
\import{naprochelibrary.tex}

\section{§12 Topological Spaces}

% from §12 Item 1: topology on a set
\begin{definition}\label{topology_on}
    $T$ is a topology on $X$ iff
    $T\subseteq\pow{X}$ and
    $\emptyset\in T$ and
    $X\in T$ and
    $T$ is closed under arbitrary unions and
    $T$ is closed under binary intersections.
\end{definition}

% from §12 Item 2: topological space with topology
\begin{definition}\label{topological_space}
    $X$ is a topological space with topology $T$ iff $T$ is a topology on $X$.
\end{definition}

% from §12 Item 3: open set
\begin{definition}\label{open_in}
    $U$ is open in $X$ with respect to $T$ iff $T$ is a topology on $X$ and $U\in T$.
\end{definition}

% from §12 Item 4: discrete topology
\begin{definition}\label{discrete_topology}
    $\discrete{X} = \pow{X}$.
\end{definition}

% from §12 Item 5: indiscrete topology
\begin{definition}\label{indiscrete_topology}
    $\indiscrete{X} = \{\emptyset, X\}$.
\end{definition}

% from §12 Item 6: finite complement topology
\begin{definition}\label{finite_complement_topology}
    $\finiteComplement{X} = \{U\in\pow{X} \mid \text{$X\setminus U$ is finite or $X\setminus U = X$}\}$.
\end{definition}

% from §12 Item 7: countable complement topology
\begin{definition}\label{countable_complement_topology}
    $\countableComplement{X} = \{U\in\pow{X} \mid \text{$X\setminus U$ is countable or $X\setminus U = X$}\}$.
\end{definition}

% from §12 Item 8: finer topology
\begin{definition}\label{finer_topology}
    $T_1$ is finer than $T$ on $X$ iff
    $T$ is a topology on $X$ and
    $T_1$ is a topology on $X$ and
    $T\subseteq T_1$.
\end{definition}

% from §12 Item 9: strictly finer topology
\begin{definition}\label{strictly_finer_topology}
    $T_1$ is strictly finer than $T$ on $X$ iff $T_1$ is finer than $T$ on $X$ and $T\neq T_1$.
\end{definition}

% from §12 Item 10: coarser topology
\begin{definition}\label{coarser_topology}
    $T$ is coarser than $T_1$ on $X$ iff $T_1$ is finer than $T$ on $X$.
\end{definition}

% from §12 Item 11: strictly coarser topology
\begin{definition}\label{strictly_coarser_topology}
    $T$ is strictly coarser than $T_1$ on $X$ iff $T$ is coarser than $T_1$ on $X$ and $T\neq T_1$.
\end{definition}

% from §12 Item 12: comparable topologies
\begin{definition}\label{comparable_topologies}
    $T$ is comparable with $T_1$ on $X$ iff
    $T$ is a topology on $X$ and
    $T_1$ is a topology on $X$ and
    either $T\subseteq T_1$ or $T_1\subseteq T$.
\end{definition}

\section{§13 Basis for a Topology}

% from §13 Item 1: basis on a set
\begin{definition}\label{basis_on}
    $B$ is a basis on $X$ iff
    $B\subseteq\pow{X}$ and
    for all $x\in X$ there exists $B_1\in B$ such that $x\in B_1$ and
    for all $B_1,B_2,x$ such that $B_1\in B$ and $B_2\in B$ and $x\in B_1\inter B_2$
    there exists $B_3\in B$ such that $x\in B_3$ and $B_3\subseteq B_1\inter B_2$.
\end{definition}

% from §13 Item 2: topology generated by a basis
\begin{definition}\label{topology_generated_by_basis}
    $\basisTopology{B}{X} = \{U\in\pow{X} \mid \text{for all $x\in U$ there exists $B_1\in B$ such that $x\in B_1$ and $B_1\subseteq U$}\}$.
\end{definition}

% from §13 Item 3: basis for a topology
\begin{definition}\label{basis_for_topology}
    $B$ is a basis for $T$ on $X$ iff $B$ is a basis on $X$ and $T = \basisTopology{B}{X}$.
\end{definition}

% from §13 Item 4: topology generated by a basis is a topology
\begin{lemma}\label{basis_topology_is_topology}
    Assume $B$ is a basis on $X$.
    Then $\basisTopology{B}{X}$ is a topology on $X$.
\end{lemma}
\begin{proof}
    Show that $\basisTopology{B}{X}\subseteq\pow{X}$.
    \begin{subproof}
        Follows by \cref{topology_generated_by_basis,subseteq}.
    \end{subproof}
    Show that $\emptyset\in\basisTopology{B}{X}$.
    \begin{subproof}
        For all $x\in\emptyset$ there exists $B_1\in B$ such that $x\in B_1$ and $B_1\subseteq\emptyset$.
        Hence $\emptyset\in\basisTopology{B}{X}$ by \cref{powerset_bottom,topology_generated_by_basis}.
    \end{subproof}
    Show that $X\in\basisTopology{B}{X}$.
    \begin{subproof}
        Follows by \cref{basis_on,subseteq,pow_iff,subseteq_refl,topology_generated_by_basis}.
    \end{subproof}
    Show that for all $M$ such that $M\subseteq\basisTopology{B}{X}$ we have $\unions{M}\in\basisTopology{B}{X}$.
    \begin{subproof}
        Follows by \cref{powerset_closed_unions,unions_iff,subseteq,topology_generated_by_basis}.
    \end{subproof}
    Show that for all $U,V$ such that $U\in\basisTopology{B}{X}$ and $V\in\basisTopology{B}{X}$ we have $U\inter V\in\basisTopology{B}{X}$.
    \begin{subproof}
        Fix $U$.
        Fix $V$.
        Assume $U\in\basisTopology{B}{X}$ and $V\in\basisTopology{B}{X}$.
        Show that for all $x$ such that $x\in U\inter V$ there exists $B_3\in B$ such that $x\in B_3$ and $B_3\subseteq U\inter V$.
        \begin{subproof}
            Fix $x$ such that $x\in U\inter V$.
            Then $x\in X$ by \cref{subseteq,pow_iff,inter_elim_left}.
            There exists $B_1\in B$ such that $x\in B_1$ and $B_1\subseteq U$ by \cref{inter_elim_left,topology_generated_by_basis}.
            There exists $B_2\in B$ such that $x\in B_2$ and $B_2\subseteq V$ by \cref{inter_elim_right,topology_generated_by_basis}.
            Then there exists $B_3\in B$ such that $x\in B_3$ and $B_3\subseteq B_1\inter B_2$ by \cref{basis_on,inter_intro}.
            Then $B_3\subseteq U\inter V$.
        \end{subproof}
        Therefore $U\inter V\in\basisTopology{B}{X}$ by \cref{subseteq,pow_iff,inter_subseteq,topology_generated_by_basis}.
    \end{subproof}
    Therefore $\basisTopology{B}{X}$ is a topology on $X$ by \cref{topology_on}.
\end{proof}

% from §13 helper lemma 1: basis elements are open in the generated topology
\begin{lemma}\label{basis_elem_open_in_generated}
    Assume $B$ is a basis on $X$.
    Suppose $B_1\in B$.
    Then $B_1\in\basisTopology{B}{X}$.
\end{lemma}
\begin{proof}
    Follows by \cref{subseteq_refl,basis_on,subseteq,pow_iff,topology_generated_by_basis}.
\end{proof}

% from §13 Item 5: unions of basis elements
\begin{lemma}\label{basis_topology_unions}
    Assume $B$ is a basis for $T$ on $X$.
    Then $T = \{U\in\pow{X} \mid \exists M\subseteq B. U = \unions{M}\}$.
\end{lemma}
\begin{proof}
    $T = \basisTopology{B}{X}$ by \cref{basis_for_topology}.
    It suffices to show that for all $U$ we have
    $U\in\basisTopology{B}{X}$ iff there exists $M\subseteq B$ such that $U = \unions{M}$.
    Fix $U$.
        Show that if $U\in\basisTopology{B}{X}$ then there exists $M\subseteq B$ such that $U = \unions{M}$.
        \begin{subproof}
            Assume $U\in\basisTopology{B}{X}$.
            Let $M = \{B_1\in B \mid B_1\subseteq U\}$.
            Hence $M\subseteq B$ and $U = \unions{M}$
                by \cref{topology_generated_by_basis,unions_iff,subseteq,subseteq_antisymmetric}.
        \end{subproof}
        Show that if there exists $M\subseteq B$ such that $U = \unions{M}$ then $U\in\basisTopology{B}{X}$.
        \begin{subproof}
            Follows by \cref{basis_topology_is_topology,basis_for_topology,basis_elem_open_in_generated,subseteq,topology_on}.
        \end{subproof}
\end{proof}

% from §13 Item 6: criterion for a collection of open sets to be a basis
\begin{lemma}\label{open_collection_basis}
    Assume $T$ is a topology on $X$.
    Assume for all $C_1\in C$ we have $C_1$ is open in $X$ with respect to $T$.
    Assume for all $U,x$ such that $U\in T$ and $x\in U$ there exists $C_1\in C$ such that $x\in C_1$ and $C_1\subseteq U$.
    Then $C$ is a basis for $T$ on $X$.
\end{lemma}
\begin{proof}
    Show that for all $x$ such that $x\in X$ there exists $C_1\in C$ such that $x\in C_1$ and $C_1\subseteq X$.
    \begin{subproof}
        Follows by \cref{topology_on}.
    \end{subproof}
    Show that for all $C_1,C_2,x$ such that $C_1\in C$ and $C_2\in C$ and $x\in C_1\inter C_2$
        there exists $C_3\in C$ such that $x\in C_3$ and $C_3\subseteq C_1\inter C_2$.
    \begin{subproof}
        Follows by \cref{open_in,topology_on}.
    \end{subproof}
    Hence $C$ is a basis on $X$ by \cref{open_in,topology_on,subseteq,basis_on}.
    Show that $T\subseteq\basisTopology{C}{X}$.
    \begin{subproof}
        Follows by \cref{topology_on,subseteq,topology_generated_by_basis}.
    \end{subproof}
    Show that $\basisTopology{C}{X}\subseteq T$.
    \begin{subproof}
        Follows by \cref{basis_topology_unions,basis_for_topology,basis_topology_is_topology,open_in,subseteq,topology_on}.
    \end{subproof}
    Hence $C$ is a basis for $T$ on $X$ by \cref{subseteq,subseteq_antisymmetric,basis_for_topology}.
\end{proof}

% from §13 Item 7: bases and finer topologies
\begin{lemma}\label{basis_finer_criterion}
    Assume $B$ is a basis for $T$ on $X$.
    Assume $B_1$ is a basis for $T_1$ on $X$.
    Then $T_1$ is finer than $T$ on $X$ iff
    for all $x,B_2$ such that $x\in X$ and $B_2\in B$ and $x\in B_2$
    there exists $B_3\in B_1$ such that $x\in B_3$ and $B_3\subseteq B_2$.
\end{lemma}
\begin{proof}
    Show that if for all $x,B_2$ such that $x\in X$ and $B_2\in B$ and $x\in B_2$
        there exists $B_3\in B_1$ such that $x\in B_3$ and $B_3\subseteq B_2$
        then $T_1$ is finer than $T$ on $X$.
    \begin{subproof}
        Assume for all $x,B_2$ such that $x\in X$ and $B_2\in B$ and $x\in B_2$
            there exists $B_3\in B_1$ such that $x\in B_3$ and $B_3\subseteq B_2$.
        $B$ is a basis on $X$ and $T = \basisTopology{B}{X}$ by \cref{basis_for_topology}.
        $B_1$ is a basis on $X$ and $T_1 = \basisTopology{B_1}{X}$ by \cref{basis_for_topology}.
        Show that for all $U\in\basisTopology{B}{X}$ we have $U\in\basisTopology{B_1}{X}$.
        \begin{subproof}
            Fix $U\in\basisTopology{B}{X}$.
            Show that for all $x\in U$ there exists $B_3\in B_1$ such that $x\in B_3$ and $B_3\subseteq U$.
            \begin{subproof}
                Fix $x\in U$.
                Then $x\in X$ by \cref{topology_generated_by_basis,pow_iff,subseteq}.
                Take $B_2\in B$ such that $x\in B_2$ and $B_2\subseteq U$ by \cref{topology_generated_by_basis}.
                Take $B_3\in B_1$ such that $x\in B_3$ and $B_3\subseteq B_2$.
                Then $B_3\subseteq U$ by \cref{subseteq_transitive}.
            \end{subproof}
            Hence $U\in\basisTopology{B_1}{X}$ by \cref{topology_generated_by_basis}.
        \end{subproof}
        Hence $\basisTopology{B}{X}\subseteq\basisTopology{B_1}{X}$ by \cref{subseteq}.
        Hence $T_1$ is finer than $T$ on $X$ by \cref{basis_topology_is_topology,finer_topology}.
    \end{subproof}
    Show that if $T_1$ is finer than $T$ on $X$ then
        for all $x,B_2$ such that $x\in X$ and $B_2\in B$ and $x\in B_2$
        there exists $B_3\in B_1$ such that $x\in B_3$ and $B_3\subseteq B_2$.
    \begin{subproof}
        Follows by \cref{basis_elem_open_in_generated,basis_for_topology,finer_topology,subseteq,topology_generated_by_basis}.
    \end{subproof}
\end{proof}

% from §13 Item 8: half-open interval on R
\begin{definition}\label{interval_closed_open}
    $\intervalclosedopen{a}{b} = \{x\in\reals \mid a \leq x \land x < b\}$.
\end{definition}

% from §13 Item 9: standard basis on R
\begin{definition}\label{standard_basis_reals}
    $\standardBasisR = \{U\in\pow{\reals} \mid \exists a,b\in\reals. a < b \land U = \intervalopen{a}{b}\}$.
\end{definition}

% from §13 Item 10: standard topology on R
\begin{definition}\label{standard_topology_reals}
    $\standardTopologyR = \basisTopology{\standardBasisR}{\reals}$.
\end{definition}

% from §13 Item 11: lower limit basis on R
\begin{definition}\label{lower_limit_basis_reals}
    $\lowerLimitBasisR = \{U\in\pow{\reals} \mid \exists a,b\in\reals. a < b \land U = \intervalclosedopen{a}{b}\}$.
\end{definition}

% from §13 Item 12: lower limit topology on R
\begin{definition}\label{lower_limit_topology_reals}
    $\lowerLimitTopologyR = \basisTopology{\lowerLimitBasisR}{\reals}$.
\end{definition}

% from §13 Item 14: K set
\begin{definition}\label{k_set}
    $\kset = \{x\in\reals \mid \exists n\in\naturalsPlus. x = 1 / n\}$.
\end{definition}

% from §13 Item 15: K basis on R
\begin{definition}\label{k_basis_reals}
    $\kBasisR = \{U\in\pow{\reals} \mid \exists a,b\in\reals. a < b \land (U = \intervalopen{a}{b} \lor U = \intervalopen{a}{b}\setminus \kset)\}$.
\end{definition}

% from §13 Item 16: K-topology on R
\begin{definition}\label{k_topology_reals}
    $\kTopologyR = \basisTopology{\kBasisR}{\reals}$.
\end{definition}

% from §13 helper lemma 2: split a weak inequality
\begin{lemma}\label{real_leq_split}
    Suppose $a \leq b$.
    Then $a < b$ or $a = b$.
\end{lemma}
\begin{proof}
    Trivial.
\end{proof}

% from §13 helper lemma 3: substitute equality on the left
\begin{lemma}\label{real_lt_subst_left}
    Suppose $a = b$ and $b < c$.
    Then $a < c$.
\end{lemma}
\begin{proof}
    Trivial.
\end{proof}

% from §13 helper lemma 4: substitute equality on the right
\begin{lemma}\label{real_lt_subst_right}
    Suppose $a < b$ and $b = c$.
    Then $a < c$.
\end{lemma}
\begin{proof}
    Trivial.
\end{proof}

% from §13 helper lemma 5: negative zero
\begin{lemma}\label{real_neg_zero}
    $\neg{0} = 0$.
\end{lemma}
\begin{proof}
    Follows by \cref{real_add_ident,real_add_inverse_unique}.
\end{proof}

% from §13 helper lemma 6: refine open intervals
\begin{lemma}\label{open_interval_refinement}
    Suppose $a,b,c,d,x\in\reals$.
    Suppose $a < x$ and $x < b$ and $c < x$ and $x < d$.
    Then there exist $e,f\in\reals$ such that $e < x$ and $x < f$ and
    $\intervalopen{e}{f} \subseteq \intervalopen{a}{b}\inter \intervalopen{c}{d}$.
\end{lemma}
    \begin{proof}
        Let $A = \{a,c\}$.
        $A\subseteq\reals$ by \cref{upair_elim,subseteq}.
        Take $m\in A$ such that $m$ is a maximum of $A$
            by \cref{pair_is_finite,upair_intro_left,emptyset,finite_real_maximum}.
        Let $B = \{b,d\}$.
        $B\subseteq\reals$ by \cref{upair_elim,subseteq}.
        Take $n\in B$ such that $n$ is a minimum of $B$
            by \cref{pair_is_finite,upair_intro_left,emptyset,finite_real_minimum}.
        Then $m < x$ and $x < n$ by \cref{upair_elim}.
        Let $e = m$.
        Let $f = n$.
        Hence $e\in\reals$ and $f\in\reals$ and $e < x$ and $x < f$ and
            $\intervalopen{e}{f}\subseteq \intervalopen{a}{b}\inter \intervalopen{c}{d}$
            by \cref{intervalopen,real_maximum,upair_intro_left,upair_intro_right,real_minimum,real_leq_split,real_lt_subst_left,real_lt_subst_right,real_lt_trans,inter_intro,subseteq}.
\end{proof}

% from §13 helper lemma 7: refine half-open intervals
\begin{lemma}\label{closedopen_interval_refinement}
    Suppose $a,b,c,d,x\in\reals$.
    Suppose $a \leq x$ and $x < b$ and $c \leq x$ and $x < d$.
    Then there exist $e,f\in\reals$ such that $e \leq x$ and $x < f$ and
    $\intervalclosedopen{e}{f} \subseteq \intervalclosedopen{a}{b}\inter \intervalclosedopen{c}{d}$.
\end{lemma}
\begin{proof}
    Let $A = \{a,c\}$.
    $A\subseteq\reals$ by \cref{upair_elim,subseteq}.
    Take $m\in A$ such that $m$ is a maximum of $A$
        by \cref{pair_is_finite,upair_intro_left,emptyset,finite_real_maximum}.
    Let $B = \{b,d\}$.
    $B\subseteq\reals$ by \cref{upair_elim,subseteq}.
    Take $n\in B$ such that $n$ is a minimum of $B$
        by \cref{pair_is_finite,upair_intro_left,emptyset,finite_real_minimum}.
    Then $m \leq x$ and $x < n$ by \cref{upair_elim}.
    Let $e = m$.
    Let $f = n$.
    Hence $e\in\reals$ and $f\in\reals$ and $e \leq x$ and $x < f$ and
        $\intervalclosedopen{e}{f}\subseteq \intervalclosedopen{a}{b}\inter \intervalclosedopen{c}{d}$
        by \cref{interval_closed_open,real_maximum,upair_intro_left,upair_intro_right,real_minimum,real_leq_trans,real_leq_split,real_lt_subst_right,real_lt_trans,inter_intro,subseteq}.
\end{proof}

% from §13 helper lemma 8: open intervals are subsets of R
\begin{lemma}\label{intervalopen_subset_reals}
    Suppose $a\in\reals$ and $b\in\reals$.
    Then $\intervalopen{a}{b}\subseteq\reals$.
\end{lemma}
\begin{proof}
    Follows by \cref{intervalopen,subseteq}.
\end{proof}

% from §13 helper lemma 9: half-open intervals are subsets of R
\begin{lemma}\label{intervalclosedopen_subset_reals}
    Suppose $a\in\reals$ and $b\in\reals$.
    Then $\intervalclosedopen{a}{b}\subseteq\reals$.
\end{lemma}
\begin{proof}
    Follows by \cref{interval_closed_open,subseteq}.
\end{proof}

% from §13 helper lemma 10: naturalsPlus elements are positive reals
\begin{lemma}\label{naturalsplus_positive_real}
    Suppose $n\in\naturalsPlus$.
    Then $0 < n$.
\end{lemma}
\begin{proof}
    $1\in\reals$ and $n\in\reals$ by \cref{real_one_in_naturals,real_naturals_subset_reals,subseteq}.
    Follows by \cref{real_one_leq_naturalsplus,real_leq_split,real_zero_lt_one,real_add_ident,real_lt_subst_right,real_lt_trans}.
\end{proof}

% from §13 helper lemma 11: reciprocals of naturalsPlus elements are reals
\begin{lemma}\label{naturalsplus_reciprocal_real}
    Suppose $n\in\naturalsPlus$.
    Then $1 / n\in\reals$.
\end{lemma}
\begin{proof}
    $1\in\reals$ and $n\in\reals$ and $0 < n$
        by \cref{real_one_in_naturals,real_naturals_subset_reals,subseteq,naturalsplus_positive_real}.
    Follows by \cref{real_pos_nonzero,real_mul_inverse,real_div,real_mul_closed}.
\end{proof}

% from §13 helper lemma 12: K-set elements are positive
\begin{lemma}\label{k_set_positive}
    Suppose $x\in\kset$.
    Then $0 < x$.
\end{lemma}
\begin{proof}
    Take $n\in\naturalsPlus$ such that $x = 1 / n$ by \cref{k_set}.
    $n\in\reals$ and $0 < n$ by \cref{real_naturals_subset_reals,subseteq,naturalsplus_positive_real}.
    Hence $0 < x$ by \cref{real_inv_pos,real_pos_nonzero,real_mul_inverse,real_div,real_mul_ident,real_lt_subst_right}.
\end{proof}

% from §13 helper lemma 13: rational between a negative real and zero
\begin{lemma}\label{rational_between_negative_and_zero}
    Suppose $a\in\reals$ and $a < 0$.
    Then there exists $p\in\rationals$ such that $p\in\intervalopen{a}{0}$.
\end{lemma}
\begin{proof}
    $0\in\reals$ and $\neg{a}\in\reals$ by \cref{real_add_ident,real_add_inverse}.
    Hence $0 < \neg{a}$ by \cref{real_lt_neg,real_neg_zero,real_lt_subst_left}.
    Take $N\in\naturalsPlus$ such that $0 < 1 / N$ and $1 / N < \neg{a}$ by \cref{real_small_reciprocal}.
    $0 < N$ and $1 / N\in\reals$ by \cref{naturalsplus_positive_real,naturalsplus_reciprocal_real}.
    Let $p = \neg{(1 / N)}$.
    Hence $p\in\intervalopen{a}{0}$ by \cref{real_add_inverse,real_lt_neg,real_neg_zero,real_lt_subst_right,real_lt_subst_left,real_neg_neg,intervalopen}.
    $1\in\naturalsPlus$ and $1\in\reals$ by \cref{real_one_in_naturals,real_mul_ident}.
    $N\in\reals$ by \cref{real_naturals_subset_reals,subseteq}.
    $\neg{1}\in\reals$ and $\neg{\neg{1}} = 1$ by \cref{real_add_inverse,real_neg_neg}.
    Then $\neg{1}\in\integers$ by \cref{real_integers}.
    Then $N\in\integers$ by \cref{real_integers}.
    Hence $N\in\integers\setminus\{0\}$ by \cref{real_pos_nonzero,singleton_iff,setminus_intro}.
    $\inv{N}\in\reals$ by \cref{real_naturals_subset_reals,subseteq,real_pos_nonzero,real_mul_inverse}.
    Then $p = \neg{1} \rmul \inv{N}$ by \cref{real_add_inverse,real_div,real_mul_neg_right,real_mul_comm}.
    Hence $p = (\neg{1}) / N$ by \cref{real_div}.
    Therefore $p\in\rationals$ by \cref{real_rationals}.
\end{proof}

% from §13 Item 17: standard basis is a basis on R
\begin{lemma}\label{standard_basis_is_basis}
    $\standardBasisR$ is a basis on $\reals$.
\end{lemma}
\begin{proof}
    Show that $\standardBasisR\subseteq\pow{\reals}$.
    \begin{subproof}
        Follows by \cref{standard_basis_reals,subseteq}.
    \end{subproof}
    Show that for all $x$ such that $x\in\reals$ we have $x\in\intervalopen{x + \neg{1}}{x + 1}$ and
        $\intervalopen{x + \neg{1}}{x + 1}\in\standardBasisR$.
    \begin{subproof}
        Follows by \cref{real_naturals_subset_reals,real_one_in_naturals,subseteq,real_add_inverse,real_add_closed,real_zero_lt_one,real_add_ident,real_lt_neg,real_neg_zero,real_lt_subst_right,real_lt_add,real_add_comm,real_lt_trans,intervalopen,intervalopen_subset_reals,powerset_intro,standard_basis_reals}.
    \end{subproof}
    Show that for all $B_1,B_2,x$ such that $B_1\in\standardBasisR$ and $B_2\in\standardBasisR$ and $x\in B_1\inter B_2$
        there exists $B_3\in\standardBasisR$ such that $x\in B_3$ and $B_3\subseteq B_1\inter B_2$.
    \begin{subproof}
        Follows by \cref{standard_basis_reals,inter_elim_left,inter_elim_right,intervalopen,open_interval_refinement,intervalopen_subset_reals,powerset_intro,real_lt_trans}.
    \end{subproof}
    Hence $\standardBasisR$ is a basis on $\reals$ by \cref{basis_on}.
\end{proof}

% from §13 Item 18: lower limit basis is a basis on R
\begin{lemma}\label{lower_limit_basis_is_basis}
    $\lowerLimitBasisR$ is a basis on $\reals$.
\end{lemma}
\begin{proof}
    Show that $\lowerLimitBasisR\subseteq\pow{\reals}$.
    \begin{subproof}
        Follows by \cref{lower_limit_basis_reals,subseteq}.
    \end{subproof}
    Show that for all $x$ such that $x\in\reals$ we have $x\in\intervalclosedopen{x}{x + 1}$ and
        $\intervalclosedopen{x}{x + 1}\in\lowerLimitBasisR$.
    \begin{subproof}
        Follows by \cref{real_naturals_subset_reals,real_one_in_naturals,subseteq,real_add_closed,real_zero_lt_one,real_add_ident,real_lt_add,real_add_comm,real_lt_subst_left,interval_closed_open,intervalclosedopen_subset_reals,powerset_intro,lower_limit_basis_reals}.
    \end{subproof}
    Show that for all $B_1,B_2,x$ such that $B_1\in\lowerLimitBasisR$ and $B_2\in\lowerLimitBasisR$ and $x\in B_1\inter B_2$
        there exists $B_3\in\lowerLimitBasisR$ such that $x\in B_3$ and $B_3\subseteq B_1\inter B_2$.
    \begin{subproof}
        Follows by \cref{lower_limit_basis_reals,inter_elim_left,inter_elim_right,interval_closed_open,closedopen_interval_refinement,intervalclosedopen_subset_reals,powerset_intro,real_leq_split,real_lt_subst_left,real_lt_trans}.
    \end{subproof}
    Hence $\lowerLimitBasisR$ is a basis on $\reals$ by \cref{basis_on}.
\end{proof}

% from §13 Item 19: K basis is a basis on R
\begin{lemma}\label{k_basis_is_basis}
    $\kBasisR$ is a basis on $\reals$.
\end{lemma}
\begin{proof}
    Show that $\kBasisR\subseteq\pow{\reals}$.
    \begin{subproof}
        Follows by \cref{k_basis_reals,subseteq}.
    \end{subproof}
    Show that for all $x$ such that $x\in\reals$ we have $x\in\intervalopen{x + \neg{1}}{x + 1}$ and
        $\intervalopen{x + \neg{1}}{x + 1}\in\kBasisR$.
    \begin{subproof}
        Follows by \cref{real_naturals_subset_reals,real_one_in_naturals,subseteq,real_add_inverse,real_add_closed,real_zero_lt_one,real_add_ident,real_lt_neg,real_neg_zero,real_lt_subst_right,real_lt_add,real_add_comm,real_lt_trans,intervalopen,intervalopen_subset_reals,powerset_intro,k_basis_reals}.
    \end{subproof}
    Show that for all $B_1,B_2,x$ such that $B_1\in\kBasisR$ and $B_2\in\kBasisR$ and $x\in B_1\inter B_2$
        there exists $B_3\in\kBasisR$ such that $x\in B_3$ and $B_3\subseteq B_1\inter B_2$.
    \begin{subproof}
        Fix $B_1,B_2$.
        Fix $x$ such that $B_1\in\kBasisR$ and $B_2\in\kBasisR$ and $x\in B_1\inter B_2$.
        Take $a,b\in\reals$ such that $a < b$ and
            $(B_1 = \intervalopen{a}{b} \lor B_1 = \intervalopen{a}{b}\setminus \kset)$
            by \cref{k_basis_reals}.
        Take $c,d\in\reals$ such that $c < d$ and
            $(B_2 = \intervalopen{c}{d} \lor B_2 = \intervalopen{c}{d}\setminus \kset)$
            by \cref{k_basis_reals}.
        \begin{byCase}
        \caseOf{$B_1 = \intervalopen{a}{b}$ and $B_2 = \intervalopen{c}{d}$.}
            Then $x\in\reals$ and $a < x$ and $x < b$ and $c < x$ and $x < d$
                by \cref{inter_elim_left,inter_elim_right,intervalopen}.
            Take $e,f\in\reals$ such that $e < x$ and $x < f$ and
                $\intervalopen{e}{f} \subseteq \intervalopen{a}{b}\inter \intervalopen{c}{d}$
                by \cref{open_interval_refinement}.
            Hence there exists $B_3\in\kBasisR$ such that $x\in B_3$ and $B_3\subseteq B_1\inter B_2$
                by \cref{intervalopen_subset_reals,powerset_intro,k_basis_reals,real_lt_trans,intervalopen}.
        \caseOf{$B_1 = \intervalopen{a}{b}\setminus \kset$ and $B_2 = \intervalopen{c}{d}$.}
            Then $x\in\reals$ and $a < x$ and $x < b$ and $c < x$ and $x < d$ and $x\notin\kset$
                by \cref{inter_elim_left,inter_elim_right,setminus_elim_left,setminus_elim_right,intervalopen}.
            Take $e,f\in\reals$ such that $e < x$ and $x < f$ and
                $\intervalopen{e}{f} \subseteq \intervalopen{a}{b}\inter \intervalopen{c}{d}$
                by \cref{open_interval_refinement}.
            Hence there exists $B_3\in\kBasisR$ such that $x\in B_3$ and $B_3\subseteq B_1\inter B_2$
                by \cref{setminus_elim_left,intervalopen,subseteq,powerset_intro,k_basis_reals,real_lt_trans,setminus_intro,setminus_elim_right,inter_elim_left,inter_elim_right,inter_intro}.
        \caseOf{$B_1 = \intervalopen{a}{b}$ and $B_2 = \intervalopen{c}{d}\setminus \kset$.}
            Then $x\in\reals$ and $a < x$ and $x < b$ and $c < x$ and $x < d$ and $x\notin\kset$
                by \cref{inter_elim_left,inter_elim_right,setminus_elim_left,setminus_elim_right,intervalopen}.
            Take $e,f\in\reals$ such that $e < x$ and $x < f$ and
                $\intervalopen{e}{f} \subseteq \intervalopen{a}{b}\inter \intervalopen{c}{d}$
                by \cref{open_interval_refinement}.
            Hence there exists $B_3\in\kBasisR$ such that $x\in B_3$ and $B_3\subseteq B_1\inter B_2$
                by \cref{setminus_elim_left,intervalopen,subseteq,powerset_intro,k_basis_reals,real_lt_trans,setminus_intro,setminus_elim_right,inter_elim_left,inter_elim_right,inter_intro}.
        \caseOf{$B_1 = \intervalopen{a}{b}\setminus \kset$ and $B_2 = \intervalopen{c}{d}\setminus \kset$.}
            Then $x\in\reals$ and $a < x$ and $x < b$ and $c < x$ and $x < d$ and $x\notin\kset$
                by \cref{inter_elim_left,inter_elim_right,setminus_elim_left,setminus_elim_right,intervalopen}.
            Take $e,f\in\reals$ such that $e < x$ and $x < f$ and
                $\intervalopen{e}{f} \subseteq \intervalopen{a}{b}\inter \intervalopen{c}{d}$
                by \cref{open_interval_refinement}.
            Let $B_3 = \intervalopen{e}{f}\setminus \kset$.
            Then $x\in\intervalopen{e}{f}$ by \cref{intervalopen}.
            Then $x\in B_3$ by \cref{setminus_intro}.
            Then $B_3\subseteq\reals$ by \cref{setminus_elim_left,intervalopen,subseteq}.
            Hence $B_3\in\pow{\reals}$ by \cref{powerset_intro}.
            Then $e < f$ by \cref{real_lt_trans}.
            Then $B_3\in\kBasisR$ by \cref{k_basis_reals}.
            Show that for all $y$ such that $y\in B_3$ we have $y\in B_1\inter B_2$.
            \begin{subproof}
                Fix $y$ such that $y\in B_3$.
                Then $y\in\intervalopen{e}{f}$ and $y\notin\kset$ by \cref{setminus_elim_left,setminus_elim_right}.
                Then $y\in\intervalopen{a}{b}\inter \intervalopen{c}{d}$ by \cref{subseteq}.
                Then $y\in\intervalopen{a}{b}$ and $y\in\intervalopen{c}{d}$ by \cref{inter_elim_left,inter_elim_right}.
                Then $y\in B_1$ and $y\in B_2$ by \cref{setminus_intro}.
                Thus $y\in B_1\inter B_2$ by \cref{inter_intro}.
            \end{subproof}
            Hence $B_3\subseteq B_1\inter B_2$ by \cref{subseteq}.
            Hence there exists $C\in\kBasisR$ such that $x\in C$ and $C\subseteq B_1\inter B_2$.
        \end{byCase}
    \end{subproof}
    Hence $\kBasisR$ is a basis on $\reals$ by \cref{basis_on}.
\end{proof}

% from §13 Item 20: lower limit topology strictly finer than standard topology
\begin{lemma}\label{lower_limit_strictly_finer_standard}
    $\lowerLimitTopologyR$ is strictly finer than $\standardTopologyR$ on $\reals$.
\end{lemma}
\begin{proof}
    Let $T = \standardTopologyR$.
    Let $T_1 = \lowerLimitTopologyR$.
    $\standardBasisR$ is a basis for $T$ on $\reals$ and
    $\lowerLimitBasisR$ is a basis for $T_1$ on $\reals$
        by \cref{standard_basis_is_basis,lower_limit_basis_is_basis,basis_for_topology,standard_topology_reals,lower_limit_topology_reals}.
    Show that for all $x,B_2$ such that $x\in\reals$ and $B_2\in\standardBasisR$ and $x\in B_2$
        there exists $B_3\in\lowerLimitBasisR$ such that $x\in B_3$ and $B_3\subseteq B_2$.
    \begin{subproof}
        Follows by \cref{standard_basis_reals,intervalopen,intervalclosedopen_subset_reals,powerset_intro,lower_limit_basis_reals,interval_closed_open,real_leq_split,real_lt_subst_right,real_lt_trans,subseteq}.
    \end{subproof}
    Hence $T_1$ is finer than $T$ on $\reals$ by \cref{basis_finer_criterion}.
    Let $U = \intervalclosedopen{0}{1}$.
    Show that $U\in T_1$.
    \begin{subproof}
        Then $0\in U$ by \cref{interval_closed_open,real_add_ident,real_naturals_subset_reals,real_one_in_naturals,subseteq,real_zero_lt_one}.
        Therefore $U\in T_1$ by \cref{real_add_ident,real_naturals_subset_reals,real_one_in_naturals,subseteq,real_zero_lt_one,intervalclosedopen_subset_reals,powerset_intro,lower_limit_basis_reals,basis_elem_open_in_generated,basis_for_topology,lower_limit_basis_is_basis,lower_limit_topology_reals}.
    \end{subproof}
    Show that $U\notin T$.
    \begin{subproof}
        Suppose not.
        Then $0\in U$ by \cref{interval_closed_open,real_add_ident,real_naturals_subset_reals,real_one_in_naturals,subseteq,real_zero_lt_one}.
        There exists $B_2\in\standardBasisR$ such that $0\in B_2$ and $B_2\subseteq U$
            by \cref{interval_closed_open,topology_generated_by_basis,standard_topology_reals}.
        Take $a,b\in\reals$ such that $a < b$ and $B_2 = \intervalopen{a}{b}$ by \cref{standard_basis_reals}.
        Hence $a < 0$ and $0 < b$ by \cref{intervalopen}.
        Take $p$ such that $p\in\intervalopen{a}{0}$ by \cref{rational_between_negative_and_zero}.
        Then $p\in\reals$ and $p < 0$ by \cref{intervalopen}.
        Hence $p\in B_2$ by \cref{intervalopen,real_add_ident,real_lt_trans}.
        $p\notin U$ by \cref{interval_closed_open,real_add_ident,real_lt_subst_left,real_lt_trans,real_lt_irreflexive}.
        Contradiction by \cref{subseteq}.
    \end{subproof}
    Hence $T_1$ is strictly finer than $T$ on $\reals$ by \cref{strictly_finer_topology}.
\end{proof}

% from §13 Item 21: K-topology strictly finer than standard topology
\begin{lemma}\label{k_strictly_finer_standard}
    $\kTopologyR$ is strictly finer than $\standardTopologyR$ on $\reals$.
\end{lemma}
\begin{proof}
    Let $T = \standardTopologyR$.
    Let $T_2 = \kTopologyR$.
    $\standardBasisR$ is a basis for $T$ on $\reals$ and
    $\kBasisR$ is a basis for $T_2$ on $\reals$
        by \cref{standard_basis_is_basis,k_basis_is_basis,basis_for_topology,standard_topology_reals,k_topology_reals}.
    Show that for all $x,B_2$ such that $x\in\reals$ and $B_2\in\standardBasisR$ and $x\in B_2$
        there exists $B_3\in\kBasisR$ such that $x\in B_3$ and $B_3\subseteq B_2$.
    \begin{subproof}
        Follows by \cref{k_basis_reals,standard_basis_reals,subseteq_refl}.
    \end{subproof}
    Hence $T_2$ is finer than $T$ on $\reals$ by \cref{basis_finer_criterion}.
    Let $U = \intervalopen{\neg{1}}{1}\setminus \kset$.
    Show that $U\in T_2$.
    \begin{subproof}
        $1\in\reals$ and $\neg{1}\in\reals$ and $0 < 1$ and $0\in\reals$ and $\neg{1} < 0$
            by \cref{real_naturals_subset_reals,real_one_in_naturals,subseteq,real_add_inverse,real_zero_lt_one,real_add_ident,real_lt_neg,real_neg_zero,real_lt_subst_right}.
        Therefore $\neg{1} < 1$ by \cref{real_lt_trans}.
        Then $0\in U$ by \cref{intervalopen,k_set_positive,real_lt_irreflexive,setminus_intro}.
        Therefore $U\in T_2$ by \cref{setminus_elim_left,intervalopen,subseteq,powerset_intro,k_basis_reals,basis_elem_open_in_generated,basis_for_topology,k_basis_is_basis,k_topology_reals}.
    \end{subproof}
    Show that $U\notin T$.
    \begin{subproof}
        Suppose not.
        $1\in\reals$ and $\neg{1}\in\reals$ and $0 < 1$ and $0\in\reals$ and $\neg{1} < 0$
            by \cref{real_naturals_subset_reals,real_one_in_naturals,subseteq,real_add_inverse,real_zero_lt_one,real_add_ident,real_lt_neg,real_neg_zero,real_lt_subst_right}.
        Therefore $\neg{1} < 1$ by \cref{real_lt_trans}.
        Then $0\in U$ by \cref{intervalopen,k_set_positive,real_lt_irreflexive,setminus_intro}.
        There exists $B_2\in\standardBasisR$ such that $0\in B_2$ and $B_2\subseteq U$
            by \cref{topology_generated_by_basis,standard_topology_reals}.
        Take $a,b\in\reals$ such that $a < b$ and $B_2 = \intervalopen{a}{b}$ by \cref{standard_basis_reals}.
        Hence $a < 0$ and $0 < b$ by \cref{intervalopen}.
        Take $N\in\naturalsPlus$ such that $0 < 1 / N$ and $1 / N < b$ by \cref{real_small_reciprocal}.
        $0 < N$ and $1 / N\in\reals$ by \cref{naturalsplus_positive_real,naturalsplus_reciprocal_real}.
        Hence $a < 1 / N$ by \cref{real_add_ident,real_lt_trans}.
        Then $1 / N\in B_2$ by \cref{intervalopen}.
        Therefore $1 / N\notin U$ by \cref{k_set,setminus_elim_right}.
        Contradiction by \cref{subseteq}.
    \end{subproof}
    Hence $T_2$ is strictly finer than $T$ on $\reals$ by \cref{strictly_finer_topology}.
\end{proof}

% from §13 Item 22: lower limit and K topologies are not comparable
\begin{lemma}\label{lower_limit_k_not_comparable}
    $\lowerLimitTopologyR$ is not comparable with $\kTopologyR$ on $\reals$.
\end{lemma}
\begin{proof}
    Let $T_1 = \lowerLimitTopologyR$.
    Let $T_2 = \kTopologyR$.
    Let $U_1 = \intervalclosedopen{0}{1}$.
    Let $U_2 = \intervalopen{\neg{1}}{1}\setminus \kset$.
    Show that $U_1\in T_1$.
    \begin{subproof}
        $0\in\reals$ and $1\in\reals$ and $0 < 1$
            by \cref{real_add_ident,real_naturals_subset_reals,real_one_in_naturals,subseteq,real_zero_lt_one}.
        Therefore $U_1\in T_1$ by \cref{intervalclosedopen_subset_reals,powerset_intro,lower_limit_basis_reals,basis_elem_open_in_generated,basis_for_topology,lower_limit_basis_is_basis,lower_limit_topology_reals}.
    \end{subproof}
    Show that $U_2\in T_2$.
    \begin{subproof}
        $1\in\reals$ and $\neg{1}\in\reals$ and $0 < 1$ and $0\in\reals$ and $\neg{1} < 0$
            by \cref{real_naturals_subset_reals,real_one_in_naturals,subseteq,real_add_inverse,real_zero_lt_one,real_add_ident,real_lt_neg,real_neg_zero,real_lt_subst_right}.
        Therefore $\neg{1} < 1$ by \cref{real_lt_trans}.
        Then $0\in U_2$ by \cref{intervalopen,k_set_positive,real_lt_irreflexive,setminus_intro}.
        Therefore $U_2\in T_2$ by \cref{setminus_elim_left,intervalopen,subseteq,powerset_intro,k_basis_reals,basis_elem_open_in_generated,basis_for_topology,k_basis_is_basis,k_topology_reals}.
    \end{subproof}
    Show that $U_1\notin T_2$.
    \begin{subproof}
        Suppose not.
        Then $0\in U_1$ by \cref{interval_closed_open,real_add_ident,real_naturals_subset_reals,real_one_in_naturals,subseteq,real_zero_lt_one}.
        There exists $B_2\in\kBasisR$ such that $0\in B_2$ and $B_2\subseteq U_1$
            by \cref{interval_closed_open,topology_generated_by_basis,k_topology_reals}.
        Take $a,b\in\reals$ such that $a < b$ and
            $(B_2 = \intervalopen{a}{b} \lor B_2 = \intervalopen{a}{b}\setminus \kset)$
            by \cref{k_basis_reals}.
        Hence $a < 0$ and $0 < b$ by \cref{intervalopen,setminus_elim_left}.
        Take $p$ such that $p\in\intervalopen{a}{0}$ by \cref{rational_between_negative_and_zero}.
        Then $p\in\reals$ and $p < 0$ by \cref{intervalopen}.
        Then $a < p$ by \cref{intervalopen}.
        $0\in\reals$ by \cref{real_add_ident}.
        Hence $p < b$ by \cref{real_lt_trans}.
        Therefore $p\in\intervalopen{a}{b}$ by \cref{intervalopen}.
        $p\notin\kset$ by \cref{k_set_positive,real_lt_trans,real_lt_irreflexive}.
        Therefore $p\in B_2$ by \cref{setminus_intro}.
        $p\notin U_1$ by \cref{interval_closed_open,real_add_ident,real_lt_subst_left,real_lt_trans,real_lt_irreflexive}.
        Contradiction by \cref{subseteq}.
    \end{subproof}
    Show that $U_2\notin T_1$.
    \begin{subproof}
        Suppose not.
        $1\in\reals$ and $\neg{1}\in\reals$ and $0 < 1$ and $0\in\reals$ and $\neg{1} < 0$
            by \cref{real_naturals_subset_reals,real_one_in_naturals,subseteq,real_add_inverse,real_zero_lt_one,real_add_ident,real_lt_neg,real_neg_zero,real_lt_subst_right}.
        Therefore $\neg{1} < 1$ by \cref{real_lt_trans}.
        Then $0\in U_2$ by \cref{intervalopen,k_set_positive,real_lt_irreflexive,setminus_intro}.
        There exists $B_1\in\lowerLimitBasisR$ such that $0\in B_1$ and $B_1\subseteq U_2$
            by \cref{topology_generated_by_basis,lower_limit_topology_reals}.
        Take $a,b\in\reals$ such that $a < b$ and $B_1 = \intervalclosedopen{a}{b}$ by \cref{lower_limit_basis_reals}.
        Hence $a \leq 0$ and $0 < b$ by \cref{interval_closed_open}.
        Take $N\in\naturalsPlus$ such that $0 < 1 / N$ and $1 / N < b$ by \cref{real_small_reciprocal}.
        $0 < N$ and $1 / N\in\reals$ by \cref{naturalsplus_positive_real,naturalsplus_reciprocal_real}.
        Hence $a < 1 / N$ by \cref{real_leq_split,real_lt_subst_left,real_lt_trans}.
        Then $1 / N\in B_1$ by \cref{interval_closed_open}.
        Therefore $1 / N\notin U_2$ by \cref{k_set,setminus_elim_right}.
        Contradiction by \cref{subseteq}.
    \end{subproof}
    $T_1$ is not comparable with $T_2$ on $\reals$ by \cref{comparable_topologies,subseteq}.
\end{proof}

% from §13 Item 23: subbasis on a set
\begin{definition}\label{subbasis_on}
    $S$ is a subbasis on $X$ iff $S\subseteq\pow{X}$ and $\unions{S} = X$.
\end{definition}

% from §13 Item 24: finite intersections of subbasis elements
\begin{definition}\label{finite_intersections}
    $\finiteIntersections{S} = \{B\in\pow{\unions{S}} \mid \text{there exists $F\subseteq S$ such that $F$ is finite and inhabited and $B = \inters{F}$}\}$.
\end{definition}

% from §13 Item 25: topology generated by a subbasis
\begin{definition}\label{topology_generated_by_subbasis}
    $\subbasisTopology{S}{X} = \basisTopology{\finiteIntersections{S}}{X}$.
\end{definition}

% from §13 Item 26: finite intersections of a subbasis form a basis
\begin{lemma}\label{subbasis_finite_intersections_basis}
    Assume $S$ is a subbasis on $X$.
    Then $\finiteIntersections{S}$ is a basis on $X$.
\end{lemma}
\begin{proof}
    Show that for all $B_1,B_2,x$ such that $B_1\in\finiteIntersections{S}$ and $B_2\in\finiteIntersections{S}$ and $x\in B_1\inter B_2$
        there exists $B_3\in\finiteIntersections{S}$ such that $x\in B_3$ and $B_3\subseteq B_1\inter B_2$.
    \begin{subproof}
        Follows by \cref{finite_intersections,union_subsets_is_subset,union_of_finite_is_finite,union_intro_left,pow_iff,inter_subseteq,powerset_intro,inters_distrib_union,subseteq_refl}.
    \end{subproof}
    Show that for all $x$ such that $x\in X$ there exists $B\in\finiteIntersections{S}$ such that $x\in B$.
    \begin{subproof}
        Fix $x$ such that $x\in X$.
        Take $U$ such that $U\in S$ and $x\in U$ by \cref{subbasis_on,unions_iff}.
        Then $U\in\finiteIntersections{S}$
            by \cref{subbasis_on,subseteq,pow_iff,powerset_intro,singleton_subset_intro,pair_is_finite,upair_intro_left,subseteq_finite_is_finite,singleton_intro,inters_singleton,finite_intersections}.
        Hence there exists $B\in\finiteIntersections{S}$ such that $x\in B$.
    \end{subproof}
    Therefore $\finiteIntersections{S}$ is a basis on $X$ by \cref{subbasis_on,finite_intersections,pow_iff,powerset_intro,subseteq,basis_on}.
\end{proof}

\section{§14 The Order Topology}

% from §14 Item 1: more than one element
\begin{definition}\label{more_than_one_element}
    $X$ has more than one element iff there exist $x, y\in X$ such that $x\neq y$.
\end{definition}

% from §14 Item 2: simple order relation
\begin{definition}\label{simple_order_relation}
    $R$ is a simple order relation on $X$ iff
    $R$ is a binary relation on $X$ and
    $R$ is transitive and
    $R$ is asymmetric and
    $R$ is connex on $X$.
\end{definition}

% from §14 Item 3: smallest element
\begin{definition}\label{smallest_element}
    $a$ is a smallest element of $X$ with respect to $R$ iff
    $a\in X$ and for all $x\in X$ we have $a\mathrel{R}x$ or $a = x$.
\end{definition}

% from §14 Item 4: largest element
\begin{definition}\label{largest_element}
    $b$ is a largest element of $X$ with respect to $R$ iff
    $b\in X$ and for all $x\in X$ we have $x\mathrel{R}b$ or $x = b$.
\end{definition}

% from §14 Item 5: open interval in an ordered set
\begin{definition}\label{intervalopen_rel}
    $\intervalopenrel{R}{X}{a}{b} = \{x\in X \mid a\mathrel{R}x \land x\mathrel{R}b\}$.
\end{definition}

% from §14 Item 6: half-open interval in an ordered set
\begin{definition}\label{intervalopenclosed_rel}
    $\intervalopenclosedrel{R}{X}{a}{b} = \{x\in X \mid a\mathrel{R}x \land (x\mathrel{R}b \lor x = b)\}$.
\end{definition}

% from §14 Item 7: half-open interval in an ordered set
\begin{definition}\label{intervalclosedopen_rel}
    $\intervalclosedopenrel{R}{X}{a}{b} = \{x\in X \mid (a\mathrel{R}x \lor x = a) \land x\mathrel{R}b\}$.
\end{definition}

% from §14 Item 8: closed interval in an ordered set
\begin{definition}\label{intervalclosed_rel}
    $\intervalclosedrel{R}{X}{a}{b} = \{x\in X \mid (a\mathrel{R}x \lor x = a) \land (x\mathrel{R}b \lor x = b)\}$.
\end{definition}

% from §14 Item 9: order basis
\begin{definition}\label{order_basis}
    $\orderBasis{R}{X} = \{B\in\pow{X} \mid
        (\exists a, b\in X. a\mathrel{R}b \land B = \intervalopenrel{R}{X}{a}{b}) \lor
        (\exists b\in X. \exists a_0\in X. (\forall x\in X. a_0\mathrel{R}x \lor a_0 = x) \land B = \intervalclosedopenrel{R}{X}{a_0}{b}) \lor
        (\exists a\in X. \exists b_0\in X. (\forall x\in X. x\mathrel{R}b_0 \lor x = b_0) \land B = \intervalopenclosedrel{R}{X}{a}{b_0})\}$.
\end{definition}

% from §14 Item 10: order topology
\begin{definition}\label{order_topology}
    $T$ is the order topology on $X$ with respect to $R$ iff
    $R$ is a simple order relation on $X$ and
    $X$ has more than one element and
    $T = \basisTopology{\orderBasis{R}{X}}{X}$.
\end{definition}

% from §14 helper lemma 1: distinct element in a set with more than one element
\begin{lemma}\label{more_than_one_element_has_other}
    Assume $X$ has more than one element.
    Suppose $x\in X$.
    Then there exists $y\in X$ such that $y\neq x$.
\end{lemma}
\begin{proof}
    Follows by \cref{more_than_one_element}.
\end{proof}

% from §14 helper lemma 2: non-smallest has a predecessor
\begin{lemma}\label{non_smallest_has_predecessor}
    Assume $R$ is a simple order relation on $X$.
    Suppose $x\in X$.
    Suppose $x$ is not a smallest element of $X$ with respect to $R$.
    Then there exists $y\in X$ such that $y\mathrel{R}x$.
\end{lemma}
\begin{proof}
    Follows by \cref{simple_order_relation,connex,smallest_element}.
\end{proof}

% from §14 helper lemma 3: non-largest has a successor
\begin{lemma}\label{non_largest_has_successor}
    Assume $R$ is a simple order relation on $X$.
    Suppose $x\in X$.
    Suppose $x$ is not a largest element of $X$ with respect to $R$.
    Then there exists $y\in X$ such that $x\mathrel{R}y$.
\end{lemma}
\begin{proof}
    Follows by \cref{simple_order_relation,connex,largest_element}.
\end{proof}

% from §14 helper lemma 4: open interval inclusion by endpoints
\begin{lemma}\label{intervalopenrel_subset_by_endpoints}
    Assume $R$ is transitive.
    Suppose $a,a_1,b,b_1\in X$.
    Suppose $a_1\mathrel{R}a$ or $a_1 = a$.
    Suppose $b\mathrel{R}b_1$ or $b = b_1$.
    Then $\intervalopenrel{R}{X}{a}{b} \subseteq \intervalopenrel{R}{X}{a_1}{b_1}$.
\end{lemma}
\begin{proof}
    Follows by \cref{intervalopen_rel,subseteq,transitive}.
\end{proof}

% from §14 helper lemma 5: closed-open interval inclusion by right endpoint
\begin{lemma}\label{intervalclosedopenrel_subset_by_right}
    Assume $R$ is transitive.
    Suppose $a,b,b_1\in X$.
    Suppose $b\mathrel{R}b_1$ or $b = b_1$.
    Then $\intervalclosedopenrel{R}{X}{a}{b} \subseteq \intervalclosedopenrel{R}{X}{a}{b_1}$.
\end{lemma}
\begin{proof}
    Follows by \cref{intervalclosedopen_rel,subseteq,transitive}.
\end{proof}

% from §14 helper lemma 6: open-closed interval inclusion by left endpoint
\begin{lemma}\label{intervalopenclosedrel_subset_by_left}
    Assume $R$ is transitive.
    Suppose $a,a_1,b\in X$.
    Suppose $a_1\mathrel{R}a$ or $a_1 = a$.
    Then $\intervalopenclosedrel{R}{X}{a}{b} \subseteq \intervalopenclosedrel{R}{X}{a_1}{b}$.
\end{lemma}
\begin{proof}
    Follows by \cref{intervalopenclosed_rel,subseteq,transitive}.
\end{proof}
% from §14 helper lemma 7: basis element at a smallest element
\begin{lemma}\label{orderbasis_elem_at_smallest}
    Assume $R$ is a simple order relation on $X$.
    Suppose $x$ is a smallest element of $X$ with respect to $R$.
    Suppose $B\in\orderBasis{R}{X}$.
    Suppose $x\in B$.
    Then there exists $b\in X$ such that $x\mathrel{R}b$ and $B = \intervalclosedopenrel{R}{X}{x}{b}$.
\end{lemma}
\begin{proof}
    $((\exists a, b\in X. a\mathrel{R}b \land B = \intervalopenrel{R}{X}{a}{b}) \lor
        (\exists b\in X. \exists a_0\in X. (\forall x\in X. a_0\mathrel{R}x \lor a_0 = x) \land B = \intervalclosedopenrel{R}{X}{a_0}{b}) \lor
        (\exists a\in X. \exists b_0\in X. (\forall x\in X. x\mathrel{R}b_0 \lor x = b_0) \land B = \intervalopenclosedrel{R}{X}{a}{b_0}))$
        by \cref{order_basis}.
    \begin{byCase}
    \caseOf{$\exists a, b\in X. a\mathrel{R}b \land B = \intervalopenrel{R}{X}{a}{b}$.}
        Take $a, b\in X$ such that $a\mathrel{R}b$ and $B = \intervalopenrel{R}{X}{a}{b}$.
        Then $a\mathrel{R}x$ by \cref{intervalopen_rel}.
        $x\mathrel{R}a$ or $x = a$ by \cref{smallest_element}.
        Suppose not.
        Contradiction by \cref{simple_order_relation,asymmetric}.
    \caseOf{$\exists b\in X. \exists a_0\in X. (\forall x\in X. a_0\mathrel{R}x \lor a_0 = x) \land B = \intervalclosedopenrel{R}{X}{a_0}{b}$.}
        Take $a_0, b\in X$ such that
            $(\forall x\in X. a_0\mathrel{R}x \lor a_0 = x)$ and $B = \intervalclosedopenrel{R}{X}{a_0}{b}$.
        Then $(a_0\mathrel{R}x \lor x = a_0)$ and $x\mathrel{R}b$ by \cref{intervalclosedopen_rel}.
        Hence $a_0 = x$ by \cref{intervalclosedopen_rel,smallest_element,simple_order_relation,asymmetric}.
        Therefore there exists $c\in X$ such that $x\mathrel{R}c$ and $B = \intervalclosedopenrel{R}{X}{x}{c}$.
    \caseOf{$\exists a\in X. \exists b_0\in X. (\forall x\in X. x\mathrel{R}b_0 \lor x = b_0) \land B = \intervalopenclosedrel{R}{X}{a}{b_0}$.}
        Take $a, b_0\in X$ such that
            $(\forall x\in X. x\mathrel{R}b_0 \lor x = b_0)$ and $B = \intervalopenclosedrel{R}{X}{a}{b_0}$.
        Then $a\mathrel{R}x$ by \cref{intervalopenclosed_rel}.
        $x\mathrel{R}a$ or $x = a$ by \cref{smallest_element}.
        Suppose not.
        Contradiction by \cref{simple_order_relation,asymmetric}.
    \end{byCase}
\end{proof}
% from §14 helper lemma 8: basis element at a largest element
\begin{lemma}\label{orderbasis_elem_at_largest}
    Assume $R$ is a simple order relation on $X$.
    Suppose $x$ is a largest element of $X$ with respect to $R$.
    Suppose $B\in\orderBasis{R}{X}$.
    Suppose $x\in B$.
    Then there exists $a\in X$ such that $a\mathrel{R}x$ and $B = \intervalopenclosedrel{R}{X}{a}{x}$.
\end{lemma}
\begin{proof}
    $((\exists a, b\in X. a\mathrel{R}b \land B = \intervalopenrel{R}{X}{a}{b}) \lor
        (\exists b\in X. \exists a_0\in X. (\forall x\in X. a_0\mathrel{R}x \lor a_0 = x) \land B = \intervalclosedopenrel{R}{X}{a_0}{b}) \lor
        (\exists a\in X. \exists b_0\in X. (\forall x\in X. x\mathrel{R}b_0 \lor x = b_0) \land B = \intervalopenclosedrel{R}{X}{a}{b_0}))$
        by \cref{order_basis}.
    \begin{byCase}
    \caseOf{$\exists a, b\in X. a\mathrel{R}b \land B = \intervalopenrel{R}{X}{a}{b}$.}
        Take $a, b\in X$ such that $a\mathrel{R}b$ and $B = \intervalopenrel{R}{X}{a}{b}$.
        Then $x\mathrel{R}b$ by \cref{intervalopen_rel}.
        $b\mathrel{R}x$ or $b = x$ by \cref{largest_element}.
        Suppose not.
        Contradiction by \cref{simple_order_relation,asymmetric}.
    \caseOf{$\exists b\in X. \exists a_0\in X. (\forall x\in X. a_0\mathrel{R}x \lor a_0 = x) \land B = \intervalclosedopenrel{R}{X}{a_0}{b}$.}
        Take $a_0, b\in X$ such that
            $(\forall x\in X. a_0\mathrel{R}x \lor a_0 = x)$ and $B = \intervalclosedopenrel{R}{X}{a_0}{b}$.
        Then $x\mathrel{R}b$ by \cref{intervalclosedopen_rel}.
        $b\mathrel{R}x$ or $b = x$ by \cref{largest_element}.
        Suppose not.
        Contradiction by \cref{simple_order_relation,asymmetric}.
    \caseOf{$\exists a\in X. \exists b_0\in X. (\forall x\in X. x\mathrel{R}b_0 \lor x = b_0) \land B = \intervalopenclosedrel{R}{X}{a}{b_0}$.}
        Take $a, b_0\in X$ such that
            $(\forall x\in X. x\mathrel{R}b_0 \lor x = b_0)$ and $B = \intervalopenclosedrel{R}{X}{a}{b_0}$.
        Then $a\mathrel{R}x$ and $(x\mathrel{R}b_0 \lor x = b_0)$ by \cref{intervalopenclosed_rel}.
        Hence $b_0 = x$ by \cref{intervalopenclosed_rel,largest_element,simple_order_relation,asymmetric}.
        Therefore there exists $c\in X$ such that $c\mathrel{R}x$ and $B = \intervalopenclosedrel{R}{X}{c}{x}$.
    \end{byCase}
\end{proof}

% from §14 helper lemma 9: open interval inside a basis element
\begin{lemma}\label{orderbasis_elem_contains_open_interval}
    Assume $R$ is a simple order relation on $X$.
    Suppose $x\in X$.
    Suppose $x$ is not a smallest element of $X$ with respect to $R$.
    Suppose $x$ is not a largest element of $X$ with respect to $R$.
    Suppose $B\in\orderBasis{R}{X}$.
    Suppose $x\in B$.
    Then there exist $a,b\in X$ such that $a\mathrel{R}x$ and $x\mathrel{R}b$ and
        $\intervalopenrel{R}{X}{a}{b} \subseteq B$.
\end{lemma}
\begin{proof}
    $((\exists a, b\in X. a\mathrel{R}b \land B = \intervalopenrel{R}{X}{a}{b}) \lor
        (\exists b\in X. \exists a_0\in X. (\forall x\in X. a_0\mathrel{R}x \lor a_0 = x) \land B = \intervalclosedopenrel{R}{X}{a_0}{b}) \lor
        (\exists a\in X. \exists b_0\in X. (\forall x\in X. x\mathrel{R}b_0 \lor x = b_0) \land B = \intervalopenclosedrel{R}{X}{a}{b_0}))$
        by \cref{order_basis}.
    \begin{byCase}
    \caseOf{$\exists a, b\in X. a\mathrel{R}b \land B = \intervalopenrel{R}{X}{a}{b}$.}
        Take $a, b\in X$ such that $a\mathrel{R}b$ and $B = \intervalopenrel{R}{X}{a}{b}$.
        Then $a\mathrel{R}x$ and $x\mathrel{R}b$ and $\intervalopenrel{R}{X}{a}{b} \subseteq B$
            by \cref{intervalopen_rel,subseteq_refl}.
    \caseOf{$\exists b\in X. \exists a_0\in X. (\forall x\in X. a_0\mathrel{R}x \lor a_0 = x) \land B = \intervalclosedopenrel{R}{X}{a_0}{b}$.}
        Take $a_0, b\in X$ such that
            $(\forall x\in X. a_0\mathrel{R}x \lor a_0 = x)$ and $B = \intervalclosedopenrel{R}{X}{a_0}{b}$.
        Hence $a_0\mathrel{R}x$ and $x\mathrel{R}b$ and $\intervalopenrel{R}{X}{a_0}{b} \subseteq B$
            by \cref{intervalopen_rel,intervalclosedopen_rel,smallest_element,subseteq}.
    \caseOf{$\exists a\in X. \exists b_0\in X. (\forall x\in X. x\mathrel{R}b_0 \lor x = b_0) \land B = \intervalopenclosedrel{R}{X}{a}{b_0}$.}
        Take $a, b_0\in X$ such that
            $(\forall x\in X. x\mathrel{R}b_0 \lor x = b_0)$ and $B = \intervalopenclosedrel{R}{X}{a}{b_0}$.
        Therefore $a\mathrel{R}x$ and $x\mathrel{R}b_0$ and $\intervalopenrel{R}{X}{a}{b_0} \subseteq B$
            by \cref{intervalopen_rel,intervalopenclosed_rel,largest_element,subseteq}.
    \end{byCase}
\end{proof}

% from §14 helper lemma 10: maximum of two elements in a simple order
\begin{lemma}\label{max_of_two_elements}
    Assume $R$ is a simple order relation on $X$.
    Suppose $a_1, a_2\in X$.
    Then there exists $a\in X$ such that
        $(a = a_1 \lor a = a_2) \land
        (a_1\mathrel{R}a \lor a_1 = a) \land
        (a_2\mathrel{R}a \lor a_2 = a)$.
\end{lemma}
\begin{proof}
    Follows by \cref{simple_order_relation,connex}.
\end{proof}

% from §14 helper lemma 11: minimum of two elements in a simple order
\begin{lemma}\label{min_of_two_elements}
    Assume $R$ is a simple order relation on $X$.
    Suppose $b_1, b_2\in X$.
    Then there exists $b\in X$ such that
        $(b = b_1 \lor b = b_2) \land
        (b\mathrel{R}b_1 \lor b = b_1) \land
        (b\mathrel{R}b_2 \lor b = b_2)$.
\end{lemma}
\begin{proof}
    Follows by \cref{simple_order_relation,connex}.
\end{proof}

% from §14 helper lemma 12: left endpoint from a disjunction
\begin{lemma}\label{rel_from_or_left}
    Suppose $a = a_1$ or $a = a_2$.
    Suppose $a_1\mathrel{R}x$ and $a_2\mathrel{R}x$.
    Then $a\mathrel{R}x$.
\end{lemma}
\begin{proof}
    Trivial.
\end{proof}

% from §14 helper lemma 13: right endpoint from a disjunction
\begin{lemma}\label{rel_from_or_right}
    Suppose $b = b_1$ or $b = b_2$.
    Suppose $x\mathrel{R}b_1$ and $x\mathrel{R}b_2$.
    Then $x\mathrel{R}b$.
\end{lemma}
\begin{proof}
    Trivial.
\end{proof}

% from §14 Item 11: order basis is a basis on X
\begin{lemma}\label{order_basis_is_basis}
    Assume $R$ is a simple order relation on $X$.
    Assume $X$ has more than one element.
    Then $\orderBasis{R}{X}$ is a basis on $X$.
\end{lemma}
\begin{proof}
    Show that $\orderBasis{R}{X}\subseteq\pow{X}$.
    \begin{subproof}
        Follows by \cref{order_basis,subseteq}.
    \end{subproof}
    Show that for all $x\in X$ there exists $B_1\in\orderBasis{R}{X}$ such that $x\in B_1$.
    \begin{subproof}
        Fix $x\in X$.
        \begin{byCase}
        \caseOf{$x$ is a smallest element of $X$ with respect to $R$.}
            Take $y\in X$ such that $y\neq x$ by \cref{more_than_one_element_has_other}.
            $x\mathrel{R}y$ by \cref{smallest_element}.
            Let $B_1 = \intervalclosedopenrel{R}{X}{x}{y}$.
            Thus $B_1\in\orderBasis{R}{X}$ and $x\in B_1$ by \cref{intervalclosedopen_rel,subseteq,powerset_intro,smallest_element,order_basis}.
        \caseOf{$x$ is not a smallest element of $X$ with respect to $R$.}
            \begin{byCase}
            \caseOf{$x$ is a largest element of $X$ with respect to $R$.}
                Take $y\in X$ such that $y\neq x$ by \cref{more_than_one_element_has_other}.
                $y\mathrel{R}x$ by \cref{largest_element}.
                Let $B_1 = \intervalopenclosedrel{R}{X}{y}{x}$.
                $B_1\in\orderBasis{R}{X}$ and $x\in B_1$ by \cref{intervalopenclosed_rel,subseteq,powerset_intro,largest_element,order_basis}.
            \caseOf{$x$ is not a largest element of $X$ with respect to $R$.}
                Take $a\in X$ such that $a\mathrel{R}x$ by \cref{non_smallest_has_predecessor}.
                Take $b\in X$ such that $x\mathrel{R}b$ by \cref{non_largest_has_successor}.
                Let $B_1 = \intervalopenrel{R}{X}{a}{b}$.
                $B_1\in\orderBasis{R}{X}$ and $x\in B_1$ by \cref{simple_order_relation,transitive,intervalopen_rel,subseteq,powerset_intro,order_basis}.
            \end{byCase}
        \end{byCase}
    \end{subproof}
    Show that for all $B_1,B_2,x$ such that $B_1\in\orderBasis{R}{X}$ and $B_2\in\orderBasis{R}{X}$ and $x\in B_1\inter B_2$
        there exists $B_3\in\orderBasis{R}{X}$ such that $x\in B_3$ and $B_3\subseteq B_1\inter B_2$.
    \begin{subproof}
        Fix $B_1,B_2$.
        Fix $x$ such that $B_1\in\orderBasis{R}{X}$ and $B_2\in\orderBasis{R}{X}$ and $x\in B_1\inter B_2$.
        $x\in B_1$ and $x\in B_2$ by \cref{inter_elim_left,inter_elim_right}.
        Thus $x\in X$ by \cref{order_basis,powerset_elim,subseteq}.
        \begin{byCase}
        \caseOf{$x$ is a smallest element of $X$ with respect to $R$.}
            Take $b_1\in X$ such that $x\mathrel{R}b_1$ and $B_1 = \intervalclosedopenrel{R}{X}{x}{b_1}$ by \cref{orderbasis_elem_at_smallest}.
            Take $b_2\in X$ such that $x\mathrel{R}b_2$ and $B_2 = \intervalclosedopenrel{R}{X}{x}{b_2}$ by \cref{orderbasis_elem_at_smallest}.
            Take $b\in X$ such that
                $(b = b_1 \lor b = b_2) \land
                (b\mathrel{R}b_1 \lor b = b_1) \land
                (b\mathrel{R}b_2 \lor b = b_2)$
                by \cref{min_of_two_elements}.
            Let $B_3 = \intervalclosedopenrel{R}{X}{x}{b}$.
            $B_3\in\orderBasis{R}{X}$ and $x\in B_3$ and $B_3\subseteq B_1\inter B_2$
                by \cref{order_basis,intervalclosedopen_rel,subseteq,powerset_intro,smallest_element,simple_order_relation,intervalclosedopenrel_subset_by_right,subseteq_inter_iff}.
        \caseOf{$x$ is not a smallest element of $X$ with respect to $R$.}
            \begin{byCase}
            \caseOf{$x$ is a largest element of $X$ with respect to $R$.}
                Take $a_1\in X$ such that $a_1\mathrel{R}x$ and $B_1 = \intervalopenclosedrel{R}{X}{a_1}{x}$ by \cref{orderbasis_elem_at_largest}.
                Take $a_2\in X$ such that $a_2\mathrel{R}x$ and $B_2 = \intervalopenclosedrel{R}{X}{a_2}{x}$ by \cref{orderbasis_elem_at_largest}.
                Take $a\in X$ such that
                    $(a = a_1 \lor a = a_2) \land
                    (a_1\mathrel{R}a \lor a_1 = a) \land
                    (a_2\mathrel{R}a \lor a_2 = a)$
                    by \cref{max_of_two_elements}.
                Let $B_3 = \intervalopenclosedrel{R}{X}{a}{x}$.
                $B_3\in\orderBasis{R}{X}$ and $x\in B_3$ and $B_3\subseteq B_1\inter B_2$
                    by \cref{order_basis,intervalopenclosed_rel,subseteq,powerset_intro,largest_element,simple_order_relation,intervalopenclosedrel_subset_by_left,subseteq_inter_iff}.
            \caseOf{$x$ is not a largest element of $X$ with respect to $R$.}
                Take $a_1,b_1\in X$ such that $a_1\mathrel{R}x$ and $x\mathrel{R}b_1$ and
                    $\intervalopenrel{R}{X}{a_1}{b_1} \subseteq B_1$ by \cref{orderbasis_elem_contains_open_interval}.
                Take $a_2,b_2\in X$ such that $a_2\mathrel{R}x$ and $x\mathrel{R}b_2$ and
                    $\intervalopenrel{R}{X}{a_2}{b_2} \subseteq B_2$ by \cref{orderbasis_elem_contains_open_interval}.
                Take $a\in X$ such that
                    $(a = a_1 \lor a = a_2) \land
                    (a_1\mathrel{R}a \lor a_1 = a) \land
                    (a_2\mathrel{R}a \lor a_2 = a)$
                    by \cref{max_of_two_elements}.
                Take $b\in X$ such that
                    $(b = b_1 \lor b = b_2) \land
                    (b\mathrel{R}b_1 \lor b = b_1) \land
                    (b\mathrel{R}b_2 \lor b = b_2)$
                    by \cref{min_of_two_elements}.
                Let $B_3 = \intervalopenrel{R}{X}{a}{b}$.
                $R$ is transitive by \cref{simple_order_relation}.
                Thus $a\mathrel{R}b$ by \hyperref[transitive]{transitivity}.
                $B_3\in\orderBasis{R}{X}$ and $x\in B_3$
                    by \cref{intervalopen_rel,subseteq,powerset_intro,order_basis,rel_from_or_left,rel_from_or_right}.
                $B_3\subseteq B_1\inter B_2$
                    by \cref{intervalopenrel_subset_by_endpoints,subseteq_transitive,subseteq_inter_iff}.
            \end{byCase}
        \end{byCase}
    \end{subproof}
    $\orderBasis{R}{X}$ is a basis on $X$ by \cref{basis_on}.
\end{proof}

% from §14 Item 12: open right ray
\begin{definition}\label{openray_right}
    $\openrayright{R}{X}{a} = \{x\in X \mid a\mathrel{R}x\}$.
\end{definition}

% from §14 Item 13: open left ray
\begin{definition}\label{openray_left}
    $\openrayleft{R}{X}{a} = \{x\in X \mid x\mathrel{R}a\}$.
\end{definition}

% from §14 Item 14: closed right ray
\begin{definition}\label{closedray_right}
    $\closedrayright{R}{X}{a} = \{x\in X \mid a\mathrel{R}x \lor x = a\}$.
\end{definition}

% from §14 Item 15: closed left ray
\begin{definition}\label{closedray_left}
    $\closedrayleft{R}{X}{a} = \{x\in X \mid x\mathrel{R}a \lor x = a\}$.
\end{definition}

% from §14 Item 16: open rays
\begin{definition}\label{open_rays}
    $\openRays{R}{X} = \{U\in\pow{X} \mid \exists a\in X. (U = \openrayright{R}{X}{a} \lor U = \openrayleft{R}{X}{a})\}$.
\end{definition}

% from §14 helper lemma 14: open interval as intersection of open rays
\begin{lemma}\label{intervalopenrel_as_intersection_openrays}
    $\intervalopenrel{R}{X}{a}{b} = \openrayright{R}{X}{a} \inter \openrayleft{R}{X}{b}$.
\end{lemma}
\begin{proof}
    Follows by \cref{intervalopen_rel,openray_right,openray_left,inter,setext}.
\end{proof}

% from §14 helper lemma 15: closed-open interval at a smallest element is a left ray
\begin{lemma}\label{intervalclosedopenrel_smallest_eq_openrayleft}
    Assume $a_0$ is a smallest element of $X$ with respect to $R$.
    Then $\intervalclosedopenrel{R}{X}{a_0}{b} = \openrayleft{R}{X}{b}$.
\end{lemma}
\begin{proof}
    Follows by \cref{intervalclosedopen_rel,openray_left,smallest_element,setext}.
\end{proof}

% from §14 helper lemma 16: open-closed interval at a largest element is a right ray
\begin{lemma}\label{intervalopenclosedrel_largest_eq_openrayright}
    Assume $b_0$ is a largest element of $X$ with respect to $R$.
    Then $\intervalopenclosedrel{R}{X}{a}{b_0} = \openrayright{R}{X}{a}$.
\end{lemma}
\begin{proof}
    Follows by \cref{intervalopenclosed_rel,openray_right,largest_element,setext}.
\end{proof}

% from §14 helper lemma 17: open rays form a subbasis on X
\begin{lemma}\label{open_rays_subbasis_on}
    Assume $R$ is a simple order relation on $X$.
    Assume $X$ has more than one element.
    Then $\openRays{R}{X}$ is a subbasis on $X$.
\end{lemma}
\begin{proof}
    Let $S = \openRays{R}{X}$.
    $S\subseteq\pow{X}$ by \cref{open_rays,subseteq}.
    $\unions{S}\subseteq X$ by \cref{unions_subseteq_of_powerset_is_subseteq}.
    For all $x\in X$ there exists $y\in X$ such that $y\neq x$ by \cref{more_than_one_element_has_other}.
    Show that for all $x$ such that $x\in X$ we have $x\in \unions{S}$.
    \begin{subproof}
        Fix $x$ such that $x\in X$.
        Take $y\in X$ such that $y\neq x$ by \cref{more_than_one_element_has_other}.
        Then $y\mathrel{R} x \lor x\mathrel{R} y$ by \cref{simple_order_relation,connex}.
        \begin{byCase}
        \caseOf{$y\mathrel{R} x$.}
            $x\in\openrayright{R}{X}{y}$ by \cref{openray_right}.
            $\openrayright{R}{X}{y}\in S$ by \cref{openray_right,subseteq,powerset_intro,open_rays}.
            $x\in\unions{S}$ by \cref{unions_iff}.
        \caseOf{$x\mathrel{R} y$.}
            $x\in\openrayleft{R}{X}{y}$ by \cref{openray_left}.
            $\openrayleft{R}{X}{y}\in S$ by \cref{openray_left,subseteq,powerset_intro,open_rays}.
            $x\in\unions{S}$ by \cref{unions_iff}.
        \end{byCase}
    \end{subproof}
    $\unions{S} = X$ by \cref{subseteq,subseteq_antisymmetric}.
    $S$ is a subbasis on $X$ by \cref{subbasis_on}.
\end{proof}

% from §14 helper lemma 18: order basis elements are finite intersections of open rays
\begin{lemma}\label{order_basis_in_finite_intersections}
    Assume $R$ is a simple order relation on $X$.
    Assume $X$ has more than one element.
    Suppose $B\in\orderBasis{R}{X}$.
    Then $B\in\finiteIntersections{\openRays{R}{X}}$.
\end{lemma}
\begin{proof}
    Let $S = \openRays{R}{X}$.
    $B\in\pow{\unions{S}}$ and
        $((\exists a, b\in X. a\mathrel{R}b \land B = \intervalopenrel{R}{X}{a}{b}) \lor
        (\exists b\in X. \exists a_0\in X. (\forall x\in X. a_0\mathrel{R}x \lor a_0 = x) \land B = \intervalclosedopenrel{R}{X}{a_0}{b}) \lor
        (\exists a\in X. \exists b_0\in X. (\forall x\in X. x\mathrel{R}b_0 \lor x = b_0) \land B = \intervalopenclosedrel{R}{X}{a}{b_0}))$
        by \cref{open_rays_subbasis_on,subbasis_on,order_basis,pow_iff,powerset_intro}.
    \begin{byCase}
    \caseOf{$\exists a, b\in X. a\mathrel{R}b \land B = \intervalopenrel{R}{X}{a}{b}$.}
        Take $a, b\in X$ such that $a\mathrel{R}b$ and $B = \intervalopenrel{R}{X}{a}{b}$.
        Let $F = \{\openrayright{R}{X}{a}, \openrayleft{R}{X}{b}\}$.
        $F\subseteq S$ by \cref{openray_right,openray_left,subseteq,powerset_intro,open_rays,upair_iff}.
        $F$ is finite and $F$ is inhabited by \cref{pair_is_finite,upair_intro_left}.
        Thus $B\in\finiteIntersections{S}$
            by \cref{intervalopenrel_as_intersection_openrays,inter_as_inters,finite_intersections}.
    \caseOf{$\exists b\in X. \exists a_0\in X. (\forall x\in X. a_0\mathrel{R}x \lor a_0 = x) \land B = \intervalclosedopenrel{R}{X}{a_0}{b}$.}
        Take $a_0, b\in X$ such that
            $(\forall x\in X. a_0\mathrel{R}x \lor a_0 = x)$ and $B = \intervalclosedopenrel{R}{X}{a_0}{b}$.
        Hence $B\in\finiteIntersections{S}$
            by \cref{openray_left,subseteq,powerset_intro,open_rays,singleton_subset_intro,pair_is_finite,upair_intro_left,subseteq_finite_is_finite,singleton_intro,smallest_element,intervalclosedopenrel_smallest_eq_openrayleft,inters_singleton,finite_intersections}.
    \caseOf{$\exists a\in X. \exists b_0\in X. (\forall x\in X. x\mathrel{R}b_0 \lor x = b_0) \land B = \intervalopenclosedrel{R}{X}{a}{b_0}$.}
        Take $a, b_0\in X$ such that
            $(\forall x\in X. x\mathrel{R}b_0 \lor x = b_0)$ and $B = \intervalopenclosedrel{R}{X}{a}{b_0}$.
        Therefore $B\in\finiteIntersections{S}$
            by \cref{openray_right,subseteq,powerset_intro,open_rays,singleton_subset_intro,pair_is_finite,upair_intro_left,subseteq_finite_is_finite,singleton_intro,largest_element,intervalopenclosedrel_largest_eq_openrayright,inters_singleton,finite_intersections}.
    \end{byCase}
\end{proof}

% from §14 helper lemma 18a: first initial segment
\begin{lemma}\label{initial_segment_one}
    $\initialSegment{1} = \{1\}$.
\end{lemma}
\begin{proof}
    For all $x$ we have $x\in\initialSegment{1}$ iff $x\in\{1\}$
        by \cref{initial_segment_naturalsplus,real_one_in_naturals,real_one_leq_naturalsplus,real_naturals_subset_reals,subseteq,singleton_iff,real_lt_asym}.
    Follows by \cref{setext}.
\end{proof}

% from §14 helper lemma 18b: initial segments are nested
\begin{lemma}\label{naturalsplus_lt_succ}
    Suppose $n\in\naturalsPlus$.
    Then $n < n + 1$.
\end{lemma}
\begin{proof}
    $n\in\reals$ and $1\in\reals$ and $0\in\reals$
        by \cref{real_naturals_subset_reals,real_one_in_naturals,subseteq,real_add_ident}.
    $0 < 1$ by \cref{real_zero_lt_one}.
    $0 + n < 1 + n$ by \cref{real_lt_add}.
    $0 + n = n$ and $1 + n = n + 1$ by \cref{real_add_ident,real_add_comm}.
    Follows by \cref{real_lt_subst_left,real_lt_subst_right}.
\end{proof}

\begin{lemma}\label{initial_segment_subset_succ}
    Suppose $n\in\naturalsPlus$.
    Then $\initialSegment{n}\subseteq \initialSegment{n + 1}$.
\end{lemma}
\begin{proof}
    Show that for all $x\in\initialSegment{n}$ we have $x\in\initialSegment{n + 1}$.
    \begin{subproof}
        Fix $x\in\initialSegment{n}$.
        $x\in\naturalsPlus$ and $x \leq n$ by \cref{initial_segment_naturalsplus}.
        $x\in\reals$ and $n\in\reals$ and $n + 1\in\reals$
            by \cref{real_naturals_subset_reals,real_naturals_closed_succ,subseteq}.
        $n < n + 1$ by \cref{naturalsplus_lt_succ}.
        $x < n + 1$ or $x = n + 1$ by \cref{real_leq_split,real_lt_trans,real_lt_subst_left}.
        Thus $x\in\initialSegment{n + 1}$ by \cref{initial_segment_naturalsplus,real_naturals_closed_succ}.
    \end{subproof}
    Follows by \cref{subseteq}.
\end{proof}

% from §14 helper lemma 18c: split a successor initial segment
\begin{lemma}\label{initial_segment_succ_split}
    Suppose $n\in\naturalsPlus$.
    Suppose $x\in\initialSegment{n + 1}$.
    Then $x\in\initialSegment{n}$ or $x = n + 1$.
\end{lemma}
\begin{proof}
    Follows by \cref{initial_segment_naturalsplus,real_succ_discrete}.
\end{proof}

% from §14 helper lemma 18d: initial segments are inhabited
\begin{lemma}\label{initial_segment_inhabited}
    Suppose $n\in\naturalsPlus$.
    Then $\initialSegment{n}$ is inhabited.
\end{lemma}
\begin{proof}
    $1\in\initialSegment{n}$ by \cref{initial_segment_naturalsplus,real_one_in_naturals,real_one_leq_naturalsplus}.
\end{proof}

% from §14 helper lemma 19: finite intersections in a topology are open
\begin{lemma}\label{finite_intersections_open_in_topology}
    Assume $T$ is a topology on $X$.
    Suppose $F\subseteq T$.
    Suppose $F$ is finite and inhabited.
    Then $\inters{F}\in T$.
\end{lemma}
\begin{proof}
    Take $k\in\naturalsPlus$ such that there exists a bijection from $\initialSegment{k}$ to $F$ by \cref{finite}.
    Let $A = \{ n\in\naturalsPlus \mid \text{for all $G,g$ such that $g$ is a bijection from $\initialSegment{n}$ to $G$ and $G\subseteq T$ and $G$ is inhabited we have $\inters{G}\in T$} \}$.
    Show that for all $G,g$ such that $g$ is a bijection from $\initialSegment{1}$ to $G$ and $G\subseteq T$ and $G$ is inhabited we have $\inters{G}\in T$.
    \begin{subproof}
        Fix $G$.
        Fix $g$ such that $g$ is a bijection from $\initialSegment{1}$ to $G$ and $G\subseteq T$ and $G$ is inhabited.
        $g$ is a function by \cref{bijection_is_function}.
        $\dom{g} = \initialSegment{1}$ by \cref{bijections_dom}.
        $1\in\dom{g}$ by \cref{initial_segment_one,singleton_intro}.
        Let $a = g(1)$.
        $(1,a)\in g$ by \cref{function_apply_elim}.
        $g\in\surj{\initialSegment{1}}{G}$ by \cref{bijections}.
        $\ran{g} = G$ by \cref{ran_of_surj}.
        $a\in G$ by \cref{ran_iff}.
        $a\in T$ by \cref{elem_subseteq}.
        Show that for all $z\in G$ we have $z\in\{a\}$.
        \begin{subproof}
            Fix $z\in G$.
            Take $x\in\initialSegment{1}$ such that $g(x)=z$ by \cref{bijections,surj}.
            $x = 1$ by \cref{initial_segment_one,singleton_iff}.
            $a = z$.
            Thus $z\in\{a\}$ by \cref{singleton_iff}.
        \end{subproof}
        $G\subseteq\{a\}$ by \cref{subseteq}.
        $\{a\}\subseteq G$ by \cref{singleton_subset_intro}.
        $G = \{a\}$ by \cref{subseteq_antisymmetric}.
        $\inters{G}\in T$ by \cref{inters_singleton}.
    \end{subproof}
    $1\in A$.
    Show that for all $n\in A$ we have $n + 1\in A$.
    \begin{subproof}
            Fix $n\in A$.
            Show that for all $G,g$ such that $g$ is a bijection from $\initialSegment{n + 1}$ to $G$ and $G\subseteq T$ and $G$ is inhabited we have $\inters{G}\in T$.
            \begin{subproof}
                Fix $G$.
                Fix $g$ such that $g$ is a bijection from $\initialSegment{n + 1}$ to $G$ and $G\subseteq T$ and $G$ is inhabited.
                $g$ is a function by \cref{bijection_is_function}.
                $\dom{g} = \initialSegment{n + 1}$ by \cref{bijections_dom}.
                $n + 1\in\dom{g}$ by \cref{initial_segment_naturalsplus,real_naturals_closed_succ}.
                Let $a = g(n + 1)$.
                Let $G_1 = \img{g}{\initialSegment{n}}$.
                $\initialSegment{n}\subseteq \initialSegment{n + 1}$ by \cref{initial_segment_subset_succ}.
                $G_1\subseteq G$ by \cref{img_subseteq,bijections_dom,img_dom,bijections,ran_of_surj}.
                $G_1\subseteq T$ by \cref{subseteq_transitive}.
                $g\in\surj{\initialSegment{n + 1}}{G}$ by \cref{bijections}.
                $\ran{g} = G$ by \cref{ran_of_surj}.
                $(n + 1,a)\in g$ by \cref{function_apply_elim}.
                $a\in G$ by \cref{ran_iff}.
                $a\in T$ by \cref{elem_subseteq}.
                Show that for all $z\in G$ we have $z\in G_1\union\{a\}$.
                    \begin{subproof}
                        Fix $z\in G$.
                        There exists $x\in\initialSegment{n + 1}$ such that $g(x)=z$ by \cref{bijections,surj}.
                        $x\in\initialSegment{n}$ or $x = n + 1$ by \cref{initial_segment_succ_split}.
                        \begin{byCase}
                        \caseOf{$x\in\initialSegment{n}$.}
                            $z\in G_1\union\{a\}$ by \cref{function_apply_elim,img_iff,union_intro_left}.
                        \caseOf{$x = n + 1$.}
                            $z\in G_1\union\{a\}$ by \cref{function_apply_elim,union_intro_right,singleton_intro}.
                        \end{byCase}
                \end{subproof}
                Thus $G = G_1\union\{a\}$
                    by \cref{subseteq,singleton_subset_intro,union_subsets_is_subset,subseteq_antisymmetric}.
                Take $x_0$ such that $x_0\in\initialSegment{n}$ by \cref{initial_segment_inhabited}.
                $x_0\in\dom{g}$ by \cref{initial_segment_subset_succ,bijections_dom,subseteq}.
                $g(x_0)\in G_1$ by \cref{function_apply_elim,img_iff}.
                $G_1$ is inhabited.
                Let $h = \restrl{g}{\initialSegment{n}}$.
                $h$ is a function by \cref{restrl_of_function_is_function}.
                For all $x\in\initialSegment{n}$ we have $x\in\dom{g}$ by \cref{initial_segment_subset_succ,bijections_dom,subseteq}.
                $\dom{h} = \initialSegment{n}$ by \cref{subseteq,dom_of_restrl_eq_inter,inter_absorb_supseteq_left}.
                $h\in\funs{\initialSegment{n}}{G_1}$ by \cref{function_apply_elim,restrl_subseteq,elem_subseteq,img_iff,funs_intro}.
                For all $y\in G_1$ there exists $x\in\initialSegment{n}$ such that $h(x) = y$
                    by \cref{img_iff,function_apply_iff,restrl_of_function_apply}.
                $h$ is a bijection from $\initialSegment{n}$ to $G_1$ by \cref{surj,restrl_of_function_apply,bijections,injective_function}.
                Therefore $\inters{G_1}\in T$.
                $G = \cons{a}{G_1}$ by \cref{union_comm,union_eq_cons}.
                $\inters{G}\in T$ by \cref{inters_cons,topology_on}.
            \end{subproof}
    \end{subproof}
    $A\subseteq\naturalsPlus$ by \cref{subseteq}.
    $k\in A$ by \cref{real_naturals_induction}.
    For all $G,g$ such that $g$ is a bijection from $\initialSegment{k}$ to $G$ and $G\subseteq T$ and $G$ is inhabited we have $\inters{G}\in T$.
    $\inters{F}\in T$.
\end{proof}

% from §14 Item 17: open rays form a subbasis
\begin{lemma}\label{open_rays_subbasis}
    Assume $R$ is a simple order relation on $X$.
    Assume $X$ has more than one element.
    Assume $T$ is the order topology on $X$ with respect to $R$.
    Then $\subbasisTopology{\openRays{R}{X}}{X} = T$.
\end{lemma}
\begin{proof}
    Let $S = \openRays{R}{X}$.
    Let $B = \orderBasis{R}{X}$.
    Let $T_1 = \subbasisTopology{S}{X}$.
    $T = \basisTopology{B}{X}$ by \cref{order_topology}.
    $B$ is a basis for $T$ on $X$ by \cref{order_basis_is_basis,basis_for_topology,order_topology}.
    $\finiteIntersections{S}$ is a basis for $T_1$ on $X$
        by \cref{open_rays_subbasis_on,subbasis_finite_intersections_basis,basis_for_topology,topology_generated_by_subbasis}.
    Show that for all $x,B_2$ such that $x\in X$ and $B_2\in B$ and $x\in B_2$
        there exists $B_3\in\finiteIntersections{S}$ such that $x\in B_3$ and $B_3\subseteq B_2$.
    \begin{subproof}
        Follows by \cref{order_basis_in_finite_intersections,subseteq_refl}.
    \end{subproof}
    Thus $T\subseteq T_1$ by \cref{basis_finer_criterion,finer_topology}.
    Show that for all $U\in S$ we have $U\in T$.
    \begin{subproof}
        Fix $U\in S$.
        Take $a\in X$ such that $U = \openrayright{R}{X}{a}$ or $U = \openrayleft{R}{X}{a}$ by \cref{open_rays}.
        \begin{byCase}
        \caseOf{$U = \openrayright{R}{X}{a}$.}
            \begin{byCase}
            \caseOf{$\exists b_0\in X. (\forall x\in X. x\mathrel{R}b_0 \lor x = b_0)$.}
                Take $b_0\in X$ such that $b_0$ is a largest element of $X$ with respect to $R$.
                $U = \intervalopenclosedrel{R}{X}{a}{b_0}$ by \cref{intervalopenclosedrel_largest_eq_openrayright}.
                $\exists a\in X. \exists b_0\in X. (\forall x\in X. x\mathrel{R}b_0 \lor x = b_0) \land U = \intervalopenclosedrel{R}{X}{a}{b_0}$.
                Show that $U\in B$ and $U\in T$.
                \begin{subproof}
                    Follows by \cref{intervalopenclosed_rel,subseteq,powerset_intro,order_basis,basis_elem_open_in_generated,order_basis_is_basis}.
                \end{subproof}
            \caseOf{it is wrong that there exists $b_0\in X$ such that for all $x\in X$ we have $x\mathrel{R}b_0$ or $x = b_0$.}
                Show that for all $x\in U$ there exists $B_1\in B$ such that $x\in B_1$ and $B_1\subseteq U$.
                \begin{subproof}
                    Fix $x\in U$.
                    $x\in X$ and $a\mathrel{R}x$ by \cref{openray_right}.
                    $x$ is not a largest element of $X$ with respect to $R$.
                    Take $b\in X$ such that $x\mathrel{R}b$ by \cref{non_largest_has_successor}.
                    Let $B_1 = \intervalopenrel{R}{X}{a}{b}$.
                    Hence $B_1\in B$ and $x\in B_1$ and $B_1\subseteq U$
                        by \cref{simple_order_relation,transitive,order_basis,intervalopen_rel,subseteq,powerset_intro,openray_right}.
                \end{subproof}
                $U\in T$ by \cref{openray_right,subseteq,powerset_intro,topology_generated_by_basis}.
            \end{byCase}
        \caseOf{$U = \openrayleft{R}{X}{a}$.}
            \begin{byCase}
            \caseOf{$\exists a_0\in X. (\forall x\in X. a_0\mathrel{R}x \lor a_0 = x)$.}
                Take $a_0\in X$ such that $a_0$ is a smallest element of $X$ with respect to $R$.
                $U = \intervalclosedopenrel{R}{X}{a_0}{a}$ by \cref{intervalclosedopenrel_smallest_eq_openrayleft}.
                $\exists b\in X. \exists a_0\in X. (\forall x\in X. a_0\mathrel{R}x \lor a_0 = x) \land U = \intervalclosedopenrel{R}{X}{a_0}{b}$.
                Show that $U\in B$ and $U\in T$.
                \begin{subproof}
                    Follows by \cref{intervalclosedopen_rel,subseteq,powerset_intro,order_basis,basis_elem_open_in_generated,order_basis_is_basis}.
                \end{subproof}
            \caseOf{it is wrong that there exists $a_0\in X$ such that for all $x\in X$ we have $a_0\mathrel{R}x$ or $a_0 = x$.}
                Show that for all $x\in U$ there exists $B_1\in B$ such that $x\in B_1$ and $B_1\subseteq U$.
                \begin{subproof}
                    Fix $x\in U$.
                    $x\in X$ and $x\mathrel{R}a$ by \cref{openray_left}.
                    $x$ is not a smallest element of $X$ with respect to $R$.
                    Take $b\in X$ such that $b\mathrel{R}x$ by \cref{non_smallest_has_predecessor}.
                    Let $B_1 = \intervalopenrel{R}{X}{b}{a}$.
                    Thus $B_1\in B$ and $x\in B_1$ and $B_1\subseteq U$
                        by \cref{simple_order_relation,transitive,order_basis,intervalopen_rel,subseteq,powerset_intro,openray_left}.
                \end{subproof}
                $U\in T$ by \cref{openray_left,subseteq,powerset_intro,topology_generated_by_basis}.
            \end{byCase}
        \end{byCase}
    \end{subproof}
    Show that for all $x,B_2$ such that $x\in X$ and $B_2\in\finiteIntersections{S}$ and $x\in B_2$
        there exists $B_3\in B$ such that $x\in B_3$ and $B_3\subseteq B_2$.
    \begin{subproof}
        Follows by \cref{finite_intersections,subseteq,order_basis_is_basis,basis_topology_is_topology,order_topology,finite_intersections_open_in_topology,topology_generated_by_basis}.
    \end{subproof}
    Hence $T_1\subseteq T$ by \cref{basis_finer_criterion,finer_topology}.
    $T_1 = T$ by \cref{subseteq_antisymmetric}.
\end{proof}

\section{§15 The Product Topology on $X \times Y$}

% from §15 Item 1: product basis on X×Y
\begin{definition}\label{product_basis}
    $\productBasis{T}{T_1}{X}{Y} = \{B\in\pow{X\times Y} \mid \exists U\in T. \exists V\in T_1. B = U\times V\}$.
\end{definition}

% from §15 Item 2: product topology on X×Y
\begin{definition}\label{product_topology}
    $\productTopology{T}{T_1}{X}{Y} = \basisTopology{\productBasis{T}{T_1}{X}{Y}}{X\times Y}$.
\end{definition}

% from §15 Item 3: product basis is a basis
\begin{lemma}\label{product_basis_is_basis}
    Assume $T$ is a topology on $X$ and $T_1$ is a topology on $Y$.
    Then $\productBasis{T}{T_1}{X}{Y}$ is a basis on $X\times Y$.
\end{lemma}
\begin{proof}
    Let $B = \productBasis{T}{T_1}{X}{Y}$.
    We have $B\subseteq\pow{X\times Y}$ by \cref{product_basis,subseteq}.
    For all $p\in X\times Y$ there exists $B_1\in B$ such that $p\in B_1$
        by \cref{topology_on,product_basis,subseteq_refl,times_subseteq_left,times_subseteq_right,subseteq_transitive,powerset_intro}.
    Show that for all $B_1,B_2,p$ such that $B_1\in B$ and $B_2\in B$ and $p\in B_1\inter B_2$
        there exists $B_3\in B$ such that $p\in B_3$ and $B_3\subseteq B_1\inter B_2$.
    \begin{subproof}
        Fix $B_1,B_2$.
        Fix $p$ such that $B_1\in B$ and $B_2\in B$ and $p\in B_1\inter B_2$.
        Take $U_1,V_1$ such that $U_1\in T$ and $V_1\in T_1$ and $B_1 = U_1\times V_1$ by \cref{product_basis}.
        Take $U_2,V_2$ such that $U_2\in T$ and $V_2\in T_1$ and $B_2 = U_2\times V_2$ by \cref{product_basis}.
        Let $B_3 = (U_1\inter U_2)\times (V_1\inter V_2)$.
        Thus $B_3\in B$ and $p\in B_3$ and $B_3\subseteq B_1\inter B_2$
            by \cref{topology_on,powerset_elim,subseteq,inter_subseteq,times_subseteq_left,times_subseteq_right,subseteq_transitive,powerset_intro,product_basis,inter_times_elim,inter_times_intro}.
    \end{subproof}
    $B$ is a basis on $X\times Y$ by \cref{basis_on}.
\end{proof}

% from §15 Item 4: bases for X and Y give a basis for X×Y
\begin{lemma}\label{product_basis_from_bases}
    Assume $B$ is a basis for $T$ on $X$ and $C$ is a basis for $T_1$ on $Y$.
    Let $D = \{D_1\in\pow{X\times Y} \mid \exists B_1\in B. \exists C_1\in C. D_1 = B_1\times C_1\}$.
    Then $D$ is a basis for $\productTopology{T}{T_1}{X}{Y}$ on $X\times Y$.
\end{lemma}
\begin{proof}
    Let $T_2 = \productTopology{T}{T_1}{X}{Y}$.
    $B$ is a basis on $X$ and $C$ is a basis on $Y$
        and $T = \basisTopology{B}{X}$ and $T_1 = \basisTopology{C}{Y}$ by \cref{basis_for_topology}.
    $T$ is a topology on $X$ and $T_1$ is a topology on $Y$
        and $\productBasis{T}{T_1}{X}{Y}$ is a basis on $X\times Y$
        and $T_2$ is a topology on $X\times Y$
        by \cref{basis_topology_is_topology,product_basis_is_basis,product_topology}.
    For all $D_1$ such that $D_1\in D$ we have $D_1$ is open in $X\times Y$ with respect to $T_2$
        by \cref{subseteq,basis_elem_open_in_generated,basis_for_topology,open_in,product_basis,product_topology}.
    Show that for all $U,p$ such that $U\in T_2$ and $p\in U$ there exists $D_1\in D$ such that $p\in D_1$ and $D_1\subseteq U$.
    \begin{subproof}
        Fix $U$.
        Fix $p$ such that $U\in T_2$ and $p\in U$.
        Take $B_0\in\productBasis{T}{T_1}{X}{Y}$ such that $p\in B_0$ and $B_0\subseteq U$
            by \cref{product_topology,topology_generated_by_basis}.
        Take $U_1,V_1$ such that $U_1\in T$ and $V_1\in T_1$ and $B_0 = U_1\times V_1$ by \cref{product_basis}.
        Take $x,y$ such that $p = (x,y)$ and $x\in U_1$ and $y\in V_1$ by \cref{times_elem_is_tuple}.
        Take $B_1\in B$ such that $x\in B_1$ and $B_1\subseteq U_1$ by \cref{basis_for_topology,topology_generated_by_basis,open_in}.
        Take $C_1\in C$ such that $y\in C_1$ and $C_1\subseteq V_1$ by \cref{basis_for_topology,topology_generated_by_basis,open_in}.
        Let $D_1 = B_1\times C_1$.
        Thus $D_1\in D$ and $p\in D_1$ and $D_1\subseteq U$
            by \cref{basis_on,subseteq,powerset_elim,times_subseteq_left,times_subseteq_right,subseteq_transitive,powerset_intro,times,inter_intro}.
    \end{subproof}
    $D$ is a basis for $T_2$ on $X\times Y$ by \cref{open_collection_basis}.
\end{proof}

% from §15 Item 5: projection onto the first factor
\begin{definition}\label{projection_first}
    $\piOne{X}{Y} = \{((x,y),x) \mid x\in X, y\in Y\}$.
\end{definition}

% from §15 Item 6: projection onto the second factor
\begin{definition}\label{projection_second}
    $\piTwo{X}{Y} = \{((x,y),y) \mid x\in X, y\in Y\}$.
\end{definition}

% from §15 helper lemma 1: preimage of the first projection
\begin{lemma}\label{preimg_pi1}
    Suppose $U\subseteq X$.
    Then $\preimg{\piOne{X}{Y}}{U} = U\times Y$.
\end{lemma}
\begin{proof}
    Show that for all $p$ we have $p\in\preimg{\piOne{X}{Y}}{U}$ iff $p\in U\times Y$.
    \begin{subproof}
        Follows by \cref{preimg_iff,projection_first,pair_eq_iff,times,times_elem_is_tuple,subseteq}.
    \end{subproof}
    Follows by \cref{setext}.
\end{proof}

% from §15 helper lemma 2: preimage of the second projection
\begin{lemma}\label{preimg_pi2}
    Suppose $V\subseteq Y$.
    Then $\preimg{\piTwo{X}{Y}}{V} = X\times V$.
\end{lemma}
\begin{proof}
    Show that for all $p$ we have $p\in\preimg{\piTwo{X}{Y}}{V}$ iff $p\in X\times V$.
    \begin{subproof}
        Follows by \cref{preimg_iff,projection_second,pair_eq_iff,times,times_elem_is_tuple,subseteq}.
    \end{subproof}
    Follows by \cref{setext}.
\end{proof}

% from §15 Item 7: left projection subbasis collection
\begin{definition}\label{product_subbasis_left}
    $\productSubbasisLeft{T}{X}{Y} = \{A \mid \exists U\in T. A = \preimg{\piOne{X}{Y}}{U}\}$.
\end{definition}

% from §15 Item 8: right projection subbasis collection
\begin{definition}\label{product_subbasis_right}
    $\productSubbasisRight{T_1}{X}{Y} = \{A \mid \exists V\in T_1. A = \preimg{\piTwo{X}{Y}}{V}\}$.
\end{definition}

% from §15 Item 9: projection subbasis for the product topology
\begin{lemma}\label{product_topology_subbasis}
    Assume $T$ is a topology on $X$ and $T_1$ is a topology on $Y$.
    Let $S = \productSubbasisLeft{T}{X}{Y}\union \productSubbasisRight{T_1}{X}{Y}$.
    Then $\subbasisTopology{S}{X\times Y} = \productTopology{T}{T_1}{X}{Y}$.
\end{lemma}
\begin{proof}
    Let $T_2 = \productTopology{T}{T_1}{X}{Y}$.
    Let $T_3 = \subbasisTopology{S}{X\times Y}$.
    Let $S_1 = \productSubbasisLeft{T}{X}{Y}$.
    Let $S_2 = \productSubbasisRight{T_1}{X}{Y}$.
    $\productBasis{T}{T_1}{X}{Y}$ is a basis on $X\times Y$ and $T_2$ is a topology on $X\times Y$
        by \cref{product_basis_is_basis,basis_topology_is_topology,product_topology}.
    For all $A$ such that $A\in S_1$ we have $A\in\pow{X\times Y}$
        by \cref{product_subbasis_left,topology_on,subseteq,powerset_elim,preimg_pi1,times_subseteq_left,powerset_intro}.
    For all $A$ such that $A\in S_2$ we have $A\in\pow{X\times Y}$
        by \cref{product_subbasis_right,topology_on,subseteq,powerset_elim,preimg_pi2,times_subseteq_right,powerset_intro}.
    $S\subseteq\pow{X\times Y}$ by \cref{union_iff,subseteq}.
    Let $A = \preimg{\piOne{X}{Y}}{X}$.
    $A\in S$ and $A = X\times Y$ by \cref{topology_on,product_subbasis_left,union_intro_left,preimg_pi1,subseteq_refl}.
    $\unions{S} = X\times Y$
        by \cref{unions_iff,subseteq,unions_subseteq_of_powerset_is_subseteq,subseteq_antisymmetric}.
    $S$ is a subbasis on $X\times Y$ and $\finiteIntersections{S}$ is a basis on $X\times Y$
        by \cref{subbasis_on,subbasis_finite_intersections_basis}.
    $\finiteIntersections{S}$ is a basis for $T_3$ on $X\times Y$ by \cref{basis_for_topology,topology_generated_by_subbasis}.
    Show that for all $U$ such that $U\in S_1$ we have $U\in T_2$.
    \begin{subproof}
        Fix $U$ such that $U\in S_1$.
        Take $V\in T$ such that $U = \preimg{\piOne{X}{Y}}{V}$ by \cref{product_subbasis_left}.
        Thus $U\in T_2$ by \cref{topology_on,subseteq,powerset_elim,preimg_pi1,product_basis,subseteq_refl,times_subseteq_left,times_subseteq_right,subseteq_transitive,powerset_intro,product_basis_is_basis,basis_elem_open_in_generated,product_topology}.
    \end{subproof}
    Show that for all $U$ such that $U\in S_2$ we have $U\in T_2$.
    \begin{subproof}
        Fix $U$ such that $U\in S_2$.
        Take $V\in T_1$ such that $U = \preimg{\piTwo{X}{Y}}{V}$ by \cref{product_subbasis_right}.
        Hence $U\in T_2$ by \cref{topology_on,subseteq,powerset_elim,preimg_pi2,product_basis,subseteq_refl,times_subseteq_left,times_subseteq_right,subseteq_transitive,powerset_intro,product_basis_is_basis,basis_elem_open_in_generated,product_topology}.
    \end{subproof}
    For all $U$ such that $U\in S$ we have $U\in T_2$ by \cref{union_iff}.
    $\finiteIntersections{S}\subseteq T_2$ by \cref{finite_intersections,subseteq,finite_intersections_open_in_topology}.
    Show that $T_3\subseteq T_2$.
    \begin{subproof}
        Follows by \cref{basis_topology_unions,topology_generated_by_subbasis,subseteq,topology_on}.
    \end{subproof}
    Show that for all $B$ such that $B\in\productBasis{T}{T_1}{X}{Y}$ we have $B\in T_3$.
    \begin{subproof}
        Fix $B$ such that $B\in\productBasis{T}{T_1}{X}{Y}$.
        Take $U,V$ such that $U\in T$ and $V\in T_1$ and $B = U\times V$ by \cref{product_basis}.
        $U\subseteq X$ and $V\subseteq Y$ by \cref{topology_on,subseteq,powerset_elim}.
        Let $A_1 = \preimg{\piOne{X}{Y}}{U}$.
        Let $A_2 = \preimg{\piTwo{X}{Y}}{V}$.
        $A_1\in S$ and $A_2\in S$ by \cref{product_subbasis_left,product_subbasis_right,union_intro_left,union_intro_right}.
        Let $F = \{A_1, A_2\}$.
        $F\subseteq S$ and $F$ is finite and $F$ is inhabited
            by \cref{subseteq,upair_elim,pair_is_finite,upair_intro_left}.
        Therefore $B\in\finiteIntersections{S}$ and $B\in T_3$
            by \cref{preimg_pi1,preimg_pi2,inter_absorb_supseteq_right,inter_absorb_supseteq_left,inter_times,inter_as_inters,subbasis_on,times_subseteq_left,times_subseteq_right,subseteq_transitive,powerset_intro,finite_intersections,basis_elem_open_in_generated,topology_generated_by_subbasis}.
    \end{subproof}
    Show that for all $U$ such that $U\in T_2$ we have $U\in T_3$.
    \begin{subproof}
        Fix $U$ such that $U\in T_2$.
        There exists $M\subseteq\productBasis{T}{T_1}{X}{Y}$ such that $U = \unions{M}$
            by \cref{product_basis_is_basis,basis_for_topology,product_topology,basis_topology_unions}.
        For all $A$ such that $A\in M$ we have $A\in T_3$ by \cref{subseteq}.
        $M\subseteq T_3$ by \cref{subseteq}.
        Thus $U\in T_3$
            by \cref{basis_topology_is_topology,basis_for_topology,topology_generated_by_subbasis,topology_on}.
    \end{subproof}
    Show that $T_2\subseteq T_3$.
    \begin{subproof}
        Follows by \cref{subseteq}.
    \end{subproof}
    $T_2 = T_3$ by \cref{subseteq_antisymmetric}.
\end{proof}

\section{§16 The Subspace Topology}

% from §16 Item 1: subspace topology
\begin{definition}\label{subspace_topology}
    $\subspaceTopology{T}{Y} = \{A \mid \exists U\in T. A = Y\inter U\}$.
\end{definition}

% from §16 Item 2: subspace of a space
\begin{definition}\label{subspace_of}
    $Y$ is a subspace of $X$ with respect to $T$ iff $T$ is a topology on $X$ and $Y\subseteq X$.
\end{definition}

% from §16 helper lemma 1: subspace topology is a topology
\begin{lemma}\label{subspace_topology_is_topology}
    Assume $T$ is a topology on $X$.
    Assume $Y\subseteq X$.
    Then $\subspaceTopology{T}{Y}$ is a topology on $Y$.
\end{lemma}
\begin{proof}
    Let $T_1 = \subspaceTopology{T}{Y}$.
    $T_1\subseteq\pow{Y}$ by \cref{subspace_topology,inter_lower_left,powerset_intro,subseteq}.
    Show that $\emptyset\in T_1$ and $Y\in T_1$.
    \begin{subproof}
        Follows by \cref{topology_on,inter_emptyset,inter_absorb_supseteq_right,subspace_topology}.
    \end{subproof}
    Show that for all $M$ such that $M\subseteq T_1$ we have $\unions{M}\in T_1$.
    \begin{subproof}
        Fix $M$ such that $M\subseteq T_1$.
        Let $N = \{U\in T \mid \exists A\in M. A = Y\inter U\}$.
        $N\subseteq T$ by \cref{subseteq}.
        For all $x$ we have $x\in\unions{M}$ iff $x\in Y\inter\unions{N}$
            by \cref{unions_iff,subseteq,subspace_topology,inter,inter_intro}.
        Thus $\unions{M} = Y\inter\unions{N}$ and $\unions{M}\in T_1$
            by \cref{setext,topology_on,subspace_topology}.
    \end{subproof}
    Show that for all $A,B$ such that $A\in T_1$ and $B\in T_1$ we have $A\inter B\in T_1$.
    \begin{subproof}
        Fix $A$.
        Fix $B$ such that $A\in T_1$ and $B\in T_1$.
        Take $U,V\in T$ such that $A = Y\inter U$ and $B = Y\inter V$ by \cref{subspace_topology}.
        For all $x$ we have $x\in A\inter B$ iff $x\in Y\inter (U\inter V)$ by \cref{inter,inter_intro}.
        Thus $A\inter B = Y\inter (U\inter V)$ and $A\inter B\in T_1$
            by \cref{setext,topology_on,subspace_topology}.
    \end{subproof}
    $T_1$ is closed under arbitrary unions and closed under binary intersections.
    Therefore $T_1$ is a topology on $Y$ by \cref{topology_on}.
\end{proof}
% from §16 Item 3: basis for the subspace topology
\begin{lemma}\label{subspace_basis}
    Assume $B$ is a basis for $T$ on $X$.
    Assume $Y\subseteq X$.
    Let $C = \{A\in\pow{Y} \mid \exists B_1\in B. A = B_1\inter Y\}$.
    Then $C$ is a basis for $\subspaceTopology{T}{Y}$ on $Y$.
\end{lemma}
\begin{proof}
    Let $T_1 = \subspaceTopology{T}{Y}$.
    $B$ is a basis on $X$ and $T = \basisTopology{B}{X}$ and $T_1$ is a topology on $Y$
        by \cref{basis_for_topology,basis_topology_is_topology,subspace_topology_is_topology}.
    For all $C_1$ such that $C_1\in C$ we have $C_1$ is open in $Y$ with respect to $T_1$
        by \cref{subseteq,basis_elem_open_in_generated,basis_for_topology,inter_comm,subspace_topology,open_in}.
    Show that for all $U,x$ such that $U\in T_1$ and $x\in U$
        there exists $C_1\in C$ such that $x\in C_1$ and $C_1\subseteq U$.
    \begin{subproof}
        Fix $U$.
        Fix $x$ such that $U\in T_1$ and $x\in U$.
        Take $U_1\in T$ such that $U = Y\inter U_1$ by \cref{subspace_topology}.
        $x\in Y$ and $x\in U_1$ by \cref{inter}.
        Take $B_1\in B$ such that $x\in B_1$ and $B_1\subseteq U_1$
            by \cref{basis_for_topology,topology_generated_by_basis}.
        Let $C_1 = B_1\inter Y$.
        Thus $C_1\in C$.
        $x\in C_1$ and $C_1\subseteq U$
            by \cref{inter_intro,inter_lower_left,subseteq_transitive,inter_lower_right,subseteq_inter_iff}.
    \end{subproof}
    $C$ is a basis for $T_1$ on $Y$ by \cref{open_collection_basis}.
\end{proof}

% from §16 Item 4: open in a subspace implies open in the space if the subspace is open
\begin{lemma}\label{open_in_subspace_open_in_space}
    Assume $Y$ is a subspace of $X$ with respect to $T$.
    Assume $U$ is open in $Y$ with respect to $\subspaceTopology{T}{Y}$.
    Assume $Y$ is open in $X$ with respect to $T$.
    Then $U$ is open in $X$ with respect to $T$.
\end{lemma}
\begin{proof}
    Take $V\in T$ such that $U = Y\inter V$ by \cref{subspace_of,open_in,subspace_topology}.
    $U$ is open in $X$ with respect to $T$ by \cref{topology_on,open_in}.
\end{proof}

% from §16 Item 5: product topology and subspace topology
\begin{theorem}\label{subspace_product_topology}
    Assume $A$ is a subspace of $X$ with respect to $T$.
    Assume $B$ is a subspace of $Y$ with respect to $T_1$.
    Then $\productTopology{\subspaceTopology{T}{A}}{\subspaceTopology{T_1}{B}}{A}{B}
        = \subspaceTopology{\productTopology{T}{T_1}{X}{Y}}{A\times B}$.
\end{theorem}
\begin{proof}
    Let $T_A = \subspaceTopology{T}{A}$.
    Let $T_B = \subspaceTopology{T_1}{B}$.
    Let $T_2 = \productTopology{T}{T_1}{X}{Y}$.
    Let $T_3 = \subspaceTopology{T_2}{A\times B}$.
    Let $C = \productBasis{T_A}{T_B}{A}{B}$.
    $T$ is a topology on $X$ and $T_1$ is a topology on $Y$ and
        $A\subseteq X$ and $B\subseteq Y$ and
        $T_A$ is a topology on $A$ and $T_B$ is a topology on $B$
        by \cref{subspace_of,subspace_topology_is_topology}.
    $\productBasis{T}{T_1}{X}{Y}$ is a basis on $X\times Y$ and $T_2$ is a topology on $X\times Y$
        by \cref{product_basis_is_basis,basis_topology_is_topology,product_topology}.
    $T_3$ is a topology on $A\times B$
        by \cref{times_subseteq_left,times_subseteq_right,subseteq_transitive,subspace_topology_is_topology}.
    Show that for all $C_1$ such that $C_1\in C$ we have $C_1$ is open in $A\times B$ with respect to $T_3$.
    \begin{subproof}
        Fix $C_1$ such that $C_1\in C$.
        Take $U_A, V_B$ such that $U_A\in T_A$ and $V_B\in T_B$ and $C_1 = U_A\times V_B$ by \cref{product_basis}.
        Take $U\in T$ such that $U_A = A\inter U$ by \cref{subspace_topology}.
        Take $V\in T_1$ such that $V_B = B\inter V$ by \cref{subspace_topology}.
        $U\times V\in T_2$
            by \cref{topology_on,subseteq,powerset_elim,times_subseteq_left,times_subseteq_right,subseteq_transitive,powerset_intro,product_basis,basis_elem_open_in_generated,product_basis_is_basis,product_topology}.
        Thus $C_1 = (A\times B)\inter(U\times V)$ by \cref{inter_times}.
        $C_1\in T_3$ by \cref{subspace_topology}.
        $C_1$ is open in $A\times B$ with respect to $T_3$ by \cref{open_in}.
    \end{subproof}
    Show that for all $U,p$ such that $U\in T_3$ and $p\in U$
        there exists $C_1\in C$ such that $p\in C_1$ and $C_1\subseteq U$.
    \begin{subproof}
        Fix $U$.
        Fix $p$ such that $U\in T_3$ and $p\in U$.
        Take $W\in T_2$ such that $U = (A\times B)\inter W$ by \cref{subspace_topology}.
        Thus $p\in A\times B$ and $p\in W$ by \cref{inter}.
        Take $W_1\in\productBasis{T}{T_1}{X}{Y}$ such that $p\in W_1$ and $W_1\subseteq W$
            by \cref{product_basis_is_basis,basis_for_topology,product_topology,topology_generated_by_basis}.
        Take $U_1, V_1$ such that $U_1\in T$ and $V_1\in T_1$ and $W_1 = U_1\times V_1$ by \cref{product_basis}.
        Let $U_A = A\inter U_1$.
        Let $V_B = B\inter V_1$.
        $U_A\in T_A$ and $V_B\in T_B$ by \cref{subspace_topology}.
        Let $C_1 = U_A\times V_B$.
        Thus $C_1\in C$
            by \cref{topology_on,subseteq,powerset_elim,times_subseteq_left,times_subseteq_right,subseteq_transitive,powerset_intro,product_basis}.
        Take $a,b$ such that $p = (a,b)$ and $a\in A$ and $b\in B$ by \cref{times_elem_is_tuple}.
        Take $a_1,b_1$ such that $p = (a_1,b_1)$ and $a_1\in U_1$ and $b_1\in V_1$ by \cref{times_elem_is_tuple}.
        $p\in C_1$ and $C_1\subseteq U$
            by \cref{pair_eq_iff,inter_intro,times,inter_times,inter_lower_left,inter_lower_right,subseteq_transitive,subseteq_inter_iff}.
    \end{subproof}
    $C$ is a basis for $T_3$ on $A\times B$ and $T_3 = \basisTopology{C}{A\times B}$
        by \cref{open_collection_basis,basis_for_topology}.
    $\productTopology{T_A}{T_B}{A}{B} = T_3$ by \cref{product_topology}.
\end{proof}
% from §16 Item 6: convex subset of an ordered set
\begin{definition}\label{convex_in}
    $Y$ is convex in $X$ with respect to $R$ iff
    $Y\subseteq X$ and
    for all $a,b$ such that $a\in Y$ and $b\in Y$ and $a\mathrel{R}b$
        for all $x$ such that $x\in X$ and $a\mathrel{R}x$ and $x\mathrel{R}b$ we have $x\in Y$.
\end{definition}

% from §16 helper lemma 2: point left of one element of a convex set is left of all of it
\begin{lemma}\label{convex_outside_left}
    Assume $Y$ is convex in $X$ with respect to $R$.
    Assume $R$ is a simple order relation on $X$.
    Suppose $c\in X$ and $c\notin Y$.
    Suppose $y_0\in Y$ and $c\mathrel{R}y_0$.
    Then for all $y$ such that $y\in Y$ we have $c\mathrel{R}y$.
\end{lemma}
\begin{proof}
    Fix $y$ such that $y\in Y$.
    Suppose not.
    $y\in X$ and $y\neq c$ by \cref{convex_in,subseteq}.
    $c\in Y$ by \cref{convex_in,subseteq,simple_order_relation,connex,transitive}.
    Contradiction.
\end{proof}

% from §16 helper lemma 3: point right of one element of a convex set is right of all of it
\begin{lemma}\label{convex_outside_right}
    Assume $Y$ is convex in $X$ with respect to $R$.
    Assume $R$ is a simple order relation on $X$.
    Suppose $c\in X$ and $c\notin Y$.
    Suppose $y_0\in Y$ and $y_0\mathrel{R}c$.
    Then for all $y$ such that $y\in Y$ we have $y\mathrel{R}c$.
\end{lemma}
\begin{proof}
    Fix $y$ such that $y\in Y$.
    Suppose not.
    $y\in X$ and $y\neq c$ by \cref{convex_in,subseteq}.
    $c\in Y$ by \cref{convex_in,subseteq,simple_order_relation,connex,transitive}.
    Contradiction.
\end{proof}

% from §16 helper lemma 4: restricting a right ray to a subset
\begin{lemma}\label{openrayright_restrict}
    Suppose $Y\subseteq X$.
    Suppose $a\in Y$.
    Then $\openrayright{R}{Y}{a} = Y\inter\openrayright{R}{X}{a}$.
\end{lemma}
\begin{proof}
    Follows by \cref{openray_right,inter,subseteq,setext}.
\end{proof}

% from §16 helper lemma 5: restricting a left ray to a subset
\begin{lemma}\label{openrayleft_restrict}
    Suppose $Y\subseteq X$.
    Suppose $a\in Y$.
    Then $\openrayleft{R}{Y}{a} = Y\inter\openrayleft{R}{X}{a}$.
\end{lemma}
\begin{proof}
    Follows by \cref{openray_left,inter,subseteq,setext}.
\end{proof}

% from §16 Item 7: convex subspace order topology equals subspace topology
\begin{theorem}\label{convex_order_subspace_topology}
    Assume $T$ is the order topology on $X$ with respect to $R$.
    Assume $Y$ is convex in $X$ with respect to $R$.
    Assume $R$ is a simple order relation on $Y$.
    Assume $Y$ has more than one element.
    Then $\basisTopology{\orderBasis{R}{Y}}{Y} = \subspaceTopology{T}{Y}$.
\end{theorem}
\begin{proof}
    Let $T_Y = \basisTopology{\orderBasis{R}{Y}}{Y}$.
    Let $T_S = \subspaceTopology{T}{Y}$.
    Let $S_Y = \openRays{R}{Y}$.
    Let $S_X = \openRays{R}{X}$.
    $R$ is a simple order relation on $X$ and $X$ has more than one element and
        $T = \basisTopology{\orderBasis{R}{X}}{X}$ by \cref{order_topology}.
    $Y\subseteq X$ by \cref{convex_in}.
    $R$ is a simple order relation on $Y$ and $Y$ has more than one element.
    $\orderBasis{R}{X}$ is a basis on $X$ and $\orderBasis{R}{X}$ is a basis for $T$ on $X$
        by \cref{order_basis_is_basis,basis_for_topology,order_topology}.
    $T$ is a topology on $X$ and $T_S$ is a topology on $Y$
        by \cref{basis_topology_is_topology,subspace_topology_is_topology}.
    $\orderBasis{R}{Y}$ is a basis on $Y$ and $T_Y$ is a topology on $Y$
        and $T_Y$ is the order topology on $Y$ with respect to $R$
        by \cref{order_basis_is_basis,basis_topology_is_topology,order_topology}.
    Show that $\subbasisTopology{S_Y}{Y} = T_Y$ and $\subbasisTopology{S_X}{X} = T$.
    \begin{subproof}
        Follows by \cref{open_rays_subbasis}.
    \end{subproof}

    Show that for all $U$ such that $U\in S_X$ we have $U\in T$.
    \begin{subproof}
        Follows by \cref{open_rays_subbasis_on,subbasis_on,open_rays,powerset_elim,subseteq_transitive,powerset_intro,singleton_subset_intro,pair_is_finite,upair_intro_left,subseteq_finite_is_finite,singleton_intro,inters_singleton,finite_intersections,subbasis_finite_intersections_basis,basis_elem_open_in_generated,topology_generated_by_subbasis}.
    \end{subproof}

    Show that for all $U$ such that $U\in S_Y$ we have $U\in T_Y$.
    \begin{subproof}
        Follows by \cref{open_rays_subbasis_on,subbasis_on,open_rays,powerset_elim,subseteq_transitive,powerset_intro,singleton_subset_intro,pair_is_finite,upair_intro_left,subseteq_finite_is_finite,singleton_intro,inters_singleton,finite_intersections,subbasis_finite_intersections_basis,basis_elem_open_in_generated,topology_generated_by_subbasis}.
    \end{subproof}

    Show that for all $U$ such that $U\in S_Y$ we have $U\in T_S$.
    \begin{subproof}
        Fix $U$ such that $U\in S_Y$.
        Take $a\in Y$ such that $U = \openrayright{R}{Y}{a}$ or $U = \openrayleft{R}{Y}{a}$ by \cref{open_rays}.
        \begin{byCase}
        \caseOf{$U = \openrayright{R}{Y}{a}$.}
            Let $U_1 = \openrayright{R}{X}{a}$.
            $a\in X$ by \cref{subseteq}.
            Thus $U_1\in S_X$ and $U_1\in T$ by \cref{openray_right,subseteq,powerset_intro,open_rays}.
            $U = Y\inter U_1$ by \cref{openrayright_restrict}.
            $U\in T_S$ by \cref{subspace_topology}.
        \caseOf{$U = \openrayleft{R}{Y}{a}$.}
            Let $U_1 = \openrayleft{R}{X}{a}$.
            $a\in X$ by \cref{subseteq}.
            Hence $U_1\in S_X$ and $U_1\in T$ by \cref{openray_left,subseteq,powerset_intro,open_rays}.
            $U = Y\inter U_1$ by \cref{openrayleft_restrict}.
            $U\in T_S$ by \cref{subspace_topology}.
        \end{byCase}
    \end{subproof}

    For all $U$ such that $U\in T_Y$ there exists $M\subseteq\finiteIntersections{S_Y}$ such that $U = \unions{M}$
        by \cref{open_rays_subbasis_on,subbasis_finite_intersections_basis,basis_for_topology,topology_generated_by_subbasis,basis_topology_unions}.
    For all $B$ such that $B\in\finiteIntersections{S_Y}$ we have $B\in T_S$
        by \cref{finite_intersections,subseteq,finite_intersections_open_in_topology}.
    Show that for all $U$ such that $U\in T_Y$ we have $U\in T_S$.
    \begin{subproof}
        Follows by \cref{subseteq,topology_on}.
    \end{subproof}
    $T_Y\subseteq T_S$ by \cref{subseteq}.

    Show that for all $c\in X$ we have $Y\inter\openrayright{R}{X}{c}\in T_Y$.
    \begin{subproof}
        Fix $c$ such that $c\in X$.
        \begin{byCase}
        \caseOf{$c\in Y$.}
            $Y\inter\openrayright{R}{X}{c} = \openrayright{R}{Y}{c}$
                by \cref{inter,openray_right,subseteq,inter_intro,subseteq_antisymmetric}.
            $Y\inter\openrayright{R}{X}{c}\in T_Y$
                by \cref{openray_right,subseteq,powerset_intro,open_rays}.
            \caseOf{$c\notin Y$.}
                Take $y_0, y_1\in Y$ such that $y_0\neq y_1$ by \cref{more_than_one_element}.
                $y_0\neq c$ and $(c\mathrel{R}y_0 \lor y_0\mathrel{R}c)$
                    by \cref{subseteq,simple_order_relation,connex}.
            \begin{byCase}
            \caseOf{$c\mathrel{R}y_0$.}
                For all $y$ such that $y\in Y$ we have $c\mathrel{R}y$
                    by \cref{convex_outside_left}.
                Thus $Y\subseteq \openrayright{R}{X}{c}$ by \cref{subseteq,elem_subseteq,openray_right}.
                $Y\inter\openrayright{R}{X}{c}\in T_Y$ by \cref{inter_absorb_supseteq_right,topology_on}.
            \caseOf{$y_0\mathrel{R}c$.}
                For all $y$ such that $y\in Y$ we have $y\mathrel{R}c$
                    by \cref{convex_outside_right}.
                Show that $Y\inter\openrayright{R}{X}{c} = \emptyset$.
                \begin{subproof}
                    Suppose not.
                    Take $x$ such that $x\in Y\inter\openrayright{R}{X}{c}$ by \cref{nonempty_inhabited}.
                    $c\mathrel{R}x$ and $x\mathrel{R}c$ by \cref{inter,openray_right}.
                    Contradiction by \cref{simple_order_relation,asymmetric}.
                \end{subproof}
                Therefore $Y\inter\openrayright{R}{X}{c}\in T_Y$ by \cref{topology_on}.
            \end{byCase}
        \end{byCase}
    \end{subproof}

    Show that for all $c\in X$ we have $Y\inter\openrayleft{R}{X}{c}\in T_Y$.
    \begin{subproof}
        Fix $c$ such that $c\in X$.
        \begin{byCase}
        \caseOf{$c\in Y$.}
            $Y\inter\openrayleft{R}{X}{c} = \openrayleft{R}{Y}{c}$
                by \cref{inter,openray_left,subseteq,inter_intro,subseteq_antisymmetric}.
            $Y\inter\openrayleft{R}{X}{c}\in T_Y$
                by \cref{openray_left,subseteq,powerset_intro,open_rays}.
        \caseOf{$c\notin Y$.}
            Take $y_0, y_1\in Y$ such that $y_0\neq y_1$ by \cref{more_than_one_element}.
            $y_0\neq c$ and $(c\mathrel{R}y_0 \lor y_0\mathrel{R}c)$
                by \cref{subseteq,simple_order_relation,connex}.
            \begin{byCase}
            \caseOf{$y_0\mathrel{R}c$.}
                For all $y$ such that $y\in Y$ we have $y\mathrel{R}c$
                    by \cref{convex_outside_right}.
                Hence $Y\subseteq \openrayleft{R}{X}{c}$ by \cref{subseteq,elem_subseteq,openray_left}.
                $Y\inter\openrayleft{R}{X}{c}\in T_Y$ by \cref{inter_absorb_supseteq_right,topology_on}.
            \caseOf{$c\mathrel{R}y_0$.}
                For all $y$ such that $y\in Y$ we have $c\mathrel{R}y$
                    by \cref{convex_outside_left}.
                Show that $Y\inter\openrayleft{R}{X}{c} = \emptyset$.
                \begin{subproof}
                    Suppose not.
                    Take $x$ such that $x\in Y\inter\openrayleft{R}{X}{c}$ by \cref{nonempty_inhabited}.
                    $x\mathrel{R}c$ and $c\mathrel{R}x$ by \cref{inter,openray_left}.
                    Contradiction by \cref{simple_order_relation,asymmetric}.
                \end{subproof}
                Therefore $Y\inter\openrayleft{R}{X}{c}\in T_Y$ by \cref{topology_on}.
            \end{byCase}
        \end{byCase}
    \end{subproof}

    Show that for all $B$ such that $B\in\orderBasis{R}{X}$ we have $Y\inter B\in T_Y$.
    \begin{subproof}
        Fix $B$ such that $B\in\orderBasis{R}{X}$.
        $((\exists a, b\in X. a\mathrel{R}b \land B = \intervalopenrel{R}{X}{a}{b}) \lor
            (\exists b\in X. \exists a_0\in X. (\forall x\in X. a_0\mathrel{R}x \lor a_0 = x) \land B = \intervalclosedopenrel{R}{X}{a_0}{b}) \lor
            (\exists a\in X. \exists b_0\in X. (\forall x\in X. x\mathrel{R}b_0 \lor x = b_0) \land B = \intervalopenclosedrel{R}{X}{a}{b_0}))$
            by \cref{order_basis}.
        \begin{byCase}
        \caseOf{$\exists a, b\in X. a\mathrel{R}b \land B = \intervalopenrel{R}{X}{a}{b}$.}
            Take $a, b\in X$ such that $a\mathrel{R}b$ and $B = \intervalopenrel{R}{X}{a}{b}$.
            $B = \openrayright{R}{X}{a} \inter \openrayleft{R}{X}{b}$ by \cref{intervalopenrel_as_intersection_openrays}.
            $Y\inter\openrayright{R}{X}{a}\in T_Y$ and $Y\inter\openrayleft{R}{X}{b}\in T_Y$
                and $T_Y$ is closed under binary intersections by \cref{topology_on}.
            Show that $Y\inter B = (Y\inter\openrayright{R}{X}{a}) \inter (Y\inter\openrayleft{R}{X}{b})$.
            \begin{subproof}
                Show that for all $x$ such that $x\in Y\inter B$ we have
                    $x\in (Y\inter\openrayright{R}{X}{a}) \inter (Y\inter\openrayleft{R}{X}{b})$.
                \begin{subproof}
                    Fix $x$ such that $x\in Y\inter B$.
                    $x\in Y$ and $x\in B$ by \cref{inter}.
                    $x\in \openrayright{R}{X}{a} \inter \openrayleft{R}{X}{b}$.
                    $x\in\openrayright{R}{X}{a}$ and $x\in\openrayleft{R}{X}{b}$ by \cref{inter}.
                    $x\in Y\inter\openrayright{R}{X}{a}$ and $x\in Y\inter\openrayleft{R}{X}{b}$ by \cref{inter_intro}.
                    $x\in (Y\inter\openrayright{R}{X}{a}) \inter (Y\inter\openrayleft{R}{X}{b})$ by \cref{inter_intro}.
                \end{subproof}
                $Y\inter B\subseteq (Y\inter\openrayright{R}{X}{a}) \inter (Y\inter\openrayleft{R}{X}{b})$ by \cref{subseteq}.
                Show that for all $x$ such that
                    $x\in (Y\inter\openrayright{R}{X}{a}) \inter (Y\inter\openrayleft{R}{X}{b})$
                    we have $x\in Y\inter B$.
                \begin{subproof}
                    Fix $x$ such that $x\in (Y\inter\openrayright{R}{X}{a}) \inter (Y\inter\openrayleft{R}{X}{b})$.
                    $x\in Y\inter\openrayright{R}{X}{a}$ and $x\in Y\inter\openrayleft{R}{X}{b}$ by \cref{inter}.
                    $x\in Y$ and $x\in\openrayright{R}{X}{a}$ and $x\in\openrayleft{R}{X}{b}$ by \cref{inter}.
                    $x\in \openrayright{R}{X}{a} \inter \openrayleft{R}{X}{b}$ by \cref{inter_intro}.
                    $x\in B$.
                    $x\in Y\inter B$ by \cref{inter_intro}.
                \end{subproof}
                $(Y\inter\openrayright{R}{X}{a}) \inter (Y\inter\openrayleft{R}{X}{b})\subseteq Y\inter B$ by \cref{subseteq}.
                $Y\inter B = (Y\inter\openrayright{R}{X}{a}) \inter (Y\inter\openrayleft{R}{X}{b})$ by \cref{subseteq_antisymmetric}.
            \end{subproof}
            Hence $Y\inter B\in T_Y$ by \cref{subseteq_antisymmetric,topology_on}.
        \caseOf{$\exists b\in X. \exists a_0\in X. (\forall x\in X. a_0\mathrel{R}x \lor a_0 = x) \land B = \intervalclosedopenrel{R}{X}{a_0}{b}$.}
            Take $a_0, b\in X$ such that
                $(\forall x\in X. a_0\mathrel{R}x \lor a_0 = x)$ and $B = \intervalclosedopenrel{R}{X}{a_0}{b}$.
            $B = \openrayleft{R}{X}{b}$ and $Y\inter B = Y\inter\openrayleft{R}{X}{b}$
                by \cref{smallest_element,intervalclosedopenrel_smallest_eq_openrayleft}.
            $Y\inter B\in T_Y$.
        \caseOf{$\exists a\in X. \exists b_0\in X. (\forall x\in X. x\mathrel{R}b_0 \lor x = b_0) \land B = \intervalopenclosedrel{R}{X}{a}{b_0}$.}
            Take $a, b_0\in X$ such that
                $(\forall x\in X. x\mathrel{R}b_0 \lor x = b_0)$ and $B = \intervalopenclosedrel{R}{X}{a}{b_0}$.
            $B = \openrayright{R}{X}{a}$ and $Y\inter B = Y\inter\openrayright{R}{X}{a}$
                by \cref{largest_element,intervalopenclosedrel_largest_eq_openrayright}.
            $Y\inter B\in T_Y$.
        \end{byCase}
    \end{subproof}

    Show that for all $U$ such that $U\in T_S$ we have $U\in T_Y$.
    \begin{subproof}
            Fix $U$ such that $U\in T_S$.
            Take $V\in T$ such that $U = Y\inter V$ by \cref{subspace_topology}.
            There exists $M\subseteq\orderBasis{R}{X}$ such that $V = \unions{M}$ by \cref{basis_topology_unions}.
            Let $N = \{A\in\pow{Y} \mid \exists B\in M. A = Y\inter B\}$.
            Show that for all $x$ such that $x\in U$ we have $x\in\unions{N}$.
            \begin{subproof}
                Fix $x$ such that $x\in U$.
                $x\in Y$ and $x\in V$ by \cref{inter}.
                Take $B\in M$ such that $x\in B$ by \cref{unions_iff}.
                $x\in Y\inter B$ by \cref{inter_intro}.
                $Y\inter B\in N$.
                $x\in\unions{N}$ by \cref{unions_iff}.
            \end{subproof}
            $U\subseteq\unions{N}$ by \cref{subseteq}.
            Show that for all $x$ such that $x\in\unions{N}$ we have $x\in U$.
            \begin{subproof}
                Follows by \cref{unions_iff,inter,inter_intro}.
            \end{subproof}
            $\unions{N}\subseteq U$ by \cref{subseteq}.
            For all $A$ such that $A\in N$ we have $A\in T_Y$ by \cref{subseteq}.
            Therefore $U\in T_Y$ by \cref{subseteq,topology_on,subseteq_antisymmetric}.
    \end{subproof}
    $T_S\subseteq T_Y$ by \cref{subseteq}.

    $T_Y = T_S$ by \cref{subseteq_antisymmetric}.
\end{proof}
\section{§17 Closed Sets and Limit Points}

% from §17 Item 1: closed set
\begin{definition}\label{closed_in}
    $A$ is closed in $X$ with respect to $T$ iff
    $A\subseteq X$ and $X\setminus A$ is open in $X$ with respect to $T$.
\end{definition}

% from §17 Item 2: empty set is closed
\begin{lemma}\label{closed_emptyset}
    Assume $T$ is a topology on $X$.
    Then $\emptyset$ is closed in $X$ with respect to $T$.
\end{lemma}
\begin{proof}
    Follows by \cref{closed_in,emptyset_subseteq,setminus_emptyset,topology_on,open_in}.
\end{proof}

% from §17 Item 3: the whole space is closed
\begin{lemma}\label{closed_whole_space}
    Assume $T$ is a topology on $X$.
    Then $X$ is closed in $X$ with respect to $T$.
\end{lemma}
\begin{proof}
    Follows by \cref{closed_in,subseteq_refl,setminus_self,topology_on,open_in}.
\end{proof}

% from §17 Item 4: intersections of closed sets are closed
\begin{lemma}\label{closed_inters}
    Assume $T$ is a topology on $X$.
    Suppose $C\subseteq\pow{X}$.
    Suppose for all $A$ such that $A\in C$ we have $A$ is closed in $X$ with respect to $T$.
    Then $\inters{C}$ is closed in $X$ with respect to $T$.
\end{lemma}
\begin{proof}
    $\inters{C}\subseteq X$ by \cref{unions_subseteq_of_powerset_is_subseteq,inters,elem_subseteq,subseteq}.
    \begin{byCase}
    \caseOf{$C = \emptyset$.}
        We have $\inters{C} = \emptyset$ by \cref{inters,unions_emptyset,emptyset,setext}.
        $\inters{C}$ is closed in $X$ with respect to $T$ by \cref{closed_emptyset}.
    \caseOf{$C \neq \emptyset$.}
        $C$ is inhabited by \cref{nonempty_inhabited}.
        Let $D = \{U\in T \mid \exists A\in C. U = X\setminus A\}$.
        Show that $X\setminus\inters{C} = \unions{D}$.
        \begin{subproof}
            Show that for all $x$ such that $x\in X\setminus\inters{C}$ we have $x\in\unions{D}$.
            \begin{subproof}
                Fix $x$ such that $x\in X\setminus\inters{C}$.
                Suppose not.
                $x\notin\unions{D}$.
                Show that for all $A$ such that $A\in C$ we have $x\in A$.
                \begin{subproof}
                    Fix $A$ such that $A\in C$.
                    Suppose not.
                    Thus $x\in\unions{D}$ by \cref{setminus_elim_left,setminus_intro,closed_in,open_in,unions_iff}.
                    Contradiction.
                \end{subproof}
                Therefore $x\in\inters{C}$ by \cref{inters_intro}.
                Contradiction by \cref{setminus_elim_right}.
            \end{subproof}
            We have $X\setminus\inters{C}\subseteq\unions{D}$ by \cref{subseteq}.
            Show that for all $x$ such that $x\in\unions{D}$ we have $x\in X\setminus\inters{C}$.
            \begin{subproof}
                Fix $x$ such that $x\in\unions{D}$.
                Take $U\in D$ such that $x\in U$ by \cref{unions_iff}.
                Take $A\in C$ such that $U = X\setminus A$.
                We have $x\in X$ and $x\notin A$ by \cref{setminus_elim_left,setminus_elim_right}.
                $x\notin\inters{C}$.
                Thus $x\in X\setminus\inters{C}$ by \cref{setminus_intro}.
            \end{subproof}
            We have $\unions{D}\subseteq X\setminus\inters{C}$ by \cref{subseteq}.
            Thus $X\setminus\inters{C} = \unions{D}$ by \cref{subseteq_antisymmetric}.
        \end{subproof}
        Therefore $X\setminus\inters{C}$ is open in $X$ with respect to $T$ by \cref{subseteq,topology_on,open_in}.
        Hence $\inters{C}$ is closed in $X$ with respect to $T$ by \cref{closed_in}.
    \end{byCase}
\end{proof}

% from §17 helper lemma 1: complements of subsets are injective
\begin{lemma}\label{complement_eq_implies_eq}
    Suppose $A\subseteq X$ and $B\subseteq X$.
    Suppose $X\setminus A = X\setminus B$.
    Then $A = B$.
\end{lemma}
\begin{proof}
    Follows by \cref{subseteq_refl,double_complement}.
\end{proof}

% from §17 Item 5: finite unions of closed sets are closed
\begin{lemma}\label{closed_finite_unions}
    Assume $T$ is a topology on $X$.
    Suppose $F\subseteq\pow{X}$.
    Suppose $F$ is finite.
    Suppose for all $A$ such that $A\in F$ we have $A$ is closed in $X$ with respect to $T$.
    Then $\unions{F}$ is closed in $X$ with respect to $T$.
\end{lemma}
\begin{proof}
    \begin{byCase}
    \caseOf{$F = \emptyset$.}
        We have $\unions{F}$ is closed in $X$ with respect to $T$ by \cref{unions_emptyset,closed_emptyset}.
    \caseOf{$F \neq \emptyset$.}
        Let $G = \{U\in T \mid \exists A\in F. U = X\setminus A\}$.
        $G$ is inhabited by \cref{nonempty_inhabited,closed_in,open_in}.
        Show that for all $x$ such that $x\in X\setminus\unions{F}$ we have $x\in\inters{G}$.
        \begin{subproof}
            Fix $x$ such that $x\in X\setminus\unions{F}$.
            $x\in X$ and $x\notin\unions{F}$ by \cref{setminus_elim_left,setminus_elim_right}.
            Show that for all $U$ such that $U\in G$ we have $x\in U$.
            \begin{subproof}
                Fix $U$ such that $U\in G$.
                Take $A\in F$ such that $U = X\setminus A$.
                Suppose not.
                We have $x\notin U$.
                Hence $x\notin X\setminus A$.
                Show that $x\in A$.
                \begin{subproof}
                    Suppose not.
                    We have $x\notin A$.
                    Hence $x\in X\setminus A$ by \cref{setminus_intro}.
                    Contradiction.
                \end{subproof}
                Therefore $x\in\unions{F}$ by \cref{unions_iff}.
                Contradiction.
            \end{subproof}
            Therefore $x\in\inters{G}$ by \cref{inters_intro}.
        \end{subproof}
        We have $X\setminus\unions{F}\subseteq\inters{G}$ by \cref{subseteq}.
        Show that for all $x$ such that $x\in\inters{G}$ we have $x\in X\setminus\unions{F}$.
        \begin{subproof}
            Fix $x$ such that $x\in\inters{G}$.
            Take $A$ such that $A\in F$.
            We have $x\in X\setminus A$ by \cref{closed_in,open_in,inters_destr}.
            Show that $x\notin\unions{F}$.
            \begin{subproof}
                Suppose not.
                We have $x\in\unions{F}$.
                Take $A_1\in F$ such that $x\in A_1$ by \cref{unions_iff}.
                Therefore $x\notin A_1$ by \cref{closed_in,open_in,inters_destr,setminus_elim_right}.
                Contradiction.
            \end{subproof}
            Hence $x\in X\setminus\unions{F}$ by \cref{inters_destr,setminus_elim_left,setminus_intro}.
        \end{subproof}
        We have $\inters{G}\subseteq X\setminus\unions{F}$ by \cref{subseteq}.
        Therefore $X\setminus\unions{F} = \inters{G}$ by \cref{subseteq_antisymmetric}.
        $G\subseteq T$ by \cref{subseteq}.
        Take $k\in\naturalsPlus$ such that there exists a bijection from $\initialSegment{k}$ to $F$ by \cref{finite}.
        Take $f$ such that $f$ is a bijection from $\initialSegment{k}$ to $F$.
        Let $g = \{(n, X\setminus f(n)) \mid n\in\initialSegment{k}\}$.
        Show that $g$ is a function to $G$ and $\dom{g} = \initialSegment{k}$.
        \begin{subproof}
            For all $w$ such that $w\in g$ there exist $x,y$ such that $w = (x,y)$.
            For all $n,y,y'$ such that $(n,y)\in g$ and $(n,y')\in g$ we have $y = y'$ by \cref{pair_eq_iff}.
            Hence $g$ is a function by \cref{relation,rightunique}.
            For all $n$ such that $n\in\initialSegment{k}$ we have $n\in\dom{g}$ by \cref{dom_intro}.
            We have $\initialSegment{k}\subseteq\dom{g}$ by \cref{subseteq}.
            For all $n$ such that $n\in\dom{g}$ we have $n\in\initialSegment{k}$ by \cref{dom_iff,pair_eq_iff}.
            $\dom{g}\subseteq\initialSegment{k}$ by \cref{subseteq}.
            Therefore $\dom{g} = \initialSegment{k}$ by \cref{subseteq_antisymmetric}.
            For all $n$ such that $n\in\dom{g}$ we have $f(n)\in F$ and $X\setminus f(n)\in G$ and
                $(n, X\setminus f(n))\in g$ by \cref{bijections,bijections_to_funs,funs_elim,closed_in,open_in}.
            For all $n$ such that $n\in\dom{g}$ we have $g(n)\in G$ by \cref{function_apply_iff}.
            Thus $g$ is a function to $G$.
        \end{subproof}
        We have $g\in\funs{\initialSegment{k}}{G}$ by \cref{funs_intro}.
        $g\in\surj{\initialSegment{k}}{G}$ by \cref{bijections,surj,function_apply_iff}.
        For all $n,m$ such that $n\in\initialSegment{k}$ and $m\in\initialSegment{k}$ and $g(n) = g(m)$ we have $n = m$
            by \cref{function_apply_iff,bijections,bijections_to_funs,funs_elim,subseteq,pow_iff,complement_eq_implies_eq,injective_function}.
        Hence $g$ is injective by \cref{injective_function}.
        Therefore $g$ is a bijection from $\initialSegment{k}$ to $G$ by \cref{bijections}.
        Hence $G$ is finite by \cref{finite}.
        Therefore $X\setminus\unions{F}$ is open in $X$ with respect to $T$
            by \cref{subseteq,finite_intersections_open_in_topology,open_in}.
        $\unions{F}$ is closed in $X$ with respect to $T$
            by \cref{unions_subseteq_of_powerset_is_subseteq,closed_in}.
    \end{byCase}
\end{proof}

% from §17 Item 6: closed sets in subspaces
\begin{theorem}\label{closed_in_subspace_iff}
    Assume $Y$ is a subspace of $X$ with respect to $T$.
    Then for all $A$ we have
    $A$ is closed in $Y$ with respect to $\subspaceTopology{T}{Y}$
    iff there exists $C$ such that $C$ is closed in $X$ with respect to $T$ and $A = C\inter Y$.
\end{theorem}
\begin{proof}
    $T$ is a topology on $X$ and $Y\subseteq X$ by \cref{subspace_of}.
    Let $T_Y = \subspaceTopology{T}{Y}$.
    Fix $A$.
    Show that if $A$ is closed in $Y$ with respect to $T_Y$ then
        there exists $C$ such that $C$ is closed in $X$ with respect to $T$ and $A = C\inter Y$.
    \begin{subproof}
        Assume $A$ is closed in $Y$ with respect to $T_Y$.
        $Y\setminus A\in T_Y$ by \cref{closed_in,open_in}.
        Take $U\in T$ such that $Y\setminus A = Y\inter U$ by \cref{subspace_topology}.
        Let $C = X\setminus U$.
        $U\subseteq X$ by \cref{open_in,topology_on,subseteq,pow_iff}.
        Therefore $C$ is closed in $X$ with respect to $T$ by
            \cref{setminus_subseteq,setminus_setminus,inter_comm,topology_on,open_in,subseteq,pow_iff,inter_absorb_supseteq_right,closed_in}.
        For all $x$ such that $x\in A$ we have $x\in Y$ by \cref{closed_in,elem_subseteq}.
        For all $x$ such that $x\in A$ we have $x\notin U$.
        We have $A\subseteq Y\setminus U$ by \cref{setminus_intro,subseteq}.
        Show that for all $x$ such that $x\in Y\setminus U$ we have $x\in A$.
        \begin{subproof}
            Fix $x$ such that $x\in Y\setminus U$.
            Suppose not.
            We have $x\notin A$.
            $x\in Y$ by \cref{setminus_elim_left}.
            Hence $x\in Y\setminus A$ by \cref{setminus_intro}.
            Thus $x\in Y\inter U$.
            Therefore $x\in U$ by \cref{inter_elim_right}.
            Contradiction by \cref{setminus_elim_right}.
        \end{subproof}
        We have $Y\setminus U\subseteq A$ by \cref{subseteq}.
        Therefore $A = Y\inter (X\setminus U)$ by \cref{subseteq_antisymmetric,setminus_eq_inter_complement}.
        Hence $A = C\inter Y$ by \cref{inter_comm}.
    \end{subproof}
    Show that if there exists $C$ such that $C$ is closed in $X$ with respect to $T$ and $A = C\inter Y$ then
        $A$ is closed in $Y$ with respect to $T_Y$.
    \begin{subproof}
        Assume there exists $C$ such that $C$ is closed in $X$ with respect to $T$ and $A = C\inter Y$.
        Take $C$ such that $C$ is closed in $X$ with respect to $T$ and $A = C\inter Y$.
        Show that for all $x$ such that $x\in Y\setminus A$ we have $x\in Y\setminus C$.
        \begin{subproof}
            Fix $x$ such that $x\in Y\setminus A$.
            $x\in Y$ and $x\notin A$ by \cref{setminus_elim_left,setminus_elim_right}.
            Suppose not.
            We have $x\in C$.
            Hence $x\in A$ by \cref{inter_intro}.
            Contradiction.
            $x\notin C$.
            Therefore $x\in Y\setminus C$ by \cref{setminus_intro}.
        \end{subproof}
        We have $Y\setminus A\subseteq Y\setminus C$ by \cref{subseteq}.
        Show that for all $x$ such that $x\in Y\setminus C$ we have $x\in Y\setminus A$.
        \begin{subproof}
            Fix $x$ such that $x\in Y\setminus C$.
            $x\in Y$ and $x\notin C$ by \cref{setminus_elim_left,setminus_elim_right}.
            Suppose not.
            We have $x\in A$.
            Hence $x\in C$ by \cref{inter_elim_left}.
            Contradiction.
            $x\notin A$.
            Therefore $x\in Y\setminus A$ by \cref{setminus_intro}.
        \end{subproof}
        We have $Y\setminus C\subseteq Y\setminus A$ by \cref{subseteq}.
        Hence $Y\setminus A = Y\setminus C$ by \cref{subseteq_antisymmetric}.
        Therefore $A$ is closed in $Y$ with respect to $T_Y$ by \cref{subseteq_antisymmetric,setminus_eq_inter_complement,closed_in,open_in,subspace_topology,inter_lower_right,subspace_topology_is_topology}.
    \end{subproof}
\end{proof}

% from §17 Item 7: closed in a closed subspace
\begin{theorem}\label{closed_in_closed_subspace}
    Assume $T$ is a topology on $X$.
    Suppose $Y\subseteq X$.
    Suppose $Y$ is closed in $X$ with respect to $T$.
    Suppose $A$ is closed in $Y$ with respect to $\subspaceTopology{T}{Y}$.
    Then $A$ is closed in $X$ with respect to $T$.
\end{theorem}
\begin{proof}
    Take $C$ such that $C$ is closed in $X$ with respect to $T$ and $A = C\inter Y$
        by \cref{closed_in_subspace_iff,subspace_of}.
    Let $D = \{C, Y\}$.
    We have $D\subseteq\pow{X}$ and for all $B$ such that $B\in D$ we have $B$ is closed in $X$ with respect to $T$
        by \cref{upair_iff,closed_in,pow_iff,subseteq}.
    Therefore $A$ is closed in $X$ with respect to $T$ by \cref{closed_inters,inter_as_inters}.
\end{proof}

% from §17 Item 8: interior of a set
\begin{definition}\label{interior}
    $\interior{A}{X}{T} = \{x\in X \mid \exists U\in T. x\in U \land U\subseteq A\}$.
\end{definition}

% from §17 Item 9: closure of a set
\begin{definition}\label{closure}
    $\closure{A}{X}{T} = \{x\in X \mid \text{for all $C$ such that $C$ is closed in $X$ with respect to $T$ and for all $y\in A$ we have $y\in C$ we have $x\in C$}\}$.
\end{definition}

% from §17 Item 10: interior is open
\begin{lemma}\label{interior_open_in}
    Assume $T$ is a topology on $X$.
    Then $\interior{A}{X}{T}$ is open in $X$ with respect to $T$.
\end{lemma}
\begin{proof}
    Let $S = \{U\in T \mid U\subseteq A\}$.
    $\interior{A}{X}{T}\subseteq\unions{S}$ by \cref{interior,unions_iff,subseteq}.
    $\unions{S}\subseteq\interior{A}{X}{T}$
        by \cref{unions_iff,topology_on,subseteq,pow_iff,elem_subseteq,interior,subseteq}.
    Therefore $\interior{A}{X}{T}$ is open in $X$ with respect to $T$ by \cref{subseteq,topology_on,subseteq_antisymmetric,open_in}.
\end{proof}

% from §17 Item 11: closure is closed
\begin{lemma}\label{closure_closed_in}
    Assume $T$ is a topology on $X$.
    Suppose $A\subseteq X$.
    Then $\closure{A}{X}{T}$ is closed in $X$ with respect to $T$.
\end{lemma}
\begin{proof}
    Let $F = \{C\in\pow{X} \mid \text{$C$ is closed in $X$ with respect to $T$ and for all $x\in A$ we have $x\in C$}\}$.
    $F$ is inhabited by \cref{closed_whole_space,powerset_top,subseteq}.
    Show that $\closure{A}{X}{T}\subseteq\inters{F}$.
    \begin{subproof}
        Follows by \cref{closure,inters_iff_forall,subseteq}.
    \end{subproof}
    Show that for all $x$ such that $x\in\inters{F}$ we have $x\in\closure{A}{X}{T}$.
    \begin{subproof}
        Fix $x$ such that $x\in\inters{F}$.
        For all $C$ such that $C\in F$ we have $x\in C$ by \cref{inters_iff_forall}.
        We have $x\in X$ by \cref{inters_iff_forall,pow_iff,elem_subseteq}.
        For all $C$ such that $C$ is closed in $X$ with respect to $T$ and
            for all $y\in A$ we have $y\in C$ we have $x\in C$.
        Therefore $x\in\closure{A}{X}{T}$ by \cref{closure}.
    \end{subproof}
    We have $\inters{F}\subseteq\closure{A}{X}{T}$ by \cref{subseteq}.
    Therefore $\closure{A}{X}{T} = \inters{F}$ by \cref{subseteq_antisymmetric}.
    $F\subseteq\pow{X}$ by \cref{subseteq}.
    For all $C$ such that $C\in F$ we have $C$ is closed in $X$ with respect to $T$.
    Therefore $\inters{F}$ is closed in $X$ with respect to $T$ by \cref{closed_inters}.
    Hence $\closure{A}{X}{T}$ is closed in $X$ with respect to $T$.
\end{proof}

% from §17 Item 12: interior subset
\begin{lemma}\label{interior_subseteq}
    Assume $T$ is a topology on $X$.
    Then $\interior{A}{X}{T}\subseteq A$.
\end{lemma}
\begin{proof}
    Follows by \cref{interior,elem_subseteq,subseteq}.
\end{proof}

% from §17 Item 13: subset of closure
\begin{lemma}\label{subseteq_closure}
    Assume $T$ is a topology on $X$.
    Suppose $A\subseteq X$.
    Then $A\subseteq\closure{A}{X}{T}$.
\end{lemma}
\begin{proof}
    Follows by \cref{elem_subseteq,closure,subseteq}.
\end{proof}

% from §17 helper lemma 2: closure contained in closed superset
\begin{lemma}\label{closure_subseteq_closed}
    Assume $T$ is a topology on $X$.
    Suppose $A\subseteq X$.
    Suppose $C$ is closed in $X$ with respect to $T$.
    Suppose $A\subseteq C$.
    Then $\closure{A}{X}{T}\subseteq C$.
\end{lemma}
\begin{proof}
    Follows by \cref{closure,subseteq}.
\end{proof}

% from §17 Item 14: open set equals its interior
\begin{lemma}\label{open_eq_interior}
    Assume $T$ is a topology on $X$.
    Suppose $A$ is open in $X$ with respect to $T$.
    Then $A = \interior{A}{X}{T}$.
\end{lemma}
\begin{proof}
    Show that for all $x$ such that $x\in A$ we have $x\in\interior{A}{X}{T}$.
    \begin{subproof}
        Fix $x$ such that $x\in A$.
        $A\in T$ by \cref{open_in}.
        We have $x\in X$ by \cref{open_in,topology_on,subseteq,pow_iff,elem_subseteq}.
        Therefore $x\in\interior{A}{X}{T}$ by \cref{interior,subseteq_refl}.
    \end{subproof}
    We have $A\subseteq\interior{A}{X}{T}$ by \cref{subseteq}.
    $\interior{A}{X}{T}\subseteq A$ by \cref{interior_subseteq}.
    Therefore $A = \interior{A}{X}{T}$ by \cref{subseteq_antisymmetric}.
\end{proof}

% from §17 Item 15: closed set equals its closure
\begin{lemma}\label{closed_eq_closure}
    Assume $T$ is a topology on $X$.
    Suppose $A$ is closed in $X$ with respect to $T$.
    Then $A = \closure{A}{X}{T}$.
\end{lemma}
\begin{proof}
    Follows by \cref{closed_in,subseteq_closure,subseteq_refl,closure_subseteq_closed,subseteq_antisymmetric}.
\end{proof}

% from §17 Item 16: closure in subspace
\begin{theorem}\label{closure_in_subspace}
    Assume $Y$ is a subspace of $X$ with respect to $T$.
    Suppose $A\subseteq Y$.
    Let $T_Y = \subspaceTopology{T}{Y}$.
    Then $\closure{A}{Y}{T_Y} = \closure{A}{X}{T}\inter Y$.
\end{theorem}
\begin{proof}
    $T$ is a topology on $X$ and $Y\subseteq X$ and $A\subseteq X$ by \cref{subspace_of,subseteq,subseteq_transitive}.
    Let $B = \closure{A}{Y}{T_Y}$.
    Let $C = \closure{A}{X}{T}$.
    We have $C\inter Y$ is closed in $Y$ with respect to $T_Y$
        by \cref{closure_closed_in,closed_in_subspace_iff,subspace_of}.
    We have $A\subseteq C\inter Y$ by \cref{subseteq_closure,subseteq_inter_iff}.
    $T_Y$ is a topology on $Y$ by \cref{subspace_topology_is_topology}.
    Show that $B\subseteq C\inter Y$.
    \begin{subproof}
        Follows by \cref{closure_subseteq_closed}.
    \end{subproof}
    $B$ is closed in $Y$ with respect to $T_Y$ by \cref{closure_closed_in}.
    Take $D$ such that $D$ is closed in $X$ with respect to $T$ and $B = D\inter Y$
        by \cref{closed_in_subspace_iff,subspace_of}.
    We have $A\subseteq D$ by \cref{inter_lower_left,subseteq_closure,subseteq_transitive}.
    Hence $C\subseteq D$ by \cref{closure_subseteq_closed}.
    Show that $C\inter Y\subseteq B$.
    \begin{subproof}
        Follows by \cref{inter,elem_subseteq,inter_intro,subseteq}.
    \end{subproof}
    Therefore $B = C\inter Y$ by \cref{subseteq_antisymmetric}.
\end{proof}

% from §17 Item 17: intersects
\begin{definition}\label{intersects}
    $A$ intersects $B$ iff $A\inter B \neq \emptyset$.
\end{definition}
% from §17 Item 18: closure and open sets
\begin{theorem}\label{closure_iff_intersects_open}
    Assume $T$ is a topology on $X$.
    Suppose $A\subseteq X$.
    Suppose $x\in X$.
    Then $x\in\closure{A}{X}{T}$ iff for all $U$ such that
        $U$ is open in $X$ with respect to $T$ and $x\in U$ we have $U$ intersects $A$.
\end{theorem}
\begin{proof}
    Show that if $x\in\closure{A}{X}{T}$ then for all $U$ such that
        $U$ is open in $X$ with respect to $T$ and $x\in U$ we have $U$ intersects $A$.
    \begin{subproof}
        Assume $x\in\closure{A}{X}{T}$.
            Fix $U$ such that $U$ is open in $X$ with respect to $T$ and $x\in U$.
        Suppose not.
        We have $A\subseteq X\setminus U$ by \cref{intersects,inter_comm,subseteq_setminus}.
        Therefore $X\setminus U$ is closed in $X$ with respect to $T$ by
            \cref{setminus_setminus,inter_comm,topology_on,open_in,subseteq,pow_iff,inter_absorb_supseteq_right,setminus_subseteq,closed_in}.
        Thus $x\notin U$ by \cref{closure_subseteq_closed,elem_subseteq,setminus_elim_right}.
        Contradiction.
    \end{subproof}
    Show that if for all $U$ such that $U$ is open in $X$ with respect to $T$ and $x\in U$ we have $U$ intersects $A$
        then for all $C$ such that $C$ is closed in $X$ with respect to $T$ and
            for all $y\in A$ we have $y\in C$ we have $x\in C$.
    \begin{subproof}
        Assume for all $U$ such that $U$ is open in $X$ with respect to $T$ and $x\in U$ we have $U$ intersects $A$.
        Fix $C$ such that $C$ is closed in $X$ with respect to $T$ and
            for all $y\in A$ we have $y\in C$.
        Suppose not.
        We have $X\setminus C$ is open in $X$ with respect to $T$ and $x\in X\setminus C$ by \cref{closed_in,setminus_intro}.
        For all $y$ such that $y\in (X\setminus C)\inter A$ we have $y\in C$ and $y\notin C$
            by \cref{inter,setminus_elim_right}.
        Therefore $X\setminus C$ does not intersect $A$ by \cref{intersects,nonempty_inhabited}.
        Contradiction.
    \end{subproof}
    Show that if for all $U$ such that $U$ is open in $X$ with respect to $T$ and $x\in U$ we have $U$ intersects $A$
        then $x\in\closure{A}{X}{T}$.
    \begin{subproof}
        Assume for all $U$ such that $U$ is open in $X$ with respect to $T$ and $x\in U$ we have $U$ intersects $A$.
        For all $C$ such that $C$ is closed in $X$ with respect to $T$ and
            for all $y\in A$ we have $y\in C$ we have $x\in C$.
        Therefore $x\in\closure{A}{X}{T}$ by \cref{closure}.
    \end{subproof}
\end{proof}

% from §17 Item 19: closure and basis elements
\begin{theorem}\label{closure_iff_intersects_basis}
    Assume $T$ is a topology on $X$.
    Suppose $B$ is a basis for $T$ on $X$.
    Suppose $A\subseteq X$.
    Suppose $x\in X$.
    Then $x\in\closure{A}{X}{T}$ iff for all $B_1$ such that $B_1\in B$ and $x\in B_1$ we have $B_1$ intersects $A$.
\end{theorem}
\begin{proof}
    Show that if $x\in\closure{A}{X}{T}$ then for all $B_1$ such that $B_1\in B$ and $x\in B_1$ we have $B_1$ intersects $A$.
    \begin{subproof}
        Follows by \cref{basis_for_topology,basis_elem_open_in_generated,open_in,closure_iff_intersects_open}.
    \end{subproof}
    Show that if for all $B_1$ such that $B_1\in B$ and $x\in B_1$ we have $B_1$ intersects $A$ then
        $x\in\closure{A}{X}{T}$.
    \begin{subproof}
        Assume for all $B_1$ such that $B_1\in B$ and $x\in B_1$ we have $B_1$ intersects $A$.
        Show that for all $U$ such that $U$ is open in $X$ with respect to $T$ and $x\in U$ we have $U$ intersects $A$.
        \begin{subproof}
            Fix $U$ such that $U$ is open in $X$ with respect to $T$ and $x\in U$.
            $U\in\basisTopology{B}{X}$ by \cref{open_in,basis_for_topology}.
            Take $B_1\in B$ such that $x\in B_1$ and $B_1\subseteq U$ by \cref{topology_generated_by_basis}.
            $B_1$ intersects $A$.
            Take $y$ such that $y\in B_1\inter A$ by \cref{intersects,nonempty_inhabited}.
            We have $y\in U$ and $y\in A$ by \cref{inter,elem_subseteq}.
            Hence $U$ intersects $A$ by \cref{inter_intro,emptyset,intersects}.
        \end{subproof}
        Therefore $x\in\closure{A}{X}{T}$ by \cref{closure_iff_intersects_open}.
    \end{subproof}
\end{proof}

% from §17 Item 20: neighborhood
\begin{definition}\label{neighborhood}
    $U$ is a neighborhood of $x$ in $X$ with respect to $T$ iff
        $U$ is open in $X$ with respect to $T$ and $x\in U$.
\end{definition}

% from §17 Item 21: limit point
\begin{definition}\label{limit_point}
    $x$ is a limit point of $A$ in $X$ with respect to $T$ iff
        for all $U$ such that $U$ is a neighborhood of $x$ in $X$ with respect to $T$
        we have $U$ intersects $A\setminus\{x\}$.
\end{definition}

% from §17 Item 22: set of limit points
\begin{definition}\label{limit_points}
    $\limitPoints{A}{X}{T} = \{x\in X \mid \text{$x$ is a limit point of $A$ in $X$ with respect to $T$}\}$.
\end{definition}
% from §17 Item 23: closure equals set plus limit points
\begin{theorem}\label{closure_union_limit_points}
    Assume $T$ is a topology on $X$.
    Suppose $A\subseteq X$.
    Then $\closure{A}{X}{T} = A\union\limitPoints{A}{X}{T}$.
\end{theorem}
\begin{proof}
    Show that for all $x$ such that $x\in A\union\limitPoints{A}{X}{T}$ we have $x\in\closure{A}{X}{T}$.
    \begin{subproof}
        Fix $x$ such that $x\in A\union\limitPoints{A}{X}{T}$.
        We have $x\in A$ or $x\in\limitPoints{A}{X}{T}$ by \cref{union_iff}.
        \begin{byCase}
        \caseOf{$x\in A$.}
            We have $x\in\closure{A}{X}{T}$ by \cref{subseteq_closure,elem_subseteq}.
        \caseOf{$x\in\limitPoints{A}{X}{T}$.}
            We have $x\in X$ and $x$ is a limit point of $A$ in $X$ with respect to $T$ by \cref{limit_points}.
            Show that for all $U$ such that $U$ is open in $X$ with respect to $T$ and $x\in U$ we have $U$ intersects $A$.
            \begin{subproof}
                Fix $U$ such that $U$ is open in $X$ with respect to $T$ and $x\in U$.
                $U$ is a neighborhood of $x$ in $X$ with respect to $T$ by \cref{neighborhood}.
                We have $U$ intersects $A\setminus\{x\}$ by \cref{limit_point}.
                Take $y$ such that $y\in U\inter(A\setminus\{x\})$ by \cref{intersects,nonempty_inhabited}.
                We have $y\in U$ and $y\in A\setminus\{x\}$ by \cref{inter}.
                $y\in A$ by \cref{setminus_elim_left}.
                Therefore $U$ intersects $A$ by \cref{inter_intro,emptyset,intersects}.
            \end{subproof}
            Therefore $x\in\closure{A}{X}{T}$ by \cref{closure_iff_intersects_open}.
        \end{byCase}
    \end{subproof}
    We have $A\union\limitPoints{A}{X}{T}\subseteq\closure{A}{X}{T}$ by \cref{subseteq}.
    Show that for all $x$ such that $x\in\closure{A}{X}{T}$ we have $x\in A\union\limitPoints{A}{X}{T}$.
    \begin{subproof}
        Fix $x$ such that $x\in\closure{A}{X}{T}$.
        \begin{byCase}
        \caseOf{$x\in A$.}
            We have $x\in A\union\limitPoints{A}{X}{T}$ by \cref{union_intro_left}.
        \caseOf{$x\notin A$.}
            Show that for all $U$ such that $U$ is a neighborhood of $x$ in $X$ with respect to $T$ we have
                $U$ intersects $A\setminus\{x\}$.
            \begin{subproof}
            Fix $U$ such that $U$ is a neighborhood of $x$ in $X$ with respect to $T$.
                We have $U$ intersects $A$ by \cref{neighborhood,closure,closure_iff_intersects_open}.
                Take $y$ such that $y\in U\inter A$ by \cref{intersects,nonempty_inhabited}.
                $y\neq x$.
                Therefore $U$ intersects $A\setminus\{x\}$ by \cref{inter,singleton_iff,setminus_intro,inter_intro,emptyset,intersects}.
            \end{subproof}
            Therefore $x\in\limitPoints{A}{X}{T}$ by \cref{closure,limit_point,limit_points}.
            Hence $x\in A\union\limitPoints{A}{X}{T}$ by \cref{union_intro_right}.
        \end{byCase}
    \end{subproof}
    We have $\closure{A}{X}{T}\subseteq A\union\limitPoints{A}{X}{T}$ by \cref{subseteq}.
    Therefore $\closure{A}{X}{T} = A\union\limitPoints{A}{X}{T}$ by \cref{subseteq_antisymmetric}.
\end{proof}

% from §17 Item 24: closed iff contains limit points
\begin{corollary}\label{closed_iff_contains_limit_points}
    Assume $T$ is a topology on $X$.
    Suppose $A\subseteq X$.
    Then $A$ is closed in $X$ with respect to $T$ iff $\limitPoints{A}{X}{T}\subseteq A$.
\end{corollary}
\begin{proof}
    Show that if $A$ is closed in $X$ with respect to $T$ then $\limitPoints{A}{X}{T}\subseteq A$.
    \begin{subproof}
        Follows by \cref{closed_eq_closure,closure_union_limit_points,union_comm,union_eq_self_implies_subseteq}.
    \end{subproof}
    Show that if $\limitPoints{A}{X}{T}\subseteq A$ then $A$ is closed in $X$ with respect to $T$.
    \begin{subproof}
        Follows by \cref{union_absorb_subseteq_right,closure_union_limit_points,closure_closed_in}.
    \end{subproof}
\end{proof}

% from §17 Item 25: Hausdorff space
\begin{definition}\label{hausdorff}
    $X$ is Hausdorff with respect to $T$ iff
        $T$ is a topology on $X$ and
        for all $x_1,x_2$ such that $x_1\in X$ and $x_2\in X$ and $x_1\neq x_2$
        there exist $U_1,U_2$ such that
            $U_1$ is a neighborhood of $x_1$ in $X$ with respect to $T$ and
            $U_2$ is a neighborhood of $x_2$ in $X$ with respect to $T$ and
            $U_1$ is disjoint from $U_2$.
\end{definition}

% from §17 helper lemma 3: finite preserved by bijection
\begin{lemma}\label{finite_bijection}
    Suppose $A$ is finite.
    Suppose $f$ is a bijection from $A$ to $B$.
    Then $B$ is finite.
\end{lemma}
\begin{proof}
    \begin{byCase}
    \caseOf{$A$ is empty.}
        Show that for all $b$ we have $b\notin B$.
        \begin{subproof}
            Fix $b$.
            Suppose not.
            Then $b\in B$.
            $f\in\surj{A}{B}$ by \cref{bijections}.
            Take $a\in A$ such that $f(a) = b$ by \cref{surj}.
            Contradiction.
        \end{subproof}
        Hence $B$ is empty.
        Thus $B$ is finite by \cref{finite}.
    \caseOf{$A$ is not empty.}
        Take $k\in\naturalsPlus$ such that there exists a bijection from $\initialSegment{k}$ to $A$ by \cref{finite}.
        Take $g$ such that $g$ is a bijection from $\initialSegment{k}$ to $A$.
        $f\circ g$ is a bijection from $\initialSegment{k}$ to $B$ by \cref{bijection_circ}.
        Thus $B$ is finite by \cref{finite}.
    \end{byCase}
\end{proof}

% from §17 helper lemma 4: bijection to singleton family
\begin{lemma}\label{singleton_family_bijection}
    Let $S = \{\{x\} \mid x\in A\}$.
    Then there exists a bijection from $A$ to $S$.
\end{lemma}
\begin{proof}
    Let $f = \{(x,\{x\}) \mid x\in A\}$.
    For all $w$ such that $w\in f$ there exist $x,y$ such that $w = (x,y)$.
    For all $x,y,z$ such that $(x,y)\in f$ and $(x,z)\in f$ we have $y = z$ by \cref{pair_eq_iff}.
    Hence $f$ is a function by \cref{relation,rightunique}.
    $f\in\funs{A}{S}$ by
        \cref{dom_iff,pair_eq_iff,subseteq,subseteq_antisymmetric,function_apply_iff,funs_intro}.
    We have $f\in\surj{A}{S}$ by \cref{ran_iff,pair_eq_iff,subseteq,subseteq_antisymmetric,funs_surj_iff}.
    Therefore there exists a bijection from $A$ to $S$ by \cref{pair_eq_iff,singleton_elim,singleton_intro,injective,bijections}.
\end{proof}

% from §17 Item 26: finite point sets are closed in Hausdorff spaces
\begin{theorem}\label{finite_point_set_closed_hausdorff}
    Assume $X$ is Hausdorff with respect to $T$.
    Suppose $A\subseteq X$.
    Suppose $A$ is finite.
    Then $A$ is closed in $X$ with respect to $T$.
\end{theorem}
\begin{proof}
    $T$ is a topology on $X$ by \cref{hausdorff}.
    Show that for all $x$ such that $x\in X$ we have $\{x\}$ is closed in $X$ with respect to $T$.
    \begin{subproof}
        Fix $x$ such that $x\in X$.
        Show that for all $y$ such that $y\in\limitPoints{\{x\}}{X}{T}$ we have $y\in\{x\}$.
        \begin{subproof}
            Fix $y$ such that $y\in\limitPoints{\{x\}}{X}{T}$.
            $y\in X$ and $y$ is a limit point of $\{x\}$ in $X$ with respect to $T$ by \cref{limit_points}.
            Suppose not.
            Take $U_1,U_2$ such that
                $U_1$ is a neighborhood of $y$ in $X$ with respect to $T$ and
                $U_2$ is a neighborhood of $x$ in $X$ with respect to $T$ and
                $U_1$ is disjoint from $U_2$ by \cref{singleton_iff,hausdorff}.
            For all $z$ such that $z\in U_1\inter (\{x\}\setminus\{y\})$ we have $x\in U_1$ and $x\in U_2$
                by \cref{inter,setminus_elim_left,singleton_elim,neighborhood}.
            Hence $U_1$ does not intersect $\{x\}\setminus\{y\}$ by \cref{intersects,nonempty_inhabited,inter_intro,emptyset,intersects,disjoint}.
            Therefore $y$ is not a limit point of $\{x\}$ in $X$ with respect to $T$ by \cref{limit_point}.
            Contradiction.
        \end{subproof}
        Therefore $\{x\}$ is closed in $X$ with respect to $T$ by \cref{subseteq,singleton_subset_intro,closed_iff_contains_limit_points}.
    \end{subproof}
    Let $F = \{\{x\} \mid x\in A\}$.
    We have $F$ is finite and $F\subseteq\pow{X}$ by
        \cref{singleton_family_bijection,finite_bijection,elem_subseteq,singleton_subset_intro,powerset_intro,subseteq}.
    Show that for all $B$ such that $B\in F$ we have $B$ is closed in $X$ with respect to $T$.
    \begin{subproof}
        Fix $B$ such that $B\in F$.
        Take $x$ such that $x\in A$ and $\{x\} = B$.
        $x\in X$ by \cref{subseteq}.
        Therefore $B$ is closed in $X$ with respect to $T$.
    \end{subproof}
    Show that for all $z$ such that $z\in A$ we have $z\in\unions{F}$.
    \begin{subproof}
        Fix $z$ such that $z\in A$.
        We have $\{z\}\in F$.
        $z\in\{z\}$ by \cref{singleton_intro}.
        Therefore $z\in\unions{F}$ by \cref{unions_iff}.
    \end{subproof}
    We have $A\subseteq\unions{F}$ by \cref{subseteq}.
    Show that for all $z$ such that $z\in\unions{F}$ we have $z\in A$.
    \begin{subproof}
        Fix $z$ such that $z\in\unions{F}$.
        Take $B\in F$ such that $z\in B$ by \cref{unions_iff}.
        Take $x$ such that $x\in A$ and $\{x\} = B$.
        We have $z = x$ by \cref{singleton_elim}.
        Therefore $z\in A$.
    \end{subproof}
    We have $\unions{F}\subseteq A$ by \cref{subseteq}.
    Therefore $A = \unions{F}$ by \cref{subseteq_antisymmetric}.
    Hence $\unions{F}$ is closed in $X$ with respect to $T$ by \cref{closed_finite_unions}.
    Therefore $A$ is closed in $X$ with respect to $T$.
\end{proof}

% from §17 Item 27: T1 axiom
\begin{definition}\label{t1_axiom}
    $X$ satisfies the T-one axiom with respect to $T$ iff
        $T$ is a topology on $X$ and
        for all $A$ such that $A\subseteq X$ and $A$ is finite
        we have $A$ is closed in $X$ with respect to $T$.
\end{definition}

% from §17 helper lemma 5: avoiding a deleted singleton
\begin{lemma}\label{inter_empty_setminus_singleton_subset_singleton}
    Suppose $U\inter (A\setminus\{x\}) = \emptyset$.
    Then $U\inter A\subseteq\{x\}$.
\end{lemma}
\begin{proof}
    Follows by \cref{subseteq,setminus_intro,inter,inter_intro,emptyset,singleton_iff}.
\end{proof}

% from §17 Item 28: limit points in $T_1$ spaces
\begin{theorem}\label{limit_point_infinite_neighborhood}
    Assume $X$ satisfies the T-one axiom with respect to $T$.
    Suppose $A\subseteq X$.
    Suppose $x\in X$.
    Then $x$ is a limit point of $A$ in $X$ with respect to $T$ iff
        for all $U$ such that $U$ is a neighborhood of $x$ in $X$ with respect to $T$ we have $U\inter A$ is infinite.
\end{theorem}
\begin{proof}
    Show that if $x$ is a limit point of $A$ in $X$ with respect to $T$ then
        for all $U$ such that $U$ is a neighborhood of $x$ in $X$ with respect to $T$ we have $U\inter A$ is infinite.
    \begin{subproof}
        Assume $x$ is a limit point of $A$ in $X$ with respect to $T$.
        Fix $U$ such that $U$ is a neighborhood of $x$ in $X$ with respect to $T$.
        Suppose not.
        $U\inter A$ is finite.
        Let $B = U\inter (A\setminus\{x\})$.
        We have $B\subseteq U\inter A$ by \cref{inter,setminus_elim_left,subseteq,subseteq_inter_iff}.
        We have $B$ is finite by \cref{subseteq_finite_is_finite}.
        $B\subseteq X$ and $T$ is a topology on $X$ and $B$ is closed in $X$ with respect to $T$ by
            \cref{inter,subseteq,subseteq_transitive,t1_axiom}.
        Therefore $U\inter (X\setminus B)$ is open in $X$ with respect to $T$ by
            \cref{neighborhood,open_in,closed_in,topology_on,t1_axiom}.
        $x\notin B$.
        We have $x\in U\inter (X\setminus B)$ by \cref{neighborhood,setminus_intro,inter_intro}.
        Therefore $U\inter (X\setminus B)$ is a neighborhood of $x$ in $X$ with respect to $T$ by \cref{neighborhood}.
        For all $y$ such that $y\in (U\inter (X\setminus B)) \inter (A\setminus\{x\})$ we have
            $y\in B$ and $y\notin B$ by \cref{inter,inter_intro,setminus_elim_right}.
        Hence $U\inter (X\setminus B)$ does not intersect $A\setminus\{x\}$ by \cref{intersects,nonempty_inhabited}.
        Contradiction by \cref{limit_point}.
    \end{subproof}
    Show that if for all $U$ such that $U$ is a neighborhood of $x$ in $X$ with respect to $T$ we have $U\inter A$ is infinite then
        for all $U$ such that $U$ is a neighborhood of $x$ in $X$ with respect to $T$ we have
            $U$ intersects $A\setminus\{x\}$.
    \begin{subproof}
        Assume for all $U$ such that $U$ is a neighborhood of $x$ in $X$ with respect to $T$ we have $U\inter A$ is infinite.
        Fix $U$ such that $U$ is a neighborhood of $x$ in $X$ with respect to $T$.
        Suppose not.
        Hence $U\inter (A\setminus\{x\}) = \emptyset$.
        Therefore $U\inter A\subseteq\{x\}$ by \cref{inter_empty_setminus_singleton_subset_singleton}.
        Therefore $U\inter A$ is finite by \cref{singleton_subset_intro,pair_is_finite,upair_intro_left,singleton_intro,subseteq_finite_is_finite}.
        Contradiction.
    \end{subproof}
    Show that if for all $U$ such that $U$ is a neighborhood of $x$ in $X$ with respect to $T$ we have $U\inter A$ is infinite then
        $x$ is a limit point of $A$ in $X$ with respect to $T$.
    \begin{subproof}
        Assume for all $U$ such that $U$ is a neighborhood of $x$ in $X$ with respect to $T$ we have $U\inter A$ is infinite.
        We have for all $U$ such that $U$ is a neighborhood of $x$ in $X$ with respect to $T$ we have
            $U$ intersects $A\setminus\{x\}$.
        Therefore $x$ is a limit point of $A$ in $X$ with respect to $T$ by \cref{limit_point}.
    \end{subproof}
\end{proof}

% from §17 Item 29: sequence in a set
\begin{definition}\label{sequence_in}
    $s$ is a sequence in $X$ iff $s\in\funs{\naturalsPlus}{X}$.
\end{definition}

% from §17 Item 30: convergence of a sequence
\begin{definition}\label{converges_to}
    $s$ converges to $x$ in $X$ with respect to $T$ iff
        $s$ is a sequence in $X$ and
        for all $U$ such that $U$ is a neighborhood of $x$ in $X$ with respect to $T$
        there exists $N\in\naturalsPlus$ such that for all $n$ such that $n\in\naturalsPlus$ if $N\leq n$ then $s(n)\in U$.
\end{definition}

% from §17 Item 31: limits of sequences are unique in Hausdorff spaces
\begin{theorem}\label{sequence_limit_unique}
    Assume $X$ is Hausdorff with respect to $T$.
    Suppose $s$ is a sequence in $X$.
    Suppose $x\in X$.
    Suppose $y\in X$.
    Suppose $s$ converges to $x$ in $X$ with respect to $T$.
    Suppose $s$ converges to $y$ in $X$ with respect to $T$.
    Then $x = y$.
\end{theorem}
\begin{proof}
    Suppose not.
    Take $U_1,U_2$ such that
        $U_1$ is a neighborhood of $x$ in $X$ with respect to $T$ and
        $U_2$ is a neighborhood of $y$ in $X$ with respect to $T$ and
        $U_1$ is disjoint from $U_2$ by \cref{hausdorff}.
    Take $N_1\in\naturalsPlus$ such that for all $n$ such that $n\in\naturalsPlus$ if $N_1\leq n$ then $s(n)\in U_1$ by \cref{converges_to}.
    Take $N_2\in\naturalsPlus$ such that for all $n$ such that $n\in\naturalsPlus$ if $N_2\leq n$ then $s(n)\in U_2$ by \cref{converges_to}.
    Take $m\in\{N_1,N_2\}$ such that for all $a\in\{N_1,N_2\}$ we have $a\leq m$ by
        \cref{upair_iff,subseteq,pair_is_finite,emptyset,finite_naturals_maximum}.
    We have $N_1\leq m$ and $N_2\leq m$ by \cref{upair_iff}.
    Therefore $s(m)\in U_1$ and $s(m)\in U_2$.
    Hence $U_1$ intersects $U_2$ by \cref{inter_intro,emptyset,intersects}.
    Contradiction by \cref{disjoint}.
\end{proof}

% from §17 helper lemma 5: open right ray is open in order topology
\begin{lemma}\label{openrayright_open_in_order_topology}
    Assume $R$ is a simple order relation on $X$.
    Assume $X$ has more than one element.
    Assume $T$ is the order topology on $X$ with respect to $R$.
    Suppose $a\in X$.
    Then $\openrayright{R}{X}{a}$ is open in $X$ with respect to $T$.
\end{lemma}
\begin{proof}
    $\openrayright{R}{X}{a}\in \openRays{R}{X}$ and $\openrayright{R}{X}{a}\in T$ by
        \cref{open_rays_subbasis_on,subbasis_on,openray_right,subseteq,powerset_intro,open_rays,singleton_subset_intro,pair_is_finite,upair_intro_left,subseteq_finite_is_finite,singleton_intro,inters_singleton,finite_intersections,subbasis_finite_intersections_basis,basis_elem_open_in_generated,topology_generated_by_subbasis,open_rays_subbasis}.
    Hence $\openrayright{R}{X}{a}$ is open in $X$ with respect to $T$ by
        \cref{order_basis_is_basis,basis_topology_is_topology,order_topology,open_in}.
\end{proof}

% from §17 helper lemma 6: open left ray is open in order topology
\begin{lemma}\label{openrayleft_open_in_order_topology}
    Assume $R$ is a simple order relation on $X$.
    Assume $X$ has more than one element.
    Assume $T$ is the order topology on $X$ with respect to $R$.
    Suppose $a\in X$.
    Then $\openrayleft{R}{X}{a}$ is open in $X$ with respect to $T$.
\end{lemma}
\begin{proof}
    $\openrayleft{R}{X}{a}\in \openRays{R}{X}$ and $\openrayleft{R}{X}{a}\in T$ by
        \cref{open_rays_subbasis_on,subbasis_on,openray_left,subseteq,powerset_intro,open_rays,singleton_subset_intro,pair_is_finite,upair_intro_left,subseteq_finite_is_finite,singleton_intro,inters_singleton,finite_intersections,subbasis_finite_intersections_basis,basis_elem_open_in_generated,topology_generated_by_subbasis,open_rays_subbasis}.
    Hence $\openrayleft{R}{X}{a}$ is open in $X$ with respect to $T$ by
        \cref{order_basis_is_basis,basis_topology_is_topology,order_topology,open_in}.
\end{proof}

% from §17 Item 32: order topology is Hausdorff
\begin{theorem}\label{order_topology_hausdorff}
    Assume $R$ is a simple order relation on $X$.
    Assume $X$ has more than one element.
    Assume $T$ is the order topology on $X$ with respect to $R$.
    Then $X$ is Hausdorff with respect to $T$.
\end{theorem}
\begin{proof}
    Show that for all $x_1,x_2$ such that $x_1\in X$ and $x_2\in X$ and $x_1\neq x_2$ there exist $U_1,U_2$ such that
        $U_1$ is a neighborhood of $x_1$ in $X$ with respect to $T$ and
        $U_2$ is a neighborhood of $x_2$ in $X$ with respect to $T$ and
        $U_1$ is disjoint from $U_2$.
    \begin{subproof}
        Fix $x_1$.
        Fix $x_2$ such that $x_1\in X$ and $x_2\in X$ and $x_1\neq x_2$.
        We have $x_1\mathrel{R}x_2$ or $x_2\mathrel{R}x_1$ by \cref{simple_order_relation,connex}.
        \begin{byCase}
        \caseOf{$x_1\mathrel{R}x_2$.}
            \begin{byCase}
            \caseOf{there exists $z\in X$ such that $x_1\mathrel{R}z$ and $z\mathrel{R}x_2$.}
                Take $z\in X$ such that $x_1\mathrel{R}z$ and $z\mathrel{R}x_2$.
                Let $U_1 = \openrayleft{R}{X}{z}$.
                Let $U_2 = \openrayright{R}{X}{z}$.
                Therefore $U_1$ is a neighborhood of $x_1$ in $X$ with respect to $T$ and
                    $U_2$ is a neighborhood of $x_2$ in $X$ with respect to $T$
                    by \cref{openrayleft_open_in_order_topology,openrayright_open_in_order_topology,openray_left,openray_right,neighborhood}.
                $U_1$ is disjoint from $U_2$ by
                    \cref{disjoint,openray_left,openray_right,simple_order_relation,asymmetric}.
            \caseOf{there exists no $z\in X$ such that $x_1\mathrel{R}z$ and $z\mathrel{R}x_2$.}
                Let $U_1 = \openrayleft{R}{X}{x_2}$.
                Let $U_2 = \openrayright{R}{X}{x_1}$.
                Hence $U_1$ is a neighborhood of $x_1$ in $X$ with respect to $T$ and
                    $U_2$ is a neighborhood of $x_2$ in $X$ with respect to $T$
                    by \cref{openrayleft_open_in_order_topology,openrayright_open_in_order_topology,openray_left,openray_right,neighborhood}.
                $U_1$ is disjoint from $U_2$ by \cref{disjoint,openray_right,openray_left}.
            \end{byCase}
        \caseOf{$x_2\mathrel{R}x_1$.}
            \begin{byCase}
            \caseOf{there exists $z\in X$ such that $x_2\mathrel{R}z$ and $z\mathrel{R}x_1$.}
                Take $z\in X$ such that $x_2\mathrel{R}z$ and $z\mathrel{R}x_1$.
                Let $U_2 = \openrayleft{R}{X}{z}$.
                Let $U_1 = \openrayright{R}{X}{z}$.
                Therefore $U_1$ is a neighborhood of $x_1$ in $X$ with respect to $T$ and
                    $U_2$ is a neighborhood of $x_2$ in $X$ with respect to $T$
                    by \cref{openrayright_open_in_order_topology,openrayleft_open_in_order_topology,openray_left,openray_right,neighborhood}.
                $U_1$ is disjoint from $U_2$ by
                    \cref{disjoint,openray_right,openray_left,simple_order_relation,asymmetric}.
            \caseOf{there exists no $z\in X$ such that $x_2\mathrel{R}z$ and $z\mathrel{R}x_1$.}
                Let $U_2 = \openrayleft{R}{X}{x_1}$.
                Let $U_1 = \openrayright{R}{X}{x_2}$.
                Hence $U_1$ is a neighborhood of $x_1$ in $X$ with respect to $T$ and
                    $U_2$ is a neighborhood of $x_2$ in $X$ with respect to $T$
                    by \cref{openrayright_open_in_order_topology,openrayleft_open_in_order_topology,openray_left,openray_right,neighborhood}.
                $U_1$ is disjoint from $U_2$ by \cref{disjoint,openray_right,openray_left}.
            \end{byCase}
        \end{byCase}
    \end{subproof}
    Hence $X$ is Hausdorff with respect to $T$ by \cref{order_basis_is_basis,basis_topology_is_topology,order_topology,hausdorff}.
\end{proof}

% from §17 Item 33: product of Hausdorff spaces is Hausdorff
\begin{theorem}\label{product_hausdorff}
    Assume $X$ is Hausdorff with respect to $T$.
    Assume $Y$ is Hausdorff with respect to $T_1$.
    Let $T_2 = \productTopology{T}{T_1}{X}{Y}$.
    Then $X\times Y$ is Hausdorff with respect to $T_2$.
\end{theorem}
\begin{proof}
    $T$ is a topology on $X$ and $T_1$ is a topology on $Y$ by \cref{hausdorff}.
    $\productBasis{T}{T_1}{X}{Y}$ is a basis on $X\times Y$ by \cref{product_basis_is_basis}.
    We have $T_2$ is a topology on $X\times Y$ by \cref{product_basis_is_basis,basis_topology_is_topology,product_topology}.
    Show that for all $p_1,p_2$ such that $p_1\in X\times Y$ and $p_2\in X\times Y$ and $p_1\neq p_2$
        there exist $U_1,U_2$ such that
            $U_1$ is a neighborhood of $p_1$ in $X\times Y$ with respect to $T_2$ and
            $U_2$ is a neighborhood of $p_2$ in $X\times Y$ with respect to $T_2$ and
            $U_1$ is disjoint from $U_2$.
    \begin{subproof}
        Fix $p_1$.
        Fix $p_2$ such that $p_1\in X\times Y$ and $p_2\in X\times Y$ and $p_1\neq p_2$.
        Take $x_1,y_1$ such that $x_1\in X$ and $y_1\in Y$ and $p_1 = (x_1,y_1)$ by \cref{times_elem_is_tuple}.
        Take $x_2,y_2$ such that $x_2\in X$ and $y_2\in Y$ and $p_2 = (x_2,y_2)$ by \cref{times_elem_is_tuple}.
        $x_1\neq x_2$ or $y_1\neq y_2$ by \cref{pair_eq_iff}.
        \begin{byCase}
        \caseOf{$x_1\neq x_2$.}
            Take $U_1,U_2$ such that
                $U_1$ is a neighborhood of $x_1$ in $X$ with respect to $T$ and
                $U_2$ is a neighborhood of $x_2$ in $X$ with respect to $T$ and
                $U_1$ is disjoint from $U_2$ by \cref{hausdorff}.
            $Y$ is open in $Y$ with respect to $T_1$ by \cref{topology_on,open_in}.
            Let $W_1 = U_1\times Y$.
            Let $W_2 = U_2\times Y$.
            Therefore $W_1\in\productBasis{T}{T_1}{X}{Y}$ and $W_2\in\productBasis{T}{T_1}{X}{Y}$
                by \cref{open_in,neighborhood,topology_on,subseteq,pow_iff,times_subseteq_left,times_subseteq_right,subseteq_transitive,powerset_intro,product_basis}.
            $W_1\in T_2$ and $W_2\in T_2$ by \cref{basis_elem_open_in_generated,product_topology}.
            Hence $W_1$ is a neighborhood of $p_1$ in $X\times Y$ with respect to $T_2$ and
                $W_2$ is a neighborhood of $p_2$ in $X\times Y$ with respect to $T_2$ by
                \cref{open_in,neighborhood,times_tuple_intro}.
            $W_1$ is disjoint from $W_2$ by \cref{disjoint,times_elem_is_tuple,times_tuple_elim}.
        \caseOf{$y_1\neq y_2$.}
            Take $V_1,V_2$ such that
                $V_1$ is a neighborhood of $y_1$ in $Y$ with respect to $T_1$ and
                $V_2$ is a neighborhood of $y_2$ in $Y$ with respect to $T_1$ and
                $V_1$ is disjoint from $V_2$ by \cref{hausdorff}.
            $X$ is open in $X$ with respect to $T$ by \cref{topology_on,open_in}.
            Let $W_1 = X\times V_1$.
            Let $W_2 = X\times V_2$.
            Therefore $W_1\in\productBasis{T}{T_1}{X}{Y}$ and $W_2\in\productBasis{T}{T_1}{X}{Y}$
                by \cref{open_in,neighborhood,topology_on,subseteq,pow_iff,times_subseteq_left,times_subseteq_right,subseteq_transitive,powerset_intro,product_basis}.
            $W_1\in T_2$ and $W_2\in T_2$ by \cref{basis_elem_open_in_generated,product_topology}.
            Hence $W_1$ is a neighborhood of $p_1$ in $X\times Y$ with respect to $T_2$ and
                $W_2$ is a neighborhood of $p_2$ in $X\times Y$ with respect to $T_2$ by
                \cref{open_in,neighborhood,times_tuple_intro}.
            $W_1$ is disjoint from $W_2$ by \cref{disjoint,times_elem_is_tuple,times_tuple_elim}.
        \end{byCase}
    \end{subproof}
    Therefore $X\times Y$ is Hausdorff with respect to $T_2$ by \cref{hausdorff}.
\end{proof}

% from §17 Item 34: subspace of a Hausdorff space is Hausdorff
\begin{theorem}\label{subspace_hausdorff}
    Assume $X$ is Hausdorff with respect to $T$.
    Suppose $Y\subseteq X$.
    Suppose $T_1 = \subspaceTopology{T}{Y}$.
    Then $Y$ is Hausdorff with respect to $T_1$.
\end{theorem}
\begin{proof}
    $T_1$ is a topology on $Y$ by \cref{hausdorff,subspace_topology_is_topology}.
    Show that for all $y_1,y_2$ such that $y_1\in Y$ and $y_2\in Y$ and $y_1\neq y_2$
        there exist $V_1,V_2$ such that
            $V_1$ is a neighborhood of $y_1$ in $Y$ with respect to $T_1$ and
            $V_2$ is a neighborhood of $y_2$ in $Y$ with respect to $T_1$ and
            $V_1$ is disjoint from $V_2$.
    \begin{subproof}
        Fix $y_1$.
        Fix $y_2$ such that $y_1\in Y$ and $y_2\in Y$ and $y_1\neq y_2$.
        Take $U_1,U_2$ such that
            $U_1$ is a neighborhood of $y_1$ in $X$ with respect to $T$ and
            $U_2$ is a neighborhood of $y_2$ in $X$ with respect to $T$ and
            $U_1$ is disjoint from $U_2$ by \cref{elem_subseteq,hausdorff}.
        Let $V_1 = Y\inter U_1$.
        Let $V_2 = Y\inter U_2$.
        Therefore $V_1$ is a neighborhood of $y_1$ in $Y$ with respect to $T_1$ and
            $V_2$ is a neighborhood of $y_2$ in $Y$ with respect to $T_1$ by
            \cref{neighborhood,open_in,subspace_topology,inter_intro}.
        $V_1$ is disjoint from $V_2$ by \cref{disjoint,inter}.
    \end{subproof}
    Therefore $Y$ is Hausdorff with respect to $T_1$ by \cref{hausdorff}.
\end{proof}

\section{§18 Continuous Functions}

% from §18 Item 1: continuous function
\begin{definition}\label{continuous}
    $f$ is continuous from $X$ to $Y$ with respect to $T$ and $T_1$ iff
        $f$ is a function from $X$ to $Y$ and
        $T$ is a topology on $X$ and
        $T_1$ is a topology on $Y$ and
        for all $V$ such that $V$ is open in $Y$ with respect to $T_1$
        we have $\preimg{f}{V}$ is open in $X$ with respect to $T$.
\end{definition}

% from §18 helper lemma 1: preimage of an intersection under a function
\begin{lemma}\label{preimg_inter_function}
    Suppose $f$ is a function.
    Then $\preimg{f}{A\inter B} = \preimg{f}{A}\inter \preimg{f}{B}$.
\end{lemma}
\begin{proof}
    Follows by \cref{preimg_inter,inter,preimg_iff,function_rightunique,inter_intro,subseteq,subseteq_antisymmetric}.
\end{proof}

% from §18 helper lemma 2: image of a preimage is contained in the set
\begin{lemma}\label{img_preimg_subseteq}
    Suppose $f$ is a function.
    Then $\img{f}{\preimg{f}{B}}\subseteq B$.
\end{lemma}
\begin{proof}
    Follows by \cref{img_of_function_elim,inter_elim_right,preimg_iff,function_apply_elim,inter_elim_left,function_rightunique,subseteq}.
\end{proof}

% from §18 helper lemma 3: preimage of a complement
\begin{lemma}\label{preimg_setminus_function}
    Assume $f$ is a function from $X$ to $Y$.
    Suppose $B\subseteq Y$.
    Then $\preimg{f}{Y\setminus B} = X\setminus \preimg{f}{B}$.
\end{lemma}
\begin{proof}
    Show that $\preimg{f}{Y\setminus B}\subseteq X\setminus \preimg{f}{B}$.
    \begin{subproof}
        Follows by \cref{preimg_iff,funs,funs_is_function,subseteq,preimg_subseteq_dom,function_rightunique,setminus_elim_right,setminus_intro}.
    \end{subproof}
    Show that $X\setminus \preimg{f}{B}\subseteq \preimg{f}{Y\setminus B}$.
    \begin{subproof}
        Show that for all $x$ such that $x\in X\setminus \preimg{f}{B}$ we have
            $x\in\preimg{f}{Y\setminus B}$.
        \begin{subproof}
            Fix $x$ such that $x\in X\setminus \preimg{f}{B}$.
            We have $x\in\dom{f}$ by \cref{funs,setminus_elim_left}.
            Take $y$ such that $x\mathrel{f} y$ by \cref{dom_iff}.
            $y\in Y$ by \cref{funs,rels_type_ran}.
            $y\notin B$ by \cref{preimg_iff,setminus_elim_right}.
            Therefore $y\in Y\setminus B$ by \cref{setminus_intro}.
            Hence $x\in\preimg{f}{Y\setminus B}$ by \cref{preimg_iff}.
        \end{subproof}
        We have $X\setminus \preimg{f}{B}\subseteq \preimg{f}{Y\setminus B}$ by \cref{subseteq}.
    \end{subproof}
    Therefore $\preimg{f}{Y\setminus B} = X\setminus \preimg{f}{B}$ by \cref{subseteq_antisymmetric}.
\end{proof}

% from §18 Item 2: continuity implies image of closure is contained in closure of image
\begin{lemma}\label{continuous_closure_image_subset}
    Assume $f$ is continuous from $X$ to $Y$ with respect to $T$ and $T_1$.
    Suppose $A\subseteq X$.
    Then $\img{f}{\closure{A}{X}{T}}\subseteq \closure{\img{f}{A}}{Y}{T_1}$.
\end{lemma}
\begin{proof}
    We have $f\in\funs{X}{Y}$ and $f$ is a function and
        $T$ is a topology on $X$ and $T_1$ is a topology on $Y$
        by \cref{continuous,funs_is_function}.
    We have $f\in\rels{X}{Y}$ by \cref{funs_subseteq_rels,subseteq}.
    Therefore $\img{f}{A}\subseteq Y$ by
        \cref{funs_subseteq_rels,subseteq,rels_ran_subseteq,img_subseteq_ran,subseteq_transitive}.
    Show that for all $y$ such that $y\in\img{f}{\closure{A}{X}{T}}$ we have
        $y\in\closure{\img{f}{A}}{Y}{T_1}$.
        \begin{subproof}
            Fix $y$ such that $y\in\img{f}{\closure{A}{X}{T}}$.
            Take $x\in\dom{f}\inter \closure{A}{X}{T}$ such that $y = f(x)$ by \cref{img_of_function_elim}.
            We have $x\in\dom{f}$ and $x\mathrel{f} y$ and $y\in Y$
                by \cref{inter_elim_left,function_apply_elim,rels_type_ran}.
            Therefore $x\in\closure{A}{X}{T}$ and $x\in X$ by \cref{inter_elim_right,closure}.
        Show that for all $V$ such that $V$ is open in $Y$ with respect to $T_1$ and $y\in V$
            we have $V$ intersects $\img{f}{A}$.
        \begin{subproof}
            Fix $V$ such that $V$ is open in $Y$ with respect to $T_1$ and $y\in V$.
            Let $U = \preimg{f}{V}$.
            We have $U$ is open in $X$ with respect to $T$ and $x\in U$ by
                \cref{continuous,function_apply_elim,inter_elim_left,preimg_iff}.
            Take $z$ such that $z\in U\inter A$ by \cref{closure_iff_intersects_open,intersects,nonempty_inhabited}.
            Take $w\in V$ such that $z\mathrel{f} w$ by \cref{inter,preimg_iff}.
            Hence $V$ intersects $\img{f}{A}$ by \cref{inter,function_apply_iff,inter_intro,img_of_function_intro,emptyset,intersects}.
        \end{subproof}
        Therefore $y\in\closure{\img{f}{A}}{Y}{T_1}$ by \cref{closure_iff_intersects_open}.
    \end{subproof}
    Hence $\img{f}{\closure{A}{X}{T}}\subseteq \closure{\img{f}{A}}{Y}{T_1}$ by \cref{subseteq}.
\end{proof}

% from §18 Item 3: closure condition implies preimages of closed sets are closed
\begin{lemma}\label{closure_image_subset_implies_preimg_closed}
    Assume $f$ is a function from $X$ to $Y$ and
        $T$ is a topology on $X$ and
        $T_1$ is a topology on $Y$.
    Suppose for all $A$ such that $A\subseteq X$ we have
        $\img{f}{\closure{A}{X}{T}}\subseteq \closure{\img{f}{A}}{Y}{T_1}$.
    Then for all $B$ such that $B$ is closed in $Y$ with respect to $T_1$
        we have $\preimg{f}{B}$ is closed in $X$ with respect to $T$.
\end{lemma}
\begin{proof}
    Fix $B$ such that $B$ is closed in $Y$ with respect to $T_1$.
    Let $A = \preimg{f}{B}$.
    We have $f$ is right-unique and $f$ is a relation and $\dom{f} = X$
        by \cref{funs,funs_is_function,function_rightunique,funs_is_relation}.
    We have $\ran{f}\subseteq Y$ by \cref{funs_subseteq_rels,subseteq,rels_ran_subseteq}.
    $A\subseteq X$ by \cref{preimg_subseteq_dom}.
    $\img{f}{A}\subseteq B$ by \cref{img_preimg_subseteq}.
    Hence $\img{f}{A}\subseteq Y$ by \cref{img_subseteq_ran,subseteq_transitive}.
    Therefore $\closure{\img{f}{A}}{Y}{T_1}\subseteq B$ by \cref{closure_subseteq_closed}.
    Show that for all $x$ such that $x\in\closure{A}{X}{T}$ we have $x\in A$.
    \begin{subproof}
        Fix $x$ such that $x\in\closure{A}{X}{T}$.
        We have $x\in X$ by \cref{closure}.
        Hence $x\in\dom{f}$ by \cref{funs}.
        We have $x\in\dom{f}\inter\closure{A}{X}{T}$ by \cref{inter_intro}.
        Hence $f(x)\in\img{f}{\closure{A}{X}{T}}$ by \cref{img_of_function_intro}.
        Therefore $f(x)\in\closure{\img{f}{A}}{Y}{T_1}$ by \cref{subseteq}.
        Hence $f(x)\in B$ by \cref{elem_subseteq}.
        We have $x\mathrel{f} f(x)$ by \cref{function_apply_elim}.
        Therefore $x\in A$ by \cref{preimg_iff}.
    \end{subproof}
    Hence $A$ is closed in $X$ with respect to $T$ by
        \cref{subseteq,subseteq_closure,subseteq_antisymmetric,closure_closed_in}.
\end{proof}

% from §18 Item 4: preimages of closed sets imply continuity
\begin{lemma}\label{preimg_closed_implies_continuous}
    Assume $f$ is a function from $X$ to $Y$ and
        $T$ is a topology on $X$ and
        $T_1$ is a topology on $Y$.
    Suppose for all $B$ such that $B$ is closed in $Y$ with respect to $T_1$
        we have $\preimg{f}{B}$ is closed in $X$ with respect to $T$.
    Then $f$ is continuous from $X$ to $Y$ with respect to $T$ and $T_1$.
\end{lemma}
\begin{proof}
    We have $f\in\funs{X}{Y}$ and $\dom{f} = X$ by \cref{funs}.
    Show that for all $V$ such that $V$ is open in $Y$ with respect to $T_1$
        we have $\preimg{f}{V}$ is open in $X$ with respect to $T$.
    \begin{subproof}
        Fix $V$ such that $V$ is open in $Y$ with respect to $T_1$.
        Let $B = Y\setminus V$.
        We have $Y\setminus B = V$ by \cref{setminus_subseteq,setminus_setminus,open_in,topology_on,subseteq,pow_iff,inter_absorb_supseteq_left}.
        Therefore $B$ is closed in $Y$ with respect to $T_1$ by \cref{setminus_subseteq,closed_in}.
        Hence $\preimg{f}{B}$ is closed in $X$ with respect to $T$ by \cref{subseteq}.
        Therefore $\preimg{f}{V}$ is open in $X$ with respect to $T$ by
            \cref{closed_in,funs_is_function,preimg_setminus_function,subseteq,open_in}.
    \end{subproof}
    Hence $f$ is continuous from $X$ to $Y$ with respect to $T$ and $T_1$ by \cref{continuous}.
\end{proof}

% from §18 Item 5: continuity at a point
\begin{definition}\label{continuous_at}
    $f$ is continuous at $x$ from $X$ to $Y$ with respect to $T$ and $T_1$ iff
        $f$ is a function from $X$ to $Y$ and
        $T$ is a topology on $X$ and
        $T_1$ is a topology on $Y$ and
        $x\in X$ and
        for all $V$ such that $V$ is a neighborhood of $f(x)$ in $Y$ with respect to $T_1$
            there exists $U$ such that
                $U$ is a neighborhood of $x$ in $X$ with respect to $T$ and
                $\img{f}{U}\subseteq V$.
\end{definition}

% from §18 Item 6: continuity implies continuity at a point
\begin{lemma}\label{continuous_implies_continuous_at}
    Assume $f$ is continuous from $X$ to $Y$ with respect to $T$ and $T_1$.
    Suppose $x\in X$.
    Then $f$ is continuous at $x$ from $X$ to $Y$ with respect to $T$ and $T_1$.
\end{lemma}
\begin{proof}
    Follows by \cref{neighborhood,continuous,funs,function_apply_elim,preimg_iff,funs_is_function,img_preimg_subseteq,continuous_at}.
\end{proof}

% from §18 Item 7: continuity at each point implies continuity
\begin{lemma}\label{continuous_at_implies_continuous}
    Assume $f$ is a function from $X$ to $Y$ and
        $T$ is a topology on $X$ and
        $T_1$ is a topology on $Y$.
    Suppose for all $x\in X$ we have $f$ is continuous at $x$ from $X$ to $Y$ with respect to $T$ and $T_1$.
    Then $f$ is continuous from $X$ to $Y$ with respect to $T$ and $T_1$.
\end{lemma}
\begin{proof}
    We have $f\in\funs{X}{Y}$ and $\dom{f} = X$ by \cref{funs}.
    Show that for all $V$ such that $V$ is open in $Y$ with respect to $T_1$
        we have $\preimg{f}{V}$ is open in $X$ with respect to $T$.
    \begin{subproof}
        Fix $V$ such that $V$ is open in $Y$ with respect to $T_1$.
        Let $M = \{U\in T \mid U\subseteq \preimg{f}{V}\}$.
        Show that for all $x$ such that $x\in\preimg{f}{V}$ we have $x\in\unions{M}$.
        \begin{subproof}
            Fix $x$ such that $x\in\preimg{f}{V}$.
            Take $U$ such that
                $U$ is a neighborhood of $x$ in $X$ with respect to $T$ and
                $\img{f}{U}\subseteq V$ by
                \cref{preimg_subseteq_dom,subseteq,funs,funs_is_function,preimg_iff,function_apply_iff,neighborhood,continuous_at}.
            Therefore $x\in\unions{M}$ by
                \cref{preimg_subseteq_dom,subseteq,funs,funs_is_function,neighborhood,open_in,topology_on,pow_iff,inter_intro,img_of_function_intro,function_apply_elim,preimg_iff,unions_iff}.
        \end{subproof}
        Hence $\preimg{f}{V}\subseteq \unions{M}$ by \cref{subseteq}.
        Therefore $\unions{M}\in T$ and $\unions{M}\subseteq \preimg{f}{V}$
            by \cref{subseteq,pow_iff,topology_on,unions_subseteq_of_powerset_is_subseteq}.
        Hence $\preimg{f}{V}$ is open in $X$ with respect to $T$ by \cref{subseteq_antisymmetric,open_in}.
    \end{subproof}
    Therefore $f$ is continuous from $X$ to $Y$ with respect to $T$ and $T_1$ by \cref{continuous}.
\end{proof}

% from §18 Item 8: homeomorphism
\begin{definition}\label{homeomorphism}
    $f$ is a homeomorphism between $X$ and $Y$ with respect to $T$ and $T_1$ iff
        $f$ is a bijection from $X$ to $Y$ and
        $f$ is continuous from $X$ to $Y$ with respect to $T$ and $T_1$ and
        $\converse{f}$ is continuous from $Y$ to $X$ with respect to $T_1$ and $T$.
\end{definition}

% from §18 Item 9: topological property
\begin{definition}\label{topological_property}
    $P$ is a topological property iff
        for all $X,Y,T,T_1,f$ such that
            $f$ is a homeomorphism between $X$ and $Y$ with respect to $T$ and $T_1$
        we have $(X,T)\in P$ iff $(Y,T_1)\in P$.
\end{definition}

% from §18 Item 10: topological embedding
\begin{definition}\label{topological_embedding}
    $f$ is a topological embedding from $X$ to $Y$ with respect to $T$ and $T_1$ iff
        $f$ is a function from $X$ to $Y$ and
        $f$ is injective and
        $f$ is continuous from $X$ to $Y$ with respect to $T$ and $T_1$ and
        $f$ is a homeomorphism between $X$ and $\img{f}{X}$ with respect to
            $T$ and $\subspaceTopology{T_1}{\img{f}{X}}$.
\end{definition}
% from §18 Item 11: constant function is continuous
\begin{lemma}\label{constant_continuous}
    Assume $T$ is a topology on $X$ and $T_1$ is a topology on $Y$.
    Suppose $f$ is a function from $X$ to $Y$.
    Suppose there exists $y_0\in Y$ such that for all $x\in X$ we have $f(x) = y_0$.
    Then $f$ is continuous from $X$ to $Y$ with respect to $T$ and $T_1$.
\end{lemma}
\begin{proof}
    Show that for all $V$ such that $V$ is open in $Y$ with respect to $T_1$
        we have $\preimg{f}{V}$ is open in $X$ with respect to $T$.
    \begin{subproof}
        Fix $V$ such that $V$ is open in $Y$ with respect to $T_1$.
        Take $y_0\in Y$ such that for all $x\in X$ we have $f(x) = y_0$.
        \begin{byCase}
        \caseOf{$y_0\in V$.}
            We have $\preimg{f}{V}$ is open in $X$ with respect to $T$ by
                \cref{funs,funs_is_function,funs_is_relation,function_apply_elim,preimg_iff,subseteq,preimg_subseteq_dom,subseteq_transitive,subseteq_antisymmetric,open_in,topology_on}.
        \caseOf{$y_0\notin V$.}
            We have $\preimg{f}{V}$ is open in $X$ with respect to $T$ by
                \cref{preimg_iff,funs,funs_is_function,function_rightunique,funs_is_relation,preimg_subseteq_dom,subseteq,function_apply_iff,emptyset_subseteq,subseteq_antisymmetric,open_in,topology_on}.
        \end{byCase}
    \end{subproof}
    Hence $f$ is continuous from $X$ to $Y$ with respect to $T$ and $T_1$ by \cref{continuous}.
\end{proof}

% from §18 Item 12: inclusion is continuous
\begin{lemma}\label{inclusion_continuous}
    Assume $A$ is a subspace of $X$ with respect to $T$.
    Let $T_A = \subspaceTopology{T}{A}$.
    Let $j = \identity{A}$.
    Then $j$ is continuous from $A$ to $X$ with respect to $T_A$ and $T$.
\end{lemma}
\begin{proof}
    $T_A$ is a topology on $A$ by \cref{subspace_topology_is_topology,subspace_of}.
    We have $T$ is a topology on $X$ and $A\subseteq X$ by \cref{subspace_of}.
    Hence $j\in\funs{A}{A}$ and $j\in\funs{A}{X}$ by \cref{id_elem_funs,funs_weaken_codom}.
    Show that for all $U$ such that $U$ is open in $X$ with respect to $T$ we have
        $\preimg{j}{U} = A\inter U$.
    \begin{subproof}
        Fix $U$ such that $U$ is open in $X$ with respect to $T$.
        Show that $\preimg{j}{U}\subseteq A\inter U$.
        \begin{subproof}
            For all $x$ such that $x\in\preimg{j}{U}$ we have $x\in A\inter U$
                by \cref{preimg_iff,id_iff,inter_intro}.
        \end{subproof}
        Show that $A\inter U\subseteq\preimg{j}{U}$.
        \begin{subproof}
            For all $x$ such that $x\in A\inter U$ we have $x\in\preimg{j}{U}$
                by \cref{preimg_iff,id_iff,inter}.
        \end{subproof}
        Hence $\preimg{j}{U} = A\inter U$ by \cref{subseteq_antisymmetric}.
    \end{subproof}
    Show that for all $U$ such that $U$ is open in $X$ with respect to $T$ we have
        $\preimg{j}{U}$ is open in $A$ with respect to $T_A$.
    \begin{subproof}
        Follows by \cref{preimg_iff,id_iff,inter_intro,inter,subseteq,subseteq_antisymmetric,subspace_topology,open_in}.
    \end{subproof}
    Hence $j$ is continuous from $A$ to $X$ with respect to $T_A$ and $T$ by \cref{continuous}.
\end{proof}

% from §18 Item 13: composites of continuous functions are continuous
\begin{lemma}\label{continuous_composite}
    Assume $f$ is continuous from $X$ to $Y$ with respect to $T$ and $T_1$.
    Assume $g$ is continuous from $Y$ to $Z$ with respect to $T_1$ and $T_2$.
    Then $g\circ f$ is continuous from $X$ to $Z$ with respect to $T$ and $T_2$.
\end{lemma}
\begin{proof}
    $g\circ f\in\funs{X}{Z}$ by \cref{continuous,funs_circ}.
    $T$ is a topology on $X$ and $T_2$ is a topology on $Z$ by \cref{continuous}.
    Show that for all $V$ such that $V$ is open in $Z$ with respect to $T_2$ we have
        $\preimg{g\circ f}{V}$ is open in $X$ with respect to $T$.
    \begin{subproof}
        Fix $V$ such that $V$ is open in $Z$ with respect to $T_2$.
        Show that $\preimg{g\circ f}{V}\subseteq\preimg{f}{\preimg{g}{V}}$.
        \begin{subproof}
            Show that for all $x$ such that $x\in\preimg{g\circ f}{V}$ we have $x\in\preimg{f}{\preimg{g}{V}}$.
            \begin{subproof}
                Fix $x$ such that $x\in\preimg{g\circ f}{V}$.
                Take $z\in V$ such that $(x,z)\in g\circ f$ by \cref{preimg_iff}.
                Take $y$ such that $(x,y)\in f$ and $(y,z)\in g$ by \cref{circ_iff}.
                Hence $y\in\preimg{g}{V}$ by \cref{preimg_iff}.
                Therefore $x\in\preimg{f}{\preimg{g}{V}}$ by \cref{preimg_iff}.
            \end{subproof}
        \end{subproof}
        Hence $\preimg{g\circ f}{V}\subseteq\preimg{f}{\preimg{g}{V}}$ by \cref{subseteq}.
        Show that $\preimg{f}{\preimg{g}{V}}\subseteq\preimg{g\circ f}{V}$.
        \begin{subproof}
            Show that for all $x$ such that $x\in\preimg{f}{\preimg{g}{V}}$ we have $x\in\preimg{g\circ f}{V}$.
            \begin{subproof}
                Fix $x$ such that $x\in\preimg{f}{\preimg{g}{V}}$.
                Take $y\in\preimg{g}{V}$ such that $(x,y)\in f$ by \cref{preimg_iff}.
                Take $z\in V$ such that $(y,z)\in g$ by \cref{preimg_iff}.
                Hence $(x,z)\in g\circ f$ by \cref{circ_iff}.
                Therefore $x\in\preimg{g\circ f}{V}$ by \cref{preimg_iff}.
            \end{subproof}
        \end{subproof}
        Hence $\preimg{f}{\preimg{g}{V}}\subseteq\preimg{g\circ f}{V}$ by \cref{subseteq}.
        Therefore $\preimg{g\circ f}{V} = \preimg{f}{\preimg{g}{V}}$ by \cref{subseteq_antisymmetric}.
        $\preimg{g}{V}$ is open in $Y$ with respect to $T_1$ by \cref{continuous}.
        Hence $\preimg{f}{\preimg{g}{V}}$ is open in $X$ with respect to $T$ by \cref{continuous}.
        Therefore $\preimg{g\circ f}{V}$ is open in $X$ with respect to $T$.
    \end{subproof}
    Therefore $g\circ f$ is continuous from $X$ to $Z$ with respect to $T$ and $T_2$ by \cref{continuous}.
\end{proof}

% from §18 Item 14: restricting the domain
\begin{lemma}\label{continuous_restrict_domain}
    Assume $f$ is continuous from $X$ to $Y$ with respect to $T$ and $T_1$.
    Suppose $A$ is a subspace of $X$ with respect to $T$.
    Let $T_A = \subspaceTopology{T}{A}$.
    Let $g = \restrl{f}{A}$.
    Then $g$ is continuous from $A$ to $Y$ with respect to $T_A$ and $T_1$.
\end{lemma}
    \begin{proof}
    $T_A$ is a topology on $A$ by \cref{subspace_topology_is_topology,subspace_of}.
    Show that for all $V$ such that $V$ is open in $Y$ with respect to $T_1$ we have
        $\preimg{g}{V} = A\inter \preimg{f}{V}$.
    \begin{subproof}
        Fix $V$ such that $V$ is open in $Y$ with respect to $T_1$.
        Show that $\preimg{g}{V}\subseteq A\inter \preimg{f}{V}$.
        \begin{subproof}
            For all $x$ such that $x\in\preimg{g}{V}$ we have $x\in A\inter \preimg{f}{V}$
                by \cref{preimg_iff,restrl_iff,inter_intro}.
        \end{subproof}
        Show that $A\inter \preimg{f}{V}\subseteq\preimg{g}{V}$.
        \begin{subproof}
            For all $x$ such that $x\in A\inter \preimg{f}{V}$ we have $x\in\preimg{g}{V}$
                by \cref{inter,preimg_iff,restrl_iff}.
        \end{subproof}
        Hence $\preimg{g}{V} = A\inter \preimg{f}{V}$ by \cref{subseteq_antisymmetric}.
    \end{subproof}
    Show that for all $V$ such that $V$ is open in $Y$ with respect to $T_1$ we have
        $\preimg{g}{V}$ is open in $A$ with respect to $T_A$.
    \begin{subproof}
        Follows by \cref{preimg_iff,restrl_iff,inter,inter_intro,subseteq,subseteq_antisymmetric,continuous,open_in,subspace_topology}.
    \end{subproof}
    We have $f\in\funs{X}{Y}$ and $g$ is a function and $\dom{g} = A$
        by \cref{continuous,funs_is_function,restrl_of_function_is_function,dom_of_restrl_eq_inter,funs,inter_absorb_supseteq_left,subspace_of}.
    Hence $g$ is right-unique and $g$ is a relation.
    $A\subseteq X$ and $\dom{f} = X$ and $A\subseteq \dom{f}$
        by \cref{subspace_of,funs,subseteq_transitive}.
    $\ran{g} = \img{f}{A}$ and $\img{f}{A}\subseteq \ran{f}$ by \cref{restrl_ran,img_subseteq_ran}.
    Therefore $f\in\rels{X}{Y}$ and $\ran{f}\subseteq Y$ by \cref{funs_subseteq_rels,subseteq,rels_ran_subseteq}.
    Hence $g\in\funs{A}{Y}$ by
        \cref{function_apply_elim,ran_iff,subseteq,funs_intro,subseteq_transitive}.
    Therefore $T_1$ is a topology on $Y$ and
        $g$ is continuous from $A$ to $Y$ with respect to $T_A$ and $T_1$ by \cref{continuous}.
\end{proof}

% from §18 Item 15: restricting the range
\begin{lemma}\label{continuous_restrict_range}
    Assume $f$ is continuous from $X$ to $Y$ with respect to $T$ and $T_1$.
    Suppose $Z$ is a subspace of $Y$ with respect to $T_1$.
    Suppose $\img{f}{X}\subseteq Z$.
    Let $T_Z = \subspaceTopology{T_1}{Z}$.
    Then $f$ is continuous from $X$ to $Z$ with respect to $T$ and $T_Z$.
\end{lemma}
\begin{proof}
    $f\in\funs{X}{Y}$ and $\dom{f} = X$ by \cref{continuous,funs}.
    We have $T_1$ is a topology on $Y$ and $T_Z$ is a topology on $Z$
        by \cref{continuous,subspace_topology_is_topology,subspace_of}.
    For all $B$ such that $B$ is open in $Z$ with respect to $T_Z$
        there exists $U\in T_1$ such that $B = Z\inter U$ by \cref{subspace_topology,open_in}.
    Show that for all $B$ such that $B$ is open in $Z$ with respect to $T_Z$
        we have $\preimg{f}{B}$ is open in $X$ with respect to $T$.
    \begin{subproof}
        Fix $B$ such that $B$ is open in $Z$ with respect to $T_Z$.
        Take $U\in T_1$ such that $B = Z\inter U$ by \cref{subspace_topology,open_in}.
        Show that $\preimg{f}{B}\subseteq\preimg{f}{U}$.
        \begin{subproof}
            For all $x$ such that $x\in\preimg{f}{B}$ we have $x\in\preimg{f}{U}$
                by \cref{preimg_iff,inter_elim_right}.
        \end{subproof}
        Show that $\preimg{f}{U}\subseteq\preimg{f}{B}$.
        \begin{subproof}
            For all $x$ such that $x\in\preimg{f}{U}$ we have $x\in\preimg{f}{B}$
                by \cref{preimg_iff,preimg_subseteq_dom,funs,funs_is_function,funs_is_relation,function_apply_iff,img_of_function_intro,inter_intro,subseteq}.
        \end{subproof}
        Therefore $\preimg{f}{B} = \preimg{f}{U}$ by \cref{subseteq_antisymmetric}.
        We have $U$ is open in $Y$ with respect to $T_1$ by \cref{open_in}.
        Therefore $\preimg{f}{B}$ is open in $X$ with respect to $T$ by \cref{continuous}.
    \end{subproof}
    $f\in\funs{X}{Z}$ by
        \cref{function_apply_elim,ran_iff,subseteq,funs_intro,funs_is_function,img_dom}.
    Therefore $f$ is continuous from $X$ to $Z$ with respect to $T$ and $T_Z$ by \cref{continuous}.
\end{proof}
% from §18 Item 16: expanding the range
\begin{lemma}\label{continuous_expand_range}
    Assume $f$ is continuous from $X$ to $Y$ with respect to $T$ and $T_1$.
    Suppose $Y$ is a subspace of $Z$ with respect to $T_2$.
    Let $T_Y = \subspaceTopology{T_2}{Y}$.
    Suppose $T_1 = T_Y$.
    Let $j = \identity{Y}$.
    Then $j\circ f$ is continuous from $X$ to $Z$ with respect to $T$ and $T_2$.
\end{lemma}
\begin{proof}
    Follows by \cref{inclusion_continuous,subseteq,continuous_composite}.
\end{proof}

% from §18 Item 17: local formulation of continuity
\begin{lemma}\label{continuous_local}
    Assume $T$ is a topology on $X$.
    Assume $T_1$ is a topology on $Y$.
    Assume $f$ is a function from $X$ to $Y$.
    Suppose $M\subseteq T$.
    Suppose $\unions{M} = X$.
    Suppose for all $U$ such that $U\in M$ we have
        $\restrl{f}{U}$ is continuous from $U$ to $Y$ with respect to
            $\subspaceTopology{T}{U}$ and $T_1$.
    Then $f$ is continuous from $X$ to $Y$ with respect to $T$ and $T_1$.
\end{lemma}
\begin{proof}
    Show that for all $V$ such that $V$ is open in $Y$ with respect to $T_1$
        we have $\preimg{f}{V}$ is open in $X$ with respect to $T$.
    \begin{subproof}
        Fix $V$ such that $V$ is open in $Y$ with respect to $T_1$.
        Let $N = \{W\in T \mid W\subseteq \preimg{f}{V}\}$.
        Show that for all $x$ such that $x\in\preimg{f}{V}$ we have $x\in\unions{N}$.
        \begin{subproof}
            Fix $x$ such that $x\in\preimg{f}{V}$.
            Take $U\in M$ such that $x\in U$ by \cref{unions_iff,subseteq,preimg_subseteq_dom,funs}.
            Let $g = \restrl{f}{U}$.
            We have $\preimg{g}{V}\in N$ by \cref{continuous,open_in,subseteq,topology_on,pow_iff,subspace_of,open_in_subspace_open_in_space,preimg_iff,restrl_iff}.
            Therefore $x\in\unions{N}$ by \cref{restrl_iff,preimg_iff,unions_iff}.
        \end{subproof}
        Hence $\preimg{f}{V}\subseteq \unions{N}$ by \cref{subseteq}.
        Therefore $\preimg{f}{V}$ is open in $X$ with respect to $T$ by
            \cref{subseteq,pow_iff,topology_on,unions_subseteq_of_powerset_is_subseteq,subseteq_antisymmetric,open_in}.
    \end{subproof}
    Therefore $f$ is continuous from $X$ to $Y$ with respect to $T$ and $T_1$ by \cref{continuous}.
\end{proof}

% from §18 helper lemma 4: union of two compatible functions is a function
\begin{lemma}\label{union_compatible_functions}
    Suppose $f$ is a function.
    Suppose $g$ is a function.
    Suppose for all $x$ such that $x\in\dom{f}\inter\dom{g}$ we have $f(x) = g(x)$.
    Then $f\union g$ is a function.
\end{lemma}
\begin{proof}
    $f$ is a relation and $g$ is a relation and $f\union g$ is a relation
        by \cref{funs_is_relation,union_relations_is_relation}.
    Show that for all $x,y,z$ such that $(x,y)\in f\union g$ and $(x,z)\in f\union g$ we have $y = z$.
    \begin{subproof}
        Fix $x$.
        Fix $y$.
        Fix $z$ such that $(x,y)\in f\union g$ and $(x,z)\in f\union g$.
        \begin{byCase}
        \caseOf{$(x,y)\in f$ and $(x,z)\in f$.}
            We have $y = z$ by \cref{function_rightunique,rightunique}.
        \caseOf{$(x,y)\in g$ and $(x,z)\in g$.}
            We have $y = z$ by \cref{function_rightunique,rightunique}.
        \caseOf{$(x,y)\in f$ and $(x,z)\in g$.}
            We have $y = z$ by \cref{dom_intro,inter_intro,subseteq,function_apply_iff}.
        \caseOf{$(x,y)\in g$ and $(x,z)\in f$.}
            We have $y = z$ by \cref{dom_intro,inter_intro,subseteq,function_apply_iff}.
        \end{byCase}
    \end{subproof}
    Hence $f\union g$ is a function by \cref{rightunique,relation}.
\end{proof}

% from §18 Item 18: pasting lemma
\begin{theorem}\label{pasting_lemma}
    Assume $X = A\union B$.
    Assume $A$ is closed in $X$ with respect to $T$.
    Assume $B$ is closed in $X$ with respect to $T$.
    Assume $f$ is continuous from $A$ to $Y$ with respect to $\subspaceTopology{T}{A}$ and $T_1$.
    Assume $g$ is continuous from $B$ to $Y$ with respect to $\subspaceTopology{T}{B}$ and $T_1$.
    Suppose for all $x$ such that $x\in A\inter B$ we have $f(x) = g(x)$.
    Let $h = f\union g$.
    Then $h$ is continuous from $X$ to $Y$ with respect to $T$ and $T_1$.
\end{theorem}
\begin{proof}
    $f\in\funs{A}{Y}$ and $g\in\funs{B}{Y}$
        and $\dom{f} = A$ and $\dom{g} = B$ by \cref{continuous,funs}.
    $h$ is a function by \cref{inter,funs_is_function,subseteq,union_compatible_functions}.
    Show that for all $C$ such that $C$ is closed in $Y$ with respect to $T_1$
        we have $\preimg{h}{C}$ is closed in $X$ with respect to $T$.
    \begin{subproof}
        Fix $C$ such that $C$ is closed in $Y$ with respect to $T_1$.
        Let $A_1 = \preimg{f}{C}$.
        Let $B_1 = \preimg{g}{C}$.
        Let $D = Y\setminus C$.
        $\preimg{f}{D}$ is open in $A$ with respect to $\subspaceTopology{T}{A}$ and
            $\preimg{g}{D}$ is open in $B$ with respect to $\subspaceTopology{T}{B}$ by \cref{closed_in,continuous}.
        $\preimg{f}{D} = A\setminus A_1$ and $\preimg{g}{D} = B\setminus B_1$ by
            \cref{continuous,closed_in,setminus_subseteq,preimg_setminus_function}.
        We have $A_1\subseteq A$ and $B_1\subseteq B$ by \cref{continuous,preimg_subseteq_dom,funs}.
        Therefore $A_1 = A\setminus \preimg{f}{D}$ and $B_1 = B\setminus \preimg{g}{D}$
            by \cref{subseteq,setminus_setminus,inter_absorb_supseteq_left}.
        Hence $A_1$ is closed in $A$ with respect to $\subspaceTopology{T}{A}$ and
            $B_1$ is closed in $B$ with respect to $\subspaceTopology{T}{B}$ and
            $X\setminus A$ is open in $X$ with respect to $T$ by \cref{closed_in}.
        $A_1$ is closed in $X$ with respect to $T$ and $B_1$ is closed in $X$ with respect to $T$
            by \cref{closed_in,open_in,topology_on,closed_in_closed_subspace}.
        For all $x$ such that $x\in\preimg{h}{C}$ we have $x\in A_1\union B_1$ by
            \cref{preimg_iff,union_iff}.
        We have $\preimg{h}{C}\subseteq A_1\union B_1$ by \cref{subseteq}.
        For all $x$ such that $x\in A_1\union B_1$ we have $x\in\preimg{h}{C}$ by
            \cref{union_iff,preimg_iff,union_intro_left,union_intro_right}.
        Therefore $A_1\union B_1\subseteq \preimg{h}{C}$ by \cref{subseteq}.
        Therefore $\preimg{h}{C} = A_1\union B_1$ by \cref{subseteq_antisymmetric}.
        Hence $X\setminus (A_1\union B_1)$ is open in $X$ with respect to $T$
            by \cref{closed_in,open_in,topology_on,setminus_union,subseteq}.
        Therefore $A_1\union B_1$ is closed in $X$ with respect to $T$ by
            \cref{closed_in,preimg_subseteq_dom,funs,subseteq_transitive,union_subsets_is_subset}.
        Hence $\preimg{h}{C}$ is closed in $X$ with respect to $T$ by \cref{subseteq}.
    \end{subproof}
    $\dom{h} = X$ by \cref{dom_union}.
    $f\in\rels{A}{Y}$ and $g\in\rels{B}{Y}$ and $\ran{h} = \ran{f}\union\ran{g}$
        by \cref{funs_subseteq_rels,subseteq,ran_union}.
    Hence $\ran{h}\subseteq Y$ by \cref{rels_ran_subseteq,union_subsets_is_subset}.
    Therefore $h\in\funs{X}{Y}$ by \cref{function_apply_elim,ran_iff,subseteq,funs_intro}.
    Hence $h$ is continuous from $X$ to $Y$ with respect to $T$ and $T_1$
        by \cref{continuous,closed_in,open_in,preimg_closed_implies_continuous}.
\end{proof}

% from §18 helper lemma 5: first projection is continuous
\begin{lemma}\label{projection_one_continuous}
    Assume $T$ is a topology on $X$ and $T_1$ is a topology on $Y$.
    Let $T_2 = \productTopology{T}{T_1}{X}{Y}$.
    Then $\piOne{X}{Y}$ is continuous from $X\times Y$ to $X$ with respect to $T_2$ and $T$.
\end{lemma}
\begin{proof}
    $T_2$ is a topology on $X\times Y$ by \cref{basis_topology_is_topology,product_basis_is_basis,product_topology}.
    For all $p$ we have for all $x$ we have for all $x'$ such that
        $p\mathrel{\piOne{X}{Y}} x$ and $p\mathrel{\piOne{X}{Y}} x'$ we have $x = x'$
        by \cref{projection_first,pair_eq_iff}.
    Therefore $\piOne{X}{Y}$ is a function
        by \cref{projection_first,pair_eq_iff,rightunique,relation}.
    For all $p\in\dom{\piOne{X}{Y}}$ we have $\apply{\piOne{X}{Y}}{p}\in X$ by
        \cref{dom_iff,projection_first,pair_eq_iff,function_apply_intro}.
    For all $p$ such that $p\in X\times Y$ we have $p\in\dom{\piOne{X}{Y}}$ by
        \cref{times_elem_is_tuple,projection_first,dom_intro}.
    Hence $X\times Y\subseteq \dom{\piOne{X}{Y}}$ by \cref{subseteq}.
    For all $p$ such that $p\in\dom{\piOne{X}{Y}}$ we have $p\in X\times Y$ by
        \cref{dom_iff,projection_first,pair_eq_iff,times}.
    Therefore $\dom{\piOne{X}{Y}}\subseteq X\times Y$ by \cref{subseteq}.
    Hence $\dom{\piOne{X}{Y}} = X\times Y$ by \cref{subseteq_antisymmetric}.
    Therefore $\piOne{X}{Y}\in\funs{X\times Y}{X}$ by \cref{funs_intro}.
    For all $U$ such that $U$ is open in $X$ with respect to $T$ we have
        $\preimg{\piOne{X}{Y}}{U}$ is open in $X\times Y$ with respect to $T_2$ by
        \cref{open_in,topology_on,pow_iff,subseteq,times_subseteq_left,product_basis,product_basis_is_basis,basis_elem_open_in_generated,product_topology,preimg_pi1}.
    Therefore $\piOne{X}{Y}$ is continuous from $X\times Y$ to $X$ with respect to $T_2$ and $T$ by \cref{continuous}.
\end{proof}

% from §18 helper lemma 6: second projection is continuous
\begin{lemma}\label{projection_two_continuous}
    Assume $T$ is a topology on $X$ and $T_1$ is a topology on $Y$.
    Let $T_2 = \productTopology{T}{T_1}{X}{Y}$.
    Then $\piTwo{X}{Y}$ is continuous from $X\times Y$ to $Y$ with respect to $T_2$ and $T_1$.
\end{lemma}
\begin{proof}
    $T_2$ is a topology on $X\times Y$ by \cref{basis_topology_is_topology,product_basis_is_basis,product_topology}.
    For all $p$ we have for all $y$ we have for all $y'$ such that
        $p\mathrel{\piTwo{X}{Y}} y$ and $p\mathrel{\piTwo{X}{Y}} y'$ we have $y = y'$
        by \cref{projection_second,pair_eq_iff}.
    Therefore $\piTwo{X}{Y}$ is a function
        by \cref{projection_second,pair_eq_iff,rightunique,relation}.
    For all $p\in\dom{\piTwo{X}{Y}}$ we have $\apply{\piTwo{X}{Y}}{p}\in Y$ by
        \cref{dom_iff,projection_second,pair_eq_iff,function_apply_intro}.
    For all $p$ such that $p\in X\times Y$ we have $p\in\dom{\piTwo{X}{Y}}$ by
        \cref{times_elem_is_tuple,projection_second,dom_intro}.
    Hence $X\times Y\subseteq \dom{\piTwo{X}{Y}}$ by \cref{subseteq}.
    For all $p$ such that $p\in\dom{\piTwo{X}{Y}}$ we have $p\in X\times Y$ by
        \cref{dom_iff,projection_second,pair_eq_iff,times}.
    Therefore $\dom{\piTwo{X}{Y}}\subseteq X\times Y$ by \cref{subseteq}.
    Hence $\dom{\piTwo{X}{Y}} = X\times Y$ by \cref{subseteq_antisymmetric}.
    Therefore $\piTwo{X}{Y}\in\funs{X\times Y}{Y}$ by \cref{funs_intro}.
    For all $V$ such that $V$ is open in $Y$ with respect to $T_1$ we have
        $\preimg{\piTwo{X}{Y}}{V}$ is open in $X\times Y$ with respect to $T_2$ by
        \cref{open_in,topology_on,pow_iff,subseteq,times_subseteq_right,product_basis,product_basis_is_basis,product_topology,basis_elem_open_in_generated,preimg_pi2}.
    Therefore $\piTwo{X}{Y}$ is continuous from $X\times Y$ to $Y$ with respect to $T_2$ and $T_1$ by \cref{continuous}.
\end{proof}

% from §18 Item 19: maps into products (forward)
\begin{lemma}\label{continuous_map_into_product_forward}
    Assume $T_1$ is a topology on $X$ and $T_2$ is a topology on $Y$.
    Assume $f$ is continuous from $A$ to $X\times Y$ with respect to $T$ and
        $\productTopology{T_1}{T_2}{X}{Y}$.
    Let $f_1 = \piOne{X}{Y}\circ f$.
    Let $f_2 = \piTwo{X}{Y}\circ f$.
    Then $f_1$ is continuous from $A$ to $X$ with respect to $T$ and $T_1$ and
        $f_2$ is continuous from $A$ to $Y$ with respect to $T$ and $T_2$.
\end{lemma}
\begin{proof}
    Follows by \cref{projection_one_continuous,projection_two_continuous,continuous_composite}.
\end{proof}

% from §18 Item 20: maps into products (backward)
\begin{lemma}\label{continuous_map_into_product_backward}
    Assume $f_1$ is continuous from $A$ to $X$ with respect to $T$ and $T_1$.
    Assume $f_2$ is continuous from $A$ to $Y$ with respect to $T$ and $T_2$.
    Let $f = \{(a,(f_1(a),f_2(a))) \mid a\in A\}$.
    Then $f$ is continuous from $A$ to $X\times Y$ with respect to $T$ and
        $\productTopology{T_1}{T_2}{X}{Y}$.
\end{lemma}
\begin{proof}
    Let $T_3 = \productTopology{T_1}{T_2}{X}{Y}$.
    Show that for all $U$ such that $U$ is open in $X\times Y$ with respect to $T_3$
        we have $\preimg{f}{U}$ is open in $A$ with respect to $T$.
    \begin{subproof}
        Fix $U$ such that $U$ is open in $X\times Y$ with respect to $T_3$.
        Take $M\subseteq \productBasis{T_1}{T_2}{X}{Y}$ such that $U = \unions{M}$ by
            \cref{open_in,continuous,basis_for_topology,product_basis_is_basis,product_topology,basis_topology_unions}.
        Let $N = \{\preimg{f}{B} \mid B\in M\}$.
        Show that for all $C$ such that $C\in N$ we have $C\in T$.
        \begin{subproof}
            Fix $C$ such that $C\in N$.
            Take $B\in M$ such that $C = \preimg{f}{B}$.
            Take $U_1,V_1$ such that $U_1\in T_1$ and $V_1\in T_2$ and $B = U_1\times V_1$
                by \cref{subseteq,product_basis}.
            For all $a$ such that $a\in\preimg{f}{B}$ we have
                $a\in\preimg{f_1}{U_1}\inter \preimg{f_2}{V_1}$ by
                \cref{preimg_iff,times_elem_is_tuple,pair_eq_iff,continuous,funs_elim,function_apply_elim,inter_intro}.
            Hence $\preimg{f}{B}\subseteq \preimg{f_1}{U_1}\inter \preimg{f_2}{V_1}$ by \cref{subseteq}.
            For all $a$ such that $a\in\preimg{f_1}{U_1}\inter \preimg{f_2}{V_1}$ we have
                $a\in\preimg{f}{B}$ by
                \cref{inter,preimg_iff,continuous,funs_elim,function_apply_intro,dom_intro,times_tuple_intro,pair_eq_iff}.
            Therefore $\preimg{f_1}{U_1}\inter \preimg{f_2}{V_1}\subseteq \preimg{f}{B}$ by \cref{subseteq}.
            Hence $C\in T$ by \cref{continuous,open_in,topology_on,subseteq_antisymmetric}.
        \end{subproof}
        $T$ is a topology on $A$ and $\unions{N}\in T$ by \cref{continuous,topology_on,subseteq}.
        Show that for all $a$ such that $a\in\preimg{f}{U}$ we have $a\in\unions{N}$.
        \begin{subproof}
            Fix $a$ such that $a\in\preimg{f}{U}$.
            Take $b\in U$ such that $(a,b)\in f$ by \cref{preimg_iff}.
            Take $B\in M$ such that $b\in B$ by \cref{unions_iff}.
            Hence $a\in\preimg{f}{B}$ by \cref{preimg_iff}.
            We have $\preimg{f}{B}\in N$.
            Therefore $a\in\unions{N}$ by \cref{unions_iff}.
        \end{subproof}
        Hence $\preimg{f}{U}\subseteq \unions{N}$ by \cref{subseteq}.
        Show that for all $a$ such that $a\in\unions{N}$ we have $a\in\preimg{f}{U}$.
        \begin{subproof}
            Fix $a$ such that $a\in\unions{N}$.
            Take $C\in N$ such that $a\in C$ by \cref{unions_iff}.
            Take $B\in M$ such that $C = \preimg{f}{B}$.
            Take $b\in B$ such that $(a,b)\in f$ by \cref{preimg_iff}.
            Hence $b\in U$ by \cref{unions_iff}.
            Therefore $a\in\preimg{f}{U}$ by \cref{preimg_iff}.
        \end{subproof}
        Hence $\unions{N}\subseteq \preimg{f}{U}$ by \cref{subseteq}.
        Therefore $\preimg{f}{U} = \unions{N}$ by \cref{subseteq_antisymmetric}.
        Hence $\preimg{f}{U}$ is open in $A$ with respect to $T$ by \cref{open_in,subseteq}.
    \end{subproof}
    For all $x,y,z$ such that $(x,y)\in f$ and $(x,z)\in f$ we have $y = z$ by \cref{pair_eq_iff}.
    Hence $f$ is a function by \cref{relation,pair_eq_iff,rightunique}.
    We have $\dom{f} = A$ by \cref{dom_iff,pair_eq_iff,dom_intro,subseteq,subseteq_antisymmetric}.
    Therefore $f\in\funs{A}{X\times Y}$ by
        \cref{continuous,funs_elim,times_tuple_intro,function_apply_intro,funs_intro}.
    Hence $f$ is continuous from $A$ to $X\times Y$ with respect to $T$ and $T_3$ by
        \cref{continuous,basis_topology_is_topology,product_basis_is_basis,product_topology}.
\end{proof}
\section{§19 The Product Topology}

% from §19 Item 1: J-tuples
\begin{definition}\label{j_tuple}
    $f$ is a $J$-tuple of elements of $X$ iff $f\in\funs{J}{X}$.
\end{definition}

% from §19 Item 2: the set of J-tuples
\begin{definition}\label{tuple_power}
    $\tuplePower{X}{J} = \funs{J}{X}$.
\end{definition}

% from §19 Item 3: product of an indexed family
\begin{definition}\label{indexed_product}
    Let $A$ be a function.
    $\productFamily{A} = \{f\in\funs{\dom{A}}{\unions{\ran{A}}} \mid \forall \alpha\in\dom{A}. \apply{f}{\alpha}\in\apply{A}{\alpha}\}$.
\end{definition}

% from §19 Item 4: box basis
\begin{definition}\label{box_basis}
    $\boxBasis{T}{X} = \{B\in\pow{\productFamily{X}} \mid \exists U.
        U\in\funs{\dom{X}}{\pow{\unions{\ran{X}}}}\land
        (\forall \alpha\in\dom{X}. \apply{U}{\alpha}\in\apply{T}{\alpha})\land
        B = \productFamily{U}\}$.
\end{definition}

% from §19 Item 5: box topology
\begin{definition}\label{box_topology}
    $\boxTopology{T}{X} = \basisTopology{\boxBasis{T}{X}}{\productFamily{X}}$.
\end{definition}

% from §19 Item 6: projection mapping
\begin{definition}\label{projection_map_indexed}
    $\piIndex{X}{\beta} = \{(x,\apply{x}{\beta}) \mid x\in\productFamily{X}\}$.
\end{definition}

% from §19 Item 7: product subbasis for a fixed index
\begin{definition}\label{product_subbasis_indexed}
    $\productSubbasisAt{T}{X}{\beta} = \{A\in\pow{\productFamily{X}} \mid
        \exists U\in\apply{T}{\beta}. A = \preimg{\piIndex{X}{\beta}}{U}\}$.
\end{definition}

% from §19 Item 8: product subbasis
\begin{definition}\label{product_subbasis_union_indexed}
    $\productSubbasis{T}{X} = \{A\in\pow{\productFamily{X}} \mid
        \exists \beta\in\dom{X}. A\in\productSubbasisAt{T}{X}{\beta}\}$.
\end{definition}

% from §19 Item 9: product topology on indexed products
\begin{definition}\label{product_topology_indexed}
    $\productTopologyIndexed{T}{X} = \subbasisTopology{\productSubbasis{T}{X}}{\productFamily{X}}$.
\end{definition}

% from §19 Item 10: box basis is a basis on the indexed product
\begin{lemma}\label{box_basis_is_basis}
    Assume $X$ is a function.
    Assume $T$ is a function.
    Assume $\dom{T} = \dom{X}$.
    Assume that for all $\alpha\in\dom{X}$ we have $\apply{T}{\alpha}$ is a topology on $\apply{X}{\alpha}$.
    Then $\boxBasis{T}{X}$ is a basis on $\productFamily{X}$.
\end{lemma}
    \begin{proof}
        Let $B = \boxBasis{T}{X}$.
        Show that $B\subseteq\pow{\productFamily{X}}$.
        \begin{subproof}
            Follows by \cref{box_basis,subseteq}.
        \end{subproof}
        We have $X\in\funs{\dom{X}}{\pow{\unions{\ran{X}}}}$ by
            \cref{funs_intro,function_apply_elim,ran_intro,subseteq_pow_unions,subseteq}.
        For all $\alpha\in\dom{X}$ we have $\apply{X}{\alpha}\in\apply{T}{\alpha}$ by \cref{topology_on}.
        Show that for all $x$ such that $x\in\productFamily{X}$ there exists $B_1\in B$ such that $x\in B_1$.
        \begin{subproof}
            Follows by \cref{subseteq_refl,powerset_intro,box_basis}.
        \end{subproof}
        Show that for all $B_1,B_2,x$ such that $B_1\in B$ and $B_2\in B$ and $x\in B_1\inter B_2$
            there exists $B_3\in B$ such that $x\in B_3$ and $B_3\subseteq B_1\inter B_2$.
    \begin{subproof}
        Fix $B_1$.
        Fix $B_2$.
        Fix $x$ such that $B_1\in B$ and $B_2\in B$ and $x\in B_1\inter B_2$.
        Take $U_1$ such that
            $U_1\in\funs{\dom{X}}{\pow{\unions{\ran{X}}}}$ and
            $B_1 = \productFamily{U_1}$ and
            for all $\alpha\in\dom{X}$ we have $\apply{U_1}{\alpha}\in\apply{T}{\alpha}$
            by \cref{box_basis}.
        $\dom{U_1} = \dom{X}$ by \cref{funs_elim}.
        Take $U_2$ such that
            $U_2\in\funs{\dom{X}}{\pow{\unions{\ran{X}}}}$ and
            $B_2 = \productFamily{U_2}$ and
            for all $\alpha\in\dom{X}$ we have $\apply{U_2}{\alpha}\in\apply{T}{\alpha}$
            by \cref{box_basis}.
        $\dom{U_2} = \dom{X}$ by \cref{funs_elim}.
        Let $R = \{(\alpha,\apply{U_1}{\alpha}\inter\apply{U_2}{\alpha}) \mid \alpha\in\dom{X}\}$.
        We have $R$ is a relation by \cref{times_elem_is_tuple,relation}.
        Let $U_3 = R$.
        Show that for all $a$ we have for all $y$ we have for all $z$ if $(a,y)\in U_3$ and $(a,z)\in U_3$ then $y = z$.
        \begin{subproof}
            Fix $a$.
            Fix $y$.
            Fix $z$ such that $(a,y)\in U_3$ and $(a,z)\in U_3$.
            We have $(a,y)\in R$ and $(a,z)\in R$.
            Take $\alpha_1\in\dom{X}$ such that $(a,y) = (\alpha_1,\apply{U_1}{\alpha_1}\inter\apply{U_2}{\alpha_1})$
                by definition.
            Take $\alpha_2\in\dom{X}$ such that $(a,z) = (\alpha_2,\apply{U_1}{\alpha_2}\inter\apply{U_2}{\alpha_2})$
                by definition.
            Hence $a = \alpha_1$ and $y = \apply{U_1}{\alpha_1}\inter\apply{U_2}{\alpha_1}$ and
                $a = \alpha_2$ and $z = \apply{U_1}{\alpha_2}\inter\apply{U_2}{\alpha_2}$ by \cref{pair_eq_iff}.
            Therefore $y = \apply{U_1}{a}\inter\apply{U_2}{a}$.
            Hence $z = \apply{U_1}{a}\inter\apply{U_2}{a}$.
            Therefore $y = z$.
        \end{subproof}
        Thus $U_3$ is a function by \cref{rightunique,relation}.
        Show that for all $\alpha$ such that $\alpha\in\dom{X}$ we have $\alpha\in\dom{U_3}$.
        \begin{subproof}
            Fix $\alpha$ such that $\alpha\in\dom{X}$.
            We have $(\alpha,\apply{U_1}{\alpha}\inter\apply{U_2}{\alpha})\in U_3$.
            Hence $\alpha\in\dom{U_3}$ by \cref{dom_intro}.
        \end{subproof}
        We have $\dom{X}\subseteq\dom{U_3}$ by \cref{subseteq}.
        Show that for all $\alpha$ such that $\alpha\in\dom{U_3}$ we have $\alpha\in\dom{X}$.
        \begin{subproof}
            Fix $\alpha$ such that $\alpha\in\dom{U_3}$.
            Take $y$ such that $(\alpha,y)\in U_3$ by \cref{dom_iff}.
            We have $(\alpha,y)\in R$.
            Take $\beta\in\dom{X}$ such that $(\alpha,y) = (\beta,\apply{U_1}{\beta}\inter\apply{U_2}{\beta})$
                by definition.
            Hence $\alpha = \beta$ by \cref{pair_eq_iff}.
        \end{subproof}
        Therefore $\dom{U_3}\subseteq\dom{X}$ by \cref{subseteq}.
        Hence $\dom{U_3} = \dom{X}$ by \cref{subseteq_antisymmetric}.
        For all $\alpha\in\dom{X}$ we have $\apply{U_3}{\alpha} = \apply{U_1}{\alpha}\inter\apply{U_2}{\alpha}$
            by \cref{function_apply_intro}.
        For all $\alpha\in\dom{X}$ we have $\apply{U_3}{\alpha}\in\pow{\unions{\ran{X}}}$ by
            \cref{funs_elim,pow_iff,inter_subseteq,powerset_intro}.
        Therefore $U_3\in\funs{\dom{X}}{\pow{\unions{\ran{X}}}}$ by
            \cref{funs_intro,funs_elim,pow_iff,inter_subseteq,powerset_intro}.
        For all $\alpha\in\dom{X}$ we have $\apply{U_3}{\alpha}\in\apply{T}{\alpha}$ by \cref{topology_on}.
        Let $B_3 = \productFamily{U_3}$.
        Show that for all $\alpha\in\dom{X}$ we have $\apply{U_3}{\alpha}\subseteq\apply{X}{\alpha}$.
        \begin{subproof}
            Fix $\alpha$ such that $\alpha\in\dom{X}$.
            $\apply{U_1}{\alpha}\in\apply{T}{\alpha}$ and $\apply{U_2}{\alpha}\in\apply{T}{\alpha}$.
            Hence $\apply{U_1}{\alpha}\subseteq\apply{X}{\alpha}$ and
                $\apply{U_2}{\alpha}\subseteq\apply{X}{\alpha}$ by \cref{topology_on,subseteq,pow_iff}.
            $\apply{U_3}{\alpha} = \apply{U_1}{\alpha}\inter\apply{U_2}{\alpha}$.
            Show that $\apply{U_3}{\alpha}\subseteq\apply{U_1}{\alpha}$.
            \begin{subproof}
                For all $a$ such that $a\in\apply{U_3}{\alpha}$ we have $a\in\apply{U_1}{\alpha}$ by \cref{inter_elim_left,subseteq}.
            \end{subproof}
            Therefore $\apply{U_3}{\alpha}\subseteq\apply{X}{\alpha}$ by \cref{subseteq_transitive}.
        \end{subproof}
        Therefore for all $z$ such that $z\in B_3$ we have $z\in\productFamily{X}$ by
            \cref{indexed_product,funs_elim,subseteq,funs_intro,function_apply_elim,ran_intro,unions_iff}.
        Hence $B_3\in\pow{\productFamily{X}}$ by \cref{subseteq,powerset_intro}.
        Therefore $B_3\in B$ by \cref{box_basis}.
        We have $x\in\productFamily{U_1}$ and $x\in\productFamily{U_2}$ by
            \cref{inter_elim_left,inter_elim_right,subseteq_refl,elem_subseteq}.
        For all $\alpha\in\dom{X}$ we have $\apply{x}{\alpha}\in\apply{U_3}{\alpha}$
            by \cref{indexed_product,inter_intro}.
        Therefore $x\in B_3$ by
            \cref{indexed_product,inter_intro,funs_intro,funs_elim,function_apply_elim,ran_intro,unions_iff}.
        For all $z$ such that $z\in B_3$ we have $z\in B_1$ and $z\in B_2$ by
            \cref{indexed_product,inter_elim_left,funs_intro,funs_elim,function_apply_elim,ran_intro,unions_iff,inter_elim_right}.
        Therefore for all $z$ such that $z\in B_3$ we have $z\in B_1\inter B_2$ by \cref{inter_intro}.
        Hence $B_3\subseteq B_1\inter B_2$ by \cref{subseteq}.
    \end{subproof}
    Hence $B$ is a basis on $\productFamily{X}$ by \cref{basis_on}.
\end{proof}

% from §19 Item 11: product subbasis is a subbasis on the indexed product
\begin{lemma}\label{product_subbasis_is_subbasis}
    Assume $X$ is a function.
    Assume $T$ is a function.
    Assume $\dom{T} = \dom{X}$.
    Assume $\dom{X}$ is inhabited.
    Assume that for all $\alpha\in\dom{X}$ we have $\apply{T}{\alpha}$ is a topology on $\apply{X}{\alpha}$.
    Then $\productSubbasis{T}{X}$ is a subbasis on $\productFamily{X}$.
\end{lemma}
\begin{proof}
    Let $S = \productSubbasis{T}{X}$.
    Show that $S\subseteq\pow{\productFamily{X}}$.
    \begin{subproof}
        Follows by \cref{product_subbasis_union_indexed,product_subbasis_indexed,subseteq}.
    \end{subproof}
    Take $\beta$ such that $\beta\in\dom{X}$.
    Let $A = \preimg{\piIndex{X}{\beta}}{\apply{X}{\beta}}$.
    Hence $A\in S$ by
        \cref{topology_on,product_subbasis_indexed,product_subbasis_union_indexed,preimg_iff,projection_map_indexed,pair_eq_iff,subseteq,powerset_intro}.
    Show that $\unions{S} = \productFamily{X}$.
    \begin{subproof}
        Follows by \cref{indexed_product,projection_map_indexed,preimg_iff,unions_iff,subseteq,unions_subseteq_of_powerset_is_subseteq,subseteq_antisymmetric}.
    \end{subproof}
    Hence $S$ is a subbasis on $\productFamily{X}$ by \cref{subbasis_on}.
\end{proof}

% from §19 Item 12: finite intersections of product subbasis elements
\begin{lemma}\label{product_subbasis_finite_intersections_basis}
    Assume $X$ is a function.
    Assume $T$ is a function.
    Assume $\dom{T} = \dom{X}$.
    Assume $\dom{X}$ is inhabited.
    Assume that for all $\alpha\in\dom{X}$ we have $\apply{T}{\alpha}$ is a topology on $\apply{X}{\alpha}$.
    Then $\finiteIntersections{\productSubbasis{T}{X}}$ is a basis for $\productTopologyIndexed{T}{X}$ on $\productFamily{X}$.
\end{lemma}
\begin{proof}
    Follows by \cref{product_subbasis_is_subbasis,product_topology_indexed,subbasis_finite_intersections_basis,basis_for_topology,topology_generated_by_subbasis}.
\end{proof}

% from §19 helper lemma 1: choice function on a relation
\begin{axiom}\label{choice_function}
    Suppose $R$ is a relation.
    Suppose for all $x\in A$ there exists $y$ such that $x\mathrel{R} y$.
    Then there exists $f$ such that $f$ is a function on $A$ and for all $x\in A$ we have $x\mathrel{R}\apply{f}{x}$.
\end{axiom}

% from §19 Item 13: basis for the box topology from bases
\begin{lemma}\label{box_basis_from_bases}
    Assume $X$ is a function.
    Assume $T$ is a function.
    Assume $B$ is a function.
    Assume $\dom{T} = \dom{X}$.
    Assume $\dom{B} = \dom{X}$.
    Assume that for all $\alpha\in\dom{X}$ we have $\apply{B}{\alpha}$ is a basis for $\apply{T}{\alpha}$ on $\apply{X}{\alpha}$.
    Let $C = \{A\in\pow{\productFamily{X}} \mid
        \exists U. U\in\funs{\dom{X}}{\pow{\unions{\ran{X}}}}\land
        (\forall \alpha\in\dom{X}. \apply{U}{\alpha}\in\apply{B}{\alpha})\land
        A = \productFamily{U}\}$.
    Then $C$ is a basis for $\boxTopology{T}{X}$ on $\productFamily{X}$.
\end{lemma}
\begin{proof}
    Let $T_0 = \boxTopology{T}{X}$.
    Let $B_0 = \boxBasis{T}{X}$.
    For all $\alpha\in\dom{X}$ we have $\apply{T}{\alpha}$ is a topology on $\apply{X}{\alpha}$
        by \cref{basis_for_topology,basis_topology_is_topology}.
    We have $T_0$ is a topology on $\productFamily{X}$ by
        \cref{box_basis_is_basis,basis_topology_is_topology,box_topology}.
    Show that for all $A$ such that $A\in C$ we have $A$ is open in $\productFamily{X}$ with respect to $T_0$.
    \begin{subproof}
        Follows by \cref{subseteq,basis_for_topology,basis_elem_open_in_generated,box_basis,box_basis_is_basis,box_topology,open_in}.
    \end{subproof}
    Show that for all $U,x$ such that $U\in T_0$ and $x\in U$ there exists $C_1\in C$ such that $x\in C_1$ and $C_1\subseteq U$.
    \begin{subproof}
        Fix $U$.
        Fix $x$ such that $U\in T_0$ and $x\in U$.
        Take $B_1\in B_0$ such that $x\in B_1$ and $B_1\subseteq U$ by \cref{box_topology,topology_generated_by_basis}.
        Take $V$ such that
            $V\in\funs{\dom{X}}{\pow{\unions{\ran{X}}}}$ and
            $B_1 = \productFamily{V}$ and
            for all $\alpha\in\dom{X}$ we have $\apply{V}{\alpha}\in\apply{T}{\alpha}$
            by \cref{box_basis}.
        $\dom{V} = \dom{X}$ by \cref{funs_elim}.
        For all $\alpha\in\dom{X}$ there exists $W$ such that
            $W\in\apply{B}{\alpha}$ and $\apply{x}{\alpha}\in W$ and $W\subseteq\apply{V}{\alpha}$
            by \cref{indexed_product,open_in,basis_for_topology,topology_generated_by_basis}.
        Let $R = \{r\in\dom{X}\times\pow{\unions{\ran{X}}} \mid \apply{x}{\fst{r}}\in\snd{r} \land \snd{r}\subseteq\apply{V}{\fst{r}} \land \snd{r}\in\apply{B}{\fst{r}}\}$.
        We have $R$ is a relation by \cref{times_elem_is_tuple,relation}.
        Let $A = \dom{X}$.
        For all $\alpha\in A$ there exists $W$ such that
            $W\in\apply{B}{\alpha}$ and $\apply{x}{\alpha}\in W$ and $W\subseteq\apply{V}{\alpha}$ by
            \cref{indexed_product,open_in,basis_for_topology,topology_generated_by_basis}.
        Thus for all $\alpha\in A$ there exists $W$ such that $\alpha\mathrel{R} W$ by
            \cref{topology_on,subseteq,pow_iff,function_apply_elim,ran_intro,unions_iff,powerset_intro,times_tuple_intro,fst_eq,snd_eq}.
        Take $U_1$ such that $U_1$ is a function on $A$ and for all $\alpha\in A$ we have $\alpha\mathrel{R}\apply{U_1}{\alpha}$
            by \cref{choice_function}.
        We have $U_1$ is a function and $\dom{U_1} = A$.
        For all $\alpha\in\dom{X}$ we have $\apply{U_1}{\alpha}\in\apply{B}{\alpha}$ and
            $\apply{x}{\alpha}\in\apply{U_1}{\alpha}$ and $\apply{U_1}{\alpha}\subseteq\apply{V}{\alpha}$
            by \cref{fst_eq,snd_eq}.
        We have $\dom{U_1} = \dom{X}$.
        Show that for all $\alpha\in\dom{X}$ we have $\apply{U_1}{\alpha}\in\pow{\unions{\ran{X}}}$.
        \begin{subproof}
            Fix $\alpha$ such that $\alpha\in\dom{X}$.
            $\apply{U_1}{\alpha}\subseteq\apply{V}{\alpha}$.
            $\apply{V}{\alpha}\in\pow{\unions{\ran{X}}}$ by \cref{funs_elim}.
            We have $\apply{V}{\alpha}\subseteq\unions{\ran{X}}$ by \cref{pow_iff}.
            Therefore $\apply{U_1}{\alpha}\subseteq\unions{\ran{X}}$ by \cref{subseteq_transitive}.
            Hence $\apply{U_1}{\alpha}\in\pow{\unions{\ran{X}}}$ by \cref{powerset_intro}.
        \end{subproof}
        Therefore $U_1\in\funs{\dom{X}}{\pow{\unions{\ran{X}}}}$ by
            \cref{funs_intro}.
        Let $C_1 = \productFamily{U_1}$.
        Show that for all $\alpha\in\dom{X}$ we have $\apply{U_1}{\alpha}\subseteq\apply{X}{\alpha}$.
        \begin{subproof}
            Fix $\alpha$ such that $\alpha\in\dom{X}$.
            $\apply{U_1}{\alpha}\subseteq\apply{V}{\alpha}$.
            $\apply{V}{\alpha}\in\apply{T}{\alpha}$.
            We have $\apply{V}{\alpha}\subseteq\apply{X}{\alpha}$ by \cref{topology_on,subseteq,pow_iff}.
            Therefore $\apply{U_1}{\alpha}\subseteq\apply{X}{\alpha}$ by \cref{subseteq_transitive}.
        \end{subproof}
        Therefore for all $z$ such that $z\in C_1$ we have $z\in\productFamily{X}$ by
            \cref{indexed_product,funs_intro,funs_elim,function_apply_elim,ran_intro,unions_iff,subseteq}.
        We have $C_1\in\pow{\productFamily{X}}$ by \cref{subseteq,powerset_intro}.
        Therefore $C_1\in C$ by \cref{powerset_intro}.
        For all $\alpha\in\dom{X}$ we have $\apply{x}{\alpha}\in\apply{U_1}{\alpha}$.
        Hence $x\in C_1$ by
            \cref{indexed_product,funs_intro,funs_elim,function_apply_elim,ran_intro,unions_iff}.
        For all $z$ such that $z\in C_1$ we have $z\in U$ by
            \cref{indexed_product,subseteq,funs_intro,funs_elim,function_apply_elim,ran_intro,unions_iff}.
        Therefore $C_1\subseteq U$ by \cref{subseteq}.
    \end{subproof}
    Hence $C$ is a basis for $T_0$ on $\productFamily{X}$ by \cref{open_collection_basis}.
\end{proof}

% from §19 Item 14: basis for the product topology from bases
\begin{lemma}\label{product_basis_from_bases_indexed}
    Assume $X$ is a function.
    Assume $T$ is a function.
    Assume $B$ is a function.
    Assume $\dom{T} = \dom{X}$.
    Assume $\dom{B} = \dom{X}$.
    Assume $\dom{X}$ is inhabited.
    Assume that for all $\alpha\in\dom{X}$ we have $\apply{B}{\alpha}$ is a basis for $\apply{T}{\alpha}$ on $\apply{X}{\alpha}$.
    Let $S$ be a set such that for all $A$ we have
        $A\in S$ iff there exists $\beta\in\dom{X}$ such that there exists $U\in\apply{B}{\beta}$ such that
            $A = \preimg{\piIndex{X}{\beta}}{U}$.
    Then $\finiteIntersections{S}$ is a basis for $\productTopologyIndexed{T}{X}$ on $\productFamily{X}$.
\end{lemma}
\begin{proof}
    Let $T_0 = \productTopologyIndexed{T}{X}$.
    Let $T_1 = \subbasisTopology{S}{\productFamily{X}}$.
    For all $A$ such that $A\in S$ we have $A\in\pow{\productFamily{X}}$ by
        \cref{preimg_iff,projection_map_indexed,pair_eq_iff,subseteq,powerset_intro}.
    Show that $S\subseteq\pow{\productFamily{X}}$.
    \begin{subproof}
        Follows by \cref{subseteq}.
    \end{subproof}
    For all $x$ such that $x\in\productFamily{X}$ we have $x\in\unions{S}$ by
        \cref{indexed_product,basis_for_topology,basis_on,projection_map_indexed,preimg_iff,unions_iff}.
    Show that $\unions{S} = \productFamily{X}$.
    \begin{subproof}
        Follows by \cref{subseteq,indexed_product,projection_map_indexed,preimg_iff,unions_iff,unions_subseteq_of_powerset_is_subseteq,subseteq_antisymmetric}.
    \end{subproof}
    Hence $S$ is a subbasis on $\productFamily{X}$ by \cref{subbasis_on}.
    Therefore $\finiteIntersections{S}$ is a basis for $T_1$ on $\productFamily{X}$
        by \cref{subbasis_finite_intersections_basis,basis_for_topology,topology_generated_by_subbasis}.
    Show that for all $C$ such that $C\in S$ we have $C\in T_1$.
    \begin{subproof}
        Follows by \cref{subbasis_on,subseteq,pow_iff,singleton_subset_intro,upair_iff,pair_is_finite,subseteq_finite_is_finite,singleton_intro,inters_singleton,finite_intersections,basis_elem_open_in_generated,basis_for_topology}.
    \end{subproof}
    Therefore $S\subseteq T_1$ by \cref{subseteq}.
    For all $\alpha\in\dom{X}$ we have $\apply{T}{\alpha}$ is a topology on $\apply{X}{\alpha}$
        by \cref{basis_for_topology,basis_topology_is_topology}.
    Let $S_0 = \productSubbasis{T}{X}$.
    $T_0 = \subbasisTopology{S_0}{\productFamily{X}}$ by \cref{product_topology_indexed}.
    Hence $\finiteIntersections{S_0}$ is a basis for $T_0$ on $\productFamily{X}$ by
        \cref{product_subbasis_finite_intersections_basis,product_topology_indexed}.
    Therefore $T_0$ is a topology on $\productFamily{X}$ and
        $T_1$ is a topology on $\productFamily{X}$ by \cref{basis_topology_is_topology,basis_for_topology}.
    Show that for all $A$ such that $A\in S$ we have $A\in T_0$.
    \begin{subproof}
        Follows by \cref{basis_for_topology,basis_elem_open_in_generated,preimg_iff,projection_map_indexed,pair_eq_iff,subseteq,powerset_intro,product_subbasis_indexed,product_subbasis_union_indexed,product_subbasis_is_subbasis,subbasis_on,pow_iff,singleton_subset_intro,upair_iff,pair_is_finite,subseteq_finite_is_finite,singleton_intro,inters_singleton,finite_intersections}.
    \end{subproof}
    Show that for all $B$ such that $B\in\finiteIntersections{S}$ we have $B\in T_0$.
    \begin{subproof}
        Follows by \cref{finite_intersections,subseteq,finite_intersections_open_in_topology}.
    \end{subproof}
    Therefore $\finiteIntersections{S}\subseteq T_0$ by \cref{subseteq}.
    Show that for all $U$ such that $U\in T_1$ we have $U\in T_0$.
    \begin{subproof}
        Follows by \cref{basis_topology_unions,subseteq,topology_on}.
    \end{subproof}
    Hence $T_1\subseteq T_0$ by \cref{subseteq}.
    Show that for all $A$ such that $A\in S_0$ we have $A\in T_1$.
    \begin{subproof}
        Fix $A$ such that $A\in S_0$.
        Take $\beta\in\dom{X}$ such that $A\in\productSubbasisAt{T}{X}{\beta}$ by \cref{product_subbasis_union_indexed}.
        Take $U\in\apply{T}{\beta}$ such that $A = \preimg{\piIndex{X}{\beta}}{U}$ by \cref{product_subbasis_indexed}.
        Take $M\subseteq\apply{B}{\beta}$ such that $U = \unions{M}$ by \cref{basis_for_topology,basis_topology_unions}.
        Let $N = \{\preimg{\piIndex{X}{\beta}}{V} \mid V\in M\}$.
        We have $N\subseteq T_1$ by \cref{subseteq,subseteq_transitive}.
        Therefore $A\subseteq\unions{N}$ by \cref{preimg_iff,unions_iff,subseteq}.
        Hence $\unions{N}\subseteq A$ by \cref{unions_iff,subseteq,preimg_subseteq}.
        Thus $A\in T_1$ by \cref{subseteq_antisymmetric,topology_on,subseteq}.
    \end{subproof}
    Show that for all $B$ such that $B\in\finiteIntersections{S_0}$ we have $B\in T_1$.
    \begin{subproof}
        Follows by \cref{finite_intersections,subseteq,finite_intersections_open_in_topology}.
    \end{subproof}
    Therefore $\finiteIntersections{S_0}\subseteq T_1$ by \cref{subseteq}.
    Show that for all $U$ such that $U\in T_0$ we have $U\in T_1$.
    \begin{subproof}
        Follows by \cref{basis_topology_unions,subseteq,topology_on}.
    \end{subproof}
    Hence $T_0\subseteq T_1$ by \cref{subseteq}.
    Therefore $T_0 = T_1$ by \cref{subseteq_antisymmetric}.
    Hence $\finiteIntersections{S}$ is a basis for $T_0$ on $\productFamily{X}$ by \cref{basis_for_topology}.
\end{proof}
% from §19 helper lemma 2: replace one coordinate in a function
\begin{lemma}\label{replace_coordinate_function}
    Assume $X$ is a function.
    Assume $\alpha\in\dom{X}$.
    Assume $U\subseteq\unions{\ran{X}}$.
    Then there exists $V$ such that
        $V\in\funs{\dom{X}}{\pow{\unions{\ran{X}}}}$ and
        $\apply{V}{\alpha} = U$ and
        for all $\beta\in\dom{X}$ such that $\beta\neq\alpha$ we have $\apply{V}{\beta} = \apply{X}{\beta}$.
\end{lemma}
\begin{proof}
    Let $R_1 = \{(\alpha,U)\}$.
    Let $R_2 = \{(\beta,\apply{X}{\beta}) \mid \beta\in\dom{X}\setminus\{\alpha\}\}$.
    Let $R = R_1\union R_2$.
    For all $w$ such that $w\in R$ there exist $x,y$ such that $w = (x,y)$ by \cref{union_iff,singleton_elim}.
    Hence $R$ is a relation by \cref{relation}.
    For all $\beta\in\dom{X}$ there exists $y$ such that $\beta\mathrel{R} y$ by
        \cref{singleton_intro,singleton_iff,setminus_intro,union_iff}.
    Take $V$ such that $V$ is a function on $\dom{X}$ and for all $\beta\in\dom{X}$ we have $\beta\mathrel{R}\apply{V}{\beta}$
        by \cref{choice_function}.
    $V$ is a function and $\dom{V} = \dom{X}$.
    $(\alpha,\apply{V}{\alpha})\notin R_2$ by \cref{pair_eq_iff,setminus_intro,setminus_elim_right,singleton_intro}.
    Hence $\apply{V}{\alpha} = U$ by \cref{union_iff,singleton_elim,pair_eq_iff}.
    For all $\beta\in\dom{X}$ such that $\beta\neq\alpha$ we have $\apply{V}{\beta} = \apply{X}{\beta}$ by
        \cref{singleton_elim,pair_eq_iff,union_iff}.
    Show that for all $\beta\in\dom{X}$ we have $\apply{V}{\beta}\in\pow{\unions{\ran{X}}}$.
    \begin{subproof}
        Fix $\beta\in\dom{X}$.
        \begin{byCase}
        \caseOf{$\beta = \alpha$.}
            $\apply{V}{\beta} = U$.
            Hence $\apply{V}{\beta}\subseteq\unions{\ran{X}}$ by \cref{subseteq}.
            Thus $\apply{V}{\beta}\in\pow{\unions{\ran{X}}}$ by \cref{powerset_intro}.
        \caseOf{$\beta\neq\alpha$.}
            $\apply{V}{\beta} = \apply{X}{\beta}$.
            $(\beta,\apply{X}{\beta})\in X$ by \cref{function_apply_elim}.
            Hence $\apply{X}{\beta}\in\ran{X}$ by \cref{ran_intro}.
            $\apply{X}{\beta}\subseteq\unions{\ran{X}}$ by \cref{subseteq,unions_iff}.
            Therefore $\apply{V}{\beta}\in\pow{\unions{\ran{X}}}$ by \cref{powerset_intro}.
        \end{byCase}
    \end{subproof}
    Therefore $V\in\funs{\dom{X}}{\pow{\unions{\ran{X}}}}$ by \cref{funs_intro}.
\end{proof}

% from §19 helper lemma 2a: coordinatewise subset implies product subset
\begin{lemma}\label{indexed_product_subset_from_coordinate_subsets}
    Assume $X$ is a function.
    Assume $U\in\funs{\dom{X}}{\pow{\unions{\ran{X}}}}$.
    Assume that for all $\alpha\in\dom{X}$ we have $\apply{U}{\alpha}\subseteq\apply{X}{\alpha}$.
    Then $\productFamily{U}\subseteq\productFamily{X}$.
\end{lemma}
\begin{proof}
    For all $x$ such that $x\in\productFamily{U}$ we have $x\in\productFamily{X}$
        by \cref{indexed_product,funs_elim,subseteq,funs_intro,function_apply_elim,ran_intro,unions_iff}.
    Hence $\productFamily{U}\subseteq\productFamily{X}$ by \cref{subseteq}.
\end{proof}

% from §19 helper lemma 2b: replacing one coordinate preserves coordinatewise subsets
\begin{lemma}\label{replace_coordinate_function_subset}
    Assume $X$ is a function.
    Assume $U\in\funs{\dom{X}}{\pow{\unions{\ran{X}}}}$.
    Assume $\alpha\in\dom{X}$.
    Assume $\apply{U}{\alpha}\subseteq\apply{X}{\alpha}$.
    Assume that for all $\beta\in\dom{X}$ such that $\beta\neq\alpha$ we have $\apply{U}{\beta} = \apply{X}{\beta}$.
    Then for all $\beta\in\dom{X}$ we have $\apply{U}{\beta}\subseteq\apply{X}{\beta}$.
\end{lemma}
\begin{proof}
    For all $\beta\in\dom{X}$ we have $\apply{U}{\beta}\subseteq\apply{X}{\beta}$ by \cref{subseteq_refl}.
\end{proof}

% from §19 helper lemma 2c: products are disjoint if one coordinate is disjoint
\begin{lemma}\label{indexed_product_disjoint_from_coordinate_disjoint}
    Assume $X$ is a function.
    Assume $U_1\in\funs{\dom{X}}{\pow{\unions{\ran{X}}}}$.
    Assume $U_2\in\funs{\dom{X}}{\pow{\unions{\ran{X}}}}$.
    Assume $\alpha\in\dom{X}$.
    Assume $\apply{U_1}{\alpha}$ is disjoint from $\apply{U_2}{\alpha}$.
    Then $\productFamily{U_1}$ is disjoint from $\productFamily{U_2}$.
\end{lemma}
\begin{proof}
    Suppose not.
    There exists $p$ such that $p\in\productFamily{U_1}$ and $p\in\productFamily{U_2}$.
    Hence $\apply{p}{\alpha}\in\apply{U_1}{\alpha}$ and $\apply{p}{\alpha}\in\apply{U_2}{\alpha}$
        by \cref{funs_elim,indexed_product}.
    Thus $\apply{U_1}{\alpha}$ intersects $\apply{U_2}{\alpha}$ by \cref{inter_intro,emptyset,intersects}.
    Contradiction by \cref{disjoint}.
\end{proof}

% from §19 helper lemma 2d: preimages under a projection preserve disjointness
\begin{lemma}\label{projection_preimgs_disjoint}
    Assume $X$ is a function.
    Assume $\alpha\in\dom{X}$.
    Assume $U_1$ is disjoint from $U_2$.
    Then $\preimg{\piIndex{X}{\alpha}}{U_1}$ is disjoint from $\preimg{\piIndex{X}{\alpha}}{U_2}$.
\end{lemma}
\begin{proof}
    Suppose not.
    There exists $p$ such that $p\in\preimg{\piIndex{X}{\alpha}}{U_1}$ and $p\in\preimg{\piIndex{X}{\alpha}}{U_2}$.
    Hence $\apply{p}{\alpha}\in U_1$ and $\apply{p}{\alpha}\in U_2$
        by \cref{preimg_iff,projection_map_indexed,pair_eq_iff}.
    Thus $U_1$ intersects $U_2$ by \cref{inter_intro,emptyset,intersects}.
    Contradiction by \cref{disjoint}.
\end{proof}
% from §19 Item 15: products of subspaces for the box topology
\begin{lemma}\label{product_subspace_box}
    Assume $X$ is a function.
    Assume $A$ is a function.
    Assume $T$ is a function.
    Assume $\dom{A} = \dom{X}$.
    Assume $\dom{T} = \dom{X}$.
    Assume that for all $\alpha\in\dom{X}$ we have $\apply{A}{\alpha}\subseteq\apply{X}{\alpha}$.
    Assume that for all $\alpha\in\dom{X}$ we have $\apply{T}{\alpha}$ is a topology on $\apply{X}{\alpha}$.
    Then $\productFamily{A}$ is a subspace of $\productFamily{X}$ with respect to $\boxTopology{T}{X}$.
\end{lemma}
\begin{proof}
    Show that for all $\alpha\in\dom{X}$ we have $\apply{A}{\alpha}\in\pow{\unions{\ran{X}}}$.
    \begin{subproof}
        Fix $\alpha$ such that $\alpha\in\dom{X}$.
        $\apply{A}{\alpha}\subseteq\apply{X}{\alpha}$.
        $\apply{X}{\alpha}\in\ran{X}$ by \cref{function_apply_elim,ran_intro}.
        Hence $\apply{X}{\alpha}\in\pow{\unions{\ran{X}}}$ by \cref{subseteq_pow_unions,subseteq}.
        $\apply{X}{\alpha}\subseteq\unions{\ran{X}}$ by \cref{pow_iff}.
        Therefore $\apply{A}{\alpha}\subseteq\unions{\ran{X}}$ by \cref{subseteq_transitive}.
        Hence $\apply{A}{\alpha}\in\pow{\unions{\ran{X}}}$ by \cref{powerset_intro}.
    \end{subproof}
    Therefore $A\in\funs{\dom{X}}{\pow{\unions{\ran{X}}}}$ by \cref{funs_intro}.
    $\productFamily{A}\subseteq\productFamily{X}$ by \cref{indexed_product_subset_from_coordinate_subsets}.
    $\boxTopology{T}{X}$ is a topology on $\productFamily{X}$ by
        \cref{box_basis_is_basis,basis_topology_is_topology,box_topology}.
    Therefore $\productFamily{A}$ is a subspace of $\productFamily{X}$ with respect to $\boxTopology{T}{X}$ by \cref{subspace_of}.
\end{proof}

% from §19 Item 16: products of subspaces for the product topology
\begin{lemma}\label{product_subspace_product}
    Assume $X$ is a function.
    Assume $A$ is a function.
    Assume $T$ is a function.
    Assume $\dom{A} = \dom{X}$.
    Assume $\dom{T} = \dom{X}$.
    Assume $\dom{X}$ is inhabited.
    Assume that for all $\alpha\in\dom{X}$ we have $\apply{A}{\alpha}\subseteq\apply{X}{\alpha}$.
    Assume that for all $\alpha\in\dom{X}$ we have $\apply{T}{\alpha}$ is a topology on $\apply{X}{\alpha}$.
    Then $\productFamily{A}$ is a subspace of $\productFamily{X}$ with respect to $\productTopologyIndexed{T}{X}$.
\end{lemma}
\begin{proof}
    Let $T_0 = \productTopologyIndexed{T}{X}$.
    Show that for all $\alpha\in\dom{X}$ we have $\apply{A}{\alpha}\in\pow{\unions{\ran{X}}}$.
    \begin{subproof}
        Fix $\alpha$ such that $\alpha\in\dom{X}$.
        $\apply{A}{\alpha}\subseteq\apply{X}{\alpha}$.
        $\apply{X}{\alpha}\in\ran{X}$ by \cref{function_apply_elim,ran_intro}.
        Hence $\apply{X}{\alpha}\in\pow{\unions{\ran{X}}}$ by \cref{subseteq_pow_unions,subseteq}.
        $\apply{X}{\alpha}\subseteq\unions{\ran{X}}$ by \cref{pow_iff}.
        Therefore $\apply{A}{\alpha}\subseteq\unions{\ran{X}}$ by \cref{subseteq_transitive}.
        Hence $\apply{A}{\alpha}\in\pow{\unions{\ran{X}}}$ by \cref{powerset_intro}.
    \end{subproof}
    Therefore $A\in\funs{\dom{X}}{\pow{\unions{\ran{X}}}}$ by \cref{funs_intro}.
    $\productFamily{A}\subseteq\productFamily{X}$ by \cref{indexed_product_subset_from_coordinate_subsets}.
    $T_0$ is a topology on $\productFamily{X}$ by
        \cref{product_subbasis_finite_intersections_basis,basis_for_topology,basis_topology_is_topology}.
    Therefore $\productFamily{A}$ is a subspace of $\productFamily{X}$ with respect to $T_0$ by \cref{subspace_of}.
\end{proof}

% from §19 Item 17: product of Hausdorff spaces with the box topology
\begin{lemma}\label{product_hausdorff_box}
    Assume $X$ is a function.
    Assume $T$ is a function.
    Assume $\dom{T} = \dom{X}$.
    Assume that for all $\alpha\in\dom{X}$ we have $\apply{X}{\alpha}$ is Hausdorff with respect to $\apply{T}{\alpha}$.
    Then $\productFamily{X}$ is Hausdorff with respect to $\boxTopology{T}{X}$.
\end{lemma}
\begin{proof}
    Let $T_0 = \boxTopology{T}{X}$.
    Let $B_0 = \boxBasis{T}{X}$.
    For all $\alpha\in\dom{X}$ we have $\apply{T}{\alpha}$ is a topology on $\apply{X}{\alpha}$ by \cref{hausdorff}.
    $\boxBasis{T}{X}$ is a basis on $\productFamily{X}$ by \cref{box_basis_is_basis}.
    Hence $T_0$ is a topology on $\productFamily{X}$ by \cref{basis_topology_is_topology,box_topology}.
    Show that for all $x_1,x_2$ such that $x_1\in\productFamily{X}$ and $x_2\in\productFamily{X}$ and $x_1\neq x_2$
        there exist $U_1,U_2$ such that
            $U_1$ is a neighborhood of $x_1$ in $\productFamily{X}$ with respect to $T_0$ and
            $U_2$ is a neighborhood of $x_2$ in $\productFamily{X}$ with respect to $T_0$ and
            $U_1$ is disjoint from $U_2$.
    \begin{subproof}
        Fix $x_1$.
        Fix $x_2$ such that $x_1\in\productFamily{X}$ and $x_2\in\productFamily{X}$ and $x_1\neq x_2$.
        Take $w$ such that either $w\in x_1$ and $w\notin x_2$ or $w\notin x_1$ and $w\in x_2$
            by \cref{neq_witness}.
        \begin{byCase}
        \caseOf{$w\in x_1$ and $w\notin x_2$.}
            We have $x_1$ is a function and $\dom{x_1} = \dom{X}$ and
                $x_2$ is a function and $\dom{x_2} = \dom{X}$ by \cref{indexed_product,funs_elim}.
            Take $\alpha\in\dom{X}$ such that $w = (\alpha,\apply{x_1}{\alpha})$ by \cref{funs_elim,function_member_elim}.
            Therefore $\apply{x_1}{\alpha}\neq\apply{x_2}{\alpha}$ by \cref{function_apply_elim,pair_eq_iff}.
            $\apply{x_1}{\alpha}\in\apply{X}{\alpha}$ and
                $\apply{x_2}{\alpha}\in\apply{X}{\alpha}$ by \cref{indexed_product}.
            Take $U_1,U_2$ such that
                $U_1$ is a neighborhood of $\apply{x_1}{\alpha}$ in $\apply{X}{\alpha}$ with respect to $\apply{T}{\alpha}$ and
                $U_2$ is a neighborhood of $\apply{x_2}{\alpha}$ in $\apply{X}{\alpha}$ with respect to $\apply{T}{\alpha}$ and
                $U_1$ is disjoint from $U_2$ by \cref{hausdorff}.
            Take $V_1$ such that
                $V_1\in\funs{\dom{X}}{\pow{\unions{\ran{X}}}}$ and
                $\apply{V_1}{\alpha} = U_1$ and
                for all $\beta\in\dom{X}$ such that $\beta\neq\alpha$ we have $\apply{V_1}{\beta} = \apply{X}{\beta}$
                by \cref{neighborhood,open_in,hausdorff,topology_on,subseteq,pow_iff,function_apply_elim,ran_intro,unions_iff,subseteq_transitive,replace_coordinate_function}.
            Take $V_2$ such that
                $V_2\in\funs{\dom{X}}{\pow{\unions{\ran{X}}}}$ and
                $\apply{V_2}{\alpha} = U_2$ and
                for all $\beta\in\dom{X}$ such that $\beta\neq\alpha$ we have $\apply{V_2}{\beta} = \apply{X}{\beta}$
                by \cref{neighborhood,open_in,hausdorff,topology_on,subseteq,pow_iff,function_apply_elim,ran_intro,unions_iff,subseteq_transitive,replace_coordinate_function}.
            Let $W_1 = \productFamily{V_1}$.
            Let $W_2 = \productFamily{V_2}$.
            For all $\beta\in\dom{X}$ we have $\apply{V_1}{\beta}\subseteq\apply{X}{\beta}$
                by \cref{neighborhood,open_in,hausdorff,topology_on,subseteq,pow_iff,replace_coordinate_function_subset}.
            Hence $W_1\subseteq\productFamily{X}$ by \cref{indexed_product_subset_from_coordinate_subsets}.
            For all $\beta\in\dom{X}$ we have $\apply{V_2}{\beta}\subseteq\apply{X}{\beta}$
                by \cref{neighborhood,open_in,hausdorff,topology_on,subseteq,pow_iff,replace_coordinate_function_subset}.
            Hence $W_2\subseteq\productFamily{X}$ by \cref{indexed_product_subset_from_coordinate_subsets}.
            Show that $W_1$ is open in $\productFamily{X}$ with respect to $T_0$ and
                $W_2$ is open in $\productFamily{X}$ with respect to $T_0$.
            \begin{subproof}
                Follows by \cref{neighborhood,open_in,topology_on,subseteq,pow_iff,box_basis,basis_elem_open_in_generated,box_topology,open_in}.
            \end{subproof}
            For all $\beta\in\dom{X}$ we have $\apply{x_1}{\beta}\in\apply{V_1}{\beta}$
                by \cref{neighborhood,indexed_product,subseteq}.
            Hence $x_1$ is a function and $\dom{x_1} = \dom{X}$ and
                $\dom{V_1} = \dom{X}$ by \cref{funs_elim,indexed_product}.
            For all $\beta\in\dom{X}$ we have $\apply{x_2}{\beta}\in\apply{V_2}{\beta}$
                by \cref{neighborhood,indexed_product,subseteq}.
            Hence $x_2$ is a function and $\dom{x_2} = \dom{X}$ and
                $\dom{V_2} = \dom{X}$ by \cref{funs_elim,indexed_product}.
            Show that $W_1$ is a neighborhood of $x_1$ in $\productFamily{X}$ with respect to $T_0$ and
                $W_2$ is a neighborhood of $x_2$ in $\productFamily{X}$ with respect to $T_0$.
            \begin{subproof}
                Follows by \cref{funs_intro,indexed_product,funs_elim,function_apply_elim,ran_intro,unions_iff,neighborhood}.
            \end{subproof}
            Hence $W_1$ is disjoint from $W_2$ by \cref{indexed_product_disjoint_from_coordinate_disjoint}.
        \caseOf{$w\notin x_1$ and $w\in x_2$.}
            We have $x_1$ is a function and $\dom{x_1} = \dom{X}$ and
                $x_2$ is a function and $\dom{x_2} = \dom{X}$ by \cref{indexed_product,funs_elim}.
            Take $\alpha\in\dom{X}$ such that $w = (\alpha,\apply{x_2}{\alpha})$ by \cref{funs_elim,function_member_elim}.
            Therefore $\apply{x_1}{\alpha}\neq\apply{x_2}{\alpha}$ by \cref{function_apply_elim,pair_eq_iff}.
            $\apply{x_1}{\alpha}\in\apply{X}{\alpha}$ and
                $\apply{x_2}{\alpha}\in\apply{X}{\alpha}$ by \cref{indexed_product}.
            Take $U_1,U_2$ such that
                $U_1$ is a neighborhood of $\apply{x_1}{\alpha}$ in $\apply{X}{\alpha}$ with respect to $\apply{T}{\alpha}$ and
                $U_2$ is a neighborhood of $\apply{x_2}{\alpha}$ in $\apply{X}{\alpha}$ with respect to $\apply{T}{\alpha}$ and
                $U_1$ is disjoint from $U_2$ by \cref{hausdorff}.
            Take $V_1$ such that
                $V_1\in\funs{\dom{X}}{\pow{\unions{\ran{X}}}}$ and
                $\apply{V_1}{\alpha} = U_1$ and
                for all $\beta\in\dom{X}$ such that $\beta\neq\alpha$ we have $\apply{V_1}{\beta} = \apply{X}{\beta}$
                by \cref{neighborhood,open_in,hausdorff,topology_on,subseteq,pow_iff,function_apply_elim,ran_intro,unions_iff,subseteq_transitive,replace_coordinate_function}.
            Take $V_2$ such that
                $V_2\in\funs{\dom{X}}{\pow{\unions{\ran{X}}}}$ and
                $\apply{V_2}{\alpha} = U_2$ and
                for all $\beta\in\dom{X}$ such that $\beta\neq\alpha$ we have $\apply{V_2}{\beta} = \apply{X}{\beta}$
                by \cref{neighborhood,open_in,hausdorff,topology_on,subseteq,pow_iff,function_apply_elim,ran_intro,unions_iff,subseteq_transitive,replace_coordinate_function}.
            Let $W_1 = \productFamily{V_1}$.
            Let $W_2 = \productFamily{V_2}$.
            For all $\beta\in\dom{X}$ we have $\apply{V_1}{\beta}\subseteq\apply{X}{\beta}$
                by \cref{neighborhood,open_in,hausdorff,topology_on,subseteq,pow_iff,replace_coordinate_function_subset}.
            Hence $W_1\subseteq\productFamily{X}$ by \cref{indexed_product_subset_from_coordinate_subsets}.
            For all $\beta\in\dom{X}$ we have $\apply{V_2}{\beta}\subseteq\apply{X}{\beta}$
                by \cref{neighborhood,open_in,hausdorff,topology_on,subseteq,pow_iff,replace_coordinate_function_subset}.
            Hence $W_2\subseteq\productFamily{X}$ by \cref{indexed_product_subset_from_coordinate_subsets}.
            Show that $W_1$ is open in $\productFamily{X}$ with respect to $T_0$ and
                $W_2$ is open in $\productFamily{X}$ with respect to $T_0$.
            \begin{subproof}
                Follows by \cref{neighborhood,open_in,topology_on,subseteq,pow_iff,box_basis,basis_elem_open_in_generated,box_topology,open_in}.
            \end{subproof}
            For all $\beta\in\dom{X}$ we have $\apply{x_1}{\beta}\in\apply{V_1}{\beta}$
                by \cref{neighborhood,indexed_product,subseteq}.
            Hence $x_1$ is a function and $\dom{x_1} = \dom{X}$ and
                $\dom{V_1} = \dom{X}$ by \cref{funs_elim,indexed_product}.
            For all $\beta\in\dom{X}$ we have $\apply{x_2}{\beta}\in\apply{V_2}{\beta}$
                by \cref{neighborhood,indexed_product,subseteq}.
            Hence $x_2$ is a function and $\dom{x_2} = \dom{X}$ and
                $\dom{V_2} = \dom{X}$ by \cref{funs_elim,indexed_product}.
            Show that $W_1$ is a neighborhood of $x_1$ in $\productFamily{X}$ with respect to $T_0$ and
                $W_2$ is a neighborhood of $x_2$ in $\productFamily{X}$ with respect to $T_0$.
            \begin{subproof}
                Follows by \cref{funs_intro,indexed_product,funs_elim,function_apply_elim,ran_intro,unions_iff,neighborhood}.
            \end{subproof}
            Hence $W_1$ is disjoint from $W_2$ by \cref{indexed_product_disjoint_from_coordinate_disjoint}.
        \end{byCase}
    \end{subproof}
    Hence $\productFamily{X}$ is Hausdorff with respect to $T_0$ by \cref{hausdorff}.
\end{proof}

% from §19 Item 18: product of Hausdorff spaces with the product topology
\begin{lemma}\label{product_hausdorff_product}
    Assume $X$ is a function.
    Assume $T$ is a function.
    Assume $\dom{T} = \dom{X}$.
    Assume $\dom{X}$ is inhabited.
    Assume that for all $\alpha\in\dom{X}$ we have $\apply{X}{\alpha}$ is Hausdorff with respect to $\apply{T}{\alpha}$.
    Then $\productFamily{X}$ is Hausdorff with respect to $\productTopologyIndexed{T}{X}$.
\end{lemma}
\begin{proof}
    Let $T_0 = \productTopologyIndexed{T}{X}$.
    Let $S = \productSubbasis{T}{X}$.
    For all $\alpha\in\dom{X}$ we have $\apply{T}{\alpha}$ is a topology on $\apply{X}{\alpha}$ by \cref{hausdorff}.
    $S$ is a subbasis on $\productFamily{X}$ by \cref{product_subbasis_is_subbasis}.
    Hence $\finiteIntersections{S}$ is a basis for $T_0$ on $\productFamily{X}$ by \cref{product_subbasis_finite_intersections_basis}.
    Therefore $T_0$ is a topology on $\productFamily{X}$ by \cref{basis_topology_is_topology,basis_for_topology}.
    Show that for all $x_1,x_2$ such that $x_1\in\productFamily{X}$ and $x_2\in\productFamily{X}$ and $x_1\neq x_2$
        there exist $U_1,U_2$ such that
            $U_1$ is a neighborhood of $x_1$ in $\productFamily{X}$ with respect to $T_0$ and
            $U_2$ is a neighborhood of $x_2$ in $\productFamily{X}$ with respect to $T_0$ and
            $U_1$ is disjoint from $U_2$.
    \begin{subproof}
        Fix $x_1$.
        Fix $x_2$ such that $x_1\in\productFamily{X}$ and $x_2\in\productFamily{X}$ and $x_1\neq x_2$.
        Take $w$ such that either $w\in x_1$ and $w\notin x_2$ or $w\notin x_1$ and $w\in x_2$
            by \cref{neq_witness}.
        \begin{byCase}
        \caseOf{$w\in x_1$ and $w\notin x_2$.}
            We have $x_1$ is a function and $\dom{x_1} = \dom{X}$ and
                $x_2$ is a function and $\dom{x_2} = \dom{X}$ by \cref{indexed_product,funs_elim}.
            Take $\alpha\in\dom{X}$ such that $w = (\alpha,\apply{x_1}{\alpha})$ by \cref{funs_elim,function_member_elim}.
            Therefore $\apply{x_1}{\alpha}\neq\apply{x_2}{\alpha}$ by \cref{function_apply_elim,pair_eq_iff}.
            $\apply{x_1}{\alpha}\in\apply{X}{\alpha}$ and
                $\apply{x_2}{\alpha}\in\apply{X}{\alpha}$ by \cref{indexed_product}.
            Take $U_1,U_2$ such that
                $U_1$ is a neighborhood of $\apply{x_1}{\alpha}$ in $\apply{X}{\alpha}$ with respect to $\apply{T}{\alpha}$ and
                $U_2$ is a neighborhood of $\apply{x_2}{\alpha}$ in $\apply{X}{\alpha}$ with respect to $\apply{T}{\alpha}$ and
                $U_1$ is disjoint from $U_2$ by \cref{hausdorff}.
            Let $W_1 = \preimg{\piIndex{X}{\alpha}}{U_1}$.
            Let $W_2 = \preimg{\piIndex{X}{\alpha}}{U_2}$.
            For all $x$ such that $x\in\dom{\piIndex{X}{\alpha}}$ we have $x\in\productFamily{X}$
                by \cref{dom,projection_map_indexed,pair_eq_iff}.
            Hence $\dom{\piIndex{X}{\alpha}}\subseteq\productFamily{X}$ by \cref{subseteq}.
            For all $x$ such that $x\in\productFamily{X}$ we have $x\in\dom{\piIndex{X}{\alpha}}$ by \cref{projection_map_indexed,dom}.
            $W_1\in\pow{\productFamily{X}}$ and $W_2\in\pow{\productFamily{X}}$
                by \cref{subseteq,subseteq_antisymmetric,preimg_subseteq_dom,subseteq_transitive,powerset_intro}.
            Hence $W_1\in S$ and $W_2\in S$ by \cref{neighborhood,open_in,product_subbasis_indexed,product_subbasis_union_indexed}.
            Show that $W_1\in\finiteIntersections{S}$.
            \begin{subproof}
                Follows by \cref{subbasis_on,subseteq,pow_iff,powerset_intro,singleton_subset_intro,pair_is_finite,upair_intro_left,subseteq_finite_is_finite,singleton_intro,inters_singleton,finite_intersections}.
            \end{subproof}
            Show that $W_2\in\finiteIntersections{S}$.
            \begin{subproof}
                Follows by \cref{subbasis_on,subseteq,pow_iff,powerset_intro,singleton_subset_intro,pair_is_finite,upair_intro_left,subseteq_finite_is_finite,singleton_intro,inters_singleton,finite_intersections}.
            \end{subproof}
            Show that $W_1$ is open in $\productFamily{X}$ with respect to $T_0$ and
                $W_2$ is open in $\productFamily{X}$ with respect to $T_0$.
            \begin{subproof}
                Follows by \cref{finite_intersections,basis_elem_open_in_generated,basis_for_topology,open_in}.
            \end{subproof}
            Show that $W_1$ is a neighborhood of $x_1$ in $\productFamily{X}$ with respect to $T_0$ and
                $W_2$ is a neighborhood of $x_2$ in $\productFamily{X}$ with respect to $T_0$.
            \begin{subproof}
                Follows by \cref{neighborhood,projection_map_indexed,preimg_iff}.
            \end{subproof}
            Hence $W_1$ is disjoint from $W_2$ by \cref{projection_preimgs_disjoint}.
        \caseOf{$w\notin x_1$ and $w\in x_2$.}
            We have $x_1$ is a function and $\dom{x_1} = \dom{X}$ and
                $x_2$ is a function and $\dom{x_2} = \dom{X}$ by \cref{indexed_product,funs_elim}.
            Take $\alpha\in\dom{X}$ such that $w = (\alpha,\apply{x_2}{\alpha})$ by \cref{funs_elim,function_member_elim}.
            Therefore $\apply{x_1}{\alpha}\neq\apply{x_2}{\alpha}$ by \cref{function_apply_elim,pair_eq_iff}.
            $\apply{x_1}{\alpha}\in\apply{X}{\alpha}$ and
                $\apply{x_2}{\alpha}\in\apply{X}{\alpha}$ by \cref{indexed_product}.
            Take $U_1,U_2$ such that
                $U_1$ is a neighborhood of $\apply{x_1}{\alpha}$ in $\apply{X}{\alpha}$ with respect to $\apply{T}{\alpha}$ and
                $U_2$ is a neighborhood of $\apply{x_2}{\alpha}$ in $\apply{X}{\alpha}$ with respect to $\apply{T}{\alpha}$ and
                $U_1$ is disjoint from $U_2$ by \cref{hausdorff}.
            Let $W_1 = \preimg{\piIndex{X}{\alpha}}{U_1}$.
            Let $W_2 = \preimg{\piIndex{X}{\alpha}}{U_2}$.
            For all $x$ such that $x\in\dom{\piIndex{X}{\alpha}}$ we have $x\in\productFamily{X}$
                by \cref{dom,projection_map_indexed,pair_eq_iff}.
            Hence $\dom{\piIndex{X}{\alpha}}\subseteq\productFamily{X}$ by \cref{subseteq}.
            For all $x$ such that $x\in\productFamily{X}$ we have $x\in\dom{\piIndex{X}{\alpha}}$ by \cref{projection_map_indexed,dom}.
            $W_1\in\pow{\productFamily{X}}$ and $W_2\in\pow{\productFamily{X}}$
                by \cref{subseteq,subseteq_antisymmetric,preimg_subseteq_dom,subseteq_transitive,powerset_intro}.
            Hence $W_1\in S$ and $W_2\in S$ by \cref{neighborhood,open_in,product_subbasis_indexed,product_subbasis_union_indexed}.
            Show that $W_1\in\finiteIntersections{S}$.
            \begin{subproof}
                Follows by \cref{subbasis_on,subseteq,pow_iff,powerset_intro,singleton_subset_intro,pair_is_finite,upair_intro_left,subseteq_finite_is_finite,singleton_intro,inters_singleton,finite_intersections}.
            \end{subproof}
            Show that $W_2\in\finiteIntersections{S}$.
            \begin{subproof}
                Follows by \cref{subbasis_on,subseteq,pow_iff,powerset_intro,singleton_subset_intro,pair_is_finite,upair_intro_left,subseteq_finite_is_finite,singleton_intro,inters_singleton,finite_intersections}.
            \end{subproof}
            Show that $W_1$ is open in $\productFamily{X}$ with respect to $T_0$ and
                $W_2$ is open in $\productFamily{X}$ with respect to $T_0$.
            \begin{subproof}
                Follows by \cref{finite_intersections,basis_elem_open_in_generated,basis_for_topology,open_in}.
            \end{subproof}
            Show that $W_1$ is a neighborhood of $x_1$ in $\productFamily{X}$ with respect to $T_0$ and
                $W_2$ is a neighborhood of $x_2$ in $\productFamily{X}$ with respect to $T_0$.
            \begin{subproof}
                Follows by \cref{neighborhood,projection_map_indexed,preimg_iff}.
            \end{subproof}
            Hence $W_1$ is disjoint from $W_2$ by \cref{projection_preimgs_disjoint}.
        \end{byCase}
    \end{subproof}
    Therefore $\productFamily{X}$ is Hausdorff with respect to $T_0$ by \cref{hausdorff}.
\end{proof}

% from §19 helper lemma 5: closure family
\begin{definition}\label{closure_family}
    $\closureFamily{A}{X}{T} =
        \{(\alpha,\closure{\apply{A}{\alpha}}{\apply{X}{\alpha}}{\apply{T}{\alpha}}) \mid \alpha\in\dom{X}\}$.
\end{definition}

% from §19 Item 19: closure of products with the box topology
\begin{lemma}\label{product_closure_box}
    Assume $X$ is a function.
    Assume $A$ is a function.
    Assume $T$ is a function.
    Assume $\dom{A} = \dom{X}$.
    Assume $\dom{T} = \dom{X}$.
    Assume that for all $\alpha\in\dom{X}$ we have $\apply{A}{\alpha}\subseteq\apply{X}{\alpha}$.
    Assume that for all $\alpha\in\dom{X}$ we have $\apply{T}{\alpha}$ is a topology on $\apply{X}{\alpha}$.
    Let $C = \closureFamily{A}{X}{T}$.
    Then $\productFamily{C} = \closure{\productFamily{A}}{\productFamily{X}}{\boxTopology{T}{X}}$.
\end{lemma}
\begin{proof}
    Let $T_0 = \boxTopology{T}{X}$.
    Let $B_0 = \boxBasis{T}{X}$.
    $B_0$ is a basis on $\productFamily{X}$ by \cref{box_basis_is_basis}.
    Hence $T_0$ is a topology on $\productFamily{X}$ by \cref{basis_topology_is_topology,box_topology}.
    $\productFamily{A}$ is a subspace of $\productFamily{X}$ with respect to $T_0$
        and $\productFamily{A}\subseteq\productFamily{X}$ by \cref{product_subspace_box,subspace_of}.
    $C$ is a relation by \cref{closure_family,relation}.
    For all $a,b_1,b_2$ such that $(a,b_1)\in C$ and $(a,b_2)\in C$ we have $b_1 = b_2$
        by \cref{closure_family,pair_eq_iff}.
    Therefore $C$ is a function by \cref{rightunique}.
    Show that for all $x$ such that $x\in\productFamily{C}$ we have
        $x\in\closure{\productFamily{A}}{\productFamily{X}}{T_0}$.
    \begin{subproof}
                Fix $x$ such that $x\in\productFamily{C}$.
                $\dom{C} = \dom{X}$ by \cref{closure_family,dom,pair_eq_iff,subseteq,subseteq_antisymmetric}.
                $x$ is a function and $\dom{x} = \dom{C}$ by \cref{indexed_product,funs_elim}.
            For all $\alpha\in\dom{X}$ we have $\apply{x}{\alpha}\in\apply{X}{\alpha}$
                by \cref{indexed_product,closure_family,function_apply_intro,closure}.
            Hence $x\in\productFamily{X}$ by \cref{funs_intro,indexed_product,function_apply_elim,ran_intro,unions_iff}.
            Show that for all $B_1$ such that $B_1\in B_0$ and $x\in B_1$ we have $B_1$ intersects $\productFamily{A}$.
            \begin{subproof}
                Fix $B_1$ such that $B_1\in B_0$ and $x\in B_1$.
                Take $U$ such that
                    $U\in\funs{\dom{X}}{\pow{\unions{\ran{X}}}}$ and
                    $B_1 = \productFamily{U}$ and
                    for all $\alpha\in\dom{X}$ we have $\apply{U}{\alpha}\in\apply{T}{\alpha}$
                    by \cref{box_basis}.
                $\dom{U} = \dom{X}$ by \cref{funs_elim}.
                For all $\alpha\in\dom{X}$ we have $\apply{x}{\alpha}\in\apply{U}{\alpha}$ by \cref{indexed_product}.
                For all $\alpha\in\dom{X}$ we have $\apply{U}{\alpha}$ intersects $\apply{A}{\alpha}$
                    by \cref{open_in,closure_family,dom,indexed_product,function_apply_intro,closure_iff_intersects_open}.
                Let $R = \{w\in\dom{X}\times\unions{\ran{X}} \mid
                    \exists \alpha\in\dom{X}.\exists a\in\apply{A}{\alpha}\inter\apply{U}{\alpha}. w = (\alpha,a)\}$.
                $R$ is a relation by \cref{relation}.
                For all $\alpha\in\dom{X}$ there exists $a$ such that $\alpha\mathrel{R} a$
                    by \cref{intersects,nonempty_inhabited,inter,elem_subseteq,function_apply_elim,ran_intro,unions_iff,times_tuple_intro}.
                Take $f$ such that $f$ is a function on $\dom{X}$ and for all $\alpha\in\dom{X}$ we have
                    $\alpha\mathrel{R}\apply{f}{\alpha}$ by \cref{choice_function}.
                Hence $f$ is a function and $\dom{f} = \dom{X}$.
                For all $\alpha\in\dom{X}$ we have $\apply{f}{\alpha}\in\apply{A}{\alpha}$ and
                    $\apply{f}{\alpha}\in\apply{U}{\alpha}$ by \cref{pair_eq_iff,inter}.
                Thus $f\in\productFamily{A}$ by \cref{funs_intro,indexed_product,function_apply_elim,ran_intro,unions_iff}.
                Therefore $B_1$ intersects $\productFamily{A}$ by \cref{funs_intro,indexed_product,funs_elim,function_apply_elim,ran_intro,unions_iff,inter_intro,emptyset,intersects}.
            \end{subproof}
            Hence $x\in\closure{\productFamily{A}}{\productFamily{X}}{T_0}$
                by \cref{box_basis_is_basis,basis_for_topology,box_topology,closure_iff_intersects_basis}.
    \end{subproof}
    Therefore $\productFamily{C}\subseteq \closure{\productFamily{A}}{\productFamily{X}}{T_0}$ by \cref{subseteq}.
    Show that for all $x$ such that $x\in\closure{\productFamily{A}}{\productFamily{X}}{T_0}$ we have
        $x\in\productFamily{C}$.
    \begin{subproof}
            Fix $x$ such that $x\in\closure{\productFamily{A}}{\productFamily{X}}{T_0}$.
            $x\in\productFamily{X}$ by \cref{closure}.
            $\dom{C} = \dom{X}$ by \cref{closure_family,dom,pair_eq_iff,subseteq,subseteq_antisymmetric}.
            Show that for all $\alpha\in\dom{X}$ we have
                $\apply{x}{\alpha}\in\closure{\apply{A}{\alpha}}{\apply{X}{\alpha}}{\apply{T}{\alpha}}$.
            \begin{subproof}
                Fix $\alpha$ such that $\alpha\in\dom{X}$.
                Show that for all $V$ such that
                    $V$ is open in $\apply{X}{\alpha}$ with respect to $\apply{T}{\alpha}$ and $\apply{x}{\alpha}\in V$
                    we have $V$ intersects $\apply{A}{\alpha}$.
                \begin{subproof}
                    Fix $V$ such that $V$ is open in $\apply{X}{\alpha}$ with respect to $\apply{T}{\alpha}$ and $\apply{x}{\alpha}\in V$.
                    Take $U$ such that
                        $U\in\funs{\dom{X}}{\pow{\unions{\ran{X}}}}$ and
                        $\apply{U}{\alpha} = V$ and
                        for all $\beta\in\dom{X}$ such that $\beta\neq\alpha$ we have $\apply{U}{\beta} = \apply{X}{\beta}$
                        by \cref{open_in,topology_on,subseteq,pow_iff,function_apply_elim,ran_intro,unions_iff,subseteq_transitive,replace_coordinate_function}.
                    For all $\beta\in\dom{X}$ we have $\apply{U}{\beta}\in\apply{T}{\beta}$ by \cref{open_in,topology_on}.
                    Let $B_1 = \productFamily{U}$.
                    For all $\beta\in\dom{X}$ we have $\apply{U}{\beta}\subseteq\apply{X}{\beta}$
                        by \cref{open_in,topology_on,subseteq,pow_iff,replace_coordinate_function_subset}.
                    Hence $B_1\subseteq\productFamily{X}$ by \cref{indexed_product_subset_from_coordinate_subsets}.
                    Therefore $B_1$ is open in $\productFamily{X}$ with respect to $T_0$ by \cref{subseteq,pow_iff,box_basis,basis_elem_open_in_generated,box_topology,open_in}.
                    For all $\beta\in\dom{X}$ we have $\apply{x}{\beta}\in\apply{U}{\beta}$ by \cref{indexed_product,subseteq}.
                    $x$ is a function and $\dom{x} = \dom{X}$ and
                        $\dom{U} = \dom{X}$ by \cref{indexed_product,funs_elim}.
                    $\dom{x} = \dom{U}$.
                    $x\in B_1$ by \cref{funs_intro,indexed_product,funs_elim,function_apply_elim,ran_intro,unions_iff}.
                    For all $U$ such that $U$ is open in $\productFamily{X}$ with respect to $T_0$ and $x\in U$
                        we have $U$ intersects $\productFamily{A}$ by \cref{closure_iff_intersects_open}.
                    Hence $B_1$ intersects $\productFamily{A}$.
                    Take $f$ such that $f\in B_1\inter\productFamily{A}$ by \cref{intersects,nonempty_inhabited}.
                    $f\in\productFamily{A}$ and $f\in B_1$ by \cref{inter}.
                    $\apply{f}{\alpha}\in\apply{A}{\alpha}$ and $\apply{f}{\alpha}\in\apply{U}{\alpha}$ by \cref{inter,indexed_product,funs_elim}.
                    Thus $V$ intersects $\apply{A}{\alpha}$ by \cref{inter_intro,emptyset,intersects}.
                \end{subproof}
                Therefore $\apply{x}{\alpha}\in\closure{\apply{A}{\alpha}}{\apply{X}{\alpha}}{\apply{T}{\alpha}}$ by \cref{indexed_product,closure_iff_intersects_open}.
            \end{subproof}
            For all $\alpha\in\dom{X}$ we have $\apply{x}{\alpha}\in\apply{C}{\alpha}$ by \cref{closure_family,function_apply_intro}.
            $x$ is a function and $\dom{x} = \dom{X}$ by \cref{indexed_product,funs_elim}.
            Therefore $x\in\productFamily{C}$ by \cref{funs_intro,indexed_product,funs_elim,closure_family,dom,pair_eq_iff,subseteq,subseteq_antisymmetric,function_apply_elim,ran_intro,unions_iff}.
    \end{subproof}
    Hence $\closure{\productFamily{A}}{\productFamily{X}}{T_0}\subseteq\productFamily{C}$ by \cref{subseteq}.
    Hence $\productFamily{C} = \closure{\productFamily{A}}{\productFamily{X}}{T_0}$ by \cref{subseteq_antisymmetric}.
\end{proof}
% from §19 Item 20: closure of products with the product topology
\begin{lemma}\label{product_closure_product}
    Assume $X$ is a function.
    Assume $A$ is a function.
    Assume $T$ is a function.
    Assume $\dom{A} = \dom{X}$.
    Assume $\dom{T} = \dom{X}$.
    Assume $\dom{X}$ is inhabited.
    Assume that for all $\alpha\in\dom{X}$ we have $\apply{A}{\alpha}\subseteq\apply{X}{\alpha}$.
    Assume that for all $\alpha\in\dom{X}$ we have $\apply{T}{\alpha}$ is a topology on $\apply{X}{\alpha}$.
    Let $C = \closureFamily{A}{X}{T}$.
    Then $\productFamily{C} = \closure{\productFamily{A}}{\productFamily{X}}{\productTopologyIndexed{T}{X}}$.
\end{lemma}
\begin{proof}
    Let $T_0 = \productTopologyIndexed{T}{X}$.
    Let $T_1 = \boxTopology{T}{X}$.
    Let $S = \productSubbasis{T}{X}$.
    $C$ is a relation by \cref{closure_family,relation}.
    For all $a,b_1,b_2$ such that $(a,b_1)\in C$ and $(a,b_2)\in C$ we have $b_1 = b_2$
        by \cref{closure_family,pair_eq_iff}.
    Therefore $C$ is a function by \cref{rightunique}.
    $\productFamily{A}$ is a subspace of $\productFamily{X}$ with respect to $T_0$
        and $\productFamily{A}\subseteq\productFamily{X}$ by \cref{product_subspace_product,subspace_of}.
    Show that for all $x$ such that $x\in\closure{\productFamily{A}}{\productFamily{X}}{T_0}$ we have
        $x\in\productFamily{C}$.
    \begin{subproof}
            Fix $x$ such that $x\in\closure{\productFamily{A}}{\productFamily{X}}{T_0}$.
            $x\in\productFamily{X}$ by \cref{closure}.
            $\dom{C} = \dom{X}$ by \cref{closure_family,dom,pair_eq_iff,subseteq,subseteq_antisymmetric}.
            Show that for all $\alpha\in\dom{X}$ we have
                $\apply{x}{\alpha}\in\closure{\apply{A}{\alpha}}{\apply{X}{\alpha}}{\apply{T}{\alpha}}$.
            \begin{subproof}
                Fix $\alpha$ such that $\alpha\in\dom{X}$.
                Show that for all $V$ such that
                    $V$ is open in $\apply{X}{\alpha}$ with respect to $\apply{T}{\alpha}$ and $\apply{x}{\alpha}\in V$
                    we have $V$ intersects $\apply{A}{\alpha}$.
                \begin{subproof}
                    Fix $V$ such that $V$ is open in $\apply{X}{\alpha}$ with respect to $\apply{T}{\alpha}$ and $\apply{x}{\alpha}\in V$.
                    Let $W = \preimg{\piIndex{X}{\alpha}}{V}$.
                    $W\in\productSubbasisAt{T}{X}{\alpha}$ by
                        \cref{preimg_iff,projection_map_indexed,pair_eq_iff,subseteq,pow_iff,open_in,product_subbasis_indexed}.
                    $W\in S$ by \cref{product_subbasis_indexed,product_subbasis_union_indexed}.
                    $W\in\finiteIntersections{S}$ by
                        \cref{unions_iff,subseteq,pow_iff,singleton_subset_intro,upair_iff,pair_is_finite,subseteq_finite_is_finite,singleton_intro,inters_singleton,finite_intersections}.
                    Therefore $W$ is open in $\productFamily{X}$ with respect to $T_0$ by
                        \cref{product_subbasis_is_subbasis,product_subbasis_finite_intersections_basis,basis_elem_open_in_generated,basis_for_topology,basis_topology_is_topology,open_in}.
                    $x\in W$ by \cref{projection_map_indexed,preimg_iff}.
                    Hence $W$ intersects $\productFamily{A}$ by \cref{basis_for_topology,basis_topology_is_topology,product_subbasis_finite_intersections_basis,subspace_of,closure_iff_intersects_open}.
                    Take $f$ such that $f\in W\inter\productFamily{A}$ by \cref{intersects,nonempty_inhabited}.
                    $f\in W$ and $f\in\productFamily{A}$ by \cref{inter}.
                    Hence $\apply{f}{\alpha}\in V$ and $\apply{f}{\alpha}\in\apply{A}{\alpha}$
                        by \cref{inter,preimg_iff,projection_map_indexed,pair_eq_iff,indexed_product}.
                    Therefore $V$ intersects $\apply{A}{\alpha}$ by \cref{inter_intro,emptyset,intersects}.
                \end{subproof}
                Therefore $\apply{x}{\alpha}\in\closure{\apply{A}{\alpha}}{\apply{X}{\alpha}}{\apply{T}{\alpha}}$
                    by \cref{indexed_product,closure_iff_intersects_open}.
            \end{subproof}
            For all $\alpha\in\dom{X}$ we have $\apply{x}{\alpha}\in\apply{C}{\alpha}$ by \cref{closure_family,function_apply_intro}.
            Therefore $x\in\productFamily{C}$ by \cref{funs_intro,indexed_product,funs_elim,closure_family,dom,pair_eq_iff,subseteq,subseteq_antisymmetric,function_apply_elim,ran_intro,unions_iff}.
    \end{subproof}
    Hence $\closure{\productFamily{A}}{\productFamily{X}}{T_0}\subseteq\productFamily{C}$ by \cref{subseteq}.
    $\productFamily{C} = \closure{\productFamily{A}}{\productFamily{X}}{T_1}$ by \cref{product_closure_box}.
    Show that for all $U$ such that $U\in T_0$ we have $U\in T_1$.
    \begin{subproof}
                Fix $U$ such that $U\in T_0$.
                $U\in\basisTopology{\finiteIntersections{S}}{\productFamily{X}}$ by
                    \cref{product_subbasis_is_subbasis,subbasis_finite_intersections_basis,product_topology_indexed,topology_generated_by_subbasis}.
                Let $M = \{B\in\finiteIntersections{S} \mid B\subseteq U\}$.
                For all $x$ such that $x\in U$ there exists $B\in\finiteIntersections{S}$ such that $x\in B$ and $B\subseteq U$
                    by \cref{topology_generated_by_basis}.
                For all $x$ such that $x\in U$ we have $x\in\unions{M}$ by \cref{unions_iff}.
                Therefore $U = \unions{M}$ by \cref{unions_iff,subseteq,subseteq_antisymmetric}.
                Show that for all $B$ such that $B\in\finiteIntersections{S}$ we have $B\in T_1$.
                \begin{subproof}
                    Fix $B$ such that $B\in\finiteIntersections{S}$.
                    Take $F\subseteq S$ such that $F$ is finite and inhabited and $B = \inters{F}$ by \cref{finite_intersections}.
                    Show that for all $A_0$ such that $A_0\in F$ we have $A_0\in T_1$.
                    \begin{subproof}
                        Fix $A_0$ such that $A_0\in F$.
                        Take $\beta\in\dom{X}$ such that $A_0\in\productSubbasisAt{T}{X}{\beta}$ by \cref{subseteq,product_subbasis_union_indexed}.
                        Take $V\in\apply{T}{\beta}$ such that $A_0 = \preimg{\piIndex{X}{\beta}}{V}$ by \cref{product_subbasis_indexed}.
                        Show that for all $x$ such that $x\in A_0$ there exists $B_1\in\boxBasis{T}{X}$ such that $x\in B_1$ and $B_1\subseteq A_0$.
                        \begin{subproof}
                            Fix $x$ such that $x\in A_0$.
                            $\apply{X}{\beta}\subseteq\unions{\ran{X}}$ by
                                \cref{function_apply_elim,ran_intro,unions_iff,subseteq}.
                            Hence $V\subseteq\unions{\ran{X}}$ by
                                \cref{open_in,topology_on,subseteq,pow_iff,subseteq_transitive}.
                            Take $W$ such that
                                $W\in\funs{\dom{X}}{\pow{\unions{\ran{X}}}}$ and
                                $\apply{W}{\beta} = V$ and
                                for all $\gamma\in\dom{X}$ such that $\gamma\neq\beta$ we have $\apply{W}{\gamma} = \apply{X}{\gamma}$
                                by \cref{replace_coordinate_function}.
                            For all $\gamma\in\dom{X}$ we have $\apply{W}{\gamma}\in\apply{T}{\gamma}$ by \cref{open_in,topology_on}.
                            Let $B_1 = \productFamily{W}$.
                            For all $\gamma\in\dom{X}$ we have $\apply{W}{\gamma}\subseteq\apply{X}{\gamma}$
                                by \cref{open_in,topology_on,subseteq,pow_iff,replace_coordinate_function_subset}.
                            Hence $B_1\subseteq\productFamily{X}$ by \cref{indexed_product_subset_from_coordinate_subsets}.
                            $B_1\in\pow{\productFamily{X}}$ by \cref{subseteq,pow_iff}.
                            $B_1\in\boxBasis{T}{X}$ by \cref{box_basis}.
                            $x\in\productFamily{X}$ and $\apply{x}{\beta}\in V$ by \cref{preimg_iff,projection_map_indexed,pair_eq_iff}.
                            For all $\gamma\in\dom{X}$ we have $\apply{x}{\gamma}\in\apply{W}{\gamma}$ by \cref{indexed_product,subseteq}.
                            $x\in\funs{\dom{X}}{\unions{\ran{X}}}$ by \cref{indexed_product}.
                            $x$ is a function and $\dom{x} = \dom{X}$ and
                                $\dom{W} = \dom{X}$ by \cref{funs_elim}.
                            $\dom{x} = \dom{W}$.
                            For all $\gamma\in\dom{W}$ we have $\apply{x}{\gamma}\in\unions{\ran{W}}$
                                by \cref{funs_elim,function_apply_elim,ran_intro,unions_iff}.
                            $x\in\funs{\dom{W}}{\unions{\ran{W}}}$ by \cref{funs_intro}.
                            $x\in B_1$ by \cref{indexed_product}.
                            For all $z$ such that $z\in B_1$ we have $z\in A_0$
                                by \cref{indexed_product,funs_elim,topology_on,elem_subseteq,pow_iff,funs_intro,function_apply_elim,ran_intro,unions_iff,projection_map_indexed,preimg_iff}.
                            Hence $B_1\subseteq A_0$ by \cref{subseteq}.
                        \end{subproof}
                        Therefore $A_0\in T_1$ by
                            \cref{preimg_iff,projection_map_indexed,pair_eq_iff,subseteq,pow_iff,topology_generated_by_basis,box_topology}.
                    \end{subproof}
                    $F\subseteq T_1$ by \cref{subseteq}.
                    Hence $B\in T_1$ by
                        \cref{box_basis_is_basis,basis_topology_is_topology,box_topology,finite_intersections_open_in_topology}.
                \end{subproof}
                Therefore $U\in T_1$ by
                    \cref{subseteq,box_basis_is_basis,basis_topology_is_topology,box_topology,topology_on}.
    \end{subproof}
    Hence $T_0\subseteq T_1$ by \cref{subseteq}.
    Show that for all $x$ such that $x\in\closure{\productFamily{A}}{\productFamily{X}}{T_1}$ we have
        $x\in\closure{\productFamily{A}}{\productFamily{X}}{T_0}$.
    \begin{subproof}
            Fix $x$ such that $x\in\closure{\productFamily{A}}{\productFamily{X}}{T_1}$.
            For all $U$ such that
                $U$ is open in $\productFamily{X}$ with respect to $T_0$ and $x\in U$
                we have $U$ intersects $\productFamily{A}$
                by \cref{subseteq,open_in,box_basis_is_basis,basis_topology_is_topology,box_topology,topology_on,pow_iff,elem_subseteq,subspace_of,closure_iff_intersects_open}.
            Hence $\finiteIntersections{S}$ is a basis for $T_0$ on $\productFamily{X}$
                by \cref{product_subbasis_is_subbasis,product_subbasis_finite_intersections_basis}.
            Therefore $T_0$ is a topology on $\productFamily{X}$ by \cref{basis_topology_is_topology,basis_for_topology}.
            $\productFamily{A}\subseteq\productFamily{X}$ by \cref{subspace_of}.
            $\boxBasis{T}{X}$ is a basis on $\productFamily{X}$ and
                $T_1$ is a topology on $\productFamily{X}$ by \cref{box_basis_is_basis,basis_topology_is_topology,box_topology}.
            $\productFamily{X}$ is closed in $\productFamily{X}$ with respect to $T_1$ by \cref{closed_whole_space}.
            $x\in\productFamily{X}$ by \cref{closure_subseteq_closed,elem_subseteq}.
            Hence $x\in\closure{\productFamily{A}}{\productFamily{X}}{T_0}$ by \cref{closure_iff_intersects_open}.
    \end{subproof}
    Hence $\closure{\productFamily{A}}{\productFamily{X}}{T_1}\subseteq\closure{\productFamily{A}}{\productFamily{X}}{T_0}$ by \cref{subseteq}.
    Therefore $\productFamily{C}\subseteq\closure{\productFamily{A}}{\productFamily{X}}{T_0}$ by \cref{subseteq}.
    Hence $\productFamily{C} = \closure{\productFamily{A}}{\productFamily{X}}{T_0}$ by \cref{subseteq_antisymmetric}.
\end{proof}

% from §19 helper lemma 3: projection map is a function
\begin{lemma}\label{projection_map_indexed_function}
    Assume $X$ is a function.
    Assume $\beta\in\dom{X}$.
    Then $\piIndex{X}{\beta}\in\funs{\productFamily{X}}{\apply{X}{\beta}}$.
\end{lemma}
\begin{proof}
    Let $P = \piIndex{X}{\beta}$.
    $P$ is a relation by \cref{projection_map_indexed,relation}.
    For all $x,y_1,y_2$ such that $(x,y_1)\in P$ and $(x,y_2)\in P$ we have $y_1 = y_2$
        by \cref{projection_map_indexed,pair_eq_iff}.
    Hence $P$ is right-unique by \cref{rightunique}.
    $P$ is a function by \cref{relation}.
    For all $x$ such that $x\in\productFamily{X}$ we have $x\in\dom{P}$ by \cref{projection_map_indexed,dom_intro}.
    $\productFamily{X}\subseteq\dom{P}$ by \cref{subseteq}.
    For all $x$ such that $x\in\dom{P}$ we have $x\in\productFamily{X}$ by \cref{dom_iff,projection_map_indexed,pair_eq_iff}.
    $\dom{P}\subseteq\productFamily{X}$ by \cref{subseteq}.
    Therefore $\dom{P} = \productFamily{X}$ by \cref{subseteq_antisymmetric}.
    For all $x$ such that $x\in\dom{P}$ we have $\apply{P}{x}\in\apply{X}{\beta}$
        by \cref{projection_map_indexed,function_apply_iff,indexed_product}.
    Hence $P\in\funs{\productFamily{X}}{\apply{X}{\beta}}$ by \cref{funs_intro}.
\end{proof}

% from §19 helper lemma 4: image of a finite set under a function is finite
\begin{lemma}\label{finite_image_function}
    Assume $F$ is finite.
    Assume $g$ is a function on $F$.
    Then $\img{g}{F}$ is finite.
\end{lemma}
\begin{proof}
    Let $B = \img{g}{F}$.
    Let $R = \{w\in B\times F \mid \exists y\in B.\exists x\in F. w = (y,x) \land \apply{g}{x} = y\}$.
    $R$ is a relation by \cref{relation}.
    For all $y\in B$ there exists $x\in F$ such that $g(x) = y$ by \cref{img_iff,function_apply_iff}.
    Hence for all $y\in B$ there exists $x$ such that $y\mathrel{R} x$ by \cref{times_tuple_intro}.
    Take $s$ such that $s$ is a function on $B$ and for all $y\in B$ we have $y\mathrel{R}\apply{s}{y}$ by \cref{choice_function}.
    Hence $s$ is a function and $\dom{s} = B$.
    For all $y\in B$ we have $\apply{s}{y}\in F$ and $\apply{g}{\apply{s}{y}} = y$ by \cref{pair_eq_iff}.
    Let $C = \ran{s}$.
    For all $y_1,y_2$ such that $y_1\in B$ and $y_2\in B$ and $\apply{s}{y_1} = \apply{s}{y_2}$ we have $y_1 = y_2$
        by \cref{subseteq}.
    Therefore $s$ is injective.
    Show that for all $x$ such that $x\in C$ we have $x\in F$.
    \begin{subproof}
        Fix $x$ such that $x\in C$.
        Take $y$ such that $y\mathrel{s} x$ by \cref{ran_iff}.
        Hence $y\in B$ by \cref{dom_intro}.
        Therefore $x\in F$ by \cref{subseteq,function_apply_iff}.
    \end{subproof}
    Hence $C$ is finite by \cref{subseteq,subseteq_finite_is_finite}.
    Therefore $s\in\bijections{B}{C}$ by \cref{function_apply_elim,ran_intro,funs_intro,funs_surj_iff,bijections}.
    Hence $\converse{s}$ is a bijection from $C$ to $B$ by \cref{bijection_converse_is_bijection}.
    Therefore $B$ is finite by \cref{finite_bijection}.
\end{proof}

% from §19 Item 21: continuity into an indexed product
\begin{lemma}\label{continuous_map_into_indexed_product}
    Assume $X$ is a function.
    Assume $T$ is a function.
    Assume $\dom{T} = \dom{X}$.
    Assume $\dom{X}$ is inhabited.
    Assume that for all $\alpha\in\dom{X}$ we have $\apply{T}{\alpha}$ is a topology on $\apply{X}{\alpha}$.
    Assume $T_A$ is a topology on $A$.
    Assume $f\in\funs{A}{\productFamily{X}}$.
    Then $f$ is continuous from $A$ to $\productFamily{X}$ with respect to $T_A$ and $\productTopologyIndexed{T}{X}$ iff
        for all $\alpha\in\dom{X}$ we have $\piIndex{X}{\alpha}\circ f$ is continuous from $A$ to $\apply{X}{\alpha}$ with respect to $T_A$ and $\apply{T}{\alpha}$.
\end{lemma}
\begin{proof}
    Let $T_0 = \productTopologyIndexed{T}{X}$.
    Let $S = \productSubbasis{T}{X}$.
    For all $\alpha\in\dom{X}$ we have $\apply{T}{\alpha}$ is a topology on $\apply{X}{\alpha}$ by \cref{subseteq}.
    Show that if $f$ is continuous from $A$ to $\productFamily{X}$ with respect to $T_A$ and $T_0$ then
        for all $\alpha\in\dom{X}$ we have $\piIndex{X}{\alpha}\circ f$ is continuous from $A$ to $\apply{X}{\alpha}$
        with respect to $T_A$ and $\apply{T}{\alpha}$.
    \begin{subproof}
        Assume $f$ is continuous from $A$ to $\productFamily{X}$ with respect to $T_A$ and $T_0$.
        Fix $\alpha$ such that $\alpha\in\dom{X}$.
        Let $g = \piIndex{X}{\alpha}\circ f$.
        $\piIndex{X}{\alpha}\in\funs{\productFamily{X}}{\apply{X}{\alpha}}$ and
            $g\in\funs{A}{\apply{X}{\alpha}}$ by \cref{projection_map_indexed_function,funs_circ,continuous}.
        Show that for all $V$ such that $V$ is open in $\apply{X}{\alpha}$ with respect to $\apply{T}{\alpha}$
            we have $\preimg{g}{V}$ is open in $A$ with respect to $T_A$.
        \begin{subproof}
            Fix $V$ such that $V$ is open in $\apply{X}{\alpha}$ with respect to $\apply{T}{\alpha}$.
            Let $W = \preimg{\piIndex{X}{\alpha}}{V}$.
            $W\in\pow{\productFamily{X}}$
                by \cref{projection_map_indexed_function,funs_elim,preimg_subseteq_dom,subseteq_transitive,pow_iff}.
            $V\in\apply{T}{\alpha}$ and $W\in\productSubbasisAt{T}{X}{\alpha}$ and $W\in S$
                by \cref{open_in,product_subbasis_indexed,product_subbasis_union_indexed}.
            Let $F = \{W\}$.
            $F\subseteq S$ and $F$ is finite and inhabited and $\inters{F} = W$ by
                \cref{singleton_subset_intro,upair_iff,pair_is_finite,subseteq_finite_is_finite,singleton_intro,inters_singleton}.
            Therefore $W\in\finiteIntersections{S}$ by \cref{unions_iff,subseteq,pow_iff,finite_intersections}.
            $W$ is open in $\productFamily{X}$ with respect to $T_0$ by
                \cref{product_subbasis_is_subbasis,product_subbasis_finite_intersections_basis,basis_elem_open_in_generated,basis_for_topology,basis_topology_is_topology,open_in}.
            Hence $\preimg{g}{V}$ is open in $A$ with respect to $T_A$ by
                \cref{continuous,circ_iff,preimg_iff,subseteq,subseteq_antisymmetric}.
        \end{subproof}
        Therefore $g$ is continuous from $A$ to $\apply{X}{\alpha}$ with respect to $T_A$ and $\apply{T}{\alpha}$ by \cref{continuous}.
    \end{subproof}
    Show that if for all $\alpha\in\dom{X}$ we have $\piIndex{X}{\alpha}\circ f$ is continuous from $A$ to $\apply{X}{\alpha}$
        with respect to $T_A$ and $\apply{T}{\alpha}$ then $f$ is continuous from $A$ to $\productFamily{X}$ with respect to $T_A$ and $T_0$.
    \begin{subproof}
        Assume for all $\alpha\in\dom{X}$ we have $\piIndex{X}{\alpha}\circ f$ is continuous from $A$ to $\apply{X}{\alpha}$
            with respect to $T_A$ and $\apply{T}{\alpha}$.
        Show that for all $U$ such that $U$ is open in $\productFamily{X}$ with respect to $T_0$
            we have $\preimg{f}{U}$ is open in $A$ with respect to $T_A$.
        \begin{subproof}
            Fix $U$ such that $U$ is open in $\productFamily{X}$ with respect to $T_0$.
            Take $M\subseteq \finiteIntersections{S}$ such that $U = \unions{M}$ by
                \cref{open_in,product_subbasis_is_subbasis,product_subbasis_finite_intersections_basis,basis_topology_unions}.
            Let $N = \{\preimg{f}{B} \mid B\in M\}$.
            Show that for all $C$ such that $C\in N$ we have $C$ is open in $A$ with respect to $T_A$.
            \begin{subproof}
                Fix $C$ such that $C\in N$.
                Take $B\in M$ such that $C = \preimg{f}{B}$.
                Take $F\subseteq S$ such that $F$ is finite and inhabited and $B = \inters{F}$ by \cref{subseteq,finite_intersections}.
                Let $H = \{\preimg{f}{A} \mid A\in F\}$.
                Show that for all $D$ such that $D\in H$ we have $D$ is open in $A$ with respect to $T_A$.
                \begin{subproof}
                    Fix $D$ such that $D\in H$.
                    Take $E\in F$ such that $D = \preimg{f}{E}$.
                    Take $\beta\in\dom{X}$ such that $E\in\productSubbasisAt{T}{X}{\beta}$ by \cref{subseteq,product_subbasis_union_indexed}.
                    Take $V\in\apply{T}{\beta}$ such that $E = \preimg{\piIndex{X}{\beta}}{V}$ by \cref{product_subbasis_indexed}.
                    Let $g = \piIndex{X}{\beta}\circ f$.
                    Therefore $D$ is open in $A$ with respect to $T_A$ by
                        \cref{subseteq,open_in,continuous,circ_iff,preimg_iff,subseteq_antisymmetric}.
                \end{subproof}
                Let $h = \{(A,\preimg{f}{A}) \mid A\in F\}$.
                $h$ is a function on $F$ by
                    \cref{relation,pair_eq_iff,rightunique,dom_iff,subseteq,subseteq_antisymmetric}.
                For all $D$ such that $D\in H$ we have $D\in\img{h}{F}$.
                For all $D$ such that $D\in\img{h}{F}$ we have $D\in H$ by \cref{img_iff,pair_eq_iff}.
                Hence $H = \img{h}{F}$ by \cref{subseteq,subseteq_antisymmetric}.
                Therefore $H$ is finite by \cref{finite_image_function,subseteq,subseteq_finite_is_finite}.
                Take $E_0$ such that $E_0\in F$.
                Hence $H$ is inhabited and for all $D$ such that $D\in H$ we have $D\in T_A$
                    by \cref{open_in}.
                Therefore $\inters{H}$ is open in $A$ with respect to $T_A$ by
                    \cref{subseteq,topology_on,finite_intersections_open_in_topology,open_in}.
                For all $a$ such that $a\in C$ we have $a\in\inters{H}$ by \cref{preimg_iff,subseteq,inters_iff_forall}.
                Hence $C\subseteq\inters{H}$ by \cref{subseteq}.
                Show that for all $a$ such that $a\in\inters{H}$ we have $a\in C$.
                \begin{subproof}
                    Fix $a$ such that $a\in\inters{H}$.
                    $a\in A$ by \cref{open_in,topology_on,subseteq,pow_iff}.
                    Show that for all $E$ such that $E\in F$ we have $\apply{f}{a}\in E$.
                    \begin{subproof}
                        Fix $E$ such that $E\in F$.
                        $(E,\preimg{f}{E})\in h$ by \cref{pair_eq_iff}.
                        Hence $\preimg{f}{E}\in\img{h}{F}$ by \cref{img_iff}.
                        $\preimg{f}{E}\in H$ by \cref{subseteq,subseteq_antisymmetric}.
                        Therefore $a\in\preimg{f}{E}$ by \cref{inters_iff_forall}.
                        Hence $\apply{f}{a}\in E$ by
                            \cref{preimg_iff,funs_is_function,function_rightunique,function_apply_iff}.
                    \end{subproof}
                    Hence $\apply{f}{a}\in\inters{F}$ by \cref{inters_iff_forall}.
                    $\apply{f}{a}\in B$ by \cref{subseteq,subseteq_antisymmetric}.
                    $f$ is a function and $\dom{f} = A$ by \cref{funs_elim}.
                    Hence $a\in\dom{f}$ by \cref{subseteq}.
                    Therefore $(a,\apply{f}{a})\in f$ by \cref{function_apply_elim}.
                    Hence $a\in C$ by \cref{preimg_iff}.
                \end{subproof}
                $\inters{H}\subseteq C$ by \cref{subseteq}.
                Therefore $C = \inters{H}$ by \cref{subseteq_antisymmetric}.
                Hence $C$ is open in $A$ with respect to $T_A$.
            \end{subproof}
            Therefore $N\subseteq T_A$ by \cref{subseteq,open_in}.
            $T_A$ is a topology on $A$ and $\unions{N}\in T_A$ by \cref{continuous,topology_on}.
            For all $a$ such that $a\in\preimg{f}{U}$ we have $a\in\unions{N}$ by \cref{preimg_iff,unions_iff,img_iff}.
            Hence $\preimg{f}{U}\subseteq\unions{N}$ by \cref{subseteq}.
            For all $a$ such that $a\in\unions{N}$ we have $a\in\preimg{f}{U}$ by \cref{unions_iff,img_iff,preimg_iff,elem_subseteq}.
            $\unions{N}\subseteq\preimg{f}{U}$ by \cref{subseteq}.
            Therefore $\preimg{f}{U} = \unions{N}$ by \cref{subseteq_antisymmetric}.
            Hence $\preimg{f}{U}$ is open in $A$ with respect to $T_A$ by \cref{open_in}.
        \end{subproof}
        Therefore $f$ is continuous from $A$ to $\productFamily{X}$ with respect to $T_A$ and $T_0$ by
            \cref{product_subbasis_is_subbasis,product_subbasis_finite_intersections_basis,basis_topology_is_topology,basis_for_topology,continuous}.
    \end{subproof}
\end{proof}
\section{§20 The Metric Topology}

% from §20 helper lemma 1: metric nonnegative
\begin{definition}\label{metric_nonnegative}
    $d$ is nonnegative on $X$ iff
    for all $x,y\in X$ we have $d((x,y)) \geq 0$.
\end{definition}

% from §20 helper lemma 2: metric zero characterization
\begin{definition}\label{metric_zero_characterization}
    $d$ is zero-characterized on $X$ iff
    for all $x,y\in X$ we have $d((x,y)) = 0$ iff $x = y$.
\end{definition}

% from §20 helper lemma 3: metric symmetry
\begin{definition}\label{metric_symmetric}
    $d$ is symmetric on $X$ iff
    for all $x,y\in X$ we have $d((x,y)) = d((y,x))$.
\end{definition}

% from §20 helper lemma 4: metric triangle inequality
\begin{definition}\label{metric_triangle_inequality}
    $d$ satisfies the triangle inequality on $X$ iff
    for all $x,y,z\in X$ we have $d((x,y)) + d((y,z)) \geq d((x,z))$.
\end{definition}

% from §20 Item 1: metric on a set
\begin{definition}\label{metric_on}
    $d$ is a metric on $X$ iff
    $d\in\funs{X\times X}{\reals}$ and
    $d$ is nonnegative on $X$ and
    $d$ is zero-characterized on $X$ and
    $d$ is symmetric on $X$ and
    $d$ satisfies the triangle inequality on $X$.
\end{definition}

% from §20 Item 2: metric ball
\begin{definition}\label{metric_ball}
    $\metricBall{d}{X}{x}{\epsilon} = \{y\in X \mid d((x,y)) < \epsilon\}$.
\end{definition}

% from §20 Item 3: metric basis
\begin{definition}\label{metric_basis}
    $\metricBasis{d}{X} = \{B\in\pow{X} \mid \exists x\in X. \exists \epsilon\in\reals. 0 < \epsilon \land B = \metricBall{d}{X}{x}{\epsilon}\}$.
\end{definition}

% from §20 Item 4: metric topology induced by a metric
\begin{definition}\label{metric_topology}
    $\metricTopology{d}{X} = \basisTopology{\metricBasis{d}{X}}{X}$.
\end{definition}

% from §20 Item 5: metric basis is a basis
\begin{lemma}\label{metric_basis_is_basis}
    Assume $d$ is a metric on $X$.
    Then $\metricBasis{d}{X}$ is a basis on $X$.
\end{lemma}
\begin{proof}
    Let $B = \metricBasis{d}{X}$.
    Show that $B\subseteq\pow{X}$.
    \begin{subproof}
        Follows by \cref{metric_basis,metric_ball,subseteq,powerset_intro}.
    \end{subproof}
    Show that for all $x$ such that $x\in X$ there exists $B_1\in B$ such that $x\in B_1$.
    \begin{subproof}
        Follows by \cref{real_naturals_subset_reals,real_one_in_naturals,subseteq,real_zero_lt_one,metric_on,metric_zero_characterization,real_lt_subst_left,metric_ball,powerset_intro,metric_basis}.
    \end{subproof}
    Show that for all $B_1,B_2,x$ such that $B_1\in B$ and $B_2\in B$ and $x\in B_1\inter B_2$
        there exists $B_3\in B$ such that $x\in B_3$ and $B_3\subseteq B_1\inter B_2$.
    \begin{subproof}
        Fix $B_1$.
        Fix $B_2$.
        Fix $x$ such that $B_1\in B$ and $B_2\in B$ and $x\in B_1\inter B_2$.
        Take $x_1,\epsilon_1$ such that $x_1\in X$ and $\epsilon_1\in\reals$ and $0 < \epsilon_1$ and $B_1 = \metricBall{d}{X}{x_1}{\epsilon_1}$
            by \cref{metric_basis}.
        Take $x_2,\epsilon_2$ such that $x_2\in X$ and $\epsilon_2\in\reals$ and $0 < \epsilon_2$ and $B_2 = \metricBall{d}{X}{x_2}{\epsilon_2}$
            by \cref{metric_basis}.
        Thus $x\in X$ and $d((x_1,x)) < \epsilon_1$ and $d((x_2,x)) < \epsilon_2$
            by \cref{inter_elim_left,inter_elim_right,metric_ball}.
        $d\in\funs{X\times X}{\reals}$ and $d((x_1,x))\in\reals$ and $d((x_2,x))\in\reals$
            by \cref{metric_on,times_tuple_intro,funs_type_apply}.
        Let $r_1 = \epsilon_1 - d((x_1,x))$.
        Let $r_2 = \epsilon_2 - d((x_2,x))$.
        $0 < r_1$ and $0 < r_2$ by \cref{metric_on,times_tuple_intro,funs_type_apply,real_lt_imp_pos_diff}.
        $r_1\in\reals$ and $r_2\in\reals$ by \cref{real_sub,real_add_inverse,real_add_closed}.
        Let $A = \{r_1, r_2\}$.
        $A$ is finite and $A\neq\emptyset$ by \cref{pair_is_finite,upair_intro_left,emptyset}.
        $A\subseteq\reals$ by \cref{upair_elim,real_sub,subseteq}.
        Take $m\in A$ such that $m$ is a minimum of $A$ by \cref{finite_real_minimum}.
        Thus $0 < m$ by \cref{upair_elim}.
        Let $B_3 = \metricBall{d}{X}{x}{m}$.
        Hence $B_3\in B$ and $x\in B_3$ by
            \cref{subseteq,metric_ball,powerset_intro,metric_basis,metric_on,metric_zero_characterization,real_lt_subst_left}.
        Show that for all $z$ such that $z\in B_3$ we have $z\in B_1\inter B_2$.
        \begin{subproof}
            Fix $z$ such that $z\in B_3$.
            Thus $z\in X$ and $d((x,z)) < m$ by \cref{metric_ball}.
            $d((x,z))\in\reals$ and $d((x_1,z))\in\reals$ and $d((x_2,z))\in\reals$
                by \cref{metric_on,times_tuple_intro,funs_type_apply}.
            $r_1 + d((x_1,x)) = \epsilon_1$ and $r_2 + d((x_2,x)) = \epsilon_2$
                by \cref{real_sub,real_add_assoc,real_add_inverse,real_add_comm,real_add_ident}.
            $d((x,z)) < r_1$ by
                \cref{real_minimum,upair_intro_left,times_tuple_intro,funs_type_apply,subseteq,real_sub,real_add_closed,real_leq_split,real_lt_trans,real_lt_subst_right}.
            $d((x,z)) + d((x_1,x)) < \epsilon_1$ by
                \cref{times_tuple_intro,funs_type_apply,real_sub,real_add_closed,real_lt_add,real_add_inverse,real_add_assoc,real_add_comm,real_add_ident,real_lt_subst_right}.
            $d((x_1,x)) + d((x,z)) < \epsilon_1$ by
                \cref{funs_type_apply,times_tuple_intro,real_add_comm,real_lt_subst_left}.
            $d((x_1,z)) \leq d((x_1,x)) + d((x,z))$ by \cref{metric_on,metric_triangle_inequality}.
            Hence $d((x_1,z)) < \epsilon_1$ by
                \cref{times_tuple_intro,funs_type_apply,real_add_closed,real_leq_split,real_lt_trans,real_lt_subst_left}.
            $z\in B_1$ by \cref{metric_ball}.
            $d((x,z)) < r_2$ by
                \cref{real_minimum,upair_intro_right,times_tuple_intro,funs_type_apply,subseteq,real_sub,real_add_inverse,real_add_closed,real_leq_split,real_lt_trans,real_lt_subst_right}.
            $d((x,z)) + d((x_2,x)) < \epsilon_2$ by
                \cref{times_tuple_intro,funs_type_apply,real_sub,real_add_inverse,real_add_closed,real_lt_add,real_add_assoc,real_add_comm,real_add_ident,real_lt_subst_right}.
            $d((x_2,x)) + d((x,z)) < \epsilon_2$ by
                \cref{funs_type_apply,times_tuple_intro,real_add_comm,real_lt_subst_left}.
            $d((x_2,z)) \leq d((x_2,x)) + d((x,z))$ by \cref{metric_on,metric_triangle_inequality}.
            Hence $d((x_2,z)) < \epsilon_2$ by
                \cref{times_tuple_intro,funs_type_apply,real_add_closed,real_leq_split,real_lt_trans,real_lt_subst_left}.
            $z\in B_2$ by \cref{metric_ball}.
            Therefore $z\in B_1\inter B_2$ by \cref{inter_intro}.
        \end{subproof}
        Hence $B_3\subseteq B_1\inter B_2$ by \cref{subseteq}.
    \end{subproof}
    Therefore $B$ is a basis on $X$ by \cref{basis_on}.
\end{proof}

% from §20 Item 6: metric open sets
\begin{lemma}\label{metric_open_iff}
    Assume $d$ is a metric on $X$.
    Let $T = \metricTopology{d}{X}$.
    Then for all $U$ such that $U\subseteq X$ we have
    $U$ is open in $X$ with respect to $T$ iff
    for all $y\in U$ there exists $\delta\in\reals$ such that $0 < \delta$ and $\metricBall{d}{X}{y}{\delta}\subseteq U$.
\end{lemma}
\begin{proof}
    Let $B = \metricBasis{d}{X}$.
    $T = \basisTopology{B}{X}$ by \cref{metric_topology}.
    Fix $U$ such that $U\subseteq X$.
    Show that if $U$ is open in $X$ with respect to $T$ then
        for all $y\in U$ there exists $\delta\in\reals$ such that $0 < \delta$ and $\metricBall{d}{X}{y}{\delta}\subseteq U$.
    \begin{subproof}
        Assume $U$ is open in $X$ with respect to $T$.
        Fix $y$ such that $y\in U$.
        There exists $B_1\in B$ such that $y\in B_1$ and $B_1\subseteq U$ by \cref{open_in,topology_generated_by_basis}.
        Take $x,\epsilon$ such that $x\in X$ and $\epsilon\in\reals$ and $0 < \epsilon$ and
            $B_1 = \metricBall{d}{X}{x}{\epsilon}$ by \cref{metric_basis}.
        Thus $y\in X$ and $d((x,y)) < \epsilon$ by \cref{metric_ball}.
        $d\in\funs{X\times X}{\reals}$ and $d((x,y))\in\reals$
            by \cref{metric_on,times_tuple_intro,funs_type_apply}.
        Let $\delta = \epsilon - d((x,y))$.
        $0 < \delta$ by \cref{metric_on,times_tuple_intro,funs_type_apply,real_lt_imp_pos_diff}.
        Show that for all $z$ such that $z\in\metricBall{d}{X}{y}{\delta}$ we have $z\in B_1$.
        \begin{subproof}
            Fix $z$ such that $z\in\metricBall{d}{X}{y}{\delta}$.
            Thus $z\in X$ and $d((y,z)) < \delta$ by \cref{metric_ball}.
            $d((x,y)) + d((y,z)) < \epsilon$ by
                \cref{times_tuple_intro,funs_type_apply,real_sub,real_add_inverse,real_add_closed,real_lt_subst_right,real_lt_add,real_add_assoc,real_add_comm,real_add_ident,real_lt_subst_left}.
            $d((x,z)) \leq d((x,y)) + d((y,z))$ by \cref{metric_on,metric_triangle_inequality}.
            Hence $d((x,z)) < \epsilon$ by
                \cref{times_tuple_intro,funs_type_apply,real_add_closed,real_leq_split,real_lt_trans,real_lt_subst_left}.
            Therefore $z\in B_1$ by \cref{metric_ball}.
        \end{subproof}
        Therefore $\metricBall{d}{X}{y}{\delta}\subseteq U$ by \cref{subseteq,subseteq_transitive}.
    \end{subproof}
    Show that if for all $y\in U$ there exists $\delta\in\reals$ such that $0 < \delta$ and $\metricBall{d}{X}{y}{\delta}\subseteq U$
        then for all $y\in U$ there exists $B_1\in B$ such that $y\in B_1$ and $B_1\subseteq U$.
    \begin{subproof}
        Follows by \cref{subseteq,metric_on,metric_zero_characterization,real_lt_subst_left,metric_ball,powerset_intro,metric_basis}.
    \end{subproof}
    Show that if for all $y\in U$ there exists $\delta\in\reals$ such that $0 < \delta$ and $\metricBall{d}{X}{y}{\delta}\subseteq U$
        then $U$ is open in $X$ with respect to $T$.
    \begin{subproof}
        Follows by \cref{metric_basis_is_basis,basis_topology_is_topology,powerset_intro,topology_generated_by_basis,open_in}.
    \end{subproof}
\end{proof}

% from §20 Item 7: discrete metric
\begin{definition}\label{discrete_metric}
    $\discreteMetric{X} = \{w \mid \exists p\in X\times X. \exists z\in\reals.
        w = (p,z) \land
        ((\fst{p} = \snd{p} \land z = 0) \lor (\fst{p}\neq \snd{p} \land z = 1))\}$.
\end{definition}
% from §20 Item 8: discrete metric is a metric
\begin{lemma}\label{discrete_metric_is_metric}
    $\discreteMetric{X}$ is a metric on $X$.
\end{lemma}
\begin{proof}
    Let $d = \discreteMetric{X}$.
    For all $p,z_1,z_2$ such that $(p,z_1)\in d$ and $(p,z_2)\in d$ we have $z_1 = z_2$
        by \cref{discrete_metric,pair_eq_iff}.
    Hence $d$ is right-unique by \cref{rightunique}.
    $d$ is a relation by \cref{discrete_metric,relation}.
    $d$ is a function.
    For all $p$ such that $p\in\dom{d}$ we have $d(p)\in\reals$
        by \cref{dom_iff,discrete_metric,pair_eq_iff,function_apply_intro}.
    Hence $d$ is a function to $\reals$.
    $\dom{d}\subseteq X\times X$ by \cref{dom_iff,discrete_metric,pair_eq_iff,subseteq}.
    For all $p$ such that $p\in X\times X$ we have $p\in\dom{d}$
        by \cref{real_add_ident,real_naturals_subset_reals,real_one_in_naturals,subseteq,discrete_metric,dom_intro}.
    Therefore $d\in\funs{X\times X}{\reals}$ by \cref{subseteq,subseteq_antisymmetric,funs_intro}.
    Show that for all $x,y$ such that $x\in X$ and $y\in X$ we have $d((x,y)) \geq 0$.
    \begin{subproof}
        Follows by \cref{times_tuple_intro,real_add_ident,real_naturals_subset_reals,real_one_in_naturals,subseteq,fst_eq,snd_eq,discrete_metric,function_apply_intro,funs_elim,real_zero_lt_one,real_lt_subst_right}.
    \end{subproof}
    Hence $d$ is nonnegative on $X$ by \cref{metric_nonnegative}.
    Show that for all $x,y$ such that $x\in X$ and $y\in X$ we have $d((x,y)) = 0$ iff $x = y$.
    \begin{subproof}
        Follows by \cref{times_tuple_intro,real_add_ident,real_naturals_subset_reals,real_one_in_naturals,subseteq,fst_eq,snd_eq,discrete_metric,function_apply_intro,funs_elim,real_zero_lt_one,real_lt_subst_right,real_lt_irreflexive}.
    \end{subproof}
    Therefore $d$ is zero-characterized on $X$ by \cref{metric_zero_characterization}.
    Show that for all $x,y$ such that $x\in X$ and $y\in X$ we have $d((x,y)) = d((y,x))$.
    \begin{subproof}
        Follows by \cref{times_tuple_intro,real_add_ident,real_naturals_subset_reals,real_one_in_naturals,subseteq,fst_eq,snd_eq,discrete_metric,function_apply_intro,funs_elim}.
    \end{subproof}
    Hence $d$ is symmetric on $X$ by \cref{metric_symmetric}.
    Show that for all $x,y,z$ such that $x\in X$ and $y\in X$ and $z\in X$ we have
        $d((x,y)) + d((y,z)) \geq d((x,z))$.
    \begin{subproof}
        Fix $x$.
        Fix $y$.
        Fix $z$ such that $x\in X$ and $y\in X$ and $z\in X$.
        \begin{byCase}
        \caseOf{$x = z$.}
            $d((x,z)) = 0$ by \cref{times_tuple_intro,real_add_ident,fst_eq,snd_eq,discrete_metric,function_apply_intro,funs_elim}.
            $0\in\reals$ by \cref{real_add_ident}.
            We have $d((x,y))\in\reals$ and $d((y,z))\in\reals$ by \cref{funs_type_apply,times_tuple_intro}.
            Therefore $d((x,y)) + d((y,z))\in\reals$ by \cref{real_add_closed}.
            We have $0 \leq d((x,y))$ and $0 \leq d((y,z))$ by \cref{metric_nonnegative}.
            Thus $0 + d((y,z)) \leq d((x,y)) + d((y,z))$ by \cref{real_leq_add}.
            Hence $0 \leq d((x,y)) + d((y,z))$ by \cref{real_add_ident,real_leq_trans}.
            $d((x,z)) \leq 0$.
            Therefore $d((x,z)) \leq d((x,y)) + d((y,z))$ by \cref{real_leq_trans}.
            Hence $d((x,y)) + d((y,z)) \geq d((x,z))$.
        \caseOf{$x \neq z$.}
            We have $d((x,z)) = 1$ by
                \cref{times_tuple_intro,real_naturals_subset_reals,real_one_in_naturals,subseteq,fst_eq,snd_eq,discrete_metric,function_apply_intro,funs_elim}.
            Show that $d((x,y)) + d((y,z)) \geq 1$.
            \begin{subproof}
                \begin{byCase}
                \caseOf{$x = y$.}
                    $d((y,z)) = 1$ by
                        \cref{times_tuple_intro,real_naturals_subset_reals,real_one_in_naturals,subseteq,fst_eq,snd_eq,discrete_metric,function_apply_intro,funs_elim}.
                    $d((x,y))\in\reals$ by \cref{funs_type_apply,times_tuple_intro}.
                    Hence $0 \leq d((x,y))$ by \cref{metric_nonnegative}.
                    Therefore $d((x,y)) + d((y,z)) \geq 1$ by
                        \cref{real_add_ident,real_naturals_subset_reals,real_one_in_naturals,subseteq,real_leq_add,real_add_comm,real_leq_trans}.
                \caseOf{$x \neq y$.}
                    $d((x,y)) = 1$ by
                        \cref{times_tuple_intro,real_naturals_subset_reals,real_one_in_naturals,subseteq,fst_eq,snd_eq,discrete_metric,function_apply_intro,funs_elim}.
                    $d((y,z))\in\reals$ by \cref{funs_type_apply,times_tuple_intro}.
                    Hence $0 \leq d((y,z))$ by \cref{metric_nonnegative}.
                    Therefore $d((x,y)) + d((y,z)) \geq 1$ by
                        \cref{real_add_ident,real_naturals_subset_reals,real_one_in_naturals,subseteq,real_leq_add,real_add_comm,real_leq_trans}.
                \end{byCase}
            \end{subproof}
            Hence $1 \leq d((x,y)) + d((y,z))$.
            $d((x,z)) \leq 1$.
            Therefore $d((x,z)) \leq d((x,y)) + d((y,z))$ by \cref{real_leq_trans}.
            Hence $d((x,y)) + d((y,z)) \geq d((x,z))$.
        \end{byCase}
    \end{subproof}
    Hence $d$ is a metric on $X$ by \cref{metric_triangle_inequality,metric_on}.
\end{proof}
% from §20 Item 9: discrete metric topology
\begin{lemma}\label{discrete_metric_topology}
    $\metricTopology{\discreteMetric{X}}{X} = \discrete{X}$.
\end{lemma}
\begin{proof}
    Let $d = \discreteMetric{X}$.
    Let $T = \metricTopology{d}{X}$.
    $d$ is a metric on $X$ by \cref{discrete_metric_is_metric}.
    Thus $\metricBasis{d}{X}$ is a basis on $X$ and $T$ is a topology on $X$
        by \cref{metric_basis_is_basis,basis_topology_is_topology,metric_topology,topology_on}.
    Show that $T\subseteq\discrete{X}$.
    \begin{subproof}
        Follows by \cref{metric_basis_is_basis,basis_topology_is_topology,metric_topology,topology_on,discrete_topology}.
    \end{subproof}
    Show that for all $U$ such that $U\in\discrete{X}$ we have $U\in T$.
    \begin{subproof}
        Fix $U$ such that $U\in\discrete{X}$.
        $U\subseteq X$ by \cref{discrete_topology,powerset_elim}.
        $1\in\reals$ and $0 < 1$ by \cref{real_naturals_subset_reals,real_one_in_naturals,subseteq,real_zero_lt_one}.
        Show that for all $y$ such that $y\in U$ we have $\metricBall{d}{X}{y}{1}\subseteq U$.
        \begin{subproof}
            Follows by \cref{metric_ball,subseteq,times_tuple_intro,fst_eq,snd_eq,discrete_metric,function_apply_intro,funs_elim,metric_on,real_lt_subst_right,real_lt_irreflexive}.
        \end{subproof}
        For all $y\in U$ there exists $\delta\in\reals$ such that
            $0 < \delta$ and $\metricBall{d}{X}{y}{\delta}\subseteq U$.
        Hence $U\in T$ by \cref{discrete_topology,powerset_elim,metric_open_iff,open_in}.
    \end{subproof}
    Therefore $T = \discrete{X}$ by \cref{subseteq,subseteq_antisymmetric}.
\end{proof}

% from §20 Item 10: standard metric on $\reals$
\begin{definition}\label{standard_metric_r}
    $\standardMetricR = \{w \mid \exists p\in\reals\times\reals. \exists z\in\reals.
        w = (p,z) \land z = \abs{\fst{p} - \snd{p}}\}$.
\end{definition}

% from §20 Item 11: standard metric is a metric
\begin{lemma}\label{standard_metric_is_metric}
    $\standardMetricR$ is a metric on $\reals$.
\end{lemma}
\begin{proof}
    Let $d = \standardMetricR$.
    For all $p,z_1,z_2$ such that $(p,z_1)\in d$ and $(p,z_2)\in d$ we have $z_1 = z_2$
        by \cref{standard_metric_r,pair_eq_iff}.
    Hence $d$ is right-unique by \cref{rightunique}.
    $d$ is a relation by \cref{standard_metric_r,relation}.
    $d$ is a function.
    For all $p$ such that $p\in\dom{d}$ we have $d(p)\in\reals$
        by \cref{dom_iff,standard_metric_r,pair_eq_iff,function_apply_intro}.
    Hence $d$ is a function to $\reals$.
    Show that $\dom{d}\subseteq\reals\times\reals$.
    \begin{subproof}
        Follows by \cref{dom_iff,standard_metric_r,pair_eq_iff,subseteq}.
    \end{subproof}
    Show that $\reals\times\reals\subseteq\dom{d}$.
    \begin{subproof}
        Follows by \cref{standard_metric_r,dom_intro,times_elem_is_tuple,fst_eq,snd_eq,real_sub,real_add_inverse,real_add_closed,real_abs_in_reals,subseteq}.
    \end{subproof}
    Hence $d\in\funs{\reals\times\reals}{\reals}$ by \cref{subseteq,subseteq_antisymmetric,funs_intro}.
    Show that for all $x,y$ such that $x\in\reals$ and $y\in\reals$ we have $d((x,y)) \geq 0$.
    \begin{subproof}
        Follows by \cref{times_tuple_intro,standard_metric_r,function_apply_intro,funs_elim,fst_eq,snd_eq,real_sub,real_add_inverse,real_add_closed,real_abs_in_reals,real_abs_nonnegative}.
    \end{subproof}
    Hence $d$ is nonnegative on $\reals$ by \cref{metric_nonnegative}.
    Show that for all $x,y$ such that $x\in\reals$ and $y\in\reals$ we have $d((x,y)) = 0$ iff $x = y$.
    \begin{subproof}
        Follows by \cref{times_tuple_intro,standard_metric_r,function_apply_intro,funs_elim,fst_eq,snd_eq,real_sub,real_add_inverse,real_add_closed,real_abs_in_reals,real_abs_eq_zero,real_add_cancel,real_add_ident,real_abs_nonneg}.
    \end{subproof}
    Therefore $d$ is zero-characterized on $\reals$ by \cref{metric_zero_characterization}.
    Show that for all $x,y$ such that $x\in\reals$ and $y\in\reals$ we have $d((x,y)) = d((y,x))$.
    \begin{subproof}
        Follows by \cref{times_tuple_intro,standard_metric_r,function_apply_intro,funs_elim,fst_eq,snd_eq,real_sub,real_add_inverse,real_add_closed,real_abs_in_reals,real_abs_sub_swap}.
    \end{subproof}
    Hence $d$ is symmetric on $\reals$ by \cref{metric_symmetric}.
    Show that for all $x,y,z$ such that $x\in\reals$ and $y\in\reals$ and $z\in\reals$ we have
        $d((x,y)) + d((y,z)) \geq d((x,z))$.
    \begin{subproof}
        Fix $x$.
        Fix $y$.
        Fix $z$ such that $x\in\reals$ and $y\in\reals$ and $z\in\reals$.
        $\neg{y}\in\reals$ and $\neg{z}\in\reals$ by \cref{real_add_inverse}.
        $x - y = x + \neg{y}$ by \cref{real_sub}.
        $y - z = y + \neg{z}$ by \cref{real_sub}.
        $x - z = x + \neg{z}$ by \cref{real_sub}.
        $\abs{(x - y)}\in\reals$ and $\abs{(y - z)}\in\reals$ and $\abs{(x - z)}\in\reals$
            by \cref{real_sub,real_add_inverse,real_add_closed,real_abs_in_reals}.
        Let $p_1 = (x,y)$.
        Let $p_2 = (y,z)$.
        Let $p_3 = (x,z)$.
        $p_1\in\reals\times\reals$ and $p_2\in\reals\times\reals$ and $p_3\in\reals\times\reals$
            by \cref{times_tuple_intro}.
        $\fst{p_1} = x$ and $\snd{p_1} = y$ and $\fst{p_2} = y$ and $\snd{p_2} = z$ and
            $\fst{p_3} = x$ and $\snd{p_3} = z$ by \cref{fst_eq,snd_eq}.
        $\fst{p_1} - \snd{p_1} = x - y$.
        $\fst{p_2} - \snd{p_2} = y - z$.
        $\fst{p_3} - \snd{p_3} = x - z$.
        $\fst{p_1} - \snd{p_1}\in\reals$.
        $\fst{p_2} - \snd{p_2}\in\reals$.
        $\fst{p_3} - \snd{p_3}\in\reals$.
        Let $z_1 = \abs{\fst{p_1} - \snd{p_1}}$.
        Let $z_2 = \abs{\fst{p_2} - \snd{p_2}}$.
        Let $z_3 = \abs{\fst{p_3} - \snd{p_3}}$.
        $z_1\in\reals$ and $z_2\in\reals$ and $z_3\in\reals$
            by \cref{fst_eq,snd_eq,real_sub,real_add_inverse,real_add_closed,real_abs_in_reals}.
        Let $w_1 = (p_1,z_1)$.
        Let $w_2 = (p_2,z_2)$.
        Let $w_3 = (p_3,z_3)$.
        $w_1\in d$ and $w_2\in d$ and $w_3\in d$ by \cref{standard_metric_r}.
        $d((x,y)) = \abs{(x - y)}$ and $d((y,z)) = \abs{(y - z)}$ and $d((x,z)) = \abs{(x - z)}$
            by \cref{function_apply_intro,funs_elim,fst_eq,snd_eq}.
        $x - y = x + \neg{y}$ and $y - z = y + \neg{z}$ and
            $(x - y) + (y - z) = (x + \neg{y}) + (y + \neg{z})$ by \cref{real_sub}.
        $(x + \neg{y}) + (y + \neg{z}) = x + (\neg{y} + (y + \neg{z}))$ and
            $\neg{y} + (y + \neg{z}) = (\neg{y} + y) + \neg{z}$ by \cref{real_add_assoc}.
        $\neg{y} + y = y + \neg{y}$ by \cref{real_add_comm}.
        $y + \neg{y} = 0$ by \cref{real_add_inverse}.
        $\neg{y} + (y + \neg{z}) = 0 + \neg{z}$.
        $0 + \neg{z} = \neg{z}$ by \cref{real_add_ident}.
        $(x - y) + (y - z) = x + \neg{z}$.
        $x + \neg{z} = x - z$ by \cref{real_sub}.
        $\abs{((x - y) + (y - z))} = \abs{(x - z)}$.
        Hence $d((x,z)) \leq \abs{(x - y)} + \abs{(y - z)}$ by \cref{real_triangle_inequality}.
        Therefore $d((x,z)) \leq d((x,y)) + d((y,z))$.
        Hence $d((x,y)) + d((y,z)) \geq d((x,z))$.
    \end{subproof}
    Hence $d$ is a metric on $\reals$ by \cref{metric_triangle_inequality,metric_on}.
\end{proof}

% from §20 helper lemma 6: evaluation of the standard metric
\begin{lemma}\label{standard_metric_apply}
    Assume $x\in\reals$ and $y\in\reals$ and $d = \standardMetricR$.
    Then $d((x,y)) = \abs{(x - y)}$.
\end{lemma}
\begin{proof}
    $d((x,y)) = \abs{(x - y)}$ by
        \cref{fst_eq,snd_eq,real_sub,real_add_inverse,real_add_closed,real_abs_in_reals,times_tuple_intro,standard_metric_r,standard_metric_is_metric,metric_on,function_apply_intro,funs_elim}.
\end{proof}

% from §20 Item 12: standard metric topology
\begin{lemma}\label{standard_metric_topology}
    $\metricTopology{\standardMetricR}{\reals} = \standardTopologyR$.
\end{lemma}
\begin{proof}
    Let $d = \standardMetricR$.
    Let $T = \metricTopology{d}{\reals}$.
    Let $T_1 = \standardTopologyR$.
    Let $B = \metricBasis{d}{\reals}$.
    Let $B_1 = \standardBasisR$.
    $d$ is a metric on $\reals$ by \cref{standard_metric_is_metric}.
    Hence $B$ is a basis for $T$ on $\reals$ by \cref{metric_basis_is_basis,basis_for_topology,metric_topology}.
    $B_1$ is a basis for $T_1$ on $\reals$ by \cref{standard_basis_is_basis,basis_for_topology,standard_topology_reals}.
    Show that for all $x,B_2$ such that $x\in\reals$ and $B_2\in B_1$ and $x\in B_2$
        there exists $B_3\in B$ such that $x\in B_3$ and $B_3\subseteq B_2$.
    \begin{subproof}
            Fix $x$.
            Fix $B_2$ such that $x\in\reals$ and $B_2\in B_1$ and $x\in B_2$.
            Take $a,b\in\reals$ such that $a < b$ and $B_2 = \intervalopen{a}{b}$ by \cref{standard_basis_reals}.
            $a < x$ and $x < b$ by \cref{intervalopen}.
            Let $r_1 = x - a$.
            Let $r_2 = b - x$.
            $0 < r_1$ and $0 < r_2$ by \cref{real_lt_imp_pos_diff}.
            Let $A = \{r_1, r_2\}$.
            $A$ is finite and $A\neq\emptyset$ by \cref{pair_is_finite,upair_intro_left,emptyset}.
            For all $y$ such that $y\in A$ we have $y\in\reals$.
            $A\subseteq\reals$ by \cref{subseteq}.
            Take $m\in A$ such that $m$ is a minimum of $A$ by \cref{finite_real_minimum}.
            $0 < m$ by \cref{upair_elim,real_lt_subst_right}.
            Let $B_3 = \metricBall{d}{\reals}{x}{m}$.
            Hence $B_3\in B$ and $x\in B_3$ by
                \cref{metric_ball,subseteq,powerset_intro,metric_basis,metric_on,metric_zero_characterization,real_lt_subst_left}.
            Show that for all $z$ such that $z\in B_3$ we have $z\in B_2$.
            \begin{subproof}
                    Fix $z$ such that $z\in B_3$.
                    $z\in\reals$ and $d((x,z)) < m$ by \cref{metric_ball}.
                    $\abs{(x - z)} < m$ by \cref{standard_metric_apply,real_lt_subst_left}.
                    $x - z\in\reals$ by \cref{real_sub,real_add_inverse,real_add_closed}.
                    $\abs{(x - z)}\in\reals$ by \cref{real_sub,real_add_inverse,real_add_closed,real_abs_in_reals}.
                    $r_1\in A$ by \cref{upair_intro_left}.
                    $m < r_1$ or $m = r_1$ by \cref{real_minimum,subseteq,real_leq_split}.
                    $\abs{(x - z)} < r_1$ by \cref{real_lt_subst_right,real_lt_trans}.
                    $x - z \leq \abs{(x - z)}$ by \cref{real_leq_abs}.
                    $x - z < \abs{(x - z)}$ or $x - z = \abs{(x - z)}$ by \cref{real_leq_split}.
                    $x - z < r_1$ by \cref{real_lt_subst_left,real_lt_trans}.
                        $x - z < x - a$ by \cref{real_lt_subst_right}.
                        $x - a\in\reals$ by \cref{real_sub,real_add_closed}.
                        $\neg{(x - a)} < \neg{(x - z)}$ by \cref{real_lt_neg}.
                        $\neg{(x - a)} = a - x$ and $\neg{(x - z)} = z - x$ by \cref{real_sub_swap_neg}.
                        $a - x < z - x$ by \cref{real_lt_subst_left,real_lt_subst_right}.
                        $\neg{x}\in\reals$ by \cref{real_add_inverse}.
                        $a - x = a + \neg{x}$ and $z - x = z + \neg{x}$ by \cref{real_sub}.
                        $a - x\in\reals$ by \cref{real_add_closed}.
                        $z - x\in\reals$ by \cref{real_add_closed}.
                        $(a - x) + x < (z - x) + x$ by \cref{real_lt_add}.
                        $\neg{x} + x = 0$ by \cref{real_add_comm,real_add_inverse}.
                        $(a - x) + x = a + 0$ and $(z - x) + x = z + 0$
                            by \cref{real_add_assoc,real_add_eq_subst}.
                        $a + 0 = a$ and $z + 0 = z$ by \cref{real_add_ident,real_add_comm}.
                        $(a - x) + x = a$ and $(z - x) + x = z$.
                        $a < z$.
                    $r_2\in A$ by \cref{upair_intro_right}.
                    $m < r_2$ or $m = r_2$ by \cref{real_minimum,subseteq,real_leq_split}.
                    $\abs{(x - z)} < r_2$ by \cref{real_lt_subst_right,real_lt_trans}.
                        $\neg{(x - z)}\in\reals$ by \cref{real_add_inverse}.
                        $\neg{(x - z)} \leq \abs{(x - z)}$ by \cref{real_neg_leq_abs}.
                        $\neg{(x - z)} < \abs{(x - z)}$ or $\neg{(x - z)} = \abs{(x - z)}$ by \cref{real_leq_split}.
                        $\neg{(x - z)} < r_2$ by \cref{real_lt_subst_left,real_lt_trans}.
                        $\neg{(x - z)} = z - x$ by \cref{real_sub_swap_neg}.
                        $z - x < r_2$ by \cref{real_lt_subst_left}.
                        $z - x < b - x$ by \cref{real_lt_subst_right}.
                        $z - x\in\reals$ and $b - x\in\reals$ by \cref{real_sub,real_add_closed}.
                        $(z - x) + x < (b - x) + x$ by \cref{real_lt_add}.
                        $z - x = z + \neg{x}$ and $b - x = b + \neg{x}$ by \cref{real_sub}.
                        $\neg{x} + x = 0$ by \cref{real_add_comm,real_add_inverse}.
                        $(z - x) + x = z + 0$ and $(b - x) + x = b + 0$
                            by \cref{real_add_assoc,real_add_eq_subst}.
                        $z + 0 = z$ and $b + 0 = b$ by \cref{real_add_ident,real_add_comm}.
                        $(z - x) + x = z$ and $(b - x) + x = b$.
                        $z < b$.
                    $z\in B_2$ by \cref{intervalopen}.
                \end{subproof}
            Therefore $B_3\subseteq B_2$ by \cref{subseteq}.
    \end{subproof}
    Hence $T$ is finer than $T_1$ on $\reals$ by \cref{basis_finer_criterion}.
    Show that for all $x,B_2$ such that $x\in\reals$ and $B_2\in B$ and $x\in B_2$
        there exists $B_3\in B_1$ such that $x\in B_3$ and $B_3\subseteq B_2$.
    \begin{subproof}
            Fix $x$.
            Fix $B_2$ such that $x\in\reals$ and $B_2\in B$ and $x\in B_2$.
            Take $x_0,\epsilon\in\reals$ such that $0 < \epsilon$ and $B_2 = \metricBall{d}{\reals}{x_0}{\epsilon}$
                by \cref{metric_basis}.
            $d((x_0,x)) < \epsilon$ by \cref{metric_ball}.
            $d\in\funs{\reals\times\reals}{\reals}$ by \cref{metric_on}.
            $d((x_0,x))\in\reals$ by \cref{times_tuple_intro,funs_type_apply}.
            $\neg{d((x_0,x))}\in\reals$ by \cref{real_add_inverse}.
            Let $\delta = \epsilon - d((x_0,x))$.
            $\delta\in\reals$ by \cref{real_sub,real_add_inverse,real_add_closed}.
            $0 < \delta$ by \cref{real_lt_imp_pos_diff}.
            $\neg{\delta}\in\reals$ by \cref{real_add_inverse}.
            $x - \delta = x + \neg{\delta}$ by \cref{real_sub}.
            $0\in\reals$ by \cref{real_add_ident}.
            $\neg{\delta} < \neg{0}$ by \cref{real_lt_neg}.
            $\neg{0} = 0$ by \cref{real_add_inverse_unique,real_add_ident}.
            $\neg{\delta} < 0$ by \cref{real_lt_subst_right}.
            $\neg{\delta} + x < 0 + x$ by \cref{real_lt_add}.
            $x + \neg{\delta} < 0 + x$ by \cref{real_add_comm,real_lt_subst_left}.
            $x - \delta < 0 + x$ by \cref{real_lt_subst_left,real_sub}.
            $0 + x = x$ by \cref{real_add_ident}.
            $x - \delta < x$ by \cref{real_lt_subst_right}.
            $0 + x < \delta + x$ by \cref{real_lt_add}.
            $\delta + x = x + \delta$ by \cref{real_add_comm}.
            $x < x + \delta$ by \cref{real_lt_subst_left,real_lt_subst_right}.
            Let $B_3 = \intervalopen{x - \delta}{x + \delta}$.
            We have $x - \delta\in\reals$ and $x + \delta\in\reals$ by \cref{real_sub,real_add_inverse,real_add_closed}.
            Therefore $x - \delta < x + \delta$ by \cref{real_lt_trans}.
            Hence $B_3\in B_1$ by \cref{intervalopen_subset_reals,powerset_intro,standard_basis_reals}.
            $x\in B_3$ by \cref{intervalopen}.
            Show that for all $z$ such that $z\in B_3$ we have $z\in B_2$.
                \begin{subproof}
                    Fix $z$ such that $z\in B_3$.
                    $x - \delta < z$ and $z < x + \delta$ and $z\in\reals$ by \cref{intervalopen}.
                    $x - z\in\reals$ and $z - x\in\reals$ by \cref{real_sub,real_add_inverse,real_add_closed}.
                    $\neg{x}\in\reals$ by \cref{real_add_inverse}.
                    $z + \neg{x} < (x + \delta) + \neg{x}$ by \cref{real_lt_add}.
                    $z + \neg{x} = z - x$ by \cref{real_sub}.
                    $(x + \delta) + \neg{x} = x + (\delta + \neg{x})$ by \cref{real_add_assoc}.
                    $\delta + \neg{x} = \neg{x} + \delta$ by \cref{real_add_comm}.
                    $x + (\neg{x} + \delta) = (x + \neg{x}) + \delta$ by \cref{real_add_assoc}.
                    $x + \neg{x} = 0$ by \cref{real_add_inverse}.
                    $(x + \delta) + \neg{x} = 0 + \delta$ by \cref{real_add_assoc,real_add_comm,real_add_inverse}.
                    $0 + \delta = \delta$ by \cref{real_add_ident}.
                    $z - x < \delta$ by \cref{real_lt_subst_left,real_lt_subst_right}.
                    $(x - \delta) + \neg{x} < z + \neg{x}$ by \cref{real_lt_add}.
                    $(x - \delta) + \neg{x} = \neg{\delta}$ by \cref{real_sub,real_add_assoc,real_add_inverse,real_add_comm,real_add_ident}.
                    $\neg{\delta} < z - x$ by \cref{real_lt_subst_left,real_lt_subst_right}.
                    $\neg{(z - x)} < \neg{(\neg{\delta})}$ by \cref{real_lt_neg}.
                    $\neg{(z - x)} = x - z$ by \cref{real_sub_swap_neg}.
                    $\delta = \neg{(\neg{\delta})}$ by \cref{real_neg_neg}.
                    $x - z < \delta$ by \cref{real_lt_subst_left,real_lt_subst_right}.
                    $x - z \leq \delta$ and $\neg{(x - z)} \leq \delta$.
                    $\abs{(x - z)} \leq \delta$ by \cref{real_abs_leq_if_pm_leq}.
                    $\abs{(x - z)}\in\reals$ by \cref{real_abs_in_reals}.
                    $\abs{(x - z)} < \delta$.
                    $d((x,z)) = \abs{(x - z)}$ by \cref{standard_metric_apply}.
                    $d((x,z)) < \delta$.
                    $d((x_0,z)) \leq d((x_0,x)) + d((x,z))$ by \cref{metric_on,metric_triangle_inequality}.
                    $d((x_0,x)) + d((x,z))\in\reals$ and $d((x_0,x)) + \delta\in\reals$
                        by \cref{funs_type_apply,real_add_closed}.
                    $d((x,z)) + d((x_0,x)) < \delta + d((x_0,x))$ by \cref{real_lt_add}.
                    $d((x,z)) + d((x_0,x)) = d((x_0,x)) + d((x,z))$ and
                        $\delta + d((x_0,x)) = d((x_0,x)) + \delta$ by \cref{real_add_comm}.
                    $d((x_0,x)) + d((x,z)) < d((x_0,x)) + \delta$
                        by \cref{real_lt_subst_left,real_lt_subst_right}.
                    $(x_0,z)\in\reals\times\reals$ by \cref{times_tuple_intro}.
                    $d((x_0,z))\in\reals$ by \cref{funs_type_apply}.
                    $d((x_0,z)) < d((x_0,x)) + d((x,z))$ or
                        $d((x_0,z)) = d((x_0,x)) + d((x,z))$
                        by \cref{times_tuple_intro,funs_type_apply,real_leq_split}.
                    $d((x_0,z)) < d((x_0,x)) + \delta$ by \cref{real_lt_subst_left,real_lt_trans}.
                    $\delta = \epsilon + \neg{d((x_0,x))}$ by \cref{real_sub}.
                    $d((x_0,x)) + \delta = \delta + d((x_0,x))$ by \cref{real_add_comm}.
                    $d((x_0,x)) + \delta = (\epsilon + \neg{d((x_0,x))}) + d((x_0,x))$
                        by \cref{real_add_eq_subst}.
                    $(\epsilon + \neg{d((x_0,x))}) + d((x_0,x)) =
                        \epsilon + (\neg{d((x_0,x))} + d((x_0,x)))$ by \cref{real_add_assoc}.
                    $\neg{d((x_0,x))} + d((x_0,x)) = 0$ by \cref{real_add_comm,real_add_inverse}.
                    $d((x_0,x)) + \delta = \epsilon + 0$ by \cref{real_add_eq_subst}.
                    $\epsilon + 0 = \epsilon$ by \cref{real_add_ident,real_add_comm}.
                    $d((x_0,x)) + \delta = \epsilon$.
                    $d((x_0,z)) < \epsilon$ by \cref{real_lt_subst_right}.
                    $z\in B_2$ by \cref{metric_ball}.
                \end{subproof}
            Therefore $B_3\subseteq B_2$ by \cref{subseteq}.
    \end{subproof}
    Hence $T_1$ is finer than $T$ on $\reals$ by \cref{basis_finer_criterion}.
    Therefore $T = T_1$ by \cref{finer_topology,subseteq_antisymmetric}.
\end{proof}

% from §20 Item 13: metrizable space
\begin{definition}\label{metrizable}
    $X$ is metrizable with respect to $T$ iff
    $T$ is a topology on $X$ and
    there exists $d$ such that $d$ is a metric on $X$ and $T = \metricTopology{d}{X}$.
\end{definition}

% from §20 Item 14: metric space
\begin{definition}\label{metric_space}
    $X$ is a metric space with metric $d$ and topology $T$ iff
    $d$ is a metric on $X$ and $T = \metricTopology{d}{X}$.
\end{definition}

% from §20 Item 15: bounded subset
\begin{definition}\label{metric_bounded}
    $A$ is bounded in $X$ with respect to $d$ iff
    $d$ is a metric on $X$ and
    $A\subseteq X$ and
    there exists $M\in\reals$ such that for all $a_1,a_2\in A$ we have $d((a_1,a_2)) \leq M$.
\end{definition}

% from §20 helper lemma 5: distance set of a subset
\begin{definition}\label{metric_distance_set}
    $\metricDistanceSet{A}{d} = \{d((a_1,a_2)) \mid a_1\in A, a_2\in A\}$.
\end{definition}

% from §20 Item 16: diameter of a set
\begin{definition}\label{metric_diameter}
    $r$ is a diameter of $A$ in $X$ with respect to $d$ iff
    $A$ is bounded in $X$ with respect to $d$ and
    $A\neq\emptyset$ and
    $r$ is a supremum of $\metricDistanceSet{A}{d}$.
\end{definition}

% from §20 Item 17: standard bounded metric
\begin{definition}\label{bounded_metric}
    $\boundedMetric{d} = \{w \mid \exists p\in\dom{d}. \exists z\in\reals.
        w = (p,z) \land ((z = d(p) \land d(p) < 1) \lor (z = 1 \land 1 \leq d(p)))\}$.
\end{definition}

% from §20 helper lemma 7: bounded metric values are at most one
\begin{lemma}\label{bounded_metric_leq_one}
    Assume $d$ is a metric on $X$ and $e = \boundedMetric{d}$ and $e\in\funs{X\times X}{\reals}$ and $p\in X\times X$.
    Then $e(p) \leq 1$.
\end{lemma}
\begin{proof}
    $p\in\dom{d}$ and $d(p)\in\reals$ by \cref{metric_on,funs_elim,subseteq,funs_type_apply}.
    $1\in\reals$ by \cref{real_naturals_subset_reals,real_one_in_naturals,subseteq}.
    \begin{byCase}
    \caseOf{$d(p) < 1$.}
        Let $w = (p,d(p))$.
        $w\in e$ by \cref{bounded_metric}.
        Thus $e(p) \leq 1$ by \cref{function_apply_intro,funs_elim,real_lt_subst_left}.
    \caseOf{it is wrong that $d(p) < 1$.}
        We have $1 \leq d(p)$.
        Let $w = (p,1)$.
        $w\in e$ by \cref{bounded_metric}.
        Thus $e(p) \leq 1$ by \cref{function_apply_intro,funs_elim}.
    \end{byCase}
\end{proof}

% from §20 helper lemma 8: bounded metric agrees with the original metric below one
\begin{lemma}\label{bounded_metric_apply_lt_one}
    Assume $d$ is a metric on $X$ and $e = \boundedMetric{d}$ and $e\in\funs{X\times X}{\reals}$ and $p\in X\times X$ and $d(p) < 1$.
    Then $e(p) = d(p)$.
\end{lemma}
\begin{proof}
    $e(p) = d(p)$ by \cref{metric_on,funs_elim,subseteq,funs_type_apply,bounded_metric,function_apply_intro}.
\end{proof}

% from §20 helper lemma 9: bounded metric is one above the cutoff
\begin{lemma}\label{bounded_metric_apply_geq_one}
    Assume $d$ is a metric on $X$ and $e = \boundedMetric{d}$ and $e\in\funs{X\times X}{\reals}$ and $p\in X\times X$ and it is wrong that $d(p) < 1$.
    Then $e(p) = 1$.
\end{lemma}
\begin{proof}
    $p\in\dom{d}$ and $d(p)\in\reals$ by \cref{metric_on,funs_elim,subseteq,funs_type_apply}.
    $1\in\reals$ by \cref{real_naturals_subset_reals,real_one_in_naturals,subseteq}.
    $1 \leq d(p)$ by \cref{real_lt_trichotomy}.
    Hence $e(p) = 1$ by \cref{metric_on,funs_elim,subseteq,bounded_metric,function_apply_intro}.
\end{proof}

% from §20 Item 18: standard bounded metric is a metric
\begin{theorem}\label{bounded_metric_is_metric}
    Assume $d$ is a metric on $X$.
    Then $\boundedMetric{d}$ is a metric on $X$.
\end{theorem}
\begin{proof}
    Let $e = \boundedMetric{d}$.
    Show that $e\in\funs{X\times X}{\reals}$.
    \begin{subproof}
        Show that for all $p,z_1,z_2$ such that $(p,z_1)\in e$ and $(p,z_2)\in e$ we have $z_1 = z_2$.
        \begin{subproof}
                    Fix $p$.
                    Fix $z_1$.
                    Fix $z_2$ such that $(p,z_1)\in e$ and $(p,z_2)\in e$.
                    Take $p_1, w_1$ such that
                        $p_1\in\dom{d}$ and $w_1\in\reals$ and
                        $(p,z_1) = (p_1,w_1)$ and
                        $((w_1 = d(p_1) \land d(p_1) < 1) \lor (w_1 = 1 \land 1 \leq d(p_1)))$
                        by \cref{bounded_metric}.
                    Take $p_2, w_2$ such that
                        $p_2\in\dom{d}$ and $w_2\in\reals$ and
                        $(p,z_2) = (p_2,w_2)$ and
                        $((w_2 = d(p_2) \land d(p_2) < 1) \lor (w_2 = 1 \land 1 \leq d(p_2)))$
                        by \cref{bounded_metric}.
                    Hence $((z_1 = d(p) \land d(p) < 1) \lor (z_1 = 1 \land 1 \leq d(p)))$
                        and $((z_2 = d(p) \land d(p) < 1) \lor (z_2 = 1 \land 1 \leq d(p)))$
                        by \cref{bounded_metric,pair_eq_iff}.
                    $d(p)\in\reals$ by \cref{pair_eq_iff,metric_on,funs_elim}.
                    $1\in\reals$ by \cref{real_naturals_subset_reals,real_one_in_naturals,subseteq}.
                    \begin{byCase}
                    \caseOf{$d(p) < 1$.}
                        We have $z_1 = z_2$.
                    \caseOf{it is wrong that $d(p) < 1$.}
                        We have $z_1 = z_2$.
                    \end{byCase}
                \end{subproof}
                Hence $e$ is a function by \cref{bounded_metric,relation,rightunique}.
        For all $p\in\dom{e}$ we have $e(p)\in\reals$
            by \cref{dom_iff,bounded_metric,pair_eq_iff,function_apply_intro}.
        Hence $e$ is a function to $\reals$.
        Show that $\dom{e}\subseteq X\times X$.
        \begin{subproof}
            Follows by \cref{dom_iff,bounded_metric,pair_eq_iff,metric_on,funs_elim,subseteq}.
        \end{subproof}
        Show that $X\times X\subseteq\dom{e}$.
        \begin{subproof}
            Follows by \cref{metric_on,funs_elim,subseteq,funs_type_apply,real_naturals_subset_reals,real_one_in_naturals,bounded_metric,dom_intro,real_lt_trichotomy}.
        \end{subproof}
        Hence $e\in\funs{X\times X}{\reals}$ by \cref{subseteq,subseteq_antisymmetric,funs_intro}.
    \end{subproof}
    Show that for all $x,y$ such that $x\in X$ and $y\in X$ we have $e((x,y)) \geq 0$.
    \begin{subproof}
        Fix $x$.
        Fix $y$ such that $x\in X$ and $y\in X$.
        Let $p = (x,y)$.
        $p\in X\times X$ by \cref{times_tuple_intro}.
        \begin{byCase}
        \caseOf{$d(p) < 1$.}
            We have $e(p) = d(p)$ by \cref{bounded_metric_apply_lt_one}.
            Hence $e((x,y)) \geq 0$ by \cref{metric_on,metric_nonnegative}.
        \caseOf{it is wrong that $d(p) < 1$.}
            We have $e(p) = 1$ by \cref{bounded_metric_apply_geq_one}.
            Hence $e((x,y)) \geq 0$ by
                \cref{real_naturals_subset_reals,real_one_in_naturals,subseteq,real_zero_lt_one,real_lt_subst_left}.
        \end{byCase}
    \end{subproof}
    Hence $e$ is nonnegative on $X$ by \cref{metric_nonnegative}.
    Show that for all $x,y$ such that $x\in X$ and $y\in X$ we have $e((x,y)) = 0$ iff $x = y$.
    \begin{subproof}
        Fix $x$.
        Fix $y$ such that $x\in X$ and $y\in X$.
        Show that if $e((x,y)) = 0$ then $x = y$.
        \begin{subproof}
            Suppose not.
            $e((x,y)) = 0$ and $x \neq y$.
            Let $p = (x,y)$.
            $p\in X\times X$ by \cref{times_tuple_intro}.
            Hence $e(p) = 0$.
            \begin{byCase}
            \caseOf{$d(p) < 1$.}
                $e(p) = d(p)$ by \cref{bounded_metric_apply_lt_one}.
                Hence $d(p) = 0$.
                Therefore $x = y$ by \cref{metric_on,metric_zero_characterization}.
            \caseOf{it is wrong that $d(p) < 1$.}
                $e(p) = 1$ by \cref{bounded_metric_apply_geq_one}.
                Hence $0 = 1$.
                $1\in\reals$ by \cref{real_naturals_subset_reals,real_one_in_naturals,subseteq}.
                Contradiction by \cref{real_zero_lt_one,real_lt_subst_right,real_lt_irreflexive}.
            \end{byCase}
        \end{subproof}
        Show that if $x = y$ then $e((x,y)) = 0$.
        \begin{subproof}
            Follows by \cref{times_tuple_intro,metric_on,metric_zero_characterization,real_zero_lt_one,real_lt_subst_left,funs_elim,subseteq,funs_type_apply,bounded_metric,function_apply_intro}.
        \end{subproof}
    \end{subproof}
    Therefore $e$ is zero-characterized on $X$ by \cref{metric_zero_characterization}.
    Show that for all $x,y$ such that $x\in X$ and $y\in X$ we have $e((x,y)) = e((y,x))$.
    \begin{subproof}
        Fix $x$.
        Fix $y$ such that $x\in X$ and $y\in X$.
        Let $p_1 = (x,y)$.
        Let $p_2 = (y,x)$.
        $d((x,y)) = d((y,x))$ by \cref{times_tuple_intro,metric_on,metric_symmetric}.
        $p_1\in X\times X$ and $p_2\in X\times X$ and $d(p_1)\in\reals$ and $d(p_2)\in\reals$
            by \cref{times_tuple_intro,metric_on,funs_type_apply}.
        $1\in\reals$ by \cref{real_naturals_subset_reals,real_one_in_naturals,subseteq}.
        \begin{byCase}
        \caseOf{$d(p_1) < 1$.}
            We have $d(p_2) < 1$ by \cref{real_lt_subst_left}.
            $e(p_1) = d(p_1)$ by \cref{bounded_metric_apply_lt_one}.
            $e(p_2) = d(p_2)$ by \cref{bounded_metric_apply_lt_one}.
            Therefore $e((x,y)) = e((y,x))$.
        \caseOf{it is wrong that $d(p_1) < 1$.}
            $1 \leq d(p_1)$.
            $1 \leq d(p_2)$ by \cref{real_lt_subst_left}.
            $e(p_1) = 1$ by \cref{bounded_metric_apply_geq_one}.
            $e(p_2) = 1$ by \cref{bounded_metric_apply_geq_one}.
            Therefore $e((x,y)) = e((y,x))$.
        \end{byCase}
    \end{subproof}
    Hence $e$ is symmetric on $X$ by \cref{metric_symmetric}.
    Show that for all $x,y,z$ such that $x\in X$ and $y\in X$ and $z\in X$ we have
        $e((x,y)) + e((y,z)) \geq e((x,z))$.
    \begin{subproof}
        Fix $x$.
        Fix $y$.
        Fix $z$ such that $x\in X$ and $y\in X$ and $z\in X$.
        Let $p_1 = (x,y)$.
        Let $p_2 = (y,z)$.
        Let $p_3 = (x,z)$.
        $p_1\in X\times X$ and $p_2\in X\times X$ and $p_3\in X\times X$ by \cref{times_tuple_intro}.
        $d(p_1)\in\reals$ and $d(p_2)\in\reals$ and $d(p_3)\in\reals$ by \cref{metric_on,funs_type_apply}.
        $1\in\reals$ by \cref{real_naturals_subset_reals,real_one_in_naturals,subseteq}.
        $d(p_3) \leq d(p_1) + d(p_2)$ by \cref{metric_on,metric_triangle_inequality}.
        \begin{byCase}
        \caseOf{$d(p_1) < 1$.}
            \begin{byCase}
            \caseOf{$d(p_2) < 1$.}
                $e(p_1) = d(p_1)$ by \cref{bounded_metric_apply_lt_one}.
                $e(p_2) = d(p_2)$ by \cref{bounded_metric_apply_lt_one}.
                $e(p_3) \leq d(p_3)$ by
                    \cref{bounded_metric_apply_lt_one,bounded_metric_apply_geq_one,real_lt_trichotomy}.
                We have $e(p_3)\in\reals$ by \cref{funs_elim,funs_type_apply}.
                $d(p_1) + d(p_2)\in\reals$ by \cref{real_add_closed}.
                Therefore $e(p_3) \leq d(p_1) + d(p_2)$ by \cref{real_leq_trans}.
                $e(p_1)\in\reals$ and $e(p_2)\in\reals$ by \cref{funs_elim}.
                Hence $e(p_1) + e(p_2) = d(p_1) + e(p_2)$ and
                    $d(p_1) + e(p_2) = d(p_1) + d(p_2)$ by \cref{real_add_eq_subst}.
                Therefore $e(p_1) + e(p_2) = d(p_1) + d(p_2)$.
                $d(p_1) + d(p_2) \leq e(p_1) + e(p_2)$.
                Hence $e(p_3) \leq e(p_1) + e(p_2)$ by \cref{real_leq_trans}.
            \caseOf{it is wrong that $d(p_2) < 1$.}
                $1 \leq d(p_2)$.
                $e(p_1) = d(p_1)$ by \cref{bounded_metric_apply_lt_one}.
                $e(p_2) = 1$ by \cref{bounded_metric_apply_geq_one}.
                We have $e(p_3) \leq 1$ by \cref{bounded_metric_leq_one}.
                $0 \leq e(p_1)$ by \cref{metric_nonnegative}.
                $e(p_1)\in\reals$ and $e(p_2)\in\reals$ and $e(p_3)\in\reals$ and
                    $e(p_1) + e(p_2)\in\reals$ by \cref{funs_type_apply,real_add_closed}.
                $0\in\reals$ and $1\in\reals$ by \cref{real_add_ident,real_naturals_subset_reals,real_one_in_naturals,subseteq}.
                $0 + 1 \leq e(p_1) + 1$ by \cref{real_leq_add}.
                $0 + 1 = 1$ and $e(p_1) + 1 = 1 + e(p_1)$ by \cref{real_add_ident,real_add_comm}.
                Therefore $1 \leq e(p_1) + e(p_2)$ and $e(p_3) \leq e(p_1) + e(p_2)$ by \cref{real_leq_trans}.
            \end{byCase}
        \caseOf{it is wrong that $d(p_1) < 1$.}
            $1 \leq d(p_1)$.
            $e(p_1) = 1$ by \cref{bounded_metric_apply_geq_one}.
            We have $e(p_3) \leq 1$ by \cref{bounded_metric_leq_one}.
            $0 \leq e(p_2)$ by \cref{metric_nonnegative}.
            $e(p_1)\in\reals$ and $e(p_2)\in\reals$ and $e(p_3)\in\reals$ and
                $e(p_1) + e(p_2)\in\reals$ by \cref{funs_type_apply,real_add_closed}.
            $0\in\reals$ and $1\in\reals$ by \cref{real_add_ident,real_naturals_subset_reals,real_one_in_naturals,subseteq}.
            $0 + 1 \leq e(p_2) + 1$ by \cref{real_leq_add}.
            $0 + 1 = 1$ and $e(p_2) + 1 = 1 + e(p_2)$ by \cref{real_add_ident,real_add_comm}.
            Therefore $1 \leq e(p_1) + e(p_2)$ and $e(p_3) \leq e(p_1) + e(p_2)$ by \cref{real_leq_trans}.
        \end{byCase}
    \end{subproof}
    Hence $e$ is a metric on $X$ by \cref{metric_triangle_inequality,metric_on}.
\end{proof}
% from §20 helper lemma 6: bounded metric agrees with the original metric below one
\begin{lemma}\label{bounded_metric_apply_below_one}
    Assume $d$ is a metric on $X$.
    Let $e = \boundedMetric{d}$.
    Suppose $p\in X\times X$.
    Suppose $d(p) < 1$.
    Then $e(p) = d(p)$.
\end{lemma}
\begin{proof}
    $e(p) = d(p)$ by \cref{bounded_metric,bounded_metric_is_metric,metric_on,funs_elim,subseteq,funs_type_apply,function_apply_intro}.
\end{proof}

% from §20 helper lemma 7: bounded metric agrees with the original metric when its value is below one
\begin{lemma}\label{bounded_metric_apply_value_lt_one}
    Assume $d$ is a metric on $X$.
    Let $e = \boundedMetric{d}$.
    Suppose $p\in X\times X$.
    Suppose $e(p) < 1$.
    Then $e(p) = d(p)$.
\end{lemma}
\begin{proof}
    $(e(p) = d(p) \land d(p) < 1) \lor (e(p) = 1 \land 1 \leq d(p))$
        by \cref{bounded_metric_is_metric,metric_on,funs_elim,subseteq,function_apply_iff,bounded_metric,pair_eq_iff}.
    Therefore $e(p) = d(p)$.
\end{proof}

% from §20 Item 19: standard bounded metric induces the same topology
\begin{theorem}\label{bounded_metric_topology}
    Assume $d$ is a metric on $X$.
    Then $\metricTopology{\boundedMetric{d}}{X} = \metricTopology{d}{X}$.
\end{theorem}
\begin{proof}
    Let $e = \boundedMetric{d}$.
    $e$ is a metric on $X$ by \cref{bounded_metric_is_metric}.
    Let $T = \metricTopology{d}{X}$.
    Let $T_1 = \metricTopology{e}{X}$.
    Show that for all $U$ such that $U\in T_1$ we have $U\in T$.
    \begin{subproof}
        Fix $U$ such that $U\in T_1$.
        Hence $U$ is open in $X$ with respect to $T_1$ and $U\subseteq X$
            by \cref{metric_basis_is_basis,metric_topology,basis_topology_is_topology,basis_for_topology,open_in,topology_on,subseteq,powerset_elim}.
        Show that for all $y\in U$ there exists $\delta\in\reals$ such that $0 < \delta$ and $\metricBall{d}{X}{y}{\delta}\subseteq U$.
        \begin{subproof}
                Fix $y$ such that $y\in U$.
                Take $\epsilon\in\reals$ such that $0 < \epsilon$ and $\metricBall{e}{X}{y}{\epsilon}\subseteq U$
                    by \cref{metric_open_iff}.
                $1\in\reals$ by \cref{real_naturals_subset_reals,real_one_in_naturals,subseteq}.
                Let $A = \{\epsilon,1\}$.
                $A$ is finite by \cref{pair_is_finite}.
                $\epsilon\in A$ and $1\in A$ by \cref{upair_intro_left,upair_intro_right}.
                Hence $A\neq\emptyset$ by \cref{upair_intro_left,emptyset}.
                Therefore $A\subseteq\reals$ by \cref{upair_elim,subseteq}.
                Take $m\in A$ such that $m$ is a minimum of $A$ by \cref{finite_real_minimum}.
                Hence $0 < m$ by \cref{upair_elim,real_zero_lt_one,real_lt_subst_right}.
                Show that $\metricBall{d}{X}{y}{m}\subseteq U$.
                \begin{subproof}
                    Follows by \cref{metric_ball,subseteq,times_tuple_intro,metric_on,funs_type_apply,real_lt_subst_left,real_minimum,real_leq_split,real_lt_subst_right,real_lt_trans,bounded_metric_apply_below_one}.
                \end{subproof}
        \end{subproof}
        Hence $U\in T$ by \cref{metric_open_iff,open_in}.
    \end{subproof}
    Hence $T_1\subseteq T$ by \cref{subseteq}.
    Show that for all $U$ such that $U\in T$ we have $U\in T_1$.
    \begin{subproof}
        Fix $U$ such that $U\in T$.
        Hence $U$ is open in $X$ with respect to $T$ and $U\subseteq X$
            by \cref{metric_basis_is_basis,metric_topology,basis_topology_is_topology,basis_for_topology,open_in,topology_on,subseteq,powerset_elim}.
        Show that for all $y\in U$ there exists $\delta\in\reals$ such that $0 < \delta$ and $\metricBall{e}{X}{y}{\delta}\subseteq U$.
        \begin{subproof}
                Fix $y$ such that $y\in U$.
                Take $\epsilon\in\reals$ such that $0 < \epsilon$ and $\metricBall{d}{X}{y}{\epsilon}\subseteq U$
                    by \cref{metric_open_iff}.
                $1\in\reals$ by \cref{real_naturals_subset_reals,real_one_in_naturals,subseteq}.
                Let $A = \{\epsilon,1\}$.
                $A$ is finite by \cref{pair_is_finite}.
                $\epsilon\in A$ and $1\in A$ by \cref{upair_intro_left,upair_intro_right}.
                Hence $A\neq\emptyset$ by \cref{upair_intro_left,emptyset}.
                Therefore $A\subseteq\reals$ by \cref{upair_elim,subseteq}.
                Take $m\in A$ such that $m$ is a minimum of $A$ by \cref{finite_real_minimum}.
                Hence $0 < m$ by \cref{upair_elim,real_zero_lt_one,real_lt_subst_right}.
                For all $z$ such that $z\in\metricBall{e}{X}{y}{m}$ we have $z\in U$
                    by \cref{metric_ball,subseteq,times_tuple_intro,metric_on,funs_elim,funs_type_apply,real_lt_subst_left,real_minimum,real_leq_split,real_lt_subst_right,real_lt_trans,bounded_metric_apply_value_lt_one}.
                Therefore $\metricBall{e}{X}{y}{m}\subseteq U$ by \cref{subseteq}.
        \end{subproof}
        Hence $U\in T_1$ by \cref{metric_open_iff,open_in}.
    \end{subproof}
    Therefore $T_1 = T$ by \cref{subseteq,subseteq_antisymmetric}.
\end{proof}

% from §20 Item 20: $\reals^n$
\begin{definition}\label{reals_power}
    $\realsPower{J} = \tuplePower{\reals}{J}$.
\end{definition}

% from §20 Item 21: euclidean norm
\begin{definition}\label{euclidean_norm}
    $r$ is a euclidean norm of $x$ with respect to $n$ iff
    $x\in\realsPower{\initialSegment{n}}$ and
    $r\in\reals$ and
    there exists $a\in\funs{\initialSegment{n}}{\reals}$ such that
        for all $i\in\initialSegment{n}$ we have $a(i) = \exp{x(i)}{1+1}$ and
        there exists $s$ such that $s$ is the sum of $a$ and $r = \sqrt{s}$.
\end{definition}
% from §20 helper lemma 6: sum of a finite sequence predicate
\begin{definition}\label{sum_finite_sequence_pred}
    $\sumFiniteSequence{a}{s}$ iff
    there exists $n\in\naturalsPlus$ such that
    $a\in\funs{\initialSegment{n}}{\reals}$ and
    there exists $t\in\funs{\initialSegment{n}}{\reals}$ such that
    $t(1) = a(1)$ and
    $s = t(n)$ and
    for all $k\in\naturalsPlus$ such that $k + 1 \leq n$ we have
    $t(k + 1) = t(k) + a(k + 1)$.
\end{definition}

% from §20 helper lemma 7: sum of a finite sequence is unique
\begin{lemma}\label{sum_finite_sequence_unique}
    Suppose $\sumFiniteSequence{a}{s_1}$.
    Suppose $\sumFiniteSequence{a}{s_2}$.
    Then $s_1 = s_2$.
\end{lemma}
\begin{proof}
    Take $n_1$ such that $n_1\in\naturalsPlus$ and $a\in\funs{\initialSegment{n_1}}{\reals}$ and
        there exists $t_1\in\funs{\initialSegment{n_1}}{\reals}$ such that
            $t_1(1) = a(1)$ and
            $s_1 = t_1(n_1)$ and
            for all $k\in\naturalsPlus$ such that $k + 1 \leq n_1$ we have $t_1(k + 1) = t_1(k) + a(k + 1)$
        by \cref{sum_finite_sequence_pred}.
    Take $n_2$ such that $n_2\in\naturalsPlus$ and $a\in\funs{\initialSegment{n_2}}{\reals}$ and
        there exists $t_2\in\funs{\initialSegment{n_2}}{\reals}$ such that
            $t_2(1) = a(1)$ and
            $s_2 = t_2(n_2)$ and
            for all $k\in\naturalsPlus$ such that $k + 1 \leq n_2$ we have $t_2(k + 1) = t_2(k) + a(k + 1)$
        by \cref{sum_finite_sequence_pred}.
    Take $t_1$ such that $t_1\in\funs{\initialSegment{n_1}}{\reals}$ and
        $t_1(1) = a(1)$ and
        $s_1 = t_1(n_1)$ and
        for all $k\in\naturalsPlus$ such that $k + 1 \leq n_1$ we have $t_1(k + 1) = t_1(k) + a(k + 1)$.
    Take $t_2$ such that $t_2\in\funs{\initialSegment{n_2}}{\reals}$ and
        $t_2(1) = a(1)$ and
        $s_2 = t_2(n_2)$ and
        for all $k\in\naturalsPlus$ such that $k + 1 \leq n_2$ we have $t_2(k + 1) = t_2(k) + a(k + 1)$.
    Hence $n_1 = n_2$
        by \cref{funs_elim,subseteq_antisymmetric,subseteq,initial_segment_naturalsplus,real_naturals_subset_reals,real_leq_antisym}.
    Let $n = n_1$.
    We have $t_1\in\funs{\initialSegment{n}}{\reals}$ and $t_2\in\funs{\initialSegment{n}}{\reals}$.
    Let $A = \{k\in\naturalsPlus \mid \text{if $k \leq n$ then $t_1(k) = t_2(k)$}\}$.
    $A\subseteq\naturalsPlus$ by \cref{subseteq}.
    Hence $1\in A$ by \cref{real_one_in_naturals,real_one_leq_naturalsplus}.
    Show that for all $k\in A$ we have $k + 1\in A$.
    \begin{subproof}
        Follows by \cref{subseteq,real_naturals_subset_reals,real_add_ident,real_one_in_naturals,real_zero_lt_one,real_lt_add,real_add_comm,real_lt_subst_left,real_add_closed,real_leq_trans,initial_segment_naturalsplus,real_naturals_closed_succ,funs_type_apply,real_add_eq_subst}.
    \end{subproof}
    Therefore for all $k\in\naturalsPlus$ we have $k\in A$ by \cref{real_naturals_induction}.
    For all $k\in\naturalsPlus$ we have if $k \leq n$ then $t_1(k) = t_2(k)$.
    We have $\dom{t_1} = \initialSegment{n}$ and $\dom{t_2} = \initialSegment{n}$ and
        $t_1$ is a function and $t_2$ is a function by \cref{funs_elim}.
    Show that $t_1\subseteq t_2$.
    \begin{subproof}
        Follows by \cref{subseteq,initial_segment_naturalsplus,fun_subseteq}.
    \end{subproof}
    Show that $t_2\subseteq t_1$.
    \begin{subproof}
        Follows by \cref{subseteq,initial_segment_naturalsplus,fun_subseteq}.
    \end{subproof}
    Hence $t_1 = t_2$ by \cref{subseteq_antisymmetric}.
    Therefore $s_1 = s_2$.
\end{proof}

% from §20 helper lemma 8: squared coordinate sequence
\begin{definition}\label{coord_square_sequence}
    $\coordSquareSequence{n}{p} = \{w \mid \exists i\in\initialSegment{n}.
        w = (i, \exp{\apply{\fst{p}}{i} - \apply{\snd{p}}{i}}{1+1})\}$.
\end{definition}

% from §20 Item 22: euclidean metric
\begin{definition}\label{euclidean_metric}
    $\euclideanMetric{n} = \{w \mid \exists p\in\realsPower{\initialSegment{n}}\times\realsPower{\initialSegment{n}}. \exists r\in\reals.
        \exists s\in\reals.
        w = (p,r) \land
        \sumFiniteSequence{\coordSquareSequence{n}{p}}{s} \land
        r = \sqrt{s}\}$.
\end{definition}

% from §20 helper lemma 9: coordinate difference set
\begin{definition}\label{coord_diff_set}
    $\coordDiffSet{n}{x}{y} = \{\abs{x(i) - y(i)} \mid i\in\initialSegment{n}\}$.
\end{definition}

% from §20 Item 23: square metric
\begin{definition}\label{square_metric}
    $\squareMetric{n} = \{w \mid \exists p\in\realsPower{\initialSegment{n}}\times\realsPower{\initialSegment{n}}.
        \exists x\in\realsPower{\initialSegment{n}}. \exists y\in\realsPower{\initialSegment{n}}. \exists r\in\reals.
        w = (p,r) \land p = (x,y) \land
        r\in\coordDiffSet{n}{x}{y} \land
        (\forall t\in\coordDiffSet{n}{x}{y}. t \leq r)\}$.
\end{definition}

% from §20 Item 24: comparison of metric topologies
\begin{lemma}\label{metric_topology_finer_iff}
    Suppose $d$ is a metric on $X$.
    Suppose $d_1$ is a metric on $X$.
    Let $T = \metricTopology{d}{X}$.
    Let $T_1 = \metricTopology{d_1}{X}$.
    Then $T_1$ is finer than $T$ on $X$ iff
    for all $x\in X$ we have
    for all $\epsilon\in\reals$ such that $0 < \epsilon$ there exists $\delta\in\reals$
        such that $0 < \delta$ and
        $\metricBall{d_1}{X}{x}{\delta} \subseteq \metricBall{d}{X}{x}{\epsilon}$.
\end{lemma}
\begin{proof}
    Show that if $T_1$ is finer than $T$ on $X$ then
        for all $x\in X$ we have
        for all $\epsilon\in\reals$ such that $0 < \epsilon$ there exists $\delta\in\reals$
            such that $0 < \delta$ and
            $\metricBall{d_1}{X}{x}{\delta} \subseteq \metricBall{d}{X}{x}{\epsilon}$.
    \begin{subproof}
        Follows by \cref{metric_ball,subseteq,powerset_intro,metric_basis,metric_basis_is_basis,basis_elem_open_in_generated,basis_for_topology,metric_topology,finer_topology,basis_topology_is_topology,open_in,metric_zero_characterization,real_lt_subst_left,metric_on,metric_open_iff}.
    \end{subproof}
    Show that if
        for all $x\in X$ we have
        for all $\epsilon\in\reals$ such that $0 < \epsilon$ there exists $\delta\in\reals$
            such that $0 < \delta$ and
            $\metricBall{d_1}{X}{x}{\delta} \subseteq \metricBall{d}{X}{x}{\epsilon}$
        then $T_1$ is finer than $T$ on $X$.
    \begin{subproof}
        Assume that for all $x\in X$ we have
            for all $\epsilon\in\reals$ such that $0 < \epsilon$ there exists $\delta\in\reals$
                such that $0 < \delta$ and
                $\metricBall{d_1}{X}{x}{\delta} \subseteq \metricBall{d}{X}{x}{\epsilon}$.
        $T$ is a topology on $X$ and $T_1$ is a topology on $X$
            by \cref{metric_topology,metric_basis_is_basis,basis_topology_is_topology}.
        Show that for all $U$ such that $U\in T$ we have $U\in T_1$.
        \begin{subproof}
            Fix $U$ such that $U\in T$.
            Hence $U$ is open in $X$ with respect to $T$ and $U\subseteq X$
                by \cref{open_in,topology_on,subseteq,powerset_elim}.
            Show that for all $y\in U$ there exists $\delta\in\reals$ such that $0 < \delta$ and
                $\metricBall{d_1}{X}{y}{\delta}\subseteq U$.
            \begin{subproof}
                Fix $y$ such that $y\in U$.
                Take $\epsilon\in\reals$ such that $0 < \epsilon$ and
                    $\metricBall{d}{X}{y}{\epsilon}\subseteq U$
                    by \cref{metric_open_iff}.
                Take $\delta\in\reals$ such that $0 < \delta$ and
                    $\metricBall{d_1}{X}{y}{\delta}\subseteq \metricBall{d}{X}{y}{\epsilon}$.
                Show that $\metricBall{d_1}{X}{y}{\delta}\subseteq U$.
                \begin{subproof}
                    Follows by \cref{subseteq_transitive}.
                \end{subproof}
            \end{subproof}
            Hence $U\in T_1$ by \cref{metric_open_iff,open_in}.
        \end{subproof}
        Therefore $T_1$ is finer than $T$ on $X$ by \cref{subseteq,finer_topology}.
    \end{subproof}
\end{proof}

% from §20 helper lemma 13: initial segment is finite
\begin{lemma}\label{initial_segment_finite}
    Suppose $n\in\naturalsPlus$.
    Then $\initialSegment{n}$ is finite.
\end{lemma}
\begin{proof}
    Let $A = \{n\in\naturalsPlus \mid \text{$\initialSegment{n}$ is finite}\}$.
    $A\subseteq\naturalsPlus$ by \cref{subseteq}.
    Let $B = \initialSegment{1}$.
    We have for all $k$ such that $k\in B$ we have $k = 1$.
    Hence $B$ is finite by \cref{upair_intro_left,subseteq,pair_is_finite,subseteq_finite_is_finite}.
    Therefore $1\in A$.
    Show that for all $m\in A$ we have $m + 1\in A$.
    \begin{subproof}
        Fix $m$ such that $m\in A$.
        We have $\initialSegment{m}$ is finite.
        $m + 1\in\naturalsPlus$ by \cref{subseteq,real_naturals_closed_succ}.
        Let $C = \initialSegment{m + 1}$.
        For all $k$ such that $k\in C$ we have $k\in \initialSegment{m}\union\{m + 1,m + 1\}$.
        Hence $C$ is finite by \cref{subseteq,pair_is_finite,union_of_finite_is_finite,subseteq_finite_is_finite}.
        Therefore $m + 1\in A$.
    \end{subproof}
    Therefore $\initialSegment{n}$ is finite by \cref{real_naturals_induction,subseteq}.
\end{proof}

% from §20 helper lemma 14: product family of constant reals
\begin{lemma}\label{product_family_reals_power}
    Suppose $J$ is inhabited.
    Let $X = \{(i,\reals) \mid i\in J\}$.
    Then $\productFamily{X} = \realsPower{J}$.
\end{lemma}
\begin{proof}
    Show that for all $i,y_1,y_2$ such that $(i,y_1)\in X$ and $(i,y_2)\in X$ we have $y_1 = y_2$.
    \begin{subproof}
        Trivial.
    \end{subproof}
    Hence $X$ is a function on $J$
        by \cref{relation,rightunique,dom_intro,dom_iff,pair_eq_iff,subseteq,subseteq_antisymmetric}.
    Show that for all $z$ such that $z\in\unions{\ran{X}}$ we have $z\in\reals$.
    \begin{subproof}
        Follows by \cref{ran_iff,pair_eq_iff,unions_iff}.
    \end{subproof}
    Show that for all $z$ such that $z\in\reals$ we have $z\in\unions{\ran{X}}$.
    \begin{subproof}
        Follows by \cref{ran_intro,unions_iff}.
    \end{subproof}
    Hence $\unions{\ran{X}} = \reals$ by \cref{subseteq,subseteq_antisymmetric}.
    Show that $\productFamily{X}\subseteq\realsPower{J}$.
    \begin{subproof}
        Follows by \cref{indexed_product,funs_elim,function_on_weaken_codom,funs_intro,tuple_power,reals_power,subseteq}.
    \end{subproof}
    Show that $\realsPower{J}\subseteq\productFamily{X}$.
    \begin{subproof}
        Follows by \cref{reals_power,tuple_power,funs_elim,function_on_weaken_codom,funs_intro,subseteq,function_apply_intro,indexed_product}.
    \end{subproof}
    Therefore $\productFamily{X} = \realsPower{J}$ by \cref{subseteq,subseteq_antisymmetric}.
\end{proof}

% from §20 helper lemma 15: coordinate difference set is finite
\begin{lemma}\label{coord_diff_set_finite}
    Suppose $n\in\naturalsPlus$.
    Let $J = \initialSegment{n}$.
    Suppose $x\in\realsPower{J}$.
    Suppose $y\in\realsPower{J}$.
    Then $\coordDiffSet{n}{x}{y}\subseteq\reals$ and
         $\coordDiffSet{n}{x}{y}$ is finite and
         $\coordDiffSet{n}{x}{y}$ is inhabited.
\end{lemma}
\begin{proof}
    Let $D = \coordDiffSet{n}{x}{y}$.
    Show that $D\subseteq\reals$.
    \begin{subproof}
        Follows by \cref{coord_diff_set,reals_power,tuple_power,funs_type_apply,real_add_inverse,real_add_closed,real_sub,real_abs_in_reals,subseteq}.
    \end{subproof}
    Show that $D$ is finite.
    \begin{subproof}
        Let $g = \{(i,\abs{x(i) - y(i)}) \mid i\in J\}$.
        Show that for all $i,t_1,t_2$ such that $(i,t_1)\in g$ and $(i,t_2)\in g$ we have $t_1 = t_2$.
        \begin{subproof}
            Follows by \cref{pair_eq_iff}.
        \end{subproof}
        Hence $g$ is a function on $J$ by \cref{coord_diff_set,relation,rightunique,dom_intro,dom_iff,pair_eq_iff,subseteq,subseteq_antisymmetric}.
        Therefore $D$ is finite by \cref{coord_diff_set,img_iff,setext,initial_segment_finite,finite_image_function}.
    \end{subproof}
    Show that $D$ is inhabited.
    \begin{subproof}
        Follows by \cref{real_one_in_naturals,real_one_leq_naturalsplus,initial_segment_naturalsplus,coord_diff_set}.
    \end{subproof}
\end{proof}

% from §20 helper lemma 16: square metric is a function
\begin{lemma}\label{square_metric_function}
    Suppose $n\in\naturalsPlus$.
    Then $\squareMetric{n}$ is a function.
\end{lemma}
\begin{proof}
    Follows by \cref{square_metric,relation,rightunique,pair_eq_iff,coord_diff_set_finite,subseteq,real_leq_antisym}.
\end{proof}

% from §20 helper lemma 17: absolute value bound from interval
\begin{lemma}\label{real_abs_lt_from_interval}
    Suppose $a\in\reals$.
    Suppose $b\in\reals$.
    Suppose $\epsilon\in\reals$.
    Suppose $0 < \epsilon$.
    Suppose $a - \epsilon < b$ and $b < a + \epsilon$.
    Then $\abs{a - b} < \epsilon$.
\end{lemma}
\begin{proof}
    $b - a\in\reals$ by \cref{real_add_inverse,real_sub,real_add_closed}.
    Hence $b + \neg{a} < (a + \epsilon) + \neg{a}$ by \cref{real_add_inverse,real_sub,real_add_closed,real_lt_add}.
    We have $(a + \epsilon) + \neg{a} = 0 + \epsilon$ by \cref{real_add_assoc,real_add_comm,real_add_inverse}.
    Thus $b - a < \epsilon$ by \cref{real_add_ident,real_sub}.
    $\neg{\epsilon}\in\reals$ by \cref{real_add_inverse}.
    We have $a + \neg{\epsilon} < b$ by \cref{real_sub,real_lt_subst_left}.
    Therefore $(a + \neg{\epsilon}) + \neg{a} < b + \neg{a}$ by \cref{real_add_inverse,real_add_closed,real_lt_add}.
    We have $(a - \epsilon) + \neg{a} = 0 + \neg{\epsilon}$ by \cref{real_sub,real_add_assoc,real_add_comm,real_add_inverse}.
    Hence $\neg{\epsilon} < b - a$ by \cref{real_add_ident,real_sub}.
    $b - a = 0$ or $b - a < 0$ or $0 < b - a$ by \cref{real_add_ident,real_lt_trichotomy}.
    Hence $\abs{b - a} < \epsilon$
        by \cref{real_abs_nonneg,real_lt_subst_left,real_abs_neg,real_neg_neg,real_add_inverse,real_lt_neg,real_lt_subst_right}.
    Therefore $\abs{a - b} < \epsilon$ by \cref{real_abs_sub_swap}.
\end{proof}

% from §20 helper lemma 17a: predecessor stays in initial segment
\begin{lemma}\label{initial_segment_pred_from_succ_bound_aux}
    Suppose $n\in\naturalsPlus$.
    Suppose $k\in\naturalsPlus$ and $k + 1 \leq n$.
    Then $k\in\initialSegment{n}$.
\end{lemma}
\begin{proof}
    $k\in\reals$ and $1\in\reals$ and $0\in\reals$
        by \cref{real_naturals_subset_reals,real_one_in_naturals,subseteq,real_add_ident}.
    Hence $k < k + 1$
        by \cref{real_zero_lt_one,real_lt_add,real_add_ident,real_add_comm,real_lt_subst_left}.
    $k + 1\in\reals$ and $n\in\reals$ by \cref{real_add_closed,real_naturals_subset_reals,subseteq}.
    Show that $k \leq n$.
    \begin{subproof}
        Suppose not.
        Hence $n < k$ by \cref{real_lt_trichotomy}.
        $k + 1 < n$ or $k + 1 = n$.
        Hence $k < n$ by \cref{real_lt_trans,real_lt_subst_right}.
        Therefore $n < n$ by \cref{real_lt_trans}.
        Contradiction by \cref{real_lt_irreflexive}.
    \end{subproof}
    Therefore $k\in\initialSegment{n}$ by \cref{initial_segment_naturalsplus}.
\end{proof}

% from §20 helper lemma 18: sum of nonnegative finite sequence is nonnegative
\begin{lemma}\label{sum_finite_sequence_nonneg}
    Suppose $\sumFiniteSequence{a}{s}$.
    Suppose for all $i\in\dom{a}$ we have $0 \leq a(i)$.
    Then $0 \leq s$.
\end{lemma}
\begin{proof}
    Take $n$ such that $n\in\naturalsPlus$ and $a\in\funs{\initialSegment{n}}{\reals}$ and
        there exists $t\in\funs{\initialSegment{n}}{\reals}$ such that
            $t(1) = a(1)$ and $s = t(n)$ and
            for all $k\in\naturalsPlus$ such that $k + 1 \leq n$ we have $t(k + 1) = t(k) + a(k + 1)$
        by \cref{sum_finite_sequence_pred}.
    Take $t$ such that $t\in\funs{\initialSegment{n}}{\reals}$ and
            $t(1) = a(1)$ and $s = t(n)$ and
            for all $k\in\naturalsPlus$ such that $k + 1 \leq n$ we have $t(k + 1) = t(k) + a(k + 1)$.
    $\dom{a} = \initialSegment{n}$ by \cref{funs_elim}.
    Let $A = \{k\in\naturalsPlus \mid \text{if $k \leq n$ then $0 \leq t(k)$}\}$.
    $A\subseteq\naturalsPlus$ by \cref{subseteq}.
    $1\in\naturalsPlus$ and $1 \leq n$ by \cref{real_one_in_naturals,real_one_leq_naturalsplus}.
    $0 \leq t(1)$.
    $1\in A$ by \cref{real_one_in_naturals}.
    Show that for all $k\in A$ we have $k + 1\in A$.
    \begin{subproof}
        Fix $k$ such that $k\in A$.
        $k\in\naturalsPlus$ and $k + 1\in\naturalsPlus$ by \cref{subseteq,real_naturals_closed_succ}.
        Show that if $k + 1 \leq n$ then $0 \leq t(k + 1)$.
        \begin{subproof}
            Assume $k + 1 \leq n$.
            $k\in\initialSegment{n}$ by \cref{initial_segment_pred_from_succ_bound_aux}.
            $k \leq n$ by \cref{initial_segment_naturalsplus}.
            $0 \leq t(k)$ by \cref{subseteq}.
            $k + 1\in\initialSegment{n}$ by \cref{initial_segment_naturalsplus}.
            $k + 1\in\dom{a}$ by \cref{funs_elim}.
            $0 \leq a(k + 1)$.
            $0\in\reals$ by \cref{real_add_ident}.
            $t(k)\in\reals$ and $a(k + 1)\in\reals$ by \cref{initial_segment_naturalsplus,funs_elim,funs_type_apply}.
            $a(k + 1) \leq t(k) + a(k + 1)$ by \cref{real_add_ident,real_leq_add}.
            $0 \leq t(k) + a(k + 1)$ by \cref{real_add_closed,real_leq_trans}.
            $0 \leq t(k + 1)$ by assumption.
        \end{subproof}
        $k + 1\in A$.
    \end{subproof}
    $0 \leq s$ by \cref{real_naturals_induction,subseteq}.
\end{proof}

% from §20 helper lemma 19: terms are bounded by the sum
\begin{lemma}\label{sum_finite_sequence_term_leq}
    Suppose $\sumFiniteSequence{a}{s}$.
    Suppose for all $i\in\dom{a}$ we have $0 \leq a(i)$.
    Then for all $i\in\dom{a}$ we have $a(i) \leq s$.
\end{lemma}
\begin{proof}
    Take $n$ such that $n\in\naturalsPlus$ and $a\in\funs{\initialSegment{n}}{\reals}$ and
        there exists $t\in\funs{\initialSegment{n}}{\reals}$ such that
            $t(1) = a(1)$ and $s = t(n)$ and
            for all $k\in\naturalsPlus$ such that $k + 1 \leq n$ we have $t(k + 1) = t(k) + a(k + 1)$
        by \cref{sum_finite_sequence_pred}.
    Take $t$ such that $t\in\funs{\initialSegment{n}}{\reals}$ and
            $t(1) = a(1)$ and $s = t(n)$ and
            for all $k\in\naturalsPlus$ such that $k + 1 \leq n$ we have $t(k + 1) = t(k) + a(k + 1)$.
    $\dom{a} = \initialSegment{n}$ and $\dom{t} = \initialSegment{n}$ by \cref{funs_elim}.
    Let $A = \{k\in\naturalsPlus \mid \text{if $k \leq n$ then
        for all $j\in\initialSegment{k}$ we have $t(j) \leq t(k)$}\}$.
    $A\subseteq\naturalsPlus$ by \cref{subseteq}.
    For all $j\in\initialSegment{1}$ we have $t(j) \leq t(1)$.
    $1\in A$.
    Show that for all $k\in A$ we have $k + 1\in A$.
    \begin{subproof}
        Fix $k$ such that $k\in A$.
        $k\in\naturalsPlus$ and $k + 1\in\naturalsPlus$ by \cref{subseteq,real_naturals_closed_succ}.
        Show that if $k + 1 \leq n$ then
            for all $j\in\initialSegment{k + 1}$ we have $t(j) \leq t(k + 1)$.
        \begin{subproof}
            Assume $k + 1 \leq n$.
            $k\in\initialSegment{n}$ by \cref{initial_segment_pred_from_succ_bound_aux}.
            $k \leq n$ by \cref{initial_segment_naturalsplus}.
            For all $j\in\initialSegment{k}$ we have $t(j) \leq t(k)$ by \cref{subseteq}.
            $0 \leq a(k + 1)$ by \cref{real_naturals_closed_succ,initial_segment_naturalsplus}.
            $a(k + 1)\in\reals$ and $t(k)\in\reals$ by \cref{initial_segment_naturalsplus,funs_elim,funs_type_apply}.
            $t(k) \leq t(k + 1)$ by \cref{real_leq_add,real_add_ident,real_add_comm,real_leq_trans}.
            Show that for all $j\in\initialSegment{k + 1}$ we have $t(j) \leq t(k + 1)$.
            \begin{subproof}
                Fix $j$ such that $j\in\initialSegment{k + 1}$.
                $j\in\naturalsPlus$ and $j \leq k + 1$ by \cref{initial_segment_naturalsplus}.
                $j = k + 1$ or $j \neq k + 1$.
                \begin{byCase}
                \caseOf{$j = k + 1$.}
                    $t(j) \leq t(k + 1)$.
                \caseOf{$j \neq k + 1$.}
                    $j \leq k$ by \cref{real_naturals_leq_succ_elim}.
                    $j\in\initialSegment{k}$ by \cref{initial_segment_naturalsplus}.
                    $t(j) \leq t(k)$.
                    $k \leq n$.
                    $j\in\reals$ and $k\in\reals$ and $n\in\reals$
                        by \cref{real_naturals_subset_reals,subseteq}.
                    $j \leq n$ by \cref{real_leq_trans}.
                    $j\in\initialSegment{n}$ and $k + 1\in\initialSegment{n}$ by \cref{initial_segment_naturalsplus}.
                    $t(j)\in\reals$ and $t(k + 1)\in\reals$ by \cref{funs_type_apply}.
                    $t(j) \leq t(k + 1)$ by \cref{real_leq_trans}.
                \end{byCase}
            \end{subproof}
        \end{subproof}
        $k + 1\in A$ by \cref{subseteq}.
    \end{subproof}
    For all $k\in\naturalsPlus$ we have $k\in A$ by \cref{real_naturals_induction}.
    Show that for all $i\in\dom{a}$ we have $a(i) \leq s$.
    \begin{subproof}
        Fix $i$ such that $i\in\dom{a}$.
        $i\in\naturalsPlus$ and $i \leq n$ by \cref{initial_segment_naturalsplus}.
        $t(i) \leq t(n)$ by \cref{subseteq,initial_segment_naturalsplus}.
        Show that $a(i) \leq t(i)$.
        \begin{subproof}
            $i = 1$ or $i \neq 1$.
            \begin{byCase}
            \caseOf{$i = 1$.}
                $a(i) = t(i)$.
            \caseOf{$i \neq 1$.}
                Take $m\in\naturalsPlus$ such that $i = m + 1$ by \cref{real_naturals_predecessor}.
                $m\in\reals$ and $i\in\reals$ and $n\in\reals$ and $1\in\reals$ and $0\in\reals$
                    by \cref{real_naturals_subset_reals,subseteq,real_one_in_naturals,real_add_ident}.
                $0 < 1$ by \cref{real_zero_lt_one}.
                $0 + m < 1 + m$ by \cref{real_lt_add}.
                $0 + m = m$ by \cref{real_add_ident}.
                $1 + m = m + 1$ by \cref{real_add_comm}.
                $m < m + 1$ by \cref{real_lt_subst_left}.
                $m < i$.
                $m \leq i$.
                $m \leq n$ by \cref{real_leq_trans}.
                $m\in A$ by \cref{subseteq}.
                $1 \leq m$ by \cref{real_one_leq_naturalsplus}.
                $1\in\initialSegment{m}$ by \cref{initial_segment_naturalsplus}.
                $t(1) \leq t(m)$.
                $1\in\dom{a}$ by \cref{real_one_in_naturals,real_one_leq_naturalsplus,initial_segment_naturalsplus}.
                $0 \leq a(1)$.
                $0 \leq t(1)$.
                $m\in\dom{t}$ by \cref{initial_segment_naturalsplus}.
                $t(1)\in\reals$ and $t(m)\in\reals$ by \cref{funs_type_apply}.
                $0 \leq t(m)$ by \cref{real_leq_trans}.
                $0 \leq a(i)$.
                $a(i)\in\reals$ and $t(m)\in\reals$ by \cref{funs_type_apply}.
                $m + 1 \leq n$.
                $t(m + 1) = t(m) + a(m + 1)$ by assumption.
                $t(i) = t(m) + a(i)$.
                $a(i) \leq t(i)$ by \cref{real_add_closed,real_leq_add,real_add_ident,real_leq_trans}.
            \end{byCase}
        \end{subproof}
        $a(i)\in\reals$ and $t(i)\in\reals$ and $t(n)\in\reals$
            by \cref{funs_type_apply,initial_segment_naturalsplus,funs_elim}.
        $a(i) \leq s$ by \cref{real_leq_trans}.
    \end{subproof}
\end{proof}

% from §20 helper lemma 20: square metric value exists
\begin{lemma}\label{square_metric_value_exists}
    Suppose $n\in\naturalsPlus$.
    Let $J = \initialSegment{n}$.
    Suppose $x\in\realsPower{J}$.
    Suppose $y\in\realsPower{J}$.
    Then there exists $r\in\reals$ such that $((x,y),r)\in\squareMetric{n}$ and
        for all $i\in J$ we have $\abs{x(i) - y(i)} \leq r$.
\end{lemma}
\begin{proof}
    Let $D = \coordDiffSet{n}{x}{y}$.
    $D\subseteq\reals$ and $D$ is finite and $D\neq\emptyset$ by \cref{coord_diff_set_finite,emptyset}.
    Take $r\in D$ such that $r$ is a maximum of $D$ by \cref{finite_real_maximum}.
    Show that $((x,y),r)\in\squareMetric{n}$ and
        for all $i\in J$ we have $\abs{x(i) - y(i)} \leq r$.
    \begin{subproof}
        Follows by \cref{real_maximum,times_tuple_intro,subseteq,square_metric,coord_diff_set}.
    \end{subproof}
\end{proof}
% from §20 helper lemma 21: bounded finite sum
\begin{lemma}\label{sum_finite_sequence_leq_n_times}
    Suppose $n\in\naturalsPlus$.
    Suppose $a\in\funs{\initialSegment{n}}{\reals}$.
    Suppose $\sumFiniteSequence{a}{s}$.
    Suppose $M\in\reals$.
    Suppose for all $i\in\initialSegment{n}$ we have $0 \leq a(i)$ and $a(i) \leq M$.
    Then $s \leq n \rmul M$.
\end{lemma}
\begin{proof}
    Take $n_0$ such that $n_0\in\naturalsPlus$ and $a\in\funs{\initialSegment{n_0}}{\reals}$ and
        there exists $t\in\funs{\initialSegment{n_0}}{\reals}$ such that
            $t(1) = a(1)$ and $s = t(n_0)$ and
            for all $k\in\naturalsPlus$ such that $k + 1 \leq n_0$ we have $t(k + 1) = t(k) + a(k + 1)$
        by \cref{sum_finite_sequence_pred}.
    Take $t$ such that $t\in\funs{\initialSegment{n_0}}{\reals}$ and
        $t(1) = a(1)$ and $s = t(n_0)$ and
        for all $k\in\naturalsPlus$ such that $k + 1 \leq n_0$ we have $t(k + 1) = t(k) + a(k + 1)$.
    $n = n_0$
        by \cref{funs_elim,subseteq_antisymmetric,subseteq,initial_segment_naturalsplus,real_naturals_subset_reals,real_leq_antisym}.
    $s = t(n)$.
    $t\in\funs{\initialSegment{n}}{\reals}$ and $\dom{t} = \initialSegment{n}$ by \cref{funs_elim}.
    Let $A = \{k\in\naturalsPlus \mid \text{if $k \leq n$ then $t(k) \leq k \rmul M$}\}$.
    $A\subseteq\naturalsPlus$ by \cref{subseteq}.
    $1\in\dom{a}$ by \cref{real_one_in_naturals,real_one_leq_naturalsplus,funs_elim,initial_segment_naturalsplus}.
    $0 \leq a(1)$ and $a(1) \leq M$.
    $1\in A$
        by \cref{real_naturals_subset_reals,real_one_in_naturals,subseteq,real_mul_ident}.
    Show that for all $k\in A$ we have $k + 1\in A$.
    \begin{subproof}
        Fix $k$ such that $k\in A$.
        $k\in\naturalsPlus$ and $k + 1\in\naturalsPlus$ by \cref{subseteq,real_naturals_closed_succ}.
        Show that if $k + 1 \leq n$ then $t(k + 1) \leq (k + 1) \rmul M$.
        \begin{subproof}
            Assume $k + 1 \leq n$.
            $k\in\initialSegment{n}$ by \cref{initial_segment_pred_from_succ_bound_aux}.
            $k \leq n$ by \cref{initial_segment_naturalsplus}.
            $t(k) \leq k \rmul M$ by \cref{subseteq}.
            $0 \leq a(k + 1)$ and $a(k + 1) \leq M$ by \cref{initial_segment_naturalsplus,funs_elim}.
            $k\in\reals$ by \cref{real_naturals_subset_reals,subseteq}.
            $a(k + 1)\in\reals$ and $t(k)\in\reals$ by \cref{initial_segment_naturalsplus,funs_elim,funs_type_apply}.
            $k \rmul M\in\reals$ by \cref{real_mul_closed}.
            $t(k) + a(k + 1)\in\reals$ by \cref{real_add_closed}.
            $t(k) + a(k + 1) \leq (k \rmul M) + a(k + 1)$ and
                $a(k + 1) + (k \rmul M) \leq M + (k \rmul M)$ by \cref{real_leq_add}.
            $a(k + 1) + (k \rmul M) = (k \rmul M) + a(k + 1)$ and
                $M + (k \rmul M) = (k \rmul M) + M$ by \cref{real_add_comm}.
            $(k \rmul M) + a(k + 1)\in\reals$ and $(k \rmul M) + M\in\reals$
                by \cref{real_add_closed}.
            $(k \rmul M) + a(k + 1) \leq (k \rmul M) + M$ by \cref{real_leq_trans}.
            $t(k) + a(k + 1) \leq (k \rmul M) + M$ by \cref{real_leq_trans}.
            $k + 1 \leq n_0$.
            $t(k + 1) = t(k) + a(k + 1)$ by assumption.
            $(k \rmul M) + M = (k + 1) \rmul M$
                by \cref{real_mul_ident,real_add_closed,real_mul_comm,real_distrib}.
            $t(k + 1) \leq (k + 1) \rmul M$ by \cref{real_leq_trans}.
        \end{subproof}
        $k + 1\in A$.
    \end{subproof}
    For all $k\in\naturalsPlus$ we have $k\in A$ by \cref{real_naturals_induction}.
    $s \leq n \rmul M$ by \cref{subseteq}.
\end{proof}

% from §20 helper lemma 22: squares preserve order on nonnegative reals
\begin{lemma}\label{real_nonneg_leq_iff_sq_leq}
    Suppose $a\in\reals$.
    Suppose $b\in\reals$.
    Suppose $0 \leq a$.
    Suppose $0 \leq b$.
    Then $a \leq b$ iff $\exp{a}{1+1} \leq \exp{b}{1+1}$.
\end{lemma}
\begin{proof}
    Show that if $a \leq b$ then $\exp{a}{1+1} \leq \exp{b}{1+1}$.
    \begin{subproof}
        Assume $a \leq b$.
        $a < b$ or $a = b$ by \cref{real_leq_split}.
        \begin{byCase}
        \caseOf{$a = b$.}
            $\exp{a}{1+1} \leq \exp{b}{1+1}$.
        \caseOf{$a < b$.}
            $0 < a$ or $a = 0$ by \cref{real_leq_split}.
            \begin{byCase}
            \caseOf{$a = 0$.}
                $0 < b$ by \cref{real_lt_subst_left}.
                $\exp{a}{1+1} \leq \exp{b}{1+1}$
                    by \cref{real_pow_one,real_one_in_naturals,real_pow_succ,real_mul_zero,real_sq_nonnegative}.
            \caseOf{$0 < a$.}
                $0 < b$ by \cref{real_lt_trans}.
                $\exp{a}{1+1} < \exp{b}{1+1}$ by \cref{real_pos_lt_iff_sq_lt}.
                $\exp{a}{1+1} \leq \exp{b}{1+1}$.
            \end{byCase}
        \end{byCase}
    \end{subproof}
    Show that if $\exp{a}{1+1} \leq \exp{b}{1+1}$ then $a \leq b$.
    \begin{subproof}
        Assume $\exp{a}{1+1} \leq \exp{b}{1+1}$.
        Suppose not.
        $b < a$.
        $0 < b$ or $b = 0$ by \cref{real_leq_split}.
        \begin{byCase}
        \caseOf{$b = 0$.}
            $0 < a$ by \cref{real_lt_subst_left}.
            $a$ is positive.
            $\exp{a}{1+1} = a \rmul a$ by \cref{real_one_in_naturals,real_pow_succ,real_pow_one}.
            $\exp{a}{1+1}\in\reals$ by \cref{real_mul_closed}.
            $\exp{a}{1+1} > 0$ by \cref{real_mul_pos_pos_gt}.
            $\exp{b}{1+1} = 0$ by \cref{real_pow_one,real_one_in_naturals,real_pow_succ,real_mul_zero}.
            $\exp{b}{1+1}\in\reals$ by \cref{real_add_ident}.
            $\exp{b}{1+1} < \exp{a}{1+1}$ by \cref{real_lt_subst_left}.
            $\exp{a}{1+1} < \exp{b}{1+1}$ or $\exp{a}{1+1} = \exp{b}{1+1}$ by \cref{real_leq_split}.
            \begin{byCase}
            \caseOf{$\exp{a}{1+1} < \exp{b}{1+1}$.}
                $\exp{a}{1+1} < \exp{a}{1+1}$ by \cref{real_lt_trans}.
                Contradiction by \cref{real_lt_irreflexive}.
            \caseOf{$\exp{a}{1+1} = \exp{b}{1+1}$.}
                $\exp{b}{1+1} < \exp{b}{1+1}$ by \cref{real_lt_subst_right}.
                Contradiction by \cref{real_lt_irreflexive}.
            \end{byCase}
        \caseOf{$0 < b$.}
            $0 < a$ by \cref{real_lt_trans}.
            $\exp{b}{1+1} < \exp{a}{1+1}$ by \cref{real_pos_lt_iff_sq_lt}.
            $\exp{a}{1+1} = a \rmul a$ by \cref{real_one_in_naturals,real_pow_succ,real_pow_one}.
            $\exp{a}{1+1}\in\reals$ by \cref{real_mul_closed}.
            $\exp{b}{1+1} = b \rmul b$ by \cref{real_one_in_naturals,real_pow_succ,real_pow_one}.
            $\exp{b}{1+1}\in\reals$ by \cref{real_mul_closed}.
            $\exp{a}{1+1} < \exp{b}{1+1}$ or $\exp{a}{1+1} = \exp{b}{1+1}$ by \cref{real_leq_split}.
            \begin{byCase}
            \caseOf{$\exp{a}{1+1} < \exp{b}{1+1}$.}
                $\exp{a}{1+1} < \exp{a}{1+1}$ by \cref{real_lt_trans}.
                Contradiction by \cref{real_lt_irreflexive}.
            \caseOf{$\exp{a}{1+1} = \exp{b}{1+1}$.}
                $\exp{b}{1+1} < \exp{b}{1+1}$ by \cref{real_lt_subst_right}.
                Contradiction by \cref{real_lt_irreflexive}.
            \end{byCase}
        \end{byCase}
    \end{subproof}
\end{proof}

% from §20 helper lemma 23: coordinate square sequence element
\begin{lemma}\label{coord_square_sequence_elem}
    Suppose $n\in\naturalsPlus$.
    Let $J = \initialSegment{n}$.
    Suppose $x\in\realsPower{J}$.
    Suppose $y\in\realsPower{J}$.
    Let $a = \coordSquareSequence{n}{(x,y)}$.
    Suppose $i\in J$.
    Then $(i,\exp{x(i) - y(i)}{1+1})\in a$.
\end{lemma}
\begin{proof}
    Follows by \cref{coord_square_sequence,fst_eq,snd_eq}.
\end{proof}

% from §20 helper lemma 24: euclidean and square metric bounds
\begin{lemma}\label{euclidean_square_metric_bounds}
    Suppose $n\in\naturalsPlus$.
    Let $J = \initialSegment{n}$.
    Suppose $x\in\realsPower{J}$.
    Suppose $y\in\realsPower{J}$.
    Suppose $((x,y),r)\in\euclideanMetric{n}$.
    Suppose $((x,y),m)\in\squareMetric{n}$.
    Then $m \leq r$ and $r \leq \sqrt{n} \rmul m$.
\end{lemma}
\begin{proof}
    $1\in\naturalsPlus$ by \cref{real_one_in_naturals}.
    Let $D = \coordDiffSet{n}{x}{y}$.
    Take $s\in\reals$ such that
        $\sumFiniteSequence{\coordSquareSequence{n}{(x,y)}}{s}$ and $r = \sqrt{s}$
        by \cref{euclidean_metric,pair_eq_iff}.
    Let $a = \coordSquareSequence{n}{(x,y)}$.
    Take $n_1$ such that $n_1\in\naturalsPlus$ and $a\in\funs{\initialSegment{n_1}}{\reals}$ and
        there exists $t\in\funs{\initialSegment{n_1}}{\reals}$ such that
            $t(1) = a(1)$ and $s = t(n_1)$ and
            for all $k\in\naturalsPlus$ such that $k + 1 \leq n_1$ we have $t(k + 1) = t(k) + a(k + 1)$
        by \cref{sum_finite_sequence_pred}.
    $a$ is a function and $\dom{a} = \initialSegment{n_1}$ by \cref{funs_elim}.
    Take $x_0,y_0\in\realsPower{J}$ such that
        $(x,y) = (x_0,y_0)$ and
        $m\in\coordDiffSet{n}{x_0}{y_0}$ and
        for all $t\in\coordDiffSet{n}{x_0}{y_0}$ we have $t \leq m$
        by \cref{square_metric,pair_eq_iff}.
    $x = x_0$ and $y = y_0$ by \cref{pair_eq_iff}.
    $m\in D$ and for all $t\in D$ we have $t \leq m$.
    Take $i_0\in J$ such that $m = \abs{x(i_0) - y(i_0)}$ by \cref{coord_diff_set}.
    Show that for all $i\in J$ we have $0 \leq a(i)$.
    \begin{subproof}
        Follows by \cref{coord_square_sequence_elem,reals_power,tuple_power,funs_type_apply,real_add_inverse,real_add_closed,real_sub,real_sq_nonnegative,function_apply_intro}.
    \end{subproof}
    $\dom{a} = J$
        by \cref{coord_square_sequence_elem,dom_intro,subseteq,dom_iff,coord_square_sequence,pair_eq_iff,subseteq_antisymmetric}.
    $a\in\funs{J}{\reals}$ by \cref{funs_elim,funs_intro}.
    $0 \leq s$ by \cref{sum_finite_sequence_nonneg}.
    $r\in\reals$ and $r \geq 0$ and $r \rmul r = s$ by \cref{positive_square_root}.
    $\exp{r}{1+1} = s$ by \cref{real_pow_succ,real_pow_one}.
    $a(i_0) \leq s$ by \cref{coord_square_sequence_elem,dom_intro,sum_finite_sequence_term_leq}.
    $x(i_0) - y(i_0)\in\reals$
        by \cref{reals_power,tuple_power,funs_type_apply,real_add_inverse,real_add_closed,real_sub}.
    $\exp{m}{1+1} \leq s$
        by \cref{real_abs_square,coord_square_sequence_elem,function_apply_intro}.
    $m \leq r$ by \cref{coord_diff_set_finite,subseteq,real_abs_nonnegative,real_nonneg_leq_iff_sq_leq}.
    Show that for all $i\in\initialSegment{n}$ we have $a(i) \leq \exp{m}{1+1}$.
    \begin{subproof}
        Follows by \cref{coord_diff_set,coord_diff_set_finite,subseteq,reals_power,tuple_power,funs_type_apply,real_add_inverse,real_add_closed,real_sub,real_abs_in_reals,real_abs_nonnegative,real_add_ident,real_leq_trans,real_nonneg_leq_iff_sq_leq,real_abs_square,coord_square_sequence_elem,function_apply_intro}.
    \end{subproof}
    Show that for all $i\in\initialSegment{n}$ we have $0 \leq a(i)$ and $a(i) \leq \exp{m}{1+1}$.
    \begin{subproof}
        Trivial.
    \end{subproof}
    $m\in\reals$ and $\exp{m}{1+1} = m \rmul m$ and $\exp{m}{1+1}\in\reals$
        by \cref{coord_diff_set_finite,subseteq,real_pow_succ,real_pow_one,real_mul_closed}.
    $s \leq n \rmul \exp{m}{1+1}$ by \cref{sum_finite_sequence_leq_n_times}.
    $n\in\reals$ by \cref{real_naturals_subset_reals,subseteq}.
    $0 < 1$ and $0\in\reals$ by \cref{real_zero_lt_one,real_add_ident}.
    $1\in\reals$ by \cref{real_naturals_subset_reals,subseteq}.
    $0 \leq 1$.
    $1 \leq n$ by \cref{real_one_leq_naturalsplus}.
    $0 \leq n$ by \cref{real_leq_trans}.
    $0 < n$.
    $\sqrt{n}\in\reals$ and $\sqrt{n} \geq 0$ and $\sqrt{n} \rmul \sqrt{n} = n$ by \cref{positive_square_root}.
    $\sqrt{n} \rmul m\in\reals$ by \cref{real_mul_closed}.
    $\exp{\sqrt{n} \rmul m}{1+1} = (\sqrt{n} \rmul m) \rmul (\sqrt{n} \rmul m)$ by \cref{real_pow_succ,real_pow_one}.
    $(\sqrt{n} \rmul m) \rmul (\sqrt{n} \rmul m) = (\sqrt{n} \rmul \sqrt{n}) \rmul (m \rmul m)$ by \cref{real_mul_assoc,real_mul_comm}.
    $\exp{\sqrt{n} \rmul m}{1+1} = n \rmul (m \rmul m)$.
    $\exp{m}{1+1} = m \rmul m$ by \cref{real_pow_succ,real_pow_one}.
    $\exp{\sqrt{n} \rmul m}{1+1} = n \rmul \exp{m}{1+1}$.
    $\exp{r}{1+1} \leq \exp{\sqrt{n} \rmul m}{1+1}$ by \cref{real_leq_trans}.
    Show that $\sqrt{n} \rmul m \geq 0$.
    \begin{subproof}
        $x(i_0) - y(i_0)\in\reals$
            by \cref{reals_power,tuple_power,funs_type_apply,real_add_inverse,real_add_closed,real_sub}.
        $m \geq 0$ by \cref{real_abs_nonnegative}.
        $0 < m$ or $m = 0$ by \cref{real_leq_split}.
        \begin{byCase}
        \caseOf{$m = 0$.}
            $\sqrt{n} \rmul m \geq 0$ by \cref{real_mul_comm,real_mul_zero}.
        \caseOf{$0 < m$.}
            $m$ is positive.
            Suppose not.
            $\sqrt{n} = 0$.
            $n = 0$ by \cref{positive_square_root,real_mul_zero}.
            $0 < 0$ by \cref{real_lt_subst_right}.
            Contradiction by \cref{real_lt_irreflexive}.
            $\sqrt{n}$ is positive.
            $\sqrt{n} \rmul m \geq 0$ by \cref{real_mul_pos_pos_gt}.
        \end{byCase}
    \end{subproof}
    Show that $m \leq r$.
    \begin{subproof}
        Trivial.
    \end{subproof}
    Show that $r \leq \sqrt{n} \rmul m$.
    \begin{subproof}
        Follows by \cref{real_nonneg_leq_iff_sq_leq}.
    \end{subproof}
\end{proof}
% from §20 helper lemma 25: interval bound from absolute value
\begin{lemma}\label{real_interval_from_abs_lt}
    Suppose $a\in\reals$.
    Suppose $b\in\reals$.
    Suppose $\epsilon\in\reals$.
    Suppose $0 < \epsilon$.
    Suppose $\abs{a - b} < \epsilon$.
    Then $a - \epsilon < b$ and $b < a + \epsilon$.
\end{lemma}
\begin{proof}
    $\abs{a - b} \leq \epsilon$ and $\epsilon$ is nonnegative.
    $b - \epsilon \leq a$ by \cref{real_abs_sub_leq_imp_left_bound}.
    $a \leq b + \epsilon$ by \cref{real_abs_sub_leq_imp_right_bound}.
    $\neg{\epsilon}\in\reals$ and $b - \epsilon\in\reals$ by \cref{real_add_inverse,real_sub,real_add_closed}.
    $b \leq a + \epsilon$
        by \cref{real_leq_add,real_sub,real_add_assoc,real_add_inverse,real_add_comm,real_add_ident,real_leq_trans}.
    $a - \epsilon \leq b$
        by \cref{real_add_closed,real_leq_add,real_sub,real_add_assoc,real_add_inverse,real_add_ident,real_add_comm}.
    Suppose not.
    $b = a + \epsilon$.
    $\abs{a - b} = \epsilon$
        by \cref{real_sub,real_add_inverse,real_add_assoc,real_add_comm,real_add_ident,real_abs_nonneg,real_abs_sub_swap}.
    $\epsilon < \epsilon$ by \cref{real_lt_subst_left}.
    Contradiction by \cref{real_lt_irreflexive}.
    Suppose not.
    $a - \epsilon = b$.
    $a - b = \epsilon$
        by \cref{real_sub,real_neg_addition,real_neg_neg,real_add_inverse,real_add_assoc,real_add_ident}.
    $\abs{a - b} = \epsilon$ by \cref{real_abs_nonneg}.
    $\epsilon < \epsilon$ by \cref{real_lt_subst_left}.
    Contradiction by \cref{real_lt_irreflexive}.
    Show that $a - \epsilon < b$.
    \begin{subproof}
        Trivial.
    \end{subproof}
    Show that $b < a + \epsilon$.
    \begin{subproof}
        Trivial.
    \end{subproof}
\end{proof}

% from §20 helper lemma 26: coordinate square sequence is a function
\begin{lemma}\label{coord_square_sequence_function}
    Suppose $n\in\naturalsPlus$.
    Let $J = \initialSegment{n}$.
    Suppose $x\in\realsPower{J}$.
    Suppose $y\in\realsPower{J}$.
    Then $\coordSquareSequence{n}{(x,y)}\in\funs{J}{\reals}$.
\end{lemma}
\begin{proof}
    Let $a = \coordSquareSequence{n}{(x,y)}$.
    $a\in\funs{J}{\reals}$
        by \cref{coord_square_sequence,pair_eq_iff,fst_eq,snd_eq,relation,rightunique,coord_square_sequence_elem,dom_intro,subseteq,dom_iff,subseteq_antisymmetric,reals_power,tuple_power,funs_type_apply,real_add_inverse,real_add_closed,real_sub,real_one_in_naturals,real_pow_succ,real_pow_one,real_mul_closed,function_apply_intro,funs_intro}.
\end{proof}

% from §20 helper lemma 26a: strict inequality implies weak inequality
\begin{lemma}\label{real_lt_imp_leq_for_sum_succ}
    Suppose $a < b$.
    Then $a \leq b$.
\end{lemma}
\begin{proof}
    Trivial.
\end{proof}

% from §20 helper lemma 27: sum finite sequence successor closure
\begin{lemma}\label{sum_finite_sequence_exists_succ}
    Let $A = \{n\in\naturalsPlus \mid \text{for all $a\in\funs{\initialSegment{n}}{\reals}$ there exists $s\in\reals$ such that $\sumFiniteSequence{a}{s}$}\}$.
    Then for all $n\in A$ we have $n + 1\in A$.
\end{lemma}
\begin{proof}
    Fix $n$ such that $n\in A$.
    For all $a\in\funs{\initialSegment{n}}{\reals}$ there exists $s\in\reals$
        such that $\sumFiniteSequence{a}{s}$.
    $n + 1\in\naturalsPlus$ by \cref{subseteq,real_naturals_closed_succ}.
    Show that for all $a\in\funs{\initialSegment{n + 1}}{\reals}$ there exists $s\in\reals$ such that $\sumFiniteSequence{a}{s}$.
    \begin{subproof}
        Fix $a$ such that $a\in\funs{\initialSegment{n + 1}}{\reals}$.
        Let $a_0 = \{(i,a(i)) \mid i\in\initialSegment{n}\}$.
        $a_0$ is a function by \cref{pair_eq_iff,relation,rightunique}.
        $\dom{a_0} = \initialSegment{n}$ by \cref{dom_intro,dom_iff,pair_eq_iff,subseteq,subseteq_antisymmetric}.
        For all $i$ such that $i\in\dom{a_0}$ we have $a_0(i)\in\reals$
            by \cref{subseteq,initial_segment_naturalsplus,real_naturals_subset_reals,real_one_in_naturals,real_add_ident,real_add_closed,real_zero_lt_one,real_lt_add,real_add_comm,real_lt_subst_left,real_lt_trans,funs_elim,funs_type_apply,function_apply_intro}.
                    $a_0\in\funs{\initialSegment{n}}{\reals}$ by \cref{funs_intro}.
                Take $s_0\in\reals$ such that $\sumFiniteSequence{a_0}{s_0}$.
                Take $n_0$ such that $n_0\in\naturalsPlus$ and $a_0\in\funs{\initialSegment{n_0}}{\reals}$ and
                    there exists $t_0\in\funs{\initialSegment{n_0}}{\reals}$ such that
                        $t_0(1) = a_0(1)$ and $s_0 = t_0(n_0)$ and
                        for all $k\in\naturalsPlus$ such that $k + 1 \leq n_0$ we have $t_0(k + 1) = t_0(k) + a_0(k + 1)$
                    by \cref{sum_finite_sequence_pred}.
                Take $t_0$ such that $t_0\in\funs{\initialSegment{n_0}}{\reals}$ and
                    $t_0(1) = a_0(1)$ and $s_0 = t_0(n_0)$ and
                    for all $k\in\naturalsPlus$ such that $k + 1 \leq n_0$ we have $t_0(k + 1) = t_0(k) + a_0(k + 1)$.
                    $\initialSegment{n_0} = \initialSegment{n}$ by \cref{funs_elim,subseteq_antisymmetric,subseteq}.
                    $\dom{t_0} = \initialSegment{n}$ by \cref{funs_elim}.
                $t_0$ is a function by \cref{funs_elim}.
                $t_0\in\funs{\initialSegment{n}}{\reals}$ by \cref{funs_intro}.
                    $t_0(n)\in\reals$ by \cref{initial_segment_naturalsplus,funs_type_apply}.
                    $a(n + 1)\in\reals$ by \cref{initial_segment_naturalsplus,funs_type_apply}.
                Let $v = t_0(n) + a(n + 1)$.
                $v\in\reals$ by \cref{real_add_closed}.
                Let $g = \{(n + 1,v)\}$.
                $g$ is a function by \cref{singleton_elim,pair_eq_iff,relation,rightunique}.
                Show that for all $x$ such that $x\in\dom{t_0}\inter\dom{g}$ we have $t_0(x) = g(x)$.
                \begin{subproof}
                    Fix $x$ such that $x\in\dom{t_0}\inter\dom{g}$.
                    Suppose not.
                        $x\in\dom{g}$.
                    Take $y$ such that $(x,y)\in g$ by \cref{dom_iff}.
                        $x = n + 1$ by \cref{singleton_elim,pair_eq_iff}.
                        $x\in\dom{t_0}$.
                        $x\in\initialSegment{n}$ by \cref{funs_elim}.
                        $x\in\naturalsPlus$ by \cref{initial_segment_naturalsplus}.
                        $x < n$ or $x = n$ by \cref{initial_segment_naturalsplus}.
                    $x\in\reals$ and $n\in\reals$ and $1\in\reals$ and $0\in\reals$
                        by \cref{real_naturals_subset_reals,subseteq,real_one_in_naturals,real_add_ident}.
                    $n < n + 1$
                        by \cref{real_zero_lt_one,real_lt_add,real_add_ident,real_add_comm,real_lt_subst_left}.
                        $n < x$ by \cref{real_lt_subst_right}.
                    $n < n$ by \cref{real_lt_trans,real_lt_subst_right}.
                    Contradiction by \cref{real_lt_irreflexive}.
                \end{subproof}
                Let $t = t_0\union g$.
                $t$ is a function and $\dom{t} = \dom{t_0}\union\dom{g}$ by \cref{union_compatible_functions,dom_union}.
                For all $i$ such that $i\in\dom{g}$ we have $i\in\{n + 1\}$
                    by \cref{dom_iff,singleton_elim,pair_eq_iff,singleton_intro}.
                For all $i$ such that $i\in\{n + 1\}$ we have $i\in\dom{g}$
                    by \cref{singleton_elim,singleton_intro,pair_eq_iff,dom_intro}.
                    $\dom{g} = \{n + 1\}$ and $\dom{t} = \initialSegment{n}\union\{n + 1\}$
                        by \cref{subseteq,subseteq_antisymmetric,funs_elim}.
                For all $i$ such that $i\in\dom{t_0}$ we have $t(i)\in\reals$
                    by \cref{funs_type_apply,function_apply_elim,union_iff,function_apply_intro}.
                For all $i$ such that $i\in\dom{g}$ we have $t(i)\in\reals$
                    by \cref{dom_iff,singleton_elim,pair_eq_iff,union_iff,function_apply_intro}.
                Show that for all $i$ such that $i\in\dom{t}$ we have $t(i)\in\reals$.
                \begin{subproof}
                    Follows by \cref{dom_union,union_iff}.
                \end{subproof}
                    $1\in\dom{t_0}$ by \cref{real_one_leq_naturalsplus,initial_segment_naturalsplus,real_one_in_naturals,funs_elim}.
                $t(1) = t_0(1)$ by \cref{function_apply_elim,union_iff,function_apply_intro}.
                $(1,a(1))\in a_0$.
                    $a_0(1) = a(1)$ by \cref{funs_elim,function_apply_intro}.
                    $t(1) = a(1)$.
                Show that for all $k\in\naturalsPlus$ such that $k + 1 \leq n + 1$ we have
                    $t(k + 1) = t(k) + a(k + 1)$.
                \begin{subproof}
                    Fix $k$ such that $k\in\naturalsPlus$ and $k + 1 \leq n + 1$.
                    $k + 1\in\naturalsPlus$ by \cref{real_naturals_closed_succ}.
                    $k + 1 = n + 1$ or $k + 1 \neq n + 1$.
                    \begin{byCase}
                    \caseOf{$k + 1 = n + 1$.}
                        $k\in\reals$ and $n\in\reals$ and $1\in\reals$ by \cref{real_naturals_subset_reals,subseteq,real_one_in_naturals}.
                            $k = n$ by \cref{real_add_cancel}.
                            $k\in\dom{t_0}$ by \cref{real_add_cancel,initial_segment_naturalsplus,funs_elim}.
                        $(k,t_0(k))\in t_0$ by \cref{function_apply_elim}.
                            $(k,t_0(k))\in t$ by \cref{union_iff}.
                            $t(k) = t_0(k)$ by \cref{function_apply_intro}.
                        $t_0(n)\in\reals$ by \cref{funs_elim,initial_segment_naturalsplus,funs_type_apply}.
                            $t(k) = t_0(n)$.
                        $(k + 1,v)\in g$ by \cref{singleton_intro,pair_eq_iff}.
                            $(k + 1,v)\in t$ by \cref{union_iff}.
                            $t(k + 1) = v$ by \cref{function_apply_intro}.
                            $t(k + 1) = t(k) + a(k + 1)$.
                    \caseOf{$k + 1 \neq n + 1$.}
                            $k + 1 \leq n$ by \cref{real_naturals_leq_succ_elim}.
                            $k + 1\in\initialSegment{n}$ by \cref{initial_segment_naturalsplus}.
                            $k\in\initialSegment{n}$ by \cref{initial_segment_pred_from_succ_bound_aux}.
                        $k\in\dom{t_0}$ and $k + 1\in\dom{t_0}$ by \cref{funs_elim,subseteq}.
                        $t(k) = t_0(k)$ and $t(k + 1) = t_0(k + 1)$
                            by \cref{function_apply_elim,union_iff,function_apply_intro}.
                            $k + 1\in\initialSegment{n_0}$.
                            $k + 1 < n_0$ or $k + 1 = n_0$ by \cref{initial_segment_naturalsplus}.
                            $k + 1 \leq n_0$ by \cref{real_lt_imp_leq_for_sum_succ}.
                            $t_0(k + 1) = t_0(k) + a_0(k + 1)$ by \cref{sum_finite_sequence_pred}.
                            $t(k + 1) = t(k) + a(k + 1)$
                                by \cref{funs_elim,function_apply_intro}.
                    \end{byCase}
                \end{subproof}
                Let $s = t(n + 1)$.
                $s\in\reals$ by \cref{singleton_intro,union_iff,dom_intro,funs_type_apply}.
                $n + 1\in\naturalsPlus$ by \cref{real_naturals_closed_succ}.
                For all $i$ such that $i\in\initialSegment{n}\union\{n + 1\}$ we have $i\in\initialSegment{n + 1}$
                    by \cref{union_iff,initial_segment_naturalsplus,real_naturals_subset_reals,subseteq,real_one_in_naturals,real_add_ident,real_add_closed,real_zero_lt_one,real_lt_add,real_add_comm,real_lt_subst_left,real_lt_trans,singleton_elim}.
                For all $i$ such that $i\in\initialSegment{n + 1}$ we have $i < n + 1$ or $i = n + 1$
                    by \cref{initial_segment_naturalsplus}.
                For all $i$ such that $i\in\initialSegment{n + 1}$ we have $i\in\initialSegment{n}\union\{n + 1\}$
                    by \cref{initial_segment_naturalsplus,real_naturals_leq_succ_elim,union_iff,singleton_intro}.
                    $\initialSegment{n}\union\{n + 1\} = \initialSegment{n + 1}$ by \cref{subseteq,subseteq_antisymmetric}.
                    $\dom{t} = \initialSegment{n + 1}$.
                    $t\in\funs{\initialSegment{n + 1}}{\reals}$ by \cref{funs_intro}.
                    $\sumFiniteSequence{a}{s}$ by \cref{sum_finite_sequence_pred}.
    \end{subproof}
    $n + 1\in A$.
\end{proof}

% from §20 helper lemma 28: sum of a finite sequence exists
\begin{lemma}\label{sum_finite_sequence_exists}
    Suppose $m\in\naturalsPlus$.
    Suppose $f\in\funs{\initialSegment{m}}{\reals}$.
    Then there exists $s\in\reals$ such that $\sumFiniteSequence{f}{s}$.
\end{lemma}
\begin{proof}
    Let $A = \{n\in\naturalsPlus \mid \text{for all $a\in\funs{\initialSegment{n}}{\reals}$ there exists $s\in\reals$ such that $\sumFiniteSequence{a}{s}$}\}$.
    $A\subseteq\naturalsPlus$ by \cref{subseteq}.
    $1\in A$
        by \cref{initial_segment_naturalsplus,real_one_in_naturals,funs_type_apply,real_one_leq_naturalsplus,real_naturals_subset_reals,subseteq,real_add_closed,real_leq_add,real_add_ident,real_zero_lt_one,real_lt_add,real_lt_subst_left,real_lt_trans,real_lt_subst_right,real_lt_irreflexive,sum_finite_sequence_pred}.
    For all $n\in A$ we have $n + 1\in A$ by \cref{sum_finite_sequence_exists_succ}.
    For all $n\in\naturalsPlus$ we have $n\in A$ by \cref{real_naturals_induction}.
    There exists $s\in\reals$ such that $\sumFiniteSequence{f}{s}$ by \cref{subseteq}.
\end{proof}

% from §20 helper definition 1: coordinate sum sequence
\begin{definition}\label{coord_sum_sequence}
    $\coordSumSequence{n}{a}{b} = \{w \mid \exists i\in\initialSegment{n}.
        w = (i, a(i) + b(i))\}$.
\end{definition}

% from §20 helper definition 2: coordinate product sequence
\begin{definition}\label{coord_product_sequence}
    $\coordProductSequence{n}{a}{b} = \{w \mid \exists i\in\initialSegment{n}.
        w = (i, a(i) \rmul b(i))\}$.
\end{definition}

% from §20 helper definition 3: coordinate square sequence
\begin{definition}\label{coord_square_sequence_mul}
    $\coordSquareSequenceMul{n}{a} = \{w \mid \exists i\in\initialSegment{n}.
        w = (i, a(i) \rmul a(i))\}$.
\end{definition}

% from §20 helper definition 4: coordinate difference sequence
\begin{definition}\label{coord_diff_sequence}
    $\coordDiffSequence{n}{x}{y} = \{w \mid \exists i\in\initialSegment{n}.
        w = (i, x(i) - y(i))\}$.
\end{definition}

% from §20 helper definition 5: coordinate scalar sequence
\begin{definition}\label{coord_scalar_sequence}
    $\coordScalarSequence{n}{c}{a} = \{w \mid \exists i\in\initialSegment{n}.
        w = (i, c \rmul a(i))\}$.
\end{definition}

% from §20 helper lemma 29: coordinate sum sequence is a function
\begin{lemma}\label{coord_sum_sequence_function}
    Suppose $n\in\naturalsPlus$.
    Let $J = \initialSegment{n}$.
    Suppose $a\in\funs{J}{\reals}$.
    Suppose $b\in\funs{J}{\reals}$.
    Then $\coordSumSequence{n}{a}{b}\in\funs{J}{\reals}$.
\end{lemma}
\begin{proof}
    Let $c = \coordSumSequence{n}{a}{b}$.
    $c\in\funs{J}{\reals}$
        by \cref{coord_sum_sequence,pair_eq_iff,relation,rightunique,dom_intro,subseteq,dom_iff,subseteq_antisymmetric,funs_type_apply,real_add_closed,function_apply_intro,funs_intro}.
\end{proof}

% from §20 helper lemma 30: coordinate product sequence is a function
\begin{lemma}\label{coord_product_sequence_function}
    Suppose $n\in\naturalsPlus$.
    Let $J = \initialSegment{n}$.
    Suppose $a\in\funs{J}{\reals}$.
    Suppose $b\in\funs{J}{\reals}$.
    Then $\coordProductSequence{n}{a}{b}\in\funs{J}{\reals}$.
\end{lemma}
\begin{proof}
    Let $c = \coordProductSequence{n}{a}{b}$.
    $c\in\funs{J}{\reals}$
        by \cref{coord_product_sequence,pair_eq_iff,relation,rightunique,dom_intro,subseteq,dom_iff,subseteq_antisymmetric,funs_type_apply,real_mul_closed,function_apply_intro,funs_intro}.
\end{proof}

% from §20 helper lemma 31: coordinate square sequence is a function
\begin{lemma}\label{coord_square_sequence_mul_function}
    Suppose $n\in\naturalsPlus$.
    Let $J = \initialSegment{n}$.
    Suppose $a\in\funs{J}{\reals}$.
    Then $\coordSquareSequenceMul{n}{a}\in\funs{J}{\reals}$.
\end{lemma}
\begin{proof}
    Let $c = \coordSquareSequenceMul{n}{a}$.
    $c\in\funs{J}{\reals}$
        by \cref{coord_square_sequence_mul,pair_eq_iff,relation,rightunique,dom_intro,subseteq,dom_iff,subseteq_antisymmetric,funs_type_apply,real_mul_closed,function_apply_intro,funs_intro}.
\end{proof}

% from §20 helper lemma 32: coordinate difference sequence is a function
\begin{lemma}\label{coord_diff_sequence_function}
    Suppose $n\in\naturalsPlus$.
    Let $J = \initialSegment{n}$.
    Suppose $x\in\realsPower{J}$.
    Suppose $y\in\realsPower{J}$.
    Then $\coordDiffSequence{n}{x}{y}\in\funs{J}{\reals}$.
\end{lemma}
\begin{proof}
    Let $c = \coordDiffSequence{n}{x}{y}$.
    For all $i,z_1,z_2$ such that $(i,z_1)\in c$ and $(i,z_2)\in c$ we have $z_1 = z_2$
        by \cref{coord_diff_sequence,pair_eq_iff}.
    $c$ is right-unique by \cref{rightunique}.
    $c$ is a relation by \cref{coord_diff_sequence,relation}.
    $c$ is a function.
    For all $i$ such that $i\in\dom{c}$ we have $c(i)\in\reals$
        by \cref{dom_iff,coord_diff_sequence,pair_eq_iff,reals_power,tuple_power,funs_type_apply,real_add_inverse,real_add_closed,real_sub,function_apply_intro}.
    $c$ is a function to $\reals$.
    $\dom{c}\subseteq J$ by \cref{dom_iff,coord_diff_sequence,pair_eq_iff,subseteq}.
    For all $i$ such that $i\in J$ we have $i\in\dom{c}$ by \cref{coord_diff_sequence,dom_intro}.
    $c\in\funs{J}{\reals}$ by \cref{subseteq,subseteq_antisymmetric,funs_intro}.
\end{proof}

% from §20 helper lemma 33: coordinate scalar sequence is a function
\begin{lemma}\label{coord_scalar_sequence_function}
    Suppose $n\in\naturalsPlus$.
    Let $J = \initialSegment{n}$.
    Suppose $c\in\reals$.
    Suppose $a\in\funs{J}{\reals}$.
    Then $\coordScalarSequence{n}{c}{a}\in\funs{J}{\reals}$.
\end{lemma}
\begin{proof}
    Let $d = \coordScalarSequence{n}{c}{a}$.
    $d\in\funs{J}{\reals}$
        by \cref{coord_scalar_sequence,pair_eq_iff,relation,rightunique,dom_intro,subseteq,dom_iff,subseteq_antisymmetric,funs_type_apply,real_mul_closed,function_apply_intro,funs_intro}.
\end{proof}

% from §20 helper lemma 34: strict inequality implies weak inequality
\begin{lemma}\label{real_lt_imp_leq}
    Suppose $a < b$.
    Then $a \leq b$.
\end{lemma}
\begin{proof}
    Trivial.
\end{proof}

% from §20 helper lemma 34a: predecessor stays in initial segment
\begin{lemma}\label{initial_segment_pred_from_succ_bound}
    Suppose $n\in\naturalsPlus$.
    Suppose $k\in\naturalsPlus$ and $k + 1 \leq n$.
    Then $k\in\initialSegment{n}$.
\end{lemma}
\begin{proof}
    Follows by \cref{initial_segment_pred_from_succ_bound_aux}.
\end{proof}

% from §20 helper lemma 34b: initial segment is contained in its successor segment
\begin{lemma}\label{initial_segment_subset_succ_helper20}
    Suppose $n\in\naturalsPlus$.
    Then $\initialSegment{n}\subseteq\initialSegment{n + 1}$.
\end{lemma}
\begin{proof}
    Follows by \cref{initial_segment_naturalsplus,real_naturals_subset_reals,subseteq,real_add_ident,real_one_in_naturals,real_zero_lt_one,real_lt_add,real_add_comm,real_add_closed,real_lt_subst_left,real_lt_imp_leq,real_leq_trans}.
\end{proof}

% from §20 helper lemma 34c: exchange middle terms in a sum of four reals
\begin{lemma}\label{real_add_four_exchange}
    Suppose $a\in\reals$.
    Suppose $b\in\reals$.
    Suppose $c\in\reals$.
    Suppose $d\in\reals$.
    Then $(a + b) + (c + d) = (a + c) + (b + d)$.
\end{lemma}
\begin{proof}
    Follows by \cref{real_add_assoc,real_add_comm,real_add_closed}.
\end{proof}

\begin{lemma}\label{real_square_sum_expand}
    Suppose $a\in\reals$.
    Suppose $b\in\reals$.
    Then $(a + b) \rmul (a + b) = (a \rmul a) + ((1 + 1) \rmul (a \rmul b)) + (b \rmul b)$.
\end{lemma}
\begin{proof}
    $1\in\naturalsPlus$ and $1\in\reals$ and $1 + 1\in\reals$
        by \cref{real_one_in_naturals,real_naturals_subset_reals,subseteq,real_add_closed}.
    $a + b\in\reals$ by \cref{real_add_closed}.
    $(a + b) \rmul (a + b) = (a + b) \rmul a + (a + b) \rmul b$ by \cref{real_distrib}.
    $(a + b) \rmul a = a \rmul (a + b)$ by \cref{real_mul_comm}.
    $a \rmul (a + b) = (a \rmul a) + (a \rmul b)$ by \cref{real_distrib}.
    $(a + b) \rmul a = (a \rmul a) + (a \rmul b)$.
    $(a + b) \rmul b = b \rmul (a + b)$ by \cref{real_mul_comm}.
    $b \rmul (a + b) = (b \rmul a) + (b \rmul b)$ by \cref{real_distrib}.
    $(a + b) \rmul b = (b \rmul a) + (b \rmul b)$.
    $b \rmul a = a \rmul b$ and $a \rmul a\in\reals$ and $a \rmul b\in\reals$ and $b \rmul b\in\reals$
        by \cref{real_mul_comm,real_mul_closed}.
    $(a + b) \rmul (a + b) = ((a \rmul a) + (a \rmul b)) + ((a \rmul b) + (b \rmul b))$
        by \cref{real_add_eq_subst}.
    $(a \rmul b) + (b \rmul b)\in\reals$ by \cref{real_add_closed}.
    $(a + b) \rmul (a + b) =
        (a \rmul a) + ((a \rmul b) + ((a \rmul b) + (b \rmul b)))$ by \cref{real_add_assoc}.
    $(a \rmul b) + ((a \rmul b) + (b \rmul b)) =
        ((a \rmul b) + (a \rmul b)) + (b \rmul b)$ by \cref{real_add_assoc}.
    $(a + b) \rmul (a + b) =
        (a \rmul a) + (((a \rmul b) + (a \rmul b)) + (b \rmul b))$ by \cref{real_add_eq_subst}.
    $(a \rmul b) + (a \rmul b)\in\reals$ by \cref{real_add_closed}.
    $(a + b) \rmul (a + b) =
        ((a \rmul a) + ((a \rmul b) + (a \rmul b))) + (b \rmul b)$ by \cref{real_add_assoc}.
    $(1 + 1) \rmul (a \rmul b) = (a \rmul b) + (a \rmul b)$
        by \cref{real_mul_comm,real_distrib,real_mul_ident,real_add_eq_subst}.
    $(a + b) \rmul (a + b) =
        ((a \rmul a) + ((1 + 1) \rmul (a \rmul b))) + (b \rmul b)$ by \cref{real_add_eq_subst}.
    $(a + b) \rmul (a + b) =
        (a \rmul a) + ((1 + 1) \rmul (a \rmul b)) + (b \rmul b)$.
\end{proof}
% from §20 helper lemma 35: sum of a coordinate sum sequence
\begin{lemma}\label{sum_finite_sequence_add}
    Suppose $n\in\naturalsPlus$.
    Let $J = \initialSegment{n}$.
    Suppose $a\in\funs{J}{\reals}$.
    Suppose $b\in\funs{J}{\reals}$.
    Suppose $\sumFiniteSequence{a}{s}$.
    Suppose $\sumFiniteSequence{b}{t}$.
    Let $c = \coordSumSequence{n}{a}{b}$.
    Then $\sumFiniteSequence{c}{s + t}$.
\end{lemma}
\begin{proof}
    Take $n_1$ such that $n_1\in\naturalsPlus$ and $a\in\funs{\initialSegment{n_1}}{\reals}$ and
        there exists $u\in\funs{\initialSegment{n_1}}{\reals}$ such that
            $u(1) = a(1)$ and $s = u(n_1)$ and
            for all $k\in\naturalsPlus$ such that $k + 1 \leq n_1$ we have $u(k + 1) = u(k) + a(k + 1)$
        by \cref{sum_finite_sequence_pred}.
    Take $n_2$ such that $n_2\in\naturalsPlus$ and $b\in\funs{\initialSegment{n_2}}{\reals}$ and
        there exists $v\in\funs{\initialSegment{n_2}}{\reals}$ such that
            $v(1) = b(1)$ and $t = v(n_2)$ and
            for all $k\in\naturalsPlus$ such that $k + 1 \leq n_2$ we have $v(k + 1) = v(k) + b(k + 1)$
        by \cref{sum_finite_sequence_pred}.
    Take $u$ such that $u\in\funs{\initialSegment{n_1}}{\reals}$ and
        $u(1) = a(1)$ and $s = u(n_1)$ and
        for all $k\in\naturalsPlus$ such that $k + 1 \leq n_1$ we have $u(k + 1) = u(k) + a(k + 1)$.
    Take $v$ such that $v\in\funs{\initialSegment{n_2}}{\reals}$ and
        $v(1) = b(1)$ and $t = v(n_2)$ and
        for all $k\in\naturalsPlus$ such that $k + 1 \leq n_2$ we have $v(k + 1) = v(k) + b(k + 1)$.
    $n_1 = n$ and $n_2 = n$
        by \cref{subseteq_antisymmetric,subseteq,funs_elim,initial_segment_naturalsplus,real_naturals_subset_reals,real_leq_antisym}.
    $u\in\funs{J}{\reals}$ and $v\in\funs{J}{\reals}$ by \cref{funs_elim}.
    $c\in\funs{J}{\reals}$ by \cref{coord_sum_sequence_function}.
    Let $w = \coordSumSequence{n}{u}{v}$.
    $w\in\funs{J}{\reals}$ and $w(1) = c(1)$
        by \cref{real_one_in_naturals,real_one_leq_naturalsplus,initial_segment_naturalsplus,coord_sum_sequence,funs_elim,function_apply_intro,coord_sum_sequence_function}.
    Show that for all $k\in\naturalsPlus$ such that $k + 1 \leq n$ we have $w(k + 1) = w(k) + c(k + 1)$.
    \begin{subproof}
        Follows by \cref{initial_segment_pred_from_succ_bound,initial_segment_naturalsplus,real_naturals_closed_succ,funs_type_apply,funs_elim,coord_sum_sequence,function_apply_intro,real_add_eq_subst,real_add_four_exchange}.
    \end{subproof}
    $s + t = w(n)$
        by \cref{initial_segment_naturalsplus,coord_sum_sequence,funs_elim,function_apply_intro}.
    $\sumFiniteSequence{c}{s + t}$ by \cref{coord_sum_sequence_function,sum_finite_sequence_pred}.
\end{proof}

% from §20 helper lemma 36: sum of a coordinate scalar sequence
\begin{lemma}\label{sum_finite_sequence_scalar}
    Suppose $n\in\naturalsPlus$.
    Let $J = \initialSegment{n}$.
    Suppose $c\in\reals$.
    Suppose $a\in\funs{J}{\reals}$.
    Suppose $\sumFiniteSequence{a}{s}$.
    Let $b = \coordScalarSequence{n}{c}{a}$.
    Then $\sumFiniteSequence{b}{c \rmul s}$.
\end{lemma}
\begin{proof}
    Take $n_1$ such that $n_1\in\naturalsPlus$ and $a\in\funs{\initialSegment{n_1}}{\reals}$ and
        there exists $u\in\funs{\initialSegment{n_1}}{\reals}$ such that
            $u(1) = a(1)$ and $s = u(n_1)$ and
            for all $k\in\naturalsPlus$ such that $k + 1 \leq n_1$ we have $u(k + 1) = u(k) + a(k + 1)$
        by \cref{sum_finite_sequence_pred}.
    Take $u$ such that $u\in\funs{\initialSegment{n_1}}{\reals}$ and
        $u(1) = a(1)$ and $s = u(n_1)$ and
        for all $k\in\naturalsPlus$ such that $k + 1 \leq n_1$ we have $u(k + 1) = u(k) + a(k + 1)$.
    $n_1 = n$
        by \cref{subseteq_antisymmetric,subseteq,funs_elim,initial_segment_naturalsplus,real_naturals_subset_reals,real_leq_antisym}.
    Let $w = \coordScalarSequence{n}{c}{u}$.
    $w\in\funs{J}{\reals}$ and $w(1) = b(1)$
        by \cref{real_one_in_naturals,real_one_leq_naturalsplus,initial_segment_naturalsplus,coord_scalar_sequence,funs_elim,function_apply_intro,coord_scalar_sequence_function}.
    Show that for all $k\in\naturalsPlus$ such that $k + 1 \leq n$ we have $w(k + 1) = w(k) + b(k + 1)$.
    \begin{subproof}
        Follows by \cref{real_naturals_closed_succ,initial_segment_naturalsplus,initial_segment_pred_from_succ_bound,coord_scalar_sequence,funs_elim,function_apply_intro,funs_type_apply,real_add_eq_subst,real_distrib,coord_scalar_sequence_function}.
    \end{subproof}
    $c \rmul s = w(n)$
        by \cref{initial_segment_naturalsplus,coord_scalar_sequence,funs_elim,function_apply_intro}.
    $\sumFiniteSequence{b}{c \rmul s}$ by \cref{coord_scalar_sequence_function,sum_finite_sequence_pred}.
\end{proof}

% from §20 helper lemma 37: split sum at the last term
\begin{lemma}\label{sum_finite_sequence_split_last}
    Suppose $n\in\naturalsPlus$.
    Let $J = \initialSegment{n}$.
    Let $J_1 = \initialSegment{n + 1}$.
    Suppose $a\in\funs{J_1}{\reals}$.
    Suppose $\sumFiniteSequence{a}{s}$.
    Let $a_0 = \restrl{a}{J}$.
    Then $a_0\in\funs{J}{\reals}$ and
        there exists $s_0\in\reals$ such that $\sumFiniteSequence{a_0}{s_0}$ and $s = s_0 + a(n + 1)$.
\end{lemma}
\begin{proof}
    $a$ is a function and $\dom{a} = J_1$ by \cref{funs_elim}.
    $a_0$ is a function and $\dom{a_0} = J_1\inter J$ by \cref{restrl_of_function_is_function,dom_of_restrl_eq_inter}.
    $J\subseteq J_1$ by \cref{initial_segment_subset_succ}.
    $\dom{a_0} = J$ by \cref{initial_segment_subset_succ,inter_absorb_supseteq_left,subseteq}.
    For all $i$ such that $i\in\dom{a_0}$ we have $a_0(i)\in\reals$
        by \cref{subseteq,funs_elim,funs_type_apply,restrl_of_function_apply}.
    Show that $a_0\in\funs{J}{\reals}$.
    \begin{subproof}
        Follows by \cref{funs_intro}.
    \end{subproof}
    Take $n_0$ such that $n_0\in\naturalsPlus$ and $a\in\funs{\initialSegment{n_0}}{\reals}$ and
        there exists $t\in\funs{\initialSegment{n_0}}{\reals}$ such that
            $t(1) = a(1)$ and $s = t(n_0)$ and
            for all $k\in\naturalsPlus$ such that $k + 1 \leq n_0$ we have $t(k + 1) = t(k) + a(k + 1)$
        by \cref{sum_finite_sequence_pred}.
    Take $t$ such that $t\in\funs{\initialSegment{n_0}}{\reals}$ and
        $t(1) = a(1)$ and $s = t(n_0)$ and
        for all $k\in\naturalsPlus$ such that $k + 1 \leq n_0$ we have $t(k + 1) = t(k) + a(k + 1)$.
    $\initialSegment{n_0} = J_1$ by \cref{subseteq_antisymmetric,subseteq,funs_elim}.
    $n_0 = n + 1$
        by \cref{initial_segment_naturalsplus,subseteq,real_naturals_closed_succ,real_naturals_subset_reals,real_leq_antisym}.
    Let $t_0 = \restrl{t}{J}$.
    $t$ is a function and $\dom{t} = J_1$ by \cref{funs_elim}.
    $t_0$ is a function and $\dom{t_0} = J_1\inter J$ by \cref{restrl_of_function_is_function,dom_of_restrl_eq_inter}.
    $J\subseteq J_1$ by \cref{initial_segment_subset_succ}.
    $\dom{t_0} = J$ by \cref{initial_segment_subset_succ,inter_absorb_supseteq_left,subseteq}.
    For all $i$ such that $i\in\dom{t_0}$ we have $t_0(i)\in\reals$
        by \cref{subseteq,funs_elim,funs_type_apply,restrl_of_function_apply}.
    $t_0\in\funs{J}{\reals}$ by \cref{funs_intro}.
    $1\in J$ and $t_0(1) = a_0(1)$
        by \cref{real_one_in_naturals,real_one_leq_naturalsplus,initial_segment_naturalsplus,restrl_of_function_apply,subseteq}.
    Show that for all $k\in\naturalsPlus$ such that $k + 1 \leq n$ we have
        $t_0(k + 1) = t_0(k) + a_0(k + 1)$.
    \begin{subproof}
        Follows by \cref{real_naturals_closed_succ,initial_segment_naturalsplus,initial_segment_pred_from_succ_bound,restrl_of_function_apply,subseteq,real_add_eq_subst}.
    \end{subproof}
    Let $s_0 = t_0(n)$.
    $\sumFiniteSequence{a_0}{s_0}$ by \cref{initial_segment_naturalsplus,funs_type_apply,sum_finite_sequence_pred}.
    $t_0(n) + a(n + 1) = s$
        by \cref{initial_segment_naturalsplus,real_naturals_closed_succ,restrl_of_function_apply,subseteq,real_add_eq_subst}.
    Show that there exists $u\in\reals$ such that $\sumFiniteSequence{a_0}{u}$ and $s = u + a(n + 1)$.
    \begin{subproof}
        Follows by \cref{initial_segment_naturalsplus,funs_type_apply,real_add_comm}.
    \end{subproof}
\end{proof}

% from §20 helper lemma 38: zero sum of a nonnegative finite sequence
\begin{lemma}\label{sum_finite_sequence_zero_terms}
    Suppose $n\in\naturalsPlus$.
    Let $J = \initialSegment{n}$.
    Suppose $a\in\funs{J}{\reals}$.
    Suppose $\sumFiniteSequence{a}{s}$.
    Suppose $s = 0$.
    Suppose for all $i\in J$ we have $0 \leq a(i)$.
    Then for all $i\in J$ we have $a(i) = 0$.
\end{lemma}
\begin{proof}
    Follows by \cref{funs_elim,sum_finite_sequence_term_leq,real_add_ident,funs_type_apply,real_leq_antisym}.
\end{proof}

% from §20 helper lemma 39: sum of squares of a scalar sequence
\begin{lemma}\label{sum_finite_sequence_square_scalar}
    Suppose $n\in\naturalsPlus$.
    Let $J = \initialSegment{n}$.
    Suppose $c\in\reals$.
    Suppose $a\in\funs{J}{\reals}$.
    Suppose $\sumFiniteSequence{\coordSquareSequenceMul{n}{a}}{s}$.
    Let $b = \coordScalarSequence{n}{c}{a}$.
    Then $\sumFiniteSequence{\coordSquareSequenceMul{n}{b}}{(c \rmul c) \rmul s}$.
\end{lemma}
\begin{proof}
    Let $u = \coordSquareSequenceMul{n}{a}$.
    Let $v = \coordSquareSequenceMul{n}{b}$.
    Let $w = \coordScalarSequence{n}{c \rmul c}{u}$.
    $b\in\funs{J}{\reals}$ and $u\in\funs{J}{\reals}$ and $v\in\funs{J}{\reals}$ and
        $c \rmul c\in\reals$ and $w\in\funs{J}{\reals}$
        by \cref{coord_scalar_sequence_function,coord_square_sequence_mul_function,real_mul_closed}.
    Show that for all $i\in J$ we have $v(i) = w(i)$.
    \begin{subproof}
        Follows by \cref{coord_scalar_sequence,coord_square_sequence_mul,funs_elim,function_apply_intro,funs_type_apply,real_mul_closed,real_mul_assoc,real_mul_comm}.
    \end{subproof}
    $\sumFiniteSequence{v}{(c \rmul c) \rmul s}$
        by \cref{funs_elim,subseteq,fun_subseteq,subseteq_antisymmetric,sum_finite_sequence_scalar}.
\end{proof}
% from §20 helper lemma 40: sum of squares of a coordinate sum sequence
\begin{lemma}\label{sum_finite_sequence_square_sum}
    Suppose $n\in\naturalsPlus$.
    Let $J = \initialSegment{n}$.
    Suppose $a\in\funs{J}{\reals}$.
    Suppose $b\in\funs{J}{\reals}$.
    Suppose $\sumFiniteSequence{\coordSquareSequenceMul{n}{a}}{s_1}$.
    Suppose $\sumFiniteSequence{\coordSquareSequenceMul{n}{b}}{s_2}$.
    Suppose $\sumFiniteSequence{\coordProductSequence{n}{a}{b}}{s}$.
    Let $c = \coordSumSequence{n}{a}{b}$.
    Then $\sumFiniteSequence{\coordSquareSequenceMul{n}{c}}{s_1 + s_2 + (1+1) \rmul s}$.
\end{lemma}
\begin{proof}
    Let $u = \coordSquareSequenceMul{n}{a}$.
    Let $v = \coordSquareSequenceMul{n}{b}$.
    Let $p = \coordProductSequence{n}{a}{b}$.
    $1\in\naturalsPlus$ and $1\in\reals$ and $1 + 1\in\reals$
        by \cref{real_one_in_naturals,real_naturals_subset_reals,subseteq,real_add_closed}.
    Let $p_2 = \coordScalarSequence{n}{1+1}{p}$.
    Let $w = \coordSumSequence{n}{u}{v}$.
    Let $d = \coordSumSequence{n}{w}{p_2}$.
    Let $q = \coordSquareSequenceMul{n}{c}$.
    $u\in\funs{J}{\reals}$ and $v\in\funs{J}{\reals}$ by \cref{coord_square_sequence_mul_function}.
    $p\in\funs{J}{\reals}$ by \cref{coord_product_sequence_function}.
    $p_2\in\funs{J}{\reals}$ and $w\in\funs{J}{\reals}$ and $d\in\funs{J}{\reals}$ and
        $c\in\funs{J}{\reals}$ and $q\in\funs{J}{\reals}$
        by \cref{coord_scalar_sequence_function,coord_sum_sequence_function,coord_square_sequence_mul_function}.
    Show that for all $i\in J$ we have $q(i) = d(i)$.
    \begin{subproof}
        Fix $i$ such that $i\in J$.
        $a(i)\in\reals$ and $b(i)\in\reals$ and $c(i)\in\reals$ and $u(i)\in\reals$ and $v(i)\in\reals$ and
            $p(i)\in\reals$ and $p_2(i)\in\reals$ and $w(i)\in\reals$
            by \cref{funs_type_apply}.
        $(a(i) \rmul a(i))\in\reals$ and $(b(i) \rmul b(i))\in\reals$ and
            $(a(i) \rmul b(i))\in\reals$ and $((1 + 1) \rmul (a(i) \rmul b(i)))\in\reals$
            by \cref{real_mul_closed}.
        $q(i) = (a(i) + b(i)) \rmul (a(i) + b(i))$
            by \cref{coord_sum_sequence,coord_square_sequence_mul,funs_elim,function_apply_intro}.
        $d(i) = ((a(i) \rmul a(i)) + (b(i) \rmul b(i))) + ((1 + 1) \rmul (a(i) \rmul b(i)))$
            by \cref{coord_sum_sequence,coord_square_sequence_mul,coord_product_sequence,coord_scalar_sequence,funs_elim,function_apply_intro,real_add_eq_subst}.
        $q(i) = ((a(i) \rmul a(i)) + ((1 + 1) \rmul (a(i) \rmul b(i)))) + (b(i) \rmul b(i))$
            by \cref{real_square_sum_expand}.
        $q(i) = (a(i) \rmul a(i)) + (((1 + 1) \rmul (a(i) \rmul b(i))) + (b(i) \rmul b(i)))$
            by \cref{real_add_assoc}.
        $d(i) = (a(i) \rmul a(i)) + (((1 + 1) \rmul (a(i) \rmul b(i))) + (b(i) \rmul b(i)))$
            by \cref{real_add_assoc,real_add_eq_subst,real_add_comm}.
        $q(i) = d(i)$.
    \end{subproof}
    $q = d$ by \cref{funs_elim,subseteq,fun_subseteq,subseteq_antisymmetric}.
    $\sumFiniteSequence{q}{(s_1 + s_2) + (1+1) \rmul s}$
        by \cref{sum_finite_sequence_add,sum_finite_sequence_scalar}.
    $\sumFiniteSequence{\coordSquareSequenceMul{n}{c}}{s_1 + s_2 + (1+1) \rmul s}$
        by \cref{real_add_assoc}.
\end{proof}

% from §20 helper lemma 41a: finite sum value is real
\begin{lemma}\label{sum_finite_sequence_value_in_reals}
    Suppose $n\in\naturalsPlus$.
    Suppose $a\in\funs{\initialSegment{n}}{\reals}$.
    Suppose $\sumFiniteSequence{a}{s}$.
    Then $s\in\reals$.
\end{lemma}
\begin{proof}
    Follows by \cref{sum_finite_sequence_pred,initial_segment_naturalsplus,funs_type_apply}.
\end{proof}

% from §20 helper lemma 41b: coordinate square sequence terms are nonnegative
\begin{lemma}\label{coord_square_sequence_mul_term_nonneg}
    Suppose $n\in\naturalsPlus$.
    Let $J = \initialSegment{n}$.
    Suppose $a\in\funs{J}{\reals}$.
    Then for all $i\in J$ we have $0 \leq \apply{\coordSquareSequenceMul{n}{a}}{i}$.
\end{lemma}
\begin{proof}
    Follows by \cref{coord_square_sequence_mul,coord_square_sequence_mul_function,function_apply_intro,funs_elim,funs_type_apply,real_sq_nonnegative,real_one_in_naturals,real_pow_succ,real_pow_one}.
\end{proof}

% from §20 helper lemma 41c: finite sum of coordinate squares is nonnegative
\begin{lemma}\label{sum_finite_sequence_square_nonneg}
    Suppose $n\in\naturalsPlus$.
    Let $J = \initialSegment{n}$.
    Suppose $a\in\funs{J}{\reals}$.
    Suppose $\sumFiniteSequence{\coordSquareSequenceMul{n}{a}}{s}$.
    Then $0 \leq s$.
\end{lemma}
\begin{proof}
    Follows by \cref{coord_square_sequence_mul_function,funs_elim,coord_square_sequence_mul_term_nonneg,sum_finite_sequence_nonneg}.
\end{proof}

% from §20 helper lemma 41: Cauchy-Schwarz for finite sequences
\begin{lemma}\label{cauchy_schwarz_finite_sequences}
    Suppose $n\in\naturalsPlus$.
    Let $J = \initialSegment{n}$.
    Suppose $a\in\funs{J}{\reals}$.
    Suppose $b\in\funs{J}{\reals}$.
    Suppose $\sumFiniteSequence{\coordProductSequence{n}{a}{b}}{s}$.
    Suppose $\sumFiniteSequence{\coordSquareSequenceMul{n}{a}}{s_1}$.
    Suppose $\sumFiniteSequence{\coordSquareSequenceMul{n}{b}}{s_2}$.
    Then $s \rmul s \leq s_1 \rmul s_2$.
\end{lemma}
\begin{proof}
    Let $a_1 = \coordScalarSequence{n}{s_2}{a}$.
    Let $b_1 = \coordScalarSequence{n}{\neg{s}}{b}$.
    Let $c = \coordSumSequence{n}{a_1}{b_1}$.
    Let $p = \coordProductSequence{n}{a}{b}$.
    Let $p_1 = \coordProductSequence{n}{a_1}{b_1}$.
    $s_2\in\reals$ and $s\in\reals$ and $s_1\in\reals$ and $\neg{s}\in\reals$
        by \cref{coord_square_sequence_mul_function,coord_product_sequence_function,sum_finite_sequence_value_in_reals,real_add_inverse}.
    $a_1\in\funs{J}{\reals}$ and $b_1\in\funs{J}{\reals}$ and
        $p\in\funs{J}{\reals}$ and $p_1\in\funs{J}{\reals}$
        by \cref{coord_scalar_sequence_function,coord_product_sequence_function}.
    $\sumFiniteSequence{\coordSquareSequenceMul{n}{a_1}}{(s_2 \rmul s_2) \rmul s_1}$ by \cref{sum_finite_sequence_square_scalar}.
    $\sumFiniteSequence{\coordSquareSequenceMul{n}{b_1}}{(s \rmul s) \rmul s_2}$
        by \cref{sum_finite_sequence_square_scalar,real_mul_neg_right,real_mul_comm,real_mul_closed,real_neg_neg}.
    Let $r = \coordScalarSequence{n}{s_2 \rmul (\neg{s})}{p}$.
    $s_2 \rmul (\neg{s})\in\reals$ and $r\in\funs{J}{\reals}$
        by \cref{real_mul_closed,coord_scalar_sequence_function}.
    For all $i\in J$ we have $a_1(i) = s_2 \rmul a(i)$ and $b_1(i) = \neg{s} \rmul b(i)$ and
        $p(i) = a(i) \rmul b(i)$
        by \cref{coord_scalar_sequence,coord_product_sequence,funs_elim,function_apply_intro}.
    For all $i\in J$ we have $p_1(i) = r(i)$
        by \cref{coord_scalar_sequence,coord_product_sequence,funs_elim,function_apply_intro,funs_type_apply,real_mul_closed,real_mul_assoc,real_mul_comm}.
    $p_1 = r$ by \cref{funs_elim,subseteq,fun_subseteq,subseteq_antisymmetric}.
    We have $\sumFiniteSequence{r}{(s_2 \rmul (\neg{s})) \rmul s}$ by \cref{sum_finite_sequence_scalar}.
    $\sumFiniteSequence{p_1}{(s_2 \rmul (\neg{s})) \rmul s}$.
    $c\in\funs{J}{\reals}$ and $\coordSquareSequenceMul{n}{c}\in\funs{J}{\reals}$
        by \cref{coord_sum_sequence_function,coord_square_sequence_mul_function}.
    Take $S\in\reals$ such that $\sumFiniteSequence{\coordSquareSequenceMul{n}{c}}{S}$ by \cref{sum_finite_sequence_exists}.
    $\sumFiniteSequence{\coordSquareSequenceMul{n}{c}}{(s_2 \rmul s_2) \rmul s_1 + (s \rmul s) \rmul s_2 + (1+1) \rmul ((s_2 \rmul (\neg{s})) \rmul s)}$
        by \cref{sum_finite_sequence_square_sum}.
    $S = (s_2 \rmul s_2) \rmul s_1 + (s \rmul s) \rmul s_2 + (1+1) \rmul ((s_2 \rmul (\neg{s})) \rmul s)$
        by \cref{sum_finite_sequence_unique}.
    We have $0 \leq S$ by \cref{sum_finite_sequence_square_nonneg}.
    $(s_2 \rmul (\neg{s})) \rmul s = s_2 \rmul ((\neg{s}) \rmul s)$ by \cref{real_mul_assoc}.
    $(\neg{s}) \rmul s = s \rmul (\neg{s})$ by \cref{real_mul_comm}.
    $(s_2 \rmul (\neg{s})) \rmul s = \neg{(s_2 \rmul (s \rmul s))}$
        by \cref{real_mul_assoc,real_mul_closed,real_mul_neg_right}.
    $(1+1) \rmul ((s_2 \rmul (\neg{s})) \rmul s) =
        (1+1) \rmul \neg{(s_2 \rmul (s \rmul s))}$.
    $1\in\naturalsPlus$ and $1\in\reals$ and $1 + 1\in\reals$ by \cref{real_one_in_naturals,real_naturals_subset_reals,subseteq,real_add_closed}.
    We have $s_2 \rmul (s \rmul s)\in\reals$ by \cref{real_mul_closed}.
    $(1+1) \rmul \neg{(s_2 \rmul (s \rmul s))} =
        \neg{((1+1) \rmul (s_2 \rmul (s \rmul s)))}$ by \cref{real_mul_neg_right}.
    $S = (s_2 \rmul s_2) \rmul s_1 + (s \rmul s) \rmul s_2 +
        \neg{((1+1) \rmul (s_2 \rmul (s \rmul s)))}$ by \cref{real_add_eq_subst}.
    $s \rmul s\in\reals$ and $(s \rmul s) \rmul s_2 = s_2 \rmul (s \rmul s)$ by \cref{real_mul_closed,real_mul_comm}.
    $S = s_2 \rmul (s_2 \rmul s_1) + s_2 \rmul (s \rmul s) +
        \neg{((1+1) \rmul (s_2 \rmul (s \rmul s)))}$ by \cref{real_mul_assoc}.
    We have $s_2 \rmul (s_2 \rmul s_1)\in\reals$ and $s_2 \rmul (s \rmul s)\in\reals$ by \cref{real_mul_closed}.
    $\neg{((1+1) \rmul (s_2 \rmul (s \rmul s)))}\in\reals$ by \cref{real_add_inverse,real_mul_closed}.
    $S = s_2 \rmul (s_2 \rmul s_1) +
        (s_2 \rmul (s \rmul s) + \neg{((1+1) \rmul (s_2 \rmul (s \rmul s)))})$ by \cref{real_add_assoc}.
    We have $(1+1) \rmul (s_2 \rmul (s \rmul s)) =
        (s_2 \rmul (s \rmul s)) + (s_2 \rmul (s \rmul s))$
        by \cref{real_mul_comm,real_distrib,real_mul_ident,real_add_eq_subst}.
    $\neg{((1+1) \rmul (s_2 \rmul (s \rmul s)))} =
        \neg{(s_2 \rmul (s \rmul s))} + \neg{(s_2 \rmul (s \rmul s))}$
        by \cref{real_add_eq_subst,real_neg_addition}.
    $s_2 \rmul (s \rmul s) + \neg{((1+1) \rmul (s_2 \rmul (s \rmul s)))} =
        \neg{(s_2 \rmul (s \rmul s))}$
        by \cref{real_add_eq_subst,real_add_assoc,real_add_inverse,real_add_ident}.
    $S = s_2 \rmul (s_2 \rmul s_1) + \neg{(s_2 \rmul (s \rmul s))}$ by \cref{real_add_eq_subst}.
    We have $s_2 \rmul s_1\in\reals$ and $s \rmul s\in\reals$ and $\neg{(s \rmul s)}\in\reals$
        by \cref{real_mul_closed,real_add_inverse}.
    $s_2 \rmul (s_2 \rmul s_1 + \neg{(s \rmul s)}) =
        s_2 \rmul (s_2 \rmul s_1) + s_2 \rmul \neg{(s \rmul s)}$ by \cref{real_distrib}.
    $s_2 \rmul \neg{(s \rmul s)} = \neg{(s_2 \rmul (s \rmul s))}$ by \cref{real_mul_neg_right}.
    Hence $s_2 \rmul (s_2 \rmul s_1 + \neg{(s \rmul s)}) =
        s_2 \rmul (s_2 \rmul s_1) + \neg{(s_2 \rmul (s \rmul s))}$ by \cref{real_add_eq_subst}.
    $S = s_2 \rmul (s_2 \rmul s_1 + \neg{(s \rmul s)})$.
    Hence $S = s_2 \rmul (s_2 \rmul s_1 - s \rmul s)$ by \cref{real_sub}.
    We have $0 \leq s_2$ by \cref{sum_finite_sequence_square_nonneg}.
    $0 < s_2$ or $s_2 = 0$ by \cref{real_leq_split}.
    \begin{byCase}
    \caseOf{$s_2 = 0$.}
            For all $i\in J$ we have $b(i) = 0$
                by \cref{coord_square_sequence_mul_term_nonneg,coord_square_sequence_mul_function,sum_finite_sequence_zero_terms,coord_square_sequence_mul,function_apply_intro,funs_elim,funs_type_apply,real_mul_zero_product}.
            For all $i\in J$ we have $p(i) = 0$
                by \cref{coord_product_sequence,funs_elim,function_apply_intro,funs_type_apply,real_mul_comm,real_mul_zero}.
            $0\in\reals$ by \cref{real_add_ident}.
            We have $\sumFiniteSequence{p}{s}$.
            $s \leq 0$ by \cref{sum_finite_sequence_leq_n_times,real_naturals_subset_reals,subseteq,real_mul_comm,real_mul_zero}.
            For all $i\in\dom{p}$ we have $0 \leq p(i)$ by \cref{funs_elim,subseteq}.
            $0 \leq s$ by \cref{sum_finite_sequence_nonneg}.
            Hence $s = 0$ by \cref{real_leq_antisym}.
            $s \rmul s = s_1 \rmul s_2$ by \cref{real_mul_comm,real_mul_zero}.
            Thus $s \rmul s \leq s_1 \rmul s_2$.
        \caseOf{$0 < s_2$.}
            We have $s \rmul s < s_2 \rmul s_1$ or $s \rmul s = s_2 \rmul s_1$ or $s_2 \rmul s_1 < s \rmul s$
                by \cref{real_lt_trichotomy}.
            \begin{byCase}
                \caseOf{$s \rmul s < s_2 \rmul s_1$.}
                    Show that $s \rmul s \leq s_1 \rmul s_2$.
                    \begin{subproof}
                        Follows by \cref{real_lt_imp_leq,real_mul_comm}.
                    \end{subproof}
                \caseOf{$s \rmul s = s_2 \rmul s_1$.}
                    Show that $s \rmul s \leq s_1 \rmul s_2$.
                    \begin{subproof}
                        Follows by \cref{real_mul_comm}.
                    \end{subproof}
                \caseOf{$s_2 \rmul s_1 < s \rmul s$.}
                    Show that $s \rmul s \leq s_1 \rmul s_2$.
                    \begin{subproof}
                        Suppose not.
                        We have $0 < s \rmul s - s_2 \rmul s_1$ by \cref{real_lt_imp_pos_diff}.
                        $s \rmul s - s_2 \rmul s_1\in\reals$ by \cref{real_sub,real_add_closed,real_add_inverse,real_mul_closed}.
                        $s_2 \rmul s_1 - s \rmul s = \neg{(s \rmul s - s_2 \rmul s_1)}$ by \cref{real_sub_swap_neg}.
                        $s_2 \rmul s_1 - s \rmul s < 0$ by \cref{real_pos_imp_neg}.
                        $0\in\reals$ and $s_2 \rmul s_1\in\reals$ and $s \rmul s\in\reals$ and
                            $\neg{(s \rmul s)}\in\reals$ and $s_2 \rmul s_1 - s \rmul s\in\reals$
                            by \cref{real_add_ident,real_mul_closed,real_add_inverse,real_sub,real_add_closed}.
                        $(s_2 \rmul s_1 - s \rmul s) \rmul s_2 < 0 \rmul s_2$ by \cref{real_lt_mul_pos}.
                        $0 \rmul s_2 = 0$ by \cref{real_mul_zero}.
                        Thus $(s_2 \rmul s_1 - s \rmul s) \rmul s_2 < 0$ by \cref{real_lt_subst_right}.
                        $(s_2 \rmul s_1 - s \rmul s) \rmul s_2 = s_2 \rmul (s_2 \rmul s_1 - s \rmul s)$ by \cref{real_mul_comm}.
                        Hence $s_2 \rmul (s_2 \rmul s_1 - s \rmul s) < 0$ by \cref{real_lt_subst_left}.
                        $S < 0$ by \cref{real_lt_subst_left}.
                        Therefore $0 < 0$ by \cref{real_lt_trans,real_lt_subst_right}.
                        Contradiction by \cref{real_lt_irreflexive}.
                    \end{subproof}
            \end{byCase}
    \end{byCase}
\end{proof}

% from §20 helper lemma 42: Minkowski for finite sequences
\begin{lemma}\label{minkowski_finite_sequences}
    Suppose $n\in\naturalsPlus$.
    Let $J = \initialSegment{n}$.
    Suppose $a\in\funs{J}{\reals}$.
    Suppose $b\in\funs{J}{\reals}$.
    Suppose $\sumFiniteSequence{\coordSquareSequenceMul{n}{a}}{s_1}$.
    Suppose $\sumFiniteSequence{\coordSquareSequenceMul{n}{b}}{s_2}$.
    Let $c = \coordSumSequence{n}{a}{b}$.
    Suppose $\sumFiniteSequence{\coordSquareSequenceMul{n}{c}}{s}$.
    Then $\sqrt{s} \leq \sqrt{s_1} + \sqrt{s_2}$.
\end{lemma}
\begin{proof}
    Let $p = \coordProductSequence{n}{a}{b}$.
    Take $t\in\reals$ such that $\sumFiniteSequence{p}{t}$
        by \cref{coord_product_sequence_function,sum_finite_sequence_exists}.
    $\sumFiniteSequence{\coordSquareSequenceMul{n}{c}}{s_1 + s_2 + (1+1) \rmul t}$ by \cref{sum_finite_sequence_square_sum}.
    $s = s_1 + s_2 + (1+1) \rmul t$ by \cref{sum_finite_sequence_unique}.
    $s_1\in\reals$ and $s_2\in\reals$ by \cref{coord_square_sequence_mul_function,sum_finite_sequence_value_in_reals}.
    $1\in\naturalsPlus$ and $1\in\reals$ and $1 + 1\in\reals$
        by \cref{real_one_in_naturals,real_naturals_subset_reals,subseteq,real_add_closed}.
    $(1+1) \rmul t\in\reals$ and $s_1 + s_2\in\reals$ and $s\in\reals$
        by \cref{real_one_in_naturals,real_naturals_subset_reals,subseteq,real_add_closed,real_mul_closed}.
    $0 \leq s_1$ and $0 \leq s_2$ by \cref{sum_finite_sequence_square_nonneg}.
    $0 \leq s$ by \cref{coord_sum_sequence_function,sum_finite_sequence_square_nonneg}.
    Let $r = \sqrt{s}$.
    Let $r_1 = \sqrt{s_1}$.
    Let $r_2 = \sqrt{s_2}$.
    $r\in\reals$ and $r \geq 0$ and $r \rmul r = s$ by \cref{positive_square_root}.
    $r_1\in\reals$ and $r_1 \geq 0$ and $r_1 \rmul r_1 = s_1$ by \cref{positive_square_root}.
    $r_2\in\reals$ and $r_2 \geq 0$ and $r_2 \rmul r_2 = s_2$ by \cref{positive_square_root}.
    $\abs{r_1} = r_1$ and $\abs{r_2} = r_2$ by \cref{real_abs_nonneg}.
    $\abs{r_1 \rmul r_2} = r_1 \rmul r_2$ by \cref{real_abs_mul}.
    $r_1 \rmul r_2\in\reals$ by \cref{real_mul_closed}.
    We have $\abs{r_1 \rmul r_2} \geq 0$ by \cref{real_abs_nonnegative}.
    $0 \leq r_1 \rmul r_2$.
    $\abs{t}\in\reals$ and $\abs{t} \geq 0$ and $r_1 \rmul r_2\in\reals$
        by \cref{real_abs_in_reals,real_abs_nonnegative,real_mul_closed}.
    $\abs{t} \leq r_1 \rmul r_2$ iff
        $\exp{\abs{t}}{1+1} \leq \exp{r_1 \rmul r_2}{1+1}$ by \cref{real_nonneg_leq_iff_sq_leq}.
    $\exp{\abs{t}}{1+1} = t \rmul t$ by \cref{real_abs_square,real_pow_succ,real_pow_one}.
    $t \rmul t \leq s_1 \rmul s_2$ by \cref{cauchy_schwarz_finite_sequences}.
    $\exp{r_1 \rmul r_2}{1+1} = s_1 \rmul s_2$
        by \cref{real_pow_succ,real_pow_one,real_mul_assoc,real_mul_comm}.
    Hence $\exp{\abs{t}}{1+1} \leq \exp{r_1 \rmul r_2}{1+1}$ by \cref{real_leq_trans}.
    $\abs{t} \leq r_1 \rmul r_2$ by \cref{real_nonneg_leq_iff_sq_leq}.
    $t \leq r_1 \rmul r_2$ by \cref{real_leq_abs,real_leq_trans}.
    We have $0 < 1 + 1$
        by \cref{real_add_ident,real_naturals_subset_reals,real_one_in_naturals,subseteq,real_zero_lt_one,real_lt_add,real_lt_subst_left,real_lt_trans}.
    $(1+1) \rmul t \leq (1+1) \rmul (r_1 \rmul r_2)$
        by \cref{real_mul_closed,real_add_closed,real_leq_split,real_lt_mul_pos,real_mul_comm,real_lt_subst_left,real_lt_subst_right,real_lt_imp_leq}.
    We have $(1+1) \rmul t\in\reals$ and $(1+1) \rmul (r_1 \rmul r_2)\in\reals$ by \cref{real_mul_closed}.
    $(1+1) \rmul t + s_1 \leq (1+1) \rmul (r_1 \rmul r_2) + s_1$ by \cref{real_leq_add}.
    We have $(1+1) \rmul t + s_1 = s_1 + (1+1) \rmul t$ and
        $(1+1) \rmul (r_1 \rmul r_2) + s_1 = s_1 + (1+1) \rmul (r_1 \rmul r_2)$ by \cref{real_add_comm}.
    Hence $s_1 + (1+1) \rmul t \leq s_1 + (1+1) \rmul (r_1 \rmul r_2)$ by \cref{real_leq_trans}.
    We have $s_1 + (1+1) \rmul t\in\reals$ and $s_1 + (1+1) \rmul (r_1 \rmul r_2)\in\reals$ by \cref{real_add_closed}.
    $(s_1 + (1+1) \rmul t) + s_2 \leq (s_1 + (1+1) \rmul (r_1 \rmul r_2)) + s_2$ by \cref{real_leq_add}.
    We have $(s_1 + (1+1) \rmul t) + s_2 = s_1 + ((1+1) \rmul t + s_2)$ by \cref{real_add_assoc}.
    We have $(1+1) \rmul t + s_2 = s_2 + (1+1) \rmul t$ by \cref{real_add_comm}.
    $s_1 + ((1+1) \rmul t + s_2) = s_1 + (s_2 + (1+1) \rmul t)$ by \cref{real_add_eq_subst}.
    Hence $s_1 + (s_2 + (1+1) \rmul t) = (s_1 + s_2) + (1+1) \rmul t$ by \cref{real_add_assoc}.
    $(s_1 + (1+1) \rmul t) + s_2 = (s_1 + s_2) + (1+1) \rmul t$.
    We have $(s_1 + (1+1) \rmul (r_1 \rmul r_2)) + s_2 = s_1 + ((1+1) \rmul (r_1 \rmul r_2) + s_2)$ and
        $(1+1) \rmul (r_1 \rmul r_2) + s_2 = s_2 + (1+1) \rmul (r_1 \rmul r_2)$ by \cref{real_add_assoc,real_add_comm}.
    $s_1 + ((1+1) \rmul (r_1 \rmul r_2) + s_2) = s_1 + (s_2 + (1+1) \rmul (r_1 \rmul r_2))$ by \cref{real_add_eq_subst}.
    Hence $s_1 + (s_2 + (1+1) \rmul (r_1 \rmul r_2)) = (s_1 + s_2) + (1+1) \rmul (r_1 \rmul r_2)$ by \cref{real_add_assoc}.
    $(s_1 + (1+1) \rmul (r_1 \rmul r_2)) + s_2 = (s_1 + s_2) + (1+1) \rmul (r_1 \rmul r_2)$.
    $(s_1 + s_2) + (1+1) \rmul t \leq (s_1 + s_2) + (1+1) \rmul (r_1 \rmul r_2)$ by \cref{real_leq_trans}.
    Hence $s \leq (s_1 + s_2) + (1+1) \rmul (r_1 \rmul r_2)$ by \cref{real_leq_trans}.
    We have $r_1 + r_2\in\reals$ by \cref{real_add_closed}.
    $(r_1 + r_2) \rmul (r_1 + r_2) =
        (r_1 \rmul r_1) + ((1 + 1) \rmul (r_1 \rmul r_2)) + (r_2 \rmul r_2)$ by \cref{real_square_sum_expand}.
    We have $r_1 \rmul r_1 = s_1$ and $r_2 \rmul r_2 = s_2$.
    Hence $(r_1 + r_2) \rmul (r_1 + r_2) =
        s_1 + ((1 + 1) \rmul (r_1 \rmul r_2)) + s_2$ by \cref{real_add_eq_subst}.
    $(r_1 + r_2) \rmul (r_1 + r_2) =
        s_1 + s_2 + (1 + 1) \rmul (r_1 \rmul r_2)$ by \cref{real_add_assoc,real_add_comm}.
    We have $r_1 + r_2\in\reals$ and $0\in\reals$ and $0 \leq r_1$
        by \cref{real_add_closed,real_add_ident}.
    $0 + r_2 \leq r_1 + r_2$ and $0 + r_2 = r_2$ by \cref{real_leq_add,real_add_ident}.
    Hence $r_2 \leq r_1 + r_2$ and $0 \leq r_2$ and $0 \leq r_1 + r_2$ by \cref{real_leq_trans}.
    We have $\exp{r}{1+1} = s$ by \cref{real_pow_succ,real_pow_one}.
    We have $\exp{r_1 + r_2}{1+1} = s_1 + s_2 + (1 + 1) \rmul (r_1 \rmul r_2)$
        by \cref{real_pow_succ,real_pow_one,real_add_eq_subst}.
    $\exp{r}{1+1} \leq \exp{r_1 + r_2}{1+1}$ by \cref{real_leq_trans}.
    Hence $r \leq r_1 + r_2$ by \cref{real_nonneg_leq_iff_sq_leq}.
    $\sqrt{s} \leq \sqrt{s_1} + \sqrt{s_2}$.
\end{proof}

% from §20 helper lemma 43: euclidean metric value exists
\begin{lemma}\label{euclidean_metric_value_exists}
    Suppose $n\in\naturalsPlus$.
    Let $J = \initialSegment{n}$.
    Suppose $x\in\realsPower{J}$.
    Suppose $y\in\realsPower{J}$.
    Then there exists $r\in\reals$ such that $((x,y),r)\in\euclideanMetric{n}$.
\end{lemma}
\begin{proof}
    Let $a = \coordSquareSequence{n}{(x,y)}$.
    $a\in\funs{J}{\reals}$ by \cref{coord_square_sequence_function}.
    Take $s\in\reals$ such that $\sumFiniteSequence{a}{s}$ by \cref{sum_finite_sequence_exists}.
    For all $i\in\dom{a}$ we have $0 \leq a(i)$
        by \cref{funs_elim,reals_power,tuple_power,funs_type_apply,real_add_inverse,real_add_closed,real_sub,real_sq_nonnegative,coord_square_sequence_elem,function_apply_intro}.
    $0 \leq s$ by \cref{sum_finite_sequence_nonneg}.
    There exists $r\in\reals$ such that $((x,y),r)\in\euclideanMetric{n}$
        by \cref{positive_square_root,times_tuple_intro,euclidean_metric}.
\end{proof}
% from §20 helper lemma 44: euclidean metric is a function
\begin{lemma}\label{euclidean_metric_function}
    Suppose $n\in\naturalsPlus$.
    Then $\euclideanMetric{n}$ is a function.
\end{lemma}
\begin{proof}
    Follows by \cref{euclidean_metric,relation,rightunique,pair_eq_iff,sum_finite_sequence_unique}.
\end{proof}

% from §20 helper lemma 44a: chained subtraction
\begin{lemma}\label{real_sub_add_chain}
    Suppose $x\in\reals$.
    Suppose $y\in\reals$.
    Suppose $z\in\reals$.
    Then $(x - y) + (y - z) = x - z$.
\end{lemma}
\begin{proof}
    Follows by \cref{real_add_inverse,real_sub,real_add_assoc,real_add_closed,real_add_comm,real_add_eq_subst,real_add_ident}.
\end{proof}

% from §20 helper lemma 45: square metric is a metric
\begin{lemma}\label{square_metric_is_metric}
    Suppose $n\in\naturalsPlus$.
    Let $J = \initialSegment{n}$.
    Let $X = \realsPower{J}$.
    Then $\squareMetric{n}$ is a metric on $X$.
\end{lemma}
\begin{proof}
    Let $d = \squareMetric{n}$.
    We have $d$ is a function and $X\times X\subseteq\dom{d}$
        by \cref{square_metric_function,times_elem_is_tuple,square_metric_value_exists,pair_eq_iff,dom_intro,subseteq}.
    $\dom{d}\subseteq X\times X$ by \cref{dom_iff,square_metric,pair_eq_iff,times_tuple_intro,subseteq}.
    Hence $\dom{d} = X\times X$ by \cref{subseteq_antisymmetric}.
    For all $p\in\dom{d}$ we have $d(p)\in\reals$
        by \cref{dom_iff,square_metric,pair_eq_iff,coord_diff_set_finite,subseteq,function_apply_intro}.
    $d\in\funs{X\times X}{\reals}$ by \cref{funs_intro}.
    Show that for all $x,y$ such that $x\in X$ and $y\in X$ we have $d((x,y)) \geq 0$.
    \begin{subproof}
        Follows by \cref{square_metric_value_exists,square_metric,pair_eq_iff,coord_diff_set,reals_power,tuple_power,funs_type_apply,real_sub,real_add_closed,real_add_inverse,real_abs_nonnegative,function_apply_intro}.
    \end{subproof}
    $d$ is nonnegative on $X$ by \cref{metric_nonnegative}.
    Show that for all $x,y$ such that $x\in X$ and $y\in X$ we have if $d((x,y)) = 0$ then $x = y$.
    \begin{subproof}
        Follows by \cref{square_metric_value_exists,function_apply_intro,reals_power,tuple_power,funs_type_apply,real_sub,real_add_closed,real_add_inverse,real_abs_in_reals,real_abs_nonnegative,real_leq_antisym,real_abs_eq_zero,real_add_cancel,funs_elim,subseteq,fun_subseteq,subseteq_antisymmetric}.
    \end{subproof}
    Show that for all $x,y$ such that $x\in X$ and $y\in X$ we have if $x = y$ then $d((x,y)) = 0$.
    \begin{subproof}
        Follows by \cref{square_metric_value_exists,coord_diff_set,reals_power,tuple_power,funs_type_apply,real_sub,real_add_inverse,real_add_eq_subst,real_add_ident,real_abs_nonneg,square_metric,pair_eq_iff,function_apply_intro}.
    \end{subproof}
    $d$ is zero-characterized on $X$ by \cref{metric_zero_characterization}.
    Show that for all $x,y$ such that $x\in X$ and $y\in X$ we have $d((x,y)) = d((y,x))$.
    \begin{subproof}
        Follows by \cref{square_metric_value_exists,coord_diff_set,reals_power,tuple_power,funs_type_apply,real_abs_sub_swap,square_metric,pair_eq_iff,real_leq_antisym,function_apply_intro}.
    \end{subproof}
    $d$ is symmetric on $X$ by \cref{metric_symmetric}.
    Show that for all $x,y,z$ such that $x\in X$ and $y\in X$ and $z\in X$ we have
        $d((x,y)) + d((y,z)) \geq d((x,z))$.
    \begin{subproof}
            Fix $x$.
            Fix $y$.
            Fix $z$ such that $x\in X$ and $y\in X$ and $z\in X$.
            Take $r_1\in\reals$ such that $((x,y),r_1)\in d$ and
                for all $i\in J$ we have $\abs{x(i) - y(i)} \leq r_1$
                by \cref{square_metric_value_exists}.
            Take $r_2\in\reals$ such that $((y,z),r_2)\in d$ and
                for all $i\in J$ we have $\abs{y(i) - z(i)} \leq r_2$
                by \cref{square_metric_value_exists}.
            Take $r_3\in\reals$ such that $((x,z),r_3)\in d$ and
                for all $i\in J$ we have $\abs{x(i) - z(i)} \leq r_3$
                by \cref{square_metric_value_exists}.
            $r_3\in\coordDiffSet{n}{x}{z}$ by \cref{square_metric,pair_eq_iff}.
            Take $i_0\in J$ such that $r_3 = \abs{x(i_0) - z(i_0)}$ by \cref{coord_diff_set}.
            $x(i_0)\in\reals$ and $y(i_0)\in\reals$ and $z(i_0)\in\reals$ and
                $\neg{y(i_0)}\in\reals$ and $\neg{z(i_0)}\in\reals$ and
                $x(i_0) - y(i_0)\in\reals$ and $y(i_0) - z(i_0)\in\reals$
                by \cref{reals_power,tuple_power,funs_type_apply,real_add_inverse,real_sub,real_add_closed}.
            $(x(i_0) - y(i_0)) + (y(i_0) - z(i_0)) = x(i_0) - z(i_0)$
                by \cref{real_sub_add_chain}.
            $\abs{(x(i_0) - y(i_0)) + (y(i_0) - z(i_0))} \leq \abs{x(i_0) - y(i_0)} + \abs{y(i_0) - z(i_0)}$ by \cref{real_triangle_inequality}.
            Hence $\abs{x(i_0) - z(i_0)} \leq \abs{x(i_0) - y(i_0)} + \abs{y(i_0) - z(i_0)}$.
            $r_3 \leq \abs{x(i_0) - y(i_0)} + \abs{y(i_0) - z(i_0)}$.
            $\abs{x(i_0) - y(i_0)} \leq r_1$.
            $\abs{x(i_0) - y(i_0)}\in\reals$ and $\abs{y(i_0) - z(i_0)}\in\reals$ by \cref{real_abs_in_reals}.
            $\abs{x(i_0) - y(i_0)} + \abs{y(i_0) - z(i_0)} \leq r_1 + \abs{y(i_0) - z(i_0)}$ and
                $\abs{x(i_0) - y(i_0)} + \abs{y(i_0) - z(i_0)}\in\reals$ and
                $r_1 + \abs{y(i_0) - z(i_0)}\in\reals$ by \cref{real_leq_add,real_add_closed}.
            Hence $r_3 \leq r_1 + \abs{y(i_0) - z(i_0)}$ by \cref{real_leq_trans}.
            $\abs{y(i_0) - z(i_0)} \leq r_2$.
            $r_1 + \abs{y(i_0) - z(i_0)} \leq r_1 + r_2$
                by \cref{real_leq_add,real_add_comm}.
            $r_1 + r_2\in\reals$ by \cref{real_add_closed}.
            Hence $d((x,z)) \leq d((x,y)) + d((y,z))$ by \cref{real_leq_trans,function_apply_intro}.
            $d((x,y)) + d((y,z)) \geq d((x,z))$.
    \end{subproof}
    $d$ satisfies the triangle inequality on $X$ by \cref{metric_triangle_inequality}.
    Hence $d$ is a metric on $X$ by \cref{metric_on}.
\end{proof}

% from §20 helper lemma 46: square sequence of a difference sequence
\begin{lemma}\label{coord_square_sequence_diff}
    Suppose $n\in\naturalsPlus$.
    Let $J = \initialSegment{n}$.
    Suppose $x\in\realsPower{J}$.
    Suppose $y\in\realsPower{J}$.
    Let $a = \coordDiffSequence{n}{x}{y}$.
    Then $\coordSquareSequenceMul{n}{a} = \coordSquareSequence{n}{(x,y)}$.
\end{lemma}
\begin{proof}
    Follows by \cref{setext,coord_square_sequence_mul,coord_diff_sequence_function,funs_elim,coord_diff_sequence,function_apply_intro,reals_power,tuple_power,funs_type_apply,real_add_inverse,real_add_closed,real_sub,real_one_in_naturals,real_pow_succ,real_pow_one,coord_square_sequence,fst_eq,snd_eq}.
\end{proof}

% from §20 helper lemma 47: difference sequence of a sum
\begin{lemma}\label{coord_diff_sequence_sum}
    Suppose $n\in\naturalsPlus$.
    Let $J = \initialSegment{n}$.
    Suppose $x\in\realsPower{J}$.
    Suppose $y\in\realsPower{J}$.
    Suppose $z\in\realsPower{J}$.
    Let $a = \coordDiffSequence{n}{x}{y}$.
    Let $b = \coordDiffSequence{n}{y}{z}$.
    Let $c = \coordSumSequence{n}{a}{b}$.
    Then $c = \coordDiffSequence{n}{x}{z}$.
\end{lemma}
\begin{proof}
    Follows by \cref{setext,coord_sum_sequence,coord_diff_sequence,coord_diff_sequence_function,funs_elim,function_apply_intro,reals_power,tuple_power,funs_type_apply,real_sub_add_chain}.
\end{proof}

% from §20 Item 25: euclidean metric is a metric
\begin{theorem}\label{euclidean_metric_is_metric}
    Suppose $n\in\naturalsPlus$.
    Let $J = \initialSegment{n}$.
    Let $X = \realsPower{J}$.
    Then $\euclideanMetric{n}$ is a metric on $X$.
\end{theorem}
\begin{proof}
    Let $d = \euclideanMetric{n}$.
    We have $d$ is a function and $X\times X\subseteq\dom{d}$
        by \cref{euclidean_metric_function,times_elem_is_tuple,euclidean_metric_value_exists,pair_eq_iff,dom_intro,subseteq}.
    $\dom{d}\subseteq X\times X$
        by \cref{dom_iff,euclidean_metric,pair_eq_iff,times_elem_is_tuple,times_tuple_intro,subseteq}.
    Thus $\dom{d} = X\times X$ by \cref{subseteq_antisymmetric}.
    For all $p\in\dom{d}$ we have $d(p)\in\reals$
        by \cref{dom_iff,euclidean_metric,pair_eq_iff,function_apply_intro}.
    $d\in\funs{X\times X}{\reals}$ by \cref{funs_intro}.
    Show that for all $x,y$ such that $x\in X$ and $y\in X$ we have $d((x,y)) \geq 0$.
    \begin{subproof}
        Follows by
        \cref{euclidean_metric_value_exists,euclidean_metric,pair_eq_iff,coord_square_sequence_diff,coord_diff_sequence_function,sum_finite_sequence_square_nonneg,positive_square_root,function_apply_intro}.
    \end{subproof}
    $d$ is nonnegative on $X$ by \cref{metric_nonnegative}.
    Show that for all $x,y$ such that $x\in X$ and $y\in X$ we have if $d((x,y)) = 0$ then $x = y$.
    \begin{subproof}
        Follows by \cref{euclidean_metric_value_exists,function_apply_intro,square_metric_value_exists,euclidean_square_metric_bounds,square_metric_is_metric,metric_on,metric_nonnegative,funs_elim,real_leq_antisym,metric_zero_characterization,real_add_ident,real_zero_lt_one,real_lt_imp_leq,real_one_leq_naturalsplus,real_naturals_subset_reals,real_one_in_naturals,subseteq,real_leq_trans,positive_square_root,real_mul_comm,real_mul_zero}.
    \end{subproof}
    Show that for all $x,y$ such that $x\in X$ and $y\in X$ we have if $x = y$ then $d((x,y)) = 0$.
    \begin{subproof}
        Follows by \cref{euclidean_metric_value_exists,function_apply_intro,square_metric_value_exists,euclidean_square_metric_bounds,square_metric_is_metric,metric_on,metric_nonnegative,funs_elim,real_leq_antisym,metric_zero_characterization,real_add_ident,real_zero_lt_one,real_lt_imp_leq,real_one_leq_naturalsplus,real_naturals_subset_reals,real_one_in_naturals,subseteq,real_leq_trans,positive_square_root,real_mul_comm,real_mul_zero}.
    \end{subproof}
    $d$ is zero-characterized on $X$ by \cref{metric_zero_characterization}.
    Show that for all $x,y$ such that $x\in X$ and $y\in X$ we have $d((x,y)) = d((y,x))$.
    \begin{subproof}
            Fix $x$.
            Fix $y$ such that $x\in X$ and $y\in X$.
            Take $r\in\reals$ such that $((x,y),r)\in d$ by \cref{euclidean_metric_value_exists}.
            Take $s\in\reals$ such that $((y,x),s)\in d$ by \cref{euclidean_metric_value_exists}.
            Take $s_1\in\reals$ such that
                $\sumFiniteSequence{\coordSquareSequence{n}{(x,y)}}{s_1}$ and $r = \sqrt{s_1}$
                by \cref{euclidean_metric,pair_eq_iff}.
            Take $s_2\in\reals$ such that
                $\sumFiniteSequence{\coordSquareSequence{n}{(y,x)}}{s_2}$ and $s = \sqrt{s_2}$
                by \cref{euclidean_metric,pair_eq_iff}.
            $\coordSquareSequence{n}{(x,y)} = \coordSquareSequence{n}{(y,x)}$
                by \cref{setext,coord_square_sequence,fst_eq,snd_eq,reals_power,tuple_power,funs_type_apply,real_add_inverse,real_add_closed,real_sub,real_abs_sub_swap,real_abs_square}.
            $d((x,y)) = d((y,x))$ by \cref{sum_finite_sequence_unique,function_apply_intro}.
    \end{subproof}
    $d$ is symmetric on $X$ by \cref{metric_symmetric}.
    Show that for all $x,y,z$ such that $x\in X$ and $y\in X$ and $z\in X$ we have
        $d((x,y)) + d((y,z)) \geq d((x,z))$.
    \begin{subproof}
            Fix $x$.
            Fix $y$.
            Fix $z$ such that $x\in X$ and $y\in X$ and $z\in X$.
            Take $r_1\in\reals$ such that $((x,y),r_1)\in d$ by \cref{euclidean_metric_value_exists}.
            Take $r_2\in\reals$ such that $((y,z),r_2)\in d$ by \cref{euclidean_metric_value_exists}.
            Take $r_3\in\reals$ such that $((x,z),r_3)\in d$ by \cref{euclidean_metric_value_exists}.
            Take $s_1\in\reals$ such that
                $\sumFiniteSequence{\coordSquareSequence{n}{(x,y)}}{s_1}$ and $r_1 = \sqrt{s_1}$
                by \cref{euclidean_metric,pair_eq_iff}.
            Take $s_2\in\reals$ such that
                $\sumFiniteSequence{\coordSquareSequence{n}{(y,z)}}{s_2}$ and $r_2 = \sqrt{s_2}$
                by \cref{euclidean_metric,pair_eq_iff}.
            Take $s_3\in\reals$ such that
                $\sumFiniteSequence{\coordSquareSequence{n}{(x,z)}}{s_3}$ and $r_3 = \sqrt{s_3}$
                by \cref{euclidean_metric,pair_eq_iff}.
            Let $a = \coordDiffSequence{n}{x}{y}$.
            Let $b = \coordDiffSequence{n}{y}{z}$.
            Let $c = \coordSumSequence{n}{a}{b}$.
            $\sqrt{s_3} \leq \sqrt{s_1} + \sqrt{s_2}$
                by \cref{coord_diff_sequence_function,coord_sum_sequence_function,coord_square_sequence_diff,coord_diff_sequence_sum,minkowski_finite_sequences}.
            $d((x,z)) \leq d((x,y)) + d((y,z))$ by \cref{function_apply_intro}.
            Hence $d((x,y)) + d((y,z)) \geq d((x,z))$.
    \end{subproof}
    $d$ satisfies the triangle inequality on $X$ by \cref{metric_triangle_inequality}.
    Hence $d$ is a metric on $X$ by \cref{metric_on}.
\end{proof}
% from §20 Item 26: euclidean metric topology equals product topology
\begin{theorem}\label{euclidean_metric_product_topology}
    Suppose $n\in\naturalsPlus$.
    Let $J = \initialSegment{n}$.
    Let $X = \{(i,\reals) \mid i\in J\}$.
    Let $T = \{(i,\standardTopologyR) \mid i\in J\}$.
    Let $T_0 = \productTopologyIndexed{T}{X}$.
    Then $\metricTopology{\euclideanMetric{n}}{\realsPower{J}} = T_0$.
\end{theorem}
\begin{proof}
    Let $X_0 = \realsPower{J}$.
    Let $d = \euclideanMetric{n}$.
    Let $d_1 = \squareMetric{n}$.
    Let $T_E = \metricTopology{d}{X_0}$.
    Let $T_S = \metricTopology{d_1}{X_0}$.
    We have $J$ is inhabited by \cref{real_one_in_naturals,real_one_leq_naturalsplus,initial_segment_naturalsplus}.
    For all $w$ such that $w\in X$ there exist $x,y$ such that $w = (x,y)$.
    For all $i,y_1,y_2$ such that $(i,y_1)\in X$ and $(i,y_2)\in X$ we have $y_1 = y_2$.
    $X$ is a function by \cref{relation,rightunique}.
    For all $w$ such that $w\in T$ there exist $x,y$ such that $w = (x,y)$.
    For all $i,y_1,y_2$ such that $(i,y_1)\in T$ and $(i,y_2)\in T$ we have $y_1 = y_2$.
    $T$ is a function by \cref{relation,rightunique}.
    We have $J\subseteq\dom{X}$ by \cref{dom_intro,subseteq}.
    For all $i$ such that $i\in\dom{X}$ we have $i\in J$.
    $\dom{X}\subseteq J$ by \cref{subseteq}.
    Hence $\dom{X} = J$ by \cref{subseteq_antisymmetric}.
    We have $J\subseteq\dom{T}$ by \cref{dom_intro,subseteq}.
    For all $i$ such that $i\in\dom{T}$ we have $i\in J$.
    $\dom{T}\subseteq J$ by \cref{subseteq}.
    Thus $\dom{T} = J$ by \cref{subseteq_antisymmetric}.
    For all $\alpha\in\dom{X}$ we have $\apply{T}{\alpha}$ is a topology on $\apply{X}{\alpha}$
        by \cref{subseteq,function_apply_intro,standard_basis_is_basis,basis_topology_is_topology,standard_topology_reals}.
    Let $S_0 = \productSubbasis{T}{X}$.
    Let $B_0 = \finiteIntersections{S_0}$.
    $B_0$ is a basis for $T_0$ on $\productFamily{X}$ by \cref{product_subbasis_finite_intersections_basis}.
    $B_0$ is a basis for $T_0$ on $X_0$ by \cref{product_family_reals_power}.
    $d$ is a metric on $X_0$ and $d_1$ is a metric on $X_0$ by \cref{euclidean_metric_is_metric,square_metric_is_metric}.
    Let $B_S = \metricBasis{d_1}{X_0}$.
    $B_S$ is a basis for $T_S$ on $X_0$
        by \cref{metric_basis_is_basis,basis_for_topology,metric_topology}.
    Show that for all $x\in X_0$ we have
        for all $\epsilon\in\reals$ such that $0 < \epsilon$ there exists $\delta\in\reals$
            such that $0 < \delta$ and
            $\metricBall{d}{X_0}{x}{\delta} \subseteq \metricBall{d_1}{X_0}{x}{\epsilon}$.
    \begin{subproof}
                Fix $x$ such that $x\in X_0$.
                Fix $\epsilon$ such that $\epsilon\in\reals$ and $0 < \epsilon$.
                Let $\delta = \epsilon$.
                $\delta\in\reals$ and $0 < \delta$.
                For all $y$ such that $y\in\metricBall{d}{X_0}{x}{\delta}$ we have $y\in\metricBall{d_1}{X_0}{x}{\epsilon}$
                    by \cref{metric_ball,euclidean_metric_value_exists,square_metric_value_exists,euclidean_metric_function,square_metric_function,function_apply_intro,real_lt_subst_left,real_lt_subst_right,euclidean_square_metric_bounds,real_leq_split,real_lt_trans}.
                $\metricBall{d}{X_0}{x}{\delta} \subseteq \metricBall{d_1}{X_0}{x}{\epsilon}$ by \cref{subseteq}.
    \end{subproof}
    $T_E$ is finer than $T_S$ on $X_0$ by \cref{metric_topology_finer_iff}.
    Show that for all $x\in X_0$ we have
        for all $\epsilon\in\reals$ such that $0 < \epsilon$ there exists $\delta\in\reals$
            such that $0 < \delta$ and
            $\metricBall{d_1}{X_0}{x}{\delta} \subseteq \metricBall{d}{X_0}{x}{\epsilon}$.
    \begin{subproof}
                Fix $x$ such that $x\in X_0$.
                Fix $\epsilon$ such that $\epsilon\in\reals$ and $0 < \epsilon$.
                $0 < 1$ and $0\in\reals$ and $1\in\reals$ and $n\in\reals$
                    by \cref{real_zero_lt_one,real_add_ident,real_naturals_subset_reals,real_one_in_naturals,subseteq}.
                $0 < n$ by \cref{real_one_leq_naturalsplus,real_lt_subst_right,real_lt_trans}.
                We have $\sqrt{n}\in\reals$ and $\sqrt{n} \geq 0$ and $\sqrt{n} \rmul \sqrt{n} = n$ by \cref{positive_square_root}.
                Hence $\sqrt{n}$ is positive.
                Let $\delta = \epsilon \rmul \inv{\sqrt{n}}$.
                $\sqrt{n} \neq 0$ by \cref{real_pos_nonzero}.
                $\inv{\sqrt{n}}\in\reals$ and $\inv{\sqrt{n}} > 0$ by \cref{real_mul_inverse,real_inv_pos}.
                $\epsilon \rmul \inv{\sqrt{n}}\in\reals$ by \cref{real_mul_closed}.
                $\delta > 0$ by \cref{real_mul_pos_pos_gt}.
                Show that for all $y$ such that $y\in\metricBall{d_1}{X_0}{x}{\delta}$ we have $y\in\metricBall{d}{X_0}{x}{\epsilon}$.
                \begin{subproof}
                    Fix $y$ such that $y\in\metricBall{d_1}{X_0}{x}{\delta}$.
                    $y\in X_0$ and $d_1((x,y)) < \delta$ by \cref{metric_ball}.
                    Take $m\in\reals$ such that $((x,y),m)\in d_1$ by \cref{square_metric_value_exists}.
                    Take $r\in\reals$ such that $((x,y),r)\in d$ by \cref{euclidean_metric_value_exists}.
                    $m < \delta$ by \cref{square_metric_function,function_apply_intro,real_lt_subst_left}.
                    $r \leq \sqrt{n} \rmul m$ by \cref{euclidean_square_metric_bounds}.
                    Hence $\sqrt{n} \rmul m < \sqrt{n} \rmul \delta$ by \cref{real_lt_mul_pos,real_mul_comm}.
                    $\sqrt{n} \rmul m < \epsilon$ by
                        \cref{real_lt_subst_right,real_mul_assoc,real_mul_comm,real_mul_inverse,real_mul_ident}.
                    Hence $d((x,y)) < \epsilon$ by
                        \cref{euclidean_metric_function,real_mul_closed,real_leq_split,real_lt_subst_left,real_lt_trans,function_apply_intro}.
                    $y\in\metricBall{d}{X_0}{x}{\epsilon}$ by \cref{metric_ball}.
                \end{subproof}
                $\metricBall{d_1}{X_0}{x}{\delta} \subseteq \metricBall{d}{X_0}{x}{\epsilon}$ by \cref{subseteq}.
    \end{subproof}
    $T_S$ is finer than $T_E$ on $X_0$ by \cref{metric_topology_finer_iff}.
    Hence $T_S\subseteq T_E$ and $T_E\subseteq T_S$ by \cref{finer_topology}.
    $T_E = T_S$ by \cref{subseteq_antisymmetric}.
    Show that for all $x,B_2$ such that $x\in X_0$ and $B_2\in B_S$ and $x\in B_2$
        there exists $B_3\in B_0$ such that $x\in B_3$ and $B_3\subseteq B_2$.
    \begin{subproof}
                Fix $x$.
                Fix $B_2$ such that $x\in X_0$ and $B_2\in B_S$ and $x\in B_2$.
                Take $x_0,\epsilon$ such that $x_0\in X_0$ and $\epsilon\in\reals$ and $0 < \epsilon$ and
                    $B_2 = \metricBall{d_1}{X_0}{x_0}{\epsilon}$ by \cref{metric_basis}.
                Let $R_g = \{w\in J\times\pow{\productFamily{X}} \mid \exists i\in J.
                    w = (i,\preimg{\piIndex{X}{i}}{\intervalopen{x_0(i) - \epsilon}{x_0(i) + \epsilon}})\}$.
                For all $w$ such that $w\in R_g$ there exist $i,A$ such that $w = (i,A)$.
                $R_g$ is a relation by \cref{relation}.
                For all $i\in J$ we have
                    $\preimg{\piIndex{X}{i}}{\intervalopen{x_0(i) - \epsilon}{x_0(i) + \epsilon}}\in\pow{\productFamily{X}}$
                    by \cref{preimg_iff,projection_map_indexed,pair_eq_iff,subseteq,powerset_intro}.
                For all $i\in J$ there exists $A$ such that $i\mathrel{R_g} A$
                    by \cref{times_tuple_intro}.
                Take $g$ such that $g$ is a function on $J$ and for all $i\in J$ we have $i\mathrel{R_g}\apply{g}{i}$
                    by \cref{choice_function}.
                Let $F = \img{g}{J}$.
                $F$ is finite by \cref{initial_segment_finite,finite_image_function}.
                Take $i$ such that $i\in J$.
                $\apply{g}{i}\in F$ by \cref{function_apply_elim,img_iff}.
                Hence $F\neq\emptyset$ by \cref{function_apply_elim,img_iff,emptyset}.
                Show that for all $A$ such that $A\in F$ we have $A\in S_0$.
                \begin{subproof}
                    Fix $A$ such that $A\in F$.
                    Take $j\in J$ such that $(j,A)\in g$ by \cref{img_iff}.
                    $j\mathrel{R_g}\apply{g}{j}$.
                    Take $i_0\in J$ such that
                        $(j,\apply{g}{j}) = (i_0,\preimg{\piIndex{X}{i_0}}{\intervalopen{x_0(i_0) - \epsilon}{x_0(i_0) + \epsilon}})$.
                    $A = \preimg{\piIndex{X}{j}}{\intervalopen{x_0(j) - \epsilon}{x_0(j) + \epsilon}}$
                        by \cref{function_apply_intro,pair_eq_iff}.
                    We have $x_0(j)\in\reals$ by \cref{reals_power,tuple_power,funs_type_apply}.
                    We have $x_0(j) - \epsilon\in\reals$ and $x_0(j) + \epsilon\in\reals$
                        by \cref{real_add_inverse,real_sub,real_add_closed}.
                    $0 < \epsilon$ and $0\in\reals$ by \cref{real_zero_lt_one,real_add_ident}.
                    $x_0(j) - \epsilon < x_0(j)$ by \cref{real_sub_pos_lt}.
                    $x_0(j) < x_0(j) + \epsilon$ by \cref{real_lt_add,real_add_comm,real_add_ident}.
                    Hence $x_0(j) - \epsilon < x_0(j) + \epsilon$ by \cref{real_lt_trans}.
                    $\intervalopen{x_0(j) - \epsilon}{x_0(j) + \epsilon}\in\standardTopologyR$
                        by \cref{intervalopen_subset_reals,powerset_intro,standard_basis_reals,basis_elem_open_in_generated,standard_basis_is_basis,standard_topology_reals}.
                    Let $U = \intervalopen{x_0(j) - \epsilon}{x_0(j) + \epsilon}$.
                    $U\in\apply{T}{j}$ by \cref{function_apply_intro}.
                    $A\in\productSubbasisAt{T}{X}{j}$
                        by \cref{preimg_iff,projection_map_indexed,pair_eq_iff,subseteq,powerset_intro,product_subbasis_indexed}.
                    Hence $A\in S_0$ by \cref{product_subbasis_union_indexed}.
                \end{subproof}
                $F\subseteq S_0$ by \cref{subseteq}.
                Show that for all $y$ such that $y\in B_2$ we have $y\in\inters{F}$.
                \begin{subproof}
                            Fix $y$ such that $y\in B_2$.
                            We have $y\in X_0$ and $d_1((x_0,y)) < \epsilon$ by \cref{metric_ball}.
                            Take $m\in\reals$ such that $((x_0,y),m)\in d_1$ and
                                for all $j\in J$ we have $\abs{x_0(j) - y(j)} \leq m$ by \cref{square_metric_value_exists}.
                            $m < \epsilon$ by \cref{square_metric_function,function_apply_intro,real_lt_subst_left}.
                            Show that for all $A$ such that $A\in F$ we have $y\in A$.
                            \begin{subproof}
                                Fix $A$ such that $A\in F$.
                                Take $j\in J$ such that $(j,A)\in g$ by \cref{img_iff}.
                                $j\mathrel{R_g}\apply{g}{j}$.
                                Take $i_0\in J$ such that
                                    $(j,\apply{g}{j}) = (i_0,\preimg{\piIndex{X}{i_0}}{\intervalopen{x_0(i_0) - \epsilon}{x_0(i_0) + \epsilon}})$.
                                $A = \preimg{\piIndex{X}{j}}{\intervalopen{x_0(j) - \epsilon}{x_0(j) + \epsilon}}$
                                    by \cref{function_apply_intro,pair_eq_iff}.
                                We have $(y,\apply{y}{j})\in\piIndex{X}{j}$ by
                                    \cref{product_family_reals_power,subseteq,projection_map_indexed}.
                                We have $x_0(j)\in\reals$ and $\apply{y}{j}\in\reals$
                                    by \cref{reals_power,tuple_power,funs_type_apply}.
                                $x_0(j) - \apply{y}{j}\in\reals$
                                    by \cref{real_add_inverse,real_sub,real_add_closed}.
                                We have $\abs{x_0(j) - \apply{y}{j}}\in\reals$ by \cref{real_abs_in_reals}.
                                We have $\abs{x_0(j) - \apply{y}{j}} \leq m$.
                                $\abs{x_0(j) - \apply{y}{j}} < \epsilon$ by \cref{real_leq_split,real_lt_subst_left,real_lt_trans}.
                                Hence $x_0(j) - \epsilon < \apply{y}{j}$ and $\apply{y}{j} < x_0(j) + \epsilon$
                                    by \cref{real_interval_from_abs_lt}.
                                $\apply{y}{j}\in\intervalopen{x_0(j) - \epsilon}{x_0(j) + \epsilon}$ by \cref{intervalopen}.
                                Thus $y\in A$ by \cref{preimg_iff}.
                            \end{subproof}
                            $y\in\inters{F}$ by \cref{inters_iff_forall}.
                \end{subproof}
                $B_2\subseteq\inters{F}$ by \cref{subseteq}.
                Show that for all $y$ such that $y\in\inters{F}$ we have $y\in B_2$.
                \begin{subproof}
                            Fix $y$ such that $y\in\inters{F}$.
                            Take $j$ such that $j\in J$.
                            $\apply{g}{j}\in F$ by \cref{function_apply_elim,img_iff}.
                            $y\in\apply{g}{j}$ by \cref{inters_iff_forall}.
                            Hence $j\mathrel{R_g}\apply{g}{j}$.
                            Take $i_0\in J$ such that
                                $(j,\apply{g}{j}) = (i_0,\preimg{\piIndex{X}{i_0}}{\intervalopen{x_0(i_0) - \epsilon}{x_0(i_0) + \epsilon}})$.
                            We have $\apply{g}{j} = \preimg{\piIndex{X}{j}}{\intervalopen{x_0(j) - \epsilon}{x_0(j) + \epsilon}}$ by \cref{pair_eq_iff}.
                            $y\in\preimg{\piIndex{X}{j}}{\intervalopen{x_0(j) - \epsilon}{x_0(j) + \epsilon}}$.
                            Take $z\in\intervalopen{x_0(j) - \epsilon}{x_0(j) + \epsilon}$ such that
                                $y\mathrel{\piIndex{X}{j}} z$ by \cref{preimg_iff}.
                            We have $(y,z)\in\piIndex{X}{j}$.
                            Take $y_0\in\productFamily{X}$ such that $(y,z) = (y_0,\apply{y_0}{j})$ by \cref{projection_map_indexed}.
                            $y = y_0$ by \cref{pair_eq_iff}.
                            Thus $y\in X_0$ by \cref{product_family_reals_power,subseteq}.
                            Take $m\in\reals$ such that $((x_0,y),m)\in d_1$ and
                                for all $j\in J$ we have $\abs{x_0(j) - y(j)} \leq m$ by \cref{square_metric_value_exists}.
                            Show that for all $t\in\coordDiffSet{n}{x_0}{y}$ we have $t < \epsilon$.
                            \begin{subproof}
                                Fix $t$ such that $t\in\coordDiffSet{n}{x_0}{y}$.
                                Take $k\in J$ such that $t = \abs{x_0(k) - y(k)}$ by \cref{coord_diff_set}.
                                We have $y\in\apply{g}{k}$ by \cref{function_apply_elim,img_iff,inters_iff_forall}.
                                $y\in\preimg{\piIndex{X}{k}}{\intervalopen{x_0(k) - \epsilon}{x_0(k) + \epsilon}}$
                                    by \cref{function_apply_intro,pair_eq_iff}.
                                Take $u\in\intervalopen{x_0(k) - \epsilon}{x_0(k) + \epsilon}$ such that
                                    $y\mathrel{\piIndex{X}{k}} u$ by \cref{preimg_iff}.
                                We have $(y,u)\in\piIndex{X}{k}$.
                                Take $y_1\in\productFamily{X}$ such that $(y,u) = (y_1,\apply{y_1}{k})$ by \cref{projection_map_indexed}.
                                $y = y_1$ and $u = \apply{y}{k}$ by \cref{pair_eq_iff}.
                                Hence $x_0(k) - \epsilon < y(k)$ and $y(k) < x_0(k) + \epsilon$ by \cref{intervalopen}.
                                $\abs{x_0(k) - y(k)} < \epsilon$
                                    by \cref{reals_power,tuple_power,funs_type_apply,real_abs_lt_from_interval}.
                                Thus $t < \epsilon$.
                            \end{subproof}
                            $d_1((x_0,y)) < \epsilon$ by
                                \cref{square_metric,pair_eq_iff,square_metric_function,function_apply_intro,real_lt_subst_left}.
                            Hence $y\in B_2$ by \cref{metric_ball}.
                \end{subproof}
                $\inters{F}\subseteq B_2$ by \cref{subseteq}.
                Hence $B_2 = \inters{F}$ by \cref{subseteq_antisymmetric}.
                $B_2\in B_0$
                    by \cref{inters_iff_forall,unions_iff,subseteq,powerset_intro,finite_intersections}.
                Thus $x\in B_2$ and $B_2\subseteq B_2$.
    \end{subproof}
    $T_0$ is finer than $T_S$ on $X_0$ by \cref{basis_finer_criterion}.
    Show that for all $x,B_2$ such that $x\in X_0$ and $B_2\in B_0$ and $x\in B_2$
        there exists $B_3\in B_S$ such that $x\in B_3$ and $B_3\subseteq B_2$.
    \begin{subproof}
                Fix $x$.
                Fix $B_2$ such that $x\in X_0$ and $B_2\in B_0$ and $x\in B_2$.
                Take $F\subseteq S_0$ such that $F$ is finite and inhabited and $B_2 = \inters{F}$ by \cref{finite_intersections}.
                Let $R = \{w\in F\times\reals \mid \exists \beta\in J.\exists U\in\standardTopologyR.
                    \fst{w} = \preimg{\piIndex{X}{\beta}}{U} \land \apply{x}{\beta}\in U \land
                    0 < \snd{w} \land \metricBall{\standardMetricR}{\reals}{\apply{x}{\beta}}{\snd{w}} \subseteq U\}$.
                For all $w$ such that $w\in R$ there exist $y,z$ such that $w = (y,z)$.
                $R$ is a relation by \cref{relation}.
                Show that for all $A\in F$ there exists $\epsilon\in\reals$ such that $A\mathrel{R}\epsilon$.
                \begin{subproof}
                    Fix $A$ such that $A\in F$.
                    $x\in A$ by \cref{inters_iff_forall}.
                    Take $\beta\in\dom{X}$ such that $A\in\productSubbasisAt{T}{X}{\beta}$
                        by \cref{subseteq,product_subbasis_union_indexed}.
                    Take $U\in\apply{T}{\beta}$ such that $A = \preimg{\piIndex{X}{\beta}}{U}$ by \cref{product_subbasis_indexed}.
                    $\beta\in J$ by \cref{subseteq}.
                    $x\in\preimg{\piIndex{X}{\beta}}{U}$.
                    Take $z\in U$ such that $x\mathrel{\piIndex{X}{\beta}} z$ by \cref{preimg_iff}.
                    We have $(x,z)\in\piIndex{X}{\beta}$.
                    Take $x_1\in\productFamily{X}$ such that $(x,z) = (x_1,\apply{x_1}{\beta})$ by \cref{projection_map_indexed}.
                    $z = \apply{x}{\beta}$ by \cref{pair_eq_iff}.
                    $\apply{x}{\beta}\in U$.
                    $U\in\standardTopologyR$ by \cref{function_apply_intro}.
                    $\standardBasisR$ is a basis on $\reals$ by \cref{standard_basis_is_basis}.
                    We have $U\subseteq\reals$
                        by \cref{basis_topology_is_topology,standard_topology_reals,topology_on,subseteq,pow_iff}.
                    $U\in\metricTopology{\standardMetricR}{\reals}$ by \cref{standard_metric_topology}.
                    $\standardMetricR$ is a metric on $\reals$ by \cref{standard_metric_is_metric}.
                    $U$ is open in $\reals$ with respect to $\metricTopology{\standardMetricR}{\reals}$ by \cref{open_in,basis_topology_is_topology,metric_basis_is_basis,metric_topology}.
                    Take $\epsilon\in\reals$ such that $0 < \epsilon$ and
                        $\metricBall{\standardMetricR}{\reals}{\apply{x}{\beta}}{\epsilon}\subseteq U$
                        by \cref{metric_open_iff}.
                    Hence $A\mathrel{R}\epsilon$ by
                        \cref{times_tuple_intro,fst_eq,snd_eq,pair_eq_iff}.
                \end{subproof}
                Take $e$ such that $e$ is a function on $F$ and for all $A\in F$ we have $A\mathrel{R}\apply{e}{A}$
                    by \cref{choice_function}.
                $e$ is a function and $\dom{e} = F$.
                Let $E = \img{e}{F}$.
                $E$ is finite by \cref{finite_image_function}.
                Take $A_0$ such that $A_0\in F$.
                $\apply{e}{A_0}\in E$ by \cref{function_apply_elim,img_iff}.
                Hence $E\neq\emptyset$ by \cref{function_apply_elim,img_iff,emptyset}.
                Show that for all $t$ such that $t\in E$ we have $t\in\reals$ and $0 < t$.
                \begin{subproof}
                    Fix $t$ such that $t\in E$.
                    Take $A_1$ such that $A_1\in F$ and $(A_1,t)\in e$ by \cref{img_iff}.
                    $t = \apply{e}{A_1}$ by \cref{subseteq,function_apply_iff}.
                    Hence $A_1\mathrel{R}t$ by \cref{subseteq}.
                    $t\in\reals$ and $0 < t$.
                \end{subproof}
                $E\subseteq\reals$ by \cref{subseteq}.
                Take $\delta\in E$ such that $\delta$ is a minimum of $E$ by \cref{finite_real_minimum}.
                $\delta\in\reals$ and $0 < \delta$ by \cref{subseteq}.
                Let $B_3 = \metricBall{d_1}{X_0}{x}{\delta}$.
                $B_3\in B_S$ by \cref{metric_ball,subseteq,powerset_intro,metric_basis}.
                Show that for all $y$ such that $y\in B_3$ we have $y\in B_2$.
                \begin{subproof}
                        Fix $y$ such that $y\in B_3$.
                        We have $y\in X_0$ and $d_1((x,y)) < \delta$ by \cref{metric_ball}.
                        Take $m\in\reals$ such that $((x,y),m)\in d_1$ and
                            for all $j\in J$ we have $\abs{x(j) - y(j)} \leq m$ by \cref{square_metric_value_exists}.
                        $m < \delta$ by \cref{square_metric_function,function_apply_intro,real_lt_subst_left}.
                        Show that for all $A\in F$ we have $y\in A$.
                        \begin{subproof}
                            Fix $A$ such that $A\in F$.
                            $(A,\apply{e}{A})\in R$.
                            Let $w = (A,\apply{e}{A})$.
                            We have $w\in R$.
                            Take $\beta,U$ such that $\beta\in J$ and $U\in\standardTopologyR$ and
                                $\fst{w} = \preimg{\piIndex{X}{\beta}}{U}$ and $\apply{x}{\beta}\in U$ and
                                $0 < \snd{w}$ and
                                $\metricBall{\standardMetricR}{\reals}{\apply{x}{\beta}}{\snd{w}} \subseteq U$.
                            $\fst{w} = A$ and $\snd{w} = \apply{e}{A}$ by \cref{fst_eq,snd_eq}.
                            $A = \preimg{\piIndex{X}{\beta}}{U}$.
                            $0 < \apply{e}{A}$.
                            $\metricBall{\standardMetricR}{\reals}{\apply{x}{\beta}}{\apply{e}{A}} \subseteq U$.
                            $\apply{e}{A}\in E$ by \cref{function_apply_elim,img_iff}.
                            $\apply{e}{A}\in\reals$ by \cref{subseteq}.
                            Hence $\delta \leq \apply{e}{A}$ by \cref{function_apply_elim,img_iff,real_minimum}.
                            $m < \apply{e}{A}$ by \cref{real_leq_split,real_lt_trans,real_lt_subst_right}.
                            We have $x(\beta)\in\reals$ and $\apply{y}{\beta}\in\reals$
                                by \cref{reals_power,tuple_power,funs_type_apply}.
                            We have $\abs{x(\beta) - \apply{y}{\beta}}\in\reals$
                                by \cref{reals_power,tuple_power,funs_type_apply,real_sub,real_add_inverse,real_add_closed,real_abs_in_reals}.
                            $\abs{x(\beta) - \apply{y}{\beta}} \leq m$.
                            Hence $\abs{x(\beta) - \apply{y}{\beta}} < \apply{e}{A}$ by \cref{real_lt_subst_left,real_lt_trans}.
                            $\apply{\standardMetricR}{(\apply{x}{\beta},\apply{y}{\beta})} < \apply{e}{A}$
                                by \cref{standard_metric_apply,real_lt_subst_left}.
                            $\apply{y}{\beta}\in\metricBall{\standardMetricR}{\reals}{\apply{x}{\beta}}{\apply{e}{A}}$ by \cref{metric_ball}.
                            $\apply{y}{\beta}\in U$ by \cref{subseteq}.
                            $y\in\preimg{\piIndex{X}{\beta}}{U}$ by
                                \cref{metric_ball,product_family_reals_power,subseteq,projection_map_indexed,preimg_iff}.
                            Thus $y\in A$.
                        \end{subproof}
                        $y\in\inters{F}$ by \cref{inters_iff_forall}.
                        Hence $y\in B_2$ by \cref{subseteq_antisymmetric}.
                \end{subproof}
                $B_3\subseteq B_2$ by \cref{subseteq}.
                Thus $x\in B_3$ and $B_3\subseteq B_2$ by \cref{metric_ball,metric_zero_characterization,metric_on}.
    \end{subproof}
    $T_S$ is finer than $T_0$ on $X_0$ by \cref{basis_finer_criterion}.
    Hence $T_0\subseteq T_S$ and $T_S\subseteq T_0$ by \cref{finer_topology}.
    $T_S = T_0$ by \cref{subseteq_antisymmetric}.
    Hence $T_E = T_0$.
\end{proof}

% from §20 helper lemma 10: uniform coordinate metric set
\begin{definition}\label{uniform_metric_set}
    $\uniformMetricSet{J}{x}{y} = \{\apply{\boundedMetric{\standardMetricR}}{(x(\alpha),y(\alpha))} \mid \alpha\in J\}$.
\end{definition}

% from §20 helper lemma 11: supremum predicate
\begin{definition}\label{real_supremum_pred}
    $\realSupremum{p}{A}$ iff $p$ is a supremum of $A$.
\end{definition}

% from §20 Item 27: uniform metric
\begin{definition}\label{uniform_metric}
    $\uniformMetric{J} = \{w \mid \exists p\in\realsPower{J}\times\realsPower{J}.
        \exists x\in\realsPower{J}. \exists y\in\realsPower{J}. \exists r\in\reals.
        w = (p,r) \land p = (x,y) \land
        \realSupremum{r}{\uniformMetricSet{J}{x}{y}}\}$.
\end{definition}

% from §20 Item 28: uniform topology
\begin{definition}\label{uniform_topology}
    $\uniformTopology{J} = \metricTopology{\uniformMetric{J}}{\realsPower{J}}$.
\end{definition}

% from §20 helper lemma 11a: uniform metric is a metric on inhabited products
\begin{lemma}\label{uniform_metric_is_metric}
    Suppose $J$ is inhabited.
    Let $X = \realsPower{J}$.
    Then $\uniformMetric{J}$ is a metric on $X$.
\end{lemma}
\begin{proof}
    Let $d = \uniformMetric{J}$.
    Let $e = \boundedMetric{\standardMetricR}$.
    $\standardMetricR$ is a metric on $\reals$ by \cref{standard_metric_is_metric}.
    $e$ is a metric on $\reals$ and $e\in\funs{\reals\times\reals}{\reals}$
        by \cref{bounded_metric_is_metric,metric_on}.
    Show that $d\in\funs{X\times X}{\reals}$.
    \begin{subproof}
        Show that for all $p,z_1,z_2$ such that $(p,z_1)\in d$ and $(p,z_2)\in d$ we have $z_1 = z_2$.
        \begin{subproof}
            Fix $p$.
            Fix $z_1$.
            Fix $z_2$ such that $(p,z_1)\in d$ and $(p,z_2)\in d$.
            Take $p_1,x_1,y_1,r_1$ such that
                $p_1\in X\times X$ and $x_1\in X$ and $y_1\in X$ and $r_1\in\reals$ and
                $(p,z_1) = (p_1,r_1)$ and $p_1 = (x_1,y_1)$ and
                $\realSupremum{r_1}{\uniformMetricSet{J}{x_1}{y_1}}$ by \cref{uniform_metric}.
            Take $p_2,x_2,y_2,r_2$ such that
                $p_2\in X\times X$ and $x_2\in X$ and $y_2\in X$ and $r_2\in\reals$ and
                $(p,z_2) = (p_2,r_2)$ and $p_2 = (x_2,y_2)$ and
                $\realSupremum{r_2}{\uniformMetricSet{J}{x_2}{y_2}}$ by \cref{uniform_metric}.
            $p = p_1$ and $z_1 = r_1$ and $p = p_2$ and $z_2 = r_2$ by \cref{pair_eq_iff}.
            $p_1 = p_2$.
            $(x_1,y_1) = (x_2,y_2)$ by \cref{pair_eq_iff}.
            We have $x_1 = x_2$ and $y_1 = y_2$ by \cref{pair_eq_iff}.
            $\uniformMetricSet{J}{x_1}{y_1} = \uniformMetricSet{J}{x_2}{y_2}$.
            Hence $r_1 = r_2$ by \cref{real_supremum_pred,real_supremum_unique}.
            $z_1 = z_2$.
        \end{subproof}
        $d$ is a function by \cref{uniform_metric,relation,rightunique}.
        Show that for all $p$ such that $p\in X\times X$ we have $p\in\dom{d}$.
        \begin{subproof}
            Fix $p$ such that $p\in X\times X$.
            Take $x,y$ such that $p = (x,y)$ and $x\in X$ and $y\in X$ by \cref{times_elem_is_tuple}.
            Take $\beta$ such that $\beta\in J$.
            Let $A = \uniformMetricSet{J}{x}{y}$.
            For all $a$ such that $a\in A$ we have $a\in\reals$
                by \cref{uniform_metric_set,reals_power,tuple_power,funs_type_apply,times_tuple_intro,metric_on}.
            $A\subseteq\reals$ by \cref{subseteq}.
            $A\neq\emptyset$ by \cref{uniform_metric_set,emptyset}.
            For all $a\in A$ we have $a \leq 1$ by
                \cref{uniform_metric_set,reals_power,tuple_power,funs_type_apply,times_tuple_intro,bounded_metric_leq_one,real_lt_subst_left,real_lt_subst_right}.
            $1\in\reals$ by \cref{real_naturals_subset_reals,real_one_in_naturals,subseteq}.
            $1$ is an upper bound of $A$ by \cref{real_upper_bound}.
            $A$ is bounded above by \cref{real_bounded_above}.
            Take $r$ such that $r$ is a supremum of $A$ by \cref{real_completeness_axiom}.
            We have $(x,y)\in X\times X$ by \cref{times_tuple_intro}.
            $r\in\reals$ by \cref{real_supremum,real_upper_bound}.
            $((x,y),r)\in d$ by \cref{real_supremum_pred,uniform_metric}.
            Hence $p\in\dom{d}$ by \cref{dom_intro,pair_eq_iff}.
        \end{subproof}
        $X\times X\subseteq\dom{d}$ by \cref{subseteq}.
        For all $p$ such that $p\in\dom{d}$ we have $p\in X\times X$
            by \cref{dom_iff,uniform_metric,pair_eq_iff}.
        $\dom{d}\subseteq X\times X$ by \cref{subseteq}.
        Thus $\dom{d} = X\times X$ by \cref{subseteq_antisymmetric}.
        For all $p\in\dom{d}$ we have $d(p)\in\reals$
            by \cref{dom_iff,uniform_metric,pair_eq_iff,function_apply_iff,real_lt_subst_left}.
        $d\in\funs{X\times X}{\reals}$ by \cref{funs_intro}.
    \end{subproof}
    Show that for all $x,y$ such that $x\in X$ and $y\in X$ we have $d((x,y)) \geq 0$.
    \begin{subproof}
        Fix $x$.
        Fix $y$ such that $x\in X$ and $y\in X$.
        Take $\beta$ such that $\beta\in J$.
        Let $r = d((x,y))$.
        $(x,y)\in\dom{d}$ by \cref{times_tuple_intro,funs_elim}.
        $((x,y),r)\in d$ by \cref{funs_elim,function_apply_elim}.
        Take $p_0,x_0,y_0,r_0$ such that
            $p_0\in X\times X$ and $x_0\in X$ and $y_0\in X$ and $r_0\in\reals$ and
            $((x,y),r) = (p_0,r_0)$ and $p_0 = (x_0,y_0)$ and
            $\realSupremum{r_0}{\uniformMetricSet{J}{x_0}{y_0}}$ by \cref{uniform_metric}.
        $p_0 = (x,y)$ and $r = r_0$ by \cref{pair_eq_iff}.
        $r\in\reals$.
        We have $x_0 = x$ and $y_0 = y$ by \cref{pair_eq_iff}.
        Let $a = \apply{e}{(x(\beta),y(\beta))}$.
        $a \leq r$ by \cref{uniform_metric_set,real_supremum_pred,real_supremum,real_upper_bound}.
        $x(\beta)\in\reals$ and $y(\beta)\in\reals$ by \cref{reals_power,tuple_power,funs_type_apply}.
        $(x(\beta),y(\beta))\in\reals\times\reals$ by \cref{times_tuple_intro}.
        We have $a\in\reals$ and $0 \leq a$ by \cref{funs_type_apply,metric_on,metric_nonnegative}.
        $0 \leq r$ by \cref{real_add_ident,real_leq_trans}.
        Hence $d((x,y)) \geq 0$.
    \end{subproof}
    $d$ is nonnegative on $X$ by \cref{metric_nonnegative}.
    Show that for all $x,y$ such that $x\in X$ and $y\in X$ we have if $d((x,y)) = 0$ then $x = y$.
    \begin{subproof}
        Fix $x$.
        Fix $y$ such that $x\in X$ and $y\in X$.
        Assume $d((x,y)) = 0$.
        Let $r = d((x,y))$.
        $(x,y)\in\dom{d}$ by \cref{times_tuple_intro,funs_elim}.
        $((x,y),r)\in d$ by \cref{funs_elim,function_apply_elim}.
        Take $p_0,x_0,y_0,r_0$ such that
            $p_0\in X\times X$ and $x_0\in X$ and $y_0\in X$ and $r_0\in\reals$ and
            $((x,y),r) = (p_0,r_0)$ and $p_0 = (x_0,y_0)$ and
            $\realSupremum{r_0}{\uniformMetricSet{J}{x_0}{y_0}}$ by \cref{uniform_metric}.
        $p_0 = (x,y)$ and $r = r_0$ by \cref{pair_eq_iff}.
        We have $x_0 = x$ and $y_0 = y$ by \cref{pair_eq_iff}.
        Show that for all $\alpha\in J$ we have $x(\alpha) = y(\alpha)$.
        \begin{subproof}
            Fix $\alpha$ such that $\alpha\in J$.
            Let $a = \apply{e}{(x(\alpha),y(\alpha))}$.
            $a \leq r$ by \cref{uniform_metric_set,real_supremum_pred,real_supremum,real_upper_bound}.
            $x(\alpha)\in\reals$ and $y(\alpha)\in\reals$ by \cref{reals_power,tuple_power,funs_type_apply}.
            $(x(\alpha),y(\alpha))\in\reals\times\reals$ by \cref{times_tuple_intro}.
            We have $a\in\reals$ and $0 \leq a$ by \cref{funs_type_apply,metric_on,metric_nonnegative}.
            $a \leq 0$ by \cref{real_supremum_pred,real_supremum,real_upper_bound,real_lt_subst_left,real_lt_subst_right}.
            Hence $a = 0$ by \cref{real_add_ident,real_leq_antisym,real_lt_subst_right}.
            $x(\alpha) = y(\alpha)$ by \cref{metric_on,metric_zero_characterization}.
        \end{subproof}
        $\dom{x} = J$ and $\dom{y} = J$ by \cref{reals_power,tuple_power,funs_elim}.
        $\dom{x}\subseteq\dom{y}$ and $\dom{y}\subseteq\dom{x}$ by \cref{subseteq_refl}.
        $x\subseteq y$ and $y\subseteq x$ by \cref{reals_power,tuple_power,funs_elim,fun_subseteq,subseteq}.
        Hence $x = y$ by \cref{subseteq_antisymmetric}.
    \end{subproof}
    Show that for all $x,y$ such that $x\in X$ and $y\in X$ we have if $x = y$ then $d((x,y)) = 0$.
    \begin{subproof}
        Fix $x$.
        Fix $y$ such that $x\in X$ and $y\in X$.
        Assume $x = y$.
        Let $A = \uniformMetricSet{J}{x}{y}$.
        For all $a\in A$ we have $a \leq 0$ by \cref{uniform_metric_set,real_lt_subst_left,reals_power,tuple_power,funs_type_apply,times_tuple_intro,metric_on,metric_zero_characterization}.
        $0\in\reals$ by \cref{real_add_ident}.
        For all $a$ such that $a\in A$ we have $a\in\reals$
            by \cref{uniform_metric_set,reals_power,tuple_power,funs_type_apply,times_tuple_intro,metric_on}.
        We have $A\subseteq\reals$ and $0$ is an upper bound of $A$ by \cref{subseteq,real_add_ident,real_upper_bound}.
        Show that for all $b$ such that $b$ is an upper bound of $A$ we have $0 \leq b$.
        \begin{subproof}
            Fix $b$.
            Assume $b$ is an upper bound of $A$.
            Take $\beta$ such that $\beta\in J$.
            Let $a = \apply{e}{(x(\beta),y(\beta))}$.
            $a \leq b$ by \cref{uniform_metric_set,real_upper_bound}.
            $x(\beta)\in\reals$ and $y(\beta)\in\reals$ by \cref{reals_power,tuple_power,funs_type_apply}.
            $(x(\beta),y(\beta))\in\reals\times\reals$ by \cref{times_tuple_intro}.
            $a = 0$ by \cref{metric_on,metric_zero_characterization}.
            Hence $0 \leq b$ by \cref{real_lt_subst_left}.
        \end{subproof}
        $0$ is a supremum of $A$ by \cref{real_supremum}.
        $((x,y),0)\in d$ by \cref{real_supremum_pred,uniform_metric,pair_eq_iff,times_tuple_intro,real_add_ident}.
        $(x,y)\in\dom{d}$ by \cref{dom_intro}.
        Hence $d((x,y)) = 0$ by \cref{funs_elim,function_apply_iff}.
    \end{subproof}
    $d$ is zero-characterized on $X$ by \cref{metric_zero_characterization}.
    Show that for all $x,y$ such that $x\in X$ and $y\in X$ we have $d((x,y)) = d((y,x))$.
    \begin{subproof}
        Fix $x$.
        Fix $y$ such that $x\in X$ and $y\in X$.
        Let $A_1 = \uniformMetricSet{J}{x}{y}$.
        Let $A_2 = \uniformMetricSet{J}{y}{x}$.
        For all $a$ such that $a\in A_1$ we have $a\in A_2$
            by \cref{uniform_metric_set,reals_power,tuple_power,funs_type_apply,times_tuple_intro,metric_on,metric_symmetric}.
        $A_1\subseteq A_2$ by \cref{subseteq}.
        For all $a$ such that $a\in A_2$ we have $a\in A_1$
            by \cref{uniform_metric_set,reals_power,tuple_power,funs_type_apply,times_tuple_intro,metric_on,metric_symmetric}.
        $A_2\subseteq A_1$ by \cref{subseteq}.
        Hence $A_1 = A_2$ by \cref{subseteq_antisymmetric}.
        Let $r_1 = d((x,y))$.
        Let $r_2 = d((y,x))$.
        $(x,y)\in\dom{d}$ and $(y,x)\in\dom{d}$ by \cref{times_tuple_intro,funs_elim}.
        $((x,y),r_1)\in d$ and $((y,x),r_2)\in d$ by \cref{funs_elim,function_apply_elim}.
        Take $p_1,x_1,y_1,s_1$ such that
            $p_1\in X\times X$ and $x_1\in X$ and $y_1\in X$ and $s_1\in\reals$ and
            $((x,y),r_1) = (p_1,s_1)$ and $p_1 = (x_1,y_1)$ and
            $\realSupremum{s_1}{\uniformMetricSet{J}{x_1}{y_1}}$ by \cref{uniform_metric}.
        Take $p_2,x_2,y_2,s_2$ such that
            $p_2\in X\times X$ and $x_2\in X$ and $y_2\in X$ and $s_2\in\reals$ and
            $((y,x),r_2) = (p_2,s_2)$ and $p_2 = (x_2,y_2)$ and
            $\realSupremum{s_2}{\uniformMetricSet{J}{x_2}{y_2}}$ by \cref{uniform_metric}.
        $p_1 = (x,y)$ and $r_1 = s_1$ and $p_2 = (y,x)$ and $r_2 = s_2$ by \cref{pair_eq_iff}.
        We have $x_1 = x$ and $y_1 = y$ and $x_2 = y$ and $y_2 = x$ by \cref{pair_eq_iff}.
        $r_1 = r_2$ by \cref{real_supremum_pred,real_supremum_unique}.
        Hence $d((x,y)) = d((y,x))$.
    \end{subproof}
    $d$ is symmetric on $X$ by \cref{metric_symmetric}.
    Show that for all $x,y,z$ such that $x\in X$ and $y\in X$ and $z\in X$ we have
        $d((x,y)) + d((y,z)) \geq d((x,z))$.
    \begin{subproof}
        Fix $x$.
        Fix $y$.
        Fix $z$ such that $x\in X$ and $y\in X$ and $z\in X$.
        Let $A_1 = \uniformMetricSet{J}{x}{y}$.
        Let $A_2 = \uniformMetricSet{J}{y}{z}$.
        Let $A_3 = \uniformMetricSet{J}{x}{z}$.
        Let $r_1 = d((x,y))$.
        Let $r_2 = d((y,z))$.
        Let $r_3 = d((x,z))$.
        $(x,y)\in\dom{d}$ and $(y,z)\in\dom{d}$ and $(x,z)\in\dom{d}$ by \cref{times_tuple_intro,funs_elim}.
        $((x,y),r_1)\in d$ and $((y,z),r_2)\in d$ and $((x,z),r_3)\in d$
            by \cref{funs_elim,function_apply_elim}.
        Take $p_1,x_1,y_1,s_1$ such that
            $p_1\in X\times X$ and $x_1\in X$ and $y_1\in X$ and $s_1\in\reals$ and
            $((x,y),r_1) = (p_1,s_1)$ and $p_1 = (x_1,y_1)$ and
            $\realSupremum{s_1}{\uniformMetricSet{J}{x_1}{y_1}}$ by \cref{uniform_metric}.
        Take $p_2,x_2,y_2,s_2$ such that
            $p_2\in X\times X$ and $x_2\in X$ and $y_2\in X$ and $s_2\in\reals$ and
            $((y,z),r_2) = (p_2,s_2)$ and $p_2 = (x_2,y_2)$ and
            $\realSupremum{s_2}{\uniformMetricSet{J}{x_2}{y_2}}$ by \cref{uniform_metric}.
        Take $p_3,x_3,y_3,s_3$ such that
            $p_3\in X\times X$ and $x_3\in X$ and $y_3\in X$ and $s_3\in\reals$ and
            $((x,z),r_3) = (p_3,s_3)$ and $p_3 = (x_3,y_3)$ and
            $\realSupremum{s_3}{\uniformMetricSet{J}{x_3}{y_3}}$ by \cref{uniform_metric}.
        $p_1 = (x,y)$ and $r_1 = s_1$ and $p_2 = (y,z)$ and $r_2 = s_2$ and
            $p_3 = (x,z)$ and $r_3 = s_3$ by \cref{pair_eq_iff}.
        $r_1\in\reals$ and $r_2\in\reals$ and $r_3\in\reals$.
        We have $x_1 = x$ and $y_1 = y$ and $x_2 = y$ and $y_2 = z$ and
            $x_3 = x$ and $y_3 = z$ by \cref{pair_eq_iff}.
        For all $a$ such that $a\in A_3$ we have $a\in\reals$
            by \cref{uniform_metric_set,reals_power,tuple_power,funs_type_apply,times_tuple_intro,metric_on}.
        $A_3\subseteq\reals$ by \cref{subseteq}.
        Show that for all $a\in A_3$ we have $a \leq r_1 + r_2$.
        \begin{subproof}
            Fix $a$ such that $a\in A_3$.
            $a\in\reals$ by \cref{subseteq}.
            Take $\alpha\in J$ such that $a = \apply{e}{(x(\alpha),z(\alpha))}$ by \cref{uniform_metric_set}.
            Let $a_1 = \apply{e}{(x(\alpha),y(\alpha))}$.
            Let $a_2 = \apply{e}{(y(\alpha),z(\alpha))}$.
            $a_1\in A_1$ and $a_2\in A_2$ by \cref{uniform_metric_set}.
            $x(\alpha)\in\reals$ and $y(\alpha)\in\reals$ and $z(\alpha)\in\reals$ by \cref{reals_power,tuple_power,funs_type_apply}.
            $(x(\alpha),y(\alpha))\in\reals\times\reals$ and
                $(y(\alpha),z(\alpha))\in\reals\times\reals$ and
                $(x(\alpha),z(\alpha))\in\reals\times\reals$ by \cref{times_tuple_intro}.
            We have $a_1\in\reals$ and $a_2\in\reals$ by \cref{funs_type_apply}.
            $a \leq a_1 + a_2$ by \cref{metric_on,metric_triangle_inequality}.
            $a_1 \leq r_1$ and $a_2 \leq r_2$ by \cref{real_supremum_pred,real_supremum,real_upper_bound}.
            $a_1 + a_2\in\reals$ and $r_1 + a_2\in\reals$ and $r_1 + r_2\in\reals$ by \cref{real_add_closed}.
            $a_1 + a_2 \leq r_1 + a_2$ by \cref{real_leq_add}.
            $r_1 + a_2 \leq r_1 + r_2$ by \cref{real_leq_add,real_add_comm}.
            Hence $a \leq r_1 + r_2$ by \cref{real_leq_trans}.
        \end{subproof}
        $r_1 + r_2$ is an upper bound of $A_3$ by \cref{real_upper_bound,real_add_closed}.
        Hence $r_3 \leq r_1 + r_2$ by \cref{real_supremum_pred,real_supremum}.
        $d((x,z)) \leq d((x,y)) + d((y,z))$.
        Hence $d((x,y)) + d((y,z)) \geq d((x,z))$.
    \end{subproof}
    $d$ satisfies the triangle inequality on $X$ by \cref{metric_triangle_inequality}.
    Hence $d$ is a metric on $X$ by \cref{metric_on}.
\end{proof}

% from §20 Item 29: uniform topology is finer than product topology
\begin{theorem}\label{uniform_finer_product}
    Suppose $J$ is inhabited.
    Let $X = \{(\alpha,\reals) \mid \alpha\in J\}$.
    Let $T = \{(\alpha,\standardTopologyR) \mid \alpha\in J\}$.
    Let $T_0 = \productTopologyIndexed{T}{X}$.
    Then $\uniformTopology{J}$ is finer than $T_0$ on $\realsPower{J}$.
\end{theorem}
\begin{proof}
    Let $X_0 = \realsPower{J}$.
    Let $d = \uniformMetric{J}$.
    Let $T_1 = \uniformTopology{J}$.
    For all $w$ such that $w\in X$ there exist $x,y$ such that $w = (x,y)$.
    For all $\alpha,y_1,y_2$ such that $(\alpha,y_1)\in X$ and $(\alpha,y_2)\in X$ we have $y_1 = y_2$.
    $X$ is a function by \cref{relation,rightunique}.
    For all $w$ such that $w\in T$ there exist $x,y$ such that $w = (x,y)$.
    For all $\alpha,y_1,y_2$ such that $(\alpha,y_1)\in T$ and $(\alpha,y_2)\in T$ we have $y_1 = y_2$.
    $T$ is a function by \cref{relation,rightunique}.
    $J\subseteq\dom{X}$ and $J\subseteq\dom{T}$ by \cref{dom_intro,subseteq}.
    For all $\alpha$ such that $\alpha\in\dom{X}$ we have $\alpha\in J$.
    For all $\alpha$ such that $\alpha\in\dom{T}$ we have $\alpha\in J$.
    $\dom{X}\subseteq J$ and $\dom{T}\subseteq J$ by \cref{subseteq}.
    $\dom{X} = J$ by \cref{subseteq_antisymmetric}.
    $\dom{T} = J$ by \cref{subseteq_antisymmetric}.
    For all $\alpha\in\dom{X}$ we have $\apply{X}{\alpha} = \reals$ and $\apply{T}{\alpha} = \standardTopologyR$
        by \cref{subseteq,function_apply_intro}.
    For all $\alpha\in\dom{X}$ we have $\apply{T}{\alpha}$ is a topology on $\apply{X}{\alpha}$
        by \cref{subseteq,function_apply_intro,standard_basis_is_basis,basis_topology_is_topology,standard_topology_reals}.
    Let $S_0 = \productSubbasis{T}{X}$.
    Let $B_0 = \finiteIntersections{S_0}$.
    $B_0$ is a basis for $T_0$ on $\productFamily{X}$ by \cref{product_subbasis_finite_intersections_basis}.
    $B_0$ is a basis for $T_0$ on $X_0$ by \cref{product_family_reals_power}.
    $d$ is a metric on $X_0$ by \cref{uniform_metric_is_metric}.
    $\dom{d} = X_0\times X_0$ by \cref{metric_on,funs_elim}.
    Let $B_U = \metricBasis{d}{X_0}$.
    $B_U$ is a basis on $X_0$ by \cref{metric_basis_is_basis}.
    $T_1 = \metricTopology{d}{X_0}$ by \cref{uniform_topology}.
    $T_1 = \basisTopology{B_U}{X_0}$ by \cref{metric_topology}.
    $B_U$ is a basis for $T_1$ on $X_0$ by \cref{basis_for_topology}.
    Show that for all $x,B_2$ such that $x\in X_0$ and $B_2\in B_0$ and $x\in B_2$
        there exists $B_3\in B_U$ such that $x\in B_3$ and $B_3\subseteq B_2$.
    \begin{subproof}
        Fix $x$.
        Fix $B_2$ such that $x\in X_0$ and $B_2\in B_0$ and $x\in B_2$.
        Take $F\subseteq S_0$ such that $F$ is finite and inhabited and $B_2 = \inters{F}$ by \cref{finite_intersections}.
        Let $R = \{w\in F\times\reals \mid \exists \beta\in J.\exists U\in\standardTopologyR.
            \fst{w} = \preimg{\piIndex{X}{\beta}}{U} \land \apply{x}{\beta}\in U \land
            0 < \snd{w} \land
            \snd{w} \leq 1 \land
            \metricBall{\standardMetricR}{\reals}{\apply{x}{\beta}}{\snd{w}} \subseteq U\}$.
        For all $w$ such that $w\in R$ there exist $y,z$ such that $w = (y,z)$.
        $R$ is a relation by \cref{relation}.
        Show that for all $A\in F$ there exists $\epsilon\in\reals$ such that $A\mathrel{R}\epsilon$.
        \begin{subproof}
            Fix $A$ such that $A\in F$.
            $x\in A$ by \cref{inters_iff_forall}.
            Take $\beta\in\dom{X}$ such that $A\in\productSubbasisAt{T}{X}{\beta}$
                by \cref{subseteq,product_subbasis_union_indexed}.
            Take $U\in\apply{T}{\beta}$ such that $A = \preimg{\piIndex{X}{\beta}}{U}$ by \cref{product_subbasis_indexed}.
            $\beta\in J$ by \cref{subseteq}.
            $x\in\preimg{\piIndex{X}{\beta}}{U}$.
            Take $z\in U$ such that $x\mathrel{\piIndex{X}{\beta}} z$ by \cref{preimg_iff}.
            $(x,z)\in\piIndex{X}{\beta}$.
            Take $x_1\in\productFamily{X}$ such that $(x,z) = (x_1,\apply{x_1}{\beta})$ by \cref{projection_map_indexed}.
            $\apply{x}{\beta}\in U$ by \cref{pair_eq_iff}.
            $U\in\standardTopologyR$ by \cref{function_apply_intro}.
            $\standardBasisR$ is a basis on $\reals$ by \cref{standard_basis_is_basis}.
            $U\subseteq\reals$
                by \cref{basis_topology_is_topology,standard_topology_reals,topology_on,subseteq,pow_iff}.
            $\standardMetricR$ is a metric on $\reals$ by \cref{standard_metric_is_metric}.
            $U$ is open in $\reals$ with respect to $\metricTopology{\standardMetricR}{\reals}$ by
                \cref{standard_metric_topology,open_in,basis_topology_is_topology,metric_basis_is_basis,metric_topology}.
            Take $\epsilon\in\reals$ such that $0 < \epsilon$ and
                $\metricBall{\standardMetricR}{\reals}{\apply{x}{\beta}}{\epsilon}\subseteq U$
                by \cref{metric_open_iff}.
            $1\in\reals$ by \cref{real_naturals_subset_reals,real_one_in_naturals,subseteq}.
            Let $C = \{\epsilon,1\}$.
            $C$ is finite and $\epsilon\in C$ and $1\in C$ by \cref{pair_is_finite,upair_intro_left,upair_intro_right}.
            $C\neq\emptyset$ and $C\subseteq\reals$ by \cref{upair_intro_left,emptyset,upair_elim,subseteq}.
            Take $\epsilon_0\in C$ such that $\epsilon_0$ is a minimum of $C$ by \cref{finite_real_minimum}.
            $0 < \epsilon_0$ by \cref{upair_elim,real_zero_lt_one,real_lt_subst_right}.
            $\epsilon_0 \leq 1$ by \cref{upair_intro_right,real_minimum}.
            Show that for all $z_0$ such that $z_0\in\metricBall{\standardMetricR}{\reals}{\apply{x}{\beta}}{\epsilon_0}$ we have $z_0\in U$.
            \begin{subproof}
                Fix $z_0$ such that $z_0\in\metricBall{\standardMetricR}{\reals}{\apply{x}{\beta}}{\epsilon_0}$.
                $z_0\in\reals$ and $\apply{\standardMetricR}{(\apply{x}{\beta},z_0)} < \epsilon_0$ by \cref{metric_ball}.
                $\apply{x}{\beta}\in\reals$ by \cref{reals_power,tuple_power,funs_type_apply}.
                $(\apply{x}{\beta},z_0)\in\reals\times\reals$ by \cref{times_tuple_intro}.
                $\apply{\standardMetricR}{(\apply{x}{\beta},z_0)}\in\reals$ by
                    \cref{standard_metric_is_metric,metric_on,funs_type_apply}.
                $\epsilon_0\in\reals$ and $\epsilon_0 \leq \epsilon$ by \cref{subseteq,upair_intro_left,real_minimum}.
                $\apply{\standardMetricR}{(\apply{x}{\beta},z_0)} < \epsilon$
                    by \cref{real_leq_split,real_lt_trans,real_lt_subst_right}.
                $z_0\in U$ by \cref{metric_ball,subseteq}.
            \end{subproof}
            $\metricBall{\standardMetricR}{\reals}{\apply{x}{\beta}}{\epsilon_0}\subseteq U$ and $\epsilon_0\in\reals$ by \cref{subseteq}.
            $A\mathrel{R}\epsilon_0$ by \cref{times_tuple_intro,fst_eq,snd_eq,pair_eq_iff}.
        \end{subproof}
        Take $e$ such that $e$ is a function on $F$ and for all $A\in F$ we have $A\mathrel{R}\apply{e}{A}$
            by \cref{choice_function}.
        $e$ is a function and $\dom{e} = F$.
        Let $E = \img{e}{F}$.
        $E$ is finite by \cref{finite_image_function}.
        Take $A_0$ such that $A_0\in F$.
        $E\neq\emptyset$ by \cref{function_apply_elim,img_iff,emptyset}.
        For all $t$ such that $t\in E$ we have $t\in\reals$ and $0 < t$
            by \cref{img_iff,subseteq,function_apply_iff,times_tuple_elim,snd_eq}.
        $E\subseteq\reals$ by \cref{subseteq}.
        Take $\delta\in E$ such that $\delta$ is a minimum of $E$ by \cref{finite_real_minimum}.
        $\delta\in\reals$ and $0 < \delta$ by \cref{subseteq}.
        Let $B_3 = \metricBall{d}{X_0}{x}{\delta}$.
        $B_3\in B_U$ by \cref{metric_ball,subseteq,powerset_intro,metric_basis}.
        Show that for all $y$ such that $y\in B_3$ we have $y\in B_2$.
        \begin{subproof}
            Fix $y$ such that $y\in B_3$.
            $y\in X_0$ and $d((x,y)) < \delta$ by \cref{metric_ball}.
            $(x,y)\in\dom{d}$ by \cref{times_tuple_intro,subseteq}.
            Take $s$ such that $((x,y),s)\in d$ by \cref{dom_iff}.
            $d((x,y)) = s$ by \cref{metric_on,funs_elim,subseteq,function_apply_intro}.
            $s < \delta$ by \cref{real_lt_subst_left}.
            Take $p_0,x_0,y_0,r_0$ such that
                $p_0\in X_0\times X_0$ and $x_0\in X_0$ and $y_0\in X_0$ and $r_0\in\reals$ and
                $((x,y),s) = (p_0,r_0)$ and $p_0 = (x_0,y_0)$ and
                $\realSupremum{r_0}{\uniformMetricSet{J}{x_0}{y_0}}$ by \cref{uniform_metric}.
            $p_0 = (x,y)$ and $s = r_0$ by \cref{pair_eq_iff}.
            $r_0 < \delta$ by \cref{real_lt_subst_left}.
            $x_0 = x$ and $y_0 = y$ by \cref{pair_eq_iff}.
            Show that for all $A\in F$ we have $y\in A$.
            \begin{subproof}
                Fix $A$ such that $A\in F$.
                $(A,\apply{e}{A})\in R$.
                Let $w = (A,\apply{e}{A})$.
                $w\in R$.
                Take $\beta,U$ such that $\beta\in J$ and $U\in\standardTopologyR$ and
                    $\fst{w} = \preimg{\piIndex{X}{\beta}}{U}$ and $\apply{x}{\beta}\in U$ and
                    $0 < \snd{w}$ and
                    $\snd{w} \leq 1$ and
                    $\metricBall{\standardMetricR}{\reals}{\apply{x}{\beta}}{\snd{w}} \subseteq U$.
                $\fst{w} = A$ and $\snd{w} = \apply{e}{A}$ by \cref{fst_eq,snd_eq}.
                $A = \preimg{\piIndex{X}{\beta}}{U}$.
                $0 < \apply{e}{A}$.
                $\apply{e}{A} \leq 1$ by \cref{real_leq_split}.
                $\metricBall{\standardMetricR}{\reals}{\apply{x}{\beta}}{\apply{e}{A}} \subseteq U$.
                $\apply{e}{A}\in E$ by \cref{function_apply_elim,img_iff}.
                $\apply{e}{A}\in\reals$ by \cref{subseteq}.
                $\delta \leq \apply{e}{A}$ by \cref{function_apply_elim,img_iff,real_minimum}.
                $r_0 < \apply{e}{A}$ by \cref{real_leq_split,real_lt_trans,real_lt_subst_right}.
                Let $m = \apply{\boundedMetric{\standardMetricR}}{(x(\beta),y(\beta))}$.
                $m \leq r_0$ by \cref{uniform_metric_set,real_supremum_pred,real_supremum,real_upper_bound}.
                $x(\beta)\in\reals$ and $y(\beta)\in\reals$ by \cref{reals_power,tuple_power,funs_type_apply}.
                $(x(\beta),y(\beta))\in\reals\times\reals$ by \cref{times_tuple_intro}.
                $m\in\reals$ by \cref{bounded_metric_is_metric,standard_metric_is_metric,metric_on,funs_type_apply}.
                $m < r_0$ or $m = r_0$ by \cref{real_leq_split}.
                $m < \apply{e}{A}$ by \cref{real_lt_subst_left,real_lt_trans}.
                $1\in\reals$ by \cref{real_naturals_subset_reals,real_one_in_naturals,subseteq}.
                $m < 1$ by \cref{real_leq_split,real_lt_trans,real_lt_subst_right}.
                $m = \apply{\standardMetricR}{(x(\beta),y(\beta))}$ by
                    \cref{bounded_metric_apply_value_lt_one,standard_metric_is_metric}.
                $\apply{\standardMetricR}{(x(\beta),y(\beta))} < \apply{e}{A}$ by \cref{real_lt_subst_left}.
                $\apply{y}{\beta}\in\metricBall{\standardMetricR}{\reals}{\apply{x}{\beta}}{\apply{e}{A}}$ by
                    \cref{metric_ball}.
                $\apply{y}{\beta}\in U$ by \cref{subseteq}.
                $y\in\preimg{\piIndex{X}{\beta}}{U}$ by
                    \cref{metric_ball,product_family_reals_power,subseteq,projection_map_indexed,preimg_iff}.
                $y\in A$.
            \end{subproof}
            $y\in B_2$ by \cref{inters_iff_forall,subseteq_antisymmetric}.
        \end{subproof}
        $B_3\subseteq B_2$ by \cref{subseteq}.
        $x\in B_3$ and $B_3\subseteq B_2$ by \cref{metric_ball,metric_zero_characterization,metric_on}.
    \end{subproof}
    Show that $T_1$ is finer than $T_0$ on $X_0$.
    \begin{subproof}
        Follows by \cref{basis_finer_criterion}.
    \end{subproof}
\end{proof}


% from §20 helper lemma 11b: positive real has a smaller positive half
\begin{lemma}\label{positive_real_smaller_half}
    Suppose $r\in\reals$ and $0 < r$.
    Let $s = r \rmul \inv{(1+1)}$.
    Then $s\in\reals$ and $0 < s$ and $s < r$.
\end{lemma}
\begin{proof}
    $1\in\reals$ and $1+1\in\reals$ and $0 < 1 + 1$ and $1 < 1 + 1$ by
        \cref{real_add_ident,real_naturals_subset_reals,real_one_in_naturals,subseteq,real_zero_lt_one,real_add_closed,real_lt_add,real_lt_subst_left,real_lt_trans}.
    $1 + 1 \neq 0$ and $\inv{(1+1)}\in\reals$ and $0 < \inv{(1+1)}$ by
        \cref{real_pos_nonzero,real_mul_inverse,real_inv_pos}.
    Show that $s\in\reals$ and $0 < s$.
    \begin{subproof}
        Follows by \cref{real_mul_closed,real_mul_pos_pos_gt}.
    \end{subproof}
    Show that $s < r$.
    \begin{subproof}
        $1 \rmul s < (1 + 1) \rmul s$ by \cref{real_lt_mul_pos}.
        $s < r$ by
            \cref{real_mul_assoc,real_mul_comm,real_mul_inverse,real_mul_ident,real_lt_subst_left,real_lt_subst_right}.
    \end{subproof}
\end{proof}

% from §20 helper lemma 11c: bounded standard metric basis generates the standard topology
\begin{lemma}\label{standard_metric_basis_standard_topology}
    $\metricBasis{\boundedMetric{\standardMetricR}}{\reals}$ is a basis for $\standardTopologyR$ on $\reals$.
\end{lemma}
\begin{proof}
    Follows by \cref{standard_metric_is_metric,bounded_metric_is_metric,metric_basis_is_basis,metric_topology,basis_for_topology,bounded_metric_topology,standard_metric_topology,subseteq}.
\end{proof}

% from §20 helper lemma 11d: the center of a metric ball belongs to the ball
\begin{lemma}\label{metric_ball_center}
    Suppose $d$ is a metric on $X$.
    Suppose $x\in X$.
    Suppose $\epsilon\in\reals$ and $0 < \epsilon$.
    Then $x\in\metricBall{d}{X}{x}{\epsilon}$.
\end{lemma}
\begin{proof}
    Follows by \cref{metric_on,metric_zero_characterization,real_lt_subst_left,metric_ball}.
\end{proof}


% from §20 Item 30: uniform topology is coarser than box topology
\begin{theorem}\label{uniform_coarser_box}
    Suppose $J$ is inhabited.
    Let $X = \{(\alpha,\reals) \mid \alpha\in J\}$.
    Let $T = \{(\alpha,\standardTopologyR) \mid \alpha\in J\}$.
    Let $T_0 = \boxTopology{T}{X}$.
    Then $\uniformTopology{J}$ is coarser than $T_0$ on $\realsPower{J}$.
\end{theorem}
\begin{proof}
    Let $X_0 = \realsPower{J}$.
    Let $d = \uniformMetric{J}$.
    Let $T_1 = \uniformTopology{J}$.
    Let $B_U = \metricBasis{d}{X_0}$.
    Let $B = \{(\alpha,\metricBasis{\boundedMetric{\standardMetricR}}{\reals}) \mid \alpha\in J\}$.
    Let $C = \{A\in\pow{\productFamily{X}} \mid
        \exists U. U\in\funs{\dom{X}}{\pow{\unions{\ran{X}}}}\land
        (\forall \alpha\in\dom{X}. \apply{U}{\alpha}\in\apply{B}{\alpha})\land
        A = \productFamily{U}\}$.
    $X$ is a function.
    $T$ is a function.
    $B$ is a function.
    For all $\alpha$ such that $\alpha\in J$ we have $\alpha\in\dom{X}$ and $\alpha\in\dom{T}$ and $\alpha\in\dom{B}$
        by \cref{dom_intro}.
    $J\subseteq\dom{X}$ and $J\subseteq\dom{T}$ and $J\subseteq\dom{B}$ by \cref{subseteq}.
    For all $\alpha$ such that $\alpha\in\dom{X}$ we have $\alpha\in J$ by \cref{dom_iff,pair_eq_iff}.
    $\dom{X}\subseteq J$ by \cref{subseteq}.
    $\dom{X} = J$ by \cref{subseteq_antisymmetric}.
    For all $\alpha$ such that $\alpha\in\dom{T}$ we have $\alpha\in J$ by \cref{dom_iff,pair_eq_iff}.
    $\dom{T}\subseteq J$ by \cref{subseteq}.
    $\dom{T} = J$ by \cref{subseteq_antisymmetric}.
    For all $\alpha$ such that $\alpha\in\dom{B}$ we have $\alpha\in J$ by \cref{dom_iff,pair_eq_iff}.
    $\dom{B}\subseteq J$ by \cref{subseteq}.
    $\dom{B} = J$ by \cref{subseteq_antisymmetric}.
    For all $\alpha$ such that $\alpha\in\dom{X}$ we have
        $\apply{X}{\alpha} = \reals$ and
        $\apply{T}{\alpha} = \standardTopologyR$ and
        $\apply{B}{\alpha} = \metricBasis{\boundedMetric{\standardMetricR}}{\reals}$
        by \cref{subseteq,function_apply_intro}.
    For all $\alpha$ such that $\alpha\in\dom{X}$ we have $\apply{B}{\alpha}$ is a basis for $\apply{T}{\alpha}$ on $\apply{X}{\alpha}$
        by \cref{subseteq,standard_metric_basis_standard_topology}.
    $C$ is a basis for $T_0$ on $\productFamily{X}$ by \cref{box_basis_from_bases}.
    $C$ is a basis for $T_0$ on $X_0$ by \cref{product_family_reals_power,basis_for_topology}.
    $d$ is a metric on $X_0$ by \cref{uniform_metric_is_metric}.
    $\dom{d} = X_0\times X_0$ by \cref{metric_on,funs_elim}.
    $B_U$ is a basis for $T_1$ on $X_0$ by \cref{metric_basis_is_basis,basis_for_topology,uniform_topology,metric_topology}.
    Show that for all $x,B_2$ such that $x\in X_0$ and $B_2\in B_U$ and $x\in B_2$
        there exists $B_3\in C$ such that $x\in B_3$ and $B_3\subseteq B_2$.
    \begin{subproof}
        Fix $x$.
        Fix $B_2$ such that $x\in X_0$ and $B_2\in B_U$ and $x\in B_2$.
        Take $z,\epsilon$ such that $z\in X_0$ and $\epsilon\in\reals$ and $0 < \epsilon$
            and $B_2 = \metricBall{d}{X_0}{z}{\epsilon}$ by \cref{metric_basis}.
        $d((z,x)) < \epsilon$ by \cref{subseteq,metric_ball}.
        $(z,x)\in X_0\times X_0$ by \cref{times_tuple_intro}.
        $d((z,x))\in\reals$ by \cref{metric_on,funs_type_apply}.
        Let $r = \epsilon - d((z,x))$.
        $0 < r$ by \cref{real_lt_imp_pos_diff}.
        $r\in\reals$ by \cref{real_sub,real_add_inverse,real_add_closed}.
        Let $\delta = r \rmul \inv{(1+1)}$.
        $\delta\in\reals$ and $0 < \delta$ and $\delta < r$ by \cref{positive_real_smaller_half}.
        Let $R = \{w\in J\times\pow{\unions{\ran{X}}} \mid \exists \alpha\in J.
            \fst{w} = \alpha \land \snd{w} = \metricBall{\boundedMetric{\standardMetricR}}{\reals}{\apply{x}{\alpha}}{\delta}\}$.
        For all $w$ such that $w\in R$ there exist $y,z$ such that $w = (y,z)$.
        $R$ is a relation by \cref{relation}.
        Show that for all $\alpha\in J$ there exists $V\in\pow{\unions{\ran{X}}}$ such that $\alpha\mathrel{R}V$.
        \begin{subproof}
            Fix $\alpha$ such that $\alpha\in J$.
            Let $V = \metricBall{\boundedMetric{\standardMetricR}}{\reals}{\apply{x}{\alpha}}{\delta}$.
            $V\in\pow{\unions{\ran{X}}}$ by \cref{metric_ball,ran_intro,unions_iff,subseteq,powerset_intro}.
            $\alpha\mathrel{R}V$ by \cref{times_tuple_intro,fst_eq,snd_eq,pair_eq_iff}.
        \end{subproof}
        Take $U$ such that $U$ is a function on $J$ and for all $\alpha\in J$ we have $\alpha\mathrel{R}\apply{U}{\alpha}$
            by \cref{choice_function}.
        $U$ is a function and $\dom{U} = J$.
        For all $\alpha\in J$ we have $(\alpha,\apply{U}{\alpha})\in R$ by \cref{subseteq}.
        For all $\alpha\in J$ we have
            $\apply{U}{\alpha} = \metricBall{\boundedMetric{\standardMetricR}}{\reals}{\apply{x}{\alpha}}{\delta}$
            by \cref{fst_eq,snd_eq,pair_eq_iff}.
        For all $\alpha\in\dom{X}$ we have $\apply{U}{\alpha}\in\apply{B}{\alpha}$ by
            \cref{subseteq,reals_power,tuple_power,funs_type_apply,metric_ball,powerset_intro,metric_basis,bounded_metric_is_metric,standard_metric_is_metric}.
        For all $\alpha$ such that $\alpha\in\dom{U}$ we have $\apply{U}{\alpha}\in\pow{\unions{\ran{X}}}$
            by \cref{subseteq,times_tuple_elim}.
        $U\in\funs{\dom{X}}{\pow{\unions{\ran{X}}}}$ by \cref{funs_intro}.
        Let $B_3 = \productFamily{U}$.
        For all $\alpha\in\dom{X}$ we have $\apply{U}{\alpha}\subseteq\apply{X}{\alpha}$ by \cref{subseteq,metric_ball}.
        $B_3\subseteq\productFamily{X}$ by \cref{indexed_product_subset_from_coordinate_subsets}.
        $B_3\in\pow{\productFamily{X}}$ by \cref{powerset_intro}.
        $B_3\in C$ by definition.
        For all $\alpha\in\dom{U}$ we have $\apply{x}{\alpha}\in\apply{U}{\alpha}$ by
            \cref{subseteq,reals_power,tuple_power,funs_type_apply,metric_ball_center,bounded_metric_is_metric,standard_metric_is_metric}.
        For all $\alpha$ such that $\alpha\in\dom{U}$ we have $\apply{x}{\alpha}\in\unions{\ran{U}}$
            by \cref{subseteq,funs_elim,function_apply_elim,ran_intro,unions_iff}.
        $x\in\funs{\dom{U}}{\unions{\ran{U}}}$ by \cref{subseteq,reals_power,tuple_power,funs_elim,funs_intro}.
        $x\in B_3$ by \cref{indexed_product}.
        Show that for all $y$ such that $y\in B_3$ we have $y\in B_2$.
        \begin{subproof}
            Fix $y$ such that $y\in B_3$.
            $y\in X_0$ by \cref{subseteq,product_family_reals_power}.
            $(x,y)\in X_0\times X_0$ by \cref{times_tuple_intro}.
            $(z,y)\in X_0\times X_0$ by \cref{times_tuple_intro}.
            $d((x,y))\in\reals$ and $d((z,y))\in\reals$ by \cref{metric_on,funs_type_apply}.
            $(x,y)\in\dom{d}$ and $(z,y)\in\dom{d}$ by \cref{subseteq}.
            Take $s$ such that $((x,y),s)\in d$ by \cref{dom_iff}.
            $d((x,y)) = s$ by \cref{metric_on,funs_elim,subseteq,function_apply_intro}.
            Take $p_0,x_0,y_0,r_0$ such that
                $p_0\in X_0\times X_0$ and $x_0\in X_0$ and $y_0\in X_0$ and $r_0\in\reals$ and
                $((x,y),s) = (p_0,r_0)$ and $p_0 = (x_0,y_0)$ and
                $\realSupremum{r_0}{\uniformMetricSet{J}{x_0}{y_0}}$ by \cref{uniform_metric}.
            $p_0 = (x,y)$ and $s = r_0$ by \cref{pair_eq_iff}.
            $x_0 = x$ and $y_0 = y$ by \cref{pair_eq_iff}.
            For all $a\in\uniformMetricSet{J}{x}{y}$ we have $a\in\reals$ by \cref{uniform_metric_set,reals_power,tuple_power,funs_type_apply,times_tuple_intro,bounded_metric_is_metric,standard_metric_is_metric,metric_on}.
            $\uniformMetricSet{J}{x}{y}\subseteq\reals$ by \cref{subseteq}.
            Show that for all $a\in\uniformMetricSet{J}{x}{y}$ we have $a \leq \delta$.
            \begin{subproof}
                Fix $a$ such that $a\in\uniformMetricSet{J}{x}{y}$.
                Take $\alpha\in J$ such that $a = \apply{\boundedMetric{\standardMetricR}}{(x(\alpha),y(\alpha))}$ by \cref{uniform_metric_set}.
                $y(\alpha)\in\metricBall{\boundedMetric{\standardMetricR}}{\reals}{\apply{x}{\alpha}}{\delta}$ by \cref{indexed_product,subseteq}.
                $a < \delta$ by \cref{metric_ball}.
                $a \leq \delta$.
            \end{subproof}
            $\delta$ is an upper bound of $\uniformMetricSet{J}{x}{y}$ by \cref{real_upper_bound}.
            $r_0 \leq \delta$ by \cref{real_supremum_pred,real_supremum,real_upper_bound}.
            $d((x,y)) \leq \delta$ by \cref{real_lt_subst_left,real_lt_subst_right}.
            $(z,x)\in X_0\times X_0$ by \cref{times_tuple_intro}.
            $d((z,y)) \leq d((z,x)) + d((x,y))$ by \cref{metric_on,metric_triangle_inequality}.
            $d((z,x)) + d((x,y))\in\reals$ and $d((z,x)) + \delta\in\reals$ and $d((z,x)) + r\in\reals$
                by \cref{real_add_closed}.
            $d((z,x)) + d((x,y)) \leq d((z,x)) + \delta$ by \cref{real_leq_add,real_add_comm}.
            $d((z,x)) + \delta < d((z,x)) + r$ by \cref{real_lt_add,real_add_comm}.
            $d((z,x)) + r = \epsilon$ by
                \cref{real_sub,real_add_assoc,real_add_inverse,real_add_comm,real_add_ident}.
            $d((z,x)) + \delta < \epsilon$ by \cref{real_lt_subst_right}.
            $y\in B_2$ by \cref{real_leq_split,real_lt_trans,metric_ball,subseteq}.
        \end{subproof}
        $B_3\subseteq B_2$ by \cref{subseteq}.
    \end{subproof}
    Show that $T_0$ is finer than $T_1$ on $X_0$.
    \begin{subproof}
        Follows by \cref{basis_finer_criterion}.
    \end{subproof}
\end{proof}

% from §20 helper lemma 13: one divided by one
\begin{lemma}\label{real_one_div_one}
    $1 / 1 = 1$.
\end{lemma}
\begin{proof}
    Follows by \cref{real_mul_ident,real_zero_lt_one,real_pos_nonzero,real_mul_inverse_unique,real_div,subseteq}.
\end{proof}

% from §20 helper lemma 14: reciprocal decreases at successors
\begin{lemma}\label{naturalsplus_reciprocal_step}
    Suppose $n\in\naturalsPlus$.
    Then $1 / (n + 1) < 1 / n$.
\end{lemma}
\begin{proof}
    $n\in\reals$ and $n + 1\in\reals$ and $0 < n$ and $0 < n + 1$ by
        \cref{real_naturals_subset_reals,subseteq,real_naturals_closed_succ,naturalsplus_positive_real}.
    $n \neq 0$ and $n + 1 \neq 0$ and
        $\inv{n}\in\reals$ and $\inv{(n + 1)}\in\reals$ and $0 < \inv{n}$ and $0 < \inv{(n + 1)}$
        by \cref{real_pos_nonzero,real_mul_inverse,real_inv_pos}.
    Let $c = \inv{n} \rmul \inv{(n + 1)}$.
    $c\in\reals$ and $0 < c$ by \cref{real_mul_closed,real_mul_pos_pos_gt}.
    $0\in\reals$ and $1\in\reals$ and $n < n + 1$ by
        \cref{real_add_ident,real_mul_ident,real_zero_lt_one,real_lt_add,real_add_comm,real_lt_subst_left,real_lt_subst_right}.
    $n \rmul c < (n + 1) \rmul c$ by \cref{real_lt_mul_pos}.
    $1 / (n + 1) < 1 / n$ by
        \cref{real_mul_assoc,real_mul_comm,real_mul_inverse,real_mul_ident,real_div,real_lt_subst_left,real_lt_subst_right}.
\end{proof}

% from §20 helper lemma 15: reciprocals are antitone on naturalsPlus
\begin{lemma}\label{naturalsplus_reciprocal_antitone}
    Suppose $m\in\naturalsPlus$ and $n\in\naturalsPlus$ and $m \leq n$.
    Then $1 / n \leq 1 / m$.
\end{lemma}
\begin{proof}
    Let $A = \{k\in\naturalsPlus \mid \text{if $m \leq k$ then $1 / k \leq 1 / m$}\}$.
    For all $k$ such that $k\in A$ we have $k\in\naturalsPlus$ by \cref{subseteq}.
    $A\subseteq\naturalsPlus$ by \cref{subseteq}.
    Show that if $m \leq 1$ then $1 / 1 \leq 1 / m$.
    \begin{subproof}
        Follows by \cref{real_mul_ident,real_naturals_subset_reals,subseteq,real_one_leq_naturalsplus,real_leq_antisym}.
    \end{subproof}
    $1\in A$.
    Show that for all $k\in A$ we have $k + 1\in A$.
    \begin{subproof}
        Fix $k$ such that $k\in A$.
        $k\in\naturalsPlus$ by \cref{subseteq}.
        Show that if $m \leq k + 1$ then $1 / (k + 1) \leq 1 / m$.
        \begin{subproof}
            Assume $m \leq k + 1$.
            $m < k + 1$ or $m = k + 1$ by \cref{real_leq_split}.
            \begin{byCase}
            \caseOf{$m < k + 1$.}
                $m < k$ or $m = k$ or $m = k + 1$ by \cref{real_succ_discrete}.
                \begin{byCase}
                \caseOf{$m < k$.}
                    $m \leq k$.
                    $1 / (k + 1)\in\reals$ and $1 / k\in\reals$ and $1 / m\in\reals$ by
                        \cref{naturalsplus_reciprocal_real,real_naturals_closed_succ,subseteq}.
                    $1 / k \leq 1 / m$ by \cref{subseteq}.
                    $1 / (k + 1) < 1 / k$ by \cref{naturalsplus_reciprocal_step}.
                    $1 / k < 1 / m$ or $1 / k = 1 / m$ by \cref{real_leq_split}.
                    $1 / (k + 1) \leq 1 / m$ by
                        \cref{real_leq_split,real_lt_trans,real_lt_subst_right}.
                \caseOf{$m = k$.}
                    $1 / k = 1 / m$ by \cref{subseteq}.
                    $1 / (k + 1) < 1 / k$ by \cref{naturalsplus_reciprocal_step}.
                    $1 / (k + 1) \leq 1 / m$ by \cref{real_lt_subst_right}.
                \caseOf{$m = k + 1$.}
                    $1 / (k + 1) \leq 1 / m$ by \cref{subseteq}.
                \end{byCase}
            \caseOf{$m = k + 1$.}
                $1 / (k + 1) \leq 1 / m$ by \cref{subseteq}.
            \end{byCase}
        \end{subproof}
        $k + 1\in A$ by \cref{real_naturals_closed_succ}.
    \end{subproof}
    For all $k\in\naturalsPlus$ we have $k\in A$ by \cref{real_naturals_induction}.
    $1 / n \leq 1 / m$ by \cref{subseteq}.
\end{proof}

% from §20 helper lemma 16: cancel a positive natural denominator in a strict inequality
\begin{lemma}\label{real_div_lt_same_denom_cancel}
    Suppose $a\in\reals$ and $b\in\reals$ and $n\in\naturalsPlus$ and $a / n < b / n$.
    Then $a < b$.
\end{lemma}
\begin{proof}
    $n\in\reals$ and $0 < n$ by \cref{real_naturals_subset_reals,subseteq,naturalsplus_positive_real}.
    $n \neq 0$ by \cref{real_pos_nonzero}.
    $\inv{n}\in\reals$ by \cref{real_mul_inverse}.
    $a / n\in\reals$ and $b / n\in\reals$ by \cref{real_div,real_mul_closed}.
    $(a / n) \rmul n < (b / n) \rmul n$ by \cref{real_lt_mul_pos}.
    $(a / n) \rmul n = a$ and $(b / n) \rmul n = b$ by
        \cref{real_div,real_mul_assoc,real_mul_inverse,real_mul_comm,real_mul_ident}.
    $a < b$ by \cref{real_lt_subst_left,real_lt_subst_right}.
\end{proof}

% from §20 helper lemma 17: division by a positive natural preserves weak order
\begin{lemma}\label{real_div_leq_same_denom}
    Suppose $a\in\reals$ and $b\in\reals$ and $n\in\naturalsPlus$ and $a \leq b$.
    Then $a / n \leq b / n$.
\end{lemma}
\begin{proof}
    $n\in\reals$ and $0 < n$ and $\inv{n}\in\reals$ and $0 < \inv{n}$ by
        \cref{real_naturals_subset_reals,subseteq,naturalsplus_positive_real,real_pos_nonzero,real_mul_inverse,real_inv_pos}.
    $a < b$ or $a = b$ by \cref{real_leq_split}.
    \begin{byCase}
    \caseOf{$a < b$.}
        $a \rmul \inv{n} < b \rmul \inv{n}$ by \cref{real_lt_mul_pos}.
        $a / n \leq b / n$ by \cref{real_div,real_lt_subst_left,real_lt_subst_right}.
    \caseOf{$a = b$.}
        $a / n \leq b / n$ by \cref{subseteq}.
    \end{byCase}
\end{proof}

% from §20 helper lemma 18: division by a positive natural distributes over sums
\begin{lemma}\label{real_div_add_same_denom}
    Suppose $a\in\reals$ and $b\in\reals$ and $n\in\naturalsPlus$.
    Then $(a + b) / n = a / n + b / n$.
\end{lemma}
\begin{proof}
    $n\in\reals$ and $\inv{n}\in\reals$ and $a + b\in\reals$ by
        \cref{real_naturals_subset_reals,subseteq,real_mul_inverse,naturalsplus_positive_real,real_pos_nonzero,real_add_closed}.
    $(a + b) / n = (a \rmul \inv{n}) + (b \rmul \inv{n})$ by
        \cref{real_div,real_mul_comm,real_distrib,real_add_eq_subst}.
    $(a + b) / n = a / n + b / n$ by \cref{real_div,real_lt_subst_left,real_lt_subst_right}.
\end{proof}

% from §20 helper lemma 19: division of a nonnegative real by a positive natural is nonnegative
\begin{lemma}\label{real_div_naturalsplus_nonnegative}
    Suppose $a\in\reals$ and $0 \leq a$ and $n\in\naturalsPlus$.
    Then $0 \leq a / n$.
\end{lemma}
\begin{proof}
    $0\in\reals$ and $0 / n = 0$ by
        \cref{real_add_ident,real_naturals_subset_reals,subseteq,naturalsplus_positive_real,real_pos_nonzero,real_mul_inverse,real_div,real_mul_zero}.
    $0 / n \leq a / n$ by \cref{real_div_leq_same_denom}.
    $0 \leq a / n$ by \cref{real_lt_subst_left}.
\end{proof}

% from §20 helper lemma 20: zero quotient with positive natural denominator
\begin{lemma}\label{real_div_naturalsplus_zero}
    Suppose $a\in\reals$ and $0 \leq a$ and $n\in\naturalsPlus$ and $a / n = 0$.
    Then $a = 0$.
\end{lemma}
\begin{proof}
    Suppose not.
    $0 < a$ by \cref{real_leq_split}.
    $n\in\reals$ and $0 < n$ and $\inv{n}\in\reals$ and $0 < \inv{n}$ and $0\in\reals$ by
        \cref{real_naturals_subset_reals,subseteq,naturalsplus_positive_real,real_pos_nonzero,real_mul_inverse,real_inv_pos,real_add_ident}.
    $0 < a \rmul \inv{n}$ by \cref{real_lt_mul_pos,real_mul_zero,real_lt_subst_left}.
    $0 < a / n$ by \cref{real_div,real_lt_subst_right}.
    $0 < 0$ by \cref{real_lt_subst_right}.
    Contradiction by \cref{real_lt_irreflexive,real_add_ident}.
\end{proof}

% from §20 helper lemma 21: quotient by a positive natural is real
\begin{lemma}\label{real_div_naturalsplus_real}
    Suppose $a\in\reals$ and $n\in\naturalsPlus$.
    Then $a / n\in\reals$.
\end{lemma}
\begin{proof}
    $a / n\in\reals$ by
        \cref{real_naturals_subset_reals,subseteq,naturalsplus_positive_real,real_pos_nonzero,real_mul_inverse,real_div,real_mul_closed}.
\end{proof}

% from §20 helper lemma 22: weak-strong transitivity for reals
\begin{lemma}\label{real_leq_lt_trans}
    Suppose $a\in\reals$ and $b\in\reals$ and $c\in\reals$ and $a \leq b$ and $b < c$.
    Then $a < c$.
\end{lemma}
\begin{proof}
    $a < c$ by \cref{real_leq_split,real_lt_trans,real_lt_subst_left}.
\end{proof}

% from §20 helper lemma 12: omega coordinate metric set
\begin{definition}\label{omega_metric_set}
    $\omegaMetricSet{x}{y} = \{\apply{\boundedMetric{\standardMetricR}}{(x(i),y(i))} / i \mid i\in\naturalsPlus\}$.
\end{definition}

% from §20 Item 31: metric on $\reals^{\omega}$
\begin{definition}\label{omega_metric}
    $\omegaMetric = \{w \mid \exists p\in\realsPower{\naturalsPlus}\times\realsPower{\naturalsPlus}.
        \exists x\in\realsPower{\naturalsPlus}. \exists y\in\realsPower{\naturalsPlus}. \exists r\in\reals.
        w = (p,r) \land p = (x,y) \land
        \realSupremum{r}{\omegaMetricSet{x}{y}}\}$.
\end{definition}

% from §20 Item 32: $\omega$-metric is a metric
\begin{theorem}\label{omega_metric_is_metric}
    $\omegaMetric$ is a metric on $\realsPower{\naturalsPlus}$.
\end{theorem}
\begin{proof}
    Let $X = \realsPower{\naturalsPlus}$.
    Let $d = \omegaMetric$.
    Let $e = \boundedMetric{\standardMetricR}$.
    $\standardMetricR$ is a metric on $\reals$ by \cref{standard_metric_is_metric}.
    Hence $e$ is a metric on $\reals$ and $e\in\funs{\reals\times\reals}{\reals}$
        by \cref{bounded_metric_is_metric,metric_on}.
    Show that $d\in\funs{X\times X}{\reals}$.
    \begin{subproof}
        Show that for all $p,z_1,z_2$ such that $(p,z_1)\in d$ and $(p,z_2)\in d$ we have $z_1 = z_2$.
        \begin{subproof}
            Fix $p$.
            Fix $z_1$.
            Fix $z_2$ such that $(p,z_1)\in d$ and $(p,z_2)\in d$.
            Take $p_1,x_1,y_1,r_1$ such that
                $p_1\in X\times X$ and $x_1\in X$ and $y_1\in X$ and $r_1\in\reals$ and
                $(p,z_1) = (p_1,r_1)$ and $p_1 = (x_1,y_1)$ and
                $\realSupremum{r_1}{\omegaMetricSet{x_1}{y_1}}$ by \cref{omega_metric}.
            Take $p_2,x_2,y_2,r_2$ such that
                $p_2\in X\times X$ and $x_2\in X$ and $y_2\in X$ and $r_2\in\reals$ and
                $(p,z_2) = (p_2,r_2)$ and $p_2 = (x_2,y_2)$ and
                $\realSupremum{r_2}{\omegaMetricSet{x_2}{y_2}}$ by \cref{omega_metric}.
            $p = p_1$ and $z_1 = r_1$ and $p = p_2$ and $z_2 = r_2$ by \cref{pair_eq_iff}.
            $p_1 = p_2$.
            $(x_1,y_1) = (x_2,y_2)$ by \cref{pair_eq_iff}.
            $x_1 = x_2$ and $y_1 = y_2$ by \cref{pair_eq_iff}.
            $\omegaMetricSet{x_1}{y_1} = \omegaMetricSet{x_2}{y_2}$.
            $r_1 = r_2$ by \cref{real_supremum_pred,real_supremum_unique}.
            $z_1 = z_2$.
        \end{subproof}
        $d$ is a function by \cref{omega_metric,relation,rightunique}.
        Show that for all $p$ such that $p\in X\times X$ we have $p\in\dom{d}$.
        \begin{subproof}
            Fix $p$ such that $p\in X\times X$.
            Take $x,y$ such that $p = (x,y)$ and $x\in X$ and $y\in X$ by \cref{times_elem_is_tuple}.
            Let $A = \omegaMetricSet{x}{y}$.
            For all $a$ such that $a\in A$ we have $a\in\reals$
                by \cref{omega_metric_set,reals_power,tuple_power,funs_type_apply,times_tuple_intro,metric_on,real_div_naturalsplus_real}.
            $A\subseteq\reals$ by \cref{subseteq}.
            $1\in\naturalsPlus$ by \cref{real_one_in_naturals}.
            $A\neq\emptyset$ by \cref{omega_metric_set,emptyset}.
            Show that for all $a\in A$ we have $a \leq 1$.
            \begin{subproof}
                Fix $a$ such that $a\in A$.
                Take $i\in\naturalsPlus$ such that $a = \apply{e}{(x(i),y(i))} / i$ by \cref{omega_metric_set}.
                Let $m = \apply{e}{(x(i),y(i))}$.
                $a = m / i$.
                $x(i)\in\reals$ and $y(i)\in\reals$ by \cref{reals_power,tuple_power,funs_type_apply}.
                $(x(i),y(i))\in\reals\times\reals$ by \cref{times_tuple_intro}.
                $1\in\reals$ by \cref{real_mul_ident}.
                $m\in\reals$ and $0 \leq m$ by \cref{metric_on,funs_type_apply,metric_nonnegative}.
                $m \leq 1$ by \cref{bounded_metric_leq_one}.
                $m / i \leq 1 / i$ by \cref{real_div_leq_same_denom}.
                $1 \leq i$ by \cref{real_one_leq_naturalsplus}.
                $1 / i \leq 1$ by \cref{naturalsplus_reciprocal_antitone,real_one_div_one,subseteq}.
                $m / i\in\reals$ and $1 / i\in\reals$ by \cref{real_div_naturalsplus_real,naturalsplus_reciprocal_real}.
                $m / i \leq 1$ by \cref{real_leq_trans}.
                $a \leq 1$ by \cref{subseteq}.
            \end{subproof}
            $1\in\reals$ by \cref{real_mul_ident}.
            $1$ is an upper bound of $A$ by \cref{real_upper_bound}.
            $A$ is bounded above by \cref{real_bounded_above}.
            Take $r$ such that $r$ is a supremum of $A$ by \cref{real_completeness_axiom}.
            $(x,y)\in X\times X$ by \cref{times_tuple_intro}.
            $r\in\reals$ by \cref{real_supremum,real_upper_bound}.
            $((x,y),r)\in d$ by \cref{real_supremum_pred,omega_metric}.
            $p\in\dom{d}$ by \cref{dom_intro,pair_eq_iff}.
        \end{subproof}
        $X\times X\subseteq\dom{d}$ by \cref{subseteq}.
        For all $p$ such that $p\in\dom{d}$ we have $p\in X\times X$
            by \cref{dom_iff,omega_metric,pair_eq_iff,subseteq}.
        $\dom{d}\subseteq X\times X$ by \cref{subseteq}.
        $\dom{d} = X\times X$ by \cref{subseteq_antisymmetric}.
        For all $p\in\dom{d}$ we have $d(p)\in\reals$
            by \cref{dom_iff,omega_metric,pair_eq_iff,function_apply_iff,real_lt_subst_left}.
        $d\in\funs{X\times X}{\reals}$ by \cref{funs_intro}.
    \end{subproof}
    Show that for all $x,y$ such that $x\in X$ and $y\in X$ we have $d((x,y)) \geq 0$.
    \begin{subproof}
        Fix $x$.
        Fix $y$ such that $x\in X$ and $y\in X$.
        Let $r = d((x,y))$.
        $(x,y)\in\dom{d}$ by \cref{times_tuple_intro,funs_elim}.
        $((x,y),r)\in d$ by \cref{funs_elim,function_apply_elim}.
        Take $p_0,x_0,y_0,r_0$ such that
            $p_0\in X\times X$ and $x_0\in X$ and $y_0\in X$ and $r_0\in\reals$ and
            $((x,y),r) = (p_0,r_0)$ and $p_0 = (x_0,y_0)$ and
            $\realSupremum{r_0}{\omegaMetricSet{x_0}{y_0}}$ by \cref{omega_metric}.
        $p_0 = (x,y)$ and $r = r_0$ by \cref{pair_eq_iff}.
        $r\in\reals$.
        $x_0 = x$ and $y_0 = y$ by \cref{pair_eq_iff}.
        Let $a = \apply{e}{(x(1),y(1))} / 1$.
        $a \leq r$ by \cref{omega_metric_set,real_one_in_naturals,real_supremum_pred,real_supremum,real_upper_bound}.
        $x(1)\in\reals$ and $y(1)\in\reals$ by \cref{reals_power,tuple_power,real_one_in_naturals,funs_type_apply}.
        $(x(1),y(1))\in\reals\times\reals$ by \cref{times_tuple_intro}.
        $a\in\reals$ and $0 \leq a$ by
            \cref{real_lt_subst_left,metric_on,funs_type_apply,metric_nonnegative,real_div_naturalsplus_real,real_div_naturalsplus_nonnegative,real_one_in_naturals,subseteq}.
        $0 \leq r$ by \cref{real_add_ident,real_leq_trans}.
        $d((x,y)) \geq 0$.
    \end{subproof}
    Show that $d$ is nonnegative on $X$.
    \begin{subproof}
        Follows by \cref{metric_nonnegative}.
    \end{subproof}
    Show that for all $x,y$ such that $x\in X$ and $y\in X$ we have if $d((x,y)) = 0$ then $x = y$.
    \begin{subproof}
        Fix $x$.
        Fix $y$ such that $x\in X$ and $y\in X$.
        Assume $d((x,y)) = 0$.
        Let $r = d((x,y))$.
        $(x,y)\in\dom{d}$ by \cref{times_tuple_intro,funs_elim}.
        $((x,y),r)\in d$ by \cref{funs_elim,function_apply_elim}.
        Take $p_0,x_0,y_0,r_0$ such that
            $p_0\in X\times X$ and $x_0\in X$ and $y_0\in X$ and $r_0\in\reals$ and
            $((x,y),r) = (p_0,r_0)$ and $p_0 = (x_0,y_0)$ and
            $\realSupremum{r_0}{\omegaMetricSet{x_0}{y_0}}$ by \cref{omega_metric}.
        $p_0 = (x,y)$ and $r = r_0$ by \cref{pair_eq_iff}.
        $x_0 = x$ and $y_0 = y$ by \cref{pair_eq_iff}.
        Show that for all $i\in\naturalsPlus$ we have $x(i) = y(i)$.
        \begin{subproof}
            Fix $i$ such that $i\in\naturalsPlus$.
            Let $a = \apply{e}{(x(i),y(i))} / i$.
            $a \leq r$ by \cref{omega_metric_set,real_supremum_pred,real_supremum,real_upper_bound}.
            Let $m = \apply{e}{(x(i),y(i))}$.
            $a = m / i$.
            $x(i)\in\reals$ and $y(i)\in\reals$ by \cref{reals_power,tuple_power,funs_type_apply}.
            $(x(i),y(i))\in\reals\times\reals$ by \cref{times_tuple_intro}.
            $m\in\reals$ and $0 \leq m$ by \cref{metric_on,funs_type_apply,metric_nonnegative}.
            $a\in\reals$ and $0 \leq a$ by
                \cref{real_div_naturalsplus_real,real_div_naturalsplus_nonnegative,subseteq}.
            $a \leq 0$ by \cref{real_supremum_pred,real_supremum,real_upper_bound,real_lt_subst_left,real_lt_subst_right}.
            $a = 0$ by \cref{real_add_ident,real_leq_antisym,real_lt_subst_right}.
            $m = 0$ by \cref{real_div_naturalsplus_zero,subseteq}.
            $x(i) = y(i)$ by \cref{metric_on,metric_zero_characterization}.
        \end{subproof}
        $\dom{x} = \naturalsPlus$ and $\dom{y} = \naturalsPlus$ by \cref{reals_power,tuple_power,funs_elim}.
        $\dom{x}\subseteq\dom{y}$ and $\dom{y}\subseteq\dom{x}$ by \cref{subseteq_refl}.
        $x\subseteq y$ and $y\subseteq x$ by \cref{reals_power,tuple_power,funs_elim,fun_subseteq,subseteq}.
        $x = y$ by \cref{subseteq_antisymmetric}.
    \end{subproof}
    Show that for all $x,y$ such that $x\in X$ and $y\in X$ we have if $x = y$ then $d((x,y)) = 0$.
    \begin{subproof}
        Fix $x$.
        Fix $y$ such that $x\in X$ and $y\in X$.
        Assume $x = y$.
        Let $A = \omegaMetricSet{x}{y}$.
        Show that for all $a\in A$ we have $a \leq 0$.
        \begin{subproof}
            Fix $a$ such that $a\in A$.
            Take $i\in\naturalsPlus$ such that $a = \apply{e}{(x(i),y(i))} / i$ by \cref{omega_metric_set}.
            $x(i) = y(i)$ by \cref{real_lt_subst_left}.
            $x(i)\in\reals$ and $y(i)\in\reals$ by \cref{reals_power,tuple_power,funs_type_apply}.
            $(x(i),y(i))\in\reals\times\reals$ by \cref{times_tuple_intro}.
            $\apply{e}{(x(i),y(i))} = 0$ by \cref{metric_on,metric_zero_characterization}.
            $a = 0 / i$ by \cref{subseteq}.
            $a = 0 \rmul \inv{i}$ by \cref{real_div}.
            $i\in\reals$ and $0 < i$ and $\inv{i}\in\reals$ by
                \cref{real_naturals_subset_reals,subseteq,naturalsplus_positive_real,real_pos_nonzero,real_mul_inverse}.
            $a = 0$ by \cref{real_mul_zero}.
            $a \leq 0$.
        \end{subproof}
        $0\in\reals$ by \cref{real_add_ident}.
        For all $a$ such that $a\in A$ we have $a\in\reals$
            by \cref{omega_metric_set,reals_power,tuple_power,funs_type_apply,times_tuple_intro,metric_on,real_div_naturalsplus_real}.
        $A\subseteq\reals$ and $0$ is an upper bound of $A$ by \cref{subseteq,real_add_ident,real_upper_bound}.
        $1\in\naturalsPlus$ by \cref{real_one_in_naturals}.
        $\apply{e}{(x(1),y(1))} / 1\in A$ by \cref{omega_metric_set}.
        $x(1) = y(1)$ by \cref{real_lt_subst_left}.
        $x(1)\in\reals$ and $y(1)\in\reals$ by \cref{reals_power,tuple_power,real_one_in_naturals,funs_type_apply}.
        $(x(1),y(1))\in\reals\times\reals$ by \cref{times_tuple_intro}.
        $\apply{e}{(x(1),y(1))} = 0$ by \cref{metric_on,metric_zero_characterization}.
        $\apply{e}{(x(1),y(1))} / 1 = 0 / 1$ by \cref{real_lt_subst_left}.
        $0 / 1 = 0$ by
            \cref{real_div,real_mul_ident,real_zero_lt_one,real_pos_nonzero,real_mul_inverse,real_mul_zero,real_lt_subst_left}.
        $0\in A$ by \cref{real_lt_subst_right,subseteq}.
        For all $b$ such that $b$ is an upper bound of $A$ we have $0 \leq b$ by \cref{real_upper_bound}.
        $0$ is a supremum of $A$ by \cref{real_supremum}.
        $((x,y),0)\in d$ by \cref{real_supremum_pred,omega_metric,pair_eq_iff,times_tuple_intro,real_add_ident}.
        $(x,y)\in\dom{d}$ by \cref{dom_intro}.
        $d((x,y)) = 0$ by \cref{funs_elim,function_apply_iff}.
    \end{subproof}
    Show that $d$ is zero-characterized on $X$.
    \begin{subproof}
        Follows by \cref{metric_zero_characterization}.
    \end{subproof}
    Show that for all $x,y$ such that $x\in X$ and $y\in X$ we have $d((x,y)) = d((y,x))$.
    \begin{subproof}
        Fix $x$.
        Fix $y$ such that $x\in X$ and $y\in X$.
        Let $A_1 = \omegaMetricSet{x}{y}$.
        Let $A_2 = \omegaMetricSet{y}{x}$.
        For all $a$ such that $a\in A_1$ we have $a\in A_2$
            by \cref{omega_metric_set,reals_power,tuple_power,funs_type_apply,times_tuple_intro,metric_on,metric_symmetric,subseteq}.
        $A_1\subseteq A_2$ by \cref{subseteq}.
        For all $a$ such that $a\in A_2$ we have $a\in A_1$
            by \cref{omega_metric_set,reals_power,tuple_power,funs_type_apply,times_tuple_intro,metric_on,metric_symmetric,subseteq}.
        $A_2\subseteq A_1$ by \cref{subseteq}.
        $A_1 = A_2$ by \cref{subseteq_antisymmetric}.
        Let $r_1 = d((x,y))$.
        Let $r_2 = d((y,x))$.
        $(x,y)\in\dom{d}$ and $(y,x)\in\dom{d}$ by \cref{times_tuple_intro,funs_elim}.
        $((x,y),r_1)\in d$ and $((y,x),r_2)\in d$ by \cref{funs_elim,function_apply_elim}.
        Take $p_1,x_1,y_1,s_1$ such that
            $p_1\in X\times X$ and $x_1\in X$ and $y_1\in X$ and $s_1\in\reals$ and
            $((x,y),r_1) = (p_1,s_1)$ and $p_1 = (x_1,y_1)$ and
            $\realSupremum{s_1}{\omegaMetricSet{x_1}{y_1}}$ by \cref{omega_metric}.
        Take $p_2,x_2,y_2,s_2$ such that
            $p_2\in X\times X$ and $x_2\in X$ and $y_2\in X$ and $s_2\in\reals$ and
            $((y,x),r_2) = (p_2,s_2)$ and $p_2 = (x_2,y_2)$ and
            $\realSupremum{s_2}{\omegaMetricSet{x_2}{y_2}}$ by \cref{omega_metric}.
        $p_1 = (x,y)$ and $r_1 = s_1$ and $p_2 = (y,x)$ and $r_2 = s_2$ by \cref{pair_eq_iff}.
        $x_1 = x$ and $y_1 = y$ and $x_2 = y$ and $y_2 = x$ by \cref{pair_eq_iff}.
        $r_1 = r_2$ by \cref{real_supremum_pred,real_supremum_unique}.
        $d((x,y)) = d((y,x))$.
    \end{subproof}
    Show that $d$ is symmetric on $X$.
    \begin{subproof}
        Follows by \cref{metric_symmetric}.
    \end{subproof}
    Show that for all $x,y,z$ such that $x\in X$ and $y\in X$ and $z\in X$ we have
        $d((x,y)) + d((y,z)) \geq d((x,z))$.
    \begin{subproof}
        Fix $x$.
        Fix $y$.
        Fix $z$ such that $x\in X$ and $y\in X$ and $z\in X$.
        Let $A_1 = \omegaMetricSet{x}{y}$.
        Let $A_2 = \omegaMetricSet{y}{z}$.
        Let $A_3 = \omegaMetricSet{x}{z}$.
        Let $r_1 = d((x,y))$.
        Let $r_2 = d((y,z))$.
        Let $r_3 = d((x,z))$.
        $(x,y)\in\dom{d}$ and $(y,z)\in\dom{d}$ and $(x,z)\in\dom{d}$ by \cref{times_tuple_intro,funs_elim}.
        $((x,y),r_1)\in d$ and $((y,z),r_2)\in d$ and $((x,z),r_3)\in d$
            by \cref{funs_elim,function_apply_elim}.
        Take $p_1,x_1,y_1,s_1$ such that
            $p_1\in X\times X$ and $x_1\in X$ and $y_1\in X$ and $s_1\in\reals$ and
            $((x,y),r_1) = (p_1,s_1)$ and $p_1 = (x_1,y_1)$ and
            $\realSupremum{s_1}{\omegaMetricSet{x_1}{y_1}}$ by \cref{omega_metric}.
        Take $p_2,x_2,y_2,s_2$ such that
            $p_2\in X\times X$ and $x_2\in X$ and $y_2\in X$ and $s_2\in\reals$ and
            $((y,z),r_2) = (p_2,s_2)$ and $p_2 = (x_2,y_2)$ and
            $\realSupremum{s_2}{\omegaMetricSet{x_2}{y_2}}$ by \cref{omega_metric}.
        Take $p_3,x_3,y_3,s_3$ such that
            $p_3\in X\times X$ and $x_3\in X$ and $y_3\in X$ and $s_3\in\reals$ and
            $((x,z),r_3) = (p_3,s_3)$ and $p_3 = (x_3,y_3)$ and
            $\realSupremum{s_3}{\omegaMetricSet{x_3}{y_3}}$ by \cref{omega_metric}.
        $p_1 = (x,y)$ and $r_1 = s_1$ and $p_2 = (y,z)$ and $r_2 = s_2$ and
            $p_3 = (x,z)$ and $r_3 = s_3$ by \cref{pair_eq_iff}.
        $r_1\in\reals$ and $r_2\in\reals$ and $r_3\in\reals$.
        $x_1 = x$ and $y_1 = y$ and $x_2 = y$ and $y_2 = z$ and
            $x_3 = x$ and $y_3 = z$ by \cref{pair_eq_iff}.
        For all $a$ such that $a\in A_3$ we have $a\in\reals$
            by \cref{omega_metric_set,reals_power,tuple_power,funs_type_apply,times_tuple_intro,metric_on,real_div_naturalsplus_real}.
        $A_3\subseteq\reals$ by \cref{subseteq}.
        Show that for all $a\in A_3$ we have $a \leq r_1 + r_2$.
        \begin{subproof}
            Fix $a$ such that $a\in A_3$.
            $a\in\reals$ by \cref{subseteq}.
            Take $i\in\naturalsPlus$ such that $a = \apply{e}{(x(i),z(i))} / i$ by \cref{omega_metric_set}.
            Let $a_1 = \apply{e}{(x(i),y(i))} / i$.
            Let $a_2 = \apply{e}{(y(i),z(i))} / i$.
            $a_1\in A_1$ and $a_2\in A_2$ by \cref{omega_metric_set}.
            Let $m = \apply{e}{(x(i),z(i))}$.
            Let $m_1 = \apply{e}{(x(i),y(i))}$.
            Let $m_2 = \apply{e}{(y(i),z(i))}$.
            $a = m / i$ and $a_1 = m_1 / i$ and $a_2 = m_2 / i$.
            $x(i)\in\reals$ and $y(i)\in\reals$ and $z(i)\in\reals$ by \cref{reals_power,tuple_power,funs_type_apply}.
            $(x(i),y(i))\in\reals\times\reals$ and
                $(y(i),z(i))\in\reals\times\reals$ and
                $(x(i),z(i))\in\reals\times\reals$ by \cref{times_tuple_intro}.
            $m\in\reals$ and $m_1\in\reals$ and $m_2\in\reals$ by \cref{funs_type_apply}.
            $m_1 + m_2\in\reals$ by \cref{real_add_closed}.
            $m \leq m_1 + m_2$ by \cref{metric_on,metric_triangle_inequality}.
            $m / i \leq (m_1 + m_2) / i$ by \cref{real_div_leq_same_denom}.
            $(m_1 + m_2) / i = m_1 / i + m_2 / i$ by \cref{real_div_add_same_denom}.
            $a \leq a_1 + a_2$ by \cref{subseteq}.
            $a_1 \leq r_1$ and $a_2 \leq r_2$ by \cref{real_supremum_pred,real_supremum,real_upper_bound}.
            $a_1\in\reals$ and $a_2\in\reals$ by \cref{real_div_naturalsplus_real,subseteq}.
            $a_1 + a_2\in\reals$ and $r_1 + a_2\in\reals$ and $r_1 + r_2\in\reals$ by \cref{real_add_closed}.
            $a_1 + a_2 \leq r_1 + a_2$ by \cref{real_leq_add}.
            $r_1 + a_2 \leq r_1 + r_2$ by \cref{real_leq_add,real_add_comm}.
            $a \leq r_1 + r_2$ by \cref{real_leq_trans}.
        \end{subproof}
        $r_1 + r_2$ is an upper bound of $A_3$ by \cref{real_upper_bound,real_add_closed}.
        $r_3 \leq r_1 + r_2$ by \cref{real_supremum_pred,real_supremum}.
        $d((x,z)) \leq d((x,y)) + d((y,z))$.
        $d((x,y)) + d((y,z)) \geq d((x,z))$.
    \end{subproof}
    Show that $d$ satisfies the triangle inequality on $X$.
    \begin{subproof}
        Follows by \cref{metric_triangle_inequality}.
    \end{subproof}
    $d$ is a metric on $X$ by \cref{metric_on}.
\end{proof}

% from §20 Item 33: $\omega$-metric induces the product topology
\begin{theorem}\label{omega_metric_product_topology}
    Let $X = \{(i,\reals) \mid i\in\naturalsPlus\}$.
    Let $T = \{(i,\standardTopologyR) \mid i\in\naturalsPlus\}$.
    Let $T_0 = \productTopologyIndexed{T}{X}$.
    Then $\metricTopology{\omegaMetric}{\realsPower{\naturalsPlus}} = T_0$.
\end{theorem}
\begin{proof}
    Let $X_0 = \realsPower{\naturalsPlus}$.
    Let $d = \omegaMetric$.
    Let $T_1 = \metricTopology{d}{X_0}$.
    For all $w$ such that $w\in X$ there exist $x,y$ such that $w = (x,y)$.
    For all $i,y_1,y_2$ such that $(i,y_1)\in X$ and $(i,y_2)\in X$ we have $y_1 = y_2$.
    $X$ is a function by \cref{relation,rightunique}.
    For all $w$ such that $w\in T$ there exist $x,y$ such that $w = (x,y)$.
    For all $i,y_1,y_2$ such that $(i,y_1)\in T$ and $(i,y_2)\in T$ we have $y_1 = y_2$.
    $T$ is a function by \cref{relation,rightunique}.
    $\naturalsPlus\subseteq\dom{X}$ by \cref{dom_intro,subseteq}.
    For all $i$ such that $i\in\dom{X}$ we have $i\in\naturalsPlus$.
    $\dom{X}\subseteq\naturalsPlus$ by \cref{subseteq}.
    $\dom{X} = \naturalsPlus$ by \cref{subseteq_antisymmetric}.
    $\naturalsPlus\subseteq\dom{T}$ by \cref{dom_intro,subseteq}.
    For all $i$ such that $i\in\dom{T}$ we have $i\in\naturalsPlus$.
    $\dom{T}\subseteq\naturalsPlus$ by \cref{subseteq}.
    $\dom{T} = \naturalsPlus$ by \cref{subseteq_antisymmetric}.
    For all $i\in\dom{X}$ we have $\apply{X}{i} = \reals$ and $\apply{T}{i} = \standardTopologyR$
        by \cref{subseteq,function_apply_intro}.
    For all $i\in\dom{X}$ we have $\apply{T}{i}$ is a topology on $\apply{X}{i}$
        by \cref{subseteq,function_apply_intro,standard_basis_is_basis,basis_topology_is_topology,standard_topology_reals}.
    $1\in\naturalsPlus$ by \cref{real_one_in_naturals}.
    $1\in\dom{X}$ by \cref{subseteq}.
    Let $S_0 = \productSubbasis{T}{X}$.
    Let $B_0 = \finiteIntersections{S_0}$.
    $B_0$ is a basis for $T_0$ on $\productFamily{X}$ by \cref{product_subbasis_finite_intersections_basis}.
    $B_0$ is a basis for $T_0$ on $X_0$ by \cref{product_family_reals_power}.
    $d$ is a metric on $X_0$ by \cref{omega_metric_is_metric}.
    $\dom{d} = X_0\times X_0$ by \cref{metric_on,funs_elim}.
    Let $B_1 = \metricBasis{d}{X_0}$.
    $B_1$ is a basis on $X_0$ by \cref{metric_basis_is_basis}.
    $B_1$ is a basis for $T_1$ on $X_0$ by \cref{metric_topology,basis_for_topology}.
    Show that for all $x,B_2$ such that $x\in X_0$ and $B_2\in B_0$ and $x\in B_2$
        there exists $B_3\in B_1$ such that $x\in B_3$ and $B_3\subseteq B_2$.
    \begin{subproof}
        Fix $x$.
        Fix $B_2$ such that $x\in X_0$ and $B_2\in B_0$ and $x\in B_2$.
        Take $F\subseteq S_0$ such that $F$ is finite and inhabited and $B_2 = \inters{F}$ by \cref{finite_intersections}.
        Let $R = \{w\in F\times\reals \mid \exists \beta\in\naturalsPlus.\exists U\in\standardTopologyR.\exists \eta\in\reals.
            \fst{w} = \preimg{\piIndex{X}{\beta}}{U} \land \apply{x}{\beta}\in U \land
            0 < \eta \land
            \metricBall{\boundedMetric{\standardMetricR}}{\reals}{\apply{x}{\beta}}{\eta} \subseteq U \land
            0 < \snd{w} \land \snd{w} = \eta / \beta\}$.
        For all $w$ such that $w\in R$ there exist $y,z$ such that $w = (y,z)$.
        $R$ is a relation by \cref{relation}.
        Show that for all $A\in F$ there exists $\epsilon\in\reals$ such that $A\mathrel{R}\epsilon$.
        \begin{subproof}
            Fix $A$ such that $A\in F$.
            $x\in A$ by \cref{inters_iff_forall}.
            Take $\beta\in\dom{X}$ such that $A\in\productSubbasisAt{T}{X}{\beta}$
                by \cref{subseteq,product_subbasis_union_indexed}.
            Take $U\in\apply{T}{\beta}$ such that $A = \preimg{\piIndex{X}{\beta}}{U}$ by \cref{product_subbasis_indexed}.
            $\beta\in\naturalsPlus$ by \cref{subseteq}.
            $x\in\preimg{\piIndex{X}{\beta}}{U}$.
            Take $z\in U$ such that $x\mathrel{\piIndex{X}{\beta}} z$ by \cref{preimg_iff}.
            $(x,z)\in\piIndex{X}{\beta}$.
            Take $x_1\in\productFamily{X}$ such that $(x,z) = (x_1,\apply{x_1}{\beta})$ by \cref{projection_map_indexed}.
            $z = \apply{x}{\beta}$ by \cref{pair_eq_iff}.
            $\apply{x}{\beta}\in U$.
            $U\in\standardTopologyR$ by \cref{function_apply_intro}.
            $\standardBasisR$ is a basis on $\reals$ by \cref{standard_basis_is_basis}.
            $U\subseteq\reals$
                by \cref{basis_topology_is_topology,standard_topology_reals,topology_on,subseteq,pow_iff}.
            $\standardMetricR$ is a metric on $\reals$ by \cref{standard_metric_is_metric}.
            $\boundedMetric{\standardMetricR}$ is a metric on $\reals$ by \cref{bounded_metric_is_metric}.
            $U\in\metricTopology{\boundedMetric{\standardMetricR}}{\reals}$
                by \cref{standard_metric_topology,bounded_metric_topology,standard_metric_is_metric,subseteq}.
            $U$ is open in $\reals$ with respect to $\metricTopology{\boundedMetric{\standardMetricR}}{\reals}$ by
                \cref{open_in,basis_topology_is_topology,metric_basis_is_basis,metric_topology}.
            Take $\eta\in\reals$ such that $0 < \eta$ and
                $\metricBall{\boundedMetric{\standardMetricR}}{\reals}{\apply{x}{\beta}}{\eta}\subseteq U$
                by \cref{metric_open_iff}.
            Let $\epsilon = \eta / \beta$.
            $\epsilon\in\reals$ by \cref{real_div_naturalsplus_real}.
            $\beta\in\reals$ and $0 < \beta$ and $\inv{\beta}\in\reals$ and $0 < \inv{\beta}$ by
                \cref{real_naturals_subset_reals,subseteq,naturalsplus_positive_real,real_pos_nonzero,real_mul_inverse,real_inv_pos}.
            $0\in\reals$ by \cref{real_add_ident}.
            $0 \rmul \inv{\beta} < \eta \rmul \inv{\beta}$ by \cref{real_lt_mul_pos}.
            $0 < \eta \rmul \inv{\beta}$ by \cref{real_mul_zero,real_lt_subst_left}.
            $0 < \epsilon$ by \cref{real_div,real_lt_subst_right}.
            $(A,\epsilon)\in F\times\reals$ by \cref{times_tuple_intro}.
            $A\mathrel{R}\epsilon$ by \cref{fst_eq,snd_eq,pair_eq_iff}.
        \end{subproof}
        Take $e$ such that $e$ is a function on $F$ and for all $A\in F$ we have $A\mathrel{R}\apply{e}{A}$
            by \cref{choice_function}.
        $e$ is a function and $\dom{e} = F$.
        Let $E = \img{e}{F}$.
        $E$ is finite by \cref{finite_image_function}.
        Take $A_0$ such that $A_0\in F$.
        $\apply{e}{A_0}\in E$ by \cref{function_apply_elim,img_iff}.
        $E\neq\emptyset$ by \cref{function_apply_elim,img_iff,emptyset}.
        For all $t$ such that $t\in E$ we have $t\in\reals$ and $0 < t$
            by \cref{img_iff,subseteq,function_apply_iff,fst_eq,snd_eq,times_tuple_elim}.
        $E\subseteq\reals$ by \cref{subseteq}.
        Take $\delta\in E$ such that $\delta$ is a minimum of $E$ by \cref{finite_real_minimum}.
        $\delta\in\reals$ and $0 < \delta$ by \cref{subseteq}.
        Let $B_3 = \metricBall{d}{X_0}{x}{\delta}$.
        $B_3\in B_1$ by \cref{metric_ball,subseteq,powerset_intro,metric_basis}.
        Show that for all $y$ such that $y\in B_3$ we have $y\in B_2$.
        \begin{subproof}
            Fix $y$ such that $y\in B_3$.
            $y\in X_0$ and $d((x,y)) < \delta$ by \cref{metric_ball}.
            $(x,y)\in X_0\times X_0$ by \cref{times_tuple_intro}.
            $(x,y)\in\dom{d}$ by \cref{subseteq}.
            Take $s$ such that $((x,y),s)\in d$ by \cref{dom_iff}.
            $d((x,y)) = s$ by \cref{metric_on,funs_elim,subseteq,function_apply_intro}.
            $s < \delta$ by \cref{real_lt_subst_left}.
            Take $p_0,x_0,y_0,r_0$ such that
                $p_0\in X_0\times X_0$ and $x_0\in X_0$ and $y_0\in X_0$ and $r_0\in\reals$ and
                $((x,y),s) = (p_0,r_0)$ and $p_0 = (x_0,y_0)$ and
                $\realSupremum{r_0}{\omegaMetricSet{x_0}{y_0}}$ by \cref{omega_metric}.
            $p_0 = (x,y)$ and $s = r_0$ by \cref{pair_eq_iff}.
            $r_0 < \delta$ by \cref{real_lt_subst_left}.
            $x_0 = x$ and $y_0 = y$ by \cref{pair_eq_iff}.
            Show that for all $A\in F$ we have $y\in A$.
            \begin{subproof}
                Fix $A$ such that $A\in F$.
                $(A,\apply{e}{A})\in R$.
                Let $w = (A,\apply{e}{A})$.
                $w\in R$.
                Take $\beta,U,\eta$ such that $\beta\in\naturalsPlus$ and $U\in\standardTopologyR$ and $\eta\in\reals$ and
                    $\fst{w} = \preimg{\piIndex{X}{\beta}}{U}$ and $\apply{x}{\beta}\in U$ and
                    $0 < \eta$ and
                    $\metricBall{\boundedMetric{\standardMetricR}}{\reals}{\apply{x}{\beta}}{\eta}\subseteq U$ and
                    $0 < \snd{w}$ and $\snd{w} = \eta / \beta$.
                $\fst{w} = A$ and $\snd{w} = \apply{e}{A}$ by \cref{fst_eq,snd_eq}.
                $A = \preimg{\piIndex{X}{\beta}}{U}$.
                $\apply{e}{A}\in E$ by \cref{function_apply_elim,img_iff}.
                $\apply{e}{A}\in\reals$ by \cref{subseteq}.
                $\delta \leq \apply{e}{A}$ by \cref{function_apply_elim,img_iff,real_minimum}.
                $s < \apply{e}{A}$ by \cref{real_leq_split,real_lt_trans,real_lt_subst_right}.
                Let $m = \apply{\boundedMetric{\standardMetricR}}{(x(\beta),y(\beta))}$.
                Let $a = m / \beta$.
                $x(\beta)\in\reals$ and $y(\beta)\in\reals$ by \cref{reals_power,tuple_power,funs_type_apply}.
                $(x(\beta),y(\beta))\in\reals\times\reals$ by \cref{times_tuple_intro}.
                $\standardMetricR$ is a metric on $\reals$ by \cref{standard_metric_is_metric}.
                $\boundedMetric{\standardMetricR}$ is a metric on $\reals$ by \cref{bounded_metric_is_metric}.
                $m\in\reals$ by \cref{metric_on,funs_type_apply}.
                $a\in\reals$ by \cref{real_div_naturalsplus_real}.
                $a\in\omegaMetricSet{x}{y}$ by \cref{omega_metric_set}.
                $a \leq s$ by \cref{real_supremum_pred,real_supremum,real_upper_bound}.
                $a < \apply{e}{A}$ by \cref{real_leq_lt_trans}.
                $a = m / \beta$ and $\apply{e}{A} = \eta / \beta$ by \cref{subseteq}.
                $m < \eta$ by \cref{real_div_lt_same_denom_cancel,real_lt_subst_left,real_lt_subst_right}.
                $\apply{y}{\beta}\in\metricBall{\boundedMetric{\standardMetricR}}{\reals}{\apply{x}{\beta}}{\eta}$ by \cref{metric_ball}.
                $\apply{y}{\beta}\in U$ by \cref{subseteq}.
                $y\in\preimg{\piIndex{X}{\beta}}{U}$ by
                    \cref{metric_ball,product_family_reals_power,subseteq,projection_map_indexed,preimg_iff}.
                $y\in A$.
            \end{subproof}
            $y\in B_2$ by \cref{inters_iff_forall,subseteq_antisymmetric}.
        \end{subproof}
        $B_3\subseteq B_2$ by \cref{subseteq}.
        $x\in B_3$ and $B_3\subseteq B_2$ by \cref{metric_ball,metric_zero_characterization,metric_on}.
    \end{subproof}
    Show that $T_1$ is finer than $T_0$ on $X_0$.
    \begin{subproof}
        Follows by \cref{basis_finer_criterion}.
    \end{subproof}
    Show that for all $x,B_2$ such that $x\in X_0$ and $B_2\in B_1$ and $x\in B_2$
        there exists $B_3\in B_0$ such that $x\in B_3$ and $B_3\subseteq B_2$.
    \begin{subproof}
        Fix $x$.
        Fix $B_2$ such that $x\in X_0$ and $B_2\in B_1$ and $x\in B_2$.
        Take $z,\epsilon$ such that $z\in X_0$ and $\epsilon\in\reals$ and $0 < \epsilon$ and
            $B_2 = \metricBall{d}{X_0}{z}{\epsilon}$ by \cref{metric_basis}.
        $x\in\metricBall{d}{X_0}{z}{\epsilon}$ by \cref{subseteq}.
        $d((z,x)) < \epsilon$ by \cref{metric_ball}.
        $(z,x)\in X_0\times X_0$ by \cref{times_tuple_intro}.
        $d((z,x))\in\reals$ by \cref{metric_on,funs_type_apply}.
        Let $r = \epsilon - d((z,x))$.
        $0 < r$ by \cref{real_lt_imp_pos_diff}.
        $r\in\reals$ by \cref{real_sub,real_add_inverse,real_add_closed}.
        Let $\delta = r \rmul \inv{(1+1)}$.
        $\delta\in\reals$ and $0 < \delta$ and $\delta < r$ by \cref{positive_real_smaller_half}.
        Take $N\in\naturalsPlus$ such that $0 < 1 / N$ and $1 / N < \delta$ by \cref{real_small_reciprocal}.
        Let $J = \initialSegment{N}$.
        Let $R_g = \{w\in J\times\pow{\productFamily{X}} \mid \exists i\in J.
            \fst{w} = i \land
            \snd{w} = \preimg{\piIndex{X}{i}}{\metricBall{\boundedMetric{\standardMetricR}}{\reals}{\apply{x}{i}}{i \rmul \delta}}\}$.
        For all $w$ such that $w\in R_g$ there exist $i,A$ such that $w = (i,A)$.
        $R_g$ is a relation by \cref{relation}.
        For all $i\in J$ we have
            $\preimg{\piIndex{X}{i}}{\metricBall{\boundedMetric{\standardMetricR}}{\reals}{\apply{x}{i}}{i \rmul \delta}}
            \in\pow{\productFamily{X}}$
            by \cref{preimg_iff,projection_map_indexed,pair_eq_iff,subseteq,powerset_intro}.
        For all $i\in J$ there exists $A$ such that $i\mathrel{R_g} A$
            by \cref{times_tuple_intro,fst_eq,snd_eq,pair_eq_iff}.
        Take $g$ such that $g$ is a function on $J$ and for all $i\in J$ we have $i\mathrel{R_g}\apply{g}{i}$
            by \cref{choice_function}.
        Let $F = \img{g}{J}$.
        $F$ is finite by \cref{initial_segment_finite,finite_image_function}.
        $1\in J$ by \cref{real_one_in_naturals,real_one_leq_naturalsplus,initial_segment_naturalsplus}.
        $\apply{g}{1}\in F$ by \cref{function_apply_elim,img_iff}.
        $F\neq\emptyset$ by \cref{function_apply_elim,img_iff,emptyset}.
        Show that for all $A$ such that $A\in F$ we have $A\in S_0$.
        \begin{subproof}
            Fix $A$ such that $A\in F$.
            Take $j\in J$ such that $(j,A)\in g$ by \cref{img_iff}.
            $j\in\dom{g}$ and $A = \apply{g}{j}$ by \cref{subseteq,function_apply_iff}.
            $j\mathrel{R_g}\apply{g}{j}$.
            $(j,\apply{g}{j})\in R_g$.
            $(j,\apply{g}{j})\in J\times\pow{\productFamily{X}}$ by \cref{subseteq}.
            $A\in\pow{\productFamily{X}}$ by \cref{times_tuple_elim,snd_eq}.
            Take $i_0\in J$ such that
                $\fst{(j,\apply{g}{j})} = i_0$ and
                $\snd{(j,\apply{g}{j})}
                    = \preimg{\piIndex{X}{i_0}}{\metricBall{\boundedMetric{\standardMetricR}}{\reals}{\apply{x}{i_0}}{i_0 \rmul \delta}}$.
            $j = i_0$ and $A = \preimg{\piIndex{X}{j}}{\metricBall{\boundedMetric{\standardMetricR}}{\reals}{\apply{x}{j}}{j \rmul \delta}}$
                by \cref{function_apply_intro,fst_eq,snd_eq,pair_eq_iff}.
            $j\in\naturalsPlus$ and $j\in\reals$ and $0 < j$ by
                \cref{initial_segment_naturalsplus,real_naturals_subset_reals,subseteq,naturalsplus_positive_real}.
            $j \rmul \delta\in\reals$ and $0 < j \rmul \delta$ by \cref{real_mul_closed,real_mul_pos_pos_gt}.
            Let $U = \metricBall{\boundedMetric{\standardMetricR}}{\reals}{\apply{x}{j}}{j \rmul \delta}$.
            $\apply{x}{j}\in\reals$ by \cref{reals_power,tuple_power,funs_type_apply}.
            $U\in\metricBasis{\boundedMetric{\standardMetricR}}{\reals}$ by
                \cref{metric_ball,subseteq,powerset_intro,metric_basis,bounded_metric_is_metric,standard_metric_is_metric}.
            $U\in\standardTopologyR$ by
                \cref{standard_metric_basis_standard_topology,basis_for_topology,basis_elem_open_in_generated}.
            $U\in\apply{T}{j}$ by \cref{function_apply_intro}.
            $A\in\productSubbasisAt{T}{X}{j}$ by
                \cref{preimg_iff,projection_map_indexed,pair_eq_iff,subseteq,powerset_intro,product_subbasis_indexed}.
            $j\in\dom{X}$ by \cref{subseteq}.
            $A\in S_0$ by \cref{product_subbasis_union_indexed}.
        \end{subproof}
        $F\subseteq S_0$ by \cref{subseteq}.
        Let $B_3 = \inters{F}$.
        For all $u$ such that $u\in B_3$ we have $u\in\unions{S_0}$.
        $B_3\subseteq\unions{S_0}$ by \cref{subseteq}.
        $B_3\in\pow{\unions{S_0}}$ by \cref{powerset_intro}.
        $B_3\in B_0$ by \cref{finite_intersections}.
        For all $A$ such that $A\in F$ there exists $j\in J$ such that
            $A = \preimg{\piIndex{X}{j}}{\metricBall{\boundedMetric{\standardMetricR}}{\reals}{\apply{x}{j}}{j \rmul \delta}}$
            by \cref{img_iff,function_apply_intro,fst_eq,snd_eq,pair_eq_iff}.
        For all $A$ such that $A\in F$ we have $x\in A$
            by \cref{initial_segment_naturalsplus,real_naturals_subset_reals,subseteq,naturalsplus_positive_real,real_mul_closed,real_mul_pos_pos_gt,reals_power,tuple_power,funs_type_apply,metric_ball_center,bounded_metric_is_metric,standard_metric_is_metric,metric_ball,product_family_reals_power,projection_map_indexed,preimg_iff}.
        $x\in\inters{F}$ by \cref{inters_iff_forall}.
        $x\in B_3$ by \cref{subseteq}.
        Show that for all $y$ such that $y\in B_3$ we have $y\in B_2$.
        \begin{subproof}
            Fix $y$ such that $y\in B_3$.
            $\apply{g}{1}\in F$ by \cref{function_apply_elim,img_iff}.
            $y\in\apply{g}{1}$ by \cref{inters_iff_forall}.
            $1\in J$ by \cref{real_one_in_naturals,real_one_leq_naturalsplus,initial_segment_naturalsplus}.
            We have $1\mathrel{R_g}\apply{g}{1}$ by \cref{subseteq}.
            Take $i_1\in J$ such that
                $\fst{(1,\apply{g}{1})} = i_1$ and
                $\snd{(1,\apply{g}{1})}
                    = \preimg{\piIndex{X}{i_1}}{\metricBall{\boundedMetric{\standardMetricR}}{\reals}{\apply{x}{i_1}}{i_1 \rmul \delta}}$.
            We have $1 = i_1$ and $\apply{g}{1}
                = \preimg{\piIndex{X}{1}}{\metricBall{\boundedMetric{\standardMetricR}}{\reals}{\apply{x}{1}}{1 \rmul \delta}}$
                by \cref{fst_eq,snd_eq,pair_eq_iff}.
            Take $u_0\in\metricBall{\boundedMetric{\standardMetricR}}{\reals}{\apply{x}{1}}{1 \rmul \delta}$ such that
                $y\mathrel{\piIndex{X}{1}} u_0$ by \cref{preimg_iff}.
            $(y,u_0)\in\piIndex{X}{1}$.
            Take $y_2\in\productFamily{X}$ such that $(y,u_0) = (y_2,\apply{y_2}{1})$ by \cref{projection_map_indexed}.
            $y = y_2$ by \cref{pair_eq_iff}.
            $y\in\productFamily{X}$ by \cref{subseteq}.
            $y\in X_0$ by \cref{product_family_reals_power,subseteq}.
            $(x,y)\in X_0\times X_0$ and $(z,x)\in X_0\times X_0$ and $(z,y)\in X_0\times X_0$ by \cref{times_tuple_intro}.
            $d((x,y))\in\reals$ and $d((z,x))\in\reals$ and $d((z,y))\in\reals$ by \cref{metric_on,funs_type_apply}.
            $(x,y)\in\dom{d}$ by \cref{subseteq}.
            Take $s$ such that $((x,y),s)\in d$ by \cref{dom_iff}.
            $d((x,y)) = s$ by \cref{metric_on,funs_elim,subseteq,function_apply_intro}.
            Take $p_0,x_0,y_0,r_0$ such that
                $p_0\in X_0\times X_0$ and $x_0\in X_0$ and $y_0\in X_0$ and $r_0\in\reals$ and
                $((x,y),s) = (p_0,r_0)$ and $p_0 = (x_0,y_0)$ and
                $\realSupremum{r_0}{\omegaMetricSet{x_0}{y_0}}$ by \cref{omega_metric}.
            $p_0 = (x,y)$ and $s = r_0$ by \cref{pair_eq_iff}.
            $x_0 = x$ and $y_0 = y$ by \cref{pair_eq_iff}.
            Show that for all $a\in\omegaMetricSet{x}{y}$ we have $a \leq \delta$.
            \begin{subproof}
                Fix $a$ such that $a\in\omegaMetricSet{x}{y}$.
                Take $i\in\naturalsPlus$ such that $a = \apply{\boundedMetric{\standardMetricR}}{(x(i),y(i))} / i$ by \cref{omega_metric_set}.
                Let $m = \apply{\boundedMetric{\standardMetricR}}{(x(i),y(i))}$.
                $a = m / i$.
                $x(i)\in\reals$ and $y(i)\in\reals$ by \cref{reals_power,tuple_power,funs_type_apply}.
                $(x(i),y(i))\in\reals\times\reals$ by \cref{times_tuple_intro}.
                $\standardMetricR$ is a metric on $\reals$ by \cref{standard_metric_is_metric}.
                $\boundedMetric{\standardMetricR}\in\funs{\reals\times\reals}{\reals}$ by \cref{bounded_metric_is_metric,metric_on,funs_elim}.
                $m\in\reals$ and $0 \leq m$ by \cref{bounded_metric_is_metric,standard_metric_is_metric,metric_on,funs_type_apply,metric_nonnegative}.
                $m \leq 1$ by \cref{bounded_metric_leq_one}.
                \begin{byCase}
                \caseOf{$i\in J$.}
                    $\apply{g}{i}\in F$ by \cref{function_apply_elim,img_iff}.
                    $y\in\apply{g}{i}$ by \cref{inters_iff_forall}.
                    $i\mathrel{R_g}\apply{g}{i}$.
                    Take $i_0\in J$ such that
                        $\fst{(i,\apply{g}{i})} = i_0$ and
                        $\snd{(i,\apply{g}{i})}
                            = \preimg{\piIndex{X}{i_0}}{\metricBall{\boundedMetric{\standardMetricR}}{\reals}{\apply{x}{i_0}}{i_0 \rmul \delta}}$.
                    $i = i_0$ and
                        $\apply{g}{i} = \preimg{\piIndex{X}{i}}{\metricBall{\boundedMetric{\standardMetricR}}{\reals}{\apply{x}{i}}{i \rmul \delta}}$
                        by \cref{fst_eq,snd_eq,pair_eq_iff,function_apply_intro}.
                    $y\in\preimg{\piIndex{X}{i}}{\metricBall{\boundedMetric{\standardMetricR}}{\reals}{\apply{x}{i}}{i \rmul \delta}}$.
                    Take $u\in\metricBall{\boundedMetric{\standardMetricR}}{\reals}{\apply{x}{i}}{i \rmul \delta}$ such that
                        $y\mathrel{\piIndex{X}{i}} u$ by \cref{preimg_iff}.
                    $(y,u)\in\piIndex{X}{i}$.
                    Take $y_1\in\productFamily{X}$ such that $(y,u) = (y_1,\apply{y_1}{i})$ by \cref{projection_map_indexed}.
                    $u = \apply{y}{i}$ by \cref{pair_eq_iff}.
                    $\apply{y}{i}\in\metricBall{\boundedMetric{\standardMetricR}}{\reals}{\apply{x}{i}}{i \rmul \delta}$.
                    $m < i \rmul \delta$ by \cref{metric_ball}.
                    $i\in\reals$ and $0 < i$ by \cref{real_naturals_subset_reals,subseteq,naturalsplus_positive_real}.
                    $i \rmul \delta\in\reals$ by \cref{real_mul_closed}.
                    $i \neq 0$ by \cref{real_pos_nonzero}.
                    $\inv{i}\in\reals$ and $0 < \inv{i}$ by \cref{real_mul_inverse,real_inv_pos}.
                    $m \rmul \inv{i} < (i \rmul \delta) \rmul \inv{i}$ by \cref{real_lt_mul_pos}.
                    $m / i < \delta$ by
                        \cref{real_div,real_lt_subst_left,real_lt_subst_right,real_mul_assoc,real_mul_comm,real_mul_inverse,real_mul_ident}.
                    $a < \delta$ by \cref{real_lt_subst_right,subseteq}.
                    $a \leq \delta$.
                \caseOf{$i\notin J$.}
                    Show that $N < i$.
                    \begin{subproof}
                        Suppose not.
                        $i\in\reals$ and $N\in\reals$ by \cref{real_naturals_subset_reals,subseteq}.
                        \begin{byCase}
                        \caseOf{$i = N$.}
                            $i\in J$ by \cref{initial_segment_naturalsplus}.
                            Contradiction.
                        \caseOf{$i \neq N$.}
                            $i < N$ or $N < i$ by \cref{real_lt_trichotomy}.
                            $i < N$.
                            $i \leq N$.
                            $i\in J$ by \cref{initial_segment_naturalsplus}.
                            Contradiction.
                        \end{byCase}
                    \end{subproof}
                    $1 / i \leq 1 / N$ by \cref{naturalsplus_reciprocal_antitone}.
                    $a\in\reals$ by \cref{real_div_naturalsplus_real}.
                    $1\in\reals$ by \cref{real_naturals_subset_reals,real_one_in_naturals,subseteq}.
                    $1 / i\in\reals$ by \cref{real_div_naturalsplus_real}.
                    $1 / N\in\reals$ by \cref{real_div_naturalsplus_real}.
                    $a \leq 1 / i$ by \cref{real_div_leq_same_denom}.
                    $1 / i < \delta$ by \cref{real_leq_lt_trans,real_lt_subst_left}.
                    $a < \delta$ by \cref{real_leq_lt_trans}.
                    $a \leq \delta$.
                \end{byCase}
            \end{subproof}
            For all $a\in\omegaMetricSet{x}{y}$ we have $a\in\reals$
                by \cref{omega_metric_set,reals_power,tuple_power,funs_type_apply,times_tuple_intro,standard_metric_is_metric,bounded_metric_is_metric,metric_on,real_div_naturalsplus_real}.
            $\omegaMetricSet{x}{y}\subseteq\reals$ by \cref{subseteq}.
            $\delta\in\reals$ by \cref{subseteq}.
            $\delta$ is an upper bound of $\omegaMetricSet{x}{y}$ by \cref{real_upper_bound}.
            $r_0 \leq \delta$ by \cref{real_supremum_pred,real_supremum,real_upper_bound}.
            $d((x,y)) \leq \delta$ by \cref{real_lt_subst_left,real_lt_subst_right}.
            $d((z,y)) \leq d((z,x)) + d((x,y))$ by \cref{metric_on,metric_triangle_inequality}.
            $d((x,y)) + d((z,x))\in\reals$ and $\delta + d((z,x))\in\reals$ and $r + d((z,x))\in\reals$
                by \cref{real_add_closed}.
            $d((z,x)) + d((x,y))\in\reals$ and $d((z,x)) + \delta\in\reals$ by \cref{real_add_closed}.
            $d((x,y)) + d((z,x)) \leq \delta + d((z,x))$ by \cref{real_leq_add}.
            $d((z,x)) + d((x,y)) \leq \delta + d((z,x))$ by \cref{real_add_comm,real_leq_trans}.
            $\delta + d((z,x)) = d((z,x)) + \delta$ by \cref{real_add_comm}.
            $d((z,y)) \leq d((z,x)) + \delta$ by \cref{real_leq_trans,subseteq}.
            $\delta + d((z,x)) < r + d((z,x))$ by \cref{real_lt_add}.
            $d((z,x)) + \delta < r + d((z,x))$ by \cref{real_add_comm,real_lt_subst_left}.
            $d((z,x)) + \delta < d((z,x)) + r$ by \cref{real_add_comm,real_lt_subst_right}.
            $d((z,x)) + r = \epsilon$ by
                \cref{real_sub,real_add_assoc,real_add_inverse,real_add_comm,real_add_ident}.
            $d((z,x)) + \delta < \epsilon$ by \cref{real_lt_subst_right}.
            $d((z,y)) < \epsilon$ by \cref{real_leq_lt_trans}.
            $y\in\metricBall{d}{X_0}{z}{\epsilon}$ by \cref{metric_ball}.
            $y\in B_2$ by \cref{subseteq}.
        \end{subproof}
        $B_3\subseteq B_2$ by \cref{subseteq}.
    \end{subproof}
    Show that $T_0$ is finer than $T_1$ on $X_0$.
    \begin{subproof}
        Follows by \cref{basis_finer_criterion}.
    \end{subproof}
    Show that $T_1\subseteq T_0$.
    \begin{subproof}
        Follows by \cref{finer_topology}.
    \end{subproof}
    Show that $T_0\subseteq T_1$.
    \begin{subproof}
        Follows by \cref{finer_topology}.
    \end{subproof}
    $T_1 = T_0$ by \cref{subseteq_antisymmetric}.
    $\metricTopology{\omegaMetric}{\realsPower{\naturalsPlus}} = T_0$.
\end{proof}
 
\section{§21 The Metric Topology (continued)}

% from §21 helper lemma 1: positive reciprocal of naturalsplus
\begin{lemma}\label{positive_reciprocal_naturalsplus}
    Suppose $n\in\naturalsPlus$.
    Then $0 < 1 / n$.
\end{lemma}
\begin{proof}
    Follows by \cref{real_zero_lt_one,real_add_ident,real_naturals_subset_reals,real_one_in_naturals,subseteq,real_one_leq_naturalsplus,real_lt_subst_right,real_lt_trans,real_inv_pos,real_pos_nonzero,real_mul_inverse,real_div,real_mul_ident}.
\end{proof}

% from §21 helper lemma 2: metric ball monotonic in radius
\begin{lemma}\label{metric_ball_mono}
    Assume $d$ is a metric on $X$.
    Suppose $x\in X$.
    Suppose $r_1\in\reals$.
    Suppose $r_2\in\reals$.
    Suppose $r_1 < r_2$.
    Then $\metricBall{d}{X}{x}{r_1} \subseteq \metricBall{d}{X}{x}{r_2}$.
\end{lemma}
\begin{proof}
    Follows by \cref{metric_ball,metric_on,times_tuple_intro,funs_type_apply,real_lt_trans,subseteq}.
\end{proof}

% from §21 helper lemma 3: epsilon-delta characterization of continuity at a point
\begin{lemma}\label{metric_continuity_at_eps_delta}
    Suppose $d_X$ is a metric on $X$.
    Suppose $d_Y$ is a metric on $Y$.
    Let $T_X = \metricTopology{d_X}{X}$.
    Let $T_Y = \metricTopology{d_Y}{Y}$.
    Suppose $f$ is a function from $X$ to $Y$.
    Suppose $x\in X$.
    Then $f$ is continuous at $x$ from $X$ to $Y$ with respect to $T_X$ and $T_Y$ iff
        for all $\epsilon\in\reals$ such that $0 < \epsilon$
        there exists $\delta\in\reals$ such that $0 < \delta$ and
        for all $y\in X$ if $d_X((x,y)) < \delta$ then $d_Y((f(x),f(y))) < \epsilon$.
\end{lemma}
\begin{proof}
    Show that if $f$ is continuous at $x$ from $X$ to $Y$ with respect to $T_X$ and $T_Y$ then
        for all $\epsilon\in\reals$ such that $0 < \epsilon$
        there exists $\delta\in\reals$ such that $0 < \delta$ and
        for all $y\in X$ if $d_X((x,y)) < \delta$ then $d_Y((f(x),f(y))) < \epsilon$.
    \begin{subproof}
        Assume $f$ is continuous at $x$ from $X$ to $Y$ with respect to $T_X$ and $T_Y$.
        $f$ is a function by \cref{continuous_at,funs_is_function}.
        Fix $\epsilon\in\reals$.
        Assume $0 < \epsilon$.
        Let $V = \metricBall{d_Y}{Y}{f(x)}{\epsilon}$.
        $f(x)\in Y$ by \cref{continuous_at,funs_type_apply}.
        $V$ is open in $Y$ with respect to $T_Y$ by
            \cref{continuous_at,funs_type_apply,metric_ball,subseteq,powerset_intro,metric_basis,metric_basis_is_basis,basis_elem_open_in_generated,metric_topology,basis_for_topology,basis_topology_is_topology,open_in}.
        $d_Y((f(x),f(x))) = 0$ by \cref{metric_zero_characterization,metric_on,continuous_at,funs_type_apply}.
        $f(x)\in V$ by \cref{real_lt_subst_left,metric_ball}.
        $V$ is a neighborhood of $f(x)$ in $Y$ with respect to $T_Y$ by \cref{neighborhood}.
        Take $U$ such that
            $U$ is a neighborhood of $x$ in $X$ with respect to $T_X$ and
            $\img{f}{U}\subseteq V$ by \cref{continuous_at}.
        $U$ is open in $X$ with respect to $T_X$ and $x\in U$ by \cref{neighborhood}.
        $U\subseteq X$ by
            \cref{open_in,basis_topology_is_topology,metric_topology,metric_basis_is_basis,basis_for_topology,topology_on,subseteq,pow_iff}.
        Take $\delta\in\reals$ such that $0 < \delta$ and
            $\metricBall{d_X}{X}{x}{\delta} \subseteq U$ by \cref{metric_open_iff}.
        For all $y\in X$ if $d_X((x,y)) < \delta$ then $d_Y((f(x),f(y))) < \epsilon$
            by \cref{metric_ball,subseteq,continuous_at,funs,inter_intro,img_of_function_intro}.
    \end{subproof}
    Show that if
        for all $\epsilon\in\reals$ such that $0 < \epsilon$
        there exists $\delta\in\reals$ such that $0 < \delta$ and
        for all $y\in X$ if $d_X((x,y)) < \delta$ then $d_Y((f(x),f(y))) < \epsilon$
        then $f$ is continuous at $x$ from $X$ to $Y$ with respect to $T_X$ and $T_Y$.
    \begin{subproof}
        Assume for all $\epsilon\in\reals$ such that $0 < \epsilon$
            there exists $\delta\in\reals$ such that $0 < \delta$ and
            for all $y\in X$ if $d_X((x,y)) < \delta$ then $d_Y((f(x),f(y))) < \epsilon$.
        $f$ is a function by \cref{funs_is_function}.
        Show that for all $V$ such that $V$ is a neighborhood of $f(x)$ in $Y$ with respect to $T_Y$
            there exists $U$ such that
            $U$ is a neighborhood of $x$ in $X$ with respect to $T_X$ and
            $\img{f}{U}\subseteq V$.
        \begin{subproof}
            Fix $V$ such that $V$ is a neighborhood of $f(x)$ in $Y$ with respect to $T_Y$.
            $V$ is open in $Y$ with respect to $T_Y$ and $f(x)\in V$ by \cref{neighborhood}.
            $V\subseteq Y$ by
                \cref{open_in,basis_topology_is_topology,metric_topology,metric_basis_is_basis,basis_for_topology,topology_on,subseteq,pow_iff}.
            Take $\epsilon\in\reals$ such that $0 < \epsilon$ and
                $\metricBall{d_Y}{Y}{f(x)}{\epsilon} \subseteq V$ by \cref{metric_open_iff}.
            Take $\delta\in\reals$ such that $0 < \delta$ and
                for all $y\in X$ if $d_X((x,y)) < \delta$ then $d_Y((f(x),f(y))) < \epsilon$.
            Let $U = \metricBall{d_X}{X}{x}{\delta}$.
            $U$ is open in $X$ with respect to $T_X$ by
                \cref{metric_ball,subseteq,powerset_intro,metric_basis,metric_basis_is_basis,basis_elem_open_in_generated,basis_for_topology,metric_topology,basis_topology_is_topology,open_in}.
            $d_X((x,x)) = 0$ by \cref{metric_zero_characterization,metric_on}.
            $x\in U$ by \cref{real_lt_subst_left,metric_ball}.
            $U$ is a neighborhood of $x$ in $X$ with respect to $T_X$ by \cref{neighborhood}.
            For all $z$ such that $z\in\img{f}{U}$ we have $z\in V$
                by \cref{img_of_function_elim,inter_elim_right,metric_ball,subseteq,funs_type_apply}.
            $\img{f}{U}\subseteq V$ by \cref{subseteq}.
            There exists $W$ such that
                $W$ is a neighborhood of $x$ in $X$ with respect to $T_X$ and
                $\img{f}{W}\subseteq V$.
        \end{subproof}
        $T_X$ is a topology on $X$ and $T_Y$ is a topology on $Y$ by
            \cref{basis_topology_is_topology,metric_topology,metric_basis_is_basis,basis_for_topology}.
        $f$ is continuous at $x$ from $X$ to $Y$ with respect to $T_X$ and $T_Y$
            by \cref{continuous_at}.
\end{subproof}
\end{proof}

% from §21 helper lemma 4: continuity implies epsilon-delta
\begin{lemma}\label{metric_continuity_eps_delta_forward}
    Suppose $d_X$ is a metric on $X$.
    Suppose $d_Y$ is a metric on $Y$.
    Let $T_X = \metricTopology{d_X}{X}$.
    Let $T_Y = \metricTopology{d_Y}{Y}$.
    Suppose $f$ is a function from $X$ to $Y$.
    Suppose $f$ is continuous from $X$ to $Y$ with respect to $T_X$ and $T_Y$.
    Then for all $x\in X$ we have
        for all $\epsilon\in\reals$ such that $0 < \epsilon$
        there exists $\delta\in\reals$ such that $0 < \delta$ and
        for all $y\in X$ if $d_X((x,y)) < \delta$ then $d_Y((f(x),f(y))) < \epsilon$.
\end{lemma}
\begin{proof}
    Show that for all $x\in X$ we have
        for all $\epsilon\in\reals$ such that $0 < \epsilon$
        there exists $\delta\in\reals$ such that $0 < \delta$ and
        for all $y\in X$ if $d_X((x,y)) < \delta$ then $d_Y((f(x),f(y))) < \epsilon$.
    \begin{subproof}
        Follows by \cref{continuous_implies_continuous_at,metric_continuity_at_eps_delta}.
    \end{subproof}
\end{proof}

% from §21 helper lemma 5: epsilon-delta implies continuity
\begin{lemma}\label{metric_continuity_eps_delta_backward}
    Suppose $d_X$ is a metric on $X$.
    Suppose $d_Y$ is a metric on $Y$.
    Let $T_X = \metricTopology{d_X}{X}$.
    Let $T_Y = \metricTopology{d_Y}{Y}$.
    Suppose $f$ is a function from $X$ to $Y$.
    Suppose for all $x\in X$ we have
        for all $\epsilon\in\reals$ such that $0 < \epsilon$
        there exists $\delta\in\reals$ such that $0 < \delta$ and
        for all $y\in X$ if $d_X((x,y)) < \delta$ then $d_Y((f(x),f(y))) < \epsilon$.
    Then $f$ is continuous from $X$ to $Y$ with respect to $T_X$ and $T_Y$.
\end{lemma}
\begin{proof}
    Show that for all $x\in X$ we have
        $f$ is continuous at $x$ from $X$ to $Y$ with respect to $T_X$ and $T_Y$.
    \begin{subproof}
        Follows by \cref{metric_continuity_at_eps_delta}.
    \end{subproof}
    $f$ is continuous from $X$ to $Y$ with respect to $T_X$ and $T_Y$ by
        \cref{continuous_at_implies_continuous,basis_topology_is_topology,metric_topology,metric_basis_is_basis,basis_for_topology}.
\end{proof}

% from §21 Item 1: epsilon-delta characterization of continuity in metric spaces
\begin{theorem}\label{metric_continuity_eps_delta}
    Suppose $d_X$ is a metric on $X$.
    Suppose $d_Y$ is a metric on $Y$.
    Let $T_X = \metricTopology{d_X}{X}$.
    Let $T_Y = \metricTopology{d_Y}{Y}$.
    Suppose $f$ is a function from $X$ to $Y$.
    Then $f$ is continuous from $X$ to $Y$ with respect to $T_X$ and $T_Y$ iff
        for all $x\in X$ we have
        for all $\epsilon\in\reals$ such that $0 < \epsilon$
        there exists $\delta\in\reals$ such that $0 < \delta$ and
        for all $y\in X$ if $d_X((x,y)) < \delta$ then $d_Y((f(x),f(y))) < \epsilon$.
\end{theorem}
\begin{proof}
    Show that if $f$ is continuous from $X$ to $Y$ with respect to $T_X$ and $T_Y$ then
        for all $x\in X$ we have
        for all $\epsilon\in\reals$ such that $0 < \epsilon$
        there exists $\delta\in\reals$ such that $0 < \delta$ and
        for all $y\in X$ if $d_X((x,y)) < \delta$ then $d_Y((f(x),f(y))) < \epsilon$.
    \begin{subproof}
        Follows by \cref{metric_continuity_eps_delta_forward}.
    \end{subproof}
    Show that if
        for all $x\in X$ we have
        for all $\epsilon\in\reals$ such that $0 < \epsilon$
        there exists $\delta\in\reals$ such that $0 < \delta$ and
        for all $y\in X$ if $d_X((x,y)) < \delta$ then $d_Y((f(x),f(y))) < \epsilon$
        then $f$ is continuous from $X$ to $Y$ with respect to $T_X$ and $T_Y$.
    \begin{subproof}
        Follows by \cref{metric_continuity_eps_delta_backward}.
    \end{subproof}
\end{proof}

% from §21 Item 2: sequence implies closure
\begin{lemma}\label{sequence_implies_closure}
    Assume $T$ is a topology on $X$.
    Suppose $A\subseteq X$.
    Suppose $s$ is a sequence in $A$.
    Suppose $x\in X$.
    Suppose $s$ converges to $x$ in $X$ with respect to $T$.
    Then $x\in\closure{A}{X}{T}$.
\end{lemma}
\begin{proof}
    Show that for all $U$ such that $U$ is open in $X$ with respect to $T$ and $x\in U$ we have $U$ intersects $A$.
    \begin{subproof}
        Follows by \cref{converges_to,neighborhood,real_naturals_closed_succ,real_naturals_subset_reals,subseteq,real_add_ident,real_one_in_naturals,real_zero_lt_one,real_lt_add,real_add_comm,real_add_closed,real_lt_subst_left,real_lt_imp_leq,sequence_in,funs_type_apply,inter_intro,emptyset,intersects}.
    \end{subproof}
    $x\in\closure{A}{X}{T}$ by \cref{closure_iff_intersects_open}.
\end{proof}

% from §21 Item 3: closure implies sequence in a metrizable space
\begin{lemma}\label{closure_implies_sequence_metrizable}
    Assume $T$ is a topology on $X$.
    Suppose $A\subseteq X$.
    Suppose $x\in X$.
    Suppose $X$ is metrizable with respect to $T$.
    Suppose $x\in\closure{A}{X}{T}$.
    Then there exists $s$ such that $s$ is a sequence in $A$ and $s$ converges to $x$ in $X$ with respect to $T$.
\end{lemma}
\begin{proof}
    Take $d$ such that $d$ is a metric on $X$ and $T = \metricTopology{d}{X}$ by \cref{metrizable}.
    Let $b = \{(k,\metricBall{d}{X}{x}{1 / k}) \mid k\in\naturalsPlus\}$.
    For all $k,U_1,U_2$ such that $(k,U_1)\in b$ and $(k,U_2)\in b$ we have $U_1 = U_2$ by \cref{pair_eq_iff}.
    $b$ is a function by \cref{rightunique,relation}.
    $\dom{b} = \naturalsPlus$ by \cref{pair_eq_iff,dom_intro,dom_iff,subseteq,subseteq_antisymmetric}.
    $b$ is a function on $\naturalsPlus$.

    Let $F = \{(n,\img{\restrl{b}{\initialSegment{n}}}{\initialSegment{n}}) \mid n\in\naturalsPlus\}$.
    Let $B = \{(n,\inters{\apply{F}{n}}) \mid n\in\naturalsPlus\}$.

    For all $n,U_1,U_2$ such that $(n,U_1)\in F$ and $(n,U_2)\in F$ we have $U_1 = U_2$
        by \cref{pair_eq_iff}.
    $F$ is a function by \cref{rightunique,relation}.
    $\dom{F} = \naturalsPlus$ by \cref{pair_eq_iff,dom_intro,dom_iff,subseteq,subseteq_antisymmetric}.
    $F$ is a function on $\naturalsPlus$.

    For all $n\in\naturalsPlus$ we have $F(n) = \img{\restrl{b}{\initialSegment{n}}}{\initialSegment{n}}$
        by \cref{pair_eq_iff,function_apply_intro}.

    For all $n,U_1,U_2$ such that $(n,U_1)\in B$ and $(n,U_2)\in B$ we have $U_1 = U_2$
        by \cref{pair_eq_iff}.
    $B$ is a function by \cref{rightunique,relation}.
    $\dom{B} = \naturalsPlus$ by \cref{pair_eq_iff,dom_intro,dom_iff,subseteq,subseteq_antisymmetric}.
    $B$ is a function on $\naturalsPlus$.

    For all $n\in\naturalsPlus$ we have $B(n) = \inters{F(n)}$ by \cref{pair_eq_iff,function_apply_intro}.

    Show that for all $n\in\naturalsPlus$ we have $B(n)$ is a neighborhood of $x$ in $X$ with respect to $T$.
    \begin{subproof}
        Fix $n\in\naturalsPlus$.
        $\initialSegment{n}$ is finite by \cref{initial_segment_finite}.
        $\initialSegment{n}\subseteq\dom{b}$
            by \cref{initial_segment_naturalsplus,pair_eq_iff,dom_intro,subseteq}.
        Let $g = \restrl{b}{\initialSegment{n}}$.
        $g$ is a function by \cref{restrl_of_function_is_function}.
        $\dom{g} = \initialSegment{n}$ by
            \cref{restrl_dom,initial_segment_naturalsplus,pair_eq_iff,restrl_iff,dom_intro,subseteq,subseteq_antisymmetric}.
        $g$ is a function on $\initialSegment{n}$.
        $F(n) = \img{g}{\initialSegment{n}}$.
        $F(n)$ is finite by \cref{finite_image_function}.
        $1\in\initialSegment{n}$ by \cref{real_one_in_naturals,real_one_leq_naturalsplus,initial_segment_naturalsplus}.
        Let $U_1 = \metricBall{d}{X}{x}{1 / 1}$.
        $(1,U_1)\in b$ by \cref{initial_segment_naturalsplus,pair_eq_iff}.
        $1\in\dom{g}\inter\initialSegment{n}$ by \cref{subseteq,restrl_iff,dom_intro,inter_intro}.
        $b(1)\in F(n)$ by \cref{restrl_of_function_apply,img_of_function_intro}.
        $F(n)$ is inhabited.
        Show that for all $U$ such that $U\in F(n)$ we have $U\in T$.
        \begin{subproof}
            Fix $U$ such that $U\in F(n)$.
            $U\in\img{g}{\initialSegment{n}}$.
            Take $k\in\initialSegment{n}$ such that $(k,U)\in g$ by \cref{img_iff}.
            $(k,U)\in b$ by \cref{restrl_iff}.
            $U = \metricBall{d}{X}{x}{1 / k}$ by \cref{pair_eq_iff}.
            $k\in\naturalsPlus$ by \cref{initial_segment_naturalsplus}.
            $0 < 1 / k$ by \cref{positive_reciprocal_naturalsplus}.
            $0\in\reals$ and $1\in\reals$ and $k\in\reals$
                by \cref{real_add_ident,real_one_in_naturals,real_naturals_subset_reals,subseteq}.
            $0 < k$ by \cref{real_zero_lt_one,real_one_leq_naturalsplus,real_lt_subst_right,real_lt_trans}.
            $k \neq 0$ by \cref{real_pos_nonzero}.
            $1 / k\in\reals$ by \cref{real_mul_inverse,real_div,real_mul_closed}.
            For all $y$ such that $y\in U$ we have $y\in X$ by \cref{metric_ball,subseteq}.
            $U\subseteq X$ by \cref{subseteq}.
            $U\in\pow{X}$ by \cref{powerset_intro}.
            $U\in\metricBasis{d}{X}$ by \cref{metric_basis}.
            $U\in T$ by \cref{metric_basis_is_basis,basis_elem_open_in_generated,metric_topology,basis_for_topology}.
        \end{subproof}
        $F(n)\subseteq T$ by \cref{subseteq}.
        $B(n)\in T$ by \cref{finite_intersections_open_in_topology}.
        Show that for all $U$ such that $U\in F(n)$ we have $x\in U$.
        \begin{subproof}
            Fix $U$ such that $U\in F(n)$.
            $U\in\img{g}{\initialSegment{n}}$.
            Take $k\in\initialSegment{n}$ such that $(k,U)\in g$ by \cref{img_iff}.
            $(k,U)\in b$ by \cref{restrl_iff}.
            $U = \metricBall{d}{X}{x}{1 / k}$ by \cref{pair_eq_iff}.
            $k\in\naturalsPlus$ by \cref{initial_segment_naturalsplus}.
            $0 < 1 / k$ by \cref{positive_reciprocal_naturalsplus}.
            $d((x,x)) < 1 / k$ by \cref{metric_zero_characterization,metric_on,real_lt_subst_left}.
            $x\in U$ by \cref{metric_ball}.
        \end{subproof}
        $B(n)$ is a neighborhood of $x$ in $X$ with respect to $T$ by \cref{inters_iff_forall,neighborhood,open_in}.
    \end{subproof}

    Show that for all $n\in\naturalsPlus$ we have $B(n)$ intersects $A$.
    \begin{subproof}
        Follows by \cref{closure_iff_intersects_open,neighborhood}.
    \end{subproof}

    Let $R = \{w\in\naturalsPlus\times X \mid \snd{w}\in A\inter \apply{B}{\fst{w}}\}$.
    $R$ is a relation by \cref{times_elem_is_tuple,relation}.
    Show that for all $n\in\naturalsPlus$ there exists $y$ such that $n\mathrel{R} y$.
    \begin{subproof}
        Follows by \cref{intersects,inter_comm,nonempty_inhabited,inter_elim_left,inter_elim_right,subseteq,times_tuple_intro,fst_eq,snd_eq,inter_intro,pair_eq_iff}.
    \end{subproof}
    Take $s$ such that $s$ is a function on $\naturalsPlus$ and
        for all $n\in\naturalsPlus$ we have $n\mathrel{R} s(n)$ by \cref{axiom_of_choice}.

    $s$ is a sequence in $A$ by \cref{funs_intro,sequence_in,subseteq,pair_eq_iff,fst_eq,snd_eq,inter}.

    $s$ is a sequence in $X$ by \cref{sequence_in,funs_weaken_codom}.
    Show that for all $U$ such that $U$ is a neighborhood of $x$ in $X$ with respect to $T$
        there exists $N\in\naturalsPlus$ such that for all $n\in\naturalsPlus$ if $N \leq n$ then $s(n)\in U$.
    \begin{subproof}
        Fix $U$ such that $U$ is a neighborhood of $x$ in $X$ with respect to $T$.
        $U$ is open in $X$ with respect to $T$ and $x\in U$ by \cref{neighborhood}.
        $U\in T$ and $U\subseteq X$ by \cref{open_in,topology_on,subseteq,pow_iff}.
        Take $\epsilon\in\reals$ such that $0 < \epsilon$ and
            $\metricBall{d}{X}{x}{\epsilon} \subseteq U$ by \cref{metric_open_iff}.
        Take $N\in\naturalsPlus$ such that $0 < 1 / N$ and $1 / N < \epsilon$ by \cref{real_small_reciprocal}.
        $\metricBall{d}{X}{x}{1 / N} \subseteq U$ by
            \cref{real_naturals_subset_reals,real_one_in_naturals,subseteq,real_add_ident,real_zero_lt_one,real_one_leq_naturalsplus,real_lt_subst_right,real_lt_trans,real_pos_nonzero,real_mul_inverse,real_div,real_mul_closed,metric_ball_mono,subseteq_transitive}.
        Show that for all $n\in\naturalsPlus$ if $N \leq n$ then $s(n)\in U$.
        \begin{subproof}
            Fix $n\in\naturalsPlus$.
            Assume $N \leq n$.
            $\initialSegment{N}\subseteq\initialSegment{n}$ by
                \cref{initial_segment_naturalsplus,real_naturals_subset_reals,subseteq,real_leq_trans}.
            Show that for all $U_0$ such that $U_0\in F(N)$ we have $U_0\in F(n)$.
            \begin{subproof}
                Fix $U_0$ such that $U_0\in F(N)$.
                $U_0\in\img{\restrl{b}{\initialSegment{N}}}{\initialSegment{N}}$.
                Take $k\in\initialSegment{N}$ such that $(k,U_0)\in\restrl{b}{\initialSegment{N}}$ by \cref{img_iff}.
                $k\in\initialSegment{n}$ by \cref{subseteq}.
                $(k,U_0)\in b$ by \cref{restrl_iff}.
                $(k,U_0)\in\restrl{b}{\initialSegment{n}}$ by \cref{restrl_iff}.
                $U_0\in\img{\restrl{b}{\initialSegment{n}}}{\initialSegment{n}}$ by \cref{img_iff}.
                $U_0\in F(n)$.
            \end{subproof}
            $F(N)\subseteq F(n)$ by \cref{subseteq}.
            $F(N) = \img{\restrl{b}{\initialSegment{N}}}{\initialSegment{N}}$.
            $N\in\initialSegment{N}$ by \cref{initial_segment_naturalsplus}.
            For all $i$ such that $i\in\initialSegment{N}$ we have $i\in\dom{b}$ by
                \cref{initial_segment_naturalsplus,pair_eq_iff,dom_intro}.
            $\initialSegment{N}\subseteq\dom{b}$ by \cref{subseteq}.
            $(N,b(N))\in b$ by \cref{subseteq,function_apply_elim}.
            $b(N)\in F(N)$ by \cref{restrl_iff,img_iff}.
            $F(N)$ is inhabited.
            For all $U_0$ such that $U_0\in F(N)$ we have $B(n)\subseteq U_0$ by
                \cref{subseteq,inters_subseteq_elem}.
            $B(n)\subseteq B(N)$ by \cref{subseteq_inters_iff}.
            $B(N)\subseteq \metricBall{d}{X}{x}{1 / N}$ by
                \cref{inters_subseteq_elem,pair_eq_iff}.
            $B(n)\subseteq \metricBall{d}{X}{x}{1 / N}$ by \cref{subseteq_transitive}.
            $s(n)\in U$ by \cref{pair_eq_iff,fst_eq,snd_eq,inter,subseteq}.
        \end{subproof}
        There exists $M\in\naturalsPlus$ such that for all $n\in\naturalsPlus$ if $M \leq n$ then $s(n)\in U$.
    \end{subproof}
    $s$ converges to $x$ in $X$ with respect to $T$ by \cref{converges_to}.

    There exists $t$ such that $t$ is a sequence in $A$ and $t$ converges to $x$ in $X$ with respect to $T$.
\end{proof}

% from §21 Item 4: continuous implies sequential continuity
\begin{theorem}\label{continuous_implies_sequential}
    Assume $T$ is a topology on $X$.
    Assume $T_1$ is a topology on $Y$.
    Suppose $f$ is a function from $X$ to $Y$.
    Suppose $f$ is continuous from $X$ to $Y$ with respect to $T$ and $T_1$.
    Then for all $s$ such that $s$ is a sequence in $X$ we have
        for all $x\in X$ if $s$ converges to $x$ in $X$ with respect to $T$ then
        $f\circ s$ converges to $f(x)$ in $Y$ with respect to $T_1$.
\end{theorem}
\begin{proof}
    Fix $s$ such that $s$ is a sequence in $X$.
    Fix $x\in X$.
    Assume $s$ converges to $x$ in $X$ with respect to $T$.
    $f\circ s$ is a sequence in $Y$ by \cref{continuous,funs_circ,sequence_in}.
    Show that for all $V$ such that $V$ is a neighborhood of $f(x)$ in $Y$ with respect to $T_1$
        there exists $N\in\naturalsPlus$ such that for all $n\in\naturalsPlus$ if $N\leq n$ then
        $(f\circ s)(n)\in V$.
    \begin{subproof}
        Follows by \cref{neighborhood,continuous,funs,funs_is_function,function_apply_iff,preimg_iff,converges_to,sequence_in,funs_elim,function_to_ran,subseteq,circ_apply}.
    \end{subproof}
    $f\circ s$ converges to $f(x)$ in $Y$ with respect to $T_1$ by \cref{converges_to}.
\end{proof}
% from §21 Item 5: sequential continuity implies continuity in metrizable spaces
\begin{theorem}\label{sequential_implies_continuous_metrizable}
    Assume $T$ is a topology on $X$.
    Assume $T_1$ is a topology on $Y$.
    Suppose $f$ is a function from $X$ to $Y$.
    Suppose $X$ is metrizable with respect to $T$.
    Suppose for all $s$ such that $s$ is a sequence in $X$ we have
        for all $x\in X$ if $s$ converges to $x$ in $X$ with respect to $T$ then
        $f\circ s$ converges to $f(x)$ in $Y$ with respect to $T_1$.
    Then $f$ is continuous from $X$ to $Y$ with respect to $T$ and $T_1$.
\end{theorem}
\begin{proof}
    Show that for all $A$ such that $A\subseteq X$ we have
        $\img{f}{\closure{A}{X}{T}}\subseteq \closure{\img{f}{A}}{Y}{T_1}$.
    \begin{subproof}
        Fix $A$ such that $A\subseteq X$.
        Show that for all $y$ such that $y\in\img{f}{\closure{A}{X}{T}}$ we have
            $y\in\closure{\img{f}{A}}{Y}{T_1}$.
        \begin{subproof}
            Fix $y$ such that $y\in\img{f}{\closure{A}{X}{T}}$.
            Take $x\in\dom{f}\inter\closure{A}{X}{T}$ such that $y = f(x)$ by \cref{funs_is_function,img_of_function_elim}.
            Take $s$ such that $s$ is a sequence in $A$ and
                $s$ converges to $x$ in $X$ with respect to $T$
                by \cref{inter_elim_right,closure,closure_implies_sequence_metrizable}.
            $x\in X$ and $s$ is a sequence in $X$ and
                $f\circ s$ converges to $f(x)$ in $Y$ with respect to $T_1$ by
                \cref{inter_elim_right,closure,sequence_in,funs_weaken_codom}.
            $f\circ s\in\funs{\naturalsPlus}{Y}$ by \cref{funs_circ,sequence_in}.
            For all $n\in\dom{(f\circ s)}$ we have $(f\circ s)(n)\in\img{f}{A}$ by
                \cref{funs,sequence_in,funs_elim,funs_type_apply,function_to_ran,subseteq,subseteq_transitive,inter_intro,img_of_function_intro,circ_apply}.
            $f\circ s$ is a sequence in $\img{f}{A}$ by \cref{funs,sequence_in,funs_elim,funs_intro}.
        $f(x)\in\closure{\img{f}{A}}{Y}{T_1}$ by
            \cref{img_subseteq_ran,funs_elim,function_to_ran,subseteq_transitive,funs_type_apply,sequence_implies_closure}.
        $y\in\closure{\img{f}{A}}{Y}{T_1}$.
        \end{subproof}
        $\img{f}{\closure{A}{X}{T}}\subseteq \closure{\img{f}{A}}{Y}{T_1}$ by \cref{subseteq}.
    \end{subproof}
    Show that for all $B$ such that $B$ is closed in $Y$ with respect to $T_1$
        we have $\preimg{f}{B}$ is closed in $X$ with respect to $T$.
    \begin{subproof}
        Follows by \cref{closure_image_subset_implies_preimg_closed}.
    \end{subproof}
    Show that $f$ is continuous from $X$ to $Y$ with respect to $T$ and $T_1$.
    \begin{subproof}
        Follows by \cref{preimg_closed_implies_continuous}.
    \end{subproof}
\end{proof}

% from §21 Item 6: countable basis at a point
\begin{definition}\label{countable_basis_at_point}
    $b$ is a countable basis at $x$ in $X$ with respect to $T$ iff
        $b$ is a sequence in $\pow{X}$ and
        for all $n\in\naturalsPlus$ we have
            $b(n)$ is a neighborhood of $x$ in $X$ with respect to $T$ and
        for all $U$ such that $U$ is a neighborhood of $x$ in $X$ with respect to $T$
            there exists $n\in\naturalsPlus$ such that $b(n)\subseteq U$.
\end{definition}

% from §21 Item 7: first countability axiom
\begin{definition}\label{first_countability_axiom}
    $X$ satisfies the first countability axiom with respect to $T$ iff
        for all $x\in X$ there exists $b$ such that
            $b$ is a countable basis at $x$ in $X$ with respect to $T$.
\end{definition}

% from §21 helper definition 1: addition operation on reals
\begin{definition}\label{real_addition_operation}
    $\realAddition = \{w \mid \exists p\in\reals\times\reals. \exists z\in\reals.
        w = (p,z) \land z = \fst{p} + \snd{p}\}$.
\end{definition}

% from §21 helper definition 2: subtraction operation on reals
\begin{definition}\label{real_subtraction_operation}
    $\realSubtraction = \{w \mid \exists p\in\reals\times\reals. \exists z\in\reals.
        w = (p,z) \land z = \fst{p} - \snd{p}\}$.
\end{definition}

% from §21 helper definition 3: multiplication operation on reals
\begin{definition}\label{real_multiplication_operation}
    $\realMultiplication = \{w \mid \exists p\in\reals\times\reals. \exists z\in\reals.
        w = (p,z) \land z = \fst{p} \rmul \snd{p}\}$.
\end{definition}

% from §21 helper definition 4: quotient operation on reals
\begin{definition}\label{real_division_operation}
    $\realDivision = \{w \mid \exists p\in\reals\times(\reals\setminus\{0\}). \exists z\in\reals.
        w = (p,z) \land z = \fst{p} \rmul \inv{\snd{p}}\}$.
\end{definition}

% from §21 helper lemma 3: standard metric value
\begin{lemma}\label{standard_metric_value}
    Suppose $x\in\reals$.
    Suppose $y\in\reals$.
    Then $\apply{\standardMetricR}{(x,y)} = \abs{x - y}$.
\end{lemma}
\begin{proof}
    $\apply{\standardMetricR}{(x,y)} = \abs{x - y}$ by
        \cref{standard_metric_is_metric,metric_on,times_tuple_intro,fst_eq,snd_eq,real_add_inverse,real_sub,real_add_closed,real_abs_in_reals,standard_metric_r,funs_elim,function_apply_intro}.
\end{proof}

% from §21 helper lemma 4: halving a positive real
\begin{lemma}\label{real_half_of_positive}
    Suppose $\epsilon\in\reals$.
    Suppose $0 < \epsilon$.
    Suppose $\delta = \epsilon / (1 + 1)$.
    Then $\delta\in\reals$ and $0 < \delta$ and $\delta + \delta = \epsilon$.
\end{lemma}
\begin{proof}
    $1\in\reals$ and $0\in\reals$ and $0 < 1$
        by \cref{real_mul_ident,real_add_ident,real_zero_lt_one}.
    $0 < 1 + 1$ and $1 + 1\in\reals$ by
        \cref{real_lt_add,real_add_ident,real_add_closed,real_lt_subst_left,real_lt_trans}.
    $(1 + 1)\neq 0$ and $\inv{(1 + 1)}\in\reals$ by \cref{real_pos_nonzero,real_mul_inverse}.
    $\delta = \epsilon \rmul \inv{(1 + 1)}$ by \cref{real_div}.
    Show that $\delta\in\reals$.
    \begin{subproof}
        Follows by \cref{real_mul_closed}.
    \end{subproof}
    Show that $0 < \delta$.
    \begin{subproof}
        $0\rmul \inv{(1 + 1)} = 0$ and $0 < \inv{(1 + 1)}$ by \cref{real_mul_zero,real_inv_pos}.
        $0 < \delta$ by \cref{real_lt_mul_pos}.
    \end{subproof}
    Show that $\delta + \delta = \epsilon$.
    \begin{subproof}
        $(1 + 1) \rmul \delta = \delta + \delta$ by
            \cref{real_mul_ident,real_right_distrib,real_add_eq_subst}.
        $\delta + \delta = \epsilon$ by
            \cref{real_div,real_mul_assoc,real_mul_comm,real_mul_inverse,real_mul_ident}.
    \end{subproof}
\end{proof}

% from §21 Item 8: addition operation is continuous
\begin{lemma}\label{real_addition_continuous}
    Let $T = \productTopology{\standardTopologyR}{\standardTopologyR}{\reals}{\reals}$.
    Then $\realAddition$ is continuous from $\reals\times\reals$ to $\reals$ with respect to $T$ and $\standardTopologyR$.
\end{lemma}
\begin{proof}
    Let $f = \realAddition$.
    For all $p,z_1,z_2$ such that $(p,z_1)\in f$ and $(p,z_2)\in f$ we have $z_1 = z_2$
        by \cref{real_addition_operation,pair_eq_iff}.
    $f$ is a function by \cref{rightunique,real_addition_operation,relation}.
    For all $p\in\dom{f}$ we have $f(p)\in\reals$
        by \cref{dom_iff,real_addition_operation,pair_eq_iff,function_apply_intro}.
    $f$ is a function to $\reals$.
    $\dom{f} = \reals\times\reals$ by
        \cref{dom_iff,real_addition_operation,pair_eq_iff,subseteq,times_elem_is_tuple,real_add_closed,fst_eq,snd_eq,dom_intro,subseteq_antisymmetric}.
    Show that $f\in\funs{\reals\times\reals}{\reals}$.
    \begin{subproof}
        Follows by \cref{funs_intro}.
    \end{subproof}
    Show that for all $p\in\reals\times\reals$ we have
        $f$ is continuous at $p$ from $\reals\times\reals$ to $\reals$ with respect to $T$ and $\standardTopologyR$.
        \begin{subproof}
            Fix $p$ such that $p\in\reals\times\reals$.
            Take $x_0,y_0$ such that $p = (x_0,y_0)$ and $x_0\in\reals$ and $y_0\in\reals$ by \cref{times_elem_is_tuple}.
            $f(p) = x_0 + y_0$ by
                \cref{fst_eq,snd_eq,real_add_closed,real_addition_operation,function_apply_intro}.
        Show that for all $V$ such that $V$ is a neighborhood of $f(p)$ in $\reals$ with respect to $\standardTopologyR$
            there exists $U$ such that $U$ is a neighborhood of $p$ in $\reals\times\reals$ with respect to $T$ and $\img{f}{U}\subseteq V$.
        \begin{subproof}
            Fix $V$ such that $V$ is a neighborhood of $f(p)$ in $\reals$ with respect to $\standardTopologyR$.
            Let $d = \standardMetricR$.
            $\metricTopology{d}{\reals} = \standardTopologyR$ by \cref{standard_metric_topology}.
            $\standardBasisR$ is a basis on $\reals$ and $\standardTopologyR$ is a topology on $\reals$
                by \cref{standard_basis_is_basis,basis_topology_is_topology,standard_topology_reals}.
            $V\subseteq\reals$ by
                \cref{neighborhood,open_in,standard_basis_is_basis,basis_topology_is_topology,standard_topology_reals,topology_on,subseteq,pow_iff}.
            $d$ is a metric on $\reals$ by \cref{standard_metric_is_metric}.
            Take $\epsilon\in\reals$ such that $0 < \epsilon$ and
                $\metricBall{d}{\reals}{f(p)}{\epsilon} \subseteq V$ by \cref{neighborhood,standard_metric_topology,metric_open_iff}.
            Let $\delta = \epsilon / (1 + 1)$.
            $\delta\in\reals$ and $0 < \delta$ and $\delta + \delta = \epsilon$ by \cref{real_half_of_positive}.
            Let $U_1 = \metricBall{d}{\reals}{x_0}{\delta}$.
            Let $U_2 = \metricBall{d}{\reals}{y_0}{\delta}$.
            $x_0\in\reals$ and $y_0\in\reals$ and $\delta\in\reals$ and $0 < \delta$.
            $U_1$ is open in $\reals$ with respect to $\standardTopologyR$ and
                $U_2$ is open in $\reals$ with respect to $\standardTopologyR$
                by \cref{metric_ball,subseteq,pow_iff,metric_basis,metric_basis_is_basis,basis_elem_open_in_generated,metric_topology,basis_for_topology,open_in}.
            Let $U = U_1\times U_2$.
            $U$ is open in $\reals\times\reals$ with respect to $T$ by
                \cref{open_in,topology_on,subseteq,pow_iff,times_subseteq_left,times_subseteq_right,subseteq_transitive,product_basis,product_basis_is_basis,basis_topology_is_topology,product_topology,basis_elem_open_in_generated}.
            $p\in U$ by
                \cref{metric_zero_characterization,standard_metric_is_metric,metric_on,real_lt_subst_left,metric_ball,times_tuple_intro}.
            Show that for all $q$ such that $q\in U$ we have
                $f(q)\in\metricBall{d}{\reals}{f(p)}{\epsilon}$.
            \begin{subproof}
                    Fix $q$ such that $q\in U$.
                    Take $x,y$ such that $q = (x,y)$ and $x\in U_1$ and $y\in U_2$ by \cref{times_elem_is_tuple}.
                    $x\in\reals$ and $y\in\reals$ by \cref{metric_ball}.
                    $\abs{x - x_0} < \delta$ and $\abs{y - y_0} < \delta$
                        by \cref{metric_ball,standard_metric_value,real_abs_sub_swap,real_lt_subst_left}.
                    $(q, x + y)\in f$ by \cref{real_add_closed,times_tuple_intro,fst_eq,snd_eq,real_addition_operation}.
                    $f(q) = x + y$ by \cref{function_apply_intro}.
                    $x - x_0 = x + \neg{x_0}$ by \cref{real_sub}.
                    $y - y_0 = y + \neg{y_0}$ by \cref{real_sub}.
                    $(x + y) - (x_0 + y_0) = (x + y) + \neg{(x_0 + y_0)}$ by \cref{real_sub}.
                    $\neg{(x_0 + y_0)} = \neg{x_0} + \neg{y_0}$ by \cref{real_neg_addition}.
                    $\neg{x_0}\in\reals$ and $\neg{y_0}\in\reals$ by \cref{real_add_inverse}.
                    $\neg{x_0} + \neg{y_0}\in\reals$ by \cref{real_add_closed}.
                    $y + \neg{x_0}\in\reals$ and $y + \neg{y_0}\in\reals$ by \cref{real_add_closed}.
                    $\neg{x_0} + y\in\reals$ and $\neg{x_0} + (y + \neg{y_0})\in\reals$ by \cref{real_add_closed}.
                    $(x + y) + (\neg{x_0} + \neg{y_0}) = x + (y + (\neg{x_0} + \neg{y_0}))$ by \cref{real_add_assoc}.
                    $y + (\neg{x_0} + \neg{y_0}) = (y + \neg{x_0}) + \neg{y_0}$ by \cref{real_add_assoc}.
                    $y + \neg{x_0} = \neg{x_0} + y$ by \cref{real_add_comm}.
                    $(\neg{x_0} + y) + \neg{y_0} = \neg{x_0} + (y + \neg{y_0})$ by \cref{real_add_assoc}.
                    $y + (\neg{x_0} + \neg{y_0}) = \neg{x_0} + (y + \neg{y_0})$.
                    $x + (\neg{x_0} + (y + \neg{y_0})) = (x + \neg{x_0}) + (y + \neg{y_0})$ by \cref{real_add_assoc}.
                    $(x + y) + (\neg{x_0} + \neg{y_0}) = (x + \neg{x_0}) + (y + \neg{y_0})$.
                    $(x + y) - (x_0 + y_0) = (x - x_0) + (y - y_0)$.
                    $x - x_0\in\reals$ and $y - y_0\in\reals$ by
                        \cref{real_sub,real_add_inverse,real_add_closed}.
                    $x + y\in\reals$ by \cref{real_add_closed}.
                    $x_0 + y_0\in\reals$ by \cref{real_add_closed}.
                    $(x + y) - (x_0 + y_0)\in\reals$ and $\abs{(x + y) - (x_0 + y_0)}\in\reals$ by
                        \cref{real_sub,real_add_closed,real_abs_in_reals}.
                    $\abs{(x + y) - (x_0 + y_0)} \leq \abs{x - x_0} + \abs{y - y_0}$ by \cref{real_triangle_inequality}.
                    $\abs{x - x_0}\in\reals$ by \cref{real_abs_in_reals}.
                    $\abs{y - y_0}\in\reals$ by \cref{real_abs_in_reals}.
                    $\abs{x - x_0} + \abs{y - y_0}\in\reals$ by \cref{real_add_closed}.
                    $\delta + \delta\in\reals$ by \cref{real_add_closed}.
                    $\abs{x - x_0} + \abs{y - y_0} < \delta + \delta$ by \cref{real_lt_add_two}.
                    $\abs{(x + y) - (x_0 + y_0)} < \delta + \delta$
                        by \cref{real_leq_split,real_lt_trans,real_lt_subst_left}.
                    $\abs{(x + y) - (x_0 + y_0)} < \epsilon$ by \cref{real_lt_subst_right}.
                    $\abs{f(q) - f(p)} = \abs{(x + y) - (x_0 + y_0)}$.
                    $\abs{f(q) - f(p)} < \epsilon$ by \cref{real_lt_subst_left}.
                $f$ is a function to $\reals$ and $\dom{f} = \reals\times\reals$ by \cref{funs_elim}.
                $f(p)\in\reals$ and $f(q)\in\reals$.
                $d((f(p),f(q))) < \epsilon$ by \cref{standard_metric_value,real_abs_sub_swap,real_lt_subst_left}.
                $f(q)\in\metricBall{d}{\reals}{f(p)}{\epsilon}$ by \cref{metric_ball}.
            \end{subproof}
            $\img{f}{U}\subseteq V$ by
                \cref{img_iff,function_apply_intro,subseteq,subseteq_transitive}.
            $U$ is a neighborhood of $p$ in $\reals\times\reals$ with respect to $T$ by \cref{neighborhood}.
        \end{subproof}
        $f$ is continuous at $p$ from $\reals\times\reals$ to $\reals$ with respect to $T$ and $\standardTopologyR$ by
            \cref{standard_basis_is_basis,basis_topology_is_topology,standard_topology_reals,product_basis_is_basis,product_topology,continuous_at}.
    \end{subproof}
    Show that $f$ is continuous from $\reals\times\reals$ to $\reals$ with respect to $T$ and $\standardTopologyR$.
    \begin{subproof}
        Follows by
            \cref{standard_basis_is_basis,basis_topology_is_topology,standard_topology_reals,product_basis_is_basis,product_topology,continuous_at_implies_continuous}.
    \end{subproof}
\end{proof}

% from §21 Item 9: subtraction operation is continuous
\begin{lemma}\label{real_subtraction_continuous}
    Let $T = \productTopology{\standardTopologyR}{\standardTopologyR}{\reals}{\reals}$.
    Then $\realSubtraction$ is continuous from $\reals\times\reals$ to $\reals$ with respect to $T$ and $\standardTopologyR$.
\end{lemma}
\begin{proof}
    Let $f = \realSubtraction$.
    For all $p,z_1,z_2$ such that $(p,z_1)\in f$ and $(p,z_2)\in f$ we have $z_1 = z_2$
        by \cref{real_subtraction_operation,pair_eq_iff}.
    $f$ is a function by \cref{rightunique,real_subtraction_operation,relation}.
    For all $p\in\dom{f}$ we have $f(p)\in\reals$
        by \cref{dom_iff,real_subtraction_operation,pair_eq_iff,function_apply_intro}.
    $f$ is a function to $\reals$.
    Show that for all $p\in\reals\times\reals$ we have $p\in\dom{f}$.
    \begin{subproof}
        Fix $p\in\reals\times\reals$.
        Take $x,y$ such that $p = (x,y)$ and $x\in\reals$ and $y\in\reals$ by \cref{times_elem_is_tuple}.
        $x - y\in\reals$ by \cref{real_sub,real_add_closed,real_add_inverse}.
        $(p,x - y)\in f$ by \cref{real_subtraction_operation,fst_eq,snd_eq,times_tuple_intro}.
        $p\in\dom{f}$ by \cref{dom_intro}.
    \end{subproof}
    $\reals\times\reals\subseteq\dom{f}$ by \cref{subseteq}.
    Show that for all $p\in\dom{f}$ we have $p\in\reals\times\reals$.
    \begin{subproof}
        Follows by \cref{dom_iff,real_subtraction_operation,pair_eq_iff}.
    \end{subproof}
    $\dom{f}\subseteq\reals\times\reals$ by \cref{subseteq}.
    $\dom{f} = \reals\times\reals$ by \cref{subseteq_antisymmetric}.
    Show that $f\in\funs{\reals\times\reals}{\reals}$.
    \begin{subproof}
        Follows by \cref{funs_intro}.
    \end{subproof}
    Show that for all $p\in\reals\times\reals$ we have
        $f$ is continuous at $p$ from $\reals\times\reals$ to $\reals$ with respect to $T$ and $\standardTopologyR$.
        \begin{subproof}
            Fix $p$ such that $p\in\reals\times\reals$.
            Take $x_0,y_0$ such that $p = (x_0,y_0)$ and $x_0\in\reals$ and $y_0\in\reals$ by \cref{times_elem_is_tuple}.
            $f(p) = x_0 - y_0$ by
                \cref{fst_eq,snd_eq,real_sub,real_add_closed,real_add_inverse,real_subtraction_operation,function_apply_intro}.
        Show that for all $V$ such that $V$ is a neighborhood of $f(p)$ in $\reals$ with respect to $\standardTopologyR$
            there exists $U$ such that $U$ is a neighborhood of $p$ in $\reals\times\reals$ with respect to $T$ and $\img{f}{U}\subseteq V$.
        \begin{subproof}
            Fix $V$ such that $V$ is a neighborhood of $f(p)$ in $\reals$ with respect to $\standardTopologyR$.
            Let $d = \standardMetricR$.
            $\metricTopology{d}{\reals} = \standardTopologyR$ by \cref{standard_metric_topology}.
            $\standardBasisR$ is a basis on $\reals$ and $\standardTopologyR$ is a topology on $\reals$
                by \cref{standard_basis_is_basis,basis_topology_is_topology,standard_topology_reals}.
            $V\subseteq\reals$ by
                \cref{neighborhood,open_in,standard_basis_is_basis,basis_topology_is_topology,standard_topology_reals,topology_on,subseteq,pow_iff}.
            $d$ is a metric on $\reals$ by \cref{standard_metric_is_metric}.
            Take $\epsilon\in\reals$ such that $0 < \epsilon$ and
                $\metricBall{d}{\reals}{f(p)}{\epsilon} \subseteq V$ by \cref{neighborhood,standard_metric_topology,metric_open_iff}.
            Let $\delta = \epsilon / (1 + 1)$.
            $\delta\in\reals$ and $0 < \delta$ and $\delta + \delta = \epsilon$ by \cref{real_half_of_positive}.
            Let $U_1 = \metricBall{d}{\reals}{x_0}{\delta}$.
            Let $U_2 = \metricBall{d}{\reals}{y_0}{\delta}$.
            $x_0\in\reals$ and $y_0\in\reals$ and $\delta\in\reals$ and $0 < \delta$.
            $U_1$ is open in $\reals$ with respect to $\standardTopologyR$ and
                $U_2$ is open in $\reals$ with respect to $\standardTopologyR$
                by \cref{metric_ball,subseteq,pow_iff,metric_basis,metric_basis_is_basis,basis_elem_open_in_generated,metric_topology,basis_for_topology,open_in}.
            Let $U = U_1\times U_2$.
            $U$ is open in $\reals\times\reals$ with respect to $T$ by
                \cref{open_in,topology_on,subseteq,pow_iff,times_subseteq_left,times_subseteq_right,subseteq_transitive,product_basis,product_basis_is_basis,basis_topology_is_topology,product_topology,basis_elem_open_in_generated}.
            $p\in U$ by
                \cref{metric_zero_characterization,standard_metric_is_metric,metric_on,real_lt_subst_left,metric_ball,times_tuple_intro}.
            Show that for all $q$ such that $q\in U$ we have
                $f(q)\in\metricBall{d}{\reals}{f(p)}{\epsilon}$.
            \begin{subproof}
                    Fix $q$ such that $q\in U$.
                    Take $x,y$ such that $q = (x,y)$ and $x\in U_1$ and $y\in U_2$ by \cref{times_elem_is_tuple}.
                    $x\in\reals$ and $y\in\reals$ by \cref{metric_ball}.
                    $\abs{x - x_0} < \delta$ and $\abs{y - y_0} < \delta$
                        by \cref{metric_ball,standard_metric_value,real_abs_sub_swap,real_lt_subst_left}.
                    $(q, x - y)\in f$ by \cref{real_sub,real_add_closed,real_add_inverse,times_tuple_intro,fst_eq,snd_eq,real_subtraction_operation}.
                    $f(q) = x - y$ by \cref{function_apply_intro}.
                    $x - y = x + \neg{y}$ by \cref{real_sub}.
                    $x_0 - y_0 = x_0 + \neg{y_0}$ by \cref{real_sub}.
                    $(x - y) - (x_0 - y_0) = (x + \neg{y}) + \neg{(x_0 + \neg{y_0})}$ by \cref{real_sub}.
                    $\neg{y_0}\in\reals$ by \cref{real_add_inverse}.
                    $\neg{(x_0 + \neg{y_0})} = \neg{x_0} + \neg{(\neg{y_0})}$ by \cref{real_neg_addition}.
                    $\neg{(\neg{y_0})} = y_0$ by \cref{real_neg_neg}.
                    $(x - y) - (x_0 - y_0) = (x + \neg{y}) + (\neg{x_0} + y_0)$.
                    $\neg{y}\in\reals$ by \cref{real_add_inverse}.
                    $\neg{x_0}\in\reals$ by \cref{real_add_inverse}.
                    $\neg{x_0} + y_0\in\reals$ by \cref{real_add_closed}.
                    $\neg{y} + \neg{x_0}\in\reals$ by \cref{real_add_closed}.
                    $\neg{y} + y_0\in\reals$ by \cref{real_add_closed}.
                    $(x + \neg{y}) + (\neg{x_0} + y_0) = x + (\neg{y} + (\neg{x_0} + y_0))$ by \cref{real_add_assoc}.
                    $\neg{y} + (\neg{x_0} + y_0) = (\neg{y} + \neg{x_0}) + y_0$ by \cref{real_add_assoc}.
                    $\neg{y} + \neg{x_0} = \neg{x_0} + \neg{y}$ by \cref{real_add_comm}.
                    $\neg{y} + (\neg{x_0} + y_0) = (\neg{x_0} + \neg{y}) + y_0$.
                    $(\neg{x_0} + \neg{y}) + y_0 = \neg{x_0} + (\neg{y} + y_0)$ by \cref{real_add_assoc}.
                    $(x - y) - (x_0 - y_0) = x + (\neg{x_0} + (\neg{y} + y_0))$ by \cref{real_add_assoc}.
                    $\neg{y} + y_0 = y_0 + \neg{y}$ by \cref{real_add_comm}.
                    $(x - y) - (x_0 - y_0) = x + (\neg{x_0} + (y_0 + \neg{y}))$.
                    $y_0 + \neg{y} = y_0 - y$ by \cref{real_sub}.
                    $x + \neg{x_0} = x - x_0$ by \cref{real_sub}.
                    $(x - y) - (x_0 - y_0) = (x - x_0) + (y_0 - y)$ by \cref{real_add_assoc}.
                    $x - x_0\in\reals$ and $y_0 - y\in\reals$ by \cref{real_sub,real_add_closed}.
                    $\abs{(x - x_0) + (y_0 - y)} \leq \abs{x - x_0} + \abs{y_0 - y}$ by \cref{real_triangle_inequality}.
                    $\abs{y_0 - y} = \abs{y - y_0}$ by \cref{real_abs_sub_swap}.
                    $\abs{y_0 - y} < \delta$ by \cref{real_lt_subst_left}.
                    $\abs{x - x_0}\in\reals$ by \cref{real_abs_in_reals}.
                    $\abs{y_0 - y}\in\reals$ by \cref{real_abs_in_reals}.
                    $\abs{x - x_0} + \abs{y_0 - y} < \delta + \delta$ by \cref{real_lt_add_two}.
                    $\abs{x - x_0} + \abs{y_0 - y}\in\reals$ by \cref{real_add_closed}.
                    $\delta + \delta\in\reals$ by \cref{real_add_closed}.
                    $(x - x_0) + (y_0 - y)\in\reals$ by \cref{real_add_closed}.
                    $\abs{(x - x_0) + (y_0 - y)}\in\reals$ by \cref{real_abs_in_reals}.
                    $\abs{(x - x_0) + (y_0 - y)} < \delta + \delta$
                        by \cref{real_leq_split,real_lt_trans,real_lt_subst_left}.
                    $\abs{(x - x_0) + (y_0 - y)} < \epsilon$ by \cref{real_lt_subst_right}.
                    $\abs{f(q) - f(p)} = \abs{(x - y) - (x_0 - y_0)}$.
                    $\abs{f(q) - f(p)} < \epsilon$ by \cref{real_lt_subst_left}.
                $f$ is a function to $\reals$ and $\dom{f} = \reals\times\reals$ by \cref{funs_elim}.
                $f(p)\in\reals$ and $f(q)\in\reals$.
                $d((f(p),f(q))) < \epsilon$ by \cref{standard_metric_value,real_abs_sub_swap,real_lt_subst_left}.
                $f(q)\in\metricBall{d}{\reals}{f(p)}{\epsilon}$ by \cref{metric_ball}.
            \end{subproof}
            $\img{f}{U}\subseteq V$ by
                \cref{img_iff,function_apply_intro,subseteq,subseteq_transitive}.
            $U$ is a neighborhood of $p$ in $\reals\times\reals$ with respect to $T$ by \cref{neighborhood}.
        \end{subproof}
        $f$ is continuous at $p$ from $\reals\times\reals$ to $\reals$ with respect to $T$ and $\standardTopologyR$ by
            \cref{standard_basis_is_basis,basis_topology_is_topology,standard_topology_reals,product_basis_is_basis,product_topology,continuous_at}.
    \end{subproof}
    Show that $f$ is continuous from $\reals\times\reals$ to $\reals$ with respect to $T$ and $\standardTopologyR$.
    \begin{subproof}
        Follows by
            \cref{standard_basis_is_basis,basis_topology_is_topology,standard_topology_reals,product_basis_is_basis,product_topology,continuous_at_implies_continuous}.
    \end{subproof}
\end{proof}

% from §21 helper lemma 4: absolute value bound by nearby point
\begin{lemma}\label{real_abs_sub_lt_bound}
    Suppose $x\in\reals$ and $y\in\reals$ and $\delta\in\reals$.
    Suppose $\abs{x - y} < \delta$.
    Then $\abs{x} < \abs{y} + \delta$.
\end{lemma}
\begin{proof}
    $x - y = x + \neg{y}$ and $\neg{y}\in\reals$ and $x - y\in\reals$ by
        \cref{real_sub,real_add_inverse,real_add_closed}.
    $(x - y) + y = x$ by \cref{real_add_assoc,real_add_inverse,real_add_comm,real_add_ident,real_add_eq_subst}.
    $\abs{x} \leq \abs{x - y} + \abs{y}$ by
        \cref{real_sub,real_add_inverse,real_add_closed,real_triangle_inequality,real_lt_subst_left}.
    $\abs{x - y} + \abs{y}\in\reals$ and $\delta + \abs{y}\in\reals$ and
        $\abs{x - y} + \abs{y} < \delta + \abs{y}$ by \cref{real_abs_in_reals,real_add_closed,real_lt_add}.
    $\abs{x} < \abs{y} + \delta$ by
        \cref{real_abs_in_reals,real_leq_split,real_lt_trans,real_lt_subst_left,real_add_comm,real_lt_subst_right}.
\end{proof}
% from §21 helper lemma 5: positive sum with nonnegative term
\begin{lemma}\label{real_pos_add_nonneg}
    Suppose $a\in\reals$ and $b\in\reals$.
    Suppose $a \geq 0$ and $0 < b$.
    Then $0 < a + b$.
\end{lemma}
\begin{proof}
    $0 < a + b$ by
        \cref{real_add_closed,real_add_ident,real_leq_split,real_add_eq_subst,real_lt_subst_right,real_lt_add,real_lt_subst_left,real_lt_trans}.
\end{proof}

% from §21 helper lemma 6: substitute equality in multiplication
\begin{lemma}\label{real_mul_eq_subst}
    Suppose $a\in\reals$ and $b\in\reals$ and $c\in\reals$.
    Suppose $a = b$.
    Then $a \rmul c = b \rmul c$.
\end{lemma}
\begin{proof}
    Trivial.
\end{proof}

% from §21 Item 10: multiplication operation is continuous
\begin{lemma}\label{real_multiplication_continuous}
    Let $T = \productTopology{\standardTopologyR}{\standardTopologyR}{\reals}{\reals}$.
    Then $\realMultiplication$ is continuous from $\reals\times\reals$ to $\reals$ with respect to $T$ and $\standardTopologyR$.
\end{lemma}
\begin{proof}
    Let $f = \realMultiplication$.
    For all $p,z_1,z_2$ such that $(p,z_1)\in f$ and $(p,z_2)\in f$ we have $z_1 = z_2$
        by \cref{real_multiplication_operation,pair_eq_iff}.
    $f$ is a function by \cref{rightunique,real_multiplication_operation,relation}.
    For all $p\in\dom{f}$ we have $f(p)\in\reals$
        by \cref{dom_iff,real_multiplication_operation,pair_eq_iff,function_apply_intro}.
    $f$ is a function to $\reals$.
    $\dom{f} = \reals\times\reals$ by
        \cref{dom_iff,real_multiplication_operation,pair_eq_iff,subseteq,times_elem_is_tuple,real_mul_closed,fst_eq,snd_eq,dom_intro,subseteq_antisymmetric}.
    Show that $f\in\funs{\reals\times\reals}{\reals}$.
    \begin{subproof}
        Follows by \cref{funs_intro}.
    \end{subproof}
    Show that for all $p\in\reals\times\reals$ we have
        $f$ is continuous at $p$ from $\reals\times\reals$ to $\reals$ with respect to $T$ and $\standardTopologyR$.
    \begin{subproof}
        Fix $p$ such that $p\in\reals\times\reals$.
        Take $x_0,y_0$ such that $p = (x_0,y_0)$ and $x_0\in\reals$ and $y_0\in\reals$ by \cref{times_elem_is_tuple}.
        $f(p) = x_0 \rmul y_0$ by
            \cref{fst_eq,snd_eq,real_mul_closed,real_multiplication_operation,function_apply_intro}.
        Show that for all $V$ such that $V$ is a neighborhood of $f(p)$ in $\reals$ with respect to $\standardTopologyR$
            there exists $U$ such that $U$ is a neighborhood of $p$ in $\reals\times\reals$ with respect to $T$ and $\img{f}{U}\subseteq V$.
        \begin{subproof}
            Fix $V$ such that $V$ is a neighborhood of $f(p)$ in $\reals$ with respect to $\standardTopologyR$.
            Let $d = \standardMetricR$.
            $\metricTopology{d}{\reals} = \standardTopologyR$ by \cref{standard_metric_topology}.
            $\standardBasisR$ is a basis on $\reals$ and $\standardTopologyR$ is a topology on $\reals$
                by \cref{standard_basis_is_basis,basis_topology_is_topology,standard_topology_reals}.
            $V\subseteq\reals$ by
                \cref{open_in,neighborhood,standard_basis_is_basis,basis_topology_is_topology,standard_topology_reals,topology_on,subseteq,pow_iff}.
            $d$ is a metric on $\reals$ by \cref{standard_metric_is_metric}.
            Take $\epsilon\in\reals$ such that $0 < \epsilon$ and
                $\metricBall{d}{\reals}{f(p)}{\epsilon} \subseteq V$ by \cref{neighborhood,standard_metric_topology,metric_open_iff}.
            Let $m = \abs{x_0} + \abs{y_0} + 1$.
            $\abs{x_0}\in\reals$ by \cref{real_abs_in_reals}.
            $\abs{y_0}\in\reals$ by \cref{real_abs_in_reals}.
            $\abs{x_0} + \abs{y_0}\in\reals$ by \cref{real_add_closed}.
            $1\in\reals$ by \cref{real_mul_ident}.
            $m\in\reals$ by \cref{real_add_closed}.
            $\abs{y_0} \geq 0$ by \cref{real_abs_nonnegative}.
            $0 < 1$ by \cref{real_zero_lt_one}.
            $\abs{y_0} + 1\in\reals$ by \cref{real_add_closed}.
            $0 < \abs{y_0} + 1$ by \cref{real_pos_add_nonneg}.
            $\abs{x_0} \geq 0$ by \cref{real_abs_nonnegative}.
            $0 < \abs{x_0} + (\abs{y_0} + 1)$ by \cref{real_pos_add_nonneg}.
            $0 < m$ by \cref{real_add_assoc,real_lt_subst_right}.
            $m \geq 0$ by \cref{real_lt_imp_leq}.
            Let $\delta = \epsilon / (m + \epsilon)$.
            $m + \epsilon\in\reals$ by \cref{real_add_closed}.
            $0 < m + \epsilon$ by \cref{real_pos_add_nonneg}.
            $(m + \epsilon)\neq 0$ by \cref{real_pos_nonzero}.
            $\inv{(m + \epsilon)}\in\reals$ by \cref{real_mul_inverse}.
            $0 < \inv{(m + \epsilon)}$ by \cref{real_inv_pos}.
            $0\in\reals$ and $0 \rmul \inv{(m + \epsilon)} = 0$ by \cref{real_add_ident,real_mul_zero}.
            $0 \rmul \inv{(m + \epsilon)} < \epsilon \rmul \inv{(m + \epsilon)}$ by \cref{real_lt_mul_pos}.
            $0 < \epsilon \rmul \inv{(m + \epsilon)}$ by \cref{real_lt_subst_left}.
            $\delta = \epsilon \rmul \inv{(m + \epsilon)}$ by \cref{real_div}.
            $0 < \delta$ by \cref{real_lt_subst_left}.
            $\epsilon < m + \epsilon$ by \cref{real_lt_add,real_add_ident,real_lt_subst_left}.
            $\epsilon \rmul \inv{(m + \epsilon)} < (m + \epsilon) \rmul \inv{(m + \epsilon)}$ by \cref{real_lt_mul_pos}.
            $(m + \epsilon) \rmul \inv{(m + \epsilon)} = 1$ by \cref{real_mul_inverse}.
            $\epsilon \rmul \inv{(m + \epsilon)} < 1$ by \cref{real_lt_subst_right}.
            $\delta < 1$ by \cref{real_lt_subst_left}.
            Let $U_1 = \metricBall{d}{\reals}{x_0}{\delta}$.
            Let $U_2 = \metricBall{d}{\reals}{y_0}{\delta}$.
            $x_0\in\reals$ and $y_0\in\reals$ and $\delta\in\reals$ and $0 < \delta$.
            $U_1$ is open in $\reals$ with respect to $\standardTopologyR$ and
                $U_2$ is open in $\reals$ with respect to $\standardTopologyR$
                by \cref{metric_ball,subseteq,pow_iff,metric_basis,metric_basis_is_basis,basis_elem_open_in_generated,metric_topology,basis_for_topology,open_in}.
            Let $U = U_1\times U_2$.
            $U$ is open in $\reals\times\reals$ with respect to $T$ by
                \cref{open_in,topology_on,subseteq,pow_iff,times_subseteq_left,times_subseteq_right,subseteq_transitive,product_basis,product_basis_is_basis,basis_topology_is_topology,product_topology,basis_elem_open_in_generated}.
            $d((x_0,x_0)) = 0$ and $d((y_0,y_0)) = 0$ by \cref{metric_zero_characterization,standard_metric_is_metric,metric_on}.
            $d((x_0,x_0)) < \delta$ and $d((y_0,y_0)) < \delta$ by \cref{real_lt_subst_left}.
            $p\in U$ by \cref{metric_ball,times_tuple_intro}.
            Show that for all $q$ such that $q\in U$ we have
                $f(q)\in\metricBall{d}{\reals}{f(p)}{\epsilon}$.
            \begin{subproof}
                Fix $q$ such that $q\in U$.
                Take $x,y$ such that $q = (x,y)$ and $x\in U_1$ and $y\in U_2$ by \cref{times_elem_is_tuple}.
                $x\in\reals$ and $y\in\reals$ by \cref{metric_ball}.
                $\abs{x - x_0} < \delta$ and $\abs{y - y_0} < \delta$
                    by \cref{metric_ball,standard_metric_value,real_abs_sub_swap,real_lt_subst_left}.
                $(q, x \rmul y)\in f$ by \cref{real_mul_closed,times_tuple_intro,fst_eq,snd_eq,real_multiplication_operation}.
                $f(q) = x \rmul y$ by \cref{function_apply_intro}.
                Show that $\abs{f(q) - f(p)} < \epsilon$.
                \begin{subproof}
                        $x - x_0 = x + \neg{x_0}$ by \cref{real_sub}.
                        $y - y_0 = y + \neg{y_0}$ by \cref{real_sub}.
                        $\neg{x_0}\in\reals$ by \cref{real_add_inverse}.
                        $\neg{y_0}\in\reals$ by \cref{real_add_inverse}.
                        $(x - x_0) \rmul y = (x + \neg{x_0}) \rmul y$.
                        $x_0 \rmul (y - y_0) = x_0 \rmul (y + \neg{y_0})$.
                        $y + \neg{y_0}\in\reals$ by \cref{real_add_closed}.
                        $(x + \neg{x_0}) \rmul y = (x \rmul y) + (\neg{x_0} \rmul y)$ by \cref{real_right_distrib}.
                        $x_0 \rmul (y + \neg{y_0}) = (y + \neg{y_0}) \rmul x_0$ by \cref{real_mul_comm}.
                        $(y + \neg{y_0}) \rmul x_0 = (y \rmul x_0) + (\neg{y_0} \rmul x_0)$ by \cref{real_right_distrib}.
                        $y \rmul x_0 = x_0 \rmul y$ by \cref{real_mul_comm}.
                        $\neg{y_0} \rmul x_0 = x_0 \rmul \neg{y_0}$ by \cref{real_mul_comm}.
                        $x_0 \rmul \neg{y_0} = \neg{(x_0 \rmul y_0)}$ by \cref{real_mul_neg_right}.
                        $(x - x_0) \rmul y + x_0 \rmul (y - y_0) = (x \rmul y) + (\neg{x_0} \rmul y) + (x_0 \rmul y) + \neg{(x_0 \rmul y_0)}$.
                        $(\neg{x_0} \rmul y) + (x_0 \rmul y) = (\neg{x_0} + x_0) \rmul y$ by \cref{real_right_distrib}.
                        $\neg{x_0} + x_0 = 0$ by \cref{real_add_inverse,real_add_comm}.
                        $0 \rmul y = 0$ by \cref{real_mul_zero}.
                        $(\neg{x_0} \rmul y) + (x_0 \rmul y) = 0$.
                        $0\in\reals$ by \cref{real_add_ident}.
                        $x \rmul y\in\reals$ by \cref{real_mul_closed}.
                        $\neg{x_0} \rmul y\in\reals$ by \cref{real_mul_closed}.
                        $x_0 \rmul y\in\reals$ by \cref{real_mul_closed}.
                        $(\neg{x_0} \rmul y) + (x_0 \rmul y)\in\reals$ by \cref{real_add_closed}.
                        $(x \rmul y) + (\neg{x_0} \rmul y) + (x_0 \rmul y) = x \rmul y$
                            by \cref{real_add_assoc,real_add_comm,real_add_eq_subst,real_add_ident}.
                        $(x \rmul y) + (\neg{x_0} \rmul y)\in\reals$ by \cref{real_add_closed}.
                        $(x \rmul y) + (\neg{x_0} \rmul y) + (x_0 \rmul y)\in\reals$ by \cref{real_add_closed}.
                        $\neg{(x_0 \rmul y_0)}\in\reals$ by \cref{real_add_inverse}.
                        $(x \rmul y) + (\neg{x_0} \rmul y) + (x_0 \rmul y) + \neg{(x_0 \rmul y_0)} = (x \rmul y) + \neg{(x_0 \rmul y_0)}$ by \cref{real_add_eq_subst}.
                        $(x - x_0) \rmul y + x_0 \rmul (y - y_0) = (x \rmul y) + \neg{(x_0 \rmul y_0)}$.
                        $(x \rmul y) + \neg{(x_0 \rmul y_0)} = (x \rmul y) - (x_0 \rmul y_0)$ by \cref{real_sub}.
                        $(x \rmul y) - (x_0 \rmul y_0) = (x - x_0) \rmul y + x_0 \rmul (y - y_0)$.
                        $x - x_0\in\reals$ by \cref{real_sub,real_add_closed}.
                        $y - y_0\in\reals$ by \cref{real_sub,real_add_closed}.
                        $(x - x_0) \rmul y\in\reals$ by \cref{real_mul_closed}.
                        $x_0 \rmul (y - y_0)\in\reals$ by \cref{real_mul_closed}.
                        $\abs{(x - x_0) \rmul y + x_0 \rmul (y - y_0)} \leq \abs{(x - x_0) \rmul y} + \abs{x_0 \rmul (y - y_0)}$ by \cref{real_triangle_inequality}.
                        $\abs{(x - x_0) \rmul y} = \abs{x - x_0} \rmul \abs{y}$ by \cref{real_abs_mul}.
                        $\abs{x_0 \rmul (y - y_0)} = \abs{x_0} \rmul \abs{y - y_0}$ by \cref{real_abs_mul}.
                        $\abs{(x - x_0) \rmul y + x_0 \rmul (y - y_0)} \leq \abs{x - x_0} \rmul \abs{y} + \abs{x_0} \rmul \abs{y - y_0}$.
                        $\delta\in\reals$ by \cref{real_mul_closed}.
                        $\abs{y} < \abs{y_0} + \delta$ by \cref{real_abs_sub_lt_bound}.
                        $0 < \abs{y_0} + \delta$ by \cref{real_abs_nonnegative,real_pos_add_nonneg}.
                        $0 < \abs{x - x_0}$ or $0 = \abs{x - x_0}$ by \cref{real_abs_nonnegative,real_leq_split}.
                        Show that $\abs{x - x_0} \rmul \abs{y} < \delta \rmul (\abs{y_0} + \delta)$.
                        \begin{subproof}
                            \begin{byCase}
                            \caseOf{$0 = \abs{x - x_0}$.}
                                $\abs{x - x_0} = 0$ by assumption.
                                $\abs{x - x_0}\in\reals$ by \cref{real_abs_in_reals}.
                                $\abs{y}\in\reals$ by \cref{real_abs_in_reals}.
                                $0\in\reals$ by \cref{real_add_ident}.
                                $\abs{x - x_0} \rmul \abs{y} = 0 \rmul \abs{y}$ by \cref{real_mul_eq_subst}.
                                $0 \rmul \abs{y} = 0$ by \cref{real_mul_zero}.
                                $\abs{x - x_0} \rmul \abs{y} = 0$.
                                $\delta$ is positive.
                                $\abs{y_0} + \delta$ is positive.
                                $\abs{y_0}\in\reals$ by \cref{real_abs_in_reals}.
                                $\abs{y_0} + \delta\in\reals$ by \cref{real_add_closed}.
                                $\delta \rmul (\abs{y_0} + \delta)$ is positive by \cref{real_mul_pos_pos}.
                                $0 < \delta \rmul (\abs{y_0} + \delta)$.
                                $\abs{x - x_0} \rmul \abs{y} < \delta \rmul (\abs{y_0} + \delta)$ by \cref{real_lt_subst_left}.
                            \caseOf{$0 < \abs{x - x_0}$.}
                                $\abs{y}\in\reals$ by \cref{real_abs_in_reals}.
                                $\abs{y_0}\in\reals$ by \cref{real_abs_in_reals}.
                                $\abs{x - x_0}\in\reals$ by \cref{real_abs_in_reals}.
                                $\abs{y_0} + \delta\in\reals$ by \cref{real_add_closed}.
                                $\abs{x - x_0} \rmul \abs{y}\in\reals$ by \cref{real_mul_closed}.
                                $\abs{x - x_0} \rmul (\abs{y_0} + \delta)\in\reals$ by \cref{real_mul_closed}.
                                $\delta \rmul (\abs{y_0} + \delta)\in\reals$ by \cref{real_mul_closed}.
                                $\abs{y} \rmul \abs{x - x_0} < (\abs{y_0} + \delta) \rmul \abs{x - x_0}$ by \cref{real_lt_mul_pos}.
                                $\abs{y} \rmul \abs{x - x_0} = \abs{x - x_0} \rmul \abs{y}$ by \cref{real_mul_comm}.
                                $(\abs{y_0} + \delta) \rmul \abs{x - x_0} = \abs{x - x_0} \rmul (\abs{y_0} + \delta)$ by \cref{real_mul_comm}.
                                $\abs{x - x_0} \rmul \abs{y} < \abs{x - x_0} \rmul (\abs{y_0} + \delta)$.
                                $\abs{x - x_0} \rmul (\abs{y_0} + \delta) < \delta \rmul (\abs{y_0} + \delta)$ by \cref{real_lt_mul_pos}.
                                $\abs{x - x_0} \rmul \abs{y} < \delta \rmul (\abs{y_0} + \delta)$ by \cref{real_lt_trans}.
                            \end{byCase}
                        \end{subproof}
                        $0 < \abs{x_0}$ or $0 = \abs{x_0}$ by \cref{real_abs_nonnegative,real_leq_split}.
                        Show that $\abs{x - x_0} \rmul \abs{y} + \abs{x_0} \rmul \abs{y - y_0} < \delta \rmul m$.
                        \begin{subproof}
                            \begin{byCase}
                            \caseOf{$0 = \abs{x_0}$.}
                                $\abs{x_0} = 0$ by assumption.
                                $\abs{x_0}\in\reals$ by \cref{real_abs_in_reals}.
                                $\abs{y - y_0}\in\reals$ by \cref{real_abs_in_reals}.
                                $0\in\reals$ by \cref{real_add_ident}.
                                $\abs{x_0} \rmul \abs{y - y_0} = 0 \rmul \abs{y - y_0}$ by \cref{real_mul_eq_subst}.
                                $0 \rmul \abs{y - y_0} = 0$ by \cref{real_mul_zero}.
                                $\abs{x_0} \rmul \abs{y - y_0} = 0$.
                                $\abs{x - x_0}\in\reals$ by \cref{real_abs_in_reals}.
                                $\abs{y}\in\reals$ by \cref{real_abs_in_reals}.
                                $\abs{x - x_0} \rmul \abs{y}\in\reals$ by \cref{real_mul_closed}.
                                $\abs{x - x_0} \rmul \abs{y} + \abs{x_0} \rmul \abs{y - y_0} = \abs{x - x_0} \rmul \abs{y} + 0$ by \cref{real_add_eq_subst}.
                                $\abs{x - x_0} \rmul \abs{y} + 0 = 0 + \abs{x - x_0} \rmul \abs{y}$ by \cref{real_add_comm}.
                                $0 + \abs{x - x_0} \rmul \abs{y} = \abs{x - x_0} \rmul \abs{y}$ by \cref{real_add_ident}.
                                $\abs{x - x_0} \rmul \abs{y} + \abs{x_0} \rmul \abs{y - y_0} = \abs{x - x_0} \rmul \abs{y}$.
                                $\abs{x - x_0} \rmul \abs{y} < \delta \rmul (\abs{y_0} + \delta)$.
                                $\delta + \abs{y_0} < 1 + \abs{y_0}$ by \cref{real_lt_add}.
                                $\abs{y_0} + \delta < \abs{y_0} + 1$ by \cref{real_add_comm,real_add_assoc}.
                                $\abs{y_0} + \delta\in\reals$ by \cref{real_add_closed}.
                                $\abs{y_0} + 1\in\reals$ by \cref{real_add_closed}.
                                $\delta\in\reals$ by \cref{real_mul_closed}.
                                $\delta$ is positive.
                                $\delta \rmul (\abs{y_0} + \delta) = (\abs{y_0} + \delta) \rmul \delta$ by \cref{real_mul_comm}.
                                $\delta \rmul (\abs{y_0} + 1) = (\abs{y_0} + 1) \rmul \delta$ by \cref{real_mul_comm}.
                                $(\abs{y_0} + \delta) \rmul \delta < (\abs{y_0} + 1) \rmul \delta$ by \cref{real_lt_mul_pos}.
                                $\delta \rmul (\abs{y_0} + \delta) < (\abs{y_0} + 1) \rmul \delta$ by \cref{real_lt_subst_left}.
                                $\delta \rmul (\abs{y_0} + \delta) < \delta \rmul (\abs{y_0} + 1)$ by \cref{real_lt_subst_right}.
                                $\delta \rmul (\abs{y_0} + \delta)\in\reals$ by \cref{real_mul_closed}.
                                $\delta \rmul (\abs{y_0} + 1)\in\reals$ by \cref{real_mul_closed}.
                                $\abs{x - x_0} \rmul \abs{y} < \delta \rmul (\abs{y_0} + 1)$ by \cref{real_lt_trans}.
                                $\abs{x_0} = 0$.
                                $\abs{x_0} + \abs{y_0} + 1 = \abs{y_0} + 1$ by \cref{real_add_ident,real_add_comm,real_add_assoc,real_add_eq_subst}.
                                $\delta \rmul (\abs{y_0} + 1) = (\abs{y_0} + 1) \rmul \delta$ by \cref{real_mul_comm}.
                                $(\abs{y_0} + 1) \rmul \delta = (\abs{x_0} + \abs{y_0} + 1) \rmul \delta$ by \cref{real_mul_eq_subst}.
                                $\delta \rmul (\abs{y_0} + 1) = \delta \rmul (\abs{x_0} + \abs{y_0} + 1)$ by \cref{real_mul_comm}.
                                $\abs{x - x_0} \rmul \abs{y} + \abs{x_0} \rmul \abs{y - y_0} < \delta \rmul (\abs{x_0} + \abs{y_0} + 1)$ by \cref{real_lt_subst_right}.
                                $\abs{x - x_0} \rmul \abs{y} + \abs{x_0} \rmul \abs{y - y_0} < \delta \rmul m$ by \cref{real_lt_subst_right}.
                            \caseOf{$0 < \abs{x_0}$.}
                                $\abs{y - y_0}\in\reals$ by \cref{real_abs_in_reals}.
                                $\abs{x_0}\in\reals$ by \cref{real_abs_in_reals}.
                                $\delta\in\reals$ by \cref{real_mul_closed}.
                                $\abs{y - y_0} \rmul \abs{x_0} < \delta \rmul \abs{x_0}$ by \cref{real_lt_mul_pos}.
                                $\abs{y - y_0} \rmul \abs{x_0} = \abs{x_0} \rmul \abs{y - y_0}$ by \cref{real_mul_comm}.
                                $\delta \rmul \abs{x_0} = \abs{x_0} \rmul \delta$ by \cref{real_mul_comm}.
                                $\abs{x_0} \rmul \abs{y - y_0} < \delta \rmul \abs{x_0}$ by \cref{real_lt_subst_left}.
                                $\abs{x_0} \rmul \abs{y - y_0} < \abs{x_0} \rmul \delta$ by \cref{real_lt_subst_right}.
                                $\abs{x - x_0}\in\reals$ by \cref{real_abs_in_reals}.
                                $\abs{y}\in\reals$ by \cref{real_abs_in_reals}.
                                $\abs{x - x_0} \rmul \abs{y}\in\reals$ by \cref{real_mul_closed}.
                                $\abs{y_0}\in\reals$ by \cref{real_abs_in_reals}.
                                $\abs{y_0} + \delta\in\reals$ by \cref{real_add_closed}.
                                $\delta \rmul (\abs{y_0} + \delta)\in\reals$ by \cref{real_mul_closed}.
                                $\abs{y - y_0}\in\reals$ by \cref{real_abs_in_reals}.
                                $\abs{x_0}\in\reals$ by \cref{real_abs_in_reals}.
                                $\abs{x_0} \rmul \abs{y - y_0}\in\reals$ by \cref{real_mul_closed}.
                                $\abs{x_0} \rmul \delta\in\reals$ by \cref{real_mul_closed}.
                                $\abs{x - x_0} \rmul \abs{y} + \abs{x_0} \rmul \abs{y - y_0} < \delta \rmul (\abs{y_0} + \delta) + \abs{x_0} \rmul \delta$ by \cref{real_lt_add_two}.
                                $(\abs{x_0} + \abs{y_0}) + \delta = \abs{x_0} + (\abs{y_0} + \delta)$ by \cref{real_add_assoc}.
                                $\delta \rmul (\abs{y_0} + \delta) = (\abs{y_0} + \delta) \rmul \delta$ by \cref{real_mul_comm}.
                                $(\abs{y_0} + \delta) \rmul \delta = (\abs{y_0} \rmul \delta) + (\delta \rmul \delta)$ by \cref{real_right_distrib}.
                                $\abs{x_0} \rmul \delta = \delta \rmul \abs{x_0}$ by \cref{real_mul_comm}.
                                $\abs{y_0} \rmul \delta = \delta \rmul \abs{y_0}$ by \cref{real_mul_comm}.
                                $\delta \rmul (\abs{y_0} + \delta) + \abs{x_0} \rmul \delta = (\delta \rmul \abs{x_0}) + (\delta \rmul \abs{y_0}) + (\delta \rmul \delta)$
                                    by \cref{real_add_eq_subst,real_right_distrib,real_mul_comm,real_add_assoc,real_add_comm}.
                                $(\delta \rmul \abs{x_0}) + (\delta \rmul \abs{y_0}) = \delta \rmul (\abs{x_0} + \abs{y_0})$
                                    by \cref{real_mul_comm,real_right_distrib,real_add_eq_subst}.
                                $\delta \rmul (\abs{y_0} + \delta) + \abs{x_0} \rmul \delta = (\delta \rmul (\abs{x_0} + \abs{y_0})) + (\delta \rmul \delta)$ by \cref{real_add_assoc}.
                                $(\delta \rmul (\abs{x_0} + \abs{y_0})) + (\delta \rmul \delta) = \delta \rmul ((\abs{x_0} + \abs{y_0}) + \delta)$ by \cref{real_abs_in_reals,real_add_closed,real_mul_closed,real_mul_comm,real_right_distrib,real_add_eq_subst}.
                                $\delta \rmul (\abs{y_0} + \delta) + \abs{x_0} \rmul \delta = \delta \rmul ((\abs{x_0} + \abs{y_0}) + \delta)$.
                                $\abs{x_0} + \abs{y_0} = \abs{y_0} + \abs{x_0}$ by \cref{real_add_comm}.
                                $(\abs{x_0} + \abs{y_0}) + \delta = (\abs{y_0} + \abs{x_0}) + \delta$ by \cref{real_add_eq_subst}.
                                $\delta \rmul ((\abs{x_0} + \abs{y_0}) + \delta) = \delta \rmul (\abs{x_0} + \abs{y_0} + \delta)$ by \cref{real_add_assoc}.
                                $\abs{x - x_0} \rmul \abs{y} + \abs{x_0} \rmul \abs{y - y_0} < \delta \rmul (\abs{x_0} + \abs{y_0} + \delta)$ by \cref{real_lt_subst_right}.
                                $\delta + (\abs{x_0} + \abs{y_0}) < 1 + (\abs{x_0} + \abs{y_0})$ by \cref{real_lt_add}.
                                $\abs{x_0} + \abs{y_0} + \delta < \abs{x_0} + \abs{y_0} + 1$ by \cref{real_add_comm,real_add_assoc}.
                                $\abs{x_0}\in\reals$ by \cref{real_abs_in_reals}.
                                $\abs{y_0}\in\reals$ by \cref{real_abs_in_reals}.
                                $\abs{x_0} + \abs{y_0}\in\reals$ by \cref{real_add_closed}.
                                $\abs{x_0} + \abs{y_0} + \delta\in\reals$ by \cref{real_add_closed}.
                                $1\in\reals$ by \cref{real_mul_ident}.
                                $\abs{x_0} + \abs{y_0} + 1\in\reals$ by \cref{real_add_closed}.
                                $(\abs{x_0} + \abs{y_0} + \delta) \rmul \delta < (\abs{x_0} + \abs{y_0} + 1) \rmul \delta$ by \cref{real_lt_mul_pos}.
                                $(\abs{x_0} + \abs{y_0} + \delta) \rmul \delta = \delta \rmul (\abs{x_0} + \abs{y_0} + \delta)$ by \cref{real_mul_comm}.
                                $\delta \rmul (\abs{x_0} + \abs{y_0} + \delta) < (\abs{x_0} + \abs{y_0} + 1) \rmul \delta$ by \cref{real_lt_subst_left}.
                                $(\abs{x_0} + \abs{y_0} + 1) \rmul \delta = \delta \rmul (\abs{x_0} + \abs{y_0} + 1)$ by \cref{real_mul_comm}.
                                $\delta \rmul (\abs{x_0} + \abs{y_0} + \delta) < \delta \rmul (\abs{x_0} + \abs{y_0} + 1)$ by \cref{real_lt_subst_right}.
                                $\abs{x - x_0} \rmul \abs{y}\in\reals$ by \cref{real_mul_closed}.
                                $\abs{x_0} \rmul \abs{y - y_0}\in\reals$ by \cref{real_mul_closed}.
                                $\abs{x - x_0} \rmul \abs{y} + \abs{x_0} \rmul \abs{y - y_0}\in\reals$ by \cref{real_add_closed}.
                                $\delta \rmul (\abs{x_0} + \abs{y_0} + \delta)\in\reals$ by \cref{real_mul_closed}.
                                $\delta \rmul (\abs{x_0} + \abs{y_0} + 1)\in\reals$ by \cref{real_mul_closed}.
                                $\abs{x - x_0} \rmul \abs{y} + \abs{x_0} \rmul \abs{y - y_0} < \delta \rmul (\abs{x_0} + \abs{y_0} + 1)$ by \cref{real_lt_trans}.
                                $\abs{x - x_0} \rmul \abs{y} + \abs{x_0} \rmul \abs{y - y_0} < \delta \rmul m$ by \cref{real_lt_subst_right}.
                            \end{byCase}
                        \end{subproof}
                        $m < \epsilon + m$ by \cref{real_lt_add,real_add_ident,real_lt_subst_left}.
                        $m < m + \epsilon$ by \cref{real_add_comm,real_lt_subst_right}.
                        $m \rmul \inv{(m + \epsilon)} < (m + \epsilon) \rmul \inv{(m + \epsilon)}$ by \cref{real_lt_mul_pos}.
                        $(m + \epsilon) \rmul \inv{(m + \epsilon)} = 1$ by \cref{real_mul_inverse}.
                        $m \rmul \inv{(m + \epsilon)} < 1$ by \cref{real_lt_subst_right}.
                        $m \rmul \inv{(m + \epsilon)}\in\reals$ by \cref{real_mul_closed}.
                        $1\in\reals$ by \cref{real_mul_ident}.
                        $(m \rmul \inv{(m + \epsilon)}) \rmul \epsilon < 1 \rmul \epsilon$ by \cref{real_lt_mul_pos}.
                        $(m \rmul \inv{(m + \epsilon)}) \rmul \epsilon = \epsilon \rmul (m \rmul \inv{(m + \epsilon)})$ by \cref{real_mul_comm}.
                        $\epsilon \rmul (m \rmul \inv{(m + \epsilon)}) < 1 \rmul \epsilon$ by \cref{real_lt_subst_left}.
                        $1 \rmul \epsilon = \epsilon \rmul 1$ by \cref{real_mul_comm}.
                        $\epsilon \rmul (m \rmul \inv{(m + \epsilon)}) < \epsilon \rmul 1$ by \cref{real_lt_subst_right}.
                        $\epsilon \rmul 1 = \epsilon$ by \cref{real_mul_comm,real_mul_ident}.
                        $\delta = \epsilon \rmul \inv{(m + \epsilon)}$ by \cref{real_div}.
                        $\delta \rmul m = \epsilon \rmul (m \rmul \inv{(m + \epsilon)})$ by \cref{real_div,real_mul_assoc,real_mul_comm}.
                        $\delta \rmul m < \epsilon$ by \cref{real_lt_subst_left}.
                        $x - x_0\in\reals$ and $y - y_0\in\reals$ by \cref{real_sub,real_add_closed}.
                        $\abs{x - x_0}\in\reals$ and $\abs{y}\in\reals$ and $\abs{y - y_0}\in\reals$ by \cref{real_abs_in_reals}.
                        $\abs{x - x_0} \rmul \abs{y}\in\reals$ by \cref{real_mul_closed}.
                        $\abs{x_0} \rmul \abs{y - y_0}\in\reals$ by \cref{real_mul_closed}.
                        $\abs{x - x_0} \rmul \abs{y} + \abs{x_0} \rmul \abs{y - y_0}\in\reals$ by \cref{real_add_closed}.
                        $\delta \rmul m\in\reals$ by \cref{real_mul_closed}.
                        $\abs{x - x_0} \rmul \abs{y} + \abs{x_0} \rmul \abs{y - y_0} < \epsilon$ by \cref{real_lt_trans}.
                        $(x - x_0) \rmul y + x_0 \rmul (y - y_0)\in\reals$ and
                            $\abs{(x - x_0) \rmul y + x_0 \rmul (y - y_0)}\in\reals$ by
                            \cref{real_add_closed,real_abs_in_reals}.
                        $\abs{(x - x_0) \rmul y + x_0 \rmul (y - y_0)} < \epsilon$
                            by \cref{real_leq_split,real_lt_trans,real_lt_subst_left}.
                        $\abs{(x \rmul y) - (x_0 \rmul y_0)} < \epsilon$ by \cref{real_lt_subst_left}.
                        $\abs{f(q) - f(p)} < \epsilon$ by \cref{real_lt_subst_left}.
                \end{subproof}
                $f$ is a function to $\reals$ and $\dom{f} = \reals\times\reals$ by \cref{funs_elim}.
                $f(p)\in\reals$ and $f(q)\in\reals$.
                $d((f(p),f(q))) < \epsilon$ by \cref{standard_metric_value,real_abs_sub_swap,real_lt_subst_left}.
                $f(q)\in\metricBall{d}{\reals}{f(p)}{\epsilon}$ by \cref{metric_ball}.
            \end{subproof}
            $\img{f}{U}\subseteq \metricBall{d}{\reals}{f(p)}{\epsilon}$ by \cref{img_iff,function_apply_intro,subseteq}.
            $\img{f}{U}\subseteq V$ by \cref{subseteq,subseteq_transitive}.
            $U$ is a neighborhood of $p$ in $\reals\times\reals$ with respect to $T$ by \cref{neighborhood}.
        \end{subproof}
        $f$ is continuous at $p$ from $\reals\times\reals$ to $\reals$ with respect to $T$ and $\standardTopologyR$ by
            \cref{standard_basis_is_basis,basis_topology_is_topology,standard_topology_reals,product_basis_is_basis,product_topology,continuous_at}.
    \end{subproof}
    Show that $f$ is continuous from $\reals\times\reals$ to $\reals$ with respect to $T$ and $\standardTopologyR$.
    \begin{subproof}
        Follows by
            \cref{standard_basis_is_basis,basis_topology_is_topology,standard_topology_reals,product_basis_is_basis,product_topology,continuous_at_implies_continuous}.
    \end{subproof}
\end{proof}

% from §21 helper definition 5: reciprocal operation on reals
\begin{definition}\label{real_inverse_operation}
    $\realInverse = \{w \mid \exists y\in\reals\setminus\{0\}. \exists z\in\reals.
        w = (y,z) \land z = \inv{y}\}$.
\end{definition}
% from §21 helper lemma 7: reciprocal operation is continuous
\begin{lemma}\label{real_inverse_continuous}
    Let $T_1 = \standardTopologyR$.
    Let $T_2 = \subspaceTopology{T_1}{\reals\setminus\{0\}}$.
    Then $\realInverse$ is continuous from $\reals\setminus\{0\}$ to $\reals$
        with respect to $T_2$ and $T_1$.
\end{lemma}
\begin{proof}
    Let $f = \realInverse$.
    For all $p,z_1,z_2$ such that $(p,z_1)\in f$ and $(p,z_2)\in f$ we have $z_1 = z_2$
        by \cref{real_inverse_operation,pair_eq_iff}.
    $f$ is a function by \cref{rightunique,real_inverse_operation,relation}.
    For all $p\in\dom{f}$ we have $f(p)\in\reals$
        by \cref{dom_iff,real_inverse_operation,pair_eq_iff,function_apply_intro}.
    $f$ is a function to $\reals$.
    $\dom{f}\subseteq\reals\setminus\{0\}$ by \cref{dom_iff,real_inverse_operation,pair_eq_iff,subseteq}.
    Show that for all $p\in\reals\setminus\{0\}$ we have $p\in\dom{f}$.
    \begin{subproof}
        Follows by \cref{setminus_elim_left,setminus_elim_right,singleton_iff,real_mul_inverse,real_inverse_operation,dom_intro}.
    \end{subproof}
    $\reals\setminus\{0\}\subseteq\dom{f}$ by \cref{subseteq}.
    Show that $f\in\funs{\reals\setminus\{0\}}{\reals}$.
    \begin{subproof}
        Follows by \cref{subseteq_antisymmetric,funs_intro}.
    \end{subproof}
    Show that for all $p\in\reals\setminus\{0\}$ we have
        $f$ is continuous at $p$ from $\reals\setminus\{0\}$ to $\reals$ with respect to $T_2$ and $T_1$.
        \begin{subproof}
            Fix $p$ such that $p\in\reals\setminus\{0\}$.
            $p\in\reals$ and $p\notin\{0\}$ by \cref{setminus_elim_left,setminus_elim_right}.
            $p\neq 0$ and $\inv{p}\in\reals$ by \cref{singleton_iff,real_mul_inverse}.
        $f(p) = \inv{p}$ by \cref{real_inverse_operation,function_apply_intro}.
        Show that for all $V$ such that $V$ is a neighborhood of $f(p)$ in $\reals$ with respect to $T_1$
            there exists $U$ such that $U$ is a neighborhood of $p$ in $\reals\setminus\{0\}$ with respect to $T_2$
            and $\img{f}{U}\subseteq V$.
        \begin{subproof}
            Fix $V$ such that $V$ is a neighborhood of $f(p)$ in $\reals$ with respect to $T_1$.
            $V$ is open in $\reals$ with respect to $T_1$ by \cref{neighborhood}.
            Let $d = \standardMetricR$.
            Let $T_d = \metricTopology{d}{\reals}$.
            $T_d = \standardTopologyR$ by \cref{standard_metric_topology}.
            $V\in T_d$ by \cref{open_in,subseteq}.
            $f(p)\in V$ by \cref{neighborhood}.
            $\standardBasisR$ is a basis on $\reals$ and $\standardTopologyR$ is a topology on $\reals$
                by \cref{standard_basis_is_basis,basis_topology_is_topology,standard_topology_reals}.
            $V\subseteq\reals$ by \cref{topology_on,subseteq,pow_iff}.
            $d$ is a metric on $\reals$ by \cref{standard_metric_is_metric}.
            Take $\epsilon\in\reals$ such that $0 < \epsilon$ and
                $\metricBall{d}{\reals}{f(p)}{\epsilon} \subseteq V$ by \cref{metric_open_iff}.
            $\abs{p}\in\reals$ by \cref{real_abs_in_reals}.
            $\abs{p} > 0$ by \cref{real_abs_eq_zero,real_lt_trichotomy,real_abs_nonnegative}.
            Let $a = \abs{p} \rmul \inv{(1 + 1)}$.
            $1\in\reals$ and $0\in\reals$ and $0 < 1$ and $1 + 1\in\reals$
                by \cref{real_mul_ident,real_add_ident,real_zero_lt_one,real_add_closed}.
            $0 + 1 < 1 + 1$ and $0 + 1 = 1$ by \cref{real_lt_add,real_add_ident}.
            $1 < 1 + 1$ by \cref{real_lt_subst_left}.
            $0 < 1 + 1$ by \cref{real_lt_trans}.
            $(1 + 1)\neq 0$ by \cref{real_pos_nonzero}.
            $\inv{(1 + 1)}\in\reals$ by \cref{real_mul_inverse}.
            $0 < \inv{(1 + 1)}$ by \cref{real_inv_pos}.
            $0 \rmul \inv{(1 + 1)} < \abs{p} \rmul \inv{(1 + 1)}$ by \cref{real_lt_mul_pos}.
            $0 \rmul \inv{(1 + 1)} < a$ by \cref{real_lt_subst_right}.
            $0 \rmul \inv{(1 + 1)} = 0$ by \cref{real_mul_zero}.
            $0 < a$ by \cref{real_lt_subst_left}.
            Let $b = (\epsilon \rmul (\abs{p} \rmul \abs{p})) \rmul \inv{(1 + 1)}$.
            $\abs{p} \rmul \abs{p}\in\reals$ and $\epsilon \rmul (\abs{p} \rmul \abs{p})\in\reals$ and $(1 + 1)\in\reals$
                by \cref{real_mul_closed,real_add_closed}.
            $\abs{p} \rmul \abs{p} > 0$ by \cref{real_mul_pos_pos_gt}.
            $0 \rmul (\abs{p} \rmul \abs{p}) < \epsilon \rmul (\abs{p} \rmul \abs{p})$ by \cref{real_lt_mul_pos}.
            $0 \rmul (\abs{p} \rmul \abs{p}) = 0$ by \cref{real_mul_zero}.
            $0 < \epsilon \rmul (\abs{p} \rmul \abs{p})$ by \cref{real_lt_subst_left}.
            $0 \rmul \inv{(1 + 1)} < (\epsilon \rmul (\abs{p} \rmul \abs{p})) \rmul \inv{(1 + 1)}$ by \cref{real_lt_mul_pos}.
            $0 \rmul \inv{(1 + 1)} < b$ by \cref{real_lt_subst_right}.
            $0 \rmul \inv{(1 + 1)} = 0$ by \cref{real_mul_zero}.
            $0 < b$ by \cref{real_lt_subst_left}.
            $a\in\reals$ and $b\in\reals$ by \cref{real_mul_closed}.
            Let $E = \{a,b\}$.
            For all $x$ such that $x\in E$ we have $x\in\reals$.
            $E\subseteq\reals$ by \cref{subseteq}.
            $E$ is finite by \cref{pair_is_finite}.
            $E\neq\emptyset$ by \cref{upair_intro_left,emptyset}.
            Take $\delta\in E$ such that $\delta$ is a minimum of $E$ by \cref{finite_real_minimum}.
            $\delta\in\reals$ by \cref{subseteq}.
            $a\in E$ and $b\in E$ by \cref{upair_intro_left,upair_intro_right}.
            $0 < \delta$ by \cref{upair_elim,real_lt_subst_right}.
            $\delta \leq a$ and $\delta \leq b$ by \cref{real_minimum}.
            Let $U = \metricBall{d}{\reals}{p}{\delta}$.
            Show that for all $y$ such that $y\in U$ we have $y\in\reals\setminus\{0\}$.
            \begin{subproof}
                Fix $y$ such that $y\in U$.
                $y\in\reals$ by \cref{metric_ball}.
                $d((p,y)) < \delta$ by \cref{metric_ball}.
                $d((p,y)) = \abs{p - y}$ by \cref{standard_metric_value}.
                $\abs{p - y} < \delta$.
                Show that $y\neq 0$.
                \begin{subproof}
                    $\neg{y}\in\reals$ by \cref{real_add_inverse}.
                    $p - y\in\reals$ by \cref{real_sub,real_add_closed}.
                    $\abs{p - y}\in\reals$ by \cref{real_abs_in_reals}.
                    $\abs{p - y} < a$ by \cref{real_leq_split,real_lt_trans,real_lt_subst_right}.
                    $a + a = \abs{p}$ by
                        \cref{real_right_distrib,real_mul_ident,real_mul_eq_subst,real_mul_assoc,real_mul_comm,real_mul_inverse}.
                    $\abs{p} < \abs{y} + a$ by \cref{real_abs_sub_lt_bound}.
                    Show that it is wrong that $\abs{y} \leq a$.
                    \begin{subproof}
                        Suppose not.
                        $\abs{y} \leq a$.
                        $\abs{y}\in\reals$ by \cref{real_abs_in_reals}.
                        $\abs{y} + a \leq a + a$ by \cref{real_leq_add}.
                        $\abs{y} + a\in\reals$ by \cref{real_add_closed}.
                        $\abs{y} + a < a + a$ or $\abs{y} + a = a + a$ by \cref{real_leq_split}.
                        \begin{byCase}
                        \caseOf{$\abs{y} + a < a + a$.}
                            $\abs{y} + a < \abs{p}$ by \cref{real_lt_subst_right}.
                            $\abs{p} < \abs{p}$ by \cref{real_lt_trans}.
                            Contradiction by \cref{real_lt_irreflexive}.
                        \caseOf{$\abs{y} + a = a + a$.}
                            $\abs{y} + a = \abs{p}$.
                            $\abs{p} < \abs{p}$ by \cref{real_lt_subst_right}.
                            Contradiction by \cref{real_lt_irreflexive}.
                        \end{byCase}
                    \end{subproof}
                    $\abs{y} > a$.
                    $\abs{y}\neq 0$ by \cref{real_lt_trichotomy,real_abs_nonnegative}.
                    Show that $y\neq 0$.
                    \begin{subproof}
                        Suppose not.
                        $y = 0$.
                        $0 < y$ or $0 = y$.
                        $0 \leq y$.
                        $\abs{y} = y$ by \cref{real_abs_nonneg}.
                        $\abs{y} = 0$.
                        Contradiction.
                    \end{subproof}
                \end{subproof}
                $y\in\reals\setminus\{0\}$ by \cref{singleton_iff,setminus_intro}.
            \end{subproof}
            $U\subseteq\reals\setminus\{0\}$ by \cref{subseteq}.
            $p\in\reals$ and $\delta\in\reals$ and $0 < \delta$.
            For all $y$ such that $y\in U$ we have $y\in\reals$ by \cref{metric_ball}.
            $U\in\metricBasis{d}{\reals}$ by \cref{subseteq,powerset_intro,metric_basis}.
            $U\in T_d$ by \cref{metric_basis_is_basis,basis_elem_open_in_generated,basis_for_topology,metric_topology}.
            $T_d\subseteq\standardTopologyR$ by \cref{subseteq}.
            $U\in\standardTopologyR$ by \cref{subseteq}.
            $\standardBasisR$ is a basis on $\reals$ and $\standardTopologyR$ is a topology on $\reals$
                by \cref{standard_basis_is_basis,basis_topology_is_topology,standard_topology_reals}.
            $U$ is open in $\reals$ with respect to $\standardTopologyR$ by \cref{open_in}.
            $U\in T_1$ by \cref{open_in}.
            $U\subseteq(\reals\setminus\{0\})\inter U$ by \cref{subseteq,subseteq_inter_iff}.
            $(\reals\setminus\{0\})\inter U\subseteq U$ by \cref{inter_lower_right}.
            $U = (\reals\setminus\{0\})\inter U$ by \cref{subseteq_antisymmetric}.
            $U\in T_2$ by \cref{subspace_topology}.
            $\standardBasisR$ is a basis on $\reals$ and $T_1$ is a topology on $\reals$
                by \cref{standard_basis_is_basis,basis_topology_is_topology,standard_topology_reals}.
            $\reals\setminus\{0\}\subseteq\reals$ and $T_2$ is a topology on $\reals\setminus\{0\}$
                by \cref{setminus_elim_left,subseteq,subspace_topology_is_topology}.
            $U$ is open in $\reals\setminus\{0\}$ with respect to $T_2$ by \cref{open_in}.
            $d((p,p)) = 0$ by \cref{metric_zero_characterization,standard_metric_is_metric,metric_on}.
            $d((p,p)) < \delta$ by \cref{real_lt_subst_left}.
            $p\in U$ by \cref{metric_ball}.
            $U$ is a neighborhood of $p$ in $\reals\setminus\{0\}$ with respect to $T_2$ by \cref{neighborhood}.
            Show that for all $q$ such that $q\in U$ we have
                $f(q)\in\metricBall{d}{\reals}{f(p)}{\epsilon}$.
            \begin{subproof}
                Fix $q$ such that $q\in U$.
                $q\in\reals$ by \cref{metric_ball}.
                $d((p,q)) < \delta$ by \cref{metric_ball}.
                $d((p,q)) = \abs{p - q}$ by \cref{standard_metric_value}.
                $\abs{p - q} < \delta$.
                $q\neq 0$ by \cref{subseteq,setminus_elim_right,metric_ball,singleton_iff}.
                $q\in\reals\setminus\{0\}$ by \cref{setminus_intro}.
                $\inv{q}\in\reals$ by \cref{real_mul_inverse}.
                $f(q) = \inv{q}$ by \cref{real_inverse_operation,function_apply_intro}.
                    $q \rmul p \neq 0$ by \cref{real_mul_zero_product}.
                    $q \rmul p\in\reals$ by \cref{real_mul_closed}.
                    $\inv{(q \rmul p)}\in\reals$ by \cref{real_mul_inverse}.
                    $q \rmul (p \rmul \inv{(q \rmul p)}) = 1$ by \cref{real_mul_assoc,real_mul_inverse}.
                    $p \rmul \inv{(q \rmul p)}\in\reals$ by \cref{real_mul_closed}.
                    $\inv{q} = p \rmul \inv{(q \rmul p)}$ by \cref{real_mul_inverse_unique}.
                    $p \rmul (q \rmul \inv{(q \rmul p)}) = 1$ by \cref{real_mul_assoc,real_mul_comm,real_mul_inverse}.
                    $q \rmul \inv{(q \rmul p)}\in\reals$ by \cref{real_mul_closed}.
                    $\inv{p} = q \rmul \inv{(q \rmul p)}$ by \cref{real_mul_inverse_unique}.
                        $\inv{q} - \inv{p} = (p \rmul \inv{(q \rmul p)}) - (q \rmul \inv{(q \rmul p)})$ by \cref{real_lt_subst_left}.
                        $(p \rmul \inv{(q \rmul p)}) - (q \rmul \inv{(q \rmul p)}) =
                            (p \rmul \inv{(q \rmul p)}) + \neg{(q \rmul \inv{(q \rmul p)})}$ by \cref{real_sub}.
                        Show that $\neg{(q \rmul \inv{(q \rmul p)})} = (\neg{q}) \rmul \inv{(q \rmul p)}$.
                        \begin{subproof}
                            $\neg{(q \rmul \inv{(q \rmul p)})} = \neg{1} \rmul (q \rmul \inv{(q \rmul p)})$ by \cref{real_neg_as_mul}.
                            $q \rmul \inv{(q \rmul p)} = \inv{(q \rmul p)} \rmul q$ by \cref{real_mul_comm}.
                            $\neg{1} \rmul (q \rmul \inv{(q \rmul p)}) = \neg{1} \rmul (\inv{(q \rmul p)} \rmul q)$ by \cref{real_mul_eq_subst}.
                            $\neg{(q \rmul \inv{(q \rmul p)})} = \neg{1} \rmul (\inv{(q \rmul p)} \rmul q)$ by \cref{real_mul_eq_subst}.
                            $\neg{1} \rmul (\inv{(q \rmul p)} \rmul q) = \neg{(\inv{(q \rmul p)} \rmul q)}$ by \cref{real_neg_as_mul}.
                            $\inv{(q \rmul p)} \rmul \neg{q} = \neg{(\inv{(q \rmul p)} \rmul q)}$ by \cref{real_mul_neg_right}.
                            $\neg{q}\in\reals$ by \cref{real_add_inverse}.
                            $\inv{(q \rmul p)} \rmul \neg{q} = (\neg{q}) \rmul \inv{(q \rmul p)}$ by \cref{real_mul_comm}.
                            $\neg{(q \rmul \inv{(q \rmul p)})} = (\neg{q}) \rmul \inv{(q \rmul p)}$.
                        \end{subproof}
                        $(p \rmul \inv{(q \rmul p)}) + \neg{(q \rmul \inv{(q \rmul p)})} =
                            (p \rmul \inv{(q \rmul p)}) + ((\neg{q}) \rmul \inv{(q \rmul p)})$ by \cref{real_add_eq_subst}.
                        $\neg{q}\in\reals$ by \cref{real_add_inverse}.
                        $(p \rmul \inv{(q \rmul p)}) + ((\neg{q}) \rmul \inv{(q \rmul p)}) =
                            (p + \neg{q}) \rmul \inv{(q \rmul p)}$ by \cref{real_right_distrib}.
                        $p + \neg{q} = p - q$ by \cref{real_sub}.
                        $(p \rmul \inv{(q \rmul p)}) + ((\neg{q}) \rmul \inv{(q \rmul p)}) =
                            (p - q) \rmul \inv{(q \rmul p)}$ by \cref{real_mul_eq_subst}.
                        $(p \rmul \inv{(q \rmul p)}) - (q \rmul \inv{(q \rmul p)}) = (p - q) \rmul \inv{(q \rmul p)}$
                            by \cref{real_add_eq_subst}.
                        $\abs{\inv{q} - \inv{p}} = \abs{(p - q) \rmul \inv{(q \rmul p)}}$ by \cref{real_lt_subst_left}.
                        $\neg{q}\in\reals$ by \cref{real_add_inverse}.
                        $p - q\in\reals$ by \cref{real_sub,real_add_closed}.
                        $\abs{(p - q) \rmul \inv{(q \rmul p)}} = \abs{p - q} \rmul \abs{\inv{(q \rmul p)}}$ by \cref{real_abs_mul}.
                        $\abs{p - q} = \abs{q - p}$ by \cref{real_abs_sub_swap}.
                        Show that $\abs{p - q} < \epsilon \rmul \abs{q \rmul p}$.
                        \begin{subproof}
                            $\neg{q}\in\reals$ by \cref{real_add_inverse}.
                            $p - q\in\reals$ by \cref{real_sub,real_add_closed}.
                            $\abs{p - q}\in\reals$ by \cref{real_abs_in_reals}.
                            $a + a = \abs{p}$ by
                                \cref{real_right_distrib,real_mul_ident,real_mul_eq_subst,real_mul_assoc,real_mul_comm,real_mul_inverse}.
                            Show that it is wrong that $\abs{q} \leq a$.
                            \begin{subproof}
                                Suppose not.
                                $\abs{q} \leq a$.
                                $\abs{q}\in\reals$ by \cref{real_abs_in_reals}.
                                $\abs{p - q} < a$ by \cref{real_leq_split,real_lt_trans,real_lt_subst_right}.
                                $\abs{p} < \abs{q} + a$ by \cref{real_abs_sub_lt_bound}.
                                $\abs{q} + a \leq a + a$ by \cref{real_leq_add}.
                                $\abs{q} + a\in\reals$ by \cref{real_add_closed}.
                                $\abs{q} + a < a + a$ or $\abs{q} + a = a + a$ by \cref{real_leq_split}.
                                \begin{byCase}
                                \caseOf{$\abs{q} + a < a + a$.}
                                    $\abs{q} + a < \abs{p}$ by \cref{real_lt_subst_right}.
                                    $\abs{p} < \abs{p}$ by \cref{real_lt_trans}.
                                    Contradiction by \cref{real_lt_irreflexive}.
                                \caseOf{$\abs{q} + a = a + a$.}
                                    $\abs{q} + a = \abs{p}$.
                                    $\abs{p} < \abs{p}$ by \cref{real_lt_subst_right}.
                                    Contradiction by \cref{real_lt_irreflexive}.
                                \end{byCase}
                            \end{subproof}
                            $\abs{q} > a$.
                            $\abs{p}\in\reals$ by \cref{real_abs_in_reals}.
                            $0 < \abs{p}$ by \cref{real_abs_eq_zero,real_abs_nonnegative,real_lt_trichotomy}.
                            $a\in\reals$ by \cref{real_mul_closed}.
                            $\abs{q}\in\reals$ by \cref{real_abs_in_reals}.
                            $a \rmul \abs{p} < \abs{q} \rmul \abs{p}$ by \cref{real_lt_mul_pos}.
                            $\abs{q \rmul p} = \abs{q} \rmul \abs{p}$ by \cref{real_abs_mul}.
                            $\abs{q \rmul p} > a \rmul \abs{p}$ by \cref{real_lt_subst_left}.
                            $\abs{p - q} < b$ by \cref{real_leq_split,real_lt_trans,real_lt_subst_right}.
                            $a \rmul \abs{p} = (\abs{p} \rmul \abs{p}) \rmul \inv{(1 + 1)}$
                                by \cref{real_mul_eq_subst,real_mul_assoc,real_mul_comm}.
                            $\epsilon \rmul (a \rmul \abs{p}) = (\epsilon \rmul (\abs{p} \rmul \abs{p})) \rmul \inv{(1 + 1)}$
                                by \cref{real_mul_assoc,real_mul_comm}.
                            $b = \epsilon \rmul (a \rmul \abs{p})$.
                            $\abs{q \rmul p} = \abs{q} \rmul \abs{p}$ by \cref{real_abs_mul}.
                            $a \rmul \abs{p} < \abs{q \rmul p}$ by \cref{real_lt_subst_right}.
                            $a \rmul \abs{p}\in\reals$ by \cref{real_mul_closed}.
                            $\abs{q \rmul p}\in\reals$ by \cref{real_abs_in_reals}.
                            $(a \rmul \abs{p}) \rmul \epsilon < \abs{q \rmul p} \rmul \epsilon$ by \cref{real_lt_mul_pos}.
                            $\epsilon \rmul (a \rmul \abs{p}) = (a \rmul \abs{p}) \rmul \epsilon$ by \cref{real_mul_comm}.
                            $\epsilon \rmul \abs{q \rmul p} = \abs{q \rmul p} \rmul \epsilon$ by \cref{real_mul_comm}.
                            $\epsilon \rmul (a \rmul \abs{p}) < \abs{q \rmul p} \rmul \epsilon$ by \cref{real_lt_subst_left}.
                            $\epsilon \rmul (a \rmul \abs{p}) < \epsilon \rmul \abs{q \rmul p}$ by \cref{real_lt_subst_right}.
                            $b < \epsilon \rmul \abs{q \rmul p}$ by \cref{real_lt_subst_left}.
                            $\epsilon \rmul \abs{q \rmul p}\in\reals$ by \cref{real_mul_closed}.
                            $\abs{p - q} < \epsilon \rmul \abs{q \rmul p}$ by \cref{real_lt_trans}.
                        \end{subproof}
                        $\inv{(q \rmul p)} \rmul (q \rmul p) = 1$ by \cref{real_mul_inverse,real_mul_comm}.
                        $\inv{(q \rmul p)} \neq 0$ by \cref{real_mul_unit_left_nonzero}.
                        $\abs{\inv{(q \rmul p)}}\neq 0$ by \cref{real_abs_eq_zero}.
                        $\abs{\inv{(q \rmul p)}} \geq 0$ by \cref{real_abs_nonnegative}.
                        $0 < \abs{\inv{(q \rmul p)}}$ by \cref{real_lt_trichotomy}.
                        $\abs{p - q}\in\reals$ by \cref{real_abs_in_reals}.
                        $\abs{\inv{(q \rmul p)}}\in\reals$ by \cref{real_abs_in_reals}.
                        $\epsilon\in\reals$.
                        $\abs{q \rmul p}\in\reals$ by \cref{real_abs_in_reals}.
                        $\epsilon \rmul \abs{q \rmul p}\in\reals$ by \cref{real_mul_closed}.
                        $\abs{p - q} \rmul \abs{\inv{(q \rmul p)}} < \epsilon \rmul \abs{q \rmul p} \rmul \abs{\inv{(q \rmul p)}}$
                            by \cref{real_lt_mul_pos}.
                        $\abs{q \rmul p} \rmul \abs{\inv{(q \rmul p)}} = \abs{(q \rmul p) \rmul \inv{(q \rmul p)}}$ by \cref{real_abs_mul,real_mul_assoc}.
                        $(q \rmul p) \rmul \inv{(q \rmul p)} = 1$ by \cref{real_mul_inverse}.
                        $\abs{1} = 1$ by \cref{real_abs_nonneg,real_zero_lt_one,real_lt_imp_leq}.
                        $\abs{q \rmul p} \rmul \abs{\inv{(q \rmul p)}} = 1$ by \cref{real_lt_subst_right}.
                        $\epsilon \rmul \abs{q \rmul p} \rmul \abs{\inv{(q \rmul p)}} =
                            \epsilon \rmul (\abs{q \rmul p} \rmul \abs{\inv{(q \rmul p)}})$ by \cref{real_mul_assoc}.
                        $\epsilon \rmul \abs{q \rmul p} \rmul \abs{\inv{(q \rmul p)}} = \epsilon \rmul 1$ by \cref{real_mul_eq_subst}.
                        $\epsilon \rmul 1 = 1 \rmul \epsilon$ by \cref{real_mul_comm}.
                        $\epsilon \rmul 1 = \epsilon$ by \cref{real_mul_ident}.
                        $\epsilon \rmul \abs{q \rmul p} \rmul \abs{\inv{(q \rmul p)}} = \epsilon$ by \cref{real_mul_eq_subst}.
                        $\abs{p - q} \rmul \abs{\inv{(q \rmul p)}} < \epsilon$ by \cref{real_lt_subst_right}.
                    $\abs{\inv{q} - \inv{p}} < \epsilon$ by \cref{real_lt_subst_left}.
                $\abs{f(q) - f(p)} < \epsilon$ by \cref{real_lt_subst_left,real_div}.
                $f(p)\in\reals$ and $f(q)\in\reals$ by \cref{funs_type_apply}.
                $d((f(p),f(q))) < \epsilon$ by \cref{standard_metric_value,real_abs_sub_swap,real_lt_subst_left}.
                $f(q)\in\metricBall{d}{\reals}{f(p)}{\epsilon}$ by \cref{metric_ball}.
            \end{subproof}
            $\img{f}{U}\subseteq \metricBall{d}{\reals}{f(p)}{\epsilon}$ by \cref{img_iff,function_apply_intro,subseteq}.
            $\metricBall{d}{\reals}{f(p)}{\epsilon}\subseteq V$ by \cref{subseteq}.
            $\img{f}{U}\subseteq V$ by \cref{subseteq_transitive}.
        \end{subproof}
        $\standardBasisR$ is a basis on $\reals$ and $T_1$ is a topology on $\reals$
            by \cref{standard_basis_is_basis,basis_topology_is_topology,standard_topology_reals}.
        $\reals\setminus\{0\}\subseteq\reals$ and $T_2$ is a topology on $\reals\setminus\{0\}$
            by \cref{setminus_elim_left,subseteq,subspace_topology_is_topology}.
        $f$ is continuous at $p$ from $\reals\setminus\{0\}$ to $\reals$ with respect to $T_2$ and $T_1$ by \cref{continuous_at}.
    \end{subproof}
    $\standardBasisR$ is a basis on $\reals$ and $T_1$ is a topology on $\reals$
        by \cref{standard_basis_is_basis,basis_topology_is_topology,standard_topology_reals}.
    $\reals\setminus\{0\}\subseteq\reals$ and $T_2$ is a topology on $\reals\setminus\{0\}$
        by \cref{setminus_elim_left,subseteq,subspace_topology_is_topology}.
    Show that $f$ is continuous from $\reals\setminus\{0\}$ to $\reals$ with respect to $T_2$ and $T_1$.
    \begin{subproof}
        Follows by \cref{continuous_at_implies_continuous}.
    \end{subproof}
\end{proof}

% from §21 Item 11: quotient operation is continuous
\begin{lemma}\label{real_division_continuous}
    Let $T_1 = \standardTopologyR$.
    Let $T_2 = \subspaceTopology{T_1}{\reals\setminus\{0\}}$.
    Let $T = \productTopology{T_1}{T_2}{\reals}{\reals\setminus\{0\}}$.
    Then $\realDivision$ is continuous from $\reals\times(\reals\setminus\{0\})$ to $\reals$
        with respect to $T$ and $\standardTopologyR$.
\end{lemma}
\begin{proof}
    Let $T_3 = \productTopology{T_1}{T_1}{\reals}{\reals}$.
    $\standardBasisR$ is a basis on $\reals$ and $T_1$ is a topology on $\reals$
        by \cref{standard_basis_is_basis,basis_topology_is_topology,standard_topology_reals}.
    $\reals\setminus\{0\}\subseteq\reals$ and $T_2$ is a topology on $\reals\setminus\{0\}$
        by \cref{setminus_elim_left,subseteq,subspace_topology_is_topology}.
    Let $f_1 = \piOne{\reals}{\reals\setminus\{0\}}$.
    Let $f_2 = \piTwo{\reals}{\reals\setminus\{0\}}$.
    $f_1$ is continuous from $\reals\times(\reals\setminus\{0\})$ to $\reals$ with respect to $T$ and $T_1$ and
        $f_2$ is continuous from $\reals\times(\reals\setminus\{0\})$ to $\reals\setminus\{0\}$ with respect to $T$ and $T_2$
        by \cref{projection_one_continuous,projection_two_continuous}.
    $\realInverse$ is continuous from $\reals\setminus\{0\}$ to $\reals$ with respect to $T_2$ and $T_1$
        by \cref{real_inverse_continuous}.
    Let $g = \realInverse\circ f_2$.
    $g$ is continuous from $\reals\times(\reals\setminus\{0\})$ to $\reals$ with respect to $T$ and $T_1$
        by \cref{continuous_composite}.
    Let $k = \{(p,(f_1(p), g(p))) \mid p\in\reals\times(\reals\setminus\{0\})\}$.
    $k$ is continuous from $\reals\times(\reals\setminus\{0\})$ to $\reals\times\reals$
        with respect to $T$ and $T_3$ by \cref{continuous_map_into_product_backward}.
    $\realMultiplication$ is continuous from $\reals\times\reals$ to $\reals$ with respect to $T_3$ and $T_1$
        by \cref{real_multiplication_continuous}.
    Let $m = \realMultiplication\circ k$.
    $m$ is continuous from $\reals\times(\reals\setminus\{0\})$ to $\reals$ with respect to $T$ and $T_1$
        by \cref{continuous_composite}.
    $m$ is a function and $\dom{m} = \reals\times(\reals\setminus\{0\})$ by \cref{continuous,funs_elim}.
    For all $p,z_1,z_2$ such that $(p,z_1)\in\realDivision$ and $(p,z_2)\in\realDivision$ we have $z_1 = z_2$.
    $\realDivision$ is a function by \cref{rightunique,real_division_operation,relation}.
    For all $p\in\dom{\realDivision}$ we have $\apply{\realDivision}{p}\in\reals$
        by \cref{dom_iff,real_division_operation,pair_eq_iff,function_apply_intro}.
    $\realDivision$ is a function to $\reals$.
    Show that for all $p$ such that $p\in\dom{\realDivision}$ we have $p\in\reals\times(\reals\setminus\{0\})$.
    \begin{subproof}
        Fix $p$ such that $p\in\dom{\realDivision}$.
        Take $z$ such that $(p,z)\in\realDivision$ by \cref{dom_iff}.
        Take $q,r$ such that $q\in\reals\times(\reals\setminus\{0\})$ and $r\in\reals$ and
            $(p,z) = (q,r)$ and $r = \fst{q} \rmul \inv{\snd{q}}$ by \cref{real_division_operation}.
        $p = q$ by \cref{pair_eq_iff}.
        $p\in\reals\times(\reals\setminus\{0\})$ by \cref{real_division_operation,pair_eq_iff}.
    \end{subproof}
    $\dom{\realDivision}\subseteq\reals\times(\reals\setminus\{0\})$ by \cref{subseteq}.
    Show that for all $p\in\reals\times(\reals\setminus\{0\})$ we have $p\in\dom{\realDivision}$.
    \begin{subproof}
        Follows by \cref{times_elem_is_tuple,fst_eq,snd_eq,setminus_elim_left,setminus_elim_right,singleton_iff,real_mul_inverse,real_mul_closed,real_division_operation,dom_intro}.
    \end{subproof}
    $\reals\times(\reals\setminus\{0\})\subseteq\dom{\realDivision}$ by \cref{subseteq}.
    $\dom{\realDivision} = \reals\times(\reals\setminus\{0\})$ by \cref{subseteq_antisymmetric}.
    Show that $\realDivision\in\funs{\reals\times(\reals\setminus\{0\})}{\reals}$.
    \begin{subproof}
        Follows by \cref{funs_intro}.
    \end{subproof}
    Show that for all $p\in\reals\times(\reals\setminus\{0\})$ we have $m(p) = \apply{\realDivision}{p}$.
    \begin{subproof}
            Fix $p\in\reals\times(\reals\setminus\{0\})$.
            $k\in\funs{\reals\times(\reals\setminus\{0\})}{\reals\times\reals}$ by \cref{continuous}.
            $(p,(f_1(p), g(p)))\in k$.
            $k$ is a function by \cref{continuous,funs_is_function}.
            $k(p) = (f_1(p), g(p))$ by \cref{continuous,function_apply_intro}.
            $\realMultiplication\in\funs{\reals\times\reals}{\reals}$ and
                $\dom{\realMultiplication} = \reals\times\reals$ by \cref{continuous,funs}.
            $\ran{k}\subseteq\reals\times\reals$ by \cref{funs_ran}.
            $\realMultiplication$ is a function by \cref{funs_is_function}.
            $p\in\dom{k}$ by \cref{funs}.
            $m(p) = \apply{\realMultiplication}{k(p)}$ by \cref{circ_apply}.
            $f_1\in\funs{\reals\times(\reals\setminus\{0\})}{\reals}$ and
                $g\in\funs{\reals\times(\reals\setminus\{0\})}{\reals}$ by \cref{continuous}.
            $f_1(p)\in\reals$ and $g(p)\in\reals$ and $(f_1(p), g(p))\in\reals\times\reals$
                by \cref{funs_type_apply,times_tuple_intro}.
            $\fst{k(p)} = f_1(p)$ and $\snd{k(p)} = g(p)$ by \cref{fst_eq,snd_eq}.
            $\apply{\realMultiplication}{k(p)} = f_1(p) \rmul g(p)$
                by \cref{continuous,real_multiplication_operation,real_multiplication_continuous,funs_is_function,function_apply_intro,real_mul_closed}.
            $m(p) = f_1(p) \rmul g(p)$.
            Take $a,b$ such that $p = (a,b)$ and $a\in\reals$ and $b\in\reals\setminus\{0\}$ by \cref{times_elem_is_tuple}.
            $f_1(p) = a$ by \cref{projection_first,continuous,funs_is_function,function_apply_intro}.
            $f_2(p) = b$ by \cref{projection_second,continuous,funs_is_function,function_apply_intro}.
            $f_2\in\funs{\reals\times(\reals\setminus\{0\})}{\reals\setminus\{0\}}$ and
                $\realInverse\in\funs{\reals\setminus\{0\}}{\reals}$ by \cref{continuous}.
            $f_2$ is right-unique and $f_2$ is a relation and
                $\realInverse$ is right-unique and $\realInverse$ is a relation.
            $\ran{f_2}\subseteq\dom{\realInverse}$ by \cref{funs_ran,funs,subseteq}.
            $p\in\dom{f_2}$ by \cref{funs}.
            $g(p) = \apply{\realInverse}{f_2(p)}$ by \cref{circ_apply}.
            $\fst{p} = a$ and $\snd{p} = b$ by \cref{fst_eq,snd_eq}.
            $f_1(p) = \fst{p}$ and $f_2(p) = \snd{p}$.
            $g(p) = \apply{\realInverse}{\snd{p}}$.
            $\snd{p}\in\reals\setminus\{0\}$.
            $\apply{\realInverse}{\snd{p}} = \inv{\snd{p}}$
                by \cref{continuous,setminus_elim_left,setminus_elim_right,singleton_iff,real_mul_inverse,real_inverse_operation,funs_is_function,function_apply_intro}.
            $g(p) = \inv{\snd{p}}$.
            $m(p) = \fst{p} \rmul \inv{\snd{p}}$.
            $m(p) = \apply{\realDivision}{p}$
                by \cref{real_mul_closed,real_mul_inverse,real_division_operation,funs_is_function,function_apply_intro}.
        \end{subproof}
        Show that $m = \realDivision$.
        \begin{subproof}
            Follows by \cref{continuous,funs_elim,subseteq,fun_subseteq,subseteq_antisymmetric}.
        \end{subproof}
    $\realDivision$ is continuous from $\reals\times(\reals\setminus\{0\})$ to $\reals$
        with respect to $T$ and $\standardTopologyR$ by \cref{continuous_composite}.
\end{proof}
% from §21 Item 12: sum of continuous functions
\begin{theorem}\label{real_function_addition_continuous}
    Suppose $f$ is continuous from $X$ to $\reals$ with respect to $T$ and $\standardTopologyR$.
    Suppose $g$ is continuous from $X$ to $\reals$ with respect to $T$ and $\standardTopologyR$.
    Let $h = \{(x, f(x) + g(x)) \mid x\in X\}$.
    Then $h$ is continuous from $X$ to $\reals$ with respect to $T$ and $\standardTopologyR$.
\end{theorem}
\begin{proof}
    Let $T_1 = \standardTopologyR$.
    Let $T_2 = \productTopology{T_1}{T_1}{\reals}{\reals}$.
    Let $k = \{(x,(f(x), g(x))) \mid x\in X\}$.
    $k$ is continuous from $X$ to $\reals\times\reals$ with respect to $T$ and $T_2$
        by \cref{continuous_map_into_product_backward}.
    $\realAddition$ is continuous from $\reals\times\reals$ to $\reals$ with respect to $T_2$ and $T_1$
        by \cref{real_addition_continuous}.
    Let $m = \realAddition\circ k$.
    $m$ is continuous from $X$ to $\reals$ with respect to $T$ and $T_1$ by \cref{continuous_composite}.
    Show that for all $x,z_1,z_2$ such that $(x,z_1)\in h$ and $(x,z_2)\in h$ we have $z_1 = z_2$.
    \begin{subproof}
        Trivial.
    \end{subproof}
    Show that for all $w$ such that $w\in h$ there exist $x,y$ such that $w = (x,y)$.
    \begin{subproof}
        Trivial.
    \end{subproof}
    $h$ is a function by \cref{rightunique,relation}.
    For all $x\in\dom{h}$ there exists $z$ such that $(x,z)\in h$ by \cref{dom_iff}.
    For all $x,z$ such that $x\in\dom{h}$ and $(x,z)\in h$ we have $x\in X$ and $z = f(x) + g(x)$.
    for all $x\in\dom{h}$ we have $h(x)\in\reals$
        by \cref{function_apply_intro,continuous,funs_type_apply,real_add_closed}.
    $h$ is a function to $\reals$.
    Show that for all $x\in\dom{h}$ we have $x\in X$.
    \begin{subproof}
        Trivial.
    \end{subproof}
    Show that for all $x\in X$ we have $(x, f(x) + g(x))\in h$.
    \begin{subproof}
        Trivial.
    \end{subproof}
    Show that $h\in\funs{X}{\reals}$.
    \begin{subproof}
        Follows by \cref{subseteq,dom_intro,subseteq_antisymmetric,funs_intro}.
    \end{subproof}
    Show that for all $x\in X$ we have $m(x) = \apply{h}{x}$.
    \begin{subproof}
        Fix $x\in X$.
        $k\in\funs{X}{\reals\times\reals}$ and $k$ is a function by \cref{continuous,funs_is_function}.
        $k(x) = (f(x), g(x))$ by \cref{continuous,function_apply_intro}.
        $f(x)\in\reals$ and $g(x)\in\reals$ and $(f(x), g(x))\in\reals\times\reals$
            by \cref{continuous,funs_type_apply,times_tuple_intro}.
        $\fst{k(x)} = f(x)$ and $\snd{k(x)} = g(x)$ by \cref{fst_eq,snd_eq}.
        $\realAddition$ is a function by \cref{continuous,real_addition_continuous,funs_is_function}.
        $\apply{\realAddition}{k(x)} = f(x) + g(x)$
            by \cref{continuous,real_addition_operation,real_addition_continuous,funs_is_function,function_apply_intro,real_add_closed}.
        $m(x) = \apply{\realAddition}{k(x)}$
            by \cref{funs_ran,continuous,real_addition_continuous,funs,subseteq,circ_apply}.
        $m(x) = \apply{h}{x}$ by \cref{funs_is_function,function_apply_intro}.
    \end{subproof}
    Show that $m = h$.
    \begin{subproof}
        Follows by \cref{continuous,funs_elim,subseteq,fun_subseteq,subseteq_antisymmetric}.
    \end{subproof}
    $h$ is continuous from $X$ to $\reals$ with respect to $T$ and $\standardTopologyR$ by \cref{continuous_composite}.
\end{proof}

% from §21 Item 13: difference of continuous functions
\begin{theorem}\label{real_function_subtraction_continuous}
    Suppose $f$ is continuous from $X$ to $\reals$ with respect to $T$ and $\standardTopologyR$.
    Suppose $g$ is continuous from $X$ to $\reals$ with respect to $T$ and $\standardTopologyR$.
    Let $h = \{(x, f(x) - g(x)) \mid x\in X\}$.
    Then $h$ is continuous from $X$ to $\reals$ with respect to $T$ and $\standardTopologyR$.
\end{theorem}
\begin{proof}
    Let $T_1 = \standardTopologyR$.
    Let $T_2 = \productTopology{T_1}{T_1}{\reals}{\reals}$.
    Let $k = \{(x,(f(x), g(x))) \mid x\in X\}$.
    $k$ is continuous from $X$ to $\reals\times\reals$ with respect to $T$ and $T_2$
        by \cref{continuous_map_into_product_backward}.
    $\realSubtraction$ is continuous from $\reals\times\reals$ to $\reals$ with respect to $T_2$ and $T_1$
        by \cref{real_subtraction_continuous}.
    Let $m = \realSubtraction\circ k$.
    $m$ is continuous from $X$ to $\reals$ with respect to $T$ and $T_1$ by \cref{continuous_composite}.
    Show that for all $x,z_1,z_2$ such that $(x,z_1)\in h$ and $(x,z_2)\in h$ we have $z_1 = z_2$.
    \begin{subproof}
        Trivial.
    \end{subproof}
    Show that for all $w$ such that $w\in h$ there exist $x,y$ such that $w = (x,y)$.
    \begin{subproof}
        Trivial.
    \end{subproof}
    $h$ is a function by \cref{rightunique,relation}.
    For all $x\in\dom{h}$ there exists $z$ such that $(x,z)\in h$ by \cref{dom_iff}.
    For all $x,z$ such that $x\in\dom{h}$ and $(x,z)\in h$ we have $x\in X$ and $z = f(x) - g(x)$.
    for all $x\in\dom{h}$ we have $h(x)\in\reals$
        by \cref{function_apply_intro,continuous,funs_type_apply,real_add_inverse,real_sub,real_add_closed}.
    $h$ is a function to $\reals$.
    Show that for all $x\in\dom{h}$ we have $x\in X$.
    \begin{subproof}
        Trivial.
    \end{subproof}
    Show that for all $x\in X$ we have $(x, f(x) - g(x))\in h$.
    \begin{subproof}
        Trivial.
    \end{subproof}
    Show that $h\in\funs{X}{\reals}$.
    \begin{subproof}
        Follows by \cref{subseteq,dom_intro,subseteq_antisymmetric,funs_intro}.
    \end{subproof}
    Show that for all $x\in X$ we have $m(x) = \apply{h}{x}$.
    \begin{subproof}
        Fix $x\in X$.
        $k\in\funs{X}{\reals\times\reals}$ and $k$ is a function by \cref{continuous,funs_is_function}.
        $k(x) = (f(x), g(x))$ by \cref{continuous,function_apply_intro}.
        $f(x)\in\reals$ and $g(x)\in\reals$ and $(f(x), g(x))\in\reals\times\reals$
            by \cref{continuous,funs_type_apply,times_tuple_intro}.
        $\fst{k(x)} = f(x)$ and $\snd{k(x)} = g(x)$ by \cref{fst_eq,snd_eq}.
        $\realSubtraction$ is a function by \cref{continuous,real_subtraction_continuous,funs_is_function}.
        $\apply{\realSubtraction}{k(x)} = f(x) - g(x)$
            by \cref{continuous,real_subtraction_operation,real_subtraction_continuous,funs_is_function,function_apply_intro,real_add_inverse,real_sub,real_add_closed}.
        $m(x) = \apply{\realSubtraction}{k(x)}$
            by \cref{funs_ran,continuous,real_subtraction_continuous,funs,subseteq,circ_apply}.
        $m(x) = \apply{h}{x}$ by \cref{funs_is_function,function_apply_intro}.
    \end{subproof}
    Show that $m = h$.
    \begin{subproof}
        Follows by \cref{continuous,funs_elim,subseteq,fun_subseteq,subseteq_antisymmetric}.
    \end{subproof}
    $h$ is continuous from $X$ to $\reals$ with respect to $T$ and $\standardTopologyR$ by \cref{continuous_composite}.
\end{proof}

% from §21 Item 14: product of continuous functions
\begin{theorem}\label{real_function_multiplication_continuous}
    Suppose $f$ is continuous from $X$ to $\reals$ with respect to $T$ and $\standardTopologyR$.
    Suppose $g$ is continuous from $X$ to $\reals$ with respect to $T$ and $\standardTopologyR$.
    Let $h = \{(x, f(x) \rmul g(x)) \mid x\in X\}$.
    Then $h$ is continuous from $X$ to $\reals$ with respect to $T$ and $\standardTopologyR$.
\end{theorem}
\begin{proof}
    Let $T_1 = \standardTopologyR$.
    Let $T_2 = \productTopology{T_1}{T_1}{\reals}{\reals}$.
    Let $k = \{(x,(f(x), g(x))) \mid x\in X\}$.
    $k$ is continuous from $X$ to $\reals\times\reals$ with respect to $T$ and $T_2$
        by \cref{continuous_map_into_product_backward}.
    $\realMultiplication$ is continuous from $\reals\times\reals$ to $\reals$ with respect to $T_2$ and $T_1$
        by \cref{real_multiplication_continuous}.
    Let $m = \realMultiplication\circ k$.
    $m$ is continuous from $X$ to $\reals$ with respect to $T$ and $T_1$ by \cref{continuous_composite}.
    Show that for all $x,z_1,z_2$ such that $(x,z_1)\in h$ and $(x,z_2)\in h$ we have $z_1 = z_2$.
    \begin{subproof}
        Trivial.
    \end{subproof}
    Show that for all $w$ such that $w\in h$ there exist $x,y$ such that $w = (x,y)$.
    \begin{subproof}
        Trivial.
    \end{subproof}
    $h$ is a function by \cref{rightunique,relation}.
    For all $x\in\dom{h}$ there exists $z$ such that $(x,z)\in h$ by \cref{dom_iff}.
    For all $x,z$ such that $x\in\dom{h}$ and $(x,z)\in h$ we have $x\in X$ and $z = f(x) \rmul g(x)$.
    for all $x\in\dom{h}$ we have $h(x)\in\reals$
        by \cref{function_apply_intro,continuous,funs_type_apply,real_mul_closed}.
    $h$ is a function to $\reals$.
    Show that for all $x\in\dom{h}$ we have $x\in X$.
    \begin{subproof}
        Trivial.
    \end{subproof}
    Show that for all $x\in X$ we have $(x, f(x) \rmul g(x))\in h$.
    \begin{subproof}
        Trivial.
    \end{subproof}
    Show that $h\in\funs{X}{\reals}$.
    \begin{subproof}
        Follows by \cref{subseteq,dom_intro,subseteq_antisymmetric,funs_intro}.
    \end{subproof}
    Show that for all $x\in X$ we have $m(x) = \apply{h}{x}$.
    \begin{subproof}
        Fix $x\in X$.
        $k\in\funs{X}{\reals\times\reals}$ and $k$ is a function by \cref{continuous,funs_is_function}.
        $k(x) = (f(x), g(x))$ by \cref{continuous,function_apply_intro}.
        $f(x)\in\reals$ and $g(x)\in\reals$ and $(f(x), g(x))\in\reals\times\reals$
            by \cref{continuous,funs_type_apply,times_tuple_intro}.
        $\fst{k(x)} = f(x)$ and $\snd{k(x)} = g(x)$ by \cref{fst_eq,snd_eq}.
        $\realMultiplication$ is a function by \cref{continuous,real_multiplication_continuous,funs_is_function}.
        $\apply{\realMultiplication}{k(x)} = f(x) \rmul g(x)$
            by \cref{continuous,real_multiplication_operation,real_multiplication_continuous,funs_is_function,function_apply_intro,real_mul_closed}.
        $m(x) = \apply{\realMultiplication}{k(x)}$
            by \cref{funs_ran,continuous,real_multiplication_continuous,funs,subseteq,circ_apply}.
        $m(x) = \apply{h}{x}$ by \cref{funs_is_function,function_apply_intro}.
    \end{subproof}
    Show that $m = h$.
    \begin{subproof}
        Follows by \cref{continuous,funs_elim,subseteq,fun_subseteq,subseteq_antisymmetric}.
    \end{subproof}
    $h$ is continuous from $X$ to $\reals$ with respect to $T$ and $\standardTopologyR$ by \cref{continuous_composite}.
\end{proof}

% from §21 Item 15: quotient of continuous functions
\begin{theorem}\label{real_function_division_continuous}
    Suppose $f$ is continuous from $X$ to $\reals$ with respect to $T$ and $\standardTopologyR$.
    Suppose $g$ is continuous from $X$ to $\reals$ with respect to $T$ and $\standardTopologyR$.
    Suppose for all $x\in X$ we have $g(x)\neq 0$.
    Let $h = \{(x, f(x) \rmul \inv{(g(x))}) \mid x\in X\}$.
    Then $h$ is continuous from $X$ to $\reals$ with respect to $T$ and $\standardTopologyR$.
\end{theorem}
\begin{proof}
    Let $T_1 = \standardTopologyR$.
    Let $T_2 = \subspaceTopology{T_1}{\reals\setminus\{0\}}$.
    For all $y$ such that $y\in\img{g}{X}$ we have $y\in\reals\setminus\{0\}$
        by \cref{img_iff,continuous,funs_is_function,function_apply_intro,funs_type_apply,singleton_iff,setminus_intro}.
    $\img{g}{X}\subseteq \reals\setminus\{0\}$ by \cref{subseteq}.
    $\standardBasisR$ is a basis on $\reals$ and $T_1$ is a topology on $\reals$
        by \cref{standard_basis_is_basis,basis_topology_is_topology,standard_topology_reals}.
    $\reals\setminus\{0\}$ is a subspace of $\reals$ with respect to $T_1$
        by \cref{setminus_elim_left,subseteq,subspace_of}.
    $g$ is continuous from $X$ to $\reals\setminus\{0\}$ with respect to $T$ and $T_2$
        by \cref{continuous_restrict_range}.
    $\realInverse$ is continuous from $\reals\setminus\{0\}$ to $\reals$ with respect to $T_2$ and $T_1$
        by \cref{real_inverse_continuous}.
    Let $j = \realInverse\circ g$.
    $j$ is continuous from $X$ to $\reals$ with respect to $T$ and $T_1$ by \cref{continuous_composite}.
    Let $k = \{(x,(f(x), j(x))) \mid x\in X\}$.
    $k$ is continuous from $X$ to $\reals\times\reals$ with respect to $T$ and
        $\productTopology{T_1}{T_1}{\reals}{\reals}$ by \cref{continuous_map_into_product_backward}.
    $\realMultiplication$ is continuous from $\reals\times\reals$ to $\reals$ with respect to
        $\productTopology{T_1}{T_1}{\reals}{\reals}$ and $T_1$ by \cref{real_multiplication_continuous}.
    Let $m = \realMultiplication\circ k$.
    $m$ is continuous from $X$ to $\reals$ with respect to $T$ and $T_1$ by \cref{continuous_composite}.
    Show that for all $x,z_1,z_2$ such that $(x,z_1)\in h$ and $(x,z_2)\in h$ we have $z_1 = z_2$.
    \begin{subproof}
        Trivial.
    \end{subproof}
    Show that for all $w$ such that $w\in h$ there exist $x,y$ such that $w = (x,y)$.
    \begin{subproof}
        Trivial.
    \end{subproof}
    $h$ is a function by \cref{rightunique,relation}.
    For all $x\in\dom{h}$ there exists $z$ such that $(x,z)\in h$ by \cref{dom_iff}.
    For all $x,z$ such that $x\in\dom{h}$ and $(x,z)\in h$ we have $x\in X$ and $z = f(x) \rmul \inv{(g(x))}$ by \cref{pair_eq_iff}.
    for all $x\in\dom{h}$ we have $h(x)\in\reals$
        by \cref{function_apply_intro,continuous,funs_type_apply,real_mul_inverse,real_mul_closed}.
    $h$ is a function to $\reals$.
    Show that for all $x\in\dom{h}$ we have $x\in X$.
    \begin{subproof}
        Trivial.
    \end{subproof}
    Show that for all $x\in X$ we have $(x, f(x) \rmul \inv{(g(x))})\in h$.
    \begin{subproof}
        Trivial.
    \end{subproof}
    Show that $h\in\funs{X}{\reals}$.
    \begin{subproof}
        Follows by \cref{subseteq,dom_intro,subseteq_antisymmetric,funs_intro}.
    \end{subproof}
    Show that for all $x\in X$ we have $m(x) = \apply{h}{x}$.
    \begin{subproof}
        Fix $x\in X$.
        $k\in\funs{X}{\reals\times\reals}$ and $k$ is a function by \cref{continuous,funs_is_function}.
        $k(x) = (f(x), j(x))$ by \cref{continuous,function_apply_intro}.
        $f(x)\in\reals$ and $j(x)\in\reals$ and $(f(x), j(x))\in\reals\times\reals$
            by \cref{continuous,funs_type_apply,times_tuple_intro}.
        $\fst{k(x)} = f(x)$ and $\snd{k(x)} = j(x)$ by \cref{fst_eq,snd_eq}.
        $\realMultiplication$ is a function by \cref{continuous,real_multiplication_continuous,funs_is_function}.
        $\apply{\realMultiplication}{k(x)} = f(x) \rmul j(x)$
            by \cref{continuous,real_multiplication_operation,real_multiplication_continuous,funs_is_function,function_apply_intro,real_mul_closed}.
        $m(x) = \apply{\realMultiplication}{k(x)}$
            by \cref{funs_ran,continuous,real_multiplication_continuous,funs,subseteq,circ_apply}.
        $g\in\funs{X}{\reals\setminus\{0\}}$ by \cref{continuous}.
        $g(x)\in\reals\setminus\{0\}$ by \cref{funs_type_apply}.
        $\realInverse\in\funs{\reals\setminus\{0\}}{\reals}$ and
            $\dom{\realInverse} = \reals\setminus\{0\}$ by \cref{continuous,real_inverse_continuous,funs}.
        $\ran{g}\subseteq\dom{\realInverse}$ by \cref{funs_ran,subseteq}.
        $g$ is a function and $\realInverse$ is a function by \cref{funs_is_function}.
        $x\in\dom{g}$ by \cref{funs}.
        $j(x) = \apply{\realInverse}{g(x)}$ by \cref{circ_apply}.
        $g(x)\neq 0$ and $g(x)\in\reals$ and $\inv{(g(x))}\in\reals$
            by \cref{singleton_iff,setminus_elim_left,setminus_elim_right,real_mul_inverse}.
        $\apply{\realInverse}{g(x)} = \inv{(g(x))}$
            by \cref{real_inverse_operation,funs_is_function,function_apply_intro}.
        $m(x) = \apply{h}{x}$ by \cref{funs_is_function,function_apply_intro}.
    \end{subproof}
    Show that $m = h$.
    \begin{subproof}
        Follows by \cref{continuous,funs_elim,subseteq,fun_subseteq,subseteq_antisymmetric}.
    \end{subproof}
    $h$ is continuous from $X$ to $\reals$ with respect to $T$ and $\standardTopologyR$ by \cref{continuous_composite}.
\end{proof}

% from §21 Item 16: uniform convergence
\begin{definition}\label{uniform_converges_to}
    $s$ converges uniformly to $f$ from $X$ to $Y$ with respect to $d$ iff
        $s$ is a sequence in $\funs{X}{Y}$ and
        $f\in\funs{X}{Y}$ and
        for all $\epsilon\in\reals$ such that $0 < \epsilon$
        there exists $N\in\naturalsPlus$ such that for all $n\in\naturalsPlus$ if $N\leq n$ then
            for all $x\in X$ we have $d((\apply{\apply{s}{n}}{x}, f(x))) < \epsilon$.
\end{definition}

% from §21 Item 17: uniform limit theorem
\begin{theorem}\label{uniform_limit_theorem}
    Suppose $T_X$ is a topology on $X$.
    Suppose $s$ is a sequence in $\funs{X}{Y}$.
    Suppose $f\in\funs{X}{Y}$.
    Suppose $d$ is a metric on $Y$.
    Let $T = \metricTopology{d}{Y}$.
    Suppose for all $n\in\naturalsPlus$ we have
        $s(n)$ is continuous from $X$ to $Y$ with respect to $T_X$ and $T$.
    Suppose $s$ converges uniformly to $f$ from $X$ to $Y$ with respect to $d$.
    Then $f$ is continuous from $X$ to $Y$ with respect to $T_X$ and $T$.
\end{theorem}
\begin{proof}
    $f$ is a function from $X$ to $Y$ by \cref{funs}.
    $d$ is a metric on $Y$.
    $\metricBasis{d}{Y}$ is a basis on $Y$ by \cref{metric_basis_is_basis}.
    $T$ is a topology on $Y$ by \cref{basis_topology_is_topology,metric_topology,basis_for_topology}.
    Show that for all $V$ such that $V$ is open in $Y$ with respect to $T$ we have
        $\preimg{f}{V}$ is open in $X$ with respect to $T_X$.
    \begin{subproof}
        Fix $V$ such that $V$ is open in $Y$ with respect to $T$.
        $V\subseteq Y$ by \cref{open_in,topology_on,subseteq,pow_iff}.
        Show that for all $x\in\preimg{f}{V}$ there exists $U$ such that
            $U$ is open in $X$ with respect to $T_X$ and $x\in U$ and $U\subseteq\preimg{f}{V}$.
        \begin{subproof}
            Fix $x$ such that $x\in\preimg{f}{V}$.
            Take $y\in V$ such that $(x,y)\in f$ by \cref{preimg_iff}.
            $f(x)\in V$ by \cref{funs_is_function,function_apply_iff}.
            Take $\epsilon\in\reals$ such that $0 < \epsilon$ and
                $\metricBall{d}{Y}{f(x)}{\epsilon} \subseteq V$ by \cref{metric_open_iff}.
            Let $\epsilon_1 = \epsilon / (1 + 1)$.
            $0 < 1$ and $1\in\reals$ by \cref{real_zero_lt_one,real_naturals_subset_reals,real_one_in_naturals,subseteq}.
            $1 + 1$ is positive by \cref{real_pos_add}.
            $1 + 1\in\reals$ and $1 + 1 \neq 0$ by \cref{real_add_closed,real_pos_nonzero}.
            $\inv{(1 + 1)}\in\reals$ by \cref{real_mul_inverse}.
            $0 < \inv{(1 + 1)}$ by \cref{real_inv_pos}.
            $\epsilon \rmul \inv{(1 + 1)}$ is positive by \cref{real_mul_pos_pos}.
            $\epsilon / (1 + 1) = \epsilon \rmul \inv{(1 + 1)}$ by \cref{real_div}.
            $0 < \epsilon_1$ by \cref{real_lt_subst_right}.
            $\epsilon_1\in\reals$ by \cref{real_div,real_mul_closed}.
            $s\in\funs{\naturalsPlus}{\funs{X}{Y}}$ by \cref{sequence_in}.
            Take $N\in\naturalsPlus$ such that for all $n\in\naturalsPlus$ if $N\leq n$ then
                for all $z\in X$ we have $d((\apply{\apply{s}{n}}{z}, f(z))) < \epsilon_1$
                by \cref{uniform_converges_to}.
            Let $g = s(N)$.
            $g\in\funs{X}{Y}$ by \cref{funs_type_apply}.
            $g$ is a function and $f$ is a function by \cref{funs_is_function}.
            $N\leq N$.
            for all $z\in X$ we have $d((g(z), f(z))) < \epsilon_1$.
            $N\in\naturalsPlus$.
            $g$ is continuous from $X$ to $Y$ with respect to $T_X$ and $T$.
            Let $W = \metricBall{d}{Y}{f(x)}{\epsilon_1}$.
            $W\in\pow{Y}$ by \cref{subseteq,powerset_intro,metric_ball}.
            $x\in X$ by \cref{preimg_subseteq_dom,elem_subseteq,funs}.
            $f(x)\in Y$ by \cref{preimg_subseteq_dom,elem_subseteq,funs,funs_type_apply}.
            $W$ is open in $Y$ with respect to $T$
                by \cref{metric_basis,metric_basis_is_basis,basis_elem_open_in_generated,metric_topology,basis_for_topology,open_in}.
            Let $U = \preimg{g}{W}$.
            $U$ is open in $X$ with respect to $T_X$ by \cref{continuous}.
            $g(x)\in Y$ and $f(x)\in Y$ by \cref{funs_type_apply}.
            $d((f(x), g(x))) = d((g(x), f(x)))$ by \cref{metric_on,metric_symmetric}.
            $d((f(x), g(x))) < \epsilon_1$ by \cref{real_lt_subst_left}.
            $g(x)\in W$ by \cref{metric_ball}.
            $x\in U$ by \cref{funs,function_apply_iff,preimg_iff}.
            Show that for all $u$ such that $u\in U$ we have $u\in\preimg{f}{V}$.
            \begin{subproof}
                Fix $u$ such that $u\in U$.
                Take $z\in W$ such that $(u,z)\in g$ by \cref{preimg_iff}.
                $g(u)\in Y$ and $d((f(x), g(u))) < \epsilon_1$ by \cref{function_apply_iff,metric_ball}.
                $d((g(u), f(x))) < \epsilon_1$
                    by \cref{subseteq,metric_on,metric_symmetric,real_lt_subst_left}.
                $u\in X$ by \cref{preimg_subseteq_dom,elem_subseteq,funs}.
                $d((g(u), f(u))) < \epsilon_1$.
                $f(u)\in Y$ by \cref{funs_type_apply}.
                $d((f(u), g(u))) < \epsilon_1$
                    by \cref{metric_on,metric_symmetric,real_lt_subst_left}.
                $d\in\funs{Y\times Y}{\reals}$ by \cref{metric_on}.
                $(f(u), g(u))\in Y\times Y$ and $(g(u), f(x))\in Y\times Y$ and $(f(u), f(x))\in Y\times Y$
                    by \cref{times_tuple_intro}.
                $d((f(u), g(u)))\in\reals$ and $d((g(u), f(x)))\in\reals$ and $d((f(u), f(x)))\in\reals$
                    by \cref{funs_type_apply}.
                $d((f(u), g(u))) + d((g(u), f(x))) \geq d((f(u), f(x)))$
                    by \cref{metric_on,metric_triangle_inequality}.
                $d((f(u), f(x))) \leq d((f(u), g(u))) + d((g(u), f(x)))$.
                $d((f(u), g(u))) + d((g(u), f(x))) < \epsilon_1 + \epsilon_1$
                    by \cref{real_add_closed,real_lt_add_two}.
                $(1 + 1) \rmul \epsilon_1 = \epsilon_1 + \epsilon_1$
                    by \cref{real_right_distrib,real_mul_ident,real_add_eq_subst}.
                $\epsilon_1 + \epsilon_1 = \epsilon$
                    by \cref{real_div,real_mul_eq_subst,real_mul_assoc,real_mul_comm,real_mul_inverse,real_mul_ident}.
                $d((f(u), g(u))) + d((g(u), f(x))) < \epsilon$ by \cref{real_lt_subst_right}.
                Show that $d((f(u), f(x))) < \epsilon$.
                \begin{subproof}
                    $d((f(u), f(x))) < d((f(u), g(u))) + d((g(u), f(x)))$ or
                        $d((f(u), f(x))) = d((f(u), g(u))) + d((g(u), f(x)))$ by \cref{real_leq_split}.
                    \begin{byCase}
                    \caseOf{$d((f(u), f(x))) < d((f(u), g(u))) + d((g(u), f(x)))$.}
                        $d((f(u), g(u))) + d((g(u), f(x)))\in\reals$ by \cref{real_add_closed}.
                        $d((f(u), f(x))) < \epsilon$ by \cref{real_lt_trans}.
                    \caseOf{$d((f(u), f(x))) = d((f(u), g(u))) + d((g(u), f(x)))$.}
                        $d((f(u), f(x))) < \epsilon$ by \cref{real_lt_subst_left}.
                    \end{byCase}
                \end{subproof}
                $d((f(x), f(u))) = d((f(u), f(x)))$ by \cref{metric_on,metric_symmetric}.
                $d((f(x), f(u))) < \epsilon$ by \cref{real_lt_subst_left}.
                $f(u)\in\metricBall{d}{Y}{f(x)}{\epsilon}$ by \cref{metric_ball}.
                $f(u)\in V$ by \cref{subseteq}.
                $u\in\preimg{f}{V}$ by \cref{funs,function_apply_iff,preimg_iff}.
            \end{subproof}
            $U\subseteq\preimg{f}{V}$ by \cref{subseteq}.
            Show that there exists $W_1$ such that
                $W_1$ is open in $X$ with respect to $T_X$ and $x\in W_1$ and $W_1\subseteq\preimg{f}{V}$.
            \begin{subproof}
                Trivial.
            \end{subproof}
        \end{subproof}
        Let $M = \{U\in T_X \mid U\subseteq\preimg{f}{V}\}$.
        $\unions{M}\in T_X$ by \cref{subseteq,topology_on}.
        Show that for all $x$ such that $x\in\preimg{f}{V}$ we have $x\in\unions{M}$.
        \begin{subproof}
            Fix $x$ such that $x\in\preimg{f}{V}$.
            Take $U$ such that $U$ is open in $X$ with respect to $T_X$ and
                $x\in U$ and $U\subseteq\preimg{f}{V}$.
            $U\in T_X$ by \cref{open_in}.
            $U\in M$.
            $x\in\unions{M}$ by \cref{unions_iff}.
        \end{subproof}
        $\preimg{f}{V}\subseteq\unions{M}$ by \cref{subseteq}.
        Show that for all $x$ such that $x\in\unions{M}$ we have $x\in\preimg{f}{V}$.
        \begin{subproof}
            Fix $x$ such that $x\in\unions{M}$.
            Take $U\in M$ such that $x\in U$ by \cref{unions_iff}.
            $U\subseteq\preimg{f}{V}$.
            $x\in\preimg{f}{V}$ by \cref{subseteq}.
        \end{subproof}
        $\unions{M}\subseteq\preimg{f}{V}$ by \cref{subseteq}.
        $\preimg{f}{V} = \unions{M}$ by \cref{subseteq_antisymmetric}.
        $\preimg{f}{V}\in T_X$.
        $\preimg{f}{V}$ is open in $X$ with respect to $T_X$ by \cref{open_in}.
    \end{subproof}
    $f$ is continuous from $X$ to $Y$ with respect to $T_X$ and $T$ by \cref{continuous}.
\end{proof}
\section{§22 The Quotient Topology}

% from §22 Item 1: quotient map
\begin{definition}\label{quotient_map}
    $p$ is a quotient map from $X$ to $Y$ with respect to $T_X$ and $T_Y$ iff
        $p$ is a surjection from $X$ to $Y$ and
        $p$ is continuous from $X$ to $Y$ with respect to $T_X$ and $T_Y$ and
        for all $U$ such that $U\subseteq Y$ if $\preimg{p}{U}$ is open in $X$ with respect to $T_X$ then
            $U$ is open in $Y$ with respect to $T_Y$.
\end{definition}

% from §22 Item 2: quotient map closed set characterization
\begin{lemma}\label{quotient_map_closed_iff}
    Suppose $p$ is a surjection from $X$ to $Y$.
    Suppose $T_X$ is a topology on $X$.
    Suppose $T_Y$ is a topology on $Y$.
    Then $p$ is a quotient map from $X$ to $Y$ with respect to $T_X$ and $T_Y$ iff
        for all $A$ such that $A\subseteq Y$ we have
            $A$ is closed in $Y$ with respect to $T_Y$ iff
            $\preimg{p}{A}$ is closed in $X$ with respect to $T_X$.
\end{lemma}
\begin{proof}
    $p$ is a function from $X$ to $Y$ by \cref{surj,funs}.
    Show that if $p$ is a quotient map from $X$ to $Y$ with respect to $T_X$ and $T_Y$ then
        for all $A$ such that $A\subseteq Y$ we have
            $A$ is closed in $Y$ with respect to $T_Y$ iff
            $\preimg{p}{A}$ is closed in $X$ with respect to $T_X$.
    \begin{subproof}
        Follows by \cref{closed_in,setminus_subseteq,quotient_map,continuous,preimg_setminus_function,preimg_subseteq_dom,funs,subseteq}.
    \end{subproof}
    Show that if for all $A$ such that $A\subseteq Y$ we have
            $A$ is closed in $Y$ with respect to $T_Y$ iff
            $\preimg{p}{A}$ is closed in $X$ with respect to $T_X$
        then $p$ is a quotient map from $X$ to $Y$ with respect to $T_X$ and $T_Y$.
    \begin{subproof}
        Assume that for all $A$ such that $A\subseteq Y$ we have
            $A$ is closed in $Y$ with respect to $T_Y$ iff
            $\preimg{p}{A}$ is closed in $X$ with respect to $T_X$.
        For all $B$ such that $B$ is closed in $Y$ with respect to $T_Y$ we have
            $\preimg{p}{B}$ is closed in $X$ with respect to $T_X$.
        $p$ is continuous from $X$ to $Y$ with respect to $T_X$ and $T_Y$
            by \cref{preimg_closed_implies_continuous}.
        Show that for all $U$ such that $U\subseteq Y$ we have
            $U$ is open in $Y$ with respect to $T_Y$ iff
            $\preimg{p}{U}$ is open in $X$ with respect to $T_X$.
        \begin{subproof}
            Follows by
                \cref{continuous,preimg_subseteq_dom,funs,subseteq,setminus_subseteq,double_relative_complement,open_in,closed_in,preimg_setminus_function}.
        \end{subproof}
        $p$ is a quotient map from $X$ to $Y$ with respect to $T_X$ and $T_Y$ by \cref{quotient_map}.
    \end{subproof}
\end{proof}

% from §22 Item 3: saturated subset
\begin{definition}\label{saturated}
    $C$ is saturated with respect to $p$ iff
        $C\subseteq\dom{p}$ and
        for all $y\in\ran{p}$ if $C$ intersects $\preimg{p}{\{y\}}$ then
            $\preimg{p}{\{y\}}\subseteq C$.
\end{definition}

% from §22 Item 4: saturated sets are complete inverse images
\begin{lemma}\label{saturated_iff_preimg}
    Suppose $p$ is a function.
    Then $C$ is saturated with respect to $p$ iff there exists $B$ such that
        $C = \preimg{p}{B}$.
\end{lemma}
\begin{proof}
    Show that if $C$ is saturated with respect to $p$ then there exists $B$ such that
        $C = \preimg{p}{B}$.
    \begin{subproof}
        Assume $C$ is saturated with respect to $p$.
        Let $B = \img{p}{C}$.
        Show that for all $x$ such that $x\in C$ we have $x\in\preimg{p}{B}$.
        \begin{subproof}
            Follows by \cref{saturated,subseteq,img_of_function_intro,inter_intro,function_apply_elim,preimg_iff}.
        \end{subproof}
        $C\subseteq \preimg{p}{B}$ by \cref{subseteq}.
        Show that for all $x$ such that $x\in\preimg{p}{B}$ we have $x\in C$.
        \begin{subproof}
            Follows by
                \cref{preimg_iff,img_of_function_elim,inter_elim_right,inter_elim_left,function_apply_iff,singleton_iff,inter_intro,emptyset,intersects,ran_iff,saturated,subseteq}.
        \end{subproof}
        $\preimg{p}{B}\subseteq C$ by \cref{subseteq}.
        $C = \preimg{p}{B}$ by \cref{subseteq_antisymmetric}.
    \end{subproof}
    Show that if there exists $B$ such that $C = \preimg{p}{B}$ then
        $C$ is saturated with respect to $p$.
    \begin{subproof}
        Assume that there exists $B$ such that $C = \preimg{p}{B}$.
        Take $B$ such that $C = \preimg{p}{B}$.
        Show that $C\subseteq\dom{p}$.
        \begin{subproof}
            Follows by \cref{preimg_subseteq_dom,subseteq}.
        \end{subproof}
        For all $y\in\ran{p}$ if $C$ intersects $\preimg{p}{\{y\}}$ then
            there exists $c$ such that $c\in C\inter \preimg{p}{\{y\}}$
            by \cref{intersects,nonempty_inhabited}.
        Show that for all $y\in\ran{p}$ if $C$ intersects $\preimg{p}{\{y\}}$ then
            $\preimg{p}{\{y\}}\subseteq C$.
        \begin{subproof}
            Follows by
                \cref{preimg_iff,subseteq,inter_elim_left,inter_elim_right,singleton_iff,function_rightunique,rightunique}.
        \end{subproof}
        $C$ is saturated with respect to $p$ by \cref{saturated}.
    \end{subproof}
\end{proof}

% from §22 helper lemma 1: image of a preimage for surjections
\begin{lemma}\label{img_preimg_surjection}
    Suppose $p$ is a surjection from $X$ to $Y$.
    Suppose $B\subseteq Y$.
    Then $\img{p}{\preimg{p}{B}} = B$.
\end{lemma}
\begin{proof}
    $p\in\surj{X}{Y}$ and $p\in\funs{X}{Y}$ and $p$ is a function
        by \cref{surj,funs_is_function}.
    Show that $\img{p}{\preimg{p}{B}}\subseteq B$.
    \begin{subproof}
        Follows by \cref{img_preimg_subseteq}.
    \end{subproof}
    Show that for all $y\in B$ we have $y\in\img{p}{\preimg{p}{B}}$.
    \begin{subproof}
        Fix $y\in B$.
        $y\in Y$ by \cref{subseteq}.
        Take $x\in X$ such that $p(x) = y$ by \cref{surj}.
        $x\in\dom{p}$ by \cref{funs}.
        $(x,y)\in p$ by \cref{function_apply_iff}.
        $x\in\preimg{p}{B}$ by \cref{preimg_iff}.
        $y\in\img{p}{\preimg{p}{B}}$ by \cref{img_iff}.
    \end{subproof}
    $B\subseteq\img{p}{\preimg{p}{B}}$ by \cref{subseteq}.
    $\img{p}{\preimg{p}{B}} = B$ by \cref{subseteq_antisymmetric}.
\end{proof}

% from §22 Item 5: quotient maps and saturated open sets
\begin{theorem}\label{quotient_map_saturated_open_iff}
    Suppose $p$ is a surjection from $X$ to $Y$.
    Suppose $T_X$ is a topology on $X$.
    Suppose $T_Y$ is a topology on $Y$.
    Then $p$ is a quotient map from $X$ to $Y$ with respect to $T_X$ and $T_Y$ iff
        $p$ is continuous from $X$ to $Y$ with respect to $T_X$ and $T_Y$ and
        for all $C$ such that
            $C$ is saturated with respect to $p$ and
            $C$ is open in $X$ with respect to $T_X$
        we have $\img{p}{C}$ is open in $Y$ with respect to $T_Y$.
\end{theorem}
\begin{proof}
    Show that if $p$ is a quotient map from $X$ to $Y$ with respect to $T_X$ and $T_Y$ then
        $p$ is continuous from $X$ to $Y$ with respect to $T_X$ and $T_Y$ and
        for all $C$ such that
            $C$ is saturated with respect to $p$ and
            $C$ is open in $X$ with respect to $T_X$
        we have $\img{p}{C}$ is open in $Y$ with respect to $T_Y$.
    \begin{subproof}
        Assume $p$ is a quotient map from $X$ to $Y$ with respect to $T_X$ and $T_Y$.
        Show that $p$ is continuous from $X$ to $Y$ with respect to $T_X$ and $T_Y$.
        \begin{subproof}
            Follows by \cref{quotient_map}.
        \end{subproof}
        Show that for all $C$ such that
            $C$ is saturated with respect to $p$ and
            $C$ is open in $X$ with respect to $T_X$
        we have $\img{p}{C}$ is open in $Y$ with respect to $T_Y$.
        \begin{subproof}
            Fix $C$ such that
                $C$ is saturated with respect to $p$ and
                $C$ is open in $X$ with respect to $T_X$.
            $p\in\funs{X}{Y}$ and $p$ is a function by \cref{surj,funs_is_function}.
            Take $B$ such that $C = \preimg{p}{B}$ by \cref{saturated_iff_preimg}.
            Let $B_1 = B\inter Y$.
            $\preimg{p}{B_1} = \preimg{p}{B}$ by
                \cref{inter_lower_left,preimg_subseteq,preimg_iff,ran_iff,surj,funs_ran,subseteq,inter_intro,subseteq_antisymmetric}.
            $\img{p}{C}$ is open in $Y$ with respect to $T_Y$
                by \cref{subseteq,quotient_map,inter_lower_right,img_preimg_surjection,open_in}.
        \end{subproof}
    \end{subproof}
    Show that if
        $p$ is continuous from $X$ to $Y$ with respect to $T_X$ and $T_Y$ and
        for all $C$ such that
            $C$ is saturated with respect to $p$ and
            $C$ is open in $X$ with respect to $T_X$
        we have $\img{p}{C}$ is open in $Y$ with respect to $T_Y$
    then $p$ is a quotient map from $X$ to $Y$ with respect to $T_X$ and $T_Y$.
    \begin{subproof}
        Assume $p$ is continuous from $X$ to $Y$ with respect to $T_X$ and $T_Y$.
        Assume that for all $C$ such that
            $C$ is saturated with respect to $p$ and
            $C$ is open in $X$ with respect to $T_X$
        we have $\img{p}{C}$ is open in $Y$ with respect to $T_Y$.
        $p$ is a function by \cref{surj,funs_is_function}.
        Show that for all $U$ such that $U\subseteq Y$ if $U$ is open in $Y$ with respect to $T_Y$ then
            $\preimg{p}{U}$ is open in $X$ with respect to $T_X$.
        \begin{subproof}
            Follows by \cref{continuous}.
        \end{subproof}
        For all $U$ such that $U\subseteq Y$ we have
            $\preimg{p}{U}$ is saturated with respect to $p$ by \cref{saturated_iff_preimg}.
        Show that for all $U$ such that $U\subseteq Y$ if $\preimg{p}{U}$ is open in $X$ with respect to $T_X$ then
            $U$ is open in $Y$ with respect to $T_Y$.
        \begin{subproof}
            For all $U$ such that $U\subseteq Y$ if $\preimg{p}{U}$ is open in $X$ with respect to $T_X$ then
                $\img{p}{\preimg{p}{U}}$ is open in $Y$ with respect to $T_Y$.
            Follows by \cref{img_preimg_surjection}.
        \end{subproof}
        Show that for all $U$ such that $U\subseteq Y$ we have
            $U$ is open in $Y$ with respect to $T_Y$ iff
            $\preimg{p}{U}$ is open in $X$ with respect to $T_X$.
        \begin{subproof}
            Trivial.
        \end{subproof}
        $p$ is a quotient map from $X$ to $Y$ with respect to $T_X$ and $T_Y$ by \cref{quotient_map}.
    \end{subproof}
\end{proof}

% from §22 Item 6: open map
\begin{definition}\label{open_map}
    $f$ is an open map from $X$ to $Y$ with respect to $T$ and $T_1$ iff
        $f$ is a function from $X$ to $Y$ and
        for all $U$ such that $U$ is open in $X$ with respect to $T$ we have
            $\img{f}{U}$ is open in $Y$ with respect to $T_1$.
\end{definition}

% from §22 Item 7: closed map
\begin{definition}\label{closed_map}
    $f$ is a closed map from $X$ to $Y$ with respect to $T$ and $T_1$ iff
        $f$ is a function from $X$ to $Y$ and
        for all $A$ such that $A$ is closed in $X$ with respect to $T$ we have
            $\img{f}{A}$ is closed in $Y$ with respect to $T_1$.
\end{definition}

% from §22 Item 8: open or closed surjections are quotient maps
\begin{lemma}\label{open_or_closed_surjection_quotient}
    Suppose $p$ is a surjection from $X$ to $Y$.
    Suppose $T_X$ is a topology on $X$.
    Suppose $T_Y$ is a topology on $Y$.
    Suppose $p$ is continuous from $X$ to $Y$ with respect to $T_X$ and $T_Y$.
    Suppose $p$ is an open map from $X$ to $Y$ with respect to $T_X$ and $T_Y$
        or $p$ is a closed map from $X$ to $Y$ with respect to $T_X$ and $T_Y$.
    Then $p$ is a quotient map from $X$ to $Y$ with respect to $T_X$ and $T_Y$.
\end{lemma}
\begin{proof}
    Show that if $p$ is an open map from $X$ to $Y$ with respect to $T_X$ and $T_Y$ then
        $p$ is a quotient map from $X$ to $Y$ with respect to $T_X$ and $T_Y$.
    \begin{subproof}
        Follows by \cref{continuous,open_map,img_preimg_surjection,quotient_map}.
    \end{subproof}
    Show that if $p$ is a closed map from $X$ to $Y$ with respect to $T_X$ and $T_Y$ then
        $p$ is a quotient map from $X$ to $Y$ with respect to $T_X$ and $T_Y$.
    \begin{subproof}
        Follows by \cref{closed_in,continuous,preimg_setminus_function,preimg_subseteq_dom,funs,subseteq,closed_map,img_preimg_surjection,quotient_map_closed_iff}.
    \end{subproof}
    $p$ is a quotient map from $X$ to $Y$ with respect to $T_X$ and $T_Y$.
\end{proof}

% from §22 Item 9: quotient topology
\begin{definition}\label{quotient_topology}
    $\quotientTopology{p}{T}{X}{A}
        = \{U\in\pow{A} \mid \text{$\preimg{p}{U}$ is open in $X$ with respect to $T$}\}$.
\end{definition}

% from §22 Item 10: quotient topology is a topology
\begin{lemma}\label{quotient_topology_is_topology}
    Assume $p$ is a function from $X$ to $A$.
    Assume $T$ is a topology on $X$.
    Then $\quotientTopology{p}{T}{X}{A}$ is a topology on $A$.
\end{lemma}
\begin{proof}
    Let $T_Q = \quotientTopology{p}{T}{X}{A}$.
    $p\in\funs{X}{A}$ and $\dom{p} = X$ and $p$ is a function by \cref{funs,funs_is_function}.
    Show that $T_Q\subseteq\pow{A}$.
    \begin{subproof}
        Follows by \cref{quotient_topology,subseteq}.
    \end{subproof}
    $\preimg{p}{\emptyset}\subseteq\emptyset$ by \cref{preimg_iff,emptyset,subseteq}.
    Show that $\emptyset\in T_Q$.
    \begin{subproof}
        Follows by \cref{emptyset_subseteq,subseteq_antisymmetric,topology_on,open_in,pow_iff,quotient_topology}.
    \end{subproof}
    $\preimg{p}{A} = X$ by
        \cref{preimg_subseteq_dom,funs,subseteq,dom_iff,ran_iff,funs_ran,subseteq_antisymmetric,preimg_iff}.
    Show that $A\in T_Q$.
    \begin{subproof}
        Follows by \cref{topology_on,open_in,subseteq_refl,pow_iff,quotient_topology}.
    \end{subproof}
    Show that for all $M$ such that $M\subseteq T_Q$ we have $\unions{M}\in T_Q$.
    \begin{subproof}
        Fix $M$ such that $M\subseteq T_Q$.
        Let $N = \{V\in T \mid \text{there exists $U\in M$ such that $V = \preimg{p}{U}$}\}$.
        $\unions{N}\in T$ by \cref{subseteq,topology_on}.
        Show that for all $x$ such that $x\in\preimg{p}{\unions{M}}$ we have $x\in\unions{N}$.
        \begin{subproof}
            Fix $x$ such that $x\in\preimg{p}{\unions{M}}$.
            Take $y\in\unions{M}$ such that $(x,y)\in p$ by \cref{preimg_iff}.
            Take $U\in M$ such that $y\in U$ by \cref{unions_iff}.
            $U\in T_Q$ by \cref{subseteq}.
            Hence $\preimg{p}{U}$ is open in $X$ with respect to $T$ by \cref{quotient_topology}.
            $\preimg{p}{U}\in T$ by \cref{open_in}.
            $\preimg{p}{U}\in N$.
            $x\in\preimg{p}{U}$ by \cref{preimg_iff}.
            Therefore $x\in\unions{N}$ by \cref{unions_iff}.
        \end{subproof}
        Hence $\preimg{p}{\unions{M}}\subseteq \unions{N}$ by \cref{subseteq}.
        Show that for all $x$ such that $x\in\unions{N}$ we have $x\in\preimg{p}{\unions{M}}$.
        \begin{subproof}
            Fix $x$ such that $x\in\unions{N}$.
            Take $V\in N$ such that $x\in V$ by \cref{unions_iff}.
            Take $U\in M$ such that $V = \preimg{p}{U}$.
            Take $y\in U$ such that $(x,y)\in p$ by \cref{preimg_iff}.
            $y\in\unions{M}$ by \cref{unions_iff}.
            Therefore $x\in\preimg{p}{\unions{M}}$ by \cref{preimg_iff}.
        \end{subproof}
        Therefore $\unions{N}\subseteq \preimg{p}{\unions{M}}$ by \cref{subseteq}.
        Hence $\preimg{p}{\unions{M}} = \unions{N}$ by \cref{subseteq,subseteq_antisymmetric}.
        Therefore $\unions{M}\in T_Q$ by \cref{open_in,subseteq,quotient_topology,unions_subseteq_of_powerset_is_subseteq,pow_iff}.
    \end{subproof}
    Show that for all $U,V$ such that $U\in T_Q$ and $V\in T_Q$ we have $U\inter V\in T_Q$.
    \begin{subproof}
        Follows by \cref{quotient_topology,open_in,topology_on,preimg_inter_function,pow_iff,inter_subseteq}.
    \end{subproof}
    Show that $T_Q$ is closed under arbitrary unions.
    \begin{subproof}
        Trivial.
    \end{subproof}
    Show that $T_Q$ is closed under binary intersections.
    \begin{subproof}
        Trivial.
    \end{subproof}
    Therefore $\quotientTopology{p}{T}{X}{A}$ is a topology on $A$ by \cref{topology_on}.
\end{proof}

% from §22 Item 11: quotient topology induces a quotient map
\begin{lemma}\label{quotient_topology_quotient_map}
    Suppose $p$ is a surjection from $X$ to $A$.
    Suppose $T$ is a topology on $X$.
    Let $T_Q = \quotientTopology{p}{T}{X}{A}$.
    Then $p$ is a quotient map from $X$ to $A$ with respect to $T$ and $T_Q$.
\end{lemma}
\begin{proof}
    $p$ is a function from $X$ to $A$ and $T_Q$ is a topology on $A$
        by \cref{surj,funs,quotient_topology_is_topology}.
    Show that for all $V$ such that $V$ is open in $A$ with respect to $T_Q$
        we have $\preimg{p}{V}$ is open in $X$ with respect to $T$.
    \begin{subproof}
        Follows by \cref{open_in,quotient_topology}.
    \end{subproof}
    $p$ is continuous from $X$ to $A$ with respect to $T$ and $T_Q$ by \cref{continuous}.
    Show that for all $U$ such that $U\subseteq A$ if $\preimg{p}{U}$ is open in $X$ with respect to $T$
        then $U$ is open in $A$ with respect to $T_Q$.
    \begin{subproof}
        Follows by \cref{pow_iff,quotient_topology,open_in}.
    \end{subproof}
    Therefore $p$ is a quotient map from $X$ to $A$ with respect to $T$ and $T_Q$
        by \cref{quotient_map}.
\end{proof}
% from §22 Item 12: quotient topology is unique
\begin{lemma}\label{quotient_topology_unique}
    Suppose $p$ is a surjection from $X$ to $A$.
    Suppose $T$ is a topology on $X$.
    Suppose $T_A$ is a topology on $A$.
    Suppose $p$ is a quotient map from $X$ to $A$ with respect to $T$ and $T_A$.
    Then $T_A = \quotientTopology{p}{T}{X}{A}$.
\end{lemma}
\begin{proof}
    Let $T_Q = \quotientTopology{p}{T}{X}{A}$.
    Show that $T_A\subseteq T_Q$.
    \begin{subproof}
        Follows by \cref{open_in,topology_on,subseteq,continuous,quotient_map,quotient_topology}.
    \end{subproof}
    Show that $T_Q\subseteq T_A$.
    \begin{subproof}
        Follows by \cref{quotient_topology,pow_iff,quotient_map,open_in,subseteq}.
    \end{subproof}
    Therefore $T_A = T_Q$ by \cref{subseteq_antisymmetric}.
\end{proof}

% from §22 helper definition 1: partition map
\begin{definition}\label{partition_map}
    $\partitionMap{P}{X} = \{w\in X\times P \mid \fst{w}\in \snd{w}\}$.
\end{definition}

% from §22 Item 13: quotient space
\begin{definition}\label{quotient_space}
    $P$ is a quotient space of $X$ with respect to $T$ and $T_P$ iff
        $P$ is a partition of $X$ and
        $T$ is a topology on $X$ and
        $T_P = \quotientTopology{\partitionMap{P}{X}}{T}{X}{P}$.
\end{definition}

% from §22 helper lemma 2: restriction preimage for saturated subsets
\begin{lemma}\label{preimg_restrl_saturated}
    Suppose $p$ is a surjection from $X$ to $Y$.
    Suppose $A\subseteq X$.
    Suppose $A$ is saturated with respect to $p$.
    Let $q = \restrl{p}{A}$.
    Suppose $V\subseteq \img{p}{A}$.
    Then $\preimg{q}{V} = \preimg{p}{V}$.
\end{lemma}
\begin{proof}
    Show that $\preimg{q}{V}\subseteq \preimg{p}{V}$.
    \begin{subproof}
        Follows by \cref{preimg_iff,restrl_iff,subseteq}.
    \end{subproof}
    Show that for all $x$ such that $x\in\preimg{p}{V}$ we have $x\in\preimg{q}{V}$.
    \begin{subproof}
        Follows by \cref{preimg_iff,img_iff,subseteq,intersects,emptyset,inter_intro,ran_iff,saturated,singleton_iff,restrl_iff}.
    \end{subproof}
    Hence $\preimg{p}{V}\subseteq \preimg{q}{V}$ by \cref{subseteq}.
    Therefore $\preimg{q}{V} = \preimg{p}{V}$ by \cref{subseteq,subseteq_antisymmetric}.
\end{proof}

% from §22 helper lemma 3: image of intersection with saturated set
\begin{lemma}\label{img_inter_saturated}
    Suppose $p$ is a function.
    Suppose $A\subseteq \dom{p}$.
    Suppose $A$ is saturated with respect to $p$.
    Suppose $U\subseteq \dom{p}$.
    Then $\img{p}{U\inter A} = \img{p}{U}\inter \img{p}{A}$.
\end{lemma}
\begin{proof}
    Show that $\img{p}{U\inter A}\subseteq \img{p}{U}\inter \img{p}{A}$.
    \begin{subproof}
        Follows by \cref{img_iff,inter,inter_intro,subseteq}.
    \end{subproof}
    Show that for all $y$ such that $y\in\img{p}{U}\inter \img{p}{A}$ we have $y\in\img{p}{U\inter A}$.
    \begin{subproof}
        Fix $y$ such that $y\in\img{p}{U}\inter \img{p}{A}$.
        $y\in\img{p}{U}$ and $y\in\img{p}{A}$ by \cref{inter}.
        Take $x$ such that $x\in U$ and $(x,y)\in p$ by \cref{img_iff}.
        Take $a$ such that $a\in A$ and $(a,y)\in p$ by \cref{img_iff}.
        $y\in\ran{p}$ by \cref{ran_iff}.
        $a\in\preimg{p}{\{y\}}$ by \cref{preimg_iff,singleton_iff}.
        Hence $A$ intersects $\preimg{p}{\{y\}}$ by \cref{inter_intro,emptyset,intersects}.
        $\preimg{p}{\{y\}}\subseteq A$ by \cref{saturated}.
        $x\in\preimg{p}{\{y\}}$ by \cref{preimg_iff,singleton_iff}.
        $x\in A$ by \cref{subseteq}.
        $x\in U\inter A$ by \cref{inter_intro}.
        Therefore $y\in\img{p}{U\inter A}$ by \cref{img_iff}.
    \end{subproof}
    Hence $\img{p}{U}\inter \img{p}{A}\subseteq \img{p}{U\inter A}$ by \cref{subseteq}.
    Therefore $\img{p}{U\inter A} = \img{p}{U}\inter \img{p}{A}$ by \cref{subseteq,subseteq_antisymmetric}.
\end{proof}

% from §22 helper lemma 4: preimage of image of a saturated set
\begin{lemma}\label{preimg_img_saturated}
    Suppose $p$ is a function.
    Suppose $A\subseteq \dom{p}$.
    Suppose $A$ is saturated with respect to $p$.
    Then $\preimg{p}{\img{p}{A}} = A$.
\end{lemma}
\begin{proof}
    Show that $A\subseteq \preimg{p}{\img{p}{A}}$.
    \begin{subproof}
        Follows by \cref{subseteq,function_apply_elim,img_iff,preimg_iff}.
    \end{subproof}
    Show that for all $x$ such that $x\in\preimg{p}{\img{p}{A}}$ we have $x\in A$.
    \begin{subproof}
        Follows by \cref{preimg_iff,img_iff,singleton_iff,inter_intro,emptyset,intersects,ran_iff,saturated,subseteq}.
    \end{subproof}
    Hence $\preimg{p}{\img{p}{A}}\subseteq A$ by \cref{subseteq}.
    Therefore $\preimg{p}{\img{p}{A}} = A$ by \cref{subseteq,subseteq_antisymmetric}.
\end{proof}

% from §22 Item 14: restriction to a saturated subspace
\begin{theorem}\label{quotient_map_restriction_open_or_closed}
    Suppose $p$ is a quotient map from $X$ to $Y$ with respect to $T_X$ and $T_Y$.
    Suppose $A$ is a subspace of $X$ with respect to $T_X$.
    Suppose $A$ is saturated with respect to $p$.
    Let $q = \restrl{p}{A}$.
    Let $T_A = \subspaceTopology{T_X}{A}$.
    Let $T_B = \subspaceTopology{T_Y}{\img{p}{A}}$.
    Suppose $A$ is open in $X$ with respect to $T_X$
        or $A$ is closed in $X$ with respect to $T_X$.
    Then $q$ is a quotient map from $A$ to $\img{p}{A}$ with respect to $T_A$ and $T_B$.
\end{theorem}
    \begin{proof}
        $p\in\surj{X}{Y}$ and $p\in\funs{X}{Y}$ and $p$ is a function
            by \cref{quotient_map,surj,funs_is_function}.
    $p$ is continuous from $X$ to $Y$ with respect to $T_X$ and $T_Y$ by \cref{quotient_map}.
    Therefore $T_X$ is a topology on $X$ and $T_Y$ is a topology on $Y$ by \cref{continuous}.
    $A\subseteq X$ and $A$ is a subspace of $X$ with respect to $T_X$ by \cref{subspace_of}.
    We have $\img{p}{A}\subseteq Y$ and $\img{p}{A}$ is a subspace of $Y$ with respect to $T_Y$
        by \cref{img_subseteq_ran,funs_ran,subseteq_transitive,subspace_of}.
    Therefore $T_A$ is a topology on $A$ and $T_B$ is a topology on $\img{p}{A}$
        by \cref{subspace_topology_is_topology,subspace_of}.
    We have $q$ is a function by \cref{restrl_of_function_is_function}.
    $q$ is right-unique and $q$ is a relation.
        $A\subseteq \dom{p}$ by \cref{funs,subseteq_transitive,subspace_of}.
        $\dom{q} = A$ and $\ran{q} = \img{p}{A}$
            by \cref{dom_of_restrl_eq_inter,funs,inter_absorb_supseteq_left,subspace_of,restrl_ran}.
        We have $\dom{q}\subseteq A$ and $\ran{q}\subseteq \img{p}{A}$ by \cref{subseteq_refl}.
        $q\in\rels{A}{\img{p}{A}}$ by \cref{rels_intro_dom_and_ran}.
        $q\in\funs{A}{\img{p}{A}}$ by \cref{funs}.
        Therefore $q\in\surj{A}{\img{p}{A}}$ by \cref{funs_surj_iff,restrl_ran}.
        We have $q$ is continuous from $A$ to $Y$ with respect to $T_A$ and $T_Y$
            by \cref{continuous_restrict_domain}.
    Therefore $q$ is continuous from $A$ to $\img{p}{A}$ with respect to $T_A$ and $T_B$
        by \cref{restrl_img,inter_absorb_supseteq_left,subseteq,subseteq_refl,continuous_restrict_range}.
    \begin{byCase}
    \caseOf{$A$ is open in $X$ with respect to $T_X$.}
        Show that for all $V$ such that $V\subseteq \img{p}{A}$ if $\preimg{q}{V}$ is open in $A$ with respect to $T_A$ then
            $V$ is open in $\img{p}{A}$ with respect to $T_B$.
        \begin{subproof}
            Follows by \cref{open_in_subspace_open_in_space,subspace_of,preimg_restrl_saturated,open_in,subseteq_transitive,quotient_map,subspace_topology,inter_absorb_supseteq_left}.
        \end{subproof}
        Therefore $q$ is a quotient map from $A$ to $\img{p}{A}$ with respect to $T_A$ and $T_B$ by \cref{quotient_map}.
    \caseOf{$A$ is closed in $X$ with respect to $T_X$.}
        Show that for all $B$ such that $B\subseteq \img{p}{A}$ we have
            $B$ is closed in $\img{p}{A}$ with respect to $T_B$ iff
            $\preimg{q}{B}$ is closed in $A$ with respect to $T_A$.
        \begin{subproof}
            Fix $B$ such that $B\subseteq \img{p}{A}$.
            Show that if $B$ is closed in $\img{p}{A}$ with respect to $T_B$ then
                $\preimg{q}{B}$ is closed in $A$ with respect to $T_A$.
            \begin{subproof}
                Assume $B$ is closed in $\img{p}{A}$ with respect to $T_B$.
                Take $C$ such that $C$ is closed in $Y$ with respect to $T_Y$ and
                    $B = C\inter \img{p}{A}$ by \cref{closed_in_subspace_iff,subspace_of}.
                Hence $\preimg{q}{B} = \preimg{p}{C}\inter A$ and
                    $\preimg{p}{C}$ is closed in $X$ with respect to $T_X$
                    by \cref{closed_in,continuous,preimg_setminus_function,preimg_subseteq_dom,funs,subseteq,preimg_restrl_saturated,preimg_inter_function,subseteq_transitive,subspace_of,preimg_img_saturated}.
                Therefore $\preimg{q}{B}$ is closed in $A$ with respect to $T_A$ by \cref{closed_in_subspace_iff,subspace_of}.
            \end{subproof}
            Show that if $\preimg{q}{B}$ is closed in $A$ with respect to $T_A$ then
                $B$ is closed in $\img{p}{A}$ with respect to $T_B$.
            \begin{subproof}
                Assume $\preimg{q}{B}$ is closed in $A$ with respect to $T_A$.
                Take $C$ such that $C$ is closed in $X$ with respect to $T_X$ and
                    $\preimg{q}{B} = C\inter A$ by \cref{closed_in_subspace_iff,subspace_of}.
                $\preimg{q}{B} = \preimg{p}{B}$ by \cref{preimg_restrl_saturated}.
                Let $M = \{X\setminus C, X\setminus A\}$.
                For all $U$ such that $U\in M$ we have $U\in T_X$ by \cref{upair_elim,closed_in,open_in}.
                $(X\setminus C)\union (X\setminus A)$ is open in $X$ with respect to $T_X$
                    by \cref{union_as_unions,subseteq,topology_on,open_in}.
                $X\setminus (C\inter A)$ is open in $X$ with respect to $T_X$
                    by \cref{setminus_inter,open_in}.
                Hence $C\inter A$ is closed in $X$ with respect to $T_X$
                    by \cref{closed_in,inter_lower_left,subseteq_transitive}.
                Therefore $\preimg{p}{B}$ is closed in $X$ with respect to $T_X$ by \cref{subseteq}.
                $B$ is closed in $Y$ with respect to $T_Y$
                    by \cref{subseteq_transitive,quotient_map_closed_iff,quotient_map}.
                Hence $B$ is closed in $\img{p}{A}$ with respect to $T_B$
                    by \cref{inter_absorb_supseteq_left,inter_comm,closed_in_subspace_iff,subspace_of}.
            \end{subproof}
            Therefore $B$ is closed in $\img{p}{A}$ with respect to $T_B$ iff
                $\preimg{q}{B}$ is closed in $A$ with respect to $T_A$.
        \end{subproof}
        Therefore $q$ is a quotient map from $A$ to $\img{p}{A}$ with respect to $T_A$ and $T_B$
            by \cref{quotient_map_closed_iff}.
    \end{byCase}
\end{proof}

% from §22 Item 15: restriction of a quotient map with open or closed map
\begin{theorem}\label{quotient_map_restriction_open_or_closed_map}
    Suppose $p$ is a quotient map from $X$ to $Y$ with respect to $T_X$ and $T_Y$.
    Suppose $A$ is a subspace of $X$ with respect to $T_X$.
    Suppose $A$ is saturated with respect to $p$.
    Let $q = \restrl{p}{A}$.
    Let $T_A = \subspaceTopology{T_X}{A}$.
    Let $T_B = \subspaceTopology{T_Y}{\img{p}{A}}$.
    Suppose $p$ is an open map from $X$ to $Y$ with respect to $T_X$ and $T_Y$
        or $p$ is a closed map from $X$ to $Y$ with respect to $T_X$ and $T_Y$.
    Then $q$ is a quotient map from $A$ to $\img{p}{A}$ with respect to $T_A$ and $T_B$.
\end{theorem}
    \begin{proof}
        $p\in\surj{X}{Y}$ and $p\in\funs{X}{Y}$ and $p$ is a function
            by \cref{quotient_map,surj,funs_is_function}.
    $p$ is continuous from $X$ to $Y$ with respect to $T_X$ and $T_Y$ by \cref{quotient_map}.
    Therefore $T_X$ is a topology on $X$ and $T_Y$ is a topology on $Y$ by \cref{continuous}.
    $A\subseteq X$ and $A$ is a subspace of $X$ with respect to $T_X$ by \cref{subspace_of}.
    We have $\img{p}{A}\subseteq Y$ and $\img{p}{A}$ is a subspace of $Y$ with respect to $T_Y$
        by \cref{img_subseteq_ran,funs_ran,subseteq_transitive,subspace_of}.
    Therefore $T_A$ is a topology on $A$ and $T_B$ is a topology on $\img{p}{A}$
        by \cref{subspace_topology_is_topology,subspace_of}.
    We have $q$ is a function by \cref{restrl_of_function_is_function}.
    $q$ is right-unique and $q$ is a relation.
        $A\subseteq \dom{p}$ by \cref{funs,subseteq_transitive,subspace_of}.
        $\dom{q} = A$ and $\ran{q} = \img{p}{A}$
            by \cref{dom_of_restrl_eq_inter,funs,inter_absorb_supseteq_left,subspace_of,restrl_ran}.
        We have $\dom{q}\subseteq A$ and $\ran{q}\subseteq \img{p}{A}$ by \cref{subseteq_refl}.
        $q\in\rels{A}{\img{p}{A}}$ by \cref{rels_intro_dom_and_ran}.
        $q\in\funs{A}{\img{p}{A}}$ by \cref{funs}.
        Therefore $q\in\surj{A}{\img{p}{A}}$ by \cref{funs_surj_iff,restrl_ran}.
        We have $q$ is continuous from $A$ to $Y$ with respect to $T_A$ and $T_Y$
            by \cref{continuous_restrict_domain}.
    Therefore $q$ is continuous from $A$ to $\img{p}{A}$ with respect to $T_A$ and $T_B$
        by \cref{restrl_img,inter_absorb_supseteq_left,subseteq,subseteq_refl,continuous_restrict_range}.
    \begin{byCase}
    \caseOf{$p$ is an open map from $X$ to $Y$ with respect to $T_X$ and $T_Y$.}
        Show that for all $U$ such that $U$ is open in $A$ with respect to $T_A$ we have
            $\img{q}{U}$ is open in $\img{p}{A}$ with respect to $T_B$.
        \begin{subproof}
            Fix $U$ such that $U$ is open in $A$ with respect to $T_A$.
            $U\in T_A$ and $U\subseteq A$ by \cref{open_in,topology_on,subseteq,pow_iff}.
            Take $W\in T_X$ such that $U = A\inter W$ by \cref{subspace_topology}.
            $\img{q}{U} = \img{p}{A\inter U}$ by \cref{funs,subseteq_transitive,subspace_of,restrl_img}.
            $\img{q}{U} = \img{p}{A\inter W}$ by \cref{inter_absorb_supseteq_left,subseteq}.
            Hence $\img{q}{U} = \img{p}{W}\inter \img{p}{A}$
                by \cref{inter_comm,topology_on,subseteq,pow_iff,funs,subseteq_transitive,img_inter_saturated}.
            Therefore $\img{q}{U}$ is open in $\img{p}{A}$ with respect to $T_B$
                by \cref{inter_comm,open_in,open_map,subspace_topology}.
        \end{subproof}
        Hence $q$ is an open map from $A$ to $\img{p}{A}$ with respect to $T_A$ and $T_B$
            by \cref{surj_to_fun,funs,open_map}.
        Therefore $q$ is a quotient map from $A$ to $\img{p}{A}$ with respect to $T_A$ and $T_B$
            by \cref{open_or_closed_surjection_quotient}.
    \caseOf{$p$ is a closed map from $X$ to $Y$ with respect to $T_X$ and $T_Y$.}
        Show that for all $U$ such that $U$ is closed in $A$ with respect to $T_A$ we have
            $\img{q}{U}$ is closed in $\img{p}{A}$ with respect to $T_B$.
        \begin{subproof}
            Follows by \cref{closed_in_subspace_iff,subspace_of,closed_in,funs,subseteq_transitive,restrl_img,inter_absorb_supseteq_left,subseteq,img_inter_saturated,closed_map}.
        \end{subproof}
        Hence $q$ is a closed map from $A$ to $\img{p}{A}$ with respect to $T_A$ and $T_B$
            by \cref{surj_to_fun,funs,closed_map}.
        Therefore $q$ is a quotient map from $A$ to $\img{p}{A}$ with respect to $T_A$ and $T_B$
            by \cref{open_or_closed_surjection_quotient}.
    \end{byCase}
\end{proof}

% from §22 Item 16: composite of quotient maps
\begin{lemma}\label{quotient_map_composite}
    Suppose $p$ is a quotient map from $X$ to $Y$ with respect to $T_X$ and $T_Y$.
    Suppose $q$ is a quotient map from $Y$ to $Z$ with respect to $T_Y$ and $T_Z$.
    Then $q\circ p$ is a quotient map from $X$ to $Z$ with respect to $T_X$ and $T_Z$.
\end{lemma}
\begin{proof}
    $q\circ p\in\surj{X}{Z}$ and $q\circ p$ is continuous from $X$ to $Z$
        with respect to $T_X$ and $T_Z$ by \cref{quotient_map,surjections_circ,continuous_composite}.
    Show that for all $U$ such that $U\subseteq Z$ we have
        $\preimg{(q\circ p)}{U} = \preimg{p}{\preimg{q}{U}}$.
    \begin{subproof}
        Fix $U$ such that $U\subseteq Z$.
        Show that $\preimg{(q\circ p)}{U}\subseteq \preimg{p}{\preimg{q}{U}}$.
        \begin{subproof}
            Show that for all $x$ such that $x\in\preimg{(q\circ p)}{U}$ we have $x\in\preimg{p}{\preimg{q}{U}}$.
            \begin{subproof}
                Fix $x$ such that $x\in\preimg{(q\circ p)}{U}$.
                Take $z\in U$ such that $x\mathrel{(q\circ p)} z$ by \cref{preimg_iff}.
                Take $y$ such that $x\mathrel{p} y\mathrel{q} z$ by \cref{circ_iff}.
                Hence $y\in\preimg{q}{U}$ by \cref{preimg_iff}.
                $x\in\preimg{p}{\preimg{q}{U}}$ by \cref{preimg_iff}.
            \end{subproof}
            Follows by \cref{subseteq}.
        \end{subproof}
        Show that $\preimg{p}{\preimg{q}{U}}\subseteq \preimg{(q\circ p)}{U}$.
        \begin{subproof}
            Show that for all $x$ such that $x\in\preimg{p}{\preimg{q}{U}}$ we have $x\in\preimg{(q\circ p)}{U}$.
            \begin{subproof}
                Fix $x$ such that $x\in\preimg{p}{\preimg{q}{U}}$.
                Take $y\in\preimg{q}{U}$ such that $x\mathrel{p} y$ by \cref{preimg_iff}.
                Take $z\in U$ such that $y\mathrel{q} z$ by \cref{preimg_iff}.
                Hence $x\mathrel{(q\circ p)} z$ by \cref{circ_iff}.
                $x\in\preimg{(q\circ p)}{U}$ by \cref{preimg_iff}.
            \end{subproof}
            Follows by \cref{subseteq}.
        \end{subproof}
        Therefore $\preimg{(q\circ p)}{U} = \preimg{p}{\preimg{q}{U}}$ by \cref{subseteq_antisymmetric}.
    \end{subproof}
    Show that for all $U$ such that $U\subseteq Z$ if $\preimg{(q\circ p)}{U}$ is open in $X$ with respect to $T_X$ then
        $\preimg{p}{\preimg{q}{U}}$ is open in $X$ with respect to $T_X$.
    \begin{subproof}
        Follows by \cref{open_in}.
    \end{subproof}
    Show that for all $U$ such that $U\subseteq Z$ if $\preimg{(q\circ p)}{U}$ is open in $X$ with respect to $T_X$ then
        $U$ is open in $Z$ with respect to $T_Z$.
    \begin{subproof}
        Follows by \cref{preimg_subseteq_dom,quotient_map,surj_to_fun,funs,subseteq}.
    \end{subproof}
    Therefore $q\circ p$ is a quotient map from $X$ to $Z$ with respect to $T_X$ and $T_Z$ by \cref{quotient_map}.
\end{proof}

% from §22 helper definition 2: induced map
\begin{definition}\label{induced_map}
    $\inducedMap{p}{g}{X}{Y}{Z} =
        \{w\in Y\times Z \mid \snd{w}\in \img{g}{\preimg{p}{\{\fst{w}\}}}\}$.
\end{definition}

% from §22 Item 17: induced map exists
\begin{lemma}\label{quotient_map_induced_map}
    Suppose $p$ is a quotient map from $X$ to $Y$ with respect to $T_X$ and $T_Y$.
    Suppose $g$ is a function from $X$ to $Z$.
    Suppose for all $y\in Y$ there exists $z$ such that for all $x\in\preimg{p}{\{y\}}$ we have $g(x) = z$.
    Then there exists $f$ such that
        $f$ is a function from $Y$ to $Z$ and
        $f\circ p = g$.
\end{lemma}
\begin{proof}
    $p\in\surj{X}{Y}$ by \cref{quotient_map}.
    Hence $p\in\funs{X}{Y}$ and $p$ is a function by \cref{surj,funs_is_function}.
    Let $f = \inducedMap{p}{g}{X}{Y}{Z}$.
    For all $y,z_1,z_2$ such that $y\mathrel{f} z_1$ and $y\mathrel{f} z_2$ we have
        $y\in Y$ and $z_1\in \img{g}{\preimg{p}{\{y\}}}$ and $z_2\in \img{g}{\preimg{p}{\{y\}}}$
        by \cref{induced_map,times_elem_is_tuple,pair_eq_iff,fst_eq,snd_eq}.
    For all $y,z_1,z_2$ such that $y\mathrel{f} z_1$ and $y\mathrel{f} z_2$ there exist $x_1, x_2$ such that
        $x_1\in \preimg{p}{\{y\}}$ and $x_2\in \preimg{p}{\{y\}}$ and $x_1\mathrel{g} z_1$ and $x_2\mathrel{g} z_2$
        by \cref{img_iff}.
    For all $y,z_1,z_2$ such that $y\mathrel{f} z_1$ and $y\mathrel{f} z_2$ we have $z_1 = z_2$ by
        \cref{funs,funs_is_function,function_apply_intro}.
    Therefore $f$ is right-unique.
    $\dom{f}\subseteq Y$ by
        \cref{dom,induced_map,times_elem_is_tuple,pair_eq_iff,subseteq}.
    For all $y$ such that $y\in Y$ there exists $x\in X$ such that $p(x) = y$ by \cref{surj}.
    For all $y$ such that $y\in Y$ we have $y\in\dom{f}$ by
        \cref{funs,function_apply_elim,funs_elim,subseteq,times_tuple_intro,singleton_intro,preimg_iff,img_iff,fst_eq,snd_eq,induced_map,dom}.
    Hence $\dom{f} = Y$ by \cref{subseteq,subseteq_antisymmetric}.
    For all $w$ such that $w\in f$ we have $w\in Y\times Z$ by \cref{induced_map}.
    Therefore $f\in\funs{Y}{Z}$ and $f$ is a function from $Y$ to $Z$
        by \cref{subseteq,rels_intro,funs}.
    Show that for all $w$ such that $w\in g$ we have $w\in f\circ p$.
    \begin{subproof}
        Fix $w$ such that $w\in g$.
        Take $x,z$ such that $w = (x,z)$ and $x\in X$ and $z\in Z$
            by \cref{funs,rels_elim,subseteq,times_elem_is_tuple}.
        $(x,z)\in g$.
        $z = g(x)$ by \cref{funs,funs_is_function,function_apply_intro}.
        Take $y$ such that $x\mathrel{p} y$ by \cref{funs,subseteq,dom_iff}.
        $y\in Y$ by \cref{ran_iff,funs_ran,subseteq}.
        $x\in\preimg{p}{\{y\}}$ by \cref{singleton_intro,preimg_iff}.
        $z\in \img{g}{\preimg{p}{\{y\}}}$ by \cref{img_iff}.
        $(y,z)\in f$ by \cref{times_tuple_intro,fst_eq,snd_eq,induced_map}.
        Therefore $w\in f\circ p$ by \cref{circ_iff}.
    \end{subproof}
    Hence $g\subseteq f\circ p$ by \cref{subseteq}.
    Show that for all $w$ such that $w\in f\circ p$ we have $w\in g$.
    \begin{subproof}
        Fix $w$ such that $w\in f\circ p$.
        Take $x,z$ such that $w = (x,z)$ by \cref{circ_is_relation,relation}.
        $(x,z)\in f\circ p$.
        Take $y$ such that $x\mathrel{p} y$ and $y\mathrel{f} z$ by \cref{circ_iff}.
        $z\in \img{g}{\preimg{p}{\{y\}}}$ by \cref{induced_map,fst_eq,snd_eq}.
        $g$ is a function by \cref{funs,funs_is_function}.
        Take $x_1\in \preimg{p}{\{y\}}$ such that $x_1\mathrel{g} z$ by \cref{img_iff}.
        $g(x_1) = z$ by \cref{funs,funs_is_function,function_apply_intro}.
        $x\in\preimg{p}{\{y\}}$ by \cref{singleton_intro,preimg_iff}.
        $y\in Y$ by \cref{funs_ran,ran_iff,subseteq}.
        Take $z_0$ such that for all $x\in\preimg{p}{\{y\}}$ we have $g(x) = z_0$.
        $g(x) = z_0$ and $g(x_1) = z_0$.
        Hence $g(x) = z$.
        Therefore $w\in g$ by \cref{preimg_subseteq_dom,funs,subseteq,function_apply_elim}.
    \end{subproof}
    Hence $f\circ p\subseteq g$ by \cref{subseteq}.
    Therefore $f\circ p = g$ by \cref{subseteq_antisymmetric}.
\end{proof}

% from §22 Item 18: induced map continuity
\begin{lemma}\label{quotient_map_induced_map_continuous}
    Suppose $p$ is a quotient map from $X$ to $Y$ with respect to $T_X$ and $T_Y$.
    Suppose $g$ is a function from $X$ to $Z$.
    Suppose for all $y\in Y$ there exists $z$ such that for all $x\in\preimg{p}{\{y\}}$ we have $g(x) = z$.
    Suppose $f$ is a function from $Y$ to $Z$.
    Suppose $f\circ p = g$.
    Then $f$ is continuous from $Y$ to $Z$ with respect to $T_Y$ and $T_Z$
        iff $g$ is continuous from $X$ to $Z$ with respect to $T_X$ and $T_Z$.
\end{lemma}
\begin{proof}
    Show that if $f$ is continuous from $Y$ to $Z$ with respect to $T_Y$ and $T_Z$ then
        $g$ is continuous from $X$ to $Z$ with respect to $T_X$ and $T_Z$.
    \begin{subproof}
        Follows by \cref{quotient_map,continuous_composite}.
    \end{subproof}
    Show that if $g$ is continuous from $X$ to $Z$ with respect to $T_X$ and $T_Z$ then
        for all $V$ such that $V$ is open in $Z$ with respect to $T_Z$
        we have $\preimg{f}{V}$ is open in $Y$ with respect to $T_Y$.
    \begin{subproof}
        Assume $g$ is continuous from $X$ to $Z$ with respect to $T_X$ and $T_Z$.
        Fix $V$ such that $V$ is open in $Z$ with respect to $T_Z$.
        $\preimg{f}{V}\subseteq Y$ by \cref{preimg_subseteq_dom,funs,subseteq}.
        Show that $\preimg{p}{\preimg{f}{V}} = \preimg{g}{V}$.
        \begin{subproof}
            Show that $\preimg{p}{\preimg{f}{V}}\subseteq \preimg{g}{V}$.
            \begin{subproof}
                For all $x$ such that $x\in\preimg{p}{\preimg{f}{V}}$ we have $x\in\preimg{g}{V}$
                    by \cref{preimg_iff,circ_iff,subseteq}.
            \end{subproof}
            Show that $\preimg{g}{V}\subseteq \preimg{p}{\preimg{f}{V}}$.
            \begin{subproof}
                For all $x$ such that $x\in\preimg{g}{V}$ we have $x\in\preimg{p}{\preimg{f}{V}}$
                    by \cref{preimg_iff,circ_iff,subseteq}.
            \end{subproof}
            Follows by \cref{subseteq_antisymmetric}.
        \end{subproof}
        Hence $\preimg{g}{V}$ is open in $X$ with respect to $T_X$ by \cref{continuous}.
        $\preimg{p}{\preimg{f}{V}}$ is open in $X$ with respect to $T_X$.
        $\preimg{f}{V}$ is open in $Y$ with respect to $T_Y$ by \cref{quotient_map}.
    \end{subproof}
    Show that if $g$ is continuous from $X$ to $Z$ with respect to $T_X$ and $T_Z$ then
        $f$ is continuous from $Y$ to $Z$ with respect to $T_Y$ and $T_Z$.
    \begin{subproof}
        Follows by \cref{continuous,quotient_map}.
    \end{subproof}
\end{proof}

% from §22 helper lemma 7: surjectivity descends along a surjective quotient factor
\begin{lemma}\label{surj_factor_of_surj_composite}
    Suppose $p$ is a surjection from $X$ to $Y$.
    Suppose $g$ is a surjection from $X$ to $Z$.
    Suppose $f$ is a function from $Y$ to $Z$.
    Suppose $f\circ p = g$.
    Then $f$ is a surjection from $Y$ to $Z$.
\end{lemma}
\begin{proof}
    $p\in\funs{X}{Y}$ and $g\in\funs{X}{Z}$ and $p$ is a function and
        $g$ is a function and $f$ is a function by \cref{surj,funs_is_function}.
    For all $z\in Z$ there exists $x\in X$ such that $g(x) = z$ by \cref{surj}.
    Hence for all $z\in Z$ there exists $y\in Y$ such that $f(y) = z$ by
        \cref{funs_type_apply,funs,function_apply_elim,function_apply_iff,circ_iff,function_rightunique}.
    Therefore $f\in\surj{Y}{Z}$ by \cref{surj}.
\end{proof}

% from §22 Item 19: induced map and quotient maps
\begin{lemma}\label{quotient_map_induced_map_quotient}
    Suppose $p$ is a quotient map from $X$ to $Y$ with respect to $T_X$ and $T_Y$.
    Suppose $g$ is a function from $X$ to $Z$.
    Suppose for all $y\in Y$ there exists $z$ such that for all $x\in\preimg{p}{\{y\}}$ we have $g(x) = z$.
    Suppose $f$ is a function from $Y$ to $Z$.
    Suppose $f\circ p = g$.
    Then $f$ is a quotient map from $Y$ to $Z$ with respect to $T_Y$ and $T_Z$
        iff $g$ is a quotient map from $X$ to $Z$ with respect to $T_X$ and $T_Z$.
\end{lemma}
\begin{proof}
    Show that if $f$ is a quotient map from $Y$ to $Z$ with respect to $T_Y$ and $T_Z$ then
        $g$ is a quotient map from $X$ to $Z$ with respect to $T_X$ and $T_Z$.
    \begin{subproof}
        Follows by \cref{quotient_map_composite}.
    \end{subproof}
    Show that if $g$ is a quotient map from $X$ to $Z$ with respect to $T_X$ and $T_Z$
        then $f$ is a quotient map from $Y$ to $Z$ with respect to $T_Y$ and $T_Z$.
    \begin{subproof}
        Assume $g$ is a quotient map from $X$ to $Z$ with respect to $T_X$ and $T_Z$.
        Hence $f\in\surj{Y}{Z}$ by \cref{quotient_map,surj_factor_of_surj_composite}.
        Therefore $f$ is continuous from $Y$ to $Z$ with respect to $T_Y$ and $T_Z$
            by \cref{quotient_map,quotient_map_induced_map_continuous}.
        Show that for all $V$ such that $V\subseteq Z$ if $\preimg{f}{V}$ is open in $Y$ with respect to $T_Y$
            then $V$ is open in $Z$ with respect to $T_Z$.
        \begin{subproof}
            Fix $V$ such that $V\subseteq Z$.
            Assume $\preimg{f}{V}$ is open in $Y$ with respect to $T_Y$.
            Hence $\preimg{g}{V}$ is open in $X$ with respect to $T_X$ by
                \cref{quotient_map,continuous,preimg_iff,circ_iff,subseteq,subseteq_antisymmetric}.
            $V$ is open in $Z$ with respect to $T_Z$ by \cref{quotient_map}.
        \end{subproof}
        Therefore $f$ is a quotient map from $Y$ to $Z$ with respect to $T_Y$ and $T_Z$ by \cref{quotient_map}.
    \end{subproof}
\end{proof}
% from §22 helper definition 3: fiber partition
\begin{definition}\label{fiber_partition}
    $\fiberPartition{g}{Z} = \{\preimg{g}{\{z\}} \mid z\in Z\}$.
\end{definition}

% from §22 helper lemma 5: fibers form a partition
\begin{lemma}\label{fiber_partition_is_partition}
    Suppose $g$ is a surjection from $X$ to $Z$.
    Let $X_s = \fiberPartition{g}{Z}$.
    Then $X_s$ is a partition of $X$.
\end{lemma}
\begin{proof}
    $g\in\funs{X}{Z}$ and $g$ is a function by \cref{surj,funs_is_function}.
    Show that $X_s$ is a partition.
    \begin{subproof}
        For all $A$ such that $A\in X_s$ we have $A\neq\emptyset$
            by \cref{fiber_partition,surj,funs,function_apply_iff,singleton_intro,preimg_iff,emptyset}.
        Hence $\emptyset\notin X_s$ by \cref{emptyset}.
        For all $A,B$ such that $A\in X_s$ and $B\in X_s$ and $A\neq B$ we have $A$ is disjoint from $B$
            by \cref{fiber_partition,preimg_iff,singleton_elim,function_rightunique,disjoint}.
        Therefore $X_s$ is a partition by \cref{partition}.
    \end{subproof}
    Show that $\unions{X_s} = X$.
    \begin{subproof}
        For all $x$ such that $x\in X$ we have $x\in\unions{X_s}$ by
            \cref{funs_type_apply,fiber_partition,funs,function_apply_elim,preimg_iff,singleton_intro,unions_intro}.
        For all $x$ such that $x\in\unions{X_s}$ we have $x\in X$
            by \cref{unions_iff,fiber_partition,preimg_subseteq_dom,funs,subseteq}.
        Hence $\unions{X_s} = X$ by \cref{subseteq,subseteq_antisymmetric}.
    \end{subproof}
    Therefore $X_s$ is a partition of $X$.
\end{proof}

% from §22 helper lemma 6: partition map is a surjection
\begin{lemma}\label{partition_map_surjection}
    Suppose $P$ is a partition of $X$.
    Let $p = \partitionMap{P}{X}$.
    Then $p$ is a surjection from $X$ to $P$.
\end{lemma}
\begin{proof}
    $P$ is a partition and $\unions{P} = X$.
    Show that $p$ is a function from $X$ to $P$.
    \begin{subproof}
        $p\in\rels{X}{P}$ and $\dom{p}\subseteq X$ by \cref{partition_map,subseteq,rels_intro,rels_dom_subseteq}.
        For all $x$ such that $x\in X$ we have $x\in\dom{p}$
            by \cref{subseteq,unions_iff,times_tuple_intro,fst_eq,snd_eq,partition_map,dom}.
        Hence $\dom{p} = X$ by \cref{subseteq,subseteq_antisymmetric}.
        For all $x,C,D$ such that $x\mathrel{p} C$ and $x\mathrel{p} D$ we have $C = D$
            by \cref{partition_map,times_tuple_elim,fst_eq,snd_eq,partition,disjoint}.
        Thus $p\in\funs{X}{P}$ and $p$ is a function from $X$ to $P$ by \cref{rightunique,funs}.
    \end{subproof}
    Show that for all $C\in P$ there exists $x\in X$ such that $p(x) = C$.
    \begin{subproof}
        Follows by \cref{partition,nonempty_inhabited,unions_intro,subseteq,times_tuple_intro,fst_eq,snd_eq,partition_map,funs_is_function,function_apply_intro}.
    \end{subproof}
    Therefore $p\in\surj{X}{P}$ by \cref{surj}.
\end{proof}

% from §22 Item 20: Corollary 22.3(a) induced bijective map
\begin{lemma}\label{corollary_quotient_space_induced_bijection}
    Suppose $g$ is a surjection from $X$ to $Z$.
    Suppose $g$ is continuous from $X$ to $Z$ with respect to $T_X$ and $T_Z$.
    Let $X_s = \fiberPartition{g}{Z}$.
    Let $p = \partitionMap{X_s}{X}$.
    Let $T_s = \quotientTopology{p}{T_X}{X}{X_s}$.
    Then there exists $f$ such that
        $f$ is a bijection from $X_s$ to $Z$ and
        $f$ is continuous from $X_s$ to $Z$ with respect to $T_s$ and $T_Z$ and
        $f\circ p = g$.
\end{lemma}
\begin{proof}
    $T_X$ is a topology on $X$ and $X_s$ is a partition of $X$
        by \cref{continuous,fiber_partition_is_partition}.
    Hence $p$ is a surjection from $X$ to $X_s$ by \cref{partition_map_surjection}.
    Therefore $p$ is a quotient map from $X$ to $X_s$ with respect to $T_X$ and $T_s$ by \cref{quotient_topology_quotient_map}.
    We have $g\in\funs{X}{Z}$ and $g$ is a function by \cref{surj,funs_is_function}.
    For all $y\in X_s$ there exists $z$ such that for all $x\in\preimg{p}{\{y\}}$ we have $g(x) = z$
        by \cref{fiber_partition,preimg_iff,singleton_elim,partition_map,fst_eq,snd_eq,function_apply_intro}.
    Take $f$ such that $f$ is a function from $X_s$ to $Z$ and $f\circ p = g$ by \cref{quotient_map_induced_map}.
    Hence $f\in\surj{X_s}{Z}$ by \cref{surj,partition_map_surjection,surj_factor_of_surj_composite}.
    Show that for all $y_1,y_2,z$ such that $y_1\mathrel{f} z$ and $y_2\mathrel{f} z$ we have $y_1 = y_2$.
    \begin{subproof}
        Fix $y_1$.
        Fix $y_2$.
        Fix $z$.
        Assume $y_1\mathrel{f} z$ and $y_2\mathrel{f} z$.
        $y_1\in X_s$ and $y_2\in X_s$
            by \cref{funs_is_function,function_apply_intro,function_apply_iff,funs}.
        Take $z_1\in Z$ such that $y_1 = \preimg{g}{\{z_1\}}$ by \cref{fiber_partition}.
        Take $z_2\in Z$ such that $y_2 = \preimg{g}{\{z_2\}}$ by \cref{fiber_partition}.
        $y_1\neq\emptyset$ and $y_2\neq\emptyset$ by \cref{partition}.
        Take $x_1$ such that $x_1\in y_1$ by \cref{nonempty_inhabited}.
        Take $x_2$ such that $x_2\in y_2$ by \cref{nonempty_inhabited}.
        $x_1\in X$ and $x_2\in X$ by \cref{unions_intro,subseteq}.
        $x_1\mathrel{p} y_1$ and $x_2\mathrel{p} y_2$ by \cref{times_tuple_intro,fst_eq,snd_eq,partition_map}.
        Hence $x_1\mathrel{(f\circ p)} z$ and $x_2\mathrel{(f\circ p)} z$ by \cref{circ_iff}.
        $x_1\mathrel{g} z$ and $x_2\mathrel{g} z$.
        $g(x_1) = z$ and $g(x_2) = z$ by \cref{function_apply_intro}.
        $x_1\in\preimg{g}{\{z_1\}}$ and $x_2\in\preimg{g}{\{z_2\}}$.
        $g(x_1) = z_1$ and $g(x_2) = z_2$
            by \cref{preimg_iff,singleton_elim,function_apply_intro}.
        $z_1 = z$ and $z_2 = z$.
        Therefore $y_1 = \preimg{g}{\{z\}}$ and $y_2 = \preimg{g}{\{z\}}$.
        $y_1 = y_2$.
    \end{subproof}
    Hence $f$ is injective by \cref{injective}.
    Thus $f\in\bijections{X_s}{Z}$ by \cref{bijections}.
    Hence $f$ is continuous from $X_s$ to $Z$ with respect to $T_s$ and $T_Z$
        by \cref{quotient_map_induced_map_continuous}.
\end{proof}

% from §22 Item 21: Corollary 22.3(a) homeomorphism criterion
\begin{lemma}\label{corollary_quotient_space_homeomorphism}
    Suppose $g$ is a surjection from $X$ to $Z$.
    Suppose $g$ is continuous from $X$ to $Z$ with respect to $T_X$ and $T_Z$.
    Let $X_s = \fiberPartition{g}{Z}$.
    Let $p = \partitionMap{X_s}{X}$.
    Let $T_s = \quotientTopology{p}{T_X}{X}{X_s}$.
    Suppose $f$ is a function from $X_s$ to $Z$.
    Suppose $f\circ p = g$.
    Suppose $f$ is a bijection from $X_s$ to $Z$.
    Then $f$ is a homeomorphism between $X_s$ and $Z$ with respect to $T_s$ and $T_Z$
        iff $g$ is a quotient map from $X$ to $Z$ with respect to $T_X$ and $T_Z$.
\end{lemma}
\begin{proof}
    $T_X$ is a topology on $X$ and $X_s$ is a partition of $X$
        by \cref{continuous,fiber_partition_is_partition}.
    Hence $p$ is a surjection from $X$ to $X_s$ by \cref{partition_map_surjection}.
    Therefore $p$ is a quotient map from $X$ to $X_s$ with respect to $T_X$ and $T_s$ by \cref{quotient_topology_quotient_map}.
    Show that if $f$ is a homeomorphism between $X_s$ and $Z$ with respect to $T_s$ and $T_Z$
        then $g$ is a quotient map from $X$ to $Z$ with respect to $T_X$ and $T_Z$.
    \begin{subproof}
        Assume $f$ is a homeomorphism between $X_s$ and $Z$ with respect to $T_s$ and $T_Z$.
        Hence $f$ is continuous from $X_s$ to $Z$ with respect to $T_s$ and $T_Z$ by \cref{homeomorphism}.
        For all $U$ such that $U$ is open in $X_s$ with respect to $T_s$
            we have $\img{f}{U}$ is open in $Z$ with respect to $T_Z$
            by \cref{homeomorphism,continuous,funs_is_relation,funs,preim_eq_img_of_converse,converse_converse_eq}.
        Therefore $f$ is an open map from $X_s$ to $Z$ with respect to $T_s$ and $T_Z$ by \cref{open_map}.
        $f$ is a surjection from $X_s$ to $Z$ by \cref{bijections}.
        $T_s$ is a topology on $X_s$ and $T_Z$ is a topology on $Z$
            by \cref{surj,quotient_topology_is_topology,continuous}.
        Therefore $f$ is a quotient map from $X_s$ to $Z$ with respect to $T_s$ and $T_Z$ by \cref{open_or_closed_surjection_quotient}.
        Thus $g$ is a quotient map from $X$ to $Z$ with respect to $T_X$ and $T_Z$
            by \cref{quotient_map_composite}.
    \end{subproof}
    Show that if $g$ is a quotient map from $X$ to $Z$ with respect to $T_X$ and $T_Z$
        then $f$ is a homeomorphism between $X_s$ and $Z$ with respect to $T_s$ and $T_Z$.
    \begin{subproof}
        Assume $g$ is a quotient map from $X$ to $Z$ with respect to $T_X$ and $T_Z$.
        We have $g\in\funs{X}{Z}$ and $g$ is a function by \cref{quotient_map,surj,funs_is_function}.
        For all $y\in X_s$ there exists $z$ such that for all $x\in\preimg{p}{\{y\}}$ we have $g(x) = z$
            by \cref{fiber_partition,preimg_iff,singleton_elim,partition_map,fst_eq,snd_eq,function_apply_intro}.
        Therefore $f$ is a quotient map from $X_s$ to $Z$ with respect to $T_s$ and $T_Z$
            by \cref{quotient_map_induced_map_quotient}.
        We have $f$ is continuous from $X_s$ to $Z$ with respect to $T_s$ and $T_Z$ by \cref{quotient_map}.
        Show that for all $U$ such that $U$ is open in $X_s$ with respect to $T_s$
            we have $\preimg{\converse{f}}{U}$ is open in $Z$ with respect to $T_Z$.
        \begin{subproof}
            Fix $U$ such that $U$ is open in $X_s$ with respect to $T_s$.
            For all $x$ such that $x\in\preimg{f}{\img{f}{U}}$ we have $x\in U$
                by \cref{bijections,preimg_iff,img_iff,injective}.
            For all $x$ such that $x\in U$ we have $x\in\preimg{f}{\img{f}{U}}$ by
                \cref{open_in,topology_on,subseteq,powerset_elim,funs,funs_is_function,function_apply_elim,img_iff,preimg_iff}.
            Therefore $\preimg{f}{\img{f}{U}} = U$ by \cref{subseteq,subseteq_antisymmetric}.
            Hence $\img{f}{U}$ is open in $Z$ with respect to $T_Z$
                by \cref{funs,rels_ran_subseteq,img_subseteq_ran,subseteq_transitive,quotient_map}.
            $\preimg{\converse{f}}{U} = \img{f}{U}$
                by \cref{funs_is_relation,funs,preim_eq_img_of_converse,converse_converse_eq}.
        \end{subproof}
        $\converse{f}\in\funs{Z}{X_s}$ by \cref{bijective_converse_are_funs}.
        $T_Z$ is a topology on $Z$ and $T_s$ is a topology on $X_s$ by \cref{continuous}.
        We have $\converse{f}$ is continuous from $Z$ to $X_s$ with respect to $T_Z$ and $T_s$ by \cref{continuous}.
        Therefore $f$ is a homeomorphism between $X_s$ and $Z$ with respect to $T_s$ and $T_Z$ by \cref{homeomorphism}.
    \end{subproof}
\end{proof}

% from §22 Item 22: Corollary 22.3(b) Hausdorff
\begin{lemma}\label{corollary_quotient_space_hausdorff}
    Suppose $g$ is a surjection from $X$ to $Z$.
    Suppose $g$ is continuous from $X$ to $Z$ with respect to $T_X$ and $T_Z$.
    Suppose $Z$ is Hausdorff with respect to $T_Z$.
    Let $X_s = \fiberPartition{g}{Z}$.
    Let $p = \partitionMap{X_s}{X}$.
    Let $T_s = \quotientTopology{p}{T_X}{X}{X_s}$.
    Suppose $f$ is a bijection from $X_s$ to $Z$.
    Suppose $f$ is continuous from $X_s$ to $Z$ with respect to $T_s$ and $T_Z$.
    Then $X_s$ is Hausdorff with respect to $T_s$.
\end{lemma}
\begin{proof}
    $T_X$ is a topology on $X$ and $X_s$ is a partition of $X$
        by \cref{continuous,fiber_partition_is_partition}.
    Hence $p$ is a surjection from $X$ to $X_s$ by \cref{partition_map_surjection}.
    Therefore $T_s$ is a topology on $X_s$ by \cref{surj,quotient_topology_is_topology}.
    Show that for all $x_1,x_2$ such that $x_1\in X_s$ and $x_2\in X_s$ and $x_1\neq x_2$ there exist $U_1,U_2$ such that
        $U_1$ is a neighborhood of $x_1$ in $X_s$ with respect to $T_s$ and
        $U_2$ is a neighborhood of $x_2$ in $X_s$ with respect to $T_s$ and
        $U_1$ is disjoint from $U_2$.
    \begin{subproof}
        Fix $x_1$.
        Fix $x_2$.
        Assume $x_1\in X_s$ and $x_2\in X_s$ and $x_1\neq x_2$.
        $f$ is injective and $f\in\funs{X_s}{Z}$ by \cref{bijections,bijections_to_funs}.
        We have $f$ is a function and $x_1\in\dom{f}$ and $x_2\in\dom{f}$
            by \cref{bijections_to_funs,funs_is_function,funs,subseteq}.
        $f(x_1) \neq f(x_2)$ and $f(x_1)\in Z$ and $f(x_2)\in Z$
            by \cref{injective_function,funs_type_apply}.
        There exist $V_1,V_2$ such that
            $V_1$ is a neighborhood of $f(x_1)$ in $Z$ with respect to $T_Z$ and
            $V_2$ is a neighborhood of $f(x_2)$ in $Z$ with respect to $T_Z$ and
            $V_1$ is disjoint from $V_2$ by \cref{hausdorff}.
        Therefore $V_1$ and $V_2$ are open in $Z$ with respect to $T_Z$ and
            $f(x_1)\in V_1$ and $f(x_2)\in V_2$ by \cref{neighborhood}.
        Let $U_1 = \preimg{f}{V_1}$.
        Let $U_2 = \preimg{f}{V_2}$.
        We have $U_1$ and $U_2$ are open in $X_s$ with respect to $T_s$ and
            $x_1\in U_1$ and $x_2\in U_2$ by \cref{continuous,preimg_iff,function_apply_elim}.
        $U_1$ is a neighborhood of $x_1$ in $X_s$ with respect to $T_s$ and
            $U_2$ is a neighborhood of $x_2$ in $X_s$ with respect to $T_s$ by \cref{neighborhood}.
        Thus $U_1$ is disjoint from $U_2$ by \cref{preimg_iff,function_apply_iff,disjoint}.
    \end{subproof}
    Therefore $X_s$ is Hausdorff with respect to $T_s$ by \cref{hausdorff}.
\end{proof}

\section{§23 Connected Spaces}

% from §23 Item 1: separation
\begin{definition}\label{separation}
    $X$ is separated by $U$ and $V$ with respect to $T$ iff
        $U$ is open in $X$ with respect to $T$ and
        $V$ is open in $X$ with respect to $T$ and
        $U\neq\emptyset$ and $V\neq\emptyset$ and
        $U$ is disjoint from $V$ and
        $U\union V = X$.
\end{definition}

% from §23 Item 2: connected
\begin{definition}\label{connected}
    $X$ is connected with respect to $T$ iff
        there exists no $U$ such that there exists $V$ such that
            $X$ is separated by $U$ and $V$ with respect to $T$.
\end{definition}
% from §23 Item 3: connected iff only clopen sets are empty or whole space
\begin{lemma}\label{connected_iff_clopen}
    Assume $T$ is a topology on $X$.
    Then $X$ is connected with respect to $T$ iff
        for all $A$ such that $A\subseteq X$ we have
            if $A$ is open in $X$ with respect to $T$ and $A$ is closed in $X$ with respect to $T$ then
                $A = \emptyset$ or $A = X$.
\end{lemma}
\begin{proof}
    Show that if $X$ is connected with respect to $T$ then
        for all $A$ such that $A\subseteq X$ we have
            if $A$ is open in $X$ with respect to $T$ and $A$ is closed in $X$ with respect to $T$ then
                $A = \emptyset$ or $A = X$.
    \begin{subproof}
        Assume $X$ is connected with respect to $T$.
        Fix $A$ such that $A\subseteq X$.
        Assume $A$ is open in $X$ with respect to $T$ and $A$ is closed in $X$ with respect to $T$.
        Suppose not.
        $A\neq\emptyset$ and $A\neq X$.
        Let $U = A$.
        Let $V = X\setminus A$.
        $U$ is open in $X$ with respect to $T$.
        $V$ is open in $X$ with respect to $T$ by \cref{closed_in}.
        $U\neq\emptyset$.
        $V\neq\emptyset$ by \cref{setminus_eq_emptyset_iff_subseteq,subseteq_antisymmetric,subseteq}.
        Hence $U$ is disjoint from $V$ by \cref{setminus_elim_right,disjoint}.
        $U\union V = X$ by \cref{setminus_partition}.
        Therefore $X$ is separated by $U$ and $V$ with respect to $T$ by \cref{separation}.
        Contradiction by \cref{connected}.
        Therefore $A = \emptyset$ or $A = X$.
    \end{subproof}
    Show that if
        for all $A$ such that $A\subseteq X$ we have
            if $A$ is open in $X$ with respect to $T$ and $A$ is closed in $X$ with respect to $T$ then
                $A = \emptyset$ or $A = X$
        then $X$ is connected with respect to $T$.
    \begin{subproof}
        Assume that for all $A$ such that $A\subseteq X$ we have
            if $A$ is open in $X$ with respect to $T$ and $A$ is closed in $X$ with respect to $T$ then
                $A = \emptyset$ or $A = X$.
        Suppose not.
        Take $U$ such that there exists $V$ such that
            $X$ is separated by $U$ and $V$ with respect to $T$.
        Take $V$ such that $X$ is separated by $U$ and $V$ with respect to $T$.
        $U$ and $V$ are open in $X$ with respect to $T$ and $U\neq\emptyset$ and $V\neq\emptyset$
            by \cref{separation}.
        $U$ is disjoint from $V$ and $U\union V = X$ by \cref{separation}.
        Hence $V\subseteq X\setminus U$ by
            \cref{disjoint,union_intro_right,union_upper_left,subseteq,setminus_intro}.
        $X\setminus U\subseteq V$ by
            \cref{setminus_elim_left,union_upper_left,subseteq,union_iff,setminus_elim_right}.
        $U$ is open in $X$ with respect to $T$ and
            $U$ is closed in $X$ with respect to $T$
            by \cref{subseteq_antisymmetric,separation,union_upper_left,subseteq,open_in,closed_in}.
        Take $y$ such that $y\in V$ by \cref{nonempty_inhabited}.
        $y\in X$ and $y\notin U$ by \cref{union_intro_right,subseteq,disjoint}.
        Hence $U\neq X$ by \cref{subseteq}.
        $U = \emptyset$ or $U = X$.
        Contradiction.
        Therefore $X$ is connected with respect to $T$ by \cref{connected}.
    \end{subproof}
\end{proof}

% from §23 Item 4: separation of a subspace
\begin{lemma}\label{separation_subspace_limit_points}
    Assume $Y$ is a subspace of $X$ with respect to $T$.
    Let $T_Y = \subspaceTopology{T}{Y}$.
    Suppose $A\subseteq Y$.
    Suppose $B\subseteq Y$.
    Then $Y$ is separated by $A$ and $B$ with respect to $T_Y$ iff
        $A\neq\emptyset$ and $B\neq\emptyset$ and
        $A$ is disjoint from $B$ and
        $A\union B = Y$ and
        $\limitPoints{A}{X}{T}$ is disjoint from $B$ and
        $\limitPoints{B}{X}{T}$ is disjoint from $A$.
\end{lemma}
\begin{proof}
    Show that if $Y$ is separated by $A$ and $B$ with respect to $T_Y$ then
        $A\neq\emptyset$ and $B\neq\emptyset$ and
        $A$ is disjoint from $B$ and
        $A\union B = Y$ and
        $\limitPoints{A}{X}{T}$ is disjoint from $B$ and
        $\limitPoints{B}{X}{T}$ is disjoint from $A$.
    \begin{subproof}
        Assume $Y$ is separated by $A$ and $B$ with respect to $T_Y$.
        $A$ and $B$ are open in $Y$ with respect to $T_Y$ and $A\neq\emptyset$ and $B\neq\emptyset$
            by \cref{separation}.
        $A$ is disjoint from $B$ and $A\union B = Y$ by \cref{separation}.
        $T$ is a topology on $X$ and $Y\subseteq X$ and $T_Y$ is a topology on $Y$
            by \cref{subspace_of,subspace_topology_is_topology}.
        Therefore $B = Y\setminus A$ by
            \cref{subseteq,disjoint,setminus_intro,setminus_elim_left,union_iff,setminus_elim_right,subseteq_antisymmetric}.
        We have $Y\setminus A$ is open in $Y$ with respect to $T_Y$ and
            $A$ is closed in $Y$ with respect to $T_Y$ by \cref{open_in,closed_in}.
        $A = Y\setminus B$ by \cref{double_complement,subseteq_refl}.
        $Y\setminus B$ is open in $Y$ with respect to $T_Y$ and
            $B$ is closed in $Y$ with respect to $T_Y$ by \cref{open_in,closed_in}.
        Therefore $\closure{A}{X}{T}\inter Y = A$ and $\closure{B}{X}{T}\inter Y = B$
            by \cref{closed_eq_closure,closure_in_subspace,subspace_of}.
        Hence $\limitPoints{A}{X}{T}$ is disjoint from $B$ by
            \cref{subseteq_transitive,subseteq,closure_union_limit_points,union_intro_right,inter_intro,disjoint}.
        Thus $\limitPoints{B}{X}{T}$ is disjoint from $A$ by
            \cref{subseteq_transitive,subseteq,closure_union_limit_points,union_intro_right,inter_intro,disjoint}.
    \end{subproof}
    Show that if
        $A\neq\emptyset$ and $B\neq\emptyset$ and
        $A$ is disjoint from $B$ and
        $A\union B = Y$ and
        $\limitPoints{A}{X}{T}$ is disjoint from $B$ and
        $\limitPoints{B}{X}{T}$ is disjoint from $A$
        then $Y$ is separated by $A$ and $B$ with respect to $T_Y$.
    \begin{subproof}
        Assume $A\neq\emptyset$ and $B\neq\emptyset$ and $A$ is disjoint from $B$ and $A\union B = Y$ and
            $\limitPoints{A}{X}{T}$ is disjoint from $B$ and $\limitPoints{B}{X}{T}$ is disjoint from $A$.
        $T$ is a topology on $X$ and $Y\subseteq X$ and $T_Y$ is a topology on $Y$
            by \cref{subspace_of,subspace_topology_is_topology}.
        Therefore $B = Y\setminus A$ by
            \cref{subseteq,disjoint,setminus_intro,setminus_elim_left,union_iff,setminus_elim_right,subseteq_antisymmetric}.
        $A = Y\setminus B$ by \cref{double_complement,subseteq_refl}.
        $\closure{A}{X}{T}\inter Y \subseteq A$
            by \cref{inter,closure_union_limit_points,subseteq,union_iff,disjoint}.
        Therefore $\closure{A}{X}{T}\inter Y = A$ by
            \cref{subseteq,subseteq_transitive,subseteq_closure,subseteq_inter_iff,subseteq_antisymmetric}.
        We have $\closure{A}{Y}{T_Y} = A$ by \cref{closure_in_subspace,subspace_of}.
        Therefore $Y\setminus A$ is open in $Y$ with respect to $T_Y$ and
            $A$ is closed in $Y$ with respect to $T_Y$ by \cref{closure_closed_in,closed_in,open_in}.
        Show that for all $x$ such that $x\in\closure{B}{X}{T}\inter Y$ we have $x\in B$.
        \begin{subproof}
            Fix $x$ such that $x\in\closure{B}{X}{T}\inter Y$.
            $x\in\closure{B}{X}{T}$ and $x\in Y$ by \cref{inter}.
            $B\subseteq X$ by \cref{subseteq_transitive}.
            Hence $x\in B$ or $x\in\limitPoints{B}{X}{T}$ by \cref{closure_union_limit_points,union_iff}.
            \begin{byCase}
            \caseOf{$x\in B$.}
                $x\in B$.
            \caseOf{$x\in\limitPoints{B}{X}{T}$.}
                Hence $x\notin A$ by \cref{disjoint}.
                Thus $x\in A\union B$ by \cref{union_iff}.
                Therefore $x\in B$ by \cref{union_iff}.
            \end{byCase}
        \end{subproof}
        Hence $\closure{B}{X}{T}\inter Y \subseteq B$ by \cref{subseteq}.
        Therefore $\closure{B}{X}{T}\inter Y = B$ by
            \cref{subseteq,subseteq_transitive,subseteq_closure,subseteq_inter_iff,subseteq_antisymmetric}.
        We have $\closure{B}{Y}{T_Y} = B$ by \cref{closure_in_subspace,subspace_of}.
        Therefore $B$ is closed in $Y$ with respect to $T_Y$ and
            $Y\setminus B$ is open in $Y$ with respect to $T_Y$ by \cref{closure_closed_in,closed_in,open_in}.
        We have $A$ and $B$ are open in $Y$ with respect to $T_Y$ by \cref{open_in}.
        Therefore $Y$ is separated by $A$ and $B$ with respect to $T_Y$ by \cref{separation}.
    \end{subproof}
\end{proof}

% from §23 Item 5: connected subspace lies in a separation set
\begin{lemma}\label{connected_subspace_in_separation}
    Assume $X$ is separated by $C$ and $D$ with respect to $T$.
    Assume $Y$ is a subspace of $X$ with respect to $T$.
    Let $T_Y = \subspaceTopology{T}{Y}$.
    Assume $Y$ is connected with respect to $T_Y$.
    Then $Y\subseteq C$ or $Y\subseteq D$.
\end{lemma}
\begin{proof}
    $C$ and $D$ are open in $X$ with respect to $T$ and $C$ is disjoint from $D$
        and $C\union D = X$ by \cref{separation}.
    $T$ is a topology on $X$ and $Y\subseteq X$ and $T_Y$ is a topology on $Y$
        by \cref{subspace_of,subspace_topology_is_topology}.
    $C\inter Y$ and $D\inter Y$ are open in $Y$ with respect to $T_Y$
        by \cref{open_in,inter_comm,subspace_topology}.
    We have $C\inter Y$ is disjoint from $D\inter Y$ by \cref{inter,disjoint}.
    $(C\inter Y)\union(D\inter Y)\subseteq Y$ by \cref{union_iff,inter,subseteq}.
    $Y\subseteq (C\inter Y)\union(D\inter Y)$
        by \cref{subseteq,union_iff,inter_intro,union_intro_left,union_intro_right}.
    We have $(C\inter Y)\union(D\inter Y) = Y$ by \cref{subseteq,subseteq_antisymmetric}.
    Hence $C\inter Y = \emptyset$ or $D\inter Y = \emptyset$ by \cref{connected,separation}.
    Therefore $Y\subseteq C$ or $Y\subseteq D$ by \cref{union_comm,union_emptyset,inter_lower_left,subseteq}.
\end{proof}

% from §23 helper lemma 1: subspace topology of a subspace
\begin{lemma}\label{subspace_subspace_topology}
    Assume $T$ is a topology on $X$.
    Suppose $A\subseteq Y$.
    Suppose $Y\subseteq X$.
    Let $T_Y = \subspaceTopology{T}{Y}$.
    Then $\subspaceTopology{T_Y}{A} = \subspaceTopology{T}{A}$.
\end{lemma}
\begin{proof}
    Show that $\subspaceTopology{T_Y}{A}\subseteq \subspaceTopology{T}{A}$.
    \begin{subproof}
        It suffices to show that for all $B$ such that $B\in\subspaceTopology{T_Y}{A}$
            we have $B\in\subspaceTopology{T}{A}$ by \cref{subseteq}.
        Fix $B$ such that $B\in\subspaceTopology{T_Y}{A}$.
        Take $V\in T_Y$ such that $B = A\inter V$ by \cref{subspace_topology}.
        Take $U\in T$ such that $V = Y\inter U$ by \cref{subspace_topology}.
        $B = A\inter U$ by \cref{inter_assoc,inter_absorb_supseteq_right,subseteq}.
        Therefore $B\in\subspaceTopology{T}{A}$ by \cref{subspace_topology}.
    \end{subproof}
    Show that $\subspaceTopology{T}{A}\subseteq \subspaceTopology{T_Y}{A}$.
    \begin{subproof}
        It suffices to show that for all $B$ such that $B\in\subspaceTopology{T}{A}$
            we have $B\in\subspaceTopology{T_Y}{A}$ by \cref{subseteq}.
        Fix $B$ such that $B\in\subspaceTopology{T}{A}$.
        Take $U\in T$ such that $B = A\inter U$ by \cref{subspace_topology}.
        $Y\inter U\in T_Y$ by \cref{subspace_topology}.
        $B = A\inter (Y\inter U)$ by \cref{inter_assoc,inter_absorb_supseteq_right,subseteq}.
        Therefore $B\in\subspaceTopology{T_Y}{A}$ by \cref{subspace_topology}.
    \end{subproof}
    Therefore $\subspaceTopology{T_Y}{A} = \subspaceTopology{T}{A}$ by \cref{subseteq_antisymmetric}.
\end{proof}

% from §23 Item 6: union of connected subspaces with a common point
\begin{theorem}\label{union_connected_common_point}
    Assume $M\neq\emptyset$.
    Assume $T$ is a topology on $X$.
    Assume for all $A$ such that $A\in M$ we have $A\subseteq X$.
    Assume for all $A$ such that $A\in M$ we have
        $A$ is connected with respect to $\subspaceTopology{T}{A}$.
    Assume $\inters{M}\neq\emptyset$.
    Then $\unions{M}$ is connected with respect to $\subspaceTopology{T}{\unions{M}}$.
\end{theorem}
\begin{proof}
    Let $Y = \unions{M}$.
    Let $T_Y = \subspaceTopology{T}{Y}$.
    $Y\subseteq X$ by \cref{subseteq,unions_iff}.
    $Y$ is a subspace of $X$ with respect to $T$ and $T_Y$ is a topology on $Y$
        by \cref{subspace_of,subspace_topology_is_topology}.
    Show that there exists no $C$ such that there exists $D$ such that
        $Y$ is separated by $C$ and $D$ with respect to $T_Y$.
    \begin{subproof}
        Suppose not.
        Take $C$ such that there exists $D$ such that
            $Y$ is separated by $C$ and $D$ with respect to $T_Y$.
        Take $D$ such that $Y$ is separated by $C$ and $D$ with respect to $T_Y$.
        $C$ and $D$ are open in $Y$ with respect to $T_Y$ and $C\neq\emptyset$ and $D\neq\emptyset$
            and $C$ is disjoint from $D$ and $C\union D = Y$ by \cref{separation}.
        Take $p$ such that $p\in\inters{M}$ by \cref{nonempty_inhabited}.
        $p\in C$ or $p\in D$ by \cref{inters_iff_forall,unions_iff,union_iff}.
        \begin{byCase}
            \caseOf{$p\in C$.}
                Show that for all $A$ such that $A\in M$ we have $A\subseteq C$.
                \begin{subproof}
                    Fix $A$ such that $A\in M$.
                    $A\subseteq X$ and $A\subseteq Y$ by \cref{subseteq,unions_iff}.
                    $A$ is a subspace of $Y$ with respect to $T_Y$ by \cref{subspace_of}.
                    $\subspaceTopology{T_Y}{A} = \subspaceTopology{T}{A}$ by \cref{subspace_subspace_topology}.
                    $A$ is connected with respect to $\subspaceTopology{T_Y}{A}$.
                    $A\subseteq C$ or $A\subseteq D$ by \cref{connected_subspace_in_separation}.
                    $p\in A$ by \cref{inters_iff_forall}.
                    \begin{byCase}
                    \caseOf{$A\subseteq C$.}
                        Straightforward.
                    \caseOf{$A\subseteq D$.}
                        $p\in D$ by \cref{subseteq}.
                        Contradiction by \cref{disjoint}.
                    \end{byCase}
                \end{subproof}
            $Y\subseteq C$ by \cref{subseteq,unions_iff}.
            $C = Y$ and $D = \emptyset$ by
                \cref{nonempty_inhabited,open_in,topology_on,subseteq,powerset_elim,subseteq_antisymmetric,disjoint}.
            Contradiction.
            \caseOf{$p\in D$.}
                Show that for all $A$ such that $A\in M$ we have $A\subseteq D$.
                \begin{subproof}
                    Fix $A$ such that $A\in M$.
                    $A\subseteq X$ and $A\subseteq Y$ by \cref{subseteq,unions_iff}.
                    $A$ is a subspace of $Y$ with respect to $T_Y$ by \cref{subspace_of}.
                    $\subspaceTopology{T_Y}{A} = \subspaceTopology{T}{A}$ by \cref{subspace_subspace_topology}.
                    $A$ is connected with respect to $\subspaceTopology{T_Y}{A}$.
                    $A\subseteq C$ or $A\subseteq D$ by \cref{connected_subspace_in_separation}.
                    $p\in A$ by \cref{inters_iff_forall}.
                    \begin{byCase}
                    \caseOf{$A\subseteq C$.}
                        $p\in C$ by \cref{subseteq}.
                        Contradiction by \cref{disjoint}.
                    \caseOf{$A\subseteq D$.}
                        Straightforward.
                    \end{byCase}
                \end{subproof}
            $Y\subseteq D$ by \cref{subseteq,unions_iff}.
            $D = Y$ and $C = \emptyset$ by
                \cref{nonempty_inhabited,open_in,topology_on,subseteq,powerset_elim,subseteq_antisymmetric,disjoint}.
            Contradiction.
        \end{byCase}
    \end{subproof}
    $Y$ is connected with respect to $T_Y$ by \cref{connected}.
\end{proof}

% from §23 Item 7: connectedness preserved under adjoining limit points
\begin{theorem}\label{connected_superset_closure}
    Assume $T$ is a topology on $X$.
    Suppose $A\subseteq X$.
    Suppose $B\subseteq X$.
    Assume $A$ is connected with respect to $\subspaceTopology{T}{A}$.
    Suppose $A\subseteq B$.
    Suppose $B\subseteq \closure{A}{X}{T}$.
    Then $B$ is connected with respect to $\subspaceTopology{T}{B}$.
\end{theorem}
\begin{proof}
    Let $T_B = \subspaceTopology{T}{B}$.
    $T$ is a topology on $X$ and $A\subseteq X$ and $B\subseteq X$.
    $T_B$ is a topology on $B$ by \cref{subspace_topology_is_topology}.
    Show that there exists no $C$ such that there exists $D$ such that
        $B$ is separated by $C$ and $D$ with respect to $T_B$.
    \begin{subproof}
        Suppose not.
        Take $C$ such that there exists $D$ such that
            $B$ is separated by $C$ and $D$ with respect to $T_B$.
        Take $D$ such that $B$ is separated by $C$ and $D$ with respect to $T_B$.
        $C$ and $D$ are open in $B$ with respect to $T_B$ and $C\neq\emptyset$ and $D\neq\emptyset$
            and $C$ is disjoint from $D$ and $C\union D = B$ by \cref{separation}.
        $C\subseteq B$ and $D\subseteq B$ and $A$ is a subspace of $B$ with respect to $T_B$
            by \cref{open_in,topology_on,subseteq,powerset_elim,subspace_of}.
        $A\subseteq C$ or $A\subseteq D$ by
            \cref{subspace_subspace_topology,subspace_of,connected_subspace_in_separation}.
        \begin{byCase}
        \caseOf{$A\subseteq C$.}
            $\limitPoints{C}{X}{T}$ is disjoint from $D$ by \cref{subspace_of,separation_subspace_limit_points}.
            $\closure{C}{X}{T}$ is disjoint from $D$ and $C\subseteq X$ by
                \cref{subseteq_transitive,subseteq,closure_union_limit_points,union_iff,disjoint}.
            $D\subseteq \closure{C}{X}{T}$ by
                \cref{subseteq_transitive,subseteq,subseteq_closure,closure_closed_in,closure_subseteq_closed}.
            $D = \emptyset$ by \cref{nonempty_inhabited,subseteq,disjoint}.
            Contradiction by \cref{emptyset}.
        \caseOf{$A\subseteq D$.}
            $\limitPoints{D}{X}{T}$ is disjoint from $C$ by \cref{subspace_of,separation_subspace_limit_points}.
            $\closure{D}{X}{T}$ is disjoint from $C$ and $D\subseteq X$ by
                \cref{subseteq_transitive,subseteq,closure_union_limit_points,union_iff,disjoint}.
            $C\subseteq \closure{D}{X}{T}$ by
                \cref{subseteq_transitive,subseteq,subseteq_closure,closure_closed_in,closure_subseteq_closed}.
            $C = \emptyset$ by \cref{nonempty_inhabited,subseteq,disjoint}.
            Contradiction by \cref{emptyset}.
        \end{byCase}
    \end{subproof}
    $B$ is connected with respect to $T_B$ by \cref{connected}.
\end{proof}
% from §23 Item 8: continuous image of connected space
\begin{theorem}\label{connected_image_continuous}
    Assume $T$ is a topology on $X$.
    Assume $T_1$ is a topology on $Y$.
    Suppose $f$ is continuous from $X$ to $Y$ with respect to $T$ and $T_1$.
    Assume $X$ is connected with respect to $T$.
    Let $Z = \img{f}{X}$.
    Let $T_Z = \subspaceTopology{T_1}{Z}$.
    Then $Z$ is connected with respect to $T_Z$.
\end{theorem}
\begin{proof}
    Suppose not.
    Take $A$ such that there exists $B$ such that
        $Z$ is separated by $A$ and $B$ with respect to $T_Z$.
    Take $B$ such that $Z$ is separated by $A$ and $B$ with respect to $T_Z$.
    $A$ and $B$ are open in $Z$ with respect to $T_Z$ and $A\neq\emptyset$ and $B\neq\emptyset$
        and $A$ is disjoint from $B$ and $A\union B = Z$ by \cref{separation}.
    $Z$ is a subspace of $Y$ with respect to $T_1$ and $T_Z$ is a topology on $Z$ and
        $A\subseteq Z$ and $B\subseteq Z$ by
        \cref{continuous,funs_ran,img_subseteq_ran,subseteq_transitive,subspace_of,subspace_topology_is_topology,open_in,topology_on,subseteq,powerset_elim}.
    $f$ is continuous from $X$ to $Z$ with respect to $T$ and $T_Z$
        by \cref{subseteq_refl,continuous_restrict_range}.
    $f\in\funs{X}{Z}$ by \cref{continuous}.
    $\preimg{f}{A}$ and $\preimg{f}{B}$ are open in $X$ with respect to $T$ by \cref{continuous}.
    $\preimg{f}{A}\neq\emptyset$ and $\preimg{f}{B}\neq\emptyset$
        by \cref{nonempty_inhabited,subseteq,img_iff,preimg_iff,emptyset}.
    $\preimg{f}{A}$ is disjoint from $\preimg{f}{B}$ by
        \cref{preimg_iff,funs_is_function,function_apply_iff,disjoint}.
    Show that for all $x$ such that $x\in X$ we have $x\in\preimg{f}{Z}$.
    \begin{subproof}
        Trivial.
    \end{subproof}
    $\preimg{f}{A}\union\preimg{f}{B} = X$
        by \cref{subseteq,preimg_subseteq_dom,funs,subseteq_antisymmetric,separation,preimg_union}.
    $X$ is separated by $\preimg{f}{A}$ and $\preimg{f}{B}$ with respect to $T$ by \cref{separation}.
    Contradiction by \cref{connected,separation}.
    $Z$ is connected with respect to $T_Z$ by \cref{connected}.
\end{proof}

% from §23 helper lemma 2: identity is continuous
\begin{lemma}\label{identity_continuous}
    Assume $T$ is a topology on $X$.
    Let $j = \identity{X}$.
    Then $j$ is continuous from $X$ to $X$ with respect to $T$ and $T$.
\end{lemma}
\begin{proof}
    Show that for all $U$ such that $U$ is open in $X$ with respect to $T$
        we have $\preimg{j}{U} = U$ and $\preimg{j}{U}$ is open in $X$ with respect to $T$.
    \begin{subproof}
        Follows by \cref{preimg_iff,id_iff,subseteq,open_in,topology_on,pow_iff,subseteq_antisymmetric}.
    \end{subproof}
    $j$ is continuous from $X$ to $X$ with respect to $T$ and $T$ by \cref{id_elem_funs,continuous}.
\end{proof}

% from §23 helper lemma 3: constant map is a function
\begin{lemma}\label{constant_map_function}
    Suppose $c\in B$.
    Let $f = \{(a,c) \mid a\in A\}$.
    Then $f$ is a function from $A$ to $B$.
\end{lemma}
\begin{proof}
    Show that for all $a,b_1,b_2$ such that $(a,b_1)\in f$ and $(a,b_2)\in f$ we have $b_1 = b_2$.
    \begin{subproof}
        Trivial.
    \end{subproof}
    Show that for all $w\in f$ there exist $x,y$ such that $w = (x,y)$.
    \begin{subproof}
        Trivial.
    \end{subproof}
    $f$ is a function by \cref{relation,rightunique}.
    Show that for all $x\in\dom{f}$ we have $f(x)\in B$.
    \begin{subproof}
        Trivial.
    \end{subproof}
    $f$ is a function to $B$.
    Show that for all $x$ such that $x\in\dom{f}$ we have $x\in A$.
    \begin{subproof}
        Trivial.
    \end{subproof}
    Show that for all $a$ such that $a\in A$ we have $a\in\dom{f}$.
    \begin{subproof}
        Trivial.
    \end{subproof}
    $\dom{f} = A$ by \cref{subseteq,subseteq_antisymmetric}.
    $f$ is a function from $A$ to $B$ by \cref{funs_intro,funs}.
\end{proof}

% from §23 Item 9: finite products of connected spaces
\begin{theorem}\label{connected_product_finite}
    Assume $T_1$ is a topology on $X$.
    Assume $T_2$ is a topology on $Y$.
    Assume $X$ is connected with respect to $T_1$.
    Assume $Y$ is connected with respect to $T_2$.
    Let $T = \productTopology{T_1}{T_2}{X}{Y}$.
    Then $X\times Y$ is connected with respect to $T$.
\end{theorem}
\begin{proof}
    $T$ is a topology on $X\times Y$ by
        \cref{product_basis_is_basis,basis_topology_is_topology,product_topology}.
    \begin{byCase}
    \caseOf{$X = \emptyset$.}
        $X\times Y = \emptyset$ by \cref{times_empty_left}.
        Show that there exists no $U$ such that there exists $V$ such that
            $X\times Y$ is separated by $U$ and $V$ with respect to $T$.
        \begin{subproof}
            Follows by \cref{separation,open_in,topology_on,pow_iff,subseteq,subseteq_emptyset_iff}.
        \end{subproof}
        $X\times Y$ is connected with respect to $T$ by \cref{connected}.
    \caseOf{$X\neq\emptyset$.}
        \begin{byCase}
        \caseOf{$Y = \emptyset$.}
            $X\times Y = \emptyset$ by \cref{times_empty_right}.
            Show that there exists no $U$ such that there exists $V$ such that
                $X\times Y$ is separated by $U$ and $V$ with respect to $T$.
            \begin{subproof}
                Follows by \cref{separation,open_in,topology_on,pow_iff,subseteq,subseteq_emptyset_iff}.
            \end{subproof}
            $X\times Y$ is connected with respect to $T$ by \cref{connected}.
        \caseOf{$Y\neq\emptyset$.}
            Take $a$ such that $a\in X$ by \cref{nonempty_inhabited}.
            Take $b$ such that $b\in Y$ by \cref{nonempty_inhabited}.
            Let $j_X = \identity{X}$.
            $j_X$ is continuous from $X$ to $X$ with respect to $T_1$ and $T_1$
                by \cref{identity_continuous}.
            Let $k = \{(x,b) \mid x\in X\}$.
            $k$ is a function from $X$ to $Y$ by \cref{constant_map_function}.
            Show that for all $x$ such that $x\in X$ we have $j_X(x) = x$.
            \begin{subproof}
                Follows by \cref{continuous,funs_is_function,function_apply_intro,id_iff}.
            \end{subproof}
            Show that for all $x$ such that $x\in X$ we have $(x,b)\in k$.
            \begin{subproof}
                Trivial.
            \end{subproof}
            Show that for all $x$ such that $x\in X$ we have $k(x) = b$.
            \begin{subproof}
                Follows by \cref{funs,funs_is_function,function_apply_intro}.
            \end{subproof}
            $k$ is continuous from $X$ to $Y$ with respect to $T_1$ and $T_2$
                by \cref{constant_continuous}.
            Let $f = \{(x,(j_X(x),k(x))) \mid x\in X\}$.
            $f$ is continuous from $X$ to $X\times Y$ with respect to $T_1$ and $T$
                by \cref{continuous_map_into_product_backward}.
            Let $Z = \img{f}{X}$.
            Show that for all $p$ such that $p\in Z$ we have $p\in X\times\{b\}$.
            \begin{subproof}
                Follows by \cref{img_iff,pair_eq_iff,singleton_intro,times_tuple_intro}.
            \end{subproof}
            Show that for all $p$ such that $p\in X\times\{b\}$ we have $p\in Z$.
            \begin{subproof}
                Follows by \cref{times_elem_is_tuple,singleton_iff,pair_eq_iff,function_apply_intro,img_iff}.
            \end{subproof}
            $Z = X\times\{b\}$ by \cref{subseteq,subseteq_antisymmetric}.
            Show that $X\times\{b\}$ is connected with respect to $\subspaceTopology{T}{X\times\{b\}}$.
            \begin{subproof}
                Follows by \cref{connected_image_continuous}.
            \end{subproof}
            Show that for all $x$ such that $x\in X$ we have
                $\{x\}\times Y$ is connected with respect to $\subspaceTopology{T}{\{x\}\times Y}$.
            \begin{subproof}
                Fix $x$ such that $x\in X$.
                Let $j_Y = \identity{Y}$.
                $j_Y$ is continuous from $Y$ to $Y$ with respect to $T_2$ and $T_2$
                    by \cref{identity_continuous}.
                Let $h = \{(y,x) \mid y\in Y\}$.
                $h$ is a function from $Y$ to $X$ by \cref{constant_map_function}.
                Show that for all $y$ such that $y\in Y$ we have $j_Y(y) = y$.
                \begin{subproof}
                    Follows by \cref{continuous,funs_is_function,function_apply_intro,id_iff}.
                \end{subproof}
                Show that for all $y$ such that $y\in Y$ we have $h(y) = x$.
                \begin{subproof}
                    Follows by \cref{funs,funs_is_function,function_apply_intro}.
                \end{subproof}
                $h$ is continuous from $Y$ to $X$ with respect to $T_2$ and $T_1$
                    by \cref{constant_continuous}.
                Let $g = \{(y,(h(y),j_Y(y))) \mid y\in Y\}$.
                $g$ is continuous from $Y$ to $X\times Y$ with respect to $T_2$ and $T$
                    by \cref{continuous_map_into_product_backward}.
                Let $Z_x = \img{g}{Y}$.
                Show that for all $p$ such that $p\in Z_x$ we have $p\in \{x\}\times Y$.
                \begin{subproof}
                    Follows by \cref{img_iff,pair_eq_iff,singleton_intro,times_tuple_intro}.
                \end{subproof}
                Show that for all $p$ such that $p\in \{x\}\times Y$ we have $p\in Z_x$.
                \begin{subproof}
                    Follows by \cref{times_elem_is_tuple,singleton_iff,pair_eq_iff,function_apply_intro,img_iff}.
                \end{subproof}
                $Z_x = \{x\}\times Y$ by \cref{subseteq,subseteq_antisymmetric}.
                $\{x\}\times Y$ is connected with respect to $\subspaceTopology{T}{\{x\}\times Y}$
                    by \cref{connected_image_continuous}.
            \end{subproof}
            Show that for all $x$ such that $x\in X$ we have
                $(X\times\{b\})\union(\{x\}\times Y)$ is connected with respect to
                $\subspaceTopology{T}{(X\times\{b\})\union(\{x\}\times Y)}$.
            \begin{subproof}
                Fix $x$ such that $x\in X$.
                Let $M = \{X\times\{b\},\{x\}\times Y\}$.
                $M\neq\emptyset$ by \cref{upair_intro_left,emptyset}.
                Show that for all $A$ such that $A\in M$ we have $A\subseteq X\times Y$.
                \begin{subproof}
                    Follows by \cref{upair_iff,singleton_iff,subseteq,times_subseteq_right,times_subseteq_left}.
                \end{subproof}
                Show that for all $A$ such that $A\in M$ we have
                    $A$ is connected with respect to $\subspaceTopology{T}{A}$.
                \begin{subproof}
                    Follows by \cref{upair_iff}.
                \end{subproof}
                $\inters{M}\neq\emptyset$
                    by \cref{upair_intro_left,upair_iff,singleton_intro,times_tuple_intro,inters_iff_forall,emptyset}.
                $(X\times\{b\})\union(\{x\}\times Y)$ is connected with respect to
                    $\subspaceTopology{T}{(X\times\{b\})\union(\{x\}\times Y)}$
                    by \cref{union_as_unions,union_connected_common_point}.
            \end{subproof}
            Show that there exists no $C$ such that there exists $D$ such that
                $X\times Y$ is separated by $C$ and $D$ with respect to $T$.
            \begin{subproof}
                Suppose not.
                Take $C$ such that there exists $D$ such that
                    $X\times Y$ is separated by $C$ and $D$ with respect to $T$.
                Take $D$ such that $X\times Y$ is separated by $C$ and $D$ with respect to $T$.
                Show that for all $x$ such that $x\in X$ we have
                    $(X\times\{b\})\union(\{x\}\times Y)\subseteq C$ or
                    $(X\times\{b\})\union(\{x\}\times Y)\subseteq D$.
                \begin{subproof}
                    Fix $x$ such that $x\in X$.
                    Let $W = (X\times\{b\})\union(\{x\}\times Y)$.
                    $W\subseteq X\times Y$ by
                        \cref{union_iff,singleton_iff,times_subseteq_right,times_subseteq_left,subseteq}.
                    $W$ is a subspace of $X\times Y$ with respect to $T$ by \cref{subspace_of}.
                    $W$ is connected with respect to $\subspaceTopology{T}{W}$.
                    $W\subseteq C$ or $W\subseteq D$ by \cref{connected_subspace_in_separation}.
                \end{subproof}
                $(a,b)\in C$ or $(a,b)\in D$ by \cref{times_tuple_intro,separation,union_iff,subseteq}.
                \begin{byCase}
                \caseOf{$(a,b)\in C$.}
                    Show that for all $x$ such that $x\in X$ we have
                        $(X\times\{b\})\union(\{x\}\times Y)\subseteq C$.
                    \begin{subproof}
                        Follows by \cref{singleton_intro,times_tuple_intro,union_intro_left,subseteq,separation,disjoint}.
                    \end{subproof}
                    Show that for all $p$ such that $p\in X\times Y$ we have $p\in C$.
                    \begin{subproof}
                        Follows by \cref{times_elem_is_tuple,singleton_intro,times_tuple_intro,union_intro_right,subseteq}.
                    \end{subproof}
                    $X\times Y\subseteq C$ by \cref{subseteq}.
                    $D = \emptyset$ by \cref{nonempty_inhabited,separation,open_in,topology_on,pow_iff,subseteq,disjoint}.
                    Contradiction by \cref{separation}.
                \caseOf{$(a,b)\in D$.}
                    Show that for all $x$ such that $x\in X$ we have
                        $(X\times\{b\})\union(\{x\}\times Y)\subseteq D$.
                    \begin{subproof}
                        Follows by \cref{singleton_intro,times_tuple_intro,union_intro_left,subseteq,separation,disjoint}.
                    \end{subproof}
                    Show that for all $p$ such that $p\in X\times Y$ we have $p\in D$.
                    \begin{subproof}
                        Follows by \cref{times_elem_is_tuple,singleton_intro,times_tuple_intro,union_intro_right,subseteq}.
                    \end{subproof}
                    $X\times Y\subseteq D$ by \cref{subseteq}.
                    $C = \emptyset$ by \cref{nonempty_inhabited,separation,open_in,topology_on,pow_iff,subseteq,disjoint}.
                    Contradiction by \cref{separation}.
                \end{byCase}
            \end{subproof}
            $X\times Y$ is connected with respect to $T$ by \cref{connected}.
        \end{byCase}
    \end{byCase}
\end{proof}

\section{§24 Connected Subspaces of the Real Line}

% from §24 helper definition 1: upper bound in an ordered set
\begin{definition}\label{upper_bound_rel}
    $b$ is an upper bound of $A$ in $X$ with respect to $R$ iff
        $b\in X$ and $A\subseteq X$ and
        for all $a\in A$ we have $a\mathrel{R}b$ or $a = b$.
\end{definition}

% from §24 helper definition 2: least upper bound in an ordered set
\begin{definition}\label{least_upper_bound_rel}
    $b$ is a least upper bound of $A$ in $X$ with respect to $R$ iff
        $b$ is an upper bound of $A$ in $X$ with respect to $R$ and
        for all $c$ such that $c$ is an upper bound of $A$ in $X$ with respect to $R$
            we have $b\mathrel{R}c$ or $b = c$.
\end{definition}

% from §24 helper definition 3: least upper bound property
\begin{definition}\label{least_upper_bound_property}
    $X$ has the least upper bound property with respect to $R$ iff
        for all $A$ such that $A\subseteq X$ and $A$ is inhabited and
            there exists $b$ such that $b$ is an upper bound of $A$ in $X$ with respect to $R$
        there exists $c$ such that $c$ is a least upper bound of $A$ in $X$ with respect to $R$.
\end{definition}

% from §24 Item 1: linear continuum
\begin{definition}\label{linear_continuum}
    $X$ is a linear continuum with respect to $R$ iff
        $R$ is a simple order relation on $X$ and
        $X$ has more than one element and
        $X$ has the least upper bound property with respect to $R$ and
        for all $x,y\in X$ if $x\mathrel{R}y$ then
            there exists $z\in X$ such that $x\mathrel{R}z$ and $z\mathrel{R}y$.
\end{definition}

% from §24 helper lemma 1: open interval is convex
\begin{lemma}\label{intervalopenrel_convex}
    Assume $R$ is transitive.
    Suppose $a,b\in X$.
    Then $\intervalopenrel{R}{X}{a}{b}$ is convex in $X$ with respect to $R$.
\end{lemma}
\begin{proof}
    Follows by \cref{intervalopen_rel,transitive,subseteq,convex_in}.
\end{proof}

% from §24 helper lemma 2: open-closed interval is convex
\begin{lemma}\label{intervalopenclosedrel_convex}
    Assume $R$ is transitive.
    Suppose $a,b\in X$.
    Then $\intervalopenclosedrel{R}{X}{a}{b}$ is convex in $X$ with respect to $R$.
\end{lemma}
\begin{proof}
    Follows by \cref{intervalopenclosed_rel,transitive,subseteq,convex_in}.
\end{proof}

% from §24 helper lemma 3: closed-open interval is convex
\begin{lemma}\label{intervalclosedopenrel_convex}
    Assume $R$ is transitive.
    Suppose $a,b\in X$.
    Then $\intervalclosedopenrel{R}{X}{a}{b}$ is convex in $X$ with respect to $R$.
\end{lemma}
\begin{proof}
    Follows by \cref{intervalclosedopen_rel,transitive,subseteq,convex_in}.
\end{proof}

% from §24 helper lemma 4: closed interval is convex
\begin{lemma}\label{intervalclosedrel_convex}
    Assume $R$ is transitive.
    Suppose $a,b\in X$.
    Then $\intervalclosedrel{R}{X}{a}{b}$ is convex in $X$ with respect to $R$.
\end{lemma}
\begin{proof}
    Follows by \cref{intervalclosed_rel,transitive,subseteq,convex_in}.
\end{proof}

% from §24 helper lemma 5: open right ray is convex
\begin{lemma}\label{openrayright_convex}
    Assume $R$ is transitive.
    Suppose $a\in X$.
    Then $\openrayright{R}{X}{a}$ is convex in $X$ with respect to $R$.
\end{lemma}
\begin{proof}
    Follows by \cref{openray_right,transitive,subseteq,convex_in}.
\end{proof}

% from §24 helper lemma 6: open left ray is convex
\begin{lemma}\label{openrayleft_convex}
    Assume $R$ is transitive.
    Suppose $a\in X$.
    Then $\openrayleft{R}{X}{a}$ is convex in $X$ with respect to $R$.
\end{lemma}
\begin{proof}
    Follows by \cref{openray_left,transitive,subseteq,convex_in}.
\end{proof}

% from §24 helper lemma 7: closed right ray is convex
\begin{lemma}\label{closedrayright_convex}
    Assume $R$ is transitive.
    Suppose $a\in X$.
    Then $\closedrayright{R}{X}{a}$ is convex in $X$ with respect to $R$.
\end{lemma}
\begin{proof}
    Follows by \cref{closedray_right,transitive,subseteq,convex_in}.
\end{proof}

% from §24 helper lemma 8: closed left ray is convex
\begin{lemma}\label{closedrayleft_convex}
    Assume $R$ is transitive.
    Suppose $a\in X$.
    Then $\closedrayleft{R}{X}{a}$ is convex in $X$ with respect to $R$.
\end{lemma}
\begin{proof}
    Follows by \cref{closedray_left,transitive,subseteq,convex_in}.
\end{proof}
% from §24 helper lemma 9: convex separation contradiction
\begin{lemma}\label{convex_separation_contradiction}
    Assume $T$ is the order topology on $L$ with respect to $R$.
    Assume $L$ is a linear continuum with respect to $R$.
    Suppose $Y$ is convex in $L$ with respect to $R$.
    Let $T_Y = \subspaceTopology{T}{Y}$.
    Suppose $Y$ is separated by $A$ and $B$ with respect to $T_Y$.
    Suppose $a\in A$ and $b\in B$ and $a\mathrel{R}b$.
    Contradiction.
\end{lemma}
\begin{proof}
    $R$ is a simple order relation on $L$ by \cref{linear_continuum}.
    $R$ is transitive and $R$ is asymmetric and $R$ is connex on $L$ by \cref{simple_order_relation}.
    $L$ has more than one element and $L$ has the least upper bound property with respect to $R$
        by \cref{linear_continuum}.
    $Y\subseteq L$ by \cref{convex_in}.

    $T = \basisTopology{\orderBasis{R}{L}}{L}$ and $\orderBasis{R}{L}$ is a basis on $L$
        by \cref{order_topology,order_basis_is_basis}.
    $T$ is a topology on $L$ by \cref{basis_topology_is_topology}.
    $T_Y$ is a topology on $Y$ by \cref{subspace_topology_is_topology,subspace_of,convex_in}.
    $A$ is open in $Y$ with respect to $T_Y$ and $B$ is open in $Y$ with respect to $T_Y$ by \cref{separation}.
    $a\in Y$ and $b\in Y$ and $a\in L$ and $b\in L$
        by \cref{open_in,separation,subspace_topology_is_topology,subspace_of,convex_in,topology_on,subseteq,powerset_elim}.

    Let $I = \intervalclosedrel{R}{L}{a}{b}$.
    $I\subseteq Y$ by \cref{intervalclosed_rel,convex_in,subseteq}.
    $I\subseteq L$ by \cref{convex_in,subseteq_transitive}.

    Let $A_0 = A\inter I$.
    Let $B_0 = B\inter I$.
    Let $T_I = \subspaceTopology{T}{I}$.

    Show that there exists $U\in T$ such that $A = Y\inter U$.
    \begin{subproof}
        Follows by \cref{subspace_topology,open_in}.
    \end{subproof}
    $A_0$ is open in $I$ with respect to $T_I$
        by \cref{subspace_topology,subspace_topology_is_topology,subspace_of,open_in,inter_assoc,inter_comm,inter_absorb_supseteq_right,subseteq}.

    Show that there exists $V\in T$ such that $B = Y\inter V$.
    \begin{subproof}
        Follows by \cref{subspace_topology,open_in}.
    \end{subproof}
    $B_0$ is open in $I$ with respect to $T_I$
        by \cref{subspace_topology,subspace_topology_is_topology,subspace_of,open_in,inter_assoc,inter_comm,inter_absorb_supseteq_right,subseteq}.

    $a\in I$ and $b\in I$ by \cref{intervalclosed_rel}.
    $a\in A_0$ and $b\in B_0$ by \cref{inter_intro}.
    $A_0\neq\emptyset$ and $B_0\neq\emptyset$ by \cref{emptyset}.

    $A_0$ is disjoint from $B_0$ by \cref{inter,disjoint,separation}.

    $A_0\subseteq I$ and $B_0\subseteq I$ by \cref{inter_elim_right,subseteq}.
    $I\subseteq A_0\union B_0$
        by \cref{subseteq,separation,union_iff,inter_intro,union_intro_left,union_intro_right}.
    $A_0\union B_0 = I$ by \cref{union_subsets_is_subset,subseteq,subseteq_antisymmetric}.

    $I$ is separated by $A_0$ and $B_0$ with respect to $T_I$ by \cref{separation}.

    $A_0\subseteq L$ by \cref{inter_elim_right,subseteq,subseteq_transitive}.
    $A_0$ is inhabited by \cref{nonempty_inhabited}.

    $b$ is an upper bound of $A_0$ in $L$ with respect to $R$
        by \cref{inter_elim_right,intervalclosed_rel,upper_bound_rel}.
    Take $c$ such that $c$ is a least upper bound of $A_0$ in $L$ with respect to $R$ by
        \cref{upper_bound_rel,least_upper_bound_property,linear_continuum}.

    $c\in I$ by \cref{least_upper_bound_rel,upper_bound_rel,intervalclosed_rel}.

    $c\in A_0$ or $c\in B_0$ by \cref{union_iff}.
    \begin{byCase}
    \caseOf{$c\in B_0$.}
        Suppose not.
        $c\notin A_0$ by \cref{disjoint}.
        $c\neq a$.
        $a\mathrel{R}c$ by \cref{intervalclosed_rel}.

        Take $U_0\in T$ such that $B_0 = I\inter U_0$ by \cref{subspace_topology,open_in}.
        $c\in U_0$ by \cref{inter_elim_right}.
        Take $B_1\in\orderBasis{R}{L}$ such that $c\in B_1$ and $B_1\subseteq U_0$
            by \cref{topology_generated_by_basis}.
        $B_1\inter I \subseteq B_0$ by \cref{subseteq_inter_iff,subseteq}.

        Show that there exists $d\in I$ such that $d\mathrel{R}c$ and
            $\intervalopenclosedrel{R}{I}{d}{c} \subseteq B_0$.
        \begin{subproof}
            $B_1\in\pow{L}$ and
                $((\exists u, v\in L. u\mathrel{R}v \land B_1 = \intervalopenrel{R}{L}{u}{v}) \lor
                (\exists v\in L. \exists u_0\in L. (\forall x\in L. u_0\mathrel{R}x \lor u_0 = x) \land B_1 = \intervalclosedopenrel{R}{L}{u_0}{v}) \lor
                (\exists u\in L. \exists v_0\in L. (\forall x\in L. x\mathrel{R}v_0 \lor x = v_0) \land B_1 = \intervalopenclosedrel{R}{L}{u}{v_0}))$
                by \cref{order_basis}.
            \begin{byCase}
            \caseOf{$\exists u, v\in L. u\mathrel{R}v \land B_1 = \intervalopenrel{R}{L}{u}{v}$.}
                Take $u, v\in L$ such that $u\mathrel{R}v$ and $B_1 = \intervalopenrel{R}{L}{u}{v}$.
                $u\mathrel{R}c$ and $c\mathrel{R}v$ by \cref{intervalopen_rel}.
                Take $d\in L$ such that
                    $(d = u \lor d = a) \land (u\mathrel{R}d \lor u = d) \land (a\mathrel{R}d \lor a = d)$
                    by \cref{max_of_two_elements}.
                $d\mathrel{R}c$ by \cref{max_of_two_elements,intervalopen_rel,transitive}.
                $d\in I$ by \cref{max_of_two_elements,intervalclosed_rel,transitive}.
                For all $x$ such that $x\in \intervalopenclosedrel{R}{I}{d}{c}$ we have $x\in B_0$
                    by \cref{intervalopenclosed_rel,max_of_two_elements,transitive,intervalopen_rel,inter_intro,subseteq}.
                $\intervalopenclosedrel{R}{I}{d}{c} \subseteq B_0$ by \cref{subseteq}.
                There exists $d_0\in I$ such that $d_0\mathrel{R}c$ and
                    $\intervalopenclosedrel{R}{I}{d_0}{c} \subseteq B_0$.
            \caseOf{$\exists v\in L. \exists u_0\in L. (\forall x\in L. u_0\mathrel{R}x \lor u_0 = x) \land B_1 = \intervalclosedopenrel{R}{L}{u_0}{v}$.}
                Take $u_0, v\in L$ such that
                    $(\forall x\in L. u_0\mathrel{R}x \lor u_0 = x)$ and $B_1 = \intervalclosedopenrel{R}{L}{u_0}{v}$.
                $c\mathrel{R}v$ by \cref{intervalclosedopen_rel}.
                Let $d = a$.
                $d\in I$ by \cref{intervalclosed_rel}.
                $d\mathrel{R}c$ by \cref{intervalclosed_rel}.
                For all $x$ such that $x\in \intervalopenclosedrel{R}{I}{d}{c}$ we have $x\in B_0$
                    by \cref{intervalopenclosed_rel,subseteq,transitive,intervalclosedopen_rel,inter_intro}.
                $\intervalopenclosedrel{R}{I}{d}{c} \subseteq B_0$ by \cref{subseteq}.
                There exists $d_0\in I$ such that $d_0\mathrel{R}c$ and
                    $\intervalopenclosedrel{R}{I}{d_0}{c} \subseteq B_0$.
            \caseOf{$\exists u\in L. \exists v_0\in L. (\forall x\in L. x\mathrel{R}v_0 \lor x = v_0) \land B_1 = \intervalopenclosedrel{R}{L}{u}{v_0}$.}
                Take $u, v_0\in L$ such that
                    $(\forall x\in L. x\mathrel{R}v_0 \lor x = v_0)$ and $B_1 = \intervalopenclosedrel{R}{L}{u}{v_0}$.
                $u\mathrel{R}c$ by \cref{intervalopenclosed_rel}.
                Take $d\in L$ such that
                    $(d = u \lor d = a) \land (u\mathrel{R}d \lor u = d) \land (a\mathrel{R}d \lor a = d)$
                    by \cref{max_of_two_elements}.
                $a\mathrel{R}d$ or $d = a$.
                $d\mathrel{R}c$ by \cref{intervalopenclosed_rel}.
                $c\mathrel{R}b$ or $c = b$ by \cref{intervalclosed_rel}.
                $d\mathrel{R}b$ or $d = b$ by \hyperref[transitive]{transitivity}.
                $d\in I$ by \cref{intervalclosed_rel}.
                For all $x$ such that $x\in \intervalopenclosedrel{R}{I}{d}{c}$ we have $x\in B_0$
                    by \cref{intervalopenclosed_rel,max_of_two_elements,transitive,inter_intro,subseteq}.
                $\intervalopenclosedrel{R}{I}{d}{c} \subseteq B_0$ by \cref{subseteq}.
                There exists $d_0\in I$ such that $d_0\mathrel{R}c$ and
                    $\intervalopenclosedrel{R}{I}{d_0}{c} \subseteq B_0$.
            \end{byCase}
        \end{subproof}

        Take $d\in I$ such that $d\mathrel{R}c$ and $\intervalopenclosedrel{R}{I}{d}{c} \subseteq B_0$.
        Show that for all $x\in A_0$ we have $x\mathrel{R}d$ or $x = d$.
        \begin{subproof}
            Fix $x$ such that $x\in A_0$.
            $x\mathrel{R}c$ or $x = c$ by \cref{least_upper_bound_rel,upper_bound_rel}.
            Show that it is wrong that $d\mathrel{R}x$.
            \begin{subproof}
                Suppose not.
                $d\mathrel{R}x$.
                \begin{byCase}
                \caseOf{$x = c$.}
                    $c\in A_0$.
                    Contradiction by \cref{disjoint}.
                \caseOf{$x\mathrel{R}c$.}
                    $x\in B_0$ by \cref{intervalopenclosed_rel,subseteq}.
                    Contradiction by \cref{disjoint}.
                \end{byCase}
            \end{subproof}
            \begin{byCase}
            \caseOf{$x = d$.}
                $x\mathrel{R}d$ or $x = d$.
            \caseOf{$x \neq d$.}
                $x\mathrel{R}d$ or $d\mathrel{R}x$ by \cref{subseteq,connex}.
                $x\mathrel{R}d$.
                $x\mathrel{R}d$ or $x = d$.
            \end{byCase}
        \end{subproof}
        $d$ is an upper bound of $A_0$ in $L$ with respect to $R$ by \cref{intervalclosed_rel,upper_bound_rel}.
        $c\mathrel{R}d$ or $c = d$ by \cref{least_upper_bound_rel}.
        Contradiction by \cref{asymmetric}.
    \caseOf{$c\in A_0$.}
        Suppose not.
        $c\notin B_0$ by \cref{disjoint}.
        $c\neq b$.

        Take $U_1\in T$ such that $A_0 = I\inter U_1$ by \cref{subspace_topology,open_in}.
        $c\in U_1$ by \cref{inter_elim_right}.
        Take $B_1\in\orderBasis{R}{L}$ such that $c\in B_1$ and $B_1\subseteq U_1$
            by \cref{topology_generated_by_basis}.
        $B_1\inter I \subseteq A_0$ by \cref{subseteq_inter_iff,subseteq}.

        Show that there exists $e\in I$ such that $c\mathrel{R}e$ and
            $\intervalclosedopenrel{R}{I}{c}{e} \subseteq A_0$.
        \begin{subproof}
            $B_1\in\pow{L}$ and
                $((\exists u, v\in L. u\mathrel{R}v \land B_1 = \intervalopenrel{R}{L}{u}{v}) \lor
                (\exists v\in L. \exists u_0\in L. (\forall x\in L. u_0\mathrel{R}x \lor u_0 = x) \land B_1 = \intervalclosedopenrel{R}{L}{u_0}{v}) \lor
                (\exists u\in L. \exists v_0\in L. (\forall x\in L. x\mathrel{R}v_0 \lor x = v_0) \land B_1 = \intervalopenclosedrel{R}{L}{u}{v_0}))$
                by \cref{order_basis}.
            \begin{byCase}
            \caseOf{$\exists u, v\in L. u\mathrel{R}v \land B_1 = \intervalopenrel{R}{L}{u}{v}$.}
                Take $u, v\in L$ such that $u\mathrel{R}v$ and $B_1 = \intervalopenrel{R}{L}{u}{v}$.
                $u\mathrel{R}c$ and $c\mathrel{R}v$ by \cref{intervalopen_rel}.
                Take $e\in L$ such that
                    $(e = v \lor e = b) \land (e\mathrel{R}v \lor e = v) \land (e\mathrel{R}b \lor e = b)$
                    by \cref{min_of_two_elements}.
                $c\mathrel{R}e$ by \cref{min_of_two_elements,intervalclosed_rel,transitive}.
                $e\in I$ by \cref{min_of_two_elements,intervalclosed_rel,transitive}.
                For all $x$ such that $x\in \intervalclosedopenrel{R}{I}{c}{e}$ we have $x\in A_0$
                    by \cref{intervalclosedopen_rel,transitive,intervalopen_rel,inter_intro,subseteq}.
                $\intervalclosedopenrel{R}{I}{c}{e} \subseteq A_0$ by \cref{subseteq}.
                There exists $e_0\in I$ such that $c\mathrel{R}e_0$ and
                    $\intervalclosedopenrel{R}{I}{c}{e_0} \subseteq A_0$.
            \caseOf{$\exists v\in L. \exists u_0\in L. (\forall x\in L. u_0\mathrel{R}x \lor u_0 = x) \land B_1 = \intervalclosedopenrel{R}{L}{u_0}{v}$.}
                Take $u_0, v\in L$ such that
                    $(\forall x\in L. u_0\mathrel{R}x \lor u_0 = x)$ and $B_1 = \intervalclosedopenrel{R}{L}{u_0}{v}$.
                $c\mathrel{R}v$ by \cref{intervalclosedopen_rel}.
                Take $e\in L$ such that
                    $(e = v \lor e = b) \land (e\mathrel{R}v \lor e = v) \land (e\mathrel{R}b \lor e = b)$
                    by \cref{min_of_two_elements}.
                $e\mathrel{R}b$ or $e = b$.
                $c\mathrel{R}b$ by \cref{intervalclosed_rel}.
                $c\mathrel{R}e$.
                $a\mathrel{R}c$ or $a = c$ by \cref{intervalclosed_rel}.
                $a\mathrel{R}e$ or $e = a$ by \hyperref[transitive]{transitivity}.
                $e\in I$ by \cref{intervalclosed_rel}.
                For all $x$ such that $x\in \intervalclosedopenrel{R}{I}{c}{e}$ we have $x\in A_0$
                    by \cref{intervalclosedopen_rel,transitive,inter_intro,subseteq}.
                $\intervalclosedopenrel{R}{I}{c}{e} \subseteq A_0$ by \cref{subseteq}.
                There exists $e_0\in I$ such that $c\mathrel{R}e_0$ and
                    $\intervalclosedopenrel{R}{I}{c}{e_0} \subseteq A_0$.
            \caseOf{$\exists u\in L. \exists v_0\in L. (\forall x\in L. x\mathrel{R}v_0 \lor x = v_0) \land B_1 = \intervalopenclosedrel{R}{L}{u}{v_0}$.}
                Take $u, v_0\in L$ such that
                    $(\forall x\in L. x\mathrel{R}v_0 \lor x = v_0)$ and $B_1 = \intervalopenclosedrel{R}{L}{u}{v_0}$.
                $u\mathrel{R}c$ by \cref{intervalopenclosed_rel}.
                Take $e\in L$ such that
                    $(e = v_0 \lor e = b) \land (e\mathrel{R}v_0 \lor e = v_0) \land (e\mathrel{R}b \lor e = b)$
                    by \cref{min_of_two_elements}.
                Show that $c\mathrel{R}e$.
                \begin{subproof}
                    \begin{byCase}
                    \caseOf{$e = b$.}
                        $c\mathrel{R}b$ or $c = b$ by \cref{intervalclosed_rel}.
                        $c \neq b$ by \cref{disjoint}.
                        $c\mathrel{R}e$.
                    \caseOf{$e = v_0$.}
                        $c\mathrel{R}v_0$ or $c = v_0$.
                        Show that $c\mathrel{R}v_0$.
                        \begin{subproof}
                            Suppose not.
                            $c = v_0$.
                            $b\mathrel{R}v_0$ or $b = v_0$.
                            \begin{byCase}
                            \caseOf{$b = v_0$.} Contradiction by \cref{disjoint}.
                            \caseOf{$b\mathrel{R}v_0$.}
                                $c\mathrel{R}b$ or $c = b$ by \cref{intervalclosed_rel}.
                                \begin{byCase}
                                \caseOf{$c = b$.} Contradiction by \cref{disjoint}.
                                \caseOf{$c\mathrel{R}b$.}
                                    $b\mathrel{R}c$.
                                    Contradiction by \cref{asymmetric}.
                                \end{byCase}
                            \end{byCase}
                        \end{subproof}
                        $c\mathrel{R}e$.
                    \end{byCase}
                \end{subproof}
                $a\mathrel{R}c$ or $a = c$ by \cref{intervalclosed_rel}.
                $a\mathrel{R}e$ or $e = a$ by \hyperref[transitive]{transitivity}.
                $e\in I$ by \cref{intervalclosed_rel}.
                For all $x$ such that $x\in \intervalclosedopenrel{R}{I}{c}{e}$ we have $x\in A_0$
                    by \cref{intervalclosedopen_rel,transitive,intervalopenclosed_rel,inter_intro,subseteq}.
                $\intervalclosedopenrel{R}{I}{c}{e} \subseteq A_0$ by \cref{subseteq}.
                There exists $e_0\in I$ such that $c\mathrel{R}e_0$ and
                    $\intervalclosedopenrel{R}{I}{c}{e_0} \subseteq A_0$.
            \end{byCase}
        \end{subproof}

        Take $e\in I$ such that $c\mathrel{R}e$ and $\intervalclosedopenrel{R}{I}{c}{e} \subseteq A_0$.
        Take $z\in L$ such that $c\mathrel{R}z$ and $z\mathrel{R}e$ by \cref{subseteq,linear_continuum}.
        $a\mathrel{R}c$ or $a = c$ by \cref{intervalclosed_rel}.
        $a\mathrel{R}z$ or $z = a$ by \hyperref[transitive]{transitivity}.
        $e\mathrel{R}b$ or $e = b$ by \cref{intervalclosed_rel}.
        $z\mathrel{R}b$ or $z = b$ by \hyperref[transitive]{transitivity}.
        $z\in I$ by \cref{intervalclosed_rel}.
        $z\in A_0$ by \cref{intervalclosedopen_rel,subseteq}.
        $z\mathrel{R}c$ or $z = c$ by \cref{least_upper_bound_rel,upper_bound_rel}.
        Contradiction by \cref{asymmetric}.
    \end{byCase}
\end{proof}

% from §24 helper lemma 10: convex subset of a linear continuum is connected
\begin{lemma}\label{convex_connected_linear_continuum}
    Assume $T$ is the order topology on $L$ with respect to $R$.
    Assume $L$ is a linear continuum with respect to $R$.
    Suppose $Y$ is convex in $L$ with respect to $R$.
    Let $T_Y = \subspaceTopology{T}{Y}$.
    Then $Y$ is connected with respect to $T_Y$.
\end{lemma}
\begin{proof}
    Suppose not.
    Take $A,B$ such that $Y$ is separated by $A$ and $B$ with respect to $T_Y$.
    Take $a,b$ such that $a\in A$ and $b\in B$ by \cref{nonempty_inhabited,separation}.
    $a\mathrel{R}b$ or $b\mathrel{R}a$
        by \cref{separation,union_upper_left,union_upper_right,convex_in,subseteq_transitive,subseteq,disjoint,linear_continuum,simple_order_relation,connex}.
    Contradiction by \cref{convex_separation_contradiction,separation,disjoint_symmetric,union_comm}.
\end{proof}

% from §24 Item 2: linear continuum connected, intervals and rays connected
\begin{theorem}\label{linear_continuum_connected}
    Assume $T$ is the order topology on $L$ with respect to $R$.
    Assume $L$ is a linear continuum with respect to $R$.
    Then $L$ is connected with respect to $T$ and
        for all $a,b\in L$ if $a\mathrel{R}b$ then
            $\intervalopenrel{R}{L}{a}{b}$ is connected with respect to $\subspaceTopology{T}{\intervalopenrel{R}{L}{a}{b}}$ and
            $\intervalopenclosedrel{R}{L}{a}{b}$ is connected with respect to $\subspaceTopology{T}{\intervalopenclosedrel{R}{L}{a}{b}}$ and
            $\intervalclosedopenrel{R}{L}{a}{b}$ is connected with respect to $\subspaceTopology{T}{\intervalclosedopenrel{R}{L}{a}{b}}$ and
            $\intervalclosedrel{R}{L}{a}{b}$ is connected with respect to $\subspaceTopology{T}{\intervalclosedrel{R}{L}{a}{b}}$ and
        for all $a\in L$ we have
            $\openrayright{R}{L}{a}$ is connected with respect to $\subspaceTopology{T}{\openrayright{R}{L}{a}}$ and
            $\openrayleft{R}{L}{a}$ is connected with respect to $\subspaceTopology{T}{\openrayleft{R}{L}{a}}$ and
            $\closedrayright{R}{L}{a}$ is connected with respect to $\subspaceTopology{T}{\closedrayright{R}{L}{a}}$ and
            $\closedrayleft{R}{L}{a}$ is connected with respect to $\subspaceTopology{T}{\closedrayleft{R}{L}{a}}$.
\end{theorem}
\begin{proof}
    $T$ is a topology on $L$ by \cref{order_topology,order_basis_is_basis,basis_topology_is_topology}.
    Let $T_L = \subspaceTopology{T}{L}$.
    $L$ is connected with respect to $T_L$ by \cref{convex_in,subseteq,convex_connected_linear_continuum}.
    Show that $T_L\subseteq T$.
    \begin{subproof}
        Follows by \cref{subspace_topology,topology_on,powerset_elim,inter_absorb_supseteq_left,subseteq}.
    \end{subproof}
    Show that $T\subseteq T_L$.
    \begin{subproof}
        Show that for all $U$ such that $U\in T$ we have $U\in T_L$.
        \begin{subproof}
            Follows by \cref{topology_on,subseteq,powerset_elim,inter_absorb_supseteq_left,subspace_topology}.
        \end{subproof}
        $T\subseteq T_L$ by \cref{subseteq}.
    \end{subproof}
    $L$ is connected with respect to $T$ by \cref{subseteq_antisymmetric,subseteq}.

    Show that for all $a,b\in L$ if $a\mathrel{R}b$ then
        $\intervalopenrel{R}{L}{a}{b}$ is connected with respect to $\subspaceTopology{T}{\intervalopenrel{R}{L}{a}{b}}$ and
        $\intervalopenclosedrel{R}{L}{a}{b}$ is connected with respect to $\subspaceTopology{T}{\intervalopenclosedrel{R}{L}{a}{b}}$ and
        $\intervalclosedopenrel{R}{L}{a}{b}$ is connected with respect to $\subspaceTopology{T}{\intervalclosedopenrel{R}{L}{a}{b}}$ and
        $\intervalclosedrel{R}{L}{a}{b}$ is connected with respect to $\subspaceTopology{T}{\intervalclosedrel{R}{L}{a}{b}}$.
    \begin{subproof}
        Follows by \cref{linear_continuum,simple_order_relation,intervalopenrel_convex,intervalopenclosedrel_convex,intervalclosedopenrel_convex,intervalclosedrel_convex,convex_connected_linear_continuum}.
    \end{subproof}

    Show that for all $a\in L$ we have
        $\openrayright{R}{L}{a}$ is connected with respect to $\subspaceTopology{T}{\openrayright{R}{L}{a}}$ and
        $\openrayleft{R}{L}{a}$ is connected with respect to $\subspaceTopology{T}{\openrayleft{R}{L}{a}}$ and
        $\closedrayright{R}{L}{a}$ is connected with respect to $\subspaceTopology{T}{\closedrayright{R}{L}{a}}$ and
        $\closedrayleft{R}{L}{a}$ is connected with respect to $\subspaceTopology{T}{\closedrayleft{R}{L}{a}}$.
    \begin{subproof}
        Follows by \cref{linear_continuum,simple_order_relation,openrayright_convex,openrayleft_convex,closedrayright_convex,closedrayleft_convex,convex_connected_linear_continuum}.
    \end{subproof}
\end{proof}
% from §24 helper definition 4: real order relation
\begin{definition}\label{real_order_relation}
    $\realOrderRel = \{w\in\reals\times\reals \mid \fst{w} < \snd{w}\}$.
\end{definition}

% from §24 helper definition 5: order topology on reals
\begin{definition}\label{real_order_topology}
    $\realOrderTopology = \basisTopology{\orderBasis{\realOrderRel}{\reals}}{\reals}$.
\end{definition}

% from §24 helper definition 6: between set in an ordered set
\begin{definition}\label{between_rel}
    $\betweenrel{R}{Y}{x}{y} = \intervalclosedrel{R}{Y}{x}{y} \union \intervalclosedrel{R}{Y}{y}{x}$.
\end{definition}

% from §24 helper lemma 11: real order relation intro
\begin{lemma}\label{real_order_relation_intro}
    Suppose $a\in\reals$ and $b\in\reals$ and $a < b$.
    Then $a\mathrel{\realOrderRel} b$.
\end{lemma}
\begin{proof}
    Follows by \cref{times_tuple_intro,fst_eq,snd_eq,real_lt_subst_left,real_lt_subst_right,real_order_relation}.
\end{proof}

% from §24 helper lemma 12: real order relation elim
\begin{lemma}\label{real_order_relation_elim}
    Suppose $a\mathrel{\realOrderRel} b$.
    Then $a\in\reals$ and $b\in\reals$ and $a < b$.
\end{lemma}
\begin{proof}
    Follows by \cref{real_order_relation,times_tuple_elim,fst_eq,snd_eq,real_lt_subst_left,real_lt_subst_right}.
\end{proof}

% from §24 helper lemma 13: real order relation is simple
\begin{lemma}\label{real_order_relation_simple}
    $\realOrderRel$ is a simple order relation on $\reals$.
\end{lemma}
\begin{proof}
    Show that $\realOrderRel$ is transitive.
    \begin{subproof}
        Follows by \cref{real_order_relation_elim,real_lt_trans,real_order_relation_intro,transitive}.
    \end{subproof}
    Show that for all $a,b$ such that $a\mathrel{\realOrderRel} b$ we have it is wrong that $b\mathrel{\realOrderRel} a$.
    \begin{subproof}
        Follows by \cref{real_order_relation_elim,real_lt_trans,real_lt_irreflexive}.
    \end{subproof}
    $\realOrderRel$ is a simple order relation on $\reals$
        by \cref{real_order_relation,subseteq,asymmetric,real_lt_trichotomy,real_order_relation_intro,connex,simple_order_relation}.
\end{proof}

% from §24 helper lemma 14: inhabited implies nonempty
\begin{lemma}\label{inhabited_nonempty}
    Suppose $A$ is inhabited.
    Then $A\neq\emptyset$.
\end{lemma}
\begin{proof}
    Take $a$ such that $a\in A$.
    $A\neq\emptyset$ by \cref{emptyset}.
\end{proof}

% from §24 helper lemma 15: real order upper bound conversion
\begin{lemma}\label{real_order_upper_bound_is_real_upper_bound}
    Suppose $A\subseteq\reals$.
    Suppose $u$ is an upper bound of $A$ in $\reals$ with respect to $\realOrderRel$.
    Then $u$ is an upper bound of $A$.
\end{lemma}
\begin{proof}
    $u\in\reals$ by \cref{upper_bound_rel}.
    Show that for all $a\in A$ we have $a \leq u$.
    \begin{subproof}
        Fix $a\in A$.
        $a\mathrel{\realOrderRel} u$ or $a = u$ by \cref{upper_bound_rel}.
        $a < u$ or $a = u$ by \cref{real_order_relation_elim}.
        $a \leq u$.
    \end{subproof}
    $u$ is an upper bound of $A$ by \cref{real_upper_bound}.
\end{proof}

% from §24 helper lemma 16: real supremum gives least upper bound
\begin{lemma}\label{real_supremum_is_order_least_upper_bound}
    Suppose $A\subseteq\reals$.
    Suppose $c$ is a supremum of $A$.
    Then $c$ is a least upper bound of $A$ in $\reals$ with respect to $\realOrderRel$.
\end{lemma}
\begin{proof}
    $c$ is an upper bound of $A$ by \cref{real_supremum}.
    $c\in\reals$ by \cref{real_upper_bound}.
    Show that for all $a\in A$ we have $a\mathrel{\realOrderRel} c$ or $a = c$.
    \begin{subproof}
        Fix $a\in A$.
        $a\in\reals$ by \cref{subseteq}.
        $a \leq c$ by \cref{real_upper_bound}.
        $a < c$ or $a = c$ by \cref{real_leq_split}.
        $a\mathrel{\realOrderRel} c$ or $a = c$ by \cref{real_order_relation_intro}.
    \end{subproof}
    $c$ is an upper bound of $A$ in $\reals$ with respect to $\realOrderRel$ by \cref{upper_bound_rel}.
    Show that for all $u$ such that $u$ is an upper bound of $A$ in $\reals$ with respect to $\realOrderRel$
        we have $c\mathrel{\realOrderRel} u$ or $c = u$.
    \begin{subproof}
        Fix $u$ such that $u$ is an upper bound of $A$ in $\reals$ with respect to $\realOrderRel$.
        $u$ is an upper bound of $A$ by \cref{real_order_upper_bound_is_real_upper_bound}.
        $u\in\reals$ by \cref{real_upper_bound}.
        $c \leq u$ by \cref{real_supremum}.
        $c < u$ or $c = u$ by \cref{real_leq_split}.
        $c\mathrel{\realOrderRel} u$ or $c = u$
            by \cref{real_order_relation_intro,real_upper_bound}.
    \end{subproof}
    $c$ is a least upper bound of $A$ in $\reals$ with respect to $\realOrderRel$
        by \cref{least_upper_bound_rel}.
\end{proof}

% from §24 helper lemma 17: real line is a linear continuum
\begin{lemma}\label{reals_linear_continuum}
    $\reals$ is a linear continuum with respect to $\realOrderRel$.
\end{lemma}
\begin{proof}
    Show that $\realOrderRel$ is a simple order relation on $\reals$.
    \begin{subproof}
        Follows by \cref{real_order_relation_simple}.
    \end{subproof}
    Show that $\reals$ has more than one element.
    \begin{subproof}
        Follows by \cref{real_add_ident,real_mul_ident,real_zero_lt_one,real_lt_irreflexive,more_than_one_element}.
    \end{subproof}
    Show that for all $A$ such that $A\subseteq\reals$ and $A$ is inhabited and
        there exists $u$ such that $u$ is an upper bound of $A$ in $\reals$ with respect to $\realOrderRel$
        there exists $c$ such that $c$ is a least upper bound of $A$ in $\reals$ with respect to $\realOrderRel$.
    \begin{subproof}
        Fix $A$ such that $A\subseteq\reals$ and $A$ is inhabited and
            there exists $u$ such that $u$ is an upper bound of $A$ in $\reals$ with respect to $\realOrderRel$.
        Take $u$ such that $u$ is an upper bound of $A$ in $\reals$ with respect to $\realOrderRel$.
        $u$ is an upper bound of $A$ by \cref{real_order_upper_bound_is_real_upper_bound}.
        $A$ is bounded above by \cref{real_bounded_above}.
        $A\neq\emptyset$ by \cref{inhabited_nonempty}.
        Take $c$ such that $c$ is a supremum of $A$ by \cref{real_completeness_axiom}.
        $c$ is a least upper bound of $A$ in $\reals$ with respect to $\realOrderRel$
            by \cref{real_supremum_is_order_least_upper_bound}.
    \end{subproof}
    Show that $\reals$ has the least upper bound property with respect to $\realOrderRel$.
    \begin{subproof}
        Follows by \cref{least_upper_bound_property}.
    \end{subproof}
    Show that for all $a,b\in\reals$ if $a\mathrel{\realOrderRel} b$ then
        there exists $z\in\reals$ such that $a\mathrel{\realOrderRel} z$ and $z\mathrel{\realOrderRel} b$.
    \begin{subproof}
        Follows by \cref{real_order_relation_elim,real_rational_open_interval,intervalopen,real_order_relation_intro}.
    \end{subproof}
    $\reals$ is a linear continuum with respect to $\realOrderRel$ by \cref{linear_continuum}.
\end{proof}

% from §24 Item 3: real line connected, intervals and rays connected
\begin{corollary}\label{reals_connected}
    $\reals$ is connected with respect to $\realOrderTopology$ and
        for all $a,b\in\reals$ if $a\mathrel{\realOrderRel} b$ then
            $\intervalopenrel{\realOrderRel}{\reals}{a}{b}$ is connected with respect to $\subspaceTopology{\realOrderTopology}{\intervalopenrel{\realOrderRel}{\reals}{a}{b}}$ and
            $\intervalopenclosedrel{\realOrderRel}{\reals}{a}{b}$ is connected with respect to $\subspaceTopology{\realOrderTopology}{\intervalopenclosedrel{\realOrderRel}{\reals}{a}{b}}$ and
            $\intervalclosedopenrel{\realOrderRel}{\reals}{a}{b}$ is connected with respect to $\subspaceTopology{\realOrderTopology}{\intervalclosedopenrel{\realOrderRel}{\reals}{a}{b}}$ and
            $\intervalclosedrel{\realOrderRel}{\reals}{a}{b}$ is connected with respect to $\subspaceTopology{\realOrderTopology}{\intervalclosedrel{\realOrderRel}{\reals}{a}{b}}$ and
        for all $a\in\reals$ we have
            $\openrayright{\realOrderRel}{\reals}{a}$ is connected with respect to $\subspaceTopology{\realOrderTopology}{\openrayright{\realOrderRel}{\reals}{a}}$ and
            $\openrayleft{\realOrderRel}{\reals}{a}$ is connected with respect to $\subspaceTopology{\realOrderTopology}{\openrayleft{\realOrderRel}{\reals}{a}}$ and
            $\closedrayright{\realOrderRel}{\reals}{a}$ is connected with respect to $\subspaceTopology{\realOrderTopology}{\closedrayright{\realOrderRel}{\reals}{a}}$ and
            $\closedrayleft{\realOrderRel}{\reals}{a}$ is connected with respect to $\subspaceTopology{\realOrderTopology}{\closedrayleft{\realOrderRel}{\reals}{a}}$.
\end{corollary}
\begin{proof}
    Follows by \cref{reals_linear_continuum,linear_continuum,real_order_topology,order_topology,linear_continuum_connected}.
\end{proof}

% from §24 Item 4: intermediate value theorem
\begin{theorem}\label{intermediate_value_theorem}
    Suppose $f$ is continuous from $X$ to $Y$ with respect to $T$ and $T_1$.
    Assume $X$ is connected with respect to $T$.
    Assume $R$ is a simple order relation on $Y$.
    Assume $Y$ has more than one element.
    Assume $T_1$ is the order topology on $Y$ with respect to $R$.
    Suppose $a\in X$.
    Suppose $b\in X$.
    Suppose $r\in Y$.
    Suppose $r\in\betweenrel{R}{Y}{f(a)}{f(b)}$.
    Then there exists $c\in X$ such that $f(c) = r$.
\end{theorem}
\begin{proof}
    $f\in\funs{X}{Y}$ and $f$ is a function by \cref{continuous,funs_elim}.
    $T$ is a topology on $X$ and $T_1$ is a topology on $Y$ by \cref{continuous}.
    Let $Z = \img{f}{X}$.
    Let $T_Z = \subspaceTopology{T_1}{Z}$.
    $Z$ is connected with respect to $T_Z$ and
        $Z$ is a subspace of $Y$ with respect to $T_1$ and
        $T_Z$ is a topology on $Z$
        by \cref{connected_image_continuous,funs_ran,img_subseteq_ran,subseteq_transitive,subspace_of,subspace_topology_is_topology}.
    Suppose not.
    There exists no $c\in X$ such that $f(c) = r$.
    Show that $r\notin Z$.
    \begin{subproof}
        Suppose not.
        $r\in Z$.
        Take $c\in X$ such that $f(c) = r$ by \cref{img_iff,function_apply_iff}.
        Contradiction.
    \end{subproof}
    $f(a)\in Z$ and $f(b)\in Z$ and $f(a)\in Y$ and $f(b)\in Y$
        by \cref{funs,function_apply_elim,img_elem_intro,funs_type_apply}.
    Let $A = Z\inter\openrayleft{R}{Y}{r}$.
    Let $B = Z\inter\openrayright{R}{Y}{r}$.
    Show that $A$ is open in $Z$ with respect to $T_Z$.
    \begin{subproof}
        Follows by \cref{openrayleft_open_in_order_topology,open_in,subspace_topology}.
    \end{subproof}
    Show that $B$ is open in $Z$ with respect to $T_Z$.
    \begin{subproof}
        Follows by \cref{openrayright_open_in_order_topology,open_in,subspace_topology}.
    \end{subproof}
    Show that $A\neq\emptyset$.
    \begin{subproof}
        Follows by \cref{between_rel,union_iff,intervalclosed_rel,openray_left,inter_intro,emptyset}.
    \end{subproof}
    Show that $B\neq\emptyset$.
    \begin{subproof}
        Follows by \cref{between_rel,union_iff,intervalclosed_rel,openray_right,inter_intro,emptyset}.
    \end{subproof}
    $A$ is disjoint from $B$ by \cref{inter,openray_left,openray_right,simple_order_relation,asymmetric,disjoint}.
    Show that $A\union B\subseteq Z$.
    \begin{subproof}
        Follows by \cref{union_iff,inter,subseteq}.
    \end{subproof}
    Show that $Z\subseteq A\union B$.
    \begin{subproof}
        Show that for all $y$ such that $y\in Z$ we have $y\in A\union B$.
        \begin{subproof}
            Follows by \cref{subspace_of,subseteq,simple_order_relation,connex,openray_left,openray_right,inter_intro,union_iff}.
        \end{subproof}
        $Z\subseteq A\union B$ by \cref{subseteq}.
    \end{subproof}
    $A\union B = Z$ by \cref{subseteq_antisymmetric}.
    $Z$ is separated by $A$ and $B$ with respect to $T_Z$ by \cref{separation}.
    Contradiction by \cref{connected}.
\end{proof}

% from §24 Item 5: path in a space
\begin{definition}\label{path_in_space}
    $f$ is a path in $X$ from $x$ to $y$ with respect to $T$ iff
        there exist $a,b\in\reals$ such that
            $a\mathrel{\realOrderRel} b$ and
            $f$ is continuous from $\intervalclosedrel{\realOrderRel}{\reals}{a}{b}$ to $X$ with respect to
                $\subspaceTopology{\realOrderTopology}{\intervalclosedrel{\realOrderRel}{\reals}{a}{b}}$ and $T$ and
            $f(a) = x$ and $f(b) = y$.
\end{definition}

% from §24 Item 6: path connected
\begin{definition}\label{path_connected}
    $X$ is path connected with respect to $T$ iff
        for all $x,y\in X$ there exists $f$ such that
            $f$ is a path in $X$ from $x$ to $y$ with respect to $T$.
\end{definition}

% from §24 Item 7: path-connected implies connected
\begin{theorem}\label{path_connected_implies_connected}
    Assume $T$ is a topology on $X$.
    Assume $X$ is path connected with respect to $T$.
    Then $X$ is connected with respect to $T$.
\end{theorem}
\begin{proof}
    Suppose not.
    Take $A,B$ such that $X$ is separated by $A$ and $B$ with respect to $T$.
    $A\neq\emptyset$ and $B\neq\emptyset$ by \cref{separation}.
    Take $x,y$ such that $x\in A$ and $y\in B$ by \cref{nonempty_inhabited}.
    Take $f$ such that $f$ is a path in $X$ from $x$ to $y$ with respect to $T$
        by \cref{separation,open_in,topology_on,powerset_elim,subseteq,path_connected}.
    Take $a,b\in\reals$ such that
        $a\mathrel{\realOrderRel} b$ and
        $f$ is continuous from $\intervalclosedrel{\realOrderRel}{\reals}{a}{b}$ to $X$ with respect to
            $\subspaceTopology{\realOrderTopology}{\intervalclosedrel{\realOrderRel}{\reals}{a}{b}}$ and $T$ and
        $f(a) = x$ and $f(b) = y$
        by \cref{path_in_space}.
    Let $I = \intervalclosedrel{\realOrderRel}{\reals}{a}{b}$.
    Let $T_I = \subspaceTopology{\realOrderTopology}{I}$.
    $I$ is connected with respect to $T_I$ by \cref{continuous,reals_connected}.
    Let $Z = \img{f}{I}$.
    Let $T_Z = \subspaceTopology{T}{Z}$.
    Show that $Z\subseteq A$ or $Z\subseteq B$.
    \begin{subproof}
        Follows by \cref{connected_image_continuous,continuous,funs_elim,funs_ran,img_subseteq_ran,subseteq_transitive,subspace_of,connected_subspace_in_separation}.
    \end{subproof}
    Show that $x\in Z$ and $y\in Z$.
    \begin{subproof}
        Follows by \cref{intervalclosed_rel,continuous,funs_elim,function_apply_elim,img_elem_intro}.
    \end{subproof}
    Contradiction by \cref{subseteq,separation,disjoint}.
\end{proof}

\section{§25 Components and Local Connectedness}

% from §25 Item 1: component relation
\begin{definition}\label{component_relation}
    $\componentRelation{T}{X} = \{w\in X\times X \mid
        \text{there exists $A$ such that
            $A$ is a subspace of $X$ with respect to $T$ and
            $\fst{w}\in A$ and $\snd{w}\in A$ and
            $A$ is connected with respect to $\subspaceTopology{T}{A}$}\}$.
\end{definition}

% from §25 Item 2: component
\begin{definition}\label{component}
    Assume $T$ is a topology on $X$.
    $C$ is a component of $X$ with respect to $T$ iff
        there exists $x\in X$ such that
            $C = \equivalenceClass{\componentRelation{T}{X}}{x}$.
\end{definition}

% from §25 helper lemma 1: component relation is a binary relation
\begin{lemma}\label{component_relation_binary}
    Assume $T$ is a topology on $X$.
    Then $\componentRelation{T}{X}$ is a binary relation on $X$.
\end{lemma}
\begin{proof}
    Follows by \cref{component_relation,subseteq}.
\end{proof}

% from §25 helper lemma 2: component relation is reflexive
\begin{lemma}\label{component_relation_reflexive}
    Assume $T$ is a topology on $X$.
    Then $\componentRelation{T}{X}$ is reflexive on $X$.
\end{lemma}
\begin{proof}
    Show that for all $x\in X$ we have $x\mathrel{\componentRelation{T}{X}}x$.
    \begin{subproof}
        Fix $x\in X$.
        Let $A = \{x\}$.
        We have $A$ is a subspace of $X$ with respect to $T$ by \cref{singleton_subset_intro,subspace_of}.
        Show that $A$ is connected with respect to $\subspaceTopology{T}{A}$.
        \begin{subproof}
            Suppose not.
            Take $U,V$ such that $A$ is separated by $U$ and $V$ with respect to $\subspaceTopology{T}{A}$.
            Take $u,v$ such that $u\in U$ and $v\in V$ by \cref{nonempty_inhabited,separation}.
            $u = x$ and $v = x$
                by \cref{separation,open_in,subspace_topology_is_topology,subspace_of,topology_on,powerset_elim,subseteq,singleton_iff}.
            Contradiction by \cref{disjoint,separation}.
        \end{subproof}
        Therefore $x\mathrel{\componentRelation{T}{X}}x$
            by \cref{singleton_intro,fst_eq,snd_eq,times_tuple_intro,component_relation}.
    \end{subproof}
    Hence $\componentRelation{T}{X}$ is reflexive on $X$ by \cref{reflexive_on}.
\end{proof}

% from §25 helper lemma 3: component relation is symmetric
\begin{lemma}\label{component_relation_symmetric}
    Assume $T$ is a topology on $X$.
    Then $\componentRelation{T}{X}$ is symmetric.
\end{lemma}
\begin{proof}
    Follows by \cref{component_relation,subspace_of,subseteq,times_tuple_intro,fst_eq,snd_eq,symmetric}.
\end{proof}

% from §25 helper lemma 4: component relation is transitive
\begin{lemma}\label{component_relation_transitive}
    Assume $T$ is a topology on $X$.
    Then $\componentRelation{T}{X}$ is transitive.
\end{lemma}
\begin{proof}
    Show that for all $x,y,z$ such that
        $x\mathrel{\componentRelation{T}{X}}y$ and $y\mathrel{\componentRelation{T}{X}}z$
        we have $x\mathrel{\componentRelation{T}{X}}z$.
    \begin{subproof}
        Fix $x$.
        Fix $y$.
        Fix $z$.
        Assume $x\mathrel{\componentRelation{T}{X}}y$ and $y\mathrel{\componentRelation{T}{X}}z$.
        Take $A$ such that
            $A$ is a subspace of $X$ with respect to $T$ and
            $\fst{(x,y)}\in A$ and $\snd{(x,y)}\in A$ and
            $A$ is connected with respect to $\subspaceTopology{T}{A}$
            by \cref{component_relation}.
        $x\in A$ and $y\in A$ by \cref{fst_eq,snd_eq}.
        Take $B$ such that
            $B$ is a subspace of $X$ with respect to $T$ and
            $\fst{(y,z)}\in B$ and $\snd{(y,z)}\in B$ and
            $B$ is connected with respect to $\subspaceTopology{T}{B}$
            by \cref{component_relation}.
        $y\in B$ and $z\in B$ by \cref{fst_eq,snd_eq}.
        Let $M = \{A,B\}$.
        $M\neq\emptyset$ by \cref{upair_intro_left,emptyset}.
        Show that for all $C$ such that $C\in M$ we have
            $C\subseteq X$ and
            $C$ is connected with respect to $\subspaceTopology{T}{C}$ and
            $y\in C$.
        \begin{subproof}
            Follows by \cref{upair_iff,subspace_of}.
        \end{subproof}
        Therefore $\inters{M}\neq\emptyset$ by \cref{upair_intro_left,nonempty_inhabited,inters_iff_forall,emptyset}.
        Hence $A\union B$ is connected with respect to $\subspaceTopology{T}{A\union B}$
            by \cref{union_connected_common_point,union_as_unions}.
        $A\union B$ is a subspace of $X$ with respect to $T$ by \cref{union_iff,subspace_of,subseteq}.
        We have $x\in A\union B$ and $z\in A\union B$ by \cref{union_iff}.
        Therefore $x\mathrel{\componentRelation{T}{X}}z$
            by \cref{subspace_of,subseteq,times_tuple_intro,fst_eq,snd_eq,component_relation}.
    \end{subproof}
    Therefore $\componentRelation{T}{X}$ is transitive by \cref{transitive}.
\end{proof}

% from §25 helper lemma 5: component relation is an equivalence
\begin{lemma}\label{component_relation_equivalence}
    Assume $T$ is a topology on $X$.
    Then $\componentRelation{T}{X}$ is an equivalence on $X$.
\end{lemma}
\begin{proof}
    Follows by \cref{component_relation_binary,component_relation_reflexive,component_relation_transitive,component_relation_symmetric}.
\end{proof}

% from §25 Item 3: components are connected
\begin{theorem}\label{component_connected}
    Assume $T$ is a topology on $X$.
    Suppose $C$ is a component of $X$ with respect to $T$.
    Then $C$ is connected with respect to $\subspaceTopology{T}{C}$.
\end{theorem}
    \begin{proof}
        Take $x\in X$ such that $C = \equivalenceClass{\componentRelation{T}{X}}{x}$ by \cref{component}.
        Let $M = \{A\in\pow{X} \mid \text{$x\in A$ and $A$ is connected with respect to $\subspaceTopology{T}{A}$}\}$.
        Take $A$ such that
            $A$ is a subspace of $X$ with respect to $T$ and
            $\fst{(x,x)}\in A$ and $\snd{(x,x)}\in A$ and
            $A$ is connected with respect to $\subspaceTopology{T}{A}$
            by \cref{component_relation_reflexive,reflexive_on,component_relation}.
        Therefore $M\neq\emptyset$ by \cref{fst_eq,snd_eq,subspace_of,powerset_intro,emptyset}.
    Show that for all $A$ such that $A\in M$ we have
        $A\subseteq X$ and
        $A$ is connected with respect to $\subspaceTopology{T}{A}$ and
        $x\in A$.
    \begin{subproof}
        Follows by \cref{powerset_elim}.
    \end{subproof}
    Hence $\inters{M}\neq\emptyset$ by \cref{fst_eq,snd_eq,subspace_of,powerset_intro,nonempty_inhabited,inters_iff_forall,emptyset}.
    Let $U = \unions{M}$.
    $U$ is connected with respect to $\subspaceTopology{T}{U}$ by \cref{union_connected_common_point}.
    Show that $U\subseteq C$.
    \begin{subproof}
        Show that for all $y$ such that $y\in U$ we have $y\in C$.
        \begin{subproof}
            Follows by \cref{unions_iff,powerset_elim,subspace_of,subseteq,times_tuple_intro,fst_eq,snd_eq,component_relation,component_relation_symmetric,symmetric,downward_closure_iff}.
        \end{subproof}
        Hence $U\subseteq C$ by \cref{subseteq}.
    \end{subproof}
    Show that $C\subseteq U$.
    \begin{subproof}
        Show that for all $y$ such that $y\in C$ we have $y\in U$.
        \begin{subproof}
            Follows by \cref{downward_closure_iff,component_relation,fst_eq,snd_eq,subspace_of,powerset_intro,unions_iff}.
        \end{subproof}
        Hence $C\subseteq U$ by \cref{subseteq}.
    \end{subproof}
    Therefore $U = C$ by \cref{subseteq_antisymmetric}.
    Hence $C$ is connected with respect to $\subspaceTopology{T}{C}$.
\end{proof}
% from §25 Item 4: distinct components are disjoint
\begin{theorem}\label{components_disjoint}
    Assume $T$ is a topology on $X$.
    Suppose $C$ is a component of $X$ with respect to $T$.
    Suppose $D$ is a component of $X$ with respect to $T$.
    Suppose $C\neq D$.
    Then $C$ is disjoint from $D$.
\end{theorem}
\begin{proof}
    Take $x\in X$ such that $C = \equivalenceClass{\componentRelation{T}{X}}{x}$ by \cref{component}.
    Take $y\in X$ such that $D = \equivalenceClass{\componentRelation{T}{X}}{y}$ by \cref{component}.
    We have $C$ is disjoint from $D$
        by \cref{component_relation_equivalence,equivalenceon_equivclasses_diseq_implies_disjoint}.
\end{proof}

% from §25 Item 5: components cover the space
\begin{theorem}\label{components_cover}
    Assume $T$ is a topology on $X$.
    Then for all $x\in X$ there exists $C$ such that
        $C$ is a component of $X$ with respect to $T$ and $x\in C$.
\end{theorem}
\begin{proof}
    Follows by \cref{component_relation_equivalence,equivclasses_inhabited_equivalenceon,component}.
\end{proof}

% from §25 Item 6: connected subspaces intersect only one component
\begin{theorem}\label{connected_subspace_single_component}
    Assume $T$ is a topology on $X$.
    Suppose $A$ is a subspace of $X$ with respect to $T$.
    Suppose $A$ is connected with respect to $\subspaceTopology{T}{A}$.
    Suppose $A\neq\emptyset$.
    Then for all $C,D$ such that
        $C$ is a component of $X$ with respect to $T$ and
        $D$ is a component of $X$ with respect to $T$ and
        $A$ intersects $C$ and $A$ intersects $D$
        we have $C = D$.
\end{theorem}
\begin{proof}
    Fix $C$.
    Fix $D$ such that
        $C$ is a component of $X$ with respect to $T$ and
        $D$ is a component of $X$ with respect to $T$ and
        $A$ intersects $C$ and $A$ intersects $D$.
    Take $x,y$ such that $x\in A\inter C$ and $y\in A\inter D$ by \cref{intersects,nonempty_inhabited}.
    Therefore $(x,y)\in \componentRelation{T}{X}$
        by \cref{inter,subspace_of,subseteq,times_tuple_intro,fst_eq,snd_eq,component_relation}.
    Take $c\in X$ such that $C = \equivalenceClass{\componentRelation{T}{X}}{c}$ by \cref{component}.
    Take $d\in X$ such that $D = \equivalenceClass{\componentRelation{T}{X}}{d}$ by \cref{component}.
    Therefore $C = D$
        by \cref{inter,downward_closure_iff,component_relation_symmetric,component_relation_transitive,symmetric,transitive,component_relation_equivalence,equiv_iff_equivclasses_eq}.
\end{proof}

% from §25 Item 7: path relation
\begin{definition}\label{path_relation}
    $\pathRelation{T}{X} = \{w\in X\times X \mid
        \text{there exists $f$ such that
            $f$ is a path in $X$ from $\fst{w}$ to $\snd{w}$ with respect to $T$}\}$.
\end{definition}

% from §25 Item 8: path component
\begin{definition}\label{path_component}
    Assume $T$ is a topology on $X$.
    $C$ is a path component of $X$ with respect to $T$ iff
        there exists $x\in X$ such that
            $C = \equivalenceClass{\pathRelation{T}{X}}{x}$.
\end{definition}

% from §25 helper lemma 6: real line has smaller elements
\begin{lemma}\label{reals_has_smaller}
    For all $a\in\reals$ there exists $b\in\reals$ such that $b < a$.
\end{lemma}
\begin{proof}
    Follows by \cref{real_add_ident,real_mul_ident,real_add_inverse,real_zero_lt_one,real_lt_neg,real_neg_zero,real_lt_subst_right,real_lt_add,real_add_comm,real_lt_subst_left,real_add_closed}.
\end{proof}

% from §25 helper lemma 7: real line has larger elements
\begin{lemma}\label{reals_has_larger}
    For all $a\in\reals$ there exists $b\in\reals$ such that $a < b$.
\end{lemma}
\begin{proof}
    Follows by \cref{real_add_ident,real_mul_ident,real_zero_lt_one,real_lt_add,real_add_comm,real_lt_subst_left,real_lt_subst_right,real_add_closed}.
\end{proof}

% from §25 helper lemma 8: order interval equals real open interval
\begin{lemma}\label{intervalopenrel_real_eq_intervalopen}
    Suppose $a\in\reals$ and $b\in\reals$.
    Then $\intervalopenrel{\realOrderRel}{\reals}{a}{b} = \intervalopen{a}{b}$.
\end{lemma}
\begin{proof}
    Follows by \cref{intervalopen_rel,real_order_relation_elim,intervalopen,real_order_relation_intro,subseteq,subseteq_antisymmetric}.
\end{proof}

% from §25 helper lemma 9: order basis on reals equals the standard basis
\begin{lemma}\label{order_basis_reals_standard}
    $\orderBasis{\realOrderRel}{\reals} = \standardBasisR$.
\end{lemma}
\begin{proof}
    Show that for all $B$ such that $B\in\orderBasis{\realOrderRel}{\reals}$ we have
        $B\in\standardBasisR$.
    \begin{subproof}
            Fix $B$ such that $B\in\orderBasis{\realOrderRel}{\reals}$.
            We have
                $(\exists a, b\in\reals. a\mathrel{\realOrderRel}b \land B = \intervalopenrel{\realOrderRel}{\reals}{a}{b}) \lor
                (\exists b\in\reals. \exists a_0\in\reals. (\forall x\in\reals. a_0\mathrel{\realOrderRel}x \lor a_0 = x) \land B = \intervalclosedopenrel{\realOrderRel}{\reals}{a_0}{b}) \lor
                (\exists a\in\reals. \exists b_0\in\reals. (\forall x\in\reals. x\mathrel{\realOrderRel}b_0 \lor x = b_0) \land B = \intervalopenclosedrel{\realOrderRel}{\reals}{a}{b_0})$
                by \cref{order_basis}.
            \begin{byCase}
            \caseOf{$\exists a, b\in\reals. a\mathrel{\realOrderRel}b \land B = \intervalopenrel{\realOrderRel}{\reals}{a}{b}$.}
                Take $a,b\in\reals$ such that $a\mathrel{\realOrderRel}b$ and
                    $B = \intervalopenrel{\realOrderRel}{\reals}{a}{b}$.
                Hence $B\in\standardBasisR$
                    by \cref{real_order_relation_elim,intervalopenrel_real_eq_intervalopen,intervalopen_subset_reals,subseteq,powerset_intro,standard_basis_reals}.
            \caseOf{$\exists b\in\reals. \exists a_0\in\reals. (\forall x\in\reals. a_0\mathrel{\realOrderRel}x \lor a_0 = x) \land B = \intervalclosedopenrel{\realOrderRel}{\reals}{a_0}{b}$.}
                Suppose not.
                Take $b,a_0\in\reals$ such that
                    $(\forall x\in\reals. a_0\mathrel{\realOrderRel}x \lor a_0 = x)$ and
                    $B = \intervalclosedopenrel{\realOrderRel}{\reals}{a_0}{b}$.
                Take $c\in\reals$ such that $c < a_0$ by \cref{reals_has_smaller}.
                $a_0\mathrel{\realOrderRel}c$ or $a_0 = c$.
                Contradiction by \cref{real_order_relation_elim,real_lt_trans,real_lt_irreflexive,real_lt_subst_left,real_lt_subst_right}.
            \caseOf{$\exists a\in\reals. \exists b_0\in\reals. (\forall x\in\reals. x\mathrel{\realOrderRel}b_0 \lor x = b_0) \land B = \intervalopenclosedrel{\realOrderRel}{\reals}{a}{b_0}$.}
                Suppose not.
                Take $a,b_0\in\reals$ such that
                    $(\forall x\in\reals. x\mathrel{\realOrderRel}b_0 \lor x = b_0)$ and
                    $B = \intervalopenclosedrel{\realOrderRel}{\reals}{a}{b_0}$.
                Take $c\in\reals$ such that $b_0 < c$ by \cref{reals_has_larger}.
                $c\mathrel{\realOrderRel}b_0$ or $c = b_0$.
                Contradiction by \cref{real_order_relation_elim,real_lt_trans,real_lt_irreflexive,real_lt_subst_left,real_lt_subst_right}.
            \end{byCase}
    \end{subproof}
    Show that $\orderBasis{\realOrderRel}{\reals} \subseteq \standardBasisR$.
    \begin{subproof}
        Follows by \cref{subseteq}.
    \end{subproof}
    Show that for all $B$ such that $B\in\standardBasisR$ we have
        $B\in\orderBasis{\realOrderRel}{\reals}$.
    \begin{subproof}
        Follows by \cref{standard_basis_reals,real_order_relation_intro,intervalopenrel_real_eq_intervalopen,intervalopen_subset_reals,subseteq,powerset_intro,order_basis}.
    \end{subproof}
    Show that $\standardBasisR \subseteq \orderBasis{\realOrderRel}{\reals}$.
    \begin{subproof}
        Follows by \cref{subseteq}.
    \end{subproof}
    Hence $\orderBasis{\realOrderRel}{\reals} = \standardBasisR$ by \cref{subseteq_antisymmetric}.
\end{proof}

% from §25 helper lemma 10: real order topology equals standard topology
\begin{lemma}\label{real_order_topology_eq_standard}
    $\realOrderTopology = \standardTopologyR$.
\end{lemma}
\begin{proof}
    Follows by \cref{real_order_topology,standard_topology_reals,order_basis_reals_standard}.
\end{proof}

% from §25 helper lemma 11: closed right ray is complement of open left ray
\begin{lemma}\label{closedrayright_complement_openrayleft}
    Assume $R$ is a simple order relation on $X$.
    Suppose $a\in X$.
    Then $\closedrayright{R}{X}{a} = X\setminus \openrayleft{R}{X}{a}$.
\end{lemma}
\begin{proof}
    Show that for all $x$ such that $x\in\closedrayright{R}{X}{a}$ we have
        $x\in X\setminus \openrayleft{R}{X}{a}$.
    \begin{subproof}
        Fix $x$ such that $x\in\closedrayright{R}{X}{a}$.
        $x\in X$ by \cref{closedray_right}.
        $a\mathrel{R}x$ or $x = a$ by \cref{closedray_right}.
        Show that $x\notin\openrayleft{R}{X}{a}$.
        \begin{subproof}
            Suppose not.
            $x\mathrel{R}a$ by \cref{openray_left}.
            Contradiction by \cref{simple_order_relation,asymmetric}.
        \end{subproof}
        Therefore $x\in X\setminus \openrayleft{R}{X}{a}$ by \cref{setminus_intro}.
    \end{subproof}
    Show that $\closedrayright{R}{X}{a}\subseteq X\setminus \openrayleft{R}{X}{a}$.
    \begin{subproof}
        Follows by \cref{subseteq}.
    \end{subproof}
    Show that for all $x$ such that $x\in X\setminus \openrayleft{R}{X}{a}$ we have
        $x\in\closedrayright{R}{X}{a}$.
    \begin{subproof}
        Fix $x$ such that $x\in X\setminus \openrayleft{R}{X}{a}$.
        $x\in X$ and $x\notin\openrayleft{R}{X}{a}$ by \cref{setminus_elim_left,setminus_elim_right}.
        $a\mathrel{R}x$ or $x = a$ by \cref{simple_order_relation,connex,openray_left}.
        Therefore $x\in\closedrayright{R}{X}{a}$ by \cref{closedray_right}.
    \end{subproof}
    Show that $X\setminus \openrayleft{R}{X}{a}\subseteq \closedrayright{R}{X}{a}$.
    \begin{subproof}
        Follows by \cref{subseteq}.
    \end{subproof}
    Therefore $\closedrayright{R}{X}{a} = X\setminus \openrayleft{R}{X}{a}$
        by \cref{subseteq_antisymmetric}.
\end{proof}

% from §25 helper lemma 12: closed left ray is complement of open right ray
\begin{lemma}\label{closedrayleft_complement_openrayright}
    Assume $R$ is a simple order relation on $X$.
    Suppose $a\in X$.
    Then $\closedrayleft{R}{X}{a} = X\setminus \openrayright{R}{X}{a}$.
\end{lemma}
\begin{proof}
    Show that for all $x$ such that $x\in\closedrayleft{R}{X}{a}$ we have
        $x\in X\setminus \openrayright{R}{X}{a}$.
    \begin{subproof}
        Fix $x$ such that $x\in\closedrayleft{R}{X}{a}$.
        $x\in X$ by \cref{closedray_left}.
        $x\mathrel{R}a$ or $x = a$ by \cref{closedray_left}.
        Show that $x\notin\openrayright{R}{X}{a}$.
        \begin{subproof}
            Suppose not.
            $a\mathrel{R}x$ by \cref{openray_right}.
            Contradiction by \cref{simple_order_relation,asymmetric}.
        \end{subproof}
        Therefore $x\in X\setminus \openrayright{R}{X}{a}$ by \cref{setminus_intro}.
    \end{subproof}
    Show that $\closedrayleft{R}{X}{a}\subseteq X\setminus \openrayright{R}{X}{a}$.
    \begin{subproof}
        Follows by \cref{subseteq}.
    \end{subproof}
    Show that for all $x$ such that $x\in X\setminus \openrayright{R}{X}{a}$ we have
        $x\in\closedrayleft{R}{X}{a}$.
    \begin{subproof}
        Fix $x$ such that $x\in X\setminus \openrayright{R}{X}{a}$.
        $x\in X$ and $x\notin\openrayright{R}{X}{a}$ by \cref{setminus_elim_left,setminus_elim_right}.
        $x\mathrel{R}a$ or $x = a$ by \cref{simple_order_relation,connex,openray_right}.
        Therefore $x\in\closedrayleft{R}{X}{a}$ by \cref{closedray_left}.
    \end{subproof}
    Show that $X\setminus \openrayright{R}{X}{a}\subseteq \closedrayleft{R}{X}{a}$.
    \begin{subproof}
        Follows by \cref{subseteq}.
    \end{subproof}
    Therefore $\closedrayleft{R}{X}{a} = X\setminus \openrayright{R}{X}{a}$
        by \cref{subseteq_antisymmetric}.
\end{proof}
% from §25 helper lemma 13: closed right ray is closed in the order topology
\begin{lemma}\label{closedrayright_closed_in_order_topology}
    Assume $R$ is a simple order relation on $X$.
    Assume $X$ has more than one element.
    Assume $T$ is the order topology on $X$ with respect to $R$.
    Suppose $a\in X$.
    Then $\closedrayright{R}{X}{a}$ is closed in $X$ with respect to $T$.
\end{lemma}
\begin{proof}
    Follows by \cref{openrayleft_open_in_order_topology,closedray_right,openray_left,subseteq,closedrayright_complement_openrayleft,setminus_setminus,inter_absorb_supseteq_left,open_in,closed_in}.
\end{proof}

% from §25 helper lemma 14: closed left ray is closed in the order topology
\begin{lemma}\label{closedrayleft_closed_in_order_topology}
    Assume $R$ is a simple order relation on $X$.
    Assume $X$ has more than one element.
    Assume $T$ is the order topology on $X$ with respect to $R$.
    Suppose $a\in X$.
    Then $\closedrayleft{R}{X}{a}$ is closed in $X$ with respect to $T$.
\end{lemma}
\begin{proof}
    Follows by \cref{openrayright_open_in_order_topology,closedray_left,openray_right,subseteq,closedrayleft_complement_openrayright,setminus_setminus,inter_absorb_supseteq_left,open_in,closed_in}.
\end{proof}

% from §25 helper lemma 15: closed interval is closed in the order topology
\begin{lemma}\label{intervalclosedrel_closed_in_order_topology}
    Assume $R$ is a simple order relation on $X$.
    Assume $X$ has more than one element.
    Assume $T$ is the order topology on $X$ with respect to $R$.
    Suppose $a\in X$.
    Suppose $b\in X$.
    Then $\intervalclosedrel{R}{X}{a}{b}$ is closed in $X$ with respect to $T$.
\end{lemma}
\begin{proof}
    Let $A = \closedrayright{R}{X}{a}$.
    Let $B = \closedrayleft{R}{X}{b}$.
    $T$ is a topology on $X$ by \cref{order_basis_is_basis,basis_topology_is_topology,order_topology}.
    Show that $A$ is closed in $X$ with respect to $T$.
    \begin{subproof}
        Follows by \cref{closedrayright_closed_in_order_topology}.
    \end{subproof}
    Show that $B$ is closed in $X$ with respect to $T$.
    \begin{subproof}
        Follows by \cref{closedrayleft_closed_in_order_topology}.
    \end{subproof}
    Show that $A\inter B$ is closed in $X$ with respect to $T$.
    \begin{subproof}
        Follows by \cref{upair_iff,closedray_right,closedray_left,subseteq,powerset_intro,closed_inters,inter_as_inters}.
    \end{subproof}
    Show that for all $x$ such that $x\in\intervalclosedrel{R}{X}{a}{b}$ we have $x\in A\inter B$.
    \begin{subproof}
        Fix $x$ such that $x\in\intervalclosedrel{R}{X}{a}{b}$.
        $x\in X$ by \cref{intervalclosed_rel}.
        $a\mathrel{R}x$ or $x = a$ by \cref{intervalclosed_rel}.
        $x\mathrel{R}b$ or $x = b$ by \cref{intervalclosed_rel}.
        Hence $x\in A$ and $x\in B$ by \cref{closedray_right,closedray_left}.
        Therefore $x\in A\inter B$ by \cref{inter}.
    \end{subproof}
    Show that $\intervalclosedrel{R}{X}{a}{b}\subseteq A\inter B$.
    \begin{subproof}
        Follows by \cref{subseteq}.
    \end{subproof}
    Show that for all $x$ such that $x\in A\inter B$ we have $x\in\intervalclosedrel{R}{X}{a}{b}$.
    \begin{subproof}
        Fix $x$ such that $x\in A\inter B$.
        $x\in A$ and $x\in B$ by \cref{inter}.
        $x\in X$ by \cref{closedray_right}.
        $a\mathrel{R}x$ or $x = a$ by \cref{closedray_right}.
        $x\mathrel{R}b$ or $x = b$ by \cref{closedray_left}.
        Therefore $x\in\intervalclosedrel{R}{X}{a}{b}$ by \cref{intervalclosed_rel}.
    \end{subproof}
    Show that $A\inter B\subseteq \intervalclosedrel{R}{X}{a}{b}$.
    \begin{subproof}
        Follows by \cref{subseteq}.
    \end{subproof}
    Therefore $\intervalclosedrel{R}{X}{a}{b} = A\inter B$ by \cref{subseteq_antisymmetric}.
    Show that $\intervalclosedrel{R}{X}{a}{b}$ is closed in $X$ with respect to $T$.
    \begin{subproof}
        Follows by \cref{subseteq}.
    \end{subproof}
\end{proof}

% from §25 helper lemma 16: reverse a path
\begin{lemma}\label{path_reverse}
    Assume $T$ is a topology on $X$.
    Suppose $f$ is a path in $X$ from $x$ to $y$ with respect to $T$.
    Then there exists $g$ such that $g$ is a path in $X$ from $y$ to $x$ with respect to $T$.
\end{lemma}
\begin{proof}
    Take $a,b\in\reals$ such that
        $a\mathrel{\realOrderRel} b$ and
        $f$ is continuous from $\intervalclosedrel{\realOrderRel}{\reals}{a}{b}$ to $X$ with respect to
            $\subspaceTopology{\realOrderTopology}{\intervalclosedrel{\realOrderRel}{\reals}{a}{b}}$ and $T$ and
        $f(a) = x$ and $f(b) = y$
        by \cref{path_in_space}.
    Let $I = \intervalclosedrel{\realOrderRel}{\reals}{a}{b}$.
    Let $T_I = \subspaceTopology{\realOrderTopology}{I}$.
    $\realOrderTopology$ is a topology on $\reals$ by
        \cref{standard_basis_is_basis,standard_topology_reals,basis_topology_is_topology,real_order_topology_eq_standard}.
    $I$ is a subspace of $\reals$ with respect to $\realOrderTopology$ by \cref{intervalclosed_rel,subseteq,subspace_of}.
    $T_I$ is a topology on $I$ by \cref{subspace_topology_is_topology,subspace_of}.
    Let $c_0 = a + b$.
    $c_0\in\reals$ by \cref{real_add_closed}.
    Let $k = \{(t,c_0) \mid t\in I\}$.
    $k$ is a function from $I$ to $\reals$ by \cref{constant_map_function}.
    Show that for all $t\in I$ we have $k(t) = c_0$.
    \begin{subproof}
        Fix $t\in I$.
        $(t,c_0)\in k$.
        Hence $k(t) = c_0$ by \cref{funs_is_function,function_apply_iff}.
    \end{subproof}
    Hence $k$ is continuous from $I$ to $\reals$ with respect to $T_I$ and $\standardTopologyR$
        by \cref{constant_continuous,real_order_topology_eq_standard}.
    Let $j = \identity{I}$.
    $j$ is continuous from $I$ to $\reals$ with respect to $T_I$ and $\standardTopologyR$
        by \cref{inclusion_continuous,real_order_topology_eq_standard}.
    Let $r = \{(t, k(t) - j(t)) \mid t\in I\}$.
    $r$ is continuous from $I$ to $\reals$ with respect to $T_I$ and $\standardTopologyR$
        by \cref{real_function_subtraction_continuous}.
    $r$ is continuous from $I$ to $\reals$ with respect to $T_I$ and $\realOrderTopology$
        by \cref{real_order_topology_eq_standard}.
    Show that for all $t$ such that $t\in I$ we have $r(t)\in I$.
    \begin{subproof}
            Fix $t$ such that $t\in I$.
            $t\in\reals$ by \cref{intervalclosed_rel}.
            $a\mathrel{\realOrderRel}t$ or $t = a$ by \cref{intervalclosed_rel}.
            $t\mathrel{\realOrderRel}b$ or $t = b$ by \cref{intervalclosed_rel}.
            $j(t) = t$ by \cref{id_elem_funs,id_iff,funs_is_function,function_apply_iff}.
            $(t, k(t) - j(t))\in r$.
            $r(t) = k(t) - j(t)$ by \cref{continuous,funs_is_function,function_apply_iff}.
            $k(t) = c_0$.
            $r(t) = c_0 - t$.
            $r(t)\in\reals$ by \cref{real_add_inverse,real_add_closed,real_sub}.
            Show that $a\mathrel{\realOrderRel}r(t)$ or $r(t) = a$.
            \begin{subproof}
                \begin{byCase}
                \caseOf{$t = b$.}
                    $b + \neg{b} = 0$ by \cref{real_add_inverse}.
                    $c_0 - b = (a + b) + \neg{b}$ by \cref{real_sub,real_add_eq_subst}.
                    $a\in\reals$ and $b\in\reals$.
                    $\neg{b}\in\reals$ by \cref{real_add_inverse}.
                    $(a + b) + \neg{b} = a + (b + \neg{b})$ by \cref{real_add_assoc}.
                    Hence $r(t) = a + 0$ by \cref{real_add_eq_subst}.
                    Therefore $r(t) = a$ by \cref{real_add_ident,real_add_comm}.
                \caseOf{$t\mathrel{\realOrderRel}b$.}
                    $t < b$ by \cref{real_order_relation_elim}.
                    $b\in\reals$.
                    $\neg{t}\in\reals$ by \cref{real_add_inverse}.
                    Hence $t + \neg{t} < b + \neg{t}$ by \cref{real_lt_add}.
                    $t + \neg{t} = 0$ by \cref{real_add_inverse}.
                    Therefore $0 < b + \neg{t}$ by \cref{real_lt_subst_left}.
                    Hence $0 < b - t$ by \cref{real_sub,real_lt_subst_right}.
                    $0\in\reals$ by \cref{real_add_ident}.
                    $a\in\reals$.
                    $b - t\in\reals$ by \cref{real_sub,real_add_closed}.
                    Hence $0 + a < (b - t) + a$ by \cref{real_lt_add}.
                    $0 + a = a$ by \cref{real_add_ident}.
                    $(b - t) + a = a + (b - t)$ by \cref{real_add_comm}.
                    Therefore $a < a + (b - t)$ by \cref{real_lt_subst_left,real_lt_subst_right}.
                    $a + (b - t) = a + (b + \neg{t})$ by \cref{real_sub,real_add_eq_subst}.
                    $a + (b + \neg{t}) = (a + b) + \neg{t}$ by \cref{real_add_assoc}.
                    $(a + b) + \neg{t} = (a + b) - t$ by \cref{real_sub}.
                    $a + (b - t) = c_0 - t$.
                    Hence $a < c_0 - t$ by \cref{real_lt_subst_right}.
                    Therefore $a\mathrel{\realOrderRel}r(t)$ by \cref{real_order_relation_intro}.
                \end{byCase}
            \end{subproof}
            Show that $r(t)\mathrel{\realOrderRel}b$ or $r(t) = b$.
            \begin{subproof}
                \begin{byCase}
                \caseOf{$t = a$.}
                    $c_0 = b + a$ by \cref{real_add_comm,real_add_eq_subst}.
                    Hence $c_0 - a = (b + a) + \neg{a}$ by \cref{real_sub,real_add_eq_subst}.
                    $a\in\reals$ and $b\in\reals$.
                    $\neg{a}\in\reals$ by \cref{real_add_inverse}.
                    $(b + a) + \neg{a} = b + (a + \neg{a})$ by \cref{real_add_assoc}.
                    $a + \neg{a} = 0$ by \cref{real_add_inverse}.
                    Therefore $r(t) = b + 0$ by \cref{real_add_eq_subst}.
                    Hence $r(t) = b$ by \cref{real_add_ident,real_add_comm}.
                \caseOf{$a\mathrel{\realOrderRel}t$.}
                    $a < t$ by \cref{real_order_relation_elim}.
                    $a\in\reals$ and $t\in\reals$ and $b\in\reals$.
                    $\neg{t}\in\reals$ by \cref{real_add_inverse}.
                    Hence $b - t\in\reals$ by \cref{real_sub,real_add_closed}.
                    $a + (b - t) < t + (b - t)$ by \cref{real_lt_add}.
                    $t + (b - t) = t + (b + \neg{t})$ by \cref{real_sub,real_add_eq_subst}.
                    $t + (b + \neg{t}) = (t + b) + \neg{t}$ by \cref{real_add_assoc}.
                    $(t + b) + \neg{t} = (b + t) + \neg{t}$ by \cref{real_add_comm}.
                    $(b + t) + \neg{t} = b + (t + \neg{t})$ by \cref{real_add_assoc}.
                    $t + \neg{t} = 0$ by \cref{real_add_inverse}.
                    Therefore $t + (b - t) = b + 0$ by \cref{real_add_eq_subst}.
                    Hence $t + (b - t) = b$ by \cref{real_add_ident,real_add_comm}.
                    Hence $a + (b - t) < b$ by \cref{real_lt_subst_right}.
                    $a + (b - t) = a + (b + \neg{t})$ by \cref{real_sub,real_add_eq_subst}.
                    $a + (b + \neg{t}) = (a + b) + \neg{t}$ by \cref{real_add_assoc}.
                    $(a + b) + \neg{t} = (a + b) - t$ by \cref{real_sub}.
                    $a + (b - t) = c_0 - t$.
                    Therefore $c_0 - t < b$ by \cref{real_lt_subst_left}.
                    Hence $r(t)\mathrel{\realOrderRel}b$ by \cref{real_order_relation_intro}.
                \end{byCase}
            \end{subproof}
            Therefore $r(t)\in I$ by \cref{intervalclosed_rel}.
    \end{subproof}
    Hence $\img{r}{I}\subseteq I$
        by \cref{img_iff,continuous,funs_is_function,function_apply_iff,subseteq}.
    Therefore $r$ is continuous from $I$ to $I$ with respect to $T_I$ and $T_I$ by \cref{continuous_restrict_range}.
    Let $g = f\circ r$.
    $f$ is composable with $r$ by \cref{continuous,funs,img_dom,subseteq,funs_is_function}.
    Hence $g$ is continuous from $I$ to $X$ with respect to $T_I$ and $T$ by \cref{continuous_composite}.
    $a\in I$ and $b\in I$ by \cref{intervalclosed_rel}.
    $(a, k(a) - j(a))\in r$ and $(b, k(b) - j(b))\in r$.
    Hence $r(a) = k(a) - j(a)$ and $r(b) = k(b) - j(b)$
        by \cref{continuous,funs_is_function,function_apply_iff}.
    We have $k(a) = c_0$ and $k(b) = c_0$.
    We have $j(a) = a$ and $j(b) = b$
        by \cref{id_elem_funs,id_iff,funs_is_function,function_apply_iff}.
    $r(a) = c_0 - a$ and $r(b) = c_0 - b$.
    We have $c_0 = b + a$ by \cref{real_add_comm,real_add_eq_subst}.
    Therefore $c_0 - a = (b + a) + \neg{a}$ by \cref{real_sub,real_add_eq_subst}.
    $\neg{a}\in\reals$ by \cref{real_add_inverse}.
    $(b + a) + \neg{a} = b + (a + \neg{a})$ by \cref{real_add_assoc}.
    $a + \neg{a} = 0$ by \cref{real_add_inverse}.
    Therefore $r(a) = b + 0$ by \cref{real_add_eq_subst}.
    Show that $r(a) = b$.
    \begin{subproof}
        Follows by \cref{real_add_ident,real_add_comm}.
    \end{subproof}
    We have $c_0 - b = (a + b) + \neg{b}$ by \cref{real_sub,real_add_eq_subst}.
    $\neg{b}\in\reals$ by \cref{real_add_inverse}.
    $(a + b) + \neg{b} = a + (b + \neg{b})$ by \cref{real_add_assoc}.
    $b + \neg{b} = 0$ by \cref{real_add_inverse}.
    Therefore $r(b) = a + 0$ by \cref{real_add_eq_subst}.
    Show that $r(b) = a$.
    \begin{subproof}
        Follows by \cref{real_add_ident,real_add_comm}.
    \end{subproof}
    Hence $g(a) = f(r(a))$ and $g(b) = f(r(b))$
        by \cref{intervalclosed_rel,continuous,funs,funs_is_function,subseteq,circ_apply}.
    Show that $g(a) = y$ and $g(b) = x$.
    \begin{subproof}
        Trivial.
    \end{subproof}
    Therefore $g$ is a path in $X$ from $y$ to $x$ with respect to $T$ by \cref{path_in_space}.
    Hence there exists a function $g$ such that $g$ is a path in $X$ from $y$ to $x$ with respect to $T$.
\end{proof}

% from §25 helper lemma 17a: add a difference on the right
\begin{lemma}\label{real_add_sub_cancel_right}
    Suppose $a\in\reals$ and $b\in\reals$.
    Then $b + (a - b) = a$.
\end{lemma}
\begin{proof}
    $\neg{b}\in\reals$ by \cref{real_add_inverse}.
    $a - b\in\reals$ by \cref{real_sub,real_add_closed}.
    $b + (a - b) = (a - b) + b$ by \cref{real_add_comm}.
    $a - b = a + \neg{b}$ by \cref{real_sub}.
    We have $b + (a - b) = (a + \neg{b}) + b$ by \cref{real_add_eq_subst}.
    $(a + \neg{b}) + b = a + (\neg{b} + b)$ by \cref{real_add_assoc}.
    $\neg{b} + b = 0$ by \cref{real_add_inverse,real_add_comm}.
    Hence $b + (a - b) = a + 0$ by \cref{real_add_eq_subst}.
    Therefore $b + (a - b) = a$ by \cref{real_add_ident,real_add_comm}.
\end{proof}

% from §25 helper lemma 17b: translate an interval endpoint sum
\begin{lemma}\label{real_add_sub_sub_cancel}
    Suppose $b\in\reals$ and $c\in\reals$ and $d\in\reals$.
    Then $(b + (d - c)) + (c - b) = d$.
\end{lemma}
\begin{proof}
    $\neg{c}\in\reals$ and $\neg{b}\in\reals$ by \cref{real_add_inverse}.
    $d - c\in\reals$ by \cref{real_sub,real_add_closed}.
    $c - b\in\reals$ by \cref{real_sub,real_add_closed}.
    $b + (d - c) = (d - c) + b$ by \cref{real_add_comm}.
    Hence $(b + (d - c)) + (c - b) = (d - c) + (b + (c - b))$ by \cref{real_add_assoc,real_add_eq_subst}.
    $b + (c - b) = c$ by \cref{real_add_sub_cancel_right}.
    We have $(b + (d - c)) + (c - b) = (d - c) + c$ by \cref{real_add_eq_subst}.
    $(d - c) + c = c + (d - c)$ by \cref{real_add_comm}.
    Hence $(b + (d - c)) + (c - b) = d$ by \cref{real_add_eq_subst,real_add_sub_cancel_right}.
\end{proof}

% from §25 helper lemma 17: concatenate paths
\begin{lemma}\label{path_concatenation}
    Assume $T$ is a topology on $X$.
    Suppose $f$ is a path in $X$ from $x$ to $y$ with respect to $T$.
    Suppose $g$ is a path in $X$ from $y$ to $z$ with respect to $T$.
    Then there exists $h$ such that $h$ is a path in $X$ from $x$ to $z$ with respect to $T$.
\end{lemma}
\begin{proof}
    Take $a,b\in\reals$ such that
        $a\mathrel{\realOrderRel} b$ and
        $f$ is continuous from $\intervalclosedrel{\realOrderRel}{\reals}{a}{b}$ to $X$ with respect to
            $\subspaceTopology{\realOrderTopology}{\intervalclosedrel{\realOrderRel}{\reals}{a}{b}}$ and $T$ and
        $f(a) = x$ and $f(b) = y$
        by \cref{path_in_space}.
    Let $I_1 = \intervalclosedrel{\realOrderRel}{\reals}{a}{b}$.
    Let $T_1 = \subspaceTopology{\realOrderTopology}{I_1}$.
    Take $c,d\in\reals$ such that
        $c\mathrel{\realOrderRel} d$ and
        $g$ is continuous from $\intervalclosedrel{\realOrderRel}{\reals}{c}{d}$ to $X$ with respect to
            $\subspaceTopology{\realOrderTopology}{\intervalclosedrel{\realOrderRel}{\reals}{c}{d}}$ and $T$ and
        $g(c) = y$ and $g(d) = z$
        by \cref{path_in_space}.
    Let $I_2 = \intervalclosedrel{\realOrderRel}{\reals}{c}{d}$.
    Let $T_2 = \subspaceTopology{\realOrderTopology}{I_2}$.
    Let $s = c - b$.
    $b\in\reals$ and $c\in\reals$.
    $\neg{b}\in\reals$ by \cref{real_add_inverse}.
    $s\in\reals$ by \cref{real_sub,real_add_closed}.
    Let $e = b + (d - c)$.
    $\neg{c}\in\reals$ by \cref{real_add_inverse}.
    $d - c\in\reals$ by \cref{real_sub,real_add_closed}.
    Hence $e\in\reals$ by \cref{real_add_closed}.
    Let $I_3 = \intervalclosedrel{\realOrderRel}{\reals}{b}{e}$.
    Let $I = \intervalclosedrel{\realOrderRel}{\reals}{a}{e}$.
    Let $T_I = \subspaceTopology{\realOrderTopology}{I}$.
    Let $T_3 = \subspaceTopology{\realOrderTopology}{I_3}$.
    $\realOrderTopology$ is a topology on $\reals$ by
        \cref{standard_basis_is_basis,standard_topology_reals,basis_topology_is_topology,real_order_topology_eq_standard}.
    $I$ is a subspace of $\reals$ with respect to $\realOrderTopology$
        by \cref{intervalclosed_rel,subseteq,subspace_of}.
    $T_I$ is a topology on $I$ by \cref{subspace_topology_is_topology,subspace_of}.
    $I_1\subseteq\reals$ by \cref{intervalclosed_rel,subseteq}.
    $I_2\subseteq\reals$ by \cref{intervalclosed_rel,subseteq}.
    $I_3\subseteq\reals$ by \cref{intervalclosed_rel,subseteq}.
    Hence $I_1$ is a subspace of $\reals$ with respect to $\realOrderTopology$ and
        $I_2$ is a subspace of $\reals$ with respect to $\realOrderTopology$ and
        $I_3$ is a subspace of $\reals$ with respect to $\realOrderTopology$ by \cref{subspace_of}.
    Therefore $T_1$ is a topology on $I_1$ and $T_2$ is a topology on $I_2$ and $T_3$ is a topology on $I_3$
        by \cref{subspace_topology_is_topology,subspace_of}.
    Let $k = \{(t,s) \mid t\in I_3\}$.
    $k$ is a function from $I_3$ to $\reals$ by \cref{constant_map_function}.
    Show that for all $t\in I_3$ we have $k(t) = s$.
    \begin{subproof}
        Fix $t\in I_3$.
        $(t,s)\in k$.
        Hence $k(t) = s$ by \cref{funs_is_function,function_apply_iff}.
    \end{subproof}
    Hence $k$ is continuous from $I_3$ to $\reals$ with respect to $T_3$ and $\standardTopologyR$
        by \cref{constant_continuous,real_order_topology_eq_standard}.
    Let $j = \identity{I_3}$.
    $j$ is continuous from $I_3$ to $\reals$ with respect to $T_3$ and $\standardTopologyR$
        by \cref{inclusion_continuous,real_order_topology_eq_standard}.
    Let $u = \{(t, j(t) + k(t)) \mid t\in I_3\}$.
    $u$ is continuous from $I_3$ to $\reals$ with respect to $T_3$ and $\standardTopologyR$
        by \cref{real_function_addition_continuous}.
    Hence $u$ is continuous from $I_3$ to $\reals$ with respect to $T_3$ and $\realOrderTopology$
        by \cref{real_order_topology_eq_standard}.
    Show that for all $t$ such that $t\in I_3$ we have $u(t)\in I_2$.
    \begin{subproof}
            Fix $t$ such that $t\in I_3$.
            $t\in\reals$ by \cref{intervalclosed_rel}.
            $b\mathrel{\realOrderRel}t$ or $t = b$ by \cref{intervalclosed_rel}.
            $t\mathrel{\realOrderRel}e$ or $t = e$ by \cref{intervalclosed_rel}.
            $j(t) = t$ by \cref{id_elem_funs,id_iff,funs_is_function,function_apply_iff}.
            $(t, j(t) + k(t))\in u$.
            $u(t) = j(t) + k(t)$ by \cref{continuous,funs_is_function,function_apply_iff}.
            $k(t) = s$.
            $u(t) = t + s$ by \cref{real_add_eq_subst}.
            $s\in\reals$.
            $u(t)\in\reals$ by \cref{real_add_closed}.
            Show that $c\mathrel{\realOrderRel}u(t)$ or $u(t) = c$.
            \begin{subproof}
                \begin{byCase}
                \caseOf{$t = b$.}
                    $s = c - b$.
                    Hence $u(t) = c$ by \cref{real_add_sub_cancel_right}.
                \caseOf{$b\mathrel{\realOrderRel}t$.}
                    $b < t$ by \cref{real_order_relation_elim}.
                    $b\in\reals$ and $t\in\reals$ and $s\in\reals$.
                    $b + s < t + s$ by \cref{real_lt_add}.
                    Hence $b + s = c$ by \cref{real_add_sub_cancel_right}.
                    Therefore $c < u(t)$ by \cref{real_lt_subst_left,real_lt_subst_right}.
                    Hence $c\mathrel{\realOrderRel}u(t)$ by \cref{real_order_relation_intro}.
                \end{byCase}
            \end{subproof}
            Show that $u(t)\mathrel{\realOrderRel}d$ or $u(t) = d$.
            \begin{subproof}
                \begin{byCase}
                \caseOf{$t = e$.}
                    $s = c - b$.
                    $e = b + (d - c)$.
                    Hence $u(t) = d$ by \cref{real_add_sub_sub_cancel}.
                \caseOf{$t\mathrel{\realOrderRel}e$.}
                    $t < e$ by \cref{real_order_relation_elim}.
                    $t\in\reals$ and $e\in\reals$ and $s\in\reals$.
                    $t + s < e + s$ by \cref{real_lt_add}.
                    $s = c - b$.
                    $e = b + (d - c)$.
                    Hence $e + s = d$ by \cref{real_add_sub_sub_cancel}.
                    Therefore $u(t) < d$ by \cref{real_lt_subst_right}.
                    Hence $u(t)\mathrel{\realOrderRel}d$ by \cref{real_order_relation_intro}.
                \end{byCase}
            \end{subproof}
            Therefore $u(t)\in I_2$ by \cref{intervalclosed_rel}.
    \end{subproof}
    Show that for all $w$ such that $w\in\img{u}{I_3}$ we have $w\in I_2$.
    \begin{subproof}
        Fix $w$ such that $w\in\img{u}{I_3}$.
        Take $t\in I_3$ such that $(t,w)\in u$ by \cref{img_iff}.
        Hence $u(t) = w$ by \cref{continuous,funs_is_function,function_apply_iff}.
        $u(t)\in I_2$.
        Hence $w\in I_2$.
    \end{subproof}
    Therefore $\img{u}{I_3}\subseteq I_2$ by \cref{subseteq}.
    Hence $u$ is continuous from $I_3$ to $I_2$ with respect to $T_3$ and $T_2$
        by \cref{continuous_restrict_range}.
    Let $g_1 = g\circ u$.
    $g$ is composable with $u$ by \cref{continuous,funs,img_dom,subseteq,funs_is_function}.
    Hence $g_1$ is continuous from $I_3$ to $X$ with respect to $T_3$ and $T$ by \cref{continuous_composite}.
    $c < d$ by \cref{real_order_relation_elim}.
    $c\in\reals$ and $d\in\reals$.
    $c + \neg{c} < d + \neg{c}$ by \cref{real_lt_add}.
    $c + \neg{c} = 0$ by \cref{real_add_inverse}.
    Therefore $0 < d + \neg{c}$ by \cref{real_lt_subst_left}.
    $d + \neg{c} = d - c$ by \cref{real_sub}.
    Hence $0 < d - c$ by \cref{real_lt_subst_right}.
    $b\in\reals$ and $d - c\in\reals$.
    $0\in\reals$ by \cref{real_add_ident}.
    Therefore $0 + b < (d - c) + b$ by \cref{real_lt_add}.
    $0 + b = b$ by \cref{real_add_comm,real_add_ident}.
    Hence $b < (d - c) + b$ by \cref{real_lt_subst_left}.
    $(d - c) + b = b + (d - c)$ by \cref{real_add_comm}.
    Hence $b < b + (d - c)$ by \cref{real_lt_subst_right}.
    Therefore $b < e$ by \cref{real_lt_subst_right}.
    Hence $b\mathrel{\realOrderRel}e$ by \cref{real_order_relation_intro}.
    $b\in I_3$ and $e\in I_3$ by \cref{intervalclosed_rel}.
    $(b, j(b) + k(b))\in u$ and $(e, j(e) + k(e))\in u$.
    $u(b) = j(b) + k(b)$ and $u(e) = j(e) + k(e)$
        by \cref{continuous,funs_is_function,function_apply_iff}.
    $j(b) = b$ and $j(e) = e$
        by \cref{id_elem_funs,id_iff,funs_is_function,function_apply_iff}.
    We have $k(b) = s$ and $k(e) = s$.
    $u(b) = b + s$ and $u(e) = e + s$ by \cref{real_add_eq_subst}.
    Show that $u(b) = c$.
    \begin{subproof}
        Follows by \cref{real_add_sub_cancel_right}.
    \end{subproof}
    Show that $u(e) = d$.
    \begin{subproof}
        Follows by \cref{real_add_sub_sub_cancel}.
    \end{subproof}
    Therefore $g_1(b) = g(u(b))$ and $g_1(e) = g(u(e))$
        by \cref{intervalclosed_rel,continuous,funs,funs_is_function,subseteq,circ_apply}.
    Show that $g_1(b) = y$ and $g_1(e) = z$.
    \begin{subproof}
        Trivial.
    \end{subproof}
    Show that for all $t$ such that $t\in I_1\union I_3$ we have $t\in I$.
    \begin{subproof}
            Fix $t$ such that $t\in I_1\union I_3$.
            $t\in I_1$ or $t\in I_3$ by \cref{union_iff}.
            \begin{byCase}
            \caseOf{$t\in I_1$.}
                $a\mathrel{\realOrderRel}t$ or $t = a$ by \cref{intervalclosed_rel}.
                $t\mathrel{\realOrderRel}b$ or $t = b$ by \cref{intervalclosed_rel}.
                \begin{byCase}
                \caseOf{$t = b$.}
                    $t\mathrel{\realOrderRel}e$ by \cref{real_lt_subst_right}.
                    $t\in I$ by \cref{intervalclosed_rel}.
                \caseOf{$t\mathrel{\realOrderRel}b$.}
                    $t < b$ by \cref{real_order_relation_elim}.
                    $b < e$ by \cref{real_order_relation_elim}.
                    $t\in\reals$ by \cref{intervalclosed_rel}.
                    $b\in\reals$ and $e\in\reals$.
                    $t < e$ by \cref{real_lt_trans}.
                    $t\mathrel{\realOrderRel}e$ by \cref{real_order_relation_intro}.
                    $t\in I$ by \cref{intervalclosed_rel}.
                \end{byCase}
            \caseOf{$t\in I_3$.}
                $b\mathrel{\realOrderRel}t$ or $t = b$ by \cref{intervalclosed_rel}.
                $t\mathrel{\realOrderRel}e$ or $t = e$ by \cref{intervalclosed_rel}.
                \begin{byCase}
                \caseOf{$t = b$.}
                    $a\mathrel{\realOrderRel}t$ by \cref{real_lt_subst_right}.
                    $t\in I$ by \cref{intervalclosed_rel}.
                \caseOf{$b\mathrel{\realOrderRel}t$.}
                    $a < b$ by \cref{real_order_relation_elim}.
                    $b < t$ by \cref{real_order_relation_elim}.
                    $a\in\reals$ and $b\in\reals$.
                    $t\in\reals$ by \cref{intervalclosed_rel}.
                    $a < t$ by \cref{real_lt_trans}.
                    $a\mathrel{\realOrderRel}t$ by \cref{real_order_relation_intro}.
                    $t\in I$ by \cref{intervalclosed_rel}.
                \end{byCase}
            \end{byCase}
    \end{subproof}
    Show that $I_1\union I_3\subseteq I$.
    \begin{subproof}
        Follows by \cref{subseteq}.
    \end{subproof}
    Show that for all $t$ such that $t\in I$ we have $t\in I_1\union I_3$.
    \begin{subproof}
            Fix $t$ such that $t\in I$.
            $t\in\reals$ by \cref{intervalclosed_rel}.
            $a\mathrel{\realOrderRel}t$ or $t = a$ by \cref{intervalclosed_rel}.
            $t\mathrel{\realOrderRel}e$ or $t = e$ by \cref{intervalclosed_rel}.
            \begin{byCase}
            \caseOf{$t = b$.}
                $t\in I_1$ by \cref{intervalclosed_rel}.
                $t\in I_1\union I_3$ by \cref{union_iff}.
            \caseOf{$t\neq b$.}
                $b\in\reals$ and $t\in\reals$.
                $t\neq b$.
                $t < b$ or $b < t$ by \cref{real_lt_trichotomy}.
                \begin{byCase}
                \caseOf{$t < b$.}
                    $t\mathrel{\realOrderRel}b$ by \cref{real_order_relation_intro}.
                    $t\in I_1$ by \cref{intervalclosed_rel}.
                    $t\in I_1\union I_3$ by \cref{union_iff}.
                \caseOf{$b < t$.}
                    $b\mathrel{\realOrderRel}t$ by \cref{real_order_relation_intro}.
                    $t\in I_3$ by \cref{intervalclosed_rel}.
                    $t\in I_1\union I_3$ by \cref{union_iff}.
                \end{byCase}
            \end{byCase}
    \end{subproof}
    Show that $I\subseteq I_1\union I_3$.
    \begin{subproof}
        Follows by \cref{subseteq}.
    \end{subproof}
    Hence $I = I_1\union I_3$ by \cref{subseteq_antisymmetric}.
    $I_1\subseteq I$ and $I_3\subseteq I$ by \cref{union_upper_left,union_upper_right,subseteq}.
    $\realOrderRel$ is a simple order relation on $\reals$ and $\reals$ has more than one element
        by \cref{reals_linear_continuum,linear_continuum}.
    Hence $\realOrderTopology$ is the order topology on $\reals$ with respect to $\realOrderRel$
        by \cref{real_order_topology,order_topology}.
    $I$ is a subspace of $\reals$ with respect to $\realOrderTopology$ by \cref{subspace_of}.
    $I_1$ is closed in $\reals$ with respect to $\realOrderTopology$ by \cref{intervalclosedrel_closed_in_order_topology}.
    Therefore there exists $C$ such that $C$ is closed in $\reals$ with respect to $\realOrderTopology$ and $I_1 = C\inter I$.
    Show that $I_1$ is closed in $I$ with respect to $T_I$.
    \begin{subproof}
        Follows by \cref{closed_in_subspace_iff,subspace_of}.
    \end{subproof}
    $I_3$ is closed in $\reals$ with respect to $\realOrderTopology$ by \cref{intervalclosedrel_closed_in_order_topology}.
    Therefore there exists a set $D$ such that $D$ is closed in $\reals$ with respect to $\realOrderTopology$ and $I_3 = D\inter I$.
    Show that $I_3$ is closed in $I$ with respect to $T_I$.
    \begin{subproof}
        Follows by \cref{closed_in_subspace_iff,subspace_of}.
    \end{subproof}
    Show that for all $t$ such that $t\in I_1\inter I_3$ we have $f(t) = g_1(t)$.
    \begin{subproof}
        Follows by \cref{inter,intervalclosed_rel,real_order_relation_elim,real_lt_asym,real_lt_subst_right}.
    \end{subproof}
    Let $T_{I1} = \subspaceTopology{T_I}{I_1}$.
    Let $T_{I3} = \subspaceTopology{T_I}{I_3}$.
    Show that $T_{I1} = T_1$.
    \begin{subproof}
        Follows by \cref{intervalclosed_rel,union_upper_left,subseteq,subspace_subspace_topology,subspace_of,subspace_topology}.
    \end{subproof}
    Show that $T_{I3} = T_3$.
    \begin{subproof}
        Follows by \cref{intervalclosed_rel,union_upper_right,subseteq,subspace_subspace_topology,subspace_of,subspace_topology}.
    \end{subproof}
    We have $f$ is continuous from $I_1$ to $X$ with respect to $T_{I1}$ and $T$.
    $g_1$ is continuous from $I_3$ to $X$ with respect to $T_{I3}$ and $T$.
    Let $h = f\union g_1$.
    Therefore $h$ is continuous from $I$ to $X$ with respect to $T_I$ and $T$ by \cref{pasting_lemma}.
    $a\in I_1$ and $e\in I_3$ by \cref{intervalclosed_rel,real_order_relation_elim}.
    $(a, f(a))\in f$ and $(e, g_1(e))\in g_1$ by \cref{continuous,funs_elim,function_apply_iff}.
    Hence $h(a) = f(a)$ and $h(e) = g_1(e)$
        by \cref{continuous,funs_is_function,function_apply_iff,union_iff}.
    Show that $h(a) = x$ and $h(e) = z$.
    \begin{subproof}
        Trivial.
    \end{subproof}
    $a < b$ by \cref{real_order_relation_elim}.
    We have $b < e$ by \cref{real_lt_subst_right}.
    $a\in\reals$ and $b\in\reals$ and $e\in\reals$.
    Therefore $a < e$ by \cref{real_lt_trans}.
    Hence $a\mathrel{\realOrderRel}e$ by \cref{real_order_relation_intro}.
    We have $h$ is continuous from $\intervalclosedrel{\realOrderRel}{\reals}{a}{e}$ to $X$ with respect to
        $\subspaceTopology{\realOrderTopology}{\intervalclosedrel{\realOrderRel}{\reals}{a}{e}}$ and $T$.
    Therefore $h$ is a path in $X$ from $x$ to $z$ with respect to $T$ by \cref{path_in_space}.
    Hence there exists a function $h$ such that $h$ is a path in $X$ from $x$ to $z$ with respect to $T$.
\end{proof}

% from §25 helper lemma 18: path relation is a binary relation
\begin{lemma}\label{path_relation_binary}
    Assume $T$ is a topology on $X$.
    Then $\pathRelation{T}{X}$ is a binary relation on $X$.
\end{lemma}
\begin{proof}
    Follows by \cref{path_relation,subseteq}.
\end{proof}

% from §25 helper lemma 19: path relation is reflexive
\begin{lemma}\label{path_relation_reflexive}
    Assume $T$ is a topology on $X$.
    Then $\pathRelation{T}{X}$ is reflexive on $X$.
\end{lemma}
\begin{proof}
    Show that for all $x\in X$ we have $x\mathrel{\pathRelation{T}{X}}x$.
    \begin{subproof}
        Fix $x\in X$.
        We have $0\mathrel{\realOrderRel}1$ by
            \cref{real_add_ident,real_mul_ident,real_zero_lt_one,real_order_relation_intro}.
        Let $I = \intervalclosedrel{\realOrderRel}{\reals}{0}{1}$.
        Let $T_I = \subspaceTopology{\realOrderTopology}{I}$.
        For all $t$ such that $t\in I$ we have $t\in\reals$ by \cref{intervalclosed_rel}.
        Hence $I\subseteq\reals$ by \cref{subseteq}.
        $\realOrderTopology$ is a topology on $\reals$ by
            \cref{standard_basis_is_basis,standard_topology_reals,basis_topology_is_topology,real_order_topology_eq_standard}.
        Therefore $I$ is a subspace of $\reals$ with respect to $\realOrderTopology$ by \cref{subspace_of}.
        Hence $T_I$ is a topology on $I$ by \cref{subspace_topology_is_topology,subspace_of}.
        Let $f = \{(t,x) \mid t\in I\}$.
        $f$ is a function from $I$ to $X$ by \cref{constant_map_function}.
        Show that for all $t\in I$ we have $f(t) = x$.
        \begin{subproof}
            Fix $t\in I$.
            $(t,x)\in f$.
            Hence $f(t) = x$ by \cref{funs_is_function,function_apply_iff}.
        \end{subproof}
        Therefore $f$ is continuous from $I$ to $X$ with respect to $T_I$ and $T$ by \cref{constant_continuous}.
        $0\in I$ and $1\in I$ by \cref{intervalclosed_rel,real_order_relation_elim}.
        $f(0) = x$ and $f(1) = x$.
        Hence $x\mathrel{\pathRelation{T}{X}}x$ by \cref{path_in_space,path_relation,times_tuple_intro,fst_eq,snd_eq}.
    \end{subproof}
    Hence $\pathRelation{T}{X}$ is reflexive on $X$ by \cref{reflexive_on}.
\end{proof}

% from §25 helper lemma 20: path relation is symmetric
\begin{lemma}\label{path_relation_symmetric}
    Assume $T$ is a topology on $X$.
    Then $\pathRelation{T}{X}$ is symmetric.
\end{lemma}
\begin{proof}
    Follows by \cref{path_relation,times_tuple_elim,fst_eq,snd_eq,path_reverse,times_tuple_intro,symmetric}.
\end{proof}

% from §25 helper lemma 21: path relation is transitive
\begin{lemma}\label{path_relation_transitive}
    Assume $T$ is a topology on $X$.
    Then $\pathRelation{T}{X}$ is transitive.
\end{lemma}
\begin{proof}
    Follows by \cref{path_relation,times_tuple_elim,fst_eq,snd_eq,path_concatenation,times_tuple_intro,transitive}.
\end{proof}
% from §25 helper lemma 22: path relation is an equivalence
\begin{lemma}\label{path_relation_equivalence}
    Assume $T$ is a topology on $X$.
    Then $\pathRelation{T}{X}$ is an equivalence on $X$.
\end{lemma}
\begin{proof}
    Follows by \cref{path_relation_binary,path_relation_reflexive,path_relation_transitive,path_relation_symmetric}.
\end{proof}

% from §25 helper lemma 22a: a path in a subspace is a path in the ambient space
\begin{lemma}\label{path_in_subspace_implies_path_in_superspace}
    Assume $T$ is a topology on $X$.
    Suppose $A$ is a subspace of $X$ with respect to $T$.
    Suppose $f$ is a path in $A$ from $x$ to $y$ with respect to $\subspaceTopology{T}{A}$.
    Then there exists $g$ such that $g$ is a path in $X$ from $x$ to $y$ with respect to $T$.
\end{lemma}
\begin{proof}
    Take $a,b\in\reals$ such that
        $a\mathrel{\realOrderRel} b$ and
        $f$ is continuous from $\intervalclosedrel{\realOrderRel}{\reals}{a}{b}$ to $A$ with respect to
            $\subspaceTopology{\realOrderTopology}{\intervalclosedrel{\realOrderRel}{\reals}{a}{b}}$ and
            $\subspaceTopology{T}{A}$ and
        $f(a) = x$ and $f(b) = y$
        by \cref{path_in_space}.
    Let $I = \intervalclosedrel{\realOrderRel}{\reals}{a}{b}$.
    Let $T_I = \subspaceTopology{\realOrderTopology}{I}$.
    Let $j = \identity{A}$.
    $j$ is continuous from $A$ to $X$ with respect to $\subspaceTopology{T}{A}$ and $T$
        by \cref{inclusion_continuous}.
    Let $g = j\circ f$.
    $j$ is composable with $f$
        by \cref{continuous,funs,funs_is_function,funs_ran,subseteq,subspace_of}.
    $g$ is continuous from $I$ to $X$ with respect to $T_I$ and $T$ by \cref{continuous_composite}.
    $g(a) = x$ and $g(b) = y$
        by \cref{intervalclosed_rel,continuous,funs,funs_is_function,funs_type_apply,subseteq,circ_apply,id_apply}.
    $g$ is continuous from $\intervalclosedrel{\realOrderRel}{\reals}{a}{b}$ to $X$ with respect to
        $\subspaceTopology{\realOrderTopology}{\intervalclosedrel{\realOrderRel}{\reals}{a}{b}}$ and $T$.
    $g$ is a path in $X$ from $x$ to $y$ with respect to $T$ by \cref{path_in_space}.
    Show that there exists a function $h$ such that $h$ is a path in $X$ from $x$ to $y$ with respect to $T$.
    \begin{subproof}
        Trivial.
    \end{subproof}
\end{proof}

% from §25 helper lemma 23: union of path-connected subspaces with a common point
\begin{lemma}\label{union_path_connected_common_point}
    Assume $M\neq\emptyset$.
    Assume $T$ is a topology on $X$.
    Assume for all $A$ such that $A\in M$ we have $A\subseteq X$.
    Assume for all $A$ such that $A\in M$ we have
        $A$ is path connected with respect to $\subspaceTopology{T}{A}$.
    Assume $\inters{M}\neq\emptyset$.
    Then $\unions{M}$ is path connected with respect to $\subspaceTopology{T}{\unions{M}}$.
\end{lemma}
\begin{proof}
    Let $U = \unions{M}$.
    Let $T_U = \subspaceTopology{T}{U}$.
    $U$ is a subspace of $X$ with respect to $T$ by \cref{unions_iff,subseteq,subspace_of}.
    $T_U$ is a topology on $U$ by \cref{subspace_topology_is_topology,subspace_of}.
    Show that for all $x,y\in U$ there exists $h$ such that
        $h$ is a path in $U$ from $x$ to $y$ with respect to $T_U$.
    \begin{subproof}
        Fix $x$.
        Fix $y$ such that $x\in U$ and $y\in U$.
        Take $A\in M$ such that $x\in A$ by \cref{unions_iff}.
        Take $B\in M$ such that $y\in B$ by \cref{unions_iff}.
        Take $p$ such that $p\in\inters{M}$ by \cref{nonempty_inhabited}.
        $p\in A$ and $p\in B$ by \cref{inters_iff_forall}.
        Take $f$ such that $f$ is a path in $A$ from $x$ to $p$ with respect to $\subspaceTopology{T}{A}$
            by \cref{path_connected}.
        Take $g$ such that $g$ is a path in $B$ from $p$ to $y$ with respect to $\subspaceTopology{T}{B}$
            by \cref{path_connected}.
        Show that $A$ is a subspace of $U$ with respect to $T_U$.
        \begin{subproof}
            Follows by \cref{unions_iff,subseteq,subspace_of}.
        \end{subproof}
        Show that $B$ is a subspace of $U$ with respect to $T_U$.
        \begin{subproof}
            Follows by \cref{unions_iff,subseteq,subspace_of}.
        \end{subproof}
        Show that $\subspaceTopology{T_U}{A} = \subspaceTopology{T}{A}$.
        \begin{subproof}
            Follows by \cref{subspace_of,subspace_subspace_topology}.
        \end{subproof}
        Show that $\subspaceTopology{T_U}{B} = \subspaceTopology{T}{B}$.
        \begin{subproof}
            Follows by \cref{subspace_of,subspace_subspace_topology}.
        \end{subproof}
        Take $f_1$ such that $f_1$ is a path in $U$ from $x$ to $p$ with respect to $T_U$
            by \cref{path_in_subspace_implies_path_in_superspace}.
        Take $g_1$ such that $g_1$ is a path in $U$ from $p$ to $y$ with respect to $T_U$
            by \cref{path_in_subspace_implies_path_in_superspace}.
        Show that there exists $h$ such that $h$ is a path in $U$ from $x$ to $y$ with respect to $T_U$.
        \begin{subproof}
            Follows by \cref{path_concatenation}.
        \end{subproof}
    \end{subproof}
    $U$ is path connected with respect to $T_U$ by \cref{path_connected}.
\end{proof}


% from §25 helper lemma 24: initial segment of a path
\begin{lemma}\label{path_initial_segment}
    Assume $T$ is a topology on $X$.
    Suppose $f$ is a path in $X$ from $x$ to $y$ with respect to $T$.
    Suppose $t\in\dom{f}$.
    Then there exists $g$ such that $g$ is a path in $X$ from $x$ to $f(t)$ with respect to $T$.
\end{lemma}
\begin{proof}
    Take $a,b\in\reals$ such that
        $a\mathrel{\realOrderRel} b$ and
        $f$ is continuous from $\intervalclosedrel{\realOrderRel}{\reals}{a}{b}$ to $X$ with respect to
            $\subspaceTopology{\realOrderTopology}{\intervalclosedrel{\realOrderRel}{\reals}{a}{b}}$ and $T$ and
        $f(a) = x$ and $f(b) = y$
        by \cref{path_in_space}.
        Let $I = \intervalclosedrel{\realOrderRel}{\reals}{a}{b}$.
        Let $T_I = \subspaceTopology{\realOrderTopology}{I}$.
        $\dom{f} = I$ by \cref{continuous,funs}.
        $t\in I$ by \cref{continuous,funs,subseteq}.
        $f$ is a function by \cref{continuous,funs_is_function}.
    $x\in X$ by \cref{intervalclosed_rel,continuous,funs_type_apply}.
    \begin{byCase}
    \caseOf{$t = a$.}
        $f(t) = x$.
        $(x,x)\in\pathRelation{T}{X}$ by \cref{path_relation_reflexive,reflexive_on}.
        Take $g$ such that $g$ is a path in $X$ from $x$ to $x$ with respect to $T$
            by \cref{path_relation,fst_eq,snd_eq}.
        $g$ is a path in $X$ from $x$ to $f(t)$ with respect to $T$.
        Show that there exists $h$ such that $h$ is a path in $X$ from $x$ to $f(t)$ with respect to $T$.
        \begin{subproof}
            Trivial.
        \end{subproof}
    \caseOf{$t\neq a$.}
        $a\mathrel{\realOrderRel} t$ by \cref{intervalclosed_rel}.
        $t\in\reals$ by \cref{real_order_relation_elim}.
        Let $J = \intervalclosedrel{\realOrderRel}{\reals}{a}{t}$.
        $J\subseteq I$ by \cref{real_order_relation_simple,simple_order_relation,intervalclosedrel_convex,intervalclosed_rel,convex_in,subseteq}.
        $J$ is a subspace of $I$ with respect to $T_I$ by \cref{continuous,subspace_of}.
        Let $T_J = \subspaceTopology{T_I}{J}$.
        Let $g = \restrl{f}{J}$.
        $g$ is continuous from $J$ to $X$ with respect to $T_J$ and $T$
            by \cref{continuous_restrict_domain}.
        $\realOrderTopology$ is a topology on $\reals$ by
            \cref{standard_basis_is_basis,standard_topology_reals,basis_topology_is_topology,real_order_topology_eq_standard}.
        $T_J = \subspaceTopology{\realOrderTopology}{J}$ by \cref{intervalclosed_rel,subseteq,subspace_subspace_topology}.
        $g$ is continuous from $J$ to $X$ with respect to $\subspaceTopology{\realOrderTopology}{J}$ and $T$.
        $g(a) = x$ and $g(t) = f(t)$ by \cref{intervalclosed_rel,funs_is_function,restrl_of_function_apply}.
        $g$ is continuous from $\intervalclosedrel{\realOrderRel}{\reals}{a}{t}$ to $X$ with respect to
            $\subspaceTopology{\realOrderTopology}{\intervalclosedrel{\realOrderRel}{\reals}{a}{t}}$ and $T$.
        $g$ is a path in $X$ from $x$ to $f(t)$ with respect to $T$ by \cref{path_in_space}.
        Show that there exists $h$ such that $h$ is a path in $X$ from $x$ to $f(t)$ with respect to $T$.
        \begin{subproof}
            Trivial.
        \end{subproof}
    \end{byCase}
\end{proof}
% from §25 Item 9: path components are path connected
\begin{theorem}\label{path_component_path_connected}
    Assume $T$ is a topology on $X$.
    Suppose $C$ is a path component of $X$ with respect to $T$.
    Then $C$ is path connected with respect to $\subspaceTopology{T}{C}$.
\end{theorem}
\begin{proof}
    Take $x\in X$ such that $C = \equivalenceClass{\pathRelation{T}{X}}{x}$ by \cref{path_component}.
    $C$ is a subspace of $X$ with respect to $T$
        by \cref{downward_closure_iff,path_relation,times_tuple_elim,subseteq,subspace_of}.
    Let $T_C = \subspaceTopology{T}{C}$.
    $T_C$ is a topology on $C$ by \cref{subspace_topology_is_topology,subspace_of}.
    Show that for all $y,z\in C$ there exists $h$ such that
        $h$ is a path in $C$ from $y$ to $z$ with respect to $T_C$.
    \begin{subproof}
        Fix $y$.
        Fix $z$ such that $y\in C$ and $z\in C$.
        $x\mathrel{\pathRelation{T}{X}}y$ and $x\mathrel{\pathRelation{T}{X}}z$
            by \cref{downward_closure_iff,path_relation_symmetric,symmetric}.
        Take $f$ such that $f$ is a path in $X$ from $x$ to $y$ with respect to $T$
            by \cref{path_relation,fst_eq,snd_eq}.
        Take $g$ such that $g$ is a path in $X$ from $x$ to $z$ with respect to $T$
            by \cref{path_relation,fst_eq,snd_eq}.
        Take $a,b\in\reals$ such that
            $a\mathrel{\realOrderRel} b$ and
            $f$ is continuous from $\intervalclosedrel{\realOrderRel}{\reals}{a}{b}$ to $X$ with respect to
                $\subspaceTopology{\realOrderTopology}{\intervalclosedrel{\realOrderRel}{\reals}{a}{b}}$ and $T$ and
            $f(a) = x$ and $f(b) = y$
            by \cref{path_in_space}.
        Let $I = \intervalclosedrel{\realOrderRel}{\reals}{a}{b}$.
        Let $T_I = \subspaceTopology{\realOrderTopology}{I}$.
        Take $c,d\in\reals$ such that
            $c\mathrel{\realOrderRel} d$ and
            $g$ is continuous from $\intervalclosedrel{\realOrderRel}{\reals}{c}{d}$ to $X$ with respect to
                $\subspaceTopology{\realOrderTopology}{\intervalclosedrel{\realOrderRel}{\reals}{c}{d}}$ and $T$ and
            $g(c) = x$ and $g(d) = z$
            by \cref{path_in_space}.
        Let $J = \intervalclosedrel{\realOrderRel}{\reals}{c}{d}$.
        Let $T_J = \subspaceTopology{\realOrderTopology}{J}$.
        $\ran{f}\subseteq C$
            by \cref{fun_ran_iff,function_apply_elim,ran_intro,continuous,funs_is_function,funs_ran,subseteq,path_initial_segment,times_tuple_intro,path_relation,fst_eq,snd_eq,path_relation_symmetric,symmetric,downward_closure_iff}.
        $\ran{g}\subseteq C$
            by \cref{fun_ran_iff,function_apply_elim,ran_intro,continuous,funs_is_function,funs_ran,subseteq,path_initial_segment,times_tuple_intro,path_relation,fst_eq,snd_eq,path_relation_symmetric,symmetric,downward_closure_iff}.
        $f$ is continuous from $I$ to $C$ with respect to $T_I$ and $T_C$
            and $g$ is continuous from $J$ to $C$ with respect to $T_J$ and $T_C$
            by \cref{img_subseteq_ran,subseteq,continuous_restrict_range}.
        $f$ is a path in $C$ from $x$ to $y$ with respect to $T_C$
            and $g$ is a path in $C$ from $x$ to $z$ with respect to $T_C$
            by \cref{path_in_space}.
        Show that there exists $h$ such that $h$ is a path in $C$ from $y$ to $z$ with respect to $T_C$.
        \begin{subproof}
            Follows by \cref{path_reverse,path_concatenation}.
        \end{subproof}
    \end{subproof}
    $C$ is path connected with respect to $T_C$ by \cref{path_connected}.
\end{proof}

% from §25 Item 10: distinct path components are disjoint
\begin{theorem}\label{path_components_disjoint}
    Assume $T$ is a topology on $X$.
    Suppose $C$ is a path component of $X$ with respect to $T$.
    Suppose $D$ is a path component of $X$ with respect to $T$.
    Suppose $C\neq D$.
    Then $C$ is disjoint from $D$.
\end{theorem}
\begin{proof}
    Follows by \cref{path_component,path_relation_equivalence,equivalenceon_equivclasses_diseq_implies_disjoint}.
\end{proof}

% from §25 Item 11: path components cover the space
\begin{theorem}\label{path_components_cover}
    Assume $T$ is a topology on $X$.
    Then for all $x\in X$ there exists $C$ such that
        $C$ is a path component of $X$ with respect to $T$ and $x\in C$.
\end{theorem}
\begin{proof}
    Fix $x\in X$.
    Follows by \cref{path_relation_equivalence,equivclasses_inhabited_equivalenceon,path_component}.
\end{proof}

% from §25 Item 12: path-connected subspaces intersect only one path component
\begin{theorem}\label{path_connected_subspace_single_component}
    Assume $T$ is a topology on $X$.
    Suppose $A$ is a subspace of $X$ with respect to $T$.
    Suppose $A$ is path connected with respect to $\subspaceTopology{T}{A}$.
    Suppose $A\neq\emptyset$.
    Then for all $C,D$ such that
        $C$ is a path component of $X$ with respect to $T$ and
        $D$ is a path component of $X$ with respect to $T$ and
        $A$ intersects $C$ and $A$ intersects $D$
        we have $C = D$.
\end{theorem}
\begin{proof}
    Fix $C$.
    Fix $D$ such that
        $C$ is a path component of $X$ with respect to $T$ and
        $D$ is a path component of $X$ with respect to $T$ and
        $A$ intersects $C$ and $A$ intersects $D$.
    Take $x$ such that $x\in A\inter C$ by \cref{intersects,nonempty_inhabited}.
    Take $y$ such that $y\in A\inter D$ by \cref{intersects,nonempty_inhabited}.
    $x\in A$ and $x\in C$ and $y\in A$ and $y\in D$ by \cref{inter}.
    Take $f$ such that $f$ is a path in $A$ from $x$ to $y$ with respect to $\subspaceTopology{T}{A}$
        by \cref{path_connected}.
    Take $f_1$ such that $f_1$ is a path in $X$ from $x$ to $y$ with respect to $T$
        by \cref{path_in_subspace_implies_path_in_superspace}.
    $x\mathrel{\pathRelation{T}{X}}y$ by \cref{subspace_of,subseteq,times_tuple_intro,path_relation,fst_eq,snd_eq}.
    Take $c\in X$ such that $C = \equivalenceClass{\pathRelation{T}{X}}{c}$ by \cref{path_component}.
    Take $d\in X$ such that $D = \equivalenceClass{\pathRelation{T}{X}}{d}$ by \cref{path_component}.
    $x\in \downward{\pathRelation{T}{X}}{c}$ and $y\in \downward{\pathRelation{T}{X}}{d}$.
    $C = D$
        by \cref{downward_closure_iff,path_relation_symmetric,path_relation_transitive,symmetric,transitive,path_relation_equivalence,equiv_iff_equivclasses_eq}.
\end{proof}

% from §25 Item 13: locally connected at a point
\begin{definition}\label{locally_connected_at}
    $X$ is locally connected at $x$ with respect to $T$ iff
        for all $U$ such that $U$ is a neighborhood of $x$ in $X$ with respect to $T$
        there exists $V$ such that
            $V$ is a neighborhood of $x$ in $X$ with respect to $T$ and
            $V\subseteq U$ and
            $V$ is connected with respect to $\subspaceTopology{T}{V}$.
\end{definition}

% from §25 Item 14: locally connected
\begin{definition}\label{locally_connected}
    $X$ is locally connected with respect to $T$ iff
        for all $x\in X$ we have $X$ is locally connected at $x$ with respect to $T$.
\end{definition}

% from §25 Item 15: locally path connected at a point
\begin{definition}\label{locally_path_connected_at}
    $X$ is locally path connected at $x$ with respect to $T$ iff
        for all $U$ such that $U$ is a neighborhood of $x$ in $X$ with respect to $T$
        there exists $V$ such that
            $V$ is a neighborhood of $x$ in $X$ with respect to $T$ and
            $V\subseteq U$ and
            $V$ is path connected with respect to $\subspaceTopology{T}{V}$.
\end{definition}

% from §25 Item 16: locally path connected
\begin{definition}\label{locally_path_connected}
    $X$ is locally path connected with respect to $T$ iff
        for all $x\in X$ we have $X$ is locally path connected at $x$ with respect to $T$.
\end{definition}

% from §25 Item 17: locally connected iff components of open sets are open
\begin{theorem}\label{locally_connected_iff_components_open}
    Assume $T$ is a topology on $X$.
    Then $X$ is locally connected with respect to $T$ iff
        for all $U$ such that $U$ is open in $X$ with respect to $T$
        for all $C$ such that $C$ is a component of $U$ with respect to $\subspaceTopology{T}{U}$
        we have $C$ is open in $X$ with respect to $T$.
\end{theorem}
\begin{proof}
    Show that if $X$ is locally connected with respect to $T$ then
        for all $U$ such that $U$ is open in $X$ with respect to $T$
        for all $C$ such that $C$ is a component of $U$ with respect to $\subspaceTopology{T}{U}$
        we have $C$ is open in $X$ with respect to $T$.
    \begin{subproof}
        Assume $X$ is locally connected with respect to $T$.
        Fix $U$ such that $U$ is open in $X$ with respect to $T$.
        $U$ is a subspace of $X$ with respect to $T$
            by \cref{topology_on,open_in,powerset_elim,subseteq,subspace_of}.
        Let $T_U = \subspaceTopology{T}{U}$.
        $T_U$ is a topology on $U$ by \cref{subspace_topology_is_topology,subspace_of}.
        Fix $C$ such that $C$ is a component of $U$ with respect to $\subspaceTopology{T}{U}$.
        $C$ is a component of $U$ with respect to $T_U$.
        Show that for all $x$ such that $x\in C$ we have
            there exists $V$ such that $V$ is open in $X$ with respect to $T$ and $x\in V$ and $V\subseteq C$.
        \begin{subproof}
            Fix $x$ such that $x\in C$.
            $C\subseteq U$ by \cref{component,downward_closure_iff,component_relation,times_tuple_elim,subseteq}.
            $x\in U$ and $x\in X$ by \cref{subseteq,subspace_of}.
            $U$ is a neighborhood of $x$ in $X$ with respect to $T$ by \cref{neighborhood}.
            Take $W$ such that
                $W$ is a neighborhood of $x$ in $X$ with respect to $T$ and
                $W\subseteq U$ and
                $W$ is connected with respect to $\subspaceTopology{T}{W}$
                by \cref{locally_connected,locally_connected_at}.
            $W$ is open in $X$ with respect to $T$ and $x\in W$ by \cref{neighborhood}.
            $W$ is a subspace of $U$ with respect to $T_U$ by \cref{subseteq,subspace_of}.
            $\subspaceTopology{T_U}{W} = \subspaceTopology{T}{W}$ by \cref{subspace_subspace_topology,subspace_of,subseteq}.
            $W$ is connected with respect to $\subspaceTopology{T_U}{W}$.
            $W\neq\emptyset$ by \cref{inhabited_nonempty}.
            Show that $W\subseteq C$.
            \begin{subproof}
                Show that for all $y$ such that $y\in W$ we have $y\in C$.
                \begin{subproof}
                    Fix $y$ such that $y\in W$.
                    $y\in U$ by \cref{subseteq}.
                    Take $D$ such that $D$ is a component of $U$ with respect to $T_U$ and $y\in D$
                        by \cref{components_cover}.
                    $x\in W\inter C$ by \cref{inter}.
                    $W$ intersects $C$ by \cref{inhabited_nonempty,intersects}.
                    $y\in W\inter D$ by \cref{inter}.
                    $W$ intersects $D$ by \cref{inhabited_nonempty,intersects}.
                    $C = D$ by \cref{connected_subspace_single_component}.
                    $y\in C$.
                \end{subproof}
                $W\subseteq C$ by \cref{subseteq}.
            \end{subproof}
            Show that there exists $V$ such that $V$ is open in $X$ with respect to $T$ and $x\in V$ and $V\subseteq C$.
            \begin{subproof}
                Trivial.
            \end{subproof}
        \end{subproof}
        Let $M = \{V\in T \mid V\subseteq C\}$.
        For all $x$ such that $x\in C$ we have $x\in \unions{M}$ by \cref{subseteq,open_in,unions_iff}.
        $C\subseteq \unions{M}$ by \cref{subseteq}.
        Show that $\unions{M}\subseteq C$.
        \begin{subproof}
            Follows by \cref{unions_iff,subseteq}.
        \end{subproof}
        $C$ is open in $X$ with respect to $T$ by \cref{subseteq_antisymmetric,topology_on,subseteq,open_in}.
    \end{subproof}
    Show that if for all $U$ such that $U$ is open in $X$ with respect to $T$
        for all $C$ such that $C$ is a component of $U$ with respect to $\subspaceTopology{T}{U}$
        we have $C$ is open in $X$ with respect to $T$ then
        $X$ is locally connected with respect to $T$.
    \begin{subproof}
        Assume for all $U$ such that $U$ is open in $X$ with respect to $T$
            for all $C$ such that $C$ is a component of $U$ with respect to $\subspaceTopology{T}{U}$
            we have $C$ is open in $X$ with respect to $T$.
        Show that for all $x\in X$ we have $X$ is locally connected at $x$ with respect to $T$.
        \begin{subproof}
            Fix $x\in X$.
            Show that for all $U$ such that $U$ is a neighborhood of $x$ in $X$ with respect to $T$
                there exists $V$ such that
                    $V$ is a neighborhood of $x$ in $X$ with respect to $T$ and
                    $V\subseteq U$ and
                    $V$ is connected with respect to $\subspaceTopology{T}{V}$.
            \begin{subproof}
                Fix $U$ such that $U$ is a neighborhood of $x$ in $X$ with respect to $T$.
                $U$ is open in $X$ with respect to $T$ and $x\in U$ by \cref{neighborhood}.
                $U$ is a subspace of $X$ with respect to $T$
                    by \cref{topology_on,open_in,powerset_elim,subseteq,subspace_of}.
                Let $T_U = \subspaceTopology{T}{U}$.
                $T_U$ is a topology on $U$ by \cref{subspace_topology_is_topology,subspace_of}.
                Take $C$ such that
                    $C$ is a component of $U$ with respect to $T_U$ and $x\in C$
                    by \cref{components_cover}.
                $C$ is open in $X$ with respect to $T$.
                $C\subseteq U$ by \cref{component,downward_closure_iff,component_relation,times_tuple_elim,subseteq}.
                Let $T_C = \subspaceTopology{T}{C}$.
                $C$ is connected with respect to $T_C$
                    by \cref{component_connected,subspace_of,subspace_subspace_topology,subseteq}.
                $C$ is a neighborhood of $x$ in $X$ with respect to $T$ by \cref{neighborhood,open_in}.
                Show that there exists $V$ such that
                    $V$ is a neighborhood of $x$ in $X$ with respect to $T$ and
                    $V\subseteq U$ and
                    $V$ is connected with respect to $\subspaceTopology{T}{V}$.
                \begin{subproof}
                    Trivial.
                \end{subproof}
            \end{subproof}
            $X$ is locally connected at $x$ with respect to $T$ by \cref{locally_connected_at}.
        \end{subproof}
        $X$ is locally connected with respect to $T$ by \cref{locally_connected}.
    \end{subproof}
\end{proof}

% from §25 Item 18: locally path connected iff path components of open sets are open
\begin{theorem}\label{locally_path_connected_iff_components_open}
    Assume $T$ is a topology on $X$.
    Then $X$ is locally path connected with respect to $T$ iff
        for all $U$ such that $U$ is open in $X$ with respect to $T$
        for all $C$ such that $C$ is a path component of $U$ with respect to $\subspaceTopology{T}{U}$
        we have $C$ is open in $X$ with respect to $T$.
\end{theorem}
\begin{proof}
    Show that if $X$ is locally path connected with respect to $T$ then
        for all $U$ such that $U$ is open in $X$ with respect to $T$
        for all $C$ such that $C$ is a path component of $U$ with respect to $\subspaceTopology{T}{U}$
        we have $C$ is open in $X$ with respect to $T$.
    \begin{subproof}
        Assume $X$ is locally path connected with respect to $T$.
        Fix $U$ such that $U$ is open in $X$ with respect to $T$.
        $U$ is a subspace of $X$ with respect to $T$
            by \cref{topology_on,open_in,powerset_elim,subseteq,subspace_of}.
        Let $T_U = \subspaceTopology{T}{U}$.
        $T_U$ is a topology on $U$ by \cref{subspace_topology_is_topology,subspace_of}.
        Fix $C$ such that $C$ is a path component of $U$ with respect to $\subspaceTopology{T}{U}$.
        $C$ is a path component of $U$ with respect to $T_U$.
        Show that for all $x$ such that $x\in C$ we have
            there exists $V$ such that $V$ is open in $X$ with respect to $T$ and $x\in V$ and $V\subseteq C$.
        \begin{subproof}
                Fix $x$ such that $x\in C$.
                $C\subseteq U$ by \cref{path_component,downward_closure_iff,path_relation,times_tuple_elim,subseteq}.
                $x\in U$ and $x\in X$ by \cref{subseteq,subspace_of}.
                $U$ is a neighborhood of $x$ in $X$ with respect to $T$ by \cref{neighborhood}.
                Take $W$ such that
                    $W$ is a neighborhood of $x$ in $X$ with respect to $T$ and
                    $W\subseteq U$ and
                    $W$ is path connected with respect to $\subspaceTopology{T}{W}$
                    by \cref{locally_path_connected,locally_path_connected_at}.
                $W$ is open in $X$ with respect to $T$ and $x\in W$ by \cref{neighborhood}.
                $W$ is a subspace of $U$ with respect to $T_U$ by \cref{subseteq,subspace_of}.
                $\subspaceTopology{T_U}{W} = \subspaceTopology{T}{W}$ by \cref{subspace_subspace_topology,subspace_of,subseteq}.
                $W$ is path connected with respect to $\subspaceTopology{T_U}{W}$.
                $W\neq\emptyset$ by \cref{inhabited_nonempty}.
                Show that $W\subseteq C$.
                \begin{subproof}
                    Show that for all $y$ such that $y\in W$ we have $y\in C$.
                    \begin{subproof}
                        Fix $y$ such that $y\in W$.
                        $y\in U$ by \cref{subseteq}.
                        Take $D$ such that $D$ is a path component of $U$ with respect to $T_U$ and $y\in D$
                            by \cref{path_components_cover}.
                        $x\in W\inter C$ by \cref{inter}.
                        $W$ intersects $C$ by \cref{inhabited_nonempty,intersects}.
                        $y\in W\inter D$ by \cref{inter}.
                        $W$ intersects $D$ by \cref{inhabited_nonempty,intersects}.
                        $C = D$ by \cref{path_connected_subspace_single_component}.
                        $y\in C$.
                    \end{subproof}
                    $W\subseteq C$ by \cref{subseteq}.
                \end{subproof}
                Show that there exists $V$ such that $V$ is open in $X$ with respect to $T$ and $x\in V$ and $V\subseteq C$.
                \begin{subproof}
                    Trivial.
                \end{subproof}
        \end{subproof}
        Let $M = \{V\in T \mid V\subseteq C\}$.
        For all $x$ such that $x\in C$ we have $x\in \unions{M}$ by \cref{subseteq,open_in,unions_iff}.
        $C\subseteq \unions{M}$ by \cref{subseteq}.
        Show that $\unions{M}\subseteq C$.
        \begin{subproof}
            Follows by \cref{unions_iff,subseteq}.
        \end{subproof}
        $C$ is open in $X$ with respect to $T$ by \cref{subseteq_antisymmetric,topology_on,subseteq,open_in}.
    \end{subproof}
    Show that if for all $U$ such that $U$ is open in $X$ with respect to $T$
        for all $C$ such that $C$ is a path component of $U$ with respect to $\subspaceTopology{T}{U}$
        we have $C$ is open in $X$ with respect to $T$ then
        $X$ is locally path connected with respect to $T$.
    \begin{subproof}
        Assume for all $U$ such that $U$ is open in $X$ with respect to $T$
            for all $C$ such that $C$ is a path component of $U$ with respect to $\subspaceTopology{T}{U}$
            we have $C$ is open in $X$ with respect to $T$.
        Show that for all $x\in X$ we have $X$ is locally path connected at $x$ with respect to $T$.
        \begin{subproof}
            Fix $x\in X$.
            Show that for all $U$ such that $U$ is a neighborhood of $x$ in $X$ with respect to $T$
                there exists $V$ such that
                    $V$ is a neighborhood of $x$ in $X$ with respect to $T$ and
                    $V\subseteq U$ and
                    $V$ is path connected with respect to $\subspaceTopology{T}{V}$.
            \begin{subproof}
                Fix $U$ such that $U$ is a neighborhood of $x$ in $X$ with respect to $T$.
                $U$ is open in $X$ with respect to $T$ and $x\in U$ by \cref{neighborhood}.
                $U$ is a subspace of $X$ with respect to $T$
                    by \cref{topology_on,open_in,powerset_elim,subseteq,subspace_of}.
                Let $T_U = \subspaceTopology{T}{U}$.
                $T_U$ is a topology on $U$ by \cref{subspace_topology_is_topology,subspace_of}.
                Take $C$ such that
                    $C$ is a path component of $U$ with respect to $T_U$ and $x\in C$
                    by \cref{path_components_cover}.
                $C$ is open in $X$ with respect to $T$.
                $C\subseteq U$ by \cref{path_component,downward_closure_iff,path_relation,times_tuple_elim,subseteq}.
                Let $T_C = \subspaceTopology{T}{C}$.
                $C$ is path connected with respect to $T_C$
                    by \cref{path_component_path_connected,subspace_of,subspace_subspace_topology,subseteq}.
                $C$ is a neighborhood of $x$ in $X$ with respect to $T$ by \cref{neighborhood,open_in}.
                Show that there exists $V$ such that
                    $V$ is a neighborhood of $x$ in $X$ with respect to $T$ and
                    $V\subseteq U$ and
                    $V$ is path connected with respect to $\subspaceTopology{T}{V}$.
                \begin{subproof}
                    Trivial.
                \end{subproof}
            \end{subproof}
            $X$ is locally path connected at $x$ with respect to $T$ by \cref{locally_path_connected_at}.
        \end{subproof}
        $X$ is locally path connected with respect to $T$ by \cref{locally_path_connected}.
    \end{subproof}
\end{proof}

% from §25 Item 19: path components lie in components
\begin{theorem}\label{path_component_subset_component}
    Assume $T$ is a topology on $X$.
    Suppose $P$ is a path component of $X$ with respect to $T$.
    Then there exists $C$ such that $C$ is a component of $X$ with respect to $T$ and $P\subseteq C$.
\end{theorem}
\begin{proof}
    Take $x\in X$ such that $P = \equivalenceClass{\pathRelation{T}{X}}{x}$ by \cref{path_component}.
    For all $y$ such that $y\in P$ we have
        $(y,x)\in X\times X$ and there exists $f$ such that
        $f$ is a path in $X$ from $\fst{(y,x)}$ to $\snd{(y,x)}$ with respect to $T$
        by \cref{downward_closure_iff,path_relation}.
    $P$ is a subspace of $X$ with respect to $T$ by \cref{times_tuple_elim,subseteq,subspace_of}.
    Let $T_P = \subspaceTopology{T}{P}$.
    $T_P$ is a topology on $P$ by \cref{subspace_topology_is_topology,subspace_of}.
    $P$ is connected with respect to $T_P$
        by \cref{path_component_path_connected,path_connected_implies_connected}.
    Take $D$ such that
        $D$ is a component of $X$ with respect to $T$ and $x\in D$
        by \cref{components_cover}.
    For all $y$ such that $y\in P$ we have $y\in D$
        by \cref{times_tuple_elim,components_cover,path_relation_equivalence,equivclasses_inhabited_equivalenceon,inter,emptyset,intersects,connected_subspace_single_component,subspace_of}.
    Show that there exists $C$ such that $C$ is a component of $X$ with respect to $T$ and $P\subseteq C$.
    \begin{subproof}
        Follows by \cref{subseteq}.
    \end{subproof}
\end{proof}


% from §25 Item 20: components equal path components in locally path connected spaces
\begin{theorem}\label{components_equal_path_components_locally_path_connected}
    Assume $T$ is a topology on $X$.
    Suppose $X$ is locally path connected with respect to $T$.
    Then for all $C$ such that $C$ is a component of $X$ with respect to $T$
        there exists $P$ such that $P$ is a path component of $X$ with respect to $T$ and $C = P$.
\end{theorem}
\begin{proof}
    Fix $C$ such that $C$ is a component of $X$ with respect to $T$.
    Take $x\in X$ such that $C = \equivalenceClass{\componentRelation{T}{X}}{x}$ by \cref{component}.
    $x\in C$ by \cref{component_relation_equivalence,equivclasses_inhabited_equivalenceon}.
    $C$ is a subspace of $X$ with respect to $T$
        by \cref{downward_closure_iff,component_relation,times_tuple_elim,subseteq,subspace_of}.
    Let $T_X = \subspaceTopology{T}{X}$.
    $T_X$ is a topology on $X$ by \cref{subseteq_refl,subspace_of,subspace_topology_is_topology}.
    Show that $T\subseteq T_X$.
    \begin{subproof}
        Follows by \cref{topology_on,powerset_elim,subseteq,inter_absorb_supseteq_left,subspace_topology}.
    \end{subproof}
    Show that $T_X\subseteq T$.
    \begin{subproof}
        Follows by \cref{subspace_topology,topology_on,powerset_elim,subseteq,inter_absorb_supseteq_left}.
    \end{subproof}
    $T_X = T$ by \cref{subseteq_antisymmetric}.
    For all $D$ such that $D$ is a path component of $X$ with respect to $T$
        we have $D$ is open in $X$ with respect to $T$
        by \cref{topology_on,open_in,locally_path_connected_iff_components_open}.
    Take $P$ such that $P$ is a path component of $X$ with respect to $T$ and $x\in P$
        by \cref{path_components_cover}.
    Suppose not.
    Take $C_1$ such that
        $C_1$ is a component of $X$ with respect to $T$ and $P\subseteq C_1$
        by \cref{path_component_subset_component}.
    $C = C_1$ by \cref{components_disjoint,subseteq,disjoint}.
    $P\subset C$ by \cref{subseteq,subset}.
    Take $y$ such that $y\in C\setminus P$
        by \cref{difference_with_proper_subset_is_inhabited,nonempty_inhabited}.
    $y\in C$ and $y\notin P$ by \cref{setminus_elim_left,setminus_elim_right}.
    Let $M = \{D\in\pow{X} \mid
        \text{$D$ is a path component of $X$ with respect to $T$ and $D\subseteq C$ and $D\neq P$}\}$.
    Let $Q = \unions{M}$.
    Show that $Q\subseteq C\setminus P$.
    \begin{subproof}
        Follows by \cref{unions_iff,subseteq,path_components_disjoint,disjoint,setminus_intro}.
    \end{subproof}
    Show that for all $z$ such that $z\in C\setminus P$ we have $z\in Q$.
    \begin{subproof}
        Fix $z$ such that $z\in C\setminus P$.
        $z\in C$ and $z\notin P$ by \cref{setminus_elim_left,setminus_elim_right}.
        Take $D$ such that $D$ is a path component of $X$ with respect to $T$ and $z\in D$
            by \cref{subspace_of,subseteq,path_components_cover}.
        Take $C_2$ such that
            $C_2$ is a component of $X$ with respect to $T$ and $D\subseteq C_2$
            by \cref{path_component_subset_component}.
        $C = C_2$ by \cref{components_disjoint,subseteq,disjoint}.
        $D\in M$ by \cref{subspace_of,subseteq_transitive,powerset_intro}.
        $z\in Q$ by \cref{unions_iff}.
    \end{subproof}
    $C\setminus P\subseteq Q$ by \cref{subseteq}.
    $Q = C\setminus P$ by \cref{subseteq_antisymmetric}.
    Let $T_C = \subspaceTopology{T}{C}$.
    $T_C$ is a topology on $C$ by \cref{subspace_topology_is_topology,subspace_of}.
    $P$ is open in $C$ with respect to $T_C$
        by \cref{open_in,inter_absorb_supseteq_left,subspace_topology}.
    $Q$ is open in $C$ with respect to $T_C$
        by \cref{open_in,subseteq,topology_on,inter_absorb_supseteq_left,setminus_subseteq,subspace_topology}.
    $x\in P$ and $y\in Q$.
    $P\neq\emptyset$ and $Q\neq\emptyset$ by \cref{emptyset}.
    $P$ is disjoint from $Q$ by \cref{setminus_disjoint,inter,emptyset,disjoint}.
    $P\union Q = C$ by \cref{setminus_partition,subseteq}.
    $C$ is separated by $P$ and $Q$ with respect to $T_C$ by \cref{separation}.
    $C$ is not connected with respect to $T_C$ by \cref{connected}.
    Contradiction by \cref{component_connected}.
    $C = P$.
    Show that there exists $R$ such that $R$ is a path component of $X$ with respect to $T$ and $C = R$.
    \begin{subproof}
        Trivial.
    \end{subproof}
\end{proof}


\section{§26 Compact Spaces}

% from §26 Item 1: covering
\begin{definition}\label{covering}
    $A$ is a covering of $X$ iff $X\subseteq\unions{A}$.
\end{definition}

% from §26 Item 2: open covering
\begin{definition}\label{open_covering}
    $A$ is an open covering of $X$ with respect to $T$ iff
        $A$ is a covering of $X$ and
        for all $U$ such that $U\in A$ we have $U$ is open in $X$ with respect to $T$.
\end{definition}

% from §26 Item 3: compact space
\begin{definition}\label{compact}
    $X$ is compact with respect to $T$ iff
        for all $A$ such that $A$ is an open covering of $X$ with respect to $T$
        there exists $B$ such that
            $B\subseteq A$ and
            $B$ is finite and
            $B$ is a covering of $X$.
\end{definition}

% from §26 Item 4: compactness via coverings open in the ambient space
\begin{lemma}\label{compact_subspace_open_covering_forward}
    Assume $Y$ is a subspace of $X$ with respect to $T$.
    Let $T_Y = \subspaceTopology{T}{Y}$.
    Suppose $Y$ is compact with respect to $T_Y$.
    Then for all $A$ such that
            $A$ is a covering of $Y$ and
            for all $U$ such that $U\in A$ we have $U$ is open in $X$ with respect to $T$
        there exists $B$ such that
            $B\subseteq A$ and
            $B$ is finite and
            $B$ is a covering of $Y$.
\end{lemma}
\begin{proof}
    $T$ is a topology on $X$ and $Y\subseteq X$ by \cref{subspace_of}.
    $T_Y$ is a topology on $Y$ by \cref{subspace_topology_is_topology}.
    Fix $A$ such that
        $A$ is a covering of $Y$ and
        for all $U$ such that $U\in A$ we have $U$ is open in $X$ with respect to $T$.
    Let $A_Y = \{V \mid \exists U\in A. V = Y\inter U\}$.
    For all $V$ such that $V\in A_Y$ we have $V$ is open in $Y$ with respect to $T_Y$
        by \cref{open_in,subspace_topology}.
    $Y\subseteq\unions{A}$ by \cref{covering}.
    Show that for all $y$ such that $y\in Y$ we have $y\in\unions{A_Y}$.
    \begin{subproof}
        Fix $y$ such that $y\in Y$.
        $y\in\unions{A}$ by \cref{subseteq}.
        Take $U$ such that $U\in A$ and $y\in U$ by \cref{unions_iff}.
        Let $V = Y\inter U$.
        $V\in A_Y$.
        $y\in V$ by \cref{inter_intro}.
        $y\in\unions{A_Y}$ by \cref{unions_iff}.
    \end{subproof}
    $A_Y$ is a covering of $Y$ by \cref{subseteq,covering}.
    $A_Y$ is an open covering of $Y$ with respect to $T_Y$ by \cref{open_covering}.
    Take $B_Y$ such that
        $B_Y\subseteq A_Y$ and $B_Y$ is finite and $B_Y$ is a covering of $Y$
        by \cref{compact}.
    Let $R = \{w\in B_Y\times A \mid \exists V\in B_Y.\exists U\in A. w = (V,U) \land V = Y\inter U\}$.
    For all $w$ such that $w\in R$ there exist $V,U$ such that $w = (V,U)$.
    $R$ is a relation by \cref{relation}.
    Show that for all $V\in B_Y$ there exists $U$ such that $V\mathrel{R} U$.
    \begin{subproof}
        Fix $V\in B_Y$.
        $V\in A_Y$ by \cref{subseteq}.
        Take $U\in A$ such that $V = Y\inter U$.
        $(V,U)\in B_Y\times A$ by \cref{times_tuple_intro}.
        $(V,U)\in R$.
        $V\mathrel{R} U$.
    \end{subproof}
    Take $g$ such that $g$ is a function on $B_Y$ and for all $V\in B_Y$ we have $V\mathrel{R}\apply{g}{V}$ by \cref{choice_function}.
    Let $B = \img{g}{B_Y}$.
    $B\subseteq A$ by \cref{img_iff,function_apply_iff,pair_eq_iff,subseteq}.
    $B$ is finite by \cref{finite_image_function}.
    For all $y$ such that $y\in Y$ we have $y\in\unions{B}$
        by \cref{covering,subseteq,function_apply_iff,pair_eq_iff,inter_intro,inter_elim_right,function_apply_elim,img_iff,unions_iff}.
    $B$ is a covering of $Y$ by \cref{subseteq,covering}.
\end{proof}

\begin{lemma}\label{compact_subspace_open_covering_backward}
    Assume $Y$ is a subspace of $X$ with respect to $T$.
    Let $T_Y = \subspaceTopology{T}{Y}$.
    Suppose for all $A$ such that
            $A$ is a covering of $Y$ and
            for all $U$ such that $U\in A$ we have $U$ is open in $X$ with respect to $T$
        there exists $B$ such that
            $B\subseteq A$ and
            $B$ is finite and
            $B$ is a covering of $Y$.
    Then $Y$ is compact with respect to $T_Y$.
\end{lemma}
\begin{proof}
    $T$ is a topology on $X$ and $Y\subseteq X$ by \cref{subspace_of}.
    $T_Y$ is a topology on $Y$ by \cref{subspace_topology_is_topology}.
    It suffices to show that for all $A$ such that $A$ is an open covering of $Y$ with respect to $T_Y$
        there exists $B$ such that $B\subseteq A$ and $B$ is finite and $B$ is a covering of $Y$
        by \cref{compact}.
    Fix $A$ such that $A$ is an open covering of $Y$ with respect to $T_Y$.
    Let $R = \{w\in A\times T \mid \exists V\in A.\exists U\in T. w = (V,U) \land V = Y\inter U\}$.
    For all $w$ such that $w\in R$ there exist $V,U$ such that $w = (V,U)$.
    $R$ is a relation by \cref{relation}.
    For all $V\in A$ there exists $U\in T$ such that $V = Y\inter U$
        by \cref{open_covering,open_in,subspace_topology}.
    For all $V\in A$ there exists $U$ such that $V\mathrel{R} U$.
    Take $g$ such that $g$ is a function on $A$ and for all $V\in A$ we have $V\mathrel{R}\apply{g}{V}$ by \cref{choice_function}.
    Let $C = \img{g}{A}$.
    For all $U$ such that $U\in C$ we have $U$ is open in $X$ with respect to $T$
        by \cref{img_iff,function_apply_iff,pair_eq_iff,open_in}.
    For all $y$ such that $y\in Y$ we have $y\in\unions{C}$
        by \cref{open_covering,covering,subseteq,unions_iff,function_apply_iff,pair_eq_iff,inter_intro,inter_elim_right,function_apply_elim,img_iff}.
    $Y\subseteq\unions{C}$ by \cref{subseteq}.
    $C$ is a covering of $Y$ by \cref{covering}.
    Take $B$ such that $B\subseteq C$ and $B$ is finite and $B$ is a covering of $Y$ by assumption.
    Let $S = \{w\in B\times A \mid \exists U\in B.\exists V\in A. w = (U,V) \land \apply{g}{V} = U\}$.
    For all $w$ such that $w\in S$ there exist $U,V$ such that $w = (U,V)$.
    $S$ is a relation by \cref{relation}.
    For all $U\in B$ there exists $V\in A$ such that $\apply{g}{V} = U$
        by \cref{subseteq,img_iff,function_apply_iff}.
    For all $U\in B$ there exists $V$ such that $U\mathrel{S} V$.
    Take $h$ such that $h$ is a function on $B$ and for all $U\in B$ we have $U\mathrel{S}\apply{h}{U}$ by \cref{choice_function}.
    Let $B_1 = \img{h}{B}$.
    $B_1\subseteq A$ by \cref{img_iff,function_apply_iff,pair_eq_iff,subseteq}.
    $B_1$ is finite by \cref{finite_image_function}.
    For all $y$ such that $y\in Y$ we have $y\in\unions{B_1}$
        by \cref{covering,subseteq,unions_iff,function_apply_iff,pair_eq_iff,inter_intro,open_covering,open_in,subspace_topology,function_apply_elim,img_iff}.
    $Y\subseteq\unions{B_1}$ by \cref{subseteq}.
    $B_1$ is a covering of $Y$ by \cref{covering}.
    Show that there exists $D$ such that $D\subseteq A$ and $D$ is finite and $D$ is a covering of $Y$.
    \begin{subproof}
        Trivial.
    \end{subproof}
\end{proof}

\begin{lemma}\label{compact_subspace_open_covering}
    Assume $Y$ is a subspace of $X$ with respect to $T$.
    Let $T_Y = \subspaceTopology{T}{Y}$.
    Then $Y$ is compact with respect to $T_Y$ iff
        for all $A$ such that
            $A$ is a covering of $Y$ and
            for all $U$ such that $U\in A$ we have $U$ is open in $X$ with respect to $T$
        there exists $B$ such that
            $B\subseteq A$ and
            $B$ is finite and
            $B$ is a covering of $Y$.
\end{lemma}
\begin{proof}
    Show that if $Y$ is compact with respect to $T_Y$ then
        for all $A$ such that
            $A$ is a covering of $Y$ and
            for all $U$ such that $U\in A$ we have $U$ is open in $X$ with respect to $T$
        there exists $B$ such that
            $B\subseteq A$ and
            $B$ is finite and
            $B$ is a covering of $Y$.
    \begin{subproof}
        Follows by \cref{compact_subspace_open_covering_forward}.
    \end{subproof}
    Show that if for all $A$ such that
            $A$ is a covering of $Y$ and
            for all $U$ such that $U\in A$ we have $U$ is open in $X$ with respect to $T$
        there exists $B$ such that
            $B\subseteq A$ and
            $B$ is finite and
            $B$ is a covering of $Y$ then
        $Y$ is compact with respect to $T_Y$.
    \begin{subproof}
        Follows by \cref{compact_subspace_open_covering_backward}.
    \end{subproof}
\end{proof}
% from §26 Item 5: closed subspace of a compact space is compact
\begin{theorem}\label{closed_subspace_compact}
    Assume $T$ is a topology on $X$.
    Assume $X$ is compact with respect to $T$.
    Suppose $Y\subseteq X$.
    Suppose $Y$ is closed in $X$ with respect to $T$.
    Let $T_Y = \subspaceTopology{T}{Y}$.
    Then $Y$ is compact with respect to $T_Y$.
\end{theorem}
\begin{proof}
    $Y$ is a subspace of $X$ with respect to $T$ by \cref{subspace_of}.
    It suffices to show that for all $A$ such that
        $A$ is a covering of $Y$ and
        for all $U$ such that $U\in A$ we have $U$ is open in $X$ with respect to $T$
        there exists $B$ such that
            $B\subseteq A$ and
            $B$ is finite and
            $B$ is a covering of $Y$
        by \cref{compact_subspace_open_covering}.
    Fix $A$ such that
        $A$ is a covering of $Y$ and
        for all $U$ such that $U\in A$ we have $U$ is open in $X$ with respect to $T$.
    Let $C = A\union \{X\setminus Y\}$.
    For all $U$ such that $U\in C$ we have $U$ is open in $X$ with respect to $T$
        by \cref{union_iff,singleton_iff,closed_in}.
    Show that for all $x$ such that $x\in X$ we have $x\in\unions{C}$.
    \begin{subproof}
        Fix $x$ such that $x\in X$.
        \begin{byCase}
        \caseOf{$x\in Y$.}
            Take $U\in A$ such that $x\in U$ by \cref{covering,subseteq,unions_iff}.
            $x\in\unions{C}$ by \cref{union_iff,unions_iff}.
        \caseOf{$x\notin Y$.}
            $x\in\unions{C}$ by \cref{setminus_intro,singleton_iff,union_iff,unions_iff}.
        \end{byCase}
    \end{subproof}
    $C$ is a covering of $X$ by \cref{subseteq,covering}.
    $C$ is an open covering of $X$ with respect to $T$ by \cref{open_covering}.
    Take $B$ such that $B\subseteq C$ and $B$ is finite and $B$ is a covering of $X$ by \cref{compact}.
    \begin{byCase}
    \caseOf{$X\setminus Y\in B$.}
        Let $B_1 = B\setminus \{X\setminus Y\}$.
        For all $y$ such that $y\in Y$ there exists $U\in B$ such that $y\in U$
            by \cref{covering,subseteq,unions_iff}.
        For all $y$ such that $y\in Y$ there exists $U\in B_1$ such that $y\in U$
            by \cref{setminus_intro,singleton_iff,setminus_elim_right}.
        For all $y$ such that $y\in Y$ we have $y\in\unions{B_1}$ by \cref{unions_iff}.
        $B_1$ is a covering of $Y$ by \cref{subseteq,covering}.
        $B_1\subseteq B$ by \cref{setminus_subseteq}.
        $B_1$ is finite by \cref{setminus_subseteq,subseteq_finite_is_finite}.
        $B_1\subseteq A$ by \cref{subseteq,setminus_elim_right,union_iff,subseteq}.
        Show that there exists $D$ such that $D\subseteq A$ and $D$ is finite and $D$ is a covering of $Y$.
        \begin{subproof}
            Trivial.
        \end{subproof}
    \caseOf{$X\setminus Y\notin B$.}
        Show that there exists $D$ such that $D\subseteq A$ and $D$ is finite and $D$ is a covering of $Y$.
        \begin{subproof}
            Follows by \cref{subseteq,union_iff,singleton_iff,covering}.
        \end{subproof}
    \end{byCase}
\end{proof}
% from §26 Item 6: compact subspace of Hausdorff is closed
\begin{theorem}\label{compact_subspace_closed}
    Assume $T$ is a topology on $X$.
    Assume $X$ is Hausdorff with respect to $T$.
    Suppose $Y\subseteq X$.
    Let $T_Y = \subspaceTopology{T}{Y}$.
    Suppose $Y$ is compact with respect to $T_Y$.
    Then $Y$ is closed in $X$ with respect to $T$.
\end{theorem}
\begin{proof}
    \begin{byCase}
    \caseOf{$Y = \emptyset$.}
        Show that $Y$ is closed in $X$ with respect to $T$.
        \begin{subproof}
            Follows by \cref{closed_emptyset}.
        \end{subproof}
    \caseOf{$Y \neq \emptyset$.}
        Let $M = \{U\in T \mid \text{$U$ is disjoint from $Y$}\}$.
        Show that for all $x_0$ such that $x_0\in X\setminus Y$ we have $x_0\in\unions{M}$.
        \begin{subproof}
                        Fix $x_0$ such that $x_0\in X\setminus Y$.
                        $x_0\in X$ and $x_0\notin Y$ by \cref{setminus_elim_left,setminus_elim_right}.
                        Let $A = \{V\in T \mid \text{there exists $y\in Y$ such that there exists $U$ such that
                                $U$ is a neighborhood of $x_0$ in $X$ with respect to $T$ and
                                $V$ is a neighborhood of $y$ in $X$ with respect to $T$ and
                                $U$ is disjoint from $V$}\}$.
                        For all $y$ such that $y\in Y$ there exists $V\in A$ such that $y\in V$
                            by \cref{subseteq,hausdorff,neighborhood,open_in}.
                        For all $y$ such that $y\in Y$ we have $y\in\unions{A}$ by \cref{unions_iff}.
                        $A$ is a covering of $Y$ by \cref{subseteq,covering}.
                        For all $U$ such that $U\in A$ we have $U$ is open in $X$ with respect to $T$
                            by \cref{subseteq,open_in}.
                        Take $B$ such that
                            $B\subseteq A$ and $B$ is finite and $B$ is a covering of $Y$
                            by \cref{subspace_of,compact_subspace_open_covering,subseteq}.
                        Let $R = \{w\in B\times T \mid \text{there exists $V\in B$ such that there exists $U\in T$ such that
                                $w = (V,U)$ and
                                $U$ is a neighborhood of $x_0$ in $X$ with respect to $T$ and
                                $U$ is disjoint from $V$}\}$.
                        For all $w$ such that $w\in R$ there exist $V,U$ such that $w = (V,U)$.
                        $R$ is a relation by \cref{relation}.
                        Show that for all $V\in B$ there exists $U$ such that $V\mathrel{R} U$.
                        \begin{subproof}
                            Fix $V$ such that $V\in B$.
                            Take $y,U$ such that $y\in Y$ and
                                $U$ is a neighborhood of $x_0$ in $X$ with respect to $T$ and
                                $V$ is a neighborhood of $y$ in $X$ with respect to $T$ and
                                $U$ is disjoint from $V$
                                by \cref{subseteq}.
                            $U\in T$ by \cref{neighborhood,open_in}.
                            $(V,U)\in B\times T$ by \cref{times_tuple_intro}.
                            $(V,U)\in R$ by \cref{times_tuple_intro}.
                            $V\mathrel{R} U$.
                        \end{subproof}
                        Take $g$ such that $g$ is a function on $B$ and for all $V\in B$ we have $V\mathrel{R}\apply{g}{V}$ by \cref{choice_function}.
                        Let $G = \img{g}{B}$.
                        $G\subseteq T$ by \cref{img_iff,function_apply_iff,pair_eq_iff,subseteq}.
                        $G$ is finite by \cref{finite_image_function}.
                        Take $y$ such that $y\in Y$ by \cref{nonempty_inhabited}.
                        Take $V\in B$ such that $y\in V$ by \cref{covering,subseteq,unions_iff}.
                        $G$ is inhabited by \cref{function_apply_elim,subseteq,img_iff}.
                        Let $U_0 = \inters{G}$.
                        $U_0\in T$ by \cref{hausdorff,finite_intersections_open_in_topology}.
                        For all $U$ such that $U\in G$ there exists $V\in B$ such that $(V,U)\in g$
                            by \cref{img_iff}.
                        For all $U$ such that $U\in G$ we have $x_0\in U$
                            by \cref{function_apply_iff,pair_eq_iff,neighborhood}.
                        $x_0\in U_0$ by \cref{inters_intro}.
                        Show that $U_0$ is disjoint from $Y$.
                        \begin{subproof}
                            Suppose not.
                            Take $z$ such that $z\in U_0$ and $z\in Y$ by \cref{disjoint}.
                            Take $W\in B$ such that $z\in W$ by \cref{covering,subseteq,unions_iff}.
                            $z\in \apply{g}{W}$ by \cref{function_apply_elim,subseteq,img_iff,inters_subseteq_elem}.
                            Take $V_1,U_1$ such that $V_1\in B$ and $U_1\in T$ and
                                $(W,\apply{g}{W}) = (V_1,U_1)$ and
                                $U_1$ is a neighborhood of $x_0$ in $X$ with respect to $T$ and
                                $U_1$ is disjoint from $V_1$
                                by \cref{subseteq,function_apply_iff}.
                            $z\notin \apply{g}{W}$ by \cref{pair_eq_iff,disjoint}.
                            Contradiction.
                        \end{subproof}
                        $U_0\in M$.
                        $x_0\in\unions{M}$ by \cref{unions_iff}.
        \end{subproof}
        $X\setminus Y\subseteq \unions{M}$ by \cref{subseteq}.
        Show that for all $x$ such that $x\in\unions{M}$ we have $x\in X\setminus Y$.
        \begin{subproof}
            Follows by \cref{unions_iff,disjoint,subseteq,topology_on,powerset_elim,setminus_intro}.
        \end{subproof}
        $X\setminus Y = \unions{M}$ by \cref{subseteq,subseteq_antisymmetric}.
        $X\setminus Y$ is open in $X$ with respect to $T$ by \cref{subseteq,topology_on,open_in}.
        Show that $Y$ is closed in $X$ with respect to $T$.
        \begin{subproof}
            Follows by \cref{closed_in}.
        \end{subproof}
    \end{byCase}
\end{proof}

% from §26 Item 7: separation of a point and a compact set in a Hausdorff space
\begin{lemma}\label{compact_hausdorff_separation}
    Assume $T$ is a topology on $X$.
    Assume $X$ is Hausdorff with respect to $T$.
    Suppose $Y\subseteq X$.
    Let $T_Y = \subspaceTopology{T}{Y}$.
    Suppose $Y$ is compact with respect to $T_Y$.
    Suppose $x_0\in X\setminus Y$.
    Then there exist $U,V$ such that
        $U$ is open in $X$ with respect to $T$ and
        $V$ is open in $X$ with respect to $T$ and
        $x_0\in U$ and
        $Y\subseteq V$ and
        $U$ is disjoint from $V$.
\end{lemma}
\begin{proof}
    \begin{byCase}
    \caseOf{$Y = \emptyset$.}
        Let $U_0 = X$.
        Let $V_0 = \emptyset$.
        Show that there exist $U,V$ such that
            $U$ is open in $X$ with respect to $T$ and
            $V$ is open in $X$ with respect to $T$ and
            $x_0\in U$ and
            $Y\subseteq V$ and
            $U$ is disjoint from $V$.
        \begin{subproof}
            Follows by \cref{open_in,topology_on,setminus_elim_left,emptyset_subseteq,emptyset,disjoint}.
        \end{subproof}
    \caseOf{$Y \neq \emptyset$.}
        $x_0\in X$ and $x_0\notin Y$ by \cref{setminus_elim_left,setminus_elim_right}.
        Let $A = \{V\in T \mid \text{there exists $y\in Y$ such that there exists $U$ such that
                $U$ is a neighborhood of $x_0$ in $X$ with respect to $T$ and
                $V$ is a neighborhood of $y$ in $X$ with respect to $T$ and
                $U$ is disjoint from $V$}\}$.
        For all $y$ such that $y\in Y$ there exists $V\in A$ such that $y\in V$
            by \cref{subseteq,hausdorff,neighborhood,open_in}.
        For all $y$ such that $y\in Y$ we have $y\in\unions{A}$ by \cref{unions_iff}.
        $A$ is a covering of $Y$ by \cref{subseteq,covering}.
        For all $U$ such that $U\in A$ we have $U$ is open in $X$ with respect to $T$
            by \cref{subseteq,open_in}.
        Take $B$ such that
            $B\subseteq A$ and $B$ is finite and $B$ is a covering of $Y$
            by \cref{subspace_of,compact_subspace_open_covering,subseteq}.
        Let $R = \{w\in B\times T \mid \text{there exists $V\in B$ such that there exists $U\in T$ such that
                $w = (V,U)$ and
                $U$ is a neighborhood of $x_0$ in $X$ with respect to $T$ and
                $U$ is disjoint from $V$}\}$.
        For all $w$ such that $w\in R$ there exist $V,U$ such that $w = (V,U)$.
        $R$ is a relation by \cref{relation}.
        Show that for all $V\in B$ there exists $U$ such that $V\mathrel{R} U$.
        \begin{subproof}
            Fix $V$ such that $V\in B$.
            Take $y,U$ such that $y\in Y$ and
                $U$ is a neighborhood of $x_0$ in $X$ with respect to $T$ and
                $V$ is a neighborhood of $y$ in $X$ with respect to $T$ and
                $U$ is disjoint from $V$
                by \cref{subseteq}.
            $U\in T$ by \cref{neighborhood,open_in}.
            $(V,U)\in B\times T$ by \cref{times_tuple_intro}.
            $(V,U)\in R$ by \cref{times_tuple_intro}.
            $V\mathrel{R} U$.
        \end{subproof}
        Take $g$ such that $g$ is a function on $B$ and for all $V\in B$ we have $V\mathrel{R}\apply{g}{V}$ by \cref{choice_function}.
        Let $G = \img{g}{B}$.
        $G\subseteq T$ by \cref{img_iff,function_apply_iff,pair_eq_iff,subseteq}.
        $G$ is finite by \cref{finite_image_function}.
        Take $y$ such that $y\in Y$ by \cref{nonempty_inhabited}.
        Take $V\in B$ such that $y\in V$ by \cref{covering,subseteq,unions_iff}.
        $G$ is inhabited by \cref{function_apply_elim,subseteq,img_iff}.
        Let $U_0 = \inters{G}$.
        $U_0\in T$ by \cref{hausdorff,finite_intersections_open_in_topology}.
        For all $U$ such that $U\in G$ there exists $V\in B$ such that $(V,U)\in g$
            by \cref{img_iff}.
        For all $U$ such that $U\in G$ we have $x_0\in U$
            by \cref{function_apply_iff,pair_eq_iff,neighborhood}.
        $x_0\in U_0$ by \cref{inters_intro}.
        Let $V_0 = \unions{B}$.
        $V_0$ is open in $X$ with respect to $T$ by \cref{subseteq,topology_on,open_in}.
        $Y\subseteq V_0$ by \cref{covering}.
        Show that $U_0$ is disjoint from $V_0$.
        \begin{subproof}
            Suppose not.
            There exists $z$ such that $z\in U_0$ and $z\in V_0$.
            Take $W\in B$ such that $z\in W$ by \cref{unions_iff}.
            $z\in \apply{g}{W}$ by \cref{function_apply_elim,subseteq,img_iff,inters_subseteq_elem}.
            Take $V_1,U_1$ such that $V_1\in B$ and $U_1\in T$ and
                $(W,\apply{g}{W}) = (V_1,U_1)$ and
                $U_1$ is a neighborhood of $x_0$ in $X$ with respect to $T$ and
                $U_1$ is disjoint from $V_1$
                by \cref{subseteq,function_apply_iff}.
            $z\notin \apply{g}{W}$ by \cref{pair_eq_iff,disjoint}.
            Contradiction.
        \end{subproof}
        Show that there exist $U,V$ such that
            $U$ is open in $X$ with respect to $T$ and
            $V$ is open in $X$ with respect to $T$ and
            $x_0\in U$ and
            $Y\subseteq V$ and
            $U$ is disjoint from $V$.
        \begin{subproof}
            Trivial.
        \end{subproof}
    \end{byCase}
\end{proof}

% from §26 Item 8: image of a compact space under a continuous map is compact
\begin{theorem}\label{continuous_image_compact}
    Assume $f$ is continuous from $X$ to $Y$ with respect to $T$ and $T_1$.
    Assume $X$ is compact with respect to $T$.
    Let $A = \img{f}{X}$.
    Let $T_A = \subspaceTopology{T_1}{A}$.
    Then $A$ is compact with respect to $T_A$.
\end{theorem}
\begin{proof}
    $A$ is a subspace of $Y$ with respect to $T_1$
        by \cref{continuous,funs,rels_ran_subseteq,img_subseteq_ran,subseteq_transitive,subspace_of}.
    It suffices to show that for all $C$ such that
        $C$ is a covering of $A$ and
        for all $U$ such that $U\in C$ we have $U$ is open in $Y$ with respect to $T_1$
        there exists $B$ such that
            $B\subseteq C$ and
            $B$ is finite and
            $B$ is a covering of $A$
        by \cref{compact_subspace_open_covering}.
    Fix $C$ such that
        $C$ is a covering of $A$ and
        for all $U$ such that $U\in C$ we have $U$ is open in $Y$ with respect to $T_1$.
    Let $D = \{W \mid \exists U\in C. W = \preimg{f}{U}\}$.
    Show that for all $x$ such that $x\in X$ there exists $U\in C$ such that $f(x)\in U$.
    \begin{subproof}
        Follows by \cref{continuous,funs_is_function,funs,function_apply_elim,img_iff,covering,subseteq,unions_iff}.
    \end{subproof}
    Show that for all $x$ such that $x\in X$ we have $x\in\unions{D}$.
    \begin{subproof}
        Follows by \cref{continuous,funs_is_function,funs,function_apply_elim,preimg_iff,unions_iff}.
    \end{subproof}
    $D$ is a covering of $X$ by \cref{subseteq,covering}.
    Show that for all $W$ such that $W\in D$ we have $W$ is open in $X$ with respect to $T$.
    \begin{subproof}
        Fix $W$ such that $W\in D$.
        Take $U\in C$ such that $W = \preimg{f}{U}$.
        We have $U$ is open in $Y$ with respect to $T_1$.
        $\preimg{f}{U}$ is open in $X$ with respect to $T$ by \cref{continuous}.
        $W$ is open in $X$ with respect to $T$.
    \end{subproof}
    Show that $D$ is an open covering of $X$ with respect to $T$.
    \begin{subproof}
        Follows by \cref{open_covering}.
    \end{subproof}
    Take $E$ such that
        $E\subseteq D$ and $E$ is finite and $E$ is a covering of $X$
        by \cref{compact}.
    Let $R = \{w\in E\times C \mid \exists W\in E.\exists U\in C. w = (W,U) \land W = \preimg{f}{U}\}$.
    For all $w$ such that $w\in R$ there exist $W,U$ such that $w = (W,U)$.
    $R$ is a relation by \cref{relation}.
    Show that for all $W\in E$ there exists $U\in C$ such that $W = \preimg{f}{U}$.
    \begin{subproof}
        Follows by \cref{subseteq}.
    \end{subproof}
    For all $W\in E$ there exists $U$ such that $W\mathrel{R} U$.
    Take $h$ such that $h$ is a function on $E$ and for all $W\in E$ we have $W\mathrel{R}\apply{h}{W}$ by \cref{choice_function}.
    Let $B = \img{h}{E}$.
    $B\subseteq C$ by \cref{img_iff,function_apply_iff,pair_eq_iff,subseteq}.
    $B$ is finite by \cref{finite_image_function}.
    Show that for all $a$ such that $a\in A$ we have $a\in\unions{B}$.
    \begin{subproof}
        Fix $a$ such that $a\in A$.
        Take $x$ such that $x\in X$ and $(x,a)\in f$ by \cref{img_iff}.
        We have $a = \apply{f}{x}$ by \cref{continuous,funs_is_function,function_apply_intro}.
        Take $W\in E$ such that $x\in W$ by \cref{covering,subseteq,unions_iff}.
        Take $W_1,U$ such that $W_1\in E$ and $U\in C$ and
            $(W,\apply{h}{W}) = (W_1,U)$ and
            $W_1 = \preimg{f}{U}$
            by \cref{subseteq,function_apply_iff}.
        $\apply{h}{W}\in C$ and $W = \preimg{f}{\apply{h}{W}}$ by \cref{pair_eq_iff}.
        $x\in \preimg{f}{\apply{h}{W}}$.
        Take $b$ such that $b\in \apply{h}{W}$ and $(x,b)\in f$ by \cref{preimg_iff}.
        We have $b = \apply{f}{x}$ by \cref{continuous,funs_is_function,function_apply_intro}.
        $a\in \apply{h}{W}$.
        $\apply{h}{W}\in B$ by \cref{function_apply_elim,subseteq,img_iff}.
        $a\in\unions{B}$ by \cref{unions_iff}.
    \end{subproof}
    $B$ is a covering of $A$ by \cref{subseteq,covering}.
\end{proof}

% from §26 Item 9: bijective continuous map from compact to Hausdorff is a homeomorphism
\begin{theorem}\label{compact_hausdorff_bijection_homeomorphism}
    Assume $f$ is a bijection from $X$ to $Y$.
    Assume $f$ is continuous from $X$ to $Y$ with respect to $T$ and $T_1$.
    Assume $X$ is compact with respect to $T$.
    Assume $Y$ is Hausdorff with respect to $T_1$.
    Then $f$ is a homeomorphism between $X$ and $Y$ with respect to $T$ and $T_1$.
\end{theorem}
\begin{proof}
    Show that $T$ is a topology on $X$ and $T_1$ is a topology on $Y$.
    \begin{subproof}
        Follows by \cref{continuous}.
    \end{subproof}
    Show that for all $A$ such that $A$ is closed in $X$ with respect to $T$
        we have $\img{f}{A}$ is closed in $Y$ with respect to $T_1$.
    \begin{subproof}
        Fix $A$ such that $A$ is closed in $X$ with respect to $T$.
        We have $A\subseteq X$ and $A$ is a subspace of $X$ with respect to $T$
            by \cref{closed_in,subspace_of}.
        Let $T_A = \subspaceTopology{T}{A}$.
        $A$ is compact with respect to $T_A$ by \cref{closed_subspace_compact}.
        Let $g = \restrl{f}{A}$.
        $g$ is continuous from $A$ to $Y$ with respect to $T_A$ and $T_1$
            by \cref{continuous_restrict_domain}.
        Let $B = \img{g}{A}$.
        $B = \img{f}{A}$ by \cref{bijections_to_funs,subspace_of,subseteq,funs,restrl_img,inter_idempotent}.
        Let $T_B = \subspaceTopology{T_1}{B}$.
        $B$ is compact with respect to $T_B$ by \cref{continuous_image_compact}.
        $B$ is closed in $Y$ with respect to $T_1$ by \cref{bijections_to_funs,funs,rels_ran_subseteq,img_subseteq_ran,subseteq_transitive,compact_subspace_closed}.
        $\img{f}{A}$ is closed in $Y$ with respect to $T_1$.
    \end{subproof}
    $f$ is a closed map from $X$ to $Y$ with respect to $T$ and $T_1$ by \cref{bijections_to_funs,closed_map}.
    $\converse{f}\in\funs{Y}{X}$ by \cref{bijective_converse_are_funs}.
    Show that for all $C$ such that $C$ is closed in $X$ with respect to $T$
        we have $\preimg{\converse{f}}{C}$ is closed in $Y$ with respect to $T_1$.
    \begin{subproof}
        Follows by \cref{closed_map,bijective_converse_are_funs,funs_is_relation,bijections_to_funs,preim_eq_img_of_converse,converse_converse_eq}.
    \end{subproof}
    $\converse{f}$ is continuous from $Y$ to $X$ with respect to $T_1$ and $T$
        by \cref{preimg_closed_implies_continuous}.
    $f$ is a homeomorphism between $X$ and $Y$ with respect to $T$ and $T_1$
        by \cref{homeomorphism}.
\end{proof}

% from §26 helper lemma 1: slice is compact
\begin{lemma}\label{slice_compact}
    Assume $T$ is a topology on $X$.
    Assume $T_1$ is a topology on $Y$.
    Assume $Y$ is compact with respect to $T_1$.
    Assume $x_0\in X$.
    Let $T_2 = \productTopology{T}{T_1}{X}{Y}$.
    Let $S = \{x_0\}\times Y$.
    Let $T_S = \subspaceTopology{T_2}{S}$.
    Then $S$ is compact with respect to $T_S$.
\end{lemma}
\begin{proof}
    Let $f_1 = \{(y,x_0) \mid y\in Y\}$.
    Show that $f_1\in\funs{Y}{X}$ and for all $y$ such that $y\in Y$ we have $f_1(y) = x_0$.
    \begin{subproof}
        Follows by \cref{constant_map_function,funs_is_function,function_apply_intro}.
    \end{subproof}
    $f_1$ is continuous from $Y$ to $X$ with respect to $T_1$ and $T$
        by \cref{constant_continuous}.
    Let $f_2 = \identity{Y}$.
    $f_2$ is continuous from $Y$ to $Y$ with respect to $T_1$ and $T_1$
        by \cref{identity_continuous}.
    Show that $f_2\in\funs{Y}{Y}$ and for all $y$ such that $y\in Y$ we have $f_2(y) = y$.
    \begin{subproof}
        Follows by \cref{id_iff,id_elem_funs,funs_is_function,function_apply_intro}.
    \end{subproof}
    Let $k = \{(y,(f_1(y),f_2(y))) \mid y\in Y\}$.
    $k$ is continuous from $Y$ to $X\times Y$ with respect to $T_1$ and $T_2$
        by \cref{continuous_map_into_product_backward}.
    Let $S_0 = \img{k}{Y}$.
    $S_0$ is compact with respect to $\subspaceTopology{T_2}{S_0}$ by \cref{continuous_image_compact}.
    Show that for all $p$ such that $p\in S_0$ we have $p\in S$.
    \begin{subproof}
        Follows by \cref{img_iff,pair_eq_iff,singleton_iff,times_tuple_intro}.
    \end{subproof}
    Show that for all $p$ such that $p\in S$ we have $p\in S_0$.
    \begin{subproof}
        Follows by \cref{times_elem_is_tuple,singleton_iff,img_iff}.
    \end{subproof}
    Show that $S_0 = S$.
    \begin{subproof}
        Follows by \cref{setext}.
    \end{subproof}
    $S$ is compact with respect to $T_S$.
\end{proof}

% from §26 helper lemma 2: tube lemma
\begin{lemma}\label{tube_lemma_helper}
    Assume $T$ is a topology on $X$.
    Assume $T_1$ is a topology on $Y$.
    Assume $Y$ is compact with respect to $T_1$.
    Let $T_2 = \productTopology{T}{T_1}{X}{Y}$.
    Assume $N$ is open in $X\times Y$ with respect to $T_2$.
    Assume $x_0\in X$.
    Suppose $\{x_0\}\times Y\subseteq N$.
    Then there exists $W$ such that
        $W$ is open in $X$ with respect to $T$ and
        $x_0\in W$ and
        $W\times Y\subseteq N$.
\end{lemma}
\begin{proof}
    Let $S = \{x_0\}\times Y$.
    Let $T_S = \subspaceTopology{T_2}{S}$.
    \begin{byCase}
    \caseOf{$Y = \emptyset$.}
        Let $W = X$.
        We have $W$ is open in $X$ with respect to $T$ and
            $x_0\in W$ and
            $W\times Y\subseteq N$
            by \cref{open_in,topology_on,subseteq_refl,times_empty_right,emptyset_subseteq}.
        \caseOf{$Y\neq\emptyset$.}
            Take $y_0$ such that $y_0\in Y$ by \cref{nonempty_inhabited}.
            We have $x_0\in\{x_0\}$ by \cref{singleton_iff}.
            $(x_0,y_0)\in S$ and $S$ is inhabited by \cref{times_tuple_intro}.
        $S$ is compact with respect to $T_S$ by \cref{slice_compact}.
        Let $B = \productBasis{T}{T_1}{X}{Y}$.
        $B$ is a basis on $X\times Y$ by \cref{product_basis_is_basis}.
        $T_2$ is a topology on $X\times Y$ by \cref{basis_topology_is_topology,product_topology}.
        For all $p$ such that $p\in S$ we have $p\in X\times Y$
            by \cref{times_elem_is_tuple,singleton_iff,times_tuple_intro}.
        $S$ is a subspace of $X\times Y$ with respect to $T_2$ by \cref{subseteq,subspace_of}.
        Let $R = \{w\in S\times B \mid \exists p\in S.\exists B_1\in B. w = (p,B_1) \land p\in B_1 \land B_1\subseteq N\}$.
        For all $w$ such that $w\in R$ there exist $p,B_1$ such that $w = (p,B_1)$.
        Show that $R$ is a relation and for all $p$ such that $p\in S$ we have $p\in N$.
        \begin{subproof}
            Follows by \cref{relation,subseteq}.
        \end{subproof}
        For all $p\in S$ there exists $B_1\in B$ such that $p\in B_1$ and $B_1\subseteq N$
            by \cref{open_in,product_topology,topology_generated_by_basis}.
        For all $p\in S$ there exists $B_1$ such that $p\mathrel{R} B_1$.
        Take $g$ such that $g$ is a function on $S$ and for all $p\in S$ we have $p\mathrel{R}\apply{g}{p}$
            by \cref{choice_function}.
        Let $C = \img{g}{S}$.
        Show that for all $p$ such that $p\in S$ we have $p\in\unions{C}$.
        \begin{subproof}
            Follows by \cref{function_apply_iff,pair_eq_iff,function_apply_elim,img_iff,unions_iff}.
        \end{subproof}
        $C$ is a covering of $S$ by \cref{subseteq,covering}.
        Show that for all $U$ such that $U\in C$ we have $U$ is open in $X\times Y$ with respect to $T_2$.
        \begin{subproof}
            Follows by \cref{img_iff,function_apply_iff,pair_eq_iff,basis_elem_open_in_generated,product_topology,open_in}.
        \end{subproof}
        Take $F$ such that
            $F\subseteq C$ and $F$ is finite and $F$ is a covering of $S$
            by \cref{compact_subspace_open_covering}.
        Let $R_1 = \{w\in F\times T \mid \exists E\in F.\exists U\in T.\exists V\in T_1. w = (E,U) \land E = U\times V\}$.
        For all $w$ such that $w\in R_1$ there exist $E,U$ such that $w = (E,U)$.
        $R_1$ is a relation by \cref{relation}.
        Show that for all $E$ such that $E\in F$ we have $E\in B$.
        \begin{subproof}
            Fix $E$ such that $E\in F$.
            We have $E\in C$ by \cref{subseteq}.
            Take $p$ such that $p\in S$ and $(p,E)\in g$ by \cref{img_iff}.
            $E = \apply{g}{p}$ by \cref{function_apply_iff}.
            $p\mathrel{R} E$.
            $(p,E)\in S\times B$ by \cref{subseteq}.
            $E\in B$ by \cref{times_tuple_elim}.
        \end{subproof}
        Show that for all $E\in F$ there exists $U\in T$ such that there exists $V\in T_1$ such that $E = U\times V$.
        \begin{subproof}
            Follows by \cref{product_basis}.
        \end{subproof}
        Show that for all $E\in F$ there exists $U$ such that $E\mathrel{R_1} U$.
        \begin{subproof}
            Trivial.
        \end{subproof}
        Take $h$ such that $h$ is a function on $F$ and for all $E\in F$ we have $E\mathrel{R_1}\apply{h}{E}$
            by \cref{choice_function}.
        Let $G = \img{h}{F}$.
        $G$ is finite by \cref{finite_image_function}.
        Take $p_0$ such that $p_0\in S$ by \cref{nonempty_inhabited}.
        Take $E\in F$ such that $p_0\in E$ by \cref{covering,subseteq,unions_iff}.
        $G$ is inhabited by \cref{function_apply_elim,img_iff}.
        $G\subseteq T$ by \cref{img_iff,function_apply_iff,pair_eq_iff,subseteq}.
        Let $W = \inters{G}$.
        $W$ is open in $X$ with respect to $T$ by \cref{finite_intersections_open_in_topology,open_in}.
        Show that for all $U$ such that $U\in G$ there exists $E\in F$ such that $(E,U)\in h$.
        \begin{subproof}
            Follows by \cref{img_iff}.
        \end{subproof}
        Show that for all $U$ such that $U\in G$ we have $x_0\in U$.
        \begin{subproof}
            Follows by \cref{function_apply_iff,subseteq,img_iff,pair_eq_iff,times_elem_is_tuple,singleton_iff,times_tuple_elim}.
        \end{subproof}
        $x_0\in W$ by \cref{inters_intro}.
        Show that for all $p$ such that $p\in W\times Y$ we have $p\in N$.
        \begin{subproof}
            Fix $p$ such that $p\in W\times Y$.
            Take $x,y$ such that $p = (x,y)$ and $x\in W$ and $y\in Y$
                by \cref{times_elem_is_tuple}.
            $(x_0,y)\in S$ by \cref{times_tuple_intro}.
            Take $E_0\in F$ such that $(x_0,y)\in E_0$ by \cref{covering,subseteq,unions_iff}.
            Take $E_1,U_1,V_1$ such that $E_1\in F$ and $U_1\in T$ and $V_1\in T_1$ and
                $(E_0,\apply{h}{E_0}) = (E_1,U_1)$ and $E_1 = U_1\times V_1$
                by \cref{function_apply_iff}.
            $E_0 = E_1$ and $\apply{h}{E_0} = U_1$ by \cref{pair_eq_iff}.
            We have $E_0 = U_1\times V_1$.
            $(x_0,y)\in U_1\times V_1$.
            $y\in V_1$ by \cref{times_tuple_elim}.
            We have $(E_0,\apply{h}{E_0})\in h$ and $U_1\in G$ by \cref{function_apply_elim,img_iff,pair_eq_iff}.
            $W\subseteq U_1$ by \cref{inters_subseteq_elem}.
            $p\in U_1\times V_1$ by \cref{subseteq,times_tuple_intro}.
            $p\in E_0$ and $E_0\in C$ by \cref{subseteq}.
            Take $p_2\in S$ such that $(p_2,E_0)\in g$ by \cref{img_iff}.
            Take $p_1,B_1$ such that $p_1\in S$ and $B_1\in B$ and
                $(p_2,E_0) = (p_1,B_1)$ and $p_1\in B_1$ and $B_1\subseteq N$
                by \cref{function_apply_iff}.
            $p\in N$ by \cref{pair_eq_iff,subseteq}.
        \end{subproof}
        $W\times Y\subseteq N$ by \cref{subseteq}.
        $W$ is open in $X$ with respect to $T$ and
            $x_0\in W$ and
            $W\times Y\subseteq N$.
    \end{byCase}
\end{proof}

% from §26 helper lemma 3: union of a finite family of finite sets is finite
\begin{lemma}\label{finite_union_of_finite_family}
    Suppose $F$ is finite.
    Suppose for all $A$ such that $A\in F$ we have $A$ is finite.
    Then $\unions{F}$ is finite.
\end{lemma}
\begin{proof}
    \begin{byCase}
    \caseOf{$F = \emptyset$.}
        $\unions{F} = \emptyset$ by \cref{unions_emptyset}.
        $\unions{F}$ is finite by \cref{emptyset_is_finite}.
    \caseOf{$F \neq \emptyset$.}
        Take $F_0$ such that $F_0\in F$ by \cref{nonempty_inhabited}.
        Take $k\in\naturalsPlus$ such that there exists a bijection from $\initialSegment{k}$ to $F$ by \cref{finite}.
        Take $f$ such that $f$ is a bijection from $\initialSegment{k}$ to $F$.
        Let $A = \{ n\in\naturalsPlus \mid \text{for all $G,g$ such that $g$ is a bijection from $\initialSegment{n}$ to $G$ and for all $B\in G$ we have $B$ is finite we have $\unions{G}$ is finite} \}$.
        Show that for all $G,g$ such that $g$ is a bijection from $\initialSegment{1}$ to $G$ and
            for all $B\in G$ we have $B$ is finite we have $\unions{G}$ is finite.
        \begin{subproof}
            Fix $G$.
            Fix $g$ such that $g$ is a bijection from $\initialSegment{1}$ to $G$ and for all $B\in G$ we have $B$ is finite.
            $g$ is a function by \cref{bijection_is_function}.
            $\dom{g} = \initialSegment{1}$ by \cref{bijections_dom}.
            $1\in\dom{g}$ by \cref{initial_segment_one,singleton_intro}.
            Let $a = g(1)$.
            $(1,a)\in g$ by \cref{function_apply_elim}.
            $g\in\surj{\initialSegment{1}}{G}$ by \cref{bijections}.
            $\ran{g} = G$ by \cref{ran_of_surj}.
            $a\in G$ by \cref{ran_iff}.
            $a$ is finite.
            Show that for all $z\in G$ we have $z\in\{a\}$.
            \begin{subproof}
                Fix $z\in G$.
                Take $x\in\initialSegment{1}$ such that $g(x)=z$ by \cref{bijections,surj}.
                $x = 1$ by \cref{initial_segment_one,singleton_iff}.
                $a = z$.
                Thus $z\in\{a\}$ by \cref{singleton_iff}.
            \end{subproof}
            $G\subseteq\{a\}$ by \cref{subseteq}.
            $\{a\}\subseteq G$ by \cref{singleton_subset_intro}.
            $G = \{a\}$ by \cref{subseteq_antisymmetric}.
            $G = \cons{a}{\emptyset}$.
            $\unions{G} = a\union\emptyset$ by \cref{unions_cons,unions_emptyset}.
            $\unions{G} = a$ by \cref{union_emptyset}.
            $\unions{G}$ is finite.
        \end{subproof}
        $1\in A$.
        Show that for all $n\in A$ we have $n + 1\in A$.
        \begin{subproof}
            Fix $n\in A$.
            For all $G,g$ such that $g$ is a bijection from $\initialSegment{n}$ to $G$ and
                for all $B\in G$ we have $B$ is finite we have $\unions{G}$ is finite.
            Show that for all $G,g$ such that $g$ is a bijection from $\initialSegment{n + 1}$ to $G$
                and for all $B\in G$ we have $B$ is finite we have $\unions{G}$ is finite.
            \begin{subproof}
                Fix $G$.
                Fix $g$ such that $g$ is a bijection from $\initialSegment{n + 1}$ to $G$ and for all $B\in G$ we have $B$ is finite.
                $g$ is a function by \cref{bijection_is_function}.
                $\dom{g} = \initialSegment{n + 1}$ by \cref{bijections_dom}.
                $n + 1\in\dom{g}$ by \cref{initial_segment_naturalsplus,real_naturals_closed_succ}.
                Let $a = g(n + 1)$.
                Let $G_1 = \img{g}{\initialSegment{n}}$.
                $\initialSegment{n}\subseteq \initialSegment{n + 1}$ by \cref{initial_segment_subset_succ}.
                $G_1\subseteq G$ by \cref{img_subseteq,bijections_dom,img_dom,bijections,ran_of_surj}.
                $g\in\surj{\initialSegment{n + 1}}{G}$ by \cref{bijections}.
                $\ran{g} = G$ by \cref{ran_of_surj}.
                $(n + 1,a)\in g$ by \cref{function_apply_elim}.
                $a\in G$ by \cref{function_apply_elim,ran_iff,bijections,ran_of_surj}.
                $a$ is finite.
                Show that for all $z\in G$ we have $z\in G_1\union\{a\}$.
                \begin{subproof}
                    Fix $z\in G$.
                    There exists $x\in\initialSegment{n + 1}$ such that $g(x)=z$ by \cref{bijections,surj}.
                    $x\in\initialSegment{n}$ or $x = n + 1$ by \cref{initial_segment_succ_split}.
                    \begin{byCase}
                    \caseOf{$x\in\initialSegment{n}$.}
                        $z\in G_1\union\{a\}$ by \cref{function_apply_elim,img_iff,union_intro_left}.
                    \caseOf{$x = n + 1$.}
                        $z\in G_1\union\{a\}$ by \cref{function_apply_elim,union_intro_right,singleton_intro}.
                    \end{byCase}
                \end{subproof}
                Thus $G = G_1\union\{a\}$
                    by \cref{subseteq,singleton_subset_intro,union_subsets_is_subset,subseteq_antisymmetric}.
                Show that for all $B$ such that $B\in G_1$ we have $B\in G$ and $B$ is finite.
                \begin{subproof}
                    Follows by \cref{subseteq}.
                \end{subproof}
                Let $h = \restrl{g}{\initialSegment{n}}$.
                $h$ is a function by \cref{restrl_of_function_is_function}.
                For all $x\in\initialSegment{n}$ we have $x\in\dom{g}$ by \cref{initial_segment_subset_succ,bijections_dom,subseteq}.
                $\dom{h} = \initialSegment{n}$ by \cref{subseteq,dom_of_restrl_eq_inter,inter_absorb_supseteq_left}.
                $h\in\funs{\initialSegment{n}}{G_1}$ by \cref{function_apply_elim,restrl_subseteq,elem_subseteq,img_iff,funs_intro}.
                For all $y\in G_1$ there exists $x\in\initialSegment{n}$ such that $h(x) = y$
                    by \cref{img_iff,function_apply_iff,restrl_of_function_apply}.
                $h\in\surj{\initialSegment{n}}{G_1}$ by \cref{surj}.
                For all $x,y$ such that $x\in\dom{h}$ and $y\in\dom{h}$ and $h(x) = h(y)$ we have $x\in\initialSegment{n}$ and $y\in\initialSegment{n}$.
                For all $x,y$ such that $x\in\dom{h}$ and $y\in\dom{h}$ and $h(x) = h(y)$
                    we have $h(x) = g(x)$ and $h(y) = g(y)$ by \cref{restrl_of_function_apply}.
                For all $x,y$ such that $x\in\dom{h}$ and $y\in\dom{h}$ and $h(x) = h(y)$ we have $g(x) = g(y)$.
                For all $x,y$ such that $x\in\dom{h}$ and $y\in\dom{h}$ and $h(x) = h(y)$
                    we have $x\in\dom{g}$ and $y\in\dom{g}$ by \cref{subseteq}.
                For all $x,y$ such that $x\in\dom{h}$ and $y\in\dom{h}$ and $h(x) = h(y)$ we have $x = y$
                    by \cref{bijections,injective_function}.
                $h$ is injective by \cref{injective_function}.
                $h$ is a bijection from $\initialSegment{n}$ to $G_1$ by \cref{bijections}.
                $\unions{G_1}$ is finite.
                $a$ is finite.
                $\unions{G} = a\union\unions{G_1}$ by \cref{union_comm,union_eq_cons,unions_cons}.
                $\unions{G}$ is finite by \cref{union_of_finite_is_finite}.
            \end{subproof}
        \end{subproof}
        $A\subseteq\naturalsPlus$ by \cref{subseteq}.
        $k\in A$ by \cref{real_naturals_induction}.
        For all $G,g$ such that $g$ is a bijection from $\initialSegment{k}$ to $G$ and for all $B\in G$ we have $B$ is finite
            we have $\unions{G}$ is finite.
        $\unions{F}$ is finite.
    \end{byCase}
\end{proof}

% from §26 Item 10: product of finitely many compact spaces is compact
\begin{theorem}\label{product_compact}
    Assume $T$ is a topology on $X$.
    Assume $T_1$ is a topology on $Y$.
    Assume $X$ is compact with respect to $T$.
    Assume $Y$ is compact with respect to $T_1$.
    Let $T_2 = \productTopology{T}{T_1}{X}{Y}$.
    Then $X\times Y$ is compact with respect to $T_2$.
\end{theorem}
\begin{proof}
    $T_2$ is a topology on $X\times Y$ by \cref{product_basis_is_basis,basis_topology_is_topology,product_topology}.
    It suffices to show that for all $A$ such that $A$ is an open covering of $X\times Y$ with respect to $T_2$
        there exists $B$ such that
            $B\subseteq A$ and
            $B$ is finite and
            $B$ is a covering of $X\times Y$
        by \cref{compact}.
    Fix $A$ such that $A$ is an open covering of $X\times Y$ with respect to $T_2$.
    Let $C = \{W\in T \mid \text{there exists $B\subseteq A$ such that $B$ is finite and $W\times Y\subseteq \unions{B}$}\}$.
    Show that for all $W$ such that $W\in C$ we have $W\in T$ and $W$ is open in $X$ with respect to $T$.
    \begin{subproof}
        Follows by \cref{open_in}.
    \end{subproof}
    Show that for all $x$ such that $x\in X$ we have $x\in\unions{C}$.
    \begin{subproof}
        Fix $x$ such that $x\in X$.
        Let $S = \{x\}\times Y$.
        Let $T_S = \subspaceTopology{T_2}{S}$.
        $S$ is compact with respect to $T_S$ by \cref{slice_compact}.
        For all $p$ such that $p\in S$ we have $p\in X\times Y$
            by \cref{times_elem_is_tuple,singleton_iff,times_tuple_intro}.
        $S$ is a subspace of $X\times Y$ with respect to $T_2$ by \cref{subseteq,subspace_of}.
        $A$ is a covering of $S$ by \cref{open_covering,covering,subseteq}.
        Show that for all $U$ such that $U\in A$ we have $U$ is open in $X\times Y$ with respect to $T_2$.
        \begin{subproof}
            Follows by \cref{open_covering}.
        \end{subproof}
        Take $B$ such that
            $B\subseteq A$ and $B$ is finite and $B$ is a covering of $S$
            by \cref{compact_subspace_open_covering}.
        Let $N = \unions{B}$.
        $N$ is open in $X\times Y$ with respect to $T_2$ by \cref{subseteq,open_covering,open_in,topology_on}.
        Take $W$ such that
            $W$ is open in $X$ with respect to $T$ and
            $x\in W$ and
            $W\times Y\subseteq N$ by \cref{covering,tube_lemma_helper}.
        $x\in\unions{C}$ by \cref{open_in,unions_iff}.
    \end{subproof}
    Show that $C$ is an open covering of $X$ with respect to $T$.
    \begin{subproof}
        Follows by \cref{subseteq,covering,open_covering}.
    \end{subproof}
    Take $F$ such that
        $F\subseteq C$ and $F$ is finite and $F$ is a covering of $X$
        by \cref{compact}.
    Let $R = \{w\in F\times\pow{A} \mid \text{there exists $W\in F$ such that there exists $B\subseteq A$ such that $w = (W,B)$ and $B$ is finite and $W\times Y\subseteq \unions{B}$}\}$.
    For all $w$ such that $w\in R$ there exist $W,B$ such that $w = (W,B)$.
    $R$ is a relation by \cref{relation}.
    Show that for all $W\in F$ there exists $B\subseteq A$ such that $B$ is finite and $W\times Y\subseteq \unions{B}$.
    \begin{subproof}
        Follows by \cref{subseteq}.
    \end{subproof}
    Show that for all $W\in F$ there exists $B$ such that $(W,B)\in R$.
    \begin{subproof}
        Fix $W$ such that $W\in F$.
        Take $B\subseteq A$ such that $B$ is finite and $W\times Y\subseteq \unions{B}$.
        $B\in\pow{A}$ by \cref{powerset_intro}.
        $(W,B)\in F\times\pow{A}$ by \cref{times_tuple_intro}.
        $(W,B)\in R$ by \cref{times_tuple_intro,powerset_intro}.
    \end{subproof}
    For all $W\in F$ there exists $B$ such that $W\mathrel{R} B$.
    Take $h$ such that $h$ is a function on $F$ and for all $W\in F$ we have $W\mathrel{R}\apply{h}{W}$
        by \cref{choice_function}.
    Let $H = \img{h}{F}$.
    Show that for all $B$ such that $B\in H$ we have $B$ is finite.
    \begin{subproof}
        Follows by \cref{img_iff,function_apply_iff,pair_eq_iff}.
    \end{subproof}
    Let $B = \unions{H}$.
    $B$ is finite by \cref{finite_image_function,finite_union_of_finite_family}.
    $B\subseteq A$ by \cref{img_iff,function_apply_iff,pair_eq_iff,subseteq,powerset_intro,unions_subseteq_of_powerset_is_subseteq}.
    Show that for all $p$ such that $p\in X\times Y$ we have $p\in\unions{B}$.
    \begin{subproof}
        Fix $p$ such that $p\in X\times Y$.
        Take $x,y$ such that $p = (x,y)$ and $x\in X$ and $y\in Y$
            by \cref{times_elem_is_tuple}.
        Take $W\in F$ such that $x\in W$ by \cref{covering,subseteq,unions_iff}.
        $p\in W\times Y$ by \cref{times_tuple_intro}.
        Take $W_1,B_1$ such that $W_1\in F$ and $B_1\subseteq A$ and
            $(W,\apply{h}{W}) = (W_1,B_1)$ and $B_1$ is finite and $W_1\times Y\subseteq \unions{B_1}$
            by \cref{function_apply_iff}.
        Take $U\in\apply{h}{W}$ such that $p\in U$ by \cref{pair_eq_iff,subseteq,unions_iff}.
        $p\in\unions{B}$ by \cref{function_apply_elim,img_iff,unions_iff}.
    \end{subproof}
    $B$ is a covering of $X\times Y$ by \cref{subseteq,covering}.
\end{proof}

% from §26 Item 11: tube lemma
\begin{lemma}\label{tube_lemma}
    Assume $T$ is a topology on $X$.
    Assume $T_1$ is a topology on $Y$.
    Assume $Y$ is compact with respect to $T_1$.
    Let $T_2 = \productTopology{T}{T_1}{X}{Y}$.
    Assume $N$ is open in $X\times Y$ with respect to $T_2$.
    Assume $x_0\in X$.
    Suppose $\{x_0\}\times Y\subseteq N$.
    Then there exists $W$ such that
        $W$ is open in $X$ with respect to $T$ and
        $x_0\in W$ and
        $W\times Y\subseteq N$.
\end{lemma}
\begin{proof}
    Follows by \cref{tube_lemma_helper}.
\end{proof}
% from §26 Item 12: finite intersection property
\begin{definition}\label{finite_intersection_property}
    $C$ has the finite intersection property iff
        $C$ is inhabited and
        for all $F$ such that $F\subseteq C$ and $F$ is finite and inhabited
        we have $\inters{F}\neq\emptyset$.
\end{definition}
% from §26 Item 13: compactness via finite intersection property
\begin{theorem}\label{compact_iff_finite_intersection_property}
    Assume $T$ is a topology on $X$.
    Then $X$ is compact with respect to $T$ iff
        for all $C$ if
            $C$ has the finite intersection property and
            for all $A$ such that $A\in C$ we have $A$ is closed in $X$ with respect to $T$
        then $\inters{C}\neq\emptyset$.
\end{theorem}
\begin{proof}
    Show that if $X$ is compact with respect to $T$ then
        for all $C$ if
            $C$ has the finite intersection property and
            for all $A$ such that $A\in C$ we have $A$ is closed in $X$ with respect to $T$
        then $\inters{C}\neq\emptyset$.
    \begin{subproof}
        Assume $X$ is compact with respect to $T$.
        Fix $C$ such that $C$ has the finite intersection property and
            for all $A$ such that $A\in C$ we have $A$ is closed in $X$ with respect to $T$.
        Suppose not.
        $\inters{C} = \emptyset$ by \cref{emptyset}.
        Take $A_0$ such that $A_0\in C$ by \cref{finite_intersection_property}.
        Let $F_0 = \{A_0\}$.
        We have $F_0$ is finite and $F_0\subseteq C$
            by \cref{pair_is_finite,upair_intro_left,singleton_subset_intro,subseteq_finite_is_finite}.
        $F_0$ is inhabited.
        $\inters{F_0}\neq\emptyset$ by \cref{finite_intersection_property}.
        $A_0\neq\emptyset$ by \cref{inters_singleton,emptyset}.
        Let $A = \{U \mid \exists D\in C. U = X\setminus D\}$.
        Show that $A$ is an open covering of $X$ with respect to $T$.
        \begin{subproof}
            Show that for all $U$ such that $U\in A$ we have $U$ is open in $X$ with respect to $T$.
            \begin{subproof}
                Follows by \cref{closed_in}.
            \end{subproof}
            Show that for all $x$ such that $x\in X$ we have $x\in\unions{A}$.
            \begin{subproof}
                Trivial.
            \end{subproof}
            Show that $X\subseteq\unions{A}$ and $A$ is a covering of $X$.
            \begin{subproof}
                Follows by \cref{subseteq,covering}.
            \end{subproof}
            Show that $A$ is an open covering of $X$ with respect to $T$.
            \begin{subproof}
                Follows by \cref{open_covering}.
            \end{subproof}
        \end{subproof}
        Take $B$ such that $B\subseteq A$ and $B$ is finite and $B$ is a covering of $X$
            by \cref{compact}.
        $B$ is inhabited by \cref{closed_in,subseteq,nonempty_inhabited,covering,unions_iff}.
        Let $R = \{w\in B\times C \mid \exists U\in B.\exists D\in C. w = (U,D) \land U = X\setminus D\}$.
        For all $w$ such that $w\in R$ there exist $U,D$ such that $w = (U,D)$ and
            $R$ is a relation by \cref{times_elem_is_tuple,relation}.
        Show that for all $U\in B$ there exists $D\in C$ such that $U = X\setminus D$.
        \begin{subproof}
            Follows by \cref{subseteq}.
        \end{subproof}
        For all $U\in B$ there exists $D$ such that $U\mathrel{R} D$.
        Take $h$ such that $h$ is a function on $B$ and for all $U\in B$ we have $U\mathrel{R}\apply{h}{U}$
            by \cref{choice_function}.
        Let $F = \img{h}{B}$.
        $F$ is inhabited and $F\subseteq C$
            by \cref{nonempty_inhabited,function_apply_elim,img_iff,function_apply_iff,pair_eq_iff,subseteq}.
        Show that for all $x$ such that $x\in X\setminus\unions{B}$ we have $x\in\inters{F}$.
        \begin{subproof}
            Follows by \cref{img_iff,function_apply_iff,setminus_elim_left,setminus_elim_right,pair_eq_iff,setminus_intro,unions_iff,inters_intro}.
        \end{subproof}
        $X\setminus\unions{B}\subseteq\inters{F}$ by \cref{subseteq}.
        Show that for all $x$ such that $x\in\inters{F}$ we have $x\in X\setminus\unions{B}$.
        \begin{subproof}
            Fix $x$ such that $x\in\inters{F}$.
            We have $x\in X$ by \cref{nonempty_inhabited,inters_destr,subseteq,closed_in}.
            Suppose not.
            $x\in\unions{B}$.
            Take $U\in B$ such that $x\in U$ by \cref{unions_iff}.
            $U\mathrel{R}\apply{h}{U}$.
            $(U,\apply{h}{U})\in R$.
            Take $U_1,D_1$ such that $U_1\in B$ and $D_1\in C$ and
                $(U,\apply{h}{U}) = (U_1,D_1)$ and $U_1 = X\setminus D_1$
                by \cref{subseteq}.
            We have $U = U_1$ and $\apply{h}{U} = D_1$ by \cref{pair_eq_iff}.
            $x\in D_1$ by \cref{img_iff,function_apply_elim,inters_destr}.
            $x\in X\setminus D_1$ by \cref{pair_eq_iff}.
            $x\notin D_1$ by \cref{setminus_elim_right}.
            Contradiction.
        \end{subproof}
        $\inters{F}\subseteq X\setminus\unions{B}$ by \cref{subseteq}.
        Show that $X\setminus\unions{B} = \inters{F}$.
        \begin{subproof}
            Follows by \cref{subseteq_antisymmetric}.
        \end{subproof}
        For all $U$ such that $U\in B$ we have $U\subseteq X$.
        $X\setminus\unions{B} = \emptyset$
            by \cref{covering,powerset_intro,subseteq,unions_subseteq_of_powerset_is_subseteq,subseteq_antisymmetric,setminus_self}.
        $\inters{F} = \emptyset$ by \cref{subseteq_antisymmetric,subseteq_refl}.
        We have $\inters{F}\neq\emptyset$ by \cref{finite_image_function,finite_intersection_property}.
        Contradiction.
    \end{subproof}
    Show that if
        for all $C$ if
            $C$ has the finite intersection property and
            for all $A$ such that $A\in C$ we have $A$ is closed in $X$ with respect to $T$
        then $\inters{C}\neq\emptyset$
        then $X$ is compact with respect to $T$.
    \begin{subproof}
        Assume for all $C$ if
            $C$ has the finite intersection property and
            for all $A$ such that $A\in C$ we have $A$ is closed in $X$ with respect to $T$
        then $\inters{C}\neq\emptyset$.
        Show that for all $A$ such that $A$ is an open covering of $X$ with respect to $T$
            there exists $B$ such that
                $B\subseteq A$ and
                $B$ is finite and
                $B$ is a covering of $X$.
        \begin{subproof}
            Fix $A$ such that $A$ is an open covering of $X$ with respect to $T$.
            \begin{byCase}
            \caseOf{$X = \emptyset$.}
                Let $B = \emptyset$.
                We have $B\subseteq A$ and $B$ is finite and $B$ is a covering of $X$
                    by \cref{emptyset_subseteq,emptyset_is_finite,unions_emptyset,subseteq_refl,covering}.
                $B\subseteq A$ and $B$ is finite and $B$ is a covering of $X$.
            \caseOf{$X \neq \emptyset$.}
                Show that there exists $B$ such that
                    $B\subseteq A$ and $B$ is finite and $B$ is a covering of $X$.
                \begin{subproof}
                    Suppose not.
                    For all $B$ such that $B\subseteq A$ and $B$ is finite we have
                        $B$ is not a covering of $X$.
                    Let $C = \{D \mid \exists U\in A. D = X\setminus U\}$.
                    Show that for all $D$ such that $D\in C$ we have $D$ is closed in $X$ with respect to $T$.
                    \begin{subproof}
                        Follows by \cref{open_covering,open_in,topology_on,subseteq,pow_iff,double_relative_complement,setminus_subseteq,closed_in}.
                    \end{subproof}
                    Show that $C$ has the finite intersection property.
                    \begin{subproof}
                        Take $x_0,U_1$ such that $x_0\in X$ and $U_1\in A$ and $x_0\in U_1$
                            by \cref{open_covering,nonempty_inhabited,covering,subseteq,unions_iff}.
                        $C$ is inhabited.
                        Show that for all $F$ such that $F\subseteq C$ and $F$ is finite and inhabited
                            we have $\inters{F}\neq\emptyset$.
                        \begin{subproof}
                            Fix $F$ such that $F\subseteq C$ and $F$ is finite and inhabited.
                            Suppose not.
                            $\inters{F} = \emptyset$ by \cref{emptyset}.
                            Let $R_0 = \{w\in F\times A \mid \exists D\in F.\exists U\in A. w = (D,U) \land D = X\setminus U\}$.
                            For all $w$ such that $w\in R_0$ there exist $D,U$ such that $w = (D,U)$ and
                                $R_0$ is a relation by \cref{times_elem_is_tuple,relation}.
                            Show that for all $D\in F$ there exists $U\in A$ such that $D = X\setminus U$.
                            \begin{subproof}
                                Follows by \cref{subseteq}.
                            \end{subproof}
                            For all $D\in F$ there exists $U$ such that $D\mathrel{R_0} U$.
                            Take $h$ such that $h$ is a function on $F$ and for all $D\in F$ we have $D\mathrel{R_0}\apply{h}{D}$
                                by \cref{choice_function}.
                            Let $B = \img{h}{F}$.
                            $B$ is finite and $B\subseteq A$
                                by \cref{finite_image_function,img_iff,function_apply_iff,pair_eq_iff,subseteq}.
                            Show that for all $x$ such that $x\in X\setminus\unions{B}$ we have $x\in X$ and $x\notin\unions{B}$.
                            \begin{subproof}
                                Follows by \cref{setminus_elim_left,setminus_elim_right}.
                            \end{subproof}
                            Show that for all $x$ such that $x\in X\setminus\unions{B}$ we have $x\in\inters{F}$.
                            \begin{subproof}
                                Follows by \cref{pair_eq_iff,setminus_intro,function_apply_iff,function_apply_elim,img_iff,unions_iff,inters_intro}.
                            \end{subproof}
                            $X\setminus\unions{B}\subseteq\inters{F}$ by \cref{subseteq}.
                            Show that for all $x$ such that $x\in\inters{F}$ we have $x\in X\setminus\unions{B}$.
                            \begin{subproof}
                                Fix $x$ such that $x\in\inters{F}$.
                                We have $x\in X$ by \cref{nonempty_inhabited,inters_destr,subseteq,closed_in}.
                                Suppose not.
                                $x\in\unions{B}$.
                                Take $U\in B$ such that $x\in U$ by \cref{unions_iff}.
                                Take $D\in F$ such that $(D,U)\in h$ by \cref{img_iff}.
                                $D\mathrel{R_0} U$ by \cref{function_apply_iff}.
                                Take $D_1,U_2$ such that $D_1\in F$ and $U_2\in A$ and
                                    $(D,U) = (D_1,U_2)$ and $D_1 = X\setminus U_2$.
                                We have $D = D_1$ and $U = U_2$ by \cref{pair_eq_iff}.
                                $x\in X\setminus D$ by \cref{open_covering,open_in,topology_on,subseteq,pow_iff,double_relative_complement,pair_eq_iff}.
                                $x\notin D$ by \cref{setminus_elim_right}.
                                $x\in D$ by \cref{inters_destr}.
                                Contradiction.
                            \end{subproof}
                            $\inters{F}\subseteq X\setminus\unions{B}$ by \cref{subseteq}.
                            Show that $X\setminus\unions{B} = \inters{F}$.
                            \begin{subproof}
                                Follows by \cref{subseteq_antisymmetric}.
                            \end{subproof}
                            $X\setminus\unions{B} = \emptyset$.
                            $B$ is a covering of $X$ by \cref{setminus_eq_emptyset_iff_subseteq,covering}.
                            Contradiction.
                        \end{subproof}
                        Show that $C$ has the finite intersection property.
                        \begin{subproof}
                            Follows by \cref{finite_intersection_property}.
                        \end{subproof}
                    \end{subproof}
                    $\inters{C}\neq\emptyset$.
                    Show that for all $x$ such that $x\in X\setminus\unions{A}$ we have $x\in\inters{C}$.
                    \begin{subproof}
                        Trivial.
                    \end{subproof}
                    $X\setminus\unions{A}\subseteq\inters{C}$ by \cref{subseteq}.
                    Show that for all $x$ such that $x\in\inters{C}$ we have $x\in X\setminus\unions{A}$.
                    \begin{subproof}
                        Trivial.
                    \end{subproof}
                    $\inters{C}\subseteq X\setminus\unions{A}$ by \cref{subseteq}.
                    Show that $X\setminus\unions{A} = \inters{C}$.
                    \begin{subproof}
                        Follows by \cref{subseteq_antisymmetric}.
                    \end{subproof}
                    $X\setminus\unions{A} = \emptyset$
                        by \cref{covering,powerset_intro,open_covering,open_in,topology_on,subseteq,pow_iff,unions_subseteq_of_powerset_is_subseteq,subseteq_antisymmetric,setminus_self}.
                    $\inters{C} = \emptyset$ by \cref{subseteq_antisymmetric,subseteq_refl}.
                    Contradiction.
                \end{subproof}
            \end{byCase}
        \end{subproof}
        $X$ is compact with respect to $T$ by \cref{compact}.
    \end{subproof}
\end{proof}

\section{§27 Compact Subspaces of the Real Line}

% from §27 helper definition 1: immediate successor
\begin{definition}\label{immediate_successor}
    $y$ is an immediate successor of $x$ in $X$ with respect to $R$ iff
        $x\in X$ and $y\in X$ and $x\mathrel{R}y$ and
        $\intervalclosedrel{R}{X}{x}{y} = \{x,y\}$.
\end{definition}

% from §27 helper lemma 1: interval of an immediate successor
\begin{lemma}\label{immediate_successor_interval}
    Assume $R$ is asymmetric.
    Suppose $y$ is an immediate successor of $x$ in $X$ with respect to $R$.
    Let $I = \intervalclosedrel{R}{X}{x}{y}$.
    Then $I = \{x,y\}$.
\end{lemma}
\begin{proof}
    Follows by \cref{immediate_successor}.
\end{proof}

% from §27 helper lemma 2: non-immediate successor gives an intermediate point
\begin{lemma}\label{no_immediate_successor_between}
    Suppose $x\in X$.
    Suppose $y\in X$.
    Suppose $x\mathrel{R}y$.
    Suppose it is wrong that $y$ is an immediate successor of $x$ in $X$ with respect to $R$.
    Then there exists $z\in X$ such that $x\mathrel{R}z$ and $z\mathrel{R}y$.
\end{lemma}
\begin{proof}
    Let $I = \intervalclosedrel{R}{X}{x}{y}$.
    Suppose not.
    Show that $\{x,y\}\subseteq I$.
    \begin{subproof}
        Follows by \cref{upair_elim,intervalclosed_rel,subseteq}.
    \end{subproof}
    Show that for all $z$ such that $z\in I$ we have $z\in\{x,y\}$.
    \begin{subproof}
        Fix $z$ such that $z\in I$.
        $(x\mathrel{R}z \lor z = x)$ and $(z\mathrel{R}y \lor z = y)$
            by \cref{intervalclosed_rel}.
        \begin{byCase}
        \caseOf{$z = x$.}
            $z\in\{x,y\}$ by \cref{upair_intro_left}.
        \caseOf{$z = y$.}
            $z\in\{x,y\}$ by \cref{upair_intro_right}.
        \caseOf{$x\mathrel{R}z$ and $z\mathrel{R}y$.}
            Contradiction.
        \end{byCase}
    \end{subproof}
    Show that $y$ is an immediate successor of $x$ in $X$ with respect to $R$.
    \begin{subproof}
        Follows by \cref{subseteq,subseteq_antisymmetric,immediate_successor}.
    \end{subproof}
    Contradiction.
\end{proof}

% from §27 Item 1: closed interval is compact in an ordered set with the least upper bound property
\begin{theorem}\label{ordered_closed_interval_compact}
    Assume $R$ is a simple order relation on $X$.
    Assume $X$ has the least upper bound property with respect to $R$.
    Assume $T$ is the order topology on $X$ with respect to $R$.
    Suppose $a\in X$.
    Suppose $b\in X$.
    Suppose $a\mathrel{R}b$ or $a = b$.
    Let $Y = \intervalclosedrel{R}{X}{a}{b}$.
    Let $T_Y = \subspaceTopology{T}{Y}$.
    Then $Y$ is compact with respect to $T_Y$.
\end{theorem}
\begin{proof}
    Show that $R$ is transitive and $R$ is asymmetric and $R$ is connex on $X$.
    \begin{subproof}
        Follows by \cref{simple_order_relation}.
    \end{subproof}
    Show that for all $A$ such that $A$ is an open covering of $Y$ with respect to $T_Y$
        there exists $B$ such that $B\subseteq A$ and $B$ is finite and $B$ is a covering of $Y$.
    \begin{subproof}
        Fix $A$ such that $A$ is an open covering of $Y$ with respect to $T_Y$.
        \begin{byCase}
        \caseOf{$a = b$.}
            Take $U\in A$ such that $a\in U$
                by \cref{open_covering,covering,intervalclosed_rel,subseteq,unions_iff}.
            Let $B = \{U\}$.
            We have $B\subseteq A$ and $B$ is finite
                by \cref{singleton_subset_intro,pair_is_finite,upair_intro_left,subseteq_finite_is_finite}.
            For all $x$ such that $x\in Y$ we have $x\in\unions{B}$
                by \cref{intervalclosed_rel,asymmetric,singleton_intro,unions_iff}.
            $B$ is a covering of $Y$ by \cref{subseteq,covering}.
            $B\subseteq A$ and $B$ is finite and $B$ is a covering of $Y$.
        \caseOf{$a\mathrel{R}b$.}
            Show that $a\in Y$ and $b\in Y$.
            \begin{subproof}
                Follows by \cref{intervalclosed_rel}.
            \end{subproof}
            $a\neq b$ by \cref{asymmetric}.
            Show that $X$ has more than one element and $Y$ has more than one element.
            \begin{subproof}
                Follows by \cref{more_than_one_element}.
            \end{subproof}
            $Y$ is convex in $X$ with respect to $R$ by \cref{intervalclosedrel_convex}.
            Show that for all $y$ such that $y\in Y$ we have $y\in X$.
            \begin{subproof}
                Follows by \cref{intervalclosed_rel}.
            \end{subproof}
            $Y\subseteq X$ by \cref{subseteq}.
            $a$ is a smallest element of $Y$ with respect to $R$ by \cref{intervalclosed_rel,smallest_element}.
            $b$ is a largest element of $Y$ with respect to $R$ by \cref{intervalclosed_rel,largest_element}.
            $T = \basisTopology{\orderBasis{R}{X}}{X}$ by \cref{order_topology}.
            $T$ is a topology on $X$ by \cref{order_basis_is_basis,basis_topology_is_topology}.
            $T_Y$ is a topology on $Y$ by \cref{subspace_topology_is_topology,subspace_of}.

            Show that for all $x$ such that $x\in Y$ and $x\neq b$ there exists $y\in Y$ such that
                $x\mathrel{R}y$ and
                there exists $B$ such that $B\subseteq A$ and $B$ is finite and
                    $B$ is a covering of $\intervalclosedrel{R}{X}{x}{y}$.
            \begin{subproof}
                Fix $x$ such that $x\in Y$ and $x\neq b$.
                Take $U\in A$ such that $x\in U$ by \cref{open_covering,covering,subseteq,unions_iff}.
                Take $V\in T$ such that $U = Y\inter V$ by \cref{open_covering,open_in,subspace_topology}.
                $x\in V$ by \cref{inter_elim_right}.
                Take $B_1\in\orderBasis{R}{X}$ such that $x\in B_1$ and $B_1\subseteq V$
                    by \cref{topology_generated_by_basis}.
                Show that there exists $c\in Y$ such that $x\mathrel{R}c$ and
                    $\intervalclosedopenrel{R}{X}{x}{c} \subseteq V$.
                \begin{subproof}
                    $B_1\in\pow{X}$ and
                        $((\exists u, v\in X. u\mathrel{R}v \land B_1 = \intervalopenrel{R}{X}{u}{v}) \lor
                        (\exists v\in X. \exists a_0\in X. (\forall z\in X. a_0\mathrel{R}z \lor a_0 = z) \land
                            B_1 = \intervalclosedopenrel{R}{X}{a_0}{v}) \lor
                        (\exists u\in X. \exists b_0\in X. (\forall z\in X. z\mathrel{R}b_0 \lor z = b_0) \land
                            B_1 = \intervalopenclosedrel{R}{X}{u}{b_0}))$
                        by \cref{order_basis}.
                    \begin{byCase}
                    \caseOf{$\exists u, v\in X. u\mathrel{R}v \land B_1 = \intervalopenrel{R}{X}{u}{v}$.}
                        Take $u, v\in X$ such that $u\mathrel{R}v$ and $B_1 = \intervalopenrel{R}{X}{u}{v}$.
                        $u\mathrel{R}x$ and $x\mathrel{R}v$ by \cref{intervalopen_rel}.
                        \begin{byCase}
                        \caseOf{$v\in Y$.}
                            Let $c = v$.
                            $\intervalclosedopenrel{R}{X}{x}{c} \subseteq B_1$
                                by \cref{intervalclosedopen_rel,transitive,intervalopen_rel,subseteq}.
                        \caseOf{it is wrong that $v\in Y$.}
                            $a\mathrel{R}v$ by \cref{intervalclosed_rel,transitive}.
                            $b\mathrel{R}v$ by \cref{connex,intervalclosed_rel}.
                            $\intervalclosedopenrel{R}{X}{x}{b} \subseteq B_1$
                                by \cref{intervalclosedopen_rel,transitive,intervalopen_rel,subseteq}.
                            There exists $c\in Y$ such that $x\mathrel{R}c$ and
                                $\intervalclosedopenrel{R}{X}{x}{c} \subseteq B_1$ by \cref{intervalclosed_rel}.
                        \end{byCase}
                    \caseOf{$\exists v\in X. \exists a_0\in X. (\forall z\in X. a_0\mathrel{R}z \lor a_0 = z) \land
                        B_1 = \intervalclosedopenrel{R}{X}{a_0}{v}$.}
                        Take $a_0, v\in X$ such that
                            $(\forall z\in X. a_0\mathrel{R}z \lor a_0 = z)$ and
                            $B_1 = \intervalclosedopenrel{R}{X}{a_0}{v}$.
                        $x\in X$ and $(a_0\mathrel{R}x \lor x = a_0)$ and $x\mathrel{R}v$
                            by \cref{intervalclosedopen_rel}.
                        \begin{byCase}
                        \caseOf{$v\in Y$.}
                            Let $c = v$.
                            $\intervalclosedopenrel{R}{X}{x}{c} \subseteq B_1$
                                by \cref{intervalclosedopen_rel,smallest_element,intervalclosed_rel,subseteq}.
                        \caseOf{it is wrong that $v\in Y$.}
                            $a\mathrel{R}v$ by \cref{intervalclosed_rel,transitive}.
                            $b\mathrel{R}v$ by \cref{connex,intervalclosed_rel}.
                            $\intervalclosedopenrel{R}{X}{x}{b} \subseteq B_1$
                                by \cref{intervalclosedopen_rel,smallest_element,transitive,subseteq}.
                            There exists $c\in Y$ such that $x\mathrel{R}c$ and
                                $\intervalclosedopenrel{R}{X}{x}{c} \subseteq B_1$ by \cref{intervalclosed_rel}.
                        \end{byCase}
                    \caseOf{$\exists u\in X. \exists b_0\in X. (\forall z\in X. z\mathrel{R}b_0 \lor z = b_0) \land
                        B_1 = \intervalopenclosedrel{R}{X}{u}{b_0}$.}
                        Take $u, b_0\in X$ such that
                            $(\forall z\in X. z\mathrel{R}b_0 \lor z = b_0)$ and
                            $B_1 = \intervalopenclosedrel{R}{X}{u}{b_0}$.
                        $u\mathrel{R}x$ and $(x\mathrel{R}b_0 \lor x = b_0)$ by \cref{intervalopenclosed_rel}.
                        $\intervalclosedopenrel{R}{X}{x}{b} \subseteq B_1$
                            by \cref{intervalclosedopen_rel,transitive,intervalopenclosed_rel,subseteq}.
                        There exists $c\in Y$ such that $x\mathrel{R}c$ and
                            $\intervalclosedopenrel{R}{X}{x}{c} \subseteq B_1$ by \cref{intervalclosed_rel}.
                    \end{byCase}
                \end{subproof}
                Take $c\in Y$ such that $x\mathrel{R}c$ and $\intervalclosedopenrel{R}{X}{x}{c} \subseteq V$.
                \begin{byCase}
                \caseOf{there exists $y\in X$ such that $y$ is an immediate successor of $x$ in $X$ with respect to $R$.}
                    Take $y\in X$ such that $y$ is an immediate successor of $x$ in $X$ with respect to $R$.
                    Show that $y\in Y$.
                    \begin{subproof}
                        Show that $y\mathrel{R}b$ or $y = b$.
                        \begin{subproof}
                            $y\mathrel{R}b$ or $y = b$ or $b\mathrel{R}y$ by \cref{connex}.
                            \begin{byCase}
                            \caseOf{$b\mathrel{R}y$.}
                                $b\in\intervalclosedrel{R}{X}{x}{y}$ by \cref{intervalclosed_rel}.
                                $b\in\{x,y\}$ by \cref{immediate_successor_interval}.
                                $b = x$ or $b = y$ by \cref{upair_elim}.
                                $y = b$.
                            \caseOf{$y\mathrel{R}b$ or $y = b$.}
                            \end{byCase}
                        \end{subproof}
                        $y\in Y$ by \cref{immediate_successor,intervalclosed_rel,transitive}.
                    \end{subproof}
                    Take $U_1\in A$ such that $y\in U_1$
                        by \cref{open_covering,covering,subseteq,unions_iff}.
                    Let $B = \{U,U_1\}$.
                    $B\subseteq A$ and $B$ is finite by \cref{upair_elim,subseteq,pair_is_finite}.
                    For all $z$ such that $z\in\intervalclosedrel{R}{X}{x}{y}$ we have $z\in\unions{B}$
                        by \cref{immediate_successor_interval,upair_elim,upair_intro_left,upair_intro_right,unions_iff}.
                    $B$ is a covering of $\intervalclosedrel{R}{X}{x}{y}$ by \cref{subseteq,covering}.
                \caseOf{it is wrong that there exists $y\in X$ such that $y$ is an immediate successor of $x$ in $X$ with respect to $R$.}
                    It is wrong that $c$ is an immediate successor of $x$ in $X$ with respect to $R$
                        by \cref{subseteq}.
                    Take $y\in X$ such that $x\mathrel{R}y$ and $y\mathrel{R}c$ by \cref{no_immediate_successor_between}.
                    $y\in Y$ by \cref{convex_in}.
                    Let $B = \{U\}$.
                    $B\subseteq A$ and $B$ is finite
                        by \cref{singleton_subset_intro,pair_is_finite,upair_intro_left,subseteq_finite_is_finite}.
                    For all $z$ such that $z\in\intervalclosedrel{R}{X}{x}{y}$ we have
                        $z\in X$ and $(x\mathrel{R}z \lor z = x)$ and $(z\mathrel{R}y \lor z = y)$
                        by \cref{intervalclosed_rel}.
                    For all $z$ such that $z\in\intervalclosedrel{R}{X}{x}{y}$ we have $z\mathrel{R}c$
                        by \hyperref[transitive]{transitivity}.
                    For all $z$ such that $z\in\intervalclosedrel{R}{X}{x}{y}$ we have $z\in\unions{B}$
                        by \cref{intervalclosedopen_rel,subseteq,convex_in,inter_intro,singleton_intro,unions_iff}.
                    $B$ is a covering of $\intervalclosedrel{R}{X}{x}{y}$ by \cref{subseteq,covering}.
                \end{byCase}
            \end{subproof}

            Let $C = \{y\in Y \mid \text{$a\mathrel{R}y$ and there exists $B\subseteq A$ such that
                $B$ is finite and $B$ is a covering of $\intervalclosedrel{R}{X}{a}{y}$}\}$.
            $a\in Y$.
            Take $y\in Y$ such that $a\mathrel{R}y$ and
                there exists $B$ such that $B\subseteq A$ and $B$ is finite and
                    $B$ is a covering of $\intervalclosedrel{R}{X}{a}{y}$.
            $y\in C$.
            $C$ is inhabited.
            We have $C\subseteq X$ by \cref{intervalclosed_rel,subseteq}.
            Show that for all $y\in C$ we have $y\mathrel{R}b$ or $y = b$.
            \begin{subproof}
                Follows by \cref{intervalclosed_rel}.
            \end{subproof}
            $b$ is an upper bound of $C$ in $X$ with respect to $R$ by \cref{upper_bound_rel}.
            Take $c$ such that $c$ is a least upper bound of $C$ in $X$ with respect to $R$
                by \cref{least_upper_bound_property}.
            Show that $a\mathrel{R}c$.
            \begin{subproof}
                Take $y_0\in C$ such that $a\mathrel{R}y_0$ by \cref{nonempty_inhabited}.
                Suppose not.
                \begin{byCase}
                \caseOf{$c = a$.}
                    $y_0\mathrel{R}c$ or $y_0 = c$ by \cref{least_upper_bound_rel,upper_bound_rel}.
                    $y_0\mathrel{R}c$ by \cref{least_upper_bound_rel,upper_bound_rel,asymmetric}.
                    $y_0\mathrel{R}a$.
                    Contradiction by \cref{asymmetric}.
                \caseOf{$c\neq a$.}
                    $c\mathrel{R}a$ or $a\mathrel{R}c$ by \cref{least_upper_bound_rel,upper_bound_rel,connex}.
                    $c\mathrel{R}a$ by \cref{least_upper_bound_rel,upper_bound_rel,connex}.
                    $c\mathrel{R}y$ by \hyperref[transitive]{transitivity}.
                    $y\mathrel{R}c$ or $y = c$ by \cref{least_upper_bound_rel,upper_bound_rel}.
                    Contradiction by \cref{asymmetric}.
                \end{byCase}
            \end{subproof}
            Show that $c\in X$ and $c\in Y$.
            \begin{subproof}
                Follows by \cref{least_upper_bound_rel,upper_bound_rel,intervalclosed_rel}.
            \end{subproof}

            Show that $c\in C$.
            \begin{subproof}
                Take $U\in A$ such that $c\in U$ by \cref{open_covering,covering,subseteq,unions_iff}.
                Take $V\in T$ such that $U = Y\inter V$ by \cref{open_covering,open_in,subspace_topology}.
                $c\in V$ by \cref{inter_elim_right}.
                Take $B_1\in\orderBasis{R}{X}$ such that $c\in B_1$ and $B_1\subseteq V$
                    by \cref{topology_generated_by_basis}.
                Show that there exists $d\in Y$ such that $d\mathrel{R}c$ and
                    $\intervalopenclosedrel{R}{X}{d}{c} \subseteq V$.
                \begin{subproof}
                    $B_1\in\pow{X}$ and
                        $((\exists u, v\in X. u\mathrel{R}v \land B_1 = \intervalopenrel{R}{X}{u}{v}) \lor
                        (\exists v\in X. \exists a_0\in X. (\forall z\in X. a_0\mathrel{R}z \lor a_0 = z) \land
                            B_1 = \intervalclosedopenrel{R}{X}{a_0}{v}) \lor
                        (\exists u\in X. \exists b_0\in X. (\forall z\in X. z\mathrel{R}b_0 \lor z = b_0) \land
                            B_1 = \intervalopenclosedrel{R}{X}{u}{b_0}))$
                        by \cref{order_basis}.
                    \begin{byCase}
                    \caseOf{$\exists u, v\in X. u\mathrel{R}v \land B_1 = \intervalopenrel{R}{X}{u}{v}$.}
                        Take $u, v\in X$ such that $u\mathrel{R}v$ and $B_1 = \intervalopenrel{R}{X}{u}{v}$.
                        $u\mathrel{R}c$ and $c\mathrel{R}v$ by \cref{intervalopen_rel}.
                        \begin{byCase}
                        \caseOf{$u\in Y$.}
                            Let $d = u$.
                            $\intervalopenclosedrel{R}{X}{d}{c} \subseteq B_1$
                                by \cref{intervalopenclosed_rel,transitive,intervalopen_rel,subseteq}.
                        \caseOf{it is wrong that $u\in Y$.}
                            $u\mathrel{R}a$ by \cref{connex,least_upper_bound_rel,transitive,intervalclosed_rel}.
                            Let $d = a$.
                            $\intervalopenclosedrel{R}{X}{d}{c} \subseteq B_1$
                                by \cref{intervalopenclosed_rel,transitive,intervalopen_rel,subseteq}.
                        \end{byCase}
                    \caseOf{$\exists v\in X. \exists a_0\in X. (\forall z\in X. a_0\mathrel{R}z \lor a_0 = z) \land
                        B_1 = \intervalclosedopenrel{R}{X}{a_0}{v}$.}
                        Take $a_0, v\in X$ such that
                            $(\forall z\in X. a_0\mathrel{R}z \lor a_0 = z)$ and
                            $B_1 = \intervalclosedopenrel{R}{X}{a_0}{v}$.
                        We have $c\in X$ and $(a_0\mathrel{R}c \lor c = a_0)$ and $c\mathrel{R}v$
                            by \cref{intervalclosedopen_rel}.
                        Let $d = a$.
                        $\intervalopenclosedrel{R}{X}{d}{c} \subseteq B_1$
                            by \cref{intervalopenclosed_rel,smallest_element,transitive,intervalclosedopen_rel,subseteq}.
                    \caseOf{$\exists u\in X. \exists b_0\in X. (\forall z\in X. z\mathrel{R}b_0 \lor z = b_0) \land
                        B_1 = \intervalopenclosedrel{R}{X}{u}{b_0}$.}
                        Take $u, b_0\in X$ such that
                            $(\forall z\in X. z\mathrel{R}b_0 \lor z = b_0)$ and
                            $B_1 = \intervalopenclosedrel{R}{X}{u}{b_0}$.
                        $u\mathrel{R}c$ and $(c\mathrel{R}b_0 \lor c = b_0)$ by \cref{intervalopenclosed_rel}.
                        \begin{byCase}
                        \caseOf{$u\in Y$.}
                            Let $d = u$.
                            $\intervalopenclosedrel{R}{X}{d}{c} \subseteq B_1$
                                by \cref{intervalopenclosed_rel,transitive,intervalopenclosed_rel,subseteq}.
                        \caseOf{it is wrong that $u\in Y$.}
                            $u\mathrel{R}a$ by \cref{connex,least_upper_bound_rel,transitive,intervalclosed_rel}.
                            Let $d = a$.
                            $\intervalopenclosedrel{R}{X}{d}{c} \subseteq B_1$
                                by \cref{intervalopenclosed_rel,transitive,intervalopenclosed_rel,subseteq}.
                        \end{byCase}
                    \end{byCase}
                \end{subproof}
                Take $d\in Y$ such that $d\mathrel{R}c$ and $\intervalopenclosedrel{R}{X}{d}{c} \subseteq V$.
                Suppose not.
                Show that there exists $z\in C$ such that $d\mathrel{R}z$.
                \begin{subproof}
                    Suppose not.
                    Show that for all $y\in C$ we have $y\mathrel{R}d$ or $y = d$.
                    \begin{subproof}
                        Follows by \cref{subseteq,connex}.
                    \end{subproof}
                    $c\mathrel{R}d$ or $c = d$ by \cref{subseteq,upper_bound_rel,least_upper_bound_rel}.
                    Contradiction by \cref{asymmetric}.
                \end{subproof}
                Take $z\in C$ such that $d\mathrel{R}z$.
                $z\mathrel{R}c$ by \cref{least_upper_bound_rel,upper_bound_rel}.
                Take $B_0$ such that $B_0\subseteq A$ and $B_0$ is finite and
                    $B_0$ is a covering of $\intervalclosedrel{R}{X}{a}{z}$.
                Let $B = B_0\union\{U\}$.
                $B$ is finite by \cref{pair_is_finite,upair_intro_left,singleton_subset_intro,subseteq_finite_is_finite,union_of_finite_is_finite}.
                Show that for all $w$ such that $w\in\intervalclosedrel{R}{X}{a}{c}$ we have $w\in\unions{B}$.
                \begin{subproof}
                    Fix $w$ such that $w\in\intervalclosedrel{R}{X}{a}{c}$.
                    We have $w\in X$ and $(a\mathrel{R}w \lor w = a)$ and $(w\mathrel{R}c \lor w = c)$
                        by \cref{intervalclosed_rel}.
                    $w\mathrel{R}z$ or $w = z$ or $z\mathrel{R}w$ by \cref{connex}.
                    \begin{byCase}
                    \caseOf{$w\mathrel{R}z$ or $w = z$.}
                        $w\in\unions{B}$ by \cref{subseteq,intervalclosed_rel,covering,unions_iff,union_intro_left}.
                    \caseOf{$z\mathrel{R}w$.}
                        $w\in V$ by \cref{transitive,intervalopenclosed_rel,subseteq}.
                        $w\in Y$ by \cref{convex_in}.
                        $w\in Y\inter V$ by \cref{inter_intro}.
                        $w\in U$.
                        $w\in\unions{B}$ by \cref{singleton_intro,union_intro_right,unions_iff}.
                    \end{byCase}
                \end{subproof}
                $B$ is a covering of $\intervalclosedrel{R}{X}{a}{c}$ by \cref{subseteq,covering}.
                We have $\{U\}\subseteq A$ by \cref{singleton_subset_intro}.
                $B\subseteq A$ by \cref{union_subsets_is_subset}.
                $B\subseteq A$ and $B$ is finite and
                    $B$ is a covering of $\intervalclosedrel{R}{X}{a}{c}$.
                $c\in C$.
                Contradiction.
            \end{subproof}

            Show that $c = b$.
            \begin{subproof}
                Suppose not.
                $c\mathrel{R}b$ by \cref{least_upper_bound_rel}.
                Take $y_1,B_1$ such that $y_1\in Y$ and $c\mathrel{R}y_1$ and $B_1\subseteq A$ and
                    $B_1$ is finite and $B_1$ is a covering of $\intervalclosedrel{R}{X}{c}{y_1}$.
                $a\mathrel{R}y_1$ by \hyperref[transitive]{transitivity}.
                Take $B_0$ such that $B_0\subseteq A$ and $B_0$ is finite and
                    $B_0$ is a covering of $\intervalclosedrel{R}{X}{a}{c}$.
                Let $B = B_0\union B_1$.
                $B\subseteq A$ by \cref{union_subsets_is_subset}.
                $B$ is finite by \cref{union_of_finite_is_finite}.
                For all $w$ such that $w\in\intervalclosedrel{R}{X}{a}{y_1}$ we have $w\in\unions{B}$
                    by \cref{intervalclosed_rel,covering,subseteq,unions_iff,union_intro_left,union_intro_right,connex}.
                $B$ is a covering of $\intervalclosedrel{R}{X}{a}{y_1}$ by \cref{subseteq,covering}.
                $y_1\in C$.
                $y_1\mathrel{R}c$ or $y_1 = c$ by \cref{least_upper_bound_rel,upper_bound_rel}.
                Contradiction by \cref{asymmetric}.
            \end{subproof}
            Take $B$ such that $B\subseteq A$ and $B$ is finite and
                $B$ is a covering of $\intervalclosedrel{R}{X}{a}{b}$.
        \end{byCase}
    \end{subproof}
    $Y$ is compact with respect to $T_Y$ by \cref{compact}.
\end{proof}
% from §27 Item 2: closed interval in reals is compact
\begin{corollary}\label{real_closed_interval_compact}
    Suppose $a\in\reals$.
    Suppose $b\in\reals$.
    Suppose $a < b$ or $a = b$.
    Let $R = \realOrderRel$.
    Let $T = \realOrderTopology$.
    Let $Y = \intervalclosedrel{R}{\reals}{a}{b}$.
    Let $T_Y = \subspaceTopology{T}{Y}$.
    Then $Y$ is compact with respect to $T_Y$.
\end{corollary}
\begin{proof}
    Let $X = \reals$.
    $R$ is a simple order relation on $X$ and $X$ has the least upper bound property with respect to $R$ and
        $X$ has more than one element by \cref{real_order_relation_simple,reals_linear_continuum,linear_continuum}.
    $Y$ is compact with respect to $T_Y$
        by \cref{real_order_topology,order_topology,real_order_relation_intro,ordered_closed_interval_compact}.
\end{proof}

% from §27 helper lemma 3: four-term real sum bound
\begin{lemma}\label{real_four_sum_bound}
    Suppose $r_1,r_2,r_3,r_4,M_0\in\reals$.
    Suppose $r_1 \leq 1$ and $r_4 \leq 1$ and $r_2 \leq M_0$ and $r_3 \leq M_0$.
    Then $(r_1 + r_2) + (r_3 + r_4) \leq (M_0 + 1) + (M_0 + 1)$.
\end{lemma}
\begin{proof}
    $1\in\reals$ and $r_2 + 1\in\reals$ and $r_3 + 1\in\reals$ and $M_0 + 1\in\reals$
        by \cref{real_naturals_subset_reals,real_one_in_naturals,subseteq,real_add_closed}.
    $r_1 + r_2\in\reals$ and $r_3 + r_4\in\reals$ and $r_4 + r_3\in\reals$ and
        $(r_1 + r_2) + (r_3 + r_4)\in\reals$ and $(M_0 + 1) + (r_3 + r_4)\in\reals$ and
        $(r_3 + r_4) + (M_0 + 1)\in\reals$ and $(M_0 + 1) + (M_0 + 1)\in\reals$
        by \cref{real_add_closed}.
    $r_1 + r_2 \leq 1 + r_2$ by \cref{real_leq_add}.
    $1 + r_2 = r_2 + 1$ by \cref{real_add_comm}.
    $r_2 + 1 \leq M_0 + 1$ by \cref{real_leq_add}.
    Hence $r_1 + r_2 \leq M_0 + 1$ by \cref{real_leq_trans}.
    $r_4 + r_3 \leq 1 + r_3$ by \cref{real_leq_add}.
    $r_4 + r_3 = r_3 + r_4$ and $1 + r_3 = r_3 + 1$ by \cref{real_add_comm}.
    $r_3 + 1 \leq M_0 + 1$ by \cref{real_leq_add}.
    Hence $r_3 + r_4 \leq M_0 + 1$ by \cref{real_leq_trans}.
    Thus $(r_1 + r_2) + (r_3 + r_4) \leq (M_0 + 1) + (r_3 + r_4)$
        by \cref{real_leq_add}.
    $(M_0 + 1) + (r_3 + r_4) = (r_3 + r_4) + (M_0 + 1)$ by \cref{real_add_comm}.
    $(r_3 + r_4) + (M_0 + 1) \leq (M_0 + 1) + (M_0 + 1)$ by \cref{real_leq_add}.
    Therefore $(r_1 + r_2) + (r_3 + r_4) \leq (M_0 + 1) + (M_0 + 1)$
        by \cref{real_leq_trans}.
\end{proof}

% from §27 helper lemma 4: compact standard closed interval
\begin{lemma}\label{real_intervalclosed_compact_standard}
    Suppose $a\in\reals$.
    Suppose $b\in\reals$.
    Suppose $a \leq b$.
    Let $I = \intervalclosed{a}{b}$.
    Then $I$ is compact with respect to $\subspaceTopology{\standardTopologyR}{I}$.
\end{lemma}
\begin{proof}
    Let $R = \realOrderRel$.
    Let $J = \intervalclosedrel{R}{\reals}{a}{b}$.
    Show that for all $x\in I$ we have $x\in J$.
    \begin{subproof}
        Follows by \cref{intervalclosed,intervalclosed_rel,real_order_relation_intro}.
    \end{subproof}
    Hence $I\subseteq J$ by \cref{subseteq}.
    Show that for all $x\in J$ we have $x\in I$.
    \begin{subproof}
        Follows by \cref{intervalclosed,intervalclosed_rel,real_order_relation_elim}.
    \end{subproof}
    $J\subseteq I$ by \cref{subseteq}.
    Thus $I = J$ by \cref{subseteq_antisymmetric}.
    Let $T = \realOrderTopology$.
    Let $T_I = \subspaceTopology{T}{I}$.
    $I$ is compact with respect to $T_I$ by \cref{real_closed_interval_compact}.
    Therefore $I$ is compact with respect to $\subspaceTopology{\standardTopologyR}{I}$
        by \cref{real_order_topology_eq_standard}.
\end{proof}

% from §27 helper lemma 5: indexed projection map is continuous
\begin{lemma}\label{projection_map_indexed_continuous}
    Assume $X$ is a function.
    Assume $T$ is a function.
    Assume $\dom{T} = \dom{X}$.
    Assume $\dom{X}$ is inhabited.
    Assume that for all $\alpha\in\dom{X}$ we have $\apply{T}{\alpha}$ is a topology on $\apply{X}{\alpha}$.
    Let $T_0 = \productTopologyIndexed{T}{X}$.
    Then for all $\alpha\in\dom{X}$ we have $\piIndex{X}{\alpha}$ is continuous from $\productFamily{X}$ to $\apply{X}{\alpha}$
        with respect to $T_0$ and $\apply{T}{\alpha}$.
\end{lemma}
\begin{proof}
    $T_0$ is a topology on $\productFamily{X}$
        by \cref{product_subbasis_is_subbasis,product_subbasis_finite_intersections_basis,basis_topology_is_topology,basis_for_topology}.
    Let $j = \identity{\productFamily{X}}$.
    $j\in\funs{\productFamily{X}}{\productFamily{X}}$ and
        $j$ is continuous from $\productFamily{X}$ to $\productFamily{X}$ with respect to $T_0$ and $T_0$
        by \cref{id_elem_funs,identity_continuous}.
    Therefore for all $\alpha\in\dom{X}$ we have $\piIndex{X}{\alpha}\circ j$ is continuous from $\productFamily{X}$
        to $\apply{X}{\alpha}$ with respect to $T_0$ and $\apply{T}{\alpha}$
        by \cref{continuous_map_into_indexed_product}.
    Fix $\alpha$ such that $\alpha\in\dom{X}$.
    Show that for all $x,z$ we have $x\mathrel{\piIndex{X}{\alpha}\circ j} z$ iff $x\mathrel{\piIndex{X}{\alpha}} z$.
    \begin{subproof}
        Follows by \cref{circ_iff,id_iff,projection_map_indexed,pair_eq_iff}.
    \end{subproof}
    Therefore $\piIndex{X}{\alpha}$ is continuous from $\productFamily{X}$ to $\apply{X}{\alpha}$
        with respect to $T_0$ and $\apply{T}{\alpha}$
        by \cref{continuous_map_into_indexed_product,circ_is_relation,projection_map_indexed_function,id_elem_funs,funs_is_relation,relext}.
\end{proof}

% from §27 helper definition 2: indexed product compact index
\begin{definition}\label{indexed_product_compact_index}
    A natural number $m$ is indexed-product-compact iff
        for all $Y_1,T_1$ such that $Y_1$ is a function and $T_1$ is a function and
            $\dom{Y_1} = \initialSegment{m}$ and $\dom{T_1} = \initialSegment{m}$,
            if for all $i\in\initialSegment{m}$ we have $\apply{T_1}{i}$ is a topology on $\apply{Y_1}{i}$ then
            if for all $i\in\initialSegment{m}$ we have $\apply{Y_1}{i}$ is compact with respect to $\apply{T_1}{i}$ then
            $\productFamily{Y_1}$ is compact with respect to $\productTopologyIndexed{T_1}{Y_1}$.
\end{definition}

% from §27 helper definition 3: finite indexed product compactness property
\begin{definition}\label{finite_indexed_product_compact_property}
    $\finiteIndexedProductCompactProp{m}$ iff
        for all $Y_1,T_1$ such that
            $Y_1$ is a function and $T_1$ is a function and
            $\dom{Y_1} = \initialSegment{m}$ and $\dom{T_1} = \initialSegment{m}$ and
            for all $i\in\dom{Y_1}$ we have $\apply{T_1}{i}$ is a topology on $\apply{Y_1}{i}$ and
                $\apply{Y_1}{i}$ is compact with respect to $\apply{T_1}{i}$,
        we have $\productFamily{Y_1}$ is compact with respect to $\productTopologyIndexed{T_1}{Y_1}$.
\end{definition}

% from §27 helper lemma 7: continuity preserved under equality
\begin{lemma}\label{continuous_eq_subst}
    Let $f,g$ be functions.
    Suppose $f = g$.
    Suppose $g$ is continuous from $X$ to $Y$ with respect to $T$ and $T_1$.
    Then $f$ is continuous from $X$ to $Y$ with respect to $T$ and $T_1$.
\end{lemma}
\begin{proof}
    Trivial.
\end{proof}

% from §27 helper lemma 8: continuity respects equal codomain/topology
\begin{lemma}\label{continuous_eq_codomain}
    Suppose $Y = Y_2$.
    Suppose $T_1 = T_2$.
    Suppose $f$ is continuous from $X$ to $Y$ with respect to $T$ and $T_1$.
    Then $f$ is continuous from $X$ to $Y_2$ with respect to $T$ and $T_2$.
\end{lemma}
\begin{proof}
    Trivial.
\end{proof}

\begin{lemma}\label{finite_indexed_product_compact_property_succ}
    Assume $k\in\naturalsPlus$.
    Assume $\finiteIndexedProductCompactProp{k}$.
    Then $\finiteIndexedProductCompactProp{k + 1}$.
\end{lemma}
\begin{proof}
    Show that for all $Y_1,T_1$ such that $Y_1$ is a function and $T_1$ is a function and
        $\dom{Y_1} = \initialSegment{k + 1}$ and $\dom{T_1} = \initialSegment{k + 1}$ and
        for all $i\in\dom{Y_1}$ we have $\apply{T_1}{i}$ is a topology on $\apply{Y_1}{i}$ and
            $\apply{Y_1}{i}$ is compact with respect to $\apply{T_1}{i}$,
        we have $\productFamily{Y_1}$ is compact with respect to $\productTopologyIndexed{T_1}{Y_1}$.
    \begin{subproof}
        Fix $Y_1$.
        Fix $T_1$.
        Assume $Y_1$ is a function.
        Assume $T_1$ is a function.
        Assume $\dom{Y_1} = \initialSegment{k + 1}$.
        Assume $\dom{T_1} = \initialSegment{k + 1}$.
        Assume that for all $i\in\dom{Y_1}$ we have $\apply{T_1}{i}$ is a topology on $\apply{Y_1}{i}$ and
            $\apply{Y_1}{i}$ is compact with respect to $\apply{T_1}{i}$.
        Let $J_1 = \initialSegment{k}$.
        Let $J = \initialSegment{k + 1}$.
        For all $i$ such that $i\in J_1$ we have $i\in J$
            by \cref{initial_segment_naturalsplus,real_naturals_subset_reals,subseteq,real_naturals_closed_succ,real_add_ident,real_one_in_naturals,real_zero_lt_one,real_lt_add,real_add_comm,real_lt_subst_left,real_lt_imp_leq,real_leq_trans}.
        Hence $J_1\subseteq J$ by \cref{subseteq}.
        $k\in\reals$ by \cref{real_naturals_subset_reals,subseteq}.
        $k + 1\in\naturalsPlus$ and $k + 1\in\reals$
            by \cref{real_naturals_closed_succ,real_naturals_subset_reals,subseteq}.
        Thus $k < k + 1$
            by \cref{real_zero_lt_one,real_add_ident,real_naturals_subset_reals,real_one_in_naturals,subseteq,real_lt_add,real_add_comm,real_lt_subst_left}.
        Show that $k \neq k + 1$.
        \begin{subproof}
            Suppose not.
            Thus $k = k + 1$.
            Hence $k < k$ by \cref{real_lt_subst_left}.
            Contradiction by \cref{real_lt_irreflexive}.
        \end{subproof}
        Let $Y_0 = \{(i,\apply{Y_1}{i}) \mid i\in J_1\}$.
        Let $T_0 = \{(i,\apply{T_1}{i}) \mid i\in J_1\}$.
        $\dom{Y_0} = J_1$ by \cref{dom_intro,dom_iff,pair_eq_iff,subseteq,subseteq_antisymmetric}.
        We have $\dom{T_0} = J_1$ by \cref{dom_intro,dom_iff,pair_eq_iff,subseteq,subseteq_antisymmetric}.
        $Y_0$ is a relation by \cref{relation}.
        Show that for all $a,b_1,b_2$ such that $(a,b_1)\in Y_0$ and $(a,b_2)\in Y_0$ we have $b_1 = b_2$.
        \begin{subproof}
            Follows by \cref{pair_eq_iff}.
        \end{subproof}
        Hence $Y_0$ is a function by \cref{rightunique}.
        $T_0$ is a relation by \cref{relation}.
        Show that for all $a,b_1,b_2$ such that $(a,b_1)\in T_0$ and $(a,b_2)\in T_0$ we have $b_1 = b_2$.
        \begin{subproof}
            Follows by \cref{pair_eq_iff}.
        \end{subproof}
        Thus $T_0$ is a function by \cref{rightunique}.
        For all $i$ such that $i\in J_1$ we have $\apply{T_0}{i}$ is a topology on $\apply{Y_0}{i}$
            by \cref{subseteq,function_apply_intro}.
        For all $i$ such that $i\in J_1$ we have $\apply{Y_0}{i}$ is compact with respect to $\apply{T_0}{i}$
            by \cref{subseteq,function_apply_intro}.
        Therefore $\productFamily{Y_0}$ is compact with respect to $\productTopologyIndexed{T_0}{Y_0}$
            by \cref{finite_indexed_product_compact_property}.
        Let $X_0 = \productFamily{Y_0}$.
        Let $T_X = \productTopologyIndexed{T_0}{Y_0}$.
        Let $X_1 = \apply{Y_1}{k + 1}$.
        Let $T_{1p} = \apply{T_1}{k + 1}$.
        $k + 1\in J$ by \cref{real_naturals_closed_succ,initial_segment_naturalsplus}.
        Show that $T_{1p}$ is a topology on $X_1$ and $X_1$ is compact with respect to $T_{1p}$.
        \begin{subproof}
            Follows by \cref{subseteq}.
        \end{subproof}
        We have $J_1$ is inhabited by \cref{real_one_leq_naturalsplus,real_one_in_naturals,initial_segment_naturalsplus}.
        Hence $T_X$ is a topology on $X_0$
            by \cref{product_subbasis_finite_intersections_basis,basis_for_topology,basis_topology_is_topology}.
        Let $T_P = \productTopology{T_X}{T_{1p}}{X_0}{X_1}$.
        Therefore $X_0\times X_1$ is compact with respect to $T_P$ by \cref{product_compact}.
        Let $g = \{(p,\fst{p}\union\{(k + 1,\snd{p})\}) \mid p\in X_0\times X_1\}$.
        $g$ is a relation by \cref{relation}.
        For all $p,u,v$ such that $(p,u)\in g$ and $(p,v)\in g$ we have $u = v$ by \cref{pair_eq_iff}.
        Hence $g$ is a function by \cref{rightunique}.
        We have $\dom{g} = X_0\times X_1$ by \cref{dom_intro,dom_iff,pair_eq_iff,subseteq,subseteq_antisymmetric}.
            Show that for all $p\in\dom{g}$ we have $g(p)\in\productFamily{Y_1}$.
            \begin{subproof}
                Fix $p$ such that $p\in\dom{g}$.
                Take $v$ such that $(p,v)\in g$ by \cref{dom_iff}.
                Hence $p\in X_0\times X_1$ by \cref{subseteq}.
                Let $u = \fst{p}\union\{(k + 1,\snd{p})\}$.
                $(p,u)\in g$.
                Hence $v = u$ by \cref{pair_eq_iff}.
                Let $x = \fst{p}$.
                    Let $y = \snd{p}$.
                    We have $x\in X_0$ and $y\in X_1$ by \cref{subseteq,times_proj_elim}.
                    Hence $x\in\funs{\dom{Y_0}}{\unions{\ran{Y_0}}}$ and $x$ is a function and $\dom{x} = \dom{Y_0}$
                        by \cref{subseteq,indexed_product,funs_elim}.
                    Thus $x\in\rels{\dom{Y_0}}{\unions{\ran{Y_0}}}$ and $x$ is a relation
                        by \cref{funs_subseteq_rels,subseteq,rels_is_relation}.
                    Show that $v\in\funs{\dom{Y_1}}{\unions{\ran{Y_1}}}$.
                    \begin{subproof}
                            For all $w$ such that $w\in v$ there exist $q_1,q_2$ such that $w = (q_1,q_2)$
                                by \cref{union_iff,singleton_iff,relation}.
                            Hence $v$ is a relation by \cref{relation}.
                            Show that for all $a,b_1,b_2$ such that $(a,b_1)\in v$ and $(a,b_2)\in v$ we have $b_1 = b_2$.
                            \begin{subproof}
                                Fix $a$.
                                Fix $b_1$.
                                Fix $b_2$ such that $(a,b_1)\in v$ and $(a,b_2)\in v$.
                                \begin{byCase}
                                \caseOf{$a = k + 1$.}
                                    Hence $(a,b_1)\in x$ or $(a,b_1) = (k + 1,y)$ by \cref{union_iff,singleton_iff}.
                                    Show that $(a,b_1)\notin x$.
                                    \begin{subproof}
                                        Suppose not.
                                        $(a,b_1)\in x$.
                                        Hence $a\in\dom{x}$ by \cref{dom_intro}.
                                        We have $\dom{x} = \dom{Y_0}$ by \cref{funs_elim}.
                                        Hence $a\in J_1$ by \cref{subseteq}.
                                        Therefore $a\in\naturalsPlus$ and $a \leq k$ by \cref{initial_segment_naturalsplus}.
                                        Hence $k + 1 \leq k$.
                                        Therefore $k + 1 < k$ or $k + 1 = k$ by \cref{real_leq_split}.
                                        Contradiction by \cref{real_lt_asym,real_lt_subst_left,real_lt_irreflexive}.
                                    \end{subproof}
                                    Hence $(a,b_1) = (k + 1,y)$.
                                    We have $(a,b_2)\in x$ or $(a,b_2) = (k + 1,y)$ by \cref{union_iff,singleton_iff}.
                                    Show that $(a,b_2)\notin x$.
                                    \begin{subproof}
                                        Suppose not.
                                        $(a,b_2)\in x$.
                                        Hence $a\in\dom{x}$ by \cref{dom_intro}.
                                        We have $\dom{x} = \dom{Y_0}$ by \cref{funs_elim}.
                                        Hence $a\in J_1$ by \cref{subseteq}.
                                        Therefore $a\in\naturalsPlus$ and $a \leq k$ by \cref{initial_segment_naturalsplus}.
                                        Hence $k + 1 \leq k$.
                                        Therefore $k + 1 < k$ or $k + 1 = k$ by \cref{real_leq_split}.
                                        Contradiction by \cref{real_lt_asym,real_lt_subst_left,real_lt_irreflexive}.
                                    \end{subproof}
                                    Hence $(a,b_2) = (k + 1,y)$.
                                    Hence $b_1 = y$ and $b_2 = y$ by \cref{pair_eq_iff}.
                                    Therefore $b_1 = b_2$.
                                \caseOf{$a \neq k + 1$.}
                                    Hence $(a,b_1)\in x$ by \cref{union_iff,singleton_iff,pair_eq_iff}.
                                    $(a,b_2)\in x$ by \cref{union_iff,singleton_iff,pair_eq_iff}.
                                    Thus $b_1 = b_2$ by \cref{rightunique,relation}.
                                \end{byCase}
                            \end{subproof}
                            Hence $v$ is a function by \cref{rightunique}.
                            Show that for all $i$ such that $i\in\dom{Y_1}$ we have $i\in\dom{v}$.
                            \begin{subproof}
                                Follows by \cref{subseteq,union_iff,singleton_iff,dom_intro,real_naturals_leq_succ_elim,initial_segment_naturalsplus,indexed_product,funs_elim,dom_iff}.
                            \end{subproof}
                            Hence $\dom{Y_1}\subseteq\dom{v}$ by \cref{subseteq}.
                            Show that for all $i$ such that $i\in\dom{v}$ we have $i\in\dom{Y_1}$.
                            \begin{subproof}
                                Fix $i$ such that $i\in\dom{v}$.
                                Take $b$ such that $(i,b)\in v$ by \cref{dom_iff}.
                                We have $(i,b)\in x$ or $(i,b) = (k + 1,y)$ by \cref{union_iff,singleton_iff}.
                                \begin{byCase}
                                \caseOf{$(i,b) = (k + 1,y)$.}
                                    Hence $i = k + 1$ by \cref{pair_eq_iff}.
                                    We have $i\in J$.
                                    Therefore $i\in\dom{Y_1}$ by \cref{subseteq}.
                                \caseOf{$(i,b)\in x$.}
                                    Hence $i\in\dom{x}$ by \cref{dom_intro}.
                                    We have $i\in J_1$ by \cref{indexed_product,funs_elim}.
                                    Therefore $i\in J$ by \cref{subseteq,initial_segment_naturalsplus}.
                                    Hence $i\in\dom{Y_1}$ by \cref{subseteq}.
                                \end{byCase}
                            \end{subproof}
                            Therefore $\dom{v}\subseteq\dom{Y_1}$ by \cref{subseteq}.
                            Hence $\dom{v} = \dom{Y_1}$ by \cref{subseteq_antisymmetric}.
                        Show that for all $i\in\dom{Y_1}$ we have $v(i)\in\unions{\ran{Y_1}}$.
                        \begin{subproof}
                            Fix $i$ such that $i\in\dom{Y_1}$.
                            Hence $i\in J$ by \cref{subseteq}.
                            \begin{byCase}
                            \caseOf{$i = k + 1$.}
                                We have $(k + 1,y)\in v$ by \cref{union_iff,singleton_iff}.
                                Hence $v(k + 1) = y$ by \cref{function_apply_intro}.
                                $(k + 1,X_1)\in Y_1$ by \cref{function_apply_elim,subseteq}.
                                Hence $X_1\in\ran{Y_1}$ by \cref{ran_intro}.
                                Therefore $y\in\unions{\ran{Y_1}}$ by \cref{unions_iff}.
                                Therefore $v(i)\in\unions{\ran{Y_1}}$.
                            \caseOf{$i \neq k + 1$.}
                                Hence $i\in J_1$ by \cref{real_naturals_leq_succ_elim,initial_segment_naturalsplus}.
                                We have $i\in\dom{x}$ by \cref{subseteq,indexed_product,funs_elim}.
                                Therefore $x(i)\in\apply{Y_0}{i}$ by \cref{indexed_product}.
                                $(i,\apply{Y_1}{i})\in Y_0$.
                                Hence $x(i)\in\apply{Y_1}{i}$ by \cref{function_apply_intro}.
                                Therefore $(i,\apply{Y_1}{i})\in Y_1$ by \cref{function_apply_elim}.
                                Hence $\apply{Y_1}{i}\in\ran{Y_1}$ by \cref{ran_intro}.
                                Therefore $x(i)\in\unions{\ran{Y_1}}$ by \cref{unions_iff}.
                                We have $(i,x(i))\in v$ by \cref{function_apply_elim,union_iff}.
                                Hence $v(i) = x(i)$ by \cref{function_apply_intro}.
                                Hence $v(i)\in\unions{\ran{Y_1}}$.
                            \end{byCase}
                        \end{subproof}
                        Therefore $v\in\funs{\dom{Y_1}}{\unions{\ran{Y_1}}}$ by \cref{funs_intro}.
                    \end{subproof}
                    Hence $v$ is right-unique and $v$ is a relation by \cref{funs,rels_is_relation}.
                    Show that for all $i\in\dom{Y_1}$ we have $v(i)\in\apply{Y_1}{i}$.
                    \begin{subproof}
                        Fix $i$ such that $i\in\dom{Y_1}$.
                        Hence $i\in J$ by \cref{subseteq}.
                        \begin{byCase}
                        \caseOf{$i = k + 1$.}
                            We have $(k + 1,y)\in v$ by \cref{union_iff,singleton_iff}.
                            Hence $v(k + 1) = y$ by \cref{function_apply_intro}.
                            Therefore $y\in X_1$ by \cref{subseteq}.
                            Hence $v(i)\in\apply{Y_1}{i}$.
                        \caseOf{$i \neq k + 1$.}
                            Hence $i\in J_1$ by \cref{real_naturals_leq_succ_elim,initial_segment_naturalsplus}.
                            We have $x(i)\in\apply{Y_0}{i}$ by \cref{indexed_product}.
                            Hence $x(i)\in\apply{Y_1}{i}$ by \cref{function_apply_intro}.
                            Therefore $(i,x(i))\in v$ by \cref{function_apply_elim,union_iff}.
                            Hence $v(i) = x(i)$ by \cref{function_apply_intro}.
                            Hence $v(i)\in\apply{Y_1}{i}$.
                        \end{byCase}
                    \end{subproof}
                    Therefore $v\in\productFamily{Y_1}$ by \cref{indexed_product}.
                Hence $g(p)\in\productFamily{Y_1}$ by \cref{function_apply_intro}.
            \end{subproof}
            Therefore $g\in\funs{X_0\times X_1}{\productFamily{Y_1}}$ by \cref{funs_intro}.
        Hence $g$ is right-unique and $g$ is a relation by \cref{funs,rels_is_relation}.
        Show that for all $i\in J$ we have $\piIndex{Y_1}{i}\circ g$ is continuous from $X_0\times X_1$
            to $\apply{Y_1}{i}$ with respect to $T_P$ and $\apply{T_1}{i}$.
        \begin{subproof}
            Fix $i$ such that $i\in J$.
            \begin{byCase}
            \caseOf{$i = k + 1$.}
                    Let $h = \piTwo{X_0}{X_1}$.
                    $h$ is continuous from $X_0\times X_1$ to $X_1$ with respect to $T_P$ and $T_{1p}$
                        by \cref{projection_two_continuous}.
                    Show that for all $p\in X_0\times X_1$ we have $\apply{\piIndex{Y_1}{i}\circ g}{p} = \apply{h}{p}$.
                    \begin{subproof}
                        Fix $p$ such that $p\in X_0\times X_1$.
                        We have $g$ is a function to $\productFamily{Y_1}$ and $g$ is a function
                            by \cref{funs_elim,funs_is_function}.
                        Hence $g(p)\in\productFamily{Y_1}$
                            by \cref{funs,function_apply_elim,ran_intro,function_to_ran,subseteq}.
                        $(p,\fst{p}\union\{(k + 1,\snd{p})\})\in g$.
                        Thus $g(p) = \fst{p}\union\{(k + 1,\snd{p})\}$ by \cref{function_apply_intro}.
                        Hence $(i,\snd{p})\in g(p)$ by \cref{union_iff,singleton_iff}.
                        We have $g(p)\in\funs{\dom{Y_1}}{\unions{\ran{Y_1}}}$
                            and $g(p)$ is a function by \cref{indexed_product,funs_is_function}.
                        Therefore $\apply{g(p)}{i} = \snd{p}$ by \cref{function_apply_intro}.
                        Hence $(g(p),\snd{p})\in\piIndex{Y_1}{i}$ by \cref{pair_eq_iff,projection_map_indexed}.
                        Therefore $(p,\snd{p})\in\piIndex{Y_1}{i}\circ g$ by \cref{circ_iff}.
                        We have $\piIndex{Y_1}{i}\in\funs{\productFamily{Y_1}}{\apply{Y_1}{i}}$
                            and $\piIndex{Y_1}{i}$ is a function by \cref{projection_map_indexed_function,funs_is_function}.
                        Hence $\piIndex{Y_1}{i}\circ g$ is a function by \cref{function_circ}.
                        Therefore $\apply{\piIndex{Y_1}{i}\circ g}{p} = \snd{p}$ by \cref{function_apply_intro}.
                        We have $h\in\funs{X_0\times X_1}{X_1}$ and $h$ is a function by \cref{continuous,funs_is_function}.
                        Take $x,y$ such that $p = (x,y)$ by \cref{times_elem_is_tuple}.
                        $(x,y)\in X_0\times X_1$.
                        Hence $x\in X_0$ and $y\in X_1$ by \cref{times_tuple_elim}.
                        Therefore $\snd{p} = y$ by \cref{snd_eq}.
                        Hence $(p,\snd{p})\in h$ by \cref{pair_eq_iff,projection_second}.
                        Therefore $\apply{\piIndex{Y_1}{i}\circ g}{p} = \apply{h}{p}$
                            by \cref{function_apply_intro}.
                    \end{subproof}
                    $\piIndex{Y_1}{i}\in\funs{\productFamily{Y_1}}{\apply{Y_1}{i}}$
                        by \cref{projection_map_indexed_function}.
                    Thus $\piIndex{Y_1}{i}\circ g\in\funs{X_0\times X_1}{\apply{Y_1}{i}}$
                        by \cref{funs_circ,funs_elim}.
                    Thus $\piIndex{Y_1}{i}\circ g$ is a function and $\dom{\piIndex{Y_1}{i}\circ g} = X_0\times X_1$
                        by \cref{funs_elim,funs_is_function}.
                    $h$ is a function and $\dom{h} = X_0\times X_1$ by \cref{continuous,funs_is_function,funs_elim}.
                    Hence $\piIndex{Y_1}{i}\circ g = h$ by \cref{funs_elim,subseteq,fun_subseteq,subseteq_antisymmetric}.
                    Thus $\piIndex{Y_1}{i}\circ g$ is continuous from $X_0\times X_1$ to $\apply{Y_1}{i}$ with respect to $T_P$ and $\apply{T_1}{i}$
                        by \cref{continuous_eq_subst}.
                \caseOf{$i \neq k + 1$.}
                    Hence $i\in J_1$ by \cref{real_naturals_leq_succ_elim,initial_segment_naturalsplus}.
                    Let $h_0 = \piOne{X_0}{X_1}$.
                    $h_0$ is continuous from $X_0\times X_1$ to $X_0$ with respect to $T_P$ and $T_X$
                        by \cref{projection_one_continuous}.
                    $\piIndex{Y_0}{i}$ is continuous from $X_0$ to $\apply{Y_0}{i}$ with respect to $T_X$ and $\apply{T_0}{i}$
                        by \cref{projection_map_indexed_continuous,subseteq,subseteq_antisymmetric}.
                    Let $h = \piIndex{Y_0}{i}\circ h_0$.
                    $h$ is continuous from $X_0\times X_1$ to $\apply{Y_0}{i}$ with respect to $T_P$ and $\apply{T_0}{i}$
                        by \cref{continuous_composite}.
                    Show that for all $p\in X_0\times X_1$ we have $\apply{\piIndex{Y_1}{i}\circ g}{p} = \apply{h}{p}$.
                    \begin{subproof}
                            Fix $p$ such that $p\in X_0\times X_1$.
                            $(p,\fst{p}\union\{(k + 1,\snd{p})\})\in g$.
                            Thus $g(p) = \fst{p}\union\{(k + 1,\snd{p})\}$ by \cref{function_apply_intro}.
                            Let $x = \fst{p}$.
                            Take $a,b$ such that $p = (a,b)$ and $a\in X_0$ and $b\in X_1$ by \cref{times_elem_is_tuple}.
                            Hence $\fst{p} = a$ by \cref{fst_eq}.
                            We have $x = a$.
                            Therefore $x\in X_0$.
                            Hence $(p,x)\in h_0$ by \cref{pair_eq_iff,projection_first}.
                            Therefore $(x,\apply{x}{i})\in\piIndex{Y_0}{i}$ by \cref{projection_map_indexed}.
                            Hence $(p,\apply{x}{i})\in\piIndex{Y_0}{i}\circ h_0$ by \cref{circ_iff}.
                            We have $h\in\funs{X_0\times X_1}{\apply{Y_0}{i}}$ and $h$ is a function
                                by \cref{continuous,funs_is_function}.
                            Therefore $\apply{h}{p} = \apply{x}{i}$ by \cref{function_apply_intro}.
                            We have $g$ is a function to $\productFamily{Y_1}$ and $g$ is a function
                                by \cref{funs_elim,funs_is_function}.
                            Hence $g(p)\in\productFamily{Y_1}$
                                by \cref{funs,function_apply_elim,ran_intro,function_to_ran,subseteq}.
                            We have $g(p)\in\funs{\dom{Y_1}}{\unions{\ran{Y_1}}}$
                                and $g(p)$ is a function by \cref{indexed_product,funs_is_function}.
                            We have $x\in\funs{\dom{Y_0}}{\unions{\ran{Y_0}}}$ and $x$ is a function
                                by \cref{indexed_product,funs_is_function}.
                            Hence $i\in\dom{x}$ by \cref{funs_elim,subseteq}.
                            Therefore $(i,\apply{x}{i})\in x$ by \cref{function_apply_elim}.
                            Hence $(i,\apply{x}{i})\in g(p)$ by \cref{union_iff}.
                            Therefore $\apply{g(p)}{i} = \apply{x}{i}$ by \cref{function_apply_intro}.
                            Hence $(g(p),\apply{x}{i})\in\piIndex{Y_1}{i}$ by \cref{projection_map_indexed}.
                            Therefore $(p,\apply{x}{i})\in\piIndex{Y_1}{i}\circ g$ by \cref{circ_iff}.
                            We have $\piIndex{Y_1}{i}\in\funs{\productFamily{Y_1}}{\apply{Y_1}{i}}$ and
                                $\piIndex{Y_1}{i}$ is a function by \cref{projection_map_indexed_function,funs_is_function}.
                            Hence $\piIndex{Y_1}{i}\circ g$ is a function by \cref{function_circ}.
                            Therefore $\apply{\piIndex{Y_1}{i}\circ g}{p} = \apply{h}{p}$
                                by \cref{function_apply_intro}.
                    \end{subproof}
                    $\piIndex{Y_1}{i}\in\funs{\productFamily{Y_1}}{\apply{Y_1}{i}}$
                        by \cref{projection_map_indexed_function}.
                    Thus $\piIndex{Y_1}{i}\circ g\in\funs{X_0\times X_1}{\apply{Y_1}{i}}$
                        by \cref{funs_circ,funs_elim}.
                    Thus $\piIndex{Y_1}{i}\circ g$ is a function and $\dom{\piIndex{Y_1}{i}\circ g} = X_0\times X_1$
                        by \cref{funs_elim,funs_is_function}.
                    $h$ is a function and $\dom{h} = X_0\times X_1$ by \cref{continuous,funs_is_function,funs_elim}.
                    Hence $\piIndex{Y_1}{i}\circ g = h$ by \cref{funs_elim,subseteq,fun_subseteq,subseteq_antisymmetric}.
                    Hence $h$ is continuous from $X_0\times X_1$ to $\apply{Y_1}{i}$ with respect to $T_P$ and $\apply{T_1}{i}$
                        by \cref{function_apply_intro,continuous_eq_codomain}.
                    Hence $h\in\funs{X_0\times X_1}{\apply{Y_1}{i}}$ and $h$ is a function
                        by \cref{continuous,funs_is_function}.
                    Thus $\piIndex{Y_1}{i}\in\funs{\productFamily{Y_1}}{\apply{Y_1}{i}}$ and
                        $\piIndex{Y_1}{i}$ is a function by \cref{projection_map_indexed_function,funs_is_function}.
                    Hence $\piIndex{Y_1}{i}\circ g$ is a function by \cref{function_circ}.
                    Therefore $\piIndex{Y_1}{i}\circ g$ is continuous from $X_0\times X_1$ to $\apply{Y_1}{i}$ with respect to $T_P$ and $\apply{T_1}{i}$
                        by \cref{continuous_eq_subst}.
            \end{byCase}
        \end{subproof}
        $J$ is inhabited.
        $T_P$ is a topology on $X_0\times X_1$.
        Hence $g$ is continuous from $X_0\times X_1$ to $\productFamily{Y_1}$ with respect to $T_P$ and $\productTopologyIndexed{T_1}{Y_1}$
            by \cref{continuous_map_into_indexed_product}.
        Show that for all $z\in\productFamily{Y_1}$ there exists $p\in X_0\times X_1$ such that $g(p) = z$.
        \begin{subproof}
                Fix $z$ such that $z\in\productFamily{Y_1}$.
                We have $z$ is a function and $\dom{z} = \dom{Y_1}$ by \cref{indexed_product,funs_elim}.
                Let $x = \{(i,z(i)) \mid i\in J_1\}$.
                $x$ is a relation by \cref{relation}.
                        For all $a,b_1,b_2$ such that $(a,b_1)\in x$ and $(a,b_2)\in x$ we have $b_1 = b_2$
                            by \cref{pair_eq_iff}.
                        Hence $x$ is a function by \cref{rightunique}.
                        We have $\dom{x} = \dom{Y_0}$ by \cref{dom_intro,dom_iff,pair_eq_iff,subseteq,subseteq_antisymmetric}.
                        For all $i$ such that $i\in\dom{Y_0}$ we have $x(i)\in\unions{\ran{Y_0}}$
                            by \cref{subseteq,function_apply_intro,indexed_product,ran_intro,unions_iff}.
                        Therefore $x\in\funs{\dom{Y_0}}{\unions{\ran{Y_0}}}$ and $x$ is a function
                            by \cref{funs_intro,funs_is_function}.
                        For all $i\in\dom{Y_0}$ we have $x(i)\in\apply{Y_0}{i}$
                            by \cref{subseteq,function_apply_intro,indexed_product}.
                Thus $x\in\productFamily{Y_0}$ by \cref{indexed_product}.
                Hence $x\in X_0$ by \cref{subseteq}.
                Let $y = z(k + 1)$.
                Hence $y\in X_1$ by \cref{indexed_product,subseteq}.
                Let $p = (x,y)$.
                Hence $p\in X_0\times X_1$ by \cref{times_tuple_intro}.
                Show that for all $i\in\dom{z}$ we have $\apply{g(p)}{i} = z(i)$.
                    \begin{subproof}
                        Fix $i$ such that $i\in\dom{z}$.
                        Hence $i\in J$ by \cref{subseteq}.
                        We have $p\in\dom{g}$ by \cref{funs}.
                        Hence $g(p)\in\productFamily{Y_1}$
                            by \cref{funs_elim,funs,function_apply_elim,ran_intro,function_to_ran,subseteq}.
                        We have $g(p)\in\funs{\dom{Y_1}}{\unions{\ran{Y_1}}}$
                            and $g(p)$ is a function by \cref{indexed_product,funs_is_function}.
                        $(p,\fst{p}\union\{(k + 1,\snd{p})\})\in g$.
                        Thus $g(p) = \fst{p}\union\{(k + 1,\snd{p})\}$ by \cref{function_apply_intro}.
                        Hence $g(p) = x\union\{(k + 1,y)\}$ by \cref{fst_eq,snd_eq}.
                        \begin{byCase}
                        \caseOf{$i = k + 1$.}
                            We have $(k + 1,y)\in\{(k + 1,y)\}$ by \cref{singleton_intro}.
                            Hence $(k + 1,y)\in g(p)$ by \cref{union_iff}.
                            Therefore $\apply{g(p)}{k + 1} = y$ by \cref{function_apply_intro}.
                            Hence $\apply{g(p)}{i} = z(i)$.
                        \caseOf{$i \neq k + 1$.}
                            Hence $i\in J_1$ by \cref{real_naturals_leq_succ_elim,initial_segment_naturalsplus}.
                            We have $(i,z(i))\in x$.
                            Hence $(i,z(i))\in g(p)$ by \cref{union_iff}.
                            Therefore $\apply{g(p)}{i} = z(i)$ by \cref{function_apply_intro}.
                        \end{byCase}
                    \end{subproof}
                We have $g$ is a function to $\productFamily{Y_1}$ by \cref{funs_elim}.
                Hence $p\in\dom{g}$ by \cref{funs}.
                Hence $g(p)\in\productFamily{Y_1}$
                    by \cref{funs_elim,funs,function_apply_elim,ran_intro,function_to_ran,subseteq}.
                Therefore $g(p)$ is a function and $\dom{g(p)} = \dom{Y_1}$ by \cref{indexed_product,funs_elim}.
                Therefore there exists $q\in X_0\times X_1$ such that $g(q) = z$
                    by \cref{subseteq,fun_subseteq,subseteq_antisymmetric}.
            \end{subproof}
        Therefore $g\in\surj{X_0\times X_1}{\productFamily{Y_1}}$ by \cref{surj}.
        Hence $\img{g}{X_0\times X_1} = \productFamily{Y_1}$ by \cref{surj,funs_elim,img_dom,funs_surj_iff}.
        Let $A = \img{g}{X_0\times X_1}$.
        Let $T_Q = \productTopologyIndexed{T_1}{Y_1}$.
        Let $T_A = \subspaceTopology{T_Q}{A}$.
        Hence $T_A = \subspaceTopology{T_Q}{\productFamily{Y_1}}$.
        $T_Q$ is a topology on $\productFamily{Y_1}$
            by \cref{product_subbasis_finite_intersections_basis,basis_for_topology,basis_topology_is_topology}.
        Therefore $T_A = T_Q$
            by \cref{subspace_topology,topology_on,subseteq,powerset_elim,inter_absorb_supseteq_left,subseteq_antisymmetric}.
        Hence $\productFamily{Y_1}$ is compact with respect to $T_Q$ by \cref{continuous_image_compact}.
    \end{subproof}
    Hence $\finiteIndexedProductCompactProp{k + 1}$ by \cref{finite_indexed_product_compact_property}.
\end{proof}

% from §27 helper lemma 6: compact indexed product on an initial segment
\begin{lemma}\label{compact_indexed_product_initial_segment}
    Suppose $n\in\naturalsPlus$.
    Let $J = \initialSegment{n}$.
    Assume $Y$ is a function.
    Assume $T_Y$ is a function.
    Assume $\dom{Y} = J$.
    Assume $\dom{T_Y} = J$.
    Assume that for all $i\in J$ we have $\apply{T_Y}{i}$ is a topology on $\apply{Y}{i}$.
    Assume that for all $i\in J$ we have $\apply{Y}{i}$ is compact with respect to $\apply{T_Y}{i}$.
    Then $\productFamily{Y}$ is compact with respect to $\productTopologyIndexed{T_Y}{Y}$.
\end{lemma}
\begin{proof}
    Let $A = \{m\in\naturalsPlus \mid \finiteIndexedProductCompactProp{m}\}$.
    $A\subseteq\naturalsPlus$ by \cref{subseteq}.
    Show that $1\in A$.
    \begin{subproof}
        It suffices to show that $1\in\naturalsPlus$ and $\finiteIndexedProductCompactProp{1}$.
        Show that for all $Y_1,T_1$ such that $Y_1$ is a function and $T_1$ is a function and
            $\dom{Y_1} = \initialSegment{1}$ and $\dom{T_1} = \initialSegment{1}$ and
            for all $i\in\dom{Y_1}$ we have $\apply{T_1}{i}$ is a topology on $\apply{Y_1}{i}$ and
                $\apply{Y_1}{i}$ is compact with respect to $\apply{T_1}{i}$,
            we have $\productFamily{Y_1}$ is compact with respect to $\productTopologyIndexed{T_1}{Y_1}$.
        \begin{subproof}
            Fix $Y_1$.
            Fix $T_1$.
            Assume $Y_1$ is a function.
            Assume $T_1$ is a function.
            Assume $\dom{Y_1} = \initialSegment{1}$.
            Assume $\dom{T_1} = \initialSegment{1}$.
            Assume that for all $i\in\dom{Y_1}$ we have $\apply{T_1}{i}$ is a topology on $\apply{Y_1}{i}$ and
                $\apply{Y_1}{i}$ is compact with respect to $\apply{T_1}{i}$.
            Let $J_1 = \initialSegment{1}$.
            We have $1\in J_1$ by \cref{real_one_in_naturals,initial_segment_naturalsplus}.
            For all $i\in J_1$ we have $i < 1$ or $i = 1$ by \cref{initial_segment_naturalsplus}.
            For all $i\in J_1$ we have $1\leq i$ by \cref{initial_segment_naturalsplus,real_one_leq_naturalsplus}.
            Thus for all $i\in J_1$ we have $i = 1$.
            Therefore $J_1 = \{1,1\}$ by \cref{subseteq_antisymmetric,subseteq,upair_iff}.
            Let $X = \apply{Y_1}{1}$.
            Let $T_X = \apply{T_1}{1}$.
            Show that $T_X$ is a topology on $X$ and $X$ is compact with respect to $T_X$.
            \begin{subproof}
                Follows by \cref{subseteq}.
            \end{subproof}
            Let $g = \{(x,\{(1,x)\}) \mid x\in X\}$.
            $g$ is a relation by \cref{relation}.
            For all $x,u,v$ such that $(x,u)\in g$ and $(x,v)\in g$ we have $u = v$ by \cref{pair_eq_iff}.
            Hence $g$ is a function by \cref{rightunique}.
            We have $\dom{g} = X$ by \cref{dom_intro,dom_iff,pair_eq_iff,subseteq,subseteq_antisymmetric}.
            Show that for all $x\in\dom{g}$ we have $g(x)\in\productFamily{Y_1}$.
            \begin{subproof}
                Fix $x$ such that $x\in\dom{g}$.
                Take $y$ such that $(x,y)\in g$ by \cref{dom_iff}.
                Take $x_0\in X$ such that $(x,y) = (x_0,\{(1,x_0)\})$.
                We have $y = \{(1,x_0)\}$ by \cref{pair_eq_iff}.
                Hence $y$ is a relation by \cref{singleton_iff,relation}.
                For all $a,b_1,b_2$ such that $(a,b_1)\in y$ and $(a,b_2)\in y$ we have $b_1 = b_2$
                    by \cref{singleton_iff,pair_eq_iff}.
                Hence $y$ is a function by \cref{rightunique}.
                Show that for all $a$ such that $a\in\dom{y}$ we have $a\in\dom{Y_1}$.
                \begin{subproof}
                    Fix $a$ such that $a\in\dom{y}$.
                    Take $b$ such that $(a,b)\in y$ by \cref{dom_iff}.
                    We have $(a,b) = (1,x_0)$ by \cref{singleton_iff}.
                    Hence $a = 1$ by \cref{pair_eq_iff}.
                    Therefore $a\in J_1$.
                    Therefore $a\in\dom{Y_1}$.
                \end{subproof}
                Hence $\dom{y}\subseteq\dom{Y_1}$ by \cref{subseteq}.
                Show that for all $a$ such that $a\in\dom{Y_1}$ we have $a\in\dom{y}$.
                \begin{subproof}
                    Fix $a$ such that $a\in\dom{Y_1}$.
                    We have $a\in J_1$.
                    Hence $a = 1$.
                    Therefore $(a,x_0)\in y$ by \cref{singleton_iff}.
                    Therefore $a\in\dom{y}$ by \cref{dom_intro}.
                \end{subproof}
                Hence $\dom{Y_1}\subseteq\dom{y}$ by \cref{subseteq}.
                Therefore $\dom{y} = \dom{Y_1}$ by \cref{subseteq_antisymmetric}.
                For all $i$ such that $i\in\dom{Y_1}$ we have $y(i)\in\unions{\ran{Y_1}}$
                    by \cref{subseteq,singleton_iff,function_apply_intro,function_apply_elim,ran_intro,unions_iff}.
                Hence $y\in\funs{\dom{Y_1}}{\unions{\ran{Y_1}}}$ by \cref{funs_intro}.
                For all $i$ such that $i\in\dom{Y_1}$ we have $y(i)\in\apply{Y_1}{i}$
                    by \cref{subseteq,singleton_iff,funs_elim,function_apply_intro}.
                Therefore $y\in\productFamily{Y_1}$ by \cref{indexed_product}.
                Hence $g(x)\in\productFamily{Y_1}$ by \cref{function_apply_intro}.
            \end{subproof}
            Therefore $g\in\funs{X}{\productFamily{Y_1}}$ by \cref{funs_intro}.
            Let $T_P = \productTopologyIndexed{T_1}{Y_1}$.
        $J_1$ is inhabited.
        Show that for all $i\in J_1$ we have $\piIndex{Y_1}{i}\circ g$ is continuous from $X$ to $\apply{Y_1}{i}$ with respect to $T_X$ and $\apply{T_1}{i}$.
        \begin{subproof}
            Fix $i$ such that $i\in J_1$.
            Hence $i = 1$.
            Let $j = \identity{X}$.
            $j$ is continuous from $X$ to $X$ with respect to $T_X$ and $T_X$ by \cref{identity_continuous}.
                    Show that for all $x\in X$ we have $\apply{\piIndex{Y_1}{i}\circ g}{x} = \apply{j}{x}$.
                    \begin{subproof}
                        Fix $x$ such that $x\in X$.
                        We have $g$ is a function to $\productFamily{Y_1}$ and $g$ is a function and $\dom{g} = X$
                            by \cref{funs_elim,funs_is_function}.
                        Hence $g(x)\in\productFamily{Y_1}$ by \cref{subseteq,funs,function_apply_elim,ran_intro,function_to_ran}.
                        $(x,\{(1,x)\})\in g$.
                        Thus $g(x) = \{(1,x)\}$ by \cref{function_apply_intro}.
                        We have $g(x)\in\funs{\dom{Y_1}}{\unions{\ran{Y_1}}}$ and $g(x)$ is a function
                            by \cref{indexed_product,funs_elim}.
                        Hence $(1,x)\in g(x)$ by \cref{singleton_iff}.
                        Hence $\apply{g(x)}{i} = x$ by \cref{function_apply_intro}.
                        Therefore $(g(x),x)\in\piIndex{Y_1}{i}$ by \cref{pair_eq_iff,projection_map_indexed}.
                        Hence $(x,x)\in\piIndex{Y_1}{i}\circ g$ by \cref{circ_iff}.
                        We have $\piIndex{Y_1}{i}\in\funs{\productFamily{Y_1}}{\apply{Y_1}{i}}$
                            and $\piIndex{Y_1}{i}$ is a function by \cref{projection_map_indexed_function,funs_is_function}.
                        Hence $\piIndex{Y_1}{i}\circ g$ is a function by \cref{function_circ}.
                        Therefore $\apply{\piIndex{Y_1}{i}\circ g}{x} = \apply{j}{x}$
                            by \cref{function_apply_intro,id_apply}.
                    \end{subproof}
            $\piIndex{Y_1}{i}\in\funs{\productFamily{Y_1}}{\apply{Y_1}{i}}$
                by \cref{projection_map_indexed_function}.
            Thus $\piIndex{Y_1}{i}\circ g\in\funs{X}{\apply{Y_1}{i}}$ by \cref{funs_circ,funs_elim}.
            Thus $\piIndex{Y_1}{i}\circ g$ is a function and $\dom{\piIndex{Y_1}{i}\circ g} = X$
                by \cref{funs_elim,funs_is_function}.
            $j$ is a function and $\dom{j} = X$ by \cref{id_is_function,id_dom}.
            Therefore $\piIndex{Y_1}{i}\circ g = j$ by \cref{subseteq,fun_subseteq,subseteq_antisymmetric}.
            Hence $\piIndex{Y_1}{i}\circ g$ is continuous from $X$ to $\apply{Y_1}{i}$ with respect to $T_X$ and $\apply{T_1}{i}$.
        \end{subproof}
        Therefore $g$ is continuous from $X$ to $\productFamily{Y_1}$ with respect to $T_X$ and $T_P$
            by \cref{continuous_map_into_indexed_product}.
        Show that for all $y\in\productFamily{Y_1}$ there exists $x\in X$ such that $g(x) = y$.
        \begin{subproof}
            Fix $y$ such that $y\in\productFamily{Y_1}$.
            We have $y$ is a function and $\dom{y} = \dom{Y_1}$ by \cref{indexed_product,funs_elim}.
            Hence $1\in\dom{y}$ by \cref{subseteq,real_one_in_naturals,initial_segment_naturalsplus}.
            Thus $y(1)\in X$ by \cref{indexed_product}.
            Take $x\in X$ such that $x = y(1)$.
            Thus $(x,\{(1,x)\})\in g$.
            For all $i$ such that $i\in\dom{y}$ we have $\apply{g(x)}{i} = y(i)$
                by \cref{subseteq,funs_elim,function_apply_intro,singleton_iff,function_apply_elim,ran_intro,function_to_ran,indexed_product}.
            Thus $g(x)$ is a function and $\dom{g(x)} = \dom{Y_1}$
                by \cref{subseteq,funs,funs_is_function,function_apply_elim,ran_intro,function_to_ran,indexed_product,funs_elim}.
            Therefore there exists $z\in X$ such that $g(z) = y$
                by \cref{subseteq,fun_subseteq,subseteq_antisymmetric}.
        \end{subproof}
        Therefore $g\in\surj{X}{\productFamily{Y_1}}$ by \cref{surj}.
        Hence $\img{g}{X} = \productFamily{Y_1}$ by \cref{surj,funs_elim,img_dom,funs_surj_iff}.
        Let $A_0 = \img{g}{X}$.
        Let $T_A = \subspaceTopology{T_P}{A_0}$.
        We have $T_A = \subspaceTopology{T_P}{\productFamily{Y_1}}$.
        $T_P$ is a topology on $\productFamily{Y_1}$
            by \cref{product_subbasis_finite_intersections_basis,basis_for_topology,basis_topology_is_topology}.
        Therefore $T_A = T_P$
            by \cref{subspace_topology,topology_on,subseteq,powerset_elim,inter_absorb_supseteq_left,subseteq_antisymmetric}.
        Hence $\productFamily{Y_1}$ is compact with respect to $T_P$ by \cref{continuous_image_compact}.
    \end{subproof}
        Hence $\finiteIndexedProductCompactProp{1}$ by \cref{finite_indexed_product_compact_property}.
        Therefore $1\in A$ by \cref{real_one_in_naturals}.
    \end{subproof}
    For all $k$ such that $k\in A$ we have $k\in\naturalsPlus$ and $\finiteIndexedProductCompactProp{k}$.
    Thus for all $k\in A$ we have $k + 1\in A$
        by \cref{real_naturals_closed_succ,finite_indexed_product_compact_property_succ}.
    Hence $n\in A$ by \cref{real_naturals_induction}.
    Therefore $\productFamily{Y}$ is compact with respect to $\productTopologyIndexed{T_Y}{Y}$
        by \cref{subseteq,finite_indexed_product_compact_property}.
\end{proof}
% from §27 helper lemma 9: square and euclidean metric topologies agree
\begin{lemma}\label{square_metric_topology_equals_euclidean}
    Suppose $n\in\naturalsPlus$.
    Let $J = \initialSegment{n}$.
    Let $X = \realsPower{J}$.
    Let $T_E = \metricTopology{\euclideanMetric{n}}{X}$.
    Let $T_S = \metricTopology{\squareMetric{n}}{X}$.
    Then $T_E = T_S$.
\end{lemma}
\begin{proof}
    Let $X_0 = X$.
    Let $d = \euclideanMetric{n}$.
    Let $d_1 = \squareMetric{n}$.
    $d$ is a metric on $X_0$ and $d_1$ is a metric on $X_0$ by \cref{euclidean_metric_is_metric,square_metric_is_metric}.
    Show that for all $x\in X_0$ we have
        for all $\epsilon\in\reals$ such that $0 < \epsilon$ there exists $\delta\in\reals$
            such that $0 < \delta$ and
            $\metricBall{d}{X_0}{x}{\delta} \subseteq \metricBall{d_1}{X_0}{x}{\epsilon}$.
    \begin{subproof}
        Fix $x$ such that $x\in X_0$.
        Fix $\epsilon$ such that $\epsilon\in\reals$ and $0 < \epsilon$.
        Let $\delta = \epsilon$.
        Show that $0 < \delta$ and
            $\metricBall{d}{X_0}{x}{\delta} \subseteq \metricBall{d_1}{X_0}{x}{\epsilon}$.
        \begin{subproof}
            Follows by \cref{subseteq,metric_ball,euclidean_metric_value_exists,square_metric_value_exists,euclidean_square_metric_bounds,euclidean_metric_function,square_metric_function,real_leq_split,real_lt_subst_left,real_lt_trans,function_apply_intro}.
        \end{subproof}
    \end{subproof}
    Show that $T_E$ is finer than $T_S$ on $X_0$.
    \begin{subproof}
        Follows by \cref{metric_topology_finer_iff}.
    \end{subproof}
    Show that for all $x\in X_0$ we have
        for all $\epsilon\in\reals$ such that $0 < \epsilon$ there exists $\delta\in\reals$
            such that $0 < \delta$ and
            $\metricBall{d_1}{X_0}{x}{\delta} \subseteq \metricBall{d}{X_0}{x}{\epsilon}$.
    \begin{subproof}
        Fix $x$ such that $x\in X_0$.
        Fix $\epsilon$ such that $\epsilon\in\reals$ and $0 < \epsilon$.
        $0 < 1$ and $0\in\reals$ and $1\in\reals$ and $n\in\reals$
            by \cref{real_zero_lt_one,real_add_ident,real_naturals_subset_reals,real_one_in_naturals,subseteq}.
        Hence $0 < n$ by \cref{real_one_leq_naturalsplus,real_lt_subst_right,real_lt_trans}.
        We have $\sqrt{n}\in\reals$ and $\sqrt{n} \geq 0$ and $\sqrt{n} \rmul \sqrt{n} = n$ by \cref{positive_square_root}.
        Hence $\sqrt{n}$ is positive.
        Let $\delta = \epsilon \rmul \inv{\sqrt{n}}$.
        $\sqrt{n} \neq 0$ and $\inv{\sqrt{n}}\in\reals$ and $\inv{\sqrt{n}} > 0$
            by \cref{real_pos_nonzero,real_mul_inverse,real_inv_pos}.
        Therefore $\delta\in\reals$ and $\delta > 0$ by \cref{real_mul_closed,real_mul_pos_pos_gt}.
        Show that for all $y$ such that $y\in\metricBall{d_1}{X_0}{x}{\delta}$ we have
            $y\in\metricBall{d}{X_0}{x}{\epsilon}$.
        \begin{subproof}
            Fix $y$ such that $y\in\metricBall{d_1}{X_0}{x}{\delta}$.
            Thus $y\in X_0$ and $d_1((x,y)) < \delta$ by \cref{metric_ball}.
            Take $m\in\reals$ such that $((x,y),m)\in d_1$ by \cref{square_metric_value_exists}.
            Take $r\in\reals$ such that $((x,y),r)\in d$ by \cref{euclidean_metric_value_exists}.
            Hence $m < \delta$ by \cref{square_metric_function,function_apply_intro,real_lt_subst_left}.
            Hence $r \leq \sqrt{n} \rmul m$ by \cref{euclidean_square_metric_bounds}.
            $m \rmul \sqrt{n} < \delta \rmul \sqrt{n}$ by \cref{real_lt_mul_pos}.
            Hence $\sqrt{n} \rmul m < \sqrt{n} \rmul \delta$ by \cref{real_mul_comm}.
            Thus $\sqrt{n} \rmul \delta = \epsilon$ by
                \cref{real_mul_assoc,real_mul_comm,real_mul_inverse,real_mul_ident}.
            Hence $\sqrt{n} \rmul m < \epsilon$ by \cref{real_lt_subst_right}.
            Therefore $d((x,y)) < \epsilon$ by
                \cref{euclidean_metric_function,real_mul_closed,real_leq_split,real_lt_subst_left,real_lt_trans,function_apply_intro}.
            Hence $y\in\metricBall{d}{X_0}{x}{\epsilon}$ by \cref{metric_ball}.
        \end{subproof}
        Therefore $\metricBall{d_1}{X_0}{x}{\delta} \subseteq \metricBall{d}{X_0}{x}{\epsilon}$ by \cref{subseteq}.
    \end{subproof}
    Show that $T_S$ is finer than $T_E$ on $X_0$.
    \begin{subproof}
        Follows by \cref{metric_topology_finer_iff}.
    \end{subproof}
    Therefore $T_E = T_S$ by \cref{finer_topology,subseteq_antisymmetric}.
\end{proof}

% from §27 helper lemma 10: indexed real/standard families are functions with domain J
\begin{lemma}\label{reals_standard_indexed_families_data}
    Suppose $n\in\naturalsPlus$.
    Let $J = \initialSegment{n}$.
    Let $X_1 = \{(i,\reals) \mid i\in J\}$.
    Let $T_1 = \{(i,\standardTopologyR) \mid i\in J\}$.
    Then $X_1$ is a function and $T_1$ is a function and $\dom{X_1} = J$ and $\dom{T_1} = J$.
\end{lemma}
\begin{proof}
    Show that for all $i,y_1,y_2$ such that $(i,y_1)\in X_1$ and $(i,y_2)\in X_1$ we have $y_1 = y_2$.
    \begin{subproof}
        Trivial.
    \end{subproof}
    Show that $X_1$ is a function.
    \begin{subproof}
        Follows by \cref{relation,rightunique}.
    \end{subproof}
    Show that for all $i,y_1,y_2$ such that $(i,y_1)\in T_1$ and $(i,y_2)\in T_1$ we have $y_1 = y_2$.
    \begin{subproof}
        Trivial.
    \end{subproof}
    Show that $T_1$ is a function.
    \begin{subproof}
        Follows by \cref{relation,rightunique}.
    \end{subproof}
    Show that $\dom{X_1} = J$.
    \begin{subproof}
        Follows by \cref{dom_intro,dom_iff,pair_eq_iff,subseteq,subseteq_antisymmetric}.
    \end{subproof}
    Show that $\dom{T_1} = J$.
    \begin{subproof}
        Follows by \cref{dom_intro,dom_iff,pair_eq_iff,subseteq,subseteq_antisymmetric}.
    \end{subproof}
\end{proof}

% from §27 Item 3: compact subsets of $\reals^n$
\begin{theorem}\label{compact_reals_power_iff_closed_bounded}
    Suppose $n\in\naturalsPlus$.
    Let $J = \initialSegment{n}$.
    Let $X = \realsPower{J}$.
    Let $d = \euclideanMetric{n}$.
    Let $\rho = \squareMetric{n}$.
    Let $T = \metricTopology{\rho}{X}$.
    Suppose $A\subseteq X$.
    Let $T_A = \subspaceTopology{T}{A}$.
    Then $A$ is compact with respect to $T_A$ iff
        $A$ is closed in $X$ with respect to $T$ and
        $A$ is bounded in $X$ with respect to $d$ or $A$ is bounded in $X$ with respect to $\rho$.
\end{theorem}
\begin{proof}
    Let $X_0 = X$.
    Let $d_1 = \rho$.
    Show that if $A$ is compact with respect to $T_A$ then
        $A$ is closed in $X$ with respect to $T$ and
        $A$ is bounded in $X$ with respect to $d$ or $A$ is bounded in $X$ with respect to $\rho$.
    \begin{subproof}
        Assume $A$ is compact with respect to $T_A$.
        $d_1$ is a metric on $X$ by \cref{square_metric_is_metric}.
        $T$ is a topology on $X$ by \cref{basis_topology_is_topology,metric_topology,metric_basis_is_basis,basis_for_topology}.
        Let $T_S = \metricTopology{d_1}{X_0}$.
        Let $X_1 = \{(i,\reals) \mid i\in J\}$.
        Let $T_1 = \{(i,\standardTopologyR) \mid i\in J\}$.
        Let $T_0 = \productTopologyIndexed{T_1}{X_1}$.
        $d$ is a metric on $X_0$ and $d_1$ is a metric on $X_0$ by \cref{euclidean_metric_is_metric,square_metric_is_metric}.
        Hence $J$ is inhabited by \cref{real_one_in_naturals,real_one_leq_naturalsplus,initial_segment_naturalsplus}.
        $X_1$ is a function and $T_1$ is a function and $\dom{X_1} = J$ and $\dom{T_1} = J$
            by \cref{reals_standard_indexed_families_data}.
        For all $\alpha\in\dom{X_1}$ we have $\apply{X_1}{\alpha} = \reals$ and $\apply{T_1}{\alpha} = \standardTopologyR$
            by \cref{subseteq,function_apply_intro}.
        $\realOrderTopology = \standardTopologyR$ and $\realOrderTopology = \basisTopology{\orderBasis{\realOrderRel}{\reals}}{\reals}$
            by \cref{real_order_topology_eq_standard,real_order_topology}.
        $\realOrderRel$ is a simple order relation on $\reals$ by \cref{real_order_relation_simple}.
        $\reals$ has more than one element by \cref{linear_continuum,reals_linear_continuum}.
        Hence $\realOrderTopology$ is the order topology on $\reals$ with respect to $\realOrderRel$ by \cref{order_topology}.
        Therefore for all $\alpha\in\dom{X_1}$ we have $\apply{X_1}{\alpha}$ is Hausdorff with respect to $\apply{T_1}{\alpha}$
            by \cref{order_topology_hausdorff,real_order_topology_eq_standard}.
        Hence $X_0$ is Hausdorff with respect to $T_0$ by \cref{product_hausdorff_product,product_family_reals_power}.
        Therefore $T_S = T_0$ by \cref{square_metric_topology_equals_euclidean,euclidean_metric_product_topology}.
        Hence $X$ is Hausdorff with respect to $T$.
        Show that $A$ is closed in $X$ with respect to $T$.
        \begin{subproof}
            Follows by \cref{compact_subspace_closed}.
        \end{subproof}
        Show that $A$ is bounded in $X$ with respect to $\rho$.
        \begin{subproof}
            \begin{byCase}
            \caseOf{$A = \emptyset$.}
                Let $M = 0$.
                For all $a_1,a_2\in A$ we have $\rho((a_1,a_2)) \leq M$
                    by \cref{subseteq_emptyset_iff,elem_subseteq,emptyset}.
                Thus $A$ is bounded in $X$ with respect to $\rho$ by \cref{real_add_ident,metric_bounded}.
            \caseOf{$A \neq \emptyset$.}
                $A$ is inhabited by \cref{nonempty_inhabited}.
                Let $C = \{\metricBall{\rho}{X}{a}{1} \mid a\in A\}$.
                For all $U$ such that $U\in C$ we have $U$ is open in $X$ with respect to $T$
                    by \cref{real_naturals_subset_reals,real_one_in_naturals,subseteq,real_zero_lt_one,metric_ball,powerset_intro,metric_basis,metric_basis_is_basis,basis_elem_open_in_generated,metric_topology,basis_for_topology,open_in}.
                Hence $C$ is a covering of $A$
                    by \cref{subseteq,covering,square_metric_is_metric,metric_zero_characterization,metric_on,real_zero_lt_one,real_lt_subst_left,metric_ball,unions_iff}.
                $A$ is a subspace of $X$ with respect to $T$ by \cref{subspace_of}.
                Thus there exists $B$ such that
                    $B\subseteq C$ and
                    $B$ is finite and
                    $B$ is a covering of $A$
                    by \cref{compact_subspace_open_covering}.
                Let $R = \{w\in B\times A \mid \fst{w} = \metricBall{\rho}{X}{\snd{w}}{1}\}$.
                Hence $R$ is a relation by \cref{times_elem_is_tuple,relation}.
                Show that for all $U\in B$ there exists $a$ such that $U\mathrel{R} a$.
                \begin{subproof}
                    Fix $U$ such that $U\in B$.
                    We have $U\in C$ by \cref{subseteq}.
                    Take $a\in A$ such that $U = \metricBall{\rho}{X}{a}{1}$.
                    Therefore $(U,a)\in B\times A$ by \cref{times_tuple_intro}.
                    Hence $\fst{(U,a)} = U$ and $\snd{(U,a)} = a$ by \cref{fst_eq,snd_eq}.
                    Therefore $(U,a)\in R$.
                    Hence $U\mathrel{R} a$.
                \end{subproof}
                Take $g$ such that $g$ is a function on $B$ and for all $U\in B$ we have $U\mathrel{R}\apply{g}{U}$
                    by \cref{choice_function}.
                Let $F = \img{g}{B}$.
                $F$ is finite by \cref{finite_image_function}.
                Hence $F\subseteq A$ by \cref{subseteq,img_iff,function_apply_intro,subseteq,times_tuple_elim}.
                Take $a_0$ such that $a_0\in A$.
                Thus $a_0\in\unions{B}$ by \cref{covering,subseteq}.
                Take $U_0\in B$ such that $a_0\in U_0$ by \cref{unions_iff}.
                Take $f_0$ such that $f_0\in F$ by \cref{subseteq,function_apply_elim,img_iff}.
                Let $h = \{(f,\rho((a_0,f))) \mid f\in F\}$.
                $h$ is a relation by \cref{relation}.
                For all $x,y_1,y_2$ such that $(x,y_1)\in h$ and $(x,y_2)\in h$ we have $y_1 = y_2$
                    by \cref{pair_eq_iff}.
                Hence $h$ is right-unique by \cref{rightunique}.
                For all $f$ such that $f\in F$ we have $f\in\dom{h}$.
                For all $f$ such that $f\in\dom{h}$ we have $f\in F$ by \cref{dom_iff,pair_eq_iff}.
                Therefore $\dom{h} = F$ by \cref{subseteq,subseteq_antisymmetric}.
                Hence $h$ is a function on $F$.
                Let $D = \img{h}{F}$.
                $D$ is finite by \cref{finite_image_function}.
                For all $r$ such that $r\in D$ we have $r\in\reals$
                    by \cref{img_iff,subseteq,times_tuple_intro,metric_on,funs_type_apply,square_metric_is_metric,pair_eq_iff}.
                Hence $D\subseteq\reals$ by \cref{subseteq}.
                Take $f_1$ such that $f_1\in F$.
                Thus $(f_1,\rho((a_0,f_1)))\in h$.
                Therefore $D\neq\emptyset$ by \cref{img_iff,nonempty_inhabited,inhabited_nonempty}.
                Take $M_0\in D$ such that $M_0$ is a maximum of $D$ by \cref{finite_real_maximum}.
                Let $M = (M_0 + 1) + (M_0 + 1)$.
                Show that for all $x,y\in A$ we have $\rho((x,y)) \leq M$.
                \begin{subproof}
                    Fix $x$.
                    Fix $y$ such that $x\in A$ and $y\in A$.
                    Take $U_x\in B$ such that $x\in U_x$ by \cref{covering,subseteq,unions_iff}.
                    Take $U_y\in B$ such that $y\in U_y$ by \cref{covering,subseteq,unions_iff}.
                    Let $f_x = g(U_x)$.
                    Let $f_y = g(U_y)$.
                    $U_x\mathrel{R} f_x$ by \cref{function_apply_elim}.
                    $U_y\mathrel{R} f_y$ by \cref{function_apply_elim}.
                    We have $(U_x,f_x)\in R$.
                    Therefore $(U_x,f_x)\in B\times A$.
                    $\fst{(U_x,f_x)} = \metricBall{\rho}{X}{\snd{(U_x,f_x)}}{1}$.
                    Hence $U_x = \metricBall{\rho}{X}{f_x}{1}$ by \cref{fst_eq,snd_eq}.
                    We have $(U_y,f_y)\in R$.
                    Therefore $(U_y,f_y)\in B\times A$.
                    $\fst{(U_y,f_y)} = \metricBall{\rho}{X}{\snd{(U_y,f_y)}}{1}$.
                    Hence $U_y = \metricBall{\rho}{X}{f_y}{1}$ by \cref{fst_eq,snd_eq}.
                    Hence $x\in \metricBall{\rho}{X}{f_x}{1}$.
                    Thus $x\in X$ and $\rho((f_x,x)) < 1$ by \cref{metric_ball}.
                    Thus $f_x\in A$ and $f_x\in X$ by \cref{subseteq,times_tuple_elim}.
                    Therefore $\rho((x,f_x)) < 1$
                        by \cref{metric_symmetric,metric_on,square_metric_is_metric,real_lt_subst_left}.
                    Hence $y\in \metricBall{\rho}{X}{f_y}{1}$.
                    Thus $y\in X$ and $\rho((f_y,y)) < 1$ by \cref{metric_ball}.
                    Thus $f_y\in A$ and $f_y\in X$ by \cref{subseteq,times_tuple_elim}.
                    Therefore $\rho((y,f_y)) < 1$
                        by \cref{metric_symmetric,metric_on,square_metric_is_metric,real_lt_subst_left}.
                    Hence $f_x\in F$ and $f_y\in F$ by \cref{subseteq,function_apply_elim,img_iff}.
                    We have $h(f_x)\in D$ and $h(f_y)\in D$ by \cref{subseteq,function_apply_elim,img_iff}.
                    Therefore $h(f_x) \leq M_0$ and $h(f_y) \leq M_0$ by \cref{real_maximum}.
                    Thus $\rho((a_0,f_x)) \leq M_0$ and $\rho((a_0,f_y)) \leq M_0$ by \cref{function_apply_intro}.
                    Show that $\rho((x,y)) \leq \rho((x,f_x)) + \rho((f_x,a_0)) + \rho((a_0,f_y)) + \rho((f_y,y))$.
                    \begin{subproof}
                        $x\in X$ and $y\in X$ and $f_x\in X$ and $f_y\in X$ and $a_0\in X$ by \cref{subseteq}.
                        Therefore $\rho((x,y)) \leq \rho((x,f_x)) + \rho((f_x,y))$
                            by \cref{square_metric_is_metric,metric_on,metric_triangle_inequality}.
                        Hence $\rho((f_x,y)) \leq \rho((f_x,a_0)) + \rho((a_0,y))$
                            by \cref{square_metric_is_metric,metric_on,metric_triangle_inequality}.
                        Therefore $\rho((a_0,y)) \leq \rho((a_0,f_y)) + \rho((f_y,y))$
                            by \cref{square_metric_is_metric,metric_on,metric_triangle_inequality}.
                        We have $x\in X$ and $f_x\in A$ and $f_x\in X$ and $f_y\in A$ by \cref{elem_subseteq}.
                        Therefore $f_y\in X$ and $y\in X$ and $a_0\in A$ and $a_0\in X$ by \cref{elem_subseteq}.
                        $(x,f_x)\in X\times X$ by \cref{times_tuple_intro}.
                        $(f_x,y)\in X\times X$ by \cref{times_tuple_intro}.
                        $(f_x,a_0)\in X\times X$ by \cref{times_tuple_intro}.
                        $(a_0,y)\in X\times X$ by \cref{times_tuple_intro}.
                        $(a_0,f_y)\in X\times X$ by \cref{times_tuple_intro}.
                        $(f_y,y)\in X\times X$ by \cref{times_tuple_intro}.
                        $(x,y)\in X\times X$ by \cref{times_tuple_intro}.
                        We have $\rho((x,y))\in\reals$ by \cref{metric_on,funs_type_apply}.
                        $\rho((x,f_x))\in\reals$ by \cref{metric_on,funs_type_apply}.
                        $\rho((f_x,y))\in\reals$ by \cref{metric_on,funs_type_apply}.
                        $\rho((f_x,a_0))\in\reals$ by \cref{metric_on,funs_type_apply}.
                        $\rho((a_0,y))\in\reals$ by \cref{metric_on,funs_type_apply}.
                        $\rho((a_0,f_y))\in\reals$ by \cref{metric_on,funs_type_apply}.
                        $\rho((f_y,y))\in\reals$ by \cref{metric_on,funs_type_apply}.
                        Let $r_1 = \rho((x,f_x))$.
                        Let $r_2 = \rho((f_x,y))$.
                        Let $r_3 = \rho((f_x,a_0))$.
                        Let $r_4 = \rho((a_0,y))$.
                        We have $r_1\in\reals$ and $r_2\in\reals$ and $r_3\in\reals$ and $r_4\in\reals$.
                        $r_3 + r_4\in\reals$ by \cref{real_add_closed}.
                        $r_1 + r_2\in\reals$ by \cref{real_add_closed}.
                        $r_1 + (r_3 + r_4)\in\reals$ by \cref{real_add_closed}.
                        $r_2 \leq r_3 + r_4$.
                        Hence $r_2 + r_1 \leq (r_3 + r_4) + r_1$ by \cref{real_leq_add}.
                        We have $r_1 + r_2 = r_2 + r_1$ by \cref{real_add_comm}.
                        $(r_3 + r_4) + r_1 = r_1 + (r_3 + r_4)$ by \cref{real_add_comm}.
                        Therefore $r_1 + r_2 \leq (r_3 + r_4) + r_1$ by \cref{real_leq_trans}.
                        Hence $r_1 + r_2 \leq r_1 + (r_3 + r_4)$ by \cref{real_leq_trans}.
                        Thus $\rho((x,y)) \leq r_1 + (r_3 + r_4)$ by \cref{real_leq_trans}.
                        Let $u = \rho((a_0,f_y)) + \rho((f_y,y))$.
                        We have $u\in\reals$ by \cref{real_add_closed}.
                        $r_1 + r_3\in\reals$ by \cref{real_add_closed}.
                        $(r_1 + r_3) + u\in\reals$ by \cref{real_add_closed}.
                        $r_4 \leq u$.
                        Therefore $r_4 + (r_1 + r_3) \leq u + (r_1 + r_3)$ by \cref{real_leq_add}.
                        Thus $(r_1 + r_3) + r_4 = r_4 + (r_1 + r_3)$ by \cref{real_add_comm}.
                        Thus $u + (r_1 + r_3) = (r_1 + r_3) + u$ by \cref{real_add_comm}.
                        Hence $(r_1 + r_3) + r_4 \leq u + (r_1 + r_3)$ by \cref{real_leq_trans}.
                        Thus $(r_1 + r_3) + r_4 \leq (r_1 + r_3) + u$ by \cref{real_leq_trans}.
                        $r_1 + (r_3 + r_4) = (r_1 + r_3) + r_4$ by \cref{real_add_assoc}.
                        Therefore $r_1 + (r_3 + r_4) \leq (r_1 + r_3) + u$ by \cref{real_leq_trans}.
                        Hence $\rho((x,y)) \leq (r_1 + r_3) + u$ by \cref{real_leq_trans}.
                        Therefore $\rho((x,y)) \leq \rho((x,f_x)) + \rho((f_x,a_0)) + \rho((a_0,f_y)) + \rho((f_y,y))$
                            by \cref{real_add_assoc,real_add_comm}.
                    \end{subproof}
                    Show that $\rho((x,y)) \leq M$.
                    \begin{subproof}
                        Let $s_1 = \rho((x,f_x))$.
                        Let $s_2 = \rho((f_x,a_0))$.
                        Let $s_3 = \rho((a_0,f_y))$.
                        Let $s_4 = \rho((f_y,y))$.
                        We have $x\in X$ and $y\in X$ and $f_x\in A$ and $f_x\in X$ and $f_y\in A$ and $f_y\in X$ and $a_0\in A$ and $a_0\in X$
                            by \cref{elem_subseteq}.
                        $(x,y)\in X\times X$ by \cref{times_tuple_intro}.
                        $(x,f_x)\in X\times X$ by \cref{times_tuple_intro}.
                        $(f_x,a_0)\in X\times X$ by \cref{times_tuple_intro}.
                        $(a_0,f_y)\in X\times X$ by \cref{times_tuple_intro}.
                        $(f_y,y)\in X\times X$ by \cref{times_tuple_intro}.
                        We have $s_1\in\reals$ and $s_2\in\reals$ and $s_3\in\reals$ and $s_4\in\reals$ by \cref{metric_on,funs_type_apply}.
                        $s_1 + s_2\in\reals$ and $s_3 + s_4\in\reals$ by \cref{real_add_closed}.
                        $1\in\reals$ by \cref{real_naturals_subset_reals,real_one_in_naturals,subseteq}.
                        Hence $s_2 = \rho((a_0,f_x))$ by \cref{metric_symmetric,metric_on,square_metric_is_metric}.
                        Therefore $s_1 \leq 1$ by \cref{real_lt_imp_leq}.
                        Hence $s_4 \leq 1$ by \cref{metric_symmetric,metric_on,square_metric_is_metric,real_lt_subst_left,real_lt_imp_leq}.
                        We have $s_2 \leq M_0$ and $s_3 \leq M_0$.
                        Therefore $M_0\in\reals$ by \cref{real_maximum,subseteq}.
                        $M_0 + 1\in\reals$ by \cref{real_add_closed}.
                        $(s_1 + s_2) + (s_3 + s_4)\in\reals$ by \cref{real_add_closed}.
                        $(M_0 + 1) + (M_0 + 1)\in\reals$ by \cref{real_add_closed}.
                        We have $\rho((x,y))\in\reals$ by \cref{metric_on,funs_type_apply}.
                        $\rho((x,y)) \leq s_1 + s_2 + s_3 + s_4$.
                        Hence $s_1 + s_2 + s_3 + s_4 = (s_1 + s_2) + (s_3 + s_4)$ by \cref{real_add_assoc}.
                        Therefore $\rho((x,y)) \leq (s_1 + s_2) + (s_3 + s_4)$ by \cref{real_leq_trans}.
                        Hence $(s_1 + s_2) + (s_3 + s_4) \leq (M_0 + 1) + (M_0 + 1)$ by \cref{real_four_sum_bound}.
                        Thus $\rho((x,y)) \leq (M_0 + 1) + (M_0 + 1)$ by \cref{real_leq_trans}.
                        Therefore $\rho((x,y)) \leq M$.
                    \end{subproof}
                \end{subproof}
                Hence $M\in\reals$ by \cref{real_maximum,subseteq,real_naturals_subset_reals,real_one_in_naturals,real_add_closed}.
                Therefore $A$ is bounded in $X$ with respect to $\rho$ by \cref{metric_bounded}.
            \end{byCase}
        \end{subproof}
        Show that $A$ is bounded in $X$ with respect to $d$ or $A$ is bounded in $X$ with respect to $\rho$.
        \begin{subproof}
            Trivial.
        \end{subproof}
    \end{subproof}
    Show that if
        $A$ is closed in $X$ with respect to $T$ and
        $A$ is bounded in $X$ with respect to $d$ or $A$ is bounded in $X$ with respect to $\rho$
        then $A$ is compact with respect to $T_A$.
    \begin{subproof}
        Assume $A$ is closed in $X$ with respect to $T$ and
            $A$ is bounded in $X$ with respect to $d$ or $A$ is bounded in $X$ with respect to $\rho$.
        \begin{byCase}
            \caseOf{$A = \emptyset$.}
                Show that for all $C$ such that $C$ is an open covering of $A$ with respect to $T_A$ there exists $B$ such that
                    $B\subseteq C$ and $B$ is finite and $B$ is a covering of $A$.
                \begin{subproof}
                    Follows by \cref{emptyset_is_finite,emptyset_subseteq,covering}.
                \end{subproof}
            Therefore $A$ is compact with respect to $T_A$ by \cref{compact}.
        \caseOf{$A \neq \emptyset$.}
            $A$ is inhabited by \cref{nonempty_inhabited}.
            Show that $A$ is bounded in $X$ with respect to $\rho$.
            \begin{subproof}
                \begin{byCase}
                \caseOf{$A$ is bounded in $X$ with respect to $\rho$.}
                    Show that $A$ is bounded in $X$ with respect to $\rho$.
                    \begin{subproof}
                        Trivial.
                    \end{subproof}
                \caseOf{$A$ is bounded in $X$ with respect to $d$.}
                    Take $M\in\reals$ such that for all $a_1,a_2\in A$ we have $d((a_1,a_2)) \leq M$
                        by \cref{metric_bounded}.
                    For all $a_1,a_2\in A$ we have $\rho((a_1,a_2)) \leq M$
                        by \cref{elem_subseteq,euclidean_metric_value_exists,square_metric_value_exists,times_tuple_intro,euclidean_square_metric_bounds,euclidean_metric_function,square_metric_function,function_apply_intro,real_leq_trans}.
                    Hence $A$ is bounded in $X$ with respect to $\rho$ by \cref{square_metric_is_metric,metric_bounded}.
                \end{byCase}
            \end{subproof}
            Take $M\in\reals$ such that for all $a_1,a_2\in A$ we have $\rho((a_1,a_2)) \leq M$
                by \cref{metric_bounded}.
            Take $x_0$ such that $x_0\in A$.
            Let $Y = \{(i,\intervalclosed{x_0(i) - M}{x_0(i) + M}) \mid i\in J\}$.
            Show that for all $i,u,v$ such that $(i,u)\in Y$ and $(i,v)\in Y$ we have $u = v$.
            \begin{subproof}
                Trivial.
            \end{subproof}
            Hence $Y$ is a function by \cref{relation,rightunique}.
            We have $\dom{Y} = J$ by \cref{dom_intro,dom_iff,pair_eq_iff,subseteq,subseteq_antisymmetric}.
            Show that for all $x$ such that $x\in A$ we have $x\in \productFamily{Y}$.
            \begin{subproof}
                Fix $x$ such that $x\in A$.
                Show that for all $i\in J$ we have $x(i)\in \intervalclosed{x_0(i) - M}{x_0(i) + M}$.
                \begin{subproof}
                    Fix $i$ such that $i\in J$.
                    We have $x_0\in\realsPower{J}$ and $x\in\realsPower{J}$ by \cref{elem_subseteq}.
                    $x_0(i)\in\reals$ and $x(i)\in\reals$ by \cref{reals_power,tuple_power,funs_type_apply}.
                    $\neg{x(i)}\in\reals$ by \cref{real_add_inverse}.
                    $x_0(i) - x(i) = x_0(i) + \neg{x(i)}$ by \cref{real_sub}.
                    Therefore $x_0(i) - x(i)\in\reals$ by \cref{real_add_closed}.
                    Hence $\abs{x_0(i) - x(i)}\in\reals$ by \cref{real_abs_in_reals}.
                    Take $m\in\reals$ such that $((x_0,x),m)\in \rho$ by \cref{square_metric_value_exists}.
                    We have $m\in\coordDiffSet{n}{x_0}{x}$ and for all $t\in\coordDiffSet{n}{x_0}{x}$ we have $t \leq m$
                        by \cref{square_metric,pair_eq_iff}.
                    Therefore $\abs{x_0(i) - x(i)} \leq m$ by \cref{coord_diff_set}.
                    Hence $m \leq M$ by \cref{square_metric_function,function_apply_intro,real_leq_trans}.
                    Therefore $\abs{x_0(i) - x(i)} \leq M$ by \cref{real_leq_trans}.
                    Hence $\abs{x(i) - x_0(i)} \leq M$ by \cref{real_abs_sub_swap,real_leq_trans}.
                    Therefore $M$ is nonnegative by \cref{real_abs_nonnegative,real_add_ident,real_leq_trans}.
                    Hence $x_0(i) - M \leq x(i)$ by \cref{real_abs_sub_leq_imp_left_bound}.
                    Therefore $x(i) \leq x_0(i) + M$ by \cref{real_abs_sub_leq_imp_right_bound}.
                    Hence $x(i)\in\intervalclosed{x_0(i) - M}{x_0(i) + M}$ by \cref{intervalclosed}.
                \end{subproof}
                Thus $x$ is a function to $\reals$ and $\dom{x} = \dom{Y}$ by \cref{subseteq,reals_power,tuple_power,funs_elim}.
                For all $i\in\dom{Y}$ we have $x(i)\in\unions{\ran{Y}}$ by \cref{subseteq,ran_intro,unions_iff}.
                Hence $x\in\funs{\dom{Y}}{\unions{\ran{Y}}}$ by \cref{funs_intro}.
                For all $i\in\dom{Y}$ we have $x(i)\in\apply{Y}{i}$ by \cref{dom_iff,pair_eq_iff,function_apply_intro}.
                Hence $x\in\productFamily{Y}$ by \cref{indexed_product}.
            \end{subproof}
            Hence $A\subseteq \productFamily{Y}$ by \cref{subseteq}.
            Let $T_Y = \{(i,\subspaceTopology{\standardTopologyR}{\intervalclosed{x_0(i) - M}{x_0(i) + M}}) \mid i\in J\}$.
            Show that for all $i,u,v$ such that $(i,u)\in T_Y$ and $(i,v)\in T_Y$ we have $u = v$.
            \begin{subproof}
                Trivial.
            \end{subproof}
            Therefore $T_Y$ is a function by \cref{relation,rightunique}.
            Hence $\dom{T_Y} = J$ by \cref{dom_intro,dom_iff,pair_eq_iff,subseteq,subseteq_antisymmetric}.
            % product of finitely many compact intervals
            Show that for all $i\in J$ we have $\intervalclosed{x_0(i) - M}{x_0(i) + M}$ is compact
                with respect to $\subspaceTopology{\standardTopologyR}{\intervalclosed{x_0(i) - M}{x_0(i) + M}}$.
            \begin{subproof}
                Fix $i$ such that $i\in J$.
                We have $x_0(i)\in\reals$ by \cref{elem_subseteq,reals_power,tuple_power,funs_type_apply}.
                $0\in\reals$ by \cref{real_add_ident}.
                $\rho$ is a metric on $X$ by \cref{square_metric_is_metric}.
                Hence $\rho$ is nonnegative on $X$ by \cref{metric_on}.
                Therefore $x_0\in X$ by \cref{elem_subseteq}.
                $(x_0,x_0)\in X\times X$ by \cref{times_tuple_intro}.
                $\rho((x_0,x_0))\in\reals$ by \cref{metric_on,funs_type_apply}.
                We have $\rho((x_0,x_0)) \geq 0$ by \cref{metric_nonnegative}.
                $\rho((x_0,x_0)) \leq M$.
                Therefore $0 \leq M$ by \cref{real_leq_trans}.
                $\neg{M}\in\reals$ and $\neg{0}\in\reals$ by \cref{real_add_inverse}.
                $0 + \neg{0} = 0$ and $0 + \neg{0} = \neg{0}$ by \cref{real_add_inverse,real_add_ident}.
                We have $\neg{0} = 0$.
                Therefore $\neg{M} \leq 0$ by \cref{real_leq_neg,real_leq_trans}.
                Hence $\neg{M} \leq M$ by \cref{real_leq_trans}.
                Therefore $\neg{M} + x_0(i) \leq M + x_0(i)$ by \cref{real_leq_add}.
                $x_0(i) + \neg{M} = \neg{M} + x_0(i)$ and $x_0(i) + M = M + x_0(i)$ by \cref{real_add_comm}.
                Hence $x_0(i) + \neg{M} \leq x_0(i) + M$ by \cref{real_leq_trans}.
                $x_0(i) + M\in\reals$ by \cref{real_add_closed}.
                $x_0(i) - M = x_0(i) + \neg{M}$ by \cref{real_sub}.
                $x_0(i) - M\in\reals$ by \cref{real_add_closed}.
                Therefore $x_0(i) - M \leq x_0(i) + M$ by \cref{real_sub}.
                Hence $\intervalclosed{x_0(i) - M}{x_0(i) + M}$ is compact
                    with respect to $\subspaceTopology{\standardTopologyR}{\intervalclosed{x_0(i) - M}{x_0(i) + M}}$
                    by \cref{real_intervalclosed_compact_standard}.
            \end{subproof}
            For all $i\in J$ we have $\apply{T_Y}{i}$ is a topology on $\apply{Y}{i}$
                by \cref{subseteq,function_apply_intro,standard_basis_is_basis,basis_topology_is_topology,standard_topology_reals,intervalclosed,subspace_topology_is_topology}.
            For all $i\in J$ we have $\apply{Y}{i}$ is compact with respect to $\apply{T_Y}{i}$
                by \cref{subseteq,function_apply_intro}.
            Therefore $\productFamily{Y}$ is compact with respect to $\productTopologyIndexed{T_Y}{Y}$ by \cref{compact_indexed_product_initial_segment}.
            Let $T_{P0} = \productTopologyIndexed{T_Y}{Y}$.
            Let $T_{Y0} = \subspaceTopology{T}{\productFamily{Y}}$.
            Let $X_1 = \{(i,\reals) \mid i\in J\}$.
            Let $T_1 = \{(i,\standardTopologyR) \mid i\in J\}$.
            Let $T_0 = \productTopologyIndexed{T_1}{X_1}$.
            Hence $T = T_0$ by \cref{euclidean_metric_product_topology,square_metric_topology_equals_euclidean}.
            $X_1$ is a function and $T_1$ is a function and $\dom{X_1} = J$ and $\dom{T_1} = J$
                by \cref{reals_standard_indexed_families_data}.
            Therefore $\dom{X_1}$ is inhabited by \cref{real_one_in_naturals,real_one_leq_naturalsplus,initial_segment_naturalsplus,subseteq}.
            For all $i\in\dom{X_1}$ we have $\apply{Y}{i}\subseteq\apply{X_1}{i}$
                by \cref{subseteq,function_apply_intro,intervalclosed,subseteq_transitive}.
            For all $i\in\dom{X_1}$ we have $\apply{T_1}{i}$ is a topology on $\apply{X_1}{i}$
                by \cref{subseteq,function_apply_intro,standard_basis_is_basis,basis_topology_is_topology,standard_topology_reals}.
            Hence $\productFamily{Y}$ is a subspace of $\productFamily{X_1}$ with respect to $T$ by \cref{product_subspace_product}.
            Let $j = \identity{\productFamily{Y}}$.
            Therefore $\dom{Y}$ is inhabited by \cref{real_one_in_naturals,real_one_leq_naturalsplus,initial_segment_naturalsplus,subseteq}.
            For all $i\in\dom{Y}$ we have $\apply{T_Y}{i}$ is a topology on $\apply{Y}{i}$
                by \cref{subseteq,function_apply_intro,standard_basis_is_basis,basis_topology_is_topology,standard_topology_reals,intervalclosed,subspace_topology_is_topology}.
            Therefore $T_{P0}$ is a topology on $\productFamily{Y}$
                by \cref{product_subbasis_finite_intersections_basis,basis_for_topology,basis_topology_is_topology}.
            Hence $\productFamily{X_1} = X$ by \cref{product_family_reals_power}.
            We have $\productFamily{Y}\subseteq X$ by \cref{subspace_of,subseteq_transitive}.
            $j\in\funs{\productFamily{Y}}{X}$ by \cref{id_elem_funs,funs_weaken_codom}.
            Show that for all $i\in\dom{X_1}$ we have $\piIndex{X_1}{i}\circ j$ is continuous from $\productFamily{Y}$
                to $\apply{X_1}{i}$ with respect to $T_{P0}$ and $\apply{T_1}{i}$.
            \begin{subproof}
                Fix $i$ such that $i\in\dom{X_1}$.
                Let $g = \piIndex{Y}{i}$.
                Let $k = \identity{\apply{Y}{i}}$.
                $k$ is continuous from $\apply{Y}{i}$ to $\apply{X_1}{i}$ with respect to $\apply{T_Y}{i}$ and $\apply{T_1}{i}$
                    by \cref{subseteq,subspace_of,function_apply_intro,inclusion_continuous}.
                $g$ is continuous from $\productFamily{Y}$ to $\apply{Y}{i}$ with respect to $T_{P0}$ and $\apply{T_Y}{i}$
                    by \cref{projection_map_indexed_continuous}.
                Therefore $k\circ g$ is continuous from $\productFamily{Y}$ to $\apply{X_1}{i}$ with respect to $T_{P0}$ and $\apply{T_1}{i}$
                    by \cref{continuous_composite}.
                Show that $k\circ g = \piIndex{X_1}{i}\circ j$.
                \begin{subproof}
                    Show that for all $y,z$ if $y\mathrel{k\circ g} z$ then $y\mathrel{\piIndex{X_1}{i}\circ j} z$.
                    \begin{subproof}
                        Follows by \cref{circ_iff,id_iff,projection_map_indexed,pair_eq_iff,subspace_of,subseteq}.
                    \end{subproof}
                    Show that for all $y,z$ if $y\mathrel{\piIndex{X_1}{i}\circ j} z$ then $y\mathrel{k\circ g} z$.
                    \begin{subproof}
                        Follows by \cref{circ_iff,id_iff,projection_map_indexed,pair_eq_iff,subseteq,indexed_product}.
                    \end{subproof}
                    Follows by set extensionality.
                \end{subproof}
                Hence $\piIndex{X_1}{i}\circ j$ is continuous from $\productFamily{Y}$ to $\apply{X_1}{i}$ with respect to $T_{P0}$ and $\apply{T_1}{i}$
                    by \cref{continuous_eq_subst}.
            \end{subproof}
            Therefore $j$ is continuous from $\productFamily{Y}$ to $\productFamily{X_1}$ with respect to $T_{P0}$ and $T_0$
                by \cref{continuous_map_into_indexed_product}.
            Hence $j$ is continuous from $\productFamily{Y}$ to $X$ with respect to $T_{P0}$ and $T$ by \cref{continuous_eq_codomain}.
            Show that for all $U$ such that $U\in T_{Y0}$ we have $U\in T_{P0}$.
            \begin{subproof}
                Fix $U$ such that $U\in T_{Y0}$.
                Take $V\in T$ such that $U = \productFamily{Y}\inter V$ by \cref{subspace_topology}.
                Let $j_0 = \identity{\productFamily{Y}}$.
                Thus $V$ is open in $X$ with respect to $T$ by \cref{continuous,open_in}.
                Hence $\preimg{j_0}{V}$ is open in $\productFamily{Y}$ with respect to $T_{P0}$ by \cref{continuous}.
                Show that $\preimg{j_0}{V} = \productFamily{Y}\inter V$.
                \begin{subproof}
                    Show that for all $x$ such that $x\in\preimg{j_0}{V}$ we have $x\in\productFamily{Y}\inter V$.
                    \begin{subproof}
                        Fix $x$ such that $x\in\preimg{j_0}{V}$.
                        Take $v$ such that $v\in V$ and $(x,v)\in j_0$ by \cref{preimg_iff}.
                        Thus $x = v$ and $v\in\productFamily{Y}$ by \cref{id_iff}.
                        Hence $x\in\productFamily{Y}$ and $x\in V$.
                        Thus $x\in\productFamily{Y}\inter V$ by \cref{inter_intro}.
                    \end{subproof}
                    Hence $\preimg{j_0}{V}\subseteq \productFamily{Y}\inter V$ by \cref{subseteq}.
                    Show that for all $x$ such that $x\in\productFamily{Y}\inter V$ we have $x\in\preimg{j_0}{V}$.
                    \begin{subproof}
                        Fix $x$ such that $x\in\productFamily{Y}\inter V$.
                        Thus $x\in\productFamily{Y}$ and $x\in V$ by \cref{inter}.
                        Hence $(x,x)\in j_0$ by \cref{id_iff}.
                        Thus $x\in\preimg{j_0}{V}$ by \cref{preimg_iff}.
                    \end{subproof}
                    Hence $\productFamily{Y}\inter V\subseteq \preimg{j_0}{V}$ by \cref{subseteq}.
                    Follows by \cref{subseteq_antisymmetric}.
                \end{subproof}
                Hence $U\in T_{P0}$ by \cref{open_in}.
            \end{subproof}
            Hence $T_{Y0}\subseteq T_{P0}$ by \cref{subseteq}.
            For all $C$ such that $C$ is an open covering of $\productFamily{Y}$ with respect to $T_{Y0}$ there exists $B$ such that
                $B\subseteq C$ and $B$ is finite and $B$ is a covering of $\productFamily{Y}$
                by \cref{open_covering,open_in,subseteq,compact}.
            Therefore $\productFamily{Y}$ is compact with respect to $T_{Y0}$ by \cref{compact}.
            We have $T$ is a topology on $X$ by \cref{square_metric_is_metric,metric_basis_is_basis,basis_topology_is_topology,metric_topology,basis_for_topology}.
            $\productFamily{Y}\subseteq X$
                by \cref{indexed_product,funs_elim,subseteq,function_apply_intro,intervalclosed,funs_intro,tuple_power,reals_power}.
            Thus $\productFamily{Y}$ is a subspace of $X$ with respect to $T$ by \cref{subspace_of}.
            We have $A = A\inter\productFamily{Y}$ by \cref{subseteq,inter_absorb_supseteq_right}.
            Thus $A$ is closed in $\productFamily{Y}$ with respect to $T_{Y0}$ by \cref{closed_in_subspace_iff,subspace_of}.
            Therefore $T$ is a topology on $X$ and $\productFamily{Y}\subseteq X$ by \cref{subspace_of}.
            Thus $T_{Y0}$ is a topology on $\productFamily{Y}$ by \cref{subspace_topology_is_topology}.
            Hence $A$ is compact with respect to $T_A$ by \cref{closed_subspace_compact,subspace_subspace_topology}.
        \end{byCase}
    \end{subproof}
\end{proof}
% from §27 helper lemma 10: finite inhabited subset has a largest element
\begin{lemma}\label{finite_inhabited_largest_element}
    Assume $R$ is a simple order relation on $X$.
    Suppose $F\subseteq X$.
    Suppose $F$ is finite.
    Suppose $F$ is inhabited.
    Then there exists $b$ such that $b$ is a largest element of $F$ with respect to $R$.
\end{lemma}
\begin{proof}
    Take $k\in\naturalsPlus$ such that there exists a bijection from $\initialSegment{k}$ to $F$ by \cref{finite}.
    Take $f$ such that $f$ is a bijection from $\initialSegment{k}$ to $F$.
    Let $A = \{ n\in\naturalsPlus \mid \text{for all $G,g$ such that $g$ is a bijection from $\initialSegment{n}$ to $G$ and $G\subseteq X$ and $G$ is inhabited
        there exists $b$ such that $b$ is a largest element of $G$ with respect to $R$}\}$.
    Show that for all $G,g$ such that $g$ is a bijection from $\initialSegment{1}$ to $G$ and $G\subseteq X$ and $G$ is inhabited
        we have there exists $b$ such that $b$ is a largest element of $G$ with respect to $R$.
    \begin{subproof}
        Fix $G$.
        Fix $g$ such that $g$ is a bijection from $\initialSegment{1}$ to $G$ and $G\subseteq X$ and $G$ is inhabited.
        $g$ is a function by \cref{bijection_is_function}.
        $\dom{g} = \initialSegment{1}$ by \cref{bijections_dom}.
        $1\in\dom{g}$ by \cref{initial_segment_one,singleton_intro}.
        Let $a = g(1)$.
        $(1,a)\in g$ by \cref{function_apply_elim}.
        $g\in\surj{\initialSegment{1}}{G}$ by \cref{bijections}.
        $\ran{g} = G$ by \cref{ran_of_surj}.
        $a\in G$ by \cref{ran_iff}.
        Show that for all $z\in G$ we have $z\in\{a\}$.
        \begin{subproof}
            Fix $z\in G$.
            Take $x\in\initialSegment{1}$ such that $g(x)=z$ by \cref{bijections,surj}.
            $x = 1$ by \cref{initial_segment_one,singleton_iff}.
            $a = z$.
            Thus $z\in\{a\}$ by \cref{singleton_iff}.
        \end{subproof}
        $G\subseteq\{a\}$ by \cref{subseteq}.
        $\{a\}\subseteq G$ by \cref{singleton_subset_intro}.
        $G = \{a\}$ by \cref{subseteq_antisymmetric}.
        $a$ is a largest element of $G$ with respect to $R$ by \cref{singleton_iff,largest_element}.
    \end{subproof}
    $1\in A$.
    Show that for all $n\in A$ we have $n + 1\in A$.
    \begin{subproof}
        Fix $n\in A$.
        Show that for all $G,g$ such that $g$ is a bijection from $\initialSegment{n + 1}$ to $G$ and $G\subseteq X$ and $G$ is inhabited
            there exists $b$ such that $b$ is a largest element of $G$ with respect to $R$.
        \begin{subproof}
            Fix $G$.
            Fix $g$ such that $g$ is a bijection from $\initialSegment{n + 1}$ to $G$ and $G\subseteq X$ and $G$ is inhabited.
            $g$ is a function by \cref{bijection_is_function}.
            $\dom{g} = \initialSegment{n + 1}$ by \cref{bijections_dom}.
            $n + 1\in\dom{g}$ by \cref{initial_segment_naturalsplus,real_naturals_closed_succ}.
            Let $a = g(n + 1)$.
            Let $G_1 = \img{g}{\initialSegment{n}}$.
            $\initialSegment{n}\subseteq \initialSegment{n + 1}$ by \cref{initial_segment_subset_succ}.
            $G_1\subseteq G$ and $G_1\subseteq X$
                by \cref{subseteq,img_subseteq,img_dom,bijections,ran_of_surj,subseteq_transitive}.
            $g\in\surj{\initialSegment{n + 1}}{G}$ by \cref{bijections}.
            $\ran{g} = G$ by \cref{ran_of_surj}.
            $(n + 1,a)\in g$ by \cref{function_apply_elim}.
            $a\in G$ by \cref{function_apply_elim,ran_iff,bijections,ran_of_surj}.
            Show that for all $z$ such that $z\in G$ we have $z\in G_1\union\{a\}$.
            \begin{subproof}
                Fix $z\in G$.
                There exists $x\in\initialSegment{n + 1}$ such that $g(x)=z$ by \cref{bijections,surj}.
                $x\in\initialSegment{n}$ or $x = n + 1$ by \cref{initial_segment_succ_split}.
                \begin{byCase}
                \caseOf{$x\in\initialSegment{n}$.}
                    $(x,z)\in g$ by \cref{function_apply_elim}.
                    $z\in G_1$ by \cref{img_iff}.
                    $z\in G_1\union\{a\}$ by \cref{union_intro_left}.
                \caseOf{$x = n + 1$.}
                    $z = a$.
                    $z\in G_1\union\{a\}$ by \cref{union_intro_right,singleton_intro}.
                \end{byCase}
            \end{subproof}
            $G\subseteq G_1\union\{a\}$ by \cref{subseteq}.
            $G_1\union\{a\}\subseteq G$ by \cref{singleton_subset_intro,union_subsets_is_subset}.
            $G = G_1\union\{a\}$ by \cref{subseteq_antisymmetric}.
            Let $h = \restrl{g}{\initialSegment{n}}$.
            $h$ is a function by \cref{restrl_of_function_is_function}.
            For all $x\in\initialSegment{n}$ we have $x\in\dom{g}$ by \cref{initial_segment_subset_succ,bijections_dom,subseteq}.
            $\dom{h} = \initialSegment{n}$ by \cref{subseteq,dom_of_restrl_eq_inter,inter_absorb_supseteq_left}.
            For all $x\in\dom{h}$ we have $h(x)\in G_1$
                by \cref{function_apply_elim,restrl_subseteq,elem_subseteq,img_iff}.
            $h\in\funs{\initialSegment{n}}{G_1}$ by \cref{funs_intro}.
            For all $y\in G_1$ there exists $x\in\initialSegment{n}$ such that $h(x) = y$
                by \cref{img_iff,function_apply_iff,restrl_of_function_apply}.
            $h\in\surj{\initialSegment{n}}{G_1}$ by \cref{surj}.
            $h$ is injective by \cref{restrl_of_function_apply,bijections,injective_function}.
            $h$ is a bijection from $\initialSegment{n}$ to $G_1$ by \cref{bijections}.
            Take $x_0$ such that $x_0\in\initialSegment{n}$ by \cref{initial_segment_inhabited}.
            $x_0\in\dom{g}$ by \cref{initial_segment_subset_succ,bijections_dom,subseteq}.
            $g(x_0)\in G_1$ by \cref{function_apply_elim,img_iff}.
            $G_1$ is inhabited.
                Take $b$ such that $b$ is a largest element of $G_1$ with respect to $R$.
                $b\mathrel{R}a$ or $a\mathrel{R}b$ or $b = a$
                    by \cref{largest_element,elem_subseteq,simple_order_relation,connex}.
                \begin{byCase}
                \caseOf{$b\mathrel{R}a$ or $b = a$.}
                    For all $y\in G$ we have $y\in G_1$ or $y = a$ by \cref{union_iff,singleton_iff}.
                    For all $y\in G_1$ we have $y\mathrel{R}b$ or $y = b$ by \cref{largest_element}.
                    For all $y\in G$ we have $y\mathrel{R}a$ or $y = a$
                        by \cref{simple_order_relation,transitive}.
                    $a$ is a largest element of $G$ with respect to $R$
                        by \cref{union_intro_right,singleton_intro,largest_element}.
                \caseOf{$a\mathrel{R}b$.}
                    For all $y\in G$ we have $y\in G_1$ or $y = a$ by \cref{union_iff,singleton_iff}.
                    For all $y\in G$ we have $y\mathrel{R}b$ or $y = b$ by \cref{largest_element}.
                    $b$ is a largest element of $G$ with respect to $R$ by \cref{elem_subseteq,largest_element}.
                \end{byCase}
        \end{subproof}
    \end{subproof}
    $A\subseteq\naturalsPlus$ by \cref{subseteq}.
    $k\in A$ by \cref{real_naturals_induction}.
    For all $G,g$ such that $g$ is a bijection from $\initialSegment{k}$ to $G$ and $G\subseteq X$ and $G$ is inhabited
        there exists $b$ such that $b$ is a largest element of $G$ with respect to $R$.
    Therefore there exists $b$ such that $b$ is a largest element of $F$ with respect to $R$.
\end{proof}

% from §27 helper lemma 11: finite inhabited subset has a smallest element
\begin{lemma}\label{finite_inhabited_smallest_element}
    Assume $R$ is a simple order relation on $X$.
    Suppose $F\subseteq X$.
    Suppose $F$ is finite.
    Suppose $F$ is inhabited.
    Then there exists $a$ such that $a$ is a smallest element of $F$ with respect to $R$.
\end{lemma}
\begin{proof}
    Let $S = \converse{R}$.
    $S$ is a simple order relation on $X$
        by \cref{simple_order_relation,converse_type,converse_iff,transitive,asymmetric,connex}.
    Take $a$ such that $a$ is a largest element of $F$ with respect to $S$
        by \cref{finite_inhabited_largest_element}.
    For all $x\in F$ we have $x\mathrel{S}a$ or $x = a$ by \cref{largest_element}.
    For all $x\in F$ we have $a\mathrel{R}x$ or $a = x$ by \cref{converse_iff}.
    $a$ is a smallest element of $F$ with respect to $R$ by \cref{largest_element,smallest_element}.
\end{proof}

% from §27 Item 4: extreme value theorem
\begin{theorem}\label{extreme_value_theorem}
    Assume $T_X$ is a topology on $X$.
    Assume $T_Y$ is the order topology on $Y$ with respect to $R$.
    Suppose $f$ is continuous from $X$ to $Y$ with respect to $T_X$ and $T_Y$.
    Suppose $X$ is compact with respect to $T_X$.
    Suppose $X\neq\emptyset$.
    Then there exist $c,d\in X$ such that
        $\apply{f}{c}$ is a smallest element of $\img{f}{X}$ with respect to $R$ and
        $\apply{f}{d}$ is a largest element of $\img{f}{X}$ with respect to $R$.
\end{theorem}
\begin{proof}
    Let $A = \img{f}{X}$.
    Let $T_A = \subspaceTopology{T_Y}{A}$.
    $A$ is compact with respect to $T_A$ by \cref{continuous_image_compact}.
    $T_Y = \basisTopology{\orderBasis{R}{Y}}{Y}$ and $R$ is a simple order relation on $Y$
        and $Y$ has more than one element by \cref{order_topology}.
    $T_Y$ is a topology on $Y$ by \cref{order_topology,order_basis_is_basis,basis_topology_is_topology}.
    $f$ is a function to $Y$ and $\dom{f} = X$ by \cref{continuous,funs_elim}.
    $A\subseteq Y$ by \cref{img_dom,function_to_ran}.
    $A$ is a subspace of $Y$ with respect to $T_Y$ by \cref{subspace_of}.
    $A$ is inhabited by \cref{nonempty_inhabited,continuous,funs_elim,function_apply_elim,img_iff}.
    Show that there exists $M$ such that $M$ is a largest element of $A$ with respect to $R$.
    \begin{subproof}
        Suppose not.
        Let $C = \{U \mid \exists a\in A. U = \openrayleft{R}{Y}{a}\}$.
        Show that for all $y$ such that $y\in A$ we have $y\in\unions{C}$.
        \begin{subproof}
            Follows by \cref{largest_element,subseteq,order_topology,simple_order_relation,connex,openray_left,unions_iff}.
        \end{subproof}
        $C$ is a covering of $A$ by \cref{subseteq,covering}.
        Take $B$ such that $B\subseteq C$ and $B$ is finite and $B$ is a covering of $A$
            by \cref{subseteq,openrayleft_open_in_order_topology,compact_subspace_open_covering,compact}.
        $B$ is inhabited by \cref{nonempty_inhabited,covering,subseteq,unions_iff}.
        Let $R_B = \{w\in B\times A \mid \exists U\in B.\exists a\in A. w = (U,a) \land U = \openrayleft{R}{Y}{a}\}$.
        For all $w$ such that $w\in R_B$ there exist $U,a$ such that $w = (U,a)$.
        $R_B$ is a relation by \cref{relation}.
        Show that for all $U\in B$ there exists $a$ such that $U\mathrel{R_B}a$.
        \begin{subproof}
            Fix $U$ such that $U\in B$.
            $U\in C$ by \cref{subseteq}.
            Take $a\in A$ such that $U = \openrayleft{R}{Y}{a}$.
            $(U,a)\in B\times A$ by \cref{times_tuple_intro}.
            $(U,a)\in R_B$ by definition.
            $U\mathrel{R_B}a$.
        \end{subproof}
        Take $g$ such that $g$ is a function on $B$ and for all $U\in B$ we have $U\mathrel{R_B} \apply{g}{U}$
            by \cref{axiom_of_choice}.
        Let $E = \img{g}{B}$.
        For every $U$ such that $U \in B$, $g(U) \in A$ by \cref{relation,pair_eq_iff}.
        $E\subseteq A$ by \cref{img_iff,function_apply_intro,subseteq}.
        $E$ is finite and $E$ is inhabited by \cref{finite_image_function,nonempty_inhabited,function_apply_elim,img_iff}.
        $E\subseteq Y$ by \cref{subseteq_transitive}.
        Take $a_0$ such that $a_0$ is a largest element of $E$ with respect to $R$
            by \cref{finite_inhabited_largest_element}.
        Show that for all $U\in B$ we have $a_0\notin U$.
        \begin{subproof}
            Follows by \cref{choice_function,relation,pair_eq_iff,function_apply_elim,img_iff,largest_element,openray_left,simple_order_relation,order_topology,asymmetric}.
        \end{subproof}
        $a_0\notin\unions{B}$ by \cref{unions_iff}.
        $a_0\in\unions{B}$ by \cref{largest_element,elem_subseteq,covering,subseteq}.
        Contradiction.
    \end{subproof}
    Show that there exists $m$ such that $m$ is a smallest element of $A$ with respect to $R$.
    \begin{subproof}
        Suppose not.
        Let $C = \{U \mid \exists a\in A. U = \openrayright{R}{Y}{a}\}$.
        Show that for all $y$ such that $y\in A$ we have $y\in\unions{C}$.
        \begin{subproof}
            Follows by \cref{smallest_element,subseteq,order_topology,simple_order_relation,connex,openray_right,unions_iff}.
        \end{subproof}
        $C$ is a covering of $A$ by \cref{subseteq,covering}.
        Take $B$ such that $B\subseteq C$ and $B$ is finite and $B$ is a covering of $A$
            by \cref{subseteq,openrayright_open_in_order_topology,compact_subspace_open_covering,compact}.
        $B$ is inhabited by \cref{nonempty_inhabited,covering,subseteq,unions_iff}.
        Let $R_B = \{w\in B\times A \mid \exists U\in B.\exists a\in A. w = (U,a) \land U = \openrayright{R}{Y}{a}\}$.
        For all $w$ such that $w\in R_B$ there exist $U,a$ such that $w = (U,a)$.
        $R_B$ is a relation by \cref{relation}.
        Show that for all $U\in B$ there exists $a$ such that $U\mathrel{R_B}a$.
        \begin{subproof}
            Fix $U$ such that $U\in B$.
            $U\in C$ by \cref{subseteq}.
            Take $a\in A$ such that $U = \openrayright{R}{Y}{a}$.
            $(U,a)\in B\times A$ by \cref{times_tuple_intro}.
            $(U,a)\in R_B$ by definition.
            $U\mathrel{R_B}a$.
        \end{subproof}
        Take $g$ such that $g$ is a function on $B$ and for all $U\in B$ we have $U\mathrel{R_B} \apply{g}{U}$
            by \cref{axiom_of_choice}.
        Let $E = \img{g}{B}$.
        For every $U$ such that $U \in B$, $g(U) \in A$ by \cref{relation,pair_eq_iff}.
        $E\subseteq A$ by \cref{img_iff,function_apply_intro,subseteq}.
        $E$ is finite and $E$ is inhabited by \cref{finite_image_function,nonempty_inhabited,function_apply_elim,img_iff}.
        $E\subseteq Y$ by \cref{subseteq_transitive}.
        Take $a_0$ such that $a_0$ is a smallest element of $E$ with respect to $R$
            by \cref{finite_inhabited_smallest_element}.
        Show that for all $U\in B$ we have $a_0\notin U$.
        \begin{subproof}
            Follows by \cref{choice_function,relation,pair_eq_iff,function_apply_elim,img_iff,smallest_element,openray_right,simple_order_relation,order_topology,asymmetric}.
        \end{subproof}
        $a_0\notin\unions{B}$ by \cref{unions_iff}.
        $a_0\in\unions{B}$ by \cref{smallest_element,elem_subseteq,covering,subseteq}.
        Contradiction.
    \end{subproof}
    Take $m$ such that $m$ is a smallest element of $A$ with respect to $R$.
    Take $c\in X$ such that $f(c) = m$ by \cref{smallest_element,img_iff,function_apply_intro}.
    Take $M$ such that $M$ is a largest element of $A$ with respect to $R$.
    Take $d\in X$ such that $f(d) = M$ by \cref{largest_element,img_iff,function_apply_intro}.
    Show that $\apply{f}{c}$ is a smallest element of $A$ with respect to $R$ and
        $\apply{f}{d}$ is a largest element of $A$ with respect to $R$.
    \begin{subproof}
        Follows by \cref{smallest_element,largest_element}.
    \end{subproof}
\end{proof}

% from §27 helper definition 4: distance set from a point to a set
\begin{definition}\label{point_distance_set}
    $\pointDistanceSet{d}{A}{x} = \{d((x,a)) \mid a\in A\}$.
\end{definition}

% from §27 Item 5: distance from a point to a set
\begin{definition}\label{point_set_distance}
    $r$ is the distance from $x$ to $A$ in $X$ with respect to $d$ iff
        $x\in X$ and $A\subseteq X$ and $r\in\reals$ and
        $r$ is an infimum of $\pointDistanceSet{d}{A}{x}$.
\end{definition}

% from §27 helper lemma 3: distance-to-set map is continuous
\begin{lemma}\label{point_set_distance_continuous}
    Suppose $d$ is a metric on $X$.
    Suppose $A\subseteq X$.
    Suppose $A\neq\emptyset$.
    Let $T = \metricTopology{d}{X}$.
    Then there exists $f$ such that
        $f$ is a function from $X$ to $\reals$ and
        for all $x\in X$ we have $f(x)$ is the distance from $x$ to $A$ in $X$ with respect to $d$ and
        $f$ is continuous from $X$ to $\reals$ with respect to $T$ and $\standardTopologyR$.
\end{lemma}
\begin{proof}
    Let $R = \{w\in X\times\reals \mid \text{there exists $x\in X$ such that there exists $r\in\reals$ such that
        $w = (x,r)$ and $r$ is the distance from $x$ to $A$ in $X$ with respect to $d$}\}$.
    For all $w$ such that $w\in R$ there exist $x,r$ such that $w = (x,r)$.
    $R$ is a relation by \cref{relation}.
    Show that for all $x\in X$ there exists $r$ such that $x\mathrel{R} r$.
    \begin{subproof}
        Fix $x\in X$.
        Let $S = \pointDistanceSet{d}{A}{x}$.
        $S\subseteq\reals$ by \cref{point_distance_set,subseteq,times_tuple_intro,metric_on,funs_type_apply}.
        $S\neq\emptyset$ by \cref{nonempty_inhabited,point_distance_set,emptyset}.
        $0$ is a lower bound of $S$ and $S$ is bounded below
            by \cref{point_distance_set,subseteq,metric_on,metric_nonnegative,real_add_ident,real_lower_bound,real_bounded_below}.
        Take $r$ such that $r$ is an infimum of $S$ by \cref{real_infimum_exists}.
        $r$ is the distance from $x$ to $A$ in $X$ with respect to $d$
            by \cref{real_infimum,real_lower_bound,point_set_distance}.
        $(x,r)\in R$.
        $x\mathrel{R} r$ by \cref{relation}.
    \end{subproof}
    Take $f$ such that $f$ is a function on $X$ and for all $x\in X$ we have $x\mathrel{R} f(x)$
        by \cref{axiom_of_choice}.
    For all $x\in X$ we have $f(x)$ is the distance from $x$ to $A$ in $X$ with respect to $d$ and $f(x)\in\reals$
        by \cref{relation,pair_eq_iff,point_set_distance}.
    $f$ is a function from $X$ to $\reals$ by \cref{funs_intro,funs}.
    Let $T_Y = \metricTopology{\standardMetricR}{\reals}$.
    Show that for all $x\in X$ we have
        for all $\epsilon\in\reals$ such that $0 < \epsilon$
        there exists $\delta\in\reals$ such that $0 < \delta$ and
        for all $y\in X$ if $d((x,y)) < \delta$ then
            $\apply{\standardMetricR}{(f(x),f(y))} < \epsilon$.
    \begin{subproof}
        Fix $x\in X$.
        Fix $\epsilon$ such that $\epsilon\in\reals$ and $0 < \epsilon$.
        Let $\delta = \epsilon$.
        $0 < \delta$.
        Show that for all $y\in X$ if $d((x,y)) < \delta$ then
            $\apply{\standardMetricR}{(f(x),f(y))} < \epsilon$.
        \begin{subproof}
            Fix $y\in X$.
            Assume $d((x,y)) < \delta$.
            $d((x,y)) < \epsilon$ by \cref{real_lt_subst_right}.
            Let $S_x = \pointDistanceSet{d}{A}{x}$.
            Let $S_y = \pointDistanceSet{d}{A}{y}$.
                $f(x)$ is an infimum of $S_x$ and $f(y)$ is an infimum of $S_y$ by \cref{point_set_distance}.
                $f(x)$ is a lower bound of $S_x$ and $f(y)$ is a lower bound of $S_y$ by \cref{real_infimum}.
                Show that $f(x) \geq f(y) - d((x,y))$.
                \begin{subproof}
                    Let $b = f(y) - d((x,y))$.
                    $f(y)\in\reals$ by \cref{point_set_distance}.
                    $(x,y)\in X\times X$ by \cref{times_tuple_intro}.
                    $d\in\funs{X\times X}{\reals}$ by \cref{metric_on}.
                    $d((x,y))\in\reals$ by \cref{funs_type_apply}.
                    $b = f(y) + \neg{(d((x,y)))}$ by \cref{real_sub}.
                    $\neg{(d((x,y)))}\in\reals$ by \cref{real_add_inverse}.
                    $b\in\reals$ by \cref{real_add_closed}.
                    Show that for all $s\in S_x$ we have $b \leq s$.
                    \begin{subproof}
                        Fix $s\in S_x$.
                        Take $a\in A$ such that $s = d((x,a))$ by \cref{point_distance_set}.
                        $a\in X$ by \cref{subseteq}.
                        $d((y,x)) = d((x,y))$ by \cref{metric_on,metric_symmetric}.
                        $d((y,x)) + d((x,a)) \geq d((y,a))$ by \cref{metric_on,metric_triangle_inequality}.
                        $d((y,a)) \leq d((y,x)) + d((x,a))$.
                        $d((y,a)) \leq d((x,y)) + d((x,a))$.
                        $(y,a)\in X\times X$ and $(x,y)\in X\times X$ and $(x,a)\in X\times X$
                            by \cref{times_tuple_intro}.
                        $d((y,a))\in\reals$ and $d((x,y))\in\reals$ and $d((x,a))\in\reals$
                            by \cref{funs_type_apply}.
                        $d((x,y)) + d((x,a))\in\reals$ by \cref{real_add_closed}.
                        $f(y) \leq d((y,a))$ by \cref{point_set_distance,real_lower_bound,point_distance_set}.
                        $f(y) \leq d((x,y)) + d((x,a))$ by \cref{real_leq_trans}.
                        Let $u = d((x,y))$.
                        Let $v = d((x,a))$.
                        $f(y) \leq u + v$ by \cref{real_leq_trans}.
                        $f(y) + \neg{u} \leq (u + v) + \neg{u}$ by \cref{real_leq_add}.
                        $u\in\reals$ and $v\in\reals$.
                        $u + \neg{u} = 0$ by \cref{real_add_inverse}.
                        $(u + v) + \neg{u} = v + (u + \neg{u})$ by \cref{real_add_assoc,real_add_comm}.
                        $(u + v) + \neg{u} = v + 0$ by \cref{real_add_inverse}.
                        $0\in\reals$ by \cref{real_add_ident}.
                        $v + 0 = 0 + v$ by \cref{real_add_comm}.
                        $0 + v = v$ by \cref{real_add_ident}.
                        $(u + v) + \neg{u} = v$ by \cref{real_add_comm,real_add_ident}.
                        $f(y) + \neg{u} \leq v$ by \cref{real_leq_trans}.
                        $f(y) - u \leq v$ by \cref{real_sub}.
                        $f(y) - d((x,y)) \leq d((x,a))$.
                        $b \leq s$.
                    \end{subproof}
                    $b$ is a lower bound of $S_x$ by \cref{real_lower_bound,subseteq}.
                    $f(x) \geq b$ by \cref{real_infimum}.
                \end{subproof}
                Show that $f(y) \geq f(x) - d((x,y))$.
                \begin{subproof}
                    Let $b = f(x) - d((x,y))$.
                    $f(x)\in\reals$ and $(x,y)\in X\times X$ and $d\in\funs{X\times X}{\reals}$
                        by \cref{point_set_distance,times_tuple_intro,metric_on}.
                    $d((x,y))\in\reals$ by \cref{funs_type_apply}.
                    $b = f(x) + \neg{(d((x,y)))}$ by \cref{real_sub}.
                    $\neg{(d((x,y)))}\in\reals$ by \cref{real_add_inverse}.
                    $b\in\reals$ by \cref{real_add_closed}.
                    Show that for all $s\in S_y$ we have $b \leq s$.
                    \begin{subproof}
                        Fix $s\in S_y$.
                        Take $a\in A$ such that $s = d((y,a))$ by \cref{point_distance_set}.
                        $a\in X$ by \cref{subseteq}.
                        $d((x,y)) = d((y,x))$ by \cref{metric_on,metric_symmetric}.
                        $d((x,y)) + d((y,a)) \geq d((x,a))$ by \cref{metric_on,metric_triangle_inequality}.
                        $d((x,a)) \leq d((x,y)) + d((y,a))$.
                        $(x,a)\in X\times X$ and $(x,y)\in X\times X$ and $(y,a)\in X\times X$
                            by \cref{times_tuple_intro}.
                        $d((x,a))\in\reals$ and $d((x,y))\in\reals$ and $d((y,a))\in\reals$
                            by \cref{funs_type_apply}.
                        $d((x,y)) + d((y,a))\in\reals$ by \cref{real_add_closed}.
                        $f(x) \leq d((x,a))$ by \cref{point_set_distance,real_lower_bound,point_distance_set}.
                        $f(x) \leq d((x,y)) + d((y,a))$ by \cref{real_leq_trans}.
                        Let $u = d((x,y))$.
                        Let $v = d((y,a))$.
                        $f(x) \leq u + v$ by \cref{real_leq_trans}.
                        $f(x) + \neg{u} \leq (u + v) + \neg{u}$ by \cref{real_leq_add}.
                        $u\in\reals$ and $v\in\reals$.
                        $u + \neg{u} = 0$ by \cref{real_add_inverse}.
                        $(u + v) + \neg{u} = v + (u + \neg{u})$ by \cref{real_add_assoc,real_add_comm}.
                        $(u + v) + \neg{u} = v + 0$ by \cref{real_add_inverse}.
                        $0\in\reals$ by \cref{real_add_ident}.
                        $v + 0 = 0 + v$ by \cref{real_add_comm}.
                        $0 + v = v$ by \cref{real_add_ident}.
                        $(u + v) + \neg{u} = v$ by \cref{real_add_comm,real_add_ident}.
                        $f(x) + \neg{u} \leq v$ by \cref{real_leq_trans}.
                        $f(x) - u \leq v$ by \cref{real_sub}.
                        $f(x) - d((x,y)) \leq d((y,a))$.
                        $b \leq s$.
                    \end{subproof}
                    $b$ is a lower bound of $S_y$ by \cref{real_lower_bound,subseteq}.
                    $f(y) \geq b$ by \cref{real_infimum}.
                \end{subproof}
                $f(x) - d((x,y)) \leq f(y)$.
                $f(x)\in\reals$ and $f(y)\in\reals$ by \cref{point_set_distance}.
                $(x,y)\in X\times X$ and $d\in\funs{X\times X}{\reals}$
                    by \cref{times_tuple_intro,metric_on}.
                $d((x,y))\in\reals$ by \cref{funs_type_apply}.
                $\neg{(d((x,y)))}\in\reals$ by \cref{real_add_inverse}.
                $f(x) + \neg{(d((x,y)))}\in\reals$ by \cref{real_add_closed}.
                $f(x) - d((x,y))\in\reals$ by \cref{real_sub}.
                $(f(x) - d((x,y))) + d((x,y)) \leq f(y) + d((x,y))$ by \cref{real_leq_add}.
                Let $s_0 = (f(x) - d((x,y))) + d((x,y))$.
                $f(x) - d((x,y)) = f(x) + \neg{(d((x,y)))}$ by \cref{real_sub}.
                $s_0 = f(x) + (\neg{(d((x,y)))} + d((x,y)))$ by \cref{real_add_assoc}.
                $\neg{(d((x,y)))} + d((x,y)) = 0$ by \cref{real_add_inverse,real_add_comm}.
                $s_0 = f(x) + 0$ by \cref{real_add_assoc}.
                $f(x) + 0 = 0 + f(x)$ by \cref{point_set_distance,real_add_ident,real_add_comm}.
                $0 + f(x) = f(x)$ by \cref{point_set_distance,real_add_ident}.
                $s_0 = f(x)$ by \cref{real_add_comm,real_add_ident}.
                $f(x) \leq s_0$.
                $f(x) \leq f(y) + d((x,y))$ by \cref{real_leq_trans}.
                $f(y) - d((x,y)) \leq f(x)$.
                $f(y)\in\reals$ and $f(x)\in\reals$ by \cref{point_set_distance}.
                $(x,y)\in X\times X$ and $d\in\funs{X\times X}{\reals}$
                    by \cref{times_tuple_intro,metric_on}.
                $d((x,y))\in\reals$ by \cref{funs_type_apply}.
                $\neg{(d((x,y)))}\in\reals$ by \cref{real_add_inverse}.
                $f(y) + \neg{(d((x,y)))}\in\reals$ by \cref{real_add_closed}.
                $f(y) - d((x,y))\in\reals$ by \cref{real_sub}.
                We have $(f(y) - d((x,y))) + d((x,y)) \leq f(x) + d((x,y))$ by \cref{real_leq_add}.
                Let $s_1 = (f(y) - d((x,y))) + d((x,y))$.
                $f(y) - d((x,y)) = f(y) + \neg{(d((x,y)))}$ by \cref{real_sub}.
                $s_1 = f(y) + (\neg{(d((x,y)))} + d((x,y)))$ by \cref{real_add_assoc}.
                $\neg{(d((x,y)))} + d((x,y)) = 0$ by \cref{real_add_inverse,real_add_comm}.
                $s_1 = f(y) + 0$ by \cref{real_add_assoc}.
                $f(y) + 0 = 0 + f(y)$ by \cref{point_set_distance,real_add_ident,real_add_comm}.
                $0 + f(y) = f(y)$ by \cref{point_set_distance,real_add_ident}.
                $s_1 = f(y)$ by \cref{real_add_comm,real_add_ident}.
                $f(y) \leq s_1$.
                $f(y) \leq f(x) + d((x,y))$ by \cref{real_leq_trans}.
                $f(x)\in\reals$ and $f(y)\in\reals$ by \cref{point_set_distance}.
                $(x,y)\in X\times X$ and $d\in\funs{X\times X}{\reals}$
                    by \cref{times_tuple_intro,metric_on}.
                $d((x,y))\in\reals$ by \cref{funs_type_apply}.
                $\neg{(f(y))}\in\reals$ by \cref{real_add_inverse}.
                $f(x) + \neg{(f(y))}\in\reals$ by \cref{real_add_closed}.
                $f(x) - f(y)\in\reals$ by \cref{real_sub}.
                $\neg{(f(x) - f(y))}\in\reals$ by \cref{real_add_inverse}.
                $f(y) + d((x,y))\in\reals$ by \cref{real_add_closed}.
                We have $f(x) + \neg{(f(y))} \leq (f(y) + d((x,y))) + \neg{(f(y))}$ by \cref{real_leq_add}.
                $(f(y) + d((x,y))) + \neg{(f(y))} = f(y) + (d((x,y)) + \neg{(f(y))})$ by \cref{real_add_assoc}.
                $d((x,y)) + \neg{(f(y))} = \neg{(f(y))} + d((x,y))$ by \cref{real_add_comm}.
                $f(y) + (\neg{(f(y))} + d((x,y))) = (f(y) + \neg{(f(y))}) + d((x,y))$ by \cref{real_add_assoc}.
                $f(y) + \neg{(f(y))} = 0$ by \cref{real_add_inverse}.
                $0\in\reals$ by \cref{real_add_ident}.
                $(f(y) + d((x,y))) + \neg{(f(y))} = d((x,y))$ by \cref{real_add_ident,real_add_assoc,real_add_comm}.
                $f(x) + \neg{(f(y))} \leq d((x,y))$ by \cref{real_leq_trans}.
                $f(x) - f(y) \leq d((x,y))$ by \cref{real_sub}.
                $f(x) + d((x,y))\in\reals$ and $\neg{(f(x))}\in\reals$
                    by \cref{real_add_closed,real_add_inverse}.
                $f(y) + \neg{(f(x))} \leq (f(x) + d((x,y))) + \neg{(f(x))}$ by \cref{real_leq_add}.
                $(f(x) + d((x,y))) + \neg{(f(x))} = f(x) + (d((x,y)) + \neg{(f(x))})$ by \cref{real_add_assoc}.
                $d((x,y)) + \neg{(f(x))} = \neg{(f(x))} + d((x,y))$ by \cref{real_add_comm}.
                $f(x) + (\neg{(f(x))} + d((x,y))) = (f(x) + \neg{(f(x))}) + d((x,y))$ by \cref{real_add_assoc}.
                $f(x) + \neg{(f(x))} = 0$ by \cref{real_add_inverse}.
                $0\in\reals$ by \cref{real_add_ident}.
                $(f(x) + d((x,y))) + \neg{(f(x))} = d((x,y))$ by \cref{real_add_ident,real_add_assoc,real_add_comm}.
                $f(y) + \neg{(f(x))} \leq d((x,y))$ by \cref{real_leq_trans}.
                $f(y) - f(x) \leq d((x,y))$ by \cref{real_sub}.
                $f(y) - f(x) = \neg{(f(x) - f(y))}$ by \cref{real_sub_swap_neg}.
                $\neg{(f(x) - f(y))} \leq d((x,y))$.
                $\abs{(f(x) - f(y))} \leq d((x,y))$ by \cref{real_abs_leq_if_pm_leq}.
                $f(x)\in\reals$ and $f(y)\in\reals$ by \cref{point_set_distance}.
                $\neg{(f(y))}\in\reals$ by \cref{real_add_inverse}.
                $f(x) + \neg{(f(y))}\in\reals$ by \cref{real_add_closed}.
                $f(x) - f(y)\in\reals$ by \cref{real_sub}.
                $(f(x),f(y))\in\reals\times\reals$ by \cref{times_tuple_intro}.
                $\abs{(f(x) - f(y))}\in\reals$ by \cref{real_abs_in_reals}.
                Let $p = (f(x),f(y))$.
                $\fst{p} = f(x)$ and $\snd{p} = f(y)$ by \cref{fst_eq,snd_eq}.
                $\abs{(\fst{p} - \snd{p})} = \abs{(f(x) - f(y))}$.
                $(p,\abs{(f(x) - f(y))})\in\standardMetricR$ by \cref{standard_metric_r}.
                $\apply{\standardMetricR}{(f(x),f(y))} = \abs{(f(x) - f(y))}$
                    by \cref{standard_metric_is_metric,metric_on,funs_elim,function_apply_intro}.
                $\apply{\standardMetricR}{(f(x),f(y))} \leq d((x,y))$.
                $(x,y)\in X\times X$ by \cref{times_tuple_intro}.
                $d\in\funs{X\times X}{\reals}$ by \cref{metric_on}.
                $d((x,y))\in\reals$ and $\apply{\standardMetricR}{(f(x),f(y))}\in\reals$
                    by \cref{funs_type_apply}.
                $\apply{\standardMetricR}{(f(x),f(y))} < \epsilon$
                    by \cref{real_leq_split,real_lt_trans,real_lt_subst_left}.
        \end{subproof}
    \end{subproof}
    $f$ is continuous from $X$ to $\reals$ with respect to $T$ and $T_Y$
        by \cref{standard_metric_is_metric,metric_continuity_eps_delta}.
    $f$ is continuous from $X$ to $\reals$ with respect to $T$ and $\standardTopologyR$
        by \cref{standard_metric_topology,continuous_eq_codomain}.
\end{proof}

from §27 Item 6: Lebesgue number lemma
\begin{lemma}\label{lebesgue_number_lemma}
    Suppose $d$ is a metric on $X$.
    Let $T = \metricTopology{d}{X}$.
    Suppose $A$ is an open covering of $X$ with respect to $T$.
    Suppose $X\neq\emptyset$.
    Suppose $X$ is compact with respect to $T$.
    Then there exists $\delta\in\reals$ such that $0 < \delta$ and
        for all $B$ such that $B\subseteq X$ and
            for all $x,y\in B$ we have $d((x,y)) < \delta$
        there exists $U\in A$ such that $B\subseteq U$.
\end{lemma}
\begin{proof}
    Take $x_0$ such that $x_0\in X$ by \cref{nonempty_inhabited}.
    Take $U_0\in A$ such that $x_0\in U_0$ by \cref{open_covering,covering,subseteq,unions_iff}.
    Let $R = \{w\in X\times\pow{X} \mid \text{there exists $x\in X$ such that there exists $D\in\pow{X}$ such that
        there exists $\delta\in\reals$ such that there exists $U\in A$ such that
        $w = (x,D)$ and
        $0 < \delta$ and
        $D = \metricBall{d}{X}{x}{\delta}$ and
        $\metricBall{d}{X}{x}{\delta + \delta} \subseteq U$}\}$.
    For all $w$ such that $w\in R$ there exist $x,D$ such that $w = (x,D)$.
    $R$ is a relation by \cref{relation}.
    Show that for all $x\in X$ there exists $D$ such that $x\mathrel{R}D$.
    \begin{subproof}
        Fix $x\in X$.
        Take $U\in A$ such that $x\in U$ by \cref{open_covering,covering,subseteq,unions_iff}.
        $U$ is open in $X$ with respect to $T$ by \cref{open_covering}.
        $T$ is a topology on $X$ by \cref{metric_basis_is_basis,basis_topology_is_topology,metric_topology}.
        $U\in T$ by \cref{open_in}.
        $U\in\pow{X}$ by \cref{topology_on,subseteq}.
        $U\subseteq X$ by \cref{powerset_elim}.
        Take $\epsilon\in\reals$ such that $0 < \epsilon$ and $\metricBall{d}{X}{x}{\epsilon}\subseteq U$
            by \cref{metric_open_iff}.
        Let $\delta = \epsilon / (1 + 1)$.
        $1\in\reals$ by \cref{real_mul_ident}.
        $0\in\reals$ by \cref{real_add_ident}.
        $0 < 1$ by \cref{real_zero_lt_one}.
        $1 < 1 + 1$ by \cref{real_lt_add,real_add_ident}.
        $(1 + 1)\in\reals$ and $0 < 1 + 1$ by \cref{real_add_closed,real_lt_trans}.
        $(1 + 1)\neq 0$ by \cref{real_pos_nonzero}.
        $\inv{(1 + 1)}\in\reals$ and $\inv{(1 + 1)} > 0$ by \cref{real_mul_inverse,real_inv_pos}.
        $\delta = \epsilon \rmul \inv{(1 + 1)}$ and $\delta\in\reals$ by \cref{real_div,real_mul_closed}.
        $0 < \delta$ by \cref{real_div,real_mul_zero,real_lt_mul_pos,real_lt_subst_left}.
        Let $D = \metricBall{d}{X}{x}{\delta}$.
        $D\subseteq X$ by \cref{metric_ball,subseteq}.
        $D\in\pow{X}$ by \cref{powerset_intro}.
        Show that $\metricBall{d}{X}{x}{\delta + \delta}\subseteq U$.
        \begin{subproof}
            $\delta = \epsilon \rmul \inv{(1 + 1)}$ by \cref{real_div}.
            $\delta + \delta =
                (\epsilon \rmul \inv{(1 + 1)}) + (\epsilon \rmul \inv{(1 + 1)})$.
            $(\epsilon + \epsilon) \rmul \inv{(1 + 1)}
                = (\epsilon \rmul \inv{(1 + 1)}) + (\epsilon \rmul \inv{(1 + 1)})$
                by \cref{real_right_distrib}.
            $\delta + \delta = (\epsilon + \epsilon) \rmul \inv{(1 + 1)}$.
            $\epsilon \rmul (1 + 1) = (\epsilon \rmul 1) + (\epsilon \rmul 1)$ by \cref{real_distrib}.
            $\epsilon \rmul 1 = \epsilon$ by \cref{real_mul_ident,real_mul_comm}.
            $\epsilon \rmul (1 + 1) = \epsilon + \epsilon$.
            $(\epsilon + \epsilon) \rmul \inv{(1 + 1)} = (\epsilon \rmul (1 + 1)) \rmul \inv{(1 + 1)}$.
            $(\epsilon \rmul (1 + 1)) \rmul \inv{(1 + 1)}
                = \epsilon \rmul ((1 + 1) \rmul \inv{(1 + 1)})$ by \cref{real_mul_assoc}.
            $(1 + 1) \rmul \inv{(1 + 1)} = 1$ by \cref{real_mul_inverse}.
            $\delta + \delta = \epsilon \rmul 1$.
            $\delta + \delta = \epsilon$ by \cref{real_mul_ident}.
            Show that for all $y$ such that $y\in\metricBall{d}{X}{x}{\delta + \delta}$ we have $y\in U$.
            \begin{subproof}
                Fix $y$ such that $y\in\metricBall{d}{X}{x}{\delta + \delta}$.
                $y\in\metricBall{d}{X}{x}{\epsilon}$ by \cref{real_lt_subst_right,metric_ball}.
                $y\in U$ by \cref{subseteq}.
            \end{subproof}
            $\metricBall{d}{X}{x}{\delta + \delta}\subseteq U$ by \cref{subseteq}.
        \end{subproof}
        $(x,D)\in R$.
        $x\mathrel{R}D$ by \cref{relation}.
    \end{subproof}
    Take $g$ such that $g$ is a function on $X$ and for all $x\in X$ we have $x\mathrel{R}\apply{g}{x}$
        by \cref{choice_function}.
    Let $C = \img{g}{X}$.
    Show that for all $D$ such that $D\in C$ we have $D$ is open in $X$ with respect to $T$.
    \begin{subproof}
        Fix $D$ such that $D\in C$.
        Take $x$ such that $x\in X$ and $(x,D)\in g$ by \cref{img_iff}.
        $x\in\dom{g}$ by \cref{choice_function}.
        $D = \apply{g}{x}$ by \cref{function_apply_iff}.
        $x\mathrel{R}D$ by \cref{choice_function}.
        Take $x_1,D_1,\delta,U$ such that $x_1\in X$ and $D_1\in\pow{X}$ and $\delta\in\reals$ and $U\in A$ and
            $(x,D) = (x_1,D_1)$ and
            $0 < \delta$ and
            $D_1 = \metricBall{d}{X}{x_1}{\delta}$ and
            $\metricBall{d}{X}{x_1}{\delta + \delta}\subseteq U$
            by \cref{relation}.
        $x = x_1$ and $D = D_1$ by \cref{pair_eq_iff}.
        $D\in\metricBasis{d}{X}$ by \cref{metric_ball,powerset_intro,metric_basis}.
        Let $B = \metricBasis{d}{X}$.
        $B$ is a basis on $X$ by \cref{metric_basis_is_basis}.
        $T = \basisTopology{B}{X}$ by \cref{metric_topology}.
        $D\in T$ by \cref{basis_elem_open_in_generated}.
        $T$ is a topology on $X$ by \cref{basis_topology_is_topology}.
        $D$ is open in $X$ with respect to $T$ by \cref{open_in}.
    \end{subproof}
    Show that for all $x$ such that $x\in X$ we have $x\in\unions{C}$.
    \begin{subproof}
        Fix $x$ such that $x\in X$.
        $x\mathrel{R}\apply{g}{x}$ by \cref{choice_function}.
        Take $x_1,D,\delta,U$ such that $x_1\in X$ and $D\in\pow{X}$ and $\delta\in\reals$ and $U\in A$ and
            $(x,\apply{g}{x}) = (x_1,D)$ and
            $0 < \delta$ and
            $D = \metricBall{d}{X}{x_1}{\delta}$ and
            $\metricBall{d}{X}{x_1}{\delta + \delta}\subseteq U$
            by \cref{relation}.
        $x = x_1$ and $\apply{g}{x} = D$ by \cref{pair_eq_iff}.
        $x\in D$
            by \cref{metric_ball,metric_on,metric_zero_characterization,real_lt_subst_left}.
        $\apply{g}{x}\in C$ by \cref{function_apply_elim,img_iff}.
        $x\in\unions{C}$ by \cref{unions_iff}.
    \end{subproof}
    $C$ is a covering of $X$ by \cref{subseteq,covering}.
    $C$ is an open covering of $X$ with respect to $T$ by \cref{open_covering}.
    Take $F$ such that $F\subseteq C$ and $F$ is finite and $F$ is a covering of $X$ by \cref{compact}.
    $F$ is inhabited by \cref{nonempty_inhabited,covering,subseteq,unions_iff}.
    Let $S = \{w\in F\times\reals \mid \text{there exists $D\in F$ such that there exists $\delta\in\reals$ such that
        there exists $x\in X$ such that there exists $U\in A$ such that
        $w = (D,\delta)$ and
        $D = \apply{g}{x}$ and
        $0 < \delta$ and
        $D = \metricBall{d}{X}{x}{\delta}$ and
        $\metricBall{d}{X}{x}{\delta + \delta}\subseteq U$}\}$.
    For all $w$ such that $w\in S$ there exist $D,\delta$ such that $w = (D,\delta)$.
    $S$ is a relation by \cref{relation}.
    Show that for all $D\in F$ there exists $\delta$ such that $D\mathrel{S}\delta$.
    \begin{subproof}
        Fix $D\in F$.
        $D\in C$ by \cref{subseteq}.
        Take $x$ such that $x\in X$ and $(x,D)\in g$ by \cref{img_iff}.
        $x\in\dom{g}$ by \cref{choice_function}.
        $D = \apply{g}{x}$ by \cref{function_apply_iff}.
        $x\mathrel{R}D$ by \cref{choice_function}.
        Take $x_1,D_1,\delta,U$ such that $x_1\in X$ and $D_1\in\pow{X}$ and $\delta\in\reals$ and $U\in A$ and
            $(x,D) = (x_1,D_1)$ and
            $0 < \delta$ and
            $D_1 = \metricBall{d}{X}{x_1}{\delta}$ and
            $\metricBall{d}{X}{x_1}{\delta + \delta}\subseteq U$
            by \cref{relation}.
        $x = x_1$ and $D = D_1$ by \cref{pair_eq_iff}.
        $(D,\delta)\in F\times\reals$ by \cref{times_tuple_intro}.
        $(D,\delta)\in S$ by \cref{times_tuple_intro}.
        $D\mathrel{S}\delta$ by \cref{relation}.
    \end{subproof}
    Take $e$ such that $e$ is a function on $F$ and for all $D\in F$ we have $D\mathrel{S}\apply{e}{D}$
        by \cref{choice_function}.
    Let $E = \img{e}{F}$.
    $E$ is finite by \cref{finite_image_function}.
    Take $D_0$ such that $D_0\in F$ by \cref{nonempty_inhabited}.
    $\apply{e}{D_0}\in E$ by \cref{function_apply_elim,img_iff}.
    $E\neq\emptyset$ by \cref{emptyset}.
    Show that for all $t$ such that $t\in E$ we have $t\in\reals$ and $0 < t$.
    \begin{subproof}
        Fix $t$ such that $t\in E$.
        Take $D_1$ such that $D_1\in F$ and $(D_1,t)\in e$ by \cref{img_iff}.
        $D_1\in\dom{e}$ by \cref{subseteq}.
        $t = \apply{e}{D_1}$ by \cref{function_apply_iff}.
        $D_1\mathrel{S}t$ by \cref{subseteq}.
        $(D_1,t)\in S$ by \cref{relation}.
        Take $D_2,\delta,x,U$ such that $D_2\in F$ and $\delta\in\reals$ and $x\in X$ and $U\in A$ and
            $(D_1,t) = (D_2,\delta)$ and
            $D_2 = \apply{g}{x}$ and
            $0 < \delta$ and
            $D_2 = \metricBall{d}{X}{x}{\delta}$ and
            $\metricBall{d}{X}{x}{\delta + \delta}\subseteq U$
            by \cref{subseteq}.
        $t = \delta$ by \cref{pair_eq_iff}.
        $t\in\reals$ and $0 < t$ by \cref{real_lt_subst_right,real_lt_subst_left}.
    \end{subproof}
    $E\subseteq\reals$ by \cref{subseteq}.
    Take $\delta\in E$ such that $\delta$ is a minimum of $E$ by \cref{finite_real_minimum}.
    $\delta\in\reals$ and $0 < \delta$ by \cref{subseteq}.
    Show that for all $B$ such that $B\subseteq X$ and
        for all $x,y\in B$ we have $d((x,y)) < \delta$
        there exists $U\in A$ such that $B\subseteq U$.
    \begin{subproof}
        Fix $B$ such that $B\subseteq X$ and
            for all $x,y\in B$ we have $d((x,y)) < \delta$.
        \begin{byCase}
        \caseOf{$B = \emptyset$.}
            $B\subseteq U_0$ by \cref{emptyset_subseteq}.
            Show that there exists $U\in A$ such that $B\subseteq U$.
            \begin{subproof}
                Follows by \cref{open_covering}.
            \end{subproof}
        \caseOf{$B \neq \emptyset$.}
            Take $z$ such that $z\in B$ by \cref{nonempty_inhabited}.
            Take $D\in F$ such that $z\in D$ by \cref{covering,subseteq,unions_iff}.
            $D\mathrel{S}\apply{e}{D}$ by \cref{choice_function}.
            $(D,\apply{e}{D})\in S$ by \cref{relation}.
            Take $D_1,\delta_1,x,U$ such that $D_1\in F$ and $\delta_1\in\reals$ and $x\in X$ and $U\in A$ and
                $(D,\apply{e}{D}) = (D_1,\delta_1)$ and
                $D_1 = \apply{g}{x}$ and
                $0 < \delta_1$ and
                $D_1 = \metricBall{d}{X}{x}{\delta_1}$ and
                $\metricBall{d}{X}{x}{\delta_1 + \delta_1}\subseteq U$
                by \cref{subseteq}.
            $D = D_1$ and $\apply{e}{D} = \delta_1$ by \cref{pair_eq_iff}.
            $\delta_1\in E$ by \cref{function_apply_elim,img_iff}.
            $\delta \leq \delta_1$ by \cref{real_minimum}.
            Show that for all $y$ such that $y\in B$ we have $y\in U$.
            \begin{subproof}
                Fix $y$ such that $y\in B$.
                $d((z,y)) < \delta$.
                $z\in X$ and $y\in X$ by \cref{subseteq}.
                $(z,y)\in X\times X$ and $(x,z)\in X\times X$ and $(x,y)\in X\times X$
                    by \cref{times_tuple_intro,metric_ball}.
                $d\in\funs{X\times X}{\reals}$ by \cref{metric_on}.
                $d((z,y))\in\reals$ and $d((x,z))\in\reals$ and $d((x,y))\in\reals$
                    by \cref{funs_type_apply}.
                $z\in\metricBall{d}{X}{x}{\delta_1}$ by \cref{pair_eq_iff,real_lt_subst_left,metric_ball}.
                $d((x,z)) < \delta_1$ by \cref{metric_ball}.
                Show that $d((z,y)) < \delta_1$.
                \begin{subproof}
                    We have $\delta < \delta_1$ or $\delta = \delta_1$ by \cref{real_leq_split}.
                    \begin{byCase}
                    \caseOf{$\delta < \delta_1$.}
                        $d((z,y)) < \delta_1$ by \cref{real_lt_trans}.
                    \caseOf{$\delta = \delta_1$.}
                        $d((z,y)) < \delta_1$ by \cref{real_lt_subst_right}.
                    \end{byCase}
                \end{subproof}
                $d((x,z)) + d((z,y)) < \delta_1 + \delta_1$ by \cref{real_lt_add_two}.
                $d((x,y)) \leq d((x,z)) + d((z,y))$ by \cref{metric_on,metric_triangle_inequality}.
                $d((x,y)) < \delta_1 + \delta_1$
                    by \cref{real_add_closed,real_leq_split,real_lt_trans,real_lt_subst_left}.
                $y\in\metricBall{d}{X}{x}{\delta_1 + \delta_1}$ by \cref{metric_ball}.
                $y\in U$ by \cref{subseteq}.
            \end{subproof}
            $B\subseteq U$ by \cref{subseteq}.
            Show that there exists $V\in A$ such that $B\subseteq V$.
            \begin{subproof}
                Follows by \cref{subseteq}.
            \end{subproof}
        \end{byCase}
    \end{subproof}
    Show that there exists $\eta\in\reals$ such that $0 < \eta$ and
        for all $B$ such that $B\subseteq X$ and
            for all $x,y\in B$ we have $d((x,y)) < \eta$
        there exists $U\in A$ such that $B\subseteq U$.
    \begin{subproof}
        Trivial.
    \end{subproof}
\end{proof}

% from §27 Item 7: Lebesgue number
\begin{definition}\label{lebesgue_number}
    $\delta$ is a Lebesgue number for $A$ in $X$ with respect to $d$ iff
        $d$ is a metric on $X$ and
        $A$ is an open covering of $X$ with respect to $\metricTopology{d}{X}$ and
        $X\neq\emptyset$ and
        $\delta\in\reals$ and $0 < \delta$ and
        for all $B$ such that $B\subseteq X$ and
            for all $x,y\in B$ we have $d((x,y)) < \delta$
        there exists $U\in A$ such that $B\subseteq U$.
\end{definition}

% from §27 Item 8: uniformly continuous function
\begin{definition}\label{uniformly_continuous}
    $f$ is uniformly continuous from $X$ to $Y$ with respect to $d_X$ and $d_Y$ iff
        $f\in\funs{X}{Y}$ and
        $d_X$ is a metric on $X$ and
        $d_Y$ is a metric on $Y$ and
        for all $\epsilon\in\reals$ such that $0 < \epsilon$
            there exists $\delta\in\reals$ such that $0 < \delta$ and
                for all $x_0,x_1\in X$ such that $d_X((x_0,x_1)) < \delta$
                    we have $d_Y((\apply{f}{x_0},\apply{f}{x_1})) < \epsilon$.
\end{definition}

% from §27 Item 9: uniform continuity theorem
\begin{theorem}\label{uniform_continuity_theorem}
    Suppose $d_X$ is a metric on $X$.
    Suppose $d_Y$ is a metric on $Y$.
    Let $T_X = \metricTopology{d_X}{X}$.
    Let $T_Y = \metricTopology{d_Y}{Y}$.
    Suppose $f$ is continuous from $X$ to $Y$ with respect to $T_X$ and $T_Y$.
    Suppose $X$ is compact with respect to $T_X$.
    Then $f$ is uniformly continuous from $X$ to $Y$ with respect to $d_X$ and $d_Y$.
\end{theorem}
\begin{proof}
    $f\in\funs{X}{Y}$ by \cref{continuous,funs}.
    Show that for all $\epsilon\in\reals$ such that $0 < \epsilon$
        there exists $\delta\in\reals$ such that $0 < \delta$ and
            for all $x_0,x_1\in X$ such that $d_X((x_0,x_1)) < \delta$
                we have $d_Y((\apply{f}{x_0},\apply{f}{x_1})) < \epsilon$.
    \begin{subproof}
        Fix $\epsilon\in\reals$.
        Assume $0 < \epsilon$.
        Let $\epsilon_0 = \epsilon / (1 + 1)$.
        $1\in\reals$ by \cref{real_mul_ident}.
        $0\in\reals$ by \cref{real_add_ident}.
        $0 < 1$ by \cref{real_zero_lt_one}.
        $1 < 1 + 1$ by \cref{real_lt_add,real_add_ident}.
        $(1 + 1)\in\reals$ and $0 < 1 + 1$ by \cref{real_add_closed,real_lt_trans}.
        $(1 + 1)\neq 0$ by \cref{real_pos_nonzero}.
        $\inv{(1 + 1)}\in\reals$ and $\inv{(1 + 1)} > 0$ by \cref{real_mul_inverse,real_inv_pos}.
        $\epsilon_0 = \epsilon \rmul \inv{(1 + 1)}$ and $\epsilon_0\in\reals$ by \cref{real_div,real_mul_closed}.
        $0 < \epsilon_0$ by \cref{real_div,real_mul_zero,real_lt_mul_pos,real_lt_subst_left}.
        \begin{byCase}
        \caseOf{$X = \emptyset$.}
            Let $\delta_0 = \epsilon_0$.
            $\delta_0\in\reals$ and $0 < \delta_0$.
            Show that for all $x_0,x_1$ such that $x_0\in X$ and $x_1\in X$ and $d_X((x_0,x_1)) < \delta_0$
                we have $d_Y((\apply{f}{x_0},\apply{f}{x_1})) < \epsilon$.
            \begin{subproof}
                Trivial.
            \end{subproof}
            $\delta_0\in\reals$ and $0 < \delta_0$ and
                for all $x_0,x_1\in X$ such that $d_X((x_0,x_1)) < \delta_0$
                    we have $d_Y((\apply{f}{x_0},\apply{f}{x_1})) < \epsilon$.
            Show that there exists $\delta\in\reals$ such that $0 < \delta$ and
                for all $x_0,x_1\in X$ such that $d_X((x_0,x_1)) < \delta$
                    we have $d_Y((\apply{f}{x_0},\apply{f}{x_1})) < \epsilon$.
            \begin{subproof}
                Trivial.
            \end{subproof}
        \caseOf{$X \neq \emptyset$.}
            Let $R = \{w\in X\times\reals \mid \text{$0 < \snd{w}$ and
                for all $y\in X$ if $d_X((\fst{w},y)) < \snd{w}$ then
                $d_Y((\apply{f}{\fst{w}},\apply{f}{y})) < \epsilon_0$}\}$.
            $R$ is a relation by \cref{times_elem_is_tuple,relation}.
            For all $x\in X$ there exists $\delta\in\reals$ such that $x\mathrel{R}\delta$
                by \cref{metric_continuity_eps_delta,times_tuple_intro,fst_eq,snd_eq,relation}.
            Take $g$ such that $g$ is a function on $X$ and for all $x\in X$ we have $x\mathrel{R}\apply{g}{x}$
                by \cref{choice_function}.
            Let $A = \{U\in\pow{X} \mid \exists x\in X. U = \metricBall{d_X}{X}{x}{\apply{g}{x}}\}$.

            For all $x\in X$ we have $0 < \apply{g}{x}$ by \cref{choice_function,fst_eq,snd_eq,relation}.
            For all $x\in X$ we have $x\in\unions{A}$
                by \cref{metric_zero_characterization,metric_on,real_lt_subst_left,metric_ball,subseteq,powerset_intro,unions_iff}.
            For all $U$ such that $U\in A$ we have $U$ is open in $X$ with respect to $T_X$
                by \cref{choice_function,fst_eq,snd_eq,times_tuple_elim,metric_ball,subseteq,powerset_intro,metric_basis,metric_basis_is_basis,basis_elem_open_in_generated,basis_for_topology,metric_topology,open_in,basis_topology_is_topology}.
            $A$ is an open covering of $X$ with respect to $T_X$ by \cref{subseteq,covering,open_covering}.

            Take $\delta_1\in\reals$ such that $0 < \delta_1$ and
                for all $B$ such that $B\subseteq X$ and
                    for all $x,y\in B$ we have $d_X((x,y)) < \delta_1$
                there exists $U\in A$ such that $B\subseteq U$
                by \cref{lebesgue_number_lemma}.

            Show that for all $x_0,x_1$ such that $x_0\in X$ and $x_1\in X$ and $d_X((x_0,x_1)) < \delta_1$
                we have $d_Y((\apply{f}{x_0},\apply{f}{x_1})) < \epsilon$.
            \begin{subproof}
                Fix $x_0$.
                Fix $x_1$.
                Assume $x_0\in X$ and $x_1\in X$ and $d_X((x_0,x_1)) < \delta_1$.
        Let $B = \{x_0,x_1\}$.
                $B\subseteq X$ by \cref{upair_iff,subseteq}.
                Show that for all $x\in B$ we have for all $y\in B$ we have $d_X((x,y)) < \delta_1$.
                \begin{subproof}
                    Fix $x$.
                    Assume $x\in B$.
                    Fix $y$.
                    Assume $y\in B$.
                    $x = x_0$ or $x = x_1$ by \cref{upair_iff}.
                    $y = x_0$ or $y = x_1$ by \cref{upair_iff}.
                    \begin{byCase}
                    \caseOf{$x = x_0$.}
                        \begin{byCase}
                        \caseOf{$y = x_0$.}
                            $d_X((x,y)) < \delta_1$ by \cref{metric_zero_characterization,metric_on,real_lt_subst_left}.
                        \caseOf{$y \neq x_0$.}
                            $y = x_1$.
                            $d_X((x,y)) < \delta_1$.
                        \end{byCase}
                    \caseOf{$x \neq x_0$.}
                        $x = x_1$.
                        \begin{byCase}
                        \caseOf{$y = x_0$.}
                            $d_X((x_1,x_0)) = d_X((x_0,x_1))$ by \cref{metric_symmetric,metric_on}.
                            $d_X((x,y)) < \delta_1$.
                        \caseOf{$y \neq x_0$.}
                            $y = x_1$.
                            $d_X((x,y)) < \delta_1$ by \cref{metric_zero_characterization,metric_on,real_lt_subst_left}.
                        \end{byCase}
                    \end{byCase}
                \end{subproof}
                $B\subseteq X$ and for all $x,y\in B$ we have $d_X((x,y)) < \delta_1$.
                Take $U$ such that $U\in A$ and $B\subseteq U$ by \cref{lebesgue_number_lemma}.
                Take $x\in X$ such that $U = \metricBall{d_X}{X}{x}{\apply{g}{x}}$.
                $d_X((x,x_0)) < \apply{g}{x}$ and
                    $d_X((x,x_1)) < \apply{g}{x}$ by \cref{upair_intro_left,upair_intro_right,subseteq,metric_ball}.
                $(x,\apply{g}{x})\in R$ by \cref{choice_function,relation}.
                For all $y\in X$ if $d_X((x,y)) < \apply{g}{x}$ then
                    $d_Y((\apply{f}{x},\apply{f}{y})) < \epsilon_0$
                    by \cref{fst_eq,snd_eq}.
                $d_Y((\apply{f}{x},\apply{f}{x_0})) < \epsilon_0$ and
                    $d_Y((\apply{f}{x},\apply{f}{x_1})) < \epsilon_0$.
                $\apply{f}{x_0}\in Y$ and $\apply{f}{x}\in Y$ and $\apply{f}{x_1}\in Y$ by \cref{funs_type_apply}.
                $d_Y((\apply{f}{x_0},\apply{f}{x})) < \epsilon_0$ by \cref{metric_symmetric,metric_on,real_lt_subst_left}.
                Let $s = d_Y((\apply{f}{x_0},\apply{f}{x})) + d_Y((\apply{f}{x},\apply{f}{x_1}))$.
                $d_Y((\apply{f}{x_0},\apply{f}{x_1})) \leq s$ by \cref{metric_on,metric_triangle_inequality}.
                $d_Y\in\funs{Y\times Y}{\reals}$ by \cref{metric_on}.
                $(\apply{f}{x_0},\apply{f}{x})\in Y\times Y$ and
                    $(\apply{f}{x},\apply{f}{x_1})\in Y\times Y$ by \cref{times_tuple_intro}.
                $d_Y((\apply{f}{x_0},\apply{f}{x}))\in\reals$ and
                    $d_Y((\apply{f}{x},\apply{f}{x_1}))\in\reals$ by \cref{funs_type_apply}.
                $s < \epsilon_0 + \epsilon_0$ by \cref{real_lt_add_two}.
                $d_Y((\apply{f}{x_0},\apply{f}{x_1})) < \epsilon_0 + \epsilon_0$
                    by \cref{times_tuple_intro,funs_type_apply,real_add_closed,real_leq_split,real_lt_trans,real_lt_subst_left}.

                $\epsilon_0 = \epsilon \rmul \inv{(1 + 1)}$ by \cref{real_div}.
                $\epsilon_0 + \epsilon_0 =
                    (\epsilon \rmul \inv{(1 + 1)}) + (\epsilon \rmul \inv{(1 + 1)})$.
                $(\epsilon + \epsilon) \rmul \inv{(1 + 1)}
                    = (\epsilon \rmul \inv{(1 + 1)}) + (\epsilon \rmul \inv{(1 + 1)})$
                    by \cref{real_right_distrib}.
                $\epsilon_0 + \epsilon_0 = (\epsilon + \epsilon) \rmul \inv{(1 + 1)}$.
                $\epsilon \rmul (1 + 1) = (\epsilon \rmul 1) + (\epsilon \rmul 1)$ by \cref{real_distrib}.
                $\epsilon \rmul 1 = \epsilon$ by \cref{real_mul_ident,real_mul_comm}.
                $\epsilon \rmul (1 + 1) = \epsilon + \epsilon$.
                $(\epsilon + \epsilon) \rmul \inv{(1 + 1)} = (\epsilon \rmul (1 + 1)) \rmul \inv{(1 + 1)}$.
                $(\epsilon \rmul (1 + 1)) \rmul \inv{(1 + 1)}
                    = \epsilon \rmul ((1 + 1) \rmul \inv{(1 + 1)})$ by \cref{real_mul_assoc}.
                $(1 + 1) \rmul \inv{(1 + 1)} = 1$ by \cref{real_mul_inverse}.
                $\epsilon_0 + \epsilon_0 = \epsilon \rmul 1$.
                $\epsilon_0 + \epsilon_0 = \epsilon$ by \cref{real_mul_ident}.
                $d_Y((\apply{f}{x_0},\apply{f}{x_1})) < \epsilon$ by \cref{real_lt_subst_right}.
            \end{subproof}
            $\delta_1\in\reals$ and $0 < \delta_1$ and
                for all $x_0,x_1\in X$ such that $d_X((x_0,x_1)) < \delta_1$
                    we have $d_Y((\apply{f}{x_0},\apply{f}{x_1})) < \epsilon$.
            Take $\eta$ such that $\eta = \delta_1$.
            $\eta\in\reals$ and $0 < \eta$ and
                for all $x_0,x_1\in X$ such that $d_X((x_0,x_1)) < \eta$
                    we have $d_Y((\apply{f}{x_0},\apply{f}{x_1})) < \epsilon$.
            Show that there exists $\delta\in\reals$ such that $0 < \delta$ and
                for all $x_0,x_1\in X$ such that $d_X((x_0,x_1)) < \delta$
                    we have $d_Y((\apply{f}{x_0},\apply{f}{x_1})) < \epsilon$.
            \begin{subproof}
                Trivial.
            \end{subproof}
        \end{byCase}
    \end{subproof}
    $f$ is uniformly continuous from $X$ to $Y$ with respect to $d_X$ and $d_Y$ by \cref{uniformly_continuous}.
\end{proof}
% from §27 Item 10: isolated point
\begin{definition}\label{isolated_point}
    $x$ is an isolated point of $X$ with respect to $T$ iff
        $x\in X$ and $\{x\}$ is open in $X$ with respect to $T$.
\end{definition}

% from §27 helper lemma 12: shrink a nonempty open set away from a point
\begin{lemma}\label{open_set_shrink_away_point}
    Assume $T$ is a topology on $X$.
    Assume $X$ is Hausdorff with respect to $T$.
    Assume $U$ is open in $X$ with respect to $T$.
    Assume $U\neq\emptyset$.
    Assume $x\in X$.
    Assume for all $z\in X$ we have $z$ is not an isolated point of $X$ with respect to $T$.
    Then there exists $V$ such that
        $V$ is open in $X$ with respect to $T$ and
        $V\neq\emptyset$ and
        $V\subseteq U$ and
        $x\notin\closure{V}{X}{T}$.
\end{lemma}
\begin{proof}
    Show that there exists $y$ such that $y\in U$ and $y\neq x$.
    \begin{subproof}
        \begin{byCase}
        \caseOf{$x\in U$.}
            Suppose not.
            Show that for all $z$ such that $z\in U$ we have $z\in\{x\}$.
            \begin{subproof}
                Fix $z$ such that $z\in U$.
                $z = x$.
                $z\in\{x\}$ by \cref{singleton_intro}.
            \end{subproof}
            $U\subseteq\{x\}$ by \cref{subseteq}.
            $\{x\}\subseteq U$ by \cref{singleton_subset_intro}.
            $U = \{x\}$ by \cref{subseteq_antisymmetric}.
            $\{x\}$ is open in $X$ with respect to $T$ by \cref{open_in}.
            $x$ is an isolated point of $X$ with respect to $T$ by \cref{isolated_point}.
            Contradiction.
        \caseOf{$x\notin U$.}
            Take $y$ such that $y\in U$ by \cref{nonempty_inhabited}.
            $y\neq x$ by \cref{open_in,subseteq}.
            Show that there exists $z$ such that $z\in U$ and $z\neq x$.
            \begin{subproof}
                Trivial.
            \end{subproof}
        \end{byCase}
    \end{subproof}
    Take $y$ such that $y\in U$ and $y\neq x$.
    $U\in T$ by \cref{open_in}.
    $U\in\pow{X}$ by \cref{topology_on,subseteq}.
    $U\subseteq X$ by \cref{powerset_elim}.
    $y\in X$ by \cref{subseteq}.
    Take $W_1,W_2$ such that
        $W_1$ is a neighborhood of $x$ in $X$ with respect to $T$ and
        $W_2$ is a neighborhood of $y$ in $X$ with respect to $T$ and
        $W_1$ is disjoint from $W_2$ by \cref{hausdorff}.
    Let $V = U\inter W_2$.
    $W_2$ is open in $X$ with respect to $T$ by \cref{neighborhood}.
    $V$ is open in $X$ with respect to $T$ by \cref{open_in,topology_on,inter}.
    $y\in V$ by \cref{inter_intro,neighborhood}.
    $V\neq\emptyset$ by \cref{emptyset}.
    $V\subseteq U$ by \cref{inter_lower_left}.
    Show that for all $z$ such that $z\in V$ we have $z\in X\setminus W_1$.
    \begin{subproof}
        Fix $z$ such that $z\in V$.
        Suppose not.
        $z\in W_1$.
        $z\in U\inter W_2$ by \cref{inter}.
        $z\in W_2$ by \cref{inter_elim_right}.
        $z\in W_1\inter W_2$ by \cref{inter_intro}.
        Contradiction by \cref{disjoint}.
    \end{subproof}
    $V\subseteq X\setminus W_1$ by \cref{subseteq}.
    $W_1$ is open in $X$ with respect to $T$ by \cref{neighborhood}.
    $W_1\in T$ by \cref{open_in}.
    $W_1\subseteq X$ by \cref{topology_on,subseteq,powerset_elim}.
    $X\setminus W_1\subseteq X$ by \cref{setminus_subseteq}.
    $X\setminus (X\setminus W_1) = W_1$ by \cref{double_relative_complement}.
    $X\setminus W_1$ is closed in $X$ with respect to $T$ by \cref{closed_in,open_in,setminus_subseteq}.
    $V\in T$ by \cref{open_in}.
    $V\in\pow{X}$ by \cref{topology_on,subseteq}.
    $V\subseteq X$ by \cref{powerset_elim}.
    $\closure{V}{X}{T}\subseteq X\setminus W_1$ by \cref{closure_subseteq_closed}.
    $x\notin X\setminus W_1$ by \cref{neighborhood,setminus_elim_right}.
    $x\notin\closure{V}{X}{T}$ by \cref{subseteq}.
\end{proof}

% from §27 helper lemma 13: countable nonempty set is an image of naturals
\begin{lemma}\label{countable_nonempty_surj_naturals}
    Suppose $X$ is countable.
    Suppose $X\neq\emptyset$.
    Then there exists $g$ such that $g$ is a surjection from $\naturals$ to $X$.
\end{lemma}
\begin{proof}
    Take $f$ such that $f\in\inj{X}{\naturals}$ by \cref{countable}.
    $f\in\funs{X}{\naturals}$ and
        for all $x,y\in X$ such that $f(x) = f(y)$ we have $x = y$
        by \cref{inj}.
    Hence $f$ is a function by \cref{funs_is_function}.
    Therefore $f$ is a function from $X$ to $\naturals$ by \cref{funs}.
    Take $x_0$ such that $x_0\in X$ by \cref{nonempty_inhabited}.
    Let $A = \img{f}{X}$.
    Let $R_1 = \{w\in\naturals\times X \mid \text{$\fst{w}\in A$ and $f(\snd{w}) = \fst{w}$}\}$.
    Let $R_2 = \{w\in\naturals\times X \mid \text{$\fst{w}\notin A$ and $\snd{w} = x_0$}\}$.
    Let $R = R_1\union R_2$.
    Show that $R_1$ is a relation.
    \begin{subproof}
        For all $w$ such that $w\in R_1$ there exist $n,x$ such that $w = (n,x)$.
        Hence $R_1$ is a relation by \cref{relation}.
    \end{subproof}
    Show that $R_2$ is a relation.
    \begin{subproof}
        For all $w$ such that $w\in R_2$ there exist $n,x$ such that $w = (n,x)$.
        Thus $R_2$ is a relation by \cref{relation}.
    \end{subproof}
    Therefore $R$ is a relation by \cref{union_relations_is_relation}.
    $R_1\subseteq\naturals\times X$ and $R_2\subseteq\naturals\times X$ by \cref{subseteq}.
    We have $R\subseteq\naturals\times X$ by \cref{union_subsets_is_subset}.
    Show that for all $n\in\naturals$ there exists $x$ such that $n\mathrel{R}x$.
    \begin{subproof}
        Fix $n\in\naturals$.
        \begin{byCase}
        \caseOf{$n\in A$.}
            Take $x$ such that $x\in X$ and $(x,n)\in f$ by \cref{img_iff}.
            $f(x) = n$ by \cref{function_apply_intro}.
            $(n,x)\in\naturals\times X$ by \cref{times_tuple_intro}.
            Hence $(n,x)\in R_1$.
            Therefore $(n,x)\in R$ by \cref{union_intro_left}.
            Thus $n\mathrel{R}x$ by \cref{relation}.
        \caseOf{$n\notin A$.}
            $(n,x_0)\in\naturals\times X$ by \cref{times_tuple_intro}.
            Hence $(n,x_0)\in R_2$.
            Therefore $(n,x_0)\in R$ by \cref{union_intro_right}.
            Thus $n\mathrel{R}x_0$ by \cref{relation}.
        \end{byCase}
    \end{subproof}
    Take $g$ such that $g$ is a function on $\naturals$ and for all $n\in\naturals$ we have $n\mathrel{R}g(n)$
        by \cref{choice_function}.
    For all $n\in\naturals$ we have $(n,g(n))\in R$ by \cref{choice_function,relation}.
    For all $n\in\naturals$ we have $(n,g(n))\in\naturals\times X$ by \cref{subseteq}.
    For all $n\in\naturals$ we have $g(n)\in X$ by \cref{times_tuple_elim}.
    Show that $g$ is a function from $\naturals$ to $X$.
    \begin{subproof}
        Follows by \cref{funs_intro,funs}.
    \end{subproof}
    Show that for all $x\in X$ there exists $n\in\naturals$ such that $g(n) = x$.
    \begin{subproof}
        Fix $x\in X$.
        Let $n = f(x)$.
        $n\in\naturals$ by \cref{funs_type_apply}.
        $x\in\dom{f}$ by \cref{funs}.
        $(x,n)\in f$ by \cref{function_apply_elim}.
        $n\in A$ by \cref{img_iff}.
        Therefore $n\mathrel{R}g(n)$ by \cref{choice_function}.
        Hence $(n,g(n))\in R$ by \cref{relation}.
        Show that $f(g(n)) = n$.
        \begin{subproof}
            Show that $(n,g(n))\in R_1$.
            \begin{subproof}
                Suppose not.
                $(n,g(n))\in R_2$ by \cref{union_iff}.
                $\fst{(n,g(n))}\notin A$.
                Hence $n\notin A$ by \cref{fst_eq}.
                Contradiction.
            \end{subproof}
            $f(\snd{(n,g(n))}) = \fst{(n,g(n))}$.
            Thus $f(g(n)) = n$ by \cref{fst_eq,snd_eq}.
        \end{subproof}
        Thus $f(g(n)) = f(x)$.
        $g(n)\in X$ by \cref{funs_type_apply}.
        Therefore $g(n) = x$.
        Therefore there exists $m\in\naturals$ such that $g(m) = x$.
    \end{subproof}
    Therefore $g$ is a surjection from $\naturals$ to $X$ by \cref{surj}.
\end{proof}

% local compatibility interface for the old set-theoretic naturals proofs
\begin{abbreviation}\label{num_inductive_set}
    $A$ is an inductive set iff $\emptyset\in A$ and for all $a\in A$ we have $\suc{a}\in A$.
\end{abbreviation}

\begin{axiom}\label{num_emptyset_in_naturals}
    $\emptyset\in\naturals$.
\end{axiom}

\begin{axiom}\label{num_naturals_inductive_set}
    $\naturals$ is an inductive set.
\end{axiom}

\begin{axiom}\label{num_naturals_smallest_inductive_set}
    Let $A$ be an inductive set.
    Then $\naturals\subseteq A$.
\end{axiom}

% from §27 helper lemma 14: naturals are subsets of naturals
\begin{lemma}\label{natural_subseteq_naturals}
    Suppose $n\in\naturals$.
    Then $n\subseteq\naturals$.
\end{lemma}
\begin{proof}
    Let $A = \{k\in\naturals \mid \text{$k\subseteq\naturals$}\}$.
    $A\subseteq\naturals$ by \cref{subseteq}.
    $\emptyset\subseteq\naturals$ by \cref{emptyset_subseteq}.
    Hence $\emptyset\in A$ by \cref{num_emptyset_in_naturals}.
    Show that for all $k\in A$ we have $\suc{k}\in A$.
    \begin{subproof}
        Fix $k\in A$.
        We have $k\in\naturals$ and $k\subseteq\naturals$ by \cref{subseteq}.
        Hence $\suc{k}\in\naturals$ by \cref{num_naturals_inductive_set}.
        Show that for all $x$ such that $x\in\suc{k}$ we have $x\in\naturals$.
        \begin{subproof}
            Fix $x$ such that $x\in\suc{k}$.
            $x\in k$ or $x = k$ by \cref{suc_iff}.
            \begin{byCase}
            \caseOf{$x\in k$.}
                $x\in\naturals$ by \cref{subseteq}.
            \caseOf{$x = k$.}
                $x\in\naturals$.
            \end{byCase}
        \end{subproof}
        Hence $\suc{k}\subseteq\naturals$ by \cref{subseteq}.
        Thus $\suc{k}\in A$.
    \end{subproof}
    Hence $A$ is an inductive set.
    Therefore $\naturals\subseteq A$ by \cref{num_naturals_smallest_inductive_set}.
    Hence $n\in A$ by \cref{subseteq}.
    Thus $n\subseteq\naturals$ by \cref{subseteq}.
\end{proof}

% from §27 helper lemma 15: naturals are transitive sets
\begin{lemma}\label{natural_is_transitiveset}
    Suppose $n\in\naturals$.
    Then for all $a\in n$ we have $a\subseteq n$.
\end{lemma}
\begin{proof}
    Let $A = \{k\in\naturals \mid \text{for all $a\in k$ we have $a\subseteq k$}\}$.
    $A\subseteq\naturals$ by \cref{subseteq}.
    For all $a$ if $a\in\emptyset$ then $a\subseteq\emptyset$ by \cref{emptyset}.
    Hence $\emptyset\in A$ by \cref{num_emptyset_in_naturals}.
    Show that for all $k\in A$ we have $\suc{k}\in A$.
    \begin{subproof}
        Fix $k\in A$.
        $k\in\naturals$ and for all $a\in k$ we have $a\subseteq k$ by \cref{subseteq}.
        $\suc{k}\in\naturals$ by \cref{num_naturals_inductive_set}.
        Show that for all $x\in\suc{k}$ we have $x\subseteq\suc{k}$.
        \begin{subproof}
            Fix $x\in\suc{k}$.
            $x\in k$ or $x = k$ by \cref{suc_iff}.
            \begin{byCase}
            \caseOf{$x\in k$.}
                $x\subseteq k$.
                $k\subseteq\suc{k}$ by \cref{subseteq_self_suc_intro}.
                Hence $x\subseteq\suc{k}$ by \cref{subseteq_transitive}.
            \caseOf{$x = k$.}
                $k\subseteq\suc{k}$ by \cref{subseteq_self_suc_intro}.
                Hence $x\subseteq\suc{k}$.
            \end{byCase}
        \end{subproof}
        Thus $\suc{k}\in A$.
    \end{subproof}
    Hence $A$ is an inductive set.
    Therefore $\naturals\subseteq A$ by \cref{num_naturals_smallest_inductive_set}.
    Hence $n\in A$ by \cref{subseteq}.
    Thus for all $a\in n$ we have $a\subseteq n$ by \cref{subseteq}.
\end{proof}

\begin{lemma}\label{natural_elem_implies_subseteq}
    Suppose $n\in\naturals$.
    Suppose $m\in\naturals$.
    Suppose $n\in m$.
    Then $n\subseteq m$.
\end{lemma}
\begin{proof}
    For all $a\in m$ we have $a\subseteq m$ by \cref{natural_is_transitiveset}.
    Hence $n\subseteq m$ by \cref{subseteq}.
\end{proof}

% from §27 helper lemma 16: naturals are comparable by membership
\begin{lemma}\label{natural_comparability}
    For all natural numbers $m,n$ we have $m\in n$ or $n\in m$ or $m = n$.
\end{lemma}
\begin{proof}[Proof by \in-induction on $m$]
    Assume $m$ is a natural number.
    Show for all natural numbers $n$ we have $m\in n$ or $n\in m$ or $m = n$.
    \begin{subproof}[Proof by \in-induction on $n$]
        Assume $n$ is a natural number.
        For all $a\in m$ we have $a\subseteq m$ by \cref{natural_is_transitiveset}.
        For all $a\in n$ we have $a\subseteq n$ by \cref{natural_is_transitiveset}.
        Suppose not.
        Then $m\notin n$ and $n\notin m$ and $m\neq n$.
        Show that for all $x\in m$ we have $x\in n$.
        \begin{subproof}
            Fix $x\in m$.
            $x\in\naturals$ by \cref{natural_subseteq_naturals,subseteq}.
            $x\in n$ or $n\in x$ or $x = n$ by assumption.
            \begin{byCase}
            \caseOf{$x\in n$.}
                Trivial.
            \caseOf{$n\in x$.}
                $x\subseteq m$.
                Hence $n\in m$ by \cref{subseteq}.
                Contradiction.
            \caseOf{$x = n$.}
                Hence $n\in m$.
                Contradiction.
            \end{byCase}
        \end{subproof}
        Hence $m\subseteq n$ by \cref{subseteq}.
        Show that for all $y\in n$ we have $y\in m$.
        \begin{subproof}
            Fix $y\in n$.
            $y\in\naturals$ by \cref{natural_subseteq_naturals,subseteq}.
            $m\in y$ or $y\in m$ or $m = y$ by assumption.
            \begin{byCase}
            \caseOf{$y\in m$.}
                Trivial.
            \caseOf{$m\in y$.}
                $y\subseteq n$.
                Hence $m\in n$ by \cref{subseteq}.
                Contradiction.
            \caseOf{$m = y$.}
                Hence $m\in n$.
                Contradiction.
            \end{byCase}
        \end{subproof}
        Hence $n\subseteq m$ by \cref{subseteq}.
        Hence $m = n$ by \cref{subseteq_antisymmetric}.
        Contradiction.
    \end{subproof}
\end{proof}

% from §27 helper lemma 17: finite inhabited subset of naturals has a largest element
\begin{lemma}\label{finite_inhabited_naturals_largest}
    Suppose $F\subseteq\naturals$.
    Suppose $F$ is finite.
    Suppose $F$ is inhabited.
    Then there exists $m\in F$ such that for all $a\in F$ we have $a\in m$ or $a = m$.
\end{lemma}
\begin{proof}
    Let $R = \{w\in\naturals\times\naturals \mid \fst{w}\in\snd{w}\}$.
    Show that $R$ is a simple order relation on $\naturals$.
    \begin{subproof}
        $R\subseteq\naturals\times\naturals$ by \cref{subseteq}.
        Show that $R$ is a binary relation on $\naturals$.
        \begin{subproof}
            Trivial.
        \end{subproof}
        Show that for all $a,b,c$ such that $a\mathrel{R}b$ and $b\mathrel{R}c$ we have $a\mathrel{R}c$.
        \begin{subproof}
            Fix $a$.
            Fix $b$.
            Fix $c$.
            Assume $a\mathrel{R}b$ and $b\mathrel{R}c$.
            $(a,b)\in R$ and $(b,c)\in R$ by \cref{relation}.
            Hence $(a,b)\in\naturals\times\naturals$ and $(b,c)\in\naturals\times\naturals$ by \cref{subseteq}.
            $a\in\naturals$ and $b\in\naturals$ and $c\in\naturals$ by \cref{times_tuple_elim}.
            $\fst{(a,b)}\in\snd{(a,b)}$ and $\fst{(b,c)}\in\snd{(b,c)}$.
            $a\in b$ and $b\in c$ by \cref{fst_eq,snd_eq}.
            For all $d\in c$ we have $d\subseteq c$ by \cref{natural_is_transitiveset}.
            $b\subseteq c$ by \cref{subseteq}.
            Therefore $a\in c$ by \cref{subseteq}.
            $(a,c)\in\naturals\times\naturals$ by \cref{times_tuple_intro}.
            $\fst{(a,c)}\in\snd{(a,c)}$ by \cref{fst_eq,snd_eq}.
            Therefore $(a,c)\in R$.
            Hence $a\mathrel{R}c$ by \cref{relation}.
        \end{subproof}
        Hence $R$ is transitive by \cref{transitive}.
        Show that for all $a,b$ such that $a\mathrel{R}b$ we have it is wrong that $b\mathrel{R}a$.
        \begin{subproof}
            Fix $a$.
            Fix $b$.
            Assume $a\mathrel{R}b$.
            Suppose not.
            Hence $b\mathrel{R}a$.
            $(a,b)\in R$ and $(b,a)\in R$ by \cref{relation}.
            $\fst{(a,b)}\in\snd{(a,b)}$ and $\fst{(b,a)}\in\snd{(b,a)}$.
            Hence $a\in b$ and $b\in a$ by \cref{fst_eq,snd_eq}.
            Contradiction by \cref{in_asymmetric}.
        \end{subproof}
        Hence $R$ is asymmetric by \cref{asymmetric}.
        Show that for all $a,b$ such that $a\in\naturals$ and $b\in\naturals$ we have
            $a\mathrel{R}b$ or $b\mathrel{R}a$ or $a = b$.
        \begin{subproof}
            Fix $a$.
            Fix $b$.
            Assume $a\in\naturals$ and $b\in\naturals$.
            $a\in b$ or $b\in a$ or $a = b$ by \cref{natural_comparability}.
            \begin{byCase}
            \caseOf{$a\in b$.}
                $(a,b)\in\naturals\times\naturals$ by \cref{times_tuple_intro}.
                Hence $(a,b)\in R$.
                Thus $a\mathrel{R}b$ or $b\mathrel{R}a$ or $a = b$ by \cref{relation}.
            \caseOf{$b\in a$.}
                $(b,a)\in\naturals\times\naturals$ by \cref{times_tuple_intro}.
                Hence $(b,a)\in R$.
                Hence $a\mathrel{R}b$ or $b\mathrel{R}a$ or $a = b$ by \cref{relation}.
            \caseOf{$a = b$.}
                $a\mathrel{R}b$ or $b\mathrel{R}a$ or $a = b$.
            \end{byCase}
        \end{subproof}
        Hence $R$ is connex on $\naturals$ by \cref{connex}.
        Therefore $R$ is a simple order relation on $\naturals$ by \cref{simple_order_relation}.
    \end{subproof}
    Take $m$ such that $m$ is a largest element of $F$ with respect to $R$
        by \cref{finite_inhabited_largest_element}.
    Thus $m\in F$ and for all $a\in F$ we have $a\mathrel{R}m$ or $a = m$ by \cref{largest_element}.
    Show that for all $a\in F$ we have $a\in m$ or $a = m$.
    \begin{subproof}
        Fix $a\in F$.
        $a\mathrel{R}m$ or $a = m$ by \cref{largest_element}.
        \begin{byCase}
        \caseOf{$a\mathrel{R}m$.}
            $(a,m)\in R$ by \cref{relation}.
            Hence $\fst{(a,m)}\in\snd{(a,m)}$.
            Thus $a\in m$ by \cref{fst_eq,snd_eq}.
        \caseOf{$a = m$.}
            Thus $a\in m$ or $a = m$.
        \end{byCase}
    \end{subproof}
    Therefore there exists $n\in F$ such that for all $a\in F$ we have $a\in n$ or $a = n$.
\end{proof}

% from §27 helper definition 2: shrinking sequence up to a stage
\begin{definition}\label{shrinking_sequence_upto}
    $s$ is a shrinking sequence up to $n$ for $g$ and $c$ in $X$ with respect to $T$ iff
        $n\in\naturals$ and
        $s$ is a function on $\suc{n}$ and
        $s(\emptyset) = c((X,g(\emptyset)))$ and
        for all $k\in\suc{n}$ we have
            $s(k)$ is open in $X$ with respect to $T$ and
            $s(k)\neq\emptyset$ and
            $g(k)\notin\closure{s(k)}{X}{T}$ and
        for all $k\in n$ we have
            $s(\suc{k}) = c((s(k),g(\suc{k})))$ and
            $s(\suc{k})\subseteq s(k)$.
\end{definition}

% from §27 helper lemma 18: existence of finite shrinking sequences
\begin{lemma}\label{shrinking_sequence_exists}
    Assume $T$ is a topology on $X$.
    Assume $X\neq\emptyset$.
    Suppose $g$ is a function from $\naturals$ to $X$.
    Suppose for all $U,x$ such that
        $U$ is open in $X$ with respect to $T$ and
        $U\neq\emptyset$ and
        $x\in X$
    we have
        $c((U,x))$ is open in $X$ with respect to $T$ and
        $c((U,x))\neq\emptyset$ and
        $c((U,x))\subseteq U$ and
        $x\notin\closure{c((U,x))}{X}{T}$.
    Then for all $n\in\naturals$ there exists $s$ such that
        $s$ is a shrinking sequence up to $n$ for $g$ and $c$ in $X$ with respect to $T$.
\end{lemma}
\begin{proof}
    Fix $m\in\naturals$.
    Let $A = \{n\in\naturals \mid \text{there exists $s$ such that
        $s$ is a shrinking sequence up to $n$ for $g$ and $c$ in $X$ with respect to $T$}\}$.
    $X\in T$ by \cref{topology_on}.
    Hence $X$ is open in $X$ with respect to $T$ by \cref{open_in}.
    $\emptyset\in\naturals$ by \cref{num_emptyset_in_naturals}.
    $\emptyset\in\dom{g}$ by \cref{funs}.
    Therefore $g(\emptyset)\in X$ by \cref{funs_type_apply}.
    Let $V = c((X,g(\emptyset)))$.
    $V$ is open in $X$ with respect to $T$ and
        $V\neq\emptyset$ and
        $V\subseteq X$ and
        $g(\emptyset)\notin\closure{V}{X}{T}$ by assumption.
    Let $s_0 = \{(\emptyset,V)\}$.
    For all $w$ such that $w\in s_0$ there exist $x, y$ such that $w = (x,y)$ by \cref{singleton_iff,relation}.
    Hence $s_0$ is a relation by \cref{relation}.
    Show that for all $a,b,b'$ such that $(a,b)\in s_0$ and $(a,b')\in s_0$ we have $b = b'$.
    \begin{subproof}
        Fix $a$.
        Fix $b$.
        Fix $b'$.
        Assume $(a,b)\in s_0$ and $(a,b')\in s_0$.
        Hence $(a,b) = (\emptyset,V)$ and $(a,b') = (\emptyset,V)$ by \cref{singleton_iff}.
        Thus $b = V$ and $b' = V$ by \cref{pair_eq_iff}.
    \end{subproof}
    Hence $s_0$ is a function by \cref{rightunique}.
    Show that for all $x$ such that $x\in\dom{s_0}$ we have $x\in\{\emptyset\}$.
    \begin{subproof}
        Fix $x$ such that $x\in\dom{s_0}$.
        Take $y$ such that $(x,y)\in s_0$ by \cref{dom_iff}.
        Hence $(x,y) = (\emptyset,V)$ by \cref{singleton_iff}.
        Thus $x = \emptyset$ by \cref{pair_eq_iff}.
        Hence $x\in\{\emptyset\}$ by \cref{singleton_intro}.
    \end{subproof}
    We have $\dom{s_0}\subseteq\{\emptyset\}$ by \cref{subseteq}.
    $(\emptyset,V)\in s_0$ by \cref{singleton_intro}.
    $\emptyset\in\dom{s_0}$ by \cref{dom_intro}.
    Thus $\{\emptyset\}\subseteq\dom{s_0}$ by \cref{singleton_subset_intro}.
    Therefore $\dom{s_0} = \{\emptyset\}$ by \cref{subseteq_antisymmetric}.
    Show that for all $x$ such that $x\in\dom{s_0}$ we have $x\in\suc{\emptyset}$.
    \begin{subproof}
        Fix $x$ such that $x\in\dom{s_0}$.
        $x = \emptyset$ by \cref{subseteq,singleton_iff}.
        $x\in\suc{\emptyset}$ by \cref{suc_intro_self}.
    \end{subproof}
    Therefore $\dom{s_0}\subseteq\suc{\emptyset}$ by \cref{subseteq}.
    Show that for all $x$ such that $x\in\suc{\emptyset}$ we have $x\in\dom{s_0}$.
    \begin{subproof}
        Fix $x$ such that $x\in\suc{\emptyset}$.
        $x\in\emptyset$ or $x = \emptyset$ by \cref{suc_iff}.
        Thus $x = \emptyset$ by \cref{emptyset}.
        Therefore $x\in\dom{s_0}$ by \cref{subseteq,singleton_iff}.
    \end{subproof}
    Therefore $\suc{\emptyset}\subseteq\dom{s_0}$ by \cref{subseteq}.
    Therefore $\dom{s_0} = \suc{\emptyset}$ by \cref{subseteq_antisymmetric}.
    Show that $s_0$ is a function on $\suc{\emptyset}$.
    \begin{subproof}
        Trivial.
    \end{subproof}
    Show that $s_0(\emptyset) = c((X,g(\emptyset)))$.
    \begin{subproof}
        $(\emptyset,V)\in s_0$ by \cref{singleton_iff}.
        Therefore $s_0(\emptyset) = c((X,g(\emptyset)))$ by \cref{function_apply_intro}.
    \end{subproof}
    Show that for all $k\in\suc{\emptyset}$ we have
        $s_0(k)$ is open in $X$ with respect to $T$ and
        $s_0(k)\neq\emptyset$ and
        $g(k)\notin\closure{s_0(k)}{X}{T}$.
    \begin{subproof}
        Fix $k\in\suc{\emptyset}$.
        $k\in\emptyset$ or $k = \emptyset$ by \cref{suc_iff}.
        Hence $k = \emptyset$ by \cref{emptyset}.
        $(\emptyset,V)\in s_0$ by \cref{singleton_iff}.
        Hence $s_0(k) = V$ by \cref{function_apply_intro}.
        Thus $s_0(k)$ is open in $X$ with respect to $T$ and
            $s_0(k)\neq\emptyset$ and
            $g(k)\notin\closure{s_0(k)}{X}{T}$.
    \end{subproof}
    Show that for all $k\in\emptyset$ we have
        $s_0(\suc{k}) = c((s_0(k),g(\suc{k})))$ and
        $s_0(\suc{k})\subseteq s_0(k)$.
    \begin{subproof}
        Trivial.
    \end{subproof}
    Therefore $s_0$ is a shrinking sequence up to $\emptyset$ for $g$ and $c$ in $X$ with respect to $T$ by \cref{shrinking_sequence_upto,num_emptyset_in_naturals}.
    Thus $\emptyset\in A$ by \cref{num_emptyset_in_naturals,subseteq}.
    Show that for all $n\in A$ we have $\suc{n}\in A$.
    \begin{subproof}
        Fix $n\in A$.
        Take $s$ such that $s$ is a shrinking sequence up to $n$ for $g$ and $c$ in $X$ with respect to $T$ by \cref{subseteq}.
        We have $n\in\naturals$ and $s$ is a function on $\suc{n}$ and
            $s(\emptyset) = c((X,g(\emptyset)))$ and
            for all $k\in\suc{n}$ we have
                $s(k)$ is open in $X$ with respect to $T$ and
                $s(k)\neq\emptyset$ and
                $g(k)\notin\closure{s(k)}{X}{T}$ and
            for all $k\in n$ we have
                $s(\suc{k}) = c((s(k),g(\suc{k})))$ and
                $s(\suc{k})\subseteq s(k)$ by \cref{shrinking_sequence_upto}.
        Hence $\dom{s} = \suc{n}$ by definition.
        $\suc{n}\in\naturals$ by \cref{num_naturals_inductive_set}.
        $\suc{n}\in\dom{g}$ by \cref{funs}.
        Therefore $g(\suc{n})\in X$ by \cref{funs_type_apply}.
        $n\in\suc{n}$ by \cref{suc_intro_self}.
        We have $s(n)$ is open in $X$ with respect to $T$ and
            $s(n)\neq\emptyset$ and
            $g(n)\notin\closure{s(n)}{X}{T}$ by \cref{shrinking_sequence_upto}.
        Let $W = c((s(n),g(\suc{n})))$.
        $W$ is open in $X$ with respect to $T$ and
            $W\neq\emptyset$ and
            $W\subseteq s(n)$ and
            $g(\suc{n})\notin\closure{W}{X}{T}$ by assumption.
        Let $t = s\union\{(\suc{n},W)\}$.
        For all $w$ such that $w\in t$ there exist $x, y$ such that $w = (x,y)$ by \cref{union_iff,singleton_iff,relation}.
        Hence $t$ is a relation by \cref{relation}.
        Show that for all $a,b,b'$ such that $(a,b)\in t$ and $(a,b')\in t$ we have $b = b'$.
        \begin{subproof}
            Fix $a$.
            Fix $b$.
            Fix $b'$.
            Assume $(a,b)\in t$ and $(a,b')\in t$.
            \begin{byCase}
            \caseOf{$a = \suc{n}$.}
                $(a,b)\in s$ or $(a,b) = (\suc{n},W)$ by \cref{union_iff,singleton_iff}.
                Show that $(a,b)\notin s$.
                \begin{subproof}
                    Suppose not.
                    $a\in\dom{s}$ by \cref{dom_intro}.
                    Hence $a\in\suc{n}$.
                    Contradiction by \cref{in_irrefl}.
                \end{subproof}
                Hence $(a,b) = (\suc{n},W)$.
                $(a,b')\in s$ or $(a,b') = (\suc{n},W)$ by \cref{union_iff,singleton_iff}.
                Show that $(a,b')\notin s$.
                \begin{subproof}
                    Suppose not.
                    $a\in\dom{s}$ by \cref{dom_intro}.
                    Hence $a\in\suc{n}$.
                    Contradiction by \cref{in_irrefl}.
                \end{subproof}
                Hence $(a,b') = (\suc{n},W)$.
                Thus $b = W$ and $b' = W$ by \cref{pair_eq_iff}.
            \caseOf{$a \neq \suc{n}$.}
                $(a,b)\in s$ by \cref{union_iff,singleton_iff,pair_eq_iff}.
                $(a,b')\in s$ by \cref{union_iff,singleton_iff,pair_eq_iff}.
                Thus $b = b'$ by \cref{function_rightunique}.
            \end{byCase}
        \end{subproof}
        Hence $t$ is a function by \cref{rightunique}.
        Show that for all $x$ such that $x\in\dom{t}$ we have $x\in\suc{\suc{n}}$.
        \begin{subproof}
            Fix $x$ such that $x\in\dom{t}$.
            Take $y$ such that $(x,y)\in t$ by \cref{dom_iff}.
            $(x,y)\in s$ or $(x,y) = (\suc{n},W)$ by \cref{union_iff,singleton_iff}.
            \begin{byCase}
            \caseOf{$(x,y) = (\suc{n},W)$.}
                $x = \suc{n}$ by \cref{pair_eq_iff}.
                Hence $x\in\suc{\suc{n}}$ by \cref{suc_intro_self}.
            \caseOf{$(x,y)\in s$.}
                $x\in\dom{s}$ by \cref{dom_intro}.
                We have $x\in\suc{n}$.
                Therefore $x\in\suc{\suc{n}}$ by \cref{suc_intro_in}.
            \end{byCase}
        \end{subproof}
        Hence $\dom{t}\subseteq\suc{\suc{n}}$ by \cref{subseteq}.
        Show that for all $x$ such that $x\in\suc{\suc{n}}$ we have $x\in\dom{t}$.
        \begin{subproof}
            Fix $x$ such that $x\in\suc{\suc{n}}$.
            $x\in\suc{n}$ or $x = \suc{n}$ by \cref{suc_iff}.
            \begin{byCase}
            \caseOf{$x\in\suc{n}$.}
                $x\in\dom{s}$.
                Take $y$ such that $(x,y)\in s$ by \cref{dom_iff}.
                Hence $(x,y)\in t$ by \cref{union_intro_left}.
                Thus $x\in\dom{t}$ by \cref{dom_intro}.
            \caseOf{$x = \suc{n}$.}
                $(x,W) = (\suc{n},W)$ by \cref{pair_eq_iff}.
                Hence $(x,W)\in t$ by \cref{union_intro_right,singleton_intro}.
                Thus $x\in\dom{t}$ by \cref{dom_intro}.
            \end{byCase}
        \end{subproof}
        Hence $\suc{\suc{n}}\subseteq\dom{t}$ by \cref{subseteq}.
        Therefore $\dom{t} = \suc{\suc{n}}$ by \cref{subseteq_antisymmetric}.
        Show that $t$ is a function on $\suc{\suc{n}}$.
        \begin{subproof}
            Trivial.
        \end{subproof}
        Show that $t(\emptyset) = c((X,g(\emptyset)))$.
        \begin{subproof}
            $\emptyset\in n$ or $n\in\emptyset$ or $\emptyset = n$
                by \cref{natural_comparability,num_emptyset_in_naturals}.
            Thus $\emptyset\in n$ or $\emptyset = n$ by \cref{emptyset}.
            $\emptyset\in\suc{n}$ by \cref{suc_iff}.
            We have $\emptyset\in\dom{s}$.
            $\emptyset\in\suc{\suc{n}}$ by \cref{suc_intro_in}.
            Therefore $\emptyset\in\dom{t}$.
            $(\emptyset,t(\emptyset))\in t$ and $(\emptyset,s(\emptyset))\in s$ by \cref{function_apply_elim}.
            Thus $(\emptyset,s(\emptyset))\in t$ by \cref{union_intro_left}.
            Therefore $t(\emptyset) = s(\emptyset)$ by \cref{function_rightunique}.
            Thus $t(\emptyset) = c((X,g(\emptyset)))$ by \cref{shrinking_sequence_upto}.
        \end{subproof}
        Show that for all $k\in\suc{\suc{n}}$ we have
            $t(k)$ is open in $X$ with respect to $T$ and
            $t(k)\neq\emptyset$ and
            $g(k)\notin\closure{t(k)}{X}{T}$.
        \begin{subproof}
            Fix $k\in\suc{\suc{n}}$.
            $k\in\suc{n}$ or $k = \suc{n}$ by \cref{suc_iff}.
            \begin{byCase}
            \caseOf{$k = \suc{n}$.}
                $(k,W) = (\suc{n},W)$ by \cref{pair_eq_iff}.
                Hence $(k,W)\in t$ by \cref{union_intro_right,singleton_intro}.
                Hence $t(k) = W$ by \cref{function_apply_intro}.
                Thus $t(k)$ is open in $X$ with respect to $T$ and
                    $t(k)\neq\emptyset$ and
                    $g(k)\notin\closure{t(k)}{X}{T}$.
            \caseOf{$k\in\suc{n}$.}
                $k\in\dom{s}$.
                $(k,s(k))\in s$ by \cref{function_apply_elim}.
                Therefore $(k,s(k))\in t$ by \cref{union_intro_left}.
                Thus $t(k) = s(k)$ by \cref{function_apply_intro}.
                We have $t(k)$ is open in $X$ with respect to $T$ and
                    $t(k)\neq\emptyset$ and
                    $g(k)\notin\closure{t(k)}{X}{T}$ by \cref{shrinking_sequence_upto}.
            \end{byCase}
        \end{subproof}
        Show that for all $k\in\suc{n}$ we have
            $t(\suc{k}) = c((t(k),g(\suc{k})))$ and
            $t(\suc{k})\subseteq t(k)$.
        \begin{subproof}
            Fix $k\in\suc{n}$.
            $k\in n$ or $k = n$ by \cref{suc_iff}.
            \begin{byCase}
            \caseOf{$k = n$.}
                $(\suc{k},W) = (\suc{n},W)$ by \cref{pair_eq_iff}.
                Hence $(\suc{k},W)\in t$ by \cref{union_intro_right,singleton_intro}.
                Hence $t(\suc{k}) = W$ by \cref{function_apply_intro}.
                $k\in\dom{s}$.
                $(k,s(k))\in s$ by \cref{function_apply_elim}.
                Hence $(k,s(k))\in t$ by \cref{union_intro_left}.
                Hence $t(k) = s(k)$ by \cref{function_apply_intro}.
                Thus $t(\suc{k}) = c((t(k),g(\suc{k})))$ and $t(\suc{k})\subseteq t(k)$.
            \caseOf{$k\in n$.}
                We have $n\subseteq\naturals$ by \cref{natural_subseteq_naturals}.
                Hence $k\in\naturals$ by \cref{subseteq}.
                Therefore $\suc{k}\in\naturals$ by \cref{num_naturals_inductive_set}.
                $\suc{k}\in n$ or $n\in\suc{k}$ or $\suc{k} = n$ by \cref{natural_comparability}.
                Show that it is wrong that $n\in\suc{k}$.
                \begin{subproof}
                    Suppose not.
                    $n\in k$ or $n = k$ by \cref{suc_iff}.
                    \begin{byCase}
                    \caseOf{$n\in k$.}
                        Contradiction by \cref{in_asymmetric}.
                    \caseOf{$n = k$.}
                        $k\in k$.
                        Contradiction by \cref{in_irrefl}.
                    \end{byCase}
                \end{subproof}
                Thus $\suc{k}\in n$ or $\suc{k} = n$.
                $\suc{k}\in\suc{n}$ by \cref{suc_iff}.
                $\suc{k}\in\dom{s}$ and $k\in\dom{s}$.
                $(\suc{k},s(\suc{k}))\in s$ and $(k,s(k))\in s$
                    by \cref{function_apply_elim}.
                Therefore $(\suc{k},s(\suc{k}))\in t$ and $(k,s(k))\in t$ by \cref{union_intro_left}.
                Thus $t(\suc{k}) = s(\suc{k})$ and $t(k) = s(k)$ by \cref{function_apply_intro}.
                We have $t(\suc{k}) = c((t(k),g(\suc{k})))$ and $t(\suc{k})\subseteq t(k)$
                    by \cref{shrinking_sequence_upto}.
            \end{byCase}
        \end{subproof}
        Therefore $t$ is a shrinking sequence up to $\suc{n}$ for $g$ and $c$ in $X$ with respect to $T$
            by \cref{shrinking_sequence_upto}.
        Thus $\suc{n}\in A$ by \cref{num_naturals_inductive_set,subseteq}.
    \end{subproof}
    Hence $A$ is an inductive set.
    Therefore $\naturals\subseteq A$ by \cref{num_naturals_smallest_inductive_set}.
    Hence $m\in A$ by \cref{subseteq}.
    Thus there exists $s$ such that $s$ is a shrinking sequence up to $m$ for $g$ and $c$ in $X$ with respect to $T$ by \cref{subseteq}.
\end{proof}

% from §27 helper lemma 19: restriction of a shrinking sequence
\begin{lemma}\label{shrinking_sequence_restriction}
    Suppose $s$ is a shrinking sequence up to $m$ for $g$ and $c$ in $X$ with respect to $T$.
    Suppose $n\in m$ or $n = m$.
    Then $\restrl{s}{\suc{n}}$ is a shrinking sequence up to $n$ for $g$ and $c$ in $X$ with respect to $T$.
\end{lemma}
\begin{proof}
    We have $m\in\naturals$ and $s$ is a function on $\suc{m}$ and
        $s(\emptyset) = c((X,g(\emptyset)))$ and
        for all $k\in\suc{m}$ we have
            $s(k)$ is open in $X$ with respect to $T$ and
            $s(k)\neq\emptyset$ and
            $g(k)\notin\closure{s(k)}{X}{T}$ and
        for all $k\in m$ we have
            $s(\suc{k}) = c((s(k),g(\suc{k})))$ and
            $s(\suc{k})\subseteq s(k)$ by \cref{shrinking_sequence_upto}.
    Hence $\dom{s} = \suc{m}$ by definition.
    Show that $n\in\naturals$ and $n\subseteq m$ and $\suc{n}\subseteq\suc{m}$.
    \begin{subproof}
        \begin{byCase}
        \caseOf{$n\in m$.}
            $n\in\naturals$ by \cref{natural_subseteq_naturals,subseteq}.
            Hence $n\subseteq m$ by \cref{natural_elem_implies_subseteq}.
            Thus $\suc{n}\subseteq m$ by \cref{suc_subseteq_intro}.
            We have $m\subseteq\suc{m}$ by \cref{suc,suc_intro_in,subseteq}.
            Therefore $\suc{n}\subseteq\suc{m}$ by \cref{subseteq_transitive}.
        \caseOf{$n = m$.}
            $n\in\naturals$.
            Hence $n\subseteq m$ by \cref{subseteq_refl}.
            Thus $\suc{n}\subseteq\suc{m}$ by \cref{subseteq_refl}.
        \end{byCase}
    \end{subproof}
    Let $t = \restrl{s}{\suc{n}}$.
    $t$ is a function by \cref{restrl_of_function_is_function}.
    We have $\dom{t} = \dom{s}\inter\suc{n}$ by \cref{dom_of_restrl_eq_inter}.
    Hence $\dom{t} = \suc{n}$ by \cref{inter_absorb_supseteq_left}.
    Show that $t$ is a function on $\suc{n}$.
    \begin{subproof}
        Trivial.
    \end{subproof}
    Show that $t(\emptyset) = c((X,g(\emptyset)))$.
    \begin{subproof}
        $\emptyset\in n$ or $n\in\emptyset$ or $\emptyset = n$
            by \cref{natural_comparability,num_emptyset_in_naturals}.
        Hence $\emptyset\in n$ or $\emptyset = n$ by \cref{emptyset}.
        $\emptyset\in\suc{n}$ by \cref{suc_iff}.
        We have $\emptyset\in\dom{t}$ and $\emptyset\in\dom{s}$.
        Thus $t(\emptyset) = s(\emptyset)$ by \cref{restrl_of_function_apply}.
        Therefore $t(\emptyset) = c((X,g(\emptyset)))$ by \cref{shrinking_sequence_upto}.
    \end{subproof}
    Show that for all $k\in\suc{n}$ we have
        $t(k)$ is open in $X$ with respect to $T$ and
        $t(k)\neq\emptyset$ and
        $g(k)\notin\closure{t(k)}{X}{T}$.
    \begin{subproof}
        Fix $k\in\suc{n}$.
        We have $k\in\dom{t}$ and $k\in\dom{s}$.
        Hence $t(k) = s(k)$ by \cref{restrl_of_function_apply}.
        Therefore $t(k)$ is open in $X$ with respect to $T$ and
            $t(k)\neq\emptyset$ and
            $g(k)\notin\closure{t(k)}{X}{T}$ by \cref{shrinking_sequence_upto}.
    \end{subproof}
    Show that for all $k\in n$ we have
        $t(\suc{k}) = c((t(k),g(\suc{k})))$ and
        $t(\suc{k})\subseteq t(k)$.
    \begin{subproof}
        Fix $k\in n$.
        $k\in\suc{n}$ by \cref{suc_intro_in}.
        $k\in\naturals$ by \cref{natural_subseteq_naturals,subseteq}.
        Hence $\suc{k}\in\naturals$ by \cref{num_naturals_inductive_set}.
        $\suc{k}\in n$ or $n\in\suc{k}$ or $\suc{k} = n$ by \cref{natural_comparability}.
        Show that it is wrong that $n\in\suc{k}$.
        \begin{subproof}
            Suppose not.
            $n\in k$ or $n = k$ by \cref{suc_iff}.
            \begin{byCase}
            \caseOf{$n\in k$.}
                Contradiction by \cref{in_asymmetric}.
            \caseOf{$n = k$.}
                $k\in k$.
                Contradiction by \cref{in_irrefl}.
            \end{byCase}
        \end{subproof}
        Hence $\suc{k}\in n$ or $\suc{k} = n$.
        $\suc{k}\in\suc{n}$ by \cref{suc_iff}.
        Therefore $k\in\dom{t}$ and $\suc{k}\in\dom{t}$.
        Thus $k\in\dom{s}$ and $\suc{k}\in\dom{s}$.
        $k\in m$ by \cref{subseteq}.
        Therefore $t(k) = s(k)$ and $t(\suc{k}) = s(\suc{k})$ by \cref{restrl_of_function_apply}.
        We have $t(\suc{k}) = c((t(k),g(\suc{k})))$ and $t(\suc{k})\subseteq t(k)$
            by \cref{shrinking_sequence_upto}.
    \end{subproof}
    Therefore $t$ is a shrinking sequence up to $n$ for $g$ and $c$ in $X$ with respect to $T$
        by \cref{shrinking_sequence_upto}.
\end{proof}

% from §27 helper lemma 20: uniqueness of shrinking sequences
\begin{lemma}\label{shrinking_sequence_unique}
    Suppose $s$ is a shrinking sequence up to $n$ for $g$ and $c$ in $X$ with respect to $T$.
    Suppose $t$ is a shrinking sequence up to $n$ for $g$ and $c$ in $X$ with respect to $T$.
    Then $s = t$.
\end{lemma}
\begin{proof}
    Let $A = \{m\in\naturals \mid \text{for all $s_0,t_0$ such that
        $s_0$ is a shrinking sequence up to $m$ for $g$ and $c$ in $X$ with respect to $T$ and
        $t_0$ is a shrinking sequence up to $m$ for $g$ and $c$ in $X$ with respect to $T$
        we have $s_0 = t_0$}\}$.
    $A\subseteq\naturals$ by \cref{subseteq}.
    Show that $\emptyset\in A$.
    \begin{subproof}
        Show that for all $s_0,t_0$ such that
            $s_0$ is a shrinking sequence up to $\emptyset$ for $g$ and $c$ in $X$ with respect to $T$ and
            $t_0$ is a shrinking sequence up to $\emptyset$ for $g$ and $c$ in $X$ with respect to $T$
        we have $s_0 = t_0$.
        \begin{subproof}
            Fix $s_0$.
            Fix $t_0$.
            Assume $s_0$ is a shrinking sequence up to $\emptyset$ for $g$ and $c$ in $X$ with respect to $T$ and
                $t_0$ is a shrinking sequence up to $\emptyset$ for $g$ and $c$ in $X$ with respect to $T$.
            We have $s_0$ is a function on $\suc{\emptyset}$ and
                $s_0(\emptyset) = c((X,g(\emptyset)))$ by \cref{shrinking_sequence_upto}.
            $t_0$ is a function on $\suc{\emptyset}$ and
                $t_0(\emptyset) = c((X,g(\emptyset)))$ by \cref{shrinking_sequence_upto}.
            Hence $\dom{s_0} = \suc{\emptyset}$ and $\dom{t_0} = \suc{\emptyset}$ by definition.
            Show that for all $x$ such that $x\in\dom{s_0}$ we have $s_0(x) = t_0(x)$.
            \begin{subproof}
                Fix $x$ such that $x\in\dom{s_0}$.
                $x\in\emptyset$ or $x = \emptyset$ by \cref{suc_iff}.
                Hence $x = \emptyset$ by \cref{emptyset}.
                Thus $s_0(x) = c((X,g(\emptyset)))$ and $t_0(x) = c((X,g(\emptyset)))$
                    by \cref{shrinking_sequence_upto}.
            \end{subproof}
            Show that $s_0\subseteq t_0$.
            \begin{subproof}
                $\dom{s_0}\subseteq\dom{t_0}$ by \cref{subseteq_refl}.
                Follows by \cref{fun_subseteq}.
            \end{subproof}
            Show that for all $x\in\dom{t_0}$ we have $t_0(x) = s_0(x)$.
            \begin{subproof}
                Fix $x\in\dom{t_0}$.
                $x\in\dom{s_0}$ by \cref{subseteq}.
                Hence $s_0(x) = t_0(x)$ by \cref{subseteq}.
            \end{subproof}
            Show that $t_0\subseteq s_0$.
            \begin{subproof}
                $\dom{t_0}\subseteq\dom{s_0}$ by \cref{subseteq_refl}.
                Follows by \cref{fun_subseteq}.
            \end{subproof}
            Therefore $s_0 = t_0$ by \cref{subseteq_antisymmetric}.
        \end{subproof}
        Hence $\emptyset\in A$ by \cref{num_emptyset_in_naturals}.
    \end{subproof}
    Show that for all $m\in A$ we have $\suc{m}\in A$.
    \begin{subproof}
        Fix $m\in A$.
        $m\in\naturals$ by \cref{subseteq}.
        Show that for all $s_0,t_0$ such that
            $s_0$ is a shrinking sequence up to $\suc{m}$ for $g$ and $c$ in $X$ with respect to $T$ and
            $t_0$ is a shrinking sequence up to $\suc{m}$ for $g$ and $c$ in $X$ with respect to $T$
        we have $s_0 = t_0$.
        \begin{subproof}
            Fix $s_0$.
            Fix $t_0$.
            Assume $s_0$ is a shrinking sequence up to $\suc{m}$ for $g$ and $c$ in $X$ with respect to $T$ and
                $t_0$ is a shrinking sequence up to $\suc{m}$ for $g$ and $c$ in $X$ with respect to $T$.
            Let $s_1 = \restrl{s_0}{\suc{m}}$.
            Let $t_1 = \restrl{t_0}{\suc{m}}$.
            We have $s_1$ is a shrinking sequence up to $m$ for $g$ and $c$ in $X$ with respect to $T$ and
                $t_1$ is a shrinking sequence up to $m$ for $g$ and $c$ in $X$ with respect to $T$
                by \cref{shrinking_sequence_restriction,suc_intro_self}.
            Hence $s_1 = t_1$ by \cref{subseteq}.
            $s_0$ is a function on $\suc{\suc{m}}$ by \cref{shrinking_sequence_upto}.
            $t_0$ is a function on $\suc{\suc{m}}$ by \cref{shrinking_sequence_upto}.
            Hence $\dom{s_0} = \suc{\suc{m}}$ and $\dom{t_0} = \suc{\suc{m}}$ by definition.
            $\suc{m}\subseteq\suc{\suc{m}}$ by \cref{subseteq_self_suc_intro}.
            Hence $\suc{m}\subseteq\dom{s_0}$ and $\suc{m}\subseteq\dom{t_0}$ by \cref{subseteq}.
            We have $\dom{s_1} = \dom{s_0}\inter\suc{m}$ and $\dom{t_1} = \dom{t_0}\inter\suc{m}$
                by \cref{dom_of_restrl_eq_inter}.
            Hence $\dom{s_1} = \suc{m}$ and $\dom{t_1} = \suc{m}$ by \cref{inter_absorb_supseteq_left}.
            Show that for all $x$ such that $x\in\dom{s_0}$ we have $s_0(x) = t_0(x)$.
            \begin{subproof}
                Fix $x$ such that $x\in\dom{s_0}$.
                $x\in\suc{m}$ or $x = \suc{m}$ by \cref{suc_iff}.
                \begin{byCase}
                \caseOf{$x\in\suc{m}$.}
                    Hence $x\in\dom{s_1}$ and $x\in\dom{t_1}$ by \cref{subseteq}.
                    Hence $s_0(x) = s_1(x)$ and $t_0(x) = t_1(x)$ by \cref{restrl_of_function_apply}.
                    Thus $s_1(x) = t_1(x)$.
                    Therefore $s_0(x) = t_0(x)$.
                \caseOf{$x = \suc{m}$.}
                    $m\in\suc{m}$ by \cref{suc_intro_self}.
                    Hence $m\in\dom{s_1}$ and $m\in\dom{t_1}$ by \cref{subseteq}.
                    Therefore $s_0(m) = s_1(m)$ and $t_0(m) = t_1(m)$ by \cref{restrl_of_function_apply}.
                    Thus $s_0(m) = t_0(m)$.
                    Therefore $s_0(\suc{m}) = c((s_0(m),g(\suc{m})))$ by \cref{shrinking_sequence_upto}.
                    Hence $t_0(\suc{m}) = c((t_0(m),g(\suc{m})))$ by \cref{shrinking_sequence_upto}.
                    $s_0(x) = c((s_0(m),g(\suc{m})))$ and
                        $t_0(x) = c((t_0(m),g(\suc{m})))$.
                    Thus $s_0(x) = t_0(x)$.
                \end{byCase}
            \end{subproof}
            Show that $s_0\subseteq t_0$.
            \begin{subproof}
                $\dom{s_0}\subseteq\dom{t_0}$ by \cref{subseteq_refl}.
                Follows by \cref{fun_subseteq}.
            \end{subproof}
            Show that for all $x\in\dom{t_0}$ we have $t_0(x) = s_0(x)$.
            \begin{subproof}
                Fix $x\in\dom{t_0}$.
                $x\in\dom{s_0}$ by \cref{subseteq}.
                Hence $s_0(x) = t_0(x)$ by \cref{subseteq}.
            \end{subproof}
            Show that $t_0\subseteq s_0$.
            \begin{subproof}
                $\dom{t_0}\subseteq\dom{s_0}$ by \cref{subseteq_refl}.
                Follows by \cref{fun_subseteq}.
            \end{subproof}
            Therefore $s_0 = t_0$ by \cref{subseteq_antisymmetric}.
        \end{subproof}
        Hence $\suc{m}\in A$ by \cref{num_naturals_inductive_set}.
    \end{subproof}
    Hence $A$ is an inductive set.
    Therefore $\naturals\subseteq A$ by \cref{num_naturals_smallest_inductive_set}.
    $n\in\naturals$ by \cref{shrinking_sequence_upto}.
    Hence $n\in A$ by \cref{subseteq}.
    Thus $s = t$ by \cref{subseteq}.
\end{proof}

% from §27 helper lemma 21: later term of a shrinking sequence lies in every earlier stage
\begin{lemma}\label{shrinking_sequence_terms_nested}
    Suppose $s$ is a shrinking sequence up to $m$ for $g$ and $c$ in $X$ with respect to $T$.
    Suppose $n\in\suc{m}$.
    Then $s(m)\subseteq s(n)$.
\end{lemma}
\begin{proof}
    Let $A = \{k\in\naturals \mid \text{for all $s_0,n_0$ we have
        if $s_0$ is a shrinking sequence up to $k$ for $g$ and $c$ in $X$ with respect to $T$ and
            $n_0\in\suc{k}$
        then $s_0(k)\subseteq s_0(n_0)$}\}$.
    $A\subseteq\naturals$ by \cref{subseteq}.
    Show that $\emptyset\in A$.
    \begin{subproof}
        Show that for all $s_0,n_0$ we have
            if $s_0$ is a shrinking sequence up to $\emptyset$ for $g$ and $c$ in $X$ with respect to $T$ and
                $n_0\in\suc{\emptyset}$
            then $s_0(\emptyset)\subseteq s_0(n_0)$.
        \begin{subproof}
            Fix $s_0$.
            Fix $n_0$.
            Assume $s_0$ is a shrinking sequence up to $\emptyset$ for $g$ and $c$ in $X$ with respect to $T$ and
                $n_0\in\suc{\emptyset}$.
            $n_0\in\emptyset$ or $n_0 = \emptyset$ by \cref{suc_iff}.
            Hence $n_0 = \emptyset$ by \cref{emptyset}.
            Thus $s_0(\emptyset)\subseteq s_0(n_0)$ by \cref{subseteq_refl}.
        \end{subproof}
        Hence $\emptyset\in A$ by \cref{num_emptyset_in_naturals}.
    \end{subproof}
    Show that for all $k\in A$ we have $\suc{k}\in A$.
    \begin{subproof}
        Fix $k\in A$.
        $k\in\naturals$ by \cref{subseteq}.
        Show that for all $s_0,n_0$ we have
            if $s_0$ is a shrinking sequence up to $\suc{k}$ for $g$ and $c$ in $X$ with respect to $T$ and
                $n_0\in\suc{\suc{k}}$
            then $s_0(\suc{k})\subseteq s_0(n_0)$.
        \begin{subproof}
            Fix $s_0$.
            Fix $n_0$.
            Assume $s_0$ is a shrinking sequence up to $\suc{k}$ for $g$ and $c$ in $X$ with respect to $T$ and
                $n_0\in\suc{\suc{k}}$.
            $n_0\in\suc{k}$ or $n_0 = \suc{k}$ by \cref{suc_iff}.
            \begin{byCase}
            \caseOf{$n_0 = \suc{k}$.}
                Thus $s_0(\suc{k})\subseteq s_0(n_0)$ by \cref{subseteq_refl}.
            \caseOf{$n_0\in\suc{k}$.}
                We have $s_0$ is a function on $\suc{\suc{k}}$ and
                    $s_0(\suc{k})\subseteq s_0(k)$ by \cref{shrinking_sequence_upto,suc_intro_self}.
                Hence $\dom{s_0} = \suc{\suc{k}}$ by definition.
                $\suc{k}\subseteq\suc{\suc{k}}$ by \cref{subseteq_self_suc_intro}.
                Hence $\suc{k}\subseteq\dom{s_0}$ by \cref{subseteq}.
                Let $t_0 = \restrl{s_0}{\suc{k}}$.
                $t_0$ is a shrinking sequence up to $k$ for $g$ and $c$ in $X$ with respect to $T$
                    by \cref{shrinking_sequence_restriction,suc_intro_self}.
                \begin{byCase}
                \caseOf{$n_0 = k$.}
                    Thus $s_0(\suc{k})\subseteq s_0(n_0)$ by \cref{shrinking_sequence_upto,suc_intro_self}.
                \caseOf{$n_0\in k$.}
                    Hence $t_0(k)\subseteq t_0(n_0)$ by \cref{subseteq}.
                    $k\in\suc{k}$ by \cref{suc_intro_self}.
                    Therefore $n_0\in\suc{k}$ by \cref{suc_intro_in}.
                    Thus $t_0(k) = s_0(k)$ and $t_0(n_0) = s_0(n_0)$ by \cref{restrl_of_function_apply}.
                    We have $s_0(k)\subseteq s_0(n_0)$ by \cref{subseteq}.
                    Therefore $s_0(\suc{k})\subseteq s_0(n_0)$ by \cref{subseteq_transitive}.
                \end{byCase}
            \end{byCase}
        \end{subproof}
        Hence $\suc{k}\in A$ by \cref{num_naturals_inductive_set}.
    \end{subproof}
    Hence $A$ is an inductive set.
    Therefore $\naturals\subseteq A$ by \cref{num_naturals_smallest_inductive_set}.
    $m\in\naturals$ by \cref{shrinking_sequence_upto}.
    Hence $m\in A$ by \cref{subseteq}.
    Thus $s(m)\subseteq s(n)$ by \cref{subseteq}.
\end{proof}

% from §27 Item 11: compact Hausdorff without isolated points is uncountable
\begin{theorem}\label{compact_hausdorff_no_isolated_uncountable}
    Assume $T$ is a topology on $X$.
    Assume $X$ is compact with respect to $T$.
    Assume $X$ is Hausdorff with respect to $T$.
    Assume $X\neq\emptyset$.
    Suppose for all $x\in X$ we have $x$ is not an isolated point of $X$ with respect to $T$.
    Then $X$ is not countable.
\end{theorem}
\begin{proof}
    Suppose not.
    Take $g$ such that $g$ is a surjection from $\naturals$ to $X$
        by \cref{countable_nonempty_surj_naturals}.
    $g$ is a function from $\naturals$ to $X$ by \cref{surj_to_fun}.

    Let $D = \{p\in\pow{X}\times X \mid \text{$\fst{p}$ is open in $X$ with respect to $T$ and $\fst{p}\neq\emptyset$}\}$.
    Let $R = \{w\in D\times\pow{X} \mid \text{there exist $p,V$ such that
        $p\in D$ and
        $V\in\pow{X}$ and
        $w = (p,V)$ and
        $V$ is open in $X$ with respect to $T$ and
        $V\neq\emptyset$ and
        $V\subseteq\fst{p}$ and
        $\snd{p}\notin\closure{V}{X}{T}$}\}$.
    For all $w$ such that $w\in R$ there exist $p,V$ such that $w = (p,V)$ by \cref{relation}.
    Hence $R$ is a relation by \cref{relation}.
    Show that for all $p\in D$ there exists $V$ such that $p\mathrel{R}V$.
    \begin{subproof}
        Fix $p\in D$.
        $p\in\pow{X}\times X$ by \cref{subseteq}.
        Take $U,x$ such that $p = (U,x)$ and $U\in\pow{X}$ and $x\in X$ by \cref{times_elem_is_tuple}.
        $U = \fst{p}$ and $x = \snd{p}$ by \cref{fst_eq,snd_eq}.
        Hence $U$ is open in $X$ with respect to $T$ and $U\neq\emptyset$ by \cref{subseteq}.
        Take $V$ such that
            $V$ is open in $X$ with respect to $T$ and
            $V\neq\emptyset$ and
            $V\subseteq U$ and
            $x\notin\closure{V}{X}{T}$ by \cref{open_set_shrink_away_point}.
        $V\subseteq X$ by \cref{open_in,topology_on,subseteq,pow_iff}.
        Hence $V\in\pow{X}$ by \cref{powerset_intro}.
        $(p,V)\in D\times\pow{X}$ by \cref{times_tuple_intro}.
        Hence $(p,V)\in R$ by \cref{times_tuple_intro,powerset_intro}.
        Thus $p\mathrel{R}V$ by \cref{relation}.
    \end{subproof}
    Take $c$ such that $c$ is a function on $D$ and for all $p\in D$ we have $p\mathrel{R}c(p)$
        by \cref{choice_function}.

    Show that for all $U,x$ such that
        $U$ is open in $X$ with respect to $T$ and
        $U\neq\emptyset$ and
        $x\in X$
    we have
        $c((U,x))$ is open in $X$ with respect to $T$ and
        $c((U,x))\neq\emptyset$ and
        $c((U,x))\subseteq U$ and
        $x\notin\closure{c((U,x))}{X}{T}$.
    \begin{subproof}
        Fix $U$.
        Fix $x$.
        Assume $U$ is open in $X$ with respect to $T$ and $U\neq\emptyset$ and $x\in X$.
        $U\subseteq X$ by \cref{open_in,topology_on,subseteq,pow_iff}.
        Therefore $U\in\pow{X}$ by \cref{powerset_intro}.
        $(U,x)\in\pow{X}\times X$ by \cref{times_tuple_intro}.
        Thus $\fst{(U,x)} = U$ by \cref{fst_eq}.
        Therefore $\fst{(U,x)}$ is open in $X$ with respect to $T$ and $\fst{(U,x)}\neq\emptyset$.
        Thus $(U,x)\in D$ by \cref{times_tuple_intro}.
        We have $\dom{c} = D$ by definition.
        Therefore $(U,x)\in\dom{c}$ by \cref{subseteq}.
        $((U,x),c((U,x)))\in c$ by \cref{function_apply_elim}.
        Thus $((U,x),c((U,x)))\in R$ by \cref{function_apply_iff}.
        Take $p,V$ such that $p\in D$ and $V\in\pow{X}$ and
            $((U,x),c((U,x))) = (p,V)$ and
            $V$ is open in $X$ with respect to $T$ and
            $V\neq\emptyset$ and
            $V\subseteq\fst{p}$ and
            $\snd{p}\notin\closure{V}{X}{T}$ by \cref{subseteq}.
        $p = (U,x)$ and $c((U,x)) = V$ by \cref{pair_eq_iff}.
        Thus $\fst{p} = U$ and $\snd{p} = x$ by \cref{fst_eq,snd_eq}.
        Therefore $c((U,x))$ is open in $X$ with respect to $T$ and
            $c((U,x))\neq\emptyset$ and
            $c((U,x))\subseteq U$ and
            $x\notin\closure{c((U,x))}{X}{T}$.
    \end{subproof}

    Show that for all $n\in\naturals$ there exists $s$ such that
        $s$ is a shrinking sequence up to $n$ for $g$ and $c$ in $X$ with respect to $T$.
    \begin{subproof}
        Follows by \cref{shrinking_sequence_exists}.
    \end{subproof}

    Let $C = \{A\in\pow{X} \mid \text{there exists $n\in\naturals$ such that there exists $s$ such that
        $s$ is a shrinking sequence up to $n$ for $g$ and $c$ in $X$ with respect to $T$ and
        $A = \closure{s(n)}{X}{T}$}\}$.
    Show that $C$ has the finite intersection property.
    \begin{subproof}
        Show that $C$ is inhabited.
        \begin{subproof}
            $\emptyset\in\naturals$ by \cref{num_emptyset_in_naturals}.
            Take $s$ such that $s$ is a shrinking sequence up to $\emptyset$ for $g$ and $c$ in $X$ with respect to $T$ by \cref{subseteq}.
            Let $A_0 = \closure{s(\emptyset)}{X}{T}$.
            $\emptyset\in\suc{\emptyset}$ by \cref{suc_intro_self}.
            Hence $s(\emptyset)$ is open in $X$ with respect to $T$ by \cref{shrinking_sequence_upto}.
            Therefore $s(\emptyset)\subseteq X$ by \cref{open_in,topology_on,subseteq,pow_iff}.
            Thus $A_0$ is closed in $X$ with respect to $T$ by \cref{closure_closed_in}.
            Therefore $A_0\subseteq X$ by \cref{closed_in}.
            Thus $A_0\in\pow{X}$ by \cref{powerset_intro}.
            $A_0\in C$ by \cref{subseteq,num_emptyset_in_naturals}.
            Therefore $C$ is inhabited.
        \end{subproof}
        Show that for all $F$ such that $F\subseteq C$ and $F$ is finite and inhabited we have $\inters{F}\neq\emptyset$.
        \begin{subproof}
            Fix $F$ such that $F\subseteq C$ and $F$ is finite and inhabited.
            Let $R_0 = \{w\in F\times\naturals \mid \text{there exists $A\in F$ such that there exists $n\in\naturals$ such that there exists $s$ such that
                $w = (A,n)$ and
                $s$ is a shrinking sequence up to $n$ for $g$ and $c$ in $X$ with respect to $T$ and
                $A = \closure{s(n)}{X}{T}$}\}$.
            For all $w$ such that $w\in R_0$ there exist $A,n$ such that $w = (A,n)$ by \cref{times_elem_is_tuple,relation}.
            Hence $R_0$ is a relation by \cref{relation}.
            Show that for all $A\in F$ there exists $n$ such that $A\mathrel{R_0}n$.
            \begin{subproof}
                Fix $A\in F$.
                $A\in C$ by \cref{subseteq}.
                Take $n$ such that $n\in\naturals$ and there exists $s$ such that
                    $s$ is a shrinking sequence up to $n$ for $g$ and $c$ in $X$ with respect to $T$ and
                    $A = \closure{s(n)}{X}{T}$ by \cref{subseteq}.
                Take $s$ such that $s$ is a shrinking sequence up to $n$ for $g$ and $c$ in $X$ with respect to $T$ and
                    $A = \closure{s(n)}{X}{T}$.
                $(A,n)\in F\times\naturals$ by \cref{times_tuple_intro}.
                Hence $(A,n)\in R_0$ by \cref{times_tuple_intro}.
                Therefore $A\mathrel{R_0}n$ by \cref{relation}.
            \end{subproof}
            Take $h$ such that $h$ is a function on $F$ and for all $A\in F$ we have $A\mathrel{R_0}h(A)$
                by \cref{choice_function}.
            Let $N = \img{h}{F}$.
            We have $N\subseteq\naturals$ and $N$ is finite and $N$ is inhabited
                by \cref{img_iff,function_apply_iff,pair_eq_iff,subseteq,finite_image_function,nonempty_inhabited,function_apply_elim}.
            Take $m\in N$ such that for all $a\in N$ we have $a\in m$ or $a = m$
                by \cref{finite_inhabited_naturals_largest}.
            Take $u$ such that $u$ is a shrinking sequence up to $m$ for $g$ and $c$ in $X$ with respect to $T$ by \cref{subseteq}.
            $m\in\suc{m}$ by \cref{suc_intro_self}.
            Hence $u(m)\neq\emptyset$ by \cref{shrinking_sequence_upto}.
            Take $y$ such that $y\in u(m)$ by \cref{nonempty_inhabited}.
            Show that $y\in\inters{F}$.
            \begin{subproof}
                Show that for all $A\in F$ we have $y\in A$.
                \begin{subproof}
                    Fix $A\in F$.
                    $A\mathrel{R_0}h(A)$ by \cref{function_apply_iff}.
                    Hence $(A,h(A))\in R_0$ by \cref{relation}.
                    $h(A)\in N$ by \cref{function_apply_elim,img_iff}.
                    Therefore $h(A)\in m$ or $h(A) = m$ by \cref{subseteq}.
                    Hence $h(A)\in\suc{m}$ by \cref{suc_iff}.
                    Take $A_1,n_1,s_A$ such that $A_1\in F$ and $n_1\in\naturals$ and
                        $(A,h(A)) = (A_1,n_1)$ and
                        $s_A$ is a shrinking sequence up to $n_1$ for $g$ and $c$ in $X$ with respect to $T$ and
                        $A_1 = \closure{s_A(n_1)}{X}{T}$ by \cref{subseteq}.
                    $A = A_1$ and $h(A) = n_1$ by \cref{pair_eq_iff}.
                    Let $v = \restrl{u}{\suc{h(A)}}$.
                    $v$ is a shrinking sequence up to $h(A)$ for $g$ and $c$ in $X$ with respect to $T$
                        by \cref{shrinking_sequence_restriction}.
                    Hence $v = s_A$ by \cref{shrinking_sequence_unique}.
                    Therefore $u(m)\subseteq u(h(A))$ by \cref{shrinking_sequence_terms_nested}.
                    Thus $y\in u(h(A))$ by \cref{subseteq}.
                    Show that $\suc{h(A)}\subseteq\dom{u}$.
                    \begin{subproof}
                        We have $m\in\naturals$ and $u$ is a function on $\suc{m}$ by \cref{shrinking_sequence_upto}.
                        Hence $\dom{u} = \suc{m}$ by definition.
                        \begin{byCase}
                        \caseOf{$h(A)\in m$.}
                            $h(A)\in\naturals$ by \cref{subseteq}.
                            Hence $h(A)\subseteq m$ by \cref{natural_elem_implies_subseteq}.
                            $\suc{h(A)}\subseteq m$ by \cref{suc_subseteq_intro}.
                            $m\subseteq\suc{m}$ by \cref{subseteq_self_suc_intro}.
                            Therefore $\suc{h(A)}\subseteq\dom{u}$ by \cref{subseteq_transitive}.
                        \caseOf{$h(A) = m$.}
                            Thus $\suc{h(A)}\subseteq\dom{u}$ by \cref{subseteq_refl}.
                        \end{byCase}
                    \end{subproof}
                    $u$ is a function by \cref{shrinking_sequence_upto}.
                    $h(A)\in\suc{h(A)}$ by \cref{suc_intro_self}.
                    Hence $v(h(A)) = u(h(A))$ by \cref{restrl_of_function_apply}.
                    $u(h(A)) = s_A(h(A))$ by \cref{subseteq}.
                    Therefore $y\in s_A(h(A))$ by \cref{subseteq}.
                    $h(A)\in\suc{h(A)}$ by \cref{suc_intro_self}.
                    Hence $s_A(h(A))$ is open in $X$ with respect to $T$ by \cref{shrinking_sequence_upto}.
                    Hence $s_A(h(A))\subseteq X$ by \cref{open_in,topology_on,subseteq,pow_iff}.
                    Therefore $y\in\closure{s_A(h(A))}{X}{T}$ by \cref{subseteq_closure,elem_subseteq}.
                    Thus $y\in A$ by \cref{subseteq}.
                \end{subproof}
                Therefore $y\in\inters{F}$ by \cref{inters_intro}.
            \end{subproof}
            Hence $\inters{F}\neq\emptyset$ by \cref{emptyset}.
        \end{subproof}
        Hence $C$ has the finite intersection property by \cref{finite_intersection_property}.
    \end{subproof}

    Show that for all $A$ such that $A\in C$ we have $A$ is closed in $X$ with respect to $T$.
    \begin{subproof}
        Fix $A$ such that $A\in C$.
        Take $n$ such that $n\in\naturals$ and there exists $s$ such that
            $s$ is a shrinking sequence up to $n$ for $g$ and $c$ in $X$ with respect to $T$ and
            $A = \closure{s(n)}{X}{T}$ by \cref{subseteq}.
        Take $s$ such that $s$ is a shrinking sequence up to $n$ for $g$ and $c$ in $X$ with respect to $T$ and
            $A = \closure{s(n)}{X}{T}$.
        $n\in\suc{n}$ by \cref{suc_intro_self}.
        Hence $s(n)$ is open in $X$ with respect to $T$ by \cref{shrinking_sequence_upto}.
        Hence $s(n)\subseteq X$ by \cref{open_in,topology_on,subseteq,pow_iff}.
        Therefore $\closure{s(n)}{X}{T}$ is closed in $X$ with respect to $T$ by \cref{closure_closed_in}.
        Thus $A$ is closed in $X$ with respect to $T$ by \cref{subseteq}.
    \end{subproof}

    Therefore $\inters{C}\neq\emptyset$ by \cref{compact_iff_finite_intersection_property}.
    Take $x$ such that $x\in\inters{C}$ by \cref{nonempty_inhabited}.
    Take $A_0$ such that $A_0\in C$ by \cref{finite_intersection_property}.
    $x\in A_0$ by \cref{inters_destr}.
    Hence $x\in X$ by \cref{subseteq,closed_in}.
    Take $n\in\naturals$ such that $g(n) = x$ by \cref{surj}.
    Take $s$ such that $s$ is a shrinking sequence up to $n$ for $g$ and $c$ in $X$ with respect to $T$ by \cref{subseteq}.
    Let $B = \closure{s(n)}{X}{T}$.
    $n\in\suc{n}$ by \cref{suc_intro_self}.
    Hence $s(n)$ is open in $X$ with respect to $T$ by \cref{shrinking_sequence_upto}.
    Therefore $s(n)\subseteq X$ by \cref{open_in,topology_on,subseteq,pow_iff}.
    Thus $B$ is closed in $X$ with respect to $T$ by \cref{closure_closed_in}.
    Therefore $B\subseteq X$ by \cref{closed_in}.
    Thus $B\in\pow{X}$ by \cref{powerset_intro}.
    $B\in C$ by \cref{subseteq}.
    Thus $x\in B$ by \cref{inters_destr}.
    Therefore $g(n)\notin\closure{s(n)}{X}{T}$ by \cref{shrinking_sequence_upto}.
    Contradiction.
\end{proof}

% from §27 Item 12: closed interval in reals is uncountable
\begin{corollary}\label{real_closed_interval_uncountable}
    Suppose $a\in\reals$.
    Suppose $b\in\reals$.
    Suppose $a < b$.
    Let $R = \realOrderRel$.
    Let $T = \realOrderTopology$.
    Let $Y = \intervalclosedrel{R}{\reals}{a}{b}$.
    Let $T_Y = \subspaceTopology{T}{Y}$.
    Then $Y$ is not countable.
\end{corollary}
\begin{proof}
    $R$ is a simple order relation on $\reals$ and
        $\reals$ is a linear continuum with respect to $\realOrderRel$
        by \cref{real_order_relation_simple,reals_linear_continuum}.
    $\reals$ has more than one element by \cref{linear_continuum}.
    $T$ is the order topology on $\reals$ with respect to $R$ and $T$ is a topology on $\reals$
        by \cref{real_order_topology,order_topology,order_basis_is_basis,basis_topology_is_topology}.

    $a\mathrel{R}b$ by \cref{real_order_relation_intro}.
    $a\neq b$ by \cref{real_lt_subst_right,real_lt_irreflexive}.
    $a\in Y$ and $b\in Y$ by \cref{intervalclosed_rel}.
    $Y$ has more than one element by \cref{more_than_one_element}.
    $Y\neq\emptyset$ by \cref{emptyset}.

    $Y\subseteq\reals$ by \cref{intervalclosed_rel,subseteq}.

    $\reals$ is Hausdorff with respect to $T$ and $Y$ is Hausdorff with respect to $T_Y$
        by \cref{order_topology_hausdorff,subspace_hausdorff}.
    $T_Y$ is a topology on $Y$ by \cref{subspace_topology_is_topology}.

    $Y$ is compact with respect to $T_Y$ and $Y$ is connected with respect to $T_Y$
        by \cref{real_closed_interval_compact,linear_continuum_connected}.

    Show that for all $x\in Y$ we have $x$ is not an isolated point of $Y$ with respect to $T_Y$.
    \begin{subproof}
        Fix $x\in Y$.
        Suppose not.
        $\{x\}$ is open in $Y$ with respect to $T_Y$ by \cref{isolated_point}.
        Take $y$ such that $y\neq x$ by \cref{more_than_one_element_has_other}.
        $\{x\}\subseteq Y$ by \cref{singleton_subset_intro}.
        $\{x\}$ is closed in $Y$ with respect to $T_Y$
            by \cref{singleton_subset_intro,pair_is_finite,upair_intro_left,subseteq_finite_is_finite,finite_point_set_closed_hausdorff}.
        $\{x\} = \emptyset$ or $\{x\} = Y$ by \cref{connected_iff_clopen}.
        $\{x\} = Y$ by \cref{singleton_iff,emptyset}.
        $y = x$ by \cref{subseteq,singleton_iff}.
        Contradiction.
    \end{subproof}

    $Y$ is not countable by \cref{compact_hausdorff_no_isolated_uncountable}.
\end{proof}
\section{§28 Limit Point Compactness}

% from §28 Item 1: limit point compact
\begin{definition}\label{limit_point_compact}
    $X$ is limit point compact with respect to $T$ iff
        $T$ is a topology on $X$ and
        for all $A$ such that $A\subseteq X$ and $A$ is infinite
            there exists $x\in X$ such that $x$ is a limit point of $A$ in $X$ with respect to $T$.
\end{definition}

% from §28 Item 2: compact implies limit point compact
\begin{theorem}\label{compact_implies_limit_point_compact}
    Assume $T$ is a topology on $X$.
    Suppose $X$ is compact with respect to $T$.
    Then $X$ is limit point compact with respect to $T$.
\end{theorem}
\begin{proof}
    Show that for all $A$ such that $A\subseteq X$ and $A$ is infinite
        there exists $x\in X$ such that $x$ is a limit point of $A$ in $X$ with respect to $T$.
    \begin{subproof}
        Fix $A$.
        Assume $A\subseteq X$ and $A$ is infinite.
        Suppose not.
        For all $x\in X$ we have $x$ is not a limit point of $A$ in $X$ with respect to $T$.

        Let $C = \{U\in\pow{X} \mid \text{there exists $a\in A$ such that
            $U$ is a neighborhood of $a$ in $X$ with respect to $T$ and
            $U\inter (A\setminus\{a\}) = \emptyset$}\}$.

        For all $U$ such that $U\in C$ we have $U$ is open in $X$ with respect to $T$
            by \cref{neighborhood}.

        For all $a$ such that $a\in A$ we have $a\in\unions{C}$
            by \cref{subseteq,limit_point,intersects,neighborhood,open_in,topology_on,unions_iff}.
        $A\subseteq\unions{C}$ by \cref{subseteq}.

        For all $x$ such that $x\in\limitPoints{A}{X}{T}$ we have $x\in\emptyset$ by \cref{limit_points}.
        $X\setminus A$ is open in $X$ with respect to $T$
            by \cref{subseteq,emptyset_subseteq,subseteq_transitive,closed_iff_contains_limit_points,closed_in}.

        Let $D = C\union\{X\setminus A\}$.

        For all $U$ such that $U\in D$ we have $U$ is open in $X$ with respect to $T$
            by \cref{union_iff,singleton_iff,closed_in}.
        For all $x$ such that $x\in X$ we have $x\in\unions{D}$
            by \cref{subseteq,union_iff,unions_iff,setminus_intro,singleton_intro}.
        $D$ is an open covering of $X$ with respect to $T$ by \cref{subseteq,covering,open_covering}.

        Take $B$ such that $B\subseteq D$ and $B$ is finite and $B$ is a covering of $X$ by \cref{compact}.
        Let $B_A = \{U\in B \mid U\neq X\setminus A\}$.
        $B_A\subseteq B$ and $B_A$ is finite by \cref{subseteq,subseteq_finite_is_finite}.

        Show that for all $a$ such that $a\in A$ we have $a\in\unions{B_A}$.
        \begin{subproof}
            Fix $a$.
            Assume $a\in A$.
            Take $U\in B$ such that $a\in U$ by \cref{subseteq,covering,unions_iff}.
            $U\neq X\setminus A$.
            $a\in\unions{B_A}$ by \cref{unions_iff}.
        \end{subproof}
        $A\subseteq\unions{B_A}$ by \cref{subseteq}.

        Let $h = \{(U, U\inter A) \mid U\in B_A\}$.
        For all $U,V_1,V_2$ such that $(U,V_1)\in h$ and $(U,V_2)\in h$ we have $V_1 = V_2$
            by \cref{pair_eq_iff}.
        $h$ is right-unique by \cref{rightunique}.
        $h$ is a function.
        Show that for all $U$ such that $U\in B_A$ we have $U\in\dom{h}$.
        \begin{subproof}
            Follows by \cref{pair_eq_iff,dom_intro}.
        \end{subproof}
        $B_A\subseteq\dom{h}$ by \cref{subseteq}.
        Show that for all $U$ such that $U\in\dom{h}$ we have $U\in B_A$.
        \begin{subproof}
            Follows by \cref{dom_iff,pair_eq_iff}.
        \end{subproof}
        $\dom{h}\subseteq B_A$ by \cref{subseteq}.
        $\dom{h} = B_A$ by \cref{subseteq_antisymmetric}.
        $h$ is a function on $B_A$.

        Let $F = \img{h}{B_A}$.

        Show that for all $V$ such that $V\in F$ we have $V$ is finite.
        \begin{subproof}
            Fix $V$.
            Assume $V\in F$.
            Take $U\in B_A$ such that $(U,V)\in h$ by \cref{img_iff}.
            Take $U_0$ such that $(U,V) = (U_0,U_0\inter A)$ by \cref{pair_eq_iff}.
            Take $a$ such that
                $U$ is a neighborhood of $a$ in $X$ with respect to $T$ and
                $U\inter (A\setminus\{a\}) = \emptyset$
                by \cref{subseteq,union_iff,singleton_iff}.
            Show that for all $x$ such that $x\in V$ we have $x\in\{a\}$.
            \begin{subproof}
                Fix $x$.
                Assume $x\in V$.
                $x\in U$ and $x\in A$ by \cref{pair_eq_iff,inter}.
                \begin{byCase}
                \caseOf{$x = a$.}
                    $x\in\{a\}$ by \cref{singleton_intro}.
                \caseOf{$x\neq a$.}
                    $x\notin\{a\}$ and $x\in A\setminus\{a\}$ by \cref{singleton_iff,setminus_intro}.
                    $x\in U\inter (A\setminus\{a\})$ by \cref{inter_intro}.
                    Contradiction by \cref{emptyset}.
                \end{byCase}
            \end{subproof}
            $V$ is finite
                by \cref{subseteq,pair_is_finite,upair_intro_left,singleton_subset_intro,subseteq_finite_is_finite}.
        \end{subproof}

        For all $a$ such that $a\in A$ we have $a\in\unions{F}$
            by \cref{subseteq,unions_iff,pair_eq_iff,img_iff,inter_intro}.
        $A\subseteq\unions{F}$ by \cref{subseteq}.
        For all $x$ such that $x\in\unions{F}$ we have $x\in A$
            by \cref{unions_iff,img_iff,pair_eq_iff,inter}.
        $\unions{F}\subseteq A$ by \cref{subseteq}.
        $\unions{F} = A$ by \cref{subseteq_antisymmetric}.

        $A$ is finite by \cref{finite_image_function,finite_union_of_finite_family,subseteq_antisymmetric}.
        Contradiction.
    \end{subproof}
    $X$ is limit point compact with respect to $T$ by \cref{limit_point_compact}.
\end{proof}

% from §28 helper lemma 1: naturalsPlus is infinite
\begin{lemma}\label{naturalsplus_infinite}
    $\naturalsPlus$ is infinite.
\end{lemma}
\begin{proof}
    Suppose not.
    $\naturalsPlus$ is finite.
    Take $m\in\naturalsPlus$ such that for all $a\in\naturalsPlus$ we have $a \leq m$
        by \cref{real_one_in_naturals,emptyset,subseteq_refl,finite_naturals_maximum}.
    $m + 1 \leq m$ and $m\in\reals$ and $m + 1\in\reals$
        by \cref{real_naturals_closed_succ,real_naturals_subset_reals,subseteq}.
    $m < m + 1$
        by \cref{real_add_ident,real_one_in_naturals,real_naturals_subset_reals,subseteq,real_zero_lt_one,real_lt_add,real_add_comm}.
    Contradiction by \cref{real_leq_antisym,real_lt_subst_right,real_lt_irreflexive}.
\end{proof}

% from §28 Item 3: limit point compact does not imply compact
\begin{theorem}\label{limit_point_compact_not_compact}
    There exists $X,T$ such that
        $X$ is limit point compact with respect to $T$ and
        $X$ is not compact with respect to $T$.
\end{theorem}
\begin{proof}
    Let $Y = \{0,1\}$.
    Let $X = \naturalsPlus\times Y$.
    Let $T_0 = \discrete{\naturalsPlus}$.
    Let $T_1 = \{\emptyset, Y\}$.
    Let $T = \productTopology{T_0}{T_1}{\naturalsPlus}{Y}$.

    $T_0 = \pow{\naturalsPlus}$ and $T_0\subseteq\pow{\naturalsPlus}$
        by \cref{discrete_topology,subseteq_refl}.
    $\emptyset\in T_0$ and $\naturalsPlus\in T_0$ by \cref{discrete_topology,powerset_bottom,powerset_top}.
    $T_0$ is closed under arbitrary unions by \cref{subseteq_transitive,powerset_closed_unions,discrete_topology}.
    $T_0$ is closed under binary intersections by \cref{discrete_topology,subseteq,inter_powerset}.
    $T_0$ is a topology on $\naturalsPlus$ by \cref{topology_on}.

    $T_1\subseteq\pow{Y}$ by \cref{upair_iff,emptyset_subseteq,subseteq_refl,powerset_intro,subseteq}.
    $\emptyset\in T_1$ and $Y\in T_1$ by \cref{upair_intro_left,upair_intro_right}.
    Show that $T_1$ is closed under arbitrary unions.
    \begin{subproof}
        Fix $M$ such that $M\subseteq T_1$.
        \begin{byCase}
        \caseOf{$Y\in M$.}
            $\unions{M}\subseteq Y$
                by \cref{subseteq_transitive,powerset_closed_unions,pow_iff}.
            $Y\subseteq\unions{M}$ by \cref{unions_iff,subseteq}.
            $\unions{M}\in T_1$ by \cref{subseteq_antisymmetric,upair_intro_right}.
        \caseOf{$Y\notin M$.}
            For all $U$ such that $U\in M$ we have $U\subseteq\emptyset$
                by \cref{subseteq,upair_iff,subseteq_refl}.
            $\unions{M}\in T_1$ by \cref{unions_family,subseteq_emptyset_iff,upair_intro_left}.
        \end{byCase}
    \end{subproof}
    Show that $T_1$ is closed under binary intersections.
    \begin{subproof}
        Follows by \cref{upair_iff,inter_comm,inter_emptyset,inter_idempotent,upair_intro_left,upair_intro_right}.
    \end{subproof}
    $T_1$ is a topology on $Y$ by \cref{topology_on}.

    $\productBasis{T_0}{T_1}{\naturalsPlus}{Y}$ is a basis on $X$ by \cref{product_basis_is_basis}.
    $T$ is a topology on $X$ by \cref{basis_topology_is_topology,product_topology}.

    Show that for all $A$ such that $A\subseteq X$ and $A$ is infinite
        there exists $x\in X$ such that $x$ is a limit point of $A$ in $X$ with respect to $T$.
        \begin{subproof}
            Fix $A$ such that $A\subseteq X$ and $A$ is infinite.
            Take $a$ such that $a\in A$ by \cref{emptyset_is_finite,nonempty_inhabited}.
            Take $n,y$ such that $a = (n,y)$ and $n\in\naturalsPlus$ and $y\in Y$ by \cref{subseteq,times_elem_is_tuple}.

        Show that there exists $y'$ such that $y'\in Y$ and $y'\neq y$.
        \begin{subproof}
            $y = 0$ or $y = 1$ by \cref{upair_iff}.
            \begin{byCase}
            \caseOf{$y = 0$.}
                Let $z = 1$.
                There exists $y'$ such that $y'\in Y$ and $y'\neq y$.
            \caseOf{$y = 1$.}
                Let $z = 0$.
                There exists $y'$ such that $y'\in Y$ and $y'\neq y$.
            \end{byCase}
        \end{subproof}
        Take $y'$ such that $y'\in Y$ and $y'\neq y$.
        Let $x = (n,y')$.
        $x\in X$ by \cref{times_tuple_intro}.

        Show that for all $U$ such that $U$ is a neighborhood of $x$ in $X$ with respect to $T$
            we have $U$ intersects $A\setminus\{x\}$.
        \begin{subproof}
            Fix $U$ such that $U$ is a neighborhood of $x$ in $X$ with respect to $T$.
            $x\in U$ and $U\in\basisTopology{\productBasis{T_0}{T_1}{\naturalsPlus}{Y}}{X}$
                by \cref{neighborhood,open_in,product_topology}.
            Take $B\in\productBasis{T_0}{T_1}{\naturalsPlus}{Y}$ such that $x\in B$ and $B\subseteq U$
                by \cref{topology_generated_by_basis}.
            Take $U_0,V_0$ such that $V_0\in T_1$ and $B = U_0\times V_0$
                by \cref{product_basis}.
            $n\in U_0$ and $y'\in V_0$ by \cref{pair_eq_iff,times_tuple_elim}.
            $V_0 = Y$ by \cref{upair_iff,emptyset}.
            $B = U_0\times Y$.
            $(n,y)\in B$ by \cref{times_tuple_intro}.
            $a\in U$ by \cref{subseteq}.
            $a\in A\setminus\{x\}$ by \cref{pair_eq_iff,setminus_intro,singleton_iff}.
            $U$ intersects $A\setminus\{x\}$ by \cref{inter_intro,emptyset,intersects}.
        \end{subproof}
        $x$ is a limit point of $A$ in $X$ with respect to $T$ by \cref{limit_point}.
    \end{subproof}
    $X$ is limit point compact with respect to $T$ by \cref{limit_point_compact}.

    Show that $X$ is not compact with respect to $T$.
    \begin{subproof}
        Suppose not.
        $X$ is compact with respect to $T$.
        Let $\pi = \piOne{\naturalsPlus}{Y}$.
        $\pi$ is continuous from $X$ to $\naturalsPlus$ with respect to $T$ and $T_0$
            by \cref{projection_one_continuous}.
        $\img{\pi}{X}\subseteq\naturalsPlus$
            by \cref{img_iff,times_elem_is_tuple,projection_first,pair_eq_iff,subseteq}.
        $\naturalsPlus\subseteq\img{\pi}{X}$
            by \cref{upair_intro_left,times_tuple_intro,projection_first,img_iff,subseteq}.
        $\img{\pi}{X} = \naturalsPlus$ by \cref{subseteq_antisymmetric}.
        Let $T_A = \subspaceTopology{T_0}{\naturalsPlus}$.
        $\naturalsPlus$ is compact with respect to $T_A$ by \cref{continuous_image_compact}.
        $T_A = T_0$
            by \cref{subspace_topology,discrete_topology,pow_iff,inter_absorb_supseteq_left,subseteq,subseteq_antisymmetric}.
        $\naturalsPlus$ is compact with respect to $T_0$.

        Let $C = \{\cons{n}{\emptyset} \mid n\in\naturalsPlus\}$.
        $C$ is an open covering of $\naturalsPlus$ with respect to $T_0$
            by \cref{emptyset_subseteq,cons_subseteq_intro,powerset_intro,discrete_topology,open_in,cons_left,unions_iff,subseteq,covering,open_covering}.
        Take $B$ such that
            $B\subseteq C$ and
            $B$ is finite and
            $B$ is a covering of $\naturalsPlus$ by \cref{compact}.
        $\naturalsPlus$ is finite
            by \cref{subseteq,pair_is_finite,upair_intro_left,emptyset_subseteq,cons_subseteq_intro,subseteq_finite_is_finite,covering,finite_union_of_finite_family}.
        Contradiction by \cref{naturalsplus_infinite}.
    \end{subproof}
    There exists $Z,S$ such that
        $Z$ is limit point compact with respect to $S$ and
        $Z$ is not compact with respect to $S$.
\end{proof}

% from §28 helper definition 1: strictly increasing on a set
\begin{definition}\label{strictly_increasing_on_set}
    $u$ is strictly increasing on $A$ iff
        $u\in\funs{A}{A}$ and
        for all $m$ such that $m\in A$ we have
            for all $n$ such that $n\in A$ and $m < n$ we have $u(m) < u(n)$.
\end{definition}

% from §28 Item 4: subsequence
\begin{definition}\label{subsequence}
    $t$ is a subsequence of $s$ iff
        there exists $u$ such that
            $u$ is strictly increasing on $\naturalsPlus$ and
            for all $n\in\naturalsPlus$ we have $t(n) = \apply{s}{\apply{u}{n}}$.
\end{definition}

% from §28 Item 5: sequentially compact
\begin{definition}\label{sequentially_compact}
    $X$ is sequentially compact with respect to $T$ iff
        $T$ is a topology on $X$ and
        for all $s$ such that $s$ is a sequence in $X$
            there exist $t,x$ such that
                $t$ is a sequence in $X$ and
                $t$ is a subsequence of $s$ and
                $x\in X$ and
                $t$ converges to $x$ in $X$ with respect to $T$.
\end{definition}

% from §28 helper lemma 2: metric spaces are Hausdorff
\begin{lemma}\label{metric_space_hausdorff}
    Suppose $d$ is a metric on $X$.
    Let $T = \metricTopology{d}{X}$.
    Then $X$ is Hausdorff with respect to $T$.
\end{lemma}
\begin{proof}
    $\metricBasis{d}{X}$ is a basis on $X$ by \cref{metric_basis_is_basis}.
    $T$ is a topology on $X$ by \cref{basis_topology_is_topology,metric_topology}.
    Show that for all $x_1,x_2$ such that $x_1\in X$ and $x_2\in X$ and $x_1\neq x_2$
        there exist $U_1,U_2$ such that
            $U_1$ is a neighborhood of $x_1$ in $X$ with respect to $T$ and
            $U_2$ is a neighborhood of $x_2$ in $X$ with respect to $T$ and
            $U_1$ is disjoint from $U_2$.
    \begin{subproof}
        Fix $x_1$.
        Fix $x_2$ such that $x_1\in X$ and $x_2\in X$ and $x_1\neq x_2$.
        Let $r = d((x_1,x_2))$.
        $r\in\reals$ by \cref{metric_on,times_tuple_intro,funs_type_apply}.
        $0 < r$ by \cref{metric_on,metric_nonnegative,metric_zero_characterization,real_leq_split}.
        Let $m_2 = 1 + 1$.
        $m_2\in\reals$ and $0 < m_2$
            by \cref{real_one_in_naturals,real_naturals_subset_reals,subseteq,real_add_closed,real_add_ident,real_zero_lt_one,real_lt_add,real_lt_trans}.
        $m_2 \neq 0$.
        Let $\epsilon = r / m_2$.
        $\epsilon\in\reals$ and $0 < \epsilon$
            by \cref{real_div,real_mul_inverse,real_mul_closed,real_inv_pos,real_add_ident,real_lt_mul_pos,real_mul_zero}.
        Let $U_1 = \metricBall{d}{X}{x_1}{\epsilon}$.
        Let $U_2 = \metricBall{d}{X}{x_2}{\epsilon}$.
        $U_1$ is a neighborhood of $x_1$ in $X$ with respect to $T$ and
            $U_2$ is a neighborhood of $x_2$ in $X$ with respect to $T$
            by \cref{metric_ball,subseteq,powerset_intro,metric_basis,metric_basis_is_basis,basis_elem_open_in_generated,basis_for_topology,basis_topology_is_topology,metric_topology,open_in,metric_on,metric_zero_characterization,neighborhood}.
        Show that for all $y$ such that $y\in U_1$ we have $y\notin U_2$.
        \begin{subproof}
            Fix $y$ such that $y\in U_1$.
            Suppose not.
            $y\in X$ and $d((x_1,y)) < \epsilon$ and $d((x_2,y)) < \epsilon$ and
                $d((x_1,y))\in\reals$ and $d((x_2,y))\in\reals$
                by \cref{metric_ball,metric_on,times_tuple_intro,funs_type_apply}.
            $r \leq d((x_1,y)) + d((x_2,y))$ by \cref{metric_on,metric_triangle_inequality,metric_symmetric}.
            $d((x_1,y)) + \epsilon\in\reals$ and
                $d((x_1,y)) + d((x_2,y))\in\reals$ and
                $\epsilon + \epsilon\in\reals$ by \cref{real_add_closed}.
            $d((x_1,y)) + d((x_2,y)) < \epsilon + \epsilon$
                by \cref{real_lt_add,real_add_comm,real_lt_trans}.
            $r < \epsilon + \epsilon$ by \cref{real_leq_split,real_lt_trans,real_lt_subst_left}.
            $r < r$
                by \cref{real_div,real_right_distrib,real_mul_ident,real_mul_inverse,real_mul_assoc,real_mul_comm,real_lt_subst_right}.
            Contradiction by \cref{real_lt_irreflexive}.
        \end{subproof}
        $U_1$ is disjoint from $U_2$ by \cref{disjoint}.
    \end{subproof}
    $X$ is Hausdorff with respect to $T$ by \cref{hausdorff}.
\end{proof}


% from §28 helper definition 1: iteration sequence predicate
\begin{definition}\label{iteration_sequence}
    $u$ is an iteration sequence for $f$ on $n$ starting at $a_0$ iff
        $u$ is a function and
        $\dom{u} = \initialSegment{n}$ and
        $u(1) = a_0$ and
        for all $k\in\naturalsPlus$ such that $k + 1 \leq n$ we have $u(k + 1) = \apply{f}{u(k)}$.
\end{definition}

% from §28 helper definition 2: iteration set
\begin{definition}\label{iteration_set}
    $\iterationSet{f}{a_0} = \{ n\in\naturalsPlus \mid \text{there exists $u$ such that
        $u$ is an iteration sequence for $f$ on $n$ starting at $a_0$} \}$.
\end{definition}

% from §28 helper lemma 3: iteration set introduction
\begin{lemma}\label{iteration_set_intro}
    Suppose $n\in\naturalsPlus$.
    Suppose there exists $u$ such that
        $u$ is an iteration sequence for $f$ on $n$ starting at $a_0$.
    Then $n\in\iterationSet{f}{a_0}$.
\end{lemma}
\begin{proof}
    $n\in\iterationSet{f}{a_0}$ by \cref{iteration_set}.
\end{proof}

% from §28 helper lemma 4: iteration set elimination
\begin{lemma}\label{iteration_set_elim}
    Suppose $n\in\iterationSet{f}{a_0}$.
    Then $n\in\naturalsPlus$ and there exists $u$ such that
        $u$ is an iteration sequence for $f$ on $n$ starting at $a_0$.
\end{lemma}
\begin{proof}
    $n\in\naturalsPlus$ and there exists $u$ such that
        $u$ is an iteration sequence for $f$ on $n$ starting at $a_0$
        by \cref{iteration_set}.
\end{proof}

% from §28 helper lemma 5: iteration sequence introduction
\begin{lemma}\label{iteration_sequence_intro}
    Suppose $u$ is a function and $\dom{u} = \initialSegment{n}$ and $u(1) = a_0$.
    Suppose for all $k\in\naturalsPlus$ such that $k + 1 \leq n$ we have $u(k + 1) = \apply{f}{u(k)}$.
    Then $u$ is an iteration sequence for $f$ on $n$ starting at $a_0$.
\end{lemma}
\begin{proof}
    $u$ is an iteration sequence for $f$ on $n$ starting at $a_0$ by \cref{iteration_sequence}.
\end{proof}

% from §28 helper lemma 5a: naturalsPlus successor is larger
\begin{lemma}\label{real_naturals_lt_succ_local}
    Suppose $n\in\naturalsPlus$.
    Then $n < n + 1$.
\end{lemma}
\begin{proof}
    $n < n + 1$
        by \cref{real_naturals_subset_reals,real_one_in_naturals,subseteq,real_add_ident,real_zero_lt_one,real_lt_add,real_add_comm,real_lt_subst_left}.
\end{proof}

% from §28 helper lemma 5b: naturalsPlus successor is not smaller
\begin{lemma}\label{real_naturals_leq_succ_local}
    Suppose $n\in\naturalsPlus$.
    Then $n \leq n + 1$.
\end{lemma}
\begin{proof}
    $n \leq n + 1$ by \cref{real_naturals_lt_succ_local}.
\end{proof}

% from §28 helper lemma 5c: predecessor bound elimination
\begin{lemma}\label{real_naturals_succ_bound_elim_local}
    Suppose $k\in\naturalsPlus$ and $n\in\naturalsPlus$.
    Suppose $k + 1 \leq n$.
    Then $k \leq n$.
\end{lemma}
\begin{proof}
    $k \leq n$
        by \cref{real_naturals_leq_succ_local,real_naturals_subset_reals,real_naturals_closed_succ,subseteq,real_leq_trans}.
\end{proof}

% from §28 helper lemma 6: iteration sequence uniqueness
\begin{lemma}\label{iteration_sequence_unique}
    Suppose $n\in\naturalsPlus$.
    Suppose $u$ and $v$ are iteration sequences for $f$ on $n$ starting at $a_0$.
    Then $u = v$.
\end{lemma}
\begin{proof}
    $u$ is a function and $v$ is a function and
        $\dom{u} = \initialSegment{n}$ and $\dom{v} = \initialSegment{n}$ by \cref{iteration_sequence}.
    Let $A = \{k\in\naturalsPlus \mid \text{if $k\leq n$ then $u(k) = v(k)$}\}$.
    $1\in A$ by \cref{real_one_in_naturals,iteration_sequence}.
    Show that for all $k\in A$ we have if $k + 1\leq n$ then $u(k + 1) = v(k + 1)$.
    \begin{subproof}
        Follows by \cref{real_naturals_succ_bound_elim_local,iteration_sequence}.
    \end{subproof}
    For all $k\in A$ we have $k + 1\in A$.
    For all $k\in\initialSegment{n}$ we have $u(k) = v(k)$
        by \cref{subseteq,real_naturals_induction,initial_segment_naturalsplus}.
    Show that for all $w$ such that $w\in u$ we have $w\in v$.
    \begin{subproof}
        Follows by \cref{function_member_elim,subseteq,function_apply_elim}.
    \end{subproof}
    Show that for all $w$ such that $w\in v$ we have $w\in u$.
    \begin{subproof}
        Follows by \cref{function_member_elim,subseteq,function_apply_elim}.
    \end{subproof}
    $u = v$ by \cref{subseteq,subseteq_antisymmetric}.
\end{proof}
% from §28 helper lemma 7: infinite subset of naturalsPlus has increasing sequence
\begin{lemma}\label{infinite_naturalsplus_has_strictly_increasing_sequence}
    Suppose $A\subseteq\naturalsPlus$.
    Suppose $A$ is infinite.
    Then there exists $u$ such that
        $u$ is strictly increasing on $\naturalsPlus$ and
        for all $n\in\naturalsPlus$ we have $u(n)\in A$.
\end{lemma}
\begin{proof}
    $A\neq\emptyset$ by \cref{emptyset_is_finite}.
    Take $a_0\in A$ such that for all $a\in A$ we have $a_0 \leq a$
        by \cref{real_naturals_least_element}.
    Show that for all $a\in A$ there exists $b\in A$ such that $a < b$.
    \begin{subproof}
        Fix $a$ such that $a\in A$.
        $a\in\naturalsPlus$ by \cref{subseteq}.
        Suppose not.
        For all $b\in A$ we have it is wrong that $a < b$.
        Show that for all $b$ such that $b\in A$ we have $b\in\initialSegment{a}$.
        \begin{subproof}
            Fix $b$ such that $b\in A$.
            $b\in\naturalsPlus$ by \cref{subseteq}.
            $a\in\reals$ and $b\in\reals$ by \cref{real_naturals_subset_reals,subseteq}.
            Show that $b\in\initialSegment{a}$.
            \begin{subproof}
                $b < a$ or $b = a$ or $a < b$ by \cref{real_lt_trichotomy}.
                \begin{byCase}
                \caseOf{$b < a$.}
                    $b \leq a$ by \cref{real_lt_imp_leq}.
                    $b\in\initialSegment{a}$ by \cref{initial_segment_naturalsplus}.
                \caseOf{$b = a$.}
                    $b \leq a$.
                    $b\in\initialSegment{a}$ by \cref{initial_segment_naturalsplus}.
                \caseOf{$a < b$.}
                    Contradiction.
                \end{byCase}
            \end{subproof}
        \end{subproof}
        $A\subseteq\initialSegment{a}$ by \cref{subseteq}.
        $A$ is finite by \cref{initial_segment_finite,subseteq_finite_is_finite}.
        Contradiction.
    \end{subproof}
    Let $R = \{w\in A\times A \mid \fst{w} < \snd{w}\}$.
    Show that for all $x\in A$ there exists $y$ such that $x\mathrel{R} y$.
    \begin{subproof}
        Fix $x\in A$.
        Take $y_0\in A$ such that $x < y_0$.
        $(x,y_0)\in A\times A$ by \cref{times_tuple_intro}.
        $(x,y_0)\in R$ by \cref{fst_eq,snd_eq}.
        There exists $y$ such that $x\mathrel{R} y$ by \cref{relation}.
    \end{subproof}
    Take $f$ such that $f$ is a function on $A$ and for all $x\in A$ we have
        $x\mathrel{R}\apply{f}{x}$ by \cref{times_elem_is_tuple,relation,choice_function}.
    Show that for all $x\in A$ we have $\apply{f}{x}\in A$ and $x < \apply{f}{x}$.
    \begin{subproof}
        Fix $x\in A$.
        $x\mathrel{R}\apply{f}{x}$.
        $(x,\apply{f}{x})\in R$ by \cref{relation}.
        $(x,\apply{f}{x})\in A\times A$ and
            $\fst{(x,\apply{f}{x})} < \snd{(x,\apply{f}{x})}$.
        $x\in A$ and $\apply{f}{x}\in A$ by \cref{times_tuple_elim}.
        $x < \apply{f}{x}$ by \cref{fst_eq,snd_eq}.
    \end{subproof}
    Let $S = \iterationSet{f}{a_0}$.
    Show that for all $n\in\naturalsPlus$ we have $n\in S$.
    \begin{subproof}
        $S\subseteq\naturalsPlus$ by \cref{iteration_set_elim,subseteq}.
        $1\in\naturalsPlus$ by \cref{real_one_in_naturals}.
        Show that $1\in S$.
        \begin{subproof}
            Let $u = \{(1,a_0)\}$.
            For all $x,y,y'$ such that $(x,y)\in u$ and $(x,y')\in u$ we have $y = y'$
                by \cref{singleton_iff,pair_eq_iff}.
            $u$ is a function by \cref{singleton_iff,relation,rightunique}.
            Show that for all $x$ such that $x\in\dom{u}$ we have $x\in\initialSegment{1}$.
            \begin{subproof}
                Follows by \cref{dom,singleton_iff,pair_eq_iff,real_one_in_naturals,initial_segment_naturalsplus}.
            \end{subproof}
            Show that for all $x$ such that $x\in\initialSegment{1}$ we have $x\in\dom{u}$.
            \begin{subproof}
                Fix $x$ such that $x\in\initialSegment{1}$.
                $x\in\naturalsPlus$ and $x \leq 1$ by \cref{initial_segment_naturalsplus}.
                $x\in\reals$ and $1\in\reals$
                    by \cref{real_naturals_subset_reals,subseteq,real_one_in_naturals}.
                $1 \leq x$ by \cref{real_one_leq_naturalsplus}.
                $x = 1$ by \cref{real_leq_antisym}.
                $(x,a_0)\in u$ by \cref{singleton_iff,pair_eq_iff}.
                $x\in\dom{u}$ by \cref{dom}.
            \end{subproof}
            $\dom{u} = \initialSegment{1}$ by \cref{subseteq,subseteq_antisymmetric}.
            $u(1) = a_0$ by \cref{singleton_intro,function_apply_intro}.
            Show that for all $k\in\naturalsPlus$ such that $k + 1 \leq 1$ we have $u(k + 1) = \apply{f}{u(k)}$.
            \begin{subproof}
                Fix $k$ such that $k\in\naturalsPlus$.
                Assume $k + 1 \leq 1$.
                Suppose not.
                $k < k + 1$ by \cref{real_naturals_lt_succ_local}.
                $1\in\reals$ and $k\in\reals$ and $k + 1\in\reals$
                    by \cref{real_naturals_subset_reals,real_one_in_naturals,real_naturals_closed_succ,subseteq}.
                $1 < k$ or $1 = k$ by \cref{real_one_leq_naturalsplus}.
                $1 < k + 1$ by \cref{real_lt_subst_left,real_lt_trans}.
                $k + 1 < 1$ or $k + 1 = 1$ by \cref{real_leq_split}.
                $1 < 1$ by \cref{real_lt_trans,real_lt_subst_right}.
                Contradiction by \cref{real_naturals_subset_reals,real_one_in_naturals,subseteq,real_lt_irreflexive}.
            \end{subproof}
            $1\in S$ by \cref{iteration_sequence,iteration_set_intro}.
        \end{subproof}
        Show that for all $n\in S$ we have $n + 1\in S$.
        \begin{subproof}
            Fix $n$ such that $n\in S$.
            $n\in\naturalsPlus$ by \cref{iteration_set_elim}.
            Take $u$ such that
                $u$ is an iteration sequence for $f$ on $n$ starting at $a_0$.
            For all $k\in\naturalsPlus$ such that $k + 1 \leq n$ we have $u(k + 1) = \apply{f}{u(k)}$
                by \cref{iteration_sequence}.
            $n + 1\in\naturalsPlus$ by \cref{real_naturals_closed_succ}.
            Let $v = u\union\{(n + 1,\apply{f}{u(n)})\}$.
            $v$ is a relation by \cref{singleton_iff,relation,iteration_sequence,union_relations_is_relation}.
            Show that for all $x,y,y'$ such that $(x,y)\in v$ and $(x,y')\in v$ we have $y = y'$.
            \begin{subproof}
                Fix $x$.
                Fix $y$.
                Fix $y'$.
                Assume $(x,y)\in v$ and $(x,y')\in v$.
                \begin{byCase}
                \caseOf{$x = n + 1$.}
                    $(n + 1,y')\in v$.
                    Show that $(n + 1,y)\in\{(n + 1,\apply{f}{u(n)})\}$.
                    \begin{subproof}
                        Suppose not.
                        $(n + 1,y)\in u$ by \cref{union_iff}.
                        $n + 1\in\dom{u}$ by \cref{dom}.
                        $\dom{u} = \initialSegment{n}$ by \cref{iteration_sequence}.
                        $n + 1\in\initialSegment{n}$.
                        $n + 1 \leq n$ by \cref{initial_segment_naturalsplus}.
                        $n < n + 1$ by \cref{real_naturals_lt_succ_local}.
                        $n\in\reals$ and $n + 1\in\reals$
                            by \cref{real_naturals_subset_reals,real_naturals_closed_succ,subseteq}.
                        $n + 1 < n$ or $n + 1 = n$ by \cref{real_leq_split}.
                        $n < n$ by \cref{real_lt_trans,real_lt_subst_right}.
                        Contradiction by \cref{real_lt_irreflexive}.
                    \end{subproof}
                    $(n + 1,y) = (n + 1,\apply{f}{u(n)})$ by \cref{singleton_iff}.
                    $y = \apply{f}{u(n)}$ by \cref{pair_eq_iff}.
                    Show that $(n + 1,y')\in\{(n + 1,\apply{f}{u(n)})\}$.
                    \begin{subproof}
                        Suppose not.
                        $(n + 1,y')\in u$ by \cref{union_iff}.
                        $n + 1\in\dom{u}$ by \cref{dom}.
                        $\dom{u} = \initialSegment{n}$ by \cref{iteration_sequence}.
                        $n + 1\in\initialSegment{n}$.
                        $n + 1 \leq n$ by \cref{initial_segment_naturalsplus}.
                        $n < n + 1$ by \cref{real_naturals_lt_succ_local}.
                        $n\in\reals$ and $n + 1\in\reals$
                            by \cref{real_naturals_subset_reals,real_naturals_closed_succ,subseteq}.
                        $n + 1 < n$ or $n + 1 = n$ by \cref{real_leq_split}.
                        $n < n$ by \cref{real_lt_trans,real_lt_subst_right}.
                        Contradiction by \cref{real_lt_irreflexive}.
                    \end{subproof}
                    $(n + 1,y') = (n + 1,\apply{f}{u(n)})$ by \cref{singleton_iff}.
                    $y' = \apply{f}{u(n)}$ by \cref{pair_eq_iff}.
                    $y = y'$.
                \caseOf{$x \neq n + 1$.}
                    $(x,y)\in u$ by \cref{union_iff,singleton_iff,pair_eq_iff}.
                    $(x,y')\in u$ by \cref{union_iff,singleton_iff,pair_eq_iff}.
                    $y = y'$ by \cref{iteration_sequence,function_rightunique}.
                \end{byCase}
            \end{subproof}
            $v$ is a function by \cref{rightunique}.
            Show that for all $x$ such that $x\in\dom{v}$ we have $x\in\initialSegment{n + 1}$.
            \begin{subproof}
                Fix $x$ such that $x\in\dom{v}$.
                Take $y$ such that $(x,y)\in v$ by \cref{dom}.
                $(x,y)\in u$ or $(x,y)\in\{(n + 1,\apply{f}{u(n)})\}$ by \cref{union_iff}.
                \begin{byCase}
                \caseOf{$(x,y)\in u$.}
                    $x\in\initialSegment{n}$ by \cref{dom,iteration_sequence}.
                    $x\in\initialSegment{n + 1}$
                        by \cref{initial_segment_naturalsplus,real_naturals_leq_succ_local,real_leq_trans,real_naturals_subset_reals,real_naturals_closed_succ,subseteq}.
                \caseOf{$(x,y)\in\{(n + 1,\apply{f}{u(n)})\}$.}
                    $(x,y) = (n + 1,\apply{f}{u(n)})$ by \cref{singleton_iff}.
                    $x = n + 1$ by \cref{pair_eq_iff}.
                    $x\in\initialSegment{n + 1}$ by \cref{initial_segment_naturalsplus}.
                \end{byCase}
            \end{subproof}
            Show that for all $x$ such that $x\in\initialSegment{n + 1}$ we have $x\in\dom{v}$.
            \begin{subproof}
                Fix $x$ such that $x\in\initialSegment{n + 1}$.
                $x\in\naturalsPlus$ and $x \leq n + 1$ by \cref{initial_segment_naturalsplus}.
                \begin{byCase}
                \caseOf{$x = n + 1$.}
                    $x\in\dom{v}$ by \cref{cons_left,union_intro_right,dom}.
                \caseOf{$x \neq n + 1$.}
                    $x \leq n$ by \cref{real_naturals_leq_succ_elim}.
                    $x\in\dom{u}$ by \cref{initial_segment_naturalsplus,iteration_sequence}.
                    $x\in\dom{v}$ by \cref{dom,union_iff}.
                \end{byCase}
            \end{subproof}
            $\dom{v} = \initialSegment{n + 1}$ by \cref{subseteq,subseteq_antisymmetric}.
            $v(1) = a_0$
                by \cref{real_one_in_naturals,real_one_leq_naturalsplus,initial_segment_naturalsplus,iteration_sequence,function_apply_elim,union_intro_left,function_apply_intro}.
            Show that for all $k\in\naturalsPlus$ such that $k + 1 \leq n + 1$ we have $v(k + 1) = \apply{f}{v(k)}$.
            \begin{subproof}
                Fix $k$ such that $k\in\naturalsPlus$.
                Assume $k + 1 \leq n + 1$.
                \begin{byCase}
                \caseOf{$k + 1 = n + 1$.}
                    $k = n$.
                    $v(k + 1) = \apply{f}{u(n)}$
                        by \cref{cons_left,union_intro_right,function_apply_intro}.
                    $v(k + 1) = \apply{f}{u(k)}$.
                    $k \leq n$.
                    $v(k) = u(k)$
                        by \cref{initial_segment_naturalsplus,iteration_sequence,function_apply_elim,union_intro_left,function_apply_intro}.
                    $v(k + 1) = \apply{f}{v(k)}$.
                \caseOf{$k + 1 \neq n + 1$.}
                    $k + 1\in\naturalsPlus$ by \cref{real_naturals_closed_succ}.
                    $k + 1 \leq n$ by \cref{real_naturals_leq_succ_elim}.
                    $k + 1\in\initialSegment{n}$ by \cref{initial_segment_naturalsplus}.
                    $(k + 1,u(k + 1))\in v$
                        by \cref{iteration_sequence,function_apply_elim,union_intro_left}.
                    $k + 1\in\dom{v}$ by \cref{dom}.
                    $(k + 1,v(k + 1))\in v$ by \cref{function_apply_elim}.
                    $u(k + 1) = v(k + 1)$ by \cref{function_rightunique}.
                    $v(k + 1) = u(k + 1)$.
                    $k\in\naturalsPlus$ and $k + 1 \leq n$.
                    $u(k + 1) = \apply{f}{u(k)}$.
                    $v(k + 1) = \apply{f}{u(k)}$.
                    $k \leq n$ by \cref{real_naturals_succ_bound_elim_local}.
                    $(k,u(k))\in v$
                        by \cref{initial_segment_naturalsplus,iteration_sequence,function_apply_elim,union_intro_left}.
                    $k\in\dom{v}$ by \cref{dom}.
                    $(k,v(k))\in v$ by \cref{function_apply_elim}.
                    $u(k) = v(k)$ by \cref{function_rightunique}.
                    $v(k + 1) = \apply{f}{v(k)}$.
                \end{byCase}
            \end{subproof}
            $v$ is an iteration sequence for $f$ on $n + 1$ starting at $a_0$ by \cref{iteration_sequence_intro}.
            $n + 1\in S$ by \cref{iteration_set_intro}.
        \end{subproof}
        For all $n\in\naturalsPlus$ we have $n\in S$ by \cref{real_naturals_induction}.
    \end{subproof}
    Let $F = \{u \mid \text{there exists $n\in\naturalsPlus$ such that
        $u$ is an iteration sequence for $f$ on $n$ starting at $a_0$} \}$.
    For all $u\in F$ we have $u$ is a function by \cref{iteration_sequence}.
    Show that for all $u,v$ such that $u\in F$ and $v\in F$ we have $u\subseteq v$ or $v\subseteq u$.
    \begin{subproof}
        Fix $u$.
        Fix $v$.
        Assume $u\in F$ and $v\in F$.
        Take $m\in\naturalsPlus$ such that
            $u$ is an iteration sequence for $f$ on $m$ starting at $a_0$.
        Take $n\in\naturalsPlus$ such that
            $v$ is an iteration sequence for $f$ on $n$ starting at $a_0$.
            $u$ is a function and $\dom{u} = \initialSegment{m}$ and $u(1) = a_0$
                by \cref{iteration_sequence}.
            For all $k\in\naturalsPlus$ such that $k + 1 \leq m$ we have $u(k + 1) = \apply{f}{u(k)}$
                by \cref{iteration_sequence}.
            $v$ is a function and $\dom{v} = \initialSegment{n}$ and $v(1) = a_0$
                by \cref{iteration_sequence}.
            For all $k\in\naturalsPlus$ such that $k + 1 \leq n$ we have $v(k + 1) = \apply{f}{v(k)}$
                by \cref{iteration_sequence}.
        $m = n$ or $m \neq n$.
        \begin{byCase}
        \caseOf{$m = n$.}
            $u\subseteq v$ or $v\subseteq u$ by \cref{iteration_sequence_unique,subseteq_refl}.
        \caseOf{$m \neq n$.}
            $m\in\reals$ and $n\in\reals$ by \cref{real_naturals_subset_reals,subseteq}.
            $m < n$ or $n < m$ by \cref{real_lt_trichotomy}.
            \begin{byCase}
        \caseOf{$m < n$.}
            Let $B = \{k\in\naturalsPlus \mid \text{if $k\leq m$ then $u(k) = v(k)$}\}$.
            $1\in\naturalsPlus$ by \cref{real_one_in_naturals}.
            Show that if $1\leq m$ then $u(1) = v(1)$.
            \begin{subproof}
                Follows by \cref{iteration_sequence}.
            \end{subproof}
            $1\in B$.
            Show that for all $k\in B$ we have if $k + 1\leq m$ then $u(k + 1) = v(k + 1)$.
            \begin{subproof}
                Follows by \cref{real_naturals_succ_bound_elim_local,real_leq_trans,real_lt_imp_leq,real_add_closed,real_naturals_subset_reals,subseteq,iteration_sequence}.
            \end{subproof}
            Show that for all $k\in B$ we have $k + 1\in B$.
            \begin{subproof}
                Fix $k\in B$.
                $k\in\naturalsPlus$ by \cref{subseteq}.
                $k + 1\in\naturalsPlus$ by \cref{real_naturals_closed_succ}.
                Show that if $k + 1\leq m$ then $u(k + 1) = v(k + 1)$.
                \begin{subproof}
                    Trivial.
                \end{subproof}
                $k + 1\in B$.
            \end{subproof}
            For all $k\in\initialSegment{m}$ we have $u(k) = v(k)$
                by \cref{subseteq,real_naturals_induction,initial_segment_naturalsplus}.
            For all $w$ such that $w\in u$ we have $w\in v$
                by \cref{iteration_sequence,function_member_elim,initial_segment_naturalsplus,real_naturals_subset_reals,subseteq,real_leq_trans,real_lt_imp_leq,function_apply_elim}.
            $u\subseteq v$ by \cref{subseteq}.
        \caseOf{$n < m$.}
            Let $C = \{k\in\naturalsPlus \mid \text{if $k\leq n$ then $v(k) = u(k)$}\}$.
            $1\in\naturalsPlus$ by \cref{real_one_in_naturals}.
            Show that if $1\leq n$ then $v(1) = u(1)$.
            \begin{subproof}
                Follows by \cref{iteration_sequence}.
            \end{subproof}
            $1\in C$.
            Show that for all $k\in C$ we have if $k + 1\leq n$ then $v(k + 1) = u(k + 1)$.
            \begin{subproof}
                Follows by \cref{real_naturals_succ_bound_elim_local,real_leq_trans,real_add_closed,real_naturals_subset_reals,subseteq,iteration_sequence}.
            \end{subproof}
            Show that for all $k\in C$ we have $k + 1\in C$.
            \begin{subproof}
                Fix $k\in C$.
                $k\in\naturalsPlus$ by \cref{subseteq}.
                $k + 1\in\naturalsPlus$ by \cref{real_naturals_closed_succ}.
                Show that if $k + 1\leq n$ then $v(k + 1) = u(k + 1)$.
                \begin{subproof}
                    Trivial.
                \end{subproof}
                $k + 1\in C$.
            \end{subproof}
            For all $k\in\initialSegment{n}$ we have $v(k) = u(k)$
                by \cref{subseteq,real_naturals_induction,initial_segment_naturalsplus}.
            For all $w$ such that $w\in v$ we have $w\in u$
                by \cref{iteration_sequence,function_member_elim,initial_segment_naturalsplus,real_naturals_subset_reals,subseteq,real_leq_trans,real_lt_imp_leq,function_apply_elim}.
            $v\subseteq u$ by \cref{subseteq}.
        \end{byCase}
        \end{byCase}
    \end{subproof}
    Let $u = \unions{F}$.
    $F$ is a family of functions.
    $u$ is a function by \cref{unions_of_compatible_family_of_function_is_function}.
    Show that for all $x$ such that $x\in\dom{u}$ we have $x\in\naturalsPlus$.
    \begin{subproof}
        Follows by \cref{dom,unions_iff,iteration_sequence,initial_segment_naturalsplus}.
    \end{subproof}
    Show that for all $n\in\naturalsPlus$ we have $n\in\dom{u}$.
    \begin{subproof}
        Follows by \cref{subseteq,iteration_set_elim,iteration_sequence,initial_segment_naturalsplus,dom,unions_iff}.
    \end{subproof}
    $\dom{u} = \naturalsPlus$ by \cref{subseteq,subseteq_antisymmetric}.
    Show that for all $n\in\naturalsPlus$ we have $u(n)\in A$.
    \begin{subproof}
        Fix $n\in\naturalsPlus$.
        Take $v\in F$ such that $(n,u(n))\in v$ by \cref{function_apply_elim,unions_iff}.
        Take $n_0\in\naturalsPlus$ such that
            $v$ is an iteration sequence for $f$ on $n_0$ starting at $a_0$.
        $n \leq n_0$ by \cref{dom,iteration_sequence,initial_segment_naturalsplus}.
        $v(1) = a_0$ and for all $k\in\naturalsPlus$ such that $k + 1 \leq n_0$ we have $v(k + 1) = \apply{f}{v(k)}$
            by \cref{iteration_sequence}.
        Let $B = \{k\in\naturalsPlus \mid \text{if $k\leq n$ then $v(k)\in A$}\}$.
        $1\in\naturalsPlus$ by \cref{real_one_in_naturals}.
        Show that if $1\leq n$ then $v(1)\in A$.
        \begin{subproof}
            Follows by \cref{iteration_sequence}.
        \end{subproof}
        $1\in B$.
        Show that for all $k\in B$ we have if $k + 1\leq n$ then $v(k + 1)\in A$.
        \begin{subproof}
            Follows by \cref{real_naturals_succ_bound_elim_local,funs_type_apply,real_leq_trans,real_naturals_closed_succ,real_naturals_subset_reals,subseteq,iteration_sequence}.
        \end{subproof}
        Show that for all $k\in B$ we have $k + 1\in B$.
        \begin{subproof}
            Fix $k\in B$.
            $k\in\naturalsPlus$ by \cref{subseteq}.
            $k + 1\in\naturalsPlus$ by \cref{real_naturals_closed_succ}.
            Show that if $k + 1\leq n$ then $v(k + 1)\in A$.
            \begin{subproof}
                Trivial.
            \end{subproof}
            $k + 1\in B$.
        \end{subproof}
        $v(n)\in A$ by \cref{subseteq,real_naturals_induction,initial_segment_naturalsplus}.
        $u(n)\in A$ by \cref{function_apply_intro}.
    \end{subproof}
    Show that for all $n\in\naturalsPlus$ we have $u(n) < u(n + 1)$.
        \begin{subproof}
            Fix $n\in\naturalsPlus$.
            $(n + 1,u(n + 1))\in\unions{F}$ by \cref{real_naturals_closed_succ,function_apply_elim}.
            Take $v\in F$ such that $(n + 1,u(n + 1))\in v$ by \cref{unions_iff}.
            $n + 1\in\dom{v}$ by \cref{dom}.
            $v(n + 1) = u(n + 1)$ by \cref{function_apply_intro}.
            Take $n_0\in\naturalsPlus$ such that
                $v$ is an iteration sequence for $f$ on $n_0$ starting at $a_0$.
            $n + 1 \leq n_0$ by \cref{iteration_sequence,initial_segment_naturalsplus}.
            $v(n + 1) = \apply{f}{v(n)}$ by \cref{iteration_sequence}.
            $u(n + 1) = \apply{f}{v(n)}$.
            $n \leq n_0$ by \cref{real_naturals_succ_bound_elim_local}.
            $(n,v(n))\in u$
                by \cref{initial_segment_naturalsplus,iteration_sequence,function_apply_elim,unions_iff}.
            $u(n) = v(n)$ by \cref{function_apply_intro}.
            $u(n + 1) = \apply{f}{u(n)}$.
            $u(n) < u(n + 1)$ by \cref{function_apply_intro}.
        \end{subproof}
        Show that for all $m,n\in\naturalsPlus$ such that $m < n$ we have $u(m) < u(n)$.
        \begin{subproof}
            Let $T = \{n\in\naturalsPlus \mid \text{for all $m\in\naturalsPlus$ such that $m < n$ we have $u(m) < u(n)$}\}$.
            Show that for all $n\in\naturalsPlus$ we have $n\in T$.
            \begin{subproof}
                $1\in\naturalsPlus$ by \cref{real_one_in_naturals}.
                For all $m\in\naturalsPlus$ such that $m < 1$ we have $u(m) < u(1)$
                    by \cref{real_naturals_subset_reals,real_one_in_naturals,subseteq,real_one_leq_naturalsplus,real_lt_subst_right,real_lt_trans,real_lt_irreflexive}.
                $1\in T$.
                Show that for all $n\in T$ we have $n + 1\in T$.
                \begin{subproof}
                    Fix $n$ such that $n\in T$.
                    Show that for all $m\in\naturalsPlus$ such that $m < n + 1$ we have $u(m) < u(n + 1)$.
                    \begin{subproof}
                        Fix $m$.
                        Assume $m\in\naturalsPlus$ and $m < n + 1$.
                        Suppose not.
                        $u(m)\in A$ and $u(n)\in A$ and $u(n + 1)\in A$.
                        $u(m)\in\reals$ and $u(n)\in\reals$ and $u(n + 1)\in\reals$
                            by \cref{subseteq,real_naturals_subset_reals}.
                        $m < n$ or $m = n$ or $m = n + 1$ by \cref{real_succ_discrete}.
                        \begin{byCase}
                        \caseOf{$m = n$.}
                            $u(m) < u(n + 1)$ by \cref{real_lt_subst_left}.
                            Contradiction.
                        \caseOf{$m < n$.}
                            $u(m) < u(n + 1)$ by \cref{subseteq,real_lt_trans}.
                            Contradiction.
                        \caseOf{$m = n + 1$.}
                            Contradiction by \cref{real_lt_subst_left,real_naturals_subset_reals,subseteq,real_lt_irreflexive}.
                        \end{byCase}
                    \end{subproof}
                    $n + 1\in T$.
                \end{subproof}
                For all $m,n\in\naturalsPlus$ such that $m < n$ we have $u(m) < u(n)$
                    by \cref{subseteq,real_naturals_induction}.
            \end{subproof}
    \end{subproof}
    For all $n\in\dom{u}$ we have $u(n)\in\naturalsPlus$.
    $u\in\funs{\naturalsPlus}{\naturalsPlus}$ by \cref{funs_intro}.
    Show that $u$ is strictly increasing on $\naturalsPlus$.
    \begin{subproof}
        Follows by \cref{strictly_increasing_on_set}.
    \end{subproof}
    Let $v = u$.
    $v$ is strictly increasing on $\naturalsPlus$ and
        for all $n\in\naturalsPlus$ we have $v(n)\in A$.
\end{proof}


% from §28 helper lemma 8: constant subsequence in a finite range
\begin{lemma}\label{finite_sequence_constant_subsequence}
    Suppose $A\subseteq X$.
    Suppose $A$ is finite.
    Suppose $s$ is a sequence in $A$.
    Then there exist $t,a$ such that
        $a\in A$ and
        $t$ is a sequence in $A$ and
        $t$ is a subsequence of $s$ and
        for all $n\in\naturalsPlus$ we have $t(n) = a$.
\end{lemma}
\begin{proof}
    $s\in\funs{\naturalsPlus}{A}$ by \cref{sequence_in}.
    $\dom{s} = \naturalsPlus$ and $s$ is right-unique and $s$ is a function by \cref{funs,funs_is_function}.
    Let $g = \{ (a,\preimg{s}{\{a\}}) \mid a\in A \}$.
    Let $F = \img{g}{A}$.
    For all $w$ such that $w\in g$ there exist $a,B$ such that $w = (a,B)$.
    For all $a,B_1,B_2$ such that $(a,B_1)\in g$ and $(a,B_2)\in g$ we have $B_1 = B_2$.
    Show that for all $a$ such that $a\in\dom{g}$ we have $a\in A$.
    \begin{subproof}
        Follows by \cref{dom_iff,pair_eq_iff}.
    \end{subproof}
    Show that for all $a$ such that $a\in A$ we have $a\in\dom{g}$.
    \begin{subproof}
        Follows by \cref{dom_intro}.
    \end{subproof}
    $\dom{g} = A$ by \cref{subseteq,subseteq_antisymmetric}.
    $g$ is a function on $A$ by \cref{relation,rightunique}.
    $F$ is finite by \cref{finite_image_function}.
    Show that for all $n$ such that $n\in\naturalsPlus$ we have $n\in\unions{F}$.
    \begin{subproof}
        Follows by \cref{funs_type_apply,img_iff,singleton_intro,preimg_iff,function_apply_iff,unions_iff}.
    \end{subproof}
    Show that for all $n$ such that $n\in\unions{F}$ we have $n\in\naturalsPlus$.
    \begin{subproof}
        Follows by \cref{unions_iff,img_iff,function_apply_iff,preimg_subseteq_dom,subseteq}.
    \end{subproof}
    $\unions{F} = \naturalsPlus$ by \cref{subseteq,subseteq_antisymmetric}.
    Show that there exists $a\in A$ such that $\preimg{s}{\{a\}}$ is infinite.
    \begin{subproof}
        Suppose not.
        For all $a\in A$ we have $\preimg{s}{\{a\}}$ is finite.
        For all $B\in F$ we have $B$ is finite by \cref{img_iff,function_apply_iff}.
        $\unions{F}$ is finite by \cref{finite_union_of_finite_family}.
        $\naturalsPlus$ is finite.
        Contradiction by \cref{naturalsplus_infinite}.
    \end{subproof}
    Take $a\in A$ such that $\preimg{s}{\{a\}}$ is infinite.
    Let $I = \preimg{s}{\{a\}}$.
    $I\subseteq\naturalsPlus$ by \cref{preimg_subseteq_dom,subseteq}.
    Take $u$ such that
        $u$ is strictly increasing on $\naturalsPlus$ and
        for all $n\in\naturalsPlus$ we have $u(n)\in I$
        by \cref{infinite_naturalsplus_has_strictly_increasing_sequence}.
    Let $t = s\circ u$.
    $t$ is a function by \cref{strictly_increasing_on_set,funs,function_circ}.
    Show that for all $n$ such that $n\in\dom{t}$ we have $n\in\naturalsPlus$.
    \begin{subproof}
        Follows by \cref{dom,circ_elem_elim,strictly_increasing_on_set,funs}.
    \end{subproof}
    Show that for all $n$ such that $n\in\naturalsPlus$ we have $n\in\dom{t}$.
    \begin{subproof}
        Follows by \cref{strictly_increasing_on_set,funs,subseteq,funs_is_function,function_apply_elim,circ_elem_intro,dom}.
    \end{subproof}
    $\dom{t} = \naturalsPlus$ by \cref{subseteq,subseteq_antisymmetric}.
    For all $n\in\dom{t}$ we have $t(n)\in A$
        by \cref{dom,strictly_increasing_on_set,funs_is_function,function_apply_iff,function_apply_elim,circ_elem_elim,subseteq,funs_type_apply}.
    $t\in\funs{\naturalsPlus}{A}$ by \cref{funs_intro}.
    Show that $t$ is a sequence in $A$.
    \begin{subproof}
        Follows by \cref{sequence_in}.
    \end{subproof}
    For all $n\in\naturalsPlus$ we have $t(n) = \apply{s}{\apply{u}{n}}$
        by \cref{sequence_in,funs,funs_is_function,function_apply_elim,circ_elem_elim,strictly_increasing_on_set,function_apply_iff}.
    Show that $t$ is a subsequence of $s$.
    \begin{subproof}
        Follows by \cref{subsequence}.
    \end{subproof}
    Show that for all $n\in\naturalsPlus$ we have $t(n) = a$.
    \begin{subproof}
        Follows by \cref{strictly_increasing_on_set,funs,preimg_iff,singleton_iff,function_apply_iff}.
    \end{subproof}
    Let $r = t$.
    Let $b = a$.
    $b\in A$ and
        $r$ is a sequence in $A$ and
        $r$ is a subsequence of $s$ and
        for all $n\in\naturalsPlus$ we have $r(n) = b$.
\end{proof}

% from §28 helper lemma: reciprocal reverses order on naturalsPlus
\begin{lemma}\label{naturalsplus_reciprocal_antitone_alt}
    Suppose $m,n\in\naturalsPlus$.
    Suppose $m \leq n$.
    Then $1 / n \leq 1 / m$.
\end{lemma}
\begin{proof}
    $m < n$ or $m = n$ by \cref{real_leq_split}.
    $m\in\reals$ and $n\in\reals$ and $0\in\reals$ and $1\in\reals$
        by \cref{real_naturals_subset_reals,subseteq,real_add_ident,real_mul_ident}.

    $0 < m$ by \cref{real_add_ident,real_zero_lt_one,real_one_leq_naturalsplus,real_lt_subst_right,real_lt_trans}.
    $0 < n$ by \cref{real_add_ident,real_zero_lt_one,real_one_leq_naturalsplus,real_lt_subst_right,real_lt_trans}.
    Suppose not.
    $1 / m < 1 / n$.
    $0 < m \rmul n$ by \cref{real_lt_mul_pos,real_mul_zero,real_lt_subst_left}.
    $m\neq 0$ and $n\neq 0$ by \cref{real_pos_nonzero}.
    $1 / m\in\reals$ and $1 / n\in\reals$ and $m \rmul n\in\reals$
        by \cref{real_mul_inverse,real_div,real_mul_closed}.
    $(1 / m) \rmul (m \rmul n) < (1 / n) \rmul (m \rmul n)$ by \cref{real_lt_mul_pos}.
    $(m \rmul n) \rmul (1 / m) < (m \rmul n) \rmul (1 / n)$
        by \cref{real_mul_comm,real_lt_subst_left,real_lt_subst_right}.
    $(m \rmul n) \rmul (1 / m) = n$
        by \cref{real_div,real_mul_assoc,real_mul_comm,real_mul_closed,real_mul_inverse,real_mul_ident}.
    $(m \rmul n) \rmul (1 / n) = m$
        by \cref{real_div,real_mul_assoc,real_mul_closed,real_mul_comm,real_mul_inverse,real_mul_ident}.
    $n < m$.
    Contradiction by \cref{real_lt_trans,real_lt_subst_left,real_lt_irreflexive}.
\end{proof}
% from §28 helper lemma 9: sequential compactness implies Lebesgue number lemma
\begin{lemma}\label{sequentially_compact_lebesgue}
    Suppose $d$ is a metric on $X$.
    Let $T = \metricTopology{d}{X}$.
    Suppose $A$ is an open covering of $X$ with respect to $T$.
    Suppose $X\neq\emptyset$.
    Suppose $X$ is sequentially compact with respect to $T$.
    Then there exists $\delta\in\reals$ such that $0 < \delta$ and
        for all $B$ such that $B\subseteq X$ and
            for all $x,y\in B$ we have $d((x,y)) < \delta$
        there exists $U\in A$ such that $B\subseteq U$.
\end{lemma}
\begin{proof}
    Suppose not.
    Thus for all $\delta\in\reals$ such that $0 < \delta$ there exists $B$ such that
        $B\subseteq X$ and
        for all $U\in A$ we have $B\not\subseteq U$ and
        for all $x,y\in B$ we have $d((x,y)) < \delta$.
    Take $x_0$ such that $x_0\in X$ by \cref{nonempty_inhabited}.
    Hence $x_0\in\unions{A}$ by \cref{open_covering,covering,subseteq}.
    Take $U_0\in A$ such that $x_0\in U_0$ by \cref{unions_iff}.

    Let $R = \{w\in\naturalsPlus\times X \mid
        \text{for all $U\in A$ we have $\metricBall{d}{X}{\snd{w}}{1 / (\fst{w})}\not\subseteq U$}\}$.
    Hence $R$ is a relation by \cref{times_elem_is_tuple,relation}.

    Show that for all $n\in\naturalsPlus$ there exists $x$ such that $n\mathrel{R} x$.
    \begin{subproof}
        Fix $n\in\naturalsPlus$.
        Let $\delta = 1 / n$.
        $0 < \delta$ by \cref{positive_reciprocal_naturalsplus}.
        $n\in\reals$ and $0 < n$
            by \cref{real_naturals_subset_reals,subseteq,real_mul_ident,real_add_ident,real_zero_lt_one,real_one_leq_naturalsplus,real_lt_subst_right,real_lt_trans}.
        Hence $\inv{n}\in\reals$ and $\delta\in\reals$
            by \cref{real_pos_nonzero,real_mul_inverse,real_div,real_mul_closed}.
        Take $B$ such that
            $B\subseteq X$ and
            for all $U\in A$ we have $B\not\subseteq U$ and
            for all $x,y\in B$ we have $d((x,y)) < \delta$.
        We have $B\neq\emptyset$ by \cref{emptyset_subseteq}.
        Take $x$ such that $x\in B$.
        For all $U\in A$ we have $\metricBall{d}{X}{x}{\delta}\not\subseteq U$
            by \cref{subseteq,metric_ball,subseteq_transitive}.
        $x\in X$ by \cref{subseteq}.
        Hence $(n,x)\in\naturalsPlus\times X$ by \cref{times_tuple_intro}.
        Therefore for all $U\in A$ we have $\metricBall{d}{X}{\snd{(n,x)}}{1 / (\fst{(n,x)})}\not\subseteq U$
            by \cref{fst_eq,snd_eq}.
        Thus $(n,x)\in R$.
        $n\mathrel{R} x$ by \cref{relation}.
    \end{subproof}

    Take $s$ such that $s$ is a function on $\naturalsPlus$ and
        for all $n\in\naturalsPlus$ we have $n\mathrel{R} \apply{s}{n}$ by \cref{choice_function}.
    For all $n\in\dom{s}$ we have $s(n)\in X$ by \cref{times_tuple_elim}.
    Hence $s\in\funs{\naturalsPlus}{X}$ by \cref{funs_intro}.
    Therefore $s$ is a sequence in $X$ by \cref{sequence_in}.

    Take $t,x$ such that
        $t$ is a sequence in $X$ and
        $t$ is a subsequence of $s$ and
        $x\in X$ and
        $t$ converges to $x$ in $X$ with respect to $T$ by \cref{sequentially_compact}.
    Take $u$ such that
        $u$ is strictly increasing on $\naturalsPlus$ and
        for all $n\in\naturalsPlus$ we have $t(n) = \apply{s}{\apply{u}{n}}$ by \cref{subsequence}.

    Show that for all $n\in\naturalsPlus$ we have $n \leq u(n)$.
    \begin{subproof}
        Let $P = \{n\in\naturalsPlus \mid n \leq u(n)\}$.
            For all $n$ such that $n\in P$ we have $n\in\naturalsPlus$.
            $P\subseteq\naturalsPlus$ by \cref{subseteq}.
            Hence $1\in P$ by \cref{real_one_in_naturals,strictly_increasing_on_set,funs_type_apply,real_one_leq_naturalsplus}.
            Show that for all $n\in P$ we have $n + 1\in P$.
            \begin{subproof}
                Fix $n$ such that $n\in P$.
                $n \leq u(n)$.
                $n\in\naturalsPlus$.
                Hence $n + 1\in\naturalsPlus$ and $n < n + 1$
                    by \cref{real_naturals_closed_succ,real_naturals_lt_succ_local}.
                $n\in\reals$ by \cref{real_naturals_subset_reals,subseteq}.
                Therefore $u(n) < u(n + 1)$ by \cref{strictly_increasing_on_set}.
                Show that it is wrong that $u(n + 1) < n + 1$.
                \begin{subproof}
                    Suppose not.
                    We have $u(n + 1) < n + 1$.
                    Hence $u(n + 1)\in\naturalsPlus$ by \cref{strictly_increasing_on_set,funs_type_apply}.
                    Therefore $u(n + 1) < n$ or $u(n + 1) = n$ or $u(n + 1) = n + 1$
                        by \cref{real_succ_discrete}.
                    \begin{byCase}
                    \caseOf{$u(n + 1) < n$.}
                        $u(n)\in\reals$ and $u(n + 1)\in\reals$
                            by \cref{strictly_increasing_on_set,funs_type_apply,real_naturals_subset_reals,subseteq}.
                        $n < u(n)$ or $n = u(n)$ by \cref{real_leq_split}.
                        \begin{byCase}
                        \caseOf{$n < u(n)$.}
                            Contradiction by \cref{real_lt_trans,real_lt_irreflexive}.
                        \caseOf{$n = u(n)$.}
                            Contradiction by \cref{real_lt_subst_right,real_lt_trans,real_lt_irreflexive}.
                        \end{byCase}
                    \caseOf{$u(n + 1) = n$.}
                        We have $u(n)\in\reals$
                            by \cref{strictly_increasing_on_set,funs_type_apply,real_naturals_subset_reals,subseteq}.
                        Hence $u(n) < n$ by \cref{real_lt_subst_right}.
                        $n < u(n)$ or $n = u(n)$ by \cref{real_leq_split}.
                        \begin{byCase}
                        \caseOf{$n < u(n)$.}
                            Contradiction by \cref{real_lt_trans,real_lt_irreflexive}.
                        \caseOf{$n = u(n)$.}
                            Contradiction by \cref{real_lt_subst_left,real_lt_irreflexive}.
                        \end{byCase}
                    \caseOf{$u(n + 1) = n + 1$.}
                        Contradiction by \cref{real_lt_subst_left,real_naturals_closed_succ,real_naturals_subset_reals,subseteq,real_lt_irreflexive}.
                    \end{byCase}
                \end{subproof}
                Hence $u(n + 1)\in\reals$ and $n + 1\in\reals$
                    by \cref{strictly_increasing_on_set,funs_type_apply,real_naturals_closed_succ,real_naturals_subset_reals,subseteq}.
                Therefore $n + 1 < u(n + 1)$ or $n + 1 = u(n + 1)$ or $u(n + 1) < n + 1$
                    by \cref{real_lt_trichotomy}.
                $n + 1 < u(n + 1)$ or $n + 1 = u(n + 1)$.
                Hence $n + 1 \leq u(n + 1)$.
                Therefore $n + 1\in P$.
            \end{subproof}
            Thus for all $n\in\naturalsPlus$ we have $n\in P$ by \cref{real_naturals_induction}.
        Thus for all $n\in\naturalsPlus$ we have $n \leq u(n)$ by \cref{subseteq}.
    \end{subproof}

    Hence $x\in\unions{A}$ by \cref{open_covering,covering,subseteq}.
    Take $U\in A$ such that $x\in U$ by \cref{unions_iff}.
    Therefore $U$ is open in $X$ with respect to $T$ and $U\subseteq X$
        by \cref{metric_topology,basis_topology_is_topology,metric_basis_is_basis,open_covering,open_in,topology_on,subseteq,powerset_elim}.
    Take $\delta\in\reals$ such that $0 < \delta$ and $\metricBall{d}{X}{x}{\delta}\subseteq U$
        by \cref{metric_open_iff}.
    Let $m_2 = 1 + 1$.
    Let $\epsilon = \delta / m_2$.
    $1\in\reals$ and $0\in\reals$ and $0 < 1$ by \cref{real_mul_ident,real_add_ident,real_zero_lt_one}.
    Hence $m_2\in\reals$ and $0 < m_2$
        by \cref{real_add_closed,real_lt_add,real_add_ident,real_lt_subst_left,real_lt_trans}.
    $m_2\neq 0$ by \cref{real_pos_nonzero}.
    Thus $\inv{m_2}\in\reals$ and $0 < \inv{m_2}$ by \cref{real_mul_inverse,real_inv_pos}.
    $\epsilon = \delta \rmul \inv{m_2}$ by \cref{real_div}.
    Hence $\epsilon\in\reals$ and $0 < \epsilon$
        by \cref{real_mul_closed,real_add_ident,real_lt_mul_pos,real_mul_zero,real_lt_subst_left}.
    Take $N\in\naturalsPlus$ such that $0 < 1 / N$ and $1 / N < \epsilon$ by \cref{real_small_reciprocal}.

            We have $\metricBall{d}{X}{x}{\epsilon}\subseteq X$ by \cref{metric_ball,subseteq}.
            Show that for all $y\in\metricBall{d}{X}{x}{\epsilon}$ there exists $\eta\in\reals$ such that
                $0 < \eta$ and $\metricBall{d}{X}{y}{\eta}\subseteq\metricBall{d}{X}{x}{\epsilon}$.
            \begin{subproof}
                Fix $y$ such that $y\in\metricBall{d}{X}{x}{\epsilon}$.
                $y\in X$ and $d((x,y)) < \epsilon$ by \cref{metric_ball}.
                Hence $d((x,y))\in\reals$ by \cref{metric_on,times_tuple_intro,funs_type_apply}.
                Let $\eta = \epsilon - d((x,y))$.
                $0 < \eta$ by \cref{real_lt_imp_pos_diff}.
                Show that for all $z$ such that $z\in\metricBall{d}{X}{y}{\eta}$ we have $z\in\metricBall{d}{X}{x}{\epsilon}$.
                \begin{subproof}
                        Fix $z$ such that $z\in\metricBall{d}{X}{y}{\eta}$.
                        $z\in X$ and $d((y,z)) < \eta$ by \cref{metric_ball}.
                        Hence $d((y,z))\in\reals$ by \cref{metric_on,times_tuple_intro,funs_type_apply}.
                        We have $\eta\in\reals$ by \cref{real_add_inverse,real_sub,real_add_closed}.
                        Therefore $d((x,z)) \leq d((x,y)) + d((y,z))$ by \cref{metric_on,metric_triangle_inequality}.
                        $d((x,z))\in\reals$ by \cref{metric_on,times_tuple_intro,funs_type_apply}.
                        Hence $\eta + d((x,y))\in\reals$ and $d((x,y)) + d((y,z))\in\reals$
                            by \cref{real_add_closed}.
                        Therefore $d((x,y)) + d((y,z)) < \eta + d((x,y))$
                            by \cref{real_lt_add,real_add_comm,real_lt_subst_left}.
                        $d((x,z)) < d((x,y)) + d((y,z))$ or
                            $d((x,z)) = d((x,y)) + d((y,z))$ by \cref{real_leq_split}.
                        Hence $d((x,z)) < \eta + d((x,y))$
                            by \cref{real_lt_trans,metric_on,times_tuple_intro,funs_type_apply,real_add_closed,real_lt_subst_left}.
                        Thus $z\in\metricBall{d}{X}{x}{\epsilon}$
                            by \cref{real_sub,real_add_assoc,real_add_comm,real_add_inverse,real_add_ident,metric_ball,real_lt_subst_right}.
                \end{subproof}
                $\metricBall{d}{X}{y}{\eta}\subseteq\metricBall{d}{X}{x}{\epsilon}$ by \cref{subseteq}.
            \end{subproof}
        Hence $\metricBall{d}{X}{x}{\epsilon}$ is open in $X$ with respect to $T$ by \cref{metric_open_iff}.
    Therefore $\metricBall{d}{X}{x}{\epsilon}$ is a neighborhood of $x$ in $X$ with respect to $T$
        by \cref{metric_zero_characterization,metric_on,metric_ball,real_lt_subst_left,neighborhood}.

    Take $N_0\in\naturalsPlus$ such that for all $n\in\naturalsPlus$ if $N_0 \leq n$ then
        $t(n)\in\metricBall{d}{X}{x}{\epsilon}$ by \cref{converges_to}.
    Show that there exists $n\in\naturalsPlus$ such that $N_0 \leq n$ and $N \leq n$.
    \begin{subproof}
        \begin{byCase}
        \caseOf{$N_0 = N$.}
            Take $n$ such that $n = N$.
            $N_0 \leq n$ and $N \leq n$.
        \caseOf{$N_0 \neq N$.}
            Thus $N_0 < N$ or $N < N_0$ by \cref{real_naturals_subset_reals,subseteq,real_lt_trichotomy}.
            \begin{byCase}
            \caseOf{$N_0 < N$.}
                Take $n$ such that $n = N$.
                $N_0 \leq n$ and $N \leq n$ by \cref{real_lt_imp_leq}.
            \caseOf{$N < N_0$.}
                Take $n$ such that $n = N_0$.
                $N_0 \leq n$ and $N \leq n$ by \cref{real_lt_imp_leq}.
            \end{byCase}
        \end{byCase}
    \end{subproof}
    Take $n\in\naturalsPlus$ such that $N_0 \leq n$ and $N \leq n$.

    Hence $t(n)\in\metricBall{d}{X}{x}{\epsilon}$.
    $d((x,t(n))) < \epsilon$ by \cref{metric_ball}.
    $u(n)\in\naturalsPlus$ by \cref{strictly_increasing_on_set,funs_type_apply}.
    Hence $n \leq u(n)$.
    Therefore $N \leq u(n)$
        by \cref{real_naturals_subset_reals,subseteq,strictly_increasing_on_set,funs_type_apply,real_leq_trans}.
    Hence $1 / u(n) \leq 1 / N$ by \cref{naturalsplus_reciprocal_antitone}.
    $u(n)\in\reals$ and $N\in\reals$ by \cref{real_naturals_subset_reals,subseteq}.
    $0 < u(n)$ and $0 < N$
        by \cref{real_zero_lt_one,real_one_in_naturals,real_one_leq_naturalsplus,real_lt_trans,strictly_increasing_on_set,funs_type_apply}.
    $1 / u(n)\in\reals$ and $1 / N\in\reals$
        by \cref{real_pos_nonzero,real_div,real_mul_closed,real_mul_inverse}.
    Therefore $1 / u(n) < \epsilon$ by \cref{real_lt_trans,real_lt_subst_left}.

    Show that for all $z$ such that $z\in\metricBall{d}{X}{t(n)}{1 / u(n)}$ we have $z\in U$.
        \begin{subproof}
            Fix $z$ such that $z\in\metricBall{d}{X}{t(n)}{1 / u(n)}$.
            $z\in X$ and $d((t(n),z)) < 1 / u(n)$ by \cref{metric_ball}.
            Hence $t(n)\in X$ by \cref{sequence_in,funs_type_apply}.
            Therefore $d((x,z)) \leq d((x,t(n))) + d((t(n),z))$ by \cref{metric_on,metric_triangle_inequality}.
            Show that $d((x,z)) < \epsilon + \epsilon$.
            \begin{subproof}
                We have $d((t(n),z))\in\reals$ by \cref{metric_on,times_tuple_intro,funs_type_apply}.
                $1 / u(n)\in\reals$.
                $\epsilon\in\reals$.
                Hence $d((t(n),z)) < \epsilon$ by \cref{real_lt_trans}.
                $d((x,t(n)))\in\reals$ by \cref{metric_on,times_tuple_intro,funs_type_apply}.
                Hence $d((x,t(n))) + d((t(n),z)) < d((x,t(n))) + \epsilon$
                    by \cref{real_lt_add,real_add_comm,real_lt_subst_left,real_lt_subst_right}.
                Therefore $d((x,t(n))) + \epsilon < \epsilon + \epsilon$ by \cref{real_lt_add}.
                $d((x,t(n))) + d((t(n),z))\in\reals$ and $d((x,t(n))) + \epsilon\in\reals$
                    and $\epsilon + \epsilon\in\reals$ by \cref{real_add_closed}.
                Hence $d((x,t(n))) + d((t(n),z)) < \epsilon + \epsilon$ by \cref{real_lt_trans}.
                $d((x,z))\in\reals$ by \cref{metric_on,times_tuple_intro,funs_type_apply}.
                $d((x,z)) < d((x,t(n))) + d((t(n),z))$ or
                    $d((x,z)) = d((x,t(n))) + d((t(n),z))$ by \cref{real_leq_split}.
                Therefore $d((x,z)) < \epsilon + \epsilon$ by \cref{real_lt_trans,real_lt_subst_left}.
            \end{subproof}
            Hence $d((x,z)) < \delta$
                by \cref{real_div,real_mul_assoc,real_mul_inverse,real_mul_comm,real_mul_ident,real_right_distrib,real_lt_subst_right}.
            Therefore $z\in U$ by \cref{metric_ball,subseteq}.
    \end{subproof}
    Hence $\metricBall{d}{X}{t(n)}{1 / u(n)}\subseteq U$ by \cref{subseteq}.

    Thus $(u(n),t(n))\in R$ by \cref{pair_eq_iff}.
    Therefore for all $V\in A$ we have $\metricBall{d}{X}{t(n)}{1 / u(n)}\not\subseteq V$.
    Hence $\metricBall{d}{X}{t(n)}{1 / u(n)}\not\subseteq U$.
    Contradiction.
\end{proof}

% from §28 helper lemma 11: iteration sequence on naturalsPlus
\begin{lemma}\label{iteration_sequence_naturalsplus}
    Suppose $f\in\funs{A}{A}$.
    Suppose $a_0\in A$.
    Then there exists $u$ such that
        $u$ is a sequence in $A$ and
        $u(1) = a_0$ and
        for all $n\in\naturalsPlus$ we have $u(n + 1) = \apply{f}{u(n)}$.
\end{lemma}
\begin{proof}
    Let $S = \iterationSet{f}{a_0}$.
    Show that for all $n\in\naturalsPlus$ we have $n\in S$.
    \begin{subproof}
        $S\subseteq\naturalsPlus$ by \cref{iteration_set_elim,subseteq}.
        $1\in\naturalsPlus$ by \cref{real_one_in_naturals}.
            Let $u = \{(1,a_0)\}$.
            For all $x,y,y'$ such that $(x,y)\in u$ and $(x,y')\in u$ we have $y = y'$
                by \cref{singleton_iff,pair_eq_iff}.
            Hence $u$ is a function by \cref{singleton_iff,relation,rightunique}.
            For all $x$ such that $x\in\dom{u}$ we have $x\in\initialSegment{1}$
                by \cref{dom,singleton_iff,pair_eq_iff,real_one_in_naturals,initial_segment_naturalsplus}.
            $\dom{u}\subseteq\initialSegment{1}$ by \cref{subseteq}.
            Show that for all $x$ such that $x\in\initialSegment{1}$ we have $x\in\dom{u}$.
            \begin{subproof}
                Fix $x$ such that $x\in\initialSegment{1}$.
                $x\in\naturalsPlus$ and $x \leq 1$ by \cref{initial_segment_naturalsplus}.
                Hence $x\in\reals$ and $1\in\reals$
                    by \cref{real_naturals_subset_reals,subseteq,real_one_in_naturals}.
                Therefore $1 \leq x$ by \cref{real_one_leq_naturalsplus}.
                $x = 1$ by \cref{real_leq_antisym}.
                Hence $(x,a_0)\in u$ by \cref{singleton_iff,pair_eq_iff}.
                $x\in\dom{u}$ by \cref{dom}.
            \end{subproof}
            Hence $\dom{u} = \initialSegment{1}$ by \cref{subseteq,subseteq_antisymmetric}.
            Therefore $u(1) = a_0$ by \cref{singleton_intro,function_apply_intro}.
            Show that for all $k\in\naturalsPlus$ such that $k + 1 \leq 1$ we have $u(k + 1) = \apply{f}{u(k)}$.
            \begin{subproof}
                Fix $k\in\naturalsPlus$.
                Assume $k + 1 \leq 1$.
                Suppose not.
                $k < k + 1$ by \cref{real_naturals_lt_succ_local}.
                $1\in\reals$ and $k\in\reals$ and $k + 1\in\reals$
                    by \cref{real_naturals_subset_reals,real_one_in_naturals,real_naturals_closed_succ,subseteq}.
                $1 < k$ or $1 = k$ by \cref{real_one_leq_naturalsplus}.
                Therefore $1 < k + 1$ by \cref{real_lt_trans,real_lt_subst_left}.
                $k + 1 < 1$ or $k + 1 = 1$ by \cref{real_leq_split}.
                Thus $1 < 1$ by \cref{real_lt_trans,real_lt_subst_right}.
                Contradiction by \cref{real_naturals_subset_reals,real_one_in_naturals,subseteq,real_lt_irreflexive}.
            \end{subproof}
            Therefore $1\in S$ by \cref{iteration_sequence,iteration_set_intro}.
        Show that for all $n\in S$ we have $n + 1\in S$.
        \begin{subproof}
            Fix $n$ such that $n\in S$.
            $n\in\naturalsPlus$ by \cref{iteration_set_elim}.
            Take $w$ such that
                $w$ is an iteration sequence for $f$ on $n$ starting at $a_0$
                by \cref{iteration_set_elim}.
            Hence $n + 1\in\naturalsPlus$ by \cref{real_naturals_closed_succ}.
            Let $v = w\union\{(n + 1,\apply{f}{w(n)})\}$.
            Therefore $v$ is a relation by \cref{singleton_iff,relation,iteration_sequence,union_relations_is_relation}.
            Show that for all $x,y,y'$ such that $(x,y)\in v$ and $(x,y')\in v$ we have $y = y'$.
            \begin{subproof}
                Fix $x$.
                Fix $y$.
                Fix $y'$.
                Assume $(x,y)\in v$ and $(x,y')\in v$.
                \begin{byCase}
                \caseOf{$x = n + 1$.}
                    $(n + 1,y)\in v$.
                    $(n + 1,y')\in v$.
                    It is wrong that $(n + 1,y)\in w$
                        by \cref{dom,iteration_sequence,initial_segment_naturalsplus,real_naturals_lt_succ_local,real_naturals_subset_reals,real_naturals_closed_succ,subseteq,real_leq_split,real_lt_trans,real_lt_subst_right,real_lt_irreflexive}.
                    Hence $(n + 1,y)\in\{(n + 1,\apply{f}{w(n)})\}$ by \cref{union_iff}.
                    Therefore $(n + 1,y) = (n + 1,\apply{f}{w(n)})$ by \cref{singleton_iff}.
                    $y = \apply{f}{w(n)}$ by \cref{pair_eq_iff}.
                    It is wrong that $(n + 1,y')\in w$
                        by \cref{dom,iteration_sequence,initial_segment_naturalsplus,real_naturals_lt_succ_local,real_naturals_subset_reals,real_naturals_closed_succ,subseteq,real_leq_split,real_lt_trans,real_lt_subst_right,real_lt_irreflexive}.
                    Hence $(n + 1,y')\in\{(n + 1,\apply{f}{w(n)})\}$ by \cref{union_iff}.
                    Therefore $(n + 1,y') = (n + 1,\apply{f}{w(n)})$ by \cref{singleton_iff}.
                    $y' = \apply{f}{w(n)}$ by \cref{pair_eq_iff}.
                    Hence $y = y'$.
                \caseOf{$x \neq n + 1$.}
                    $(x,y)\in w$ by \cref{union_iff,singleton_iff,pair_eq_iff}.
                    Hence $(x,y')\in w$ by \cref{union_iff,singleton_iff,pair_eq_iff}.
                    Therefore $y = y'$ by \cref{iteration_sequence,function_rightunique}.
                \end{byCase}
            \end{subproof}
            Therefore $v$ is a function by \cref{rightunique}.
            Show that for all $x$ such that $x\in\dom{v}$ we have $x\in\initialSegment{n + 1}$.
            \begin{subproof}
                Fix $x$ such that $x\in\dom{v}$.
                Take $y$ such that $(x,y)\in v$ by \cref{dom}.
                Thus $(x,y)\in w$ or $(x,y)\in\{(n + 1,\apply{f}{w(n)})\}$ by \cref{union_iff}.
                \begin{byCase}
                \caseOf{$(x,y)\in w$.}
                    $x\in\initialSegment{n + 1}$
                        by \cref{dom,iteration_sequence,initial_segment_naturalsplus,real_naturals_subset_reals,real_naturals_closed_succ,subseteq,real_leq_trans,real_naturals_lt_succ_local}.
                \caseOf{$(x,y)\in\{(n + 1,\apply{f}{w(n)})\}$.}
                    Hence $x\in\initialSegment{n + 1}$
                        by \cref{singleton_iff,pair_eq_iff,initial_segment_naturalsplus}.
                \end{byCase}
            \end{subproof}
            $\dom{v}\subseteq\initialSegment{n + 1}$ by \cref{subseteq}.
            Show that for all $x$ such that $x\in\initialSegment{n + 1}$ we have $x\in\dom{v}$.
            \begin{subproof}
                Fix $x$ such that $x\in\initialSegment{n + 1}$.
                $x\in\naturalsPlus$ and $x \leq n + 1$ by \cref{initial_segment_naturalsplus}.
                \begin{byCase}
                \caseOf{$x = n + 1$.}
                    Hence $x\in\dom{v}$ by \cref{cons_left,union_intro_right,dom}.
                \caseOf{$x \neq n + 1$.}
                    Therefore $x \leq n$ by \cref{real_naturals_leq_succ_elim}.
                    $x\in\dom{w}$ by \cref{initial_segment_naturalsplus,iteration_sequence}.
                    Hence $x\in\dom{v}$ by \cref{dom,union_iff}.
                \end{byCase}
            \end{subproof}
            Hence $\dom{v} = \initialSegment{n + 1}$ by \cref{subseteq,subseteq_antisymmetric}.
            Thus $v(1) = a_0$
                by \cref{real_one_in_naturals,real_one_leq_naturalsplus,initial_segment_naturalsplus,iteration_sequence,function_apply_elim,union_intro_left,function_apply_intro}.
            Show that for all $k\in\naturalsPlus$ such that $k + 1 \leq n + 1$ we have $v(k + 1) = \apply{f}{v(k)}$.
            \begin{subproof}
                Fix $k\in\naturalsPlus$.
                Assume $k + 1 \leq n + 1$.
                \begin{byCase}
                \caseOf{$k + 1 = n + 1$.}
                    $k = n$.
                    Hence $v(k + 1) = \apply{f}{w(n)}$
                        by \cref{cons_left,union_intro_right,function_apply_intro}.
                    $v(k + 1) = \apply{f}{w(k)}$.
                    Therefore $k \leq n$.
                    $v(k) = w(k)$
                        by \cref{initial_segment_naturalsplus,iteration_sequence,function_apply_elim,union_intro_left,function_apply_intro}.
                    Hence $v(k + 1) = \apply{f}{v(k)}$.
                \caseOf{$k + 1 \neq n + 1$.}
                    $k + 1\in\naturalsPlus$ by \cref{real_naturals_closed_succ}.
                    Hence $k + 1 \leq n$ by \cref{real_naturals_leq_succ_elim}.
                    Therefore $k + 1\in\initialSegment{n}$ by \cref{initial_segment_naturalsplus}.
                    $(k + 1,w(k + 1))\in w$
                        by \cref{iteration_sequence,function_apply_elim}.
                    Hence $v(k + 1) = w(k + 1)$
                        by \cref{union_intro_left,dom,function_apply_elim,function_rightunique}.
                    $k\in\naturalsPlus$ and $k + 1 \leq n$.
                    Therefore $w(k + 1) = \apply{f}{w(k)}$ by \cref{iteration_sequence}.
                    $v(k + 1) = \apply{f}{w(k)}$.
                    Hence $k \leq n$ by \cref{real_naturals_succ_bound_elim_local}.
                    Therefore $(k,w(k))\in v$
                        by \cref{initial_segment_naturalsplus,iteration_sequence,function_apply_elim,union_intro_left}.
                    $w(k) = v(k)$ by \cref{dom,function_apply_elim,function_rightunique}.
                    Hence $v(k + 1) = \apply{f}{v(k)}$.
                \end{byCase}
            \end{subproof}
            Hence $v$ is an iteration sequence for $f$ on $n + 1$ starting at $a_0$ by \cref{iteration_sequence_intro}.
            Therefore $n + 1\in S$ by \cref{iteration_set_intro}.
        \end{subproof}
        Thus for all $n\in\naturalsPlus$ we have $n\in S$ by \cref{real_naturals_induction}.
    \end{subproof}
    Let $F = \{u \mid \text{there exists $n\in\naturalsPlus$ such that
        $u$ is an iteration sequence for $f$ on $n$ starting at $a_0$} \}$.
    For all $u\in F$ we have $u$ is a function by \cref{iteration_sequence}.
    Show that for all $u,v$ such that $u\in F$ and $v\in F$ we have $u\subseteq v$ or $v\subseteq u$.
    \begin{subproof}
        Fix $u$.
        Fix $v$.
        Assume $u\in F$ and $v\in F$.
        Take $m\in\naturalsPlus$ such that
            $u$ is an iteration sequence for $f$ on $m$ starting at $a_0$.
        Take $n\in\naturalsPlus$ such that
            $v$ is an iteration sequence for $f$ on $n$ starting at $a_0$.
        $u$ is a function and $\dom{u} = \initialSegment{m}$ and $u(1) = a_0$
            by \cref{iteration_sequence}.
        Hence for all $k\in\naturalsPlus$ such that $k + 1 \leq m$ we have $u(k + 1) = \apply{f}{u(k)}$
            by \cref{iteration_sequence}.
        $v$ is a function and $\dom{v} = \initialSegment{n}$ and $v(1) = a_0$
            by \cref{iteration_sequence}.
        Therefore for all $k\in\naturalsPlus$ such that $k + 1 \leq n$ we have $v(k + 1) = \apply{f}{v(k)}$
            by \cref{iteration_sequence}.
        $m = n$ or $m \neq n$.
        \begin{byCase}
        \caseOf{$m = n$.}
            Thus $u\subseteq v$ or $v\subseteq u$ by \cref{iteration_sequence_unique,subseteq_refl}.
        \caseOf{$m \neq n$.}
            $m\in\reals$ and $n\in\reals$ by \cref{real_naturals_subset_reals,subseteq}.
            Thus $m < n$ or $n < m$ by \cref{real_lt_trichotomy}.
            \begin{byCase}
            \caseOf{$m < n$.}
                Let $B = \{k\in\naturalsPlus \mid \text{if $k\leq m$ then $u(k) = v(k)$}\}$.
                $1\in\naturalsPlus$ by \cref{real_one_in_naturals}.
                Show that if $1\leq m$ then $u(1) = v(1)$.
                \begin{subproof}
                    Follows by \cref{iteration_sequence}.
                \end{subproof}
                Hence $1\in B$.
                Show that for all $k\in B$ we have $k + 1\in B$.
                \begin{subproof}
                    Fix $k\in B$.
                    $k\in\naturalsPlus$.
                    Hence $k + 1\in\naturalsPlus$ by \cref{real_naturals_closed_succ}.
                    Show that if $k + 1\leq m$ then $u(k + 1) = v(k + 1)$.
                    \begin{subproof}
                        Suppose $k + 1\leq m$.
                        $k \leq m$ by \cref{real_naturals_succ_bound_elim_local}.
                        Hence $u(k) = v(k)$.
                        Therefore $m\leq n$ by \cref{real_lt_imp_leq}.
                        $k + 1\leq n$ by \cref{real_leq_trans,real_add_closed,real_naturals_subset_reals,subseteq}.
                        Hence $u(k + 1) = \apply{f}{u(k)}$ by \cref{iteration_sequence}.
                        Therefore $v(k + 1) = \apply{f}{v(k)}$ by \cref{iteration_sequence}.
                        $u(k + 1) = v(k + 1)$.
                    \end{subproof}
                    Hence $k + 1\in B$.
                \end{subproof}
                Therefore for all $k\in\initialSegment{m}$ we have $u(k) = v(k)$
                    by \cref{subseteq,real_naturals_induction,initial_segment_naturalsplus}.
                For all $w$ such that $w\in u$ we have $w\in v$
                    by \cref{iteration_sequence,function_member_elim,initial_segment_naturalsplus,real_naturals_subset_reals,subseteq,real_leq_trans,real_lt_imp_leq,function_apply_elim}.
                Hence $u\subseteq v$ by \cref{subseteq}.
            \caseOf{$n < m$.}
                Let $C = \{k\in\naturalsPlus \mid \text{if $k\leq n$ then $v(k) = u(k)$}\}$.
                $1\in\naturalsPlus$ by \cref{real_one_in_naturals}.
                Show that if $1\leq n$ then $v(1) = u(1)$.
                \begin{subproof}
                    Follows by \cref{iteration_sequence}.
                \end{subproof}
                Hence $1\in C$.
                Show that for all $k\in C$ we have $k + 1\in C$.
                \begin{subproof}
                    Fix $k\in C$.
                    $k\in\naturalsPlus$.
                    Hence $k + 1\in\naturalsPlus$ by \cref{real_naturals_closed_succ}.
                    Show that if $k + 1\leq n$ then $v(k + 1) = u(k + 1)$.
                    \begin{subproof}
                        Suppose $k + 1\leq n$.
                        $k \leq n$ by \cref{real_naturals_succ_bound_elim_local}.
                        Hence $v(k) = u(k)$.
                        Therefore $n\leq m$ by \cref{real_lt_imp_leq}.
                        $k + 1\leq m$ by \cref{real_leq_trans,real_add_closed,real_naturals_subset_reals,subseteq}.
                        Hence $v(k + 1) = \apply{f}{v(k)}$ by \cref{iteration_sequence}.
                        Therefore $u(k + 1) = \apply{f}{u(k)}$ by \cref{iteration_sequence}.
                        $v(k + 1) = u(k + 1)$.
                    \end{subproof}
                    Hence $k + 1\in C$.
                \end{subproof}
                Therefore for all $k\in\initialSegment{n}$ we have $v(k) = u(k)$
                    by \cref{subseteq,real_naturals_induction,initial_segment_naturalsplus}.
                For all $w$ such that $w\in v$ we have $w\in u$
                    by \cref{iteration_sequence,function_member_elim,initial_segment_naturalsplus,real_naturals_subset_reals,subseteq,real_leq_trans,real_lt_imp_leq,function_apply_elim}.
                Hence $v\subseteq u$ by \cref{subseteq}.
            \end{byCase}
        \end{byCase}
    \end{subproof}
    Let $u = \unions{F}$.
    $F$ is a family of functions.
    Hence $u$ is a function by \cref{unions_of_compatible_family_of_function_is_function}.
    Show that for all $x$ such that $x\in\dom{u}$ we have $x\in\naturalsPlus$.
    \begin{subproof}
        Follows by \cref{dom,unions_iff,iteration_sequence,initial_segment_naturalsplus}.
    \end{subproof}
    Show that for all $n\in\naturalsPlus$ we have $n\in\dom{u}$.
    \begin{subproof}
        Follows by \cref{subseteq,iteration_set_elim,iteration_sequence,initial_segment_naturalsplus,dom,unions_iff}.
    \end{subproof}
    Therefore $\dom{u} = \naturalsPlus$ by \cref{subseteq,subseteq_antisymmetric}.
    Show that for all $n\in\naturalsPlus$ we have $u(n)\in A$.
    \begin{subproof}
        Fix $n\in\naturalsPlus$.
        Take $v\in F$ such that $(n,u(n))\in v$
            by \cref{function_apply_elim,unions_iff}.
        Take $n_0\in\naturalsPlus$ such that
            $v$ is an iteration sequence for $f$ on $n_0$ starting at $a_0$.
        $n \leq n_0$ by \cref{dom,iteration_sequence,initial_segment_naturalsplus}.
        Let $B = \{k\in\naturalsPlus \mid \text{if $k\leq n$ then $v(k)\in A$}\}$.
        $B\subseteq\naturalsPlus$.
        Hence $1\in B$ by \cref{real_one_in_naturals,iteration_sequence}.
        Show that for all $k\in B$ we have $k + 1\in B$.
        \begin{subproof}
            Fix $k\in B$.
            Show that if $k + 1\leq n$ then $v(k + 1)\in A$.
            \begin{subproof}
                Suppose $k + 1\leq n$.
                $k \leq n$ by \cref{real_naturals_succ_bound_elim_local}.
                Hence $v(k)\in A$.
                Therefore $v(k + 1) = \apply{f}{v(k)}$
                    by \cref{real_naturals_closed_succ,real_naturals_subset_reals,subseteq,real_leq_trans,iteration_sequence}.
                $\apply{f}{v(k)}\in A$ by \cref{funs_type_apply}.
                Hence $v(k + 1)\in A$.
            \end{subproof}
            $k + 1\in B$.
        \end{subproof}
        Thus for all $k\in\naturalsPlus$ we have $k\in B$ by \cref{real_naturals_induction}.
        Thus $v(n)\in A$ by \cref{initial_segment_naturalsplus}.
        Thus $u(n)\in A$ by \cref{function_apply_intro}.
    \end{subproof}
    Take $v\in F$ such that $(1,u(1))\in v$
        by \cref{real_one_in_naturals,function_apply_elim,unions_iff}.
    Take $n_0\in\naturalsPlus$ such that
        $v$ is an iteration sequence for $f$ on $n_0$ starting at $a_0$.
    Show that $u(1) = a_0$.
    \begin{subproof}
        Follows by \cref{iteration_sequence,dom,function_apply_intro}.
    \end{subproof}
    Show that for all $n\in\naturalsPlus$ we have $u(n + 1) = \apply{f}{u(n)}$.
    \begin{subproof}
        Fix $n\in\naturalsPlus$.
        Take $w\in F$ such that $(n + 1,u(n + 1))\in w$
            by \cref{real_naturals_closed_succ,function_apply_elim,unions_iff}.
        Take $m\in\naturalsPlus$ such that
            $w$ is an iteration sequence for $f$ on $m$ starting at $a_0$.
        $n + 1\in\dom{w}$ by \cref{dom}.
        Hence $w(n + 1) = u(n + 1)$ by \cref{function_apply_intro}.
        Therefore $n + 1 \leq m$ by \cref{iteration_sequence,initial_segment_naturalsplus}.
        $u(n + 1) = \apply{f}{w(n)}$ by \cref{function_apply_intro,iteration_sequence}.
        Hence $u(n) = w(n)$
            by \cref{real_naturals_succ_bound_elim_local,initial_segment_naturalsplus,iteration_sequence,function_apply_elim,unions_iff,function_apply_intro}.
        Therefore $u(n + 1) = \apply{f}{u(n)}$.
    \end{subproof}
    Hence $u\in\funs{\naturalsPlus}{A}$ by \cref{funs_intro}.
    Show that $u$ is a sequence in $A$.
    \begin{subproof}
        Follows by \cref{sequence_in}.
    \end{subproof}
    Show that there exists $w$ such that
        $w$ is a sequence in $A$ and
        $w(1) = a_0$ and
        for all $n\in\naturalsPlus$ we have $w(n + 1) = \apply{f}{w(n)}$.
    \begin{subproof}
        Trivial.
    \end{subproof}
\end{proof}
% from §28 helper lemma 10: sequential compactness implies finite epsilon-ball cover
\begin{lemma}\label{sequentially_compact_finite_ball_cover}
    Suppose $d$ is a metric on $X$.
    Let $T = \metricTopology{d}{X}$.
    Suppose $X$ is sequentially compact with respect to $T$.
    Suppose $\epsilon\in\reals$.
    Suppose $0 < \epsilon$.
    Then there exists $B$ such that
        $B$ is finite and
        $B$ is a covering of $X$ and
        for all $U\in B$ there exists $x\in X$ such that $U = \metricBall{d}{X}{x}{\epsilon}$.
\end{lemma}
\begin{proof}
    \begin{byCase}
    \caseOf{$X = \emptyset$.}
        Let $B = \emptyset$.
        $B$ is finite by \cref{emptyset_is_finite}.
        For all $U\in B$ we have $\exists x\in X. U = \metricBall{d}{X}{x}{\epsilon}$.
        Hence $B$ is a covering of $X$ by \cref{emptyset_subseteq,unions_emptyset,covering}.
        Show that there exists $C$ such that
            $C$ is finite and
            $C$ is a covering of $X$ and
            for all $U\in C$ there exists $x\in X$ such that $U = \metricBall{d}{X}{x}{\epsilon}$.
        \begin{subproof}
            Trivial.
        \end{subproof}
    \caseOf{$X \neq \emptyset$.}
        Suppose not.
        Thus for all $B$ such that
            $B$ is finite and
            for all $U\in B$ we have $\exists x\in X. U = \metricBall{d}{X}{x}{\epsilon}$
        we have $B$ is not a covering of $X$.
        Let $F = \{B\in\pow{X} \mid \text{$B$ is finite}\}$.
        Show that for all $B\in F$ there exists $x\in X$ such that
            for all $y\in B$ we have $x\notin\metricBall{d}{X}{y}{\epsilon}$.
        \begin{subproof}
            Fix $B$ such that $B\in F$.
            $B$ is finite and $B\subseteq X$ by \cref{powerset_elim}.
            Let $g = \{(x,\metricBall{d}{X}{x}{\epsilon}) \mid x\in B\}$.
                $g$ is a relation by \cref{relation}.
                For all $x,U_1,U_2$ such that $(x,U_1)\in g$ and $(x,U_2)\in g$ we have $U_1 = U_2$
                    by \cref{pair_eq_iff}.
                Hence $g$ is right-unique by \cref{rightunique}.
                    Show that for all $x$ such that $x\in\dom{g}$ we have $x\in B$.
                    \begin{subproof}
                        Follows by \cref{dom,pair_eq_iff}.
                    \end{subproof}
                    $\dom{g}\subseteq B$ by \cref{subseteq}.
                    Show that for all $x$ such that $x\in B$ we have $x\in\dom{g}$.
                    \begin{subproof}
                        Follows by \cref{dom}.
                    \end{subproof}
                    Hence $\dom{g} = B$ by \cref{subseteq,subseteq_antisymmetric}.
                $g$ is a function on $B$.
            Let $G = \img{g}{B}$.
            Hence $G$ is finite by \cref{finite_image_function}.
            For all $U\in G$ we have $\exists x\in X. U = \metricBall{d}{X}{x}{\epsilon}$
                by \cref{img_iff,pair_eq_iff,subseteq}.
            Hence $X\not\subseteq\unions{G}$ by \cref{covering}.
            For all $U\in G$ we have $U\subseteq X$.
            $\unions{G}\subseteq X$ by \cref{unions_family}.
            Hence $\unions{G}\neq X$ by \cref{subseteq_refl}.
            Therefore $\unions{G}\subset X$ by \cref{subset}.
            Take $x\in X$ such that $x\notin\unions{G}$ by \cref{subset_witness}.
            For all $y\in B$ we have $(y,\metricBall{d}{X}{y}{\epsilon})\in g$.
            Hence for all $y\in B$ we have $\metricBall{d}{X}{y}{\epsilon}\in G$
                by \cref{function_apply_iff,img_iff}.
            Therefore for all $y\in B$ we have $x\notin\metricBall{d}{X}{y}{\epsilon}$ by \cref{unions_iff}.
            There exists $z\in X$ such that for all $y\in B$ we have $z\notin\metricBall{d}{X}{y}{\epsilon}$.
        \end{subproof}
        Let $R = \{ w\in F\times X \mid
            \text{for all $y\in\fst{w}$ we have $\snd{w}\notin\metricBall{d}{X}{y}{\epsilon}$} \}$.
        $R$ is a relation by \cref{relation,subseteq,times_elem_is_tuple}.
        For all $B\in F$ there exists $x$ such that $B\mathrel{R} x$
            by \cref{times_tuple_intro,fst_eq,snd_eq}.
        Take $c$ such that $c$ is a function on $F$ and for all $B\in F$ we have $B\mathrel{R} \apply{c}{B}$ by \cref{choice_function}.
        For all $B\in F$ we have $\apply{c}{B}\in X$ and for all $y\in B$ we have $\apply{c}{B}\notin\metricBall{d}{X}{y}{\epsilon}$
            by \cref{subseteq,times_tuple_elim,fst_eq,snd_eq}.
        $c\in\funs{F}{X}$ by \cref{funs_intro,subseteq}.
        Let $S = \{ w\in F\times F \mid
            \text{$\snd{w} = \fst{w}\union\{\apply{c}{\fst{w}},\apply{c}{\fst{w}}\}$} \}$.
        $S$ is a relation by \cref{relation,subseteq,times_elem_is_tuple}.
        For all $B\in F$ there exists $D\in F$ such that $D = B\union\{\apply{c}{B},\apply{c}{B}\}$
            by \cref{powerset_elim,upair_elim,subseteq,pair_is_finite,union_of_finite_is_finite,union_subsets_is_subset,powerset_intro}.
        Thus for all $B\in F$ there exists $D$ such that $B\mathrel{S} D$
            by \cref{times_tuple_intro,fst_eq,snd_eq}.
        Take $g$ such that $g$ is a function on $F$ and for all $B\in F$ we have $B\mathrel{S}\apply{g}{B}$
            by \cref{choice_function}.
        For all $B\in F$ we have $g(B)\in F$ and $g(B) = B\union\{\apply{c}{B},\apply{c}{B}\}$
            by \cref{times_tuple_elim,fst_eq,snd_eq}.
        Hence $g\in\funs{F}{F}$ by \cref{subseteq,funs_intro}.
        Take $x_0$ such that $x_0\in X$ by \cref{nonempty_inhabited}.
        Let $a_0 = \{x_0,x_0\}$.
        $a_0\in F$ by \cref{pair_is_finite,upair_elim,subseteq,powerset_intro}.
        Take $u$ such that
            $u$ is a sequence in $F$ and
            $u(1) = a_0$ and
            for all $n\in\naturalsPlus$ we have $u(n + 1) = \apply{g}{u(n)}$ by \cref{iteration_sequence_naturalsplus}.
        Let $s = c\circ u$.
        $u\in\funs{\naturalsPlus}{F}$ and $c\in\funs{F}{X}$ by \cref{sequence_in,funs}.
        Hence $s$ is a sequence in $X$ by \cref{funs_circ,sequence_in}.
        Therefore $\ran{u}\subseteq\dom{c}$ by \cref{funs_ran,funs,subseteq}.
        We have $u$ is right-unique and $u$ is a relation and $c$ is right-unique and $c$ is a relation
            by \cref{funs_is_function}.
        Show that for all $n\in\naturalsPlus$ we have
            for all $y\in u(n)$ we have $s(n)\notin\metricBall{d}{X}{y}{\epsilon}$.
        \begin{subproof}
            Fix $n\in\naturalsPlus$.
            $u(n)\in F$ by \cref{funs_type_apply,sequence_in}.
            Hence $u(n)\mathrel{R} c(u(n))$.
            $n\in\dom{u}$ by \cref{funs,subseteq}.
            Therefore for all $y\in u(n)$ we have $c(u(n))\notin\metricBall{d}{X}{y}{\epsilon}$ and $s(n) = c(u(n))$
                by \cref{circ_apply}.
            Hence for all $y\in u(n)$ we have $s(n)\notin\metricBall{d}{X}{y}{\epsilon}$.
        \end{subproof}
        For all $n\in\naturalsPlus$ we have $u(n + 1) = u(n)\union\{c(u(n)),c(u(n))\}$
            by \cref{iteration_sequence_naturalsplus,sequence_in,funs_type_apply,funs_is_function,function_apply_intro}.
        Thus for all $n\in\naturalsPlus$ we have $u(n)\subseteq u(n + 1)$ by \cref{union_upper_left}.
        For all $n\in\naturalsPlus$ we have $s(n)\in u(n + 1)$
            by \cref{iteration_sequence_naturalsplus,sequence_in,funs_type_apply,funs_is_function,function_apply_intro,funs,subseteq,circ_apply,upair_intro_left,union_iff}.
            Let $P = \{n\in\naturalsPlus \mid \text{for all $m\in\naturalsPlus$ such that $m < n$ we have $s(m)\in u(n)$}\}$.
                $P\subseteq\naturalsPlus$ by \cref{subseteq}.
                Show that $1\in P$.
                \begin{subproof}
                    Follows by \cref{real_one_in_naturals,real_one_leq_naturalsplus,real_naturals_subset_reals,subseteq,real_lt_subst_right,real_lt_trans,real_lt_irreflexive}.
                \end{subproof}
                Show that for all $n\in P$ we have $n + 1\in P$.
                \begin{subproof}
                    Fix $n$ such that $n\in P$.
                    Show that for all $m\in\naturalsPlus$ such that $m < n + 1$ we have $s(m)\in u(n + 1)$.
                    \begin{subproof}
                        Fix $m$ such that $m\in\naturalsPlus$ and $m < n + 1$.
                        $m \neq n + 1$.
                        Hence $m \leq n$ by \cref{real_naturals_leq_succ_elim}.
                        $m < n$ or $m = n$ by \cref{real_leq_split}.
                        \begin{byCase}
                        \caseOf{$m < n$.}
                            $s(m)\in u(n)$.
                            Hence $s(m)\in u(n + 1)$ by \cref{union_upper_left,subseteq}.
                        \caseOf{$m = n$.}
                            $s(n)\in u(n + 1)$.
                        \end{byCase}
                    \end{subproof}
                    Hence $n + 1\in P$.
                \end{subproof}
                Therefore for all $n\in\naturalsPlus$ we have $n\in P$ by \cref{real_naturals_induction}.
            For all $n\in\naturalsPlus$ we have
                for all $m\in\naturalsPlus$ such that $m < n$ we have $s(m)\in u(n)$ by \cref{subseteq}.
        For all $m\in\naturalsPlus$ we have
            for all $n\in\naturalsPlus$ such that $m < n$ we have $s(m)\in u(n)$
            by \cref{subseteq}.
        Therefore for all $m\in\naturalsPlus$ we have
            for all $n\in\naturalsPlus$ such that $m < n$ we have
            $s(n)\notin\metricBall{d}{X}{s(m)}{\epsilon}$.
        Show that for all $m\in\naturalsPlus$ we have
            for all $n\in\naturalsPlus$ such that $m < n$ we have
            it is wrong that $d((s(m),s(n))) < \epsilon$.
        \begin{subproof}
            Follows by \cref{sequence_in,funs_type_apply,metric_ball}.
        \end{subproof}
        Take $t,x$ such that
            $t$ is a sequence in $X$ and
            $t$ is a subsequence of $s$ and
            $x\in X$ and
            $t$ converges to $x$ in $X$ with respect to $T$ by \cref{sequentially_compact}.
        Take $v$ such that
            $v$ is strictly increasing on $\naturalsPlus$ and
            for all $n\in\naturalsPlus$ we have $t(n) = \apply{s}{\apply{v}{n}}$ by \cref{subsequence}.
        Let $m_2 = 1 + 1$.
        Let $\delta = \epsilon / m_2$.
        $1\in\reals$ and $m_2\in\reals$ and $0\in\reals$ and $0 < 1$
            by \cref{real_naturals_subset_reals,real_one_in_naturals,subseteq,real_add_closed,real_add_ident,real_zero_lt_one}.
        Hence $0 < m_2$ by \cref{real_lt_add,real_add_ident,real_lt_subst_right,real_lt_trans}.
        $m_2\neq 0$ by \cref{real_pos_nonzero}.
        Thus $\inv{m_2}\in\reals$ and $0 < \inv{m_2}$ by \cref{real_mul_inverse,real_inv_pos}.
        $\delta = \epsilon \rmul \inv{m_2}$ by \cref{real_div}.
        Hence $0 < \delta$ by \cref{real_lt_mul_pos,real_mul_zero,real_lt_subst_right}.
        $\metricBall{d}{X}{x}{\delta}\in\pow{X}$ by \cref{metric_ball,subseteq,powerset_intro}.
        $x\in X$ and $\delta\in\reals$ and $0 < \delta$.
        Therefore $\metricBall{d}{X}{x}{\delta}\in\metricBasis{d}{X}$ by \cref{metric_basis}.
        Hence $\metricBall{d}{X}{x}{\delta}\in T$
            by \cref{metric_basis_is_basis,basis_elem_open_in_generated,metric_topology}.
        Hence $\metricBall{d}{X}{x}{\delta}$ is open in $X$ with respect to $T$
            by \cref{open_in,basis_topology_is_topology,metric_topology,metric_basis_is_basis,basis_for_topology}.
        Therefore $\metricBall{d}{X}{x}{\delta}$ is a neighborhood of $x$ in $X$ with respect to $T$
            by \cref{metric_ball,metric_zero_characterization,real_lt_subst_left,metric_on,neighborhood}.
        Take $N\in\naturalsPlus$ such that for all $n\in\naturalsPlus$ if $N\leq n$ then $t(n)\in\metricBall{d}{X}{x}{\delta}$ by \cref{converges_to}.
        $N + 1\in\naturalsPlus$ by \cref{real_naturals_closed_succ}.
        $N\leq N + 1$.
        Hence $t(N)\in\metricBall{d}{X}{x}{\delta}$ and $t(N + 1)\in\metricBall{d}{X}{x}{\delta}$ by \cref{converges_to}.
        Therefore $d((x,t(N))) < \delta$ and $d((x,t(N + 1))) < \delta$ by \cref{metric_ball}.
        $t(N)\in X$ and $t(N + 1)\in X$ by \cref{sequence_in,funs_type_apply}.
        Hence $d((t(N),x)) < \delta$ by \cref{metric_on,metric_symmetric,real_lt_subst_left}.
        $d$ satisfies the triangle inequality on $X$ by \cref{metric_on}.
        Therefore $d((t(N),t(N + 1))) \leq d((t(N),x)) + d((x,t(N + 1)))$ by \cref{metric_triangle_inequality}.
        $(t(N),x)\in X\times X$ and $(x,t(N + 1))\in X\times X$ and
            $d\in\funs{X\times X}{\reals}$ by \cref{times_tuple_intro,metric_on}.
        $d((t(N),x))\in\reals$ and $d((x,t(N + 1)))\in\reals$ by \cref{funs_type_apply}.
        $\delta\in\reals$.
        Hence $d((x,t(N + 1))) + d((t(N),x)) < \delta + d((t(N),x))$ by \cref{real_lt_add}.
        $d((x,t(N + 1))) + d((t(N),x)) = d((t(N),x)) + d((x,t(N + 1)))$ and
            $\delta + d((t(N),x)) = d((t(N),x)) + \delta$ by \cref{real_add_comm}.
        Therefore $d((t(N),x)) + d((x,t(N + 1))) < d((t(N),x)) + \delta$
            by \cref{real_lt_subst_left,real_lt_subst_right}.
        $d((t(N),x)) + \delta < \delta + \delta$ by \cref{real_lt_add}.
        Hence $d((t(N),x)) + d((x,t(N + 1))) < \delta + \delta$
            by \cref{real_add_closed,real_lt_trans}.
        Therefore $d((t(N),x)) + d((x,t(N + 1))) < \epsilon$
            by \cref{real_mul_assoc,real_mul_inverse,real_mul_comm,real_mul_ident,real_right_distrib,real_lt_subst_right}.
        Thus $d((t(N),t(N + 1))) < \epsilon$
            by \cref{real_leq_split,times_tuple_intro,metric_on,funs_type_apply,real_add_closed,real_lt_trans,real_lt_subst_left}.
        $v(N)\in\naturalsPlus$ and $v(N + 1)\in\naturalsPlus$ by \cref{strictly_increasing_on_set,funs_type_apply}.
        Hence $N < N + 1$ by \cref{real_naturals_lt_succ_local}.
        Therefore $v(N) < v(N + 1)$ by \cref{strictly_increasing_on_set}.
        $t(N) = \apply{s}{\apply{v}{N}}$ and $t(N + 1) = \apply{s}{\apply{v}{(N + 1)}}$.
        Hence it is wrong that $d((\apply{s}{\apply{v}{N}},\apply{s}{\apply{v}{(N + 1)}})) < \epsilon$.
        Therefore it is wrong that $d((t(N),t(N + 1))) < \epsilon$
            by \cref{pair_eq_iff,real_lt_subst_left}.
        Contradiction.
    \end{byCase}
\end{proof}

% from §28 Item 6: compactness, limit point compactness, and sequential compactness
\begin{theorem}\label{metrizable_compact_iff_limit_point_compact_iff_sequentially_compact}
    Assume $T$ is a topology on $X$.
    Suppose $X$ is metrizable with respect to $T$.
    Then $X$ is compact with respect to $T$ iff
        $X$ is limit point compact with respect to $T$ and
        $X$ is sequentially compact with respect to $T$.
\end{theorem}
\begin{proof}
    Take $d$ such that $d$ is a metric on $X$ and $T = \metricTopology{d}{X}$.

    Show that if $X$ is compact with respect to $T$ then $X$ is limit point compact with respect to $T$.
    \begin{subproof}
        Follows by \cref{compact_implies_limit_point_compact}.
    \end{subproof}

    Show that if $X$ is sequentially compact with respect to $T$ then $X$ is compact with respect to $T$.
    \begin{subproof}
        Assume $X$ is sequentially compact with respect to $T$.
        Show that for all $A$ if $A$ is an open covering of $X$ with respect to $T$ then
            there exists $B$ such that $B\subseteq A$ and $B$ is finite and $B$ is a covering of $X$.
        \begin{subproof}
            Fix $A$ such that $A$ is an open covering of $X$ with respect to $T$.
            \begin{byCase}
            \caseOf{$X = \emptyset$.}
                Let $B = \emptyset$.
                Show that there exists $C$ such that $C\subseteq A$ and $C$ is finite and $C$ is a covering of $X$.
                \begin{subproof}
                    Follows by \cref{emptyset_subseteq,emptyset_is_finite,unions_emptyset,covering}.
                \end{subproof}
            \caseOf{$X\neq\emptyset$.}
                Take $\delta\in\reals$ such that $0 < \delta$ and
                    for all $B$ such that $B\subseteq X$ and
                        for all $x,y\in B$ we have $d((x,y)) < \delta$
                    there exists $U\in A$ such that $B\subseteq U$
                    by \cref{sequentially_compact_lebesgue}.
                Let $m_2 = 1 + 1$.
                Let $\epsilon = \delta / m_2$.
                $1\in\reals$ and $0\in\reals$ and $0 < 1$
                    by \cref{real_mul_ident,real_add_ident,real_zero_lt_one}.
                Hence $m_2\in\reals$ by \cref{real_add_closed}.
                Therefore $0 < m_2$ by \cref{real_lt_add,real_add_ident,real_lt_subst_left,real_lt_trans}.
                $m_2\neq 0$ by \cref{real_pos_nonzero}.
                Hence $\inv{m_2}\in\reals$ by \cref{real_mul_inverse}.
                Therefore $0 < \inv{m_2}$ by \cref{real_inv_pos}.
                $\epsilon = \delta \rmul \inv{m_2}$ by \cref{real_div}.
                $\epsilon\in\reals$ by \cref{real_mul_closed}.
                Hence $0 < \epsilon$ by \cref{real_add_ident,real_lt_mul_pos,real_mul_zero,real_lt_subst_left}.
                Take $B$ such that
                    $B$ is finite and
                    $B$ is a covering of $X$ and
                    for all $U\in B$ there exists $x\in X$ such that $U = \metricBall{d}{X}{x}{\epsilon}$
                    by \cref{sequentially_compact_finite_ball_cover}.

                Show that for all $V\in B$ there exists $U\in A$ such that $V\subseteq U$.
                \begin{subproof}
                    Fix $V$ such that $V\in B$.
                    Take $x_0\in X$ such that $V = \metricBall{d}{X}{x_0}{\epsilon}$.
                    Show that for all $x,y\in V$ we have $d((x,y)) < \delta$.
                    \begin{subproof}
                        Fix $x$.
                        Fix $y$ such that $x\in V$ and $y\in V$.
                        $x\in X$ and $y\in X$ by \cref{metric_ball,subseteq}.
                        Hence $d((x_0,x)) < \epsilon$ and $d((x_0,y)) < \epsilon$ by \cref{metric_ball}.
                        Therefore $d((x,y)) \leq d((x,x_0)) + d((x_0,y))$ by \cref{metric_on,metric_triangle_inequality}.
                        $d((x,x_0)) < \epsilon$ by \cref{metric_on,metric_symmetric,real_lt_subst_left}.
                        Show that $d((x,y)) < \epsilon + \epsilon$.
                        \begin{subproof}
                            $(x,x_0)\in X\times X$ and $(x_0,y)\in X\times X$ and $(x,y)\in X\times X$
                                by \cref{times_tuple_intro}.
                            $d((x,x_0))\in\reals$ and $d((x_0,y))\in\reals$ and $d((x,y))\in\reals$
                                by \cref{metric_on,funs_type_apply}.
                            Hence $d((x,x_0)) + d((x_0,y)) < \epsilon + d((x_0,y))$ by \cref{real_lt_add}.
                            $d((x,x_0)) + d((x_0,y)) = d((x_0,y)) + d((x,x_0))$ and
                                $\epsilon + d((x_0,y)) = d((x_0,y)) + \epsilon$ by \cref{real_add_comm}.
                            Therefore $d((x,x_0)) + d((x_0,y)) < d((x_0,y)) + \epsilon$
                                by \cref{real_lt_subst_left,real_lt_subst_right}.
                            $d((x_0,y)) + \epsilon < \epsilon + \epsilon$ by \cref{metric_ball,real_lt_add}.
                            Hence $d((x,x_0)) + d((x_0,y))\in\reals$ and $d((x_0,y)) + \epsilon\in\reals$ and
                                $\epsilon + \epsilon\in\reals$ by \cref{real_add_closed}.
                            Therefore $d((x,x_0)) + d((x_0,y)) < \epsilon + \epsilon$ by \cref{real_lt_trans}.
                            $d((x,y)) < \epsilon + \epsilon$
                                by \cref{real_leq_split,real_lt_trans,real_lt_subst_left}.
                        \end{subproof}
                        Hence $d((x,y)) < \delta$
                            by \cref{real_div,real_mul_assoc,real_mul_inverse,real_mul_comm,real_mul_ident,real_right_distrib,real_lt_subst_right}.
                    \end{subproof}
                    $V\subseteq X$ by \cref{metric_ball,subseteq}.
                    Hence there exists $U\in A$ such that $V\subseteq U$.
                \end{subproof}

                Let $R = \{w\in B\times A \mid \fst{w}\subseteq \snd{w}\}$.
                Show that for all $V\in B$ there exists $U$ such that $V\mathrel{R} U$.
                \begin{subproof}
                    Fix $V\in B$.
                    Take $U\in A$ such that $V\subseteq U$.
                    $(V,U)\in B\times A$ by \cref{times_tuple_intro}.
                    Hence $\fst{(V,U)}\subseteq\snd{(V,U)}$ by \cref{fst_eq,snd_eq}.
                    Therefore $V\mathrel{R} U$.
                \end{subproof}
                Take $h$ such that $h$ is a function on $B$ and for all $V\in B$ we have $V\mathrel{R}\apply{h}{V}$
                    by \cref{subseteq,times_elem_is_tuple,relation,choice_function}.
                Let $C = \img{h}{B}$.
                Hence $C$ is finite by \cref{finite_image_function}.
                Therefore $C\subseteq A$ by \cref{img_iff,function_apply_intro,subseteq,times_tuple_elim}.
                For all $x$ such that $x\in X$ we have $x\in\unions{C}$
                    by \cref{covering,subseteq,unions_iff,fst_eq,snd_eq,function_apply_elim,img_iff}.
                Hence $C$ is a covering of $X$ by \cref{subseteq,covering}.
                Show that there exists $D$ such that $D\subseteq A$ and $D$ is finite and $D$ is a covering of $X$.
                \begin{subproof}
                    Trivial.
                \end{subproof}
            \end{byCase}
        \end{subproof}
        Hence $X$ is compact with respect to $T$ by \cref{compact}.
    \end{subproof}

    Show that if $X$ is limit point compact with respect to $T$ then $X$ is sequentially compact with respect to $T$.
    \begin{subproof}
        Assume $X$ is limit point compact with respect to $T$.
        Show that for all $s$ such that $s$ is a sequence in $X$
            there exist $t,x$ such that
                $t$ is a sequence in $X$ and
                $t$ is a subsequence of $s$ and
                $x\in X$ and
                $t$ converges to $x$ in $X$ with respect to $T$.
        \begin{subproof}
            Fix $s$ such that $s$ is a sequence in $X$.
            $s\in\funs{\naturalsPlus}{X}$ by \cref{sequence_in}.
            Let $A = \ran{s}$.
            $A\subseteq X$ by \cref{funs_ran}.
            \begin{byCase}
            \caseOf{$A$ is finite.}
                For all $n\in\dom{s}$ we have $s(n)\in A$
                    by \cref{funs_is_function,function_apply_elim,ran_iff}.
                $s\in\funs{\naturalsPlus}{A}$.
                Hence $s$ is a sequence in $A$ by \cref{sequence_in}.
                Take $t,a$ such that
                    $a\in A$ and
                    $t$ is a sequence in $A$ and
                    $t$ is a subsequence of $s$ and
                    for all $n\in\naturalsPlus$ we have $t(n) = a$
                    by \cref{finite_sequence_constant_subsequence}.
                $a\in X$ by \cref{subseteq}.
                For all $n\in\dom{t}$ we have $t(n)\in X$
                    by \cref{sequence_in,funs,funs_type_apply,subseteq}.
                $t\in\funs{\naturalsPlus}{X}$.
                Hence $t$ is a sequence in $X$ by \cref{sequence_in}.
                Hence $t$ converges to $a$ in $X$ with respect to $T$
                    by \cref{neighborhood,real_one_in_naturals,subseteq,converges_to}.
                Show that there exist $u,b$ such that
                    $u$ is a sequence in $X$ and
                    $u$ is a subsequence of $s$ and
                    $b\in X$ and
                    $u$ converges to $b$ in $X$ with respect to $T$.
                \begin{subproof}
                    Trivial.
                \end{subproof}
            \caseOf{$A$ is infinite.}
                Take $x\in X$ such that $x$ is a limit point of $A$ in $X$ with respect to $T$
                    by \cref{limit_point_compact}.

                Hence $X$ satisfies the T-one axiom with respect to $T$
                    by \cref{metrizable,metric_space_hausdorff,finite_point_set_closed_hausdorff,t1_axiom}.

                For all $U$ such that $U$ is a neighborhood of $x$ in $X$ with respect to $T$
                    we have $U\inter A$ is infinite by \cref{limit_point_infinite_neighborhood}.

                For all $n\in\naturalsPlus$ we have
                    $\metricBall{d}{X}{x}{1 / n}$ is a neighborhood of $x$ in $X$ with respect to $T$
                    by \cref{positive_reciprocal_naturalsplus,real_naturals_subset_reals,subseteq,real_one_in_naturals,real_add_ident,real_zero_lt_one,real_one_leq_naturalsplus,real_lt_subst_right,real_lt_trans,real_pos_nonzero,real_mul_inverse,real_div,real_mul_closed,metric_ball,powerset_intro,metric_basis,metric_basis_is_basis,basis_elem_open_in_generated,metric_topology,basis_for_topology,basis_topology_is_topology,open_in,metric_zero_characterization,metric_on,real_lt_subst_left,neighborhood}.

                Let $D = \{w\in\naturalsPlus\times\naturalsPlus \mid
                    \apply{s}{\snd{w}}\in\metricBall{d}{X}{x}{1 / (\fst{w})}\}$.

                Let $U_1 = \metricBall{d}{X}{x}{1 / 1}$.
                Hence $U_1\inter A$ is infinite by \cref{real_one_in_naturals,subseteq}.
                Take $y$ such that $y\in U_1\inter A$ by \cref{emptyset_is_finite,nonempty_inhabited}.
                Take $k$ such that $(k,y)\in s$ by \cref{inter_elim_right,ran_iff}.
                $k\in\naturalsPlus$ and $s(k) = y$
                    by \cref{dom,funs,subseteq,sequence_in,function_apply_iff,funs_is_function}.
                Let $a_0 = (1,k)$.
                Hence $a_0\in D$ and $\fst{a_0} = 1$
                    by \cref{real_one_in_naturals,times_tuple_intro,inter_elim_left,fst_eq,snd_eq}.

                Let $R = \{w\in D\times D \mid
                    \fst{\snd{w}} = \fst{\fst{w}} + 1 \land \snd{\fst{w}} < \snd{\snd{w}}\}$.

                Show that for all $p\in D$ there exists $q$ such that $p\mathrel{R} q$.
                \begin{subproof}
                    Fix $p\in D$.
                    Take $n_0,m_0$ such that $p = (n_0,m_0)$ by \cref{times_elem_is_tuple}.
                    $n_0\in\naturalsPlus$ and $m_0\in\naturalsPlus$ by \cref{times_tuple_elim}.
                    Let $V = \metricBall{d}{X}{x}{1 / (n_0 + 1)}$.
                    Hence $V$ is a neighborhood of $x$ in $X$ with respect to $T$ by \cref{real_naturals_closed_succ,subseteq}.
                    Hence $V\inter A$ is infinite.
                    Let $J = \{m\in\naturalsPlus \mid \apply{s}{m}\in V\}$.
                    Show that $J$ is infinite.
                    \begin{subproof}
                        Suppose not.
                        $J$ is finite.
                        For all $y$ such that $y\in\img{s}{J}$ we have $y\in V\inter A$
                            by \cref{img_iff,function_apply_intro,funs_is_function,subseteq,ran_iff,inter_intro}.
                        Hence $\img{s}{J}\subseteq V\inter A$ by \cref{subseteq}.
                        For all $y$ such that $y\in V\inter A$ we have $y\in\img{s}{J}$
                            by \cref{inter_elim_right,ran_iff,dom,funs,subseteq,sequence_in,function_apply_intro,funs_is_function,inter_elim_left,img_iff}.
                        Therefore $V\inter A\subseteq \img{s}{J}$ by \cref{subseteq}.
                        $\img{s}{J} = V\inter A$ by \cref{subseteq_antisymmetric}.
                        Let $g = \restrl{s}{J}$.
                        Hence $g$ is a function on $J$
                            by \cref{sequence_in,funs_is_function,restrl_of_function_is_function,dom_of_restrl_eq_inter,funs_elim,subseteq,inter_absorb_supseteq_left}.
                        Therefore $\img{g}{J} = V\inter A$
                            by \cref{sequence_in,funs_elim,subseteq,restrl_img,subseteq_refl,inter_absorb_supseteq_right}.
                        $\img{g}{J}$ is finite by \cref{finite_image_function}.
                        Hence $V\inter A$ is finite.
                        Contradiction.
                    \end{subproof}
                    Show that there exists $k_1\in J$ such that $m_0 < k_1$.
                    \begin{subproof}
                        Suppose not.
                        For all $m\in J$ we have $m \leq m_0$.
                        Hence $J$ is finite
                            by \cref{initial_segment_naturalsplus,subseteq,initial_segment_finite,subseteq_finite_is_finite}.
                        Contradiction.
                    \end{subproof}
                    Take $k_1\in J$ such that $m_0 < k_1$.
                    $\apply{s}{k_1}\in V$.
                    Let $q = (n_0 + 1,k_1)$.
                    Hence $q\in\naturalsPlus\times\naturalsPlus$ by \cref{real_naturals_closed_succ,subseteq,times_tuple_intro}.
                    Therefore $\fst{q} = n_0 + 1$ and $\snd{q} = k_1$ by \cref{fst_eq,snd_eq}.
                    $\apply{s}{\snd{q}}\in\metricBall{d}{X}{x}{1 / (n_0 + 1)}$ by \cref{real_lt_subst_left}.
                    Hence $\apply{s}{\snd{q}}\in\metricBall{d}{X}{x}{1 / (\fst{q})}$ by \cref{real_lt_subst_right}.
                    Therefore $q\in D$ by \cref{subseteq}.
                    $k_1\in\naturalsPlus$ by \cref{subseteq}.
                    Hence $n_0 + 1\in\naturalsPlus$ by \cref{real_naturals_closed_succ}.
                    Therefore $(p,q) = ((n_0,m_0),(n_0 + 1,k_1))$ by \cref{real_naturals_closed_succ,pair_eq_iff}.
                    $(p,q)\in D\times D$ by \cref{times_tuple_intro}.
                    Hence $\fst{\snd{(p,q)}} = n_0 + 1$ and $\fst{\fst{(p,q)}} = n_0$
                        by \cref{fst_eq,snd_eq}.
                    Therefore $\snd{\fst{(p,q)}} = m_0$ and $\snd{\snd{(p,q)}} = k_1$
                        by \cref{fst_eq,snd_eq}.
                    $\fst{\snd{(p,q)}} = \fst{\fst{(p,q)}} + 1$ and
                        $\snd{\fst{(p,q)}} < \snd{\snd{(p,q)}}$
                        by \cref{real_lt_subst_left,real_lt_subst_right,real_add_eq_subst}.
                    Hence $(p,q)\in R$.
                    Therefore $p\mathrel{R} q$.
                    There exists $r$ such that $p\mathrel{R} r$.
                \end{subproof}

                Take $f$ such that $f$ is a function on $D$ and for all $p\in D$ we have
                    $p\mathrel{R}\apply{f}{p}$ by \cref{subseteq,times_elem_is_tuple,relation,choice_function}.
                Hence $f\in\funs{D}{D}$ by \cref{funs_intro,subseteq,times_tuple_elim}.
                Take $v$ such that
                    $v$ is a sequence in $D$ and
                    $v(1) = a_0$ and
                    for all $n\in\naturalsPlus$ we have $v(n + 1) = \apply{f}{v(n)}$
                    by \cref{iteration_sequence_naturalsplus}.

                Let $g = \{(p,\snd{p}) \mid p\in D\}$.
                For all $p,y_1,y_2$ such that $(p,y_1)\in g$ and $(p,y_2)\in g$ we have $y_1 = y_2$.
                Therefore $g$ is a function by \cref{rightunique,relation}.
                $\dom{g} = D$ by \cref{dom_iff,pair_eq_iff,subseteq,subseteq_antisymmetric}.
                For all $p\in\dom{g}$ we have $g(p)\in\naturalsPlus$
                    by \cref{pair_eq_iff,function_apply_intro,subseteq,times_elem_is_tuple,snd_eq}.
                Hence $g\in\funs{D}{\naturalsPlus}$ by \cref{funs_intro}.

                Let $u = g\circ v$.
                Hence $u$ is a sequence in $\naturalsPlus$ by \cref{sequence_in,funs_circ}.

                For all $n\in\naturalsPlus$ we have $u(n) = \snd{v(n)}$
                    by \cref{funs,sequence_in,funs_is_function,function_apply_elim,funs_type_apply,circ_iff,function_apply_intro}.

                Let $P = \{n\in\naturalsPlus \mid \fst{v(n)} = n\}$.
                Show that for all $n\in\naturalsPlus$ we have $n\in P$.
                \begin{subproof}
                    $P\subseteq\naturalsPlus$ by \cref{subseteq}.
                    Show that $1\in P$.
                    \begin{subproof}
                        Follows by \cref{real_one_in_naturals,fst_eq,pair_eq_pair_of_proj}.
                    \end{subproof}
                    Show that for all $n\in P$ we have $n + 1\in P$.
                    \begin{subproof}
                        Fix $n$ such that $n\in P$.
                        $\fst{v(n)} = n$.
                        Let $p = v(n)$.
                        Let $q = \apply{f}{p}$.
                        Hence $v(n + 1) = q$ by \cref{real_naturals_closed_succ,iteration_sequence_naturalsplus}.
                        $p\in D$ by \cref{sequence_in,funs_type_apply}.
                        Hence $q\in D$ by \cref{funs_type_apply}.
                        Take $n_0,k_0$ such that $p = (n_0,k_0)$ by \cref{times_elem_is_tuple}.
                        Take $n_1,k_1$ such that $q = (n_1,k_1)$ by \cref{times_elem_is_tuple}.
                        $n_0\in\naturalsPlus$ and $k_0\in\naturalsPlus$ and
                            $n_1\in\naturalsPlus$ and $k_1\in\naturalsPlus$ by \cref{times_tuple_elim}.
                        Hence $p\mathrel{R} q$.
                        Therefore $\fst{q} = \fst{p} + 1$ and $\snd{p} < \snd{q}$ by \cref{fst_eq,snd_eq}.
                        $(v(n),v(n + 1)) = ((n_0,k_0),(n_1,k_1))$
                            by \cref{pair_eq_iff}.
                        Hence $n_1 = n_0 + 1$ and $k_0 < k_1$
                            by \cref{fst_eq,snd_eq,real_lt_subst_left,real_lt_subst_right,real_add_eq_subst}.
                        Therefore $\fst{v(n + 1)} = n + 1$
                            by \cref{pair_eq_iff,fst_eq,real_add_comm,real_add_ident}.
                        Hence $n + 1\in P$.
                    \end{subproof}
                    Therefore for all $n\in\naturalsPlus$ we have $n\in P$ by \cref{real_naturals_induction}.
                \end{subproof}
                For all $n\in\naturalsPlus$ we have $\apply{s}{u(n)}\in\metricBall{d}{X}{x}{1 / n}$
                    by \cref{funs_type_apply,sequence_in,snd_eq,subseteq}.

                Show that for all $n\in\naturalsPlus$ we have $u(n) < u(n + 1)$.
                \begin{subproof}
                    Fix $n\in\naturalsPlus$.
                    Let $p = v(n)$.
                    Let $q = \apply{f}{p}$.
                    Hence $v(n + 1) = q$ by \cref{iteration_sequence_naturalsplus}.
                    $p\in D$ by \cref{sequence_in,funs_type_apply}.
                    Hence $q\in D$ by \cref{funs_type_apply}.
                    Take $n_0,k_0$ such that $p = (n_0,k_0)$ by \cref{times_elem_is_tuple}.
                    Take $n_1,k_1$ such that $q = (n_1,k_1)$ by \cref{times_elem_is_tuple}.
                    $n_0\in\naturalsPlus$ and $k_0\in\naturalsPlus$ and
                        $n_1\in\naturalsPlus$ and $k_1\in\naturalsPlus$ by \cref{times_tuple_elim}.
                    Hence $p\mathrel{R} q$.
                    Therefore $\fst{q} = \fst{p} + 1$ and $\snd{p} < \snd{q}$ by \cref{fst_eq,snd_eq}.
                    $(v(n),v(n + 1)) = ((n_0,k_0),(n_1,k_1))$
                        by \cref{pair_eq_iff}.
                    Hence $n_1 = n_0 + 1$ and $k_0 < k_1$
                        by \cref{fst_eq,snd_eq,real_lt_subst_left,real_lt_subst_right,real_add_eq_subst}.
                    Therefore $u(n) < u(n + 1)$
                        by \cref{pair_eq_iff,snd_eq,real_naturals_closed_succ,real_lt_subst_left,real_lt_subst_right}.
                \end{subproof}

                Show that for all $m,n\in\naturalsPlus$ such that $m < n$ we have $u(m) < u(n)$.
                \begin{subproof}
                    Let $T_0 = \{n\in\naturalsPlus \mid \text{for all $m\in\naturalsPlus$ such that $m < n$ we have $u(m) < u(n)$}\}$.
                    Show that for all $n\in\naturalsPlus$ we have $n\in T_0$.
                    \begin{subproof}
                        $T_0\subseteq\naturalsPlus$ by \cref{subseteq}.
                        Show that $1\in T_0$.
                        \begin{subproof}
                            Follows by \cref{real_one_in_naturals,real_one_leq_naturalsplus,real_naturals_subset_reals,subseteq,real_lt_subst_right,real_lt_trans,real_lt_irreflexive}.
                        \end{subproof}
                        Show that for all $n\in T_0$ we have $n + 1\in T_0$.
                        \begin{subproof}
                            Fix $n$ such that $n\in T_0$.
                            Show that for all $m\in\naturalsPlus$ such that $m < n + 1$ we have $u(m) < u(n + 1)$.
                            \begin{subproof}
                                Fix $m$ such that $m\in\naturalsPlus$ and $m < n + 1$.
                                Suppose not.
                                Thus $m < n$ or $m = n$ or $m = n + 1$ by \cref{real_succ_discrete}.
                                \begin{byCase}
                                \caseOf{$m = n$.}
                                    Contradiction by \cref{real_lt_subst_left}.
                                \caseOf{$m < n$.}
                                    $u(m) < u(n)$ by \cref{subseteq}.
                                    Hence $u(n) < u(n + 1)$.
                                    Therefore $u(m) < u(n + 1)$
                                        by \cref{sequence_in,real_naturals_closed_succ,funs_type_apply,real_naturals_subset_reals,subseteq,real_lt_trans}.
                                    Contradiction.
                                \caseOf{$m = n + 1$.}
                                    Contradiction
                                        by \cref{real_lt_subst_left,real_naturals_subset_reals,subseteq,real_lt_irreflexive}.
                                \end{byCase}
                            \end{subproof}
                            Hence $n + 1\in T_0$.
                        \end{subproof}
                        Therefore for all $n\in\naturalsPlus$ we have $n\in T_0$ by \cref{real_naturals_induction}.
                    \end{subproof}
                    Thus for all $m,n\in\naturalsPlus$ such that $m < n$ we have $u(m) < u(n)$ by \cref{subseteq}.
                \end{subproof}
                Hence $u$ is strictly increasing on $\naturalsPlus$ by \cref{sequence_in,strictly_increasing_on_set}.

                Let $t = s\circ u$.
                Hence $t$ is a sequence in $X$ by \cref{sequence_in,funs_circ}.

                For all $n\in\naturalsPlus$ we have $t(n) = \apply{s}{\apply{u}{n}}$
                    by \cref{funs,sequence_in,funs_is_function,function_apply_elim,funs_type_apply,circ_iff,function_apply_intro}.
                Hence $t$ is a subsequence of $s$ by \cref{subsequence}.

                Show that for all $U$ if $U$ is a neighborhood of $x$ in $X$ with respect to $T$ then
                    there exists $N\in\naturalsPlus$ such that for all $n\in\naturalsPlus$ if $N\leq n$ then $t(n)\in U$.
                \begin{subproof}
                    Fix $U$ such that $U$ is a neighborhood of $x$ in $X$ with respect to $T$.
                    $U$ is open in $X$ with respect to $T$ and $x\in U$ by \cref{neighborhood}.
                    Take $\epsilon\in\reals$ such that $0 < \epsilon$ and $\metricBall{d}{X}{x}{\epsilon}\subseteq U$
                        by \cref{open_in,topology_on,subseteq,powerset_elim,metric_open_iff}.
                    Take $N\in\naturalsPlus$ such that $0 < 1 / N$ and $1 / N < \epsilon$ by \cref{real_small_reciprocal}.

                        Show that for all $n\in\naturalsPlus$ if $N\leq n$ then $t(n)\in U$.
                        \begin{subproof}
                            Follows by \cref{naturalsplus_reciprocal_antitone,real_leq_split,real_naturals_subset_reals,real_one_in_naturals,subseteq,real_add_ident,real_zero_lt_one,real_one_leq_naturalsplus,real_lt_subst_right,real_lt_trans,real_pos_nonzero,real_mul_inverse,real_div,real_mul_closed,metric_ball_mono,subseteq_transitive,funs_type_apply,sequence_in,snd_eq}.
                        \end{subproof}
                        Hence there exists $M\in\naturalsPlus$ such that for all $n\in\naturalsPlus$ if $M\leq n$ then $t(n)\in U$.
                \end{subproof}
                Hence $t$ converges to $x$ in $X$ with respect to $T$ by \cref{converges_to}.
                Show that there exist $r,b$ such that
                    $r$ is a sequence in $X$ and
                    $r$ is a subsequence of $s$ and
                    $b\in X$ and
                    $r$ converges to $b$ in $X$ with respect to $T$.
                \begin{subproof}
                    Trivial.
                \end{subproof}
            \end{byCase}
        \end{subproof}
        Hence $X$ is sequentially compact with respect to $T$ by \cref{sequentially_compact}.
    \end{subproof}

    Show that if $X$ is compact with respect to $T$ then
        $X$ is limit point compact with respect to $T$ and
        $X$ is sequentially compact with respect to $T$.
    \begin{subproof}
        Trivial.
    \end{subproof}

    Show that if $X$ is limit point compact with respect to $T$ and $X$ is sequentially compact with respect to $T$ then $X$ is compact with respect to $T$.
    \begin{subproof}
        Trivial.
    \end{subproof}
\end{proof}
    
\section{§29 Local Compactness}

% from §29 helper definition 1: compact subspace
\begin{definition}\label{compact_subspace}
    $C$ is a compact subspace of $X$ with respect to $T$ iff
        $C\subseteq X$ and
        $C$ is compact with respect to $\subspaceTopology{T}{C}$.
\end{definition}

% from §29 Item 1: locally compact at a point
\begin{definition}\label{locally_compact_at_point}
    $X$ is locally compact at $x$ with respect to $T$ iff
        $x\in X$ and
        there exists $C$ such that
            $C$ is a compact subspace of $X$ with respect to $T$ and
            there exists $U$ such that
                $U$ is a neighborhood of $x$ in $X$ with respect to $T$ and
                $U\subseteq C$.
\end{definition}

% from §29 Item 2: locally compact
\begin{definition}\label{locally_compact}
    $X$ is locally compact with respect to $T$ iff
        for all $x\in X$ we have $X$ is locally compact at $x$ with respect to $T$.
\end{definition}

% from §29 Item 3: compact implies locally compact
\begin{lemma}\label{compact_implies_locally_compact}
    Assume $T$ is a topology on $X$.
    Suppose $X$ is compact with respect to $T$.
    Then $X$ is locally compact with respect to $T$.
\end{lemma}
\begin{proof}
    Let $T_X = \subspaceTopology{T}{X}$.
    Show that $T\subseteq T_X$.
    \begin{subproof}
        Follows by \cref{topology_on,powerset_elim,subseteq,inter_absorb_supseteq_left,subspace_topology}.
    \end{subproof}
    Show that $T_X\subseteq T$.
    \begin{subproof}
        Follows by \cref{subspace_topology,topology_on,powerset_elim,subseteq,inter_absorb_supseteq_left}.
    \end{subproof}
    Therefore $T_X = T$ by \cref{subseteq,subseteq_antisymmetric}.
    Hence $X$ is a compact subspace of $X$ with respect to $T$ by \cref{subseteq_refl,compact_subspace}.
    Show that for all $x\in X$ we have $X$ is locally compact at $x$ with respect to $T$.
    \begin{subproof}
        Follows by \cref{subseteq_refl,topology_on,open_in,neighborhood,locally_compact_at_point}.
    \end{subproof}
    Hence $X$ is locally compact with respect to $T$ by \cref{locally_compact}.
\end{proof}

% from §29 helper lemma 1: set not member of itself
\begin{lemma}\label{set_notin_self}
    For all $x$ we have $x\notin x$.
\end{lemma}
\begin{proof}
    Fix $x$.
    Suppose not.
    Take $a\in\{x\}$ such that for all $y\in a$ we have $y\notin\{x\}$ by \cref{singleton_intro,emptyset,foundation}.
    Hence $x\notin\{x\}$ by \cref{singleton_elim}.
    Contradiction by \cref{singleton_intro}.
\end{proof}

% from §29 helper lemma 2: empty set is compact
\begin{lemma}\label{emptyset_compact}
    Assume $T$ is a topology on $X$.
    Then $\emptyset$ is compact with respect to $T$.
\end{lemma}
\begin{proof}
    Show that for all $A$ such that $A$ is an open covering of $\emptyset$ with respect to $T$
        there exists $B$ such that $B\subseteq A$ and $B$ is finite and $B$ is a covering of $\emptyset$.
    \begin{subproof}
        Fix $A$ such that $A$ is an open covering of $\emptyset$ with respect to $T$.
        Let $B = \emptyset$.
        We have $B\subseteq A$ and $B$ is finite and $B$ is a covering of $\emptyset$
            by \cref{emptyset_subseteq,emptyset_is_finite,unions_emptyset,covering}.
    \end{subproof}
    Hence $\emptyset$ is compact with respect to $T$ by \cref{compact}.
\end{proof}

% from §29 helper lemma 3: intersection of compact subspaces is compact
\begin{lemma}\label{compact_subspace_inters}
    Assume $T$ is a topology on $X$.
    Assume $X$ is Hausdorff with respect to $T$.
    Suppose $C\subseteq\pow{X}$.
    Suppose $C$ is inhabited.
    Suppose for all $K$ such that $K\in C$ we have $K$ is a compact subspace of $X$ with respect to $T$.
    Then $\inters{C}$ is a compact subspace of $X$ with respect to $T$.
\end{lemma}
\begin{proof}
    Take $K_0$ such that $K_0\in C$.
    For all $K$ such that $K\in C$ we have $K$ is closed in $X$ with respect to $T$
        by \cref{compact_subspace,compact_subspace_closed}.
    Hence $\inters{C}$ is closed in $X$ with respect to $T$ by \cref{closed_inters}.
    $\inters{C}\subseteq K_0$ by \cref{inters,subseteq}.
    Let $T_{K0} = \subspaceTopology{T}{K_0}$.
    $K_0$ is compact with respect to $T_{K0}$ and $K_0\subseteq X$ by \cref{compact_subspace}.
    $\inters{C}\inter K_0 = \inters{C}$ by \cref{closed_in_subspace_iff,inter_absorb_supseteq_right,subseteq}.
    Hence $\inters{C}$ is compact with respect to $\subspaceTopology{T_{K0}}{\inters{C}}$
        by \cref{compact_subspace,subspace_of,subspace_topology_is_topology,closed_in_subspace_iff,closed_subspace_compact}.
    Therefore $\inters{C}$ is compact with respect to $\subspaceTopology{T}{\inters{C}}$
        by \cref{closed_in_subspace_iff,closed_subspace_compact,subspace_subspace_topology,subseteq}.
    Finally $\inters{C}$ is a compact subspace of $X$ with respect to $T$ by \cref{subseteq_transitive,compact_subspace}.
\end{proof}

% from §29 helper lemma 4: union of compact subspaces is compact
\begin{lemma}\label{compact_subspace_union}
    Assume $T$ is a topology on $X$.
    Suppose $A$ is a compact subspace of $X$ with respect to $T$.
    Suppose $B$ is a compact subspace of $X$ with respect to $T$.
    Let $C = A\union B$.
    Then $C$ is a compact subspace of $X$ with respect to $T$.
\end{lemma}
\begin{proof}
    We have $C\subseteq X$ by \cref{compact_subspace,union_subsets_is_subset}.
    Let $T_C = \subspaceTopology{T}{C}$.
    Show that for all $U$ such that $U$ is an open covering of $C$ with respect to $T_C$
        there exists $D$ such that
            $D\subseteq U$ and
            $D$ is finite and
            $D$ is a covering of $C$.
    \begin{subproof}
        Fix $U$ such that $U$ is an open covering of $C$ with respect to $T_C$.
        Let $V = \{W\in T \mid \text{there exists $U_0\in U$ such that $U_0 = C\inter W$}\}$.
        For all $a$ such that $a\in A$ we have $a\in\unions{V}$
            by \cref{union_intro_left,covering,elem_subseteq,unions_iff,open_covering,open_in,subspace_topology,inter_lower_right}.
        Hence $V$ is a covering of $A$ by \cref{subseteq,covering}.
        For all $V_0$ such that $V_0\in V$ we have $V_0$ is open in $X$ with respect to $T$ by \cref{open_in}.
        Let $T_A = \subspaceTopology{T}{A}$.
        $A$ is a subspace of $X$ with respect to $T$ and $A$ is compact with respect to $T_A$
            by \cref{subspace_of,compact_subspace}.
        Take $V_A$ such that
            $V_A\subseteq V$ and
            $V_A$ is finite and
            $V_A$ is a covering of $A$
            by \cref{compact_subspace_open_covering}.

        For all $b$ such that $b\in B$ we have $b\in\unions{V}$
            by \cref{union_intro_right,covering,elem_subseteq,unions_iff,open_covering,open_in,subspace_topology,inter_lower_right}.
        Hence $V$ is a covering of $B$ by \cref{subseteq,covering}.
        Let $T_B = \subspaceTopology{T}{B}$.
        $B$ is a subspace of $X$ with respect to $T$ and $B$ is compact with respect to $T_B$
            by \cref{subspace_of,compact_subspace}.
        Take $V_B$ such that
            $V_B\subseteq V$ and
            $V_B$ is finite and
            $V_B$ is a covering of $B$
            by \cref{compact_subspace_open_covering}.

        Let $V_0 = V_A\union V_B$.
        $V_0$ is finite by \cref{union_of_finite_is_finite}.
        For all $x$ such that $x\in C$ we have $x\in\unions{V_0}$
            by \cref{union_iff,covering,elem_subseteq,unions_iff,union_intro_left,union_intro_right}.
        Therefore $V_0$ is a covering of $C$ by \cref{subseteq,covering}.

        Let $g = \{(v, C\inter v) \mid v\in V_0\}$.
        For all $w$ such that $w\in g$ there exist $x,y$ such that $w = (x,y)$.
        For all $x,y,z$ such that $(x,y)\in g$ and $(x,z)\in g$ we have $y = z$
            by \cref{pair_eq_iff}.
        Hence $g$ is a function by \cref{rightunique,relation}.
        Therefore $\dom{g} = V_0$ by \cref{dom_iff,pair_eq_iff,subseteq,subseteq_antisymmetric}.
        For all $v$ such that $v\in\dom{g}$ we have $g(v)\in\pow{C}$
            by \cref{function_apply_iff,inter_lower_left,powerset_intro}.
        Hence $g\in\funs{V_0}{\pow{C}}$ by \cref{function_apply_iff,inter_lower_left,powerset_intro,funs_intro}.
        Let $E = \{w \mid \exists v\in V_0. w = C\inter v\}$.
        Hence $E$ is finite by \cref{funs_is_function,funs,finite_image_function,img_iff,subseteq,subseteq_finite_is_finite}.
        $E\subseteq U$ by \cref{union_subsets_is_subset,elem_subseteq,subseteq}.
        Hence $E$ is a covering of $C$ by \cref{covering,elem_subseteq,unions_iff,inter_intro,subseteq}.
        There exists $D$ such that
            $D\subseteq U$ and
            $D$ is finite and
            $D$ is a covering of $C$.
    \end{subproof}
    Finally $C$ is a compact subspace of $X$ with respect to $T$ by \cref{compact,compact_subspace}.
\end{proof}

% from §29 helper lemma 5: compact subspace minus open set is compact
\begin{lemma}\label{compact_subspace_minus_open}
    Assume $T$ is a topology on $X$.
    Suppose $C$ is a compact subspace of $X$ with respect to $T$.
    Suppose $U$ is open in $X$ with respect to $T$.
    Let $D = C\setminus U$.
    Then $D$ is a compact subspace of $X$ with respect to $T$.
\end{lemma}
\begin{proof}
    We have $D\subseteq X$ by \cref{compact_subspace,setminus_subseteq,subseteq_transitive}.
    Let $T_C = \subspaceTopology{T}{C}$.
    Let $T_D = \subspaceTopology{T}{D}$.
    $X\setminus U$ is closed in $X$ with respect to $T$
        by \cref{open_in,topology_on,subseteq,pow_iff,setminus_subseteq,setminus_setminus,inter_absorb_supseteq_left,closed_in}.
    Hence $D$ is closed in $C$ with respect to $T_C$
        by \cref{open_in,topology_on,subseteq,pow_iff,setminus_eq_inter_complement,inter_comm,closed_in_subspace_iff,subspace_of,compact_subspace}.
    Hence $D$ is compact with respect to $T_D$
        by \cref{compact_subspace,setminus_subseteq,subspace_topology_is_topology,closed_subspace_compact,subspace_subspace_topology}.
    Therefore $D$ is a compact subspace of $X$ with respect to $T$ by \cref{compact_subspace}.
\end{proof}

% from §29 helper definition 2: open sets from X
\begin{definition}\label{one_point_open_sets}
    $\onePointOpenSets{T}{X}{p}
        = \{U\in\pow{X\union\{p\}} \mid \text{$U\subseteq X$ and $U\in T$}\}$.
\end{definition}

% from §29 helper definition 3: complements of compact sets
\begin{definition}\label{one_point_compact_complements}
    $\onePointCompactComplements{T}{X}{p}
        = \{U\in\pow{X\union\{p\}} \mid \text{there exists $C$ such that
            $C$ is a compact subspace of $X$ with respect to $T$ and
            $U = (X\union\{p\})\setminus C$}\}$.
\end{definition}

% from §29 helper definition 4: one-point compactification topology
\begin{definition}\label{one_point_compactification_topology}
    $\onePointTopology{T}{X}{p}
        = \onePointOpenSets{T}{X}{p}\union \onePointCompactComplements{T}{X}{p}$.
\end{definition}

% from §29 helper lemma 6: one-point topology is contained in powerset
\begin{lemma}\label{one_point_topology_subseteq_pow}
    Assume $p\notin X$.
    Let $Y = X\union\{p\}$.
    Let $S = \onePointTopology{T}{X}{p}$.
    Then $S\subseteq\pow{Y}$.
\end{lemma}
\begin{proof}
    We have $S\subseteq\pow{Y}$
        by \cref{one_point_compactification_topology,union_iff,one_point_open_sets,one_point_compact_complements,subseteq}.
\end{proof}

% from §29 helper lemma 7: open set from X is open in one-point topology
\begin{lemma}\label{one_point_topology_intro_left}
    Assume $p\notin X$.
    Let $Y = X\union\{p\}$.
    Let $S = \onePointTopology{T}{X}{p}$.
    Suppose $U\subseteq X$.
    Suppose $U\in T$.
    Then $U\in S$.
\end{lemma}
\begin{proof}
    We have $U\in S$ by \cref{one_point_compactification_topology,one_point_open_sets,union_upper_left,subseteq_transitive,pow_iff,union_intro_left}.
\end{proof}

% from §29 helper lemma 8: complement of compact set is open in one-point topology
\begin{lemma}\label{one_point_topology_intro_right}
    Assume $p\notin X$.
    Let $Y = X\union\{p\}$.
    Let $S = \onePointTopology{T}{X}{p}$.
    Suppose there exists $C$ such that
        $C$ is a compact subspace of $X$ with respect to $T$ and
        $U = Y\setminus C$.
    Then $U\in S$.
\end{lemma}
\begin{proof}
    Take $C$ such that
        $C$ is a compact subspace of $X$ with respect to $T$ and
        $U = Y\setminus C$.
    Hence $U\in S$ by \cref{one_point_compactification_topology,one_point_compact_complements,setminus_subseteq,pow_iff,union_intro_right}.
\end{proof}

% from §29 helper lemma 9: open set without added point
\begin{lemma}\label{one_point_topology_without_p}
    Assume $p\notin X$.
    Let $Y = X\union\{p\}$.
    Let $S = \onePointTopology{T}{X}{p}$.
    Suppose $U\in S$.
    Suppose $p\notin U$.
    Then $U\subseteq X$ and $U\in T$.
\end{lemma}
\begin{proof}
    Show that $U\in \onePointOpenSets{T}{X}{p}$.
    \begin{subproof}
        Suppose not.
        Hence $U\in \onePointCompactComplements{T}{X}{p}$ by \cref{one_point_compactification_topology,union_iff}.
        Take $C$ such that
            $C$ is a compact subspace of $X$ with respect to $T$ and
            $U = Y\setminus C$
            by \cref{one_point_compact_complements}.
        Thus $p\in U$ by \cref{compact_subspace,subseteq,singleton_intro,union_intro_right,setminus_intro}.
        Contradiction.
    \end{subproof}
    Show that $U\subseteq X$ and $U\in T$.
    \begin{subproof}
        Follows by \cref{one_point_open_sets}.
    \end{subproof}
\end{proof}

% from §29 helper lemma 10: open set containing added point
\begin{lemma}\label{one_point_topology_with_p}
    Assume $p\notin X$.
    Let $Y = X\union\{p\}$.
    Let $S = \onePointTopology{T}{X}{p}$.
    Suppose $U\in S$.
    Suppose $p\in U$.
    Then there exists $C$ such that $C$ is a compact subspace of $X$ with respect to $T$
        and $U = (X\union\{p\})\setminus C$.
\end{lemma}
\begin{proof}
    Show that $U\in \onePointCompactComplements{T}{X}{p}$.
    \begin{subproof}
        Suppose not.
        Hence $U\in \onePointOpenSets{T}{X}{p}$ by \cref{one_point_compactification_topology,union_iff}.
        Therefore $p\in X$ by \cref{one_point_open_sets,subseteq}.
        Contradiction.
    \end{subproof}
    Show that there exists $C$ such that $C$ is a compact subspace of $X$ with respect to $T$
        and $U = Y\setminus C$.
    \begin{subproof}
        Follows by \cref{one_point_compact_complements}.
    \end{subproof}
\end{proof}

% from §29 helper lemma 11: one-point topology is a topology
\begin{lemma}\label{one_point_topology_is_topology}
    Assume $T$ is a topology on $X$.
    Assume $X$ is Hausdorff with respect to $T$.
    Assume $p\notin X$.
    Let $Y = X\union\{p\}$.
    Let $S = \onePointTopology{T}{X}{p}$.
    Then $S$ is a topology on $Y$.
\end{lemma}
\begin{proof}
    Show that $S\subseteq\pow{Y}$.
    \begin{subproof}
        Follows by \cref{one_point_topology_subseteq_pow}.
    \end{subproof}
    Show that $\emptyset\in S$.
    \begin{subproof}
        Follows by \cref{emptyset_subseteq,topology_on,one_point_topology_intro_left}.
    \end{subproof}
    Let $T_0 = \subspaceTopology{T}{\emptyset}$.
    We have $\emptyset$ is a compact subspace of $X$ with respect to $T$
        by \cref{emptyset_subseteq,subspace_topology_is_topology,emptyset_compact,compact_subspace}.
    Show that $Y\in S$.
    \begin{subproof}
        Follows by \cref{setminus_emptyset,one_point_topology_intro_right}.
    \end{subproof}
    Show that for all $A$ such that $A\subseteq S$ we have $\unions{A}\in S$.
    \begin{subproof}
        Fix $A$ such that $A\subseteq S$.
        \begin{byCase}
        \caseOf{$p\notin\unions{A}$.}
            We have $A\subseteq T$ by \cref{unions_iff,subseteq,one_point_topology_without_p}.
            $\unions{A}\subseteq X$
                by \cref{topology_on,subseteq_transitive,unions_subseteq_of_powerset_is_subseteq}.
            Hence $\unions{A}\in S$ by \cref{topology_on,one_point_topology_intro_left}.
        \caseOf{$p\in\unions{A}$.}
            Take $U_0\in A$ such that $p\in U_0$ by \cref{unions_iff}.
            Take $C_0$ such that
                $C_0$ is a compact subspace of $X$ with respect to $T$ and
                $U_0 = (X\union\{p\})\setminus C_0$
                by \cref{subseteq,one_point_topology_with_p}.
            Let $A_1 = \{U\in A \mid p\notin U\}$.
            We have $A_1\subseteq T$ by \cref{subseteq,one_point_topology_without_p}.
            Let $U = \unions{A_1}$.
            Therefore $U\subseteq X$ and $U\in T$ by \cref{topology_on,subseteq,pow_iff}.
            Let $A_2 = \{U\in A \mid p\in U\}$.
            Let $F = \{Y\setminus U \mid U\in A_2\}$.
            $U_0\in A_2$.
            Hence $Y\setminus U_0\in F$.
            $F$ is inhabited.
            Show that for all $C$ such that $C\in F$ we have $C$ is a compact subspace of $X$ with respect to $T$.
            \begin{subproof}
                Fix $C$ such that $C\in F$.
                Take $U_1\in A_2$ such that $C = Y\setminus U_1$.
                Hence $U_1\in A$ and $p\in U_1$.
                Take $K$ such that
                    $K$ is a compact subspace of $X$ with respect to $T$ and
                    $U_1 = (X\union\{p\})\setminus K$
                    by \cref{subseteq,one_point_topology_with_p}.
                We have $K\subseteq X$ by \cref{compact_subspace}.
                $K\subseteq Y$ by \cref{union_upper_left,subseteq_transitive}.
                Hence $C = K$ by \cref{setminus_setminus,inter_absorb_supseteq_left}.
                Therefore $C$ is a compact subspace of $X$ with respect to $T$.
            \end{subproof}
            Hence $F\subseteq\pow{X}$ by \cref{subseteq,compact_subspace,pow_iff}.
            Let $C = \inters{F}$.
            $C$ is a compact subspace of $X$ with respect to $T$ by \cref{compact_subspace_inters}.
            Show that $\unions{A_2}\subseteq (X\union\{p\})\setminus C$.
            \begin{subproof}
                Follows by \cref{unions_iff,one_point_topology_subseteq_pow,subseteq,pow_iff,setminus_setminus,inter_absorb_supseteq_left,setminus_elim_right,inters_iff_forall,setminus_elim_left,setminus_intro}.
            \end{subproof}
            Show that for all $x$ such that $x\in (X\union\{p\})\setminus C$ we have
                $x\in\unions{A_2}$.
            \begin{subproof}
                Fix $x$ such that $x\in (X\union\{p\})\setminus C$.
                Hence $x\notin C$ by \cref{setminus_elim_right}.
                Suppose not.
                For all $C_1$ such that $C_1\in F$ we have $x\in C_1$
                    by \cref{subseteq,one_point_topology_subseteq_pow,pow_iff,setminus_elim_left,setminus_intro,setminus_setminus,inter_absorb_supseteq_left,unions_iff}.
                Therefore $x\in C$ by \cref{inters_iff_forall}.
                Contradiction.
                Hence $x\in\unions{A_2}$.
            \end{subproof}
            Therefore $(X\union\{p\})\setminus C\subseteq \unions{A_2}$ by \cref{subseteq}.
            Hence $\unions{A_2} = (X\union\{p\})\setminus C$ by \cref{subseteq_antisymmetric}.
            For all $x$ such that $x\in\unions{A}$ we have $x\in\unions{A_1}\union\unions{A_2}$.
            Therefore $\unions{A}\subseteq \unions{A_1}\union\unions{A_2}$ by \cref{subseteq}.
            We have $\unions{A_1}\union\unions{A_2}\subseteq \unions{A}$
                by \cref{union_iff,unions_iff,subseteq}.
            Hence $\unions{A} = \unions{A_1}\union\unions{A_2}$ by \cref{subseteq_antisymmetric}.
            Let $D = C\setminus U$.
            $D$ is a compact subspace of $X$ with respect to $T$
                by \cref{open_in,compact_subspace_minus_open}.
            Show that for all $x$ such that $x\in\unions{A}$ we have
                $x\in (X\union\{p\})\setminus D$.
            \begin{subproof}
                Fix $x$ such that $x\in\unions{A}$.
                Hence $x\in\unions{A_1}\union\unions{A_2}$.
                Therefore $x\in\unions{A_1}$ or $x\in\unions{A_2}$ by \cref{union_iff}.
                \begin{byCase}
                \caseOf{$x\in\unions{A_1}$.}
                    We have $x\in U$.
                    Hence $x\notin D$ by \cref{setminus_elim_right}.
                    Therefore $x\in (X\union\{p\})\setminus D$
                        by \cref{subseteq,union_intro_left,setminus_intro}.
                \caseOf{$x\in\unions{A_2}$.}
                    We have $x\in (X\union\{p\})\setminus C$.
                    Hence $x\notin C$ by \cref{setminus_elim_right}.
                    Therefore $x\in (X\union\{p\})\setminus D$
                        by \cref{setminus_elim_left,setminus_subseteq,subseteq,setminus_intro}.
                \end{byCase}
            \end{subproof}
            Hence $\unions{A}\subseteq (X\union\{p\})\setminus D$ by \cref{subseteq}.
            Show that for all $x$ such that $x\in (X\union\{p\})\setminus D$ we have
                $x\in\unions{A}$.
            \begin{subproof}
                Fix $x$ such that $x\in (X\union\{p\})\setminus D$.
                Hence $x\notin D$ by \cref{setminus_elim_right}.
                \begin{byCase}
                \caseOf{$x\in U$.}
                    Take $U_1\in A_1$ such that $x\in U_1$ by \cref{unions_iff}.
                    We have $U_1\in A$.
                    Therefore $x\in\unions{A}$ by \cref{unions_iff}.
                \caseOf{$x\notin U$.}
                    Show that $x\in (X\union\{p\})\setminus C$.
                    \begin{subproof}
                        Suppose not.
                        We have $x\in X\union\{p\}$ by \cref{setminus_elim_left}.
                        It is wrong that $x\notin C$ by \cref{setminus_intro}.
                        Therefore $x\in C$.
                        We have $x\in C$ and $x\notin U$.
                        Hence $x\in D$ by \cref{setminus_intro}.
                        Contradiction.
                    \end{subproof}
                    Hence $x\in\unions{A_2}$.
                    Take $U_1\in A_2$ such that $x\in U_1$ by \cref{unions_iff}.
                    We have $U_1\in A$.
                    Hence $x\in\unions{A}$ by \cref{unions_iff}.
                \end{byCase}
            \end{subproof}
            Therefore $(X\union\{p\})\setminus D\subseteq \unions{A}$ by \cref{subseteq}.
            Hence $\unions{A} = (X\union\{p\})\setminus D$ by \cref{subseteq_antisymmetric}.
            Show that there exists $K$ such that $K$ is a compact subspace of $X$ with respect to $T$
                and $\unions{A} = Y\setminus K$.
            \begin{subproof}
                Trivial.
            \end{subproof}
            Hence $\unions{A}\in S$ by \cref{one_point_topology_intro_right}.
        \end{byCase}
    \end{subproof}
    Show that for all $U,V$ such that $U\in S$ and $V\in S$ we have $U\inter V\in S$.
    \begin{subproof}
        Fix $U$.
        Fix $V$ such that $U\in S$ and $V\in S$.
        \begin{byCase}
        \caseOf{$p\notin U$ and $p\notin V$.}
            We have $U\subseteq X$ and $U\in T$ and $V\subseteq X$ and $V\in T$
                by \cref{one_point_topology_without_p}.
            Therefore $U\inter V\in S$
                by \cref{inter_lower_left,subseteq_transitive,topology_on,one_point_topology_intro_left}.
        \caseOf{$p\in U$ and $p\in V$.}
            Take $C_1$ such that
                $C_1$ is a compact subspace of $X$ with respect to $T$ and
                $U = (X\union\{p\})\setminus C_1$
                by \cref{one_point_topology_with_p}.
            Take $C_2$ such that
                $C_2$ is a compact subspace of $X$ with respect to $T$ and
                $V = (X\union\{p\})\setminus C_2$
                by \cref{one_point_topology_with_p}.
            Let $C_3 = C_1\union C_2$.
            Therefore $U\inter V\in S$ by \cref{compact_subspace_union,setminus_union,one_point_topology_intro_right}.
        \caseOf{$p\in U$ and $p\notin V$.}
            Take $C_1$ such that
                $C_1$ is a compact subspace of $X$ with respect to $T$ and
                $U = (X\union\{p\})\setminus C_1$
                by \cref{one_point_topology_with_p}.
            We have $V\subseteq X$ and $V\in T$ by \cref{one_point_topology_without_p}.
            $C_1\subseteq X$ by \cref{compact_subspace}.
            $C_1$ is closed in $X$ with respect to $T$ by \cref{compact_subspace,compact_subspace_closed}.
            Therefore $X\setminus C_1\in T$ by \cref{closed_in,setminus_subseteq,open_in}.
            Hence $V\setminus C_1\in T$ by \cref{setminus_eq_inter_complement,topology_on}.
            $V\subseteq X\union\{p\}$ and $C_1\subseteq X\union\{p\}$
                by \cref{subseteq_transitive,union_upper_left}.
            Hence $V\setminus C_1 = V\inter U$ by \cref{setminus_eq_inter_complement}.
            Therefore $U\inter V = V\setminus C_1$ by \cref{inter_comm}.
            Hence $U\inter V\subseteq X$ by \cref{inter_lower_left,subseteq_transitive}.
            Hence $U\inter V\in T$.
            Therefore $U\inter V\in S$ by \cref{one_point_topology_intro_left}.
        \caseOf{$p\notin U$ and $p\in V$.}
            Take $C_1$ such that
                $C_1$ is a compact subspace of $X$ with respect to $T$ and
                $V = (X\union\{p\})\setminus C_1$
                by \cref{one_point_topology_with_p}.
            We have $U\subseteq X$ and $U\in T$ by \cref{one_point_topology_without_p}.
            $C_1\subseteq X$ by \cref{compact_subspace}.
            $C_1$ is closed in $X$ with respect to $T$ by \cref{compact_subspace,compact_subspace_closed}.
            Therefore $X\setminus C_1\in T$ by \cref{closed_in,setminus_subseteq,open_in}.
            Hence $U\setminus C_1\in T$ by \cref{setminus_eq_inter_complement,topology_on}.
            $U\subseteq X\union\{p\}$ and $C_1\subseteq X\union\{p\}$
                by \cref{subseteq_transitive,union_upper_left}.
            Hence $U\setminus C_1 = U\inter V$ by \cref{setminus_eq_inter_complement}.
            Therefore $U\inter V = U\setminus C_1$.
            Hence $U\inter V\subseteq X$ by \cref{inter_lower_left,subseteq_transitive}.
            Hence $U\inter V\in T$.
            Therefore $U\inter V\in S$ by \cref{one_point_topology_intro_left}.
        \end{byCase}
    \end{subproof}
    Finally $S$ is a topology on $Y$ by \cref{topology_on}.
\end{proof}

% from §29 helper lemma 12: subspace topology from one-point topology
\begin{lemma}\label{one_point_topology_subspace}
    Assume $T$ is a topology on $X$.
    Assume $X$ is Hausdorff with respect to $T$.
    Assume $p\notin X$.
    Let $Y = X\union\{p\}$.
    Let $S = \onePointTopology{T}{X}{p}$.
    Let $T_X = \subspaceTopology{S}{X}$.
    Then $T_X = T$.
\end{lemma}
\begin{proof}
    Show that $T\subseteq T_X$.
    \begin{subproof}
        Follows by \cref{topology_on,pow_iff,subseteq,one_point_topology_intro_left,inter_absorb_supseteq_left,subspace_topology}.
    \end{subproof}
    Show that for all $U$ such that $U\in T_X$ we have $U\in T$.
    \begin{subproof}
        Fix $U$ such that $U\in T_X$.
        Take $V\in S$ such that $U = X\inter V$ by \cref{subspace_topology}.
        \begin{byCase}
        \caseOf{$p\notin V$.}
            Hence $U\in T$ by \cref{one_point_topology_without_p,inter_absorb_supseteq_left}.
        \caseOf{$p\in V$.}
            Take $C$ such that
                $C$ is a compact subspace of $X$ with respect to $T$ and
                $V = (X\union\{p\})\setminus C$
                by \cref{one_point_topology_with_p}.
            $C\subseteq X$ by \cref{compact_subspace}.
            $C$ is closed in $X$ with respect to $T$ by \cref{compact_subspace,compact_subspace_closed}.
            Therefore $X\setminus C\in T$ by \cref{closed_in,setminus_subseteq,open_in}.
            Therefore $U = X\setminus C$
                by \cref{compact_subspace,union_upper_left,subseteq_transitive,setminus_eq_inter_complement}.
            Hence $U\in T$.
        \end{byCase}
    \end{subproof}
    Therefore $T_X\subseteq T$ by \cref{subseteq}.
    Hence $T_X = T$ by \cref{subseteq_antisymmetric}.
\end{proof}

% from §29 helper lemma 13: one-point topology is compact
\begin{lemma}\label{one_point_topology_compact}
    Assume $T$ is a topology on $X$.
    Assume $X$ is Hausdorff with respect to $T$.
    Assume $p\notin X$.
    Let $Y = X\union\{p\}$.
    Let $S = \onePointTopology{T}{X}{p}$.
    Then $Y$ is compact with respect to $S$.
\end{lemma}
\begin{proof}
    Show that for all $A$ such that $A$ is an open covering of $Y$ with respect to $S$
        there exists $B$ such that $B\subseteq A$ and $B$ is finite and $B$ is a covering of $Y$.
    \begin{subproof}
        Fix $A$ such that $A$ is an open covering of $Y$ with respect to $S$.
        We have $A$ is a covering of $Y$ and
            for all $U$ such that $U\in A$ we have $U$ is open in $Y$ with respect to $S$
            by \cref{open_covering}.
        Hence $p\in\unions{A}$ by \cref{singleton_intro,union_intro_right,covering,subseteq}.
        Take $U_0\in A$ such that $p\in U_0$ by \cref{unions_iff}.
        Take $C_0$ such that
            $C_0$ is a compact subspace of $X$ with respect to $T$ and
            $U_0 = (X\union\{p\})\setminus C_0$
            by \cref{open_covering,open_in,one_point_topology_with_p}.
        Let $A_0 = A\setminus\{U_0\}$.
        Therefore $A_0$ is a covering of $C_0$
            by \cref{setminus_elim_right,compact_subspace,subseteq,union_upper_left,covering,unions_iff,singleton_iff,setminus_intro}.
        $T = \subspaceTopology{S}{X}$ and $S$ is a topology on $Y$
            by \cref{one_point_topology_subspace,one_point_topology_is_topology}.
        Hence $C_0$ is compact with respect to $\subspaceTopology{S}{C_0}$
            by \cref{compact_subspace,union_upper_left,subspace_subspace_topology}.
        Hence $C_0$ is a subspace of $Y$ with respect to $S$
            by \cref{compact_subspace,union_upper_left,subseteq_transitive,subspace_of}.
        For all $U$ such that $U\in A_0$ we have $U$ is open in $Y$ with respect to $S$
            by \cref{setminus_elim_left,open_covering}.
        Take $B$ such that
            $B\subseteq A_0$ and
            $B$ is finite and
            $B$ is a covering of $C_0$
            by \cref{compact_subspace_open_covering}.
        Let $E = B\union\{U_0\}$.
        $E$ is finite by \cref{pair_is_finite,upair_intro_left,singleton_subset_intro,subseteq_finite_is_finite,union_of_finite_is_finite}.
        We have $E\subseteq A$ by \cref{setminus_subseteq,subseteq_transitive,singleton_subset_intro,union_subsets_is_subset}.
        Hence $E$ is a covering of $Y$
            by \cref{singleton_intro,union_intro_right,unions_iff,setminus_intro,compact_subspace,subseteq_transitive,setminus_setminus,inter_absorb_supseteq_left,subseteq,covering,union_intro_left}.
        There exists $B_0$ such that
            $B_0\subseteq A$ and
            $B_0$ is finite and
            $B_0$ is a covering of $Y$.
    \end{subproof}
    Finally $Y$ is compact with respect to $S$ by \cref{compact}.
\end{proof}

% from §29 helper lemma 14: one-point topology is Hausdorff
\begin{lemma}\label{one_point_topology_hausdorff}
    Assume $T$ is a topology on $X$.
    Assume $X$ is Hausdorff with respect to $T$.
    Assume $X$ is locally compact with respect to $T$.
    Assume $p\notin X$.
    Let $Y = X\union\{p\}$.
    Let $S = \onePointTopology{T}{X}{p}$.
    Then $Y$ is Hausdorff with respect to $S$.
\end{lemma}
\begin{proof}
    $S$ is a topology on $Y$ by \cref{one_point_topology_is_topology}.
    Show that for all $x,y$ such that $x\in Y$ and $y\in Y$ and $x\neq y$ there exist $U,V$ such that
        $U$ is a neighborhood of $x$ in $Y$ with respect to $S$ and
        $V$ is a neighborhood of $y$ in $Y$ with respect to $S$ and
        $U$ is disjoint from $V$.
    \begin{subproof}
        Fix $x$.
        Fix $y$ such that $x\in Y$ and $y\in Y$ and $x\neq y$.
        Show that $x\in X$ or $y\in X$.
        \begin{subproof}
            Suppose not.
            Hence $x\notin X$ and $y\notin X$.
            Therefore $x\in\{p\}$ and $y\in\{p\}$ by \cref{union_iff}.
            Hence $x = p$ and $y = p$ by \cref{singleton_iff}.
            Thus $x = y$.
            Contradiction.
        \end{subproof}
        \begin{byCase}
        \caseOf{$x\in X$ and $y\in X$.}
            Take $U_0,V_0$ such that
                $U_0$ is a neighborhood of $x$ in $X$ with respect to $T$ and
                $V_0$ is a neighborhood of $y$ in $X$ with respect to $T$ and
                $U_0$ is disjoint from $V_0$
                by \cref{hausdorff}.
            Hence $U_0$ is a neighborhood of $x$ in $Y$ with respect to $S$ and
                $V_0$ is a neighborhood of $y$ in $Y$ with respect to $S$
                by \cref{neighborhood,open_in,topology_on,subseteq,pow_iff,one_point_topology_intro_left}.
            Show that there exist $U,V$ such that
                $U$ is a neighborhood of $x$ in $Y$ with respect to $S$ and
                $V$ is a neighborhood of $y$ in $Y$ with respect to $S$ and
                $U$ is disjoint from $V$.
            \begin{subproof}
                Trivial.
            \end{subproof}
        \caseOf{$x\in X$ and $y\notin X$.}
            Take $C$ such that
                $C$ is a compact subspace of $X$ with respect to $T$ and
                there exists $U_0$ such that
                    $U_0$ is a neighborhood of $x$ in $X$ with respect to $T$ and
                    $U_0\subseteq C$
                by \cref{locally_compact,locally_compact_at_point}.
            Take $U_0$ such that
                $U_0$ is a neighborhood of $x$ in $X$ with respect to $T$ and
                $U_0\subseteq C$.
            Let $V_0 = (X\union\{p\})\setminus C$.
            Therefore $V_0\in S$ by \cref{one_point_topology_intro_right}.
            Hence $y\in V_0$ by \cref{compact_subspace,subseteq,union_iff,singleton_iff,setminus_intro}.
            Hence $U_0$ is disjoint from $V_0$ by \cref{subseteq,setminus_elim_right,disjoint}.
            Therefore $U_0$ is a neighborhood of $x$ in $Y$ with respect to $S$ and
                $V_0$ is a neighborhood of $y$ in $Y$ with respect to $S$
                by \cref{neighborhood,open_in,topology_on,subseteq,pow_iff,one_point_topology_intro_left,compact_subspace,setminus_intro}.
            Show that there exist $U,V$ such that
                $U$ is a neighborhood of $x$ in $Y$ with respect to $S$ and
                $V$ is a neighborhood of $y$ in $Y$ with respect to $S$ and
                $U$ is disjoint from $V$.
            \begin{subproof}
                Trivial.
            \end{subproof}
        \caseOf{$x\notin X$ and $y\in X$.}
            Take $C$ such that
                $C$ is a compact subspace of $X$ with respect to $T$ and
                there exists $U_0$ such that
                    $U_0$ is a neighborhood of $y$ in $X$ with respect to $T$ and
                    $U_0\subseteq C$
                by \cref{locally_compact,locally_compact_at_point}.
            Take $U_0$ such that
                $U_0$ is a neighborhood of $y$ in $X$ with respect to $T$ and
                $U_0\subseteq C$.
            Let $V_0 = (X\union\{p\})\setminus C$.
            Therefore $V_0\in S$ by \cref{one_point_topology_intro_right}.
            Hence $x\in V_0$ by \cref{compact_subspace,subseteq,union_iff,singleton_iff,setminus_intro}.
            Hence $U_0$ is disjoint from $V_0$ by \cref{subseteq,setminus_elim_right,disjoint}.
            Therefore $U_0$ is a neighborhood of $y$ in $Y$ with respect to $S$ and
                $V_0$ is a neighborhood of $x$ in $Y$ with respect to $S$
                by \cref{neighborhood,open_in,topology_on,subseteq,pow_iff,one_point_topology_intro_left,compact_subspace,setminus_intro}.
            Therefore $V_0$ is disjoint from $U_0$ by \cref{disjoint_symmetric}.
            Show that there exist $U,V$ such that
                $U$ is a neighborhood of $x$ in $Y$ with respect to $S$ and
                $V$ is a neighborhood of $y$ in $Y$ with respect to $S$ and
                $U$ is disjoint from $V$.
            \begin{subproof}
                Trivial.
            \end{subproof}
        \end{byCase}
    \end{subproof}
    Finally $Y$ is Hausdorff with respect to $S$ by \cref{hausdorff}.
\end{proof}
% from §29 Item 4: one-point compactification existence (forward)
\begin{theorem}\label{one_point_compactification_existence_forward}
    Assume $T$ is a topology on $X$.
    If $X$ is locally compact with respect to $T$ and $X$ is Hausdorff with respect to $T$, then
        there exist $Y,S,p$ such that
            $S$ is a topology on $Y$ and
            $X\subseteq Y$ and
            $T = \subspaceTopology{S}{X}$ and
            $Y\setminus X = \{p\}$ and
            $Y$ is compact with respect to $S$ and
            $Y$ is Hausdorff with respect to $S$.
\end{theorem}
\begin{proof}
    Assume $X$ is locally compact with respect to $T$ and $X$ is Hausdorff with respect to $T$.
    Let $p_0 = X$.
    $p_0\notin X$ by \cref{set_notin_self}.
    Let $Y_0 = X\union\{p_0\}$.
    Let $S_0 = \onePointTopology{T}{X}{p_0}$.
    $S_0$ is a topology on $Y_0$ by \cref{one_point_topology_is_topology}.
    $X\subseteq Y_0$ by \cref{union_upper_left}.
    $T = \subspaceTopology{S_0}{X}$ by \cref{one_point_topology_subspace}.
    Hence $Y_0\setminus X = \{p_0\}$
        by \cref{setminus_elim_left,setminus_elim_right,union_iff,singleton_iff,union_intro_right,setminus_intro,subseteq,subseteq_antisymmetric}.
    $Y_0$ is compact with respect to $S_0$ and $Y_0$ is Hausdorff with respect to $S_0$
        by \cref{one_point_topology_compact,one_point_topology_hausdorff}.
    Show that there exist $Y,S,p$ such that
        $S$ is a topology on $Y$ and
        $X\subseteq Y$ and
        $T = \subspaceTopology{S}{X}$ and
        $Y\setminus X = \{p\}$ and
        $Y$ is compact with respect to $S$ and
        $Y$ is Hausdorff with respect to $S$.
    \begin{subproof}
        Trivial.
    \end{subproof}
\end{proof}

% from §29 Item 4: one-point compactification existence (reverse)
\begin{theorem}\label{one_point_compactification_existence_reverse}
    Assume $T$ is a topology on $X$.
    If there exist $Y,S,p$ such that
            $S$ is a topology on $Y$ and
            $X\subseteq Y$ and
            $T = \subspaceTopology{S}{X}$ and
            $Y\setminus X = \{p\}$ and
            $Y$ is compact with respect to $S$ and
            $Y$ is Hausdorff with respect to $S$
        then $X$ is locally compact with respect to $T$ and $X$ is Hausdorff with respect to $T$.
\end{theorem}
\begin{proof}
    Assume there exist $Y,S,p$ such that
        $S$ is a topology on $Y$ and
        $X\subseteq Y$ and
        $T = \subspaceTopology{S}{X}$ and
        $Y\setminus X = \{p\}$ and
        $Y$ is compact with respect to $S$ and
        $Y$ is Hausdorff with respect to $S$.
    Take $Y_0,S_0,p_0$ such that
        $S_0$ is a topology on $Y_0$ and
        $X\subseteq Y_0$ and
        $T = \subspaceTopology{S_0}{X}$ and
        $Y_0\setminus X = \{p_0\}$ and
        $Y_0$ is compact with respect to $S_0$ and
        $Y_0$ is Hausdorff with respect to $S_0$.
    Hence $Y_0 = X\union\{p_0\}$ by \cref{setminus_partition}.
    Show that $X$ is Hausdorff with respect to $T$.
    \begin{subproof}
        Follows by \cref{subspace_of,subspace_hausdorff}.
    \end{subproof}
    Show that for all $x\in X$ we have $X$ is locally compact at $x$ with respect to $T$.
    \begin{subproof}
            Fix $x$ such that $x\in X$.
            We have $p_0\in Y_0$ and $p_0\notin X$ by \cref{singleton_intro,setminus_elim_left,setminus_elim_right}.
            $x\in Y_0$ by \cref{subseteq}.
            Therefore $x\neq p_0$.
            Take $U,V$ such that
                $U$ is a neighborhood of $x$ in $Y_0$ with respect to $S_0$ and
                $V$ is a neighborhood of $p_0$ in $Y_0$ with respect to $S_0$ and
                $U$ is disjoint from $V$
                by \cref{hausdorff}.
            Let $C = Y_0\setminus V$.
            $C\subseteq Y_0$ and $V\subseteq Y_0$ and $V$ is open in $Y_0$ with respect to $S_0$
                by \cref{setminus_subseteq,neighborhood,open_in,topology_on,subseteq,pow_iff}.
            Hence $Y_0\setminus C = V$ by \cref{setminus_setminus,inter_absorb_supseteq_left}.
            Therefore $C$ is closed in $Y_0$ with respect to $S_0$ and
                $C$ is compact with respect to $\subspaceTopology{S_0}{C}$
                by \cref{neighborhood,closed_in,closed_subspace_compact}.
            Hence $p_0\notin C$ by \cref{neighborhood,setminus_elim_right}.
            Hence $C\subseteq X$ by \cref{setminus_elim_left,union_iff,singleton_iff,setminus_elim_right,subseteq}.
            Therefore $C$ is a compact subspace of $X$ with respect to $T$
                by \cref{subspace_subspace_topology,compact_subspace}.
            Hence $U\inter X$ is a neighborhood of $x$ in $X$ with respect to $T$
                by \cref{neighborhood,inter_comm,open_in,topology_on,subseteq,pow_iff,subspace_topology,inter_intro}.
            Therefore $U\inter X\subseteq C$
                by \cref{disjoint,inter_lower_left,inter_lower_right,setminus_intro,subseteq}.
            Hence $X$ is locally compact at $x$ with respect to $T$ by \cref{locally_compact_at_point}.
    \end{subproof}
    Therefore $X$ is locally compact with respect to $T$ by \cref{locally_compact}.
    Show that $X$ is locally compact with respect to $T$ and $X$ is Hausdorff with respect to $T$.
    \begin{subproof}
        Trivial.
    \end{subproof}
\end{proof}

 
% from §29 Item 5: one-point compactification uniqueness
\begin{theorem}\label{one_point_compactification_uniqueness}
    Assume $T$ is a topology on $X$.
    Assume $S$ is a topology on $Y$.
    Assume $S_1$ is a topology on $Y_1$.
    Suppose $X$ is a subspace of $Y$ with respect to $S$.
    Suppose $X$ is a subspace of $Y_1$ with respect to $S_1$.
    Suppose $T = \subspaceTopology{S}{X}$.
    Suppose $T = \subspaceTopology{S_1}{X}$.
    Suppose $Y\setminus X = \{p\}$.
    Suppose $Y_1\setminus X = \{q\}$.
    Suppose $Y$ is compact with respect to $S$ and $Y$ is Hausdorff with respect to $S$.
    Suppose $Y_1$ is compact with respect to $S_1$ and $Y_1$ is Hausdorff with respect to $S_1$.
    Then there exists $h$ such that
        $h$ is a homeomorphism between $Y$ and $Y_1$ with respect to $S$ and $S_1$ and
        $\restrl{h}{X} = \identity{X}$.
\end{theorem}
\begin{proof}
    We have $X\subseteq Y$ and $X\subseteq Y_1$ by \cref{subspace_of}.
    Therefore $Y = X\union\{p\}$ and $Y_1 = X\union\{q\}$ by \cref{setminus_partition}.
    Hence $p\in Y$ and $q\in Y_1$ and $p\notin X$ and $q\notin X$
        by \cref{singleton_intro,setminus_elim_left,setminus_elim_right}.

    Let $f = \identity{X}$.
    Let $g = \{(p,q)\}$.
    Let $h = f\union g$.

    $f$ is a function by \cref{id_is_function}.
    We have $g$ is a function by \cref{singleton_iff,relation,pair_eq_iff,rightunique}.

    $\dom{f} = X$ by \cref{id_dom}.
    $g = \cons{(p,q)}{\emptyset}$.
    Hence $\dom{g} = \{p\}$ by \cref{dom_cons,dom_emptyset}.
    For all $x$ such that $x\in\dom{f}\inter\dom{g}$ we have $f(x) = g(x)$.
    Therefore $h$ is a function by \cref{union_compatible_functions}.

    Hence $\dom{h} = Y$ by \cref{dom_union}.

    For all $x$ such that $x\in\dom{h}$ we have $h(x)\in Y_1$
        by \cref{subseteq,union_iff,id_elem_intro,union_intro_left,function_apply_intro,singleton_iff,singleton_intro,union_intro_right}.
    Therefore $h\in\funs{Y}{Y_1}$ by \cref{funs_intro}.

    For all $y$ such that $y\in Y_1$ there exists $x$ such that $x\in Y$ and $h(x) = y$
        by \cref{union_iff,union_intro_left,id_elem_intro,union_intro_left,function_apply_intro,singleton_iff,union_intro_right,singleton_intro}.
    Hence $h\in\surj{Y}{Y_1}$ by \cref{surj}.

    Show that for all $x_1,x_2$ such that $x_1\in\dom{h}$ and $x_2\in\dom{h}$ and $h(x_1) = h(x_2)$
        we have $x_1 = x_2$.
    \begin{subproof}
        Fix $x_1$.
        Fix $x_2$ such that $x_1\in\dom{h}$ and $x_2\in\dom{h}$ and $h(x_1) = h(x_2)$.
        We have $x_1\in Y$ and $x_2\in Y$.
        Hence $x_1\in X$ or $x_1\in\{p\}$ by \cref{union_iff}.
        Therefore $x_2\in X$ or $x_2\in\{p\}$ by \cref{union_iff}.
        \begin{byCase}
        \caseOf{$x_1\in X$ and $x_2\in X$.}
            We have $x_1 = x_2$
                by \cref{id_elem_intro,union_intro_left,function_apply_intro}.
        \caseOf{$x_1\in X$ and $x_2\in\{p\}$.}
            Suppose not.
            Hence $h(x_1) = x_1$ and $h(x_2) = q$
                by \cref{id_elem_intro,union_intro_left,function_apply_intro,singleton_iff,singleton_intro,union_intro_right}.
            Therefore $x_1 = q$.
            Thus $x_1\notin X$.
            Contradiction.
            Therefore $x_1 = x_2$.
        \caseOf{$x_1\in\{p\}$ and $x_2\in X$.}
            Suppose not.
            Hence $h(x_1) = q$ and $h(x_2) = x_2$
                by \cref{singleton_iff,singleton_intro,union_intro_right,function_apply_intro,id_elem_intro,union_intro_left}.
            Therefore $x_2 = q$.
            Thus $x_2\notin X$.
            Contradiction.
            Therefore $x_1 = x_2$.
        \caseOf{$x_1\in\{p\}$ and $x_2\in\{p\}$.}
            We have $x_1 = x_2$ by \cref{singleton_iff}.
        \end{byCase}
    \end{subproof}
    Therefore $h$ is injective by \cref{injective_function}.
    Hence $h\in\bijections{Y}{Y_1}$ by \cref{bijections}.

    Show that for all $U$ such that $U$ is open in $Y_1$ with respect to $S_1$
        we have $\preimg{h}{U}$ is open in $Y$ with respect to $S$.
    \begin{subproof}
        Fix $U$ such that $U$ is open in $Y_1$ with respect to $S_1$.
        \begin{byCase}
        \caseOf{$q\notin U$.}
            We have $U\subseteq X$
                by \cref{open_in,topology_on,subseteq,pow_iff,union_iff,singleton_iff,subseteq}.
            For all $y$ such that $y\in\preimg{h}{U}$ we have $y\in U$
                by \cref{preimg_iff,preimg_subseteq_dom,subseteq,union_iff,singleton_iff,singleton_intro,union_intro_right,function_apply_intro,id_elem_intro,union_intro_left}.
            Hence $\preimg{h}{U}\subseteq U$ by \cref{subseteq}.
            Therefore $U\subseteq\preimg{h}{U}$
                by \cref{subseteq,id_elem_intro,union_intro_left,function_apply_intro,preimg_iff}.
            Hence $\preimg{h}{U} = U$ by \cref{subseteq_antisymmetric}.
            Therefore $U$ is open in $X$ with respect to $T$
                by \cref{open_in,inter_absorb_supseteq_left,subspace_topology}.
            Hence $\{p\}$ is closed in $Y$ with respect to $S$
                by \cref{singleton_subset_intro,pair_is_finite,upair_intro_left,subseteq_finite_is_finite,finite_point_set_closed_hausdorff}.
            Therefore $Y\setminus\{p\}$ is open in $Y$ with respect to $S$ by \cref{closed_in}.
            Hence $X = Y\setminus\{p\}$
                by \cref{union_intro_left,singleton_iff,setminus_intro,setminus_elim_left,union_iff,setminus_elim_right,subseteq,subseteq_antisymmetric}.
            Hence $U$ is open in $Y$ with respect to $S$ by \cref{open_in_subspace_open_in_space,subspace_of}.
            Therefore $\preimg{h}{U}$ is open in $Y$ with respect to $S$.
        \caseOf{$q\in U$.}
            Let $C = Y_1\setminus U$.
            $C\subseteq Y_1$ and $U\in S_1$ by \cref{setminus_subseteq,open_in}.
            Hence $U\subseteq Y_1$ by \cref{topology_on,subseteq,pow_iff}.
            Therefore $Y_1\setminus C = U$ by \cref{setminus_setminus,inter_absorb_supseteq_left}.
            Hence $C$ is closed in $Y_1$ with respect to $S_1$ by \cref{open_in,closed_in}.
            Therefore $C\subseteq X$
                by \cref{setminus_elim_left,union_iff,singleton_iff,setminus_elim_right,subseteq}.
            Hence $C$ is compact with respect to $\subspaceTopology{S_1}{C}$ by \cref{closed_subspace_compact}.
            Therefore $C$ is a compact subspace of $X$ with respect to $T$
                by \cref{subspace_subspace_topology,compact_subspace}.
            Hence $C$ is compact with respect to $\subspaceTopology{S}{C}$ by \cref{subspace_subspace_topology}.
            Therefore $C$ is closed in $Y$ with respect to $S$
                by \cref{subseteq_transitive,compact_subspace,compact_subspace_closed}.
            Therefore $Y\setminus C$ is open in $Y$ with respect to $S$ by \cref{closed_in}.
            For all $y$ such that $y\in\preimg{h}{U}$ we have $y\in Y\setminus C$
                by \cref{preimg_subseteq_dom,subseteq,preimg_iff,id_elem_intro,union_intro_left,function_apply_intro,setminus_elim_right,setminus_intro}.
            Therefore $\preimg{h}{U}\subseteq Y\setminus C$ by \cref{subseteq}.
            Hence $Y\setminus C\subseteq\preimg{h}{U}$
                by \cref{setminus_elim_left,setminus_elim_right,union_iff,union_intro_left,setminus_intro,id_elem_intro,function_apply_intro,preimg_iff,singleton_intro,union_intro_right,singleton_iff,subseteq}.
            Therefore $\preimg{h}{U} = Y\setminus C$ by \cref{subseteq_antisymmetric}.
            Hence $\preimg{h}{U}$ is open in $Y$ with respect to $S$ by \cref{subseteq}.
        \end{byCase}
    \end{subproof}
    Therefore $h$ is continuous from $Y$ to $Y_1$ with respect to $S$ and $S_1$ by \cref{continuous}.

    Hence $h$ is a homeomorphism between $Y$ and $Y_1$ with respect to $S$ and $S_1$
        by \cref{bijections,compact_hausdorff_bijection_homeomorphism}.

    Show that $\restrl{h}{X} = \identity{X}$.
    \begin{subproof}
        We have $\restrl{h}{X}\subseteq \identity{X}$
            by \cref{restrl,id_elem_intro,union_intro_left,function_apply_intro,subseteq}.
        Therefore $\identity{X}\subseteq\restrl{h}{X}$
            by \cref{id_elem_inspect,id_elem_intro,union_intro_left,restrl,subseteq}.
        Hence $\restrl{h}{X} = \identity{X}$ by \cref{subseteq_antisymmetric}.
    \end{subproof}
    Show that there exists $k$ such that
        $k$ is a homeomorphism between $Y$ and $Y_1$ with respect to $S$ and $S_1$ and
        $\restrl{k}{X} = \identity{X}$.
    \begin{subproof}
        Trivial.
    \end{subproof}
\end{proof}


% from §29 Item 6: compactification
\begin{definition}\label{compactification}
    $Y$ is a compactification of $X$ with respect to $S$ and $T$ iff
        $S$ is a topology on $Y$ and
        $T$ is a topology on $X$ and
        $X$ is a subspace of $Y$ with respect to $S$ and
        $T = \subspaceTopology{S}{X}$ and
        $X\neq Y$ and
        $\closure{X}{Y}{S} = Y$ and
        $Y$ is compact with respect to $S$ and
        $Y$ is Hausdorff with respect to $S$.
\end{definition}

% from §29 Item 7: one-point compactification
\begin{definition}\label{one_point_compactification}
    $Y$ is a one-point compactification of $X$ with respect to $S$ and $T$ iff
        $Y$ is a compactification of $X$ with respect to $S$ and $T$ and
        there exists $p$ such that $Y\setminus X = \{p\}$.
\end{definition}

% from §29 Item 8: one-point compactification characterization
\begin{theorem}\label{one_point_compactification_characterization}
    Assume $T$ is a topology on $X$.
    Then there exist $Y,S$ such that $Y$ is a one-point compactification of $X$ with respect to $S$ and $T$
        iff $X$ is locally compact with respect to $T$ and $X$ is Hausdorff with respect to $T$
            and $X$ is not compact with respect to $T$.
\end{theorem}
\begin{proof}
    Show that if there exist $Y,S$ such that $Y$ is a one-point compactification of $X$ with respect to $S$ and $T$ then
        $X$ is locally compact with respect to $T$ and $X$ is Hausdorff with respect to $T$ and
            $X$ is not compact with respect to $T$.
    \begin{subproof}
        Assume there exist $Y,S$ such that $Y$ is a one-point compactification of $X$ with respect to $S$ and $T$.
        Take $Y,S$ such that $Y$ is a one-point compactification of $X$ with respect to $S$ and $T$.
        Take $p$ such that $Y\setminus X = \{p\}$ and
            $Y$ is a compactification of $X$ with respect to $S$ and $T$
            by \cref{one_point_compactification}.
        $T = \subspaceTopology{S}{X}$ by \cref{compactification}.
        $Y$ is compact with respect to $S$ and $Y$ is Hausdorff with respect to $S$ by \cref{compactification}.
        $X\subseteq Y$ and $S$ is a topology on $Y$ by \cref{subspace_of,compactification}.
        Show that $X$ is locally compact with respect to $T$ and $X$ is Hausdorff with respect to $T$.
        \begin{subproof}
            Follows by \cref{one_point_compactification_existence_reverse}.
        \end{subproof}

        Show that $X$ is not compact with respect to $T$.
        \begin{subproof}
            Suppose not.
            Therefore $X$ is closed in $Y$ with respect to $S$
                by \cref{compact_subspace,compactification,compact_subspace_closed}.
            Hence $\closure{X}{Y}{S} = X$
                by \cref{subseteq_refl,closure_subseteq_closed,subseteq_closure,subseteq_antisymmetric}.
            Therefore $X = Y$ by \cref{compactification}.
            Contradiction by \cref{compactification}.
        \end{subproof}
        Show that $X$ is locally compact with respect to $T$ and $X$ is Hausdorff with respect to $T$
            and $X$ is not compact with respect to $T$.
        \begin{subproof}
            Trivial.
        \end{subproof}
    \end{subproof}

    Show that if $X$ is locally compact with respect to $T$ and $X$ is Hausdorff with respect to $T$
        and $X$ is not compact with respect to $T$
        then there exist $Y,S$ such that $Y$ is a one-point compactification of $X$ with respect to $S$ and $T$.
    \begin{subproof}
        Assume $X$ is locally compact with respect to $T$ and $X$ is Hausdorff with respect to $T$
            and $X$ is not compact with respect to $T$.
        Take $Y,S,p$ such that
            $S$ is a topology on $Y$ and
            $X\subseteq Y$ and
            $T = \subspaceTopology{S}{X}$ and
            $Y\setminus X = \{p\}$ and
            $Y$ is compact with respect to $S$ and
            $Y$ is Hausdorff with respect to $S$
            by \cref{one_point_compactification_existence_forward}.
        We have $Y = X\union\{p\}$ by \cref{setminus_partition}.

        Show that $p\in \closure{X}{Y}{S}$.
        \begin{subproof}
            Suppose not.
            Show that for all $y$ such that $y\in\closure{X}{Y}{S}$ we have $y\in X$.
            \begin{subproof}
                Fix $y$ such that $y\in\closure{X}{Y}{S}$.
                Hence $y\in Y$ by \cref{closure}.
                Therefore $y\in X$ or $y\in\{p\}$ by \cref{union_iff}.
                \begin{byCase}
                \caseOf{$y\in X$.}
                    Trivial.
                \caseOf{$y\in\{p\}$.}
                    We have $y = p$ by \cref{singleton_iff}.
                    Therefore $p\in\closure{X}{Y}{S}$.
                    Contradiction.
                    Thus $y\in X$.
                \end{byCase}
            \end{subproof}
            Therefore $\closure{X}{Y}{S}\subseteq X$ by \cref{subseteq}.
            Therefore $X$ is closed in $Y$ with respect to $S$
                by \cref{subseteq_closure,subseteq_antisymmetric,closure_closed_in}.
            Hence $X$ is compact with respect to $T$ by \cref{closed_subspace_compact}.
            Contradiction.
        \end{subproof}
        Hence $\closure{X}{Y}{S} = Y$
            by \cref{closure_closed_in,subseteq_closure,singleton_subset_intro,union_subsets_is_subset,closed_in,subseteq_antisymmetric}.

        Therefore $X\neq Y$ by \cref{singleton_intro,setminus_elim_right}.
        Hence $X$ is a subspace of $Y$ with respect to $S$ by \cref{subspace_of}.
        Therefore $Y$ is a compactification of $X$ with respect to $S$ and $T$ by \cref{compactification}.
        Hence $Y$ is a one-point compactification of $X$ with respect to $S$ and $T$
            by \cref{one_point_compactification}.
        Show that there exist $Y_1,S_1$ such that $Y_1$ is a one-point compactification of $X$ with respect to $S_1$ and $T$.
        \begin{subproof}
            Trivial.
        \end{subproof}
    \end{subproof}
\end{proof}

% from §29 Item 9: local compactness via neighborhoods
\begin{theorem}\label{locally_compact_neighborhood_characterization}
    Assume $T$ is a topology on $X$.
    Assume $X$ is Hausdorff with respect to $T$.
    Then $X$ is locally compact with respect to $T$ iff
        for all $x\in X$ we have
            for all $U$ such that $U$ is a neighborhood of $x$ in $X$ with respect to $T$ we have
                there exists $V$ such that
                    $V$ is a neighborhood of $x$ in $X$ with respect to $T$ and
                    $\closure{V}{X}{T}$ is a compact subspace of $X$ with respect to $T$ and
                    $\closure{V}{X}{T}\subseteq U$.
\end{theorem}
\begin{proof}
    Show that if for all $x\in X$ we have
        for all $U$ such that $U$ is a neighborhood of $x$ in $X$ with respect to $T$ we have
            there exists $V$ such that
                $V$ is a neighborhood of $x$ in $X$ with respect to $T$ and
                $\closure{V}{X}{T}$ is a compact subspace of $X$ with respect to $T$ and
                $\closure{V}{X}{T}\subseteq U$
        then $X$ is locally compact with respect to $T$.
    \begin{subproof}
        Assume for all $x\in X$ we have
            for all $U$ such that $U$ is a neighborhood of $x$ in $X$ with respect to $T$ we have
                there exists $V$ such that
                    $V$ is a neighborhood of $x$ in $X$ with respect to $T$ and
                    $\closure{V}{X}{T}$ is a compact subspace of $X$ with respect to $T$ and
                    $\closure{V}{X}{T}\subseteq U$.
        For all $x$ such that $x\in X$ we have $X$ is a neighborhood of $x$ in $X$ with respect to $T$
            by \cref{topology_on,open_in,neighborhood}.
        For all $x$ such that $x\in X$ there exists $V$ such that
            $V$ is a neighborhood of $x$ in $X$ with respect to $T$ and
            $\closure{V}{X}{T}$ is a compact subspace of $X$ with respect to $T$ and
            $\closure{V}{X}{T}\subseteq X$.
        For all $x\in X$ we have $X$ is locally compact at $x$ with respect to $T$
            by \cref{neighborhood,open_in,topology_on,subseteq,pow_iff,subseteq_closure,locally_compact_at_point}.
        $X$ is locally compact with respect to $T$ by \cref{locally_compact}.
    \end{subproof}

    Show that if $X$ is locally compact with respect to $T$ then
        for all $x\in X$ we have
            for all $U$ such that $U$ is a neighborhood of $x$ in $X$ with respect to $T$ we have
                there exists $V$ such that
                    $V$ is a neighborhood of $x$ in $X$ with respect to $T$ and
                    $\closure{V}{X}{T}$ is a compact subspace of $X$ with respect to $T$ and
                    $\closure{V}{X}{T}\subseteq U$.
    \begin{subproof}
        Assume $X$ is locally compact with respect to $T$.
        Let $p = X$.
        $p\notin X$ by \cref{set_notin_self}.
        Let $Y = X\union\{p\}$.
        Let $S = \onePointTopology{T}{X}{p}$.
        $S$ is a topology on $Y$ and $Y$ is compact with respect to $S$ and $Y$ is Hausdorff with respect to $S$
            by \cref{one_point_topology_is_topology,one_point_topology_compact,one_point_topology_hausdorff}.
        $T = \subspaceTopology{S}{X}$ by \cref{one_point_topology_subspace}.
        $X\subseteq Y$ by \cref{union_upper_left}.
        $X$ is a subspace of $Y$ with respect to $S$ by \cref{subspace_of}.

        Show that for all $x\in X$ we have
            for all $U$ such that $U$ is a neighborhood of $x$ in $X$ with respect to $T$ we have
                there exists $V$ such that
                    $V$ is a neighborhood of $x$ in $X$ with respect to $T$ and
                    $\closure{V}{X}{T}$ is a compact subspace of $X$ with respect to $T$ and
                    $\closure{V}{X}{T}\subseteq U$.
        \begin{subproof}
            Fix $x$ such that $x\in X$.
            Fix $U$ such that $U$ is a neighborhood of $x$ in $X$ with respect to $T$.
            $U$ is open in $X$ with respect to $T$ by \cref{neighborhood}.
            $U\subseteq X$ by \cref{open_in,topology_on,subseteq,pow_iff}.

            $X$ is open in $Y$ with respect to $S$
                by \cref{topology_on,subseteq_refl,one_point_topology_intro_left,open_in}.
            $U$ is open in $Y$ with respect to $S$ by
                \cref{open_in_subspace_open_in_space,subspace_of}.

            Let $C = Y\setminus U$.
            $C\subseteq Y$ by \cref{setminus_subseteq}.
            $U\in S$ and $U\subseteq Y$ by \cref{open_in,topology_on,subseteq,pow_iff}.
            $Y\setminus C = U$ by \cref{setminus_setminus,inter_absorb_supseteq_left}.
            $C$ is closed in $Y$ with respect to $S$ and
                $C$ is compact with respect to $\subspaceTopology{S}{C}$
                by \cref{setminus_subseteq,open_in,closed_in,closed_subspace_compact}.

            $x\in U$ by \cref{neighborhood}.
            $x\notin C$ by \cref{setminus_elim_right}.
            Take $V,W$ such that
                $V$ is open in $Y$ with respect to $S$ and
                $W$ is open in $Y$ with respect to $S$ and
                $x\in V$ and
                $C\subseteq W$ and
                $V$ is disjoint from $W$
                by \cref{compact_hausdorff_separation}.

            $p\notin U$ by \cref{subseteq}.
            $p\in W$ by \cref{singleton_intro,union_intro_right,setminus_intro,subseteq}.
            $p\notin V$ by \cref{disjoint}.

            $V\subseteq X$ by \cref{open_in,topology_on,subseteq,pow_iff,union_iff,singleton_iff}.

            $V$ is a neighborhood of $x$ in $X$ with respect to $T$
                by \cref{inter_absorb_supseteq_left,subspace_topology,open_in,neighborhood}.
            $V\subseteq Y$ by \cref{subseteq_transitive}.

            Let $B = \closure{V}{Y}{S}$.
            $B$ is closed in $Y$ with respect to $S$ and
                $B$ is compact with respect to $\subspaceTopology{S}{B}$
                by \cref{closure_closed_in,closed_in,closed_subspace_compact}.

            $Y\setminus W$ is closed in $Y$ with respect to $S$ by \cref{setminus_subseteq,open_in,topology_on,subseteq,pow_iff,setminus_setminus,inter_absorb_supseteq_left,closed_in}.
            $V\subseteq Y\setminus W$ by \cref{disjoint,subseteq,setminus_intro,subseteq}.
            $B\subseteq Y\setminus W$ by \cref{closure_subseteq_closed}.

            For all $y$ such that $y\in Y\setminus W$ we have $y\in U$.
            $Y\setminus W\subseteq U$ by \cref{subseteq}.
            $B\subseteq U$ and $B\subseteq X$ by \cref{subseteq_transitive}.

            $B$ is a compact subspace of $X$ with respect to $T$
                by \cref{subspace_subspace_topology,compact_subspace}.

            $\closure{V}{X}{T} = B$ by
                \cref{closure_in_subspace,subspace_of,inter_comm,inter_absorb_supseteq_left}.
            $\closure{V}{X}{T}$ is a compact subspace of $X$ with respect to $T$.
            $\closure{V}{X}{T}\subseteq U$ by \cref{subseteq}.
        \end{subproof}
    \end{subproof}
\end{proof}

% from §29 Item 10: locally compact subspaces
\begin{corollary}\label{locally_compact_subspace}
    Assume $T$ is a topology on $X$.
    Assume $X$ is locally compact with respect to $T$.
    Assume $X$ is Hausdorff with respect to $T$.
    Suppose $A$ is a subspace of $X$ with respect to $T$.
    If $A$ is closed in $X$ with respect to $T$ or $A$ is open in $X$ with respect to $T$ then
        $A$ is locally compact with respect to $\subspaceTopology{T}{A}$.
\end{corollary}
\begin{proof}
    Let $T_A = \subspaceTopology{T}{A}$.
    Assume $A$ is closed in $X$ with respect to $T$ or $A$ is open in $X$ with respect to $T$.
    \begin{byCase}
    \caseOf{$A$ is closed in $X$ with respect to $T$.}
        Show that for all $x\in A$ we have $A$ is locally compact at $x$ with respect to $T_A$.
        \begin{subproof}
            Fix $x$ such that $x\in A$.
            $A\subseteq X$ by \cref{subspace_of}.
            $x\in X$ by \cref{subseteq}.
            $X$ is locally compact at $x$ with respect to $T$ by \cref{locally_compact}.
            Take $C$ such that
                $C$ is a compact subspace of $X$ with respect to $T$ and
                there exists $U$ such that
                    $U$ is a neighborhood of $x$ in $X$ with respect to $T$ and
                    $U\subseteq C$
                by \cref{locally_compact_at_point}.
            Take $U$ such that
                $U$ is a neighborhood of $x$ in $X$ with respect to $T$ and
                $U\subseteq C$.
            Let $D = C\inter A$.
            $D\subseteq A$ by \cref{inter_lower_right}.
            Let $T_C = \subspaceTopology{T}{C}$.
            $C$ is compact with respect to $T_C$ by \cref{compact_subspace}.
            $C$ is a subspace of $X$ with respect to $T$ by \cref{subspace_of,compact_subspace}.
            $D$ is closed in $C$ with respect to $T_C$ by \cref{inter_comm,closed_in_subspace_iff}.
            $D\subseteq C$ by \cref{inter_lower_left}.
            $T_C$ is a topology on $C$ by \cref{subspace_topology_is_topology,subspace_of}.
            $C\subseteq X$ by \cref{compact_subspace}.
            $D$ is compact with respect to $\subspaceTopology{T_C}{D}$ by \cref{closed_subspace_compact}.
            $D$ is a compact subspace of $A$ with respect to $T_A$
                by \cref{subspace_subspace_topology,compact_subspace}.

            $U\inter A$ is a neighborhood of $x$ in $A$ with respect to $T_A$
                by \cref{inter_comm,neighborhood,open_in,subspace_topology,subspace_topology_is_topology,subspace_of,inter_intro}.
            $U\inter A\subseteq C\inter A$
                by \cref{inter_lower_left,subseteq_transitive,inter_lower_right,subseteq_inter_iff}.
            $A$ is locally compact at $x$ with respect to $T_A$ by \cref{locally_compact_at_point}.
        \end{subproof}
        Hence $A$ is locally compact with respect to $T_A$ by \cref{locally_compact}.
    \caseOf{$A$ is open in $X$ with respect to $T$.}
        Show that for all $x\in A$ we have $A$ is locally compact at $x$ with respect to $T_A$.
        \begin{subproof}
            Fix $x$ such that $x\in A$.
            $A\subseteq X$ by \cref{subspace_of}.
            $x\in X$ by \cref{subseteq}.
            $A$ is a neighborhood of $x$ in $X$ with respect to $T$ by \cref{neighborhood}.
            Take $V$ such that
                $V$ is a neighborhood of $x$ in $X$ with respect to $T$ and
                $\closure{V}{X}{T}$ is a compact subspace of $X$ with respect to $T$ and
                $\closure{V}{X}{T}\subseteq A$
                by \cref{locally_compact_neighborhood_characterization}.

            $V$ is open in $X$ with respect to $T$ by \cref{neighborhood}.
            $V\subseteq X$ by \cref{open_in,topology_on,subseteq,pow_iff}.
            $V\subseteq \closure{V}{X}{T}$ by \cref{subseteq_closure}.
            $V\subseteq A$ by \cref{subseteq_transitive}.
            $V$ is open in $A$ with respect to $T_A$ and
                $V$ is a neighborhood of $x$ in $A$ with respect to $T_A$
                by \cref{inter_absorb_supseteq_left,neighborhood,open_in,subspace_topology,subspace_topology_is_topology,subspace_of}.

            Let $D = \closure{V}{X}{T}$.
            $D\subseteq A$ and $D$ is a compact subspace of $A$ with respect to $T_A$
                by \cref{subseteq,subspace_subspace_topology,compact_subspace}.
            $A$ is locally compact at $x$ with respect to $T_A$ by \cref{locally_compact_at_point}.
        \end{subproof}
        $A$ is locally compact with respect to $T_A$ by \cref{locally_compact}.
    \end{byCase}
\end{proof}

% from §29 helper lemma 15: homeomorphism preserves local compactness
\begin{lemma}\label{homeomorphism_preserves_local_compactness}
    Assume $T$ is a topology on $X$.
    Assume $S$ is a topology on $Y$.
    Suppose $f$ is a homeomorphism between $X$ and $Y$ with respect to $T$ and $S$.
    If $Y$ is locally compact with respect to $S$ then $X$ is locally compact with respect to $T$.
\end{lemma}
\begin{proof}
    Assume $Y$ is locally compact with respect to $S$.
    Let $g = \converse{f}$.
    $f\in\funs{X}{Y}$ and $f$ is a function and $\dom{f} = X$
        by \cref{homeomorphism,bijections,bijections_to_funs,funs_is_function,funs}.
    We have $g$ is a bijection from $Y$ to $X$ and $g\in\bijections{Y}{X}$ and $\dom{g} = Y$ and $g$ is a function
        by \cref{homeomorphism,bijection_converse_is_bijection,bijections,bijections_dom,bijection_is_function}.
    $g$ is continuous from $Y$ to $X$ with respect to $S$ and $T$ by \cref{homeomorphism}.

    Show that for all $x\in X$ we have $X$ is locally compact at $x$ with respect to $T$.
        \begin{subproof}
        Fix $x$ such that $x\in X$.
        Let $y = f(x)$.
        $Y$ is locally compact at $y$ with respect to $S$ by \cref{funs_type_apply,locally_compact}.
        Take $C$ such that
            $C$ is a compact subspace of $Y$ with respect to $S$ and
            there exists $U$ such that
                $U$ is a neighborhood of $y$ in $Y$ with respect to $S$ and
                $U\subseteq C$
            by \cref{locally_compact_at_point}.
        Take $U$ such that
            $U$ is a neighborhood of $y$ in $Y$ with respect to $S$ and
            $U\subseteq C$.

        Let $g_C = \restrl{g}{C}$.
        $C\subseteq \dom{g}$ and $C$ is a subspace of $Y$ with respect to $S$
            by \cref{compact_subspace,bijections_dom,subspace_of}.
        Let $T_C = \subspaceTopology{S}{C}$.
        $g_C$ is continuous from $C$ to $X$ with respect to $T_C$ and $T$ by \cref{continuous_restrict_domain}.
        $C$ is compact with respect to $T_C$ by \cref{compact_subspace}.
        Let $A = \img{g_C}{C}$.
        $A$ is compact with respect to $\subspaceTopology{T}{A}$ by \cref{continuous_image_compact}.
        $g_C\in\funs{C}{X}$ by \cref{continuous}.
        $A\subseteq X$
            by \cref{continuous,funs_subseteq_rels,subseteq,rels_ran_subseteq,img_subseteq_ran,subseteq_transitive}.
        $A$ is a compact subspace of $X$ with respect to $T$ by \cref{compact_subspace}.

        Let $V = \img{g}{U}$.
        $V = \preimg{f}{U}$ by \cref{preim_eq_img_of_converse}.
        $U$ is open in $Y$ with respect to $S$ and $y\in U$ by \cref{neighborhood}.
        $f$ is continuous from $X$ to $Y$ with respect to $T$ and $S$ by \cref{homeomorphism}.
        $V$ is a neighborhood of $x$ in $X$ with respect to $T$
            by \cref{continuous,function_apply_elim,preimg_iff,neighborhood}.

        $V\subseteq A$
            by \cref{subseteq,img_subseteq,continuous_restrict_domain,funs,img_dom,restrl_ran}.
        $X$ is locally compact at $x$ with respect to $T$ by \cref{locally_compact_at_point}.
    \end{subproof}
    Hence $X$ is locally compact with respect to $T$ by \cref{locally_compact}.
\end{proof}

% from §29 Item 11: locally compact Hausdorff embeds as open subspace
\begin{corollary}\label{locally_compact_open_subspace_compact_hausdorff_forward}
    Assume $T$ is a topology on $X$.
    If $X$ is locally compact with respect to $T$ and $X$ is Hausdorff with respect to $T$ then
        there exist $Y,S,A,f$ such that
            $S$ is a topology on $Y$ and
            $Y$ is compact with respect to $S$ and
            $Y$ is Hausdorff with respect to $S$ and
            $A$ is open in $Y$ with respect to $S$ and
            $f$ is a homeomorphism between $X$ and $A$ with respect to $T$ and $\subspaceTopology{S}{A}$.
\end{corollary}
\begin{proof}
    Assume $X$ is locally compact with respect to $T$ and $X$ is Hausdorff with respect to $T$.
    Let $p = X$.
    $p\notin X$ by \cref{set_notin_self}.
    Let $Y = X\union\{p\}$.
    Let $S = \onePointTopology{T}{X}{p}$.
    $S$ is a topology on $Y$ and $Y$ is compact with respect to $S$ and $Y$ is Hausdorff with respect to $S$
        by \cref{one_point_topology_is_topology,one_point_topology_compact,one_point_topology_hausdorff}.
    Let $A = X$.
    We have $A$ is open in $Y$ with respect to $S$ and $T = \subspaceTopology{S}{A}$
        by \cref{topology_on,subseteq,one_point_topology_intro_left,open_in,one_point_topology_subspace}.
    Let $f = \identity{A}$.
    $f$ is a bijection from $X$ to $A$ by \cref{id_elem_bijections}.
    $f$ is continuous from $X$ to $A$ with respect to $T$ and $\subspaceTopology{S}{A}$ by
        \cref{identity_continuous}.
    $\converse{f}$ is a function from $A$ to $X$ by \cref{bijective_converse_are_funs}.
    $A$ is a subspace of $Y$ with respect to $S$ and $\subspaceTopology{S}{A}$ is a topology on $A$
        by \cref{open_in,topology_on,subseteq,pow_iff,subspace_of,subspace_topology_is_topology}.
    For all $U$ such that $U$ is open in $X$ with respect to $T$ we have
        $\preimg{\converse{f}}{U}$ is open in $A$ with respect to $\subspaceTopology{S}{A}$
        by \cref{open_in,topology_on,subseteq,pow_iff,id_is_relation,preim_eq_img_of_converse,converse_converse_eq,id_img,inter_absorb_supseteq_left}.
    $\converse{f}$ is continuous from $A$ to $X$ with respect to $\subspaceTopology{S}{A}$ and $T$
        by \cref{continuous}.
    $f$ is a homeomorphism between $X$ and $A$ with respect to $T$ and $\subspaceTopology{S}{A}$
        by \cref{homeomorphism}.
    Show that there exist $Y_1,S_1,A_1,f_1$ such that
        $S_1$ is a topology on $Y_1$ and
        $Y_1$ is compact with respect to $S_1$ and
        $Y_1$ is Hausdorff with respect to $S_1$ and
        $A_1$ is open in $Y_1$ with respect to $S_1$ and
        $f_1$ is a homeomorphism between $X$ and $A_1$ with respect to $T$ and $\subspaceTopology{S_1}{A_1}$.
    \begin{subproof}
        Trivial.
    \end{subproof}
\end{proof}

% from §29 Item 12: open subspace of compact Hausdorff is locally compact Hausdorff
\begin{corollary}\label{open_subspace_compact_hausdorff_locally_compact_hausdorff}
    Assume $T$ is a topology on $X$.
    If there exist $Y,S,A,f$ such that
        $S$ is a topology on $Y$ and
        $Y$ is compact with respect to $S$ and
        $Y$ is Hausdorff with respect to $S$ and
        $A$ is open in $Y$ with respect to $S$ and
        $f$ is a homeomorphism between $X$ and $A$ with respect to $T$ and $\subspaceTopology{S}{A}$
    then $X$ is locally compact with respect to $T$ and $X$ is Hausdorff with respect to $T$.
\end{corollary}
\begin{proof}
    Assume there exist $Y,S,A,f$ such that
        $S$ is a topology on $Y$ and
        $Y$ is compact with respect to $S$ and
        $Y$ is Hausdorff with respect to $S$ and
        $A$ is open in $Y$ with respect to $S$ and
        $f$ is a homeomorphism between $X$ and $A$ with respect to $T$ and $\subspaceTopology{S}{A}$.
    Take $Y,S,A,f$ such that
        $S$ is a topology on $Y$ and
        $Y$ is compact with respect to $S$ and
        $Y$ is Hausdorff with respect to $S$ and
        $A$ is open in $Y$ with respect to $S$ and
    $f$ is a homeomorphism between $X$ and $A$ with respect to $T$ and $\subspaceTopology{S}{A}$.

    $Y$ is locally compact with respect to $S$ by \cref{compact_implies_locally_compact}.
    $A\subseteq Y$ by \cref{open_in,topology_on,subseteq,pow_iff}.
    Let $T_A = \subspaceTopology{S}{A}$.
    $A$ is a subspace of $Y$ with respect to $S$ and $T_A$ is a topology on $A$
        by \cref{subspace_of,subspace_topology_is_topology}.
    $A$ is Hausdorff with respect to $T_A$ by \cref{subspace_hausdorff}.
    $A$ is locally compact with respect to $T_A$ by \cref{locally_compact_subspace}.
    Show that $X$ is locally compact with respect to $T$.
    \begin{subproof}
        Follows by \cref{homeomorphism_preserves_local_compactness}.
    \end{subproof}

    Show that for all $x_1,x_2$ such that $x_1\in X$ and $x_2\in X$ and $x_1\neq x_2$ there exist $U_1,U_2$ such that
        $U_1$ is a neighborhood of $x_1$ in $X$ with respect to $T$ and
        $U_2$ is a neighborhood of $x_2$ in $X$ with respect to $T$ and
        $U_1$ is disjoint from $U_2$.
    \begin{subproof}
            Fix $x_1$.
            Fix $x_2$ such that $x_1\in X$ and $x_2\in X$ and $x_1\neq x_2$.
            $f\in\funs{X}{A}$ and $f$ is a function and $f$ is injective
                by \cref{bijections_to_funs,homeomorphism,bijection_is_function,bijections}.
            Let $y_1 = f(x_1)$.
            Let $y_2 = f(x_2)$.
            $y_1\in A$ and $y_2\in A$ by \cref{funs_type_apply,homeomorphism}.
            $y_1\neq y_2$ by \cref{funs,subseteq,injective_function}.
            Take $V_1,V_2$ such that
                $V_1$ is a neighborhood of $y_1$ in $A$ with respect to $T_A$ and
                $V_2$ is a neighborhood of $y_2$ in $A$ with respect to $T_A$ and
                $V_1$ is disjoint from $V_2$
                by \cref{hausdorff}.
            Let $U_1 = \preimg{f}{V_1}$.
            Let $U_2 = \preimg{f}{V_2}$.
            $U_1$ is a neighborhood of $x_1$ in $X$ with respect to $T$ and
                $U_2$ is a neighborhood of $x_2$ in $X$ with respect to $T$
                by \cref{neighborhood,continuous,homeomorphism,funs,subseteq,function_apply_elim,preimg_iff}.
            $U_1$ is disjoint from $U_2$
                by \cref{preimg_iff,funs_is_function,function_apply_iff,disjoint}.
            Show that there exist $W_1,W_2$ such that
                $W_1$ is a neighborhood of $x_1$ in $X$ with respect to $T$ and
                $W_2$ is a neighborhood of $x_2$ in $X$ with respect to $T$ and
                $W_1$ is disjoint from $W_2$.
            \begin{subproof}
                Trivial.
            \end{subproof}
    \end{subproof}
    Hence $X$ is Hausdorff with respect to $T$ by \cref{hausdorff}.
    Show that $X$ is locally compact with respect to $T$ and $X$ is Hausdorff with respect to $T$.
    \begin{subproof}
        Trivial.
    \end{subproof}
\end{proof}
\section{§30 The Countability Axioms}

% from §30 Item 1: sequence implies closure
\begin{lemma}\label{sequence_implies_closure_first_countable}
    Assume $T$ is a topology on $X$.
    Suppose $A\subseteq X$.
    Suppose $s$ is a sequence in $A$.
    Suppose $x\in X$.
    Suppose $s$ converges to $x$ in $X$ with respect to $T$.
    Then $x\in\closure{A}{X}{T}$.
\end{lemma}
\begin{proof}
    Follows by \cref{sequence_implies_closure}.
\end{proof}

% from §30 helper lemma 1: nested countable basis at a point
\begin{lemma}\label{nested_countable_basis}
    Suppose $b$ is a countable basis at $x$ in $X$ with respect to $T$.
    Then there exists $c$ such that
        $c$ is a countable basis at $x$ in $X$ with respect to $T$ and
        for all $m,n\in\naturalsPlus$ if $m\leq n$ then $c(n)\subseteq c(m)$.
\end{lemma}
\begin{proof}
    $1\in\naturalsPlus$ by \cref{real_one_in_naturals}.
    $b$ is a sequence in $\pow{X}$ by \cref{countable_basis_at_point}. $b(1)$ is a neighborhood of $x$ in $X$ with respect to $T$ by \cref{countable_basis_at_point}.
    $T$ is a topology on $X$ and $b(1)$ is open in $X$ with respect to $T$ by \cref{neighborhood,open_in}.
    Let $F = \{(n,\img{\restrl{b}{\initialSegment{n}}}{\initialSegment{n}}) \mid n\in\naturalsPlus\}$.
    Let $C = \{(n,\inters{\apply{F}{n}}) \mid n\in\naturalsPlus\}$.

    For all $n,U_1,U_2$ such that $(n,U_1)\in F$ and $(n,U_2)\in F$ we have $U_1 = U_2$
        by \cref{pair_eq_iff}.
    $F$ is right-unique by \cref{rightunique}.
    $F$ is a function.
    For all $n$ such that $n\in\naturalsPlus$ we have $n\in\dom{F}$
        by \cref{pair_eq_iff,dom_intro}.
    For all $n$ such that $n\in\dom{F}$ we have $n\in\naturalsPlus$
        by \cref{dom_iff,pair_eq_iff}.
    $F$ is a function on $\naturalsPlus$ by \cref{subseteq,subseteq_antisymmetric}.

    For all $n\in\naturalsPlus$ we have
        $F(n) = \img{\restrl{b}{\initialSegment{n}}}{\initialSegment{n}}$
        by \cref{pair_eq_iff,function_apply_intro}.

    For all $n,U_1,U_2$ such that $(n,U_1)\in C$ and $(n,U_2)\in C$ we have $U_1 = U_2$
        by \cref{pair_eq_iff}.
    $C$ is right-unique by \cref{rightunique}.
    $C$ is a function.
    For all $n$ such that $n\in\naturalsPlus$ we have $n\in\dom{C}$
        by \cref{pair_eq_iff,dom_intro}.
    For all $n$ such that $n\in\dom{C}$ we have $n\in\naturalsPlus$
        by \cref{dom_iff,pair_eq_iff}.
    $C$ is a function on $\naturalsPlus$ by \cref{subseteq,subseteq_antisymmetric}.

    For all $n\in\naturalsPlus$ we have $C(n) = \inters{F(n)}$
        by \cref{pair_eq_iff,function_apply_intro}.

    Show that for all $n\in\naturalsPlus$ we have $C(n)$ is a neighborhood of $x$ in $X$ with respect to $T$.
    \begin{subproof}
        Fix $n\in\naturalsPlus$.
        $\initialSegment{n}$ is finite by \cref{initial_segment_finite}.
        $b\in\funs{\naturalsPlus}{\pow{X}}$ and $\dom{b} = \naturalsPlus$
            by \cref{sequence_in,funs}.
        $\initialSegment{n}\subseteq\dom{b}$ by \cref{initial_segment_naturalsplus,subseteq}.
        Let $g = \restrl{b}{\initialSegment{n}}$.
        $g$ is a function by \cref{sequence_in,funs_is_function,restrl_of_function_is_function}.
        $\dom{g}\subseteq\initialSegment{n}$ by \cref{restrl_dom}.
        For all $i$ such that $i\in\initialSegment{n}$ we have $i\in\dom{g}$
            by \cref{initial_segment_naturalsplus,subseteq,dom_iff,restrl_iff,dom_intro}.
        $\initialSegment{n}\subseteq\dom{g}$ by \cref{subseteq}.
        $\dom{g} = \initialSegment{n}$ by \cref{subseteq_antisymmetric}.
        $g$ is a function on $\initialSegment{n}$.
        $F(n) = \img{g}{\initialSegment{n}}$.
        $F(n)$ is finite by \cref{finite_image_function}.
        $1\in\initialSegment{n}$ by \cref{real_one_in_naturals,real_one_leq_naturalsplus,initial_segment_naturalsplus}.
        Take $U_1$ such that $1\mathrel{b} U_1$ by \cref{dom_iff,funs}.
        $F(n)$ is inhabited by \cref{img_iff,restrl_iff,nonempty_inhabited}.
        For all $U$ such that $U\in F(n)$ we have $U\in T$
            by \cref{img_iff,restrl_iff,sequence_in,funs,initial_segment_naturalsplus,subseteq,funs_is_function,function_apply_iff,countable_basis_at_point,neighborhood,open_in}.
        $C(n)\in T$ by \cref{subseteq,finite_intersections_open_in_topology}.
        For all $U$ such that $U\in F(n)$ we have $x\in U$
            by \cref{img_iff,restrl_iff,sequence_in,funs,initial_segment_naturalsplus,subseteq,funs_is_function,function_apply_iff,countable_basis_at_point,neighborhood}.
        $x\in C(n)$ by \cref{inters_iff_forall}.
        $C(n)$ is a neighborhood of $x$ in $X$ with respect to $T$ by \cref{neighborhood,open_in}.
    \end{subproof}

    For all $n\in\naturalsPlus$ we have $C(n)\in\pow{X}$
        by \cref{subseteq,neighborhood,open_in,topology_on}.
    $C$ is a sequence in $\pow{X}$ by \cref{funs_intro,sequence_in}.
    For all $U$ such that $U$ is a neighborhood of $x$ in $X$ with respect to $T$
        there exists $n\in\naturalsPlus$ such that $b(n)\subseteq U$ by \cref{countable_basis_at_point}.
    For all $U$ such that $U$ is a neighborhood of $x$ in $X$ with respect to $T$
        there exists $n\in\naturalsPlus$ such that $C(n)\subseteq U$
        by \cref{initial_segment_naturalsplus,sequence_in,funs,subseteq,funs_is_function,function_apply_iff,restrl_iff,img_iff,inters_subseteq_elem,subseteq_transitive}.
    Show that $C$ is a countable basis at $x$ in $X$ with respect to $T$.
    \begin{subproof}
        Follows by \cref{countable_basis_at_point}.
    \end{subproof}

    Show that for all $m,n$ such that $m\in\naturalsPlus$ and $n\in\naturalsPlus$ and $m\leq n$ we have $C(n)\subseteq C(m)$.
        \begin{subproof}
            Fix $m$.
            Fix $n$ such that $m\in\naturalsPlus$ and $n\in\naturalsPlus$ and $m\leq n$.
            For all $z$ such that $z\in C(n)$ we have for all $U$ such that $U\in F(n)$ we have $z\in U$
                by \cref{inters_iff_forall}.
            For all $i$ such that $i\in\initialSegment{m}$ we have $i\in\initialSegment{n}$
                by \cref{initial_segment_naturalsplus,real_naturals_subset_reals,subseteq,real_leq_trans}.
            $\initialSegment{m}\subseteq\initialSegment{n}$ by \cref{subseteq}.
            For all $y$ such that $y\in\img{\restrl{b}{\initialSegment{m}}}{\initialSegment{m}}$ we have
                $y\in\img{\restrl{b}{\initialSegment{n}}}{\initialSegment{n}}$
                by \cref{img_iff,restrl_iff,subseteq}.
            $\img{\restrl{b}{\initialSegment{m}}}{\initialSegment{m}}\subseteq
                \img{\restrl{b}{\initialSegment{n}}}{\initialSegment{n}}$ by \cref{subseteq}.
            For all $U$ such that $U\in F(m)$ we have $U\in F(n)$ by \cref{subseteq}.
            For all $z$ such that $z\in C(n)$ we have for all $U$ such that $U\in F(m)$ we have $z\in U$.
            $1\in\initialSegment{m}$ by \cref{real_one_in_naturals,real_one_leq_naturalsplus,initial_segment_naturalsplus}.
            $(1,b(1))\in\restrl{b}{\initialSegment{m}}$ by \cref{sequence_in,funs,subseteq,funs_is_function,function_apply_iff,restrl_iff}.
            $F(m)$ is inhabited by \cref{img_iff,nonempty_inhabited}.
            For all $z$ such that $z\in C(n)$ we have $z\in C(m)$ by \cref{inters_iff_forall}.
        $C(n)\subseteq C(m)$ by \cref{subseteq}.
    \end{subproof}

    Show that there exists $c$ such that
        $c$ is a countable basis at $x$ in $X$ with respect to $T$ and
        for all $m,n\in\naturalsPlus$ if $m\leq n$ then $c(n)\subseteq c(m)$.
    \begin{subproof}
        Trivial.
    \end{subproof}
\end{proof}

% from §30 Item 2: closure implies sequence in a first-countable space
\begin{lemma}\label{closure_implies_sequence_first_countable}
    Assume $T$ is a topology on $X$.
    Suppose $A\subseteq X$.
    Suppose $x\in X$.
    Suppose $X$ satisfies the first countability axiom with respect to $T$.
    Suppose $x\in\closure{A}{X}{T}$.
    Then there exists $s$ such that $s$ is a sequence in $A$ and $s$ converges to $x$ in $X$ with respect to $T$.
\end{lemma}
\begin{proof}
    Take $c$ such that
        $c$ is a countable basis at $x$ in $X$ with respect to $T$ and
        for all $m,n\in\naturalsPlus$ if $m\leq n$ then $c(n)\subseteq c(m)$
        by \cref{first_countability_axiom,nested_countable_basis}.

    Let $R = \{w\in\naturalsPlus\times X \mid \snd{w}\in A\inter \apply{c}{\fst{w}}\}$.
    $R$ is a relation by \cref{times_elem_is_tuple,relation}.
    For all $n\in\naturalsPlus$ there exists $y$ such that $n\mathrel{R} y$
        by \cref{countable_basis_at_point,closure_iff_intersects_open,neighborhood,intersects,nonempty_inhabited,inter_elim_right,inter_elim_left,subseteq,times_tuple_intro,fst_eq,snd_eq,inter_intro,pair_eq_iff}.
    Take $s$ such that $s$ is a function on $\naturalsPlus$ and
        for all $n\in\naturalsPlus$ we have $n\mathrel{R} s(n)$ by \cref{axiom_of_choice}.

    For all $n$ such that $n\in\dom{s}$ we have $s(n)\in A$
        by \cref{choice_function,subseteq,pair_eq_iff,fst_eq,snd_eq,inter}.
    Show that $s$ is a sequence in $A$.
    \begin{subproof}
        Follows by \cref{funs_intro,sequence_in}.
    \end{subproof}

    $s$ is a sequence in $X$ by \cref{sequence_in,funs_weaken_codom}.
    $1\in\naturalsPlus$ by \cref{real_one_in_naturals}.
    For all $U$ such that $U$ is a neighborhood of $x$ in $X$ with respect to $T$
        there exists $N\in\naturalsPlus$ such that $c(N)\subseteq U$
        by \cref{countable_basis_at_point}.
    For all $U$ such that $U$ is a neighborhood of $x$ in $X$ with respect to $T$
        there exists $N\in\naturalsPlus$ such that for all $n\in\naturalsPlus$ if $N\leq n$ then $s(n)\in U$
        by \cref{nested_countable_basis,subseteq_transitive,pair_eq_iff,fst_eq,snd_eq,inter,subseteq}.
    $s$ converges to $x$ in $X$ with respect to $T$ by \cref{converges_to}.

    Show that there exists $t$ such that $t$ is a sequence in $A$ and $t$ converges to $x$ in $X$ with respect to $T$.
    \begin{subproof}
        Trivial.
    \end{subproof}
\end{proof}

% from §30 Item 3: continuous implies sequential continuity
\begin{theorem}\label{continuous_implies_sequential_first_countable}
    Assume $T$ is a topology on $X$.
    Assume $T_1$ is a topology on $Y$.
    Suppose $f$ is a function from $X$ to $Y$.
    Suppose $f$ is continuous from $X$ to $Y$ with respect to $T$ and $T_1$.
    Then for all $s$ such that $s$ is a sequence in $X$ we have
        for all $x\in X$ if $s$ converges to $x$ in $X$ with respect to $T$ then
        $f\circ s$ converges to $f(x)$ in $Y$ with respect to $T_1$.
\end{theorem}
\begin{proof}
    Follows by \cref{continuous_implies_sequential}.
\end{proof}

% from §30 Item 4: sequential continuity implies continuity in a first-countable space
\begin{theorem}\label{sequential_implies_continuous_first_countable}
    Assume $T$ is a topology on $X$.
    Assume $T_1$ is a topology on $Y$.
    Suppose $f$ is a function from $X$ to $Y$.
    Suppose $X$ satisfies the first countability axiom with respect to $T$.
    Suppose for all $s$ such that $s$ is a sequence in $X$ we have
        for all $x\in X$ if $s$ converges to $x$ in $X$ with respect to $T$ then
        $f\circ s$ converges to $f(x)$ in $Y$ with respect to $T_1$.
    Then $f$ is continuous from $X$ to $Y$ with respect to $T$ and $T_1$.
\end{theorem}
\begin{proof}
    Show that for all $A$ such that $A\subseteq X$ we have
        $\img{f}{\closure{A}{X}{T}}\subseteq \closure{\img{f}{A}}{Y}{T_1}$.
    \begin{subproof}
        Fix $A$ such that $A\subseteq X$.
        Show that for all $y$ such that $y\in\img{f}{\closure{A}{X}{T}}$ we have
            $y\in\closure{\img{f}{A}}{Y}{T_1}$.
        \begin{subproof}
            Fix $y$ such that $y\in\img{f}{\closure{A}{X}{T}}$.
            Take $x\in\dom{f}\inter\closure{A}{X}{T}$ such that $y = f(x)$ by \cref{funs_is_function,img_of_function_elim}.
            Take $s$ such that $s$ is a sequence in $A$ and $s$ converges to $x$ in $X$ with respect to $T$ by
                \cref{inter_elim_right,closure,closure_implies_sequence_first_countable}.
            We have $x\in X$ and $s$ is a sequence in $X$ and $f\circ s$ converges to $f(x)$ in $Y$ with respect to $T_1$ by
                \cref{inter_elim_right,closure,sequence_in,funs_weaken_codom}.
            $f\circ s$ converges to $y$ in $Y$ with respect to $T_1$.
            $f\circ s\in\funs{\naturalsPlus}{Y}$ by \cref{funs_circ,sequence_in}.
            For all $n\in\dom{(f\circ s)}$ we have $(f\circ s)(n)\in\img{f}{A}$
                by \cref{funs,sequence_in,subseteq,inter_intro,img_of_function_intro,funs_elim,function_to_ran,subseteq_transitive,circ_apply}.
            $f\circ s$ is a sequence in $\img{f}{A}$ by \cref{funs,sequence_in,funs_elim,funs_intro}.
            $f(x)\in\closure{\img{f}{A}}{Y}{T_1}$ by
                \cref{img_subseteq_ran,funs_elim,function_to_ran,subseteq_transitive,funs_type_apply,sequence_implies_closure}.
            $y\in\closure{\img{f}{A}}{Y}{T_1}$.
        \end{subproof}
        Hence $\img{f}{\closure{A}{X}{T}}\subseteq \closure{\img{f}{A}}{Y}{T_1}$ by \cref{subseteq}.
    \end{subproof}
    For all $B$ such that $B$ is closed in $Y$ with respect to $T_1$
        we have $\preimg{f}{B}$ is closed in $X$ with respect to $T$ by
        \cref{closure_image_subset_implies_preimg_closed}.
    $f$ is continuous from $X$ to $Y$ with respect to $T$ and $T_1$ by
        \cref{preimg_closed_implies_continuous}.
\end{proof}
% from §30 Item 5: second countability axiom
\begin{definition}\label{second_countability_axiom}
    $X$ satisfies the second countability axiom with respect to $T$ iff
        there exists $b$ such that
            $b$ is a sequence in $\pow{X}$ and
            for all $n\in\naturalsPlus$ we have $b(n)$ is open in $X$ with respect to $T$ and
            for all $x\in X$ we have
                for all $U$ such that $U$ is open in $X$ with respect to $T$ and $x\in U$
                    there exists $n\in\naturalsPlus$ such that $x\in b(n)$ and $b(n)\subseteq U$.
\end{definition}

% from §30 Item 6: second countability implies first countability
\begin{lemma}\label{second_countability_implies_first}
    Assume $T$ is a topology on $X$.
    Suppose $X$ satisfies the second countability axiom with respect to $T$.
    Then $X$ satisfies the first countability axiom with respect to $T$.
\end{lemma}
\begin{proof}
    Take $b$ such that
        $b$ is a sequence in $\pow{X}$ and
        for all $n\in\naturalsPlus$ we have $b(n)$ is open in $X$ with respect to $T$ and
        for all $x\in X$ we have
            for all $U$ such that $U$ is open in $X$ with respect to $T$ and $x\in U$
                there exists $n\in\naturalsPlus$ such that $x\in b(n)$ and $b(n)\subseteq U$
        by \cref{second_countability_axiom}.
    Show that for all $x\in X$ there exists $c$ such that
        $c$ is a countable basis at $x$ in $X$ with respect to $T$.
    \begin{subproof}
        Fix $x\in X$.
        Let $R = \{w \mid \exists n\in\naturalsPlus. \exists y\in\pow{X}.
            w = (n,y) \land ((x\in b(n)\land y = b(n)) \lor (x\notin b(n)\land y = X))\}$.
        $R$ is a relation by \cref{pair_eq_iff,relation}.
        Show that for all $n\in\naturalsPlus$ there exists $y$ such that $n\mathrel{R} y$.
        \begin{subproof}
            Fix $n\in\naturalsPlus$.
            \begin{byCase}
                \caseOf{$x\in b(n)$.}
                    Let $y = b(n)$.
                    $y\in\pow{X}$ by \cref{sequence_in,funs_type_apply}.
                    $(n,y)\in R$ and $n\mathrel{R} y$ by \cref{pair_eq_iff}.
                \caseOf{$x\notin b(n)$.}
                    Let $y = X$.
                    $y\in\pow{X}$ by \cref{subseteq_refl,powerset_intro}.
                    $(n,y)\in R$ and $n\mathrel{R} y$ by \cref{pair_eq_iff}.
            \end{byCase}
        \end{subproof}
        Take $c$ such that $c$ is a function on $\naturalsPlus$ and for all $n\in\naturalsPlus$ we have
            $n\mathrel{R} c(n)$ by \cref{axiom_of_choice}.
        For all $n\in\dom{c}$ we have $(x\in b(n)\land c(n) = b(n))\lor(x\notin b(n)\land c(n) = X)$
            by \cref{subseteq,choice_function,pair_eq_iff}.
        For all $n\in\dom{c}$ we have $c(n)\in\pow{X}$ by \cref{sequence_in,funs_type_apply,subseteq_refl,powerset_intro}.
        $c$ is a sequence in $\pow{X}$ by \cref{funs_intro,sequence_in}.
        For all $n\in\naturalsPlus$ we have $(x\in b(n)\land c(n) = b(n))\lor(x\notin b(n)\land c(n) = X)$
            by \cref{choice_function,pair_eq_iff}.
        Show that for all $n\in\naturalsPlus$ we have $c(n)$ is a neighborhood of $x$ in $X$ with respect to $T$.
        \begin{subproof}
            Follows by \cref{neighborhood,topology_on,open_in}.
        \end{subproof}
        For all $U$ such that $U$ is a neighborhood of $x$ in $X$ with respect to $T$
            we have $U$ is open in $X$ with respect to $T$ and $x\in U$ by \cref{neighborhood}.
        For all $U$ such that $U$ is a neighborhood of $x$ in $X$ with respect to $T$
            there exists $n\in\naturalsPlus$ such that $x\in b(n)$ and $b(n)\subseteq U$.
        For all $U$ such that $U$ is a neighborhood of $x$ in $X$ with respect to $T$
            there exists $n\in\naturalsPlus$ such that $c(n)\subseteq U$
            by \cref{choice_function,pair_eq_iff,subseteq}.
        $c$ is a countable basis at $x$ in $X$ with respect to $T$ by \cref{countable_basis_at_point}.
    \end{subproof}
    Hence $X$ satisfies the first countability axiom with respect to $T$ by \cref{first_countability_axiom}.
\end{proof}
% from §30 Item 7: subspace of a first-countable space is first-countable
\begin{lemma}\label{subspace_first_countable}
    Assume $T$ is a topology on $X$.
    Suppose $A\subseteq X$.
    Suppose $X$ satisfies the first countability axiom with respect to $T$.
    Then $A$ satisfies the first countability axiom with respect to $\subspaceTopology{T}{A}$.
\end{lemma}
\begin{proof}
    Let $T_1 = \subspaceTopology{T}{A}$.
    $T_1$ is a topology on $A$ by \cref{subspace_topology_is_topology}.
    Show that for all $x\in A$ there exists $c$ such that
        $c$ is a countable basis at $x$ in $A$ with respect to $T_1$.
    \begin{subproof}
        Fix $x\in A$.
        Take $b$ such that $b$ is a countable basis at $x$ in $X$ with respect to $T$
            by \cref{subseteq,first_countability_axiom}.
        Let $c = \{(n, A\inter b(n)) \mid n\in\naturalsPlus\}$.
        For all $n,U_1,U_2$ such that $(n,U_1)\in c$ and $(n,U_2)\in c$ we have $U_1 = U_2$
            by \cref{pair_eq_iff}.
        $c$ is right-unique by \cref{rightunique}.
        $c$ is a function.
        For all $n\in\dom{c}$ we have $c(n)\in\pow{A}$
            by \cref{dom_iff,pair_eq_iff,inter_elim_left,subseteq,powerset_intro,function_apply_iff}.
        For all $n$ such that $n\in\naturalsPlus$ we have $n\in\dom{c}$
            by \cref{pair_eq_iff,dom_intro}.
        For all $n$ such that $n\in\dom{c}$ we have $n\in\naturalsPlus$ by \cref{dom_iff,pair_eq_iff}.
        $\dom{c} = \naturalsPlus$ by \cref{subseteq,subseteq_antisymmetric}.
        $c$ is a sequence in $\pow{A}$ by \cref{funs_intro,sequence_in}.
        Show that for all $n\in\naturalsPlus$ we have $c(n)$ is a neighborhood of $x$ in $A$ with respect to $T_1$.
        \begin{subproof}
            Follows by \cref{countable_basis_at_point,neighborhood,open_in,subspace_topology,inter_intro,sequence_in,funs_elim,subseteq,pair_eq_iff,function_apply_iff}.
        \end{subproof}
        Show that for all $U$ such that $U$ is a neighborhood of $x$ in $A$ with respect to $T_1$
            there exists $n\in\naturalsPlus$ such that $c(n)\subseteq U$.
        \begin{subproof}
            Fix $U$ such that $U$ is a neighborhood of $x$ in $A$ with respect to $T_1$.
            $U\in T_1$ and $x\in U$ by \cref{neighborhood,open_in}.
            Take $V\in T$ such that $U = A\inter V$ by \cref{subspace_topology}.
            $V$ is a neighborhood of $x$ in $X$ with respect to $T$ by \cref{inter_elim_right,open_in,neighborhood}.
            $1\in\naturalsPlus$ by \cref{real_one_in_naturals}.
            For all $U$ such that $U$ is a neighborhood of $x$ in $X$ with respect to $T$
                there exists $n\in\naturalsPlus$ such that $b(n)\subseteq U$
                by \cref{countable_basis_at_point}.
            Take $n\in\naturalsPlus$ such that $b(n)\subseteq V$.
            For all $z$ such that $z\in A\inter b(n)$ we have $z\in A\inter V$
                by \cref{inter_elim_left,inter_elim_right,subseteq,inter_intro}.
            $A\inter b(n)\subseteq A\inter V$ by \cref{subseteq}.
            $c(n)\subseteq A\inter V$ by \cref{sequence_in,funs_elim,subseteq,pair_eq_iff,function_apply_iff}.
            $c(n)\subseteq U$ by \cref{subseteq}.
        \end{subproof}
        $c$ is a countable basis at $x$ in $A$ with respect to $T_1$ by \cref{countable_basis_at_point}.
    \end{subproof}
    $A$ satisfies the first countability axiom with respect to $T_1$ by \cref{first_countability_axiom}.
\end{proof}
% from §30 helper lemma 2: basis elements are neighborhoods
\begin{lemma}\label{countable_basis_elem_neighborhood}
    Assume $b$ is a countable basis at $x$ in $X$ with respect to $T$.
    Assume $n\in\naturalsPlus$.
    Then $b(n)$ is a neighborhood of $x$ in $X$ with respect to $T$.
\end{lemma}
\begin{proof}
    Follows by \cref{countable_basis_at_point}.
\end{proof}
% from §30 Item 8: countable product of first-countable spaces is first-countable
\begin{lemma}\label{countable_product_first_countable}
    Assume $X$ is a function.
    Assume $T$ is a function.
    Assume $\dom{X} = \naturalsPlus$.
    Assume $\dom{T} = \naturalsPlus$.
    Assume that for all $n\in\naturalsPlus$ we have $\apply{T}{n}$ is a topology on $\apply{X}{n}$.
    Suppose for all $n\in\naturalsPlus$ we have $\apply{X}{n}$ satisfies the first countability axiom with respect to $\apply{T}{n}$.
    Then $\productFamily{X}$ satisfies the first countability axiom with respect to $\productTopologyIndexed{T}{X}$.
\end{lemma}
\begin{proof}
    Let $T_0 = \productTopologyIndexed{T}{X}$.
    Show that for all $x\in\productFamily{X}$ there exists $d$ such that
        $d$ is a countable basis at $x$ in $\productFamily{X}$ with respect to $T_0$.
    \begin{subproof}
        Fix $x\in\productFamily{X}$.
        Let $B_0 = \funs{\naturalsPlus}{\pow{\unions{\ran{X}}}}$.
        Let $R = \{w\in\naturalsPlus\times B_0 \mid \text{there exists $n\in\naturalsPlus$ such that there exists $b\in B_0$ such that
            $w = (n,b)$ and $b$ is a countable basis at $\apply{x}{n}$ in $\apply{X}{n}$ with respect to $\apply{T}{n}$}\}$.
        For all $w$ such that $w\in R$ there exist $n,b$ such that $w = (n,b)$.
        $R$ is a relation by \cref{relation}.
        Show that for all $n\in\naturalsPlus$ there exists $b$ such that $n\mathrel{R} b$.
        \begin{subproof}
            Fix $n\in\naturalsPlus$.
            $\apply{x}{n}\in\apply{X}{n}$ by \cref{indexed_product,funs_elim,subseteq}.
            Take $b$ such that $b$ is a countable basis at $\apply{x}{n}$ in $\apply{X}{n}$ with respect to $\apply{T}{n}$
                by \cref{first_countability_axiom}.
            $n\in\dom{X}$ by \cref{subseteq}.
            $b$ is a sequence in $\pow{\apply{X}{n}}$ by \cref{countable_basis_at_point}.
            $b$ is a function and $\dom{b} = \naturalsPlus$ by \cref{sequence_in,funs_elim}.
            For all $m\in\naturalsPlus$ we have $\apply{b}{m}\in\pow{\unions{\ran{X}}}$
                by \cref{countable_basis_at_point,sequence_in,funs_elim,pow_iff,function_apply_iff,ran_intro,subseteq_pow_unions,subseteq,subseteq_transitive,powerset_intro}.
            $b\in B_0$ by \cref{funs_intro}.
            $(n,b)\in\naturalsPlus\times B_0$ by \cref{times_tuple_intro}.
            $(n,b)\in R$.
            $n\mathrel{R} b$.
        \end{subproof}
        Take $b$ such that $b$ is a function on $\naturalsPlus$ and for all $n\in\naturalsPlus$ we have
            $n\mathrel{R}\apply{b}{n}$ by \cref{choice_function}.
        Let $R_1 = \{w\in R \mid \text{for all $m,k\in\naturalsPlus$ if $m\leq k$ then
            $\apply{\snd{w}}{k}\subseteq \apply{\snd{w}}{m}$}\}$.
        For all $w$ such that $w\in R_1$ there exist $n,c$ such that $w = (n,c)$.
        $R_1$ is a relation by \cref{relation}.
        Show that for all $n\in\naturalsPlus$ there exists $c$ such that $n\mathrel{R_1} c$.
        \begin{subproof}
            Fix $n\in\naturalsPlus$.
            Take $c$ such that
                $c$ is a countable basis at $\apply{x}{n}$ in $\apply{X}{n}$ with respect to $\apply{T}{n}$ and
                for all $m,k\in\naturalsPlus$ if $m\leq k$ then $c(k)\subseteq c(m)$
                by \cref{nested_countable_basis,choice_function,pair_eq_iff}.
            $n\in\dom{X}$ by \cref{subseteq}.
            $c$ is a sequence in $\pow{\apply{X}{n}}$ by \cref{countable_basis_at_point}.
            $c$ is a function and $\dom{c} = \naturalsPlus$ by \cref{sequence_in,funs_elim}.
            For all $m\in\naturalsPlus$ we have $\apply{c}{m}\in\pow{\unions{\ran{X}}}$
                by \cref{countable_basis_at_point,sequence_in,funs_elim,pow_iff,function_apply_iff,ran_intro,subseteq_pow_unions,subseteq,subseteq_transitive,powerset_intro}.
            $c\in B_0$ by \cref{funs_intro}.
            $(n,c)\in\naturalsPlus\times B_0$ by \cref{times_tuple_intro}.
            $(n,c)\in R$.
            $(n,c)\in R_1$ by \cref{snd_eq}.
            $n\mathrel{R_1} c$.
        \end{subproof}
        Take $c$ such that $c$ is a function on $\naturalsPlus$ and for all $n\in\naturalsPlus$ we have
            $n\mathrel{R_1}c(n)$ by \cref{choice_function}.
        For all $n\in\naturalsPlus$ we have $c(n)$ is a countable basis at
            $\apply{x}{n}$ in $\apply{X}{n}$ with respect to $\apply{T}{n}$
            by \cref{choice_function,pair_eq_iff}.
        For all $n\in\naturalsPlus$ we have for all $m,k\in\naturalsPlus$ if $m\leq k$ then
            $\apply{c(n)}{k}\subseteq \apply{c(n)}{m}$ by \cref{choice_function,snd_eq}.

        Let $G_0 = \pow{\naturalsPlus\times\pow{\productFamily{X}}}$.
        Let $R_d = \{w\in\naturalsPlus\times G_0 \mid \text{there exists $n\in\naturalsPlus$ such that there exists $g\in G_0$ such that
            $w = (n,g)$ and $g$ is a function on $\initialSegment{n}$ and for all $i\in\initialSegment{n}$ we have
            $g(i) = \preimg{\piIndex{X}{i}}{\apply{c(i)}{n}}$}\}$.
        For all $w$ such that $w\in R_d$ there exist $n,g$ such that $w = (n,g)$.
        $R_d$ is a relation by \cref{relation}.
        Show that for all $n\in\naturalsPlus$ there exists $g$ such that $n\mathrel{R_d} g$.
        \begin{subproof}
            Fix $n\in\naturalsPlus$.
            Let $g = \{(i, \preimg{\piIndex{X}{i}}{\apply{c(i)}{n}}) \mid i\in\initialSegment{n}\}$.
            $g$ is a function by \cref{rightunique,relation,pair_eq_iff}.
            For all $i$ such that $i\in\initialSegment{n}$ we have $i\in\dom{g}$
                by \cref{pair_eq_iff,dom_intro}.
            $\initialSegment{n}\subseteq\dom{g}$ by \cref{subseteq}.
            For all $i$ such that $i\in\dom{g}$ we have $i\in\initialSegment{n}$ by \cref{dom_iff,pair_eq_iff}.
            $\dom{g}\subseteq\initialSegment{n}$ by \cref{subseteq}.
            $\dom{g} = \initialSegment{n}$ by \cref{subseteq_antisymmetric}.
            $g$ is a function on $\initialSegment{n}$.
            $g$ is right-unique and $g$ is a relation and $\dom{g} = \initialSegment{n}$
                by \cref{relation,function_rightunique}.
            Show that for all $w$ such that $w\in g$ we have $w\in\naturalsPlus\times\pow{\productFamily{X}}$.
            \begin{subproof}
                Fix $w$ such that $w\in g$.
                Take $i$ such that $i\in\initialSegment{n}$ and
                    $w = (i, \preimg{\piIndex{X}{i}}{\apply{c(i)}{n}})$ by \cref{pair_eq_iff}.
                $i\in\naturalsPlus$ by \cref{initial_segment_naturalsplus,subseteq}.
                $i\in\dom{X}$ by \cref{subseteq}.
                $c(i)$ is a countable basis at $\apply{x}{i}$ in $\apply{X}{i}$ with respect to $\apply{T}{i}$.
                $\apply{c(i)}{n}\in\apply{T}{i}$ by \cref{countable_basis_at_point,neighborhood,open_in}.
                $\preimg{\piIndex{X}{i}}{\apply{c(i)}{n}}\in\pow{\productFamily{X}}$ by
                    \cref{preimg_iff,projection_map_indexed,pair_eq_iff,subseteq,powerset_intro}.
                $w\in\naturalsPlus\times\pow{\productFamily{X}}$ by \cref{pair_eq_iff,times_tuple_intro}.
            \end{subproof}
            $g\subseteq\naturalsPlus\times\pow{\productFamily{X}}$ by \cref{subseteq}.
            $g\in G_0$ by \cref{powerset_intro}.
            $(n,g)\in\naturalsPlus\times G_0$ by \cref{times_tuple_intro}.
            For all $i\in\initialSegment{n}$ we have $\apply{g}{i} = \preimg{\piIndex{X}{i}}{\apply{c(i)}{n}}$
                by \cref{pair_eq_iff,dom_intro,function_apply_iff}.
            $(n,g)\in R_d$ by \cref{pair_eq_iff}.
            $n\mathrel{R_d} g$.
        \end{subproof}
        Take $h$ such that $h$ is a function on $\naturalsPlus$ and for all $n\in\naturalsPlus$ we have
            $n\mathrel{R_d}\apply{h}{n}$ by \cref{choice_function}.
        Let $d = \{(n, \inters{\img{\apply{h}{n}}{\initialSegment{n}}}) \mid n\in\naturalsPlus\}$.

        Show that $d$ is a countable basis at $x$ in $\productFamily{X}$ with respect to $T_0$.
        \begin{subproof}
            $d$ is a function by \cref{rightunique,relation,pair_eq_iff}.
            Show that for all $n\in\dom{d}$ we have $d(n)\in\pow{\productFamily{X}}$.
            \begin{subproof}
                        Fix $n\in\dom{d}$.
                        Take $U$ such that $(n,U)\in d$ by \cref{dom_iff}.
                        $U = \inters{\img{\apply{h}{n}}{\initialSegment{n}}}$ by \cref{pair_eq_iff}.
                        $n\in\naturalsPlus$ by \cref{pair_eq_iff}.
                        Let $F = \img{\apply{h}{n}}{\initialSegment{n}}$.
                        $n\mathrel{R_d}\apply{h}{n}$ by \cref{choice_function}.
                        $\apply{h}{n}$ is right-unique and $\apply{h}{n}$ is a relation and
                            $\initialSegment{n} = \dom{\apply{h}{n}}$ and
                            for all $j\in\initialSegment{n}$ we have
                            $\apply{\apply{h}{n}}{j} = \preimg{\piIndex{X}{j}}{\apply{c(j)}{n}}$ by \cref{pair_eq_iff}.
                        Show that for all $A$ such that $A\in F$ we have $A\in\productSubbasis{T}{X}$.
                        \begin{subproof}
                            Fix $A$ such that $A\in F$.
                            Take $i$ such that $i\in\initialSegment{n}$ and $(i,A)\in\apply{h}{n}$ by \cref{img_iff}.
                            $i\in\dom{\apply{h}{n}}$ by \cref{subseteq}.
                            $A = \apply{\apply{h}{n}}{i}$ by \cref{pair_eq_iff,function_apply_iff}.
                            $\apply{\apply{h}{n}}{i} = \preimg{\piIndex{X}{i}}{\apply{c(i)}{n}}$.
                            $A = \preimg{\piIndex{X}{i}}{\apply{c(i)}{n}}$.
                            $i\in\dom{X}$ by \cref{initial_segment_naturalsplus,subseteq}.
                            $c(i)$ is a countable basis at $\apply{x}{i}$ in $\apply{X}{i}$ with respect to $\apply{T}{i}$.
                            $\apply{c(i)}{n}\in\apply{T}{i}$ by \cref{countable_basis_at_point,neighborhood,open_in}.
                            For all $z$ such that $z\in A$ we have $z\in\productFamily{X}$
                                by \cref{preimg_iff,projection_map_indexed,pair_eq_iff}.
                            $A\subseteq\productFamily{X}$ by \cref{subseteq}.
                            $A\in\pow{\productFamily{X}}$ by \cref{powerset_intro}.
                            $A\in\productSubbasisAt{T}{X}{i}$ by \cref{product_subbasis_indexed}.
                            $A\in\productSubbasis{T}{X}$ by \cref{product_subbasis_union_indexed}.
                        \end{subproof}
                        $F\subseteq\productSubbasis{T}{X}$ by \cref{subseteq}.
                        $F$ is finite by \cref{initial_segment_finite,choice_function,pair_eq_iff,finite_image_function}.
                        $n\in\initialSegment{n}$ by \cref{initial_segment_naturalsplus}.
                        $\apply{h}{n}$ is a function and $\dom{\apply{h}{n}} = \initialSegment{n}$ by \cref{pair_eq_iff}.
                        $n\in\dom{\apply{h}{n}}$ by \cref{subseteq}.
                        For all $i\in\initialSegment{n}$ we have
                            $\apply{\apply{h}{n}}{i} = \preimg{\piIndex{X}{i}}{\apply{c(i)}{n}}$ by \cref{pair_eq_iff}.
                        $(n, \preimg{\piIndex{X}{n}}{\apply{c(n)}{n}})\in\apply{h}{n}$ by \cref{pair_eq_iff,function_apply_iff}.
                        $F$ is inhabited by \cref{img_iff,nonempty_inhabited}.
                        Take $A$ such that $A\in F$.
                        $A\in\productSubbasis{T}{X}$ by \cref{subseteq}.
                        For all $z$ such that $z\in A$ we have $z\in\unions{\productSubbasis{T}{X}}$ by \cref{unions_intro}.
                        $A\subseteq\unions{\productSubbasis{T}{X}}$ by \cref{subseteq}.
                        $U\subseteq A$ by \cref{inters_subseteq_elem}.
                        $U\subseteq\unions{\productSubbasis{T}{X}}$ by \cref{subseteq_transitive}.
                        $\unions{\productSubbasis{T}{X}} = \productFamily{X}$ by \cref{product_subbasis_is_subbasis,subbasis_on}.
                        $U\subseteq\productFamily{X}$ by \cref{subseteq}.
                        $U\in\pow{\productFamily{X}}$ by \cref{powerset_intro}.
                        $d(n)\in\pow{\productFamily{X}}$ by \cref{function_apply_iff}.
            \end{subproof}
            For all $n$ such that $n\in\naturalsPlus$ we have $n\in\dom{d}$
                by \cref{pair_eq_iff,dom_intro}.
            $\naturalsPlus\subseteq\dom{d}$ by \cref{subseteq}.
            For all $n$ such that $n\in\dom{d}$ we have $n\in\naturalsPlus$ by \cref{dom_iff,pair_eq_iff}.
            $\dom{d}\subseteq\naturalsPlus$ by \cref{subseteq}.
            $\dom{d} = \naturalsPlus$ by \cref{subseteq_antisymmetric}.
            $d$ is a sequence in $\pow{\productFamily{X}}$ by \cref{funs_intro,sequence_in}.

            Show that for all $n\in\naturalsPlus$ we have $d(n)$ is a neighborhood of $x$ in $\productFamily{X}$ with respect to $T_0$.
            \begin{subproof}
                Fix $n\in\naturalsPlus$.
                Let $F = \img{\apply{h}{n}}{\initialSegment{n}}$.
                $n\mathrel{R_d}\apply{h}{n}$ by \cref{choice_function}.
                $\apply{h}{n}$ is right-unique and $\apply{h}{n}$ is a relation and
                    $\initialSegment{n} = \dom{\apply{h}{n}}$ and
                    for all $j\in\initialSegment{n}$ we have
                    $\apply{\apply{h}{n}}{j} = \preimg{\piIndex{X}{j}}{\apply{c(j)}{n}}$ by \cref{pair_eq_iff}.
                Show that for all $A$ such that $A\in F$ we have $x\in A$.
                \begin{subproof}
                    Fix $A$ such that $A\in F$.
                    Take $i$ such that $i\in\initialSegment{n}$ and $(i,A)\in\apply{h}{n}$ by \cref{img_iff}.
                    $i\in\dom{\apply{h}{n}}$ by \cref{subseteq}.
                    $A = \apply{\apply{h}{n}}{i}$ by \cref{pair_eq_iff,function_apply_iff}.
                    $\apply{\apply{h}{n}}{i} = \preimg{\piIndex{X}{i}}{\apply{c(i)}{n}}$.
                    $A = \preimg{\piIndex{X}{i}}{\apply{c(i)}{n}}$.
                    $c(i)$ is a countable basis at $\apply{x}{i}$ in $\apply{X}{i}$ with respect to $\apply{T}{i}$.
                    $\apply{c(i)}{n}$ is a neighborhood of $\apply{x}{i}$ in $\apply{X}{i}$ with respect to $\apply{T}{i}$
                        by \cref{countable_basis_elem_neighborhood}.
                    $\apply{x}{i}\in\apply{c(i)}{n}$ by \cref{neighborhood}.
                    $x\in A$ by \cref{preimg_iff,projection_map_indexed}.
                \end{subproof}
                $n\in\initialSegment{n}$ by \cref{initial_segment_naturalsplus}.
                $n\mathrel{R_d}\apply{h}{n}$ by \cref{choice_function}.
                $\apply{h}{n}$ is a function and $\dom{\apply{h}{n}} = \initialSegment{n}$ by \cref{pair_eq_iff}.
                $n\in\dom{\apply{h}{n}}$ by \cref{subseteq}.
                For all $j\in\initialSegment{n}$ we have
                    $\apply{\apply{h}{n}}{j} = \preimg{\piIndex{X}{j}}{\apply{c(j)}{n}}$ by \cref{pair_eq_iff}.
                $(n, \preimg{\piIndex{X}{n}}{\apply{c(n)}{n}})\in\apply{h}{n}$ by \cref{pair_eq_iff,function_apply_iff}.
                $\preimg{\piIndex{X}{n}}{\apply{c(n)}{n}}\in F$ by \cref{img_iff}.
                $F$ is inhabited by \cref{nonempty_inhabited}.
                $x\in\inters{F}$ by \cref{inters_intro}.
                Show that for all $A$ such that $A\in F$ we have $A\in\productSubbasis{T}{X}$.
                \begin{subproof}
                    Fix $A$ such that $A\in F$.
                    Take $i$ such that $i\in\initialSegment{n}$ and $(i,A)\in\apply{h}{n}$ by \cref{img_iff}.
                    $i\in\dom{\apply{h}{n}}$ by \cref{subseteq}.
                    $A = \apply{\apply{h}{n}}{i}$ by \cref{pair_eq_iff,function_apply_iff}.
                    $\apply{\apply{h}{n}}{i} = \preimg{\piIndex{X}{i}}{\apply{c(i)}{n}}$.
                    $A = \preimg{\piIndex{X}{i}}{\apply{c(i)}{n}}$.
                    $i\in\dom{X}$ by \cref{initial_segment_naturalsplus,subseteq}.
                    $c(i)$ is a countable basis at $\apply{x}{i}$ in $\apply{X}{i}$ with respect to $\apply{T}{i}$.
                    $\apply{c(i)}{n}\in\apply{T}{i}$ by \cref{countable_basis_elem_neighborhood,neighborhood,open_in}.
                    For all $z$ such that $z\in A$ we have $z\in\productFamily{X}$
                        by \cref{preimg_iff,projection_map_indexed,pair_eq_iff}.
                    $A\subseteq\productFamily{X}$ by \cref{subseteq}.
                    $A\in\pow{\productFamily{X}}$ by \cref{powerset_intro}.
                    $A\in\productSubbasisAt{T}{X}{i}$ by \cref{product_subbasis_indexed}.
                    $A\in\productSubbasis{T}{X}$ by \cref{product_subbasis_union_indexed}.
                \end{subproof}
                $F\subseteq\productSubbasis{T}{X}$ by \cref{subseteq}.
                $F$ is finite by \cref{initial_segment_finite,choice_function,pair_eq_iff,finite_image_function}.
                $n\in\initialSegment{n}$ by \cref{initial_segment_naturalsplus}.
                $\apply{h}{n}$ is a function and $\dom{\apply{h}{n}} = \initialSegment{n}$ by \cref{pair_eq_iff}.
                $n\in\dom{\apply{h}{n}}$ by \cref{subseteq}.
                For all $j\in\initialSegment{n}$ we have
                    $\apply{\apply{h}{n}}{j} = \preimg{\piIndex{X}{j}}{\apply{c(j)}{n}}$ by \cref{pair_eq_iff}.
                $(n, \preimg{\piIndex{X}{n}}{\apply{c(n)}{n}})\in\apply{h}{n}$ by \cref{pair_eq_iff,function_apply_iff}.
                $F$ is inhabited by \cref{img_iff,nonempty_inhabited}.
                Take $A$ such that $A\in F$.
                $A\in\productSubbasis{T}{X}$ by \cref{subseteq}.
                For all $z$ such that $z\in A$ we have $z\in\unions{\productSubbasis{T}{X}}$ by \cref{unions_intro}.
                $A\subseteq\unions{\productSubbasis{T}{X}}$ by \cref{subseteq}.
                $\inters{F}\subseteq A$ by \cref{inters_subseteq_elem}.
                $\inters{F}\subseteq\unions{\productSubbasis{T}{X}}$ by \cref{subseteq_transitive}.
                $\inters{F}\in\pow{\unions{\productSubbasis{T}{X}}}$ by \cref{powerset_intro}.
                $\inters{F}\in\finiteIntersections{\productSubbasis{T}{X}}$ by \cref{finite_intersections}.
                $\finiteIntersections{\productSubbasis{T}{X}}$ is a basis for $T_0$ on $\productFamily{X}$ by
                    \cref{product_subbasis_finite_intersections_basis,product_topology_indexed,basis_for_topology,topology_generated_by_subbasis}.
                $T_0 = \basisTopology{\finiteIntersections{\productSubbasis{T}{X}}}{\productFamily{X}}$ by \cref{basis_for_topology}.
                $\inters{F}\in\basisTopology{\finiteIntersections{\productSubbasis{T}{X}}}{\productFamily{X}}$
                    by \cref{basis_elem_open_in_generated,basis_for_topology}.
                $\inters{F}\in T_0$ by \cref{basis_for_topology}.
                $\basisTopology{\finiteIntersections{\productSubbasis{T}{X}}}{\productFamily{X}}$ is a topology on $\productFamily{X}$
                    by \cref{basis_topology_is_topology,basis_for_topology}.
                $T_0$ is a topology on $\productFamily{X}$ by \cref{basis_for_topology}.
                $\inters{F}$ is open in $\productFamily{X}$ with respect to $T_0$ by \cref{open_in}.
                $\inters{F}$ is a neighborhood of $x$ in $\productFamily{X}$ with respect to $T_0$ by \cref{neighborhood}.
                $n\in\dom{d}$ by \cref{sequence_in,funs_elim,subseteq}.
                $(n, \inters{F})\in d$ by \cref{pair_eq_iff}.
                $d(n) = \inters{F}$ by \cref{function_apply_iff}.
                $d(n)$ is a neighborhood of $x$ in $\productFamily{X}$ with respect to $T_0$.
            \end{subproof}

            Show that for all $U$ such that $U$ is a neighborhood of $x$ in $\productFamily{X}$ with respect to $T_0$
                there exists $n\in\naturalsPlus$ such that $d(n)\subseteq U$.
            \begin{subproof}
                Fix $U$ such that $U$ is a neighborhood of $x$ in $\productFamily{X}$ with respect to $T_0$.
                $U$ is open in $\productFamily{X}$ with respect to $T_0$ and $x\in U$ by \cref{neighborhood}.
                $1\in\naturalsPlus$ by \cref{real_one_in_naturals}.
                $1\in\dom{X}$ by \cref{subseteq}.
                $\dom{X}$ is inhabited.
                For all $n\in\dom{X}$ we have $\apply{T}{n}$ is a topology on $\apply{X}{n}$.
                $\finiteIntersections{\productSubbasis{T}{X}}$ is a basis for $\productTopologyIndexed{T}{X}$ on $\productFamily{X}$
                    by \cref{product_subbasis_finite_intersections_basis}.
                $T_0 = \productTopologyIndexed{T}{X}$.
                $\finiteIntersections{\productSubbasis{T}{X}}$ is a basis for $T_0$ on $\productFamily{X}$.
                $U\in T_0$ by \cref{open_in}.
                $T_0 = \basisTopology{\finiteIntersections{\productSubbasis{T}{X}}}{\productFamily{X}}$
                    by \cref{basis_for_topology}.
                $U\in\basisTopology{\finiteIntersections{\productSubbasis{T}{X}}}{\productFamily{X}}$.
                Take $B\in\finiteIntersections{\productSubbasis{T}{X}}$ such that $x\in B$ and $B\subseteq U$
                    by \cref{topology_generated_by_basis}.
                Take $F\subseteq\productSubbasis{T}{X}$ such that $F$ is finite and inhabited and $B = \inters{F}$ by \cref{finite_intersections}.

                Let $R_2 = \{w\in F\times\dom{X} \mid \text{there exists $A\in F$ such that there exists $i\in\dom{X}$ such that
                    $w = (A,i)$ and there exists $V\in\apply{T}{i}$ such that
                    $A = \preimg{\piIndex{X}{i}}{V}$}\}$.
                For all $w$ such that $w\in R_2$ there exist $A,i$ such that $w = (A,i)$.
                $R_2$ is a relation by \cref{relation}.
                Show that for all $A\in F$ there exists $i$ such that $A\mathrel{R_2} i$.
                \begin{subproof}
                    Fix $A\in F$.
                    Take $i$ such that $i\in\dom{X}$ and $A\in\productSubbasisAt{T}{X}{i}$
                        by \cref{subseteq,product_subbasis_union_indexed}.
                    Take $V\in\apply{T}{i}$ such that $A = \preimg{\piIndex{X}{i}}{V}$ by \cref{product_subbasis_indexed}.
                    $(A,i)\in F\times\dom{X}$ by \cref{times_tuple_intro}.
                    $(A,i)\in R_2$ by definition.
                    $A\mathrel{R_2} i$.
                \end{subproof}
                Take $p$ such that $p$ is a function on $F$ and for all $A\in F$ we have $A\mathrel{R_2}\apply{p}{A}$
                    by \cref{choice_function}.
                Let $J = \img{p}{F}$.
                $J$ is finite by \cref{choice_function,finite_image_function}.
                Take $A_0$ such that $A_0\in F$.
                $J$ is inhabited by \cref{choice_function,subseteq,function_apply_iff,img_iff,nonempty_inhabited}.

                Let $R_3 = \{w\in F\times\naturalsPlus \mid \text{there exists $A\in F$ such that there exists $k\in\naturalsPlus$ such that
                    $w = (A,k)$ and there exists $V\in\apply{T}{\apply{p}{A}}$ such that
                    $A = \preimg{\piIndex{X}{\apply{p}{A}}}{V}$ and
                    $\apply{c(\apply{p}{A})}{k}\subseteq V$}\}$.
                For all $w$ such that $w\in R_3$ there exist $A,k$ such that $w = (A,k)$.
                $R_3$ is a relation by \cref{relation}.
                Show that for all $A\in F$ there exists $k$ such that $A\mathrel{R_3} k$.
                \begin{subproof}
                    Fix $A\in F$.
                    $\apply{p}{A}\in\dom{X}$ and there exists $V\in\apply{T}{\apply{p}{A}}$ such that
                        $A = \preimg{\piIndex{X}{\apply{p}{A}}}{V}$ by \cref{choice_function,pair_eq_iff}.
                    Take $V\in\apply{T}{\apply{p}{A}}$ such that $A = \preimg{\piIndex{X}{\apply{p}{A}}}{V}$.
                    $B\subseteq A$ by \cref{inters_subseteq_elem}.
                    $x\in A$ and $x\in \preimg{\piIndex{X}{\apply{p}{A}}}{V}$ by \cref{subseteq}.
                    Take $y\in V$ such that $(x,y)\in\piIndex{X}{\apply{p}{A}}$ by \cref{preimg_iff}.
                    Take $z\in\productFamily{X}$ such that $(z,\apply{z}{\apply{p}{A}}) = (x,y)$ by \cref{projection_map_indexed}.
                    $\apply{x}{\apply{p}{A}}\in V$ by \cref{pair_eq_iff}.
                    $\dom{X} = \naturalsPlus$.
                    $\apply{p}{A}\in\naturalsPlus$ by \cref{subseteq}.
                    $\apply{T}{\apply{p}{A}}$ is a topology on $\apply{X}{\apply{p}{A}}$.
                    $V$ is a neighborhood of $\apply{x}{\apply{p}{A}}$ in $\apply{X}{\apply{p}{A}}$ with respect to $\apply{T}{\apply{p}{A}}$ by \cref{open_in,neighborhood}.
                    $c(\apply{p}{A})$ is a countable basis at $\apply{x}{\apply{p}{A}}$ in $\apply{X}{\apply{p}{A}}$ with respect to $\apply{T}{\apply{p}{A}}$.
                    Take $k\in\naturalsPlus$ such that $\apply{c(\apply{p}{A})}{k}\subseteq V$ by \cref{countable_basis_at_point}.
                    $(A,k)\in F\times\naturalsPlus$ by \cref{times_tuple_intro}.
                    There exists $V_0\in\apply{T}{\apply{p}{A}}$ such that
                        $A = \preimg{\piIndex{X}{\apply{p}{A}}}{V_0}$ and $\apply{c(\apply{p}{A})}{k}\subseteq V_0$.
                    $(A,k)\in R_3$ by definition.
                    $A\mathrel{R_3} k$.
                \end{subproof}
                Take $q$ such that $q$ is a function on $F$ and for all $A\in F$ we have $A\mathrel{R_3}\apply{q}{A}$
                    by \cref{choice_function}.
                Let $K = \img{q}{F}$.
                $K$ is finite by \cref{choice_function,finite_image_function}.
                Take $A$ such that $A\in F$.
                $K$ is inhabited by \cref{choice_function,subseteq,function_apply_iff,img_iff,nonempty_inhabited}.

                For all $j$ such that $j\in J$ we have $j\in\naturalsPlus$
                    by \cref{img_iff,choice_function,subseteq,function_apply_iff,pair_eq_iff}.
                $J\subseteq\naturalsPlus$ by \cref{subseteq}.

                For all $k$ such that $k\in K$ we have $k\in\naturalsPlus$
                    by \cref{img_iff,choice_function,subseteq,function_apply_iff,pair_eq_iff}.
                $K\subseteq\naturalsPlus$ by \cref{subseteq}.

                Let $G = J\union K$.
                $G\subseteq\naturalsPlus$ by \cref{union_subsets_is_subset}.
                $\naturalsPlus\subseteq\reals$ by \cref{real_naturals_subset_reals}.
                $G\subseteq\reals$ by \cref{subseteq_transitive}.
                $G$ is finite by \cref{union_of_finite_is_finite}.
                $G$ is inhabited.
                Take $m$ such that $m$ is a largest element of $G$ with respect to $\realOrderRel$
                    by \cref{finite_inhabited_largest_element,real_order_relation_simple}.
                $m\in G$ by \cref{largest_element}.
                $m\in\naturalsPlus$ by \cref{subseteq}.

                Show that for all $A\in F$ we have $d(m)\subseteq A$.
                \begin{subproof}
                    Fix $A_1\in F$.
                    $\apply{p}{A_1}\in J$ by \cref{choice_function,subseteq,function_apply_iff,img_iff}.
                    $\apply{p}{A_1}\in G$ by \cref{union_intro_left}.
                    $\apply{p}{A_1}\leq m$ by \cref{largest_element,real_order_relation_elim}.
                    $\apply{p}{A_1}\in\dom{X}$ by \cref{pair_eq_iff}.
                    $\dom{X} = \naturalsPlus$.
                    $\apply{p}{A_1}\in\naturalsPlus$.
                    $\apply{p}{A_1}\in\initialSegment{m}$ by \cref{initial_segment_naturalsplus}.

                    $\apply{q}{A_1}\in K$ by \cref{choice_function,subseteq,function_apply_iff,img_iff}.
                    $\apply{q}{A_1}\in G$ by \cref{union_intro_right}.
                    $\apply{q}{A_1}\leq m$ by \cref{largest_element,real_order_relation_elim}.

                    Let $i = \apply{p}{A_1}$.
                    $i\in\initialSegment{m}$.
                    $A_1\mathrel{R_2} i$ by \cref{choice_function}.
                    $i\in\dom{X}$ by \cref{pair_eq_iff}.
                    Let $k = \apply{q}{A_1}$.
                    $A_1\mathrel{R_3} k$ by \cref{choice_function}.
                    There exists $W_{17}\in\apply{T}{i}$ such that
                        $A_1 = \preimg{\piIndex{X}{i}}{W_{17}}$ and $k\in\naturalsPlus$ and $\apply{c(i)}{k}\subseteq W_{17}$
                        by \cref{pair_eq_iff}.
                    Take $Z_{17}\in\apply{T}{i}$ such that
                        $A_1 = \preimg{\piIndex{X}{i}}{Z_{17}}$ and $k\in\naturalsPlus$ and $\apply{c(i)}{k}\subseteq Z_{17}$.

                    $k\leq m$.
                    $i\mathrel{R_1}\apply{c}{i}$ by \cref{choice_function}.
                    For all $m_1,k_1\in\naturalsPlus$ if $m_1\leq k_1$ then
                        $\apply{c(i)}{k_1}\subseteq \apply{c(i)}{m_1}$ by \cref{pair_eq_iff}.
                    $\apply{c(i)}{m}\subseteq \apply{c(i)}{k}$.
                    $\apply{c(i)}{m}\subseteq Z_{17}$ by \cref{subseteq}.

                    Let $F_m = \img{\apply{h}{m}}{\initialSegment{m}}$.
                    $m\mathrel{R_d}\apply{h}{m}$ by \cref{choice_function}.
                    $\apply{h}{m}$ is right-unique and $\apply{h}{m}$ is a relation and
                        $\initialSegment{m} = \dom{\apply{h}{m}}$ and
                        for all $j\in\initialSegment{m}$ we have
                        $\apply{\apply{h}{m}}{j} = \preimg{\piIndex{X}{j}}{\apply{c(j)}{m}}$ by \cref{pair_eq_iff}.
                    $\apply{h}{m}$ is a function.
                    $i\in\dom{\apply{h}{m}}$ by \cref{subseteq}.
                    $\preimg{\piIndex{X}{i}}{\apply{c(i)}{m}}\in F_m$ by \cref{pair_eq_iff,function_apply_iff,img_iff}.
                    $m\in\dom{d}$ by \cref{sequence_in,funs_elim,subseteq}.
                    $(m, \inters{F_m})\in d$ by \cref{pair_eq_iff}.
                    $d(m) = \inters{F_m}$ by \cref{function_apply_iff}.
                    $d(m)\subseteq \preimg{\piIndex{X}{i}}{\apply{c(i)}{m}}$ by \cref{inters_subseteq_elem}.
                    $d(m)\subseteq A_1$ by \cref{subseteq_transitive,preimg_subseteq,subseteq}.
                \end{subproof}
                $d(m)\subseteq U$ by \cref{subseteq_inters_iff,subseteq,subseteq_transitive}.
            \end{subproof}
            Hence $d$ is a countable basis at $x$ in $\productFamily{X}$ with respect to $T_0$ by \cref{countable_basis_at_point}.
        \end{subproof}
    \end{subproof}
    Therefore $\productFamily{X}$ satisfies the first countability axiom with respect to $T_0$ by \cref{first_countability_axiom}.
\end{proof}
% from §30 Item 9: subspace of a second-countable space is second-countable
\begin{lemma}\label{subspace_second_countable}
    Assume $T$ is a topology on $X$.
    Suppose $A\subseteq X$.
    Suppose $X$ satisfies the second countability axiom with respect to $T$.
    Then $A$ satisfies the second countability axiom with respect to $\subspaceTopology{T}{A}$.
\end{lemma}
\begin{proof}
    Let $T_1 = \subspaceTopology{T}{A}$.
    $T_1$ is a topology on $A$ by \cref{subspace_topology_is_topology}.
    Take $b$ such that
        $b$ is a sequence in $\pow{X}$ and
        for all $n\in\naturalsPlus$ we have $b(n)$ is open in $X$ with respect to $T$ and
        for all $x\in X$ we have
            for all $U$ such that $U$ is open in $X$ with respect to $T$ and $x\in U$
                there exists $n\in\naturalsPlus$ such that $x\in b(n)$ and $b(n)\subseteq U$
        by \cref{second_countability_axiom}.
    Let $c = \{(n, A\inter b(n)) \mid n\in\naturalsPlus\}$.
    For all $n,U_1,U_2$ such that $(n,U_1)\in c$ and $(n,U_2)\in c$ we have $U_1 = U_2$ by \cref{pair_eq_iff}.
    $c$ is a function by \cref{rightunique,relation,pair_eq_iff}.
    $c\in\funs{\naturalsPlus}{\pow{A}}$
        by \cref{funs_intro,dom_iff,pair_eq_iff,inter_elim_left,subseteq,powerset_intro,function_apply_iff,dom_intro,subseteq_antisymmetric}.
    Show that $c$ is a sequence in $\pow{A}$.
    \begin{subproof}
        Follows by \cref{sequence_in}.
    \end{subproof}
    Show that for all $n\in\naturalsPlus$ we have $c(n)$ is open in $A$ with respect to $T_1$.
    \begin{subproof}
        Follows by \cref{sequence_in,funs_is_function,funs_elim,subseteq,pair_eq_iff,function_apply_iff,open_in,subspace_topology}.
    \end{subproof}
    Show that for all $x\in A$ we have
        for all $U$ such that $U$ is open in $A$ with respect to $T_1$ and $x\in U$
            there exists $n\in\naturalsPlus$ such that $x\in c(n)$ and $c(n)\subseteq U$.
    \begin{subproof}
        Fix $x\in A$.
        Fix $U$ such that $U$ is open in $A$ with respect to $T_1$ and $x\in U$.
        Take $V\in T$ such that $U = A\inter V$ by \cref{open_in,subspace_topology}.
        $V$ is open in $X$ with respect to $T$ and $x\in V$ by \cref{inter_elim_right,open_in}.
        $x\in X$ by \cref{subseteq}.
        Take $n\in\naturalsPlus$ such that $x\in b(n)$ and $b(n)\subseteq V$.
        $x\in c(n)$ and $c(n)\subseteq U$ by
            \cref{sequence_in,funs_is_function,funs_elim,subseteq,pair_eq_iff,function_apply_iff,inter_elim_left,inter_elim_right,inter_intro}.
    \end{subproof}
    $A$ satisfies the second countability axiom with respect to $T_1$ by \cref{second_countability_axiom}.
\end{proof}

% from §30 helper lemma 3: naturalsPlus closed under addition
\begin{lemma}\label{naturalsplus_add_closed}
    Suppose $a\in\naturalsPlus$ and $b\in\naturalsPlus$.
    Then $a + b\in\naturalsPlus$.
\end{lemma}
\begin{proof}
    Let $A = \{b\in\naturalsPlus \mid a + b\in\naturalsPlus\}$.
    $1\in A$ by \cref{real_one_in_naturals,real_naturals_closed_succ}.
    For all $n\in A$ we have $n + 1\in A$
        by \cref{real_naturals_closed_succ,real_naturals_subset_reals,subseteq,real_one_in_naturals,real_add_assoc}.
    $a + b\in\naturalsPlus$ by \cref{subseteq,real_naturals_induction}.
\end{proof}

% from §30 helper lemma 4: naturalsPlus closed under multiplication
\begin{lemma}\label{naturalsplus_mul_closed}
    Suppose $a\in\naturalsPlus$ and $b\in\naturalsPlus$.
    Then $a \rmul b\in\naturalsPlus$.
\end{lemma}
\begin{proof}
    Let $A = \{b\in\naturalsPlus \mid a \rmul b\in\naturalsPlus\}$.
    $1\in A$ by \cref{real_one_in_naturals,real_naturals_subset_reals,subseteq,real_mul_comm,real_mul_ident}.
    For all $n\in A$ we have $n + 1\in A$
        by \cref{naturalsplus_add_closed,real_naturals_subset_reals,subseteq,real_one_in_naturals,real_distrib,real_mul_comm,real_mul_ident}.
    $a \rmul b\in\naturalsPlus$ by \cref{subseteq,real_naturals_induction}.
\end{proof}

% from §30 helper lemma 5: successor product below next square
\begin{lemma}\label{naturalsplus_mul_succ_lt_square}
    Suppose $s\in\naturalsPlus$.
    Then $s \rmul (s + 1) < (s + 1) \rmul (s + 1)$.
\end{lemma}
\begin{proof}
    $s\in\reals$ and $1\in\reals$ and $0\in\reals$ and $0 < 1$ by \cref{real_naturals_subset_reals,real_one_in_naturals,subseteq,real_add_ident,real_zero_lt_one}.
    $s < s + 1$ and $s + 1\in\reals$
        by \cref{real_naturals_closed_succ,real_naturals_subset_reals,subseteq,real_lt_add,real_add_ident,real_add_comm}.
    $0 < s + 1$ by \cref{real_one_leq_naturalsplus,real_lt_trans,real_lt_subst_right,real_zero_lt_one}.
    $s \rmul (s + 1) < (s + 1) \rmul (s + 1)$ by \cref{real_lt_mul_pos}.
\end{proof}

% from §30 helper lemma 6: bound for square against successor product
\begin{lemma}\label{naturalsplus_mul_succ_lt_square_bound}
    Suppose $s\in\naturalsPlus$ and $t\in\naturalsPlus$ and $s < t$.
    Then $s \rmul (s + 1) < t \rmul t$.
\end{lemma}
\begin{proof}
    $t = 1$ or there exists $m\in\naturalsPlus$ such that $t = m + 1$ by \cref{real_naturals_predecessor}.
    \begin{byCase}
    \caseOf{$t = 1$.}
        Show that $s \rmul (s + 1) < t \rmul t$.
        \begin{subproof}
            Suppose not.
            $1 < s$ or $1 = s$ by \cref{real_one_leq_naturalsplus}.
            $s < 1$.
            $1\in\reals$ and $s\in\reals$ by \cref{real_naturals_subset_reals,real_one_in_naturals,subseteq}.
            $1 < 1$ by \cref{real_lt_trans,real_lt_subst_left}.
            Contradiction by \cref{real_lt_irreflexive}.
        \end{subproof}
    \caseOf{there exists $m\in\naturalsPlus$ such that $t = m + 1$.}
        Take $m\in\naturalsPlus$ such that $t = m + 1$.
        $s < m + 1$.
        $s \leq m + 1$.
        $s \neq m + 1$.
        $s \leq m$ by \cref{real_naturals_leq_succ_elim}.
        $s < m$ or $s = m$.
        \begin{byCase}
        \caseOf{$s = m$.}
            Show that $s \rmul (s + 1) < t \rmul t$.
            \begin{subproof}
                Follows by \cref{naturalsplus_mul_succ_lt_square}.
            \end{subproof}
        \caseOf{$s < m$.}
            $s\in\reals$ and $t\in\reals$ by \cref{real_naturals_subset_reals,subseteq}.
            $s + 1\in\naturalsPlus$ and $s + 1\in\reals$ by \cref{real_naturals_closed_succ,real_naturals_subset_reals,subseteq}.
            $1\in\reals$ and $0\in\reals$ and $0 < 1$ by \cref{real_naturals_subset_reals,real_one_in_naturals,subseteq,real_add_ident,real_zero_lt_one}.
            $1 < s + 1$ or $1 = s + 1$ by \cref{real_one_leq_naturalsplus}.
            $0 < s + 1$ by \cref{real_lt_trans,real_lt_subst_right}.
            $s \rmul (s + 1) < t \rmul (s + 1)$ by \cref{real_lt_mul_pos}.
            $s + 1 < m + 1$ by \cref{real_naturals_subset_reals,real_one_in_naturals,subseteq,real_add_ident,real_zero_lt_one,real_lt_add}.
            $s + 1 < t$.
            $0 < t$ by \cref{real_one_leq_naturalsplus,real_lt_trans,real_lt_subst_right,real_zero_lt_one}.
            $(s + 1) \rmul t < t \rmul t$ by \cref{real_lt_mul_pos}.
            $t \rmul (s + 1) < t \rmul t$ by \cref{real_mul_comm}.
            $s \rmul (s + 1)\in\reals$ and $t \rmul (s + 1)\in\reals$ and $t \rmul t\in\reals$ by \cref{real_mul_closed}.
            $s \rmul (s + 1) < t \rmul t$ by \cref{real_lt_trans}.
        \end{byCase}
    \end{byCase}
\end{proof}
% from §30 helper lemma 7: injective pairing of naturalsPlus
\begin{lemma}\label{naturalsplus_pairing_injection}
    There exists $f\in\inj{\naturalsPlus\times\naturalsPlus}{\naturalsPlus}$.
\end{lemma}
    \begin{proof}
        Let $g = \{(p,(\fst{p} + \snd{p}) \rmul (\fst{p} + \snd{p}) + \fst{p}) \mid p\in\naturalsPlus\times\naturalsPlus\}$.
        For all $p,a,b$ such that $(p,a)\in g$ and $(p,b)\in g$ we have $a = b$ by \cref{pair_eq_iff}.
        $g$ is a function by \cref{rightunique,relation,pair_eq_iff}.
        Show that for all $p\in\dom{g}$ we have $g(p)\in\naturalsPlus$.
        \begin{subproof}
            Fix $p\in\dom{g}$.
            Take $a$ such that $(p,a)\in g$ by \cref{dom_iff}.
            $p\in\naturalsPlus\times\naturalsPlus$ and
                $a = (\fst{p} + \snd{p}) \rmul (\fst{p} + \snd{p}) + \fst{p}$ by \cref{pair_eq_iff}.
            Take $i,j$ such that $i\in\naturalsPlus$ and $j\in\naturalsPlus$ and $p = (i,j)$ by \cref{times}.
            $\fst{p} = i$ and $\snd{p} = j$ by \cref{fst_eq,snd_eq}.
            $\fst{p} + \snd{p}\in\naturalsPlus$ by \cref{naturalsplus_add_closed}.
            $(\fst{p} + \snd{p}) \rmul (\fst{p} + \snd{p})\in\naturalsPlus$ by \cref{naturalsplus_mul_closed}.
            $a\in\naturalsPlus$ by \cref{naturalsplus_add_closed}.
            $g(p) = a$ by \cref{function_apply_iff}.
            $g(p)\in\naturalsPlus$.
        \end{subproof}
        For all $p$ such that $p\in\naturalsPlus\times\naturalsPlus$ we have $p\in\dom{g}$
            by \cref{times,pair_eq_iff,dom_intro}.
        $\dom{g} = \naturalsPlus\times\naturalsPlus$ by \cref{subseteq,dom_iff,pair_eq_iff,times_tuple_intro,dom_intro,subseteq_antisymmetric}.
        Show that $g\in\funs{\naturalsPlus\times\naturalsPlus}{\naturalsPlus}$.
        \begin{subproof}
            Follows by \cref{funs_intro}.
        \end{subproof}
        Show that for all $p_1,p_2$ such that $p_1\in\naturalsPlus\times\naturalsPlus$ and $p_2\in\naturalsPlus\times\naturalsPlus$ and $g(p_1) = g(p_2)$ we have $p_1 = p_2$.
        \begin{subproof}
            Fix $p_1$.
            Fix $p_2$ such that $p_1\in\naturalsPlus\times\naturalsPlus$ and $p_2\in\naturalsPlus\times\naturalsPlus$ and $g(p_1) = g(p_2)$.
            $p_1\in\dom{g}$ and $p_2\in\dom{g}$ by \cref{funs_elim,subseteq}.
            $(p_1,g(p_1))\in g$ and $(p_2,g(p_2))\in g$ by \cref{function_apply_iff}.
            Take $i_1,j_1$ such that $i_1\in\naturalsPlus$ and $j_1\in\naturalsPlus$ and $p_1 = (i_1,j_1)$ by \cref{times}.
            $g(p_1) = (\fst{p_1} + \snd{p_1}) \rmul (\fst{p_1} + \snd{p_1}) + \fst{p_1}$ by \cref{pair_eq_iff}.
            $g(p_1) = (i_1 + j_1) \rmul (i_1 + j_1) + i_1$ by \cref{fst_eq,snd_eq}.
            Take $i_2,j_2$ such that $i_2\in\naturalsPlus$ and $j_2\in\naturalsPlus$ and $p_2 = (i_2,j_2)$ by \cref{times}.
            $g(p_2) = (\fst{p_2} + \snd{p_2}) \rmul (\fst{p_2} + \snd{p_2}) + \fst{p_2}$ by \cref{pair_eq_iff}.
            $g(p_2) = (i_2 + j_2) \rmul (i_2 + j_2) + i_2$ by \cref{fst_eq,snd_eq}.
            $(i_1 + j_1) \rmul (i_1 + j_1) + i_1 = (i_2 + j_2) \rmul (i_2 + j_2) + i_2$.
            Let $s = i_1 + j_1$.
            Let $t = i_2 + j_2$.
            $s \rmul s + i_1 = t \rmul t + i_2$.
            \begin{byCase}
            \caseOf{$s = t$.}
            $i_1\in\reals$ and $j_1\in\reals$ and $i_2\in\reals$ and $j_2\in\reals$ and $s\in\reals$ and $t\in\reals$ by \cref{naturalsplus_add_closed,real_naturals_subset_reals,subseteq}.
            $t \rmul t\in\reals$ by \cref{real_mul_closed}.
            $i_1 + t \rmul t = i_2 + t \rmul t$ by \cref{real_add_comm}.
            $i_1 = i_2$ by \cref{real_add_cancel}.
            $i_1 + j_1 = i_2 + j_2$.
            $j_1 = j_2$ by \cref{real_add_comm,real_add_cancel}.
            $p_1 = p_2$ by \cref{pair_eq_iff}.
            \caseOf{$s < t$.}
                Suppose not.
                $j_1\in\naturalsPlus$ and $j_1\in\reals$ by \cref{subseteq,real_naturals_subset_reals}.
                $1\in\reals$ and $0\in\reals$ and $0 < 1$ by \cref{real_naturals_subset_reals,real_one_in_naturals,subseteq,real_add_ident,real_zero_lt_one}.
                $0 + j_1 < 1 + j_1$ by \cref{real_lt_add}.
                $0 + j_1 = j_1$ by \cref{real_add_ident}.
                $1 + j_1 = j_1 + 1$ by \cref{real_add_comm}.
                $0 < j_1$.
                $i_1\in\reals$ by \cref{real_naturals_subset_reals,subseteq}.
                $s\in\naturalsPlus$ and $s\in\reals$ by \cref{naturalsplus_add_closed,real_naturals_subset_reals,subseteq}.
                $i_1 < s$ by \cref{real_lt_add,real_add_ident,real_add_comm}.
                $s \rmul s\in\reals$ by \cref{real_mul_closed}.
                $s \rmul s + i_1 < s \rmul s + s$ by \cref{real_lt_add,real_add_comm}.
                $s \rmul s + i_1 < s \rmul (s + 1)$ by \cref{real_lt_subst_right,real_distrib,real_mul_comm,real_mul_ident}.
                $s \rmul s + i_1\in\reals$ and $s + 1\in\reals$ and $s \rmul (s + 1)\in\reals$ by \cref{real_add_closed,real_mul_closed}.
                $t\in\naturalsPlus$ and $t\in\reals$ and $t \rmul t\in\reals$
                    by \cref{naturalsplus_add_closed,real_naturals_subset_reals,subseteq,real_mul_closed}.
                $s \rmul (s + 1) < t \rmul t$ by \cref{naturalsplus_mul_succ_lt_square_bound}.
                $s \rmul s + i_1 < t \rmul t$ by \cref{real_lt_trans}.
                $i_2\in\reals$ by \cref{real_naturals_subset_reals,subseteq}.
                $0 < i_2$
                    by \cref{real_one_leq_naturalsplus,real_one_in_naturals,real_naturals_subset_reals,subseteq,real_lt_subst_right,real_lt_trans,real_zero_lt_one}.
                $t \rmul t < t \rmul t + i_2$ by \cref{real_add_ident,real_lt_add,real_lt_subst_left,real_add_comm}.
                $s \rmul s + i_1 < t \rmul t + i_2$ by \cref{real_lt_trans}.
                It is wrong that $s \rmul s + i_1 = t \rmul t + i_2$ by \cref{real_lt_irreflexive}.
                Contradiction.
            \caseOf{$t < s$.}
                Suppose not.
                $j_2\in\reals$ and $1\in\reals$ and $0\in\reals$ and $0 < 1$ by \cref{subseteq,real_naturals_subset_reals,real_one_in_naturals,real_add_ident,real_zero_lt_one}.
                $0 + j_2 < 1 + j_2$ by \cref{real_lt_add}.
                $0 + j_2 = j_2$ by \cref{real_add_ident}.
                $1 + j_2 = j_2 + 1$ by \cref{real_add_comm}.
                $0 < j_2$.
                $i_2\in\reals$ by \cref{real_naturals_subset_reals,subseteq}.
                $i_2 < t$ by \cref{real_lt_add,real_add_ident,real_add_comm}.
                $t\in\naturalsPlus$ and $t\in\reals$ and $t \rmul t\in\reals$
                    by \cref{naturalsplus_add_closed,real_naturals_subset_reals,subseteq,real_mul_closed}.
                $t \rmul t + i_2 < t \rmul t + t$ by \cref{real_lt_add,real_add_comm}.
                $t \rmul t + i_2 < t \rmul (t + 1)$ by \cref{real_lt_subst_right,real_distrib,real_mul_comm,real_mul_ident}.
                $t \rmul t + i_2\in\reals$ and $t + 1\in\reals$ and $t \rmul (t + 1)\in\reals$ by \cref{real_add_closed,real_mul_closed}.
                $s\in\naturalsPlus$ and $s\in\reals$ by \cref{naturalsplus_add_closed,real_naturals_subset_reals,subseteq}.
                $s \rmul s\in\reals$ by \cref{real_mul_closed}.
                $t \rmul (t + 1) < s \rmul s$ by \cref{naturalsplus_mul_succ_lt_square_bound}.
                $t \rmul t + i_2 < s \rmul s$ by \cref{real_lt_trans}.
                $i_1\in\reals$ by \cref{subseteq,real_naturals_subset_reals}.
                $0 < i_1$
                    by \cref{real_one_leq_naturalsplus,real_one_in_naturals,real_naturals_subset_reals,subseteq,real_lt_subst_right,real_lt_trans,real_zero_lt_one}.
                $s \rmul s < s \rmul s + i_1$ by \cref{real_add_ident,real_lt_add,real_lt_subst_left,real_add_comm}.
                $t \rmul t + i_2 < s \rmul s + i_1$ by \cref{real_lt_trans}.
                It is wrong that $s \rmul s + i_1 = t \rmul t + i_2$ by \cref{real_lt_irreflexive}.
                Contradiction.
            \end{byCase}
        \end{subproof}
        Show that there exists $f\in\inj{\naturalsPlus\times\naturalsPlus}{\naturalsPlus}$.
        \begin{subproof}
            Follows by \cref{inj}.
        \end{subproof}
    \end{proof}

% from §30 helper lemma 8: function equality by pointwise equality
\begin{lemma}\label{functions_eq_iff}
    Suppose $f\in\funs{X}{Y}$ and $g\in\funs{X}{Y}$.
    Suppose for all $x\in X$ we have $f(x) = g(x)$.
    Then $f = g$.
\end{lemma}
\begin{proof}
    Follows by \cref{funs_elim,fun_subseteq,subseteq_refl,subseteq_antisymmetric}.
\end{proof}

% from §30 helper lemma 8a: restriction of a sequence to a shorter initial segment
\begin{lemma}\label{restrl_initialsegment_funs}
    Suppose $n\in\naturalsPlus$ and $u\in\funs{\initialSegment{n + 1}}{\naturalsPlus}$.
    Then $\restrl{u}{\initialSegment{n}}\in\funs{\initialSegment{n}}{\naturalsPlus}$.
\end{lemma}
\begin{proof}
    $u$ is a function and $\dom{u} = \initialSegment{n + 1}$ by \cref{funs_elim}.
    $\restrl{u}{\initialSegment{n}}$ is a function by \cref{restrl_of_function_is_function}.
    $\dom{\restrl{u}{\initialSegment{n}}} = \initialSegment{n}$
        by \cref{initial_segment_subset_succ,subseteq,elem_dom_and_restr_implies_elem_of_restr,restrl_dom,subseteq_antisymmetric}.
    $\restrl{u}{\initialSegment{n}}\in\funs{\initialSegment{n}}{\naturalsPlus}$
        by \cref{elem_dom_of_restrl_implies_elem_dom_and_restr,funs_elim,initial_segment_subset_succ,subseteq,restrl_of_function_apply,funs_intro}.
\end{proof}

% from §30 helper lemma 9: injective coding of finite sequences of naturalsPlus
\begin{lemma}\label{naturalsplus_finite_sequences_inj}
    Suppose $n\in\naturalsPlus$.
    Then there exists $f\in\inj{\funs{\initialSegment{n}}{\naturalsPlus}}{\naturalsPlus}$.
\end{lemma}
\begin{proof}
    Let $A = \{n\in\naturalsPlus \mid \exists f\in\inj{\funs{\initialSegment{n}}{\naturalsPlus}}{\naturalsPlus}. f = f\}$.
    $1\in\naturalsPlus$ by \cref{real_one_in_naturals}.
    Let $f_1 = \{(u, \apply{u}{1}) \mid u\in\funs{\initialSegment{1}}{\naturalsPlus}\}$.
    $f_1$ is a function by \cref{rightunique,relation,pair_eq_iff}.
    $f_1\in\funs{\funs{\initialSegment{1}}{\naturalsPlus}}{\naturalsPlus}$
        by \cref{funs_intro,subseteq,dom_iff,pair_eq_iff,dom_intro,subseteq_antisymmetric,initial_segment_naturalsplus,real_one_in_naturals,funs_type_apply,function_apply_iff}.
    For all $u_1,u_2$ such that $u_1\in\funs{\initialSegment{1}}{\naturalsPlus}$ and $u_2\in\funs{\initialSegment{1}}{\naturalsPlus}$ and $f_1(u_1) = f_1(u_2)$
        we have for all $k\in\initialSegment{1}$ we have $u_1(k) = u_2(k)$
        by \cref{pair_eq_iff,funs_elim,function_apply_iff,initial_segment_naturalsplus,real_naturals_subset_reals,real_one_in_naturals,subseteq,real_one_leq_naturalsplus,real_leq_antisym}.
    For all $u_1,u_2$ such that $u_1\in\funs{\initialSegment{1}}{\naturalsPlus}$ and $u_2\in\funs{\initialSegment{1}}{\naturalsPlus}$ and $f_1(u_1) = f_1(u_2)$
        we have $u_1 = u_2$ by \cref{functions_eq_iff}.
    $f_1\in\inj{\funs{\initialSegment{1}}{\naturalsPlus}}{\naturalsPlus}$ by \cref{inj}.
        Show that $1\in A$.
        \begin{subproof}
            Trivial.
        \end{subproof}
        Show that for all $m\in A$ we have $m + 1\in A$.
        \begin{subproof}
            Fix $m\in A$.
            Take $f$ such that $f\in\inj{\funs{\initialSegment{m}}{\naturalsPlus}}{\naturalsPlus}$.
            Take $g$ such that $g\in\inj{\naturalsPlus\times\naturalsPlus}{\naturalsPlus}$ by \cref{naturalsplus_pairing_injection}.
            Let $F = \{(u, \apply{g}{(\apply{f}{\restrl{u}{\initialSegment{m}}}, \apply{u}{m + 1})}) \mid u\in\funs{\initialSegment{m + 1}}{\naturalsPlus}\}$.
            $F$ is a function by \cref{rightunique,relation,pair_eq_iff}.
            For all $u\in\dom{F}$ there exists $u_0\in\funs{\initialSegment{m + 1}}{\naturalsPlus}$ such that
                $(u,F(u)) = (u_0,\apply{g}{(\apply{f}{\restrl{u_0}{\initialSegment{m}}}, \apply{u_0}{m + 1})})$
                by \cref{dom_iff,pair_eq_iff,function_apply_iff}.
            For all $u\in\dom{F}$ we have $F(u)\in\naturalsPlus$
                by \cref{pair_eq_iff,funs_elim,funs_intro,restrl_initialsegment_funs,inj,funs_type_apply,function_apply_iff,real_naturals_closed_succ,initial_segment_naturalsplus,times_tuple_intro}.
            $F\in\funs{\funs{\initialSegment{m + 1}}{\naturalsPlus}}{\naturalsPlus}$
                by \cref{funs_intro,subseteq,dom_iff,pair_eq_iff,dom_intro,subseteq_antisymmetric}.
            Show that for all $u_1,u_2$ such that $u_1\in\funs{\initialSegment{m + 1}}{\naturalsPlus}$ and $u_2\in\funs{\initialSegment{m + 1}}{\naturalsPlus}$ and $F(u_1) = F(u_2)$ we have $u_1 = u_2$.
            \begin{subproof}
                    Fix $u_1$.
                    Fix $u_2$ such that $u_1\in\funs{\initialSegment{m + 1}}{\naturalsPlus}$ and $u_2\in\funs{\initialSegment{m + 1}}{\naturalsPlus}$ and $F(u_1) = F(u_2)$.
                    $\restrl{u_1}{\initialSegment{m}}\in\funs{\initialSegment{m}}{\naturalsPlus}$ and
                        $\restrl{u_2}{\initialSegment{m}}\in\funs{\initialSegment{m}}{\naturalsPlus}$
                        by \cref{funs_elim,restrl_initialsegment_funs}.
                    $(u_1,\apply{g}{(\apply{f}{\restrl{u_1}{\initialSegment{m}}}, \apply{u_1}{m + 1})})\in F$ and
                        $(u_2,\apply{g}{(\apply{f}{\restrl{u_2}{\initialSegment{m}}}, \apply{u_2}{m + 1})})\in F$ by \cref{pair_eq_iff}.
                    $\apply{g}{(\apply{f}{\restrl{u_1}{\initialSegment{m}}}, \apply{u_1}{m + 1})} = \apply{g}{(\apply{f}{\restrl{u_2}{\initialSegment{m}}}, \apply{u_2}{m + 1})}$ by \cref{funs_elim,function_apply_iff}.
                    $\apply{f}{\restrl{u_1}{\initialSegment{m}}}\in\naturalsPlus$ and
                        $\apply{f}{\restrl{u_2}{\initialSegment{m}}}\in\naturalsPlus$ by \cref{inj,funs_type_apply}.
                    $m + 1\in\naturalsPlus$ and $m + 1\in\initialSegment{m + 1}$ by \cref{real_naturals_closed_succ,initial_segment_naturalsplus}.
                    $\apply{u_1}{m + 1}\in\naturalsPlus$ and $\apply{u_2}{m + 1}\in\naturalsPlus$ by \cref{funs_type_apply}.
                    $(\apply{f}{\restrl{u_1}{\initialSegment{m}}}, \apply{u_1}{m + 1})\in\naturalsPlus\times\naturalsPlus$ and
                        $(\apply{f}{\restrl{u_2}{\initialSegment{m}}}, \apply{u_2}{m + 1})\in\naturalsPlus\times\naturalsPlus$ by \cref{times_tuple_intro}.
                    $(\apply{f}{\restrl{u_1}{\initialSegment{m}}}, \apply{u_1}{m + 1}) = (\apply{f}{\restrl{u_2}{\initialSegment{m}}}, \apply{u_2}{m + 1})$ by \cref{inj}.
                    $\apply{f}{\restrl{u_1}{\initialSegment{m}}} = \apply{f}{\restrl{u_2}{\initialSegment{m}}}$ and
                        $\apply{u_1}{m + 1} = \apply{u_2}{m + 1}$ by \cref{pair_eq_iff}.
                    $\restrl{u_1}{\initialSegment{m}} = \restrl{u_2}{\initialSegment{m}}$ by \cref{inj}.
                    $u_1$ is a function and $\dom{u_1} = \initialSegment{m + 1}$ and
                        $u_2$ is a function and $\dom{u_2} = \initialSegment{m + 1}$ by \cref{funs_elim}.
                    For all $k\in\initialSegment{m + 1}$ we have $u_1(k) = u_2(k)$
                        by \cref{real_naturals_leq_succ_elim,initial_segment_naturalsplus,initial_segment_subset_succ,funs_elim,subseteq,restrl_of_function_apply}.
                    $u_1 = u_2$ by \cref{functions_eq_iff}.
                \end{subproof}
            $F\in\inj{\funs{\initialSegment{m + 1}}{\naturalsPlus}}{\naturalsPlus}$ by \cref{inj}.
            $m + 1\in A$.
        \end{subproof}
    Show that there exists $f\in\inj{\funs{\initialSegment{n}}{\naturalsPlus}}{\naturalsPlus}$.
    \begin{subproof}
        Follows by \cref{subseteq,real_naturals_induction,pair_eq_iff}.
    \end{subproof}
\end{proof}

% from §30 helper lemma 10: injective coding of all finite sequences of naturalsPlus
\begin{lemma}\label{naturalsplus_finite_sequences_all_inj}
    Let $F = \{ \funs{\initialSegment{k}}{\naturalsPlus} \mid k\in\naturalsPlus \}$.
    Let $S = \unions{F}$.
    Then there exists $h\in\inj{S}{\naturalsPlus}$.
\end{lemma}
\begin{proof}
    Let $B = \pow{S\times\naturalsPlus}$.
    Let $R = \{w\in\naturalsPlus\times B \mid \text{there exists $k\in\naturalsPlus$ such that there exists
        $f\in\inj{\funs{\initialSegment{k}}{\naturalsPlus}}{\naturalsPlus}$ such that $w = (k,f)$}\}$.
    Show that for all $k\in\naturalsPlus$ there exists $f$ such that $k\mathrel{R} f$.
        \begin{subproof}
            Fix $k\in\naturalsPlus$.
        Take $f$ such that $f\in\inj{\funs{\initialSegment{k}}{\naturalsPlus}}{\naturalsPlus}$ by \cref{naturalsplus_finite_sequences_inj}.
        $f\in\funs{\funs{\initialSegment{k}}{\naturalsPlus}}{\naturalsPlus}$ by \cref{inj}.
        $f$ is a function and $\dom{f} = \funs{\initialSegment{k}}{\naturalsPlus}$ by \cref{funs_elim}.
        For all $w$ such that $w\in f$ we have $w\in S\times\naturalsPlus$
            by \cref{funs_elim,function_member_elim,pair_eq_iff,unions_intro,times_tuple_intro}.
        $f\subseteq S\times\naturalsPlus$ by \cref{subseteq}.
        $f\in B$ by \cref{powerset_intro}.
        $(k,f)\in R$.
        $k\mathrel{R} f$.
    \end{subproof}
    For all $w$ such that $w\in R$ there exist $k,f$ such that $w = (k,f)$.
    $R$ is a relation by \cref{relation}.
    Take $C$ such that $C$ is a function on $\naturalsPlus$ and for all $k\in\naturalsPlus$ we have $k\mathrel{R} \apply{C}{k}$
        by \cref{choice_function}.
    Take $g$ such that $g\in\inj{\naturalsPlus\times\naturalsPlus}{\naturalsPlus}$ by \cref{naturalsplus_pairing_injection}.
    Let $h = \{w\in S\times\naturalsPlus \mid \text{there exists $k\in\naturalsPlus$ such that there exists
        $u\in\funs{\initialSegment{k}}{\naturalsPlus}$ such that
        $w = (u, \apply{g}{(k, \apply{\apply{C}{k}}{u})})$}\}$.
    Show that $h\in\inj{S}{\naturalsPlus}$.
    \begin{subproof}
            Show that for all $u,a,b$ such that $(u,a)\in h$ and $(u,b)\in h$ we have $a = b$.
            \begin{subproof}
                Fix $u$.
                Fix $a$.
                Fix $b$ such that $(u,a)\in h$ and $(u,b)\in h$.
                Take $k_1\in\naturalsPlus$ such that $u\in\funs{\initialSegment{k_1}}{\naturalsPlus}$ and
                    $a = \apply{g}{(k_1, \apply{\apply{C}{k_1}}{u})}$ by \cref{pair_eq_iff}.
                Take $k_2\in\naturalsPlus$ such that $u\in\funs{\initialSegment{k_2}}{\naturalsPlus}$ and
                    $b = \apply{g}{(k_2, \apply{\apply{C}{k_2}}{u})}$ by \cref{pair_eq_iff}.
                $\initialSegment{k_1} = \initialSegment{k_2}$ by \cref{funs_elim,subseteq_antisymmetric,subseteq}.
                $k_1 \leq k_1$.
                $k_1\in\initialSegment{k_1}$ by \cref{initial_segment_naturalsplus}.
                $k_2 \leq k_2$.
                $k_2\in\initialSegment{k_2}$ by \cref{initial_segment_naturalsplus}.
                $k_1\in\initialSegment{k_2}$ and $k_2\in\initialSegment{k_1}$ by \cref{subseteq}.
                $k_1 \leq k_2$ and $k_2 \leq k_1$ by \cref{initial_segment_naturalsplus}.
                $k_1\in\reals$ and $k_2\in\reals$ by \cref{real_naturals_subset_reals,subseteq}.
                $k_1 = k_2$ by \cref{real_leq_antisym}.
                $a = b$.
            \end{subproof}
            Show that $h$ is a function.
            \begin{subproof}
                Follows by \cref{rightunique,relation,pair_eq_iff}.
            \end{subproof}
            Show that for all $u\in\dom{h}$ we have $h(u)\in\naturalsPlus$.
            \begin{subproof}
                Follows by \cref{dom_iff,pair_eq_iff,inj,funs_type_apply,function_apply_iff,times_tuple_intro}.
            \end{subproof}
            Show that for all $u$ such that $u\in\dom{h}$ we have $u\in S$.
            \begin{subproof}
                Follows by \cref{dom_iff,pair_eq_iff,unions_iff}.
            \end{subproof}
            Show that for all $u$ such that $u\in S$ we have $u\in\dom{h}$.
            \begin{subproof}
                Fix $u$ such that $u\in S$.
                Take $k\in\naturalsPlus$ such that $u\in\funs{\initialSegment{k}}{\naturalsPlus}$ by \cref{unions_iff,pair_eq_iff}.
                $\apply{C}{k}\in\inj{\funs{\initialSegment{k}}{\naturalsPlus}}{\naturalsPlus}$ by \cref{choice_function,pair_eq_iff}.
                $\apply{\apply{C}{k}}{u}\in\naturalsPlus$ by \cref{inj,funs_type_apply}.
                $\apply{g}{(k,\apply{\apply{C}{k}}{u})}\in\naturalsPlus$ by \cref{inj,funs_type_apply,times_tuple_intro}.
                $(u,\apply{g}{(k,\apply{\apply{C}{k}}{u})})\in S\times\naturalsPlus$ by \cref{times_tuple_intro}.
                $(u,\apply{g}{(k,\apply{\apply{C}{k}}{u})})\in h$ by \cref{pair_eq_iff}.
                $u\in\dom{h}$ by \cref{dom_intro}.
            \end{subproof}
            Show that $h\in\funs{S}{\naturalsPlus}$.
            \begin{subproof}
                Follows by \cref{funs_intro,dom_iff,pair_eq_iff,inj,funs_type_apply,function_apply_iff,times_tuple_intro,unions_iff,subseteq,dom_intro,subseteq_antisymmetric}.
            \end{subproof}
        Show that for all $u_1,u_2$ such that $u_1\in S$ and $u_2\in S$ and $h(u_1) = h(u_2)$ we have $u_1 = u_2$.
        \begin{subproof}
            Fix $u_1$.
            Fix $u_2$ such that $u_1\in S$ and $u_2\in S$ and $h(u_1) = h(u_2)$.
            $(u_1,\apply{h}{u_1})\in h$ and $(u_2,\apply{h}{u_2})\in h$ by \cref{funs_elim,subseteq,function_apply_iff}.
            Take $k_1\in\naturalsPlus$ such that $u_1\in\funs{\initialSegment{k_1}}{\naturalsPlus}$ and
                $\apply{h}{u_1} = \apply{g}{(k_1, \apply{\apply{C}{k_1}}{u_1})}$ by \cref{pair_eq_iff}.
            Take $k_2\in\naturalsPlus$ such that $u_2\in\funs{\initialSegment{k_2}}{\naturalsPlus}$ and
                $\apply{h}{u_2} = \apply{g}{(k_2, \apply{\apply{C}{k_2}}{u_2})}$ by \cref{pair_eq_iff}.
            $\apply{g}{(k_1, \apply{\apply{C}{k_1}}{u_1})} = \apply{g}{(k_2, \apply{\apply{C}{k_2}}{u_2})}$ by \cref{funs_elim,function_apply_iff}.
            $\apply{C}{k_1}\in\inj{\funs{\initialSegment{k_1}}{\naturalsPlus}}{\naturalsPlus}$ and
                $\apply{C}{k_2}\in\inj{\funs{\initialSegment{k_2}}{\naturalsPlus}}{\naturalsPlus}$ by \cref{pair_eq_iff}.
            $\apply{\apply{C}{k_1}}{u_1}\in\naturalsPlus$ and $\apply{\apply{C}{k_2}}{u_2}\in\naturalsPlus$ by \cref{inj,funs_type_apply}.
            $(k_1, \apply{\apply{C}{k_1}}{u_1})\in\naturalsPlus\times\naturalsPlus$ and
                $(k_2, \apply{\apply{C}{k_2}}{u_2})\in\naturalsPlus\times\naturalsPlus$ by \cref{times_tuple_intro}.
            $(k_1, \apply{\apply{C}{k_1}}{u_1}) = (k_2, \apply{\apply{C}{k_2}}{u_2})$ by \cref{inj}.
            $u_1 = u_2$ by \cref{pair_eq_iff,inj}.
        \end{subproof}
        $h\in\inj{S}{\naturalsPlus}$ by \cref{inj}.
    \end{subproof}
    Show that there exists $h_0\in\inj{S}{\naturalsPlus}$.
    \begin{subproof}
        Follows by \cref{inj}.
    \end{subproof}
\end{proof}

% from §30 helper lemma 11: enumerate finite sequences of naturalsPlus
\begin{lemma}\label{naturalsplus_finite_sequence_enumeration}
    Let $F = \{ \funs{\initialSegment{k}}{\naturalsPlus} \mid k\in\naturalsPlus \}$.
    Let $S = \unions{F}$.
    There exists $q$ such that
        $q$ is a function on $\naturalsPlus$ and
        for all $n\in\naturalsPlus$ we have $q(n)\in S$ and
        for all $u$ such that $u\in S$
            there exists $n\in\naturalsPlus$ such that $q(n) = u$.
\end{lemma}
\begin{proof}
    Take $h$ such that $h\in\inj{S}{\naturalsPlus}$ by \cref{naturalsplus_finite_sequences_all_inj}.
    Let $u_0 = \{(1,1)\}$.
    $1\in\naturalsPlus$ by \cref{real_one_in_naturals}.
        $u_0$ is a function by \cref{singleton_elim,pair_eq_iff,rightunique,relation}.
        For all $x$ such that $x\in\dom{u_0}$ we have $x\in\initialSegment{1}$
            by \cref{dom,singleton_elim,pair_eq_iff,real_one_in_naturals,initial_segment_naturalsplus}.
        For all $x$ such that $x\in\initialSegment{1}$ we have $x\in\dom{u_0}$
            by \cref{initial_segment_naturalsplus,real_naturals_subset_reals,real_one_in_naturals,subseteq,real_one_leq_naturalsplus,real_leq_antisym,pair_eq_iff,singleton_intro,dom}.
        $\dom{u_0} = \initialSegment{1}$ by \cref{subseteq,subseteq_antisymmetric}.
    For all $x\in\dom{u_0}$ we have $u_0(x)\in\naturalsPlus$
        by \cref{dom,singleton_elim,pair_eq_iff,function_apply_intro,real_one_in_naturals}.
    $u_0\in\funs{\initialSegment{1}}{\naturalsPlus}$ by \cref{funs_intro}.
    $u_0\in S$ by \cref{unions_iff}.
    Let $R_a = \{w\in\ran{h}\times S \mid \text{there exists $u\in S$ such that $w = (h(u),u)$}\}$.
    Let $R_b = (\naturalsPlus\setminus\ran{h})\times\{u_0\}$.
    Let $R_1 = R_a\union R_b$.
    For all $w$ such that $w\in R_a$ there exist $n,u$ such that $w = (n,u)$.
    $R_a$ is a relation by \cref{relation}.
    $R_b$ is a relation by \cref{relation,times_elem_is_tuple}.
    $R_1$ is a relation by \cref{union_relations_is_relation}.
    Show that for all $n\in\naturalsPlus$ there exists $u$ such that $n\mathrel{R_1} u$.
    \begin{subproof}
        Fix $n\in\naturalsPlus$.
        \begin{byCase}
        \caseOf{$n\in\ran{h}$.}
        Take $u$ such that $(u,n)\in h$ by \cref{ran_iff}.
        $u\in S$ and $h(u) = n$ by \cref{dom,inj,funs_elim,function_apply_intro}.
        $(n,u)\in\ran{h}\times S$ by \cref{times_tuple_intro}.
        $(n,u)\in R_a$ by \cref{pair_eq_iff}.
        $n\mathrel{R_1} u$ by \cref{union_intro_left}.
        \caseOf{$n\notin\ran{h}$.}
        Show that $n\mathrel{R_1} u_0$.
        \begin{subproof}
            Follows by \cref{setminus_intro,singleton_intro,times_tuple_intro,union_intro_right}.
        \end{subproof}
        \end{byCase}
    \end{subproof}
    Take $q$ such that $q$ is a function on $\naturalsPlus$ and for all $n\in\naturalsPlus$ we have $n\mathrel{R_1} \apply{q}{n}$
        by \cref{choice_function}.
    For all $n\in\naturalsPlus$ we have $q(n)\in S$
        by \cref{choice_function,union_iff,pair_eq_iff,times_tuple_elim,singleton_elim}.
    Show that for all $u\in S$ there exists $n\in\naturalsPlus$ such that $q(n) = u$.
    \begin{subproof}
        Fix $u\in S$.
        Let $n = h(u)$.
        $n\in\naturalsPlus$ and $n\in\ran{h}$
            by \cref{inj,funs_elim,funs_type_apply,function_apply_elim,ran_intro}.
        $(n,\apply{q}{n})\in R_a$ or $(n,\apply{q}{n})\in R_b$ by \cref{choice_function,union_iff}.
        $\apply{q}{n} = u$ by \cref{pair_eq_iff,inj,times_tuple_elim,setminus_elim_right}.
    \end{subproof}
    Let $q_0 = q$.
    Show that $q_0$ is a function on $\naturalsPlus$.
    \begin{subproof}
        Follows by \cref{choice_function}.
    \end{subproof}
    Show that for all $n\in\naturalsPlus$ we have $q_0(n)\in S$.
    \begin{subproof}
        Trivial.
    \end{subproof}
    Show that for all $u\in S$ there exists $n\in\naturalsPlus$ such that $q_0(n) = u$.
    \begin{subproof}
        Trivial.
    \end{subproof}
    Show that there exists $q_0$ such that
        $q_0$ is a function on $\naturalsPlus$ and
        for all $n\in\naturalsPlus$ we have $q_0(n)\in S$ and
        for all $u\in S$ there exists $n\in\naturalsPlus$ such that $q_0(n) = u$.
    \begin{subproof}
        Trivial.
    \end{subproof}
\end{proof}
% from §30 Item 10: countable product of second-countable spaces is second-countable
\begin{lemma}\label{countable_product_second_countable}
    Assume $X$ is a function.
    Assume $T$ is a function.
    Assume $\dom{X} = \naturalsPlus$.
    Assume $\dom{T} = \naturalsPlus$.
    Assume that for all $n\in\naturalsPlus$ we have $\apply{T}{n}$ is a topology on $\apply{X}{n}$.
    Suppose for all $n\in\naturalsPlus$ we have $\apply{X}{n}$ satisfies the second countability axiom with respect to $\apply{T}{n}$.
    Then $\productFamily{X}$ satisfies the second countability axiom with respect to $\productTopologyIndexed{T}{X}$.
\end{lemma}
\begin{proof}
    Let $T_0 = \productTopologyIndexed{T}{X}$.
    Let $R = \{(n,b) \mid n\in\naturalsPlus, b\in\funs{\naturalsPlus}{\pow{\unions{\ran{X}}}} \mid
        \text{$b$ is a sequence in $\pow{\apply{X}{n}}$ and
        for all $k\in\naturalsPlus$ we have $b(k)$ is open in $\apply{X}{n}$ with respect to $\apply{T}{n}$ and
        for all $x\in\apply{X}{n}$ we have
            for all $U$ such that $U$ is open in $\apply{X}{n}$ with respect to $\apply{T}{n}$ and $x\in U$
                there exists $k\in\naturalsPlus$ such that $x\in b(k)$ and $b(k)\subseteq U$}\}$.
    $R$ is a relation by \cref{pair_eq_iff,relation}.
    Show that for all $n\in\naturalsPlus$ there exists $b$ such that $n\mathrel{R} b$.
    \begin{subproof}
        Fix $n\in\naturalsPlus$.
        Take $b$ such that
            $b$ is a sequence in $\pow{\apply{X}{n}}$ and
            for all $k\in\naturalsPlus$ we have $b(k)$ is open in $\apply{X}{n}$ with respect to $\apply{T}{n}$ and
            for all $x\in\apply{X}{n}$ we have
                for all $U$ such that $U$ is open in $\apply{X}{n}$ with respect to $\apply{T}{n}$ and $x\in U$
                    there exists $k\in\naturalsPlus$ such that $x\in b(k)$ and $b(k)\subseteq U$
            by \cref{second_countability_axiom}.
        $b$ is a function and $\dom{b} = \naturalsPlus$ by \cref{sequence_in,funs_elim}.
        For all $k\in\naturalsPlus$ we have $b(k)\in\pow{\unions{\ran{X}}}$
            by \cref{sequence_in,funs_elim,function_apply_elim,ran_iff,function_to_ran,function_apply_elim,ran_intro,subseteq,unions_iff,pow_iff,subseteq_transitive,powerset_intro}.
        $b\in\funs{\naturalsPlus}{\pow{\unions{\ran{X}}}}$ by \cref{funs_intro}.
        $n\mathrel{R} b$ by \cref{pair_eq_iff}.
    \end{subproof}
    Take $b$ such that $b$ is a function on $\naturalsPlus$ and for all $n\in\naturalsPlus$ we have
        $n\mathrel{R}\apply{b}{n}$ by \cref{choice_function}.

    Let $F = \{ \funs{\initialSegment{k}}{\naturalsPlus} \mid k\in\naturalsPlus \}$.
    Let $S = \unions{F}$.

    Take $q$ such that
        $q$ is a function on $\naturalsPlus$ and
        for all $n\in\naturalsPlus$ we have $q(n)\in S$ and
        for all $u$ such that $u\in S$
            there exists $n\in\naturalsPlus$ such that $q(n) = u$
        by \cref{naturalsplus_finite_sequence_enumeration}.

    Let $R_d = \{(n,g) \mid n\in\naturalsPlus, g\in\pow{\naturalsPlus\times\pow{\productFamily{X}}} \mid \text{$g$ is a function on $\dom{\apply{q}{n}}$ and
        for all $i\in\dom{\apply{q}{n}}$ we have
            $(i, \preimg{\piIndex{X}{i}}{\apply{\apply{b}{i}}{\apply{\apply{q}{n}}{i}}})\in g$}\}$.
    $R_d$ is a relation by \cref{pair_eq_iff,relation}.
    Show that for all $n\in\naturalsPlus$ there exists $g$ such that $n\mathrel{R_d} g$.
    \begin{subproof}
        Fix $n\in\naturalsPlus$.
        Take $m\in\naturalsPlus$ such that $\apply{q}{n}\in\funs{\initialSegment{m}}{\naturalsPlus}$
            by \cref{naturalsplus_finite_sequence_enumeration,unions_iff}.
        $\apply{q}{n}$ is a function and $\dom{\apply{q}{n}} = \initialSegment{m}$ by \cref{funs_elim}.
        Let $g = \{(i, \preimg{\piIndex{X}{i}}{\apply{\apply{b}{i}}{\apply{\apply{q}{n}}{i}}}) \mid i\in\dom{\apply{q}{n}}\}$.
            For all $i,A_1,A_2$ such that $(i,A_1)\in g$ and $(i,A_2)\in g$ we have $A_1 = A_2$ by \cref{pair_eq_iff}.
            $g$ is a function by \cref{rightunique,relation,pair_eq_iff}.
            For all $i$ such that $i\in\dom{\apply{q}{n}}$ we have $i\in\dom{g}$ by \cref{pair_eq_iff,dom_intro}.
            $\dom{\apply{q}{n}}\subseteq\dom{g}$ by \cref{subseteq}.
            For all $i$ such that $i\in\dom{g}$ we have $i\in\dom{\apply{q}{n}}$ by \cref{dom_iff,pair_eq_iff}.
            $\dom{g}\subseteq\dom{\apply{q}{n}}$ by \cref{subseteq}.
            $\dom{\apply{q}{n}} = \dom{g}$ by \cref{subseteq_antisymmetric}.
            $g$ is a function on $\dom{\apply{q}{n}}$.
        Show that $g\in\pow{\naturalsPlus\times\pow{\productFamily{X}}}$.
        \begin{subproof}
            For all $w$ such that $w\in g$ we have $w\in\naturalsPlus\times\pow{\productFamily{X}}$
                by \cref{pair_eq_iff,subseteq,initial_segment_naturalsplus,projection_map_indexed_function,funs_elim,preimg_subseteq_dom,powerset_intro,times_tuple_intro}.
            $g\subseteq\naturalsPlus\times\pow{\productFamily{X}}$ by \cref{subseteq}.
            $g\in\pow{\naturalsPlus\times\pow{\productFamily{X}}}$ by \cref{powerset_intro}.
        \end{subproof}
        For all $i\in\dom{\apply{q}{n}}$ we have
            $\apply{g}{i} = \preimg{\piIndex{X}{i}}{\apply{\apply{b}{i}}{\apply{\apply{q}{n}}{i}}}$
            by \cref{pair_eq_iff,dom_intro,function_apply_iff}.
        $n\mathrel{R_d} g$ by \cref{pair_eq_iff}.
    \end{subproof}
    Take $h$ such that $h$ is a function on $\naturalsPlus$ and for all $n\in\naturalsPlus$ we have
        $n\mathrel{R_d}\apply{h}{n}$ by \cref{choice_function}.

    Let $d = \{(n, \inters{\img{\apply{h}{n}}{\dom{\apply{q}{n}}}}) \mid n\in\naturalsPlus\}$.

    $1\in\dom{X}$ by \cref{real_one_in_naturals,subseteq}.
    There exists $a$ such that $a\in\dom{X}$.
    For all $n\in\dom{X}$ we have $\apply{T}{n}$ is a topology on $\apply{X}{n}$.
    $\productSubbasis{T}{X}$ is a subbasis on $\productFamily{X}$ by \cref{product_subbasis_is_subbasis}.

    For all $n,U_1,U_2$ such that $(n,U_1)\in d$ and $(n,U_2)\in d$ we have $U_1 = U_2$ by \cref{pair_eq_iff}.
    $d$ is a function by \cref{rightunique,relation,pair_eq_iff}.
    Show that for all $n\in\naturalsPlus$ we have $d(n)\in\pow{\productFamily{X}}$.
    \begin{subproof}
        Fix $n\in\naturalsPlus$.
        Let $U = \inters{\img{\apply{h}{n}}{\dom{\apply{q}{n}}}}$.
        $d(n) = U$ by \cref{pair_eq_iff,dom_intro,function_apply_iff}.
        Let $F_0 = \img{\apply{h}{n}}{\dom{\apply{q}{n}}}$.
        Show that for all $A$ such that $A\in F_0$ we have $A\in\productSubbasis{T}{X}$.
        \begin{subproof}
            Fix $A$ such that $A\in F_0$.
            Take $i$ such that $i\in\dom{\apply{q}{n}}$ and $(i,A)\in\apply{h}{n}$ by \cref{img_iff}.
            $\apply{h}{n}$ is a function on $\dom{\apply{q}{n}}$ and
                for all $j\in\dom{\apply{q}{n}}$ we have
                    $(j, \preimg{\piIndex{X}{j}}{\apply{\apply{b}{j}}{\apply{\apply{q}{n}}{j}}})\in\apply{h}{n}$
                by \cref{choice_function,pair_eq_iff}.
            $\apply{h}{n}$ is a function and $\dom{\apply{h}{n}} = \dom{\apply{q}{n}}$.
            $i\in\dom{\apply{h}{n}}$.
            $\apply{\apply{h}{n}}{i} = A$ by \cref{function_apply_iff}.
            $(i, \preimg{\piIndex{X}{i}}{\apply{\apply{b}{i}}{\apply{\apply{q}{n}}{i}}})\in\apply{h}{n}$.
            $\apply{\apply{h}{n}}{i} = \preimg{\piIndex{X}{i}}{\apply{\apply{b}{i}}{\apply{\apply{q}{n}}{i}}}$
                by \cref{function_apply_iff}.
            $A = \preimg{\piIndex{X}{i}}{\apply{\apply{b}{i}}{\apply{\apply{q}{n}}{i}}}$.
            Take $m\in\naturalsPlus$ such that $\apply{q}{n}\in\funs{\initialSegment{m}}{\naturalsPlus}$
                by \cref{naturalsplus_finite_sequence_enumeration,unions_iff}.
            $\apply{q}{n}$ is a function and $\dom{\apply{q}{n}} = \initialSegment{m}$ by \cref{funs_elim}.
            $i\in\dom{X}$ by \cref{funs_elim,subseteq,initial_segment_naturalsplus}.
            $\apply{b}{i}$ is a sequence in $\pow{\apply{X}{i}}$ and
                for all $k\in\naturalsPlus$ we have $\apply{\apply{b}{i}}{k}$ is open in $\apply{X}{i}$ with respect to $\apply{T}{i}$
                by \cref{choice_function,pair_eq_iff}.
            $\apply{\apply{b}{i}}{\apply{\apply{q}{n}}{i}}\in\apply{T}{i}$
                by \cref{funs_elim,function_apply_elim,ran_iff,function_to_ran,subseteq,open_in}.
            $A\in\pow{\productFamily{X}}$ by \cref{projection_map_indexed_function,funs_elim,preimg_subseteq_dom,powerset_intro}.
            $A\in\productSubbasisAt{T}{X}{i}$ by \cref{product_subbasis_indexed}.
            $A\in\productSubbasis{T}{X}$ by \cref{product_subbasis_union_indexed}.
        \end{subproof}
        $F_0\subseteq\productSubbasis{T}{X}$ by \cref{subseteq}.
        Take $m\in\naturalsPlus$ such that $\apply{q}{n}\in\funs{\initialSegment{m}}{\naturalsPlus}$
            by \cref{naturalsplus_finite_sequence_enumeration,unions_iff}.
        $\apply{q}{n}$ is a function and $\dom{\apply{q}{n}} = \initialSegment{m}$ by \cref{funs_elim}.
        $\dom{\apply{q}{n}}$ is finite by \cref{initial_segment_finite}.
        $\apply{h}{n}$ is a function on $\dom{\apply{q}{n}}$ by \cref{choice_function,pair_eq_iff}.
        $\apply{h}{n}$ is a function and $\dom{\apply{h}{n}} = \dom{\apply{q}{n}}$.
        $F_0$ is finite by \cref{finite_image_function}.
        $\apply{\apply{h}{n}}{m}\in F_0$
            by \cref{funs_elim,subseteq,initial_segment_naturalsplus,choice_function,pair_eq_iff,function_apply_iff,img_iff}.
        For all $x$ such that $x\in U$ we have $x\in\unions{F_0}$ by \cref{inters}.
        $U\subseteq\unions{F_0}$ by \cref{subseteq}.
        For all $x$ such that $x\in\unions{F_0}$ we have $x\in\unions{\productSubbasis{T}{X}}$.
        $\unions{F_0}\subseteq\unions{\productSubbasis{T}{X}}$ by \cref{subseteq}.
        $U\subseteq\unions{\productSubbasis{T}{X}}$ by \cref{subseteq_transitive}.
        $U\in\pow{\unions{\productSubbasis{T}{X}}}$ by \cref{powerset_intro}.
        $\unions{\productSubbasis{T}{X}} = \productFamily{X}$ by \cref{subbasis_on}.
        $U\in\pow{\productFamily{X}}$ by \cref{subseteq}.
        $d(n)\in\pow{\productFamily{X}}$.
    \end{subproof}
    $d\in\funs{\naturalsPlus}{\pow{\productFamily{X}}}$
        by \cref{funs_intro,subseteq,pair_eq_iff,dom_intro,dom_iff,subseteq_antisymmetric}.
    Show that $d$ is a sequence in $\pow{\productFamily{X}}$.
    \begin{subproof}
        Follows by \cref{sequence_in}.
    \end{subproof}

    Show that for all $n\in\naturalsPlus$ we have $d(n)$ is open in $\productFamily{X}$ with respect to $T_0$.
    \begin{subproof}
        Fix $n\in\naturalsPlus$.
        Let $F_1 = \img{\apply{h}{n}}{\dom{\apply{q}{n}}}$.
        Show that for all $A$ such that $A\in F_1$ we have $A\in\productSubbasis{T}{X}$.
        \begin{subproof}
            Fix $A$ such that $A\in F_1$.
            Take $i$ such that $i\in\dom{\apply{q}{n}}$ and $(i,A)\in\apply{h}{n}$ by \cref{img_iff}.
            $\apply{h}{n}$ is a function on $\dom{\apply{q}{n}}$ and
                for all $j\in\dom{\apply{q}{n}}$ we have
                    $(j, \preimg{\piIndex{X}{j}}{\apply{\apply{b}{j}}{\apply{\apply{q}{n}}{j}}})\in\apply{h}{n}$
                by \cref{choice_function,pair_eq_iff}.
            $\apply{h}{n}$ is a function and $\dom{\apply{h}{n}} = \dom{\apply{q}{n}}$.
            $i\in\dom{\apply{h}{n}}$.
            $\apply{\apply{h}{n}}{i} = A$ by \cref{function_apply_iff}.
            $(i, \preimg{\piIndex{X}{i}}{\apply{\apply{b}{i}}{\apply{\apply{q}{n}}{i}}})\in\apply{h}{n}$.
            $\apply{\apply{h}{n}}{i} = \preimg{\piIndex{X}{i}}{\apply{\apply{b}{i}}{\apply{\apply{q}{n}}{i}}}$
                by \cref{function_apply_iff}.
            $A = \preimg{\piIndex{X}{i}}{\apply{\apply{b}{i}}{\apply{\apply{q}{n}}{i}}}$.
            Take $m\in\naturalsPlus$ such that $\apply{q}{n}\in\funs{\initialSegment{m}}{\naturalsPlus}$
                by \cref{naturalsplus_finite_sequence_enumeration,unions_iff}.
            $\apply{q}{n}$ is a function to $\naturalsPlus$ and $\dom{\apply{q}{n}} = \initialSegment{m}$ by \cref{funs_elim}.
            $i\in\dom{X}$ by \cref{funs_elim,subseteq,initial_segment_naturalsplus}.
            $\apply{b}{i}$ is a sequence in $\pow{\apply{X}{i}}$ and
                for all $k\in\naturalsPlus$ we have $\apply{\apply{b}{i}}{k}$ is open in $\apply{X}{i}$ with respect to $\apply{T}{i}$
                by \cref{choice_function,pair_eq_iff}.
            $\apply{\apply{b}{i}}{\apply{\apply{q}{n}}{i}}\in\apply{T}{i}$
                by \cref{function_to_ran,function_apply_elim,ran_iff,subseteq,open_in}.
            $A\in\pow{\productFamily{X}}$ by \cref{projection_map_indexed_function,funs_elim,preimg_subseteq_dom,powerset_intro}.
            $A\in\productSubbasisAt{T}{X}{i}$ by \cref{product_subbasis_indexed}.
            $A\in\productSubbasis{T}{X}$ by \cref{product_subbasis_union_indexed}.
        \end{subproof}
        $F_1\subseteq\productSubbasis{T}{X}$ by \cref{subseteq}.
        Take $m\in\naturalsPlus$ such that $\apply{q}{n}\in\funs{\initialSegment{m}}{\naturalsPlus}$
            by \cref{naturalsplus_finite_sequence_enumeration,unions_iff}.
        $\apply{q}{n}$ is a function and $\dom{\apply{q}{n}} = \initialSegment{m}$ by \cref{funs_elim}.
        $\dom{\apply{q}{n}}$ is finite by \cref{initial_segment_finite}.
        $\apply{h}{n}$ is a function on $\dom{\apply{q}{n}}$ by \cref{choice_function,pair_eq_iff}.
        $\apply{h}{n}$ is a function and $\dom{\apply{h}{n}} = \dom{\apply{q}{n}}$.
        $F_1$ is finite by \cref{finite_image_function}.
        $\apply{\apply{h}{n}}{m}\in F_1$
            by \cref{funs_elim,subseteq,initial_segment_naturalsplus,choice_function,pair_eq_iff,function_apply_iff,img_iff}.
        For all $x$ such that $x\in\inters{F_1}$ we have $x\in\unions{\productSubbasis{T}{X}}$.
        $\inters{F_1}\subseteq\unions{\productSubbasis{T}{X}}$ by \cref{subseteq}.
        $\inters{F_1}\in\finiteIntersections{\productSubbasis{T}{X}}$ by \cref{powerset_intro,finite_intersections}.
        $\finiteIntersections{\productSubbasis{T}{X}}$ is a basis for $T_0$ on $\productFamily{X}$
            by \cref{product_subbasis_finite_intersections_basis}.
        $T_0 = \basisTopology{\finiteIntersections{\productSubbasis{T}{X}}}{\productFamily{X}}$
            by \cref{product_topology_indexed,topology_generated_by_subbasis,basis_for_topology}.
        $\inters{F_1}\in\basisTopology{\finiteIntersections{\productSubbasis{T}{X}}}{\productFamily{X}}$
            by \cref{basis_elem_open_in_generated,basis_for_topology}.
        $\inters{F_1}\in T_0$ by \cref{basis_for_topology}.
        $T_0$ is a topology on $\productFamily{X}$ by \cref{basis_for_topology,basis_topology_is_topology}.
        $d(n)$ is open in $\productFamily{X}$ with respect to $T_0$
            by \cref{basis_for_topology,basis_topology_is_topology,sequence_in,funs_is_function,funs_elim,subseteq,pair_eq_iff,function_apply_iff,open_in}.
    \end{subproof}

    Show that for all $x\in\productFamily{X}$ we have
        for all $U$ such that $U$ is open in $\productFamily{X}$ with respect to $T_0$ and $x\in U$
            there exists $n\in\naturalsPlus$ such that $x\in d(n)$ and $d(n)\subseteq U$.
    \begin{subproof}
        Fix $x\in\productFamily{X}$.
        Fix $U$ such that $U$ is open in $\productFamily{X}$ with respect to $T_0$ and $x\in U$.
        $U\in T_0$ by \cref{open_in}.
        $1\in\dom{X}$ by \cref{real_one_in_naturals,subseteq}.
        There exists $a_0$ such that $a_0\in\dom{X}$.
        For all $n\in\dom{X}$ we have $\apply{T}{n}$ is a topology on $\apply{X}{n}$.
        $\finiteIntersections{\productSubbasis{T}{X}}$ is a basis for $T_0$ on $\productFamily{X}$
            by \cref{product_subbasis_finite_intersections_basis}.
        $T_0 = \basisTopology{\finiteIntersections{\productSubbasis{T}{X}}}{\productFamily{X}}$
            by \cref{product_topology_indexed,topology_generated_by_subbasis,basis_for_topology}.
        $U\in\basisTopology{\finiteIntersections{\productSubbasis{T}{X}}}{\productFamily{X}}$.
        Take $B\in\finiteIntersections{\productSubbasis{T}{X}}$ such that $x\in B$ and $B\subseteq U$
            by \cref{topology_generated_by_basis}.
        Take $F_2\subseteq\productSubbasis{T}{X}$ such that $F_2$ is finite and inhabited and $B = \inters{F_2}$ by \cref{finite_intersections}.

        Let $R_2 = \{(A,i) \mid A\in F_2, i\in\dom{X} \mid
            \exists V\in\apply{T}{i}. A = \preimg{\piIndex{X}{i}}{V}\}$.
        $R_2$ is a relation by \cref{pair_eq_iff,relation}.
        Show that for all $A\in F_2$ there exists $i$ such that $A\mathrel{R_2} i$.
        \begin{subproof}
            Fix $A\in F_2$.
            Take $i$ such that $i\in\dom{X}$ and $A\in\productSubbasisAt{T}{X}{i}$
                by \cref{subseteq,product_subbasis_union_indexed}.
            Take $V\in\apply{T}{i}$ such that $A = \preimg{\piIndex{X}{i}}{V}$ by \cref{product_subbasis_indexed}.
            $A\mathrel{R_2} i$ by \cref{pair_eq_iff}.
        \end{subproof}
        Take $p$ such that $p$ is a function on $F_2$ and for all $A\in F_2$ we have $A\mathrel{R_2}\apply{p}{A}$
            by \cref{choice_function}.
        Let $J = \img{p}{F_2}$.
        $J$ is finite and inhabited by \cref{finite_image_function,choice_function,subseteq,function_apply_iff,img_iff}.
        For all $j$ such that $j\in J$ we have $j\in\naturalsPlus$
            by \cref{img_iff,choice_function,subseteq,function_apply_iff,pair_eq_iff}.
        $J\subseteq\naturalsPlus$ by \cref{subseteq}.

        Let $R_3 = \{(A,V) \mid A\in F_2, V\in\pow{\unions{\ran{X}}} \mid
            V\in\apply{T}{\apply{p}{A}} \land A = \preimg{\piIndex{X}{\apply{p}{A}}}{V}\}$.
        $R_3$ is a relation by \cref{pair_eq_iff,relation}.
        Show that for all $A\in F_2$ there exists $V$ such that $A\mathrel{R_3} V$.
        \begin{subproof}
            Fix $A\in F_2$.
            Take $V\in\apply{T}{\apply{p}{A}}$ such that $A = \preimg{\piIndex{X}{\apply{p}{A}}}{V}$
                by \cref{choice_function,pair_eq_iff}.
            $\apply{p}{A}\in\dom{X}$ by \cref{choice_function,pair_eq_iff}.
            $V$ is open in $\apply{X}{\apply{p}{A}}$ with respect to $\apply{T}{\apply{p}{A}}$ by \cref{open_in}.
            $V\subseteq\apply{X}{\apply{p}{A}}$ by \cref{topology_on,open_in,powerset_elim,subseteq}.
            $\apply{X}{\apply{p}{A}}\in\ran{X}$ by \cref{function_apply_elim,ran_intro}.
            $V\subseteq\unions{\ran{X}}$ by \cref{subseteq_transitive,subseteq,unions_iff}.
            $V\in\pow{\unions{\ran{X}}}$ by \cref{powerset_intro}.
            $A\mathrel{R_3} V$ by \cref{pair_eq_iff}.
        \end{subproof}
        Take $r$ such that $r$ is a function on $F_2$ and for all $A\in F_2$ we have $A\mathrel{R_3}\apply{r}{A}$
            by \cref{choice_function}.

        Let $R_4 = \{(i,\inters{\img{r}{\preimg{p}{\{i\}}}}) \mid i\in J\}$.
        $R_4$ is a relation by \cref{pair_eq_iff,relation}.
        Show that for all $i\in J$ there exists $V$ such that $i\mathrel{R_4} V$.
        \begin{subproof}
            Fix $i\in J$.
            Let $H = \img{r}{\preimg{p}{\{i\}}}$.
            Let $V = \inters{H}$.
            $i\mathrel{R_4} V$ by \cref{pair_eq_iff}.
        \end{subproof}
        Take $v$ such that $v$ is a function on $J$ and for all $i\in J$ we have $i\mathrel{R_4}\apply{v}{i}$
            by \cref{choice_function}.

        Show that for all $i\in J$ we have $\apply{v}{i}\in\apply{T}{i}$.
        \begin{subproof}
            Fix $i\in J$.
            Let $H = \img{r}{\preimg{p}{\{i\}}}$.
            $\apply{v}{i} = \inters{H}$ by \cref{choice_function,pair_eq_iff}.
            For all $V$ such that $V\in H$ we have $V\in\apply{T}{i}$
                by \cref{img_iff,preimg_iff,choice_function,pair_eq_iff,dom_intro,subseteq,singleton_elim,function_apply_iff}.
            $H\subseteq\apply{T}{i}$ by \cref{subseteq}.
            $\img{r}{F_2}$ is finite by \cref{finite_image_function}.
            For all $V$ such that $V\in H$ we have $V\in\img{r}{F_2}$
                by \cref{img_iff,preimg_iff,choice_function,dom_intro,subseteq,singleton_elim,function_apply_iff}.
            $H$ is finite by \cref{subseteq,subseteq_finite_is_finite,finite_image_function}.
            $H$ is inhabited
                by \cref{img_iff,singleton_intro,preimg_iff,choice_function,subseteq,function_apply_iff}.
            $\apply{v}{i}\in\apply{T}{i}$ by \cref{subseteq,finite_intersections_open_in_topology}.
        \end{subproof}

        Show that for all $i\in J$ we have $\apply{x}{i}\in\apply{v}{i}$.
        \begin{subproof}
            Fix $i\in J$.
            Let $H = \img{r}{\preimg{p}{\{i\}}}$.
            $\apply{v}{i} = \inters{H}$ by \cref{choice_function,pair_eq_iff}.
            Show that for all $V\in H$ we have $\apply{x}{i}\in V$.
            \begin{subproof}
                Fix $V\in H$.
                Take $A$ such that $A\in\preimg{p}{\{i\}}$ and $(A,V)\in r$ by \cref{img_iff}.
                Take $y\in\{i\}$ such that $A\mathrel{p} y$ by \cref{preimg_iff}.
                $(A,y)\in p$.
                $A\in F_2$ and $\apply{p}{A} = i$ by \cref{choice_function,dom_intro,subseteq,singleton_elim,function_apply_iff}.
                $\apply{r}{A} = V$ by \cref{choice_function,subseteq,function_apply_iff}.
                $A = \preimg{\piIndex{X}{i}}{V}$ by \cref{choice_function,pair_eq_iff}.
                $x\in A$ by \cref{inters_subseteq_elem,subseteq}.
                Take $y_0\in V$ such that $(x,y_0)\in\piIndex{X}{i}$ by \cref{preimg_iff}.
                Take $z\in\productFamily{X}$ such that $(z,\apply{z}{i}) = (x,y_0)$ by \cref{projection_map_indexed}.
                $z = x$ and $\apply{z}{i} = y_0$ by \cref{pair_eq_iff}.
                $\apply{x}{i} = y_0$.
                $\apply{x}{i}\in V$.
            \end{subproof}
            $H$ is inhabited
                by \cref{img_iff,singleton_intro,preimg_iff,choice_function,subseteq,function_apply_iff}.
            $\apply{x}{i}\in\apply{v}{i}$ by \cref{choice_function,pair_eq_iff,inters_intro}.
        \end{subproof}

        Let $R_5 = \{(i,k) \mid i\in J, k\in\naturalsPlus \mid
            \apply{x}{i}\in\apply{\apply{b}{i}}{k} \land \apply{\apply{b}{i}}{k}\subseteq\apply{v}{i}\}$.
        $R_5$ is a relation by \cref{pair_eq_iff,relation}.
        Show that for all $i\in J$ there exists $k$ such that $i\mathrel{R_5} k$.
        \begin{subproof}
            Fix $i\in J$.
            $\apply{v}{i}$ is open in $\apply{X}{i}$ with respect to $\apply{T}{i}$ by \cref{open_in}.
            $\apply{x}{i}\in\apply{v}{i}$.
            $\apply{v}{i}\subseteq\apply{X}{i}$ by \cref{topology_on,open_in,powerset_elim,subseteq}.
            $\apply{x}{i}\in\apply{X}{i}$ by \cref{subseteq}.
            $i\in\naturalsPlus$ by \cref{subseteq}.
            $i\mathrel{R}\apply{b}{i}$ by \cref{choice_function,pair_eq_iff}.
            $\apply{b}{i}$ is a sequence in $\pow{\apply{X}{i}}$ by \cref{pair_eq_iff}.
            For all $k\in\naturalsPlus$ we have $\apply{\apply{b}{i}}{k}$ is open in $\apply{X}{i}$ with respect to $\apply{T}{i}$
                by \cref{pair_eq_iff}.
            For all $y\in\apply{X}{i}$ we have
                for all $U$ such that $U$ is open in $\apply{X}{i}$ with respect to $\apply{T}{i}$ and $y\in U$
                    there exists $k\in\naturalsPlus$ such that $y\in\apply{\apply{b}{i}}{k}$ and $\apply{\apply{b}{i}}{k}\subseteq U$
                by \cref{pair_eq_iff}.
            Take $k\in\naturalsPlus$ such that
                $\apply{x}{i}\in\apply{\apply{b}{i}}{k}$ and $\apply{\apply{b}{i}}{k}\subseteq\apply{v}{i}$.
            $i\mathrel{R_5} k$ by \cref{pair_eq_iff}.
        \end{subproof}
        Take $s$ such that $s$ is a function on $J$ and for all $i\in J$ we have $i\mathrel{R_5}\apply{s}{i}$
            by \cref{choice_function}.

        $J\subseteq\reals$ by \cref{real_naturals_subset_reals,subseteq_transitive}.
        Take $m$ such that $m$ is a largest element of $J$ with respect to $\realOrderRel$
            by \cref{finite_inhabited_largest_element,real_order_relation_simple}.
        $m\in\naturalsPlus$ by \cref{largest_element,subseteq}.

        Let $R_6 = \{(i,k) \mid i\in\initialSegment{m}, k\in\naturalsPlus \mid
            (i\in J \land \apply{x}{i}\in\apply{\apply{b}{i}}{k} \land \apply{\apply{b}{i}}{k}\subseteq\apply{v}{i})
            \lor (i\notin J \land \apply{x}{i}\in\apply{\apply{b}{i}}{k})\}$.
        $R_6$ is a relation by \cref{pair_eq_iff,relation}.
        Show that for all $i\in\initialSegment{m}$ there exists $k$ such that $i\mathrel{R_6} k$.
        \begin{subproof}
            Fix $i\in\initialSegment{m}$.
            \begin{byCase}
            \caseOf{$i\in J$.}
                $i\mathrel{R_6}\apply{s}{i}$ by \cref{choice_function,pair_eq_iff}.
            \caseOf{$i\notin J$.}
                $i\in\dom{X}$ by \cref{initial_segment_naturalsplus,subseteq}.
                $i\in\naturalsPlus$ by \cref{initial_segment_naturalsplus,subseteq}.
                $\apply{X}{i}$ is open in $\apply{X}{i}$ with respect to $\apply{T}{i}$ by \cref{topology_on,open_in}.
                $\apply{x}{i}\in\apply{X}{i}$ by \cref{indexed_product}.
                $i\mathrel{R}\apply{b}{i}$ by \cref{choice_function,pair_eq_iff}.
                $\apply{b}{i}$ is a sequence in $\pow{\apply{X}{i}}$ by \cref{pair_eq_iff}.
                For all $k\in\naturalsPlus$ we have $\apply{\apply{b}{i}}{k}$ is open in $\apply{X}{i}$ with respect to $\apply{T}{i}$
                    by \cref{pair_eq_iff}.
                For all $y\in\apply{X}{i}$ we have
                    for all $U$ such that $U$ is open in $\apply{X}{i}$ with respect to $\apply{T}{i}$ and $y\in U$
                        there exists $k\in\naturalsPlus$ such that $y\in\apply{\apply{b}{i}}{k}$ and $\apply{\apply{b}{i}}{k}\subseteq U$
                    by \cref{pair_eq_iff}.
                Take $k\in\naturalsPlus$ such that $\apply{x}{i}\in\apply{\apply{b}{i}}{k}$ and $\apply{\apply{b}{i}}{k}\subseteq\apply{X}{i}$.
                $i\mathrel{R_6} k$ by \cref{pair_eq_iff}.
            \end{byCase}
        \end{subproof}
        Take $u$ such that $u$ is a function on $\initialSegment{m}$ and for all $i\in\initialSegment{m}$ we have
            $i\mathrel{R_6}\apply{u}{i}$ by \cref{choice_function}.
        $u\in\funs{\initialSegment{m}}{\naturalsPlus}$ by \cref{funs_intro,subseteq,choice_function,pair_eq_iff}.
        $u\in S$ by \cref{unions_intro}.
        Take $n\in\naturalsPlus$ such that $\apply{q}{n} = u$ by \cref{naturalsplus_finite_sequence_enumeration}.

        Let $F_n = \img{\apply{h}{n}}{\dom{\apply{q}{n}}}$.
        $\apply{q}{n} = u$.
        $\dom{\apply{q}{n}} = \initialSegment{m}$ by \cref{funs_elim}.
        For all $i\in\initialSegment{m}$ we have
            $x\in\preimg{\piIndex{X}{i}}{\apply{\apply{b}{i}}{\apply{u}{i}}}$
            by \cref{choice_function,pair_eq_iff,preimg_iff,projection_map_indexed}.
        Show that for all $j\in\dom{\apply{q}{n}}$ we have
            $(j,\preimg{\piIndex{X}{j}}{\apply{\apply{b}{j}}{\apply{\apply{q}{n}}{j}}})\in\apply{h}{n}$.
        \begin{subproof}
            Fix $j\in\dom{\apply{q}{n}}$.
            $n\mathrel{R_d}\apply{h}{n}$ by \cref{choice_function}.
            Take $n_0,g$ such that $n_0\in\naturalsPlus$ and
                $g\in\pow{\naturalsPlus\times\pow{\productFamily{X}}}$ and
                $(n,\apply{h}{n}) = (n_0,g)$ and
                $g$ is a function on $\dom{\apply{q}{n_0}}$ and
                for all $i\in\dom{\apply{q}{n_0}}$ we have
                    $(i,\preimg{\piIndex{X}{i}}{\apply{\apply{b}{i}}{\apply{\apply{q}{n_0}}{i}}})\in g$
                by \cref{pair_eq_iff}.
            $n_0 = n$ and $g = \apply{h}{n}$ by \cref{pair_eq_iff}.
            $(j,\preimg{\piIndex{X}{j}}{\apply{\apply{b}{j}}{\apply{\apply{q}{n}}{j}}})\in\apply{h}{n}$.
        \end{subproof}
        Show that for all $A\in F_n$ we have $x\in A$.
        \begin{subproof}
            Fix $A\in F_n$.
            Take $i$ such that $i\in\dom{\apply{q}{n}}$ and $(i,A)\in\apply{h}{n}$ by \cref{img_iff}.
            $i\in\initialSegment{m}$ by \cref{subseteq}.
            $n\mathrel{R_d}\apply{h}{n}$ by \cref{choice_function}.
            $\apply{h}{n}$ is a function on $\dom{\apply{q}{n}}$ by \cref{choice_function,pair_eq_iff}.
            $(i,\preimg{\piIndex{X}{i}}{\apply{\apply{b}{i}}{\apply{\apply{q}{n}}{i}}})\in\apply{h}{n}$.
            $i\in\dom{\apply{h}{n}}$ and
                $\apply{\apply{h}{n}}{i} = \preimg{\piIndex{X}{i}}{\apply{\apply{b}{i}}{\apply{\apply{q}{n}}{i}}}$
                by \cref{dom_intro,function_apply_iff}.
            $A = \apply{\apply{h}{n}}{i}$ by \cref{pair_eq_iff,function_apply_iff}.
            $A = \preimg{\piIndex{X}{i}}{\apply{\apply{b}{i}}{\apply{\apply{q}{n}}{i}}}$.
            $\apply{\apply{q}{n}}{i} = \apply{u}{i}$.
            $x\in A$ by \cref{subseteq}.
        \end{subproof}
        $m\in\dom{\apply{q}{n}}$ by \cref{subseteq,initial_segment_naturalsplus}.
        $n\mathrel{R_d}\apply{h}{n}$ by \cref{choice_function}.
        $\apply{h}{n}$ is a function on $\dom{\apply{q}{n}}$ and
            for all $j\in\dom{\apply{q}{n}}$ we have
                $(j,\preimg{\piIndex{X}{j}}{\apply{\apply{b}{j}}{\apply{\apply{q}{n}}{j}}})\in\apply{h}{n}$
            by \cref{pair_eq_iff}.
        $\apply{h}{n}$ is a function and $\dom{\apply{h}{n}} = \dom{\apply{q}{n}}$.
        $(m,\apply{\apply{h}{n}}{m})\in\apply{h}{n}$ by \cref{function_apply_iff}.
        $\apply{\apply{h}{n}}{m}\in F_n$ by \cref{img_iff}.
        $F_n$ is inhabited.
        $d(n) = \inters{F_n}$ by \cref{pair_eq_iff,dom_intro,function_apply_iff}.
        $x\in d(n)$ by \cref{inters_intro}.

        Show that for all $A\in F_2$ we have $d(n)\subseteq A$.
        \begin{subproof}
            Fix $A\in F_2$.
            $(A,\apply{p}{A})\in p$ and $\apply{p}{A}\in J$
                by \cref{choice_function,subseteq,function_apply_elim,img_iff}.
            Let $i = \apply{p}{A}$.
            $i\in\naturalsPlus$ by \cref{subseteq}.
            $i\leq m$ by \cref{largest_element,real_order_relation_elim}.
            $i\in\initialSegment{m}$ by \cref{initial_segment_naturalsplus}.
            $\apply{u}{i}\in\naturalsPlus$ and
                $\apply{x}{i}\in\apply{\apply{b}{i}}{\apply{u}{i}}$ and
                $\apply{\apply{b}{i}}{\apply{u}{i}}\subseteq\apply{v}{i}$ by \cref{choice_function,pair_eq_iff}.
            $A = \preimg{\piIndex{X}{i}}{\apply{r}{A}}$ by \cref{choice_function,pair_eq_iff}.
            Let $H = \img{r}{\preimg{p}{\{i\}}}$.
            $\apply{v}{i} = \inters{H}$ by \cref{choice_function,pair_eq_iff}.
            $A\in\preimg{p}{\{i\}}$ by \cref{singleton_intro,preimg_iff}.
            $\apply{r}{A}\in H$ by \cref{choice_function,subseteq,function_apply_iff,img_iff}.
            $\apply{v}{i}\subseteq\apply{r}{A}$ by \cref{inters_subseteq_elem}.
            $\apply{\apply{b}{i}}{\apply{u}{i}}\subseteq\apply{r}{A}$ by \cref{subseteq_transitive}.
            $\preimg{\piIndex{X}{i}}{\apply{\apply{b}{i}}{\apply{u}{i}}}\subseteq A$ by \cref{preimg_subseteq,subseteq}.
            $(i,\preimg{\piIndex{X}{i}}{\apply{\apply{b}{i}}{\apply{\apply{q}{n}}{i}}})\in\apply{h}{n}$ by \cref{subseteq}.
            $\apply{\apply{q}{n}}{i} = \apply{u}{i}$.
            $\preimg{\piIndex{X}{i}}{\apply{\apply{b}{i}}{\apply{u}{i}}}\in F_n$ by \cref{img_iff}.
            $d(n) = \inters{F_n}$ by \cref{sequence_in,funs_is_function,funs_elim,subseteq,pair_eq_iff,function_apply_iff}.
            $d(n)\subseteq \preimg{\piIndex{X}{i}}{\apply{\apply{b}{i}}{\apply{u}{i}}}$ by \cref{inters_subseteq_elem}.
            $d(n)\subseteq A$ by \cref{subseteq_transitive}.
        \end{subproof}
        $d(n)\subseteq U$ by \cref{subseteq_inters_iff,subseteq,subseteq_transitive}.
    \end{subproof}
    $\productFamily{X}$ satisfies the second countability axiom with respect to $T_0$ by \cref{second_countability_axiom}.
\end{proof}

% from §30 Item 11: dense subset
\begin{definition}\label{dense_in}
    $A$ is dense in $X$ with respect to $T$ iff $\closure{A}{X}{T} = X$.
\end{definition}

% from §30 Item 12: countable subcover in a second-countable space
\begin{lemma}\label{second_countable_open_covering_countable_subcover}
    Assume $T$ is a topology on $X$.
    Suppose $X$ satisfies the second countability axiom with respect to $T$.
    Then for all $C$ such that $C$ is an open covering of $X$ with respect to $T$
        there exists $J$ such that
            $J\subseteq\naturalsPlus$ and
            there exists $f$ such that
                $f$ is a function on $J$ and
                for all $j\in J$ we have $\apply{f}{j}\in C$ and
                $X\subseteq\unions{\img{f}{J}}$.
\end{lemma}
\begin{proof}
    Take $b$ such that
        $b$ is a sequence in $\pow{X}$ and
        for all $n\in\naturalsPlus$ we have $b(n)$ is open in $X$ with respect to $T$ and
        for all $x\in X$ we have
            for all $U$ such that $U$ is open in $X$ with respect to $T$ and $x\in U$
                there exists $n\in\naturalsPlus$ such that $x\in b(n)$ and $b(n)\subseteq U$
        by \cref{second_countability_axiom}.
    Fix $C$ such that $C$ is an open covering of $X$ with respect to $T$.
    Let $J_0 = \{n\in\naturalsPlus \mid \exists A\in C. b(n)\subseteq A\}$.
    Let $R = \{(n,A) \mid n\in J_0, A\in C \mid b(n)\subseteq A\}$.
    $R$ is a relation by \cref{pair_eq_iff,relation}.
    For all $n\in J_0$ there exists $A$ such that $n\mathrel{R} A$ by \cref{pair_eq_iff}.
    Take $f_0$ such that $f_0$ is a function on $J_0$ and for all $n\in J_0$ we have $n\mathrel{R}\apply{f_0}{n}$
        by \cref{choice_function}.
    Show that for all $j\in J_0$ we have $\apply{f_0}{j}\in C$.
    \begin{subproof}
        Follows by \cref{choice_function,pair_eq_iff}.
    \end{subproof}
    Show that for all $x$ such that $x\in X$ we have $x\in\unions{\img{f_0}{J_0}}$.
    \begin{subproof}
        Fix $x$ such that $x\in X$.
        Take $A_0$ such that $A_0\in C$ and $x\in A_0$ by \cref{open_covering,covering,subseteq,unions_iff}.
        $A_0$ is open in $X$ with respect to $T$ by \cref{open_covering}.
        Take $n\in\naturalsPlus$ such that $x\in b(n)$ and $b(n)\subseteq A_0$.
        $\apply{f_0}{n}\in C$ and $b(n)\subseteq\apply{f_0}{n}$ by \cref{choice_function,pair_eq_iff}.
        $x\in\unions{\img{f_0}{J_0}}$ by \cref{choice_function,pair_eq_iff,subseteq,function_apply_iff,img_iff,unions_intro}.
    \end{subproof}
    $X\subseteq\unions{\img{f_0}{J_0}}$ by \cref{subseteq}.
    Show that there exists $J$ such that
        $J\subseteq\naturalsPlus$ and
        there exists $f$ such that
            $f$ is a function on $J$ and
            for all $j\in J$ we have $\apply{f}{j}\in C$ and
            $X\subseteq\unions{\img{f}{J}}$.
    \begin{subproof}
        Trivial.
    \end{subproof}
\end{proof}

% from §30 Item 13: countable dense subset in a second-countable space
\begin{lemma}\label{second_countable_countable_dense_subset}
    Assume $T$ is a topology on $X$.
    Suppose $X$ satisfies the second countability axiom with respect to $T$.
    Then there exists $J$ such that
        $J\subseteq\naturalsPlus$ and
        there exists $f$ such that
            $f$ is a function on $J$ and
            for all $j\in J$ we have $\apply{f}{j}\in X$ and
            $\img{f}{J}$ is dense in $X$ with respect to $T$.
\end{lemma}
\begin{proof}
    Take $b$ such that
        $b$ is a sequence in $\pow{X}$ and
        for all $n\in\naturalsPlus$ we have $b(n)$ is open in $X$ with respect to $T$ and
        for all $x\in X$ we have
            for all $U$ such that $U$ is open in $X$ with respect to $T$ and $x\in U$
                there exists $n\in\naturalsPlus$ such that $x\in b(n)$ and $b(n)\subseteq U$
        by \cref{second_countability_axiom}.
    Let $J_1 = \{n\in\naturalsPlus \mid b(n)\neq\emptyset\}$.
    Let $R = \{(n,x) \mid n\in J_1, x\in X \mid x\in b(n)\}$.
    $R$ is a relation by \cref{pair_eq_iff,relation}.
    For all $n\in J_1$ there exists $x$ such that $x\in b(n)$ and $x\in X$
        by \cref{pair_eq_iff,nonempty_inhabited,sequence_in,funs_elim,subseteq,function_to_ran,ran_iff,powerset_elim}.
    For all $n\in J_1$ there exists $x$ such that $n\mathrel{R} x$ by \cref{pair_eq_iff}.
    Take $f_1$ such that $f_1$ is a function on $J_1$ and for all $n\in J_1$ we have $n\mathrel{R}\apply{f_1}{n}$
        by \cref{choice_function}.
    $\img{f_1}{J_1}\subseteq X$
        by \cref{img_iff,choice_function,pair_eq_iff,sequence_in,funs_elim,subseteq,function_apply_iff,function_to_ran,ran_iff,powerset_elim}.
        Show that for all $x\in X$ we have $x\in\closure{\img{f_1}{J_1}}{X}{T}$.
        \begin{subproof}
            Fix $x\in X$.
            For all $U$ such that $U$ is open in $X$ with respect to $T$ and $x\in U$
                there exists $n\in\naturalsPlus$ such that $x\in b(n)$ and $b(n)\subseteq U$.
            Thus for all $U$ such that $U$ is open in $X$ with respect to $T$ and $x\in U$
                we have $U$ intersects $\img{f_1}{J_1}$
                by \cref{choice_function,pair_eq_iff,subseteq,function_apply_iff,img_iff,inter_intro,emptyset,intersects}.
            Therefore $x\in\closure{\img{f_1}{J_1}}{X}{T}$ by \cref{closure_iff_intersects_open}.
        \end{subproof}
        Hence $\img{f_1}{J_1}$ is dense in $X$ with respect to $T$
            by \cref{subseteq,closed_whole_space,closure_subseteq_closed,subseteq_antisymmetric,dense_in}.
    Show that there exists $J$ such that
        $J\subseteq\naturalsPlus$ and
        there exists $f$ such that
            $f$ is a function on $J$ and
            for all $j\in J$ we have $\apply{f}{j}\in X$ and
            $\img{f}{J}$ is dense in $X$ with respect to $T$.
    \begin{subproof}
        Trivial.
    \end{subproof}
\end{proof}

% from §30 Item 14: Lindelöf space
\begin{definition}\label{lindelof_space}
    $X$ is Lindelöf with respect to $T$ iff
        for all $C$ such that $C$ is an open covering of $X$ with respect to $T$
            there exists $J$ such that
                $J\subseteq\naturalsPlus$ and
                there exists $f$ such that
                    $f$ is a function on $J$ and
                    for all $j\in J$ we have $\apply{f}{j}\in C$ and
                    $X\subseteq\unions{\img{f}{J}}$.
\end{definition}

% from §30 Item 15: separable space
\begin{definition}\label{separable_space}
    $X$ is separable with respect to $T$ iff
        there exists $J$ such that
            $J\subseteq\naturalsPlus$ and
            there exists $f$ such that
                $f$ is a function on $J$ and
                for all $j\in J$ we have $\apply{f}{j}\in X$ and
                $\img{f}{J}$ is dense in $X$ with respect to $T$.
\end{definition}


\section{§31 The Separation Axioms}

% from §31 helper definition 1: closed singletons
\begin{definition}\label{one_point_closed}
    $X$ has closed singletons with respect to $T$ iff
        for all $x$ such that $x\in X$ we have $\{x\}$ is closed in $X$ with respect to $T$.
\end{definition}

% from §31 Item 1: regular space
\begin{definition}\label{regular_space}
    $X$ is regular with respect to $T$ iff
        $T$ is a topology on $X$ and
        $X$ has closed singletons with respect to $T$ and
        for all $x,B$ such that $x\in X$ and $B$ is closed in $X$ with respect to $T$ and $x\notin B$
            there exist $U,V$ such that
                $U$ is open in $X$ with respect to $T$ and
                $V$ is open in $X$ with respect to $T$ and
                $x\in U$ and
                $B\subseteq V$ and
                $U$ is disjoint from $V$.
\end{definition}

% from §31 Item 2: normal space
\begin{definition}\label{normal_space}
    $X$ is normal with respect to $T$ iff
        $T$ is a topology on $X$ and
        $X$ has closed singletons with respect to $T$ and
        for all $A,B$ such that $A$ is closed in $X$ with respect to $T$ and $B$ is closed in $X$ with respect to $T$
            and $A$ is disjoint from $B$
            there exist $U,V$ such that
                $U$ is open in $X$ with respect to $T$ and
                $V$ is open in $X$ with respect to $T$ and
                $A\subseteq U$ and
                $B\subseteq V$ and
                $U$ is disjoint from $V$.
\end{definition}

% from §31 helper lemma 1: regular implies Hausdorff
\begin{lemma}\label{regular_implies_hausdorff}
    Assume $X$ is regular with respect to $T$.
    Then $X$ is Hausdorff with respect to $T$.
\end{lemma}
\begin{proof}
    Show that for all $x_1,x_2$ such that $x_1\in X$ and $x_2\in X$ and $x_1\neq x_2$
        there exist $U_1,U_2$ such that
            $U_1$ is a neighborhood of $x_1$ in $X$ with respect to $T$ and
            $U_2$ is a neighborhood of $x_2$ in $X$ with respect to $T$ and
            $U_1$ is disjoint from $U_2$.
    \begin{subproof}
        Follows by \cref{regular_space,one_point_closed,singleton_elim,singleton_intro,subseteq,neighborhood}.
    \end{subproof}
    Hence $X$ is Hausdorff with respect to $T$ by \cref{regular_space,hausdorff}.
\end{proof}

% from §31 helper lemma 2: normal implies regular
\begin{lemma}\label{normal_implies_regular}
    Assume $X$ is normal with respect to $T$.
    Then $X$ is regular with respect to $T$.
\end{lemma}
\begin{proof}
    Show that for all $x,B$ such that $x\in X$ and $B$ is closed in $X$ with respect to $T$ and $x\notin B$
        there exist $U,V$ such that
            $U$ is open in $X$ with respect to $T$ and
            $V$ is open in $X$ with respect to $T$ and
            $x\in U$ and
            $B\subseteq V$ and
            $U$ is disjoint from $V$.
    \begin{subproof}
        Follows by \cref{normal_space,one_point_closed,singleton_elim,disjoint,singleton_intro,subseteq}.
    \end{subproof}
    Hence $X$ is regular with respect to $T$ by \cref{normal_space,regular_space}.
\end{proof}

% from §31 Item 3: regular neighborhood closure property
\begin{lemma}\label{regular_closure_neighborhood}
    Assume $X$ is regular with respect to $T$.
    Then for all $x,U$ such that $x\in X$ and $U$ is a neighborhood of $x$ in $X$ with respect to $T$
        there exists $V$ such that
            $V$ is a neighborhood of $x$ in $X$ with respect to $T$ and
            $\closure{V}{X}{T}\subseteq U$.
\end{lemma}
\begin{proof}
    $T$ is a topology on $X$ by \cref{regular_space}.
    Fix $x$.
    Fix $U$ such that $x\in X$ and $U$ is a neighborhood of $x$ in $X$ with respect to $T$.
    Let $B = X\setminus U$.
    $U$ is open in $X$ with respect to $T$ and $U\subseteq X$
        by \cref{neighborhood,open_in,topology_on,subseteq,pow_iff}.
    $B$ is closed in $X$ with respect to $T$ by \cref{setminus_subseteq,double_relative_complement,subseteq,open_in,closed_in}.
    Take $U_1,V_1$ such that
        $U_1$ is open in $X$ with respect to $T$ and
        $V_1$ is open in $X$ with respect to $T$ and
        $x\in U_1$ and
        $B\subseteq V_1$ and
        $U_1$ is disjoint from $V_1$
        by \cref{setminus_elim_right,neighborhood,regular_space}.
    Let $V = U_1$.
    $V$ is a neighborhood of $x$ in $X$ with respect to $T$ by \cref{open_in,neighborhood}.
    $V$ is open in $X$ with respect to $T$ and $V\subseteq X$
        by \cref{open_in,topology_on,subseteq,pow_iff}.
    We have $V\subseteq X\setminus V_1$ by \cref{disjoint,subseteq,setminus_intro}.
    $X\setminus V_1$ is closed in $X$ with respect to $T$
        by \cref{open_in,topology_on,subseteq,pow_iff,setminus_subseteq,double_relative_complement,closed_in}.
    Hence $\closure{V}{X}{T}\subseteq X\setminus V_1$ by \cref{closure_subseteq_closed}.
    We have $X\setminus V_1\subseteq X\setminus B$ by \cref{subseteq_implies_setminus_supseteq,subseteq}.
    Finally $X\setminus V_1\subseteq U$ and $\closure{V}{X}{T}\subseteq U$
        by \cref{subseteq_transitive,double_relative_complement,subseteq}.
\end{proof}

% from §31 Item 4: regularity from neighborhood closure property
\begin{lemma}\label{regular_from_closure_neighborhood}
    Suppose $T$ is a topology on $X$.
    Suppose for all $x$ such that $x\in X$ we have $\{x\}$ is closed in $X$ with respect to $T$.
    Suppose for all $x,U$ such that $x\in X$ and $U$ is a neighborhood of $x$ in $X$ with respect to $T$
        there exists $V$ such that
            $V$ is a neighborhood of $x$ in $X$ with respect to $T$ and
            $\closure{V}{X}{T}\subseteq U$.
    Then $X$ is regular with respect to $T$.
\end{lemma}
\begin{proof}
    Show that $T$ is a topology on $X$.
    \begin{subproof}
        Trivial.
    \end{subproof}
    Show that $X$ has closed singletons with respect to $T$.
    \begin{subproof}
        Trivial.
    \end{subproof}
    Show that for all $x,B$ such that $x\in X$ and $B$ is closed in $X$ with respect to $T$ and $x\notin B$
        there exist $U,V$ such that
            $U$ is open in $X$ with respect to $T$ and
            $V$ is open in $X$ with respect to $T$ and
            $x\in U$ and
            $B\subseteq V$ and
            $U$ is disjoint from $V$.
    \begin{subproof}
        Fix $x$.
        Fix $B$ such that $x\in X$ and $B$ is closed in $X$ with respect to $T$ and $x\notin B$.
        Let $U = X\setminus B$.
        $U$ is a neighborhood of $x$ in $X$ with respect to $T$ by \cref{closed_in,setminus_intro,neighborhood}.
        Take $V$ such that
            $V$ is a neighborhood of $x$ in $X$ with respect to $T$ and
            $\closure{V}{X}{T}\subseteq U$.
        Let $W = X\setminus \closure{V}{X}{T}$.
        $V$ is open in $X$ with respect to $T$ and $V\subseteq X$
            by \cref{neighborhood,open_in,topology_on,subseteq,pow_iff}.
        $W$ is open in $X$ with respect to $T$ and $B\subseteq X$ by \cref{closure_closed_in,closed_in}.
        Thus $B\subseteq W$ by \cref{subseteq_implies_setminus_supseteq,double_relative_complement,subseteq}.
        We have $V\subseteq \closure{V}{X}{T}$ by \cref{subseteq_closure}.
        Hence $V$ is disjoint from $W$ by \cref{subseteq,setminus_elim_right,disjoint}.
        Show that there exist $U_0,V_0$ such that
            $U_0$ is open in $X$ with respect to $T$ and
            $V_0$ is open in $X$ with respect to $T$ and
            $x\in U_0$ and
            $B\subseteq V_0$ and
            $U_0$ is disjoint from $V_0$.
        \begin{subproof}
            Trivial.
        \end{subproof}
    \end{subproof}
    Therefore $X$ is regular with respect to $T$ by \cref{one_point_closed,regular_space}.
\end{proof}

% from §31 Item 5: normal neighborhood closure property
\begin{lemma}\label{normal_closure_neighborhood}
    Assume $X$ is normal with respect to $T$.
    Then for all $A,U$ such that $A$ is closed in $X$ with respect to $T$ and $U$ is open in $X$ with respect to $T$
        and $A\subseteq U$
        there exists $V$ such that
            $V$ is open in $X$ with respect to $T$ and
            $A\subseteq V$ and
            $\closure{V}{X}{T}\subseteq U$.
\end{lemma}
\begin{proof}
    Fix $A$.
    Fix $U$ such that $A$ is closed in $X$ with respect to $T$ and $U$ is open in $X$ with respect to $T$ and $A\subseteq U$.
    Let $B = X\setminus U$.
    $U\in T$ and $U\subseteq X$ by \cref{open_in,normal_space,topology_on,subseteq,pow_iff}.
    $B$ is closed in $X$ with respect to $T$ by \cref{setminus_subseteq,double_relative_complement,subseteq,open_in,closed_in}.
    Thus $A$ is disjoint from $B$ by \cref{disjoint,subseteq,setminus_elim_right}.
    Take $V,W$ such that
        $V$ is open in $X$ with respect to $T$ and
        $W$ is open in $X$ with respect to $T$ and
        $A\subseteq V$ and
        $B\subseteq W$ and
        $V$ is disjoint from $W$
        by \cref{normal_space}.

    We have $V\subseteq X$ by \cref{open_in,normal_space,topology_on,subseteq,pow_iff}.
    We have $V\subseteq X\setminus W$ by \cref{disjoint,subseteq,setminus_intro}.
    $X\setminus W$ is closed in $X$ with respect to $T$
        by \cref{open_in,normal_space,topology_on,subseteq,pow_iff,setminus_subseteq,double_relative_complement,closed_in}.
    Hence $\closure{V}{X}{T}\subseteq X\setminus W$ by \cref{normal_space,closure_subseteq_closed}.
    We have $X\setminus W\subseteq X\setminus B$ by \cref{subseteq,subseteq_implies_setminus_supseteq}.
    Thus $\closure{V}{X}{T}\subseteq U$ by \cref{subseteq_transitive,double_relative_complement,subseteq}.
    Show that there exists $V_0$ such that
        $V_0$ is open in $X$ with respect to $T$ and
        $A\subseteq V_0$ and
        $\closure{V_0}{X}{T}\subseteq U$.
    \begin{subproof}
        Trivial.
    \end{subproof}
\end{proof}
% from §31 Item 6: normality from neighborhood closure property
\begin{lemma}\label{normal_from_closure_neighborhood}
    Suppose $T$ is a topology on $X$.
    Suppose for all $x$ such that $x\in X$ we have $\{x\}$ is closed in $X$ with respect to $T$.
    Suppose for all $A,U$ such that $A$ is closed in $X$ with respect to $T$ and $U$ is open in $X$ with respect to $T$
        and $A\subseteq U$
        there exists $V$ such that
            $V$ is open in $X$ with respect to $T$ and
            $A\subseteq V$ and
            $\closure{V}{X}{T}\subseteq U$.
    Then $X$ is normal with respect to $T$.
\end{lemma}
\begin{proof}
    Show that $T$ is a topology on $X$.
    \begin{subproof}
        Trivial.
    \end{subproof}
    Show that $X$ has closed singletons with respect to $T$.
    \begin{subproof}
        Trivial.
    \end{subproof}
    Show that for all $A,B$ such that $A$ is closed in $X$ with respect to $T$ and $B$ is closed in $X$ with respect to $T$
        and $A$ is disjoint from $B$
        there exist $U,V$ such that
            $U$ is open in $X$ with respect to $T$ and
            $V$ is open in $X$ with respect to $T$ and
            $A\subseteq U$ and
            $B\subseteq V$ and
            $U$ is disjoint from $V$.
    \begin{subproof}
        Fix $A$.
        Fix $B$ such that $A$ is closed in $X$ with respect to $T$ and $B$ is closed in $X$ with respect to $T$
            and $A$ is disjoint from $B$.
        Let $U = X\setminus B$.
        $A\subseteq U$ by \cref{disjoint,closed_in,subseteq,setminus_intro}.
        $U$ is open in $X$ with respect to $T$ by \cref{closed_in}.
        Take $V$ such that
            $V$ is open in $X$ with respect to $T$ and
            $A\subseteq V$ and
            $\closure{V}{X}{T}\subseteq U$.
        Let $W = X\setminus \closure{V}{X}{T}$.
        $V\subseteq X$ by \cref{open_in,topology_on,subseteq,pow_iff}.
        $\closure{V}{X}{T}$ is closed in $X$ with respect to $T$
            by \cref{open_in,topology_on,subseteq,pow_iff,closure_closed_in}.
        $W$ is open in $X$ with respect to $T$ and $B\subseteq X$ by \cref{closed_in}.
        Thus $B\subseteq W$
            by \cref{subseteq,double_relative_complement,subseteq_implies_setminus_supseteq}.
        We have $V\subseteq \closure{V}{X}{T}$ by \cref{subseteq_closure}.
        Hence $V$ is disjoint from $W$ by \cref{subseteq,setminus_elim_right,disjoint}.
        Show that there exist $U_0,V_0$ such that
            $U_0$ is open in $X$ with respect to $T$ and
            $V_0$ is open in $X$ with respect to $T$ and
            $A\subseteq U_0$ and
            $B\subseteq V_0$ and
            $U_0$ is disjoint from $V_0$.
        \begin{subproof}
            Trivial.
        \end{subproof}
    \end{subproof}
    Therefore $X$ is normal with respect to $T$ by \cref{one_point_closed,normal_space}.
\end{proof}

% from §31 Item 7: subspace of a Hausdorff space is Hausdorff
\begin{lemma}\label{subspace_hausdorff_31}
    Assume $X$ is Hausdorff with respect to $T$.
    Assume $Y$ is a subspace of $X$ with respect to $T$.
    Let $T_Y = \subspaceTopology{T}{Y}$.
    Then $Y$ is Hausdorff with respect to $T_Y$.
\end{lemma}
\begin{proof}
    Follows by \cref{subspace_of,subspace_hausdorff}.
\end{proof}

% from §31 Item 8: product of Hausdorff spaces is Hausdorff
\begin{lemma}\label{product_hausdorff_31}
    Assume $X$ is a function.
    Assume $T$ is a function.
    Assume $\dom{T} = \dom{X}$.
    Assume $\dom{X}$ is inhabited.
    Suppose for all $\alpha\in\dom{X}$ we have $\apply{X}{\alpha}$ is Hausdorff with respect to $\apply{T}{\alpha}$.
    Then $\productFamily{X}$ is Hausdorff with respect to $\productTopologyIndexed{T}{X}$.
\end{lemma}
\begin{proof}
    Follows by \cref{product_hausdorff_product}.
\end{proof}

% from §31 Item 9: subspace of a regular space is regular
\begin{lemma}\label{subspace_regular}
    Assume $X$ is regular with respect to $T$.
    Assume $Y$ is a subspace of $X$ with respect to $T$.
    Let $T_Y = \subspaceTopology{T}{Y}$.
    Then $Y$ is regular with respect to $T_Y$.
\end{lemma}
    \begin{proof}
    $T$ is a topology on $X$ and $Y\subseteq X$ by \cref{regular_space,subspace_of}.
    Show that $T_Y$ is a topology on $Y$.
    \begin{subproof}
        Follows by \cref{subspace_topology_is_topology,subspace_of}.
    \end{subproof}
    Show that $Y$ has closed singletons with respect to $T_Y$.
    \begin{subproof}
        Follows by \cref{regular_space,elem_subseteq,one_point_closed,singleton_subset_intro,inter_absorb_supseteq_right,closed_in_subspace_iff,subspace_of}.
    \end{subproof}
    Show that for all $x,B$ such that $x\in Y$ and $B$ is closed in $Y$ with respect to $T_Y$ and $x\notin B$
        there exist $U,V$ such that
            $U$ is open in $Y$ with respect to $T_Y$ and
            $V$ is open in $Y$ with respect to $T_Y$ and
            $x\in U$ and
            $B\subseteq V$ and
            $U$ is disjoint from $V$.
    \begin{subproof}
        Fix $x$.
        Fix $B$ such that $x\in Y$ and $B$ is closed in $Y$ with respect to $T_Y$ and $x\notin B$.
        Take $C$ such that $C$ is closed in $X$ with respect to $T$ and $B = C\inter Y$
            by \cref{closed_in_subspace_iff,subspace_of}.
        Take $U,V$ such that
            $U$ is open in $X$ with respect to $T$ and
            $V$ is open in $X$ with respect to $T$ and
            $x\in U$ and
            $C\subseteq V$ and
            $U$ is disjoint from $V$
            by \cref{inter_intro,elem_subseteq,regular_space}.
        Let $U_1 = U\inter Y$.
        Let $V_1 = V\inter Y$.
        $U_1$ is open in $Y$ with respect to $T_Y$ and $V_1$ is open in $Y$ with respect to $T_Y$
            and $x\in U_1$ by \cref{open_in,inter_comm,subspace_topology,inter_intro}.
        Thus $B\subseteq V_1$ and $U_1$ is disjoint from $V_1$
            by \cref{inter_lower_left,subseteq_transitive,inter_lower_right,subseteq_inter_iff,inter,disjoint}.
    \end{subproof}
    Hence $Y$ is regular with respect to $T_Y$ by \cref{regular_space}.
\end{proof}

% from §31 Item 10: product of regular spaces is regular
\begin{lemma}\label{product_regular}
    Assume $X$ is a function.
    Assume $T$ is a function.
    Assume $\dom{T} = \dom{X}$.
    Assume $\dom{X}$ is inhabited.
    Suppose for all $\alpha\in\dom{X}$ we have $\apply{X}{\alpha}$ is regular with respect to $\apply{T}{\alpha}$.
    Then $\productFamily{X}$ is regular with respect to $\productTopologyIndexed{T}{X}$.
\end{lemma}
\begin{proof}
    Let $T_0 = \productTopologyIndexed{T}{X}$.
    Let $S = \productSubbasis{T}{X}$.
    For all $\alpha\in\dom{X}$ we have $\apply{T}{\alpha}$ is a topology on $\apply{X}{\alpha}$ by \cref{regular_space}.
    Hence for all $\alpha\in\dom{X}$ we have $\apply{X}{\alpha}$ is Hausdorff with respect to $\apply{T}{\alpha}$
        by \cref{regular_implies_hausdorff}.
    Therefore $\productFamily{X}$ is Hausdorff with respect to $T_0$ by \cref{product_hausdorff_31}.
    Show that $T_0$ is a topology on $\productFamily{X}$.
    \begin{subproof}
        Follows by \cref{hausdorff}.
    \end{subproof}
    Show that for all $x$ such that $x\in\productFamily{X}$ we have $\{x\}$ is closed in $\productFamily{X}$ with respect to $T_0$.
    \begin{subproof}
        Follows by \cref{singleton_subset_intro,pair_is_finite,upair_intro_left,subseteq_finite_is_finite,finite_point_set_closed_hausdorff}.
    \end{subproof}
    Show that for all $x,U$ such that $x\in\productFamily{X}$ and
        $U$ is a neighborhood of $x$ in $\productFamily{X}$ with respect to $T_0$
        there exists $V$ such that
            $V$ is a neighborhood of $x$ in $\productFamily{X}$ with respect to $T_0$ and
            $\closure{V}{\productFamily{X}}{T_0}\subseteq U$.
    \begin{subproof}
        Fix $x$.
        Fix $U$.
        Assume $x\in\productFamily{X}$ and
            $U$ is a neighborhood of $x$ in $\productFamily{X}$ with respect to $T_0$.
        We have $U\in T_0$ and $x\in U$ by \cref{neighborhood,open_in}.
        $\finiteIntersections{S}$ is a basis for $T_0$ on $\productFamily{X}$ by \cref{product_subbasis_finite_intersections_basis}.
        Thus $T_0 = \basisTopology{\finiteIntersections{S}}{\productFamily{X}}$
            by \cref{product_topology_indexed,topology_generated_by_subbasis,basis_for_topology}.
        $U\in\basisTopology{\finiteIntersections{S}}{\productFamily{X}}$ by \cref{basis_for_topology}.
        Take $B_1\in\finiteIntersections{S}$ such that $x\in B_1$ and $B_1\subseteq U$ by \cref{topology_generated_by_basis}.
        Take $F\subseteq S$ such that $F$ is finite and inhabited and $B_1 = \inters{F}$ by \cref{finite_intersections}.
        Let $R = \{(A,B) \mid A\in F, B\in\pow{\productFamily{X}} \mid \exists \beta\in\dom{X}. \exists U_0\in\apply{T}{\beta}. \exists V_0\in\apply{T}{\beta}.
            A = \preimg{\piIndex{X}{\beta}}{U_0} \land
            \apply{x}{\beta}\in V_0 \land
            \closure{V_0}{\apply{X}{\beta}}{\apply{T}{\beta}}\subseteq U_0 \land
            B = \preimg{\piIndex{X}{\beta}}{V_0}\}$.
        $R$ is a relation by \cref{pair_eq_iff,relation}.
        Show that for all $A\in F$ there exists $B$ such that $A\mathrel{R} B$.
        \begin{subproof}
            Fix $A$ such that $A\in F$.
            Take $\beta\in\dom{X}$ such that $A\in\productSubbasisAt{T}{X}{\beta}$ by \cref{subseteq,product_subbasis_union_indexed}.
            Take $U_0\in\apply{T}{\beta}$ such that $A = \preimg{\piIndex{X}{\beta}}{U_0}$ by \cref{product_subbasis_indexed}.
            $x\in A$ by \cref{inters_subseteq_elem,subseteq}.
            Take $y\in U_0$ such that $x\mathrel{\piIndex{X}{\beta}} y$ by \cref{preimg_iff}.
            Take $x_0\in\productFamily{X}$ such that $(x,y) = (x_0,\apply{x_0}{\beta})$ by \cref{projection_map_indexed}.
            Thus $\apply{x}{\beta}\in U_0$ by \cref{pair_eq_iff}.
            Take $V_0$ such that
                $V_0$ is a neighborhood of $\apply{x}{\beta}$ in $\apply{X}{\beta}$ with respect to $\apply{T}{\beta}$ and
                $\closure{V_0}{\apply{X}{\beta}}{\apply{T}{\beta}}\subseteq U_0$
                by \cref{open_in,neighborhood,indexed_product,regular_space,regular_closure_neighborhood}.
            $V_0\in\apply{T}{\beta}$ and $\apply{x}{\beta}\in V_0$ by \cref{neighborhood,open_in}.
            Let $B = \preimg{\piIndex{X}{\beta}}{V_0}$.
            We have $B\subseteq\productFamily{X}$
                by \cref{projection_map_indexed_function,funs_is_relation,preimg_subseteq_dom,funs,subseteq_transitive}.
            Thus $B\in\pow{\productFamily{X}}$ by \cref{powerset_intro}.
            Hence $A\mathrel{R}B$ by \cref{pair_eq_iff}.
        \end{subproof}
        Take $g$ such that $g$ is a function on $F$ and for all $A\in F$ we have $A\mathrel{R}\apply{g}{A}$
            by \cref{choice_function}.
        Let $F_1 = \img{g}{F}$.
        Hence $F_1$ is finite and inhabited by \cref{finite_image_function,function_apply_elim,img_iff}.
        Show that for all $B$ such that $B\in F_1$ we have $B\in T_0$.
        \begin{subproof}
            Fix $B$ such that $B\in F_1$.
            Take $A$ such that $A\in F$ and $\apply{g}{A} = B$ by \cref{img_iff,function_apply_iff}.
            $A\mathrel{R}\apply{g}{A}$ by \cref{choice_function}.
            Thus $A\mathrel{R}B$.
            Take $\beta\in\dom{X}$ such that $B\in\pow{\productFamily{X}}$ and there exists $U_0\in\apply{T}{\beta}$ such that there exists $V_0\in\apply{T}{\beta}$ such that
                $A = \preimg{\piIndex{X}{\beta}}{U_0}$ and
                $\apply{x}{\beta}\in V_0$ and
                $\closure{V_0}{\apply{X}{\beta}}{\apply{T}{\beta}}\subseteq U_0$ and
                $B = \preimg{\piIndex{X}{\beta}}{V_0}$
                by \cref{pair_eq_iff}.
            Take $U_0\in\apply{T}{\beta}$ such that there exists $V_0\in\apply{T}{\beta}$ such that
                $A = \preimg{\piIndex{X}{\beta}}{U_0}$ and
                $\apply{x}{\beta}\in V_0$ and
                $\closure{V_0}{\apply{X}{\beta}}{\apply{T}{\beta}}\subseteq U_0$ and
                $B = \preimg{\piIndex{X}{\beta}}{V_0}$.
            Take $V_0\in\apply{T}{\beta}$ such that
                $A = \preimg{\piIndex{X}{\beta}}{U_0}$ and
                $\apply{x}{\beta}\in V_0$ and
                $\closure{V_0}{\apply{X}{\beta}}{\apply{T}{\beta}}\subseteq U_0$ and
                $B = \preimg{\piIndex{X}{\beta}}{V_0}$.
            Therefore $B\in T_0$ by \cref{projection_map_indexed_continuous,continuous,open_in}.
        \end{subproof}
        $F_1\subseteq T_0$ by \cref{subseteq}.
        Let $V = \inters{F_1}$.
        We have $V\in T_0$ and $V$ is open in $\productFamily{X}$ with respect to $T_0$
            by \cref{finite_intersections_open_in_topology,open_in}.
        Show that for all $B$ such that $B\in F_1$ we have $x\in B$.
        \begin{subproof}
            Fix $B$ such that $B\in F_1$.
            Take $A$ such that $A\in F$ and $\apply{g}{A} = B$ by \cref{img_iff,function_apply_iff}.
            $A\mathrel{R}\apply{g}{A}$ by \cref{choice_function}.
            Thus $A\mathrel{R}B$.
            Take $\beta\in\dom{X}$ such that $B\in\pow{\productFamily{X}}$ and there exists $U_0$ such that there exists $V_0$ such that
                $U_0\in\apply{T}{\beta}$ and
                $V_0\in\apply{T}{\beta}$ and
                $A = \preimg{\piIndex{X}{\beta}}{U_0}$ and
                $\apply{x}{\beta}\in V_0$ and
                $\closure{V_0}{\apply{X}{\beta}}{\apply{T}{\beta}}\subseteq U_0$ and
                $B = \preimg{\piIndex{X}{\beta}}{V_0}$
                by \cref{pair_eq_iff}.
            Take $U_0$ such that
                $U_0\in\apply{T}{\beta}$ and
                there exists $V_0$ such that
                    $V_0\in\apply{T}{\beta}$ and
                    $A = \preimg{\piIndex{X}{\beta}}{U_0}$ and
                    $\apply{x}{\beta}\in V_0$ and
                    $\closure{V_0}{\apply{X}{\beta}}{\apply{T}{\beta}}\subseteq U_0$ and
                    $B = \preimg{\piIndex{X}{\beta}}{V_0}$.
            Take $V_0$ such that
                $V_0\in\apply{T}{\beta}$ and
                $A = \preimg{\piIndex{X}{\beta}}{U_0}$ and
                $\apply{x}{\beta}\in V_0$ and
                $\closure{V_0}{\apply{X}{\beta}}{\apply{T}{\beta}}\subseteq U_0$ and
                $B = \preimg{\piIndex{X}{\beta}}{V_0}$.
            Therefore $x\in B$ by \cref{projection_map_indexed,preimg_iff}.
        \end{subproof}
        $x\in V$ by \cref{inters_iff_forall}.
        Thus $V$ is a neighborhood of $x$ in $\productFamily{X}$ with respect to $T_0$ by \cref{neighborhood}.
        Show that for all $A\in F$ we have $\closure{V}{\productFamily{X}}{T_0}\subseteq A$.
        \begin{subproof}
            Fix $A$ such that $A\in F$.
            Let $B = \apply{g}{A}$.
            We have $A\mathrel{R} B$ by \cref{choice_function}.
            Take $\beta\in\dom{X}$ such that $B\in\pow{\productFamily{X}}$ and there exists $U_0\in\apply{T}{\beta}$ such that there exists $V_0\in\apply{T}{\beta}$ such that
                $A = \preimg{\piIndex{X}{\beta}}{U_0}$ and
                $\apply{x}{\beta}\in V_0$ and
                $\closure{V_0}{\apply{X}{\beta}}{\apply{T}{\beta}}\subseteq U_0$ and
                $B = \preimg{\piIndex{X}{\beta}}{V_0}$
                by \cref{pair_eq_iff}.
            Take $U_0\in\apply{T}{\beta}$ such that there exists $V_0\in\apply{T}{\beta}$ such that
                $A = \preimg{\piIndex{X}{\beta}}{U_0}$ and
                $\apply{x}{\beta}\in V_0$ and
                $\closure{V_0}{\apply{X}{\beta}}{\apply{T}{\beta}}\subseteq U_0$ and
                $B = \preimg{\piIndex{X}{\beta}}{V_0}$.
            Take $V_0\in\apply{T}{\beta}$ such that
                $A = \preimg{\piIndex{X}{\beta}}{U_0}$ and
                $\apply{x}{\beta}\in V_0$ and
                $\closure{V_0}{\apply{X}{\beta}}{\apply{T}{\beta}}\subseteq U_0$ and
                $B = \preimg{\piIndex{X}{\beta}}{V_0}$.
            $B\in F_1$ and $V\subseteq B$ by \cref{img_iff,function_apply_elim,inters_subseteq_elem}.
            We have $V_0\subseteq\apply{X}{\beta}$ by \cref{regular_space,topology_on,subseteq,powerset_elim}.
            $\closure{V_0}{\apply{X}{\beta}}{\apply{T}{\beta}}$ is closed in $\apply{X}{\beta}$ with respect to $\apply{T}{\beta}$
                by \cref{regular_space,topology_on,subseteq,powerset_elim,closure_closed_in}.
            Let $D = \apply{X}{\beta}\setminus \closure{V_0}{\apply{X}{\beta}}{\apply{T}{\beta}}$.
            Therefore $\preimg{\piIndex{X}{\beta}}{D}$ is open in $\productFamily{X}$ with respect to $T_0$
                by \cref{closed_in,projection_map_indexed_continuous,continuous}.
            We have $\closure{V_0}{\apply{X}{\beta}}{\apply{T}{\beta}}\subseteq \apply{X}{\beta}$
                by \cref{closed_whole_space,closure_subseteq_closed}.
            Thus $\preimg{\piIndex{X}{\beta}}{D} = \productFamily{X}\setminus \preimg{\piIndex{X}{\beta}}{\closure{V_0}{\apply{X}{\beta}}{\apply{T}{\beta}}}$
                by \cref{projection_map_indexed_function,preimg_setminus_function}.
            We have $\productFamily{X}\setminus \preimg{\piIndex{X}{\beta}}{\closure{V_0}{\apply{X}{\beta}}{\apply{T}{\beta}}}$
                is open in $\productFamily{X}$ with respect to $T_0$ by \cref{open_in}.
            $\preimg{\piIndex{X}{\beta}}{\closure{V_0}{\apply{X}{\beta}}{\apply{T}{\beta}}}$
                is closed in $\productFamily{X}$ with respect to $T_0$
                by \cref{projection_map_indexed_function,funs_is_relation,preimg_subseteq_dom,funs,subseteq_transitive,closed_in}.
            Hence $B\subseteq \preimg{\piIndex{X}{\beta}}{\closure{V_0}{\apply{X}{\beta}}{\apply{T}{\beta}}}$
                by \cref{subseteq_closure,preimg_subseteq}.
            We have $V\subseteq \preimg{\piIndex{X}{\beta}}{\closure{V_0}{\apply{X}{\beta}}{\apply{T}{\beta}}}$
                by \cref{subseteq_transitive}.
            Hence $\closure{V}{\productFamily{X}}{T_0}\subseteq
                \preimg{\piIndex{X}{\beta}}{\closure{V_0}{\apply{X}{\beta}}{\apply{T}{\beta}}}$
                by \cref{hausdorff,topology_on,subseteq,powerset_elim,closure_subseteq_closed}.
            Therefore $\closure{V}{\productFamily{X}}{T_0}\subseteq \preimg{\piIndex{X}{\beta}}{U_0}$
                by \cref{preimg_subseteq,subseteq_transitive}.
            $\closure{V}{\productFamily{X}}{T_0}\subseteq A$ by \cref{subseteq}.
        \end{subproof}
        Hence $\closure{V}{\productFamily{X}}{T_0}\subseteq U$ by \cref{subseteq_inters_iff,subseteq,subseteq_transitive}.
    \end{subproof}
    Finally $\productFamily{X}$ is regular with respect to $T_0$ by \cref{regular_from_closure_neighborhood}.
\end{proof}

\section{§32 Normal Spaces}

% from §32 Item 1: regular space with countable basis is normal
\begin{theorem}\label{regular_second_countable_normal}
    Assume $X$ is regular with respect to $T$.
    Suppose $X$ satisfies the second countability axiom with respect to $T$.
    Then $X$ is normal with respect to $T$.
\end{theorem}
\begin{proof}
    Show that $T$ is a topology on $X$.
    \begin{subproof}
        Follows by \cref{regular_space}.
    \end{subproof}
    Take $b$ such that
        $b$ is a sequence in $\pow{X}$ and
        for all $n\in\naturalsPlus$ we have $b(n)$ is open in $X$ with respect to $T$ and
        for all $x\in X$ we have
            for all $U$ such that $U$ is open in $X$ with respect to $T$ and $x\in U$
                there exists $n\in\naturalsPlus$ such that $x\in b(n)$ and $b(n)\subseteq U$
        by \cref{second_countability_axiom}.
    Show that for all $A,B$ such that $A$ is closed in $X$ with respect to $T$ and $B$ is closed in $X$ with respect to $T$
        and $A$ is disjoint from $B$
        there exist $U,V$ such that
            $U$ is open in $X$ with respect to $T$ and
            $V$ is open in $X$ with respect to $T$ and
            $A\subseteq U$ and
            $B\subseteq V$ and
            $U$ is disjoint from $V$.
    \begin{subproof}
        Fix $A$.
        Fix $B$.
        Assume $A$ is closed in $X$ with respect to $T$ and $B$ is closed in $X$ with respect to $T$
            and $A$ is disjoint from $B$.
        Let $P = \{n\in\naturalsPlus \mid b(n)\inter A \neq \emptyset \land \closure{b(n)}{X}{T}\inter B = \emptyset\}$.
        Let $Q = \{n\in\naturalsPlus \mid b(n)\inter B \neq \emptyset \land \closure{b(n)}{X}{T}\inter A = \emptyset\}$.
        Let $U_1 = \{w \mid \exists n\in P. w = (n,b(n))\}$. Let $U_2 = \{w \mid \exists n\in\naturalsPlus\setminus P. w = (n,\emptyset)\}$.
        Let $U = U_1 \union U_2$.
        Let $V_1 = \{w \mid \exists n\in Q. w = (n,b(n))\}$. Let $V_2 = \{w \mid \exists n\in\naturalsPlus\setminus Q. w = (n,\emptyset)\}$.
        Let $V = V_1 \union V_2$.

        Show that $U$ is a function.
        \begin{subproof}
                For all $w$ such that $w\in U$ there exist $x,y$ such that $w = (x,y)$.
                Hence $U$ is a relation by \cref{relation}.
                Show that for all $n,y_1,y_2$ such that $(n,y_1)\in U$ and $(n,y_2)\in U$ we have $y_1 = y_2$.
                \begin{subproof}
                    Fix $n$.
                    Fix $y_1$.
                    Fix $y_2$ such that $(n,y_1)\in U$ and $(n,y_2)\in U$.
                    \begin{byCase}
                    \caseOf{$n\in P$.}
                        Show that $(n,y_1)\notin U_2$.
                        \begin{subproof}
                            Suppose not.
                            Take $n_1\in\naturalsPlus\setminus P$ such that $(n,y_1) = (n_1,\emptyset)$ by definition.
                            Therefore $n\notin P$ by \cref{pair_eq_iff,setminus_elim_right}.
                            Contradiction.
                        \end{subproof}
                        We have $(n,y_1)\in U_1$ by \cref{union_iff}.
                        Take $n_1\in P$ such that $(n,y_1) = (n_1,b(n_1))$ by definition.
                        Show that $(n,y_2)\notin U_2$.
                        \begin{subproof}
                            Suppose not.
                            Take $n_2\in\naturalsPlus\setminus P$ such that $(n,y_2) = (n_2,\emptyset)$ by definition.
                            Therefore $n\notin P$ by \cref{pair_eq_iff,setminus_elim_right}.
                            Contradiction.
                        \end{subproof}
                        Hence $(n,y_2)\in U_1$ by \cref{union_iff}.
                        Take $n_2\in P$ such that $(n,y_2) = (n_2,b(n_2))$ by definition.
                        Hence $y_1 = y_2$ by \cref{pair_eq_iff}.
                    \caseOf{$n\notin P$.}
                        Show that $(n,y_1)\notin U_1$.
                        \begin{subproof}
                            Suppose not.
                            Take $n_1\in P$ such that $(n,y_1) = (n_1,b(n_1))$ by definition.
                            Therefore $n\in P$ by \cref{pair_eq_iff}.
                            Contradiction.
                        \end{subproof}
                        We have $(n,y_1)\in U_2$ by \cref{union_iff}.
                        Take $n_1\in\naturalsPlus\setminus P$ such that $(n,y_1) = (n_1,\emptyset)$ by definition.
                        Show that $(n,y_2)\notin U_1$.
                        \begin{subproof}
                            Suppose not.
                            Take $n_2\in P$ such that $(n,y_2) = (n_2,b(n_2))$ by definition.
                            Therefore $n\in P$ by \cref{pair_eq_iff}.
                            Contradiction.
                        \end{subproof}
                        Hence $(n,y_2)\in U_2$ by \cref{union_iff}.
                        Take $n_2\in\naturalsPlus\setminus P$ such that $(n,y_2) = (n_2,\emptyset)$ by definition.
                        Hence $y_1 = y_2$ by \cref{pair_eq_iff}.
                    \end{byCase}
                \end{subproof}
                $U$ is right-unique by \cref{rightunique}.
        \end{subproof}
        For all $n$ such that $n\in\dom{U}$ we have $n\in\naturalsPlus$
            by \cref{dom_iff,union_iff,pair_eq_iff,setminus_elim_left}.
        We have $\dom{U}\subseteq\naturalsPlus$ by \cref{subseteq}.
        $\naturalsPlus\subseteq\dom{U}$
            by \cref{union_intro_left,setminus_intro,union_intro_right,dom_intro,subseteq}.
        Thus $\dom{U} = \naturalsPlus$ by \cref{subseteq_antisymmetric}.
        Hence $U$ is a function on $\naturalsPlus$.

        Show that for all $n\in\dom{U}$ we have $U(n)\in\pow{X}$.
        \begin{subproof}
            Fix $n$ such that $n\in\dom{U}$.
            $\dom{U} = \naturalsPlus$.
            Hence $n\in\naturalsPlus$ by \cref{subseteq}.
            \begin{byCase}
            \caseOf{$n\in P$.}
                $(n,U(n))\in U$ by \cref{function_apply_elim}.
                Show that $(n,U(n))\notin U_2$.
                \begin{subproof}
                    Suppose not.
                    Take $n_0\in\naturalsPlus\setminus P$ such that $(n,U(n)) = (n_0,\emptyset)$ by definition.
                    Therefore $n\notin P$ by \cref{pair_eq_iff,setminus_elim_right}.
                    Contradiction.
                \end{subproof}
                We have $(n,U(n))\in U_1$ by \cref{union_iff}.
                Take $n_0\in P$ such that $(n,U(n)) = (n_0,b(n_0))$ by definition.
                Therefore $U(n)\in\pow{X}$ by \cref{pair_eq_iff,sequence_in,funs_type_apply}.
            \caseOf{$n\notin P$.}
                $(n,U(n))\in U$ by \cref{function_apply_elim}.
                Show that $(n,U(n))\notin U_1$.
                \begin{subproof}
                    Suppose not.
                    Take $n_0\in P$ such that $(n,U(n)) = (n_0,b(n_0))$ by definition.
                    Therefore $n\in P$ by \cref{pair_eq_iff}.
                    Contradiction.
                \end{subproof}
                Hence $(n,U(n))\in U_2$ by \cref{union_iff}.
                Take $n_0\in\naturalsPlus\setminus P$ such that $(n,U(n)) = (n_0,\emptyset)$ by definition.
                Therefore $U(n)\in\pow{X}$ by \cref{pair_eq_iff,emptyset_subseteq,powerset_intro}.
            \end{byCase}
        \end{subproof}
        $U$ is a function to $\pow{X}$.
        Therefore $U\in\funs{\naturalsPlus}{\pow{X}}$ by \cref{funs_intro}.

        Show that $V$ is a function.
        \begin{subproof}
                For all $w$ such that $w\in V$ there exist $x,y$ such that $w = (x,y)$.
                Hence $V$ is a relation by \cref{relation}.
                Show that for all $n,y_1,y_2$ such that $(n,y_1)\in V$ and $(n,y_2)\in V$ we have $y_1 = y_2$.
                \begin{subproof}
                    Fix $n$.
                    Fix $y_1$.
                    Fix $y_2$ such that $(n,y_1)\in V$ and $(n,y_2)\in V$.
                    \begin{byCase}
                    \caseOf{$n\in Q$.}
                        Show that $(n,y_1)\notin V_2$.
                        \begin{subproof}
                            Suppose not.
                            Take $n_1\in\naturalsPlus\setminus Q$ such that $(n,y_1) = (n_1,\emptyset)$ by definition.
                            Therefore $n\notin Q$ by \cref{pair_eq_iff,setminus_elim_right}.
                            Contradiction.
                        \end{subproof}
                        We have $(n,y_1)\in V_1$ by \cref{union_iff}.
                        Take $n_1\in Q$ such that $(n,y_1) = (n_1,b(n_1))$ by definition.
                        Show that $(n,y_2)\notin V_2$.
                        \begin{subproof}
                            Suppose not.
                            Take $n_2\in\naturalsPlus\setminus Q$ such that $(n,y_2) = (n_2,\emptyset)$ by definition.
                            Therefore $n\notin Q$ by \cref{pair_eq_iff,setminus_elim_right}.
                            Contradiction.
                        \end{subproof}
                        Hence $(n,y_2)\in V_1$ by \cref{union_iff}.
                        Take $n_2\in Q$ such that $(n,y_2) = (n_2,b(n_2))$ by definition.
                        Hence $y_1 = y_2$ by \cref{pair_eq_iff}.
                    \caseOf{$n\notin Q$.}
                        Show that $(n,y_1)\notin V_1$.
                        \begin{subproof}
                            Suppose not.
                            Take $n_1\in Q$ such that $(n,y_1) = (n_1,b(n_1))$ by definition.
                            Therefore $n\in Q$ by \cref{pair_eq_iff}.
                            Contradiction.
                        \end{subproof}
                        We have $(n,y_1)\in V_2$ by \cref{union_iff}.
                        Take $n_1\in\naturalsPlus\setminus Q$ such that $(n,y_1) = (n_1,\emptyset)$ by definition.
                        Show that $(n,y_2)\notin V_1$.
                        \begin{subproof}
                            Suppose not.
                            Take $n_2\in Q$ such that $(n,y_2) = (n_2,b(n_2))$ by definition.
                            Therefore $n\in Q$ by \cref{pair_eq_iff}.
                            Contradiction.
                        \end{subproof}
                        Hence $(n,y_2)\in V_2$ by \cref{union_iff}.
                        Take $n_2\in\naturalsPlus\setminus Q$ such that $(n,y_2) = (n_2,\emptyset)$ by definition.
                        Hence $y_1 = y_2$ by \cref{pair_eq_iff}.
                    \end{byCase}
                \end{subproof}
                $V$ is right-unique by \cref{rightunique}.
        \end{subproof}
        For all $n\in\dom{V}$ we have $n\in\naturalsPlus$
            by \cref{dom_iff,union_iff,pair_eq_iff,setminus_elim_left}.
        We have $\dom{V}\subseteq\naturalsPlus$ by \cref{subseteq}.
        $\naturalsPlus\subseteq\dom{V}$
            by \cref{union_intro_left,setminus_intro,union_intro_right,dom_intro,subseteq}.
        Thus $\dom{V} = \naturalsPlus$ by \cref{subseteq_antisymmetric}.
        Hence $V$ is a function on $\naturalsPlus$.

        Show that for all $n\in\dom{V}$ we have $V(n)\in\pow{X}$.
        \begin{subproof}
            Fix $n$ such that $n\in\dom{V}$.
            $\dom{V} = \naturalsPlus$.
            Hence $n\in\naturalsPlus$ by \cref{subseteq}.
            \begin{byCase}
            \caseOf{$n\in Q$.}
                $(n,V(n))\in V$ by \cref{function_apply_elim}.
                Show that $(n,V(n))\notin V_2$.
                \begin{subproof}
                    Suppose not.
                    Take $n_0\in\naturalsPlus\setminus Q$ such that $(n,V(n)) = (n_0,\emptyset)$ by definition.
                    Therefore $n\notin Q$ by \cref{pair_eq_iff,setminus_elim_right}.
                    Contradiction.
                \end{subproof}
                We have $(n,V(n))\in V_1$ by \cref{union_iff}.
                Take $n_0\in Q$ such that $(n,V(n)) = (n_0,b(n_0))$ by definition.
                Therefore $V(n)\in\pow{X}$ by \cref{pair_eq_iff,sequence_in,funs_type_apply}.
            \caseOf{$n\notin Q$.}
                $(n,V(n))\in V$ by \cref{function_apply_elim}.
                Show that $(n,V(n))\notin V_1$.
                \begin{subproof}
                    Suppose not.
                    Take $n_0\in Q$ such that $(n,V(n)) = (n_0,b(n_0))$ by definition.
                    Therefore $n\in Q$ by \cref{pair_eq_iff}.
                    Contradiction.
                \end{subproof}
                Hence $(n,V(n))\in V_2$ by \cref{union_iff}.
                Take $n_0\in\naturalsPlus\setminus Q$ such that $(n,V(n)) = (n_0,\emptyset)$ by definition.
                Therefore $V(n)\in\pow{X}$ by \cref{pair_eq_iff,emptyset_subseteq,powerset_intro}.
            \end{byCase}
        \end{subproof}
        $V$ is a function to $\pow{X}$.
        Therefore $V\in\funs{\naturalsPlus}{\pow{X}}$ by \cref{funs_intro}.

        Show that for all $n\in\naturalsPlus$ we have $U(n)$ is open in $X$ with respect to $T$ and
            $\closure{U(n)}{X}{T}$ is disjoint from $B$.
        \begin{subproof}
            Fix $n\in\naturalsPlus$.
            \begin{byCase}
            \caseOf{$n\in P$.}
                $b(n)$ is open in $X$ with respect to $T$ by \cref{second_countability_axiom}.
                We have $U(n) = b(n)$ by \cref{union_intro_left,function_apply_elim,function_rightunique}.
                $U(n)$ is open in $X$ with respect to $T$.
                Hence $\closure{U(n)}{X}{T}$ is disjoint from $B$ by \cref{inter,emptyset,disjoint}.
            \caseOf{$n\notin P$.}
                We have $U(n) = \emptyset$ by \cref{setminus_intro,union_intro_right,function_apply_elim,function_rightunique}.
                $U(n)$ is open in $X$ with respect to $T$ by \cref{rightunique,topology_on,open_in}.
                We have $\closure{U(n)}{X}{T} = \emptyset$
                    by \cref{closed_emptyset,emptyset_subseteq,subseteq_emptyset_iff,closure_subseteq_closed}.
                Hence $\closure{U(n)}{X}{T}$ is disjoint from $B$ by \cref{emptyset,disjoint}.
            \end{byCase}
        \end{subproof}

        Show that for all $n\in\naturalsPlus$ we have $V(n)$ is open in $X$ with respect to $T$ and
            $\closure{V(n)}{X}{T}$ is disjoint from $A$.
        \begin{subproof}
            Fix $n\in\naturalsPlus$.
            \begin{byCase}
            \caseOf{$n\in Q$.}
                $b(n)$ is open in $X$ with respect to $T$ by \cref{second_countability_axiom}.
                We have $V(n) = b(n)$ by \cref{union_intro_left,function_apply_elim,function_rightunique}.
                $V(n)$ is open in $X$ with respect to $T$.
                Hence $\closure{V(n)}{X}{T}$ is disjoint from $A$ by \cref{inter,emptyset,disjoint}.
            \caseOf{$n\notin Q$.}
                We have $V(n) = \emptyset$ by \cref{setminus_intro,union_intro_right,function_apply_elim,function_rightunique}.
                $V(n)$ is open in $X$ with respect to $T$ by \cref{rightunique,topology_on,open_in}.
                We have $\closure{V(n)}{X}{T} = \emptyset$
                    by \cref{closed_emptyset,emptyset_subseteq,subseteq_emptyset_iff,closure_subseteq_closed}.
                Hence $\closure{V(n)}{X}{T}$ is disjoint from $A$ by \cref{emptyset,disjoint}.
            \end{byCase}
        \end{subproof}

        Show that for all $x$ such that $x\in A$ we have $x\in\unions{\img{U}{\naturalsPlus}}$.
        \begin{subproof}
                Fix $x$ such that $x\in A$.
                Take $U_0,V_0$ such that
                    $U_0$ is open in $X$ with respect to $T$ and
                    $V_0$ is open in $X$ with respect to $T$ and
                    $x\in U_0$ and
                    $B\subseteq V_0$ and
                    $U_0$ is disjoint from $V_0$
                    by \cref{closed_in,subseteq,disjoint,regular_space}.
                We have $x\in X$ by \cref{closed_in,subseteq}.
                $1\in\naturalsPlus$ by \cref{real_one_in_naturals}.
                For all $y\in X$ we have
                    for all $W$ such that $W$ is open in $X$ with respect to $T$ and $y\in W$
                        there exists $n\in\naturalsPlus$ such that $y\in b(n)$ and $b(n)\subseteq W$
                    by \cref{second_countability_axiom}.
                Take $n\in\naturalsPlus$ such that $x\in b(n)$ and $b(n)\subseteq U_0$
                    by \cref{second_countability_axiom}.
                Hence $b(n)\inter A\neq\emptyset$ by \cref{inter_intro,emptyset}.
                Therefore $X\setminus V_0$ is closed in $X$ with respect to $T$
                    by \cref{open_in,topology_on,subseteq,powerset_elim,setminus_subseteq,double_relative_complement,closed_in}.
                Hence $U_0\subseteq X\setminus V_0$ by \cref{disjoint,subseteq,open_in,topology_on,powerset_elim,setminus_intro,subseteq}.
                Therefore $\closure{b(n)}{X}{T}\subseteq X\setminus V_0$
                    by \cref{subseteq_transitive,sequence_in,funs_type_apply,powerset_elim,closure_subseteq_closed}.
                Hence $\closure{b(n)}{X}{T}$ is disjoint from $B$ by \cref{subseteq,setminus_elim_right,disjoint}.
                Show that $\closure{b(n)}{X}{T}\inter B = \emptyset$.
                \begin{subproof}
                    Suppose not.
                    Take $y$ such that $y\in\closure{b(n)}{X}{T}\inter B$ by \cref{nonempty_inhabited}.
                    Contradiction by \cref{inter,disjoint}.
                \end{subproof}
                Thus $n\in P$.
                Hence $(n,b(n))\in U$ by \cref{union_intro_left}.
                $(n,U(n))\in U$ by \cref{function_apply_elim}.
                $x\in U(n)$ by \cref{function_rightunique}.
                Hence $x\in\unions{\img{U}{\naturalsPlus}}$ by \cref{unions_intro,img_iff}.
        \end{subproof}
        Therefore $A\subseteq\unions{\img{U}{\naturalsPlus}}$ by \cref{subseteq}.

        Show that for all $x$ such that $x\in B$ we have $x\in\unions{\img{V}{\naturalsPlus}}$.
        \begin{subproof}
                Fix $x$ such that $x\in B$.
                Take $U_0,V_0$ such that
                    $U_0$ is open in $X$ with respect to $T$ and
                    $V_0$ is open in $X$ with respect to $T$ and
                    $x\in U_0$ and
                    $A\subseteq V_0$ and
                    $U_0$ is disjoint from $V_0$
                    by \cref{closed_in,subseteq,disjoint,regular_space}.
                We have $x\in X$ by \cref{closed_in,subseteq}.
                $1\in\naturalsPlus$ by \cref{real_one_in_naturals}.
                For all $y\in X$ we have
                    for all $W$ such that $W$ is open in $X$ with respect to $T$ and $y\in W$
                        there exists $n\in\naturalsPlus$ such that $y\in b(n)$ and $b(n)\subseteq W$
                    by \cref{second_countability_axiom}.
                Take $n\in\naturalsPlus$ such that $x\in b(n)$ and $b(n)\subseteq U_0$
                    by \cref{second_countability_axiom}.
                Hence $b(n)\inter B\neq\emptyset$ by \cref{inter_intro,emptyset}.
                Therefore $X\setminus V_0$ is closed in $X$ with respect to $T$
                    by \cref{open_in,topology_on,subseteq,powerset_elim,setminus_subseteq,double_relative_complement,closed_in}.
                Hence $U_0\subseteq X\setminus V_0$ by \cref{disjoint,subseteq,open_in,topology_on,powerset_elim,setminus_intro,subseteq}.
                Therefore $\closure{b(n)}{X}{T}\subseteq X\setminus V_0$
                    by \cref{subseteq_transitive,sequence_in,funs_type_apply,powerset_elim,closure_subseteq_closed}.
                Hence $\closure{b(n)}{X}{T}$ is disjoint from $A$ by \cref{subseteq,setminus_elim_right,disjoint}.
                Show that $\closure{b(n)}{X}{T}\inter A = \emptyset$.
                \begin{subproof}
                    Suppose not.
                    Take $y$ such that $y\in\closure{b(n)}{X}{T}\inter A$ by \cref{nonempty_inhabited}.
                    Contradiction by \cref{inter,disjoint}.
                \end{subproof}
                Thus $n\in Q$.
                Hence $(n,b(n))\in V$ by \cref{union_intro_left}.
                $(n,V(n))\in V$ by \cref{function_apply_elim}.
                $x\in V(n)$ by \cref{function_rightunique}.
                Hence $x\in\unions{\img{V}{\naturalsPlus}}$ by \cref{unions_intro,img_iff}.
        \end{subproof}
        Therefore $B\subseteq\unions{\img{V}{\naturalsPlus}}$ by \cref{subseteq}.

        Let $c = \{w \mid \exists n\in\naturalsPlus. w = (n,\closure{V(n)}{X}{T})\}$. Let $d = \{w \mid \exists n\in\naturalsPlus. w = (n,\closure{U(n)}{X}{T})\}$.
        Show that $c$ is a function on $\naturalsPlus$.
        \begin{subproof}
            For all $w$ such that $w\in c$ there exist $x,y$ such that $w = (x,y)$.
            Hence $c$ is a relation by \cref{relation}.
            Show that for all $n,y_1,y_2$ such that $(n,y_1)\in c$ and $(n,y_2)\in c$ we have $y_1 = y_2$.
            \begin{subproof}
                Fix $n$.
                Fix $y_1$.
                Fix $y_2$ such that $(n,y_1)\in c$ and $(n,y_2)\in c$.
                Take $n_1\in\naturalsPlus$ such that $(n,y_1) = (n_1,\closure{V(n_1)}{X}{T})$ by definition.
                Take $n_2\in\naturalsPlus$ such that $(n,y_2) = (n_2,\closure{V(n_2)}{X}{T})$ by definition.
                Hence $y_1 = \closure{V(n)}{X}{T}$ and $y_2 = \closure{V(n)}{X}{T}$ by \cref{pair_eq_iff}.
                Therefore $y_1 = y_2$.
            \end{subproof}
            Hence $c$ is right-unique by \cref{rightunique}.
            Show that for all $n$ such that $n\in\dom{c}$ we have $n\in\naturalsPlus$.
            \begin{subproof}
                Fix $n$ such that $n\in\dom{c}$.
                Take $y$ such that $(n,y)\in c$ by \cref{dom_iff}.
                Take $n_0\in\naturalsPlus$ such that $(n,y) = (n_0,\closure{V(n_0)}{X}{T})$ by definition.
                Thus $n\in\naturalsPlus$ by \cref{pair_eq_iff}.
            \end{subproof}
            Hence $\dom{c}\subseteq\naturalsPlus$ by \cref{subseteq}.
            Show that for all $n$ such that $n\in\naturalsPlus$ we have $n\in\dom{c}$.
            \begin{subproof}
                Fix $n$ such that $n\in\naturalsPlus$.
                We have $(n,\closure{V(n)}{X}{T})\in c$ by definition.
                Thus $n\in\dom{c}$ by \cref{dom_intro}.
            \end{subproof}
            Hence $\naturalsPlus\subseteq\dom{c}$ by \cref{subseteq}.
            Therefore $\dom{c} = \naturalsPlus$ by \cref{subseteq_antisymmetric}.
        \end{subproof}

        Show that $d$ is a function on $\naturalsPlus$.
        \begin{subproof}
            For all $w$ such that $w\in d$ there exist $x,y$ such that $w = (x,y)$.
            Hence $d$ is a relation by \cref{relation}.
            Show that for all $n,y_1,y_2$ such that $(n,y_1)\in d$ and $(n,y_2)\in d$ we have $y_1 = y_2$.
            \begin{subproof}
                Fix $n$.
                Fix $y_1$.
                Fix $y_2$ such that $(n,y_1)\in d$ and $(n,y_2)\in d$.
                Take $n_1\in\naturalsPlus$ such that $(n,y_1) = (n_1,\closure{U(n_1)}{X}{T})$ by definition.
                Take $n_2\in\naturalsPlus$ such that $(n,y_2) = (n_2,\closure{U(n_2)}{X}{T})$ by definition.
                Hence $y_1 = \closure{U(n)}{X}{T}$ and $y_2 = \closure{U(n)}{X}{T}$ by \cref{pair_eq_iff}.
                Therefore $y_1 = y_2$.
            \end{subproof}
            Hence $d$ is right-unique by \cref{rightunique}.
            Show that for all $n$ such that $n\in\dom{d}$ we have $n\in\naturalsPlus$.
            \begin{subproof}
                Fix $n$ such that $n\in\dom{d}$.
                Take $y$ such that $(n,y)\in d$ by \cref{dom_iff}.
                Take $n_0\in\naturalsPlus$ such that $(n,y) = (n_0,\closure{U(n_0)}{X}{T})$ by definition.
                Thus $n\in\naturalsPlus$ by \cref{pair_eq_iff}.
            \end{subproof}
            Hence $\dom{d}\subseteq\naturalsPlus$ by \cref{subseteq}.
            Show that for all $n$ such that $n\in\naturalsPlus$ we have $n\in\dom{d}$.
            \begin{subproof}
                Fix $n$ such that $n\in\naturalsPlus$.
                We have $(n,\closure{U(n)}{X}{T})\in d$ by definition.
                Thus $n\in\dom{d}$ by \cref{dom_intro}.
            \end{subproof}
            Hence $\naturalsPlus\subseteq\dom{d}$ by \cref{subseteq}.
            Therefore $\dom{d} = \naturalsPlus$ by \cref{subseteq_antisymmetric}.
        \end{subproof}

        Let $F = \{W \mid \exists n\in\naturalsPlus. W = U(n)\setminus \unions{\img{\restrl{c}{\initialSegment{n}}}{\initialSegment{n}}}\}$. Let $G = \{W \mid \exists n\in\naturalsPlus. W = V(n)\setminus \unions{\img{\restrl{d}{\initialSegment{n}}}{\initialSegment{n}}}\}$.
        Show that for all $W$ such that $W\in F$ we have $W\in T$.
        \begin{subproof}
            Fix $W$ such that $W\in F$.
                Take $n\in\naturalsPlus$ such that $W = U(n)\setminus \unions{\img{\restrl{c}{\initialSegment{n}}}{\initialSegment{n}}}$ by definition.
                $\restrl{c}{\initialSegment{n}}$ is a function by \cref{restrl_of_function_is_function}.
                $\initialSegment{n}$ is finite by \cref{initial_segment_finite}.
                Hence $\dom{\restrl{c}{\initialSegment{n}}} = \initialSegment{n}$
                    by \cref{dom_of_restrl_eq_inter,initial_segment_naturalsplus,subseteq,inter_absorb_supseteq_left}.
                Therefore $\img{\restrl{c}{\initialSegment{n}}}{\initialSegment{n}}$ is finite by \cref{finite_image_function}.
                For all $C$ such that $C\in \img{\restrl{c}{\initialSegment{n}}}{\initialSegment{n}}$ we have
                    $C$ is closed in $X$ with respect to $T$ by
                    \cref{img_iff,pair_eq_iff,restrl_iff,funs_type_apply,powerset_elim,closure_closed_in}.
                Hence $\img{\restrl{c}{\initialSegment{n}}}{\initialSegment{n}}\subseteq\pow{X}$
                    by \cref{closed_in,powerset_intro,subseteq}.
                Therefore $\unions{\img{\restrl{c}{\initialSegment{n}}}{\initialSegment{n}}}$ is closed in $X$ with respect to $T$
                    by \cref{closed_finite_unions}.
                Hence $U(n)\in T$ and $U(n)\subseteq X$ by \cref{open_in,topology_on,subseteq,powerset_elim}.
                Therefore $W\in T$
                    by \cref{closed_in,setminus_eq_inter_complement,open_in,topology_on}.
        \end{subproof}
        Therefore $F\subseteq T$ by \cref{subseteq}.

        Show that for all $W$ such that $W\in G$ we have $W\in T$.
        \begin{subproof}
            Fix $W$ such that $W\in G$.
                Take $n\in\naturalsPlus$ such that $W = V(n)\setminus \unions{\img{\restrl{d}{\initialSegment{n}}}{\initialSegment{n}}}$ by definition.
                $\restrl{d}{\initialSegment{n}}$ is a function by \cref{restrl_of_function_is_function}.
                $\initialSegment{n}$ is finite by \cref{initial_segment_finite}.
                $d$ is a function on $\naturalsPlus$.
                Hence $\dom{\restrl{d}{\initialSegment{n}}} = \initialSegment{n}$
                    by \cref{dom_of_restrl_eq_inter,initial_segment_naturalsplus,subseteq,inter_absorb_supseteq_left}.
                Therefore $\img{\restrl{d}{\initialSegment{n}}}{\initialSegment{n}}$ is finite by \cref{finite_image_function}.
                For all $C$ such that $C\in \img{\restrl{d}{\initialSegment{n}}}{\initialSegment{n}}$ we have
                    $C$ is closed in $X$ with respect to $T$ by
                    \cref{img_iff,pair_eq_iff,restrl_iff,funs_type_apply,powerset_elim,closure_closed_in}.
                Hence $\img{\restrl{d}{\initialSegment{n}}}{\initialSegment{n}}\subseteq\pow{X}$
                    by \cref{closed_in,powerset_intro,subseteq}.
                Therefore $\unions{\img{\restrl{d}{\initialSegment{n}}}{\initialSegment{n}}}$ is closed in $X$ with respect to $T$
                    by \cref{closed_finite_unions}.
                Hence $V(n)\in T$ and $V(n)\subseteq X$ by \cref{open_in,topology_on,subseteq,powerset_elim}.
                Therefore $W\in T$
                    by \cref{closed_in,setminus_eq_inter_complement,open_in,topology_on}.
        \end{subproof}
        Therefore $G\subseteq T$ by \cref{subseteq}.

        Let $U_3 = \unions{F}$. Let $V_3 = \unions{G}$.
        Show that $U_3$ is open in $X$ with respect to $T$.
        \begin{subproof}
            Follows by \cref{topology_on,subseteq,open_in}.
        \end{subproof}
        Show that $V_3$ is open in $X$ with respect to $T$.
        \begin{subproof}
            Follows by \cref{topology_on,subseteq,open_in}.
        \end{subproof}

        Show that for all $x$ such that $x\in A$ we have $x\in U_3$.
        \begin{subproof}
            Fix $x$ such that $x\in A$.
                Take $W$ such that $W\in \img{U}{\naturalsPlus}$ and $x\in W$ by \cref{subseteq,unions_iff}.
                Take $n\in\naturalsPlus$ such that $n\mathrel{U} W$ by \cref{img_iff}.
                $(n,U(n))\in U$ by \cref{function_apply_elim}.
                Hence $x\in U(n)$ by \cref{img_iff,pair_eq_iff,function_rightunique}.
                Show that $x\notin \unions{\img{\restrl{c}{\initialSegment{n}}}{\initialSegment{n}}}$.
                \begin{subproof}
                    Suppose not.
                    There exists $C$ such that $C\in \img{\restrl{c}{\initialSegment{n}}}{\initialSegment{n}}$ and $x\in C$ by \cref{unions_iff}.
                    Take $k\in\initialSegment{n}$ such that $k\mathrel{\restrl{c}{\initialSegment{n}}} C$ by \cref{img_iff}.
                    Thus $(k,C)\in c$ by \cref{pair_eq_iff,restrl_iff}.
                    Take $k_0\in\naturalsPlus$ such that $(k,C) = (k_0,\closure{V(k_0)}{X}{T})$ by definition.
                    Contradiction by \cref{pair_eq_iff,disjoint}.
                \end{subproof}
                Therefore $x\in U(n)\setminus \unions{\img{\restrl{c}{\initialSegment{n}}}{\initialSegment{n}}}$ by \cref{setminus_intro}.
                Hence $U(n)\setminus \unions{\img{\restrl{c}{\initialSegment{n}}}{\initialSegment{n}}}\in F$ by definition.
                Therefore $x\in U_3$ by \cref{unions_intro}.
        \end{subproof}
        Hence $A\subseteq U_3$ by \cref{subseteq}.

        Show that for all $x$ such that $x\in B$ we have $x\in V_3$.
        \begin{subproof}
            Fix $x$ such that $x\in B$.
                Take $W$ such that $W\in \img{V}{\naturalsPlus}$ and $x\in W$ by \cref{subseteq,unions_iff}.
                Take $n\in\naturalsPlus$ such that $n\mathrel{V} W$ by \cref{img_iff}.
                $(n,V(n))\in V$ by \cref{function_apply_elim}.
                Hence $x\in V(n)$ by \cref{img_iff,pair_eq_iff,function_rightunique}.
                Show that $x\notin \unions{\img{\restrl{d}{\initialSegment{n}}}{\initialSegment{n}}}$.
                \begin{subproof}
                    Suppose not.
                    There exists $C$ such that $C\in \img{\restrl{d}{\initialSegment{n}}}{\initialSegment{n}}$ and $x\in C$ by \cref{unions_iff}.
                    Take $k\in\initialSegment{n}$ such that $k\mathrel{\restrl{d}{\initialSegment{n}}} C$ by \cref{img_iff}.
                    Thus $(k,C)\in d$ by \cref{pair_eq_iff,restrl_iff}.
                    Take $k_0\in\naturalsPlus$ such that $(k,C) = (k_0,\closure{U(k_0)}{X}{T})$ by definition.
                    Contradiction by \cref{pair_eq_iff,disjoint}.
                \end{subproof}
                Therefore $x\in V(n)\setminus \unions{\img{\restrl{d}{\initialSegment{n}}}{\initialSegment{n}}}$ by \cref{setminus_intro}.
                Hence $V(n)\setminus \unions{\img{\restrl{d}{\initialSegment{n}}}{\initialSegment{n}}}\in G$ by definition.
                Therefore $x\in V_3$ by \cref{unions_intro}.
        \end{subproof}
        Hence $B\subseteq V_3$ by \cref{subseteq}.

        Show that $U_3$ is disjoint from $V_3$.
        \begin{subproof}
            Suppose not.
            Therefore there exists $x$ such that $x\in U_3$ and $x\in V_3$ by \cref{disjoint}.
            Take $W_1$ such that $W_1\in F$ and $x\in W_1$ by \cref{unions_iff}.
            Take $W_2$ such that $W_2\in G$ and $x\in W_2$ by \cref{unions_iff}.
            Take $j\in\naturalsPlus$ such that $W_1 = U(j)\setminus \unions{\img{\restrl{c}{\initialSegment{j}}}{\initialSegment{j}}}$.
            Take $k\in\naturalsPlus$ such that $W_2 = V(k)\setminus \unions{\img{\restrl{d}{\initialSegment{k}}}{\initialSegment{k}}}$.
            $x\in U(j)$ and $x\in V(k)$ by \cref{setminus_elim_left}.
            $x\notin \unions{\img{\restrl{c}{\initialSegment{j}}}{\initialSegment{j}}}$ and
                $x\notin \unions{\img{\restrl{d}{\initialSegment{k}}}{\initialSegment{k}}}$ by \cref{setminus_elim_right}.
            \begin{byCase}
            \caseOf{$j\leq k$.}
                $j\in\initialSegment{k}$ by \cref{initial_segment_naturalsplus}.
                Take $C$ such that $C = \closure{U(j)}{X}{T}$.
                $j\mathrel{\restrl{d}{\initialSegment{k}}} C$ by \cref{restrl_iff}.
                Hence $C\in \img{\restrl{d}{\initialSegment{k}}}{\initialSegment{k}}$ by \cref{img_iff}.
                Therefore $x\in C$ by \cref{funs_type_apply,powerset_elim,subseteq_closure,elem_subseteq}.
                Hence $x\in \unions{\img{\restrl{d}{\initialSegment{k}}}{\initialSegment{k}}}$ by \cref{unions_intro}.
                Contradiction.
            \caseOf{$k\leq j$.}
                $k\in\initialSegment{j}$ by \cref{initial_segment_naturalsplus}.
                Take $C$ such that $C = \closure{V(k)}{X}{T}$.
                $k\mathrel{\restrl{c}{\initialSegment{j}}} C$ by \cref{restrl_iff}.
                Hence $C\in \img{\restrl{c}{\initialSegment{j}}}{\initialSegment{j}}$ by \cref{img_iff}.
                Therefore $x\in C$ by \cref{funs_type_apply,powerset_elim,subseteq_closure,elem_subseteq}.
                Hence $x\in \unions{\img{\restrl{c}{\initialSegment{j}}}{\initialSegment{j}}}$ by \cref{unions_intro}.
                Contradiction.
            \end{byCase}
        \end{subproof}

        Show that there exist $U_4,V_4$ such that
            $U_4$ is open in $X$ with respect to $T$ and
            $V_4$ is open in $X$ with respect to $T$ and
            $A\subseteq U_4$ and
            $B\subseteq V_4$ and
            $U_4$ is disjoint from $V_4$.
        \begin{subproof}
            Trivial.
        \end{subproof}
    \end{subproof}
    Show that $X$ has closed singletons with respect to $T$.
    \begin{subproof}
        Follows by \cref{regular_space}.
    \end{subproof}
    Therefore $X$ is normal with respect to $T$ by \cref{normal_space}.
\end{proof}
% from §32 helper definition 1: metric ball family avoiding a set
\begin{definition}\label{metric_ball_family}
    $\metricBallFamily{d}{X}{A}{B} = \{W\in\pow{X} \mid \exists a\in A. \exists \epsilon\in\reals.
        0 < \epsilon \land W = \metricBall{d}{X}{a}{\epsilon / (1 + 1)} \land
        \metricBall{d}{X}{a}{\epsilon}\inter B = \emptyset\}$.
\end{definition}

\begin{lemma}\label{subseteq_setminus_disjoint}
    Suppose $C\subseteq X\setminus B$.
    Then $C\inter B = \emptyset$.
\end{lemma}
\begin{proof}
    Suppose not.
    Take $x$ such that $x\in C\inter B$ by \cref{nonempty_inhabited}.
    $x\in B$ and $x\notin B$ by \cref{inter,subseteq,setminus_elim_right}.
    Contradiction.
\end{proof}

\begin{lemma}\label{real_double_as_mul_two}
    Suppose $\epsilon\in\reals$.
    Then $\epsilon + \epsilon = \epsilon \rmul (1 + 1)$.
\end{lemma}
\begin{proof}
    $1\in\reals$ by \cref{real_mul_ident}.
    $\epsilon \rmul 1 = 1 \rmul \epsilon$ by \cref{real_mul_comm}.
    $\epsilon \rmul 1 = \epsilon$ by \cref{real_mul_ident}.
    $(\epsilon \rmul 1) + (\epsilon \rmul 1) = \epsilon + \epsilon$ by \cref{real_add_eq_subst}.
    $\epsilon \rmul (1 + 1) = (\epsilon \rmul 1) + (\epsilon \rmul 1)$ by \cref{real_distrib}.
    $\epsilon + \epsilon = \epsilon \rmul (1 + 1)$.
\end{proof}

\begin{lemma}\label{real_half_plus_half}
    Suppose $\epsilon\in\reals$.
    Then $\epsilon / (1 + 1) + \epsilon / (1 + 1) = \epsilon$.
\end{lemma}
\begin{proof}
    $1 + 1\in\reals$ and $0 < 1 + 1$
        by \cref{real_naturals_subset_reals,real_one_in_naturals,subseteq,real_add_ident,real_add_closed,real_zero_lt_one,real_lt_add,real_lt_trans}.
    $1 + 1\neq 0$ and $\inv{(1 + 1)}\in\reals$ by \cref{real_add_closed,real_pos_nonzero,real_mul_inverse}.
    $\epsilon / (1 + 1) + \epsilon / (1 + 1) =
        (\epsilon + \epsilon) \rmul \inv{(1 + 1)}$ by \cref{real_div,real_right_distrib}.
    $\epsilon + \epsilon = \epsilon \rmul (1 + 1)$ by \cref{real_double_as_mul_two}.
    $(\epsilon + \epsilon) \rmul \inv{(1 + 1)} =
        \epsilon \rmul ((1 + 1) \rmul \inv{(1 + 1)})$ by \cref{real_mul_assoc}.
    $(\epsilon + \epsilon) \rmul \inv{(1 + 1)} = \epsilon$ by \cref{real_mul_inverse,real_mul_ident,real_mul_comm}.
    $\epsilon / (1 + 1) + \epsilon / (1 + 1) = \epsilon$.
\end{proof}

\begin{lemma}\label{metric_ball_family_open}
    Assume $d$ is a metric on $X$ and $T = \metricTopology{d}{X}$.
    Assume $1 + 1\in\reals$ and $0 < 1 + 1$ and $\inv{(1 + 1)}\in\reals$.
    Suppose $A$ is closed in $X$ with respect to $T$.
    Suppose $W\in\metricBallFamily{d}{X}{A}{B}$.
    Then $W\in T$.
\end{lemma}
\begin{proof}
    Take $a,\epsilon$ such that
        $a\in A$ and $\epsilon\in\reals$ and $0 < \epsilon$ and
        $W = \metricBall{d}{X}{a}{\epsilon / (1 + 1)}$ and
        $\metricBall{d}{X}{a}{\epsilon}\inter B = \emptyset$
        by \cref{metric_ball_family}.
    $W\in\pow{X}$ and $a\in X$ by \cref{metric_ball_family,closed_in,subseteq}.
    $\epsilon / (1 + 1)\in\reals$ by \cref{real_div,real_mul_closed}.
    $0 < \epsilon / (1 + 1)$ by \cref{real_add_ident,real_lt_mul_pos,real_inv_pos,real_mul_zero,real_lt_subst_left,real_div}.
    $W\in\metricBasis{d}{X}$ by \cref{metric_basis}.
    $W\in T$ by \cref{metric_on,metric_basis_is_basis,basis_elem_open_in_generated,metric_topology,basis_for_topology}.
\end{proof}

\begin{lemma}\label{metric_ball_family_center_in_union}
    Assume $d$ is a metric on $X$.
    Assume $1 + 1\in\reals$ and $0 < 1 + 1$ and $\inv{(1 + 1)}\in\reals$.
    Suppose $a\in A$ and $a\in X$.
    Suppose $\epsilon\in\reals$ and $0 < \epsilon$ and $\metricBall{d}{X}{a}{\epsilon}\subseteq X\setminus B$.
    Then $a\in\unions{\metricBallFamily{d}{X}{A}{B}}$.
\end{lemma}
\begin{proof}
    $d((a,a)) = 0$ by \cref{metric_on,metric_zero_characterization}.
    $0 < \epsilon / (1 + 1)$ by \cref{real_add_ident,real_lt_mul_pos,real_inv_pos,real_mul_zero,real_lt_subst_left,real_div}.
    $a\in \metricBall{d}{X}{a}{\epsilon / (1 + 1)}$ by \cref{metric_ball,real_lt_subst_left}.
    Let $W = \metricBall{d}{X}{a}{\epsilon / (1 + 1)}$.
    $W\in \metricBallFamily{d}{X}{A}{B}$ by \cref{metric_ball,subseteq,powerset_intro,subseteq_setminus_disjoint,metric_ball_family}.
    $a\in\unions{\metricBallFamily{d}{X}{A}{B}}$ by \cref{unions_intro}.
\end{proof}

% from §32 Item 2: metrizable space is normal
\begin{theorem}\label{metrizable_normal}
    Suppose $X$ is metrizable with respect to $T$.
    Then $X$ is normal with respect to $T$.
\end{theorem}
\begin{proof}
    Take $d$ such that $d$ is a metric on $X$ and $T = \metricTopology{d}{X}$ by \cref{metrizable}.
    $1 + 1\in\reals$ and $0 < 1 + 1$
        by \cref{real_naturals_subset_reals,real_one_in_naturals,subseteq,real_add_ident,real_add_closed,real_zero_lt_one,real_lt_add,real_lt_trans}.
    $1 + 1\neq 0$ and $\inv{(1 + 1)}\in\reals$ by \cref{real_add_closed,real_pos_nonzero,real_mul_inverse}.
    Show that for all $A,B$ such that $A$ is closed in $X$ with respect to $T$ and $B$ is closed in $X$ with respect to $T$
        and $A$ is disjoint from $B$
        there exist $U,V$ such that
            $U$ is open in $X$ with respect to $T$ and
            $V$ is open in $X$ with respect to $T$ and
            $A\subseteq U$ and
            $B\subseteq V$ and
            $U$ is disjoint from $V$.
    \begin{subproof}
        Fix $A$.
        Fix $B$.
        Assume $A$ is closed in $X$ with respect to $T$ and $B$ is closed in $X$ with respect to $T$
            and $A$ is disjoint from $B$.
        Let $F = \metricBallFamily{d}{X}{A}{B}$. Let $G = \metricBallFamily{d}{X}{B}{A}$.
        Let $U = \unions{F}$. Let $V = \unions{G}$.
        $T$ is a topology on $X$ by \cref{metrizable}.

        Show that for all $W$ such that $W\in F$ we have $W\in T$.
        \begin{subproof}
            Follows by \cref{metric_ball_family_open}.
        \end{subproof}
        $U$ is open in $X$ with respect to $T$ by \cref{topology_on,subseteq,open_in}.

        Show that for all $W$ such that $W\in G$ we have $W\in T$.
        \begin{subproof}
            Follows by \cref{metric_ball_family_open}.
        \end{subproof}
        $V$ is open in $X$ with respect to $T$ by \cref{topology_on,subseteq,open_in}.

        Show that for all $a$ such that $a\in A$ there exists $\epsilon\in\reals$ such that
            $0 < \epsilon$ and $\metricBall{d}{X}{a}{\epsilon}\subseteq X\setminus B$.
        \begin{subproof}
            Follows by \cref{closed_in,subseteq,setminus_intro,disjoint,setminus_subseteq,metric_open_iff}.
        \end{subproof}
        Show that for all $a$ such that $a\in A$ we have $a\in U$.
        \begin{subproof}
            Follows by \cref{closed_in,subseteq,metric_ball_family_center_in_union}.
        \end{subproof}
        $A\subseteq U$ by \cref{subseteq}.

        Show that for all $b$ such that $b\in B$ there exists $\epsilon\in\reals$ such that
            $0 < \epsilon$ and $\metricBall{d}{X}{b}{\epsilon}\subseteq X\setminus A$.
        \begin{subproof}
            Follows by \cref{closed_in,subseteq,setminus_intro,disjoint,setminus_subseteq,metric_open_iff}.
        \end{subproof}
        Show that for all $b$ such that $b\in B$ we have $b\in V$.
        \begin{subproof}
            Follows by \cref{closed_in,subseteq,metric_ball_family_center_in_union}.
        \end{subproof}
        $B\subseteq V$ by \cref{subseteq}.

        Show that $U$ is disjoint from $V$.
        \begin{subproof}
            Suppose not.
            Take $z$ such that $z\in U$ and $z\in V$ by \cref{disjoint}.
            Take $W_1$ such that $W_1\in F$ and $z\in W_1$ by \cref{unions_iff}.
            Take $W_2$ such that $W_2\in G$ and $z\in W_2$ by \cref{unions_iff}.
            Take $a,\epsilon_a$ such that
                $a\in A$ and $\epsilon_a\in\reals$ and $0 < \epsilon_a$ and
                $W_1 = \metricBall{d}{X}{a}{\epsilon_a / (1 + 1)}$ and
                $\metricBall{d}{X}{a}{\epsilon_a}\inter B = \emptyset$
                by \cref{metric_ball_family}.
            Take $b,\epsilon_b$ such that
                $b\in B$ and $\epsilon_b\in\reals$ and $0 < \epsilon_b$ and
                $W_2 = \metricBall{d}{X}{b}{\epsilon_b / (1 + 1)}$ and
                $\metricBall{d}{X}{b}{\epsilon_b}\inter A = \emptyset$
                by \cref{metric_ball_family}.
            $d((a,z)) < \epsilon_a / (1 + 1)$ and $d((b,z)) < \epsilon_b / (1 + 1)$ by \cref{metric_ball}.
            $d$ satisfies the triangle inequality on $X$ and $d$ is symmetric on $X$ and $z\in X$
                by \cref{metric_on,metric_symmetric,metric_ball}.
            $a\in X$ and $b\in X$ by \cref{closed_in,subseteq}.
            $d((z,b)) = d((b,z))$ by \cref{metric_symmetric}.
            $(a,z)\in X\times X$ and $(z,b)\in X\times X$ by \cref{times_tuple_intro}.
            $d((a,z))\in\reals$ and $d((z,b))\in\reals$ by \cref{funs_type_apply,metric_on}.
            $\epsilon_a / (1 + 1)\in\reals$ and $\epsilon_b / (1 + 1)\in\reals$ by \cref{real_div,real_mul_closed}.
            $d((a,b))\in\reals$ by \cref{times_tuple_intro,funs_type_apply,metric_on}.
            $d((a,z)) + d((z,b))\in\reals$ by \cref{real_add_closed}.
            $\epsilon_a / (1 + 1) + \epsilon_b / (1 + 1)\in\reals$ by \cref{real_add_closed}.
            $d((a,b)) < d((a,z)) + d((z,b))$ or $d((a,b)) = d((a,z)) + d((z,b))$
                by \cref{metric_triangle_inequality}.
            $d((a,z)) + d((z,b)) < \epsilon_a / (1 + 1) + \epsilon_b / (1 + 1)$ by \cref{real_lt_subst_left,real_lt_add_two}.
            $d((a,b)) < \epsilon_a / (1 + 1) + \epsilon_b / (1 + 1)$
                by \cref{real_lt_trans,real_lt_subst_left}.
            \begin{byCase}
            \caseOf{$\epsilon_a < \epsilon_b$.}
                $\epsilon_a / (1 + 1)\in\reals$ and $\epsilon_b / (1 + 1)\in\reals$ by \cref{real_div,real_mul_closed}.
                $\epsilon_a / (1 + 1) < \epsilon_b / (1 + 1)$ by \cref{real_lt_mul_pos,real_inv_pos,real_div}.
                $\epsilon_a / (1 + 1) + \epsilon_b / (1 + 1) < \epsilon_b / (1 + 1) + \epsilon_b / (1 + 1)$ by \cref{real_lt_add}.
                $\epsilon_b / (1 + 1) + \epsilon_b / (1 + 1) = \epsilon_b$ by \cref{real_half_plus_half}.
                $d((a,b)) < \epsilon_b / (1 + 1) + \epsilon_b / (1 + 1)$ by \cref{real_lt_trans}.
                $d((a,b)) < \epsilon_b$ by \cref{real_lt_subst_right}.
                $a\in\metricBall{d}{X}{b}{\epsilon_b}$
                    by \cref{metric_symmetric,metric_on,real_lt_subst_left,metric_ball}.
                Contradiction by \cref{inter_intro,subseteq_emptyset_iff,subseteq,emptyset}.
            \caseOf{$\epsilon_a = \epsilon_b$.}
                $\epsilon_b = \epsilon_a$.
                $\epsilon_b / (1 + 1) = \epsilon_a / (1 + 1)$.
                $\epsilon_a / (1 + 1) + \epsilon_b / (1 + 1) = \epsilon_a / (1 + 1) + \epsilon_a / (1 + 1)$ by \cref{real_add_eq_subst}.
                $\epsilon_a / (1 + 1) + \epsilon_a / (1 + 1) = \epsilon_a$ by \cref{real_half_plus_half}.
                $d((a,b)) < \epsilon_a / (1 + 1) + \epsilon_a / (1 + 1)$ by \cref{real_lt_subst_right}.
                $d((a,b)) < \epsilon_a$ by \cref{real_lt_subst_right}.
                $b\in\metricBall{d}{X}{a}{\epsilon_a}$ by \cref{metric_ball}.
                Contradiction by \cref{inter_intro,subseteq_emptyset_iff,subseteq,emptyset}.
            \caseOf{$\epsilon_b < \epsilon_a$.}
                $\epsilon_a / (1 + 1)\in\reals$ and $\epsilon_b / (1 + 1)\in\reals$ by \cref{real_div,real_mul_closed}.
                $\epsilon_b / (1 + 1) < \epsilon_a / (1 + 1)$ by \cref{real_lt_mul_pos,real_inv_pos,real_div}.
                $\epsilon_b / (1 + 1) + \epsilon_a / (1 + 1) < \epsilon_a / (1 + 1) + \epsilon_a / (1 + 1)$ by \cref{real_lt_add}.
                $\epsilon_a / (1 + 1) + \epsilon_b / (1 + 1) < \epsilon_a / (1 + 1) + \epsilon_a / (1 + 1)$ by \cref{real_add_comm,real_lt_subst_left}.
                $\epsilon_a / (1 + 1) + \epsilon_a / (1 + 1) = \epsilon_a$ by \cref{real_half_plus_half}.
                $d((a,b)) < \epsilon_a / (1 + 1) + \epsilon_a / (1 + 1)$ by \cref{real_lt_trans}.
                $d((a,b)) < \epsilon_a$ by \cref{real_lt_subst_right}.
                $b\in\metricBall{d}{X}{a}{\epsilon_a}$ by \cref{metric_ball}.
                Contradiction by \cref{inter_intro,subseteq_emptyset_iff,subseteq,emptyset}.
            \end{byCase}
        \end{subproof}
    \end{subproof}
    $T$ is a topology on $X$ and $X$ is Hausdorff with respect to $T$ by \cref{metrizable,metric_space_hausdorff}.
    Show that $X$ has closed singletons with respect to $T$.
    \begin{subproof}
        Follows by \cref{singleton_subset_intro,pair_is_finite,singleton_iff,upair_intro_left,subseteq,subseteq_finite_is_finite,finite_point_set_closed_hausdorff,one_point_closed}.
    \end{subproof}
    $X$ is normal with respect to $T$ by \cref{normal_space}.
\end{proof}

% from §32 Item 3: compact Hausdorff space is normal
\begin{theorem}\label{compact_hausdorff_normal}
    Suppose $X$ is compact with respect to $T$ and $X$ is Hausdorff with respect to $T$.
    Then $X$ is normal with respect to $T$.
\end{theorem}
\begin{proof}
    $T$ is a topology on $X$ by \cref{hausdorff}.
    Show that for all $A,B$ such that $A$ is closed in $X$ with respect to $T$ and $B$ is closed in $X$ with respect to $T$
        and $A$ is disjoint from $B$
        there exist $U,V$ such that
            $U$ is open in $X$ with respect to $T$ and
            $V$ is open in $X$ with respect to $T$ and
            $A\subseteq U$ and
            $B\subseteq V$ and
            $U$ is disjoint from $V$.
    \begin{subproof}
        Fix $A$.
        Fix $B$.
        Assume $A$ is closed in $X$ with respect to $T$ and $B$ is closed in $X$ with respect to $T$
            and $A$ is disjoint from $B$.
        \begin{byCase}
        \caseOf{$A = \emptyset$.}
            Show that there exist $U,V$ such that
                $U$ is open in $X$ with respect to $T$ and
                $V$ is open in $X$ with respect to $T$ and
                $A\subseteq U$ and
                $B\subseteq V$ and
                $U$ is disjoint from $V$.
            \begin{subproof}
                Follows by \cref{topology_on,open_in,emptyset_subseteq,closed_in,subseteq,disjoint,emptyset}.
            \end{subproof}
        \caseOf{$A\neq\emptyset$.}
            Show that there exist $U,V$ such that
                $U$ is open in $X$ with respect to $T$ and
                $V$ is open in $X$ with respect to $T$ and
                $A\subseteq U$ and
                $B\subseteq V$ and
                $U$ is disjoint from $V$.
            \begin{subproof}
                $A\subseteq X$ and $B\subseteq X$ by \cref{closed_in}.
                Let $T_A = \subspaceTopology{T}{A}$. Let $T_B = \subspaceTopology{T}{B}$.
                $A$ is compact with respect to $T_A$ and $B$ is compact with respect to $T_B$ by \cref{closed_subspace_compact}.
                Let $F = \{U\in T \mid \text{there exists $a\in A$ such that there exists $V$ such that
                    $V$ is open in $X$ with respect to $T$ and
                    $a\in U$ and
                    $B\subseteq V$ and
                    $U$ is disjoint from $V$}\}$.
                Show that for all $a$ such that $a\in A$ we have $a\in\unions{F}$.
                \begin{subproof}
                    Fix $a$ such that $a\in A$.
                    $a\in X$ and $a\notin B$ by \cref{subseteq,disjoint}.
                    $a\in X\setminus B$ by \cref{setminus_intro}.
                    Take $U,V$ such that
                    $U$ is open in $X$ with respect to $T$ and
                    $V$ is open in $X$ with respect to $T$ and
                    $a\in U$ and
                        $B\subseteq V$ and
                        $U$ is disjoint from $V$
                        by \cref{compact_hausdorff_separation}.
                    $U\in T$ by \cref{open_in}.
                    $U\in F$ by definition.
                    $a\in\unions{F}$ by \cref{unions_intro}.
                \end{subproof}
                $F$ is a covering of $A$ by \cref{subseteq,covering}.
                Take $G$ such that $G\subseteq F$ and $G$ is finite and $G$ is a covering of $A$
                    by \cref{open_in,subspace_of,compact_subspace_open_covering,subseteq}.
                $G$ is inhabited by \cref{nonempty_inhabited,covering,subseteq,unions_iff}.
                Let $R = \{w\in G\times T \mid \text{there exists $U_1\in G$ such that there exists $V_1\in T$ such that
                    $w = (U_1,V_1)$ and
                    $B\subseteq V_1$ and
                    $U_1$ is disjoint from $V_1$}\}$.
                $R$ is a relation by \cref{subseteq,times_elem_is_tuple,relation}.
                Show that for all $U\in G$ there exists $V$ such that $U\mathrel{R} V$.
                \begin{subproof}
                    Fix $U$ such that $U\in G$.
                    Take $a,V$ such that $a\in A$ and
                        $V$ is open in $X$ with respect to $T$ and
                        $a\in U$ and
                        $B\subseteq V$ and
                        $U$ is disjoint from $V$
                        by \cref{subseteq}.
                    $V\in T$ by \cref{open_in}.
                    $(U,V)\in G\times T$ by \cref{times_tuple_intro}.
                    $(U,V)\in R$ by definition.
                    $U\mathrel{R} V$.
                \end{subproof}
                Take $g$ such that $g$ is a function on $G$ and for all $U\in G$ we have $U\mathrel{R}\apply{g}{U}$
                    by \cref{choice_function}.
                Let $H = \img{g}{G}$.
                $H\subseteq T$ by \cref{img_iff,function_apply_iff,pair_eq_iff,subseteq}.
                $H$ is finite by \cref{finite_image_function}.
                $H$ is inhabited by \cref{nonempty_inhabited,inhabited_nonempty,function_apply_elim,subseteq,img_iff}.
                Let $U_0 = \unions{G}$.
                Let $V_0 = \inters{H}$.
                $U_0$ is open in $X$ with respect to $T$ by \cref{topology_on,subseteq,open_in}.
                $V_0$ is open in $X$ with respect to $T$ by \cref{subseteq,finite_intersections_open_in_topology,open_in}.
                $A\subseteq U_0$ by \cref{covering}.
                $B\subseteq V_0$ by \cref{img_iff,function_apply_iff,pair_eq_iff,inters_greatest}.
                Show that for all $z$ such that $z\in U_0$ we have $z\notin V_0$.
                \begin{subproof}
                    Fix $z$ such that $z\in U_0$.
                    Take $W\in G$ such that $z\in W$ by \cref{unions_iff}.
                    $\apply{g}{W}\in H$ by \cref{function_apply_elim,subseteq,img_iff}.
                    $V_0\subseteq \apply{g}{W}$ by \cref{inters_subseteq_elem}.
                    Take $U_1,V_1$ such that $U_1\in G$ and $V_1\in T$ and
                        $(W,\apply{g}{W}) = (U_1,V_1)$ and
                        $B\subseteq V_1$ and
                        $U_1$ is disjoint from $V_1$
                        by \cref{subseteq,function_apply_iff}.
                    $W$ is disjoint from $\apply{g}{W}$ by \cref{pair_eq_iff}.
                    $W$ is disjoint from $V_0$ by \cref{subseteq,disjoint}.
                    $z\notin V_0$ by \cref{disjoint}.
                \end{subproof}
                $U_0$ is disjoint from $V_0$ by \cref{disjoint}.
                Show that there exist $W,Z$ such that
                    $W$ is open in $X$ with respect to $T$ and
                    $Z$ is open in $X$ with respect to $T$ and
                    $A\subseteq W$ and
                    $B\subseteq Z$ and
                    $W$ is disjoint from $Z$.
                \begin{subproof}
                    Trivial.
                \end{subproof}
            \end{subproof}
        \end{byCase}
    \end{subproof}
    Show that $X$ has closed singletons with respect to $T$.
    \begin{subproof}
        Follows by \cref{singleton_subset_intro,pair_is_finite,singleton_iff,upair_intro_left,subseteq,subseteq_finite_is_finite,finite_point_set_closed_hausdorff,one_point_closed}.
    \end{subproof}
    $X$ is normal with respect to $T$ by \cref{normal_space}.
\end{proof}

% from §32 helper definition 2: well order relation
\begin{definition}\label{well_order_relation}
    $R$ is a well order relation on $X$ iff
        $R$ is a simple order relation on $X$ and
        for all $A$ such that $A\subseteq X$ and $A\neq\emptyset$
            there exists $a$ such that $a$ is a smallest element of $A$ with respect to $R$.
\end{definition}

% from §32 helper lemma 1: non-largest element has an immediate successor
\begin{lemma}\label{well_order_non_largest_has_immediate_successor}
    Assume $R$ is a well order relation on $X$.
    Suppose $x\in X$.
    Suppose $x$ is not a largest element of $X$ with respect to $R$.
    Then there exists $y\in X$ such that $y$ is an immediate successor of $x$ in $X$ with respect to $R$.
\end{lemma}
\begin{proof}
    $R$ is a simple order relation on $X$ by \cref{well_order_relation}.
    Take $z\in X$ such that $x\mathrel{R}z$ by \cref{non_largest_has_successor}.
    Let $A = \openrayright{R}{X}{x}$.
    $A\subseteq X$ by \cref{openray_right,subseteq}.
    $z\in A$ by \cref{openray_right}.
    $A\neq\emptyset$ by \cref{emptyset}.
    Take $y$ such that $y$ is a smallest element of $A$ with respect to $R$ by \cref{well_order_relation}.
    $y\in X$ and $x\mathrel{R}y$ by \cref{smallest_element,openray_right,subseteq}.
    Show that $\intervalclosedrel{R}{X}{x}{y} = \{x,y\}$.
    \begin{subproof}
        $\{x,y\}\subseteq\intervalclosedrel{R}{X}{x}{y}$ by \cref{upair_elim,intervalclosed_rel,subseteq}.
        Show that for all $u$ such that $u\in\intervalclosedrel{R}{X}{x}{y}$ we have $u\in\{x,y\}$.
        \begin{subproof}
            Fix $u$ such that $u\in\intervalclosedrel{R}{X}{x}{y}$.
            \begin{byCase}
            \caseOf{$u = x$.}
                $u\in\{x,y\}$ by \cref{upair_intro_left}.
            \caseOf{$u\neq x$.}
                $u\in X$ and $x\mathrel{R}u$ by \cref{intervalclosed_rel}.
                $u\in A$ by \cref{openray_right}.
                $y\mathrel{R}u$ or $y = u$ by \cref{smallest_element}.
                \begin{byCase}
                \caseOf{$y = u$.}
                    $u\in\{x,y\}$ by \cref{upair_intro_right}.
                \caseOf{$y\mathrel{R}u$.}
                    Suppose not.
                    $u\mathrel{R}y$ by \cref{intervalclosed_rel}.
                    Contradiction by \cref{simple_order_relation,asymmetric}.
                    $u = y$.
                    $u\in\{x,y\}$ by \cref{upair_intro_right}.
                \end{byCase}
            \end{byCase}
        \end{subproof}
        $\intervalclosedrel{R}{X}{x}{y}\subseteq\{x,y\}$ by \cref{subseteq}.
        $\intervalclosedrel{R}{X}{x}{y} = \{x,y\}$ by \cref{subseteq_antisymmetric}.
    \end{subproof}
    $y$ is an immediate successor of $x$ in $X$ with respect to $R$ by \cref{immediate_successor}.
\end{proof}

% from §32 helper lemma 2: left half-open intervals are open in a well order
\begin{lemma}\label{well_order_intervalopenclosed_open}
    Assume $R$ is a well order relation on $X$.
    Assume $T$ is the order topology on $X$ with respect to $R$.
    Suppose $x,y\in X$.
    Suppose $x\mathrel{R}y$.
    Then $\intervalopenclosedrel{R}{X}{x}{y}$ is open in $X$ with respect to $T$.
\end{lemma}
\begin{proof}
    $R$ is a simple order relation on $X$ and $X$ has more than one element by \cref{well_order_relation,order_topology}.
    \begin{byCase}
    \caseOf{$y$ is a largest element of $X$ with respect to $R$.}
        $\intervalopenclosedrel{R}{X}{x}{y} = \openrayright{R}{X}{x}$ by \cref{intervalopenclosedrel_largest_eq_openrayright}.
        $\intervalopenclosedrel{R}{X}{x}{y}$ is open in $X$ with respect to $T$ by \cref{openrayright_open_in_order_topology}.
    \caseOf{$y$ is not a largest element of $X$ with respect to $R$.}
        Take $s\in X$ such that $s$ is an immediate successor of $y$ in $X$ with respect to $R$
            by \cref{well_order_non_largest_has_immediate_successor}.
        Let $U = \openrayright{R}{X}{x}\inter \openrayleft{R}{X}{s}$.
        Show that $\intervalopenclosedrel{R}{X}{x}{y}\subseteq U$.
        \begin{subproof}
            Show that for all $z$ such that $z\in\intervalopenclosedrel{R}{X}{x}{y}$ we have $z\in U$.
            \begin{subproof}
                Fix $z$ such that $z\in\intervalopenclosedrel{R}{X}{x}{y}$.
                $z\in X$ and $x\mathrel{R}z$ and $(z\mathrel{R}y \lor z = y)$ by \cref{intervalopenclosed_rel}.
                $y\mathrel{R}s$ by \cref{immediate_successor}.
                Show that $z\mathrel{R}s$.
                \begin{subproof}
                    \begin{byCase}
                    \caseOf{$z\mathrel{R}y$.}
                        $z\mathrel{R}s$ by \cref{simple_order_relation,transitive}.
                    \caseOf{$z = y$.}
                        $z\mathrel{R}s$ by \cref{immediate_successor}.
                    \end{byCase}
                \end{subproof}
                $z\in U$ by \cref{openray_right,openray_left,inter_intro}.
            \end{subproof}
            $\intervalopenclosedrel{R}{X}{x}{y}\subseteq U$ by \cref{subseteq}.
        \end{subproof}
        Show that $U\subseteq\intervalopenclosedrel{R}{X}{x}{y}$.
        \begin{subproof}
            Show that for all $z$ such that $z\in U$ we have $z\in\intervalopenclosedrel{R}{X}{x}{y}$.
            \begin{subproof}
                Fix $z$ such that $z\in U$.
                $x\mathrel{R}z$ and $z\mathrel{R}s$ and $z\in X$ by \cref{inter,openray_right,openray_left}.
                \begin{byCase}
                \caseOf{$z = y$.}
                    $z\in\intervalopenclosedrel{R}{X}{x}{y}$ by \cref{intervalopenclosed_rel}.
                \caseOf{$z \neq y$.}
                    $z\mathrel{R}y$ or $y\mathrel{R}z$ by \cref{simple_order_relation,connex}.
                    \begin{byCase}
                    \caseOf{$z\mathrel{R}y$.}
                        $z\in\intervalopenclosedrel{R}{X}{x}{y}$ by \cref{intervalopenclosed_rel}.
                    \caseOf{$y\mathrel{R}z$.}
                        Show that $z\in\intervalopenclosedrel{R}{X}{x}{y}$.
                        \begin{subproof}
                            Suppose not.
                            $z\in\intervalclosedrel{R}{X}{y}{s}$ by \cref{intervalclosed_rel}.
                            $z\in\{y,s\}$ by \cref{simple_order_relation,asymmetric,immediate_successor_interval}.
                            $z = y$ or $z = s$ by \cref{upair_elim}.
                            $z = s$ by \cref{simple_order_relation,asymmetric}.
                            $s\mathrel{R}s$ by \cref{openray_left}.
                            Contradiction by \cref{simple_order_relation,asymmetric}.
                        \end{subproof}
                    \end{byCase}
                \end{byCase}
            \end{subproof}
            $U\subseteq\intervalopenclosedrel{R}{X}{x}{y}$ by \cref{subseteq}.
        \end{subproof}
        $U = \intervalopenclosedrel{R}{X}{x}{y}$ by \cref{subseteq_antisymmetric}.
        $\openrayright{R}{X}{x}$ is open in $X$ with respect to $T$ and $\openrayleft{R}{X}{s}$ is open in $X$ with respect to $T$
            by \cref{openrayright_open_in_order_topology,openrayleft_open_in_order_topology}.
        $U$ is open in $X$ with respect to $T$ by \cref{order_basis_is_basis,basis_topology_is_topology,order_topology,open_in,topology_on}.
        $\intervalopenclosedrel{R}{X}{x}{y}$ is open in $X$ with respect to $T$.
    \end{byCase}
\end{proof}

% from §32 helper lemma 3: neighborhoods contain left half-open intervals
\begin{lemma}\label{well_order_left_interval_neighborhood}
    Assume $R$ is a well order relation on $X$.
    Assume $T$ is the order topology on $X$ with respect to $R$.
    Suppose $a\in X$.
    Suppose $a$ is not a smallest element of $X$ with respect to $R$.
    Suppose $U$ is open in $X$ with respect to $T$.
    Suppose $a\in U$.
    Then there exists $x\in X$ such that $x\mathrel{R}a$ and $\intervalopenclosedrel{R}{X}{x}{a}\subseteq U$.
\end{lemma}
\begin{proof}
    $R$ is a simple order relation on $X$ by \cref{well_order_relation}.
    Let $B = \orderBasis{R}{X}$.
    $T = \basisTopology{B}{X}$ by \cref{order_topology}.
    $B$ is a basis for $T$ on $X$ by \cref{order_basis_is_basis,basis_for_topology,order_topology}.
    $U\in\basisTopology{B}{X}$ by \cref{open_in,basis_for_topology}.
    Take $B_1\in B$ such that $a\in B_1$ and $B_1\subseteq U$ by \cref{topology_generated_by_basis}.
    $((\exists u, v\in X. u\mathrel{R}v \land B_1 = \intervalopenrel{R}{X}{u}{v}) \lor
        (\exists v\in X. \exists a_0\in X. (\forall z\in X. a_0\mathrel{R}z \lor a_0 = z) \land B_1 = \intervalclosedopenrel{R}{X}{a_0}{v}) \lor
        (\exists u\in X. \exists b_0\in X. (\forall z\in X. z\mathrel{R}b_0 \lor z = b_0) \land B_1 = \intervalopenclosedrel{R}{X}{u}{b_0}))$
        by \cref{order_basis}.
    \begin{byCase}
    \caseOf{$\exists u, v\in X. u\mathrel{R}v \land B_1 = \intervalopenrel{R}{X}{u}{v}$.}
        Take $u, v\in X$ such that $u\mathrel{R}v$ and $B_1 = \intervalopenrel{R}{X}{u}{v}$.
        $u\mathrel{R}a$ and $a\mathrel{R}v$ by \cref{intervalopen_rel}.
        Show that $\intervalopenclosedrel{R}{X}{u}{a}\subseteq B_1$.
        \begin{subproof}
            Show that for all $z$ such that $z\in\intervalopenclosedrel{R}{X}{u}{a}$ we have $z\in B_1$.
            \begin{subproof}
                Fix $z$ such that $z\in\intervalopenclosedrel{R}{X}{u}{a}$.
                $z\in X$ and $u\mathrel{R}z$ and $(z\mathrel{R}a \lor z = a)$ by \cref{intervalopenclosed_rel}.
                \begin{byCase}
                \caseOf{$z\mathrel{R}a$.}
                    $z\mathrel{R}v$ by \cref{simple_order_relation,transitive}.
                    $z\in B_1$ by \cref{intervalopen_rel}.
                \caseOf{$z = a$.}
                    $z\mathrel{R}v$.
                    $z\in B_1$ by \cref{intervalopen_rel}.
                \end{byCase}
            \end{subproof}
            $\intervalopenclosedrel{R}{X}{u}{a}\subseteq B_1$ by \cref{subseteq}.
        \end{subproof}
        Show that there exists $x\in X$ such that $x\mathrel{R}a$ and $\intervalopenclosedrel{R}{X}{x}{a}\subseteq U$.
        \begin{subproof}
            Follows by \cref{subseteq_transitive}.
        \end{subproof}
    \caseOf{$\exists v\in X. \exists a_0\in X. (\forall z\in X. a_0\mathrel{R}z \lor a_0 = z) \land B_1 = \intervalclosedopenrel{R}{X}{a_0}{v}$.}
        Take $a_0, v\in X$ such that
            $(\forall z\in X. a_0\mathrel{R}z \lor a_0 = z)$ and
            $B_1 = \intervalclosedopenrel{R}{X}{a_0}{v}$.
        $a\mathrel{R}v$ by \cref{intervalclosedopen_rel}.
        Take $x\in X$ such that $x\mathrel{R}a$ by \cref{non_smallest_has_predecessor}.
        Show that $\intervalopenclosedrel{R}{X}{x}{a}\subseteq B_1$.
        \begin{subproof}
            Show that for all $z$ such that $z\in\intervalopenclosedrel{R}{X}{x}{a}$ we have $z\in B_1$.
            \begin{subproof}
                Fix $z$ such that $z\in\intervalopenclosedrel{R}{X}{x}{a}$.
                $z\in X$ and $x\mathrel{R}z$ and $(z\mathrel{R}a \lor z = a)$ by \cref{intervalopenclosed_rel}.
                $a_0\mathrel{R}z$ or $z = a_0$ by \cref{smallest_element}.
                \begin{byCase}
                \caseOf{$z\mathrel{R}a$.}
                    $z\mathrel{R}v$ by \cref{simple_order_relation,transitive}.
                    $z\in B_1$ by \cref{intervalclosedopen_rel}.
                \caseOf{$z = a$.}
                    $z\mathrel{R}v$.
                    $z\in B_1$ by \cref{intervalclosedopen_rel}.
                \end{byCase}
            \end{subproof}
            $\intervalopenclosedrel{R}{X}{x}{a}\subseteq B_1$ by \cref{subseteq}.
        \end{subproof}
        $\intervalopenclosedrel{R}{X}{x}{a}\subseteq U$ by \cref{subseteq_transitive}.
        Show that there exists $y\in X$ such that $y\mathrel{R}a$ and $\intervalopenclosedrel{R}{X}{y}{a}\subseteq U$.
        \begin{subproof}
            Trivial.
        \end{subproof}
    \caseOf{$\exists u\in X. \exists b_0\in X. (\forall z\in X. z\mathrel{R}b_0 \lor z = b_0) \land B_1 = \intervalopenclosedrel{R}{X}{u}{b_0}$.}
        Take $u, b_0\in X$ such that
            $(\forall z\in X. z\mathrel{R}b_0 \lor z = b_0)$ and
            $B_1 = \intervalopenclosedrel{R}{X}{u}{b_0}$.
        $u\mathrel{R}a$ and $(a\mathrel{R}b_0 \lor a = b_0)$ by \cref{intervalopenclosed_rel}.
        Show that $\intervalopenclosedrel{R}{X}{u}{a}\subseteq B_1$.
        \begin{subproof}
            Show that for all $z$ such that $z\in\intervalopenclosedrel{R}{X}{u}{a}$ we have $z\in B_1$.
            \begin{subproof}
                Fix $z$ such that $z\in\intervalopenclosedrel{R}{X}{u}{a}$.
                $z\in X$ and $u\mathrel{R}z$ and $(z\mathrel{R}a \lor z = a)$ by \cref{intervalopenclosed_rel}.
                \begin{byCase}
                \caseOf{$z\mathrel{R}a$.}
                    $z\mathrel{R}b_0$ or $z = b_0$ by \cref{simple_order_relation,transitive,intervalopenclosed_rel}.
                    $z\in B_1$ by \cref{intervalopenclosed_rel}.
                \caseOf{$z = a$.}
                    $z\mathrel{R}b_0$ or $z = b_0$ by \cref{intervalopenclosed_rel}.
                    $z\in B_1$ by \cref{intervalopenclosed_rel}.
                \end{byCase}
            \end{subproof}
            $\intervalopenclosedrel{R}{X}{u}{a}\subseteq B_1$ by \cref{subseteq}.
        \end{subproof}
        Show that there exists $x\in X$ such that $x\mathrel{R}a$ and $\intervalopenclosedrel{R}{X}{x}{a}\subseteq U$.
        \begin{subproof}
            Follows by \cref{subseteq_transitive}.
        \end{subproof}
    \end{byCase}
\end{proof}

% from §32 helper lemma 4: the smallest singleton is open
\begin{lemma}\label{well_order_smallest_singleton_open}
    Assume $R$ is a well order relation on $X$.
    Assume $T$ is the order topology on $X$ with respect to $R$.
    Suppose $a_0$ is a smallest element of $X$ with respect to $R$.
    Then $\{a_0\}$ is open in $X$ with respect to $T$.
\end{lemma}
\begin{proof}
    $R$ is a simple order relation on $X$ and $X$ has more than one element by \cref{well_order_relation,order_topology}.
    $a_0\in X$ by \cref{smallest_element}.
    Take $y\in X$ such that $y\neq a_0$ by \cref{more_than_one_element_has_other}.
    $a_0\mathrel{R}y$ by \cref{smallest_element}.
    Show that $a_0$ is not a largest element of $X$ with respect to $R$.
    \begin{subproof}
        Suppose not.
        $y\mathrel{R}a_0$ or $y = a_0$ by \cref{largest_element}.
        \begin{byCase}
        \caseOf{$y\mathrel{R}a_0$.}
            Contradiction by \cref{simple_order_relation,asymmetric}.
        \caseOf{$y = a_0$.}
            Contradiction.
        \end{byCase}
    \end{subproof}
    Take $s\in X$ such that $s$ is an immediate successor of $a_0$ in $X$ with respect to $R$
        by \cref{well_order_non_largest_has_immediate_successor}.
    Let $U = \openrayleft{R}{X}{s}$.
    Show that $\{a_0\}\subseteq U$.
    \begin{subproof}
        $a_0\mathrel{R}s$ by \cref{immediate_successor}.
        $a_0\in U$ by \cref{openray_left}.
        $\{a_0\}\subseteq U$ by \cref{singleton_subset_intro,subseteq}.
    \end{subproof}
    Show that $U\subseteq\{a_0\}$.
    \begin{subproof}
        Show that for all $z$ such that $z\in U$ we have $z\in\{a_0\}$.
        \begin{subproof}
            Fix $z$ such that $z\in U$.
            $z\in X$ and $z\mathrel{R}s$ by \cref{openray_left}.
            $a_0\mathrel{R}z$ or $z = a_0$ by \cref{smallest_element}.
            \begin{byCase}
            \caseOf{$z = a_0$.}
                $z\in\{a_0\}$ by \cref{singleton_intro}.
            \caseOf{$a_0\mathrel{R}z$.}
                $z\in\intervalclosedrel{R}{X}{a_0}{s}$ by \cref{intervalclosed_rel}.
                $z\in\{a_0,s\}$ by \cref{simple_order_relation,asymmetric,immediate_successor_interval}.
                $z = a_0$ or $z = s$ by \cref{upair_elim}.
                Suppose not.
                $z = s$.
                $s\mathrel{R}s$ by \cref{openray_left}.
                Contradiction by \cref{simple_order_relation,asymmetric}.
                $z = a_0$.
                $z\in\{a_0\}$ by \cref{singleton_intro}.
            \end{byCase}
        \end{subproof}
        $U\subseteq\{a_0\}$ by \cref{subseteq}.
    \end{subproof}
    $U = \{a_0\}$ by \cref{subseteq_antisymmetric}.
    $U$ is open in $X$ with respect to $T$ by \cref{openrayleft_open_in_order_topology}.
    $\{a_0\}$ is open in $X$ with respect to $T$.
\end{proof}

% from §32 Item 4: well-ordered set is normal in the order topology
\begin{theorem}\label{well_ordered_normal}
    Assume $R$ is a well order relation on $X$ and $T$ is the order topology on $X$ with respect to $R$.
    Then $X$ is normal with respect to $T$.
\end{theorem}
\begin{proof}
    $R$ is a simple order relation on $X$ and $X$ has more than one element by \cref{well_order_relation,order_topology}.
    $T$ is a topology on $X$ and $X$ is Hausdorff with respect to $T$ by \cref{order_basis_is_basis,basis_topology_is_topology,order_topology,order_topology_hausdorff}.
    Show that $X$ has closed singletons with respect to $T$.
    \begin{subproof}
        Follows by \cref{singleton_subset_intro,pair_is_finite,singleton_iff,upair_intro_left,subseteq,subseteq_finite_is_finite,finite_point_set_closed_hausdorff,one_point_closed}.
    \end{subproof}
    Take $p\in X$ such that $p\in X$ by \cref{more_than_one_element}.
    $X\neq\emptyset$ by \cref{emptyset}.
    $X\subseteq X$ by \cref{subseteq_refl}.
    Take $a_0$ such that $a_0$ is a smallest element of $X$ with respect to $R$ by \cref{well_order_relation}.
    Show that for all $C,D$ such that $C$ is closed in $X$ with respect to $T$ and $D$ is closed in $X$ with respect to $T$
        and $C$ is disjoint from $D$ and $a_0\notin C$ and $a_0\notin D$
        there exist $U,V$ such that
            $U$ is open in $X$ with respect to $T$ and
            $V$ is open in $X$ with respect to $T$ and
            $C\subseteq U$ and
            $D\subseteq V$ and
            $U$ is disjoint from $V$ and
            $a_0\notin U$ and
            $a_0\notin V$.
    \begin{subproof}
        Fix $C$.
        Fix $D$ such that $C$ is closed in $X$ with respect to $T$ and $D$ is closed in $X$ with respect to $T$
            and $C$ is disjoint from $D$ and $a_0\notin C$ and $a_0\notin D$.
        Let $F = \{W\in\pow{X} \mid \text{there exists $c\in C$ such that there exists $x\in X$ such that
            $x\mathrel{R}c$ and $W = \intervalopenclosedrel{R}{X}{x}{c}$ and $W$ is disjoint from $D$}\}$.
        Let $G = \{W\in\pow{X} \mid \text{there exists $d\in D$ such that there exists $y\in X$ such that
            $y\mathrel{R}d$ and $W = \intervalopenclosedrel{R}{X}{y}{d}$ and $W$ is disjoint from $C$}\}$.
        Let $U = \unions{F}$.
        Let $V = \unions{G}$.

        Show that $U$ is open in $X$ with respect to $T$.
        \begin{subproof}
            Show that for all $W$ such that $W\in F$ we have $W\in T$.
            \begin{subproof}
                Fix $W$ such that $W\in F$.
                Take $c, x$ such that $c\in C$ and $x\in X$ and
                    $x\mathrel{R}c$ and $W = \intervalopenclosedrel{R}{X}{x}{c}$ and $W$ is disjoint from $D$.
                $c\in X$ by \cref{closed_in,subseteq}.
                $W$ is open in $X$ with respect to $T$ by \cref{well_order_intervalopenclosed_open}.
                $W\in T$ by \cref{open_in}.
            \end{subproof}
            $F\subseteq T$ by \cref{subseteq}.
            $U\in T$ and $U$ is open in $X$ with respect to $T$ by \cref{topology_on,open_in}.
        \end{subproof}

        Show that $V$ is open in $X$ with respect to $T$.
        \begin{subproof}
            Show that for all $W$ such that $W\in G$ we have $W\in T$.
            \begin{subproof}
                Fix $W$ such that $W\in G$.
                Take $d, y$ such that $d\in D$ and $y\in X$ and
                    $y\mathrel{R}d$ and $W = \intervalopenclosedrel{R}{X}{y}{d}$ and $W$ is disjoint from $C$.
                $d\in X$ by \cref{closed_in,subseteq}.
                $W$ is open in $X$ with respect to $T$ by \cref{well_order_intervalopenclosed_open}.
                $W\in T$ by \cref{open_in}.
            \end{subproof}
            $G\subseteq T$ by \cref{subseteq}.
            $V\in T$ and $V$ is open in $X$ with respect to $T$ by \cref{topology_on,open_in}.
        \end{subproof}

        Show that $C\subseteq U$.
        \begin{subproof}
            Show that for all $c$ such that $c\in C$ we have $c\in U$.
            \begin{subproof}
                Fix $c$ such that $c\in C$.
                $c\in X$ by \cref{closed_in,subseteq}.
                $c\notin D$ by \cref{disjoint}.
                $c\neq a_0$.
                $c$ is not a smallest element of $X$ with respect to $R$ by \cref{smallest_element,simple_order_relation,asymmetric}.
                $c\in X\setminus D$ by \cref{setminus_intro}.
                $X\setminus D$ is open in $X$ with respect to $T$ by \cref{closed_in}.
                Take $x\in X$ such that $x\mathrel{R}c$ and $\intervalopenclosedrel{R}{X}{x}{c}\subseteq X\setminus D$
                    by \cref{well_order_left_interval_neighborhood}.
                Let $W = \intervalopenclosedrel{R}{X}{x}{c}$.
                Show that $W$ is disjoint from $D$.
                \begin{subproof}
                    Suppose not.
                    Take $z$ such that $z\in W$ and $z\in D$ by \cref{disjoint}.
                    $z\in X\setminus D$ by \cref{subseteq}.
                    Contradiction by \cref{setminus_elim_right}.
                \end{subproof}
                $W\in F$ by \cref{intervalopenclosed_rel,subseteq,powerset_intro}.
                $c\in W$ by \cref{intervalopenclosed_rel}.
                $c\in U$ by \cref{unions_intro}.
            \end{subproof}
            $C\subseteq U$ by \cref{subseteq}.
        \end{subproof}

        Show that $D\subseteq V$.
        \begin{subproof}
            Show that for all $d$ such that $d\in D$ we have $d\in V$.
            \begin{subproof}
                Fix $d$ such that $d\in D$.
                $d\in X$ by \cref{closed_in,subseteq}.
                $d\notin C$ by \cref{disjoint}.
                $d\neq a_0$.
                $d$ is not a smallest element of $X$ with respect to $R$ by \cref{smallest_element,simple_order_relation,asymmetric}.
                $d\in X\setminus C$ by \cref{setminus_intro}.
                $X\setminus C$ is open in $X$ with respect to $T$ by \cref{closed_in}.
                Take $y\in X$ such that $y\mathrel{R}d$ and $\intervalopenclosedrel{R}{X}{y}{d}\subseteq X\setminus C$
                    by \cref{well_order_left_interval_neighborhood}.
                Let $W = \intervalopenclosedrel{R}{X}{y}{d}$.
                Show that $W$ is disjoint from $C$.
                \begin{subproof}
                    Suppose not.
                    Take $z$ such that $z\in W$ and $z\in C$ by \cref{disjoint}.
                    $z\in X\setminus C$ by \cref{subseteq}.
                    Contradiction by \cref{setminus_elim_right}.
                \end{subproof}
                $W\in G$ by \cref{intervalopenclosed_rel,subseteq,powerset_intro}.
                $d\in W$ by \cref{intervalopenclosed_rel}.
                $d\in V$ by \cref{unions_intro}.
            \end{subproof}
            $D\subseteq V$ by \cref{subseteq}.
        \end{subproof}

        Show that $a_0\notin U$.
        \begin{subproof}
            Suppose not.
            Take $W\in F$ such that $a_0\in W$ by \cref{unions_iff}.
            Take $c, x$ such that $c\in C$ and $x\in X$ and
                $x\mathrel{R}c$ and $W = \intervalopenclosedrel{R}{X}{x}{c}$ and $W$ is disjoint from $D$.
            $x\mathrel{R}a_0$ by \cref{intervalopenclosed_rel}.
            Contradiction by \cref{smallest_element,simple_order_relation,asymmetric}.
        \end{subproof}

        Show that $a_0\notin V$.
        \begin{subproof}
            Suppose not.
            Take $W\in G$ such that $a_0\in W$ by \cref{unions_iff}.
            Take $d, y$ such that $d\in D$ and $y\in X$ and
                $y\mathrel{R}d$ and $W = \intervalopenclosedrel{R}{X}{y}{d}$ and $W$ is disjoint from $C$.
            $y\mathrel{R}a_0$ by \cref{intervalopenclosed_rel}.
            Contradiction by \cref{smallest_element,simple_order_relation,asymmetric}.
        \end{subproof}

        Show that $U$ is disjoint from $V$.
        \begin{subproof}
            Suppose not.
            Take $z$ such that $z\in U$ and $z\in V$ by \cref{disjoint}.
            Take $W_1\in F$ such that $z\in W_1$ by \cref{unions_iff}.
            Take $W_2\in G$ such that $z\in W_2$ by \cref{unions_iff}.
            Take $a, x$ such that $a\in C$ and $x\in X$ and
                $x\mathrel{R}a$ and $W_1 = \intervalopenclosedrel{R}{X}{x}{a}$ and $W_1$ is disjoint from $D$.
            Take $b, y$ such that $b\in D$ and $y\in X$ and
                $y\mathrel{R}b$ and $W_2 = \intervalopenclosedrel{R}{X}{y}{b}$ and $W_2$ is disjoint from $C$.
            $a\in X$ and $b\in X$ by \cref{closed_in,subseteq}.
            $a\neq b$ by \cref{disjoint}.
            $a\mathrel{R}b$ or $b\mathrel{R}a$ by \cref{simple_order_relation,connex}.
            \begin{byCase}
            \caseOf{$a\mathrel{R}b$.}
                \begin{byCase}
                \caseOf{$y = a$.}
                    $a\mathrel{R}z$ by \cref{intervalopenclosed_rel}.
                    $z\mathrel{R}a$ or $z = a$ by \cref{intervalopenclosed_rel}.
                    Contradiction by \cref{simple_order_relation,asymmetric}.
                \caseOf{$y \neq a$.}
                    $y\mathrel{R}a$ or $a\mathrel{R}y$ by \cref{simple_order_relation,connex}.
                    \begin{byCase}
                    \caseOf{$y\mathrel{R}a$.}
                        $a\in W_2$ by \cref{intervalopenclosed_rel}.
                        Contradiction by \cref{disjoint}.
                    \caseOf{$a\mathrel{R}y$.}
                        $y\mathrel{R}z$ by \cref{intervalopenclosed_rel}.
                        $z\mathrel{R}a$ or $z = a$ by \cref{intervalopenclosed_rel}.
                        \begin{byCase}
                        \caseOf{$z\mathrel{R}a$.}
                            $a\mathrel{R}z$ by \cref{simple_order_relation,transitive}.
                            Contradiction by \cref{simple_order_relation,asymmetric}.
                        \caseOf{$z = a$.}
                            $y\mathrel{R}a$.
                            Contradiction by \cref{simple_order_relation,asymmetric}.
                        \end{byCase}
                    \end{byCase}
                \end{byCase}
            \caseOf{$b\mathrel{R}a$.}
                \begin{byCase}
                \caseOf{$x = b$.}
                    $b\mathrel{R}z$ by \cref{intervalopenclosed_rel}.
                    $z\mathrel{R}b$ or $z = b$ by \cref{intervalopenclosed_rel}.
                    Contradiction by \cref{simple_order_relation,asymmetric}.
                \caseOf{$x \neq b$.}
                    $x\mathrel{R}b$ or $b\mathrel{R}x$ by \cref{simple_order_relation,connex}.
                    \begin{byCase}
                    \caseOf{$x\mathrel{R}b$.}
                        $b\in W_1$ by \cref{intervalopenclosed_rel}.
                        Contradiction by \cref{disjoint}.
                    \caseOf{$b\mathrel{R}x$.}
                        $x\mathrel{R}z$ by \cref{intervalopenclosed_rel}.
                        $z\mathrel{R}b$ or $z = b$ by \cref{intervalopenclosed_rel}.
                        \begin{byCase}
                        \caseOf{$z\mathrel{R}b$.}
                            $b\mathrel{R}z$ by \cref{simple_order_relation,transitive}.
                            Contradiction by \cref{simple_order_relation,asymmetric}.
                        \caseOf{$z = b$.}
                            $x\mathrel{R}b$.
                            Contradiction by \cref{simple_order_relation,asymmetric}.
                        \end{byCase}
                    \end{byCase}
                \end{byCase}
            \end{byCase}
        \end{subproof}

        Show that there exist $U_0,V_0$ such that
            $U_0$ is open in $X$ with respect to $T$ and
            $V_0$ is open in $X$ with respect to $T$ and
            $C\subseteq U_0$ and
            $D\subseteq V_0$ and
            $U_0$ is disjoint from $V_0$ and
            $a_0\notin U_0$ and
            $a_0\notin V_0$.
        \begin{subproof}
            Trivial.
        \end{subproof}
    \end{subproof}
    Show that for all $A,B$ such that $A$ is closed in $X$ with respect to $T$ and $B$ is closed in $X$ with respect to $T$
        and $A$ is disjoint from $B$
        there exist $U,V$ such that
            $U$ is open in $X$ with respect to $T$ and
            $V$ is open in $X$ with respect to $T$ and
            $A\subseteq U$ and
            $B\subseteq V$ and
            $U$ is disjoint from $V$.
    \begin{subproof}
        Fix $A$.
        Fix $B$ such that $A$ is closed in $X$ with respect to $T$ and $B$ is closed in $X$ with respect to $T$
            and $A$ is disjoint from $B$.
        \begin{byCase}
        \caseOf{$a_0\in A$.}
            $a_0\notin B$ by \cref{disjoint}.
            Let $A_1 = A\setminus\{a_0\}$.
            $\{a_0\}$ is open in $X$ with respect to $T$ by \cref{well_order_smallest_singleton_open}.
            $X\setminus\{a_0\}$ is closed in $X$ with respect to $T$
                by \cref{singleton_subset_intro,setminus_subseteq,double_relative_complement,subseteq,open_in,closed_in}.
            $A\subseteq X$ and $\{a_0\}\subseteq X$ by \cref{closed_in,singleton_subset_intro,subseteq}.
            $A_1 = A\inter (X\setminus\{a_0\})$ by \cref{setminus_eq_inter_complement}.
            Let $H = \{A, X\setminus\{a_0\}\}$.
            $H\subseteq\pow{X}$ by \cref{upair_iff,closed_in,powerset_intro,subseteq}.
            Show that for all $K$ such that $K\in H$ we have $K$ is closed in $X$ with respect to $T$.
            \begin{subproof}
                Follows by \cref{upair_iff}.
            \end{subproof}
            $\inters{H}$ is closed in $X$ with respect to $T$ by \cref{topology_on,closed_inters}.
            $A_1$ is closed in $X$ with respect to $T$ by \cref{inter_as_inters}.
            Show that $a_0\notin A_1$.
            \begin{subproof}
                Suppose not.
                $a_0\in A_1$.
                $a_0\notin\{a_0\}$ by \cref{setminus_elim_right}.
                Contradiction by \cref{singleton_intro}.
            \end{subproof}
            Show that $A_1$ is disjoint from $B$.
            \begin{subproof}
                Suppose not.
                Take $z$ such that $z\in A_1$ and $z\in B$ by \cref{disjoint}.
                $z\in A$ by \cref{setminus_elim_left}.
                Contradiction by \cref{disjoint}.
            \end{subproof}
            Take $U_1,V_1$ such that
                $U_1$ is open in $X$ with respect to $T$ and
                $V_1$ is open in $X$ with respect to $T$ and
                $A_1\subseteq U_1$ and
                $B\subseteq V_1$ and
                $U_1$ is disjoint from $V_1$ and
                $a_0\notin U_1$ and
                $a_0\notin V_1$.
            Let $U = U_1\union\{a_0\}$.
            $\{a_0\}\subseteq U$ by \cref{union_upper_right,subseteq_refl}.
            Show that $A\subseteq U$.
            \begin{subproof}
                Show that for all $z$ such that $z\in A$ we have $z\in U$.
                \begin{subproof}
                    Fix $z$.
                    Assume $z\in A$.
                    \begin{byCase}
                    \caseOf{$z = a_0$.}
                        $z\in\{a_0\}$ by \cref{singleton_intro}.
                        $z\in U$ by \cref{union_intro_right}.
                    \caseOf{$z \neq a_0$.}
                        $z\notin\{a_0\}$ by \cref{singleton_iff}.
                        $z\in A_1$ by \cref{setminus_intro}.
                        $z\in U_1$ by \cref{subseteq}.
                        $z\in U$ by \cref{union_intro_left}.
                    \end{byCase}
                \end{subproof}
            \end{subproof}
            $\{a_0\}$ is open in $X$ with respect to $T$ and $U_1$ is open in $X$ with respect to $T$.
            $\{a_0\}\in T$ and $U_1\in T$ by \cref{open_in}.
            Let $M = \{U_1, \{a_0\}\}$.
            Show that for all $W$ such that $W\in M$ we have $W\in T$.
            \begin{subproof}
                Follows by \cref{upair_iff}.
            \end{subproof}
            $\unions{M}\in T$ by \cref{subseteq,topology_on}.
            $U = \unions{M}$ by \cref{union_as_unions}.
            $U$ is open in $X$ with respect to $T$ by \cref{open_in}.
            Show that $U$ is disjoint from $V_1$.
            \begin{subproof}
                Suppose not.
                Take $z$ such that $z\in U$ and $z\in V_1$ by \cref{disjoint}.
                $z\in U_1$ or $z\in\{a_0\}$ by \cref{union_iff}.
                \begin{byCase}
                \caseOf{$z\in U_1$.}
                    Contradiction by \cref{disjoint}.
                \caseOf{$z\in\{a_0\}$.}
                    $z = a_0$ by \cref{singleton_elim}.
                    Contradiction.
                \end{byCase}
            \end{subproof}
            Show that there exist $W,Z$ such that
                $W$ is open in $X$ with respect to $T$ and
                $Z$ is open in $X$ with respect to $T$ and
                $A\subseteq W$ and
                $B\subseteq Z$ and
                $W$ is disjoint from $Z$.
            \begin{subproof}
                Trivial.
            \end{subproof}
        \caseOf{$a_0\notin A$.}
            \begin{byCase}
            \caseOf{$a_0\in B$.}
                Let $B_1 = B\setminus\{a_0\}$.
                $\{a_0\}$ is open in $X$ with respect to $T$ by \cref{well_order_smallest_singleton_open}.
                $X\setminus\{a_0\}$ is closed in $X$ with respect to $T$
                    by \cref{singleton_subset_intro,setminus_subseteq,double_relative_complement,subseteq,open_in,closed_in}.
                $B\subseteq X$ and $\{a_0\}\subseteq X$ by \cref{closed_in,singleton_subset_intro,subseteq}.
                $B_1 = B\inter (X\setminus\{a_0\})$ by \cref{setminus_eq_inter_complement}.
                Let $H = \{B, X\setminus\{a_0\}\}$.
                $H\subseteq\pow{X}$ by \cref{upair_iff,closed_in,powerset_intro,subseteq}.
                Show that for all $K$ such that $K\in H$ we have $K$ is closed in $X$ with respect to $T$.
                \begin{subproof}
                    Follows by \cref{upair_iff}.
                \end{subproof}
                $\inters{H}$ is closed in $X$ with respect to $T$ by \cref{topology_on,closed_inters}.
                $B_1$ is closed in $X$ with respect to $T$ by \cref{inter_as_inters}.
                Show that $a_0\notin B_1$.
                \begin{subproof}
                    Suppose not.
                    $a_0\in B_1$.
                    $a_0\notin\{a_0\}$ by \cref{setminus_elim_right}.
                    Contradiction by \cref{singleton_intro}.
                \end{subproof}
                Show that $A$ is disjoint from $B_1$.
                \begin{subproof}
                    Suppose not.
                    Take $z$ such that $z\in A$ and $z\in B_1$ by \cref{disjoint}.
                    $z\in B$ by \cref{setminus_elim_left}.
                    Contradiction by \cref{disjoint}.
                \end{subproof}
                Take $U_1,V_1$ such that
                    $U_1$ is open in $X$ with respect to $T$ and
                    $V_1$ is open in $X$ with respect to $T$ and
                    $A\subseteq U_1$ and
                    $B_1\subseteq V_1$ and
                    $U_1$ is disjoint from $V_1$ and
                    $a_0\notin U_1$ and
                    $a_0\notin V_1$.
                Let $V_0 = V_1\union\{a_0\}$.
                $\{a_0\}\subseteq V_0$ by \cref{union_upper_right,subseteq_refl}.
                Show that $B\subseteq V_0$.
                \begin{subproof}
                    Show that for all $z$ such that $z\in B$ we have $z\in V_0$.
                    \begin{subproof}
                        Fix $z$.
                        Assume $z\in B$.
                        \begin{byCase}
                        \caseOf{$z = a_0$.}
                            $z\in\{a_0\}$ by \cref{singleton_intro}.
                            $z\in V_0$ by \cref{union_intro_right}.
                        \caseOf{$z \neq a_0$.}
                            $z\notin\{a_0\}$ by \cref{singleton_iff}.
                            $z\in B_1$ by \cref{setminus_intro}.
                            $z\in V_1$ by \cref{subseteq}.
                            $z\in V_0$ by \cref{union_intro_left}.
                        \end{byCase}
                    \end{subproof}
                \end{subproof}
                $\{a_0\}$ is open in $X$ with respect to $T$ and $V_1$ is open in $X$ with respect to $T$.
                $\{a_0\}\in T$ and $V_1\in T$ by \cref{open_in}.
                Let $M = \{V_1, \{a_0\}\}$.
                Show that for all $W$ such that $W\in M$ we have $W\in T$.
                \begin{subproof}
                    Follows by \cref{upair_iff}.
                \end{subproof}
                $\unions{M}\in T$ by \cref{subseteq,topology_on}.
                $V_0 = \unions{M}$ by \cref{union_as_unions}.
                $V_0$ is open in $X$ with respect to $T$ by \cref{open_in}.
                Show that $U_1$ is disjoint from $V_0$.
                \begin{subproof}
                    Suppose not.
                    Take $z$ such that $z\in U_1$ and $z\in V_0$ by \cref{disjoint}.
                    $z\in V_1$ or $z\in\{a_0\}$ by \cref{union_iff}.
                    \begin{byCase}
                    \caseOf{$z\in V_1$.}
                        Contradiction by \cref{disjoint}.
                    \caseOf{$z\in\{a_0\}$.}
                        $z = a_0$ by \cref{singleton_elim}.
                        Contradiction.
                    \end{byCase}
                \end{subproof}
                Show that there exist $W,Z$ such that
                    $W$ is open in $X$ with respect to $T$ and
                    $Z$ is open in $X$ with respect to $T$ and
                    $A\subseteq W$ and
                    $B\subseteq Z$ and
                    $W$ is disjoint from $Z$.
                \begin{subproof}
                    Trivial.
                \end{subproof}
            \caseOf{$a_0\notin B$.}
                Show that there exist $U,V$ such that
                    $U$ is open in $X$ with respect to $T$ and
                    $V$ is open in $X$ with respect to $T$ and
                    $A\subseteq U$ and
                    $B\subseteq V$ and
                    $U$ is disjoint from $V$ and
                    $a_0\notin U$ and
                    $a_0\notin V$.
                \begin{subproof}
                    Trivial.
                \end{subproof}
                Show that there exist $W,Z$ such that
                    $W$ is open in $X$ with respect to $T$ and
                    $Z$ is open in $X$ with respect to $T$ and
                    $A\subseteq W$ and
                    $B\subseteq Z$ and
                    $W$ is disjoint from $Z$.
                \begin{subproof}
                    Trivial.
                \end{subproof}
            \end{byCase}
        \end{byCase}
    \end{subproof}
    $X$ is normal with respect to $T$ by \cref{normal_space}.
\end{proof}
\section{§33 The Urysohn Lemma}

% from §33 helper definition 1: unit rational set
\begin{definition}\label{unit_rationals}
    $\unitRationals =
        \{\fst{q} \rmul \inv{\snd{q}} \mid q\in\naturalsPlus\times\naturalsPlus\}$.
\end{definition}

% from §33 helper definition 2: unit rational set with zero
\begin{definition}\label{unit_rationals_zero}
    $\unitRationalsZero = \unitRationals \union \{0\}$.
\end{definition}

% from §33 helper definition 3: Urysohn admissible functions
\begin{definition}\label{urysohn_admissible}
    $\urysohnAdmissible{s}{X}{T}{A}{B}{n} =
        \{f\in\funs{\{0,1\}\union \img{s}{\initialSegment{n}}}{\pow{X}} \mid
            \text{for all $p\in\dom{f}$ we have $f(p)\in T$ and
            $f(1) = X\setminus B$ and
            $A\subseteq f(0)$ and
            $\closure{f(0)}{X}{T}\subseteq f(1)$ and
            for all $p\in\dom{f}$ for all $q\in\dom{f}$ if $p < q$ then
                $\closure{f(p)}{X}{T}\subseteq f(q)$ and
            for all $p\in\dom{f}$ if $1 < p$ then $f(p) = X$}\}$.
\end{definition}

% from §33 helper definition 4: Urysohn good pair predicate
\begin{definition}\label{urysohn_good_pair}
    $\urysohnGoodPair{s}{X}{T}{A}{B}{k}{g}$ iff
        $k\in\naturalsPlus$ and $g\in\urysohnAdmissible{s}{X}{T}{A}{B}{k}$.
\end{definition}

% from §33 helper definition 5: Urysohn good pair set
\begin{definition}\label{urysohn_good_pairs}
    $\urysohnGoodPairs{s}{X}{T}{A}{B} =
        \{(k,g) \mid k\in\naturalsPlus,
            g\in\rels{\{0,1\}\union \img{s}{\naturalsPlus}}{\pow{X}} \mid
            \text{$\urysohnGoodPair{s}{X}{T}{A}{B}{k}{g}$}\}$.
\end{definition}

% from §33 helper definition 6: Urysohn level set
\begin{definition}\label{urysohn_level_set}
    $\urysohnLevelSet{U}{x} = \{q\in\unitRationalsZero \mid x\in U(q)\}$.
\end{definition}

% from §33 helper lemma 1: unit rationals are dense in nonnegative intervals
\begin{lemma}\label{unit_rationals_dense_nonneg}
    Suppose $a\in\reals$ and $b\in\reals$ and $0 \leq a$ and $a < b$.
    Then there exists $p$ such that $p\in\unitRationals$ and $p\in\intervalopen{a}{b}$.
\end{lemma}
\begin{proof}
    Take $n, N$ such that $n\in\naturalsPlus$ and $N\in\naturalsPlus$ and $n / N\in\intervalopen{a}{b}$
        by \cref{real_rationals_nonneg_interval}.
    Let $p = n / N$. Then $p\in\unitRationals$ and $p\in\intervalopen{a}{b}$
        by \cref{times_tuple_intro,real_div,fst_eq,snd_eq,unit_rationals}.
\end{proof}

% from §33 helper lemma 2: sequence enumerating unit rationals
\begin{lemma}\label{unit_rationals_sequence}
    There exists $s$ such that
        $s$ is a sequence in $\unitRationals$ and
        for all $p$ such that $p\in\unitRationals$ there exists $n\in\naturalsPlus$ such that $s(n) = p$.
\end{lemma}
\begin{proof}
    Take $g$ such that $g\in\inj{\naturalsPlus\times\naturalsPlus}{\naturalsPlus}$
        by \cref{naturalsplus_pairing_injection}.
    Let $p_0 = (1,1)$. Let $R = \converse{g} \union ((\naturalsPlus\setminus\ran{g})\times\{p_0\})$.
    Show that for all $x,y_1,y_2$ such that $(x,y_1)\in R$ and $(x,y_2)\in R$ we have $y_1 = y_2$.
    \begin{subproof}
        Fix $x$.
        Fix $y_1$.
        Fix $y_2$ such that $(x,y_1)\in R$ and $(x,y_2)\in R$.
        \begin{byCase}
        \caseOf{$x\in\ran{g}$.}
            We have $(x,y_1)\in\converse{g}$ and $(x,y_2)\in\converse{g}$
                by \cref{union_iff,times_tuple_elim,setminus_elim_right}.
            $y_1\mathrel{g} x$ and $y_2\mathrel{g} x$ by \cref{converse_iff}.
            Therefore $g(y_1) = x$ and $g(y_2) = x$ by \cref{inj,funs_is_function,function_apply_intro}.
            Thus $y_1 = y_2$ by \cref{dom_iff,funs,subseteq,inj}.
        \caseOf{$x\notin\ran{g}$.}
            We have $(x,y_1)\in(\naturalsPlus\setminus\ran{g})\times\{p_0\}$ and
                $(x,y_2)\in(\naturalsPlus\setminus\ran{g})\times\{p_0\}$ by \cref{union_iff,converse_iff,ran_iff}.
            Therefore $y_1 = y_2$ by \cref{times_tuple_elim,singleton_iff}.
        \end{byCase}
    \end{subproof}
    $R$ is a function by \cref{union_iff,converse_relation,times_elem_is_tuple,relation,rightunique}.
    For all $x\in\naturalsPlus$ we have $x\in\dom{R}$
        by \cref{ran_iff,converse_iff,union_iff,setminus_intro,singleton_intro,times_tuple_intro,dom_iff}.
    For all $x\in\dom{R}$ we have $x\in\naturalsPlus$
        by \cref{dom_iff,union_iff,converse_iff,inj,funs,funs_is_function,funs_type_apply,function_apply_intro,times_tuple_elim,setminus_elim_left}.
    $\dom{R} = \naturalsPlus$ by \cref{subseteq,subseteq_antisymmetric}.
    Let $h = \{(n, \fst{R(n)} \rmul \inv{\snd{R(n)}}) \mid n\in\naturalsPlus\}$.
    For all $n,y_1,y_2$ such that $(n,y_1)\in h$ and $(n,y_2)\in h$ we have $y_1 = y_2$
        by \cref{pair_eq_iff}.
    $h$ is a function by \cref{pair_eq_iff,relation,rightunique}.
    For all $n\in\dom{h}$ we have $h(n)\in\unitRationals$ by
        \cref{dom_iff,pair_eq_iff,subseteq,function_apply_elim,union_iff,converse_iff,inj,funs,times_tuple_elim,real_one_in_naturals,times_tuple_intro,singleton_subset_intro,fst_type,snd_type,unit_rationals,function_apply_intro}.
    For all $n\in\naturalsPlus$ we have $n\in\dom{h}$ by \cref{pair_eq_iff,dom_iff}.
    For all $n\in\dom{h}$ we have $n\in\naturalsPlus$ by \cref{dom_iff,pair_eq_iff}.
    We have $\dom{h} = \naturalsPlus$ by \cref{subseteq,subseteq_antisymmetric}.
    Show that $h$ is a sequence in $\unitRationals$.
    \begin{subproof}
        Follows by \cref{funs_intro,sequence_in}.
    \end{subproof}
    Show that for all $p$ such that $p\in\unitRationals$ there exists $n\in\naturalsPlus$ such that $h(n) = p$.
    \begin{subproof}
        Follows by
            \cref{unit_rationals,inj,funs_is_function,funs,subseteq,funs_type_apply,function_apply_elim,converse_iff,union_iff,dom_iff,function_apply_intro,pair_eq_iff,funs_is_relation}.
    \end{subproof}
    There exists $s$ such that
        $s$ is a sequence in $\unitRationals$ and
        for all $p$ such that $p\in\unitRationals$ there exists $n\in\naturalsPlus$ such that $s(n) = p$.
\end{proof}


% from §33 helper lemma 3: good pair decomposition
\begin{lemma}\label{urysohn_good_pair_decomp}
    Let $G = \urysohnGoodPairs{s}{X}{T}{A}{B}$.
    Suppose $z\in G$.
    Then there exist $n,f$ such that
        $n\in\naturalsPlus$ and
        $f\in\urysohnAdmissible{s}{X}{T}{A}{B}{n}$ and
        $z = (n,f)$.
\end{lemma}
\begin{proof}
    Take $n,f$ such that
        $n\in\naturalsPlus$ and
        $\urysohnGoodPair{s}{X}{T}{A}{B}{n}{f}$ and
        $z = (n,f)$.
    $f\in\urysohnAdmissible{s}{X}{T}{A}{B}{n}$ by \cref{urysohn_good_pair}.
\end{proof}

% from §33 helper lemma 4: good pair introduction
\begin{lemma}\label{urysohn_good_pair_intro}
    Let $G = \urysohnGoodPairs{s}{X}{T}{A}{B}$.
    Suppose $n\in\naturalsPlus$ and $f\in\urysohnAdmissible{s}{X}{T}{A}{B}{n}$.
    Then $(n,f)\in G$.
\end{lemma}
\begin{proof}
    $f\in\funs{\{0,1\}\union \img{s}{\initialSegment{n}}}{\pow{X}}$ by \cref{urysohn_admissible}.
    We have $\dom{f} = \{0,1\}\union \img{s}{\initialSegment{n}}$ by \cref{funs_elim}.
    $f\in\rels{\{0,1\}\union \img{s}{\initialSegment{n}}}{\pow{X}}$ by \cref{funs_subseteq_rels,elem_subseteq}.
    $f$ is a relation by \cref{rels_is_relation}.
    We have $\ran{f}\subseteq \pow{X}$ by \cref{rels_ran_subseteq}.
    $\initialSegment{n}\subseteq \naturalsPlus$ by \cref{subseteq,initial_segment_naturalsplus}.
    We have $\img{s}{\initialSegment{n}}\subseteq \img{s}{\naturalsPlus}$ by \cref{img_subseteq}.
    $\dom{f}\subseteq \{0,1\}\union \img{s}{\naturalsPlus}$
        by \cref{union_upper_left,union_upper_right,subseteq_transitive,union_subsets_is_subset}.
    Therefore $f\in\rels{\{0,1\}\union \img{s}{\naturalsPlus}}{\pow{X}}$ by \cref{rels_intro_dom_and_ran}.
    $(n,f)\in G$.
\end{proof}

% from §33 helper lemma 6: continuity at a point via open neighborhoods
\begin{lemma}\label{continuous_at_open}
    Assume $f$ is a function from $X$ to $Y$.
    Assume $T$ is a topology on $X$.
    Assume $T_1$ is a topology on $Y$.
    Assume $x\in X$.
    Suppose for all $V$ such that $V$ is open in $Y$ with respect to $T_1$ and $f(x)\in V$ there exists $U$ such that
        $U$ is open in $X$ with respect to $T$ and $x\in U$ and $\img{f}{U}\subseteq V$.
    Then $f$ is continuous at $x$ from $X$ to $Y$ with respect to $T$ and $T_1$.
\end{lemma}
\begin{proof}
    Follows by \cref{neighborhood,continuous_at}.
\end{proof}

% from §33 helper lemma 7: continuity to a subspace via range inclusion
\begin{lemma}\label{continuous_subspace_range}
    Assume $f\in\funs{X}{\reals}$.
    Assume $T$ is a topology on $X$.
    Assume $I_0\subseteq\reals$.
    Let $T_0 = \subspaceTopology{\standardTopologyR}{I_0}$.
    Assume $f$ is continuous from $X$ to $\reals$ with respect to $T$ and $\standardTopologyR$.
    Assume $\img{f}{X}\subseteq I_0$.
    Then $f$ is continuous from $X$ to $I_0$ with respect to $T$ and $T_0$.
\end{lemma}
    \begin{proof}
        We have $T_0$ is a topology on $I_0$
            by \cref{standard_basis_is_basis,basis_topology_is_topology,standard_topology_reals,subspace_of,subspace_topology_is_topology}.
        Show that for all $B$ such that $B$ is open in $I_0$ with respect to $T_0$
            we have $\preimg{f}{B}$ is open in $X$ with respect to $T$.
        \begin{subproof}
            Fix $B$ such that $B$ is open in $I_0$ with respect to $T_0$.
            Take $U\in\standardTopologyR$ such that $B = I_0\inter U$ by \cref{subspace_topology,open_in}.
            Show that for all $x$ such that $x\in\preimg{f}{B}$ we have $x\in\preimg{f}{U}$.
            \begin{subproof}
                Follows by \cref{preimg_iff,inter_elim_right,subseteq}.
            \end{subproof}
            Show that $\preimg{f}{B}\subseteq \preimg{f}{U}$.
            \begin{subproof}
                Follows by \cref{subseteq}.
            \end{subproof}
            Show that for all $x$ such that $x\in\preimg{f}{U}$ we have $x\in\preimg{f}{B}$.
            \begin{subproof}
                Fix $x$ such that $x\in\preimg{f}{U}$.
                We have $x\in\dom{f}$ by \cref{preimg_subseteq_dom,subseteq}.
                Hence $x\in X$ by \cref{funs,subseteq}.
                Thus $x\in\dom{f}\inter X$ by \cref{inter_intro}.
                Therefore $f(x)\in \img{f}{X}$ by \cref{funs,funs_is_function,img_of_function_intro}.
                $f(x)\in I_0$ by \cref{subseteq}.
                $f(x)\in U$ by \cref{preimg_iff,funs,funs_is_function,function_apply_iff}.
                Thus $f(x)\in I_0\inter U$ by \cref{inter_intro}.
        $x\in\preimg{f}{B}$ by \cref{preimg_iff,funs,funs_is_function,function_apply_intro}.
            \end{subproof}
            Show that $\preimg{f}{U}\subseteq \preimg{f}{B}$.
            \begin{subproof}
                Follows by \cref{subseteq}.
            \end{subproof}
            We have $\preimg{f}{B} = \preimg{f}{U}$ by \cref{subseteq_antisymmetric}.
            $\preimg{f}{B}$ is open in $X$ with respect to $T$ by \cref{open_in,continuous}.
        \end{subproof}
        $f$ is continuous from $X$ to $I_0$ with respect to $T$ and $T_0$
            by \cref{funs_is_function,funs_elim,img_dom,funs,subseteq,function_apply_elim,ran_iff,funs_intro,continuous}.
    \end{proof}

% from §33 helper lemma 8: initial good pair when $r_1 = 1$
\begin{lemma}\label{urysohn_u1_eq1}
    Assume $U_1 = X\setminus B$.
    Assume $U_1$ is open in $X$ with respect to $T$.
    Assume $U_0$ is open in $X$ with respect to $T$.
    Assume $A\subseteq U_0$.
    Assume $\closure{U_0}{X}{T}\subseteq U_1$.
    Assume $r_1 = 1$.
    Assume $\img{s}{\initialSegment{1}} = \{r_1\}$.
    Let $G = \urysohnGoodPairs{s}{X}{T}{A}{B}$.
    Then there exists $u_1$ such that $(1,u_1)\in G$.
\end{lemma}
\begin{proof}
    Let $u_1 = \{(0,U_0),(1,U_1)\}$.
    For all $a,b,b'$ such that $(a,b)\in u_1$ and $(a,b')\in u_1$ we have $b = b'$
        by \cref{upair_elim,pair_eq_iff,real_mul_ident}.
    $u_1$ is a function by \cref{upair_elim,relation,rightunique}.
    For all $p$ such that $p\in\dom{u_1}$ we have $p\in\{0,1\}$
        by \cref{dom_iff,upair_elim,pair_eq_iff,upair_intro_left,upair_intro_right}.
    Show that $\dom{u_1}\subseteq \{0,1\}$.
    \begin{subproof}
        Follows by \cref{subseteq}.
    \end{subproof}
    For all $p$ such that $p\in\{0,1\}$ we have $p\in\dom{u_1}$
        by \cref{upair_elim,upair_intro_left,upair_intro_right,dom_intro}.
    Show that $\{0,1\}\subseteq \dom{u_1}$.
    \begin{subproof}
        Follows by \cref{subseteq}.
    \end{subproof}
    $\dom{u_1} = \{0,1\}$ by \cref{subseteq_antisymmetric}.
    We have $\img{s}{\initialSegment{1}} = \{1\}$.
    Show that $\{0,1\}\union \img{s}{\initialSegment{1}}\subseteq \{0,1\}$.
    \begin{subproof}
        Follows by \cref{upair_intro_right,subseteq_refl,singleton_subset_intro,subseteq_transitive,union_subsets_is_subset,union_intro_left,subseteq}.
    \end{subproof}
    Show that $\{0,1\}\subseteq \{0,1\}\union \img{s}{\initialSegment{1}}$.
    \begin{subproof}
        Follows by \cref{subseteq_refl,union_intro_left,subseteq}.
    \end{subproof}
    $\{0,1\}\union \img{s}{\initialSegment{1}} = \{0,1\}$ by \cref{subseteq_antisymmetric}.
    $\dom{u_1} = \{0,1\}\union \img{s}{\initialSegment{1}}$ by \cref{subseteq_antisymmetric}.
    For all $p\in\dom{u_1}$ we have $u_1(p)\in\pow{X}$
        by \cref{subseteq,upair_elim,upair_intro_left,upair_intro_right,function_apply_intro,open_in,topology_on,elem_subseteq}.
    $u_1\in\funs{\{0,1\}\union \img{s}{\initialSegment{1}}}{\pow{X}}$ by \cref{funs_intro}.
    Show that for all $p\in\dom{u_1}$ we have
                $u_1(p)\in T$ and
                $u_1(1) = X\setminus B$ and
                $A\subseteq u_1(0)$ and
                $\closure{u_1(0)}{X}{T}\subseteq u_1(1)$ and
                for all $p\in\dom{u_1}$ for all $q\in\dom{u_1}$ if $p < q$ then
                    $\closure{u_1(p)}{X}{T}\subseteq u_1(q)$ and
                for all $p\in\dom{u_1}$ if $1 < p$ then $u_1(p) = X$.
    \begin{subproof}
        Fix $p$ such that $p\in\dom{u_1}$.
        We have $u_1$ is a function and $u_1$ is right-unique and $u_1$ is a relation
            by \cref{funs_is_function,function_rightunique,funs_is_relation}.
        $u_1(p)\in T$ by \cref{subseteq,upair_elim,upair_intro_left,upair_intro_right,function_apply_intro,open_in}.
        $u_1(1) = X\setminus B$ and $A\subseteq u_1(0)$ and
            $\closure{u_1(0)}{X}{T}\subseteq u_1(1)$
            by \cref{upair_intro_right,upair_intro_left,function_apply_intro,subseteq}.
        Show that for all $p_1\in\dom{u_1}$ for all $q_1\in\dom{u_1}$ if $p_1 < q_1$ then
            $\closure{u_1(p_1)}{X}{T}\subseteq u_1(q_1)$.
        \begin{subproof}
            Fix $p_1\in\dom{u_1}$.
            Fix $q_1\in\dom{u_1}$.
            Assume $p_1 < q_1$.
            We have $p_1\in\{0,1\}$ and $q_1\in\{0,1\}$.
            Hence $p_1 = 0$ and $q_1 = 1$
                by \cref{upair_elim,real_add_ident,real_naturals_subset_reals,real_one_in_naturals,subseteq,real_zero_lt_one,real_lt_irreflexive,real_lt_subst_left,real_lt_subst_right,real_lt_trans}.
            Thus $u_1(p_1) = U_0$ and $u_1(q_1) = U_1$
                by \cref{pair_eq_iff,upair_intro_left,upair_intro_right,function_apply_intro,funs_is_function,function_rightunique,funs_is_relation}.
            Therefore $\closure{u_1(p_1)}{X}{T}\subseteq u_1(q_1)$ by \cref{subseteq}.
        \end{subproof}
        Show that for all $p\in\dom{u_1}$ if $1 < p$ then $u_1(p) = X$.
        \begin{subproof}
            Follows by \cref{subseteq,upair_elim,real_add_ident,real_naturals_subset_reals,real_one_in_naturals,subseteq,real_zero_lt_one,real_lt_irreflexive,real_lt_subst_right,real_lt_trans}.
        \end{subproof}
    \end{subproof}
    $u_1\in\urysohnAdmissible{s}{X}{T}{A}{B}{1}$ by \cref{urysohn_admissible}.
    $(1,u_1)\in G$
        by \cref{urysohn_good_pair_intro,real_one_in_naturals}.
    There exists $u$ such that $(1,u)\in G$.
\end{proof}
% from §33 helper lemma 9: initial good pair when $1 < r_1$
\begin{lemma}\label{urysohn_u1_gt1}
    Assume $U_1 = X\setminus B$.
    Assume $U_1$ is open in $X$ with respect to $T$.
    Assume $U_0$ is open in $X$ with respect to $T$.
    Assume $A\subseteq U_0$.
    Assume $\closure{U_0}{X}{T}\subseteq U_1$.
    Assume $r_1\in\reals$.
    Assume $1 < r_1$.
    Assume $\img{s}{\initialSegment{1}} = \{r_1\}$.
    Let $G = \urysohnGoodPairs{s}{X}{T}{A}{B}$.
    Then there exists $u_1$ such that $(1,u_1)\in G$.
\end{lemma}
\begin{proof}
    Let $U_r = X$.
    Let $u_1 = \{(0,U_0),(1,U_1)\}\union\{(r_1,U_r)\}$.
    Show that $u_1\in\urysohnAdmissible{s}{X}{T}{A}{B}{1}$.
    \begin{subproof}
            Show that $u_1\in\funs{\{0,1\}\union \img{s}{\initialSegment{1}}}{\pow{X}}$.
            \begin{subproof}
                We have $r_1 \neq 0$ by \cref{real_add_ident,real_mul_ident,real_zero_lt_one,real_lt_trans,real_pos_nonzero}.
                Hence $r_1 \neq 1$ by \cref{real_lt_irreflexive,real_lt_subst_right}.
                For all $a,b,b'$ such that $(a,b)\in u_1$ and $(a,b')\in u_1$ we have $b = b'$
                    by \cref{union_iff,upair_elim,singleton_elim,pair_eq_iff,real_mul_ident}.
                $u_1$ is a function by \cref{union_iff,upair_elim,singleton_elim,relation,rightunique}.
                    Show that $\dom{u_1}\subseteq \{0,1\}\union\{r_1\}$.
                    \begin{subproof}
                        Follows by \cref{dom_iff,union_iff,upair_elim,singleton_elim,pair_eq_iff,upair_intro_left,upair_intro_right,singleton_intro,union_intro_left,union_intro_right,subseteq}.
                    \end{subproof}
                    Show that $\{0,1\}\union\{r_1\}\subseteq \dom{u_1}$.
                    \begin{subproof}
                        Follows by \cref{union_iff,upair_elim,singleton_elim,upair_intro_left,upair_intro_right,singleton_intro,union_intro_left,union_intro_right,dom_intro,subseteq}.
                    \end{subproof}
                    Hence $\dom{u_1} = \{0,1\}\union\{r_1\}$ by \cref{subseteq_antisymmetric}.
                    Show that $\{0,1\}\union \img{s}{\initialSegment{1}}\subseteq \{0,1\}\union\{r_1\}$.
                    \begin{subproof}
                        Follows by \cref{upair_intro_right,subseteq_refl,singleton_subset_intro,subseteq_transitive,union_subsets_is_subset,union_intro_left,subseteq}.
                    \end{subproof}
                    Show that $\{0,1\}\union\{r_1\}\subseteq \{0,1\}\union \img{s}{\initialSegment{1}}$.
                    \begin{subproof}
                        Follows by \cref{upair_intro_right,subseteq_refl,singleton_subset_intro,subseteq_transitive,union_subsets_is_subset,union_intro_right,subseteq}.
                    \end{subproof}
                    Thus $\{0,1\}\union \img{s}{\initialSegment{1}} = \{0,1\}\union\{r_1\}$ by \cref{subseteq_antisymmetric}.
                    Hence $\dom{u_1} = \{0,1\}\union \img{s}{\initialSegment{1}}$ by \cref{subseteq_antisymmetric}.
                For all $p\in\dom{u_1}$ we have $u_1(p)\in\pow{X}$
                    by \cref{subseteq,union_iff,upair_elim,upair_intro_left,upair_intro_right,singleton_elim,singleton_intro,union_intro_left,union_intro_right,function_apply_intro,open_in,topology_on,elem_subseteq}.
                Hence $u_1\in\funs{\{0,1\}\union \img{s}{\initialSegment{1}}}{\pow{X}}$ by \cref{funs_intro}.
            \end{subproof}
            Show that for all $p\in\dom{u_1}$ we have
                $u_1(p)\in T$ and
                $u_1(1) = X\setminus B$ and
                $A\subseteq u_1(0)$ and
                $\closure{u_1(0)}{X}{T}\subseteq u_1(1)$ and
                for all $p\in\dom{u_1}$ for all $q\in\dom{u_1}$ if $p < q$ then
                    $\closure{u_1(p)}{X}{T}\subseteq u_1(q)$ and
                for all $p\in\dom{u_1}$ if $1 < p$ then $u_1(p) = X$.
            \begin{subproof}
                Fix $p$ such that $p\in\dom{u_1}$.
                We have $u_1$ is a function and $u_1$ is right-unique and $u_1$ is a relation
                    by \cref{funs_is_function,function_rightunique,funs_is_relation}.
                Show that $u_1(p)\in T$.
                \begin{subproof}
                    We have $p\in\{0,1\}\union\{r_1\}$.
                    Hence $p\in\{0,1\}$ or $p\in\{r_1\}$ by \cref{union_iff}.
                    \begin{byCase}
                    \caseOf{$p\in\{0,1\}$.}
                        Hence $p = 0$ or $p = 1$ by \cref{upair_elim}.
                        \begin{byCase}
                        \caseOf{$p = 0$.}
                            $u_1(p)\in T$ by \cref{upair_intro_left,union_intro_left,function_apply_intro,open_in}.
                        \caseOf{$p = 1$.}
                            $u_1(p)\in T$ by \cref{upair_intro_right,union_intro_left,function_apply_intro,open_in}.
                        \end{byCase}
                    \caseOf{$p\in\{r_1\}$.}
                        $u_1(p)\in T$ by \cref{singleton_elim,singleton_intro,union_intro_right,function_apply_intro,open_in,topology_on}.
                    \end{byCase}
                \end{subproof}
                $u_1(1) = X\setminus B$ and $A\subseteq u_1(0)$ and
                    $\closure{u_1(0)}{X}{T}\subseteq u_1(1)$
                    by \cref{upair_intro_right,upair_intro_left,union_intro_left,function_apply_intro,subseteq}.
                Show that for all $p_1\in\dom{u_1}$ for all $q_1\in\dom{u_1}$ if $p_1 < q_1$ then
                    $\closure{u_1(p_1)}{X}{T}\subseteq u_1(q_1)$.
                \begin{subproof}
                    Fix $p_1\in\dom{u_1}$.
                    Fix $q_1\in\dom{u_1}$.
                    Assume $p_1 < q_1$.
                    We have $p_1\in\{0,1\}\union\{r_1\}$ and $q_1\in\{0,1\}\union\{r_1\}$.
                    Show that $q_1 \neq 0$.
                    \begin{subproof}
                        Suppose not.
                        Hence $q_1 = 0$.
                        Hence $p_1 < 0$ by \cref{real_lt_subst_right}.
                        We have $p_1\in\{0,1\}$ or $p_1\in\{r_1\}$ by \cref{union_iff}.
                        \begin{byCase}
                        \caseOf{$p_1\in\{0,1\}$.}
                            Hence $p_1 = 0$ or $p_1 = 1$ by \cref{upair_elim}.
                            \begin{byCase}
                            \caseOf{$p_1 = 0$.}
                                Contradiction by \cref{real_lt_subst_left,real_add_ident,real_lt_irreflexive}.
                            \caseOf{$p_1 = 1$.}
                                Hence $1 < 0$ by \cref{real_lt_subst_left}.
                                $0\in\reals$ and $1\in\reals$ and $0 < 1$ by \cref{real_add_ident,real_mul_ident,real_zero_lt_one}.
                                Hence $1 < 1$ by \cref{real_lt_trans}.
                                Contradiction by \cref{real_lt_irreflexive}.
                            \end{byCase}
                        \caseOf{$p_1\in\{r_1\}$.}
                            Hence $r_1 < 0$ by \cref{singleton_elim,real_lt_subst_left}.
                            $0\in\reals$ and $1\in\reals$ and $0 < 1$ by \cref{real_add_ident,real_mul_ident,real_zero_lt_one}.
                            Hence $1 < 0$ by \cref{real_lt_trans}.
                            Hence $1 < 1$ by \cref{real_lt_trans}.
                            Contradiction by \cref{real_lt_irreflexive}.
                        \end{byCase}
                    \end{subproof}
                    Show that $q_1\in\{1\}\union\{r_1\}$.
                    \begin{subproof}
                        We have $q_1\in\{0,1\}$ or $q_1\in\{r_1\}$ by \cref{union_iff}.
                        \begin{byCase}
                        \caseOf{$q_1\in\{0,1\}$.}
                            $q_1 = 0$ or $q_1 = 1$ by \cref{upair_elim}.
                            \begin{byCase}
                            \caseOf{$q_1 = 0$.}
                                Suppose not.
                                Contradiction.
                            \caseOf{$q_1 = 1$.}
                                Hence $q_1\in\{1\}$ by \cref{singleton_intro}.
                                Therefore $q_1\in\{1\}\union\{r_1\}$ by \cref{union_intro_left}.
                            \end{byCase}
                        \caseOf{$q_1\in\{r_1\}$.}
                            Hence $q_1\in\{1\}\union\{r_1\}$ by \cref{union_intro_right}.
                        \end{byCase}
                    \end{subproof}
                    Hence $q_1\in\{1\}$ or $q_1\in\{r_1\}$ by \cref{union_iff}.
                    \begin{byCase}
                    \caseOf{$q_1\in\{r_1\}$.}
                        We have $u_1(q_1) = X$ by \cref{singleton_elim,singleton_intro,union_intro_right,function_apply_intro}.
                        We have $u_1(p_1)\subseteq X$ by \cref{funs_type_apply,pow_iff}.
                        $T$ is a topology on $X$ by \cref{open_in}.
                        Hence $X$ is closed in $X$ with respect to $T$ by \cref{closed_whole_space}.
                        Hence $\closure{u_1(p_1)}{X}{T}\subseteq u_1(q_1)$
                            by \cref{closure_subseteq_closed,subseteq}.
                    \caseOf{$q_1\in\{1\}$.}
                        Hence $q_1 = 1$ by \cref{singleton_elim}.
                        Show that $p_1 = 0$.
                        \begin{subproof}
                            Suppose not.
                            We have $p_1\in\{0,1\}$ or $p_1\in\{r_1\}$ by \cref{union_iff}.
                            \begin{byCase}
                            \caseOf{$p_1\in\{0,1\}$.}
                                Hence $p_1 = 0$ or $p_1 = 1$ by \cref{upair_elim}.
                                \begin{byCase}
                                \caseOf{$p_1 = 1$.}
                                    Contradiction by \cref{real_lt_subst_left,real_mul_ident,real_lt_irreflexive}.
                                \caseOf{$p_1 = 0$.}
                                    Contradiction.
                                \end{byCase}
                            \caseOf{$p_1\in\{r_1\}$.}
                                Hence $r_1 < 1$ by \cref{singleton_elim,real_lt_subst_left}.
                                Contradiction.
                            \end{byCase}
                        \end{subproof}
                        Hence $p_1 = 0$.
                        Hence $\closure{u_1(p_1)}{X}{T}\subseteq u_1(q_1)$
                            by \cref{singleton_elim,pair_eq_iff,union_intro_left,upair_intro_left,upair_intro_right,function_apply_intro,subseteq}.
                    \end{byCase}
                \end{subproof}
                Show that for all $p_1\in\dom{u_1}$ if $1 < p_1$ then $u_1(p_1) = X$.
                \begin{subproof}
                    Fix $p_1\in\dom{u_1}$.
                    Assume $1 < p_1$.
                    Suppose not.
                    We have $p_1\in\{0,1\}\union\{r_1\}$.
                    Hence $p_1\in\{0,1\}$ or $p_1\in\{r_1\}$ by \cref{union_iff}.
                    \begin{byCase}
                    \caseOf{$p_1\in\{0,1\}$.}
                        Hence $p_1 = 0$ or $p_1 = 1$ by \cref{upair_elim}.
                        \begin{byCase}
                        \caseOf{$p_1 = 0$.}
                            Hence $1 < 0$ by \cref{real_lt_subst_left}.
                            $0\in\reals$ and $1\in\reals$ and $0 < 1$ by \cref{real_add_ident,real_mul_ident,real_zero_lt_one}.
                            Hence $1 < 1$ by \cref{real_lt_trans}.
                            Contradiction by \cref{real_lt_irreflexive}.
                        \caseOf{$p_1 = 1$.}
                            Contradiction by \cref{real_lt_subst_left,real_mul_ident,real_lt_irreflexive}.
                        \end{byCase}
                    \caseOf{$p_1\in\{r_1\}$.}
                        Hence $1 < r_1$ by \cref{singleton_elim,real_lt_subst_left}.
                        We have $u_1(p_1) = X$ by \cref{singleton_elim,singleton_intro,union_intro_right,function_apply_intro}.
                        Contradiction.
                    \end{byCase}
                \end{subproof}
            \end{subproof}
            Thus $u_1\in\urysohnAdmissible{s}{X}{T}{A}{B}{1}$ by \cref{urysohn_admissible}.
        \end{subproof}
        Therefore $(1,u_1)\in G$
            by \cref{urysohn_good_pair_intro,real_one_in_naturals}.
    There exists $u$ such that $(1,u)\in G$.
\end{proof}

% from §33 helper lemma 12: unit rational is positive
\begin{lemma}\label{unit_rational_pos}
    Assume $r_1\in\unitRationals$.
    Then $0 < r_1$.
\end{lemma}
\begin{proof}
    Take $q$ such that
        $q\in\naturalsPlus\times\naturalsPlus$ and
        $r_1 = \fst{q} \rmul \inv{\snd{q}}$
        by \cref{unit_rationals}.
    $r_1 > 0$
        by \cref{fst_type,snd_type,real_naturals_subset_reals,subseteq,real_add_ident,real_zero_lt_one,real_mul_ident,real_one_leq_naturalsplus,real_lt_subst_right,real_lt_trans,real_pos_nonzero,real_mul_inverse,real_inv_pos,real_mul_pos_pos}.
\end{proof}

% from §33 helper lemma 10: admissible good pair for $r_1 < 1$ given $U_r$
\begin{lemma}\label{urysohn_u1_lt1_admissible}
    Assume $U_1 = X\setminus B$.
    Assume $U_1$ is open in $X$ with respect to $T$.
    Assume $U_0$ is open in $X$ with respect to $T$.
    Assume $A\subseteq U_0$.
    Assume $\closure{U_0}{X}{T}\subseteq U_1$.
    Assume $r_1\in\reals$.
    Assume $r_1\in\unitRationals$.
    Assume $r_1 < 1$.
    Assume $\img{s}{\initialSegment{1}} = \{r_1\}$.
    Assume $U_r$ is open in $X$ with respect to $T$.
    Assume $\closure{U_0}{X}{T}\subseteq U_r$.
    Assume $\closure{U_r}{X}{T}\subseteq U_1$.
    Let $G = \urysohnGoodPairs{s}{X}{T}{A}{B}$.
    Then there exists $u_1$ such that $(1,u_1)\in G$.
\end{lemma}
\begin{proof}
    Let $u_1 = \{(0,U_0),(1,U_1)\}\union\{(r_1,U_r)\}$.
    Show that $(1,u_1)\in G$.
    \begin{subproof}
        Show that $u_1\in\urysohnAdmissible{s}{X}{T}{A}{B}{1}$.
        \begin{subproof}
            Show that $u_1\in\funs{\{0,1\}\union \img{s}{\initialSegment{1}}}{\pow{X}}$.
            \begin{subproof}
                We have $r_1 \neq 0$ by \cref{unit_rational_pos,real_pos_nonzero}.
                Hence $r_1 \neq 1$ by \cref{real_lt_irreflexive,real_lt_subst_right}.
                For all $a,b,b'$ such that $(a,b)\in u_1$ and $(a,b')\in u_1$ we have $b = b'$
                    by \cref{union_iff,upair_elim,singleton_elim,pair_eq_iff,real_mul_ident}.
                Hence $u_1$ is right-unique by \cref{rightunique}.
                Therefore $u_1$ is a function by \cref{union_iff,upair_elim,singleton_elim,relation,rightunique}.
                Show that $\dom{u_1} = \{0,1\}\union\{r_1\}$.
                \begin{subproof}
                    For all $x\in\dom{u_1}$ we have $x\in\{0,1\}\union\{r_1\}$
                        by \cref{dom_iff,union_iff,upair_elim,singleton_elim,pair_eq_iff,upair_intro_left,upair_intro_right,singleton_intro,union_intro_left,union_intro_right}.
                    Show that $\dom{u_1}\subseteq \{0,1\}\union\{r_1\}$.
                    \begin{subproof}
                        Follows by \cref{subseteq}.
                    \end{subproof}
                    For all $x\in\{0,1\}\union\{r_1\}$ we have $x\in\dom{u_1}$
                        by \cref{union_iff,upair_elim,singleton_elim,upair_intro_left,upair_intro_right,singleton_intro,union_intro_left,union_intro_right,dom_intro}.
                    Show that $\{0,1\}\union\{r_1\}\subseteq \dom{u_1}$.
                    \begin{subproof}
                        Follows by \cref{subseteq}.
                    \end{subproof}
                    Therefore $\dom{u_1} = \{0,1\}\union\{r_1\}$ by \cref{subseteq_antisymmetric}.
                \end{subproof}
                Hence $\dom{u_1} = \{0,1\}\union \img{s}{\initialSegment{1}}$.
                For all $p\in\dom{u_1}$ we have $u_1(p)\in\pow{X}$
                    by \cref{subseteq,union_iff,upair_elim,upair_intro_left,upair_intro_right,singleton_elim,singleton_intro,union_intro_left,union_intro_right,function_apply_intro,open_in,topology_on,elem_subseteq}.
                Therefore $u_1\in\funs{\{0,1\}\union \img{s}{\initialSegment{1}}}{\pow{X}}$ by \cref{funs_intro}.
            \end{subproof}
            Show that for all $p\in\dom{u_1}$ we have
                $u_1(p)\in T$ and
                $u_1(1) = X\setminus B$ and
                $A\subseteq u_1(0)$ and
                $\closure{u_1(0)}{X}{T}\subseteq u_1(1)$ and
                for all $p\in\dom{u_1}$ for all $q\in\dom{u_1}$ if $p < q$ then
                    $\closure{u_1(p)}{X}{T}\subseteq u_1(q)$ and
                for all $p\in\dom{u_1}$ if $1 < p$ then $u_1(p) = X$.
            \begin{subproof}
                Fix $p$ such that $p\in\dom{u_1}$.
                We have $u_1$ is a function and $u_1$ is right-unique and $u_1$ is a relation
                    by \cref{funs_is_function,function_rightunique,funs_is_relation}.
                Hence $u_1(p)\in T$
                    by \cref{funs_elim,subseteq,union_iff,upair_elim,upair_intro_left,upair_intro_right,singleton_elim,singleton_intro,union_intro_left,union_intro_right,function_apply_intro,open_in}.
                $u_1(1) = X\setminus B$ and $A\subseteq u_1(0)$ and
                    $\closure{u_1(0)}{X}{T}\subseteq u_1(1)$
                    by \cref{upair_intro_right,upair_intro_left,union_intro_left,function_apply_intro,subseteq}.
                Show that for all $p_1\in\dom{u_1}$ for all $q_1\in\dom{u_1}$ if $p_1 < q_1$ then
                    $\closure{u_1(p_1)}{X}{T}\subseteq u_1(q_1)$.
                \begin{subproof}
                    Fix $p_1\in\dom{u_1}$.
                    Fix $q_1\in\dom{u_1}$.
                    Assume $p_1 < q_1$.
                    We have $p_1\in\{0,1\}\union\{r_1\}$ and $q_1\in\{0,1\}\union\{r_1\}$.
                    Hence $q_1 \neq 0$
                        by \cref{union_iff,upair_elim,singleton_elim,real_lt_subst_right,real_lt_subst_left,real_add_ident,real_mul_ident,real_zero_lt_one,unit_rational_pos,real_lt_trans,real_lt_irreflexive}.
                    Hence $q_1\in\{1\}\union\{r_1\}$.
                    Hence $q_1\in\{1\}$ or $q_1\in\{r_1\}$ by \cref{union_iff}.
                    \begin{byCase}
                    \caseOf{$q_1\in\{r_1\}$.}
                        Hence $q_1 = r_1$ by \cref{singleton_elim}.
                        We have $p_1\in\{0,1\}$ or $p_1\in\{r_1\}$ by \cref{union_iff}.
                        Hence $p_1 = 0$
                            by \cref{upair_elim,singleton_elim,real_lt_subst_left,real_mul_ident,real_lt_trans,real_lt_irreflexive}.
                        Hence $\closure{u_1(p_1)}{X}{T}\subseteq u_1(q_1)$
                            by \cref{singleton_elim,pair_eq_iff,union_intro_left,upair_intro_left,union_intro_right,singleton_intro,function_apply_intro,subseteq}.
                    \caseOf{$q_1\in\{1\}$.}
                        Hence $q_1 = 1$ by \cref{singleton_elim}.
                        We have $p_1\in\{0,1\}$ or $p_1\in\{r_1\}$ by \cref{union_iff}.
                        \begin{byCase}
                        \caseOf{$p_1\in\{0,1\}$.}
                            Hence $p_1 = 0$
                                by \cref{upair_elim,real_lt_subst_left,real_mul_ident,real_lt_irreflexive}.
                            Hence $\closure{u_1(p_1)}{X}{T}\subseteq u_1(q_1)$
                                by \cref{pair_eq_iff,union_intro_left,upair_intro_left,upair_intro_right,function_apply_intro,subseteq}.
                        \caseOf{$p_1\in\{r_1\}$.}
                            Hence $\closure{u_1(p_1)}{X}{T}\subseteq u_1(q_1)$
                                by \cref{singleton_elim,pair_eq_iff,union_intro_right,singleton_intro,union_intro_left,upair_intro_right,function_apply_intro,subseteq}.
                        \end{byCase}
                    \end{byCase}
                \end{subproof}
                Show that for all $p\in\dom{u_1}$ if $1 < p$ then $u_1(p) = X$.
                \begin{subproof}
                    For all $p\in\dom{u_1}$ we have $p = 0$ or $p = 1$ or $p = r_1$
                        by \cref{funs_elim,subseteq,union_iff,upair_elim,singleton_elim}.
                    For all $p\in\dom{u_1}$ if $1 < p$ then $u_1(p) = X$
                        by \cref{real_lt_subst_left,real_add_ident,real_mul_ident,real_zero_lt_one,real_lt_trans,real_lt_irreflexive}.
                \end{subproof}
            \end{subproof}
            Hence $u_1\in\urysohnAdmissible{s}{X}{T}{A}{B}{1}$ by \cref{urysohn_admissible}.
        \end{subproof}
        $(1,u_1)\in G$ by \cref{urysohn_good_pair_intro,real_one_in_naturals}.
    \end{subproof}
\end{proof}

% from §33 helper lemma 11: initial good pair when $r_1 < 1$
\begin{lemma}\label{urysohn_u1_lt1}
    Assume $X$ is normal with respect to $T$.
    Assume $U_1 = X\setminus B$.
    Assume $U_1$ is open in $X$ with respect to $T$.
    Assume $U_0$ is open in $X$ with respect to $T$.
    Assume $A\subseteq U_0$.
    Assume $\closure{U_0}{X}{T}\subseteq U_1$.
    Assume $r_1\in\reals$.
    Assume $r_1\in\unitRationals$.
    Assume $r_1 < 1$.
    Assume $\img{s}{\initialSegment{1}} = \{r_1\}$.
    Let $G = \urysohnGoodPairs{s}{X}{T}{A}{B}$.
    Then there exists $u_1$ such that $(1,u_1)\in G$.
\end{lemma}
\begin{proof}
    We have $\closure{U_0}{X}{T}$ is closed in $X$ with respect to $T$ by \cref{open_in,topology_on,elem_subseteq,pow_iff,closure_closed_in}.
    Take $U_r$ such that
        $U_r$ is open in $X$ with respect to $T$ and
        $\closure{U_0}{X}{T}\subseteq U_r$ and
        $\closure{U_r}{X}{T}\subseteq U_1$
        by \cref{normal_closure_neighborhood}.
    There exists $u_1$ such that $(1,u_1)\in G$ by \cref{urysohn_u1_lt1_admissible}.
\end{proof}

% from §33 helper lemma 13: admissible chain is linearly ordered by inclusion
\begin{lemma}\label{urysohn_unit_interval_chain}
    Assume $s$ is a sequence in $\unitRationals$.
    Let $G = \urysohnGoodPairs{s}{X}{T}{A}{B}$.
    Let $R = \{(z,w) \mid z\in G, w\in G \mid
        \fst{w} = \fst{z} + 1 \land \snd{z}\subseteq\snd{w}\}$.
    Assume $c$ is a function on $G$.
    Assume for all $z\in G$ we have $z\mathrel{R} \apply{c}{z}$.
    Assume $u$ is a sequence in $G$.
    Assume for all $n\in\naturalsPlus$ we have $u(n + 1) = \apply{c}{u(n)}$.
    Let $F = \{\snd{u(n)} \mid n\in\naturalsPlus\}$.
    Then for all $f,g$ such that $f\in F$ and $g\in F$ we have $f\subseteq g$ or $g\subseteq f$.
\end{lemma}
\begin{proof}
    For all $n\in\naturalsPlus$ we have $\snd{u(n)}\subseteq \snd{u(n + 1)}$
        by \cref{sequence_in,funs_type_apply,pair_eq_iff}.
    Let $A_0 = \{m\in\naturalsPlus \mid \text{for all $n\in\naturalsPlus$ if $n \leq m$ then
        $\snd{u(n)}\subseteq \snd{u(m)}$}\}$.
    $A_0\subseteq\naturalsPlus$ by \cref{subseteq}.
    We have $1\in A_0$
        by \cref{real_one_in_naturals,real_naturals_subset_reals,subseteq,real_one_leq_naturalsplus,real_leq_antisym,subseteq_refl}.
    Show that for all $k\in A_0$ we have $k + 1\in A_0$.
    \begin{subproof}
        Fix $k$ such that $k\in A_0$.
        We have $k\in\naturalsPlus$ by \cref{subseteq}.
        Hence $k + 1\in\naturalsPlus$ by \cref{real_naturals_closed_succ}.
        Show that for all $n\in\naturalsPlus$ if $n \leq k + 1$ then $\snd{u(n)}\subseteq \snd{u(k + 1)}$.
        \begin{subproof}
            Fix $n\in\naturalsPlus$.
            Assume $n \leq k + 1$.
            Hence $n \leq k$ or $n = k + 1$ by \cref{real_naturals_leq_succ_elim}.
            \begin{byCase}
            \caseOf{$n \leq k$.}
                Hence $\snd{u(n)}\subseteq \snd{u(k)}$ by definition.
                $\snd{u(k)}\subseteq \snd{u(k + 1)}$ by \cref{sequence_in,funs_type_apply,pair_eq_iff}.
                Therefore $\snd{u(n)}\subseteq \snd{u(k + 1)}$ by \cref{subseteq_transitive}.
            \caseOf{$n = k + 1$.}
                We have $\snd{u(n)} = \snd{u(k + 1)}$.
                Therefore $\snd{u(n)}\subseteq \snd{u(k + 1)}$ by \cref{subseteq_refl}.
            \end{byCase}
        \end{subproof}
        $k + 1\in A_0$ by definition.
    \end{subproof}
    For all $m\in\naturalsPlus$ we have $m\in A_0$ by \cref{real_naturals_induction}.
    Fix $f$.
    Fix $g$ such that $f\in F$ and $g\in F$.
    Take $n$ such that $n\in\naturalsPlus$ and $f = \snd{u(n)}$.
    Take $m$ such that $m\in\naturalsPlus$ and $g = \snd{u(m)}$.
    We have $n < m$ or $n = m$ or $m < n$ by \cref{real_naturals_subset_reals,subseteq,real_lt_trichotomy}.
    Hence $f\subseteq g$ or $g\subseteq f$ by \cref{subseteq_refl,real_lt_imp_leq}.
\end{proof}

% from §33 helper lemma 14: unions of the admissible chain are a function
\begin{lemma}\label{urysohn_unit_interval_union_function}
    Assume $s$ is a sequence in $\unitRationals$.
    Let $G = \urysohnGoodPairs{s}{X}{T}{A}{B}$.
    Let $R = \{(z,w) \mid z\in G, w\in G \mid
        \fst{w} = \fst{z} + 1 \land \snd{z}\subseteq\snd{w}\}$.
    Assume $c$ is a function on $G$.
    Assume for all $z\in G$ we have $z\mathrel{R} \apply{c}{z}$.
    Assume $u$ is a sequence in $G$.
    Assume for all $n\in\naturalsPlus$ we have $u(n + 1) = \apply{c}{u(n)}$.
    Assume for all $n\in\naturalsPlus$ we have $\fst{u(n)} = n$.
    Let $F = \{\snd{u(n)} \mid n\in\naturalsPlus\}$.
    Let $U = \unions{F}$.
    Then $U$ is a function.
\end{lemma}
\begin{proof}
    We have $U$ is a function
        by \cref{sequence_in,funs_type_apply,urysohn_good_pair_decomp,snd_eq,urysohn_admissible,funs_is_function,urysohn_unit_interval_chain,unions_of_compatible_family_of_function_is_function}.
\end{proof}

% from §33 helper lemma 20d: sequence values occur in the initial image
\begin{lemma}\label{sequence_value_in_initial_image}
    Assume $s$ is a sequence in $Y$.
    Assume $k\in\naturalsPlus$.
    Then $s(k)\in\img{s}{\initialSegment{k}}$.
\end{lemma}
\begin{proof}
    $s(k)\in\img{s}{\initialSegment{k}}$ by \cref{sequence_in,funs_is_function,funs_elim,initial_segment_naturalsplus,subseteq,function_apply_elim,img_iff}.
\end{proof}

% from §33 helper lemma 20a: unit rationals are real (local helper)
\begin{lemma}\label{unit_rational_real_aux}
    Assume $r\in\unitRationals$.
    Then $r\in\reals$.
\end{lemma}
\begin{proof}
    Take $q$ such that
        $q\in\naturalsPlus\times\naturalsPlus$ and
        $r = \fst{q} \rmul \inv{\snd{q}}$
        by \cref{unit_rationals}.
    $r\in\reals$
        by \cref{fst_type,snd_type,real_naturals_subset_reals,subseteq,real_add_ident,real_mul_ident,real_zero_lt_one,real_one_leq_naturalsplus,real_lt_subst_right,real_lt_trans,real_pos_nonzero,real_mul_inverse,real_mul_closed}.
\end{proof}

% from §33 helper lemma 20b: unit rationals zero are real
\begin{lemma}\label{unit_rational_zero_real}
    Assume $r\in\unitRationalsZero$.
    Then $r\in\reals$.
\end{lemma}
\begin{proof}
    We have $r\in\reals$
        by \cref{unit_rationals_zero,union_iff,unit_rational_real_aux,singleton_elim,real_add_ident}.
\end{proof}

% from §33 helper lemma 18a: Urysohn level values are real
\begin{lemma}\label{urysohn_level_value_real}
    Assume $r\in\urysohnLevelSet{U}{x}$.
    Then $r\in\reals$.
\end{lemma}
\begin{proof}
    We have $r\in\reals$ by \cref{urysohn_level_set,unit_rational_zero_real}.
\end{proof}

% from §33 helper lemma 18aa: the infimum graph definition can be eliminated
\begin{lemma}\label{urysohn_infimum_graph_elim}
    Let $f = \{(x,y) \mid x\in X, y\in\reals \mid
        \text{$y$ is an infimum of $\urysohnLevelSet{U}{x}$}\}$.
    Assume $(x,y)\in f$.
    Then $y$ is an infimum of $\urysohnLevelSet{U}{x}$.
\end{lemma}
\begin{proof}
    Follows by \cref{pair_eq_iff}.
\end{proof}

% from §33 helper lemma 15: Urysohn function is a function to reals
\begin{lemma}\label{urysohn_unit_interval_f_funs}
    Assume $s$ is a sequence in $\unitRationals$.
    Assume for all $p$ such that $p\in\unitRationals$ there exists $n\in\naturalsPlus$ such that $s(n) = p$.
    Let $G = \urysohnGoodPairs{s}{X}{T}{A}{B}$.
    Assume $u$ is a sequence in $G$.
    Assume for all $n\in\naturalsPlus$ we have $\fst{u(n)} = n$.
    Let $F = \{\snd{u(n)} \mid n\in\naturalsPlus\}$.
    Let $U = \unions{F}$.
    Assume $U$ is a function.
    Let $f = \{(x,y) \mid x\in X, y\in\reals \mid
        \text{$y$ is an infimum of $\urysohnLevelSet{U}{x}$}\}$.
    Then $f$ is a function from $X$ to $\reals$.
\end{lemma}
\begin{proof}
        For all $x,y_1,y_2$ such that $(x,y_1)\in f$ and $(x,y_2)\in f$ we have $y_1 = y_2$
            by \cref{urysohn_infimum_graph_elim,real_infimum_unique}.
        $f$ is right-unique by \cref{rightunique}.
        $f$ is a function by \cref{relation,rightunique}.
            Show that for all $x\in X$ we have $x\in\dom{f}$.
            \begin{subproof}
                Fix $x\in X$.
                Let $Q_x = \urysohnLevelSet{U}{x}$.
                Show that $Q_x\neq \emptyset$.
                \begin{subproof}
                    $0\in\reals$ and $1\in\reals$ and $0 < 1$
                        by \cref{real_add_ident,real_mul_ident,real_zero_lt_one}.
                    $0 \leq 1$ by \cref{real_lt_imp_leq}.
                    Let $b = 1 + 1$.
                    We have $b\in\reals$ and $1 < b$
                        by \cref{real_add_closed,real_lt_add,real_add_ident,real_lt_subst_left}.
                    Take $p$ such that $p\in\unitRationals$ and $p\in\intervalopen{1}{b}$
                        by \cref{unit_rationals_dense_nonneg}.
                    $1 < p$ by \cref{intervalopen}.
                    Take $k\in\naturalsPlus$ such that $s(k) = p$ by assumption.
                    $u(k)\in G$ by \cref{sequence_in,funs_type_apply}.
                    Take $n,t$ such that
                        $t\in\urysohnAdmissible{s}{X}{T}{A}{B}{n}$ and
                        $u(k) = (n,t)$
                        by \cref{urysohn_good_pair_decomp}.
                    $k = n$ and $\snd{u(k)} = t$ by \cref{fst_eq,snd_eq}.
                    We have $t$ is a function and $\dom{t} = \{0,1\}\union \img{s}{\initialSegment{k}}$
                        by \cref{urysohn_admissible,funs_is_function,funs_elim}.
                    $p\in\dom{t}$ by \cref{sequence_value_in_initial_image,urysohn_admissible,funs_elim,union_intro_right}.
                    $1\in\dom{t}$ by \cref{upair_intro_right,union_intro_left,subseteq}.
                    We have for all $q\in\dom{t}$ if $1 < q$ then $t(q) = X$
                        by \cref{urysohn_admissible}.
                    $t(p) = X$.
                    $(p,t(p))\in U$ by \cref{function_apply_elim,unions_intro}.
                    $U(p) = t(p)$ by \cref{function_apply_intro}.
                    $x\in U(p)$.
                    $p\in Q_x$ by \cref{unit_rationals_zero,union_intro_left,urysohn_level_set}.
                    $Q_x\neq \emptyset$.
                \end{subproof}
                Show that $Q_x\subseteq\reals$ and $Q_x$ is bounded below.
                \begin{subproof}
                    For all $p\in Q_x$ we have $p\in\reals$ by \cref{urysohn_level_value_real}.
                    $Q_x\subseteq\reals$ by \cref{subseteq}.
                    For all $p\in Q_x$ we have $0 \leq p$
                        by \cref{urysohn_level_set,unit_rationals_zero,union_iff,singleton_elim,real_add_ident,unit_rational_pos,real_lt_imp_leq}.
                    $0$ is a lower bound of $Q_x$
                        by \cref{real_add_ident,subseteq,real_lower_bound}.
                    $Q_x$ is bounded below.
                \end{subproof}
                Take $p$ such that $p$ is an infimum of $Q_x$ by \cref{real_infimum_exists}.
                $x\in\dom{f}$ by \cref{real_infimum,real_lower_bound,dom_intro}.
            \end{subproof}
            Show that $X\subseteq\dom{f}$.
            \begin{subproof}
                Follows by \cref{subseteq}.
            \end{subproof}
        Show that $\dom{f}\subseteq X$.
        \begin{subproof}
            Follows by \cref{dom_iff,pair_eq_iff,subseteq}.
        \end{subproof}
        $\dom{f} = X$ by \cref{subseteq_antisymmetric}.
        For all $x\in\dom{f}$ we have $f(x)\in\reals$
            by \cref{dom_iff,urysohn_infimum_graph_elim,real_infimum,real_lower_bound,function_apply_intro}.
    $f\in\funs{X}{\reals}$ by \cref{funs_intro}.
\end{proof}

% from §33 helper lemma 16: existence of a first good pair
\begin{lemma}\label{urysohn_unit_interval_u1}
    Assume $X$ is normal with respect to $T$.
    Assume $A$ is closed in $X$ with respect to $T$.
    Assume $B$ is closed in $X$ with respect to $T$.
    Assume $A$ is disjoint from $B$.
    Assume $s$ is a sequence in $\unitRationals$.
    Let $G = \urysohnGoodPairs{s}{X}{T}{A}{B}$.
    Then there exists $u_1$ such that $(1,u_1)\in G$.
\end{lemma}
\begin{proof}
    Let $U_1 = X\setminus B$.
    $U_1$ is open in $X$ with respect to $T$ by \cref{closed_in}.
    $A\subseteq U_1$ by \cref{disjoint,closed_in,subseteq,setminus_intro}.
    Take $U_0$ such that
        $U_0$ is open in $X$ with respect to $T$ and
        $A\subseteq U_0$ and
        $\closure{U_0}{X}{T}\subseteq U_1$
        by \cref{normal_closure_neighborhood}.
    Let $r_1 = s(1)$.
    $s\in\funs{\naturalsPlus}{\unitRationals}$ and $1\in\naturalsPlus$
        by \cref{sequence_in,real_one_in_naturals}.
    $r_1\in\unitRationals$ by \cref{funs_type_apply}.
    Show that $\img{s}{\initialSegment{1}}\subseteq \{r_1\}$.
    \begin{subproof}
        Follows by \cref{img_iff,initial_segment_naturalsplus,real_one_leq_naturalsplus,real_naturals_subset_reals,subseteq,real_leq_antisym,funs_is_function,function_apply_intro,singleton_intro}.
    \end{subproof}
    Show that $\{r_1\}\subseteq \img{s}{\initialSegment{1}}$.
    \begin{subproof}
        Follows by \cref{sequence_value_in_initial_image,real_one_in_naturals,singleton_subset_intro}.
    \end{subproof}
    $\img{s}{\initialSegment{1}} = \{r_1\}$ by \cref{subseteq_antisymmetric}.
    \begin{byCase}
    \caseOf{$r_1 = 1$.}
        Show that there exists $u_1$ such that $(1,u_1)\in G$.
        \begin{subproof}
            Follows by \cref{urysohn_u1_eq1}.
        \end{subproof}
    \caseOf{$r_1 \neq 1$.}
        Show that there exists $u_1$ such that $(1,u_1)\in G$.
        \begin{subproof}
            Follows by \cref{unit_rational_real_aux,real_mul_ident,real_lt_trichotomy,urysohn_u1_gt1,urysohn_u1_lt1}.
        \end{subproof}
    \end{byCase}
\end{proof}

% from §33 helper lemma 17: indices of the iteration sequence
\begin{lemma}\label{urysohn_unit_interval_fst_u}
    Assume $s$ is a sequence in $\unitRationals$.
    Let $G = \urysohnGoodPairs{s}{X}{T}{A}{B}$.
    Let $R = \{(z,w) \mid z\in G, w\in G \mid
        \fst{w} = \fst{z} + 1 \land \snd{z}\subseteq\snd{w}\}$.
    Assume $c$ is a function on $G$.
    Assume for all $z\in G$ we have $z\mathrel{R} \apply{c}{z}$.
    Assume $u$ is a sequence in $G$.
    Assume $u(1) = a_0$.
    Assume $a_0 = (1,u_1)$.
    Assume for all $n\in\naturalsPlus$ we have $u(n + 1) = \apply{c}{u(n)}$.
    Then for all $n\in\naturalsPlus$ we have $\fst{u(n)} = n$.
\end{lemma}
\begin{proof}
    Let $A_0 = \{n\in\naturalsPlus \mid \fst{u(n)} = n\}$.
    $A_0\subseteq\naturalsPlus$ by \cref{subseteq}.
    We have $1\in A_0$ by \cref{real_one_in_naturals,fst_eq}.
    Show that for all $k\in A_0$ we have $k + 1\in A_0$.
    \begin{subproof}
        Follows by \cref{subseteq,sequence_in,funs_type_apply,pair_eq_iff,real_naturals_closed_succ}.
    \end{subproof}
    For all $n\in\naturalsPlus$ we have $n\in A_0$ by \cref{real_naturals_induction}.
    For all $n\in\naturalsPlus$ we have $\fst{u(n)} = n$.
\end{proof}

% from §33 helper lemma 18: continuous subspace range
\begin{lemma}\label{urysohn_unit_interval_continuous_subspace}
    Assume $f\in\funs{X}{\reals}$.
    Assume $T$ is a topology on $X$.
    Assume $I_0\subseteq\reals$.
    Let $T_0 = \subspaceTopology{\standardTopologyR}{I_0}$.
    Assume $f$ is continuous from $X$ to $\reals$ with respect to $T$ and $\standardTopologyR$.
    Assume $\img{f}{X}\subseteq I_0$.
    Then $f$ is continuous from $X$ to $I_0$ with respect to $T$ and $T_0$.
\end{lemma}
\begin{proof}
    We have $f$ is continuous from $X$ to $I_0$ with respect to $T$ and $T_0$
        by \cref{continuous_subspace_range}.
\end{proof}

% from §33 helper lemma 20da: rational values of U come from an admissible stage
\begin{lemma}\label{urysohn_union_rational_value_aux}
    Assume $s$ is a sequence in $\unitRationals$.
    Assume for all $p$ such that $p\in\unitRationals$ there exists $n\in\naturalsPlus$ such that $s(n) = p$.
    Let $G = \urysohnGoodPairs{s}{X}{T}{A}{B}$.
    Assume $u$ is a sequence in $G$.
    Assume for all $n\in\naturalsPlus$ we have $\fst{u(n)} = n$.
    Let $F = \{\snd{u(n)} \mid n\in\naturalsPlus\}$.
    Let $U = \unions{F}$.
    Assume $U$ is a function.
    Assume $q\in\unitRationals$.
    Then there exist $n,t$ such that
        $t\in\urysohnAdmissible{s}{X}{T}{A}{B}{n}$ and
        $t$ is a function and
        $q\in\dom{t}$ and
        $t\in F$ and
        $U(q) = t(q)$.
\end{lemma}
\begin{proof}
    Take $k$ such that $k\in\naturalsPlus$ and $s(k) = q$
        by \cref{unit_rationals_sequence}.
    $u(k)\in G$ by \cref{sequence_in,funs_type_apply}.
    Take $n,t$ such that
        $t\in\urysohnAdmissible{s}{X}{T}{A}{B}{n}$ and
        $u(k) = (n,t)$
        by \cref{urysohn_good_pair_decomp}.
    $\fst{u(k)} = n$ and $\snd{u(k)} = t$ by \cref{fst_eq,snd_eq}.
    $n\in\naturalsPlus$ and $q = s(n)$ and $t\in F$ by \cref{subseteq}.
    $t$ is a function by \cref{urysohn_admissible,funs_is_function}.
    $q\in\dom{t}$ by \cref{urysohn_admissible,sequence_value_in_initial_image,funs_elim,union_intro_right}.
    $U(q) = t(q)$ by \cref{function_apply_elim,unions_intro,function_apply_intro}.
\end{proof}

% from §33 helper lemma 18b: rational U-values are open
\begin{lemma}\label{urysohn_union_rational_open}
    Assume $X$ is normal with respect to $T$.
    Assume $A$ is closed in $X$ with respect to $T$.
    Assume $B$ is closed in $X$ with respect to $T$.
    Assume $A$ is disjoint from $B$.
    Assume $s$ is a sequence in $\unitRationals$.
    Assume for all $p$ such that $p\in\unitRationals$ there exists $n\in\naturalsPlus$ such that $s(n) = p$.
    Let $G = \urysohnGoodPairs{s}{X}{T}{A}{B}$.
    Assume $u$ is a sequence in $G$.
    Assume for all $n\in\naturalsPlus$ we have $\fst{u(n)} = n$.
    Let $F = \{\snd{u(n)} \mid n\in\naturalsPlus\}$.
    Let $U = \unions{F}$.
    Assume $U$ is a function.
    Assume $q\in\unitRationals$.
    Then $U(q)$ is open in $X$ with respect to $T$.
\end{lemma}
\begin{proof}
    Take $n,t$ such that
        $t\in\urysohnAdmissible{s}{X}{T}{A}{B}{n}$ and
        $q\in\dom{t}$ and
        $U(q) = t(q)$
        by \cref{urysohn_union_rational_value_aux}.
    $U(q)$ is open in $X$ with respect to $T$
        by \cref{urysohn_admissible,subseteq,function_apply_elim,unions_intro,function_apply_intro,normal_space,open_in}.
\end{proof}

% from §33 helper lemma 18c: zero-or-rational U-values lie in X
\begin{lemma}\label{urysohn_union_zero_subset_X_aux}
    Assume $s$ is a sequence in $\unitRationals$.
    Assume for all $p$ such that $p\in\unitRationals$ there exists $n\in\naturalsPlus$ such that $s(n) = p$.
    Let $G = \urysohnGoodPairs{s}{X}{T}{A}{B}$.
    Assume $u$ is a sequence in $G$.
    Assume for all $n\in\naturalsPlus$ we have $\fst{u(n)} = n$.
    Let $F = \{\snd{u(n)} \mid n\in\naturalsPlus\}$.
    Let $U = \unions{F}$.
    Assume $U$ is a function.
    Assume $r\in\unitRationalsZero$.
    Then $U(r)\subseteq X$.
\end{lemma}
\begin{proof}
    We have $r\in\unitRationals$ or $r\in\{0\}$ by \cref{unit_rationals_zero,union_iff}.
    \begin{byCase}
    \caseOf{$r\in\unitRationals$.}
        Take $n,t$ such that
            $t\in\urysohnAdmissible{s}{X}{T}{A}{B}{n}$ and
            $r\in\dom{t}$ and
            $U(r) = t(r)$
            by \cref{urysohn_union_rational_value_aux}.
        We have $t\in\funs{\{0,1\}\union \img{s}{\initialSegment{n}}}{\pow{X}}$ and
            $\dom{t} = \{0,1\}\union \img{s}{\initialSegment{n}}$
            by \cref{urysohn_admissible,funs_elim}.
        $r\in\{0,1\}\union \img{s}{\initialSegment{n}}$ by \cref{subseteq}.
        $U(r)\subseteq X$ by \cref{subseteq,funs_type_apply,powerset_elim}.
    \caseOf{$r\in\{0\}$.}
        $r = 0$ by \cref{singleton_elim}.
        $0\in\reals$ and $1\in\reals$ and $0 < 1$
            by \cref{real_add_ident,real_mul_ident,real_zero_lt_one}.
        Take $q$ such that $q\in\unitRationals$ and $q\in\intervalopen{0}{1}$
            by \cref{unit_rationals_dense_nonneg}.
        Take $k$ such that $k\in\naturalsPlus$ and $s(k) = q$.
        $u(k)\in G$ by \cref{sequence_in,funs_type_apply}.
        Take $n,t$ such that
            $t\in\urysohnAdmissible{s}{X}{T}{A}{B}{n}$ and
            $u(k) = (n,t)$
            by \cref{urysohn_good_pair_decomp}.
        $\fst{u(k)} = n$ and $\snd{u(k)} = t$ by \cref{fst_eq,snd_eq}.
        $k = n$.
        $t\in\funs{\{0,1\}\union \img{s}{\initialSegment{k}}}{\pow{X}}$
            by \cref{urysohn_admissible}.
        $t$ is a function and $\dom{t} = \{0,1\}\union \img{s}{\initialSegment{k}}$
            by \cref{funs_is_function,funs_elim}.
        $0\in\dom{t}$ by \cref{upair_intro_left,union_intro_left,subseteq}.
        $U(r)\subseteq X$
            by \cref{function_apply_elim,unions_intro,function_apply_intro,funs_type_apply,powerset_elim}.
    \end{byCase}
\end{proof}

% from §33 helper lemma 18d: closures of lower levels lie in higher rational levels
\begin{lemma}\label{urysohn_union_closure_subset}
    Assume $A$ is closed in $X$ with respect to $T$.
    Assume $B$ is closed in $X$ with respect to $T$.
    Assume $A$ is disjoint from $B$.
    Assume $s$ is a sequence in $\unitRationals$.
    Assume for all $p$ such that $p\in\unitRationals$ there exists $n\in\naturalsPlus$ such that $s(n) = p$.
    Let $G = \urysohnGoodPairs{s}{X}{T}{A}{B}$.
    Let $R = \{(z,w) \mid z\in G, w\in G \mid
        \fst{w} = \fst{z} + 1 \land \snd{z}\subseteq\snd{w}\}$.
    Assume $c$ is a function on $G$.
    Assume for all $z\in G$ we have $z\mathrel{R} \apply{c}{z}$.
    Assume $u$ is a sequence in $G$.
    Assume for all $n\in\naturalsPlus$ we have $u(n + 1) = \apply{c}{u(n)}$.
    Assume for all $n\in\naturalsPlus$ we have $\fst{u(n)} = n$.
    Let $F = \{\snd{u(n)} \mid n\in\naturalsPlus\}$.
    Let $U = \unions{F}$.
    Assume $U$ is a function.
    Assume $p\in\unitRationalsZero$.
    Assume $q\in\unitRationals$.
    Assume $p < q$.
    Then $\closure{U(p)}{X}{T}\subseteq U(q)$.
\end{lemma}
\begin{proof}
    Take $n_q,t_q$ such that
        $t_q\in\urysohnAdmissible{s}{X}{T}{A}{B}{n_q}$ and
        $t_q$ is a function and
        $q\in\dom{t_q}$ and
        $t_q\in F$ and
        $U(q) = t_q(q)$
        by \cref{urysohn_union_rational_value_aux}.
    $(q,t_q(q))\in t_q$ by \cref{function_apply_elim}.
    We have $p\in\unitRationals$ or $p\in\{0\}$ by \cref{unit_rationals_zero,union_iff}.
    \begin{byCase}
    \caseOf{$p\in\unitRationals$.}
        Take $n_p,t_p$ such that
            $t_p\in\urysohnAdmissible{s}{X}{T}{A}{B}{n_p}$ and
            $t_p$ is a function and
            $p\in\dom{t_p}$ and
            $t_p\in F$ and
            $U(p) = t_p(p)$
            by \cref{urysohn_union_rational_value_aux}.
        $(p,t_p(p))\in t_p$ by \cref{function_apply_elim}.
        We have $t_p\subseteq t_q$ or $t_q\subseteq t_p$
            by \cref{urysohn_unit_interval_chain}.
        \begin{byCase}
        \caseOf{$t_p\subseteq t_q$.}
            $(p,t_p(p))\in t_q$ by \cref{subseteq}.
            Hence $\closure{t_q(p)}{X}{T}\subseteq t_q(q)$ by \cref{dom_intro,urysohn_admissible}.
            We have $U(p) = t_q(p)$ by \cref{function_apply_elim,function_apply_intro}.
            Therefore $\closure{U(p)}{X}{T}\subseteq U(q)$.
        \caseOf{$t_q\subseteq t_p$.}
            $(q,t_q(q))\in t_p$ by \cref{subseteq}.
            Hence $\closure{t_p(p)}{X}{T}\subseteq t_p(q)$ by \cref{dom_intro,urysohn_admissible}.
            We have $U(q) = t_p(q)$ by \cref{function_apply_elim,function_apply_intro}.
            Therefore $\closure{U(p)}{X}{T}\subseteq U(q)$.
        \end{byCase}
    \caseOf{$p\in\{0\}$.}
        $p = 0$ by \cref{singleton_elim}.
        Show that $\closure{U(p)}{X}{T}\subseteq U(q)$.
        \begin{subproof}
            Follows by \cref{urysohn_admissible,funs_elim,upair_intro_left,union_intro_left,function_apply_elim,unions_intro,function_apply_intro,subseteq}.
        \end{subproof}
    \end{byCase}
\end{proof}

% from §33 helper lemma 20c: value of the Urysohn infimum function
\begin{lemma}\label{urysohn_infimum_value}
    Let $f = \{(x,y) \mid x\in X, y\in\reals \mid
        \text{$y$ is an infimum of $\urysohnLevelSet{U}{x}$}\}$.
    Assume $f\in\funs{X}{\reals}$.
    Assume $x\in X$.
    Then $f(x)$ is an infimum of $\urysohnLevelSet{U}{x}$.
\end{lemma}
\begin{proof}
    Follows by \cref{funs_is_function,funs_elim,subseteq,function_apply_elim,pair_eq_iff}.
\end{proof}

% from §33 helper lemma 19a: infimum below a rational level implies membership
\begin{lemma}\label{urysohn_infimum_lt_rational_in_union}
    Assume $X$ is normal with respect to $T$.
    Assume $A$ is closed in $X$ with respect to $T$.
    Assume $B$ is closed in $X$ with respect to $T$.
    Assume $A$ is disjoint from $B$.
    Assume $s$ is a sequence in $\unitRationals$.
    Assume for all $p$ such that $p\in\unitRationals$ there exists $n\in\naturalsPlus$ such that $s(n) = p$.
    Let $G = \urysohnGoodPairs{s}{X}{T}{A}{B}$.
    Let $R = \{(z,w) \mid z\in G, w\in G \mid
        \fst{w} = \fst{z} + 1 \land \snd{z}\subseteq\snd{w}\}$.
    Assume $c$ is a function on $G$.
    Assume for all $z\in G$ we have $z\mathrel{R} \apply{c}{z}$.
    Assume $u$ is a sequence in $G$.
    Assume $u(1) = a_0$.
    Assume $a_0 = (1,u_1)$.
    Assume for all $n\in\naturalsPlus$ we have $u(n + 1) = \apply{c}{u(n)}$.
    Assume for all $n\in\naturalsPlus$ we have $\fst{u(n)} = n$.
    Let $F = \{\snd{u(n)} \mid n\in\naturalsPlus\}$.
    Let $U = \unions{F}$.
    Assume $U$ is a function.
    Let $Q_x = \urysohnLevelSet{U}{x}$.
    Assume $m$ is an infimum of $Q_x$.
    Assume $q\in\unitRationals$.
    Assume $q\in\reals$.
    Assume $m < q$.
    Then $x\in U(q)$.
\end{lemma}
\begin{proof}
    $m$ is a lower bound of $Q_x$ and $m \leq q$ by \cref{real_infimum,real_lower_bound}.
    Show that there exists $r$ such that $r\in Q_x$ and $r < q$.
    \begin{subproof}
        Suppose not.
        For all $r\in Q_x$ we have $q \leq r$.
        $q$ is a lower bound of $Q_x$ by \cref{real_lower_bound}.
        We have $m\in\reals$ by \cref{real_infimum,real_lower_bound}.
        Hence $q \leq m$ by \cref{real_infimum}.
        Therefore $q < m$ or $q = m$ by \cref{real_leq_split}.
        \begin{byCase}
        \caseOf{$q < m$.}
            $q < q$ by \cref{real_lt_trans}.
            Contradiction by \cref{real_lt_irreflexive}.
        \caseOf{$q = m$.}
            $q < q$ by \cref{real_lt_subst_left}.
            Contradiction by \cref{real_lt_irreflexive}.
        \end{byCase}
    \end{subproof}
    Take $r$ such that $r\in Q_x$ and $r < q$.
    $x\in U(r)$ and $r\in\unitRationalsZero$ by \cref{urysohn_level_set}.
    Hence $\closure{U(r)}{X}{T}\subseteq U(q)$ by \cref{urysohn_union_closure_subset}.
    Therefore $x\in U(q)$ by \cref{urysohn_union_zero_subset_X_aux,subseteq_closure,normal_space,subseteq}.
\end{proof}


% from §33 helper lemma 19: continuity at zero
\begin{lemma}\label{urysohn_unit_interval_continuous_zero}
    Assume $X$ is normal with respect to $T$ and $A$ is closed in $X$ with respect to $T$ and
        $B$ is closed in $X$ with respect to $T$ and $A$ is disjoint from $B$.
    Assume $s$ is a sequence in $\unitRationals$ and for all $p$ such that $p\in\unitRationals$
        there exists $n\in\naturalsPlus$ such that $s(n) = p$.
    Let $G = \urysohnGoodPairs{s}{X}{T}{A}{B}$.
    Let $R = \{(z,w) \mid z\in G, w\in G \mid
        \fst{w} = \fst{z} + 1 \land \snd{z}\subseteq\snd{w}\}$.
    Assume $c$ is a function on $G$ and for all $z\in G$ we have $z\mathrel{R} \apply{c}{z}$.
    Assume $u$ is a sequence in $G$ and $u(1) = a_0$ and $a_0 = (1,u_1)$ and
        for all $n\in\naturalsPlus$ we have $u(n + 1) = \apply{c}{u(n)}$ and $\fst{u(n)} = n$.
    Let $F = \{\snd{u(n)} \mid n\in\naturalsPlus\}$.
    Let $U = \unions{F}$.
    Assume $U$ is a function.
    Let $I_0 = \intervalclosed{0}{1}$.
    Let $f = \{(x,y) \mid x\in X, y\in\reals \mid
        \text{$y$ is an infimum of $\urysohnLevelSet{U}{x}$}\}$.
    Assume $f\in\funs{X}{\reals}$.
    Assume $\img{f}{X}\subseteq I_0$.
    Assume $x\in X$.
    Assume $V$ is open in $\reals$ with respect to $\standardTopologyR$.
    Assume $a\in\reals$ and $b\in\reals$ and $a < b$.
    Assume $f(x)\in\intervalopen{a}{b}$.
    Assume $\intervalopen{a}{b}\subseteq V$.
    Assume $f(x) = 0$.
    Then there exists $U$ such that
        $U$ is open in $X$ with respect to $T$ and $x\in U$ and $\img{f}{U}\subseteq V$.
\end{lemma}
\begin{proof}
    Let $Q_x = \urysohnLevelSet{U}{x}$.
    $f(x) < b$ by \cref{intervalopen}.
    Hence $0 < b$ by \cref{real_lt_subst_left}.
    $0\in\reals$ by \cref{real_add_ident}.
    Take $q$ such that $q\in\unitRationals$ and $q\in\intervalopen{0}{b}$
        by \cref{unit_rationals_dense_nonneg}.
    $q\in\reals$ and $0 < q$ and $q < b$ by \cref{intervalopen}.
    Therefore $q\in\unitRationalsZero$ by \cref{unit_rationals_zero,union_intro_left}.
    Let $W = U(q)$.
    $W$ is open in $X$ with respect to $T$
        by \cref{urysohn_union_rational_open}.
    $f(x)$ is an infimum of $Q_x$ by \cref{urysohn_infimum_value}.
    Hence $x\in U(q)$ by \cref{urysohn_infimum_lt_rational_in_union,real_lt_subst_right}.
    $x\in W$.
    Show that for all $y\in W$ we have $f(y)\in V$.
    \begin{subproof}
            Fix $y\in W$.
            Let $Q_y = \urysohnLevelSet{U}{y}$.
            $q\in Q_y$ by \cref{urysohn_level_set}.
            We have $y\in X$ and $f(y)$ is an infimum of $Q_y$ by
                \cref{open_in,topology_on,subseteq,powerset_elim,urysohn_infimum_value}.
            Hence $f(y) \leq q$ by \cref{real_infimum,real_lower_bound}.
            $f(y)\in\reals$ and $0 \leq f(y)$
                by \cref{funs_is_function,funs_elim,function_apply_elim,subseteq,img_iff,intervalclosed}.
            $a < 0$ by \cref{intervalopen,real_lt_subst_right}.
            Hence $a < f(y)$ by \cref{real_lt_subst_right,real_leq_split,real_lt_trans}.
            Therefore $f(y) < b$ by \cref{real_lt_trans}.
            Hence $f(y)\in\intervalopen{a}{b}$ by \cref{intervalopen}.
            Therefore $f(y)\in V$ by \cref{subseteq}.
    \end{subproof}
    Hence $\img{f}{W}\subseteq V$
        by \cref{img_iff,funs_is_function,function_apply_intro,subseteq}.
    Therefore there exists $W_0$ such that
        $W_0$ is open in $X$ with respect to $T$ and $x\in W_0$ and $\img{f}{W_0}\subseteq V$.
\end{proof}

% from §33 helper lemma 20: positive case produces open neighborhood
\begin{lemma}\label{urysohn_unit_interval_continuous_pos_open}
    Assume $X$ is normal with respect to $T$.
    Assume $A$ is closed in $X$ with respect to $T$.
    Assume $B$ is closed in $X$ with respect to $T$.
    Assume $A$ is disjoint from $B$.
    Assume $s$ is a sequence in $\unitRationals$.
    Assume for all $p$ such that $p\in\unitRationals$ there exists $n\in\naturalsPlus$ such that $s(n) = p$.
    Let $G = \urysohnGoodPairs{s}{X}{T}{A}{B}$.
    Let $R = \{(z,w) \mid z\in G, w\in G \mid
        \fst{w} = \fst{z} + 1 \land \snd{z}\subseteq\snd{w}\}$.
    Assume $c$ is a function on $G$.
    Assume for all $z\in G$ we have $z\mathrel{R} \apply{c}{z}$.
    Assume $u$ is a sequence in $G$.
    Assume $u(1) = a_0$.
    Assume $a_0 = (1,u_1)$.
    Assume for all $n\in\naturalsPlus$ we have $u(n + 1) = \apply{c}{u(n)}$.
    Assume for all $n\in\naturalsPlus$ we have $\fst{u(n)} = n$.
    Let $F = \{\snd{u(n)} \mid n\in\naturalsPlus\}$.
    Let $U = \unions{F}$.
    Assume $U$ is a function.
    Let $I_0 = \intervalclosed{0}{1}$.
    Let $f = \{(x,y) \mid x\in X, y\in\reals \mid
        \text{$y$ is an infimum of $\urysohnLevelSet{U}{x}$}\}$.
    Assume $f\in\funs{X}{\reals}$.
    Assume $\img{f}{X}\subseteq I_0$.
    Assume $x\in X$.
    Assume $V$ is open in $\reals$ with respect to $\standardTopologyR$.
    Assume $a\in\reals$ and $b\in\reals$ and $a < b$.
    Assume $f(x)\in\intervalopen{a}{b}$.
    Assume $\intervalopen{a}{b}\subseteq V$.
    Assume $0 < f(x)$.
    Then there exist $p,q,r_0,W$ such that
        $p\in\unitRationals$ and
        $p\in\intervalopen{a}{f(x)}$ and
        $p\in\intervalopen{0}{f(x)}$ and
        $q\in\unitRationals$ and
        $q\in\intervalopen{f(x)}{b}$ and
        $r_0\in\unitRationals$ and
        $r_0\in\intervalopen{p}{f(x)}$ and
        $p\in\unitRationalsZero$ and
        $q\in\unitRationalsZero$ and
        $r_0\in\unitRationalsZero$ and
        $W = U(q)\setminus \closure{U(p)}{X}{T}$ and
        $W$ is open in $X$ with respect to $T$.
\end{lemma}
\begin{proof}
    Show that there exists $p$ such that
        $p\in\unitRationals$ and
        $p\in\intervalopen{a}{f(x)}$ and
        $p\in\intervalopen{0}{f(x)}$.
    \begin{subproof}
        $0\in\reals$ by \cref{real_add_ident}.
        $f(x)\in\reals$ by \cref{funs_type_apply}.
        Hence $a < 0$ or $a = 0$ or $0 < a$ by \cref{real_lt_trichotomy}.
        \begin{byCase}
        \caseOf{$a < 0$.}
            Take $p_0$ such that $p_0\in\unitRationals$ and $p_0\in\intervalopen{0}{f(x)}$
                by \cref{unit_rationals_dense_nonneg}.
            $p_0\in\reals$ and $0 < p_0$ and $p_0 < f(x)$ by \cref{intervalopen}.
            Hence $p_0\in\intervalopen{a}{f(x)}$ by \cref{intervalopen,real_lt_trans}.
            Therefore there exists $p$ such that
                $p\in\unitRationals$ and
                $p\in\intervalopen{a}{f(x)}$ and
                $p\in\intervalopen{0}{f(x)}$.
        \caseOf{$a = 0$.}
            Take $p_0$ such that $p_0\in\unitRationals$ and $p_0\in\intervalopen{0}{f(x)}$
                by \cref{unit_rationals_dense_nonneg}.
            $p_0\in\reals$ and $0 < p_0$ and $p_0 < f(x)$ by \cref{intervalopen}.
            Hence $p_0\in\intervalopen{a}{f(x)}$ by \cref{intervalopen,real_lt_subst_left}.
            Therefore there exists $p$ such that
                $p\in\unitRationals$ and
                $p\in\intervalopen{a}{f(x)}$ and
                $p\in\intervalopen{0}{f(x)}$.
        \caseOf{$0 < a$.}
            $a < f(x)$ by \cref{intervalopen}.
            Take $p_0$ such that $p_0\in\unitRationals$ and $p_0\in\intervalopen{a}{f(x)}$
                by \cref{unit_rationals_dense_nonneg}.
            $p_0\in\reals$ and $a < p_0$ and $p_0 < f(x)$ by \cref{intervalopen}.
            Hence $p_0\in\intervalopen{0}{f(x)}$ by \cref{intervalopen,real_lt_trans}.
            Therefore there exists $p$ such that
                $p\in\unitRationals$ and
                $p\in\intervalopen{a}{f(x)}$ and
                $p\in\intervalopen{0}{f(x)}$.
        \end{byCase}
    \end{subproof}
    Take $p$ such that
        $p\in\unitRationals$ and
        $p\in\intervalopen{a}{f(x)}$ and
        $p\in\intervalopen{0}{f(x)}$.
    $f(x)\in\reals$ and $f(x) < b$ by \cref{funs_type_apply,intervalopen}.
    Take $q$ such that $q\in\unitRationals$ and $q\in\intervalopen{f(x)}{b}$
        by \cref{unit_rationals_dense_nonneg}.
    $p\in\reals$ and $a < p$ and $0 < p$ and $p < f(x)$ by \cref{intervalopen}.
    Take $r_0$ such that $r_0\in\unitRationals$ and $r_0\in\intervalopen{p}{f(x)}$
        by \cref{unit_rationals_dense_nonneg}.
    $p < r_0$ and $r_0 < f(x)$ by \cref{intervalopen}.
    Therefore $p\in\unitRationalsZero$ and $q\in\unitRationalsZero$ and $r_0\in\unitRationalsZero$
        by \cref{unit_rationals_zero,union_intro_left}.
    Let $W = U(q)\setminus \closure{U(p)}{X}{T}$.
    Show that $W$ is open in $X$ with respect to $T$.
    \begin{subproof}
        $U(q)\subseteq X$ and $U(q)\in T$
            by \cref{urysohn_union_rational_open,normal_space,open_in,topology_on,subseteq,powerset_elim}.
        $U(p)\subseteq X$ by \cref{urysohn_union_zero_subset_X_aux}.
        $T$ is a topology on $X$ by \cref{normal_space}.
        We have $\closure{U(p)}{X}{T}$ is closed in $X$ with respect to $T$
            by \cref{closure_closed_in}.
        Therefore $X\setminus \closure{U(p)}{X}{T}\in T$ by \cref{closed_in,open_in}.
        $U(q)\inter (X\setminus \closure{U(p)}{X}{T})\in T$ by \cref{topology_on}.
        $\closure{U(p)}{X}{T}\subseteq X$ by \cref{closed_in}.
        We have $W = U(q)\inter (X\setminus \closure{U(p)}{X}{T})$
            by \cref{setminus_eq_inter_complement}.
        Hence $W\in T$.
        Thus $W$ is open in $X$ with respect to $T$ by \cref{open_in}.
    \end{subproof}
    Therefore there exist $p_1,q_1,r_1,W_1$ such that
        $p_1\in\unitRationals$ and
        $p_1\in\intervalopen{a}{f(x)}$ and
        $p_1\in\intervalopen{0}{f(x)}$ and
        $q_1\in\unitRationals$ and
        $q_1\in\intervalopen{f(x)}{b}$ and
        $r_1\in\unitRationals$ and
        $r_1\in\intervalopen{p_1}{f(x)}$ and
        $p_1\in\unitRationalsZero$ and
        $q_1\in\unitRationalsZero$ and
        $r_1\in\unitRationalsZero$ and
        $W_1 = U(q_1)\setminus \closure{U(p_1)}{X}{T}$ and
        $W_1$ is open in $X$ with respect to $T$.
\end{proof}

% from §33 helper lemma 20e: rational U-values lie in X
\begin{lemma}\label{urysohn_union_rational_subset_X}
    Assume $s$ is a sequence in $\unitRationals$.
    Assume for all $p$ such that $p\in\unitRationals$ there exists $n\in\naturalsPlus$ such that $s(n) = p$.
    Let $G = \urysohnGoodPairs{s}{X}{T}{A}{B}$.
    Assume $u$ is a sequence in $G$.
    Assume for all $n\in\naturalsPlus$ we have $\fst{u(n)} = n$.
    Let $F = \{\snd{u(n)} \mid n\in\naturalsPlus\}$.
    Let $U = \unions{F}$.
    Assume $U$ is a function.
    Assume $q\in\unitRationals$.
    Then $U(q)\subseteq X$.
\end{lemma}
\begin{proof}
    Follows by \cref{unit_rationals_zero,union_intro_left,urysohn_union_zero_subset_X_aux}.
\end{proof}
% from §33 helper lemma 21: positive case point membership
\begin{lemma}\label{urysohn_unit_interval_continuous_pos_x_in_W}
    Assume $X$ is normal with respect to $T$.
    Assume $A$ is closed in $X$ with respect to $T$.
    Assume $B$ is closed in $X$ with respect to $T$.
    Assume $A$ is disjoint from $B$.
    Assume $s$ is a sequence in $\unitRationals$.
    Assume for all $p$ such that $p\in\unitRationals$ there exists $n\in\naturalsPlus$ such that $s(n) = p$.
    Let $G = \urysohnGoodPairs{s}{X}{T}{A}{B}$.
    Let $R = \{(z,w) \mid z\in G, w\in G \mid
        \fst{w} = \fst{z} + 1 \land \snd{z}\subseteq\snd{w}\}$.
    Assume $c$ is a function on $G$.
    Assume for all $z\in G$ we have $z\mathrel{R} \apply{c}{z}$.
    Assume $u$ is a sequence in $G$.
    Assume $u(1) = a_0$.
    Assume $a_0 = (1,u_1)$.
    Assume for all $n\in\naturalsPlus$ we have $u(n + 1) = \apply{c}{u(n)}$.
    Assume for all $n\in\naturalsPlus$ we have $\fst{u(n)} = n$.
    Let $F = \{\snd{u(n)} \mid n\in\naturalsPlus\}$.
    Let $U = \unions{F}$.
    Assume $U$ is a function.
    Let $I_0 = \intervalclosed{0}{1}$.
    Let $f = \{(x,y) \mid x\in X, y\in\reals \mid
        \text{$y$ is an infimum of $\urysohnLevelSet{U}{x}$}\}$.
    Assume $f\in\funs{X}{\reals}$ and $\img{f}{X}\subseteq I_0$.
    Assume $x\in X$ and $V$ is open in $\reals$ with respect to $\standardTopologyR$ and
        $a\in\reals$ and $b\in\reals$ and $a < b$ and
        $f(x)\in\intervalopen{a}{b}$ and $\intervalopen{a}{b}\subseteq V$ and $0 < f(x)$.
    Assume $p\in\unitRationals$ and $p\in\intervalopen{a}{f(x)}$ and $p\in\intervalopen{0}{f(x)}$ and
        $q\in\unitRationals$ and $q\in\intervalopen{f(x)}{b}$ and
        $r_0\in\unitRationals$ and $r_0\in\intervalopen{p}{f(x)}$ and
        $p\in\unitRationalsZero$ and $q\in\unitRationalsZero$ and $r_0\in\unitRationalsZero$.
    Let $W = U(q)\setminus \closure{U(p)}{X}{T}$.
    Then $x\in W$.
\end{lemma}
\begin{proof}
    Let $Q_x = \urysohnLevelSet{U}{x}$.
    $f(x)$ is an infimum of $Q_x$ and $f(x)$ is a lower bound of $Q_x$
        by \cref{urysohn_infimum_value,real_infimum}.
    $f(x) < q$ by \cref{intervalopen}.
    Hence $x\in U(q)$ by \cref{urysohn_infimum_lt_rational_in_union,intervalopen}.
    $p < r_0$ by \cref{intervalopen}.
    Hence $\closure{U(p)}{X}{T}\subseteq U(r_0)$ by \cref{urysohn_union_closure_subset}.
    Show that $x\notin U(r_0)$.
    \begin{subproof}
        Suppose not.
        $x\in U(r_0)$.
        Hence $r_0\in Q_x$ by \cref{urysohn_level_set}.
        $f(x) \leq r_0$ by \cref{real_lower_bound}.
        Hence $f(x)\in\reals$ and $r_0\in\reals$ by \cref{funs_type_apply,intervalopen}.
        We have $f(x) < r_0$ or $f(x) = r_0$ by \cref{real_leq_split}.
        $r_0 < f(x)$ by \cref{intervalopen}.
        Contradiction by \cref{real_lt_trans,real_lt_irreflexive,real_lt_subst_left}.
    \end{subproof}
    Therefore $x\notin\closure{U(p)}{X}{T}$ by \cref{not_in_subseteq}.
    Hence $x\in W$ by \cref{setminus_intro}.
\end{proof}
% from §33 helper lemma 22: positive case image control
\begin{lemma}\label{urysohn_unit_interval_continuous_pos_img}
    Assume $X$ is normal with respect to $T$ and $A$ is closed in $X$ with respect to $T$ and
        $B$ is closed in $X$ with respect to $T$ and $A$ is disjoint from $B$.
    Assume $s$ is a sequence in $\unitRationals$ and for all $p$ such that $p\in\unitRationals$
        there exists $n\in\naturalsPlus$ such that $s(n) = p$.
    Let $G = \urysohnGoodPairs{s}{X}{T}{A}{B}$.
    Let $R = \{(z,w) \mid z\in G, w\in G \mid
        \fst{w} = \fst{z} + 1 \land \snd{z}\subseteq\snd{w}\}$.
    Assume $c$ is a function on $G$ and for all $z\in G$ we have $z\mathrel{R} \apply{c}{z}$.
    Assume $u$ is a sequence in $G$ and $u(1) = a_0$ and $a_0 = (1,u_1)$ and
        for all $n\in\naturalsPlus$ we have $u(n + 1) = \apply{c}{u(n)}$ and $\fst{u(n)} = n$.
    Let $F = \{\snd{u(n)} \mid n\in\naturalsPlus\}$.
    Let $U = \unions{F}$.
    Assume $U$ is a function.
    Let $I_0 = \intervalclosed{0}{1}$.
    Let $f = \{(x,y) \mid x\in X, y\in\reals \mid
        \text{$y$ is an infimum of $\urysohnLevelSet{U}{x}$}\}$.
    Assume $f\in\funs{X}{\reals}$ and $\img{f}{X}\subseteq I_0$.
    Assume $x\in X$ and $V$ is open in $\reals$ with respect to $\standardTopologyR$ and
        $a\in\reals$ and $b\in\reals$ and $a < b$ and
        $f(x)\in\intervalopen{a}{b}$ and $\intervalopen{a}{b}\subseteq V$ and $0 < f(x)$.
    Assume $p\in\unitRationals$ and $p\in\intervalopen{a}{f(x)}$ and $p\in\intervalopen{0}{f(x)}$ and
        $q\in\unitRationals$ and $q\in\intervalopen{f(x)}{b}$ and
        $r_0\in\unitRationals$ and $r_0\in\intervalopen{p}{f(x)}$ and
        $p\in\unitRationalsZero$ and $q\in\unitRationalsZero$ and $r_0\in\unitRationalsZero$.
    Let $W = U(q)\setminus \closure{U(p)}{X}{T}$.
    Assume $W\subseteq X$.
    Then $\img{f}{W}\subseteq V$.
\end{lemma}
\begin{proof}
    Show that for all $y\in W$ we have $f(y)\in V$.
    \begin{subproof}
        Fix $y$ such that $y\in W$.
        $y\in X$ and $f(y)\in\reals$ by \cref{subseteq,funs_type_apply}.
        Let $Q_y = \urysohnLevelSet{U}{y}$.
        $y\notin\closure{U(p)}{X}{T}$ by \cref{setminus_elim_right}.
        $q\in Q_y$ by \cref{setminus_elim_left,urysohn_level_set}.
        Therefore $f(y) \leq q$ by \cref{urysohn_infimum_value,real_infimum,real_lower_bound}.
        Show that $p \leq f(y)$.
        \begin{subproof}
            Show that $p$ is a lower bound of $Q_y$.
            \begin{subproof}
        $p\in\reals$ and $Q_y\subseteq\reals$ by \cref{intervalopen,urysohn_level_value_real,subseteq}.
                Show that for all $r\in Q_y$ we have $p \leq r$.
                \begin{subproof}
                    Fix $r$ such that $r\in Q_y$.
                    $r\in\reals$ by \cref{subseteq}.
                    Hence $y\in U(r)$ and $r\in\unitRationalsZero$ by \cref{urysohn_level_set}.
                    Show that $p \neq r$.
                    \begin{subproof}
                        Suppose not.
                        $p = r$.
                        Hence $y\in U(p)$.
                        Therefore $y\in\closure{U(p)}{X}{T}$
                            by \cref{urysohn_union_zero_subset_X_aux,subseteq_closure,normal_space,subseteq}.
                        Contradiction.
                    \end{subproof}
                    Hence $p < r$ or $r < p$ by \cref{real_lt_trichotomy}.
                    \begin{byCase}
                    \caseOf{$p < r$.}
                        Show that $p \leq r$.
                        \begin{subproof}
                            Follows by \cref{real_lt_imp_leq}.
                        \end{subproof}
                    \caseOf{$r < p$.}
                        Show that $p \leq r$.
                        \begin{subproof}
                            Suppose not.
                            We have $\closure{U(r)}{X}{T}\subseteq U(p)$ by \cref{urysohn_union_closure_subset}.
                            We have $y\in\closure{U(r)}{X}{T}$
                                by \cref{urysohn_union_zero_subset_X_aux,subseteq_closure,normal_space,subseteq}.
                            $y\in U(p)$ by \cref{subseteq}.
                            Therefore $y\in\closure{U(p)}{X}{T}$
                                by \cref{urysohn_union_zero_subset_X_aux,subseteq_closure,normal_space,subseteq}.
                            Contradiction.
                        \end{subproof}
                    \end{byCase}
                \end{subproof}
                Therefore $p$ is a lower bound of $Q_y$ by \cref{real_lower_bound}.
            \end{subproof}
            Hence $p \leq f(y)$ by \cref{urysohn_infimum_value,real_infimum}.
        \end{subproof}
        $a < f(y)$ by \cref{intervalopen,real_leq_split,real_lt_subst_right,real_lt_trans}.
        $f(y) < b$ by \cref{intervalopen,real_leq_split,real_lt_subst_left,real_lt_trans}.
        Therefore $f(y)\in V$ by \cref{intervalopen,subseteq}.
    \end{subproof}
    Hence $\img{f}{W}\subseteq V$ by \cref{img_iff,funs_is_function,function_apply_intro,subseteq}.
\end{proof}

% from §33 helper lemma 23: continuity at a point
\begin{lemma}\label{urysohn_unit_interval_continuous_at_point}
    Assume $X$ is normal with respect to $T$.
    Assume $A$ is closed in $X$ with respect to $T$.
    Assume $B$ is closed in $X$ with respect to $T$.
    Assume $A$ is disjoint from $B$.
    Assume $s$ is a sequence in $\unitRationals$.
    Assume for all $p$ such that $p\in\unitRationals$ there exists $n\in\naturalsPlus$ such that $s(n) = p$.
    Let $G = \urysohnGoodPairs{s}{X}{T}{A}{B}$.
    Let $R = \{(z,w) \mid z\in G, w\in G \mid
        \fst{w} = \fst{z} + 1 \land \snd{z}\subseteq\snd{w}\}$.
    Assume $c$ is a function on $G$.
    Assume for all $z\in G$ we have $z\mathrel{R} \apply{c}{z}$.
    Assume $u$ is a sequence in $G$.
    Assume $u(1) = a_0$.
    Assume $a_0 = (1,u_1)$.
    Assume for all $n\in\naturalsPlus$ we have $u(n + 1) = \apply{c}{u(n)}$.
    Assume for all $n\in\naturalsPlus$ we have $\fst{u(n)} = n$.
    Let $F = \{\snd{u(n)} \mid n\in\naturalsPlus\}$.
    Let $U = \unions{F}$.
    Assume $U$ is a function.
    Let $I_0 = \intervalclosed{0}{1}$.
    Let $f = \{(x,y) \mid x\in X, y\in\reals \mid
        \text{$y$ is an infimum of $\urysohnLevelSet{U}{x}$}\}$.
    Assume $f\in\funs{X}{\reals}$.
    Assume $\img{f}{X}\subseteq I_0$.
    Assume $x\in X$.
    Assume $V$ is open in $\reals$ with respect to $\standardTopologyR$.
    Assume $f(x)\in V$.
    Then there exists $U_0$ such that
        $U_0$ is open in $X$ with respect to $T$ and $x\in U_0$ and $\img{f}{U_0}\subseteq V$.
\end{lemma}
\begin{proof}
    $V\in\basisTopology{\standardBasisR}{\reals}$ by \cref{open_in,standard_topology_reals}.
    Take $B_1$ such that $B_1\in\standardBasisR$ and $f(x)\in B_1$ and $B_1\subseteq V$
        by \cref{topology_generated_by_basis}.
    Take $a,b\in\reals$ such that $a < b$ and $B_1 = \intervalopen{a}{b}$
        by \cref{standard_basis_reals}.
    $a < f(x)$ and $f(x) < b$ by \cref{intervalopen}.
    Hence $0 < f(x)$ or $0 = f(x)$
        by \cref{funs_elim,funs_is_function,subseteq,function_apply_elim,img_iff,funs_type_apply,intervalclosed,real_add_ident,real_leq_split}.
    \begin{byCase}
    \caseOf{$0 = f(x)$.}
        Show that there exists $U_0$ such that
            $U_0$ is open in $X$ with respect to $T$ and $x\in U_0$ and $\img{f}{U_0}\subseteq V$
        .
        \begin{subproof}
            Follows by \cref{urysohn_unit_interval_continuous_zero}.
        \end{subproof}
    \caseOf{$0 < f(x)$.}
        Show that there exists $U_0$ such that
            $U_0$ is open in $X$ with respect to $T$ and $x\in U_0$ and $\img{f}{U_0}\subseteq V$.
        \begin{subproof}
            Take $p,q,r_0,W$ such that
                $p\in\unitRationals$ and
                $p\in\intervalopen{a}{f(x)}$ and
                $p\in\intervalopen{0}{f(x)}$ and
                $q\in\unitRationals$ and
                $q\in\intervalopen{f(x)}{b}$ and
                $r_0\in\unitRationals$ and
                $r_0\in\intervalopen{p}{f(x)}$ and
                $p\in\unitRationalsZero$ and
                $q\in\unitRationalsZero$ and
                $r_0\in\unitRationalsZero$ and
                $W = U(q)\setminus \closure{U(p)}{X}{T}$ and
                $W$ is open in $X$ with respect to $T$
                by \cref{urysohn_unit_interval_continuous_pos_open}.
            $x\in W$ by \cref{urysohn_unit_interval_continuous_pos_x_in_W}.
            We have $W\subseteq X$ by \cref{open_in,normal_space,topology_on,subseteq,powerset_elim}.
            We have $\img{f}{W}\subseteq V$ by \cref{urysohn_unit_interval_continuous_pos_img}.
            Therefore there exists $U_0$ such that
                $U_0$ is open in $X$ with respect to $T$ and $x\in U_0$ and $\img{f}{U_0}\subseteq V$.
        \end{subproof}
    \end{byCase}
\end{proof}

% from §33 helper lemma 24: Urysohn function is continuous to reals
\begin{lemma}\label{urysohn_unit_interval_continuous_reals}
    Assume $X$ is normal with respect to $T$.
    Assume $A$ is closed in $X$ with respect to $T$.
    Assume $B$ is closed in $X$ with respect to $T$.
    Assume $A$ is disjoint from $B$.
    Assume $s$ is a sequence in $\unitRationals$.
    Assume for all $p$ such that $p\in\unitRationals$ there exists $n\in\naturalsPlus$ such that $s(n) = p$.
    Let $G = \urysohnGoodPairs{s}{X}{T}{A}{B}$.
    Let $R = \{(z,w) \mid z\in G, w\in G \mid
        \fst{w} = \fst{z} + 1 \land \snd{z}\subseteq\snd{w}\}$.
    Assume $c$ is a function on $G$.
    Assume for all $z\in G$ we have $z\mathrel{R} \apply{c}{z}$.
    Assume $u$ is a sequence in $G$.
    Assume $u(1) = a_0$.
    Assume $a_0 = (1,u_1)$.
    Assume for all $n\in\naturalsPlus$ we have $u(n + 1) = \apply{c}{u(n)}$.
    Assume for all $n\in\naturalsPlus$ we have $\fst{u(n)} = n$.
    Let $F = \{\snd{u(n)} \mid n\in\naturalsPlus\}$.
    Let $U = \unions{F}$.
    Assume $U$ is a function.
    Let $I_0 = \intervalclosed{0}{1}$.
    Let $f = \{(x,y) \mid x\in X, y\in\reals \mid
        \text{$y$ is an infimum of $\urysohnLevelSet{U}{x}$}\}$.
    Assume $f\in\funs{X}{\reals}$.
    Assume $\img{f}{X}\subseteq I_0$.
    Then $f$ is continuous from $X$ to $\reals$ with respect to $T$ and $\standardTopologyR$.
\end{lemma}
\begin{proof}
    Follows by \cref{urysohn_unit_interval_continuous_at_point,standard_basis_is_basis,basis_topology_is_topology,standard_topology_reals,normal_space,continuous_at_open,continuous_at_implies_continuous}.
\end{proof}

% from §33 helper lemma 34: admissible values above $1$
\begin{lemma}\label{urysohn_admissible_above_one}
    Assume $f\in\urysohnAdmissible{s}{X}{T}{A}{B}{n}$.
    Assume $p\in\dom{f}$.
    Assume $1 < p$.
    Then $f(p) = X$.
\end{lemma}
\begin{proof}
    Follows by \cref{urysohn_admissible,funs_elim,upair_intro_right,union_intro_left}.
\end{proof}

% from §33 helper lemma 35: unit rationals are real
\begin{lemma}\label{unit_rational_real}
    Assume $r\in\unitRationals$.
    Then $r\in\reals$.
\end{lemma}
\begin{proof}
    Follows by \cref{unit_rational_real_aux}.
\end{proof}

% from §33 helper lemma 25: extension when $r\in\dom{f}$
\begin{lemma}\label{urysohn_unit_interval_extend_in_dom}
    Assume $X$ is normal with respect to $T$.
    Assume $A$ is closed in $X$ with respect to $T$.
    Assume $B$ is closed in $X$ with respect to $T$.
    Assume $A$ is disjoint from $B$.
    Assume $s$ is a sequence in $\unitRationals$.
    Let $G = \urysohnGoodPairs{s}{X}{T}{A}{B}$.
    Let $R = \{(z,w) \mid z\in G, w\in G \mid
        \fst{w} = \fst{z} + 1 \land \snd{z}\subseteq\snd{w}\}$.
    Suppose $n\in\naturalsPlus$ and $f\in\urysohnAdmissible{s}{X}{T}{A}{B}{n}$ and $z = (n,f)$.
    Let $r = s(n + 1)$.
    Assume $r\in\dom{f}$.
    Then there exists $w$ such that $z\mathrel{R} w$.
\end{lemma}
\begin{proof}
    $f\in\funs{\{0,1\}\union \img{s}{\initialSegment{n}}}{\pow{X}}$ and
        for all $p\in\dom{f}$ we have $f(p)\in T$ and
        $f(1) = X\setminus B$ and
        $A\subseteq f(0)$ and
        $\closure{f(0)}{X}{T}\subseteq f(1)$ and
        for all $p\in\dom{f}$ for all $q\in\dom{f}$ if $p < q$ then
            $\closure{f(p)}{X}{T}\subseteq f(q)$ and
        for all $p\in\dom{f}$ if $1 < p$ then $f(p) = X$
        by \cref{urysohn_admissible}.
    $\{0,1\}\union \img{s}{\initialSegment{n}}\subseteq \{0,1\}\union \img{s}{\initialSegment{n + 1}}$
        by \cref{initial_segment_subset_succ,img_subseteq,union_intro_left,union_intro_right,subseteq,union_subsets_is_subset}.
    Show that for all $y\in \img{s}{\initialSegment{n + 1}}$ we have $y\in\dom{f}$.
    \begin{subproof}
        Fix $y$ such that $y\in \img{s}{\initialSegment{n + 1}}$.
        Take $k$ such that $k\in\initialSegment{n + 1}$ and $(k,y)\in s$ by \cref{img_iff}.
        $s(k) = y$ by \cref{sequence_in,funs_is_function,function_apply_intro}.
        \begin{byCase}
        \caseOf{$k\in\initialSegment{n}$.}
            Show that $y\in\dom{f}$.
            \begin{subproof}
                Follows by \cref{img_iff,union_intro_right,funs_elim}.
            \end{subproof}
        \caseOf{$k\notin\initialSegment{n}$.}
            $k\in\naturalsPlus$ and $k\leq n + 1$ by \cref{initial_segment_naturalsplus}.
            Hence $k = n + 1$ by \cref{real_naturals_leq_succ_elim,initial_segment_naturalsplus}.
            Therefore $y = r$ and $y\in\dom{f}$ by \cref{subseteq}.
        \end{byCase}
    \end{subproof}
    Hence $\img{s}{\initialSegment{n + 1}}\subseteq \dom{f}$ by \cref{subseteq}.
    Therefore $\{0,1\}\union \img{s}{\initialSegment{n + 1}}\subseteq \dom{f}$
        by \cref{funs_elim,union_intro_left,subseteq,union_subsets_is_subset}.
    Show that for all $y$ such that $y\in\dom{f}$ we have $y\in \{0,1\}\union \img{s}{\initialSegment{n + 1}}$.
    \begin{subproof}
        Follows by \cref{funs_elim,subseteq}.
    \end{subproof}
    $\dom{f}\subseteq \{0,1\}\union \img{s}{\initialSegment{n + 1}}$ by \cref{subseteq}.
    Therefore $\dom{f} = \{0,1\}\union \img{s}{\initialSegment{n + 1}}$ by \cref{subseteq_antisymmetric}.
    Let $w = (n + 1,f)$.
    $n + 1\in\naturalsPlus$ and $f\in\urysohnAdmissible{s}{X}{T}{A}{B}{n + 1}$
        by \cref{naturalsplus_add_closed,real_one_in_naturals,funs_elim,funs_intro,urysohn_admissible}.
    Hence $z\in G$ and $w\in G$ and $\fst{w} = \fst{z} + 1$ and $\snd{z}\subseteq \snd{w}$
        by \cref{urysohn_good_pair_intro,fst_eq,snd_eq,subseteq_refl}.
    Therefore $z\mathrel{R} w$.
\end{proof}

% from §33 helper lemma 26: extension for $1 < r$ (setup)
\begin{lemma}\label{urysohn_unit_interval_extend_gt1_setup}
    Assume $X$ is normal with respect to $T$.
    Assume $A$ is closed in $X$ with respect to $T$.
    Assume $B$ is closed in $X$ with respect to $T$.
    Assume $A$ is disjoint from $B$.
    Assume $s$ is a sequence in $\unitRationals$.
    Suppose $n\in\naturalsPlus$ and $f\in\urysohnAdmissible{s}{X}{T}{A}{B}{n}$ and $z = (n,f)$.
    Let $r = s(n + 1)$.
    Assume $r\notin\dom{f}$.
    Let $U_r = X$.
    Let $h = \{(r,U_r)\}$.
    Let $g = f\union h$.
    Then $g$ is a function.
\end{lemma}
\begin{proof}
    Show that for all $x$ such that $x\in\dom{f}\inter\dom{h}$ we have $f(x) = h(x)$.
    \begin{subproof}
        Fix $x$ such that $x\in\dom{f}\inter\dom{h}$.
        Suppose not.
        $x\in\dom{h}$ by \cref{inter_elim_right}.
        $x\in\{r\}$ by \cref{dom_iff,singleton_elim,pair_eq_iff,singleton_intro}.
        Thus $x = r$ by \cref{singleton_elim}.
        Therefore $r\in\dom{f}$ by \cref{inter_elim_left}.
        Contradiction.
    \end{subproof}
    Show that $h$ is a function.
    \begin{subproof}
        Follows by \cref{singleton_elim,pair_eq_iff,relation,rightunique}.
    \end{subproof}
    Show that $g$ is a function.
    \begin{subproof}
        Follows by \cref{urysohn_admissible,funs_is_function,union_compatible_functions}.
    \end{subproof}
\end{proof}

% from §33 helper lemma 27: extension for $1 < r$
\begin{lemma}\label{urysohn_unit_interval_extend_gt1}
    Assume $X$ is normal with respect to $T$.
    Assume $A$ is closed in $X$ with respect to $T$.
    Assume $B$ is closed in $X$ with respect to $T$.
    Assume $A$ is disjoint from $B$.
    Assume $s$ is a sequence in $\unitRationals$.
    Let $G = \urysohnGoodPairs{s}{X}{T}{A}{B}$.
    Let $R = \{(z,w) \mid z\in G, w\in G \mid
        \fst{w} = \fst{z} + 1 \land \snd{z}\subseteq\snd{w}\}$.
    Suppose $n\in\naturalsPlus$ and $f\in\urysohnAdmissible{s}{X}{T}{A}{B}{n}$ and $z = (n,f)$.
    Let $r = s(n + 1)$.
    Assume $r\notin\dom{f}$.
    Assume $1 < r$.
    Then there exists $w$ such that $z\mathrel{R} w$.
\end{lemma}
\begin{proof}
    $f\in\funs{\{0,1\}\union \img{s}{\initialSegment{n}}}{\pow{X}}$ and
        for all $p\in\dom{f}$ we have $f(p)\in T$ and
        $f(1) = X\setminus B$ and
        $A\subseteq f(0)$ and
        $\closure{f(0)}{X}{T}\subseteq f(1)$ and
        for all $p\in\dom{f}$ for all $q\in\dom{f}$ if $p < q$ then
            $\closure{f(p)}{X}{T}\subseteq f(q)$ and
        for all $p\in\dom{f}$ if $1 < p$ then $f(p) = X$
        by \cref{urysohn_admissible}.
    Let $U_r = X$.
    Let $h = \{(r,U_r)\}$.
    Let $g = f\union h$.
    $g$ is a function by \cref{urysohn_unit_interval_extend_gt1_setup}.
    $\dom{g} = \dom{f}\union\dom{h}$ by \cref{dom_union}.
    For all $x\in\dom{h}$ we have $x\in\{r\}$ by \cref{dom_iff,singleton_elim,pair_eq_iff,singleton_intro}.
    For all $x\in\{r\}$ we have $x\in\dom{h}$
        by \cref{singleton_elim,singleton_intro,pair_eq_iff,dom_intro}.
    Hence $\dom{h} = \{r\}$ by \cref{subseteq,subseteq_antisymmetric}.
    Therefore $\dom{g} = \dom{f}\union\{r\}$.
    $\dom{f} = \{0,1\}\union \img{s}{\initialSegment{n}}$ by \cref{funs_elim}.
    Show that for all $y\in\dom{g}$ we have $y\in \{0,1\}\union \img{s}{\initialSegment{n + 1}}$.
        \begin{subproof}
            Fix $y$ such that $y\in\dom{g}$.
            $y\in\dom{f}$ or $y\in\{r\}$ by \cref{union_iff}.
            \begin{byCase}
            \caseOf{$y\in\dom{f}$.}
                $y\in \{0,1\}$ or $y\in \img{s}{\initialSegment{n}}$ by \cref{funs_elim,union_iff}.
                \begin{byCase}
                \caseOf{$y\in \{0,1\}$.}
                    Show that $y\in \{0,1\}\union \img{s}{\initialSegment{n + 1}}$.
                    \begin{subproof}
                        Follows by \cref{union_intro_left}.
                    \end{subproof}
                \caseOf{$y\in \img{s}{\initialSegment{n}}$.}
                    Take $k$ such that $k\in\initialSegment{n}$ and $(k,y)\in s$ by \cref{img_iff}.
                    $y\in \{0,1\}\union \img{s}{\initialSegment{n + 1}}$
                        by \cref{initial_segment_subset_succ,subseteq,img_iff,union_intro_right}.
                \end{byCase}
            \caseOf{$y\in\{r\}$.}
                $y\in \{0,1\}\union \img{s}{\initialSegment{n + 1}}$
                    by \cref{singleton_elim,naturalsplus_add_closed,real_one_in_naturals,sequence_in,sequence_value_in_initial_image,union_intro_right}.
            \end{byCase}
    \end{subproof}
    Hence $\dom{g}\subseteq \{0,1\}\union \img{s}{\initialSegment{n + 1}}$ by \cref{subseteq}.
    Show that for all $y\in \{0,1\}\union \img{s}{\initialSegment{n + 1}}$ we have $y\in\dom{g}$.
    \begin{subproof}
            Fix $y$ such that $y\in \{0,1\}\union \img{s}{\initialSegment{n + 1}}$.
            $y\in \{0,1\}$ or $y\in \img{s}{\initialSegment{n + 1}}$ by \cref{union_iff}.
            \begin{byCase}
            \caseOf{$y\in \{0,1\}$.}
                Show that $y\in\dom{g}$.
                \begin{subproof}
                    Follows by \cref{union_intro_left,funs_elim}.
                \end{subproof}
            \caseOf{$y\in \img{s}{\initialSegment{n + 1}}$.}
                Take $k$ such that $k\in\initialSegment{n + 1}$ and $(k,y)\in s$ by \cref{img_iff}.
                \begin{byCase}
                \caseOf{$k\in\initialSegment{n}$.}
                    $y\in\dom{g}$ by \cref{img_iff,union_intro_right,funs_elim,union_intro_left}.
                \caseOf{$k\notin\initialSegment{n}$.}
                    $k\in\naturalsPlus$ and $k\leq n + 1$ by \cref{initial_segment_naturalsplus}.
                    Hence $k = n + 1$ by \cref{real_naturals_leq_succ_elim,initial_segment_naturalsplus}.
                    Thus $y = r$.
                    Therefore $y\in\dom{g}$ by \cref{singleton_intro,union_intro_right}.
                \end{byCase}
            \end{byCase}
    \end{subproof}
    Thus $\{0,1\}\union \img{s}{\initialSegment{n + 1}}\subseteq \dom{g}$ by \cref{subseteq}.
    Therefore $\dom{g} = \{0,1\}\union \img{s}{\initialSegment{n + 1}}$ by \cref{subseteq_antisymmetric}.
    Show that for all $p\in\dom{g}$ we have $g(p)\in\pow{X}$.
    \begin{subproof}
        Fix $p$ such that $p\in\dom{g}$.
        $p\in\dom{f}$ or $p\in\{r\}$ by \cref{union_iff}.
        \begin{byCase}
        \caseOf{$p\in\dom{f}$.}
            $g(p)\in\pow{X}$
                by \cref{urysohn_admissible,funs_type_apply,funs_is_function,function_apply_elim,union_intro_left,function_apply_intro}.
        \caseOf{$p\in\{r\}$.}
            $g(p)\in\pow{X}$
                by \cref{singleton_elim,singleton_intro,union_intro_right,function_apply_intro,normal_space,topology_on,subseteq}.
        \end{byCase}
    \end{subproof}
    Let $w = (n + 1,g)$.
    $n + 1\in\naturalsPlus$ by \cref{naturalsplus_add_closed,real_one_in_naturals}.
    Therefore $g\in\funs{\{0,1\}\union \img{s}{\initialSegment{n + 1}}}{\pow{X}}$
        by \cref{urysohn_unit_interval_extend_gt1_setup,funs_intro}.
    Show that for all $p\in\dom{g}$ we have
        $g(p)\in T$ and
        $g(1) = X\setminus B$ and
        $A\subseteq g(0)$ and
        $\closure{g(0)}{X}{T}\subseteq g(1)$ and
        for all $p\in\dom{g}$ for all $q\in\dom{g}$ if $p < q$ then
            $\closure{g(p)}{X}{T}\subseteq g(q)$ and
        for all $p\in\dom{g}$ if $1 < p$ then $g(p) = X$.
    \begin{subproof}
        Fix $p$ such that $p\in\dom{g}$.
        Show that $g(p)\in T$.
        \begin{subproof}
            $p\in\dom{f}$ or $p\in\dom{h}$ by \cref{dom_union,union_iff}.
            \begin{byCase}
            \caseOf{$p\in\dom{f}$.}
                Show that $g(p)\in T$.
                \begin{subproof}
                    Follows by \cref{urysohn_admissible,funs_is_function,function_apply_elim,union_intro_left,function_apply_intro}.
                \end{subproof}
            \caseOf{$p\in\dom{h}$.}
                Take $y$ such that $(p,y)\in h$ by \cref{dom_iff}.
                Hence $(p,y) = (r,U_r)$ by \cref{singleton_elim}.
                Therefore $p = r$ by \cref{pair_eq_iff}.
                We have $g(p) = U_r$ by \cref{union_intro_right,singleton_intro,function_apply_intro}.
                Therefore $g(p)\in T$ by \cref{normal_space,topology_on}.
            \end{byCase}
        \end{subproof}
        Hence $T$ is a topology on $X$ and $g(p)$ is open in $X$ with respect to $T$ by \cref{normal_space,open_in}.
        $1\in\dom{f}$ by \cref{urysohn_admissible,funs_elim,upair_intro_right,union_intro_left}.
        Hence $(1,f(1))\in f$ and $(1,f(1))\in g$
            by \cref{urysohn_admissible,funs_is_function,function_apply_elim,union_intro_left}.
        Therefore $g(1) = X\setminus B$ by \cref{function_apply_intro}.
        $(0,f(0))\in f$ and $(0,f(0))\in g$
            by \cref{urysohn_admissible,funs_is_function,funs_elim,function_apply_elim,upair_intro_left,union_intro_left}.
        Hence $g(0) = f(0)$ by \cref{function_apply_intro}.
        We have $A\subseteq g(0)$ and $\closure{g(0)}{X}{T}\subseteq g(1)$ by \cref{subseteq}.
        Show that for all $p\in\dom{g}$ for all $q\in\dom{g}$ if $p < q$ then
            $\closure{g(p)}{X}{T}\subseteq g(q)$.
        \begin{subproof}
            Fix $p_1$ such that $p_1\in\dom{g}$.
            Fix $q$ such that $q\in\dom{g}$.
            Assume $p_1 < q$.
            $p_1\in\dom{f}$ or $p_1\in\dom{h}$ by \cref{dom_union,union_iff}.
            We have $q\in\dom{f}$ or $q\in\dom{h}$ by \cref{dom_union,union_iff}.
            \begin{byCase}
            \caseOf{$q\in\dom{h}$.}
                Take $y$ such that $(q,y)\in h$ by \cref{dom_iff}.
                Hence $(q,y) = (r,U_r)$ by \cref{singleton_elim}.
                Therefore $q = r$ by \cref{pair_eq_iff}.
                $g(q) = X$ by \cref{union_intro_right,singleton_intro,function_apply_intro}.
                We have $\closure{g(p_1)}{X}{T}\subseteq X$
                    by \cref{closed_whole_space,funs_type_apply,powerset_elim,closure_subseteq_closed}.
                Therefore $\closure{g(p_1)}{X}{T}\subseteq g(q)$ by \cref{subseteq}.
            \caseOf{$q\in\dom{f}$.}
                \begin{byCase}
                \caseOf{$p_1\in\dom{f}$.}
                    $g(p_1) = f(p_1)$ and $g(q) = f(q)$
                        by \cref{funs_is_function,function_apply_elim,union_intro_left,function_apply_intro}.
                    Show that $\closure{g(p_1)}{X}{T}\subseteq g(q)$.
                    \begin{subproof}
                        Follows by \cref{urysohn_admissible,subseteq}.
                    \end{subproof}
                \caseOf{$p_1\in\dom{h}$.}
                    Take $y$ such that $(p_1,y)\in h$ by \cref{dom_iff}.
                    Therefore $p_1 = r$ by \cref{singleton_elim,pair_eq_iff}.
                    $r < q$ by \cref{real_lt_subst_left}.
                    We have $r\in\unitRationals$ and $r\in\reals$
                        by \cref{sequence_in,naturalsplus_add_closed,real_one_in_naturals,funs_type_apply,unit_rational_real}.
                    $q\in \{0,1\}\union \img{s}{\initialSegment{n}}$ by \cref{urysohn_admissible,funs_elim}.
                    Hence $q\in\reals$ by
                        \cref{union_iff,upair_elim,real_add_ident,real_mul_ident,img_iff,sequence_in,initial_segment_naturalsplus,funs_is_function,function_apply_intro,funs_type_apply,unit_rational_real}.
                    We have $1\in\reals$ and $1 < q$ by \cref{real_mul_ident,real_lt_trans}.
                    Therefore $f(q) = X$ by \cref{urysohn_admissible_above_one}.
                    $g(q) = f(q)$
                        by \cref{funs_is_function,function_apply_elim,union_intro_left,function_apply_intro}.
                    $g(p_1) = X$ by \cref{union_intro_right,singleton_intro,function_apply_intro}.
                    Show that $\closure{g(p_1)}{X}{T}\subseteq g(q)$.
                    \begin{subproof}
                        $\closure{g(p_1)}{X}{T}\subseteq X$
                            by \cref{closed_whole_space,funs_type_apply,powerset_elim,closure_subseteq_closed}.
                        Therefore $\closure{g(p_1)}{X}{T}\subseteq g(q)$ by \cref{subseteq}.
                    \end{subproof}
                \end{byCase}
            \end{byCase}
        \end{subproof}
        Show that for all $p\in\dom{g}$ if $1 < p$ then $g(p) = X$.
        \begin{subproof}
            Fix $p_1$ such that $p_1\in\dom{g}$.
            Assume $1 < p_1$.
            $p_1\in\dom{f}$ or $p_1\in\{r\}$ by \cref{union_iff}.
            \begin{byCase}
            \caseOf{$p_1\in\{r\}$.}
                $g(p_1) = X$ by \cref{singleton_elim,union_intro_right,singleton_intro,function_apply_intro}.
            \caseOf{$p_1\in\dom{f}$.}
                $g(p_1) = X$
                    by \cref{funs_is_function,function_apply_elim,union_intro_left,function_apply_intro,urysohn_admissible_above_one}.
            \end{byCase}
        \end{subproof}
    \end{subproof}
    Therefore $g\in\urysohnAdmissible{s}{X}{T}{A}{B}{n + 1}$ by \cref{urysohn_admissible}.
    Hence $z\in G$ and $w\in G$ and $\fst{w} = \fst{z} + 1$ and $\snd{z}\subseteq \snd{w}$
        by \cref{urysohn_good_pair_intro,fst_eq,snd_eq,union_upper_left}.
    Thus $(z,w)\in R$.
    Therefore $z\mathrel{R} w$.
\end{proof}
% from §33 helper lemma 28: extension for $r < 1$ (setup)
\begin{lemma}\label{urysohn_unit_interval_extend_lt1_setup}
    Assume $X$ is normal with respect to $T$.
    Assume $A$ is closed in $X$ with respect to $T$.
    Assume $B$ is closed in $X$ with respect to $T$.
    Assume $A$ is disjoint from $B$.
    Assume $s$ is a sequence in $\unitRationals$.
    Let $G = \urysohnGoodPairs{s}{X}{T}{A}{B}$.
    Let $R = \{(z,w) \mid z\in G, w\in G \mid
        \fst{w} = \fst{z} + 1 \land \snd{z}\subseteq\snd{w}\}$.
    Suppose $n\in\naturalsPlus$ and $f\in\urysohnAdmissible{s}{X}{T}{A}{B}{n}$ and $z = (n,f)$.
    Let $r = s(n + 1)$.
    Assume $r\notin\dom{f}$.
    Assume $r < 1$.
    Let $S = \{p\in\dom{f} \mid p < r\}$.
    Let $T_1 = \{q\in\dom{f} \mid r < q\}$.
    Then there exist $p_0,q_0,U_r,h,g$ such that
        $p_0$ is a largest element of $S$ with respect to $\realOrderRel$ and
        $q_0$ is a smallest element of $T_1$ with respect to $\realOrderRel$ and
        $U_r$ is open in $X$ with respect to $T$ and
        $\closure{f(p_0)}{X}{T}\subseteq U_r$ and
        $\closure{U_r}{X}{T}\subseteq f(q_0)$ and
        $h = \{(r,U_r)\}$ and
        $g = f\union h$ and
        $g$ is a function.
\end{lemma}
\begin{proof}
    $f\in\funs{\{0,1\}\union \img{s}{\initialSegment{n}}}{\pow{X}}$ and
        for all $p\in\dom{f}$ we have $f(p)\in T$ and
        $f(1) = X\setminus B$ and
        $A\subseteq f(0)$ and
        $\closure{f(0)}{X}{T}\subseteq f(1)$ and
        for all $p\in\dom{f}$ for all $q\in\dom{f}$ if $p < q$ then
            $\closure{f(p)}{X}{T}\subseteq f(q)$ and
        for all $p\in\dom{f}$ if $1 < p$ then $f(p) = X$
        by \cref{urysohn_admissible}.
    $0\in S$ and $1\in T_1$
        by \cref{urysohn_admissible,funs_elim,upair_intro_left,upair_intro_right,union_intro_left,sequence_in,naturalsplus_add_closed,real_one_in_naturals,funs_type_apply,unit_rational_pos}.
    Let $D = \{0,1\}\union \img{s}{\initialSegment{n}}$.
    $\dom{f} = D$ by \cref{funs_elim}.
    $\{0,1\}$ is finite and $\initialSegment{n}$ is finite by \cref{pair_is_finite,initial_segment_finite}.
    $s\in\funs{\naturalsPlus}{\unitRationals}$ and $s$ is a function by \cref{sequence_in,funs_is_function}.
    Let $g_0 = \restrl{s}{\initialSegment{n}}$.
    $g_0$ is a function and $\dom{s} = \naturalsPlus$ by \cref{restrl_of_function_is_function,funs_elim}.
    $\dom{g_0} = \dom{s}\inter \initialSegment{n}$ by \cref{dom_of_restrl_eq_inter}.
    Hence $\dom{g_0}\subseteq\initialSegment{n}$ by \cref{restrl_dom}.
    Thus $\initialSegment{n}\subseteq\dom{g_0}$
        by \cref{initial_segment_naturalsplus,funs_elim,inter_intro,subseteq}.
    Therefore $\dom{g_0} = \initialSegment{n}$ by \cref{subseteq_antisymmetric}.
    Hence $g_0$ is a function on $\initialSegment{n}$.
    Thus $\img{g_0}{\initialSegment{n}}$ is finite by \cref{finite_image_function}.
    Therefore $\img{s}{\initialSegment{n}}$ is finite
        by \cref{funs_elim,initial_segment_naturalsplus,subseteq,restrl_img,inter_idempotent}.
    Thus $D$ is finite and $\dom{f}$ is finite by \cref{union_of_finite_is_finite}.
    We have $S\subseteq\dom{f}$ and $T_1\subseteq\dom{f}$ by \cref{subseteq}.
    Hence $S$ is finite and $T_1$ is finite by \cref{subseteq_finite_is_finite}.
    $\{0,1\}\subseteq\reals$
        by \cref{upair_elim,real_add_ident,real_naturals_subset_reals,real_one_in_naturals,subseteq}.
    $\img{s}{\initialSegment{n}}\subseteq\reals$
        by \cref{img_iff,initial_segment_naturalsplus,sequence_in,funs_is_function,function_apply_intro,funs_type_apply,unit_rational_real,subseteq}.
    Therefore $\dom{f}\subseteq\reals$ by \cref{union_subsets_is_subset}.
    Thus $S\subseteq\reals$ and $T_1\subseteq\reals$ by \cref{subseteq_transitive}.
    $\realOrderRel$ is a simple order relation on $\reals$ by \cref{real_order_relation_simple}.
    Take $p_1$ such that $p_1$ is a largest element of $S$ with respect to $\realOrderRel$
        by \cref{finite_inhabited_largest_element}.
    Take $q_1$ such that $q_1$ is a smallest element of $T_1$ with respect to $\realOrderRel$
        by \cref{finite_inhabited_smallest_element}.
    $p_1\in\dom{f}$ and $p_1 < r$ and $q_1\in\dom{f}$ and $r < q_1$
        by \cref{largest_element,smallest_element,subseteq}.
    Hence $p_1 < q_1$
        by \cref{subseteq,sequence_in,naturalsplus_add_closed,real_one_in_naturals,funs_type_apply,unit_rational_real,real_lt_trans}.
    Thus $\closure{f(p_1)}{X}{T}\subseteq f(q_1)$ by \cref{urysohn_admissible}.
    Show that $\closure{f(p_1)}{X}{T}$ is closed in $X$ with respect to $T$.
    \begin{subproof}
        Follows by \cref{sequence_in,naturalsplus_add_closed,real_one_in_naturals,funs_type_apply,unit_rational_real,subseteq,real_lt_trans,normal_space,urysohn_admissible,topology_on,elem_subseteq,powerset_elim,closure_closed_in}.
    \end{subproof}
    Show that $f(q_1)$ is open in $X$ with respect to $T$.
    \begin{subproof}
        Follows by \cref{sequence_in,naturalsplus_add_closed,real_one_in_naturals,funs_type_apply,unit_rational_real,subseteq,real_lt_trans,normal_space,urysohn_admissible,topology_on,elem_subseteq,powerset_elim,open_in}.
    \end{subproof}
    Take $U_1$ such that
        $U_1$ is open in $X$ with respect to $T$ and
        $\closure{f(p_1)}{X}{T}\subseteq U_1$ and
        $\closure{U_1}{X}{T}\subseteq f(q_1)$
        by \cref{normal_closure_neighborhood}.
    Let $h_1 = \{(r,U_1)\}$.
    Let $g_1 = f\union h_1$.
    Show that for all $x$ such that $x\in\dom{f}\inter\dom{h_1}$ we have $f(x) = h_1(x)$.
    \begin{subproof}
        Fix $x$ such that $x\in\dom{f}\inter\dom{h_1}$.
        Suppose not.
        $x\in\dom{h_1}$ by \cref{inter_elim_right}.
        $x\in\{r\}$ by \cref{dom_iff,singleton_elim,pair_eq_iff,singleton_intro}.
        Thus $x = r$ by \cref{singleton_elim}.
        Therefore $r\in\dom{f}$ by \cref{inter_elim_left}.
        Contradiction.
    \end{subproof}
    Show that $h_1$ is a relation.
    \begin{subproof}
        Follows by \cref{singleton_elim,relation}.
    \end{subproof}
    Show that $h_1$ is right-unique.
    \begin{subproof}
        Follows by \cref{singleton_elim,pair_eq_iff,rightunique}.
    \end{subproof}
    Hence $g_1$ is a function by \cref{urysohn_admissible,funs_is_function,union_compatible_functions}.
    Therefore there exist $p_0,q_0,U_r,h,g$ such that
        $p_0$ is a largest element of $S$ with respect to $\realOrderRel$ and
        $q_0$ is a smallest element of $T_1$ with respect to $\realOrderRel$ and
        $U_r$ is open in $X$ with respect to $T$ and
        $\closure{f(p_0)}{X}{T}\subseteq U_r$ and
        $\closure{U_r}{X}{T}\subseteq f(q_0)$ and
        $h = \{(r,U_r)\}$ and
        $g = f\union h$ and
        $g$ is a function
        by \cref{largest_element,smallest_element,subseteq}.
\end{proof}
% from §33 helper lemma 29: extension for $r < 1$
\begin{lemma}\label{urysohn_unit_interval_extend_lt1}
    Assume $X$ is normal with respect to $T$.
    Assume $A$ is closed in $X$ with respect to $T$.
    Assume $B$ is closed in $X$ with respect to $T$.
    Assume $A$ is disjoint from $B$.
    Assume $s$ is a sequence in $\unitRationals$.
    Let $G = \urysohnGoodPairs{s}{X}{T}{A}{B}$.
    Let $R = \{(z,w) \mid z\in G, w\in G \mid
        \fst{w} = \fst{z} + 1 \land \snd{z}\subseteq\snd{w}\}$.
    Suppose $n\in\naturalsPlus$ and $f\in\urysohnAdmissible{s}{X}{T}{A}{B}{n}$ and $z = (n,f)$.
    Let $r = s(n + 1)$.
    Assume $r\notin\dom{f}$.
    Assume $r < 1$.
    Let $S = \{p\in\dom{f} \mid p < r\}$.
    Let $T_1 = \{q\in\dom{f} \mid r < q\}$.
    Then there exists $w$ such that $z\mathrel{R} w$.
\end{lemma}
\begin{proof}
    We have $f\in\funs{\{0,1\}\union \img{s}{\initialSegment{n}}}{\pow{X}}$ and
        for all $p\in\dom{f}$ we have $f(p)\in T$ and
        $f(1) = X\setminus B$ and
        $A\subseteq f(0)$ and
        $\closure{f(0)}{X}{T}\subseteq f(1)$ and
        for all $p\in\dom{f}$ for all $q\in\dom{f}$ if $p < q$ then
            $\closure{f(p)}{X}{T}\subseteq f(q)$ and
        for all $p\in\dom{f}$ if $1 < p$ then $f(p) = X$
        by \cref{urysohn_admissible}.
    Take $p_0,q_0,U_r,h,g$ such that
        $p_0$ is a largest element of $S$ with respect to $\realOrderRel$ and
        $q_0$ is a smallest element of $T_1$ with respect to $\realOrderRel$ and
        $U_r$ is open in $X$ with respect to $T$ and
        $\closure{f(p_0)}{X}{T}\subseteq U_r$ and
        $\closure{U_r}{X}{T}\subseteq f(q_0)$ and
        $h = \{(r,U_r)\}$ and
        $g = f\union h$ and
        $g$ is a function
        by \cref{urysohn_unit_interval_extend_lt1_setup}.
    $\dom{g} = \dom{f}\union\dom{h}$ by \cref{dom_union}.
    $\dom{h} = \{r\}$
        by \cref{dom_iff,singleton_elim,pair_eq_iff,singleton_intro,dom_intro,subseteq,subseteq_antisymmetric}.
    $\dom{g} = \dom{f}\union\{r\}$.
    Show that for all $y\in\dom{g}$ we have $y\in \{0,1\}\union \img{s}{\initialSegment{n + 1}}$.
    \begin{subproof}
        Follows by \cref{union_iff,funs_elim,initial_segment_subset_succ,subseteq,img_iff,union_intro_left,singleton_elim,naturalsplus_add_closed,real_one_in_naturals,sequence_in,sequence_value_in_initial_image,union_intro_right}.
    \end{subproof}
    Thus $\dom{g}\subseteq \{0,1\}\union \img{s}{\initialSegment{n + 1}}$ by \cref{subseteq}.
    Show that for all $y\in \{0,1\}\union \img{s}{\initialSegment{n + 1}}$ we have $y\in\dom{g}$.
    \begin{subproof}
        Fix $y$ such that $y\in \{0,1\}\union \img{s}{\initialSegment{n + 1}}$.
        $y\in \{0,1\}$ or $y\in \img{s}{\initialSegment{n + 1}}$ by \cref{union_iff}.
        \begin{byCase}
        \caseOf{$y\in \{0,1\}$.}
            We have $y\in\dom{g}$ by \cref{union_intro_left,funs_elim}.
        \caseOf{$y\in \img{s}{\initialSegment{n + 1}}$.}
            Take $k$ such that $k\in\initialSegment{n + 1}$ and $(k,y)\in s$ by \cref{img_iff}.
            \begin{byCase}
            \caseOf{$k\in\initialSegment{n}$.}
                $y\in\dom{g}$ by \cref{img_iff,union_intro_right,funs_elim,union_intro_left}.
            \caseOf{$k\notin\initialSegment{n}$.}
                $k\in\naturalsPlus$ and $k\leq n + 1$ by \cref{initial_segment_naturalsplus}.
                $k = n + 1$ by \cref{real_naturals_leq_succ_elim,initial_segment_naturalsplus}.
                $y = r$.
                $y\in\dom{g}$ by \cref{singleton_intro,union_intro_right}.
            \end{byCase}
        \end{byCase}
    \end{subproof}
    Hence $\{0,1\}\union \img{s}{\initialSegment{n + 1}}\subseteq \dom{g}$ by \cref{subseteq}.
    Therefore $\dom{g} = \{0,1\}\union \img{s}{\initialSegment{n + 1}}$ by \cref{subseteq_antisymmetric}.
    For all $p\in\dom{g}$ we have $g(p)$ is open in $X$ with respect to $T$
        by \cref{union_iff,urysohn_admissible,open_in,topology_on,normal_space,funs_is_function,function_apply_elim,union_intro_left,function_apply_intro,singleton_elim,singleton_intro,union_intro_right,pair_eq_iff}.
    Show that for all $p,q$ such that $p\in\dom{g}$ and $q\in\dom{g}$ and $p < q$ we have
        $\closure{g(p)}{X}{T}\subseteq g(q)$.
    \begin{subproof}
        Fix $p$.
        Fix $q$.
        Assume $p\in\dom{g}$ and $q\in\dom{g}$ and $p < q$.
        $p\in\dom{f}\union\{r\}$ and $q\in\dom{f}\union\{r\}$.
        $p\in\dom{f}$ or $p\in\{r\}$ by \cref{union_iff}.
        We have $q\in\dom{f}$ or $q\in\{r\}$ by \cref{union_iff}.
        \begin{byCase}
        \caseOf{$p\in\dom{f}$.}
            \begin{byCase}
            \caseOf{$q\in\dom{f}$.}
                $g(p) = f(p)$ and $g(q) = f(q)$
                    by \cref{funs_is_function,function_apply_elim,union_intro_left,function_apply_intro}.
                $\closure{g(p)}{X}{T}\subseteq g(q)$
                    by \cref{urysohn_admissible,funs_is_function,function_apply_elim,union_intro_left,function_apply_intro}.
            \caseOf{$q\in\{r\}$.}
                $q = r$ by \cref{singleton_elim}.
                $p < r$.
                $g(p) = f(p)$
                    by \cref{funs_is_function,function_apply_elim,union_intro_left,function_apply_intro}.
                $p\in S$.
                $\closure{f(p)}{X}{T}\subseteq U_r$
                    by \cref{real_order_relation_elim,largest_element,subseteq,urysohn_admissible,normal_space,topology_on,elem_subseteq,powerset_elim,subseteq_closure,subseteq_transitive}.
                $g(q) = U_r$ by \cref{singleton_elim,singleton_intro,union_intro_right,function_apply_intro}.
                $\closure{g(p)}{X}{T}\subseteq g(q)$.
            \end{byCase}
        \caseOf{$p\in\{r\}$.}
            \begin{byCase}
            \caseOf{$q\in\dom{f}$.}
                $p = r$ by \cref{singleton_elim}.
                $r < q$.
                $g(p) = U_r$ by \cref{singleton_elim,singleton_intro,union_intro_right,function_apply_intro}.
                $g(q) = f(q)$
                    by \cref{funs_is_function,function_apply_elim,union_intro_left,function_apply_intro}.
                $q\in T_1$.
                $\closure{U_r}{X}{T}\subseteq f(q)$
                    by \cref{real_order_relation_elim,smallest_element,subseteq,urysohn_admissible,normal_space,topology_on,elem_subseteq,powerset_elim,subseteq_closure,subseteq_transitive}.
                $\closure{g(p)}{X}{T}\subseteq g(q)$.
            \caseOf{$q\in\{r\}$.}
                Suppose not.
                $p = r$ and $q = r$ by \cref{singleton_elim}.
                $p\in\reals$ and $p < p$
                    by \cref{sequence_in,naturalsplus_add_closed,real_one_in_naturals,funs_type_apply,unit_rational_real,real_lt_subst_left,real_lt_subst_right}.
                Contradiction by \cref{real_lt_irreflexive}.
            \end{byCase}
        \end{byCase}
    \end{subproof}
    Show that for all $p\in\dom{g}$ if $1 < p$ then $g(p) = X$.
    \begin{subproof}
        Fix $p$ such that $p\in\dom{g}$.
        Assume $1 < p$.
        $p\in\dom{f}\union\{r\}$.
        $p\in\dom{f}$ or $p = r$ by \cref{union_iff,singleton_elim}.
        $g(p) = X$
            by \cref{urysohn_admissible_above_one,funs_is_function,function_apply_elim,union_intro_left,function_apply_intro,real_lt_subst_left,sequence_in,naturalsplus_add_closed,real_one_in_naturals,funs_type_apply,unit_rational_real,real_naturals_subset_reals,subseteq,real_lt_trans,real_lt_irreflexive}.
    \end{subproof}
    We have for all $p\in\dom{g}$ we have $g(p)\in\pow{X}$
        by \cref{open_in,normal_space,topology_on,elem_subseteq}.
    $g\in\funs{\{0,1\}\union \img{s}{\initialSegment{n + 1}}}{\pow{X}}$
        by \cref{urysohn_unit_interval_extend_lt1_setup,funs_intro}.
    $0\in\dom{f}$ and $1\in\dom{f}$
        by \cref{funs_elim,upair_intro_left,upair_intro_right,union_intro_left}.
    $g(1) = X\setminus B$ and $A\subseteq g(0)$ and
        $\closure{g(0)}{X}{T}\subseteq g(1)$
        by \cref{urysohn_admissible,funs_is_function,function_apply_elim,union_intro_left,function_apply_intro}.
    For all $p\in\dom{g}$ we have $g(p)\in T$ by \cref{open_in}.
    $g\in\urysohnAdmissible{s}{X}{T}{A}{B}{n + 1}$ by \cref{urysohn_admissible}.
    Let $w = (n + 1,g)$.
    $z\in G$ and $w\in G$ and $\fst{w} = \fst{z} + 1$ and $\snd{z}\subseteq \snd{w}$
        by \cref{naturalsplus_add_closed,real_one_in_naturals,urysohn_good_pair_intro,fst_eq,snd_eq,union_upper_left}.
    Hence $(z,w)\in R$.
    Therefore $z\mathrel{R} w$.
\end{proof}

% from §33 helper lemma 30: extension relation on good pairs
\begin{lemma}\label{urysohn_unit_interval_extend}
    Assume $X$ is normal with respect to $T$.
    Assume $A$ is closed in $X$ with respect to $T$.
    Assume $B$ is closed in $X$ with respect to $T$.
    Assume $A$ is disjoint from $B$.
    Assume $s$ is a sequence in $\unitRationals$.
    Assume for all $p$ such that $p\in\unitRationals$ there exists $n\in\naturalsPlus$ such that $s(n) = p$.
    Let $G = \urysohnGoodPairs{s}{X}{T}{A}{B}$.
    Let $R = \{(z,w) \mid z\in G, w\in G \mid
        \fst{w} = \fst{z} + 1 \land \snd{z}\subseteq\snd{w}\}$.
    Then for all $z\in G$ there exists $w$ such that $z\mathrel{R} w$.
\end{lemma}
\begin{proof}
    Fix $z$ such that $z\in G$.
    Take $n,f$ such that
        $n\in\naturalsPlus$ and
        $f\in\urysohnAdmissible{s}{X}{T}{A}{B}{n}$ and
        $z = (n,f)$
        by \cref{urysohn_good_pair_decomp}.
    Let $r = s(n + 1)$.
    $r\in\unitRationals$ and $r\in\reals$ and $0 < r$
        by \cref{sequence_in,funs_type_apply,naturalsplus_add_closed,real_one_in_naturals,unit_rational_real,unit_rational_pos}.
    Show that there exists $w$ such that $z\mathrel{R} w$.
    \begin{subproof}
        \begin{byCase}
        \caseOf{$r\in\dom{f}$.}
            Follows by \cref{urysohn_unit_interval_extend_in_dom}.
        \caseOf{$r\notin\dom{f}$.}
            $1\in\dom{f}$ by \cref{urysohn_admissible,funs_elim,upair_intro_right,union_intro_left}.
            $r \neq 1$ by \cref{real_mul_ident}.
            $r < 1$ or $1 < r$
                by \cref{real_naturals_subset_reals,real_one_in_naturals,subseteq,real_lt_trichotomy}.
            \begin{byCase}
            \caseOf{$1 < r$.}
                Follows by \cref{urysohn_unit_interval_extend_gt1}.
            \caseOf{$r < 1$.}
                Let $S = \{p\in\dom{f} \mid p < r\}$.
                Let $T_1 = \{q\in\dom{f} \mid r < q\}$.
                Follows by \cref{urysohn_unit_interval_extend_lt1}.
            \end{byCase}
        \end{byCase}
    \end{subproof}
\end{proof}

% from §33 helper lemma 31: Urysohn function vanishes on $A$
\begin{lemma}\label{urysohn_unit_interval_value_A}
    Assume $X$ is normal with respect to $T$.
    Assume $A$ is closed in $X$ with respect to $T$.
    Assume $B$ is closed in $X$ with respect to $T$.
    Assume $A$ is disjoint from $B$.
    Assume $s$ is a sequence in $\unitRationals$.
    Assume for all $p$ such that $p\in\unitRationals$ there exists $n\in\naturalsPlus$ such that $s(n) = p$.
    Let $G = \urysohnGoodPairs{s}{X}{T}{A}{B}$.
    Let $R = \{(z,w) \mid z\in G, w\in G \mid
        \fst{w} = \fst{z} + 1 \land \snd{z}\subseteq\snd{w}\}$.
    Assume $c$ is a function on $G$.
    Assume for all $z\in G$ we have $z\mathrel{R} \apply{c}{z}$.
    Assume $u$ is a sequence in $G$.
    Assume $u(1) = a_0$.
    Assume $a_0 = (1,u_1)$.
    Assume for all $n\in\naturalsPlus$ we have $u(n + 1) = \apply{c}{u(n)}$.
    Assume for all $n\in\naturalsPlus$ we have $\fst{u(n)} = n$.
    Let $F = \{\snd{u(n)} \mid n\in\naturalsPlus\}$.
    Let $U = \unions{F}$.
    Assume $U$ is a function.
    Let $I_0 = \intervalclosed{0}{1}$.
    Let $f = \{(x,y) \mid x\in X, y\in\reals \mid
        \text{$y$ is an infimum of $\urysohnLevelSet{U}{x}$}\}$.
    Assume $f\in\funs{X}{\reals}$.
    Assume $\img{f}{X}\subseteq I_0$.
    Then for all $x\in A$ we have $f(x) = 0$.
\end{lemma}
\begin{proof}
    Fix $x\in A$.
    Let $Q_x = \urysohnLevelSet{U}{x}$.
    $1\in\naturalsPlus$ and $u$ is a sequence in $G$ and $u\in\funs{\naturalsPlus}{G}$
        by \cref{real_one_in_naturals,sequence_in}.
    Take $n,t$ such that
        $n\in\naturalsPlus$ and
        $t\in\urysohnAdmissible{s}{X}{T}{A}{B}{n}$ and
        $u(1) = (n,t)$
        by \cref{funs_type_apply,urysohn_good_pair_decomp}.
    $t$ is a function and $0\in\dom{t}$
        by \cref{urysohn_admissible,funs_is_function,funs_elim,upair_intro_left,union_intro_left,subseteq}.
    $0\in Q_x$
        by \cref{snd_eq,function_apply_elim,unions_intro,function_apply_intro,urysohn_admissible,subseteq,unit_rationals_zero,singleton_intro,union_intro_right,urysohn_level_set}.
    $x\in X$ and $x\in\dom{f}$ and $(x,f(x))\in f$
        by \cref{funs_is_function,funs_elim,closed_in,subseteq,function_apply_elim}.
    $f(x)$ is an infimum of $Q_x$ by \cref{urysohn_infimum_value}.
    Thus $f(x) \leq 0$ by \cref{real_infimum,real_lower_bound}.
    $0 \leq f(x)$ by \cref{img_iff,subseteq,intervalclosed}.
    Hence $f(x) = 0$ by \cref{real_add_ident,funs_type_apply,real_leq_antisym}.
\end{proof}

% from §33 helper lemma 32: Urysohn function equals $1$ on $B$
\begin{lemma}\label{urysohn_unit_interval_value_B}
    Assume $X$ is normal with respect to $T$.
    Assume $A$ is closed in $X$ with respect to $T$.
    Assume $B$ is closed in $X$ with respect to $T$.
    Assume $A$ is disjoint from $B$.
    Assume $s$ is a sequence in $\unitRationals$.
    Assume for all $p$ such that $p\in\unitRationals$ there exists $n\in\naturalsPlus$ such that $s(n) = p$.
    Let $G = \urysohnGoodPairs{s}{X}{T}{A}{B}$.
    Let $R = \{(z,w) \mid z\in G, w\in G \mid
        \fst{w} = \fst{z} + 1 \land \snd{z}\subseteq\snd{w}\}$.
    Assume $c$ is a function on $G$.
    Assume for all $z\in G$ we have $z\mathrel{R} \apply{c}{z}$.
    Assume $u$ is a sequence in $G$.
    Assume $u(1) = a_0$.
    Assume $a_0 = (1,u_1)$.
    Assume for all $n\in\naturalsPlus$ we have $u(n + 1) = \apply{c}{u(n)}$.
    Assume for all $n\in\naturalsPlus$ we have $\fst{u(n)} = n$.
    Let $F = \{\snd{u(n)} \mid n\in\naturalsPlus\}$.
    Let $U = \unions{F}$.
    Assume $U$ is a function.
    Let $I_0 = \intervalclosed{0}{1}$.
    Let $f = \{(x,y) \mid x\in X, y\in\reals \mid
        \text{$y$ is an infimum of $\urysohnLevelSet{U}{x}$}\}$.
    Assume $f\in\funs{X}{\reals}$.
    Assume $\img{f}{X}\subseteq I_0$.
    Then for all $x\in B$ we have $f(x) = 1$.
\end{lemma}
\begin{proof}
    Fix $x\in B$.
    Let $Q_x = \urysohnLevelSet{U}{x}$.
    $1\in\naturalsPlus$ and $u$ is a sequence in $G$ and $u\in\funs{\naturalsPlus}{G}$
        by \cref{real_one_in_naturals,sequence_in}.
    Take $n,t$ such that
        $n\in\naturalsPlus$ and
        $t\in\urysohnAdmissible{s}{X}{T}{A}{B}{n}$ and
        $u(1) = (n,t)$
        by \cref{funs_type_apply,urysohn_good_pair_decomp}.
    $t$ is a function and $1\in\dom{t}$
        by \cref{urysohn_admissible,funs_is_function,funs_elim,upair_intro_right,union_intro_left,subseteq}.
    $x\notin U(1)$
        by \cref{snd_eq,function_apply_elim,unions_intro,function_apply_intro,urysohn_admissible,subseteq_refl,subseteq,setminus_elim_right}.
    $x\in X$ and $x\in\dom{f}$ and $(x,f(x))\in f$
        by \cref{funs_is_function,funs_elim,closed_in,subseteq,function_apply_elim}.
    $f(x)$ is an infimum of $Q_x$ by \cref{urysohn_infimum_value}.
    $f(x)$ is a lower bound of $Q_x$ by \cref{real_infimum}.
    Show that for all $q\in Q_x$ we have $1 \leq q$.
    \begin{subproof}
            Fix $q$ such that $q\in Q_x$.
            We have $q\in\unitRationalsZero$ and $x\in U(q)$ by \cref{urysohn_level_set}.
            Suppose not.
            $q < 1$.
            $1\in\unitRationals$
                by \cref{times_tuple_intro,real_one_in_naturals,real_naturals_subset_reals,subseteq,real_zero_lt_one,real_pos_nonzero,real_mul_inverse,real_div,real_mul_ident,fst_eq,snd_eq,unit_rationals}.
            $\closure{U(q)}{X}{T}\subseteq U(1)$ by \cref{urysohn_union_closure_subset}.
            $U(q)\subseteq X$ by \cref{urysohn_union_zero_subset_X_aux}.
            $T$ is a topology on $X$ by \cref{normal_space}.
            We have $U(q)\subseteq \closure{U(q)}{X}{T}$ by \cref{subseteq_closure}.
            $x\in U(1)$ by \cref{subseteq,subseteq_transitive}.
            Contradiction.
    \end{subproof}
    Thus $1$ is a lower bound of $Q_x$
        by \cref{real_mul_ident,urysohn_level_set,unit_rational_zero_real,subseteq,real_lower_bound}.
    $1 \leq f(x)$ by \cref{real_infimum}.
    $f(x) \leq 1$ by \cref{img_iff,subseteq,intervalclosed}.
    Hence $f(x) = 1$ by \cref{real_mul_ident,funs_type_apply,real_leq_antisym}.
\end{proof}

% from §33 helper theorem 33: Urysohn lemma for the unit interval
\begin{lemma}\label{urysohn_unit_interval}
    Assume $X$ is normal with respect to $T$.
    Assume $A$ is closed in $X$ with respect to $T$.
    Assume $B$ is closed in $X$ with respect to $T$.
    Assume $A$ is disjoint from $B$.
    Let $I_0 = \intervalclosed{0}{1}$.
    Let $T_0 = \subspaceTopology{\standardTopologyR}{I_0}$.
    Then there exists $f$ such that
        $f$ is continuous from $X$ to $I_0$ with respect to $T$ and $T_0$ and
        $\img{f}{A}\subseteq \{0\}$ and
        $\img{f}{B}\subseteq \{1\}$.
\end{lemma}
\begin{proof}
    Take $s$ such that
        $s$ is a sequence in $\unitRationals$ and
        for all $p$ such that $p\in\unitRationals$ there exists $n\in\naturalsPlus$ such that $s(n) = p$
        by \cref{unit_rationals_sequence}.
    Let $G = \urysohnGoodPairs{s}{X}{T}{A}{B}$.

    Let $R = \{(z,w) \mid z\in G, w\in G \mid
        \fst{w} = \fst{z} + 1 \land \snd{z}\subseteq\snd{w}\}$.

    Show that for all $z\in G$ there exists $w$ such that $z\mathrel{R} w$.
    \begin{subproof}
        Follows by \cref{urysohn_unit_interval_extend}.
    \end{subproof}

    $R$ is a relation by \cref{relation}.

    Take $c$ such that $c$ is a function on $G$ and for all $z\in G$ we have $z\mathrel{R} \apply{c}{z}$
        by \cref{choice_function}.

    Take $u_1$ such that $(1,u_1)\in G$ by \cref{urysohn_unit_interval_u1}.
    Let $a_0 = (1,u_1)$.

    Show that for all $x\in\dom{c}$ we have $c(x)\in G$.
    \begin{subproof}
        Fix $x$ such that $x\in\dom{c}$.
        We have $x\in G$ by \cref{funs_elim}.
        $x\mathrel{R} c(x)$ by \cref{choice_function}.
        Take $z,w$ such that $z\in G$ and $w\in G$ and $(z,w) = (x,c(x))$ and
            $\fst{w} = \fst{z} + 1$ and $\snd{z}\subseteq\snd{w}$ by \cref{relation}.
        $w = c(x)$ by \cref{pair_eq_iff}.
        $c(x)\in G$ by \cref{subseteq}.
    \end{subproof}
    $c\in\funs{G}{G}$ by \cref{funs_intro}.
    Take $u$ such that
        $u$ is a sequence in $G$ and
        $u(1) = a_0$ and
        for all $n\in\naturalsPlus$ we have $u(n + 1) = \apply{c}{u(n)}$
        by \cref{iteration_sequence_naturalsplus}.
    For all $n\in\naturalsPlus$ we have $\fst{u(n)} = n$
        by \cref{urysohn_unit_interval_fst_u}.

    Let $F = \{\snd{u(n)} \mid n\in\naturalsPlus\}$.
    Let $U = \unions{F}$.
    $U$ is a function by \cref{urysohn_unit_interval_union_function}.

    Let $f = \{(x,y) \mid x\in X, y\in\reals \mid
        \text{$y$ is an infimum of $\urysohnLevelSet{U}{x}$}\}$.
    $f$ is a function from $X$ to $\reals$ by \cref{urysohn_unit_interval_f_funs}.

    Show that for all $x\in X$ we have $f(x)\in I_0$.
    \begin{subproof}
        Fix $x\in X$.
        Let $Q_x = \urysohnLevelSet{U}{x}$.
        We have $f(x)$ is an infimum of $Q_x$ by \cref{urysohn_infimum_value}.
        $0\in\reals$ by \cref{real_add_ident}.
        For all $p\in Q_x$ we have $p\in\reals$ and $0 \leq p$
            by \cref{urysohn_level_set,unit_rational_zero_real,unit_rationals_zero,union_iff,singleton_elim,real_add_ident,unit_rational_pos,real_lt_imp_leq}.
        $Q_x\subseteq\reals$ by \cref{subseteq}.
        $0$ is a lower bound of $Q_x$ by \cref{real_lower_bound}.
        $0 < f(x)$ or $0 = f(x)$ by \cref{real_infimum}.
        $0 \leq f(x)$.
        Show that $f(x) \leq 1$.
        \begin{subproof}
            Suppose not.
            $1 < f(x)$.
            $1\in\reals$ and $0 < 1$ by \cref{real_mul_ident,real_zero_lt_one}.
            $0 \leq 1$ by \cref{real_lt_imp_leq}.
            $f(x)\in\reals$ by \cref{funs_type_apply}.
            Take $p$ such that $p\in\unitRationals$ and $p\in\intervalopen{1}{f(x)}$
                by \cref{unit_rationals_dense_nonneg}.
            $p\in\reals$ by \cref{intervalopen}.
            $p > 1$.
            $p < f(x)$ by \cref{intervalopen}.
            $p\in\unitRationalsZero$ by \cref{unit_rationals_zero,union_iff}.
            Take $n,t$ such that
                $t\in\urysohnAdmissible{s}{X}{T}{A}{B}{n}$ and
                $t$ is a function and
                $p\in\dom{t}$ and
                $U(p) = t(p)$
                by \cref{urysohn_union_rational_value_aux}.
            $t(p) = X$ by \cref{urysohn_admissible_above_one}.
            $U(p) = X$ by \cref{function_apply_intro}.
            $x\in U(p)$.
            $p\in Q_x$ by \cref{urysohn_level_set}.
            $f(x)$ is a lower bound of $Q_x$ by \cref{real_infimum}.
            $f(x) < p$ or $f(x) = p$ by \cref{real_lower_bound}.
            $f(x) < f(x)$ by \cref{real_lt_trans,real_lt_subst_left}.
            Contradiction by \cref{real_lt_irreflexive}.
        \end{subproof}
        $f(x)\in I_0$ by \cref{funs_type_apply,intervalclosed}.
    \end{subproof}
    Hence for all $x\in\dom{f}$ we have $f(x)\in I_0$ by \cref{funs_elim}.
    $f\in\funs{X}{I_0}$
        by \cref{function_rightunique,funs_is_relation,funs_elim,funs_intro}.
    $\img{f}{X}\subseteq I_0$ by \cref{funs_ran,img_subseteq_ran,subseteq_transitive}.

    Show that for all $x\in A$ we have $f(x) = 0$.
    \begin{subproof}
        Follows by \cref{urysohn_unit_interval_value_A}.
    \end{subproof}

    Show that for all $x\in B$ we have $f(x) = 1$.
    \begin{subproof}
        Follows by \cref{urysohn_unit_interval_value_B}.
    \end{subproof}

    We have $\img{f}{A}\subseteq \{0\}$ by \cref{img_iff,funs_is_function,function_apply_iff,singleton_intro,subseteq}.
    Thus $\img{f}{B}\subseteq \{1\}$ by \cref{img_iff,funs_is_function,function_apply_iff,singleton_intro,subseteq}.

    $I_0\subseteq\reals$ by \cref{intervalclosed,subseteq}.
    $f$ is continuous from $X$ to $\reals$ with respect to $T$ and $\standardTopologyR$
        by \cref{urysohn_unit_interval_continuous_reals}.
    $T$ is a topology on $X$ by \cref{normal_space}.
    We have $T_0 = \subspaceTopology{\standardTopologyR}{I_0}$.
    $f$ is continuous from $X$ to $I_0$ with respect to $T$ and $T_0$
        by \cref{urysohn_unit_interval_f_funs,urysohn_unit_interval_continuous_subspace}.
    Therefore there exists $g$ such that
        $g$ is continuous from $X$ to $I_0$ with respect to $T$ and $T_0$ and
        $\img{g}{A}\subseteq \{0\}$ and
        $\img{g}{B}\subseteq \{1\}$.
\end{proof}

% from §33 Item 1: Urysohn lemma
\begin{theorem}\label{urysohn_lemma}
    Assume $X$ is normal with respect to $T$.
    Assume $A$ is closed in $X$ with respect to $T$.
    Assume $B$ is closed in $X$ with respect to $T$.
    Assume $A$ is disjoint from $B$.
    Assume $a\in\reals$ and $b\in\reals$ and $a < b$.
    Let $I = \intervalclosed{a}{b}$.
    Let $T_I = \subspaceTopology{\standardTopologyR}{I}$.
    Then there exists $f$ such that
        $f$ is continuous from $X$ to $I$ with respect to $T$ and $T_I$ and
        for all $x\in A$ we have $f(x) = a$ and
        for all $x\in B$ we have $f(x) = b$.
\end{theorem}
\begin{proof}
    Let $I_0 = \intervalclosed{0}{1}$.
    Let $T_0 = \subspaceTopology{\standardTopologyR}{I_0}$.
    Take $g$ such that
        $g$ is continuous from $X$ to $I_0$ with respect to $T$ and $T_0$ and
        $\img{g}{A}\subseteq \{0\}$ and
        $\img{g}{B}\subseteq \{1\}$
        by \cref{urysohn_unit_interval}.
    Let $c = b - a$.
    Let $j = \identity{I_0}$.
    Let $g_1 = j\circ g$.
    $\standardTopologyR$ is a topology on $\reals$
        by \cref{standard_basis_is_basis,basis_topology_is_topology,standard_topology_reals}.
    $1\in\reals$ and $1 + 1\in\reals$
        by \cref{real_one_in_naturals,real_naturals_subset_reals,subseteq,real_add_closed}.
    $I_0\subseteq\reals$ by \cref{intervalclosed,subseteq}.
    $I_0$ is a subspace of $\reals$ with respect to $\standardTopologyR$ by \cref{subspace_of}.
    $g_1$ is continuous from $X$ to $\reals$ with respect to $T$ and $\standardTopologyR$
        by \cref{continuous_expand_range}.
    $g\in\funs{X}{I_0}$ and $g$ is a function by \cref{continuous,funs_is_function}.
    $g$ is right-unique and $g$ is a relation and $\dom{g} = X$
        by \cref{function_rightunique,funs_is_relation,funs_elim}.
    $g\in\rels{X}{I_0}$ by \cref{funs_subseteq_rels,subseteq}.
    $\ran{g}\subseteq I_0$ by \cref{rels_ran_subseteq}.
    We have $\dom{j} = I_0$ by \cref{id_dom}.
    Hence $\ran{g}\subseteq\dom{j}$.
    $j$ is right-unique and $j$ is a relation by \cref{id_rightunique,id_is_relation}.
    For all $x\in X$ we have $g_1(x) = g(x)$ by \cref{circ_apply,funs_elim,id_apply}.
    Let $k = \{(x,a) \mid x\in X\}$.
    Let $h = \{(x,c) \mid x\in X\}$.
    $T$ is a topology on $X$ by \cref{normal_space}.
    $c\in\reals$ by \cref{real_add_inverse,real_add_closed,real_sub}.
    $a + c = b$ by
        \cref{real_sub,real_add_assoc,real_add_comm,real_add_inverse,real_add_ident}.
    $I\subseteq\reals$ by \cref{intervalclosed,subseteq}.
    $I$ is a subspace of $\reals$ with respect to $\standardTopologyR$ by \cref{subspace_of}.
    $T_I$ is a topology on $I$ by \cref{subspace_topology_is_topology}.
    $T_0$ is a topology on $I_0$ by \cref{intervalclosed,subseteq,subspace_topology_is_topology}.
    $k$ is a function from $X$ to $\reals$ and $h$ is a function from $X$ to $\reals$
        by \cref{constant_map_function}.
    For all $x\in X$ we have $k(x) = a$ and $h(x) = c$
        by \cref{constant_map_function,funs_is_function,function_apply_intro}.
    $k$ is continuous from $X$ to $\reals$ with respect to $T$ and $\standardTopologyR$ and
        $h$ is continuous from $X$ to $\reals$ with respect to $T$ and $\standardTopologyR$
        by \cref{constant_continuous}.
    Let $m = \{(x, h(x) \rmul g_1(x)) \mid x\in X\}$.
    Let $f_0 = \{(x, k(x) + m(x)) \mid x\in X\}$.
    $m$ is continuous from $X$ to $\reals$ with respect to $T$ and $\standardTopologyR$
        by \cref{real_function_multiplication_continuous}.
    $f_0$ is continuous from $X$ to $\reals$ with respect to $T$ and $\standardTopologyR$
        by \cref{real_function_addition_continuous}.
    $m\in\funs{X}{\reals}$ and $m$ is a function and
        $f_0\in\funs{X}{\reals}$ and $f_0$ is a function
        by \cref{continuous,funs_is_function}.

    Show that for all $x\in X$ we have $f_0(x)\in I$.
        \begin{subproof}
            Fix $x\in X$.
            $g_1(x)\in I_0$ by \cref{circ_apply,funs_elim,id_apply}.
            $g_1(x)\in\reals$ by \cref{intervalclosed}.
            $0 \leq g_1(x)$ and $g_1(x) \leq 1$ by \cref{intervalclosed}.
            $(x, h(x) \rmul g_1(x))\in m$.
            $m(x) = h(x) \rmul g_1(x)$ by \cref{function_apply_iff}.
            $h(x) = c$.
            $m(x) = c \rmul g_1(x)$.
            $g_1(x) = 0$ or $0 < g_1(x)$ by \cref{real_lt_trichotomy}.
            $0 \leq c \rmul g_1(x)$
                by \cref{real_mul_zero,real_mul_comm,real_lt_imp_pos_diff,real_mul_pos_pos_gt,real_lt_imp_leq}.
            $0 \leq m(x)$.
            $g_1(x) = 1$ or $g_1(x) < 1$ by \cref{real_lt_trichotomy}.
            $c \rmul g_1(x) \leq c$
                by \cref{real_mul_comm,real_mul_ident,real_lt_imp_pos_diff,real_naturals_subset_reals,real_one_in_naturals,subseteq,real_lt_mul_pos}.
            $m(x) \leq c$.
            $m(x)\in\reals$ by \cref{real_mul_closed}.
            $(x, k(x) + m(x))\in f_0$.
            $f_0(x) = a + m(x)$ by \cref{function_apply_iff,real_add_eq_subst}.
            $a + m(x)\in\reals$ and $a \leq a + m(x)$
                by \cref{real_add_closed,real_add_ident,real_leq_add,real_add_comm}.
            $a + m(x) \leq b$ by \cref{real_leq_add,real_add_comm}.
            $a + m(x)\in I$ by \cref{intervalclosed}.
            $f_0(x)\in I$ by \cref{function_apply_intro}.
    \end{subproof}
    Thus $\img{f_0}{X}\subseteq I$
        by \cref{subseteq,img_iff,continuous,funs_is_function,function_apply_iff}.

    $f_0$ is continuous from $X$ to $I$ with respect to $T$ and $T_I$
        by \cref{continuous_restrict_range}.

    Show that for all $x\in A$ we have $f_0(x) = a$.
    \begin{subproof}
        Fix $x\in A$.
        $x\in X$ and $g(x)\in\img{g}{A}$ and $g(x) = 0$ and $g_1(x) = 0$ and $g_1(x)\in\reals$ and
            $(x, h(x) \rmul g_1(x))\in m$ and $m(x) = 0$ and $(x, k(x) + m(x))\in f_0$
            by \cref{closed_in,subseteq,function_apply_elim,img_iff,singleton_elim,circ_apply,funs_elim,id_apply,real_add_ident,function_apply_iff,real_mul_comm,real_mul_zero}.
        $f_0(x) = a$ by \cref{function_apply_iff,real_add_eq_subst,real_add_comm,real_add_ident}.
    \end{subproof}
    Show that for all $x\in B$ we have $f_0(x) = b$.
    \begin{subproof}
        Fix $x\in B$.
        $x\in X$ and $g(x)\in\img{g}{B}$ and $g(x) = 1$ and $g_1(x) = 1$ and $g_1(x)\in\reals$ and
            $(x, h(x) \rmul g_1(x))\in m$ and $m(x) = c$ and $(x, k(x) + m(x))\in f_0$ and
            $f_0(x) = a + c$
            by \cref{closed_in,subseteq,function_apply_elim,img_iff,singleton_elim,circ_apply,funs_elim,id_apply,real_mul_ident,function_apply_iff,real_mul_comm,real_add_eq_subst}.
        $f_0(x) = b$.
    \end{subproof}
    Therefore there exists $f$ such that
        $f$ is continuous from $X$ to $I$ with respect to $T$ and $T_I$ and
        for all $x\in A$ we have $f(x) = a$ and
        for all $x\in B$ we have $f(x) = b$.
\end{proof}

% from §33 Item 2: separation by a continuous function
\begin{definition}\label{separated_by_continuous_function}
    $A$ is functionally separable from $B$ in $X$ with respect to $T$ iff
        $A\subseteq X$ and
        $B\subseteq X$ and
        there exists $f$ such that
            $f$ is continuous from $X$ to $\intervalclosed{0}{1}$ with respect to $T$ and
                $\subspaceTopology{\standardTopologyR}{\intervalclosed{0}{1}}$ and
            for all $x\in A$ we have $f(x) = 0$ and
            for all $x\in B$ we have $f(x) = 1$.
\end{definition}

% from §33 Item 3: completely regular space
\begin{definition}\label{completely_regular_space}
    $X$ is completely regular with respect to $T$ iff
        $T$ is a topology on $X$ and
        $X$ has closed singletons with respect to $T$ and
        for all $x,A$ such that $x\in X$ and $A$ is closed in $X$ with respect to $T$ and $x\notin A$
            there exists $f$ such that
                $f$ is continuous from $X$ to $\intervalclosed{0}{1}$ with respect to $T$ and
                    $\subspaceTopology{\standardTopologyR}{\intervalclosed{0}{1}}$ and
                $f(x) = 1$ and
                for all $y\in A$ we have $f(y) = 0$.
\end{definition}

% from §33 Item 4: subspace of a completely regular space is completely regular
\begin{theorem}\label{subspace_completely_regular}
    Assume $X$ is completely regular with respect to $T$.
    Assume $Y$ is a subspace of $X$ with respect to $T$.
    Let $T_Y = \subspaceTopology{T}{Y}$.
    Then $Y$ is completely regular with respect to $T_Y$.
\end{theorem}
\begin{proof}
    $T$ is a topology on $X$ and $Y\subseteq X$ by \cref{completely_regular_space,subspace_of}.
    $T_Y$ is a topology on $Y$ and $Y$ has closed singletons with respect to $T_Y$
        by \cref{subspace_topology_is_topology,completely_regular_space,elem_subseteq,one_point_closed,singleton_subset_intro,inter_absorb_supseteq_right,closed_in_subspace_iff,subspace_of}.
    Let $I_0 = \intervalclosed{0}{1}$.
    Let $T_0 = \subspaceTopology{\standardTopologyR}{I_0}$.
    Show that for all $x,A$ such that $x\in Y$ and $A$ is closed in $Y$ with respect to $T_Y$ and $x\notin A$
        there exists $g$ such that
            $g$ is continuous from $Y$ to $I_0$ with respect to $T_Y$ and $T_0$ and
            $g(x) = 1$ and
            for all $y\in A$ we have $g(y) = 0$.
    \begin{subproof}
        Fix $x$.
        Fix $A$ such that $x\in Y$ and $A$ is closed in $Y$ with respect to $T_Y$ and $x\notin A$.
        We have $A\subseteq Y$ by \cref{closed_in}.
        Let $C = \closure{A}{X}{T}$.
        $\closure{A}{Y}{T_Y} = C\inter Y$ by \cref{closure_in_subspace,subspace_of}.
        $A = C\inter Y$ by \cref{subspace_topology_is_topology,closed_eq_closure}.
        Show that $x\notin C$.
        \begin{subproof}
            Suppose not.
            $x\in C\inter Y$ by \cref{inter_intro}.
            $x\in A$ by \cref{subseteq}.
            Contradiction.
        \end{subproof}
        $A\subseteq X$ by \cref{subseteq_transitive,subspace_of}.
        $x\in X$ by \cref{subseteq}.
        $C$ is closed in $X$ with respect to $T$ by \cref{closure_closed_in}.
        $A\subseteq C$ by \cref{subseteq_closure}.
        Take $f$ such that
            $f$ is continuous from $X$ to $I_0$ with respect to $T$ and $T_0$ and
            $f(x) = 1$ and
            for all $y\in C$ we have $f(y) = 0$
            by \cref{completely_regular_space}.
        Let $g = \restrl{f}{Y}$.
        $g$ is continuous from $Y$ to $I_0$ with respect to $T_Y$ and $T_0$
            by \cref{continuous_restrict_domain}.
        We have $f\in\funs{X}{I_0}$ by \cref{continuous}.
        $f$ is right-unique and $f$ is a relation and $\dom{f} = X$
            by \cref{funs_is_function,function_rightunique,funs_is_relation,funs}.
        $Y\subseteq\dom{f}$ by \cref{subspace_of,subseteq}.
        $g(x) = f(x)$ by \cref{restrl_of_function_apply,subspace_of,subseteq}.
        Show that for all $y\in A$ we have $g(y) = 0$.
        \begin{subproof}
            Fix $y\in A$.
            $y\in C$ by \cref{subseteq}.
            $f(y) = 0$.
            $g(y) = f(y)$ by \cref{restrl_of_function_apply,subspace_of,subseteq,closed_in}.
            $g(y) = 0$.
        \end{subproof}
        There exists $h$ such that
            $h$ is continuous from $Y$ to $I_0$ with respect to $T_Y$ and $T_0$ and
            $h(x) = 1$ and
            for all $y\in A$ we have $h(y) = 0$.
    \end{subproof}
    Thus $Y$ is completely regular with respect to $T_Y$ by \cref{completely_regular_space}.
\end{proof}

% from §33 helper lemma 36: completely regular spaces are regular
\begin{lemma}\label{completely_regular_implies_regular}
    Assume $X$ is completely regular with respect to $T$.
    Then $X$ is regular with respect to $T$.
\end{lemma}
\begin{proof}
    Let $I_0 = \intervalclosed{0}{1}$.
    Let $T_0 = \subspaceTopology{\standardTopologyR}{I_0}$.
    $T$ is a topology on $X$ and $X$ has closed singletons with respect to $T$
        by \cref{completely_regular_space}.
    Show that for all $x,B$ such that $x\in X$ and $B$ is closed in $X$ with respect to $T$ and $x\notin B$
        there exist $U,V$ such that
            $U$ is open in $X$ with respect to $T$ and
            $V$ is open in $X$ with respect to $T$ and
            $x\in U$ and
            $B\subseteq V$ and
            $U$ is disjoint from $V$.
    \begin{subproof}
        Fix $x$.
        Fix $B$ such that $x\in X$ and $B$ is closed in $X$ with respect to $T$ and $x\notin B$.
        Take $f$ such that
            $f$ is continuous from $X$ to $I_0$ with respect to $T$ and $T_0$ and
            $f(x) = 1$ and
            for all $y\in B$ we have $f(y) = 0$
            by \cref{completely_regular_space}.
        Let $W_1 = \intervalopen{1 / (1 + 1)}{1 + 1}$.
        Let $W_2 = \intervalopen{\neg{1}}{1 / (1 + 1)}$.
        Let $U = \preimg{f}{I_0\inter W_1}$.
        Let $V = \preimg{f}{I_0\inter W_2}$.
        $0\in\reals$ and $1\in\reals$ and $1 + 1\in\reals$ and $0 < 1$
            by \cref{real_naturals_subset_reals,real_one_in_naturals,subseteq,real_add_closed,real_add_ident,real_zero_lt_one}.
        $0 < 1 + 1$ by \cref{real_lt_add,real_add_ident,real_lt_trans}.
        $1 < 1 + 1$ by \cref{real_lt_add,real_add_ident}.
        $(1 + 1)\neq 0$ and $\inv{(1 + 1)}\in\reals$
            by \cref{real_pos_nonzero,real_mul_inverse}.
        $0 < \inv{(1 + 1)}$ by \cref{real_inv_pos}.
        $0 \rmul \inv{(1 + 1)} < 1 \rmul \inv{(1 + 1)}$
            by \cref{real_lt_mul_pos}.
        $1 / (1 + 1)\in\reals$ by \cref{real_div,real_mul_closed}.
        $0 \rmul \inv{(1 + 1)} = 0$ and $1 \rmul \inv{(1 + 1)} = 1 / (1 + 1)$
            by \cref{real_mul_zero,real_mul_ident,real_div}.
        $0 < 1 / (1 + 1)$ by \cref{real_lt_subst_left,real_lt_subst_right}.
        We have $1 / (1 + 1) + 1 / (1 + 1) = 1$ by \cref{real_half_plus_half}.
        $1 / (1 + 1) < 1$ by \cref{real_lt_add,real_add_ident,real_lt_subst_left}.
        $1 / (1 + 1) < 1 + 1$ by \cref{real_lt_trans}.
        $\neg{1}\in\reals$ by \cref{real_add_inverse}.
        $\neg{1} < 0$ by \cref{real_lt_neg,real_neg_zero,real_lt_subst_right}.
        $\neg{1} < 1 / (1 + 1)$ by \cref{real_lt_trans}.
        $W_1\in\standardTopologyR$ and $W_2\in\standardTopologyR$
            by \cref{intervalopen_subset_reals,powerset_intro,standard_basis_reals,basis_elem_open_in_generated,basis_for_topology,standard_basis_is_basis,standard_topology_reals}.
        $T_0$ is a topology on $I_0$
            by \cref{intervalclosed,subseteq,standard_basis_is_basis,basis_topology_is_topology,standard_topology_reals,subspace_topology_is_topology}.
        $I_0\inter W_1\in T_0$ by \cref{subspace_topology}.
        $I_0\inter W_2\in T_0$ by \cref{subspace_topology}.
        $I_0\inter W_1$ is open in $I_0$ with respect to $T_0$ by \cref{open_in}.
        $I_0\inter W_2$ is open in $I_0$ with respect to $T_0$ by \cref{open_in}.
        $U$ is open in $X$ with respect to $T$ and $V$ is open in $X$ with respect to $T$
            by \cref{continuous}.
        $f\in\funs{X}{I_0}$ and $f$ is a function and $\dom{f} = X$
            by \cref{continuous,funs_is_function,funs}.
        $1\in W_1$ and $0\in W_2$
            by \cref{intervalopen,real_lt_neg,real_neg_zero,real_lt_subst_right}.
        $1\in I_0\inter W_1$ and $0\in I_0\inter W_2$ by \cref{intervalclosed,inter_intro}.
        $x\in\dom{f}$ by \cref{subseteq}.
        $(x,1)\in f$ by \cref{function_apply_elim}.
        $x\in U$ by \cref{preimg_iff}.
        Show that $B\subseteq V$.
        \begin{subproof}
            Show that for all $y$ such that $y\in B$ we have $y\in V$.
            \begin{subproof}
                Fix $y$ such that $y\in B$.
                We have $f(y) = 0$.
                $y\in X$ by \cref{closed_in,subseteq}.
                $y\in\dom{f}$ by \cref{subseteq}.
                $(y,0)\in f$ by \cref{function_apply_elim}.
                $y\in V$ by \cref{preimg_iff}.
            \end{subproof}
            $B\subseteq V$ by \cref{subseteq}.
        \end{subproof}
        Show that $U$ is disjoint from $V$.
        \begin{subproof}
            Suppose not.
            Take $y$ such that $y\in U$ and $y\in V$ by \cref{disjoint}.
            $f(y)\in I_0\inter W_1$ and $f(y)\in I_0\inter W_2$
                by \cref{preimg_iff,continuous,funs_is_function,function_apply_iff}.
            $f(y)\in W_1$ and $f(y)\in W_2$ by \cref{inter}.
            $f(y)\in\reals$ by \cref{inter_elim_left,intervalclosed}.
            $1 / (1 + 1) < f(y)$ and $f(y) < 1 / (1 + 1)$ by \cref{intervalopen}.
            Contradiction by \cref{real_lt_irreflexive,real_lt_trans}.
        \end{subproof}
    \end{subproof}
    Thus $X$ is regular with respect to $T$ by \cref{one_point_closed,regular_space}.
\end{proof}

% from §33 helper lemma 37: product of unit-interval valued continuous maps
\begin{lemma}\label{unit_interval_product_continuous}
    Let $I_0 = \intervalclosed{0}{1}$.
    Let $T_0 = \subspaceTopology{\standardTopologyR}{I_0}$.
    Assume $T$ is a topology on $X$.
    Assume $f$ is continuous from $X$ to $I_0$ with respect to $T$ and $T_0$.
    Assume $g$ is continuous from $X$ to $I_0$ with respect to $T$ and $T_0$.
    Let $h = \{(x, f(x) \rmul g(x)) \mid x\in X\}$.
    Then $h$ is continuous from $X$ to $I_0$ with respect to $T$ and $T_0$.
\end{lemma}
\begin{proof}
    $I_0\subseteq\reals$ by \cref{intervalclosed,subseteq}.
    $I_0$ is a subspace of $\reals$ with respect to $\standardTopologyR$
        by \cref{subspace_of,standard_basis_is_basis,basis_topology_is_topology,standard_topology_reals}.
    Let $j = \identity{I_0}$.
    Let $f_0 = j\circ f$.
    Let $g_0 = j\circ g$.
    $f\in\funs{X}{I_0}$ and $g\in\funs{X}{I_0}$ by \cref{continuous}.
    $f$ is right-unique and $f$ is a relation and $\dom{f} = X$ and
        $g$ is right-unique and $g$ is a relation and $\dom{g} = X$
        by \cref{funs_is_function,function_rightunique,funs_is_relation,funs}.
    $f\in\rels{X}{I_0}$ and $g\in\rels{X}{I_0}$ by \cref{funs_subseteq_rels,subseteq}.
    $\ran{f}\subseteq I_0$ and $\ran{g}\subseteq I_0$ by \cref{rels_ran_subseteq}.
    $\dom{j} = I_0$ and $j$ is right-unique and $j$ is a relation
        by \cref{id_dom,id_rightunique,id_is_relation}.
    $f_0$ is continuous from $X$ to $\reals$ with respect to $T$ and $\standardTopologyR$ and
        $g_0$ is continuous from $X$ to $\reals$ with respect to $T$ and $\standardTopologyR$
        by \cref{continuous_expand_range}.
    For all $x\in X$ we have $f_0(x) = f(x)$
        by \cref{circ_apply,funs_elim,id_apply}.
    For all $x\in X$ we have $g_0(x) = g(x)$
        by \cref{circ_apply,funs_elim,id_apply}.
    $h$ is continuous from $X$ to $\reals$ with respect to $T$ and $\standardTopologyR$
        by \cref{real_function_multiplication_continuous}.
    $h\in\funs{X}{\reals}$ and $h$ is a function by \cref{continuous,funs_is_function}.
    Show that for all $x\in X$ we have $h(x)\in I_0$.
    \begin{subproof}
        Fix $x\in X$.
        We have $f(x)\in I_0$ and $g(x)\in I_0$ by \cref{continuous,funs_type_apply}.
        $f(x)\in\reals$ and $g(x)\in\reals$ by \cref{intervalclosed}.
        $0 \leq f(x)$ and $f(x) \leq 1$ and $0 \leq g(x)$ and $g(x) \leq 1$
            by \cref{intervalclosed}.
        We have $(x, f(x) \rmul g(x))\in h$.
        $h(x) = f(x) \rmul g(x)$ by \cref{function_apply_intro}.
        $h(x)\in\reals$ by \cref{funs_type_apply}.
        $f(x) = 0$ or $0 < f(x)$ by \cref{real_lt_trichotomy}.
        $0 \leq f(x) \rmul g(x)$
            by \cref{real_mul_zero,real_mul_comm,real_lt_imp_leq,real_mul_pos_pos_gt}.
        $0 \leq h(x)$.
        Show that $h(x) \leq 1$.
        \begin{subproof}
            \begin{byCase}
            \caseOf{$f(x) = 0$.}
                $h(x) = 0$ by \cref{real_mul_zero,real_add_eq_subst}.
                $h(x) \leq 1$ by \cref{real_zero_lt_one,real_lt_imp_leq}.
            \caseOf{$0 < f(x)$.}
                Show that $h(x) < f(x)$ or $h(x) = f(x)$.
                \begin{subproof}
                    $g(x) = 1$ or $g(x) < 1$ by \cref{real_lt_trichotomy}.
                    \begin{byCase}
                    \caseOf{$g(x) = 1$.}
                        $h(x) = f(x) \rmul 1$ by \cref{real_add_eq_subst}.
                        $h(x) = f(x)$ by \cref{real_mul_comm,real_mul_ident}.
                    \caseOf{$g(x) < 1$.}
                        $f(x) \rmul g(x) < f(x) \rmul 1$
                            by \cref{real_lt_mul_pos,real_mul_comm,real_naturals_subset_reals,real_one_in_naturals,subseteq}.
                        $f(x) \rmul 1 = f(x)$ by \cref{real_mul_comm,real_mul_ident}.
                        $f(x) \rmul g(x) < f(x)$ by \cref{real_lt_subst_right}.
                        $h(x) < f(x)$ by \cref{real_lt_subst_left}.
                    \end{byCase}
                \end{subproof}
                $f(x) < 1$ or $f(x) = 1$ by \cref{intervalclosed}.
                $1\in\reals$ by \cref{real_mul_ident}.
                Show that $h(x) < 1$ or $h(x) = 1$.
                \begin{subproof}
                    \begin{byCase}
                    \caseOf{$h(x) < f(x)$.}
                        \begin{byCase}
                        \caseOf{$f(x) < 1$.}
                            $h(x) < 1$ by \cref{real_lt_trans}.
                        \caseOf{$f(x) = 1$.}
                            $h(x) < 1$ by \cref{real_lt_subst_right}.
                        \end{byCase}
                    \caseOf{$h(x) = f(x)$.}
                        \begin{byCase}
                        \caseOf{$f(x) < 1$.}
                            $h(x) < 1$ by \cref{real_lt_subst_left}.
                        \caseOf{$f(x) = 1$.}
                            $h(x) = 1$ by \cref{real_add_eq_subst}.
                        \end{byCase}
                    \end{byCase}
                \end{subproof}
            \end{byCase}
        \end{subproof}
        $h(x)\in I_0$ by \cref{intervalclosed}.
    \end{subproof}
    Hence $\img{h}{X}\subseteq I_0$
        by \cref{subseteq,img_iff,continuous,funs_is_function,function_apply_iff}.
    $h$ is continuous from $X$ to $I_0$ with respect to $T$ and $T_0$
        by \cref{continuous_subspace_range}.
\end{proof}

% from §33 helper definition 1: finite subbasic families admitting product maps
\begin{definition}\label{product_unit_interval_property}
    $n$ is a product unit interval index with respect to $P$ and $T_0$ and $T_I$ and $S$ iff
        for all $G,g,p_1$ such that
            $g$ is a bijection from $n$ to $G$ and
            $G\subseteq S$ and
            $G$ is inhabited and
            $p_1\in\inters{G}$
        there exists $h$ such that
            $h$ is continuous from $P$ to $\intervalclosed{0}{1}$ with respect to $T_0$ and $T_I$ and
            $h(p_1) = 1$ and
            for all $x\in P$ if $x\notin\inters{G}$ then $h(x) = 0$.
\end{definition}

% from §33 helper lemma 38: packaging a prepared continuous witness
\begin{lemma}\label{prepared_product_unit_interval_witness}
    Assume $f$ is continuous from $P$ to $I_0$ with respect to $T_0$ and $T_I$.
    Assume $f(p_1) = 1$.
    Assume for all $x\in P$ if $x\notin A$ then $f(x) = 0$.
    Then there exists $h$ such that
        $h$ is continuous from $P$ to $I_0$ with respect to $T_0$ and $T_I$ and
        $h(p_1) = 1$ and
        for all $x\in P$ if $x\notin A$ then $h(x) = 0$.
\end{lemma}
\begin{proof}
    Take $h$ such that $h = f$.
    $h$ is continuous from $P$ to $I_0$ with respect to $T_0$ and $T_I$.
    $h(p_1) = 1$.
    For all $x\in P$ if $x\notin A$ then $h(x) = 0$ by assumption.
\end{proof}

% from §33 helper lemma 38a: singleton family witness
\begin{lemma}\label{singleton_product_unit_interval_witness}
    Assume $\phi$ is continuous from $P$ to $I_0$ with respect to $T_0$ and $T_I$.
    Assume $\phi(p_1) = 1$.
    Assume for all $x\in P$ if $x\notin A_0$ then $\phi(x) = 0$.
    Assume $G = \{A_0\}$.
    Then there exists $h$ such that
        $h$ is continuous from $P$ to $I_0$ with respect to $T_0$ and $T_I$ and
        $h(p_1) = 1$ and
        for all $x\in P$ if $x\notin\inters{G}$ then $h(x) = 0$.
\end{lemma}
\begin{proof}
    Follows by \cref{inters_singleton,subseteq,prepared_product_unit_interval_witness}.
\end{proof}

% from §33 helper lemma 39: finite products of unit-interval valued continuous maps
\begin{lemma}\label{finite_unit_interval_product_continuous}
    Let $P = \productFamily{X}$.
    Let $T_0 = \productTopologyIndexed{T}{X}$.
    Let $S = \productSubbasis{T}{X}$.
    Let $I_0 = \intervalclosed{0}{1}$.
    Let $T_I = \subspaceTopology{\standardTopologyR}{I_0}$.
    Assume $X$ is a function.
    Assume $T$ is a function.
    Assume $\dom{T} = \dom{X}$.
    Assume $\dom{X}$ is inhabited.
    Suppose for all $\alpha\in\dom{X}$ we have $\apply{X}{\alpha}$ is completely regular with respect to $\apply{T}{\alpha}$.
    Suppose $F\subseteq S$.
    Suppose $F$ is finite and inhabited.
    Assume $p\in\inters{F}$.
    Then there exists $h$ such that
        $h$ is continuous from $P$ to $I_0$ with respect to $T_0$ and $T_I$ and
        $h(p) = 1$ and
        for all $x\in P$ if $x\notin\inters{F}$ then $h(x) = 0$.
\end{lemma}
\begin{proof}
    Let $A = \{ n\in\naturalsPlus \mid \text{$\initialSegment{n}$ is a product unit interval index with respect to $P$ and $T_0$ and $T_I$ and $S$} \}$.
    Show that for all $G,g,p_1$ such that
        $g$ is a bijection from $\initialSegment{1}$ to $G$ and
        $G\subseteq S$ and
        $G$ is inhabited and
        $p_1\in\inters{G}$
        there exists $h$ such that
            $h$ is continuous from $P$ to $I_0$ with respect to $T_0$ and $T_I$ and
            $h(p_1) = 1$ and
            for all $x\in P$ if $x\notin\inters{G}$ then $h(x) = 0$.
    \begin{subproof}
        Fix $G$.
        Fix $g$.
        Fix $p_1$ such that
            $g$ is a bijection from $\initialSegment{1}$ to $G$ and
            $G\subseteq S$ and
            $G$ is inhabited and
            $p_1\in\inters{G}$.
        $g$ is a function and $\dom{g} = \initialSegment{1}$ by \cref{bijection_is_function,bijections_dom}.
        $1\in\dom{g}$ by \cref{initial_segment_one,singleton_intro}.
        Let $A_0 = g(1)$.
        $(1,A_0)\in g$ by \cref{function_apply_elim}.
        $g\in\surj{\initialSegment{1}}{G}$ by \cref{bijections}.
        $\ran{g} = G$ by \cref{ran_of_surj}.
        $A_0\in G$ by \cref{ran_iff}.
        Show that for all $z\in G$ we have $z\in\{A_0\}$.
        \begin{subproof}
            Fix $z\in G$.
            Take $x\in\initialSegment{1}$ such that $g(x)=z$ by \cref{bijections,surj}.
            $x = 1$ by \cref{initial_segment_one,singleton_iff}.
            $A_0 = z$.
            Thus $z\in\{A_0\}$ by \cref{singleton_iff}.
        \end{subproof}
        $G\subseteq\{A_0\}$ by \cref{subseteq}.
        $\{A_0\}\subseteq G$ by \cref{singleton_subset_intro}.
        $G = \{A_0\}$ by \cref{subseteq_antisymmetric}.
        Take $\beta\in\dom{X}$ such that $A_0\in\productSubbasisAt{T}{X}{\beta}$
            by \cref{subseteq,product_subbasis_union_indexed}.
        Take $U_0\in\apply{T}{\beta}$ such that $A_0 = \preimg{\piIndex{X}{\beta}}{U_0}$ by \cref{product_subbasis_indexed}.
        $\piIndex{X}{\beta}\in\funs{P}{\apply{X}{\beta}}$ by \cref{projection_map_indexed_function}.
        $\ran{\piIndex{X}{\beta}}\subseteq \apply{X}{\beta}$ by \cref{funs_ran}.
        $A_0\subseteq P$ by \cref{funs_elim,preimg_subseteq_dom,subseteq}.
        $p_1\in A_0$ by \cref{inters_iff_forall}.
        $\apply{p_1}{\beta}\in\apply{X}{\beta}$ by \cref{subseteq,indexed_product}.
        Take $y\in U_0$ such that $p_1\mathrel{\piIndex{X}{\beta}} y$ by \cref{preimg_iff}.
        Take $z\in P$ such that $(p_1,y) = (z,\apply{z}{\beta})$ by \cref{projection_map_indexed}.
        $\apply{p_1}{\beta} = y$ by \cref{pair_eq_iff}.
        $\apply{p_1}{\beta}\in U_0$.
        For all $\alpha\in\dom{X}$ we have $\apply{T}{\alpha}$ is a topology on $\apply{X}{\alpha}$
            by \cref{completely_regular_space}.
        Let $C = \apply{X}{\beta}\setminus U_0$.
        $C$ is closed in $\apply{X}{\beta}$ with respect to $\apply{T}{\beta}$
            by \cref{open_in,topology_on,subseteq,pow_iff,setminus_subseteq,setminus_setminus,inter_absorb_supseteq_left,closed_in}.
        $\apply{p_1}{\beta}\notin C$ by \cref{setminus_elim_right}.
        Take $f$ such that
            $f$ is continuous from $\apply{X}{\beta}$ to $I_0$ with respect to $\apply{T}{\beta}$ and $T_I$ and
            $f(\apply{p_1}{\beta}) = 1$ and
            for all $u\in C$ we have $f(u) = 0$
            by \cref{completely_regular_space}.
        $f\in\funs{\apply{X}{\beta}}{I_0}$ and $f$ is a function and $\dom{f} = \apply{X}{\beta}$
            by \cref{continuous,funs_is_function,funs}.
        $\ran{\piIndex{X}{\beta}}\subseteq\dom{f}$ by \cref{funs,subseteq}.
        Let $\phi = f\circ\piIndex{X}{\beta}$.
        $\piIndex{X}{\beta}$ is continuous from $P$ to $\apply{X}{\beta}$ with respect to $T_0$ and $\apply{T}{\beta}$
            by \cref{projection_map_indexed_continuous}.
        $\phi$ is continuous from $P$ to $I_0$ with respect to $T_0$ and $T_I$
            by \cref{continuous_composite}.
        $\phi$ is a function by \cref{continuous,funs_is_function}.
        $\apply{p_1}{\beta}\in\dom{f}$ by \cref{funs,subseteq}.
        $(\apply{p_1}{\beta}, f(\apply{p_1}{\beta}))\in f$ by \cref{function_apply_elim}.
        $(p_1,\apply{p_1}{\beta})\in\piIndex{X}{\beta}$ by \cref{projection_map_indexed}.
        $(p_1, f(\apply{p_1}{\beta}))\in\phi$ by \cref{circ_iff}.
        $\phi(p_1) = f(\apply{p_1}{\beta})$ by \cref{function_apply_intro}.
        $\phi(p_1) = 1$.
        For all $x\in P$ if $x\notin A_0$ then $\phi(x) = 0$
            by \cref{projection_map_indexed,preimg_iff,indexed_product,setminus_intro,continuous,funs_is_function,funs,subseteq,function_apply_elim,circ_iff,function_apply_intro}.
        Show that there exists $h$ such that
            $h$ is continuous from $P$ to $I_0$ with respect to $T_0$ and $T_I$ and
            $h(p_1) = 1$ and
            for all $x\in P$ if $x\notin\inters{G}$ then $h(x) = 0$.
        \begin{subproof}
            Follows by \cref{singleton_product_unit_interval_witness}.
        \end{subproof}
    \end{subproof}
    $\initialSegment{1}$ is a product unit interval index with respect to $P$ and $T_0$ and $T_I$ and $S$
        by \cref{product_unit_interval_property}.
    $1\in\naturalsPlus$ by \cref{real_one_in_naturals}.
    $1\in A$ by definition.
    Show that for all $n\in A$ we have $n + 1\in A$.
    \begin{subproof}
        Fix $n\in A$.
        $n\in\naturalsPlus$ by definition.
        Show that for all $G,g,p_1$ such that
            $g$ is a bijection from $\initialSegment{n + 1}$ to $G$ and
            $G\subseteq S$ and
            $G$ is inhabited and
            $p_1\in\inters{G}$
            there exists $h$ such that
                $h$ is continuous from $P$ to $I_0$ with respect to $T_0$ and $T_I$ and
                $h(p_1) = 1$ and
                for all $x\in P$ if $x\notin\inters{G}$ then $h(x) = 0$.
        \begin{subproof}
            Fix $G$.
            Fix $g$.
            Fix $p_1$ such that
                $g$ is a bijection from $\initialSegment{n + 1}$ to $G$ and
                $G\subseteq S$ and
                $G$ is inhabited and
                $p_1\in\inters{G}$.
            $g$ is a function and $\dom{g} = \initialSegment{n + 1}$ by \cref{bijection_is_function,bijections_dom}.
            $n + 1\in\dom{g}$ by \cref{initial_segment_naturalsplus,real_naturals_closed_succ}.
            Let $A_0 = g(n + 1)$.
            Let $G_1 = \img{g}{\initialSegment{n}}$.
            $\initialSegment{n}\subseteq \initialSegment{n + 1}$ by \cref{initial_segment_subset_succ}.
            We have $G_1\subseteq G$
                by \cref{subseteq,img_subseteq,img_dom,bijections,ran_of_surj}.
            $G_1\subseteq S$ by \cref{subseteq_transitive}.
            $A_0\in G$ by \cref{function_apply_elim,ran_iff,bijections,ran_of_surj}.
            For all $z$ such that $z\in G$ we have $z\in G_1\union\{A_0\}$
                by \cref{bijections,surj,initial_segment_succ_split,function_apply_elim,img_iff,union_intro_left,union_intro_right,singleton_intro}.
            $G\subseteq G_1\union\{A_0\}$ by \cref{subseteq}.
            $G_1\union\{A_0\}\subseteq G$ by \cref{singleton_subset_intro,union_subsets_is_subset}.
            Hence $G = G_1\union\{A_0\}$ by \cref{subseteq_antisymmetric}.
            Take $\beta\in\dom{X}$ such that $A_0\in\productSubbasisAt{T}{X}{\beta}$
                by \cref{subseteq,product_subbasis_union_indexed}.
            Take $U_0\in\apply{T}{\beta}$ such that $A_0 = \preimg{\piIndex{X}{\beta}}{U_0}$ by \cref{product_subbasis_indexed}.
            $\piIndex{X}{\beta}\in\funs{P}{\apply{X}{\beta}}$ by \cref{projection_map_indexed_function}.
            $\ran{\piIndex{X}{\beta}}\subseteq \apply{X}{\beta}$ by \cref{funs_ran}.
            $A_0\subseteq P$ by \cref{funs_elim,preimg_subseteq_dom,subseteq}.
            $p_1\in A_0$ by \cref{inters_iff_forall}.
            $\apply{p_1}{\beta}\in\apply{X}{\beta}$ by \cref{subseteq,indexed_product}.
            Take $y\in U_0$ such that $p_1\mathrel{\piIndex{X}{\beta}} y$ by \cref{preimg_iff}.
            Take $z\in P$ such that $(p_1,y) = (z,\apply{z}{\beta})$ by \cref{projection_map_indexed}.
            $\apply{p_1}{\beta} = y$ by \cref{pair_eq_iff}.
            $\apply{p_1}{\beta}\in U_0$.
            For all $\alpha\in\dom{X}$ we have $\apply{T}{\alpha}$ is a topology on $\apply{X}{\alpha}$
                by \cref{completely_regular_space}.
            Let $C = \apply{X}{\beta}\setminus U_0$.
            $C$ is closed in $\apply{X}{\beta}$ with respect to $\apply{T}{\beta}$
                by \cref{open_in,topology_on,subseteq,pow_iff,setminus_subseteq,setminus_setminus,inter_absorb_supseteq_left,closed_in}.
            $\apply{p_1}{\beta}\notin C$ by \cref{setminus_elim_right}.
            Take $f$ such that
                $f$ is continuous from $\apply{X}{\beta}$ to $I_0$ with respect to $\apply{T}{\beta}$ and $T_I$ and
                $f(\apply{p_1}{\beta}) = 1$ and
                for all $u\in C$ we have $f(u) = 0$
                by \cref{completely_regular_space}.
            $f\in\funs{\apply{X}{\beta}}{I_0}$ and $f$ is a function and $\dom{f} = \apply{X}{\beta}$
                by \cref{continuous,funs_is_function,funs}.
            $\ran{\piIndex{X}{\beta}}\subseteq\dom{f}$ by \cref{funs,subseteq}.
            Let $\phi = f\circ\piIndex{X}{\beta}$.
            $\piIndex{X}{\beta}$ is continuous from $P$ to $\apply{X}{\beta}$ with respect to $T_0$ and $\apply{T}{\beta}$
                by \cref{projection_map_indexed_continuous}.
            $\phi$ is continuous from $P$ to $I_0$ with respect to $T_0$ and $T_I$
                by \cref{continuous_composite}.
            $\phi$ is a function by \cref{continuous,funs_is_function}.
            $\apply{p_1}{\beta}\in\dom{f}$ by \cref{funs,subseteq}.
            $(\apply{p_1}{\beta}, f(\apply{p_1}{\beta}))\in f$ by \cref{function_apply_elim}.
            $(p_1,\apply{p_1}{\beta})\in\piIndex{X}{\beta}$ by \cref{projection_map_indexed}.
            $(p_1, f(\apply{p_1}{\beta}))\in\phi$ by \cref{circ_iff}.
            $\phi(p_1) = f(\apply{p_1}{\beta})$ by \cref{function_apply_intro}.
            $\phi(p_1) = 1$.
            For all $x\in P$ if $x\notin A_0$ then $\phi(x) = 0$
                by \cref{projection_map_indexed,preimg_iff,indexed_product,setminus_intro,continuous,funs_is_function,funs,subseteq,function_apply_elim,circ_iff,function_apply_intro}.
            Take $x_0$ such that $x_0\in\initialSegment{n}$ by \cref{initial_segment_inhabited}.
            $x_0\in\dom{g}$ by \cref{initial_segment_subset_succ,bijections_dom,subseteq}.
            $g(x_0)\in G_1$ by \cref{function_apply_elim,img_iff}.
            $G_1$ is inhabited.
            For all $B$ such that $B\in G_1$ we have $p_1\in B$
                by \cref{img_iff,bijections,function_apply_intro,function_apply_elim,ran_iff,ran_of_surj,inters_iff_forall}.
            $p_1\in\inters{G_1}$ by \cref{inters_iff_forall}.
            Let $q = \restrl{g}{\initialSegment{n}}$.
            $q$ is a function by \cref{restrl_of_function_is_function}.
            For all $x\in\initialSegment{n}$ we have $x\in\dom{g}$ by \cref{initial_segment_subset_succ,bijections_dom,subseteq}.
            $\initialSegment{n}\subseteq\dom{g}$ by \cref{subseteq}.
            $\dom{q} = \initialSegment{n}$ by \cref{subseteq,dom_of_restrl_eq_inter,inter_absorb_supseteq_left}.
            $q\in\funs{\initialSegment{n}}{G_1}$ by \cref{function_apply_elim,restrl_subseteq,elem_subseteq,img_iff,funs_intro}.
            For all $w\in G_1$ there exists $x\in\initialSegment{n}$ such that $q(x) = w$
                by \cref{img_iff,function_apply_intro,restrl_of_function_apply}.
            Thus $q$ is a bijection from $\initialSegment{n}$ to $G_1$ by \cref{surj,restrl_of_function_apply,bijections,injective_function}.
            $\initialSegment{n}$ is a product unit interval index with respect to $P$ and $T_0$ and $T_I$ and $S$ by definition.
            Take $h_0$ such that
                $h_0$ is continuous from $P$ to $I_0$ with respect to $T_0$ and $T_I$ and
                $h_0(p_1) = 1$ and
                for all $x\in P$ if $x\notin\inters{G_1}$ then $h_0(x) = 0$
                by \cref{subseteq,product_unit_interval_property}.
            Let $h = \{(x, h_0(x) \rmul \phi(x)) \mid x\in P\}$.
            For all $\alpha\in\dom{X}$ we have $\apply{T}{\alpha}$ is a topology on $\apply{X}{\alpha}$
                by \cref{completely_regular_space}.
            $T_0$ is a topology on $P$ by
                \cref{product_subbasis_finite_intersections_basis,product_topology_indexed,basis_topology_is_topology,basis_for_topology}.
            $h$ is continuous from $P$ to $I_0$ with respect to $T_0$ and $T_I$
                by \cref{unit_interval_product_continuous}.
            $h$ is a function by \cref{continuous,funs_is_function}.
            $h_0\in\funs{P}{I_0}$ and $\phi\in\funs{P}{I_0}$ by \cref{continuous}.
            We have $(p_1, h_0(p_1) \rmul \phi(p_1))\in h$.
            $h(p_1) = h_0(p_1) \rmul \phi(p_1)$ by \cref{function_apply_intro}.
            $h(p_1) = 1$ by \cref{real_mul_ident}.
            Show that for all $x\in P$ if $x\notin\inters{G}$ then $h(x) = 0$.
            \begin{subproof}
                Fix $x\in P$.
                Assume $x\notin\inters{G}$.
                $h$ is a function by \cref{continuous,funs_is_function}.
                $h_0\in\funs{P}{I_0}$ and $\phi\in\funs{P}{I_0}$ by \cref{continuous}.
                We have $(x, h_0(x) \rmul \phi(x))\in h$.
                $h(x) = h_0(x) \rmul \phi(x)$ by \cref{function_apply_intro}.
                We have $h_0(x)\in\reals$ by \cref{funs_type_apply,intervalclosed,subseteq}.
                $\phi(x)\in\reals$ by \cref{funs_type_apply,intervalclosed,subseteq}.
                \begin{byCase}
                \caseOf{$x\in A_0$.}
                    Show that $x\notin\inters{G_1}$.
                    \begin{subproof}
                        Suppose not.
                        For all $B$ such that $B\in G$ we have $x\in B$
                            by \cref{subseteq,union_iff,singleton_iff,inters_iff_forall}.
                        $x\in\inters{G}$ by \cref{inters_iff_forall}.
                        Contradiction.
                    \end{subproof}
                    $h_0(x) = 0$.
                    $h(x) = 0$ by \cref{real_mul_zero}.
                \caseOf{$x\notin A_0$.}
                    $\phi(x) = 0$.
                    $h(x) = 0$ by \cref{real_mul_comm,real_mul_zero}.
                \end{byCase}
            \end{subproof}
            Follows by definition.
        \end{subproof}
        $\initialSegment{n + 1}$ is a product unit interval index with respect to $P$ and $T_0$ and $T_I$ and $S$
            by \cref{product_unit_interval_property}.
        $n + 1\in\naturalsPlus$ by \cref{real_naturals_closed_succ}.
        $n + 1\in A$ by definition.
    \end{subproof}
    $A\subseteq\naturalsPlus$ by \cref{subseteq}.
    Take $k\in\naturalsPlus$ such that there exists a bijection from $\initialSegment{k}$ to $F$ by \cref{finite}.
    Take $g$ such that $g$ is a bijection from $\initialSegment{k}$ to $F$.
    $k\in A$ by \cref{real_naturals_induction}.
    $\initialSegment{k}$ is a product unit interval index with respect to $P$ and $T_0$ and $T_I$ and $S$ by definition.
    Therefore there exists $h$ such that
        $h$ is continuous from $P$ to $I_0$ with respect to $T_0$ and $T_I$ and
        $h(p) = 1$ and
        for all $x\in P$ if $x\notin\inters{F}$ then $h(x) = 0$
        by \cref{product_unit_interval_property}.
\end{proof}

% from §33 Item 5: product of completely regular spaces is completely regular
\begin{theorem}\label{product_completely_regular}
    Assume $X$ is a function.
    Assume $T$ is a function.
    Assume $\dom{T} = \dom{X}$.
    Assume $\dom{X}$ is inhabited.
    Suppose for all $\alpha\in\dom{X}$ we have $\apply{X}{\alpha}$ is completely regular with respect to $\apply{T}{\alpha}$.
    Then $\productFamily{X}$ is completely regular with respect to $\productTopologyIndexed{T}{X}$.
\end{theorem}
\begin{proof}
    Let $P = \productFamily{X}$.
    Let $T_0 = \productTopologyIndexed{T}{X}$.
    Let $S = \productSubbasis{T}{X}$.
    Let $I_0 = \intervalclosed{0}{1}$.
    Let $T_I = \subspaceTopology{\standardTopologyR}{I_0}$.
    For all $\alpha\in\dom{X}$ we have $\apply{X}{\alpha}$ is regular with respect to $\apply{T}{\alpha}$
        by \cref{completely_regular_implies_regular}.
    Hence $P$ is regular with respect to $T_0$ by \cref{product_regular}.
    Therefore $T_0$ is a topology on $P$ and $P$ has closed singletons with respect to $T_0$
        by \cref{regular_space}.
    Show that for all $x,A$ such that $x\in P$ and $A$ is closed in $P$ with respect to $T_0$ and $x\notin A$
        there exists $f$ such that
            $f$ is continuous from $P$ to $I_0$ with respect to $T_0$ and $T_I$ and
            $f(x) = 1$ and
            for all $y\in A$ we have $f(y) = 0$.
    \begin{subproof}
        Fix $x$.
        Fix $A$ such that $x\in P$ and $A$ is closed in $P$ with respect to $T_0$ and $x\notin A$.
        Let $U = P\setminus A$.
        Hence $U$ is open in $P$ with respect to $T_0$ and $x\in U$ by \cref{closed_in,setminus_intro}.
        For all $\alpha\in\dom{X}$ we have $\apply{T}{\alpha}$ is a topology on $\apply{X}{\alpha}$
            by \cref{completely_regular_space}.
        Therefore $U\in\basisTopology{\finiteIntersections{S}}{P}$
            by \cref{open_in,product_subbasis_finite_intersections_basis,product_topology_indexed,topology_generated_by_subbasis,basis_for_topology,completely_regular_space}.
        Take $B\in\finiteIntersections{S}$ such that $x\in B$ and $B\subseteq U$ by \cref{topology_generated_by_basis}.
        Take $F\subseteq S$ such that $F$ is finite and inhabited and $B = \inters{F}$ by \cref{finite_intersections}.
        Take $h$ such that
            $h$ is continuous from $P$ to $I_0$ with respect to $T_0$ and $T_I$ and
            $h(x) = 1$ and
            for all $y\in P$ if $y\notin\inters{F}$ then $h(y) = 0$
            by \cref{finite_unit_interval_product_continuous}.
        Show that for all $y\in A$ we have $h(y) = 0$.
        \begin{subproof}
            Fix $y\in A$.
            $A\subseteq P$ by \cref{closed_in}.
            $y\in P$ by \cref{subseteq}.
            $y\notin U$ by \cref{setminus}.
            $y\notin B$ by \cref{subseteq}.
            $y\notin\inters{F}$.
            $h(y) = 0$.
        \end{subproof}
        Follows by definition.
    \end{subproof}
    Hence $P$ is completely regular with respect to $T_0$ by \cref{completely_regular_space}.
\end{proof}
\section{§34 The Urysohn Metrization Theorem}

% from §34 helper lemma 2: subspaces of metrizable spaces are metrizable
\begin{lemma}\label{subspace_metrizable}
    Assume $X$ is metrizable with respect to $T$.
    Assume $Y$ is a subspace of $X$ with respect to $T$.
    Let $T_Y = \subspaceTopology{T}{Y}$.
    Then $Y$ is metrizable with respect to $T_Y$.
\end{lemma}
\begin{proof}
    $T$ is a topology on $X$ and $Y\subseteq X$ by \cref{metrizable,subspace_of}.
    Take $d$ such that $d$ is a metric on $X$ and $T = \metricTopology{d}{X}$ by \cref{metrizable}.
    Let $e = \restrl{d}{Y\times Y}$.

    Show that $e$ is a metric on $Y$.
    \begin{subproof}
        $d\in\funs{X\times X}{\reals}$ and $d$ is a function and $d$ is a relation and $\dom{d} = X\times X$
            by \cref{metric_on,funs_elim,funs_is_relation}.
        Hence $Y\times Y\subseteq X\times X$ by \cref{times_subseteq_left,times_subseteq_right,subseteq_transitive}.
        $e$ is a function by \cref{restrl_of_function_is_function}.
        $\dom{e} = (X\times X)\inter(Y\times Y)$ by \cref{dom_of_restrl_eq_inter,pair_eq_iff}.
        $\dom{e} = Y\times Y$ by \cref{subseteq,inter_intro,restrl_dom,subseteq_antisymmetric}.
        For all $p\in\dom{e}$ we have $e(p)\in\reals$
            by \cref{inter_elim_left,dom_of_restrl_eq_inter,funs_type_apply,restrl_of_function_apply}.
        Hence $e\in\funs{Y\times Y}{\reals}$ by \cref{funs_intro}.

        For all $x,y$ such that $x\in Y$ and $y\in Y$ we have $e((x,y)) = d((x,y))$
            by \cref{times_tuple_intro,subseteq,restrl_of_function_apply}.

        Show that for all $x,y$ such that $x\in Y$ and $y\in Y$ we have $e((x,y)) \geq 0$.
        \begin{subproof}
            Follows by \cref{subseteq,metric_on,metric_nonnegative}.
        \end{subproof}
        Hence $e$ is nonnegative on $Y$ by \cref{metric_nonnegative}.

        Show that for all $x,y$ such that $x\in Y$ and $y\in Y$ we have $e((x,y)) = 0$ iff $x = y$.
        \begin{subproof}
            Follows by \cref{subseteq,metric_on,metric_zero_characterization}.
        \end{subproof}
        Thus $e$ is zero-characterized on $Y$ by \cref{metric_zero_characterization}.

        Show that for all $x,y$ such that $x\in Y$ and $y\in Y$ we have $e((x,y)) = e((y,x))$.
        \begin{subproof}
            Follows by \cref{subseteq,metric_on,metric_symmetric}.
        \end{subproof}
        Hence $e$ is symmetric on $Y$ by \cref{metric_symmetric}.

        Show that for all $x,y,z$ such that $x\in Y$ and $y\in Y$ and $z\in Y$
            we have $e((x,y)) + e((y,z)) \geq e((x,z))$.
        \begin{subproof}
            Follows by \cref{subseteq,metric_on,metric_triangle_inequality}.
        \end{subproof}
        Thus $e$ satisfies the triangle inequality on $Y$ by \cref{metric_triangle_inequality}.
        Therefore $e$ is a metric on $Y$ by \cref{metric_on}.
    \end{subproof}

    Let $T_e = \metricTopology{e}{Y}$.
    $T_Y$ is a topology on $Y$ and $T_e$ is a topology on $Y$ and
        $e\in\funs{Y\times Y}{\reals}$ and $\dom{e} = Y\times Y$
        by \cref{subspace_topology_is_topology,metric_basis_is_basis,basis_topology_is_topology,metric_topology,metric_on,funs_elim}.
    $d\in\funs{X\times X}{\reals}$ and $d$ is a function and $d$ is a relation and $\dom{d} = X\times X$
        by \cref{metric_on,funs_elim,funs_is_relation}.
    Hence $Y\times Y\subseteq\dom{d}$ by \cref{times_subseteq_left,times_subseteq_right,subseteq_transitive,pair_eq_iff}.

    Show that for all $U$ such that $U$ is open in $Y$ with respect to $T_Y$ we have $U$ is open in $Y$ with respect to $T_e$.
    \begin{subproof}
        Fix $U$ such that $U$ is open in $Y$ with respect to $T_Y$.
        Hence $U\subseteq Y$ by \cref{open_in,topology_on,subseteq,pow_iff}.
        Take $V\in T$ such that $U = Y\inter V$ by \cref{subspace_topology,open_in}.
        $V$ is open in $X$ with respect to $T$ by \cref{open_in}.
        $V\subseteq X$ by \cref{open_in,topology_on,subseteq,pow_iff}.
        Show that for all $y\in U$ there exists $\delta\in\reals$ such that $0 < \delta$ and $\metricBall{e}{Y}{y}{\delta}\subseteq U$.
        \begin{subproof}
            Fix $y\in U$.
            Hence $y\in Y$ and $y\in V$ by \cref{inter}.
            Take $\delta\in\reals$ such that $0 < \delta$ and $\metricBall{d}{X}{y}{\delta}\subseteq V$
                by \cref{metric_open_iff}.
            Show that for all $z$ such that $z\in\metricBall{e}{Y}{y}{\delta}$ we have $z\in U$.
            \begin{subproof}
                Fix $z$ such that $z\in\metricBall{e}{Y}{y}{\delta}$.
                Hence $z\in Y$ and $e((y,z)) < \delta$ by \cref{metric_ball}.
                $e((y,z)) = d((y,z))$.
                $d((y,z)) < \delta$ by \cref{real_lt_subst_left}.
                $z\in\metricBall{d}{X}{y}{\delta}$ by \cref{metric_ball,subseteq}.
                Thus $z\in U$ by \cref{subseteq,inter_intro}.
            \end{subproof}
            Hence $\metricBall{e}{Y}{y}{\delta}\subseteq U$ by \cref{subseteq}.
        \end{subproof}
        Hence $U$ is open in $Y$ with respect to $T_e$ by \cref{metric_open_iff,open_in}.
    \end{subproof}

    Show that for all $U$ such that $U$ is open in $Y$ with respect to $T_e$ we have $U$ is open in $Y$ with respect to $T_Y$.
    \begin{subproof}
        Fix $U$ such that $U$ is open in $Y$ with respect to $T_e$.
        Hence $U\subseteq Y$ by \cref{open_in,topology_on,subseteq,pow_iff}.
        Let $B = \metricBasis{e}{Y}$.
        $B$ is a basis on $Y$ by \cref{metric_basis_is_basis}.
        Show that for all $W\in B$ we have $W$ is open in $Y$ with respect to $T_Y$.
            \begin{subproof}
                Fix $W\in B$.
                Take $y\in Y$ such that there exists $\delta\in\reals$ such that $0 < \delta$ and $W = \metricBall{e}{Y}{y}{\delta}$
                    by \cref{metric_basis}.
                Take $\delta\in\reals$ such that $0 < \delta$ and $W = \metricBall{e}{Y}{y}{\delta}$ by \cref{metric_basis}.
            Hence $y\in X$ by \cref{subseteq}.
            $\metricBall{d}{X}{y}{\delta}\subseteq X$ by \cref{metric_ball,subseteq}.
            Show that for all $z$ such that $z\in W$ we have $z\in Y\inter\metricBall{d}{X}{y}{\delta}$.
            \begin{subproof}
                Fix $z$ such that $z\in W$.
                Hence $z\in\metricBall{e}{Y}{y}{\delta}$ by \cref{pair_eq_iff}.
                $z\in Y$ and $e((y,z)) < \delta$ by \cref{metric_ball}.
                $e((y,z)) = d((y,z))$.
                $d((y,z)) < \delta$ by \cref{real_lt_subst_left}.
                $z\in\metricBall{d}{X}{y}{\delta}$ by \cref{metric_ball,subseteq}.
                Thus $z\in Y\inter\metricBall{d}{X}{y}{\delta}$ by \cref{inter_intro}.
            \end{subproof}
            Hence $W\subseteq Y\inter\metricBall{d}{X}{y}{\delta}$ by \cref{subseteq}.
            Show that for all $z$ such that $z\in Y\inter\metricBall{d}{X}{y}{\delta}$ we have $z\in W$.
            \begin{subproof}
                Fix $z$ such that $z\in Y\inter\metricBall{d}{X}{y}{\delta}$.
                Hence $z\in Y$ and $z\in\metricBall{d}{X}{y}{\delta}$ by \cref{inter}.
                $d((y,z)) < \delta$ by \cref{metric_ball}.
                $e((y,z)) = d((y,z))$.
                $e((y,z)) < \delta$ by \cref{real_lt_subst_left}.
                $z\in\metricBall{e}{Y}{y}{\delta}$ by \cref{metric_ball}.
                $z\in W$ by \cref{pair_eq_iff}.
            \end{subproof}
            Hence $Y\inter\metricBall{d}{X}{y}{\delta}\subseteq W$ by \cref{subseteq}.
            Therefore $W = Y\inter\metricBall{d}{X}{y}{\delta}$ by \cref{subseteq_antisymmetric}.
            $\metricBall{d}{X}{y}{\delta}\in\metricBasis{d}{X}$ by \cref{metric_ball,powerset_intro,metric_basis}.
            Hence $\metricBall{d}{X}{y}{\delta}$ is open in $X$ with respect to $T$
                by \cref{metric_topology,metric_basis_is_basis,basis_elem_open_in_generated,open_in}.
            Therefore $W$ is open in $Y$ with respect to $T_Y$ by \cref{subspace_topology,open_in}.
        \end{subproof}

        Let $M = \{W\in B \mid W\subseteq U\}$.
        Show that for all $y$ such that $y\in U$ we have $y\in\unions{M}$.
        \begin{subproof}
            Fix $y$ such that $y\in U$.
            Take $\delta\in\reals$ such that $0 < \delta$ and $\metricBall{e}{Y}{y}{\delta}\subseteq U$
                by \cref{metric_open_iff}.
            Let $W = \metricBall{e}{Y}{y}{\delta}$.
            Hence $y\in Y$ by \cref{subseteq}.
            $W\in B$ by \cref{metric_ball,subseteq,powerset_intro,metric_basis,pair_eq_iff}.
            Therefore $y\in\unions{M}$
                by \cref{metric_ball,real_lt_subst_left,metric_on,metric_zero_characterization,pair_eq_iff,unions_intro}.
        \end{subproof}
        Hence $U\subseteq\unions{M}$ by \cref{subseteq}.
        Show that for all $y$ such that $y\in\unions{M}$ we have $y\in U$.
        \begin{subproof}
            Trivial.
        \end{subproof}
        Hence $\unions{M}\subseteq U$ by \cref{subseteq}.
        Therefore $U = \unions{M}$ by \cref{subseteq_antisymmetric}.
        $M\subseteq T_Y$ by \cref{pair_eq_iff,subseteq,open_in}.
        $\unions{M}\in T_Y$ by \cref{topology_on,subseteq,pow_iff}.
        Therefore $U$ is open in $Y$ with respect to $T_Y$ by \cref{subseteq_antisymmetric,open_in}.
    \end{subproof}

    Hence $T_Y\subseteq T_e$ and $T_e\subseteq T_Y$ by \cref{subseteq,open_in}.
    Thus $T_Y = T_e$ by \cref{subseteq_antisymmetric}.
    Therefore $Y$ is metrizable with respect to $T_Y$ by \cref{metrizable}.
\end{proof}

% from §34 helper lemma 3: homeomorphism preserves metrizability
\begin{lemma}\label{homeomorphism_preserves_metrizable}
    Assume $T$ is a topology on $X$.
    Assume $S$ is a topology on $Y$.
    Suppose $f$ is a homeomorphism between $X$ and $Y$ with respect to $T$ and $S$.
    If $Y$ is metrizable with respect to $S$ then $X$ is metrizable with respect to $T$.
\end{lemma}
\begin{proof}
    Assume $Y$ is metrizable with respect to $S$.
    Let $g = \converse{f}$.
    $f\in\funs{X}{Y}$ and $f$ is a function and $\dom{f} = X$
        by \cref{homeomorphism,bijections,bijections_to_funs,funs_is_function,funs}.
    Hence $f$ is injective by \cref{homeomorphism,bijections}.
    $g$ is a bijection from $Y$ to $X$ and $g\in\bijections{Y}{X}$ and $\dom{g} = Y$ and $g$ is a function
        by \cref{homeomorphism,bijection_converse_is_bijection,bijections,bijections_dom,bijection_is_function}.
    $g\in\funs{Y}{X}$ by \cref{bijections_to_funs}.
    Take $d$ such that $d$ is a metric on $Y$ and $S = \metricTopology{d}{Y}$ by \cref{metrizable}.
    Let $e = \{((x_1,x_2),d((f(x_1),f(x_2)))) \mid x_1\in X, x_2\in X\}$.

    Show that $e$ is a metric on $X$.
    \begin{subproof}
        For all $w$ such that $w\in e$ there exist $p,r$ such that $w = (p,r)$ by \cref{relation,pair_eq_iff}.
        Hence $e$ is a relation by \cref{relation}.
        For all $p,r_1,r_2$ such that $(p,r_1)\in e$ and $(p,r_2)\in e$ we have $r_1 = r_2$
            by \cref{pair_eq_iff}.
        $e$ is right-unique by \cref{rightunique}.
        $e$ is a function by \cref{relation}.

        For all $p$ such that $p\in X\times X$ we have $p\in\dom{e}$
            by \cref{times,funs_type_apply,times_tuple_intro,metric_on,pair_eq_iff,dom_intro}.
        Hence $X\times X\subseteq\dom{e}$ by \cref{subseteq}.
        For all $p$ such that $p\in\dom{e}$ we have $p\in X\times X$
            by \cref{dom_iff,pair_eq_iff,times_tuple_intro}.
        $\dom{e}\subseteq X\times X$ by \cref{subseteq}.
        Hence $\dom{e} = X\times X$ by \cref{subseteq_antisymmetric}.

        For all $p\in\dom{e}$ we have $e(p)\in\reals$
            by \cref{dom_iff,function_apply_intro,pair_eq_iff,funs_type_apply,times_tuple_intro,metric_on}.
        Hence $e\in\funs{X\times X}{\reals}$ by \cref{funs_intro}.

        For all $x,y$ such that $x\in X$ and $y\in X$ we have $e((x,y)) = d((f(x),f(y)))$.

        Show that for all $x,y$ such that $x\in X$ and $y\in X$ we have $e((x,y)) \geq 0$.
        \begin{subproof}
            Follows by \cref{funs_type_apply,metric_on,metric_nonnegative}.
        \end{subproof}
        Hence $e$ is nonnegative on $X$ by \cref{metric_nonnegative}.

        Show that for all $x,y$ such that $x\in X$ and $y\in X$ we have $e((x,y)) = 0$ iff $x = y$.
        \begin{subproof}
            Follows by \cref{funs_type_apply,metric_on,metric_zero_characterization,injective_function}.
        \end{subproof}
        Thus $e$ is zero-characterized on $X$ by \cref{metric_zero_characterization}.

        Show that for all $x,y$ such that $x\in X$ and $y\in X$ we have $e((x,y)) = e((y,x))$.
        \begin{subproof}
            Follows by \cref{funs_type_apply,metric_on,metric_symmetric}.
        \end{subproof}
        Hence $e$ is symmetric on $X$ by \cref{metric_symmetric}.

        Show that for all $x,y,z$ such that $x\in X$ and $y\in X$ and $z\in X$
            we have $e((x,y)) + e((y,z)) \geq e((x,z))$.
        \begin{subproof}
            Follows by \cref{funs_type_apply,metric_on,metric_triangle_inequality}.
        \end{subproof}
        Thus $e$ satisfies the triangle inequality on $X$ by \cref{metric_triangle_inequality}.
        Therefore $e$ is a metric on $X$ by \cref{metric_on}.
    \end{subproof}

    Let $T_e = \metricTopology{e}{X}$.
    Hence $T_e$ is a topology on $X$ and
        $e\in\funs{X\times X}{\reals}$ and $e$ is a function and $e$ is a relation and $\dom{e} = X\times X$
        by \cref{metric_basis_is_basis,basis_topology_is_topology,metric_topology,metric_on,funs_elim,funs_is_function,funs_is_relation,funs}.

    Show that for all $x\in X$ we have $f$ is continuous at $x$ from $X$ to $Y$ with respect to $T_e$ and $S$.
    \begin{subproof}
        Fix $x\in X$.
        Show that for all $\epsilon\in\reals$ such that $0 < \epsilon$
            there exists $\delta\in\reals$ such that $0 < \delta$ and
            for all $y\in X$ if $e((x,y)) < \delta$ then $d((f(x),f(y))) < \epsilon$.
        \begin{subproof}
            Fix $\epsilon\in\reals$.
            Assume $0 < \epsilon$.
            Let $\delta = \epsilon$.
            Hence $\delta\in\reals$ and $0 < \delta$ by \cref{real_lt_subst_left}.
            Show that for all $y\in X$ if $e((x,y)) < \delta$ then $d((f(x),f(y))) < \epsilon$.
            \begin{subproof}
                Fix $y\in X$.
                Assume $e((x,y)) < \delta$.
                $e((x,y)) = d((f(x),f(y)))$.
                $d((f(x),f(y))) < \epsilon$ by \cref{real_lt_subst_left}.
            \end{subproof}
            Follows by definition.
        \end{subproof}
        Hence $f$ is continuous at $x$ from $X$ to $Y$ with respect to $T_e$ and $S$
            by \cref{metric_continuity_at_eps_delta,funs,metrizable}.
    \end{subproof}
    Show that $f$ is continuous from $X$ to $Y$ with respect to $T_e$ and $S$.
    \begin{subproof}
        Follows by \cref{continuous_at_implies_continuous,funs,metrizable}.
    \end{subproof}

    Show that for all $y\in Y$ we have $g$ is continuous at $y$ from $Y$ to $X$ with respect to $S$ and $T_e$.
    \begin{subproof}
        Fix $y\in Y$.
        Show that for all $\epsilon\in\reals$ such that $0 < \epsilon$
            there exists $\delta\in\reals$ such that $0 < \delta$ and
            for all $z\in Y$ if $d((y,z)) < \delta$ then $e((g(y),g(z))) < \epsilon$.
        \begin{subproof}
            Follows by \cref{funs_type_apply,function_apply_elim,converse_iff,function_apply_intro,real_lt_subst_left,pair_eq_iff}.
        \end{subproof}
        Hence $g$ is continuous at $y$ from $Y$ to $X$ with respect to $S$ and $T_e$
            by \cref{metric_continuity_at_eps_delta,funs,metrizable}.
    \end{subproof}
    Show that $g$ is continuous from $Y$ to $X$ with respect to $S$ and $T_e$.
    \begin{subproof}
        Follows by \cref{continuous_at_implies_continuous,funs,metrizable}.
    \end{subproof}

    Show that for all $U$ such that $U$ is open in $X$ with respect to $T$ we have $U$ is open in $X$ with respect to $T_e$.
    \begin{subproof}
        Fix $U$ such that $U$ is open in $X$ with respect to $T$.
        Hence $\img{f}{U}$ is open in $Y$ with respect to $S$
            by \cref{homeomorphism,continuous,funs_is_relation,funs,preim_eq_img_of_converse,converse_converse_eq}.
        For all $x$ such that $x\in\preimg{f}{\img{f}{U}}$ we have $x\in U$
            by \cref{homeomorphism,bijections,preimg_iff,img_iff,injective}.
        $\preimg{f}{\img{f}{U}}\subseteq U$ by \cref{subseteq}.
        For all $x$ such that $x\in U$ we have $x\in\preimg{f}{\img{f}{U}}$ by
            \cref{open_in,topology_on,subseteq,powerset_elim,funs,funs_is_function,function_apply_elim,img_iff,preimg_iff}.
        $U\subseteq\preimg{f}{\img{f}{U}}$ by \cref{subseteq}.
        Thus $\preimg{f}{\img{f}{U}} = U$ by \cref{subseteq_antisymmetric}.
        Therefore $U$ is open in $X$ with respect to $T_e$ by \cref{subseteq_antisymmetric,continuous,open_in}.
    \end{subproof}

    Show that for all $U$ such that $U$ is open in $X$ with respect to $T_e$ we have $U$ is open in $X$ with respect to $T$.
    \begin{subproof}
        Fix $U$ such that $U$ is open in $X$ with respect to $T_e$.
        Hence $\preimg{g}{U}$ is open in $Y$ with respect to $S$ by \cref{continuous}.
        $\img{f}{U}$ is open in $Y$ with respect to $S$
            by \cref{funs_is_relation,funs,preim_eq_img_of_converse,converse_converse_eq}.
        For all $x$ such that $x\in\preimg{f}{\img{f}{U}}$ we have $x\in U$
            by \cref{homeomorphism,bijections,preimg_iff,img_iff,injective}.
        Hence $\preimg{f}{\img{f}{U}}\subseteq U$ by \cref{subseteq}.
        For all $x$ such that $x\in U$ we have $x\in\preimg{f}{\img{f}{U}}$ by
            \cref{open_in,topology_on,subseteq,powerset_elim,funs,funs_is_function,function_apply_elim,img_iff,preimg_iff}.
        $U\subseteq\preimg{f}{\img{f}{U}}$ by \cref{subseteq}.
        Thus $\preimg{f}{\img{f}{U}} = U$ by \cref{subseteq_antisymmetric}.
        Therefore $U$ is open in $X$ with respect to $T$ by \cref{subseteq_antisymmetric,homeomorphism,continuous,open_in}.
    \end{subproof}

    Hence $T\subseteq T_e$ and $T_e\subseteq T$ by \cref{subseteq,open_in}.
    Thus $T = T_e$ by \cref{subseteq_antisymmetric}.
    Therefore $X$ is metrizable with respect to $T$ by \cref{metrizable}.
\end{proof}

% from §34 helper lemma 1: countable embedding into real power from unit-interval valued maps
\begin{lemma}\label{countable_embedding_theorem_real_power_unit_interval}
    Assume $T$ is a topology on $X$.
    Assume $X$ has closed singletons with respect to $T$.
    Let $I = \intervalclosed{0}{1}$.
    Let $T_I = \subspaceTopology{\standardTopologyR}{I}$.
    Assume $f$ is a function on $\naturalsPlus$.
    Suppose for all $n\in\naturalsPlus$ we have $f(n)$ is continuous from $X$ to $I$ with respect to $T$ and $T_I$.
    Suppose for all $x,U$ such that $x\in X$ and $U$ is a neighborhood of $x$ in $X$ with respect to $T$
        there exists $n\in\naturalsPlus$ such that
            $0 < \apply{\apply{f}{n}}{x}$ and
            for all $y\in X\setminus U$ we have $\apply{\apply{f}{n}}{y} = 0$.
    Let $R = \{(n,\reals) \mid n\in\naturalsPlus\}$.
    Let $S = \{(n,\standardTopologyR) \mid n\in\naturalsPlus\}$.
    Let $P = \productTopologyIndexed{S}{R}$.
    Let $F = \{(x,z) \mid x\in X, z\in\productFamily{R} \mid
        \forall n\in\naturalsPlus. z(n) = \apply{\apply{f}{n}}{x}\}$.
    Then $F$ is a topological embedding from $X$ to $\productFamily{R}$ with respect to $T$ and $P$.
\end{lemma}
\begin{proof}
    We have $\naturalsPlus$ is inhabited by \cref{real_one_in_naturals}.
    $\standardTopologyR$ is a topology on $\reals$
        by \cref{standard_basis_is_basis,basis_topology_is_topology,standard_topology_reals}.
    $I\subseteq\reals$ by \cref{intervalclosed,subseteq}.
    Hence $I$ is a subspace of $\reals$ with respect to $\standardTopologyR$ and $T_I$ is a topology on $I$
        by \cref{subspace_of,subspace_topology_is_topology}.
    $R$ is a function on $\naturalsPlus$ by \cref{relation,rightunique,dom_intro,dom_iff,pair_eq_iff,subseteq,subseteq_antisymmetric}.
    $S$ is a function on $\naturalsPlus$ by \cref{relation,rightunique,dom_intro,dom_iff,pair_eq_iff,subseteq,subseteq_antisymmetric}.
    Thus $\dom{S} = \naturalsPlus$ and $\dom{R} = \naturalsPlus$ by \cref{funs}.
    For all $n\in\naturalsPlus$ we have $S(n) = \standardTopologyR$ and $R(n) = \reals$ by \cref{function_apply_intro}.
    Hence $\productFamily{R} = \realsPower{\naturalsPlus}$ by \cref{product_family_reals_power}.
    $P$ is a topology on $\productFamily{R}$
        by \cref{product_subbasis_is_subbasis,product_subbasis_finite_intersections_basis,basis_topology_is_topology,basis_for_topology,standard_topology_reals}.

    Show that $F$ is a function from $X$ to $\productFamily{R}$.
    \begin{subproof}
        For all $w$ such that $w\in F$ there exist $x,z$ such that $w = (x,z)$ by \cref{relation}.
        Hence $F$ is a relation by \cref{relation}.
        For all $x,z_1,z_2$ such that $(x,z_1)\in F$ and $(x,z_2)\in F$ we have $z_1 = z_2$
            by \cref{pair_eq_iff,product_family_reals_power,subseteq,reals_power,tuple_power,functions_eq_iff}.
        $F$ is right-unique by \cref{rightunique}.
        $F$ is a function by \cref{relation}.
        Show that for all $x\in X$ there exists $z$ such that $(x,z)\in F$.
        \begin{subproof}
            Fix $x\in X$.
            Let $z = \{(n,\apply{\apply{f}{n}}{x}) \mid n\in\naturalsPlus\}$.
            For all $n,y_1,y_2$ such that $(n,y_1)\in z$ and $(n,y_2)\in z$ we have $y_1 = y_2$ by \cref{pair_eq_iff}.
            Hence $z$ is a function by \cref{relation,rightunique,pair_eq_iff}.
            $\dom{z} = \naturalsPlus$ by \cref{subseteq_antisymmetric,subseteq,dom_intro,dom_iff,pair_eq_iff}.
            For all $n\in\dom{z}$ we have $z(n)\in\reals$
                by \cref{dom_iff,pair_eq_iff,continuous,funs_type_apply,intervalclosed,subseteq,function_apply_intro}.
            $z\in\productFamily{R}$ by \cref{funs_intro,reals_power,tuple_power,subseteq}.
            For all $n\in\naturalsPlus$ we have $z(n) = \apply{\apply{f}{n}}{x}$ by \cref{function_apply_intro}.
            $(x,z)\in F$ by \cref{pair_eq_iff}.
        \end{subproof}
        Hence $\dom{F} = X$ by \cref{subseteq,dom_iff,pair_eq_iff,subseteq_antisymmetric}.
        $F\in\funs{X}{\productFamily{R}}$ by \cref{funs_intro,function_apply_iff,pair_eq_iff}.
    \end{subproof}

    Show that $F$ is injective.
    \begin{subproof}
        $F$ is a function and $\dom{F} = X$ by \cref{funs_elim}.
        Show that for all $x_1,x_2$ such that $x_1\in\dom{F}$ and $x_2\in\dom{F}$ and $F(x_1) = F(x_2)$ we have $x_1 = x_2$.
        \begin{subproof}
            Fix $x_1$.
            Fix $x_2$ such that $x_1\in\dom{F}$ and $x_2\in\dom{F}$ and $F(x_1) = F(x_2)$.
            Hence $x_1\in X$ and $x_2\in X$.
            Suppose not.
            Hence $\{x_2\}$ is closed in $X$ with respect to $T$ by \cref{one_point_closed,singleton_subset_intro}.
            Let $U = X\setminus\{x_2\}$.
            Therefore $U$ is open in $X$ with respect to $T$
                by \cref{closed_in,topology_on,setminus_subseteq,double_relative_complement}.
            Hence $x_1\in U$ by \cref{singleton_elim,setminus_intro}.
            $U$ is a neighborhood of $x_1$ in $X$ with respect to $T$ by \cref{neighborhood}.
            Take $n\in\naturalsPlus$ such that
                $0 < \apply{\apply{f}{n}}{x_1}$ and
                for all $y\in X\setminus U$ we have $\apply{\apply{f}{n}}{y} = 0$
                by assumption.
            Hence $\apply{\apply{f}{n}}{x_2} = 0$
                by \cref{setminus_intro,setminus_elim_right,singleton_intro}.
            $\apply{\apply{F}{x_1}}{n} = \apply{\apply{f}{n}}{x_1}$ and $\apply{\apply{F}{x_2}}{n} = \apply{\apply{f}{n}}{x_2}$
                by \cref{funs_elim,function_apply_elim,pair_eq_iff}.
            $0 < \apply{\apply{F}{x_1}}{n}$ by \cref{real_lt_subst_right}.
            $\apply{\apply{F}{x_2}}{n} = 0$ by \cref{real_lt_subst_right}.
            $\apply{\apply{F}{x_1}}{n} \neq \apply{\apply{F}{x_2}}{n}$ by \cref{real_lt_subst_right,real_add_ident,real_lt_irreflexive}.
            $\apply{\apply{F}{x_1}}{n} = \apply{\apply{F}{x_2}}{n}$.
            Contradiction.
        \end{subproof}
        Therefore $F$ is injective by \cref{injective_function}.
    \end{subproof}

    Show that $F$ is continuous from $X$ to $\productFamily{R}$ with respect to $T$ and $P$.
    \begin{subproof}
        Show that for all $n\in\naturalsPlus$ we have $\piIndex{R}{n}\circ F$ is continuous from $X$ to $\reals$ with respect to $T$ and $\standardTopologyR$.
        \begin{subproof}
            Fix $n\in\naturalsPlus$.
            Let $j = \identity{I}$.
            Let $g = j\circ f(n)$.
            $g$ is continuous from $X$ to $\reals$ with respect to $T$ and $\standardTopologyR$
                by \cref{continuous_expand_range}.
            $g$ is a function and $\dom{g} = X$ by \cref{continuous,funs_elim}.
            $\piIndex{R}{n}\circ F\in\funs{X}{\reals}$ by \cref{projection_map_indexed_function,funs_circ,funs_elim}.
            $F$ is a function and $\dom{F} = X$ by \cref{funs_elim}.
            For all $x\in X$ we have $\apply{\piIndex{R}{n}\circ F}{x} = g(x)$
                by \cref{function_apply_elim,pair_eq_iff,continuous,funs_elim,funs_type_apply,id_iff,circ_elem_intro,projection_map_indexed,function_apply_intro}.
            Hence $\piIndex{R}{n}\circ F = g$ by \cref{functions_eq_iff,continuous}.
            $\piIndex{R}{n}\circ F$ is continuous from $X$ to $\reals$ with respect to $T$ and $\standardTopologyR$
                by \cref{continuous_eq_subst}.
        \end{subproof}
        Hence $F$ is continuous from $X$ to $\productFamily{R}$ with respect to $T$ and $P$
            by \cref{continuous_map_into_indexed_product,standard_basis_is_basis,basis_topology_is_topology,standard_topology_reals,product_family_reals_power,funs,subseteq}.
    \end{subproof}

    Show that $F$ is a homeomorphism between $X$ and $\img{F}{X}$ with respect to $T$ and $\subspaceTopology{P}{\img{F}{X}}$.
    \begin{subproof}
        $F$ is a function and $\dom{F} = X$ by \cref{funs_elim}.
        $F\in\funs{X}{\img{F}{X}}$ by \cref{funs_intro,function_apply_elim,ran_intro,img_dom}.
        Hence $F\in\surj{X}{\img{F}{X}}$ by \cref{funs_surj_iff,img_dom}.
        $F$ is a bijection from $X$ to $\img{F}{X}$ by \cref{bijections,injective_function}.
        $\converse{F}\in\funs{\img{F}{X}}{X}$ by \cref{bijective_converse_are_funs}.
        Show that for all $z\in\img{F}{X}$ we have
            $\converse{F}$ is continuous at $z$ from $\img{F}{X}$ to $X$ with respect to $\subspaceTopology{P}{\img{F}{X}}$ and $T$.
        \begin{subproof}
            Fix $z\in\img{F}{X}$.
            Take $x\in X$ such that $(x,z)\in F$ by \cref{img_iff}.
            Hence $(z,x)\in\converse{F}$ by \cref{converse_iff}.
            $\apply{\converse{F}}{z} = x$ by \cref{funs_elim,function_apply_intro}.
            Show that for all $U$ such that $U$ is a neighborhood of $\apply{\converse{F}}{z}$ in $X$ with respect to $T$
                there exists $W$ such that
                    $W$ is a neighborhood of $z$ in $\img{F}{X}$ with respect to $\subspaceTopology{P}{\img{F}{X}}$ and
                    $\img{\converse{F}}{W}\subseteq U$.
            \begin{subproof}
                Fix $U$ such that $U$ is a neighborhood of $\apply{\converse{F}}{z}$ in $X$ with respect to $T$.
                Take $n\in\naturalsPlus$ such that
                    $0 < \apply{\apply{f}{n}}{x}$ and
                    for all $y\in X\setminus U$ we have $\apply{\apply{f}{n}}{y} = 0$
                    by assumption.
                Let $V = \intervalopen{0}{1 + 1}$.
                Let $D = \preimg{\piIndex{R}{n}}{V}$.
                $0\in\reals$ and $1 + 1\in\reals$ and $0 < 1 + 1$
                    by \cref{real_naturals_subset_reals,real_one_in_naturals,subseteq,real_add_ident,real_add_closed,real_zero_lt_one,real_lt_add,real_lt_trans}.
                Therefore $D$ is open in $\productFamily{R}$ with respect to $P$
                    by \cref{intervalopen_subset_reals,powerset_intro,standard_basis_reals,basis_elem_open_in_generated,basis_for_topology,standard_basis_is_basis,standard_topology_reals,projection_map_indexed_continuous,continuous,open_in,product_family_reals_power,subseteq,basis_topology_is_topology}.
                $\apply{\apply{F}{x}}{n} = \apply{\apply{f}{n}}{x}$
                    by \cref{funs_elim,function_apply_elim,pair_eq_iff}.
                $\apply{\apply{f}{n}}{x}\in I$ by \cref{continuous,funs_type_apply}.
                $\apply{\apply{f}{n}}{x} < 1 + 1$ by \cref{intervalclosed,real_lt_add,real_add_ident,real_lt_trans,real_one_leq_naturalsplus,real_one_in_naturals,real_naturals_subset_reals,subseteq,real_lt_imp_leq}.
                $0 < \apply{\apply{F}{x}}{n}$ by \cref{real_lt_subst_right}.
                $\apply{\apply{F}{x}}{n} < 1 + 1$ by \cref{real_lt_subst_left}.
                $F(x)\in\productFamily{R}$ by \cref{funs_type_apply}.
                $\apply{\apply{F}{x}}{n}\in\reals$ by \cref{product_family_reals_power,reals_power,tuple_power,subseteq,funs_type_apply}.
                $\apply{\apply{F}{x}}{n}\in V$ by \cref{intervalopen}.
                $(z,\apply{\apply{F}{x}}{n})\in\piIndex{R}{n}$
                    by \cref{projection_map_indexed,function_apply_elim,pair_eq_iff}.
                $z\in D$ by \cref{preimg_iff}.
                Let $W = \img{F}{X}\inter D$.
                $\img{F}{X}\subseteq\productFamily{R}$ by \cref{subseteq,img_iff,pair_eq_iff}.
                $\subspaceTopology{P}{\img{F}{X}}$ is a topology on $\img{F}{X}$ by \cref{subspace_topology_is_topology}.
                Hence $W$ is open in $\img{F}{X}$ with respect to $\subspaceTopology{P}{\img{F}{X}}$ by \cref{subspace_topology,open_in}.
                $z\in W$ by \cref{img_iff,inter_intro}.
                $W$ is a neighborhood of $z$ in $\img{F}{X}$ with respect to $\subspaceTopology{P}{\img{F}{X}}$ by \cref{neighborhood}.
                $\converse{F}$ is a function and $\dom{\converse{F}} = \img{F}{X}$ by \cref{funs_elim}.
                Show that for all $y$ such that $y\in\img{\converse{F}}{W}$ we have $y\in U$.
                \begin{subproof}
                    Fix $y$ such that $y\in\img{\converse{F}}{W}$.
                    Take $p\in\dom{\converse{F}}\inter W$ such that $y = \apply{\converse{F}}{p}$ by \cref{img_of_function_elim}.
                    Hence $p\in W$ by \cref{inter_elim_right}.
                    $p\in D$ and $p\in\img{F}{X}$ by \cref{inter_elim_left,inter_elim_right}.
                    $p\in\dom{\converse{F}}$ by \cref{subseteq}.
                    $(y,p)\in F$ by \cref{converse_iff,function_apply_elim}.
                    $y\in\dom{F}$ by \cref{dom_iff}.
                    $y\in X$ by \cref{funs_elim}.
                    $p = F(y)$ by \cref{function_apply_intro}.
                    $p\in\productFamily{R}$ by \cref{subseteq}.
                    $R$ is a function and $n\in\dom{R}$ by \cref{funs_elim}.
                    $\piIndex{R}{n}\in\funs{\productFamily{R}}{\reals}$ by \cref{projection_map_indexed_function}.
                    $\piIndex{R}{n}$ is a function by \cref{funs_elim}.
                    Take $q\in V$ such that $(p,q)\in\piIndex{R}{n}$ by \cref{preimg_iff}.
                    $F(y)\in\productFamily{R}$ by \cref{funs_type_apply}.
                    $(p,\apply{\apply{F}{y}}{n})\in\piIndex{R}{n}$
                        by \cref{projection_map_indexed,function_apply_elim,pair_eq_iff}.
                    $q = \apply{\apply{F}{y}}{n}$ by \cref{function_rightunique}.
                    $0 < q$ by \cref{intervalopen}.
                    $0 < \apply{\apply{F}{y}}{n}$ by \cref{real_lt_subst_right}.
                    $\apply{\apply{F}{y}}{n} = \apply{\apply{f}{n}}{y}$ by \cref{pair_eq_iff}.
                    $0 < \apply{\apply{f}{n}}{y}$ by \cref{real_lt_subst_right}.
                    $y\in U$ by \cref{setminus_intro,real_lt_subst_right,real_lt_irreflexive}.
                \end{subproof}
                Hence $\img{\converse{F}}{W}\subseteq U$ by \cref{subseteq}.
            \end{subproof}
            $\img{F}{X}\subseteq\productFamily{R}$ by \cref{subseteq,img_iff,pair_eq_iff}.
            $\subspaceTopology{P}{\img{F}{X}}$ is a topology on $\img{F}{X}$ by \cref{subspace_topology_is_topology}.
            Hence $\converse{F}$ is continuous at $z$ from $\img{F}{X}$ to $X$ with respect to $\subspaceTopology{P}{\img{F}{X}}$ and $T$
                by \cref{continuous_at}.
        \end{subproof}
        Hence $\img{F}{X}\subseteq\productFamily{R}$ by \cref{subseteq,img_iff,pair_eq_iff}.
        $\subspaceTopology{P}{\img{F}{X}}$ is a topology on $\img{F}{X}$ by \cref{subspace_topology_is_topology}.
        Show that $\converse{F}$ is continuous from $\img{F}{X}$ to $X$ with respect to $\subspaceTopology{P}{\img{F}{X}}$ and $T$.
        \begin{subproof}
            Follows by \cref{continuous_at_implies_continuous}.
        \end{subproof}
        $\img{F}{X}$ is a subspace of $\productFamily{R}$ with respect to $P$ by \cref{subspace_of}.
        Show that $F$ is continuous from $X$ to $\img{F}{X}$ with respect to $T$ and $\subspaceTopology{P}{\img{F}{X}}$.
        \begin{subproof}
            Follows by \cref{subseteq_refl,continuous_restrict_range}.
        \end{subproof}
        Show that $F$ is a homeomorphism between $X$ and $\img{F}{X}$ with respect to $T$ and $\subspaceTopology{P}{\img{F}{X}}$.
        \begin{subproof}
            Follows by \cref{homeomorphism}.
        \end{subproof}
    \end{subproof}
    Show that $F$ is a topological embedding from $X$ to $\productFamily{R}$ with respect to $T$ and $P$.
    \begin{subproof}
        Follows by \cref{topological_embedding}.
    \end{subproof}
\end{proof}

% from §34 helper lemma 4: countable unit-interval family from a countable basis
\begin{lemma}\label{regular_second_countable_unit_interval_family}
    Assume $X$ is regular with respect to $T$.
    Let $I = \intervalclosed{0}{1}$.
    Let $T_I = \subspaceTopology{\standardTopologyR}{I}$.
    Suppose $X$ satisfies the second countability axiom with respect to $T$.
    Then there exists $f$ such that
        $f$ is a function on $\naturalsPlus$ and
        for all $n\in\naturalsPlus$ we have $f(n)$ is continuous from $X$ to $I$ with respect to $T$ and $T_I$ and
        for all $x,U$ such that $x\in X$ and $U$ is a neighborhood of $x$ in $X$ with respect to $T$
            there exists $n\in\naturalsPlus$ such that
                $0 < \apply{\apply{f}{n}}{x}$ and
                for all $y\in X\setminus U$ we have $\apply{\apply{f}{n}}{y} = 0$.
\end{lemma}
\begin{proof}
    $T$ is a topology on $X$ and $X$ has closed singletons with respect to $T$ by \cref{regular_space}.
    Hence $X$ is normal with respect to $T$ by \cref{regular_second_countable_normal}.
    $I\subseteq\reals$ by \cref{intervalclosed,subseteq}.
    Therefore $T_I$ is a topology on $I$ by \cref{subspace_of,subspace_topology_is_topology,standard_topology_reals,standard_basis_is_basis,basis_topology_is_topology}.
    We have $0\in I$ and $1\in I$
        by \cref{intervalclosed,real_add_ident,real_one_in_naturals,real_naturals_subset_reals,subseteq,real_zero_lt_one}.

    Take $b$ such that
        $b$ is a sequence in $\pow{X}$ and
        for all $n\in\naturalsPlus$ we have $b(n)$ is open in $X$ with respect to $T$ and
        for all $x\in X$ we have
            for all $U$ such that $U$ is open in $X$ with respect to $T$ and $x\in U$
                there exists $n\in\naturalsPlus$ such that $x\in b(n)$ and $b(n)\subseteq U$
        by \cref{second_countability_axiom}.

    Let $B = \pow{X\times I}$.
    Let $J = \{p\in\naturalsPlus\times\naturalsPlus \mid \closure{b(\fst{p})}{X}{T}\subseteq b(\snd{p})\}$.
    Let $R = \{w\in J\times B \mid \text{there exists $p\in J$ such that there exists $u\in B$ such that
        $w = (p,u)$ and
        $u$ is continuous from $X$ to $I$ with respect to $T$ and $T_I$ and
        $\img{u}{\closure{b(\fst{p})}{X}{T}}\subseteq \{1\}$ and
        $\img{u}{X\setminus b(\snd{p})}\subseteq \{0\}$}\}$.

    Show that for all $p\in J$ there exists $u$ such that $p\mathrel{R} u$.
    \begin{subproof}
        Fix $p\in J$.
        Hence $p\in\naturalsPlus\times\naturalsPlus$ and $\closure{b(\fst{p})}{X}{T}\subseteq b(\snd{p})$ by \cref{pair_eq_iff}.
        $\fst{p}\in\naturalsPlus$ and $\snd{p}\in\naturalsPlus$ by \cref{fst_type,snd_type}.
        $b(\fst{p})$ is open in $X$ with respect to $T$ and $b(\snd{p})$ is open in $X$ with respect to $T$ by assumption.
        $b(\fst{p})\subseteq X$ and $b(\snd{p})\subseteq X$ by \cref{open_in,topology_on,subseteq,pow_iff}.
        $\closure{b(\fst{p})}{X}{T}$ is closed in $X$ with respect to $T$ by \cref{closure_closed_in}.
        Let $A = X\setminus b(\snd{p})$.
        Hence $A$ is closed in $X$ with respect to $T$ by \cref{open_in,closed_in,setminus_subseteq,double_relative_complement,subseteq}.
        $A$ is disjoint from $\closure{b(\fst{p})}{X}{T}$ by \cref{disjoint,subseteq,setminus_elim_right}.
        Take $u$ such that
            $u$ is continuous from $X$ to $I$ with respect to $T$ and $T_I$ and
            $\img{u}{A}\subseteq \{0\}$ and
            $\img{u}{\closure{b(\fst{p})}{X}{T}}\subseteq \{1\}$
            by \cref{urysohn_unit_interval}.
        Hence $u\in B$ by \cref{continuous,funs_subseteq_rels,subseteq,rels_ran_subseteq,rels_elim,powerset_intro}.
        $\img{u}{X\setminus b(\snd{p})}\subseteq \{0\}$ by \cref{pair_eq_iff}.
        Thus $p\mathrel{R} u$ by \cref{times_tuple_intro,pair_eq_iff}.
    \end{subproof}

    $R$ is a relation by \cref{relation}.
    Take $c$ such that $c$ is a function on $J$ and for all $p\in J$ we have $p\mathrel{R} c(p)$
        by \cref{choice_function}.
    For all $p\in\dom{c}$ we have $c(p)\in B$ by \cref{subseteq,choice_function,pair_eq_iff}.
    Hence $c$ is a function and $\dom{c} = J$ by \cref{choice_function}.
    Therefore $c\in\funs{J}{B}$ by \cref{funs_intro}.

    Take $g$ such that $g\in\inj{\naturalsPlus\times\naturalsPlus}{\naturalsPlus}$ by \cref{naturalsplus_pairing_injection}.
    Let $h = \restrl{g}{J}$.
    $g\in\funs{\naturalsPlus\times\naturalsPlus}{\naturalsPlus}$ by \cref{inj}.
    Hence $g$ is a function and $g$ is a relation and $\dom{g} = \naturalsPlus\times\naturalsPlus$
        by \cref{funs_is_function,funs_is_relation,funs}.
    $g$ is injective by \cref{inj,injective_function}.
    $h$ is a function by \cref{restrl_of_function_is_function}.
    $J\subseteq\naturalsPlus\times\naturalsPlus$ by \cref{subseteq}.
    Thus $\dom{h} = J$ by \cref{funs,dom_of_restrl_eq_inter,inter_absorb_supseteq_left}.
    For all $p\in\dom{h}$ we have $h(p)\in\naturalsPlus$
        by \cref{subseteq,funs,restrl_of_function_apply,inj,funs_type_apply,pair_eq_iff}.
    Hence $h\in\funs{J}{\naturalsPlus}$ by \cref{funs_intro}.
    For all $p_1,p_2$ such that $p_1\in J$ and $p_2\in J$ and $h(p_1) = h(p_2)$ we have $p_1 = p_2$
        by \cref{subseteq,funs,restrl_of_function_apply,pair_eq_iff,injective_function}.
    Therefore $h\in\inj{J}{\naturalsPlus}$ by \cref{inj}.

    Let $A = \img{h}{J}$.
    Let $z = \{(x,0) \mid x\in X\}$.
    $z$ is a function from $X$ to $I$ by \cref{constant_map_function}.
    Hence for all $x\in X$ we have $z(x) = 0$ by \cref{constant_map_function,funs_is_function,function_apply_intro}.
    $z$ is continuous from $X$ to $I$ with respect to $T$ and $T_I$ by \cref{constant_continuous}.
    $z\in B$ by \cref{continuous,funs_subseteq_rels,subseteq,rels_ran_subseteq,rels_elim,powerset_intro}.

    Let $f_1 = \{(h(p),c(p)) \mid p\in J\}$.
    For all $w$ such that $w\in f_1$ there exist $n,u$ such that $w = (n,u)$ by \cref{relation,pair_eq_iff}.
    Hence $f_1$ is a relation by \cref{relation}.
    For all $n,u_1,u_2$ such that $(n,u_1)\in f_1$ and $(n,u_2)\in f_1$ we have $u_1 = u_2$
        by \cref{pair_eq_iff,inj}.
    $f_1$ is right-unique by \cref{rightunique}.
    $f_1$ is a function by \cref{relation}.
    For all $n$ such that $n\in A$ we have $n\in\dom{f_1}$
        by \cref{inj,funs,funs_is_function,function_rightunique,funs_is_relation,img_iff,function_apply_intro,pair_eq_iff,dom_intro}.
    Hence $A\subseteq\dom{f_1}$ by \cref{subseteq}.
    For all $n$ such that $n\in\dom{f_1}$ we have $n\in A$
        by \cref{dom_iff,pair_eq_iff,inj,funs,funs_is_function,function_rightunique,funs_is_relation,subseteq,function_apply_elim,img_iff}.
    $\dom{f_1}\subseteq A$ by \cref{subseteq}.
    Therefore $\dom{f_1} = A$ by \cref{subseteq_antisymmetric}.

    Let $f_2 = \{(n,z) \mid n\in\naturalsPlus\setminus A\}$.
    $f_2$ is a function from $\naturalsPlus\setminus A$ to $B$ by \cref{constant_map_function}.
    Hence $\dom{f_2} = \naturalsPlus\setminus A$ by \cref{constant_map_function,funs}.
    Let $f = f_1\union f_2$.
    For all $n$ such that $n\in\dom{f_1}\inter\dom{f_2}$ we have $f_1(n) = f_2(n)$
        by \cref{inter_elim_left,subseteq,inter_elim_right,setminus_elim_right}.
    Hence $f$ is a function by \cref{union_compatible_functions,funs_is_function}.
    $\dom{f} = \dom{f_1}\union\dom{f_2}$ by \cref{dom_union}.
    $\dom{f} = A\union(\naturalsPlus\setminus A)$ by \cref{funs,subseteq,real_one_in_naturals,subseteq_antisymmetric}.
    $A\subseteq\naturalsPlus$
        by \cref{inj,funs_subseteq_rels,subseteq,rels_ran_subseteq,img_subseteq_ran,subseteq_transitive}.
    Hence $\dom{f} = \naturalsPlus$ by \cref{setminus_partition}.

    Show that for all $n\in\naturalsPlus$ we have $f(n)$ is continuous from $X$ to $I$ with respect to $T$ and $T_I$.
    \begin{subproof}
        Fix $n\in\naturalsPlus$.
        \begin{byCase}
        \caseOf{$n\in A$.}
            $h\in\funs{J}{\naturalsPlus}$ by \cref{inj}.
            We have $h$ is right-unique and $h$ is a relation and $\dom{h} = J$
                by \cref{funs,funs_is_function,function_rightunique,funs_is_relation}.
            Take $p\in J$ such that $(p,n)\in h$ by \cref{img_iff}.
            Hence $(n,c(p))\in f_1$ by \cref{function_apply_intro,pair_eq_iff}.
            $(n,c(p))\in f$ by \cref{union_iff,union_intro_left}.
            $(p,c(p))\in R$ by \cref{choice_function}.
            Thus $f(n)$ is continuous from $X$ to $I$ with respect to $T$ and $T_I$
                by \cref{function_apply_intro,pair_eq_iff,continuous_eq_subst}.
        \caseOf{$n\notin A$.}
            We have $(n,z)\in f_2$ by \cref{setminus_intro,pair_eq_iff}.
            $(n,z)\in f$ by \cref{union_iff,union_intro_right}.
            Thus $f(n)$ is continuous from $X$ to $I$ with respect to $T$ and $T_I$
                by \cref{function_apply_intro,continuous_eq_subst}.
        \end{byCase}
    \end{subproof}

    Show that for all $x,U$ such that $x\in X$ and $U$ is a neighborhood of $x$ in $X$ with respect to $T$
        there exists $n\in\naturalsPlus$ such that
            $0 < \apply{\apply{f}{n}}{x}$ and
            for all $y\in X\setminus U$ we have $\apply{\apply{f}{n}}{y} = 0$.
    \begin{subproof}
        Fix $x$.
        Fix $U$ such that $x\in X$ and $U$ is a neighborhood of $x$ in $X$ with respect to $T$.
        Hence $U$ is open in $X$ with respect to $T$ by \cref{neighborhood}.
        $1\in\naturalsPlus$ by \cref{real_one_in_naturals}.
        Therefore for all $W$ such that $W$ is open in $X$ with respect to $T$ and $x\in W$
            there exists $n\in\naturalsPlus$ such that $x\in b(n)$ and $b(n)\subseteq W$
            by assumption.
        Take $m\in\naturalsPlus$ such that $x\in b(m)$ and $b(m)\subseteq U$ by \cref{neighborhood}.
        $b(m)$ is a neighborhood of $x$ in $X$ with respect to $T$ by \cref{open_in,neighborhood}.
        Take $V$ such that
            $V$ is a neighborhood of $x$ in $X$ with respect to $T$ and
            $\closure{V}{X}{T}\subseteq b(m)$ by \cref{regular_closure_neighborhood}.
        Thus $V$ is open in $X$ with respect to $T$ and $x\in V$ by \cref{neighborhood}.
        Take $k\in\naturalsPlus$ such that $x\in b(k)$ and $b(k)\subseteq V$ by assumption.
        Let $p = (k,m)$.
        $p\in\naturalsPlus\times\naturalsPlus$ by \cref{times_tuple_intro}.
        $b(k)$ is open in $X$ with respect to $T$ by assumption.
        $b(k)\subseteq X$ by \cref{open_in,topology_on,subseteq,pow_iff}.
        $V\subseteq X$ by \cref{neighborhood,open_in,topology_on,subseteq,pow_iff}.
        $V\subseteq\closure{V}{X}{T}$ by \cref{subseteq_closure}.
        $b(k)\subseteq\closure{V}{X}{T}$ by \cref{subseteq_transitive}.
        $\closure{V}{X}{T}$ is closed in $X$ with respect to $T$ by \cref{closure_closed_in}.
        $\closure{b(k)}{X}{T}\subseteq\closure{V}{X}{T}$ by \cref{closure_subseteq_closed}.
        $\closure{b(k)}{X}{T}\subseteq b(m)$ by \cref{subseteq_transitive}.
        $\fst{p} = k$ and $\snd{p} = m$ by \cref{fst_eq,snd_eq}.
        Hence $p\in J$ by \cref{pair_eq_iff}.
        Let $n = h(p)$.
        $h\in\funs{J}{\naturalsPlus}$ by \cref{inj}.
        Hence $h$ is right-unique and $h$ is a relation and $\dom{h} = J$
            by \cref{funs,funs_is_function,function_rightunique,funs_is_relation}.
        $(p,n)\in h$ by \cref{subseteq,function_apply_elim}.
        $(n,c(p))\in f_1$ by \cref{img_iff,pair_eq_iff}.
        $(n,c(p))\in f$ by \cref{union_iff,union_intro_left}.
        $f(n) = c(p)$ by \cref{function_apply_intro}.
        $(p,c(p))\in R$ by \cref{choice_function}.
        $\img{\apply{c}{p}}{\closure{b(\fst{p})}{X}{T}}\subseteq \{1\}$ and
            $\img{\apply{c}{p}}{X\setminus b(\snd{p})}\subseteq \{0\}$ by \cref{pair_eq_iff}.
        $\fst{p} = k$ and $\snd{p} = m$ by \cref{fst_eq,snd_eq}.
        $x\in\closure{b(k)}{X}{T}$ by \cref{subseteq_closure,subseteq}.
        $\apply{c}{p}$ is continuous from $X$ to $I$ with respect to $T$ and $T_I$ by \cref{pair_eq_iff}.
        $\apply{c}{p}\in\funs{X}{I}$ and $\apply{c}{p}$ is a function and $\dom{\apply{c}{p}} = X$
            by \cref{continuous,funs_is_function,funs}.
        $x\in\dom{\apply{c}{p}}$ by \cref{subseteq}.
        $(x,\apply{\apply{c}{p}}{x})\in\apply{c}{p}$ by \cref{function_apply_elim}.
        $\apply{\apply{c}{p}}{x}\in\img{\apply{c}{p}}{\closure{b(k)}{X}{T}}$ by \cref{img_iff}.
        $\apply{\apply{c}{p}}{x}\in\{1\}$ by \cref{subseteq}.
        $\apply{\apply{c}{p}}{x} = 1$ by \cref{singleton_elim}.
        $0 < \apply{\apply{c}{p}}{x}$
            by \cref{real_one_in_naturals,real_naturals_subset_reals,subseteq,real_zero_lt_one,real_lt_subst_right}.
        Hence $0 < \apply{\apply{f}{n}}{x}$ by \cref{real_lt_subst_right}.
        Show that for all $y\in X\setminus U$ we have $\apply{\apply{f}{n}}{y} = 0$.
        \begin{subproof}
            Follows by \cref{setminus_intro,setminus_elim_left,subseteq,setminus_elim_right,pair_eq_iff,continuous,funs_is_function,funs,function_apply_elim,img_iff,singleton_elim}.
        \end{subproof}
        Follows by definition.
    \end{subproof}
    Follows by definition.
\end{proof}

% from §34 Item 1: Urysohn metrization theorem
\begin{theorem}\label{urysohn_metrization_theorem}
    Assume $X$ is regular with respect to $T$.
    Suppose $X$ satisfies the second countability axiom with respect to $T$.
    Then $X$ is metrizable with respect to $T$.
\end{theorem}
\begin{proof}
    Let $I = \intervalclosed{0}{1}$.
    Let $T_I = \subspaceTopology{\standardTopologyR}{I}$.
    Take $f$ such that
        $f$ is a function on $\naturalsPlus$ and
        for all $n\in\naturalsPlus$ we have $f(n)$ is continuous from $X$ to $I$ with respect to $T$ and $T_I$ and
        for all $x,U$ such that $x\in X$ and $U$ is a neighborhood of $x$ in $X$ with respect to $T$
            there exists $n\in\naturalsPlus$ such that
                $0 < \apply{\apply{f}{n}}{x}$ and
                for all $y\in X\setminus U$ we have $\apply{\apply{f}{n}}{y} = 0$
        by \cref{regular_second_countable_unit_interval_family}.
    $1\in\naturalsPlus$ by \cref{real_one_in_naturals}.
    Hence for all $x,U$ such that $x\in X$ and $U$ is a neighborhood of $x$ in $X$ with respect to $T$
        there exists $n\in\naturalsPlus$ such that
            $0 < \apply{\apply{f}{n}}{x}$ and
            for all $y\in X\setminus U$ we have $\apply{\apply{f}{n}}{y} = 0$ by assumption.
    $T$ is a topology on $X$ and $X$ has closed singletons with respect to $T$ by \cref{regular_space}.
    Let $R = \{(n,\reals) \mid n\in\naturalsPlus\}$.
    Let $S = \{(n,\standardTopologyR) \mid n\in\naturalsPlus\}$.
    Let $P = \productTopologyIndexed{S}{R}$.
    Let $F = \{(x,z) \mid x\in X, z\in\productFamily{R} \mid
        \forall n\in\naturalsPlus. z(n) = \apply{\apply{f}{n}}{x}\}$.
    Hence $F$ is a topological embedding from $X$ to $\productFamily{R}$ with respect to $T$ and $P$
        by \cref{countable_embedding_theorem_real_power_unit_interval}.
    Show that $F$ is a homeomorphism between $X$ and $\img{F}{X}$ with respect to
        $T$ and $\subspaceTopology{P}{\img{F}{X}}$.
    \begin{subproof}
        Follows by \cref{topological_embedding}.
    \end{subproof}
    $\productFamily{R} = \realsPower{\naturalsPlus}$ by \cref{product_family_reals_power}.
    $P = \metricTopology{\omegaMetric}{\productFamily{R}}$
        by \cref{omega_metric_product_topology,product_family_reals_power}.
    $\omegaMetric$ is a metric on $\productFamily{R}$ by \cref{omega_metric_is_metric,metric_on}.
    $P$ is a topology on $\productFamily{R}$
        by \cref{basis_topology_is_topology,metric_topology,metric_basis_is_basis,basis_for_topology}.
    Therefore $\productFamily{R}$ is metrizable with respect to $P$ by \cref{metrizable}.
    $\img{F}{X}$ is a subspace of $\productFamily{R}$ with respect to $P$ and
        $\subspaceTopology{P}{\img{F}{X}}$ is a topology on $\img{F}{X}$
        by \cref{topological_embedding,funs_subseteq_rels,subseteq,rels_ran_subseteq,img_subseteq_ran,subseteq_transitive,subspace_of,subspace_topology_is_topology}.
    Hence $\img{F}{X}$ is metrizable with respect to $\subspaceTopology{P}{\img{F}{X}}$ by \cref{subspace_metrizable}.
    Therefore $X$ is metrizable with respect to $T$ by \cref{homeomorphism_preserves_metrizable}.
\end{proof}

% from §34 helper lemma 5: empty function embeds the empty space
\begin{lemma}\label{empty_domain_topological_embedding}
    Assume $T$ is a topology on $X$.
    Assume $T_1$ is a topology on $Y$.
    Assume $X = \emptyset$.
    Let $f = \emptyset$.
    Then $f$ is a topological embedding from $X$ to $Y$ with respect to $T$ and $T_1$.
\end{lemma}
\begin{proof}
    Show that for all $x,y_1,y_2$ such that $(x,y_1)\in f$ and $(x,y_2)\in f$ we have $y_1 = y_2$.
    \begin{subproof}
        Follows by \cref{emptyset}.
    \end{subproof}
    Hence $f$ is right-unique and $f$ is a function by \cref{emptyset,rightunique,relation}.
    $\dom{f} = X$ by \cref{dom_iff,emptyset,subseteq,emptyset_subseteq,subseteq_antisymmetric}.
    Hence $f\in\funs{X}{Y}$ by \cref{funs_intro,subseteq,emptyset}.

    Show that for all $x_1,x_2$ such that $x_1\in\dom{f}$ and $x_2\in\dom{f}$ and $f(x_1) = f(x_2)$ we have $x_1 = x_2$.
    \begin{subproof}
        Follows by \cref{subseteq,emptyset}.
    \end{subproof}
    Show that $f$ is injective.
    \begin{subproof}
        Follows by \cref{injective_function}.
    \end{subproof}

    Show that for all $V$ such that $V$ is open in $Y$ with respect to $T_1$ we have $\preimg{f}{V}$ is open in $X$ with respect to $T$.
    \begin{subproof}
        Follows by \cref{preimg_subseteq_dom,subseteq_transitive,subseteq_emptyset_iff,open_in,topology_on}.
    \end{subproof}
    Show that $f$ is continuous from $X$ to $Y$ with respect to $T$ and $T_1$.
    \begin{subproof}
        Follows by \cref{continuous}.
    \end{subproof}

    $\img{f}{X} = \emptyset$ by \cref{img_emptyset}.
    Hence $f\in\surj{X}{\img{f}{X}}$ by \cref{funs_intro,img_emptyset,subseteq,emptyset,funs_surj_iff,img_dom}.
    $f$ is a bijection from $X$ to $\img{f}{X}$ by \cref{bijections,injective_function}.
    $\converse{f}\in\funs{\img{f}{X}}{X}$ by \cref{bijective_converse_are_funs}.
    Hence $\img{f}{X}$ is a subspace of $Y$ with respect to $T_1$ and
        $\subspaceTopology{T_1}{\img{f}{X}}$ is a topology on $\img{f}{X}$
        by \cref{img_emptyset,emptyset_subseteq,subspace_of,subspace_topology_is_topology}.
    Show that for all $U$ such that $U$ is open in $X$ with respect to $T$ we have
        $\preimg{\converse{f}}{U}$ is open in $\img{f}{X}$ with respect to $\subspaceTopology{T_1}{\img{f}{X}}$.
    \begin{subproof}
        Follows by \cref{funs_elim,preimg_subseteq_dom,subseteq_transitive,subseteq_emptyset_iff,topology_on,subspace_topology_is_topology,subspace_of,open_in}.
    \end{subproof}
    Show that $\converse{f}$ is continuous from $\img{f}{X}$ to $X$ with respect to $\subspaceTopology{T_1}{\img{f}{X}}$ and $T$.
    \begin{subproof}
        Follows by \cref{continuous}.
    \end{subproof}
    Show that $f$ is continuous from $X$ to $\img{f}{X}$ with respect to $T$ and $\subspaceTopology{T_1}{\img{f}{X}}$.
    \begin{subproof}
        Follows by \cref{subseteq_refl,continuous_restrict_range}.
    \end{subproof}
    Show that $f$ is a homeomorphism between $X$ and $\img{f}{X}$ with respect to $T$ and $\subspaceTopology{T_1}{\img{f}{X}}$.
    \begin{subproof}
        Follows by \cref{homeomorphism}.
    \end{subproof}
    Show that $f$ is a topological embedding from $X$ to $Y$ with respect to $T$ and $T_1$.
    \begin{subproof}
        Follows by \cref{topological_embedding}.
    \end{subproof}
\end{proof}

% from §34 Item 2: imbedding theorem into real power
\begin{theorem}\label{urysohn_embedding_theorem_real_power}
    Assume $T$ is a topology on $X$.
    Assume $X$ has closed singletons with respect to $T$.
    Assume $J$ is inhabited.
    Assume $f$ is a function on $J$.
    Suppose for all $\alpha\in J$ we have $\apply{f}{\alpha}$ is continuous from $X$ to $\reals$ with respect to $T$ and $\standardTopologyR$.
    Suppose for all $x,U$ such that $x\in X$ and $U$ is a neighborhood of $x$ in $X$ with respect to $T$
        there exists $\alpha\in J$ such that
            $0 < \apply{\apply{f}{\alpha}}{x}$ and
            for all $y\in X\setminus U$ we have $\apply{\apply{f}{\alpha}}{y} = 0$.
    Let $R = \{(\alpha,\reals) \mid \alpha\in J\}$.
    Let $S = \{(\alpha,\standardTopologyR) \mid \alpha\in J\}$.
    Let $P = \productTopologyIndexed{S}{R}$.
    Let $F = \{(x,y) \mid x\in X, y\in\productFamily{R} \mid
        \forall \alpha\in J. \apply{y}{\alpha} = \apply{\apply{f}{\alpha}}{x}\}$.
    Then $F$ is a topological embedding from $X$ to $\productFamily{R}$ with respect to $T$ and $P$.
\end{theorem}
\begin{proof}
    $\standardTopologyR$ is a topology on $\reals$
        by \cref{standard_basis_is_basis,basis_topology_is_topology,standard_topology_reals}.
    $R$ is a function on $J$ by \cref{relation,rightunique,dom_intro,dom_iff,pair_eq_iff,subseteq,subseteq_antisymmetric}.
    $S$ is a function on $J$ by \cref{relation,rightunique,dom_intro,dom_iff,pair_eq_iff,subseteq,subseteq_antisymmetric}.
    Thus $\dom{S} = J$ and $\dom{R} = J$ by \cref{funs}.
    For all $\alpha\in J$ we have $S(\alpha) = \standardTopologyR$ and $R(\alpha) = \reals$ by \cref{function_apply_intro}.
    Hence $\productFamily{R} = \realsPower{J}$ by \cref{product_family_reals_power}.
    $P$ is a topology on $\productFamily{R}$
        by \cref{product_subbasis_is_subbasis,product_subbasis_finite_intersections_basis,basis_topology_is_topology,basis_for_topology,standard_topology_reals}.

    Show that for all $x,z_1,z_2$ such that $(x,z_1)\in F$ and $(x,z_2)\in F$ we have $z_1 = z_2$.
    \begin{subproof}
        Follows by \cref{pair_eq_iff,product_family_reals_power,subseteq,reals_power,tuple_power,functions_eq_iff}.
    \end{subproof}
    Hence $F$ is a function by \cref{pair_eq_iff,product_family_reals_power,subseteq,reals_power,tuple_power,functions_eq_iff,rightunique,relation}.
    Show that for all $x\in X$ there exists $z$ such that $(x,z)\in F$.
    \begin{subproof}
        Fix $x\in X$.
        Let $z = \{(\alpha,\apply{\apply{f}{\alpha}}{x}) \mid \alpha\in J\}$.
        For all $\alpha,y_1,y_2$ such that $(\alpha,y_1)\in z$ and $(\alpha,y_2)\in z$ we have $y_1 = y_2$ by \cref{pair_eq_iff}.
        Hence $z$ is a function by \cref{relation,rightunique,pair_eq_iff}.
        $\dom{z} = J$ by \cref{subseteq_antisymmetric,subseteq,dom_intro,dom_iff,pair_eq_iff}.
        For all $\alpha\in\dom{z}$ we have $z(\alpha)\in\reals$
            by \cref{dom_iff,pair_eq_iff,continuous,funs_type_apply,function_apply_intro}.
        $z\in\productFamily{R}$ by \cref{funs_intro,reals_power,tuple_power,subseteq}.
        For all $\alpha\in J$ we have $z(\alpha) = \apply{\apply{f}{\alpha}}{x}$ by \cref{function_apply_intro}.
        $(x,z)\in F$ by \cref{pair_eq_iff}.
    \end{subproof}
    Hence $\dom{F} = X$ by \cref{subseteq,dom_iff,pair_eq_iff,subseteq_antisymmetric}.
    Show that $F$ is a function from $X$ to $\productFamily{R}$.
    \begin{subproof}
        Follows by \cref{funs_intro,function_apply_iff,pair_eq_iff}.
    \end{subproof}

    $F$ is a function and $\dom{F} = X$ by \cref{funs_elim}.
    Show that for all $x_1,x_2$ such that $x_1\in\dom{F}$ and $x_2\in\dom{F}$ and $F(x_1) = F(x_2)$ we have $x_1 = x_2$.
    \begin{subproof}
        Fix $x_1$.
        Fix $x_2$ such that $x_1\in\dom{F}$ and $x_2\in\dom{F}$ and $F(x_1) = F(x_2)$.
        Hence $x_1\in X$ and $x_2\in X$.
        Suppose not.
        Hence $\{x_2\}$ is closed in $X$ with respect to $T$ by \cref{one_point_closed,singleton_subset_intro}.
        Let $U = X\setminus\{x_2\}$.
        Therefore $U$ is open in $X$ with respect to $T$
            by \cref{closed_in,topology_on,setminus_subseteq,double_relative_complement}.
        Hence $U$ is a neighborhood of $x_1$ in $X$ with respect to $T$
            by \cref{closed_in,topology_on,setminus_subseteq,double_relative_complement,singleton_elim,setminus_intro,neighborhood}.
        Take $\alpha\in J$ such that
            $0 < \apply{\apply{f}{\alpha}}{x_1}$ and
            for all $y\in X\setminus U$ we have $\apply{\apply{f}{\alpha}}{y} = 0$
            by assumption.
        Hence $\apply{\apply{f}{\alpha}}{x_2} = 0$
            by \cref{setminus_elim_right,singleton_intro,setminus_intro}.
        $\apply{\apply{F}{x_1}}{\alpha} = \apply{\apply{f}{\alpha}}{x_1}$ and
            $\apply{\apply{F}{x_2}}{\alpha} = \apply{\apply{f}{\alpha}}{x_2}$ by \cref{function_apply_elim,pair_eq_iff}.
        $0 < \apply{\apply{F}{x_1}}{\alpha}$ by \cref{real_lt_subst_right}.
        $\apply{\apply{F}{x_2}}{\alpha} = 0$ by \cref{real_lt_subst_right}.
        $\apply{\apply{F}{x_1}}{\alpha} \neq \apply{\apply{F}{x_2}}{\alpha}$ by \cref{real_lt_subst_right,real_add_ident,real_lt_irreflexive}.
        $\apply{\apply{F}{x_1}}{\alpha} = \apply{\apply{F}{x_2}}{\alpha}$.
        Contradiction.
    \end{subproof}
    Show that $F$ is injective.
    \begin{subproof}
        Follows by \cref{injective_function}.
    \end{subproof}

    Show that for all $\alpha\in J$ we have $\piIndex{R}{\alpha}\circ F = f(\alpha)$.
    \begin{subproof}
        Follows by \cref{functions_eq_iff,function_apply_elim,pair_eq_iff,continuous,funs_elim,funs_type_apply,circ_elem_intro,projection_map_indexed,function_apply_intro,projection_map_indexed_function,funs_circ,funs_is_function}.
    \end{subproof}
    Hence for all $\alpha\in J$ we have $\piIndex{R}{\alpha}\circ F$ is continuous from $X$ to $\reals$ with respect to $T$ and $\standardTopologyR$
        by \cref{continuous,continuous_eq_subst}.
    Show that $F$ is continuous from $X$ to $\productFamily{R}$ with respect to $T$ and $P$.
    \begin{subproof}
        Follows by \cref{continuous_map_into_indexed_product,standard_basis_is_basis,basis_topology_is_topology,standard_topology_reals,product_family_reals_power,funs,subseteq}.
    \end{subproof}

    $F$ is a function and $\dom{F} = X$ by \cref{funs_elim}.
    Hence $F\in\surj{X}{\img{F}{X}}$ by \cref{funs_intro,function_apply_elim,ran_intro,img_dom,funs_surj_iff}.
    $F$ is a bijection from $X$ to $\img{F}{X}$ by \cref{bijections,injective_function}.
    $\converse{F}\in\funs{\img{F}{X}}{X}$ by \cref{bijective_converse_are_funs}.
    Show that for all $z\in\img{F}{X}$ we have
        $\converse{F}$ is continuous at $z$ from $\img{F}{X}$ to $X$ with respect to $\subspaceTopology{P}{\img{F}{X}}$ and $T$.
    \begin{subproof}
        Fix $z\in\img{F}{X}$.
        Take $x\in X$ such that $(x,z)\in F$ by \cref{img_iff}.
        Hence $F(x) = z$ by \cref{function_apply_intro}.
        $(z,x)\in\converse{F}$ by \cref{converse_iff}.
        $\apply{\converse{F}}{z} = x$ by \cref{funs_elim,function_apply_intro}.
        Show that for all $U$ such that $U$ is a neighborhood of $\apply{\converse{F}}{z}$ in $X$ with respect to $T$
            there exists $W$ such that
                $W$ is a neighborhood of $z$ in $\img{F}{X}$ with respect to $\subspaceTopology{P}{\img{F}{X}}$ and
                $\img{\converse{F}}{W}\subseteq U$.
        \begin{subproof}
            Fix $U$ such that $U$ is a neighborhood of $\apply{\converse{F}}{z}$ in $X$ with respect to $T$.
            Take $\alpha\in J$ such that
                $0 < \apply{\apply{f}{\alpha}}{x}$ and
                for all $y\in X\setminus U$ we have $\apply{\apply{f}{\alpha}}{y} = 0$
                by assumption.
            Let $V = \intervalopen{0}{\apply{\apply{f}{\alpha}}{x} + 1}$.
            Let $D = \preimg{\piIndex{R}{\alpha}}{V}$.
            $\apply{\apply{f}{\alpha}}{x}\in\reals$ by \cref{continuous,funs_type_apply}.
            $0\in\reals$ by \cref{real_add_ident}.
            $1\in\reals$ and $0 < 1$ by \cref{real_one_in_naturals,real_naturals_subset_reals,subseteq,real_zero_lt_one}.
            Hence $\apply{\apply{f}{\alpha}}{x} + 1\in\reals$ by \cref{real_add_closed}.
            $0 + \apply{\apply{f}{\alpha}}{x} < 1 + \apply{\apply{f}{\alpha}}{x}$ by \cref{real_lt_add}.
            $0 + \apply{\apply{f}{\alpha}}{x} = \apply{\apply{f}{\alpha}}{x}$ and
                $1 + \apply{\apply{f}{\alpha}}{x} = \apply{\apply{f}{\alpha}}{x} + 1$
                by \cref{real_add_ident,real_add_comm}.
            $\apply{\apply{f}{\alpha}}{x} < \apply{\apply{f}{\alpha}}{x} + 1$
                by \cref{real_lt_subst_left,real_lt_subst_right}.
            $0 < \apply{\apply{f}{\alpha}}{x} + 1$ by \cref{real_lt_trans}.
            Therefore $D$ is open in $\productFamily{R}$ with respect to $P$
                by \cref{intervalopen_subset_reals,powerset_intro,standard_basis_reals,basis_elem_open_in_generated,basis_for_topology,standard_basis_is_basis,standard_topology_reals,projection_map_indexed_continuous,continuous,open_in,product_family_reals_power,subseteq,basis_topology_is_topology}.
            $\apply{\apply{F}{x}}{\alpha} = \apply{\apply{f}{\alpha}}{x}$ by \cref{function_apply_elim,pair_eq_iff}.
            $0 < \apply{\apply{F}{x}}{\alpha}$ by \cref{real_lt_subst_right}.
            $\apply{\apply{F}{x}}{\alpha} < \apply{\apply{f}{\alpha}}{x} + 1$ by \cref{real_lt_subst_left}.
            $F(x)\in\productFamily{R}$ by \cref{funs_type_apply}.
            $\apply{\apply{F}{x}}{\alpha}\in\reals$ by \cref{product_family_reals_power,reals_power,tuple_power,subseteq,funs_type_apply}.
            $\apply{\apply{F}{x}}{\alpha}\in V$ by \cref{intervalopen}.
            $(F(x),\apply{\apply{F}{x}}{\alpha})\in\piIndex{R}{\alpha}$ by \cref{projection_map_indexed,function_apply_elim}.
            $(z,\apply{\apply{F}{x}}{\alpha})\in\piIndex{R}{\alpha}$ by \cref{pair_eq_iff}.
            $z\in D$ by \cref{preimg_iff}.
            Let $W = \img{F}{X}\inter D$.
            $\img{F}{X}\subseteq\productFamily{R}$ by \cref{subseteq,img_iff,pair_eq_iff}.
            $\subspaceTopology{P}{\img{F}{X}}$ is a topology on $\img{F}{X}$ by \cref{subspace_topology_is_topology}.
            Hence $W$ is open in $\img{F}{X}$ with respect to $\subspaceTopology{P}{\img{F}{X}}$ by \cref{subspace_topology,open_in}.
            $z\in W$ by \cref{img_iff,inter_intro}.
            $W$ is a neighborhood of $z$ in $\img{F}{X}$ with respect to $\subspaceTopology{P}{\img{F}{X}}$ by \cref{neighborhood}.
            $\converse{F}$ is a function and $\dom{\converse{F}} = \img{F}{X}$ by \cref{funs_elim}.
            Show that for all $y$ such that $y\in\img{\converse{F}}{W}$ we have $y\in U$.
            \begin{subproof}
                Fix $y$ such that $y\in\img{\converse{F}}{W}$.
                Take $p\in\dom{\converse{F}}\inter W$ such that $y = \apply{\converse{F}}{p}$ by \cref{img_of_function_elim}.
                Hence $p\in W$ by \cref{inter_elim_right}.
                $p\in D$ and $p\in\img{F}{X}$ by \cref{inter_elim_left,inter_elim_right}.
                $p\in\dom{\converse{F}}$ by \cref{subseteq}.
                $(y,p)\in F$ by \cref{converse_iff,function_apply_elim}.
                $y\in X$ by \cref{dom_iff,funs_elim}.
                $p = F(y)$ by \cref{function_apply_intro}.
                $p\in\productFamily{R}$ by \cref{subseteq}.
                $R$ is a function and $\alpha\in\dom{R}$ by \cref{funs_elim}.
                $\piIndex{R}{\alpha}\in\funs{\productFamily{R}}{\reals}$ by \cref{projection_map_indexed_function}.
                $\piIndex{R}{\alpha}$ is a function by \cref{funs_elim}.
                Take $q\in V$ such that $(p,q)\in\piIndex{R}{\alpha}$ by \cref{preimg_iff}.
                $F(y)\in\productFamily{R}$ by \cref{funs_type_apply}.
                $(F(y),\apply{\apply{F}{y}}{\alpha})\in\piIndex{R}{\alpha}$ by \cref{projection_map_indexed,function_apply_elim}.
                $(p,\apply{\apply{F}{y}}{\alpha})\in\piIndex{R}{\alpha}$ by \cref{pair_eq_iff}.
                $q = \apply{\apply{F}{y}}{\alpha}$ by \cref{function_rightunique}.
                $0 < q$ by \cref{intervalopen}.
                $0 < \apply{\apply{F}{y}}{\alpha}$ by \cref{real_lt_subst_right}.
                $\apply{\apply{F}{y}}{\alpha} = \apply{\apply{f}{\alpha}}{y}$ by \cref{pair_eq_iff}.
                $0 < \apply{\apply{f}{\alpha}}{y}$ by \cref{real_lt_subst_right}.
                $y\in U$ by \cref{setminus_intro,real_lt_subst_right,real_lt_irreflexive}.
            \end{subproof}
            Hence $\img{\converse{F}}{W}\subseteq U$ by \cref{subseteq}.
        \end{subproof}
        Hence $\converse{F}$ is continuous at $z$ from $\img{F}{X}$ to $X$ with respect to $\subspaceTopology{P}{\img{F}{X}}$ and $T$
            by \cref{subseteq,img_iff,pair_eq_iff,subspace_topology_is_topology,continuous_at}.
    \end{subproof}
    $\img{F}{X}\subseteq\productFamily{R}$ by \cref{subseteq,img_iff,pair_eq_iff}.
    Show that $\converse{F}$ is continuous from $\img{F}{X}$ to $X$ with respect to $\subspaceTopology{P}{\img{F}{X}}$ and $T$.
    \begin{subproof}
        Follows by \cref{subspace_topology_is_topology,continuous_at_implies_continuous}.
    \end{subproof}
    $\img{F}{X}$ is a subspace of $\productFamily{R}$ with respect to $P$ by \cref{subspace_of}.
    Show that $F$ is continuous from $X$ to $\img{F}{X}$ with respect to $T$ and $\subspaceTopology{P}{\img{F}{X}}$.
    \begin{subproof}
        Follows by \cref{subseteq_refl,continuous_restrict_range}.
    \end{subproof}
    Show that $F$ is a homeomorphism between $X$ and $\img{F}{X}$ with respect to $T$ and $\subspaceTopology{P}{\img{F}{X}}$.
    \begin{subproof}
        Follows by \cref{homeomorphism}.
    \end{subproof}
    Show that $F$ is a topological embedding from $X$ to $\productFamily{R}$ with respect to $T$ and $P$.
    \begin{subproof}
        Follows by \cref{topological_embedding}.
    \end{subproof}
\end{proof}

% from §34 Item 2: imbedding theorem into unit cube
\begin{theorem}\label{urysohn_embedding_theorem_unit_cube}
    Assume $T$ is a topology on $X$.
    Assume $X$ has closed singletons with respect to $T$.
    Assume $J$ is inhabited.
    Assume $f$ is a function on $J$.
    Suppose for all $\alpha\in J$ we have $\apply{f}{\alpha}$ is continuous from $X$ to $\reals$ with respect to $T$ and $\standardTopologyR$.
    Suppose for all $x,U$ such that $x\in X$ and $U$ is a neighborhood of $x$ in $X$ with respect to $T$
        there exists $\alpha\in J$ such that
            $0 < \apply{\apply{f}{\alpha}}{x}$ and
            for all $y\in X\setminus U$ we have $\apply{\apply{f}{\alpha}}{y} = 0$.
    Let $I = \intervalclosed{0}{1}$.
    Suppose for all $\alpha\in J$ we have $\img{\apply{f}{\alpha}}{X}\subseteq I$.
    Let $C = \{(\alpha,I) \mid \alpha\in J\}$.
    Let $S = \{(\alpha,\subspaceTopology{\standardTopologyR}{I}) \mid \alpha\in J\}$.
    Let $P = \productTopologyIndexed{S}{C}$.
    Let $F = \{(x,y) \mid x\in X, y\in\productFamily{C} \mid
        \forall \alpha\in J. \apply{y}{\alpha} = \apply{\apply{f}{\alpha}}{x}\}$.
    Then $F$ is a topological embedding from $X$ to $\productFamily{C}$ with respect to $T$ and $P$.
\end{theorem}
\begin{proof}
    Let $T_I = \subspaceTopology{\standardTopologyR}{I}$.
    $\standardTopologyR$ is a topology on $\reals$
        by \cref{standard_basis_is_basis,basis_topology_is_topology,standard_topology_reals}.
    $I\subseteq\reals$ by \cref{intervalclosed,subseteq}.
    Hence $I$ is a subspace of $\reals$ with respect to $\standardTopologyR$ and
        $T_I$ is a topology on $I$ by \cref{subspace_of,subspace_topology_is_topology}.
    For all $\alpha\in J$ we have $\apply{f}{\alpha}$ is continuous from $X$ to $I$ with respect to $T$ and $T_I$
        by \cref{continuous,continuous_subspace_range}.
    For all $\alpha,y_1,y_2$ such that $(\alpha,y_1)\in C$ and $(\alpha,y_2)\in C$ we have $y_1 = y_2$ by \cref{pair_eq_iff}.
    Hence $C$ is a function on $J$ by \cref{rightunique,relation,dom_intro,dom_iff,pair_eq_iff,subseteq,subseteq_antisymmetric}.
    For all $\alpha,y_1,y_2$ such that $(\alpha,y_1)\in S$ and $(\alpha,y_2)\in S$ we have $y_1 = y_2$ by \cref{pair_eq_iff}.
    We have $S$ is a function on $J$ by \cref{rightunique,relation,dom_intro,dom_iff,pair_eq_iff,subseteq,subseteq_antisymmetric}.
    Thus $\dom{S} = J$ and $\dom{C} = J$ by \cref{funs}.
    For all $\alpha\in J$ we have $S(\alpha) = T_I$ and $C(\alpha) = I$ by \cref{function_apply_intro}.
    $P$ is a topology on $\productFamily{C}$
        by \cref{product_subbasis_is_subbasis,product_subbasis_finite_intersections_basis,basis_topology_is_topology,basis_for_topology,subspace_topology_is_topology}.

    For all $x,z_1,z_2$ such that $(x,z_1)\in F$ and $(x,z_2)\in F$ we have $z_1 = z_2$
        by \cref{pair_eq_iff,indexed_product,funs_is_function,funs,functions_eq_iff}.
    Hence $F$ is a function by \cref{pair_eq_iff,indexed_product,funs_is_function,funs,functions_eq_iff,rightunique,relation}.
    Show that for all $x\in X$ there exists $z$ such that $(x,z)\in F$.
    \begin{subproof}
        Fix $x\in X$.
        Let $z = \{(\alpha,\apply{\apply{f}{\alpha}}{x}) \mid \alpha\in J\}$.
        For all $\alpha,y_1,y_2$ such that $(\alpha,y_1)\in z$ and $(\alpha,y_2)\in z$ we have $y_1 = y_2$ by \cref{pair_eq_iff}.
        Hence $z$ is a function by \cref{relation,rightunique,pair_eq_iff}.
        We have $\dom{z} = J$ by \cref{subseteq_antisymmetric,subseteq,dom_intro,dom_iff,pair_eq_iff}.
        For all $\alpha\in\dom{z}$ we have $z(\alpha)\in I$
            by \cref{dom_iff,pair_eq_iff,continuous,funs_type_apply,function_apply_intro}.
        $z\in\funs{J}{I}$ by \cref{funs_intro}.
        Therefore $z\in\productFamily{C}$ by \cref{funs_intro,indexed_product,funs_elim,function_apply_intro,ran_intro,unions_iff}.
        For all $\alpha\in J$ we have $z(\alpha) = \apply{\apply{f}{\alpha}}{x}$ by \cref{function_apply_intro}.
        Hence $(x,z)\in F$ by \cref{pair_eq_iff}.
    \end{subproof}
    Hence $\dom{F} = X$ by \cref{subseteq,dom_iff,pair_eq_iff,subseteq_antisymmetric}.
    $F\in\funs{X}{\productFamily{C}}$ by \cref{funs_intro,function_apply_iff,pair_eq_iff}.

    $F$ is a function and $\dom{F} = X$ by \cref{funs_elim}.
    Show that for all $x_1,x_2$ such that $x_1\in\dom{F}$ and $x_2\in\dom{F}$ and $F(x_1) = F(x_2)$ we have $x_1 = x_2$.
    \begin{subproof}
        Fix $x_1$.
        Fix $x_2$ such that $x_1\in\dom{F}$ and $x_2\in\dom{F}$ and $F(x_1) = F(x_2)$.
        Hence $x_1\in X$ and $x_2\in X$.
        Suppose not.
        We have $\{x_2\}$ is closed in $X$ with respect to $T$ by \cref{one_point_closed,singleton_subset_intro}.
        Let $U = X\setminus\{x_2\}$.
        Hence $U$ is a neighborhood of $x_1$ in $X$ with respect to $T$
            by \cref{closed_in,topology_on,setminus_subseteq,double_relative_complement,singleton_elim,setminus_intro,neighborhood}.
        Take $\alpha\in J$ such that
            $0 < \apply{\apply{f}{\alpha}}{x_1}$ and
            for all $y\in X\setminus U$ we have $\apply{\apply{f}{\alpha}}{y} = 0$
            by assumption.
        $\apply{\apply{f}{\alpha}}{x_2} = 0$
            by \cref{setminus_elim_right,singleton_intro,setminus_intro}.
        We have $\apply{\apply{F}{x_1}}{\alpha} = \apply{\apply{f}{\alpha}}{x_1}$ and
            $\apply{\apply{F}{x_2}}{\alpha} = \apply{\apply{f}{\alpha}}{x_2}$ by \cref{function_apply_elim,pair_eq_iff}.
        $0 < \apply{\apply{F}{x_1}}{\alpha}$ by \cref{real_lt_subst_right}.
        $\apply{\apply{F}{x_2}}{\alpha} = 0$ by \cref{real_lt_subst_right}.
        Therefore $\apply{\apply{F}{x_1}}{\alpha} \neq \apply{\apply{F}{x_2}}{\alpha}$ by \cref{real_lt_subst_right,real_add_ident,real_lt_irreflexive}.
        $\apply{\apply{F}{x_1}}{\alpha} = \apply{\apply{F}{x_2}}{\alpha}$.
        Contradiction.
    \end{subproof}
    Hence $F$ is injective by \cref{injective_function}.

    Show that for all $\alpha\in J$ we have $\piIndex{C}{\alpha}\circ F = f(\alpha)$.
    \begin{subproof}
        Fix $\alpha\in J$.
        $\piIndex{C}{\alpha}\in\funs{\productFamily{C}}{I}$ by \cref{projection_map_indexed_function}.
        $\piIndex{C}{\alpha}\circ F\in\funs{X}{I}$ by \cref{funs_circ}.
        $\apply{f}{\alpha}\in\funs{X}{I}$ by \cref{continuous}.
        Show that for all $x\in X$ we have $\apply{\piIndex{C}{\alpha}\circ F}{x} = \apply{\apply{f}{\alpha}}{x}$.
        \begin{subproof}
            Fix $x\in X$.
            $\apply{F}{x}\in\productFamily{C}$ by \cref{funs_type_apply}.
            $(x,\apply{F}{x})\in F$ by \cref{function_apply_elim}.
            $\apply{\apply{F}{x}}{\alpha} = \apply{\apply{f}{\alpha}}{x}$ by \cref{pair_eq_iff}.
            $(\apply{F}{x},\apply{\apply{F}{x}}{\alpha})\in\piIndex{C}{\alpha}$ by \cref{projection_map_indexed}.
            $(x,\apply{\apply{F}{x}}{\alpha})\in\piIndex{C}{\alpha}\circ F$ by \cref{circ_elem_intro}.
            $\apply{\piIndex{C}{\alpha}\circ F}{x} = \apply{\apply{F}{x}}{\alpha}$ by \cref{function_apply_intro,funs_is_function}.
            $\apply{\piIndex{C}{\alpha}\circ F}{x} = \apply{\apply{f}{\alpha}}{x}$ by \cref{real_lt_subst_right}.
        \end{subproof}
        $\piIndex{C}{\alpha}\circ F = f(\alpha)$ by \cref{functions_eq_iff}.
    \end{subproof}
    Hence for all $\alpha\in J$ we have $\piIndex{C}{\alpha}\circ F$ is continuous from $X$ to $I$ with respect to $T$ and $T_I$
        by \cref{continuous,continuous_eq_subst}.
    Show that $F$ is continuous from $X$ to $\productFamily{C}$ with respect to $T$ and $P$.
    \begin{subproof}
        Follows by \cref{continuous_map_into_indexed_product,subspace_topology_is_topology,funs,subseteq}.
    \end{subproof}

    $F$ is a function and $\dom{F} = X$ by \cref{funs_elim}.
    Hence $F\in\surj{X}{\img{F}{X}}$ by \cref{funs_intro,function_apply_elim,ran_intro,img_dom,funs_surj_iff}.
    Thus $F$ is a bijection from $X$ to $\img{F}{X}$ by \cref{bijections,injective_function}.
    Hence $\converse{F}\in\funs{\img{F}{X}}{X}$ by \cref{bijective_converse_are_funs}.
    Show that for all $z\in\img{F}{X}$ we have
        $\converse{F}$ is continuous at $z$ from $\img{F}{X}$ to $X$ with respect to $\subspaceTopology{P}{\img{F}{X}}$ and $T$.
    \begin{subproof}
        Fix $z\in\img{F}{X}$.
        Take $x\in X$ such that $(x,z)\in F$ by \cref{img_iff}.
        Hence $F(x) = z$ by \cref{function_apply_intro}.
        Therefore $(z,x)\in\converse{F}$ by \cref{converse_iff}.
        $\apply{\converse{F}}{z} = x$ by \cref{funs_elim,function_apply_intro}.
        Show that for all $U$ such that $U$ is a neighborhood of $\apply{\converse{F}}{z}$ in $X$ with respect to $T$
            there exists $W$ such that
                $W$ is a neighborhood of $z$ in $\img{F}{X}$ with respect to $\subspaceTopology{P}{\img{F}{X}}$ and
                $\img{\converse{F}}{W}\subseteq U$.
        \begin{subproof}
            Fix $U$ such that $U$ is a neighborhood of $\apply{\converse{F}}{z}$ in $X$ with respect to $T$.
            Take $\alpha\in J$ such that
                $0 < \apply{\apply{f}{\alpha}}{x}$ and
                for all $y\in X\setminus U$ we have $\apply{\apply{f}{\alpha}}{y} = 0$
                by assumption.
            Let $V_0 = \intervalopen{0}{1 + 1}$.
            Let $V = I\inter V_0$.
            Let $D = \preimg{\piIndex{C}{\alpha}}{V}$.
            $0\in\reals$ and $1 + 1\in\reals$ and $0 < 1 + 1$
                by \cref{real_naturals_subset_reals,real_one_in_naturals,subseteq,real_add_ident,real_add_closed,real_zero_lt_one,real_lt_add,real_lt_trans}.
            Hence $V_0\in\standardTopologyR$
                by \cref{intervalopen_subset_reals,powerset_intro,standard_basis_reals,basis_elem_open_in_generated,basis_for_topology,standard_basis_is_basis,standard_topology_reals}.
            $V$ is open in $I$ with respect to $T_I$ by \cref{subspace_topology,open_in}.
            Thus $D$ is open in $\productFamily{C}$ with respect to $P$
                by \cref{projection_map_indexed_continuous,continuous,open_in,subspace_topology_is_topology,subseteq}.
            We have $\apply{\apply{F}{x}}{\alpha} = \apply{\apply{f}{\alpha}}{x}$ by \cref{function_apply_elim,pair_eq_iff}.
            $\apply{\apply{f}{\alpha}}{x}\in I$ by \cref{continuous,funs_type_apply}.
            Hence $\apply{\apply{f}{\alpha}}{x} < 1 + 1$ by \cref{intervalclosed,real_lt_add,real_add_ident,real_lt_trans,real_one_leq_naturalsplus,real_one_in_naturals,real_naturals_subset_reals,subseteq,real_lt_imp_leq}.
            $0 < \apply{\apply{F}{x}}{\alpha}$ by \cref{real_lt_subst_right}.
            $\apply{\apply{F}{x}}{\alpha} < 1 + 1$ by \cref{real_lt_subst_left}.
            $F(x)\in\productFamily{C}$ by \cref{funs_type_apply}.
            Hence $\apply{\apply{F}{x}}{\alpha}\in I$ by \cref{funs_type_apply,indexed_product}.
            We have $\apply{\apply{F}{x}}{\alpha}\in\reals$ by \cref{intervalclosed,subseteq}.
            $\apply{\apply{F}{x}}{\alpha}\in V_0$ by \cref{intervalopen}.
            Thus $\apply{\apply{F}{x}}{\alpha}\in V$ by \cref{inter_intro}.
            Hence $(F(x),\apply{\apply{F}{x}}{\alpha})\in\piIndex{C}{\alpha}$ by \cref{projection_map_indexed,function_apply_elim}.
            Therefore $(z,\apply{\apply{F}{x}}{\alpha})\in\piIndex{C}{\alpha}$ by \cref{pair_eq_iff}.
            $z\in D$ by \cref{preimg_iff}.
            Let $W = \img{F}{X}\inter D$.
            $W\in\subspaceTopology{P}{\img{F}{X}}$ by \cref{subspace_topology,open_in}.
            $\img{F}{X}\subseteq\productFamily{C}$ by \cref{subseteq,img_iff,pair_eq_iff}.
            $\subspaceTopology{P}{\img{F}{X}}$ is a topology on $\img{F}{X}$ by \cref{subspace_topology_is_topology}.
            Hence $W$ is open in $\img{F}{X}$ with respect to $\subspaceTopology{P}{\img{F}{X}}$ by \cref{open_in}.
            $z\in W$ by \cref{img_iff,inter_intro}.
            Thus $W$ is a neighborhood of $z$ in $\img{F}{X}$ with respect to $\subspaceTopology{P}{\img{F}{X}}$ by \cref{neighborhood}.
            $\converse{F}$ is a function and $\dom{\converse{F}} = \img{F}{X}$ by \cref{funs_elim}.
            Show that for all $y$ such that $y\in\img{\converse{F}}{W}$ we have $y\in U$.
            \begin{subproof}
                Fix $y$ such that $y\in\img{\converse{F}}{W}$.
                Take $p\in\dom{\converse{F}}\inter W$ such that $y = \apply{\converse{F}}{p}$ by \cref{img_of_function_elim}.
                Hence $p\in W$ by \cref{inter_elim_right}.
                We have $p\in D$ and $p\in\img{F}{X}$ by \cref{inter_elim_left,inter_elim_right}.
                $p\in\dom{\converse{F}}$ by \cref{subseteq}.
                Thus $(y,p)\in F$ by \cref{converse_iff,function_apply_elim}.
                Hence $y\in X$ by \cref{dom_iff,funs_elim}.
                Therefore $p = F(y)$ by \cref{function_apply_intro}.
                $p\in\productFamily{C}$ by \cref{subseteq}.
                $C$ is a function and $\alpha\in\dom{C}$ by \cref{funs_elim}.
                $\piIndex{C}{\alpha}\in\funs{\productFamily{C}}{I}$ by \cref{projection_map_indexed_function}.
                Hence $\piIndex{C}{\alpha}$ is a function by \cref{funs_elim}.
                Take $q\in V$ such that $(p,q)\in\piIndex{C}{\alpha}$ by \cref{preimg_iff}.
                $F(y)\in\productFamily{C}$ by \cref{funs_type_apply}.
                Hence $(F(y),\apply{\apply{F}{y}}{\alpha})\in\piIndex{C}{\alpha}$ by \cref{projection_map_indexed,function_apply_elim}.
                $(p,\apply{\apply{F}{y}}{\alpha})\in\piIndex{C}{\alpha}$ by \cref{pair_eq_iff}.
                Therefore $q = \apply{\apply{F}{y}}{\alpha}$ by \cref{function_rightunique}.
                Hence $0 < q$ by \cref{inter_elim_right,intervalopen}.
                $0 < \apply{\apply{F}{y}}{\alpha}$ by \cref{real_lt_subst_right}.
                Thus $\apply{\apply{F}{y}}{\alpha} = \apply{\apply{f}{\alpha}}{y}$ by \cref{pair_eq_iff}.
                Hence $0 < \apply{\apply{f}{\alpha}}{y}$ by \cref{real_lt_subst_right}.
                $y\in U$ by \cref{setminus_intro,real_lt_subst_right,real_lt_irreflexive}.
            \end{subproof}
            Hence $\img{\converse{F}}{W}\subseteq U$ by \cref{subseteq}.
        \end{subproof}
        Hence $\converse{F}$ is continuous at $z$ from $\img{F}{X}$ to $X$ with respect to $\subspaceTopology{P}{\img{F}{X}}$ and $T$
            by \cref{subseteq,img_iff,pair_eq_iff,subspace_topology_is_topology,continuous_at}.
    \end{subproof}
    $\img{F}{X}\subseteq\productFamily{C}$ by \cref{subseteq,img_iff,pair_eq_iff}.
    Show that $\converse{F}$ is continuous from $\img{F}{X}$ to $X$ with respect to $\subspaceTopology{P}{\img{F}{X}}$ and $T$.
    \begin{subproof}
        Follows by \cref{subspace_topology_is_topology,continuous_at_implies_continuous}.
    \end{subproof}
    $\img{F}{X}$ is a subspace of $\productFamily{C}$ with respect to $P$ by \cref{subspace_of}.
    Show that $F$ is continuous from $X$ to $\img{F}{X}$ with respect to $T$ and $\subspaceTopology{P}{\img{F}{X}}$.
    \begin{subproof}
        Follows by \cref{subseteq_refl,continuous_restrict_range}.
    \end{subproof}
    Show that $F$ is a homeomorphism between $X$ and $\img{F}{X}$ with respect to $T$ and $\subspaceTopology{P}{\img{F}{X}}$.
    \begin{subproof}
        Follows by \cref{homeomorphism}.
    \end{subproof}
    Show that $F$ is a topological embedding from $X$ to $\productFamily{C}$ with respect to $T$ and $P$.
    \begin{subproof}
        Follows by \cref{topological_embedding}.
    \end{subproof}
\end{proof}

% from §34 Item 3: family separates points from closed sets
\begin{definition}\label{separates_points_from_closed_sets}
    $f$ separates points from closed sets in $X$ with respect to $T$ and $J$ iff
        $T$ is a topology on $X$ and
        $f$ is a function on $J$ and
        for all $\alpha\in J$ we have $\apply{f}{\alpha}$ is continuous from $X$ to $\reals$ with respect to $T$ and $\standardTopologyR$ and
        for all $x,U$ such that $x\in X$ and $U$ is a neighborhood of $x$ in $X$ with respect to $T$
            there exists $\alpha\in J$ such that
                $0 < \apply{\apply{f}{\alpha}}{x}$ and
                for all $y\in X\setminus U$ we have $\apply{\apply{f}{\alpha}}{y} = 0$.
\end{definition}

% from §34 helper lemma 6: the unit interval is completely regular
\begin{lemma}\label{unit_interval_completely_regular}
    Let $I = \intervalclosed{0}{1}$.
    Let $T_I = \subspaceTopology{\standardTopologyR}{I}$.
    Then $I$ is completely regular with respect to $T_I$.
\end{lemma}
\begin{proof}
    $\standardTopologyR$ is a topology on $\reals$
        by \cref{standard_basis_is_basis,basis_topology_is_topology,standard_topology_reals}.
    $I\subseteq\reals$ by \cref{intervalclosed,subseteq}.
    Hence $I$ is a subspace of $\reals$ with respect to $\standardTopologyR$ and
        $T_I$ is a topology on $I$ by \cref{subspace_of,subspace_topology_is_topology}.

    $I$ is compact with respect to $T_I$
        by \cref{real_add_ident,real_one_in_naturals,real_naturals_subset_reals,subseteq,real_zero_lt_one,real_lt_imp_leq,real_intervalclosed_compact_standard}.

    $\realOrderTopology = \standardTopologyR$ and
        $\realOrderTopology = \basisTopology{\orderBasis{\realOrderRel}{\reals}}{\reals}$
        by \cref{real_order_topology_eq_standard,real_order_topology}.
    $\realOrderRel$ is a simple order relation on $\reals$ and
        $\reals$ has more than one element by \cref{real_order_relation_simple,linear_continuum,reals_linear_continuum}.
    Thus $\realOrderTopology$ is the order topology on $\reals$ with respect to $\realOrderRel$ by \cref{order_topology}.
    Hence $I$ is Hausdorff with respect to $T_I$
        by \cref{order_topology_hausdorff,real_order_topology_eq_standard,subspace_hausdorff}.
    Show that $I$ is normal with respect to $T_I$ and $I$ has closed singletons with respect to $T_I$.
    \begin{subproof}
        Follows by \cref{compact_hausdorff_normal,normal_space}.
    \end{subproof}

    Show that for all $x,A$ such that $x\in I$ and $A$ is closed in $I$ with respect to $T_I$ and $x\notin A$
        there exists $f$ such that
            $f$ is continuous from $I$ to $I$ with respect to $T_I$ and $T_I$ and
            $f(x) = 1$ and
            for all $y\in A$ we have $f(y) = 0$.
    \begin{subproof}
        Fix $x$.
        Fix $A$ such that $x\in I$ and $A$ is closed in $I$ with respect to $T_I$ and $x\notin A$.
        We have $\{x\}$ is closed in $I$ with respect to $T_I$ and $A$ is disjoint from $\{x\}$
            by \cref{normal_space,one_point_closed,singleton_subset_intro,disjoint,singleton_elim,pair_eq_iff}.
        Take $g$ such that
            $g$ is continuous from $I$ to $I$ with respect to $T_I$ and $T_I$ and
            $\img{g}{A}\subseteq \{0\}$ and
            $\img{g}{\{x\}}\subseteq \{1\}$
            by \cref{urysohn_unit_interval}.
        Show that $g(x) = 1$ and for all $y\in A$ we have $g(y) = 0$.
        \begin{subproof}
            Follows by \cref{continuous,funs_is_function,funs,subseteq,inter_intro,singleton_intro,img_of_function_intro,singleton_elim,closed_in}.
        \end{subproof}
        Follows by definition.
    \end{subproof}
    Show that $I$ is completely regular with respect to $T_I$.
    \begin{subproof}
        Follows by \cref{completely_regular_space}.
    \end{subproof}
\end{proof}

% from §34 helper lemma 7: homeomorphism preserves complete regularity
\begin{lemma}\label{homeomorphism_preserves_completely_regular}
    Assume $T$ is a topology on $X$.
    Assume $S$ is a topology on $Y$.
    Suppose $f$ is a homeomorphism between $X$ and $Y$ with respect to $T$ and $S$.
    If $Y$ is completely regular with respect to $S$ then $X$ is completely regular with respect to $T$.
\end{lemma}
\begin{proof}
    Assume $Y$ is completely regular with respect to $S$.
    $f\in\funs{X}{Y}$ and $f$ is a function and $\dom{f} = X$
        by \cref{homeomorphism,bijections,bijections_to_funs,funs_is_function,funs}.
    We have $f$ is injective by \cref{homeomorphism,bijections}.

    Show that for all $x$ such that $x\in X$ we have $\{x\}$ is closed in $X$ with respect to $T$.
    \begin{subproof}
        Fix $x$ such that $x\in X$.
        Let $y = f(x)$.
        We have $\{y\}$ is closed in $Y$ with respect to $S$
            by \cref{funs_type_apply,completely_regular_space,one_point_closed,singleton_subset_intro}.
        $\preimg{f}{\{y\}}\subseteq\{x\}$
            by \cref{subseteq,preimg_iff,function_apply_iff,singleton_elim,subseteq,pair_eq_iff,injective_function,singleton_intro}.
        $\{x\}\subseteq\preimg{f}{\{y\}}$
            by \cref{subseteq,singleton_elim,subseteq,pair_eq_iff,funs,function_apply_elim,singleton_intro,preimg_iff}.
        Hence $\{x\}$ is closed in $X$ with respect to $T$
            by \cref{subseteq_antisymmetric,closed_in,homeomorphism,continuous,preimg_setminus_function,preimg_subseteq_dom,funs,subseteq}.
    \end{subproof}
    Hence $X$ has closed singletons with respect to $T$ by \cref{one_point_closed}.

    Let $I = \intervalclosed{0}{1}$.
    Let $T_I = \subspaceTopology{\standardTopologyR}{I}$.
    Show that for all $x,A$ such that $x\in X$ and $A$ is closed in $X$ with respect to $T$ and $x\notin A$
        there exists $k$ such that
            $k$ is continuous from $X$ to $I$ with respect to $T$ and $T_I$ and
            $k(x) = 1$ and
            for all $a\in A$ we have $k(a) = 0$.
    \begin{subproof}
        Fix $x$.
        Fix $A$ such that $x\in X$ and $A$ is closed in $X$ with respect to $T$ and $x\notin A$.
        Let $y = f(x)$.
        Let $B = \img{f}{A}$.
        Hence $y\in Y$ by \cref{funs_type_apply}.
        We have $B$ is closed in $Y$ with respect to $S$
            by \cref{homeomorphism,continuous,closed_in,preimg_setminus_function,preimg_subseteq_dom,funs,subseteq,funs_is_relation,preim_eq_img_of_converse,converse_converse_eq}.
        We have $y\notin B$ by \cref{img_iff,closed_in,subseteq,function_apply_intro,pair_eq_iff,injective_function}.
        Take $q$ such that
            $q$ is continuous from $Y$ to $I$ with respect to $S$ and $T_I$ and
            $q(y) = 1$ and
            for all $b\in B$ we have $q(b) = 0$
            by \cref{completely_regular_space}.
        $q\in\funs{Y}{I}$ and $q$ is a function and $\dom{q} = Y$ by \cref{continuous,funs_is_function,funs}.
        $f$ is continuous from $X$ to $Y$ with respect to $T$ and $S$ and
            $f\in\funs{X}{Y}$ and $f$ is a function and $\dom{f} = X$
            by \cref{homeomorphism,continuous,funs_is_function,funs}.
        Let $k = q\circ f$.
        $q$ is composable with $f$ by \cref{continuous,funs,funs_is_function,funs_ran,subseteq}.
        Show that $k$ is continuous from $X$ to $I$ with respect to $T$ and $T_I$ and $k(x) = 1$.
        \begin{subproof}
            Follows by \cref{continuous_composite,circ_apply,funs_elim,pair_eq_iff}.
        \end{subproof}
        Show that for all $a\in A$ we have $k(a) = 0$.
        \begin{subproof}
            Follows by \cref{closed_in,subseteq,function_apply_elim,img_iff,circ_apply,funs_elim}.
        \end{subproof}
        Follows by definition.
    \end{subproof}
    Show that $X$ is completely regular with respect to $T$.
    \begin{subproof}
        Follows by \cref{completely_regular_space}.
    \end{subproof}
\end{proof}

% from §34 helper lemma 8: a subsingleton space is completely regular
\begin{lemma}\label{subsingleton_space_completely_regular}
    Assume $T$ is a topology on $X$.
    Suppose $X\subseteq \{p\}$.
    Then $X$ is completely regular with respect to $T$.
\end{lemma}
\begin{proof}
    Show that for all $x$ such that $x\in X$ we have $\{x\}$ is closed in $X$ with respect to $T$.
    \begin{subproof}
        Follows by \cref{subseteq,singleton_elim,pair_eq_iff,subseteq_antisymmetric,closed_whole_space}.
    \end{subproof}
    Hence $X$ has closed singletons with respect to $T$ by \cref{one_point_closed}.

    Let $I = \intervalclosed{0}{1}$.
    Let $T_I = \subspaceTopology{\standardTopologyR}{I}$.
    $\standardTopologyR$ is a topology on $\reals$
        by \cref{standard_basis_is_basis,basis_topology_is_topology,standard_topology_reals}.
    Hence $I$ is a subspace of $\reals$ with respect to $\standardTopologyR$ and
        $T_I$ is a topology on $I$ by \cref{intervalclosed,subseteq,subspace_of,subspace_topology_is_topology}.
    We have $1\in I$ by \cref{real_add_ident,real_one_in_naturals,real_naturals_subset_reals,subseteq,real_zero_lt_one,real_lt_imp_leq,intervalclosed}.

    Show that for all $x,A$ such that $x\in X$ and $A$ is closed in $X$ with respect to $T$ and $x\notin A$
        there exists $f$ such that
            $f$ is continuous from $X$ to $I$ with respect to $T$ and $T_I$ and
            $f(x) = 1$ and
            for all $y\in A$ we have $f(y) = 0$.
    \begin{subproof}
        Fix $x$.
        Fix $A$ such that $x\in X$ and $A$ is closed in $X$ with respect to $T$ and $x\notin A$.
        Hence $A = \emptyset$ by \cref{subseteq,closed_in,singleton_elim,pair_eq_iff,subseteq_emptyset_iff}.
        Let $f = \{(y,1) \mid y\in X\}$.
        $f$ is a function from $X$ to $I$ and for all $y\in X$ we have $f(y) = 1$
            by \cref{constant_map_function,funs_is_function,function_apply_intro}.
        Show that $f$ is continuous from $X$ to $I$ with respect to $T$ and $T_I$ and $f(x) = 1$.
        \begin{subproof}
            Follows by \cref{constant_continuous,constant_map_function,funs_is_function,function_apply_intro}.
        \end{subproof}
        Show that for all $y\in A$ we have $f(y) = 0$.
        \begin{subproof}
            Follows by \cref{pair_eq_iff,emptyset,subseteq_emptyset_iff}.
        \end{subproof}
        Follows by definition.
    \end{subproof}
    Show that $X$ is completely regular with respect to $T$.
    \begin{subproof}
        Follows by \cref{completely_regular_space}.
    \end{subproof}
\end{proof}

% from §34 Item 4: completely regular spaces embed in cubes
\begin{theorem}\label{completely_regular_embedding_characterization}
    Assume $T$ is a topology on $X$.
    Then $X$ is completely regular with respect to $T$ iff
        there exist $J,C,S,Y,h$ such that
            $C = \{(\alpha,\intervalclosed{0}{1}) \mid \alpha\in J\}$ and
            $S = \{(\alpha,\subspaceTopology{\standardTopologyR}{\intervalclosed{0}{1}}) \mid \alpha\in J\}$ and
            $Y$ is a subspace of $\productFamily{C}$ with respect to $\productTopologyIndexed{S}{C}$ and
            $h$ is a homeomorphism between $X$ and $Y$ with respect to $T$ and $\subspaceTopology{\productTopologyIndexed{S}{C}}{Y}$.
\end{theorem}
\begin{proof}
    Show that if $X$ is completely regular with respect to $T$ then
        there exist $J,C,S,Y,h$ such that
            $C = \{(\alpha,\intervalclosed{0}{1}) \mid \alpha\in J\}$ and
            $S = \{(\alpha,\subspaceTopology{\standardTopologyR}{\intervalclosed{0}{1}}) \mid \alpha\in J\}$ and
            $Y$ is a subspace of $\productFamily{C}$ with respect to $\productTopologyIndexed{S}{C}$ and
            $h$ is a homeomorphism between $X$ and $Y$ with respect to $T$ and $\subspaceTopology{\productTopologyIndexed{S}{C}}{Y}$.
    \begin{subproof}
        Assume $X$ is completely regular with respect to $T$.
        Let $I = \intervalclosed{0}{1}$.
        Let $T_I = \subspaceTopology{\standardTopologyR}{I}$.
        Let $j = \identity{I}$.
        $\standardTopologyR$ is a topology on $\reals$
            by \cref{standard_basis_is_basis,basis_topology_is_topology,standard_topology_reals}.
        $I\subseteq\reals$ by \cref{intervalclosed,subseteq}.
        Hence $I$ is a subspace of $\reals$ with respect to $\standardTopologyR$ and
            $T_I$ is a topology on $I$ by \cref{subspace_of,subspace_topology_is_topology}.
        Let $J_1 = \{g\in\pow{X\times\reals} \mid \text{$g$ is continuous from $X$ to $\reals$ with respect to $T$ and $\standardTopologyR$ and
            $\img{g}{X}\subseteq I$}\}$.
        Let $z = \{(x,0) \mid x\in X\}$.
        $0\in\reals$ and $1\in\reals$ and $0 < 1$
            by \cref{real_add_ident,real_one_in_naturals,real_naturals_subset_reals,subseteq,real_zero_lt_one}.
        Thus $0\in I$ by \cref{real_lt_imp_leq,intervalclosed}.
        $z$ is a function from $X$ to $I$ and for all $x\in X$ we have $z(x) = 0$
            by \cref{constant_map_function,funs_is_function,function_apply_intro}.
        Hence $z$ is continuous from $X$ to $I$ with respect to $T$ and $T_I$ by \cref{constant_continuous}.
        Let $z_1 = j\circ z$.
        $z_1$ is continuous from $X$ to $\reals$ with respect to $T$ and $\standardTopologyR$
            by \cref{continuous_expand_range}.
        $z_1\in\funs{X}{\reals}$ and $z_1$ is a function and $\dom{z_1} = X$ and
            $z_1\in\pow{X\times\reals}$
            by \cref{continuous,funs_is_function,funs,funs_subseteq_rels,subseteq,rels_elim,powerset_intro}.
        $z\in\funs{X}{I}$ and $z$ is a function by \cref{continuous,funs_is_function}.
        Hence $\ran{z}\subseteq\dom{j}$ by \cref{funs_subseteq_rels,subseteq,rels_ran_subseteq,id_dom}.
        $j$ is right-unique and $j$ is a relation by \cref{id_rightunique,id_is_relation}.
        For all $x\in X$ we have $z_1(x) = z(x)$ by \cref{circ_apply,funs_elim,id_apply}.
        Show that for all $u$ such that $u\in\img{z_1}{X}$ we have $u\in I$.
        \begin{subproof}
            Trivial.
        \end{subproof}
        Hence $\img{z_1}{X}\subseteq I$ by \cref{subseteq}.
        Therefore $J_1$ is inhabited.
        Let $f = \identity{J_1}$.
        Hence $f$ is a function on $J_1$ by \cref{id_is_relation,id_rightunique,id_dom,relation,rightunique,pair_eq_iff,dom_iff,subseteq,dom_intro,subseteq_antisymmetric}.
        For all $\alpha\in J_1$ we have $\apply{f}{\alpha}$ is continuous from $X$ to $\reals$ with respect to $T$ and $\standardTopologyR$
            by \cref{id_apply,id_elem_funs}.
        Show that for all $x,U$ such that $x\in X$ and $U$ is a neighborhood of $x$ in $X$ with respect to $T$
            there exists $\alpha\in J_1$ such that
                $0 < \apply{\apply{f}{\alpha}}{x}$ and
                for all $y\in X\setminus U$ we have $\apply{\apply{f}{\alpha}}{y} = 0$.
        \begin{subproof}
            Fix $x$.
            Fix $U$ such that $x\in X$ and $U$ is a neighborhood of $x$ in $X$ with respect to $T$.
            We have $U$ is open in $X$ with respect to $T$ and $x\in U$ by \cref{neighborhood}.
            Let $A = X\setminus U$.
            Hence $A$ is closed in $X$ with respect to $T$
                by \cref{open_in,topology_on,subseteq,powerset_elim,double_relative_complement,pair_eq_iff,setminus_subseteq,closed_in}.
            $x\notin A$ by \cref{setminus_elim_right}.
            Take $g$ such that
                $g$ is continuous from $X$ to $I$ with respect to $T$ and $T_I$ and
                $g(x) = 1$ and
                for all $y\in A$ we have $g(y) = 0$
                by \cref{completely_regular_space}.
            Let $g_1 = j\circ g$.
            $g_1$ is continuous from $X$ to $\reals$ with respect to $T$ and $\standardTopologyR$
                by \cref{continuous_expand_range}.
            $g_1\in\funs{X}{\reals}$ and $g_1$ is a function and $\dom{g_1} = X$ and
                $g_1\in\pow{X\times\reals}$
                by \cref{continuous,funs_is_function,funs,funs_subseteq_rels,subseteq,rels_elim,powerset_intro}.
            $g\in\funs{X}{I}$ and $g$ is a function by \cref{continuous,funs_is_function}.
            Hence $\ran{g}\subseteq\dom{j}$ by \cref{funs_subseteq_rels,subseteq,rels_ran_subseteq,id_dom}.
            For all $y\in X$ we have $g_1(y) = g(y)$ by \cref{circ_apply,funs_elim,id_apply}.
            $g_1$ is a function and $\dom{g_1} = X$ by \cref{funs_elim}.
            Show that for all $u$ such that $u\in\img{g_1}{X}$ we have $u\in I$.
            \begin{subproof}
                Fix $u$ such that $u\in\img{g_1}{X}$.
                Take $y\in\dom{g_1}\inter X$ such that $u = g_1(y)$ by \cref{img_of_function_elim}.
                Hence $y\in X$ by \cref{inter_elim_right}.
                $g(y)\in I$ by \cref{continuous,funs_type_apply}.
                We have $u = g(y)$.
                Therefore $u\in I$ by \cref{continuous,funs_type_apply}.
            \end{subproof}
            Hence $\img{g_1}{X}\subseteq I$ by \cref{subseteq}.
            Therefore $g_1\in J_1$.
            Take $\alpha\in J_1$ such that $\alpha = g_1$.
            Hence $\apply{f}{\alpha} = g_1$ by \cref{id_apply,id_elem_funs,pair_eq_iff}.
            $0 < \apply{\apply{f}{\alpha}}{x}$ by \cref{pair_eq_iff}.
            For all $y\in X\setminus U$ we have $\apply{\apply{f}{\alpha}}{y} = 0$.
        \end{subproof}
        For all $\alpha\in J_1$ we have $\img{\apply{f}{\alpha}}{X}\subseteq I$ by \cref{id_apply,id_elem_funs}.
        Let $C_1 = \{(\alpha,\intervalclosed{0}{1}) \mid \alpha\in J_1\}$.
        Let $S_1 = \{(\alpha,\subspaceTopology{\standardTopologyR}{\intervalclosed{0}{1}}) \mid \alpha\in J_1\}$.
        For all $\alpha,y_1,y_2$ such that $(\alpha,y_1)\in C_1$ and $(\alpha,y_2)\in C_1$ we have $y_1 = y_2$ by \cref{pair_eq_iff}.
        Hence $C_1$ is a function on $J_1$ by \cref{rightunique,relation,dom_intro,dom_iff,pair_eq_iff,subseteq,subseteq_antisymmetric}.
        For all $\alpha,y_1,y_2$ such that $(\alpha,y_1)\in S_1$ and $(\alpha,y_2)\in S_1$ we have $y_1 = y_2$ by \cref{pair_eq_iff}.
        We have $S_1$ is a function on $J_1$ by \cref{rightunique,relation,dom_intro,dom_iff,pair_eq_iff,subseteq,subseteq_antisymmetric}.
        Hence $\dom{C_1} = J_1$ and $\dom{S_1} = J_1$ by \cref{funs}.
        For all $\alpha\in J_1$ we have $C_1(\alpha) = I$ and $S_1(\alpha) = T_I$ by \cref{function_apply_intro}.
        Let $P = \productTopologyIndexed{S_1}{C_1}$.
        $P$ is a topology on $\productFamily{C_1}$
            by \cref{product_subbasis_is_subbasis,product_subbasis_finite_intersections_basis,basis_topology_is_topology,basis_for_topology,subspace_topology_is_topology}.
        Let $F = \{(x,y) \mid x\in X, y\in\productFamily{C_1} \mid
            \forall \alpha\in J_1. \apply{y}{\alpha} = \apply{\apply{f}{\alpha}}{x}\}$.
        Show that $F$ is a topological embedding from $X$ to $\productFamily{C_1}$ with respect to $T$ and $P$.
        \begin{subproof}
            Follows by \cref{completely_regular_space,urysohn_embedding_theorem_unit_cube}.
        \end{subproof}
        Therefore $\img{F}{X}\subseteq\productFamily{C_1}$
            by \cref{topological_embedding,funs_subseteq_rels,subseteq,rels_ran_subseteq,img_subseteq_ran,subseteq_transitive}.
        Let $Y_1 = \img{F}{X}$.
        Let $h_1 = F$.
        Show that $h_1$ is a homeomorphism between $X$ and $Y_1$ with respect to $T$ and $\subspaceTopology{P}{Y_1}$.
        \begin{subproof}
            Follows by \cref{topological_embedding,pair_eq_iff}.
        \end{subproof}
        Show that $Y_1$ is a subspace of $\productFamily{C_1}$ with respect to $\productTopologyIndexed{S_1}{C_1}$.
        \begin{subproof}
            Follows by \cref{subspace_of,pair_eq_iff}.
        \end{subproof}
        Show that $h_1$ is a homeomorphism between $X$ and $Y_1$ with respect to $T$ and $\subspaceTopology{\productTopologyIndexed{S_1}{C_1}}{Y_1}$.
        \begin{subproof}
            Follows by \cref{pair_eq_iff}.
        \end{subproof}
        Take $J,C,S,Y,h$ such that
            $J = J_1$ and $C = C_1$ and $S = S_1$ and $Y = Y_1$ and $h = h_1$.
        Follows by \cref{pair_eq_iff}.
    \end{subproof}

    Show that if there exist $J,C,S,Y,h$ such that
        $C = \{(\alpha,\intervalclosed{0}{1}) \mid \alpha\in J\}$ and
        $S = \{(\alpha,\subspaceTopology{\standardTopologyR}{\intervalclosed{0}{1}}) \mid \alpha\in J\}$ and
        $Y$ is a subspace of $\productFamily{C}$ with respect to $\productTopologyIndexed{S}{C}$ and
        $h$ is a homeomorphism between $X$ and $Y$ with respect to $T$ and $\subspaceTopology{\productTopologyIndexed{S}{C}}{Y}$
    then $X$ is completely regular with respect to $T$.
    \begin{subproof}
        Assume there exist $J,C,S,Y,h$ such that
            $C = \{(\alpha,\intervalclosed{0}{1}) \mid \alpha\in J\}$ and
            $S = \{(\alpha,\subspaceTopology{\standardTopologyR}{\intervalclosed{0}{1}}) \mid \alpha\in J\}$ and
            $Y$ is a subspace of $\productFamily{C}$ with respect to $\productTopologyIndexed{S}{C}$ and
            $h$ is a homeomorphism between $X$ and $Y$ with respect to $T$ and $\subspaceTopology{\productTopologyIndexed{S}{C}}{Y}$.
        Take $J,C,S,Y,h$ such that
            $C = \{(\alpha,\intervalclosed{0}{1}) \mid \alpha\in J\}$ and
            $S = \{(\alpha,\subspaceTopology{\standardTopologyR}{\intervalclosed{0}{1}}) \mid \alpha\in J\}$ and
            $Y$ is a subspace of $\productFamily{C}$ with respect to $\productTopologyIndexed{S}{C}$ and
            $h$ is a homeomorphism between $X$ and $Y$ with respect to $T$ and $\subspaceTopology{\productTopologyIndexed{S}{C}}{Y}$
            by assumption.
        Let $I = \intervalclosed{0}{1}$.
        Let $T_I = \subspaceTopology{\standardTopologyR}{I}$.
        Let $P = \productTopologyIndexed{S}{C}$.
        For all $\alpha,y_1,y_2$ such that $(\alpha,y_1)\in C$ and $(\alpha,y_2)\in C$ we have $y_1 = y_2$ by \cref{pair_eq_iff}.
        Hence $C$ is a function on $J$ by \cref{rightunique,relation,dom_intro,dom_iff,pair_eq_iff,subseteq,subseteq_antisymmetric}.
        For all $\alpha,y_1,y_2$ such that $(\alpha,y_1)\in S$ and $(\alpha,y_2)\in S$ we have $y_1 = y_2$ by \cref{pair_eq_iff}.
        We have $S$ is a function on $J$ by \cref{rightunique,relation,dom_intro,dom_iff,pair_eq_iff,subseteq,subseteq_antisymmetric}.
        Hence $\dom{C} = J$ and $\dom{S} = J$ by \cref{funs}.
        For all $\alpha\in J$ we have $C(\alpha) = I$ and $S(\alpha) = T_I$ by \cref{function_apply_intro}.

        Show that $Y$ is completely regular with respect to $\subspaceTopology{P}{Y}$.
        \begin{subproof}
            \begin{byCase}
            \caseOf{$J = \emptyset$.}
                Show that for all $u$ such that $u\in\productFamily{C}$ we have $u\in\{\emptyset\}$.
                \begin{subproof}
                    Trivial.
                \end{subproof}
                Show that $Y$ is completely regular with respect to $\subspaceTopology{P}{Y}$.
                \begin{subproof}
                    Follows by \cref{subseteq,subspace_of,subseteq_transitive,subspace_topology_is_topology,subsingleton_space_completely_regular}.
                \end{subproof}
            \caseOf{$J\neq\emptyset$.}
                Show that $Y$ is completely regular with respect to $\subspaceTopology{P}{Y}$.
                \begin{subproof}
                    Follows by \cref{nonempty_inhabited,function_apply_intro,unit_interval_completely_regular,pair_eq_iff,product_completely_regular,subspace_completely_regular}.
                \end{subproof}
            \end{byCase}
        \end{subproof}
        Hence $X$ is completely regular with respect to $T$
            by \cref{completely_regular_space,pair_eq_iff,homeomorphism_preserves_completely_regular}.
    \end{subproof}
    Follows by definition.
\end{proof}
\section{§35 The Tietze Extension Theorem}

% from §35 helper lemma 1: closed interval is closed in the standard topology on reals
\begin{lemma}\label{real_intervalclosed_closed_standard}
    Suppose $a\in\reals$.
    Suppose $b\in\reals$.
    Suppose $a \leq b$.
    Let $I = \intervalclosed{a}{b}$.
    Then $I$ is closed in $\reals$ with respect to $\standardTopologyR$.
\end{lemma}
\begin{proof}
    Follows by \cref{real_intervalclosed_compact_standard,intervalclosed,subseteq,standard_metric_is_metric,metric_space_hausdorff,standard_metric_topology,standard_basis_is_basis,basis_topology_is_topology,standard_topology_reals,compact_subspace_closed}.
\end{proof}

% from §35 helper lemma 2: interval-valued maps can be regarded as real-valued maps
\begin{lemma}\label{continuous_interval_expand_range}
    Assume $I = \intervalclosed{a}{b}$.
    Assume $T_I = \subspaceTopology{\standardTopologyR}{I}$.
    Assume $f$ is continuous from $X$ to $I$ with respect to $T$ and $T_I$.
    Let $j = \identity{I}$.
    Then $j\circ f$ is continuous from $X$ to $\reals$ with respect to $T$ and $\standardTopologyR$.
\end{lemma}
\begin{proof}
    Follows by \cref{intervalclosed,subseteq,subspace_of,standard_basis_is_basis,basis_topology_is_topology,standard_topology_reals,continuous_expand_range}.
\end{proof}

% from §35 helper lemma 6: twice the step width is bounded by the next error
\begin{lemma}\label{tietze_interval_double_step_width_leq_error}
    Assume $n\in\naturalsPlus$.
    Let $r = (1 + 1) / (n + 1)$.
    Let $e = (1 + 1) / ((n + 1) + 1)$.
    Let $d = r - e$.
    Then $d + d \leq e$.
\end{lemma}
\begin{proof}
    $n + 1\in\naturalsPlus$ and $((n + 1) + 1)\in\naturalsPlus$
        by \cref{real_naturals_closed_succ}.
    $0\in\reals$ and $1\in\reals$ and $n\in\reals$ and $(n + 1)\in\reals$ and $((n + 1) + 1)\in\reals$
        by \cref{real_add_ident,real_one_in_naturals,real_naturals_subset_reals,subseteq,real_naturals_closed_succ}.
    $0 < n$ and $0 < n + 1$ and $0 < ((n + 1) + 1)$
        by \cref{naturalsplus_positive_real,real_naturals_closed_succ,subseteq}.
    Hence $n + 1 \neq 0$ and $((n + 1) + 1) \neq 0$ by \cref{real_pos_nonzero}.
    We have $\inv{(n + 1)}\in\reals$ and $\inv{((n + 1) + 1)}\in\reals$
        by \cref{real_mul_inverse}.
    Hence $0 < \inv{(n + 1)}$ and $0 < \inv{((n + 1) + 1)}$ by \cref{real_inv_pos}.
    Let $c = \inv{(n + 1)} \rmul \inv{((n + 1) + 1)}$.
    $c\in\reals$ and $0 < c$ by \cref{real_mul_closed,real_mul_pos_pos_gt}.
    $1 + 1\in\reals$ and $0 < 1 + 1$
        by \cref{real_add_closed,real_zero_lt_one,real_pos_add,real_mul_ident}.
    Let $u = (1 + 1) \rmul c$.
    $u\in\reals$ and $0 < u$ by \cref{real_mul_closed,real_mul_pos_pos_gt}.
    Thus $e\in\reals$ by \cref{real_div,real_mul_closed}.

    $((n + 1) \rmul c) = ((n + 1) \rmul \inv{(n + 1)}) \rmul \inv{((n + 1) + 1)}$
        by \cref{real_mul_assoc}.
    Hence $(n + 1) \rmul c = 1 \rmul \inv{((n + 1) + 1)}$ by \cref{real_mul_inverse}.
    We have $(n + 1) \rmul c = \inv{((n + 1) + 1)}$ by \cref{real_mul_ident}.
    Therefore $(n + 1) \rmul c = 1 / ((n + 1) + 1)$ by \cref{real_div,real_mul_ident}.
    $((n + 1) \rmul u) = ((n + 1) \rmul c) \rmul (1 + 1)$
        by \cref{real_mul_assoc,real_mul_comm,real_mul_eq_subst}.
    $(n + 1) \rmul u = (1 / ((n + 1) + 1)) \rmul (1 + 1)$ by \cref{real_mul_eq_subst}.
    Hence $(n + 1) \rmul u = e$ by \cref{real_mul_comm,real_div}.

    $(((n + 1) + 1) \rmul c) = (((n + 1) + 1) \rmul \inv{((n + 1) + 1)}) \rmul \inv{(n + 1)}$
        by \cref{real_mul_assoc,real_mul_comm}.
    Hence $((n + 1) + 1) \rmul c = 1 \rmul \inv{(n + 1)}$ by \cref{real_mul_inverse}.
    We have $((n + 1) + 1) \rmul c = \inv{(n + 1)}$ by \cref{real_mul_ident}.
    Therefore $((n + 1) + 1) \rmul c = 1 / (n + 1)$ by \cref{real_div,real_mul_ident}.
    $(((n + 1) + 1) \rmul u) = (((n + 1) + 1) \rmul c) \rmul (1 + 1)$
        by \cref{real_mul_assoc,real_mul_comm,real_mul_eq_subst}.
    $((n + 1) + 1) \rmul u = (1 / (n + 1)) \rmul (1 + 1)$ by \cref{real_mul_eq_subst}.
    Hence $((n + 1) + 1) \rmul u = r$ by \cref{real_mul_comm,real_div}.

    $(((n + 1) + 1) \rmul u) = ((n + 1) \rmul u) + (1 \rmul u)$ by \cref{real_right_distrib}.
    Hence $r = e + u$ by \cref{real_mul_eq_subst,real_mul_ident}.
    We have $d = u$
        by \cref{real_sub,real_add_eq_subst,real_add_assoc,real_add_comm,real_add_inverse,real_add_ident}.

    We have $1 + 1 < n + 1$ or $1 + 1 = n + 1$
        by \cref{real_one_leq_naturalsplus,real_leq_add,real_leq_split}.
    Hence $(1 + 1) \rmul u = u + u$ by \cref{real_right_distrib,real_mul_ident,real_add_eq_subst}.
    Therefore $d + d \leq e$
        by \cref{real_lt_mul_pos,real_mul_eq_subst,real_lt_subst_left,real_lt_subst_right,real_add_eq_subst,real_lt_imp_leq,subseteq}.
\end{proof}

% from §35 helper lemma 5: one Tietze approximation step
\begin{lemma}\label{tietze_interval_approximation_step}
    Assume $X$ is normal with respect to $T$.
    Assume $A$ is closed in $X$ with respect to $T$.
    Assume $n\in\naturalsPlus$.
    Let $r = (1 + 1) / (n + 1)$.
    Let $e = (1 + 1) / ((n + 1) + 1)$.
    Let $d = r - e$.
    Let $I = \intervalclosed{\neg{r}}{r}$.
    Let $J = \intervalclosed{\neg{d}}{d}$.
    Let $K = \intervalclosed{\neg{e}}{e}$.
    Let $T_A = \subspaceTopology{T}{A}$.
    Let $T_I = \subspaceTopology{\standardTopologyR}{I}$.
    Let $T_J = \subspaceTopology{\standardTopologyR}{J}$.
    Assume $f$ is continuous from $A$ to $I$ with respect to $T_A$ and $T_I$.
    Then there exists $g$ such that
        $g$ is continuous from $X$ to $J$ with respect to $T$ and $T_J$ and
        for all $a\in A$ we have $f(a) - g(a)\in K$.
\end{lemma}
\begin{proof}
    $T$ is a topology on $X$ by \cref{normal_space}.
    $n + 1\in\naturalsPlus$ and $((n + 1) + 1)\in\naturalsPlus$
        by \cref{real_naturals_closed_succ}.
    $0\in\reals$ and $1\in\reals$ and $n\in\reals$ and $(n + 1)\in\reals$ and $((n + 1) + 1)\in\reals$
        by \cref{real_add_ident,real_one_in_naturals,real_naturals_subset_reals,subseteq,real_naturals_closed_succ}.
    $0 < n$ and $0 < n + 1$ and $0 < ((n + 1) + 1)$
        by \cref{naturalsplus_positive_real,real_naturals_closed_succ,subseteq}.
    Hence $n + 1 \neq 0$ and $((n + 1) + 1) \neq 0$ by \cref{real_pos_nonzero}.
    $\inv{(n + 1)}\in\reals$ and $\inv{((n + 1) + 1)}\in\reals$
        by \cref{real_mul_inverse}.
    Hence $0 < \inv{(n + 1)}$ and $0 < \inv{((n + 1) + 1)}$ by \cref{real_inv_pos}.
    $1 + 1\in\reals$ and $0 < 1 + 1$
        by \cref{real_add_closed,real_zero_lt_one,real_pos_add,real_mul_ident}.
    $1 / (n + 1)\in\reals$ and $1 / ((n + 1) + 1)\in\reals$
        by \cref{real_div,real_mul_closed}.
    Thus $r\in\reals$ and $e\in\reals$ by \cref{real_div,real_mul_closed}.

    We have $1 / ((n + 1) + 1) < 1 / (n + 1)$ by \cref{naturalsplus_reciprocal_step}.
    $(1 / ((n + 1) + 1)) \rmul (1 + 1) < (1 / (n + 1)) \rmul (1 + 1)$
        by \cref{real_lt_mul_pos}.
    $(1 / ((n + 1) + 1)) \rmul (1 + 1) = (1 \rmul \inv{((n + 1) + 1)}) \rmul (1 + 1)$
        by \cref{real_div}.
    Thus $(1 / ((n + 1) + 1)) \rmul (1 + 1) = (1 + 1) \rmul \inv{((n + 1) + 1)}$
        by \cref{real_mul_assoc,real_mul_comm,real_mul_ident}.
    Therefore $(1 / ((n + 1) + 1)) \rmul (1 + 1) = e$ by \cref{real_div}.
    $(1 / (n + 1)) \rmul (1 + 1) = (1 \rmul \inv{(n + 1)}) \rmul (1 + 1)$ by \cref{real_div}.
    Hence $(1 / (n + 1)) \rmul (1 + 1) = (1 + 1) \rmul \inv{(n + 1)}$
        by \cref{real_mul_assoc,real_mul_comm,real_mul_ident}.
    Therefore $(1 / (n + 1)) \rmul (1 + 1) = r$ by \cref{real_div}.
    Hence $e < r$ by \cref{real_lt_subst_left,real_lt_subst_right}.
    We have $d$ is positive by \cref{real_lt_imp_pos_diff}.
    $\neg{e}\in\reals$ by \cref{real_add_inverse}.
    $d\in\reals$ by \cref{real_sub,real_add_closed}.
    We have $d \neq 0$ by \cref{real_pos_nonzero}.
    We have $e$ is positive by \cref{real_div,real_mul_zero,real_lt_mul_pos,real_lt_subst_left}.
    Therefore $0 < r$ by \cref{real_lt_trans}.
    We have $d < r$ by \cref{real_sub_pos_lt}.
    We have $d \leq r$ by \cref{real_lt_imp_leq}.
    Hence $\neg{r} < \neg{d}$ by \cref{real_lt_neg}.
    $\neg{r} \leq \neg{d}$ by \cref{real_lt_imp_leq}.
    We have $\neg{d}\in\reals$ and $\neg{r}\in\reals$ by \cref{real_add_inverse}.
    We have $\neg{d} < 0$ by \cref{real_lt_neg,real_neg_zero,real_lt_subst_right}.
    Thus $\neg{d} < r$ by \cref{real_lt_trans}.
    We have $\neg{d} \leq r$ by \cref{real_lt_imp_leq}.
    We have $\neg{r} < 0$ by \cref{real_lt_neg,real_neg_zero,real_lt_subst_right}.
    Thus $\neg{r} < d$ by \cref{real_lt_trans}.
    We have $\neg{r} \leq d$ by \cref{real_lt_imp_leq}.
    Hence $I\subseteq\reals$ and $J\subseteq\reals$ and $K\subseteq\reals$ by \cref{intervalclosed,subseteq}.
    Hence $I$ is a subspace of $\reals$ with respect to $\standardTopologyR$ and
        $J$ is a subspace of $\reals$ with respect to $\standardTopologyR$
        by \cref{subspace_of,standard_basis_is_basis,basis_topology_is_topology,standard_topology_reals}.

    Let $I_1 = \intervalclosed{\neg{r}}{\neg{d}}$.
    Let $I_3 = \intervalclosed{d}{r}$.
    $I_1$ is closed in $\reals$ with respect to $\standardTopologyR$ and
        $I_3$ is closed in $\reals$ with respect to $\standardTopologyR$
        by \cref{real_intervalclosed_closed_standard}.

    $I_1\subseteq I$ by \cref{intervalclosed,real_leq_trans,subseteq}.

    $I_3\subseteq I$ by \cref{intervalclosed,real_leq_trans,subseteq}.

    $I_1$ is closed in $I$ with respect to $T_I$ and
        $I_3$ is closed in $I$ with respect to $T_I$
        by \cref{inter_absorb_supseteq_right,subseteq,closed_in_subspace_iff,subspace_of}.

    Let $B = \preimg{f}{I_1}$.
    Let $C = \preimg{f}{I_3}$.
    Let $D_1 = I\setminus I_1$.
    Let $D_3 = I\setminus I_3$.
    $D_1$ is open in $I$ with respect to $T_I$ and $D_3$ is open in $I$ with respect to $T_I$
        by \cref{closed_in}.
    Hence $\preimg{f}{D_1}$ is open in $A$ with respect to $T_A$ and
        $\preimg{f}{D_3}$ is open in $A$ with respect to $T_A$ by \cref{continuous}.
    $\preimg{f}{D_1} = A\setminus B$ and $\preimg{f}{D_3} = A\setminus C$ by
        \cref{continuous,closed_in,setminus_subseteq,preimg_setminus_function}.
    Hence $B\subseteq A$ and $C\subseteq A$ by \cref{continuous,preimg_subseteq_dom,funs}.
    Thus $B = A\setminus \preimg{f}{D_1}$ and $C = A\setminus \preimg{f}{D_3}$
        by \cref{subseteq,setminus_setminus,inter_absorb_supseteq_left}.
    Thus $B$ is closed in $A$ with respect to $T_A$ and $C$ is closed in $A$ with respect to $T_A$
        by \cref{closed_in}.
    Therefore $B$ is closed in $X$ with respect to $T$ and $C$ is closed in $X$ with respect to $T$
        by \cref{closed_in_closed_subspace,normal_space,closed_in}.

    Show that $B$ is disjoint from $C$.
    \begin{subproof}
        Suppose not.
        Take $a$ such that $a\in B$ and $a\in C$ by \cref{disjoint}.
        Hence $a\in A$ by \cref{subseteq}.
        Thus $f(a)\in I_1$ and $f(a)\in I_3$
            by \cref{continuous,funs_is_function,funs,preimg_iff,function_apply_iff}.
        Hence $d \leq f(a)$ and $f(a) \leq \neg{d}$ by \cref{intervalclosed,real_leq_trans}.
        Hence $f(a)\in\reals$ by \cref{intervalclosed}.
        We have $d \leq \neg{d}$ by \cref{real_leq_trans}.
        Thus $d < \neg{d}$ or $d = \neg{d}$ by \cref{real_leq_split}.
        Hence $d < 0$ by \cref{real_lt_trans,real_lt_subst_right}.
        Contradiction by \cref{real_lt_trans,real_lt_irreflexive}.
    \end{subproof}

    Hence $\neg{d} < d$ by \cref{real_lt_neg,real_neg_zero,real_lt_subst_left,real_lt_trans}.
    Let $I_0 = \intervalclosed{0}{1}$.
    Let $T_0 = \subspaceTopology{\standardTopologyR}{I_0}$.
    Take $u$ such that
        $u$ is continuous from $X$ to $I_0$ with respect to $T$ and $T_0$ and
        $\img{u}{B}\subseteq \{0\}$ and
        $\img{u}{C}\subseteq \{1\}$
        by \cref{urysohn_unit_interval}.
    Let $c = d + d$.
    Let $j = \identity{I_0}$.
    Let $u_1 = j\circ u$.
    $u_1$ is continuous from $X$ to $\reals$ with respect to $T$ and $\standardTopologyR$
        by \cref{standard_basis_is_basis,basis_topology_is_topology,standard_topology_reals,continuous_interval_expand_range}.
    $u\in\funs{X}{I_0}$ and $u$ is a function by \cref{continuous,funs_is_function}.
    Hence $u$ is right-unique and $u$ is a relation and $\dom{u} = X$
        by \cref{function_rightunique,funs_is_relation,funs_elim}.
    Therefore $u\in\rels{X}{I_0}$ by \cref{funs_subseteq_rels,subseteq}.
    We have $\ran{u}\subseteq I_0$ by \cref{rels_ran_subseteq}.
    We have $\dom{j} = I_0$ by \cref{id_dom}.
    Hence $\ran{u}\subseteq\dom{j}$.
    $j$ is right-unique and $j$ is a relation by \cref{id_rightunique,id_is_relation}.
    For all $x\in X$ we have $u_1(x) = u(x)$ by \cref{circ_apply,funs_elim,id_apply}.
    Let $k = \{(x,\neg{d}) \mid x\in X\}$.
    Let $h = \{(x,c) \mid x\in X\}$.
    We have $c\in\reals$ and $c$ is positive by \cref{real_add_closed,real_pos_add}.
    $\standardTopologyR$ is a topology on $\reals$
        by \cref{standard_basis_is_basis,basis_topology_is_topology,standard_topology_reals}.
    $k$ is a function from $X$ to $\reals$ and $h$ is a function from $X$ to $\reals$
        by \cref{constant_map_function}.
    Hence for all $x\in X$ we have $k(x) = \neg{d}$ and $h(x) = c$
        by \cref{constant_map_function,funs_is_function,function_apply_intro}.
    $k$ is continuous from $X$ to $\reals$ with respect to $T$ and $\standardTopologyR$ and
        $h$ is continuous from $X$ to $\reals$ with respect to $T$ and $\standardTopologyR$
        by \cref{constant_continuous}.
    Let $m = \{(x, h(x) \rmul u_1(x)) \mid x\in X\}$.
    Let $g = \{(x, k(x) + m(x)) \mid x\in X\}$.
    $m$ is continuous from $X$ to $\reals$ with respect to $T$ and $\standardTopologyR$
        by \cref{real_function_multiplication_continuous}.
    $g$ is continuous from $X$ to $\reals$ with respect to $T$ and $\standardTopologyR$
        by \cref{real_function_addition_continuous}.
    $m\in\funs{X}{\reals}$ and $m$ is a function and
        $g\in\funs{X}{\reals}$ and $g$ is a function
        by \cref{continuous,funs_is_function}.

    Show that for all $x\in X$ we have $g(x)\in J$.
    \begin{subproof}
        Fix $x\in X$.
        We have $u_1(x)\in I_0$ by \cref{circ_apply,funs_elim,id_apply}.
        Hence $u_1(x)\in\reals$ by \cref{intervalclosed}.
        We have $0 \leq u_1(x)$ and $u_1(x) \leq 1$ by \cref{intervalclosed}.
        Thus $(x, h(x) \rmul u_1(x))\in m$.
        Hence $m(x) = c \rmul u_1(x)$ by \cref{function_apply_iff}.
        $u_1(x) = 0$ or $0 < u_1(x)$ by \cref{real_lt_trichotomy}.
        Hence $0 \leq c \rmul u_1(x)$
            by \cref{real_mul_zero,real_mul_comm,real_lt_imp_pos_diff,real_mul_pos_pos_gt,real_lt_imp_leq}.
        $0 \leq m(x)$.
        $u_1(x) = 1$ or $u_1(x) < 1$ by \cref{real_lt_trichotomy}.
        Hence $c \rmul u_1(x) \leq c$
            by \cref{real_mul_comm,real_mul_ident,real_lt_imp_pos_diff,real_naturals_subset_reals,real_one_in_naturals,subseteq,real_lt_mul_pos}.
        $m(x) \leq c$.
        We have $m(x)\in\reals$ by \cref{real_mul_closed}.
        Thus $(x, k(x) + m(x))\in g$.
        Hence $g(x) = m(x) + \neg{d}$ by \cref{function_apply_iff,real_add_eq_subst,real_add_comm}.
        We have $m(x) + \neg{d}\in\reals$ and $0 + \neg{d} \leq m(x) + \neg{d}$ by
            \cref{real_add_closed,real_leq_add}.
        Hence $\neg{d} \leq m(x) + \neg{d}$ by \cref{real_add_ident}.
        Hence $c + \neg{d} = d$
            by \cref{real_add_eq_subst,real_add_assoc,real_add_inverse,real_add_comm,real_add_ident}.
        Hence $m(x) + \neg{d} \leq c + \neg{d}$ by \cref{real_leq_add}.
        Therefore $m(x) + \neg{d} \leq d$ by \cref{real_leq_trans}.
        Hence $m(x) + \neg{d}\in J$ by \cref{intervalclosed}.
        $g(x)\in J$ by \cref{function_apply_intro}.
    \end{subproof}
    Hence $\img{g}{X}\subseteq J$
        by \cref{subseteq,img_iff,continuous,funs_is_function,function_apply_iff}.
    Show that $g$ is continuous from $X$ to $J$ with respect to $T$ and $T_J$.
    \begin{subproof}
        Follows by \cref{continuous_restrict_range}.
    \end{subproof}

    Show that for all $x\in B$ we have $g(x) = \neg{d}$.
    \begin{subproof}
        Fix $x\in B$.
        We have $x\in X$ by \cref{closed_in,subseteq}.
        Hence $u(x)\in\img{u}{B}$ by \cref{function_apply_elim,img_iff,continuous,funs_elim}.
        $u(x)\in\{0\}$ by \cref{subseteq}.
        Thus $u(x) = 0$ by \cref{singleton_elim}.
        Hence $u_1(x) = u(x)$ by \cref{circ_apply,funs_elim,id_apply}.
        $u_1(x) = 0$.
        Therefore $(x, h(x) \rmul u_1(x))\in m$.
        Hence $m(x) = c \rmul 0$ by \cref{function_apply_iff}.
        $m(x) = 0$ by \cref{real_mul_comm,real_mul_zero}.
        Therefore $(x, k(x) + m(x))\in g$.
        Hence $g(x) = \neg{d} + 0$ by \cref{function_apply_iff}.
        $g(x) = \neg{d}$ by \cref{real_add_comm,real_add_ident}.
    \end{subproof}
    Show that for all $x\in C$ we have $g(x) = d$.
    \begin{subproof}
        Fix $x\in C$.
        We have $x\in X$ by \cref{closed_in,subseteq}.
        Hence $u(x)\in\img{u}{C}$ by \cref{function_apply_elim,img_iff,continuous,funs_elim}.
        $u(x)\in\{1\}$ by \cref{subseteq}.
        Thus $u(x) = 1$ by \cref{singleton_elim}.
        Hence $u_1(x) = u(x)$ by \cref{circ_apply,funs_elim,id_apply}.
        $u_1(x) = 1$.
        Therefore $(x, h(x) \rmul u_1(x))\in m$.
        Hence $m(x) = c \rmul 1$ by \cref{function_apply_iff}.
        $m(x) = c$ by \cref{real_mul_comm,real_mul_ident}.
        Therefore $(x, k(x) + m(x))\in g$.
        Hence $g(x) = \neg{d} + c$ by \cref{function_apply_iff}.
        $g(x) = c + \neg{d}$ by \cref{real_add_comm}.
        Hence $c + \neg{d} = d$
            by \cref{real_add_eq_subst,real_add_assoc,real_add_inverse,real_add_comm,real_add_ident}.
        Hence $g(x) = d$.
    \end{subproof}
    We have $g\in\funs{X}{J}$ and $g$ is a function and $\dom{g} = X$ by \cref{continuous,funs_is_function,funs}.
    We have $f\in\funs{A}{I}$ and $f$ is a function and $\dom{f} = A$ by \cref{continuous,funs_is_function,funs}.

    Show that for all $a\in A$ we have $f(a) - g(a)\in K$.
    \begin{subproof}
        Fix $a\in A$.
        Hence $a\in X$ by \cref{subseteq,closed_in}.
        $f(a)\in I$ and $g(a)\in J$ by \cref{funs_type_apply}.
        $f(a)\in\reals$ and $g(a)\in\reals$ by \cref{intervalclosed}.
        Thus $f(a) - g(a)\in\reals$ by \cref{real_sub,real_add_inverse,real_add_closed}.
        \begin{byCase}
        \caseOf{$a\in B$.}
            We have $g(a) = \neg{d}$.
            $f(a)\in I_1$ by \cref{preimg_iff,function_apply_iff}.
            Thus $\neg{r} \leq f(a)$ and $f(a) \leq \neg{d}$ by \cref{intervalclosed}.
        Hence $\neg{r} + d \leq f(a) + d$ and $f(a) + d \leq \neg{d} + d$
            by \cref{real_leq_add}.
        $f(a) - g(a) = f(a) + d$ by \cref{real_sub,real_neg_neg,real_add_eq_subst}.
        $\neg{r} + d = d + \neg{r}$ by \cref{real_add_comm}.
        We have $d + \neg{r} = d - r$ by \cref{real_sub}.
        Hence $d - r = \neg{e}$ by \cref{real_sub_sub_cancel,real_sub_swap_neg,real_neg_neg,real_add_eq_subst}.
        Hence $\neg{e} \leq f(a) - g(a)$ by \cref{real_leq_trans,real_leq_add,real_add_inverse,real_add_comm,real_add_ident,real_add_eq_subst}.
        $\neg{d} + d = 0$ by \cref{real_add_inverse,real_add_comm}.
        Hence $f(a) - g(a) \leq 0$ by \cref{real_leq_add,real_add_eq_subst,real_add_ident}.
            We have $0 \leq e$ by \cref{real_lt_imp_leq,real_lt_subst_right}.
            Hence $f(a) - g(a) \leq e$ by \cref{real_leq_trans}.
            Hence $f(a) - g(a)\in K$ by \cref{intervalclosed}.
        \caseOf{$a\in C$.}
            We have $g(a) = d$.
            $f(a)\in I_3$ by \cref{preimg_iff,function_apply_iff}.
            Thus $d \leq f(a)$ and $f(a) \leq r$ by \cref{intervalclosed}.
        Hence $d + \neg{d} \leq f(a) + \neg{d}$ and $f(a) + \neg{d} \leq r + \neg{d}$
            by \cref{real_leq_add}.
        $f(a) - g(a) = f(a) + \neg{d}$ by \cref{real_sub,real_add_eq_subst}.
        $d + \neg{d} = 0$ by \cref{real_add_inverse,real_add_comm}.
        Hence $0 \leq f(a) - g(a)$ by \cref{real_add_eq_subst,real_add_ident}.
        We have $\neg{e} \leq 0$ by \cref{real_lt_imp_leq,real_lt_neg,real_neg_zero,real_lt_subst_right}.
        Hence $\neg{e} \leq f(a) - g(a)$ by \cref{real_leq_trans}.
        We have $r + \neg{d} = r - d$ by \cref{real_sub}.
        Hence $f(a) - g(a) \leq e$ by \cref{real_sub_sub_cancel,real_leq_trans,real_add_eq_subst}.
        Hence $f(a) - g(a)\in K$ by \cref{intervalclosed}.
        \caseOf{$a\notin B$ and $a\notin C$.}
            Hence $f(a) < d$ by \cref{intervalclosed,preimg_iff,function_apply_elim,real_lt_trichotomy,real_lt_imp_leq}.
            We have $\neg{d} < f(a)$ by \cref{intervalclosed,preimg_iff,function_apply_elim,real_lt_trichotomy,real_lt_imp_leq}.
            Hence $\neg{d} \leq f(a)$ and $f(a) \leq d$ by \cref{real_lt_imp_leq}.
            $\neg{d} \leq g(a)$ and $g(a) \leq d$ by \cref{intervalclosed}.
            We have $\neg{(g(a))}\in\reals$ by \cref{real_add_inverse}.
            Hence $\neg{(g(a))} \leq d$ by \cref{real_leq_neg,real_neg_neg,real_lt_imp_leq}.
            Thus $\neg{d} \leq \neg{(g(a))}$ by \cref{real_leq_neg,real_neg_neg}.
            Hence $f(a) + \neg{(g(a))}\in\reals$ and $d + \neg{(g(a))}\in\reals$ and
                $\neg{d} + \neg{(g(a))}\in\reals$ and $d + d\in\reals$ and $\neg{d} + \neg{d}\in\reals$
                by \cref{real_add_closed}.
            Hence $f(a) + \neg{(g(a))} \leq d + \neg{(g(a))}$ and
                $\neg{d} + \neg{(g(a))} \leq f(a) + \neg{(g(a))}$ by \cref{real_leq_add}.
            $d + \neg{(g(a))} = \neg{(g(a))} + d$ and
                $\neg{(g(a))} + \neg{d} = \neg{d} + \neg{(g(a))}$ by \cref{real_add_comm}.
            We have $\neg{(g(a))} + d \leq d + d$ by \cref{real_leq_add}.
            Hence $\neg{d} + \neg{d} \leq \neg{(g(a))} + \neg{d}$ by \cref{real_leq_add}.
            Therefore $d + \neg{(g(a))} \leq d + d$ and
                $\neg{d} + \neg{d} \leq \neg{d} + \neg{(g(a))}$ by \cref{real_leq_trans}.
            Hence $f(a) + \neg{(g(a))} \leq d + d$ by \cref{real_leq_trans}.
            Hence $\neg{d} + \neg{d} \leq f(a) + \neg{(g(a))}$ by \cref{real_leq_trans}.
            We have $d + d \leq e$ by \cref{tietze_interval_double_step_width_leq_error}.
            Hence $f(a) + \neg{(g(a))} \leq e$ by \cref{real_leq_trans}.
            $\neg{(d + d)} = \neg{d} + \neg{d}$ by \cref{real_neg_addition}.
            We have $\neg{e} \leq \neg{(d + d)}$ by \cref{real_leq_neg}.
            Hence $\neg{e} \leq \neg{d} + \neg{d}$ by \cref{real_add_eq_subst}.
            Thus $\neg{e} \leq f(a) + \neg{(g(a))}$ by \cref{real_leq_trans}.
            Hence $f(a) - g(a) = f(a) + \neg{(g(a))}$ by \cref{real_sub}.
            Therefore $f(a) - g(a)\in K$ by \cref{intervalclosed,real_leq_trans,real_add_eq_subst}.
        \end{byCase}
    \end{subproof}

    Take $h_0$ such that $h_0 = g$.
    Follows by \cref{pair_eq_iff}.
\end{proof}

% from §35 helper definition 1: unit-interval approximation state
\begin{definition}\label{tietze_unit_state}
    $s$ is a Tietze unit state on $n$ for $f$ over $A$ in $X$ with respect to $T$ iff
        $n\in\naturalsPlus$ and
        $s$ is continuous from $X$ to $\reals$ with respect to $T$ and $\standardTopologyR$ and
        for all $a\in A$ we have
            $f(a) - s(a)\in \intervalclosed{\neg{((1 + 1) / (n + 1))}}{(1 + 1) / (n + 1)}$.
\end{definition}

% from §35 helper lemma 7: a unit-interval approximation state has a successor
\begin{lemma}\label{tietze_unit_state_successor}
    Assume $X$ is normal with respect to $T$.
    Assume $A$ is closed in $X$ with respect to $T$.
    Let $T_A = \subspaceTopology{T}{A}$.
    Assume $f$ is continuous from $A$ to $\reals$ with respect to $T_A$ and $\standardTopologyR$.
    Assume $s$ is a Tietze unit state on $n$ for $f$ over $A$ in $X$ with respect to $T$.
    Let $r = (1 + 1) / (n + 1)$.
    Let $e = (1 + 1) / ((n + 1) + 1)$.
    Let $d = r - e$.
    Let $I = \intervalclosed{\neg{r}}{r}$.
    Let $J = \intervalclosed{\neg{d}}{d}$.
    Then there exists $t$ such that
        $t$ is a Tietze unit state on $n + 1$ for $f$ over $A$ in $X$ with respect to $T$ and
        for all $x\in X$ we have $t(x) - s(x)\in J$.
\end{lemma}
\begin{proof}
    $T$ is a topology on $X$ by \cref{normal_space}.
    Hence $A\subseteq X$ by \cref{closed_in}.
    Thus $A$ is a subspace of $X$ with respect to $T$ by \cref{subspace_of}.
    Therefore $n\in\naturalsPlus$ and
        $s$ is continuous from $X$ to $\reals$ with respect to $T$ and $\standardTopologyR$ and
        for all $a\in A$ we have $f(a) - s(a)\in I$ by \cref{tietze_unit_state}.
    We have $s\in\funs{X}{\reals}$ and $s$ is a function and $\dom{s} = X$
        by \cref{continuous,funs_is_function,funs}.
    Let $s_A = \restrl{s}{A}$.
    Hence $s_A$ is continuous from $A$ to $\reals$ with respect to $T_A$ and $\standardTopologyR$
        by \cref{continuous_restrict_domain}.
    We have $s_A\in\funs{A}{\reals}$ and $s_A$ is a function and $\dom{s_A} = A$
        by \cref{continuous,funs_is_function,funs}.
    Let $h = \{(a, f(a) - s_A(a)) \mid a\in A\}$.
    Hence $h$ is continuous from $A$ to $\reals$ with respect to $T_A$ and $\standardTopologyR$
        by \cref{real_function_subtraction_continuous}.
    We have $h\in\funs{A}{\reals}$ and $h$ is a function and $\dom{h} = A$ by \cref{continuous,funs_is_function,funs}.
    For all $a\in A$ we have $s_A(a) = s(a)$ and $h(a) = f(a) - s(a)$
        by \cref{continuous,funs,subseteq,function_apply_elim,restrl_iff,function_apply_iff,function_apply_intro}.
    Hence $h$ is continuous from $A$ to $I$ with respect to $T_A$ and $\subspaceTopology{\standardTopologyR}{I}$
        by \cref{tietze_unit_state,img_iff,funs,dom,function_apply_iff,real_add_eq_subst,subseteq,intervalclosed,continuous_subspace_range,subspace_topology_is_topology,standard_basis_is_basis,basis_topology_is_topology,standard_topology_reals}.
    Let $T_I = \subspaceTopology{\standardTopologyR}{I}$.
    Let $T_J = \subspaceTopology{\standardTopologyR}{J}$.
    Take $g$ such that
        $g$ is continuous from $X$ to $J$ with respect to $T$ and $T_J$ and
        for all $a\in A$ we have $h(a) - g(a)\in \intervalclosed{\neg{e}}{e}$
        by \cref{tietze_interval_approximation_step}.
    Let $j = \identity{J}$.
    Let $g_1 = j\circ g$.
    Hence $g_1$ is continuous from $X$ to $\reals$ with respect to $T$ and $\standardTopologyR$
        by \cref{continuous_interval_expand_range}.
    We have $g\in\funs{X}{J}$ and $g$ is a function and $\dom{g} = X$ and $\ran{g}\subseteq J$
        by \cref{continuous,funs_is_function,funs,funs_ran}.
    We have $\dom{j} = J$ by \cref{id_dom}.
    We have $j$ is right-unique and $j$ is a relation by \cref{id_rightunique,id_is_relation}.
    Let $t = \{(x, s(x) + g_1(x)) \mid x\in X\}$.
    Hence $t$ is continuous from $X$ to $\reals$ with respect to $T$ and $\standardTopologyR$
        by \cref{real_function_addition_continuous}.
    We have $t\in\funs{X}{\reals}$ and $t$ is a function and $\dom{t} = X$ by \cref{continuous,funs_is_function,funs}.
    For all $a\in A$ we have $s_A(a) = s(a)$ and $h(a) = f(a) - s(a)$ and
        $g_1(a) = g(a)$ and $t(a) = s(a) + g(a)$ and $f(a) - t(a) = h(a) - g(a)$
        by \cref{continuous,funs,subseteq,function_apply_elim,restrl_iff,function_apply_iff,function_apply_intro,circ_apply,funs_elim,id_apply,real_add_eq_subst,real_sub,real_add_inverse,real_neg_addition,real_add_assoc,funs_type_apply,intervalclosed}.
    Show that $t$ is a Tietze unit state on $n + 1$ for $f$ over $A$ in $X$ with respect to $T$.
    \begin{subproof}
        Follows by \cref{real_add_eq_subst,tietze_unit_state,real_naturals_closed_succ}.
    \end{subproof}
    Show that for all $x\in X$ we have $t(x) - s(x)\in J$.
    \begin{subproof}
        Follows by \cref{continuous,funs_type_apply,intervalclosed,real_sub,real_add_inverse,real_add_closed,function_apply_intro,circ_apply,funs_elim,id_apply,real_add_eq_subst,real_add_sub_cancel_right,real_add_cancel,real_add_comm}.
    \end{subproof}
\end{proof}

% from §35 helper lemma 9: zero map is the initial unit state
\begin{lemma}\label{tietze_unit_state_zero}
    Assume $X$ is normal with respect to $T$.
    Assume $A$ is closed in $X$ with respect to $T$.
    Let $I = \intervalclosed{\neg{1}}{1}$.
    Let $T_A = \subspaceTopology{T}{A}$.
    Assume $f$ is continuous from $A$ to $\reals$ with respect to $T_A$ and $\standardTopologyR$.
    Assume for all $a\in A$ we have $f(a)\in I$.
    Let $z = \{(x,0) \mid x\in X\}$.
    Then $z$ is a Tietze unit state on $1$ for $f$ over $A$ in $X$ with respect to $T$.
\end{lemma}
\begin{proof}
    $0\in\reals$ by \cref{real_add_ident}.
    $z$ is a function from $X$ to $\reals$ by \cref{constant_map_function}.
    For all $x\in X$ we have $z(x) = 0$ by \cref{constant_map_function,funs_is_function,function_apply_intro}.
    $z$ is continuous from $X$ to $\reals$ with respect to $T$ and $\standardTopologyR$
        by \cref{normal_space,standard_basis_is_basis,basis_topology_is_topology,standard_topology_reals,constant_continuous}.
    $((1 + 1) / (1 + 1)) = 1$ by
        \cref{real_div,real_mul_assoc,real_mul_comm,real_mul_inverse,real_mul_ident,real_add_closed,real_zero_lt_one,real_pos_add,real_pos_nonzero,real_naturals_subset_reals,real_one_in_naturals,subseteq}.
    For all $a\in A$ we have $f(a) - z(a)\in \intervalclosed{\neg{((1 + 1) / (1 + 1))}}{(1 + 1) / (1 + 1)}$
        by \cref{subseteq,closed_in,continuous,funs_type_apply,intervalclosed,real_neg_zero,real_sub,real_add_eq_subst,real_add_comm,real_add_ident}.
    $z$ is a Tietze unit state on $1$ for $f$ over $A$ in $X$ with respect to $T$
        by \cref{tietze_unit_state,real_one_in_naturals}.
\end{proof}

% from §35 helper lemma 10: recursive sequence of Tietze unit states
\begin{lemma}\label{tietze_unit_state_iteration}
    Assume $X$ is normal with respect to $T$.
    Assume $A$ is closed in $X$ with respect to $T$.
    Let $I = \intervalclosed{\neg{1}}{1}$.
    Let $T_A = \subspaceTopology{T}{A}$.
    Assume $f$ is continuous from $A$ to $\reals$ with respect to $T_A$ and $\standardTopologyR$.
    Assume for all $a\in A$ we have $f(a)\in I$.
    Then there exists $u$ such that
        $u$ is a sequence in $\naturalsPlus\times\funs{X}{\reals}$ and
        for all $x\in X$ we have $\apply{\snd{u(1)}}{x} = 0$ and
        for all $n\in\naturalsPlus$ we have
            $\fst{u(n)} = n$ and
            $\snd{u(n)}$ is a Tietze unit state on $n$ for $f$ over $A$ in $X$ with respect to $T$ and
            for all $x\in X$ we have
                $\apply{\snd{u(n + 1)}}{x} - \apply{\snd{u(n)}}{x}\in
                    \intervalclosed{\neg{(((1 + 1) / (n + 1)) - ((1 + 1) / ((n + 1) + 1)))}}{((1 + 1) / (n + 1)) - ((1 + 1) / ((n + 1) + 1))}$.
\end{lemma}
\begin{proof}
    $T$ is a topology on $X$ by \cref{normal_space}.
    $f\in\funs{A}{\reals}$ and $\dom{f} = A$ by \cref{continuous,funs}.
    We have $I$ is a subspace of $\reals$ with respect to $\standardTopologyR$
        by \cref{continuous,intervalclosed,subseteq,subspace_of}.
    $\img{f}{A}\subseteq I$
        by \cref{img_iff,dom,continuous,funs_is_function,function_rightunique,funs_is_relation,function_apply_iff,real_add_eq_subst,subseteq}.
    Let $T_I = \subspaceTopology{\standardTopologyR}{I}$.
    $f$ is continuous from $A$ to $I$ with respect to $T_A$ and $T_I$
        by \cref{continuous_restrict_range}.

    Let $G = \{w\in\naturalsPlus\times\funs{X}{\reals} \mid
        \text{$\snd{w}$ is a Tietze unit state on $\fst{w}$ for $f$ over $A$ in $X$ with respect to $T$}\}$.

    Let $R = \{(z,w) \mid z\in G, w\in G \mid \text{$\fst{w} = \fst{z} + 1$ and
        for all $x\in X$ we have
            $\apply{\snd{w}}{x} - \apply{\snd{z}}{x}\in
                \intervalclosed{\neg{(((1 + 1) / (\fst{z} + 1)) - ((1 + 1) / ((\fst{z} + 1) + 1)))}}{((1 + 1) / (\fst{z} + 1)) - ((1 + 1) / ((\fst{z} + 1) + 1))}$}\}$.

    Show that for all $z\in G$ there exists $w$ such that $z\mathrel{R} w$.
    \begin{subproof}
        Fix $z\in G$.
        Take $n,s$ such that $z = (n,s)$ and $n\in\naturalsPlus$ and $s\in\funs{X}{\reals}$
            by \cref{subseteq,times_elem_is_tuple}.
        $\fst{z} = n$ and $\snd{z} = s$ by \cref{pair_eq_iff,fst_eq,snd_eq}.
        $s$ is a Tietze unit state on $n$ for $f$ over $A$ in $X$ with respect to $T$
            by \cref{subseteq,real_add_eq_subst}.
        Let $r = (1 + 1) / (n + 1)$.
        Let $e = (1 + 1) / ((n + 1) + 1)$.
        Let $d = r - e$.
        Let $J = \intervalclosed{\neg{d}}{d}$.
        Take $t$ such that
            $t$ is a Tietze unit state on $n + 1$ for $f$ over $A$ in $X$ with respect to $T$ and
            for all $x\in X$ we have $t(x) - s(x)\in J$
            by \cref{tietze_unit_state_successor}.
        Take $w$ such that $w = (n + 1,t)$.
        $(n + 1,t)\in\naturalsPlus\times\funs{X}{\reals}$
            by \cref{real_naturals_closed_succ,tietze_unit_state,continuous,funs,times_tuple_intro}.
        $\fst{w} = n + 1$ and $\snd{w} = t$ by \cref{pair_eq_iff,fst_eq,snd_eq}.
        $\snd{w}$ is a Tietze unit state on $\fst{w}$ for $f$ over $A$ in $X$ with respect to $T$
            by \cref{tietze_unit_state,real_add_eq_subst}.
        $w\in G$ by \cref{subseteq}.
        $z\mathrel{R} w$ by \cref{relation,pair_eq_iff,subseteq}.
    \end{subproof}

    We have $R$ is a relation by \cref{relation}.

    Take $c$ such that $c$ is a function on $G$ and for all $z\in G$ we have $z\mathrel{R}\apply{c}{z}$
        by \cref{choice_function}.

    Let $z_0 = \{(x,0) \mid x\in X\}$.
    Let $a_0 = (1,z_0)$.
    $z_0$ is a Tietze unit state on $1$ for $f$ over $A$ in $X$ with respect to $T$
        by \cref{tietze_unit_state_zero}.
    $(1,z_0)\in\naturalsPlus\times\funs{X}{\reals}$
        by \cref{tietze_unit_state,continuous,funs,times_tuple_intro,real_one_in_naturals}.
    $a_0\in G$ by \cref{subseteq,pair_eq_iff,fst_eq,snd_eq,tietze_unit_state,real_add_eq_subst}.

    $c\in\funs{G}{G}$ by \cref{choice_function,funs_intro,relation,pair_eq_iff,subseteq}.

    For all $z,w$ such that $z\mathrel{R} w$ we have $\fst{w} = \fst{z} + 1$
        by \cref{relation,pair_eq_iff,real_add_eq_subst}.

    Take $u$ such that
        $u$ is a sequence in $G$ and
        $u(1) = a_0$ and
        for all $n\in\naturalsPlus$ we have $u(n + 1) = \apply{c}{u(n)}$
        by \cref{iteration_sequence_naturalsplus}.

    Let $A_0 = \{n\in\naturalsPlus \mid \fst{u(n)} = n\}$.
    $A_0\subseteq\naturalsPlus$ by \cref{subseteq}.
    $1\in A_0$ by \cref{iteration_sequence_naturalsplus,real_one_in_naturals,pair_eq_iff,fst_eq}.
    For all $k\in A_0$ we have $k + 1\in A_0$
        by \cref{subseteq,sequence_in,funs_type_apply,choice_function,real_add_eq_subst,relation,pair_eq_iff,real_naturals_closed_succ}.
    For all $n\in\naturalsPlus$ we have $n\in A_0$ by \cref{real_naturals_induction}.
    For all $n\in\naturalsPlus$ we have $\fst{u(n)} = n$.
    Show that $u$ is a sequence in $\naturalsPlus\times\funs{X}{\reals}$.
    \begin{subproof}
        Follows by \cref{sequence_in,funs_weaken_codom,subseteq}.
    \end{subproof}
    Show that for all $x\in X$ we have $\apply{\snd{u(1)}}{x} = 0$.
    \begin{subproof}
        Follows by \cref{iteration_sequence_naturalsplus,pair_eq_iff,snd_eq,subseteq,continuous,funs,tietze_unit_state,real_add_eq_subst,funs_is_function,function_rightunique,funs_is_relation,function_apply_iff}.
    \end{subproof}

    Show that for all $n\in\naturalsPlus$ we have
        $\fst{u(n)} = n$ and
        $\snd{u(n)}$ is a Tietze unit state on $n$ for $f$ over $A$ in $X$ with respect to $T$ and
        for all $x\in X$ we have
            $\apply{\snd{u(n + 1)}}{x} - \apply{\snd{u(n)}}{x}\in
                \intervalclosed{\neg{(((1 + 1) / (n + 1)) - ((1 + 1) / ((n + 1) + 1)))}}{((1 + 1) / (n + 1)) - ((1 + 1) / ((n + 1) + 1))}$.
    \begin{subproof}
        Fix $n\in\naturalsPlus$.
        We have $\fst{u(n)} = n$.
        $u(n)\in G$ by \cref{sequence_in,funs_type_apply}.
        $\snd{u(n)}$ is a Tietze unit state on $n$ for $f$ over $A$ in $X$ with respect to $T$
            by \cref{subseteq,real_add_eq_subst}.
        $u(n + 1) = c(u(n))$ by \cref{subseteq}.
        $u(n)\mathrel{R} u(n + 1)$ by \cref{choice_function,sequence_in,funs_type_apply}.
        Take $z,w$ such that $z\in G$ and $w\in G$ and $(z,w) = (u(n),u(n + 1))$ and
            $\fst{w} = \fst{z} + 1$ and
            for all $y\in X$ we have
                $\apply{\snd{w}}{y} - \apply{\snd{z}}{y}\in
                    \intervalclosed{\neg{(((1 + 1) / (\fst{z} + 1)) - ((1 + 1) / ((\fst{z} + 1) + 1)))}}{((1 + 1) / (\fst{z} + 1)) - ((1 + 1) / ((\fst{z} + 1) + 1))}$
            by \cref{relation}.
        For all $x\in X$ we have
            $\apply{\snd{u(n + 1)}}{x} - \apply{\snd{u(n)}}{x}\in
                \intervalclosed{\neg{(((1 + 1) / (n + 1)) - ((1 + 1) / ((n + 1) + 1)))}}{((1 + 1) / (n + 1)) - ((1 + 1) / ((n + 1) + 1))}$
                by \cref{pair_eq_iff,subseteq,snd_eq,real_add_eq_subst}.
    \end{subproof}
    Follows by \cref{pair_eq_iff}.
\end{proof}

% from §35 helper definition 2: lower endpoint set of the Tietze iteration
\begin{definition}\label{tietze_lower_set}
    $\tietzeLowerSet{u}{x} =
        \{\apply{\snd{u(n)}}{x} - ((1 + 1) / (n + 1)) \mid n\in\naturalsPlus\}$.
\end{definition}

% from §35 helper definition 3: upper endpoint set of the Tietze iteration
\begin{definition}\label{tietze_upper_set}
    $\tietzeUpperSet{u}{x} =
        \{\apply{\snd{u(n)}}{x} + ((1 + 1) / (n + 1)) \mid n\in\naturalsPlus\}$.
\end{definition}

% from §35 helper lemma 11: approximation errors can be made arbitrarily small
\begin{lemma}\label{tietze_small_error}
    Assume $\epsilon\in\reals$ and $0 < \epsilon$.
    Then there exists $N\in\naturalsPlus$ such that $((1 + 1) / (N + 1)) < \epsilon$.
\end{lemma}
\begin{proof}
    We have $1\in\reals$ and $0 < 1$ by \cref{real_naturals_subset_reals,real_one_in_naturals,subseteq,real_zero_lt_one}.
    $1 + 1\in\reals$ and $0 < 1 + 1$ and $(1 + 1)\neq 0$
        by \cref{real_add_closed,real_pos_add,real_pos_nonzero}.
    Let $\delta = \epsilon / (1 + 1)$.
    $\delta\in\reals$ and $0 < \delta$ by \cref{real_div,real_mul_closed,real_mul_inverse,real_inv_pos,real_mul_pos_pos}.
    Take $N\in\naturalsPlus$ such that $0 < 1 / N$ and $1 / N < \delta$ by \cref{real_small_reciprocal}.
    $N + 1\in\naturalsPlus$ and $N + 1\in\reals$ and $0 < N + 1$ and $N + 1\neq 0$
        by \cref{real_naturals_closed_succ,real_naturals_subset_reals,subseteq,naturalsplus_positive_real,real_pos_nonzero}.
    $N < N + 1$ by \cref{real_naturals_lt_succ_local}.
    $1 / (N + 1) < 1 / N$ by \cref{naturalsplus_reciprocal_step}.
    $1 / (N + 1)\in\reals$ and $1 / N\in\reals$ and $\inv{(N + 1)}\in\reals$
        by \cref{naturalsplus_reciprocal_real,subseteq,real_mul_inverse}.
    $1 / (N + 1) < \delta$ by \cref{real_lt_trans}.
    $(\inv{(N + 1)} \rmul (1 + 1)) < (\delta \rmul (1 + 1))$
        by \cref{real_div,real_mul_ident,real_lt_subst_left,real_lt_mul_pos}.
    $(\inv{(N + 1)} \rmul (1 + 1)) = ((1 + 1) / (N + 1))$ by \cref{real_div,real_mul_comm}.
    $(\delta \rmul (1 + 1)) = \epsilon$ by
        \cref{real_div,real_mul_assoc,real_mul_comm,real_mul_inverse,real_mul_ident}.
    $((1 + 1) / (N + 1)) < \epsilon$ by \cref{real_lt_subst_left,real_lt_subst_right}.
\end{proof}

% from §35 helper lemma 12: doubled approximation errors can be made arbitrarily small
\begin{lemma}\label{tietze_small_double_error}
    Assume $\epsilon\in\reals$ and $0 < \epsilon$.
    Then there exists $N\in\naturalsPlus$ such that
        $((1 + 1) / (N + 1)) + ((1 + 1) / (N + 1)) < \epsilon$.
\end{lemma}
\begin{proof}
    Take $\delta$ such that $\delta\in\reals$ and $0 < \delta$ and $\delta + \delta = \epsilon$
        by \cref{real_half_of_positive}.
    Take $N\in\naturalsPlus$ such that $((1 + 1) / (N + 1)) < \delta$ by \cref{tietze_small_error}.
    $((1 + 1) / (N + 1)) + ((1 + 1) / (N + 1)) < \epsilon$
        by \cref{real_naturals_closed_succ,real_naturals_subset_reals,subseteq,naturalsplus_positive_real,real_pos_nonzero,real_one_in_naturals,real_add_closed,real_div,real_mul_closed,real_mul_inverse,real_lt_add_two,real_lt_subst_right}.
\end{proof}

% from §35 helper lemma 13: successive Tietze endpoint intervals are nested
\begin{lemma}\label{tietze_endpoint_successive_nesting}
    Assume $u$ is a sequence in $\naturalsPlus\times\funs{X}{\reals}$.
    Assume for all $n\in\naturalsPlus$ we have
        $\fst{u(n)} = n$ and
        $\snd{u(n)}$ is a Tietze unit state on $n$ for $f$ over $A$ in $X$ with respect to $T$ and
        for all $x\in X$ we have
            $\apply{\snd{u(n + 1)}}{x} - \apply{\snd{u(n)}}{x}\in
                \intervalclosed{\neg{(((1 + 1) / (n + 1)) - ((1 + 1) / ((n + 1) + 1)))}}{((1 + 1) / (n + 1)) - ((1 + 1) / ((n + 1) + 1))}$.
    Then for all $n\in\naturalsPlus$ we have
        for all $x\in X$ we have
            $\apply{\snd{u(n)}}{x} - ((1 + 1) / (n + 1))
                \leq \apply{\snd{u(n + 1)}}{x} - ((1 + 1) / ((n + 1) + 1))$
            and
            $\apply{\snd{u(n + 1)}}{x} + ((1 + 1) / ((n + 1) + 1))
                \leq \apply{\snd{u(n)}}{x} + ((1 + 1) / (n + 1))$.
\end{lemma}
\begin{proof}
    Fix $n\in\naturalsPlus$.
    Fix $x\in X$.
    Let $r = (1 + 1) / (n + 1)$.
    Let $e = (1 + 1) / ((n + 1) + 1)$.
    Let $d = r - e$.
    We have $\snd{u(n)}$ is continuous from $X$ to $\reals$ with respect to $T$ and $\standardTopologyR$
        by \cref{tietze_unit_state}.
    $\snd{u(n + 1)}$ is continuous from $X$ to $\reals$ with respect to $T$ and $\standardTopologyR$
        by \cref{real_naturals_closed_succ,tietze_unit_state}.
    $\apply{\snd{u(n)}}{x}\in\reals$ and $\apply{\snd{u(n + 1)}}{x}\in\reals$
        by \cref{continuous,funs_type_apply}.
    $n + 1\in\naturalsPlus$ and $((n + 1) + 1)\in\naturalsPlus$
        by \cref{real_naturals_closed_succ}.
    $n + 1\in\reals$ and $((n + 1) + 1)\in\reals$ by \cref{real_naturals_subset_reals,subseteq}.
    $0 < n + 1$ and $0 < ((n + 1) + 1)$ by \cref{naturalsplus_positive_real,subseteq}.
    $n + 1\neq 0$ and $((n + 1) + 1)\neq 0$ by \cref{real_pos_nonzero}.
    $1 + 1\in\reals$ by \cref{real_one_in_naturals,real_naturals_subset_reals,subseteq,real_add_closed}.
    $r\in\reals$ and $e\in\reals$ and $d\in\reals$ by \cref{real_div,real_mul_closed,real_mul_inverse,real_sub,real_add_inverse,real_add_closed}.
    $\neg{d}\in\reals$ and $\neg{e}\in\reals$ and $\neg{r}\in\reals$ by \cref{real_add_inverse}.
    $\apply{\snd{u(n)}}{x} + \neg{d}\in\reals$ and $\apply{\snd{u(n + 1)}}{x} + \neg{e}\in\reals$ and
        $(\apply{\snd{u(n)}}{x} + \neg{d}) + \neg{e}\in\reals$ and
        $d + \apply{\snd{u(n)}}{x}\in\reals$ and $\apply{\snd{u(n)}}{x} + d\in\reals$ and
        $(\apply{\snd{u(n)}}{x} + d) + e\in\reals$ by \cref{real_add_closed}.
    $\apply{\snd{u(n + 1)}}{x} - \apply{\snd{u(n)}}{x}\in \intervalclosed{\neg{d}}{d}$ by \cref{subseteq}.
    $\apply{\snd{u(n + 1)}}{x} - \apply{\snd{u(n)}}{x}\in\reals$ and
        $\neg{d} \leq \apply{\snd{u(n + 1)}}{x} - \apply{\snd{u(n)}}{x}$ and
        $\apply{\snd{u(n + 1)}}{x} - \apply{\snd{u(n)}}{x} \leq d$ by \cref{intervalclosed}.

    Show that $\apply{\snd{u(n)}}{x} - r \leq \apply{\snd{u(n + 1)}}{x} - e$.
    \begin{subproof}
        $\apply{\snd{u(n)}}{x} + \neg{d} \leq \apply{\snd{u(n + 1)}}{x}$ by
            \cref{real_leq_add,real_add_comm,real_add_sub_cancel_right,real_add_eq_subst}.
        $(\apply{\snd{u(n)}}{x} + \neg{d}) + \neg{e} \leq \apply{\snd{u(n + 1)}}{x} + \neg{e}$ by \cref{real_leq_add}.
        We have $\neg{d} + \neg{e} = \neg{(d + e)}$ by \cref{real_neg_addition}.
        $d + e = r$
            by \cref{real_sub,real_add_assoc,real_add_inverse,real_add_comm,real_add_ident,real_add_eq_subst}.
        $\neg{d} + \neg{e} = \neg{r}$ by \cref{real_add_eq_subst}.
        $\apply{\snd{u(n)}}{x} + \neg{r} \leq \apply{\snd{u(n + 1)}}{x} + \neg{e}$
            by \cref{real_add_assoc,real_add_eq_subst}.
        $\apply{\snd{u(n)}}{x} - r \leq \apply{\snd{u(n + 1)}}{x} - e$ by \cref{real_sub}.
    \end{subproof}

    Show that $\apply{\snd{u(n + 1)}}{x} + e \leq \apply{\snd{u(n)}}{x} + r$.
    \begin{subproof}
        $\apply{\snd{u(n + 1)}}{x} \leq \apply{\snd{u(n)}}{x} + d$ by
            \cref{real_leq_add,real_add_sub_cancel_right,real_add_comm,real_add_eq_subst}.
        $\apply{\snd{u(n + 1)}}{x} + e \leq (\apply{\snd{u(n)}}{x} + d) + e$ by \cref{real_leq_add}.
        We have $d + e = r$
            by \cref{real_sub,real_add_assoc,real_add_inverse,real_add_comm,real_add_ident,real_add_eq_subst}.
        $(\apply{\snd{u(n)}}{x} + d) + e = \apply{\snd{u(n)}}{x} + r$ by \cref{real_add_assoc,real_add_eq_subst}.
        $\apply{\snd{u(n + 1)}}{x} + e \leq \apply{\snd{u(n)}}{x} + r$
            by \cref{real_add_eq_subst}.
    \end{subproof}
\end{proof}

% from §35 helper lemma 14: Tietze lower endpoints increase and upper endpoints decrease
\begin{lemma}\label{tietze_endpoint_monotone}
    Assume $u$ is a sequence in $\naturalsPlus\times\funs{X}{\reals}$.
    Assume for all $n\in\naturalsPlus$ we have
        $\fst{u(n)} = n$ and
        $\snd{u(n)}$ is a Tietze unit state on $n$ for $f$ over $A$ in $X$ with respect to $T$ and
        for all $x\in X$ we have
            $\apply{\snd{u(n + 1)}}{x} - \apply{\snd{u(n)}}{x}\in
                \intervalclosed{\neg{(((1 + 1) / (n + 1)) - ((1 + 1) / ((n + 1) + 1)))}}{((1 + 1) / (n + 1)) - ((1 + 1) / ((n + 1) + 1))}$.
    Then for all $m,n\in\naturalsPlus$ if $m \leq n$ then
        for all $x\in X$ we have
            $\apply{\snd{u(m)}}{x} - ((1 + 1) / (m + 1))
                \leq \apply{\snd{u(n)}}{x} - ((1 + 1) / (n + 1))$
            and
            $\apply{\snd{u(n)}}{x} + ((1 + 1) / (n + 1))
                \leq \apply{\snd{u(m)}}{x} + ((1 + 1) / (m + 1))$.
\end{lemma}
\begin{proof}
    Fix $m$.
    Fix $n$.
    Assume $m\in\naturalsPlus$ and $n\in\naturalsPlus$.
    Assume $m \leq n$.
    Fix $x\in X$.
    Let $A_0 = \{k\in\naturalsPlus \mid \text{if $m \leq k$ then
        $\apply{\snd{u(m)}}{x} - ((1 + 1) / (m + 1))
            \leq \apply{\snd{u(k)}}{x} - ((1 + 1) / (k + 1))$ and
        $\apply{\snd{u(k)}}{x} + ((1 + 1) / (k + 1))
            \leq \apply{\snd{u(m)}}{x} + ((1 + 1) / (m + 1))$}\}$.
    $A_0\subseteq\naturalsPlus$ by \cref{subseteq}.
    Show that if $m \leq 1$ then
        $\apply{\snd{u(m)}}{x} - ((1 + 1) / (m + 1))
            \leq \apply{\snd{u(1)}}{x} - ((1 + 1) / (1 + 1))$ and
        $\apply{\snd{u(1)}}{x} + ((1 + 1) / (1 + 1))
            \leq \apply{\snd{u(m)}}{x} + ((1 + 1) / (m + 1))$.
    \begin{subproof}
        Follows by \cref{real_naturals_subset_reals,subseteq,real_one_in_naturals,real_one_leq_naturalsplus,real_leq_antisym,real_add_eq_subst}.
    \end{subproof}
    $1\in A_0$ by \cref{real_one_in_naturals}.
    Show that for all $k\in A_0$ we have $k + 1\in A_0$.
    \begin{subproof}
        Fix $k$ such that $k\in A_0$.
        $k + 1\in\naturalsPlus$ by \cref{subseteq,real_naturals_closed_succ}.
        Show that if $m \leq k + 1$ then
            $\apply{\snd{u(m)}}{x} - ((1 + 1) / (m + 1))
                \leq \apply{\snd{u(k + 1)}}{x} - ((1 + 1) / ((k + 1) + 1))$ and
            $\apply{\snd{u(k + 1)}}{x} + ((1 + 1) / ((k + 1) + 1))
                \leq \apply{\snd{u(m)}}{x} + ((1 + 1) / (m + 1))$.
        \begin{subproof}
            Assume $m \leq k + 1$.
            $m \leq k$ or $m = k + 1$ by \cref{real_naturals_leq_succ_elim}.
            \begin{byCase}
            \caseOf{$m \leq k$.}
                We have $\apply{\snd{u(m)}}{x} - ((1 + 1) / (m + 1))
                    \leq \apply{\snd{u(k)}}{x} - ((1 + 1) / (k + 1))$ and
                    $\apply{\snd{u(k)}}{x} + ((1 + 1) / (k + 1))
                        \leq \apply{\snd{u(m)}}{x} + ((1 + 1) / (m + 1))$ by definition.
                $\apply{\snd{u(k)}}{x} - ((1 + 1) / (k + 1))
                    \leq \apply{\snd{u(k + 1)}}{x} - ((1 + 1) / ((k + 1) + 1))$ and
                    $\apply{\snd{u(k + 1)}}{x} + ((1 + 1) / ((k + 1) + 1))
                        \leq \apply{\snd{u(k)}}{x} + ((1 + 1) / (k + 1))$
                    by \cref{tietze_endpoint_successive_nesting}.
                $\apply{\snd{u(m)}}{x} - ((1 + 1) / (m + 1))\in\reals$ and
                    $\apply{\snd{u(k)}}{x} - ((1 + 1) / (k + 1))\in\reals$ and
                    $\apply{\snd{u(k + 1)}}{x} - ((1 + 1) / ((k + 1) + 1))\in\reals$ and
                    $\apply{\snd{u(k + 1)}}{x} + ((1 + 1) / ((k + 1) + 1))\in\reals$ and
                    $\apply{\snd{u(k)}}{x} + ((1 + 1) / (k + 1))\in\reals$ and
                    $\apply{\snd{u(m)}}{x} + ((1 + 1) / (m + 1))\in\reals$ by
                    \cref{tietze_unit_state,continuous,funs_type_apply,real_naturals_closed_succ,real_one_in_naturals,real_naturals_subset_reals,subseteq,naturalsplus_positive_real,real_pos_nonzero,real_add_closed,real_div,real_mul_closed,real_mul_inverse,real_sub,real_add_inverse}.
                $\apply{\snd{u(m)}}{x} - ((1 + 1) / (m + 1))
                    \leq \apply{\snd{u(k + 1)}}{x} - ((1 + 1) / ((k + 1) + 1))$ and
                    $\apply{\snd{u(k + 1)}}{x} + ((1 + 1) / ((k + 1) + 1))
                        \leq \apply{\snd{u(m)}}{x} + ((1 + 1) / (m + 1))$
                    by \cref{real_leq_trans}.
            \caseOf{$m = k + 1$.}
                $\apply{\snd{u(m)}}{x} - ((1 + 1) / (m + 1))
                    \leq \apply{\snd{u(k + 1)}}{x} - ((1 + 1) / ((k + 1) + 1))$ and
                    $\apply{\snd{u(k + 1)}}{x} + ((1 + 1) / ((k + 1) + 1))
                        \leq \apply{\snd{u(m)}}{x} + ((1 + 1) / (m + 1))$ by \cref{real_add_eq_subst}.
            \end{byCase}
        \end{subproof}
        $k + 1\in A_0$ by definition.
    \end{subproof}
    For all $k\in\naturalsPlus$ we have $k\in A_0$ by \cref{real_naturals_induction}.
    $\apply{\snd{u(m)}}{x} - ((1 + 1) / (m + 1))
        \leq \apply{\snd{u(n)}}{x} - ((1 + 1) / (n + 1))$ and
        $\apply{\snd{u(n)}}{x} + ((1 + 1) / (n + 1))
            \leq \apply{\snd{u(m)}}{x} + ((1 + 1) / (m + 1))$ by \cref{subseteq}.
\end{proof}

% from §35 helper lemma 15: each lower endpoint lies below the matching upper endpoint
\begin{lemma}\label{tietze_lower_leq_upper}
    Assume $u$ is a sequence in $\naturalsPlus\times\funs{X}{\reals}$.
    Assume for all $n\in\naturalsPlus$ we have
        $\fst{u(n)} = n$ and
        $\snd{u(n)}$ is a Tietze unit state on $n$ for $f$ over $A$ in $X$ with respect to $T$ and
        for all $x\in X$ we have
            $\apply{\snd{u(n + 1)}}{x} - \apply{\snd{u(n)}}{x}\in
                \intervalclosed{\neg{(((1 + 1) / (n + 1)) - ((1 + 1) / ((n + 1) + 1)))}}{((1 + 1) / (n + 1)) - ((1 + 1) / ((n + 1) + 1))}$.
    Then for all $n\in\naturalsPlus$ we have
        for all $x\in X$ we have
            $\apply{\snd{u(n)}}{x} - ((1 + 1) / (n + 1))
                \leq \apply{\snd{u(n)}}{x} + ((1 + 1) / (n + 1))$.
\end{lemma}
\begin{proof}
    Fix $n\in\naturalsPlus$.
    Fix $x\in X$.
    Let $r = (1 + 1) / (n + 1)$.
    We have $\apply{\snd{u(n)}}{x}\in\reals$ and $r\in\reals$ and $0 < r$ by
        \cref{tietze_unit_state,continuous,funs_type_apply,real_naturals_closed_succ,real_naturals_subset_reals,subseteq,naturalsplus_positive_real,real_pos_nonzero,real_one_in_naturals,real_add_closed,real_pos_add,real_zero_lt_one,real_div,real_mul_closed,real_mul_inverse,real_inv_pos,real_mul_pos_pos}.
    $\neg{r} \leq r$ by
        \cref{real_add_inverse,real_add_ident,real_lt_neg,real_neg_zero,real_lt_subst_right,real_lt_imp_leq,real_leq_trans}.
    $\apply{\snd{u(n)}}{x} - r \leq \apply{\snd{u(n)}}{x} + r$ by
        \cref{real_add_inverse,real_add_closed,real_leq_add,real_add_comm,real_add_eq_subst,real_sub}.
\end{proof}

% from §35 helper lemma 15a: every lower endpoint lies below every upper endpoint
\begin{lemma}\label{tietze_lower_leq_upper_cross}
    Assume $u$ is a sequence in $\naturalsPlus\times\funs{X}{\reals}$.
    Assume for all $n\in\naturalsPlus$ we have
        $\fst{u(n)} = n$ and
        $\snd{u(n)}$ is a Tietze unit state on $n$ for $f$ over $A$ in $X$ with respect to $T$ and
        for all $x\in X$ we have
            $\apply{\snd{u(n + 1)}}{x} - \apply{\snd{u(n)}}{x}\in
                \intervalclosed{\neg{(((1 + 1) / (n + 1)) - ((1 + 1) / ((n + 1) + 1)))}}{((1 + 1) / (n + 1)) - ((1 + 1) / ((n + 1) + 1))}$.
    Then for all $m,n\in\naturalsPlus$ we have
        for all $x\in X$ we have
            $\apply{\snd{u(m)}}{x} - ((1 + 1) / (m + 1))
                \leq \apply{\snd{u(n)}}{x} + ((1 + 1) / (n + 1))$.
\end{lemma}
\begin{proof}
    Fix $m$.
    Fix $n$.
    Assume $m\in\naturalsPlus$ and $n\in\naturalsPlus$.
    Fix $x\in X$.
    We have $\apply{\snd{u(m)}}{x} - ((1 + 1) / (m + 1))\in\reals$ and
        $\apply{\snd{u(m)}}{x} + ((1 + 1) / (m + 1))\in\reals$ and
        $\apply{\snd{u(n)}}{x} - ((1 + 1) / (n + 1))\in\reals$ and
        $\apply{\snd{u(n)}}{x} + ((1 + 1) / (n + 1))\in\reals$ by
        \cref{tietze_unit_state,continuous,funs_type_apply,real_naturals_closed_succ,real_naturals_subset_reals,subseteq,naturalsplus_positive_real,real_pos_nonzero,real_one_in_naturals,real_add_closed,real_div,real_mul_closed,real_mul_inverse,real_sub,real_add_inverse}.
    $m < n$ or $m = n$ or $n < m$ by \cref{real_naturals_subset_reals,subseteq,real_lt_trichotomy}.
    \begin{byCase}
    \caseOf{$m < n$.}
        $\apply{\snd{u(m)}}{x} - ((1 + 1) / (m + 1))
            \leq \apply{\snd{u(n)}}{x} - ((1 + 1) / (n + 1))$
            by \cref{real_lt_imp_leq,tietze_endpoint_monotone}.
        $\apply{\snd{u(m)}}{x} - ((1 + 1) / (m + 1))
            \leq \apply{\snd{u(n)}}{x} + ((1 + 1) / (n + 1))$ by \cref{real_leq_trans,tietze_lower_leq_upper}.
    \caseOf{$m = n$.}
        $\apply{\snd{u(m)}}{x} - ((1 + 1) / (m + 1))
            \leq \apply{\snd{u(n)}}{x} + ((1 + 1) / (n + 1))$ by \cref{tietze_lower_leq_upper,real_add_eq_subst}.
    \caseOf{$n < m$.}
        $\apply{\snd{u(m)}}{x} + ((1 + 1) / (m + 1))
            \leq \apply{\snd{u(n)}}{x} + ((1 + 1) / (n + 1))$
            by \cref{real_lt_imp_leq,tietze_endpoint_monotone}.
        $\apply{\snd{u(m)}}{x} - ((1 + 1) / (m + 1))
            \leq \apply{\snd{u(n)}}{x} + ((1 + 1) / (n + 1))$ by \cref{real_leq_trans,tietze_lower_leq_upper}.
    \end{byCase}
\end{proof}

% from §35 helper lemma 16: lower endpoint sets consist of real numbers
\begin{lemma}\label{tietze_lower_set_subset_reals}
    Assume $u$ is a sequence in $\naturalsPlus\times\funs{X}{\reals}$.
    Assume for all $n\in\naturalsPlus$ we have
        $\fst{u(n)} = n$ and
        $\snd{u(n)}$ is a Tietze unit state on $n$ for $f$ over $A$ in $X$ with respect to $T$ and
        for all $x\in X$ we have
            $\apply{\snd{u(n + 1)}}{x} - \apply{\snd{u(n)}}{x}\in
                \intervalclosed{\neg{(((1 + 1) / (n + 1)) - ((1 + 1) / ((n + 1) + 1)))}}{((1 + 1) / (n + 1)) - ((1 + 1) / ((n + 1) + 1))}$.
    Then for all $x\in X$ we have $\tietzeLowerSet{u}{x}\subseteq\reals$.
\end{lemma}
\begin{proof}
    Fix $x\in X$.
    We have $\tietzeLowerSet{u}{x}\subseteq\reals$
        by \cref{tietze_lower_set,tietze_unit_state,continuous,funs_type_apply,real_naturals_closed_succ,real_naturals_subset_reals,subseteq,naturalsplus_positive_real,real_pos_nonzero,real_one_in_naturals,real_add_closed,real_div,real_mul_closed,real_mul_inverse,real_sub,real_add_inverse}.
\end{proof}

% from §35 helper lemma 16a: two over two equals one
\begin{lemma}\label{real_two_over_two_local}
    $((1 + 1) / (1 + 1)) = 1$.
\end{lemma}
\begin{proof}
    Follows by \cref{real_div,real_mul_assoc,real_mul_comm,real_mul_inverse,real_mul_ident,real_add_closed,real_zero_lt_one,real_pos_add,real_pos_nonzero,real_naturals_subset_reals,real_one_in_naturals,subseteq}.
\end{proof}

% from §35 helper lemma 17: each lower endpoint set has a supremum
\begin{lemma}\label{tietze_lower_set_has_supremum}
    Assume $u$ is a sequence in $\naturalsPlus\times\funs{X}{\reals}$.
    Assume for all $x\in X$ we have $\apply{\snd{u(1)}}{x} = 0$.
    Assume for all $n\in\naturalsPlus$ we have
        $\fst{u(n)} = n$ and
        $\snd{u(n)}$ is a Tietze unit state on $n$ for $f$ over $A$ in $X$ with respect to $T$ and
        for all $x\in X$ we have
            $\apply{\snd{u(n + 1)}}{x} - \apply{\snd{u(n)}}{x}\in
                \intervalclosed{\neg{(((1 + 1) / (n + 1)) - ((1 + 1) / ((n + 1) + 1)))}}{((1 + 1) / (n + 1)) - ((1 + 1) / ((n + 1) + 1))}$.
    Then for all $x\in X$ there exists $r$ such that $r$ is a supremum of $\tietzeLowerSet{u}{x}$.
\end{lemma}
\begin{proof}
    Fix $x\in X$.
    We have $\tietzeLowerSet{u}{x}\subseteq\reals$ by \cref{tietze_lower_set_subset_reals}.
    $\tietzeLowerSet{u}{x}\neq\emptyset$ by \cref{real_one_in_naturals,tietze_lower_set,emptyset}.
    Show that for all $y\in\tietzeLowerSet{u}{x}$ we have $y \leq 1$.
    \begin{subproof}
        Fix $y\in\tietzeLowerSet{u}{x}$.
        Take $n\in\naturalsPlus$ such that
            $y = \apply{\snd{u(n)}}{x} - ((1 + 1) / (n + 1))$ by \cref{tietze_lower_set}.
        $y \leq \apply{\snd{u(1)}}{x} + ((1 + 1) / (1 + 1))$
            by \cref{tietze_lower_leq_upper_cross,real_one_in_naturals,real_add_eq_subst}.
        $\apply{\snd{u(1)}}{x} = 0$ by assumption.
        $y \leq 1$ by
            \cref{real_two_over_two_local,real_naturals_subset_reals,real_one_in_naturals,subseteq,real_add_ident,real_leq_split,real_lt_subst_right,real_lt_imp_leq,real_add_eq_subst}.
    \end{subproof}
    $\tietzeLowerSet{u}{x}$ is bounded above
        by \cref{real_naturals_subset_reals,real_one_in_naturals,subseteq,real_upper_bound,real_bounded_above}.
    There exists $r$ such that $r$ is a supremum of $\tietzeLowerSet{u}{x}$ by \cref{real_completeness_axiom}.
\end{proof}

% from §35 helper lemma 18: choose the supremum function of the lower endpoint sets
\begin{lemma}\label{tietze_lower_set_supremum_function}
    Assume $u$ is a sequence in $\naturalsPlus\times\funs{X}{\reals}$.
    Assume for all $x\in X$ we have $\apply{\snd{u(1)}}{x} = 0$.
    Assume for all $n\in\naturalsPlus$ we have
        $\fst{u(n)} = n$ and
        $\snd{u(n)}$ is a Tietze unit state on $n$ for $f$ over $A$ in $X$ with respect to $T$ and
        for all $x\in X$ we have
            $\apply{\snd{u(n + 1)}}{x} - \apply{\snd{u(n)}}{x}\in
                \intervalclosed{\neg{(((1 + 1) / (n + 1)) - ((1 + 1) / ((n + 1) + 1)))}}{((1 + 1) / (n + 1)) - ((1 + 1) / ((n + 1) + 1))}$.
    Then there exists $g$ such that
        $g\in\funs{X}{\reals}$ and
        for all $x\in X$ we have $g(x)$ is a supremum of $\tietzeLowerSet{u}{x}$.
\end{lemma}
\begin{proof}
    Let $R = \{(x,r) \mid x\in X, r\in\reals \mid \text{$r$ is a supremum of $\tietzeLowerSet{u}{x}$}\}$.
    We have $R$ is a relation by \cref{relation}.
    For all $x\in X$ there exists $r$ such that $x\mathrel{R} r$
        by \cref{tietze_lower_set_has_supremum,tietze_lower_set_subset_reals,real_supremum,real_upper_bound,relation,pair_eq_iff,subseteq}.
    Take $g$ such that $g$ is a function on $X$ and for all $x\in X$ we have $x\mathrel{R}\apply{g}{x}$
        by \cref{choice_function}.
    For all $x\in X$ we have $g(x)$ is a supremum of $\tietzeLowerSet{u}{x}$ and $g(x)\in\reals$ by
        \cref{relation,pair_eq_iff,tietze_lower_set_subset_reals,real_supremum,real_upper_bound}.
    $g\in\funs{X}{\reals}$ by \cref{funs_intro,funs}.
    Follows by \cref{subseteq}.
\end{proof}

% from §35 helper lemma 19: the supremum stays inside every endpoint interval
\begin{lemma}\label{tietze_supremum_between_endpoints}
    Assume $u$ is a sequence in $\naturalsPlus\times\funs{X}{\reals}$.
    Assume for all $n\in\naturalsPlus$ we have
        $\fst{u(n)} = n$ and
        $\snd{u(n)}$ is a Tietze unit state on $n$ for $f$ over $A$ in $X$ with respect to $T$ and
        for all $x\in X$ we have
            $\apply{\snd{u(n + 1)}}{x} - \apply{\snd{u(n)}}{x}\in
                \intervalclosed{\neg{(((1 + 1) / (n + 1)) - ((1 + 1) / ((n + 1) + 1)))}}{((1 + 1) / (n + 1)) - ((1 + 1) / ((n + 1) + 1))}$.
    Assume $g\in\funs{X}{\reals}$.
    Assume for all $x\in X$ we have $g(x)$ is a supremum of $\tietzeLowerSet{u}{x}$.
    Then for all $n\in\naturalsPlus$ we have
        for all $x\in X$ we have
            $\apply{\snd{u(n)}}{x} - ((1 + 1) / (n + 1)) \leq g(x)$ and
            $g(x) \leq \apply{\snd{u(n)}}{x} + ((1 + 1) / (n + 1))$.
\end{lemma}
\begin{proof}
    Fix $n\in\naturalsPlus$.
    Fix $x\in X$.
    We have $g(x)$ is a supremum of $\tietzeLowerSet{u}{x}$ and $g(x)\in\reals$
        by \cref{funs_type_apply}.
    $\apply{\snd{u(n)}}{x}\in\reals$ and $((1 + 1) / (n + 1))\in\reals$ by
        \cref{tietze_unit_state,continuous,funs_type_apply,real_one_in_naturals,real_naturals_subset_reals,subseteq,real_add_closed,real_naturals_closed_succ,real_div_naturalsplus_real}.
    $\apply{\snd{u(n)}}{x} - ((1 + 1) / (n + 1))\in\reals$ and
        $\apply{\snd{u(n)}}{x} + ((1 + 1) / (n + 1))\in\reals$ by
        \cref{real_sub,real_add_inverse,real_add_closed}.
    Show that $\apply{\snd{u(n)}}{x} - ((1 + 1) / (n + 1)) \leq g(x)$.
    \begin{subproof}
        Follows by \cref{tietze_lower_set,real_supremum_pred,real_supremum,real_upper_bound}.
    \end{subproof}
    Show that for all $y\in\tietzeLowerSet{u}{x}$ we have
        $y \leq \apply{\snd{u(n)}}{x} + ((1 + 1) / (n + 1))$.
    \begin{subproof}
        Fix $y\in\tietzeLowerSet{u}{x}$.
        Take $m\in\naturalsPlus$ such that
            $y = \apply{\snd{u(m)}}{x} - ((1 + 1) / (m + 1))$ by \cref{tietze_lower_set}.
        $y \leq \apply{\snd{u(n)}}{x} + ((1 + 1) / (n + 1))$
            by \cref{tietze_lower_leq_upper_cross,real_add_eq_subst}.
    \end{subproof}
    $g(x) \leq \apply{\snd{u(n)}}{x} + ((1 + 1) / (n + 1))$
        by \cref{tietze_lower_set_subset_reals,real_upper_bound,real_supremum}.
\end{proof}

% from §35 helper lemma 20: substitute equality on the left of a weak inequality
\begin{lemma}\label{real_leq_subst_left_local}
    Suppose $a = b$ and $b \leq c$.
    Then $a \leq c$.
\end{lemma}
\begin{proof}
    Follows by \cref{real_leq_split,real_lt_subst_left,real_lt_imp_leq,subseteq}.
\end{proof}

% from §35 helper lemma 21: substitute equality on the right of a weak inequality
\begin{lemma}\label{real_leq_subst_right_local}
    Suppose $a \leq b$ and $b = c$.
    Then $a \leq c$.
\end{lemma}
\begin{proof}
    Follows by \cref{real_leq_split,real_lt_subst_right,real_lt_imp_leq,subseteq}.
\end{proof}

% from §35 helper lemma 22: strong-weak transitivity for reals
\begin{lemma}\label{real_lt_leq_trans_local}
    Suppose $a\in\reals$ and $b\in\reals$ and $c\in\reals$ and $a < b$ and $b \leq c$.
    Then $a < c$.
\end{lemma}
\begin{proof}
    Follows by \cref{real_leq_split,real_lt_trans,real_lt_subst_right}.
\end{proof}

% from §35 helper lemma 23: uniform convergence of Tietze lower-set approximants
\begin{lemma}\label{tietze_lower_set_uniform_convergence}
    Suppose $u$ is a sequence in $\naturalsPlus\times\funs{X}{\reals}$.
    Suppose for all $n\in\naturalsPlus$ we have
        $\fst{u(n)} = n$ and
        $\snd{u(n)}$ is a Tietze unit state on $n$ for $f$ over $A$ in $X$ with respect to $T$ and
        for all $x\in X$ we have
            $\apply{\snd{u(n + 1)}}{x} - \apply{\snd{u(n)}}{x}\in
                \intervalclosed{\neg{(((1 + 1) / (n + 1)) - ((1 + 1) / ((n + 1) + 1)))}}{((1 + 1) / (n + 1)) - ((1 + 1) / ((n + 1) + 1))}$.
    Suppose $g\in\funs{X}{\reals}$.
    Suppose for all $x\in X$ we have $g(x)$ is a supremum of $\tietzeLowerSet{u}{x}$.
    Let $s = \{(n,\snd{u(n)}) \mid n\in\naturalsPlus\}$.
    Then $s$ converges uniformly to $g$ from $X$ to $\reals$ with respect to $\standardMetricR$.
\end{lemma}
\begin{proof}
    $s$ is a relation by \cref{relation}.
    For all $n,v,w$ such that $(n,v)\in s$ and $(n,w)\in s$ we have $v = w$ by \cref{pair_eq_iff}.
    $s$ is a function by \cref{rightunique,relation}.
    $\dom{s} = \naturalsPlus$ by \cref{dom_intro,dom_iff,pair_eq_iff,subseteq,subseteq_antisymmetric}.
    For all $n\in\dom{s}$ we have $s(n)\in\funs{X}{\reals}$
        by \cref{function_apply_iff,subseteq,pair_eq_iff,tietze_unit_state,continuous,funs}.
    $s$ is a sequence in $\funs{X}{\reals}$ by \cref{funs_intro,sequence_in}.

    Show that for all $\epsilon\in\reals$ such that $0 < \epsilon$ there exists $N\in\naturalsPlus$ such that
        for all $n\in\naturalsPlus$ if $N\leq n$ then
            for all $x\in X$ we have $\apply{\standardMetricR}{(\apply{\apply{s}{n}}{x}, g(x))} < \epsilon$.
    \begin{subproof}
        Fix $\epsilon\in\reals$.
        Assume $0 < \epsilon$.
        Take $N\in\naturalsPlus$ such that $((1 + 1) / (N + 1)) < \epsilon$ by \cref{tietze_small_error}.
        Show that for all $n\in\naturalsPlus$ if $N\leq n$ then
            for all $x\in X$ we have $\apply{\standardMetricR}{(\apply{\apply{s}{n}}{x}, g(x))} < \epsilon$.
        \begin{subproof}
            Fix $n\in\naturalsPlus$.
            Assume $N \leq n$.
            Fix $x\in X$.
            We have $\snd{u(n)}$ is a Tietze unit state on $n$ for $f$ over $A$ in $X$ with respect to $T$ by assumption.
            $\apply{\snd{u(n)}}{x}\in\reals$ by \cref{tietze_unit_state,continuous,funs_type_apply}.
            $\apply{\apply{s}{n}}{x} = \apply{\snd{u(n)}}{x}$ by \cref{function_apply_intro,function_apply_iff}.
            $\apply{\apply{s}{n}}{x}\in\reals$ by \cref{sequence_in,funs_type_apply}.
            $g(x)\in\reals$ by \cref{funs_type_apply}.
            $N\in\reals$ and $n\in\reals$ by \cref{real_naturals_subset_reals,subseteq}.
            $N + 1\in\naturalsPlus$ and $n + 1\in\naturalsPlus$ by \cref{real_naturals_closed_succ}.
            $1\in\reals$ and $N + 1\in\reals$ and $n + 1\in\reals$ and $1 + 1\in\reals$ by
                \cref{real_naturals_subset_reals,subseteq,real_one_in_naturals,real_add_closed}.
            $N + 1\neq 0$ and $n + 1\neq 0$ by \cref{naturalsplus_positive_real,subseteq,real_pos_nonzero}.
            $1 / (N + 1)\in\reals$ and $1 / (n + 1)\in\reals$ and
                $(1 + 1) / (N + 1)\in\reals$ and $((1 + 1) / (n + 1))\in\reals$ by
                \cref{real_div,real_mul_closed,real_mul_inverse}.
            $N + 1 \leq n + 1$ by \cref{real_leq_add}.
            $1 / (n + 1) \leq 1 / (N + 1)$ by \cref{naturalsplus_reciprocal_antitone}.
            $(1 / (n + 1)) \rmul (1 + 1) \leq (1 / (N + 1)) \rmul (1 + 1)$
                by \cref{real_zero_lt_one,real_pos_add,real_leq_split,real_lt_mul_pos,real_lt_imp_leq,real_mul_eq_subst,subseteq}.
            $(1 / (n + 1)) \rmul (1 + 1) = (1 + 1) \rmul (1 / (n + 1))$ and
                $(1 / (N + 1)) \rmul (1 + 1) = (1 + 1) \rmul (1 / (N + 1))$ by \cref{real_mul_comm}.
            $1 / (n + 1) + 1 / (n + 1)\in\reals$ and
                $1 / (N + 1) + 1 / (n + 1)\in\reals$ and
                $1 / (n + 1) + 1 / (N + 1)\in\reals$ and
                $1 / (N + 1) + 1 / (N + 1)\in\reals$ by \cref{real_add_closed}.
            $1 / (n + 1) + 1 / (n + 1) \leq 1 / (N + 1) + 1 / (n + 1)$ by \cref{real_leq_add}.
            $1 / (n + 1) + 1 / (N + 1) \leq 1 / (N + 1) + 1 / (N + 1)$ by \cref{real_leq_add}.
            $1 / (n + 1) + 1 / (n + 1) \leq 1 / (n + 1) + 1 / (N + 1)$
                by \cref{real_add_comm,real_leq_subst_right_local}.
            $1 / (n + 1) + 1 / (n + 1) \leq 1 / (N + 1) + 1 / (N + 1)$ by \cref{real_leq_trans}.
            $(1 + 1) / (n + 1) = 1 / (n + 1) + 1 / (n + 1)$ and
                $(1 + 1) / (N + 1) = 1 / (N + 1) + 1 / (N + 1)$ by
                \cref{real_div_add_same_denom,real_mul_ident,real_lt_subst_left,real_lt_subst_right}.
            $((1 + 1) / (n + 1)) \leq ((1 + 1) / (N + 1))$ by
                \cref{real_leq_subst_left_local,real_leq_subst_right_local}.
            $\apply{\snd{u(n)}}{x} - ((1 + 1) / (n + 1)) \leq g(x)$ and
                $g(x) \leq \apply{\snd{u(n)}}{x} + ((1 + 1) / (n + 1))$ by \cref{tietze_supremum_between_endpoints}.
            $\neg{(((1 + 1) / (N + 1)))}\in\reals$ and $\neg{(((1 + 1) / (n + 1)))}\in\reals$
                by \cref{real_add_inverse}.
            $\neg{(((1 + 1) / (N + 1)))} \leq \neg{(((1 + 1) / (n + 1)))}$ by \cref{real_leq_neg}.
            $\apply{\snd{u(n)}}{x} + \neg{(((1 + 1) / (N + 1)))}\in\reals$ and
                $\apply{\snd{u(n)}}{x} + \neg{(((1 + 1) / (n + 1)))}\in\reals$ by \cref{real_add_closed}.
            $\neg{(((1 + 1) / (N + 1)))} + \apply{\snd{u(n)}}{x} \leq
                \neg{(((1 + 1) / (n + 1)))} + \apply{\snd{u(n)}}{x}$ by \cref{real_leq_add}.
            $\apply{\snd{u(n)}}{x} + \neg{(((1 + 1) / (N + 1)))} \leq
                \apply{\snd{u(n)}}{x} + \neg{(((1 + 1) / (n + 1)))}$ by
                \cref{real_add_comm,real_leq_subst_left_local,real_leq_subst_right_local}.
            $\apply{\snd{u(n)}}{x} - ((1 + 1) / (N + 1))
                \leq \apply{\snd{u(n)}}{x} - ((1 + 1) / (n + 1))$ by \cref{real_sub}.
            $\apply{\snd{u(n)}}{x} - ((1 + 1) / (N + 1))\in\reals$ and
                $\apply{\snd{u(n)}}{x} - ((1 + 1) / (n + 1))\in\reals$ by \cref{real_sub}.
            $\apply{\apply{s}{n}}{x} - ((1 + 1) / (N + 1)) =
                \apply{\snd{u(n)}}{x} - ((1 + 1) / (N + 1))$ by \cref{real_sub}.
            $\apply{\snd{u(n)}}{x} - ((1 + 1) / (N + 1)) \leq g(x)$ by \cref{real_leq_trans}.
            $\apply{\apply{s}{n}}{x} - ((1 + 1) / (N + 1)) \leq g(x)$ by \cref{real_leq_subst_left_local}.
            $((1 + 1) / (n + 1)) + \apply{\snd{u(n)}}{x} \leq
                ((1 + 1) / (N + 1)) + \apply{\snd{u(n)}}{x}$ by \cref{real_leq_add}.
            $\apply{\snd{u(n)}}{x} + ((1 + 1) / (n + 1)) \leq
                \apply{\snd{u(n)}}{x} + ((1 + 1) / (N + 1))$ by
                \cref{real_add_comm,real_leq_subst_left_local,real_leq_subst_right_local}.
            $\apply{\snd{u(n)}}{x} + ((1 + 1) / (n + 1))\in\reals$ and
                $\apply{\snd{u(n)}}{x} + ((1 + 1) / (N + 1))\in\reals$ by \cref{real_add_closed}.
            $g(x) \leq \apply{\snd{u(n)}}{x} + ((1 + 1) / (N + 1))$ by \cref{real_leq_trans}.
            $g(x) \leq \apply{\apply{s}{n}}{x} + ((1 + 1) / (N + 1))$
                by \cref{real_add_eq_subst,real_leq_subst_right_local}.
            $\apply{\apply{s}{n}}{x} - ((1 + 1) / (N + 1)) \leq g(x)$ and
                $g(x) \leq \apply{\apply{s}{n}}{x} + ((1 + 1) / (N + 1))$ by \cref{subseteq}.
            $\neg{\epsilon}\in\reals$ by \cref{real_add_inverse}.
            $\neg{\epsilon} < \neg{(((1 + 1) / (N + 1)))}$ by \cref{real_lt_neg,real_lt_subst_right}.
            $\neg{\epsilon} + \apply{\apply{s}{n}}{x} < \neg{(((1 + 1) / (N + 1)))} + \apply{\apply{s}{n}}{x}$ by \cref{real_lt_add}.
            $\apply{\apply{s}{n}}{x} + \neg{\epsilon}\in\reals$ and
                $\apply{\apply{s}{n}}{x} + \neg{(((1 + 1) / (N + 1)))}\in\reals$ and
                $\apply{\apply{s}{n}}{x} + \epsilon\in\reals$ by \cref{real_add_closed}.
            $\apply{\apply{s}{n}}{x} + \neg{\epsilon} < \apply{\apply{s}{n}}{x} + \neg{(((1 + 1) / (N + 1)))}$
                by \cref{real_add_comm,real_lt_subst_left,real_lt_subst_right}.
            $\apply{\apply{s}{n}}{x} + \neg{(((1 + 1) / (N + 1)))} \leq g(x)$
                by \cref{real_sub,subseteq,real_leq_subst_left_local}.
            $\apply{\apply{s}{n}}{x} + \neg{\epsilon} < g(x)$ by \cref{real_lt_leq_trans_local}.
            $\apply{\apply{s}{n}}{x} - \epsilon < g(x)$ by \cref{real_sub,real_lt_subst_left}.
            $((1 + 1) / (N + 1)) + \apply{\apply{s}{n}}{x} < \epsilon + \apply{\apply{s}{n}}{x}$ by \cref{real_lt_add}.
            $\apply{\apply{s}{n}}{x} + ((1 + 1) / (N + 1)) < \apply{\apply{s}{n}}{x} + \epsilon$
                by \cref{real_add_comm,real_lt_subst_left,real_lt_subst_right}.
            $g(x) < \apply{\apply{s}{n}}{x} + \epsilon$ by \cref{real_leq_lt_trans}.
            $\abs{\apply{\apply{s}{n}}{x} - g(x)} < \epsilon$ by \cref{real_abs_lt_from_interval}.
            $\apply{\standardMetricR}{(\apply{\apply{s}{n}}{x}, g(x))} < \epsilon$ by \cref{standard_metric_value,real_lt_subst_left}.
        \end{subproof}
        Follows by \cref{subseteq}.
    \end{subproof}
    We have $s$ is a sequence in $\funs{X}{\reals}$ and
        $g\in\funs{X}{\reals}$ and
        for all $\epsilon\in\reals$ such that $0 < \epsilon$ there exists $N\in\naturalsPlus$ such that
            for all $n\in\naturalsPlus$ if $N\leq n$ then
                for all $x\in X$ we have $\apply{\standardMetricR}{(\apply{\apply{s}{n}}{x}, g(x))} < \epsilon$
        by \cref{funs,subseteq}.
    $s$ converges uniformly to $g$ from $X$ to $\reals$ with respect to $\standardMetricR$
        by \cref{uniform_converges_to}.
\end{proof}

% from §35 helper lemma 24: Tietze extension theorem for the symmetric unit interval
\begin{lemma}\label{tietze_extension_theorem_unit_interval}
    Assume $X$ is normal with respect to $T$.
    Assume $A$ is closed in $X$ with respect to $T$.
    Let $I = \intervalclosed{\neg{1}}{1}$.
    Let $T_A = \subspaceTopology{T}{A}$.
    Let $T_I = \subspaceTopology{\standardTopologyR}{I}$.
    Assume $f$ is continuous from $A$ to $I$ with respect to $T_A$ and $T_I$.
    Then there exists $g$ such that
        $g$ is continuous from $X$ to $I$ with respect to $T$ and $T_I$ and
        for all $x\in A$ we have $g(x) = f(x)$.
\end{lemma}
\begin{proof}
    \begin{byCase}
    \caseOf{$X = \emptyset$.}
        $T$ is a topology on $X$ by \cref{normal_space}.
        We have $\standardTopologyR$ is a topology on $\reals$
            by \cref{standard_basis_is_basis,basis_topology_is_topology,standard_topology_reals}.
        $T_I$ is a topology on $I$ by \cref{intervalclosed,subseteq,subspace_topology_is_topology,subspace_of}.
        $A = \emptyset$ by \cref{closed_in,real_add_eq_subst,subseteq_emptyset_iff}.
        $0\in\reals$ and $1\in\reals$ and $0 < 1$
            by \cref{real_add_ident,real_one_in_naturals,real_naturals_subset_reals,subseteq,real_zero_lt_one}.
        $\neg{1}\in\reals$ and $\neg{1} < 0$ by \cref{real_add_inverse,real_lt_neg,real_neg_zero,real_lt_subst_right}.
        $\neg{1} \leq 0$ and $0 \leq 1$ by \cref{real_lt_imp_leq}.
        $0\in I$ by \cref{intervalclosed}.
        Let $g = \{(x,0) \mid x\in X\}$.
        $g$ is a function from $X$ to $I$ by \cref{constant_map_function}.
        For all $x\in X$ we have $g(x) = 0$ by \cref{constant_map_function,funs_is_function,function_apply_intro}.
        Show that $g$ is continuous from $X$ to $I$ with respect to $T$ and $T_I$.
        \begin{subproof}
            Follows by \cref{constant_continuous}.
        \end{subproof}
        Show that for all $x\in A$ we have $g(x) = f(x)$.
        \begin{subproof}
            Trivial.
        \end{subproof}
        There exists $h$ such that
            $h$ is continuous from $X$ to $I$ with respect to $T$ and $T_I$ and
            for all $x\in A$ we have $h(x) = f(x)$ by \cref{subseteq}.
    \caseOf{$X \neq \emptyset$.}
        We have $X$ is inhabited by \cref{nonempty_inhabited}.
    Let $j = \identity{I}$.
    Let $f_1 = j\circ f$.
    $f_1$ is continuous from $A$ to $\reals$ with respect to $T_A$ and $\standardTopologyR$
        by \cref{continuous_interval_expand_range}.
    For all $a\in A$ we have $f_1(a)\in I$ by
        \cref{continuous,funs_type_apply,funs_is_function,funs,function_rightunique,funs_is_relation,id_rightunique,id_is_relation,id_dom,circ_apply,funs_ran,subseteq,id_apply,real_add_eq_subst}.
    Take $u$ such that
        $u$ is a sequence in $\naturalsPlus\times\funs{X}{\reals}$ and
        for all $n\in\naturalsPlus$ we have
            $\fst{u(n)} = n$ and
            $\snd{u(n)}$ is a Tietze unit state on $n$ for $f_1$ over $A$ in $X$ with respect to $T$ and
            for all $x\in X$ we have
                $\apply{\snd{u(1)}}{x} = 0$ and
                $\apply{\snd{u(n + 1)}}{x} - \apply{\snd{u(n)}}{x}\in
                    \intervalclosed{\neg{(((1 + 1) / (n + 1)) - ((1 + 1) / ((n + 1) + 1)))}}{((1 + 1) / (n + 1)) - ((1 + 1) / ((n + 1) + 1))}$
        by \cref{tietze_unit_state_iteration}.
    For all $x\in X$ we have $\apply{\snd{u(1)}}{x} = 0$ by \cref{real_one_in_naturals,subseteq}.
    Take $g$ such that
        $g\in\funs{X}{\reals}$ and
        for all $x\in X$ we have $g(x)$ is a supremum of $\tietzeLowerSet{u}{x}$
        by \cref{tietze_lower_set_supremum_function}.
    Let $s = \{(n,\snd{u(n)}) \mid n\in\naturalsPlus\}$.
    $s$ converges uniformly to $g$ from $X$ to $\reals$ with respect to $\standardMetricR$
        by \cref{tietze_lower_set_uniform_convergence}.
    $s\in\funs{\naturalsPlus}{\funs{X}{\reals}}$ by \cref{uniform_converges_to,sequence_in}.
    $s$ is a function and $s$ is a relation and $\dom{s} = \naturalsPlus$ by
        \cref{funs_is_function,funs_is_relation,funs}.
    For all $n\in\dom{s}$ we have $s(n)\in\funs{X}{\reals}$ by \cref{funs_type_apply}.
    $s$ is a sequence in $\funs{X}{\reals}$ by \cref{sequence_in}.
    $T$ is a topology on $X$ by \cref{normal_space}.
    $\standardMetricR$ is a metric on $\reals$ by \cref{standard_metric_is_metric}.
    Let $T_M = \metricTopology{\standardMetricR}{\reals}$.
    $T_M = \standardTopologyR$ by \cref{standard_metric_topology}.
    $\standardTopologyR = T_M$ by \cref{subseteq}.
    For all $n\in\naturalsPlus$ we have
        $s(n)$ is continuous from $X$ to $\reals$ with respect to $T$ and $T_M$
        by \cref{tietze_unit_state,function_apply_intro,continuous_eq_subst,continuous_eq_codomain}.
    $g$ is continuous from $X$ to $\reals$ with respect to $T$ and $T_M$
        by \cref{uniform_limit_theorem}.
    $g$ is continuous from $X$ to $\reals$ with respect to $T$ and $\standardTopologyR$
        by \cref{continuous_eq_codomain}.

    For all $x\in X$ we have $\apply{\snd{u(1)}}{x} = 0$ by assumption.
    For all $x\in X$ we have $g(x)\in I$ by
        \cref{tietze_supremum_between_endpoints,real_one_in_naturals,real_naturals_subset_reals,subseteq,real_two_over_two_local,real_add_inverse,real_sub,real_add_eq_subst,real_add_ident,real_leq_subst_left_local,real_leq_subst_right_local,intervalclosed,funs_type_apply}.
    Show that $g$ is continuous from $X$ to $I$ with respect to $T$ and $T_I$.
    \begin{subproof}
        Follows by \cref{img_iff,funs_is_function,function_rightunique,funs_is_relation,funs,function_apply_iff,pair_eq_iff,subseteq,continuous_subspace_range,intervalclosed,normal_space}.
    \end{subproof}

    Show that for all $a\in A$ we have $g(a) = f_1(a)$.
    \begin{subproof}
        Fix $a\in A$.
        We have $a\in X$ and $f_1(a)\in\reals$ and $g(a)\in\reals$
            by \cref{closed_in,subseteq,continuous,funs_type_apply}.
        $f_1(a) < g(a)$ or $f_1(a) = g(a)$ or $g(a) < f_1(a)$ by \cref{real_lt_trichotomy}.
        \begin{byCase}
        \caseOf{$f_1(a) = g(a)$.}
            $g(a) = f_1(a)$.
        \caseOf{$g(a) < f_1(a)$.}
            Suppose not.
            Let $\epsilon = f_1(a) - g(a)$.
            $\epsilon\in\reals$ and $0 < \epsilon$ by \cref{real_sub,real_add_inverse,real_add_closed,real_lt_imp_pos_diff}.
            Take $N\in\naturalsPlus$ such that
                $((1 + 1) / (N + 1)) + ((1 + 1) / (N + 1)) < \epsilon$ by \cref{tietze_small_double_error}.
            $\snd{u(N)}$ is a Tietze unit state on $N$ for $f_1$ over $A$ in $X$ with respect to $T$ by assumption.
            $f_1(a) - \apply{\snd{u(N)}}{a}\in
                \intervalclosed{\neg{((1 + 1) / (N + 1))}}{(1 + 1) / (N + 1)}$ by \cref{tietze_unit_state}.
            $f_1(a) - \apply{\snd{u(N)}}{a}\in\reals$ and
                $\neg{((1 + 1) / (N + 1))} \leq f_1(a) - \apply{\snd{u(N)}}{a}$ and
                $f_1(a) - \apply{\snd{u(N)}}{a} \leq ((1 + 1) / (N + 1))$ by \cref{intervalclosed}.
            $g(a)\in\reals$ and $\apply{\snd{u(N)}}{a}\in\reals$ by
                \cref{funs_type_apply,tietze_unit_state,continuous}.
            $N\in\reals$ by \cref{real_naturals_subset_reals,subseteq}.
            $N + 1\in\naturalsPlus$ by \cref{real_naturals_closed_succ}.
            $N + 1\in\reals$ and $1 + 1\in\reals$ by
                \cref{real_naturals_subset_reals,subseteq,real_one_in_naturals,real_add_closed}.
            $N + 1\neq 0$ by \cref{naturalsplus_positive_real,subseteq,real_pos_nonzero}.
            $((1 + 1) / (N + 1))\in\reals$ by \cref{real_div,real_mul_closed,real_mul_inverse}.
            $\neg{(((1 + 1) / (N + 1)))}\in\reals$ by \cref{real_add_inverse}.
            $\apply{\snd{u(N)}}{a} + \neg{(((1 + 1) / (N + 1)))}\in\reals$ by \cref{real_add_closed}.
            $\apply{\snd{u(N)}}{a} - ((1 + 1) / (N + 1))\in\reals$ by \cref{real_sub,real_add_eq_subst}.
            $\apply{\snd{u(N)}}{a} + ((1 + 1) / (N + 1))\in\reals$ by \cref{real_add_closed}.
            $0 < 1 + 1$ by \cref{real_zero_lt_one,real_pos_add,real_naturals_subset_reals,real_one_in_naturals,subseteq}.
            $0 \leq 1 + 1$ by \cref{real_lt_imp_leq}.
            $0 \leq ((1 + 1) / (N + 1))$ by \cref{real_div_naturalsplus_nonnegative,subseteq}.
            $\neg{(f_1(a) - \apply{\snd{u(N)}}{a})} \leq \neg{\neg{(((1 + 1) / (N + 1)))}}$ by \cref{real_leq_neg}.
            $\neg{(f_1(a) - \apply{\snd{u(N)}}{a})} \leq ((1 + 1) / (N + 1))$
                by \cref{real_neg_neg,real_leq_subst_right_local}.
            $\abs{f_1(a) - \apply{\snd{u(N)}}{a}} \leq ((1 + 1) / (N + 1))$ by \cref{real_abs_leq_if_pm_leq}.
            $\apply{\snd{u(N)}}{a} - ((1 + 1) / (N + 1)) \leq g(a)$ and
                $g(a) \leq \apply{\snd{u(N)}}{a} + ((1 + 1) / (N + 1))$ by \cref{tietze_supremum_between_endpoints}.
            $\neg{g(a)} \leq \neg{(\apply{\snd{u(N)}}{a} - ((1 + 1) / (N + 1)))}$ by \cref{real_leq_neg}.
            $\neg{g(a)}\in\reals$ and
                $\neg{(\apply{\snd{u(N)}}{a} - ((1 + 1) / (N + 1)))}\in\reals$ by \cref{real_add_inverse}.
            $\neg{g(a)} + f_1(a) \leq \neg{(\apply{\snd{u(N)}}{a} - ((1 + 1) / (N + 1)))} + f_1(a)$
                by \cref{real_leq_add}.
            $f_1(a) + \neg{g(a)} \leq f_1(a) + \neg{(\apply{\snd{u(N)}}{a} - ((1 + 1) / (N + 1)))}$
                by \cref{real_add_comm,real_leq_subst_left_local,real_leq_subst_right_local}.
            $f_1(a) \leq \apply{\snd{u(N)}}{a} + ((1 + 1) / (N + 1))$ by \cref{real_abs_sub_leq_imp_right_bound}.
            $f_1(a) + \neg{(\apply{\snd{u(N)}}{a} - ((1 + 1) / (N + 1)))}
                \leq (\apply{\snd{u(N)}}{a} + ((1 + 1) / (N + 1))) + \neg{(\apply{\snd{u(N)}}{a} - ((1 + 1) / (N + 1)))}$
                by \cref{real_leq_add}.
            $f_1(a) + \neg{(\apply{\snd{u(N)}}{a} - ((1 + 1) / (N + 1)))}\in\reals$ and
                $(\apply{\snd{u(N)}}{a} + ((1 + 1) / (N + 1))) + \neg{(\apply{\snd{u(N)}}{a} - ((1 + 1) / (N + 1)))}\in\reals$
                by \cref{real_add_closed}.
            $\epsilon \leq f_1(a) + \neg{(\apply{\snd{u(N)}}{a} - ((1 + 1) / (N + 1)))}$
                by \cref{real_sub,real_leq_subst_left_local}.
            $\epsilon \leq (\apply{\snd{u(N)}}{a} + ((1 + 1) / (N + 1))) +
                \neg{(\apply{\snd{u(N)}}{a} - ((1 + 1) / (N + 1)))}$ by \cref{real_leq_trans}.
            $\neg{(\apply{\snd{u(N)}}{a} - ((1 + 1) / (N + 1)))}\in\reals$ by \cref{real_sub,real_add_inverse,real_add_closed}.
            $\apply{\snd{u(N)}}{a} - ((1 + 1) / (N + 1)) =
                \apply{\snd{u(N)}}{a} + \neg{(((1 + 1) / (N + 1)))}$ by \cref{real_sub}.
            $\neg{(\apply{\snd{u(N)}}{a} - ((1 + 1) / (N + 1)))} =
                \neg{(\apply{\snd{u(N)}}{a} + \neg{(((1 + 1) / (N + 1)))})}$ by \cref{real_add_eq_subst}.
            $(\apply{\snd{u(N)}}{a} + ((1 + 1) / (N + 1))) +
                \neg{(\apply{\snd{u(N)}}{a} - ((1 + 1) / (N + 1)))} =
                (\apply{\snd{u(N)}}{a} + ((1 + 1) / (N + 1))) +
                \neg{(\apply{\snd{u(N)}}{a} + \neg{(((1 + 1) / (N + 1)))})}$ by \cref{real_add_eq_subst}.
            $(\apply{\snd{u(N)}}{a} + ((1 + 1) / (N + 1))) +
                \neg{(\apply{\snd{u(N)}}{a} + \neg{(((1 + 1) / (N + 1)))})}
                = ((1 + 1) / (N + 1)) + ((1 + 1) / (N + 1))$
                by \cref{real_neg_addition,real_add_assoc,real_add_comm,real_add_inverse,real_add_ident,real_add_eq_subst}.
            $(\apply{\snd{u(N)}}{a} + ((1 + 1) / (N + 1))) +
                \neg{(\apply{\snd{u(N)}}{a} - ((1 + 1) / (N + 1)))} = ((1 + 1) / (N + 1)) + ((1 + 1) / (N + 1))$
                by \cref{real_add_eq_subst}.
            $\epsilon \leq ((1 + 1) / (N + 1)) + ((1 + 1) / (N + 1))$ by \cref{real_leq_trans,real_add_eq_subst}.
            $\epsilon < \epsilon$ by \cref{real_leq_lt_trans}.
            Contradiction by \cref{real_lt_irreflexive}.
        \caseOf{$f_1(a) < g(a)$.}
            Suppose not.
            Let $\epsilon = g(a) - f_1(a)$.
            $\epsilon\in\reals$ and $0 < \epsilon$ by \cref{real_sub,real_add_inverse,real_add_closed,real_lt_imp_pos_diff}.
            Take $N\in\naturalsPlus$ such that
                $((1 + 1) / (N + 1)) + ((1 + 1) / (N + 1)) < \epsilon$ by \cref{tietze_small_double_error}.
            $\snd{u(N)}$ is a Tietze unit state on $N$ for $f_1$ over $A$ in $X$ with respect to $T$ by assumption.
            $f_1(a) - \apply{\snd{u(N)}}{a}\in
                \intervalclosed{\neg{((1 + 1) / (N + 1))}}{(1 + 1) / (N + 1)}$ by \cref{tietze_unit_state}.
            $f_1(a) - \apply{\snd{u(N)}}{a}\in\reals$ and
                $\neg{((1 + 1) / (N + 1))} \leq f_1(a) - \apply{\snd{u(N)}}{a}$ and
                $f_1(a) - \apply{\snd{u(N)}}{a} \leq ((1 + 1) / (N + 1))$ by \cref{intervalclosed}.
            $g(a)\in\reals$ and $\apply{\snd{u(N)}}{a}\in\reals$ by
                \cref{funs_type_apply,tietze_unit_state,continuous}.
            $N\in\reals$ by \cref{real_naturals_subset_reals,subseteq}.
            $N + 1\in\naturalsPlus$ by \cref{real_naturals_closed_succ}.
            $N + 1\in\reals$ and $1 + 1\in\reals$ by
                \cref{real_naturals_subset_reals,subseteq,real_one_in_naturals,real_add_closed}.
            $N + 1\neq 0$ by \cref{naturalsplus_positive_real,subseteq,real_pos_nonzero}.
            $((1 + 1) / (N + 1))\in\reals$ by \cref{real_div,real_mul_closed,real_mul_inverse}.
            $\neg{(((1 + 1) / (N + 1)))}\in\reals$ by \cref{real_add_inverse}.
            $\apply{\snd{u(N)}}{a} + \neg{(((1 + 1) / (N + 1)))}\in\reals$ and
                $\apply{\snd{u(N)}}{a} + ((1 + 1) / (N + 1))\in\reals$ by \cref{real_add_closed,real_add_inverse}.
            $0 < 1 + 1$ by \cref{real_zero_lt_one,real_pos_add,real_naturals_subset_reals,real_one_in_naturals,subseteq}.
            $0 \leq 1 + 1$ by \cref{real_lt_imp_leq}.
            $0 \leq ((1 + 1) / (N + 1))$ by \cref{real_div_naturalsplus_nonnegative,subseteq}.
            $\neg{(f_1(a) - \apply{\snd{u(N)}}{a})} \leq \neg{\neg{(((1 + 1) / (N + 1)))}}$ by \cref{real_leq_neg}.
            $\neg{(f_1(a) - \apply{\snd{u(N)}}{a})} \leq ((1 + 1) / (N + 1))$
                by \cref{real_neg_neg,real_leq_subst_right_local}.
            $\abs{f_1(a) - \apply{\snd{u(N)}}{a}} \leq ((1 + 1) / (N + 1))$ by \cref{real_abs_leq_if_pm_leq}.
            $\apply{\snd{u(N)}}{a} - ((1 + 1) / (N + 1)) \leq g(a)$ and
                $g(a) \leq \apply{\snd{u(N)}}{a} + ((1 + 1) / (N + 1))$ by \cref{tietze_supremum_between_endpoints}.
            $\apply{\snd{u(N)}}{a} - ((1 + 1) / (N + 1)) \leq f_1(a)$ by \cref{real_abs_sub_leq_imp_left_bound}.
            $\apply{\snd{u(N)}}{a} + \neg{(((1 + 1) / (N + 1)))} \leq f_1(a)$
                by \cref{real_sub,real_leq_subst_left_local}.
            $\neg{f_1(a)}\in\reals$ and
                $\neg{(\apply{\snd{u(N)}}{a} + \neg{(((1 + 1) / (N + 1)))})}\in\reals$ by \cref{real_add_inverse}.
            $\neg{f_1(a)} \leq \neg{(\apply{\snd{u(N)}}{a} + \neg{(((1 + 1) / (N + 1)))})}$ by \cref{real_leq_neg}.
            $\neg{f_1(a)} + g(a) \leq
                \neg{(\apply{\snd{u(N)}}{a} + \neg{(((1 + 1) / (N + 1)))})} + g(a)$ by \cref{real_leq_add}.
            $g(a) + \neg{f_1(a)} \leq
                g(a) + \neg{(\apply{\snd{u(N)}}{a} + \neg{(((1 + 1) / (N + 1)))})}$ by
                \cref{real_add_comm,real_leq_subst_left_local,real_leq_subst_right_local}.
            $g(a) + \neg{f_1(a)}\in\reals$ and
                $g(a) + \neg{(\apply{\snd{u(N)}}{a} + \neg{(((1 + 1) / (N + 1)))})}\in\reals$ by \cref{real_add_closed}.
            $g(a) + \neg{(\apply{\snd{u(N)}}{a} + \neg{(((1 + 1) / (N + 1)))})} \leq
                (\apply{\snd{u(N)}}{a} + ((1 + 1) / (N + 1))) + \neg{(\apply{\snd{u(N)}}{a} + \neg{(((1 + 1) / (N + 1)))})}$
                by \cref{real_leq_add}.
            $(\apply{\snd{u(N)}}{a} + ((1 + 1) / (N + 1))) + \neg{(\apply{\snd{u(N)}}{a} + \neg{(((1 + 1) / (N + 1)))})}\in\reals$
                by \cref{real_add_closed}.
            $g(a) + \neg{f_1(a)} \leq
                (\apply{\snd{u(N)}}{a} + ((1 + 1) / (N + 1))) + \neg{(\apply{\snd{u(N)}}{a} + \neg{(((1 + 1) / (N + 1)))})}$
                by \cref{real_leq_trans}.
            $\neg{(\apply{\snd{u(N)}}{a} + \neg{(((1 + 1) / (N + 1)))})}\in\reals$ by
                \cref{real_add_inverse,real_add_closed}.
            $(\apply{\snd{u(N)}}{a} + ((1 + 1) / (N + 1))) + \neg{(\apply{\snd{u(N)}}{a} + \neg{(((1 + 1) / (N + 1)))})}
                = ((1 + 1) / (N + 1)) + ((1 + 1) / (N + 1))$
                by \cref{real_neg_addition,real_add_assoc,real_add_comm,real_add_inverse,real_add_ident,real_add_eq_subst}.
            $\epsilon = g(a) + \neg{f_1(a)}$ by \cref{real_sub}.
            $\epsilon \leq
                (\apply{\snd{u(N)}}{a} + ((1 + 1) / (N + 1))) + \neg{(\apply{\snd{u(N)}}{a} + \neg{(((1 + 1) / (N + 1)))})}$
                by \cref{real_leq_subst_left_local}.
            $\epsilon \leq ((1 + 1) / (N + 1)) + ((1 + 1) / (N + 1))$ by \cref{real_leq_trans,real_add_eq_subst}.
            $\epsilon < \epsilon$ by \cref{real_leq_lt_trans}.
            Contradiction by \cref{real_lt_irreflexive}.
        \end{byCase}
    \end{subproof}
    For all $a\in A$ we have $g(a) = f(a)$ by
        \cref{subseteq,continuous,funs_type_apply,funs,id_rightunique,id_is_relation,id_dom,funs_elim,funs_ran,circ_apply,id_apply}.
    There exists $h$ such that
        $h$ is continuous from $X$ to $I$ with respect to $T$ and $T_I$ and
        for all $x\in A$ we have $h(x) = f(x)$ by \cref{subseteq}.
\end{byCase}
\end{proof}

% from §35 Item 1: Tietze extension theorem for closed intervals
\begin{theorem}\label{tietze_extension_theorem_closed_interval}
    Assume $X$ is normal with respect to $T$.
    Assume $A$ is closed in $X$ with respect to $T$.
    Assume $a\in\reals$ and $b\in\reals$ and $a < b$.
    Let $I = \intervalclosed{a}{b}$.
    Let $T_A = \subspaceTopology{T}{A}$.
    Let $T_I = \subspaceTopology{\standardTopologyR}{I}$.
    Assume $f$ is continuous from $A$ to $I$ with respect to $T_A$ and $T_I$.
    Then there exists $g$ such that
        $g$ is continuous from $X$ to $I$ with respect to $T$ and $T_I$ and
        for all $x\in A$ we have $g(x) = f(x)$.
\end{theorem}
\begin{proof}
    Let $J = \intervalclosed{\neg{1}}{1}$.
    Let $T_J = \subspaceTopology{\standardTopologyR}{J}$.
    Let $j = \identity{I}$.
    Let $f_1 = j\circ f$.
    $f_1$ is continuous from $A$ to $\reals$ with respect to $T_A$ and $\standardTopologyR$
        by \cref{continuous_interval_expand_range}.
    For all $x\in A$ we have $f_1(x) = f(x)$ by
        \cref{continuous,funs_type_apply,funs,funs_is_function,funs_is_relation,id_rightunique,id_is_relation,id_dom,circ_apply,funs_ran,id_apply,subseteq}.

    Let $c = b - a$.
    $c\in\reals$ and $0 < c$ by \cref{real_sub,real_add_inverse,real_add_closed,real_lt_imp_pos_diff}.
    Hence $c \neq 0$ and $\inv{c}\in\reals$ and $0 < \inv{c}$ by \cref{real_pos_nonzero,real_mul_inverse,real_inv_pos}.

    $T$ is a topology on $X$ by \cref{normal_space}.
    Hence $A\subseteq X$ by \cref{closed_in}.
    $T_A$ is a topology on $A$ by \cref{closed_in,subspace_topology_is_topology,subspace_of}.
    $\standardTopologyR$ is a topology on $\reals$
        by \cref{standard_basis_is_basis,basis_topology_is_topology,standard_topology_reals}.

    Let $k_a = \{(x,a) \mid x\in A\}$.
    Let $k_c = \{(x,c) \mid x\in A\}$.
    $k_a$ is a function from $A$ to $\reals$ and $k_c$ is a function from $A$ to $\reals$
        by \cref{constant_map_function}.
    For all $x\in A$ we have $k_a(x) = a$ and $k_c(x) = c$
        by \cref{constant_map_function,funs_is_function,function_apply_intro}.
    $k_a$ is continuous from $A$ to $\reals$ with respect to $T_A$ and $\standardTopologyR$ and
        $k_c$ is continuous from $A$ to $\reals$ with respect to $T_A$ and $\standardTopologyR$
        by \cref{constant_continuous}.

    Let $s = \{(x, f_1(x) - k_a(x)) \mid x\in A\}$.
    $s$ is continuous from $A$ to $\reals$ with respect to $T_A$ and $\standardTopologyR$
        by \cref{real_function_subtraction_continuous}.
    Let $u = \{(x, s(x) \rmul \inv{(k_c(x))}) \mid x\in A\}$.
    For all $x\in A$ we have $k_c(x)\neq 0$ by \cref{subseteq,real_lt_subst_left}.
    Hence $u$ is continuous from $A$ to $\reals$ with respect to $T_A$ and $\standardTopologyR$
        by \cref{real_function_division_continuous}.
    Let $r = \{(x, u(x) + u(x)) \mid x\in A\}$.
    $r$ is continuous from $A$ to $\reals$ with respect to $T_A$ and $\standardTopologyR$
        by \cref{real_function_addition_continuous}.
    $1\in\reals$ by \cref{real_naturals_subset_reals,real_one_in_naturals,subseteq}.
    Let $k_1 = \{(x,1) \mid x\in A\}$.
    $k_1$ is a function from $A$ to $\reals$ by \cref{constant_map_function}.
    For all $x\in A$ we have $k_1(x) = 1$
        by \cref{constant_map_function,funs_is_function,function_apply_intro}.
    $k_1$ is continuous from $A$ to $\reals$ with respect to $T_A$ and $\standardTopologyR$
        by \cref{constant_continuous}.
    Let $h_0 = \{(x, r(x) - k_1(x)) \mid x\in A\}$.
    $h_0$ is continuous from $A$ to $\reals$ with respect to $T_A$ and $\standardTopologyR$
        by \cref{real_function_subtraction_continuous}.
    We have $h_0\in\funs{A}{\reals}$ and $h_0$ is a function and
        $u\in\funs{A}{\reals}$ and $u$ is a function and
        $s\in\funs{A}{\reals}$ and $s$ is a function and
        $r\in\funs{A}{\reals}$ and $r$ is a function
        by \cref{continuous,funs_is_function}.

    Show that for all $x\in A$ we have $h_0(x)\in J$.
    \begin{subproof}
        Fix $x\in A$.
        We have $f_1(x) = f(x)$ by \cref{subseteq}.
        Hence $f_1(x)\in I$ by \cref{continuous,funs_type_apply,real_add_eq_subst}.
        We have $f_1(x)\in\reals$ and $a \leq f_1(x)$ and $f_1(x) \leq b$ by \cref{intervalclosed}.
        $k_a(x) = a$ and $k_c(x) = c$ and $k_1(x) = 1$ by \cref{subseteq}.
        $s(x) = f_1(x) - a$ and $s(x)\in\reals$ by \cref{function_apply_intro,continuous,funs_type_apply}.
        $\neg{a}\in\reals$ by \cref{real_add_inverse}.
        We have $a + \neg{a} \leq f_1(x) + \neg{a}$ by \cref{real_leq_add}.
        Hence $0 \leq s(x)$ by
            \cref{real_add_inverse,real_sub,real_leq_subst_left_local,real_leq_subst_right_local}.
        We have $f_1(x) + \neg{a} \leq b + \neg{a}$ by \cref{real_leq_add}.
        Hence $s(x) \leq c$ by \cref{real_sub,real_leq_subst_left_local,real_leq_subst_right_local}.

        We have $u(x) = s(x) \rmul \inv{c}$ and $u(x)\in\reals$ by \cref{function_apply_intro,continuous,funs_type_apply}.
        $0\in\reals$ by \cref{real_add_ident}.
        We have $0 < s(x)$ or $s(x) = 0$ by \cref{real_leq_split}.
        Therefore $0 \leq u(x)$ by
            \cref{real_lt_mul_pos,real_mul_zero,real_lt_subst_left,real_lt_subst_right,real_lt_imp_leq,real_mul_eq_subst,subseteq}.

        We have $s(x) < c$ or $s(x) = c$ by \cref{real_leq_split}.
        Hence $u(x) \leq 1$ by
            \cref{real_lt_mul_pos,real_lt_subst_left,real_lt_subst_right,real_mul_inverse,real_lt_imp_leq,real_mul_eq_subst,subseteq}.

        We have $r(x) = u(x) + u(x)$ and $r(x)\in\reals$ by \cref{function_apply_intro,continuous,funs_type_apply}.
        We have $u(x) = 0$ or $0 < u(x)$ by \cref{real_leq_split}.
        Therefore $0 \leq r(x)$ by
            \cref{real_add_eq_subst,real_add_ident,real_pos_add,real_lt_subst_right,real_lt_imp_leq,subseteq}.

        We have $u(x) + u(x) \leq 1 + u(x)$ by \cref{real_leq_add}.
        $1 + u(x) \leq 1 + 1$ by \cref{real_leq_add,real_add_comm,real_leq_subst_left_local}.
        Hence $r(x) \leq 1 + u(x)$ by \cref{real_leq_subst_left_local}.
        We have $1 + u(x)\in\reals$ and $1 + 1\in\reals$ by \cref{real_add_closed}.
        Hence $r(x) \leq 1 + 1$ by \cref{real_leq_trans}.

        We have $h_0(x) = r(x) - 1$ and $h_0(x)\in\reals$ by \cref{function_apply_intro,continuous,funs_type_apply}.
        $\neg{1}\in\reals$ by \cref{real_add_inverse}.
        We have $0 + \neg{1} \leq r(x) + \neg{1}$ by \cref{real_leq_add}.
        Hence $\neg{1} \leq h_0(x)$ by
            \cref{real_add_ident,real_sub,real_leq_subst_left_local,real_leq_subst_right_local}.
        We have $r(x) + \neg{1} \leq (1 + 1) + \neg{1}$ by \cref{real_leq_add}.
        Hence $h_0(x) \leq 1$ by
            \cref{real_add_assoc,real_add_inverse,real_add_comm,real_add_ident,real_sub,real_leq_subst_left_local,real_leq_subst_right_local}.
        Therefore $h_0(x)\in J$ by \cref{intervalclosed}.
    \end{subproof}
    $h_0$ is continuous from $A$ to $J$ with respect to $T_A$ and $T_J$
        by \cref{img_iff,continuous,funs,funs_is_function,function_apply_iff,pair_eq_iff,subseteq,continuous_subspace_range,intervalclosed}.

    Take $e$ such that
        $e$ is continuous from $X$ to $J$ with respect to $T$ and $T_J$ and
        for all $x\in A$ we have $e(x) = h_0(x)$
        by \cref{tietze_extension_theorem_unit_interval}.

    Let $q = \{(x,a) \mid x\in X\}$.
    Let $v = \{(x,c) \mid x\in X\}$.
    Let $w = \{(x,1) \mid x\in X\}$.
    Let $z = \{(x,1 + 1) \mid x\in X\}$.
    $1\in\reals$ and $1 + 1\in\reals$ by
        \cref{real_naturals_subset_reals,real_one_in_naturals,subseteq,real_add_closed}.
    $q$ is a function from $X$ to $\reals$ and
        $v$ is a function from $X$ to $\reals$ and
        $w$ is a function from $X$ to $\reals$ and
        $z$ is a function from $X$ to $\reals$
        by \cref{constant_map_function}.
    For all $x\in X$ we have
        $q(x) = a$ and $v(x) = c$ and $w(x) = 1$ and $z(x) = 1 + 1$
        by \cref{constant_map_function,funs_is_function,function_apply_intro}.
    $q$ is continuous from $X$ to $\reals$ with respect to $T$ and $\standardTopologyR$ and
        $v$ is continuous from $X$ to $\reals$ with respect to $T$ and $\standardTopologyR$ and
        $w$ is continuous from $X$ to $\reals$ with respect to $T$ and $\standardTopologyR$ and
        $z$ is continuous from $X$ to $\reals$ with respect to $T$ and $\standardTopologyR$
        by \cref{constant_continuous}.

    Let $j_J = \identity{J}$.
    Let $e_1 = j_J\circ e$.
    $e_1$ is continuous from $X$ to $\reals$ with respect to $T$ and $\standardTopologyR$
        by \cref{continuous_interval_expand_range}.

    Let $p = \{(x, e_1(x) + w(x)) \mid x\in X\}$.
    $p$ is continuous from $X$ to $\reals$ with respect to $T$ and $\standardTopologyR$
        by \cref{real_function_addition_continuous}.
    For all $x\in X$ we have $z(x)\neq 0$ by
        \cref{subseteq,real_one_in_naturals,real_naturals_subset_reals,naturalsplus_positive_real,real_pos_add,real_pos_nonzero}.
    Let $t = \{(x, p(x) \rmul \inv{(z(x))}) \mid x\in X\}$.
    $t$ is continuous from $X$ to $\reals$ with respect to $T$ and $\standardTopologyR$
        by \cref{real_function_division_continuous}.
    Let $m = \{(x, v(x) \rmul t(x)) \mid x\in X\}$.
    $m$ is continuous from $X$ to $\reals$ with respect to $T$ and $\standardTopologyR$
        by \cref{real_function_multiplication_continuous}.
    Let $g_0 = \{(x, q(x) + m(x)) \mid x\in X\}$.
    $g_0$ is continuous from $X$ to $\reals$ with respect to $T$ and $\standardTopologyR$
        by \cref{real_function_addition_continuous}.
    We have $g_0\in\funs{X}{\reals}$ and $g_0$ is a function and
        $p\in\funs{X}{\reals}$ and $p$ is a function and
        $t\in\funs{X}{\reals}$ and $t$ is a function and
        $m\in\funs{X}{\reals}$ and $m$ is a function and
        $e_1\in\funs{X}{\reals}$ and $e_1$ is a function
        by \cref{continuous,funs_is_function}.

    Show that for all $x\in X$ we have $g_0(x)\in I$.
    \begin{subproof}
        Fix $x\in X$.
        We have $e(x)\in J$ by \cref{continuous,funs_type_apply}.
        Hence $e_1(x)\in J$ by
            \cref{continuous,funs,funs_is_function,funs_is_relation,id_rightunique,id_is_relation,id_dom,circ_apply,funs_ran,id_apply,subseteq,real_add_eq_subst}.
        We have $e_1(x)\in\reals$ and $\neg{1} \leq e_1(x)$ and $e_1(x) \leq 1$ by \cref{intervalclosed}.
        $w(x) = 1$ and $z(x) = 1 + 1$ and $v(x) = c$ and $q(x) = a$ by \cref{subseteq}.
        $p(x) = e_1(x) + 1$ and $p(x)\in\reals$ by \cref{function_apply_intro,continuous,funs_type_apply}.
        $1\in\reals$ by \cref{real_naturals_subset_reals,real_one_in_naturals,subseteq}.
        We have $\neg{1} + 1 \leq e_1(x) + 1$ by \cref{real_add_inverse,real_leq_add}.
        Hence $0 \leq p(x)$ by \cref{real_add_inverse,real_add_comm,real_leq_subst_left_local}.
        We have $p(x) \leq 1 + 1$ by \cref{real_leq_add,real_leq_subst_left_local}.

        $t(x) = p(x) \rmul \inv{(1 + 1)}$ by \cref{function_apply_intro}.
        Hence $1 + 1\in\naturalsPlus$ by \cref{real_one_in_naturals,real_naturals_closed_succ}.
        We have $p(x) / (1 + 1)\in\reals$ and $0 \leq p(x) / (1 + 1)$
            by \cref{real_div_naturalsplus_real,real_div_naturalsplus_nonnegative}.
        Hence $t(x)\in\reals$ and $0 \leq t(x)$ by \cref{real_div,real_lt_subst_left,real_lt_subst_right}.
        We have $p(x) / (1 + 1) \leq (1 + 1) / (1 + 1)$ by \cref{real_div_leq_same_denom}.
        Hence $t(x) \leq 1$ by
            \cref{real_div,real_two_over_two_local,real_leq_subst_left_local,real_leq_subst_right_local}.

        We have $m(x) = c \rmul t(x)$ and $m(x)\in\reals$ by \cref{function_apply_intro,continuous,funs_type_apply}.
        $0\in\reals$ by \cref{real_add_ident}.
        We have $t(x) = 0$ or $0 < t(x)$ by \cref{real_leq_split}.
        Hence $0 \leq m(x)$ by
            \cref{real_mul_eq_subst,real_mul_comm,real_mul_zero,subseteq,real_lt_mul_pos,real_lt_subst_left,real_lt_imp_leq}.

        We have $t(x) < 1$ or $t(x) = 1$ by \cref{real_leq_split}.
        Therefore $m(x) \leq c$ by
            \cref{real_lt_mul_pos,real_mul_comm,real_lt_subst_left,real_lt_subst_right,real_mul_ident,real_lt_imp_leq,real_mul_eq_subst,subseteq}.

        We have $g_0(x) = a + m(x)$ and $a + m(x)\in\reals$ by \cref{function_apply_intro,real_add_closed}.
        $0\in\reals$ by \cref{real_add_ident}.
        Hence $a \leq a + m(x)$ by
            \cref{real_leq_add,real_add_ident,real_add_comm,real_leq_subst_left_local,real_leq_subst_right_local}.
        We have $a \leq g_0(x)$ by \cref{real_leq_subst_right_local}.
        We have $a + m(x) \leq a + c$ by
            \cref{real_leq_add,real_add_comm,real_leq_subst_left_local,real_leq_subst_right_local}.
        We have $a + (b - a) = b$ by
            \cref{real_sub,real_add_assoc,real_add_comm,real_add_inverse,real_add_ident}.
        Hence $g_0(x) \leq b$ by \cref{real_sub,real_leq_subst_left_local,real_leq_subst_right_local}.
        Therefore $g_0(x)\in I$ by \cref{intervalclosed}.
    \end{subproof}
    $g_0$ is continuous from $X$ to $I$ with respect to $T$ and $T_I$
        by \cref{img_iff,continuous,funs,funs_is_function,function_apply_iff,pair_eq_iff,subseteq,continuous_subspace_range,intervalclosed}.

    Show that for all $x\in A$ we have $g_0(x) = f(x)$.
    \begin{subproof}
        Fix $x\in A$.
        We have $x\in X$ by \cref{subseteq}.
        $e(x) = h_0(x)$ by assumption.
        Hence $e_1(x) = h_0(x)$ by
            \cref{continuous,funs,funs_is_function,funs_is_relation,id_rightunique,id_is_relation,id_dom,circ_apply,funs_ran,id_apply,subseteq,closed_in}.
        $w(x) = 1$ and $z(x) = 1 + 1$ and $v(x) = c$ and $q(x) = a$
            by \cref{subseteq,closed_in}.
        We have $p\in\funs{X}{\reals}$ and $p$ is a function and $\dom{p} = X$ and
            $t\in\funs{X}{\reals}$ and $t$ is a function and $\dom{t} = X$ and
            $m\in\funs{X}{\reals}$ and $m$ is a function and $\dom{m} = X$
            by \cref{continuous,funs_is_function,funs}.
        $p(x) = e_1(x) + 1$ and $h_0(x) = r(x) - 1$ by \cref{function_apply_intro}.
        $r(x)\in\reals$ and $1\in\reals$ by
            \cref{continuous,funs_type_apply,real_naturals_subset_reals,real_one_in_naturals,subseteq}.
        Hence $p(x) = r(x)$ by
            \cref{real_add_eq_subst,real_sub,real_add_inverse,real_add_closed,real_add_assoc,real_add_comm,real_add_ident,subseteq}.
        $p(x) = u(x) + u(x)$ by \cref{function_apply_intro,real_add_eq_subst}.
        $u(x)\in\reals$ by \cref{continuous,funs_type_apply}.
        Hence $(1 + 1) \rmul u(x) = u(x) + u(x)$ by
            \cref{real_mul_ident,real_right_distrib,real_add_eq_subst}.
        We have $p(x) = (1 + 1) \rmul u(x)$ by \cref{real_add_eq_subst}.
        Hence $t(x) = u(x)$ by
            \cref{function_apply_intro,real_div,real_mul_eq_subst,real_mul_assoc,real_mul_comm,real_mul_inverse,real_mul_ident}.
        We have $s(x)\in\reals$ by \cref{continuous,funs_type_apply}.
        Hence $m(x) = c \rmul (s(x) / c)$ by \cref{real_div,real_mul_eq_subst,function_apply_intro}.
        We have $m(x) = s(x)$ by \cref{real_div,real_mul_assoc,real_mul_comm,real_mul_inverse,real_mul_ident}.
        We have $s(x) = f_1(x) - a$ by \cref{function_apply_intro}.
        Hence $g_0(x) = f_1(x)$ by
            \cref{function_apply_intro,continuous,funs_type_apply,real_add_eq_subst,real_sub,real_add_assoc,real_add_comm,real_add_inverse,real_add_ident}.
        Therefore $g_0(x) = f(x)$ by \cref{subseteq}.
    \end{subproof}
    Therefore there exists $g$ such that
        $g$ is continuous from $X$ to $I$ with respect to $T$ and $T_I$ and
        for all $x\in A$ we have $g(x) = f(x)$ by \cref{subseteq}.
\end{proof}

% from §35 helper lemma 1: extension theorem for maps into the open unit interval
\begin{lemma}\label{tietze_extension_theorem_open_interval}
    Assume $X$ is normal with respect to $T$.
    Assume $A$ is closed in $X$ with respect to $T$.
    Let $I = \intervalopen{\neg{1}}{1}$.
    Let $J = \intervalclosed{\neg{1}}{1}$.
    Let $T_A = \subspaceTopology{T}{A}$.
    Let $T_I = \subspaceTopology{\standardTopologyR}{I}$.
    Let $T_J = \subspaceTopology{\standardTopologyR}{J}$.
    Assume $f$ is continuous from $A$ to $I$ with respect to $T_A$ and $T_I$.
    Then there exists $g$ such that
        $g$ is continuous from $X$ to $I$ with respect to $T$ and $T_I$ and
        for all $x\in A$ we have $g(x) = f(x)$.
\end{lemma}
\begin{proof}
    We have $I\subseteq J$ by \cref{intervalopen,real_lt_imp_leq,intervalclosed,subseteq}.

    Hence $\img{f}{A}\subseteq J$ by
        \cref{img_iff,continuous,funs,funs_is_function,function_apply_iff,funs_type_apply,subseteq,pair_eq_iff}.
    We have $I\subseteq\reals$ and $J\subseteq\reals$ by \cref{intervalopen,intervalclosed,subseteq}.
    $\standardTopologyR$ is a topology on $\reals$
        by \cref{standard_basis_is_basis,basis_topology_is_topology,standard_topology_reals}.
    Hence $I$ is a subspace of $\reals$ with respect to $\standardTopologyR$ by \cref{subspace_of}.
    Let $j_I = \identity{I}$.
    $j_I$ is continuous from $I$ to $\reals$ with respect to $T_I$ and $\standardTopologyR$
        by \cref{inclusion_continuous}.
    Let $f_1 = j_I\circ f$.
    $f_1$ is continuous from $A$ to $\reals$ with respect to $T_A$ and $\standardTopologyR$
        by \cref{continuous_composite}.
    We have $f_1\in\funs{A}{\reals}$ and $T_A$ is a topology on $A$ by \cref{continuous}.
    Show that for all $y$ such that $y\in\img{f_1}{A}$ we have $y\in J$.
    \begin{subproof}
        Follows by \cref{img_iff,continuous,funs,funs_is_function,function_apply_iff,funs_type_apply,function_rightunique,funs_is_relation,id_rightunique,id_is_relation,id_dom,circ_apply,funs_ran,id_apply,real_add_eq_subst,subseteq,pair_eq_iff}.
    \end{subproof}
    $f_1$ is continuous from $A$ to $J$ with respect to $T_A$ and $T_J$
        by \cref{img_iff,continuous,funs,funs_is_function,function_apply_iff,funs_type_apply,function_rightunique,funs_is_relation,id_rightunique,id_is_relation,id_dom,circ_apply,funs_ran,id_apply,real_add_eq_subst,subseteq,pair_eq_iff,continuous_subspace_range}.
    $1\in\reals$ by \cref{real_naturals_subset_reals,real_one_in_naturals,subseteq}.
    $0\in\reals$ and $0 < 1$ by \cref{real_add_ident,real_zero_lt_one}.
    Hence $\neg{1}\in\reals$ by \cref{real_add_inverse}.
    We have $\neg{1} < 0$ by \cref{real_lt_neg,real_neg_zero,real_lt_subst_right}.
    Therefore $\neg{1} < 1$ by \cref{real_lt_trans,real_zero_lt_one}.

    Take $g_0$ such that
        $g_0$ is continuous from $X$ to $J$ with respect to $T$ and $T_J$ and
        for all $x\in A$ we have $g_0(x) = f_1(x)$
        by \cref{tietze_extension_theorem_closed_interval}.

    Let $j_J = \identity{J}$.
    Let $g_1 = j_J\circ g_0$.
    $g_1$ is continuous from $X$ to $\reals$ with respect to $T$ and $\standardTopologyR$
        by \cref{continuous_interval_expand_range}.

    $\standardMetricR$ is a metric on $\reals$ by \cref{standard_metric_is_metric}.
    Hence $\reals$ is Hausdorff with respect to $\metricTopology{\standardMetricR}{\reals}$
        by \cref{metric_space_hausdorff}.
    Therefore $\reals$ is Hausdorff with respect to $\standardTopologyR$
        by \cref{standard_metric_topology}.
    $1\in\reals$ by \cref{real_naturals_subset_reals,real_one_in_naturals,subseteq}.
    Hence $\neg{1}\in\reals$ by \cref{real_add_inverse}.
    $\{\neg{1}\}\subseteq\reals$ and $\{1\}\subseteq\reals$ by \cref{singleton_subset_intro}.
    $\{\neg{1}\}$ is closed in $\reals$ with respect to $\standardTopologyR$ and
        $\{1\}$ is closed in $\reals$ with respect to $\standardTopologyR$
        by \cref{singleton_subset_intro,pair_is_finite,upair_intro_left,subseteq_finite_is_finite,finite_point_set_closed_hausdorff}.

    Let $D_1 = \preimg{g_1}{\{\neg{1}\}}$.
    Let $D_2 = \preimg{g_1}{\{1\}}$.
    Let $D = D_1\union D_2$.
    Let $U_1 = \reals\setminus \{\neg{1}\}$.
    Let $U_2 = \reals\setminus \{1\}$.
    Hence $\preimg{g_1}{U_1}$ is open in $X$ with respect to $T$ and
        $\preimg{g_1}{U_2}$ is open in $X$ with respect to $T$ by \cref{closed_in,continuous}.
    $T$ is a topology on $X$ by \cref{normal_space}.
    $g_1\in\funs{X}{\reals}$ and $g_1$ is a function by \cref{continuous,funs_is_function}.
    Hence $\preimg{g_1}{U_1} = X\setminus D_1$ and $\preimg{g_1}{U_2} = X\setminus D_2$
        by \cref{preimg_setminus_function,subseteq}.
    Therefore $D_1$ is closed in $X$ with respect to $T$ and $D_2$ is closed in $X$ with respect to $T$
        by \cref{continuous,preimg_subseteq_dom,funs,closed_in}.
    Hence $(X\setminus D_1)\in T$ and $(X\setminus D_2)\in T$ by \cref{closed_in,open_in,subseteq}.
    We have $(X\setminus D_1)\inter (X\setminus D_2)$ is open in $X$ with respect to $T$ by \cref{topology_on,open_in,subseteq}.
    $(X\setminus D_1)\inter (X\setminus D_2) = X\setminus (D_1\union D_2)$ by \cref{setminus_union}.
    Hence $X\setminus D$ is open in $X$ with respect to $T$ by \cref{real_add_eq_subst}.
    Thus $D$ is closed in $X$ with respect to $T$ by \cref{union_subsets_is_subset,closed_in}.

    Show that $A$ is disjoint from $D$.
    \begin{subproof}
        Suppose not.
        Take $a$ such that $a\in A$ and $a\in D$ by \cref{disjoint}.
        We have $a\in D_1$ or $a\in D_2$ by \cref{union_iff}.
        We have $a\in X$ and $g_0(a)\in J$ and $f(a)\in I$
            by \cref{closed_in,subseteq,continuous,funs_type_apply}.
        We have $g_0(a) = f_1(a)$ by assumption.
        Hence $g_1(a) = g_0(a)$ and $f_1(a) = f(a)$ by
            \cref{continuous,funs,funs_is_function,funs_is_relation,id_rightunique,id_is_relation,id_dom,circ_apply,funs_ran,id_apply,subseteq}.
        $g_1(a) = f(a)$ by \cref{subseteq}.
        $f(a)\in\reals$ and $\neg{1} < f(a)$ and $f(a) < 1$ by \cref{intervalopen}.
        Hence $a\notin D_1$ and $a\notin D_2$ by
            \cref{preimg_iff,singleton_elim,continuous,funs_is_function,function_apply_iff,subseteq,real_lt_subst_right,real_lt_subst_left,real_lt_irreflexive}.
        Contradiction by \cref{union_iff}.
    \end{subproof}

    Let $I_0 = \intervalclosed{0}{1}$.
    Let $T_0 = \subspaceTopology{\standardTopologyR}{I_0}$.
    $D$ is disjoint from $A$ by \cref{disjoint_symmetric}.
    Take $\phi$ such that
        $\phi$ is continuous from $X$ to $I_0$ with respect to $T$ and $T_0$ and
        $\img{\phi}{D}\subseteq \{0\}$ and
        $\img{\phi}{A}\subseteq \{1\}$
        by \cref{urysohn_unit_interval}.
    Let $j_0 = \identity{I_0}$.
    Let $\phi_1 = j_0\circ \phi$.
    $\phi_1$ is continuous from $X$ to $\reals$ with respect to $T$ and $\standardTopologyR$
        by \cref{continuous_interval_expand_range}.
    Let $h_0 = \{(x, \phi_1(x) \rmul g_1(x)) \mid x\in X\}$.
    $h_0$ is continuous from $X$ to $\reals$ with respect to $T$ and $\standardTopologyR$
        by \cref{real_function_multiplication_continuous}.
    We have $h_0\in\funs{X}{\reals}$ and $h_0$ is a function and
        $\phi_1\in\funs{X}{\reals}$ and $\phi_1$ is a function
        by \cref{continuous,funs_is_function}.

    Show that for all $x\in X$ we have $h_0(x)\in I$.
    \begin{subproof}
        Fix $x\in X$.
        We have $g_0(x)\in J$ by \cref{continuous,funs_type_apply}.
        Hence $g_1(x)\in J$ by
            \cref{continuous,funs,funs_is_function,funs_is_relation,id_rightunique,id_is_relation,id_dom,circ_apply,funs_ran,id_apply,subseteq,real_add_eq_subst}.
        We have $g_1(x)\in\reals$ and $\neg{1} \leq g_1(x)$ and $g_1(x) \leq 1$ by \cref{intervalclosed}.

        We have $\phi(x)\in I_0$ by \cref{continuous,funs_type_apply}.
        $\phi\in\funs{X}{I_0}$ and $\phi$ is right-unique and $\phi$ is a relation and $\dom{\phi} = X$
            by \cref{continuous,funs,funs_is_function,funs_is_relation}.
        $j_0$ is right-unique and $j_0$ is a relation and $\dom{j_0} = I_0$
            by \cref{id_rightunique,id_is_relation,id_dom}.
        $\phi_1(x) = \phi(x)$ and $\phi_1(x)\in I_0$
            by \cref{circ_apply,funs_ran,subseteq,id_apply,real_add_eq_subst}.
        $\phi_1(x)\in\reals$ and $0 \leq \phi_1(x)$ and $\phi_1(x) \leq 1$ by \cref{intervalclosed}.

        We have $h_0(x) = \phi_1(x) \rmul g_1(x)$ and $h_0(x)\in\reals$ by \cref{function_apply_intro,continuous,funs_type_apply}.

        \begin{byCase}
        \caseOf{$x\in D$.}
            We have $\phi(x) = 0$ by \cref{subseteq,inter_intro,singleton_intro,img_of_function_intro,singleton_elim}.
            Hence $h_0(x) = 0$ by \cref{subseteq,real_mul_eq_subst,real_mul_zero}.
            Therefore $h_0(x)\in I$ by
                \cref{intervalopen,real_add_inverse,real_zero_lt_one,real_lt_neg,real_neg_zero,real_lt_subst_right,real_lt_subst_left}.
        \caseOf{$x\notin D$.}
            We have $x\notin D_1$ and $x\notin D_2$ by \cref{union_iff}.
            We have $g_1\in\funs{X}{\reals}$ and $g_1$ is a function and $\dom{g_1} = X$
                by \cref{continuous,funs_is_function,funs}.
            Hence $g_1(x)\neq \neg{1}$ by \cref{continuous,funs_is_function,function_apply_iff,singleton_intro,preimg_iff}.
            We have $g_1(x)\neq 1$ by \cref{continuous,funs_is_function,function_apply_iff,singleton_intro,preimg_iff}.
            We have $g_1(x) < 0$ or $g_1(x) = 0$ or $0 < g_1(x)$ by \cref{real_lt_trichotomy}.
            \begin{byCase}
            \caseOf{$g_1(x) = 0$.}
                We have $h_0(x) = 0$ by \cref{real_mul_eq_subst,real_mul_comm,real_mul_zero}.
                Therefore $h_0(x)\in I$ by
                    \cref{intervalopen,real_add_inverse,real_zero_lt_one,real_lt_neg,real_neg_zero,real_lt_subst_right,real_lt_subst_left}.
            \caseOf{$0 < g_1(x)$.}
                We have $\phi_1(x) = 0$ or $0 < \phi_1(x)$ by \cref{real_leq_split}.
                \begin{byCase}
                \caseOf{$\phi_1(x) = 0$.}
                    We have $h_0(x) = 0$ by \cref{real_mul_eq_subst,real_mul_zero}.
                    Therefore $h_0(x)\in I$ by
                        \cref{intervalopen,real_add_inverse,real_zero_lt_one,real_lt_neg,real_neg_zero,real_lt_subst_right,real_lt_subst_left}.
                \caseOf{$0 < \phi_1(x)$.}
                    $\phi_1(x) < 1$ or $\phi_1(x) = 1$ by \cref{real_leq_split}.
                    \begin{byCase}
                    \caseOf{$\phi_1(x) < 1$.}
                        We have $\phi_1(x) \rmul g_1(x) < 1 \rmul g_1(x)$ by \cref{real_lt_mul_pos}.
                        We have $h_0(x) < g_1(x)$ by \cref{real_mul_ident,real_lt_subst_right,real_lt_subst_left}.
                        Hence $h_0(x) < 1$ by \cref{real_lt_trans}.
                        We have $0 \rmul g_1(x) < \phi_1(x) \rmul g_1(x)$ by \cref{real_lt_mul_pos}.
                        Hence $0 < h_0(x)$ by \cref{real_mul_zero,real_lt_subst_left}.
                        Thus $\neg{1} < h_0(x)$ by
                            \cref{real_add_inverse,real_zero_lt_one,real_lt_neg,real_neg_zero,real_lt_subst_right,real_lt_trans}.
                        Therefore $h_0(x)\in I$ by \cref{intervalopen}.
                    \caseOf{$\phi_1(x) = 1$.}
                        We have $h_0(x) = g_1(x)$ by \cref{real_mul_eq_subst,real_mul_ident}.
                        We have $h_0(x) < 1$ by \cref{real_lt_subst_left}.
                        Hence $0 \rmul g_1(x) < 1 \rmul g_1(x)$ by \cref{real_lt_mul_pos}.
                        We have $0 < h_0(x)$ by \cref{real_mul_zero,real_mul_ident,real_lt_subst_left,real_lt_subst_right}.
                        Thus $\neg{1} < h_0(x)$ by
                            \cref{real_add_inverse,real_zero_lt_one,real_lt_neg,real_neg_zero,real_lt_subst_right,real_lt_trans}.
                        Therefore $h_0(x)\in I$ by \cref{intervalopen}.
                    \end{byCase}
                \end{byCase}
            \caseOf{$g_1(x) < 0$.}
                Let $u = \neg{(g_1(x))}$.
                Hence $u\in\reals$ and $0 < u$ by \cref{real_add_inverse,real_lt_neg,real_neg_zero,real_lt_subst_left}.
                We have $g_1(x) = \neg{u}$ by \cref{real_neg_neg}.
                We have $\phi_1(x) = 0$ or $0 < \phi_1(x)$ by \cref{real_leq_split}.
                \begin{byCase}
                \caseOf{$\phi_1(x) = 0$.}
                    We have $h_0(x) = 0$ by \cref{real_mul_eq_subst,real_mul_zero}.
                    Therefore $h_0(x)\in I$ by
                        \cref{intervalopen,real_add_inverse,real_zero_lt_one,real_lt_neg,real_neg_zero,real_lt_subst_right,real_lt_subst_left}.
                \caseOf{$0 < \phi_1(x)$.}
                    $\phi_1(x) < 1$ or $\phi_1(x) = 1$ by \cref{real_leq_split}.
                    \begin{byCase}
                    \caseOf{$\phi_1(x) < 1$.}
                        We have $\phi_1(x) \rmul u < 1 \rmul u$ by \cref{real_lt_mul_pos}.
                        We have $\phi_1(x) \rmul u < u$ by \cref{real_mul_ident,real_lt_subst_right}.
                        We have $\phi_1(x) \rmul u\in\reals$ by \cref{real_mul_closed}.
                        Hence $\neg{u} < \neg{(\phi_1(x) \rmul u)}$ by \cref{real_lt_neg}.
                        We have $g_1(x) < \neg{(\phi_1(x) \rmul u)}$ by \cref{real_lt_subst_left}.
                        Hence $\phi_1(x) \rmul \neg{u} = \neg{(\phi_1(x) \rmul u)}$ by \cref{real_mul_neg_right}.
                        Thus $g_1(x) < \phi_1(x) \rmul g_1(x)$ by \cref{real_mul_eq_subst,real_lt_subst_right}.
                        Hence $\neg{1} < h_0(x)$ by \cref{real_lt_subst_left,real_lt_trans}.
                        We have $\phi_1(x) \rmul u > 0 \rmul u$ by \cref{real_lt_mul_pos}.
                        Hence $\phi_1(x) \rmul u > 0$ by \cref{real_mul_zero,real_lt_subst_right}.
                        Hence $\neg{(\phi_1(x) \rmul u)} < \neg{0}$ by \cref{real_lt_neg}.
                        We have $\neg{(\phi_1(x) \rmul u)} < 0$ by \cref{real_neg_zero,real_lt_subst_right}.
                        Hence $\phi_1(x) \rmul \neg{u} < 0$ by \cref{real_mul_neg_right,real_lt_subst_left}.
                        Hence $h_0(x) < 0$ by \cref{real_mul_eq_subst,real_lt_subst_left}.
                        Hence $h_0(x) < 1$ by \cref{real_lt_trans,real_zero_lt_one}.
                        Therefore $h_0(x)\in I$ by \cref{intervalopen}.
                    \caseOf{$\phi_1(x) = 1$.}
                        We have $h_0(x) = g_1(x)$ by \cref{real_mul_eq_subst,real_mul_ident}.
                        We have $\neg{1} < h_0(x)$ by \cref{real_lt_subst_left}.
                        Hence $1 \rmul u > 0 \rmul u$ by \cref{real_lt_mul_pos}.
                        We have $u > 0$ by \cref{real_mul_zero,real_mul_ident,real_lt_subst_left,real_lt_subst_right}.
                        Hence $\neg{u} < 0$ by \cref{real_lt_neg,real_neg_zero,real_lt_subst_right}.
                        We have $h_0(x) < 0$ by \cref{real_lt_subst_left}.
                        Hence $h_0(x) < 1$ by \cref{real_lt_trans,real_zero_lt_one}.
                        Therefore $h_0(x)\in I$ by \cref{intervalopen}.
                    \end{byCase}
                \end{byCase}
            \end{byCase}
        \end{byCase}
    \end{subproof}
    $h_0$ is continuous from $X$ to $I$ with respect to $T$ and $T_I$
        by \cref{img_iff,continuous,funs,funs_is_function,function_apply_iff,pair_eq_iff,subseteq,continuous_subspace_range,intervalopen}.

    Show that for all $a\in A$ we have $h_0(a) = f(a)$.
    \begin{subproof}
        Follows by \cref{closed_in,subseteq,continuous,funs_is_function,funs,inter_intro,singleton_intro,img_of_function_intro,singleton_elim,funs_type_apply,id_rightunique,id_is_relation,id_dom,circ_apply,funs_ran,id_apply,intervalopen,function_rightunique,funs_is_relation,function_apply_intro,real_mul_eq_subst,real_mul_ident}.
    \end{subproof}
    Therefore there exists $g$ such that
        $g$ is continuous from $X$ to $I$ with respect to $T$ and $T_I$ and
        for all $x\in A$ we have $g(x) = f(x)$ by \cref{subseteq}.
\end{proof}

% from §35 helper lemma 2: absolute value is continuous on the real line
\begin{lemma}\label{real_abs_function_continuous}
    Let $h = \{(x,\abs{x}) \mid x\in\reals\}$.
    Then $h$ is continuous from $\reals$ to $\reals$ with respect to $\standardTopologyR$ and $\standardTopologyR$.
\end{lemma}
\begin{proof}
    For all $x,z_1,z_2$ such that $(x,z_1)\in h$ and $(x,z_2)\in h$ we have $z_1 = z_2$ and
        for all $w$ such that $w\in h$ there exist $x,y$ such that $w = (x,y)$.
    Hence $h$ is a function by \cref{rightunique,relation}.
    For all $x,z$ such that $x\in\dom{h}$ and $(x,z)\in h$ we have $x\in\reals$ and $z = \abs{x}$ by \cref{pair_eq_iff}.
    We have for all $x\in\dom{h}$ we have $h(x)\in\reals$
        by \cref{dom_iff,function_apply_intro,real_abs_in_reals}.
    Hence $h$ is a function to $\reals$.
    For all $x\in\dom{h}$ we have $x\in\reals$ and for all $x\in\reals$ we have $(x,\abs{x})\in h$.
    Therefore $h\in\funs{\reals}{\reals}$ by \cref{subseteq,dom_intro,subseteq_antisymmetric,funs_intro}.

    Show that for all $p\in\reals$ we have
        $h$ is continuous at $p$ from $\reals$ to $\reals$ with respect to $\standardTopologyR$ and $\standardTopologyR$.
    \begin{subproof}
        Fix $p\in\reals$.
        Let $d = \standardMetricR$.
        Let $T_d = \metricTopology{d}{\reals}$.
        Show that for all $V$ such that $V$ is a neighborhood of $h(p)$ in $\reals$ with respect to $\standardTopologyR$
            there exists a neighborhood $U$ of $p$ in $\reals$ with respect to $\standardTopologyR$
            such that $\img{h}{U}\subseteq V$.
        \begin{subproof}
            Fix $V$ such that $V$ is a neighborhood of $h(p)$ in $\reals$ with respect to $\standardTopologyR$.
            $T_d = \standardTopologyR$ and $T_d$ is a topology on $\reals$
                by \cref{standard_metric_topology,standard_basis_is_basis,basis_topology_is_topology,standard_topology_reals}.
            Hence $V\subseteq\reals$ by \cref{neighborhood,open_in,topology_on,subseteq,pow_iff}.
            $d$ is a metric on $\reals$ by \cref{standard_metric_is_metric}.
            Take $\epsilon\in\reals$ such that $0 < \epsilon$ and
                $\metricBall{d}{\reals}{h(p)}{\epsilon} \subseteq V$ by \cref{neighborhood,metric_open_iff}.
            Let $U = \metricBall{d}{\reals}{p}{\epsilon}$.
            For all $q$ such that $q\in U$ we have $h(q)\in \metricBall{d}{\reals}{h(p)}{\epsilon}$ by
                \cref{metric_ball,standard_metric_value,real_abs_in_reals,real_sub,real_add_closed,real_add_inverse,real_reverse_triangle_inequality,real_leq_lt_trans,funs_is_function,function_apply_intro,real_abs_sub_swap,real_lt_subst_left,real_lt_subst_right}.
            We have $\img{h}{U}\subseteq V$ by \cref{img_iff,function_apply_intro,subseteq,subseteq_transitive}.
            For all $q$ such that $q\in U$ we have $q\in\reals$ by \cref{metric_ball}.
            Hence $U\subseteq\reals$ by \cref{subseteq}.
            $p\in\reals$ and $\epsilon\in\reals$ and $0 < \epsilon$.
            Thus $U$ is open in $\reals$ with respect to $\standardTopologyR$ by
                \cref{subseteq,powerset_intro,metric_basis,metric_basis_is_basis,basis_elem_open_in_generated,basis_for_topology,metric_topology,standard_metric_topology,open_in}.
            Therefore $U$ is a neighborhood of $p$ in $\reals$ with respect to $\standardTopologyR$
                by \cref{metric_zero_characterization,standard_metric_is_metric,metric_on,real_lt_subst_left,metric_ball,neighborhood}.
        \end{subproof}
        Therefore $h$ is continuous at $p$ from $\reals$ to $\reals$ with respect to $\standardTopologyR$ and $\standardTopologyR$
            by \cref{standard_basis_is_basis,basis_topology_is_topology,standard_topology_reals,continuous_at}.
    \end{subproof}
    Therefore $h$ is continuous from $\reals$ to $\reals$ with respect to $\standardTopologyR$ and $\standardTopologyR$
        by \cref{standard_basis_is_basis,basis_topology_is_topology,standard_topology_reals,continuous_at_implies_continuous}.
\end{proof}

% from §35 helper lemma 3: bounded transform lands in the open unit interval
\begin{lemma}\label{real_interval_compression}
    Suppose $r\in\reals$.
    Let $s = r \rmul \inv{(1 + \abs{r})}$.
    Then $s\in\intervalopen{\neg{1}}{1}$.
\end{lemma}
\begin{proof}
    $\abs{r}\in\reals$ and $\abs{r}\geq 0$ by \cref{real_abs_in_reals,real_abs_nonnegative}.
    $1\in\reals$ and $0 < 1$ by \cref{real_naturals_subset_reals,real_one_in_naturals,subseteq,real_zero_lt_one}.
    Hence $1 + \abs{r}\in\reals$ and $0 < 1 + \abs{r}$
        by \cref{real_add_closed,real_pos_add_nonneg,real_add_comm}.
    We have $\inv{(1 + \abs{r})}\in\reals$ by \cref{real_mul_inverse,real_pos_nonzero}.
    Hence $s\in\reals$ by \cref{real_mul_closed}.
    $0\in\reals$ by \cref{real_add_ident}.
    We have $r < 0$ or $r = 0$ or $0 < r$ by \cref{real_add_ident,real_lt_trichotomy}.
    \begin{byCase}
    \caseOf{$r = 0$.}
        $s\in\intervalopen{\neg{1}}{1}$ by
            \cref{real_abs_nonneg,real_add_ident,real_mul_eq_subst,real_mul_zero,intervalopen,real_add_inverse,real_zero_lt_one,real_lt_neg,real_neg_zero,real_lt_subst_right,real_lt_subst_left}.
    \caseOf{$0 < r$.}
        We have $\abs{r} = r$ by \cref{real_abs_nonneg,real_lt_imp_leq}.
        Thus $0 < \inv{(1 + \abs{r})}$ by \cref{real_add_eq_subst,real_lt_add,real_add_ident,real_inv_pos,real_lt_subst_left}.
        We have $0 < s$ by
            \cref{real_lt_mul_pos,real_mul_zero,real_lt_subst_right,real_lt_subst_left}.
        Thus $r < 1 + r$ by \cref{real_lt_add,real_add_ident}.
        Hence $s < 1$ by
            \cref{real_lt_mul_pos,real_add_eq_subst,real_mul_inverse,real_lt_subst_left,real_lt_subst_right,real_pos_nonzero}.
        Hence $\neg{1} < s$ by \cref{real_add_inverse,real_zero_lt_one,real_lt_neg,real_neg_zero,real_lt_subst_right,real_lt_trans}.
        Therefore $s\in\intervalopen{\neg{1}}{1}$ by \cref{intervalopen}.
    \caseOf{$r < 0$.}
        Let $u = \neg{r}$.
        Hence $u\in\reals$ and $0 < u$ by \cref{real_add_inverse,real_lt_neg,real_neg_zero,real_lt_subst_left}.
        We have $\abs{r} = u$ by \cref{real_abs_neg,real_lt_imp_leq}.
        Thus $0 < \inv{(1 + \abs{r})}$ by \cref{real_add_eq_subst,real_lt_add,real_add_ident,real_inv_pos,real_lt_subst_left}.
        $\neg{u}\in\reals$ by \cref{real_add_inverse}.
        Hence $s < 0$ by
            \cref{real_add_inverse,real_lt_neg,real_neg_zero,real_inv_pos,real_lt_mul_pos,real_mul_zero,real_lt_subst_left,real_lt_subst_right,real_neg_neg}.
        We have $u < 1 + u$ by \cref{real_lt_add,real_add_ident}.
        Hence $u \rmul \inv{(1 + \abs{r})} < (1 + u) \rmul \inv{(1 + \abs{r})}$ by \cref{real_lt_mul_pos,real_lt_subst_left}.
        Thus $u \rmul \inv{(1 + \abs{r})} < 1$ by \cref{real_add_eq_subst,real_mul_inverse,real_lt_subst_right,real_pos_nonzero}.
        $\neg{1} < \neg{(u \rmul \inv{(1 + \abs{r})})}$ by \cref{real_mul_closed,real_lt_neg,real_neg_zero,real_lt_subst_right}.
        $\neg{u} \rmul \inv{(1 + \abs{r})} = \inv{(1 + \abs{r})} \rmul \neg{u}$ by \cref{real_mul_comm}.
        $\inv{(1 + \abs{r})} \rmul \neg{u} = \neg{(\inv{(1 + \abs{r})} \rmul u)}$ by \cref{real_mul_neg_right}.
        $\inv{(1 + \abs{r})} \rmul u = u \rmul \inv{(1 + \abs{r})}$ by \cref{real_mul_comm}.
        $\neg{u} \rmul \inv{(1 + \abs{r})} = \neg{(u \rmul \inv{(1 + \abs{r})})}$ by \cref{real_lt_subst_right}.
        Thus $\neg{1} < \neg{u} \rmul \inv{(1 + \abs{r})}$ by \cref{real_lt_subst_right}.
        Hence $\neg{1} < s$ by \cref{real_neg_neg,real_lt_subst_right}.
        Hence $s < 1$ by \cref{real_lt_trans,real_zero_lt_one}.
        Therefore $s\in\intervalopen{\neg{1}}{1}$ by \cref{intervalopen}.
    \end{byCase}
\end{proof}

% from §35 helper lemma 4: bounded transform inverts to the original value
\begin{lemma}\label{real_interval_decompression}
    Suppose $r\in\reals$.
    Let $s = r \rmul \inv{(1 + \abs{r})}$.
    Then $s \rmul \inv{(1 - \abs{s})} = r$.
\end{lemma}
\begin{proof}
    $\abs{r}\in\reals$ by \cref{real_abs_in_reals}.
    Hence $1 + \abs{r}\in\reals$ by \cref{real_naturals_subset_reals,real_one_in_naturals,subseteq,real_add_closed}.
    $1\in\reals$ and $0\in\reals$ and $0 < 1$ by \cref{real_naturals_subset_reals,real_one_in_naturals,subseteq,real_add_ident,real_zero_lt_one}.
    We have $0 < 1 + \abs{r}$ by \cref{real_abs_nonnegative,real_pos_add_nonneg,real_add_comm,real_lt_subst_right}.
    We have $\inv{(1 + \abs{r})}\in\reals$ by \cref{real_mul_inverse,real_pos_nonzero}.
    Hence $s\in\reals$ by \cref{real_mul_closed}.
    We have $r < 0$ or $r = 0$ or $0 < r$ by \cref{real_lt_trichotomy}.
    \begin{byCase}
    \caseOf{$r = 0$.}
        We have $s = 0$ by \cref{real_abs_nonneg,real_add_ident,real_mul_eq_subst,real_mul_zero}.
        We have $\abs{s} = 0$ by \cref{real_abs_nonneg,real_lt_subst_right}.
        Hence $1 - \abs{s} = 1$ by
            \cref{real_sub,real_add_eq_subst,real_neg_zero,real_add_comm,real_add_ident,real_lt_subst_right}.
        Thus $\inv{1}\in\reals$ by \cref{real_mul_inverse,real_zero_lt_one,real_pos_nonzero}.
        Therefore $s \rmul \inv{(1 - \abs{s})} = r$ by \cref{real_mul_eq_subst,real_mul_zero,subseteq}.
    \caseOf{$0 < r$.}
        We have $\abs{r} = r$ by \cref{real_abs_nonneg,real_lt_imp_leq}.
        Thus $0 < \inv{(1 + \abs{r})}$ by \cref{real_add_eq_subst,real_lt_add,real_add_ident,real_inv_pos,real_lt_subst_left}.
        We have $0 \rmul \inv{(1 + \abs{r})} < r \rmul \inv{(1 + \abs{r})}$ by \cref{real_lt_mul_pos}.
        Hence $0 < s$ by \cref{real_mul_zero,real_lt_subst_right,real_lt_subst_left}.
        $\abs{s} = s$ by \cref{real_abs_nonneg,real_lt_imp_leq}.
        $(1 + r) \rmul \inv{(1 + \abs{r})} = 1$ by \cref{real_add_eq_subst,real_mul_inverse,real_pos_nonzero}.
        $1 - s = ((1 + r) \rmul \inv{(1 + \abs{r})}) - s$ by \cref{real_lt_subst_left}.
        Hence $1 - s = ((1 + r) \rmul \inv{(1 + \abs{r})}) - (r \rmul \inv{(1 + \abs{r})})$ by \cref{real_mul_eq_subst,real_add_eq_subst}.
        Thus $1 - s = ((1 + r) \rmul \inv{(1 + \abs{r})}) + \neg{(r \rmul \inv{(1 + \abs{r})})}$ by \cref{real_sub}.
        $\neg{r}\in\reals$ by \cref{real_add_inverse}.
        $\neg{(r \rmul \inv{(1 + \abs{r})})} = \neg{(\inv{(1 + \abs{r})} \rmul r)}$ by \cref{real_mul_comm,real_lt_subst_right}.
        We have $\inv{(1 + \abs{r})} \rmul \neg{r} = \neg{(\inv{(1 + \abs{r})} \rmul r)}$ by \cref{real_mul_neg_right}.
        We have $\inv{(1 + \abs{r})} \rmul \neg{r} = (\neg{r}) \rmul \inv{(1 + \abs{r})}$ by \cref{real_mul_comm}.
        We have $\neg{(r \rmul \inv{(1 + \abs{r})})} = (\neg{r}) \rmul \inv{(1 + \abs{r})}$ by \cref{real_lt_subst_left}.
        Thus $1 - s = ((1 + r) \rmul \inv{(1 + \abs{r})}) + ((\neg{r}) \rmul \inv{(1 + \abs{r})})$ by \cref{real_add_eq_subst}.
        We have $((1 + r) \rmul \inv{(1 + \abs{r})}) + ((\neg{r}) \rmul \inv{(1 + \abs{r})})
            = ((1 + r) + \neg{r}) \rmul \inv{(1 + \abs{r})}$ by \cref{real_right_distrib}.
        We have $1 - s = ((1 + r) + \neg{r}) \rmul \inv{(1 + \abs{r})}$ by \cref{real_add_eq_subst}.
        Hence $(1 + r) + \neg{r} = 1$
            by \cref{real_add_assoc,real_add_inverse,real_add_eq_subst,real_add_comm,real_add_ident,real_lt_subst_right}.
        Thus $1 - s = 1 \rmul \inv{(1 + \abs{r})}$ by \cref{real_mul_eq_subst}.
        Hence $1 - s = \inv{(1 + \abs{r})}$ by \cref{real_mul_ident}.
        We have $1 - \abs{s} = \inv{(1 + \abs{r})}$ by \cref{subseteq}.
        We have $\inv{(1 + \abs{r})} \rmul (1 + \abs{r}) = 1$ by \cref{real_mul_inverse,real_pos_nonzero,real_mul_comm}.
        Thus $\inv{(1 + \abs{r})}\neq 0$ by \cref{real_inv_pos,real_pos_nonzero}.
        Hence $\inv{\inv{(1 + \abs{r})}} = 1 + \abs{r}$ by \cref{real_mul_inverse_unique}.
        We have $\inv{(1 - \abs{s})} = \inv{\inv{(1 + \abs{r})}}$ by \cref{real_mul_eq_subst}.
        Hence $\inv{(1 - \abs{s})} = 1 + \abs{r}$ by \cref{real_lt_subst_right}.
        Thus $s \rmul \inv{(1 - \abs{s})} = s \rmul (1 + \abs{r})$ by \cref{real_mul_eq_subst}.
        Hence $s \rmul \inv{(1 - \abs{s})} = (r \rmul \inv{(1 + \abs{r})}) \rmul (1 + \abs{r})$ by \cref{real_mul_eq_subst}.
        We have $s \rmul \inv{(1 - \abs{s})} = r \rmul (\inv{(1 + \abs{r})} \rmul (1 + \abs{r}))$ by \cref{real_mul_assoc}.
        Hence $s \rmul \inv{(1 - \abs{s})} = r \rmul 1$ by \cref{real_mul_inverse,real_pos_nonzero,real_mul_eq_subst,real_mul_comm}.
        Therefore $s \rmul \inv{(1 - \abs{s})} = r$ by \cref{real_mul_comm,real_mul_ident,real_lt_subst_right}.
    \caseOf{$r < 0$.}
        Let $u = \neg{r}$.
        Hence $u\in\reals$ and $0 < u$ by \cref{real_add_inverse,real_lt_neg,real_neg_zero,real_lt_subst_left}.
        We have $\abs{r} = u$ by \cref{real_abs_neg,real_lt_imp_leq}.
        Thus $0 < \inv{(1 + \abs{r})}$ by \cref{real_add_eq_subst,real_lt_add,real_add_ident,real_inv_pos,real_lt_subst_left}.
        $\neg{u}\in\reals$ by \cref{real_add_inverse}.
        Hence $\neg{u} < 0$ by \cref{real_lt_neg,real_neg_zero,real_lt_subst_right}.
        We have $\neg{u} \rmul \inv{(1 + \abs{r})} < 0 \rmul \inv{(1 + \abs{r})}$ by \cref{real_lt_mul_pos}.
        Hence $s < 0$ by \cref{real_mul_zero,real_lt_subst_right,real_lt_subst_left,real_neg_neg}.
        Thus $\abs{s} = \neg{s}$ by \cref{real_abs_neg,real_lt_imp_leq}.
        We have $(1 + u) \rmul \inv{(1 + \abs{r})} = 1$ by \cref{real_add_eq_subst,real_mul_inverse,real_pos_nonzero}.
        Hence $1 - \abs{s} = 1 + \neg{\abs{s}}$ by \cref{real_sub}.
        We have $1 + \neg{\abs{s}} = 1 + \neg{\neg{s}}$ by \cref{real_add_eq_subst}.
        Hence $\neg{\neg{s}} = s$ by \cref{real_neg_neg}.
        Thus $1 + \neg{\neg{s}} = 1 + s$ by \cref{real_add_eq_subst}.
        Hence $1 - \abs{s} = 1 + s$ by \cref{real_lt_subst_right}.
        We have $1 + s = ((1 + u) \rmul \inv{(1 + \abs{r})}) + s$ by \cref{real_lt_subst_left}.
        Hence $1 + s = ((1 + u) \rmul \inv{(1 + \abs{r})}) + (\neg{u} \rmul \inv{(1 + \abs{r})})$ by \cref{real_mul_eq_subst,real_add_eq_subst,real_neg_neg}.
        We have $((1 + u) \rmul \inv{(1 + \abs{r})}) + (\neg{u} \rmul \inv{(1 + \abs{r})})
            = ((1 + u) + \neg{u}) \rmul \inv{(1 + \abs{r})}$ by \cref{real_right_distrib}.
        Thus $1 + s = ((1 + u) + \neg{u}) \rmul \inv{(1 + \abs{r})}$ by \cref{real_add_eq_subst}.
        Hence $(1 + u) + \neg{u} = 1$
            by \cref{real_add_assoc,real_add_inverse,real_add_eq_subst,real_add_comm,real_add_ident,real_lt_subst_right}.
        We have $1 + s = 1 \rmul \inv{(1 + \abs{r})}$ by \cref{real_mul_eq_subst}.
        Hence $1 + s = \inv{(1 + \abs{r})}$ by \cref{real_mul_ident}.
        Thus $1 - \abs{s} = \inv{(1 + \abs{r})}$ by \cref{subseteq}.
        We have $\inv{(1 + \abs{r})} \rmul (1 + \abs{r}) = 1$ by \cref{real_mul_inverse,real_pos_nonzero,real_mul_comm}.
        We have $\inv{(1 + \abs{r})}\neq 0$ by \cref{real_inv_pos,real_pos_nonzero}.
        Hence $\inv{\inv{(1 + \abs{r})}} = 1 + \abs{r}$ by \cref{real_mul_inverse_unique}.
        Thus $\inv{(1 - \abs{s})} = \inv{\inv{(1 + \abs{r})}}$ by \cref{real_mul_eq_subst}.
        Hence $\inv{(1 - \abs{s})} = 1 + \abs{r}$ by \cref{real_lt_subst_right}.
        We have $s \rmul \inv{(1 - \abs{s})} = s \rmul (1 + \abs{r})$ by \cref{real_mul_eq_subst}.
        Hence $s \rmul \inv{(1 - \abs{s})} = (\neg{u} \rmul \inv{(1 + \abs{r})}) \rmul (1 + \abs{r})$ by \cref{real_mul_eq_subst,real_neg_neg}.
        Thus $s \rmul \inv{(1 - \abs{s})} = \neg{u} \rmul (\inv{(1 + \abs{r})} \rmul (1 + \abs{r}))$ by \cref{real_mul_assoc}.
        Hence $s \rmul \inv{(1 - \abs{s})} = \neg{u} \rmul 1$ by \cref{real_mul_inverse,real_pos_nonzero,real_mul_eq_subst,real_mul_comm}.
        We have $s \rmul \inv{(1 - \abs{s})} = \neg{u}$ by \cref{real_mul_comm,real_mul_ident,real_lt_subst_right}.
        Therefore $s \rmul \inv{(1 - \abs{s})} = r$ by \cref{real_neg_neg}.
    \end{byCase}
\end{proof}

\begin{theorem}\label{tietze_extension_theorem_real_valued}
    Assume $X$ is normal with respect to $T$.
    Assume $A$ is closed in $X$ with respect to $T$.
    Let $T_A = \subspaceTopology{T}{A}$.
    Assume $f$ is continuous from $A$ to $\reals$ with respect to $T_A$ and $\standardTopologyR$.
    Then there exists $g$ such that
        $g$ is continuous from $X$ to $\reals$ with respect to $T$ and $\standardTopologyR$ and
        for all $x\in A$ we have $g(x) = f(x)$.
\end{theorem}
\begin{proof}
    Let $I = \intervalopen{\neg{1}}{1}$.
    Let $T_I = \subspaceTopology{\standardTopologyR}{I}$.
    Let $a = \{(x,\abs{x}) \mid x\in\reals\}$.
    $a$ is continuous from $\reals$ to $\reals$ with respect to $\standardTopologyR$ and $\standardTopologyR$
        by \cref{real_abs_function_continuous}.
    Hence $a\in\funs{\reals}{\reals}$ and $a$ is a function and $\dom{a} = \reals$
        by \cref{continuous,funs_is_function,funs}.

    $T$ is a topology on $X$ by \cref{normal_space}.
    We have $A\subseteq X$ by \cref{closed_in}.
    $T_A$ is a topology on $A$ by \cref{subspace_topology_is_topology,subspace_of}.
    $\standardTopologyR$ is a topology on $\reals$
        by \cref{standard_basis_is_basis,basis_topology_is_topology,standard_topology_reals}.
    We have $I\subseteq\reals$ and $I$ is a subspace of $\reals$ with respect to $\standardTopologyR$
        by \cref{intervalopen,subseteq,subspace_of}.
    Hence $f\in\funs{A}{\reals}$ and $f$ is a function and $\dom{f} = A$
        by \cref{continuous,funs_is_function,funs}.
    We have $f$ is right-unique and $f$ is a relation and $\ran{f}\subseteq\dom{a}$
        by \cref{function_rightunique,funs_is_relation,funs_ran,subseteq}.

    Let $b = a\circ f$.
    $b$ is continuous from $A$ to $\reals$ with respect to $T_A$ and $\standardTopologyR$
        by \cref{continuous_composite}.
    $1\in\reals$ by \cref{real_naturals_subset_reals,real_one_in_naturals,subseteq}.
    Let $w_A = \{(x,1) \mid x\in A\}$.
    For all $x\in A$ we have $w_A(x) = 1$
        by \cref{constant_map_function,funs_is_function,function_apply_intro}.
    $w_A$ is continuous from $A$ to $\reals$ with respect to $T_A$ and $\standardTopologyR$
        by \cref{constant_map_function,constant_continuous}.
    Let $d = \{(x, w_A(x) + b(x)) \mid x\in A\}$.
    $d$ is continuous from $A$ to $\reals$ with respect to $T_A$ and $\standardTopologyR$
        by \cref{real_function_addition_continuous}.
    We have $d\in\funs{A}{\reals}$ and $d$ is a function and $\dom{d} = A$
        by \cref{continuous,funs_is_function,funs}.

    Show that for all $x\in A$ we have $d(x) = 1 + \abs{f(x)}$ and $0 < d(x)$.
    \begin{subproof}
        Follows by \cref{funs_type_apply,circ_apply,function_apply_intro,real_lt_subst_right,real_add_eq_subst,real_abs_in_reals,real_abs_nonnegative,real_pos_add_nonneg,real_zero_lt_one,real_add_comm}.
    \end{subproof}
    For all $x\in A$ we have $d(x)\neq 0$ by \cref{continuous,funs_type_apply,real_pos_nonzero}.

    Let $c = \{(x, f(x) \rmul \inv{(d(x))}) \mid x\in A\}$.
    $c$ is continuous from $A$ to $\reals$ with respect to $T_A$ and $\standardTopologyR$
        by \cref{real_function_division_continuous}.
    We have $c\in\funs{A}{\reals}$ and $c$ is a function and $\dom{c} = A$
        by \cref{continuous,funs_is_function,funs}.

    Hence $\img{c}{A}\subseteq I$ by
        \cref{img_iff,continuous,funs,subseteq,funs_is_function,function_apply_iff,funs_type_apply,circ_apply,function_apply_intro,real_lt_subst_right,real_add_eq_subst,real_mul_eq_subst,real_interval_compression,real_lt_subst_left}.
    $c$ is continuous from $A$ to $I$ with respect to $T_A$ and $T_I$
        by \cref{subseteq,continuous_restrict_range}.

    Take $u$ such that
        $u$ is continuous from $X$ to $I$ with respect to $T$ and $T_I$ and
        for all $x\in A$ we have $u(x) = c(x)$
        by \cref{tietze_extension_theorem_open_interval}.
    We have $u\in\funs{X}{I}$ and $u$ is a function and $\dom{u} = X$
        by \cref{continuous,funs_is_function,funs}.

    Let $j_I = \identity{I}$.
    $j_I$ is right-unique and $j_I$ is a relation and $\dom{j_I} = I$
        by \cref{id_rightunique,id_is_relation,id_dom}.
    Hence $\ran{u}\subseteq\dom{j_I}$ by \cref{funs_ran,id_dom,subseteq}.
    Let $u_1 = j_I\circ u$.
    $u_1$ is continuous from $X$ to $\reals$ with respect to $T$ and $\standardTopologyR$
        by \cref{continuous_expand_range}.
    We have $u_1\in\funs{X}{\reals}$ and $u_1$ is a function and $\dom{u_1} = X$
        by \cref{continuous,funs_is_function,funs}.
    Let $e = a\circ u_1$.
    $e$ is continuous from $X$ to $\reals$ with respect to $T$ and $\standardTopologyR$
        by \cref{continuous_composite}.
    We have $e\in\funs{X}{\reals}$ and $e$ is a function and $\dom{e} = X$
        by \cref{continuous,funs_is_function,funs}.
    Let $w_X = \{(x,1) \mid x\in X\}$.
    We have for all $x\in X$ we have $w_X(x) = 1$
        by \cref{constant_map_function,funs_is_function,function_apply_intro}.
    $w_X$ is continuous from $X$ to $\reals$ with respect to $T$ and $\standardTopologyR$
        by \cref{constant_map_function,constant_continuous}.
    Let $m = \{(x, w_X(x) - e(x)) \mid x\in X\}$.
    $m$ is continuous from $X$ to $\reals$ with respect to $T$ and $\standardTopologyR$
        by \cref{real_function_subtraction_continuous}.
    We have $m\in\funs{X}{\reals}$ and $m$ is a function and $\dom{m} = X$
        by \cref{continuous,funs_is_function,funs}.

    Show that for all $x\in X$ we have $u_1(x)\in\reals$ and $\neg{1} < u_1(x)$ and $u_1(x) < 1$.
    \begin{subproof}
        Follows by \cref{continuous,funs_type_apply,circ_apply,id_apply,real_lt_subst_left,intervalopen}.
    \end{subproof}
    For all $x\in X$ we have $\abs{u_1(x)} < 1$ by
        \cref{real_add_ident,real_naturals_subset_reals,real_one_in_naturals,subseteq,real_zero_lt_one,real_add_inverse,real_lt_add,real_sub,real_lt_subst_left,real_lt_subst_right,real_add_comm,real_abs_lt_from_interval,real_neg_zero}.
    Show that for all $x\in X$ we have $e(x) = \abs{u_1(x)}$ and $m(x) = 1 - \abs{u_1(x)}$.
    \begin{subproof}
        Follows by \cref{circ_apply,funs_ran,subseteq,function_apply_intro,real_lt_subst_right,real_lt_subst_left}.
    \end{subproof}
    We have for all $x\in X$ we have $0 < m(x)$ by
        \cref{real_abs_in_reals,real_add_inverse,real_lt_add,real_sub,real_lt_subst_left,real_lt_subst_right}.
    For all $x\in X$ we have $m(x)\neq 0$ by \cref{continuous,funs_type_apply,real_pos_nonzero}.

    Let $h = \{(x, u_1(x) \rmul \inv{(m(x))}) \mid x\in X\}$.
    $h$ is continuous from $X$ to $\reals$ with respect to $T$ and $\standardTopologyR$
        by \cref{real_function_division_continuous}.
    We have $h\in\funs{X}{\reals}$ and $h$ is a function and $\dom{h} = X$
        by \cref{continuous,funs_is_function,funs}.

    Show that for all $x\in A$ we have $f(x)\in\reals$ and $u_1(x) = f(x) \rmul \inv{(1 + \abs{f(x)})}$.
    \begin{subproof}
        Follows by \cref{subseteq,funs_type_apply,circ_apply,id_apply,real_lt_subst_right,function_apply_intro,real_add_eq_subst,real_mul_eq_subst}.
    \end{subproof}
    Show that for all $x\in A$ we have $h(x) = u_1(x) \rmul \inv{(1 - \abs{u_1(x)})}$.
    \begin{subproof}
        Follows by \cref{continuous,funs_type_apply,funs_ran,subseteq,circ_apply,function_apply_intro,real_lt_subst_right,real_lt_subst_left,real_mul_eq_subst}.
    \end{subproof}
    Show that for all $x\in A$ we have $h(x) = f(x)$.
    \begin{subproof}
        Follows by \cref{real_interval_decompression,real_lt_subst_right}.
    \end{subproof}
    Follows by \cref{subseteq}.
\end{proof}
\section{§36 Imbeddings of Manifolds}

% from §36 Item 1: $m$-manifold
\begin{definition}\label{m_manifold}
    $X$ is a manifold of dimension $m$ with respect to $T$ iff
        $X$ is Hausdorff with respect to $T$ and
        $X$ satisfies the second countability axiom with respect to $T$ and
        for all $x\in X$ there exist $U, V, R, S, h$ such that
            $U$ is a neighborhood of $x$ in $X$ with respect to $T$ and
            $R = \{(i,\reals) \mid i\in\initialSegment{m}\}$ and
            $S = \{(i,\standardTopologyR) \mid i\in\initialSegment{m}\}$ and
            $V$ is open in $\productFamily{R}$ with respect to $\productTopologyIndexed{S}{R}$ and
            $h$ is a homeomorphism between $U$ and $V$ with respect to
                $\subspaceTopology{T}{U}$ and $\subspaceTopology{\productTopologyIndexed{S}{R}}{V}$.
\end{definition}

% from §36 Item 2: support
\begin{definition}\label{support}
    $\supp{f}{X}{T} = \closure{X\setminus\preimg{f}{\{0\}}}{X}{T}$.
\end{definition}

% from §36 helper definition 5: component condition for a finite partition of unity
\begin{definition}\label{finite_partition_component_condition}
    $\finitePartitionComponentCondition{f}{X}{U}{T}{n}$ iff
        for all $i\in\initialSegment{n}$ there exists $I$ such that
            $I = \intervalclosed{0}{1}$ and
            $\apply{f}{i}$ is continuous from $X$ to $I$ with respect to $T$ and $\subspaceTopology{\standardTopologyR}{I}$ and
            $\supp{\apply{f}{i}}{X}{T}\subseteq \apply{U}{i}$.
\end{definition}

% from §36 helper definition 6: sum condition for a finite partition of unity
\begin{definition}\label{finite_partition_sum_condition}
    $\finitePartitionSumCondition{f}{X}{n}$ iff
        for all $x\in X$ there exists $a$ such that
            $a\in\funs{\initialSegment{n}}{\reals}$ and
            for all $i\in\initialSegment{n}$ we have $a(i) = \apply{\apply{f}{i}}{x}$ and
            $\sumFiniteSequence{a}{1}$.
\end{definition}

% from §36 Item 3: partition of unity dominated by a finite indexed open covering
\begin{definition}\label{finite_partition_of_unity}
    $f$ is a partition of unity on $X$ dominated by $U$ with respect to $T$ and $n$ iff
        $n\in\naturalsPlus$ and
        $U$ is a function on $\initialSegment{n}$ and
        for all $i\in\initialSegment{n}$ we have $\apply{U}{i}$ is open in $X$ with respect to $T$ and
        $\img{U}{\initialSegment{n}}$ is a covering of $X$ and
        $f$ is a function on $\initialSegment{n}$ and
        $\finitePartitionComponentCondition{f}{X}{U}{T}{n}$ and
        $\finitePartitionSumCondition{f}{X}{n}$.
\end{definition}

% from §36 helper definition 7: witness for continuous finite real sums
\begin{definition}\label{finite_real_sum_witness}
    $g$ is a finite real sum witness for $f$ on $X$ with respect to $T$ and $n$ iff
        $g$ is continuous from $X$ to $\reals$ with respect to $T$ and $\standardTopologyR$ and
        for all $x\in X$ there exists $a$ such that
            $a\in\funs{\initialSegment{n}}{\reals}$ and
            $\sumFiniteSequence{a}{g(x)}$ and
            for all $i\in\initialSegment{n}$ we have $a(i) = \apply{\apply{f}{i}}{x}$.
\end{definition}

% from §36 helper definition 1: partial shrinking stage
\begin{definition}\label{finite_shrink_stage}
    $\finiteShrinkStage{U}{X}{T}{n}{k}$ iff
        if $k \leq n$ then there exists $V$ such that
            $V$ is a function on $\initialSegment{k}$ and
            for all $i\in\initialSegment{k}$ we have
                $\apply{V}{i}$ is open in $X$ with respect to $T$ and
                $\closure{\apply{V}{i}}{X}{T}\subseteq \apply{U}{i}$
            and
            $\finiteShrinkCoverCondition{V}{U}{X}{n}{k}$.
\end{definition}

% from §36 helper definition 2: stage covering condition
\begin{definition}\label{finite_shrink_cover_condition}
    $\finiteShrinkCoverCondition{V}{U}{X}{n}{k}$ iff
        for all $x\in X$ there exists $i\in\initialSegment{n}$ such that
            either $i \leq k$ and $x\in \apply{V}{i}$
            or $k < i$ and $x\in \apply{U}{i}$.
\end{definition}

% from §36 helper definition 3: tail indices of a finite cover
\begin{definition}\label{finite_shrink_tail_indices}
    $\finiteShrinkTailIndices{n}{k} = \{i\in\initialSegment{n} \mid k < i\}$.
\end{definition}

% from §36 helper definition 4: set of shrinking stages
\begin{definition}\label{finite_shrink_stage_set}
    $\finiteShrinkStageSet{U}{X}{T}{n} = \{k\in\naturalsPlus \mid \finiteShrinkStage{U}{X}{T}{n}{k}\}$.
\end{definition}

% from §36 helper lemma 2: first shrinking stage
\begin{lemma}\label{finite_shrink_stage_one}
    Assume $X$ is normal with respect to $T$.
    Assume $n\in\naturalsPlus$.
    Assume $U$ is a function on $\initialSegment{n}$.
    Suppose for all $i\in\initialSegment{n}$ we have $\apply{U}{i}$ is open in $X$ with respect to $T$.
    Suppose $\img{U}{\initialSegment{n}}$ is a covering of $X$.
    Then $\finiteShrinkStage{U}{X}{T}{n}{1}$.
\end{lemma}
\begin{proof}
    Let $J = \finiteShrinkTailIndices{n}{1}$.
    Let $C = X\setminus\unions{\img{U}{J}}$.
    Show that if $1 \leq n$ then there exists $V$ such that
        $V$ is a function on $\initialSegment{1}$ and
        for all $i\in\initialSegment{1}$ we have
            $\apply{V}{i}$ is open in $X$ with respect to $T$ and
            $\closure{\apply{V}{i}}{X}{T}\subseteq \apply{U}{i}$ and
        $\finiteShrinkCoverCondition{V}{U}{X}{n}{1}$.
    \begin{subproof}
        Assume $1 \leq n$.
        For all $W$ such that $W\in\img{U}{J}$ we have $W\in T$
            by \cref{img_iff,finite_shrink_tail_indices,function_apply_intro,subseteq,open_in}.
        We have $\unions{\img{U}{J}}\in T$
            by \cref{img_iff,finite_shrink_tail_indices,function_apply_intro,subseteq,open_in,normal_space,topology_on,pow_iff}.
        Hence $\unions{\img{U}{J}}$ is open in $X$ with respect to $T$ by \cref{normal_space,open_in}.
        We have $\unions{\img{U}{J}}\subseteq X$ by \cref{open_in,normal_space,topology_on,subseteq,pow_iff}.
        Therefore $C$ is closed in $X$ with respect to $T$
            by \cref{setminus_subseteq,double_relative_complement,subseteq,open_in,closed_in}.
        For all $x$ such that $x\in C$ we have $x\in \apply{U}{1}$
            by \cref{setminus_elim_left,covering,unions_iff,subseteq,img_iff,function_apply_intro,function_apply_elim,initial_segment_naturalsplus,real_one_leq_naturalsplus,finite_shrink_tail_indices,real_leq_split,setminus_elim_right}.
        We have $C\subseteq \apply{U}{1}$ by \cref{subseteq}.
        $\apply{U}{1}$ is open in $X$ with respect to $T$
            by \cref{real_one_in_naturals,initial_segment_naturalsplus,subseteq}.
        Take $W$ such that
            $W$ is open in $X$ with respect to $T$ and
            $C\subseteq W$ and
            $\closure{W}{X}{T}\subseteq \apply{U}{1}$
            by \cref{normal_closure_neighborhood}.
        Let $V = \{(1,W)\}$.
        $V$ is a function by \cref{singleton_elim,pair_eq_iff,rightunique,relation}.
        We have $\dom{V} = \initialSegment{1}$ by
            \cref{dom_iff,singleton_elim,pair_eq_iff,real_one_in_naturals,initial_segment_naturalsplus,real_naturals_subset_reals,subseteq,real_one_leq_naturalsplus,real_leq_antisym,singleton_intro,dom_intro,subseteq_antisymmetric}.
        We have $V$ is a function on $\initialSegment{1}$ by \cref{funs}.
        For all $i\in\initialSegment{1}$ we have
            $\apply{V}{i}$ is open in $X$ with respect to $T$ and
            $\closure{\apply{V}{i}}{X}{T}\subseteq \apply{U}{i}$
            by \cref{initial_segment_naturalsplus,real_one_leq_naturalsplus,real_one_in_naturals,real_naturals_subset_reals,subseteq,real_leq_split,real_lt_trans,real_lt_irreflexive,singleton_intro,pair_eq_iff,function_apply_intro}.
        Show that for all $x\in X$ there exists $i$ such that
            $i\in\initialSegment{n}$ and
            either $i \leq 1$ and $x\in \apply{V}{i}$
            or $1 < i$ and $x\in \apply{U}{i}$.
        \begin{subproof}
            Fix $x\in X$.
            $x\in W$ or $x\notin W$.
            \begin{byCase}
            \caseOf{$x\in W$.}
                Take $i$ such that $i = 1$.
                We have $i\in\initialSegment{n}$ and $i \leq 1$ and $x\in \apply{V}{i}$ by
                    \cref{real_one_in_naturals,initial_segment_naturalsplus,subseteq,singleton_intro,pair_eq_iff,function_apply_intro}.
            \caseOf{$x\notin W$.}
                Take $W_0\in\img{U}{J}$ such that $x\in W_0$ by \cref{subseteq,setminus_intro,unions_iff}.
                Take $i$ such that $i\in J$ and $(i,W_0)\in U$ by \cref{img_iff}.
                We have $i\in\initialSegment{n}$ and $1 < i$ and $x\in \apply{U}{i}$
                    by \cref{finite_shrink_tail_indices,function_apply_intro}.
            \end{byCase}
        \end{subproof}
        Therefore $\finiteShrinkCoverCondition{V}{U}{X}{n}{1}$ by \cref{finite_shrink_cover_condition}.
        Follows by \cref{pair_eq_iff}.
    \end{subproof}
    Therefore $\finiteShrinkStage{U}{X}{T}{n}{1}$ by \cref{finite_shrink_stage}.
\end{proof}

% from §36 helper lemma 3: last element of a successor initial segment
\begin{lemma}\label{initial_segment_succ_last}
    Suppose $k\in\naturalsPlus$.
    Suppose $i\in\initialSegment{k + 1}$.
    Suppose $k < i$.
    Then $i = k + 1$.
\end{lemma}
\begin{proof}
    We have $i = k + 1$
        by \cref{real_naturals_leq_succ_elim,initial_segment_naturalsplus,real_naturals_subset_reals,subseteq,real_lt_trans,real_lt_irreflexive}.
\end{proof}

% from §36 helper lemma 1: shrinking a finite indexed open cover
\begin{lemma}\label{finite_shrink_cover_condition_final}
    Assume $n\in\naturalsPlus$.
    Assume $V$ is a function on $\initialSegment{n}$.
    Suppose $\finiteShrinkCoverCondition{V}{U}{X}{n}{n}$.
    Then $\img{V}{\initialSegment{n}}$ is a covering of $X$.
\end{lemma}
\begin{proof}
    For all $x$ such that $x\in X$ we have $x\in\unions{\img{V}{\initialSegment{n}}}$
        by \cref{finite_shrink_cover_condition,initial_segment_naturalsplus,real_naturals_subset_reals,subseteq,real_lt_trans,real_lt_irreflexive,function_apply_elim,img_iff,unions_iff}.
    Therefore $\img{V}{\initialSegment{n}}$ is a covering of $X$ by \cref{subseteq,covering}.
\end{proof}

% from §36 helper lemma 5: shrinking stage successor
\begin{lemma}\label{finite_shrink_stage_succ}
    Assume $X$ is normal with respect to $T$.
    Assume $n\in\naturalsPlus$.
    Assume $U$ is a function on $\initialSegment{n}$.
    Suppose for all $i\in\initialSegment{n}$ we have $\apply{U}{i}$ is open in $X$ with respect to $T$.
    Suppose $\img{U}{\initialSegment{n}}$ is a covering of $X$.
    Assume $k\in\naturalsPlus$.
    Suppose $\finiteShrinkStage{U}{X}{T}{n}{k}$.
    Then $\finiteShrinkStage{U}{X}{T}{n}{k + 1}$.
\end{lemma}
\begin{proof}
    Show that if $k + 1 \leq n$ then there exists $V$ such that
        $V$ is a function on $\initialSegment{k + 1}$ and
        for all $i\in\initialSegment{k + 1}$ we have
            $\apply{V}{i}$ is open in $X$ with respect to $T$ and
            $\closure{\apply{V}{i}}{X}{T}\subseteq \apply{U}{i}$ and
            $\finiteShrinkCoverCondition{V}{U}{X}{n}{k + 1}$.
    \begin{subproof}
        Assume $k + 1 \leq n$.
        $k\in\reals$ and $n\in\reals$ and $k + 1\in\reals$
            by \cref{real_naturals_subset_reals,real_naturals_closed_succ,subseteq,real_add_closed,real_one_in_naturals}.
        We have $k < k + 1$ by \cref{real_naturals_lt_succ_local}.
        Hence $k \leq n$ by \cref{real_lt_imp_leq,real_leq_trans}.
        Take $V_0$ such that
            $V_0$ is a function on $\initialSegment{k}$ and
            for all $i\in\initialSegment{k}$ we have
                $\apply{V_0}{i}$ is open in $X$ with respect to $T$ and
                $\closure{\apply{V_0}{i}}{X}{T}\subseteq \apply{U}{i}$ and
            $\finiteShrinkCoverCondition{V_0}{U}{X}{n}{k}$
            by \cref{finite_shrink_stage}.
        Let $J = \finiteShrinkTailIndices{n}{k + 1}$.
        Let $A = \unions{\img{V_0}{\initialSegment{k}}}$.
        Let $B = \unions{\img{U}{J}}$.
        Let $F = \{A,B\}$.
        Let $D = \unions{F}$.
        Let $C = X\setminus D$.

        For all $W$ such that $W\in\img{V_0}{\initialSegment{k}}$ we have $W\in T$
            by \cref{img_iff,function_apply_intro,subseteq,open_in}.
        $T$ is a topology on $X$ by \cref{normal_space}.
        We have $\img{V_0}{\initialSegment{k}}\subseteq T$ by \cref{subseteq}.
        $A\in T$ by \cref{topology_on,pow_iff}.
        Therefore $A$ is open in $X$ with respect to $T$ by \cref{open_in}.

        For all $W$ such that $W\in\img{U}{J}$ we have $W\in T$
            by \cref{img_iff,finite_shrink_tail_indices,function_apply_intro,subseteq,open_in}.
        We have $\img{U}{J}\subseteq T$ by \cref{subseteq}.
        $B\in T$ by \cref{normal_space,topology_on,pow_iff}.
        Hence $B$ is open in $X$ with respect to $T$ by \cref{open_in}.

        For all $W$ such that $W\in F$ we have $W\in T$ by \cref{upair_iff,open_in}.
        We have $F\subseteq T$ by \cref{subseteq}.
        $D\in T$ by \cref{normal_space,topology_on,pow_iff}.
        Therefore $D$ is open in $X$ with respect to $T$ by \cref{open_in}.
        $D\subseteq X$ by \cref{open_in,normal_space,topology_on,subseteq,pow_iff}.
        Hence $C$ is closed in $X$ with respect to $T$
            by \cref{setminus_subseteq,double_relative_complement,subseteq,open_in,closed_in}.

        For all $x$ such that $x\in C$ we have $x\in \apply{U}{k + 1}$
            by \cref{setminus_elim_left,finite_shrink_cover_condition,initial_segment_naturalsplus,funs_elim,subseteq,function_apply_elim,img_iff,unions_iff,upair_intro_left,upair_intro_right,setminus_elim_right,finite_shrink_tail_indices,real_naturals_subset_reals,real_naturals_closed_succ,real_lt_trichotomy,real_lt_imp_leq,real_naturals_leq_succ_elim,real_lt_trans,real_lt_irreflexive,real_lt_subst_right}.
        We have $C\subseteq \apply{U}{k + 1}$ by \cref{subseteq}.

        $k + 1\in\initialSegment{n}$ by \cref{real_naturals_closed_succ,initial_segment_naturalsplus}.
        $\apply{U}{k + 1}$ is open in $X$ with respect to $T$ by \cref{subseteq}.
        Take $W$ such that
            $W$ is open in $X$ with respect to $T$ and
            $C\subseteq W$ and
            $\closure{W}{X}{T}\subseteq \apply{U}{k + 1}$
            by \cref{normal_closure_neighborhood}.
        Let $G = \{(k + 1,W)\}$.
        $G$ is a function and $\dom{G} = \{k + 1\}$
            by \cref{singleton_elim,pair_eq_iff,rightunique,relation,dom_iff,singleton_intro,dom_intro,subseteq,subseteq_antisymmetric}.

        For all $a$ such that $a\in\dom{V_0}\inter\dom{G}$ we have $\apply{V_0}{a} = \apply{G}{a}$.
        Let $V = V_0\union G$.
        We have $V$ is a function by \cref{union_compatible_functions}.
        $\dom{V} = \dom{V_0}\union\dom{G}$ by \cref{dom_union}.
        Hence $\dom{V} = \initialSegment{k}\union \{k + 1\}$ by \cref{funs_elim}.

        For all $a$ such that $a\in\initialSegment{k}\union\{k + 1\}$ we have $a\in\initialSegment{k + 1}$
            by \cref{union_iff,initial_segment_naturalsplus,real_naturals_subset_reals,subseteq,real_one_in_naturals,real_add_ident,real_add_closed,real_zero_lt_one,real_lt_add,real_add_comm,real_lt_subst_left,real_lt_trans,singleton_elim}.
        For all $a$ such that $a\in\initialSegment{k + 1}$ we have $a\in\initialSegment{k}\union\{k + 1\}$
            by \cref{initial_segment_naturalsplus,real_naturals_leq_succ_elim,union_iff,singleton_intro}.
        Hence $\initialSegment{k}\union\{k + 1\} = \initialSegment{k + 1}$ by \cref{subseteq,subseteq_antisymmetric}.
        Therefore $V$ is a function on $\initialSegment{k + 1}$ by \cref{subseteq,subseteq_antisymmetric,funs}.

        For all $i\in\initialSegment{k + 1}$ we have
            $\apply{V}{i}$ is open in $X$ with respect to $T$ and
            $\closure{\apply{V}{i}}{X}{T}\subseteq \apply{U}{i}$
            by \cref{singleton_intro,pair_eq_iff,union_iff,function_apply_intro,initial_segment_naturalsplus,real_naturals_leq_succ_elim,funs_elim,subseteq,function_apply_elim}.

        Show that for all $x\in X$ there exists $i$ such that
            $i\in\initialSegment{n}$ and
            either $i \leq k + 1$ and $x\in \apply{V}{i}$
            or $k + 1 < i$ and $x\in \apply{U}{i}$.
        \begin{subproof}
            Fix $x\in X$.
            $x\in W$ or $x\notin W$.
            \begin{byCase}
            \caseOf{$x\in W$.}
                Take $i$ such that $i = k + 1$.
                Hence $i\in\initialSegment{n}$ and $i \leq k + 1$ and $x\in \apply{V}{i}$ by
                    \cref{real_naturals_closed_succ,initial_segment_naturalsplus,singleton_intro,pair_eq_iff,union_iff,function_apply_intro}.
            \caseOf{$x\notin W$.}
                We have $x\in D$ by \cref{subseteq,setminus_intro}.
                Take $Y$ such that $Y\in F$ and $x\in Y$ by \cref{unions_iff}.
                $Y = A$ or $Y = B$ by \cref{upair_iff}.
                \begin{byCase}
                \caseOf{$Y = A$.}
                    $x\in A$.
                    Take $W_0\in\img{V_0}{\initialSegment{k}}$ such that $x\in W_0$ by \cref{unions_iff}.
                    Take $i$ such that $i\in\initialSegment{k}$ and $(i,W_0)\in V_0$ by \cref{img_iff}.
                    $i\in\naturalsPlus$ and $i \leq k$ by \cref{initial_segment_naturalsplus}.
                    Hence $i\in\initialSegment{n}$ and $i \leq k + 1$ and $x\in \apply{V}{i}$ by
                        \cref{real_naturals_subset_reals,subseteq,real_leq_trans,initial_segment_naturalsplus,union_iff,function_apply_intro}.
                \caseOf{$Y = B$.}
                    $x\in B$.
                    Take $W_0\in\img{U}{J}$ such that $x\in W_0$ by \cref{unions_iff}.
                    Take $i$ such that $i\in J$ and $(i,W_0)\in U$ by \cref{img_iff}.
                    Hence $i\in\initialSegment{n}$ and $k + 1 < i$ and $x\in \apply{U}{i}$
                        by \cref{finite_shrink_tail_indices,function_apply_intro}.
                \end{byCase}
            \end{byCase}
        \end{subproof}
        Therefore $\finiteShrinkCoverCondition{V}{U}{X}{n}{k + 1}$ by \cref{finite_shrink_cover_condition}.
        Follows by \cref{pair_eq_iff}.
    \end{subproof}
    Hence $\finiteShrinkStage{U}{X}{T}{n}{k + 1}$ by \cref{finite_shrink_stage}.
\end{proof}

\begin{lemma}\label{finite_shrinking_cover}
    Assume $X$ is normal with respect to $T$.
    Assume $n\in\naturalsPlus$.
    Assume $U$ is a function on $\initialSegment{n}$.
    Suppose for all $i\in\initialSegment{n}$ we have $\apply{U}{i}$ is open in $X$ with respect to $T$.
    Suppose $\img{U}{\initialSegment{n}}$ is a covering of $X$.
    Then there exists $V$ such that
        $V$ is a function on $\initialSegment{n}$ and
        for all $i\in\initialSegment{n}$ we have
            $\apply{V}{i}$ is open in $X$ with respect to $T$ and
            $\closure{\apply{V}{i}}{X}{T}\subseteq \apply{U}{i}$ and
        $\img{V}{\initialSegment{n}}$ is a covering of $X$.
\end{lemma}
\begin{proof}
    Let $S = \finiteShrinkStageSet{U}{X}{T}{n}$.
    $S\subseteq\naturalsPlus$ by \cref{finite_shrink_stage_set,subseteq}.
    We have $1\in S$ by \cref{finite_shrink_stage_one,finite_shrink_stage_set,real_one_in_naturals}.
    For all $k\in S$ we have $k + 1\in S$
        by \cref{finite_shrink_stage_set,finite_shrink_stage_succ,real_naturals_closed_succ}.
    Hence $\finiteShrinkStage{U}{X}{T}{n}{n}$
        by \cref{real_naturals_induction,subseteq,finite_shrink_stage_set}.
    $n \leq n$.
    Take $V$ such that
        $V$ is a function on $\initialSegment{n}$ and
        for all $i\in\initialSegment{n}$ we have
            $\apply{V}{i}$ is open in $X$ with respect to $T$ and
            $\closure{\apply{V}{i}}{X}{T}\subseteq \apply{U}{i}$ and
            $\finiteShrinkCoverCondition{V}{U}{X}{n}{n}$
        by \cref{finite_shrink_stage}.
    Hence $\img{V}{\initialSegment{n}}$ is a covering of $X$
        by \cref{initial_segment_naturalsplus,finite_shrink_cover_condition_final}.
    Follows by \cref{pair_eq_iff}.
\end{proof}

% from §36 helper lemma 6: continuous finite sums of real-valued indexed families
\begin{lemma}\label{finite_real_sum_continuous}
    Assume $n\in\naturalsPlus$.
    Assume $f$ is a function on $\initialSegment{n}$.
    Suppose for all $i\in\initialSegment{n}$ we have
        $\apply{f}{i}$ is continuous from $X$ to $\reals$ with respect to $T$ and $\standardTopologyR$.
    Then there exists $g$ such that
        $g$ is a finite real sum witness for $f$ on $X$ with respect to $T$ and $n$.
\end{lemma}
\begin{proof}
    Let $A = \{k\in\naturalsPlus \mid \text{if $k \leq n$ then there exists $g$ such that
        $g$ is a finite real sum witness for $f$ on $X$ with respect to $T$ and $k$}\}$.
    Show that if $1 \leq n$ then there exists $g$ such that
        $g$ is a finite real sum witness for $f$ on $X$ with respect to $T$ and $1$.
    \begin{subproof}
        Assume $1 \leq n$.
        Let $g = \apply{f}{1}$.
        We have $g$ is continuous from $X$ to $\reals$ with respect to $T$ and $\standardTopologyR$ by \cref{real_one_in_naturals,initial_segment_naturalsplus,subseteq}.
        Show that for all $x\in X$ there exists $a$ such that
            $a\in\funs{\initialSegment{1}}{\reals}$ and
            $\sumFiniteSequence{a}{g(x)}$ and
            for all $i\in\initialSegment{1}$ we have $a(i) = \apply{\apply{f}{i}}{x}$.
        \begin{subproof}
            Fix $x\in X$.
            Let $a = \{(1,\apply{g}{x})\}$.
            $a$ is a function by \cref{singleton_elim,pair_eq_iff,rightunique,relation}.
            We have $\dom{a} = \initialSegment{1}$ by
                \cref{dom_iff,singleton_elim,pair_eq_iff,real_one_in_naturals,initial_segment_naturalsplus,real_naturals_subset_reals,subseteq,real_one_leq_naturalsplus,real_leq_antisym,singleton_intro,dom_intro,subseteq_antisymmetric}.
            For all $u\in\dom{a}$ we have $a(u)\in\reals$
                by \cref{dom_iff,singleton_elim,pair_eq_iff,singleton_intro,function_apply_intro,subseteq,continuous,funs_type_apply}.
            $a\in\funs{\initialSegment{1}}{\reals}$ by \cref{funs_intro}.
            For all $i\in\initialSegment{1}$ we have $a(i) = \apply{\apply{f}{i}}{x}$
                by \cref{initial_segment_naturalsplus,real_naturals_subset_reals,subseteq,real_one_in_naturals,real_one_leq_naturalsplus,real_leq_antisym,singleton_intro,pair_eq_iff,function_apply_intro,real_lt_subst_left}.
            $\sumFiniteSequence{a}{g(x)}$
                by \cref{initial_segment_naturalsplus,real_one_in_naturals,funs_type_apply,real_one_leq_naturalsplus,real_naturals_subset_reals,subseteq,real_add_closed,real_leq_add,real_add_ident,real_zero_lt_one,real_lt_add,real_lt_subst_left,real_lt_trans,real_lt_subst_right,real_lt_irreflexive,sum_finite_sequence_pred}.
        \end{subproof}
        Follows by \cref{finite_real_sum_witness}.
    \end{subproof}
    Hence $1\in A$ by \cref{real_one_in_naturals}.
    Show that for all $k\in A$ we have $k + 1\in A$.
    \begin{subproof}
        Fix $k$ such that $k\in A$.
        Show that if $k + 1 \leq n$ then there exists $g$ such that
            $g$ is a finite real sum witness for $f$ on $X$ with respect to $T$ and $k + 1$.
        \begin{subproof}
            Assume $k + 1 \leq n$.
            $k\in\naturalsPlus$ by \cref{subseteq}.
            $k\in\reals$ and $k + 1\in\reals$ and $n\in\reals$
                by \cref{real_naturals_subset_reals,real_naturals_closed_succ,subseteq,real_add_closed,real_one_in_naturals}.
            We have $k < k + 1$ by \cref{real_naturals_lt_succ_local}.
            Hence $k \leq n$ by \cref{real_lt_imp_leq,real_leq_trans}.
            Take $g_0$ such that
                $g_0$ is a finite real sum witness for $f$ on $X$ with respect to $T$ and $k$
                by \cref{subseteq}.
            We have $g_0$ is continuous from $X$ to $\reals$ with respect to $T$ and $\standardTopologyR$ and
                for all $x\in X$ there exists $a$ such that
                    $a\in\funs{\initialSegment{k}}{\reals}$ and
                    $\sumFiniteSequence{a}{g_0(x)}$ and
                    for all $i\in\initialSegment{k}$ we have $a(i) = \apply{\apply{f}{i}}{x}$
                by \cref{finite_real_sum_witness}.
            $k + 1\in\initialSegment{n}$ by \cref{real_naturals_closed_succ,initial_segment_naturalsplus}.
            Let $g_1 = \apply{f}{k + 1}$.
            We have $g_1$ is continuous from $X$ to $\reals$ with respect to $T$ and $\standardTopologyR$ by \cref{subseteq}.
            Let $g = \{(x, \apply{g_0}{x} + \apply{g_1}{x}) \mid x\in X\}$.
            Therefore $g$ is continuous from $X$ to $\reals$ with respect to $T$ and $\standardTopologyR$
                by \cref{real_function_addition_continuous}.
            Show that for all $x\in X$ there exists $a$ such that
                $a\in\funs{\initialSegment{k + 1}}{\reals}$ and
                $\sumFiniteSequence{a}{g(x)}$ and
                for all $i\in\initialSegment{k + 1}$ we have $a(i) = \apply{\apply{f}{i}}{x}$.
            \begin{subproof}
                Fix $x\in X$.
                Take $a_0$ such that
                    $a_0\in\funs{\initialSegment{k}}{\reals}$ and
                    $\sumFiniteSequence{a_0}{g_0(x)}$ and
                    for all $i\in\initialSegment{k}$ we have $a_0(i) = \apply{\apply{f}{i}}{x}$.
                Let $b = \{(k + 1,\apply{g_1}{x})\}$.
                $b$ is a function by \cref{singleton_elim,pair_eq_iff,rightunique,relation}.
                For all $u$ such that $u\in\dom{a_0}\inter\dom{b}$ we have $\apply{a_0}{u} = \apply{b}{u}$
                    by \cref{inter_elim_left,inter_elim_right,funs_elim,dom_iff,singleton_elim,pair_eq_iff,initial_segment_naturalsplus,real_naturals_lt_succ_local,real_lt_subst_right,real_lt_trans,real_lt_irreflexive}.
                We have $a_0$ is right-unique and $a_0$ is a relation
                    by \cref{funs_is_function,function_rightunique,funs_is_relation}.
                $b$ is right-unique and $b$ is a relation by \cref{rightunique,relation}.
                Let $a = a_0\union b$.
                $a$ is right-unique and $a$ is a relation by \cref{union_compatible_functions}.
                Therefore $a$ is a function by \cref{relation,rightunique}.
                $\dom{a} = \dom{a_0}\union\dom{b}$ by \cref{dom_union}.
                Hence $\dom{a} = \initialSegment{k}\union\dom{b}$ by \cref{funs_elim}.
                We have $\dom{b} = \{k + 1\}$
                    by \cref{dom_iff,singleton_elim,pair_eq_iff,singleton_intro,dom_intro,subseteq,subseteq_antisymmetric}.
                $\dom{a} = \initialSegment{k}\union \{k + 1\}$.
                For all $u$ such that $u\in\initialSegment{k}\union\{k + 1\}$ we have $u\in\initialSegment{k + 1}$
                    by \cref{union_iff,initial_segment_naturalsplus,real_naturals_subset_reals,subseteq,real_one_in_naturals,real_add_ident,real_add_closed,real_zero_lt_one,real_lt_add,real_add_comm,real_lt_subst_left,real_lt_trans,singleton_elim}.
                For all $u$ such that $u\in\initialSegment{k + 1}$ we have $u\in\initialSegment{k}\union\{k + 1\}$
                    by \cref{initial_segment_naturalsplus,real_naturals_leq_succ_elim,union_iff,singleton_intro}.
                Hence $\initialSegment{k}\union\{k + 1\} = \initialSegment{k + 1}$ by \cref{subseteq,subseteq_antisymmetric}.
                Therefore $\dom{a} = \initialSegment{k + 1}$.
                For all $u\in\dom{a}$ we have $a(u)\in\reals$
                    by \cref{dom_union,union_iff,function_apply_elim,function_apply_intro,funs_elim,funs_type_apply,subseteq,singleton_elim,singleton_intro,pair_eq_iff,continuous}.
                We have $a\in\funs{\initialSegment{k + 1}}{\reals}$ by \cref{funs_intro}.
                For all $i\in\initialSegment{k + 1}$ we have $a(i) = \apply{\apply{f}{i}}{x}$
                    by \cref{singleton_intro,pair_eq_iff,union_iff,function_apply_intro,real_lt_subst_left,initial_segment_naturalsplus,real_naturals_leq_succ_elim,funs_elim,subseteq,function_apply_elim}.
                Take $s\in\reals$ such that $\sumFiniteSequence{a}{s}$ by \cref{real_naturals_closed_succ,sum_finite_sequence_exists}.
                Let $a_1 = \restrl{a}{\initialSegment{k}}$.
                We have $a_1\in\funs{\initialSegment{k}}{\reals}$ and
                    there exists $s_0\in\reals$ such that $\sumFiniteSequence{a_1}{s_0}$ and $s = s_0 + a(k + 1)$
                    by \cref{sum_finite_sequence_split_last}.
                Take $s_0\in\reals$ such that $\sumFiniteSequence{a_1}{s_0}$ and $s = s_0 + a(k + 1)$.
                For all $w$ such that $w\in a_1$ we have $w\in a_0$
                    by \cref{restrl_subseteq,subseteq,relation,pair_eq_iff,dom_intro,restrl_dom,union_iff,singleton_elim,initial_segment_naturalsplus,real_naturals_lt_succ_local,real_lt_subst_right,real_lt_trans,real_lt_irreflexive}.
                $a_1\subseteq a_0$ by \cref{subseteq}.
                For all $w$ such that $w\in a_0$ we have $w\in a_1$
                    by \cref{relation,pair_eq_iff,dom_iff,funs_elim,union_iff,restrl_iff,subseteq}.
                $a_0\subseteq a_1$ by \cref{subseteq}.
                Hence $a_1 = a_0$ by \cref{subseteq_antisymmetric}.
                $s_0 = g_0(x)$ by \cref{sum_finite_sequence_unique}.
                Therefore $a(k + 1) = \apply{g_1}{x}$ by \cref{singleton_intro,pair_eq_iff,union_iff,function_apply_intro}.
                Hence $s = g_0(x) + g_1(x)$ by \cref{real_add_eq_subst}.
                Thus $g(x) = g_0(x) + g_1(x)$ by \cref{continuous,funs_is_function,function_apply_intro}.
                We have $\sumFiniteSequence{a}{g(x)}$ by \cref{real_lt_subst_right}.
            \end{subproof}
            Follows by \cref{finite_real_sum_witness}.
        \end{subproof}
        Therefore $k + 1\in A$ by \cref{real_naturals_closed_succ}.
    \end{subproof}
    Hence $n\in A$ by \cref{subseteq,real_naturals_induction}.
    $n \leq n$.
    Therefore there exists $g$ such that
        $g$ is a finite real sum witness for $f$ on $X$ with respect to $T$ and $n$
        by \cref{subseteq}.
\end{proof}

% from §36 helper lemma 7: vanishing outside an open set bounds the support
\begin{lemma}\label{support_subseteq_from_zero_outside}
    Assume $V$ is open in $X$ with respect to $T$.
    Assume $\closure{V}{X}{T}\subseteq U$.
    Suppose $f$ is continuous from $X$ to $\reals$ with respect to $T$ and $\standardTopologyR$.
    Suppose for all $x\in X\setminus V$ we have $f(x) = 0$.
    Then $\supp{f}{X}{T}\subseteq U$.
\end{lemma}
\begin{proof}
    Let $S = X\setminus\preimg{f}{\{0\}}$.
    For all $x$ such that $x\in S$ we have $x\in V$
        by \cref{setminus,setminus_intro,continuous,funs_is_function,funs,function_apply_elim,real_lt_subst_right,preimg_iff,singleton_intro,subseteq}.
    We have $S\subseteq\closure{V}{X}{T}$ by \cref{subseteq,open_in,topology_on,subseteq,pow_iff,subseteq_closure,subseteq_transitive}.
    Therefore $\closure{S}{X}{T}\subseteq\closure{V}{X}{T}$ by
        \cref{setminus_subseteq,open_in,topology_on,subseteq,pow_iff,closure_closed_in,closure_subseteq_closed}.
    Hence $\supp{f}{X}{T}\subseteq U$ by \cref{support,subseteq_transitive}.
\end{proof}

% from §36 helper lemma 8: reciprocal of a real-valued map bounded below by one
\begin{lemma}\label{reciprocal_geq_one_unit_interval_continuous}
    Let $I_0 = \intervalclosed{0}{1}$.
    Let $T_0 = \subspaceTopology{\standardTopologyR}{I_0}$.
    Suppose $g$ is continuous from $X$ to $\reals$ with respect to $T$ and $\standardTopologyR$.
    Suppose for all $x\in X$ we have $1 \leq g(x)$.
    Let $c = \{(x,1) \mid x\in X\}$.
    Let $q = \{(x, c(x) \rmul \inv{(g(x))}) \mid x\in X\}$.
    Then $q$ is continuous from $X$ to $I_0$ with respect to $T$ and $T_0$.
\end{lemma}
\begin{proof}
    $1\in\reals$ by \cref{real_one_in_naturals,real_naturals_subset_reals,subseteq}.
    $T$ is a topology on $X$ by \cref{continuous}.
    $\standardTopologyR$ is a topology on $\reals$
        by \cref{standard_basis_is_basis,basis_topology_is_topology,standard_topology_reals}.
    We have for all $x\in X$ $c(x) = 1$
        by \cref{constant_map_function,funs_is_function,function_apply_intro}.
    $c$ is continuous from $X$ to $\reals$ with respect to $T$ and $\standardTopologyR$
        by \cref{constant_map_function,constant_continuous}.
    For all $x\in X$ we have $g(x)\neq 0$
        by \cref{continuous,funs_type_apply,real_one_in_naturals,real_naturals_subset_reals,subseteq,real_add_ident,real_zero_lt_one,real_lt_imp_leq,real_leq_antisym,real_lt_irreflexive,real_lt_subst_left}.
    Therefore $q$ is continuous from $X$ to $\reals$ with respect to $T$ and $\standardTopologyR$
        by \cref{real_function_division_continuous}.
    For all $x\in X$ we have $q(x)\in I_0$
        by \cref{continuous,funs_type_apply,real_leq_split,real_add_ident,real_one_in_naturals,real_naturals_subset_reals,subseteq,real_zero_lt_one,real_lt_trans,real_lt_subst_right,real_mul_inverse,real_inv_pos,funs_is_function,function_apply_intro,real_mul_ident,real_lt_subst_left,real_lt_imp_leq,real_lt_mul_pos,real_mul_comm,real_pos_nonzero,intervalclosed}.
    Hence $q$ is continuous from $X$ to $I_0$ with respect to $T$ and $T_0$
        by \cref{subseteq,img_iff,continuous,funs_is_function,function_apply_iff,intervalclosed,continuous_subspace_range}.
\end{proof}

% from §36 Item 4: existence of finite partitions of unity
\begin{theorem}\label{finite_partition_of_unity_existence}
    Assume $X$ is normal with respect to $T$.
    Assume $n\in\naturalsPlus$.
    Assume $U$ is a function on $\initialSegment{n}$.
    Suppose for all $i\in\initialSegment{n}$ we have $\apply{U}{i}$ is open in $X$ with respect to $T$.
    Suppose $\img{U}{\initialSegment{n}}$ is a covering of $X$.
    Then there exists $f$ such that
        $f$ is a partition of unity on $X$ dominated by $U$ with respect to $T$ and $n$.
\end{theorem}
\begin{proof}
    Let $I_0 = \intervalclosed{0}{1}$.
    Let $T_0 = \subspaceTopology{\standardTopologyR}{I_0}$.
    Take $V$ such that
        $V$ is a function on $\initialSegment{n}$ and
        for all $i\in\initialSegment{n}$ we have
            $\apply{V}{i}$ is open in $X$ with respect to $T$ and
            $\closure{\apply{V}{i}}{X}{T}\subseteq \apply{U}{i}$ and
        $\img{V}{\initialSegment{n}}$ is a covering of $X$
        by \cref{finite_shrinking_cover}.
    Take $W$ such that
        $W$ is a function on $\initialSegment{n}$ and
        for all $i\in\initialSegment{n}$ we have
            $\apply{W}{i}$ is open in $X$ with respect to $T$ and
            $\closure{\apply{W}{i}}{X}{T}\subseteq \apply{V}{i}$ and
        $\img{W}{\initialSegment{n}}$ is a covering of $X$
        by \cref{finite_shrinking_cover}.

    Let $B = \pow{X\times I_0}$.
    Let $R = \{w\in\initialSegment{n}\times B \mid \text{there exists $i\in\initialSegment{n}$ such that there exists $u\in B$ such that
        $w = (i,u)$ and
        $u$ is continuous from $X$ to $I_0$ with respect to $T$ and $T_0$ and
        $\img{u}{\closure{\apply{W}{i}}{X}{T}}\subseteq \{1\}$ and
        $\img{u}{X\setminus \apply{V}{i}}\subseteq \{0\}$}\}$.
    Show that for all $i\in\initialSegment{n}$ there exists $u$ such that $i\mathrel{R} u$.
    \begin{subproof}
        Fix $i\in\initialSegment{n}$.
        $\apply{W}{i}$ is open in $X$ with respect to $T$ and
            $\closure{\apply{W}{i}}{X}{T}\subseteq \apply{V}{i}$ by assumption.
        We have $\apply{W}{i}\subseteq X$ by \cref{open_in,normal_space,topology_on,subseteq,pow_iff}.
        $T$ is a topology on $X$ by \cref{normal_space}.
        $\apply{V}{i}$ is open in $X$ with respect to $T$ by assumption.
        Hence $X\setminus \apply{V}{i}$ is closed in $X$ with respect to $T$
            by \cref{open_in,normal_space,topology_on,subseteq,pow_iff,closed_in,setminus_subseteq,double_relative_complement}.
        $\closure{\apply{W}{i}}{X}{T}$ is closed in $X$ with respect to $T$
            by \cref{closure_closed_in}.
        $X\setminus \apply{V}{i}$ is disjoint from $\closure{\apply{W}{i}}{X}{T}$
            by \cref{subseteq,setminus_elim_right,disjoint}.
        Take $u$ such that
            $u$ is continuous from $X$ to $I_0$ with respect to $T$ and $T_0$ and
            $\img{u}{X\setminus \apply{V}{i}}\subseteq \{0\}$ and
            $\img{u}{\closure{\apply{W}{i}}{X}{T}}\subseteq \{1\}$
            by \cref{urysohn_unit_interval}.
        Hence $i\mathrel{R} u$
            by \cref{continuous,funs_is_relation,funs,funs_subseteq_rels,subseteq,rels_ran_subseteq,rels_elim,powerset_intro,times_tuple_intro,pair_eq_iff}.
    \end{subproof}
    $R$ is a relation by \cref{relation}.
    Take $c$ such that $c$ is a function on $\initialSegment{n}$ and for all $i\in\initialSegment{n}$ we have $i\mathrel{R} c(i)$
        by \cref{choice_function}.
    For all $i\in\initialSegment{n}$ we have
        $\apply{c}{i}$ is continuous from $X$ to $I_0$ with respect to $T$ and $T_0$ and
        $\img{\apply{c}{i}}{\closure{\apply{W}{i}}{X}{T}}\subseteq \{1\}$ and
        $\img{\apply{c}{i}}{X\setminus \apply{V}{i}}\subseteq \{0\}$
        by \cref{choice_function,pair_eq_iff}.

    Let $j = \identity{I_0}$.
    For all $i\in\initialSegment{n}$ we have
        $j\circ \apply{c}{i}$ is continuous from $X$ to $\reals$ with respect to $T$ and $\standardTopologyR$
        by \cref{continuous_interval_expand_range}.
    Let $p = \{(i,j\circ \apply{c}{i}) \mid i\in\initialSegment{n}\}$.
    For all $i,u_1,u_2$ such that $(i,u_1)\in p$ and $(i,u_2)\in p$ we have $u_1 = u_2$ by \cref{pair_eq_iff}.
    We have $p$ is a function by \cref{relation,pair_eq_iff,rightunique}.
    For all $i$ such that $i\in\dom{p}$ we have $i\in\initialSegment{n}$ by \cref{dom_iff,pair_eq_iff}.
    For all $i$ such that $i\in\initialSegment{n}$ we have $i\in\dom{p}$ by \cref{pair_eq_iff,dom_intro}.
    Therefore $\dom{p} = \initialSegment{n}$ by \cref{subseteq,subseteq_antisymmetric}.
    Hence $p$ is a function on $\initialSegment{n}$ by \cref{funs}.
    For all $i\in\initialSegment{n}$ we have $\apply{p}{i} = j\circ \apply{c}{i}$ by \cref{pair_eq_iff,function_apply_intro}.
    For all $i\in\initialSegment{n}$ we have
        $\apply{p}{i}$ is continuous from $X$ to $\reals$ with respect to $T$ and $\standardTopologyR$
        by \cref{continuous_eq_subst}.
    For all $i\in\initialSegment{n}$ for all $x\in X$ we have $\apply{\apply{p}{i}}{x} = \apply{\apply{c}{i}}{x}$
        by \cref{continuous,funs_is_function,funs,funs_is_relation,function_rightunique,funs_subseteq_rels,subseteq,rels_ran_subseteq,id_dom,id_rightunique,id_is_relation,circ_apply,funs_type_apply,id_apply}.
    Take $g$ such that
        $g$ is a finite real sum witness for $p$ on $X$ with respect to $T$ and $n$
        by \cref{finite_real_sum_continuous}.
    Therefore $g$ is continuous from $X$ to $\reals$ with respect to $T$ and $\standardTopologyR$ and
        for all $x\in X$ there exists $a$ such that
            $a\in\funs{\initialSegment{n}}{\reals}$ and
            $\sumFiniteSequence{a}{g(x)}$ and
            for all $i\in\initialSegment{n}$ we have $a(i) = \apply{\apply{p}{i}}{x}$
        by \cref{finite_real_sum_witness}.

    Show that for all $x\in X$ we have $1 \leq g(x)$.
    \begin{subproof}
        Fix $x\in X$.
        $1\in\initialSegment{n}$ by \cref{real_one_leq_naturalsplus,real_one_in_naturals,initial_segment_naturalsplus}.
        Take $Z\in\img{W}{\initialSegment{n}}$ such that $x\in Z$ by \cref{covering,unions_iff,subseteq}.
        Take $i$ such that $i\in\initialSegment{n}$ and $(i,Z)\in W$ by \cref{img_iff}.
        We have $x\in \apply{W}{i}$ by \cref{function_apply_intro}.
        Take $a$ such that
            $a\in\funs{\initialSegment{n}}{\reals}$ and
            $\sumFiniteSequence{a}{g(x)}$ and
            for all $k\in\initialSegment{n}$ we have $a(k) = \apply{\apply{p}{k}}{x}$.
        For all $k\in\dom{a}$ we have $0 \leq a(k)$
            by \cref{funs_elim,continuous,funs_is_function,funs,subseteq,continuous,funs_type_apply,intervalclosed}.
        Therefore $a(i) \leq g(x)$ by \cref{funs_elim,subseteq,sum_finite_sequence_term_leq}.
        Hence $x\in \closure{\apply{W}{i}}{X}{T}$ by \cref{open_in,normal_space,topology_on,subseteq,pow_iff,subseteq_closure}.
        $\apply{c}{i}\in\funs{X}{I_0}$ and $\apply{c}{i}$ is a function and $\dom{\apply{c}{i}} = X$
            by \cref{continuous,funs_is_function,funs}.
        Therefore $\apply{\apply{c}{i}}{x}\in\img{\apply{c}{i}}{\closure{\apply{W}{i}}{X}{T}}$
            by \cref{subseteq,inter_intro,img_of_function_intro}.
        Thus $\apply{\apply{c}{i}}{x} = 1$ by \cref{subseteq,singleton_elim}.
        We have $a(i) = 1$ by \cref{funs_elim,subseteq,real_lt_subst_right}.
        $1\in\reals$ by \cref{real_one_in_naturals,real_naturals_subset_reals,subseteq}.
        $1 \leq g(x)$.
    \end{subproof}
    For all $x\in X$ we have $g(x)\neq 0$
        by \cref{continuous,funs_type_apply,real_one_in_naturals,real_naturals_subset_reals,subseteq,real_add_ident,real_zero_lt_one,real_lt_imp_leq,real_leq_antisym,real_lt_irreflexive,real_lt_subst_left}.

    Let $d = \{(x,1) \mid x\in X\}$.
    Let $q = \{(x, d(x) \rmul \inv{(g(x))}) \mid x\in X\}$.
    Therefore $q$ is continuous from $X$ to $I_0$ with respect to $T$ and $T_0$
        by \cref{reciprocal_geq_one_unit_interval_continuous}.
    Let $q_1 = j\circ q$.
    Hence $q_1$ is continuous from $X$ to $\reals$ with respect to $T$ and $\standardTopologyR$
        by \cref{continuous_interval_expand_range}.
    For all $x\in X$ we have $q_1(x) = q(x)$
        by \cref{continuous,funs_is_function,funs,function_rightunique,funs_is_relation,funs_subseteq_rels,subseteq,rels_ran_subseteq,id_dom,id_rightunique,id_is_relation,circ_apply,funs_type_apply,id_apply}.
    $1\in\reals$ by \cref{real_one_in_naturals,real_naturals_subset_reals,subseteq}.
    $d$ is a function from $X$ to $\reals$ by \cref{constant_map_function}.
    We have for all $x\in X$ $d(x) = 1$
        by \cref{constant_map_function,funs_is_function,function_apply_intro}.

    Let $P = \{w\in\initialSegment{n}\times B \mid \text{there exists $i\in\initialSegment{n}$ such that there exists $u\in B$ such that
        $w = (i,u)$ and
        $u$ is continuous from $X$ to $I_0$ with respect to $T$ and $T_0$ and
        $\supp{u}{X}{T}\subseteq \apply{U}{i}$ and
        for all $x\in X$ we have $u(x) = \apply{q_1}{x} \rmul \apply{\apply{c}{i}}{x}$}\}$.
    Show that for all $i\in\initialSegment{n}$ there exists $u$ such that $i\mathrel{P} u$.
    \begin{subproof}
        Fix $i\in\initialSegment{n}$.
        Let $c_i = \apply{c}{i}$.
        We have $c_i$ is continuous from $X$ to $I_0$ with respect to $T$ and $T_0$ and
            $\img{c_i}{\closure{\apply{W}{i}}{X}{T}}\subseteq \{1\}$ and
            $\img{c_i}{X\setminus \apply{V}{i}}\subseteq \{0\}$ by assumption.
        Let $u = \{(x,\apply{q_1}{x} \rmul \apply{c_i}{x}) \mid x\in X\}$.
        Therefore $u$ is continuous from $X$ to $\reals$ with respect to $T$ and $\standardTopologyR$
            by \cref{real_function_multiplication_continuous}.
        $T$ is a topology on $X$ by \cref{normal_space}.
        Let $v = \{(x,\apply{q}{x} \rmul \apply{c_i}{x}) \mid x\in X\}$.
        Hence $v$ is continuous from $X$ to $I_0$ with respect to $T$ and $T_0$
            by \cref{unit_interval_product_continuous}.

        For all $x\in X$ we have $u(x) = v(x)$
            by \cref{pair_eq_iff,continuous,funs_is_function,funs,function_apply_intro,real_mul_eq_subst}.
        We have for all $x\in X$ we have $u(x)\in I_0$ by \cref{continuous,funs_type_apply}.
        $\img{u}{X}\subseteq I_0$
            by \cref{subseteq,img_iff,continuous,funs_is_function,funs,function_apply_intro}.
        $u\in\funs{X}{\reals}$ by \cref{continuous}.
        Therefore $u$ is continuous from $X$ to $I_0$ with respect to $T$ and $T_0$
            by \cref{intervalclosed,subseteq,continuous_subspace_range}.

        Show that for all $x\in X\setminus \apply{V}{i}$ we have $u(x) = 0$.
        \begin{subproof}
            Fix $x\in X\setminus \apply{V}{i}$.
            $c_i\in\funs{X}{I_0}$ and $c_i$ is a function and $\dom{c_i} = X$
                by \cref{continuous,funs_is_function,funs}.
            We have $x\in\dom{c_i}\inter(X\setminus \apply{V}{i})$ by \cref{setminus_elim_left,subseteq,inter_intro}.
            $c_i(x)\in\img{c_i}{X\setminus \apply{V}{i}}$ by \cref{img_of_function_intro}.
            Therefore $c_i(x) = 0$ by \cref{subseteq,singleton_elim}.
            Hence $u(x) = \apply{q_1}{x} \rmul c_i(x)$
                by \cref{setminus_elim_left,pair_eq_iff,continuous,funs_is_function,funs,function_apply_intro}.
            Thus $u(x) = \apply{q_1}{x} \rmul 0$ by \cref{real_lt_subst_right}.
            $\apply{q_1}{x}\in\reals$ by \cref{setminus_elim_left,continuous,funs_type_apply}.
            $u(x) = 0$ by \cref{real_add_ident,real_mul_comm,real_mul_zero}.
        \end{subproof}
        Hence $\supp{u}{X}{T}\subseteq \apply{U}{i}$ by \cref{support_subseteq_from_zero_outside}.
        For all $x\in X$ we have $u(x) = \apply{q_1}{x} \rmul \apply{\apply{c}{i}}{x}$
            by \cref{pair_eq_iff,continuous,funs_is_function,funs,function_apply_intro,real_mul_eq_subst}.
        Therefore $i\mathrel{P} u$
            by \cref{continuous,funs_is_relation,funs,funs_subseteq_rels,subseteq,rels_ran_subseteq,rels_elim,powerset_intro,times_tuple_intro,pair_eq_iff}.
    \end{subproof}
    $P$ is a relation by \cref{relation}.
    Take $f$ such that $f$ is a function on $\initialSegment{n}$ and for all $i\in\initialSegment{n}$ we have $i\mathrel{P} f(i)$
        by \cref{choice_function}.
    Show that for all $i\in\initialSegment{n}$ we have
        $\apply{f}{i}$ is continuous from $X$ to $I_0$ with respect to $T$ and $T_0$ and
        $\supp{\apply{f}{i}}{X}{T}\subseteq \apply{U}{i}$.
    \begin{subproof}
        Follows by \cref{choice_function,pair_eq_iff}.
    \end{subproof}
    Therefore $\finitePartitionComponentCondition{f}{X}{U}{T}{n}$
        by \cref{finite_partition_component_condition,pair_eq_iff}.
    Show that for all $x\in X$ there exists $a$ such that
        $a\in\funs{\initialSegment{n}}{\reals}$ and
        for all $i\in\initialSegment{n}$ we have $a(i) = \apply{\apply{f}{i}}{x}$ and
        $\sumFiniteSequence{a}{1}$.
    \begin{subproof}
        Fix $x\in X$.
        Take $a_0$ such that
            $a_0\in\funs{\initialSegment{n}}{\reals}$ and
            $\sumFiniteSequence{a_0}{g(x)}$ and
            for all $i\in\initialSegment{n}$ we have $a_0(i) = \apply{\apply{p}{i}}{x}$.
        Let $a = \coordScalarSequence{n}{q_1(x)}{a_0}$.
        We have $a\in\funs{\initialSegment{n}}{\reals}$ by \cref{continuous,funs_type_apply,coord_scalar_sequence_function}.
        For all $i\in\initialSegment{n}$ we have $a(i) = \apply{\apply{f}{i}}{x}$
            by \cref{funs_elim,subseteq,funs_is_function,function_apply_elim,coord_scalar_sequence,function_apply_intro,real_mul_eq_subst,choice_function,pair_eq_iff,real_lt_subst_right}.
        $q_1(x)\in\reals$ by \cref{continuous,funs_type_apply}.
        $\sumFiniteSequence{a}{q_1(x) \rmul g(x)}$ by \cref{sum_finite_sequence_scalar}.
        $q_1(x) = q(x)$ by assumption.
        We have $g(x)\neq 0$ by assumption.
        $g(x)\in\reals$ by \cref{continuous,funs_type_apply}.
        $q$ is a function by \cref{continuous,funs_is_function}.
        We have $(x,d(x) \rmul \inv{(g(x))})\in q$.
        $q(x) = d(x) \rmul \inv{(g(x))}$ by \cref{function_apply_intro}.
        Hence $q_1(x) = 1 \rmul \inv{(g(x))}$ by \cref{real_lt_subst_left}.
        Therefore $q_1(x) = \inv{(g(x))}$ by \cref{real_mul_inverse,real_mul_ident}.
        Thus $q_1(x) \rmul g(x) = 1$
            by \cref{real_mul_eq_subst,real_mul_comm,real_mul_inverse,real_lt_subst_right}.
        We have $\sumFiniteSequence{a}{1}$ by \cref{real_lt_subst_right}.
        Therefore there exists $b$ such that
            $b\in\funs{\initialSegment{n}}{\reals}$ and
            for all $i\in\initialSegment{n}$ we have $b(i) = \apply{\apply{f}{i}}{x}$ and
            $\sumFiniteSequence{b}{1}$.
    \end{subproof}
    Hence $\finitePartitionSumCondition{f}{X}{n}$ by \cref{finite_partition_sum_condition}.
    $U$ is a function on $\initialSegment{n}$ by assumption.
    For all $i\in\initialSegment{n}$ we have $\apply{U}{i}$ is open in $X$ with respect to $T$ by assumption.
    $\img{U}{\initialSegment{n}}$ is a covering of $X$ by assumption.
    $f$ is a function on $\initialSegment{n}$ by \cref{choice_function}.
    Therefore $f$ is a partition of unity on $X$ dominated by $U$ with respect to $T$ and $n$
        by \cref{finite_partition_of_unity}.
\end{proof}

% from §36 helper lemma 11: nonempty compact manifold has a finite chart cover
\begin{lemma}\label{compact_manifold_finite_chart_cover}
    Assume $m\in\naturalsPlus$.
    Assume $X$ is inhabited.
    Assume $X$ is a manifold of dimension $m$ with respect to $T$.
    Assume $X$ is compact with respect to $T$.
    Let $R_0 = \{(j,\reals) \mid j\in\initialSegment{m}\}$.
    Let $S_0 = \{(j,\standardTopologyR) \mid j\in\initialSegment{m}\}$.
    Let $P_0 = \productTopologyIndexed{S_0}{R_0}$.
    Then there exist $n,U$ such that
        $n\in\naturalsPlus$ and
        $U$ is a function on $\initialSegment{n}$ and
        for all $i\in\initialSegment{n}$ there exist $V,h$ such that
            $\apply{U}{i}$ is open in $X$ with respect to $T$ and
            $V$ is open in $\productFamily{R_0}$ with respect to $P_0$ and
            $h$ is a homeomorphism between $\apply{U}{i}$ and $V$ with respect to
                $\subspaceTopology{T}{\apply{U}{i}}$ and $\subspaceTopology{P_0}{V}$ and
        $\img{U}{\initialSegment{n}}$ is a covering of $X$.
\end{lemma}
\begin{proof}
    $X$ is Hausdorff with respect to $T$ and $X$ satisfies the second countability axiom with respect to $T$
        by \cref{m_manifold}.
    We have $T$ is a topology on $X$ by \cref{hausdorff}.
    Let $C = \{U\in\pow{X} \mid \text{there exist $x\in X$ such that there exist $V,h$ such that
        $U$ is a neighborhood of $x$ in $X$ with respect to $T$ and
        $V$ is open in $\productFamily{R_0}$ with respect to $P_0$ and
        $h$ is a homeomorphism between $U$ and $V$ with respect to
            $\subspaceTopology{T}{U}$ and $\subspaceTopology{P_0}{V}$}\}$.
    Show that for all $x$ such that $x\in X$ we have $x\in\unions{C}$.
    \begin{subproof}
        Fix $x$ such that $x\in X$.
        Take $U,V,R,S,h$ such that
            $U$ is a neighborhood of $x$ in $X$ with respect to $T$ and
            $R = \{(i,\reals) \mid i\in\initialSegment{m}\}$ and
            $S = \{(i,\standardTopologyR) \mid i\in\initialSegment{m}\}$ and
            $V$ is open in $\productFamily{R}$ with respect to $\productTopologyIndexed{S}{R}$ and
            $h$ is a homeomorphism between $U$ and $V$ with respect to
                $\subspaceTopology{T}{U}$ and $\subspaceTopology{\productTopologyIndexed{S}{R}}{V}$
            by \cref{m_manifold}.
        For all $u$ such that $u\in R$ we have $u\in R_0$ by \cref{pair_eq_iff}.
        For all $u$ such that $u\in R_0$ we have $u\in R$ by \cref{pair_eq_iff}.
        Therefore $R = R_0$ by \cref{subseteq,subseteq_antisymmetric}.
        For all $u$ such that $u\in S$ we have $u\in S_0$ by \cref{pair_eq_iff}.
        For all $u$ such that $u\in S_0$ we have $u\in S$ by \cref{pair_eq_iff}.
        Thus $S = S_0$ by \cref{subseteq,subseteq_antisymmetric}.
        Hence $V$ is open in $\productFamily{R_0}$ with respect to $P_0$ and
            $h$ is a homeomorphism between $U$ and $V$ with respect to
                $\subspaceTopology{T}{U}$ and $\subspaceTopology{P_0}{V}$ by \cref{pair_eq_iff}.
        Therefore $x\in\unions{C}$ by \cref{neighborhood,open_in,topology_on,subseteq,pow_iff,pair_eq_iff,unions_iff}.
    \end{subproof}
    Therefore $C$ is a covering of $X$ by \cref{subseteq,covering}.
    For all $W$ such that $W\in C$ we have $W$ is open in $X$ with respect to $T$
        by \cref{powerset_elim,neighborhood}.
    Hence $C$ is an open covering of $X$ with respect to $T$ by \cref{open_covering}.
    Take $d$ such that
        $d$ is a sequence in $\pow{X}$ and
        for all $n\in\naturalsPlus$ we have $d(n)$ is open in $X$ with respect to $T$ and
        for all $x\in X$ we have
            for all $W$ such that $W$ is open in $X$ with respect to $T$ and $x\in W$
                there exists $n\in\naturalsPlus$ such that $x\in d(n)$ and $d(n)\subseteq W$
        by \cref{second_countability_axiom}.
    Let $J = \{n\in\naturalsPlus \mid \exists A\in C. d(n)\subseteq A\}$.
    Let $R_1 = \{(n,A) \mid n\in J, A\in C \mid d(n)\subseteq A\}$.
    $R_1$ is a relation by \cref{pair_eq_iff,relation}.
    For all $n\in J$ there exists $A$ such that $n\mathrel{R_1} A$ by \cref{pair_eq_iff}.
    Take $b$ such that $b$ is a function on $J$ and for all $n\in J$ we have $n\mathrel{R_1}\apply{b}{n}$
        by \cref{choice_function}.
    For all $j\in J$ we have $\apply{b}{j}\in C$ by \cref{choice_function,pair_eq_iff}.
    For all $x$ such that $x\in X$ we have $x\in\unions{\img{b}{J}}$
        by \cref{covering,subseteq,unions_iff,choice_function,pair_eq_iff,function_apply_iff,img_iff,unions_intro}.
    Therefore $\img{b}{J}$ is a covering of $X$ by \cref{subseteq,covering}.
    For all $W$ such that $W\in\img{b}{J}$ we have $W$ is open in $X$ with respect to $T$
        by \cref{img_iff,function_apply_intro,subseteq}.
    Hence $\img{b}{J}$ is an open covering of $X$ with respect to $T$ by \cref{open_covering}.
    Take $B$ such that
        $B\subseteq\img{b}{J}$ and
        $B$ is finite and
        $B$ is a covering of $X$ by \cref{compact}.
    Take $W_0\in B$ such that $W_0\in B$ by \cref{covering,subseteq,unions_iff}.
    Take $j_0\in J$ such that $(j_0,W_0)\in b$ by \cref{subseteq,img_iff}.
    We have $W_0 = \apply{b}{j_0}$ by \cref{function_apply_intro}.
    $W_0\subseteq X$ by \cref{subseteq,function_apply_intro,pair_eq_iff,powerset_elim}.
    Let $Q = \{w\in B\times\naturalsPlus \mid \exists W\in B.\exists j\in J. w = (W,j) \land \apply{b}{j} = W\}$.
    We have $Q$ is a relation by \cref{relation}.
    For all $W\in B$ there exists $j$ such that $W\mathrel{Q} j$
        by \cref{subseteq,img_iff,function_apply_intro,times_tuple_intro,pair_eq_iff}.
    Take $s$ such that $s$ is a function on $B$ and for all $W\in B$ we have $W\mathrel{Q}\apply{s}{W}$ by \cref{choice_function}.
    Let $F = \img{s}{B}$.
    $F$ is finite by \cref{finite_image_function}.
    We have $F\subseteq\naturalsPlus$ by \cref{img_iff,function_apply_intro,choice_function,times_tuple_elim,subseteq}.
    $F\neq\emptyset$ by \cref{subseteq,function_apply_elim,img_iff,inhabited_nonempty,nonempty_inhabited}.
    Hence $F\subseteq\reals$ by \cref{subseteq,real_naturals_subset_reals,subseteq_transitive}.
    Take $n\in F$ such that $n$ is a maximum of $F$ by \cref{finite_real_maximum}.
    Hence $F\subseteq\initialSegment{n}$ by \cref{subseteq,real_maximum,initial_segment_naturalsplus}.
    Let $U = \{(i,W) \mid i\in\initialSegment{n}, W\in\pow{X} \mid
        (i\in F \land W = \apply{b}{i}) \lor (i\notin F \land W = W_0)\}$.
    For all $i,W_1,W_2$ such that $(i,W_1)\in U$ and $(i,W_2)\in U$ we have $W_1 = W_2$ by \cref{pair_eq_iff}.
    We have $U$ is a function by \cref{rightunique,relation}.
    Show that for all $i$ such that $i\in\initialSegment{n}$ we have $i\in\dom{U}$.
    \begin{subproof}
        Fix $i$ such that $i\in\initialSegment{n}$.
        \begin{byCase}
        \caseOf{$i\in F$.}
            Take $W\in B$ such that $(W,i)\in s$ by \cref{img_iff}.
            We have $(W,i)\in Q$ by \cref{choice_function,function_apply_intro}.
            $i\in J$ by \cref{pair_eq_iff}.
            Hence $i\in\dom{U}$ by \cref{subseteq,powerset_elim,pair_eq_iff,dom_intro}.
        \caseOf{$i\notin F$.}
            We have $i\in\dom{U}$ by \cref{subseteq,img_iff,powerset_elim,pair_eq_iff,dom_intro}.
        \end{byCase}
    \end{subproof}
    Therefore $U$ is a function on $\initialSegment{n}$
        by \cref{dom_iff,pair_eq_iff,subseteq,subseteq_antisymmetric,funs}.
    Show that for all $i\in\initialSegment{n}$ there exist $V,h$ such that
        $\apply{U}{i}$ is open in $X$ with respect to $T$ and
        $V$ is open in $\productFamily{R_0}$ with respect to $P_0$ and
        $h$ is a homeomorphism between $\apply{U}{i}$ and $V$ with respect to
            $\subspaceTopology{T}{\apply{U}{i}}$ and $\subspaceTopology{P_0}{V}$.
    \begin{subproof}
        Fix $i\in\initialSegment{n}$.
        \begin{byCase}
        \caseOf{$i\in F$.}
            Take $W\in B$ such that $(W,i)\in s$ by \cref{img_iff}.
            We have $(W,i)\in Q$ by \cref{choice_function,function_apply_intro}.
            $i\in J$ by \cref{pair_eq_iff}.
            Hence $(i,\apply{U}{i})\in U$ and $(i,\apply{b}{i})\in b$ by \cref{subseteq,function_apply_elim}.
            Therefore $\apply{U}{i} = \apply{b}{i}$ by \cref{function_apply_intro}.
            Take $x\in X$ such that there exist $V,h$ such that
                $\apply{b}{i}$ is a neighborhood of $x$ in $X$ with respect to $T$ and
                $V$ is open in $\productFamily{R_0}$ with respect to $P_0$ and
                $h$ is a homeomorphism between $\apply{b}{i}$ and $V$ with respect to
                    $\subspaceTopology{T}{\apply{b}{i}}$ and $\subspaceTopology{P_0}{V}$ by \cref{subseteq,powerset_elim}.
            Take $V,h$ such that
                $\apply{b}{i}$ is a neighborhood of $x$ in $X$ with respect to $T$ and
                $V$ is open in $\productFamily{R_0}$ with respect to $P_0$ and
                $h$ is a homeomorphism between $\apply{b}{i}$ and $V$ with respect to
                    $\subspaceTopology{T}{\apply{b}{i}}$ and $\subspaceTopology{P_0}{V}$.
            Therefore there exist $V_2,h_2$ such that
                $\apply{U}{i}$ is open in $X$ with respect to $T$ and
                $V_2$ is open in $\productFamily{R_0}$ with respect to $P_0$ and
                $h_2$ is a homeomorphism between $\apply{U}{i}$ and $V_2$ with respect to
                    $\subspaceTopology{T}{\apply{U}{i}}$ and $\subspaceTopology{P_0}{V_2}$
                by \cref{neighborhood,open_in,real_lt_subst_right}.
        \caseOf{$i\notin F$.}
            We have $(i,\apply{U}{i})\in U$ by \cref{subseteq,function_apply_elim}.
            Hence $\apply{U}{i} = W_0$ by \cref{pair_eq_iff}.
            Take $x\in X$ such that there exist $V,h$ such that
                $\apply{b}{j_0}$ is a neighborhood of $x$ in $X$ with respect to $T$ and
                $V$ is open in $\productFamily{R_0}$ with respect to $P_0$ and
                $h$ is a homeomorphism between $\apply{b}{j_0}$ and $V$ with respect to
                    $\subspaceTopology{T}{\apply{b}{j_0}}$ and $\subspaceTopology{P_0}{V}$ by \cref{subseteq,powerset_elim}.
            Take $V,h$ such that
                $\apply{b}{j_0}$ is a neighborhood of $x$ in $X$ with respect to $T$ and
                $V$ is open in $\productFamily{R_0}$ with respect to $P_0$ and
                $h$ is a homeomorphism between $\apply{b}{j_0}$ and $V$ with respect to
                    $\subspaceTopology{T}{\apply{b}{j_0}}$ and $\subspaceTopology{P_0}{V}$.
            Therefore there exist $V_2,h_2$ such that
                $\apply{U}{i}$ is open in $X$ with respect to $T$ and
                $V_2$ is open in $\productFamily{R_0}$ with respect to $P_0$ and
                $h_2$ is a homeomorphism between $\apply{U}{i}$ and $V_2$ with respect to
                    $\subspaceTopology{T}{\apply{U}{i}}$ and $\subspaceTopology{P_0}{V_2}$
                by \cref{pair_eq_iff,neighborhood,open_in}.
        \end{byCase}
    \end{subproof}
    Show that for all $x$ such that $x\in X$ we have $x\in\unions{\img{U}{\initialSegment{n}}}$.
    \begin{subproof}
        Follows by \cref{covering,subseteq,unions_iff,function_apply_elim,img_iff,choice_function,pair_eq_iff,function_apply_intro,real_lt_subst_right,img_of_function_intro,inter_intro}.
    \end{subproof}
    Hence $\img{U}{\initialSegment{n}}$ is a covering of $X$ by \cref{subseteq,covering}.
\end{proof}

% from §36 helper lemma 12: a fixed row in a finite grid is finite
\begin{lemma}\label{initial_segment_row_finite}
    Assume $t\in\naturalsPlus$.
    Assume $n\in\naturalsPlus$.
    Let $A = \{(t,i) \mid i\in\initialSegment{n}\}$.
    Then $A$ is finite.
\end{lemma}
\begin{proof}
    Let $f = \{(i,(t,i)) \mid i\in\initialSegment{n}\}$.
    For all $i,b_1,b_2$ such that $(i,b_1)\in f$ and $(i,b_2)\in f$ we have $b_1 = b_2$ by \cref{pair_eq_iff}.
    We have $f$ is a function on $\initialSegment{n}$
        by \cref{relation,pair_eq_iff,rightunique,dom_intro,dom_iff,subseteq,subseteq_antisymmetric,funs}.
    For all $a\in A$ there exists $i\in\initialSegment{n}$ such that $f(i) = a$
        by \cref{pair_eq_iff,function_apply_intro}.
    Hence $A\subseteq\img{f}{\initialSegment{n}}$ by \cref{subseteq,img_iff}.
    For all $a\in\img{f}{\initialSegment{n}}$ we have $a\in A$
        by \cref{img_iff,function_apply_intro,pair_eq_iff}.
    Hence $A = \img{f}{\initialSegment{n}}$ by \cref{subseteq,subseteq_antisymmetric}.
    Therefore $A$ is finite by \cref{initial_segment_finite,finite_image_function}.
\end{proof}

% from §36 helper lemma 13: finite grid of positive integer pairs
\begin{lemma}\label{initial_segment_grid_finite}
    Assume $m\in\naturalsPlus$.
    Assume $n\in\naturalsPlus$.
    Let $J = \initialSegment{m}\times\initialSegment{n}$.
    Then $J$ is finite.
\end{lemma}
\begin{proof}
    Let $A = \{k\in\naturalsPlus \mid \text{$\initialSegment{k}\times\initialSegment{n}$ is finite}\}$.
    $A\subseteq\naturalsPlus$ by \cref{subseteq}.
    Let $B = \{(1,i) \mid i\in\initialSegment{n}\}$.
    $B$ is finite by \cref{real_one_in_naturals,initial_segment_row_finite}.
    For all $u$ such that $u\in\initialSegment{1}\times\initialSegment{n}$ we have $u\in B$
        by \cref{times_elem_is_tuple,initial_segment_naturalsplus,real_naturals_subset_reals,subseteq,real_one_in_naturals,real_one_leq_naturalsplus,real_leq_antisym,pair_eq_iff}.
    Hence $\initialSegment{1}\times\initialSegment{n}$ is finite by \cref{subseteq,subseteq_finite_is_finite}.
    Therefore $1\in A$ by \cref{real_one_in_naturals}.
    Show that for all $k\in A$ we have $k + 1\in A$.
    \begin{subproof}
        Fix $k$ such that $k\in A$.
        We have $\initialSegment{k}\times\initialSegment{n}$ is finite.
        $k + 1\in\naturalsPlus$ by \cref{subseteq,real_naturals_closed_succ}.
        Let $C = \initialSegment{k + 1}\times\initialSegment{n}$.
        For all $u$ such that $u\in C$ we have $u\in (\initialSegment{k}\times\initialSegment{n})\union(\{k + 1\}\times\initialSegment{n})$
            by \cref{times_elem_is_tuple,initial_segment_naturalsplus,real_naturals_leq_succ_elim,union_iff,singleton_intro,times_tuple_intro,union_intro_left,union_intro_right}.
        Let $D = \{(k + 1,i) \mid i\in\initialSegment{n}\}$.
        For all $u$ such that $u\in\{k + 1\}\times\initialSegment{n}$ we have $u\in D$
            by \cref{times_elem_is_tuple,singleton_elim,pair_eq_iff}.
        Hence $\{k + 1\}\times\initialSegment{n}$ is finite
            by \cref{initial_segment_row_finite,subseteq,subseteq_finite_is_finite}.
        Therefore $C$ is finite by \cref{subseteq,union_of_finite_is_finite,subseteq_finite_is_finite}.
        Therefore $k + 1\in A$ by \cref{subseteq,real_naturals_closed_succ}.
    \end{subproof}
    Therefore $J$ is finite by \cref{real_naturals_induction,subseteq}.
\end{proof}

% from §36 helper lemma 14: functions vanish outside their support
\begin{lemma}\label{outside_support_zero}
    Assume $f$ is continuous from $X$ to $\reals$ with respect to $T$ and $\standardTopologyR$.
    Assume $x\in X\setminus \supp{f}{X}{T}$.
    Then $f(x) = 0$.
\end{lemma}
\begin{proof}
    Let $A = X\setminus \preimg{f}{\{0\}}$.
    We have $A\subseteq \supp{f}{X}{T}$
        by \cref{continuous,preimg_subseteq_dom,funs,subseteq_transitive,setminus_subseteq,support,subseteq_closure}.
    Hence $x\in\preimg{f}{\{0\}}$ by \cref{subseteq,setminus_elim_right,setminus,setminus_intro}.
    Take $r\in\{0\}$ such that $(x,r)\in f$ by \cref{preimg_iff}.
    Therefore $f(x) = 0$ by \cref{singleton_elim,continuous,funs_is_function,funs,subseteq,function_apply_intro}.
\end{proof}

% from §36 helper lemma 15: extending a weighted chart coordinate
\begin{lemma}\label{weighted_chart_coordinate_extension}
    Assume $m\in\naturalsPlus$.
    Assume $n\in\naturalsPlus$.
    Assume $U$ is a function on $\initialSegment{n}$.
    Assume $p$ is a partition of unity on $X$ dominated by $U$ with respect to $T$ and $n$.
    Assume $i\in\initialSegment{n}$.
    Assume $j\in\initialSegment{m}$.
    Let $R_0 = \{(k,\reals) \mid k\in\initialSegment{m}\}$.
    Let $S_0 = \{(k,\standardTopologyR) \mid k\in\initialSegment{m}\}$.
    Let $P_0 = \productTopologyIndexed{S_0}{R_0}$.
    Suppose $V$ is open in $\productFamily{R_0}$ with respect to $P_0$.
    Suppose $h$ is a homeomorphism between $\apply{U}{i}$ and $V$ with respect to
        $\subspaceTopology{T}{\apply{U}{i}}$ and $\subspaceTopology{P_0}{V}$.
    Let $A = \supp{\apply{p}{i}}{X}{T}$.
    Then there exists $u$ such that
        $u$ is continuous from $X$ to $\reals$ with respect to $T$ and $\standardTopologyR$ and
        for all $x\in A$ we have
            $u(x) = \apply{\apply{p}{i}}{x} \rmul \apply{\piIndex{R_0}{j}\circ h}{x}$ and
        for all $x\in X\setminus A$ we have $u(x) = 0$.
\end{lemma}
\begin{proof}
    Let $I_0 = \intervalclosed{0}{1}$.
    Let $T_0 = \subspaceTopology{\standardTopologyR}{I_0}$.
    Let $p_i = \apply{p}{i}$.
    Let $T_U = \subspaceTopology{T}{\apply{U}{i}}$.
    $p_i$ is continuous from $X$ to $I_0$ with respect to $T$ and $T_0$ and
        $A\subseteq \apply{U}{i}$ by \cref{finite_partition_of_unity,finite_partition_component_condition}.
    $T$ is a topology on $X$ by \cref{continuous}.
    $\apply{U}{i}$ is open in $X$ with respect to $T$ by \cref{finite_partition_of_unity}.
    We have $\apply{U}{i}\subseteq X$ by \cref{open_in,topology_on,subseteq,pow_iff}.
    $\apply{U}{i}$ is a subspace of $X$ with respect to $T$ by \cref{subspace_of}.
    $T_U$ is a topology on $\apply{U}{i}$ by \cref{subspace_topology_is_topology}.
    $R_0$ is a relation by \cref{relation}.
    For all $k,y_1,y_2$ such that $(k,y_1)\in R_0$ and $(k,y_2)\in R_0$ we have $y_1 = y_2$ by \cref{pair_eq_iff}.
    We have $R_0$ is a function on $\initialSegment{m}$ by \cref{rightunique,relation,dom_intro,dom_iff,pair_eq_iff,subseteq,subseteq_antisymmetric}.
    $S_0$ is a relation by \cref{relation}.
    For all $k,y_1,y_2$ such that $(k,y_1)\in S_0$ and $(k,y_2)\in S_0$ we have $y_1 = y_2$ by \cref{pair_eq_iff}.
    We have $S_0$ is a function on $\initialSegment{m}$ by \cref{rightunique,relation,dom_intro,dom_iff,pair_eq_iff,subseteq,subseteq_antisymmetric}.
    For all $k\in\initialSegment{m}$ we have $R_0(k) = \reals$ and $S_0(k) = \standardTopologyR$ by \cref{function_apply_intro}.
    $\standardTopologyR$ is a topology on $\reals$
        by \cref{standard_metric_is_metric,metric_space_hausdorff,standard_metric_topology,hausdorff}.
    $P_0$ is a topology on $\productFamily{R_0}$
        by \cref{product_subbasis_is_subbasis,product_subbasis_finite_intersections_basis,basis_topology_is_topology,basis_for_topology,hausdorff}.
    We have $V\subseteq \productFamily{R_0}$ by \cref{open_in,topology_on,subseteq,pow_iff}.
    $V$ is a subspace of $\productFamily{R_0}$ with respect to $P_0$ by \cref{subspace_of}.
    Let $h_0 = \identity{V}\circ h$.
    Therefore $h_0$ is continuous from $\apply{U}{i}$ to $\productFamily{R_0}$ with respect to $T_U$ and $P_0$
        by \cref{homeomorphism,continuous_expand_range}.
    Let $q_0 = \piIndex{R_0}{j}\circ h_0$.
    Hence $q_0$ is continuous from $\apply{U}{i}$ to $\reals$ with respect to $T_U$ and $\standardTopologyR$
        by \cref{projection_map_indexed_continuous,continuous_composite}.
    Let $r_0 = \identity{I_0}$.
    Let $r = r_0\circ p_i$.
    Therefore $r$ is continuous from $X$ to $\reals$ with respect to $T$ and $\standardTopologyR$
        by \cref{continuous_interval_expand_range}.
    $p_i\in\funs{X}{I_0}$ and $p_i$ is a function and $\dom{p_i} = X$
        by \cref{continuous,funs_is_function,funs}.
    We have $p_i$ is right-unique and $p_i$ is a relation by \cref{function_rightunique,funs_is_relation}.
    $\dom{r_0} = I_0$ and $r_0$ is right-unique and $r_0$ is a relation
        by \cref{id_dom,id_rightunique,id_is_relation}.
    We have $p_i\in\rels{X}{I_0}$ by \cref{funs_subseteq_rels,subseteq}.
    $\ran{p_i}\subseteq I_0$ by \cref{rels_ran_subseteq}.
    Therefore $\ran{p_i}\subseteq\dom{r_0}$ by \cref{id_dom}.
    For all $x\in X$ we have $r(x) = p_i(x)$ by \cref{subseteq,circ_apply,continuous,funs_type_apply,id_apply}.
    Let $r_U = \restrl{r}{\apply{U}{i}}$.
    Hence $r_U$ is continuous from $\apply{U}{i}$ to $\reals$ with respect to $T_U$ and $\standardTopologyR$
        by \cref{continuous_restrict_domain}.
    Let $w = \{(x, r_U(x) \rmul q_0(x)) \mid x\in \apply{U}{i}\}$.
    Therefore $w$ is continuous from $\apply{U}{i}$ to $\reals$ with respect to $T_U$ and $\standardTopologyR$
        by \cref{real_function_multiplication_continuous}.
    Let $z = \{(x,0) \mid x\in X\setminus A\}$.
    $0\in\reals$ by \cref{real_zero_lt_one,real_add_ident}.
    We have $z$ is a function from $X\setminus A$ to $\reals$ by \cref{constant_map_function}.
    Hence $\subspaceTopology{T}{X\setminus A}$ is a topology on $X\setminus A$
        by \cref{setminus_subseteq,subspace_of,subspace_topology_is_topology}.
    For all $x\in X\setminus A$ we have $z(x) = 0$ by
        \cref{constant_map_function,funs_elim,subseteq,funs_is_function,function_apply_intro}.
    Therefore $z$ is continuous from $X\setminus A$ to $\reals$ with respect to $\subspaceTopology{T}{X\setminus A}$ and $\standardTopologyR$
        by \cref{constant_continuous}.
    Show that for all $x\in \apply{U}{i}\inter (X\setminus A)$ we have $w(x) = z(x)$.
    \begin{subproof}
        Fix $x\in \apply{U}{i}\inter (X\setminus A)$.
        We have $x\in \apply{U}{i}$ and $x\in X\setminus A$ by \cref{inter}.
        Let $B = X\setminus \preimg{p_i}{\{0\}}$.
        $B\subseteq A$
            by \cref{continuous,preimg_subseteq_dom,funs,subseteq_transitive,setminus_subseteq,support,subseteq_closure}.
        Hence $x\notin B$ by \cref{subseteq,setminus_elim_right}.
        $x\in X$ and $x\notin X\setminus \preimg{p_i}{\{0\}}$ by \cref{setminus}.
        Therefore $x\in\preimg{p_i}{\{0\}}$ by \cref{setminus_intro}.
        Take $s\in\{0\}$ such that $(x,s)\in p_i$ by \cref{preimg_iff}.
        $p_i$ is a function and $x\in\dom{p_i}$ by \cref{funs_is_function,funs,subseteq}.
        Hence $p_i(x) = 0$ by \cref{singleton_elim,function_apply_intro}.
        $r_U$ is a function and $\dom{r_U} = \apply{U}{i}$ by \cref{continuous,funs_is_function,funs}.
        We have $w$ is a function and $\dom{w} = \apply{U}{i}$ by \cref{continuous,funs_is_function,funs}.
        $x\in\dom{r_U}$ by \cref{subseteq}.
        $r$ is a function and $\dom{r} = X$ by \cref{continuous,funs_is_function,funs}.
        Therefore $\apply{U}{i}\subseteq\dom{r}$ by \cref{subseteq}.
        Hence $r_U(x) = r(x)$ by \cref{restrl_of_function_apply}.
        Therefore $r(x) = 0$ by \cref{real_lt_subst_right}.
        We have $(x, r_U(x) \rmul q_0(x))\in w$ by \cref{pair_eq_iff}.
        $q_0(x)\in\reals$ by \cref{continuous,funs_type_apply}.
        Hence $w(x) = 0$ by \cref{function_apply_intro,real_lt_subst_left,real_mul_zero}.
        We have $z(x) = 0$
            by \cref{funs_elim,subseteq,constant_map_function,funs_is_function,function_apply_intro}.
        Therefore $w(x) = z(x)$.
    \end{subproof}
    Let $u = w\union z$.
    $w\in\funs{\apply{U}{i}}{\reals}$ and $z\in\funs{X\setminus A}{\reals}$ and
        $\dom{w} = \apply{U}{i}$ and $\dom{z} = X\setminus A$ by \cref{continuous,funs}.
    $u$ is a function and $\dom{u} = \dom{w}\union\dom{z}$ by \cref{inter,funs_is_function,subseteq,union_compatible_functions,dom_union}.
    $\apply{U}{i}\subseteq X$ by \cref{open_in,topology_on,subseteq,pow_iff,finite_partition_of_unity}.
    $X\setminus A\subseteq X$ by \cref{setminus_subseteq}.
    For all $x$ such that $x\in X$ we have $x\in \apply{U}{i}\union (X\setminus A)$.
    Therefore $X = \apply{U}{i}\union (X\setminus A)$ by \cref{subseteq,union_subsets_is_subset,subseteq_antisymmetric}.
    $A$ is closed in $X$ with respect to $T$ by \cref{support,closure_closed_in,setminus_subseteq,preimg_subseteq_dom,continuous,funs}.
    Hence $X\setminus A$ is open in $X$ with respect to $T$ by \cref{closed_in,topology_on,double_relative_complement,setminus_subseteq}.
    Let $M = \{\apply{U}{i}, X\setminus A\}$.
    For all $W$ such that $W\in M$ we have $W\in T$ by \cref{upair_iff,finite_partition_of_unity,open_in}.
    We have $M\subseteq T$ by \cref{subseteq}.
    For all $x$ such that $x\in\unions{M}$ we have $x\in \apply{U}{i}\union (X\setminus A)$.
    For all $x$ such that $x\in \apply{U}{i}\union (X\setminus A)$ we have $x\in\unions{M}$
        by \cref{union_iff,upair_intro_left,upair_intro_right,unions_iff}.
    We have $\unions{M} = \apply{U}{i}\union (X\setminus A)$ by \cref{subseteq,subseteq_antisymmetric}.
    Hence $\unions{M} = X$ by \cref{pair_eq_iff}.
    Therefore $\dom{u} = X$ by \cref{pair_eq_iff}.
    Show that for all $W$ such that $W\in M$ we have
        $\restrl{u}{W}$ is continuous from $W$ to $\reals$ with respect to
            $\subspaceTopology{T}{W}$ and $\standardTopologyR$.
    \begin{subproof}
        Fix $W$ such that $W\in M$.
        We have $W = \apply{U}{i}$ or $W = X\setminus A$ by \cref{upair_iff}.
        \begin{byCase}
        \caseOf{$W = \apply{U}{i}$.}
            Show that for all $v$ such that $v\in\restrl{u}{W}$ we have $v\in w$.
            \begin{subproof}
                Fix $v$ such that $v\in\restrl{u}{W}$.
                Hence $v\in u$ by \cref{restrl_subseteq,relation,subseteq}.
                Take $x,y$ such that $v = (x,y)$ by \cref{relation,pair_eq_iff}.
                We have $x\in W$ and $(x,y)\in u$ by \cref{restrl_iff,pair_eq_iff}.
                Therefore $x\in \apply{U}{i}$ and $(x,y)\in w\union z$ by \cref{pair_eq_iff,union_iff}.
                $(x,y)\in w$ or $(x,y)\in z$ by \cref{union_iff}.
                \begin{byCase}
                \caseOf{$(x,y)\in w$.}
                    Hence $v\in w$ by \cref{pair_eq_iff}.
                \caseOf{$(x,y)\in z$.}
                    $z$ is a function and $\dom{z} = X\setminus A$ by \cref{funs_elim}.
                    Take $t\in X\setminus A$ such that $(x,y) = (t,z(t))$ by \cref{funs_elim,function_member_elim,subseteq}.
                    Hence $x = t$ by \cref{pair_eq_iff}.
                    Therefore $x\in X\setminus A$ by \cref{subseteq}.
                    Thus $x\in \apply{U}{i}\inter (X\setminus A)$ by \cref{inter_intro}.
                    Hence $w(x) = z(x)$ by \cref{subseteq}.
                    $w$ is a function and $z$ is a function by \cref{continuous,funs_is_function,funs}.
                    We have $(x,w(x))\in w$ and $(x,z(x))\in z$ by \cref{function_apply_elim}.
                    Therefore $y = z(x)$ by \cref{function_rightunique}.
                    Hence $y = w(x)$ by \cref{real_lt_subst_right}.
                    Thus $v\in w$ by \cref{pair_eq_iff,function_apply_iff}.
                \end{byCase}
            \end{subproof}
            For all $v$ such that $v\in w$ we have $v\in \restrl{u}{W}$
                by \cref{relation,pair_eq_iff,continuous,funs_is_function,funs,function_apply_iff,union_iff,union_intro_left,restrl_iff}.
            Therefore $\restrl{u}{W} = w$ by \cref{subseteq,subseteq_antisymmetric}.
            Hence $\restrl{u}{W}$ is continuous from $W$ to $\reals$ with respect to
                $\subspaceTopology{T}{W}$ and $\standardTopologyR$ by \cref{continuous_eq_subst}.
        \caseOf{$W = X\setminus A$.}
            Show that for all $v$ such that $v\in\restrl{u}{W}$ we have $v\in z$.
            \begin{subproof}
                Fix $v$ such that $v\in\restrl{u}{W}$.
                Hence $v\in u$ by \cref{restrl_subseteq,relation,subseteq}.
                Take $x,y$ such that $v = (x,y)$ by \cref{relation,pair_eq_iff}.
                We have $x\in W$ and $(x,y)\in u$ by \cref{restrl_iff,pair_eq_iff}.
                Therefore $x\in X\setminus A$ and $(x,y)\in w\union z$ by \cref{pair_eq_iff,union_iff}.
                $(x,y)\in w$ or $(x,y)\in z$ by \cref{union_iff}.
                \begin{byCase}
                \caseOf{$(x,y)\in z$.}
                    Hence $v\in z$ by \cref{pair_eq_iff}.
                \caseOf{$(x,y)\in w$.}
                    $w$ is a function and $\dom{w} = \apply{U}{i}$ by \cref{continuous,funs_is_function,funs}.
                    Take $t\in \apply{U}{i}$ such that $(x,y) = (t,w(t))$ by \cref{funs_elim,function_member_elim,subseteq}.
                    Hence $x = t$ by \cref{pair_eq_iff}.
                    Therefore $x\in \apply{U}{i}$ by \cref{subseteq}.
                    Thus $x\in \apply{U}{i}\inter (X\setminus A)$ by \cref{inter_intro}.
                    Hence $w(x) = z(x)$ by \cref{subseteq}.
                    $w$ is a function and $z$ is a function by \cref{continuous,funs_is_function,funs}.
                    We have $(x,w(x))\in w$ and $(x,z(x))\in z$ by \cref{function_apply_elim}.
                    Therefore $y = w(x)$ by \cref{function_rightunique}.
                    Hence $y = z(x)$ by \cref{real_lt_subst_right}.
                    Thus $v\in z$ by \cref{pair_eq_iff,function_apply_iff}.
                \end{byCase}
            \end{subproof}
            For all $v$ such that $v\in z$ we have $v\in \restrl{u}{W}$
                by \cref{relation,pair_eq_iff,funs_elim,function_apply_iff,union_iff,union_intro_right,restrl_iff}.
            Therefore $\restrl{u}{W} = z$ by \cref{subseteq,subseteq_antisymmetric}.
            Hence $\restrl{u}{W}$ is continuous from $W$ to $\reals$ with respect to
                $\subspaceTopology{T}{W}$ and $\standardTopologyR$ by \cref{continuous_eq_subst}.
        \end{byCase}
    \end{subproof}
    $w\in\rels{\apply{U}{i}}{\reals}$ and $z\in\rels{X\setminus A}{\reals}$ and $\ran{u} = \ran{w}\union\ran{z}$
        by \cref{funs_subseteq_rels,subseteq,ran_union}.
    Hence $\ran{u}\subseteq\reals$ by \cref{rels_ran_subseteq,union_subsets_is_subset}.
    Therefore $u\in\funs{X}{\reals}$ by \cref{function_apply_elim,ran_iff,subseteq,funs_intro}.
    Hence $u$ is continuous from $X$ to $\reals$ with respect to $T$ and $\standardTopologyR$
        by \cref{continuous_local}.
    Show that for all $x\in A$ we have
        $u(x) = p_i(x) \rmul \apply{\piIndex{R_0}{j}\circ h}{x}$.
    \begin{subproof}
        Fix $x\in A$.
        We have $x\in \apply{U}{i}$ by \cref{subseteq}.
        $x\notin X\setminus A$ by \cref{setminus_elim_right}.
        $u$ is a function and $w$ is a function by \cref{continuous,funs_is_function,funs}.
        Hence $(x,w(x))\in w$ by \cref{function_apply_elim}.
        Therefore $(x,w(x))\in u$ by \cref{union_iff}.
        Hence $u(x) = w(x)$ by \cref{function_apply_intro}.
        $w$ is a function and $\dom{w} = \apply{U}{i}$ by \cref{continuous,funs_is_function,funs}.
        $x\in\dom{w}$ by \cref{subseteq}.
        Therefore $w(x) = r_U(x) \rmul q_0(x)$ by \cref{function_apply_intro}.
        $r_U$ is a function and $\dom{r_U} = \apply{U}{i}$ by \cref{continuous,funs_is_function,funs}.
        $x\in\dom{r_U}$ by \cref{subseteq}.
        $r$ is a function and $\dom{r} = X$ by \cref{continuous,funs_is_function,funs}.
        Therefore $\apply{U}{i}\subseteq\dom{r}$ by \cref{subseteq}.
        Hence $r_U(x) = r(x)$ by \cref{restrl_of_function_apply}.
        Therefore $r(x) = p_i(x)$ by \cref{subseteq}.
        $h\in\funs{\apply{U}{i}}{V}$ and $h$ is a function and $\dom{h} = \apply{U}{i}$
            by \cref{homeomorphism,bijections,bijections_to_funs,funs_is_function,funs}.
        $\identity{V}$ is right-unique and $\identity{V}$ is a relation and $\dom{\identity{V}} = V$
            by \cref{id_rightunique,id_is_relation,id_dom}.
        We have $h_0(x) = h(x)$ by \cref{circ_apply,funs_ran,subseteq,funs_type_apply,id_apply}.
        $q_0$ is a function and $\dom{q_0} = \apply{U}{i}$ by \cref{continuous,funs_is_function,funs}.
        $h_0$ is a function and $\dom{h_0} = \apply{U}{i}$ by \cref{continuous,funs_is_function,funs}.
        Hence $(x,h_0(x))\in h_0$ by \cref{function_apply_elim}.
        We have $h_0(x)\in\productFamily{R_0}$ by \cref{continuous,funs_type_apply}.
        Therefore $(h_0(x),\apply{h_0(x)}{j})\in \piIndex{R_0}{j}$ by \cref{projection_map_indexed}.
        Hence $(x,\apply{h_0(x)}{j})\in q_0$ by \cref{circ_iff}.
        $q_0(x) = \apply{h_0(x)}{j}$ by \cref{function_apply_intro}.
        Therefore $q_0(x) = \apply{h(x)}{j}$ by \cref{real_lt_subst_right}.
        $h\in\funs{\apply{U}{i}}{\productFamily{R_0}}$ by \cref{funs_weaken_codom,subseteq}.
        $\piIndex{R_0}{j}\circ h\in\funs{\apply{U}{i}}{\reals}$ by \cref{projection_map_indexed_function,funs_circ,funs_elim}.
        We have $\piIndex{R_0}{j}\circ h$ is a function by \cref{funs_is_function}.
        Hence $(x,h(x))\in h$ by \cref{function_apply_elim}.
        Therefore $(h(x),\apply{h(x)}{j})\in \piIndex{R_0}{j}$ by \cref{projection_map_indexed,subseteq}.
        Hence $(x,\apply{h(x)}{j})\in \piIndex{R_0}{j}\circ h$ by \cref{circ_iff}.
        $\apply{\piIndex{R_0}{j}\circ h}{x} = \apply{h(x)}{j}$ by \cref{function_apply_intro}.
        Therefore $u(x) = p_i(x) \rmul \apply{\piIndex{R_0}{j}\circ h}{x}$ by \cref{real_lt_subst_left,real_lt_subst_right}.
    \end{subproof}
    Show that for all $x\in X\setminus A$ we have $u(x) = 0$.
    \begin{subproof}
        Follows by \cref{continuous,funs_is_function,funs,function_apply_elim,union_iff,function_apply_intro,funs_elim,subseteq,constant_map_function}.
    \end{subproof}
\end{proof}

% from §36 helper lemma 16: canceling a positive real factor
\begin{lemma}\label{positive_real_factor_cancel}
    Assume $r\in\reals$.
    Assume $a\in\reals$.
    Assume $b\in\reals$.
    Assume $0 < r$.
    Suppose $r \rmul a = r \rmul b$.
    Then $a = b$.
\end{lemma}
\begin{proof}
    Follows by \cref{real_pos_nonzero,real_mul_cancel}.
\end{proof}

% from §36 helper lemma 17: coordinatewise equality in a finite real product
\begin{lemma}\label{finite_real_product_coordinate_extensionality}
    Assume $m\in\naturalsPlus$.
    Let $R_0 = \{(k,\reals) \mid k\in\initialSegment{m}\}$.
    Assume $x\in\productFamily{R_0}$.
    Assume $y\in\productFamily{R_0}$.
    Suppose for all $j\in\initialSegment{m}$ we have $x(j) = y(j)$.
    Then $x = y$.
\end{lemma}
\begin{proof}
    $R_0$ is a function on $\initialSegment{m}$ by \cref{relation,rightunique,dom_intro,dom_iff,pair_eq_iff,subseteq,subseteq_antisymmetric}.
    Hence $x\in\funs{\initialSegment{m}}{\reals}$ and $y\in\funs{\initialSegment{m}}{\reals}$
        by \cref{real_one_in_naturals,real_one_leq_naturalsplus,initial_segment_naturalsplus,product_family_reals_power,reals_power,tuple_power,subseteq}.
    Therefore $x = y$ by \cref{reals_power,tuple_power,subseteq,funs_is_function,funs,functions_eq_iff}.
\end{proof}

% from §36 helper lemma 18: a nonnegative finite real sum equal to one has a positive term
\begin{lemma}\label{finite_real_sum_one_positive_term}
    Assume $n\in\naturalsPlus$.
    Assume $a\in\funs{\initialSegment{n}}{\reals}$.
    Assume $\sumFiniteSequence{a}{1}$.
    Suppose for all $i\in\initialSegment{n}$ we have $0 \leq a(i)$.
    Then there exists $i\in\initialSegment{n}$ such that $0 < a(i)$.
\end{lemma}
\begin{proof}
    Suppose not.
    For all $i\in\initialSegment{n}$ we have $a(i) = 0$ by \cref{real_leq_split}.
    $0\in\reals$ by \cref{real_add_ident}.
    Therefore $1 \leq n \rmul 0$ by \cref{sum_finite_sequence_leq_n_times}.
    $n\in\reals$ and $1\in\reals$ by \cref{real_naturals_subset_reals,subseteq,real_one_in_naturals}.
    Thus $1 \leq 0$ by \cref{real_mul_zero,real_mul_comm,real_lt_subst_right}.
    Hence $1 = 0$ by \cref{real_zero_lt_one,real_lt_imp_leq,real_leq_antisym}.
    Contradiction by \cref{real_zero_lt_one,real_lt_irreflexive,real_lt_subst_right}.
\end{proof}

% from §36 helper definition 9: agreement of a reindexed finite family
\begin{definition}\label{finite_reindex_agreement}
    $\finiteReindexAgreement{f}{h}{c}{J}$ iff
        for all $p\in J$ we have $\apply{f}{\apply{h}{p}} = \apply{c}{p}$.
\end{definition}

% from §36 helper lemma 19: reindexing a finite grid of continuous real-valued maps
\begin{lemma}\label{finite_grid_continuous_reindex}
    Assume $m\in\naturalsPlus$.
    Assume $n\in\naturalsPlus$.
    Assume $T$ is a topology on $X$.
    Let $J = \initialSegment{n}\times\initialSegment{m}$.
    Assume $c$ is a function on $J$.
    Suppose for all $p\in J$ we have $\apply{c}{p}$ is continuous from $X$ to $\reals$ with respect to $T$ and $\standardTopologyR$.
    Then there exist $N,h,f$ such that
        $N\in\naturalsPlus$ and
        $h\in\inj{J}{\initialSegment{N}}$ and
        $f$ is a function on $\initialSegment{N}$ and
        for all $k\in\initialSegment{N}$ we have $\apply{f}{k}$ is continuous from $X$ to $\reals$ with respect to $T$ and $\standardTopologyR$ and
        $\finiteReindexAgreement{f}{h}{c}{J}$.
\end{lemma}
\begin{proof}
    $J$ is finite by \cref{initial_segment_grid_finite}.
    Take $g$ such that $g\in\inj{\naturalsPlus\times\naturalsPlus}{\naturalsPlus}$ by \cref{naturalsplus_pairing_injection}.
    Let $h = \restrl{g}{J}$.
    $g\in\funs{\naturalsPlus\times\naturalsPlus}{\naturalsPlus}$ by \cref{inj}.
    We have $g$ is a function and $g$ is a relation and $\dom{g} = \naturalsPlus\times\naturalsPlus$
        by \cref{funs_is_function,funs_is_relation,funs}.
    For all $p_1,p_2$ such that $p_1\in\dom{g}$ and $p_2\in\dom{g}$ and $g(p_1) = g(p_2)$ we have $p_1 = p_2$
        by \cref{subseteq,inj}.
    Hence $g$ is injective by \cref{injective_function}.
    $h$ is a function by \cref{restrl_of_function_is_function}.
    For all $p$ such that $p\in J$ we have $p\in\naturalsPlus\times\naturalsPlus$ by
        \cref{times_elem_is_tuple,initial_segment_naturalsplus,times_tuple_intro,pair_eq_iff}.
    Therefore $J\subseteq\naturalsPlus\times\naturalsPlus$ by \cref{subseteq}.
    Hence $\dom{h} = J$ by \cref{funs,dom_of_restrl_eq_inter,inter_absorb_supseteq_left}.
    For all $p\in\dom{h}$ we have $h(p)\in\naturalsPlus$ by
        \cref{subseteq,initial_segment_naturalsplus,funs,restrl_of_function_apply,inj,funs_type_apply}.
    We have $h\in\funs{J}{\naturalsPlus}$ by \cref{funs_intro}.
    For all $p_1,p_2$ such that $p_1\in J$ and $p_2\in J$ and $h(p_1) = h(p_2)$ we have $p_1 = p_2$ by
        \cref{subseteq,initial_segment_naturalsplus,funs,restrl_of_function_apply,pair_eq_iff,injective_function}.
    Hence $h\in\inj{J}{\naturalsPlus}$ by \cref{inj}.
    Let $A = \img{h}{J}$.
    Therefore $A$ is finite by \cref{inj,funs_is_function,funs_is_relation,funs,finite_image_function}.
    We have $(1,1)\in J$ by \cref{real_one_in_naturals,real_one_leq_naturalsplus,initial_segment_naturalsplus,times_tuple_intro}.
    Hence $h((1,1))\in A$ by \cref{inj,img_of_function_intro,inter_intro}.
    We have $A\neq\emptyset$ by \cref{inhabited_nonempty,nonempty_inhabited}.
    We have $A\subseteq\naturalsPlus$ by \cref{inj,funs_subseteq_rels,subseteq,rels_ran_subseteq,img_subseteq_ran,subseteq_transitive}.
    Hence $A\subseteq\reals$ by \cref{subseteq,real_naturals_subset_reals,subseteq_transitive}.
    Take $N\in A$ such that $N$ is a maximum of $A$ by \cref{finite_real_maximum}.
    For all $k$ such that $k\in A$ we have $k\in\initialSegment{N}$ by
        \cref{subseteq,real_maximum,initial_segment_naturalsplus}.
    Therefore $A\subseteq\initialSegment{N}$ by \cref{subseteq}.
    For all $p\in\dom{h}$ we have $h(p)\in\initialSegment{N}$ by \cref{subseteq,inj,img_of_function_intro,inter_intro}.
    We have $h\in\funs{J}{\initialSegment{N}}$ by \cref{funs_intro}.
    For all $p_1,p_2$ such that $p_1\in J$ and $p_2\in J$ and $h(p_1) = h(p_2)$ we have $p_1 = p_2$
        by \cref{inj}.
    Hence $h\in\inj{J}{\initialSegment{N}}$ by \cref{inj}.
    Let $z = \{(x,0) \mid x\in X\}$.
    $0\in\reals$ by \cref{real_add_ident}.
    $\standardTopologyR$ is a topology on $\reals$
        by \cref{standard_basis_is_basis,basis_topology_is_topology,standard_topology_reals}.
    $z$ is a function from $X$ to $\reals$ by \cref{constant_map_function}.
    Therefore for all $x\in X$ we have $z(x) = 0$ by \cref{constant_map_function,funs_is_function,function_apply_intro}.
    $z$ is continuous from $X$ to $\reals$ with respect to $T$ and $\standardTopologyR$ by \cref{constant_continuous,standard_topology_reals}.
    Let $f_1 = \{(h(p),c(p)) \mid p\in J\}$.
    For all $k,u_1,u_2$ such that $(k,u_1)\in f_1$ and $(k,u_2)\in f_1$ we have $u_1 = u_2$ by
        \cref{pair_eq_iff,inj}.
    Therefore $f_1$ is a relation and $f_1$ is right-unique by \cref{relation,pair_eq_iff,inj,rightunique}.
    For all $k$ such that $k\in A$ we have $k\in\dom{f_1}$
        by \cref{inj,funs,funs_is_function,function_rightunique,funs_is_relation,img_iff,function_apply_intro,pair_eq_iff,dom_intro}.
    For all $k$ such that $k\in\dom{f_1}$ we have $k\in A$
        by \cref{dom_iff,pair_eq_iff,inj,funs_elim,subseteq,function_apply_elim,img_iff}.
    Hence $\dom{f_1} = A$ by \cref{subseteq,subseteq_antisymmetric}.
    Let $f_2 = \{(k,z) \mid k\in\initialSegment{N}\setminus A\}$.
    For all $k,u_1,u_2$ such that $(k,u_1)\in f_2$ and $(k,u_2)\in f_2$ we have $u_1 = u_2$ by \cref{pair_eq_iff}.
    We have $f_2$ is a relation and $f_2$ is right-unique by \cref{relation,pair_eq_iff,rightunique}.
    Hence $\dom{f_2} = \initialSegment{N}\setminus A$ by \cref{dom_iff,subseteq,subseteq_antisymmetric,pair_eq_iff}.
    Let $f = f_1\union f_2$.
    For all $k$ such that $k\in\dom{f_1}\inter\dom{f_2}$ we have $f_1(k) = f_2(k)$.
    Therefore $f$ is a function on $\initialSegment{N}$
        by \cref{union_compatible_functions,dom_union,pair_eq_iff,setminus_partition,subseteq,funs}.
    Show that for all $k\in\initialSegment{N}$ we have $f(k)$ is continuous from $X$ to $\reals$ with respect to $T$ and $\standardTopologyR$.
    \begin{subproof}
        Fix $k\in\initialSegment{N}$.
        \begin{byCase}
        \caseOf{$k\in A$.}
            Take $p\in J$ such that $(p,k)\in h$ by \cref{inj,img_iff}.
            We have $k = h(p)$ by \cref{function_apply_intro}.
            Hence $(k,c(p))\in f$ by \cref{pair_eq_iff,union_iff,union_intro_left}.
            $f(k) = c(p)$ by \cref{function_apply_intro}.
            Therefore $f(k)$ is continuous from $X$ to $\reals$ with respect to $T$ and $\standardTopologyR$ by \cref{continuous_eq_subst}.
        \caseOf{$k\notin A$.}
            We have $k\in\initialSegment{N}\setminus A$ by \cref{setminus_intro}.
            Hence $(k,z)\in f$ by \cref{pair_eq_iff,union_iff,union_intro_right}.
            $f(k) = z$ by \cref{function_apply_intro}.
            Hence $f(k)$ is continuous from $X$ to $\reals$ with respect to $T$ and $\standardTopologyR$ by \cref{continuous_eq_subst}.
        \end{byCase}
    \end{subproof}
    Show that for all $p\in J$ we have $\apply{f}{\apply{h}{p}} = \apply{c}{p}$.
    \begin{subproof}
        Follows by \cref{inj,funs_elim,function_apply_elim,img_iff,pair_eq_iff,union_iff,union_intro_left,function_apply_intro,funs_type_apply}.
    \end{subproof}
    Therefore $\finiteReindexAgreement{f}{h}{c}{J}$ by \cref{finite_reindex_agreement}.
\end{proof}

% from §36 helper lemma 20: evaluating a finite family of continuous real-valued maps
\begin{lemma}\label{finite_real_family_evaluation_map}
    Assume $n\in\naturalsPlus$.
    Assume $T$ is a topology on $X$.
    Assume $f$ is a function on $\initialSegment{n}$.
    Suppose for all $i\in\initialSegment{n}$ we have $\apply{f}{i}$ is continuous from $X$ to $\reals$ with respect to $T$ and $\standardTopologyR$.
    Let $R = \{(i,\reals) \mid i\in\initialSegment{n}\}$.
    Let $S = \{(i,\standardTopologyR) \mid i\in\initialSegment{n}\}$.
    Let $P = \productTopologyIndexed{S}{R}$.
    Let $F = \{(x,z) \mid x\in X, z\in\productFamily{R} \mid
        \forall i\in\initialSegment{n}. z(i) = \apply{\apply{f}{i}}{x}\}$.
    Then $F$ is a function from $X$ to $\productFamily{R}$ and
        $F$ is continuous from $X$ to $\productFamily{R}$ with respect to $T$ and $P$ and
        for all $x\in X$ for all $i\in\initialSegment{n}$ we have $\apply{\apply{F}{x}}{i} = \apply{\apply{f}{i}}{x}$.
\end{lemma}
\begin{proof}
    $\standardTopologyR$ is a topology on $\reals$
        by \cref{standard_basis_is_basis,basis_topology_is_topology,standard_topology_reals}.
    For all $i,y_1,y_2$ such that $(i,y_1)\in R$ and $(i,y_2)\in R$ we have $y_1 = y_2$ by \cref{pair_eq_iff}.
    Hence $R$ is a function on $\initialSegment{n}$ by \cref{relation,rightunique,dom_intro,dom_iff,pair_eq_iff,subseteq,subseteq_antisymmetric}.
    For all $i,y_1,y_2$ such that $(i,y_1)\in S$ and $(i,y_2)\in S$ we have $y_1 = y_2$ by \cref{pair_eq_iff}.
    Therefore $S$ is a function on $\initialSegment{n}$ by \cref{relation,rightunique,dom_intro,dom_iff,pair_eq_iff,subseteq,subseteq_antisymmetric}.
    For all $i\in\initialSegment{n}$ we have $S(i) = \standardTopologyR$ and $R(i) = \reals$
        by \cref{funs,function_apply_intro}.
    We have $1\in\initialSegment{n}$ by \cref{real_one_in_naturals,real_one_leq_naturalsplus,initial_segment_naturalsplus}.
    $\productFamily{R} = \realsPower{\initialSegment{n}}$ by \cref{product_family_reals_power}.
    $P$ is a topology on $\productFamily{R}$
        by \cref{product_subbasis_is_subbasis,product_subbasis_finite_intersections_basis,basis_topology_is_topology,basis_for_topology,standard_topology_reals}.

    For all $x,z_1,z_2$ such that $(x,z_1)\in F$ and $(x,z_2)\in F$ we have $z_1 = z_2$
        by \cref{pair_eq_iff,finite_real_product_coordinate_extensionality}.
    Hence $F$ is a function by \cref{relation,rightunique,pair_eq_iff,finite_real_product_coordinate_extensionality}.
    Show that for all $x\in X$ there exists $z$ such that $(x,z)\in F$.
    \begin{subproof}
        Fix $x\in X$.
        Let $z = \{(i,\apply{\apply{f}{i}}{x}) \mid i\in\initialSegment{n}\}$.
        For all $i,y_1,y_2$ such that $(i,y_1)\in z$ and $(i,y_2)\in z$ we have $y_1 = y_2$ by \cref{pair_eq_iff}.
        Therefore $z$ is a function by \cref{relation,rightunique,pair_eq_iff}.
        We have $\dom{z} = \initialSegment{n}$ by \cref{dom_intro,dom_iff,pair_eq_iff,subseteq,subseteq_antisymmetric}.
        For all $i\in\dom{z}$ we have $z(i)\in\reals$
            by \cref{dom_iff,pair_eq_iff,continuous,funs_type_apply,function_apply_intro}.
        Hence $z\in\productFamily{R}$ by \cref{funs_intro,reals_power,tuple_power,subseteq}.
        For all $i\in\initialSegment{n}$ we have $z(i) = \apply{\apply{f}{i}}{x}$ by \cref{function_apply_intro}.
        Therefore $(x,z)\in F$ by \cref{pair_eq_iff}.
    \end{subproof}
    Hence $\dom{F} = X$ by \cref{subseteq,dom_iff,pair_eq_iff,subseteq_antisymmetric}.
    For all $x\in\dom{F}$ we have $F(x)\in\productFamily{R}$ by \cref{function_apply_iff,pair_eq_iff}.
    Therefore $F\in\funs{X}{\productFamily{R}}$ by \cref{funs_intro}.

    Show that for all $x\in X$ for all $i\in\initialSegment{n}$ we have $\apply{\apply{F}{x}}{i} = \apply{\apply{f}{i}}{x}$.
    \begin{subproof}
        Follows by \cref{funs_elim,function_apply_elim,pair_eq_iff}.
    \end{subproof}

    Show that for all $i\in\initialSegment{n}$ we have $\piIndex{R}{i}\circ F$ is continuous from $X$ to $\reals$ with respect to $T$ and $\standardTopologyR$.
    \begin{subproof}
        Follows by \cref{projection_map_indexed_function,funs_circ,funs_elim,continuous,funs_is_function,funs_type_apply,projection_map_indexed,function_apply_elim,circ_elem_intro,function_apply_intro,pair_eq_iff,functions_eq_iff,continuous_eq_subst}.
    \end{subproof}
    Hence $F$ is continuous from $X$ to $\productFamily{R}$ with respect to $T$ and $P$
        by \cref{continuous_map_into_indexed_product,standard_basis_is_basis,basis_topology_is_topology,standard_topology_reals,product_family_reals_power,funs,subseteq}.
\end{proof}

% from §36 helper lemma 21: empty compact manifolds embed in euclidean spaces
\begin{lemma}\label{compact_manifold_empty_embeds_in_euclidean_space}
    Assume $m\in\naturalsPlus$.
    Assume $X$ is a manifold of dimension $m$ with respect to $T$.
    Assume $X$ is compact with respect to $T$.
    Assume $X = \emptyset$.
    Let $n = 1$.
    Let $R = \{(i,\reals) \mid i\in\initialSegment{n}\}$.
    Let $S = \{(i,\standardTopologyR) \mid i\in\initialSegment{n}\}$.
    Let $f = \emptyset$.
    Then $n\in\naturalsPlus$ and
        $f$ is a topological embedding from $X$ to $\productFamily{R}$ with respect to $T$ and $\productTopologyIndexed{S}{R}$.
\end{lemma}
\begin{proof}
    $X$ is Hausdorff with respect to $T$ by \cref{m_manifold}.
    Hence $T$ is a topology on $X$ by \cref{hausdorff}.
    Let $P = \productTopologyIndexed{S}{R}$.
    $\standardTopologyR$ is a topology on $\reals$
        by \cref{standard_basis_is_basis,basis_topology_is_topology,standard_topology_reals}.
    For all $i,y_1,y_2$ such that $(i,y_1)\in R$ and $(i,y_2)\in R$ we have $y_1 = y_2$ by \cref{pair_eq_iff}.
    Therefore $R$ is a function on $\initialSegment{n}$ by \cref{relation,rightunique,dom_intro,dom_iff,pair_eq_iff,subseteq,subseteq_antisymmetric}.
    For all $i,y_1,y_2$ such that $(i,y_1)\in S$ and $(i,y_2)\in S$ we have $y_1 = y_2$ by \cref{pair_eq_iff}.
    $S$ is a function on $\initialSegment{n}$ by \cref{relation,rightunique,dom_intro,dom_iff,pair_eq_iff,subseteq,subseteq_antisymmetric}.
    For all $i\in\initialSegment{n}$ we have $\apply{R}{i} = \reals$ and $\apply{S}{i} = \standardTopologyR$
        by \cref{funs,function_apply_intro}.
    We have $\initialSegment{n}$ is inhabited by \cref{real_one_in_naturals,real_one_leq_naturalsplus,initial_segment_naturalsplus}.
    $P$ is a topology on $\productFamily{R}$
        by \cref{product_subbasis_is_subbasis,product_subbasis_finite_intersections_basis,basis_topology_is_topology,basis_for_topology,standard_topology_reals}.
    Show that $f$ is a topological embedding from $X$ to $\productFamily{R}$ with respect to $T$ and $P$.
    \begin{subproof}
        Follows by \cref{empty_domain_topological_embedding}.
    \end{subproof}
    Hence $n\in\naturalsPlus$ and
        $f$ is a topological embedding from $X$ to $\productFamily{R}$ with respect to $T$ and $\productTopologyIndexed{S}{R}$.
\end{proof}

% from §36 helper lemma 22: extracting the covering from a chart-cover witness
\begin{lemma}\label{chart_cover_witness_is_covering}
    Assume $n\in\naturalsPlus$.
    Assume $U$ is a function on $\initialSegment{n}$.
    Assume for all $i\in\initialSegment{n}$ there exist $V,h$ such that
        $\apply{U}{i}$ is open in $X$ with respect to $T$ and
        $V$ is open in $\productFamily{R_0}$ with respect to $P_0$ and
        $h$ is a homeomorphism between $\apply{U}{i}$ and $V$ with respect to
            $\subspaceTopology{T}{\apply{U}{i}}$ and $\subspaceTopology{P_0}{V}$ and
        $\img{U}{\initialSegment{n}}$ is a covering of $X$.
    Then $\img{U}{\initialSegment{n}}$ is a covering of $X$.
\end{lemma}
\begin{proof}
    Follows by \cref{real_one_in_naturals,real_one_leq_naturalsplus,initial_segment_naturalsplus}.
\end{proof}

% from §36 helper lemma 25: a partition of unity yields a sum witness at a point
\begin{lemma}\label{finite_partition_of_unity_sum_witness}
    Assume $p$ is a partition of unity on $X$ dominated by $U$ with respect to $T$ and $n$.
    Assume $x\in X$.
    Then there exists $a$ such that
        $a\in\funs{\initialSegment{n}}{\reals}$ and
        for all $i\in\initialSegment{n}$ we have
            $a(i) = \apply{\apply{p}{i}}{x}$ and
            $\sumFiniteSequence{a}{1}$.
\end{lemma}
\begin{proof}
    Follows by \cref{finite_partition_of_unity,real_one_in_naturals,real_one_leq_naturalsplus,initial_segment_naturalsplus,finite_partition_sum_condition}.
\end{proof}

% from §36 helper lemma 26a: positivity forces support membership for a partition component
\begin{lemma}\label{positive_partition_component_in_support}
    Assume $p$ is a partition of unity on $X$ dominated by $U$ with respect to $T$ and $n$.
    Assume $i\in\initialSegment{n}$.
    Assume $x\in X$.
    Suppose $0 < \apply{\apply{p}{i}}{x}$.
    Then $x\in \supp{\apply{p}{i}}{X}{T}$.
\end{lemma}
\begin{proof}
    Let $A = \supp{\apply{p}{i}}{X}{T}$.
    Let $B = X\setminus \preimg{\apply{p}{i}}{\{0\}}$.
    $\apply{p}{i}$ is continuous from $X$ to $\intervalclosed{0}{1}$ with respect to $T$ and
        $\subspaceTopology{\standardTopologyR}{\intervalclosed{0}{1}}$ by \cref{finite_partition_of_unity,finite_partition_component_condition}.
    Therefore $B\subseteq A$ by \cref{continuous,setminus_subseteq,support,subseteq_closure}.
    Hence $x\in A$
        by \cref{continuous,funs_is_function,funs,preimg_iff,singleton_elim,function_apply_intro,real_lt_subst_right,real_add_ident,real_lt_irreflexive,setminus_intro,subseteq}.
\end{proof}

% from §36 helper lemma 26: nonempty compact manifolds embed in euclidean spaces
\begin{lemma}\label{compact_manifold_nonempty_embeds_in_euclidean_space}
    Assume $m\in\naturalsPlus$.
    Assume $X$ is a manifold of dimension $m$ with respect to $T$.
    Assume $X$ is compact with respect to $T$.
    Assume $X\neq\emptyset$.
    Then there exist $n, R, S, f$ such that
        $n\in\naturalsPlus$ and
        $R = \{(i,\reals) \mid i\in\initialSegment{n}\}$ and
        $S = \{(i,\standardTopologyR) \mid i\in\initialSegment{n}\}$ and
        $f$ is a topological embedding from $X$ to $\productFamily{R}$ with respect to $T$ and $\productTopologyIndexed{S}{R}$.
\end{lemma}
\begin{proof}
    We have $X$ is inhabited by \cref{nonempty_inhabited}.
    Let $I_0 = \intervalclosed{0}{1}$.
    Let $T_0 = \subspaceTopology{\standardTopologyR}{I_0}$.
    Let $j_0 = \identity{I_0}$.
    Let $R_0 = \{(j,\reals) \mid j\in\initialSegment{m}\}$.
    Let $S_0 = \{(j,\standardTopologyR) \mid j\in\initialSegment{m}\}$.
    Let $P_0 = \productTopologyIndexed{S_0}{R_0}$.
    Take $n,U$ such that
        $n\in\naturalsPlus$ and
        $U$ is a function on $\initialSegment{n}$ and
        for all $i\in\initialSegment{n}$ there exist $V,h$ such that
            $\apply{U}{i}$ is open in $X$ with respect to $T$ and
            $V$ is open in $\productFamily{R_0}$ with respect to $P_0$ and
            $h$ is a homeomorphism between $\apply{U}{i}$ and $V$ with respect to
                $\subspaceTopology{T}{\apply{U}{i}}$ and $\subspaceTopology{P_0}{V}$ and
        $\img{U}{\initialSegment{n}}$ is a covering of $X$
        by \cref{compact_manifold_finite_chart_cover}.
    $X$ is Hausdorff with respect to $T$ by \cref{m_manifold}.
    Hence $T$ is a topology on $X$ by \cref{hausdorff}.
    Therefore $X$ is normal with respect to $T$ by \cref{compact_hausdorff_normal}.
    Hence $\img{U}{\initialSegment{n}}$ is a covering of $X$ by \cref{chart_cover_witness_is_covering}.
    Take $p$ such that
        $p$ is a partition of unity on $X$ dominated by $U$ with respect to $T$ and $n$
        by \cref{finite_partition_of_unity_existence}.

    Let $B_0 = \pow{\productFamily{R_0}}\times\pow{X\times\productFamily{R_0}}$.
    Let $A = \{w\in\initialSegment{n}\times B_0 \mid
        \text{there exist $i,V,h$ such that
            $w = (i,(V,h))$ and
            $\apply{U}{i}$ is open in $X$ with respect to $T$ and
            $V$ is open in $\productFamily{R_0}$ with respect to $P_0$ and
            $h$ is a homeomorphism between $\apply{U}{i}$ and $V$ with respect to
                $\subspaceTopology{T}{\apply{U}{i}}$ and $\subspaceTopology{P_0}{V}$}\}$.
    Show that for all $i\in\initialSegment{n}$ there exists $q$ such that $i\mathrel{A} q$.
    \begin{subproof}
        Fix $i\in\initialSegment{n}$.
        Take $V,h$ such that
            $\apply{U}{i}$ is open in $X$ with respect to $T$ and
            $V$ is open in $\productFamily{R_0}$ with respect to $P_0$ and
            $h$ is a homeomorphism between $\apply{U}{i}$ and $V$ with respect to
                $\subspaceTopology{T}{\apply{U}{i}}$ and $\subspaceTopology{P_0}{V}$.
        Hence $V\subseteq \productFamily{R_0}$ by \cref{open_in,topology_on,subseteq,pow_iff}.
        $V\in\pow{\productFamily{R_0}}$ by \cref{powerset_intro}.
        Therefore $h\subseteq \apply{U}{i}\times V$
            by \cref{homeomorphism,bijections,bijections_to_funs,funs_subseteq_rels,subseteq,rels_elim}.
        $\apply{U}{i}\subseteq X$ by \cref{open_in,topology_on,subseteq,pow_iff}.
        Hence $\apply{U}{i}\times V\subseteq X\times V$ by \cref{times_subseteq_left}.
        $X\times V\subseteq X\times \productFamily{R_0}$ by \cref{times_subseteq_right,subseteq}.
        Hence $h\subseteq X\times \productFamily{R_0}$ by \cref{subseteq_transitive}.
        Therefore $h\in\pow{X\times\productFamily{R_0}}$ by \cref{powerset_intro}.
        Take $q$ such that $q = (V,h)$.
        Hence $q\in B_0$ by \cref{times_tuple_intro,pair_eq_iff}.
        $(i,q)\in \initialSegment{n}\times B_0$ by \cref{times_tuple_intro}.
        Therefore $i\mathrel{A} q$ by \cref{pair_eq_iff}.
    \end{subproof}
    Take $a$ such that $a$ is a function on $\initialSegment{n}$ and
        for all $i\in\initialSegment{n}$ we have $i\mathrel{A}\apply{a}{i}$
        by \cref{relation,choice_function}.
    Show that for all $i\in\initialSegment{n}$ we have
        $\apply{U}{i}$ is open in $X$ with respect to $T$ and
        $\fst{\apply{a}{i}}$ is open in $\productFamily{R_0}$ with respect to $P_0$ and
        $\snd{\apply{a}{i}}$ is a homeomorphism between $\apply{U}{i}$ and $\fst{\apply{a}{i}}$ with respect to
            $\subspaceTopology{T}{\apply{U}{i}}$ and $\subspaceTopology{P_0}{\fst{\apply{a}{i}}}$.
    \begin{subproof}
        Follows by \cref{choice_function,pair_eq_iff,fst_eq,snd_eq}.
    \end{subproof}

    Show that for all $i\in\initialSegment{n}$ we have
        $j_0\circ \apply{p}{i}$ is continuous from $X$ to $\reals$ with respect to $T$ and $\standardTopologyR$ and
        for all $x\in X$ we have $\apply{j_0\circ \apply{p}{i}}{x} = \apply{\apply{p}{i}}{x}$.
    \begin{subproof}
        Follows by \cref{finite_partition_of_unity,finite_partition_component_condition,continuous_interval_expand_range,continuous,funs_is_function,funs,function_rightunique,funs_is_relation,funs_subseteq_rels,subseteq,rels_ran_subseteq,id_dom,id_rightunique,id_is_relation,circ_apply,funs_type_apply,id_apply}.
    \end{subproof}

    $m + 1\in\naturalsPlus$ by \cref{real_naturals_closed_succ}.
    Let $J = \initialSegment{n}\times\initialSegment{m + 1}$.
    Let $C_0 = \pow{X\times\reals}$.
    Let $Q_1 = \{w\in J\times C_0 \mid
        \text{there exist $s,u,i,j$ such that
            $w = (s,u)$ and
            $s = (i,j)$ and
            $j\in\initialSegment{m}$ and
            $u$ is continuous from $X$ to $\reals$ with respect to $T$ and $\standardTopologyR$ and
            for all $x\in \supp{\apply{p}{i}}{X}{T}$ we have
                $u(x) = \apply{\apply{p}{i}}{x} \rmul \apply{\piIndex{R_0}{j}\circ \snd{\apply{a}{i}}}{x}$ and
            for all $x\in X\setminus \supp{\apply{p}{i}}{X}{T}$ we have $u(x) = 0$}\}$.
    Let $Q_2 = \{w\in J\times C_0 \mid
        \text{there exist $s,u,i,j$ such that
            $w = (s,u)$ and
            $s = (i,j)$ and
            $j = m + 1$ and
            $u = j_0\circ \apply{p}{i}$}\}$.
    Let $Q = Q_1\union Q_2$.
    Show that for all $s\in J$ there exists $u$ such that $s\mathrel{Q} u$.
    \begin{subproof}
        Fix $s\in J$.
        Take $i,j$ such that $i\in\initialSegment{n}$ and $j\in\initialSegment{m + 1}$ and $s = (i,j)$
            by \cref{times_elem_is_tuple}.
        $j \leq m$ or it is wrong that $j \leq m$.
        \begin{byCase}
        \caseOf{$j \leq m$.}
            $j\in\naturalsPlus$ by \cref{initial_segment_naturalsplus}.
            Hence $j\in\initialSegment{m}$ by \cref{initial_segment_naturalsplus}.
            Let $V_i = \fst{\apply{a}{i}}$.
            Let $h_i = \snd{\apply{a}{i}}$.
            Let $A_i = \supp{\apply{p}{i}}{X}{T}$.
            $\apply{U}{i}$ is open in $X$ with respect to $T$ and
                $V_i$ is open in $\productFamily{R_0}$ with respect to $P_0$ and
                $h_i$ is a homeomorphism between $\apply{U}{i}$ and $V_i$ with respect to
                    $\subspaceTopology{T}{\apply{U}{i}}$ and $\subspaceTopology{P_0}{V_i}$ by assumption.
            Take $u$ such that
                $u$ is continuous from $X$ to $\reals$ with respect to $T$ and $\standardTopologyR$ and
                for all $x\in A_i$ we have
                    $u(x) = \apply{\apply{p}{i}}{x} \rmul \apply{\piIndex{R_0}{j}\circ h_i}{x}$ and
                for all $x\in X\setminus A_i$ we have $u(x) = 0$
                by \cref{weighted_chart_coordinate_extension}.
            Hence $u\subseteq X\times \reals$
                by \cref{continuous,funs,funs_subseteq_rels,subseteq,rels_ran_subseteq,rels_elim}.
            Therefore $u\in C_0$ by \cref{powerset_intro}.
            $(s,u)\in J\times C_0$ by \cref{times_tuple_intro}.
            Hence $(s,u)\in Q$ by \cref{pair_eq_iff,union_iff,union_intro_left}.
            Therefore $s\mathrel{Q} u$ by \cref{pair_eq_iff}.
        \caseOf{it is wrong that $j \leq m$.}
            We have $j = m + 1$ by \cref{initial_segment_naturalsplus,real_naturals_leq_succ_elim}.
            Take $u$ such that $u = j_0\circ \apply{p}{i}$.
            Hence $u\subseteq X\times \reals$
                by \cref{continuous,funs,subseteq,funs_subseteq_rels,rels_ran_subseteq,rels_elim}.
            Therefore $u\in C_0$ by \cref{powerset_intro}.
            $(s,u)\in J\times C_0$ by \cref{times_tuple_intro}.
            Hence $(s,u)\in Q$ by \cref{pair_eq_iff,union_iff,union_intro_right}.
            Therefore $s\mathrel{Q} u$ by \cref{pair_eq_iff}.
        \end{byCase}
    \end{subproof}
    $Q$ is a relation by \cref{relation,union_iff}.
    Take $c$ such that $c$ is a function on $J$ and for all $s\in J$ we have $s\mathrel{Q}\apply{c}{s}$
        by \cref{choice_function}.
    Show that for all $s\in J$ we have $\apply{c}{s}$ is continuous from $X$ to $\reals$ with respect to $T$ and $\standardTopologyR$.
    \begin{subproof}
        Fix $s\in J$.
        Hence $(s,\apply{c}{s})\in Q$ by \cref{choice_function}.
        Therefore $(s,\apply{c}{s})\in Q_1$ or $(s,\apply{c}{s})\in Q_2$ by \cref{union_iff}.
        \begin{byCase}
        \caseOf{$(s,\apply{c}{s})\in Q_1$.}
            $\apply{c}{s}$ is continuous from $X$ to $\reals$ with respect to $T$ and $\standardTopologyR$
                by \cref{pair_eq_iff}.
        \caseOf{$(s,\apply{c}{s})\in Q_2$.}
            There exist $i,j$ such that
                $s = (i,j)$ and
                $j = m + 1$ and
                $\apply{c}{s} = j_0\circ \apply{p}{i}$
                by \cref{pair_eq_iff}.
            Hence $s\in\initialSegment{n}\times\initialSegment{m + 1}$ by \cref{pair_eq_iff}.
            Therefore $i\in\initialSegment{n}$ and $j\in\initialSegment{m + 1}$ by \cref{pair_eq_iff,times_tuple_elim}.
            $j_0\circ \apply{p}{i}$ is continuous from $X$ to $\reals$ with respect to $T$ and $\standardTopologyR$ by \cref{subseteq}.
            Hence $\apply{c}{s}$ is continuous from $X$ to $\reals$ with respect to $T$ and $\standardTopologyR$
                by \cref{continuous_eq_subst}.
        \end{byCase}
    \end{subproof}

    Take $N,h,f$ such that
        $N\in\naturalsPlus$ and
        $h\in\inj{J}{\initialSegment{N}}$ and
        $f$ is a function on $\initialSegment{N}$ and
        for all $k\in\initialSegment{N}$ we have $\apply{f}{k}$ is continuous from $X$ to $\reals$ with respect to $T$ and $\standardTopologyR$ and
        $\finiteReindexAgreement{f}{h}{c}{J}$
        by \cref{finite_grid_continuous_reindex}.

    Let $R = \{(k,\reals) \mid k\in\initialSegment{N}\}$.
    Let $S = \{(k,\standardTopologyR) \mid k\in\initialSegment{N}\}$.
    Let $P = \productTopologyIndexed{S}{R}$.
    Let $F = \{(x,z) \mid x\in X, z\in\productFamily{R} \mid
        \forall k\in\initialSegment{N}. z(k) = \apply{\apply{f}{k}}{x}\}$.
    Show that $F$ is a function from $X$ to $\productFamily{R}$ and
        $F$ is continuous from $X$ to $\productFamily{R}$ with respect to $T$ and $P$ and
        for all $x\in X$ for all $k\in\initialSegment{N}$ we have $\apply{\apply{F}{x}}{k} = \apply{\apply{f}{k}}{x}$.
    \begin{subproof}
        Follows by \cref{finite_real_family_evaluation_map}.
    \end{subproof}

    $F$ is a function and $\dom{F} = X$ by \cref{funs_elim}.
    Show that for all $x,y$ such that $x\in\dom{F}$ and $y\in\dom{F}$ and $F(x) = F(y)$ we have $x = y$.
    \begin{subproof}
            Fix $x$.
            Fix $y$ such that $x\in\dom{F}$ and $y\in\dom{F}$ and $F(x) = F(y)$.
            Hence $x\in X$ and $y\in X$.
            Take $b$ such that
                $b\in\funs{\initialSegment{n}}{\reals}$ and
                for all $k\in\initialSegment{n}$ we have $b(k) = \apply{\apply{p}{k}}{x}$ and
                $\sumFiniteSequence{b}{1}$
                by \cref{finite_partition_of_unity_sum_witness}.
            Hence $b(1) = \apply{\apply{p}{1}}{x}$ and $\sumFiniteSequence{b}{1}$
                by \cref{real_one_in_naturals,real_one_leq_naturalsplus,initial_segment_naturalsplus}.
            For all $k\in\initialSegment{n}$ we have $0 \leq b(k)$
                by \cref{finite_partition_of_unity,finite_partition_component_condition,continuous,funs,funs_type_apply,intervalclosed}.
            Take $i\in\initialSegment{n}$ such that $0 < b(i)$ by \cref{finite_real_sum_one_positive_term}.
            Hence $0 < \apply{\apply{p}{i}}{x}$ by \cref{real_lt_subst_right}.
            Let $q = (i,m + 1)$.
            Therefore $q\in J$ by \cref{real_naturals_closed_succ,initial_segment_naturalsplus,times_tuple_intro}.
            $h(q)\in\initialSegment{N}$ by \cref{inj,funs_elim,subseteq,funs_type_apply}.
            We have $\apply{\apply{F}{x}}{h(q)} = \apply{\apply{f}{h(q)}}{x}$ and
                $\apply{\apply{F}{y}}{h(q)} = \apply{\apply{f}{h(q)}}{y}$ by assumption.
            Therefore $\apply{\apply{F}{x}}{h(q)} = \apply{\apply{F}{y}}{h(q)}$.
            We have $\apply{\apply{f}{h(q)}}{x} = \apply{\apply{f}{h(q)}}{y}$ by \cref{pair_eq_iff}.
            Hence $\apply{f}{h(q)} = \apply{c}{q}$ by \cref{finite_reindex_agreement}.
            Therefore $\apply{\apply{c}{q}}{x} = \apply{\apply{c}{q}}{y}$ by \cref{pair_eq_iff}.
            $(q,\apply{c}{q})\in Q$ by \cref{choice_function}.
            Therefore $(q,\apply{c}{q})\in Q_1$ or $(q,\apply{c}{q})\in Q_2$ by \cref{union_iff}.
            $(q,\apply{c}{q})\notin Q_1$
                by \cref{pair_eq_iff,initial_segment_naturalsplus,real_naturals_subset_reals,subseteq,real_naturals_closed_succ,real_add_closed,real_one_in_naturals,real_naturals_lt_succ_local,real_lt_imp_leq,real_leq_antisym,real_lt_subst_right,real_lt_irreflexive}.
            Hence $(q,\apply{c}{q})\in Q_2$ by \cref{union_iff}.
            Therefore there exist $i_0,j_2$ such that
                $q = (i_0,j_2)$ and
                $j_2 = m + 1$ and
                $\apply{c}{q} = j_0\circ \apply{p}{i_0}$
                by \cref{pair_eq_iff}.
            $\apply{c}{q} = j_0\circ \apply{p}{i}$ by \cref{pair_eq_iff}.
            Therefore $j_0\circ \apply{p}{i}\in\funs{X}{\reals}$ and
                $j_0\circ \apply{p}{i}$ is a function and
                $j_0\circ \apply{p}{i}$ is a relation and
                $\dom{j_0\circ \apply{p}{i}} = X$
                by \cref{continuous,funs_is_function,funs_is_relation,funs,subseteq}.
            $\apply{p}{i}$ is continuous from $X$ to $I_0$ with respect to $T$ and $T_0$
                by \cref{finite_partition_of_unity,finite_partition_component_condition}.
            $\apply{p}{i}\in\funs{X}{I_0}$ and
                $\apply{p}{i}$ is a function and
                $\apply{p}{i}$ is a relation and
                $\dom{\apply{p}{i}} = X$
                by \cref{continuous,funs_is_function,funs_is_relation,funs}.
            $j_0\circ \apply{p}{i}$ is continuous from $X$ to $\reals$ with respect to $T$ and $\standardTopologyR$ and
                for all $z\in X$ we have $\apply{j_0\circ \apply{p}{i}}{z} = \apply{\apply{p}{i}}{z}$ by \cref{subseteq}.
            Therefore $\apply{j_0\circ \apply{p}{i}}{x} = \apply{j_0\circ \apply{p}{i}}{y}$ by \cref{pair_eq_iff}.
            We have $\apply{j_0\circ \apply{p}{i}}{x} = \apply{\apply{p}{i}}{x}$ and
                $\apply{j_0\circ \apply{p}{i}}{y} = \apply{\apply{p}{i}}{y}$ by \cref{subseteq}.
            Hence $\apply{\apply{p}{i}}{x} = \apply{\apply{p}{i}}{y}$ by \cref{pair_eq_iff}.
            Therefore $0 < \apply{\apply{p}{i}}{y}$ by \cref{real_lt_subst_right}.

            Let $A_i = \supp{\apply{p}{i}}{X}{T}$.
            Therefore $x\in A_i$ and $y\in A_i$ by \cref{positive_partition_component_in_support}.
            $A_i\subseteq \apply{U}{i}$ by \cref{finite_partition_of_unity,finite_partition_component_condition}.
            Hence $x\in \apply{U}{i}$ and $y\in \apply{U}{i}$ by \cref{subseteq}.
            $\snd{\apply{a}{i}}$ is a homeomorphism between $\apply{U}{i}$ and $\fst{\apply{a}{i}}$ with respect to
                $\subspaceTopology{T}{\apply{U}{i}}$ and $\subspaceTopology{P_0}{\fst{\apply{a}{i}}}$ by assumption.
            Hence $\snd{\apply{a}{i}}\in\funs{\apply{U}{i}}{\fst{\apply{a}{i}}}$ and
                $\snd{\apply{a}{i}}$ is a function and
                $\dom{\snd{\apply{a}{i}}} = \apply{U}{i}$ and
                $\snd{\apply{a}{i}}$ is injective
                by \cref{homeomorphism,bijections,bijections_to_funs,funs_is_function,funs,injective_function}.
            $\fst{\apply{a}{i}}$ is open in $\productFamily{R_0}$ with respect to $P_0$ by assumption.
            Therefore $\fst{\apply{a}{i}}\subseteq \productFamily{R_0}$ by \cref{open_in,topology_on,subseteq,pow_iff}.
            We have $\apply{\snd{\apply{a}{i}}}{x}\in \productFamily{R_0}$ and $\apply{\snd{\apply{a}{i}}}{y}\in \productFamily{R_0}$
                by \cref{funs_type_apply,subseteq}.

            Show that for all $j\in\initialSegment{m}$ we have
                $\apply{\apply{\snd{\apply{a}{i}}}{x}}{j} = \apply{\apply{\snd{\apply{a}{i}}}{y}}{j}$.
            \begin{subproof}
                Fix $j\in\initialSegment{m}$.
                Let $s = (i,j)$.
                Therefore $s\in J$ by \cref{initial_segment_subset_succ,subseteq,times_tuple_intro,pair_eq_iff}.
                We have $h(s)\in\initialSegment{N}$ by \cref{inj,funs_elim,subseteq,funs_type_apply}.
                We have $\apply{\apply{F}{x}}{h(s)} = \apply{\apply{f}{h(s)}}{x}$ and
                    $\apply{\apply{F}{y}}{h(s)} = \apply{\apply{f}{h(s)}}{y}$ by assumption.
                Therefore $\apply{\apply{F}{x}}{h(s)} = \apply{\apply{F}{y}}{h(s)}$.
                We have $\apply{\apply{f}{h(s)}}{x} = \apply{\apply{f}{h(s)}}{y}$ by \cref{pair_eq_iff}.
                Hence $\apply{f}{h(s)} = \apply{c}{s}$ by \cref{finite_reindex_agreement}.
                Therefore $\apply{\apply{c}{s}}{x} = \apply{\apply{c}{s}}{y}$ by \cref{pair_eq_iff}.
                We have $(s,\apply{c}{s})\in Q$ by \cref{choice_function}.
                Therefore $(s,\apply{c}{s})\in Q_1$ or $(s,\apply{c}{s})\in Q_2$ by \cref{union_iff}.
                $(s,\apply{c}{s})\notin Q_2$
                    by \cref{pair_eq_iff,initial_segment_naturalsplus,real_naturals_subset_reals,subseteq,real_naturals_closed_succ,real_add_closed,real_one_in_naturals,real_naturals_lt_succ_local,real_lt_imp_leq,real_leq_antisym,real_lt_subst_right,real_lt_irreflexive}.
                Hence $(s,\apply{c}{s})\in Q_1$ by \cref{union_iff}.
                Therefore there exist $i_1,j_1$ such that
                    $s = (i_1,j_1)$ and
                    $j_1\in\initialSegment{m}$ and
                    $\apply{c}{s}$ is continuous from $X$ to $\reals$ with respect to $T$ and $\standardTopologyR$ and
                    for all $z\in \supp{\apply{p}{i_1}}{X}{T}$ we have
                        $\apply{\apply{c}{s}}{z} = \apply{\apply{p}{i_1}}{z} \rmul \apply{\piIndex{R_0}{j_1}\circ \snd{\apply{a}{i_1}}}{z}$ and
                    for all $z\in X\setminus \supp{\apply{p}{i_1}}{X}{T}$ we have $\apply{\apply{c}{s}}{z} = 0$
                    by \cref{pair_eq_iff}.
                We have $i_1 = i$ and $j_1 = j$ by \cref{pair_eq_iff}.
                Therefore $\apply{c}{s}$ is continuous from $X$ to $\reals$ with respect to $T$ and $\standardTopologyR$ and
                    for all $z\in \supp{\apply{p}{i}}{X}{T}$ we have
                        $\apply{\apply{c}{s}}{z} = \apply{\apply{p}{i}}{z} \rmul \apply{\piIndex{R_0}{j}\circ \snd{\apply{a}{i}}}{z}$ and
                    for all $z\in X\setminus \supp{\apply{p}{i}}{X}{T}$ we have $\apply{\apply{c}{s}}{z} = 0$
                    by \cref{pair_eq_iff}.
                We have $\apply{\apply{c}{s}}{x} = \apply{\apply{p}{i}}{x} \rmul \apply{\piIndex{R_0}{j}\circ \snd{\apply{a}{i}}}{x}$ and
                    $\apply{\apply{c}{s}}{y} = \apply{\apply{p}{i}}{y} \rmul \apply{\piIndex{R_0}{j}\circ \snd{\apply{a}{i}}}{y}$ by \cref{subseteq}.
                $\apply{\apply{p}{i}}{x}\in I_0$ and $\apply{\apply{p}{i}}{y}\in I_0$ by \cref{continuous,funs_type_apply}.
                Hence $\apply{\apply{p}{i}}{x}\in\reals$ and $\apply{\apply{p}{i}}{y}\in\reals$ by \cref{intervalclosed}.
                $R_0$ is a relation by \cref{relation}.
                For all $k,y_1,y_2$ such that $(k,y_1)\in R_0$ and $(k,y_2)\in R_0$ we have $y_1 = y_2$ by \cref{pair_eq_iff}.
                For all $k\in\dom{R_0}$ we have $k\in\initialSegment{m}$ by \cref{dom_iff,pair_eq_iff}.
                For all $k\in\initialSegment{m}$ we have $k\in\dom{R_0}$ by \cref{pair_eq_iff,dom_intro}.
                Therefore $j\in\dom{R_0}$.
                We have $R_0$ is right-unique by \cref{rightunique}.
                Hence $\piIndex{R_0}{j}\in\funs{\productFamily{R_0}}{\apply{R_0}{j}}$ by \cref{projection_map_indexed_function}.
                Therefore $\apply{R_0}{j} = \reals$ by \cref{function_apply_intro}.
                We have $\piIndex{R_0}{j}\in\funs{\productFamily{R_0}}{\reals}$ by \cref{pair_eq_iff}.
                $\snd{\apply{a}{i}}$ is a function and $\dom{\snd{\apply{a}{i}}} = \apply{U}{i}$ by \cref{homeomorphism,bijections,bijections_to_funs,funs_is_function,funs}.
                For all $z\in\dom{\snd{\apply{a}{i}}}$ we have $\apply{\snd{\apply{a}{i}}}{z}\in\productFamily{R_0}$
                    by \cref{homeomorphism,bijections,bijections_to_funs,funs_type_apply,subseteq}.
                Therefore $\snd{\apply{a}{i}}\in\funs{\apply{U}{i}}{\productFamily{R_0}}$ by \cref{funs_intro}.
                We have $\piIndex{R_0}{j}\circ \snd{\apply{a}{i}}\in\funs{\apply{U}{i}}{\reals}$ by \cref{funs_circ}.
                Hence $\piIndex{R_0}{j}\circ \snd{\apply{a}{i}}$ is a function and
                    $\dom{\piIndex{R_0}{j}\circ \snd{\apply{a}{i}}} = \apply{U}{i}$ by \cref{funs_is_function,funs}.
                Therefore $\apply{\piIndex{R_0}{j}\circ \snd{\apply{a}{i}}}{x}\in\reals$ and
                    $\apply{\piIndex{R_0}{j}\circ \snd{\apply{a}{i}}}{y}\in\reals$
                    by \cref{funs_type_apply,subseteq}.
                Hence $\apply{\apply{p}{i}}{x} = \apply{\apply{p}{i}}{y}$ by \cref{pair_eq_iff}.
                Therefore $\apply{\piIndex{R_0}{j}\circ \snd{\apply{a}{i}}}{x} =
                    \apply{\piIndex{R_0}{j}\circ \snd{\apply{a}{i}}}{y}$
                    by \cref{positive_real_factor_cancel}.
                $\snd{\apply{a}{i}}\in\funs{\apply{U}{i}}{\fst{\apply{a}{i}}}$ by \cref{homeomorphism,bijections,bijections_to_funs}.
                Therefore $(x,\apply{\snd{\apply{a}{i}}}{x})\in \snd{\apply{a}{i}}$ and
                    $(y,\apply{\snd{\apply{a}{i}}}{y})\in \snd{\apply{a}{i}}$ by \cref{function_apply_elim}.
                We have $(\apply{\snd{\apply{a}{i}}}{x},\apply{\apply{\snd{\apply{a}{i}}}{x}}{j})\in \piIndex{R_0}{j}$ and
                    $(\apply{\snd{\apply{a}{i}}}{y},\apply{\apply{\snd{\apply{a}{i}}}{y}}{j})\in \piIndex{R_0}{j}$
                    by \cref{projection_map_indexed,function_apply_elim,subseteq}.
                Hence $\apply{\piIndex{R_0}{j}\circ \snd{\apply{a}{i}}}{x} = \apply{\apply{\snd{\apply{a}{i}}}{x}}{j}$ and
                    $\apply{\piIndex{R_0}{j}\circ \snd{\apply{a}{i}}}{y} = \apply{\apply{\snd{\apply{a}{i}}}{y}}{j}$
                    by \cref{circ_iff,function_apply_intro}.
                Therefore $\apply{\apply{\snd{\apply{a}{i}}}{x}}{j} = \apply{\apply{\snd{\apply{a}{i}}}{y}}{j}$ by \cref{pair_eq_iff}.
            \end{subproof}
            Hence $\apply{\snd{\apply{a}{i}}}{x} = \apply{\snd{\apply{a}{i}}}{y}$ by \cref{finite_real_product_coordinate_extensionality}.
            Therefore $x = y$ by \cref{injective_function}.
    \end{subproof}
    Hence $F$ is injective by \cref{injective_function}.

    $F$ is a function and $\dom{F} = X$ by \cref{funs_elim}.
    $F\in\funs{X}{\img{F}{X}}$ by \cref{funs_intro,function_apply_elim,ran_intro,img_dom}.
    Hence $F\in\surj{X}{\img{F}{X}}$ by \cref{funs_surj_iff,img_dom}.
    Therefore $F$ is a bijection from $X$ to $\img{F}{X}$ by \cref{bijections,injective_function}.

    For all $k,y_1,y_2$ such that $(k,y_1)\in R$ and $(k,y_2)\in R$ we have $y_1 = y_2$ by \cref{pair_eq_iff}.
    Hence $R$ is a function on $\initialSegment{N}$ by
        \cref{relation,rightunique,dom_iff,pair_eq_iff,dom_intro,subseteq,subseteq_antisymmetric,funs}.
    $\standardTopologyR\in \{\standardTopologyR\}$ by \cref{singleton_intro}.
    Hence $S$ is a function from $\initialSegment{N}$ to $\{\standardTopologyR\}$ by \cref{constant_map_function}.
    Therefore $S$ is a function and $\dom{S} = \initialSegment{N}$ by \cref{funs_elim}.
    For all $k\in\initialSegment{N}$ we have $R(k) = \reals$ and $S(k) = \standardTopologyR$ by \cref{function_apply_intro}.
    We have $\reals$ is Hausdorff with respect to $\standardTopologyR$ by \cref{standard_metric_is_metric,metric_space_hausdorff,standard_metric_topology}.
    Hence $\productFamily{R}$ is Hausdorff with respect to $P$ by \cref{real_one_in_naturals,real_one_leq_naturalsplus,initial_segment_naturalsplus,subseteq,product_hausdorff_product}.
    We have $\img{F}{X}\subseteq \productFamily{R}$ by \cref{funs_elim,funs_subseteq_rels,subseteq,rels_ran_subseteq,img_subseteq_ran,subseteq_transitive}.
    $P$ is a topology on $\productFamily{R}$ by \cref{hausdorff}.
    Hence $\img{F}{X}$ is a subspace of $\productFamily{R}$ with respect to $P$ by \cref{subspace_of}.
    Therefore $\img{F}{X}$ is Hausdorff with respect to $\subspaceTopology{P}{\img{F}{X}}$ by \cref{subspace_hausdorff}.
    We have $F$ is continuous from $X$ to $\img{F}{X}$ with respect to $T$ and $\subspaceTopology{P}{\img{F}{X}}$
        by \cref{continuous_restrict_range,subseteq_refl}.
    Therefore $F$ is a homeomorphism between $X$ and $\img{F}{X}$ with respect to
        $T$ and $\subspaceTopology{P}{\img{F}{X}}$
        by \cref{compact_hausdorff_bijection_homeomorphism}.
    Hence $F$ is a topological embedding from $X$ to $\productFamily{R}$ with respect to $T$ and $P$
        by \cref{topological_embedding}.
    Take $n_1$ such that $n_1 = N$.
    Take $R_1$ such that $R_1 = R$.
    Take $S_1$ such that $S_1 = S$.
    Take $f_1$ such that $f_1 = F$.
    Follows by \cref{pair_eq_iff}.
\end{proof}

% from §36 Item 5: compact manifolds embed in euclidean spaces
\begin{theorem}\label{compact_manifold_embeds_in_euclidean_space}
    Assume $m\in\naturalsPlus$.
    Assume $X$ is a manifold of dimension $m$ with respect to $T$.
    Assume $X$ is compact with respect to $T$.
    Then there exist $n, R, S, f$ such that
        $n\in\naturalsPlus$ and
        $R = \{(i,\reals) \mid i\in\initialSegment{n}\}$ and
        $S = \{(i,\standardTopologyR) \mid i\in\initialSegment{n}\}$ and
        $f$ is a topological embedding from $X$ to $\productFamily{R}$ with respect to $T$ and $\productTopologyIndexed{S}{R}$.
\end{theorem}
\begin{proof}
    \begin{byCase}
        \caseOf{$X = \emptyset$.}
            Let $n = 1$.
            Let $R = \{(i,\reals) \mid i\in\initialSegment{n}\}$.
            Let $S = \{(i,\standardTopologyR) \mid i\in\initialSegment{n}\}$.
            Let $f = \emptyset$.
            Follows by \cref{compact_manifold_empty_embeds_in_euclidean_space}.
        \caseOf{$X\neq\emptyset$.}
            Follows by \cref{compact_manifold_nonempty_embeds_in_euclidean_space}.
    \end{byCase}
\end{proof}

% from §36 helper definition 8: maximal finite-intersection family on a set
\begin{definition}\label{maximal_finite_intersection_property_on}
    $D$ is maximal with respect to the finite intersection property on $X$ iff
        $D\subseteq\pow{X}$ and
        $D$ has the finite intersection property and
        for all $E$ such that $D\subseteq E\subseteq\pow{X}$ and $E$ has the finite intersection property
            we have $E = D$.
\end{definition}

% imported from §11: Zorn's lemma specialized to families of subsets
\begin{theorem}\label{zorn_subset_family}
    Assume $A\subseteq\pow{X}$.
    Suppose $A$ has the finite intersection property.
    Then there exists $D$ such that
        $A\subseteq D$ and
        $D$ is maximal with respect to the finite intersection property on $X$.
\end{theorem}
\begin{proof}
    Omitted. % do not touch
\end{proof}
\section{§37 The Tychonoff Theorem}

% from §37 helper lemma 1: basic consequences of maximality
\begin{lemma}\label{maximal_finite_intersection_property_basic}
    Assume $D$ is maximal with respect to the finite intersection property on $X$.
    Then $D\subseteq\pow{X}$ and $D$ has the finite intersection property.
\end{lemma}
\begin{proof}
    Show that $D\subseteq\pow{X}$.
    \begin{subproof}
        Follows by \cref{maximal_finite_intersection_property_on}.
    \end{subproof}
    Show that $D$ has the finite intersection property.
    \begin{subproof}
        Follows by \cref{maximal_finite_intersection_property_on}.
    \end{subproof}
\end{proof}

% from §37 Item 1: maximal extension of a family with the finite intersection property
\begin{lemma}\label{maximal_finite_intersection_property_extension}
    Assume $A\subseteq\pow{X}$.
    Suppose $A$ has the finite intersection property.
    Then there exists $D$ such that
        $A\subseteq D$ and
        $D$ is maximal with respect to the finite intersection property on $X$.
\end{lemma}
\begin{proof}
    Follows by \cref{zorn_subset_family}.
\end{proof}

% from §37 Item 2: finite intersections belong to maximal families with the finite intersection property
\begin{lemma}\label{maximal_finite_intersection_property_finite_intersections}
    Assume $D$ is maximal with respect to the finite intersection property on $X$.
    Assume $F\subseteq D$.
    Assume $F$ is finite and inhabited.
    Then $\inters{F}\in D$.
\end{lemma}
\begin{proof}
    Let $B = \inters{F}$.
    $D\subseteq\pow{X}$ and $D$ has the finite intersection property
        by \cref{maximal_finite_intersection_property_basic}.
    Take $D_0$ such that $D_0\in F$.
    Let $E = D\union\{B\}$.
    Hence $E\subseteq\pow{X}$
        by \cref{inters_subseteq_elem,subseteq,pow_iff,subseteq_transitive,powerset_intro,singleton_subset_intro,union_subsets_is_subset}.
    Take $D_1$ such that $D_1\in D$ by \cref{finite_intersection_property}.
    Therefore $E$ is inhabited by \cref{union_intro_left}.
    Show that for all $G$ such that $G\subseteq E$ and $G$ is finite and inhabited
        we have $\inters{G}\neq\emptyset$.
    \begin{subproof}
        Fix $G$ such that $G\subseteq E$ and $G$ is finite and inhabited.
        \begin{byCase}
        \caseOf{$B\notin G$.}
            $\inters{G}\neq\emptyset$
                by \cref{subseteq,union_iff,singleton_iff,finite_intersection_property}.
        \caseOf{$B\in G$.}
            Let $H = (G\setminus\{B\})\union F$.
            Hence $G\setminus\{B\}$ is finite by \cref{setminus_subseteq,subseteq_finite_is_finite}.
            Therefore $H$ is finite by \cref{union_of_finite_is_finite}.
            Take $C_0$ such that $C_0\in F$.
            We have $H$ is inhabited by \cref{union_intro_right}.
            Therefore $\inters{H}\neq\emptyset$
                by \cref{union_iff,setminus_elim_left,setminus_elim_right,subseteq,finite_intersection_property}.
            Take $x$ such that $x\in\inters{H}$ by \cref{nonempty_inhabited}.
            For all $C$ such that $C\in G$ we have $x\in C$
                by \cref{union_intro_right,inters_destr,inters_intro,singleton_iff,setminus_intro,union_intro_left}.
            Hence $x\in\inters{G}$ by \cref{inters_intro}.
            Therefore $\inters{G}\neq\emptyset$ by \cref{emptyset}.
        \end{byCase}
    \end{subproof}
    Hence $E$ has the finite intersection property by \cref{finite_intersection_property}.
    Therefore $E = D$ by \cref{maximal_finite_intersection_property_on,union_upper_left}.
    Hence $B\in D$ by \cref{singleton_intro,union_intro_right}.
\end{proof}

% from §37 Item 2: sets intersecting every member belong to maximal families with the finite intersection property
\begin{lemma}\label{maximal_finite_intersection_property_intersection_test}
    Assume $D$ is maximal with respect to the finite intersection property on $X$.
    Assume $A\subseteq X$.
    Suppose for all $E\in D$ we have $A$ intersects $E$.
    Then $A\in D$.
\end{lemma}
\begin{proof}
    $D\subseteq\pow{X}$ and $D$ has the finite intersection property
        by \cref{maximal_finite_intersection_property_basic}.
    Let $E = D\union\{A\}$.
    Hence $E\subseteq\pow{X}$ by \cref{powerset_intro,singleton_subset_intro,union_subsets_is_subset}.
    Take $D_0$ such that $D_0\in D$ by \cref{finite_intersection_property}.
    Therefore $A$ intersects $D_0$.
    Take $x$ such that $x\in A\inter D_0$ by \cref{intersects,nonempty_inhabited}.
    Hence $A$ is inhabited by \cref{inter}.
    Therefore $E$ is inhabited by \cref{union_intro_left}.
    Show that for all $G$ such that $G\subseteq E$ and $G$ is finite and inhabited
        we have $\inters{G}\neq\emptyset$.
    \begin{subproof}
        Fix $G$ such that $G\subseteq E$ and $G$ is finite and inhabited.
        \begin{byCase}
        \caseOf{$A\notin G$.}
            For all $C$ such that $C\in G$ we have $C\in D$.
            $\inters{G}\neq\emptyset$ by \cref{subseteq,finite_intersection_property}.
        \caseOf{$A\in G$.}
            Let $H = G\setminus\{A\}$.
            Hence $H$ is finite by \cref{setminus_subseteq,subseteq_finite_is_finite}.
            Show that for all $C$ such that $C\in H$ we have $C\in D$.
            \begin{subproof}
                Fix $C$ such that $C\in H$.
                $C\in G$ and $C\notin\{A\}$ by \cref{setminus_elim_left,setminus_elim_right}.
                $C\in E$ by \cref{subseteq}.
                $C\in D$ or $C\in\{A\}$ by \cref{union_iff}.
                $C\in D$ by \cref{union_iff}.
            \end{subproof}
            \begin{byCase}
            \caseOf{$H = \emptyset$.}
                For all $C$ such that $C\in G$ we have $C\in\{A\}$.
                Hence $G\subseteq\{A\}$ by \cref{subseteq}.
                Therefore $G = \{A\}$ by \cref{singleton_subset_intro,subseteq_antisymmetric}.
                We have $\inters{G} = A$ by \cref{inters_singleton}.
                Therefore $\inters{G}\neq\emptyset$ by \cref{emptyset}.
            \caseOf{$H\neq\emptyset$.}
                $\inters{H}\in D$ by \cref{nonempty_inhabited,subseteq,maximal_finite_intersection_property_finite_intersections}.
                Hence $A$ intersects $\inters{H}$.
                Take $y$ such that $y\in A\inter\inters{H}$ by \cref{intersects,nonempty_inhabited}.
                For all $C$ such that $C\in G$ we have $y\in C$
                    by \cref{inter,singleton_iff,setminus_intro,inters_destr}.
                Hence $y\in\inters{G}$ by \cref{inters_intro}.
                Therefore $\inters{G}\neq\emptyset$ by \cref{emptyset}.
            \end{byCase}
        \end{byCase}
    \end{subproof}
    Hence $E$ has the finite intersection property by \cref{finite_intersection_property}.
    Therefore $E = D$ by \cref{maximal_finite_intersection_property_on,union_upper_left}.
    Hence $A\in D$ by \cref{singleton_intro,union_intro_right}.
\end{proof}

% from §37 helper lemma 2: coordinate closure family has the finite intersection property
\begin{lemma}\label{coordinate_closure_family_fip}
    Assume $X$ is a function.
    Assume $T$ is a function.
    Assume $\dom{T} = \dom{X}$.
    Assume $\alpha\in\dom{X}$.
    Assume that for all $\beta\in\dom{X}$ we have $\apply{T}{\beta}$ is a topology on $\apply{X}{\beta}$.
    Assume $D$ is maximal with respect to the finite intersection property on $\productFamily{X}$.
    Let $C = \{B\in\pow{\apply{X}{\alpha}} \mid \exists E\in D. B = \closure{\img{\piIndex{X}{\alpha}}{E}}{\apply{X}{\alpha}}{\apply{T}{\alpha}}\}$.
    Then $C$ has the finite intersection property.
\end{lemma}
\begin{proof}
    $D\subseteq\pow{\productFamily{X}}$ and $D$ has the finite intersection property
        by \cref{maximal_finite_intersection_property_basic}.
    Let $P = \piIndex{X}{\alpha}$.
    $P\in\funs{\productFamily{X}}{\apply{X}{\alpha}}$ by \cref{projection_map_indexed_function}.
    Take $E_0$ such that $E_0\in D$ by \cref{finite_intersection_property}.
    Let $B_0 = \closure{\img{P}{E_0}}{\apply{X}{\alpha}}{\apply{T}{\alpha}}$.
    Hence $B_0\in\pow{\apply{X}{\alpha}}$
        by \cref{img_subseteq_ran,funs_ran,subseteq_transitive,closure_closed_in,closed_in,powerset_intro}.
    Therefore $C$ is inhabited by \cref{powerset_intro}.
    Show that for all $F$ such that $F\subseteq C$ and $F$ is finite and inhabited
        we have $\inters{F}\neq\emptyset$.
    \begin{subproof}
        Fix $F$ such that $F\subseteq C$ and $F$ is finite and inhabited.
        Let $R = \{w\in F\times D \mid
            \fst{w} = \closure{\img{P}{\snd{w}}}{\apply{X}{\alpha}}{\apply{T}{\alpha}}\}$.
        For all $B\in F$ there exists $E$ such that $B\mathrel{R}E$
            by \cref{subseteq,times_tuple_intro,fst_eq,snd_eq}.
        Take $h$ such that $h$ is a function on $F$ and for all $B\in F$ we have $B\mathrel{R}\apply{h}{B}$
            by \cref{times_elem_is_tuple,relation,choice_function}.
        Let $G = \img{h}{F}$.
        $G$ is finite by \cref{finite_image_function}.
        Hence $G\subseteq D$ by \cref{img_iff,function_apply_iff,subseteq,times_tuple_elim}.
        Take $B_1$ such that $B_1\in F$.
        We have $G$ is inhabited by \cref{function_apply_iff,img_iff}.
        Therefore $\inters{G}\neq\emptyset$ by \cref{finite_intersection_property}.
        Take $y$ such that $y\in\inters{G}$ by \cref{nonempty_inhabited}.
        Take $E_1$ such that $E_1\in G$.
        Hence $y\in\productFamily{X}$ by \cref{inters_destr,subseteq,pow_iff}.
        For all $B$ such that $B\in F$ we have $\apply{y}{\alpha}\in B$
            by \cref{choice_function,subseteq,fst_eq,snd_eq,img_iff,function_apply_iff,inters_destr,projection_map_indexed,pair_eq_iff,img_subseteq_ran,funs_ran,subseteq_transitive,subseteq_closure}.
        Hence $\apply{y}{\alpha}\in\inters{F}$ by \cref{inters_intro}.
        Therefore $\inters{F}\neq\emptyset$ by \cref{emptyset}.
    \end{subproof}
    Hence $C$ has the finite intersection property by \cref{finite_intersection_property}.
\end{proof}

% from §37 Item 3: Tychonoff theorem
\begin{theorem}\label{tychonoff_theorem}
    Assume $X$ is a function.
    Assume $T$ is a function.
    Assume $\dom{T} = \dom{X}$.
    Assume that for all $\alpha\in\dom{X}$ we have
        $\apply{T}{\alpha}$ is a topology on $\apply{X}{\alpha}$ and
        $\apply{X}{\alpha}$ is compact with respect to $\apply{T}{\alpha}$.
    Then $\productFamily{X}$ is compact with respect to $\productTopologyIndexed{T}{X}$.
\end{theorem}
\begin{proof}
    Let $T_0 = \productTopologyIndexed{T}{X}$.
    \begin{byCase}
        \caseOf{$\dom{X} = \emptyset$.}
            $\emptyset$ is a function to $\unions{\ran{X}}$ by \cref{codom_of_emptyset_can_be_anything}.
            Hence $\emptyset\in\funs{\dom{X}}{\unions{\ran{X}}}$ by \cref{dom_emptyset,funs_intro}.
            For all $\alpha$ such that $\alpha\in\dom{X}$ we have $\apply{\emptyset}{\alpha}\in\apply{X}{\alpha}$ by \cref{emptyset}.
            Therefore $\emptyset\in\productFamily{X}$ by \cref{indexed_product}.
            We have $\{\emptyset\}\subseteq\productFamily{X}$ by \cref{singleton_subset_intro}.
            For all $y$ such that $y\in\productFamily{X}$ we have $y\in\{\emptyset\}$
                by \cref{indexed_product,funs_elim,funs_is_relation,relation,dom_intro,subseteq,emptyset,emptyset_subseteq,subseteq_antisymmetric,singleton_intro}.
            Therefore $\productFamily{X} = \{\emptyset\}$ by \cref{subseteq,singleton_subset_intro,subseteq_antisymmetric}.
            Show that for all $A$ such that $A$ is an open covering of $\productFamily{X}$ with respect to $T_0$
                there exists $B$ such that $B\subseteq A$ and $B$ is finite and $B$ is a covering of $\productFamily{X}$.
            \begin{subproof}
                Fix $A$ such that $A$ is an open covering of $\productFamily{X}$ with respect to $T_0$.
                Hence $\emptyset\in\unions{A}$ by \cref{open_covering,covering,subseteq,singleton_iff}.
                Take $U$ such that $U\in A$ and $\emptyset\in U$ by \cref{unions_iff}.
                Let $C = \{U\}$.
                Show that $C\subseteq A$ and $C$ is finite.
                \begin{subproof}
                    Follows by \cref{singleton_subset_intro,pair_is_finite,upair_intro_left,subseteq_finite_is_finite}.
                \end{subproof}
                Show that $C$ is a covering of $\productFamily{X}$.
                \begin{subproof}
                    Follows by \cref{subseteq,singleton_iff,singleton_intro,unions_iff,covering}.
                \end{subproof}
                Thus there exists $B$ such that $B\subseteq A$ and $B$ is finite and $B$ is a covering of $\productFamily{X}$.
            \end{subproof}
            Therefore $\productFamily{X}$ is compact with respect to $T_0$ by \cref{compact}.
        \caseOf{$\dom{X}\neq\emptyset$.}
            We have $\dom{X}$ is inhabited by \cref{nonempty_inhabited}.
            Let $S = \productSubbasis{T}{X}$.
            Let $B = \finiteIntersections{S}$.
            $B$ is a basis for $T_0$ on $\productFamily{X}$ by \cref{product_subbasis_finite_intersections_basis}.
            Hence $T_0$ is a topology on $\productFamily{X}$ by \cref{basis_for_topology,basis_topology_is_topology}.
            Show that for all $A$ if
                $A$ has the finite intersection property and
                for all $E$ such that $E\in A$ we have $E$ is closed in $\productFamily{X}$ with respect to $T_0$
                then $\inters{A}\neq\emptyset$.
            \begin{subproof}
                Fix $A$ such that $A$ has the finite intersection property and
                    for all $E$ such that $E\in A$ we have $E$ is closed in $\productFamily{X}$ with respect to $T_0$.
                Hence $A\subseteq\pow{\productFamily{X}}$ by \cref{closed_in,powerset_intro,subseteq}.
                Take $D$ such that $A\subseteq D$ and $D$ is maximal with respect to the finite intersection property on $\productFamily{X}$
                    by \cref{maximal_finite_intersection_property_extension}.
                $D\subseteq\pow{\productFamily{X}}$ and $D$ has the finite intersection property
                    by \cref{maximal_finite_intersection_property_basic}.
                Let $R = \{w\in\dom{X}\times\unions{\ran{X}} \mid \text{$\snd{w}\in\apply{X}{\fst{w}}$ and
                    for all $E\in D$ we have
                    $\snd{w}\in\closure{\img{\piIndex{X}{\fst{w}}}{E}}{\apply{X}{\fst{w}}}{\apply{T}{\fst{w}}}$}\}$.
                Show that for all $\alpha\in\dom{X}$ there exists $a$ such that $\alpha\mathrel{R} a$.
                \begin{subproof}
                    Fix $\alpha$ such that $\alpha\in\dom{X}$.
                    Let $C = \{B\in\pow{\apply{X}{\alpha}} \mid
                        \exists E\in D. B = \closure{\img{\piIndex{X}{\alpha}}{E}}{\apply{X}{\alpha}}{\apply{T}{\alpha}}\}$.
                    $C$ has the finite intersection property by \cref{coordinate_closure_family_fip}.
                    Show that for all $B_0$ such that $B_0\in C$ we have
                        $B_0$ is closed in $\apply{X}{\alpha}$ with respect to $\apply{T}{\alpha}$.
                    \begin{subproof}
                        Follows by \cref{subseteq,projection_map_indexed_function,img_subseteq_ran,funs_ran,subseteq_transitive,closure_closed_in}.
                    \end{subproof}
                    Therefore $\inters{C}\neq\emptyset$ by \cref{compact_iff_finite_intersection_property}.
                    Take $a$ such that $a\in\inters{C}$ by \cref{nonempty_inhabited}.
                    Hence $a\in\apply{X}{\alpha}$ by \cref{finite_intersection_property,inters_destr,subseteq,pow_iff}.
                    For all $E\in D$ we have
                        $a\in\closure{\img{\piIndex{X}{\alpha}}{E}}{\apply{X}{\alpha}}{\apply{T}{\alpha}}$
                        by \cref{projection_map_indexed_function,img_subseteq_ran,funs_ran,subseteq_transitive,closure_closed_in,closed_in,powerset_intro,inters_destr}.
                    Hence $a\in\unions{\ran{X}}$ by \cref{function_apply_elim,ran_intro,unions_iff}.
                    Hence $\alpha\mathrel{R} a$ by \cref{times_tuple_intro,fst_eq,snd_eq}.
                \end{subproof}
                Take $x$ such that $x$ is a function on $\dom{X}$ and for all $\alpha\in\dom{X}$ we have
                    $\alpha\mathrel{R}\apply{x}{\alpha}$ by \cref{times_elem_is_tuple,relation,choice_function}.
                For all $\alpha\in\dom{X}$ we have $\apply{x}{\alpha}\in\apply{X}{\alpha}$
                    by \cref{choice_function,subseteq,fst_eq,snd_eq}.
                $x$ is a function and $\dom{x} = \dom{X}$.
                For all $\alpha$ such that $\alpha\in\dom{x}$ we have $\apply{x}{\alpha}\in\unions{\ran{X}}$
                    by \cref{subseteq,function_apply_elim,ran_intro,unions_iff}.
                Therefore $x\in\funs{\dom{X}}{\unions{\ran{X}}}$ by \cref{funs_intro}.
                Hence $x\in\productFamily{X}$ by \cref{indexed_product}.
                Show that for all $E$ such that $E\in D$ we have $x\in\closure{E}{\productFamily{X}}{T_0}$.
                \begin{subproof}
                    Fix $E$ such that $E\in D$.
                    Therefore $E\subseteq\productFamily{X}$ by \cref{subseteq,pow_iff}.
                    Show that for all $M$ such that $M\in B$ and $x\in M$ we have $M$ intersects $E$.
                    \begin{subproof}
                        Fix $M$ such that $M\in B$ and $x\in M$.
                        Take $F\subseteq S$ such that $F$ is finite and inhabited and $M = \inters{F}$ by \cref{finite_intersections}.
                        Show that for all $W$ such that $W\in F$ we have $W\in D$.
                        \begin{subproof}
                            Fix $W$ such that $W\in F$.
                            Take $\alpha\in\dom{X}$ such that $W\in\productSubbasisAt{T}{X}{\alpha}$
                                by \cref{subseteq,product_subbasis_union_indexed}.
                            Take $U\in\apply{T}{\alpha}$ such that $W = \preimg{\piIndex{X}{\alpha}}{U}$
                                by \cref{product_subbasis_indexed}.
                            Hence $\apply{x}{\alpha}\in U$ by \cref{inters_destr,preimg_iff,projection_map_indexed,pair_eq_iff}.
                            Show that for all $E_0\in D$ we have $W$ intersects $E_0$.
                            \begin{subproof}
                                Fix $E_0$ such that $E_0\in D$.
                                We have $\apply{x}{\alpha}\in
                                    \closure{\img{\piIndex{X}{\alpha}}{E_0}}{\apply{X}{\alpha}}{\apply{T}{\alpha}}$
                                    by \cref{choice_function,subseteq,fst_eq,snd_eq}.
                                Hence $U$ intersects $\img{\piIndex{X}{\alpha}}{E_0}$ by
                                    \cref{open_in,topology_on,projection_map_indexed_function,img_subseteq_ran,funs_ran,subseteq_transitive,closure_iff_intersects_open}.
                                Take $a$ such that $a\in U\inter\img{\piIndex{X}{\alpha}}{E_0}$ by \cref{intersects,nonempty_inhabited}.
                                Take $y$ such that $y\in E_0$ and $(y,a)\in\piIndex{X}{\alpha}$ by \cref{inter,img_iff}.
                                Therefore $\apply{y}{\alpha}\in U$ by \cref{projection_map_indexed,pair_eq_iff,inter}.
                                Hence $y\in W$ by \cref{preimg_iff,projection_map_indexed,pair_eq_iff}.
                                Hence $y\in W\inter E_0$ by \cref{inter_intro}.
                                Therefore $W$ intersects $E_0$ by \cref{emptyset,intersects}.
                            \end{subproof}
                            Therefore $W\in D$ by \cref{product_subbasis_indexed,subseteq,pow_iff,maximal_finite_intersection_property_intersection_test}.
                        \end{subproof}
                        Hence $M\in D$ by \cref{subseteq,maximal_finite_intersection_property_finite_intersections}.
                        Let $H = \{M,E\}$.
                        $H$ is finite and inhabited by \cref{pair_is_finite,upair_intro_left}.
                        Hence $\inters{H}\neq\emptyset$ by \cref{upair_iff,subseteq,finite_intersection_property}.
                        Therefore $M$ intersects $E$ by \cref{inter_as_inters,emptyset,intersects}.
                    \end{subproof}
                    Therefore $x\in\closure{E}{\productFamily{X}}{T_0}$
                        by \cref{closure_iff_intersects_basis}.
                \end{subproof}
                Show that for all $E$ such that $E\in A$ we have $x\in E$.
                \begin{subproof}
                    Follows by \cref{subseteq,closed_eq_closure}.
                \end{subproof}
                Hence $x\in\inters{A}$ by \cref{finite_intersection_property,inters_intro}.
                Therefore $\inters{A}\neq\emptyset$ by \cref{emptyset}.
            \end{subproof}
            Therefore $\productFamily{X}$ is compact with respect to $T_0$ by \cref{compact_iff_finite_intersection_property}.
    \end{byCase}
\end{proof}
\section{§38 The Stone-Cech Compactification}

% from §38 helper definition 1: compactification in the sense of §38
\begin{definition}\label{compactifies}
    $Y$ compactifies $X$ with respect to $S$ and $T$ iff
        $S$ is a topology on $Y$ and
        $T$ is a topology on $X$ and
        $X$ is a subspace of $Y$ with respect to $S$ and
        $T = \subspaceTopology{S}{X}$ and
        $\closure{X}{Y}{S} = Y$ and
        $Y$ is compact with respect to $S$ and
        $Y$ is Hausdorff with respect to $S$.
\end{definition}

% from §38 Item 1: equivalence of compactifications
\begin{definition}\label{compactification_equivalence_over_x}
    $Y_1$ is equivalent to $Y_2$ over $X$ with respect to $S_1$ and $S_2$ and $T$ iff
        $Y_1$ compactifies $X$ with respect to $S_1$ and $T$ and
        $Y_2$ compactifies $X$ with respect to $S_2$ and $T$ and
        there exists $h$ such that
            $h$ is a homeomorphism between $Y_1$ and $Y_2$ with respect to $S_1$ and $S_2$ and
            for all $x\in X$ we have $h(x) = x$.
\end{definition}

% from §38 helper lemma 1: no two-element membership cycle
\begin{lemma}\label{no_two_cycle}
    Suppose $x\in y$.
    Then $y\notin x$.
\end{lemma}
\begin{proof}
    Suppose not.
    Let $A = \{x,y\}$.
    Take $a$ such that $a$ is an \in-minimal element of $A$ by \cref{upair_iff,regularity}.
    Hence $a = x$ or $a = y$ by \cref{upair_iff}.
    \begin{byCase}
    \caseOf{$a = x$.}
        We have $y\in a$.
        Therefore $y\in A$ by \cref{upair_iff}.
        Thus $a$ intersects $A$.
        Contradiction.
    \caseOf{$a = y$.}
        We have $x\in a$.
        Therefore $x\in A$ by \cref{upair_iff}.
        Thus $a$ intersects $A$.
        Contradiction.
    \end{byCase}
\end{proof}

% from §38 helper lemma 2: tagging by an unordered pair is injective
\begin{lemma}\label{tag_with_base_injective}
    Suppose $\{a,p\} = \{b,p\}$.
    Then $a = b$.
\end{lemma}
\begin{proof}
    \begin{byCase}
    \caseOf{$a = p$.}
        We have $a = b$ by \cref{upair_iff,subseteq,pair_eq_iff}.
    \caseOf{$a \neq p$.}
        We have $a = b$ by \cref{upair_iff,subseteq}.
    \end{byCase}
\end{proof}

% from §38 Item 2: existence of the compactification induced by an embedding
\begin{lemma}\label{induced_compactification_existence}
    Assume $h$ is a topological embedding from $X$ to $Z$ with respect to $T$ and $U$.
    Assume $Z$ is compact with respect to $U$.
    Assume $Z$ is Hausdorff with respect to $U$.
    Then there exist $Y, S, H$ such that
        $Y$ compactifies $X$ with respect to $S$ and $T$ and
        $H$ is a topological embedding from $Y$ to $Z$ with respect to $S$ and $U$ and
        for all $x\in X$ we have $H(x) = h(x)$.
\end{lemma}
\begin{proof}
    Let $X_0 = \img{h}{X}$.
    Let $Y_0 = \closure{X_0}{Z}{U}$.
    Let $T_0 = \subspaceTopology{U}{Y_0}$.
    Let $p = X$.
    Hence $p\notin X$ by \cref{set_notin_self}.
    Let $A = \{\{z,p\} \mid z\in Y_0\setminus X_0\}$.
    Let $Y = X\union A$.
    Let $g = \{w\in A\times Y_0 \mid \text{there exists $z\in Y_0\setminus X_0$ such that $w = (\{z,p\},z)$}\}$.
    Let $H = h\union g$.

    $T$ is a topology on $X$ and $U$ is a topology on $Z$ by \cref{topological_embedding,continuous}.
    $h\in\funs{X}{Z}$ and $h$ is a function and $h$ is injective by \cref{topological_embedding,funs_is_function}.
    Hence $X_0\subseteq Z$ by \cref{img_subseteq_ran,funs_ran,subseteq_transitive}.
    $Y_0\subseteq Z$ by \cref{closure,subseteq}.
    $Y_0$ is closed in $Z$ with respect to $U$ by \cref{closure_closed_in}.
    Hence $Y_0$ is compact with respect to $T_0$ by \cref{closed_subspace_compact}.
    $Y_0$ is Hausdorff with respect to $T_0$ by \cref{subspace_hausdorff}.
    $Y_0$ is a subspace of $Z$ with respect to $U$ by \cref{subspace_of,closure}.
    $T_0$ is a topology on $Y_0$ by \cref{subspace_topology_is_topology,subspace_of,closure}.
    $X_0\subseteq Y_0$ by \cref{subseteq_closure}.
    Hence $X_0$ is a subspace of $Y_0$ with respect to $T_0$ by \cref{subspace_of,subseteq_transitive}.
    Therefore $h$ is a homeomorphism between $X$ and $X_0$ with respect to $T$ and $\subspaceTopology{T_0}{X_0}$
        by \cref{topological_embedding,subspace_subspace_topology,subspace_of}.

    For all $u,v_1,v_2$ such that $(u,v_1)\in g$ and $(u,v_2)\in g$ we have $v_1 = v_2$
        by \cref{pair_eq_iff,tag_with_base_injective}.
    Hence $g$ is right-unique.
    Show that for all $u$ such that $u\in\dom{g}$ we have $u\in A$.
    \begin{subproof}
        Follows by \cref{dom_iff,times,times_tuple_elim}.
    \end{subproof}
    We have $\dom{g}\subseteq A$ by \cref{subseteq}.
    Show that for all $u$ such that $u\in A$ we have $u\in\dom{g}$.
    \begin{subproof}
        Follows by \cref{subseteq,setminus_elim_left,times_tuple_intro,dom_iff}.
    \end{subproof}
    We have $A\subseteq \dom{g}$ by \cref{subseteq}.
    Therefore $\dom{g} = A$ by \cref{subseteq,subseteq_antisymmetric}.
    For all $w$ such that $w\in g$ we have $w\in A\times Y_0$ by \cref{times}.
    $g$ is a relation by \cref{subseteq,relation}.
    For all $u$ such that $u\in\dom{g}$ we have $g(u)\in Y_0$
        by \cref{dom_iff,times,times_tuple_elim,function_apply_iff,rightunique,relation}.
    Therefore $g\in\funs{A}{Y_0}$ by \cref{funs_intro}.

    Show that $A$ is disjoint from $X$.
    \begin{subproof}
        Suppose not.
        Take $x$ such that $x\in A$ and $x\in X$ by \cref{disjoint}.
        Take $z\in Y_0\setminus X_0$ such that $x = \{z,p\}$ by \cref{subseteq}.
        Contradiction by \cref{upair_iff,no_two_cycle}.
    \end{subproof}

    For all $u$ such that $u\in\dom{h}\inter\dom{g}$ we have $h(u) = g(u)$
        by \cref{funs,inter,subseteq,dom_iff,disjoint}.
    Hence $H$ is a function by \cref{funs_is_function,union_compatible_functions}.
    Therefore $\dom{H} = Y$ by \cref{funs,subseteq_antisymmetric,dom_union}.
    For all $w$ such that $w\in h$ we have $w\in Y\times Y_0$
        by \cref{funs,rels_elim,subseteq,times_elem_is_tuple,img_iff,function_apply_intro,topological_embedding,funs_is_function,union_intro_left,times_tuple_intro}.
    For all $w$ such that $w\in g$ we have $w\in Y\times Y_0$
        by \cref{times,times_elem_is_tuple,union_intro_right,times_tuple_intro}.
    Thus for all $w$ such that $w\in H$ we have $w\in Y\times Y_0$ by \cref{union_iff}.
    For all $y$ such that $y\in\dom{H}$ we have $H(y)\in Y_0$
        by \cref{dom_iff,times,times_tuple_elim,function_apply_iff,rightunique,relation}.
    Therefore $H\in\funs{Y}{Y_0}$ by \cref{funs_intro,funs}.

    For all $x\in X$ we have $H(x) = h(x)$
        by \cref{topological_embedding,funs,function_apply_iff,funs_is_function,union_iff,subseteq}.

    Show that for all $y_1,y_2$ such that $y_1\in\dom{H}$ and $y_2\in\dom{H}$ and $H(y_1) = H(y_2)$
        we have $y_1 = y_2$.
    \begin{subproof}
        Fix $y_1$.
        Fix $y_2$ such that $y_1\in\dom{H}$ and $y_2\in\dom{H}$ and $H(y_1) = H(y_2)$.
        We have $y_1\in Y$ and $y_2\in Y$ by \cref{funs,subseteq}.
        Hence $y_1\in X$ or $y_1\in A$ by \cref{union_iff}.
        $y_2\in X$ or $y_2\in A$ by \cref{union_iff}.
        \begin{byCase}
        \caseOf{$y_1\in X$ and $y_2\in X$.}
            We have $H(y_1) = h(y_1)$ and $H(y_2) = h(y_2)$ by \cref{subseteq}.
            Hence $h(y_1) = h(y_2)$.
            Hence $y_1\in\dom{h}$ and $y_2\in\dom{h}$ by \cref{topological_embedding,funs}.
            Therefore $y_1 = y_2$ by \cref{injective_function,topological_embedding}.
        \caseOf{$y_1\in X$ and $y_2\in A$.}
            Suppose not.
            Hence $H(y_1)\in X_0$
                by \cref{subseteq,function_apply_iff,topological_embedding,funs,funs_is_function,img_iff}.
            Take $z\in Y_0\setminus X_0$ such that $y_2 = \{z,p\}$ by \cref{subseteq}.
            Therefore $(y_2,z)\in H$ by \cref{setminus_elim_left,times_tuple_intro,union_iff}.
            Hence $H(y_2) = z$ by \cref{function_apply_iff,funs_is_function}.
            Hence $H(y_2)\in Y_0\setminus X_0$.
            Contradiction by \cref{setminus_elim_right}.
        \caseOf{$y_1\in A$ and $y_2\in X$.}
            Suppose not.
            Hence $H(y_2)\in X_0$
                by \cref{subseteq,function_apply_iff,topological_embedding,funs,funs_is_function,img_iff}.
            Take $z\in Y_0\setminus X_0$ such that $y_1 = \{z,p\}$ by \cref{subseteq}.
            Therefore $(y_1,z)\in H$ by \cref{setminus_elim_left,times_tuple_intro,union_iff}.
            Hence $H(y_1) = z$ by \cref{function_apply_iff,funs_is_function}.
            Hence $H(y_1)\in Y_0\setminus X_0$.
            Contradiction by \cref{setminus_elim_right}.
        \caseOf{$y_1\in A$ and $y_2\in A$.}
            Take $z_1\in Y_0\setminus X_0$ such that $y_1 = \{z_1,p\}$ by \cref{subseteq}.
            Take $z_2\in Y_0\setminus X_0$ such that $y_2 = \{z_2,p\}$ by \cref{subseteq}.
            We have $(y_1,z_1)\in H$ and $(y_2,z_2)\in H$
                by \cref{setminus_elim_left,times_tuple_intro,union_iff}.
            Hence $H(y_1) = z_1$ and $H(y_2) = z_2$ by \cref{function_apply_iff,funs_is_function}.
            Hence $z_1 = z_2$.
            Thus $y_1 = y_2$.
        \end{byCase}
    \end{subproof}
    Hence $H$ is injective by \cref{funs_is_function,funs,injective_function}.

    Show that $\img{H}{Y}\subseteq Y_0$.
    \begin{subproof}
        Follows by \cref{img_subseteq_ran,funs_ran,subseteq_transitive,funs}.
    \end{subproof}
    Show that for all $z$ such that $z\in Y_0$ we have $z\in \img{H}{Y}$.
    \begin{subproof}
        Fix $z$ such that $z\in Y_0$.
        \begin{byCase}
        \caseOf{$z\in X_0$.}
            Take $x\in X$ such that $(x,z)\in h$ by \cref{img_iff}.
            Therefore $z\in \img{H}{Y}$ by \cref{img_iff,union_iff,union_intro_left}.
        \caseOf{$z\in Y_0\setminus X_0$.}
            $z\in \img{H}{Y}$ by \cref{subseteq,times_tuple_intro,union_iff,union_intro_right,img_iff}.
        \end{byCase}
    \end{subproof}
    Therefore $\img{H}{Y} = Y_0$ by \cref{subseteq,subseteq_antisymmetric}.

    Therefore $H$ is a bijection from $Y$ to $Y_0$ by \cref{funs,img_dom,funs_surj_iff,funs,bijections,injective}.
    Let $q = \converse{H}$.
    Hence $q$ is a function from $Y_0$ to $Y$ by \cref{bijective_converse_are_funs}.
    Let $S = \quotientTopology{q}{T_0}{Y_0}{Y}$.
    $S$ is a topology on $Y$ by \cref{quotient_topology_is_topology}.
    Hence $q\in\surj{Y_0}{Y}$ by \cref{bijective_converse_are_surjs}.
    Therefore $q$ is a quotient map from $Y_0$ to $Y$ with respect to $T_0$ and $S$
        by \cref{quotient_topology_quotient_map}.
    Hence $q$ is continuous from $Y_0$ to $Y$ with respect to $T_0$ and $S$ by \cref{quotient_map}.

    $H\circ q\in\funs{Y_0}{Y_0}$ by \cref{funs_circ}.
    Let $i_0 = \identity{Y_0}$.
    Hence $i_0\in\funs{Y_0}{Y_0}$ by \cref{id_elem_funs}.
    For all $z\in Y_0$ we have $(H\circ q)(z) = i_0(z)$
        by \cref{funs,function_apply_elim,funs_is_function,funs_type_apply,converse_iff,function_apply_iff,funs_ran,circ_apply,id_apply}.
    Therefore $H\circ q = \identity{Y_0}$ by \cref{functions_eq_iff}.

    $i_0$ is continuous from $Y_0$ to $Y_0$ with respect to $T_0$ and $T_0$
        by \cref{identity_continuous}.
    We have $H\circ q = i_0$.
    For all $V$ such that $V$ is open in $Y_0$ with respect to $T_0$ we have $\preimg{H}{V}\subseteq Y$
        by \cref{preimg_subseteq_dom,funs,subseteq}.
    Therefore for all $V$ such that $V$ is open in $Y_0$ with respect to $T_0$ we have
        $\preimg{q}{\preimg{H}{V}} = \preimg{i_0}{V}$
        by \cref{preimg_iff,circ_iff,subseteq,subseteq_antisymmetric}.
    Hence for all $V$ such that $V$ is open in $Y_0$ with respect to $T_0$ we have
        $\preimg{i_0}{V}$ is open in $Y_0$ with respect to $T_0$ by \cref{continuous}.
    Thus for all $V$ such that $V$ is open in $Y_0$ with respect to $T_0$ we have
        $\preimg{q}{\preimg{H}{V}}$ is open in $Y_0$ with respect to $T_0$.
    Hence for all $V$ such that $V$ is open in $Y_0$ with respect to $T_0$ we have
        $\preimg{H}{V}$ is open in $Y$ with respect to $S$ by \cref{quotient_map}.
    Therefore $H$ is continuous from $Y$ to $Y_0$ with respect to $S$ and $T_0$ by \cref{continuous}.
    We have $\converse{H} = q$ by \cref{converse_converse_eq}.
    Hence $\converse{H}$ is continuous from $Y_0$ to $Y$ with respect to $T_0$ and $S$ by \cref{quotient_map}.
    Therefore $H$ is a homeomorphism between $Y$ and $Y_0$ with respect to $S$ and $T_0$
        by \cref{homeomorphism,bijections}.

    Show that for all $y$ such that $y\in \img{q}{X_0}$ we have $y\in X$.
    \begin{subproof}
        Fix $y$ such that $y\in \img{q}{X_0}$.
        Take $z\in X_0$ such that $(z,y)\in q$ by \cref{img_iff}.
        Hence $(y,z)\in H$ by \cref{converse_iff}.
        \begin{byCase}
        \caseOf{$y\in X$.}
            Trivial.
        \caseOf{$y\in A$.}
            Suppose not.
            Take $u\in Y_0\setminus X_0$ such that $y = \{u,p\}$ by \cref{subseteq}.
            Therefore $(y,u)\in H$ by \cref{setminus_elim_left,times_tuple_intro,union_iff}.
            Hence $H(y) = u$ and $H(y) = z$ by \cref{function_apply_iff,funs_is_function}.
            Hence $z\in Y_0\setminus X_0$.
            Contradiction by \cref{setminus_elim_right}.
        \end{byCase}
    \end{subproof}
    Show that for all $x\in X$ we have $x\in \img{q}{X_0}$.
    \begin{subproof}
        Follows by \cref{topological_embedding,funs,function_apply_elim,funs_is_function,img_iff,subseteq,union_iff,converse_iff}.
    \end{subproof}
    Therefore $\img{q}{X_0} = X$ by \cref{subseteq,subseteq_antisymmetric}.

    Hence $Y = \img{q}{Y_0}$ by \cref{surj,funs,quotient_map,img_dom,ran_of_surj}.
    Therefore $Y$ is compact with respect to $\subspaceTopology{S}{Y}$ by \cref{continuous_image_compact}.
    Show that for all $V$ such that $V\in\subspaceTopology{S}{Y}$ we have $V\in S$.
    \begin{subproof}
        Follows by \cref{subspace_topology,topology_on,pow_iff,subseteq,inter_absorb_supseteq_left}.
    \end{subproof}
    Show that for all $V$ such that $V\in S$ we have $V\in\subspaceTopology{S}{Y}$.
    \begin{subproof}
        Follows by \cref{topology_on,pow_iff,subseteq,inter_absorb_supseteq_left,subspace_topology}.
    \end{subproof}
    Therefore $\subspaceTopology{S}{Y} = S$ by \cref{subseteq,subseteq_antisymmetric}.
    Hence $Y$ is compact with respect to $S$.

    Show that for all $y_1,y_2$ such that $y_1\in Y$ and $y_2\in Y$ and $y_1\neq y_2$
        there exist $N_1,N_2$ such that
        $N_1$ is a neighborhood of $y_1$ in $Y$ with respect to $S$ and
        $N_2$ is a neighborhood of $y_2$ in $Y$ with respect to $S$ and
        $N_1$ is disjoint from $N_2$.
    \begin{subproof}
        Fix $y_1$.
        Fix $y_2$.
        Assume $y_1\in Y$ and $y_2\in Y$ and $y_1\neq y_2$.
        Hence $H(y_1)\neq H(y_2)$ by \cref{funs,subseteq,injective_function,funs_is_function}.
        Therefore $H(y_1)\in Y_0$ and $H(y_2)\in Y_0$ by \cref{funs_type_apply}.
        Take $V_1,V_2$ such that
            $V_1$ is a neighborhood of $H(y_1)$ in $Y_0$ with respect to $T_0$ and
            $V_2$ is a neighborhood of $H(y_2)$ in $Y_0$ with respect to $T_0$ and
            $V_1$ is disjoint from $V_2$
            by \cref{hausdorff}.
        Let $W_1 = \preimg{H}{V_1}$.
        Let $W_2 = \preimg{H}{V_2}$.
        Hence $W_1$ is open in $Y$ with respect to $S$ and $W_2$ is open in $Y$ with respect to $S$
            by \cref{continuous,neighborhood}.
        Therefore $H(y_1)\in V_1$ and $H(y_2)\in V_2$ by \cref{neighborhood}.
        Hence $y_1\in\dom{H}$ and $y_2\in\dom{H}$ by \cref{funs,subseteq}.
        Therefore $(y_1,H(y_1))\in H$ and $(y_2,H(y_2))\in H$ by \cref{function_apply_elim,funs_is_function}.
        Hence $y_1\in W_1$ and $y_2\in W_2$ by \cref{preimg_iff}.
        $W_1$ is disjoint from $W_2$ by \cref{preimg_iff,funs_is_function,function_apply_iff,disjoint}.
        Therefore there exist $N_1,N_2$ such that
            $N_1$ is a neighborhood of $y_1$ in $Y$ with respect to $S$ and
            $N_2$ is a neighborhood of $y_2$ in $Y$ with respect to $S$ and
            $N_1$ is disjoint from $N_2$.
    \end{subproof}
    Hence $Y$ is Hausdorff with respect to $S$ by \cref{hausdorff}.

    Show that $X$ is a subspace of $Y$ with respect to $S$ and $T = \subspaceTopology{S}{X}$.
    \begin{subproof}
        Let $T_X = \subspaceTopology{S}{X}$.
        We have $T_X$ is a topology on $X$ and $X$ is a subspace of $Y$ with respect to $S$
            by \cref{union_upper_left,subspace_topology_is_topology,subspace_of}.
        Let $r = \restrl{H}{X}$.
        We have $r$ is continuous from $X$ to $Y_0$ with respect to $T_X$ and $T_0$
            by \cref{continuous_restrict_domain}.
        Show that for all $z$ such that $z\in\img{r}{X}$ we have $z\in X_0$.
        \begin{subproof}
            Follows by \cref{img_iff,restrl_iff,function_apply_iff,funs_is_function,subseteq,topological_embedding,funs,function_apply_elim}.
        \end{subproof}
        Show that for all $z$ such that $z\in X_0$ we have $z\in \img{r}{X}$.
        \begin{subproof}
            Follows by \cref{img_iff,function_apply_iff,topological_embedding,funs_is_function,subseteq,funs,restrl_iff}.
        \end{subproof}
        Therefore $\img{r}{X} = X_0$ by \cref{subseteq,subseteq_antisymmetric}.
        Hence $r$ is continuous from $X$ to $X_0$ with respect to $T_X$ and $\subspaceTopology{T_0}{X_0}$
            by \cref{continuous_restrict_range,subspace_of,subseteq_refl}.
        Hence $\subspaceTopology{T_0}{X_0} = \subspaceTopology{U}{X_0}$ by \cref{subspace_subspace_topology,subspace_of}.
        Therefore $h$ is continuous from $X$ to $X_0$ with respect to $T$ and $\subspaceTopology{U}{X_0}$
            by \cref{homeomorphism}.
        $h\in\funs{X}{X_0}$ by \cref{continuous,funs}.
        $r\in\funs{X}{X_0}$ by \cref{continuous,funs}.
        For all $x\in X$ we have $h(x) = r(x)$
            by \cref{subseteq,topological_embedding,funs,function_apply_iff,funs_is_function,restrl_iff,continuous}.
        Therefore $h = r$ by \cref{functions_eq_iff}.
        Hence $h$ is continuous from $X$ to $X_0$ with respect to $T_X$ and $\subspaceTopology{U}{X_0}$
            by \cref{continuous}.
        Let $q_0 = \restrl{q}{X_0}$.
        Hence $q_0$ is continuous from $X_0$ to $Y$ with respect to $\subspaceTopology{U}{X_0}$ and $S$
            by \cref{continuous_restrict_domain,subspace_of,subspace_subspace_topology}.
        Therefore $\img{q_0}{X_0} = X$ by \cref{funs,subseteq,restrl_img,inter_idempotent}.
        Hence $q_0$ is continuous from $X_0$ to $X$ with respect to $\subspaceTopology{U}{X_0}$ and $T_X$
            by \cref{continuous_restrict_range,subseteq_refl,subspace_of}.
        Hence $h$ is continuous from $X$ to $X_0$ with respect to $T$ and $\subspaceTopology{U}{X_0}$
            and $\converse{h}$ is continuous from $X_0$ to $X$ with respect to $\subspaceTopology{U}{X_0}$ and $T$
            by \cref{homeomorphism}.
        Let $j = \identity{X}$.
        $j\in\funs{X}{X}$ by \cref{id_elem_funs}.
        Hence $q_0\circ h$ is continuous from $X$ to $X$ with respect to $T$ and $T_X$
            by \cref{continuous_composite}.
        Therefore $q_0\circ h\in\funs{X}{X}$ by \cref{continuous,funs}.
        For all $x\in X$ we have $j(x) = (q_0\circ h)(x)$
            by \cref{continuous,funs,function_apply_elim,funs_is_function,union_intro_left,converse_iff,funs_type_apply,subseteq,restrl_of_function_apply,function_apply_iff,funs_ran,subseteq_transitive,circ_apply,id_apply}.
        Therefore $j = q_0\circ h$ by \cref{functions_eq_iff}.
        Hence $\converse{h}\circ h$ is continuous from $X$ to $X$ with respect to $T$ and $T$
            by \cref{continuous_composite}.
        Therefore $\converse{h}\circ h\in\funs{X}{X}$ by \cref{continuous,funs}.
        For all $x\in X$ we have $j(x) = (\converse{h}\circ h)(x)$
            by \cref{continuous,funs,function_apply_elim,funs_is_function,converse_iff,funs_type_apply,function_apply_iff,funs_ran,subseteq_transitive,subseteq,circ_apply,id_apply}.
        Therefore $j = \converse{h}\circ h$ by \cref{functions_eq_iff}.
        Hence $j$ is continuous from $X$ to $X$ with respect to $T$ and $T_X$
            and $j$ is continuous from $X$ to $X$ with respect to $T_X$ and $T$
            by \cref{continuous_composite}.
        For all $V$ such that $V\in T_X$ we have $V\subseteq X$
            by \cref{subspace_topology_is_topology,topology_on,subseteq,pow_iff}.
        Hence for all $V$ such that $V\in T_X$ we have $\preimg{j}{V} = V$
            by \cref{preimg_iff,id_iff,subseteq,subseteq_antisymmetric}.
        Therefore for all $V$ such that $V\in T_X$ we have $\preimg{j}{V}$ is open in $X$ with respect to $T$
            by \cref{continuous,open_in}.
        Therefore for all $V$ such that $V\in T_X$ we have $V\in T$ by \cref{open_in}.
        Hence $T_X\subseteq T$ by \cref{subseteq}.
        For all $V$ such that $V\in T$ we have $V\subseteq X$
            by \cref{topological_embedding,topology_on,subseteq,pow_iff}.
        Hence for all $V$ such that $V\in T$ we have $\preimg{j}{V} = V$
            by \cref{preimg_iff,id_iff,subseteq,subseteq_antisymmetric}.
        Therefore for all $V$ such that $V\in T$ we have $\preimg{j}{V}$ is open in $X$ with respect to $T_X$
            by \cref{continuous,open_in}.
        Therefore for all $V$ such that $V\in T$ we have $V\in T_X$ by \cref{open_in}.
        Hence $T\subseteq T_X$ by \cref{subseteq}.
        Therefore $T = T_X$ by \cref{subseteq_antisymmetric}.
    \end{subproof}

    Hence $\closure{X_0}{Y_0}{T_0} = Y_0$ by \cref{closure_in_subspace,subspace_of,inter_idempotent}.
    Therefore $\img{q}{Y_0}\subseteq \closure{\img{q}{X_0}}{Y}{S}$ by \cref{continuous_closure_image_subset}.
    Thus $Y\subseteq \closure{X}{Y}{S}$ by \cref{subseteq}.
    For all $y$ such that $y\in\closure{X}{Y}{S}$ we have $y\in Y$ by \cref{closure}.
    Hence $\closure{X}{Y}{S}\subseteq Y$ by \cref{subseteq}.
    Therefore $\closure{X}{Y}{S} = Y$ by \cref{subseteq_antisymmetric}.

    Hence $Y$ compactifies $X$ with respect to $S$ and $T$ by \cref{compactifies}.
    Let $j = \identity{Y_0}$.
    We have $j\circ H$ is continuous from $Y$ to $Z$ with respect to $S$ and $U$
        by \cref{continuous_expand_range,subspace_of}.
    Therefore $j\circ H\in\funs{Y}{Z}$ by \cref{continuous,funs}.
    $H\in\funs{Y}{Y_0}$ by \cref{homeomorphism,funs}.
    $H\in\funs{Y}{Z}$ by \cref{funs_weaken_codom,subseteq}.
    For all $y\in Y$ we have $(j\circ H)(y) = H(y)$
        by \cref{funs_type_apply,homeomorphism,funs,funs_ran,subseteq_transitive,id_elem_funs,subseteq,circ_apply,funs_is_function,id_apply}.
    Therefore $j\circ H = H$ by \cref{functions_eq_iff}.
    Hence $H$ is continuous from $Y$ to $Z$ with respect to $S$ and $U$ by \cref{continuous}.
    Therefore $H$ is a homeomorphism between $Y$ and $Y_0$ with respect to $S$ and $\subspaceTopology{U}{Y_0}$
        by \cref{homeomorphism,subspace_of,closure}.
    $\img{H}{Y} = Y_0$.
    Hence $H$ is a homeomorphism between $Y$ and $\img{H}{Y}$ with respect to $S$ and $\subspaceTopology{U}{\img{H}{Y}}$
        by \cref{homeomorphism}.
    Therefore $H$ is a topological embedding from $Y$ to $Z$ with respect to $S$ and $U$
        by \cref{topological_embedding}.
\end{proof}

% from §38 Item 2: uniqueness of the compactification induced by an embedding
\begin{lemma}\label{induced_compactification_uniqueness}
    Assume $h$ is a topological embedding from $X$ to $Z$ with respect to $T$ and $U$.
    Assume $Z$ is Hausdorff with respect to $U$.
    Assume $Y_1$ compactifies $X$ with respect to $S_1$ and $T$.
    Assume $Y_2$ compactifies $X$ with respect to $S_2$ and $T$.
    Assume $H_1$ is a topological embedding from $Y_1$ to $Z$ with respect to $S_1$ and $U$.
    Assume $H_2$ is a topological embedding from $Y_2$ to $Z$ with respect to $S_2$ and $U$.
    Suppose for all $x\in X$ we have $H_1(x) = h(x)$.
    Suppose for all $x\in X$ we have $H_2(x) = h(x)$.
    Then $Y_1$ is equivalent to $Y_2$ over $X$ with respect to $S_1$ and $S_2$ and $T$.
\end{lemma}
\begin{proof}
    Let $X_0 = \img{h}{X}$.

    Show that for all $z$ such that $z\in\img{H_1}{X}$ we have $z\in X_0$.
    \begin{subproof}
        Follows by \cref{img_iff,function_apply_intro,topological_embedding,funs_is_function,pair_eq_iff,function_apply_elim,funs,subseteq}.
    \end{subproof}
    Show that for all $z$ such that $z\in X_0$ we have $z\in\img{H_1}{X}$.
    \begin{subproof}
        Follows by \cref{img_iff,function_apply_intro,topological_embedding,funs_is_function,pair_eq_iff,compactifies,subspace_of,subseteq,function_apply_elim,funs}.
    \end{subproof}
    Therefore $\img{H_1}{X} = X_0$ by \cref{subseteq,subseteq_antisymmetric}.

    Show that for all $z$ such that $z\in\img{H_2}{X}$ we have $z\in X_0$.
    \begin{subproof}
        Follows by \cref{img_iff,function_apply_intro,topological_embedding,funs_is_function,pair_eq_iff,function_apply_elim,funs,subseteq}.
    \end{subproof}
    Show that for all $z$ such that $z\in X_0$ we have $z\in\img{H_2}{X}$.
    \begin{subproof}
        Follows by \cref{img_iff,function_apply_intro,topological_embedding,funs_is_function,pair_eq_iff,compactifies,subspace_of,subseteq,function_apply_elim,funs}.
    \end{subproof}
    Therefore $\img{H_2}{X} = X_0$ by \cref{subseteq,subseteq_antisymmetric}.

    We have $H_1$ is continuous from $Y_1$ to $Z$ with respect to $S_1$ and $U$ by \cref{topological_embedding}.
    Hence $X\subseteq Y_1$ and $\closure{X}{Y_1}{S_1} = Y_1$ by \cref{compactifies,subspace_of}.
    $\img{H_1}{Y_1}\subseteq \closure{\img{H_1}{X}}{Z}{U}$ by \cref{continuous_closure_image_subset}.
    Therefore $\img{H_1}{Y_1}\subseteq \closure{X_0}{Z}{U}$ by \cref{pair_eq_iff}.
    We have $U$ is a topology on $Z$ by \cref{continuous}.
    Hence $Y_1$ is compact with respect to $S_1$ by \cref{compactifies}.
    Therefore $X_0\subseteq Z$ by \cref{img_subseteq_ran,funs_ran,topological_embedding,subseteq_transitive}.
    Hence $\img{H_1}{Y_1}$ is compact with respect to $\subspaceTopology{U}{\img{H_1}{Y_1}}$
        by \cref{continuous_image_compact}.
    Therefore $\img{H_1}{Y_1}$ is closed in $Z$ with respect to $U$
        by \cref{img_subseteq_ran,funs_ran,topological_embedding,subseteq_transitive,compact_subspace_closed}.
    For all $z$ such that $z\in X_0$ we have $z\in\img{H_1}{Y_1}$
        by \cref{pair_eq_iff,img_iff,compactifies,subspace_of,subseteq}.
    Therefore $\closure{X_0}{Z}{U}\subseteq \img{H_1}{Y_1}$ by \cref{subseteq,closure_subseteq_closed}.
    Hence $\img{H_1}{Y_1} = \closure{X_0}{Z}{U}$ by \cref{subseteq,subseteq_antisymmetric}.

    We have $H_2$ is continuous from $Y_2$ to $Z$ with respect to $S_2$ and $U$ by \cref{topological_embedding}.
    Hence $X\subseteq Y_2$ and $\closure{X}{Y_2}{S_2} = Y_2$ by \cref{compactifies,subspace_of}.
    $\img{H_2}{Y_2}\subseteq \closure{\img{H_2}{X}}{Z}{U}$ by \cref{continuous_closure_image_subset}.
    Therefore $\img{H_2}{Y_2}\subseteq \closure{X_0}{Z}{U}$ by \cref{pair_eq_iff}.
    We have $U$ is a topology on $Z$ by \cref{continuous}.
    Hence $Y_2$ is compact with respect to $S_2$ by \cref{compactifies}.
    Therefore $X_0\subseteq Z$ by \cref{img_subseteq_ran,funs_ran,topological_embedding,subseteq_transitive}.
    Hence $\img{H_2}{Y_2}$ is compact with respect to $\subspaceTopology{U}{\img{H_2}{Y_2}}$
        by \cref{continuous_image_compact}.
    Therefore $\img{H_2}{Y_2}$ is closed in $Z$ with respect to $U$
        by \cref{img_subseteq_ran,funs_ran,topological_embedding,subseteq_transitive,compact_subspace_closed}.
    For all $z$ such that $z\in X_0$ we have $z\in\img{H_2}{Y_2}$
        by \cref{pair_eq_iff,img_iff,compactifies,subspace_of,subseteq}.
    Therefore $\closure{X_0}{Z}{U}\subseteq \img{H_2}{Y_2}$ by \cref{subseteq,closure_subseteq_closed}.
    Hence $\img{H_2}{Y_2} = \closure{X_0}{Z}{U}$ by \cref{subseteq,subseteq_antisymmetric}.

    Let $K = \converse{H_2}\circ H_1$.
    $H_2$ is a homeomorphism between $Y_2$ and $\img{H_2}{Y_2}$ with respect to $S_2$ and $\subspaceTopology{U}{\img{H_2}{Y_2}}$
        by \cref{topological_embedding}.
    Hence $H_2$ is a bijection from $Y_2$ to $\img{H_2}{Y_2}$ by \cref{homeomorphism,bijections}.
    Therefore $\converse{H_2}$ is a bijection from $\img{H_2}{Y_2}$ to $Y_2$
        by \cref{bijection_converse_is_bijection}.
    Hence $\converse{H_2}\in\funs{\img{H_2}{Y_2}}{Y_2}$ and $\dom{\converse{H_2}} = \img{H_2}{Y_2}$
        by \cref{bijective_converse_are_funs,funs}.
    $H_1$ is a homeomorphism between $Y_1$ and $\img{H_1}{Y_1}$ with respect to $S_1$ and $\subspaceTopology{U}{\img{H_1}{Y_1}}$
        by \cref{topological_embedding}.
    Hence $H_1$ is a bijection from $Y_1$ to $\img{H_1}{Y_1}$ by \cref{homeomorphism,bijections}.
    Therefore $H_1$ is a bijection from $Y_1$ to $\img{H_2}{Y_2}$ by \cref{pair_eq_iff}.
    We have $H_1\in\funs{Y_1}{\img{H_2}{Y_2}}$ by \cref{bijections_to_funs}.
    Hence $\ran{H_1}\subseteq \dom{\converse{H_2}}$ by \cref{funs_ran,pair_eq_iff}.
    $H_1$ is a homeomorphism between $Y_1$ and $\img{H_1}{Y_1}$ with respect to $S_1$ and $\subspaceTopology{U}{\img{H_1}{Y_1}}$
        by \cref{topological_embedding}.
    Hence $H_1$ is continuous from $Y_1$ to $\img{H_1}{Y_1}$ with respect to $S_1$ and $\subspaceTopology{U}{\img{H_1}{Y_1}}$
        by \cref{homeomorphism}.
    Hence $H_1$ is continuous from $Y_1$ to $\img{H_2}{Y_2}$ with respect to $S_1$ and $\subspaceTopology{U}{\img{H_2}{Y_2}}$
        by \cref{continuous_eq_codomain,pair_eq_iff}.
    $H_2$ is a homeomorphism between $Y_2$ and $\img{H_2}{Y_2}$ with respect to $S_2$ and $\subspaceTopology{U}{\img{H_2}{Y_2}}$
        by \cref{topological_embedding}.
    Hence $\converse{H_2}$ is continuous from $\img{H_2}{Y_2}$ to $Y_2$ with respect to $\subspaceTopology{U}{\img{H_2}{Y_2}}$ and $S_2$
        by \cref{homeomorphism}.
    Therefore $K$ is continuous from $Y_1$ to $Y_2$ with respect to $S_1$ and $S_2$ by \cref{continuous_composite}.
    Hence $K$ is a bijection from $Y_1$ to $Y_2$ by \cref{bijection_circ}.
    Therefore $K\in\funs{Y_1}{Y_2}$ by \cref{bijections_to_funs}.
    $Y_1$ is compact with respect to $S_1$ and $Y_2$ is Hausdorff with respect to $S_2$ by \cref{compactifies}.
    Hence $K$ is a homeomorphism between $Y_1$ and $Y_2$ with respect to $S_1$ and $S_2$
        by \cref{compact_hausdorff_bijection_homeomorphism}.

    For all $x\in X$ we have $K(x) = x$
        by \cref{compactifies,subspace_of,subseteq,topological_embedding,funs,function_apply_elim,pair_eq_iff,converse_iff,img_iff,bijections_dom,bijections_to_funs,circ_apply,function_apply_intro,funs_is_function}.

    Therefore $Y_1$ is equivalent to $Y_2$ over $X$ with respect to $S_1$ and $S_2$ and $T$
        by \cref{compactification_equivalence_over_x}.
\end{proof}

% from §38 helper definition 2: bounded real extension property
\begin{definition}\label{bounded_real_extension_property}
    $Y$ has the bounded real extension property over $X$ with respect to $S$ and $T$ iff
        $Y$ compactifies $X$ with respect to $S$ and $T$ and
        for all $f$ if
            $f$ is continuous from $X$ to $\reals$ with respect to $T$ and $\standardTopologyR$ and
            $\img{f}{X}$ is bounded in $\reals$ with respect to $\standardMetricR$
        then there exists $g$ such that
            $g$ is continuous from $Y$ to $\reals$ with respect to $S$ and $\standardTopologyR$ and
            for all $x\in X$ we have $g(x) = f(x)$ and
            for all $h$ if
                $h$ is continuous from $Y$ to $\reals$ with respect to $S$ and $\standardTopologyR$ and
                for all $x\in X$ we have $h(x) = f(x)$
            then $h = g$.
\end{definition}

% from §38 helper lemma 3: bounded subset of reals lies in a closed interval
\begin{lemma}\label{bounded_reals_subset_closed_interval}
    Assume $A$ is bounded in $\reals$ with respect to $\standardMetricR$.
    Then there exist $a,b$ such that
        $a\in\reals$ and
        $b\in\reals$ and
        $a \leq b$ and
        $A\subseteq \intervalclosed{a}{b}$.
\end{lemma}
\begin{proof}
    We have $\standardMetricR$ is a metric on $\reals$ and $A\subseteq\reals$ by \cref{standard_metric_is_metric,metric_bounded}.
    \begin{byCase}
    \caseOf{$A = \emptyset$.}
        Let $a = 0$.
        Let $b = 0$.
        $a\in\reals$ and $b\in\reals$ by \cref{real_add_ident}.
        $a \leq b$.
        $A\subseteq \intervalclosed{a}{b}$ by \cref{emptyset_subseteq}.
        Thus there exist $u,v$ such that
            $u\in\reals$ and
            $v\in\reals$ and
            $u \leq v$ and
            $A\subseteq \intervalclosed{u}{v}$.
    \caseOf{$A \neq \emptyset$.}
        Take $x_0$ such that $x_0\in A$ by \cref{nonempty_inhabited}.
        Take $M$ such that
            $M\in\reals$ and
            for all $u,v$ such that $u\in A$ and $v\in A$ we have $\apply{\standardMetricR}{(u,v)} \leq M$
            by \cref{metric_bounded}.
        Let $a = x_0 - M$.
        Let $b = x_0 + M$.
        $x_0\in\reals$ by \cref{subseteq}.
        We have $\neg{M}\in\reals$ by \cref{real_add_inverse}.
        Hence $a\in\reals$ by \cref{real_sub,real_add_closed}.
        $b\in\reals$ by \cref{real_add_closed}.
        We have $\apply{\standardMetricR}{(x_0,x_0)} \leq M$.
        $\apply{\standardMetricR}{(x_0,x_0)}\in\reals$ by
            \cref{standard_metric_value,real_sub,real_add_inverse,real_add_closed,real_abs_in_reals}.
        $\apply{\standardMetricR}{(x_0,x_0)} \geq 0$ by \cref{metric_nonnegative,metric_on,subseteq}.
        Hence $0 \leq M$ by \cref{real_add_ident,real_leq_trans}.
        For all $y$ such that $y\in A$ we have $y\in \intervalclosed{a}{b}$
            by \cref{subseteq,standard_metric_value,real_abs_sub_leq_imp_left_bound,real_abs_sub_swap,real_abs_sub_leq_imp_right_bound,intervalclosed}.
        Hence $A\subseteq \intervalclosed{a}{b}$ by \cref{subseteq}.
        We have $a \leq b$ by \cref{subseteq,intervalclosed,real_leq_trans}.
        Thus there exist $u,v$ such that
            $u\in\reals$ and
            $v\in\reals$ and
            $u \leq v$ and
            $A\subseteq \intervalclosed{u}{v}$.
    \end{byCase}
\end{proof}

% from §38 helper lemma 5: uniqueness of continuous real-valued extensions over a compactification
\begin{lemma}\label{compactification_real_extension_unique}
    Assume $Y$ compactifies $X$ with respect to $S$ and $T$.
    Assume $g$ is continuous from $Y$ to $\reals$ with respect to $S$ and $\standardTopologyR$.
    Assume $h$ is continuous from $Y$ to $\reals$ with respect to $S$ and $\standardTopologyR$.
    Suppose for all $x\in X$ we have $g(x) = h(x)$.
    Then $g = h$.
\end{lemma}
\begin{proof}
    We have $\reals$ is Hausdorff with respect to $\standardTopologyR$
        by \cref{standard_metric_is_metric,metric_space_hausdorff,standard_metric_topology}.
    $g\in\funs{Y}{\reals}$ and $h\in\funs{Y}{\reals}$ by \cref{continuous,funs}.

    Show that for all $y\in Y$ we have $g(y) = h(y)$.
    \begin{subproof}
        Fix $y\in Y$.
        Suppose not.
        We have $g(y)\in\reals$ and $h(y)\in\reals$ by \cref{funs_type_apply,continuous}.
        Take $N_1,N_2$ such that
            $N_1$ is a neighborhood of $g(y)$ in $\reals$ with respect to $\standardTopologyR$ and
            $N_2$ is a neighborhood of $h(y)$ in $\reals$ with respect to $\standardTopologyR$ and
            $N_1$ is disjoint from $N_2$
            by \cref{hausdorff}.
        Let $W_1 = \preimg{g}{N_1}$.
        Let $W_2 = \preimg{h}{N_2}$.
        We have $W_1$ is open in $Y$ with respect to $S$ and $W_2$ is open in $Y$ with respect to $S$
            by \cref{continuous,neighborhood}.
        Therefore $(y,g(y))\in g$ and $(y,h(y))\in h$
            by \cref{function_apply_elim,continuous,funs,funs_is_function}.
        Hence $y\in W_1$ and $y\in W_2$ by \cref{neighborhood,preimg_iff}.
        Let $C = \{W_1,W_2\}$.
        We have $S$ is a topology on $Y$ and $C\subseteq S$ by \cref{compactifies,upair_iff,open_in,subseteq}.
        Let $W = \inters{C}$.
        Hence $W$ is open in $Y$ with respect to $S$
            by \cref{pair_is_finite,upair_intro_left,inhabited_nonempty,finite_intersections_open_in_topology,open_in}.
        Therefore $y\in W$ by \cref{inters_intro,upair_iff}.
        Hence $W$ intersects $X$ by \cref{compactifies,subspace_of,pair_eq_iff,closure_iff_intersects_open}.
        Take $x$ such that $x\in W\inter X$ by \cref{intersects,nonempty_inhabited}.
        We have $x\in W$ and $x\in X$ by \cref{inter}.
        Hence $x\in W_1$ and $x\in W_2$ by \cref{upair_iff,inters_subseteq_elem,subseteq}.
        Take $z_1\in N_1$ such that $(x,z_1)\in g$ by \cref{subseteq,preimg_iff,preimg_subseteq_dom}.
        Take $z_2\in N_2$ such that $(x,z_2)\in h$ by \cref{subseteq,preimg_iff,preimg_subseteq_dom}.
        Therefore $g(x) = z_1$ and $h(x) = z_2$ by \cref{subseteq,function_apply_iff,funs_is_function}.
        Thus $h(x)\in N_1$ and $h(x)\in N_2$ by \cref{pair_eq_iff}.
        Contradiction by \cref{disjoint}.
    \end{subproof}
    Therefore $g = h$ by \cref{functions_eq_iff}.
\end{proof}

% from §38 helper lemma 6: indexed product of real subspaces carries the subspace product topology
\begin{lemma}\label{indexed_real_product_subspace_topology}
    Assume $C$ is a function.
    Assume $S$ is a function.
    Assume $R$ is a function.
    Assume $T$ is a function.
    Assume $\dom{S} = \dom{C}$.
    Assume $\dom{R} = \dom{C}$.
    Assume $\dom{T} = \dom{C}$.
    Assume $\dom{C}$ is inhabited.
    Assume that for all $\alpha\in\dom{C}$ we have
        $\apply{R}{\alpha} = \reals$ and
        $\apply{T}{\alpha} = \standardTopologyR$ and
        $\apply{C}{\alpha}\subseteq\reals$ and
        $\apply{S}{\alpha} = \subspaceTopology{\standardTopologyR}{\apply{C}{\alpha}}$.
    Then $\productTopologyIndexed{S}{C}
        = \subspaceTopology{\productTopologyIndexed{T}{R}}{\productFamily{C}}$.
\end{lemma}
\begin{proof}
    Let $P = \productTopologyIndexed{S}{C}$.
    Let $Q = \productTopologyIndexed{T}{R}$.
    Let $A = \productFamily{C}$.
    Let $T_A = \subspaceTopology{Q}{A}$.

    For all $\alpha\in\dom{C}$ we have $\apply{C}{\alpha}\subseteq\apply{R}{\alpha}$ by \cref{pair_eq_iff}.
    $\standardTopologyR$ is a topology on $\reals$
        by \cref{standard_basis_is_basis,basis_topology_is_topology,standard_topology_reals}.
    For all $\alpha\in\dom{R}$ we have $\apply{T}{\alpha}$ is a topology on $\apply{R}{\alpha}$
        by \cref{pair_eq_iff}.
    For all $\alpha\in\dom{C}$ we have $\apply{S}{\alpha}$ is a topology on $\apply{C}{\alpha}$
        by \cref{subspace_of,subspace_topology_is_topology,pair_eq_iff}.
    Hence $A$ is a subspace of $\productFamily{R}$ with respect to $Q$ by \cref{product_subspace_product}.
    Therefore $A\subseteq\productFamily{R}$ by \cref{subspace_of}.
    $P$ is a topology on $A$ by
        \cref{product_subbasis_is_subbasis,product_subbasis_finite_intersections_basis,basis_for_topology,basis_topology_is_topology}.
    $Q$ is a topology on $\productFamily{R}$ by
        \cref{product_subbasis_is_subbasis,product_subbasis_finite_intersections_basis,basis_for_topology,basis_topology_is_topology}.
    Hence $T_A$ is a topology on $A$ by \cref{subspace_topology_is_topology,subspace_of}.

    Let $j = \identity{A}$.
    We have $j\in\funs{A}{A}$ by \cref{id_elem_funs}.
    Hence $j\in\funs{A}{\productFamily{R}}$ by \cref{id_elem_funs,funs_weaken_codom,subspace_of}.

    Show that for all $\alpha\in\dom{R}$ we have $\piIndex{R}{\alpha}\circ j$ is continuous from $A$ to $\apply{R}{\alpha}$
        with respect to $P$ and $\apply{T}{\alpha}$.
    \begin{subproof}
        Fix $\alpha\in\dom{R}$.
        Let $p = \piIndex{C}{\alpha}$.
        Let $q = \identity{\apply{C}{\alpha}}\circ p$.
        We have $p$ is continuous from $A$ to $\apply{C}{\alpha}$ with respect to $P$ and $\apply{S}{\alpha}$
            by \cref{projection_map_indexed_continuous,pair_eq_iff}.
        We have $\apply{C}{\alpha}$ is a subspace of $\reals$ with respect to $\standardTopologyR$ by \cref{subspace_of}.
        Hence $q$ is continuous from $A$ to $\reals$ with respect to $P$ and $\standardTopologyR$
            by \cref{continuous_expand_range,pair_eq_iff}.
        $q\in\funs{A}{\reals}$ by \cref{continuous,funs}.
        $p\in\funs{A}{\apply{C}{\alpha}}$ by \cref{continuous,funs}.
        Hence $p\in\funs{A}{\reals}$ by \cref{funs_weaken_codom,subseteq}.
        For all $x\in A$ we have $q(x) = p(x)$
            by \cref{continuous,funs,subseteq,funs_type_apply,pair_eq_iff,function_apply_elim,funs_is_function,id_iff,circ_elem_intro,function_apply_intro}.
        Therefore $q = p$ by \cref{functions_eq_iff}.
        Hence $p$ is continuous from $A$ to $\reals$ with respect to $P$ and $\standardTopologyR$
            by \cref{continuous_eq_subst}.
        Hence $p$ is continuous from $A$ to $\apply{R}{\alpha}$ with respect to $P$ and $\apply{T}{\alpha}$
            by \cref{continuous_eq_codomain,pair_eq_iff}.
        $\piIndex{R}{\alpha}\circ j\in\funs{A}{\apply{R}{\alpha}}$ by \cref{projection_map_indexed_function,funs_circ,subseteq}.
        $p\in\funs{A}{\apply{R}{\alpha}}$ by \cref{continuous,funs}.
        For all $x\in A$ we have $(\piIndex{R}{\alpha}\circ j)(x) = p(x)$
            by \cref{subspace_of,subseteq,projection_map_indexed_function,funs,funs_ran,id_elem_funs,subseteq_transitive,circ_apply,funs_is_function,id_apply,pair_eq_iff,projection_map_indexed,function_apply_elim,function_apply_intro,continuous}.
        Therefore $\piIndex{R}{\alpha}\circ j = p$ by \cref{functions_eq_iff}.
        Therefore $\piIndex{R}{\alpha}\circ j$ is continuous from $A$ to $\apply{R}{\alpha}$
            with respect to $P$ and $\apply{T}{\alpha}$ by \cref{continuous_eq_subst}.
    \end{subproof}
    Therefore $j$ is continuous from $A$ to $\productFamily{R}$ with respect to $P$ and $Q$
        by \cref{continuous_map_into_indexed_product,subseteq,projection_map_indexed_function,funs}.
    Hence $j$ is continuous from $A$ to $A$ with respect to $P$ and $T_A$
        by \cref{img_iff,id_iff,pair_eq_iff,subseteq,continuous_restrict_range,subspace_of}.

    Show that for all $\alpha\in\dom{C}$ we have $\piIndex{C}{\alpha}\circ j$ is continuous from $A$ to $\apply{C}{\alpha}$
        with respect to $T_A$ and $\apply{S}{\alpha}$.
    \begin{subproof}
        Fix $\alpha\in\dom{C}$.
        Let $p = \piIndex{R}{\alpha}$.
        Let $g = \restrl{p}{A}$.
        We have $p$ is continuous from $\productFamily{R}$ to $\apply{R}{\alpha}$ with respect to $Q$ and $\apply{T}{\alpha}$
            by \cref{projection_map_indexed_continuous}.
        Hence $g$ is continuous from $A$ to $\apply{R}{\alpha}$ with respect to $T_A$ and $\apply{T}{\alpha}$
            by \cref{continuous_restrict_domain,subspace_of}.
        Hence $g$ is continuous from $A$ to $\reals$ with respect to $T_A$ and $\standardTopologyR$
            by \cref{continuous_eq_codomain,pair_eq_iff}.
        $\img{g}{A}\subseteq\apply{C}{\alpha}$
            by \cref{img_iff,restrl_iff,pair_eq_iff,projection_map_indexed_function,funs,subspace_of,subseteq,function_apply_intro,funs_is_function,projection_map_indexed,function_apply_elim,indexed_product}.
        Hence $g$ is continuous from $A$ to $\apply{C}{\alpha}$ with respect to $T_A$ and $\apply{S}{\alpha}$
            by \cref{continuous_subspace_range,continuous,funs,pair_eq_iff,subseteq}.
        $\piIndex{C}{\alpha}\circ j\in\funs{A}{\apply{C}{\alpha}}$ by \cref{projection_map_indexed_function,funs_circ,subseteq}.
        $g\in\funs{A}{\apply{C}{\alpha}}$ by \cref{continuous,funs}.
        For all $x\in A$ we have $(\piIndex{C}{\alpha}\circ j)(x) = g(x)$
            by \cref{projection_map_indexed_function,funs_circ,subseteq,continuous,funs,pair_eq_iff,funs_ran,id_elem_funs,circ_apply,funs_is_function,id_apply,projection_map_indexed,function_apply_elim,function_apply_intro,restrl_iff}.
        Therefore $\piIndex{C}{\alpha}\circ j = g$ by \cref{functions_eq_iff}.
        Therefore $\piIndex{C}{\alpha}\circ j$ is continuous from $A$ to $\apply{C}{\alpha}$
            with respect to $T_A$ and $\apply{S}{\alpha}$ by \cref{continuous_eq_subst}.
    \end{subproof}
    Hence $j$ is continuous from $A$ to $A$ with respect to $T_A$ and $P$
        by \cref{continuous_map_into_indexed_product,id_elem_funs,pair_eq_iff,continuous_eq_codomain}.

    For all $U$ such that $U\in T_A$ we have $U\subseteq A$
        by \cref{subspace_topology_is_topology,topology_on,subseteq,pow_iff}.
    Hence for all $U$ such that $U\in T_A$ we have $\preimg{j}{U} = U$
        by \cref{preimg_iff,id_iff,subseteq,subseteq_antisymmetric}.
    Therefore for all $U$ such that $U\in T_A$ we have $\preimg{j}{U}$ is open in $A$ with respect to $P$
        by \cref{continuous,open_in}.
    Therefore for all $U$ such that $U\in T_A$ we have $U\in P$ by \cref{open_in}.
    Show that for all $U$ such that $U\in P$ we have $U\in T_A$.
    \begin{subproof}
        Fix $U$ such that $U\in P$.
        We have $U\subseteq A$ by \cref{topology_on,subseteq,pow_iff}.
        Show that $\preimg{j}{U} = U$.
        \begin{subproof}
            Show that $\preimg{j}{U}\subseteq U$.
            \begin{subproof}
                It suffices to show that for all $x\in\preimg{j}{U}$ we have $x\in U$.
                Fix $x\in\preimg{j}{U}$.
                Take $y\in U$ such that $(x,y)\in j$ by \cref{preimg_iff}.
                Hence $x = y$ by \cref{id_iff}.
                Therefore $x\in U$.
            \end{subproof}
            Hence $\preimg{j}{U}\subseteq U$ by \cref{subseteq}.

            Show that $U\subseteq\preimg{j}{U}$.
            \begin{subproof}
                It suffices to show that for all $z\in U$ we have $z\in\preimg{j}{U}$.
                Fix $z\in U$.
                Hence $z\in A$ by \cref{subseteq}.
                $(z,z)\in j$ by \cref{id_iff}.
                Therefore $z\in\preimg{j}{U}$ by \cref{preimg_iff}.
            \end{subproof}
            Hence $U\subseteq\preimg{j}{U}$ by \cref{subseteq}.
            Therefore $\preimg{j}{U} = U$ by \cref{subseteq_antisymmetric}.
        \end{subproof}
        Hence $\preimg{j}{U}$ is open in $A$ with respect to $T_A$ by \cref{continuous,open_in}.
        Therefore $U\in T_A$ by \cref{open_in}.
    \end{subproof}
    Therefore $P = T_A$ by \cref{subseteq,subseteq_antisymmetric}.
\end{proof}

% from §38 helper lemma 7: restricting the codomain preserves topological embeddings
\begin{lemma}\label{topological_embedding_restrict_codomain}
    Assume $f$ is a topological embedding from $X$ to $Y$ with respect to $T$ and $T_1$.
    Assume $Z$ is a subspace of $Y$ with respect to $T_1$.
    Assume $S = \subspaceTopology{T_1}{Z}$.
    Suppose $\img{f}{X}\subseteq Z$.
    Then $f$ is a topological embedding from $X$ to $Z$ with respect to $T$ and $S$.
\end{lemma}
\begin{proof}
    $f$ is injective and
        $f$ is continuous from $X$ to $Y$ with respect to $T$ and $T_1$ and
        $f$ is a homeomorphism between $X$ and $\img{f}{X}$ with respect to
            $T$ and $\subspaceTopology{T_1}{\img{f}{X}}$ by \cref{topological_embedding}.
    Hence $f$ is continuous from $X$ to $Z$ with respect to $T$ and $S$ by \cref{continuous_restrict_range}.
    $f\in\funs{X}{Z}$ by \cref{continuous,funs}.
    Therefore $\img{f}{X}$ is a subspace of $Z$ with respect to $S$
        by \cref{subspace_topology_is_topology,subspace_of}.
    Therefore $\subspaceTopology{S}{\img{f}{X}} = \subspaceTopology{T_1}{\img{f}{X}}$
        by \cref{subspace_subspace_topology,continuous,subspace_of,pair_eq_iff}.
    Hence $f$ is a homeomorphism between $X$ and $\img{f}{X}$ with respect to
        $T$ and $\subspaceTopology{S}{\img{f}{X}}$ by \cref{pair_eq_iff}.
    Therefore $f$ is a topological embedding from $X$ to $Z$ with respect to $T$ and $S$
        by \cref{topological_embedding}.
\end{proof}

% from §38 helper lemma 8: bounded continuous real maps belong to the bounded family
\begin{lemma}\label{bounded_real_function_family_member}
    Let $J = \{g\in\pow{X\times\reals} \mid \text{$g$ is continuous from $X$ to $\reals$ with respect to $T$ and $\standardTopologyR$ and
        there exist $a,b\in\reals$ such that $a \leq b$ and $\img{g}{X}\subseteq\intervalclosed{a}{b}$}\}$.
    Suppose $g\in\pow{X\times\reals}$.
    Suppose $g$ is continuous from $X$ to $\reals$ with respect to $T$ and $\standardTopologyR$.
    Suppose there exist $a,b\in\reals$ such that $a \leq b$ and $\img{g}{X}\subseteq\intervalclosed{a}{b}$.
    Then $g\in J$.
\end{lemma}
\begin{proof}
    Straightforward.
\end{proof}

% from §38 Item 3: existence of a compactification with the bounded real extension property
\begin{theorem}\label{stone_cech_existence_bounded_real_property}
    Assume $X$ is completely regular with respect to $T$.
    Then there exist $Y, S$ such that
        $Y$ has the bounded real extension property over $X$ with respect to $S$ and $T$.
\end{theorem}
\begin{proof}
    $T$ is a topology on $X$ and $X$ has closed singletons with respect to $T$ by \cref{completely_regular_space}.
    Let $J = \{g\in\pow{X\times\reals} \mid \text{$g$ is continuous from $X$ to $\reals$ with respect to $T$ and $\standardTopologyR$ and
        there exist $a,b\in\reals$ such that $a \leq b$ and $\img{g}{X}\subseteq\intervalclosed{a}{b}$}\}$.
    Let $z = \{(x,0) \mid x\in X\}$.
    $z$ is a function from $X$ to $\reals$ by \cref{constant_map_function,real_add_ident}.
    Hence for all $x\in X$ we have $z(x) = 0$ by \cref{constant_map_function,funs_is_function,function_apply_intro}.
    $\standardTopologyR$ is a topology on $\reals$
        by \cref{standard_basis_is_basis,basis_topology_is_topology,standard_topology_reals}.
    $0\in\reals$ by \cref{real_add_ident}.
    Hence $z$ is continuous from $X$ to $\reals$ with respect to $T$ and $\standardTopologyR$
        by \cref{real_add_ident,constant_continuous}.
    We have $z\in\pow{X\times\reals}$
        by \cref{continuous,funs_is_function,funs,funs_subseteq_rels,subseteq,rels_elim,powerset_intro}.
    $0 \leq 0$.
    Hence $\img{z}{X}\subseteq \intervalclosed{0}{0}$
        by \cref{img_iff,function_apply_intro,funs_is_function,pair_eq_iff,intervalclosed,subseteq}.
    Therefore there exist $c,d\in\reals$ such that $c \leq d$ and $\img{z}{X}\subseteq\intervalclosed{c}{d}$.
    Hence $z\in J$ by \cref{bounded_real_function_family_member}.
    Therefore $J$ is inhabited.

    Let $f = \identity{J}$.
    $f$ is a function on $J$ by \cref{id_is_relation,id_rightunique,id_dom,relation,rightunique,subseteq,subseteq_antisymmetric}.
    For all $\alpha\in J$ we have $\apply{f}{\alpha}$ is continuous from $X$ to $\reals$ with respect to $T$ and $\standardTopologyR$
        by \cref{id_apply,id_elem_funs}.
    Show that for all $x,U$ such that $x\in X$ and $U$ is a neighborhood of $x$ in $X$ with respect to $T$
        there exists $\alpha\in J$ such that
            $0 < \apply{\apply{f}{\alpha}}{x}$ and
            for all $y\in X\setminus U$ we have $\apply{\apply{f}{\alpha}}{y} = 0$.
    \begin{subproof}
        Fix $x$.
        Fix $U$ such that $x\in X$ and $U$ is a neighborhood of $x$ in $X$ with respect to $T$.
        Let $I_0 = \intervalclosed{0}{1}$.
        Let $T_I = \subspaceTopology{\standardTopologyR}{I_0}$.
        Let $j_0 = \identity{I_0}$.
        $0\in\reals$ and $1\in\reals$ and $0 < 1$
            by \cref{real_add_ident,real_one_in_naturals,real_naturals_subset_reals,subseteq,real_zero_lt_one}.
        Hence $I_0\subseteq\reals$ and $I_0$ is a subspace of $\reals$ with respect to $\standardTopologyR$
            by \cref{real_lt_imp_leq,intervalclosed,subseteq,subspace_of}.
        Therefore $T_I$ is a topology on $I_0$ by \cref{subspace_topology_is_topology}.
        We have $U$ is open in $X$ with respect to $T$ and $x\in U$ by \cref{neighborhood}.
        Let $A = X\setminus U$.
        Hence $X\setminus A$ is open in $X$ with respect to $T$
            by \cref{open_in,topology_on,subseteq,pow_iff,double_relative_complement,pair_eq_iff}.
        Hence $A$ is closed in $X$ with respect to $T$ by \cref{closed_in,setminus_subseteq}.
        Therefore $x\notin A$ by \cref{setminus_elim_right}.
        Take $g$ such that
            $g$ is continuous from $X$ to $I_0$ with respect to $T$ and $T_I$ and
            $g(x) = 1$ and
            for all $y\in A$ we have $g(y) = 0$
            by \cref{completely_regular_space}.
        Let $g_1 = j_0\circ g$.
        $g_1$ is continuous from $X$ to $\reals$ with respect to $T$ and $\standardTopologyR$
            by \cref{continuous_expand_range}.
        $g_1\in\funs{X}{\reals}$ and $g_1$ is a function and $\dom{g_1} = X$
            by \cref{continuous,funs_is_function,funs}.
        Hence $g_1\in\pow{X\times\reals}$ by \cref{funs_subseteq_rels,subseteq,rels_elim,powerset_intro}.
        For all $y\in X$ we have $g_1(y) = g(y)$
            by \cref{continuous,funs,continuous,funs_type_apply,function_apply_elim,id_iff,circ_elem_intro,pair_eq_iff,function_apply_intro,funs_is_function}.
        Therefore $\img{g_1}{X}\subseteq \intervalclosed{0}{1}$
            by \cref{img_iff,function_apply_intro,subseteq,continuous,funs_type_apply,pair_eq_iff}.
        Hence there exist $a,b\in\reals$ such that $a \leq b$ and $\img{g_1}{X}\subseteq\intervalclosed{a}{b}$
            by \cref{pair_eq_iff}.
        Therefore $g_1\in J$ by \cref{bounded_real_function_family_member}.
        Take $\alpha\in J$ such that $\alpha = g_1$.
        We have $\apply{f}{\alpha} = g_1$ by \cref{id_apply,id_elem_funs,pair_eq_iff}.
        Therefore $\apply{\apply{f}{\alpha}}{x} = g_1(x)$ by \cref{pair_eq_iff}.
        Therefore $g_1(x) = g(x)$.
        We have $g(x) = 1$.
        Hence $0 < \apply{\apply{f}{\alpha}}{x}$ by \cref{pair_eq_iff,real_zero_lt_one}.
        For all $y\in X\setminus U$ we have $\apply{\apply{f}{\alpha}}{y} = 0$.
    \end{subproof}

    Let $R = \{(\alpha,\reals) \mid \alpha\in J\}$.
    Let $T_0 = \{(\alpha,\standardTopologyR) \mid \alpha\in J\}$.
    For all $\alpha,y_1,y_2$ such that $(\alpha,y_1)\in R$ and $(\alpha,y_2)\in R$ we have $y_1 = y_2$ by \cref{pair_eq_iff}.
    Hence $R$ is a function on $J$ by \cref{relation,rightunique,dom_intro,dom_iff,pair_eq_iff,subseteq,subseteq_antisymmetric}.
    For all $\alpha,y_1,y_2$ such that $(\alpha,y_1)\in T_0$ and $(\alpha,y_2)\in T_0$ we have $y_1 = y_2$ by \cref{pair_eq_iff}.
    Therefore $T_0$ is a function on $J$ by \cref{relation,rightunique,dom_intro,dom_iff,pair_eq_iff,subseteq,subseteq_antisymmetric}.
    We have $\dom{R} = J$ and $\dom{T_0} = J$ by \cref{funs}.
    For all $\alpha\in J$ we have $R(\alpha) = \reals$ and $T_0(\alpha) = \standardTopologyR$ by \cref{function_apply_intro}.

    Let $P_0 = \productTopologyIndexed{T_0}{R}$.
    Let $F = \{(x,y) \mid x\in X, y\in\productFamily{R} \mid
        \forall \alpha\in J. \apply{y}{\alpha} = \apply{\apply{f}{\alpha}}{x}\}$.
    Therefore $F$ is a topological embedding from $X$ to $\productFamily{R}$ with respect to $T$ and $P_0$
        by \cref{urysohn_embedding_theorem_real_power}.

    Let $M = \{w\in J\times\pow{\reals} \mid \text{there exists $\alpha\in J$ such that
        there exist $a,b\in\reals$ such that
        $w = (\alpha,\intervalclosed{a}{b})$ and $a \leq b$ and $\img{\alpha}{X}\subseteq\intervalclosed{a}{b}$}\}$.
    $M$ is a relation by \cref{times_elem_is_tuple,relation}.
    For all $\alpha\in J$ there exists $I$ such that $\alpha\mathrel{M} I$
        by \cref{subseteq,times_tuple_intro,intervalclosed,powerset_intro,pair_eq_iff}.
    Take $C$ such that $C$ is a function on $J$ and for all $\alpha\in J$ we have $\alpha\mathrel{M}\apply{C}{\alpha}$
        by \cref{choice_function}.
    Hence $C$ is a function and $\dom{C} = J$ by \cref{funs}.
    Let $S_0 = \{(\alpha,\subspaceTopology{\standardTopologyR}{C(\alpha)}) \mid \alpha\in J\}$.
    For all $\alpha,y_1,y_2$ such that $(\alpha,y_1)\in S_0$ and $(\alpha,y_2)\in S_0$ we have $y_1 = y_2$
        by \cref{pair_eq_iff}.
    Hence $S_0$ is a function on $J$ by \cref{relation,rightunique,dom_intro,dom_iff,pair_eq_iff,subseteq,subseteq_antisymmetric}.
    We have $\dom{S_0} = J$ and for all $\alpha\in J$ we have $S_0(\alpha) = \subspaceTopology{\standardTopologyR}{C(\alpha)}$
        by \cref{funs,function_apply_intro}.

    For all $\alpha\in J$ we have $C(\alpha)\subseteq\reals$
        by \cref{choice_function,pair_eq_iff,intervalclosed,subseteq}.

    Show that for all $u$ such that $u\in\img{F}{X}$ we have $u\in \productFamily{C}$.
    \begin{subproof}
        Fix $u$ such that $u\in\img{F}{X}$.
        Take $x\in X$ such that $(x,u)\in F$ by \cref{img_iff}.
        Hence $u\in\funs{\dom{R}}{\unions{\ran{R}}}$ by \cref{pair_eq_iff,indexed_product}.
        Therefore $u$ is a function and $\dom{u} = J$ by \cref{funs_is_function,funs,pair_eq_iff}.
        For all $\alpha\in J$ we have $u(\alpha)\in C(\alpha)$
            by \cref{pair_eq_iff,id_apply,id_elem_funs,choice_function,continuous,funs,subseteq,function_apply_elim,funs_is_function,img_iff}.
        For all $\alpha\in\dom{u}$ we have $u(\alpha)\in\unions{\ran{C}}$
            by \cref{subseteq,function_apply_elim,ran_intro,unions_iff,pair_eq_iff,id_apply,id_elem_funs,choice_function,continuous,funs,img_iff,funs_is_function}.
        Hence $u\in\funs{J}{\unions{\ran{C}}}$ by \cref{funs_intro}.
        Therefore $u\in\productFamily{C}$ by \cref{indexed_product,pair_eq_iff}.
    \end{subproof}
    Hence $\img{F}{X}\subseteq \productFamily{C}$ by \cref{subseteq}.

    Let $P = \productTopologyIndexed{S_0}{C}$.
    We have $\dom{S_0} = \dom{C}$ and $\dom{R} = \dom{C}$ and $\dom{T_0} = \dom{C}$ and $\dom{C}$ is inhabited by \cref{pair_eq_iff}.
    Hence $P = \subspaceTopology{P_0}{\productFamily{C}}$
        by \cref{indexed_real_product_subspace_topology,pair_eq_iff}.
    Therefore $\productFamily{C}$ is a subspace of $\productFamily{R}$ with respect to $P_0$
        by \cref{product_subspace_product,pair_eq_iff,subseteq}.
    Hence $F$ is a topological embedding from $X$ to $\productFamily{C}$ with respect to $T$ and $P$
        by \cref{topological_embedding_restrict_codomain,pair_eq_iff,subseteq}.

    For all $\alpha\in J$ we have $C(\alpha)$ is compact with respect to $S_0(\alpha)$
        by \cref{choice_function,pair_eq_iff,function_apply_intro,real_intervalclosed_compact_standard}.

    We have $\reals$ is Hausdorff with respect to $\standardTopologyR$
        by \cref{standard_metric_is_metric,metric_space_hausdorff,standard_metric_topology}.
    For all $\alpha\in J$ we have $C(\alpha)$ is Hausdorff with respect to $S_0(\alpha)$
        by \cref{subseteq,function_apply_intro,subspace_hausdorff,pair_eq_iff}.
    For all $\alpha\in\dom{C}$ we have $S_0(\alpha)$ is a topology on $C(\alpha)$
        by \cref{pair_eq_iff,subseteq,function_apply_intro,standard_basis_is_basis,basis_topology_is_topology,standard_topology_reals,subspace_of,subspace_topology_is_topology}.
    Hence $\productFamily{C}$ is compact with respect to $P$ by \cref{tychonoff_theorem,pair_eq_iff}.
    Therefore $\productFamily{C}$ is Hausdorff with respect to $P$ by \cref{product_hausdorff_product,pair_eq_iff}.

    Take $Y,S,H$ such that
        $Y$ compactifies $X$ with respect to $S$ and $T$ and
        $H$ is a topological embedding from $Y$ to $\productFamily{C}$ with respect to $S$ and $P$ and
        for all $x\in X$ we have $H(x) = F(x)$
        by \cref{induced_compactification_existence}.

    Show that for all $g$ if
        $g$ is continuous from $X$ to $\reals$ with respect to $T$ and $\standardTopologyR$ and
        $\img{g}{X}$ is bounded in $\reals$ with respect to $\standardMetricR$
    then there exists $h$ such that
        $h$ is continuous from $Y$ to $\reals$ with respect to $S$ and $\standardTopologyR$ and
        for all $x\in X$ we have $h(x) = g(x)$ and
        for all $k$ if
            $k$ is continuous from $Y$ to $\reals$ with respect to $S$ and $\standardTopologyR$ and
            for all $x\in X$ we have $k(x) = g(x)$
        then $k = h$.
    \begin{subproof}
        Fix $g$ such that
            $g$ is continuous from $X$ to $\reals$ with respect to $T$ and $\standardTopologyR$ and
            $\img{g}{X}$ is bounded in $\reals$ with respect to $\standardMetricR$.
            Hence $g\in\pow{X\times\reals}$
                by \cref{continuous,funs_is_function,funs,funs_subseteq_rels,subseteq,rels_elim,powerset_intro}.
            Take $a,b$ such that
                $a\in\reals$ and
                $b\in\reals$ and
                $a \leq b$ and
                $\img{g}{X}\subseteq\intervalclosed{a}{b}$ by \cref{bounded_reals_subset_closed_interval}.
            Therefore $g\in J$ by \cref{bounded_real_function_family_member}.
            We have $g\in\dom{C}$ and $g\in\dom{S_0}$ by \cref{funs,subseteq}.
            Let $p = \piIndex{C}{g}\circ H$.
            We have $\piIndex{C}{g}$ is continuous from $\productFamily{C}$ to $C(g)$ with respect to $P$ and $S_0(g)$
                by \cref{projection_map_indexed_continuous}.
            Hence $p$ is continuous from $Y$ to $C(g)$ with respect to $S$ and $S_0(g)$ by \cref{continuous_composite,topological_embedding}.
            Let $j = \identity{C(g)}$.
            We have $C(g)$ is a subspace of $\reals$ with respect to $\standardTopologyR$ by \cref{subspace_of,subseteq}.
            Let $h = j\circ p$.
            Hence $\ran{p}\subseteq C(g)$ by \cref{funs_ran,continuous,subseteq}.
            Therefore $\ran{p}\subseteq \dom{j}$ by \cref{id_dom,pair_eq_iff}.
            Therefore $h$ is continuous from $Y$ to $\reals$ with respect to $S$ and $\standardTopologyR$
                by \cref{function_apply_intro,continuous_expand_range,continuous,pair_eq_iff}.
            Show that for all $x\in X$ we have $h(x) = g(x)$.
            \begin{subproof}
                Fix $x\in X$.
                Hence $x\in Y$ by \cref{compactifies,subspace_of,subseteq}.
                Therefore $(x,p(x))\in p$ by \cref{function_apply_elim,continuous,funs,funs_is_function,subseteq}.
                Thus $(p(x),p(x))\in j$ by \cref{continuous,funs_type_apply,id_iff}.
                Therefore $(x,p(x))\in h$ by \cref{circ_elem_intro,pair_eq_iff}.
                Therefore $h(x) = p(x)$ by \cref{function_apply_intro,continuous,funs_is_function}.
                Let $u = H(x)$.
                $\piIndex{C}{g}\in\funs{\productFamily{C}}{C(g)}$ by \cref{projection_map_indexed_function}.
                Hence $\ran{H}\subseteq \dom{\piIndex{C}{g}}$
                    by \cref{funs,topological_embedding,funs_ran,pair_eq_iff}.
                We have $(x,u)\in H$ by \cref{function_apply_elim,topological_embedding,funs,funs_is_function,subseteq}.
                Therefore $(u,\apply{u}{g})\in\piIndex{C}{g}$
                    by \cref{projection_map_indexed,function_apply_elim,topological_embedding,funs_type_apply}.
                Thus $(x,\apply{u}{g})\in p$ by \cref{circ_elem_intro,pair_eq_iff}.
                Hence $p(x) = \apply{u}{g}$ by \cref{function_apply_intro,continuous,funs_is_function}.
                We have $p(x) = \apply{\apply{F}{x}}{g}$ by \cref{pair_eq_iff}.
                We have $(x,F(x))\in F$ by \cref{function_apply_elim,topological_embedding,funs,funs_is_function}.
                Hence $\apply{\apply{F}{x}}{g} = \apply{\apply{f}{g}}{x}$ by \cref{pair_eq_iff}.
                Thus $\apply{\apply{F}{x}}{g} = g(x)$ by \cref{id_apply,id_elem_funs,pair_eq_iff}.
                Therefore $p(x) = g(x)$ by \cref{pair_eq_iff}.
                Hence $h(x) = g(x)$ by \cref{pair_eq_iff}.
            \end{subproof}
            Show that for all $k$ if
                $k$ is continuous from $Y$ to $\reals$ with respect to $S$ and $\standardTopologyR$ and
                for all $x\in X$ we have $k(x) = g(x)$
            then $k = h$.
            \begin{subproof}
                Follows by \cref{subseteq,pair_eq_iff,compactification_real_extension_unique}.
            \end{subproof}
            Thus there exists $u$ such that
                $u$ is continuous from $Y$ to $\reals$ with respect to $S$ and $\standardTopologyR$ and
                for all $x\in X$ we have $u(x) = g(x)$ and
                for all $k$ if
                    $k$ is continuous from $Y$ to $\reals$ with respect to $S$ and $\standardTopologyR$ and
                    for all $x\in X$ we have $k(x) = g(x)$
                then $k = u$.
    \end{subproof}
    Therefore $Y$ has the bounded real extension property over $X$ with respect to $S$ and $T$
        by \cref{bounded_real_extension_property,compactifies}.
    Hence there exist $Y_0,S_1$ such that $Y_0$ has the bounded real extension property over $X$ with respect to $S_1$ and $T$.
\end{proof}

% from §38 Item 4: uniqueness of extensions to the closure
\begin{lemma}\label{extension_to_closure_unique}
    Assume $T$ is a topology on $X$.
    Assume $A\subseteq X$.
    Assume $f$ is continuous from $A$ to $Z$ with respect to $\subspaceTopology{T}{A}$ and $U$.
    Assume $Z$ is Hausdorff with respect to $U$.
    Assume $g$ is continuous from $\closure{A}{X}{T}$ to $Z$ with respect to $\subspaceTopology{T}{\closure{A}{X}{T}}$ and $U$.
    Assume $h$ is continuous from $\closure{A}{X}{T}$ to $Z$ with respect to $\subspaceTopology{T}{\closure{A}{X}{T}}$ and $U$.
    Suppose for all $x\in A$ we have $g(x) = f(x)$.
    Suppose for all $x\in A$ we have $h(x) = f(x)$.
    Then $g = h$.
\end{lemma}
\begin{proof}
    Let $B = \closure{A}{X}{T}$.
    Let $T_B = \subspaceTopology{T}{B}$.
    $B$ is a subspace of $X$ with respect to $T$ by \cref{closure,subseteq,subspace_of}.
    $T_B$ is a topology on $B$ by \cref{subspace_topology_is_topology,subspace_of,closure}.
    $A\subseteq B$ by \cref{subseteq_closure}.
    $\closure{A}{B}{T_B} = B$ by \cref{closure_in_subspace,subspace_of,inter_idempotent}.
    $g\in\funs{B}{Z}$ and $h\in\funs{B}{Z}$ by \cref{continuous,funs}.

    Show that for all $x\in B$ we have $g(x) = h(x)$.
    \begin{subproof}
        Fix $x\in B$.
        Suppose not.
        $g(x)\in Z$ and $h(x)\in Z$ by \cref{funs_type_apply,continuous}.
        Take $N_1,N_2$ such that
            $N_1$ is a neighborhood of $g(x)$ in $Z$ with respect to $U$ and
            $N_2$ is a neighborhood of $h(x)$ in $Z$ with respect to $U$ and
            $N_1$ is disjoint from $N_2$
            by \cref{hausdorff}.
        Let $W_1 = \preimg{g}{N_1}$.
        Let $W_2 = \preimg{h}{N_2}$.
        $W_1$ is open in $B$ with respect to $T_B$ and $W_2$ is open in $B$ with respect to $T_B$
            by \cref{continuous,neighborhood}.
        $x\in W_1$ and $x\in W_2$
            by \cref{function_apply_elim,continuous,funs,funs_is_function,neighborhood,preimg_iff}.
        Let $C = \{W_1,W_2\}$.
        $C\subseteq T_B$ and $C$ is finite and $C$ is inhabited
            by \cref{upair_iff,open_in,subseteq,pair_is_finite,upair_intro_left,inhabited_nonempty}.
        Let $W = \inters{C}$.
        $W$ is open in $B$ with respect to $T_B$
            by \cref{finite_intersections_open_in_topology,open_in}.
        $x\in W$ by \cref{inters_intro,upair_iff}.
        $W$ intersects $A$ by \cref{closure_iff_intersects_open}.
        Take $y$ such that $y\in W\inter A$ by \cref{intersects,nonempty_inhabited}.
        $y\in W$ and $y\in A$ by \cref{inter}.
        $y\in W_1$ and $y\in W_2$ by \cref{upair_iff,inters_subseteq_elem,subseteq}.
        Take $z_1\in N_1$ such that $(y,z_1)\in g$ by \cref{preimg_iff}.
        Take $z_2\in N_2$ such that $(y,z_2)\in h$ by \cref{preimg_iff}.
        $g(y) = z_1$ and $h(y) = z_2$
            by \cref{subseteq,preimg_subseteq_dom,function_apply_iff,funs_is_function}.
        $g(y)\in N_1$ and $h(y)\in N_2$.
        $f(y)\in N_1$ and $f(y)\in N_2$ by \cref{subseteq}.
        Contradiction by \cref{disjoint}.
    \end{subproof}
    $g = h$ by \cref{functions_eq_iff}.
\end{proof}

% from §38 helper definition 3: compact Hausdorff extension property
\begin{definition}\label{compact_hausdorff_extension_property}
    $Y$ has the compact Hausdorff extension property over $X$ with respect to $S$ and $T$ iff
        $Y$ compactifies $X$ with respect to $S$ and $T$ and
        for all $C, U, f$ if
            $U$ is a topology on $C$ and
            $C$ is compact with respect to $U$ and
            $C$ is Hausdorff with respect to $U$ and
            $f$ is continuous from $X$ to $C$ with respect to $T$ and $U$
        then there exists $g$ such that
            $g$ is continuous from $Y$ to $C$ with respect to $S$ and $U$ and
            for all $x\in X$ we have $g(x) = f(x)$ and
            for all $h$ if
                $h$ is continuous from $Y$ to $C$ with respect to $S$ and $U$ and
                for all $x\in X$ we have $h(x) = f(x)$
            then $h = g$.
\end{definition}

% from §38 helper lemma 4: compact Hausdorff spaces are completely regular
\begin{lemma}\label{compact_hausdorff_implies_completely_regular}
    Assume $U$ is a topology on $C$.
    Assume $C$ is compact with respect to $U$.
    Assume $C$ is Hausdorff with respect to $U$.
    Then $C$ is completely regular with respect to $U$.
\end{lemma}
\begin{proof}
    $C$ has closed singletons with respect to $U$
        by \cref{one_point_closed,singleton_subset_intro,singleton_iff,upair_iff,subseteq,pair_is_finite,subseteq_finite_is_finite,finite_point_set_closed_hausdorff}.
    Let $I = \intervalclosed{0}{1}$.
    Let $T_I = \subspaceTopology{\standardTopologyR}{I}$.
    Show that for all $x,A$ such that $x\in C$ and $A$ is closed in $C$ with respect to $U$ and $x\notin A$
        there exists $f$ such that
            $f$ is continuous from $C$ to $I$ with respect to $U$ and $T_I$ and
            $f(x) = 1$ and
            for all $y\in A$ we have $f(y) = 0$.
    \begin{subproof}
        Fix $x$.
        Fix $A$ such that $x\in C$ and $A$ is closed in $C$ with respect to $U$ and $x\notin A$.
        Take $f$ such that
            $f$ is continuous from $C$ to $I$ with respect to $U$ and $T_I$ and
            $\img{f}{A}\subseteq \{0\}$ and
            $\img{f}{\{x\}}\subseteq \{1\}$
            by \cref{one_point_closed,singleton_subset_intro,singleton_iff,disjoint,compact_hausdorff_normal,urysohn_unit_interval}.
        $f(x) = 1$ and for all $y\in A$ we have $f(y) = 0$
            by \cref{closed_in,singleton_intro,continuous,funs,continuous,funs_is_function,function_apply_elim,img_iff,subseteq,singleton_elim}.
    \end{subproof}
    $C$ is completely regular with respect to $U$ by \cref{completely_regular_space}.
\end{proof}

% from §38 helper lemma 10: coordinatewise continuous families induce continuous maps into indexed real products
\begin{lemma}\label{continuous_real_family_map}
    Assume $J$ is inhabited.
    Let $R = \{(\alpha,\reals) \mid \alpha\in J\}$.
    Let $T_0 = \{(\alpha,\standardTopologyR) \mid \alpha\in J\}$.
    Let $P = \productTopologyIndexed{T_0}{R}$.
    Assume $b$ is a function on $J$.
    Suppose for all $\alpha\in J$ we have $\apply{b}{\alpha}$ is continuous from $Y$ to $\reals$ with respect to $S$ and $\standardTopologyR$.
    Let $G = \{(y,z) \mid y\in Y, z\in\productFamily{R} \mid
        \forall \alpha\in J. z(\alpha) = \apply{\apply{b}{\alpha}}{y}\}$.
    Then $G$ is continuous from $Y$ to $\productFamily{R}$ with respect to $S$ and $P$.
\end{lemma}
\begin{proof}
    $\standardTopologyR$ is a topology on $\reals$
        by \cref{standard_basis_is_basis,basis_topology_is_topology,standard_topology_reals}.
    For all $\alpha,y_1,y_2$ such that $(\alpha,y_1)\in R$ and $(\alpha,y_2)\in R$ we have $y_1 = y_2$ by \cref{pair_eq_iff}.
    $R$ is a function on $J$ by \cref{rightunique,relation,dom_intro,dom_iff,pair_eq_iff,subseteq,subseteq_antisymmetric}.
    For all $\alpha,y_1,y_2$ such that $(\alpha,y_1)\in T_0$ and $(\alpha,y_2)\in T_0$ we have $y_1 = y_2$ by \cref{pair_eq_iff}.
    $T_0$ is a function on $J$ by \cref{rightunique,relation,dom_intro,dom_iff,pair_eq_iff,subseteq,subseteq_antisymmetric}.
    $\dom{R} = J$ and $\dom{T_0} = J$ by \cref{funs}.
    For all $\alpha\in J$ we have $R(\alpha) = \reals$ and $T_0(\alpha) = \standardTopologyR$ by \cref{function_apply_intro}.
    $\productFamily{R} = \realsPower{J}$ by \cref{product_family_reals_power}.
    $P$ is a topology on $\productFamily{R}$
        by \cref{product_subbasis_is_subbasis,product_subbasis_finite_intersections_basis,basis_topology_is_topology,basis_for_topology,standard_topology_reals}.

    For all $y,z_1,z_2$ such that $(y,z_1)\in G$ and $(y,z_2)\in G$ we have $z_1 = z_2$
        by \cref{pair_eq_iff,product_family_reals_power,subseteq,reals_power,tuple_power,funs_is_function,funs,functions_eq_iff}.
    $G$ is a function by \cref{rightunique,relation}.
    Show that for all $y\in Y$ there exists $z$ such that $(y,z)\in G$.
    \begin{subproof}
        Fix $y\in Y$.
        Let $z = \{(\alpha,\apply{\apply{b}{\alpha}}{y}) \mid \alpha\in J\}$.
        For all $\alpha,r_1,r_2$ such that $(\alpha,r_1)\in z$ and $(\alpha,r_2)\in z$ we have $r_1 = r_2$ by \cref{pair_eq_iff}.
        $z$ is a function by \cref{relation,rightunique,pair_eq_iff}.
        $\dom{z} = J$ by \cref{subseteq_antisymmetric,subseteq,dom_intro,dom_iff,pair_eq_iff}.
        For all $\alpha\in\dom{z}$ we have $z(\alpha)\in\reals$
            by \cref{subseteq,continuous,funs_type_apply,function_apply_intro,pair_eq_iff}.
        $z\in\productFamily{R}$ by \cref{funs_intro,reals_power,tuple_power,subseteq}.
        For all $\alpha\in J$ we have $z(\alpha) = \apply{\apply{b}{\alpha}}{y}$ by \cref{function_apply_intro}.
        $(y,z)\in G$ by \cref{pair_eq_iff}.
    \end{subproof}
    $\dom{G} = Y$ by \cref{subseteq,dom_iff,pair_eq_iff,subseteq_antisymmetric}.
    $G\in\funs{Y}{\productFamily{R}}$ by \cref{funs_intro,function_apply_iff,pair_eq_iff}.

    Show that for all $\alpha\in\dom{R}$ we have $\piIndex{R}{\alpha}\circ G$ is continuous from $Y$ to $\reals$ with respect to $S$ and $\standardTopologyR$.
    \begin{subproof}
        Fix $\alpha\in\dom{R}$.
        $\piIndex{R}{\alpha}\circ G\in\funs{Y}{\reals}$ by \cref{projection_map_indexed_function,funs_circ,funs_elim}.
        $\apply{b}{\alpha}\in\funs{Y}{\reals}$ and $\apply{b}{\alpha}$ is a function and $\dom{\apply{b}{\alpha}} = Y$
            by \cref{continuous,funs_is_function,funs,subseteq,pair_eq_iff}.
        For all $y\in Y$ we have $\apply{\piIndex{R}{\alpha}\circ G}{y} = \apply{\apply{b}{\alpha}}{y}$ by
            \cref{function_apply_elim,funs_is_function,funs,pair_eq_iff,funs_type_apply,projection_map_indexed,circ_elem_intro,funs_elim,function_apply_intro}.
        $\piIndex{R}{\alpha}\circ G = \apply{b}{\alpha}$ by \cref{functions_eq_iff}.
        $\piIndex{R}{\alpha}\circ G$ is continuous from $Y$ to $\reals$ with respect to $S$ and $\standardTopologyR$
            by \cref{continuous_eq_subst,subseteq}.
    \end{subproof}
    For all $\alpha\in\dom{R}$ we have $T_0(\alpha)$ is a topology on $R(\alpha)$
        by \cref{pair_eq_iff,function_apply_intro}.
    Take $\alpha_0$ such that $\alpha_0\in J$.
    $\apply{b}{\alpha_0}$ is continuous from $Y$ to $\reals$ with respect to $S$ and $\standardTopologyR$ by assumption.
    $S$ is a topology on $Y$ by \cref{continuous}.
    $G$ is continuous from $Y$ to $\productFamily{R}$ with respect to $S$ and $P$
        by \cref{continuous_map_into_indexed_product,pair_eq_iff}.
\end{proof}

% from §38 helper lemma 11: subsets of bounded sets are bounded
\begin{lemma}\label{metric_bounded_subset}
    Assume $A$ is bounded in $X$ with respect to $d$.
    Suppose $B\subseteq A$.
    Then $B$ is bounded in $X$ with respect to $d$.
\end{lemma}
\begin{proof}
    Follows by \cref{metric_bounded,subseteq,subseteq_transitive}.
\end{proof}

% from §38 helper lemma 12: subsets of closed real intervals are bounded
\begin{lemma}\label{intervalclosed_subset_bounded_standard}
    Assume $a\in\reals$.
    Assume $b\in\reals$.
    Assume $a \leq b$.
    Suppose $A\subseteq\intervalclosed{a}{b}$.
    Then $A$ is bounded in $\reals$ with respect to $\standardMetricR$.
\end{lemma}
\begin{proof}
    $\standardMetricR$ is a metric on $\reals$ by \cref{standard_metric_is_metric}.
    $A\subseteq\reals$ by \cref{intervalclosed,subseteq,subseteq_transitive}.
    Let $M = b - a$.
    $M\in\reals$ by \cref{real_add_inverse,real_sub,real_add_closed}.
    Show that for all $u, v$ such that $u\in A$ and $v\in A$ we have $\apply{\standardMetricR}{(u,v)} \leq M$.
    \begin{subproof}
        Fix $u$.
        Fix $v$ such that $u\in A$ and $v\in A$.
        $u\in\intervalclosed{a}{b}$ and $v\in\intervalclosed{a}{b}$ by \cref{subseteq}.
        $u\in\reals$ and $a \leq u$ and $u \leq b$ and
            $v\in\reals$ and $a \leq v$ and $v \leq b$ by \cref{intervalclosed,subseteq}.
        $b \leq v + (b - a)$
            by \cref{real_add_closed,real_leq_add,real_sub,real_add_assoc,real_add_inverse,real_add_ident,real_add_comm}.
        $v + (b - a)\in\reals$ by \cref{real_add_closed,real_sub}.
        $u \leq v + (b - a)$ by \cref{real_leq_trans}.
        $v + M = v + (b - a)$ by \cref{pair_eq_iff}.
        $v + M\in\reals$ by \cref{real_add_closed}.
        $u \leq v + M$ by \cref{pair_eq_iff}.
        $\neg{v}\in\reals$ by \cref{real_add_inverse}.
        $u + \neg{v}\in\reals$ by \cref{real_add_closed}.
        $u + \neg{v} \leq (v + M) + \neg{v}$ by \cref{real_leq_add}.
        $(v + M) + \neg{v} = M$ by \cref{real_add_assoc,real_add_inverse,real_add_comm,real_add_ident}.
        $u + \neg{v} \leq M$ by \cref{real_leq_trans,pair_eq_iff}.
        $u - v \leq M$ by \cref{real_sub}.
        $u - v\in\reals$ by \cref{real_sub,real_add_closed}.
        $b \leq u + (b - a)$
            by \cref{real_add_closed,real_leq_add,real_sub,real_add_assoc,real_add_inverse,real_add_ident,real_add_comm}.
        $u + M = u + (b - a)$ by \cref{pair_eq_iff}.
        $u + M\in\reals$ by \cref{real_add_closed}.
        $b \leq u + M$ by \cref{pair_eq_iff}.
        $v \leq u + M$ by \cref{real_leq_trans}.
        $\neg{u}\in\reals$ by \cref{real_add_inverse}.
        $v + \neg{u}\in\reals$ by \cref{real_add_closed}.
        $v + \neg{u} \leq (u + M) + \neg{u}$ by \cref{real_leq_add}.
        $(u + M) + \neg{u} = M$ by \cref{real_add_assoc,real_add_inverse,real_add_comm,real_add_ident}.
        $v + \neg{u} \leq M$ by \cref{real_leq_trans,pair_eq_iff}.
        $v - u \leq M$ by \cref{real_sub}.
        $v - u = \neg{(u - v)}$ by \cref{real_sub_swap_neg}.
        $\neg{(u - v)} \leq M$.
        $\abs{u - v} \leq M$ by \cref{real_abs_leq_if_pm_leq}.
        $\apply{\standardMetricR}{(u,v)} = \abs{u - v}$ by \cref{standard_metric_value}.
        $\apply{\standardMetricR}{(u,v)} \leq M$ by \cref{pair_eq_iff}.
    \end{subproof}
    $A$ is bounded in $\reals$ with respect to $\standardMetricR$ by \cref{metric_bounded}.
\end{proof}

\begin{theorem}\label{bounded_real_property_implies_compact_hausdorff_property}
    Assume $Y$ has the bounded real extension property over $X$ with respect to $S$ and $T$.
    Then $Y$ has the compact Hausdorff extension property over $X$ with respect to $S$ and $T$.
\end{theorem}
\begin{proof}
    We have $Y$ compactifies $X$ with respect to $S$ and $T$ by \cref{bounded_real_extension_property}.
    Show that for all $C, U, f$ if
        $U$ is a topology on $C$ and
        $C$ is compact with respect to $U$ and
        $C$ is Hausdorff with respect to $U$ and
        $f$ is continuous from $X$ to $C$ with respect to $T$ and $U$
    then there exists $g$ such that
        $g$ is continuous from $Y$ to $C$ with respect to $S$ and $U$ and
        for all $x\in X$ we have $g(x) = f(x)$ and
        for all $h$ if
            $h$ is continuous from $Y$ to $C$ with respect to $S$ and $U$ and
            for all $x\in X$ we have $h(x) = f(x)$
        then $h = g$.
    \begin{subproof}
        Fix $C$.
        Fix $U$.
        Fix $f$ such that
            $U$ is a topology on $C$ and
            $C$ is compact with respect to $U$ and
            $C$ is Hausdorff with respect to $U$ and
            $f$ is continuous from $X$ to $C$ with respect to $T$ and $U$.
        $C$ is completely regular with respect to $U$
            and $C$ has closed singletons with respect to $U$
            by \cref{compact_hausdorff_implies_completely_regular,completely_regular_space}.
        Let $J = \{g\in\pow{C\times\reals} \mid \text{$g$ is continuous from $C$ to $\reals$ with respect to $U$ and $\standardTopologyR$ and
            $\img{g}{C}$ is bounded in $\reals$ with respect to $\standardMetricR$}\}$.
        Let $z = \{(c,0) \mid c\in C\}$.
        $z$ is a function from $C$ to $\reals$ by \cref{constant_map_function,real_add_ident}.
        For all $c\in C$ we have $z(c) = 0$ by \cref{constant_map_function,funs_is_function,function_apply_intro}.
        $\standardTopologyR$ is a topology on $\reals$
            by \cref{standard_basis_is_basis,basis_topology_is_topology,standard_topology_reals}.
        $z$ is continuous from $C$ to $\reals$ with respect to $U$ and $\standardTopologyR$
            by \cref{real_add_ident,constant_continuous}.
        $0\in\reals$ by \cref{real_add_ident}.
        $z\in\pow{C\times\reals}$
            by \cref{continuous,funs_is_function,funs,funs_subseteq_rels,subseteq,rels_elim,powerset_intro}.
        $\img{z}{C}\subseteq \intervalclosed{0}{0}$
            by \cref{img_iff,function_apply_intro,funs_is_function,pair_eq_iff,intervalclosed,subseteq}.
        $\img{z}{C}$ is bounded in $\reals$ with respect to $\standardMetricR$
            by \cref{intervalclosed_subset_bounded_standard,real_add_ident}.
        $z\in J$ by \cref{pair_eq_iff}.
        $J$ is inhabited.

        Let $e = \identity{J}$.
        $e$ is a function on $J$ by \cref{id_is_relation,id_rightunique,id_dom,relation,rightunique,subseteq,subseteq_antisymmetric}.
        For all $\alpha\in J$ we have $\apply{e}{\alpha}$ is continuous from $C$ to $\reals$ with respect to $U$ and $\standardTopologyR$
            by \cref{id_apply,id_elem_funs}.
        Show that for all $c,V$ such that $c\in C$ and $V$ is a neighborhood of $c$ in $C$ with respect to $U$
            there exists $\alpha\in J$ such that
                $0 < \apply{\apply{e}{\alpha}}{c}$ and
                for all $y\in C\setminus V$ we have $\apply{\apply{e}{\alpha}}{y} = 0$.
        \begin{subproof}
            Fix $c$.
            Fix $V$ such that $c\in C$ and $V$ is a neighborhood of $c$ in $C$ with respect to $U$.
            Let $I_0 = \intervalclosed{0}{1}$.
            Let $T_I = \subspaceTopology{\standardTopologyR}{I_0}$.
            Let $j_0 = \identity{I_0}$.
            $0\in\reals$ and $1\in\reals$ and $0 < 1$
                by \cref{real_add_ident,real_one_in_naturals,real_naturals_subset_reals,subseteq,real_zero_lt_one}.
            $I_0\subseteq\reals$ and $I_0$ is a subspace of $\reals$ with respect to $\standardTopologyR$
                by \cref{real_lt_imp_leq,intervalclosed,subseteq,subspace_of}.
            $T_I$ is a topology on $I_0$ by \cref{subspace_topology_is_topology}.
            $V$ is open in $C$ with respect to $U$ and $c\in V$ by \cref{neighborhood}.
            Let $A = C\setminus V$.
            $V\subseteq C$ by \cref{open_in,topology_on,subseteq,pow_iff}.
            $C\setminus A = V$ by \cref{double_relative_complement,pair_eq_iff}.
            $C\setminus A$ is open in $C$ with respect to $U$ by \cref{open_in,pair_eq_iff}.
            $A$ is closed in $C$ with respect to $U$ by \cref{closed_in,setminus_subseteq}.
            $c\notin A$ by \cref{setminus_elim_right}.
            Take $g$ such that
                $g$ is continuous from $C$ to $I_0$ with respect to $U$ and $T_I$ and
                $g(c) = 1$ and
                for all $y\in A$ we have $g(y) = 0$
                by \cref{completely_regular_space}.
            Let $g_1 = j_0\circ g$.
            $g_1$ is continuous from $C$ to $\reals$ with respect to $U$ and $\standardTopologyR$
                by \cref{continuous_expand_range}.
            $g_1\in\funs{C}{\reals}$ and $g_1$ is a function and $\dom{g_1} = C$
                by \cref{continuous,funs_is_function,funs}.
            $g_1\in\pow{C\times\reals}$
                by \cref{funs_subseteq_rels,subseteq,rels_elim,powerset_intro}.
            For all $y\in C$ we have $g_1(y) = g(y)$
                by \cref{continuous,funs,continuous,funs_type_apply,function_apply_elim,id_iff,circ_elem_intro,pair_eq_iff,function_apply_intro,continuous,funs_is_function}.
            $\img{g_1}{C}\subseteq \intervalclosed{0}{1}$
                by \cref{img_iff,function_apply_intro,subseteq,continuous,funs_type_apply,pair_eq_iff}.
            $\img{g_1}{C}$ is bounded in $\reals$ with respect to $\standardMetricR$
                by \cref{intervalclosed_subset_bounded_standard,real_add_ident,real_one_in_naturals,real_naturals_subset_reals,subseteq,real_zero_lt_one,real_lt_imp_leq}.
            $g_1\in J$ by \cref{pair_eq_iff}.
            Take $\alpha\in J$ such that $\alpha = g_1$.
            $\apply{e}{\alpha} = g_1$ by \cref{id_apply,id_elem_funs,pair_eq_iff}.
            $0 < \apply{\apply{e}{\alpha}}{c}$ by \cref{pair_eq_iff,real_zero_lt_one}.
            For all $y\in C\setminus V$ we have $\apply{\apply{e}{\alpha}}{y} = 0$
                by \cref{pair_eq_iff,setminus_elim_left,subseteq}.
        \end{subproof}

        Let $R = \{(\alpha,\reals) \mid \alpha\in J\}$.
        Let $T_0 = \{(\alpha,\standardTopologyR) \mid \alpha\in J\}$.
        For all $\alpha,y_1,y_2$ such that $(\alpha,y_1)\in R$ and $(\alpha,y_2)\in R$ we have $y_1 = y_2$ by \cref{pair_eq_iff}.
        $R$ is a function on $J$ by \cref{rightunique,relation,dom_intro,dom_iff,pair_eq_iff,subseteq,subseteq_antisymmetric}.
        For all $\alpha,y_1,y_2$ such that $(\alpha,y_1)\in T_0$ and $(\alpha,y_2)\in T_0$ we have $y_1 = y_2$ by \cref{pair_eq_iff}.
        $T_0$ is a function on $J$ by \cref{rightunique,relation,dom_intro,dom_iff,pair_eq_iff,subseteq,subseteq_antisymmetric}.
        $\dom{R} = J$ and $\dom{T_0} = J$ by \cref{funs}.
        For all $\alpha\in J$ we have $R(\alpha) = \reals$ and $T_0(\alpha) = \standardTopologyR$ by \cref{function_apply_intro}.

        Let $P = \productTopologyIndexed{T_0}{R}$.
        Let $F = \{(c,y) \mid c\in C, y\in\productFamily{R} \mid
            \forall \alpha\in J. \apply{y}{\alpha} = \apply{\apply{e}{\alpha}}{c}\}$.
        $F$ is a topological embedding from $C$ to $\productFamily{R}$ with respect to $U$ and $P$
            by \cref{urysohn_embedding_theorem_real_power}.
        $F$ is continuous from $C$ to $\productFamily{R}$ with respect to $U$ and $P$
            and $P$ is a topology on $\productFamily{R}$ by \cref{topological_embedding,continuous}.

        $\reals$ is Hausdorff with respect to $\standardTopologyR$
            by \cref{standard_metric_is_metric,metric_space_hausdorff,standard_metric_topology}.
        For all $\alpha\in\dom{R}$ we have $T_0(\alpha)$ is a topology on $R(\alpha)$
            by \cref{pair_eq_iff,function_apply_intro}.
        For all $\alpha\in\dom{R}$ we have $R(\alpha)$ is Hausdorff with respect to $T_0(\alpha)$
            by \cref{pair_eq_iff,function_apply_intro}.
        $\productFamily{R}$ is Hausdorff with respect to $P$ by \cref{product_hausdorff_product,pair_eq_iff}.

        $\img{F}{C}$ is compact with respect to $\subspaceTopology{P}{\img{F}{C}}$ by \cref{continuous_image_compact}.
        $\img{F}{C}\subseteq\productFamily{R}$ by \cref{img_subseteq_ran,funs_ran,topological_embedding,subseteq_transitive}.
        $\img{F}{C}$ is closed in $\productFamily{R}$ with respect to $P$ by \cref{compact_subspace_closed}.

        Let $M = \{w\in J\times\pow{Y\times\reals} \mid \text{there exists $\alpha\in J$ such that there exists $h$ such that
            $w = (\alpha,h)$ and
            $h$ is continuous from $Y$ to $\reals$ with respect to $S$ and $\standardTopologyR$ and
            for all $x\in X$ we have $h(x) = \apply{\alpha}{f(x)}$}\}$.
        $M$ is a relation by \cref{times_elem_is_tuple,relation}.
        Show that for all $\alpha\in J$ there exists $h$ such that $\alpha\mathrel{M} h$.
        \begin{subproof}
            Fix $\alpha\in J$.
            $\alpha$ is continuous from $C$ to $\reals$ with respect to $U$ and $\standardTopologyR$ and
                $\img{\alpha}{C}$ is bounded in $\reals$ with respect to $\standardMetricR$ by \cref{pair_eq_iff}.
            Let $k = \alpha\circ f$.
            $k$ is continuous from $X$ to $\reals$ with respect to $T$ and $\standardTopologyR$
                by \cref{continuous_composite}.
            Show that for all $u$ such that $u\in\img{k}{X}$ we have $u\in\img{\alpha}{C}$.
            \begin{subproof}
                Fix $u$ such that $u\in\img{k}{X}$.
                Take $x\in X$ such that $(x,u)\in k$ by \cref{img_iff}.
                $u = k(x)$ by \cref{function_apply_intro,continuous,funs_is_function}.
                $f(x)\in C$ by \cref{continuous,funs_type_apply}.
                $k(x) = \apply{\alpha}{f(x)}$
                    by \cref{circ_apply,continuous,funs_is_function,continuous,funs,subseteq,continuous,funs_type_apply,funs_ran,pair_eq_iff}.
                $(f(x),\apply{\alpha}{f(x)})\in \alpha$
                    by \cref{function_apply_elim,continuous,funs_is_function,funs,subseteq}.
                $(f(x),u)\in \alpha$ by \cref{pair_eq_iff}.
                $u\in\img{\alpha}{C}$ by \cref{img_iff}.
            \end{subproof}
            $\img{k}{X}\subseteq \img{\alpha}{C}$ by \cref{subseteq}.
            $\img{k}{X}$ is bounded in $\reals$ with respect to $\standardMetricR$ by \cref{metric_bounded_subset}.
            Take $h$ such that
                $h$ is continuous from $Y$ to $\reals$ with respect to $S$ and $\standardTopologyR$ and
                for all $x\in X$ we have $h(x) = k(x)$
                by \cref{bounded_real_extension_property}.
            $h\in\pow{Y\times\reals}$
                by \cref{continuous,funs_is_function,funs,funs_subseteq_rels,subseteq,rels_elim,powerset_intro}.
            For all $x\in X$ we have $h(x) = \alpha(f(x))$
                by \cref{subseteq,circ_apply,continuous,funs,funs_is_function,funs_ran,pair_eq_iff}.
            $(\alpha,h)\in J\times\pow{Y\times\reals}$ by \cref{times_tuple_intro}.
            $(\alpha,h)\in M$ by \cref{pair_eq_iff}.
            $\alpha\mathrel{M} h$.
        \end{subproof}
        Take $b$ such that $b$ is a function on $J$ and for all $\alpha\in J$ we have $\alpha\mathrel{M}\apply{b}{\alpha}$
            by \cref{choice_function}.
        $b$ is a function and $\dom{b} = J$ by \cref{funs}.
        For all $\alpha\in J$ we have $\apply{b}{\alpha}$ is continuous from $Y$ to $\reals$ with respect to $S$ and $\standardTopologyR$
            by \cref{choice_function,pair_eq_iff}.

        Let $G = \{(y,z) \mid y\in Y, z\in\productFamily{R} \mid
            \forall \alpha\in J. z(\alpha) = \apply{\apply{b}{\alpha}}{y}\}$.
        $G$ is continuous from $Y$ to $\productFamily{R}$ with respect to $S$ and $P$
            by \cref{continuous_real_family_map}.

        Show that for all $x\in X$ we have $G(x) = F(f(x))$.
        \begin{subproof}
            Fix $x\in X$.
            $x\in Y$ by \cref{compactifies,subspace_of,subseteq}.
            $f(x)\in C$ by \cref{continuous,funs_type_apply}.
            $G(x)\in\productFamily{R}$ and $F(f(x))\in\productFamily{R}$
                by \cref{continuous,funs_type_apply,topological_embedding,funs_type_apply}.
            $G(x)\in\funs{J}{\reals}$ and $F(f(x))\in\funs{J}{\reals}$
                by \cref{product_family_reals_power,reals_power,tuple_power,subseteq}.
            $(x,G(x))\in G$ by \cref{function_apply_elim,continuous,funs_is_function,funs,subseteq}.
            $(f(x),F(f(x)))\in F$ by \cref{function_apply_elim,topological_embedding,funs_is_function,funs,subseteq}.
            For all $\alpha\in J$ we have $\apply{\apply{G}{x}}{\alpha} = \apply{\apply{F}{f(x)}}{\alpha}$
                by \cref{pair_eq_iff,choice_function,subseteq,id_apply,id_elem_funs}.
            $G(x) = F(f(x))$ by \cref{functions_eq_iff}.
        \end{subproof}

        For all $u$ such that $u\in\img{G}{X}$ we have $u\in \img{F}{C}$
            by \cref{img_iff,function_apply_intro,continuous,funs_is_function,subseteq,continuous,funs_type_apply,pair_eq_iff,topological_embedding,funs,function_apply_elim,topological_embedding,funs_is_function}.
        $\img{G}{X}\subseteq \img{F}{C}$ by \cref{subseteq}.

        $X\subseteq Y$ by \cref{compactifies,subspace_of}.
        $\closure{X}{Y}{S} = Y$ by \cref{compactifies}.
        $\img{G}{Y}\subseteq \closure{\img{G}{X}}{\productFamily{R}}{P}$ by \cref{continuous_closure_image_subset,pair_eq_iff}.
        $\closure{\img{G}{X}}{\productFamily{R}}{P}\subseteq \img{F}{C}$ by \cref{img_subseteq_ran,funs_ran,continuous,subseteq_transitive,closure_subseteq_closed}.
        $\img{G}{Y}\subseteq \img{F}{C}$ by \cref{subseteq_transitive}.

        Let $P_C = \subspaceTopology{P}{\img{F}{C}}$.
        $G$ is continuous from $Y$ to $\img{F}{C}$ with respect to $S$ and $P_C$ by \cref{subspace_of,subseteq,continuous_restrict_range}.
        $F$ is a homeomorphism between $C$ and $\img{F}{C}$ with respect to $U$ and $P_C$ by \cref{topological_embedding}.
        $\converse{F}$ is continuous from $\img{F}{C}$ to $C$ with respect to $P_C$ and $U$ by \cref{homeomorphism}.
        Let $g = \converse{F}\circ G$.
        $g$ is continuous from $Y$ to $C$ with respect to $S$ and $U$ by \cref{continuous_composite}.

        Show that for all $x\in X$ we have $g(x) = f(x)$.
        \begin{subproof}
            Fix $x\in X$.
            $x\in Y$ by \cref{compactifies,subspace_of,subseteq}.
            $G(x) = F(f(x))$ by \cref{subseteq}.
            $f(x)\in C$ by \cref{continuous,funs_type_apply}.
            $\converse{F}\in\funs{\img{F}{C}}{C}$ and $\dom{\converse{F}} = \img{F}{C}$
                by \cref{homeomorphism,bijections,bijective_converse_are_funs,funs}.
            $(x,G(x))\in G$ by \cref{function_apply_elim,continuous,funs_is_function,funs,subseteq}.
            $G(x)\in\dom{\converse{F}}$ by \cref{continuous,funs_type_apply,pair_eq_iff}.
            $(G(x),\apply{\converse{F}}{G(x)})\in\converse{F}$ by \cref{function_apply_elim,bijective_converse_are_funs,funs_is_function,homeomorphism,bijections}.
            $(x,\apply{\converse{F}}{G(x)})\in g$ by \cref{circ_elem_intro,pair_eq_iff,continuous,funs,bijective_converse_are_funs,homeomorphism,bijections,funs_is_function,funs_ran,subseteq}.
            $g(x) = \apply{\converse{F}}{G(x)}$ by \cref{function_apply_intro,continuous,funs_is_function}.
            $g(x) = \apply{\converse{F}}{F(f(x))}$ by \cref{pair_eq_iff}.
            $(f(x),F(f(x)))\in F$ by \cref{function_apply_elim,topological_embedding,funs_is_function,funs,subseteq}.
            $(F(f(x)),f(x))\in\converse{F}$ by \cref{converse_iff}.
            $\apply{\converse{F}}{F(f(x))} = f(x)$
                by \cref{function_apply_intro,bijective_converse_are_funs,funs_is_function,homeomorphism,bijections}.
            $g(x) = f(x)$ by \cref{pair_eq_iff}.
        \end{subproof}

        Let $S_Y = \subspaceTopology{S}{Y}$.
        $S_Y = S$
            by \cref{compactifies,subspace_topology,topology_on,pow_iff,subseteq,inter_absorb_supseteq_left,subseteq_antisymmetric}.
        $f$ is continuous from $X$ to $C$ with respect to $\subspaceTopology{S}{X}$ and $U$
            and $g$ is continuous from $Y$ to $C$ with respect to $S_Y$ and $U$ by \cref{compactifies,pair_eq_iff}.
        $S$ is a topology on $Y$ and $X\subseteq Y$ and $\closure{X}{Y}{S} = Y$ by \cref{compactifies,subspace_of}.
        For all $x\in X$ we have $g(x) = f(x)$ by \cref{subseteq}.
        Show that for all $h$ if
            $h$ is continuous from $Y$ to $C$ with respect to $S$ and $U$ and
            for all $x\in X$ we have $h(x) = f(x)$
        then $h = g$.
        \begin{subproof}
            Follows by \cref{pair_eq_iff,extension_to_closure_unique}.
        \end{subproof}
        There exists $g_1$ such that
            $g_1$ is continuous from $Y$ to $C$ with respect to $S$ and $U$ and
            for all $x\in X$ we have $g_1(x) = f(x)$ and
            for all $h$ if
                $h$ is continuous from $Y$ to $C$ with respect to $S$ and $U$ and
                for all $x\in X$ we have $h(x) = f(x)$
            then $h = g_1$.
    \end{subproof}
    $Y$ has the compact Hausdorff extension property over $X$ with respect to $S$ and $T$
        by \cref{compact_hausdorff_extension_property,bounded_real_extension_property}.
\end{proof}

% from §38 Item 6: uniqueness up to equivalence
\begin{theorem}\label{bounded_real_property_unique_up_to_equivalence}
    Assume $Y_1$ has the bounded real extension property over $X$ with respect to $S_1$ and $T$.
    Assume $Y_2$ has the bounded real extension property over $X$ with respect to $S_2$ and $T$.
    Then $Y_1$ is equivalent to $Y_2$ over $X$ with respect to $S_1$ and $S_2$ and $T$.
\end{theorem}
\begin{proof}
    We have $Y_1$ has the compact Hausdorff extension property over $X$ with respect to $S_1$ and $T$
        by \cref{bounded_real_property_implies_compact_hausdorff_property}.
    $Y_2$ has the compact Hausdorff extension property over $X$ with respect to $S_2$ and $T$
        by \cref{bounded_real_property_implies_compact_hausdorff_property}.
    $Y_1$ compactifies $X$ with respect to $S_1$ and $T$ and
        $Y_2$ compactifies $X$ with respect to $S_2$ and $T$
        by \cref{compact_hausdorff_extension_property}.
    $S_1$ is a topology on $Y_1$ and
        $X$ is a subspace of $Y_1$ with respect to $S_1$ and
        $Y_1$ is compact with respect to $S_1$ and
        $Y_1$ is Hausdorff with respect to $S_1$
        by \cref{compactifies}.
    $S_2$ is a topology on $Y_2$ and
        $X$ is a subspace of $Y_2$ with respect to $S_2$ and
        $Y_2$ is compact with respect to $S_2$ and
        $Y_2$ is Hausdorff with respect to $S_2$
        by \cref{compactifies}.

    Let $j = \identity{X}$.
    $j$ is continuous from $X$ to $Y_1$ with respect to $T$ and $S_1$
        and $j$ is continuous from $X$ to $Y_2$ with respect to $T$ and $S_2$
        by \cref{inclusion_continuous,compactifies,pair_eq_iff}.

    Take $h_{12}$ such that
        $h_{12}$ is continuous from $Y_1$ to $Y_2$ with respect to $S_1$ and $S_2$ and
        for all $x\in X$ we have $h_{12}(x) = j(x)$ and
        for all $h$ if
            $h$ is continuous from $Y_1$ to $Y_2$ with respect to $S_1$ and $S_2$ and
            for all $x\in X$ we have $h(x) = j(x)$
        then $h = h_{12}$
        by \cref{compact_hausdorff_extension_property}.
    Take $h_{21}$ such that
        $h_{21}$ is continuous from $Y_2$ to $Y_1$ with respect to $S_2$ and $S_1$ and
        for all $x\in X$ we have $h_{21}(x) = j(x)$ and
        for all $h$ if
            $h$ is continuous from $Y_2$ to $Y_1$ with respect to $S_2$ and $S_1$ and
            for all $x\in X$ we have $h(x) = j(x)$
        then $h = h_{21}$
        by \cref{compact_hausdorff_extension_property}.

    $\subspaceTopology{S_1}{Y_1} = S_1$
        by \cref{subspace_topology,topology_on,pow_iff,subseteq,inter_absorb_supseteq_left,subseteq_antisymmetric}.

    Let $i_1 = \identity{Y_1}$.
    $i_1$ is continuous from $Y_1$ to $Y_1$ with respect to $S_1$ and $S_1$
        and $h_{21}\circ h_{12}$ is continuous from $Y_1$ to $Y_1$ with respect to $S_1$ and $S_1$
        by \cref{identity_continuous,continuous_composite}.
    For all $x\in X$ we have $i_1(x) = j(x)$
        by \cref{compactifies,subspace_of,subseteq,id_apply,id_elem_funs,pair_eq_iff}.
    $h_{12}\in\funs{Y_1}{Y_2}$ and $h_{21}\in\funs{Y_2}{Y_1}$ and $h_{12}$ is a function and $h_{21}$ is a function
        by \cref{continuous,funs,funs_is_function}.
    For all $x\in X$ we have $h_{12}(x) = j(x)$ by assumption.
    For all $x\in X$ we have $(h_{21}\circ h_{12})(x) = h_{21}(h_{12}(x))$
        by \cref{circ_apply,continuous,funs,funs_is_function,funs_ran,compactifies,subspace_of,subseteq}.
    For all $x\in X$ we have $h_{21}(x) = j(x)$ by assumption.
    For all $x\in X$ we have $(h_{21}\circ h_{12})(x) = j(x)$ by \cref{pair_eq_iff,id_apply,id_elem_funs}.
    $i_1$ is continuous from $\closure{X}{Y_1}{S_1}$ to $Y_1$
        with respect to $\subspaceTopology{S_1}{\closure{X}{Y_1}{S_1}}$ and $S_1$
        and $h_{21}\circ h_{12}$ is continuous from $\closure{X}{Y_1}{S_1}$ to $Y_1$
        with respect to $\subspaceTopology{S_1}{\closure{X}{Y_1}{S_1}}$ and $S_1$
        by \cref{compactifies,pair_eq_iff}.
    Show that $h_{21}\circ h_{12} = i_1$.
    \begin{subproof}
        Follows by \cref{compactifies,subspace_of,extension_to_closure_unique,subseteq}.
    \end{subproof}

    $\subspaceTopology{S_2}{Y_2} = S_2$
        by \cref{subspace_topology,topology_on,pow_iff,subseteq,inter_absorb_supseteq_left,subseteq_antisymmetric}.

    Let $i_2 = \identity{Y_2}$.
    $i_2$ is continuous from $Y_2$ to $Y_2$ with respect to $S_2$ and $S_2$
        and $h_{12}\circ h_{21}$ is continuous from $Y_2$ to $Y_2$ with respect to $S_2$ and $S_2$
        by \cref{identity_continuous,continuous_composite}.
    For all $x\in X$ we have $i_2(x) = j(x)$
        by \cref{compactifies,subspace_of,subseteq,id_apply,id_elem_funs,pair_eq_iff}.
    For all $x\in X$ we have $h_{21}(x) = j(x)$ by assumption.
    For all $x\in X$ we have $(h_{12}\circ h_{21})(x) = h_{12}(h_{21}(x))$
        by \cref{circ_apply,continuous,funs,funs_is_function,funs_ran,compactifies,subspace_of,subseteq}.
    For all $x\in X$ we have $h_{12}(x) = j(x)$ by assumption.
    For all $x\in X$ we have $(h_{12}\circ h_{21})(x) = j(x)$ by \cref{pair_eq_iff,id_apply,id_elem_funs}.
    $i_2$ is continuous from $\closure{X}{Y_2}{S_2}$ to $Y_2$
        with respect to $\subspaceTopology{S_2}{\closure{X}{Y_2}{S_2}}$ and $S_2$
        and $h_{12}\circ h_{21}$ is continuous from $\closure{X}{Y_2}{S_2}$ to $Y_2$
        with respect to $\subspaceTopology{S_2}{\closure{X}{Y_2}{S_2}}$ and $S_2$
        by \cref{compactifies,pair_eq_iff}.
    Show that $h_{12}\circ h_{21} = i_2$.
    \begin{subproof}
        Follows by \cref{compactifies,subspace_of,extension_to_closure_unique,subseteq}.
    \end{subproof}

    For all $y_1,y_2$ such that $y_1\in\dom{h_{12}}$ and $y_2\in\dom{h_{12}}$ and $h_{12}(y_1) = h_{12}(y_2)$
        we have $y_1 = y_2$
        by \cref{circ_apply,continuous,funs,funs_is_function,funs_ran,pair_eq_iff,id_apply,id_elem_funs,subseteq}.
    $h_{12}$ is injective by \cref{injective_function}.

    For all $y$ such that $y\in Y_2$ we have $y\in\ran{h_{12}}$
        by \cref{continuous,funs_type_apply,continuous,funs,subseteq,circ_apply,pair_eq_iff,id_apply,id_elem_funs,function_apply_elim,funs_is_function,funs_ran,ran_intro}.
    $\ran{h_{12}} = Y_2$ by \cref{subseteq,funs_ran,continuous,subseteq_antisymmetric}.
    $h_{12}\in\surj{Y_1}{Y_2}$ by \cref{funs_surj_iff,continuous,funs}.

    $h_{12}$ is a homeomorphism between $Y_1$ and $Y_2$ with respect to $S_1$ and $S_2$
        by \cref{bijections,injective_function,compact_hausdorff_bijection_homeomorphism}.

    For all $x\in X$ we have $h_{12}(x) = x$ by \cref{id_apply,id_elem_funs,pair_eq_iff}.

    $Y_1$ is equivalent to $Y_2$ over $X$ with respect to $S_1$ and $S_2$ and $T$
        by \cref{compactification_equivalence_over_x}.
\end{proof}

% from §38 Item 7: Stone-Cech compactification
\begin{definition}\label{stone_cech_compactification}
    $Y$ is a Stone-Cech compactification of $X$ with respect to $S$ and $T$ iff
        $Y$ has the compact Hausdorff extension property over $X$ with respect to $S$ and $T$.
\end{definition}
\section{§39 Local Finiteness}

% from §39 Item 1: locally finite collection
\begin{definition}\label{locally_finite_collection}
    $A$ is locally finite in $X$ with respect to $T$ iff
        $A\subseteq\pow{X}$ and
        for all $x\in X$ there exist $U, F$ such that
            $U$ is a neighborhood of $x$ in $X$ with respect to $T$ and
            $F\subseteq A$ and
            $F$ is finite and
            for all $B\in A$ if $U$ intersects $B$ then $B\in F$.
\end{definition}

% from §39 Item 2: subcollections of locally finite collections
\begin{lemma}\label{locally_finite_subcollection}
    Assume $A$ is locally finite in $X$ with respect to $T$.
    Assume $B\subseteq A$.
    Then $B$ is locally finite in $X$ with respect to $T$.
\end{lemma}
\begin{proof}
    Follows by \cref{locally_finite_collection,subseteq,inter_lower_right,inter_lower_left,subseteq_finite_is_finite,inter_intro,subseteq_transitive}.
\end{proof}

% from §39 helper definition 1: collection of closures of a collection
\begin{definition}\label{closure_collection_in}
    $B$ is the closure collection of $A$ in $X$ with respect to $T$ iff
        $B = \{C\in\pow{X} \mid \exists D\in A. C = \closure{D}{X}{T}\}$.
\end{definition}

% from §39 Item 3: closures of a locally finite collection are locally finite
\begin{lemma}\label{locally_finite_closure_collection}
    Assume $A$ is locally finite in $X$ with respect to $T$.
    Assume $B$ is the closure collection of $A$ in $X$ with respect to $T$.
    Then $B$ is locally finite in $X$ with respect to $T$.
\end{lemma}
\begin{proof}
    Show that for all $x\in X$ there exist $U, G$ such that
        $U$ is a neighborhood of $x$ in $X$ with respect to $T$ and
        $G\subseteq B$ and
        $G$ is finite and
        for all $C\in B$ if $U$ intersects $C$ then $C\in G$.
    \begin{subproof}
        Fix $x\in X$.
        Take $U, F$ such that
            $U$ is a neighborhood of $x$ in $X$ with respect to $T$ and
            $F\subseteq A$ and
            $F$ is finite and
            for all $D\in A$ if $U$ intersects $D$ then $D\in F$
            by \cref{locally_finite_collection}.
        Let $h = \{(D,\closure{D}{X}{T}) \mid D\in F\}$.
        $h$ is a function on $F$ by \cref{relation,pair_eq_iff,rightunique,dom_iff,subseteq,subseteq_antisymmetric}.
        Let $H = \img{h}{F}$.
        Let $G = B\inter H$.
        $G\subseteq B$ and $G$ is finite
            by \cref{finite_image_function,inter_lower_left,inter_lower_right,subseteq_finite_is_finite}.
        For all $C\in B$ if $U$ intersects $C$ then $C\in G$
            by \cref{closure_collection_in,intersects,nonempty_inhabited,inter,neighborhood,open_in,topology_on,subseteq,pow_iff,elem_subseteq,powerset_elim,closure_iff_intersects_open,locally_finite_collection,img_iff,function_apply_iff,pair_eq_iff,inter_intro}.
    \end{subproof}
    $B$ is locally finite in $X$ with respect to $T$ by \cref{closure_collection_in,subseteq,locally_finite_collection}.
\end{proof}

% from §39 Item 4: closure of the union of a locally finite collection
\begin{lemma}\label{closure_of_union_locally_finite}
    Assume $A$ is locally finite in $X$ with respect to $T$.
    Assume $B$ is the closure collection of $A$ in $X$ with respect to $T$.
    Then $\closure{\unions{A}}{X}{T} = \unions{B}$.
\end{lemma}
\begin{proof}
    $A\subseteq\pow{X}$ by \cref{locally_finite_collection}.
    $B\subseteq\pow{X}$ by \cref{closure_collection_in,subseteq}.
    $\unions{A}\subseteq X$ by \cref{unions_subseteq_of_powerset_is_subseteq}.
    Show that for all $x$ such that $x\in\unions{B}$ we have $x\in\closure{\unions{A}}{X}{T}$.
        \begin{subproof}
            Fix $x$ such that $x\in\unions{B}$.
            Take $C\in B$ such that $x\in C$ by \cref{unions_iff}.
            Take $D\in A$ such that $C = \closure{D}{X}{T}$ by \cref{closure_collection_in}.
            $x\in X$ and $x\in\closure{D}{X}{T}$ by \cref{closure}.
            $T$ is a topology on $X$ by \cref{locally_finite_collection,neighborhood,open_in}.
            $D\subseteq X$ by \cref{subseteq,powerset_elim}.
            For all $U$ such that $U$ is open in $X$ with respect to $T$ and $x\in U$ we have $U$ intersects $D$
                by \cref{open_in,closure_iff_intersects_open}.
        For all $U$ such that $U$ is open in $X$ with respect to $T$ and $x\in U$
            we have $U$ intersects $\unions{A}$ by \cref{intersects,nonempty_inhabited,inter,unions_iff,inter_intro,emptyset}.
        $x\in\closure{\unions{A}}{X}{T}$ by \cref{closure_iff_intersects_open}.
    \end{subproof}
    $\unions{B}\subseteq\closure{\unions{A}}{X}{T}$ by \cref{subseteq}.
    Show that for all $x$ such that $x\in\closure{\unions{A}}{X}{T}$ we have $x\in\unions{B}$.
    \begin{subproof}
        Fix $x$ such that $x\in\closure{\unions{A}}{X}{T}$.
        $x\in X$ by \cref{closure}.
        Take $U, F$ such that
            $U$ is a neighborhood of $x$ in $X$ with respect to $T$ and
            $F\subseteq A$ and
            $F$ is finite and
            for all $D\in A$ if $U$ intersects $D$ then $D\in F$
            by \cref{locally_finite_collection}.
        $T$ is a topology on $X$ by \cref{neighborhood,open_in}.
        Let $h = \{(D,\closure{D}{X}{T}) \mid D\in F\}$.
        $h$ is a function on $F$ by \cref{relation,pair_eq_iff,rightunique,dom_iff,subseteq,subseteq_antisymmetric}.
        Let $H = \img{h}{F}$.
        $H\subseteq B$
            by \cref{img_iff,function_apply_iff,subseteq,pair_eq_iff,closure,powerset_intro,closure_collection_in}.
        $H\subseteq\pow{X}$ by \cref{subseteq_transitive}.
        For all $C$ such that $C\in H$ we have $C$ is closed in $X$ with respect to $T$
            by \cref{img_iff,function_apply_iff,subseteq,powerset_elim,pair_eq_iff,closure_closed_in}.
        $\unions{H}$ is closed in $X$ with respect to $T$ by \cref{finite_image_function,closed_finite_unions}.
        Suppose not.
        $x\notin\unions{B}$.
        $x\notin\unions{H}$ by \cref{unions_iff,subseteq,not_in_subseteq}.
        $U\inter (X\setminus\unions{H})$ is a neighborhood of $x$ in $X$ with respect to $T$ by
            \cref{neighborhood,open_in,closed_in,topology_on,setminus_intro,inter_intro}.
        Show that $U\inter (X\setminus\unions{H})$ does not intersect $\unions{A}$.
        \begin{subproof}
            Suppose not.
            $(U\inter (X\setminus\unions{H}))\inter\unions{A}\neq\emptyset$ by \cref{intersects}.
            Take $y$ such that $y\in (U\inter (X\setminus\unions{H}))\inter\unions{A}$ by \cref{nonempty_inhabited}.
            $y\in U\inter (X\setminus\unions{H})$ and $y\in\unions{A}$ by \cref{inter}.
            Take $D\in A$ such that $y\in D$ by \cref{unions_iff}.
            $y\in U$ and $y\in X\setminus\unions{H}$ by \cref{inter}.
            $y\in U\inter D$ by \cref{inter_intro}.
            $U$ intersects $D$ by \cref{inter_intro,emptyset,intersects}.
            $D\in F$ by \cref{locally_finite_collection}.
            $y\in\closure{D}{X}{T}$ by \cref{subseteq,powerset_elim,subseteq_closure,elem_subseteq}.
            $\closure{D}{X}{T}\in H$ by \cref{img_iff,function_apply_iff,pair_eq_iff}.
            $y\in\unions{H}$ by \cref{unions_iff}.
            Contradiction by \cref{setminus_elim_right}.
        \end{subproof}
        Contradiction by \cref{neighborhood,closure,closure_iff_intersects_open}.
    \end{subproof}
    $\closure{\unions{A}}{X}{T}\subseteq\unions{B}$ by \cref{subseteq}.
    $\closure{\unions{A}}{X}{T} = \unions{B}$ by \cref{subseteq_antisymmetric}.
\end{proof}

% from §39 helper definition 2: indexed family in a space
\begin{definition}\label{indexed_family_on_in}
    $A$ is an indexed family on $J$ in $X$ iff
        $A$ is a function on $J$ and
        for all $\alpha\in J$ we have $\apply{A}{\alpha}\subseteq X$.
\end{definition}

% from §39 Item 5: locally finite indexed family
\begin{definition}\label{locally_finite_indexed_family}
    $A$ is a locally finite indexed family on $J$ in $X$ with respect to $T$ iff
        $A$ is an indexed family on $J$ in $X$ and
        for all $x\in X$ there exist $U, F$ such that
            $U$ is a neighborhood of $x$ in $X$ with respect to $T$ and
            $F\subseteq J$ and
            $F$ is finite and
            for all $\alpha\in J$ if $U$ intersects $\apply{A}{\alpha}$ then $\alpha\in F$.
\end{definition}

% from §39 helper lemma 2: indexed local finiteness implies local finiteness of the image collection
\begin{lemma}\label{locally_finite_indexed_family_implies_locally_finite_collection}
    Assume $A$ is a function on $J$.
    Suppose for all $\alpha\in J$ we have $\apply{A}{\alpha}\subseteq X$.
    Suppose $A$ is a locally finite indexed family on $J$ in $X$ with respect to $T$.
    Then $\img{A}{J}$ is locally finite in $X$ with respect to $T$.
\end{lemma}
\begin{proof}
    $\img{A}{J}\subseteq\pow{X}$ by \cref{img_iff,function_apply_iff,powerset_intro,subseteq}.
    Show that for all $x\in X$ there exist $U, G$ such that
        $U$ is a neighborhood of $x$ in $X$ with respect to $T$ and
        $G\subseteq\img{A}{J}$ and
        $G$ is finite and
        for all $C\in\img{A}{J}$ if $U$ intersects $C$ then $C\in G$.
    \begin{subproof}
        Fix $x\in X$.
        Take $U, F$ such that
            $U$ is a neighborhood of $x$ in $X$ with respect to $T$ and
            $F\subseteq J$ and
            $F$ is finite and
            for all $\alpha\in J$ if $U$ intersects $\apply{A}{\alpha}$ then $\alpha\in F$
            by \cref{locally_finite_indexed_family}.
        Let $g = \restrl{A}{F}$.
        $g$ is a function and $\dom{g} = F$
            by \cref{restrl_of_function_is_function,dom_of_restrl_eq_inter,funs,subseteq,inter_absorb_supseteq_left}.
        Let $G = \img{g}{F}$.
        $G\subseteq\img{A}{J}$ and $G$ is finite
            by \cref{img_iff,restrl_iff,function_apply_iff,subseteq,finite_image_function}.
        For all $C\in\img{A}{J}$ if $U$ intersects $C$ then $C\in G$
            by \cref{img_iff,restrl_iff,function_apply_iff,locally_finite_indexed_family}.
    \end{subproof}
    $\img{A}{J}$ is locally finite in $X$ with respect to $T$ by \cref{locally_finite_collection}.
\end{proof}

% from §39 helper lemma 3: indexed local finiteness implies finite multiplicity of each nonempty value
\begin{lemma}\label{locally_finite_indexed_family_implies_finite_nonempty_fibres}
    Assume $A$ is a function on $J$.
    Suppose for all $\alpha\in J$ we have $\apply{A}{\alpha}\subseteq X$.
    Suppose $A$ is a locally finite indexed family on $J$ in $X$ with respect to $T$.
    Then for all $B$ such that $B\subseteq X$ and $B\neq\emptyset$
        there exists $F$ such that
            $F\subseteq J$ and
            $F$ is finite and
            for all $\alpha\in J$ if $\apply{A}{\alpha} = B$ then $\alpha\in F$.
\end{lemma}
\begin{proof}
    Fix $B$ such that $B\subseteq X$ and $B\neq\emptyset$.
    Take $x$ such that $x\in B$ by \cref{nonempty_inhabited}.
    Take $U, F$ such that
        $U$ is a neighborhood of $x$ in $X$ with respect to $T$ and
        $F\subseteq J$ and
        $F$ is finite and
        for all $\alpha\in J$ if $U$ intersects $\apply{A}{\alpha}$ then $\alpha\in F$
        by \cref{elem_subseteq,locally_finite_indexed_family}.
    For all $\alpha\in J$ if $\apply{A}{\alpha} = B$ then $\alpha\in F$
        by \cref{elem_subseteq,neighborhood,inter_intro,emptyset,intersects,locally_finite_indexed_family}.
\end{proof}

% from §39 helper lemma 4: collection formulation implies indexed local finiteness
\begin{lemma}\label{locally_finite_collection_and_finite_nonempty_fibres_imply_indexed_family}
    Assume $A$ is a function on $J$.
    Suppose for all $\alpha\in J$ we have $\apply{A}{\alpha}\subseteq X$.
    Suppose $\img{A}{J}$ is locally finite in $X$ with respect to $T$.
    Suppose for all $B$ such that $B\subseteq X$ and $B\neq\emptyset$
        there exists $F$ such that
            $F\subseteq J$ and
            $F$ is finite and
            for all $\alpha\in J$ if $\apply{A}{\alpha} = B$ then $\alpha\in F$.
    Then $A$ is a locally finite indexed family on $J$ in $X$ with respect to $T$.
\end{lemma}
\begin{proof}
    Show that for all $x\in X$ there exist $U, F$ such that
        $U$ is a neighborhood of $x$ in $X$ with respect to $T$ and
        $F\subseteq J$ and
        $F$ is finite and
        for all $\alpha\in J$ if $U$ intersects $\apply{A}{\alpha}$ then $\alpha\in F$.
    \begin{subproof}
        Fix $x\in X$.
        Take $U, G$ such that
            $U$ is a neighborhood of $x$ in $X$ with respect to $T$ and
            $G\subseteq\img{A}{J}$ and
            $G$ is finite and
            for all $C\in\img{A}{J}$ if $U$ intersects $C$ then $C\in G$
            by \cref{locally_finite_collection}.
        Let $G_0 = G\setminus\{\emptyset\}$.
        $G_0\subseteq G$ and $G_0$ is finite
            by \cref{setminus_subseteq,subseteq_finite_is_finite}.
        Let $R = \{ w\in G_0\times\pow{J} \mid
            \text{$\snd{w}$ is finite and
            for all $\alpha\in J$ if $\apply{A}{\alpha} = \fst{w}$ then $\alpha\in\snd{w}$} \}$.
        $R$ is a relation by \cref{relation,subseteq,times_elem_is_tuple}.
        For all $B\in G_0$ there exists $F$ such that $B\mathrel{R} F$
            by \cref{setminus_elim_left,subseteq,setminus_elim_right,singleton_iff,img_iff,function_apply_iff,times_tuple_intro,powerset_intro,fst_eq,snd_eq}.
        Take $c$ such that $c$ is a function on $G_0$ and for all $B\in G_0$ we have $B\mathrel{R}\apply{c}{B}$
            by \cref{choice_function}.
        Let $M = \img{c}{G_0}$.
        For all $B\in G_0$ we have $(B,\apply{c}{B})\in R$ by \cref{choice_function,relation}.
        For all $B\in G_0$ we have $(B,\apply{c}{B})\in G_0\times\pow{J}$ and $\apply{c}{B}$ is finite.
        For all $B\in G_0$ we have $\apply{c}{B}\subseteq J$ and $\apply{c}{B}$ is finite
            by \cref{times_tuple_elim,powerset_elim}.
        For all $F_0$ such that $F_0\in M$ we have $F_0\subseteq J$ and $F_0$ is finite
            by \cref{img_iff,function_apply_iff}.
        $M\subseteq\pow{J}$ by \cref{subseteq,powerset_intro}.
        Let $F = \unions{M}$.
        $F$ is finite and $F\subseteq J$
            by \cref{finite_image_function,finite_union_of_finite_family,unions_subseteq_of_powerset_is_subseteq}.
        For all $\alpha\in J$ if $U$ intersects $\apply{A}{\alpha}$ then $\alpha\in F$
            by \cref{img_iff,function_apply_iff,locally_finite_collection,intersects,inter_emptyset,setminus_intro,singleton_iff,choice_function,fst_eq,snd_eq,unions_iff}.
    \end{subproof}
    $A$ is a locally finite indexed family on $J$ in $X$ with respect to $T$
        by \cref{indexed_family_on_in,locally_finite_indexed_family}.
\end{proof}

% from §39 Item 6: relation between indexed and collection formulations of local finiteness
\begin{lemma}\label{locally_finite_indexed_family_iff_collection}
    Assume $A$ is a function on $J$.
    Suppose for all $\alpha\in J$ we have $\apply{A}{\alpha}\subseteq X$.
    Then $A$ is a locally finite indexed family on $J$ in $X$ with respect to $T$ iff
        $\img{A}{J}$ is locally finite in $X$ with respect to $T$ and
        for all $B$ such that $B\subseteq X$ and $B\neq\emptyset$
        there exists $F$ such that
            $F\subseteq J$ and
            $F$ is finite and
            for all $\alpha\in J$ if $\apply{A}{\alpha} = B$ then $\alpha\in F$.
\end{lemma}
\begin{proof}
    Follows by \cref{locally_finite_indexed_family_implies_locally_finite_collection,locally_finite_indexed_family_implies_finite_nonempty_fibres,locally_finite_collection_and_finite_nonempty_fibres_imply_indexed_family}.
\end{proof}

% from §39 Item 7: countably locally finite collection
\begin{definition}\label{countably_locally_finite_collection}
    $A$ is countably locally finite in $X$ with respect to $T$ iff
        $A\subseteq\pow{X}$ and
        there exists $B$ such that
            $B$ is a function on $\naturalsPlus$ and
            $A = \unions{\img{B}{\naturalsPlus}}$ and
            for all $n\in\naturalsPlus$ we have $\apply{B}{n}$ is locally finite in $X$ with respect to $T$.
\end{definition}

% from §39 Item 8: refinement
\begin{definition}\label{refinement_on}
    $B$ is a refinement of $A$ on $X$ iff
        $A\subseteq\pow{X}$ and
        $B\subseteq\pow{X}$ and
        for all $C\in B$ there exists $D\in A$ such that $C\subseteq D$.
\end{definition}

% from §39 Item 9: open refinement
\begin{definition}\label{open_refinement_on}
    $B$ is an open refinement of $A$ on $X$ with respect to $T$ iff
        $B$ is a refinement of $A$ on $X$ and
        $B$ is a covering of $X$ and
        for all $U\in B$ we have $U$ is open in $X$ with respect to $T$.
\end{definition}

% from §39 Item 10: closed refinement
\begin{definition}\label{closed_refinement_on}
    $B$ is a closed refinement of $A$ on $X$ with respect to $T$ iff
        $B$ is a refinement of $A$ on $X$ and
        for all $U\in B$ we have $U$ is closed in $X$ with respect to $T$.
\end{definition}

% from §39 helper lemma 5: well-ordering theorem for the cover
\begin{lemma}\label{well_ordering_theorem_local_finiteness}
    There exists $R$ such that $R$ is a well order relation on $A$.
\end{lemma}
\begin{proof}
    Omitted. % do not touch
\end{proof}

% from §39 helper definition 5: metric refinement sequence openness
\begin{definition}\label{metric_refinement_sequence_open}
    $b$ is an open sequence on $X$ with respect to $T$ iff
        for all $n\in\naturalsPlus$ for all $C\in\apply{b}{n}$ we have $C$ is open in $X$ with respect to $T$.
\end{definition}

% from §39 helper definition 6: metric refinement sequence cover property
\begin{definition}\label{metric_refinement_sequence_point_covering}
    $b$ is a covering sequence of $X$ iff
        for all $x\in X$ there exist $n, C$ such that
            $n\in\naturalsPlus$ and
            $C\in\apply{b}{n}$ and
            $x\in C$.
\end{definition}

% from §39 helper definition 7: metric refinement sequence refinement property
\begin{definition}\label{metric_refinement_sequence_refines}
    $b$ is a refinement sequence of $A$ iff
        for all $n\in\naturalsPlus$ for all $C\in\apply{b}{n}$ there exists $U\in A$ such that $C\subseteq U$.
\end{definition}

% from §39 helper definition 3: metric refinement sequence
\begin{definition}\label{metric_refinement_sequence}
    $b$ is a metric refinement sequence of $A$ on $X$ with respect to $T$ iff
        $b$ is a function on $\naturalsPlus$ and
        for all $n\in\naturalsPlus$ we have $\apply{b}{n}$ is locally finite in $X$ with respect to $T$ and
        $b$ is an open sequence on $X$ with respect to $T$ and
        $b$ is a covering sequence of $X$ and
        $b$ is a refinement sequence of $A$.
\end{definition}

% from §39 helper lemma 6: existence of a metric refinement sequence
% from §39 helper lemma 8: doubling the reciprocal of twice a positive natural
\begin{lemma}\label{double_half_reciprocal}
    Suppose $n\in\naturalsPlus$.
    Then $1 / ((1 + 1) \rmul n) + 1 / ((1 + 1) \rmul n) = 1 / n$.
\end{lemma}
\begin{proof}
    $1\in\naturalsPlus$ by \cref{real_one_in_naturals}.
    $1 + 1\in\naturalsPlus$ by \cref{naturalsplus_add_closed}.
    Let $m = (1 + 1) \rmul n$.
    $m\in\naturalsPlus$ by \cref{naturalsplus_mul_closed}.
    $m\in\reals$ and $0 < m$ and $m\neq 0$
        by \cref{real_naturals_subset_reals,subseteq,naturalsplus_positive_real,real_pos_nonzero}.
    $n\in\reals$ and $0 < n$ and $n\neq 0$
        by \cref{real_naturals_subset_reals,subseteq,naturalsplus_positive_real,real_pos_nonzero}.
    $1 + 1\in\reals$ by \cref{real_one_in_naturals,real_naturals_subset_reals,subseteq,real_add_closed}.
    $1 / m + 1 / m = (1 + 1) / m$ by
        \cref{real_div_add_same_denom,real_mul_ident,real_lt_subst_left,real_lt_subst_right}.
    Let $a = (1 + 1) / m$.
    $a\in\reals$ by \cref{real_div_naturalsplus_real}.
    $\inv{m}\in\reals$ by \cref{real_mul_inverse}.
    $a = (1 + 1) \rmul \inv{m}$ by \cref{real_div}.
    $a \rmul n = ((1 + 1) \rmul n) \rmul \inv{m}$ by \cref{subseteq,real_mul_assoc,real_mul_comm}.
    $a \rmul n = m \rmul \inv{m}$ by \cref{real_mul_eq_subst}.
    $a \rmul n = 1$ by \cref{real_mul_inverse}.
    $a = 1 / n$ by \cref{real_mul_comm,real_mul_inverse_unique,real_div,real_mul_ident}.
    $1 / ((1 + 1) \rmul n) + 1 / ((1 + 1) \rmul n) = 1 / n$
        by \cref{real_lt_subst_right}.
\end{proof}

% from §39 helper lemma 9: nearby core points belong to the same cover element
\begin{lemma}\label{nearby_core_points_same_cover_element}
    Assume $R$ is a well order relation on $A$.
    Suppose $d$ is a metric on $X$.
    Suppose $U\in A$ and $V\in A$ and $n\in\naturalsPlus$.
    Suppose $x\in X$ and $y\in X$.
    Suppose $\metricBall{d}{X}{x}{1 / n} \subseteq U$.
    Suppose $\metricBall{d}{X}{y}{1 / n} \subseteq V$.
    Suppose for all $W\in A$ if $W\mathrel{R}U$ then $x\notin W$.
    Suppose for all $W\in A$ if $W\mathrel{R}V$ then $y\notin W$.
    Suppose $d((x,y)) < 1 / n$.
    Then $U = V$.
\end{lemma}
\begin{proof}
    Suppose not.
    We have $U\mathrel{R}V$ or $V\mathrel{R}U$ by \cref{well_order_relation,simple_order_relation,connex}.
    \begin{byCase}
    \caseOf{$U\mathrel{R}V$.}
        We have $y\notin U$.
        Contradiction by \cref{metric_ball,subseteq}.
    \caseOf{$V\mathrel{R}U$.}
        We have $x\notin V$.
        Contradiction by \cref{metric_ball,metric_symmetric,metric_on,real_lt_subst_left,times_tuple_intro,subseteq}.
    \end{byCase}
\end{proof}

% from §39 helper lemma 10: quarter reciprocal is bounded by the original reciprocal
\begin{lemma}\label{quarter_reciprocal_leq_reciprocal}
    Suppose $n\in\naturalsPlus$.
    Then $1 / ((1 + 1) \rmul ((1 + 1) \rmul n)) \leq 1 / n$.
\end{lemma}
\begin{proof}
    $1\in\naturalsPlus$ by \cref{real_one_in_naturals}.
    $1\in\reals$ by \cref{real_mul_ident}.
    We have $1 + 1\in\naturalsPlus$ and $1 + 1\in\reals$
        by \cref{real_one_in_naturals,naturalsplus_add_closed,real_naturals_subset_reals,subseteq}.
    Let $m = (1 + 1) \rmul n$.
    $m\in\naturalsPlus$ by \cref{naturalsplus_mul_closed}.
    Let $k = (1 + 1) \rmul m$.
    $k\in\naturalsPlus$ by \cref{naturalsplus_mul_closed}.
    We have $n\in\reals$ and $0 < n$ by \cref{real_naturals_subset_reals,subseteq,naturalsplus_positive_real}.
    $n < n + n$ by \cref{real_add_ident,real_lt_add}.
    Thus $m = n + n$ by \cref{real_mul_comm,real_distrib,real_mul_ident,real_add_eq_subst,real_mul_eq_subst}.
    Therefore $n \leq m$ and $1 / m \leq 1 / n$
        by \cref{real_lt_subst_right,real_lt_imp_leq,naturalsplus_reciprocal_antitone}.
    We have $m\in\reals$ and $0 < m$ by \cref{naturalsplus_mul_closed,real_naturals_subset_reals,subseteq,naturalsplus_positive_real}.
    $m < m + m$ by \cref{real_add_ident,real_lt_add}.
    Thus $k = m + m$ by \cref{real_mul_comm,real_distrib,real_mul_ident,real_add_eq_subst,real_mul_eq_subst}.
    Therefore $m \leq k$ and $1 / k \leq 1 / m$
        by \cref{real_lt_subst_right,real_lt_imp_leq,naturalsplus_reciprocal_antitone}.
    We have $1 / k\in\reals$ and $1 / m\in\reals$ and $1 / n\in\reals$ by \cref{real_div_naturalsplus_real}.
    Hence $1 / k \leq 1 / n$ by \cref{naturalsplus_mul_closed,real_div_naturalsplus_real,real_leq_trans}.
    Therefore $1 / ((1 + 1) \rmul ((1 + 1) \rmul n)) \leq 1 / n$
        by \cref{real_leq_trans,real_mul_eq_subst,subseteq,real_leq_subst_left_local}.
\end{proof}

% from §39 helper definition 4: nth metric refinement piece
\begin{definition}\label{metric_refinement_piece}
    $C$ is an $n$-metric refinement piece of $U$ from $A$ on $X$ with respect to $R$ and $d$ iff
        $n\in\naturalsPlus$ and
        $U\in A$ and
        $C = \{y\in X \mid \text{there exists $x\in X$ such that
            $d((x,y)) < 1 / ((1 + 1) \rmul ((1 + 1) \rmul n))$ and
            $\metricBall{d}{X}{x}{1 / n} \subseteq U$ and
            for all $W\in A$ if $W\mathrel{R}U$ then $x\notin W$}\}$.
\end{definition}

% from §39 helper lemma 6: existence of a metric refinement sequence
\begin{lemma}\label{metric_refinement_sequence_exists}
    Assume $R$ is a well order relation on $A$.
    Suppose $d$ is a metric on $X$.
    Suppose $T = \metricTopology{d}{X}$.
    Suppose $A$ is an open covering of $X$ with respect to $T$.
    Then there exists $b$ such that $b$ is a metric refinement sequence of $A$ on $X$ with respect to $T$.
\end{lemma}
\begin{proof}
    $R$ is a simple order relation on $A$ by \cref{well_order_relation}.
    Let $S = \{w\in\naturalsPlus\times\pow{\pow{X}} \mid \text{there exist $n, B$ such that
        $w = (n,B)$ and
        $B = \{C\in\pow{X} \mid \text{there exists $U\in A$ such that
            $C$ is an $n$-metric refinement piece of $U$ from $A$ on $X$ with respect to $R$ and $d$}\}$}\}$.
    $S$ is a relation by \cref{relation,times_elem_is_tuple}.
    Show that for all $n\in\naturalsPlus$ there exists $B$ such that $n\mathrel{S}B$.
    \begin{subproof}
        Fix $n\in\naturalsPlus$.
        Let $B = \{C\in\pow{X} \mid \text{there exists $U\in A$ such that
            $C$ is an $n$-metric refinement piece of $U$ from $A$ on $X$ with respect to $R$ and $d$}\}$.
        We have $(n,B)\in S$ by \cref{subseteq,powerset_intro,times_tuple_intro}.
    \end{subproof}
    Take $b$ such that $b$ is a function on $\naturalsPlus$ and for all $n\in\naturalsPlus$ we have $n\mathrel{S}\apply{b}{n}$
        by \cref{choice_function}.
    Show that for all $n\in\naturalsPlus$ we have $\apply{b}{n}$ is locally finite in $X$ with respect to $T$.
    \begin{subproof}
        Fix $n\in\naturalsPlus$.
        Take $n_0, B_0$ such that
            $(n,\apply{b}{n}) = (n_0,B_0)$ and
            $B_0 = \{C\in\pow{X} \mid \text{there exists $U\in A$ such that
                $C$ is an $n_0$-metric refinement piece of $U$ from $A$ on $X$ with respect to $R$ and $d$}\}$
            by \cref{choice_function,relation,subseteq}.
        Hence $n = n_0$ and $\apply{b}{n} = B_0$ by \cref{pair_eq_iff}.
        We have $\apply{b}{n}\subseteq\pow{X}$ by \cref{subseteq}.
        Show that for all $x\in X$ there exist $U, F$ such that
            $U$ is a neighborhood of $x$ in $X$ with respect to $T$ and
            $F\subseteq\apply{b}{n}$ and
            $F$ is finite and
            for all $C\in\apply{b}{n}$ if $U$ intersects $C$ then $C\in F$.
        \begin{subproof}
            Fix $x\in X$.
            $1\in\naturalsPlus$ by \cref{real_one_in_naturals}.
            We have $1 + 1\in\naturalsPlus$ by \cref{naturalsplus_add_closed}.
            Let $m = (1 + 1) \rmul n$.
            $m\in\naturalsPlus$ by \cref{naturalsplus_mul_closed}.
            Let $k = (1 + 1) \rmul m$.
            $k\in\naturalsPlus$ by \cref{naturalsplus_mul_closed}.
            $1\in\reals$ by \cref{real_one_in_naturals,real_naturals_subset_reals,subseteq}.
            We have $1 / m\in\reals$ and $1 / n\in\reals$ by \cref{real_div_naturalsplus_real}.
            Let $q = 1 / k$.
            $q\in\reals$ and $0 < q$ by \cref{real_div_naturalsplus_real,positive_reciprocal_naturalsplus}.
            We have $q = 1 / ((1 + 1) \rmul ((1 + 1) \rmul n))$ by \cref{real_mul_eq_subst}.
            Let $U = \metricBall{d}{X}{x}{q}$.
            We have $U\in\metricBasis{d}{X}$ by \cref{metric_ball,subseteq,powerset_intro,metric_basis}.
            Hence $U$ is open in $X$ with respect to $T$ by
                \cref{metric_basis_is_basis,basis_for_topology,basis_elem_open_in_generated,metric_topology,basis_topology_is_topology,open_in}.
            We have $x\in U$ by \cref{metric_ball,metric_on,metric_zero_characterization,real_lt_subst_left}.
            Hence $U$ is a neighborhood of $x$ in $X$ with respect to $T$ by \cref{neighborhood}.
            Show that for all $C, D$ such that $C\in\apply{b}{n}$ and $D\in\apply{b}{n}$ and
                $U$ intersects $C$ and $U$ intersects $D$ we have $C = D$.
            \begin{subproof}
                Fix $C$.
                Fix $D$.
                Assume $C\in\apply{b}{n}$ and $D\in\apply{b}{n}$ and
                    $U$ intersects $C$ and $U$ intersects $D$.
                Take $V\in A$ such that $C$ is an $n$-metric refinement piece of $V$ from $A$ on $X$ with respect to $R$ and $d$
                    by \cref{subseteq,metric_refinement_piece}.
                Take $W\in A$ such that $D$ is an $n$-metric refinement piece of $W$ from $A$ on $X$ with respect to $R$ and $d$
                    by \cref{subseteq,metric_refinement_piece}.
                Take $z$ such that $z\in U\inter C$ by \cref{intersects,nonempty_inhabited}.
                Take $t$ such that $t\in U\inter D$ by \cref{intersects,nonempty_inhabited}.
                We have $z\in U$ and $z\in C$ and $t\in U$ and $t\in D$ by \cref{inter}.
                Take $x_0\in X$ such that
                    $d((x_0,z)) < q$ and
                    $\metricBall{d}{X}{x_0}{1 / n} \subseteq V$ and
                    for all $Y\in A$ if $Y\mathrel{R}V$ then $x_0\notin Y$
                    by \cref{metric_refinement_piece}.
                Take $y_0\in X$ such that
                    $d((y_0,t)) < q$ and
                    $\metricBall{d}{X}{y_0}{1 / n} \subseteq W$ and
                    for all $Y\in A$ if $Y\mathrel{R}W$ then $y_0\notin Y$
                    by \cref{metric_refinement_piece}.
                $d\in\funs{X\times X}{\reals}$ by \cref{metric_on}.
                We have $z\in X$ and $t\in X$ by \cref{metric_ball,subseteq}.
                $d((x_0,z))\in\reals$ and $d((y_0,t))\in\reals$ by \cref{times_tuple_intro,funs_type_apply}.
                We have $d((x,z)) < q$ and $d((x,t)) < q$ by \cref{metric_ball}.
                Thus $d((z,x)) = d((x,z))$ and $d((t,x)) = d((x,t))$ and $d((t,y_0)) = d((y_0,t))$
                    by \cref{metric_symmetric,metric_on,real_lt_subst_left,times_tuple_intro}.
                We have $d((z,x))\in\reals$ and $d((t,x))\in\reals$ by \cref{times_tuple_intro,funs_type_apply}.
                $d((x_0,x))\in\reals$ and $d((x,y_0))\in\reals$ by \cref{times_tuple_intro,funs_type_apply}.
                $d((x,t))\in\reals$ and $d((t,y_0))\in\reals$ by \cref{times_tuple_intro,funs_type_apply}.
                Let $s_0 = d((x_0,z)) + d((z,x))$.
                $d((x_0,x)) \leq s_0$ by \cref{metric_on,metric_triangle_inequality}.
                We have $s_0\in\reals$ by \cref{real_add_closed}.
                $d((x_0,z)) + d((z,x)) < q + q$ by \cref{real_lt_add_two}.
                We have $q + q = 1 / m$ by \cref{double_half_reciprocal,real_mul_eq_subst,real_lt_subst_left}.
                Therefore $s_0 < 1 / m$ by \cref{real_lt_subst_right}.
                Hence $d((x_0,x)) < 1 / m$ by \cref{real_leq_lt_trans}.
                Let $s_1 = d((x,t)) + d((t,y_0))$.
                $d((x,y_0)) \leq s_1$ by \cref{metric_on,metric_triangle_inequality}.
                We have $s_1\in\reals$ by \cref{real_add_closed}.
                $d((x,t)) + d((t,y_0)) < q + q$ by \cref{real_lt_add_two}.
                Thus $s_1 < 1 / m$ by \cref{double_half_reciprocal,real_mul_eq_subst,real_lt_subst_left,real_lt_subst_right}.
                Therefore $d((x,y_0)) < 1 / m$ by \cref{real_leq_lt_trans}.
                We have $d((x_0,y_0))\in\reals$ by \cref{times_tuple_intro,funs_type_apply}.
                Let $s = d((x_0,x)) + d((x,y_0))$.
                $d((x_0,y_0)) \leq s$ by \cref{metric_on,metric_triangle_inequality}.
                $s\in\reals$ by \cref{real_add_closed}.
                $d((x_0,x)) + d((x,y_0)) < 1 / m + 1 / m$ by \cref{real_lt_add_two}.
                We have $1 / m + 1 / m = 1 / n$ by \cref{double_half_reciprocal}.
                Therefore $s < 1 / n$ by \cref{real_lt_subst_right}.
                Hence $d((x_0,y_0)) < 1 / n$ by \cref{real_leq_lt_trans}.
                Therefore $V = W$ by \cref{nearby_core_points_same_cover_element}.
                $C\subseteq D$ by \cref{metric_refinement_piece,subseteq}.
                $D\subseteq C$ by \cref{metric_refinement_piece,subseteq}.
                Therefore $C = D$ by \cref{subseteq_antisymmetric}.
            \end{subproof}
            \begin{byCase}
            \caseOf{$U$ intersects $\unions{\apply{b}{n}}$.}
                Take $y$ such that $y\in U\inter\unions{\apply{b}{n}}$ by \cref{intersects,nonempty_inhabited}.
                Take $C_0\in\apply{b}{n}$ such that $y\in C_0$ by \cref{inter,unions_iff}.
                Hence $U$ intersects $C_0$ by \cref{inter,inter_intro,emptyset,intersects}.
                Let $F = \{C_0\}$.
                $F\subseteq\apply{b}{n}$ by \cref{singleton_subset_intro}.
                $F$ is finite by \cref{singleton_subset_intro,pair_is_finite,upair_intro_left,subseteq_finite_is_finite}.
                For all $C\in\apply{b}{n}$ if $U$ intersects $C$ then $C\in F$
                    by \cref{subseteq,singleton_intro}.
            \caseOf{$U$ does not intersect $\unions{\apply{b}{n}}$.}
                Let $F = \emptyset$.
                $F\subseteq\apply{b}{n}$ by \cref{emptyset_subseteq}.
                $F$ is finite by \cref{emptyset_is_finite}.
                For all $C\in\apply{b}{n}$ we have $U$ does not intersect $C$
                    by \cref{intersects,nonempty_inhabited,inter,unions_iff,inter_intro,emptyset}.
            \end{byCase}
        \end{subproof}
        Hence $\apply{b}{n}$ is locally finite in $X$ with respect to $T$ by \cref{locally_finite_collection}.
    \end{subproof}
    Show that for all $n\in\naturalsPlus$ for all $C\in\apply{b}{n}$ we have $C$ is open in $X$ with respect to $T$.
    \begin{subproof}
        Fix $n\in\naturalsPlus$.
        Fix $C\in\apply{b}{n}$.
        Take $n_0, B_0$ such that
            $(n,\apply{b}{n}) = (n_0,B_0)$ and
            $B_0 = \{D\in\pow{X} \mid \text{there exists $U\in A$ such that
                $D$ is an $n_0$-metric refinement piece of $U$ from $A$ on $X$ with respect to $R$ and $d$}\}$
            by \cref{choice_function,relation,subseteq}.
        Hence $n = n_0$ and $\apply{b}{n} = B_0$ by \cref{pair_eq_iff}.
        Take $U\in A$ such that $C$ is an $n$-metric refinement piece of $U$ from $A$ on $X$ with respect to $R$ and $d$
            by \cref{subseteq,metric_refinement_piece}.
        Hence $C\subseteq X$ by \cref{metric_refinement_piece,subseteq}.
        $1\in\naturalsPlus$ by \cref{real_one_in_naturals}.
        We have $1 + 1\in\naturalsPlus$ by \cref{naturalsplus_add_closed}.
        Let $m = (1 + 1) \rmul n$.
        $m\in\naturalsPlus$ by \cref{naturalsplus_mul_closed}.
        Let $k = (1 + 1) \rmul m$.
        $k\in\naturalsPlus$ by \cref{naturalsplus_mul_closed}.
        $1\in\reals$ by \cref{real_one_in_naturals,real_naturals_subset_reals,subseteq}.
        Let $q = 1 / k$.
        $q\in\reals$ and $0 < q$ by \cref{real_div_naturalsplus_real,positive_reciprocal_naturalsplus}.
        We have $q = 1 / ((1 + 1) \rmul ((1 + 1) \rmul n))$ by \cref{real_mul_eq_subst}.
        Show that for all $y\in C$ there exists $\delta\in\reals$ such that $0 < \delta$ and $\metricBall{d}{X}{y}{\delta}\subseteq C$.
        \begin{subproof}
            Fix $y\in C$.
            Take $x\in X$ such that
                $d((x,y)) < q$ and
                $\metricBall{d}{X}{x}{1 / n} \subseteq U$ and
                for all $W\in A$ if $W\mathrel{R}U$ then $x\notin W$
                by \cref{metric_refinement_piece}.
            $d\in\funs{X\times X}{\reals}$ by \cref{metric_on}.
            We have $y\in X$ by \cref{metric_refinement_piece,subseteq}.
            Hence $d((x,y))\in\reals$ by \cref{times_tuple_intro,funs_type_apply}.
            Let $\delta = q - d((x,y))$.
            $0 < \delta$ by \cref{real_lt_imp_pos_diff}.
            Show that for all $z$ such that $z\in\metricBall{d}{X}{y}{\delta}$ we have $z\in C$.
            \begin{subproof}
                Fix $z$ such that $z\in\metricBall{d}{X}{y}{\delta}$.
                We have $z\in X$ and $d((y,z)) < \delta$ by \cref{metric_ball}.
                $d((y,z))\in\reals$ by \cref{times_tuple_intro,funs_type_apply}.
                $d((x,z))\in\reals$ by \cref{times_tuple_intro,funs_type_apply}.
                Thus $d((x,y)) + d((y,z)) < q$ by
                    \cref{real_sub,real_add_inverse,real_add_closed,real_lt_subst_right,real_lt_add,real_add_assoc,real_add_comm,real_add_ident}.
                We have $d((x,z)) \leq d((x,y)) + d((y,z))$ by \cref{metric_on,metric_triangle_inequality}.
                Hence $d((x,z)) < q$ by \cref{real_add_closed,real_leq_lt_trans}.
                Thus $z\in C$ by \cref{metric_refinement_piece}.
            \end{subproof}
            Therefore $\metricBall{d}{X}{y}{\delta}\subseteq C$ by \cref{subseteq}.
        \end{subproof}
        Therefore $C$ is open in $X$ with respect to $T$ by \cref{metric_open_iff}.
    \end{subproof}
    Show that for all $n\in\naturalsPlus$ for all $C\in\apply{b}{n}$ there exists $U\in A$ such that $C\subseteq U$.
    \begin{subproof}
        Fix $n\in\naturalsPlus$.
        Fix $C\in\apply{b}{n}$.
        Take $n_0, B_0$ such that
            $(n,\apply{b}{n}) = (n_0,B_0)$ and
            $B_0 = \{D\in\pow{X} \mid \text{there exists $U\in A$ such that
                $D$ is an $n_0$-metric refinement piece of $U$ from $A$ on $X$ with respect to $R$ and $d$}\}$
            by \cref{choice_function,relation,subseteq}.
        We have $n = n_0$ and $\apply{b}{n} = B_0$ by \cref{pair_eq_iff}.
        Take $U\in A$ such that $C$ is an $n$-metric refinement piece of $U$ from $A$ on $X$ with respect to $R$ and $d$
            by \cref{subseteq,metric_refinement_piece}.
        We have $C\subseteq X$ by \cref{metric_refinement_piece,subseteq}.
        Hence $C\subseteq U$
            by \cref{metric_refinement_piece,quarter_reciprocal_leq_reciprocal,metric_on,subseteq,times_tuple_intro,funs_type_apply,real_one_in_naturals,naturalsplus_add_closed,naturalsplus_mul_closed,real_naturals_subset_reals,real_div_naturalsplus_real,real_lt_leq_trans_local,metric_ball}.
    \end{subproof}
    Show that for all $x\in X$ there exist $n, C$ such that
        $n\in\naturalsPlus$ and
        $C\in\apply{b}{n}$ and
        $x\in C$.
    \begin{subproof}
        Fix $x\in X$.
        We have $x\in\unions{A}$ by \cref{open_covering,covering,subseteq}.
        Take $U_0\in A$ such that $x\in U_0$ by \cref{unions_iff}.
        Let $B = \{U\in A \mid x\in U\}$.
        $B\subseteq A$ by \cref{subseteq}.
        We have $U_0\in B$ by \cref{subseteq}.
        Therefore $B\neq\emptyset$ by \cref{emptyset}.
        Take $U$ such that $U$ is a smallest element of $B$ with respect to $R$ by \cref{well_order_relation}.
        We have $U\in B$ by \cref{smallest_element}.
        Hence $U\in A$ and $x\in U$ by \cref{subseteq}.
        Thus $U$ is open in $X$ with respect to $T$ by \cref{open_covering}.
        Therefore $U\subseteq X$ by
            \cref{open_in,metric_basis_is_basis,basis_topology_is_topology,metric_topology,basis_for_topology,topology_on,subseteq,pow_iff}.
        Take $\epsilon\in\reals$ such that $0 < \epsilon$ and $\metricBall{d}{X}{x}{\epsilon}\subseteq U$
            by \cref{metric_open_iff}.
        Take $n\in\naturalsPlus$ such that $0 < 1 / n$ and $1 / n < \epsilon$ by \cref{real_small_reciprocal}.
        We have $1 / n\in\reals$ by \cref{real_one_in_naturals,real_naturals_subset_reals,subseteq,real_div_naturalsplus_real}.
        Hence $\metricBall{d}{X}{x}{1 / n}\subseteq U$ by \cref{metric_ball_mono,subseteq_transitive}.
        For all $W\in A$ if $W\mathrel{R}U$ then $x\notin W$
            by \cref{subseteq,smallest_element,simple_order_relation,asymmetric,real_lt_subst_left}.
        Let $C = \{y\in X \mid \text{there exists $z\in X$ such that
            $d((z,y)) < 1 / ((1 + 1) \rmul ((1 + 1) \rmul n))$ and
            $\metricBall{d}{X}{z}{1 / n} \subseteq U$ and
            for all $W\in A$ if $W\mathrel{R}U$ then $z\notin W$}\}$.
        Therefore $C$ is an $n$-metric refinement piece of $U$ from $A$ on $X$ with respect to $R$ and $d$
            by \cref{metric_refinement_piece}.
        Take $n_0, B_0$ such that
            $(n,\apply{b}{n}) = (n_0,B_0)$ and
            $B_0 = \{D\in\pow{X} \mid \text{there exists $V\in A$ such that
                $D$ is an $n_0$-metric refinement piece of $V$ from $A$ on $X$ with respect to $R$ and $d$}\}$
            by \cref{choice_function,relation,subseteq}.
        Hence $n = n_0$ and $\apply{b}{n} = B_0$ by \cref{pair_eq_iff}.
        Thus $C\in\apply{b}{n}$ by \cref{metric_refinement_piece,subseteq,powerset_intro}.
        $1\in\naturalsPlus$ by \cref{real_one_in_naturals}.
        We have $1 + 1\in\naturalsPlus$ by \cref{naturalsplus_add_closed}.
        Let $k = (1 + 1) \rmul ((1 + 1) \rmul n)$.
        $k\in\naturalsPlus$ by \cref{naturalsplus_mul_closed}.
        $1\in\reals$ by \cref{real_one_in_naturals,real_naturals_subset_reals,subseteq}.
        We have $1 / k\in\reals$ and $0 < 1 / k$ by \cref{real_div_naturalsplus_real,positive_reciprocal_naturalsplus}.
        Therefore $d((x,x)) < 1 / k$
            by \cref{metric_on,metric_zero_characterization,times_tuple_intro,real_add_ident,real_lt_subst_left}.
        Hence $d((x,x)) < 1 / ((1 + 1) \rmul ((1 + 1) \rmul n))$ by \cref{real_mul_eq_subst,real_lt_subst_right}.
        Thus $x\in C$ by \cref{metric_refinement_piece}.
    \end{subproof}
    We have $b$ is an open sequence on $X$ with respect to $T$ by \cref{metric_refinement_sequence_open}.
    Hence $b$ is a refinement sequence of $A$ by \cref{metric_refinement_sequence_refines}.
    Thus $b$ is a covering sequence of $X$ by \cref{metric_refinement_sequence_point_covering}.
    Therefore $b$ is a metric refinement sequence of $A$ on $X$ with respect to $T$
        by \cref{metric_refinement_sequence}.
\end{proof}

% from §39 helper lemma 7: the induced union covers the space
\begin{lemma}\label{metric_refinement_sequence_covers}
    Assume $b$ is a metric refinement sequence of $A$ on $X$ with respect to $T$.
    Let $B = \unions{\img{b}{\naturalsPlus}}$.
    Then $B$ is a covering of $X$.
\end{lemma}
\begin{proof}
    $B$ is a covering of $X$
        by \cref{real_one_in_naturals,metric_refinement_sequence_point_covering,metric_refinement_sequence,img_iff,function_apply_iff,unions_iff,subseteq,covering}.
\end{proof}

% from §39 Item 11: countably locally finite open refinement in metrizable spaces
\begin{lemma}\label{metrizable_open_cover_countably_locally_finite_refinement}
    Assume $X$ is metrizable with respect to $T$.
    Suppose $A$ is an open covering of $X$ with respect to $T$.
    Then there exists $B$ such that
        $B$ is an open refinement of $A$ on $X$ with respect to $T$ and
        $B$ is countably locally finite in $X$ with respect to $T$.
\end{lemma}
\begin{proof}
    Take $d$ such that $d$ is a metric on $X$ and $T = \metricTopology{d}{X}$ by \cref{metrizable}.
    Take $R$ such that $R$ is a well order relation on $A$ by \cref{well_ordering_theorem_local_finiteness}.
    Take $b$ such that $b$ is a metric refinement sequence of $A$ on $X$ with respect to $T$
        by \cref{metric_refinement_sequence_exists}.
    Let $B = \unions{\img{b}{\naturalsPlus}}$.
    $A\subseteq\pow{X}$
        by \cref{open_covering,metric_basis_is_basis,basis_topology_is_topology,metric_topology,basis_for_topology,open_in,topology_on,subseteq}.
    $B\subseteq\pow{X}$
        by \cref{unions_iff,metric_refinement_sequence,img_iff,function_apply_iff,metric_refinement_sequence_refines,subseteq,powerset_elim,subseteq_transitive,powerset_intro}.
    Show that for all $C\in B$ there exists $U\in A$ such that $C\subseteq U$.
    \begin{subproof}
        Follows by \cref{unions_iff,metric_refinement_sequence,img_iff,function_apply_iff,metric_refinement_sequence_refines}.
    \end{subproof}
    Hence $B$ is a refinement of $A$ on $X$ by \cref{refinement_on}.
    $B$ is a covering of $X$ by \cref{metric_refinement_sequence_covers}.
    Show that for all $C\in B$ we have $C$ is open in $X$ with respect to $T$.
    \begin{subproof}
        Follows by \cref{unions_iff,metric_refinement_sequence,img_iff,function_apply_iff,metric_refinement_sequence_open}.
    \end{subproof}
    Hence $B$ is an open refinement of $A$ on $X$ with respect to $T$ by \cref{open_refinement_on}.
    We have $B$ is countably locally finite in $X$ with respect to $T$
        by \cref{metric_refinement_sequence,countably_locally_finite_collection}.
    Therefore there exists $C$ such that
        $C$ is an open refinement of $A$ on $X$ with respect to $T$ and
        $C$ is countably locally finite in $X$ with respect to $T$.
\end{proof}
\section{§40 The Nagata-Smirnov Metrization Theorem}

% from §40 Item 1: $G_\delta$ set
\begin{definition}\label{g_delta_set}
    $A$ is a Gdelta set in $X$ with respect to $T$ iff
        there exists $s$ such that
            $s$ is a sequence in $\pow{X}$ and
            $A = \inters{\img{s}{\naturalsPlus}}$ and
            for all $n\in\naturalsPlus$ we have $\apply{s}{n}$ is open in $X$ with respect to $T$.
\end{definition}

% from §40 helper definition 1: shrink layer inside an open set
\begin{definition}\label{regular_shrink_layer}
    $\regularShrinkLayer{b}{W}{n}{X}{T} = \{C\in\apply{b}{n} \mid \closure{C}{X}{T}\subseteq W\}$.
\end{definition}

% from §40 helper definition 2: Nagata-Smirnov basis index set
\begin{definition}\label{nagata_smirnov_index_set}
    $\nagataSmirnovIndex{b}{B} = \{p\in\naturalsPlus\times B \mid \snd{p}\in\apply{b}{\fst{p}}\}$.
\end{definition}

% from §40 helper lemma 1: open sets as countable unions with closures inside
\begin{lemma}\label{regular_countably_locally_finite_basis_open_shrink}
    Assume $X$ is regular with respect to $T$.
    Assume $B$ is a basis for $T$ on $X$.
    Assume $B$ is countably locally finite in $X$ with respect to $T$.
    Assume $W$ is open in $X$ with respect to $T$.
    Then there exists $u$ such that
        $u$ is a sequence in $\pow{X}$ and
        for all $n\in\naturalsPlus$ we have
            $\apply{u}{n}$ is open in $X$ with respect to $T$ and
            $\closure{\apply{u}{n}}{X}{T}\subseteq W$ and
        $W = \unions{\img{u}{\naturalsPlus}}$.
\end{lemma}
\begin{proof}
    $T$ is a topology on $X$ by \cref{regular_space}.
    Take $b$ such that
        $b$ is a function on $\naturalsPlus$ and
        $B = \unions{\img{b}{\naturalsPlus}}$ and
        for all $n\in\naturalsPlus$ we have $\apply{b}{n}$ is locally finite in $X$ with respect to $T$
        by \cref{countably_locally_finite_collection}.
    Let $u = \{(n,\unions{\regularShrinkLayer{b}{W}{n}{X}{T}}) \mid n\in\naturalsPlus\}$.

    For all $w$ such that $w\in u$ there exist $x,y$ such that $w = (x,y)$.
    We have $u$ is a relation by \cref{relation}.
    For all $n,y_1,y_2$ such that $(n,y_1)\in u$ and $(n,y_2)\in u$ we have $y_1 = y_2$
        by \cref{pair_eq_iff}.
    Thus $u$ is a function on $\naturalsPlus$
        by \cref{rightunique,relation,dom_intro,dom_iff,pair_eq_iff,subseteq,subseteq_antisymmetric}.

    For all $n\in\dom{u}$ we have $\apply{u}{n}\in\pow{X}$
        by \cref{subseteq,pair_eq_iff,function_apply_intro,regular_shrink_layer,locally_finite_collection,funs_type_apply,powerset_closed_unions}.
    Hence $u\in\funs{\naturalsPlus}{\pow{X}}$ by \cref{funs_intro}.
    Therefore $u$ is a sequence in $\pow{X}$ by \cref{sequence_in}.

    Show that for all $n\in\naturalsPlus$ we have
        $\apply{u}{n}$ is open in $X$ with respect to $T$ and
        $\closure{\apply{u}{n}}{X}{T}\subseteq W$.
    \begin{subproof}
        Fix $n\in\naturalsPlus$.
        Let $M = \regularShrinkLayer{b}{W}{n}{X}{T}$.
        Let $H = \{D\in\pow{X} \mid \exists E\in M. D = \closure{E}{X}{T}\}$.
        We have $\apply{u}{n} = \unions{M}$ by \cref{pair_eq_iff,function_apply_intro}.
        Hence $\apply{u}{n}$ is open in $X$ with respect to $T$
            by \cref{regular_shrink_layer,subseteq,function_apply_elim,img_iff,countably_locally_finite_collection,unions_intro,basis_topology_unions,pair_eq_iff,open_in}.

        We have $M$ is locally finite in $X$ with respect to $T$
            by \cref{regular_shrink_layer,subseteq,locally_finite_subcollection}.
        Thus $H$ is the closure collection of $M$ in $X$ with respect to $T$ by \cref{closure_collection_in}.
        Hence $\closure{\unions{M}}{X}{T} = \unions{H}$ by \cref{closure_of_union_locally_finite}.
        $\unions{H}\subseteq W$
            by \cref{unions_iff,closure_collection_in,regular_shrink_layer,subseteq,pair_eq_iff}.
        Therefore $\closure{\apply{u}{n}}{X}{T}\subseteq W$
            by \cref{subseteq,pair_eq_iff,function_apply_intro}.
    \end{subproof}

    Show that $\unions{\img{u}{\naturalsPlus}}\subseteq W$.
    \begin{subproof}
        Follows by \cref{unions_iff,img_iff,function_apply_iff,sequence_in,funs_type_apply,powerset_elim,regular_space,subseteq_closure,subseteq}.
    \end{subproof}

    Show that for all $x$ such that $x\in W$ we have $x\in\unions{\img{u}{\naturalsPlus}}$.
    \begin{subproof}
        Fix $x$ such that $x\in W$.
        We have $x\in X$ by \cref{open_in,topology_on,subseteq,powerset_elim}.
        $X\setminus W$ is closed in $X$ with respect to $T$
            by \cref{open_in,topology_on,subseteq,powerset_elim,setminus_subseteq,double_relative_complement,closed_in}.
        Take $U_0,V_0$ such that
            $U_0$ is open in $X$ with respect to $T$ and
            $V_0$ is open in $X$ with respect to $T$ and
            $x\in U_0$ and
            $X\setminus W\subseteq V_0$ and
            $U_0$ is disjoint from $V_0$
            by \cref{regular_space,setminus_elim_right}.
        Take $C\in B$ such that $x\in C$ and $C\subseteq U_0$
            by \cref{basis_for_topology,topology_generated_by_basis,open_in}.
        We have $X\setminus V_0$ is closed in $X$ with respect to $T$
            by \cref{open_in,topology_on,subseteq,powerset_elim,setminus_subseteq,double_relative_complement,closed_in}.
        Hence $C\subseteq X\setminus V_0$
            by \cref{disjoint,subseteq,open_in,topology_on,powerset_elim,setminus_intro,subseteq_transitive}.
        Therefore $\closure{C}{X}{T}\subseteq X\setminus V_0$
            by \cref{regular_space,basis_for_topology,basis_on,subseteq,powerset_elim,closure_subseteq_closed}.
        Thus $\closure{C}{X}{T}\subseteq W$
            by \cref{setminus_elim_left,setminus_intro,subseteq,setminus_elim_right,subseteq_transitive}.
        Take $D$ such that $D\in\img{b}{\naturalsPlus}$ and $C\in D$ by \cref{countably_locally_finite_collection,unions_iff}.
        Take $n\in\naturalsPlus$ such that $n\mathrel{b} D$ by \cref{img_iff}.
        Let $M = \regularShrinkLayer{b}{W}{n}{X}{T}$.
        We have $C\in M$ by \cref{function_apply_iff,subseteq,regular_shrink_layer}.
        Hence $x\in\apply{u}{n}$ by \cref{subseteq,unions_intro,pair_eq_iff,function_apply_intro}.
        Therefore $x\in\unions{\img{u}{\naturalsPlus}}$ by \cref{img_iff,unions_intro}.
    \end{subproof}
    Hence $W\subseteq\unions{\img{u}{\naturalsPlus}}$ by \cref{subseteq}.
    Therefore $W = \unions{\img{u}{\naturalsPlus}}$ by \cref{subseteq_antisymmetric}.
\end{proof}

% from §40 Item 2: regular space with countably locally finite basis is normal
\begin{lemma}\label{regular_countably_locally_finite_basis_normal}
    Assume $X$ is regular with respect to $T$.
    Assume $B$ is a basis for $T$ on $X$.
    Assume $B$ is countably locally finite in $X$ with respect to $T$.
    Then $X$ is normal with respect to $T$.
\end{lemma}
\begin{proof}
    $T$ is a topology on $X$ by \cref{regular_space}.
    Show that for all $A_0,B_0$ such that $A_0$ is closed in $X$ with respect to $T$ and $B_0$ is closed in $X$ with respect to $T$
        and $A_0$ is disjoint from $B_0$
        there exist $U,V$ such that
            $U$ is open in $X$ with respect to $T$ and
            $V$ is open in $X$ with respect to $T$ and
            $A_0\subseteq U$ and
            $B_0\subseteq V$ and
            $U$ is disjoint from $V$.
    \begin{subproof}
        Fix $A_0$.
        Fix $B_0$ such that $A_0$ is closed in $X$ with respect to $T$ and $B_0$ is closed in $X$ with respect to $T$
            and $A_0$ is disjoint from $B_0$.
        Let $W_A = X\setminus B_0$.
        Let $W_B = X\setminus A_0$.
        $W_A$ is open in $X$ with respect to $T$ and $W_B$ is open in $X$ with respect to $T$
            by \cref{closed_in}.
        Take $u$ such that
            $u$ is a sequence in $\pow{X}$ and
            for all $n\in\naturalsPlus$ we have
                $u(n)$ is open in $X$ with respect to $T$ and
                $\closure{u(n)}{X}{T}\subseteq W_A$ and
            $W_A = \unions{\img{u}{\naturalsPlus}}$
            by \cref{regular_countably_locally_finite_basis_open_shrink}.
        Take $v$ such that
            $v$ is a sequence in $\pow{X}$ and
            for all $n\in\naturalsPlus$ we have
                $v(n)$ is open in $X$ with respect to $T$ and
                $\closure{v(n)}{X}{T}\subseteq W_B$ and
            $W_B = \unions{\img{v}{\naturalsPlus}}$
            by \cref{regular_countably_locally_finite_basis_open_shrink}.

        We have $W_A = \unions{\img{u}{\naturalsPlus}}$.
        $A_0\subseteq\unions{\img{u}{\naturalsPlus}}$
            by \cref{closed_in,subseteq,disjoint,setminus_intro}.

        We have $W_B = \unions{\img{v}{\naturalsPlus}}$.
        $B_0\subseteq\unions{\img{v}{\naturalsPlus}}$
            by \cref{closed_in,subseteq,disjoint,setminus_intro}.

        Let $c = \{w \mid \exists n\in\naturalsPlus. w = (n,\closure{v(n)}{X}{T})\}$.
        Let $d = \{w \mid \exists n\in\naturalsPlus. w = (n,\closure{u(n)}{X}{T})\}$.
        For all $w$ such that $w\in c$ there exist $x,y$ such that $w = (x,y)$.
        For all $n,y_1,y_2$ such that $(n,y_1)\in c$ and $(n,y_2)\in c$ we have $y_1 = y_2$.
        Hence $c$ is a function on $\naturalsPlus$
            by \cref{relation,rightunique,pair_eq_iff,dom_iff,subseteq,dom_intro,subseteq_antisymmetric}.

        For all $w$ such that $w\in d$ there exist $x,y$ such that $w = (x,y)$.
        For all $n,y_1,y_2$ such that $(n,y_1)\in d$ and $(n,y_2)\in d$ we have $y_1 = y_2$.
        Thus $d$ is a function on $\naturalsPlus$
            by \cref{relation,rightunique,pair_eq_iff,dom_iff,subseteq,dom_intro,subseteq_antisymmetric}.

        Let $F = \{W \mid \exists n\in\naturalsPlus. W = u(n)\setminus \unions{\img{\restrl{c}{\initialSegment{n}}}{\initialSegment{n}}}\}$.
        Let $G = \{W \mid \exists n\in\naturalsPlus. W = v(n)\setminus \unions{\img{\restrl{d}{\initialSegment{n}}}{\initialSegment{n}}}\}$.
        Show that for all $W$ such that $W\in F$ we have $W\in T$.
        \begin{subproof}
            Fix $W$ such that $W\in F$.
            Take $n\in\naturalsPlus$ such that $W = u(n)\setminus \unions{\img{\restrl{c}{\initialSegment{n}}}{\initialSegment{n}}}$ by \cref{pair_eq_iff}.
            $\restrl{c}{\initialSegment{n}}$ is a function by \cref{restrl_of_function_is_function}.
            $\initialSegment{n}$ is finite by \cref{initial_segment_finite}.
            We have $\dom{\restrl{c}{\initialSegment{n}}} = \initialSegment{n}$
                by \cref{dom_of_restrl_eq_inter,initial_segment_naturalsplus,subseteq,inter_absorb_supseteq_left}.
            Hence $\img{\restrl{c}{\initialSegment{n}}}{\initialSegment{n}}$ is finite by \cref{finite_image_function}.
            For all $C$ such that $C\in \img{\restrl{c}{\initialSegment{n}}}{\initialSegment{n}}$ we have
                $C$ is closed in $X$ with respect to $T$ by
                \cref{img_iff,pair_eq_iff,restrl_iff,sequence_in,funs_type_apply,powerset_elim,closure_closed_in}.
            We have $\img{\restrl{c}{\initialSegment{n}}}{\initialSegment{n}}\subseteq\pow{X}$
                by \cref{closed_in,powerset_intro,subseteq}.
            Hence $\unions{\img{\restrl{c}{\initialSegment{n}}}{\initialSegment{n}}}$ is closed in $X$ with respect to $T$
                by \cref{closed_finite_unions}.
            We have $u(n)\in T$ and $u(n)\subseteq X$ by \cref{open_in,topology_on,subseteq,powerset_elim}.
            Therefore $W\in T$ by \cref{closed_in,setminus_eq_inter_complement,open_in,topology_on}.
        \end{subproof}
        Hence $F\subseteq T$ by \cref{subseteq}.

        Show that for all $W$ such that $W\in G$ we have $W\in T$.
        \begin{subproof}
            Fix $W$ such that $W\in G$.
            Take $n\in\naturalsPlus$ such that $W = v(n)\setminus \unions{\img{\restrl{d}{\initialSegment{n}}}{\initialSegment{n}}}$ by \cref{pair_eq_iff}.
            $\restrl{d}{\initialSegment{n}}$ is a function by \cref{restrl_of_function_is_function}.
            $\initialSegment{n}$ is finite by \cref{initial_segment_finite}.
            We have $\dom{\restrl{d}{\initialSegment{n}}} = \initialSegment{n}$
                by \cref{dom_of_restrl_eq_inter,initial_segment_naturalsplus,subseteq,inter_absorb_supseteq_left}.
            Hence $\img{\restrl{d}{\initialSegment{n}}}{\initialSegment{n}}$ is finite by \cref{finite_image_function}.
            For all $C$ such that $C\in \img{\restrl{d}{\initialSegment{n}}}{\initialSegment{n}}$ we have
                $C$ is closed in $X$ with respect to $T$ by
                \cref{img_iff,pair_eq_iff,restrl_iff,sequence_in,funs_type_apply,powerset_elim,closure_closed_in}.
            We have $\img{\restrl{d}{\initialSegment{n}}}{\initialSegment{n}}\subseteq\pow{X}$
                by \cref{closed_in,powerset_intro,subseteq}.
            Hence $\unions{\img{\restrl{d}{\initialSegment{n}}}{\initialSegment{n}}}$ is closed in $X$ with respect to $T$
                by \cref{closed_finite_unions}.
            We have $v(n)\in T$ and $v(n)\subseteq X$ by \cref{open_in,topology_on,subseteq,powerset_elim}.
            Therefore $W\in T$ by \cref{closed_in,setminus_eq_inter_complement,open_in,topology_on}.
        \end{subproof}
        Hence $G\subseteq T$ by \cref{subseteq}.

        Let $U_0 = \unions{F}$.
        Let $V_0 = \unions{G}$.
        $U_0$ is open in $X$ with respect to $T$ and $V_0$ is open in $X$ with respect to $T$
            by \cref{topology_on,subseteq,open_in}.

        Show that for all $x$ such that $x\in A_0$ we have $x\in U_0$.
        \begin{subproof}
            Fix $x$ such that $x\in A_0$.
            Take $W\in\img{u}{\naturalsPlus}$ such that $x\in W$ by \cref{subseteq,unions_iff}.
            Take $n\in\naturalsPlus$ such that $n\mathrel{u} W$ by \cref{img_iff}.
            Hence $x\in u(n)$ by \cref{sequence_in,funs_is_function,funs,function_apply_iff,pair_eq_iff}.
            Therefore $x\in u(n)\setminus \unions{\img{\restrl{c}{\initialSegment{n}}}{\initialSegment{n}}}$
                by \cref{unions_iff,img_iff,pair_eq_iff,restrl_iff,subseteq,setminus_elim_right,setminus_intro}.
            Hence $x\in U_0$ by \cref{unions_intro}.
        \end{subproof}
        Therefore $A_0\subseteq U_0$ by \cref{subseteq}.

        Show that for all $x$ such that $x\in B_0$ we have $x\in V_0$.
        \begin{subproof}
            Fix $x$ such that $x\in B_0$.
            Take $W\in\img{v}{\naturalsPlus}$ such that $x\in W$ by \cref{subseteq,unions_iff}.
            Take $n\in\naturalsPlus$ such that $n\mathrel{v} W$ by \cref{img_iff}.
            Hence $x\in v(n)$ by \cref{sequence_in,funs_is_function,funs,function_apply_iff,pair_eq_iff}.
            Therefore $x\in v(n)\setminus \unions{\img{\restrl{d}{\initialSegment{n}}}{\initialSegment{n}}}$
                by \cref{unions_iff,img_iff,pair_eq_iff,restrl_iff,subseteq,setminus_elim_right,setminus_intro}.
            Hence $x\in V_0$ by \cref{unions_intro}.
        \end{subproof}
        Therefore $B_0\subseteq V_0$ by \cref{subseteq}.

        Show that $U_0$ is disjoint from $V_0$.
        \begin{subproof}
            Suppose not.
            There exists $x$ such that $x\in U_0$ and $x\in V_0$ by \cref{disjoint}.
            Take $W_1$ such that $W_1\in F$ and $x\in W_1$ by \cref{unions_iff}.
            Take $W_2$ such that $W_2\in G$ and $x\in W_2$ by \cref{unions_iff}.
            Take $j\in\naturalsPlus$ such that $W_1 = u(j)\setminus \unions{\img{\restrl{c}{\initialSegment{j}}}{\initialSegment{j}}}$.
            Take $k\in\naturalsPlus$ such that $W_2 = v(k)\setminus \unions{\img{\restrl{d}{\initialSegment{k}}}{\initialSegment{k}}}$.
            $x\in u(j)$ and $x\in v(k)$ by \cref{setminus_elim_left}.
            $x\notin \unions{\img{\restrl{c}{\initialSegment{j}}}{\initialSegment{j}}}$ and
                $x\notin \unions{\img{\restrl{d}{\initialSegment{k}}}{\initialSegment{k}}}$ by \cref{setminus_elim_right}.
            \begin{byCase}
            \caseOf{$j\leq k$.}
                $j\in\initialSegment{k}$ by \cref{initial_segment_naturalsplus}.
                $\closure{u(j)}{X}{T}\in \img{\restrl{d}{\initialSegment{k}}}{\initialSegment{k}}$
                    by \cref{img_iff,restrl_iff}.
                We have $x\in \closure{u(j)}{X}{T}$ by \cref{sequence_in,funs_type_apply,powerset_elim,subseteq_closure,elem_subseteq}.
                Therefore $x\in \unions{\img{\restrl{d}{\initialSegment{k}}}{\initialSegment{k}}}$ by \cref{unions_intro}.
                Contradiction.
            \caseOf{$k\leq j$.}
                $k\in\initialSegment{j}$ by \cref{initial_segment_naturalsplus}.
                $\closure{v(k)}{X}{T}\in \img{\restrl{c}{\initialSegment{j}}}{\initialSegment{j}}$
                    by \cref{img_iff,restrl_iff}.
                We have $x\in \closure{v(k)}{X}{T}$ by \cref{sequence_in,funs_type_apply,powerset_elim,subseteq_closure,elem_subseteq}.
                Therefore $x\in \unions{\img{\restrl{c}{\initialSegment{j}}}{\initialSegment{j}}}$ by \cref{unions_intro}.
                Contradiction.
            \end{byCase}
        \end{subproof}

        Therefore there exist $U_1,V_1$ such that
            $U_1$ is open in $X$ with respect to $T$ and
            $V_1$ is open in $X$ with respect to $T$ and
            $A_0\subseteq U_1$ and
            $B_0\subseteq V_1$ and
            $U_1$ is disjoint from $V_1$.
    \end{subproof}
    We have $X$ has closed singletons with respect to $T$ by \cref{regular_space}.
    Therefore $X$ is normal with respect to $T$ by \cref{normal_space}.
\end{proof}

% from §40 Item 2: closed sets are $G_\delta$ sets in regular spaces with countably locally finite bases
\begin{lemma}\label{regular_countably_locally_finite_basis_closed_g_delta}
    Assume $X$ is regular with respect to $T$.
    Assume $B$ is a basis for $T$ on $X$.
    Assume $B$ is countably locally finite in $X$ with respect to $T$.
    Assume $A$ is closed in $X$ with respect to $T$.
    Then $A$ is a Gdelta set in $X$ with respect to $T$.
\end{lemma}
\begin{proof}
    $T$ is a topology on $X$ by \cref{regular_space}.
    Let $W = X\setminus A$.
    $W$ is open in $X$ with respect to $T$ by \cref{closed_in}.
    Take $u$ such that
        $u$ is a sequence in $\pow{X}$ and
        for all $n\in\naturalsPlus$ we have
            $u(n)$ is open in $X$ with respect to $T$ and
            $\closure{u(n)}{X}{T}\subseteq W$ and
        $W = \unions{\img{u}{\naturalsPlus}}$
        by \cref{regular_countably_locally_finite_basis_open_shrink}.
    Let $s = \{(n,X\setminus \closure{u(n)}{X}{T}) \mid n\in\naturalsPlus\}$.

    We have $s$ is a relation by \cref{relation}.
    For all $n,y_1,y_2$ such that $(n,y_1)\in s$ and $(n,y_2)\in s$ we have $y_1 = y_2$
        by \cref{pair_eq_iff}.
    Hence $s$ is a function on $\naturalsPlus$
        by \cref{rightunique,relation,dom_intro,dom_iff,pair_eq_iff,subseteq,subseteq_antisymmetric}.

    For all $n\in\dom{s}$ we have $s(n)\in\pow{X}$
        by \cref{subseteq,sequence_in,funs_type_apply,powerset_elim,closure_closed_in,closed_in,setminus_subseteq,powerset_intro,pair_eq_iff,function_apply_intro}.
    Hence $s\in\funs{\naturalsPlus}{\pow{X}}$ by \cref{funs_intro}.
    Therefore $s$ is a sequence in $\pow{X}$ by \cref{sequence_in}.

    For all $n\in\naturalsPlus$ we have $s(n)$ is open in $X$ with respect to $T$
        by \cref{sequence_in,funs_type_apply,powerset_elim,closure_closed_in,pair_eq_iff,function_apply_intro,closed_in}.

    Show that for all $x$ such that $x\in A$ we have $x\in\inters{\img{s}{\naturalsPlus}}$.
    \begin{subproof}
        Follows by \cref{closed_in,subseteq,real_one_in_naturals,funs_is_function,funs_elim,img_iff,function_apply_elim,setminus_elim_right,function_apply_iff,setminus_intro,pair_eq_iff,function_apply_intro,inters_intro}.
    \end{subproof}
    Hence $A\subseteq\inters{\img{s}{\naturalsPlus}}$ by \cref{subseteq}.

    Show that for all $x$ such that $x\in\inters{\img{s}{\naturalsPlus}}$ we have $x\in A$.
    \begin{subproof}
        Fix $x$ such that $x\in\inters{\img{s}{\naturalsPlus}}$.
        We have $x\in\unions{\img{s}{\naturalsPlus}}$ by \cref{inters}.
        Take $S_0$ such that $S_0\in\img{s}{\naturalsPlus}$ and $x\in S_0$ by \cref{unions_iff}.
        Take $m\in\naturalsPlus$ such that $m\mathrel{s} S_0$ by \cref{img_iff}.
        Hence $x\in X\setminus \closure{u(m)}{X}{T}$ by \cref{function_apply_iff,funs_is_function,funs_elim,pair_eq_iff,function_apply_intro}.
        Suppose not.
        We have $x\in\unions{\img{u}{\naturalsPlus}}$ by \cref{setminus_elim_left,setminus_intro,pair_eq_iff,subseteq}.
        Take $U_0$ such that $U_0\in\img{u}{\naturalsPlus}$ and $x\in U_0$ by \cref{unions_iff}.
        Take $n\in\naturalsPlus$ such that $n\mathrel{u} U_0$ by \cref{img_iff}.
        Hence $x\in u(n)$ by \cref{sequence_in,funs_is_function,funs,function_apply_iff,pair_eq_iff}.
        Thus $x\in\closure{u(n)}{X}{T}$ by \cref{sequence_in,funs_type_apply,powerset_elim,subseteq_closure,subseteq}.
        Therefore $x\in s(n)$ by \cref{funs_is_function,funs_elim,img_iff,function_apply_elim,inters}.
        Hence $x\in X\setminus \closure{u(n)}{X}{T}$ by \cref{pair_eq_iff,function_apply_intro}.
        Contradiction by \cref{setminus_elim_right}.
    \end{subproof}
    Hence $\inters{\img{s}{\naturalsPlus}}\subseteq A$ by \cref{subseteq}.
    Thus $A = \inters{\img{s}{\naturalsPlus}}$ by \cref{subseteq_antisymmetric}.
    Therefore $A$ is a Gdelta set in $X$ with respect to $T$ by \cref{g_delta_set}.
\end{proof}

% from §40 helper lemma 2: distance to a singleton equals the metric value
\begin{lemma}\label{point_set_distance_singleton}
    Assume $d$ is a metric on $X$.
    Assume $x\in X$.
    Assume $a\in X$.
    Assume $r$ is the distance from $x$ to $\{a\}$ in $X$ with respect to $d$.
    Then $r = d((x,a))$.
\end{lemma}
\begin{proof}
    Let $S = \pointDistanceSet{d}{\{a\}}{x}$.
    $S\subseteq \{d((x,a))\}$ by \cref{point_distance_set,singleton_elim,pair_eq_iff,singleton_intro,subseteq}.
    $\{d((x,a))\}\subseteq S$ by \cref{singleton_elim,point_distance_set,singleton_intro,subseteq}.
    We have $S = \{d((x,a))\}$ by \cref{subseteq_antisymmetric}.
    Hence $d((x,a))\in S$ by \cref{singleton_intro}.
    Thus $r \leq d((x,a))$ by \cref{singleton_intro,point_set_distance,real_infimum,real_lower_bound}.
    We have $S\subseteq\reals$ by \cref{times_tuple_intro,pair_eq_iff,metric_on,funs_type_apply,singleton_subset_intro}.
    $d((x,a))$ is a lower bound of $S$ by \cref{subseteq,singleton_elim,real_lower_bound}.
    Hence $d((x,a)) \leq r$ by \cref{point_set_distance,real_infimum}.
    We have $r\in\reals$ and $d((x,a))\in\reals$ by \cref{point_set_distance,metric_on,funs_type_apply,times_tuple_intro}.
    Therefore $r = d((x,a))$ by \cref{real_leq_antisym}.
\end{proof}

% from §40 helper lemma 3: Urysohn family for a closed Gdelta witness
\begin{lemma}\label{normal_closed_g_delta_urysohn_family}
    Assume $X$ is normal with respect to $T$.
    Assume $A$ is closed in $X$ with respect to $T$.
    Assume $u$ is a sequence in $\pow{X}$.
    Assume $A = \inters{\img{u}{\naturalsPlus}}$.
    Assume for all $n\in\naturalsPlus$ we have $u(n)$ is open in $X$ with respect to $T$.
    Then there exists $f$ such that
        $f$ is a function on $\naturalsPlus$ and
        for all $n\in\naturalsPlus$ we have $f(n)$ is continuous from $X$ to $\reals$ with respect to $T$ and $\standardTopologyR$ and
        for all $n\in\naturalsPlus$ we have $\img{f(n)}{X}\subseteq \intervalclosed{0}{1 / n} $ and
        for all $n\in\naturalsPlus$ we have $\img{f(n)}{A}\subseteq \{0\}$ and
        for all $n\in\naturalsPlus$ we have $\img{f(n)}{X\setminus u(n)}\subseteq \{1 / n\}$.
\end{lemma}
\begin{proof}
    $T$ is a topology on $X$ by \cref{normal_space}.
    Let $R = \{w\in\naturalsPlus\times\pow{X\times\reals} \mid
        \text{there exists $n\in\naturalsPlus$ such that there exists $g$ such that
        $w = (n,g)$ and
        $g$ is continuous from $X$ to $\reals$ with respect to $T$ and $\standardTopologyR$ and
        $\img{g}{X}\subseteq \intervalclosed{0}{1 / n}$ and
        $\img{g}{A}\subseteq \{0\}$ and
        $\img{g}{X\setminus u(n)}\subseteq \{1 / n\}$}\}$.
    $R$ is a relation by \cref{relation,subseteq,times_elem_is_tuple}.

    Show that for all $n\in\naturalsPlus$ there exists $g$ such that $n\mathrel{R} g$.
    \begin{subproof}
        Fix $n\in\naturalsPlus$.
        Let $I_n = \intervalclosed{0}{1 / n}$.
        For all $x$ such that $x\in A$ we have $x\in u(n)$
            by \cref{pair_eq_iff,sequence_in,funs_is_function,funs,img_iff,function_apply_elim,inters}.
        Hence $A\subseteq u(n)$ by \cref{subseteq}.
        Let $B = X\setminus u(n)$.
        $B$ is closed in $X$ with respect to $T$
            by \cref{open_in,topology_on,subseteq,powerset_elim,setminus_subseteq,double_relative_complement,closed_in}.
        We have $A$ is disjoint from $B$ by \cref{setminus_elim_right,subseteq,disjoint}.
        $0 < 1$ and $0\in\reals$ and $1\in\reals$ and $n\in\reals$
            by \cref{real_zero_lt_one,real_add_ident,real_naturals_subset_reals,real_one_in_naturals,subseteq}.
        Thus $0 < n$ by \cref{real_one_leq_naturalsplus,real_lt_subst_right,real_lt_trans}.
        Therefore $n\neq 0$ by \cref{real_pos_nonzero}.
        Hence $\inv{n}\in\reals$ by \cref{real_mul_inverse}.
        We have $1 / n\in\reals$ by \cref{real_div,real_mul_closed}.
        $0 < 1 / n$ by \cref{positive_reciprocal_naturalsplus}.
        Let $I_0 = \intervalclosed{0}{1}$.
        Take $p$ such that
            $p$ is continuous from $X$ to $I_0$ with respect to $T$ and $\subspaceTopology{\standardTopologyR}{I_0}$ and
            $\img{p}{A}\subseteq \{0\}$ and
            $\img{p}{B}\subseteq \{1\}$
            by \cref{urysohn_unit_interval}.
        Let $j_0 = \identity{I_0}$.
        Let $p_0 = j_0\circ p$.
        $\standardTopologyR$ is a topology on $\reals$
            by \cref{standard_basis_is_basis,basis_topology_is_topology,standard_topology_reals}.
        $I_0\subseteq\reals$ by \cref{intervalclosed,subseteq}.
        Hence $I_0$ is a subspace of $\reals$ with respect to $\standardTopologyR$ by \cref{subspace_of}.
        $p_0$ is continuous from $X$ to $\reals$ with respect to $T$ and $\standardTopologyR$
            by \cref{continuous_expand_range}.
        $p\in\funs{X}{I_0}$ and $p$ is a function by \cref{continuous,funs_is_function}.
        Hence $p$ is right-unique and $p$ is a relation and $\dom{p} = X$
            by \cref{function_rightunique,funs_is_relation,funs_elim}.
        We have $p\in\rels{X}{I_0}$ by \cref{funs_subseteq_rels,subseteq}.
        Hence $\ran{p}\subseteq I_0$ by \cref{rels_ran_subseteq}.
        Therefore $\ran{p}\subseteq\dom{j_0}$ by \cref{rels_ran_subseteq,id_dom}.
        $j_0$ is right-unique and $j_0$ is a relation by \cref{id_rightunique,id_is_relation}.
        For all $x\in X$ we have $p_0(x) = p(x)$ by \cref{circ_apply,funs_elim,id_apply}.
        Let $h = \{(x,1 / n) \mid x\in X\}$.
        $h$ is a function from $X$ to $\reals$ by \cref{constant_map_function}.
        Hence for all $x\in X$ we have $h(x) = 1 / n$
            by \cref{constant_map_function,funs_is_function,function_apply_intro}.
        $h$ is continuous from $X$ to $\reals$ with respect to $T$ and $\standardTopologyR$
            by \cref{constant_continuous}.
        Let $g = \{(x, h(x) \rmul p_0(x)) \mid x\in X\}$.
        $g$ is continuous from $X$ to $\reals$ with respect to $T$ and $\standardTopologyR$
            by \cref{real_function_multiplication_continuous}.
        $g\in\funs{X}{\reals}$ and $g$ is a function by \cref{continuous,funs_is_function}.
        Show that for all $x\in X$ we have $g(x)\in I_n$.
        \begin{subproof}
            Follows by \cref{circ_apply,funs_elim,id_apply,intervalclosed,function_apply_iff,real_lt_trichotomy,real_mul_zero,real_mul_comm,real_lt_imp_pos_diff,real_mul_pos_pos_gt,real_lt_imp_leq,real_mul_ident,real_naturals_subset_reals,real_one_in_naturals,subseteq,real_lt_mul_pos,continuous}.
        \end{subproof}
        Therefore $\img{g}{X}\subseteq I_n$
            by \cref{subseteq,img_iff,continuous,funs_is_function,function_apply_iff}.
        Show that for all $x\in A$ we have
            $x\in X$ and $p(x)\in\img{p}{A}$ and $p(x) = 0$ and $p_0(x) = 0$ and $p_0(x)\in\reals$ and
            $(x, h(x) \rmul p_0(x))\in g$ and $g(x) = 0$.
        \begin{subproof}
            Follows by \cref{closed_in,subseteq,function_apply_elim,img_iff,singleton_elim,circ_apply,funs_elim,id_apply,real_mul_comm,real_mul_zero,function_apply_iff}.
        \end{subproof}
        Show that for all $x\in B$ we have
            $x\in X$ and $p(x)\in\img{p}{B}$ and $p(x) = 1$ and $p_0(x) = 1$ and $p_0(x)\in\reals$ and
            $h(x) = 1 / n$ and $(x, h(x) \rmul p_0(x))\in g$ and $g(x) = 1 / n$.
        \begin{subproof}
            Follows by \cref{closed_in,subseteq,function_apply_elim,img_iff,singleton_elim,circ_apply,funs_elim,id_apply,real_mul_ident,function_apply_iff,real_mul_comm}.
        \end{subproof}
        $\img{g}{A}\subseteq \{0\}$ by \cref{img_iff,funs_is_function,function_apply_iff,singleton_intro,subseteq}.
        $\img{g}{X\setminus u(n)}\subseteq \{1 / n\}$ by \cref{img_iff,funs_is_function,function_apply_iff,singleton_intro,subseteq}.
        $g\subseteq X\times\reals$
            by \cref{funs_subseteq_rels,funs,rels_elim,subseteq,times_elem_is_tuple}.
        Hence $g\in\pow{X\times\reals}$ by \cref{powerset_intro}.
        We have $(n,g)\in \naturalsPlus\times\pow{X\times\reals}$ by \cref{times_tuple_intro}.
        Therefore $n\mathrel{R} g$ by \cref{pair_eq_iff}.
    \end{subproof}

    Take $f$ such that
        $f$ is a function on $\naturalsPlus$ and
        for all $n\in\naturalsPlus$ we have $n\mathrel{R} f(n)$
        by \cref{axiom_of_choice}.
    Show that for all $n\in\naturalsPlus$ we have $f(n)$ is continuous from $X$ to $\reals$ with respect to $T$ and $\standardTopologyR$.
    \begin{subproof}
        Follows by \cref{function_apply_elim,pair_eq_iff}.
    \end{subproof}
    Show that for all $n\in\naturalsPlus$ we have $\img{f(n)}{X}\subseteq \intervalclosed{0}{1 / n} $.
    \begin{subproof}
        Follows by \cref{function_apply_elim,pair_eq_iff}.
    \end{subproof}
    Show that for all $n\in\naturalsPlus$ we have $\img{f(n)}{A}\subseteq \{0\}$.
    \begin{subproof}
        Follows by \cref{function_apply_elim,pair_eq_iff}.
    \end{subproof}
    Show that for all $n\in\naturalsPlus$ we have $\img{f(n)}{X\setminus u(n)}\subseteq \{1 / n\}$.
    \begin{subproof}
        Follows by \cref{function_apply_elim,pair_eq_iff}.
    \end{subproof}
    Therefore there exists $g$ such that
        $g$ is a function on $\naturalsPlus$ and
        for all $n\in\naturalsPlus$ we have $g(n)$ is continuous from $X$ to $\reals$ with respect to $T$ and $\standardTopologyR$ and
        for all $n\in\naturalsPlus$ we have $\img{g(n)}{X}\subseteq \intervalclosed{0}{1 / n} $ and
        for all $n\in\naturalsPlus$ we have $\img{g(n)}{A}\subseteq \{0\}$ and
        for all $n\in\naturalsPlus$ we have $\img{g(n)}{X\setminus u(n)}\subseteq \{1 / n\}$.
\end{proof}

% from §40 Item 3: continuous function positive off a closed $G_\delta$ set
\begin{lemma}\label{normal_closed_g_delta_positive_function}
    Assume $X$ is normal with respect to $T$.
    Assume $A$ is closed in $X$ with respect to $T$.
    Assume $A$ is a Gdelta set in $X$ with respect to $T$.
    Let $I = \intervalclosed{0}{1}$.
    Let $T_I = \subspaceTopology{\standardTopologyR}{I}$.
    Then there exists $f$ such that
        $f$ is continuous from $X$ to $I$ with respect to $T$ and $T_I$ and
        $\img{f}{A}\subseteq \{0\}$ and
        for all $x\in X\setminus A$ we have $0 < f(x)$.
\end{lemma}
\begin{proof}
    Take $u$ such that
        $u$ is a sequence in $\pow{X}$ and
        $A = \inters{\img{u}{\naturalsPlus}}$ and
        for all $n\in\naturalsPlus$ we have $u(n)$ is open in $X$ with respect to $T$
        by \cref{g_delta_set}.
    Take $h$ such that
        $h$ is a function on $\naturalsPlus$ and
        for all $n\in\naturalsPlus$ we have $h(n)$ is continuous from $X$ to $\reals$ with respect to $T$ and $\standardTopologyR$ and
        for all $n\in\naturalsPlus$ we have $\img{h(n)}{X}\subseteq \intervalclosed{0}{1 / n} $ and
        for all $n\in\naturalsPlus$ we have $\img{h(n)}{A}\subseteq \{0\}$ and
        for all $n\in\naturalsPlus$ we have $\img{h(n)}{X\setminus u(n)}\subseteq \{1 / n\}$
        by \cref{normal_closed_g_delta_urysohn_family}.
    Let $R = \{(n,\reals) \mid n\in\naturalsPlus\}$.
    Let $S = \{(n,\standardTopologyR) \mid n\in\naturalsPlus\}$.
    Let $P = \productTopologyIndexed{S}{R}$.
    Let $F = \{(x,z) \mid x\in X, z\in\productFamily{R} \mid
        \forall n\in\naturalsPlus. z(n) = \apply{\apply{h}{n}}{x}\}$.
    $1\in\naturalsPlus$ by \cref{real_one_in_naturals}.
    Hence $\naturalsPlus$ is inhabited.
    $\standardTopologyR$ is a topology on $\reals$
        by \cref{standard_basis_is_basis,basis_topology_is_topology,standard_topology_reals}.
    $R$ is a relation by \cref{relation}.
    For all $n,y_1,y_2$ such that $(n,y_1)\in R$ and $(n,y_2)\in R$ we have $y_1 = y_2$ by \cref{pair_eq_iff}.
    Hence $R$ is a function on $\naturalsPlus$ by \cref{rightunique,relation,dom_intro,dom_iff,pair_eq_iff,subseteq,subseteq_antisymmetric}.
    $S$ is a relation by \cref{relation}.
    For all $n,y_1,y_2$ such that $(n,y_1)\in S$ and $(n,y_2)\in S$ we have $y_1 = y_2$ by \cref{pair_eq_iff}.
    Hence $S$ is a function on $\naturalsPlus$ by \cref{rightunique,relation,dom_intro,dom_iff,pair_eq_iff,subseteq,subseteq_antisymmetric}.
    We have $\dom{S} = \naturalsPlus$ and $\dom{R} = \naturalsPlus$ by \cref{funs}.
    For all $n\in\naturalsPlus$ we have $S(n) = \standardTopologyR$ and $R(n) = \reals$ by \cref{function_apply_intro}.
    Hence $\productFamily{R} = \realsPower{\naturalsPlus}$ by \cref{product_family_reals_power}.
    $P$ is a topology on $\productFamily{R}$
        by \cref{product_subbasis_is_subbasis,product_subbasis_finite_intersections_basis,basis_topology_is_topology,basis_for_topology,standard_topology_reals}.

    $F$ is continuous from $X$ to $\productFamily{R}$ with respect to $T$ and $P$
        by \cref{continuous_real_family_map}.
    $F$ is a function and $\dom{F} = X$ by \cref{continuous,funs_elim}.

    Let $z = \{(n,0) \mid n\in\naturalsPlus\}$.
    $0\in\reals$ by \cref{real_add_ident}.
    Hence $z\in\productFamily{R}$
        by \cref{constant_map_function,reals_power,tuple_power,product_family_reals_power,subseteq}.
    $\boundedMetric{\omegaMetric}$ is a metric on $\realsPower{\naturalsPlus}$
        by \cref{bounded_metric_is_metric,omega_metric_is_metric}.
    $\metricTopology{\boundedMetric{\omegaMetric}}{\realsPower{\naturalsPlus}} = P$
        by \cref{bounded_metric_topology,omega_metric_is_metric,omega_metric_product_topology,product_family_reals_power}.
    We have $\{z\}\subseteq\realsPower{\naturalsPlus}$ by \cref{singleton_subset_intro,product_family_reals_power,subseteq}.
    $z\in\{z\}$ by \cref{singleton_intro}.
    Therefore $\{z\}\neq\emptyset$ by \cref{nonempty_inhabited,emptyset}.
    Take $q$ such that
        $q$ is a function from $\realsPower{\naturalsPlus}$ to $\reals$ and
        for all $y\in\realsPower{\naturalsPlus}$ we have $q(y)$ is the distance from $y$ to $\{z\}$ in $\realsPower{\naturalsPlus}$ with respect to $\boundedMetric{\omegaMetric}$ and
        $q$ is continuous from $\realsPower{\naturalsPlus}$ to $\reals$ with respect to $\metricTopology{\boundedMetric{\omegaMetric}}{\realsPower{\naturalsPlus}}$ and $\standardTopologyR$
        by \cref{point_set_distance_continuous,singleton_subset_intro,singleton_inhabited}.
    We have $q$ is continuous from $\realsPower{\naturalsPlus}$ to $\reals$ with respect to $P$ and $\standardTopologyR$
        by \cref{bounded_metric_topology,omega_metric_product_topology,product_family_reals_power}.
    Thus $q$ is a function and $\dom{q} = \realsPower{\naturalsPlus}$ by \cref{funs_elim}.
    $F\in\funs{X}{\realsPower{\naturalsPlus}}$ by \cref{continuous,product_family_reals_power,subseteq}.
    Therefore $\ran{F}\subseteq\dom{q}$ by \cref{funs_ran,funs_elim,subseteq}.
    Let $g = q\circ F$.
    Hence $g$ is continuous from $X$ to $\reals$ with respect to $T$ and $\standardTopologyR$ by \cref{continuous_composite}.
    Therefore $g\in\funs{X}{\reals}$ by \cref{continuous}.

    Show that for all $x$ such that $x\in A$ we have $g(x) = 0$.
    \begin{subproof}
        Fix $x$ such that $x\in A$.
        We have $x\in X$ by \cref{closed_in,subseteq}.
        Show that for all $n\in\naturalsPlus$ we have $\apply{\apply{F}{x}}{n} = z(n)$.
        \begin{subproof}
            Fix $n\in\naturalsPlus$.
            $(x,F(x))\in F$ by \cref{function_apply_elim}.
            Hence $\apply{\apply{F}{x}}{n} = z(n)$
                by \cref{pair_eq_iff,continuous,funs_elim,function_apply_elim,img_iff,singleton_elim,subseteq,constant_map_function,funs_is_function,function_apply_intro}.
        \end{subproof}
        $F(x)\in\funs{\naturalsPlus}{\reals}$ and $z\in\funs{\naturalsPlus}{\reals}$
            by \cref{continuous,funs_type_apply,product_family_reals_power,subseteq,reals_power,tuple_power}.
        We have $F(x) = z$ by \cref{functions_eq_iff}.
        Hence $g(x) = \apply{q}{F(x)}$ and $F(x)\in\realsPower{\naturalsPlus}$
            by \cref{circ_apply,continuous,funs_elim,subseteq,product_family_reals_power,funs_type_apply}.
        Thus $\apply{q}{F(x)}$ is the distance from $F(x)$ to $\{z\}$ in $\realsPower{\naturalsPlus}$ with respect to $\boundedMetric{\omegaMetric}$ by assumption.
        Therefore $g(x) = \apply{\boundedMetric{\omegaMetric}}{(F(x),z)}$ by \cref{point_set_distance_singleton,product_family_reals_power,subseteq}.
        Hence $g(x) = \apply{\boundedMetric{\omegaMetric}}{(z,z)}$ by \cref{real_lt_subst_right}.
        Therefore $g(x) = 0$ by \cref{metric_on,metric_zero_characterization,bounded_metric_is_metric,omega_metric_is_metric}.
    \end{subproof}

    $\img{g}{A}\subseteq \{0\}$ by \cref{img_iff,continuous,funs_is_function,function_apply_iff,singleton_intro,subseteq}.

    Show that for all $x$ such that $x\in X\setminus A$ we have $0 < g(x)$.
    \begin{subproof}
        Fix $x$ such that $x\in X\setminus A$.
        We have $x\in X$ by \cref{setminus_elim_left}.
        Show that there exists $n\in\naturalsPlus$ such that $x\in X\setminus u(n)$.
        \begin{subproof}
            Suppose not.
            We have for all $n\in\naturalsPlus$ it is wrong that $x\in X\setminus u(n)$.
            For all $n\in\naturalsPlus$ we have $x\in u(n)$.
            Hence $u(1)\in\img{u}{\naturalsPlus}$ by \cref{img_iff,function_apply_elim,sequence_in,funs_is_function,funs,real_one_in_naturals}.
            Therefore $x\in\unions{\img{u}{\naturalsPlus}}$ by \cref{unions_iff,real_one_in_naturals}.
            For all $a$ such that $a\in\img{u}{\naturalsPlus}$ we have $x\in a$
                by \cref{img_iff,function_apply_iff,funs_is_function,sequence_in,real_lt_subst_right}.
            Therefore $x\in\inters{\img{u}{\naturalsPlus}}$ by \cref{inters}.
            Hence $x\in A$ by \cref{real_lt_subst_right}.
            Contradiction by \cref{setminus_elim_right}.
        \end{subproof}
        Take $n\in\naturalsPlus$ such that $x\in X\setminus u(n)$.
        $(x,F(x))\in F$ by \cref{function_apply_elim}.
        Hence $\apply{\apply{F}{x}}{n} = 1 / n$
            by \cref{pair_eq_iff,continuous,funs_elim,function_apply_elim,img_iff,singleton_elim,subseteq}.
        $z(n) = 0$ by \cref{constant_map_function,funs_is_function,function_apply_intro}.
        $0 < 1 / n$ by \cref{positive_reciprocal_naturalsplus}.
        We have $F(x) \neq z$.
        Hence $g(x) = \apply{q}{F(x)}$ and $F(x)\in\realsPower{\naturalsPlus}$
            by \cref{circ_apply,continuous,funs_elim,subseteq,product_family_reals_power,funs_type_apply}.
        Thus $\apply{q}{F(x)}$ is the distance from $F(x)$ to $\{z\}$ in $\realsPower{\naturalsPlus}$ with respect to $\boundedMetric{\omegaMetric}$ by assumption.
        Therefore $g(x) = \apply{\boundedMetric{\omegaMetric}}{(F(x),z)}$ by \cref{point_set_distance_singleton,product_family_reals_power,subseteq}.
        Hence $0 < g(x)$ by \cref{metric_on,metric_nonnegative,metric_zero_characterization,real_leq_split,bounded_metric_is_metric,omega_metric_is_metric}.
    \end{subproof}

    For all $x\in X$ we have $g(x)\in I$
        by \cref{continuous,funs_type_apply,times_tuple_intro,circ_apply,funs_elim,subseteq,product_family_reals_power,point_set_distance_singleton,metric_on,metric_nonnegative,bounded_metric_is_metric,omega_metric_is_metric,bounded_metric_leq_one,intervalclosed}.
    Therefore $\img{g}{X}\subseteq I$ by \cref{subseteq,img_iff,continuous,funs_is_function,function_apply_iff}.
    Hence $g$ is continuous from $X$ to $I$ with respect to $T$ and $T_I$
        by \cref{intervalclosed,subseteq,subspace_of,continuous_restrict_range}.

    Therefore there exists $f$ such that
        $f$ is continuous from $X$ to $I$ with respect to $T$ and $T_I$ and
        $\img{f}{A}\subseteq \{0\}$ and
        for all $x\in X\setminus A$ we have $0 < f(x)$.
\end{proof}

% from §40 helper lemma 4: uniform metric on real powers
\begin{lemma}\label{uniform_metric_is_metric_late}
    Assume $J$ is inhabited.
    Then $\uniformMetric{J}$ is a metric on $\realsPower{J}$.
\end{lemma}
\begin{proof}
    Follows by \cref{uniform_metric_is_metric}.
\end{proof}

% from §40 helper lemma 5: flatten countably many countably locally finite collections
\begin{lemma}\label{countable_union_countably_locally_finite}
    Assume $c$ is a function on $\naturalsPlus$.
    Suppose for all $n\in\naturalsPlus$ we have $\apply{c}{n}\subseteq\pow{X}$.
    Suppose for all $n\in\naturalsPlus$ we have $\apply{c}{n}$ is countably locally finite in $X$ with respect to $T$.
    Then $\unions{\img{c}{\naturalsPlus}}$ is countably locally finite in $X$ with respect to $T$.
\end{lemma}
\begin{proof}
    Let $P = \pow{\naturalsPlus\times\pow{\pow{X}}}$.
    Let $R = \{w\in\naturalsPlus\times P \mid \text{there exists $n\in\naturalsPlus$ such that there exists $b$ such that
        $w = (n,b)$ and
        $b$ is a function on $\naturalsPlus$ and
        $\apply{c}{n} = \unions{\img{b}{\naturalsPlus}}$ and
        for all $k\in\naturalsPlus$ we have $\apply{b}{k}$ is locally finite in $X$ with respect to $T$}\}$.
    Show that for all $n\in\naturalsPlus$ there exists $b$ such that $n\mathrel{R} b$.
    \begin{subproof}
        Fix $n\in\naturalsPlus$.
        Take $b$ such that
            $b$ is a function on $\naturalsPlus$ and
            $\apply{c}{n} = \unions{\img{b}{\naturalsPlus}}$ and
            for all $k\in\naturalsPlus$ we have $\apply{b}{k}$ is locally finite in $X$ with respect to $T$
            by \cref{countably_locally_finite_collection}.
        $b\subseteq\naturalsPlus\times\pow{\pow{X}}$
            by \cref{function_member_elim,funs_elim,funs,subseteq,locally_finite_collection,powerset_intro,pair_eq_iff,times_tuple_intro}.
        Therefore $n\mathrel{R} b$ by \cref{powerset_intro,times_tuple_intro,pair_eq_iff}.
    \end{subproof}
    $R$ is a relation by \cref{relation,subseteq,times_elem_is_tuple}.
    Take $s$ such that
        $s$ is a function on $\naturalsPlus$ and
        for all $n\in\naturalsPlus$ we have $n\mathrel{R}\apply{s}{n}$
        by \cref{choice_function}.
    Take $g$ such that $g\in\inj{\naturalsPlus\times\naturalsPlus}{\naturalsPlus}$
        by \cref{naturalsplus_pairing_injection}.
    Let $p_0 = (1,1)$.
    Let $h = \converse{g} \union ((\naturalsPlus\setminus\ran{g})\times\{p_0\})$.
    Show that $h\in\funs{\naturalsPlus}{\naturalsPlus\times\naturalsPlus}$.
    \begin{subproof}
        Show that for all $m,y_1,y_2$ such that $(m,y_1)\in h$ and $(m,y_2)\in h$ we have $y_1 = y_2$.
        \begin{subproof}
            Fix $m$.
            Fix $y_1$.
            Fix $y_2$ such that $(m,y_1)\in h$ and $(m,y_2)\in h$.
            \begin{byCase}
            \caseOf{$m\in\ran{g}$.}
                We have $(m,y_1)\in\converse{g}$ and $(m,y_2)\in\converse{g}$
                    by \cref{union_iff,times_tuple_elim,setminus_elim_right}.
                $y_1\mathrel{g} m$ and $y_2\mathrel{g} m$ by \cref{converse_iff}.
                Thus $g(y_1) = m$ and $g(y_2) = m$ by \cref{inj,funs_is_function,function_apply_intro}.
                Therefore $y_1 = y_2$ by \cref{dom_iff,funs,subseteq,inj}.
            \caseOf{$m\notin\ran{g}$.}
                We have $(m,y_1)\in(\naturalsPlus\setminus\ran{g})\times\{p_0\}$ and
                    $(m,y_2)\in(\naturalsPlus\setminus\ran{g})\times\{p_0\}$ by \cref{union_iff,converse_iff,ran_iff}.
                Therefore $y_1 = y_2$ by \cref{times_tuple_elim,singleton_iff}.
            \end{byCase}
        \end{subproof}
        $h$ is a function by \cref{union_iff,converse_relation,times_elem_is_tuple,relation,rightunique}.
        Show that $\dom{h} = \naturalsPlus$.
        \begin{subproof}
            Show that for all $m\in\naturalsPlus$ we have $m\in\dom{h}$.
            \begin{subproof}
                Fix $m\in\naturalsPlus$.
                \begin{byCase}
                \caseOf{$m\in\ran{g}$.}
                    There exists $p$ such that $p\mathrel{g} m$ by \cref{ran_iff}.
                    Thus $m\mathrel{\converse{g}} p$ by \cref{converse_iff}.
                    Therefore $m\in\dom{h}$ by \cref{union_iff,dom_intro}.
                \caseOf{$m\notin\ran{g}$.}
                    Hence $m\in\dom{h}$ by \cref{setminus_intro,singleton_intro,times_tuple_intro,union_iff,dom_intro}.
                \end{byCase}
            \end{subproof}
            $\naturalsPlus\subseteq\dom{h}$ by \cref{subseteq}.
            Show that for all $m\in\dom{h}$ we have $m\in\naturalsPlus$.
            \begin{subproof}
                Fix $m\in\dom{h}$.
                Take $y$ such that $(m,y)\in h$ by \cref{dom_iff}.
                $(m,y)\in\converse{g}$ or $(m,y)\in(\naturalsPlus\setminus\ran{g})\times\{p_0\}$
                    by \cref{pair_eq_iff,union_iff}.
                \begin{byCase}
                \caseOf{$(m,y)\in\converse{g}$.}
                    Thus $y\mathrel{g} m$ by \cref{converse_iff}.
                    $y\in\dom{g}$ and $g(y) = m$ by \cref{function_apply_iff,inj,funs_is_function,funs_elim}.
                    $g\in\funs{\naturalsPlus\times\naturalsPlus}{\naturalsPlus}$ by \cref{inj}.
                    Hence $\dom{g} = \naturalsPlus\times\naturalsPlus$ by \cref{funs}.
                    $y\in\naturalsPlus\times\naturalsPlus$ by \cref{pair_eq_iff}.
                    Therefore $m\in\naturalsPlus$ by \cref{funs_type_apply,pair_eq_iff}.
                \caseOf{$(m,y)\in(\naturalsPlus\setminus\ran{g})\times\{p_0\}$.}
                    Thus $m\in\naturalsPlus$ by \cref{times_tuple_elim,setminus_elim_left}.
                \end{byCase}
            \end{subproof}
            $\dom{h}\subseteq\naturalsPlus$ by \cref{subseteq}.
            Therefore $\dom{h} = \naturalsPlus$ by \cref{subseteq_antisymmetric}.
        \end{subproof}
        Show that for all $m\in\dom{h}$ we have $h(m)\in\naturalsPlus\times\naturalsPlus$.
        \begin{subproof}
            Fix $m\in\dom{h}$.
            Take $y$ such that $(m,y)\in h$ by \cref{dom_iff}.
            We have $h(m) = y$ by \cref{function_apply_intro}.
            $(m,y)\in\converse{g}$ or $(m,y)\in(\naturalsPlus\setminus\ran{g})\times\{p_0\}$
                by \cref{pair_eq_iff,union_iff}.
            \begin{byCase}
            \caseOf{$(m,y)\in\converse{g}$.}
                Thus $y\mathrel{g} m$ by \cref{converse_iff}.
                $y\in\dom{g}$ by \cref{function_apply_iff,inj,funs_is_function,funs_elim}.
                $g\in\funs{\naturalsPlus\times\naturalsPlus}{\naturalsPlus}$ by \cref{inj}.
                Hence $\dom{g} = \naturalsPlus\times\naturalsPlus$ by \cref{funs}.
                Therefore $h(m)\in\naturalsPlus\times\naturalsPlus$ by \cref{pair_eq_iff}.
            \caseOf{$(m,y)\in(\naturalsPlus\setminus\ran{g})\times\{p_0\}$.}
                We have $y = p_0$ by \cref{times_tuple_elim,singleton_iff}.
                $1\in\naturalsPlus$ by \cref{real_one_in_naturals}.
                Thus $p_0\in\naturalsPlus\times\naturalsPlus$ by \cref{pair_eq_iff,times_tuple_intro}.
                Therefore $h(m)\in\naturalsPlus\times\naturalsPlus$ by \cref{pair_eq_iff}.
            \end{byCase}
        \end{subproof}
        Therefore $h\in\funs{\naturalsPlus}{\naturalsPlus\times\naturalsPlus}$ by \cref{funs_intro}.
    \end{subproof}
    $h$ is a function and $\dom{h} = \naturalsPlus$ by \cref{funs_is_function,funs}.
    For all $m\in\dom{h}$ we have $h(m)\in\naturalsPlus\times\naturalsPlus$ by \cref{pair_eq_iff,funs_type_apply}.

    Let $b = \{(m,\apply{\apply{s}{\fst{h(m)}}}{\snd{h(m)}}) \mid m\in\naturalsPlus\}$.
    $b$ is a function on $\naturalsPlus$
        by \cref{relation,pair_eq_iff,rightunique,subseteq,dom_intro,dom_iff,subseteq_antisymmetric}.

    For all $m\in\naturalsPlus$ we have $\apply{b}{m}$ is locally finite in $X$ with respect to $T$
        by \cref{pair_eq_iff,fst_type,snd_type,choice_function,function_apply_intro}.

    Show that for all $D$ such that $D\in\unions{\img{b}{\naturalsPlus}}$ we have $D\in\unions{\img{c}{\naturalsPlus}}$.
    \begin{subproof}
        Fix $D$ such that $D\in\unions{\img{b}{\naturalsPlus}}$.
        Take $B_0\in\img{b}{\naturalsPlus}$ such that $D\in B_0$ by \cref{unions_iff}.
        Take $m\in\naturalsPlus$ such that $m\mathrel{b} B_0$ by \cref{img_iff}.
        $B_0 = b(m)$ by \cref{function_apply_iff}.
        Let $n = \fst{h(m)}$.
        Let $k = \snd{h(m)}$.
        $h(m)\in\naturalsPlus\times\naturalsPlus$ by \cref{funs_type_apply}.
        Thus $n\in\naturalsPlus$ and $k\in\naturalsPlus$ by \cref{fst_type,snd_type}.
        We have $\apply{s}{n}$ is a function on $\naturalsPlus$ and
            $\apply{c}{n} = \unions{\img{\apply{s}{n}}{\naturalsPlus}}$
            by \cref{choice_function,pair_eq_iff}.
        Thus $(m,\apply{\apply{s}{n}}{k})\in b$ by \cref{pair_eq_iff}.
        $b(m) = \apply{\apply{s}{n}}{k}$ by \cref{function_apply_intro}.
        $\apply{\apply{s}{n}}{k}\in\img{\apply{s}{n}}{\naturalsPlus}$ by \cref{img_iff,function_apply_elim}.
        Hence $D\in\apply{c}{n}$ by \cref{pair_eq_iff,unions_iff}.
        $\apply{c}{n}\in\img{c}{\naturalsPlus}$ by \cref{img_iff,function_apply_elim}.
        Therefore $D\in\unions{\img{c}{\naturalsPlus}}$ by \cref{unions_iff}.
    \end{subproof}
    $\unions{\img{b}{\naturalsPlus}}\subseteq\unions{\img{c}{\naturalsPlus}}$ by \cref{subseteq}.

    Show that for all $D$ such that $D\in\unions{\img{c}{\naturalsPlus}}$ we have $D\in\unions{\img{b}{\naturalsPlus}}$.
    \begin{subproof}
        Fix $D$ such that $D\in\unions{\img{c}{\naturalsPlus}}$.
        Take $C_0\in\img{c}{\naturalsPlus}$ such that $D\in C_0$ by \cref{unions_iff}.
        Take $n\in\naturalsPlus$ such that $n\mathrel{c} C_0$ by \cref{img_iff}.
        $C_0 = c(n)$ by \cref{function_apply_iff}.
        We have $\apply{s}{n}$ is a function on $\naturalsPlus$ and
            $\apply{c}{n} = \unions{\img{\apply{s}{n}}{\naturalsPlus}}$
            by \cref{choice_function,pair_eq_iff}.
        Take $A_0$ such that $A_0\in\img{\apply{s}{n}}{\naturalsPlus}$ and $D\in A_0$ by \cref{pair_eq_iff,unions_iff}.
        Take $k\in\naturalsPlus$ such that $k\mathrel{\apply{s}{n}} A_0$ by \cref{img_iff}.
        $A_0 = \apply{\apply{s}{n}}{k}$ by \cref{function_apply_iff}.
        Let $m = \apply{g}{(n,k)}$.
        $g\in\funs{\naturalsPlus\times\naturalsPlus}{\naturalsPlus}$ by \cref{inj}.
        Thus $g$ is a function and $\dom{g} = \naturalsPlus\times\naturalsPlus$ by \cref{funs_is_function,funs}.
        $(n,k)\in\dom{g}$ by \cref{times_tuple_intro,pair_eq_iff}.
        $((n,k), m) \in g$ by \cref{function_apply_elim}.
        Hence $m\in\naturalsPlus$ by \cref{funs_type_apply,times_tuple_intro,pair_eq_iff}.
        Thus $(m,(n,k))\in h$ by \cref{converse_iff,union_iff}.
        $h(m) = (n,k)$ by \cref{function_apply_intro}.
        Hence $\fst{h(m)} = n$ and $\snd{h(m)} = k$ by \cref{fst_eq,snd_eq}.
        Thus $(m,\apply{\apply{s}{n}}{k})\in b$ by \cref{pair_eq_iff}.
        $b(m) = \apply{\apply{s}{n}}{k}$ by \cref{function_apply_intro}.
        We have $b(m)\in\img{b}{\naturalsPlus}$ by \cref{img_iff,function_apply_elim}.
        Therefore $D\in\unions{\img{b}{\naturalsPlus}}$ by \cref{pair_eq_iff,unions_iff}.
    \end{subproof}
    $\unions{\img{c}{\naturalsPlus}}\subseteq\unions{\img{b}{\naturalsPlus}}$ by \cref{subseteq}.
    Hence $\unions{\img{c}{\naturalsPlus}} = \unions{\img{b}{\naturalsPlus}}$ by \cref{subseteq_antisymmetric}.
    $\unions{\img{c}{\naturalsPlus}}\subseteq\pow{X}$
        by \cref{pair_eq_iff,unions_iff,img_iff,function_apply_iff,locally_finite_collection,subseteq}.
    Thus $\unions{\img{c}{\naturalsPlus}}$ is countably locally finite in $X$ with respect to $T$
        by \cref{countably_locally_finite_collection}.
\end{proof}

% from §40 helper lemma 6: values in a common initial interval are uniformly close
\begin{lemma}\label{standard_metric_intervalclosed_bound}
    Suppose $a\in\reals$.
    Suppose $0 \leq a$.
    Suppose $u\in\intervalclosed{0}{a}$.
    Suppose $v\in\intervalclosed{0}{a}$.
    Then $\apply{\standardMetricR}{(u,v)} \leq a$.
\end{lemma}
\begin{proof}
    We have $u\in\reals$ and $v\in\reals$ and
        $0 \leq u$ and $u \leq a$ and $0 \leq v$ and $v \leq a$ by \cref{intervalclosed}.
    $\neg{v}\in\reals$ and $\neg{u}\in\reals$ and
        $u - v\in\reals$ and $v - u\in\reals$ by \cref{real_add_inverse,real_sub,real_add_closed}.
    $a + 0\in\reals$ and $a + v\in\reals$ and $a + u\in\reals$ by \cref{real_add_ident,real_add_closed}.
    Hence $a \leq a + v$ and $a \leq a + u$
        by \cref{real_leq_add,real_add_ident,real_add_comm,real_lt_subst_left,real_lt_subst_right}.
    Thus $u \leq a + v$ and $v \leq a + u$ by \cref{real_leq_trans}.
    $u + \neg{v} \leq (a + v) + \neg{v}$ and $v + \neg{u} \leq (a + u) + \neg{u}$
        by \cref{real_add_closed,real_leq_add}.
    We have $(a + v) + \neg{v} = a$ and $(a + u) + \neg{u} = a$
        by \cref{real_add_assoc,real_add_inverse,real_add_ident,real_add_comm}.
    Show that $u - v \leq a$.
    \begin{subproof}
        Follows by \cref{real_sub,real_leq_trans}.
    \end{subproof}
    Show that $v - u \leq a$.
    \begin{subproof}
        Follows by \cref{real_sub,real_leq_trans}.
    \end{subproof}
    Thus $\abs{u - v} \leq a$ by \cref{real_sub_swap_neg,real_abs_leq_if_pm_leq}.
    Therefore $\apply{\standardMetricR}{(u,v)} \leq a$ by \cref{standard_metric_value,pair_eq_iff}.
\end{proof}

% from §40 helper lemma 7: finite continuous real families admit a common control neighborhood
\begin{lemma}\label{finite_continuous_real_family_common_neighborhood}
    Assume $A$ is finite.
    Assume $h$ is a function on $A$.
    Assume $x_0\in X$.
    Assume $T$ is a topology on $X$.
    Suppose for all $\alpha\in A$ we have $\apply{h}{\alpha}$ is continuous from $X$ to $\reals$ with respect to $T$ and $\standardTopologyR$.
    Suppose $\epsilon\in\reals$.
    Suppose $0 < \epsilon$.
    Then there exists $W$ such that
        $W$ is a neighborhood of $x_0$ in $X$ with respect to $T$ and
        for all $\alpha\in A$ for all $x\in W$ we have
            $\apply{\standardMetricR}{(\apply{\apply{h}{\alpha}}{x},\apply{\apply{h}{\alpha}}{x_0})} < \epsilon$.
\end{lemma}
\begin{proof}
    $\standardMetricR$ is a metric on $\reals$ by \cref{standard_metric_is_metric}.
    Let $T_0 = \metricTopology{\standardMetricR}{\reals}$.
    $T_0 = \standardTopologyR$ by \cref{standard_metric_topology}.
    Thus $T_0$ is a topology on $\reals$
        by \cref{metric_basis_is_basis,basis_topology_is_topology,metric_topology,standard_metric_is_metric}.
    \begin{byCase}
    \caseOf{$A = \emptyset$.}
        Let $W = X$.
        Hence $W$ is a neighborhood of $x_0$ in $X$ with respect to $T$ by \cref{open_in,topology_on,neighborhood}.
        Show that for all $\alpha\in A$ for all $x\in W$ we have
            $\apply{\standardMetricR}{(\apply{\apply{h}{\alpha}}{x},\apply{\apply{h}{\alpha}}{x_0})} < \epsilon$.
        \begin{subproof}
            Follows by \cref{emptyset}.
        \end{subproof}
    \caseOf{$A\neq\emptyset$.}
        Let $R = \{w\in A\times\pow{X} \mid \text{there exists $\alpha\in A$ such that there exists $U$ such that
            $w = (\alpha,U)$ and
            $U$ is a neighborhood of $x_0$ in $X$ with respect to $T$ and
            $\img{\apply{h}{\alpha}}{U}\subseteq \metricBall{\standardMetricR}{\reals}{\apply{\apply{h}{\alpha}}{x_0}}{\epsilon}$}\}$.
        $R$ is a relation by \cref{relation,subseteq,times_elem_is_tuple}.

        Show that for all $\alpha\in A$ there exists $U$ such that $\alpha\mathrel{R} U$.
        \begin{subproof}
            Fix $\alpha\in A$.
            We have $\apply{h}{\alpha}$ is continuous from $X$ to $\reals$ with respect to $T$ and $\standardTopologyR$ by assumption.
            Thus $\apply{\apply{h}{\alpha}}{x_0}\in\reals$ by \cref{continuous,funs_type_apply}.
            Let $V = \metricBall{\standardMetricR}{\reals}{\apply{\apply{h}{\alpha}}{x_0}}{\epsilon}$.
            $V$ is open in $\reals$ with respect to $T_0$
                by \cref{metric_ball,subseteq,powerset_intro,metric_basis,metric_basis_is_basis,basis_elem_open_in_generated,metric_topology,basis_for_topology,open_in}.
            Thus $V$ is a neighborhood of $\apply{\apply{h}{\alpha}}{x_0}$ in $\reals$ with respect to $T_0$
                by \cref{metric_on,metric_zero_characterization,times_tuple_intro,metric_ball,real_lt_subst_left,neighborhood,open_in}.
            Hence $\apply{h}{\alpha}$ is continuous at $x_0$ from $X$ to $\reals$ with respect to $T$ and $T_0$
                by \cref{continuous_implies_continuous_at,pair_eq_iff}.
            Take $U$ such that
                $U$ is a neighborhood of $x_0$ in $X$ with respect to $T$ and
                $\img{\apply{h}{\alpha}}{U}\subseteq V$
                by \cref{continuous_at,standard_metric_topology,pair_eq_iff}.
            Therefore $\alpha\mathrel{R} U$
                by \cref{times_tuple_intro,powerset_intro,neighborhood,open_in,topology_on,subseteq,pair_eq_iff}.
        \end{subproof}

        Take $u$ such that $u$ is a function on $A$ and for all $\alpha\in A$ we have $\alpha\mathrel{R}\apply{u}{\alpha}$
            by \cref{choice_function}.
        Let $M = \img{u}{A}$.
        $M$ is finite by \cref{finite_image_function}.
        Take $\alpha_0$ such that $\alpha_0\in A$ by \cref{nonempty_inhabited}.
        $M$ is inhabited by \cref{img_iff,function_apply_elim,nonempty_inhabited,emptyset}.

        $M\subseteq T$
            by \cref{img_iff,function_apply_iff,function_apply_elim,relation,fst_eq,snd_eq,neighborhood,open_in,pair_eq_iff,subseteq}.

        Let $W = \inters{M}$.
        Thus $W\in T$ by \cref{finite_intersections_open_in_topology}.
        For all $U$ such that $U\in M$ we have $x_0\in U$
            by \cref{img_iff,function_apply_iff,function_apply_elim,relation,fst_eq,snd_eq,neighborhood,pair_eq_iff}.
        Hence $W$ is a neighborhood of $x_0$ in $X$ with respect to $T$ by \cref{inters_iff_forall,neighborhood,open_in}.

        Show that for all $\alpha\in A$ for all $x\in W$ we have
            $\apply{\standardMetricR}{(\apply{\apply{h}{\alpha}}{x},\apply{\apply{h}{\alpha}}{x_0})} < \epsilon$.
        \begin{subproof}
            Fix $\alpha\in A$.
            Fix $x\in W$.
            Thus $x\in\apply{u}{\alpha}$
                by \cref{img_iff,function_apply_elim,subseteq,open_in,topology_on,powerset_intro,inters_subseteq_elem}.
            $\alpha\mathrel{R}\apply{u}{\alpha}$ by \cref{function_apply_elim}.
            Therefore $(\alpha,\apply{u}{\alpha})\in R$ by \cref{relation}.
            Hence $\img{\apply{h}{\alpha}}{\apply{u}{\alpha}}\subseteq
                \metricBall{\standardMetricR}{\reals}{\apply{\apply{h}{\alpha}}{x_0}}{\epsilon}$ by \cref{fst_eq,snd_eq}.
            $x\in X$ by \cref{neighborhood,open_in,topology_on,subseteq,powerset_elim}.
            $\apply{h}{\alpha}\in\funs{X}{\reals}$ by \cref{continuous}.
            Thus $\apply{\apply{h}{\alpha}}{x}\in\img{\apply{h}{\alpha}}{\apply{u}{\alpha}}$
                by \cref{funs_elim,function_apply_elim,img_iff,subseteq}.
            Hence $\apply{\apply{h}{\alpha}}{x}\in
                \metricBall{\standardMetricR}{\reals}{\apply{\apply{h}{\alpha}}{x_0}}{\epsilon}$ by \cref{subseteq}.
            $\apply{\apply{h}{\alpha}}{x}\in\reals$ and $\apply{\apply{h}{\alpha}}{x_0}\in\reals$ by \cref{continuous,funs_type_apply}.
            Thus $\apply{\standardMetricR}{(\apply{\apply{h}{\alpha}}{x_0},\apply{\apply{h}{\alpha}}{x})} < \epsilon$
                by \cref{metric_ball}.
            $(\apply{\apply{h}{\alpha}}{x},\apply{\apply{h}{\alpha}}{x_0})\in\reals\times\reals$ and
                $(\apply{\apply{h}{\alpha}}{x_0},\apply{\apply{h}{\alpha}}{x})\in\reals\times\reals$
                by \cref{times_tuple_intro}.
            Hence $\apply{\standardMetricR}{(\apply{\apply{h}{\alpha}}{x},\apply{\apply{h}{\alpha}}{x_0})}
                = \apply{\standardMetricR}{(\apply{\apply{h}{\alpha}}{x_0},\apply{\apply{h}{\alpha}}{x})}$
                by \cref{metric_on,metric_symmetric}.
            Therefore $\apply{\standardMetricR}{(\apply{\apply{h}{\alpha}}{x},\apply{\apply{h}{\alpha}}{x_0})} < \epsilon$
                by \cref{real_lt_subst_left}.
        \end{subproof}
    \end{byCase}
\end{proof}

% from §40 helper lemma 8: Nagata-Smirnov coordinate family
\begin{lemma}\label{nagata_smirnov_coordinate_family}
    Assume $X$ is regular with respect to $T$.
    Assume $B$ is a basis for $T$ on $X$.
    Assume $B$ is countably locally finite in $X$ with respect to $T$.
    Assume $b$ is a function on $\naturalsPlus$.
    Assume $B = \unions{\img{b}{\naturalsPlus}}$.
    Suppose for all $n\in\naturalsPlus$ we have $\apply{b}{n}$ is locally finite in $X$ with respect to $T$.
    Let $J = \nagataSmirnovIndex{b}{B}$.
    Then there exists $h$ such that
        $h$ is a function on $J$ and
        for all $p\in J$ we have $h(p)$ is continuous from $X$ to $\reals$ with respect to $T$ and $\standardTopologyR$ and
        for all $p,n$ such that $p\in J$ and $n = \fst{p}$ we have $\img{h(p)}{X}\subseteq \intervalclosed{0}{1 / n}$ and
        for all $p\in J$ we have $\img{h(p)}{X\setminus\snd{p}}\subseteq \{0\}$ and
        for all $p\in J$ for all $x\in\snd{p}$ we have $0 < \apply{\apply{h}{p}}{x}$.
\end{lemma}
\begin{proof}
    $T$ is a topology on $X$ by \cref{regular_space}.
    $X$ is normal with respect to $T$ by \cref{regular_countably_locally_finite_basis_normal}.
    Let $R = \{w\in J\times\pow{X\times\reals} \mid \text{there exists $p\in J$ such that there exists $g$ such that
        $w = (p,g)$ and
        $g$ is continuous from $X$ to $\reals$ with respect to $T$ and $\standardTopologyR$ and
        there exists $n\in\naturalsPlus$ such that
        $n = \fst{p}$ and
        $\img{g}{X}\subseteq \intervalclosed{0}{1 / n}$ and
        $\img{g}{X\setminus\snd{p}}\subseteq \{0\}$ and
        for all $x\in\snd{p}$ we have $0 < g(x)$}\}$.
    $R$ is a relation by \cref{relation,subseteq,times_elem_is_tuple}.

    Show that for all $p\in J$ there exists $g$ such that $p\mathrel{R} g$.
    \begin{subproof}
        Fix $p\in J$.
        $p\in\naturalsPlus\times B$ and $\snd{p}\in\apply{b}{\fst{p}}$ by \cref{nagata_smirnov_index_set}.
        Thus $\fst{p}\in\naturalsPlus$ and $\snd{p}\in B$ by \cref{fst_type,snd_type}.
        Let $n = \fst{p}$.
        $\snd{p}$ is open in $X$ with respect to $T$ by \cref{basis_for_topology,basis_elem_open_in_generated,open_in}.
        We have $\snd{p}\subseteq X$ by \cref{open_in,topology_on,subseteq,powerset_elim}.
        Let $A = X\setminus\snd{p}$.
        $A$ is closed in $X$ with respect to $T$
            by \cref{open_in,topology_on,subseteq,powerset_elim,setminus_subseteq,double_relative_complement,closed_in}.
        $A$ is a Gdelta set in $X$ with respect to $T$
            by \cref{regular_countably_locally_finite_basis_closed_g_delta}.
        Let $I_0 = \intervalclosed{0}{1}$.
        Take $q$ such that
            $q$ is continuous from $X$ to $I_0$ with respect to $T$ and $\subspaceTopology{\standardTopologyR}{I_0}$ and
            $\img{q}{A}\subseteq \{0\}$ and
            for all $x\in X\setminus A$ we have $0 < q(x)$
            by \cref{normal_closed_g_delta_positive_function}.
        Let $j_0 = \identity{I_0}$.
        Let $q_0 = j_0\circ q$.
        $\standardTopologyR$ is a topology on $\reals$
            by \cref{standard_basis_is_basis,basis_topology_is_topology,standard_topology_reals}.
        Thus $I_0$ is a subspace of $\reals$ with respect to $\standardTopologyR$ by \cref{intervalclosed,subseteq,subspace_of}.
        $q_0$ is continuous from $X$ to $\reals$ with respect to $T$ and $\standardTopologyR$
            by \cref{continuous_expand_range}.
        $q\in\funs{X}{I_0}$ by \cref{continuous}.
        Hence $\ran{q}\subseteq\dom{j_0}$ by \cref{funs_ran,subseteq,id_dom}.
        For all $x\in X$ we have $q_0(x) = q(x)$ by \cref{circ_apply,funs_elim,id_apply,id_rightunique,id_is_relation}.

        $1 / n\in\reals$ by \cref{naturalsplus_reciprocal_real}.
        $0 < 1 / n$ by \cref{positive_reciprocal_naturalsplus}.

        Let $c = \{(x,1 / n) \mid x\in X\}$.
        $c$ is a function from $X$ to $\reals$ by \cref{constant_map_function}.
        For all $x\in X$ we have $c(x) = 1 / n$
            by \cref{constant_map_function,funs_is_function,function_apply_intro}.
        $c$ is continuous from $X$ to $\reals$ with respect to $T$ and $\standardTopologyR$
            by \cref{constant_continuous}.

        Let $g = \{(x, c(x) \rmul q_0(x)) \mid x\in X\}$.
        $g$ is continuous from $X$ to $\reals$ with respect to $T$ and $\standardTopologyR$
            by \cref{real_function_multiplication_continuous}.
        $g\in\funs{X}{\reals}$ and $g$ is a function by \cref{continuous,funs_is_function}.

        For all $x\in X$ we have $g(x)\in\intervalclosed{0}{1 / n}$
            by \cref{circ_apply,funs_elim,id_apply,intervalclosed,function_apply_iff,real_lt_trichotomy,real_mul_zero,real_mul_comm,real_lt_imp_pos_diff,real_mul_pos_pos_gt,real_lt_imp_leq,real_mul_ident,real_naturals_subset_reals,real_one_in_naturals,subseteq,real_lt_mul_pos,continuous}.
        Hence $\img{g}{X}\subseteq \intervalclosed{0}{1 / n}$
            by \cref{subseteq,img_iff,continuous,funs_is_function,function_apply_iff}.

        Show that for all $x\in\snd{p}$ we have $0 < g(x)$.
        \begin{subproof}
            Follows by \cref{subseteq,pair_eq_iff,setminus_elim_right,setminus_intro,circ_apply,funs_elim,id_apply,function_apply_iff,continuous,funs_type_apply,positive_reciprocal_naturalsplus,real_mul_pos_pos_gt}.
        \end{subproof}

        Show that for all $x$ such that $x\in X\setminus\snd{p}$ we have $g(x) = 0$.
        \begin{subproof}
            Follows by \cref{pair_eq_iff,setminus_elim_left,continuous,funs_is_function,function_apply_elim,img_iff,subseteq,singleton_elim,circ_apply,funs_elim,id_apply,function_apply_iff,real_mul_comm,real_mul_zero}.
        \end{subproof}
        Thus $\img{g}{X\setminus\snd{p}}\subseteq \{0\}$
            by \cref{img_iff,funs_is_function,function_apply_iff,singleton_intro,subseteq}.

        Therefore $p\mathrel{R} g$
            by \cref{funs_subseteq_rels,funs,rels_elim,subseteq,times_elem_is_tuple,powerset_intro,times_tuple_intro,pair_eq_iff}.
    \end{subproof}

    Take $h$ such that
        $h$ is a function on $J$ and
        for all $p\in J$ we have $p\mathrel{R} h(p)$
        by \cref{choice_function}.
    Show that for all $p\in J$ we have $h(p)$ is continuous from $X$ to $\reals$ with respect to $T$ and $\standardTopologyR$.
    \begin{subproof}
        Follows by \cref{function_apply_elim,pair_eq_iff}.
    \end{subproof}
    Show that for all $p,n$ such that $p\in J$ and $n = \fst{p}$ we have $\img{h(p)}{X}\subseteq \intervalclosed{0}{1 / n}$.
    \begin{subproof}
        Follows by \cref{function_apply_elim,pair_eq_iff}.
    \end{subproof}
    Show that for all $p\in J$ we have $\img{h(p)}{X\setminus\snd{p}}\subseteq \{0\}$.
    \begin{subproof}
        Follows by \cref{function_apply_elim,pair_eq_iff}.
    \end{subproof}
    Show that for all $p\in J$ for all $x\in\snd{p}$ we have $0 < \apply{\apply{h}{p}}{x}$.
    \begin{subproof}
        Follows by \cref{function_apply_elim,pair_eq_iff}.
    \end{subproof}
\end{proof}

% from §40 helper lemma 1: Nagata-Smirnov backward implication
\begin{lemma}\label{nagata_smirnov_metrization_backward}
    If $X$ is regular with respect to $T$ and there exists $B$ such that
        $B$ is a basis for $T$ on $X$ and
        $B$ is countably locally finite in $X$ with respect to $T$
    then $X$ is metrizable with respect to $T$.
\end{lemma}
\begin{proof}
    Assume $X$ is regular with respect to $T$ and there exists $B$ such that
            $B$ is a basis for $T$ on $X$ and
            $B$ is countably locally finite in $X$ with respect to $T$.
        $T$ is a topology on $X$ by \cref{regular_space}.
        Take $G$ such that
            $G$ is a basis for $T$ on $X$ and
            $G$ is countably locally finite in $X$ with respect to $T$.
        \begin{byCase}
        \caseOf{$X = \emptyset$.}
            $\standardBasisR$ is a basis on $\reals$ by \cref{standard_basis_is_basis}.
            $\basisTopology{\standardBasisR}{\reals}$ is a topology on $\reals$
                by \cref{basis_topology_is_topology}.
            Thus $\standardTopologyR$ is a topology on $\reals$ by \cref{standard_topology_reals}.
            Let $f = \emptyset$.
            $f$ is a topological embedding from $X$ to $\reals$ with respect to $T$ and $\standardTopologyR$
                by \cref{empty_domain_topological_embedding}.
            $\standardMetricR$ is a metric on $\reals$ by \cref{standard_metric_is_metric}.
            Hence $\reals$ is metrizable with respect to $\standardTopologyR$
                by \cref{metrizable,standard_metric_topology,standard_topology_reals}.
            $f$ is a function from $X$ to $\reals$ by \cref{topological_embedding}.
            $f\in\rels{X}{\reals}$ and $\ran{f}\subseteq\reals$
                by \cref{funs_subseteq_rels,subseteq,rels_ran_subseteq}.
            Thus $\img{f}{X}\subseteq\reals$ by \cref{img_subseteq_ran,subseteq_transitive}.
            Hence $\img{f}{X}$ is a subspace of $\reals$ with respect to $\standardTopologyR$ by \cref{subspace_of}.
            $\subspaceTopology{\standardTopologyR}{\img{f}{X}}$ is a topology on $\img{f}{X}$
                by \cref{subspace_topology_is_topology,subspace_of}.
            Thus $\img{f}{X}$ is metrizable with respect to $\subspaceTopology{\standardTopologyR}{\img{f}{X}}$
                by \cref{subspace_metrizable}.
            Hence $f$ is a homeomorphism between $X$ and $\img{f}{X}$ with respect to
                $T$ and $\subspaceTopology{\standardTopologyR}{\img{f}{X}}$
                by \cref{topological_embedding}.
            Therefore $X$ is metrizable with respect to $T$ by \cref{homeomorphism_preserves_metrizable,subspace_of}.
        \caseOf{$X\neq\emptyset$.}
            Take $b$ such that
                $b$ is a function on $\naturalsPlus$ and
                $G = \unions{\img{b}{\naturalsPlus}}$ and
                for all $n\in\naturalsPlus$ we have $\apply{b}{n}$ is locally finite in $X$ with respect to $T$
                by \cref{countably_locally_finite_collection}.
            Let $J = \nagataSmirnovIndex{b}{G}$.
            Take $f$ such that
                $f$ is a function on $J$ and
                for all $p\in J$ we have $\apply{f}{p}$ is continuous from $X$ to $\reals$ with respect to $T$ and $\standardTopologyR$ and
                for all $p,n$ such that $p\in J$ and $n = \fst{p}$ we have $\img{\apply{f}{p}}{X}\subseteq \intervalclosed{0}{1 / n}$ and
                for all $p\in J$ we have $\img{\apply{f}{p}}{X\setminus \snd{p}}\subseteq \{0\}$ and
                for all $p\in J$ for all $x\in\snd{p}$ we have $0 < \apply{\apply{f}{p}}{x}$
                by \cref{nagata_smirnov_coordinate_family}.

            Take $x_0$ such that $x_0\in X$ by \cref{nonempty_inhabited}.
            Take $C_0\in G$ such that $x_0\in C_0$ by \cref{basis_for_topology,basis_on}.
            Take $D_0$ such that $D_0\in\img{b}{\naturalsPlus}$ and $C_0\in D_0$
                by \cref{countably_locally_finite_collection,unions_iff}.
            Take $n_0\in\naturalsPlus$ such that $n_0\mathrel{b} D_0$ by \cref{img_iff}.
            Hence $C_0\in\apply{b}{n_0}$ by \cref{function_apply_iff,subseteq}.
            Let $\alpha_0 = (n_0,C_0)$.
            $\fst{\alpha_0} = n_0$ and $\snd{\alpha_0} = C_0$ by \cref{fst_eq,snd_eq}.
            $\alpha_0\in\naturalsPlus\times G$ by \cref{times_tuple_intro}.
            Thus $\snd{\alpha_0}\in\apply{b}{\fst{\alpha_0}}$ by \cref{subseteq,pair_eq_iff}.
            Hence $\alpha_0\in J$ by \cref{nagata_smirnov_index_set}.
            Therefore $J$ is inhabited.

            Let $R = \{(\alpha,\reals) \mid \alpha\in J\}$.
            Let $S = \{(\alpha,\standardTopologyR) \mid \alpha\in J\}$.
            Let $P = \productTopologyIndexed{S}{R}$.
            Let $F = \{(x,z) \mid x\in X, z\in\productFamily{R} \mid
                \forall \alpha\in J. \apply{z}{\alpha} = \apply{\apply{f}{\alpha}}{x}\}$.

            Show that for all $x,U$ such that $x\in X$ and $U$ is a neighborhood of $x$ in $X$ with respect to $T$
                there exists $\alpha\in J$ such that
                    $0 < \apply{\apply{f}{\alpha}}{x}$ and
                    for all $y\in X\setminus U$ we have $\apply{\apply{f}{\alpha}}{y} = 0$.
            \begin{subproof}
                Fix $x$.
                Fix $U$ such that $x\in X$ and $U$ is a neighborhood of $x$ in $X$ with respect to $T$.
                $x\in U$ by \cref{neighborhood}.
                Take $C\in G$ such that $x\in C$ and $C\subseteq U$
                    by \cref{neighborhood,open_in,basis_for_topology,topology_generated_by_basis}.
                Take $D\in\img{b}{\naturalsPlus}$ such that $C\in D$
                    by \cref{countably_locally_finite_collection,unions_iff}.
                Take $n\in\naturalsPlus$ such that $n\mathrel{b} D$ by \cref{img_iff}.
                Let $\alpha = (n,C)$.
                We have $C\in\apply{b}{n}$ by \cref{function_apply_iff,subseteq}.
                $\fst{\alpha} = n$ and $\snd{\alpha} = C$ by \cref{fst_eq,snd_eq}.
                $\alpha\in\naturalsPlus\times G$ by \cref{times_tuple_intro}.
                Thus $\snd{\alpha}\in\apply{b}{\fst{\alpha}}$ by \cref{subseteq,pair_eq_iff}.
                Hence $\alpha\in J$ by \cref{nagata_smirnov_index_set}.
                Show that $0 < \apply{\apply{f}{\alpha}}{x}$.
                \begin{subproof}
                    Follows by \cref{subseteq}.
                \end{subproof}
                Show that for all $y\in X\setminus U$ we have $\apply{\apply{f}{\alpha}}{y} = 0$.
                \begin{subproof}
                    Follows by \cref{setminus_elim_left,setminus_elim_right,subseteq,setminus_intro,continuous,funs,funs_is_function,inter_intro,img_of_function_intro,singleton_elim}.
                \end{subproof}
            \end{subproof}

            Thus $F$ is a topological embedding from $X$ to $\productFamily{R}$ with respect to $T$ and $P$
                by \cref{urysohn_embedding_theorem_real_power,regular_space}.
            Hence $F$ is a function from $X$ to $\productFamily{R}$ and
                $F$ is injective and
                $F$ is continuous from $X$ to $\productFamily{R}$ with respect to $T$ and $P$ and
                $F$ is a homeomorphism between $X$ and $\img{F}{X}$ with respect to
                    $T$ and $\subspaceTopology{P}{\img{F}{X}}$
                by \cref{topological_embedding}.
            We have $\productFamily{R} = \realsPower{J}$ by \cref{product_family_reals_power}.
            Let $U_0 = \metricTopology{\uniformMetric{J}}{\realsPower{J}}$.
            $\uniformMetric{J}$ is a metric on $\realsPower{J}$ by \cref{uniform_metric_is_metric_late}.
            Thus $U_0$ is a topology on $\realsPower{J}$
                by \cref{metric_topology,metric_basis_is_basis,basis_topology_is_topology}.

            Show that for all $x\in X$ we have $F$ is continuous at $x$ from $X$ to $\realsPower{J}$ with respect to $T$ and $U_0$.
            \begin{subproof}
                    Fix $x\in X$.
                    Show that for all $V$ such that $V$ is a neighborhood of $F(x)$ in $\realsPower{J}$ with respect to $U_0$
                        there exists $W$ such that
                            $W$ is a neighborhood of $x$ in $X$ with respect to $T$ and
                            $\img{F}{W}\subseteq V$.
                    \begin{subproof}
                        Fix $V$ such that $V$ is a neighborhood of $F(x)$ in $\realsPower{J}$ with respect to $U_0$.
                        Let $d_J = \uniformMetric{J}$.
                        Let $e = \boundedMetric{\standardMetricR}$.
                        $\metricTopology{d_J}{\realsPower{J}} = U_0$ by definition.
                        Hence $V\subseteq\realsPower{J}$ by
                            \cref{neighborhood,open_in,metric_basis_is_basis,basis_topology_is_topology,metric_topology,topology_on,subseteq,pow_iff}.
                        Take $\epsilon\in\reals$ such that $0 < \epsilon$ and
                            $\metricBall{d_J}{\realsPower{J}}{F(x)}{\epsilon}\subseteq V$
                            by \cref{neighborhood,metric_open_iff}.
                        \begin{byCase}
                        \caseOf{$\epsilon \leq 1$.}
                            Let $\eta = \epsilon / (1 + 1)$.
                            $\eta\in\reals$ and $0 < \eta$ and $\eta + \eta = \epsilon$
                                by \cref{real_half_of_positive}.
                            Take $N\in\naturalsPlus$ such that $0 < 1 / N$ and $1 / N < \eta$
                                by \cref{real_small_reciprocal}.
                            Let $A = \{w\in\initialSegment{N}\times\pow{X} \mid \text{there exists $n\in\initialSegment{N}$ such that there exists $W$ such that there exists $H$ such that there exists $M$ such that
                                $w = (n,W)$ and
                                $W$ is a neighborhood of $x$ in $X$ with respect to $T$ and
                                $H\subseteq\apply{b}{n}$ and
                                $H$ is finite and
                                $M = \{C\in\apply{b}{n} \mid \img{\apply{f}{(n,C)}}{W}\subseteq \metricBall{\standardMetricR}{\reals}{\apply{\apply{f}{(n,C)}}{x}}{\eta}\}$ and
                                $H\subseteq M$ and
                                for all $C\in\apply{b}{n}$ such that $W$ intersects $C$ we have $C\in H$}\}$.
                            $A$ is a relation by \cref{relation,subseteq,times_elem_is_tuple}.
                            Show that for all $n\in\initialSegment{N}$ there exists $W$ such that $n\mathrel{A} W$.
                            \begin{subproof}
                                Fix $n\in\initialSegment{N}$.
                                $n\in\naturalsPlus$ by \cref{initial_segment_naturalsplus}.
                                Take $U,H$ such that
                                    $U$ is a neighborhood of $x$ in $X$ with respect to $T$ and
                                    $H\subseteq\apply{b}{n}$ and
                                    $H$ is finite and
                                    for all $C\in\apply{b}{n}$ if $U$ intersects $C$ then $C\in H$
                                    by \cref{locally_finite_collection}.
                                Let $Q = \{w\in H\times\pow{X} \mid \text{there exists $C\in H$ such that there exists $W$ such that
                                    $w = (C,W)$ and
                                    $W$ is a neighborhood of $x$ in $X$ with respect to $T$ and
                                    $\img{\apply{f}{(n,C)}}{W}\subseteq \metricBall{\standardMetricR}{\reals}{\apply{\apply{f}{(n,C)}}{x}}{\eta}$}\}$.
                                $Q$ is a relation by \cref{relation,subseteq,times_elem_is_tuple}.
                                Show that for all $C\in H$ there exists $W$ such that $C\mathrel{Q} W$.
                                \begin{subproof}
                                    Fix $C\in H$.
                                    We have $C\in\apply{b}{n}$ by \cref{subseteq}.
                                    $C\in G$ by \cref{subseteq,img_iff,function_apply_elim,unions_iff}.
                                    Let $p = (n,C)$.
                                    $\fst{p} = n$ and $\snd{p} = C$ by \cref{fst_eq,snd_eq}.
                                    $p\in\naturalsPlus\times G$ by \cref{times_tuple_intro}.
                                    Thus $\snd{p}\in\apply{b}{\fst{p}}$ by \cref{subseteq,pair_eq_iff}.
                                    Hence $p\in J$ by \cref{nagata_smirnov_index_set}.
                                    Let $g = \apply{f}{(n,C)}$.
                                    Hence $g$ is continuous at $x$ from $X$ to $\reals$ with respect to $T$ and $\standardTopologyR$
                                        by \cref{subseteq,pair_eq_iff,continuous_implies_continuous_at}.
                                    Let $V_0 = \metricBall{\standardMetricR}{\reals}{g(x)}{\eta}$.
                                    $g(x)\in\reals$ by \cref{continuous,funs_type_apply}.
                                    We have $V_0\in\metricBasis{\standardMetricR}{\reals}$
                                        by \cref{metric_ball,subseteq,powerset_intro,metric_basis}.
                                    $V_0\in\standardTopologyR$
                                        by \cref{standard_metric_is_metric,metric_basis_is_basis,basis_elem_open_in_generated,metric_topology,standard_metric_topology}.
                                    $\standardTopologyR$ is a topology on $\reals$
                                        by \cref{standard_metric_is_metric,metric_basis_is_basis,basis_topology_is_topology,metric_topology,standard_metric_topology}.
                                    Thus $V_0$ is open in $\reals$ with respect to $\standardTopologyR$ by \cref{open_in}.
                                    $\apply{\standardMetricR}{(g(x),g(x))} = 0$
                                        by \cref{standard_metric_is_metric,metric_on,metric_zero_characterization,times_tuple_intro}.
                                    Hence $V_0$ is a neighborhood of $g(x)$ in $\reals$ with respect to $\standardTopologyR$ by \cref{metric_ball,real_lt_subst_left,neighborhood}.
                                    Take $W$ such that
                                        $W$ is a neighborhood of $x$ in $X$ with respect to $T$ and
                                        $\img{g}{W}\subseteq V_0$
                                        by \cref{continuous_at}.
                                    $(C,W)\in H\times\pow{X}$ by \cref{times_tuple_intro,powerset_intro,neighborhood,open_in,topology_on,subseteq,pow_iff}.
                                    Thus $(C,W)\in Q$ by \cref{pair_eq_iff}.
                                    Therefore $C\mathrel{Q} W$ by \cref{pair_eq_iff}.
                                \end{subproof}
                                Take $s$ such that
                                    $s$ is a function on $H$ and
                                    for all $C\in H$ we have $C\mathrel{Q}\apply{s}{C}$
                                    by \cref{choice_function}.
                                Let $L = \img{s}{H}\union\{U\}$.
                                $\img{s}{H}$ is finite by \cref{choice_function,finite_image_function}.
                                $\{U\}$ is finite by \cref{singleton_subset_intro,pair_is_finite,upair_intro_left,subseteq_finite_is_finite}.
                                Thus $L$ is finite by \cref{union_of_finite_is_finite}.
                                $L$ is inhabited by \cref{union_intro_right,singleton_intro}.
                                For all $W$ such that $W\in L$ we have $W\in T$
                                    by \cref{union_iff,singleton_iff,img_iff,function_apply_iff,choice_function,fst_eq,snd_eq,pair_eq_iff,neighborhood,open_in}.
                                $L\subseteq T$ by \cref{subseteq}.
                                Let $Z = \inters{L}$.
                                Hence $Z\in T$ by \cref{finite_intersections_open_in_topology}.
                                Show that for all $W$ such that $W\in L$ we have $x\in W$.
                                \begin{subproof}
                                    Fix $W$ such that $W\in L$.
                                    \begin{byCase}
                                    \caseOf{$W\in\img{s}{H}$.}
                                        Take $C\in H$ such that $(C,W)\in s$ by \cref{img_iff}.
                                        $C\in\dom{s}$ and $W = \apply{s}{C}$
                                            by \cref{choice_function,function_apply_iff}.
                                        Thus $(C,\apply{s}{C})\in Q$ by \cref{choice_function}.
                                        We have $\apply{s}{C}$ is a neighborhood of $x$ in $X$ with respect to $T$
                                            by \cref{fst_eq,snd_eq,pair_eq_iff}.
                                        Hence $x\in \apply{s}{C}$ by \cref{neighborhood}.
                                        Therefore $x\in W$ by \cref{pair_eq_iff}.
                                    \caseOf{$W\in\{U\}$.}
                                        We have $W = U$ by \cref{singleton_iff}.
                                        $U$ is a neighborhood of $x$ in $X$ with respect to $T$
                                            by \cref{neighborhood}.
                                        Therefore $x\in W$ by \cref{neighborhood}.
                                    \end{byCase}
                                \end{subproof}
                                Hence $x\in Z$ by \cref{inters_iff_forall,union_intro_right,singleton_intro}.
                                Therefore $Z$ is a neighborhood of $x$ in $X$ with respect to $T$ by \cref{neighborhood,open_in}.
                                $U\in L$ by \cref{union_iff,singleton_iff}.
                                For all $z$ such that $z\in Z$ we have $z\in U$ by \cref{inters_iff_forall}.
                                $Z\subseteq U$ by \cref{subseteq}.
                                Hence for all $C\in\apply{b}{n}$ such that $Z$ intersects $C$ we have $C\in H$
                                    by \cref{subseteq,intersects,nonempty_inhabited,inter_elim_left,inter_elim_right,inter_intro,emptyset,locally_finite_collection}.
                                Show that for all $C\in H$ for all $y\in Z$ we have
                                    $\apply{\standardMetricR}{(\apply{\apply{f}{(n,C)}}{x},\apply{\apply{f}{(n,C)}}{y})} < \eta$.
                                \begin{subproof}
                                    Fix $C\in H$.
                                    Fix $y\in Z$.
                                    $(C,\apply{s}{C})\in Q$ by \cref{choice_function}.
                                    We have $\apply{s}{C}$ is a neighborhood of $x$ in $X$ with respect to $T$ and
                                        $\img{\apply{f}{(n,C)}}{\apply{s}{C}}\subseteq \metricBall{\standardMetricR}{\reals}{\apply{\apply{f}{(n,C)}}{x}}{\eta}$
                                        by \cref{fst_eq,snd_eq,pair_eq_iff}.
                                    Thus $\apply{s}{C}\in\img{s}{H}$ by \cref{img_iff,function_apply_elim}.
                                    $\apply{s}{C}\in L$ by \cref{union_iff}.
                                    For all $z$ such that $z\in Z$ we have $z\in\apply{s}{C}$ by \cref{inters_iff_forall}.
                                    $Z\subseteq \apply{s}{C}$ by \cref{subseteq}.
                                    Hence $y\in\apply{s}{C}$ by \cref{subseteq}.
                                    $C\in\apply{b}{n}$ by \cref{subseteq}.
                                    We have $C\in G$ by \cref{subseteq,img_iff,function_apply_elim,unions_iff}.
                                    Let $p = (n,C)$.
                                    $\fst{p} = n$ and $\snd{p} = C$ by \cref{fst_eq,snd_eq}.
                                    $p\in\naturalsPlus\times G$ by \cref{times_tuple_intro}.
                                    Thus $\snd{p}\in\apply{b}{\fst{p}}$ by \cref{subseteq,pair_eq_iff}.
                                    Hence $p\in J$ by \cref{nagata_smirnov_index_set}.
                                    $\apply{f}{(n,C)}\in\funs{X}{\reals}$ and $\apply{f}{(n,C)}$ is a function and
                                        $\dom{\apply{f}{(n,C)}} = X$
                                        by \cref{continuous,funs,funs_is_function,pair_eq_iff}.
                                    $\apply{s}{C}\subseteq X$ by \cref{neighborhood,open_in,topology_on,subseteq,pow_iff}.
                                    We have $y\in X$ by \cref{subseteq}.
                                    $y\in\dom{\apply{f}{(n,C)}}$ by \cref{subseteq}.
                                    Hence $y\in\dom{\apply{f}{(n,C)}}\inter\apply{s}{C}$ by \cref{inter_intro}.
                                    $\apply{\apply{f}{(n,C)}}{y}\in\img{\apply{f}{(n,C)}}{\apply{s}{C}}$ by \cref{img_of_function_intro}.
                                    Thus $\apply{\apply{f}{(n,C)}}{y}\in\metricBall{\standardMetricR}{\reals}{\apply{\apply{f}{(n,C)}}{x}}{\eta}$ by \cref{subseteq}.
                                    Therefore
                                        $\apply{\standardMetricR}{(\apply{\apply{f}{(n,C)}}{x},\apply{\apply{f}{(n,C)}}{y})} < \eta$
                                        by \cref{metric_ball}.
                                \end{subproof}
                                Show that for all $C\in H$ for all $a$ such that $a\in\img{\apply{f}{(n,C)}}{Z}$ we have
                                    $a\in \metricBall{\standardMetricR}{\reals}{\apply{\apply{f}{(n,C)}}{x}}{\eta}$.
                                \begin{subproof}
                                    Fix $C\in H$.
                                    Fix $a$ such that $a\in\img{\apply{f}{(n,C)}}{Z}$.
                                    We have $C\in\apply{b}{n}$ by \cref{subseteq}.
                                    $C\in G$ by \cref{subseteq,img_iff,function_apply_elim,unions_iff}.
                                    Let $p = (n,C)$.
                                    $\fst{p} = n$ and $\snd{p} = C$ by \cref{fst_eq,snd_eq}.
                                    $p\in\naturalsPlus\times G$ by \cref{times_tuple_intro}.
                                    Thus $\snd{p}\in\apply{b}{\fst{p}}$ by \cref{subseteq,pair_eq_iff}.
                                    Hence $p\in J$ by \cref{nagata_smirnov_index_set}.
                                    We have $\apply{f}{(n,C)}\in\funs{X}{\reals}$ by \cref{continuous,pair_eq_iff}.
                                    Take $y\in\dom{\apply{f}{(n,C)}}\inter Z$ such that $a = \apply{\apply{f}{(n,C)}}{y}$
                                        by \cref{funs_is_function,img_of_function_elim}.
                                    Thus $y\in X$ and $y\in Z$ by \cref{continuous,funs,pair_eq_iff,inter_elim_left,inter_elim_right}.
                                    Hence $\apply{\standardMetricR}{(\apply{\apply{f}{(n,C)}}{x},\apply{\apply{f}{(n,C)}}{y})} < \eta$
                                        by \cref{subseteq}.
                                    $a\in\reals$ by \cref{funs_type_apply,pair_eq_iff}.
                                    Therefore $a\in \metricBall{\standardMetricR}{\reals}{\apply{\apply{f}{(n,C)}}{x}}{\eta}$
                                        by \cref{metric_ball,pair_eq_iff}.
                                \end{subproof}
                                Therefore for all $C\in H$ we have
                                    $\img{\apply{f}{(n,C)}}{Z}\subseteq \metricBall{\standardMetricR}{\reals}{\apply{\apply{f}{(n,C)}}{x}}{\eta}$
                                    by \cref{subseteq}.
                                Let $M = \{C\in\apply{b}{n} \mid \img{\apply{f}{(n,C)}}{Z}\subseteq
                                    \metricBall{\standardMetricR}{\reals}{\apply{\apply{f}{(n,C)}}{x}}{\eta}\}$.
                                $H\subseteq M$ by \cref{subseteq}.
                                $Z\in T$ by \cref{neighborhood,open_in}.
                                $(n,Z)\in\initialSegment{N}\times\pow{X}$ by \cref{times_tuple_intro,topology_on,subseteq,pow_iff}.
                                Thus $(n,Z)\in A$ by \cref{pair_eq_iff}.
                                Therefore $n\mathrel{A} Z$ by \cref{pair_eq_iff}.
                            \end{subproof}
                            Take $v$ such that
                                $v$ is a function on $\initialSegment{N}$ and
                                for all $n\in\initialSegment{N}$ we have $n\mathrel{A}\apply{v}{n}$
                                by \cref{choice_function}.
                            Let $K = \img{v}{\initialSegment{N}}$.
                            $K$ is finite by \cref{initial_segment_finite,finite_image_function}.
                            $1\in\initialSegment{N}$ by \cref{real_one_in_naturals,real_one_leq_naturalsplus,initial_segment_naturalsplus}.
                            $\apply{v}{1}\in K$ by \cref{img_iff,function_apply_elim}.
                            Thus $K$ is inhabited by \cref{nonempty_inhabited,emptyset}.
                            For all $W$ such that $W\in K$ we have $W\in T$
                                by \cref{img_iff,function_apply_iff,choice_function,pair_eq_iff,neighborhood,open_in}.
                            $K\subseteq T$ by \cref{subseteq}.
                            Let $W = \inters{K}$.
                            Hence $W\in T$ by \cref{finite_intersections_open_in_topology}.
                            For all $W_0$ such that $W_0\in K$ we have $x\in W_0$
                                by \cref{img_iff,function_apply_iff,choice_function,pair_eq_iff,neighborhood}.
                            Therefore $x\in W$ by \cref{inters_iff_forall,img_iff,function_apply_elim,real_one_in_naturals}.
                            Thus $W$ is a neighborhood of $x$ in $X$ with respect to $T$ by \cref{neighborhood,open_in}.
                            Show that for all $z$ such that $z\in\img{F}{W}$ we have $z\in V$.
                            \begin{subproof}
                                Fix $z$ such that $z\in\img{F}{W}$.
                                $F\in\funs{X}{\productFamily{R}}$ and $F$ is a function and $\dom{F} = X$
                                    by \cref{topological_embedding,continuous,funs,funs_is_function}.
                                Take $y\in\dom{F}\inter W$ such that $z = F(y)$ by \cref{funs_is_function,img_of_function_elim}.
                                We have $y\in X$ and $y\in W$ and $z = F(y)$ by \cref{inter_elim_left,inter_elim_right,subseteq}.
                                Let $p_0 = (F(x),F(y))$.
                                $F(x)\in\realsPower{J}$ and $F(y)\in\realsPower{J}$ by \cref{funs_type_apply,product_family_reals_power,subseteq}.
                                Thus $p_0\in\realsPower{J}\times\realsPower{J}$ by \cref{times_tuple_intro}.
                                $d_J$ is a metric on $\realsPower{J}$ by \cref{uniform_metric_is_metric_late}.
                                Hence $d_J\in\funs{\realsPower{J}\times\realsPower{J}}{\reals}$ by \cref{metric_on}.
                                Therefore $d_J$ is a function and $\dom{d_J} = \realsPower{J}\times\realsPower{J}$ by \cref{funs_is_function,funs}.
                                $p_0\in\dom{d_J}$ by \cref{subseteq}.
                                $(p_0,d_J(p_0))\in d_J$ by \cref{function_apply_elim}.
                                Let $S_0 = \uniformMetricSet{J}{F(x)}{F(y)}$.
                                Hence $d_J(p_0)$ is a supremum of $S_0$ by \cref{uniform_metric,pair_eq_iff,real_supremum_pred}.
                                For all $a$ such that $a\in S_0$ there exists $q\in J$ such that
                                    $a = \apply{e}{(\apply{F(x)}{q},\apply{F(y)}{q})}$ by \cref{uniform_metric_set}.
                                For all $q\in J$ we have $\apply{e}{(\apply{F(x)}{q},\apply{F(y)}{q})}\in\reals$
                                    by \cref{topological_embedding,continuous,funs,funs_type_apply,product_family_reals_power,subseteq,reals_power,tuple_power,times_tuple_intro,bounded_metric_is_metric,standard_metric_is_metric,metric_on}.
                                Thus for all $a$ such that $a\in S_0$ we have $a\in\reals$
                                    by \cref{uniform_metric_set,subseteq}.
                                Hence $S_0\subseteq\reals$ by \cref{subseteq}.
                                Show that for all $a\in S_0$ we have $a \leq \eta$.
                                \begin{subproof}
                                    Fix $a\in S_0$.
                                    Take $q\in J$ such that $a = \apply{e}{(\apply{F(x)}{q},\apply{F(y)}{q})}$ by \cref{uniform_metric_set}.
                                    Let $n = \fst{q}$.
                                    Let $C = \snd{q}$.
                                    $q\in\naturalsPlus\times G$ and $C\in\apply{b}{n}$ by \cref{nagata_smirnov_index_set}.
                                    $q = (\fst{q},\snd{q})$ by \cref{times_elem_is_tuple,fst_eq,snd_eq}.
                                    Thus $\fst{q}\in\naturalsPlus$ and $\snd{q}\in G$ by \cref{times_tuple_elim,pair_eq_iff}.
                                    Hence $n\in\naturalsPlus$ by \cref{pair_eq_iff}.
                                    $F\in\funs{X}{\productFamily{R}}$ and $F$ is a function and $\dom{F} = X$
                                        by \cref{topological_embedding,continuous,funs,funs_is_function}.
                                    Hence $x\in\dom{F}$ and $y\in\dom{F}$ by \cref{subseteq}.
                                    $(x,F(x))\in F$ and $(y,F(y))\in F$ by \cref{function_apply_elim}.
                                    $\apply{F(x)}{q} = \apply{\apply{f}{q}}{x}$ and $\apply{F(y)}{q} = \apply{\apply{f}{q}}{y}$ by \cref{pair_eq_iff}.
                                    \begin{byCase}
                                            \caseOf{$n\in\initialSegment{N}$.}
                                                $(n,v(n))\in A$ by \cref{choice_function}.
                                                Take $H,M$ such that
                                                    $v(n)$ is a neighborhood of $x$ in $X$ with respect to $T$ and
                                                    $H\subseteq\apply{b}{n}$ and
                                                    $H$ is finite and
                                                    $M = \{D\in\apply{b}{n} \mid \img{\apply{f}{(n,D)}}{v(n)}\subseteq
                                                        \metricBall{\standardMetricR}{\reals}{\apply{\apply{f}{(n,D)}}{x}}{\eta}\}$ and
                                                    $H\subseteq M$ and
                                                    for all $D\in\apply{b}{n}$ such that $v(n)$ intersects $D$ we have $D\in H$
                                                    by \cref{pair_eq_iff}.
                                                $v(n)\in K$ by \cref{img_iff,function_apply_elim}.
                                                For all $z_0$ such that $z_0\in W$ we have $z_0\in v(n)$ by \cref{inters_iff_forall}.
                                                $W\subseteq v(n)$ by \cref{subseteq}.
                                                Hence $y\in v(n)$ by \cref{subseteq}.
                                                \begin{byCase}
                                                \caseOf{$C\in H$.}
                                                    We have $q = (n,C)$ by \cref{pair_eq_iff}.
                                                    $\apply{f}{q}\in\funs{X}{\reals}$ and $\apply{f}{q}$ is a function and $\dom{\apply{f}{q}} = X$
                                                        by \cref{continuous,funs,funs_is_function}.
                                                    Thus $y\in\dom{\apply{f}{q}}\inter v(n)$ by \cref{subseteq,inter_intro}.
                                                    $\apply{\apply{f}{q}}{y}\in\img{\apply{f}{q}}{v(n)}$ by \cref{img_of_function_intro}.
                                                    Therefore $C\in M$ by \cref{subseteq}.
                                                    We have $\img{\apply{f}{(n,C)}}{v(n)}\subseteq \metricBall{\standardMetricR}{\reals}{\apply{\apply{f}{(n,C)}}{x}}{\eta}$
                                                        by \cref{subseteq}.
                                                    Thus $\img{\apply{f}{q}}{v(n)}\subseteq \metricBall{\standardMetricR}{\reals}{\apply{\apply{f}{q}}{x}}{\eta}$
                                                        by \cref{pair_eq_iff}.
                                                    $\apply{\apply{f}{q}}{y}\in\metricBall{\standardMetricR}{\reals}{\apply{\apply{f}{q}}{x}}{\eta}$ by \cref{subseteq}.
                                                    Hence $\apply{\standardMetricR}{(\apply{\apply{f}{q}}{x},\apply{\apply{f}{q}}{y})} < \eta$ by \cref{metric_ball}.
                                                    We have $\eta < \epsilon$ by \cref{real_lt_add,real_add_ident,real_lt_subst_left}.
                                                    $1\in\reals$ by \cref{real_one_in_naturals,real_naturals_subset_reals,subseteq}.
                                                    Thus $\eta < 1$ by \cref{real_half_of_positive,real_lt_leq_trans_local}.
                                                    $\apply{\standardMetricR}{(\apply{F(x)}{q},\apply{F(y)}{q})} =
                                                        \apply{\standardMetricR}{(\apply{\apply{f}{q}}{x},\apply{\apply{f}{q}}{y})}$ by \cref{pair_eq_iff}.
                                                    Hence $\apply{\standardMetricR}{(\apply{F(x)}{q},\apply{F(y)}{q})} < \eta$
                                                        by \cref{real_lt_subst_right}.
                                                    $\apply{F(x)}{q}\in\reals$ and $\apply{F(y)}{q}\in\reals$
                                                        by \cref{continuous,funs_type_apply,pair_eq_iff}.
                                                    $(\apply{F(x)}{q},\apply{F(y)}{q})\in\reals\times\reals$ by \cref{times_tuple_intro}.
                                                    Thus $\apply{\standardMetricR}{(\apply{F(x)}{q},\apply{F(y)}{q})}\in\reals$
                                                        by \cref{standard_metric_is_metric,metric_on,funs_type_apply}.
                                                    Hence $\apply{\standardMetricR}{(\apply{F(x)}{q},\apply{F(y)}{q})} < 1$ by \cref{real_lt_trans}.
                                                    We have $a = \apply{\standardMetricR}{(\apply{F(x)}{q},\apply{F(y)}{q})}$
                                                        by \cref{bounded_metric_apply_below_one,standard_metric_is_metric,pair_eq_iff,times_tuple_intro}.
                                                    Therefore $a \leq \eta$ by \cref{real_lt_imp_leq,real_lt_subst_right}.
                                                \caseOf{$C\notin H$.}
                                                    Thus it is wrong that $v(n)$ intersects $C$ by \cref{subseteq}.
                                                    $x\in X\setminus C$ and $y\in X\setminus C$
                                                        by \cref{neighborhood,subseteq,inter_intro,emptyset,intersects,setminus_intro}.
                                                    We have $q = (n,C)$ by \cref{pair_eq_iff}.
                                                    $\apply{f}{q}\in\funs{X}{\reals}$ and $\apply{f}{q}$ is a function and $\dom{\apply{f}{q}} = X$
                                                        by \cref{continuous,funs,funs_is_function}.
                                                    $\fst{q} = n$ and $\snd{q} = C$ by \cref{fst_eq,snd_eq,pair_eq_iff}.
                                                    Thus $x\in\dom{\apply{f}{q}}\inter(X\setminus C)$ and
                                                        $y\in\dom{\apply{f}{q}}\inter(X\setminus C)$ by \cref{subseteq,inter_intro}.
                                                    $\apply{\apply{f}{q}}{x}\in\img{\apply{f}{q}}{X\setminus C}$ and
                                                        $\apply{\apply{f}{q}}{y}\in\img{\apply{f}{q}}{X\setminus C}$
                                                        by \cref{img_of_function_intro}.
                                                    We have $\img{\apply{f}{q}}{X\setminus \snd{q}}\subseteq \{0\}$ by \cref{subseteq}.
                                                    $\img{\apply{f}{q}}{X\setminus C}\subseteq \{0\}$ by \cref{pair_eq_iff}.
                                                    Hence $\apply{\apply{f}{q}}{x} = 0$ and $\apply{\apply{f}{q}}{y} = 0$
                                                        by \cref{singleton_elim,subseteq}.
                                                    Thus $\apply{F(x)}{q} = 0$ and $\apply{F(y)}{q} = 0$ by \cref{pair_eq_iff}.
                                                    $\apply{F(x)}{q}\in\reals$ and $\apply{F(y)}{q}\in\reals$
                                                        by \cref{continuous,funs_type_apply,pair_eq_iff}.
                                                    Therefore $\apply{F(x)}{q} = \apply{F(y)}{q}$ by \cref{pair_eq_iff}.
                                                    $(\apply{F(x)}{q},\apply{F(y)}{q})\in\reals\times\reals$ by \cref{times_tuple_intro}.
                                                    Hence $\apply{\standardMetricR}{(\apply{F(x)}{q},\apply{F(y)}{q})} = 0$
                                                        by \cref{standard_metric_is_metric,metric_on,metric_zero_characterization}.
                                                    We have $a = 0$ by \cref{bounded_metric_apply_below_one,standard_metric_is_metric,pair_eq_iff,times_tuple_intro,real_zero_lt_one}.
                                                    Therefore $a \leq \eta$ by \cref{positive_reciprocal_naturalsplus,real_lt_imp_leq,real_add_ident,real_leq_trans}.
                                                \end{byCase}
                                            \caseOf{$n\notin\initialSegment{N}$.}
                                                $N < n$
                                                    by \cref{initial_segment_naturalsplus,real_lt_trichotomy,real_naturals_subset_reals,subseteq}.
                                                Hence $N \leq n$ by \cref{real_lt_imp_leq}.
                                                Thus $1 / n \leq 1 / N$ by \cref{naturalsplus_reciprocal_antitone}.
                                                $\apply{f}{q}$ is continuous from $X$ to $\reals$ with respect to $T$ and $\standardTopologyR$ by \cref{subseteq}.
                                                $\img{\apply{f}{q}}{X}\subseteq \intervalclosed{0}{1 / n}$ by \cref{subseteq,pair_eq_iff}.
                                                $\apply{f}{q}\in\funs{X}{\reals}$ and $\apply{f}{q}$ is a function and $\dom{\apply{f}{q}} = X$
                                                    by \cref{continuous,funs,funs_is_function}.
                                                We have $x\in\dom{\apply{f}{q}}\inter X$ and $y\in\dom{\apply{f}{q}}\inter X$ by \cref{subseteq,inter_intro}.
                                                $\apply{\apply{f}{q}}{x}\in\img{\apply{f}{q}}{X}$ and $\apply{\apply{f}{q}}{y}\in\img{\apply{f}{q}}{X}$
                                                    by \cref{img_of_function_intro}.
                                                Thus $\apply{F(x)}{q}\in\intervalclosed{0}{1 / n}$ and $\apply{F(y)}{q}\in\intervalclosed{0}{1 / n}$
                                                    by \cref{subseteq,pair_eq_iff}.
                                                $1 / n\in\reals$ and $0 \leq 1 / n$
                                                    by \cref{naturalsplus_reciprocal_real,positive_reciprocal_naturalsplus,real_lt_imp_leq}.
                                                Hence $\apply{\standardMetricR}{(\apply{F(x)}{q},\apply{F(y)}{q})} \leq 1 / n$
                                                    by \cref{standard_metric_intervalclosed_bound}.
                                                $\apply{F(x)}{q}\in\reals$ and $\apply{F(y)}{q}\in\reals$ by \cref{intervalclosed}.
                                                $(\apply{F(x)}{q},\apply{F(y)}{q})\in\reals\times\reals$ by \cref{times_tuple_intro}.
                                                Thus $\apply{\standardMetricR}{(\apply{F(x)}{q},\apply{F(y)}{q})}\in\reals$
                                                    by \cref{standard_metric_is_metric,metric_on,funs_type_apply}.
                                                $1 / N\in\reals$ by \cref{naturalsplus_reciprocal_real}.
                                                $\eta\in\reals$ by \cref{real_half_of_positive}.
                                                We have $1 / N < \eta$ by \cref{pair_eq_iff}.
                                                Thus $1 / n < \eta$ by \cref{real_leq_lt_trans}.
                                                $\apply{\standardMetricR}{(\apply{F(x)}{q},\apply{F(y)}{q})} < \eta$
                                                    by \cref{real_leq_lt_trans}.
                                                Hence $\eta < \epsilon$ by \cref{real_lt_add,real_add_ident,real_lt_subst_left}.
                                                $1\in\reals$ by \cref{real_one_in_naturals,real_naturals_subset_reals,subseteq}.
                                                Thus $\eta < 1$ by \cref{real_half_of_positive,real_lt_leq_trans_local}.
                                                $\apply{\standardMetricR}{(\apply{F(x)}{q},\apply{F(y)}{q})} < 1$ by \cref{real_lt_trans}.
                                                We have $a = \apply{\standardMetricR}{(\apply{F(x)}{q},\apply{F(y)}{q})}$
                                                    by \cref{bounded_metric_apply_below_one,standard_metric_is_metric,pair_eq_iff,times_tuple_intro}.
                                                Thus $a \leq 1 / n$ by \cref{real_leq_trans,pair_eq_iff}.
                                                $a \leq 1 / N$ by \cref{real_leq_trans}.
                                                Therefore $a \leq \eta$ by \cref{real_lt_imp_leq,real_leq_trans}.
                                    \end{byCase}
                                \end{subproof}
                                Thus $\eta$ is an upper bound of $S_0$ by \cref{real_upper_bound}.
                                $d_J(p_0) \leq \eta$ by \cref{real_supremum}.
                                We have $0 < \eta$ and $\eta + \eta = \epsilon$ by \cref{real_half_of_positive}.
                                $0\in\reals$ by \cref{real_add_ident}.
                                Thus $0 + \eta < \eta + \eta$ by \cref{real_lt_add}.
                                $\eta < \eta + \eta$ by \cref{real_add_ident}.
                                Hence $\eta < \epsilon$ by \cref{real_lt_subst_right}.
                                $d_J(p_0)\in\reals$ by \cref{funs_type_apply}.
                                Thus $d_J(p_0) < \epsilon$ by \cref{real_leq_lt_trans}.
                                Hence $F(y)\in\metricBall{d_J}{\realsPower{J}}{F(x)}{\epsilon}$ by \cref{metric_ball}.
                                Therefore $z\in V$ by \cref{subseteq,real_lt_subst_right}.
                            \end{subproof}
                            Hence $\img{F}{W}\subseteq V$ by \cref{subseteq}.
                        \caseOf{$1 < \epsilon$.}
                            Let $\eta = 1$.
                            $0 < \eta$ and $\eta\in\reals$ and $\eta < \epsilon$
                                by \cref{real_one_in_naturals,real_naturals_subset_reals,subseteq,real_zero_lt_one,real_lt_subst_left}.
                            Take $N\in\naturalsPlus$ such that $0 < 1 / N$ and $1 / N < \eta$
                                by \cref{real_small_reciprocal}.
                            Let $W = X$.
                            Thus $W$ is a neighborhood of $x$ in $X$ with respect to $T$
                                by \cref{open_in,topology_on,subseteq,neighborhood}.
                            Show that for all $z$ such that $z\in\img{F}{W}$ we have $z\in V$.
                            \begin{subproof}
                                Fix $z$ such that $z\in\img{F}{W}$.
                                Take $y\in X$ such that $y\mathrel{F} z$ by \cref{img_iff,pair_eq_iff}.
                                $F\in\funs{X}{\productFamily{R}}$ and $z = F(y)$
                                    by \cref{topological_embedding,continuous,funs,funs_is_function,subseteq,function_apply_iff}.
                                Let $p_0 = (F(x),F(y))$.
                                $p_0\in\realsPower{J}\times\realsPower{J}$ by \cref{funs_type_apply,product_family_reals_power,subseteq,times_tuple_intro}.
                                $d_J$ is a metric on $\realsPower{J}$ by \cref{uniform_metric_is_metric_late}.
                                Thus $d_J\in\funs{\realsPower{J}\times\realsPower{J}}{\reals}$ by \cref{metric_on}.
                                Therefore $d_J$ is a function and $\dom{d_J} = \realsPower{J}\times\realsPower{J}$ by \cref{funs_is_function,funs}.
                                $p_0\in\dom{d_J}$ by \cref{subseteq}.
                                $(p_0,d_J(p_0))\in d_J$ by \cref{function_apply_elim}.
                                Let $S_0 = \uniformMetricSet{J}{F(x)}{F(y)}$.
                                Thus $d_J(p_0)$ is a supremum of $S_0$ by \cref{uniform_metric,pair_eq_iff,real_supremum_pred}.
                                For all $a$ such that $a\in S_0$ there exists $q\in J$ such that
                                    $a = \apply{e}{(\apply{F(x)}{q},\apply{F(y)}{q})}$ by \cref{uniform_metric_set}.
                                For all $q\in J$ we have $\apply{e}{(\apply{F(x)}{q},\apply{F(y)}{q})}\in\reals$
                                    by \cref{times_tuple_elim,pair_eq_iff,reals_power,tuple_power,funs_type_apply,times_tuple_intro,bounded_metric_is_metric,standard_metric_is_metric,metric_on}.
                                For all $a$ such that $a\in S_0$ we have $a\in\reals$
                                    by \cref{uniform_metric_set,subseteq}.
                                Hence $S_0\subseteq\reals$ by \cref{subseteq}.
                                Show that for all $a\in S_0$ we have $a \leq \eta$.
                                \begin{subproof}
                                    Follows by \cref{topological_embedding,continuous,funs_type_apply,product_family_reals_power,subseteq,reals_power,tuple_power,times_tuple_intro,bounded_metric_is_metric,standard_metric_is_metric,metric_on,bounded_metric_leq_one,real_leq_subst_right_local}.
                                \end{subproof}
                                Hence $\eta$ is an upper bound of $S_0$ by \cref{real_upper_bound}.
                                Thus $d_J(p_0) \leq \eta$ by \cref{real_supremum}.
                                We have $d_J(p_0)\in\reals$ by \cref{funs_type_apply}.
                                Therefore $d_J(p_0) < \epsilon$ by \cref{real_leq_lt_trans}.
                                $F(x)\in\realsPower{J}$ and $F(y)\in\realsPower{J}$ by \cref{times_tuple_elim,pair_eq_iff}.
                                Thus $d_J((F(x),F(y))) < \epsilon$ by \cref{pair_eq_iff}.
                                Hence $F(y)\in\metricBall{d_J}{\realsPower{J}}{F(x)}{\epsilon}$ by \cref{metric_ball}.
                                Therefore $z\in V$ by \cref{subseteq,real_lt_subst_right}.
                            \end{subproof}
                            Hence $\img{F}{W}\subseteq V$ by \cref{subseteq}.
                        \end{byCase}
                    \end{subproof}
                    Hence $F$ is continuous at $x$ from $X$ to $\realsPower{J}$ with respect to $T$ and $U_0$ by \cref{continuous_at,topological_embedding}.
            \end{subproof}
            Thus $F$ is continuous from $X$ to $\realsPower{J}$ with respect to $T$ and $U_0$
                by \cref{continuous_at_implies_continuous,topological_embedding}.

            Let $Z = \img{F}{X}$.
            $Z\subseteq\realsPower{J}$ by \cref{img_iff,pair_eq_iff,product_family_reals_power,subseteq}.
            Thus $Z$ is a subspace of $\realsPower{J}$ with respect to $U_0$ by \cref{subspace_of}.
            Let $U_Z = \subspaceTopology{U_0}{Z}$.

            $F$ is a homeomorphism between $X$ and $Z$ with respect to $T$ and $\subspaceTopology{P}{Z}$
                by \cref{topological_embedding}.
            Hence $\converse{F}$ is continuous from $Z$ to $X$ with respect to $\subspaceTopology{P}{Z}$ and $T$
                by \cref{homeomorphism}.
            Show that for all $O$ such that $O$ is open in $X$ with respect to $T$
                we have $\preimg{\converse{F}}{O}\in\subspaceTopology{P}{Z}$.
            \begin{subproof}
                Follows by \cref{continuous,open_in}.
            \end{subproof}
            $U_0 = \uniformTopology{J}$ by \cref{uniform_topology}.
            Thus $U_0$ is finer than $P$ on $\realsPower{J}$ by \cref{uniform_finer_product,product_family_reals_power,pair_eq_iff}.
            Show that for all $O$ such that $O$ is open in $X$ with respect to $T$
                we have $\preimg{\converse{F}}{O}\in U_Z$.
            \begin{subproof}
                Follows by \cref{subspace_topology,finer_topology,subseteq}.
            \end{subproof}
            $U_Z$ is a topology on $Z$ by \cref{subspace_topology_is_topology,subspace_of}.
            Show that for all $O$ such that $O$ is open in $X$ with respect to $T$
                we have $\preimg{\converse{F}}{O}$ is open in $Z$ with respect to $U_Z$.
            \begin{subproof}
                Follows by \cref{open_in}.
            \end{subproof}
            $\converse{F}\in\funs{Z}{X}$ by \cref{continuous}.
            Hence $\converse{F}$ is continuous from $Z$ to $X$ with respect to $U_Z$ and $T$ by \cref{continuous}.

            $F$ is continuous from $X$ to $Z$ with respect to $T$ and $U_Z$
                by \cref{continuous_restrict_range,subseteq_refl}.
            $F\in\funs{X}{Z}$ by \cref{continuous}.
            Thus $F$ is a function and $\dom{F} = X$ by \cref{funs,funs_is_function}.
            $F\in\surj{X}{Z}$ by \cref{funs_surj_iff,img_dom}.
            Hence $F$ is a bijection from $X$ to $Z$ by \cref{bijections,injective_function,topological_embedding}.
            Thus $F$ is a homeomorphism between $X$ and $Z$ with respect to $T$ and $U_Z$ by \cref{homeomorphism}.

            Hence $\realsPower{J}$ is metrizable with respect to $U_0$ by \cref{metrizable,uniform_metric_is_metric_late}.
            Thus $Z$ is metrizable with respect to $U_Z$ by \cref{subspace_metrizable}.
            Therefore $X$ is metrizable with respect to $T$ by \cref{homeomorphism_preserves_metrizable}.
        \end{byCase}
\end{proof}

% from §40 helper lemma 2: Nagata-Smirnov forward implication
\begin{lemma}\label{nagata_smirnov_metrization_forward}
    If $X$ is metrizable with respect to $T$ then
        $X$ is regular with respect to $T$ and
        there exists $B$ such that
            $B$ is a basis for $T$ on $X$ and
            $B$ is countably locally finite in $X$ with respect to $T$.
\end{lemma}
\begin{proof}
    Assume $X$ is metrizable with respect to $T$.
        Show that $X$ is regular with respect to $T$.
        \begin{subproof}
            Follows by \cref{metrizable_normal,normal_implies_regular}.
        \end{subproof}
        Take $d$ such that $d$ is a metric on $X$ and $T = \metricTopology{d}{X}$ by \cref{metrizable}.
        Let $P = \pow{\pow{X}}$.
        Let $R = \{w\in\naturalsPlus\times P \mid \text{there exists $m\in\naturalsPlus$ such that there exists $C$ such that there exists $A$ such that
            $w = (m,C)$ and
            $A = \{U\in\pow{X} \mid \exists x\in X. U = \metricBall{d}{X}{x}{1 / m}\}$ and
            $C$ is an open refinement of $A$ on $X$ with respect to $T$ and
            $C$ is countably locally finite in $X$ with respect to $T$}\}$.
        $R$ is a relation by \cref{relation,subseteq,times_elem_is_tuple}.
        Show that for all $m\in\naturalsPlus$ there exists $C$ such that $m\mathrel{R} C$.
        \begin{subproof}
            Fix $m\in\naturalsPlus$.
            Let $A = \{U\in\pow{X} \mid \exists x\in X. U = \metricBall{d}{X}{x}{1 / m}\}$.
            For all $U\in A$ we have $U$ is open in $X$ with respect to $T$
                by \cref{pair_eq_iff,naturalsplus_reciprocal_real,positive_reciprocal_naturalsplus,metric_ball,subseteq,powerset_intro,metric_basis,metric_basis_is_basis,basis_elem_open_in_generated,basis_for_topology,metric_topology,open_in,basis_topology_is_topology}.
            $X\subseteq\unions{A}$
                by \cref{metric_on,metric_zero_characterization,positive_reciprocal_naturalsplus,metric_ball,real_lt_subst_left,subseteq,powerset_intro,pair_eq_iff,unions_iff}.
            We have $A$ is an open covering of $X$ with respect to $T$ by \cref{covering,open_covering}.
            Take $C$ such that
                $C$ is an open refinement of $A$ on $X$ with respect to $T$ and
                $C$ is countably locally finite in $X$ with respect to $T$
                by \cref{metrizable_open_cover_countably_locally_finite_refinement}.
            We have $m\mathrel{R} C$
                by \cref{open_refinement_on,refinement_on,powerset_intro,times_tuple_intro,pair_eq_iff}.
        \end{subproof}
        Take $c$ such that
            $c$ is a function on $\naturalsPlus$ and
            for all $m\in\naturalsPlus$ we have $m\mathrel{R}\apply{c}{m}$
            by \cref{choice_function}.
        Let $B = \unions{\img{c}{\naturalsPlus}}$.
        For all $C\in B$ we have $C$ is open in $X$ with respect to $T$
            by \cref{unions_iff,img_iff,function_apply_iff,choice_function,pair_eq_iff,open_refinement_on,subseteq}.

        Show that for all $U\in T$ for all $x\in U$
            there exists $C\in B$ such that $x\in C$ and $C\subseteq U$.
        \begin{subproof}
            Fix $U\in T$.
            Fix $x\in U$.
                We have $U$ is open in $X$ with respect to $T$ and $U\subseteq X$
                    by \cref{open_in,basis_topology_is_topology,metric_topology,metric_basis_is_basis,basis_for_topology,topology_on,subseteq,powerset_elim}.
                Take $\epsilon\in\reals$ such that $0 < \epsilon$ and
                    $\metricBall{d}{X}{x}{\epsilon}\subseteq U$
                    by \cref{metric_open_iff}.
                Let $\eta = \epsilon / (1 + 1)$.
                We have $\eta\in\reals$ and $0 < \eta$ and $\eta + \eta = \epsilon$
                    by \cref{real_half_of_positive}.
                Take $m\in\naturalsPlus$ such that $0 < 1 / m$ and $1 / m < \eta$
                    by \cref{real_small_reciprocal}.
                Take $A$ such that
                    $A = \{V\in\pow{X} \mid \exists a\in X. V = \metricBall{d}{X}{a}{1 / m}\}$ and
                    $\apply{c}{m}$ is an open refinement of $A$ on $X$ with respect to $T$ and
                    $\apply{c}{m}$ is countably locally finite in $X$ with respect to $T$
                    by \cref{choice_function,pair_eq_iff}.
                Take $C\in\apply{c}{m}$ such that $x\in C$ by \cref{open_refinement_on,covering,subseteq,unions_iff}.
                Take $D$ such that
                    $D\in A$ and
                    $C\subseteq D$
                    by \cref{open_refinement_on,refinement_on}.
                Take $a\in X$ such that $D = \metricBall{d}{X}{a}{1 / m}$ by \cref{pair_eq_iff}.
                Show that for all $y$ such that $y\in C$ we have $y\in U$.
                \begin{subproof}
                    Fix $y$ such that $y\in C$.
                    We have $d((y,a)) < 1 / m$ and $d((a,x)) < 1 / m$
                        by \cref{subseteq,metric_ball,metric_on,metric_symmetric,times_tuple_intro,real_lt_subst_left}.
                    Hence $y\in X$ and $x\in X$ by \cref{subseteq,metric_ball}.
                    We have $d((y,a))\in\reals$ and $d((a,x))\in\reals$
                        by \cref{metric_on,funs_type_apply,times_tuple_intro}.
                    $1 / m\in\reals$ by \cref{naturalsplus_reciprocal_real}.
                    We have $d((y,a)) + d((a,x))\in\reals$ by \cref{real_add_closed}.
                    Hence $d((y,x)) \leq d((y,a)) + d((a,x))$ by \cref{metric_on,metric_triangle_inequality,times_tuple_intro}.
                    We have $d((y,a)) < \eta$ and $d((a,x)) < \eta$ by \cref{real_lt_trans}.
                    Hence $d((y,a)) + d((a,x)) < \eta + \eta$ by \cref{real_lt_add_two}.
                    We have $d((y,x))\in\reals$ by \cref{metric_on,funs_type_apply,times_tuple_intro}.
                    Hence $d((y,x)) < d((y,a)) + d((a,x))$ or
                        $d((y,x)) = d((y,a)) + d((a,x))$
                        by \cref{real_leq_split}.
                    We have $d((y,x)) < \eta + \eta$ by \cref{real_lt_trans,real_lt_subst_left}.
                    Hence $d((y,x)) < \epsilon$ by \cref{real_half_of_positive,real_lt_subst_right}.
                    We have $d((x,y)) < \epsilon$
                        by \cref{metric_on,metric_symmetric,times_tuple_intro,real_lt_subst_left}.
                    Hence $y\in\metricBall{d}{X}{x}{\epsilon}$ by \cref{metric_ball}.
                    Therefore $y\in U$ by \cref{subseteq}.
                \end{subproof}
                Hence $C\subseteq U$ by \cref{subseteq}.
                We have $C\in B$ by \cref{img_iff,function_apply_elim,unions_iff}.
        \end{subproof}
        Hence $B$ is a basis for $T$ on $X$
            by \cref{metric_topology,basis_topology_is_topology,metric_basis_is_basis,basis_for_topology,open_collection_basis}.

        For all $m\in\naturalsPlus$ we have $\apply{c}{m}\subseteq\pow{X}$
            by \cref{choice_function,open_refinement_on,refinement_on,fst_eq,snd_eq,pair_eq_iff}.
        For all $m\in\naturalsPlus$ we have $\apply{c}{m}$ is countably locally finite in $X$ with respect to $T$
            by \cref{choice_function,fst_eq,snd_eq,pair_eq_iff}.
        Hence $B$ is countably locally finite in $X$ with respect to $T$
            by \cref{countable_union_countably_locally_finite}.
    Hence $X$ is regular with respect to $T$ and
        there exists $B$ such that
            $B$ is a basis for $T$ on $X$ and
            $B$ is countably locally finite in $X$ with respect to $T$
        by \cref{pair_eq_iff}.
\end{proof}

% from §40 Item 4: Nagata-Smirnov metrization theorem
\begin{theorem}\label{nagata_smirnov_metrization_theorem}
    $X$ is metrizable with respect to $T$ iff
        $X$ is regular with respect to $T$ and
        there exists $B$ such that
            $B$ is a basis for $T$ on $X$ and
            $B$ is countably locally finite in $X$ with respect to $T$.
\end{theorem}
\begin{proof}
    Follows by \cref{nagata_smirnov_metrization_forward,nagata_smirnov_metrization_backward}.
\end{proof}
\section{§41 Paracompactness}

% from §41 Item 1: paracompact space
\begin{definition}\label{paracompact_space}
    $X$ is paracompact with respect to $T$ iff
        $T$ is a topology on $X$ and
        for all $A$ if
            $A$ is an open covering of $X$ with respect to $T$
        then there exists $B$ such that
            $B$ is an open refinement of $A$ on $X$ with respect to $T$ and
            $B$ is locally finite in $X$ with respect to $T$ and
            $B$ is a covering of $X$.
\end{definition}

% from §41 Item 2: paracompact Hausdorff spaces are normal
\begin{theorem}\label{paracompact_hausdorff_normal}
    Assume $X$ is paracompact with respect to $T$.
    Assume $X$ is Hausdorff with respect to $T$.
    Then $X$ is normal with respect to $T$.
\end{theorem}
\begin{proof}
    $T$ is a topology on $X$ by \cref{paracompact_space}.
    $X$ has closed singletons with respect to $T$
        by \cref{singleton_subset_intro,pair_is_finite,singleton_iff,upair_intro_left,subseteq,subseteq_finite_is_finite,finite_point_set_closed_hausdorff,one_point_closed}.
    Show that for all $x,B$ such that $x\in X$ and $B$ is closed in $X$ with respect to $T$ and $x\notin B$
        there exist $U,V$ such that
            $U$ is open in $X$ with respect to $T$ and
            $V$ is open in $X$ with respect to $T$ and
            $x\in U$ and
            $B\subseteq V$ and
            $U$ is disjoint from $V$.
    \begin{subproof}
            Fix $x$.
            Fix $B$ such that $x\in X$ and $B$ is closed in $X$ with respect to $T$ and $x\notin B$.
            Let $S = \{W\in T \mid x\notin\closure{W}{X}{T}\}$.
            Let $A = S\union\{X\setminus B\}$.
            For all $W\in A$ we have $W$ is open in $X$ with respect to $T$
                by \cref{subseteq,open_in,union_iff,singleton_iff,closed_in}.
            Show that for all $y$ such that $y\in X$ we have $y\in\unions{A}$.
            \begin{subproof}
                Fix $y$ such that $y\in X$.
                \begin{byCase}
                \caseOf{$y\in B$.}
                    We have $y\neq x$.
                    Take $U_0,V_0$ such that
                        $U_0$ is a neighborhood of $x$ in $X$ with respect to $T$ and
                        $V_0$ is a neighborhood of $y$ in $X$ with respect to $T$ and
                        $U_0$ is disjoint from $V_0$
                        by \cref{hausdorff}.
                    $V_0$ is open in $X$ with respect to $T$ and $V_0\subseteq X$
                        by \cref{neighborhood,open_in,topology_on,subseteq,pow_iff}.
                    Therefore $V_0\subseteq X\setminus U_0$ by \cref{disjoint,subseteq,setminus_intro}.
                    We have $\closure{V_0}{X}{T}\subseteq X\setminus U_0$
                        by \cref{neighborhood,open_in,topology_on,subseteq,pow_iff,setminus_subseteq,double_relative_complement,closed_in,closure_subseteq_closed}.
                    Therefore $x\notin\closure{V_0}{X}{T}$ by \cref{neighborhood,subseteq,setminus_elim_right}.
                    We have $V_0\in S$ by \cref{open_in,subseteq}.
                    Therefore $y\in\unions{A}$ by \cref{union_iff,neighborhood,unions_iff}.
                \caseOf{$y\notin B$.}
                    We have $y\in\unions{A}$ by \cref{setminus_intro,singleton_iff,union_iff,unions_iff}.
                \end{byCase}
            \end{subproof}
            We have $A$ is an open covering of $X$ with respect to $T$ by \cref{subseteq,covering,open_covering}.
            Take $C$ such that
                $C$ is an open refinement of $A$ on $X$ with respect to $T$ and
                $C$ is locally finite in $X$ with respect to $T$ and
                $C$ is a covering of $X$
                by \cref{paracompact_space}.
            Let $D = \{W\in C \mid W\inter B \neq \emptyset\}$.
            We have $D$ is locally finite in $X$ with respect to $T$ by \cref{subseteq,locally_finite_subcollection}.
            Let $H = \{E\in\pow{X} \mid \exists F\in D. E = \closure{F}{X}{T}\}$.
            We have $H$ is the closure collection of $D$ in $X$ with respect to $T$ by \cref{closure_collection_in}.
            $B\subseteq\unions{D}$
                by \cref{closed_in,subseteq,covering,unions_iff,inter_intro,emptyset,intersects}.
            For all $W$ such that $W\in D$ we have $W\in T$ by \cref{subseteq,open_refinement_on,open_in}.
            Hence $\unions{D}$ is open in $X$ with respect to $T$ by \cref{subseteq,topology_on,open_in}.
            Therefore $\closure{\unions{D}}{X}{T} = \unions{H}$ by \cref{closure_of_union_locally_finite}.
            Show that $x\notin\closure{\unions{D}}{X}{T}$.
            \begin{subproof}
                Suppose not.
                Take $E\in H$ such that $x\in E$ by \cref{subseteq,unions_iff}.
                Take $F\in D$ such that $E = \closure{F}{X}{T}$ by \cref{closure_collection_in}.
                Take $G\in A$ such that $F\subseteq G$ by \cref{open_refinement_on,refinement_on}.
                We have $F$ intersects $B$.
                We have $G\neq X\setminus B$ by \cref{subseteq,intersects,nonempty_inhabited,inter,setminus_elim_right}.
                We have $G\in T$ and $x\notin\closure{G}{X}{T}$ by \cref{union_iff,singleton_iff,subseteq}.
                Therefore $\closure{F}{X}{T}\subseteq\closure{G}{X}{T}$
                    by \cref{open_in,topology_on,subseteq,pow_iff,open_refinement_on,refinement_on,powerset_elim,subseteq_closure,subseteq_transitive,closure_subseteq_closed,closure_closed_in}.
                We have $x\notin\closure{F}{X}{T}$ by \cref{subseteq}.
                Contradiction by \cref{closure_collection_in}.
            \end{subproof}
            Let $U = X\setminus\closure{\unions{D}}{X}{T}$.
            Let $V = \unions{D}$.
            We have $V$ is open in $X$ with respect to $T$ and $V\subseteq X$
                by \cref{open_in,topology_on,subseteq,pow_iff}.
            $U$ is open in $X$ with respect to $T$ and $x\in U$ by \cref{closure_closed_in,closed_in,setminus_intro}.
            Hence $U$ is disjoint from $V$ by \cref{subseteq_closure,subseteq,setminus_elim_right,disjoint}.
            We have $U$ is open in $X$ with respect to $T$ and
                $V$ is open in $X$ with respect to $T$ and
                $x\in U$ and
                $B\subseteq V$ and
                $U$ is disjoint from $V$.
    \end{subproof}
    Hence $X$ is regular with respect to $T$ by \cref{regular_space}.
    Show that for all $A,U$ such that $A$ is closed in $X$ with respect to $T$ and $U$ is open in $X$ with respect to $T$
        and $A\subseteq U$
        there exists $V$ such that
            $V$ is open in $X$ with respect to $T$ and
            $A\subseteq V$ and
            $\closure{V}{X}{T}\subseteq U$.
    \begin{subproof}
        Fix $A$.
        Fix $U$ such that $A$ is closed in $X$ with respect to $T$ and $U$ is open in $X$ with respect to $T$
            and $A\subseteq U$.
        Let $S = \{W\in T \mid \closure{W}{X}{T}\subseteq U\}$.
        Let $B = S\union\{X\setminus A\}$.
        For all $W\in B$ we have $W$ is open in $X$ with respect to $T$
            by \cref{subseteq,open_in,union_iff,singleton_iff,closed_in}.
        Show that for all $x$ such that $x\in X$ we have $x\in\unions{B}$.
        \begin{subproof}
            Fix $x$ such that $x\in X$.
            \begin{byCase}
            \caseOf{$x\in A$.}
                We have $U$ is a neighborhood of $x$ in $X$ with respect to $T$ by \cref{subseteq,neighborhood}.
                Take $W$ such that
                    $W$ is a neighborhood of $x$ in $X$ with respect to $T$ and
                    $\closure{W}{X}{T}\subseteq U$
                    by \cref{regular_closure_neighborhood}.
                We have $W\in S$ by \cref{neighborhood,open_in,subseteq}.
                Therefore $x\in\unions{B}$ by \cref{union_iff,neighborhood,unions_iff}.
            \caseOf{$x\notin A$.}
                We have $x\in\unions{B}$ by \cref{setminus_intro,singleton_iff,union_iff,unions_iff}.
            \end{byCase}
        \end{subproof}
        We have $B$ is an open covering of $X$ with respect to $T$ by \cref{subseteq,covering,open_covering}.
        Take $C$ such that
            $C$ is an open refinement of $B$ on $X$ with respect to $T$ and
            $C$ is locally finite in $X$ with respect to $T$ and
            $C$ is a covering of $X$
            by \cref{paracompact_space}.
        Let $D = \{W\in C \mid W\inter A \neq \emptyset\}$.
        We have $D$ is locally finite in $X$ with respect to $T$ by \cref{subseteq,locally_finite_subcollection}.
        Let $H = \{E\in\pow{X} \mid \exists F\in D. E = \closure{F}{X}{T}\}$.
        We have $H$ is the closure collection of $D$ in $X$ with respect to $T$ by \cref{closure_collection_in}.
        $A\subseteq\unions{D}$
            by \cref{closed_in,subseteq,covering,unions_iff,inter_intro,emptyset,intersects}.
        For all $W$ such that $W\in D$ we have $W\in T$ by \cref{subseteq,open_refinement_on,open_in}.
        Hence $\unions{D}$ is open in $X$ with respect to $T$ by \cref{subseteq,topology_on,open_in}.
        Therefore $\closure{\unions{D}}{X}{T} = \unions{H}$ by \cref{closure_of_union_locally_finite}.
        Show that for all $x$ such that $x\in\closure{\unions{D}}{X}{T}$ we have $x\in U$.
        \begin{subproof}
            Fix $x$ such that $x\in\closure{\unions{D}}{X}{T}$.
            Take $E\in H$ such that $x\in E$ by \cref{subseteq,unions_iff}.
            Take $F\in D$ such that $E = \closure{F}{X}{T}$ by \cref{closure_collection_in}.
            Take $G\in B$ such that $F\subseteq G$ by \cref{open_refinement_on,refinement_on}.
            We have $G\neq X\setminus A$ by \cref{subseteq,intersects,nonempty_inhabited,inter,setminus_elim_right}.
            We have $G\in T$ and $\closure{G}{X}{T}\subseteq U$ by \cref{union_iff,singleton_iff,subseteq}.
            Therefore $\closure{F}{X}{T}\subseteq\closure{G}{X}{T}$
                by \cref{open_in,topology_on,subseteq,pow_iff,open_refinement_on,refinement_on,powerset_elim,subseteq_closure,subseteq_transitive,closure_subseteq_closed,closure_closed_in}.
            Therefore $x\in U$ by \cref{closure_collection_in,subseteq}.
        \end{subproof}
        Hence $\closure{\unions{D}}{X}{T}\subseteq U$ by \cref{subseteq}.
        Let $V = \unions{D}$.
        We have $V$ is open in $X$ with respect to $T$ and $A\subseteq V$ and $\closure{V}{X}{T}\subseteq U$.
    \end{subproof}
    For all $x$ such that $x\in X$ we have $\{x\}$ is closed in $X$ with respect to $T$
        by \cref{regular_space,one_point_closed,singleton_intro,subseteq}.
    Therefore $X$ is normal with respect to $T$ by \cref{normal_from_closure_neighborhood}.
\end{proof}

% from §41 Item 3: closed subspaces of paracompact spaces are paracompact
\begin{theorem}\label{closed_subspace_paracompact}
    Assume $Y$ is a subspace of $X$ with respect to $T$.
    Assume $Y$ is closed in $X$ with respect to $T$.
    Assume $X$ is paracompact with respect to $T$.
    Let $T_Y = \subspaceTopology{T}{Y}$.
    Then $Y$ is paracompact with respect to $T_Y$.
\end{theorem}
\begin{proof}
    $T$ is a topology on $X$ and $Y\subseteq X$ by \cref{subspace_of}.
    Show that $T_Y$ is a topology on $Y$.
    \begin{subproof}
        Follows by \cref{subspace_topology_is_topology,subspace_of}.
    \end{subproof}
    \begin{byCase}
    \caseOf{$Y = \emptyset$.}
        Show that for all $A$ if
            $A$ is an open covering of $Y$ with respect to $T_Y$
        then there exists $B$ such that
            $B$ is an open refinement of $A$ on $Y$ with respect to $T_Y$ and
            $B$ is locally finite in $Y$ with respect to $T_Y$ and
            $B$ is a covering of $Y$.
        \begin{subproof}
            Fix $A$ such that $A$ is an open covering of $Y$ with respect to $T_Y$.
            We have $A\subseteq\pow{Y}$
                by \cref{open_covering,open_in,topology_on,subseteq,pow_iff,powerset_intro}.
            Let $B = A$.
            Show that $B$ is a refinement of $A$ on $Y$.
            \begin{subproof}
                Follows by \cref{open_covering,open_in,topology_on,subseteq,pow_iff,powerset_intro,real_lt_subst_left,subseteq_refl,refinement_on}.
            \end{subproof}
            $B$ is a covering of $Y$ and for all $U\in B$ we have $U$ is open in $Y$ with respect to $T_Y$
                by \cref{open_covering,real_lt_subst_left}.
            Hence $B$ is an open refinement of $A$ on $Y$ with respect to $T_Y$ and
                $B$ is locally finite in $Y$ with respect to $T_Y$ and
                $B$ is a covering of $Y$
                by \cref{open_refinement_on,emptyset,locally_finite_collection,real_lt_subst_left}.
        \end{subproof}
        Hence $Y$ is paracompact with respect to $T_Y$ by \cref{paracompact_space}.
    \caseOf{$Y \neq \emptyset$.}
        Show that for all $A$ if
            $A$ is an open covering of $Y$ with respect to $T_Y$
        then there exists $C$ such that
            $C$ is an open refinement of $A$ on $Y$ with respect to $T_Y$ and
            $C$ is locally finite in $Y$ with respect to $T_Y$ and
            $C$ is a covering of $Y$.
        \begin{subproof}
            Fix $A$ such that $A$ is an open covering of $Y$ with respect to $T_Y$.
            We have $A\subseteq\pow{Y}$
                by \cref{open_covering,open_in,topology_on,subseteq,pow_iff,powerset_intro}.
            Take $y_0$ such that $y_0\in Y$ by \cref{nonempty_inhabited}.
            Take $A_0\in A$ such that $y_0\in A_0$ by \cref{open_covering,covering,subseteq,unions_iff}.
            Let $E = \{V\in T \mid Y\inter V \in A\}$.
            Let $D = E\union\{X\setminus Y\}$.
            Show that for all $x$ such that $x\in X$ we have $x\in\unions{D}$.
            \begin{subproof}
                Fix $x$ such that $x\in X$.
                \begin{byCase}
                \caseOf{$x\in Y$.}
                    Take $U\in A$ such that $x\in U$ by \cref{open_covering,covering,subseteq,unions_iff}.
                    Take $V\in T$ such that $U = Y\inter V$ by \cref{open_covering,subspace_topology,open_in}.
                    Hence $V\in D$ and $x\in Y\inter V$ by \cref{subseteq,union_iff,real_lt_subst_right}.
                    Therefore $x\in\unions{D}$ by \cref{real_lt_subst_right,inter,subseteq,unions_iff}.
                \caseOf{$x\notin Y$.}
                    We have $x\in\unions{D}$ by \cref{setminus_intro,singleton_iff,union_iff,unions_iff}.
                \end{byCase}
            \end{subproof}
            Therefore $D$ is a covering of $X$ by \cref{subseteq,covering}.
            For all $W$ such that $W\in D$ we have $W$ is open in $X$ with respect to $T$
                by \cref{subseteq,open_in,union_iff,singleton_iff,closed_in}.
            Hence $D$ is an open covering of $X$ with respect to $T$ by \cref{open_covering}.
            Take $B$ such that
                $B$ is an open refinement of $D$ on $X$ with respect to $T$ and
                $B$ is locally finite in $X$ with respect to $T$ and
                $B$ is a covering of $X$
                by \cref{paracompact_space}.
            Let $C = \{U\in\pow{Y} \mid \exists V\in B. U = Y\inter V\}$.
            Show that for all $U\in C$ there exists $V\in A$ such that $U\subseteq V$.
            \begin{subproof}
                Fix $U\in C$.
                Take $V\in B$ such that $U = Y\inter V$ by \cref{subseteq}.
                Take $W\in D$ such that $V\subseteq W$ by \cref{open_refinement_on,refinement_on}.
                \begin{byCase}
                \caseOf{$W\in E$.}
                    We have $Y\inter W\in A$ by \cref{subseteq}.
                    We have $U\subseteq Y$ and $U\subseteq W$
                        by \cref{real_lt_subst_left,inter_lower_left,inter_lower_right,subseteq_transitive}.
                    Therefore $U\subseteq Y\inter W$ by \cref{subseteq_inter_iff}.
                \caseOf{$W\notin E$.}
                    We have $W = X\setminus Y$ by \cref{union_iff,singleton_iff}.
                    We have $U\subseteq Y$ and $U\subseteq X\setminus Y$
                        by \cref{real_lt_subst_left,inter_lower_left,inter_lower_right,real_lt_subst_right,subseteq_transitive}.
                    We have $U\subseteq Y\inter(X\setminus Y)$ by \cref{subseteq_inter_iff}.
                    Hence $Y\inter(X\setminus Y) = \emptyset$ by \cref{setminus_disjoint}.
                    Therefore $U\subseteq A_0$ by \cref{real_lt_subst_right,emptyset_subseteq,subseteq_transitive}.
                \end{byCase}
            \end{subproof}
            $C\subseteq\pow{Y}$ by \cref{subseteq}.
            Show that $C$ is a refinement of $A$ on $Y$.
            \begin{subproof}
                Follows by \cref{refinement_on}.
            \end{subproof}
            $Y\subseteq\unions{C}$
                by \cref{covering,subseteq,unions_iff,inter_intro,inter_lower_left,powerset_intro}.
            Therefore $C$ is a covering of $Y$ by \cref{covering}.
            For all $U\in C$ we have $U$ is open in $Y$ with respect to $T_Y$
                by \cref{subseteq,open_refinement_on,open_in,subspace_topology,real_lt_subst_left}.
            Therefore $C$ is an open refinement of $A$ on $Y$ with respect to $T_Y$ by \cref{open_refinement_on}.
            Show that for all $y\in Y$ there exist $N, G$ such that
                $N$ is a neighborhood of $y$ in $Y$ with respect to $T_Y$ and
                $G\subseteq C$ and
                $G$ is finite and
                for all $U\in C$ if $N$ intersects $U$ then $U\in G$.
            \begin{subproof}
                Fix $y\in Y$.
                Take $N_0, F$ such that
                    $N_0$ is a neighborhood of $y$ in $X$ with respect to $T$ and
                    $F\subseteq B$ and
                    $F$ is finite and
                    for all $V\in B$ if $N_0$ intersects $V$ then $V\in F$
                    by \cref{subseteq,locally_finite_collection}.
                Let $N = Y\inter N_0$.
                We have $N$ is a neighborhood of $y$ in $Y$ with respect to $T_Y$
                    by \cref{subspace_topology,open_in,inter_intro,neighborhood}.
                Let $g = \{(V,Y\inter V) \mid V\in F\}$.
                For all $V,U_1,U_2$ such that $(V,U_1)\in g$ and $(V,U_2)\in g$ we have $U_1 = U_2$
                    by \cref{pair_eq_iff}.
                $g$ is a function on $F$ by \cref{relation,subseteq,times_elem_is_tuple,rightunique,dom_intro,dom_iff,pair_eq_iff,subseteq,subseteq_antisymmetric}.
                Let $G = \img{g}{F}$.
                $G$ is finite by \cref{finite_image_function}.
                $G\subseteq C$
                    by \cref{img_iff,function_apply_iff,function_apply_elim,function_apply_intro,real_lt_subst_left,real_lt_subst_right,subseteq,inter_lower_left,powerset_intro}.
                For all $U\in C$ if $N$ intersects $U$ then $U\in G$
                    by \cref{subseteq,intersects,nonempty_inhabited,inter,real_lt_subst_right,inter_intro,emptyset,locally_finite_collection,pair_eq_iff,function_apply_intro,img_iff,function_apply_iff}.
                We have $N$ is a neighborhood of $y$ in $Y$ with respect to $T_Y$ and
                    $G\subseteq C$ and
                    $G$ is finite and
                    for all $U\in C$ if $N$ intersects $U$ then $U\in G$.
            \end{subproof}
            Therefore $C$ is locally finite in $Y$ with respect to $T_Y$ by \cref{locally_finite_collection}.
            We have $C$ is an open refinement of $A$ on $Y$ with respect to $T_Y$ and
                $C$ is locally finite in $Y$ with respect to $T_Y$ and
                $C$ is a covering of $Y$ by \cref{open_refinement_on}.
        \end{subproof}
        Hence $Y$ is paracompact with respect to $T_Y$ by \cref{paracompact_space}.
    \end{byCase}
\end{proof}

% from §41 helper definition 1: countably locally finite open refinement property
\begin{definition}\label{countably_locally_finite_open_refinement_property}
    $X$ has the countably locally finite open refinement property with respect to $T$ iff
        $T$ is a topology on $X$ and
        for all $A$ if
            $A$ is an open covering of $X$ with respect to $T$
        then there exists $B$ such that
            $B$ is an open refinement of $A$ on $X$ with respect to $T$ and
            $B$ is countably locally finite in $X$ with respect to $T$ and
            $B$ is a covering of $X$.
\end{definition}

% from §41 helper definition 2: locally finite refinement property
\begin{definition}\label{locally_finite_refinement_property}
    $X$ has the locally finite refinement property with respect to $T$ iff
        $T$ is a topology on $X$ and
        for all $A$ if
            $A$ is an open covering of $X$ with respect to $T$
        then there exists $B$ such that
            $B$ is a refinement of $A$ on $X$ and
            $B$ is locally finite in $X$ with respect to $T$ and
            $B$ is a covering of $X$.
\end{definition}

% from §41 helper definition 3: locally finite closed refinement property
\begin{definition}\label{locally_finite_closed_refinement_property}
    $X$ has the locally finite closed refinement property with respect to $T$ iff
        $T$ is a topology on $X$ and
        for all $A$ if
            $A$ is an open covering of $X$ with respect to $T$
        then there exists $B$ such that
            $B$ is a closed refinement of $A$ on $X$ with respect to $T$ and
            $B$ is locally finite in $X$ with respect to $T$ and
            $B$ is a covering of $X$.
\end{definition}

% from §41 helper definition 7: Michael prefix indices
\begin{definition}\label{michael_prefix_indices}
    $\michaelPrefixIndices{n} = \{i\in\naturalsPlus \mid i < n\}$.
\end{definition}

% from §41 helper definition 8: Michael difference layer
\begin{definition}\label{michael_difference_layer}
    $\michaelDifferenceLayer{b}{v}{n}{X} =
        \{W\in\pow{X} \mid
            \exists U\in\apply{b}{n}.
            W = U\setminus \unions{\img{\restrl{v}{\michaelPrefixIndices{n}}}{\michaelPrefixIndices{n}}}\}$.
\end{definition}

% from §41 helper lemma 2: first countable layer containing a point
\begin{lemma}\label{michael_first_layer_point}
    Assume $b$ is a function on $\naturalsPlus$.
    Assume $B = \unions{\img{b}{\naturalsPlus}}$.
    Assume for all $n\in\naturalsPlus$ we have $\apply{b}{n}\subseteq\pow{X}$.
    Let $v = \{(n,\unions{\apply{b}{n}}) \mid n\in\naturalsPlus\}$.
    Suppose $B$ is a covering of $X$.
    Suppose $x\in X$.
    Then there exist $m, U$ such that
        $m\in\naturalsPlus$ and
        $U\in\apply{b}{m}$ and
        $x\in U$ and
        $x\notin\unions{\img{\restrl{v}{\michaelPrefixIndices{m}}}{\michaelPrefixIndices{m}}}$.
\end{lemma}
\begin{proof}
    For all $w$ such that $w\in v$ there exist $a, c$ such that $w = (a,c)$.
    For all $n, c_1, c_2$ such that $(n,c_1)\in v$ and $(n,c_2)\in v$ we have $c_1 = c_2$
        by \cref{pair_eq_iff}.
    We have $v$ is a function on $\naturalsPlus$ by \cref{relation,rightunique,dom_intro,dom_iff,pair_eq_iff,subseteq,subseteq_antisymmetric}.

    Take $D\in B$ such that $x\in D$ by \cref{covering,subseteq,unions_iff}.
    Take $E$ such that $E\in\img{b}{\naturalsPlus}$ and $D\in E$ by \cref{unions_iff,real_lt_subst_right}.
    Take $n\in\naturalsPlus$ such that $\apply{b}{n} = E$ by \cref{img_iff,function_apply_iff}.
    Let $P = \{i\in\initialSegment{n} \mid x\in\unions{\apply{b}{i}}\}$.
        We have $P$ is inhabited by \cref{initial_segment_naturalsplus,real_lt_subst_left,unions_iff,nonempty_inhabited}.
        $P\subseteq\reals$ by \cref{subseteq,initial_segment_naturalsplus,real_naturals_subset_reals}.
        Hence $P$ is finite by \cref{subseteq,initial_segment_finite,subseteq_finite_is_finite}.
    Take $m$ such that $m$ is a smallest element of $P$ with respect to $\realOrderRel$
        by \cref{finite_inhabited_smallest_element,real_order_relation_simple}.
    Therefore $m\in P$ by \cref{smallest_element}.
    Hence $m\in\naturalsPlus$ and $m\in\initialSegment{n}$ and $x\in\unions{\apply{b}{m}}$
        by \cref{subseteq,initial_segment_naturalsplus}.
    Take $U\in\apply{b}{m}$ such that $x\in U$ by \cref{unions_iff}.
    Show that $x\notin\unions{\img{\restrl{v}{\michaelPrefixIndices{m}}}{\michaelPrefixIndices{m}}}$.
    \begin{subproof}
        Suppose not.
        Take $V$ such that $V\in\img{\restrl{v}{\michaelPrefixIndices{m}}}{\michaelPrefixIndices{m}}$ and $x\in V$ by \cref{unions_iff}.
        Take $i\in\michaelPrefixIndices{m}$ such that $(i,V)\in\restrl{v}{\michaelPrefixIndices{m}}$ by \cref{img_iff}.
        Hence $i\in\naturalsPlus$ and $i < m$ by \cref{michael_prefix_indices}.
        We have $x\in\unions{\apply{b}{i}}$ by \cref{restrl_iff,function_apply_iff,pair_eq_iff}.
        $i\in\reals$ and $m\in\reals$ and $n\in\reals$ by \cref{real_naturals_subset_reals,subseteq,initial_segment_naturalsplus}.
        Therefore $i\in\initialSegment{n}$ by
            \cref{initial_segment_naturalsplus,real_lt_trans,real_lt_subst_right,michael_prefix_indices}.
        We have $i\in P$.
        Hence $m\mathrel{\realOrderRel} i$ or $m = i$ by \cref{smallest_element}.
        \begin{byCase}
        \caseOf{$m\mathrel{\realOrderRel} i$.}
            We have $m < i$ by \cref{real_order_relation_elim}.
            $m\in\reals$ and $i\in\reals$ by \cref{real_naturals_subset_reals,subseteq}.
            Thus $m \neq i$ by \cref{real_lt_subst_left,real_lt_irreflexive}.
            Contradiction by \cref{real_lt_asym}.
        \caseOf{$m = i$.}
            We have $i < i$ by \cref{real_lt_subst_right}.
            Contradiction by \cref{real_naturals_subset_reals,subseteq,real_lt_irreflexive}.
        \end{byCase}
    \end{subproof}
    Thus $m\in\naturalsPlus$ and $U\in\apply{b}{m}$ and $x\in U$ and
        $x\notin\unions{\img{\restrl{v}{\michaelPrefixIndices{m}}}{\michaelPrefixIndices{m}}}$.
\end{proof}

% from §41 helper lemma 3: shrinking a locally finite collection to an indexed family
\begin{lemma}\label{locally_finite_collection_shrinking_indexed_family}
    Assume $A$ is locally finite in $X$ with respect to $T$.
    Assume $h$ is a function on $A$.
    Suppose for all $E\in A$ we have $\apply{h}{E}\subseteq E$.
    Then $h$ is a locally finite indexed family on $A$ in $X$ with respect to $T$.
\end{lemma}
\begin{proof}
    Show that $h$ is an indexed family on $A$ in $X$.
    \begin{subproof}
        Follows by \cref{locally_finite_collection,subseteq,powerset_elim,subseteq_transitive,indexed_family_on_in}.
    \end{subproof}
    For all $x\in X$ there exist $W_0, G$ such that
        $W_0$ is a neighborhood of $x$ in $X$ with respect to $T$ and
        $G\subseteq A$ and
        $G$ is finite and
        for all $E\in A$ if $W_0$ intersects $E$ then $E\in G$
        by \cref{locally_finite_collection}.
    We have for all $x\in X$ there exist $W, F$ such that
        $W$ is a neighborhood of $x$ in $X$ with respect to $T$ and
        $F\subseteq A$ and
        $F$ is finite and
        for all $E\in A$ if $W$ intersects $\apply{h}{E}$ then $E\in F$
        by \cref{intersects,nonempty_inhabited,inter,subseteq,real_lt_subst_left,inter_intro,emptyset}.
    Therefore $h$ is a locally finite indexed family on $A$ in $X$ with respect to $T$
        by \cref{locally_finite_indexed_family}.
\end{proof}

% from §41 Item 4: open locally finite refinements imply countably locally finite open refinements
\begin{lemma}\label{paracompact_implies_countably_locally_finite_open_refinement_property}
    Assume $X$ is paracompact with respect to $T$.
    Then $X$ has the countably locally finite open refinement property with respect to $T$.
\end{lemma}
\begin{proof}
    Show that $T$ is a topology on $X$.
    \begin{subproof}
        Follows by \cref{paracompact_space}.
    \end{subproof}
    Show that for all $A$ if
        $A$ is an open covering of $X$ with respect to $T$
    then there exists $B$ such that
        $B$ is an open refinement of $A$ on $X$ with respect to $T$ and
        $B$ is countably locally finite in $X$ with respect to $T$ and
        $B$ is a covering of $X$.
    \begin{subproof}
        Fix $A$ such that $A$ is an open covering of $X$ with respect to $T$.
        Take $B$ such that
            $B$ is an open refinement of $A$ on $X$ with respect to $T$ and
            $B$ is locally finite in $X$ with respect to $T$ and
            $B$ is a covering of $X$
            by \cref{paracompact_space}.
        We have $B\subseteq\pow{X}$ by \cref{open_refinement_on,refinement_on}.
        Let $b = \{(n,B) \mid n\in\naturalsPlus\}$.
        Thus $b$ is a function from $\naturalsPlus$ to $\pow{\pow{X}}$
            by \cref{open_refinement_on,refinement_on,powerset_intro,constant_map_function}.
        We have $b$ is a function on $\naturalsPlus$.
        $\unions{\img{b}{\naturalsPlus}}\subseteq B$
            by \cref{unions_iff,img_iff,function_apply_iff,constant_map_function,funs_is_function,function_apply_intro,real_lt_subst_left,subseteq}.
        We have $B\in\img{b}{\naturalsPlus}$
            by \cref{real_one_in_naturals,constant_map_function,funs_is_function,function_apply_intro,img_iff,function_apply_elim}.
        $B\subseteq\unions{\img{b}{\naturalsPlus}}$ by \cref{unions_iff,subseteq}.
        Thus $\unions{\img{b}{\naturalsPlus}} = B$ by \cref{subseteq_antisymmetric}.
        For all $n\in\naturalsPlus$ we have $\apply{b}{n}$ is locally finite in $X$ with respect to $T$
            by \cref{constant_map_function,funs_is_function,function_apply_intro,real_lt_subst_left}.
        Hence $B$ is countably locally finite in $X$ with respect to $T$ by \cref{countably_locally_finite_collection}.
        We have $B$ is an open refinement of $A$ on $X$ with respect to $T$ and
            $B$ is countably locally finite in $X$ with respect to $T$ and
            $B$ is a covering of $X$.
    \end{subproof}
    Hence $X$ has the countably locally finite open refinement property with respect to $T$
        by \cref{countably_locally_finite_open_refinement_property}.
\end{proof}

% from §41 Item 4: countably locally finite open refinements imply locally finite refinements
\begin{lemma}\label{countably_locally_finite_open_refinement_implies_locally_finite_refinement_property}
    Assume $X$ is regular with respect to $T$.
    Assume $X$ has the countably locally finite open refinement property with respect to $T$.
    Then $X$ has the locally finite refinement property with respect to $T$.
\end{lemma}
\begin{proof}
    Show that $T$ is a topology on $X$.
    \begin{subproof}
        Follows by \cref{countably_locally_finite_open_refinement_property,regular_space}.
    \end{subproof}
    Show that for all $A$ if
        $A$ is an open covering of $X$ with respect to $T$
    then there exists $C$ such that
        $C$ is a refinement of $A$ on $X$ and
        $C$ is locally finite in $X$ with respect to $T$ and
        $C$ is a covering of $X$.
    \begin{subproof}
        Fix $A$ such that $A$ is an open covering of $X$ with respect to $T$.
        Take $B$ such that
            $B$ is an open refinement of $A$ on $X$ with respect to $T$ and
            $B$ is countably locally finite in $X$ with respect to $T$ and
            $B$ is a covering of $X$
            by \cref{countably_locally_finite_open_refinement_property}.
        We have $B\subseteq\pow{X}$ by \cref{open_refinement_on,refinement_on}.
        Take $b$ such that
            $b$ is a function on $\naturalsPlus$ and
            $B = \unions{\img{b}{\naturalsPlus}}$ and
            for all $n\in\naturalsPlus$ we have $\apply{b}{n}$ is locally finite in $X$ with respect to $T$
            by \cref{countably_locally_finite_collection}.
        Let $v = \{(n,\unions{\apply{b}{n}}) \mid n\in\naturalsPlus\}$.
        For all $n, c_1, c_2$ such that $(n,c_1)\in v$ and $(n,c_2)\in v$ we have $c_1 = c_2$
            by \cref{pair_eq_iff}.
        We have $v$ is a function on $\naturalsPlus$ by \cref{relation,rightunique,dom_intro,dom_iff,pair_eq_iff,subseteq,subseteq_antisymmetric}.
        Let $c = \{(n,\michaelDifferenceLayer{b}{v}{n}{X}) \mid n\in\naturalsPlus\}$.
        For all $n, d_1, d_2$ such that $(n,d_1)\in c$ and $(n,d_2)\in c$ we have $d_1 = d_2$
            by \cref{pair_eq_iff}.
        We have $c$ is a function on $\naturalsPlus$ by \cref{relation,rightunique,dom_intro,dom_iff,pair_eq_iff,subseteq,subseteq_antisymmetric}.
        Let $C = \unions{\img{c}{\naturalsPlus}}$.
        We have $A\subseteq\pow{X}$
            by \cref{open_covering,open_in,topology_on,subseteq,pow_iff,powerset_intro}.
        For all $W$ such that $W\in C$ we have $W\in\pow{X}$ by
            \cref{unions_iff,img_iff,function_apply_iff,function_apply_elim,pair_eq_iff,real_lt_subst_left,michael_difference_layer,subseteq}.
        We have $C\subseteq\pow{X}$ by \cref{subseteq}.
        Show that for all $W\in C$ there exists $U\in A$ such that $W\subseteq U$.
        \begin{subproof}
            Fix $W\in C$.
            Take $D$ such that $D\in\img{c}{\naturalsPlus}$ and $W\in D$ by \cref{unions_iff}.
            Take $n\in\naturalsPlus$ such that $\apply{c}{n} = D$ by \cref{img_iff,function_apply_iff}.
            We have $W\in\michaelDifferenceLayer{b}{v}{n}{X}$
                by \cref{function_apply_elim,pair_eq_iff,real_lt_subst_left}.
            Take $E\in\apply{b}{n}$ such that
                $W = E\setminus \unions{\img{\restrl{v}{\michaelPrefixIndices{n}}}{\michaelPrefixIndices{n}}}$
                by \cref{michael_difference_layer,subseteq}.
            We have $E\in B$ by \cref{countably_locally_finite_collection,unions_iff,img_iff,function_apply_elim}.
            Take $U\in A$ such that $E\subseteq U$ by \cref{open_refinement_on,refinement_on}.
            Therefore $W\subseteq U$ by \cref{setminus_subseteq,subseteq_transitive}.
        \end{subproof}
        Show that $C$ is a refinement of $A$ on $X$.
        \begin{subproof}
            Follows by \cref{refinement_on}.
        \end{subproof}

        For all $n\in\naturalsPlus$ we have $\apply{b}{n}\subseteq\pow{X}$ by \cref{locally_finite_collection}.

        Show that for all $x$ such that $x\in X$ we have $x\in\unions{C}$.
        \begin{subproof}
            Fix $x$ such that $x\in X$.
            Take $m, U$ such that
                $m\in\naturalsPlus$ and
                $U\in\apply{b}{m}$ and
                $x\in U$ and
                $x\notin\unions{\img{\restrl{v}{\michaelPrefixIndices{m}}}{\michaelPrefixIndices{m}}}$
                by \cref{michael_first_layer_point,subseteq}.
            Let $W = U\setminus \unions{\img{\restrl{v}{\michaelPrefixIndices{m}}}{\michaelPrefixIndices{m}}}$.
            We have $W\in\michaelDifferenceLayer{b}{v}{m}{X}$ by
                \cref{subseteq,powerset_elim,setminus_subseteq,subseteq_transitive,powerset_intro,michael_difference_layer}.
            $(m,\apply{c}{m})\in c$ by \cref{function_apply_elim}.
            Therefore $\apply{c}{m} = \michaelDifferenceLayer{b}{v}{m}{X}$ by \cref{pair_eq_iff}.
            We have $\apply{c}{m}\in\img{c}{\naturalsPlus}$ and $W\in C$
                by \cref{real_lt_subst_right,img_iff,function_apply_elim,subseteq,unions_iff}.
            $x\in W$ by \cref{setminus_intro}.
            Therefore $x\in\unions{C}$ by \cref{unions_iff}.
        \end{subproof}
        Therefore $C$ is a covering of $X$ by \cref{subseteq,covering}.

        Show that for all $n\in\naturalsPlus$ we have $\apply{c}{n}$ is locally finite in $X$ with respect to $T$.
        \begin{subproof}
            Fix $n\in\naturalsPlus$.
            Let $h = \{(E,E\setminus \unions{\img{\restrl{v}{\michaelPrefixIndices{n}}}{\michaelPrefixIndices{n}}}) \mid E\in\apply{b}{n}\}$.
            For all $E, d_1, d_2$ such that $(E,d_1)\in h$ and $(E,d_2)\in h$ we have $d_1 = d_2$
                by \cref{subseteq,pair_eq_iff,real_lt_subst_left}.
            We have $h$ is a function on $\apply{b}{n}$ by \cref{relation,rightunique,dom_intro,dom_iff,pair_eq_iff,subseteq,subseteq_antisymmetric}.
            For all $E\in\apply{b}{n}$ we have $\apply{h}{E}\subseteq X$
                by \cref{function_apply_iff,pair_eq_iff,subseteq,powerset_elim,real_lt_subst_left,setminus_subseteq,subseteq_transitive}.
            For all $D$ such that $D\in\apply{c}{n}$ we have $D\in\img{h}{\apply{b}{n}}$
                by \cref{pair_eq_iff,function_apply_intro,michael_difference_layer,subseteq,img_iff,function_apply_iff}.
            We have $\apply{c}{n}\subseteq\img{h}{\apply{b}{n}}$ by \cref{subseteq}.
            For all $D$ such that $D\in\img{h}{\apply{b}{n}}$ we have $D\in\apply{c}{n}$
                by \cref{img_iff,function_apply_iff,pair_eq_iff,subseteq,powerset_elim,real_lt_subst_right,setminus_subseteq,subseteq_transitive,function_apply_intro,michael_difference_layer,powerset_intro}.
            We have $\img{h}{\apply{b}{n}}\subseteq\apply{c}{n}$ by \cref{subseteq}.
            Therefore $\apply{c}{n} = \img{h}{\apply{b}{n}}$ by \cref{subseteq_antisymmetric}.
            For all $E\in\apply{b}{n}$ we have $\apply{h}{E}\subseteq E$ by
                \cref{function_apply_iff,pair_eq_iff,real_lt_subst_left,setminus_subseteq}.
            We have $h$ is a locally finite indexed family on $\apply{b}{n}$ in $X$ with respect to $T$ by
                \cref{locally_finite_collection_shrinking_indexed_family}.
            Hence $\img{h}{\apply{b}{n}}$ is locally finite in $X$ with respect to $T$ by
                \cref{locally_finite_indexed_family_implies_locally_finite_collection}.
            Therefore $\apply{c}{n}$ is locally finite in $X$ with respect to $T$ by
                \cref{locally_finite_indexed_family_implies_locally_finite_collection,real_lt_subst_left}.
        \end{subproof}
        Show that for all $x\in X$ there exist $W, F$ such that
            $W$ is a neighborhood of $x$ in $X$ with respect to $T$ and
            $F\subseteq C$ and
            $F$ is finite and
            for all $D\in C$ if $W$ intersects $D$ then $D\in F$.
        \begin{subproof}
            Fix $x\in X$.
            Take $m, U$ such that
                $m\in\naturalsPlus$ and
                $U\in\apply{b}{m}$ and
                $x\in U$ and
                $x\notin\unions{\img{\restrl{v}{\michaelPrefixIndices{m}}}{\michaelPrefixIndices{m}}}$
                by \cref{michael_first_layer_point,subseteq}.
            Let $R = \{w\in\initialSegment{m}\times(\pow{X}\times\pow{\pow{X}}) \mid
                \text{$\fst{\snd{w}}$ is a neighborhood of $x$ in $X$ with respect to $T$ and
                $\snd{\snd{w}}\subseteq\apply{c}{\fst{w}}$ and
                $\snd{\snd{w}}$ is finite and
                for all $D\in\apply{c}{\fst{w}}$ if $\fst{\snd{w}}$ intersects $D$ then $D\in\snd{\snd{w}}$}\}$.
            $R$ is a relation by \cref{times_elem_is_tuple,relation}.
            Show that for all $i\in\initialSegment{m}$ there exists $p$ such that $i\mathrel{R} p$.
            \begin{subproof}
                Fix $i\in\initialSegment{m}$.
                Hence $i\in\naturalsPlus$ by \cref{initial_segment_naturalsplus}.
                Take $W_0, G_0$ such that
                    $W_0$ is a neighborhood of $x$ in $X$ with respect to $T$ and
                    $G_0\subseteq\apply{c}{i}$ and
                    $G_0$ is finite and
                    for all $D\in\apply{c}{i}$ if $W_0$ intersects $D$ then $D\in G_0$
                    by \cref{locally_finite_collection}.
                $W_0\subseteq X$ by \cref{neighborhood,open_in,topology_on,subseteq,pow_iff}.
                We have $W_0\in\pow{X}$ by \cref{powerset_intro}.
                $\apply{c}{i}\subseteq\pow{X}$ by \cref{locally_finite_collection}.
                We have $G_0\subseteq\pow{X}$ by \cref{subseteq_transitive}.
                Therefore $G_0\in\pow{\pow{X}}$ by \cref{powerset_intro}.
                Thus $(W_0,G_0)\in\pow{X}\times\pow{\pow{X}}$ by \cref{times_tuple_intro}.
                Hence $(i,(W_0,G_0))\in\initialSegment{m}\times(\pow{X}\times\pow{\pow{X}})$
                    by \cref{times_tuple_intro}.
                Hence $i\mathrel{R}(W_0,G_0)$ by \cref{times_tuple_intro,fst_eq,snd_eq}.
            \end{subproof}
            Take $g$ such that $g$ is a function on $\initialSegment{m}$ and for all $i\in\initialSegment{m}$ we have $i\mathrel{R}\apply{g}{i}$
                by \cref{choice_function}.
            Let $J = \img{g}{\initialSegment{m}}$.
            $J$ is finite by \cref{initial_segment_finite,finite_image_function}.
            Let $q = \{(p,\fst{p}) \mid p\in J\}$.
            Let $h = \{(p,\snd{p}) \mid p\in J\}$.
            For all $p, d_1, d_2$ such that $(p,d_1)\in q$ and $(p,d_2)\in q$ we have $d_1 = d_2$
                by \cref{pair_eq_iff}.
            We have $q$ is a function on $J$ by \cref{relation,rightunique,dom_intro,dom_iff,pair_eq_iff,subseteq,subseteq_antisymmetric}.
            For all $p, d_1, d_2$ such that $(p,d_1)\in h$ and $(p,d_2)\in h$ we have $d_1 = d_2$
                by \cref{pair_eq_iff}.
            We have $h$ is a function on $J$ by \cref{relation,rightunique,dom_intro,dom_iff,pair_eq_iff,subseteq,subseteq_antisymmetric}.
            Let $P = \img{q}{J}$.
            Let $K = \img{h}{J}$.
            $P$ is finite and $K$ is finite by \cref{finite_image_function}.
            For all $W_0$ such that $W_0\in P$ we have $W_0$ is open in $X$ with respect to $T$ and $x\in W_0$
                by \cref{img_iff,function_apply_iff,choice_function,relation,snd_eq,real_lt_subst_left,neighborhood,pair_eq_iff,function_apply_intro}.
            We have $U$ is open in $X$ with respect to $T$
                by \cref{countably_locally_finite_collection,unions_iff,img_iff,function_apply_elim,open_refinement_on}.
            Let $H = P\union\{U\}$.
            Hence $H$ is finite by \cref{singleton_subset_intro,upair_intro_left,pair_is_finite,subseteq_finite_is_finite,union_of_finite_is_finite}.
            Thus $H$ is inhabited by \cref{union_intro_right,singleton_intro,nonempty_inhabited}.
            For all $W_0$ such that $W_0\in H$ we have $W_0\in T$
                by \cref{subseteq,union_iff,singleton_iff,open_in,real_lt_subst_left}.
            We have $H\subseteq T$ by \cref{subseteq}.
            Let $W = \inters{H}$.
            Hence $W$ is open in $X$ with respect to $T$ by \cref{finite_intersections_open_in_topology,open_in}.
            For all $W_0$ such that $W_0\in H$ we have $x\in W_0$
                by \cref{subseteq,union_iff,singleton_iff,real_lt_subst_left}.
            Thus $x\in W$ by \cref{inters_iff_forall}.
            Hence $W$ is a neighborhood of $x$ in $X$ with respect to $T$ by \cref{neighborhood}.
            For all $G_0$ such that $G_0\in K$ we have $G_0$ is finite by \cref{img_iff,function_apply_iff,choice_function,relation,snd_eq,real_lt_subst_left,pair_eq_iff,function_apply_intro}.
            Therefore $\unions{K}$ is finite by \cref{finite_union_of_finite_family}.
            Let $F = \unions{K}$.
            Show that for all $G_0$ such that $G_0\in K$ we have $G_0\subseteq C$.
            \begin{subproof}
                Fix $G_0$ such that $G_0\in K$.
                Take $p\in J$ such that $\apply{h}{p} = G_0$ by \cref{img_iff,function_apply_iff}.
                We have $(p,G_0)\in h$ by \cref{function_apply_elim}.
                Therefore $G_0 = \snd{p}$ by \cref{pair_eq_iff}.
                Take $i\in\initialSegment{m}$ such that $\apply{g}{i} = p$ by \cref{img_iff,function_apply_iff}.
                Hence $i\in\naturalsPlus$ by \cref{initial_segment_naturalsplus}.
                We have $i\mathrel{R}\apply{g}{i}$ by \cref{choice_function}.
                We have $G_0\subseteq\apply{c}{i}$ by \cref{relation,fst_eq,snd_eq,real_lt_subst_left}.
                Hence $\apply{c}{i}\in\img{c}{\naturalsPlus}$ by \cref{img_iff,function_apply_elim,subseteq}.
                Therefore $G_0\subseteq C$ by \cref{subseteq,unions_iff}.
            \end{subproof}
            Hence $F\subseteq C$ by \cref{subseteq,unions_iff}.
            Show that for all $D\in C$ if $W$ intersects $D$ then $D\in F$.
            \begin{subproof}
                Fix $D\in C$.
                Assume $W$ intersects $D$.
                Take $M$ such that $M\in\img{c}{\naturalsPlus}$ and $D\in M$ by \cref{unions_iff}.
                Take $n\in\naturalsPlus$ such that $\apply{c}{n} = M$ by \cref{img_iff,function_apply_iff}.
                \begin{byCase}
                \caseOf{$n\in\initialSegment{m}$.}
                    We have $\fst{\apply{g}{n}}\in H$
                        by \cref{img_iff,function_apply_elim,pair_eq_iff,union_iff}.
                    Therefore $W\subseteq\fst{\apply{g}{n}}$ by \cref{inters_subseteq_elem}.
                    Take $z$ such that $z\in W\inter D$ by \cref{intersects,nonempty_inhabited}.
                    We have $z\in\fst{\apply{g}{n}}$ and $z\in D$ by \cref{inter,subseteq}.
                    Thus $\fst{\apply{g}{n}}$ intersects $D$ by \cref{inter_intro,emptyset,intersects}.
                    We have $D\in\apply{c}{n}$ by \cref{real_lt_subst_left}.
                    Therefore $n\mathrel{R}\apply{g}{n}$ by \cref{choice_function}.
                    We have for all $E\in\apply{c}{n}$ if $\fst{\apply{g}{n}}$ intersects $E$ then $E\in\snd{\snd{(n,\apply{g}{n})}}$
                        by \cref{choice_function,relation,fst_eq,snd_eq}.
                    We have $D\in\snd{\snd{(n,\apply{g}{n})}}$ by \cref{real_lt_subst_left}.
                    Therefore $\snd{(n,\apply{g}{n})} = \apply{g}{n}$ by \cref{snd_eq}.
                    We have $\snd{\snd{(n,\apply{g}{n})}} = \snd{\apply{g}{n}}$ by \cref{real_lt_subst_left}.
                    Hence $D\in\snd{\apply{g}{n}}$ by \cref{real_lt_subst_left}.
                    Therefore $\snd{\apply{g}{n}}\in K$ by \cref{img_iff,function_apply_elim,pair_eq_iff}.
                    Therefore $D\in F$ by \cref{unions_iff}.
                \caseOf{$n\notin\initialSegment{m}$.}
                    Suppose not.
                    $n\in\reals$ and $m\in\reals$ by \cref{real_naturals_subset_reals,subseteq}.
                    Hence $m < n$ by \cref{real_lt_trichotomy,initial_segment_naturalsplus}.
                    Therefore $m\in\michaelPrefixIndices{n}$ by \cref{michael_prefix_indices,subseteq}.
                    Take $z$ such that $z\in W\inter D$ by \cref{intersects,nonempty_inhabited}.
                    We have $z\in W$ and $z\in D$ by \cref{inter}.
                    Thus $z\in U$ by \cref{union_intro_right,singleton_intro,inters_subseteq_elem,subseteq}.
                    Hence $z\in\unions{\apply{b}{m}}$ by \cref{unions_iff}.
                    Therefore $z\in\apply{v}{m}$ by \cref{pair_eq_iff,function_apply_intro}.
                    We have $\michaelPrefixIndices{n}\subseteq\dom{v}$ by \cref{michael_prefix_indices,subseteq,real_lt_subst_left}.
                    $\restrl{v}{\michaelPrefixIndices{n}}$ is a function by \cref{restrl_of_function_is_function}.
                    Hence $\dom{\restrl{v}{\michaelPrefixIndices{n}}} = \michaelPrefixIndices{n}$ by
                        \cref{dom_of_restrl_eq_inter,inter_absorb_supseteq_left}.
                    Therefore $\apply{\restrl{v}{\michaelPrefixIndices{n}}}{m} = \apply{v}{m}$ by \cref{restrl_of_function_apply,subseteq}.
                    We have $m\in\dom{\restrl{v}{\michaelPrefixIndices{n}}}$ by \cref{real_lt_subst_left}.
                    We have $\apply{v}{m}\in\img{\restrl{v}{\michaelPrefixIndices{n}}}{\michaelPrefixIndices{n}}$
                        by \cref{function_apply_elim,real_lt_subst_right,img_iff}.
                    Therefore $z\in\unions{\img{\restrl{v}{\michaelPrefixIndices{n}}}{\michaelPrefixIndices{n}}}$ by \cref{unions_iff}.
                    Hence $D\in\michaelDifferenceLayer{b}{v}{n}{X}$ by \cref{real_lt_subst_left,pair_eq_iff,function_apply_intro}.
                    Take $E\in\apply{b}{n}$ such that
                        $D = E\setminus \unions{\img{\restrl{v}{\michaelPrefixIndices{n}}}{\michaelPrefixIndices{n}}}$
                        by \cref{michael_difference_layer,subseteq}.
                    Thus $z\in E\setminus \unions{\img{\restrl{v}{\michaelPrefixIndices{n}}}{\michaelPrefixIndices{n}}}$ by \cref{real_lt_subst_right}.
                    Contradiction by \cref{setminus_elim_right}.
                \end{byCase}
            \end{subproof}
            We have there exist $W_1, F_1$ such that
                $W_1$ is a neighborhood of $x$ in $X$ with respect to $T$ and
                $F_1\subseteq C$ and
                $F_1$ is finite and
                for all $D\in C$ if $W_1$ intersects $D$ then $D\in F_1$.
        \end{subproof}
        $C\subseteq\pow{X}$ by \cref{refinement_on}.
        Therefore $C$ is locally finite in $X$ with respect to $T$ by \cref{locally_finite_collection}.
        We have $C$ is a refinement of $A$ on $X$ and
            $C$ is locally finite in $X$ with respect to $T$ and
            $C$ is a covering of $X$.
    \end{subproof}
    Therefore $X$ has the locally finite refinement property with respect to $T$
        by \cref{locally_finite_refinement_property}.
\end{proof}

% from §41 Item 4: locally finite refinements imply locally finite closed refinements.
\begin{lemma}\label{locally_finite_refinement_implies_locally_finite_closed_refinement_property}
    Assume $X$ is regular with respect to $T$.
    Assume $X$ has the locally finite refinement property with respect to $T$.
    Then $X$ has the locally finite closed refinement property with respect to $T$.
\end{lemma}
\begin{proof}
    Show that $T$ is a topology on $X$.
    \begin{subproof}
        Follows by \cref{locally_finite_refinement_property}.
    \end{subproof}
    Show that for all $A$ if $A$ is an open covering of $X$ with respect to $T$ then there exists $B$ such that $B$ is a closed refinement of $A$ on $X$ with respect to $T$ and $B$ is locally finite in $X$ with respect to $T$ and $B$ is a covering of $X$.
    \begin{subproof}
        Fix $A$ such that $A$ is an open covering of $X$ with respect to $T$.
        Let $E = \{V\in\pow{X} \mid \text{$V$ is open in $X$ with respect to $T$ and there exists $U\in A$ such that $\closure{V}{X}{T}\subseteq U$}\}$.
        Show that $E$ is an open covering of $X$ with respect to $T$.
        \begin{subproof}
            Show that $E$ is a covering of $X$.
            \begin{subproof}
                Show that for all $x\in X$ there exists $V\in E$ such that $x\in V$.
                \begin{subproof}
                    Fix $x\in X$.
                    We have $x\in\unions{A}$ by \cref{open_covering,covering,subseteq}.
                    Take $U\in A$ such that $x\in U$ by \cref{unions_iff}.
                    Hence $U$ is open in $X$ with respect to $T$ by \cref{open_covering}.
                    We have $U$ is a neighborhood of $x$ in $X$ with respect to $T$ by \cref{open_in,neighborhood}.
                    Take $V$ such that $V$ is a neighborhood of $x$ in $X$ with respect to $T$ and $\closure{V}{X}{T}\subseteq U$ by \cref{regular_closure_neighborhood}.
                    $V$ is open in $X$ with respect to $T$ and $x\in V$ and $V\subseteq X$ by \cref{neighborhood,open_in,topology_on,subseteq,pow_iff}.
                    We have $V\in\pow{X}$ by \cref{powerset_intro}.
                    Therefore $V\in E$.
                \end{subproof}
                Therefore for all $x\in X$ we have $x\in\unions{E}$ by \cref{unions_iff}.
                Hence $X\subseteq\unions{E}$ by \cref{subseteq}.
            \end{subproof}
            We have $E$ is a covering of $X$ by \cref{covering}.
            For all $V\in E$ we have $V$ is open in $X$ with respect to $T$ by \cref{subseteq}.
            Therefore $E$ is an open covering of $X$ with respect to $T$ by \cref{open_covering}.
        \end{subproof}
        Take $C$ such that $C$ is a refinement of $E$ on $X$ and $C$ is locally finite in $X$ with respect to $T$ and $C$ is a covering of $X$ by \cref{locally_finite_refinement_property}.
        Let $B = \{D\in\pow{X} \mid \exists C_0\in C. D = \closure{C_0}{X}{T}\}$.
        We have $B$ is the closure collection of $C$ in $X$ with respect to $T$ by \cref{closure_collection_in}.
        Hence $B$ is locally finite in $X$ with respect to $T$ by \cref{locally_finite_closure_collection}.
        Show that $B$ is a closed refinement of $A$ on $X$ with respect to $T$.
        \begin{subproof}
            Show that $B$ is a refinement of $A$ on $X$.
            \begin{subproof}
                For all $U\in A$ we have $U\in\pow{X}$ by \cref{open_covering,open_in,topology_on,subseteq,pow_iff,powerset_intro}.
                We have $A\subseteq\pow{X}$ by \cref{subseteq}.
                $B\subseteq\pow{X}$ by \cref{closure_collection_in,subseteq}.
                Show that for all $D\in B$ there exists $U\in A$ such that $D\subseteq U$.
                \begin{subproof}
                    Fix $D\in B$.
                    Take $C_0\in C$ such that $D = \closure{C_0}{X}{T}$ by \cref{closure_collection_in}.
                    Take $V\in E$ such that $C_0\subseteq V$ by \cref{refinement_on}.
                    Take $U\in A$ such that $\closure{V}{X}{T}\subseteq U$ by \cref{subseteq}.
                    $C\subseteq\pow{X}$ by \cref{refinement_on}.
                    We have $C_0\in\pow{X}$ and $V\in\pow{X}$ by \cref{subseteq}.
                    We have $C_0\subseteq X$ and $V\subseteq X$ by \cref{powerset_elim}.
                    Therefore $\closure{V}{X}{T}$ is closed in $X$ with respect to $T$ by \cref{closure_closed_in}.
                    We have $V\subseteq\closure{V}{X}{T}$ by \cref{subseteq_closure}.
                    Therefore $C_0\subseteq\closure{V}{X}{T}$ by \cref{subseteq_transitive}.
                    Hence $\closure{C_0}{X}{T}\subseteq\closure{V}{X}{T}$ by \cref{closure_subseteq_closed}.
                    Thus $D\subseteq U$ by \cref{real_lt_subst_left,subseteq_transitive}.
                \end{subproof}
                Therefore $B$ is a refinement of $A$ on $X$ by \cref{refinement_on}.
            \end{subproof}
            For all $D\in B$ we have $D$ is closed in $X$ with respect to $T$ by \cref{closure_collection_in,refinement_on,subseteq,powerset_elim,closure_closed_in,real_lt_subst_left}.
            Therefore $B$ is a closed refinement of $A$ on $X$ with respect to $T$ by \cref{closed_refinement_on}.
        \end{subproof}
        Show that $B$ is a covering of $X$.
        \begin{subproof}
            Show that for all $x$ such that $x\in X$ we have $x\in\unions{B}$.
            \begin{subproof}
                Fix $x$ such that $x\in X$.
                We have $x\in\unions{C}$ by \cref{covering,subseteq}.
                Take $C_0\in C$ such that $x\in C_0$ by \cref{unions_iff}.
                $C\subseteq\pow{X}$ by \cref{refinement_on}.
                We have $C_0\in\pow{X}$ by \cref{subseteq}.
                $C_0\subseteq X$ by \cref{powerset_elim}.
                $x\in\closure{C_0}{X}{T}$ by \cref{subseteq_closure,subseteq}.
                $\closure{C_0}{X}{T}\subseteq X$ by \cref{closed_whole_space,closure_subseteq_closed}.
                $\closure{C_0}{X}{T}\in\pow{X}$ by \cref{powerset_intro}.
                $\closure{C_0}{X}{T}\in B$ by \cref{closure_collection_in}.
                Therefore $x\in\unions{B}$ by \cref{unions_iff}.
            \end{subproof}
            Hence $X\subseteq\unions{B}$ by \cref{subseteq}.
            Therefore $B$ is a covering of $X$ by \cref{covering}.
        \end{subproof}
    \end{subproof}
    Therefore $X$ has the locally finite closed refinement property with respect to $T$ by \cref{locally_finite_closed_refinement_property}.
\end{proof}

% from §41 Item 4: locally finite closed refinements imply open locally finite refinements
% from §41 helper definition 9: expansion subcollection relative to a closed set
\begin{definition}\label{paracompact_expansion_subcollection}
    $\paracompactExpansionSubcollection{C}{B}{X} = \{D\in C \mid D\subseteq X\setminus B\}$.
\end{definition}

% from §41 helper definition 10: expansion open set relative to a closed set
\begin{definition}\label{paracompact_expansion_open_set}
    $\paracompactExpansionOpenSet{C}{B}{X} = X\setminus\unions{\paracompactExpansionSubcollection{C}{B}{X}}$.
\end{definition}

\begin{lemma}\label{locally_finite_closed_subcollection_union_closed}
    Assume $T$ is a topology on $X$.
    Assume $A$ is locally finite in $X$ with respect to $T$.
    Suppose for all $D\in A$ we have $D$ is closed in $X$ with respect to $T$.
    Assume $C\subseteq A$.
    Then $\unions{C}$ is closed in $X$ with respect to $T$.
\end{lemma}
\begin{proof}
    $C$ is locally finite in $X$ with respect to $T$ by \cref{locally_finite_subcollection}.
    Let $H = \{D\in\pow{X} \mid \exists E\in C. D = \closure{E}{X}{T}\}$.
    We have $H$ is the closure collection of $C$ in $X$ with respect to $T$ by \cref{closure_collection_in}.
    Therefore $\closure{\unions{C}}{X}{T} = \unions{H}$ by \cref{closure_of_union_locally_finite}.
    Show that $H\subseteq C$.
    \begin{subproof}
        Follows by \cref{closure_collection_in,subseteq,locally_finite_collection,powerset_elim,subseteq_refl,closure_subseteq_closed,subseteq_closure,real_lt_subst_left,real_lt_subst_right,subseteq_antisymmetric}.
    \end{subproof}
    Show that $C\subseteq H$.
    \begin{subproof}
        Follows by \cref{closure_collection_in,subseteq,locally_finite_collection,powerset_elim,subseteq_refl,closure_subseteq_closed,subseteq_closure,subseteq_antisymmetric,powerset_intro,real_lt_subst_left}.
    \end{subproof}
    Hence $\closure{\unions{C}}{X}{T} = \unions{C}$ by \cref{subseteq_antisymmetric,real_lt_subst_right}.
    Therefore $\unions{C}$ is closed in $X$ with respect to $T$ by
        \cref{locally_finite_collection,unions_subseteq_of_powerset_is_subseteq,closure_closed_in,real_lt_subst_left}.
\end{proof}

\begin{lemma}\label{locally_finite_closed_refinement_implies_paracompact}
    Assume $X$ is regular with respect to $T$.
    Assume $X$ has the locally finite closed refinement property with respect to $T$.
    Then $X$ is paracompact with respect to $T$.
\end{lemma}
\begin{proof}
    Show that $T$ is a topology on $X$.
    \begin{subproof}
        Follows by \cref{locally_finite_closed_refinement_property}.
    \end{subproof}
    Show that for all $A$ if
        $A$ is an open covering of $X$ with respect to $T$
    then there exists $D$ such that
        $D$ is an open refinement of $A$ on $X$ with respect to $T$ and
        $D$ is locally finite in $X$ with respect to $T$ and
        $D$ is a covering of $X$.
    \begin{subproof}
        Fix $A$ such that $A$ is an open covering of $X$ with respect to $T$.
        Take $B$ such that
            $B$ is a closed refinement of $A$ on $X$ with respect to $T$ and
            $B$ is locally finite in $X$ with respect to $T$ and
            $B$ is a covering of $X$
            by \cref{locally_finite_closed_refinement_property}.
        Let $M = \{U\in\pow{X} \mid \text{$U$ is open in $X$ with respect to $T$ and
            there exists $F\subseteq B$ such that
                $F$ is finite and
                for all $E\in B$ if $U$ intersects $E$ then $E\in F$}\}$.
        For all $x\in X$ we have $x\in\unions{M}$ by
            \cref{locally_finite_collection,neighborhood,open_in,topology_on,subseteq,pow_iff,powerset_intro,unions_iff}.
        Hence $M$ is an open covering of $X$ with respect to $T$ by \cref{subseteq,covering,open_covering}.
        Take $C$ such that
            $C$ is a closed refinement of $M$ on $X$ with respect to $T$ and
            $C$ is locally finite in $X$ with respect to $T$ and
            $C$ is a covering of $X$
            by \cref{locally_finite_closed_refinement_property}.
        For all $E\in C$ there exists $U\in M$ such that $E\subseteq U$
            by \cref{closed_refinement_on,refinement_on}.
        For all $E\in C$ there exists $F$ such that
            $F\subseteq B$ and
            $F$ is finite and
            for all $D\in B$ if $E$ intersects $D$ then $D\in F$
            by \cref{subseteq,intersects,nonempty_inhabited,inter,inter_intro,emptyset}.
        Let $R = \{w\in B\times A \mid \fst{w}\subseteq \snd{w}\}$.
        For all $w$ such that $w\in R$ there exist $D, U$ such that $w = (D,U)$.
        $R$ is a relation by \cref{relation}.
        $R\subseteq B\times A$ by \cref{subseteq}.
        For all $D\in B$ there exists $U$ such that $D\mathrel{R} U$
            by \cref{closed_refinement_on,refinement_on,times_tuple_intro,fst_eq,snd_eq}.
        Take $f$ such that $f$ is a function on $B$ and for all $D\in B$ we have $D\mathrel{R}\apply{f}{D}$ by \cref{choice_function}.
        Let $D = \{U\in\pow{X} \mid \exists E\in B. U = \paracompactExpansionOpenSet{C}{E}{X}\inter \apply{f}{E}\}$.
        Show that $D$ is a refinement of $A$ on $X$.
        \begin{subproof}
            Follows by
                \cref{open_covering,open_in,topology_on,subseteq,pow_iff,powerset_intro,choice_function,relation,times_tuple_elim,real_lt_subst_left,inter_lower_right,refinement_on}.
        \end{subproof}
        Show that for all $U\in D$ we have $U$ is open in $X$ with respect to $T$.
        \begin{subproof}
            Fix $U\in D$.
            Take $E\in B$ such that $U = \paracompactExpansionOpenSet{C}{E}{X}\inter \apply{f}{E}$ by \cref{subseteq}.
            Let $P = \paracompactExpansionSubcollection{C}{E}{X}$.
            $P\subseteq C$ by \cref{paracompact_expansion_subcollection,subseteq}.
            For all $V\in C$ we have $V$ is closed in $X$ with respect to $T$ by \cref{closed_refinement_on}.
            $\unions{P}$ is closed in $X$ with respect to $T$ by \cref{locally_finite_closed_subcollection_union_closed}.
            Hence $\paracompactExpansionOpenSet{C}{E}{X}$ is open in $X$ with respect to $T$ by \cref{paracompact_expansion_open_set,closed_in}.
            $\apply{f}{E}$ is open in $X$ with respect to $T$
                by \cref{choice_function,relation,subseteq,times_tuple_elim,open_covering}.
            Therefore $U$ is open in $X$ with respect to $T$ by \cref{real_lt_subst_left,topology_on,inter_subseteq,open_in,pow_iff}.
        \end{subproof}
        Show that for all $x\in X$ we have $x\in\unions{D}$.
        \begin{subproof}
            Fix $x\in X$.
            Take $E\in B$ such that $x\in E$ by \cref{covering,subseteq,unions_iff}.
            Let $U = \paracompactExpansionOpenSet{C}{E}{X}\inter \apply{f}{E}$.
            $U\in D$ by \cref{choice_function,relation,subseteq,times_tuple_elim,open_covering,open_in,topology_on,pow_iff,paracompact_expansion_open_set,setminus_subseteq,inter_subseteq,subseteq_transitive,powerset_intro}.
            $x\in \paracompactExpansionOpenSet{C}{E}{X}$ by
                \cref{unions_iff,paracompact_expansion_subcollection,subseteq,setminus_elim_right,paracompact_expansion_open_set,setminus_intro}.
            $x\in U$ by \cref{choice_function,relation,fst_eq,snd_eq,subseteq,inter_intro}.
            Therefore $x\in\unions{D}$ by \cref{unions_iff}.
        \end{subproof}
        $D$ is a covering of $X$ by \cref{subseteq,covering}.
        Hence $D$ is an open refinement of $A$ on $X$ with respect to $T$ by \cref{open_refinement_on}.
        Show that for all $x\in X$ there exist $W, G$ such that
            $W$ is a neighborhood of $x$ in $X$ with respect to $T$ and
            $G\subseteq D$ and
            $G$ is finite and
            for all $U\in D$ if $W$ intersects $U$ then $U\in G$.
        \begin{subproof}
            Fix $x\in X$.
            Take $W, K$ such that
                $W$ is a neighborhood of $x$ in $X$ with respect to $T$ and
                $K\subseteq C$ and
                $K$ is finite and
                for all $E\in C$ if $W$ intersects $E$ then $E\in K$
                by \cref{locally_finite_collection}.
            Let $S = \{w\in K\times\pow{B} \mid
                \text{$\snd{w}$ is finite and
                for all $E\in B$ if $\fst{w}$ intersects $E$ then $E\in\snd{w}$}\}$.
            For all $w$ such that $w\in S$ there exist $E, F$ such that $w = (E,F)$.
            $S$ is a relation by \cref{relation}.
            For all $E\in K$ there exists $F$ such that $E\mathrel{S} F$
                by \cref{subseteq,powerset_intro,times_tuple_intro,fst_eq,snd_eq}.
            Take $g$ such that $g$ is a function on $K$ and for all $E\in K$ we have $E\mathrel{S}\apply{g}{E}$ by \cref{choice_function}.
            Let $g_0 = \restrl{g}{K}$.
            $g_0$ is a function on $K$ by \cref{restrl_of_function_is_function,dom_of_restrl_eq_inter,inter_absorb_supseteq_left,subseteq,funs}.
            Let $L = \img{g_0}{K}$.
            For all $F\in L$ we have $F$ is finite by
                \cref{img_iff,function_apply_iff,restrl_of_function_apply,subseteq,choice_function,relation,fst_eq,snd_eq,real_lt_subst_left}.
            $L\subseteq\pow{B}$ by
                \cref{img_iff,function_apply_iff,restrl_of_function_apply,subseteq,choice_function,relation,times_tuple_elim,real_lt_subst_left}.
            $L$ is finite by \cref{finite_image_function}.
            $\unions{L}$ is finite by \cref{finite_union_of_finite_family}.
            Let $F = \unions{L}$.
            $F\subseteq B$ by \cref{unions_subseteq_of_powerset_is_subseteq}.
            Let $h = \{w\in F\times\pow{X} \mid \snd{w} = \paracompactExpansionOpenSet{C}{\fst{w}}{X}\inter \apply{f}{\fst{w}}\}$.
            For all $w$ such that $w\in h$ there exist $E, U$ such that $w = (E,U)$.
            $h$ is a relation by \cref{relation}.
            Show that for all $E\in F$ there exists $U$ such that $E\mathrel{h} U$.
            \begin{subproof}
                Fix $E\in F$.
                $E\in B$ by \cref{subseteq}.
                Let $U = \paracompactExpansionOpenSet{C}{E}{X}\inter \apply{f}{E}$.
                $(E,\apply{f}{E})\in R$ by \cref{choice_function}.
                $\apply{f}{E}\in A$ by \cref{subseteq,times_tuple_elim}.
                $\apply{f}{E}$ is open in $X$ with respect to $T$ by \cref{open_covering}.
                $\apply{f}{E}\subseteq X$ by \cref{open_in,topology_on,pow_iff,subseteq}.
                $\paracompactExpansionOpenSet{C}{E}{X}\subseteq X$
                    by \cref{paracompact_expansion_open_set,setminus_subseteq}.
                $U\subseteq X$ by \cref{inter_subseteq}.
                $U\in\pow{X}$ by \cref{powerset_intro}.
                $(E,U)\in F\times\pow{X}$ by \cref{times_tuple_intro}.
                $(E,U)\in h$ by \cref{fst_eq,snd_eq}.
                Therefore $E\mathrel{h} U$ by \cref{pair_eq_iff}.
            \end{subproof}
            For all $E,U_1,U_2$ such that $(E,U_1)\in h$ and $(E,U_2)\in h$ we have $U_1 = U_2$
                by \cref{fst_eq,snd_eq}.
            $h$ is right-unique by \cref{rightunique}.
            $F\subseteq\dom{h}$ by \cref{dom_iff,relation,subseteq}.
            $\dom{h}\subseteq F$ by \cref{dom_iff,subseteq,times_tuple_elim,subseteq}.
            $\dom{h} = F$ by \cref{subseteq_antisymmetric}.
            Thus $h$ is a function on $F$ by \cref{funs,subseteq}.
            Let $G = \img{h}{F}$.
            $G$ is finite by \cref{finite_image_function}.
            $G\subseteq D$ by \cref{img_iff,function_apply_iff,function_apply_elim,fst_eq,snd_eq,real_lt_subst_left,subseteq,choice_function,relation,times_tuple_elim,open_covering,open_in,topology_on,pow_iff,paracompact_expansion_open_set,setminus_subseteq,real_lt_subst_right,inter_subseteq,subseteq_transitive,powerset_intro}.
            Show that for all $U\in D$ if $W$ intersects $U$ then $U\in G$.
            \begin{subproof}
                Fix $U\in D$.
                Assume $W$ intersects $U$.
                Take $D_0\in B$ such that $U = \paracompactExpansionOpenSet{C}{D_0}{X}\inter \apply{f}{D_0}$ by \cref{subseteq}.
                Take $z$ such that $z\in W\inter U$ by \cref{intersects,nonempty_inhabited}.
                $z\in W$ and $z\in U$ by \cref{inter}.
                $z\in X$ by \cref{neighborhood,open_in,topology_on,subseteq,pow_iff}.
                Take $E\in C$ such that $z\in E$ by \cref{covering,subseteq,unions_iff}.
                $W$ intersects $E$ by \cref{intersects,inter_intro,emptyset}.
                $E\in K$ by \cref{locally_finite_collection}.
                Show that $E$ intersects $D_0$.
                \begin{subproof}
                    Suppose not.
                    $E\in\paracompactExpansionSubcollection{C}{D_0}{X}$ by
                        \cref{closed_refinement_on,refinement_on,subseteq,powerset_elim,intersects,inter_intro,emptyset,setminus_intro,paracompact_expansion_subcollection}.
                    $z\in\unions{\paracompactExpansionSubcollection{C}{D_0}{X}}$ by \cref{subseteq,unions_iff}.
                    $z\notin\paracompactExpansionOpenSet{C}{D_0}{X}$ by \cref{paracompact_expansion_open_set,setminus_elim_right}.
                    $z\in\paracompactExpansionOpenSet{C}{D_0}{X}$ by \cref{inter,real_lt_subst_right}.
                    Contradiction.
                \end{subproof}
                $E\mathrel{S}\apply{g}{E}$ by \cref{choice_function}.
                $D_0\in\apply{g}{E}$ by \cref{relation,fst_eq,snd_eq}.
                $\apply{g_0}{E} = \apply{g}{E}$ by \cref{restrl_of_function_apply,subseteq}.
                $\apply{g_0}{E}\in L$ by \cref{img_iff,function_apply_elim}.
                $D_0\in F$ by \cref{unions_iff}.
                $\apply{h}{D_0}\in G$ by \cref{img_iff,function_apply_elim}.
                $\apply{h}{D_0} = \paracompactExpansionOpenSet{C}{D_0}{X}\inter \apply{f}{D_0}$ by \cref{function_apply_elim,fst_eq,snd_eq}.
                Therefore $U\in G$ by \cref{real_lt_subst_left,real_lt_subst_right}.
            \end{subproof}
            $W$ is a neighborhood of $x$ in $X$ with respect to $T$ and
                $G\subseteq D$ and
                $G$ is finite and
                for all $U\in D$ if $W$ intersects $U$ then $U\in G$.
        \end{subproof}
        Hence $D$ is locally finite in $X$ with respect to $T$ by \cref{open_refinement_on,refinement_on,locally_finite_collection}.
        $D$ is an open refinement of $A$ on $X$ with respect to $T$ and
            $D$ is locally finite in $X$ with respect to $T$ and
            $D$ is a covering of $X$ by \cref{open_refinement_on}.
    \end{subproof}
    Therefore $X$ is paracompact with respect to $T$ by \cref{paracompact_space}.
\end{proof}

% from §41 helper lemma 4: subcollections of countably locally finite collections
\begin{lemma}\label{countably_locally_finite_subcollection}
    Assume $A$ is countably locally finite in $X$ with respect to $T$.
    Assume $B\subseteq A$.
    Then $B$ is countably locally finite in $X$ with respect to $T$.
\end{lemma}
\begin{proof}
    $B\subseteq\pow{X}$ by \cref{countably_locally_finite_collection,subseteq_transitive}.
    Take $c$ such that
        $c$ is a function on $\naturalsPlus$ and
        $A = \unions{\img{c}{\naturalsPlus}}$ and
        for all $n\in\naturalsPlus$ we have $\apply{c}{n}$ is locally finite in $X$ with respect to $T$
        by \cref{countably_locally_finite_collection}.
    Let $R = \{w\in\naturalsPlus\times\pow{\pow{X}} \mid \snd{w} = \apply{c}{\fst{w}}\inter B\}$.
    For all $w$ such that $w\in R$ there exist $m, D$ such that $w = (m,D)$.
    $R$ is a relation by \cref{relation}.
    For all $n\in\naturalsPlus$ there exists $D$ such that $n\mathrel{R} D$
        by \cref{locally_finite_collection,inter_subseteq,powerset_intro,times_tuple_intro,fst_eq,snd_eq}.
    Take $d$ such that $d$ is a function on $\naturalsPlus$ and for all $n\in\naturalsPlus$ we have $n\mathrel{R}\apply{d}{n}$ by \cref{choice_function}.
    Show that $\unions{\img{d}{\naturalsPlus}}\subseteq B$.
    \begin{subproof}
        Follows by \cref{unions_iff,img_iff,function_apply_iff,choice_function,relation,fst_eq,snd_eq,real_lt_subst_left,inter_elim_right,subseteq}.
    \end{subproof}
    Show that $B\subseteq\unions{\img{d}{\naturalsPlus}}$.
    \begin{subproof}
        Follows by \cref{subseteq,real_lt_subst_left,unions_iff,img_iff,function_apply_iff,choice_function,relation,fst_eq,snd_eq,inter_intro,function_apply_elim}.
    \end{subproof}
    We have $B = \unions{\img{d}{\naturalsPlus}}$ by \cref{subseteq_antisymmetric}.
    Hence $B$ is countably locally finite in $X$ with respect to $T$ by
        \cref{choice_function,relation,fst_eq,snd_eq,inter_lower_left,real_lt_subst_left,locally_finite_subcollection,countably_locally_finite_collection}.
\end{proof}

% from §41 helper lemma 5: countably locally finite basis yields open refinement property
\begin{lemma}\label{countably_locally_finite_basis_open_refinement_property}
    Assume $B$ is a basis for $T$ on $X$.
    Assume $B$ is countably locally finite in $X$ with respect to $T$.
    Then $X$ has the countably locally finite open refinement property with respect to $T$.
\end{lemma}
\begin{proof}
    $T$ is a topology on $X$ by \cref{basis_topology_is_topology,basis_for_topology}.
    Show that for all $A$ if
        $A$ is an open covering of $X$ with respect to $T$
    then there exists $C$ such that
        $C$ is an open refinement of $A$ on $X$ with respect to $T$ and
        $C$ is countably locally finite in $X$ with respect to $T$ and
        $C$ is a covering of $X$.
    \begin{subproof}
        Fix $A$ such that $A$ is an open covering of $X$ with respect to $T$.
        Let $C = \{U\in B \mid \exists V\in A. U\subseteq V\}$.
        $C$ is countably locally finite in $X$ with respect to $T$
            by \cref{subseteq,countably_locally_finite_subcollection}.
        Show that $C$ is a covering of $X$.
        \begin{subproof}
            Follows by \cref{open_covering,covering,subseteq,unions_iff,basis_for_topology,topology_generated_by_basis,open_in}.
        \end{subproof}
        Hence $C$ is an open refinement of $A$ on $X$ with respect to $T$ by \cref{subseteq,open_covering,open_in,topology_on,pow_iff,basis_elem_open_in_generated,basis_for_topology,powerset_intro,refinement_on,open_refinement_on}.
        $C$ is an open refinement of $A$ on $X$ with respect to $T$ and
            $C$ is countably locally finite in $X$ with respect to $T$ and
            $C$ is a covering of $X$ by \cref{open_refinement_on}.
    \end{subproof}
    Therefore $X$ has the countably locally finite open refinement property with respect to $T$
        by \cref{countably_locally_finite_open_refinement_property}.
\end{proof}

% from §41 Item 5: metrizable spaces are paracompact
\begin{theorem}\label{metrizable_implies_paracompact}
    Assume $X$ is metrizable with respect to $T$.
    Then $X$ is paracompact with respect to $T$.
\end{theorem}
\begin{proof}
    $X$ is regular with respect to $T$ by \cref{nagata_smirnov_metrization_theorem}.
    Take $B$ such that
        $B$ is a basis for $T$ on $X$ and
        $B$ is countably locally finite in $X$ with respect to $T$
        by \cref{nagata_smirnov_metrization_theorem}.
    Therefore $X$ is paracompact with respect to $T$
        by \cref{countably_locally_finite_basis_open_refinement_property,countably_locally_finite_open_refinement_implies_locally_finite_refinement_property,locally_finite_refinement_implies_locally_finite_closed_refinement_property,locally_finite_closed_refinement_implies_paracompact}.
\end{proof}

% from §41 helper lemma 6: countable indexed families are countably locally finite collections
\begin{lemma}\label{countable_indexed_family_countably_locally_finite}
    Assume $T$ is a topology on $X$.
    Assume $J\subseteq\naturalsPlus$.
    Assume $f$ is a function on $J$.
    Suppose for all $j\in J$ we have $\apply{f}{j}\subseteq X$.
    Then $\img{f}{J}$ is countably locally finite in $X$ with respect to $T$.
\end{lemma}
\begin{proof}
    Let $R = \{w\in\naturalsPlus\times\pow{\pow{X}} \mid
        ((\fst{w}\in J \land \snd{w} = \cons{\apply{f}{\fst{w}}}{\emptyset}) \lor
        (\fst{w}\notin J \land \snd{w} = \emptyset))\}$.
    For all $w$ such that $w\in R$ there exist $n, C$ such that $w = (n,C)$.
    $R$ is a relation by \cref{relation}.
    Show that for all $n\in\naturalsPlus$ there exists $C$ such that $n\mathrel{R} C$.
    \begin{subproof}
        Fix $n\in\naturalsPlus$.
        \begin{byCase}
        \caseOf{$n\in J$.}
            We have $n\mathrel{R} \cons{\apply{f}{n}}{\emptyset}$ by
                \cref{powerset_intro,singleton_subset_intro,times_tuple_intro,fst_eq,snd_eq}.
        \caseOf{$n\notin J$.}
            We have $n\mathrel{R}\emptyset$ by \cref{emptyset_subseteq,powerset_intro,times_tuple_intro,fst_eq,snd_eq}.
        \end{byCase}
    \end{subproof}
    Take $b$ such that $b$ is a function on $\naturalsPlus$ and for all $n\in\naturalsPlus$ we have $n\mathrel{R}\apply{b}{n}$ by \cref{choice_function}.
    Show that for all $U$ such that $U\in\img{f}{J}$ we have $U\in\unions{\img{b}{\naturalsPlus}}$.
    \begin{subproof}
        Follows by \cref{img_iff,function_apply_iff,subseteq,choice_function,relation,fst_eq,snd_eq,real_lt_subst_left,singleton_intro,function_apply_elim,unions_iff}.
    \end{subproof}
    $\img{f}{J}\subseteq\unions{\img{b}{\naturalsPlus}}$ by \cref{subseteq}.
    Show that for all $U$ such that $U\in\unions{\img{b}{\naturalsPlus}}$ we have $U\in\img{f}{J}$.
    \begin{subproof}
        Follows by \cref{unions_iff,img_iff,function_apply_iff,choice_function,relation,fst_eq,snd_eq,real_lt_subst_left,emptyset,singleton_iff}.
    \end{subproof}
    $\unions{\img{b}{\naturalsPlus}}\subseteq\img{f}{J}$ by \cref{subseteq}.
    Therefore $\img{f}{J} = \unions{\img{b}{\naturalsPlus}}$ by \cref{subseteq_antisymmetric}.
    Show that for all $n\in\naturalsPlus$ we have $\apply{b}{n}$ is locally finite in $X$ with respect to $T$.
    \begin{subproof}
        Fix $n\in\naturalsPlus$.
        We have $n\mathrel{R}\apply{b}{n}$ by \cref{choice_function}.
        \begin{byCase}
        \caseOf{$n\in J$.}
            We have $\apply{b}{n} = \cons{\apply{f}{n}}{\emptyset}$ by \cref{relation,fst_eq,snd_eq}.
            $\apply{b}{n}\subseteq\pow{X}$ by \cref{singleton_subset_intro,powerset_intro,real_lt_subst_left}.
            For all $x\in X$ there exist $U, F$ such that
                $U$ is a neighborhood of $x$ in $X$ with respect to $T$ and
                $F\subseteq\apply{b}{n}$ and
                $F$ is finite and
                for all $C\in\apply{b}{n}$ if $U$ intersects $C$ then $C\in F$
                by \cref{topology_on,open_in,neighborhood,subseteq_refl,real_lt_subst_left,singleton_subset_intro,pair_is_finite,upair_intro_left,subseteq_finite_is_finite}.
            Hence $\apply{b}{n}$ is locally finite in $X$ with respect to $T$ by \cref{locally_finite_collection}.
        \caseOf{$n\notin J$.}
            We have $\apply{b}{n} = \emptyset$ by \cref{relation,fst_eq,snd_eq}.
            $\apply{b}{n}\subseteq\pow{X}$ by \cref{emptyset_subseteq}.
            For all $x\in X$ there exist $U, F$ such that
                $U$ is a neighborhood of $x$ in $X$ with respect to $T$ and
                $F\subseteq\apply{b}{n}$ and
                $F$ is finite and
                for all $C\in\apply{b}{n}$ if $U$ intersects $C$ then $C\in F$
                by \cref{topology_on,open_in,neighborhood,emptyset_subseteq,emptyset_is_finite,real_lt_subst_left}.
            Therefore $\apply{b}{n}$ is locally finite in $X$ with respect to $T$ by \cref{locally_finite_collection}.
        \end{byCase}
    \end{subproof}
    $\img{f}{J}\subseteq\pow{X}$
        by \cref{subseteq,img_iff,function_apply_iff,powerset_intro,real_lt_subst_left}.
    Therefore $\img{f}{J}$ is countably locally finite in $X$ with respect to $T$ by \cref{countably_locally_finite_collection}.
\end{proof}

% from §41 Item 6: regular Lindelof spaces are paracompact
\begin{theorem}\label{regular_lindelof_implies_paracompact}
    Assume $X$ is regular with respect to $T$.
    Assume for all $A$ if
        $A$ is an open covering of $X$ with respect to $T$
    then there exists $J$ such that
        $J\subseteq\naturalsPlus$ and
        there exists $f$ such that
            $f$ is a function on $J$ and
            $X\subseteq\unions{\img{f}{J}}$ and
            for all $j\in J$ we have $\apply{f}{j}\in A$.
    Then $X$ is paracompact with respect to $T$.
\end{theorem}
\begin{proof}
    $T$ is a topology on $X$ by \cref{regular_space}.
    Show that for all $A$ if
        $A$ is an open covering of $X$ with respect to $T$
    then there exists $B$ such that
        $B$ is an open refinement of $A$ on $X$ with respect to $T$ and
        $B$ is countably locally finite in $X$ with respect to $T$ and
        $B$ is a covering of $X$.
    \begin{subproof}
        Fix $A$ such that $A$ is an open covering of $X$ with respect to $T$.
        Take $J$ such that
            $J\subseteq\naturalsPlus$ and
            there exists $f$ such that
                $f$ is a function on $J$ and
                $X\subseteq\unions{\img{f}{J}}$ and
                for all $j\in J$ we have $\apply{f}{j}\in A$.
        Take $f$ such that
            $f$ is a function on $J$ and
            $X\subseteq\unions{\img{f}{J}}$ and
            for all $j\in J$ we have $\apply{f}{j}\in A$.
        For all $j\in J$ we have $\apply{f}{j}\subseteq X$
            by \cref{open_covering,open_in,topology_on,subseteq,pow_iff}.
        $\img{f}{J}$ is countably locally finite in $X$ with respect to $T$
            by \cref{countable_indexed_family_countably_locally_finite}.
        Show that $\img{f}{J}$ is a covering of $X$.
        \begin{subproof}
            Follows by \cref{covering}.
        \end{subproof}
        Hence $\img{f}{J}$ is an open refinement of $A$ on $X$ with respect to $T$ by
            \cref{open_covering,open_in,topology_on,subseteq,pow_iff,img_iff,function_apply_iff,powerset_intro,real_lt_subst_left,subseteq_refl,refinement_on,open_refinement_on}.
        $\img{f}{J}$ is an open refinement of $A$ on $X$ with respect to $T$ and
            $\img{f}{J}$ is countably locally finite in $X$ with respect to $T$ and
            $\img{f}{J}$ is a covering of $X$ by \cref{open_refinement_on}.
    \end{subproof}
    Therefore $X$ is paracompact with respect to $T$
        by \cref{countably_locally_finite_open_refinement_property,countably_locally_finite_open_refinement_implies_locally_finite_refinement_property,locally_finite_refinement_implies_locally_finite_closed_refinement_property,locally_finite_closed_refinement_implies_paracompact}.
\end{proof}

% from §41 helper definition 4: real sum sequence along an enumeration
\begin{definition}\label{real_sum_sequence_along_enumeration}
    $a$ is a real sum sequence of $f$ at $x$ along $g$ with respect to $n$ and result $s$ iff
        $a\in\funs{\initialSegment{n}}{\reals}$ and
        $\sumFiniteSequence{a}{s}$ and
        for all $i\in\initialSegment{n}$ we have $a(i) = \apply{\apply{f}{g(i)}}{x}$.
\end{definition}

% from §41 helper definition 5: pointwise sum of an indexed family
\begin{definition}\label{pointwise_sum_of_indexed_family}
    $s$ is the pointwise sum of $f$ at $x$ on $X$ with respect to $T$ and $J$ iff
        $f$ is a function on $J$ and
        there exist $F, n, g, a$ such that
            $F\subseteq J$ and
            $F$ is finite and
            $F$ is inhabited and
            $n\in\naturalsPlus$ and
            $g$ is a bijection from $\initialSegment{n}$ to $F$ and
            $a$ is a real sum sequence of $f$ at $x$ along $g$ with respect to $n$ and result $s$ and
            for all $\alpha\in J$ if $x\in\supp{\apply{f}{\alpha}}{X}{T}$ then $\alpha\in F$.
\end{definition}

% from §41 helper definition 6: support family of an indexed family
\begin{definition}\label{support_family}
    $S$ is the support family of $f$ on $X$ with respect to $T$ and $J$ iff
        $S = \{(\alpha,\supp{\apply{f}{\alpha}}{X}{T}) \mid \alpha\in J\}$.
\end{definition}

% from §41 helper lemma 7: closure-bounded subfamilies of locally finite indexed families are locally finite
\begin{lemma}\label{locally_finite_indexed_family_closure_bound}
    Assume $A$ is a function on $J$.
    Assume $B$ is a function on $J$.
    Suppose for all $\alpha\in J$ we have $\apply{A}{\alpha}\subseteq X$ and $\apply{B}{\alpha}\subseteq X$.
    Suppose $A$ is a locally finite indexed family on $J$ in $X$ with respect to $T$.
    Suppose for all $\alpha\in J$ we have $\apply{B}{\alpha}\subseteq \closure{\apply{A}{\alpha}}{X}{T}$.
    Then $B$ is a locally finite indexed family on $J$ in $X$ with respect to $T$.
\end{lemma}
\begin{proof}
    For all $x\in X$ there exist $U,F$ such that
        $U$ is a neighborhood of $x$ in $X$ with respect to $T$ and
        $F\subseteq J$ and
        $F$ is finite and
        for all $\alpha\in J$ if $U$ intersects $\apply{A}{\alpha}$ then $\alpha\in F$
        by \cref{locally_finite_indexed_family}.
    Show that for all $x\in X$ there exist $U,F$ such that
        $U$ is a neighborhood of $x$ in $X$ with respect to $T$ and
        $F\subseteq J$ and
        $F$ is finite and
        for all $\alpha\in J$ if $U$ intersects $\apply{B}{\alpha}$ then $\alpha\in F$.
    \begin{subproof}
        Follows by \cref{intersects,nonempty_inhabited,inter,subseteq,neighborhood,open_in,topology_on,pow_iff,closure_iff_intersects_open,locally_finite_indexed_family}.
    \end{subproof}
    Hence $B$ is a locally finite indexed family on $J$ in $X$ with respect to $T$
        by \cref{indexed_family_on_in,locally_finite_indexed_family}.
\end{proof}

% from §41 helper lemma 8: locally finite open supports from local vanishing
\begin{lemma}\label{support_family_locally_finite_from_zero_outside}
    Assume $f$ is a function on $J$.
    Assume $V$ is a function on $J$.
    Suppose for all $\alpha\in J$ we have
        $\apply{f}{\alpha}$ is continuous from $X$ to $\reals$ with respect to $T$ and $\standardTopologyR$ and
        $\apply{V}{\alpha}$ is open in $X$ with respect to $T$ and
        for all $x\in X\setminus \apply{V}{\alpha}$ we have $\apply{\apply{f}{\alpha}}{x} = 0$.
    Suppose $V$ is a locally finite indexed family on $J$ in $X$ with respect to $T$.
    Assume $S$ is the support family of $f$ on $X$ with respect to $T$ and $J$.
    Then $S$ is a locally finite indexed family on $J$ in $X$ with respect to $T$.
\end{lemma}
\begin{proof}
    Show that for all $w$ such that $w\in S$ there exist $\alpha, U$ such that $w = (\alpha,U)$.
    \begin{subproof}
        Fix $w$ such that $w\in S$.
        Take $i\in J$ such that $w = (i,\supp{\apply{f}{i}}{X}{T})$
            by \cref{support_family}.
        There exist $\alpha, U$ such that $w = (\alpha,U)$.
    \end{subproof}
    $S$ is a relation by \cref{relation}.
    Show that for all $\alpha,U_1,U_2$ such that $(\alpha,U_1)\in S$ and $(\alpha,U_2)\in S$ we have $U_1 = U_2$.
    \begin{subproof}
        Fix $\alpha$.
        Fix $U_1$.
        Fix $U_2$.
        Assume $(\alpha,U_1)\in S$ and $(\alpha,U_2)\in S$.
        Take $\beta\in J$ such that $(\alpha,U_1) = (\beta,\supp{\apply{f}{\beta}}{X}{T})$
            by \cref{support_family}.
        Take $\gamma\in J$ such that $(\alpha,U_2) = (\gamma,\supp{\apply{f}{\gamma}}{X}{T})$
            by \cref{support_family}.
        Hence $U_1 = U_2$ by \cref{pair_eq_iff}.
    \end{subproof}
    $S$ is right-unique by \cref{rightunique}.
    For all $\alpha\in J$ we have $(\alpha,\supp{\apply{f}{\alpha}}{X}{T})\in S$
        by \cref{support_family,pair_eq_iff}.
    $J\subseteq\dom{S}$ by \cref{dom_intro,subseteq}.
    Show that for all $\alpha\in\dom{S}$ we have $\alpha\in J$.
    \begin{subproof}
        Fix $\alpha\in\dom{S}$.
        Take $U$ such that $(\alpha,U)\in S$ by \cref{dom_iff}.
        Take $\beta\in J$ such that $(\alpha,U) = (\beta,\supp{\apply{f}{\beta}}{X}{T})$
            by \cref{support_family}.
        Hence $\alpha\in J$ by \cref{pair_eq_iff}.
    \end{subproof}
    $\dom{S}\subseteq J$ by \cref{subseteq}.
    $\dom{S} = J$ by \cref{subseteq_antisymmetric}.
    Thus $S$ is a function on $J$
        by \cref{relation,rightunique,dom_iff,dom_intro,subseteq,subseteq_antisymmetric}.
    For all $\alpha\in J$ we have $\apply{S}{\alpha}\subseteq X$
        by \cref{support_family,pair_eq_iff,function_apply_intro,support,closure,subseteq,real_lt_subst_left}.
    For all $\alpha\in J$ we have $\apply{S}{\alpha}\subseteq \closure{\apply{V}{\alpha}}{X}{T}$
        by \cref{support_family,pair_eq_iff,function_apply_intro,support_subseteq_from_zero_outside,subseteq_refl,real_lt_subst_left}.
    Therefore $S$ is a locally finite indexed family on $J$ in $X$ with respect to $T$
        by \cref{locally_finite_indexed_family,indexed_family_on_in,locally_finite_indexed_family_closure_bound}.
\end{proof}

% from §41 helper lemma 9: finite inhabited sets admit positive initial-segment enumerations
\begin{lemma}\label{finite_inhabited_initial_segment_enumeration}
    Suppose $F$ is finite.
    Suppose $F$ is inhabited.
    Then there exist $n,g$ such that
        $n\in\naturalsPlus$ and
        $g$ is a bijection from $\initialSegment{n}$ to $F$.
\end{lemma}
\begin{proof}
    Take $n\in\naturalsPlus$ such that there exists a bijection from $\initialSegment{n}$ to $F$
        by \cref{finite,nonempty_inhabited}.
    Take $g$ such that $g$ is a bijection from $\initialSegment{n}$ to $F$.
    Therefore there exist $n_1,g_1$ such that
        $n_1\in\naturalsPlus$ and
        $g_1$ is a bijection from $\initialSegment{n_1}$ to $F$.
\end{proof}
% from §41 Item 7: partition of unity
\begin{definition}\label{indexed_partition_of_unity}
    $f$ is an indexed partition of unity on $X$ dominated by $U$ with respect to $T$ and $J$ iff
        $U$ is a function on $J$ and
        for all $\alpha\in J$ we have $\apply{U}{\alpha}$ is open in $X$ with respect to $T$ and
        $\img{U}{J}$ is a covering of $X$ and
        $f$ is a function on $J$ and
        for all $\alpha\in J$ there exists $I$ such that
            $I = \intervalclosed{0}{1}$ and
            $\apply{f}{\alpha}$ is continuous from $X$ to $I$ with respect to $T$ and $\subspaceTopology{\standardTopologyR}{I}$ and
            $\supp{\apply{f}{\alpha}}{X}{T}\subseteq \apply{U}{\alpha}$ and
        there exists $S$ such that
            $S$ is the support family of $f$ on $X$ with respect to $T$ and $J$ and
            $S$ is a locally finite indexed family on $J$ in $X$ with respect to $T$ and
        for all $x\in X$ we have $1$ is the pointwise sum of $f$ at $x$ on $X$ with respect to $T$ and $J$.
\end{definition}

% from §41 Item 8: shrinking lemma for paracompact Hausdorff spaces
\begin{lemma}\label{paracompact_shrinking_lemma}
    Assume $X$ is paracompact with respect to $T$.
    Assume $X$ is Hausdorff with respect to $T$.
    Assume $U$ is a function on $J$.
    Suppose for all $\alpha\in J$ we have $\apply{U}{\alpha}$ is open in $X$ with respect to $T$.
    Suppose $\img{U}{J}$ is a covering of $X$.
    Then there exists $V$ such that
        $V$ is a function on $J$ and
        $\img{V}{J}$ is a covering of $X$ and
        $V$ is a locally finite indexed family on $J$ in $X$ with respect to $T$ and
        for all $\alpha\in J$ we have $\apply{V}{\alpha}$ is open in $X$ with respect to $T$ and
        for all $\alpha\in J$ we have $\closure{\apply{V}{\alpha}}{X}{T}\subseteq \apply{U}{\alpha}$.
\end{lemma}
\begin{proof}
    $X$ is normal with respect to $T$ by \cref{paracompact_hausdorff_normal}.
    $X$ is regular with respect to $T$ by \cref{normal_implies_regular}.
    Let $E = \{V\in\pow{X} \mid \text{$V$ is open in $X$ with respect to $T$ and there exists $\alpha\in J$ such that $\closure{V}{X}{T}\subseteq \apply{U}{\alpha}$}\}$.
    For all $x\in X$ there exists $V\in E$ such that $x\in V$ by
        \cref{covering,subseteq,unions_iff,img_iff,function_apply_iff,open_in,neighborhood,regular_closure_neighborhood,topology_on,pow_iff,powerset_intro}.
    Hence $E$ is an open covering of $X$ with respect to $T$ by \cref{subseteq,covering,open_covering,unions_iff}.
    Take $D$ such that
        $D$ is an open refinement of $E$ on $X$ with respect to $T$ and
        $D$ is locally finite in $X$ with respect to $T$ and
        $D$ is a covering of $X$
        by \cref{paracompact_space}.
    Let $R = \{w\in D\times J \mid \closure{\fst{w}}{X}{T}\subseteq \apply{U}{\snd{w}}\}$.
    For all $w$ such that $w\in R$ there exist $A,\alpha$ such that $w = (A,\alpha)$.
    $R$ is a relation by \cref{relation}.
    For all $A\in D$ there exists $\alpha\in J$ such that $\closure{A}{X}{T}\subseteq \apply{U}{\alpha}$
        by \cref{open_refinement_on,refinement_on,subseteq,powerset_elim,regular_space,closure_closed_in,subseteq_closure,subseteq_transitive,closure_subseteq_closed}.
    For all $A\in D$ there exists $\alpha$ such that $A\mathrel{R}\alpha$ by \cref{times_tuple_intro,fst_eq,snd_eq}.
    Take $g$ such that $g$ is a function on $D$ and for all $A\in D$ we have $A\mathrel{R}\apply{g}{A}$ by \cref{choice_function}.
    Let $Q = \{w\in J\times\pow{D} \mid \text{for all $A\in D$ we have $A\in\snd{w}$ iff $\apply{g}{A} = \fst{w}$}\}$.
    For all $w$ such that $w\in Q$ there exist $\alpha,S$ such that $w = (\alpha,S)$.
    $Q$ is a relation by \cref{relation}.
    Show that for all $\alpha\in J$ there exists $S$ such that $\alpha\mathrel{Q}S$.
    \begin{subproof}
        Fix $\alpha\in J$.
        Let $S = \{A\in D \mid \apply{g}{A} = \alpha\}$.
        $\alpha\mathrel{Q}S$ by \cref{subseteq,powerset_intro,times_tuple_intro,fst_eq,snd_eq}.
    \end{subproof}
    Take $c$ such that $c$ is a function on $J$ and for all $\alpha\in J$ we have $\alpha\mathrel{Q}\apply{c}{\alpha}$ by \cref{choice_function}.
    Let $P = \{w\in J\times\pow{X} \mid \snd{w} = \unions{\apply{c}{\fst{w}}}\}$.
    For all $w$ such that $w\in P$ there exist $\alpha,V$ such that $w = (\alpha,V)$.
    $P$ is a relation by \cref{relation}.
    For all $\alpha\in J$ there exists $V$ such that $\alpha\mathrel{P}V$
        by \cref{choice_function,relation,times_tuple_elim,open_refinement_on,refinement_on,powerset_elim,subseteq_transitive,unions_subseteq_of_powerset_is_subseteq,powerset_intro,times_tuple_intro,fst_eq,snd_eq}.
    Take $V$ such that $V$ is a function on $J$ and for all $\alpha\in J$ we have $\alpha\mathrel{P}\apply{V}{\alpha}$ by \cref{choice_function}.
    For all $\alpha\in J$ we have $\apply{V}{\alpha}\subseteq X$ by \cref{choice_function,relation,times_tuple_elim,powerset_elim}.
    For all $\alpha\in J$ we have $\apply{c}{\alpha}\subseteq D$
        by \cref{choice_function,relation,times_tuple_elim,powerset_elim}.
    For all $\alpha\in J$ we have $\apply{c}{\alpha}\subseteq T$
        by \cref{subseteq,open_refinement_on,open_in}.
    $T$ is a topology on $X$ by \cref{regular_space}.
    For all $\alpha\in J$ we have $\unions{\apply{c}{\alpha}}\in T$ by \cref{topology_on}.
    For all $\alpha\in J$ we have $\apply{V}{\alpha} = \unions{\apply{c}{\alpha}}$
        by \cref{choice_function,relation,fst_eq,snd_eq}.
    Show that for all $\alpha\in J$ we have $\apply{V}{\alpha}$ is open in $X$ with respect to $T$.
    \begin{subproof}
        Follows by \cref{open_in,real_lt_subst_left}.
    \end{subproof}
    For all $A\in D$ we have $\apply{g}{A}\in J$ by \cref{choice_function,relation,times_tuple_elim}.
    For all $A\in D$ we have $A\in\apply{c}{\apply{g}{A}}$
        by \cref{choice_function,relation,fst_eq,snd_eq,subseteq_refl}.
    For all $A\in D$ we have $\apply{V}{\apply{g}{A}} = \unions{\apply{c}{\apply{g}{A}}}$
        by \cref{choice_function,relation,fst_eq,snd_eq}.
    For all $x\in X$ there exists $\alpha\in J$ such that $x\in\apply{V}{\alpha}$
        by \cref{covering,subseteq,unions_iff,real_lt_subst_right}.
    Show that $\img{V}{J}$ is a covering of $X$.
    \begin{subproof}
        Follows by \cref{subseteq,covering,unions_iff,img_iff,function_apply_elim}.
    \end{subproof}
    Show that for all $\alpha\in J$ we have $\closure{\apply{V}{\alpha}}{X}{T}\subseteq \apply{U}{\alpha}$.
    \begin{subproof}
        Fix $\alpha\in J$.
        $\apply{c}{\alpha}$ is locally finite in $X$ with respect to $T$ by
            \cref{choice_function,relation,times_tuple_elim,powerset_elim,locally_finite_subcollection}.
        Let $H = \{C\in\pow{X} \mid \exists A\in\apply{c}{\alpha}. C = \closure{A}{X}{T}\}$.
        $H$ is the closure collection of $\apply{c}{\alpha}$ in $X$ with respect to $T$ by \cref{closure_collection_in}.
        $\closure{\unions{\apply{c}{\alpha}}}{X}{T} = \unions{H}$ by \cref{closure_of_union_locally_finite}.
        For all $x$ such that $x\in\unions{H}$ we have $x\in \apply{U}{\alpha}$
            by \cref{unions_iff,closure_collection_in,choice_function,relation,fst_eq,snd_eq,subseteq,real_lt_subst_right,real_lt_subst_left}.
        Hence $\closure{\apply{V}{\alpha}}{X}{T}\subseteq \apply{U}{\alpha}$ by
            \cref{choice_function,relation,fst_eq,snd_eq,real_lt_subst_left,subseteq}.
    \end{subproof}
    Show that for all $x\in X$ there exist $W,F$ such that
        $W$ is a neighborhood of $x$ in $X$ with respect to $T$ and
        $F\subseteq J$ and
        $F$ is finite and
        for all $\alpha\in J$ if $W$ intersects $\apply{V}{\alpha}$ then $\alpha\in F$.
    \begin{subproof}
        Fix $x\in X$.
        Take $W,G_0$ such that
            $W$ is a neighborhood of $x$ in $X$ with respect to $T$ and
            $G_0\subseteq D$ and
            $G_0$ is finite and
            for all $A\in D$ if $W$ intersects $A$ then $A\in G_0$
            by \cref{locally_finite_collection}.
        Let $g_0 = \restrl{g}{G_0}$.
        $g_0$ is a function on $G_0$ by \cref{restrl_of_function_is_function,dom_of_restrl_eq_inter,subseteq,inter_absorb_supseteq_left,funs}.
        Let $F = \img{g_0}{G_0}$.
        $F$ is finite by \cref{finite_image_function}.
        $F\subseteq J$ by
            \cref{img_iff,function_apply_iff,subseteq,restrl_of_function_apply,choice_function,relation,times_tuple_elim,real_lt_subst_left}.
        Show that for all $\alpha\in J$ if $W$ intersects $\apply{V}{\alpha}$ then $\alpha\in F$.
        \begin{subproof}
            Fix $\alpha\in J$.
            Assume $W$ intersects $\apply{V}{\alpha}$.
            Take $z$ such that $z\in W\inter\apply{V}{\alpha}$ by \cref{intersects,nonempty_inhabited}.
            $z\in W$ and $z\in\unions{\apply{c}{\alpha}}$ by \cref{inter,choice_function,relation,fst_eq,snd_eq,real_lt_subst_left}.
            Take $A\in\apply{c}{\alpha}$ such that $z\in A$ by \cref{unions_iff}.
            $A\in D$ by \cref{subseteq}.
            $W$ intersects $A$ by \cref{inter_intro,emptyset,intersects}.
            $A\in G_0$ by \cref{locally_finite_collection}.
            $\alpha\in F$ by
                \cref{img_iff,function_apply_iff,restrl_of_function_apply,subseteq,choice_function,relation,fst_eq,snd_eq,real_lt_subst_right}.
        \end{subproof}
        $W$ is a neighborhood of $x$ in $X$ with respect to $T$ and
            $F\subseteq J$ and
            $F$ is finite and
            for all $\alpha\in J$ if $W$ intersects $\apply{V}{\alpha}$ then $\alpha\in F$.
    \end{subproof}
    Therefore $V$ is a locally finite indexed family on $J$ in $X$ with respect to $T$
        by \cref{indexed_family_on_in,locally_finite_indexed_family}.
\end{proof}

% from §41 helper lemma 10: local finite support control implies local vanishing
\begin{lemma}\label{support_family_local_zero_control}
    Assume $f$ is a function on $J$.
    Assume $S$ is the support family of $f$ on $X$ with respect to $T$ and $J$.
    Suppose $S$ is a locally finite indexed family on $J$ in $X$ with respect to $T$.
    Suppose for all $\alpha\in J$ we have $\apply{f}{\alpha}$ is continuous from $X$ to $\reals$ with respect to $T$ and $\standardTopologyR$.
    Suppose $x\in X$.
    Then there exist $W,F$ such that
        $W$ is a neighborhood of $x$ in $X$ with respect to $T$ and
        $F\subseteq J$ and
        $F$ is finite and
        for all $\alpha\in J$ if $\alpha\notin F$ then for all $y\in W$ we have $\apply{\apply{f}{\alpha}}{y} = 0$.
\end{lemma}
\begin{proof}
    Take $W,F$ such that
        $W$ is a neighborhood of $x$ in $X$ with respect to $T$ and
        $F\subseteq J$ and
        $F$ is finite and
        for all $\alpha\in J$ if $W$ intersects $\apply{S}{\alpha}$ then $\alpha\in F$
        by \cref{locally_finite_indexed_family}.
    Show that for all $\alpha\in J$ if $\alpha\notin F$ then for all $y\in W$ we have $\apply{\apply{f}{\alpha}}{y} = 0$.
    \begin{subproof}
        Follows by \cref{locally_finite_indexed_family,indexed_family_on_in,neighborhood,open_in,topology_on,subseteq,pow_iff,support_family,pair_eq_iff,function_apply_intro,real_lt_subst_right,inter_intro,intersects,emptyset,setminus_intro,outside_support_zero}.
    \end{subproof}
    $W$ is a neighborhood of $x$ in $X$ with respect to $T$ and
        $F\subseteq J$ and
        $F$ is finite and
        for all $\alpha\in J$ if $\alpha\notin F$ then for all $y\in W$ we have $\apply{\apply{f}{\alpha}}{y} = 0$.
\end{proof}

% from §41 helper lemma 11: appending a zero last term preserves a finite sum
\begin{lemma}\label{sum_finite_sequence_append_zero}
    Suppose $n\in\naturalsPlus$.
    Let $J = \initialSegment{n}$.
    Let $J_1 = \initialSegment{n + 1}$.
    Suppose $a\in\funs{J}{\reals}$.
    Suppose $\sumFiniteSequence{a}{s}$.
    Suppose $b\in\funs{J_1}{\reals}$.
    Suppose $\restrl{b}{J} = a$.
    Suppose $b(n + 1) = 0$.
    Then $\sumFiniteSequence{b}{s}$.
\end{lemma}
\begin{proof}
    Take $t\in\reals$ such that $\sumFiniteSequence{b}{t}$ by \cref{real_naturals_closed_succ,sum_finite_sequence_exists}.
    Take $t_0\in\reals$ such that $\sumFiniteSequence{\restrl{b}{J}}{t_0}$ and $t = t_0 + b(n + 1)$
        by \cref{sum_finite_sequence_split_last}.
    $t_0 = s$ by \cref{real_lt_subst_left,sum_finite_sequence_unique}.
    $t = s$ by \cref{real_lt_subst_left,sum_finite_sequence_value_in_reals,real_add_comm,real_add_ident,real_lt_subst_right}.
    Hence $\sumFiniteSequence{b}{s}$ by \cref{real_lt_subst_right}.
\end{proof}

% from §41 helper lemma 12: finite inhabited sets admit enumerations with prescribed last element
\begin{lemma}\label{finite_inhabited_initial_segment_enumeration_last}
    Suppose $F$ is finite.
    Suppose $\beta\in F$.
    Then there exist $n,g$ such that
        $n\in\naturalsPlus$ and
        $g$ is a bijection from $\initialSegment{n}$ to $F$ and
        $g(n) = \beta$.
\end{lemma}
\begin{proof}
    Let $E = F\setminus \{\beta\}$.
    \begin{byCase}
    \caseOf{$E = \emptyset$.}
        We have $F = \{\beta\}$ by \cref{setminus_intro,emptyset,subseteq,singleton_subset_intro,subseteq_antisymmetric}.
        Let $g = \{(1,\beta)\}$.
        $1\in\naturalsPlus$ by \cref{real_one_in_naturals}.
        $g$ is a function by \cref{singleton_elim,pair_eq_iff,rightunique,relation}.
        For all $u$ such that $u\in\dom{g}$ we have $u\in\initialSegment{1}$
            by \cref{dom_iff,singleton_elim,pair_eq_iff,real_one_in_naturals,initial_segment_naturalsplus}.
        For all $u$ such that $u\in\initialSegment{1}$ we have $u\in\dom{g}$
            by \cref{initial_segment_naturalsplus,real_one_leq_naturalsplus,real_leq_split,real_naturals_subset_reals,real_one_in_naturals,subseteq,real_lt_trans,real_lt_irreflexive,singleton_intro,pair_eq_iff,dom_intro}.
        Thus $\dom{g} = \initialSegment{1}$ by \cref{subseteq,subseteq_antisymmetric}.
        For all $u\in\dom{g}$ we have $g(u)\in F$
            by \cref{dom_iff,singleton_elim,pair_eq_iff,singleton_intro,function_apply_intro,subseteq}.
        $g\in\funs{\initialSegment{1}}{F}$ by \cref{funs_intro}.
        For all $y\in F$ there exists $u\in\initialSegment{1}$ such that $g(u) = y$
            by \cref{singleton_elim,subseteq,real_one_in_naturals,initial_segment_naturalsplus,real_one_leq_naturalsplus,singleton_intro,pair_eq_iff,function_apply_intro,real_lt_subst_left}.
        Hence $g\in\surj{\initialSegment{1}}{F}$ by \cref{surj}.
        For all $u_1,u_2$ such that $u_1\in\dom{g}$ and $u_2\in\dom{g}$ and $g(u_1) = g(u_2)$
            we have $u_1 = u_2$
            by \cref{initial_segment_naturalsplus,real_one_leq_naturalsplus,real_naturals_subset_reals,real_one_in_naturals,subseteq,real_lt_trans,real_lt_irreflexive}.
        $g$ is injective by \cref{funs_is_function,function_rightunique,relation,injective_function,subseteq}.
        Thus there exist $n,h$ such that
            $n\in\naturalsPlus$ and
            $h$ is a bijection from $\initialSegment{n}$ to $F$ and
            $h(n) = \beta$ by \cref{real_one_in_naturals,bijections,singleton_intro,pair_eq_iff,function_apply_intro}.
    \caseOf{$E\neq\emptyset$.}
        Take $m,p$ such that
            $m\in\naturalsPlus$ and
            $p$ is a bijection from $\initialSegment{m}$ to $E$
            by \cref{setminus_subseteq,subseteq_finite_is_finite,nonempty_inhabited,finite_inhabited_initial_segment_enumeration}.
        $\beta\notin E$ by \cref{setminus,singleton_intro}.
        Let $r = \{(m + 1,\beta)\}$.
        For all $u$ such that $u\in\dom{p}\inter\dom{r}$ we have $p(u) = r(u)$
            by \cref{inter_elim_left,inter_elim_right,bijections_dom,dom_iff,singleton_elim,pair_eq_iff,initial_segment_naturalsplus,real_naturals_subset_reals,subseteq,real_naturals_closed_succ,real_add_closed,real_one_in_naturals,real_naturals_lt_succ_local,real_lt_subst_right,real_lt_trans,real_lt_irreflexive}.
        $p$ is right-unique and $p$ is a relation by \cref{bijection_is_function,function_rightunique,relation}.
        $r$ is right-unique and $r$ is a relation by \cref{singleton_elim,pair_eq_iff,rightunique,relation}.
        Let $h = p\union r$.
        $h$ is right-unique and $h$ is a relation by \cref{union_compatible_functions}.
        $h$ is a function by \cref{relation,rightunique}.
        $\dom{h} = \initialSegment{m}\union \{m + 1\}$ by
            \cref{dom_union,bijections_dom,dom_iff,singleton_elim,pair_eq_iff,singleton_intro,dom_intro,subseteq,subseteq_antisymmetric}.
        For all $u$ such that $u\in\dom{h}$ we have $u\in\initialSegment{m + 1}$
            by \cref{union_iff,initial_segment_subset_succ,subseteq,singleton_elim,real_naturals_closed_succ,initial_segment_naturalsplus}.
        For all $u$ such that $u\in\initialSegment{m + 1}$ we have $u\in\dom{h}$
            by \cref{initial_segment_naturalsplus,real_naturals_leq_succ_elim,singleton_intro,union_iff,subseteq}.
        Thus $\dom{h} = \initialSegment{m + 1}$ by \cref{subseteq,subseteq_antisymmetric}.
        For all $u\in\dom{h}$ we have $h(u)\in F$
            by \cref{dom_union,union_iff,function_apply_elim,function_apply_intro,bijections,surj,funs_type_apply,setminus,subseteq,singleton_elim,singleton_intro,pair_eq_iff}.
        $h\in\funs{\initialSegment{m + 1}}{F}$ by \cref{funs_intro}.
        Show that for all $y\in F$ there exists $u\in\initialSegment{m + 1}$ such that $h(u) = y$.
        \begin{subproof}
            Fix $y\in F$.
            $y\in E\union \{\beta\}$.
            $y\in E$ or $y\in \{\beta\}$ by \cref{union_iff}.
            \begin{byCase}
            \caseOf{$y\in E$.}
                There exists $u\in\initialSegment{m}$ such that $p(u) = y$ by \cref{bijections,surj}.
                $u\in\initialSegment{m + 1}$ by \cref{initial_segment_subset_succ,subseteq}.
                $u\in\dom{p}$ by \cref{bijections_dom,subseteq}.
                $h(u) = p(u)$ by \cref{function_apply_elim,union_iff,function_apply_intro}.
                Therefore there exists $v\in\initialSegment{m + 1}$ such that $h(v) = y$ by \cref{real_lt_subst_right}.
            \caseOf{$y\in \{\beta\}$.}
                We have $y = \beta$ by \cref{singleton_elim}.
                $m + 1\in\initialSegment{m + 1}$ by \cref{real_naturals_closed_succ,initial_segment_naturalsplus}.
                $h(m + 1) = \beta$ by
                    \cref{real_naturals_closed_succ,initial_segment_naturalsplus,singleton_intro,pair_eq_iff,union_iff,function_apply_intro}.
                Hence there exists $v\in\initialSegment{m + 1}$ such that $h(v) = y$ by \cref{real_lt_subst_left}.
            \end{byCase}
        \end{subproof}
        $h\in\surj{\initialSegment{m + 1}}{F}$ by \cref{surj}.
        Show that for all $u_1,u_2$ such that $u_1\in\dom{h}$ and $u_2\in\dom{h}$ and $h(u_1) = h(u_2)$
            we have $u_1 = u_2$.
        \begin{subproof}
            Fix $u_1$.
            Fix $u_2$ such that $u_1\in\dom{h}$ and $u_2\in\dom{h}$ and $h(u_1) = h(u_2)$.
            \begin{byCase}
            \caseOf{$u_1 = m + 1$.}
                \begin{byCase}
                \caseOf{$u_2 = m + 1$.}
                    We have $u_1 = u_2$.
                \caseOf{$u_2 \neq m + 1$.}
                    $u_1 = u_2$
                        by \cref{initial_segment_naturalsplus,real_naturals_leq_succ_elim,bijections_dom,subseteq,function_apply_elim,union_iff,function_apply_intro,singleton_intro,pair_eq_iff,real_lt_subst_left,bijections,surj,funs_type_apply}.
                \end{byCase}
            \caseOf{$u_1 \neq m + 1$.}
                $u_1\in\naturalsPlus$ and $u_1 \leq m + 1$ by \cref{initial_segment_naturalsplus}.
                $u_1 \leq m$ by \cref{real_naturals_leq_succ_elim}.
                $u_1\in\initialSegment{m}$ by \cref{initial_segment_naturalsplus}.
                $u_1\in\dom{p}$ by \cref{bijections_dom,subseteq}.
                $h(u_1) = p(u_1)$ by \cref{function_apply_elim,union_iff,function_apply_intro}.
                \begin{byCase}
                \caseOf{$u_2 = m + 1$.}
                    Suppose not.
                    $h(u_2) = \beta$ by
                        \cref{singleton_intro,pair_eq_iff,union_iff,function_apply_intro,real_lt_subst_left}.
                    $p(u_1) = \beta$ by \cref{real_lt_subst_right}.
                    $\beta\in E$ by
                        \cref{bijections,surj,funs_type_apply,subseteq,real_lt_subst_left}.
                    Contradiction.
                \caseOf{$u_2 \neq m + 1$.}
                    $u_2\in\naturalsPlus$ and $u_2 \leq m + 1$ by \cref{initial_segment_naturalsplus}.
                    $u_2 \leq m$ by \cref{real_naturals_leq_succ_elim}.
                    $u_2\in\initialSegment{m}$ by \cref{initial_segment_naturalsplus}.
                    $u_2\in\dom{p}$ by \cref{bijections_dom,subseteq}.
                    $h(u_2) = p(u_2)$ by \cref{function_apply_elim,union_iff,function_apply_intro}.
                    $p(u_1) = p(u_2)$ by \cref{real_lt_subst_left,real_lt_subst_right}.
                    Hence $u_1 = u_2$ by \cref{bijections,injective_function}.
                \end{byCase}
            \end{byCase}
        \end{subproof}
        $h$ is injective by \cref{function_rightunique,relation,injective_function}.
        Therefore there exist $n,g$ such that
            $n\in\naturalsPlus$ and
            $g$ is a bijection from $\initialSegment{n}$ to $F$ and
            $g(n) = \beta$ by
            \cref{real_naturals_closed_succ,bijections,singleton_intro,pair_eq_iff,union_iff,function_apply_intro}.
    \end{byCase}
\end{proof}

% from §41 helper lemma 13: appending a last term adds that term to the sum
\begin{lemma}\label{sum_finite_sequence_append_last}
    Suppose $n\in\naturalsPlus$.
    Let $J = \initialSegment{n}$.
    Let $J_1 = \initialSegment{n + 1}$.
    Suppose $a\in\funs{J}{\reals}$.
    Suppose $\sumFiniteSequence{a}{s}$.
    Suppose $b\in\funs{J_1}{\reals}$.
    Suppose $\restrl{b}{J} = a$.
    Then $\sumFiniteSequence{b}{s + b(n + 1)}$.
\end{lemma}
\begin{proof}
    $n + 1\in\naturalsPlus$ by \cref{real_naturals_closed_succ}.
    Take $t\in\reals$ such that $\sumFiniteSequence{b}{t}$ by \cref{sum_finite_sequence_exists}.
    Take $t_0\in\reals$ such that $\sumFiniteSequence{\restrl{b}{J}}{t_0}$ and $t = t_0 + b(n + 1)$
        by \cref{sum_finite_sequence_split_last}.
    $t = s + b(n + 1)$ by \cref{real_lt_subst_left,sum_finite_sequence_unique}.
    Therefore $\sumFiniteSequence{b}{s + b(n + 1)}$
        by \cref{real_naturals_closed_succ,initial_segment_naturalsplus,sum_finite_sequence_value_in_reals,funs_type_apply,real_add_closed,real_lt_subst_right}.
\end{proof}

% from §41 helper lemma 14: the sum of a one-term sequence is its value
\begin{lemma}\label{sum_finite_sequence_one_term}
    Suppose $a\in\funs{\initialSegment{1}}{\reals}$.
    Then $\sumFiniteSequence{a}{a(1)}$.
\end{lemma}
\begin{proof}
    For all $k\in\naturalsPlus$ such that $k + 1 \leq 1$ we have $a(k + 1) = a(k) + a(k + 1)$
        by \cref{real_naturals_subset_reals,subseteq,real_one_in_naturals,real_add_closed,real_naturals_lt_succ_local,real_leq_split,real_lt_trans,real_lt_subst_right,real_one_leq_naturalsplus,real_lt_subst_left,real_lt_irreflexive}.
    Therefore $\sumFiniteSequence{a}{a(1)}$ by \cref{real_one_in_naturals,sum_finite_sequence_pred}.
\end{proof}

% from §41 helper lemma 15: vanishing on a neighborhood excludes the support
\begin{lemma}\label{local_zero_neighborhood_excludes_support}
    Assume $W$ is a neighborhood of $x$ in $X$ with respect to $T$.
    Suppose $f$ is continuous from $X$ to $\reals$ with respect to $T$ and $\standardTopologyR$.
    Suppose for all $y\in W$ we have $f(y) = 0$.
    Then $x\notin \supp{f}{X}{T}$.
\end{lemma}
\begin{proof}
    $W$ is open in $X$ with respect to $T$ and $x\in W$ by \cref{neighborhood}.
    $T$ is a topology on $X$ by \cref{open_in}.
    Let $A = X\setminus W$.
    We have $A$ is closed in $X$ with respect to $T$
        by \cref{open_in,topology_on,subseteq,pow_iff,setminus_subseteq,double_relative_complement,closed_in}.
    Let $S = X\setminus \preimg{f}{\{0\}}$.
    For all $y$ such that $y\in S$ we have $y\in A$
        by \cref{setminus,setminus_intro,continuous,funs_is_function,funs,function_apply_elim,preimg_iff,real_lt_subst_right,singleton_intro}.
    Thus $\closure{S}{X}{T}\subseteq A$ by \cref{subseteq,setminus_subseteq,closure_subseteq_closed}.
    Therefore $\supp{f}{X}{T}\subseteq A$ by \cref{support}.
    Finally $x\notin \supp{f}{X}{T}$ by \cref{setminus_elim_right,subseteq}.
\end{proof}

% from §41 helper lemma 16: local support coverage yields a pointwise sum witness
\begin{lemma}\label{pointwise_sum_exists_from_local_support_cover}
    Assume $f$ is a function on $J$.
    Assume $S$ is the support family of $f$ on $X$ with respect to $T$ and $J$.
    Suppose $S$ is a locally finite indexed family on $J$ in $X$ with respect to $T$.
    Suppose for all $\alpha\in J$ we have $\apply{f}{\alpha}$ is continuous from $X$ to $\reals$ with respect to $T$ and $\standardTopologyR$.
    Suppose $x\in X$.
    Suppose there exists $\beta\in J$ such that $x\in\apply{S}{\beta}$.
    Then there exists $s$ such that $s$ is the pointwise sum of $f$ at $x$ on $X$ with respect to $T$ and $J$.
\end{lemma}
\begin{proof}
    Take $\beta\in J$ such that $x\in\apply{S}{\beta}$.
    Take $W,F$ such that
        $W$ is a neighborhood of $x$ in $X$ with respect to $T$ and
        $F\subseteq J$ and
        $F$ is finite and
        for all $\alpha\in J$ if $\alpha\notin F$ then for all $y\in W$ we have $\apply{\apply{f}{\alpha}}{y} = 0$
        by \cref{support_family_local_zero_control}.
    Therefore $\beta\in F$
        by \cref{locally_finite_indexed_family,indexed_family_on_in,support_family_local_zero_control,support_family,pair_eq_iff,function_apply_intro,local_zero_neighborhood_excludes_support,real_lt_subst_right}.
    Take $n,g$ such that
        $n\in\naturalsPlus$ and
        $g$ is a bijection from $\initialSegment{n}$ to $F$
        by \cref{nonempty_inhabited,finite_inhabited_initial_segment_enumeration}.
    Let $R = \{w\in\initialSegment{n}\times\reals \mid w = ( \fst{w}, \apply{\apply{f}{g(\fst{w})}}{x})\}$.
    For all $i\in\initialSegment{n}$ we have $i\mathrel{R}\apply{\apply{f}{g(i)}}{x}$
        by \cref{bijections,surj,funs_type_apply,subseteq,continuous,times_tuple_intro,fst_eq}.
    For all $i\in\initialSegment{n}$ there exists $r$ such that $i\mathrel{R}r$.
    Take $a$ such that $a$ is a function on $\initialSegment{n}$ and for all $i\in\initialSegment{n}$ we have $i\mathrel{R}\apply{a}{i}$ by \cref{relation,choice_function}.
    For all $i\in\initialSegment{n}$ we have $a(i) = \apply{\apply{f}{g(i)}}{x}$
        by \cref{choice_function,fst_eq,snd_eq,pair_eq_iff}.
    We have $a\in\funs{\initialSegment{n}}{\reals}$
        by \cref{funs_intro,choice_function,funs_elim,bijections,surj,funs_type_apply,subseteq,continuous,real_lt_subst_left}.
    Take $s\in\reals$ such that $\sumFiniteSequence{a}{s}$ by \cref{sum_finite_sequence_exists}.
    Therefore $a$ is a real sum sequence of $f$ at $x$ along $g$ with respect to $n$ and result $s$
        by \cref{real_sum_sequence_along_enumeration}.
    Show that for all $\alpha\in J$ if $x\in\supp{\apply{f}{\alpha}}{X}{T}$ then $\alpha\in F$.
    \begin{subproof}
        Follows by \cref{support_family_local_zero_control,local_zero_neighborhood_excludes_support,real_lt_subst_left}.
    \end{subproof}
    Thus $s$ is the pointwise sum of $f$ at $x$ on $X$ with respect to $T$ and $J$
        by \cref{pointwise_sum_of_indexed_family}.
    Therefore there exists $t$ such that $t$ is the pointwise sum of $f$ at $x$ on $X$ with respect to $T$ and $J$
        by \cref{real_lt_subst_left}.
\end{proof}

% from §41 helper lemma 17: a zero-valued finite sequence sums to zero
\begin{lemma}\label{sum_finite_sequence_all_zero}
    Suppose $n\in\naturalsPlus$.
    Let $J = \initialSegment{n}$.
    Suppose $a\in\funs{J}{\reals}$.
    Suppose $\sumFiniteSequence{a}{s}$.
    Suppose for all $i\in J$ we have $a(i) = 0$.
    Then $s = 0$.
\end{lemma}
\begin{proof}
    $0\in\reals$ by \cref{real_add_ident}.
    Let $b = \coordScalarSequence{n}{0}{a}$.
    We have $b\in\funs{J}{\reals}$ by \cref{coord_scalar_sequence_function}.
    Hence $\sumFiniteSequence{b}{0 \rmul s}$ by \cref{sum_finite_sequence_scalar}.
    $a$ is a function and $\dom{a} = J$ by \cref{funs_is_function,funs}.
    $b$ is a function and $\dom{b} = J$ by \cref{funs_is_function,funs}.
    For all $i\in J$ we have $b(i) = a(i)$
        by \cref{funs_type_apply,coord_scalar_sequence,real_mul_zero,function_apply_intro,funs,real_lt_subst_left,real_lt_subst_right}.
    Show that $b\subseteq a$.
    \begin{subproof}
        Follows by \cref{fun_subseteq,subseteq}.
    \end{subproof}
    Show that $a\subseteq b$.
    \begin{subproof}
        Follows by \cref{fun_subseteq,subseteq,real_lt_subst_left}.
    \end{subproof}
    Thus $b = a$ by \cref{subseteq_antisymmetric}.
    Therefore $\sumFiniteSequence{a}{0}$ by \cref{real_lt_subst_left,sum_finite_sequence_value_in_reals,real_mul_zero,real_lt_subst_right}.
    Finally $s = 0$ by \cref{sum_finite_sequence_unique}.
\end{proof}

% from §41 helper lemma 18: removing the last element from a finite enumeration
\begin{lemma}\label{finite_enumeration_remove_last}
    Suppose $m\in\naturalsPlus$.
    Let $J = \initialSegment{m}$.
    Let $J_1 = \initialSegment{m + 1}$.
    Suppose $g$ is a bijection from $J_1$ to $F$.
    Suppose $g(m + 1) = \beta$.
    Let $h = \restrl{g}{J}$.
    Let $E = F\setminus \{\beta\}$.
    Then $h$ is a bijection from $J$ to $E$.
\end{lemma}
\begin{proof}
    $g$ is a function and $\dom{g} = J_1$ by \cref{bijection_is_function,bijections_dom}.
    We have $J\subseteq J_1$ by \cref{initial_segment_subset_succ}.
    Therefore $h$ is a function and $\dom{h} = J_1\inter J$ by \cref{restrl_of_function_is_function,dom_of_restrl_eq_inter}.
    Thus $\dom{h} = J$ by \cref{inter_absorb_supseteq_left}.
    For all $u\in\dom{h}$ we have $h(u)\in E$
        by \cref{subseteq,initial_segment_subset_succ,restrl_of_function_apply,bijections,surj,funs_type_apply,real_lt_subst_left,injective_function,real_naturals_closed_succ,initial_segment_naturalsplus,real_naturals_subset_reals,real_add_closed,real_one_in_naturals,real_naturals_lt_succ_local,real_lt_subst_right,real_lt_trans,real_lt_irreflexive,singleton_elim,setminus_intro}.
    $h\in\funs{J}{E}$ by \cref{funs_intro}.
    We have for all $y\in E$ there exists $u\in J_1$ such that $g(u) = y$ by \cref{setminus,bijections,surj}.
    Show that for all $y\in E$ there exists $u\in J$ such that $h(u) = y$.
    \begin{subproof}
        Follows by \cref{setminus,initial_segment_naturalsplus,real_naturals_leq_succ_elim,restrl_of_function_apply,subseteq,real_lt_subst_left,singleton_intro}.
    \end{subproof}
    Therefore $h\in\surj{J}{E}$ by \cref{surj}.
    Show that for all $u_1,u_2$ such that $u_1\in\dom{h}$ and $u_2\in\dom{h}$ and $h(u_1) = h(u_2)$
        we have $u_1 = u_2$.
    \begin{subproof}
        Follows by \cref{subseteq,initial_segment_subset_succ,restrl_of_function_apply,bijections,injective_function,pair_eq_iff}.
    \end{subproof}
    Thus $h$ is injective by \cref{funs_is_function,injective_function}.
    Finally $h$ is a bijection from $J$ to $E$ by \cref{bijections}.
\end{proof}

% from §41 helper lemma 19: positive values lie in the support
\begin{lemma}\label{positive_real_function_value_in_support}
    Assume $f$ is continuous from $X$ to $\reals$ with respect to $T$ and $\standardTopologyR$.
    Assume $x\in X$.
    Suppose $0 < f(x)$.
    Then $x\in \supp{f}{X}{T}$.
\end{lemma}
\begin{proof}
    Let $A = \supp{f}{X}{T}$.
    Let $B = X\setminus \preimg{f}{\{0\}}$.
    $B\subseteq A$ by \cref{continuous,setminus_subseteq,support,subseteq_closure}.
    Therefore $x\in A$
        by \cref{continuous,funs_is_function,funs,preimg_iff,singleton_elim,function_apply_intro,real_lt_subst_right,real_add_ident,real_lt_irreflexive,setminus_intro,subseteq}.
\end{proof}

% from §41 helper lemma 20: removing a designated value from a finite enumeration
\begin{lemma}\label{finite_enumeration_remove_value}
    Suppose $n\in\naturalsPlus$.
    Let $J = \initialSegment{n}$.
    Let $J_1 = \initialSegment{n + 1}$.
    Suppose $g$ is a bijection from $J_1$ to $F$.
    Suppose $k\in J_1$.
    Suppose $g(k) = \beta$.
    Then there exists $h$ such that
        $h$ is a bijection from $J$ to $F\setminus \{\beta\}$ and
        for all $i\in J$ we have
            $(i < k \land h(i) = g(i)) \lor (k \leq i \land h(i) = g(i + 1))$.
\end{lemma}
\begin{proof}
    $g$ is a function and $\dom{g} = J_1$ by \cref{bijection_is_function,bijections_dom}.
    Let $E = F\setminus \{\beta\}$.
    Let $R = \{(i,y) \mid i\in J, y\in E \mid
        (i < k \land y = g(i)) \lor (k \leq i \land y = g(i + 1))\}$.
    $R$ is a relation by \cref{relation,pair_eq_iff}.
    Show that for all $i\in J$ there exists $y$ such that $i\mathrel{R} y$.
    \begin{subproof}
        Fix $i\in J$.
        $i\in\naturalsPlus$ and $i \leq n$ by \cref{initial_segment_naturalsplus}.
        $k\in\naturalsPlus$ and $k \leq n + 1$ by \cref{initial_segment_naturalsplus}.
        \begin{byCase}
        \caseOf{$i < k$.}
            $i\in\dom{g}$ by \cref{bijections_dom,initial_segment_subset_succ,subseteq}.
            $g(i)\in F$ by \cref{bijections,surj,funs_type_apply,subseteq}.
            Show that $g(i)\neq \beta$.
            \begin{subproof}
                Suppose not.
                $g(i) = g(k)$ by \cref{real_lt_subst_right}.
                $i = k$ by \cref{bijections,injective_function,subseteq}.
                Therefore $k < k$ by \cref{real_lt_subst_left,real_lt_subst_right}.
                $k\in\reals$ by \cref{real_naturals_subset_reals,subseteq}.
                Contradiction by \cref{real_lt_irreflexive}.
            \end{subproof}
            $g(i)\notin\{\beta\}$ by \cref{singleton_iff}.
            Thus $g(i)\in E$ by \cref{setminus_intro}.
            Therefore there exists $y$ such that $i\mathrel{R} y$ by \cref{pair_eq_iff}.
        \caseOf{$i = k$.}
            $k \leq i$ by \cref{real_lt_subst_right}.
            $i + 1\in\naturalsPlus$ by \cref{real_naturals_closed_succ}.
            $i\in\reals$ and $n\in\reals$ by \cref{real_naturals_subset_reals,subseteq}.
            $1\in\reals$ by \cref{real_one_in_naturals,real_naturals_subset_reals,subseteq}.
            Therefore $i + 1 \leq n + 1$ by \cref{real_leq_add}.
            $i + 1\in J_1$ by \cref{initial_segment_naturalsplus}.
            $i + 1\in\dom{g}$ by \cref{bijections_dom,subseteq}.
            $g(i + 1)\in F$ by \cref{bijections,surj,funs_type_apply,subseteq}.
            Show that $g(i + 1)\neq \beta$.
            \begin{subproof}
                Suppose not.
                $g(i + 1) = g(k)$ by \cref{real_lt_subst_right}.
                $i + 1 = k$ by \cref{bijections,injective_function,subseteq}.
                Therefore $i < k$ by \cref{real_naturals_lt_succ_local,real_lt_subst_right}.
                Contradiction by \cref{real_lt_irreflexive,real_lt_subst_right}.
            \end{subproof}
            $g(i + 1)\notin\{\beta\}$ by \cref{singleton_iff}.
            Thus $g(i + 1)\in E$ by \cref{setminus_intro}.
            Therefore there exists $y$ such that $i\mathrel{R} y$ by \cref{pair_eq_iff}.
        \caseOf{$k < i$.}
            $k \leq i$ by \cref{real_lt_imp_leq}.
            $i + 1\in\naturalsPlus$ by \cref{real_naturals_closed_succ}.
            $i\in\reals$ and $n\in\reals$ by \cref{real_naturals_subset_reals,subseteq}.
            $1\in\reals$ by \cref{real_one_in_naturals,real_naturals_subset_reals,subseteq}.
            Therefore $i + 1 \leq n + 1$ by \cref{real_leq_add}.
            $i + 1\in J_1$ by \cref{initial_segment_naturalsplus}.
            $i + 1\in\dom{g}$ by \cref{bijections_dom,subseteq}.
            $g(i + 1)\in F$ by \cref{bijections,surj,funs_type_apply,subseteq}.
            Show that $g(i + 1)\neq \beta$.
            \begin{subproof}
                Suppose not.
                $g(i + 1) = g(k)$ by \cref{real_lt_subst_right}.
                $i + 1 = k$ by \cref{bijections,injective_function,subseteq}.
                Therefore $i < k$ by \cref{real_naturals_lt_succ_local,real_lt_subst_right}.
                $k\in\reals$ and $i\in\reals$ by \cref{real_naturals_subset_reals,subseteq}.
                Therefore $k < k$ by \cref{real_lt_trans}.
                Contradiction by \cref{real_lt_irreflexive}.
            \end{subproof}
            $g(i + 1)\notin\{\beta\}$ by \cref{singleton_iff}.
            Thus $g(i + 1)\in E$ by \cref{setminus_intro}.
            Therefore there exists $y$ such that $i\mathrel{R} y$ by \cref{pair_eq_iff}.
        \end{byCase}
    \end{subproof}
    Take $h$ such that $h$ is a function on $J$ and for all $i\in J$ we have $i\mathrel{R} h(i)$ by \cref{choice_function}.
    Show that for all $i\in J$ we have
        $(i < k \land h(i) = g(i)) \lor (k \leq i \land h(i) = g(i + 1))$.
    \begin{subproof}
        Follows by \cref{choice_function,pair_eq_iff}.
    \end{subproof}
    Show that for all $i$ such that $i\in J$ and $i < k$ we have $h(i) = g(i)$.
    \begin{subproof}
        Fix $i$ such that $i\in J$ and $i < k$.
        We have $(i < k \land h(i) = g(i)) \lor (k \leq i \land h(i) = g(i + 1))$
            by \cref{choice_function,pair_eq_iff}.
        Show that $h(i) = g(i)$.
        \begin{subproof}
            Suppose not.
            \begin{byCase}
            \caseOf{$i < k \land h(i) = g(i)$.}
                Contradiction.
            \caseOf{$k \leq i \land h(i) = g(i + 1)$.}
                We have $k < i$ or $k = i$ by \cref{real_leq_split}.
                \begin{byCase}
                \caseOf{$k < i$.}
                    $i\in\reals$ and $k\in\reals$ by \cref{real_naturals_subset_reals,subseteq,initial_segment_naturalsplus}.
                    $i < i$ by \cref{real_lt_trans}.
                    Contradiction by \cref{real_lt_irreflexive}.
                \caseOf{$k = i$.}
                    $i\in\reals$ by \cref{real_naturals_subset_reals,subseteq,initial_segment_naturalsplus}.
                    $i < i$ by \cref{real_lt_subst_right}.
                    Contradiction by \cref{real_lt_irreflexive}.
                \end{byCase}
            \end{byCase}
        \end{subproof}
    \end{subproof}
    Show that for all $i$ such that $i\in J$ and $k \leq i$ we have $h(i) = g(i + 1)$.
    \begin{subproof}
        Fix $i$ such that $i\in J$ and $k \leq i$.
        We have $(i < k \land h(i) = g(i)) \lor (k \leq i \land h(i) = g(i + 1))$
            by \cref{choice_function,pair_eq_iff}.
        Show that $h(i) = g(i + 1)$.
        \begin{subproof}
            Suppose not.
            \begin{byCase}
            \caseOf{$i < k \land h(i) = g(i)$.}
                We have $k < i$ or $k = i$ by \cref{real_leq_split}.
                \begin{byCase}
                \caseOf{$k < i$.}
                    $i\in\reals$ and $k\in\reals$ by \cref{real_naturals_subset_reals,subseteq,initial_segment_naturalsplus}.
                    $i < i$ by \cref{real_lt_trans}.
                    Contradiction by \cref{real_lt_irreflexive}.
                \caseOf{$k = i$.}
                    $i\in\reals$ by \cref{real_naturals_subset_reals,subseteq,initial_segment_naturalsplus}.
                    $i < i$ by \cref{real_lt_subst_right}.
                    Contradiction by \cref{real_lt_irreflexive}.
                \end{byCase}
            \caseOf{$k \leq i \land h(i) = g(i + 1)$.}
                Contradiction.
            \end{byCase}
        \end{subproof}
    \end{subproof}
    Thus for all $i\in J$ we have $h(i)\in E$ by \cref{choice_function,pair_eq_iff}.
    $h\in\funs{J}{E}$ by \cref{funs_intro}.
    Show that for all $y\in E$ there exists $i$ such that $i\in J$ and $h(i) = y$.
    \begin{subproof}
        Fix $y\in E$.
        $y\in F$ and $y\notin\{\beta\}$ by \cref{setminus_elim_left,setminus_elim_right}.
        Take $j\in J_1$ such that $g(j) = y$ by \cref{bijections,surj}.
        $j\in\naturalsPlus$ and $j \leq n + 1$ by \cref{initial_segment_naturalsplus}.
        $j\in\reals$ by \cref{real_naturals_subset_reals,subseteq}.
        $k\in\naturalsPlus$ by \cref{initial_segment_naturalsplus}.
        $k \leq n + 1$ by \cref{initial_segment_naturalsplus}.
        $k\in\reals$ by \cref{real_naturals_subset_reals,subseteq}.
        $n\in\reals$ by \cref{real_naturals_subset_reals,subseteq}.
        $1\in\reals$ by \cref{real_one_in_naturals,real_naturals_subset_reals,subseteq}.
        $n + 1\in\reals$ by \cref{real_add_closed}.
        Show that $j \neq k$.
        \begin{subproof}
            Suppose not.
            $y = \beta$ by \cref{real_lt_subst_left,real_lt_subst_right}.
            $y\in\{\beta\}$ by \cref{singleton_intro}.
            Contradiction by \cref{setminus_elim_right}.
        \end{subproof}
        Hence $j < k$ or $k < j$ by \cref{real_lt_trichotomy}.
        \begin{byCase}
        \caseOf{$j < k$.}
            Show that $j \neq n + 1$.
            \begin{subproof}
                Suppose not.
                $n + 1 < k$ by \cref{real_lt_subst_left}.
                $k < n + 1$ or $k = n + 1$ by \cref{initial_segment_naturalsplus}.
                \begin{byCase}
                \caseOf{$k < n + 1$.}
                    Contradiction by \cref{real_lt_asym}.
                \caseOf{$k = n + 1$.}
                    $n + 1 < n + 1$ by \cref{real_lt_subst_right}.
                    Contradiction by \cref{real_lt_irreflexive}.
                \end{byCase}
            \end{subproof}
            Therefore $j \leq n$ by \cref{real_naturals_leq_succ_elim}.
            $j\in J$ by \cref{initial_segment_naturalsplus}.
            We have $(j < k \land h(j) = g(j)) \lor (k \leq j \land h(j) = g(j + 1))$
                by \cref{choice_function,pair_eq_iff}.
            Show that $h(j) = g(j)$.
            \begin{subproof}
                Suppose not.
                \begin{byCase}
                \caseOf{$j < k \land h(j) = g(j)$.}
                    Contradiction.
                \caseOf{$k \leq j \land h(j) = g(j + 1)$.}
                    We have $k < j$ or $k = j$ by \cref{real_leq_split}.
                    \begin{byCase}
                    \caseOf{$k < j$.}
                        Contradiction by \cref{real_lt_asym}.
                    \caseOf{$k = j$.}
                        $j < j$ by \cref{real_lt_subst_right}.
                        Contradiction by \cref{real_lt_irreflexive}.
                    \end{byCase}
                \end{byCase}
            \end{subproof}
            $h(j) = y$.
            Therefore there exists $i\in J$ such that $h(i) = y$ by \cref{real_lt_subst_right}.
        \caseOf{$k < j$.}
            Show that $j \neq 1$.
            \begin{subproof}
                Suppose not.
                $k < 1$ by \cref{real_lt_subst_right}.
                Contradiction by \cref{real_one_leq_naturalsplus,real_lt_asym}.
            \end{subproof}
            Take $u\in\naturalsPlus$ such that $j = u + 1$ by \cref{real_naturals_predecessor}.
            $u\in\reals$ by \cref{real_naturals_subset_reals,subseteq}.
            $u < u + 1$ by \cref{real_naturals_lt_succ_local}.
            $u < j$ by \cref{real_lt_subst_right}.
            $u \leq j$ by \cref{real_lt_imp_leq}.
            Therefore $u \leq n + 1$ by \cref{real_leq_trans}.
            Show that $u \neq n + 1$.
            \begin{subproof}
                Suppose not.
                $n + 1 < j$ by \cref{real_lt_subst_left}.
                $j < n + 1$ or $j = n + 1$ by \cref{initial_segment_naturalsplus}.
                \begin{byCase}
                \caseOf{$j < n + 1$.}
                    $n + 1 < n + 1$ by \cref{real_lt_trans}.
                    Contradiction by \cref{real_lt_irreflexive}.
                \caseOf{$j = n + 1$.}
                    $n + 1 < n + 1$ by \cref{real_lt_subst_right}.
                    Contradiction by \cref{real_lt_irreflexive}.
                \end{byCase}
            \end{subproof}
            Therefore $u \leq n$ by \cref{real_naturals_leq_succ_elim}.
            $u\in J$ by \cref{initial_segment_naturalsplus}.
            $k < u + 1$ by \cref{real_lt_subst_right}.
            Show that $k \neq u + 1$.
            \begin{subproof}
                Suppose not.
                $k < k$ by \cref{real_lt_subst_right}.
                Contradiction by \cref{real_lt_irreflexive}.
            \end{subproof}
            Therefore $k \leq u$ by \cref{real_naturals_leq_succ_elim}.
            We have $(u < k \land h(u) = g(u)) \lor (k \leq u \land h(u) = g(u + 1))$
                by \cref{choice_function,pair_eq_iff}.
            Show that $h(u) = g(u + 1)$.
            \begin{subproof}
                Suppose not.
                \begin{byCase}
                \caseOf{$u < k \land h(u) = g(u)$.}
                    $u < k$.
                    We have $k < u$ or $k = u$ by \cref{real_leq_split}.
                    \begin{byCase}
                    \caseOf{$k < u$.}
                        $u < u$ by \cref{real_lt_trans}.
                        Contradiction by \cref{real_lt_irreflexive}.
                    \caseOf{$k = u$.}
                        $u < u$ by \cref{real_lt_subst_right}.
                        Contradiction by \cref{real_lt_irreflexive}.
                    \end{byCase}
                \caseOf{$k \leq u \land h(u) = g(u + 1)$.}
                    Contradiction.
                \end{byCase}
            \end{subproof}
            $h(u) = y$.
            Therefore there exists $i\in J$ such that $h(i) = y$ by \cref{real_lt_subst_right}.
        \end{byCase}
    \end{subproof}
    Thus $h\in\surj{J}{E}$ by \cref{surj}.
    Show that for all $i_1,i_2$ such that $i_1\in J$ and $i_2\in J$ and $h(i_1) = h(i_2)$
        we have $i_1 = i_2$.
    \begin{subproof}
        Fix $i_1$.
        Fix $i_2$ such that $i_1\in J$ and $i_2\in J$ and $h(i_1) = h(i_2)$.
        $i_1\in\naturalsPlus$ and $i_2\in\naturalsPlus$ and $i_1 \leq n$ and $i_2 \leq n$
            by \cref{initial_segment_naturalsplus}.
        $k\in\naturalsPlus$ and $k \leq n + 1$ by \cref{initial_segment_naturalsplus}.
        $J\subseteq J_1$ by \cref{initial_segment_subset_succ}.
        $i_1\in\reals$ and $i_2\in\reals$ and $n\in\reals$ by \cref{real_naturals_subset_reals,subseteq}.
        $k\in\reals$ by \cref{real_naturals_subset_reals,subseteq,initial_segment_naturalsplus}.
        $1\in\reals$ by \cref{real_one_in_naturals,real_naturals_subset_reals,subseteq}.
        We have $i_1 < k$ or $k \leq i_1$ by \cref{real_lt_trichotomy,real_lt_imp_leq}.
        $i_2 < k$ or $k \leq i_2$ by \cref{real_lt_trichotomy,real_lt_imp_leq}.
        \begin{byCase}
        \caseOf{$i_1 < k$ and $i_2 < k$.}
            $h(i_1) = g(i_1)$ and $h(i_2) = g(i_2)$.
            $g(i_1) = g(i_2)$ by \cref{real_lt_subst_left,real_lt_subst_right}.
            $i_1\in\dom{g}$ and $i_2\in\dom{g}$ by \cref{bijections_dom,subseteq}.
            Therefore $i_1 = i_2$ by \cref{bijections,injective_function,subseteq}.
        \caseOf{$i_1 < k$ and $k \leq i_2$.}
            $h(i_1) = g(i_1)$ and $h(i_2) = g(i_2 + 1)$.
            $g(i_1) = g(i_2 + 1)$ by \cref{real_lt_subst_left,real_lt_subst_right}.
            $i_1\in\dom{g}$ by \cref{bijections_dom,subseteq}.
            $i_2 + 1\in\naturalsPlus$ by \cref{real_naturals_closed_succ}.
            Therefore $i_2 + 1 \leq n + 1$ by \cref{real_leq_add}.
            $i_2 + 1\in J_1$ by \cref{initial_segment_naturalsplus}.
            $i_2 + 1\in\dom{g}$ by \cref{bijections_dom,subseteq}.
            Thus $i_1 = i_2 + 1$ by \cref{bijections,injective_function,subseteq,real_naturals_closed_succ,initial_segment_naturalsplus}.
            $i_2 < i_1$ by \cref{real_naturals_lt_succ_local,real_lt_subst_right}.
            Show that $k < i_1$.
            \begin{subproof}
                We have $k < i_2$ or $k = i_2$ by \cref{real_leq_split}.
                \begin{byCase}
                \caseOf{$k < i_2$.}
                    $k < i_1$ by \cref{real_lt_trans}.
                \caseOf{$k = i_2$.}
                    $k < i_1$ by \cref{real_lt_subst_left}.
                \end{byCase}
            \end{subproof}
            Show that $i_1 = i_2$.
            \begin{subproof}
                Suppose not.
                $i_1 < i_1$ by \cref{real_lt_trans}.
                Contradiction by \cref{real_lt_irreflexive}.
            \end{subproof}
        \caseOf{$k \leq i_1$ and $i_2 < k$.}
            $h(i_1) = g(i_1 + 1)$ and $h(i_2) = g(i_2)$.
            $g(i_1 + 1) = g(i_2)$ by \cref{real_lt_subst_left,real_lt_subst_right}.
            $i_1 + 1\in\naturalsPlus$ by \cref{real_naturals_closed_succ}.
            Therefore $i_1 + 1 \leq n + 1$ by \cref{real_leq_add}.
            $i_1 + 1\in J_1$ by \cref{initial_segment_naturalsplus}.
            $i_1 + 1\in\dom{g}$ by \cref{bijections_dom,subseteq}.
            $i_2\in\dom{g}$ by \cref{bijections_dom,subseteq}.
            Thus $i_1 + 1 = i_2$ by \cref{bijections,injective_function,subseteq,real_naturals_closed_succ,initial_segment_naturalsplus}.
            $i_1 < i_2$ by \cref{real_naturals_lt_succ_local,real_lt_subst_left}.
            Show that $k < i_2$.
            \begin{subproof}
                We have $k < i_1$ or $k = i_1$ by \cref{real_leq_split}.
                \begin{byCase}
                \caseOf{$k < i_1$.}
                    $k < i_2$ by \cref{real_lt_trans}.
                \caseOf{$k = i_1$.}
                    $k < i_2$ by \cref{real_lt_subst_left}.
                \end{byCase}
            \end{subproof}
            Show that $i_1 = i_2$.
            \begin{subproof}
                Suppose not.
                $i_2 < i_2$ by \cref{real_lt_trans}.
                Contradiction by \cref{real_lt_irreflexive}.
            \end{subproof}
        \caseOf{$k \leq i_1$ and $k \leq i_2$.}
            $h(i_1) = g(i_1 + 1)$ and $h(i_2) = g(i_2 + 1)$.
            $g(i_1 + 1) = g(i_2 + 1)$ by \cref{real_lt_subst_left,real_lt_subst_right}.
            $i_1 + 1\in\naturalsPlus$ and $i_2 + 1\in\naturalsPlus$ by \cref{real_naturals_closed_succ}.
            Therefore $i_1 + 1 \leq n + 1$ and $i_2 + 1 \leq n + 1$ by \cref{real_leq_add}.
            $i_1 + 1\in J_1$ and $i_2 + 1\in J_1$ by \cref{initial_segment_naturalsplus}.
            $i_1 + 1\in\dom{g}$ and $i_2 + 1\in\dom{g}$ by \cref{bijections_dom,subseteq}.
            Thus $i_1 + 1 = i_2 + 1$ by \cref{bijections,injective_function,subseteq,real_naturals_closed_succ,initial_segment_naturalsplus}.
            $i_1\in\reals$ by \cref{real_naturals_subset_reals,subseteq,initial_segment_naturalsplus}.
            $i_2\in\reals$ by \cref{real_naturals_subset_reals,subseteq,initial_segment_naturalsplus}.
            $1\in\reals$ by \cref{real_one_in_naturals,real_naturals_subset_reals,subseteq}.
            Therefore $i_1 = i_2$ by \cref{real_add_cancel}.
        \end{byCase}
    \end{subproof}
    Thus $h$ is injective by \cref{funs_is_function,injective_function}.
    Finally $h$ is a bijection from $J$ to $E$ by \cref{bijections}.
\end{proof}

% from §41 helper definition 1: increasing on an initial segment with respect to a relation
\begin{definition}\label{increasing_on_initial_segment_with_respect_to_relation}
    $g$ is increasing on $A$ with respect to $R$ iff
        for all $i,j$ such that $i\in A$ and $j\in A$ and $i < j$ we have
            $g(i)\mathrel{R} g(j)$.
\end{definition}

% from §41 helper definition 2: finite cardinality admits a well-ordered enumeration
\begin{definition}\label{well_order_enumerable_cardinality}
    $n$ is well-order-enumerable in $J$ with respect to $R$ iff
        for all $E,h$ such that $E\subseteq J$ and $E$ is inhabited and $h$ is a bijection from $\initialSegment{n}$ to $E$
            there exists $q$ such that $q$ is a bijection from $\initialSegment{n}$ to $E$ and $q$ is increasing on $\initialSegment{n}$ with respect to $R$.
\end{definition}

% from §41 helper lemma 21: finite subsets admit well-ordered enumerations
\begin{lemma}\label{finite_well_ordered_enumeration}
    Assume $R$ is a well order relation on $J$.
    Suppose $F\subseteq J$.
    Suppose $F$ is finite.
    Suppose $F$ is inhabited.
    Then there exist $n,g$ such that
        $n\in\naturalsPlus$ and
        $g$ is a bijection from $\initialSegment{n}$ to $F$ and
        $g$ is increasing on $\initialSegment{n}$ with respect to $R$.
\end{lemma}
    \begin{proof}
        Let $A = \{n\in\naturalsPlus \mid \text{$n$ is well-order-enumerable in $J$ with respect to $R$}\}$.
        Show that $1\in A$.
        \begin{subproof}
            Show that for all $E,h$ such that $E\subseteq J$ and $E$ is inhabited and $h$ is a bijection from $\initialSegment{1}$ to $E$
                we have there exists $q$ such that $q$ is a bijection from $\initialSegment{1}$ to $E$ and $q$ is increasing on $\initialSegment{1}$ with respect to $R$.
            \begin{subproof}
                Fix $E$.
                Fix $h$ such that $E\subseteq J$ and $E$ is inhabited and $h$ is a bijection from $\initialSegment{1}$ to $E$.
                Take $\beta$ such that $\beta\in E$.
                There exists $u\in\initialSegment{1}$ such that $h(u) = \beta$ by \cref{bijections,surj}.
                $u = 1$ by \cref{initial_segment_naturalsplus,real_one_in_naturals,real_one_leq_naturalsplus,real_naturals_subset_reals,subseteq,real_lt_trans,real_lt_irreflexive}.
                Thus $h(1) = \beta$ by \cref{real_lt_subst_left}.
                For all $y\in E$ there exists $u\in\initialSegment{1}$ such that $h(u) = y$ by \cref{bijections,surj}.
                Therefore for all $y\in E$ we have $y = h(1)$ by
                    \cref{initial_segment_naturalsplus,real_one_in_naturals,real_one_leq_naturalsplus,real_naturals_subset_reals,subseteq,real_lt_trans,real_lt_irreflexive,real_lt_subst_left}.
                Thus $h$ is a bijection from $\initialSegment{1}$ to $E$ and
                    $h$ is increasing on $\initialSegment{1}$ with respect to $R$
                    by \cref{initial_segment_naturalsplus,real_one_in_naturals,real_one_leq_naturalsplus,real_naturals_subset_reals,subseteq,real_lt_trans,real_lt_irreflexive,increasing_on_initial_segment_with_respect_to_relation}.
                Therefore there exists $q$ such that $q$ is a bijection from $\initialSegment{1}$ to $E$ and $q$ is increasing on $\initialSegment{1}$ with respect to $R$
                    by \cref{real_lt_subst_left}.
            \end{subproof}
            Therefore $1$ is well-order-enumerable in $J$ with respect to $R$ by \cref{well_order_enumerable_cardinality}.
        \end{subproof}
        Thus $1\in A$ by \cref{real_one_in_naturals,real_one_leq_naturalsplus}.
        Show that for all $n\in A$ we have $n + 1\in A$.
        \begin{subproof}
            Fix $n$ such that $n\in A$.
            $n$ is well-order-enumerable in $J$ with respect to $R$ by definition.
            Show that for all $E,h$ such that $E\subseteq J$ and $E$ is inhabited and $h$ is a bijection from $\initialSegment{n + 1}$ to $E$
                we have there exists $g$ such that $g$ is a bijection from $\initialSegment{n + 1}$ to $E$ and $g$ is increasing on $\initialSegment{n + 1}$ with respect to $R$.
            \begin{subproof}
                Fix $E$.
                Fix $h$ such that $E\subseteq J$ and $E$ is inhabited and $h$ is a bijection from $\initialSegment{n + 1}$ to $E$.
                $R$ is a simple order relation on $J$ by \cref{well_order_relation}.
                We have $n\in\naturalsPlus$ by \cref{subseteq}.
                Hence $n + 1\in\naturalsPlus$ by \cref{real_naturals_closed_succ}.
                $\initialSegment{n + 1}$ is finite by \cref{initial_segment_finite}.
                Hence $E$ is finite by \cref{finite_bijection}.
                Take $\beta$ such that $\beta$ is a largest element of $E$ with respect to $R$
                    by \cref{finite_inhabited_largest_element,bijections,surj}.
                Hence $\beta\in E$ by \cref{largest_element}.
                Take $k$ such that $k\in\initialSegment{n + 1}$ and $h(k) = \beta$ by \cref{bijections,surj}.
                Take $p$ such that
                    $p$ is a bijection from $\initialSegment{n}$ to $E\setminus \{\beta\}$ and
                    for all $i\in\initialSegment{n}$ we have
                        $(i < k \land p(i) = h(i)) \lor (k \leq i \land p(i) = h(i + 1))$
                    by \cref{finite_enumeration_remove_value}.
                $1\in\initialSegment{n}$ by \cref{real_one_in_naturals,real_one_leq_naturalsplus,initial_segment_naturalsplus}.
                $p(1)\in E\setminus \{\beta\}$ by \cref{bijections,surj,funs_type_apply,subseteq}.
                Thus $E\setminus \{\beta\}$ is inhabited by \cref{nonempty_inhabited}.
                Therefore $E\setminus \{\beta\}\subseteq J$ by \cref{setminus_subseteq,subseteq_transitive}.
                Take $q$ such that
                    $q$ is a bijection from $\initialSegment{n}$ to $E\setminus \{\beta\}$ and
                    $q$ is increasing on $\initialSegment{n}$ with respect to $R$
                    by \cref{well_order_enumerable_cardinality}.
                Let $r = q\union \{(n + 1,\beta)\}$.
        Show that for all $u$ such that $u\in\dom{q}\inter \{n + 1\}$ we have $q(u) = \beta$.
        \begin{subproof}
            Fix $u$ such that $u\in\dom{q}\inter \{n + 1\}$.
            Show that $q(u) = \beta$.
            \begin{subproof}
                Suppose not.
                $u\in\dom{q}$ by \cref{inter_elim_left}.
                $u\in \{n + 1\}$ by \cref{inter_elim_right}.
                $u = n + 1$ by \cref{singleton_elim}.
                $u\in\initialSegment{n}$ by \cref{bijections_dom,subseteq}.
                $u\in\initialSegment{n + 1}$ by \cref{real_naturals_closed_succ,initial_segment_naturalsplus,real_lt_subst_left}.
                Therefore $n + 1 \leq n$ by \cref{initial_segment_naturalsplus,real_lt_subst_left}.
                $n\in\reals$ by \cref{real_naturals_subset_reals,subseteq}.
                $n + 1\in\reals$ by \cref{real_naturals_closed_succ,real_naturals_subset_reals,subseteq}.
                We have $n < n + 1$ by \cref{real_naturals_lt_succ_local}.
                Thus $n + 1 < n$ or $n + 1 = n$ by \cref{real_leq_split}.
                \begin{byCase}
                \caseOf{$n + 1 < n$.}
                    $n + 1 < n + 1$ by \cref{real_lt_trans}.
                    Contradiction by \cref{real_lt_irreflexive}.
                \caseOf{$n + 1 = n$.}
                    $n < n$ by \cref{real_lt_subst_right}.
                    Contradiction by \cref{real_lt_irreflexive}.
                \end{byCase}
            \end{subproof}
        \end{subproof}
        $q$ is right-unique and $q$ is a relation by \cref{bijection_is_function,function_rightunique,relation}.
        $ \{(n + 1,\beta)\}$ is right-unique and $\{(n + 1,\beta)\}$ is a relation by \cref{singleton_elim,pair_eq_iff,rightunique,relation}.
        Thus $\dom{\{(n + 1,\beta)\}} = \{n + 1\}$
            by \cref{dom_iff,singleton_elim,pair_eq_iff,singleton_intro,dom_intro,subseteq,subseteq_antisymmetric}.
        Show that for all $u$ such that $u\in\dom{q}\inter \dom{\{(n + 1,\beta)\}}$ we have $q(u) = \apply{\{(n + 1,\beta)\}}{u}$.
        \begin{subproof}
            Fix $u$ such that $u\in\dom{q}\inter \dom{\{(n + 1,\beta)\}}$.
            We have $u\in\dom{q}\inter \{n + 1\}$ by \cref{real_lt_subst_right}.
            $q(u) = \beta$ by \cref{subseteq}.
            $u\in \{n + 1\}$ by \cref{inter_elim_right}.
            $u = n + 1$ by \cref{singleton_elim}.
            Thus $\apply{\{(n + 1,\beta)\}}{u} = \beta$
                by \cref{singleton_intro,pair_eq_iff,function_apply_intro,real_lt_subst_left}.
            Therefore $q(u) = \apply{\{(n + 1,\beta)\}}{u}$ by \cref{real_lt_subst_right}.
        \end{subproof}
        Therefore $r$ is a function by \cref{union_compatible_functions,relation,rightunique}.
        We have $\dom{r} = \dom{q}\union \{n + 1\}$ by \cref{dom_union,dom_iff,singleton_elim,pair_eq_iff,singleton_intro,dom_intro,subseteq,subseteq_antisymmetric}.
        Thus $\dom{r} = \initialSegment{n}\union \{n + 1\}$ by \cref{bijections_dom}.
        Show that for all $u$ such that $u\in\initialSegment{n}\union\{n + 1\}$ we have $u\in\initialSegment{n + 1}$.
        \begin{subproof}
            Follows by \cref{union_iff,initial_segment_naturalsplus,real_naturals_subset_reals,subseteq,real_one_in_naturals,real_add_ident,real_add_closed,real_zero_lt_one,real_lt_add,real_add_comm,real_lt_subst_left,real_lt_trans,singleton_elim}.
        \end{subproof}
        Show that for all $u$ such that $u\in\initialSegment{n + 1}$ we have $u\in\initialSegment{n}\union\{n + 1\}$.
        \begin{subproof}
            Follows by \cref{initial_segment_naturalsplus,real_naturals_leq_succ_elim,union_iff,singleton_intro}.
        \end{subproof}
        Therefore $\initialSegment{n}\union\{n + 1\} = \initialSegment{n + 1}$ by \cref{subseteq,subseteq_antisymmetric}.
        $\dom{r} = \initialSegment{n + 1}$.
        Show that for all $u\in\dom{r}$ we have $r(u)\in E$.
        \begin{subproof}
            Fix $u$ such that $u\in\dom{r}$.
            We have $u\in\initialSegment{n}\union\{n + 1\}$ by \cref{real_lt_subst_right}.
            \begin{byCase}
            \caseOf{$u\in\initialSegment{n}$.}
                $u\in\dom{q}$ by \cref{bijections_dom,subseteq}.
                $r(u) = q(u)$ by \cref{function_apply_elim,union_iff,function_apply_intro}.
                $q(u)\in E\setminus \{\beta\}$ by \cref{bijections,surj,funs_type_apply,subseteq}.
                Therefore $r(u)\in E$ by \cref{setminus_elim_left,real_lt_subst_left}.
            \caseOf{$u\in\{n + 1\}$.}
                $u = n + 1$ by \cref{singleton_elim}.
                Thus $r(u) = \beta$ by \cref{singleton_intro,pair_eq_iff,union_iff,function_apply_intro}.
                Therefore $r(u)\in E$ by \cref{largest_element}.
            \end{byCase}
        \end{subproof}
        Thus $r\in\funs{\initialSegment{n + 1}}{E}$ by \cref{funs_intro}.
        Show that for all $y\in E$ there exists $u\in\initialSegment{n + 1}$ such that $r(u) = y$.
        \begin{subproof}
            Fix $y$ such that $y\in E$.
            \begin{byCase}
            \caseOf{$y = \beta$.}
                $n + 1\in\initialSegment{n + 1}$ by \cref{real_naturals_closed_succ,initial_segment_naturalsplus}.
                Hence $r(n + 1) = \beta$ by \cref{singleton_intro,pair_eq_iff,union_iff,function_apply_intro}.
                Therefore there exists $v\in\initialSegment{n + 1}$ such that $r(v) = y$ by \cref{real_lt_subst_left}.
            \caseOf{$y \neq \beta$.}
                $y\notin \{\beta\}$ by \cref{singleton_iff}.
                Thus $y\in E\setminus \{\beta\}$ by \cref{setminus_intro}.
                There exists $u\in\initialSegment{n}$ such that $q(u) = y$ by \cref{bijections,surj}.
                $u\in\initialSegment{n + 1}$ by \cref{initial_segment_subset_succ,subseteq}.
                $u\in\dom{q}$ by \cref{bijections_dom,subseteq}.
                $r(u) = q(u)$ by \cref{function_apply_elim,union_iff,function_apply_intro}.
                Therefore there exists $v\in\initialSegment{n + 1}$ such that $r(v) = y$ by \cref{real_lt_subst_right}.
            \end{byCase}
        \end{subproof}
        Therefore $r\in\surj{\initialSegment{n + 1}}{E}$ by \cref{surj}.
        Show that for all $u_1,u_2$ such that $u_1\in\dom{r}$ and $u_2\in\dom{r}$ and $r(u_1) = r(u_2)$ we have $u_1 = u_2$.
        \begin{subproof}
            Fix $u_1$.
            Fix $u_2$ such that $u_1\in\dom{r}$ and $u_2\in\dom{r}$ and $r(u_1) = r(u_2)$.
            \begin{byCase}
            \caseOf{$u_1 = n + 1$.}
                Show that $u_2 = n + 1$.
                \begin{subproof}
                    Suppose not.
                    $u_2\in\initialSegment{n}$ by \cref{initial_segment_naturalsplus,real_naturals_leq_succ_elim,real_lt_subst_left,bijections_dom,subseteq}.
                    $u_2\in\dom{q}$ by \cref{bijections_dom,subseteq}.
                    $r(u_2) = q(u_2)$ by \cref{function_apply_elim,union_iff,function_apply_intro}.
                    $q(u_2)\in E\setminus \{\beta\}$ by \cref{bijections,surj,funs_type_apply,subseteq}.
                    Thus $r(u_2)\neq \beta$ by \cref{setminus_elim_right,singleton_iff,real_lt_subst_left}.
                    $r(u_1) = \beta$ by \cref{singleton_intro,pair_eq_iff,union_iff,function_apply_intro,real_lt_subst_left}.
                    Therefore $r(u_1)\neq r(u_2)$ by \cref{real_lt_subst_left}.
                    Contradiction.
                \end{subproof}
                Thus $u_1 = u_2$ by \cref{real_lt_subst_left}.
            \caseOf{$u_1 \neq n + 1$.}
                $u_1\in\initialSegment{n}$ by \cref{initial_segment_naturalsplus,real_naturals_leq_succ_elim,real_lt_subst_left,bijections_dom,subseteq}.
                $u_1\in\dom{q}$ by \cref{bijections_dom,subseteq}.
                $r(u_1) = q(u_1)$ by \cref{function_apply_elim,union_iff,function_apply_intro}.
                Show that $u_2 \neq n + 1$.
                \begin{subproof}
                    Suppose not.
                    $r(u_2) = \beta$ by \cref{singleton_intro,pair_eq_iff,union_iff,function_apply_intro}.
                    $q(u_1) = \beta$ by \cref{real_lt_subst_right}.
                    Therefore $\beta\in E\setminus \{\beta\}$ by \cref{bijections,surj,funs_type_apply,subseteq,real_lt_subst_left}.
                    Contradiction by \cref{setminus_elim_right,singleton_intro}.
                \end{subproof}
                $u_2\in\initialSegment{n}$ by \cref{initial_segment_naturalsplus,real_naturals_leq_succ_elim,real_lt_subst_left,bijections_dom,subseteq}.
                $u_2\in\dom{q}$ by \cref{bijections_dom,subseteq}.
                $r(u_2) = q(u_2)$ by \cref{function_apply_elim,union_iff,function_apply_intro}.
                Thus $q(u_1) = q(u_2)$ by \cref{real_lt_subst_left,real_lt_subst_right}.
                Therefore $u_1 = u_2$ by \cref{bijections,injective_function,subseteq}.
            \end{byCase}
        \end{subproof}
        Thus $r$ is injective by \cref{function_rightunique,relation,injective_function}.
        Show that for all $i,j$ such that $i\in\initialSegment{n + 1}$ and $j\in\initialSegment{n + 1}$ and $i < j$ we have $r(i)\mathrel{R} r(j)$.
        \begin{subproof}
            Fix $i$.
            Fix $j$ such that $i\in\initialSegment{n + 1}$ and $j\in\initialSegment{n + 1}$ and $i < j$.
            \begin{byCase}
            \caseOf{$j = n + 1$.}
                $i\in\naturalsPlus$ and $i \leq n + 1$ by \cref{initial_segment_naturalsplus}.
                Show that $i \neq n + 1$.
                \begin{subproof}
                    Suppose not.
                    $i\in\reals$ by \cref{real_naturals_subset_reals,subseteq}.
                    $i < i$ by \cref{real_lt_subst_right,real_lt_subst_left}.
                    Contradiction by \cref{real_lt_irreflexive}.
                \end{subproof}
                Therefore $i \leq n$ by \cref{real_naturals_leq_succ_elim}.
                $i\in\initialSegment{n}$ by \cref{initial_segment_naturalsplus}.
                $i\in\dom{q}$ by \cref{bijections_dom,subseteq}.
                $r(i) = q(i)$ by \cref{function_apply_elim,union_iff,function_apply_intro}.
                $q(i)\in E\setminus \{\beta\}$ by \cref{bijections,surj,funs_type_apply,subseteq}.
                $q(i)\in E$ and $q(i)\neq \beta$ by \cref{setminus_elim_left,setminus_elim_right,singleton_iff}.
                We have $q(i)\mathrel{R}\beta$ or $q(i) = \beta$ by \cref{largest_element}.
                $q(i)\mathrel{R}\beta$ by \cref{real_lt_subst_left}.
                $r(i)\mathrel{R}\beta$ by \cref{real_lt_subst_left}.
                Therefore $r(i)\mathrel{R} r(j)$ by \cref{singleton_intro,pair_eq_iff,union_iff,function_apply_intro,real_lt_subst_left}.
            \caseOf{$j \neq n + 1$.}
                $j\in\naturalsPlus$ and $j \leq n + 1$ by \cref{initial_segment_naturalsplus}.
                $j \leq n$ by \cref{real_naturals_leq_succ_elim}.
                $j\in\initialSegment{n}$ by \cref{initial_segment_naturalsplus}.
                $i\in\naturalsPlus$ and $i \leq j$ by \cref{initial_segment_naturalsplus,real_lt_imp_leq}.
                $i\in\reals$ and $j\in\reals$ and $n\in\reals$ by \cref{real_naturals_subset_reals,subseteq}.
                Therefore $i \leq n$ by \cref{real_leq_trans}.
                $i\in\initialSegment{n}$ by \cref{initial_segment_naturalsplus}.
                $r(i) = q(i)$ and $r(j) = q(j)$ by \cref{bijections_dom,subseteq,function_apply_elim,union_iff,function_apply_intro}.
                Thus $q(i)\mathrel{R} q(j)$ by \cref{increasing_on_initial_segment_with_respect_to_relation}.
                Therefore $r(i)\mathrel{R} r(j)$ by \cref{real_lt_subst_left,real_lt_subst_right}.
            \end{byCase}
        \end{subproof}
        Therefore $r$ is increasing on $\initialSegment{n + 1}$ with respect to $R$ by \cref{increasing_on_initial_segment_with_respect_to_relation}.
        Thus $r$ is a bijection from $\initialSegment{n + 1}$ to $E$ by \cref{bijections}.
        Let $g_2 = r$.
        Thus there exists $h_0$ such that
            $h_0$ is a bijection from $\initialSegment{n + 1}$ to $E$ and
            $h_0$ is increasing on $\initialSegment{n + 1}$ with respect to $R$
            by \cref{real_lt_subst_left}.
            \end{subproof}
            Therefore $n + 1$ is well-order-enumerable in $J$ with respect to $R$ by \cref{well_order_enumerable_cardinality}.
        \end{subproof}
        Therefore for all $n\in A$ we have $n + 1\in A$ by \cref{real_naturals_closed_succ}.
    Take $m,h$ such that $m\in\naturalsPlus$ and $h$ is a bijection from $\initialSegment{m}$ to $F$
        by \cref{finite_inhabited_initial_segment_enumeration}.
    Therefore $m\in A$ by \cref{subseteq,real_naturals_induction}.
    Thus $m$ is well-order-enumerable in $J$ with respect to $R$ by definition.
    Take $g_0$ such that
        $g_0$ is a bijection from $\initialSegment{m}$ to $F$ and
        $g_0$ is increasing on $\initialSegment{m}$ with respect to $R$
        by \cref{well_order_enumerable_cardinality}.
    Therefore there exist $n_0,g_1$ such that
        $n_0 = m$ and
        $g_1 = g_0$ and
        $g_0$ is a bijection from $\initialSegment{m}$ to $F$ and
        $g_0$ is increasing on $\initialSegment{m}$ with respect to $R$
        by \cref{real_lt_subst_left}.
    Therefore there exist $p,s$ such that
        $p\in\naturalsPlus$ and
        $s$ is a bijection from $\initialSegment{p}$ to $F$ and
        $s$ is increasing on $\initialSegment{p}$ with respect to $R$
        by \cref{real_lt_subst_left}.
\end{proof}

% from §41 helper lemma 22: ordered finite enumerations end at the largest element
\begin{lemma}\label{increasing_enumeration_last_largest}
    Assume $R$ is a simple order relation on $J$.
    Suppose $F\subseteq J$.
    Suppose $n\in\naturalsPlus$.
    Suppose $g$ is a bijection from $\initialSegment{n}$ to $F$.
    Suppose $g$ is increasing on $\initialSegment{n}$ with respect to $R$.
    Then $g(n)$ is a largest element of $F$ with respect to $R$.
\end{lemma}
\begin{proof}
    $n\in\initialSegment{n}$ by \cref{initial_segment_naturalsplus}.
    $n\in\dom{g}$ by \cref{bijections_dom,subseteq}.
    Thus $g(n)\in F$ by \cref{bijections,surj,funs_type_apply,subseteq}.
    Show that for all $y\in F$ we have $y\mathrel{R} g(n)$ or $y = g(n)$.
    \begin{subproof}
        Fix $y\in F$.
        Take $i\in\initialSegment{n}$ such that $g(i) = y$ by \cref{bijections,surj}.
        $i\in\naturalsPlus$ and $i \leq n$ by \cref{initial_segment_naturalsplus}.
        $i\in\reals$ and $n\in\reals$ by \cref{real_naturals_subset_reals,subseteq}.
        We have $i < n$ or $i = n$ by \cref{real_leq_split}.
        \begin{byCase}
        \caseOf{$i < n$.}
            Thus $g(i)\mathrel{R} g(n)$ by \cref{increasing_on_initial_segment_with_respect_to_relation}.
            Therefore $y\mathrel{R} g(n)$ by \cref{real_lt_subst_left}.
        \caseOf{$i = n$.}
            Thus $y = g(n)$ by \cref{real_lt_subst_left}.
        \end{byCase}
    \end{subproof}
    Therefore $g(n)$ is a largest element of $F$ with respect to $R$ by \cref{largest_element}.
\end{proof}

% from §41 helper lemma 23: support index sets are finite and inhabited locally
\begin{lemma}\label{local_support_index_set_finite_inhabited}
    Assume $f$ is a function on $J$.
    Assume $S$ is the support family of $f$ on $X$ with respect to $T$ and $J$.
    Suppose $S$ is a locally finite indexed family on $J$ in $X$ with respect to $T$.
    Suppose $x\in X$.
    Suppose there exists $\beta\in J$ such that $x\in\apply{S}{\beta}$.
    Let $E = \{\alpha\in J \mid x\in\apply{S}{\alpha}\}$.
    Then $E\subseteq J$ and $E$ is finite and $E$ is inhabited.
\end{lemma}
\begin{proof}
    Take $\beta\in J$ such that $x\in\apply{S}{\beta}$.
    Take $W,F$ such that
        $W$ is a neighborhood of $x$ in $X$ with respect to $T$ and
        $F\subseteq J$ and
        $F$ is finite and
        for all $\alpha\in J$ if $W$ intersects $\apply{S}{\alpha}$ then $\alpha\in F$
        by \cref{locally_finite_indexed_family}.
    Show that for all $\alpha\in E$ we have $\alpha\in F$.
    \begin{subproof}
        Fix $\alpha\in E$.
        $\alpha\in J$ and $x\in\apply{S}{\alpha}$ by definition.
        $x\in W$ by \cref{neighborhood}.
        Thus $W$ intersects $\apply{S}{\alpha}$ by \cref{inter_intro,intersects,emptyset}.
        Therefore $\alpha\in F$ by \cref{subseteq}.
    \end{subproof}
    Thus $E\subseteq F$ by \cref{subseteq}.
    Show that $E\subseteq J$.
    \begin{subproof}
        Trivial.
    \end{subproof}
    Show that $E$ is finite.
    \begin{subproof}
        Follows by \cref{subseteq_finite_is_finite}.
    \end{subproof}
    We have $\beta\in E$ by definition.
    Hence $E\neq\emptyset$ by \cref{emptyset}.
    Show that $E$ is inhabited.
    \begin{subproof}
        Follows by \cref{nonempty_inhabited}.
    \end{subproof}
    Therefore $E\subseteq J$ and $E$ is finite and $E$ is inhabited.
\end{proof}

% from §41 helper lemma 24: sums over singleton enumerations equal the unique value
\begin{lemma}\label{singleton_enumeration_sum_value}
    Suppose $n\in\naturalsPlus$.
    Suppose $g$ is a bijection from $\initialSegment{n}$ to $\{\beta\}$.
    Suppose $a\in\funs{\initialSegment{n}}{\reals}$.
    Suppose $\sumFiniteSequence{a}{s}$.
    Suppose for all $i\in\initialSegment{n}$ we have $a(i) = r$.
    Then $s = r$.
\end{lemma}
\begin{proof}
    $1\in\initialSegment{n}$ by \cref{real_one_in_naturals,real_one_leq_naturalsplus,initial_segment_naturalsplus}.
    Show that for all $i\in\initialSegment{n}$ we have $i = 1$.
    \begin{subproof}
        Fix $i\in\initialSegment{n}$.
        $i\in\dom{g}$ and $1\in\dom{g}$ by \cref{bijections_dom,subseteq}.
        $g(i)\in \{\beta\}$ and $g(1)\in \{\beta\}$ by \cref{bijections,surj,funs_type_apply,subseteq}.
        $g(i) = \beta$ and $g(1) = \beta$ by \cref{singleton_elim}.
        Thus $g(i) = g(1)$ by \cref{real_lt_subst_left}.
        Therefore $i = 1$ by \cref{bijections,bijection_is_function,function_rightunique,relation,injective_function,subseteq}.
    \end{subproof}
    Show that $\initialSegment{n}\subseteq \initialSegment{1}$.
    \begin{subproof}
        Follows by \cref{real_one_in_naturals,real_one_leq_naturalsplus,initial_segment_naturalsplus,subseteq}.
    \end{subproof}
    Show that $\initialSegment{1}\subseteq \initialSegment{n}$.
    \begin{subproof}
        Follows by \cref{real_one_in_naturals,real_one_leq_naturalsplus,initial_segment_naturalsplus,subseteq}.
    \end{subproof}
    Thus $\initialSegment{n} = \initialSegment{1}$ by \cref{subseteq_antisymmetric}.
    $a\in\funs{\initialSegment{1}}{\reals}$ by \cref{funs,real_lt_subst_left}.
    Hence $\sumFiniteSequence{a}{a(1)}$ by \cref{sum_finite_sequence_one_term}.
    Therefore $s = a(1)$ by \cref{sum_finite_sequence_unique}.
    Thus $a(1) = r$ by \cref{real_one_in_naturals,real_one_leq_naturalsplus,initial_segment_naturalsplus}.
    Finally $s = r$ by \cref{real_lt_subst_right}.
\end{proof}

% from §41 helper definition 11: well-ordered enumeration of a finite subset
\begin{definition}\label{well_ordered_enumeration_of_subset}
    $g$ is a well-ordered enumeration of $F$ in $J$ with respect to $R$ and $n$ iff
        $F\subseteq J$ and
        $n\in\naturalsPlus$ and
        $g$ is a bijection from $\initialSegment{n}$ to $F$ and
        $g$ is increasing on $\initialSegment{n}$ with respect to $R$.
\end{definition}

% from §41 helper lemma 25: removing the last entry preserves well-ordered enumeration
\begin{lemma}\label{well_ordered_enumeration_remove_last}
    Assume $R$ is a simple order relation on $J$.
    Suppose $n\in\naturalsPlus$.
    Suppose $h$ is a well-ordered enumeration of $F$ in $J$ with respect to $R$ and $n + 1$.
    Let $\beta = h(n + 1)$.
    Let $G = F\setminus \{\beta\}$.
    Let $h_1 = \restrl{h}{\initialSegment{n}}$.
    Then $h_1$ is a well-ordered enumeration of $G$ in $J$ with respect to $R$ and $n$.
\end{lemma}
\begin{proof}
    $F\subseteq J$ and $n + 1\in\naturalsPlus$ and $h$ is a bijection from $\initialSegment{n + 1}$ to $F$ and
        $h$ is increasing on $\initialSegment{n + 1}$ with respect to $R$
        by \cref{well_ordered_enumeration_of_subset}.
    Show that $h_1$ is a bijection from $\initialSegment{n}$ to $G$.
    \begin{subproof}
        Follows by \cref{finite_enumeration_remove_last}.
    \end{subproof}
    Show that $h_1$ is increasing on $\initialSegment{n}$ with respect to $R$.
    \begin{subproof}
        Show that for all $i,j$ such that $i\in\initialSegment{n}$ and $j\in\initialSegment{n}$ and $i < j$ we have $h_1(i)\mathrel{R} h_1(j)$.
        \begin{subproof}
            Fix $i$.
            Fix $j$ such that $i\in\initialSegment{n}$ and $j\in\initialSegment{n}$ and $i < j$.
            $i\in\dom{h_1}$ and $j\in\dom{h_1}$ by \cref{bijections_dom,subseteq}.
            $\initialSegment{n}\subseteq \dom{h}$ by \cref{well_ordered_enumeration_of_subset,bijections_dom,initial_segment_subset_succ,subseteq,subseteq_transitive}.
            $h_1(i) = h(i)$ and $h_1(j) = h(j)$ by \cref{restrl_of_function_apply,bijection_is_function}.
            $i\in\initialSegment{n + 1}$ and $j\in\initialSegment{n + 1}$ by \cref{initial_segment_subset_succ,subseteq}.
            Therefore $h(i)\mathrel{R} h(j)$ by \cref{increasing_on_initial_segment_with_respect_to_relation}.
            Thus $h_1(i)\mathrel{R} h_1(j)$ by \cref{real_lt_subst_left,real_lt_subst_right}.
        \end{subproof}
        Therefore $h_1$ is increasing on $\initialSegment{n}$ with respect to $R$
            by \cref{increasing_on_initial_segment_with_respect_to_relation}.
    \end{subproof}
    Show that $G\subseteq J$.
    \begin{subproof}
        Follows by \cref{setminus_subseteq,subseteq_transitive}.
    \end{subproof}
    Therefore $h_1$ is a well-ordered enumeration of $G$ in $J$ with respect to $R$ and $n$
        by \cref{well_ordered_enumeration_of_subset}.
\end{proof}

% from §41 helper lemma 26: zero-extension comparison for well-ordered local sums
\begin{lemma}\label{well_ordered_zero_extension_sum_comparison}
    Assume $R$ is a simple order relation on $J$.
    Suppose $g$ is a well-ordered enumeration of $F$ in $J$ with respect to $R$ and $n$.
    Suppose $h$ is a well-ordered enumeration of $G$ in $J$ with respect to $R$ and $m$.
    Suppose $a$ is a real sum sequence of $f$ at $x$ along $g$ with respect to $n$ and result $s$.
    Suppose $b$ is a real sum sequence of $f$ at $x$ along $h$ with respect to $m$ and result $t$.
    Suppose $G\subseteq F$.
    Suppose for all $\alpha\in F\setminus G$ we have $\apply{\apply{f}{\alpha}}{x} = 0$.
    Then $s = t$.
\end{lemma}
\begin{proof}
    Let $A = \{k\in\naturalsPlus \mid \text{for all $F_0,G_0,g_0,h_0,m_0,a_0,b_0,s_0,t_0$ such that
        $g_0$ is a well-ordered enumeration of $F_0$ in $J$ with respect to $R$ and $k$ and
        $h_0$ is a well-ordered enumeration of $G_0$ in $J$ with respect to $R$ and $m_0$ and
        $a_0$ is a real sum sequence of $f$ at $x$ along $g_0$ with respect to $k$ and result $s_0$ and
        $b_0$ is a real sum sequence of $f$ at $x$ along $h_0$ with respect to $m_0$ and result $t_0$ and
        $G_0\subseteq F_0$ and
        for all $\alpha\in F_0\setminus G_0$ we have $\apply{\apply{f}{\alpha}}{x} = 0$
        we have $s_0 = t_0$}\}$.
    Show that $1\in A$.
    \begin{subproof}
        Show that for all $F_0,G_0,g_0,h_0,m_0,a_0,b_0,s_0,t_0$ such that
            $g_0$ is a well-ordered enumeration of $F_0$ in $J$ with respect to $R$ and $1$ and
            $h_0$ is a well-ordered enumeration of $G_0$ in $J$ with respect to $R$ and $m_0$ and
            $a_0$ is a real sum sequence of $f$ at $x$ along $g_0$ with respect to $1$ and result $s_0$ and
            $b_0$ is a real sum sequence of $f$ at $x$ along $h_0$ with respect to $m_0$ and result $t_0$ and
            $G_0\subseteq F_0$ and
            for all $\alpha\in F_0\setminus G_0$ we have $\apply{\apply{f}{\alpha}}{x} = 0$
            we have $s_0 = t_0$.
        \begin{subproof}
            Fix $F_0$.
            Fix $G_0$.
            Fix $g_0$.
            Fix $h_0$.
            Fix $m_0$.
            Fix $a_0$.
            Fix $b_0$.
            Fix $s_0$.
            Fix $t_0$ such that
                $g_0$ is a well-ordered enumeration of $F_0$ in $J$ with respect to $R$ and $1$ and
                $h_0$ is a well-ordered enumeration of $G_0$ in $J$ with respect to $R$ and $m_0$ and
                $a_0$ is a real sum sequence of $f$ at $x$ along $g_0$ with respect to $1$ and result $s_0$ and
                $b_0$ is a real sum sequence of $f$ at $x$ along $h_0$ with respect to $m_0$ and result $t_0$ and
                $G_0\subseteq F_0$ and
                for all $\alpha\in F_0\setminus G_0$ we have $\apply{\apply{f}{\alpha}}{x} = 0$.
            $F_0\subseteq J$ and $g_0$ is a bijection from $\initialSegment{1}$ to $F_0$
                by \cref{well_ordered_enumeration_of_subset}.
            Let $\beta = g_0(1)$.
            $\beta\in F_0$
                by \cref{real_one_in_naturals,real_one_leq_naturalsplus,initial_segment_naturalsplus,bijections,surj,funs_type_apply,subseteq}.
            Show that for all $y\in F_0$ we have $y = \beta$.
            \begin{subproof}
                Fix $y\in F_0$.
                Take $u\in\initialSegment{1}$ such that $g_0(u) = y$ by \cref{bijections,surj}.
                $u = 1$ by \cref{initial_segment_naturalsplus,real_one_in_naturals,real_one_leq_naturalsplus,real_naturals_subset_reals,subseteq,real_lt_trans,real_lt_irreflexive}.
                $y = \beta$ by \cref{real_lt_subst_left}.
            \end{subproof}
            $F_0\subseteq \{\beta\}$ by \cref{subseteq,singleton_intro}.
            $\{\beta\}\subseteq F_0$ by \cref{singleton_subset_intro}.
            $F_0 = \{\beta\}$ by \cref{subseteq_antisymmetric}.
            $g_0$ is a bijection from $\initialSegment{1}$ to $\{\beta\}$
                by \cref{well_ordered_enumeration_of_subset,real_lt_subst_right}.
            $G_0\subseteq J$ and $m_0\in\naturalsPlus$ and $h_0$ is a bijection from $\initialSegment{m_0}$ to $G_0$
                by \cref{well_ordered_enumeration_of_subset}.
            $G_0$ is inhabited by \cref{real_one_in_naturals,real_one_leq_naturalsplus,initial_segment_naturalsplus,bijections,surj,funs_type_apply,subseteq,nonempty_inhabited}.
            $G_0 = \{\beta\}$ by \cref{subseteq,singleton_subset_intro,subseteq_antisymmetric}.
            $h_0$ is a bijection from $\initialSegment{m_0}$ to $\{\beta\}$
                by \cref{well_ordered_enumeration_of_subset,real_lt_subst_right}.
            $a_0\in\funs{\initialSegment{1}}{\reals}$ and $\sumFiniteSequence{a_0}{s_0}$ and
                for all $i\in\initialSegment{1}$ we have $a_0(i) = \apply{\apply{f}{g_0(i)}}{x}$
                by \cref{real_sum_sequence_along_enumeration}.
            Show that for all $i\in\initialSegment{1}$ we have $a_0(i) = \apply{\apply{f}{\beta}}{x}$.
            \begin{subproof}
                Fix $i\in\initialSegment{1}$.
                $i = 1$ by \cref{initial_segment_naturalsplus,real_one_in_naturals,real_one_leq_naturalsplus,real_naturals_subset_reals,subseteq,real_lt_trans,real_lt_irreflexive}.
                $a_0(i) = \apply{\apply{f}{\beta}}{x}$ by \cref{real_lt_subst_left}.
            \end{subproof}
            $s_0 = \apply{\apply{f}{\beta}}{x}$
                by \cref{real_one_in_naturals,singleton_enumeration_sum_value,real_lt_subst_right}.
            $b_0\in\funs{\initialSegment{m_0}}{\reals}$ and $\sumFiniteSequence{b_0}{t_0}$ and
                for all $i\in\initialSegment{m_0}$ we have $b_0(i) = \apply{\apply{f}{h_0(i)}}{x}$
                by \cref{real_sum_sequence_along_enumeration}.
            Show that for all $i\in\initialSegment{m_0}$ we have $b_0(i) = \apply{\apply{f}{\beta}}{x}$.
            \begin{subproof}
                Fix $i\in\initialSegment{m_0}$.
                $i\in\dom{h_0}$ by \cref{well_ordered_enumeration_of_subset,bijections_dom,subseteq}.
                $h_0(i)\in G_0$ by \cref{well_ordered_enumeration_of_subset,bijections,surj,funs_type_apply,subseteq}.
                $h_0(i) = \beta$ by \cref{real_lt_subst_left,singleton_elim}.
                $b_0(i) = \apply{\apply{f}{\beta}}{x}$ by \cref{real_lt_subst_left}.
            \end{subproof}
            $t_0 = \apply{\apply{f}{\beta}}{x}$
                by \cref{singleton_enumeration_sum_value,real_lt_subst_right}.
            $s_0 = t_0$ by \cref{real_lt_subst_left}.
        \end{subproof}
        $1\in A$ by \cref{real_one_in_naturals}.
    \end{subproof}
    Show that for all $k\in A$ we have $k + 1\in A$.
    \begin{subproof}
        Fix $k\in A$.
        Show that for all $F_0,G_0,g_0,h_0,m_0,a_0,b_0,s_0,t_0$ such that
            $g_0$ is a well-ordered enumeration of $F_0$ in $J$ with respect to $R$ and $k + 1$ and
            $h_0$ is a well-ordered enumeration of $G_0$ in $J$ with respect to $R$ and $m_0$ and
            $a_0$ is a real sum sequence of $f$ at $x$ along $g_0$ with respect to $k + 1$ and result $s_0$ and
            $b_0$ is a real sum sequence of $f$ at $x$ along $h_0$ with respect to $m_0$ and result $t_0$ and
            $G_0\subseteq F_0$ and
            for all $\alpha\in F_0\setminus G_0$ we have $\apply{\apply{f}{\alpha}}{x} = 0$
            we have $s_0 = t_0$.
        \begin{subproof}
            Fix $F_0$.
            Fix $G_0$.
            Fix $g_0$.
            Fix $h_0$.
            Fix $m_0$.
            Fix $a_0$.
            Fix $b_0$.
            Fix $s_0$.
            Fix $t_0$ such that
                $g_0$ is a well-ordered enumeration of $F_0$ in $J$ with respect to $R$ and $k + 1$ and
                $h_0$ is a well-ordered enumeration of $G_0$ in $J$ with respect to $R$ and $m_0$ and
                $a_0$ is a real sum sequence of $f$ at $x$ along $g_0$ with respect to $k + 1$ and result $s_0$ and
                $b_0$ is a real sum sequence of $f$ at $x$ along $h_0$ with respect to $m_0$ and result $t_0$ and
                $G_0\subseteq F_0$ and
                for all $\alpha\in F_0\setminus G_0$ we have $\apply{\apply{f}{\alpha}}{x} = 0$.
            $F_0\subseteq J$ and $g_0$ is a bijection from $\initialSegment{k + 1}$ to $F_0$ and
                $g_0$ is increasing on $\initialSegment{k + 1}$ with respect to $R$
                by \cref{well_ordered_enumeration_of_subset}.
            Let $\beta = g_0(k + 1)$.
            Let $E = F_0\setminus \{\beta\}$.
            Let $g_1 = \restrl{g_0}{\initialSegment{k}}$.
            $k\in\naturalsPlus$ by \cref{subseteq}.
            $g_1$ is a well-ordered enumeration of $E$ in $J$ with respect to $R$ and $k$
                by \cref{well_ordered_enumeration_remove_last}.
            $a_0\in\funs{\initialSegment{k + 1}}{\reals}$ and $\sumFiniteSequence{a_0}{s_0}$ and
                for all $i\in\initialSegment{k + 1}$ we have $a_0(i) = \apply{\apply{f}{g_0(i)}}{x}$
                by \cref{real_sum_sequence_along_enumeration}.
            Let $a_1 = \restrl{a_0}{\initialSegment{k}}$.
            $a_1\in\funs{\initialSegment{k}}{\reals}$ and
                there exists $u_0\in\reals$ such that $\sumFiniteSequence{a_1}{u_0}$ and $s_0 = u_0 + a_0(k + 1)$
                by \cref{sum_finite_sequence_split_last}.
            Take $u_0\in\reals$ such that $\sumFiniteSequence{a_1}{u_0}$ and $s_0 = u_0 + a_0(k + 1)$.
            Show that for all $i\in\initialSegment{k}$ we have $a_1(i) = \apply{\apply{f}{g_1(i)}}{x}$.
            \begin{subproof}
                Fix $i\in\initialSegment{k}$.
                $\initialSegment{k}\subseteq \dom{a_0}$ and $\initialSegment{k}\subseteq \dom{g_0}$
                    by \cref{real_sum_sequence_along_enumeration,funs,well_ordered_enumeration_of_subset,bijections_dom,initial_segment_subset_succ,subseteq,subseteq_transitive}.
                $a_1(i) = a_0(i)$ and $g_1(i) = g_0(i)$
                    by \cref{restrl_of_function_apply,funs_is_function,bijection_is_function}.
                $i\in\initialSegment{k + 1}$ by \cref{initial_segment_subset_succ,subseteq}.
                $a_0(i) = \apply{\apply{f}{g_0(i)}}{x}$ by \cref{real_sum_sequence_along_enumeration}.
                $a_1(i) = \apply{\apply{f}{g_1(i)}}{x}$ by \cref{real_lt_subst_left,real_lt_subst_right}.
            \end{subproof}
            $a_1$ is a real sum sequence of $f$ at $x$ along $g_1$ with respect to $k$ and result $u_0$
                by \cref{real_sum_sequence_along_enumeration}.
            \begin{byCase}
            \caseOf{$\beta\notin G_0$.}
                $\beta\in F_0$ by \cref{real_naturals_closed_succ,initial_segment_naturalsplus,well_ordered_enumeration_of_subset,bijections,surj,funs_type_apply,subseteq}.
                $\beta\in F_0\setminus G_0$ by \cref{setminus_intro}.
                $\apply{\apply{f}{\beta}}{x} = 0$ by assumption.
                $k + 1\in\initialSegment{k + 1}$ by \cref{real_naturals_closed_succ,initial_segment_naturalsplus}.
                $a_0(k + 1) = \apply{\apply{f}{\beta}}{x}$ by \cref{real_sum_sequence_along_enumeration,real_lt_subst_left}.
                $a_0(k + 1) = 0$ by \cref{real_lt_subst_right}.
                $s_0 = u_0$ by \cref{real_add_ident,real_add_comm,real_lt_subst_right}.
                $a_1$ is a real sum sequence of $f$ at $x$ along $g_1$ with respect to $k$ and result $s_0$
                    by \cref{real_lt_subst_right}.
                Show that for all $y\in G_0$ we have $y\in E$.
                \begin{subproof}
                    Fix $y\in G_0$.
                    $y\in F_0$ by \cref{subseteq}.
                    $y\neq \beta$ by \cref{real_lt_subst_left}.
                    $y\notin \{\beta\}$ by \cref{singleton_iff}.
                    $y\in E$ by \cref{setminus_intro}.
                \end{subproof}
                $G_0\subseteq E$ by \cref{subseteq}.
                Show that for all $\alpha\in E\setminus G_0$ we have $\apply{\apply{f}{\alpha}}{x} = 0$.
                \begin{subproof}
                    Fix $\alpha\in E\setminus G_0$.
                    $\alpha\in E$ and $\alpha\notin G_0$ by \cref{setminus}.
                    $\alpha\in F_0$ by \cref{setminus_subseteq,subseteq}.
                    $\alpha\in F_0\setminus G_0$ by \cref{setminus_intro}.
                    $\apply{\apply{f}{\alpha}}{x} = 0$ by assumption.
                \end{subproof}
                $s_0 = t_0$ by \cref{subseteq}.
            \caseOf{$\beta\in G_0$.}
                $m_0\in\naturalsPlus$ and $G_0\subseteq J$ and $h_0$ is a bijection from $\initialSegment{m_0}$ to $G_0$ and
                    $h_0$ is increasing on $\initialSegment{m_0}$ with respect to $R$
                    by \cref{well_ordered_enumeration_of_subset}.
                $\beta$ is a largest element of $F_0$ with respect to $R$
                    by \cref{simple_order_relation,real_naturals_closed_succ,increasing_enumeration_last_largest}.
                Show that $\beta$ is a largest element of $G_0$ with respect to $R$.
                \begin{subproof}
                    Show that for all $y\in G_0$ we have $y\mathrel{R}\beta$ or $y = \beta$.
                    \begin{subproof}
                        Fix $y\in G_0$.
                        $y\in F_0$ by \cref{subseteq}.
                        $y\mathrel{R}\beta$ or $y = \beta$ by \cref{largest_element}.
                    \end{subproof}
                    $\beta$ is a largest element of $G_0$ with respect to $R$ by \cref{largest_element,subseteq}.
                \end{subproof}
                Let $\gamma = h_0(m_0)$.
                $\gamma$ is a largest element of $G_0$ with respect to $R$
                    by \cref{simple_order_relation,increasing_enumeration_last_largest}.
                $\gamma\in G_0$ by \cref{largest_element}.
                $\gamma\mathrel{R}\beta$ or $\gamma = \beta$ by \cref{largest_element}.
                $\beta\mathrel{R}\gamma$ or $\beta = \gamma$ by \cref{largest_element}.
                $\gamma = \beta$ by \cref{simple_order_relation,asymmetric}.
                $h_0(m_0) = \beta$ by \cref{real_lt_subst_right}.
                \begin{byCase}
                \caseOf{$G_0 = \{\beta\}$.}
                    Show that for all $i\in\initialSegment{k}$ we have $a_1(i) = 0$.
                    \begin{subproof}
                        Fix $i\in\initialSegment{k}$.
                        $i\in\dom{g_1}$ by \cref{well_ordered_enumeration_of_subset,bijections_dom,subseteq}.
                        $g_1(i)\in E$ by \cref{well_ordered_enumeration_of_subset,bijections,surj,funs_type_apply,subseteq}.
                        $g_1(i)\in F_0\setminus G_0$ by \cref{setminus_intro,real_lt_subst_right}.
                        $\apply{\apply{f}{g_1(i)}}{x} = 0$ by assumption.
                        $a_1(i) = 0$ by \cref{real_lt_subst_right}.
                    \end{subproof}
                    $u_0 = 0$ by \cref{sum_finite_sequence_all_zero}.
                    $a_0(k + 1)\in\reals$
                        by \cref{real_naturals_closed_succ,initial_segment_naturalsplus,real_sum_sequence_along_enumeration,funs_type_apply,subseteq}.
                    $u_0 + a_0(k + 1) = 0 + a_0(k + 1)$ by \cref{real_lt_subst_left}.
                    $0 + a_0(k + 1) = a_0(k + 1)$ by \cref{real_add_ident}.
                    $s_0 = 0 + a_0(k + 1)$ by \cref{real_lt_subst_right}.
                    $s_0 = a_0(k + 1)$ by \cref{real_lt_subst_right}.
                    $a_0(k + 1) = \apply{\apply{f}{\beta}}{x}$ by \cref{real_naturals_closed_succ,initial_segment_naturalsplus,real_lt_subst_left}.
                    $s_0 = \apply{\apply{f}{\beta}}{x}$ by \cref{real_lt_subst_right}.
                    $b_0\in\funs{\initialSegment{m_0}}{\reals}$ and $\sumFiniteSequence{b_0}{t_0}$ and
                        for all $i\in\initialSegment{m_0}$ we have $b_0(i) = \apply{\apply{f}{h_0(i)}}{x}$
                        by \cref{real_sum_sequence_along_enumeration}.
                    Show that for all $i\in\initialSegment{m_0}$ we have $b_0(i) = \apply{\apply{f}{\beta}}{x}$.
                    \begin{subproof}
                        Fix $i\in\initialSegment{m_0}$.
                        $i\in\dom{h_0}$ by \cref{well_ordered_enumeration_of_subset,bijections_dom,subseteq}.
                        $h_0(i)\in G_0$ by \cref{well_ordered_enumeration_of_subset,bijections,surj,funs_type_apply,subseteq}.
                        $h_0(i) = \beta$ by \cref{real_lt_subst_left,singleton_elim}.
                        $b_0(i) = \apply{\apply{f}{\beta}}{x}$ by \cref{real_lt_subst_left}.
                    \end{subproof}
                    $t_0 = \apply{\apply{f}{\beta}}{x}$
                        by \cref{singleton_enumeration_sum_value,real_lt_subst_right}.
                    $s_0 = t_0$ by \cref{real_lt_subst_left}.
                \caseOf{$G_0\neq \{\beta\}$.}
                    Show that $m_0\neq 1$.
                    \begin{subproof}
                        Suppose not.
                        $m_0 = 1$.
                        Let $\gamma_0 = h_0(1)$.
                        Show that for all $y\in G_0$ we have $y = \gamma_0$.
                        \begin{subproof}
                            Fix $y\in G_0$.
                            Take $u\in\initialSegment{1}$ such that $h_0(u) = y$ by \cref{bijections,surj,real_lt_subst_left}.
                            $u = 1$ by \cref{initial_segment_naturalsplus,real_one_in_naturals,real_one_leq_naturalsplus,real_naturals_subset_reals,subseteq,real_lt_trans,real_lt_irreflexive}.
                            $y = \gamma_0$ by \cref{real_lt_subst_left}.
                        \end{subproof}
                        $G_0\subseteq \{\gamma_0\}$ by \cref{subseteq,singleton_intro}.
                        $\{\gamma_0\}\subseteq G_0$ by \cref{singleton_subset_intro,real_one_in_naturals,real_one_leq_naturalsplus,initial_segment_naturalsplus,bijections,surj,funs_type_apply,subseteq}.
                        $G_0 = \{\gamma_0\}$ by \cref{subseteq_antisymmetric}.
                        $\gamma_0 = \beta$ by \cref{real_lt_subst_left}.
                        Contradiction.
                    \end{subproof}
                    Take $p\in\naturalsPlus$ such that $m_0 = p + 1$ by \cref{real_naturals_predecessor}.
                    Let $H = G_0\setminus \{\beta\}$.
                    Let $h_1 = \restrl{h_0}{\initialSegment{p}}$.
                    $h_1$ is a well-ordered enumeration of $H$ in $J$ with respect to $R$ and $p$
                        by \cref{real_lt_subst_left,well_ordered_enumeration_remove_last}.
                    $b_0\in\funs{\initialSegment{m_0}}{\reals}$ and $\sumFiniteSequence{b_0}{t_0}$ and
                        for all $i\in\initialSegment{m_0}$ we have $b_0(i) = \apply{\apply{f}{h_0(i)}}{x}$
                        by \cref{real_sum_sequence_along_enumeration}.
                    Let $b_1 = \restrl{b_0}{\initialSegment{p}}$.
                    $b_1\in\funs{\initialSegment{p}}{\reals}$ and
                        there exists $v_0\in\reals$ such that $\sumFiniteSequence{b_1}{v_0}$ and $t_0 = v_0 + b_0(m_0)$
                        by \cref{real_lt_subst_left,sum_finite_sequence_split_last}.
                    Take $v_0\in\reals$ such that $\sumFiniteSequence{b_1}{v_0}$ and $t_0 = v_0 + b_0(m_0)$.
                    Show that for all $i\in\initialSegment{p}$ we have $b_1(i) = \apply{\apply{f}{h_1(i)}}{x}$.
                    \begin{subproof}
                        Fix $i\in\initialSegment{p}$.
                        $\initialSegment{p}\subseteq \dom{b_0}$ and $\initialSegment{p}\subseteq \dom{h_0}$
                            by \cref{real_sum_sequence_along_enumeration,funs,well_ordered_enumeration_of_subset,bijections_dom,initial_segment_subset_succ,subseteq,subseteq_transitive,real_lt_subst_left}.
                        $b_1(i) = b_0(i)$ and $h_1(i) = h_0(i)$
                            by \cref{restrl_of_function_apply,funs_is_function,bijection_is_function}.
                        $i\in\initialSegment{p + 1}$ by \cref{initial_segment_subset_succ,subseteq}.
                        $b_0(i) = \apply{\apply{f}{h_0(i)}}{x}$ by \cref{real_sum_sequence_along_enumeration,real_lt_subst_left}.
                        $b_1(i) = \apply{\apply{f}{h_1(i)}}{x}$ by \cref{real_lt_subst_left,real_lt_subst_right}.
                    \end{subproof}
                    $b_1$ is a real sum sequence of $f$ at $x$ along $h_1$ with respect to $p$ and result $v_0$
                        by \cref{real_sum_sequence_along_enumeration}.
                    $a_0(k + 1) = \apply{\apply{f}{\beta}}{x}$
                        by \cref{real_naturals_closed_succ,initial_segment_naturalsplus,real_lt_subst_left}.
                    $m_0\in\initialSegment{m_0}$ by \cref{initial_segment_naturalsplus}.
                    $b_0(m_0) = \apply{\apply{f}{h_0(m_0)}}{x}$ by \cref{real_sum_sequence_along_enumeration}.
                    $b_0(m_0) = \apply{\apply{f}{\beta}}{x}$ by \cref{real_lt_subst_right}.
                    $a_0(k + 1) = b_0(m_0)$ by \cref{real_lt_subst_left}.
                    Show that for all $\alpha\in E\setminus H$ we have $\apply{\apply{f}{\alpha}}{x} = 0$.
                    \begin{subproof}
                        Fix $\alpha\in E\setminus H$.
                        $\alpha\in E$ and $\alpha\notin H$ by \cref{setminus}.
                        $\alpha\in F_0$ and $\alpha\notin \{\beta\}$ by \cref{setminus}.
                        Show that $\alpha\notin G_0$.
                        \begin{subproof}
                            Suppose not.
                            $\alpha\in G_0$.
                            $\alpha\in G_0\setminus \{\beta\}$ by \cref{setminus_intro}.
                            $\alpha\in H$ by \cref{real_lt_subst_right}.
                            Contradiction.
                        \end{subproof}
                        $\alpha\in F_0\setminus G_0$ by \cref{setminus_intro}.
                        $\apply{\apply{f}{\alpha}}{x} = 0$ by assumption.
                    \end{subproof}
                    Show that for all $y\in H$ we have $y\in E$.
                    \begin{subproof}
                        Fix $y\in H$.
                        $y\in G_0$ and $y\notin \{\beta\}$ by \cref{setminus}.
                        $y\in F_0$ by \cref{subseteq}.
                        $y\in E$ by \cref{setminus_intro}.
                    \end{subproof}
                    $H\subseteq E$ by \cref{subseteq}.
                    $u_0 = v_0$ by \cref{subseteq}.
                    $s_0 = v_0 + b_0(m_0)$ by \cref{real_lt_subst_left,real_lt_subst_right}.
                    $s_0 = t_0$ by \cref{real_lt_subst_left}.
                \end{byCase}
            \end{byCase}
        \end{subproof}
        $k + 1\in A$ by \cref{real_naturals_closed_succ}.
    \end{subproof}
    $A\subseteq\naturalsPlus$ by \cref{subseteq}.
    $n\in\naturalsPlus$ by \cref{well_ordered_enumeration_of_subset}.
    $n\in A$ by \cref{subseteq,real_naturals_induction}.
    For all $F_0,G_0,g_0,h_0,m_0,a_0,b_0,s_0,t_0$ such that
        $g_0$ is a well-ordered enumeration of $F_0$ in $J$ with respect to $R$ and $n$ and
        $h_0$ is a well-ordered enumeration of $G_0$ in $J$ with respect to $R$ and $m_0$ and
        $a_0$ is a real sum sequence of $f$ at $x$ along $g_0$ with respect to $n$ and result $s_0$ and
        $b_0$ is a real sum sequence of $f$ at $x$ along $h_0$ with respect to $m_0$ and result $t_0$ and
        $G_0\subseteq F_0$ and
        for all $\alpha\in F_0\setminus G_0$ we have $\apply{\apply{f}{\alpha}}{x} = 0$
        we have $s_0 = t_0$ by \cref{subseteq}.
    $s = t$ by assumption.
\end{proof}

% from §41 helper lemma 27: support inclusion from pointwise zero implication
\begin{lemma}\label{support_subseteq_from_zero_implication}
    Suppose $f$ is continuous from $X$ to $\reals$ with respect to $T$ and $\standardTopologyR$.
    Suppose $g$ is continuous from $X$ to $\reals$ with respect to $T$ and $\standardTopologyR$.
    Suppose for all $x\in X$ if $f(x) = 0$ then $g(x) = 0$.
    Then $\supp{g}{X}{T}\subseteq \supp{f}{X}{T}$.
\end{lemma}
\begin{proof}
    Let $A = X\setminus \preimg{g}{\{0\}}$.
    Let $B = X\setminus \preimg{f}{\{0\}}$.
    Show that $A\subseteq B$.
    \begin{subproof}
        Show that for all $x$ such that $x\in A$ we have $x\in B$.
        \begin{subproof}
            Fix $x$ such that $x\in A$.
            Show that $x\in B$.
            \begin{subproof}
                Suppose not.
                $x\in X$ by \cref{setminus_elim_left}.
                $x\in \preimg{f}{\{0\}}$ by \cref{setminus_elim_left,setminus_intro}.
                Take $r\in\{0\}$ such that $(x,r)\in f$ by \cref{preimg_iff}.
                $f(x) = 0$ by \cref{singleton_elim,continuous,funs_is_function,funs,function_apply_intro}.
                $g(x) = 0$ by assumption.
                $g$ is a function and $\dom{g} = X$ by \cref{continuous,funs,funs_is_function}.
                $x\in\dom{g}$ by \cref{setminus_elim_left,subseteq}.
                $(x,g(x))\in g$ by \cref{function_apply_iff}.
                $x\in \preimg{g}{\{0\}}$
                    by \cref{singleton_intro,preimg_iff,real_lt_subst_right}.
                Contradiction by \cref{setminus_elim_right}.
            \end{subproof}
        \end{subproof}
    \end{subproof}
    $T$ is a topology on $X$ by \cref{continuous}.
    $A\subseteq X$ by \cref{setminus_subseteq}.
    $B\subseteq X$ by \cref{setminus_subseteq}.
    $B\subseteq \supp{f}{X}{T}$ by \cref{subseteq_closure,support}.
    $A\subseteq \supp{f}{X}{T}$ by \cref{subseteq_transitive}.
    $\supp{f}{X}{T}$ is closed in $X$ with respect to $T$ by \cref{support,closure_closed_in}.
    $\closure{A}{X}{T}\subseteq \supp{f}{X}{T}$ by \cref{closure_subseteq_closed}.
    $\supp{g}{X}{T}\subseteq \supp{f}{X}{T}$ by \cref{support,real_lt_subst_left}.
\end{proof}

% from §41 helper lemma 28: continuous sums of locally finite real-valued indexed families
\begin{lemma}\label{locally_finite_real_sum_continuous}
    Assume $T$ is a topology on $X$.
    Assume $R$ is a well order relation on $J$.
    Assume $f$ is a function on $J$.
    Assume $S$ is the support family of $f$ on $X$ with respect to $T$ and $J$.
    Suppose $S$ is a locally finite indexed family on $J$ in $X$ with respect to $T$.
    Suppose for all $\alpha\in J$ we have $\apply{f}{\alpha}$ is continuous from $X$ to $\reals$ with respect to $T$ and $\standardTopologyR$.
    Suppose for all $x\in X$ there exists $\beta\in J$ such that $x\in\apply{S}{\beta}$.
    Then there exists $g$ such that
        $g$ is continuous from $X$ to $\reals$ with respect to $T$ and $\standardTopologyR$ and
        for all $x\in X$ we have $g(x)$ is the pointwise sum of $f$ at $x$ on $X$ with respect to $T$ and $J$.
\end{lemma}
\begin{proof}
    Show that for all $w$ such that $w\in S$ there exist $\alpha, U$ such that $w = (\alpha,U)$.
    \begin{subproof}
        Fix $w$ such that $w\in S$.
        Take $i\in J$ such that $w = (i,\supp{\apply{f}{i}}{X}{T})$
            by \cref{support_family}.
        There exist $\alpha, U$ such that $w = (\alpha,U)$.
    \end{subproof}
    $S$ is a relation by \cref{relation}.
    Show that for all $\alpha,U_1,U_2$ such that $(\alpha,U_1)\in S$ and $(\alpha,U_2)\in S$ we have $U_1 = U_2$.
    \begin{subproof}
        Fix $\alpha$.
        Fix $U_1$.
        Fix $U_2$.
        Assume $(\alpha,U_1)\in S$ and $(\alpha,U_2)\in S$.
        Take $\beta\in J$ such that $(\alpha,U_1) = (\beta,\supp{\apply{f}{\beta}}{X}{T})$
            by \cref{support_family}.
        Take $\gamma\in J$ such that $(\alpha,U_2) = (\gamma,\supp{\apply{f}{\gamma}}{X}{T})$
            by \cref{support_family}.
        Hence $U_1 = U_2$ by \cref{pair_eq_iff}.
    \end{subproof}
    $S$ is right-unique by \cref{rightunique}.
    For all $\alpha\in J$ we have $(\alpha,\supp{\apply{f}{\alpha}}{X}{T})\in S$
        by \cref{support_family,pair_eq_iff}.
    $J\subseteq\dom{S}$ by \cref{dom_intro,subseteq}.
    Show that for all $\alpha\in\dom{S}$ we have $\alpha\in J$.
    \begin{subproof}
        Fix $\alpha\in\dom{S}$.
        Take $U$ such that $(\alpha,U)\in S$ by \cref{dom_iff}.
        Take $\beta\in J$ such that $(\alpha,U) = (\beta,\supp{\apply{f}{\beta}}{X}{T})$
            by \cref{support_family}.
        Hence $\alpha\in J$ by \cref{pair_eq_iff}.
    \end{subproof}
    $\dom{S}\subseteq J$ by \cref{subseteq}.
    $\dom{S} = J$ by \cref{subseteq_antisymmetric}.
    Thus $S$ is a function on $J$
        by \cref{relation,rightunique,dom_iff,dom_intro,subseteq,subseteq_antisymmetric}.
    Let $P = \{w\in X\times\reals \mid \text{there exist $x\in X$ such that there exist $s\in\reals$ such that there exist $E,n,h,a$ such that
        $w = (x,s)$ and
        $E = \{\alpha\in J \mid x\in\apply{S}{\alpha}\}$ and
        $h$ is a well-ordered enumeration of $E$ in $J$ with respect to $R$ and $n$ and
        $a$ is a real sum sequence of $f$ at $x$ along $h$ with respect to $n$ and result $s$}\}$.
    Show that for all $x\in X$ there exists $s$ such that $x\mathrel{P}s$.
    \begin{subproof}
        Fix $x\in X$.
        Take $\beta\in J$ such that $x\in\apply{S}{\beta}$.
        Let $E = \{\alpha\in J \mid x\in\apply{S}{\alpha}\}$.
        $E\subseteq J$ and $E$ is finite and $E$ is inhabited
            by \cref{local_support_index_set_finite_inhabited}.
        Take $n,h$ such that
            $n\in\naturalsPlus$ and
            $h$ is a bijection from $\initialSegment{n}$ to $E$ and
            $h$ is increasing on $\initialSegment{n}$ with respect to $R$
            by \cref{finite_well_ordered_enumeration}.
        $h$ is a well-ordered enumeration of $E$ in $J$ with respect to $R$ and $n$
            by \cref{well_ordered_enumeration_of_subset}.
        Let $a = \{(i,\apply{\apply{f}{h(i)}}{x}) \mid i\in\initialSegment{n}\}$.
        For all $i,r_1,r_2$ such that $(i,r_1)\in a$ and $(i,r_2)\in a$ we have $r_1 = r_2$ by \cref{pair_eq_iff}.
        $a$ is a function by \cref{relation,rightunique}.
        For all $i$ such that $i\in\dom{a}$ we have $i\in\initialSegment{n}$ by \cref{dom_iff,pair_eq_iff}.
        For all $i$ such that $i\in\initialSegment{n}$ we have $i\in\dom{a}$ by \cref{pair_eq_iff,dom_intro}.
        $\dom{a} = \initialSegment{n}$ by \cref{subseteq,subseteq_antisymmetric}.
        Show that for all $i\in\dom{a}$ we have $a(i)\in\reals$.
        \begin{subproof}
            Fix $i\in\dom{a}$.
            $i\in\initialSegment{n}$ by \cref{real_lt_subst_left}.
            $i\in\dom{h}$ by \cref{well_ordered_enumeration_of_subset,bijections_dom,subseteq}.
            $h(i)\in E$ by \cref{well_ordered_enumeration_of_subset,bijections,surj,funs_type_apply,subseteq}.
            $h(i)\in J$ by \cref{subseteq}.
            $\apply{f}{h(i)}$ is continuous from $X$ to $\reals$ with respect to $T$ and $\standardTopologyR$ by assumption.
            $a(i)\in\reals$ by \cref{dom_iff,pair_eq_iff,function_apply_intro,continuous,funs_type_apply,real_lt_subst_right}.
        \end{subproof}
        $a\in\funs{\initialSegment{n}}{\reals}$ by \cref{funs_intro}.
        For all $i\in\initialSegment{n}$ we have $a(i) = \apply{\apply{f}{h(i)}}{x}$
            by \cref{pair_eq_iff,function_apply_intro}.
        Take $s\in\reals$ such that $\sumFiniteSequence{a}{s}$ by \cref{sum_finite_sequence_exists}.
        $a$ is a real sum sequence of $f$ at $x$ along $h$ with respect to $n$ and result $s$
            by \cref{real_sum_sequence_along_enumeration}.
        $x\mathrel{P}s$ by \cref{pair_eq_iff,times_tuple_intro,powerset_intro}.
    \end{subproof}
    $P$ is a relation by \cref{relation}.
    Take $g$ such that $g$ is a function on $X$ and for all $x\in X$ we have $x\mathrel{P}g(x)$
        by \cref{choice_function}.
    For all $x\in\dom{g}$ we have $g(x)\in\reals$ by \cref{choice_function,pair_eq_iff}.
    $g\in\funs{X}{\reals}$ by \cref{choice_function,funs_intro}.

    Show that for all $x\in X$ we have $g(x)$ is the pointwise sum of $f$ at $x$ on $X$ with respect to $T$ and $J$.
    \begin{subproof}
        Fix $x\in X$.
        Take $E,n,h,a$ such that
            $E = \{\alpha\in J \mid x\in\apply{S}{\alpha}\}$ and
            $h$ is a well-ordered enumeration of $E$ in $J$ with respect to $R$ and $n$ and
            $a$ is a real sum sequence of $f$ at $x$ along $h$ with respect to $n$ and result $g(x)$
            by \cref{choice_function,pair_eq_iff}.
        $E\subseteq J$ and $n\in\naturalsPlus$ and $h$ is a bijection from $\initialSegment{n}$ to $E$
            by \cref{well_ordered_enumeration_of_subset}.
        $E$ is finite by \cref{initial_segment_finite,finite_bijection}.
        $1\in\initialSegment{n}$ by \cref{real_one_in_naturals,real_one_leq_naturalsplus,initial_segment_naturalsplus,well_ordered_enumeration_of_subset}.
        $E$ is inhabited by \cref{well_ordered_enumeration_of_subset,bijections,surj,funs_type_apply,subseteq,nonempty_inhabited}.
        $a\in\funs{\initialSegment{n}}{\reals}$ and $\sumFiniteSequence{a}{g(x)}$ and
            for all $i\in\initialSegment{n}$ we have $a(i) = \apply{\apply{f}{h(i)}}{x}$
            by \cref{real_sum_sequence_along_enumeration}.
        Show that for all $\alpha\in J$ if $x\in\supp{\apply{f}{\alpha}}{X}{T}$ then $\alpha\in E$.
        \begin{subproof}
            Fix $\alpha\in J$.
            Assume $x\in\supp{\apply{f}{\alpha}}{X}{T}$.
            $x\in\apply{S}{\alpha}$ by \cref{support_family,pair_eq_iff,function_apply_intro}.
            $\alpha\in E$ by \cref{subseteq,real_lt_subst_left}.
        \end{subproof}
        $g(x)$ is the pointwise sum of $f$ at $x$ on $X$ with respect to $T$ and $J$
            by \cref{pointwise_sum_of_indexed_family}.
    \end{subproof}

    Show that for all $x\in X$ $g$ is continuous at $x$ from $X$ to $\reals$ with respect to $T$ and $\standardTopologyR$.
    \begin{subproof}
        Fix $x\in X$.
        Take $\beta\in J$ such that $x\in\apply{S}{\beta}$.
        Take $W,F$ such that
            $W$ is a neighborhood of $x$ in $X$ with respect to $T$ and
            $F\subseteq J$ and
            $F$ is finite and
            for all $\alpha\in J$ if $\alpha\notin F$ then for all $y\in W$ we have $\apply{\apply{f}{\alpha}}{y} = 0$
            by \cref{support_family_local_zero_control}.
        $W\subseteq X$ by \cref{neighborhood,open_in,topology_on,subseteq,pow_iff}.
        $\beta\in F$
            by \cref{locally_finite_indexed_family,indexed_family_on_in,support_family_local_zero_control,support_family,pair_eq_iff,function_apply_intro,local_zero_neighborhood_excludes_support,real_lt_subst_right}.
        Take $n,h$ such that
            $n\in\naturalsPlus$ and
            $h$ is a bijection from $\initialSegment{n}$ to $F$ and
            $h$ is increasing on $\initialSegment{n}$ with respect to $R$
            by \cref{subseteq,nonempty_inhabited,finite_well_ordered_enumeration}.
        $h$ is a well-ordered enumeration of $F$ in $J$ with respect to $R$ and $n$
            by \cref{well_ordered_enumeration_of_subset}.
        Let $p = \{(i,\apply{f}{h(i)}) \mid i\in\initialSegment{n}\}$.
        For all $i,u_1,u_2$ such that $(i,u_1)\in p$ and $(i,u_2)\in p$ we have $u_1 = u_2$ by \cref{pair_eq_iff}.
        $p$ is a function by \cref{relation,rightunique}.
        For all $i$ such that $i\in\dom{p}$ we have $i\in\initialSegment{n}$ by \cref{dom_iff,pair_eq_iff}.
        For all $i$ such that $i\in\initialSegment{n}$ we have $i\in\dom{p}$ by \cref{pair_eq_iff,dom_intro}.
        $\dom{p} = \initialSegment{n}$ by \cref{subseteq,subseteq_antisymmetric}.
        $p$ is a function on $\initialSegment{n}$ by \cref{funs}.
        For all $i\in\initialSegment{n}$ we have $p(i) = \apply{f}{h(i)}$ by \cref{pair_eq_iff,function_apply_intro}.
        Show that for all $i\in\initialSegment{n}$ we have $\apply{p}{i}$ is continuous from $X$ to $\reals$ with respect to $T$ and $\standardTopologyR$.
        \begin{subproof}
            Fix $i\in\initialSegment{n}$.
            $i\in\dom{h}$ by \cref{well_ordered_enumeration_of_subset,bijections_dom,subseteq}.
            $h(i)\in F$ by \cref{well_ordered_enumeration_of_subset,bijections,surj,funs_type_apply,subseteq}.
            $h(i)\in J$ by \cref{subseteq}.
            $\apply{f}{h(i)}$ is continuous from $X$ to $\reals$ with respect to $T$ and $\standardTopologyR$ by assumption.
            $\apply{p}{i}$ is continuous from $X$ to $\reals$ with respect to $T$ and $\standardTopologyR$ by \cref{continuous_eq_subst}.
        \end{subproof}
        Take $u$ such that
            $u$ is a finite real sum witness for $p$ on $X$ with respect to $T$ and $n$
            by \cref{finite_real_sum_continuous}.
        $u$ is continuous from $X$ to $\reals$ with respect to $T$ and $\standardTopologyR$ and
            for all $y\in X$ there exists $c$ such that
                $c\in\funs{\initialSegment{n}}{\reals}$ and
                $\sumFiniteSequence{c}{u(y)}$ and
                for all $i\in\initialSegment{n}$ we have $c(i) = \apply{\apply{p}{i}}{y}$
            by \cref{finite_real_sum_witness}.

        Show that for all $y\in W$ we have $u(y) = g(y)$.
        \begin{subproof}
            Fix $y\in W$.
            $y\in X$ by \cref{subseteq}.
            $(y,g(y))\in P$ by \cref{choice_function}.
            Take $E,m,k,b$ such that
                $E = \{\alpha\in J \mid y\in\apply{S}{\alpha}\}$ and
                $k$ is a well-ordered enumeration of $E$ in $J$ with respect to $R$ and $m$ and
                $b$ is a real sum sequence of $f$ at $y$ along $k$ with respect to $m$ and result $g(y)$
                by \cref{pair_eq_iff}.
            Show that $E\subseteq F$.
            \begin{subproof}
                Show that for all $\alpha\in E$ we have $\alpha\in F$.
                \begin{subproof}
                    Fix $\alpha\in E$.
                    $\alpha\in J$ and $y\in\apply{S}{\alpha}$ by \cref{subseteq}.
                    Show that $\alpha\in F$.
                    \begin{subproof}
                        Suppose not.
                        $W$ is open in $X$ with respect to $T$ and $x\in W$ by \cref{neighborhood}.
                        $W$ is a neighborhood of $y$ in $X$ with respect to $T$ by \cref{neighborhood}.
                        For all $z\in W$ we have $\apply{\apply{f}{\alpha}}{z} = 0$ by \cref{subseteq}.
                        $y\notin \supp{\apply{f}{\alpha}}{X}{T}$ by \cref{local_zero_neighborhood_excludes_support}.
                        $y\notin \apply{S}{\alpha}$ by \cref{support_family,pair_eq_iff,function_apply_intro}.
                        Contradiction.
                    \end{subproof}
                \end{subproof}
                $E\subseteq F$ by \cref{subseteq}.
            \end{subproof}
            Take $c$ such that
                $c\in\funs{\initialSegment{n}}{\reals}$ and
                $\sumFiniteSequence{c}{u(y)}$ and
                for all $i\in\initialSegment{n}$ we have $c(i) = \apply{\apply{p}{i}}{y}$
                by \cref{neighborhood,subseteq,finite_real_sum_witness}.
            For all $i\in\initialSegment{n}$ we have $c(i) = \apply{\apply{f}{h(i)}}{y}$
                by \cref{real_lt_subst_right}.
            $c$ is a real sum sequence of $f$ at $y$ along $h$ with respect to $n$ and result $u(y)$
                by \cref{real_sum_sequence_along_enumeration}.
            Show that for all $\alpha\in F\setminus E$ we have $\apply{\apply{f}{\alpha}}{y} = 0$.
            \begin{subproof}
                Fix $\alpha\in F\setminus E$.
                $\alpha\in F$ and $\alpha\notin E$ by \cref{setminus_elim_left,setminus_elim_right}.
                $\alpha\in J$ by \cref{subseteq}.
                $y\notin \apply{S}{\alpha}$ by \cref{subseteq,setminus_intro}.
                $y\notin \supp{\apply{f}{\alpha}}{X}{T}$ by \cref{support_family,pair_eq_iff,function_apply_intro}.
                $\apply{f}{\alpha}$ is continuous from $X$ to $\reals$ with respect to $T$ and $\standardTopologyR$ by assumption.
                $y\in X\setminus \supp{\apply{f}{\alpha}}{X}{T}$ by \cref{setminus_intro}.
                $\apply{\apply{f}{\alpha}}{y} = 0$ by \cref{outside_support_zero}.
            \end{subproof}
            $u(y) = g(y)$
                by \cref{well_order_relation,well_ordered_zero_extension_sum_comparison,real_lt_subst_right}.
        \end{subproof}

        Show that $g$ is continuous at $x$ from $X$ to $\reals$ with respect to $T$ and $\standardTopologyR$.
        \begin{subproof}
            $u$ is continuous at $x$ from $X$ to $\reals$ with respect to $T$ and $\standardTopologyR$
                by \cref{continuous_implies_continuous_at}.
            Show that for all $V$ such that $V$ is a neighborhood of $g(x)$ in $\reals$ with respect to $\standardTopologyR$
                there exists $U$ such that
                    $U$ is a neighborhood of $x$ in $X$ with respect to $T$ and
                    $\img{g}{U}\subseteq V$.
            \begin{subproof}
                Fix $V$ such that $V$ is a neighborhood of $g(x)$ in $\reals$ with respect to $\standardTopologyR$.
                $x\in W$ by \cref{neighborhood}.
                $u(x) = g(x)$ by \cref{subseteq}.
                Take $U_0$ such that
                    $U_0$ is a neighborhood of $x$ in $X$ with respect to $T$ and
                    $\img{u}{U_0}\subseteq V$
                    by \cref{continuous_at,real_lt_subst_left}.
                Let $U = U_0\inter W$.
                $U_0$ is open in $X$ with respect to $T$ and $x\in U_0$ by \cref{neighborhood}.
                $W$ is open in $X$ with respect to $T$ and $x\in W$ by \cref{neighborhood}.
                $U$ is open in $X$ with respect to $T$ by \cref{open_in,topology_on,pow_iff}.
                $x\in U$ by \cref{inter_intro}.
                $U$ is a neighborhood of $x$ in $X$ with respect to $T$ by \cref{neighborhood}.
                Show that $\img{g}{U}\subseteq V$.
                \begin{subproof}
                    Show that for all $z\in\img{g}{U}$ we have $z\in V$.
                    \begin{subproof}
                        Fix $z\in\img{g}{U}$.
                        Take $y$ such that $y\in U$ and $(y,z)\in g$ by \cref{img_iff}.
                        $g(y) = z$ by \cref{funs_is_function,function_apply_intro}.
                        $y\in W$ and $y\in U_0$ by \cref{inter_elim_right,inter_elim_left}.
                        $u(y) = g(y)$ by \cref{subseteq}.
                        $u(y) = z$ by \cref{real_lt_subst_right}.
                        $u$ is a function and $\dom{u} = X$ by \cref{continuous,funs,funs_is_function}.
                        $U_0\subseteq X$ by \cref{neighborhood,open_in,topology_on,subseteq,pow_iff}.
                        $y\in\dom{u}$ by \cref{subseteq}.
                        $(y,z)\in u$ by \cref{function_apply_iff,real_lt_subst_right}.
                        $z\in\img{u}{U_0}$ by \cref{img_iff}.
                        $z\in V$ by \cref{subseteq}.
                    \end{subproof}
                \end{subproof}
            \end{subproof}
            $g$ is continuous at $x$ from $X$ to $\reals$ with respect to $T$ and $\standardTopologyR$
                by \cref{continuous_at}.
        \end{subproof}
    \end{subproof}
    $\standardTopologyR$ is a topology on $\reals$
        by \cref{standard_basis_is_basis,basis_topology_is_topology,standard_topology_reals}.
    $g$ is continuous from $X$ to $\reals$ with respect to $T$ and $\standardTopologyR$
        by \cref{funs,continuous_at_implies_continuous}.
\end{proof}

% from §41 helper definition 13: continuous locally finite real sum function
\begin{definition}\label{locally_finite_real_sum_function}
    $g$ is a locally finite real sum function of $f$ on $X$ with respect to $T$ and $J$ iff
        $g$ is continuous from $X$ to $\reals$ with respect to $T$ and $\standardTopologyR$ and
        for all $x\in X$ we have $g(x)$ is the pointwise sum of $f$ at $x$ on $X$ with respect to $T$ and $J$.
\end{definition}

% from §41 helper lemma 30: existence of continuous locally finite real sum functions
\begin{lemma}\label{locally_finite_real_sum_function_exists}
    Assume $T$ is a topology on $X$.
    Assume $R$ is a well order relation on $J$.
    Assume $f$ is a function on $J$.
    Assume $S$ is the support family of $f$ on $X$ with respect to $T$ and $J$.
    Suppose $S$ is a locally finite indexed family on $J$ in $X$ with respect to $T$.
    Suppose for all $\alpha\in J$ we have $\apply{f}{\alpha}$ is continuous from $X$ to $\reals$ with respect to $T$ and $\standardTopologyR$.
    Suppose for all $x\in X$ there exists $\beta\in J$ such that $x\in\apply{S}{\beta}$.
    Then there exists $g$ such that $g$ is a locally finite real sum function of $f$ on $X$ with respect to $T$ and $J$.
\end{lemma}
\begin{proof}
    Follows by \cref{locally_finite_real_sum_continuous,locally_finite_real_sum_function}.
\end{proof}

% from §41 helper lemma 29: nonzero values lie in the support
\begin{lemma}\label{nonzero_value_in_support}
    Assume $f$ is continuous from $X$ to $\reals$ with respect to $T$ and $\standardTopologyR$.
    Assume $x\in X$.
    Suppose $f(x)\neq 0$.
    Then $x\in \supp{f}{X}{T}$.
\end{lemma}
\begin{proof}
    Show that $x\notin \preimg{f}{\{0\}}$.
    \begin{subproof}
        Suppose not.
        Take $r\in\{0\}$ such that $(x,r)\in f$ by \cref{preimg_iff}.
        $f(x) = 0$ by \cref{singleton_elim,continuous,funs_is_function,funs,function_apply_intro}.
        Contradiction.
    \end{subproof}
    $x\in X\setminus \preimg{f}{\{0\}}$ by \cref{setminus_intro}.
    $T$ is a topology on $X$ by \cref{continuous}.
    $X\setminus \preimg{f}{\{0\}}\subseteq X$ by \cref{setminus_subseteq}.
    $X\setminus \preimg{f}{\{0\}}\subseteq \supp{f}{X}{T}$ by \cref{support,subseteq_closure}.
    $x\in \supp{f}{X}{T}$ by \cref{subseteq}.
\end{proof}

% from §41 Item 9: partitions of unity on paracompact Hausdorff spaces
\begin{theorem}\label{paracompact_partition_of_unity}
    Assume $X$ is paracompact with respect to $T$.
    Assume $X$ is Hausdorff with respect to $T$.
    Assume $U$ is a function on $J$.
    Suppose for all $\alpha\in J$ we have $\apply{U}{\alpha}$ is open in $X$ with respect to $T$.
    Suppose $\img{U}{J}$ is a covering of $X$.
    Then there exists $f$ such that
        $f$ is an indexed partition of unity on $X$ dominated by $U$ with respect to $T$ and $J$.
\end{theorem}
\begin{proof}
    Let $I_0 = \intervalclosed{0}{1}$.
    Let $T_0 = \subspaceTopology{\standardTopologyR}{I_0}$.
    Take $V$ such that
        $V$ is a function on $J$ and
        $\img{V}{J}$ is a covering of $X$ and
        $V$ is a locally finite indexed family on $J$ in $X$ with respect to $T$ and
        for all $\alpha\in J$ we have $\apply{V}{\alpha}$ is open in $X$ with respect to $T$ and
        for all $\alpha\in J$ we have $\closure{\apply{V}{\alpha}}{X}{T}\subseteq \apply{U}{\alpha}$
        by \cref{paracompact_shrinking_lemma}.
    Take $W$ such that
        $W$ is a function on $J$ and
        $\img{W}{J}$ is a covering of $X$ and
        $W$ is a locally finite indexed family on $J$ in $X$ with respect to $T$ and
        for all $\alpha\in J$ we have $\apply{W}{\alpha}$ is open in $X$ with respect to $T$ and
        for all $\alpha\in J$ we have $\closure{\apply{W}{\alpha}}{X}{T}\subseteq \apply{V}{\alpha}$
        by \cref{paracompact_shrinking_lemma}.

    $X$ is normal with respect to $T$ by \cref{paracompact_hausdorff_normal}.
    Let $B = \pow{X\times I_0}$.
    Let $M = \{w\in J\times B \mid \text{there exist $\alpha\in J$ such that there exist $u\in B$ such that
        $w = (\alpha,u)$ and
        $u$ is continuous from $X$ to $I_0$ with respect to $T$ and $T_0$ and
        $\img{u}{\closure{\apply{W}{\alpha}}{X}{T}}\subseteq \{1\}$ and
        $\img{u}{X\setminus \apply{V}{\alpha}}\subseteq \{0\}$}\}$.
    Show that for all $\alpha\in J$ there exists $u$ such that $\alpha\mathrel{M}u$.
    \begin{subproof}
        Fix $\alpha\in J$.
        $\apply{W}{\alpha}$ is open in $X$ with respect to $T$ and
            $\closure{\apply{W}{\alpha}}{X}{T}\subseteq \apply{V}{\alpha}$
            by assumption.
        $\apply{V}{\alpha}$ is open in $X$ with respect to $T$ by assumption.
        $T$ is a topology on $X$ by \cref{normal_space}.
        $X\setminus \apply{V}{\alpha}$ is closed in $X$ with respect to $T$
            by \cref{open_in,topology_on,subseteq,pow_iff,closed_in,setminus_subseteq,double_relative_complement}.
        $\closure{\apply{W}{\alpha}}{X}{T}$ is closed in $X$ with respect to $T$
            by \cref{open_in,topology_on,subseteq,pow_iff,closure_closed_in}.
        $X\setminus \apply{V}{\alpha}$ is disjoint from $\closure{\apply{W}{\alpha}}{X}{T}$
            by \cref{subseteq,setminus_elim_right,disjoint}.
        Take $u$ such that
            $u$ is continuous from $X$ to $I_0$ with respect to $T$ and $T_0$ and
            $\img{u}{X\setminus \apply{V}{\alpha}}\subseteq \{0\}$ and
            $\img{u}{\closure{\apply{W}{\alpha}}{X}{T}}\subseteq \{1\}$
            by \cref{urysohn_unit_interval}.
        $\alpha\mathrel{M}u$
            by \cref{continuous,funs_is_relation,funs,funs_subseteq_rels,subseteq,rels_ran_subseteq,rels_elim,powerset_intro,times_tuple_intro,pair_eq_iff}.
    \end{subproof}
    $M$ is a relation by \cref{relation}.
    Take $c$ such that $c$ is a function on $J$ and for all $\alpha\in J$ we have $\alpha\mathrel{M}\apply{c}{\alpha}$
        by \cref{choice_function}.
    Show that for all $\alpha\in J$ we have
        $\apply{c}{\alpha}$ is continuous from $X$ to $I_0$ with respect to $T$ and $T_0$ and
        $\img{\apply{c}{\alpha}}{\closure{\apply{W}{\alpha}}{X}{T}}\subseteq \{1\}$ and
        $\img{\apply{c}{\alpha}}{X\setminus \apply{V}{\alpha}}\subseteq \{0\}$.
    \begin{subproof}
        Follows by \cref{choice_function,pair_eq_iff}.
    \end{subproof}

    Let $j = \identity{I_0}$.
    Let $p = \{(\alpha,j\circ \apply{c}{\alpha}) \mid \alpha\in J\}$.
    For all $\alpha,u_1,u_2$ such that $(\alpha,u_1)\in p$ and $(\alpha,u_2)\in p$ we have $u_1 = u_2$ by \cref{pair_eq_iff}.
    $p$ is a function by \cref{relation,rightunique}.
    For all $\alpha$ such that $\alpha\in\dom{p}$ we have $\alpha\in J$ by \cref{dom_iff,pair_eq_iff}.
    For all $\alpha$ such that $\alpha\in J$ we have $\alpha\in\dom{p}$ by \cref{pair_eq_iff,dom_intro}.
    $\dom{p} = J$ by \cref{subseteq,subseteq_antisymmetric}.
    $p$ is a function on $J$ by \cref{funs}.
    For all $\alpha\in J$ we have $\apply{p}{\alpha} = j\circ \apply{c}{\alpha}$ by \cref{pair_eq_iff,function_apply_intro}.
    For all $\alpha\in J$ we have
        $\apply{p}{\alpha}$ is continuous from $X$ to $\reals$ with respect to $T$ and $\standardTopologyR$
        by \cref{continuous_interval_expand_range,continuous_eq_subst}.
    For all $\alpha\in J$ for all $x\in X$ we have $\apply{\apply{p}{\alpha}}{x} = \apply{\apply{c}{\alpha}}{x}$
        by \cref{continuous,funs_is_function,funs,function_rightunique,funs_is_relation,funs_subseteq_rels,subseteq,rels_ran_subseteq,id_dom,id_rightunique,id_is_relation,circ_apply,funs_type_apply,id_apply}.

    Show that for all $\alpha\in J$ for all $x\in X\setminus \apply{V}{\alpha}$ we have $\apply{\apply{p}{\alpha}}{x} = 0$.
    \begin{subproof}
        Fix $\alpha\in J$.
        Fix $x\in X\setminus \apply{V}{\alpha}$.
        $\apply{c}{\alpha}\in\funs{X}{I_0}$ and $\apply{c}{\alpha}$ is a function and $\dom{\apply{c}{\alpha}} = X$
            by \cref{continuous,funs_is_function,funs}.
        $x\in\dom{\apply{c}{\alpha}}\inter(X\setminus \apply{V}{\alpha})$
            by \cref{setminus_elim_left,subseteq,inter_intro}.
        $\apply{\apply{c}{\alpha}}{x}\in\img{\apply{c}{\alpha}}{X\setminus \apply{V}{\alpha}}$
            by \cref{img_of_function_intro}.
        $\apply{\apply{c}{\alpha}}{x} = 0$ by \cref{subseteq,singleton_elim}.
        $\apply{\apply{p}{\alpha}}{x} = \apply{\apply{c}{\alpha}}{x}$
            by \cref{setminus_elim_left,continuous,funs_is_function,funs,function_rightunique,funs_is_relation,funs_subseteq_rels,subseteq,rels_ran_subseteq,id_dom,id_rightunique,id_is_relation,circ_apply,funs_type_apply,id_apply}.
        $\apply{\apply{p}{\alpha}}{x} = 0$ by \cref{real_lt_subst_right}.
    \end{subproof}

    Let $S = \{(\alpha,\supp{\apply{p}{\alpha}}{X}{T}) \mid \alpha\in J\}$.
    For all $\alpha,A_1,A_2$ such that $(\alpha,A_1)\in S$ and $(\alpha,A_2)\in S$ we have $A_1 = A_2$
        by \cref{pair_eq_iff}.
    $S$ is a function by \cref{relation,rightunique}.
    For all $\alpha$ such that $\alpha\in\dom{S}$ we have $\alpha\in J$ by \cref{dom_iff,pair_eq_iff}.
    For all $\alpha$ such that $\alpha\in J$ we have $\alpha\in\dom{S}$ by \cref{pair_eq_iff,dom_intro}.
    $\dom{S} = J$ by \cref{subseteq,subseteq_antisymmetric}.
    $S$ is a function on $J$ by \cref{funs}.
    $S$ is the support family of $p$ on $X$ with respect to $T$ and $J$ by \cref{support_family}.
    $T$ is a topology on $X$ by \cref{normal_space}.
    $S$ is a locally finite indexed family on $J$ in $X$ with respect to $T$
        by \cref{support_family_locally_finite_from_zero_outside}.

    Show that for all $x\in X$ there exists $\beta\in J$ such that $x\in\apply{S}{\beta}$.
    \begin{subproof}
        Fix $x\in X$.
        Take $Z\in\img{W}{J}$ such that $x\in Z$ by \cref{covering,subseteq,unions_iff}.
        Take $\beta$ such that $\beta\in J$ and $(\beta,Z)\in W$ by \cref{img_iff}.
        $x\in \apply{W}{\beta}$ by \cref{function_apply_intro}.
        $x\in \closure{\apply{W}{\beta}}{X}{T}$
            by \cref{open_in,topology_on,subseteq,pow_iff,subseteq_closure}.
        $\apply{c}{\beta}\in\funs{X}{I_0}$ and $\apply{c}{\beta}$ is a function and $\dom{\apply{c}{\beta}} = X$
            by \cref{continuous,funs_is_function,funs}.
        $\apply{\apply{c}{\beta}}{x}\in\img{\apply{c}{\beta}}{\closure{\apply{W}{\beta}}{X}{T}}$
            by \cref{subseteq,inter_intro,img_of_function_intro}.
        $\apply{\apply{c}{\beta}}{x} = 1$ by \cref{subseteq,singleton_elim}.
        $\apply{\apply{p}{\beta}}{x} = 1$ by \cref{real_lt_subst_right}.
        $\apply{\apply{p}{\beta}}{x}\in\reals$ by \cref{continuous,funs_type_apply}.
        $0 < \apply{\apply{p}{\beta}}{x}$ by \cref{real_zero_lt_one,real_lt_subst_right}.
        $\apply{\apply{p}{\beta}}{x}\neq 0$ by \cref{real_pos_nonzero}.
        $x\in \supp{\apply{p}{\beta}}{X}{T}$ by \cref{nonzero_value_in_support}.
        $S$ is a function on $J$
            by \cref{support_family,relation,pair_eq_iff,rightunique,dom_iff,subseteq,dom_intro,subseteq_antisymmetric,funs}.
        $(\beta,\supp{\apply{p}{\beta}}{X}{T})\in S$ by \cref{support_family,pair_eq_iff}.
        $\apply{S}{\beta} = \supp{\apply{p}{\beta}}{X}{T}$ by \cref{function_apply_intro}.
        $x\in \apply{S}{\beta}$ by \cref{real_lt_subst_left}.
        There exists $\gamma\in J$ such that $x\in\apply{S}{\gamma}$ by \cref{real_lt_subst_left}.
    \end{subproof}

    Take $L$ such that $L$ is a well order relation on $J$ by \cref{well_ordering_theorem_local_finiteness}.
    Take $g$ such that $g$ is a locally finite real sum function of $p$ on $X$ with respect to $T$ and $J$
        by \cref{locally_finite_real_sum_function_exists}.
    $g$ is continuous from $X$ to $\reals$ with respect to $T$ and $\standardTopologyR$ by \cref{locally_finite_real_sum_function}.

    Show that for all $x\in X$ we have $1 \leq g(x)$.
    \begin{subproof}
        Fix $x\in X$.
        Take $Z\in\img{W}{J}$ such that $x\in Z$ by \cref{covering,subseteq,unions_iff}.
        Take $\beta$ such that $\beta\in J$ and $(\beta,Z)\in W$ by \cref{img_iff}.
        $x\in \apply{W}{\beta}$ by \cref{function_apply_intro}.
        $x\in \closure{\apply{W}{\beta}}{X}{T}$
            by \cref{open_in,topology_on,subseteq,pow_iff,subseteq_closure}.
        $\apply{c}{\beta}\in\funs{X}{I_0}$ and $\apply{c}{\beta}$ is a function and $\dom{\apply{c}{\beta}} = X$
            by \cref{continuous,funs_is_function,funs}.
        $\apply{\apply{c}{\beta}}{x}\in\img{\apply{c}{\beta}}{\closure{\apply{W}{\beta}}{X}{T}}$
            by \cref{subseteq,inter_intro,img_of_function_intro}.
        $\apply{\apply{c}{\beta}}{x} = 1$ by \cref{subseteq,singleton_elim}.
        $\apply{\apply{p}{\beta}}{x} = 1$ by \cref{real_lt_subst_right}.
        $\apply{\apply{p}{\beta}}{x}\in\reals$ by \cref{continuous,funs_type_apply}.
        $0 < \apply{\apply{p}{\beta}}{x}$ by \cref{real_zero_lt_one,real_lt_subst_right}.
        $\apply{\apply{p}{\beta}}{x}\neq 0$ by \cref{real_pos_nonzero}.
        $x\in \supp{\apply{p}{\beta}}{X}{T}$ by \cref{nonzero_value_in_support}.
        $g(x)$ is the pointwise sum of $p$ at $x$ on $X$ with respect to $T$ and $J$
            by \cref{subseteq,locally_finite_real_sum_function}.
        Take $F,n,h,a$ such that
            $F\subseteq J$ and
            $F$ is finite and
            $F$ is inhabited and
            $n\in\naturalsPlus$ and
            $h$ is a bijection from $\initialSegment{n}$ to $F$ and
            $a$ is a real sum sequence of $p$ at $x$ along $h$ with respect to $n$ and result $g(x)$ and
            for all $\alpha\in J$ if $x\in\supp{\apply{p}{\alpha}}{X}{T}$ then $\alpha\in F$
            by \cref{pointwise_sum_of_indexed_family}.
        $a\in\funs{\initialSegment{n}}{\reals}$ and $\sumFiniteSequence{a}{g(x)}$
            by \cref{real_sum_sequence_along_enumeration}.
        $\sumFiniteSequence{a}{g(x)}$ by \cref{subseteq}.
        For all $k\in\initialSegment{n}$ we have $a(k) = \apply{\apply{p}{h(k)}}{x}$
            by \cref{real_sum_sequence_along_enumeration}.
        For all $\alpha\in J$ if $x\in\supp{\apply{p}{\alpha}}{X}{T}$ then $\alpha\in F$ by \cref{subseteq}.
        $\beta\in F$ by \cref{subseteq}.
        Take $i\in\initialSegment{n}$ such that $h(i) = \beta$ by \cref{bijections,surj}.
        Show that for all $k\in\dom{a}$ we have $0 < a(k)$ or $0 = a(k)$.
        \begin{subproof}
            Fix $k\in\dom{a}$.
            $k\in\initialSegment{n}$ by \cref{funs,real_lt_subst_left}.
            $k\in\dom{h}$ by \cref{bijections_dom,subseteq}.
            $h(k)\in F$ by \cref{bijections,surj,funs_type_apply,subseteq}.
            $h(k)\in J$ by \cref{subseteq}.
            $\apply{c}{h(k)}\in\funs{X}{I_0}$ and $\apply{c}{h(k)}$ is a function and $\dom{\apply{c}{h(k)}} = X$
                by \cref{continuous,funs_is_function,funs}.
            $\apply{\apply{c}{h(k)}}{x}\in I_0$ by \cref{continuous,funs_type_apply}.
            $0 \leq \apply{\apply{c}{h(k)}}{x}$ by \cref{intervalclosed}.
            $a(k) = \apply{\apply{p}{h(k)}}{x}$ by \cref{subseteq}.
            $\apply{\apply{p}{h(k)}}{x} = \apply{\apply{c}{h(k)}}{x}$ by \cref{subseteq}.
            $a(k) = \apply{\apply{c}{h(k)}}{x}$ by \cref{real_lt_subst_right}.
            $0 \leq a(k)$ by \cref{real_lt_subst_right}.
            $0 < a(k)$ or $0 = a(k)$ by \cref{real_leq_split}.
        \end{subproof}
        $i\in\dom{a}$ by \cref{funs,subseteq}.
        $a(i) \leq g(x)$ by \cref{sum_finite_sequence_term_leq}.
        $a(i) = 1$ by \cref{real_lt_subst_left,real_lt_subst_right}.
        $1 \leq g(x)$ by \cref{real_lt_subst_left}.
    \end{subproof}
    For all $x\in X$ we have $g(x)\neq 0$
        by \cref{continuous,funs_type_apply,real_one_in_naturals,real_naturals_subset_reals,subseteq,real_add_ident,real_zero_lt_one,real_lt_imp_leq,real_leq_antisym,real_lt_irreflexive,real_lt_subst_left}.

    Let $d = \{(x,1) \mid x\in X\}$.
    Let $q = \{(x,d(x) \rmul \inv{(g(x))}) \mid x\in X\}$.
    $q$ is continuous from $X$ to $I_0$ with respect to $T$ and $T_0$
        by \cref{reciprocal_geq_one_unit_interval_continuous}.
    Let $q_1 = j\circ q$.
    $q_1$ is continuous from $X$ to $\reals$ with respect to $T$ and $\standardTopologyR$
        by \cref{continuous_interval_expand_range}.
    For all $x\in X$ we have $q_1(x) = q(x)$
        by \cref{continuous,funs_is_function,funs,function_rightunique,funs_is_relation,funs_subseteq_rels,subseteq,rels_ran_subseteq,id_dom,id_rightunique,id_is_relation,circ_apply,funs_type_apply,id_apply}.
    $1\in\reals$ by \cref{real_one_in_naturals,real_naturals_subset_reals,subseteq}.
    $d$ is a function from $X$ to $\reals$ by \cref{constant_map_function}.
    For all $x\in X$ we have $d(x) = 1$
        by \cref{constant_map_function,funs_is_function,function_apply_intro}.

    Let $P = \{w\in J\times B \mid \text{there exist $\alpha\in J$ such that there exists $u\in B$ such that
        $w = (\alpha,u)$ and
        $u$ is continuous from $X$ to $I_0$ with respect to $T$ and $T_0$ and
        $u$ is continuous from $X$ to $\reals$ with respect to $T$ and $\standardTopologyR$ and
        $\supp{u}{X}{T}\subseteq \apply{U}{\alpha}$ and
        for all $x\in X$ we have $u(x) = \apply{q_1}{x} \rmul \apply{\apply{c}{\alpha}}{x}$}\}$.
    Show that for all $\alpha\in J$ there exists $u$ such that $\alpha\mathrel{P}u$.
    \begin{subproof}
        Fix $\alpha\in J$.
        Let $c_1 = \apply{c}{\alpha}$.
        $c_1$ is continuous from $X$ to $I_0$ with respect to $T$ and $T_0$ and
            $\img{c_1}{\closure{\apply{W}{\alpha}}{X}{T}}\subseteq \{1\}$ and
            $\img{c_1}{X\setminus \apply{V}{\alpha}}\subseteq \{0\}$ by \cref{choice_function,pair_eq_iff}.
        Let $u = \{(x,\apply{q_1}{x} \rmul \apply{c_1}{x}) \mid x\in X\}$.
        $u$ is continuous from $X$ to $\reals$ with respect to $T$ and $\standardTopologyR$
            by \cref{real_function_multiplication_continuous}.
        Let $v = \{(x,\apply{q}{x} \rmul \apply{c_1}{x}) \mid x\in X\}$.
        $v$ is continuous from $X$ to $I_0$ with respect to $T$ and $T_0$
            by \cref{unit_interval_product_continuous}.
        For all $x\in X$ we have $u(x) = v(x)$
            by \cref{pair_eq_iff,continuous,funs_is_function,funs,function_apply_intro,real_mul_eq_subst}.
        For all $x\in X$ we have $u(x)\in I_0$ by \cref{continuous,funs_type_apply}.
        $\img{u}{X}\subseteq I_0$
            by \cref{subseteq,img_iff,continuous,funs_is_function,funs,function_apply_intro}.
        $u\in\funs{X}{\reals}$ by \cref{continuous}.
        $u$ is continuous from $X$ to $I_0$ with respect to $T$ and $T_0$
            by \cref{subseteq,continuous_subspace_range,intervalclosed}.
        Show that for all $x\in X\setminus \apply{V}{\alpha}$ we have $u(x) = 0$.
        \begin{subproof}
            Fix $x\in X\setminus \apply{V}{\alpha}$.
            $c_1\in\funs{X}{I_0}$ and $c_1$ is a function and $\dom{c_1} = X$
                by \cref{continuous,funs_is_function,funs}.
            $x\in\dom{c_1}\inter(X\setminus \apply{V}{\alpha})$ by \cref{setminus_elim_left,subseteq,inter_intro}.
            $c_1(x)\in\img{c_1}{X\setminus \apply{V}{\alpha}}$ by \cref{img_of_function_intro}.
            $c_1(x) = 0$ by \cref{subseteq,singleton_elim}.
            $u(x) = \apply{q_1}{x} \rmul c_1(x)$
                by \cref{setminus_elim_left,pair_eq_iff,continuous,funs_is_function,funs,function_apply_intro}.
            $u(x) = \apply{q_1}{x} \rmul 0$ by \cref{real_lt_subst_right}.
            $\apply{q_1}{x}\in\reals$ by \cref{setminus_elim_left,continuous,funs_type_apply}.
            $u(x) = 0$ by \cref{real_add_ident,real_mul_comm,real_mul_zero}.
        \end{subproof}
        $\supp{u}{X}{T}\subseteq \apply{U}{\alpha}$ by \cref{support_subseteq_from_zero_outside}.
        For all $x\in X$ we have $u(x) = \apply{q_1}{x} \rmul \apply{\apply{c}{\alpha}}{x}$
            by \cref{pair_eq_iff,continuous,funs_is_function,funs,function_apply_intro,real_mul_eq_subst}.
        $\alpha\mathrel{P}u$
            by \cref{continuous,funs_is_relation,funs,funs_subseteq_rels,subseteq,rels_ran_subseteq,rels_elim,powerset_intro,times_tuple_intro,pair_eq_iff}.
    \end{subproof}
    $P$ is a relation by \cref{relation}.
    Take $f$ such that $f$ is a function on $J$ and for all $\alpha\in J$ we have $\alpha\mathrel{P}\apply{f}{\alpha}$
        by \cref{choice_function}.
    Show that for all $\alpha\in J$ we have
        $\apply{f}{\alpha}$ is continuous from $X$ to $I_0$ with respect to $T$ and $T_0$ and
        $\apply{f}{\alpha}$ is continuous from $X$ to $\reals$ with respect to $T$ and $\standardTopologyR$ and
        $\supp{\apply{f}{\alpha}}{X}{T}\subseteq \apply{U}{\alpha}$.
    \begin{subproof}
        Follows by \cref{choice_function,pair_eq_iff}.
    \end{subproof}
    For all $\alpha\in J$ for all $x\in X$ we have
        $\apply{\apply{f}{\alpha}}{x} = \apply{q_1}{x} \rmul \apply{\apply{c}{\alpha}}{x}$
        by \cref{choice_function,pair_eq_iff}.
    Show that for all $\alpha\in J$ for all $x\in X\setminus \apply{V}{\alpha}$ we have $\apply{\apply{f}{\alpha}}{x} = 0$.
    \begin{subproof}
        Fix $\alpha\in J$.
        Fix $x\in X\setminus \apply{V}{\alpha}$.
        $\apply{c}{\alpha}\in\funs{X}{I_0}$ and $\apply{c}{\alpha}$ is a function and $\dom{\apply{c}{\alpha}} = X$
            by \cref{continuous,funs_is_function,funs}.
        $x\in\dom{\apply{c}{\alpha}}\inter(X\setminus \apply{V}{\alpha})$
            by \cref{setminus_elim_left,subseteq,inter_intro}.
        $\apply{\apply{c}{\alpha}}{x}\in\img{\apply{c}{\alpha}}{X\setminus \apply{V}{\alpha}}$
            by \cref{img_of_function_intro}.
        $\apply{\apply{c}{\alpha}}{x} = 0$ by \cref{subseteq,singleton_elim}.
        $x\in X$ by \cref{setminus_elim_left}.
        $\apply{\apply{f}{\alpha}}{x} = \apply{q_1}{x} \rmul \apply{\apply{c}{\alpha}}{x}$ by \cref{subseteq}.
        $\apply{\apply{f}{\alpha}}{x} = \apply{q_1}{x} \rmul 0$ by \cref{real_lt_subst_right}.
        $\apply{q_1}{x}\in\reals$ by \cref{setminus_elim_left,continuous,funs_type_apply}.
        $\apply{\apply{f}{\alpha}}{x} = 0$ by \cref{real_add_ident,real_mul_comm,real_mul_zero}.
    \end{subproof}

    Let $S_1 = \{(\alpha,\supp{\apply{f}{\alpha}}{X}{T}) \mid \alpha\in J\}$.
    $S_1$ is the support family of $f$ on $X$ with respect to $T$ and $J$ by \cref{support_family}.
    $S_1$ is a locally finite indexed family on $J$ in $X$ with respect to $T$
        by \cref{support_family_locally_finite_from_zero_outside}.

    Show that for all $x\in X$ we have $1$ is the pointwise sum of $f$ at $x$ on $X$ with respect to $T$ and $J$.
    \begin{subproof}
        Fix $x\in X$.
        Take $\beta\in J$ such that $x\in\apply{S}{\beta}$ by \cref{subseteq}.
        $S$ is a function on $J$ by \cref{subseteq}.
        $(\beta,\supp{\apply{p}{\beta}}{X}{T})\in S$ by \cref{support_family,pair_eq_iff}.
        $\apply{S}{\beta} = \supp{\apply{p}{\beta}}{X}{T}$ by \cref{function_apply_intro}.
        $x\in \supp{\apply{p}{\beta}}{X}{T}$ by \cref{real_lt_subst_left}.
        $g(x)$ is the pointwise sum of $p$ at $x$ on $X$ with respect to $T$ and $J$
            by \cref{subseteq,locally_finite_real_sum_function}.
        Take $F,n,h,a_0$ such that
            $F\subseteq J$ and
            $F$ is finite and
            $F$ is inhabited and
            $n\in\naturalsPlus$ and
            $h$ is a bijection from $\initialSegment{n}$ to $F$ and
            $a_0$ is a real sum sequence of $p$ at $x$ along $h$ with respect to $n$ and result $g(x)$ and
            for all $\alpha\in J$ if $x\in\supp{\apply{p}{\alpha}}{X}{T}$ then $\alpha\in F$
            by \cref{pointwise_sum_of_indexed_family}.
        $a_0\in\funs{\initialSegment{n}}{\reals}$ and $\sumFiniteSequence{a_0}{g(x)}$
            by \cref{real_sum_sequence_along_enumeration}.
        For all $i\in\initialSegment{n}$ we have $a_0(i) = \apply{\apply{p}{h(i)}}{x}$
            by \cref{real_sum_sequence_along_enumeration}.
        $\beta\in F$ by \cref{subseteq}.
        Let $a = \coordScalarSequence{n}{q_1(x)}{a_0}$.
        $a\in\funs{\initialSegment{n}}{\reals}$ by \cref{continuous,funs_type_apply,coord_scalar_sequence_function}.
        Show that for all $i\in\initialSegment{n}$ we have $a(i) = \apply{\apply{f}{h(i)}}{x}$.
        \begin{subproof}
            Fix $i\in\initialSegment{n}$.
            $i\in\dom{a}$ by \cref{funs,subseteq}.
            $(i,a(i))\in a$ by \cref{funs_elim,function_apply_elim}.
            $a(i) = \apply{q_1}{x} \rmul a_0(i)$ by \cref{coord_scalar_sequence,pair_eq_iff}.
            $a(i) = \apply{q_1}{x} \rmul \apply{\apply{p}{h(i)}}{x}$ by \cref{real_mul_eq_subst,subseteq}.
            $i\in\dom{h}$ by \cref{bijections_dom,subseteq}.
            $h(i)\in F$ by \cref{bijections,surj,funs_type_apply,subseteq}.
            $h(i)\in J$ by \cref{subseteq}.
            $\apply{\apply{f}{h(i)}}{x} = \apply{q_1}{x} \rmul \apply{\apply{c}{h(i)}}{x}$ by \cref{subseteq}.
            $\apply{\apply{p}{h(i)}}{x} = \apply{\apply{c}{h(i)}}{x}$ by \cref{subseteq}.
            $a(i) = \apply{\apply{f}{h(i)}}{x}$ by \cref{real_mul_eq_subst,real_lt_subst_right}.
        \end{subproof}
        $q_1(x)\in\reals$ by \cref{continuous,funs_type_apply}.
        $\sumFiniteSequence{a}{q_1(x) \rmul g(x)}$ by \cref{sum_finite_sequence_scalar}.
        $q_1(x) = q(x)$ by \cref{subseteq}.
        $g(x)\neq 0$ by \cref{subseteq}.
        $g(x)\in\reals$ by \cref{continuous,funs_type_apply}.
        $q$ is a function by \cref{continuous,funs_is_function}.
        $(x,d(x) \rmul \inv{(g(x))})\in q$.
        $q(x) = d(x) \rmul \inv{(g(x))}$ by \cref{function_apply_intro}.
        $q_1(x) = 1 \rmul \inv{(g(x))}$ by \cref{real_lt_subst_left}.
        $q_1(x) = \inv{(g(x))}$ by \cref{real_mul_inverse,real_mul_ident}.
        $q_1(x) \rmul g(x) = 1$
            by \cref{real_mul_eq_subst,real_mul_comm,real_mul_inverse,real_lt_subst_right}.
        $\sumFiniteSequence{a}{1}$ by \cref{real_lt_subst_right}.
        Show that for all $\alpha\in J$ if $x\in\supp{\apply{f}{\alpha}}{X}{T}$ then $\alpha\in F$.
        \begin{subproof}
            Fix $\alpha\in J$.
            Assume $x\in\supp{\apply{f}{\alpha}}{X}{T}$.
            Show that for all $y\in X$ if $\apply{\apply{p}{\alpha}}{y} = 0$ then $\apply{\apply{f}{\alpha}}{y} = 0$.
            \begin{subproof}
                Fix $y\in X$.
                Assume $\apply{\apply{p}{\alpha}}{y} = 0$.
                $\apply{\apply{c}{\alpha}}{y} = 0$ by \cref{real_lt_subst_left}.
                $\apply{\apply{f}{\alpha}}{y} = \apply{q_1}{y} \rmul \apply{\apply{c}{\alpha}}{y}$ by \cref{subseteq}.
                $\apply{\apply{f}{\alpha}}{y} = \apply{q_1}{y} \rmul 0$ by \cref{real_lt_subst_right}.
                $\apply{q_1}{y}\in\reals$ by \cref{continuous,funs_type_apply}.
                $\apply{\apply{f}{\alpha}}{y} = 0$ by \cref{real_add_ident,real_mul_comm,real_mul_zero}.
            \end{subproof}
            $\supp{\apply{f}{\alpha}}{X}{T}\subseteq \supp{\apply{p}{\alpha}}{X}{T}$ by \cref{support_subseteq_from_zero_implication}.
            $x\in \supp{\apply{p}{\alpha}}{X}{T}$ by \cref{subseteq}.
            $\alpha\in F$ by \cref{subseteq}.
        \end{subproof}
        $a$ is a real sum sequence of $f$ at $x$ along $h$ with respect to $n$ and result $1$
            by \cref{real_sum_sequence_along_enumeration}.
        $1$ is the pointwise sum of $f$ at $x$ on $X$ with respect to $T$ and $J$
            by \cref{pointwise_sum_of_indexed_family}.
    \end{subproof}

    $U$ is a function on $J$ by assumption.
    Show that for all $\alpha\in J$ there exists $I$ such that
        $I = \intervalclosed{0}{1}$ and
        $\apply{f}{\alpha}$ is continuous from $X$ to $I$ with respect to $T$ and $\subspaceTopology{\standardTopologyR}{I}$ and
        $\supp{\apply{f}{\alpha}}{X}{T}\subseteq \apply{U}{\alpha}$.
    \begin{subproof}
        Follows by \cref{pair_eq_iff}.
    \end{subproof}
    $f$ is an indexed partition of unity on $X$ dominated by $U$ with respect to $T$ and $J$
        by \cref{indexed_partition_of_unity}.
\end{proof}

% from §41 helper lemma 31: coordinatewise comparison for finite real sums
\begin{lemma}\label{sum_finite_sequence_coordwise_leq}
    Suppose $n\in\naturalsPlus$.
    Let $J = \initialSegment{n}$.
    Suppose $a\in\funs{J}{\reals}$.
    Suppose $b\in\funs{J}{\reals}$.
    Suppose $\sumFiniteSequence{a}{s}$.
    Suppose $\sumFiniteSequence{b}{t}$.
    Suppose for all $i\in J$ we have $a(i) \leq b(i)$.
    Then $s \leq t$.
\end{lemma}
\begin{proof}
    Let $c = \coordDiffSequence{n}{b}{a}$.
    $a\in\realsPower{J}$ and $b\in\realsPower{J}$ by \cref{reals_power,tuple_power}.
    $c\in\funs{J}{\reals}$ by \cref{coord_diff_sequence_function}.
    Show that for all $i\in J$ we have $0 \leq c(i)$.
    \begin{subproof}
        Fix $i\in J$.
        $a(i)\in\reals$ and $b(i)\in\reals$ by \cref{funs_type_apply}.
        $\neg{a(i)}\in\reals$ by \cref{real_add_inverse}.
        $a(i) + \neg{a(i)} \leq b(i) + \neg{a(i)}$ by \cref{real_leq_add}.
        $0 \leq b(i) - a(i)$ by
            \cref{real_add_inverse,real_add_ident,real_sub,real_leq_subst_left_local,real_leq_subst_right_local}.
        $(i,b(i) - a(i))\in c$ by \cref{coord_diff_sequence}.
        $c(i) = b(i) - a(i)$ by \cref{funs_is_function,function_apply_intro}.
        $0 \leq c(i)$ by \cref{real_leq_subst_right_local}.
    \end{subproof}
    Take $u\in\reals$ such that $\sumFiniteSequence{c}{u}$ by \cref{sum_finite_sequence_exists}.
    Let $z = \coordScalarSequence{n}{0}{a}$.
    $z\in\funs{J}{\reals}$ by \cref{real_add_ident,coord_scalar_sequence_function}.
    $z\in\realsPower{J}$ by \cref{reals_power,tuple_power}.
    Show that for all $i\in J$ we have $z(i) = 0$.
    \begin{subproof}
        Fix $i\in J$.
        $a(i)\in\reals$ by \cref{funs_type_apply}.
        $(i,0 \rmul a(i))\in z$ by \cref{coord_scalar_sequence}.
        $z(i) = 0 \rmul a(i)$ by \cref{funs_is_function,function_apply_intro}.
        $z(i) = 0$ by \cref{real_mul_zero,real_lt_subst_right}.
    \end{subproof}
    Let $e_0 = \coordDiffSequence{n}{a}{z}$.
    $e_0\in\funs{J}{\reals}$ by \cref{coord_diff_sequence_function}.
    Show that for all $i\in J$ we have $e_0(i) = a(i)$.
    \begin{subproof}
        Fix $i\in J$.
        $a(i)\in\reals$ by \cref{funs_type_apply}.
        $0\in\reals$ by \cref{real_add_ident}.
        $(i,a(i) - z(i))\in e_0$ by \cref{coord_diff_sequence}.
        $e_0(i) = a(i) - z(i)$ by \cref{funs_is_function,function_apply_intro}.
        $e_0(i) = a(i) - 0$ by \cref{subseteq}.
        We have $a(i) - 0 = a(i) + \neg{0}$ by \cref{real_sub}.
        $a(i) - 0 = a(i) + 0$ by \cref{real_neg_zero}.
        $a(i) - 0 = 0 + a(i)$ by \cref{real_add_comm}.
        $a(i) - 0 = a(i)$ by \cref{real_add_ident}.
        $e_0(i) = a(i)$.
    \end{subproof}
    $e_0 = a$ by \cref{funs_elim,subseteq,fun_subseteq,subseteq_antisymmetric}.
    Let $e_1 = \coordDiffSequence{n}{b}{z}$.
    $e_1\in\funs{J}{\reals}$ by \cref{coord_diff_sequence_function}.
    Show that for all $i\in J$ we have $e_1(i) = b(i)$.
    \begin{subproof}
        Fix $i\in J$.
        $b(i)\in\reals$ by \cref{funs_type_apply}.
        $0\in\reals$ by \cref{real_add_ident}.
        $(i,b(i) - z(i))\in e_1$ by \cref{coord_diff_sequence}.
        $e_1(i) = b(i) - z(i)$ by \cref{funs_is_function,function_apply_intro}.
        $e_1(i) = b(i) - 0$ by \cref{subseteq}.
        We have $b(i) - 0 = b(i) + \neg{0}$ by \cref{real_sub}.
        $b(i) - 0 = b(i) + 0$ by \cref{real_neg_zero}.
        $b(i) - 0 = 0 + b(i)$ by \cref{real_add_comm}.
        $b(i) - 0 = b(i)$ by \cref{real_add_ident}.
        $e_1(i) = b(i)$.
    \end{subproof}
    $e_1 = b$ by \cref{funs_elim,subseteq,fun_subseteq,subseteq_antisymmetric}.
    Let $d = \coordSumSequence{n}{c}{a}$.
    $d\in\funs{J}{\reals}$ by \cref{coord_sum_sequence_function}.
    $d = e_1$ by \cref{coord_diff_sequence_sum,real_lt_subst_left}.
    $d = b$ by \cref{real_lt_subst_right}.
    $\sumFiniteSequence{d}{u + s}$ by \cref{sum_finite_sequence_add}.
    $\sumFiniteSequence{b}{u + s}$ by \cref{real_lt_subst_left}.
    $t = u + s$ by \cref{sum_finite_sequence_unique}.
    For all $i\in\dom{c}$ we have $0 \leq c(i)$ by \cref{funs_elim}.
    $0 \leq u$ by \cref{sum_finite_sequence_nonneg}.
    $s\in\reals$ by \cref{sum_finite_sequence_value_in_reals}.
    $0\in\reals$ by \cref{real_add_ident}.
    $0 + s \leq u + s$ by \cref{real_leq_add}.
    $s \leq t$ by
        \cref{real_add_ident,real_leq_subst_left_local,real_leq_subst_right_local}.
\end{proof}

% from §41 helper lemma 32: deleting a designated term from a finite real sum
\begin{lemma}\label{sum_finite_sequence_remove_value}
    Suppose $n\in\naturalsPlus$.
    Let $J = \initialSegment{n}$.
    Let $J_1 = \initialSegment{n + 1}$.
    Suppose $a\in\funs{J_1}{\reals}$.
    Suppose $\sumFiniteSequence{a}{s}$.
    Suppose $k\in J_1$.
    Then there exist $b,u$ such that
        $b\in\funs{J}{\reals}$ and
        $\sumFiniteSequence{b}{u}$ and
        $s = u + a(k)$ and
        for all $i\in J$ we have
            $(i < k \land b(i) = a(i)) \lor (k \leq i \land b(i) = a(i + 1))$.
\end{lemma}
\begin{proof}
    Let $A = \{m\in\naturalsPlus \mid \text{for all $a_0,s_0,k_0$ such that
        $a_0\in\funs{\initialSegment{m + 1}}{\reals}$ and
        $\sumFiniteSequence{a_0}{s_0}$ and
        $k_0\in\initialSegment{m + 1}$
        there exist $b_0,u_0$ such that
            $b_0\in\funs{\initialSegment{m}}{\reals}$ and
            $\sumFiniteSequence{b_0}{u_0}$ and
            $s_0 = u_0 + a_0(k_0)$ and
            for all $i\in\initialSegment{m}$ we have
                $(i < k_0 \land b_0(i) = a_0(i)) \lor (k_0 \leq i \land b_0(i) = a_0(i + 1))$}\}$.
    Show that $1\in A$.
    \begin{subproof}
        Show that for all $a_0,s_0,k_0$ such that
            $a_0\in\funs{\initialSegment{1 + 1}}{\reals}$ and
            $\sumFiniteSequence{a_0}{s_0}$ and
            $k_0\in\initialSegment{1 + 1}$
            there exist $b_0,u_0$ such that
                $b_0\in\funs{\initialSegment{1}}{\reals}$ and
                $\sumFiniteSequence{b_0}{u_0}$ and
                $s_0 = u_0 + a_0(k_0)$ and
                for all $i\in\initialSegment{1}$ we have
                    $(i < k_0 \land b_0(i) = a_0(i)) \lor (k_0 \leq i \land b_0(i) = a_0(i + 1))$.
        \begin{subproof}
            Fix $a_0$.
            Fix $s_0$.
            Fix $k_0$ such that
                $a_0\in\funs{\initialSegment{1 + 1}}{\reals}$ and
                $\sumFiniteSequence{a_0}{s_0}$ and
                $k_0\in\initialSegment{1 + 1}$.
            $1\in\naturalsPlus$ by \cref{real_one_in_naturals}.
            Let $c_0 = \restrl{a_0}{\initialSegment{1}}$.
            Show that $c_0\in\funs{\initialSegment{1}}{\reals}$ and
                there exists $t_0\in\reals$ such that $\sumFiniteSequence{c_0}{t_0}$ and $s_0 = t_0 + a_0(1 + 1)$.
            \begin{subproof}
                Follows by \cref{sum_finite_sequence_split_last}.
            \end{subproof}
            Take $t_0\in\reals$ such that $\sumFiniteSequence{c_0}{t_0}$ and $s_0 = t_0 + a_0(1 + 1)$.
            \begin{byCase}
            \caseOf{$k_0 = 1 + 1$.}
                Show that for all $i\in\initialSegment{1}$ we have
                    $(i < k_0 \land c_0(i) = a_0(i)) \lor (k_0 \leq i \land c_0(i) = a_0(i + 1))$.
                \begin{subproof}
                    Fix $i\in\initialSegment{1}$.
                    $i\in\dom{c_0}$ by \cref{funs,subseteq}.
                    $\initialSegment{1}\subseteq \dom{a_0}$ by \cref{funs,initial_segment_subset_succ,subseteq}.
                    $c_0(i) = a_0(i)$ by \cref{restrl_of_function_apply,funs_is_function,subseteq}.
                    $i\in\naturalsPlus$ by \cref{initial_segment_naturalsplus}.
                    $i\in\reals$ by \cref{real_naturals_subset_reals,subseteq}.
                    $i = 1$ by \cref{initial_segment_naturalsplus,real_one_in_naturals,real_one_leq_naturalsplus,real_lt_trans,real_lt_irreflexive,real_naturals_subset_reals,subseteq}.
                    $i < 1 + 1$ by \cref{real_naturals_lt_succ_local,real_lt_subst_left}.
                    $(i < k_0 \land c_0(i) = a_0(i)) \lor (k_0 \leq i \land c_0(i) = a_0(i + 1))$
                        by \cref{real_lt_subst_right}.
                \end{subproof}
                Show that there exist $b_0,u_0$ such that
                    $b_0\in\funs{\initialSegment{1}}{\reals}$ and
                    $\sumFiniteSequence{b_0}{u_0}$ and
                    $s_0 = u_0 + a_0(k_0)$ and
                    for all $i\in\initialSegment{1}$ we have
                        $(i < k_0 \land b_0(i) = a_0(i)) \lor (k_0 \leq i \land b_0(i) = a_0(i + 1))$.
                \begin{subproof}
                    Follows by \cref{real_lt_subst_right}.
                \end{subproof}
            \caseOf{$k_0 \neq 1 + 1$.}
                $k_0\in\initialSegment{1}$
                    by \cref{initial_segment_naturalsplus,real_naturals_leq_succ_elim,real_lt_subst_right}.
                $k_0 = 1$
                    by \cref{initial_segment_naturalsplus,real_one_in_naturals,real_one_leq_naturalsplus,real_lt_trans,real_lt_irreflexive,real_naturals_subset_reals,subseteq}.
                Let $b_0 = \{(1,a_0(1 + 1))\}$.
                $b_0$ is a function by \cref{singleton_elim,pair_eq_iff,rightunique,relation}.
                $\dom{b_0} = \initialSegment{1}$
                    by \cref{dom_iff,singleton_elim,pair_eq_iff,singleton_intro,dom_intro,subseteq,subseteq_antisymmetric,real_one_in_naturals,real_one_leq_naturalsplus,initial_segment_naturalsplus,real_naturals_subset_reals,real_leq_antisym}.
                $a_0(1 + 1)\in\reals$ by \cref{funs_type_apply,real_naturals_closed_succ,initial_segment_naturalsplus}.
                For all $i\in\dom{b_0}$ we have $b_0(i)\in\reals$
                    by \cref{dom_iff,singleton_elim,pair_eq_iff,function_apply_intro,real_lt_subst_left}.
                $b_0\in\funs{\initialSegment{1}}{\reals}$ by \cref{funs_intro}.
                $\sumFiniteSequence{b_0}{b_0(1)}$ by \cref{sum_finite_sequence_one_term}.
                $b_0(1) = a_0(1 + 1)$ by \cref{function_apply_intro,singleton_intro}.
                $\sumFiniteSequence{b_0}{a_0(1 + 1)}$ by \cref{real_lt_subst_right}.
                $c_0$ is a function and $\dom{c_0} = \initialSegment{1}$ by \cref{funs_is_function,funs}.
                $\initialSegment{1}\subseteq \dom{a_0}$ by \cref{funs,initial_segment_subset_succ,subseteq}.
                $c_0(1) = a_0(1)$
                    by \cref{restrl_of_function_apply,funs_is_function,real_one_in_naturals,real_one_leq_naturalsplus,initial_segment_naturalsplus,subseteq}.
                $\sumFiniteSequence{c_0}{c_0(1)}$ by \cref{sum_finite_sequence_one_term}.
                $t_0 = c_0(1)$ by \cref{sum_finite_sequence_unique}.
                $t_0 = a_0(1)$ by \cref{real_lt_subst_right}.
                $s_0 = a_0(1) + a_0(1 + 1)$ by \cref{real_lt_subst_right}.
                $s_0 = a_0(1 + 1) + a_0(1)$ by \cref{real_add_comm,real_lt_subst_right}.
                $s_0 = a_0(1 + 1) + a_0(k_0)$ by \cref{real_lt_subst_right}.
                Show that for all $i\in\initialSegment{1}$ we have
                    $(i < k_0 \land b_0(i) = a_0(i)) \lor (k_0 \leq i \land b_0(i) = a_0(i + 1))$.
                \begin{subproof}
                    Fix $i\in\initialSegment{1}$.
                    $i = 1$
                        by \cref{initial_segment_naturalsplus,real_one_in_naturals,real_one_leq_naturalsplus,real_lt_trans,real_lt_irreflexive,real_naturals_subset_reals,subseteq}.
                    $k_0 \leq i$ by \cref{real_lt_subst_left,real_lt_subst_right}.
                    $b_0(i) = b_0(1)$ by \cref{real_lt_subst_left}.
                    $b_0(i) = a_0(1 + 1)$ by \cref{real_lt_subst_right}.
                    $i + 1 = 1 + 1$ by \cref{real_lt_subst_left}.
                    $(i < k_0 \land b_0(i) = a_0(i)) \lor (k_0 \leq i \land b_0(i) = a_0(i + 1))$
                        by \cref{real_lt_subst_right}.
                \end{subproof}
                Show that there exist $b_1,u_0$ such that
                    $b_1\in\funs{\initialSegment{1}}{\reals}$ and
                    $\sumFiniteSequence{b_1}{u_0}$ and
                    $s_0 = u_0 + a_0(k_0)$ and
                    for all $i\in\initialSegment{1}$ we have
                        $(i < k_0 \land b_1(i) = a_0(i)) \lor (k_0 \leq i \land b_1(i) = a_0(i + 1))$.
                \begin{subproof}
                    Follows by \cref{real_lt_subst_left,real_lt_subst_right}.
                \end{subproof}
            \end{byCase}
        \end{subproof}
        $1\in A$ by \cref{real_one_in_naturals}.
    \end{subproof}
    Show that for all $m\in A$ we have $m + 1\in A$.
    \begin{subproof}
        Fix $m\in A$.
        $m\in\naturalsPlus$ by \cref{subseteq}.
        $m + 1\in\naturalsPlus$ by \cref{real_naturals_closed_succ}.
        Show that for all $a_0,s_0,k_0$ such that
            $a_0\in\funs{\initialSegment{(m + 1) + 1}}{\reals}$ and
            $\sumFiniteSequence{a_0}{s_0}$ and
            $k_0\in\initialSegment{(m + 1) + 1}$
            there exist $b_0,u_0$ such that
                $b_0\in\funs{\initialSegment{m + 1}}{\reals}$ and
                $\sumFiniteSequence{b_0}{u_0}$ and
                $s_0 = u_0 + a_0(k_0)$ and
                for all $i\in\initialSegment{m + 1}$ we have
                    $(i < k_0 \land b_0(i) = a_0(i)) \lor (k_0 \leq i \land b_0(i) = a_0(i + 1))$.
        \begin{subproof}
            Fix $a_0$.
            Fix $s_0$.
            Fix $k_0$ such that
                $a_0\in\funs{\initialSegment{(m + 1) + 1}}{\reals}$ and
                $\sumFiniteSequence{a_0}{s_0}$ and
                $k_0\in\initialSegment{(m + 1) + 1}$.
            Let $a_1 = \restrl{a_0}{\initialSegment{m + 1}}$.
            Show that $a_1\in\funs{\initialSegment{m + 1}}{\reals}$ and
                there exists $t_0\in\reals$ such that $\sumFiniteSequence{a_1}{t_0}$ and $s_0 = t_0 + a_0((m + 1) + 1)$.
            \begin{subproof}
                Follows by \cref{sum_finite_sequence_split_last}.
            \end{subproof}
            Take $t_0\in\reals$ such that $\sumFiniteSequence{a_1}{t_0}$ and $s_0 = t_0 + a_0((m + 1) + 1)$.
            \begin{byCase}
            \caseOf{$k_0 = (m + 1) + 1$.}
                Show that for all $i\in\initialSegment{m + 1}$ we have
                    $(i < k_0 \land a_1(i) = a_0(i)) \lor (k_0 \leq i \land a_1(i) = a_0(i + 1))$.
                \begin{subproof}
                    Fix $i\in\initialSegment{m + 1}$.
                    $i\in\dom{a_1}$ by \cref{funs,subseteq}.
                    $\initialSegment{m + 1}\subseteq \dom{a_0}$ by \cref{funs,initial_segment_subset_succ,subseteq}.
                    $a_1(i) = a_0(i)$ by \cref{restrl_of_function_apply,funs_is_function,subseteq}.
                    $i\in\naturalsPlus$ by \cref{initial_segment_naturalsplus}.
                    Show that $i < (m + 1) + 1$.
                    \begin{subproof}
                        We have $i < m + 1$ or $i = m + 1$ by \cref{initial_segment_naturalsplus}.
                        \begin{byCase}
                        \caseOf{$i < m + 1$.}
                            $i\in\reals$ and $m + 1\in\reals$ and $(m + 1) + 1\in\reals$
                                by \cref{real_naturals_subset_reals,subseteq,real_naturals_closed_succ,real_one_in_naturals,real_add_closed}.
                            $m + 1 < (m + 1) + 1$ by \cref{real_naturals_lt_succ_local,real_naturals_closed_succ}.
                            $i < (m + 1) + 1$ by \cref{real_lt_trans}.
                        \caseOf{$i = m + 1$.}
                            $i < (m + 1) + 1$ by \cref{real_naturals_lt_succ_local,real_naturals_closed_succ,real_lt_subst_left}.
                        \end{byCase}
                    \end{subproof}
                    $(i < k_0 \land a_1(i) = a_0(i)) \lor (k_0 \leq i \land a_1(i) = a_0(i + 1))$
                        by \cref{real_lt_subst_right}.
                \end{subproof}
                Let $b_1 = a_1$.
                Let $u_1 = t_0$.
                $b_1\in\funs{\initialSegment{m + 1}}{\reals}$ by \cref{real_lt_subst_left}.
                $\sumFiniteSequence{b_1}{u_1}$ by \cref{real_lt_subst_left,real_lt_subst_right}.
                $s_0 = u_1 + a_0(k_0)$ by \cref{real_lt_subst_left,real_lt_subst_right}.
                Show that for all $i\in\initialSegment{m + 1}$ we have
                    $(i < k_0 \land b_1(i) = a_0(i)) \lor (k_0 \leq i \land b_1(i) = a_0(i + 1))$.
                \begin{subproof}
                    Fix $i\in\initialSegment{m + 1}$.
                    We have $(i < k_0 \land a_1(i) = a_0(i)) \lor (k_0 \leq i \land a_1(i) = a_0(i + 1))$
                        by assumption.
                    $(i < k_0 \land b_1(i) = a_0(i)) \lor (k_0 \leq i \land b_1(i) = a_0(i + 1))$
                        by \cref{real_lt_subst_left}.
                \end{subproof}
                Follows by definition.
            \caseOf{$k_0 \neq (m + 1) + 1$.}
                $k_0\in\initialSegment{m + 1}$
                    by \cref{initial_segment_naturalsplus,real_naturals_leq_succ_elim,real_lt_subst_right}.
                Take $c_0,v_0$ such that
                    $c_0\in\funs{\initialSegment{m}}{\reals}$ and
                    $\sumFiniteSequence{c_0}{v_0}$ and
                    $t_0 = v_0 + a_1(k_0)$ and
                    for all $i\in\initialSegment{m}$ we have
                        $(i < k_0 \land c_0(i) = a_1(i)) \lor (k_0 \leq i \land c_0(i) = a_1(i + 1))$
                    by \cref{subseteq}.
                Let $d_0 = \{((m + 1),a_0((m + 1) + 1))\}$.
                $d_0$ is a function by \cref{singleton_elim,pair_eq_iff,rightunique,relation}.
                Show that for all $u$ such that $u\in\dom{c_0}\inter\dom{d_0}$ we have $c_0(u) = d_0(u)$.
                \begin{subproof}
                    Fix $u$ such that $u\in\dom{c_0}\inter\dom{d_0}$.
                    Show that $c_0(u) = d_0(u)$.
                    \begin{subproof}
                        Suppose not.
                        $u\in\dom{c_0}$ by \cref{inter_elim_left}.
                        $u\in\initialSegment{m}$ by \cref{funs,subseteq}.
                        $u\in\dom{d_0}$ by \cref{inter_elim_right}.
                        $u = m + 1$ by \cref{dom_iff,singleton_elim,pair_eq_iff}.
                        $u\in\naturalsPlus$ and $m + 1\in\naturalsPlus$ by \cref{initial_segment_naturalsplus,real_naturals_closed_succ,subseteq}.
                        $u\in\reals$ and $m\in\reals$ and $1\in\reals$ by \cref{real_naturals_subset_reals,subseteq,real_one_in_naturals}.
                        Show that $u < m + 1$.
                        \begin{subproof}
                            We have $u < m$ or $u = m$ by \cref{initial_segment_naturalsplus}.
                            \begin{byCase}
                            \caseOf{$u < m$.}
                                $m < m + 1$ by \cref{real_naturals_lt_succ_local}.
                                $u < m + 1$ by \cref{real_lt_trans}.
                            \caseOf{$u = m$.}
                                $u < m + 1$ by \cref{real_naturals_lt_succ_local,real_lt_subst_left}.
                            \end{byCase}
                        \end{subproof}
                        $u < u$ by \cref{real_lt_subst_right}.
                        Contradiction by \cref{real_lt_irreflexive}.
                    \end{subproof}
                \end{subproof}
                $c_0$ is right-unique and $c_0$ is a relation by \cref{funs_is_function,function_rightunique,funs_is_relation}.
                $d_0$ is right-unique and $d_0$ is a relation by \cref{rightunique,relation}.
                Let $b_0 = c_0\union d_0$.
                $b_0$ is a function by \cref{union_compatible_functions,relation,rightunique}.
                $\dom{b_0} = \dom{c_0}\union\dom{d_0}$ by \cref{dom_union}.
                $\dom{b_0} = \initialSegment{m}\union\dom{d_0}$ by \cref{funs}.
                $\dom{d_0} = \{m + 1\}$ by \cref{dom_iff,singleton_elim,pair_eq_iff,singleton_intro,dom_intro,subseteq,subseteq_antisymmetric}.
                $\dom{b_0} = \initialSegment{m}\union \{m + 1\}$.
                For all $u$ such that $u\in\initialSegment{m}\union\{m + 1\}$ we have $u\in\initialSegment{m + 1}$
                    by \cref{union_iff,initial_segment_naturalsplus,real_naturals_subset_reals,subseteq,real_one_in_naturals,real_add_ident,real_add_closed,real_zero_lt_one,real_lt_add,real_add_comm,real_lt_subst_left,real_lt_trans,singleton_elim}.
                For all $u$ such that $u\in\initialSegment{m + 1}$ we have $u\in\initialSegment{m}\union\{m + 1\}$
                    by \cref{initial_segment_naturalsplus,real_naturals_leq_succ_elim,union_iff,singleton_intro}.
                $\initialSegment{m}\union\{m + 1\} = \initialSegment{m + 1}$ by \cref{subseteq,subseteq_antisymmetric}.
                $\dom{b_0} = \initialSegment{m + 1}$.
                Show that for all $u\in\dom{b_0}$ we have $b_0(u)\in\reals$.
                \begin{subproof}
                    Fix $u\in\dom{b_0}$.
                    $u\in\dom{c_0}\union\dom{d_0}$ by \cref{dom_union,real_lt_subst_left}.
                    \begin{byCase}
                    \caseOf{$u\in\dom{c_0}$.}
                        $u\in\initialSegment{m}$ by \cref{funs,subseteq}.
                        $b_0(u) = c_0(u)$ by \cref{union_iff,restrl_iff,dom_iff,function_apply_intro,funs,subseteq,real_lt_subst_left}.
                        $b_0(u)\in\reals$ by \cref{funs_type_apply,real_lt_subst_right}.
                    \caseOf{$u\in\dom{d_0}$.}
                        $u = m + 1$ by \cref{dom_iff,singleton_elim,pair_eq_iff}.
                        $b_0(u) = b_0(m + 1)$ by \cref{real_lt_subst_left}.
                        $b_0(m + 1) = a_0((m + 1) + 1)$ by \cref{function_apply_intro,singleton_intro,union_iff}.
                        $b_0(u)\in\reals$ by \cref{real_lt_subst_left,funs_type_apply,real_naturals_closed_succ,initial_segment_naturalsplus}.
                    \end{byCase}
                \end{subproof}
                $b_0\in\funs{\initialSegment{m + 1}}{\reals}$ by \cref{funs_intro}.
                Show that $\restrl{b_0}{\initialSegment{m}} = c_0$.
                \begin{subproof}
                    Show that for all $w$ such that $w\in\restrl{b_0}{\initialSegment{m}}$ we have $w\in c_0$.
                    \begin{subproof}
                        Fix $w\in\restrl{b_0}{\initialSegment{m}}$.
                        $w\in b_0$ by \cref{restrl_subseteq,subseteq}.
                        Take $u,v$ such that $w = (u,v)$ by \cref{relation}.
                        $(u,v)\in\restrl{b_0}{\initialSegment{m}}$ and $(u,v)\in b_0$ by \cref{pair_eq_iff}.
                        $u\in\dom{\restrl{b_0}{\initialSegment{m}}}$ by \cref{dom_intro}.
                        $u\in\initialSegment{m}$ by \cref{restrl_dom,subseteq}.
                        $u\in\dom{c_0}$ by \cref{funs}.
                        $(u,c_0(u))\in c_0$ by \cref{function_apply_iff,funs_is_function}.
                        $(u,c_0(u))\in b_0$ by \cref{union_iff,real_lt_subst_right}.
                        $b_0(u) = v$ and $b_0(u) = c_0(u)$ by \cref{function_apply_intro}.
                        $v = c_0(u)$ by \cref{real_lt_subst_left}.
                        $w = (u,c_0(u))$ by \cref{pair_eq_iff}.
                        $w\in c_0$ by \cref{pair_eq_iff}.
                    \end{subproof}
                    $\restrl{b_0}{\initialSegment{m}}\subseteq c_0$ by \cref{subseteq}.
                    For all $w$ such that $w\in c_0$ we have $w\in\restrl{b_0}{\initialSegment{m}}$
                        by \cref{union_iff,restrl_iff,dom_iff,funs,subseteq,dom_intro,relation,pair_eq_iff}.
                    $c_0\subseteq \restrl{b_0}{\initialSegment{m}}$ by \cref{subseteq}.
                    $\restrl{b_0}{\initialSegment{m}} = c_0$ by \cref{subseteq_antisymmetric}.
                \end{subproof}
                $\sumFiniteSequence{b_0}{v_0 + b_0(m + 1)}$ by \cref{sum_finite_sequence_append_last}.
                $b_0(m + 1) = a_0((m + 1) + 1)$ by \cref{function_apply_intro,singleton_intro,union_iff}.
                $\sumFiniteSequence{b_0}{v_0 + a_0((m + 1) + 1)}$ by \cref{real_lt_subst_right}.
                $\initialSegment{m + 1}\subseteq \dom{a_0}$ by \cref{funs,initial_segment_subset_succ,subseteq}.
                $a_1(k_0) = a_0(k_0)$ by \cref{restrl_of_function_apply,funs_is_function,subseteq}.
                $v_0\in\reals$ by \cref{sum_finite_sequence_value_in_reals}.
                $a_0(k_0)\in\reals$ by \cref{funs_type_apply,subseteq}.
                $a_0((m + 1) + 1)\in\reals$ by \cref{funs_type_apply,real_naturals_closed_succ,initial_segment_naturalsplus}.
                $s_0 = (v_0 + a_0(k_0)) + a_0((m + 1) + 1)$ by \cref{real_lt_subst_left,real_lt_subst_right}.
                $s_0 = v_0 + (a_0(k_0) + a_0((m + 1) + 1))$ by \cref{real_add_assoc,real_lt_subst_right}.
                $s_0 = v_0 + (a_0((m + 1) + 1) + a_0(k_0))$ by \cref{real_add_comm,real_lt_subst_right}.
                $s_0 = (v_0 + a_0((m + 1) + 1)) + a_0(k_0)$ by \cref{real_add_assoc,real_lt_subst_right}.
                Show that for all $i\in\initialSegment{m + 1}$ we have
                    $(i < k_0 \land b_0(i) = a_0(i)) \lor (k_0 \leq i \land b_0(i) = a_0(i + 1))$.
                \begin{subproof}
                    Fix $i\in\initialSegment{m + 1}$.
                    \begin{byCase}
                    \caseOf{$i = m + 1$.}
                        $k_0 \leq i$ by \cref{initial_segment_naturalsplus,real_lt_imp_leq,subseteq}.
                        $b_0(i) = b_0(m + 1)$ by \cref{real_lt_subst_left}.
                        $b_0(i) = a_0((m + 1) + 1)$ by \cref{real_lt_subst_right}.
                        $i + 1 = (m + 1) + 1$ by \cref{real_lt_subst_left}.
                        $(i < k_0 \land b_0(i) = a_0(i)) \lor (k_0 \leq i \land b_0(i) = a_0(i + 1))$
                            by \cref{real_lt_subst_right}.
                    \caseOf{$i \neq m + 1$.}
                        $i\in\initialSegment{m}$ by \cref{initial_segment_naturalsplus,real_naturals_leq_succ_elim,real_lt_subst_right}.
                        We have $(i < k_0 \land c_0(i) = a_1(i)) \lor (k_0 \leq i \land c_0(i) = a_1(i + 1))$ by assumption.
                        $b_0(i) = c_0(i)$ by \cref{union_iff,restrl_iff,dom_iff,function_apply_intro,funs,subseteq,real_lt_subst_left}.
                        $a_1(i) = a_0(i)$ by \cref{restrl_of_function_apply,funs_is_function,subseteq}.
                        $a_1(i + 1) = a_0(i + 1)$ by
                            \cref{restrl_of_function_apply,funs_is_function,initial_segment_subset_succ,subseteq,initial_segment_naturalsplus,real_naturals_closed_succ,real_lt_imp_leq,real_leq_add,real_one_in_naturals,real_naturals_subset_reals,subseteq,real_add_closed}.
                        $(i < k_0 \land b_0(i) = a_0(i)) \lor (k_0 \leq i \land b_0(i) = a_0(i + 1))$
                            by \cref{real_lt_subst_left,real_lt_subst_right}.
                    \end{byCase}
                \end{subproof}
                Let $b_1 = b_0$.
                Let $u_1 = v_0 + a_0((m + 1) + 1)$.
                $b_1\in\funs{\initialSegment{m + 1}}{\reals}$ by \cref{real_lt_subst_left}.
                $\sumFiniteSequence{b_1}{u_1}$ by \cref{real_lt_subst_left,real_lt_subst_right}.
                $s_0 = u_1 + a_0(k_0)$ by \cref{real_lt_subst_left,real_lt_subst_right}.
                For all $i\in\initialSegment{m + 1}$ we have
                    $(i < k_0 \land b_1(i) = a_0(i)) \lor (k_0 \leq i \land b_1(i) = a_0(i + 1))$
                    by \cref{real_lt_subst_left}.
                Follows by definition.
            \end{byCase}
        \end{subproof}
        $m + 1\in A$ by \cref{real_naturals_closed_succ}.
    \end{subproof}
    $A\subseteq\naturalsPlus$ by \cref{subseteq}.
    $n\in A$ by \cref{subseteq,real_naturals_induction}.
    Show that for all $a_0,s_0,k_0$ such that
        $a_0\in\funs{\initialSegment{n + 1}}{\reals}$ and
        $\sumFiniteSequence{a_0}{s_0}$ and
        $k_0\in\initialSegment{n + 1}$
        there exist $b_0,u_0$ such that
            $b_0\in\funs{\initialSegment{n}}{\reals}$ and
            $\sumFiniteSequence{b_0}{u_0}$ and
            $s_0 = u_0 + a_0(k_0)$ and
            for all $i\in\initialSegment{n}$ we have
                $(i < k_0 \land b_0(i) = a_0(i)) \lor (k_0 \leq i \land b_0(i) = a_0(i + 1))$.
    \begin{subproof}
        Follows by \cref{subseteq}.
    \end{subproof}
    Take $b_0,u_0$ such that
        $b_0\in\funs{\initialSegment{n}}{\reals}$ and
        $\sumFiniteSequence{b_0}{u_0}$ and
        $s = u_0 + a(k)$ and
        for all $i\in\initialSegment{n}$ we have
            $(i < k \land b_0(i) = a(i)) \lor (k \leq i \land b_0(i) = a(i + 1))$
        by \cref{subseteq,real_lt_subst_left,real_lt_subst_right}.
    $b_0\in\funs{J}{\reals}$ by \cref{real_lt_subst_left}.
    For all $i\in J$ we have
        $(i < k \land b_0(i) = a(i)) \lor (k \leq i \land b_0(i) = a(i + 1))$
        by \cref{real_lt_subst_left}.
    Let $b = b_0$.
    Let $u = u_0$.
    $b\in\funs{J}{\reals}$ by \cref{real_lt_subst_left}.
    $\sumFiniteSequence{b}{u}$ by \cref{real_lt_subst_left,real_lt_subst_right}.
    $s = u + a(k)$ by \cref{real_lt_subst_left,real_lt_subst_right}.
    For all $i\in J$ we have
        $(i < k \land b(i) = a(i)) \lor (k \leq i \land b(i) = a(i + 1))$
        by \cref{real_lt_subst_left,real_lt_subst_right}.
    Therefore there exist $b_1,u_1$ such that
        $b_1\in\funs{J}{\reals}$ and
        $\sumFiniteSequence{b_1}{u_1}$ and
        $s = u_1 + a(k)$ and
        for all $i\in J$ we have
            $(i < k \land b_1(i) = a(i)) \lor (k \leq i \land b_1(i) = a(i + 1))$
        by \cref{real_lt_subst_left}.
\end{proof}

% from §41 helper lemma 33: comparing sums along two well-ordered enumerations of the same finite set
\begin{lemma}\label{well_ordered_same_set_sum_comparison}
    Assume $R$ is a well order relation on $J$.
    Suppose $n\in\naturalsPlus$.
    Suppose $g$ is a bijection from $\initialSegment{n}$ to $F$.
    Suppose $h$ is a well-ordered enumeration of $F$ in $J$ with respect to $R$ and $m$.
    Suppose $a$ is a real sum sequence of $f$ at $x$ along $g$ with respect to $n$ and result $s$.
    Suppose $b$ is a real sum sequence of $f$ at $x$ along $h$ with respect to $m$ and result $t$.
    Then $s = t$.
\end{lemma}
\begin{proof}
    Let $A = \{m\in\naturalsPlus \mid \text{for all $F_0,g_0,n_0,a_0,s_0,h_0,b_0,t_0$ such that
        $n_0\in\naturalsPlus$ and
        $g_0$ is a bijection from $\initialSegment{n_0}$ to $F_0$ and
        $h_0$ is a well-ordered enumeration of $F_0$ in $J$ with respect to $R$ and $m$ and
        $a_0$ is a real sum sequence of $f$ at $x$ along $g_0$ with respect to $n_0$ and result $s_0$ and
        $b_0$ is a real sum sequence of $f$ at $x$ along $h_0$ with respect to $m$ and result $t_0$
        we have $s_0 = t_0$}\}$.
    Show that $1\in A$.
    \begin{subproof}
        Show that for all $F_0,g_0,n_0,a_0,s_0,h_0,b_0,t_0$ such that
            $n_0\in\naturalsPlus$ and
            $g_0$ is a bijection from $\initialSegment{n_0}$ to $F_0$ and
            $h_0$ is a well-ordered enumeration of $F_0$ in $J$ with respect to $R$ and $1$ and
            $a_0$ is a real sum sequence of $f$ at $x$ along $g_0$ with respect to $n_0$ and result $s_0$ and
            $b_0$ is a real sum sequence of $f$ at $x$ along $h_0$ with respect to $1$ and result $t_0$
            we have $s_0 = t_0$.
        \begin{subproof}
            Fix $F_0$.
            Fix $g_0$.
            Fix $n_0$.
            Fix $a_0$.
            Fix $s_0$.
            Fix $h_0$.
            Fix $b_0$.
            Fix $t_0$ such that
                $n_0\in\naturalsPlus$ and
                $g_0$ is a bijection from $\initialSegment{n_0}$ to $F_0$ and
                $h_0$ is a well-ordered enumeration of $F_0$ in $J$ with respect to $R$ and $1$ and
                $a_0$ is a real sum sequence of $f$ at $x$ along $g_0$ with respect to $n_0$ and result $s_0$ and
                $b_0$ is a real sum sequence of $f$ at $x$ along $h_0$ with respect to $1$ and result $t_0$.
            $F_0\subseteq J$ and $h_0$ is a bijection from $\initialSegment{1}$ to $F_0$
                by \cref{well_ordered_enumeration_of_subset}.
            Let $\beta = h_0(1)$.
            $\beta\in F_0$
                by \cref{real_one_in_naturals,real_one_leq_naturalsplus,initial_segment_naturalsplus,bijections,surj,funs_type_apply,subseteq}.
            Show that for all $y\in F_0$ we have $y = \beta$.
            \begin{subproof}
                Fix $y\in F_0$.
                Take $u\in\initialSegment{1}$ such that $h_0(u) = y$ by \cref{bijections,surj}.
                $u = 1$ by \cref{initial_segment_naturalsplus,real_one_in_naturals,real_one_leq_naturalsplus,real_naturals_subset_reals,subseteq,real_lt_trans,real_lt_irreflexive}.
                $y = \beta$ by \cref{real_lt_subst_left}.
            \end{subproof}
            $F_0\subseteq \{\beta\}$ by \cref{subseteq,singleton_intro}.
            $\{\beta\}\subseteq F_0$ by \cref{singleton_subset_intro}.
            $F_0 = \{\beta\}$ by \cref{subseteq_antisymmetric}.
            $g_0$ is a bijection from $\initialSegment{n_0}$ to $\{\beta\}$
                by \cref{real_lt_subst_right}.
            $a_0\in\funs{\initialSegment{n_0}}{\reals}$ and $\sumFiniteSequence{a_0}{s_0}$ and
                for all $i\in\initialSegment{n_0}$ we have $a_0(i) = \apply{\apply{f}{g_0(i)}}{x}$
                by \cref{real_sum_sequence_along_enumeration}.
            Show that for all $i\in\initialSegment{n_0}$ we have $a_0(i) = \apply{\apply{f}{\beta}}{x}$.
            \begin{subproof}
                Fix $i\in\initialSegment{n_0}$.
                $i\in\dom{g_0}$ by \cref{bijections_dom,subseteq}.
                $g_0(i)\in F_0$ by \cref{bijections,surj,funs_type_apply,subseteq}.
                $g_0(i) = \beta$ by \cref{real_lt_subst_left,singleton_elim}.
                $a_0(i) = \apply{\apply{f}{\beta}}{x}$ by \cref{real_lt_subst_left}.
            \end{subproof}
            $s_0 = \apply{\apply{f}{\beta}}{x}$
                by \cref{singleton_enumeration_sum_value,real_lt_subst_right}.
            $h_0$ is a bijection from $\initialSegment{1}$ to $\{\beta\}$
                by \cref{well_ordered_enumeration_of_subset,real_lt_subst_right}.
            $b_0\in\funs{\initialSegment{1}}{\reals}$ and $\sumFiniteSequence{b_0}{t_0}$ and
                for all $i\in\initialSegment{1}$ we have $b_0(i) = \apply{\apply{f}{h_0(i)}}{x}$
                by \cref{real_sum_sequence_along_enumeration}.
            Show that for all $i\in\initialSegment{1}$ we have $b_0(i) = \apply{\apply{f}{\beta}}{x}$.
            \begin{subproof}
                Fix $i\in\initialSegment{1}$.
                $i = 1$ by \cref{initial_segment_naturalsplus,real_one_in_naturals,real_one_leq_naturalsplus,real_naturals_subset_reals,subseteq,real_lt_trans,real_lt_irreflexive}.
                $b_0(i) = \apply{\apply{f}{\beta}}{x}$ by \cref{real_lt_subst_left}.
            \end{subproof}
            $t_0 = \apply{\apply{f}{\beta}}{x}$
                by \cref{real_one_in_naturals,singleton_enumeration_sum_value,real_lt_subst_right}.
            $s_0 = t_0$ by \cref{real_lt_subst_left}.
        \end{subproof}
        $1\in A$ by \cref{real_one_in_naturals}.
    \end{subproof}
    Show that for all $k\in A$ we have $k + 1\in A$.
    \begin{subproof}
        Fix $k\in A$.
        Show that for all $F_1,g_1,n_1,a_1,s_1,h_2,b_1,t_1$ such that
            $n_1\in\naturalsPlus$ and
            $g_1$ is a bijection from $\initialSegment{n_1}$ to $F_1$ and
            $h_2$ is a well-ordered enumeration of $F_1$ in $J$ with respect to $R$ and $k + 1$ and
            $a_1$ is a real sum sequence of $f$ at $x$ along $g_1$ with respect to $n_1$ and result $s_1$ and
            $b_1$ is a real sum sequence of $f$ at $x$ along $h_2$ with respect to $k + 1$ and result $t_1$
            we have $s_1 = t_1$.
        \begin{subproof}
            Fix $F_1$.
            Fix $g_1$.
            Fix $n_1$.
            Fix $a_1$.
            Fix $s_1$.
            Fix $h_2$.
            Fix $b_1$.
            Fix $t_1$ such that
                $n_1\in\naturalsPlus$ and
                $g_1$ is a bijection from $\initialSegment{n_1}$ to $F_1$ and
                $h_2$ is a well-ordered enumeration of $F_1$ in $J$ with respect to $R$ and $k + 1$ and
                $a_1$ is a real sum sequence of $f$ at $x$ along $g_1$ with respect to $n_1$ and result $s_1$ and
                $b_1$ is a real sum sequence of $f$ at $x$ along $h_2$ with respect to $k + 1$ and result $t_1$.
            $R$ is a simple order relation on $J$ by \cref{well_order_relation}.
            $F_1\subseteq J$ and $h_2$ is a bijection from $\initialSegment{k + 1}$ to $F_1$ and
                $h_2$ is increasing on $\initialSegment{k + 1}$ with respect to $R$
                by \cref{well_ordered_enumeration_of_subset}.
            Let $\beta = h_2(k + 1)$.
            Let $G = F_1\setminus \{\beta\}$.
            Let $h_1 = \restrl{h_2}{\initialSegment{k}}$.
            $k\in\naturalsPlus$ by \cref{subseteq}.
            $h_1$ is a well-ordered enumeration of $G$ in $J$ with respect to $R$ and $k$
                by \cref{well_ordered_enumeration_remove_last}.
            $1\in\initialSegment{k}$ by \cref{real_one_in_naturals,real_one_leq_naturalsplus,initial_segment_naturalsplus}.
            $1\in\dom{h_1}$ by \cref{well_ordered_enumeration_of_subset,bijections_dom,subseteq}.
            Let $\gamma = h_1(1)$.
            $\gamma\in G$ by \cref{well_ordered_enumeration_of_subset,bijections,surj,funs_type_apply,subseteq}.
            $\gamma\in F_1$ and $\gamma\neq \beta$ by \cref{setminus_elim_left,setminus_elim_right,singleton_iff}.
            $k + 1\in\initialSegment{k + 1}$ by \cref{real_naturals_closed_succ,initial_segment_naturalsplus}.
            $\beta\in F_1$ by \cref{bijections,surj,funs_type_apply,subseteq}.
            Take $u\in\initialSegment{n_1}$ such that $g_1(u) = \beta$ by \cref{bijections,surj}.
            Show that $n_1\neq 1$.
            \begin{subproof}
                Suppose not.
                $n_1 = 1$.
                Show that for all $y\in F_1$ we have $y = g_1(1)$.
                \begin{subproof}
                    Fix $y\in F_1$.
                    Take $r\in\initialSegment{1}$ such that $g_1(r) = y$ by \cref{bijections,surj,real_lt_subst_left}.
                    $r = 1$ by \cref{initial_segment_naturalsplus,real_one_in_naturals,real_one_leq_naturalsplus,real_naturals_subset_reals,subseteq,real_lt_trans,real_lt_irreflexive}.
                    $y = g_1(1)$ by \cref{real_lt_subst_left}.
                \end{subproof}
                $\beta = g_1(1)$ and $\gamma = g_1(1)$ by \cref{real_lt_subst_left}.
                $\gamma = \beta$ by \cref{real_lt_subst_left}.
                $\gamma\in \{\beta\}$ by \cref{singleton_intro}.
                Contradiction by \cref{setminus_elim_right}.
            \end{subproof}
            Take $q\in\naturalsPlus$ such that $n_1 = q + 1$ by \cref{real_naturals_predecessor}.
            $g_1$ is a bijection from $\initialSegment{q + 1}$ to $F_1$ by \cref{real_lt_subst_left}.
            $u\in\initialSegment{q + 1}$ and $g_1(u) = \beta$ by \cref{real_lt_subst_right}.
            Take $p$ such that
                $p$ is a bijection from $\initialSegment{q}$ to $G$ and
                for all $i\in\initialSegment{q}$ we have
                    $(i < u \land p(i) = g_1(i)) \lor (u \leq i \land p(i) = g_1(i + 1))$
                by \cref{finite_enumeration_remove_value}.
            $a_1\in\funs{\initialSegment{n_1}}{\reals}$ and $\sumFiniteSequence{a_1}{s_1}$ and
                for all $i\in\initialSegment{n_1}$ we have $a_1(i) = \apply{\apply{f}{g_1(i)}}{x}$
                by \cref{real_sum_sequence_along_enumeration}.
            $a_1\in\funs{\initialSegment{q + 1}}{\reals}$ by \cref{real_lt_subst_left}.
            Take $c,v$ such that
                $c\in\funs{\initialSegment{q}}{\reals}$ and
                $\sumFiniteSequence{c}{v}$ and
                $s_1 = v + a_1(u)$ and
                for all $i\in\initialSegment{q}$ we have
                    $(i < u \land c(i) = a_1(i)) \lor (u \leq i \land c(i) = a_1(i + 1))$
                by \cref{sum_finite_sequence_remove_value}.
            Show that for all $i\in\initialSegment{q}$ we have $c(i) = \apply{\apply{f}{p(i)}}{x}$.
            \begin{subproof}
                Fix $i\in\initialSegment{q}$.
                $\initialSegment{q}\subseteq \dom{a_1}$ and $\initialSegment{q}\subseteq \dom{g_1}$
                    by \cref{real_lt_subst_left,real_sum_sequence_along_enumeration,funs,bijections_dom,initial_segment_subset_succ,subseteq,subseteq_transitive}.
                $i\in\initialSegment{q + 1}$ by \cref{initial_segment_subset_succ,subseteq}.
                $i + 1\in\initialSegment{q + 1}$ by \cref{real_naturals_closed_succ,initial_segment_naturalsplus,real_lt_imp_leq,real_leq_add,real_one_in_naturals,real_naturals_subset_reals,subseteq,real_add_closed}.
                \begin{byCase}
                \caseOf{$i < u$.}
                    We have $(i < u \land c(i) = a_1(i)) \lor (u \leq i \land c(i) = a_1(i + 1))$
                        by assumption.
                    Show that $c(i) = a_1(i)$.
                    \begin{subproof}
                        Suppose not.
                        \begin{byCase}
                        \caseOf{$i < u \land c(i) = a_1(i)$.}
                            Contradiction.
                        \caseOf{$u \leq i \land c(i) = a_1(i + 1)$.}
                            $u < i$ or $u = i$ by \cref{real_leq_split}.
                            \begin{byCase}
                            \caseOf{$u < i$.}
                                $i\in\reals$ and $u\in\reals$ by \cref{initial_segment_naturalsplus,real_naturals_subset_reals,subseteq,real_lt_subst_left}.
                                Show that $u \neq i$.
                                \begin{subproof}
                                    Suppose not.
                                    $u < u$ by \cref{real_lt_subst_right}.
                                    Contradiction by \cref{real_lt_irreflexive}.
                                \end{subproof}
                                Contradiction by \cref{real_lt_asym}.
                            \caseOf{$u = i$.}
                                $i < i$ by \cref{real_lt_subst_right}.
                                $i\in\reals$ by \cref{initial_segment_naturalsplus,real_naturals_subset_reals,subseteq}.
                                Contradiction by \cref{real_lt_irreflexive}.
                            \end{byCase}
                        \end{byCase}
                    \end{subproof}
                    We have $(i < u \land p(i) = g_1(i)) \lor (u \leq i \land p(i) = g_1(i + 1))$
                        by assumption.
                    Show that $p(i) = g_1(i)$.
                    \begin{subproof}
                        Suppose not.
                        \begin{byCase}
                        \caseOf{$i < u \land p(i) = g_1(i)$.}
                            Contradiction.
                        \caseOf{$u \leq i \land p(i) = g_1(i + 1)$.}
                            $u < i$ or $u = i$ by \cref{real_leq_split}.
                            \begin{byCase}
                            \caseOf{$u < i$.}
                                $i\in\reals$ and $u\in\reals$ by \cref{initial_segment_naturalsplus,real_naturals_subset_reals,subseteq,real_lt_subst_left}.
                                Show that $u \neq i$.
                                \begin{subproof}
                                    Suppose not.
                                    $u < u$ by \cref{real_lt_subst_right}.
                                    Contradiction by \cref{real_lt_irreflexive}.
                                \end{subproof}
                                Contradiction by \cref{real_lt_asym}.
                            \caseOf{$u = i$.}
                                $i < i$ by \cref{real_lt_subst_right}.
                                $i\in\reals$ by \cref{initial_segment_naturalsplus,real_naturals_subset_reals,subseteq}.
                                Contradiction by \cref{real_lt_irreflexive}.
                            \end{byCase}
                        \end{byCase}
                    \end{subproof}
                    $a_1(i) = \apply{\apply{f}{g_1(i)}}{x}$ by \cref{real_sum_sequence_along_enumeration,real_lt_subst_left}.
                    $c(i) = \apply{\apply{f}{p(i)}}{x}$ by \cref{real_lt_subst_left,real_lt_subst_right}.
                \caseOf{$u \leq i$.}
                    We have $(i < u \land c(i) = a_1(i)) \lor (u \leq i \land c(i) = a_1(i + 1))$
                        by assumption.
                    Show that $c(i) = a_1(i + 1)$.
                    \begin{subproof}
                        Suppose not.
                        \begin{byCase}
                        \caseOf{$i < u \land c(i) = a_1(i)$.}
                            $u < i$ or $u = i$ by \cref{real_leq_split}.
                            \begin{byCase}
                            \caseOf{$u < i$.}
                                $i\in\reals$ and $u\in\reals$ by \cref{initial_segment_naturalsplus,real_naturals_subset_reals,subseteq,real_lt_subst_left}.
                                Show that $u \neq i$.
                                \begin{subproof}
                                    Suppose not.
                                    $u < u$ by \cref{real_lt_subst_right}.
                                    Contradiction by \cref{real_lt_irreflexive}.
                                \end{subproof}
                                Contradiction by \cref{real_lt_asym}.
                            \caseOf{$u = i$.}
                                $u < u$ by \cref{real_lt_subst_left}.
                                $u\in\reals$ by \cref{initial_segment_naturalsplus,real_naturals_subset_reals,subseteq,real_lt_subst_left}.
                                Contradiction by \cref{real_lt_irreflexive}.
                            \end{byCase}
                        \caseOf{$u \leq i \land c(i) = a_1(i + 1)$.}
                            Contradiction.
                        \end{byCase}
                    \end{subproof}
                    We have $(i < u \land p(i) = g_1(i)) \lor (u \leq i \land p(i) = g_1(i + 1))$
                        by assumption.
                    Show that $p(i) = g_1(i + 1)$.
                    \begin{subproof}
                        Suppose not.
                        \begin{byCase}
                        \caseOf{$i < u \land p(i) = g_1(i)$.}
                            $u < i$ or $u = i$ by \cref{real_leq_split}.
                            \begin{byCase}
                            \caseOf{$u < i$.}
                                $i\in\reals$ and $u\in\reals$ by \cref{initial_segment_naturalsplus,real_naturals_subset_reals,subseteq,real_lt_subst_left}.
                                Show that $u \neq i$.
                                \begin{subproof}
                                    Suppose not.
                                    $u < u$ by \cref{real_lt_subst_right}.
                                    Contradiction by \cref{real_lt_irreflexive}.
                                \end{subproof}
                                Contradiction by \cref{real_lt_asym}.
                            \caseOf{$u = i$.}
                                $u < u$ by \cref{real_lt_subst_left}.
                                $u\in\reals$ by \cref{initial_segment_naturalsplus,real_naturals_subset_reals,subseteq,real_lt_subst_left}.
                                Contradiction by \cref{real_lt_irreflexive}.
                            \end{byCase}
                        \caseOf{$u \leq i \land p(i) = g_1(i + 1)$.}
                            Contradiction.
                        \end{byCase}
                    \end{subproof}
                    $a_1(i + 1) = \apply{\apply{f}{g_1(i + 1)}}{x}$ by \cref{real_sum_sequence_along_enumeration,real_lt_subst_left}.
                    $c(i) = \apply{\apply{f}{p(i)}}{x}$ by \cref{real_lt_subst_left,real_lt_subst_right}.
                \end{byCase}
            \end{subproof}
            $c$ is a real sum sequence of $f$ at $x$ along $p$ with respect to $q$ and result $v$
                by \cref{real_sum_sequence_along_enumeration}.
            $b_1\in\funs{\initialSegment{k + 1}}{\reals}$ and $\sumFiniteSequence{b_1}{t_1}$ and
                for all $i\in\initialSegment{k + 1}$ we have $b_1(i) = \apply{\apply{f}{h_2(i)}}{x}$
                by \cref{real_sum_sequence_along_enumeration}.
            Let $b_2 = \restrl{b_1}{\initialSegment{k}}$.
            $b_2\in\funs{\initialSegment{k}}{\reals}$ by \cref{sum_finite_sequence_split_last}.
            Take $w\in\reals$ such that $\sumFiniteSequence{b_2}{w}$ and $t_1 = w + b_1(k + 1)$
                by \cref{sum_finite_sequence_split_last}.
            Show that for all $i\in\initialSegment{k}$ we have $b_2(i) = \apply{\apply{f}{h_1(i)}}{x}$.
            \begin{subproof}
                Fix $i\in\initialSegment{k}$.
                $\initialSegment{k}\subseteq \dom{b_1}$ and $\initialSegment{k}\subseteq \dom{h_2}$
                    by \cref{real_sum_sequence_along_enumeration,funs,well_ordered_enumeration_of_subset,bijections_dom,initial_segment_subset_succ,subseteq,subseteq_transitive}.
                $b_2(i) = b_1(i)$ and $h_1(i) = h_2(i)$
                    by \cref{restrl_of_function_apply,funs_is_function,bijection_is_function}.
                $i\in\initialSegment{k + 1}$ by \cref{initial_segment_subset_succ,subseteq}.
                $b_1(i) = \apply{\apply{f}{h_2(i)}}{x}$ by \cref{real_sum_sequence_along_enumeration,real_lt_subst_left}.
                $b_2(i) = \apply{\apply{f}{h_1(i)}}{x}$ by \cref{real_lt_subst_left,real_lt_subst_right}.
            \end{subproof}
            $b_2$ is a real sum sequence of $f$ at $x$ along $h_1$ with respect to $k$ and result $w$
                by \cref{real_sum_sequence_along_enumeration}.
            $v = w$ by \cref{subseteq}.
            $u\in\dom{a_1}$ by \cref{funs,real_lt_subst_left}.
            $a_1(u) = \apply{\apply{f}{g_1(u)}}{x}$ by \cref{real_sum_sequence_along_enumeration,real_lt_subst_left}.
            $a_1(u) = \apply{\apply{f}{\beta}}{x}$ by \cref{real_lt_subst_right}.
            $k + 1\in\dom{h_2}$ by \cref{well_ordered_enumeration_of_subset,bijections_dom,subseteq,real_naturals_closed_succ,initial_segment_naturalsplus}.
            $b_1(k + 1) = \apply{\apply{f}{h_2(k + 1)}}{x}$ by \cref{real_sum_sequence_along_enumeration,real_lt_subst_left}.
            $b_1(k + 1) = \apply{\apply{f}{\beta}}{x}$ by \cref{real_lt_subst_right}.
            $a_1(u) = b_1(k + 1)$ by \cref{real_lt_subst_left}.
            $s_1 = v + b_1(k + 1)$ by \cref{real_lt_subst_right}.
            $s_1 = w + b_1(k + 1)$ by \cref{real_lt_subst_left,real_lt_subst_right}.
            $s_1 = t_1$ by \cref{real_lt_subst_right}.
        \end{subproof}
        $k + 1\in A$ by \cref{real_naturals_closed_succ}.
    \end{subproof}
    $A\subseteq \naturalsPlus$ by \cref{subseteq}.
    $m\in\naturalsPlus$ by \cref{well_ordered_enumeration_of_subset}.
    $m\in A$ by \cref{subseteq,real_naturals_induction}.
    $s = t$ by \cref{subseteq}.
\end{proof}

% from §41 Item 10: positive function bounded by local bounds
\begin{theorem}\label{positive_function_with_local_bounds}
    Assume $X$ is paracompact with respect to $T$.
    Assume $X$ is Hausdorff with respect to $T$.
    Assume $C$ is locally finite in $X$ with respect to $T$.
    Assume $e$ is a function on $C$.
    Suppose for all $A\in C$ we have $\apply{e}{A}\in\reals$ and $0 < \apply{e}{A}$.
    Then there exists $f$ such that
        $f$ is continuous from $X$ to $\reals$ with respect to $T$ and $\standardTopologyR$ and
        for all $x\in X$ we have $0 < f(x)$ and
        for all $A\in C$ for all $x\in A$ we have $f(x) \leq \apply{e}{A}$.
\end{theorem}
\begin{proof}
    \begin{byCase}
    \caseOf{$C = \emptyset$.}
        Let $f = \{(x,1) \mid x\in X\}$.
        $1\in\reals$ by \cref{real_one_in_naturals,real_naturals_subset_reals,subseteq}.
        $f$ is a function from $X$ to $\reals$ by \cref{constant_map_function}.
        For all $x\in X$ we have $f(x) = 1$ by \cref{constant_map_function,funs_is_function,function_apply_intro}.
        $T$ is a topology on $X$ by \cref{paracompact_space}.
        $\standardTopologyR$ is a topology on $\reals$
            by \cref{standard_basis_is_basis,basis_topology_is_topology,standard_topology_reals}.
        Show that $f$ is continuous from $X$ to $\reals$ with respect to $T$ and $\standardTopologyR$.
        \begin{subproof}
            Follows by \cref{constant_continuous}.
        \end{subproof}
        Show that for all $x\in X$ we have $0 < f(x)$.
        \begin{subproof}
            Follows by \cref{real_zero_lt_one,real_lt_subst_right}.
        \end{subproof}
        Show that for all $A\in C$ for all $x\in A$ we have $f(x) \leq \apply{e}{A}$.
        \begin{subproof}
            Suppose not.
            There exist $A,x$ such that $A\in C$ and $x\in A$.
            $C\subseteq\emptyset$ by \cref{subseteq_emptyset_iff}.
            $A\in\emptyset$ by \cref{elem_subseteq}.
            Contradiction by \cref{emptyset}.
        \end{subproof}
        Let $g_0 = f$.
        Show that $g_0$ is continuous from $X$ to $\reals$ with respect to $T$ and $\standardTopologyR$.
        \begin{subproof}
            Follows by \cref{real_lt_subst_left}.
        \end{subproof}
        Show that for all $x\in X$ we have $0 < g_0(x)$.
        \begin{subproof}
            Follows by \cref{real_lt_subst_left,real_lt_subst_right}.
        \end{subproof}
        Show that for all $A\in C$ for all $x\in A$ we have $g_0(x) \leq \apply{e}{A}$.
        \begin{subproof}
            Follows by \cref{real_lt_subst_left,real_lt_subst_right}.
        \end{subproof}
        There exists $g$ such that
            $g$ is continuous from $X$ to $\reals$ with respect to $T$ and $\standardTopologyR$ and
            for all $x\in X$ we have $0 < g(x)$ and
            for all $A\in C$ for all $x\in A$ we have $g(x) \leq \apply{e}{A}$
            by \cref{real_lt_subst_left,real_lt_subst_right}.
    \caseOf{$C \neq \emptyset$.}
        Let $R_U = \{w\in C\times\pow{X} \mid \text{there exist $B\in C$ such that there exist $V\in\pow{X}$ such that
            $w = (B,V)$ and
            there exist $C_0,H_0$ such that
                $C_0 = \{A\in C \mid \apply{e}{A} < \apply{e}{B}\}$ and
                $H_0$ is the closure collection of $C_0$ in $X$ with respect to $T$ and
                $V = X\setminus \unions{H_0}$}\}$.
        Show that for all $B\in C$ there exists $V$ such that $B\mathrel{R_U}V$.
        \begin{subproof}
            Fix $B\in C$.
            Let $C_0 = \{A\in C \mid \apply{e}{A} < \apply{e}{B}\}$.
            Let $H_0 = \{D\in\pow{X} \mid \exists A\in C_0. D = \closure{A}{X}{T}\}$.
            $H_0$ is the closure collection of $C_0$ in $X$ with respect to $T$ by \cref{closure_collection_in}.
            Let $V = X\setminus \unions{H_0}$.
            $V\subseteq X$ by \cref{setminus_subseteq}.
            $V\in\pow{X}$ by \cref{powerset_intro}.
            $(B,V)\in C\times\pow{X}$ by \cref{times_tuple_intro}.
            $B\mathrel{R_U}V$ by \cref{subseteq,pair_eq_iff}.
        \end{subproof}
        $R_U$ is a relation by \cref{relation}.
        Take $U$ such that $U$ is a function on $C$ and for all $B\in C$ we have $B\mathrel{R_U}\apply{U}{B}$
            by \cref{choice_function}.

        Show that for all $B\in C$ we have $\apply{U}{B}$ is open in $X$ with respect to $T$.
        \begin{subproof}
            Fix $B\in C$.
            Take $C_B,H_B$ such that
                $C_B = \{A\in C \mid \apply{e}{A} < \apply{e}{B}\}$ and
                $H_B$ is the closure collection of $C_B$ in $X$ with respect to $T$ and
                $\apply{U}{B} = X\setminus \unions{H_B}$
                by \cref{choice_function,pair_eq_iff}.
            $C_B\subseteq C$ by \cref{subseteq}.
            $C_B$ is locally finite in $X$ with respect to $T$ by \cref{locally_finite_subcollection}.
            $H_B$ is locally finite in $X$ with respect to $T$ by \cref{locally_finite_closure_collection}.
            $T$ is a topology on $X$ by \cref{paracompact_space}.
            For all $D\in H_B$ we have $D$ is closed in $X$ with respect to $T$
                by \cref{closure_collection_in,subseteq,locally_finite_collection,powerset_elim,closure_closed_in,real_lt_subst_left}.
            $H_B\subseteq H_B$ by \cref{subseteq_refl}.
            $\unions{H_B}$ is closed in $X$ with respect to $T$
                by \cref{locally_finite_closed_subcollection_union_closed}.
            $X\setminus \unions{H_B}$ is open in $X$ with respect to $T$
                by \cref{closed_in,topology_on,subseteq,pow_iff,setminus_subseteq,double_relative_complement,open_in}.
            $\apply{U}{B}$ is open in $X$ with respect to $T$ by \cref{real_lt_subst_left}.
        \end{subproof}

        Show that $\img{U}{C}$ is a covering of $X$.
        \begin{subproof}
            Show that for all $x\in X$ there exists $Z\in\img{U}{C}$ such that $x\in Z$.
            \begin{subproof}
                Fix $x\in X$.
                Take $W,F$ such that
                    $W$ is a neighborhood of $x$ in $X$ with respect to $T$ and
                    $F\subseteq C$ and
                    $F$ is finite and
                    for all $A\in C$ if $W$ intersects $A$ then $A\in F$
                    by \cref{locally_finite_collection}.
                Let $E = \{A\in F \mid x\in\closure{A}{X}{T}\}$.
                $E\subseteq F$ by \cref{subseteq}.
                $E$ is finite by \cref{subseteq_finite_is_finite}.
                \begin{byCase}
                \caseOf{$E = \emptyset$.}
                    Take $B$ such that $B\in C$.
                    Take $C_B,H_B$ such that
                        $C_B = \{A\in C \mid \apply{e}{A} < \apply{e}{B}\}$ and
                        $H_B$ is the closure collection of $C_B$ in $X$ with respect to $T$ and
                        $\apply{U}{B} = X\setminus \unions{H_B}$
                        by \cref{choice_function,pair_eq_iff}.
                    Show that $x\in\apply{U}{B}$.
                    \begin{subproof}
                        Suppose not.
                        $x\in\unions{H_B}$
                            by \cref{pair_eq_iff,function_apply_intro,setminus}.
                        Take $D$ such that $D\in H_B$ and $x\in D$
                            by \cref{unions_iff}.
                        Take $A$ such that $A\in C_B$ and $D = \closure{A}{X}{T}$ by \cref{closure_collection_in}.
                        $A\in C$ and $\apply{e}{A} < \apply{e}{B}$ by \cref{subseteq}.
                        $x\in\closure{A}{X}{T}$ by \cref{real_lt_subst_left}.
                        $C$ is locally finite in $X$ with respect to $T$ by assumption.
                        $C\subseteq\pow{X}$ by \cref{locally_finite_collection}.
                        $A\in\pow{X}$ by \cref{subseteq}.
                        $A\subseteq X$ by \cref{powerset_elim}.
                        $T$ is a topology on $X$ by \cref{paracompact_space}.
                        $W$ intersects $A$ by \cref{neighborhood,closure_iff_intersects_open}.
                        $A\in F$ by \cref{subseteq}.
                        $A\in E$ by \cref{subseteq}.
                        Contradiction by \cref{emptyset}.
                    \end{subproof}
                    $\apply{U}{B}\in\img{U}{C}$ by \cref{img_iff,function_apply_elim}.
                    There exists $Z\in\img{U}{C}$ such that $x\in Z$ by \cref{real_lt_subst_left}.
                \caseOf{$E \neq \emptyset$.}
                    $E$ is inhabited by \cref{nonempty_inhabited}.
                    Let $g_0 = \restrl{e}{E}$.
                    $g_0$ is a function on $E$ by \cref{restrl_of_function_is_function,dom_of_restrl_eq_inter,funs,subseteq,inter_absorb_supseteq_left}.
                    Let $M = \img{g_0}{E}$.
                    $M$ is finite by \cref{finite_image_function}.
                    Take $A_0$ such that $A_0\in E$.
                    $\apply{g_0}{A_0}\in M$ by \cref{function_apply_elim,img_iff}.
                    $M\neq\emptyset$ by \cref{emptyset}.
                    Show that for all $t\in M$ we have $t\in\reals$.
                    \begin{subproof}
                        Fix $t\in M$.
                        Take $A_1$ such that $A_1\in E$ and $(A_1,t)\in g_0$ by \cref{img_iff}.
                        $A_1\in\dom{g_0}$ by \cref{subseteq}.
                        $t = \apply{g_0}{A_1}$ by \cref{function_apply_iff}.
                        $A_1\in C$ by \cref{subseteq,subseteq_transitive}.
                        $\apply{e}{A_1}\in\reals$ by assumption.
                        $\apply{g_0}{A_1} = \apply{e}{A_1}$ by \cref{restrl_of_function_apply,subseteq}.
                        $t\in\reals$ by \cref{real_lt_subst_right}.
                    \end{subproof}
                    $M\subseteq\reals$ by \cref{subseteq}.
                    Take $m\in M$ such that $m$ is a minimum of $M$ by \cref{finite_real_minimum}.
                    Take $B$ such that $B\in E$ and $(B,m)\in g_0$ by \cref{img_iff}.
                    $B\in\dom{g_0}$ by \cref{subseteq}.
                    $B\in E$ and $m = \apply{g_0}{B}$ by \cref{function_apply_iff}.
                    $B\in C$ by \cref{subseteq,subseteq_transitive}.
                    $\apply{g_0}{B} = \apply{e}{B}$ by \cref{restrl_of_function_apply,subseteq}.
                    $m = \apply{e}{B}$ by \cref{real_lt_subst_right}.
                    Take $C_B,H_B$ such that
                        $C_B = \{A\in C \mid \apply{e}{A} < \apply{e}{B}\}$ and
                        $H_B$ is the closure collection of $C_B$ in $X$ with respect to $T$ and
                        $\apply{U}{B} = X\setminus \unions{H_B}$
                        by \cref{choice_function,pair_eq_iff}.
                    Show that $x\in\apply{U}{B}$.
                    \begin{subproof}
                        Suppose not.
                        $x\in\unions{H_B}$
                            by \cref{pair_eq_iff,function_apply_intro,setminus}.
                        Take $D$ such that $D\in H_B$ and $x\in D$
                            by \cref{unions_iff}.
                        Take $A$ such that $A\in C_B$ and $D = \closure{A}{X}{T}$ by \cref{closure_collection_in}.
                        $A\in C$ and $\apply{e}{A} < \apply{e}{B}$ by \cref{subseteq}.
                        $x\in\closure{A}{X}{T}$ by \cref{real_lt_subst_left}.
                        $C$ is locally finite in $X$ with respect to $T$ by assumption.
                        $C\subseteq\pow{X}$ by \cref{locally_finite_collection}.
                        $A\in\pow{X}$ by \cref{subseteq}.
                        $A\subseteq X$ by \cref{powerset_elim}.
                        $T$ is a topology on $X$ by \cref{paracompact_space}.
                        $W$ intersects $A$ by \cref{neighborhood,closure_iff_intersects_open}.
                        $A\in F$ by \cref{subseteq}.
                        $A\in E$ by \cref{subseteq}.
                        Let $r = \apply{g_0}{A}$.
                        $r = \apply{e}{A}$ by \cref{restrl_of_function_apply,subseteq}.
                        $r\in M$ by \cref{function_apply_elim,img_iff}.
                        $m \leq r$ by \cref{real_minimum}.
                        $m \leq \apply{e}{A}$ by \cref{real_lt_subst_right}.
                        $m < m$ by \cref{real_lt_subst_left,real_lt_trans}.
                        Contradiction by \cref{real_lt_irreflexive}.
                    \end{subproof}
                    $\apply{U}{B}\in\img{U}{C}$ by \cref{img_iff,function_apply_elim}.
                    There exists $Z\in\img{U}{C}$ such that $x\in Z$ by \cref{real_lt_subst_left}.
                \end{byCase}
            \end{subproof}
            For all $x\in X$ we have $x\in\unions{\img{U}{C}}$ by \cref{unions_iff}.
            $X\subseteq\unions{\img{U}{C}}$ by \cref{subseteq}.
            $\img{U}{C}$ is a covering of $X$ by \cref{covering}.
        \end{subproof}

        Take $p$ such that $p$ is an indexed partition of unity on $X$ dominated by $U$ with respect to $T$ and $C$
            by \cref{paracompact_partition_of_unity}.
        Take $B_p$ such that $B_p\in C$ by \cref{nonempty_inhabited}.
        Take $I_p$ such that
            $I_p = \intervalclosed{0}{1}$ and
            $\apply{p}{B_p}$ is continuous from $X$ to $I_p$ with respect to $T$ and $\subspaceTopology{\standardTopologyR}{I_p}$ and
            $\supp{\apply{p}{B_p}}{X}{T}\subseteq \apply{U}{B_p}$ and
            there exists $S$ such that
                $S$ is the support family of $p$ on $X$ with respect to $T$ and $C$ and
                $S$ is a locally finite indexed family on $C$ in $X$ with respect to $T$ and
                for all $x\in X$ we have $1$ is the pointwise sum of $p$ at $x$ on $X$ with respect to $T$ and $C$
            by \cref{indexed_partition_of_unity,pair_eq_iff}.
        Take $S$ such that
            $S$ is the support family of $p$ on $X$ with respect to $T$ and $C$ and
            $S$ is a locally finite indexed family on $C$ in $X$ with respect to $T$ and
            for all $x\in X$ we have $1$ is the pointwise sum of $p$ at $x$ on $X$ with respect to $T$ and $C$
            by \cref{subseteq}.
        $S$ is a relation by \cref{support_family,relation,pair_eq_iff}.
        Show that for all $A,Y_1,Y_2$ if $(A,Y_1)\in S$ and $(A,Y_2)\in S$ then $Y_1 = Y_2$.
        \begin{subproof}
            Fix $A$.
            Fix $Y_1$.
            Fix $Y_2$.
            Assume $(A,Y_1)\in S$ and $(A,Y_2)\in S$.
            Take $\alpha_1\in C$ such that $(\alpha_1,\supp{\apply{p}{\alpha_1}}{X}{T}) = (A,Y_1)$ by \cref{support_family}.
            Take $\alpha_2\in C$ such that $(\alpha_2,\supp{\apply{p}{\alpha_2}}{X}{T}) = (A,Y_2)$ by \cref{support_family}.
            $\alpha_1 = A$ and $\supp{\apply{p}{\alpha_1}}{X}{T} = Y_1$ by \cref{pair_eq_iff}.
            $\alpha_2 = A$ and $\supp{\apply{p}{\alpha_2}}{X}{T} = Y_2$ by \cref{pair_eq_iff}.
            $\alpha_1 = \alpha_2$ by \cref{real_lt_subst_left}.
            $\supp{\apply{p}{\alpha_1}}{X}{T} = \supp{\apply{p}{\alpha_2}}{X}{T}$ by \cref{real_lt_subst_left}.
            $Y_1 = Y_2$ by \cref{real_lt_subst_left,real_lt_subst_right}.
        \end{subproof}
        $S$ is right-unique by \cref{rightunique}.
        Show that for all $A\in\dom{S}$ we have $A\in C$.
        \begin{subproof}
            Fix $A\in\dom{S}$.
            Take $Y$ such that $(A,Y)\in S$ by \cref{dom_iff}.
            Take $\alpha\in C$ such that $(\alpha,\supp{\apply{p}{\alpha}}{X}{T}) = (A,Y)$ by \cref{support_family}.
            $\alpha = A$ by \cref{pair_eq_iff}.
            $A\in C$ by \cref{real_lt_subst_left}.
        \end{subproof}
        Show that for all $A\in C$ we have $A\in\dom{S}$.
        \begin{subproof}
            Fix $A\in C$.
            $(A,\supp{\apply{p}{A}}{X}{T})\in S$ by \cref{support_family}.
            $A\in\dom{S}$ by \cref{dom_intro}.
        \end{subproof}
        $\dom{S}\subseteq C$ by \cref{subseteq}.
        $C\subseteq\dom{S}$ by \cref{subseteq}.
        $\dom{S} = C$ by \cref{subseteq_antisymmetric}.
        $S$ is a function on $C$ by \cref{funs}.

        Let $B_0 = \pow{X\times\reals}$.
        Let $Q = \{w\in C\times B_0 \mid \text{there exist $A\in C$ such that there exists $u\in B_0$ such that
            $w = (A,u)$ and
            $u$ is continuous from $X$ to $\reals$ with respect to $T$ and $\standardTopologyR$ and
            for all $x\in X$ we have $u(x) = \apply{e}{A} \rmul \apply{\apply{p}{A}}{x}$}\}$.
        Show that for all $A\in C$ there exists $u$ such that $A\mathrel{Q}u$.
        \begin{subproof}
            Fix $A\in C$.
            Take $I_A$ such that
                $I_A = \intervalclosed{0}{1}$ and
                $\apply{p}{A}$ is continuous from $X$ to $I_A$ with respect to $T$ and $\subspaceTopology{\standardTopologyR}{I_A}$ and
                $\supp{\apply{p}{A}}{X}{T}\subseteq \apply{U}{A}$
                by \cref{indexed_partition_of_unity,pair_eq_iff}.
            Let $j_A = \identity{I_A}$.
            Let $p_A = j_A\circ \apply{p}{A}$.
            $p_A$ is continuous from $X$ to $\reals$ with respect to $T$ and $\standardTopologyR$
                by \cref{continuous_interval_expand_range}.
            For all $x\in X$ we have $p_A(x) = \apply{\apply{p}{A}}{x}$
                by \cref{continuous,funs_is_function,funs,function_rightunique,funs_is_relation,funs_subseteq_rels,subseteq,rels_ran_subseteq,id_dom,id_rightunique,id_is_relation,circ_apply,funs_type_apply,id_apply}.
            Let $d_A = \{(x,\apply{e}{A}) \mid x\in X\}$.
            $d_A$ is a function from $X$ to $\reals$ by \cref{constant_map_function}.
            For all $x\in X$ we have $d_A(x) = \apply{e}{A}$ by \cref{constant_map_function,funs_is_function,function_apply_intro}.
            $T$ is a topology on $X$ by \cref{paracompact_space}.
            $\standardTopologyR$ is a topology on $\reals$
                by \cref{standard_basis_is_basis,basis_topology_is_topology,standard_topology_reals}.
            $d_A$ is continuous from $X$ to $\reals$ with respect to $T$ and $\standardTopologyR$ by \cref{constant_continuous}.
            Let $u = \{(x,d_A(x) \rmul p_A(x)) \mid x\in X\}$.
            $u$ is continuous from $X$ to $\reals$ with respect to $T$ and $\standardTopologyR$
                by \cref{real_function_multiplication_continuous}.
            For all $x\in X$ we have $u(x) = \apply{e}{A} \rmul \apply{\apply{p}{A}}{x}$
                by \cref{pair_eq_iff,continuous,funs_is_function,funs,function_apply_intro,real_mul_eq_subst,subseteq}.
            $A\mathrel{Q}u$
                by \cref{continuous,funs_is_relation,funs,funs_subseteq_rels,subseteq,rels_ran_subseteq,rels_elim,powerset_intro,times_tuple_intro,pair_eq_iff}.
        \end{subproof}
        $Q$ is a relation by \cref{relation}.
        Take $q$ such that $q$ is a function on $C$ and for all $A\in C$ we have $A\mathrel{Q}\apply{q}{A}$
            by \cref{choice_function}.
        Show that for all $A\in C$ we have
            $\apply{q}{A}$ is continuous from $X$ to $\reals$ with respect to $T$ and $\standardTopologyR$ and
            for all $x\in X$ we have $\apply{\apply{q}{A}}{x} = \apply{e}{A} \rmul \apply{\apply{p}{A}}{x}$.
        \begin{subproof}
            Follows by \cref{choice_function,pair_eq_iff}.
        \end{subproof}

        Let $S_1 = \{(A,\supp{\apply{q}{A}}{X}{T}) \mid A\in C\}$.
        $S_1$ is the support family of $q$ on $X$ with respect to $T$ and $C$ by \cref{support_family}.
        $S_1$ is a relation by \cref{support_family,relation,pair_eq_iff}.
        Show that for all $A,Y_1,Y_2$ if $(A,Y_1)\in S_1$ and $(A,Y_2)\in S_1$ then $Y_1 = Y_2$.
        \begin{subproof}
            Fix $A$.
            Fix $Y_1$.
            Fix $Y_2$.
            Assume $(A,Y_1)\in S_1$ and $(A,Y_2)\in S_1$.
            Take $\alpha_1\in C$ such that $(\alpha_1,\supp{\apply{q}{\alpha_1}}{X}{T}) = (A,Y_1)$ by \cref{support_family}.
            Take $\alpha_2\in C$ such that $(\alpha_2,\supp{\apply{q}{\alpha_2}}{X}{T}) = (A,Y_2)$ by \cref{support_family}.
            $\alpha_1 = A$ and $\supp{\apply{q}{\alpha_1}}{X}{T} = Y_1$ by \cref{pair_eq_iff}.
            $\alpha_2 = A$ and $\supp{\apply{q}{\alpha_2}}{X}{T} = Y_2$ by \cref{pair_eq_iff}.
            $\alpha_1 = \alpha_2$ by \cref{real_lt_subst_left}.
            $\supp{\apply{q}{\alpha_1}}{X}{T} = \supp{\apply{q}{\alpha_2}}{X}{T}$ by \cref{real_lt_subst_left}.
            $Y_1 = Y_2$ by \cref{real_lt_subst_left,real_lt_subst_right}.
        \end{subproof}
        $S_1$ is right-unique by \cref{rightunique}.
        Show that for all $A\in\dom{S_1}$ we have $A\in C$.
        \begin{subproof}
            Fix $A\in\dom{S_1}$.
            Take $Y$ such that $(A,Y)\in S_1$ by \cref{dom_iff}.
            Take $\alpha\in C$ such that $(\alpha,\supp{\apply{q}{\alpha}}{X}{T}) = (A,Y)$ by \cref{support_family}.
            $\alpha = A$ by \cref{pair_eq_iff}.
            $A\in C$ by \cref{real_lt_subst_left}.
        \end{subproof}
        Show that for all $A\in C$ we have $A\in\dom{S_1}$.
        \begin{subproof}
            Fix $A\in C$.
            $(A,\supp{\apply{q}{A}}{X}{T})\in S_1$ by \cref{support_family}.
            $A\in\dom{S_1}$ by \cref{dom_intro}.
        \end{subproof}
        $\dom{S_1}\subseteq C$ by \cref{subseteq}.
        $C\subseteq\dom{S_1}$ by \cref{subseteq}.
        $\dom{S_1} = C$ by \cref{subseteq_antisymmetric}.
        $S_1$ is a function on $C$ by \cref{funs}.
        Show that $S_1 = S$.
        \begin{subproof}
            Show that for all $A\in C$ we have $\apply{S_1}{A} = \apply{S}{A}$.
            \begin{subproof}
                Fix $A\in C$.
                Take $I_A$ such that
                    $I_A = \intervalclosed{0}{1}$ and
                    $\apply{p}{A}$ is continuous from $X$ to $I_A$ with respect to $T$ and $\subspaceTopology{\standardTopologyR}{I_A}$ and
                    $\supp{\apply{p}{A}}{X}{T}\subseteq \apply{U}{A}$
                    by \cref{indexed_partition_of_unity,pair_eq_iff}.
                Let $j_A = \identity{I_A}$.
                Let $p_A = j_A\circ \apply{p}{A}$.
                $p_A$ is continuous from $X$ to $\reals$ with respect to $T$ and $\standardTopologyR$
                    by \cref{continuous_interval_expand_range}.
                For all $x\in X$ we have $p_A(x) = \apply{\apply{p}{A}}{x}$
                    by \cref{continuous,funs_is_function,funs,function_rightunique,funs_is_relation,funs_subseteq_rels,subseteq,rels_ran_subseteq,id_dom,id_rightunique,id_is_relation,circ_apply,funs_type_apply,id_apply}.
                Show that for all $x\in X$ if $p_A(x) = 0$ then $\apply{\apply{q}{A}}{x} = 0$.
                \begin{subproof}
                    Fix $x\in X$.
                    Assume $p_A(x) = 0$.
                    $\apply{\apply{q}{A}}{x} = \apply{e}{A} \rmul p_A(x)$ by \cref{subseteq}.
                    $\apply{\apply{q}{A}}{x} = \apply{e}{A} \rmul 0$ by \cref{real_lt_subst_right}.
                    $\apply{e}{A}\in\reals$ by assumption.
                    $0\in\reals$ by \cref{real_add_ident}.
                    $\apply{e}{A} \rmul 0 = 0 \rmul \apply{e}{A}$ by \cref{real_mul_comm}.
                    $\apply{\apply{q}{A}}{x} = 0$ by \cref{real_mul_zero,real_lt_subst_right}.
                \end{subproof}
                $\supp{\apply{q}{A}}{X}{T}\subseteq \supp{p_A}{X}{T}$ by \cref{support_subseteq_from_zero_implication}.
                Show that for all $x\in X$ if $\apply{\apply{q}{A}}{x} = 0$ then $p_A(x) = 0$.
                \begin{subproof}
                    Fix $x\in X$.
                    Assume $\apply{\apply{q}{A}}{x} = 0$.
                    $\apply{\apply{q}{A}}{x} = \apply{e}{A} \rmul p_A(x)$ by \cref{subseteq}.
                    $\apply{e}{A} \rmul p_A(x) = 0$ by \cref{real_lt_subst_left}.
                    $\apply{e}{A}\in\reals$ and $0 < \apply{e}{A}$ by assumption.
                    $\apply{e}{A}\neq 0$ by \cref{real_pos_nonzero}.
                    Show that $p_A(x) = 0$.
                    \begin{subproof}
                        Suppose not.
                        $p_A(x)\neq 0$.
                        $p_A(x)\in\reals$ by \cref{continuous,funs_type_apply,subseteq}.
                        $\apply{e}{A} = 0$ or $p_A(x) = 0$ by \cref{real_mul_zero_product}.
                        \begin{byCase}
                        \caseOf{$\apply{e}{A} = 0$.}
                            Contradiction.
                        \caseOf{$p_A(x) = 0$.}
                            Contradiction.
                        \end{byCase}
                    \end{subproof}
                \end{subproof}
                $\supp{p_A}{X}{T}\subseteq \supp{\apply{q}{A}}{X}{T}$ by \cref{support_subseteq_from_zero_implication}.
                $\supp{\apply{q}{A}}{X}{T} = \supp{p_A}{X}{T}$ by \cref{subseteq_antisymmetric}.
                Show that $\supp{p_A}{X}{T} = \supp{\apply{p}{A}}{X}{T}$.
                \begin{subproof}
                    Show that $p_A = \apply{p}{A}$.
                    \begin{subproof}
                        $p_A$ is a function and $\dom{p_A} = X$ by \cref{continuous,funs,funs_is_function}.
                        $\apply{p}{A}$ is a function and $\dom{\apply{p}{A}} = X$ by \cref{continuous,funs,funs_is_function}.
                        For all $x\in X$ we have $p_A(x) = \apply{\apply{p}{A}}{x}$ by \cref{subseteq}.
                        $p_A = \apply{p}{A}$ by \cref{fun_subseteq,subseteq,subseteq_antisymmetric}.
                    \end{subproof}
                    $\supp{p_A}{X}{T} = \supp{\apply{p}{A}}{X}{T}$ by \cref{real_lt_subst_left}.
                \end{subproof}
                $\supp{\apply{q}{A}}{X}{T} = \supp{\apply{p}{A}}{X}{T}$ by \cref{real_lt_subst_right}.
                $(A,\supp{\apply{q}{A}}{X}{T})\in S_1$ and $(A,\supp{\apply{p}{A}}{X}{T})\in S$
                    by \cref{support_family,pair_eq_iff}.
                $\apply{S_1}{A} = \supp{\apply{q}{A}}{X}{T}$ by \cref{function_apply_intro}.
                $\apply{S}{A} = \supp{\apply{p}{A}}{X}{T}$ by \cref{function_apply_intro}.
                $\apply{S_1}{A} = \apply{S}{A}$ by \cref{real_lt_subst_left,real_lt_subst_right}.
            \end{subproof}
            $S_1 = S$ by \cref{funs_elim,subseteq,fun_subseteq,subseteq_antisymmetric}.
        \end{subproof}
        $S_1$ is a locally finite indexed family on $C$ in $X$ with respect to $T$ by \cref{real_lt_subst_left}.

        Show that for all $x\in X$ there exists $\beta\in C$ such that $x\in\apply{S_1}{\beta}$.
        \begin{subproof}
            Fix $x\in X$.
            $1$ is the pointwise sum of $p$ at $x$ on $X$ with respect to $T$ and $C$
                by \cref{indexed_partition_of_unity}.
            Take $F,n,h,a$ such that
                $F\subseteq C$ and
                $F$ is finite and
                $F$ is inhabited and
                $n\in\naturalsPlus$ and
                $h$ is a bijection from $\initialSegment{n}$ to $F$ and
                $a$ is a real sum sequence of $p$ at $x$ along $h$ with respect to $n$ and result $1$ and
                for all $\alpha\in C$ if $x\in\supp{\apply{p}{\alpha}}{X}{T}$ then $\alpha\in F$
                by \cref{pointwise_sum_of_indexed_family}.
            $a\in\funs{\initialSegment{n}}{\reals}$ and $\sumFiniteSequence{a}{1}$
                by \cref{real_sum_sequence_along_enumeration}.
            for all $i\in\initialSegment{n}$ we have $a(i) = \apply{\apply{p}{h(i)}}{x}$
                by \cref{real_sum_sequence_along_enumeration}.
            Show that for all $i\in\initialSegment{n}$ we have $0 < a(i)$ or $0 = a(i)$.
            \begin{subproof}
                Fix $i\in\initialSegment{n}$.
                $i\in\dom{h}$ by \cref{bijections_dom,subseteq}.
                $h(i)\in F$ by \cref{bijections,surj,funs_type_apply,subseteq}.
                $h(i)\in C$ by \cref{subseteq}.
                Take $I_i$ such that
                    $I_i = \intervalclosed{0}{1}$ and
                    $\apply{p}{h(i)}$ is continuous from $X$ to $I_i$ with respect to $T$ and $\subspaceTopology{\standardTopologyR}{I_i}$ and
                    $\supp{\apply{p}{h(i)}}{X}{T}\subseteq \apply{U}{h(i)}$
                    by \cref{indexed_partition_of_unity,pair_eq_iff}.
                $a(i) = \apply{\apply{p}{h(i)}}{x}$ by \cref{subseteq}.
                $\apply{\apply{p}{h(i)}}{x}\in I_i$ by \cref{continuous,funs_type_apply}.
                $0 \leq \apply{\apply{p}{h(i)}}{x}$ by \cref{intervalclosed}.
                $0 < \apply{\apply{p}{h(i)}}{x}$ or $0 = \apply{\apply{p}{h(i)}}{x}$ by \cref{real_leq_split}.
                $0 < a(i)$ or $0 = a(i)$ by \cref{real_lt_subst_left,real_lt_subst_right}.
            \end{subproof}
            Take $i\in\initialSegment{n}$ such that $0 < a(i)$ by \cref{finite_real_sum_one_positive_term}.
            $i\in\dom{h}$ by \cref{bijections_dom,subseteq}.
            Let $\beta = h(i)$.
            $\beta\in F$ by \cref{bijections,surj,funs_type_apply,subseteq}.
            $\beta\in C$ by \cref{subseteq}.
            $a(i) = \apply{\apply{p}{\beta}}{x}$ by \cref{real_lt_subst_right}.
            $0 < \apply{\apply{p}{\beta}}{x}$ by \cref{real_lt_subst_left}.
            Take $I_b$ such that
                $I_b = \intervalclosed{0}{1}$ and
                $\apply{p}{\beta}$ is continuous from $X$ to $I_b$ with respect to $T$ and $\subspaceTopology{\standardTopologyR}{I_b}$ and
                $\supp{\apply{p}{\beta}}{X}{T}\subseteq \apply{U}{\beta}$
                by \cref{indexed_partition_of_unity,pair_eq_iff}.
            Let $j_b = \identity{I_b}$.
            Let $p_b = j_b\circ \apply{p}{\beta}$.
            $p_b$ is continuous from $X$ to $\reals$ with respect to $T$ and $\standardTopologyR$
                by \cref{continuous_interval_expand_range}.
            For all $y\in X$ we have $p_b(y) = \apply{\apply{p}{\beta}}{y}$
                by \cref{continuous,funs_is_function,funs,function_rightunique,funs_is_relation,funs_subseteq_rels,subseteq,rels_ran_subseteq,id_dom,id_rightunique,id_is_relation,circ_apply,funs_type_apply,id_apply}.
            $p_b(x) = \apply{\apply{p}{\beta}}{x}$ by \cref{subseteq}.
            $0 < p_b(x)$ by \cref{real_lt_subst_right}.
            $p_b(x)\in\reals$ by \cref{continuous,funs_type_apply,subseteq}.
            $p_b(x)\neq 0$ by \cref{real_pos_nonzero}.
            $x\in\supp{p_b}{X}{T}$ by \cref{nonzero_value_in_support}.
            Show that $p_b = \apply{p}{\beta}$.
            \begin{subproof}
                $p_b$ is a function and $\dom{p_b} = X$ by \cref{continuous,funs,funs_is_function}.
                $\apply{p}{\beta}$ is a function and $\dom{\apply{p}{\beta}} = X$ by \cref{continuous,funs,funs_is_function}.
                For all $y\in X$ we have $p_b(y) = \apply{\apply{p}{\beta}}{y}$ by \cref{subseteq}.
                $p_b = \apply{p}{\beta}$ by \cref{fun_subseteq,subseteq,subseteq_antisymmetric}.
            \end{subproof}
            $\supp{p_b}{X}{T} = \supp{\apply{p}{\beta}}{X}{T}$ by \cref{real_lt_subst_left}.
            $x\in\supp{\apply{p}{\beta}}{X}{T}$ by \cref{real_lt_subst_left}.
            $x\in\apply{S}{\beta}$ by \cref{support_family,pair_eq_iff,function_apply_intro}.
            $\apply{\apply{q}{\beta}}{x} = \apply{e}{\beta} \rmul \apply{\apply{p}{\beta}}{x}$ by \cref{subseteq}.
            $0 < \apply{\apply{q}{\beta}}{x}$ by \cref{continuous,funs_type_apply,subseteq,real_mul_pos_pos_gt}.
            $\apply{\apply{q}{\beta}}{x}\neq 0$ by \cref{continuous,funs_type_apply,real_pos_nonzero}.
            $x\in\supp{\apply{q}{\beta}}{X}{T}$ by \cref{nonzero_value_in_support}.
            $x\in\apply{S_1}{\beta}$ by \cref{support_family,pair_eq_iff,function_apply_intro}.
            There exists $\gamma\in C$ such that $x\in\apply{S_1}{\gamma}$ by \cref{real_lt_subst_left}.
        \end{subproof}

        Take $L$ such that $L$ is a well order relation on $C$ by \cref{well_ordering_theorem_local_finiteness}.
        $T$ is a topology on $X$ by \cref{paracompact_space}.
        Take $g$ such that $g$ is a locally finite real sum function of $q$ on $X$ with respect to $T$ and $C$
            by \cref{locally_finite_real_sum_function_exists}.
        $g$ is continuous from $X$ to $\reals$ with respect to $T$ and $\standardTopologyR$ by \cref{locally_finite_real_sum_function}.

        Show that for all $x\in X$ we have $0 < g(x)$.
        \begin{subproof}
            Fix $x\in X$.
            $1$ is the pointwise sum of $p$ at $x$ on $X$ with respect to $T$ and $C$ by \cref{indexed_partition_of_unity}.
            Take $F_0,m_0,g_0,b_0$ such that
                $F_0\subseteq C$ and
                $F_0$ is finite and
                $F_0$ is inhabited and
                $m_0\in\naturalsPlus$ and
                $g_0$ is a bijection from $\initialSegment{m_0}$ to $F_0$ and
                $b_0$ is a real sum sequence of $p$ at $x$ along $g_0$ with respect to $m_0$ and result $1$ and
                for all $\alpha\in C$ if $x\in\supp{\apply{p}{\alpha}}{X}{T}$ then $\alpha\in F_0$
                by \cref{pointwise_sum_of_indexed_family}.
            $b_0\in\funs{\initialSegment{m_0}}{\reals}$ and $\sumFiniteSequence{b_0}{1}$ and
                for all $i\in\initialSegment{m_0}$ we have $b_0(i) = \apply{\apply{p}{g_0(i)}}{x}$
                by \cref{real_sum_sequence_along_enumeration}.
            Show that for all $i\in\initialSegment{m_0}$ we have $0 < b_0(i)$ or $0 = b_0(i)$.
            \begin{subproof}
                Fix $i\in\initialSegment{m_0}$.
                $i\in\dom{g_0}$ by \cref{bijections_dom,subseteq}.
                $g_0(i)\in F_0$ by \cref{bijections,surj,funs_type_apply,subseteq}.
                $g_0(i)\in C$ by \cref{subseteq}.
                Take $I_i$ such that
                    $I_i = \intervalclosed{0}{1}$ and
                    $\apply{p}{g_0(i)}$ is continuous from $X$ to $I_i$ with respect to $T$ and $\subspaceTopology{\standardTopologyR}{I_i}$ and
                    $\supp{\apply{p}{g_0(i)}}{X}{T}\subseteq \apply{U}{g_0(i)}$
                    by \cref{indexed_partition_of_unity,pair_eq_iff}.
                $\apply{\apply{p}{g_0(i)}}{x}\in I_i$ by \cref{continuous,funs_type_apply}.
                $0 \leq \apply{\apply{p}{g_0(i)}}{x}$ by \cref{intervalclosed}.
                $0 \leq b_0(i)$ by \cref{real_lt_subst_right}.
                $0 < b_0(i)$ or $0 = b_0(i)$ by \cref{real_leq_split}.
            \end{subproof}
            Take $i_0\in\initialSegment{m_0}$ such that $0 < b_0(i_0)$ by \cref{finite_real_sum_one_positive_term}.
            Let $\beta_0 = g_0(i_0)$.
            $\beta_0\in F_0$ by \cref{bijections,surj,funs_type_apply,subseteq}.
            $\beta_0\in C$ by \cref{subseteq}.
            $i_0\in\dom{g_0}$ by \cref{bijections_dom,subseteq}.
            $b_0(i_0)\in\reals$ by \cref{funs_type_apply}.
            $b_0(i_0) = \apply{\apply{p}{\beta_0}}{x}$ by \cref{real_sum_sequence_along_enumeration,real_lt_subst_left}.
            $0 < \apply{\apply{p}{\beta_0}}{x}$ by \cref{real_lt_subst_right}.
            $\apply{\apply{p}{\beta_0}}{x}\in\reals$ by \cref{real_lt_subst_right}.
            $\apply{e}{\beta_0}\in\reals$ and $0 < \apply{e}{\beta_0}$ by assumption.
            $\apply{\apply{q}{\beta_0}}{x} = \apply{e}{\beta_0} \rmul \apply{\apply{p}{\beta_0}}{x}$ by \cref{subseteq}.
            $0 < \apply{e}{\beta_0} \rmul \apply{\apply{p}{\beta_0}}{x}$ by \cref{real_mul_pos_pos_gt}.
            $0 < \apply{\apply{q}{\beta_0}}{x}$ by \cref{real_lt_subst_right}.
            $\apply{q}{\beta_0}$ is continuous from $X$ to $\reals$ with respect to $T$ and $\standardTopologyR$ by \cref{subseteq}.
            $x\in\supp{\apply{q}{\beta_0}}{X}{T}$ by \cref{positive_real_function_value_in_support}.
            $g(x)$ is the pointwise sum of $q$ at $x$ on $X$ with respect to $T$ and $C$
                by \cref{locally_finite_real_sum_function}.
            Take $F_1,m_1,g_1,b_1$ such that
                $F_1\subseteq C$ and
                $F_1$ is finite and
                $F_1$ is inhabited and
                $m_1\in\naturalsPlus$ and
                $g_1$ is a bijection from $\initialSegment{m_1}$ to $F_1$ and
                $b_1$ is a real sum sequence of $q$ at $x$ along $g_1$ with respect to $m_1$ and result $g(x)$ and
                for all $\alpha\in C$ if $x\in\supp{\apply{q}{\alpha}}{X}{T}$ then $\alpha\in F_1$
                by \cref{pointwise_sum_of_indexed_family}.
            $\beta_0\in F_1$ by \cref{subseteq}.
            Take $j_0\in\initialSegment{m_1}$ such that $g_1(j_0) = \beta_0$ by \cref{bijections,surj}.
            $b_1\in\funs{\initialSegment{m_1}}{\reals}$ and $\sumFiniteSequence{b_1}{g(x)}$ and
                for all $i\in\initialSegment{m_1}$ we have $b_1(i) = \apply{\apply{q}{g_1(i)}}{x}$
                by \cref{real_sum_sequence_along_enumeration}.
            Show that for all $i\in\dom{b_1}$ we have $0 \leq b_1(i)$.
            \begin{subproof}
                Fix $i\in\dom{b_1}$.
                $i\in\initialSegment{m_1}$ by \cref{funs,subseteq}.
                $i\in\dom{g_1}$ by \cref{bijections_dom,subseteq}.
                $g_1(i)\in F_1$ by \cref{bijections,surj,funs_type_apply,subseteq}.
                $g_1(i)\in C$ by \cref{subseteq}.
                Take $I_i$ such that
                    $I_i = \intervalclosed{0}{1}$ and
                    $\apply{p}{g_1(i)}$ is continuous from $X$ to $I_i$ with respect to $T$ and $\subspaceTopology{\standardTopologyR}{I_i}$ and
                    $\supp{\apply{p}{g_1(i)}}{X}{T}\subseteq \apply{U}{g_1(i)}$
                    by \cref{indexed_partition_of_unity,pair_eq_iff}.
                $\apply{\apply{p}{g_1(i)}}{x}\in I_i$ by \cref{continuous,funs_type_apply}.
                $\apply{\apply{p}{g_1(i)}}{x}\in\reals$ and $0 \leq \apply{\apply{p}{g_1(i)}}{x}$ by \cref{intervalclosed}.
                $\apply{e}{g_1(i)}\in\reals$ and $0 < \apply{e}{g_1(i)}$ by assumption.
                $\apply{\apply{q}{g_1(i)}}{x} = \apply{e}{g_1(i)} \rmul \apply{\apply{p}{g_1(i)}}{x}$ by \cref{subseteq}.
                $\apply{\apply{p}{g_1(i)}}{x} = 0$ or $0 < \apply{\apply{p}{g_1(i)}}{x}$ by \cref{real_lt_trichotomy}.
                $0 \leq \apply{e}{g_1(i)} \rmul \apply{\apply{p}{g_1(i)}}{x}$
                    by \cref{real_mul_zero,real_mul_comm,real_lt_imp_leq,real_mul_pos_pos_gt}.
                $0 \leq b_1(i)$ by \cref{real_lt_subst_right,real_lt_subst_left}.
            \end{subproof}
            $j_0\in\dom{b_1}$ by \cref{funs,subseteq}.
            $b_1(j_0)\in\reals$ by \cref{funs_type_apply}.
            $b_1(j_0) = \apply{\apply{q}{\beta_0}}{x}$ by \cref{real_sum_sequence_along_enumeration,real_lt_subst_left}.
            $0 < b_1(j_0)$ by \cref{real_lt_subst_right}.
            $b_1(j_0) \leq g(x)$ by \cref{sum_finite_sequence_term_leq}.
            $0\in\reals$ by \cref{real_add_ident}.
            $g(x)\in\reals$ by \cref{continuous,funs_type_apply}.
            $b_1(j_0) < g(x)$ or $b_1(j_0) = g(x)$ by \cref{real_leq_split}.
            \begin{byCase}
            \caseOf{$b_1(j_0) < g(x)$.}
                $0 < g(x)$ by \cref{real_lt_trans}.
            \caseOf{$b_1(j_0) = g(x)$.}
                $0 < g(x)$ by \cref{real_lt_subst_left}.
            \end{byCase}
        \end{subproof}

        Show that for all $A\in C$ for all $x\in A$ we have $g(x) \leq \apply{e}{A}$.
        \begin{subproof}
            Fix $A\in C$.
            Fix $x\in A$.
            $C\subseteq\pow{X}$ by \cref{locally_finite_collection}.
            $A\subseteq X$ by \cref{subseteq,powerset_elim}.
            $x\in X$ by \cref{subseteq}.
            Take $\beta$ such that $\beta\in C$ and $x\in\apply{S_1}{\beta}$ by \cref{subseteq}.
            $x\in\apply{S}{\beta}$ by \cref{real_lt_subst_left}.
            Let $E = \{\alpha\in C \mid x\in\apply{S}{\alpha}\}$.
            $E\subseteq C$ and $E$ is finite and $E$ is inhabited by \cref{local_support_index_set_finite_inhabited}.
            Take $n,h$ such that
                $n\in\naturalsPlus$ and
                $h$ is a bijection from $\initialSegment{n}$ to $E$ and
                $h$ is increasing on $\initialSegment{n}$ with respect to $L$
                by \cref{finite_well_ordered_enumeration}.
            $h$ is a well-ordered enumeration of $E$ in $C$ with respect to $L$ and $n$
                by \cref{well_ordered_enumeration_of_subset}.
            Let $a_0 = \{(i,\apply{\apply{p}{h(i)}}{x}) \mid i\in\initialSegment{n}\}$.
            Let $a_1 = \{(i,\apply{e}{h(i)} \rmul \apply{\apply{p}{h(i)}}{x}) \mid i\in\initialSegment{n}\}$.
            Let $b = \coordScalarSequence{n}{\apply{e}{A}}{a_0}$.
            For all $i,r_1,r_2$ such that $(i,r_1)\in a_0$ and $(i,r_2)\in a_0$ we have $r_1 = r_2$ by \cref{pair_eq_iff}.
            $a_0$ is a function by \cref{relation,rightunique}.
            For all $i$ such that $i\in\dom{a_0}$ we have $i\in\initialSegment{n}$ by \cref{dom_iff,pair_eq_iff}.
            For all $i$ such that $i\in\initialSegment{n}$ we have $i\in\dom{a_0}$ by \cref{pair_eq_iff,dom_intro}.
            $\dom{a_0} = \initialSegment{n}$ by \cref{subseteq,subseteq_antisymmetric}.
            Show that for all $i\in\dom{a_0}$ we have $a_0(i)\in\reals$.
            \begin{subproof}
                Fix $i\in\dom{a_0}$.
                $i\in\initialSegment{n}$ by \cref{real_lt_subst_left}.
                $i\in\dom{h}$ by \cref{well_ordered_enumeration_of_subset,bijections_dom,subseteq}.
                $h(i)\in E$ by \cref{well_ordered_enumeration_of_subset,bijections,surj,funs_type_apply,subseteq}.
                $h(i)\in C$ by \cref{subseteq}.
                Take $I_i$ such that
                    $I_i = \intervalclosed{0}{1}$ and
                    $\apply{p}{h(i)}$ is continuous from $X$ to $I_i$ with respect to $T$ and $\subspaceTopology{\standardTopologyR}{I_i}$ and
                    $\supp{\apply{p}{h(i)}}{X}{T}\subseteq \apply{U}{h(i)}$
                    by \cref{indexed_partition_of_unity,pair_eq_iff}.
                $a_0(i)\in\reals$ by \cref{dom_iff,pair_eq_iff,function_apply_intro,continuous,funs_type_apply,intervalclosed,subseteq}.
            \end{subproof}
            $a_0\in\funs{\initialSegment{n}}{\reals}$ by \cref{funs_intro}.
            For all $i\in\initialSegment{n}$ we have $a_0(i) = \apply{\apply{p}{h(i)}}{x}$ by \cref{pair_eq_iff,function_apply_intro}.
            Take $s_0\in\reals$ such that $\sumFiniteSequence{a_0}{s_0}$ by \cref{sum_finite_sequence_exists}.
            $a_0$ is a real sum sequence of $p$ at $x$ along $h$ with respect to $n$ and result $s_0$
                by \cref{real_sum_sequence_along_enumeration}.
            $1$ is the pointwise sum of $p$ at $x$ on $X$ with respect to $T$ and $C$ by \cref{indexed_partition_of_unity}.
            Take $F_0,m_0,g_0,c_0$ such that
                $F_0\subseteq C$ and
                $F_0$ is finite and
                $F_0$ is inhabited and
                $m_0\in\naturalsPlus$ and
                $g_0$ is a bijection from $\initialSegment{m_0}$ to $F_0$ and
                $c_0$ is a real sum sequence of $p$ at $x$ along $g_0$ with respect to $m_0$ and result $1$ and
                for all $\alpha\in C$ if $x\in\supp{\apply{p}{\alpha}}{X}{T}$ then $\alpha\in F_0$
                by \cref{pointwise_sum_of_indexed_family}.
            Take $m_1,h_0$ such that
                $m_1\in\naturalsPlus$ and
                $h_0$ is a bijection from $\initialSegment{m_1}$ to $F_0$ and
                $h_0$ is increasing on $\initialSegment{m_1}$ with respect to $L$
                by \cref{subseteq,nonempty_inhabited,finite_well_ordered_enumeration}.
            $h_0$ is a well-ordered enumeration of $F_0$ in $C$ with respect to $L$ and $m_1$
                by \cref{well_ordered_enumeration_of_subset}.
            Let $d_0 = \{(i,\apply{\apply{p}{h_0(i)}}{x}) \mid i\in\initialSegment{m_1}\}$.
            For all $i,r_1,r_2$ such that $(i,r_1)\in d_0$ and $(i,r_2)\in d_0$ we have $r_1 = r_2$ by \cref{pair_eq_iff}.
            $d_0$ is a function by \cref{relation,rightunique}.
            For all $i$ such that $i\in\dom{d_0}$ we have $i\in\initialSegment{m_1}$ by \cref{dom_iff,pair_eq_iff}.
            For all $i$ such that $i\in\initialSegment{m_1}$ we have $i\in\dom{d_0}$ by \cref{pair_eq_iff,dom_intro}.
            $\dom{d_0} = \initialSegment{m_1}$ by \cref{subseteq,subseteq_antisymmetric}.
            Show that for all $i\in\dom{d_0}$ we have $d_0(i)\in\reals$.
            \begin{subproof}
                Fix $i\in\dom{d_0}$.
                $i\in\initialSegment{m_1}$ by \cref{real_lt_subst_left}.
                $i\in\dom{h_0}$ by \cref{well_ordered_enumeration_of_subset,bijections_dom,subseteq}.
                $h_0(i)\in F_0$ by \cref{well_ordered_enumeration_of_subset,bijections,surj,funs_type_apply,subseteq}.
                $h_0(i)\in C$ by \cref{subseteq}.
                Take $I_i$ such that
                    $I_i = \intervalclosed{0}{1}$ and
                    $\apply{p}{h_0(i)}$ is continuous from $X$ to $I_i$ with respect to $T$ and $\subspaceTopology{\standardTopologyR}{I_i}$ and
                    $\supp{\apply{p}{h_0(i)}}{X}{T}\subseteq \apply{U}{h_0(i)}$
                    by \cref{indexed_partition_of_unity,pair_eq_iff}.
                $d_0(i)\in\reals$ by \cref{dom_iff,pair_eq_iff,function_apply_intro,continuous,funs_type_apply,intervalclosed,subseteq}.
            \end{subproof}
            $d_0\in\funs{\initialSegment{m_1}}{\reals}$ by \cref{funs_intro}.
            For all $i\in\initialSegment{m_1}$ we have $d_0(i) = \apply{\apply{p}{h_0(i)}}{x}$ by \cref{pair_eq_iff,function_apply_intro}.
            Take $t_0\in\reals$ such that $\sumFiniteSequence{d_0}{t_0}$ by \cref{sum_finite_sequence_exists}.
            $d_0$ is a real sum sequence of $p$ at $x$ along $h_0$ with respect to $m_1$ and result $t_0$
                by \cref{real_sum_sequence_along_enumeration}.
            $t_0 = 1$ by \cref{well_order_relation,simple_order_relation,well_ordered_same_set_sum_comparison}.
            $d_0$ is a real sum sequence of $p$ at $x$ along $h_0$ with respect to $m_1$ and result $1$
                by \cref{real_lt_subst_right}.
            Show that for all $\alpha\in F_0\setminus E$ we have $\apply{\apply{p}{\alpha}}{x} = 0$.
            \begin{subproof}
                Fix $\alpha\in F_0\setminus E$.
                $\alpha\in C$ and $\alpha\notin E$ by \cref{setminus,subseteq}.
                $x\notin\apply{S}{\alpha}$ by definition.
                $x\notin\supp{\apply{p}{\alpha}}{X}{T}$ by \cref{support_family,pair_eq_iff,function_apply_intro,real_lt_subst_right}.
                Take $I_a$ such that
                    $I_a = \intervalclosed{0}{1}$ and
                    $\apply{p}{\alpha}$ is continuous from $X$ to $I_a$ with respect to $T$ and $\subspaceTopology{\standardTopologyR}{I_a}$ and
                    $\supp{\apply{p}{\alpha}}{X}{T}\subseteq \apply{U}{\alpha}$
                    by \cref{indexed_partition_of_unity,pair_eq_iff}.
                Let $j_a = \identity{I_a}$.
                Let $p_a = j_a\circ \apply{p}{\alpha}$.
                $p_a$ is continuous from $X$ to $\reals$ with respect to $T$ and $\standardTopologyR$
                    by \cref{continuous_interval_expand_range}.
                $p_a\in\funs{X}{\reals}$ by \cref{continuous,funs}.
                $\apply{p}{\alpha}\in\funs{X}{I_a}$ by \cref{continuous,funs}.
                $I_a\subseteq\reals$ by \cref{intervalclosed,subseteq}.
                $\apply{p}{\alpha}\in\funs{X}{\reals}$ by \cref{funs_weaken_codom,subseteq}.
                For all $y\in X$ we have $p_a(y) = \apply{\apply{p}{\alpha}}{y}$
                    by \cref{continuous,funs_is_function,funs,function_rightunique,funs_is_relation,funs_subseteq_rels,subseteq,rels_ran_subseteq,id_dom,id_rightunique,id_is_relation,circ_apply,funs_type_apply,id_apply}.
                $p_a = \apply{p}{\alpha}$ by \cref{functions_eq_iff}.
                $\apply{p}{\alpha}$ is continuous from $X$ to $\reals$ with respect to $T$ and $\standardTopologyR$
                    by \cref{continuous_eq_subst}.
                $x\in X\setminus \supp{\apply{p}{\alpha}}{X}{T}$ by \cref{setminus_intro}.
                $\apply{\apply{p}{\alpha}}{x} = 0$ by \cref{outside_support_zero}.
            \end{subproof}
            Show that for all $\alpha\in E$ we have $\alpha\in F_0$.
            \begin{subproof}
                Fix $\alpha\in E$.
                $\alpha\in C$ and $x\in\apply{S}{\alpha}$ by definition.
                $x\in\supp{\apply{p}{\alpha}}{X}{T}$ by \cref{support_family,pair_eq_iff,function_apply_intro,real_lt_subst_right}.
                $\alpha\in F_0$ by \cref{subseteq}.
            \end{subproof}
            $E\subseteq F_0$ by \cref{subseteq}.
            $s_0 = 1$ by \cref{well_order_relation,simple_order_relation,well_ordered_zero_extension_sum_comparison}.
            $b\in\funs{\initialSegment{n}}{\reals}$ by \cref{coord_scalar_sequence_function}.
            $\sumFiniteSequence{b}{\apply{e}{A} \rmul s_0}$ by \cref{sum_finite_sequence_scalar}.
            $\apply{e}{A}\in\reals$ by assumption.
            $\apply{e}{A} \rmul s_0 = \apply{e}{A} \rmul 1$ by \cref{real_mul_eq_subst}.
            $\apply{e}{A} \rmul s_0 = 1 \rmul \apply{e}{A}$ by \cref{real_mul_comm,real_mul_eq_subst}.
            $\apply{e}{A} \rmul s_0 = \apply{e}{A}$ by \cref{real_mul_ident}.
            $\sumFiniteSequence{b}{\apply{e}{A}}$ by \cref{real_lt_subst_right}.
            For all $i,r_1,r_2$ such that $(i,r_1)\in a_1$ and $(i,r_2)\in a_1$ we have $r_1 = r_2$ by \cref{pair_eq_iff}.
            $a_1$ is a function by \cref{relation,rightunique}.
            For all $i$ such that $i\in\dom{a_1}$ we have $i\in\initialSegment{n}$ by \cref{dom_iff,pair_eq_iff}.
            For all $i$ such that $i\in\initialSegment{n}$ we have $i\in\dom{a_1}$ by \cref{pair_eq_iff,dom_intro}.
            $\dom{a_1} = \initialSegment{n}$ by \cref{subseteq,subseteq_antisymmetric}.
            Show that for all $i\in\dom{a_1}$ we have $a_1(i)\in\reals$.
            \begin{subproof}
                Fix $i\in\dom{a_1}$.
                $i\in\initialSegment{n}$ by \cref{real_lt_subst_left}.
                $i\in\dom{h}$ by \cref{well_ordered_enumeration_of_subset,bijections_dom,subseteq}.
                $h(i)\in E$ by \cref{well_ordered_enumeration_of_subset,bijections,surj,funs_type_apply,subseteq}.
                $h(i)\in C$ by \cref{subseteq}.
                $\apply{e}{h(i)}\in\reals$ by assumption.
                $\apply{\apply{p}{h(i)}}{x}\in\reals$ by \cref{funs_type_apply,subseteq}.
                $a_1(i)\in\reals$ by \cref{dom_iff,pair_eq_iff,function_apply_intro,real_mul_closed}.
            \end{subproof}
            $a_1\in\funs{\initialSegment{n}}{\reals}$ by \cref{funs_intro}.
            Show that for all $i\in\initialSegment{n}$ we have $a_1(i) = \apply{\apply{q}{h(i)}}{x}$.
            \begin{subproof}
                Fix $i\in\initialSegment{n}$.
                $i\in\dom{h}$ by \cref{well_ordered_enumeration_of_subset,bijections_dom,subseteq}.
                $h(i)\in E$ by \cref{well_ordered_enumeration_of_subset,bijections,surj,funs_type_apply,subseteq}.
                $h(i)\in C$ by \cref{subseteq}.
                $a_1(i) = \apply{e}{h(i)} \rmul \apply{\apply{p}{h(i)}}{x}$ by \cref{pair_eq_iff,function_apply_intro}.
                $\apply{\apply{q}{h(i)}}{x} = \apply{e}{h(i)} \rmul \apply{\apply{p}{h(i)}}{x}$ by \cref{subseteq}.
                $a_1(i) = \apply{\apply{q}{h(i)}}{x}$ by \cref{real_lt_subst_right}.
            \end{subproof}
            Take $s_1\in\reals$ such that $\sumFiniteSequence{a_1}{s_1}$ by \cref{sum_finite_sequence_exists}.
            $a_1$ is a real sum sequence of $q$ at $x$ along $h$ with respect to $n$ and result $s_1$
                by \cref{real_sum_sequence_along_enumeration}.
            $g(x)$ is the pointwise sum of $q$ at $x$ on $X$ with respect to $T$ and $C$
                by \cref{locally_finite_real_sum_function}.
            Take $F_1,m_2,g_1,c_1$ such that
                $F_1\subseteq C$ and
                $F_1$ is finite and
                $F_1$ is inhabited and
                $m_2\in\naturalsPlus$ and
                $g_1$ is a bijection from $\initialSegment{m_2}$ to $F_1$ and
                $c_1$ is a real sum sequence of $q$ at $x$ along $g_1$ with respect to $m_2$ and result $g(x)$ and
                for all $\alpha\in C$ if $x\in\supp{\apply{q}{\alpha}}{X}{T}$ then $\alpha\in F_1$
                by \cref{pointwise_sum_of_indexed_family}.
            Take $m_3,h_1$ such that
                $m_3\in\naturalsPlus$ and
                $h_1$ is a bijection from $\initialSegment{m_3}$ to $F_1$ and
                $h_1$ is increasing on $\initialSegment{m_3}$ with respect to $L$
                by \cref{subseteq,nonempty_inhabited,finite_well_ordered_enumeration}.
            $h_1$ is a well-ordered enumeration of $F_1$ in $C$ with respect to $L$ and $m_3$
                by \cref{well_ordered_enumeration_of_subset}.
            Let $d_1 = \{(i,\apply{\apply{q}{h_1(i)}}{x}) \mid i\in\initialSegment{m_3}\}$.
            For all $i,r_1,r_2$ such that $(i,r_1)\in d_1$ and $(i,r_2)\in d_1$ we have $r_1 = r_2$ by \cref{pair_eq_iff}.
            $d_1$ is a function by \cref{relation,rightunique}.
            For all $i$ such that $i\in\dom{d_1}$ we have $i\in\initialSegment{m_3}$ by \cref{dom_iff,pair_eq_iff}.
            For all $i$ such that $i\in\initialSegment{m_3}$ we have $i\in\dom{d_1}$ by \cref{pair_eq_iff,dom_intro}.
            $\dom{d_1} = \initialSegment{m_3}$ by \cref{subseteq,subseteq_antisymmetric}.
            Show that for all $i\in\dom{d_1}$ we have $d_1(i)\in\reals$.
            \begin{subproof}
                Fix $i\in\dom{d_1}$.
                $i\in\initialSegment{m_3}$ by \cref{real_lt_subst_left}.
                $i\in\dom{h_1}$ by \cref{well_ordered_enumeration_of_subset,bijections_dom,subseteq}.
                $h_1(i)\in F_1$ by \cref{well_ordered_enumeration_of_subset,bijections,surj,funs_type_apply,subseteq}.
                $h_1(i)\in C$ by \cref{subseteq}.
                $d_1(i)\in\reals$ by \cref{dom_iff,pair_eq_iff,function_apply_intro,subseteq,continuous,funs_type_apply}.
            \end{subproof}
            $d_1\in\funs{\initialSegment{m_3}}{\reals}$ by \cref{funs_intro}.
            For all $i\in\initialSegment{m_3}$ we have $d_1(i) = \apply{\apply{q}{h_1(i)}}{x}$ by \cref{pair_eq_iff,function_apply_intro}.
            Take $t_1\in\reals$ such that $\sumFiniteSequence{d_1}{t_1}$ by \cref{sum_finite_sequence_exists}.
            $d_1$ is a real sum sequence of $q$ at $x$ along $h_1$ with respect to $m_3$ and result $t_1$
                by \cref{real_sum_sequence_along_enumeration}.
            $t_1 = g(x)$ by \cref{well_order_relation,simple_order_relation,well_ordered_same_set_sum_comparison}.
            $d_1$ is a real sum sequence of $q$ at $x$ along $h_1$ with respect to $m_3$ and result $g(x)$
                by \cref{real_lt_subst_right}.
            Show that for all $\alpha\in F_1\setminus E$ we have $\apply{\apply{q}{\alpha}}{x} = 0$.
            \begin{subproof}
                Fix $\alpha\in F_1\setminus E$.
                $\alpha\in C$ and $\alpha\notin E$ by \cref{setminus,subseteq}.
                $x\notin\apply{S}{\alpha}$ by definition.
                $x\notin\supp{\apply{q}{\alpha}}{X}{T}$ by \cref{real_lt_subst_left,support_family,pair_eq_iff,function_apply_intro,real_lt_subst_right}.
                $\apply{q}{\alpha}$ is continuous from $X$ to $\reals$ with respect to $T$ and $\standardTopologyR$ by \cref{subseteq}.
                $x\in X\setminus \supp{\apply{q}{\alpha}}{X}{T}$ by \cref{setminus_intro}.
                $\apply{\apply{q}{\alpha}}{x} = 0$ by \cref{outside_support_zero}.
            \end{subproof}
            Show that for all $\alpha\in E$ we have $\alpha\in F_1$.
            \begin{subproof}
                Fix $\alpha\in E$.
                $\alpha\in C$ and $x\in\apply{S}{\alpha}$ by definition.
                $x\in\apply{S_1}{\alpha}$ by \cref{real_lt_subst_right}.
                $x\in\supp{\apply{q}{\alpha}}{X}{T}$ by \cref{support_family,pair_eq_iff,function_apply_intro,real_lt_subst_right}.
                $\alpha\in F_1$ by \cref{subseteq}.
            \end{subproof}
            $E\subseteq F_1$ by \cref{subseteq}.
            $s_1 = g(x)$ by \cref{well_order_relation,simple_order_relation,well_ordered_zero_extension_sum_comparison}.
            Show that for all $i\in\initialSegment{n}$ we have $a_1(i) \leq b(i)$.
            \begin{subproof}
                Fix $i\in\initialSegment{n}$.
                $i\in\dom{h}$ by \cref{well_ordered_enumeration_of_subset,bijections_dom,subseteq}.
                $h(i)\in E$ by \cref{well_ordered_enumeration_of_subset,bijections,surj,funs_type_apply,subseteq}.
                $x\in\apply{S}{h(i)}$ by definition.
                $x\in\supp{\apply{p}{h(i)}}{X}{T}$ by \cref{support_family,pair_eq_iff,function_apply_intro,real_lt_subst_right}.
                $x\in\apply{U}{h(i)}$ by \cref{indexed_partition_of_unity,pair_eq_iff,subseteq}.
                Take $C_i,H_i$ such that
                    $C_i = \{A\in C \mid \apply{e}{A} < \apply{e}{h(i)}\}$ and
                    $H_i$ is the closure collection of $C_i$ in $X$ with respect to $T$ and
                    $\apply{U}{h(i)} = X\setminus \unions{H_i}$
                    by \cref{choice_function,pair_eq_iff}.
                Show that $\apply{e}{h(i)} \leq \apply{e}{A}$.
                \begin{subproof}
                    Suppose not.
                    $\apply{e}{A} < \apply{e}{h(i)}$ by \cref{real_lt_trichotomy}.
                    $A\in C_i$ by \cref{subseteq}.
                    $A\in C$ by \cref{subseteq}.
                    $A\subseteq X$ by \cref{subseteq,powerset_elim}.
                    $X$ is closed in $X$ with respect to $T$ by \cref{closed_whole_space}.
                    $\closure{A}{X}{T}\subseteq X$ by \cref{closure_subseteq_closed}.
                    $\closure{A}{X}{T}\in\pow{X}$ by \cref{powerset_intro}.
                    $\closure{A}{X}{T}\in H_i$ by \cref{closure_collection_in,subseteq}.
                    $x\in\unions{H_i}$
                        by \cref{subseteq_closure,subseteq,unions_iff}.
                    $x\notin\apply{U}{h(i)}$ by \cref{pair_eq_iff,function_apply_intro,setminus,setminus_elim_right}.
                    Contradiction.
                \end{subproof}
                Take $I_i$ such that
                    $I_i = \intervalclosed{0}{1}$ and
                    $\apply{p}{h(i)}$ is continuous from $X$ to $I_i$ with respect to $T$ and $\subspaceTopology{\standardTopologyR}{I_i}$ and
                    $\supp{\apply{p}{h(i)}}{X}{T}\subseteq \apply{U}{h(i)}$
                    by \cref{indexed_partition_of_unity,pair_eq_iff}.
                $a_0(i) = \apply{\apply{p}{h(i)}}{x}$ by \cref{pair_eq_iff,function_apply_intro}.
                $\apply{\apply{p}{h(i)}}{x}\in I_i$ by \cref{continuous,funs_type_apply}.
                $0 \leq \apply{\apply{p}{h(i)}}{x}$ by \cref{intervalclosed}.
                $0 \leq a_0(i)$ by \cref{real_leq_subst_right_local}.
                $a_0(i) = 0$ or $0 < a_0(i)$ by \cref{real_leq_split}.
                \begin{byCase}
                \caseOf{$a_0(i) = 0$.}
                    $a_0(i)\in\reals$ by \cref{funs_type_apply,subseteq}.
                    $a_1(i) = \apply{e}{h(i)} \rmul a_0(i)$
                        by \cref{pair_eq_iff,function_apply_intro,real_mul_eq_subst}.
                    $a_1(i) = \apply{e}{h(i)} \rmul 0$ by \cref{real_mul_eq_subst}.
                    $a_1(i) = 0 \rmul \apply{e}{h(i)}$ by \cref{real_mul_comm,real_lt_subst_right}.
                    $a_1(i) = 0$ by \cref{real_mul_zero}.
                    $b(i) = \apply{e}{A} \rmul a_0(i)$ by \cref{coord_scalar_sequence,funs_elim,function_apply_intro}.
                    $b(i) = \apply{e}{A} \rmul 0$ by \cref{real_mul_eq_subst}.
                    $b(i) = 0 \rmul \apply{e}{A}$ by \cref{real_mul_comm,real_lt_subst_right}.
                    $b(i) = 0$ by \cref{real_mul_zero}.
                    $0 \leq 0$.
                    $a_1(i) \leq b(i)$ by \cref{real_leq_subst_left_local,real_leq_subst_right_local}.
                \caseOf{$0 < a_0(i)$.}
                    $\apply{e}{h(i)} < \apply{e}{A}$ or $\apply{e}{h(i)} = \apply{e}{A}$ by \cref{real_leq_split}.
                    \begin{byCase}
                    \caseOf{$\apply{e}{h(i)} < \apply{e}{A}$.}
                        $\apply{e}{h(i)} \rmul a_0(i) < \apply{e}{A} \rmul a_0(i)$ by \cref{real_lt_mul_pos}.
                        $a_1(i) < \apply{e}{A} \rmul a_0(i)$ by \cref{pair_eq_iff,function_apply_intro,real_lt_subst_left}.
                        $b(i) = \apply{e}{A} \rmul a_0(i)$ by \cref{coord_scalar_sequence,funs_elim,function_apply_intro}.
                        $a_1(i) \leq b(i)$ by \cref{real_lt_imp_leq,real_lt_subst_right}.
                    \caseOf{$\apply{e}{h(i)} = \apply{e}{A}$.}
                        $a_1(i) = \apply{e}{A} \rmul a_0(i)$ by \cref{pair_eq_iff,function_apply_intro,real_mul_eq_subst,real_lt_subst_right}.
                        $b(i) = \apply{e}{A} \rmul a_0(i)$ by \cref{coord_scalar_sequence,funs_elim,function_apply_intro}.
                        $a_1(i) \leq b(i)$ by \cref{real_lt_subst_right}.
                    \end{byCase}
                \end{byCase}
            \end{subproof}
            $s_1 < \apply{e}{A}$ or $s_1 = \apply{e}{A}$ by \cref{sum_finite_sequence_coordwise_leq}.
            \begin{byCase}
            \caseOf{$s_1 < \apply{e}{A}$.}
                $g(x) < \apply{e}{A}$ by \cref{real_lt_subst_left}.
                $g(x) \leq \apply{e}{A}$ by \cref{real_lt_imp_leq}.
            \caseOf{$s_1 = \apply{e}{A}$.}
                $g(x) = \apply{e}{A}$ by \cref{real_lt_subst_left}.
                $g(x) \leq \apply{e}{A}$ by \cref{real_lt_subst_right}.
            \end{byCase}
        \end{subproof}
        Let $f_0 = g$.
        Show that $f_0$ is continuous from $X$ to $\reals$ with respect to $T$ and $\standardTopologyR$.
        \begin{subproof}
            Follows by \cref{real_lt_subst_left}.
        \end{subproof}
        Show that for all $x\in X$ we have $0 < f_0(x)$.
        \begin{subproof}
            Follows by \cref{real_lt_subst_left,real_lt_subst_right}.
        \end{subproof}
        Show that for all $A\in C$ for all $x\in A$ we have $f_0(x) \leq \apply{e}{A}$.
        \begin{subproof}
            Follows by \cref{real_lt_subst_left,real_lt_subst_right}.
        \end{subproof}
        There exists $f$ such that
            $f$ is continuous from $X$ to $\reals$ with respect to $T$ and $\standardTopologyR$ and
            for all $x\in X$ we have $0 < f(x)$ and
            for all $A\in C$ for all $x\in A$ we have $f(x) \leq \apply{e}{A}$
            by \cref{real_lt_subst_left}.
    \end{byCase}
\end{proof}

% from §41 helper definition 6: locally metrizable space
\begin{definition}\label{locally_metrizable}
    $X$ is locally metrizable with respect to $T$ iff
        $T$ is a topology on $X$ and
        for all $x\in X$ there exists $U$ such that
            $U$ is a neighborhood of $x$ in $X$ with respect to $T$ and
            $U$ is metrizable with respect to $\subspaceTopology{T}{U}$.
\end{definition}

% from §41 helper lemma 1: basis from paracompact Hausdorff local metrizability
\begin{lemma}\label{paracompact_hausdorff_locally_metrizable_basis}
    Assume $X$ is paracompact with respect to $T$.
    Assume $X$ is Hausdorff with respect to $T$.
    Assume $X$ is locally metrizable with respect to $T$.
    Then there exists $B$ such that
        $B$ is a basis for $T$ on $X$ and
        $B$ is countably locally finite in $X$ with respect to $T$.
\end{lemma}
\begin{proof}
    $T$ is a topology on $X$ by \cref{locally_metrizable}.
    Let $A = \{U\in\pow{X} \mid \text{$U$ is open in $X$ with respect to $T$ and $U$ is metrizable with respect to $\subspaceTopology{T}{U}$}\}$.
    For all $x\in X$ there exists $U$ such that
        $U$ is a neighborhood of $x$ in $X$ with respect to $T$ and
        $U$ is metrizable with respect to $\subspaceTopology{T}{U}$
        by \cref{locally_metrizable}.
    For all $x\in X$ we have $x\in\unions{A}$ by \cref{neighborhood,open_in,topology_on,subseteq,pow_iff,powerset_intro,unions_iff}.
    $A$ is a covering of $X$ by \cref{subseteq,covering}.
    $A$ is an open covering of $X$ with respect to $T$ by \cref{subseteq,open_covering}.
    Take $M$ such that
        $M$ is an open refinement of $A$ on $X$ with respect to $T$ and
        $M$ is locally finite in $X$ with respect to $T$ and
        $M$ is a covering of $X$
        by \cref{paracompact_space}.

    Let $P = \{w\in M\times\pow{((X\times X)\times\reals)} \mid
        \text{$\snd{w}$ is a metric on $\fst{w}$ and $\subspaceTopology{T}{\fst{w}} = \metricTopology{\snd{w}}{\fst{w}}$}\}$.
    For all $w$ such that $w\in P$ there exist $C,d$ such that $w = (C,d)$.
    $P$ is a relation by \cref{relation}.
    Show that for all $C\in M$ there exists $d$ such that $C\mathrel{P} d$.
    \begin{subproof}
        Fix $C\in M$.
        Take $U\in A$ such that $C\subseteq U$ by \cref{open_refinement_on,refinement_on}.
        $U$ is open in $X$ with respect to $T$ and $U$ is metrizable with respect to $\subspaceTopology{T}{U}$ by \cref{subseteq}.
        Let $T_U = \subspaceTopology{T}{U}$.
        $U$ is a subspace of $X$ with respect to $T$ by \cref{open_in,topology_on,subseteq,pow_iff,subspace_of}.
        $T_U$ is a topology on $U$ by \cref{subspace_of,subspace_topology_is_topology}.
        $C$ is a subspace of $U$ with respect to $T_U$ by \cref{subseteq,subspace_of}.
        Hence $C$ is metrizable with respect to $\subspaceTopology{T_U}{C}$ by \cref{subspace_metrizable}.
        We have $\subspaceTopology{T_U}{C} = \subspaceTopology{T}{C}$ by \cref{subspace_of,subspace_subspace_topology,subseteq}.
        Take $d$ such that $d$ is a metric on $C$ and $\subspaceTopology{T}{C} = \metricTopology{d}{C}$ by \cref{real_lt_subst_left,metrizable}.
        $d\in\funs{C\times C}{\reals}$ and $d$ is a relation by \cref{metric_on,funs_is_relation}.
        $d\subseteq (C\times C)\times\reals$ by \cref{funs_subseteq_rels,subseteq,rels_elim}.
        $C\times C\subseteq X\times X$ by \cref{open_refinement_on,refinement_on,subseteq,powerset_elim,times_subseteq_left,times_subseteq_right,subseteq_transitive}.
        Thus $(C\times C)\times\reals\subseteq (X\times X)\times\reals$ by \cref{times_subseteq_left}.
        $d\subseteq (X\times X)\times\reals$ by \cref{subseteq_transitive}.
        Hence $d\in\pow{((X\times X)\times\reals)}$ by \cref{powerset_intro}.
        $(C,d)\in M\times\pow{((X\times X)\times\reals)}$ by \cref{times_tuple_intro}.
        Thus $C\mathrel{P} d$ by \cref{times_tuple_intro,fst_eq,snd_eq}.
    \end{subproof}
    Take $e$ such that $e$ is a function on $M$ and for all $C\in M$ we have $C\mathrel{P}\apply{e}{C}$ by \cref{choice_function}.

    Let $Q = \{w\in\naturalsPlus\times\pow{\pow{X}} \mid
        \text{$\snd{w} = \{B\in\pow{X} \mid \exists C\in M. \exists x\in C. B = \metricBall{\apply{e}{C}}{C}{x}{1 / (\fst{w})}\}$}\}$.
    For all $w$ such that $w\in Q$ there exist $m,D$ such that $w = (m,D)$.
    $Q$ is a relation by \cref{relation}.
    Show that for all $m\in\naturalsPlus$ there exists $D$ such that $m\mathrel{Q} D$.
    \begin{subproof}
        Fix $m\in\naturalsPlus$.
        Let $D = \{B\in\pow{X} \mid \exists C\in M. \exists x\in C. B = \metricBall{\apply{e}{C}}{C}{x}{1 / m}\}$.
        $m\mathrel{Q} D$ by \cref{subseteq,powerset_intro,times_tuple_intro,fst_eq,snd_eq}.
    \end{subproof}
    Take $a$ such that $a$ is a function on $\naturalsPlus$ and for all $m\in\naturalsPlus$ we have $m\mathrel{Q}\apply{a}{m}$ by \cref{choice_function}.

    Show that for all $m\in\naturalsPlus$ we have $\apply{a}{m}$ is a covering of $X$.
    \begin{subproof}
        Fix $m\in\naturalsPlus$.
        We have $m\mathrel{Q}\apply{a}{m}$ by \cref{choice_function}.
        Hence $\apply{a}{m} = \{B\in\pow{X} \mid \exists C\in M. \exists x\in C. B = \metricBall{\apply{e}{C}}{C}{x}{1 / m}\}$
            by \cref{relation,fst_eq,snd_eq}.
        Show that for all $x\in X$ we have $x\in\unions{\apply{a}{m}}$.
        \begin{subproof}
            Fix $x\in X$.
            Take $C\in M$ such that $x\in C$ by \cref{covering,subseteq,unions_iff}.
            We have $C\mathrel{P}\apply{e}{C}$ by \cref{choice_function}.
            Hence $\apply{e}{C}$ is a metric on $C$ by \cref{relation,fst_eq,snd_eq}.
            $0 < 1 / m$ by \cref{positive_reciprocal_naturalsplus}.
            Let $B = \metricBall{\apply{e}{C}}{C}{x}{1 / m}$.
            $B\in\pow{X}$ by \cref{open_refinement_on,refinement_on,subseteq,powerset_elim,metric_ball,powerset_intro}.
            $B\in\apply{a}{m}$ by \cref{real_lt_subst_left}.
            Hence $x\in B$ by \cref{metric_ball,metric_on,metric_zero_characterization,real_lt_subst_left}.
            $x\in\unions{\apply{a}{m}}$ by \cref{unions_iff}.
        \end{subproof}
        $\apply{a}{m}$ is a covering of $X$ by \cref{subseteq,covering}.
    \end{subproof}
    Show that for all $m\in\naturalsPlus$ for all $B\in\apply{a}{m}$ we have $B$ is open in $X$ with respect to $T$.
    \begin{subproof}
        Fix $m\in\naturalsPlus$.
        We have $m\mathrel{Q}\apply{a}{m}$ by \cref{choice_function}.
        Hence $\apply{a}{m} = \{B\in\pow{X} \mid \exists C\in M. \exists x\in C. B = \metricBall{\apply{e}{C}}{C}{x}{1 / m}\}$
            by \cref{relation,fst_eq,snd_eq}.
        Fix $B\in\apply{a}{m}$.
        Take $C\in M$ such that there exists $x\in C$ such that $B = \metricBall{\apply{e}{C}}{C}{x}{1 / m}$ by \cref{real_lt_subst_left}.
        Take $x\in C$ such that $B = \metricBall{\apply{e}{C}}{C}{x}{1 / m}$ by \cref{real_lt_subst_left}.
        We have $C\mathrel{P}\apply{e}{C}$ by \cref{choice_function}.
        Hence $\apply{e}{C}$ is a metric on $C$ and $\subspaceTopology{T}{C} = \metricTopology{\apply{e}{C}}{C}$
            by \cref{relation,fst_eq,snd_eq}.
        $B\in\metricBasis{\apply{e}{C}}{C}$
            by \cref{positive_reciprocal_naturalsplus,naturalsplus_reciprocal_real,metric_ball,subseteq,powerset_intro,metric_basis}.
        $\metricTopology{\apply{e}{C}}{C}$ is a topology on $C$
            by \cref{metric_basis_is_basis,metric_topology,basis_topology_is_topology}.
        Thus $B\in\metricTopology{\apply{e}{C}}{C}$
            by \cref{metric_basis_is_basis,basis_elem_open_in_generated,metric_topology}.
        $B$ is open in $C$ with respect to $\metricTopology{\apply{e}{C}}{C}$ by \cref{open_in}.
        Therefore $B$ is open in $C$ with respect to $\subspaceTopology{T}{C}$ by \cref{real_lt_subst_right,open_in}.
        $C$ is open in $X$ with respect to $T$ by \cref{open_refinement_on}.
        $C$ is a subspace of $X$ with respect to $T$ by \cref{open_refinement_on,refinement_on,subseteq,powerset_elim,subspace_of}.
        $B$ is open in $X$ with respect to $T$ by \cref{open_in_subspace_open_in_space}.
    \end{subproof}
    Show that for all $m\in\naturalsPlus$ we have $\apply{a}{m}$ is an open covering of $X$ with respect to $T$.
    \begin{subproof}
        Follows by \cref{open_covering}.
    \end{subproof}

    Let $R = \{w\in\naturalsPlus\times\pow{\pow{X}} \mid
        \text{$\snd{w}$ is an open refinement of $\apply{a}{\fst{w}}$ on $X$ with respect to $T$ and
        $\snd{w}$ is locally finite in $X$ with respect to $T$ and
        $\snd{w}$ is a covering of $X$}\}$.
    For all $w$ such that $w\in R$ there exist $m,D$ such that $w = (m,D)$.
    $R$ is a relation by \cref{relation}.
    Show that for all $m\in\naturalsPlus$ there exists $D$ such that $m\mathrel{R} D$.
    \begin{subproof}
        Follows by \cref{paracompact_space,open_refinement_on,refinement_on,subseteq,powerset_intro,times_tuple_intro,fst_eq,snd_eq}.
    \end{subproof}
    Take $d$ such that $d$ is a function on $\naturalsPlus$ and for all $m\in\naturalsPlus$ we have $m\mathrel{R}\apply{d}{m}$ by \cref{choice_function}.
    Let $B = \unions{\img{d}{\naturalsPlus}}$.

    Show that for all $U$ such that $U\in B$ we have $U\in\pow{X}$.
    \begin{subproof}
        Follows by \cref{unions_iff,img_iff,function_apply_iff,choice_function,relation,fst_eq,snd_eq,open_refinement_on,refinement_on,real_lt_subst_left,subseteq}.
    \end{subproof}
    $B\subseteq\pow{X}$ by \cref{subseteq}.
    For all $m\in\naturalsPlus$ we have $\apply{d}{m}$ is locally finite in $X$ with respect to $T$
        by \cref{choice_function,relation,fst_eq,snd_eq}.
    Hence $B$ is countably locally finite in $X$ with respect to $T$ by \cref{countably_locally_finite_collection}.

    Show that for all $D\in B$ we have $D$ is open in $X$ with respect to $T$.
    \begin{subproof}
        Follows by \cref{unions_iff,img_iff,function_apply_iff,choice_function,relation,fst_eq,snd_eq,open_refinement_on,real_lt_subst_left}.
    \end{subproof}
    Show that for all $U,x$ such that $U\in T$ and $x\in U$ there exists $D\in B$ such that $x\in D$ and $D\subseteq U$.
    \begin{subproof}
        Fix $U$.
        Fix $x$ such that $U\in T$ and $x\in U$.
        $x\in X$ by \cref{topology_on,subseteq,pow_iff}.
        Take $W,F$ such that
            $W$ is a neighborhood of $x$ in $X$ with respect to $T$ and
            $F\subseteq M$ and
            $F$ is finite and
            for all $C\in M$ if $W$ intersects $C$ then $C\in F$
            by \cref{locally_finite_collection}.
        Let $G = \{C\in F \mid x\in C\}$.
        $G$ is finite by \cref{subseteq,subseteq_finite_is_finite}.
        Take $M_0\in M$ such that $x\in M_0$ by \cref{covering,subseteq,unions_iff}.
        $W$ intersects $M_0$ by \cref{neighborhood,inter_intro,emptyset,intersects}.
        Hence $M_0\in F$ by \cref{locally_finite_collection}.
        $M_0\in G$ by \cref{inter_intro}.
        $G\neq\emptyset$ by \cref{emptyset}.
        Let $S = \{w\in G\times\reals \mid 0 < \snd{w} \land \metricBall{\apply{e}{\fst{w}}}{\fst{w}}{x}{\snd{w} + \snd{w}}\subseteq U\inter\fst{w}\}$.
        For all $w$ such that $w\in S$ there exist $C,r$ such that $w = (C,r)$.
        $S$ is a relation by \cref{relation}.
        Show that for all $C\in G$ there exists $r$ such that $C\mathrel{S} r$.
        \begin{subproof}
            Fix $C\in G$.
            We have $C\in F$ and $x\in C$ by \cref{inter}.
            $C\mathrel{P}\apply{e}{C}$ by \cref{subseteq,choice_function}.
            $\apply{e}{C}$ is a metric on $C$ and $\subspaceTopology{T}{C} = \metricTopology{\apply{e}{C}}{C}$ by \cref{relation,fst_eq,snd_eq}.
            $C\subseteq X$ by \cref{open_refinement_on,refinement_on,subseteq,powerset_elim}.
            $C$ is a subspace of $X$ with respect to $T$ by \cref{subspace_of}.
            $\subspaceTopology{T}{C}$ is a topology on $C$ by \cref{subspace_topology_is_topology}.
            $U\inter C\subseteq C$ by \cref{inter_lower_right}.
            $U\inter C = C\inter U$ by \cref{inter_comm}.
            $U\inter C\in\subspaceTopology{T}{C}$ by \cref{subspace_topology}.
            Hence $U\inter C$ is open in $C$ with respect to $\subspaceTopology{T}{C}$ by \cref{open_in}.
            $U\inter C$ is open in $C$ with respect to $\metricTopology{\apply{e}{C}}{C}$ by \cref{real_lt_subst_left,open_in}.
            $x\in U\inter C$ by \cref{inter_intro}.
            Take $\epsilon\in\reals$ such that $0 < \epsilon$ and
                $\metricBall{\apply{e}{C}}{C}{x}{\epsilon}\subseteq U\inter C$
                by \cref{metric_open_iff}.
            Let $r = \epsilon / (1 + 1)$.
            We have $r\in\reals$ and $0 < r$ and $r + r = \epsilon$ by \cref{real_half_of_positive}.
            $\metricBall{\apply{e}{C}}{C}{x}{r + r}\subseteq U\inter C$
                by \cref{metric_ball,real_div,real_mul_inverse,real_mul_closed,real_add_assoc,real_distrib,real_mul_ident,real_add_ident,real_lt_subst_right,subseteq}.
            Hence $(C,r)\in G\times\reals$ by \cref{times_tuple_intro}.
            $C\mathrel{S} r$ by \cref{times_tuple_intro,fst_eq,snd_eq}.
        \end{subproof}
        Take $s$ such that $s$ is a function on $G$ and for all $C\in G$ we have $C\mathrel{S}\apply{s}{C}$ by \cref{choice_function}.
        Let $E = \img{s}{G}$.
        $E$ is finite by \cref{finite_image_function}.
        For all $r$ such that $r\in E$ we have $r\in\reals$
            by \cref{img_iff,function_apply_iff,choice_function,relation,fst_eq,snd_eq,times_tuple_elim,real_lt_subst_left}.
        $E\subseteq\reals$ by \cref{subseteq}.
        Take $C_0$ such that $C_0\in G$ by \cref{nonempty_inhabited}.
        $E\neq\emptyset$ by \cref{img_iff,function_apply_elim,emptyset}.
        Take $\delta\in E$ such that $\delta$ is a minimum of $E$ by \cref{finite_real_minimum}.
        $\delta\in\reals$ by \cref{subseteq}.
        $0 < \delta$ by \cref{img_iff,function_apply_iff,choice_function,relation,fst_eq,snd_eq,subseteq}.
        Take $m\in\naturalsPlus$ such that $0 < 1 / m$ and $1 / m < \delta$ by \cref{real_small_reciprocal}.
        Take $D_0\in\apply{d}{m}$ such that $x\in D_0$ by \cref{choice_function,relation,fst_eq,snd_eq,covering,subseteq,unions_iff}.
        Take $B_0\in\apply{a}{m}$ such that $D_0\subseteq B_0$ by \cref{choice_function,relation,fst_eq,snd_eq,open_refinement_on,refinement_on}.
        Take $C\in M$ such that there exists $y\in C$ such that $B_0 = \metricBall{\apply{e}{C}}{C}{y}{1 / m}$ by \cref{choice_function,relation,fst_eq,snd_eq,real_lt_subst_left}.
        Take $y\in C$ such that $B_0 = \metricBall{\apply{e}{C}}{C}{y}{1 / m}$ by \cref{choice_function,relation,fst_eq,snd_eq,real_lt_subst_left}.
        $x\in B_0$ by \cref{subseteq}.
        $x\in C$ by \cref{subseteq,metric_ball}.
        $W$ intersects $C$ by \cref{neighborhood,inter_intro,emptyset,intersects}.
        Hence $C\in F$ by \cref{locally_finite_collection}.
        $C\in G$ by \cref{inter_intro}.
        We have $C\mathrel{S}\apply{s}{C}$ by \cref{choice_function}.
        Hence $0 < \apply{s}{C}$ and $\metricBall{\apply{e}{C}}{C}{x}{\apply{s}{C} + \apply{s}{C}}\subseteq U\inter C$
            by \cref{relation,fst_eq,snd_eq}.
        $\apply{s}{C}\in\reals$ by \cref{relation,fst_eq,snd_eq,times_tuple_elim}.
        Hence $\delta \leq \apply{s}{C}$ by \cref{real_minimum,img_iff,function_apply_iff}.
        Show that for all $z$ such that $z\in D_0$ we have $z\in U$.
        \begin{subproof}
            Fix $z$ such that $z\in D_0$.
            $z\in B_0$ by \cref{subseteq}.
            $z\in C$ and $\apply{\apply{e}{C}}{(y,z)} < 1 / m$ by \cref{metric_ball}.
            $C\mathrel{P}\apply{e}{C}$ by \cref{choice_function}.
            $\apply{e}{C}$ is a metric on $C$ by \cref{relation,fst_eq,snd_eq}.
            We have $\apply{\apply{e}{C}}{(y,z)}\in\reals$ and $\apply{\apply{e}{C}}{(y,x)}\in\reals$
                by \cref{metric_on,times_tuple_intro,funs_type_apply}.
            $\apply{\apply{e}{C}}{(y,x)} < 1 / m$ by \cref{metric_ball}.
            Hence $\apply{\apply{e}{C}}{(x,y)} = \apply{\apply{e}{C}}{(y,x)}$ by \cref{metric_on,metric_symmetric}.
            $\apply{\apply{e}{C}}{(x,y)}\in\reals$ by \cref{metric_on,times_tuple_intro,funs_type_apply}.
            $\apply{\apply{e}{C}}{(x,y)} < 1 / m$ by \cref{real_lt_subst_left}.
            $1 / m\in\reals$ by \cref{naturalsplus_reciprocal_real}.
            Thus $\apply{\apply{e}{C}}{(y,z)} + \apply{\apply{e}{C}}{(x,y)} < (1 / m) + \apply{\apply{e}{C}}{(x,y)}$ by \cref{real_lt_add}.
            $\apply{\apply{e}{C}}{(x,y)} + \apply{\apply{e}{C}}{(y,z)} < \apply{\apply{e}{C}}{(x,y)} + (1 / m)$
                by \cref{real_add_comm,real_lt_subst_left,real_lt_subst_right}.
            Thus $\apply{\apply{e}{C}}{(x,y)} + (1 / m) < (1 / m) + (1 / m)$ by \cref{real_lt_add}.
            We have $\apply{\apply{e}{C}}{(x,y)} + \apply{\apply{e}{C}}{(y,z)}\in\reals$ and
                $\apply{\apply{e}{C}}{(x,y)} + (1 / m)\in\reals$ and
                $(1 / m) + (1 / m)\in\reals$
                by \cref{real_add_closed}.
            Therefore $\apply{\apply{e}{C}}{(x,y)} + \apply{\apply{e}{C}}{(y,z)} < (1 / m) + (1 / m)$ by \cref{real_lt_trans}.
            $\apply{\apply{e}{C}}{(x,z)}\in\reals$ by \cref{metric_on,times_tuple_intro,funs_type_apply}.
            We have $\apply{\apply{e}{C}}{(x,z)} \leq \apply{\apply{e}{C}}{(x,y)} + \apply{\apply{e}{C}}{(y,z)}$ by \cref{metric_on,metric_triangle_inequality}.
            $\apply{\apply{e}{C}}{(x,z)} < (1 / m) + (1 / m)$ by \cref{real_leq_split,real_lt_subst_left,real_lt_trans}.
            $1 / m < \apply{s}{C}$ by \cref{real_leq_split,real_lt_trans}.
            Thus $(1 / m) + (1 / m) < \apply{s}{C} + (1 / m)$ by \cref{real_lt_add}.
            Therefore $(1 / m) + \apply{s}{C} < \apply{s}{C} + \apply{s}{C}$ by \cref{real_lt_add}.
            Thus $(1 / m) + (1 / m) < (1 / m) + \apply{s}{C}$ by \cref{real_add_comm,real_lt_subst_left,real_lt_subst_right}.
            We have $(1 / m) + (1 / m)\in\reals$ and
                $(1 / m) + \apply{s}{C}\in\reals$ and
                $\apply{s}{C} + \apply{s}{C}\in\reals$
                by \cref{real_add_closed}.
            Therefore $(1 / m) + (1 / m) < \apply{s}{C} + \apply{s}{C}$ by \cref{real_lt_trans}.
            $\apply{\apply{e}{C}}{(x,z)} < \apply{s}{C} + \apply{s}{C}$ by \cref{real_lt_trans}.
            $z\in\metricBall{\apply{e}{C}}{C}{x}{\apply{s}{C} + \apply{s}{C}}$ by \cref{metric_ball}.
            Hence $z\in U\inter C$ by \cref{subseteq}.
            $z\in U$ by \cref{inter}.
        \end{subproof}
        $D_0\subseteq U$ by \cref{subseteq}.
        $D_0\in B$ and $x\in D_0$ and $D_0\subseteq U$ by \cref{unions_iff,img_iff,function_apply_elim}.
    \end{subproof}
    $B$ is a basis for $T$ on $X$ by \cref{open_collection_basis}.
\end{proof}
\section{§42 The Smirnov Metrization Theorem}

% from §42 Item 1: locally metrizable space
\begin{definition}\label{locally_metrizable_space}
    $X$ is a locally metrizable space with respect to $T$ iff
        $T$ is a topology on $X$ and
        for all $x\in X$ there exists $U$ such that
            $U$ is a neighborhood of $x$ in $X$ with respect to $T$ and
            $U$ is metrizable with respect to $\subspaceTopology{T}{U}$.
\end{definition}

% from §42 helper lemma 1: metrizable spaces are locally metrizable
\begin{lemma}\label{metrizable_implies_locally_metrizable}
    Assume $X$ is metrizable with respect to $T$.
    Then $X$ is a locally metrizable space with respect to $T$.
\end{lemma}
\begin{proof}
    Follows by \cref{metrizable,open_in,topology_on,neighborhood,powerset_elim,subseteq,inter_absorb_supseteq_left,subspace_topology,subseteq_antisymmetric,locally_metrizable_space}.
\end{proof}

% from §42 helper lemma 2: paracompact Hausdorff locally metrizable spaces are regular
\begin{lemma}\label{paracompact_hausdorff_locally_metrizable_regular}
    Assume $X$ is paracompact with respect to $T$.
    Assume $X$ is Hausdorff with respect to $T$.
    Assume $X$ is a locally metrizable space with respect to $T$.
    Then $X$ is regular with respect to $T$.
\end{lemma}
\begin{proof}
    Follows by \cref{paracompact_hausdorff_normal,normal_implies_regular}.
\end{proof}

% from §42 helper lemma 3: locally metrizable spaces are locally metrizable
\begin{lemma}\label{locally_metrizable_space_implies_locally_metrizable}
    Assume $X$ is a locally metrizable space with respect to $T$.
    Then $X$ is locally metrizable with respect to $T$.
\end{lemma}
\begin{proof}
    Follows by \cref{locally_metrizable_space,locally_metrizable}.
\end{proof}

% from §42 Item 2: Smirnov metrization theorem
\begin{theorem}\label{smirnov_metrization_theorem}
    $X$ is metrizable with respect to $T$ iff
        $X$ is paracompact with respect to $T$ and
        $X$ is Hausdorff with respect to $T$ and
        $X$ is a locally metrizable space with respect to $T$.
\end{theorem}
\begin{proof}
    Show that if $X$ is metrizable with respect to $T$ then
        $X$ is paracompact with respect to $T$ and
        $X$ is Hausdorff with respect to $T$ and
        $X$ is a locally metrizable space with respect to $T$.
    \begin{subproof}
        Follows by \cref{metrizable_implies_paracompact,metrizable,metric_space_hausdorff,metrizable_implies_locally_metrizable}.
    \end{subproof}

    Show that if
        $X$ is paracompact with respect to $T$ and
        $X$ is Hausdorff with respect to $T$ and
        $X$ is a locally metrizable space with respect to $T$
    then $X$ is metrizable with respect to $T$.
    \begin{subproof}
        Follows by \cref{paracompact_hausdorff_locally_metrizable_regular,locally_metrizable_space_implies_locally_metrizable,paracompact_hausdorff_locally_metrizable_basis,nagata_smirnov_metrization_theorem}.
    \end{subproof}
\end{proof}
\section{§43 Complete Metric Spaces}

% from §43 Item 1: Cauchy sequence
\begin{definition}\label{cauchy_sequence}
    $s$ is a Cauchy sequence in $X$ with respect to $d$ iff
        $s$ is a sequence in $X$ and
        $d$ is a metric on $X$ and
        for all $\epsilon\in\reals$ such that $0 < \epsilon$
        there exists $N\in\naturalsPlus$ such that
            for all $n,m\in\naturalsPlus$ if $N \leq n$ and $N \leq m$ then
                $d((s(n),s(m))) < \epsilon$.
\end{definition}

% from §43 Item 1: complete metric space
\begin{definition}\label{complete_metric_space}
    $X$ is complete with respect to $d$ iff
        $d$ is a metric on $X$ and
        for all $s$ if
            $s$ is a Cauchy sequence in $X$ with respect to $d$
        then there exists $x\in X$ such that
            $s$ converges to $x$ in $X$ with respect to $\metricTopology{d}{X}$.
\end{definition}

% from §43 helper lemma 1: strictly increasing maps on naturalsPlus dominate the identity
\begin{lemma}\label{strictly_increasing_naturalsplus_geq_identity}
    Assume $u$ is strictly increasing on $\naturalsPlus$.
    Suppose $n\in\naturalsPlus$.
    Then $n \leq u(n)$.
\end{lemma}
\begin{proof}
    Let $P = \{k\in\naturalsPlus \mid k \leq u(k)\}$.
    We have $P\subseteq\naturalsPlus$ by \cref{subseteq}.
    $1\in P$ by \cref{real_one_in_naturals,strictly_increasing_on_set,funs_type_apply,real_one_leq_naturalsplus}.
    Show that for all $k\in P$ we have $k + 1\in P$.
    \begin{subproof}
        Fix $k$ such that $k\in P$.
        We have $k \leq u(k)$.
        $k + 1\in\naturalsPlus$ and $k < k + 1$
            by \cref{subseteq,real_naturals_closed_succ,real_naturals_lt_succ_local}.
        Hence $u(k) < u(k + 1)$ by \cref{strictly_increasing_on_set}.
        Show that it is wrong that $u(k + 1) < k + 1$.
        \begin{subproof}
            Suppose not.
            Then $u(k + 1) < k + 1$.
            $u(k + 1)\in\naturalsPlus$ by \cref{strictly_increasing_on_set,funs_type_apply}.
            Hence $u(k + 1)\leq k + 1$ by \cref{real_lt_imp_leq}.
            $u(k + 1)\neq k + 1$ by \cref{real_lt_subst_left,real_naturals_closed_succ,real_naturals_subset_reals,subseteq,real_lt_irreflexive}.
            $u(k + 1)\leq k$ by \cref{real_naturals_leq_succ_elim}.
            We have $k\in\reals$ and $u(k)\in\reals$ and $u(k + 1)\in\reals$
                by \cref{strictly_increasing_on_set,funs_type_apply,real_naturals_subset_reals,subseteq}.
            Hence $k < u(k + 1)$ by \cref{real_leq_lt_trans}.
            $k < k$ by \cref{real_lt_leq_trans_local}.
            Contradiction by \cref{real_lt_irreflexive}.
        \end{subproof}
        $u(k + 1)\in\reals$ and $k + 1\in\reals$
            by \cref{strictly_increasing_on_set,funs_type_apply,subseteq,real_naturals_closed_succ,real_naturals_subset_reals}.
        $k + 1 < u(k + 1)$ or $k + 1 = u(k + 1)$ or $u(k + 1) < k + 1$
            by \cref{real_lt_trichotomy}.
        $k + 1 \leq u(k + 1)$ by \cref{real_leq_split}.
        $k + 1\in P$.
    \end{subproof}
    For all $k\in\naturalsPlus$ we have $k \leq u(k)$ by \cref{real_naturals_induction,subseteq}.
\end{proof}

% from §43 Item 2: Cauchy sequence criterion for completeness
\begin{lemma}\label{cauchy_subsequence_completeness_criterion}
    Assume $d$ is a metric on $X$.
    Suppose for all $s$ if
        $s$ is a Cauchy sequence in $X$ with respect to $d$
    then there exist $t,x$ such that
        $t$ is a sequence in $X$ and
        $t$ is a subsequence of $s$ and
        $x\in X$ and
        $t$ converges to $x$ in $X$ with respect to $\metricTopology{d}{X}$.
    Then $X$ is complete with respect to $d$.
\end{lemma}
\begin{proof}
    Show that for all $s$ if
        $s$ is a Cauchy sequence in $X$ with respect to $d$
    then there exists $x\in X$ such that
        $s$ converges to $x$ in $X$ with respect to $\metricTopology{d}{X}$.
    \begin{subproof}
        Fix $s$ such that $s$ is a Cauchy sequence in $X$ with respect to $d$.
        $s$ is a sequence in $X$ by \cref{cauchy_sequence}.
        Take $t,x$ such that
            $t$ is a sequence in $X$ and
            $t$ is a subsequence of $s$ and
            $x\in X$ and
            $t$ converges to $x$ in $X$ with respect to $\metricTopology{d}{X}$
            by assumption.
        Take $u$ such that
            $u$ is strictly increasing on $\naturalsPlus$ and
            for all $n\in\naturalsPlus$ we have $t(n) = \apply{s}{\apply{u}{n}}$ by \cref{subsequence}.
        Show that for all $U$ if $U$ is a neighborhood of $x$ in $X$ with respect to $\metricTopology{d}{X}$ then
            there exists $N\in\naturalsPlus$ such that for all $n\in\naturalsPlus$ if $N\leq n$ then $s(n)\in U$.
        \begin{subproof}
            Fix $U$ such that $U$ is a neighborhood of $x$ in $X$ with respect to $\metricTopology{d}{X}$.
            We have $U$ is open in $X$ with respect to $\metricTopology{d}{X}$ and $x\in U$ by \cref{neighborhood}.
            Take $\epsilon\in\reals$ such that $0 < \epsilon$ and $\metricBall{d}{X}{x}{\epsilon}\subseteq U$
                by \cref{open_in,topology_on,subseteq,powerset_elim,metric_open_iff}.
            Let $m_2 = 1 + 1$.
            Let $\delta = \epsilon / m_2$.
            We have $\delta\in\reals$ and $0 < \delta$ and $\delta + \delta = \epsilon$ by \cref{real_half_of_positive}.
            $1\in\reals$ and $0\in\reals$ and $0 < 1$
                by \cref{real_naturals_subset_reals,real_one_in_naturals,subseteq,real_add_ident,real_zero_lt_one}.
            $\metricBall{d}{X}{x}{\delta}\in\metricBasis{d}{X}$ by \cref{metric_ball,subseteq,powerset_intro,metric_basis}.
            Hence $\metricBall{d}{X}{x}{\delta}\in\metricTopology{d}{X}$
                by \cref{metric_basis_is_basis,basis_elem_open_in_generated,metric_topology}.
            $\metricBall{d}{X}{x}{\delta}$ is open in $X$ with respect to $\metricTopology{d}{X}$
                by \cref{open_in,basis_topology_is_topology,metric_topology,metric_basis_is_basis,basis_for_topology}.
            Hence $\metricBall{d}{X}{x}{\delta}$ is a neighborhood of $x$ in $X$ with respect to $\metricTopology{d}{X}$
                by \cref{metric_ball,metric_zero_characterization,real_lt_subst_left,metric_on,neighborhood}.
            Take $N_1\in\naturalsPlus$ such that for all $n\in\naturalsPlus$ if $N_1\leq n$ then
                $t(n)\in\metricBall{d}{X}{x}{\delta}$ by \cref{converges_to}.
            Take $N_0\in\naturalsPlus$ such that for all $n,m\in\naturalsPlus$ if $N_0\leq n$ and $N_0\leq m$ then
                $d((s(n),s(m))) < \delta$ by \cref{cauchy_sequence}.
            Let $N = N_0 + N_1$.
            $N\in\naturalsPlus$ by \cref{naturalsplus_add_closed}.
            Show that for all $n\in\naturalsPlus$ if $N\leq n$ then $s(n)\in U$.
            \begin{subproof}
                Fix $n\in\naturalsPlus$.
                Suppose $N\leq n$.
                $N_0\in\reals$ and $N_1\in\reals$ and $N\in\reals$ by \cref{real_naturals_subset_reals,subseteq}.
                $n\in\reals$ by \cref{real_naturals_subset_reals,subseteq}.
                $0 < N_0$ and $0 < N_1$ by \cref{real_zero_lt_one,real_one_leq_naturalsplus,real_lt_trans}.
                $0 + N_0 < N_1 + N_0$ by \cref{real_lt_add}.
                $0 + N_0 = N_0$ and $N_1 + N_0 = N_0 + N_1$ by \cref{real_add_ident,real_add_comm}.
                Hence $N_0 < N_0 + N_1$ by \cref{real_lt_subst_left,real_lt_subst_right}.
                $N_0 < N$ by \cref{real_lt_subst_right}.
                $N_0 \leq N$ by \cref{real_lt_imp_leq}.
                $0 + N_1 < N_0 + N_1$ by \cref{real_lt_add}.
                $0 + N_1 = N_1$ by \cref{real_add_ident}.
                Hence $N_1 < N_0 + N_1$ by \cref{real_lt_subst_left,real_lt_subst_right}.
                $N_1 < N$ by \cref{real_lt_subst_right}.
                $N_1 \leq N$ by \cref{real_lt_imp_leq}.
                $N_0\leq n$ and $N_1\leq n$ by \cref{real_leq_trans}.
                $u(n)\in\naturalsPlus$ by \cref{strictly_increasing_on_set,funs_type_apply}.
                Hence $n \leq u(n)$ by \cref{strictly_increasing_naturalsplus_geq_identity}.
                We have $N_0\leq u(n)$ and $N_1\leq u(n)$
                    by \cref{real_naturals_subset_reals,subseteq,real_leq_trans}.
                $d((s(n),s(u(n)))) < \delta$ by \cref{cauchy_sequence}.
                $d((s(u(n)),x)) < \delta$
                    by \cref{converges_to,metric_ball,metric_on,metric_symmetric,real_lt_subst_left}.
                $s(n)\in X$ and $s(u(n))\in X$ by \cref{cauchy_sequence,sequence_in,funs_type_apply}.
                $d((s(n),x)) \leq d((s(n),s(u(n)))) + d((s(u(n)),x))$ by \cref{metric_on,metric_triangle_inequality}.
                $d((s(n),s(u(n))))\in\reals$ and $d((s(u(n)),x))\in\reals$ and $d((s(n),x))\in\reals$
                    by \cref{times_tuple_intro,metric_on,funs_type_apply}.
                Hence $d((s(u(n)),x)) + d((s(n),s(u(n)))) < \delta + d((s(n),s(u(n))))$ by \cref{real_lt_add}.
                $d((s(u(n)),x)) + d((s(n),s(u(n)))) = d((s(n),s(u(n)))) + d((s(u(n)),x))$ and
                    $\delta + d((s(n),s(u(n)))) = d((s(n),s(u(n)))) + \delta$ by \cref{real_add_comm}.
                $d((s(n),s(u(n)))) + d((s(u(n)),x)) < d((s(n),s(u(n)))) + \delta$
                    by \cref{real_lt_subst_left,real_lt_subst_right}.
                $d((s(n),s(u(n)))) + \delta < \delta + \delta$ by \cref{real_lt_add}.
                $d((s(n),s(u(n)))) + d((s(u(n)),x)) < \delta + \delta$
                    by \cref{real_add_closed,real_lt_trans}.
                Hence $d((s(n),s(u(n)))) + d((s(u(n)),x)) < \epsilon$ by \cref{real_lt_subst_right}.
                $d((s(n),s(u(n)))) + d((s(u(n)),x))\in\reals$ by \cref{real_add_closed}.
                $d((s(n),x)) < d((s(n),s(u(n)))) + d((s(u(n)),x))$ or
                    $d((s(n),x)) = d((s(n),s(u(n)))) + d((s(u(n)),x))$ by \cref{real_leq_split}.
                Hence $d((s(n),x)) < \epsilon$ by \cref{real_lt_trans,real_lt_subst_left}.
                $d((x,s(n))) < \epsilon$ by \cref{metric_on,metric_symmetric,real_lt_subst_left}.
                $s(n)\in\metricBall{d}{X}{x}{\epsilon}$ by \cref{metric_ball}.
                $s(n)\in U$ by \cref{subseteq}.
            \end{subproof}
            For all $n\in\naturalsPlus$ if $N\leq n$ then $s(n)\in U$.
        \end{subproof}
        $s$ converges to $x$ in $X$ with respect to $\metricTopology{d}{X}$ by \cref{converges_to}.
    \end{subproof}
    $X$ is complete with respect to $d$ by \cref{complete_metric_space}.
\end{proof}

% from §43 helper lemma 2: Cauchy sequences in metric spaces have bounded range
\begin{lemma}\label{cauchy_sequence_bounded_range}
    Assume $d$ is a metric on $X$.
    Suppose $s$ is a Cauchy sequence in $X$ with respect to $d$.
    Then $\ran{s}$ is bounded in $X$ with respect to $d$.
\end{lemma}
\begin{proof}
    We have $s\in\funs{\naturalsPlus}{X}$ and $s$ is a function by \cref{cauchy_sequence,sequence_in,funs_is_function}.
    Take $N\in\naturalsPlus$ such that for all $n,m\in\naturalsPlus$ if $N\leq n$ and $N\leq m$ then
        $d((s(n),s(m))) < 1$ by \cref{real_zero_lt_one,real_naturals_subset_reals,real_one_in_naturals,subseteq,cauchy_sequence}.
    Let $a = s(N)$.
    Let $g = \{(i,d((a,s(i)))) \mid i\in\initialSegment{N}\}$.
    For all $i,r_1,r_2$ such that $(i,r_1)\in g$ and $(i,r_2)\in g$ we have $r_1 = r_2$ by \cref{pair_eq_iff}.
    $g$ is a function by \cref{relation,rightunique}.
    For all $i$ such that $i\in\initialSegment{N}$ we have $i\in\dom{g}$ by \cref{dom_intro}.
    For all $i$ such that $i\in\dom{g}$ we have $i\in\initialSegment{N}$ by \cref{dom_iff,pair_eq_iff}.
    $\dom{g} = \initialSegment{N}$ by \cref{subseteq,subseteq_antisymmetric}.
    For all $i\in\dom{g}$ we have $g(i)\in\reals$
        by \cref{subseteq,initial_segment_naturalsplus,sequence_in,funs_type_apply,metric_on,times_tuple_intro,function_apply_intro}.
    $g\in\funs{\initialSegment{N}}{\reals}$ by \cref{funs_intro}.
    Let $D = \img{g}{\initialSegment{N}}$.
    $D$ is finite by \cref{initial_segment_finite,finite_image_function}.
    $N\in\initialSegment{N}$ by \cref{initial_segment_naturalsplus}.
    $g(N)\in D$ by \cref{function_apply_elim,img_iff}.
    $D\neq\emptyset$ by \cref{inhabited_nonempty,nonempty_inhabited}.
    $D\subseteq\reals$ by \cref{subseteq,img_iff,function_apply_elim}.
    Take $M_0\in D$ such that $M_0$ is a maximum of $D$ by \cref{finite_real_maximum}.
    $a\in X$ and $s(N)\in X$ by \cref{sequence_in,funs_type_apply}.
    $d((a,s(N))) \geq 0$ by \cref{metric_on,metric_nonnegative}.
    $g(N) = d((a,s(N)))$ by \cref{function_apply_intro}.
    $d((a,s(N)))\in\reals$ by \cref{times_tuple_intro,metric_on,funs_type_apply}.
    Hence $0 < d((a,s(N)))$ or $0 = d((a,s(N)))$ by \cref{real_leq_split}.
    Thus $0 < g(N)$ or $0 = g(N)$ by \cref{real_lt_subst_right,real_lt_subst_left}.
    Therefore $g(N) \leq M_0$ by \cref{real_maximum}.
    $g(N)\in\reals$ and $M_0\in\reals$ by \cref{function_apply_elim,subseteq,real_maximum}.
    $0\in\reals$ by \cref{real_add_ident}.
    $0 \leq g(N)$ by \cref{real_leq_split}.
    $0 \leq M_0$ by \cref{real_leq_trans}.
    Let $M = (M_0 + 1) + (M_0 + 1)$.
    Show that for all $z\in\ran{s}$ we have $d((a,z)) \leq M_0 + 1$.
    \begin{subproof}
        Fix $z\in\ran{s}$.
        Take $k$ such that $(k,z)\in s$ by \cref{ran_iff}.
        Hence $k\in\dom{s}$ by \cref{dom}.
        Thus $k\in\naturalsPlus$ by \cref{sequence_in,funs_elim,subseteq}.
        We have $s(k) = z$ by \cref{function_apply_iff}.
        $z\in X$ by \cref{sequence_in,funs_type_apply}.
        Thus $k < N$ or $k = N$ or $N < k$ by \cref{real_naturals_subset_reals,subseteq,real_lt_trichotomy}.
        \begin{byCase}
        \caseOf{$k < N$.}
            We have $k \leq N$ by \cref{real_lt_imp_leq}.
            Hence $k\in\initialSegment{N}$ by \cref{initial_segment_naturalsplus}.
            Hence $g(k)\in D$ by \cref{function_apply_elim,img_iff}.
            Therefore $g(k) \leq M_0$ by \cref{real_maximum}.
            Hence $d((a,s(k))) \leq M_0$ by \cref{function_apply_intro}.
            Thus $d((a,z)) \leq M_0$ by \cref{real_lt_subst_right}.
            $1\in\reals$ and $0 < 1$ by \cref{real_naturals_subset_reals,real_one_in_naturals,subseteq,real_zero_lt_one}.
            Thus $d((a,z))\in\reals$ by \cref{times_tuple_intro,metric_on,funs_type_apply}.
            We have $M_0 + 1\in\reals$ by \cref{real_add_closed}.
            $0 + M_0 < 1 + M_0$ by \cref{real_lt_add}.
            Hence $M_0 < M_0 + 1$ by \cref{real_add_ident,real_add_comm,real_lt_subst_left,real_lt_subst_right}.
            $d((a,z)) < M_0$ or $d((a,z)) = M_0$ by \cref{real_leq_split}.
            Thus $d((a,z)) < M_0 + 1$ by \cref{real_lt_trans,real_lt_subst_right}.
            Therefore $d((a,z)) \leq M_0 + 1$ by \cref{real_lt_imp_leq}.
        \caseOf{$k = N$.}
            We have $a = z$ by \cref{real_lt_subst_left}.
            Thus $d((a,z)) = 0$ by \cref{metric_on,metric_zero_characterization}.
            $1\in\reals$ and $0 < 1$ by \cref{real_naturals_subset_reals,real_one_in_naturals,subseteq,real_zero_lt_one}.
            Hence $d((a,z))\in\reals$ by \cref{times_tuple_intro,metric_on,funs_type_apply}.
            Thus $M_0 + 1\in\reals$ by \cref{real_add_closed}.
            We have $0 + M_0 < 1 + M_0$ by \cref{real_lt_add}.
            Hence $M_0 < M_0 + 1$ by \cref{real_add_ident,real_add_comm,real_lt_subst_left}.
            Thus $0 < M_0 + 1$ by \cref{real_leq_lt_trans}.
            Therefore $d((a,z)) < M_0 + 1$ by \cref{real_lt_subst_left}.
            We have $d((a,z)) \leq M_0 + 1$ by \cref{real_lt_imp_leq}.
        \caseOf{$N < k$.}
            We have $N \leq k$ by \cref{real_lt_imp_leq}.
            Hence $d((s(N),s(k))) < 1$ by \cref{cauchy_sequence}.
            Thus $d((a,z)) < 1$ by \cref{real_lt_subst_left,real_lt_subst_right}.
            $1\in\reals$ by \cref{real_naturals_subset_reals,real_one_in_naturals,subseteq}.
            Hence $d((a,z))\in\reals$ by \cref{times_tuple_intro,metric_on,funs_type_apply}.
            Thus $M_0 + 1\in\reals$ by \cref{real_add_closed}.
            $0 < M_0$ or $0 = M_0$.
            Therefore $d((a,z)) < M_0 + 1$ by \cref{real_lt_add,real_add_ident,real_lt_subst_left,real_lt_trans,real_lt_subst_right}.
            We have $d((a,z)) \leq M_0 + 1$ by \cref{real_lt_imp_leq}.
        \end{byCase}
    \end{subproof}
    Show that for all $x,y\in\ran{s}$ we have $d((x,y)) \leq M$.
    \begin{subproof}
        Fix $x$.
        Fix $y$ such that $x\in\ran{s}$ and $y\in\ran{s}$.
        Thus $x\in X$ and $y\in X$ by \cref{sequence_in,ran_iff,dom,funs_elim,subseteq,function_apply_iff,funs_type_apply}.
        We have $d((a,x)) \leq M_0 + 1$.
        $d((a,y)) \leq M_0 + 1$.
        We have $a\in X$ by \cref{sequence_in,funs_type_apply}.
        Thus $d((x,y)) \leq d((x,a)) + d((a,y))$ by \cref{metric_on,metric_triangle_inequality}.
        Hence $d((x,a)) \leq M_0 + 1$ by \cref{metric_on,metric_symmetric}.
        $(x,y)\in X\times X$ and $(x,a)\in X\times X$ and $(a,y)\in X\times X$ by \cref{times_tuple_intro}.
        $d((x,y))\in\reals$ and $d((x,a))\in\reals$ and $d((a,y))\in\reals$
            by \cref{metric_on,funs_type_apply}.
        We have $M_0 + 1\in\reals$ by \cref{real_naturals_subset_reals,real_one_in_naturals,subseteq,real_add_closed}.
        $d((x,a)) + d((a,y))\in\reals$ and $(M_0 + 1) + d((a,y))\in\reals$ and $(M_0 + 1) + (M_0 + 1)\in\reals$
            by \cref{real_add_closed}.
        Thus $d((x,a)) + d((a,y)) \leq (M_0 + 1) + d((a,y))$ by \cref{real_leq_add}.
        We have $d((a,y)) + (M_0 + 1) \leq (M_0 + 1) + (M_0 + 1)$ by \cref{real_leq_add}.
        $(M_0 + 1) + d((a,y)) \leq (M_0 + 1) + (M_0 + 1)$ by \cref{real_add_comm}.
        Therefore $d((x,a)) + d((a,y)) \leq (M_0 + 1) + (M_0 + 1)$ by \cref{real_leq_trans}.
        Thus $d((x,a)) + d((a,y)) \leq M$ by \cref{real_add_comm}.
        Therefore $d((x,y)) \leq M$ by \cref{real_leq_trans}.
    \end{subproof}
    Hence $M\in\reals$ by \cref{real_naturals_subset_reals,real_one_in_naturals,subseteq,real_add_closed}.
    $\ran{s}\subseteq X$ by \cref{sequence_in,funs_subseteq_rels,subseteq,rels_ran_subseteq}.
    Hence $\ran{s}$ is bounded in $X$ with respect to $d$ by \cref{metric_bounded}.
\end{proof}

% from §43 Item 3: euclidean space is complete in the square metric
\begin{theorem}\label{euclidean_space_square_metric_complete}
    Assume $n\in\naturalsPlus$.
    Let $X = \realsPower{\initialSegment{n}}$.
    Then $X$ is complete with respect to $\squareMetric{n}$.
\end{theorem}
\begin{proof}
    Let $d = \squareMetric{n}$.
    Let $T = \metricTopology{d}{X}$.
    $d$ is a metric on $X$ by \cref{square_metric_is_metric}.
    We have $T$ is a topology on $X$ by \cref{metric_basis_is_basis,basis_topology_is_topology,metric_topology,topology_on}.
    Show that for all $s$ if
        $s$ is a Cauchy sequence in $X$ with respect to $d$
    then there exist $t,x$ such that
        $t$ is a sequence in $X$ and
        $t$ is a subsequence of $s$ and
        $x\in X$ and
        $t$ converges to $x$ in $X$ with respect to $T$.
    \begin{subproof}
        Fix $s$ such that $s$ is a Cauchy sequence in $X$ with respect to $d$.
        Let $A = \closure{\ran{s}}{X}{T}$.
        Let $T_A = \subspaceTopology{T}{A}$.
        $s\in\funs{\naturalsPlus}{X}$ by \cref{cauchy_sequence,sequence_in}.
        $\ran{s}\subseteq X$ by \cref{funs_subseteq_rels,subseteq,rels_ran_subseteq}.
        $A$ is closed in $X$ with respect to $T$ and $A\subseteq X$ by \cref{closure_closed_in,closed_in,subseteq}.
        Take $M_1\in\reals$ such that for all $a_1,a_2\in\ran{s}$ we have $d((a_1,a_2)) \leq M_1$
            by \cref{cauchy_sequence_bounded_range,metric_bounded}.
        Let $M = 1 + (M_1 + 1)$.
        Show that for all $p\in A$ there exists $y\in\ran{s}$ such that $y\in X$ and $d((p,y)) < 1$.
        \begin{subproof}
            Fix $p\in A$.
            We have $p\in X$ by \cref{subseteq}.
            $\metricBall{d}{X}{p}{1}$ is open in $X$ with respect to $T$
                by \cref{metric_ball,subseteq,powerset_intro,real_naturals_subset_reals,real_one_in_naturals,real_zero_lt_one,metric_basis,metric_basis_is_basis,basis_elem_open_in_generated,metric_topology,basis_for_topology,open_in}.
            Hence $p\in\metricBall{d}{X}{p}{1}$
                by \cref{metric_ball,metric_zero_characterization,real_naturals_subset_reals,real_one_in_naturals,subseteq,real_zero_lt_one,metric_on,real_lt_subst_left}.
            Thus $\metricBall{d}{X}{p}{1}$ intersects $\ran{s}$ by \cref{closure_iff_intersects_open}.
            Take $y$ such that $y\in\metricBall{d}{X}{p}{1}\inter\ran{s}$ by \cref{intersects,nonempty_inhabited}.
            We have $y\in\ran{s}$ and $y\in X$ and $d((p,y)) < 1$ by \cref{inter_elim_right,inter_elim_left,metric_ball}.
        \end{subproof}
        Show that for all $x,z\in A$ we have $d((x,z)) \leq M$.
        \begin{subproof}
            Fix $x$.
            Fix $z$ such that $x\in A$ and $z\in A$.
            Hence $x\in X$ and $z\in X$ by \cref{subseteq}.
            Take $y_1\in\ran{s}$ such that $y_1\in X$ and $d((x,y_1)) < 1$.
            Take $y_2\in\ran{s}$ such that $y_2\in X$ and $d((z,y_2)) < 1$.
            Therefore $d((y_1,y_2)) \leq M_1$.
            Thus $d((x,z)) \leq d((x,y_1)) + d((y_1,z))$ by \cref{metric_on,metric_triangle_inequality}.
            We have $d((y_1,z)) \leq d((y_1,y_2)) + d((y_2,z))$ by \cref{metric_on,metric_triangle_inequality}.
            Hence $d((y_2,z)) < 1$ by \cref{metric_on,metric_symmetric,real_lt_subst_left}.
            Therefore $d((x,y_1)) \leq 1$ and $d((y_2,z)) \leq 1$ by \cref{real_lt_imp_leq}.
            $(x,z)\in X\times X$ and $(x,y_1)\in X\times X$ and $(y_1,z)\in X\times X$ and
                $(y_1,y_2)\in X\times X$ and $(y_2,z)\in X\times X$ by \cref{times_tuple_intro}.
            We have $d((x,z))\in\reals$ and $d((x,y_1))\in\reals$ and $d((y_1,z))\in\reals$ and
                $d((y_1,y_2))\in\reals$ and $d((y_2,z))\in\reals$ by \cref{metric_on,funs_type_apply}.
            $1\in\reals$ by \cref{real_naturals_subset_reals,real_one_in_naturals,subseteq}.
            $M_1 + 1\in\reals$ by \cref{real_naturals_subset_reals,real_one_in_naturals,subseteq,real_add_closed}.
            We have $d((y_1,y_2)) + d((y_2,z))\in\reals$ and $M_1 + d((y_2,z))\in\reals$ and $1 + M_1\in\reals$ and
                $d((x,y_1)) + d((y_1,z))\in\reals$ and $1 + d((y_1,z))\in\reals$ and $1 + (M_1 + 1)\in\reals$
                by \cref{real_add_closed}.
            Thus $d((y_1,y_2)) + d((y_2,z)) \leq M_1 + d((y_2,z))$ by \cref{real_leq_add}.
            We have $d((y_2,z)) + M_1 \leq 1 + M_1$ by \cref{real_leq_add}.
            $M_1 + d((y_2,z)) \leq 1 + M_1$ by \cref{real_add_comm}.
            Therefore $M_1 + d((y_2,z)) \leq M_1 + 1$ by \cref{real_add_comm}.
            Thus $d((y_1,z)) \leq M_1 + 1$ by \cref{real_leq_trans}.
            We have $d((x,y_1)) + d((y_1,z)) \leq 1 + d((y_1,z))$ by \cref{real_leq_add}.
            $d((y_1,z)) + 1 \leq (M_1 + 1) + 1$ by \cref{real_leq_add}.
            Thus $1 + d((y_1,z)) \leq (M_1 + 1) + 1$ by \cref{real_add_comm}.
            We have $1 + d((y_1,z)) \leq 1 + (M_1 + 1)$ by \cref{real_add_comm}.
            Therefore $d((x,y_1)) + d((y_1,z)) \leq 1 + (M_1 + 1)$ by \cref{real_leq_trans}.
            Therefore $d((x,z)) \leq M$ by \cref{real_leq_trans}.
        \end{subproof}
        $M\in\reals$ by \cref{real_naturals_subset_reals,real_one_in_naturals,subseteq,real_add_closed}.
        $A$ is bounded in $X$ with respect to $d$ by \cref{metric_bounded}.
        $A$ is compact with respect to $T_A$ by \cref{compact_reals_power_iff_closed_bounded}.
        Thus $X$ is metrizable with respect to $T$ by \cref{metrizable}.
        $A$ is a subspace of $X$ with respect to $T$ by \cref{subspace_of}.
        $T_A$ is a topology on $A$ by \cref{subspace_topology_is_topology,subspace_of}.
        $A$ is metrizable with respect to $T_A$ by \cref{subspace_of,subspace_metrizable}.
        Therefore $A$ is sequentially compact with respect to $T_A$
            by \cref{metrizable_compact_iff_limit_point_compact_iff_sequentially_compact}.
        For all $m$ such that $m\in\dom{s}$ we have $s(m)\in A$
            by \cref{cauchy_sequence,sequence_in,funs_elim,function_apply_elim,ran_iff,subseteq_closure,subseteq}.
        We have $s\in\funs{\naturalsPlus}{A}$ by \cref{cauchy_sequence,sequence_in,funs_elim,funs_intro}.
        Thus $s$ is a sequence in $A$ by \cref{sequence_in}.
        Take $t,x$ such that
            $t$ is a sequence in $A$ and
            $t$ is a subsequence of $s$ and
            $x\in A$ and
            $t$ converges to $x$ in $A$ with respect to $T_A$ by \cref{sequentially_compact}.
        $t$ is a sequence in $X$ by \cref{sequence_in,funs_weaken_codom}.
        For all $U$ if $U$ is a neighborhood of $x$ in $X$ with respect to $T$ then
            there exists $N\in\naturalsPlus$ such that for all $n\in\naturalsPlus$ if $N\leq n$ then $t(n)\in U$
            by \cref{subspace_topology,open_in,inter_intro,subseteq,neighborhood,converges_to,inter_elim_right}.
        $t$ converges to $x$ in $X$ with respect to $T$ by \cref{converges_to}.
        Thus $t$ is a sequence in $X$ and
            $t$ is a subsequence of $s$ and
            $x\in X$ and
            $t$ converges to $x$ in $X$ with respect to $T$
            by \cref{sequence_in,subseteq}.
    \end{subproof}
    $X$ is complete with respect to $d$ by \cref{cauchy_subsequence_completeness_criterion}.
\end{proof}

% from §43 Item 3: euclidean space is complete in the euclidean metric
\begin{theorem}\label{euclidean_space_euclidean_metric_complete}
    Assume $n\in\naturalsPlus$.
    Let $X = \realsPower{\initialSegment{n}}$.
    Then $X$ is complete with respect to $\euclideanMetric{n}$.
\end{theorem}
\begin{proof}
    Let $d = \euclideanMetric{n}$.
    Let $d_1 = \squareMetric{n}$.
    $d$ is a metric on $X$ and $d_1$ is a metric on $X$
        by \cref{euclidean_metric_is_metric,square_metric_is_metric}.
    Show that for all $s$ if
        $s$ is a Cauchy sequence in $X$ with respect to $d$
    then there exists $x\in X$ such that
        $s$ converges to $x$ in $X$ with respect to $\metricTopology{d}{X}$.
    \begin{subproof}
        Fix $s$ such that $s$ is a Cauchy sequence in $X$ with respect to $d$.
        Show that for all $\epsilon\in\reals$ such that $0 < \epsilon$
        there exists $N\in\naturalsPlus$ such that
            for all $k,l\in\naturalsPlus$ if $N \leq k$ and $N \leq l$ then
                $d_1((s(k),s(l))) < \epsilon$.
        \begin{subproof}
            Fix $\epsilon$ such that $\epsilon\in\reals$ and $0 < \epsilon$.
            Take $N\in\naturalsPlus$ such that
                for all $k,l\in\naturalsPlus$ if $N \leq k$ and $N \leq l$ then
                    $d((s(k),s(l))) < \epsilon$
                by \cref{cauchy_sequence}.
            Show that for all $k,l\in\naturalsPlus$ if $N \leq k$ and $N \leq l$ then
                $d_1((s(k),s(l))) < \epsilon$.
            \begin{subproof}
                Fix $k$.
                Fix $l$.
                Assume $k\in\naturalsPlus$ and $l\in\naturalsPlus$.
                Assume $N \leq k$ and $N \leq l$.
                $s(k)\in X$ and $s(l)\in X$ by \cref{cauchy_sequence,sequence_in,funs_type_apply}.
                Take $r\in\reals$ such that $((s(k),s(l)),r)\in d$ by \cref{euclidean_metric_value_exists}.
                Take $m\in\reals$ such that $((s(k),s(l)),m)\in d_1$ by \cref{square_metric_value_exists}.
                We have $d((s(k),s(l))) = r$ and $d_1((s(k),s(l))) = m$
                    by \cref{euclidean_metric_function,square_metric_function,function_apply_intro}.
                Hence $r < \epsilon$ by \cref{cauchy_sequence,real_lt_subst_left}.
                Thus $m \leq r$ by \cref{euclidean_square_metric_bounds}.
                Therefore $m < \epsilon$ by \cref{real_leq_lt_trans}.
                We have $d_1((s(k),s(l))) < \epsilon$ by \cref{real_lt_subst_left}.
            \end{subproof}
        \end{subproof}
        Hence $s$ is a Cauchy sequence in $X$ with respect to $d_1$ by \cref{cauchy_sequence}.
        Take $x\in X$ such that
            $s$ converges to $x$ in $X$ with respect to $\metricTopology{d_1}{X}$
            by \cref{euclidean_space_square_metric_complete,complete_metric_space}.
        Hence $s$ converges to $x$ in $X$ with respect to $\metricTopology{d}{X}$
            by \cref{square_metric_topology_equals_euclidean,converges_to}.
    \end{subproof}
    Hence $X$ is complete with respect to $d$ by \cref{complete_metric_space}.
\end{proof}

% from §43 helper lemma 3: a function from the empty set is empty
\begin{lemma}\label{function_from_emptyset_is_empty}
    Assume $f\in\funs{\emptyset}{Y}$.
    Then $f = \emptyset$.
\end{lemma}
\begin{proof}
    Follows by \cref{funs_is_relation,funs_elim,relation,dom_intro,emptyset,subseteq,emptyset_subseteq,subseteq_antisymmetric}.
\end{proof}

% from §43 helper lemma 4: indexed products over the empty domain are singletons
\begin{lemma}\label{empty_indexed_product_singleton}
    Assume $X$ is a function.
    Assume $\dom{X} = \emptyset$.
    Then $\productFamily{X} = \{\emptyset\}$.
\end{lemma}
\begin{proof}
    $\emptyset\in\productFamily{X}$
        by \cref{codom_of_emptyset_can_be_anything,emptyset_is_function_on_emptyset,funs_intro,indexed_product,emptyset}.
    For all $f$ if $f\in\productFamily{X}$ then $f\in\{\emptyset\}$
        by \cref{indexed_product,funs_elim,function_from_emptyset_is_empty,singleton_intro}.
    $\productFamily{X} = \{\emptyset\}$ by \cref{subseteq,singleton_subset_intro,subseteq_antisymmetric}.
\end{proof}

% from §43 Item 4: convergence in product spaces is coordinatewise
\begin{lemma}\label{product_sequence_converges_iff_coordinatewise}
    Assume $X$ is a function.
    Assume $T$ is a function.
    Assume $\dom{T} = \dom{X}$.
    Suppose for all $\alpha\in\dom{X}$ we have $\apply{T}{\alpha}$ is a topology on $\apply{X}{\alpha}$.
    Suppose $s$ is a sequence in $\productFamily{X}$.
    Suppose $x\in\productFamily{X}$.
    Then $s$ converges to $x$ in $\productFamily{X}$ with respect to $\productTopologyIndexed{T}{X}$ iff
        for all $\alpha\in\dom{X}$ we have
            $\piIndex{X}{\alpha}\circ s$ converges to $\apply{x}{\alpha}$ in $\apply{X}{\alpha}$ with respect to $\apply{T}{\alpha}$.
\end{lemma}
\begin{proof}
    Let $T_0 = \productTopologyIndexed{T}{X}$.
    Show that if $s$ converges to $x$ in $\productFamily{X}$ with respect to $T_0$ then
        for all $\alpha\in\dom{X}$ we have
            $\piIndex{X}{\alpha}\circ s$ converges to $\apply{x}{\alpha}$ in $\apply{X}{\alpha}$ with respect to $\apply{T}{\alpha}$.
    \begin{subproof}
        Follows by \cref{product_subbasis_finite_intersections_basis,basis_for_topology,basis_topology_is_topology,projection_map_indexed_function,projection_map_indexed_continuous,subseteq,continuous_implies_sequential,funs_is_function,projection_map_indexed,function_apply_intro,converges_to}.
    \end{subproof}
    Show that if for all $\alpha\in\dom{X}$ we have
        $\piIndex{X}{\alpha}\circ s$ converges to $\apply{x}{\alpha}$ in $\apply{X}{\alpha}$ with respect to $\apply{T}{\alpha}$
    then $s$ converges to $x$ in $\productFamily{X}$ with respect to $T_0$.
    \begin{subproof}
        Assume for all $\alpha\in\dom{X}$ we have
            $\piIndex{X}{\alpha}\circ s$ converges to $\apply{x}{\alpha}$ in $\apply{X}{\alpha}$ with respect to $\apply{T}{\alpha}$.
        \begin{byCase}
        \caseOf{$\dom{X}$ is not inhabited.}
            $\dom{X} = \emptyset$ by \cref{emptyset,nonempty_inhabited}.
            $x = \emptyset$ by \cref{empty_indexed_product_singleton,subseteq,singleton_elim}.
            Show that for all $U$ if $U$ is a neighborhood of $x$ in $\productFamily{X}$ with respect to $T_0$ then
                there exists $N\in\naturalsPlus$ such that for all $n\in\naturalsPlus$ if $N\leq n$ then $s(n)\in U$.
            \begin{subproof}
                Follows by \cref{neighborhood,real_one_in_naturals,sequence_in,funs_type_apply,empty_indexed_product_singleton,subseteq,singleton_elim,emptyset}.
            \end{subproof}
            $s$ converges to $x$ in $\productFamily{X}$ with respect to $T_0$ by \cref{converges_to}.
        \caseOf{$\dom{X}$ is inhabited.}
            Let $S = \productSubbasis{T}{X}$.
            Let $B = \finiteIntersections{S}$.
            $B$ is a basis for $T_0$ on $\productFamily{X}$ by \cref{product_subbasis_finite_intersections_basis}.
            Show that for all $U$ if $U$ is a neighborhood of $x$ in $\productFamily{X}$ with respect to $T_0$ then
                there exists $N\in\naturalsPlus$ such that for all $n\in\naturalsPlus$ if $N\leq n$ then $s(n)\in U$.
            \begin{subproof}
                Fix $U$ such that $U$ is a neighborhood of $x$ in $\productFamily{X}$ with respect to $T_0$.
                We have $U$ is open in $\productFamily{X}$ with respect to $T_0$ and $x\in U$ by \cref{neighborhood}.
                Hence $U\in T_0$ and $U\in\pow{\productFamily{X}}$
                    by \cref{product_subbasis_finite_intersections_basis,basis_for_topology,basis_topology_is_topology,open_in,topology_on,subseteq,pow_iff}.
                Take $M\subseteq B$ such that $U = \unions{M}$ by \cref{basis_for_topology,basis_topology_unions}.
                Take $C\in M$ such that $x\in C$ by \cref{unions_iff}.
                Thus $C\in B$ by \cref{subseteq}.
                Take $F\subseteq S$ such that $F$ is finite and inhabited and $C = \inters{F}$ by \cref{finite_intersections}.
                Let $R = \{ w\in F\times\naturalsPlus \mid
                    \text{for all $n\in\naturalsPlus$ if $\snd{w}\leq n$ then $s(n)\in\fst{w}$} \}$.
                $R$ is a relation by \cref{relation,subseteq,times_elem_is_tuple}.
                Show that for all $A\in F$ there exists $N\in\naturalsPlus$ such that $(A,N)\in R$.
                \begin{subproof}
                    Fix $A\in F$.
                    We have $A\in S$ by \cref{subseteq}.
                    Take $\beta\in\dom{X}$ such that $A\in\productSubbasisAt{T}{X}{\beta}$ by \cref{product_subbasis_union_indexed}.
                    Take $V\in\apply{T}{\beta}$ such that $A = \preimg{\piIndex{X}{\beta}}{V}$ by \cref{product_subbasis_indexed}.
                    Take $y\in V$ such that $(x,y)\in\piIndex{X}{\beta}$ by \cref{real_lt_subst_left,inters_destr,preimg_iff}.
                    Thus $\apply{x}{\beta}\in V$ by \cref{projection_map_indexed,pair_eq_iff,real_lt_subst_left}.
                    Therefore $V$ is a neighborhood of $\apply{x}{\beta}$ in $\apply{X}{\beta}$ with respect to $\apply{T}{\beta}$
                        by \cref{neighborhood,open_in,subseteq}.
                    Take $N\in\naturalsPlus$ such that for all $n\in\naturalsPlus$ if $N\leq n$ then
                        $(\piIndex{X}{\beta}\circ s)(n)\in V$ by \cref{converges_to}.
                    For all $n\in\naturalsPlus$ if $N\leq n$ then $s(n)\in A$
                        by \cref{sequence_in,projection_map_indexed_function,funs_circ,funs_elim,funs_is_function,function_apply_elim,circ_iff,function_apply_iff,pair_eq_iff,preimg_iff}.
                    Hence $(A,N)\in R$ by \cref{times_tuple_intro,fst_eq,snd_eq}.
                \end{subproof}
                Take $g$ such that $g$ is a function on $F$ and for all $A\in F$ we have $A\mathrel{R}\apply{g}{A}$ by \cref{choice_function}.
                Let $D = \img{g}{F}$.
                $D$ is finite by \cref{finite_image_function,choice_function}.
                Take $A_0$ such that $A_0\in F$.
                $D\neq\emptyset$ by \cref{function_apply_elim,img_iff,inhabited_nonempty}.
                For all $r$ such that $r\in D$ we have $r\in\naturalsPlus$
                    by \cref{img_iff,choice_function,function_apply_elim,function_rightunique,subseteq,times_tuple_elim}.
                Hence $D\subseteq\reals$
                    by \cref{real_naturals_subset_reals,subseteq}.
                Take $N\in D$ such that $N$ is a maximum of $D$ by \cref{finite_real_maximum}.
                Thus $N\in\naturalsPlus$ by \cref{subseteq}.
                For all $n\in\naturalsPlus$ if $N\leq n$ then $s(n)\in U$
                    by \cref{function_apply_elim,img_iff,subseteq,real_naturals_subset_reals,real_maximum,real_leq_trans,choice_function,fst_eq,snd_eq,inters_intro,real_lt_subst_right,unions_iff,real_lt_subst_left}.
                Follows by definition.
            \end{subproof}
            $s$ converges to $x$ in $\productFamily{X}$ with respect to $T_0$ by \cref{converges_to}.
        \end{byCase}
    \end{subproof}
    $s$ converges to $x$ in $\productFamily{X}$ with respect to $T_0$ iff
        for all $\alpha\in\dom{X}$ we have
            $\piIndex{X}{\alpha}\circ s$ converges to $\apply{x}{\alpha}$ in $\apply{X}{\alpha}$ with respect to $\apply{T}{\alpha}$.
\end{proof}

% from §43 helper lemma 5: the real line is complete in the standard metric
\begin{lemma}\label{reals_complete_standard_metric}
    $\reals$ is complete with respect to $\standardMetricR$.
\end{lemma}
\begin{proof}
    $\standardMetricR$ is a metric on $\reals$ by \cref{standard_metric_is_metric}.
    Show that for all $s$ if
        $s$ is a Cauchy sequence in $\reals$ with respect to $\standardMetricR$
    then there exist $t,x$ such that
        $t$ is a sequence in $\reals$ and
        $t$ is a subsequence of $s$ and
        $x\in\reals$ and
        $t$ converges to $x$ in $\reals$ with respect to $\metricTopology{\standardMetricR}{\reals}$.
    \begin{subproof}
        Fix $s$ such that $s$ is a Cauchy sequence in $\reals$ with respect to $\standardMetricR$.
        $s$ is a sequence in $\reals$ by \cref{cauchy_sequence}.
        $\ran{s}$ is bounded in $\reals$ with respect to $\standardMetricR$
            by \cref{standard_metric_is_metric,cauchy_sequence_bounded_range}.
        Take $a,b$ such that
            $a\in\reals$ and
            $b\in\reals$ and
            $a \leq b$ and
            $\ran{s}\subseteq\intervalclosed{a}{b}$ by \cref{bounded_reals_subset_closed_interval}.
        Let $I = \intervalclosed{a}{b}$.
        Let $T_I = \subspaceTopology{\standardTopologyR}{I}$.
        $I$ is compact with respect to $T_I$ by \cref{real_intervalclosed_compact_standard}.
        $\standardTopologyR$ is a topology on $\reals$
            by \cref{standard_basis_is_basis,basis_topology_is_topology,standard_topology_reals}.
        $\metricTopology{\standardMetricR}{\reals} = \standardTopologyR$ by \cref{standard_metric_topology}.
        $\reals$ is metrizable with respect to $\standardTopologyR$ by \cref{standard_metric_is_metric,metrizable}.
        $I\subseteq\reals$ by \cref{intervalclosed,subseteq}.
        $I$ is a subspace of $\reals$ with respect to $\standardTopologyR$ by \cref{subspace_of}.
        $T_I$ is a topology on $I$ by \cref{subspace_of,subspace_topology_is_topology}.
        $I$ is metrizable with respect to $T_I$ by \cref{subspace_metrizable}.
        Hence $I$ is sequentially compact with respect to $T_I$
            by \cref{metrizable_compact_iff_limit_point_compact_iff_sequentially_compact}.
        $s$ is a function and $\dom{s} = \naturalsPlus$ by \cref{sequence_in,funs_is_function,funs_elim}.
        For all $n\in\naturalsPlus$ we have $s(n)\in I$ by \cref{subseteq,function_apply_elim,ran_intro}.
        $s\in\funs{\naturalsPlus}{I}$ by \cref{funs_intro}.
        $s$ is a sequence in $I$ by \cref{sequence_in}.
        Take $t,x$ such that
            $t$ is a sequence in $I$ and
            $t$ is a subsequence of $s$ and
            $x\in I$ and
            $t$ converges to $x$ in $I$ with respect to $T_I$ by \cref{sequentially_compact}.
        $x\in\reals$ by \cref{subseteq}.
        $t$ is a sequence in $\reals$ by \cref{sequence_in,funs_weaken_codom,subseteq}.
        For all $U$ if $U$ is a neighborhood of $x$ in $\reals$ with respect to $\standardTopologyR$ then
            there exists $N\in\naturalsPlus$ such that for all $n\in\naturalsPlus$ if $N\leq n$ then $t(n)\in U$
            by \cref{neighborhood,open_in,subspace_of,subspace_topology,subspace_topology_is_topology,inter_intro,converges_to,inter}.
        $t$ converges to $x$ in $\reals$ with respect to $\standardTopologyR$ by \cref{converges_to}.
        $t$ converges to $x$ in $\reals$ with respect to $\metricTopology{\standardMetricR}{\reals}$
            by \cref{standard_metric_topology}.
        Follows by definition.
    \end{subproof}
    $\reals$ is complete with respect to $\standardMetricR$ by \cref{cauchy_subsequence_completeness_criterion,complete_metric_space}.
\end{proof}

% from §43 helper lemma 6: omega metric dominates each weighted coordinate distance
\begin{lemma}\label{omega_metric_coordinate_bounded_distance_leq}
    Suppose $x\in\realsPower{\naturalsPlus}$.
    Suppose $y\in\realsPower{\naturalsPlus}$.
    Suppose $i\in\naturalsPlus$.
    Then $\apply{\boundedMetric{\standardMetricR}}{(x(i),y(i))} / i \leq \apply{\omegaMetric}{(x,y)}$.
\end{lemma}
\begin{proof}
    Let $A = \omegaMetricSet{x}{y}$.
    $\omegaMetric$ is a function and
        $\dom{\omegaMetric} = \realsPower{\naturalsPlus}\times\realsPower{\naturalsPlus}$
        by \cref{omega_metric_is_metric,metric_on,funs_is_function,funs_elim}.
    $((x,y),\apply{\omegaMetric}{(x,y)})\in\omegaMetric$
        by \cref{times_tuple_intro,subseteq,function_apply_elim}.
    There exist $p,x_0,y_0,r$ such that
        $p\in\realsPower{\naturalsPlus}\times\realsPower{\naturalsPlus}$ and
        $x_0\in\realsPower{\naturalsPlus}$ and
        $y_0\in\realsPower{\naturalsPlus}$ and
        $r\in\reals$ and
        $((x,y),\apply{\omegaMetric}{(x,y)}) = (p,r)$ and
        $p = (x_0,y_0)$ and
        $\realSupremum{r}{\omegaMetricSet{x_0}{y_0}}$ by \cref{omega_metric}.
    We have $p = (x,y)$ and $r = \apply{\omegaMetric}{(x,y)}$ by \cref{pair_eq_iff}.
    $(x,y) = (x_0,y_0)$ by \cref{pair_eq_iff}.
    $x = x_0$ and $y = y_0$ by \cref{pair_eq_iff}.
    $\apply{\omegaMetric}{(x,y)}$ is an upper bound of $A$
        by \cref{real_supremum_pred,real_lt_subst_right,real_supremum}.
    $\apply{\boundedMetric{\standardMetricR}}{(x(i),y(i))} / i\in A$ by \cref{omega_metric_set}.
    Hence $\apply{\boundedMetric{\standardMetricR}}{(x(i),y(i))} / i \leq \apply{\omegaMetric}{(x,y)}$
        by \cref{real_upper_bound}.
\end{proof}

% from §43 helper lemma 7: small omega distance implies small coordinate distance
\begin{lemma}\label{omega_metric_small_implies_coordinate_small}
    Suppose $x\in\realsPower{\naturalsPlus}$.
    Suppose $y\in\realsPower{\naturalsPlus}$.
    Suppose $i\in\naturalsPlus$.
    Suppose $\epsilon\in\reals$.
    Suppose $0 < \epsilon$.
    Suppose $\epsilon < 1$.
    Suppose $\apply{\omegaMetric}{(x,y)} < \epsilon / i$.
    Then $\apply{\standardMetricR}{(x(i),y(i))} < \epsilon$.
\end{lemma}
\begin{proof}
    Let $a = \apply{\boundedMetric{\standardMetricR}}{(x(i),y(i))}$.
    $x\in\funs{\naturalsPlus}{\reals}$ and $y\in\funs{\naturalsPlus}{\reals}$ by \cref{reals_power,tuple_power}.
    $x(i)\in\reals$ and $y(i)\in\reals$ by \cref{funs_type_apply}.
    $\boundedMetric{\standardMetricR}$ is a metric on $\reals$ by \cref{bounded_metric_is_metric,standard_metric_is_metric}.
    $\boundedMetric{\standardMetricR}\in\funs{\reals\times\reals}{\reals}$ by \cref{metric_on}.
    $a\in\reals$ by \cref{funs_elim,times_tuple_intro}.
    $0 < i$ by \cref{naturalsplus_positive_real}.
    $i\in\reals$ by \cref{real_naturals_subset_reals,subseteq}.
    $i\neq 0$ by \cref{real_pos_nonzero}.
    $a / i\in\reals$ and $\epsilon / i\in\reals$ by \cref{real_div,real_mul_inverse,real_mul_closed}.
    $a / i < \epsilon / i$
        by \cref{omega_metric_coordinate_bounded_distance_leq,omega_metric_is_metric,metric_on,funs_type_apply,times_tuple_intro,real_leq_lt_trans}.
    Hence $(a / i) \rmul i < (\epsilon / i) \rmul i$ by \cref{real_lt_mul_pos}.
    $a < \epsilon$ by \cref{real_div,real_mul_assoc,real_mul_inverse,real_mul_ident,real_mul_comm,real_lt_subst_left,real_lt_subst_right}.
    $a < 1$ by \cref{real_naturals_subset_reals,real_one_in_naturals,subseteq,real_lt_trans}.
    $a = \apply{\standardMetricR}{(x(i),y(i))}$
        by \cref{bounded_metric_apply_value_lt_one,standard_metric_is_metric,times_tuple_intro}.
    Hence $\apply{\standardMetricR}{(x(i),y(i))} < \epsilon$ by \cref{real_lt_subst_left}.
\end{proof}

% from §43 helper lemma 8: coordinate projections of omega-Cauchy sequences are Cauchy
\begin{lemma}\label{omega_cauchy_coordinatewise}
    Assume $s$ is a Cauchy sequence in $\realsPower{\naturalsPlus}$ with respect to $\omegaMetric$.
    Suppose $i\in\naturalsPlus$.
    Let $X = \{(j,\reals) \mid j\in\naturalsPlus\}$.
    Then $\piIndex{X}{i}\circ s$ is a Cauchy sequence in $\reals$ with respect to $\standardMetricR$.
\end{lemma}
\begin{proof}
    $s$ is a sequence in $\realsPower{\naturalsPlus}$ and $\omegaMetric$ is a metric on $\realsPower{\naturalsPlus}$
        by \cref{cauchy_sequence}.
    $1\in\naturalsPlus$ by \cref{real_one_in_naturals}.
    $\naturalsPlus$ is inhabited.
    $\productFamily{X} = \realsPower{\naturalsPlus}$ by \cref{product_family_reals_power}.
    $s$ is a function and $\dom{s} = \naturalsPlus$ by \cref{sequence_in,funs_is_function,funs_elim}.
    For all $n\in\naturalsPlus$ we have $s(n)\in\productFamily{X}$
        by \cref{funs_type_apply,sequence_in,product_family_reals_power,subseteq}.
    $s\in\funs{\naturalsPlus}{\productFamily{X}}$ by \cref{funs_intro}.
    $s$ is a sequence in $\productFamily{X}$ by \cref{sequence_in}.
    For all $j,y_1,y_2$ such that $(j,y_1)\in X$ and $(j,y_2)\in X$ we have $y_1 = y_2$ by \cref{pair_eq_iff}.
    $X$ is a function by \cref{relation,rightunique}.
    $(i,\reals)\in X$ by \cref{pair_eq_iff}.
    $i\in\dom{X}$ by \cref{dom_intro}.
    $X(i) = \reals$ by \cref{function_apply_intro}.
    $\piIndex{X}{i}\in\funs{\productFamily{X}}{\apply{X}{i}}$ by \cref{projection_map_indexed_function}.
    $\piIndex{X}{i}\in\funs{\productFamily{X}}{\reals}$ by \cref{real_lt_subst_right}.
    $\piIndex{X}{i}$ is a function and $\dom{\piIndex{X}{i}} = \productFamily{X}$
        by \cref{funs_is_function,funs_elim}.
    $\ran{s}\subseteq\productFamily{X}$ by \cref{funs_ran}.
    $\ran{s}\subseteq\dom{\piIndex{X}{i}}$ by \cref{funs_elim,subseteq}.
    $\piIndex{X}{i}\circ s$ is a sequence in $\reals$ by \cref{funs_circ,sequence_in}.
    $\standardMetricR$ is a metric on $\reals$ by \cref{standard_metric_is_metric}.
    Show that for all $\epsilon\in\reals$ such that $0 < \epsilon$
        there exists $N\in\naturalsPlus$ such that
            for all $n,m\in\naturalsPlus$ if $N \leq n$ and $N \leq m$ then
                $\apply{\standardMetricR}{((\piIndex{X}{i}\circ s)(n),(\piIndex{X}{i}\circ s)(m))} < \epsilon$.
    \begin{subproof}
        Fix $\epsilon\in\reals$.
        Assume $0 < \epsilon$.
        $i\in\reals$ by \cref{real_naturals_subset_reals,subseteq}.
        $0 < i$ by \cref{naturalsplus_positive_real}.
        $i\neq 0$ by \cref{real_pos_nonzero}.
        $\inv{i}\in\reals$ by \cref{real_mul_inverse}.
        $0 < \inv{i}$ by \cref{real_inv_pos}.
        $0\in\reals$ by \cref{real_add_ident}.
        $1\in\reals$ by \cref{real_naturals_subset_reals,real_one_in_naturals,subseteq}.
        Let $\eta = 1 / (1 + 1)$.
        $\eta\in\reals$ and $0 < \eta$ and $\eta + \eta = 1$
            by \cref{real_naturals_subset_reals,real_one_in_naturals,subseteq,real_zero_lt_one,real_half_of_positive}.
        $\eta < 1$ by \cref{real_lt_add,real_add_ident,real_lt_subst_left}.
        \begin{byCase}
        \caseOf{$\epsilon < 1$.}
            $\epsilon / i\in\reals$ by \cref{real_div,real_mul_inverse,real_mul_closed}.
            $0 \rmul \inv{i} < \epsilon \rmul \inv{i}$ by \cref{real_lt_mul_pos}.
            $0 \rmul \inv{i} = 0$ and $\epsilon \rmul \inv{i} = \epsilon / i$ by \cref{real_mul_zero,real_div}.
            $0 < \epsilon / i$ by \cref{real_lt_subst_left,real_lt_subst_right}.
            Take $N\in\naturalsPlus$ such that
                for all $n,m\in\naturalsPlus$ if $N \leq n$ and $N \leq m$ then
                    $\apply{\omegaMetric}{(s(n),s(m))} < \epsilon / i$ by \cref{cauchy_sequence}.
            For all $n,m$ such that $n\in\naturalsPlus$ and $m\in\naturalsPlus$ and $N \leq n$ and $N \leq m$ we have
                $\apply{\standardMetricR}{((\piIndex{X}{i}\circ s)(n),(\piIndex{X}{i}\circ s)(m))} < \epsilon$
                by \cref{cauchy_sequence,sequence_in,funs_type_apply,omega_metric_small_implies_coordinate_small,circ_apply,subseteq,projection_map_indexed,product_family_reals_power,function_apply_intro,real_lt_subst_left,real_lt_subst_right}.
            Follows by definition.
        \caseOf{it is wrong that $\epsilon < 1$.}
            $1 \leq \epsilon$ by \cref{real_lt_trichotomy}.
            $\eta < \epsilon$ by \cref{real_lt_leq_trans_local}.
            $\eta / i\in\reals$ by \cref{real_div,real_mul_inverse,real_mul_closed}.
            $0 \rmul \inv{i} < \eta \rmul \inv{i}$ by \cref{real_lt_mul_pos}.
            $0 \rmul \inv{i} = 0$ and $\eta \rmul \inv{i} = \eta / i$ by \cref{real_mul_zero,real_div}.
            $0 < \eta / i$ by \cref{real_lt_subst_left,real_lt_subst_right}.
            Take $N\in\naturalsPlus$ such that
                for all $n,m\in\naturalsPlus$ if $N \leq n$ and $N \leq m$ then
                    $\apply{\omegaMetric}{(s(n),s(m))} < \eta / i$ by \cref{cauchy_sequence}.
            For all $n,m$ such that $n\in\naturalsPlus$ and $m\in\naturalsPlus$ and $N \leq n$ and $N \leq m$ we have
                $\apply{\standardMetricR}{((\piIndex{X}{i}\circ s)(n),(\piIndex{X}{i}\circ s)(m))} < \epsilon$
                by \cref{cauchy_sequence,sequence_in,funs_type_apply,omega_metric_small_implies_coordinate_small,circ_apply,subseteq,projection_map_indexed,product_family_reals_power,function_apply_intro,reals_power,tuple_power,times_tuple_intro,metric_on,standard_metric_is_metric,real_lt_trans,real_lt_subst_left,real_lt_subst_right}.
            Follows by definition.
        \end{byCase}
    \end{subproof}
    $\piIndex{X}{i}\circ s$ is a Cauchy sequence in $\reals$ with respect to $\standardMetricR$
        by \cref{cauchy_sequence}.
\end{proof}

% from §43 Item 5: $\reals^\omega$ is complete
\begin{theorem}\label{reals_omega_complete}
    $\realsPower{\naturalsPlus}$ is complete with respect to $\omegaMetric$.
\end{theorem}
\begin{proof}
    $\omegaMetric$ is a metric on $\realsPower{\naturalsPlus}$ by \cref{omega_metric_is_metric}.
    Show that for all $s$ if $s$ is a Cauchy sequence in $\realsPower{\naturalsPlus}$ with respect to $\omegaMetric$
        then there exists $x\in\realsPower{\naturalsPlus}$ such that
            $s$ converges to $x$ in $\realsPower{\naturalsPlus}$ with respect to $\metricTopology{\omegaMetric}{\realsPower{\naturalsPlus}}$.
    \begin{subproof}
        Fix $s$ such that $s$ is a Cauchy sequence in $\realsPower{\naturalsPlus}$ with respect to $\omegaMetric$.
        Let $X = \{(j,\reals) \mid j\in\naturalsPlus\}$.
        Let $T = \{(j,\standardTopologyR) \mid j\in\naturalsPlus\}$.
        Let $T_0 = \productTopologyIndexed{T}{X}$.
        $1\in\naturalsPlus$ by \cref{real_one_in_naturals}.
        $\naturalsPlus$ is inhabited.
        $\productFamily{X} = \realsPower{\naturalsPlus}$ by \cref{product_family_reals_power}.
        $s$ is a function and $\dom{s} = \naturalsPlus$ by \cref{cauchy_sequence,sequence_in,funs_is_function,funs_elim}.
        For all $n\in\naturalsPlus$ we have $s(n)\in\productFamily{X}$
            by \cref{cauchy_sequence,sequence_in,funs_type_apply,product_family_reals_power,subseteq}.
        $s\in\funs{\naturalsPlus}{\productFamily{X}}$ by \cref{funs_intro}.
        $s$ is a sequence in $\productFamily{X}$ by \cref{sequence_in}.

        For all $j,y_1,y_2$ such that $(j,y_1)\in X$ and $(j,y_2)\in X$ we have $y_1 = y_2$ by \cref{pair_eq_iff}.
        $X$ is a function by \cref{relation,rightunique}.
        $\dom{X} = \naturalsPlus$ by \cref{dom_intro,dom_iff,pair_eq_iff,subseteq,subseteq_antisymmetric}.

        For all $j,U_1,U_2$ such that $(j,U_1)\in T$ and $(j,U_2)\in T$ we have $U_1 = U_2$ by \cref{pair_eq_iff}.
        $T$ is a function by \cref{relation,rightunique}.
        $\dom{T} = \naturalsPlus$ by \cref{dom_intro,dom_iff,pair_eq_iff,subseteq,subseteq_antisymmetric}.

        For all $i\in\dom{X}$ we have $\apply{T}{i}$ is a topology on $\apply{X}{i}$
            by \cref{function_apply_intro,standard_basis_is_basis,standard_topology_reals,basis_topology_is_topology,real_lt_subst_left,real_lt_subst_right}.

        Let $R = \{w\in\naturalsPlus\times\reals \mid
            \text{$\piIndex{X}{\fst{w}}\circ s$ converges to $\snd{w}$ in $\reals$ with respect to $\standardTopologyR$}\}$.
        $R$ is a relation by \cref{relation,subseteq,times_elem_is_tuple}.

        Show that for all $i\in\naturalsPlus$ there exists $a$ such that $i\mathrel{R} a$.
        \begin{subproof}
            Follows by \cref{omega_cauchy_coordinatewise,reals_complete_standard_metric,complete_metric_space,standard_metric_topology,times_tuple_intro,fst_eq,snd_eq}.
        \end{subproof}

        Take $x$ such that $x$ is a function on $\naturalsPlus$ and for all $i\in\naturalsPlus$ we have $i\mathrel{R}\apply{x}{i}$
            by \cref{choice_function}.
        $x\in\productFamily{X}$ by \cref{choice_function,relation,subseteq,times_tuple_elim,funs_intro,reals_power,tuple_power,product_family_reals_power}.

        Show that for all $i\in\dom{X}$ we have
            $\piIndex{X}{i}\circ s$ converges to $\apply{x}{i}$ in $\apply{X}{i}$ with respect to $\apply{T}{i}$.
        \begin{subproof}
            Follows by \cref{choice_function,fst_eq,snd_eq,function_apply_intro,converges_to,real_lt_subst_left,real_lt_subst_right}.
        \end{subproof}

        $s$ converges to $x$ in $\realsPower{\naturalsPlus}$ with respect to $\metricTopology{\omegaMetric}{\realsPower{\naturalsPlus}}$
            by \cref{product_sequence_converges_iff_coordinatewise,product_family_reals_power,omega_metric_product_topology}.
        Follows by definition.
    \end{subproof}
    $\realsPower{\naturalsPlus}$ is complete with respect to $\omegaMetric$ by \cref{complete_metric_space}.
\end{proof}

% from §43 Item 6: uniform metric on a function space
\begin{definition}\label{uniform_metric_on_function_space}
    $\rho$ is the uniform metric on functions from $J$ to $Y$ with respect to $d$ iff
        $d$ is a metric on $Y$ and
        $\rho$ is a metric on $\funs{J}{Y}$ and
        for all $f,g\in\funs{J}{Y}$ there exists $A$ such that
            $A = \{r\in\reals \mid \exists \alpha\in J.
                ((r = d((\apply{f}{\alpha},\apply{g}{\alpha})) \land d((\apply{f}{\alpha},\apply{g}{\alpha})) < 1) \lor
                (r = 1 \land 1 \leq d((\apply{f}{\alpha},\apply{g}{\alpha}))))\}$ and
            $\rho((f,g))$ is a supremum of $A$.
\end{definition}

% from §43 helper lemma 8: completeness is preserved by the bounded metric
\begin{lemma}\label{complete_bounded_metric}
    Assume $d$ is a metric on $Y$.
    Suppose $Y$ is complete with respect to $d$.
    Then $Y$ is complete with respect to $\boundedMetric{d}$.
\end{lemma}
\begin{proof}
    Let $e = \boundedMetric{d}$.
    $e$ is a metric on $Y$ by \cref{bounded_metric_is_metric}.
    Show that for all $s$ if
        $s$ is a Cauchy sequence in $Y$ with respect to $e$
    then $s$ is a Cauchy sequence in $Y$ with respect to $d$.
    \begin{subproof}
        Fix $s$ such that $s$ is a Cauchy sequence in $Y$ with respect to $e$.
        Show that for all $\epsilon$ if $\epsilon\in\reals$ and $0 < \epsilon$ then
            there exists $N\in\naturalsPlus$ such that
                for all $n,m\in\naturalsPlus$ if $N \leq n$ and $N \leq m$ then
                    $d((s(n),s(m))) < \epsilon$.
        \begin{subproof}
            Fix $\epsilon$ such that $\epsilon\in\reals$ and $0 < \epsilon$.
            \begin{byCase}
            \caseOf{$\epsilon < 1$.}
                Take $N\in\naturalsPlus$ such that
                    for all $n,m\in\naturalsPlus$ if $N \leq n$ and $N \leq m$ then
                        $e((s(n),s(m))) < \epsilon$ by \cref{cauchy_sequence}.
                Show that for all $n,m\in\naturalsPlus$ if $N \leq n$ and $N \leq m$ then
                    $d((s(n),s(m))) < \epsilon$.
                \begin{subproof}
                    Follows by \cref{cauchy_sequence,sequence_in,funs_type_apply,times_tuple_intro,metric_on,real_naturals_subset_reals,real_one_in_naturals,subseteq,real_lt_trans,bounded_metric_apply_value_lt_one,real_lt_subst_left}.
                \end{subproof}
            \caseOf{it is wrong that $\epsilon < 1$.}
                $1\in\reals$ by \cref{real_naturals_subset_reals,real_one_in_naturals,subseteq}.
                $1 \leq \epsilon$ by \cref{real_lt_trichotomy}.
                Take $N\in\naturalsPlus$ such that
                    for all $n,m\in\naturalsPlus$ if $N \leq n$ and $N \leq m$ then
                        $e((s(n),s(m))) < 1$ by \cref{real_zero_lt_one,cauchy_sequence}.
                Show that for all $n,m\in\naturalsPlus$ if $N \leq n$ and $N \leq m$ then
                    $d((s(n),s(m))) < \epsilon$.
                \begin{subproof}
                    Follows by \cref{cauchy_sequence,sequence_in,funs_type_apply,times_tuple_intro,bounded_metric_apply_value_lt_one,real_lt_subst_left,metric_on,real_lt_trans}.
                \end{subproof}
            \end{byCase}
        \end{subproof}
        $s$ is a Cauchy sequence in $Y$ with respect to $d$ by \cref{cauchy_sequence}.
    \end{subproof}
    Show that for all $s$ if
        $s$ is a Cauchy sequence in $Y$ with respect to $e$
    then there exists $x\in Y$ such that
        $s$ converges to $x$ in $Y$ with respect to $\metricTopology{e}{Y}$.
    \begin{subproof}
        Follows by \cref{complete_metric_space,bounded_metric_topology,converges_to}.
    \end{subproof}
    $Y$ is complete with respect to $\boundedMetric{d}$ by \cref{complete_metric_space}.
\end{proof}

% from §43 helper lemma 9: each coordinate bounded distance is bounded by the uniform metric
\begin{lemma}\label{uniform_function_metric_coordinate_leq}
    Assume $\rho$ is the uniform metric on functions from $J$ to $Y$ with respect to $d$.
    Suppose $f\in\funs{J}{Y}$.
    Suppose $g\in\funs{J}{Y}$.
    Suppose $\alpha\in J$.
    Then $\apply{\boundedMetric{d}}{(f(\alpha),g(\alpha))} \leq \rho((f,g))$.
\end{lemma}
\begin{proof}
    $d$ is a metric on $Y$ and $\rho$ is a metric on $\funs{J}{Y}$ by \cref{uniform_metric_on_function_space}.
    Take $A$ such that
        $A = \{r\in\reals \mid \exists \beta\in J.
            ((r = d((\apply{f}{\beta},\apply{g}{\beta})) \land d((\apply{f}{\beta},\apply{g}{\beta})) < 1) \lor
            (r = 1 \land 1 \leq d((\apply{f}{\beta},\apply{g}{\beta}))))\}$ and
        $\rho((f,g))$ is a supremum of $A$ by \cref{uniform_metric_on_function_space}.
    Let $r = \apply{\boundedMetric{d}}{(f(\alpha),g(\alpha))}$.
    $f(\alpha)\in Y$ and $g(\alpha)\in Y$ by \cref{funs_type_apply}.
    $(f(\alpha),g(\alpha))\in Y\times Y$ by \cref{times_tuple_intro}.
    $\boundedMetric{d}$ is a metric on $Y$ by \cref{bounded_metric_is_metric}.
    $\boundedMetric{d}\in\funs{Y\times Y}{\reals}$ by \cref{metric_on}.
    $r\in\reals$ by \cref{bounded_metric_is_metric,metric_on,funs_type_apply}.
    Show that $r\in A$.
    \begin{subproof}
        \begin{byCase}
        \caseOf{$d((f(\alpha),g(\alpha))) < 1$.}
            $r = d((f(\alpha),g(\alpha)))$ by \cref{bounded_metric_apply_lt_one}.
            $r\in A$ by \cref{uniform_metric_on_function_space}.
        \caseOf{it is wrong that $d((f(\alpha),g(\alpha))) < 1$.}
            $r = 1$ by \cref{bounded_metric_apply_geq_one}.
            $1 \leq d((f(\alpha),g(\alpha)))$
                by \cref{metric_on,funs_type_apply,real_naturals_subset_reals,real_one_in_naturals,subseteq,real_lt_trichotomy,real_lt_imp_leq,real_lt_subst_right}.
            $r\in A$ by \cref{uniform_metric_on_function_space}.
        \end{byCase}
    \end{subproof}
    $\rho((f,g))$ is an upper bound of $A$ by \cref{real_supremum}.
    $r \leq \rho((f,g))$ by \cref{real_upper_bound}.
\end{proof}

% from §43 helper lemma 10a: positive metric balls are neighborhoods
\begin{lemma}\label{metric_ball_neighborhood_metric_topology}
    Assume $d$ is a metric on $X$.
    Suppose $x\in X$.
    Suppose $\epsilon\in\reals$.
    Suppose $0 < \epsilon$.
    Then $\metricBall{d}{X}{x}{\epsilon}$ is a neighborhood of $x$ in $X$ with respect to $\metricTopology{d}{X}$.
\end{lemma}
\begin{proof}
    $\metricBall{d}{X}{x}{\epsilon}\in\metricBasis{d}{X}$ by \cref{metric_ball,subseteq,powerset_intro,metric_basis}.
    $\metricBall{d}{X}{x}{\epsilon}$ is open in $X$ with respect to $\metricTopology{d}{X}$
        by \cref{metric_basis_is_basis,basis_elem_open_in_generated,metric_topology,basis_for_topology,basis_topology_is_topology,open_in}.
    $x\in\metricBall{d}{X}{x}{\epsilon}$ by \cref{metric_zero_characterization,metric_on,real_lt_subst_left,metric_ball}.
    $\metricBall{d}{X}{x}{\epsilon}$ is a neighborhood of $x$ in $X$ with respect to $\metricTopology{d}{X}$
        by \cref{neighborhood}.
\end{proof}

% from §43 helper lemma 10: eventual metric bounds pass to the limit
\begin{lemma}\label{metric_limit_eventual_bound}
    Assume $e$ is a metric on $Y$.
    Suppose $y\in Y$.
    Suppose $z\in Y$.
    Suppose $t$ is a sequence in $Y$.
    Suppose $t$ converges to $z$ in $Y$ with respect to $\metricTopology{e}{Y}$.
    Suppose $\delta\in\reals$.
    Suppose $0 < \delta$.
    Suppose there exists $N\in\naturalsPlus$ such that
        for all $n\in\naturalsPlus$ if $N \leq n$ then
            $e((y,t(n))) < \delta$.
    Then $e((y,z)) \leq \delta$.
\end{lemma}
\begin{proof}
    Suppose not.
    $e((y,z))\in\reals$ by \cref{times_tuple_intro,metric_on,funs_type_apply}.
    $\delta < e((y,z))$ by \cref{real_lt_trichotomy}.
    Let $\epsilon = e((y,z)) - \delta$.
    $\epsilon\in\reals$ by \cref{real_sub,real_add_inverse,real_add_closed}.
    $\delta + \neg{\delta} < e((y,z)) + \neg{\delta}$ by \cref{real_add_inverse,real_lt_add}.
    $\delta + \neg{\delta} = 0$ and $e((y,z)) + \neg{\delta} = e((y,z)) - \delta$
        by \cref{real_add_inverse,real_add_ident,real_sub}.
    $0 < \epsilon$ by \cref{real_lt_subst_left,real_lt_subst_right}.
    Take $N_0\in\naturalsPlus$ such that
        for all $n\in\naturalsPlus$ if $N_0 \leq n$ then
            $e((y,t(n))) < \delta$ by assumption.
    $\metricBall{e}{Y}{z}{\epsilon}$ is a neighborhood of $z$ in $Y$ with respect to $\metricTopology{e}{Y}$
        by \cref{metric_ball_neighborhood_metric_topology}.
    Take $N_1\in\naturalsPlus$ such that
        for all $n\in\naturalsPlus$ if $N_1 \leq n$ then
            $t(n)\in\metricBall{e}{Y}{z}{\epsilon}$ by \cref{converges_to}.
    Let $n = N_0 + N_1$.
    $n\in\naturalsPlus$ by \cref{naturalsplus_add_closed}.
    $N_0\in\reals$ and $N_1\in\reals$ and $n\in\reals$ by \cref{real_naturals_subset_reals,subseteq}.
    $0\in\reals$ by \cref{real_add_ident}.
    $1\in\reals$ by \cref{real_naturals_subset_reals,real_one_in_naturals,subseteq}.
    $0 < N_0$ and $0 < N_1$ by \cref{naturalsplus_positive_real}.
    $0 + N_0 < N_1 + N_0$ by \cref{real_lt_add}.
    $0 + N_0 = N_0$ and $N_1 + N_0 = N_0 + N_1$ by \cref{real_add_ident,real_add_comm}.
    $N_0 < N_0 + N_1$ by \cref{real_lt_subst_left,real_lt_subst_right}.
    $N_0 < n$ by \cref{real_lt_subst_right}.
    $N_0 \leq n$ by \cref{real_lt_imp_leq}.
    $0 + N_1 < N_0 + N_1$ by \cref{real_lt_add}.
    $0 + N_1 = N_1$ by \cref{real_add_ident}.
    $N_1 < N_0 + N_1$ by \cref{real_lt_subst_left,real_lt_subst_right}.
    $N_1 < n$ by \cref{real_lt_subst_right}.
    $N_1 \leq n$ by \cref{real_lt_imp_leq}.
    $e((y,t(n))) < \delta$ and $t(n)\in\metricBall{e}{Y}{z}{\epsilon}$ by \cref{subseteq}.
    $t(n)\in Y$ and $e((z,t(n))) < \epsilon$ by \cref{metric_ball}.
    $e((t(n),z)) < \epsilon$ by \cref{metric_symmetric,metric_on,real_lt_subst_left}.
    $(y,z)\in Y\times Y$ and $(y,t(n))\in Y\times Y$ and $(t(n),z)\in Y\times Y$ by \cref{times_tuple_intro}.
    $e((y,t(n)))\in\reals$ and $e((t(n),z))\in\reals$ by \cref{metric_on,funs_type_apply}.
    $e((y,z)) \leq e((y,t(n))) + e((t(n),z))$ by \cref{metric_on,metric_triangle_inequality}.
    $e((y,t(n))) + e((t(n),z))\in\reals$ and $\delta + e((t(n),z))\in\reals$ and $\delta + \epsilon\in\reals$
        by \cref{real_add_closed}.
    $e((y,t(n))) + e((t(n),z)) < \delta + e((t(n),z))$ by \cref{real_lt_add}.
    $e((t(n),z)) + \delta < \epsilon + \delta$ by \cref{real_lt_add}.
    $e((t(n),z)) + \delta = \delta + e((t(n),z))$ and $\epsilon + \delta = \delta + \epsilon$ by \cref{real_add_comm}.
    $\delta + e((t(n),z)) < \delta + \epsilon$ by \cref{real_lt_subst_left,real_lt_subst_right}.
    $e((y,t(n))) + e((t(n),z)) < \delta + \epsilon$ by \cref{real_lt_trans}.
    $\delta + \epsilon = \delta + (e((y,z)) - \delta)$ by \cref{real_sub}.
    $\delta + \epsilon = e((y,z))$ by \cref{real_add_assoc,real_add_inverse,real_add_ident,real_add_comm}.
    $e((y,t(n))) + e((t(n),z)) < e((y,z))$ by \cref{real_lt_subst_right}.
    $e((y,z)) < e((y,z))$ by \cref{real_leq_lt_trans}.
    Contradiction by \cref{real_lt_irreflexive}.
\end{proof}

% from §43 helper lemma 11: convergence in the uniform metric is uniform convergence for the bounded metric
\begin{lemma}\label{uniform_metric_convergence_uniform}
    Assume $d$ is a metric on $Y$.
    Assume $\rho$ is the uniform metric on functions from $J$ to $Y$ with respect to $d$.
    Suppose $s$ is a sequence in $\funs{J}{Y}$.
    Suppose $f\in\funs{J}{Y}$.
    Suppose $s$ converges to $f$ in $\funs{J}{Y}$ with respect to $\metricTopology{\rho}{\funs{J}{Y}}$.
    Then $s$ converges uniformly to $f$ from $J$ to $Y$ with respect to $\boundedMetric{d}$.
\end{lemma}
\begin{proof}
    Let $e = \boundedMetric{d}$.
    $e$ is a metric on $Y$ by \cref{bounded_metric_is_metric}.
    Show that for all $\epsilon$ if $\epsilon\in\reals$ and $0 < \epsilon$ then
        there exists $N\in\naturalsPlus$ such that for all $n\in\naturalsPlus$ if $N\leq n$ then
            for all $\alpha\in J$ we have $e((\apply{\apply{s}{n}}{\alpha}, f(\alpha))) < \epsilon$.
    \begin{subproof}
        Fix $\epsilon$ such that $\epsilon\in\reals$ and $0 < \epsilon$.
        Let $F = \funs{J}{Y}$.
        Take $N\in\naturalsPlus$ such that for all $n\in\naturalsPlus$ if $N\leq n$ then
            $\rho((s(n),f)) < \epsilon$
            by \cref{uniform_metric_on_function_space,metric_ball_neighborhood_metric_topology,converges_to,metric_ball,subseteq,metric_on,metric_symmetric,real_lt_subst_left}.
        Show that for all $n\in\naturalsPlus$ if $N\leq n$ then
            for all $\alpha\in J$ we have $e((\apply{\apply{s}{n}}{\alpha}, f(\alpha))) < \epsilon$.
        \begin{subproof}
            Follows by \cref{sequence_in,funs_type_apply,uniform_function_metric_coordinate_leq,times_tuple_intro,metric_on,uniform_metric_on_function_space,real_leq_lt_trans}.
        \end{subproof}
        Follows by definition.
    \end{subproof}
    $s$ converges uniformly to $f$ from $J$ to $Y$ with respect to $e$ by \cref{uniform_converges_to}.
\end{proof}

% from §43 helper lemma 12: Cauchy sequences in the uniform metric are coordinatewise Cauchy
\begin{lemma}\label{uniform_metric_cauchy_coordinate}
    Assume $d$ is a metric on $Y$.
    Assume $\rho$ is the uniform metric on functions from $J$ to $Y$ with respect to $d$.
    Suppose $s$ is a Cauchy sequence in $\funs{J}{Y}$ with respect to $\rho$.
    Suppose $\alpha\in J$.
    Let $t = \{(n,\apply{\apply{s}{n}}{\alpha}) \mid n\in\naturalsPlus\}$.
    Then $t$ is a Cauchy sequence in $Y$ with respect to $\boundedMetric{d}$.
\end{lemma}
\begin{proof}
    Let $e = \boundedMetric{d}$.
    $e$ is a metric on $Y$ by \cref{bounded_metric_is_metric}.
    For all $n,y_1,y_2$ such that $(n,y_1)\in t$ and $(n,y_2)\in t$ we have $y_1 = y_2$ by \cref{pair_eq_iff}.
    $t$ is a function by \cref{relation,rightunique}.
    $\dom{t} = \naturalsPlus$ by \cref{dom_intro,dom_iff,pair_eq_iff,subseteq,subseteq_antisymmetric}.
    For all $n$ such that $n\in\dom{t}$ we have $\apply{\apply{s}{n}}{\alpha}\in Y$
        by \cref{subseteq,cauchy_sequence,sequence_in,funs_type_apply}.
    For all $n$ such that $n\in\dom{t}$ we have $(n,\apply{\apply{s}{n}}{\alpha})\in t$.
    For all $n$ such that $n\in\dom{t}$ we have $t(n) = \apply{\apply{s}{n}}{\alpha}$ by \cref{function_apply_intro}.
    For all $n\in\dom{t}$ we have $t(n)\in Y$.
    $t$ is a sequence in $Y$ by \cref{funs_intro,sequence_in}.
    Show that for all $\epsilon$ if $\epsilon\in\reals$ and $0 < \epsilon$ then
        there exists $N\in\naturalsPlus$ such that
            for all $n,m\in\naturalsPlus$ if $N \leq n$ and $N \leq m$ then
                $e((t(n),t(m))) < \epsilon$.
    \begin{subproof}
        Follows by \cref{cauchy_sequence,sequence_in,funs_type_apply,uniform_function_metric_coordinate_leq,uniform_metric_on_function_space,metric_on,times_tuple_intro,real_leq_lt_trans,function_apply_intro,real_lt_subst_left,real_lt_subst_right}.
    \end{subproof}
    $t$ is a Cauchy sequence in $Y$ with respect to $e$ by \cref{cauchy_sequence}.
\end{proof}

% from §43 helper lemma 13: pointwise bounded metric bounds imply a uniform metric bound
\begin{lemma}\label{uniform_metric_leq_from_coordinate_bounds}
    Assume $\rho$ is the uniform metric on functions from $J$ to $Y$ with respect to $d$.
    Suppose $f\in\funs{J}{Y}$.
    Suppose $g\in\funs{J}{Y}$.
    Suppose $\delta\in\reals$.
    Suppose for all $\alpha\in J$ we have
        $\apply{\boundedMetric{d}}{(f(\alpha),g(\alpha))} \leq \delta$.
    Then $\rho((f,g)) \leq \delta$.
\end{lemma}
\begin{proof}
    Let $e = \boundedMetric{d}$.
    $d$ is a metric on $Y$ by \cref{uniform_metric_on_function_space}.
    $e$ is a metric on $Y$ by \cref{bounded_metric_is_metric}.
    Take $A$ such that
        $A = \{r\in\reals \mid \exists \alpha\in J.
            ((r = d((\apply{f}{\alpha},\apply{g}{\alpha})) \land d((\apply{f}{\alpha},\apply{g}{\alpha})) < 1) \lor
            (r = 1 \land 1 \leq d((\apply{f}{\alpha},\apply{g}{\alpha}))))\}$ and
        $\rho((f,g))$ is a supremum of $A$ by \cref{uniform_metric_on_function_space}.
    Show that for all $r$ if $r\in A$ then $r \leq \delta$.
    \begin{subproof}
        Fix $r$.
        Assume $r\in A$.
        Take $\alpha\in J$ such that
            $(r = d((\apply{f}{\alpha},\apply{g}{\alpha})) \land d((\apply{f}{\alpha},\apply{g}{\alpha})) < 1) \lor
            (r = 1 \land 1 \leq d((\apply{f}{\alpha},\apply{g}{\alpha})))$ by \cref{uniform_metric_on_function_space}.
        $f(\alpha)\in Y$ and $g(\alpha)\in Y$ by \cref{funs_type_apply}.
        Let $p = (f(\alpha),g(\alpha))$.
        $p\in Y\times Y$ by \cref{times_tuple_intro}.
        Show that $e((f(\alpha),g(\alpha))) = r$.
        \begin{subproof}
            \begin{byCase}
            \caseOf{$r = d((\apply{f}{\alpha},\apply{g}{\alpha})) \land d((\apply{f}{\alpha},\apply{g}{\alpha})) < 1$.}
                $e(p) = d(p)$ by \cref{bounded_metric_apply_below_one}.
                $e((f(\alpha),g(\alpha))) = r$ by \cref{real_lt_subst_right}.
            \caseOf{$r = 1 \land 1 \leq d((\apply{f}{\alpha},\apply{g}{\alpha}))$.}
                It is wrong that $d(p) < 1$ by \cref{real_leq_lt_trans,metric_on,funs_type_apply,real_lt_irreflexive}.
                $e\in\funs{Y\times Y}{\reals}$ by \cref{metric_on}.
                $e(p) = 1$ by \cref{bounded_metric_apply_geq_one}.
                $e((f(\alpha),g(\alpha))) = r$ by \cref{real_lt_subst_right}.
            \end{byCase}
        \end{subproof}
        $e((f(\alpha),g(\alpha))) \leq \delta$ by assumption.
        $r \leq \delta$.
    \end{subproof}
    $A\subseteq\reals$ by \cref{uniform_metric_on_function_space,subseteq}.
    For all $r$ if $r\in A$ then $r\mathrel{\realOrderRel}\delta$ or $r = \delta$.
    $\delta$ is an upper bound of $A$ in $\reals$ with respect to $\realOrderRel$ by \cref{upper_bound_rel}.
    $\delta$ is an upper bound of $A$ by \cref{real_order_upper_bound_is_real_upper_bound}.
    $\rho((f,g)) \leq \delta$ by \cref{real_supremum}.
\end{proof}

% from §43 Item 7: function spaces over complete metric spaces are complete in the uniform metric
\begin{theorem}\label{uniform_function_space_complete}
    Assume $d$ is a metric on $Y$.
    Assume $Y$ is complete with respect to $d$.
    Assume $\rho$ is the uniform metric on functions from $J$ to $Y$ with respect to $d$.
    Then $\funs{J}{Y}$ is complete with respect to $\rho$.
\end{theorem}
\begin{proof}
    Let $e = \boundedMetric{d}$.
    Let $F = \funs{J}{Y}$.
    $e$ is a metric on $Y$ by \cref{bounded_metric_is_metric}.
    $Y$ is complete with respect to $e$ by \cref{complete_bounded_metric}.
    $\rho$ is a metric on $F$ by \cref{uniform_metric_on_function_space}.
    Show that for all $s$ if $s$ is a Cauchy sequence in $F$ with respect to $\rho$
        then there exists $f\in F$ such that
            $s$ converges to $f$ in $F$ with respect to $\metricTopology{\rho}{F}$.
    \begin{subproof}
        Fix $s$ such that $s$ is a Cauchy sequence in $F$ with respect to $\rho$.
        Let $R = \{w\in J\times Y \mid
            \text{there exists $t$ such that
                $t$ is a sequence in $Y$ and
                for all $n\in\naturalsPlus$ we have $t(n) = \apply{\apply{s}{n}}{\fst{w}}$ and
                $t$ converges to $\snd{w}$ in $Y$ with respect to $\metricTopology{e}{Y}$}\}$.
        $R$ is a relation by \cref{relation,subseteq,times_elem_is_tuple}.
        Show that for all $\alpha\in J$ there exists $y$ such that $\alpha\mathrel{R} y$.
        \begin{subproof}
            Fix $\alpha\in J$.
            Let $t = \{(n,\apply{\apply{s}{n}}{\alpha}) \mid n\in\naturalsPlus\}$.
            $t$ is a Cauchy sequence in $Y$ with respect to $e$ by \cref{uniform_metric_cauchy_coordinate}.
            $t$ is a sequence in $Y$ by \cref{cauchy_sequence}.
            Take $y\in Y$ such that
                $t$ converges to $y$ in $Y$ with respect to $\metricTopology{e}{Y}$
                by \cref{complete_metric_space}.
            For all $n\in\naturalsPlus$ we have $t(n) = \apply{\apply{s}{n}}{\alpha}$
                by \cref{function_apply_intro,cauchy_sequence,sequence_in,funs_is_function,funs_elim}.
            $\alpha\mathrel{R} y$ by \cref{fst_eq,snd_eq,relation,times_tuple_intro}.
            Follows by definition.
        \end{subproof}
        Take $f$ such that
            $f$ is a function on $J$ and
            for all $\alpha\in J$ we have $\alpha\mathrel{R}\apply{f}{\alpha}$ by \cref{choice_function}.
        $f\in F$ by \cref{choice_function,relation,subseteq,times_tuple_elim,funs_intro}.
        Show that for all $U$ if $U$ is a neighborhood of $f$ in $F$ with respect to $\metricTopology{\rho}{F}$ then
            there exists $N\in\naturalsPlus$ such that for all $n\in\naturalsPlus$ if $N\leq n$ then $s(n)\in U$.
        \begin{subproof}
            Fix $U$ such that $U$ is a neighborhood of $f$ in $F$ with respect to $\metricTopology{\rho}{F}$.
            $U$ is open in $F$ with respect to $\metricTopology{\rho}{F}$ and $f\in U$ by \cref{neighborhood}.
            Take $\epsilon\in\reals$ such that $0 < \epsilon$ and $\metricBall{\rho}{F}{f}{\epsilon}\subseteq U$
                by \cref{open_in,topology_on,subseteq,powerset_elim,metric_open_iff}.
            Let $\delta = \epsilon / (1 + 1)$.
            $\delta\in\reals$ and $0 < \delta$ and $\delta + \delta = \epsilon$ by \cref{real_half_of_positive}.
            $0\in\reals$ by \cref{real_add_ident}.
            $0 + \delta < \delta + \delta$ by \cref{real_lt_add}.
            $0 + \delta = \delta$ and $\delta + \delta = \epsilon$ by \cref{real_add_ident}.
            $\delta < \epsilon$ by \cref{real_lt_subst_left,real_lt_subst_right}.
            Take $N\in\naturalsPlus$ such that
                for all $n,m\in\naturalsPlus$ if $N \leq n$ and $N \leq m$ then
                    $\rho((s(n),s(m))) < \delta$ by \cref{cauchy_sequence}.
            Show that for all $n\in\naturalsPlus$ if $N\leq n$ then $s(n)\in U$.
            \begin{subproof}
                Fix $n\in\naturalsPlus$.
                Assume $N\leq n$.
                $s(n)\in F$ by \cref{cauchy_sequence,sequence_in,funs_type_apply}.
                Show that for all $\alpha\in J$ we have $e((\apply{\apply{s}{n}}{\alpha},f(\alpha))) \leq \delta$.
                \begin{subproof}
                    Fix $\alpha\in J$.
                    $\alpha\mathrel{R} f(\alpha)$ by \cref{choice_function}.
                    Take $t$ such that
                        $t$ is a sequence in $Y$ and
                        for all $m\in\naturalsPlus$ we have $t(m) = \apply{\apply{s}{m}}{\alpha}$ and
                        $t$ converges to $f(\alpha)$ in $Y$ with respect to $\metricTopology{e}{Y}$ by
                        \cref{fst_eq,snd_eq,relation}.
                    Show that for all $m\in\naturalsPlus$ if $N \leq m$ then
                        $e((\apply{\apply{s}{n}}{\alpha},t(m))) < \delta$.
                    \begin{subproof}
                        Follows by \cref{cauchy_sequence,sequence_in,funs_type_apply,uniform_function_metric_coordinate_leq,metric_on,times_tuple_intro,real_leq_lt_trans,function_apply_intro,real_lt_subst_right}.
                    \end{subproof}
                    $e((\apply{\apply{s}{n}}{\alpha},f(\alpha))) \leq \delta$
                        by \cref{metric_limit_eventual_bound,funs_type_apply}.
                \end{subproof}
                $\rho((s(n),f)) < \epsilon$
                    by \cref{uniform_metric_leq_from_coordinate_bounds,metric_on,funs_type_apply,times_tuple_intro,real_leq_lt_trans}.
                $\rho((f,s(n))) < \epsilon$ by \cref{metric_on,metric_symmetric,real_lt_subst_left,times_tuple_intro}.
                $s(n)\in\metricBall{\rho}{F}{f}{\epsilon}$ by \cref{metric_ball}.
                $s(n)\in U$ by \cref{subseteq}.
            \end{subproof}
            Follows by definition.
        \end{subproof}
        $s$ converges to $f$ in $F$ with respect to $\metricTopology{\rho}{F}$ by \cref{cauchy_sequence,converges_to}.
        Follows by definition.
    \end{subproof}
    $F$ is complete with respect to $\rho$ by \cref{complete_metric_space}.
\end{proof}

% from §43 helper definition 1: set of continuous functions
\begin{definition}\label{continuous_function_set}
    $C$ is the set of continuous functions from $X$ to $Y$ with respect to $T$ and $U$ iff
        $C = \{f\in\funs{X}{Y} \mid \text{$f$ is continuous from $X$ to $Y$ with respect to $T$ and $U$}\}$.
\end{definition}

% from §43 helper definition 2: set of bounded functions
\begin{definition}\label{bounded_function_set}
    $B$ is the set of bounded functions from $X$ to $Y$ with respect to $d$ iff
        $B = \{f\in\funs{X}{Y} \mid \text{$\img{f}{X}$ is bounded in $Y$ with respect to $d$}\}$.
\end{definition}

% from §43 helper definition 3: distance set of two functions
\begin{definition}\label{function_distance_set}
    $A$ is the distance set of $f$ and $g$ on $X$ with respect to $d$ iff
        $A = \{d((\apply{f}{x},\apply{g}{x})) \mid x\in X\}$.
\end{definition}

% from §43 helper lemma 11: pointwise-close functions to bounded functions are bounded
\begin{lemma}\label{pointwise_close_to_bounded_function_bounded}
    Assume $d$ is a metric on $Y$.
    Suppose $f\in\funs{X}{Y}$.
    Suppose $g\in\funs{X}{Y}$.
    Assume $\img{g}{X}$ is bounded in $Y$ with respect to $d$.
    Suppose $\delta\in\reals$.
    Suppose for all $x\in X$ we have $d((g(x),f(x))) \leq \delta$.
    Then $\img{f}{X}$ is bounded in $Y$ with respect to $d$.
\end{lemma}
\begin{proof}
    Take $M\in\reals$ such that for all $z_1,z_2\in\img{g}{X}$ we have $d((z_1,z_2)) \leq M$
        by \cref{metric_bounded}.
    Show that for all $y_1,y_2\in\img{f}{X}$ we have $d((y_1,y_2)) \leq \delta + (M + \delta)$.
    \begin{subproof}
        Fix $y_1$.
        Fix $y_2$ such that $y_1\in\img{f}{X}$ and $y_2\in\img{f}{X}$.
        Take $x_1\in X$ such that $(x_1,y_1)\in f$ by \cref{img_iff}.
        Take $x_2\in X$ such that $(x_2,y_2)\in f$ by \cref{img_iff}.
        $f(x_1) = y_1$ and $f(x_2) = y_2$ by \cref{function_apply_intro,funs_is_function}.
        $d((g(x_1),y_1)) \leq \delta$ and $d((g(x_2),y_2)) \leq \delta$
            by \cref{subseteq,real_leq_subst_right_local}.
        $g(x_1)\in Y$ and $g(x_2)\in Y$ and $y_1\in Y$ and $y_2\in Y$ by \cref{funs_type_apply}.
        $(g(x_1),y_1)\in Y\times Y$ and $(g(x_2),y_2)\in Y\times Y$ and
            $(y_1,g(x_1))\in Y\times Y$ and $(g(x_1),y_2)\in Y\times Y$ and
            $(g(x_1),g(x_2))\in Y\times Y$ and $(y_1,y_2)\in Y\times Y$
            by \cref{times_tuple_intro}.
        $d((y_1,g(x_1))) \leq \delta$ by \cref{metric_on,metric_symmetric,real_leq_subst_left_local}.
        $g(x_1)\in\img{g}{X}$ and $g(x_2)\in\img{g}{X}$ by \cref{funs_is_function,funs_elim,inter_intro,img_of_function_intro}.
        $d((g(x_1),g(x_2))) \leq M$ by \cref{metric_bounded}.
        $d((y_1,y_2))\in\reals$ and $d((y_1,g(x_1)))\in\reals$ and
            $d((g(x_1),y_2))\in\reals$ and $d((g(x_1),g(x_2)))\in\reals$ and
            $d((g(x_2),y_2))\in\reals$ by \cref{metric_on,funs_type_apply}.
        $d((y_1,y_2)) \leq d((y_1,g(x_1))) + d((g(x_1),y_2))$
            by \cref{metric_on,metric_triangle_inequality}.
        $d((g(x_1),y_2)) \leq d((g(x_1),g(x_2))) + d((g(x_2),y_2))$
            by \cref{metric_on,metric_triangle_inequality}.
        $M + \delta\in\reals$ and $\delta + (M + \delta)\in\reals$ by \cref{real_add_closed}.
        $d((y_1,g(x_1))) + d((g(x_1),y_2))\in\reals$ and
            $d((g(x_1),g(x_2))) + d((g(x_2),y_2))\in\reals$ and
            $d((y_1,g(x_1))) + (d((g(x_1),g(x_2))) + d((g(x_2),y_2)))\in\reals$ and
            $d((y_1,g(x_1))) + (M + \delta)\in\reals$ and
            $M + d((g(x_2),y_2))\in\reals$ by \cref{real_add_closed}.
        $d((g(x_1),y_2)) + d((y_1,g(x_1))) \leq
            (d((g(x_1),g(x_2))) + d((g(x_2),y_2))) + d((y_1,g(x_1)))$
            by \cref{real_leq_add}.
        $d((y_1,g(x_1))) + d((g(x_1),y_2)) \leq
            d((y_1,g(x_1))) + (d((g(x_1),g(x_2))) + d((g(x_2),y_2)))$
            by \cref{real_add_comm}.
        $d((g(x_1),g(x_2))) + d((g(x_2),y_2)) \leq M + d((g(x_2),y_2))$
            by \cref{real_leq_add}.
        $d((g(x_2),y_2)) + M \leq \delta + M$ by \cref{real_leq_add}.
        $M + d((g(x_2),y_2)) \leq M + \delta$ by \cref{real_add_comm}.
        $d((g(x_1),g(x_2))) + d((g(x_2),y_2)) \leq M + \delta$ by \cref{real_leq_trans}.
        $(d((g(x_1),g(x_2))) + d((g(x_2),y_2))) + d((y_1,g(x_1))) \leq
            (M + \delta) + d((y_1,g(x_1)))$ by \cref{real_leq_add}.
        $d((y_1,g(x_1))) + (d((g(x_1),g(x_2))) + d((g(x_2),y_2))) \leq
            d((y_1,g(x_1))) + (M + \delta)$ by \cref{real_add_comm}.
        $d((y_1,g(x_1))) + (M + \delta) \leq \delta + (M + \delta)$ by \cref{real_leq_add}.
        $d((y_1,g(x_1))) + (d((g(x_1),g(x_2))) + d((g(x_2),y_2))) \leq
            \delta + (M + \delta)$ by \cref{real_leq_trans}.
        $d((y_1,y_2)) \leq \delta + (M + \delta)$ by \cref{real_leq_trans}.
    \end{subproof}
    $\img{f}{X}\subseteq\ran{f}$ and $\ran{f}\subseteq Y$ by \cref{img_subseteq_ran,funs_subseteq_rels,subseteq,rels_ran_subseteq}.
    $\img{f}{X}$ is bounded in $Y$ with respect to $d$ by \cref{metric_bounded,subseteq_transitive,real_add_closed}.
\end{proof}

% from §43 Item 8: continuous functions form a closed subset of the function space
\begin{theorem}\label{continuous_function_space_closed_uniform}
    Assume $T$ is a topology on $X$.
    Assume $d$ is a metric on $Y$.
    Assume $\rho$ is the uniform metric on functions from $X$ to $Y$ with respect to $d$.
    Assume $C$ is the set of continuous functions from $X$ to $Y$ with respect to $T$ and $\metricTopology{d}{Y}$.
    Then $C$ is closed in $\funs{X}{Y}$ with respect to $\metricTopology{\rho}{\funs{X}{Y}}$.
\end{theorem}
\begin{proof}
    Let $F = \funs{X}{Y}$.
    Let $e = \boundedMetric{d}$.
    $e$ is a metric on $Y$ by \cref{bounded_metric_is_metric}.
    $\rho$ is a metric on $F$ by \cref{uniform_metric_on_function_space}.
    $\metricTopology{\rho}{F}$ is a topology on $F$
        by \cref{metric_basis_is_basis,basis_topology_is_topology,metric_topology,basis_for_topology}.
    $F$ is metrizable with respect to $\metricTopology{\rho}{F}$ by \cref{metrizable}.
    $C\subseteq F$ by \cref{continuous_function_set,subseteq}.
    Show that for all $f$ if $f\in\closure{C}{F}{\metricTopology{\rho}{F}}$ then $f\in C$.
    \begin{subproof}
        Follows by \cref{closure,subseteq,closure_implies_sequence_metrizable,sequence_in,funs_weaken_codom,uniform_metric_convergence_uniform,bounded_metric_topology,funs_type_apply,continuous_function_set,continuous_eq_codomain,uniform_limit_theorem}.
    \end{subproof}
    $\closure{C}{F}{\metricTopology{\rho}{F}}\subseteq C$ by \cref{subseteq}.
    $C$ is closed in $F$ with respect to $\metricTopology{\rho}{F}$ by
        \cref{subseteq,subseteq_closure,subseteq_antisymmetric,closure_closed_in}.
\end{proof}

% from §43 Item 8: bounded functions form a closed subset of the function space
\begin{theorem}\label{bounded_function_space_closed_uniform}
    Assume $d$ is a metric on $Y$.
    Assume $\rho$ is the uniform metric on functions from $X$ to $Y$ with respect to $d$.
    Assume $B$ is the set of bounded functions from $X$ to $Y$ with respect to $d$.
    Then $B$ is closed in $\funs{X}{Y}$ with respect to $\metricTopology{\rho}{\funs{X}{Y}}$.
\end{theorem}
\begin{proof}
    Let $F = \funs{X}{Y}$.
    Let $e = \boundedMetric{d}$.
    $e$ is a metric on $Y$ by \cref{bounded_metric_is_metric}.
    $\rho$ is a metric on $F$ by \cref{uniform_metric_on_function_space}.
    $\metricTopology{\rho}{F}$ is a topology on $F$
        by \cref{metric_basis_is_basis,basis_topology_is_topology,metric_topology,basis_for_topology}.
    $F$ is metrizable with respect to $\metricTopology{\rho}{F}$ by \cref{metrizable}.
    $B\subseteq F$ by \cref{bounded_function_set,subseteq}.
    Show that for all $f$ if $f\in\closure{B}{F}{\metricTopology{\rho}{F}}$ then $f\in B$.
    \begin{subproof}
        Fix $f$.
        Assume $f\in\closure{B}{F}{\metricTopology{\rho}{F}}$.
        $f\in F$ by \cref{closure,subseteq}.
        Take $s$ such that
            $s$ is a sequence in $B$ and
            $s$ converges to $f$ in $F$ with respect to $\metricTopology{\rho}{F}$
            by \cref{closure_implies_sequence_metrizable}.
        $s$ converges uniformly to $f$ from $X$ to $Y$ with respect to $e$
            by \cref{sequence_in,funs_weaken_codom,uniform_metric_convergence_uniform}.
        Take $N\in\naturalsPlus$ such that for all $n\in\naturalsPlus$ if $N\leq n$ then
            for all $x\in X$ we have $e((\apply{\apply{s}{n}}{x}, f(x))) < 1$
            by \cref{uniform_converges_to,real_naturals_subset_reals,real_one_in_naturals,subseteq,real_zero_lt_one}.
        $\apply{s}{N}\in F$ by \cref{sequence_in,funs_type_apply,subseteq}.
        $\img{\apply{s}{N}}{X}$ is bounded in $Y$ with respect to $d$
            by \cref{sequence_in,funs_type_apply,bounded_function_set}.
        $1\in\reals$ by \cref{real_naturals_subset_reals,real_one_in_naturals,subseteq}.
        For all $x\in X$ we have $d((\apply{\apply{s}{N}}{x}, f(x))) \leq 1$
            by \cref{subseteq,funs_type_apply,times_tuple_intro,bounded_metric_apply_value_lt_one,real_lt_subst_left,real_lt_imp_leq,real_naturals_subset_reals,real_one_in_naturals}.
        $\img{f}{X}$ is bounded in $Y$ with respect to $d$
            by \cref{pointwise_close_to_bounded_function_bounded}.
        $f\in B$ by \cref{bounded_function_set}.
    \end{subproof}
    $\closure{B}{F}{\metricTopology{\rho}{F}}\subseteq B$ by \cref{subseteq}.
    $B$ is closed in $F$ with respect to $\metricTopology{\rho}{F}$ by
        \cref{subseteq,subseteq_closure,subseteq_antisymmetric,closure_closed_in}.
\end{proof}

% from §43 helper lemma 12: closed subsets of complete metric spaces are complete under the restricted metric
\begin{lemma}\label{closed_subset_complete_restricted_metric}
    Assume $d$ is a metric on $X$.
    Assume $X$ is complete with respect to $d$.
    Assume $A$ is closed in $X$ with respect to $\metricTopology{d}{X}$.
    Let $e = \restrl{d}{A\times A}$.
    Then $A$ is complete with respect to $e$.
\end{lemma}
\begin{proof}
    $A\subseteq X$ by \cref{closed_in,subseteq}.
    $d\in\funs{X\times X}{\reals}$ and $d$ is a function and $d$ is a relation and $\dom{d} = X\times X$
        by \cref{metric_on,funs_elim,funs_is_relation}.
    $A\times A\subseteq X\times X$
        by \cref{times_subseteq_left,times_subseteq_right,subseteq_transitive}.
    $e$ is a function by \cref{restrl_of_function_is_function}.
    $\dom{e} = A\times A$ by \cref{dom_of_restrl_eq_inter,pair_eq_iff,subseteq,inter_intro,restrl_dom,subseteq_antisymmetric}.
    For all $p\in\dom{e}$ we have $e(p)\in\reals$
        by \cref{inter_elim_left,dom_of_restrl_eq_inter,funs_type_apply,restrl_of_function_apply,subseteq,pair_eq_iff}.
    $e\in\funs{A\times A}{\reals}$ by \cref{funs_intro}.

    Show that $e$ is nonnegative on $A$.
    \begin{subproof}
        Follows by \cref{subseteq,times_tuple_intro,funs_elim,restrl_of_function_apply,metric_on,metric_nonnegative,pair_eq_iff}.
    \end{subproof}

    Show that $e$ is zero-characterized on $A$.
    \begin{subproof}
        Follows by \cref{subseteq,times_tuple_intro,funs_elim,restrl_of_function_apply,metric_on,metric_zero_characterization,pair_eq_iff}.
    \end{subproof}

    Show that $e$ is symmetric on $A$.
    \begin{subproof}
        Follows by \cref{subseteq,times_tuple_intro,funs_elim,restrl_of_function_apply,metric_on,metric_symmetric,pair_eq_iff}.
    \end{subproof}

    Show that $e$ satisfies the triangle inequality on $A$.
    \begin{subproof}
        Follows by \cref{subseteq,times_tuple_intro,funs_elim,restrl_of_function_apply,metric_on,metric_triangle_inequality,pair_eq_iff}.
    \end{subproof}
    $e$ is a metric on $A$ by \cref{metric_on}.

    Show that for all $s$ if
        $s$ is a Cauchy sequence in $A$ with respect to $e$
    then there exists $x\in A$ such that
        $s$ converges to $x$ in $A$ with respect to $\metricTopology{e}{A}$.
    \begin{subproof}
        Fix $s$ such that $s$ is a Cauchy sequence in $A$ with respect to $e$.
        $s$ is a sequence in $A$ by \cref{cauchy_sequence}.
        For all $\epsilon\in\reals$ such that $0 < \epsilon$ there exists $N\in\naturalsPlus$ such that
            for all $n,m\in\naturalsPlus$ if $N \leq n$ and $N \leq m$ then
                $e((s(n),s(m))) < \epsilon$ by \cref{cauchy_sequence}.
        $s$ is a sequence in $X$ by \cref{cauchy_sequence,sequence_in,funs_weaken_codom,subseteq}.
        For all $\epsilon\in\reals$ such that $0 < \epsilon$ there exists $N\in\naturalsPlus$ such that
            for all $n,m\in\naturalsPlus$ if $N \leq n$ and $N \leq m$ then
                $d((s(n),s(m))) < \epsilon$
            by \cref{cauchy_sequence,sequence_in,funs_type_apply,times_tuple_intro,restrl_of_function_apply,subseteq,real_lt_subst_left}.
        $s$ is a Cauchy sequence in $X$ with respect to $d$ by \cref{cauchy_sequence}.
        Take $x\in X$ such that
            $s$ converges to $x$ in $X$ with respect to $\metricTopology{d}{X}$
            by \cref{complete_metric_space}.
        $x\in A$ by \cref{metric_basis_is_basis,basis_topology_is_topology,metric_topology,basis_for_topology,sequence_implies_closure,closure_subseteq_closed,subseteq_refl,subseteq}.
        Show that for all $U$ if $U$ is a neighborhood of $x$ in $A$ with respect to $\metricTopology{e}{A}$ then
            there exists $N\in\naturalsPlus$ such that for all $n\in\naturalsPlus$ if $N\leq n$ then $s(n)\in U$.
        \begin{subproof}
            Fix $U$ such that $U$ is a neighborhood of $x$ in $A$ with respect to $\metricTopology{e}{A}$.
            Take $\epsilon\in\reals$ such that $0 < \epsilon$ and $\metricBall{e}{A}{x}{\epsilon}\subseteq U$
                by \cref{neighborhood,open_in,topology_on,subseteq,powerset_elim,metric_open_iff}.
            Take $N\in\naturalsPlus$ such that for all $n\in\naturalsPlus$ if $N\leq n$ then
                $s(n)\in\metricBall{d}{X}{x}{\epsilon}$ by \cref{converges_to,metric_ball_neighborhood_metric_topology}.
            Show that for all $n\in\naturalsPlus$ if $N\leq n$ then $s(n)\in U$.
            \begin{subproof}
                Follows by \cref{sequence_in,metric_ball,funs_type_apply,times_tuple_intro,restrl_of_function_apply,subseteq,real_lt_subst_left}.
            \end{subproof}
        \end{subproof}
        $s$ converges to $x$ in $A$ with respect to $\metricTopology{e}{A}$ by \cref{converges_to}.
        Follows by definition.
    \end{subproof}
    $A$ is complete with respect to $e$ by \cref{complete_metric_space}.
\end{proof}

% from §43 Item 8: continuous function spaces over complete metric spaces are complete
\begin{theorem}\label{continuous_function_space_complete_uniform}
    Assume $T$ is a topology on $X$.
    Assume $d$ is a metric on $Y$.
    Assume $Y$ is complete with respect to $d$.
    Assume $\rho$ is the uniform metric on functions from $X$ to $Y$ with respect to $d$.
    Assume $C$ is the set of continuous functions from $X$ to $Y$ with respect to $T$ and $\metricTopology{d}{Y}$.
    Then $C$ is complete with respect to $\restrl{\rho}{C\times C}$.
\end{theorem}
\begin{proof}
    Follows by \cref{uniform_metric_on_function_space,continuous_function_space_closed_uniform,uniform_function_space_complete,closed_subset_complete_restricted_metric}.
\end{proof}

% from §43 Item 8: bounded function spaces over complete metric spaces are complete
\begin{theorem}\label{bounded_function_space_complete_uniform}
    Assume $d$ is a metric on $Y$.
    Assume $Y$ is complete with respect to $d$.
    Assume $\rho$ is the uniform metric on functions from $X$ to $Y$ with respect to $d$.
    Assume $B$ is the set of bounded functions from $X$ to $Y$ with respect to $d$.
    Then $B$ is complete with respect to $\restrl{\rho}{B\times B}$.
\end{theorem}
\begin{proof}
    Follows by \cref{uniform_metric_on_function_space,bounded_function_space_closed_uniform,uniform_function_space_complete,closed_subset_complete_restricted_metric}.
\end{proof}

% from §43 Item 9: sup metric
\begin{definition}\label{sup_metric_on_bounded_function_space}
    $\rho$ is the sup metric on $B$ with respect to $X$ and $Y$ and $d$ iff
        $d$ is a metric on $Y$ and
        $B\subseteq\funs{X}{Y}$ and
        for all $f\in B$ there exists $M\in\reals$ such that
            for all $x_1,x_2\in X$ we have $d((\apply{f}{x_1},\apply{f}{x_2})) \leq M$ and
        $\rho$ is a metric on $B$ and
        for all $f,g\in B$ there exists $A$ such that
            $A$ is the distance set of $f$ and $g$ on $X$ with respect to $d$ and
            $\rho((f,g))$ is a supremum of $A$.
\end{definition}

% from §43 helper definition 4: isometric embedding
\begin{definition}\label{isometric_embedding}
    $h$ is an isometric embedding from $X$ to $Y$ with respect to $d_X$ and $d_Y$ iff
        $d_X$ is a metric on $X$ and
        $d_Y$ is a metric on $Y$ and
        $h\in\funs{X}{Y}$ and
        for all $a,b\in X$ we have $d_Y((h(a),h(b))) = d_X((a,b))$.
\end{definition}

% from §43 helper lemma 15: distance sets of bounded real-valued functions have suprema on inhabited domains
\begin{lemma}\label{bounded_reals_distance_set_has_supremum}
    Assume $X$ is inhabited.
    Assume $B$ is the set of bounded functions from $X$ to $\reals$ with respect to $\standardMetricR$.
    Suppose $f\in B$.
    Suppose $g\in B$.
    Then there exist $A, p$ such that
        $A$ is the distance set of $f$ and $g$ on $X$ with respect to $\standardMetricR$ and
        $p$ is a supremum of $A$.
\end{lemma}
\begin{proof}
    $\standardMetricR$ is a metric on $\reals$ by \cref{standard_metric_is_metric}.
    $f\in\funs{X}{\reals}$ and $\img{f}{X}$ is bounded in $\reals$ with respect to $\standardMetricR$
        by \cref{bounded_function_set}.
    $g\in\funs{X}{\reals}$ and $\img{g}{X}$ is bounded in $\reals$ with respect to $\standardMetricR$
        by \cref{bounded_function_set}.
    Take $x_0$ such that $x_0\in X$.
    Let $A = \{\apply{\standardMetricR}{(f(x),g(x))} \mid x\in X\}$.
    $A$ is the distance set of $f$ and $g$ on $X$ with respect to $\standardMetricR$
        by \cref{function_distance_set}.
    For all $r$ such that $r\in A$ we have $r\in\reals$
        by \cref{function_distance_set,funs_type_apply,times_tuple_intro,standard_metric_is_metric,metric_on,pair_eq_iff}.
    $A\subseteq\reals$ by \cref{subseteq}.
    $f$ is a function and $\dom{f} = X$ and $g$ is a function and $\dom{g} = X$
        by \cref{funs_is_function,funs_elim}.
    $x_0\in\dom{f}\inter X$ and $x_0\in\dom{g}\inter X$ by \cref{funs_elim,inter_intro}.
    $f(x_0)\in\img{f}{X}$ and $g(x_0)\in\img{g}{X}$ by \cref{img_of_function_intro}.
    Take $M_f\in\reals$ such that for all $u,v\in\img{f}{X}$ we have $\apply{\standardMetricR}{(u,v)} \leq M_f$
        by \cref{metric_bounded}.
    Take $M_g\in\reals$ such that for all $u,v\in\img{g}{X}$ we have $\apply{\standardMetricR}{(u,v)} \leq M_g$
        by \cref{metric_bounded}.
    $f(x_0)\in\reals$ and $g(x_0)\in\reals$ by \cref{funs_type_apply}.
    Let $c = \apply{\standardMetricR}{(f(x_0),g(x_0))}$.
    $A\neq\emptyset$
        by \cref{funs_type_apply,times_tuple_intro,function_distance_set,inhabited_nonempty,nonempty_inhabited}.
    Let $M = M_f + (c + M_g)$.
    $c + M_g\in\reals$ and $M\in\reals$
        by \cref{standard_metric_is_metric,metric_on,funs_type_apply,real_add_closed}.
    Show that for all $r\in A$ we have $r \leq M$.
    \begin{subproof}
        Fix $r$ such that $r\in A$.
        Take $x\in X$ such that $r = \apply{\standardMetricR}{(f(x),g(x))}$ by \cref{function_distance_set}.
        $x\in\dom{f}\inter X$ and $x\in\dom{g}\inter X$ by \cref{inter_intro}.
        $f(x)\in\img{f}{X}$ and $f(x_0)\in\img{f}{X}$ and
            $g(x_0)\in\img{g}{X}$ and $g(x)\in\img{g}{X}$ by \cref{img_of_function_intro}.
        Let $u = \apply{\standardMetricR}{(f(x),f(x_0))}$.
        Let $v = \apply{\standardMetricR}{(f(x_0),g(x))}$.
        Let $w = \apply{\standardMetricR}{(g(x_0),g(x))}$.
        $u \leq M_f$ and $w \leq M_g$ by \cref{metric_bounded}.
        $f(x)\in\reals$ and $f(x_0)\in\reals$ and $g(x_0)\in\reals$ and $g(x)\in\reals$
            by \cref{funs_type_apply}.
        $(f(x),f(x_0))\in\reals\times\reals$ and
            $(f(x_0),g(x))\in\reals\times\reals$ and
            $(g(x_0),g(x))\in\reals\times\reals$ and
            $(f(x),g(x))\in\reals\times\reals$ by \cref{times_tuple_intro}.
        $u\in\reals$ and $v\in\reals$ and $w\in\reals$ and
            $\apply{\standardMetricR}{(f(x),g(x))}\in\reals$ by \cref{standard_metric_is_metric,metric_on,funs_type_apply}.
        $\apply{\standardMetricR}{(f(x),g(x))} \leq u + v$
            by \cref{standard_metric_is_metric,metric_on,metric_triangle_inequality}.
        $v \leq c + w$
            by \cref{standard_metric_is_metric,metric_on,metric_triangle_inequality}.
        $c + w\in\reals$ by \cref{real_add_closed}.
        $v + u \leq (c + w) + u$ by \cref{real_leq_add}.
        $u + v \leq u + (c + w)$ by \cref{real_add_comm}.
        $u + (c + w) \leq M_f + (c + w)$
            by \cref{real_leq_add}.
        $w + c \leq M_g + c$ by \cref{real_leq_add}.
        $c + w \leq c + M_g$ by \cref{real_add_comm}.
        $(c + w) + M_f \leq (c + M_g) + M_f$ by \cref{real_leq_add}.
        $M_f + (c + w) \leq M_f + (c + M_g)$ by \cref{real_add_comm}.
        $u + (c + w)\in\reals$ and $M_f + (c + w)\in\reals$ and $M_f + (c + M_g)\in\reals$
            by \cref{real_add_closed}.
        $u + (c + w) \leq M_f + (c + M_g)$
            by \cref{real_leq_trans}.
        $u + (c + w) \leq M$.
        $u + v\in\reals$ by \cref{real_add_closed}.
        $u + v \leq M$
            by \cref{real_leq_trans}.
        $\apply{\standardMetricR}{(f(x),g(x))} \leq M$ by \cref{real_leq_trans}.
        $r \leq M$ by \cref{pair_eq_iff}.
    \end{subproof}
    $A$ is bounded above by \cref{real_upper_bound,real_bounded_above}.
    Take $p$ such that $p$ is a supremum of $A$ by \cref{real_completeness_axiom}.
    Follows by definition.
\end{proof}

% from §43 helper lemma 16: each coordinate distance is bounded by the sup metric
\begin{lemma}\label{sup_metric_coordinate_leq}
    Assume $X$ is inhabited.
    Assume $B$ is the set of bounded functions from $X$ to $\reals$ with respect to $\standardMetricR$.
    Assume $\rho$ is the sup metric on $B$ with respect to $X$ and $\reals$ and $\standardMetricR$.
    Suppose $f\in B$.
    Suppose $g\in B$.
    Suppose $x\in X$.
    Then $\apply{\standardMetricR}{(f(x),g(x))} \leq \rho((f,g))$.
\end{lemma}
\begin{proof}
    Take $A$ such that
        $A$ is the distance set of $f$ and $g$ on $X$ with respect to $\standardMetricR$ and
        $\rho((f,g))$ is a supremum of $A$ by \cref{sup_metric_on_bounded_function_space}.
    $\apply{\standardMetricR}{(f(x),g(x))} \leq \rho((f,g))$
        by \cref{sup_metric_on_bounded_function_space,bounded_function_set,funs_type_apply,times_tuple_intro,function_distance_set,standard_metric_is_metric,metric_on,real_supremum,real_upper_bound}.
\end{proof}

% from §43 helper lemma 17: pointwise bounds imply a sup-metric bound
\begin{lemma}\label{sup_metric_leq_from_coordinate_bounds}
    Assume $X$ is inhabited.
    Assume $B$ is the set of bounded functions from $X$ to $\reals$ with respect to $\standardMetricR$.
    Assume $\rho$ is the sup metric on $B$ with respect to $X$ and $\reals$ and $\standardMetricR$.
    Suppose $f\in B$.
    Suppose $g\in B$.
    Suppose $\delta\in\reals$.
    Suppose for all $x\in X$ we have $\apply{\standardMetricR}{(f(x),g(x))} \leq \delta$.
    Then $\rho((f,g)) \leq \delta$.
\end{lemma}
\begin{proof}
    Take $A$ such that
        $A$ is the distance set of $f$ and $g$ on $X$ with respect to $\standardMetricR$ and
        $\rho((f,g))$ is a supremum of $A$ by \cref{sup_metric_on_bounded_function_space}.
    For all $r$ such that $r\in A$ we have $r\in\reals$
        by \cref{function_distance_set,bounded_function_set,funs_type_apply,times_tuple_intro,standard_metric_is_metric,metric_on,pair_eq_iff}.
    $A\subseteq\reals$ by \cref{subseteq}.
    $\delta$ is an upper bound of $A$ by \cref{function_distance_set,pair_eq_iff,real_upper_bound}.
    $\rho((f,g)) \leq \delta$ by \cref{real_supremum}.
\end{proof}

% from §43 helper lemma 18: Cauchy sequences in the sup metric are coordinatewise Cauchy
\begin{lemma}\label{sup_metric_cauchy_coordinate}
    Assume $X$ is inhabited.
    Assume $B$ is the set of bounded functions from $X$ to $\reals$ with respect to $\standardMetricR$.
    Assume $\rho$ is the sup metric on $B$ with respect to $X$ and $\reals$ and $\standardMetricR$.
    Suppose $s$ is a Cauchy sequence in $B$ with respect to $\rho$.
    Suppose $x\in X$.
    Let $t = \{(n,\apply{\apply{s}{n}}{x}) \mid n\in\naturalsPlus\}$.
    Then $t$ is a Cauchy sequence in $\reals$ with respect to $\standardMetricR$.
\end{lemma}
\begin{proof}
    $t$ is a function by \cref{pair_eq_iff,relation,rightunique}.
    $\dom{t} = \naturalsPlus$ by \cref{dom_intro,dom_iff,pair_eq_iff,subseteq,subseteq_antisymmetric}.
    For all $n\in\dom{t}$ we have $t(n)\in\reals$
        by \cref{subseteq,cauchy_sequence,sequence_in,funs_type_apply,bounded_function_set,function_apply_intro}.
    $t$ is a sequence in $\reals$ by \cref{funs_intro,sequence_in}.
    Show that for all $\epsilon$ if $\epsilon\in\reals$ and $0 < \epsilon$ then
        there exists $N\in\naturalsPlus$ such that
            for all $n,m\in\naturalsPlus$ if $N \leq n$ and $N \leq m$ then
                $\apply{\standardMetricR}{(t(n),t(m))} < \epsilon$.
    \begin{subproof}
        Follows by \cref{cauchy_sequence,sequence_in,funs_type_apply,sup_metric_coordinate_leq,sup_metric_on_bounded_function_space,metric_on,times_tuple_intro,bounded_function_set,subseteq,standard_metric_is_metric,real_leq_lt_trans,function_apply_intro,real_lt_subst_left,real_lt_subst_right}.
    \end{subproof}
    $t$ is a Cauchy sequence in $\reals$ with respect to $\standardMetricR$ by \cref{standard_metric_is_metric,cauchy_sequence}.
\end{proof}

% from §43 helper lemma 14: bounded real-valued functions are complete in the sup metric on inhabited domains
\begin{lemma}\label{bounded_reals_function_space_complete_sup}
    Assume $X$ is inhabited.
    Assume $B$ is the set of bounded functions from $X$ to $\reals$ with respect to $\standardMetricR$.
    Then there exists $\rho$ such that
        $\rho$ is the sup metric on $B$ with respect to $X$ and $\reals$ and $\standardMetricR$ and
        $B$ is complete with respect to $\rho$.
\end{lemma}
\begin{proof}
    $\standardMetricR$ is a metric on $\reals$ by \cref{standard_metric_is_metric}.
    $B\subseteq\funs{X}{\reals}$ by \cref{bounded_function_set,subseteq}.
    Let $R = \{w\in(B\times B)\times\reals \mid
        \text{there exists $A$ such that
            $A$ is the distance set of $\fst{\fst{w}}$ and $\snd{\fst{w}}$ on $X$ with respect to $\standardMetricR$ and
            $\snd{w}$ is a supremum of $A$}\}$.
    $R$ is a relation by \cref{relation,subseteq,times_elem_is_tuple}.
    For all $p\in B\times B$ there exists $r$ such that $p\mathrel{R} r$
        by \cref{times,bounded_reals_distance_set_has_supremum,real_supremum,real_upper_bound,times_tuple_intro,fst_eq,snd_eq,relation}.
    Take $\rho$ such that
        $\rho$ is a function on $B\times B$ and
        for all $p\in B\times B$ we have $p\mathrel{R}\rho(p)$ by \cref{choice_function}.
        For all $p\in B\times B$ we have $\rho(p)\in\reals$
            by \cref{choice_function,relation,subseteq,times_tuple_elim}.
        Hence $\dom{\rho} = B\times B$ by \cref{choice_function}.
        Therefore $\rho\in\funs{B\times B}{\reals}$ by \cref{funs_intro}.

        For all $f,g$ such that $f\in B$ and $g\in B$ we have $\rho((f,g)) \geq 0$
            by \cref{times_tuple_intro,choice_function,relation,fst_eq,snd_eq,bounded_function_set,funs_type_apply,function_distance_set,standard_metric_is_metric,metric_on,real_supremum,real_upper_bound,metric_nonnegative,real_add_ident,real_leq_trans}.
        Therefore $\rho$ is nonnegative on $B$ by \cref{metric_nonnegative}.

        Show that for all $f,g$ such that $f\in B$ and $g\in B$ we have $\rho((f,g)) = 0$ iff $f = g$.
        \begin{subproof}
                Fix $f$.
                Fix $g$ such that $f\in B$ and $g\in B$.
                Show that if $\rho((f,g)) = 0$ then $f = g$.
                \begin{subproof}
                    Assume $\rho((f,g)) = 0$.
                    We have $f\in\funs{X}{\reals}$ and $g\in\funs{X}{\reals}$ by \cref{bounded_function_set}.
                    $f$ is a function and $\dom{f} = X$ and $g$ is a function and $\dom{g} = X$
                        by \cref{funs_is_function,funs_elim}.
                    Let $p = (f,g)$.
                    Hence $(p,\rho(p))\in R$ by \cref{times_tuple_intro,choice_function,relation}.
                    Take $A$ such that
                        $A$ is the distance set of $\fst{p}$ and $\snd{p}$ on $X$ with respect to $\standardMetricR$ and
                        $\rho(p)$ is a supremum of $A$ by \cref{fst_eq,snd_eq,relation}.
                    For all $x\in X$ we have $f(x) = g(x)$
                        by \cref{fst_eq,snd_eq,function_distance_set,real_supremum,real_upper_bound,funs_type_apply,times_tuple_intro,standard_metric_is_metric,metric_on,metric_nonnegative,real_add_ident,real_leq_antisym,metric_zero_characterization}.
                    Therefore $f = g$ by \cref{functions_eq_iff}.
                \end{subproof}

                Show that if $f = g$ then $\rho((f,g)) = 0$.
                \begin{subproof}
                    Assume $f = g$.
                    We have $f\in\funs{X}{\reals}$ and $g\in\funs{X}{\reals}$ by \cref{bounded_function_set}.
                    $f$ is a function and $\dom{f} = X$ and $g$ is a function and $\dom{g} = X$
                        by \cref{funs_is_function,funs_elim}.
                    Let $p = (f,g)$.
                    $p\in B\times B$ by \cref{times_tuple_intro}.
                    Hence $(p,\rho(p))\in R$ by \cref{choice_function,relation}.
                    Take $A$ such that
                        $A$ is the distance set of $\fst{p}$ and $\snd{p}$ on $X$ with respect to $\standardMetricR$ and
                        $\rho(p)$ is a supremum of $A$ by \cref{fst_eq,snd_eq,relation}.
                    Show that $0$ is an upper bound of $A$.
                    \begin{subproof}
                        $0\in\reals$ by \cref{real_add_ident}.
                        $A\subseteq\reals$
                            by \cref{function_distance_set,fst_eq,snd_eq,bounded_function_set,funs_type_apply,standard_metric_is_metric,metric_on,subseteq,times_tuple_intro}.
                        Show that for all $r\in A$ we have $r < 0$ or $r = 0$.
                        \begin{subproof}
                            Fix $r\in A$.
                            Take $x\in X$ such that $\apply{\standardMetricR}{(\apply{f}{x},\apply{g}{x})} = r$
                                by \cref{function_distance_set,fst_eq,snd_eq}.
                            $\apply{f}{x}\in\reals$ and $\apply{g}{x}\in\reals$ by \cref{funs_type_apply,subseteq}.
                            $\apply{f}{x} = \apply{g}{x}$ by \cref{real_lt_subst_left}.
                            $\apply{\standardMetricR}{(\apply{f}{x},\apply{g}{x})} = 0$
                                by \cref{standard_metric_is_metric,metric_on,metric_zero_characterization,times_tuple_intro,real_lt_subst_right}.
                            Hence $r = 0$ by \cref{real_lt_subst_left}.
                        \end{subproof}
                        Thus $0$ is an upper bound of $A$ by \cref{real_upper_bound}.
                    \end{subproof}
                    $\rho(p) < 0$ or $\rho(p) = 0$ by \cref{real_supremum}.
                    $0 < \rho(p)$ or $0 = \rho(p)$
                        by \cref{metric_nonnegative,real_lt_subst_left,times_tuple_intro}.
                    $\rho(p)\in\reals$ by \cref{funs_type_apply}.
                    Hence $\rho(p) = 0$ by \cref{real_add_ident,real_leq_antisym}.
                    Therefore $\rho((f,g)) = 0$ by \cref{real_lt_subst_left}.
                \end{subproof}
        \end{subproof}
        Therefore $\rho$ is zero-characterized on $B$ by \cref{metric_zero_characterization}.

        Show that for all $f,g$ such that $f\in B$ and $g\in B$ we have $\rho((f,g)) = \rho((g,f))$.
        \begin{subproof}
                Fix $f$.
                Fix $g$ such that $f\in B$ and $g\in B$.
                Let $p = (f,g)$.
                Let $q = (g,f)$.
                We have $(p,\rho(p))\in R$ and $(q,\rho(q))\in R$
                    by \cref{times_tuple_intro,choice_function,relation}.
                Take $A$ such that
                    $A$ is the distance set of $\fst{p}$ and $\snd{p}$ on $X$ with respect to $\standardMetricR$ and
                    $\rho(p)$ is a supremum of $A$ by \cref{fst_eq,snd_eq,relation}.
                Take $C$ such that
                    $C$ is the distance set of $\fst{q}$ and $\snd{q}$ on $X$ with respect to $\standardMetricR$ and
                    $\rho(q)$ is a supremum of $C$ by \cref{fst_eq,snd_eq,relation}.
                For all $r$ if $r\in A$ then $r\in C$
                    by \cref{fst_eq,snd_eq,function_distance_set,bounded_function_set,funs_type_apply,standard_metric_is_metric,metric_on,metric_symmetric,times_tuple_intro,pair_eq_iff}.
                Thus $A\subseteq C$ by \cref{subseteq}.
                For all $r$ if $r\in C$ then $r\in A$
                    by \cref{fst_eq,snd_eq,function_distance_set,bounded_function_set,funs_type_apply,standard_metric_is_metric,metric_on,metric_symmetric,times_tuple_intro,pair_eq_iff}.
                Hence $C\subseteq A$ by \cref{subseteq}.
                Therefore $\rho((f,g)) = \rho((g,f))$ by \cref{subseteq_antisymmetric,real_supremum_unique}.
        \end{subproof}
        Therefore $\rho$ is symmetric on $B$ by \cref{metric_symmetric}.

        Show that for all $f,g,h$ such that $f\in B$ and $g\in B$ and $h\in B$ we have
            $\rho((f,g)) + \rho((g,h)) \geq \rho((f,h))$.
        \begin{subproof}
                Fix $f$.
                Fix $g$.
                Fix $h$ such that $f\in B$ and $g\in B$ and $h\in B$.
                Let $p = (f,h)$.
                Let $q = (f,g)$.
                Let $s = (g,h)$.
                $p\in B\times B$ and $q\in B\times B$ and $s\in B\times B$ by \cref{times_tuple_intro}.
                Hence $(p,\rho(p))\in R$ and $(q,\rho(q))\in R$ and $(s,\rho(s))\in R$
                    by \cref{choice_function,relation}.
                Take $A$ such that
                    $A$ is the distance set of $\fst{p}$ and $\snd{p}$ on $X$ with respect to $\standardMetricR$ and
                    $\rho(p)$ is a supremum of $A$ by \cref{fst_eq,snd_eq,relation}.
                Take $C$ such that
                    $C$ is the distance set of $\fst{q}$ and $\snd{q}$ on $X$ with respect to $\standardMetricR$ and
                    $\rho(q)$ is a supremum of $C$ by \cref{fst_eq,snd_eq,relation}.
                Take $D$ such that
                    $D$ is the distance set of $\fst{s}$ and $\snd{s}$ on $X$ with respect to $\standardMetricR$ and
                    $\rho(s)$ is a supremum of $D$ by \cref{fst_eq,snd_eq,relation}.
                For all $r$ if $r\in A$ then $r\in\reals$.
                Hence $A\subseteq\reals$ by \cref{subseteq}.
                Show that for all $r\in A$ we have $r \leq \rho((f,g)) + \rho((g,h))$.
                \begin{subproof}
                    Fix $r$ such that $r\in A$.
                    Take $x\in X$ such that $r = \apply{\standardMetricR}{(f(x),h(x))}$ by \cref{fst_eq,snd_eq,function_distance_set}.
                    We have $f(x)\in\reals$ and $g(x)\in\reals$ and $h(x)\in\reals$
                        by \cref{bounded_function_set,funs_type_apply}.
                    We have $\apply{\standardMetricR}{(f(x),g(x))}\in\reals$ and
                        $\apply{\standardMetricR}{(g(x),h(x))}\in\reals$ and
                        $\apply{\standardMetricR}{(f(x),h(x))}\in\reals$ by \cref{standard_metric_is_metric,metric_on,funs_type_apply,times_tuple_intro}.
                    Thus $\apply{\standardMetricR}{(f(x),h(x))} \leq
                        \apply{\standardMetricR}{(f(x),g(x))} + \apply{\standardMetricR}{(g(x),h(x))}$
                        by \cref{standard_metric_is_metric,metric_on,metric_triangle_inequality}.
                    $\apply{\standardMetricR}{(f(x),g(x))}\in C$ and
                        $\apply{\standardMetricR}{(g(x),h(x))}\in D$ by \cref{fst_eq,snd_eq,function_distance_set}.
                    Therefore $\apply{\standardMetricR}{(f(x),g(x))} \leq \rho((f,g))$ and
                        $\apply{\standardMetricR}{(g(x),h(x))} \leq \rho((g,h))$
                        by \cref{real_supremum,real_upper_bound}.
                    $\apply{\standardMetricR}{(f(x),g(x))} + \apply{\standardMetricR}{(g(x),h(x))} \leq
                        \rho((f,g)) + \apply{\standardMetricR}{(g(x),h(x))}$ by \cref{real_leq_add}.
                    $\apply{\standardMetricR}{(g(x),h(x))} + \rho((f,g)) \leq
                        \rho((g,h)) + \rho((f,g))$ by \cref{real_leq_add}.
                    Thus $\rho((f,g)) + \apply{\standardMetricR}{(g(x),h(x))} \leq
                        \rho((f,g)) + \rho((g,h))$ by \cref{real_add_comm}.
                    We have $\apply{\standardMetricR}{(f(x),g(x))} + \apply{\standardMetricR}{(g(x),h(x))}\in\reals$ and
                        $\rho((f,g)) + \apply{\standardMetricR}{(g(x),h(x))}\in\reals$ and
                        $\rho((f,g)) + \rho((g,h))\in\reals$ by \cref{real_add_closed}.
                    Therefore $\apply{\standardMetricR}{(f(x),g(x))} + \apply{\standardMetricR}{(g(x),h(x))} \leq
                        \rho((f,g)) + \rho((g,h))$ by \cref{real_leq_trans}.
                    Thus $\apply{\standardMetricR}{(f(x),h(x))} \leq \rho((f,g)) + \rho((g,h))$
                        by \cref{real_leq_trans}.
                    Therefore $r \leq \rho((f,g)) + \rho((g,h))$ by \cref{pair_eq_iff}.
                \end{subproof}
                Thus $\rho((f,g)) + \rho((g,h))$ is an upper bound of $A$ by \cref{times_tuple_intro,funs_type_apply,real_upper_bound,real_add_closed}.
                Therefore $\rho((f,h)) \leq \rho((f,g)) + \rho((g,h))$ by \cref{real_supremum}.
                Hence $\rho((f,g)) + \rho((g,h)) \geq \rho((f,h))$.
        \end{subproof}
        Therefore $\rho$ satisfies the triangle inequality on $B$ by \cref{metric_triangle_inequality}.
        Hence $\rho$ is a metric on $B$ by \cref{metric_on}.
    For all $f,g\in B$ there exists $A$ such that
        $A$ is the distance set of $f$ and $g$ on $X$ with respect to $\standardMetricR$ and
        $\rho((f,g))$ is a supremum of $A$
        by \cref{times_tuple_intro,choice_function,relation,fst_eq,snd_eq,real_add_eq_subst}.
    For all $f\in B$ there exists $M\in\reals$ such that
        for all $x_1,x_2\in X$ we have $\apply{\standardMetricR}{(f(x_1),f(x_2))} \leq M$
        by \cref{subseteq,bounded_function_set,funs_is_function,funs_elim,inter_intro,img_of_function_intro,metric_bounded}.
    Hence $\rho$ is the sup metric on $B$ with respect to $X$ and $\reals$ and $\standardMetricR$
        by \cref{sup_metric_on_bounded_function_space}.
    Show that for all $s$ if
        $s$ is a Cauchy sequence in $B$ with respect to $\rho$
    then there exists $f\in B$ such that
        $s$ converges to $f$ in $B$ with respect to $\metricTopology{\rho}{B}$.
    \begin{subproof}
            Fix $s$ such that $s$ is a Cauchy sequence in $B$ with respect to $\rho$.
            Let $Q = \{w\in X\times\reals \mid
                \text{there exists $t$ such that
                    $t$ is a sequence in $\reals$ and
                    for all $n\in\naturalsPlus$ we have $t(n) = \apply{\apply{s}{n}}{\fst{w}}$ and
                    $t$ converges to $\snd{w}$ in $\reals$ with respect to $\metricTopology{\standardMetricR}{\reals}$}\}$.
            $Q$ is a relation by \cref{relation,subseteq,times_elem_is_tuple}.
            Show that for all $x\in X$ there exists $y$ such that $x\mathrel{Q} y$.
            \begin{subproof}
                Fix $x\in X$.
                Let $t = \{(n,\apply{\apply{s}{n}}{x}) \mid n\in\naturalsPlus\}$.
                $t$ is a Cauchy sequence in $\reals$ with respect to $\standardMetricR$
                    by \cref{sup_metric_cauchy_coordinate}.
                Take $y\in\reals$ such that
                    $t$ converges to $y$ in $\reals$ with respect to $\metricTopology{\standardMetricR}{\reals}$
                    by \cref{reals_complete_standard_metric,complete_metric_space}.
                Therefore $x\mathrel{Q} y$
                    by \cref{times_tuple_intro,fst_eq,snd_eq,relation,function_apply_intro,cauchy_sequence,sequence_in,funs_is_function,funs_elim}.
            \end{subproof}
            Take $f$ such that
                $f$ is a function on $X$ and
                for all $x\in X$ we have $x\mathrel{Q}\apply{f}{x}$ by \cref{choice_function}.
            Hence $f\in\funs{X}{\reals}$ by \cref{choice_function,relation,subseteq,times_tuple_elim,funs_intro}.
                Take $N\in\naturalsPlus$ such that
                    for all $n,m\in\naturalsPlus$ if $N \leq n$ and $N \leq m$ then
                        $\rho((s(n),s(m))) < 1$
                    by \cref{real_zero_lt_one,real_naturals_subset_reals,real_one_in_naturals,subseteq,cauchy_sequence}.
                We have $\apply{s}{N}\in\funs{X}{\reals}$ and
                    $\img{\apply{s}{N}}{X}$ is bounded in $\reals$ with respect to $\standardMetricR$
                    by \cref{sequence_in,cauchy_sequence,funs_type_apply,bounded_function_set}.
                $1\in\reals$ and $0 < 1$
                    by \cref{real_naturals_subset_reals,real_one_in_naturals,subseteq,real_zero_lt_one}.
                Show that for all $x\in X$ we have $\apply{\standardMetricR}{(\apply{\apply{s}{N}}{x},f(x))} \leq 1$.
                \begin{subproof}
                        Fix $x\in X$.
                        Take $t$ such that
                            $t$ is a sequence in $\reals$ and
                            for all $m\in\naturalsPlus$ we have $t(m) = \apply{\apply{s}{m}}{x}$ and
                            $t$ converges to $f(x)$ in $\reals$ with respect to $\metricTopology{\standardMetricR}{\reals}$
                            by \cref{choice_function,fst_eq,snd_eq,relation}.
                        For all $m\in\naturalsPlus$ if $N \leq m$ then
                            $\apply{\standardMetricR}{(\apply{\apply{s}{N}}{x},t(m))} < 1$
                            by \cref{cauchy_sequence,sequence_in,funs_type_apply,sup_metric_coordinate_leq,metric_on,times_tuple_intro,bounded_function_set,subseteq,standard_metric_is_metric,real_naturals_subset_reals,real_one_in_naturals,real_zero_lt_one,real_leq_lt_trans,function_apply_intro,real_lt_subst_right}.
                        Therefore $\apply{\standardMetricR}{(\apply{\apply{s}{N}}{x},f(x))} \leq 1$
                            by \cref{bounded_function_set,subseteq,funs_type_apply,standard_metric_is_metric,real_naturals_subset_reals,real_one_in_naturals,real_zero_lt_one,metric_limit_eventual_bound}.
                \end{subproof}
                Thus $\img{f}{X}$ is bounded in $\reals$ with respect to $\standardMetricR$
                    by \cref{pointwise_close_to_bounded_function_bounded,standard_metric_is_metric}.
                Therefore $f\in B$ by \cref{bounded_function_set}.

            Show that for all $U$ if $U$ is a neighborhood of $f$ in $B$ with respect to $\metricTopology{\rho}{B}$ then
                there exists $K\in\naturalsPlus$ such that for all $n\in\naturalsPlus$ if $K\leq n$ then $s(n)\in U$.
            \begin{subproof}
                    Fix $U$ such that $U$ is a neighborhood of $f$ in $B$ with respect to $\metricTopology{\rho}{B}$.
                    Take $\epsilon\in\reals$ such that $0 < \epsilon$ and $\metricBall{\rho}{B}{f}{\epsilon}\subseteq U$
                        by \cref{neighborhood,open_in,topology_on,subseteq,powerset_elim,metric_open_iff}.
                    Let $\delta = \epsilon / (1 + 1)$.
                    We have $\delta\in\reals$ and $0 < \delta$ and $\delta + \delta = \epsilon$ by \cref{real_half_of_positive}.
                    Hence $\delta < \epsilon$ by \cref{real_half_of_positive,real_add_ident,real_lt_add,real_lt_subst_left,real_lt_subst_right}.
                    Take $K\in\naturalsPlus$ such that
                        for all $n,m\in\naturalsPlus$ if $K \leq n$ and $K \leq m$ then
                            $\rho((s(n),s(m))) < \delta$ by \cref{cauchy_sequence}.
                    Show that for all $n\in\naturalsPlus$ if $K\leq n$ then $s(n)\in U$.
                    \begin{subproof}
                        Fix $n\in\naturalsPlus$.
                        Assume $K\leq n$.
                        We have $s(n)\in B$ by \cref{cauchy_sequence,sequence_in,funs_type_apply}.
                        Show that for all $x\in X$ we have $\apply{\standardMetricR}{(\apply{\apply{s}{n}}{x},f(x))} \leq \delta$.
                        \begin{subproof}
                            Fix $x\in X$.
                            Take $t$ such that
                                $t$ is a sequence in $\reals$ and
                                for all $m\in\naturalsPlus$ we have $t(m) = \apply{\apply{s}{m}}{x}$ and
                                $t$ converges to $f(x)$ in $\reals$ with respect to $\metricTopology{\standardMetricR}{\reals}$
                                by \cref{choice_function,fst_eq,snd_eq,relation}.
                            For all $m\in\naturalsPlus$ if $K \leq m$ then
                                $\apply{\standardMetricR}{(\apply{\apply{s}{n}}{x},\apply{\apply{s}{m}}{x})} < \delta$
                                by \cref{cauchy_sequence,sequence_in,funs_type_apply,sup_metric_coordinate_leq,metric_on,times_tuple_intro,bounded_function_set,subseteq,standard_metric_is_metric,real_leq_lt_trans}.
                            Thus for all $m\in\naturalsPlus$ if $K \leq m$ then
                                $\apply{\standardMetricR}{(\apply{\apply{s}{n}}{x},t(m))} < \delta$
                                by \cref{function_apply_intro,real_lt_subst_right}.
                            Therefore $\apply{\standardMetricR}{(\apply{\apply{s}{n}}{x},f(x))} \leq \delta$
                                by \cref{bounded_function_set,subseteq,funs_type_apply,standard_metric_is_metric,metric_limit_eventual_bound}.
                        \end{subproof}
                        Thus $\rho((s(n),f)) \leq \delta$ by \cref{sup_metric_leq_from_coordinate_bounds}.
                        Therefore $\rho((s(n),f)) < \epsilon$ by \cref{times_tuple_intro,metric_on,funs_type_apply,real_leq_lt_trans}.
                        Hence $s(n)\in U$
                            by \cref{metric_on,metric_symmetric,real_lt_subst_left,times_tuple_intro,metric_ball,subseteq}.
                    \end{subproof}
                    Hence $K\in\naturalsPlus$ and
                        for all $n\in\naturalsPlus$ if $K\leq n$ then $s(n)\in U$.
            \end{subproof}
            Hence $f\in B$ and
                $s$ converges to $f$ in $B$ with respect to $\metricTopology{\rho}{B}$
                by \cref{cauchy_sequence,converges_to}.
    \end{subproof}
    Hence $B$ is complete with respect to $\rho$ by \cref{complete_metric_space}.
    Show that $\rho$ is the sup metric on $B$ with respect to $X$ and $\reals$ and $\standardMetricR$ and
        $B$ is complete with respect to $\rho$.
    \begin{subproof}
        Trivial.
    \end{subproof}
\end{proof}

% from §43 helper definition 5: Frechet distance-difference function
\begin{definition}\label{frechet_distance_difference_function}
    $f$ is the Frechet distance-difference function of $a$ with basepoint $x_0$ in $X$ with respect to $d$ iff
        $a\in X$ and
        $x_0\in X$ and
        $f = \{(x, d((x,a)) - d((x,x_0))) \mid x\in X\}$.
\end{definition}

% from §43 helper lemma 19: point-distance values differ by at most the distance of the points
\begin{lemma}\label{metric_point_distance_difference_bound}
    Assume $d$ is a metric on $X$.
    Suppose $x\in X$.
    Suppose $a\in X$.
    Suppose $b\in X$.
    Then $\apply{\standardMetricR}{(d((x,a)),d((x,b)))} \leq d((a,b))$.
\end{lemma}
\begin{proof}
    Let $p = d((x,a))$.
    Let $q = d((x,b))$.
    Let $r = d((a,b))$.
    $(x,a)\in X\times X$ and $(x,b)\in X\times X$ and $(a,b)\in X\times X$ by \cref{times_tuple_intro}.
    Hence $p\in\reals$ and $q\in\reals$ and $r\in\reals$ by \cref{metric_on,funs_type_apply}.
    Thus $d((b,a)) = r$ by \cref{metric_on,metric_symmetric,times_tuple_intro}.
    Therefore $p \leq q + r$ by \cref{metric_on,metric_triangle_inequality,times_tuple_intro}.
    Thus $q \leq p + r$ by \cref{metric_on,metric_triangle_inequality,times_tuple_intro}.
    $\neg{q}\in\reals$ and $\neg{p}\in\reals$ by \cref{real_add_inverse}.
    We have $q + r\in\reals$ and $p + r\in\reals$ and $p + \neg{q}\in\reals$ and $q + \neg{p}\in\reals$ and
        $(q + r) + \neg{q}\in\reals$ and
        $(p + r) + \neg{p}\in\reals$ by \cref{real_add_inverse,real_add_closed}.
    Hence $p + \neg{q} \leq (q + r) + \neg{q}$ by \cref{real_leq_add}.
    Therefore $q + \neg{p} \leq (p + r) + \neg{p}$ by \cref{real_leq_add}.
    Thus $(q + r) + \neg{q} = r$ and $(p + r) + \neg{p} = r$
        by \cref{real_add_ident,real_add_assoc,real_add_comm,real_add_inverse}.
    Therefore $p - q \leq r$ and $q - p \leq r$ and $p - q\in\reals$ by \cref{real_sub}.
    Hence $\abs{(p - q)} \leq r$ by \cref{real_sub_swap_neg,real_lt_subst_left,real_abs_leq_if_pm_leq}.
    Thus $\apply{\standardMetricR}{(p,q)} = \abs{(p - q)}$ by \cref{times_tuple_intro,standard_metric_value}.
    Therefore $\apply{\standardMetricR}{(d((x,a)),d((x,b)))} \leq d((a,b))$.
\end{proof}

% from §43 helper lemma 20: Frechet distance-difference functions are bounded
\begin{lemma}\label{frechet_distance_difference_function_bounded}
    Assume $d$ is a metric on $X$.
    Suppose $f$ is the Frechet distance-difference function of $a$ with basepoint $x_0$ in $X$ with respect to $d$.
    Assume $B$ is the set of bounded functions from $X$ to $\reals$ with respect to $\standardMetricR$.
    Then $f\in B$.
\end{lemma}
\begin{proof}
    We have $a\in X$ and $x_0\in X$ and
        $f = \{(x, d((x,a)) - d((x,x_0))) \mid x\in X\}$ by \cref{frechet_distance_difference_function}.
    Hence $f$ is a function by \cref{pair_eq_iff,relation,rightunique}.
    Thus $\dom{f} = X$ by \cref{dom_intro,dom_iff,pair_eq_iff,frechet_distance_difference_function,subseteq,subseteq_antisymmetric}.
    For all $x\in\dom{f}$ we have $f(x)\in\reals$
        by \cref{subseteq,times_tuple_intro,metric_on,funs_type_apply,real_sub,real_add_inverse,real_add_closed,frechet_distance_difference_function,function_apply_intro}.
    Therefore $f\in\funs{X}{\reals}$ by \cref{funs_intro}.
    Let $M = d((a,x_0))$.
    Hence $M\in\reals$ by \cref{times_tuple_intro,metric_on,funs_type_apply}.
    Show that for all $y_1,y_2\in\img{f}{X}$ we have $\apply{\standardMetricR}{(y_1,y_2)} \leq M + M$.
    \begin{subproof}
            Fix $y_1$.
            Fix $y_2$ such that $y_1\in\img{f}{X}$ and $y_2\in\img{f}{X}$.
            Take $x_1\in X$ such that $(x_1,y_1)\in f$ by \cref{img_iff}.
            Take $x_2\in X$ such that $(x_2,y_2)\in f$ by \cref{img_iff}.
            Hence $y_1 = f(x_1)$ and $y_2 = f(x_2)$ by \cref{function_apply_intro}.
            Thus $f(x_1) = d((x_1,a)) - d((x_1,x_0))$ and
                $f(x_2) = d((x_2,a)) - d((x_2,x_0))$ by \cref{frechet_distance_difference_function,function_apply_intro}.
            $(x_1,a)\in X\times X$ and $(x_1,x_0)\in X\times X$ and
                $(x_2,a)\in X\times X$ and $(x_2,x_0)\in X\times X$ by \cref{times_tuple_intro}.
            We have $d((x_1,a))\in\reals$ and $d((x_1,x_0))\in\reals$ and
                $d((x_2,a))\in\reals$ and $d((x_2,x_0))\in\reals$ by \cref{metric_on,funs_type_apply}.
            Therefore $f(x_1)\in\reals$ and $f(x_2)\in\reals$ by \cref{real_sub,real_add_inverse,real_add_closed}.
            $0\in\reals$ by \cref{real_add_ident}.
            Thus $\apply{\standardMetricR}{(f(x_1),0)} = \abs{f(x_1)}$ and
                $\apply{\standardMetricR}{(f(x_2),0)} = \abs{f(x_2)}$
                by \cref{real_add_ident,times_tuple_intro,standard_metric_value,real_sub,real_neg_zero,real_add_comm,real_lt_subst_right}.
            Therefore $\apply{\standardMetricR}{(f(x_1),0)} \leq M$ and
                $\apply{\standardMetricR}{(f(x_2),0)} \leq M$
                by \cref{metric_point_distance_difference_bound,times_tuple_intro,standard_metric_value,real_lt_subst_right,frechet_distance_difference_function,function_apply_intro}.
            We have $\apply{\standardMetricR}{(y_1,0)}\in\reals$ and
                $\apply{\standardMetricR}{(0,y_2)}\in\reals$ and
                $\apply{\standardMetricR}{(y_1,y_2)}\in\reals$
                by \cref{times_tuple_intro,standard_metric_is_metric,metric_on,funs_type_apply}.
            Thus $\apply{\standardMetricR}{(0,y_2)} \leq M$
                by \cref{standard_metric_is_metric,metric_on,metric_symmetric,times_tuple_intro,real_lt_subst_right}.
            Therefore $\apply{\standardMetricR}{(y_1,y_2)} \leq
                \apply{\standardMetricR}{(y_1,0)} + \apply{\standardMetricR}{(0,y_2)}$
                by \cref{standard_metric_is_metric,metric_on,metric_triangle_inequality}.
            $\apply{\standardMetricR}{(y_1,0)} + \apply{\standardMetricR}{(0,y_2)} \leq
                M + \apply{\standardMetricR}{(0,y_2)}$ by \cref{real_leq_add}.
            We have $\apply{\standardMetricR}{(0,y_2)} + M \leq M + M$ by \cref{real_leq_add}.
            Thus $M + \apply{\standardMetricR}{(0,y_2)} \leq M + M$ by \cref{real_add_comm}.
            $\apply{\standardMetricR}{(y_1,0)} + \apply{\standardMetricR}{(0,y_2)}\in\reals$ and
                $M + \apply{\standardMetricR}{(0,y_2)}\in\reals$ and $M + M\in\reals$ by \cref{real_add_closed}.
            Therefore $\apply{\standardMetricR}{(y_1,0)} + \apply{\standardMetricR}{(0,y_2)} \leq M + M$
                by \cref{real_leq_trans}.
            Hence $\apply{\standardMetricR}{(y_1,y_2)} \leq M + M$ by \cref{real_leq_trans}.
    \end{subproof}
    Hence $\img{f}{X}$ is bounded in $\reals$ with respect to $\standardMetricR$
        by \cref{real_add_closed,img_subseteq_ran,funs_subseteq_rels,subseteq,rels_ran_subseteq,subseteq_transitive,standard_metric_is_metric,metric_bounded}.
    Thus $f\in B$ by \cref{bounded_function_set}.
\end{proof}

% from §43 helper lemma 21: Frechet distance-difference functions are pointwise distance-bounded
\begin{lemma}\label{frechet_distance_difference_pointwise_bound}
    Assume $d$ is a metric on $X$.
    Suppose $f$ is the Frechet distance-difference function of $a$ with basepoint $x_0$ in $X$ with respect to $d$.
    Suppose $g$ is the Frechet distance-difference function of $b$ with basepoint $x_0$ in $X$ with respect to $d$.
    Suppose $x\in X$.
    Then $\apply{\standardMetricR}{(f(x),g(x))} \leq d((a,b))$.
\end{lemma}
\begin{proof}
    Let $B = \{u\in\funs{X}{\reals} \mid \text{$\img{u}{X}$ is bounded in $\reals$ with respect to $\standardMetricR$}\}$.
    We have $a\in X$ and $b\in X$ and $x_0\in X$ by \cref{frechet_distance_difference_function}.
    Thus $f(x) = d((x,a)) - d((x,x_0))$ and $g(x) = d((x,b)) - d((x,x_0))$
        by \cref{frechet_distance_difference_function,frechet_distance_difference_function_bounded,bounded_function_set,subseteq,funs_is_function,function_apply_intro}.
    $(x,a)\in X\times X$ and $(x,b)\in X\times X$ and $(x,x_0)\in X\times X$ by \cref{times_tuple_intro}.
    We have $d((x,a))\in\reals$ and $d((x,b))\in\reals$ and $d((x,x_0))\in\reals$ by \cref{metric_on,funs_type_apply}.
    Therefore $f(x)\in\reals$ and $g(x)\in\reals$ by \cref{real_sub,real_add_inverse,real_add_closed}.
    Thus $\apply{\standardMetricR}{(f(x),g(x))} = \apply{\standardMetricR}{(d((x,a)),d((x,b)))}$
        by \cref{times_tuple_intro,standard_metric_value,real_sub,real_sub_swap_neg,real_lt_subst_left,real_lt_subst_right,real_sub_add_chain}.
    Hence $\apply{\standardMetricR}{(f(x),g(x))} \leq d((a,b))$ by \cref{metric_point_distance_difference_bound}.
\end{proof}

% from §43 helper lemma 22: Frechet distance-difference functions realize the distance at one endpoint
\begin{lemma}\label{frechet_distance_difference_realizes_distance}
    Assume $d$ is a metric on $X$.
    Suppose $f$ is the Frechet distance-difference function of $a$ with basepoint $x_0$ in $X$ with respect to $d$.
    Suppose $g$ is the Frechet distance-difference function of $b$ with basepoint $x_0$ in $X$ with respect to $d$.
    Then $\apply{\standardMetricR}{(f(a),g(a))} = d((a,b))$.
\end{lemma}
\begin{proof}
    Let $B = \{u\in\funs{X}{\reals} \mid \text{$\img{u}{X}$ is bounded in $\reals$ with respect to $\standardMetricR$}\}$.
    We have $a\in X$ and $b\in X$ and $x_0\in X$ by \cref{frechet_distance_difference_function}.
    Thus $f(a) = d((a,a)) - d((a,x_0))$ and $g(a) = d((a,b)) - d((a,x_0))$
        by \cref{frechet_distance_difference_function,frechet_distance_difference_function_bounded,bounded_function_set,subseteq,funs_is_function,function_apply_intro}.
    $(a,a)\in X\times X$ and $(a,b)\in X\times X$ and $(a,x_0)\in X\times X$ by \cref{times_tuple_intro}.
    We have $d((a,a))\in\reals$ and $d((a,b))\in\reals$ and $d((a,x_0))\in\reals$ by \cref{metric_on,funs_type_apply}.
    Therefore $f(a)\in\reals$ and $g(a)\in\reals$ by \cref{real_sub,real_add_inverse,real_add_closed}.
    Thus $f(a) - g(a) = 0 - d((a,b))$
        by \cref{metric_on,metric_zero_characterization,times_tuple_intro,real_sub,real_sub_swap_neg,real_lt_subst_left,real_lt_subst_right,real_sub_add_chain}.
    Therefore $\apply{\standardMetricR}{(f(a),g(a))} = \apply{\standardMetricR}{(0,d((a,b)))}$
        by \cref{times_tuple_intro,real_add_ident,standard_metric_value,real_lt_subst_right}.
    Thus $\apply{\standardMetricR}{(0,d((a,b)))} = d((a,b))$
        by \cref{standard_metric_is_metric,metric_on,metric_symmetric,times_tuple_intro,real_add_ident,standard_metric_value,metric_nonnegative,real_sub,real_neg_zero,real_add_comm,real_lt_subst_right,real_abs_nonneg}.
    Hence $\apply{\standardMetricR}{(f(a),g(a))} = d((a,b))$.
\end{proof}

% from §43 Item 10: isometric embedding into a complete metric space
\begin{theorem}\label{metric_space_isometric_embedding_into_complete}
    Assume $d$ is a metric on $X$.
    Then there exist $Y, D, h$ such that
        $D$ is a metric on $Y$ and
        $Y$ is complete with respect to $D$ and
        $h$ is an isometric embedding from $X$ to $Y$ with respect to $d$ and $D$.
\end{theorem}
\begin{proof}
    \begin{byCase}
    \caseOf{$X = \emptyset$.}
        Let $Y = \reals$.
        Let $D = \standardMetricR$.
        Let $h = \emptyset$.
        $D$ is a metric on $Y$ by \cref{standard_metric_is_metric}.
        $Y$ is complete with respect to $D$ by \cref{reals_complete_standard_metric}.
        Hence $h\in\funs{X}{Y}$ by \cref{codom_of_emptyset_can_be_anything,emptyset_is_function_on_emptyset,funs_intro}.
        For all $a,b$ such that $a\in X$ and $b\in X$ we have $D((h(a),h(b))) = d((a,b))$.
        Show that $D$ is a metric on $Y$ and
            $Y$ is complete with respect to $D$ and
            $h$ is an isometric embedding from $X$ to $Y$ with respect to $d$ and $D$.
        \begin{subproof}
            Follows by \cref{isometric_embedding}.
        \end{subproof}
    \caseOf{$X \neq \emptyset$.}
        Take $x_0$ such that $x_0\in X$ by \cref{nonempty_inhabited}.
        Let $B = \{u\in\funs{X}{\reals} \mid \text{$\img{u}{X}$ is bounded in $\reals$ with respect to $\standardMetricR$}\}$.
        We have $B$ is the set of bounded functions from $X$ to $\reals$ with respect to $\standardMetricR$ by \cref{bounded_function_set}.
        Take $\rho$ such that
            $\rho$ is the sup metric on $B$ with respect to $X$ and $\reals$ and $\standardMetricR$ and
            $B$ is complete with respect to $\rho$ by \cref{bounded_reals_function_space_complete_sup}.
        Hence $\rho$ is a metric on $B$ by \cref{complete_metric_space}.
        Let $Q = \{w\in X\times B \mid
            \text{$\snd{w}$ is the Frechet distance-difference function of $\fst{w}$ with basepoint $x_0$ in $X$ with respect to $d$}\}$.
        $Q$ is a relation by \cref{relation,subseteq,times_elem_is_tuple}.
        Show that for all $a\in X$ there exists $y$ such that $a\mathrel{Q} y$.
        \begin{subproof}
            Fix $a\in X$.
            Let $f = \{(x, d((x,a)) - d((x,x_0))) \mid x\in X\}$.
            Therefore $a\mathrel{Q} f$
                by \cref{frechet_distance_difference_function,frechet_distance_difference_function_bounded,times_tuple_intro,fst_eq,snd_eq,relation}.
        \end{subproof}
        Take $h$ such that
            $h$ is a function on $X$ and
            for all $a\in X$ we have $a\mathrel{Q}\apply{h}{a}$ by \cref{choice_function}.
        Hence $h\in\funs{X}{B}$ by \cref{choice_function,relation,subseteq,times_tuple_elim,funs_intro}.
        Show that for all $a,b\in X$ we have $\rho((h(a),h(b))) = d((a,b))$.
        \begin{subproof}
            Fix $a$.
            Fix $b$ such that $a\in X$ and $b\in X$.
            We have $h(a)$ is the Frechet distance-difference function of $a$ with basepoint $x_0$ in $X$ with respect to $d$ and
                $h(b)$ is the Frechet distance-difference function of $b$ with basepoint $x_0$ in $X$ with respect to $d$
                by \cref{choice_function,relation,fst_eq,snd_eq}.
            Thus $d((a,b))\in\reals$ by \cref{times_tuple_intro,metric_on,funs_type_apply}.
            For all $x\in X$ we have
                $\apply{\standardMetricR}{(\apply{\apply{h}{a}}{x},\apply{\apply{h}{b}}{x})} \leq d((a,b))$
                by \cref{frechet_distance_difference_pointwise_bound}.
            Hence $\rho((h(a),h(b))) \leq d((a,b))$
                by \cref{frechet_distance_difference_function_bounded,sup_metric_leq_from_coordinate_bounds}.
            We have $d((a,b)) \leq \rho((h(a),h(b)))$
                by \cref{frechet_distance_difference_realizes_distance,frechet_distance_difference_function_bounded,sup_metric_coordinate_leq,times_tuple_intro,metric_on,funs_type_apply,real_lt_subst_left}.
            Therefore $\rho((h(a),h(b))) = d((a,b))$ by \cref{complete_metric_space,times_tuple_intro,metric_on,funs_type_apply,real_leq_antisym}.
        \end{subproof}
        Show that $\rho$ is a metric on $B$ and
            $B$ is complete with respect to $\rho$ and
            $h$ is an isometric embedding from $X$ to $B$ with respect to $d$ and $\rho$.
        \begin{subproof}
            Follows by \cref{isometric_embedding}.
        \end{subproof}
    \end{byCase}
\end{proof}

% from §43 Item 11: completion of a metric space
\begin{definition}\label{completion_of_metric_space}
    $Z$ is a completion of $X$ with respect to $d$ iff
        there exist $Y, D, h$ such that
            $D$ is a metric on $Y$ and
            $Y$ is complete with respect to $D$ and
            $h$ is an isometric embedding from $X$ to $Y$ with respect to $d$ and $D$ and
            $Z = \closure{\img{h}{X}}{Y}{\metricTopology{D}{Y}}$.
\end{definition}

\section{§44 A Space-Filling Curve}

% from §44 helper lemma 1: compactness of the domain interval
\begin{lemma}\label{space_filling_domain_compact}
    Let $I = \intervalclosed{0}{1}$.
    Let $T_I = \subspaceTopology{\standardTopologyR}{I}$.
    Then $I$ is compact with respect to $T_I$.
\end{lemma}
\begin{proof}
    Follows by \cref{real_add_ident,real_one_in_naturals,real_naturals_subset_reals,subseteq,real_zero_lt_one,real_lt_imp_leq,real_intervalclosed_compact_standard}.
\end{proof}

% from §44 helper lemma 2: Hausdorffness of the unit square
\begin{lemma}\label{space_filling_target_hausdorff}
    Let $I = \intervalclosed{0}{1}$.
    Let $T_I = \subspaceTopology{\standardTopologyR}{I}$.
    Let $C = \{(i,I) \mid i\in\initialSegment{1 + 1}\}$.
    Let $S = \{(i,T_I) \mid i\in\initialSegment{1 + 1}\}$.
    Let $P = \productTopologyIndexed{S}{C}$.
    Then $\productFamily{C}$ is Hausdorff with respect to $P$.
\end{lemma}
\begin{proof}
    We have $C$ is a function on $\initialSegment{1 + 1}$ and
        $S$ is a function on $\initialSegment{1 + 1}$
        by \cref{relation,rightunique,dom_intro,dom_iff,pair_eq_iff,subseteq,subseteq_antisymmetric}.
    $I$ is Hausdorff with respect to $T_I$
        by \cref{real_add_ident,real_one_in_naturals,real_naturals_subset_reals,subseteq,real_zero_lt_one,real_lt_imp_leq,intervalclosed,standard_metric_is_metric,metric_space_hausdorff,standard_metric_topology,subspace_hausdorff}.
    For all $i\in\initialSegment{1 + 1}$ we have $\apply{C}{i}$ is Hausdorff with respect to $\apply{S}{i}$
        by \cref{function_apply_intro,real_lt_subst_left,real_lt_subst_right}.
    Therefore $\productFamily{C}$ is Hausdorff with respect to $P$
        by \cref{real_one_in_naturals,real_naturals_closed_succ,real_one_leq_naturalsplus,initial_segment_naturalsplus,real_add_eq_subst,product_hausdorff_31}.
\end{proof}

% from §44 helper lemma 3: existence of a dense curve in the unit square
\begin{lemma}\label{space_filling_domain_normal}
    Let $I = \intervalclosed{0}{1}$.
    Let $T_I = \subspaceTopology{\standardTopologyR}{I}$.
    Then $I$ is normal with respect to $T_I$.
\end{lemma}
\begin{proof}
    Follows by \cref{standard_basis_is_basis,basis_topology_is_topology,standard_topology_reals,standard_metric_topology,standard_metric_is_metric,metrizable,intervalclosed,subseteq,subspace_of,subspace_metrizable,metrizable_normal}.
\end{proof}

% from §44 helper lemma 5: points of the unit square are determined by their two coordinates
\begin{lemma}\label{space_filling_square_point_eq}
    Let $I = \intervalclosed{0}{1}$.
    Let $m_2 = 1 + 1$.
    Let $C = \{(i,I) \mid i\in\initialSegment{m_2}\}$.
    Assume $y_1\in\productFamily{C}$.
    Assume $y_2\in\productFamily{C}$.
    Assume $y_1(1) = y_2(1)$.
    Assume $y_1(m_2) = y_2(m_2)$.
    Then $y_1 = y_2$.
\end{lemma}
\begin{proof}
    Let $J = \initialSegment{m_2}$.
    Therefore $C$ is a function on $J$
        by \cref{relation,rightunique,dom_intro,dom_iff,pair_eq_iff,subseteq,subseteq_antisymmetric}.
    For all $i\in J$ we have $i = 1$ or $i = m_2$
        by \cref{initial_segment_naturalsplus,real_naturals_leq_succ_elim,real_one_leq_naturalsplus,real_naturals_subset_reals,real_one_in_naturals,real_leq_antisym,subseteq}.
    For all $i\in\dom{C}$ we have $y_1(i) = y_2(i)$ by \cref{indexed_product,subseteq,pair_eq_iff}.
    Therefore $y_1 = y_2$ by \cref{indexed_product,functions_eq_iff}.
\end{proof}

% from §44 helper lemma 7: binary product space over the unit interval endpoints is compact and Hausdorff
\begin{lemma}\label{space_filling_binary_product_compact_hausdorff}
    Let $I = \intervalclosed{0}{1}$.
    Let $T_I = \subspaceTopology{\standardTopologyR}{I}$.
    Let $D = \{0,1\}$.
    Let $T_D = \subspaceTopology{T_I}{D}$.
    Let $E = \{(n,D) \mid n\in\naturalsPlus\}$.
    Let $U = \{(n,T_D) \mid n\in\naturalsPlus\}$.
    Let $P_B = \productTopologyIndexed{U}{E}$.
    Then $\productFamily{E}$ is compact with respect to $P_B$ and
        $\productFamily{E}$ is Hausdorff with respect to $P_B$.
\end{lemma}
\begin{proof}
    $0\in\reals$ and $1\in\reals$ and $0 < 1$
        by \cref{real_add_ident,real_one_in_naturals,real_naturals_subset_reals,subseteq,real_zero_lt_one}.
    Hence $0\in I$ and $1\in I$ by \cref{intervalclosed,real_lt_imp_leq}.
    Thus $D\subseteq I$ by \cref{upair_iff,subseteq}.
    We have $T_I$ is a topology on $I$
        by \cref{intervalclosed,subseteq,standard_basis_is_basis,basis_topology_is_topology,standard_topology_reals,subspace_of,subspace_topology_is_topology}.
    Therefore $I$ is Hausdorff with respect to $T_I$
        by \cref{intervalclosed,subseteq,standard_metric_is_metric,metric_space_hausdorff,standard_metric_topology,subspace_hausdorff}.
    Thus $D$ is closed in $I$ with respect to $T_I$
        by \cref{singleton_subset_intro,upair_intro_left,pair_is_finite,subseteq_finite_is_finite,finite_point_set_closed_hausdorff}.
    Therefore $D$ is compact with respect to $T_D$ and $D$ is Hausdorff with respect to $T_D$
        by \cref{space_filling_domain_compact,closed_subspace_compact,subspace_hausdorff}.
    We have $E$ is a function on $\naturalsPlus$ and
        $U$ is a function on $\naturalsPlus$
        by \cref{relation,rightunique,dom_intro,dom_iff,pair_eq_iff,subseteq,subseteq_antisymmetric}.
    Hence $\dom{E}$ is inhabited by \cref{funs,real_one_in_naturals,subseteq}.
    For all $n\in\naturalsPlus$ we have $E(n) = D$ and $U(n) = T_D$ by \cref{function_apply_intro}.
    For all $n\in\dom{E}$ we have $U(n)$ is a topology on $E(n)$ by \cref{function_apply_intro,subspace_topology_is_topology,subspace_of,subseteq}.
    For all $n\in\dom{E}$ we have $E(n)$ is compact with respect to $U(n)$ and
        $E(n)$ is Hausdorff with respect to $U(n)$
        by \cref{function_apply_intro,subseteq}.
    Show that $\productFamily{E}$ is compact with respect to $P_B$ and
        $\productFamily{E}$ is Hausdorff with respect to $P_B$.
    \begin{subproof}
        Follows by \cref{funs,real_one_in_naturals,subseteq,tychonoff_theorem,product_hausdorff_product}.
    \end{subproof}
\end{proof}

% from §44 helper lemma 8: the endpoint subspace of the unit interval is discrete
\begin{lemma}\label{space_filling_binary_endpoint_discrete}
    Let $I = \intervalclosed{0}{1}$.
    Let $T_I = \subspaceTopology{\standardTopologyR}{I}$.
    Let $D = \{0,1\}$.
    Let $T_D = \subspaceTopology{T_I}{D}$.
    Then $\{0\}$ is open in $D$ with respect to $T_D$ and
        $\{1\}$ is open in $D$ with respect to $T_D$.
\end{lemma}
\begin{proof}
    We have $T_I$ is a topology on $I$
        by \cref{intervalclosed,subseteq,standard_basis_is_basis,basis_topology_is_topology,standard_topology_reals,subspace_of,subspace_topology_is_topology}.
    $0\in\reals$ and $1\in\reals$ and $0 < 1$
        by \cref{real_add_ident,real_one_in_naturals,real_naturals_subset_reals,subseteq,real_zero_lt_one}.
    Hence $0\in I$ and $1\in I$ by \cref{intervalclosed,real_lt_imp_leq}.
    Thus $0\in D$ and $1\in D$ by \cref{upair_intro_left,upair_intro_right}.
    Therefore $D\subseteq I$ by \cref{upair_iff,subseteq}.
    Thus $I$ is Hausdorff with respect to $T_I$
        by \cref{intervalclosed,subseteq,standard_metric_is_metric,metric_space_hausdorff,standard_metric_topology,subspace_hausdorff}.
    Therefore $D$ has closed singletons with respect to $T_D$
        by \cref{subspace_hausdorff,subspace_topology_is_topology,subspace_of,singleton_subset_intro,pair_is_finite,singleton_iff,upair_intro_left,subseteq,subseteq_finite_is_finite,finite_point_set_closed_hausdorff,one_point_closed}.
    Thus $D\setminus\{0\}$ is open in $D$ with respect to $T_D$ and
        $D\setminus\{1\}$ is open in $D$ with respect to $T_D$
        by \cref{one_point_closed,closed_in,setminus_subseteq}.
    We have $D\setminus\{0\}\subseteq\{1\}$
        by \cref{upair_iff,setminus_elim_left,setminus_elim_right,singleton_intro,singleton_elim,subseteq}.
    Therefore $\{1\}\subseteq D\setminus\{0\}$
        by \cref{singleton_elim,upair_iff,setminus_intro,singleton_iff,real_lt_irreflexive,real_lt_subst_right,subseteq}.
    We have $D\setminus\{1\}\subseteq\{0\}$
        by \cref{upair_iff,setminus_elim_left,setminus_elim_right,singleton_intro,singleton_elim,subseteq}.
    Therefore $\{0\}\subseteq D\setminus\{1\}$
        by \cref{singleton_elim,upair_iff,setminus_intro,singleton_iff,real_lt_irreflexive,real_lt_subst_left,subseteq}.
    Show that $\{0\}$ is open in $D$ with respect to $T_D$ and
        $\{1\}$ is open in $D$ with respect to $T_D$.
    \begin{subproof}
        Follows by \cref{subseteq_antisymmetric,open_in}.
    \end{subproof}
\end{proof}

% from §44 helper lemma 9: one-coordinate cylinders in the binary product are open
\begin{lemma}\label{space_filling_binary_coordinate_cylinders_open}
    Let $I = \intervalclosed{0}{1}$.
    Let $T_I = \subspaceTopology{\standardTopologyR}{I}$.
    Let $D = \{0,1\}$.
    Let $T_D = \subspaceTopology{T_I}{D}$.
    Let $E = \{(n,D) \mid n\in\naturalsPlus\}$.
    Let $U = \{(n,T_D) \mid n\in\naturalsPlus\}$.
    Let $P_B = \productTopologyIndexed{U}{E}$.
    Assume $n\in\naturalsPlus$.
    Then $\preimg{\piIndex{E}{n}}{\{0\}}$ is open in $\productFamily{E}$ with respect to $P_B$ and
        $\preimg{\piIndex{E}{n}}{\{1\}}$ is open in $\productFamily{E}$ with respect to $P_B$.
\end{lemma}
\begin{proof}
    We have $E$ is a function on $\naturalsPlus$ and
        $U$ is a function on $\naturalsPlus$
        by \cref{relation,rightunique,dom_intro,dom_iff,pair_eq_iff,subseteq,subseteq_antisymmetric}.
    For all $m\in\naturalsPlus$ we have $E(m) = D$ and $U(m) = T_D$ by \cref{function_apply_intro}.
    $T_I$ is a topology on $I$
        by \cref{intervalclosed,subseteq,standard_basis_is_basis,basis_topology_is_topology,standard_topology_reals,subspace_of,subspace_topology_is_topology}.
    $0\in\reals$ and $1\in\reals$ and $0 < 1$
        by \cref{real_add_ident,real_one_in_naturals,real_naturals_subset_reals,subseteq,real_zero_lt_one}.
    Hence $0\in I$ and $1\in I$ by \cref{intervalclosed,real_lt_imp_leq}.
    Therefore $D\subseteq I$ by \cref{upair_iff,subseteq}.
    For all $m\in\dom{E}$ we have $U(m)$ is a topology on $E(m)$
        by \cref{function_apply_intro,subspace_of,subseteq,subspace_topology_is_topology}.
    Thus $\piIndex{E}{n}$ is continuous from $\productFamily{E}$ to $D$ with respect to $P_B$ and $T_D$
        by \cref{funs,real_one_in_naturals,subseteq,projection_map_indexed_continuous}.
    Show that $\preimg{\piIndex{E}{n}}{\{0\}}$ is open in $\productFamily{E}$ with respect to $P_B$ and
        $\preimg{\piIndex{E}{n}}{\{1\}}$ is open in $\productFamily{E}$ with respect to $P_B$.
    \begin{subproof}
        Follows by \cref{space_filling_binary_endpoint_discrete,continuous}.
    \end{subproof}
\end{proof}

% from §44 helper lemma 10: finite binary digit assignments extend to product points
\begin{lemma}\label{space_filling_binary_prefix_extension}
    Let $I = \intervalclosed{0}{1}$.
    Let $D = \{0,1\}$.
    Let $E = \{(n,D) \mid n\in\naturalsPlus\}$.
    Assume $k\in\naturalsPlus$.
    Assume $a\in\funs{\initialSegment{k}}{D}$.
    Then there exists $x\in\productFamily{E}$ such that
        for all $i\in\initialSegment{k}$ we have $x(i) = a(i)$.
\end{lemma}
\begin{proof}
    We have $E$ is a function on $\naturalsPlus$ and $\dom{E} = \naturalsPlus$
        by \cref{relation,rightunique,dom_intro,dom_iff,pair_eq_iff,subseteq,subseteq_antisymmetric,funs}.
    For all $m\in\naturalsPlus$ we have $E(m) = D$ by \cref{function_apply_intro}.
    Let $t = \{(i,0) \mid i\in\naturalsPlus\setminus\initialSegment{k}\}$.
    For all $i,y_1,y_2$ such that $(i,y_1)\in t$ and $(i,y_2)\in t$ we have $y_1 = y_2$ by \cref{pair_eq_iff}.
    Hence $t$ is a function and $\dom{t} = \naturalsPlus\setminus\initialSegment{k}$
        by \cref{relation,rightunique,dom_intro,dom_iff,pair_eq_iff,subseteq,subseteq_antisymmetric}.
    For all $i\in\dom{t}$ we have $t(i)\in D$ by \cref{function_apply_intro,subseteq,upair_intro_left}.
    Thus $t\in\funs{\naturalsPlus\setminus\initialSegment{k}}{D}$ by \cref{funs_intro}.
    For all $u$ such that $u\in\dom{a}\inter\dom{t}$ we have $a(u) = t(u)$
        by \cref{funs,inter,setminus_elim_right}.
    Therefore $a$ is right-unique and $a$ is a relation and $t$ is right-unique and $t$ is a relation
        by \cref{funs_is_function,function_rightunique,funs_is_relation}.
    Let $x = a\union t$.
    Thus $x$ is a function by \cref{union_compatible_functions,relation,rightunique}.
    Hence $\dom{x} = \initialSegment{k}\union(\naturalsPlus\setminus\initialSegment{k})$
        by \cref{dom_union,funs,real_lt_subst_right}.
    Therefore $\naturalsPlus\subseteq\dom{x}$ by \cref{union_iff,setminus_intro,subseteq}.
    $\dom{a}\subseteq\naturalsPlus$ and $\dom{t}\subseteq\naturalsPlus$
        by \cref{funs,initial_segment_naturalsplus,setminus_elim_left,subseteq}.
    We have $\dom{x}\subseteq\naturalsPlus$ by \cref{dom_union,union_subsets_is_subset,subseteq}.
    Thus $\dom{x} = \naturalsPlus$ by \cref{subseteq_antisymmetric}.
    Therefore $\dom{x} = \dom{E}$ by \cref{subseteq}.
    For all $i\in\dom{x}$ we have $x(i)\in D$
        by \cref{dom_union,union_iff,function_apply_elim,function_apply_intro,funs_elim,funs_type_apply,subseteq,setminus_elim_left}.
    Thus $x\in\funs{\naturalsPlus}{D}$ by \cref{funs_intro}.
    For all $i\in\dom{E}$ we have $x(i)\in E(i)$ by \cref{function_apply_intro,funs_type_apply,subseteq}.
    For all $i\in\dom{x}$ we have $x(i)\in\unions{\ran{E}}$
        by \cref{subseteq,function_apply_elim,ran_intro,unions_iff,upair_intro_left}.
    Therefore $x\in\funs{\dom{E}}{\unions{\ran{E}}}$ by \cref{funs_intro}.
    Hence $x\in\productFamily{E}$ by \cref{indexed_product}.
    For all $i\in\initialSegment{k}$ we have $x(i) = a(i)$
        by \cref{union_upper_left,funs,function_apply_elim,subseteq,funs_is_function,function_apply_intro}.
    Thus $x\in\productFamily{E}$ and
        for all $i\in\initialSegment{k}$ we have $x(i) = a(i)$ by \cref{subseteq}.
\end{proof}

% from §44 helper definition 1: binary prefix cylinder in the product space
\begin{definition}\label{space_filling_binary_prefix_cylinder}
    $C$ is the binary prefix cylinder at $a$ in $E$ iff
        $C = \{x\in\productFamily{E} \mid \text{for all $i\in\dom{a}$ we have $x(i) = a(i)$}\}$.
\end{definition}

% from §44 helper lemma 11: every finite binary prefix cylinder is inhabited
\begin{lemma}\label{space_filling_binary_prefix_cylinder_inhabited}
    Let $D = \{0,1\}$.
    Let $E = \{(n,D) \mid n\in\naturalsPlus\}$.
    Assume $k\in\naturalsPlus$.
    Assume $a\in\funs{\initialSegment{k}}{D}$.
    Assume $C$ is the binary prefix cylinder at $a$ in $E$.
    Then $C$ is inhabited.
\end{lemma}
\begin{proof}
    Follows by \cref{space_filling_binary_prefix_extension,funs,subseteq,space_filling_binary_prefix_cylinder}.
\end{proof}

% from §44 helper lemma 12: every finite binary prefix cylinder is open
\begin{lemma}\label{space_filling_binary_prefix_cylinder_open}
    Let $I = \intervalclosed{0}{1}$.
    Let $T_I = \subspaceTopology{\standardTopologyR}{I}$.
    Let $D = \{0,1\}$.
    Let $T_D = \subspaceTopology{T_I}{D}$.
    Let $E = \{(n,D) \mid n\in\naturalsPlus\}$.
    Let $U = \{(n,T_D) \mid n\in\naturalsPlus\}$.
    Let $P_B = \productTopologyIndexed{U}{E}$.
    Assume $k\in\naturalsPlus$.
    Assume $a\in\funs{\initialSegment{k}}{D}$.
    Assume $C$ is the binary prefix cylinder at $a$ in $E$.
    Then $C$ is open in $\productFamily{E}$ with respect to $P_B$.
\end{lemma}
\begin{proof}
    $\productFamily{E}$ is Hausdorff with respect to $P_B$ by \cref{space_filling_binary_product_compact_hausdorff}.
    Hence $P_B$ is a topology on $\productFamily{E}$ by \cref{hausdorff}.
    Let $g = \{(i,\preimg{\piIndex{E}{i}}{\{a(i)\}}) \mid i\in\initialSegment{k}\}$.
    For all $i,V_1,V_2$ such that $(i,V_1)\in g$ and $(i,V_2)\in g$ we have $V_1 = V_2$ by \cref{pair_eq_iff}.
    Thus $g$ is a function and $\dom{g} = \initialSegment{k}$
        by \cref{relation,rightunique,pair_eq_iff,dom_intro,dom_iff,subseteq,subseteq_antisymmetric}.
    Show that for all $i\in\dom{g}$ we have $g(i)$ is open in $\productFamily{E}$ with respect to $P_B$.
    \begin{subproof}
        Fix $i\in\dom{g}$.
        We have $i\in\naturalsPlus$ by \cref{subseteq,initial_segment_naturalsplus}.
        Thus $a(i)\in D$ by \cref{funs_type_apply}.
        Therefore $a(i) = 0$ or $a(i) = 1$ by \cref{upair_iff}.
        \begin{byCase}
        \caseOf{$a(i) = 0$.}
            We have $g(i)$ is open in $\productFamily{E}$ with respect to $P_B$
                by \cref{space_filling_binary_coordinate_cylinders_open,pair_eq_iff,function_apply_intro,open_in}.
        \caseOf{$a(i) = 1$.}
            $g(i)$ is open in $\productFamily{E}$ with respect to $P_B$
                by \cref{space_filling_binary_coordinate_cylinders_open,pair_eq_iff,function_apply_intro,open_in}.
        \end{byCase}
    \end{subproof}
    We have $\img{g}{\initialSegment{k}}$ is finite and $\img{g}{\initialSegment{k}}$ is inhabited
        by \cref{initial_segment_finite,finite_image_function,real_one_in_naturals,real_one_leq_naturalsplus,initial_segment_naturalsplus,subseteq,function_apply_elim,img_iff}.
    Thus $\inters{\img{g}{\initialSegment{k}}}$ is open in $\productFamily{E}$ with respect to $P_B$
        by \cref{img_iff,function_apply_elim,open_in,subseteq,finite_intersections_open_in_topology}.
    For all $x$ such that $x\in C$ we have $x\in\inters{\img{g}{\initialSegment{k}}}$
        by \cref{img_iff,pair_eq_iff,funs,subseteq,space_filling_binary_prefix_cylinder,singleton_intro,projection_map_indexed,preimg_iff,funs_is_function,initial_segment_naturalsplus,real_one_in_naturals,inters_intro}.
    Therefore $C\subseteq \inters{\img{g}{\initialSegment{k}}}$ by \cref{subseteq}.
    Show that for all $x$ such that $x\in\inters{\img{g}{\initialSegment{k}}}$ we have $x\in C$.
    \begin{subproof}
        Fix $x$ such that $x\in\inters{\img{g}{\initialSegment{k}}}$.
        Thus $x\in \preimg{\piIndex{E}{1}}{\{a(1)\}}$
            by \cref{real_one_in_naturals,real_one_leq_naturalsplus,initial_segment_naturalsplus,subseteq,pair_eq_iff,function_apply_intro,img_iff,function_apply_elim,inters_destr}.
        We have $x\in\productFamily{E}$ by \cref{projection_map_indexed,preimg_iff,pair_eq_iff}.
        For all $i\in\dom{a}$ we have $x(i) = a(i)$
            by \cref{funs,pair_eq_iff,img_iff,inters_destr,projection_map_indexed,indexed_product,funs_is_function,function_apply_iff,preimg_iff,singleton_elim}.
        Therefore $x\in C$ by \cref{space_filling_binary_prefix_cylinder}.
    \end{subproof}
    Thus $\inters{\img{g}{\initialSegment{k}}}\subseteq C$ by \cref{subseteq}.
    Therefore $C$ is open in $\productFamily{E}$ with respect to $P_B$ by \cref{subseteq_antisymmetric,open_in}.
\end{proof}

% from §44 helper lemma 13: one-coordinate cylinders in the binary product are closed
\begin{lemma}\label{space_filling_binary_coordinate_cylinders_closed}
    Let $I = \intervalclosed{0}{1}$.
    Let $T_I = \subspaceTopology{\standardTopologyR}{I}$.
    Let $D = \{0,1\}$.
    Let $T_D = \subspaceTopology{T_I}{D}$.
    Let $E = \{(n,D) \mid n\in\naturalsPlus\}$.
    Let $U = \{(n,T_D) \mid n\in\naturalsPlus\}$.
    Let $P_B = \productTopologyIndexed{U}{E}$.
    Assume $n\in\naturalsPlus$.
    Then $\preimg{\piIndex{E}{n}}{\{0\}}$ is closed in $\productFamily{E}$ with respect to $P_B$ and
        $\preimg{\piIndex{E}{n}}{\{1\}}$ is closed in $\productFamily{E}$ with respect to $P_B$.
\end{lemma}
\begin{proof}
    Let $Z_0 = \preimg{\piIndex{E}{n}}{\{0\}}$.
    Let $Z_1 = \preimg{\piIndex{E}{n}}{\{1\}}$.
    We have $Z_0$ is open in $\productFamily{E}$ with respect to $P_B$ and
        $Z_1$ is open in $\productFamily{E}$ with respect to $P_B$
        by \cref{space_filling_binary_coordinate_cylinders_open}.
    $\productFamily{E}$ is Hausdorff with respect to $P_B$ by \cref{space_filling_binary_product_compact_hausdorff}.
    Hence $P_B$ is a topology on $\productFamily{E}$ by \cref{hausdorff}.
    $Z_0\subseteq\productFamily{E}$ and $Z_1\subseteq\productFamily{E}$
        by \cref{preimg_iff,projection_map_indexed,pair_eq_iff,subseteq}.
    We have $E$ is a function on $\naturalsPlus$ and $n\in\dom{E}$
        by \cref{relation,rightunique,dom_intro,dom_iff,pair_eq_iff,subseteq,subseteq_antisymmetric,funs}.

    For all $x$ such that $x\in\productFamily{E}\setminus Z_0$ we have $x\in Z_1$
        by \cref{setminus_elim_left,setminus_elim_right,indexed_product,subseteq,function_apply_intro,projection_map_indexed,preimg_iff,pair_eq_iff,upair_iff,singleton_iff,singleton_intro}.
    Hence $\productFamily{E}\setminus Z_0\subseteq Z_1$ by \cref{subseteq}.
    For all $x$ such that $x\in Z_1$ we have $x\in\productFamily{E}\setminus Z_0$
        by \cref{preimg_iff,singleton_elim,projection_map_indexed,pair_eq_iff,real_add_ident,real_one_in_naturals,real_naturals_subset_reals,subseteq,real_zero_lt_one,singleton_iff,real_lt_irreflexive,real_lt_subst_left,setminus_intro}.
    Thus $Z_1\subseteq\productFamily{E}\setminus Z_0$ by \cref{subseteq}.
    Therefore $\productFamily{E}\setminus Z_0$ is open in $\productFamily{E}$ with respect to $P_B$
        by \cref{subseteq_antisymmetric,open_in}.
    Thus $Z_0$ is closed in $\productFamily{E}$ with respect to $P_B$ by \cref{closed_in}.

    For all $x$ such that $x\in\productFamily{E}\setminus Z_1$ we have $x\in Z_0$
        by \cref{setminus_elim_left,setminus_elim_right,indexed_product,subseteq,function_apply_intro,projection_map_indexed,preimg_iff,pair_eq_iff,upair_iff,singleton_iff,singleton_intro}.
    Hence $\productFamily{E}\setminus Z_1\subseteq Z_0$ by \cref{subseteq}.
    For all $x$ such that $x\in Z_0$ we have $x\in\productFamily{E}\setminus Z_1$
        by \cref{preimg_iff,singleton_elim,projection_map_indexed,pair_eq_iff,real_add_ident,real_one_in_naturals,real_naturals_subset_reals,subseteq,real_zero_lt_one,singleton_iff,real_lt_irreflexive,real_lt_subst_right,setminus_intro}.
    Therefore $Z_0\subseteq\productFamily{E}\setminus Z_1$ by \cref{subseteq}.
    Thus $\productFamily{E}\setminus Z_1$ is open in $\productFamily{E}$ with respect to $P_B$
        by \cref{subseteq_antisymmetric,open_in}.
    Therefore $Z_1$ is closed in $\productFamily{E}$ with respect to $P_B$ by \cref{closed_in}.
\end{proof}

% from §44 helper lemma 14: every finite binary prefix cylinder is closed
\begin{lemma}\label{space_filling_binary_prefix_cylinder_closed}
    Let $I = \intervalclosed{0}{1}$.
    Let $T_I = \subspaceTopology{\standardTopologyR}{I}$.
    Let $D = \{0,1\}$.
    Let $T_D = \subspaceTopology{T_I}{D}$.
    Let $E = \{(n,D) \mid n\in\naturalsPlus\}$.
    Let $U = \{(n,T_D) \mid n\in\naturalsPlus\}$.
    Let $P_B = \productTopologyIndexed{U}{E}$.
    Assume $k\in\naturalsPlus$.
    Assume $a\in\funs{\initialSegment{k}}{D}$.
    Assume $C$ is the binary prefix cylinder at $a$ in $E$.
    Then $C$ is closed in $\productFamily{E}$ with respect to $P_B$.
\end{lemma}
\begin{proof}
    $\productFamily{E}$ is Hausdorff with respect to $P_B$ by \cref{space_filling_binary_product_compact_hausdorff}.
    Hence $P_B$ is a topology on $\productFamily{E}$ by \cref{hausdorff}.
    Let $g = \{(i,\preimg{\piIndex{E}{i}}{\{a(i)\}}) \mid i\in\initialSegment{k}\}$.
    For all $i,V_1,V_2$ such that $(i,V_1)\in g$ and $(i,V_2)\in g$ we have $V_1 = V_2$ by \cref{pair_eq_iff}.
    Thus $g$ is a function and $\dom{g} = \initialSegment{k}$
        by \cref{relation,rightunique,pair_eq_iff,dom_intro,dom_iff,subseteq,subseteq_antisymmetric}.
    Show that for all $i\in\dom{g}$ we have $g(i)$ is closed in $\productFamily{E}$ with respect to $P_B$.
    \begin{subproof}
        Fix $i\in\dom{g}$.
        We have $i\in\naturalsPlus$ by \cref{subseteq,initial_segment_naturalsplus}.
        Thus $a(i)\in D$ by \cref{funs_type_apply}.
        Therefore $a(i) = 0$ or $a(i) = 1$ by \cref{upair_iff}.
        \begin{byCase}
        \caseOf{$a(i) = 0$.}
            We have $g(i)$ is closed in $\productFamily{E}$ with respect to $P_B$
                by \cref{space_filling_binary_coordinate_cylinders_closed,pair_eq_iff,function_apply_intro,closed_in}.
        \caseOf{$a(i) = 1$.}
            $g(i)$ is closed in $\productFamily{E}$ with respect to $P_B$
                by \cref{space_filling_binary_coordinate_cylinders_closed,pair_eq_iff,function_apply_intro,closed_in}.
        \end{byCase}
    \end{subproof}
    For all $W\in\img{g}{\initialSegment{k}}$ we have $W$ is closed in $\productFamily{E}$ with respect to $P_B$
        by \cref{img_iff,function_apply_elim}.
    Thus $\inters{\img{g}{\initialSegment{k}}}$ is closed in $\productFamily{E}$ with respect to $P_B$
        by \cref{img_iff,function_apply_elim,closed_in,subseteq,powerset_intro,closed_inters}.
    Show that $C = \inters{\img{g}{\initialSegment{k}}}$.
    \begin{subproof}
        For all $x$ such that $x\in C$ we have $x\in\inters{\img{g}{\initialSegment{k}}}$
            by \cref{img_iff,pair_eq_iff,funs,subseteq,space_filling_binary_prefix_cylinder,singleton_intro,projection_map_indexed,preimg_iff,funs_is_function,initial_segment_naturalsplus,real_one_in_naturals,inters_intro}.
        Therefore $C\subseteq \inters{\img{g}{\initialSegment{k}}}$ by \cref{subseteq}.
        Show that for all $x$ such that $x\in\inters{\img{g}{\initialSegment{k}}}$ we have $x\in C$.
        \begin{subproof}
            Fix $x$ such that $x\in\inters{\img{g}{\initialSegment{k}}}$.
            Thus $x\in \preimg{\piIndex{E}{1}}{\{a(1)\}}$
                by \cref{real_one_in_naturals,real_one_leq_naturalsplus,initial_segment_naturalsplus,subseteq,pair_eq_iff,function_apply_intro,img_iff,function_apply_elim,inters_destr}.
            We have $x\in\productFamily{E}$ by \cref{projection_map_indexed,preimg_iff,pair_eq_iff}.
            For all $i\in\dom{a}$ we have $x(i) = a(i)$
                by \cref{funs,pair_eq_iff,img_iff,inters_destr,projection_map_indexed,indexed_product,funs_is_function,function_apply_iff,preimg_iff,singleton_elim}.
            Therefore $x\in C$ by \cref{space_filling_binary_prefix_cylinder}.
        \end{subproof}
        Thus $\inters{\img{g}{\initialSegment{k}}}\subseteq C$ by \cref{subseteq}.
    \end{subproof}
    Therefore $C$ is closed in $\productFamily{E}$ with respect to $P_B$ by \cref{subseteq_antisymmetric,closed_in}.
\end{proof}

% from §44 helper lemma 15: every finite binary prefix cylinder is compact
\begin{lemma}\label{space_filling_binary_prefix_cylinder_compact}
    Let $I = \intervalclosed{0}{1}$.
    Let $T_I = \subspaceTopology{\standardTopologyR}{I}$.
    Let $D = \{0,1\}$.
    Let $T_D = \subspaceTopology{T_I}{D}$.
    Let $E = \{(n,D) \mid n\in\naturalsPlus\}$.
    Let $U = \{(n,T_D) \mid n\in\naturalsPlus\}$.
    Let $P_B = \productTopologyIndexed{U}{E}$.
    Assume $k\in\naturalsPlus$.
    Assume $a\in\funs{\initialSegment{k}}{D}$.
    Assume $C$ is the binary prefix cylinder at $a$ in $E$.
    Then $C$ is compact with respect to $\subspaceTopology{P_B}{C}$.
\end{lemma}
\begin{proof}
    Follows by \cref{space_filling_binary_product_compact_hausdorff,hausdorff,space_filling_binary_prefix_cylinder,subseteq,space_filling_binary_prefix_cylinder_closed,closed_subspace_compact}.
\end{proof}

% from §44 helper lemma 16: compact finite-intersection families have inhabited total intersection
\begin{lemma}\label{space_filling_compact_fip_inhabited_intersection}
    Assume $T$ is a topology on $X$.
    Assume $X$ is compact with respect to $T$.
    Suppose $C$ has the finite intersection property.
    Suppose for all $A$ such that $A\in C$ we have $A$ is closed in $X$ with respect to $T$.
    Then $\inters{C}$ is inhabited.
\end{lemma}
\begin{proof}
    Follows by \cref{compact_iff_finite_intersection_property,nonempty_inhabited}.
\end{proof}

% from §44 helper lemma 17: doubling and halving a real gives the original real
\begin{lemma}\label{space_filling_double_half_real}
    Suppose $a\in\reals$.
    Then $(a + a) / (1 + 1) = a$.
\end{lemma}
\begin{proof}
    Follows by \cref{real_one_in_naturals,real_naturals_closed_succ,real_div_add_same_denom,real_half_plus_half}.
\end{proof}

% from §44 helper lemma 18: the midpoint is not below the left endpoint
\begin{lemma}\label{space_filling_left_endpoint_leq_midpoint}
    Suppose $a\in\reals$ and $b\in\reals$ and $a \leq b$.
    Then $a \leq (a + b) / (1 + 1)$.
\end{lemma}
\begin{proof}
    Follows by \cref{real_one_in_naturals,real_naturals_closed_succ,real_leq_add,real_add_comm,real_leq_subst_right_local,real_add_closed,real_div_leq_same_denom,space_filling_double_half_real,real_leq_subst_left_local}.
\end{proof}

% from §44 helper lemma 19: the midpoint does not exceed the right endpoint
\begin{lemma}\label{space_filling_midpoint_leq_right_endpoint}
    Suppose $a\in\reals$ and $b\in\reals$ and $a \leq b$.
    Then $(a + b) / (1 + 1) \leq b$.
\end{lemma}
\begin{proof}
    Follows by \cref{real_one_in_naturals,real_naturals_closed_succ,real_leq_add,real_add_closed,real_div_leq_same_denom,space_filling_double_half_real,real_leq_subst_right_local}.
\end{proof}

% from §44 helper lemma 20: the left dyadic half lies in the parent interval
\begin{lemma}\label{space_filling_left_half_subset_parent}
    Suppose $a\in\reals$ and $b\in\reals$ and $a \leq b$.
    Let $J = \intervalclosed{a}{(a + b) / (1 + 1)}$.
    Then $J\subseteq \intervalclosed{a}{b}$.
\end{lemma}
\begin{proof}
    Follows by \cref{intervalclosed,subseteq,real_add_closed,real_one_in_naturals,real_naturals_closed_succ,real_div_naturalsplus_real,space_filling_midpoint_leq_right_endpoint,real_leq_trans}.
\end{proof}

% from §44 helper lemma 21: the right dyadic half lies in the parent interval
\begin{lemma}\label{space_filling_right_half_subset_parent}
    Suppose $a\in\reals$ and $b\in\reals$ and $a \leq b$.
    Let $J = \intervalclosed{(a + b) / (1 + 1)}{b}$.
    Then $J\subseteq \intervalclosed{a}{b}$.
\end{lemma}
\begin{proof}
    Follows by \cref{intervalclosed,subseteq,real_add_closed,real_one_in_naturals,real_naturals_closed_succ,real_div_naturalsplus_real,space_filling_left_endpoint_leq_midpoint,real_leq_trans}.
\end{proof}

% from §44 helper lemma 22: the two dyadic halves cover the parent interval
\begin{lemma}\label{space_filling_halves_cover_parent}
    Suppose $a\in\reals$ and $b\in\reals$ and $a \leq b$.
    Let $m = (a + b) / (1 + 1)$.
    Let $J_0 = \intervalclosed{a}{m}$.
    Let $J_1 = \intervalclosed{m}{b}$.
    Then $\intervalclosed{a}{b}\subseteq J_0\union J_1$.
\end{lemma}
\begin{proof}
    Follows by \cref{subseteq,intervalclosed,real_add_closed,real_one_in_naturals,real_naturals_closed_succ,real_div_naturalsplus_real,real_lt_trichotomy,real_lt_imp_leq,union_intro_left,union_intro_right}.
\end{proof}

% from §44 helper lemma 23: a closed interval is the union of its two dyadic halves
\begin{lemma}\label{space_filling_parent_eq_union_halves}
    Suppose $a\in\reals$ and $b\in\reals$ and $a \leq b$.
    Let $m = (a + b) / (1 + 1)$.
    Let $J_0 = \intervalclosed{a}{m}$.
    Let $J_1 = \intervalclosed{m}{b}$.
    Then $\intervalclosed{a}{b} = J_0\union J_1$.
\end{lemma}
\begin{proof}
    Follows by \cref{space_filling_left_half_subset_parent,space_filling_right_half_subset_parent,union_subsets_is_subset,space_filling_halves_cover_parent,subseteq_antisymmetric}.
\end{proof}

% from §44 helper lemma 24: every point of a closed interval belongs to a dyadic half
\begin{lemma}\label{space_filling_point_in_some_half}
    Suppose $a\in\reals$ and $b\in\reals$ and $a \leq b$.
    Let $m = (a + b) / (1 + 1)$.
    Let $J_0 = \intervalclosed{a}{m}$.
    Let $J_1 = \intervalclosed{m}{b}$.
    Assume $y\in\intervalclosed{a}{b}$.
    Then $y\in J_0$ or $y\in J_1$.
\end{lemma}
\begin{proof}
    Follows by \cref{space_filling_parent_eq_union_halves,subseteq,union_iff}.
\end{proof}

% from §44 helper lemma 25: one dyadic half containing a given point can be chosen with a binary digit
\begin{lemma}\label{space_filling_choose_half_with_digit}
    Let $D = \{0,1\}$.
    Suppose $a\in\reals$ and $b\in\reals$ and $a \leq b$.
    Let $m = (a + b) / (1 + 1)$.
    Let $J = \intervalclosed{a}{b}$.
    Let $J_0 = \intervalclosed{a}{m}$.
    Let $J_1 = \intervalclosed{m}{b}$.
    Assume $y\in J$.
    Then there exist $d,K$ such that
        $d\in D$ and
        $y\in K$ and
        $K\subseteq J$ and
        if $d = 0$ then $K = J_0$ and
        if $d = 1$ then $K = J_1$.
\end{lemma}
\begin{proof}
    $y\in J_0$ or $y\in J_1$ by \cref{space_filling_point_in_some_half}.
    \begin{byCase}
    \caseOf{$y\in J_0$.}
        Take $d,K$ such that $d = 0$ and $K = J_0$.
        We have $d\in D$ and $y\in K$ and $K\subseteq J$
            by \cref{singleton_intro,upair_intro_left,subseteq,space_filling_left_half_subset_parent}.
        Show that if $d = 0$ then $K = J_0$ and if $d = 1$ then $K = J_1$.
        \begin{subproof}
            Trivial.
        \end{subproof}
    \caseOf{$y\in J_1$.}
        Take $d,K$ such that $d = 1$ and $K = J_1$.
        We have $d\in D$ and $y\in K$ and $K\subseteq J$
            by \cref{singleton_intro,upair_intro_right,subseteq,space_filling_right_half_subset_parent}.
        Show that if $d = 0$ then $K = J_0$ and if $d = 1$ then $K = J_1$.
        \begin{subproof}
            Trivial.
        \end{subproof}
    \end{byCase}
\end{proof}

% from §44 helper lemma 26: every interval in the point-family has a dyadic child in the same family
\begin{lemma}\label{space_filling_point_interval_family_child}
    Let $I = \intervalclosed{0}{1}$.
    Let $D = \{0,1\}$.
    Assume $y\in I$.
    Let $K = \{J\in\pow{I} \mid \text{there exist $a,b\in\reals$ such that
        $J = \intervalclosed{a}{b}$ and
        $a \leq b$ and
        $y\in J$}\}$.
    Suppose $J\in K$.
    Then there exist $a,b,d,H$ such that
        $a\in\reals$ and
        $b\in\reals$ and
        $a \leq b$ and
        $J = \intervalclosed{a}{b}$ and
        $d\in D$ and
        $H\in K$ and
        $H\subseteq J$ and
        if $d = 0$ then $H = \intervalclosed{a}{(a + b) / (1 + 1)}$ and
        if $d = 1$ then $H = \intervalclosed{(a + b) / (1 + 1)}{b}$.
\end{lemma}
\begin{proof}
    Take $a,b$ such that
        $a\in\reals$ and
        $b\in\reals$ and
        $J = \intervalclosed{a}{b}$ and
        $a \leq b$ and
        $y\in J$ by \cref{subseteq}.
    Let $m = (a + b) / (1 + 1)$.
    Let $J_0 = \intervalclosed{a}{m}$.
    Let $J_1 = \intervalclosed{m}{b}$.
    $y\in J_0$ or $y\in J_1$ by \cref{space_filling_point_in_some_half}.
    \begin{byCase}
    \caseOf{$y\in J_0$.}
        Take $d,H$ such that $d = 0$ and $H = J_0$.
        We have $d\in D$ and $H\subseteq J$
            by \cref{singleton_intro,upair_intro_left,subseteq,space_filling_left_half_subset_parent}.
        Hence $H\in K$ by \cref{subseteq,subseteq_transitive,powerset_elim,powerset_intro,real_add_closed,real_one_in_naturals,real_naturals_closed_succ,real_div_naturalsplus_real,space_filling_left_endpoint_leq_midpoint}.
        Show that if $d = 0$ then $H = \intervalclosed{a}{(a + b) / (1 + 1)}$ and
            if $d = 1$ then $H = \intervalclosed{(a + b) / (1 + 1)}{b}$.
        \begin{subproof}
            Trivial.
        \end{subproof}
    \caseOf{$y\in J_1$.}
        Take $d,H$ such that $d = 1$ and $H = J_1$.
        We have $d\in D$ and $H\subseteq J$
            by \cref{singleton_intro,upair_intro_right,subseteq,space_filling_right_half_subset_parent}.
        Hence $H\in K$ by \cref{subseteq,subseteq_transitive,powerset_elim,powerset_intro,real_add_closed,real_one_in_naturals,real_naturals_closed_succ,real_div_naturalsplus_real,space_filling_midpoint_leq_right_endpoint}.
        Show that if $d = 0$ then $H = \intervalclosed{a}{(a + b) / (1 + 1)}$ and
            if $d = 1$ then $H = \intervalclosed{(a + b) / (1 + 1)}{b}$.
        \begin{subproof}
            Trivial.
        \end{subproof}
    \end{byCase}
\end{proof}

% from §44 helper lemma 27: every interval in the point-family has a dyadic child given by one of its two halves
\begin{lemma}\label{space_filling_point_interval_family_child_disj}
    Let $I = \intervalclosed{0}{1}$.
    Assume $y\in I$.
    Let $K = \{J\in\pow{I} \mid \text{there exist $a,b\in\reals$ such that
        $J = \intervalclosed{a}{b}$ and
        $a \leq b$ and
        $y\in J$}\}$.
    Suppose $J\in K$.
    Then there exist $a,b,H$ such that
        $a\in\reals$ and
        $b\in\reals$ and
        $a \leq b$ and
        $J = \intervalclosed{a}{b}$ and
        $H\in K$ and
        $H\subseteq J$ and
        $H\in\{\intervalclosed{a}{(a + b) / (1 + 1)}, \intervalclosed{(a + b) / (1 + 1)}{b}\}$.
\end{lemma}
\begin{proof}
    Take $a,b$ such that
        $a\in\reals$ and
        $b\in\reals$ and
        $J = \intervalclosed{a}{b}$ and
        $a \leq b$ and
        $y\in J$ by \cref{subseteq}.
    Let $m = (a + b) / (1 + 1)$.
    Let $J_0 = \intervalclosed{a}{m}$.
    Let $J_1 = \intervalclosed{m}{b}$.
    $y\in J_0$ or $y\in J_1$ by \cref{space_filling_point_in_some_half}.
    \begin{byCase}
    \caseOf{$y\in J_0$.}
        Let $H = J_0$.
        $H\subseteq J$ by \cref{space_filling_left_half_subset_parent}.
        Hence $H\in K$ by \cref{subseteq,subseteq_transitive,powerset_elim,powerset_intro,real_add_closed,real_one_in_naturals,real_naturals_closed_succ,real_div_naturalsplus_real,space_filling_left_endpoint_leq_midpoint}.
        Therefore $H\in\{\intervalclosed{a}{(a + b) / (1 + 1)}, \intervalclosed{(a + b) / (1 + 1)}{b}\}$
            by \cref{upair_intro_left,subseteq}.
    \caseOf{$y\in J_1$.}
        Let $H = J_1$.
        $H\subseteq J$ by \cref{space_filling_right_half_subset_parent}.
        Thus $H\in K$ by \cref{subseteq,subseteq_transitive,powerset_elim,powerset_intro,real_add_closed,real_one_in_naturals,real_naturals_closed_succ,real_div_naturalsplus_real,space_filling_midpoint_leq_right_endpoint}.
        Hence $H\in\{\intervalclosed{a}{(a + b) / (1 + 1)}, \intervalclosed{(a + b) / (1 + 1)}{b}\}$
            by \cref{upair_intro_right,subseteq}.
    \end{byCase}
\end{proof}

% from §44 helper lemma 28: the unit interval belongs to the point-family
\begin{lemma}\label{space_filling_point_interval_family_root}
    Let $I = \intervalclosed{0}{1}$.
    Assume $y\in I$.
    Let $K = \{J\in\pow{I} \mid \text{there exist $a,b\in\reals$ such that
        $J = \intervalclosed{a}{b}$ and
        $a \leq b$ and
        $y\in J$}\}$.
    Then $I\in K$.
\end{lemma}
\begin{proof}
    Follows by \cref{real_add_ident,real_one_in_naturals,real_naturals_subset_reals,subseteq,real_zero_lt_one,real_lt_imp_leq,subseteq_refl,powerset_intro}.
\end{proof}

% from §44 helper lemma 29: a choice function selects a dyadic child relation pointwise
\begin{lemma}\label{space_filling_point_interval_child_relation_selector}
    Let $I = \intervalclosed{0}{1}$.
    Assume $y\in I$.
    Let $K = \{J\in\pow{I} \mid \text{there exist $a,b\in\reals$ such that
        $J = \intervalclosed{a}{b}$ and
        $a \leq b$ and
        $y\in J$}\}$.
    Let $R = \{w\in K\times K \mid \text{there exist $a,b$ such that
        $a\in\reals$ and
        $b\in\reals$ and
        $a \leq b$ and
        $\fst{w} = \intervalclosed{a}{b}$ and
        $\snd{w}\subseteq\fst{w}$ and
        $\snd{w}\in\{\intervalclosed{a}{(a + b) / (1 + 1)}, \intervalclosed{(a + b) / (1 + 1)}{b}\}$}\}$.
    Then there exists $g$ such that
        $g\in\funs{K}{K}$ and
        for all $J\in K$ we have $J\mathrel{R}\apply{g}{J}$.
\end{lemma}
\begin{proof}
    For all $J\in K$ there exists $H$ such that $J\mathrel{R} H$
        by \cref{space_filling_point_interval_family_child_disj,times_tuple_intro,fst_eq,snd_eq,pair_eq_iff}.
    Take $g$ such that
        $g$ is a function on $K$ and
        for all $J\in K$ we have $J\mathrel{R}\apply{g}{J}$
        by \cref{choice_function,relation,subseteq,times_elem_is_tuple}.
    Hence $g\in\funs{K}{K}$ by \cref{relation,subseteq,times_tuple_elim,funs_intro}.
\end{proof}

% from §44 helper lemma 30: every self-map of the point-family has an iterated interval sequence starting from the unit interval
\begin{lemma}\label{space_filling_point_interval_family_iteration}
    Let $I = \intervalclosed{0}{1}$.
    Assume $y\in I$.
    Let $K = \{J\in\pow{I} \mid \text{there exist $a,b\in\reals$ such that
        $J = \intervalclosed{a}{b}$ and
        $a \leq b$ and
        $y\in J$}\}$.
    Suppose $g\in\funs{K}{K}$.
    Then there exists $u$ such that
        $u$ is a sequence in $K$ and
        $u(1) = I$ and
        for all $n\in\naturalsPlus$ we have $u(n + 1) = \apply{g}{u(n)}$.
\end{lemma}
\begin{proof}
    Follows by \cref{space_filling_point_interval_family_root,iteration_sequence_naturalsplus}.
\end{proof}

% from §44 helper lemma 31: an iterated selector orbit is nested
\begin{lemma}\label{space_filling_point_interval_family_iteration_nested}
    Let $I = \intervalclosed{0}{1}$.
    Assume $y\in I$.
    Let $K = \{J\in\pow{I} \mid \text{there exist $a,b\in\reals$ such that
        $J = \intervalclosed{a}{b}$ and
        $a \leq b$ and
        $y\in J$}\}$.
    Let $R = \{w\in K\times K \mid \text{there exist $a,b$ such that
        $a\in\reals$ and
        $b\in\reals$ and
        $a \leq b$ and
        $\fst{w} = \intervalclosed{a}{b}$ and
        $\snd{w}\subseteq\fst{w}$ and
        $\snd{w}\in\{\intervalclosed{a}{(a + b) / (1 + 1)}, \intervalclosed{(a + b) / (1 + 1)}{b}\}$}\}$.
    Suppose $g\in\funs{K}{K}$.
    Suppose for all $J\in K$ we have $J\mathrel{R}\apply{g}{J}$.
    Suppose $u$ is a sequence in $K$.
    Suppose for all $n\in\naturalsPlus$ we have $u(n + 1) = \apply{g}{u(n)}$.
    Then for all $n\in\naturalsPlus$ we have $u(n + 1)\subseteq u(n)$.
\end{lemma}
\begin{proof}
    Follows by \cref{relation,subseteq,times_elem_is_tuple,sequence_in,funs_type_apply,fst_eq,snd_eq,pair_eq_iff}.
\end{proof}

% from §44 helper lemma 32: a successively nested sequence is monotone with respect to the natural order
\begin{lemma}\label{space_filling_nested_sequence_monotone}
    Suppose $u$ is a sequence in $\pow{X}$.
    Suppose for all $n\in\naturalsPlus$ we have $u(n + 1)\subseteq u(n)$.
    Then for all $m,n\in\naturalsPlus$ if $m \leq n$ then $u(n)\subseteq u(m)$.
\end{lemma}
\begin{proof}
    Let $P = \{n\in\naturalsPlus \mid \text{for all $m\in\initialSegment{n}$ we have $u(n)\subseteq u(m)$}\}$.
    $P\subseteq\naturalsPlus$ by \cref{subseteq}.
    For all $m\in\initialSegment{1}$ we have $u(1)\subseteq u(m)$
        by \cref{initial_segment_naturalsplus,real_naturals_subset_reals,subseteq,real_one_in_naturals,real_one_leq_naturalsplus,real_leq_antisym,subseteq_refl}.
    Therefore $1\in P$ by \cref{real_one_in_naturals,subseteq}.
    Show that for all $n\in P$ we have $n + 1\in P$.
    \begin{subproof}
        Fix $n\in P$.
        Show that for all $m\in\initialSegment{n + 1}$ we have $u(n + 1)\subseteq u(m)$.
        \begin{subproof}
            Fix $m\in\initialSegment{n + 1}$.
            $u(n + 1)\subseteq u(n)$ by assumption.
            \begin{byCase}
            \caseOf{$m = n + 1$.}
                We have $u(n + 1)\subseteq u(m)$ by \cref{subseteq_refl}.
            \caseOf{$m \neq n + 1$.}
                We have $u(n + 1)\subseteq u(m)$
                    by \cref{subseteq,initial_segment_naturalsplus,real_naturals_leq_succ_elim,subseteq_transitive}.
            \end{byCase}
        \end{subproof}
        Thus $n + 1\in P$ by \cref{real_naturals_closed_succ,subseteq}.
    \end{subproof}
    Hence for all $n\in\naturalsPlus$ we have $n\in P$ by \cref{real_naturals_induction}.
    Fix $m$.
    Fix $n$ such that $m\in\naturalsPlus$ and $n\in\naturalsPlus$.
    Assume $m \leq n$.
    Therefore $u(n)\subseteq u(m)$ by \cref{subseteq,initial_segment_naturalsplus}.
\end{proof}

% from §44 helper lemma 32a: three is a positive natural
\begin{lemma}\label{space_filling_three_in_naturalsplus}
    $1 + 1 + 1\in\naturalsPlus$.
\end{lemma}
\begin{proof}
    Follows by \cref{real_one_in_naturals,real_naturals_closed_succ}.
\end{proof}

% from §44 helper lemma 33: tripling and dividing by three gives the original real
\begin{lemma}\label{space_filling_triple_third_real}
    Suppose $a\in\reals$.
    Then $(a + a + a) / (1 + 1 + 1) = a$.
\end{lemma}
\begin{proof}
    We have $1 + 1 + 1\in\reals$ and $0 < 1 + 1 + 1$ and $1 + 1 + 1\neq 0$
        by \cref{space_filling_three_in_naturalsplus,real_naturals_subset_reals,subseteq,naturalsplus_positive_real,real_pos_nonzero}.
    Thus $\inv{(1 + 1 + 1)}\in\reals$ by \cref{real_mul_inverse}.
    We have $1\in\reals$ and $1 + 1\in\reals$
        by \cref{real_one_in_naturals,real_naturals_closed_succ,real_naturals_subset_reals,subseteq}.
    Hence $((1 + 1) + 1) \rmul a = (a + a) + a$
        by \cref{real_right_distrib,real_add_eq_subst,real_mul_ident}.
    Therefore $(1 + 1 + 1) \rmul a = (a + a) + a$ by \cref{real_add_assoc}.
    Thus $(1 + 1 + 1) \rmul a = a + a + a$ by \cref{real_add_assoc}.
    We have $(a + a + a) / (1 + 1 + 1) = (((1 + 1 + 1) \rmul a) \rmul \inv{(1 + 1 + 1)})$
        by \cref{real_div,real_mul_eq_subst}.
    Hence $(a + a + a) / (1 + 1 + 1) = a \rmul ((1 + 1 + 1) \rmul \inv{(1 + 1 + 1)})$
        by \cref{real_mul_comm,real_mul_eq_subst,real_mul_assoc}.
    Therefore $(a + a + a) / (1 + 1 + 1) = a \rmul 1$
        by \cref{real_mul_inverse,real_mul_comm,real_mul_eq_subst}.
    Thus $(a + a + a) / (1 + 1 + 1) = a$
        by \cref{real_mul_comm,real_mul_eq_subst,real_mul_ident}.
\end{proof}

% from §44 helper lemma 34: the left retained third lies in the parent interval
\begin{lemma}\label{space_filling_left_third_subset_parent}
    Suppose $a\in\reals$ and $b\in\reals$ and $a \leq b$.
    Let $J = \intervalclosed{a}{(a + a + b) / (1 + 1 + 1)}$.
    Then $J\subseteq \intervalclosed{a}{b}$.
\end{lemma}
\begin{proof}
    For all $x$ such that $x\in J$ we have $x\in\intervalclosed{a}{b}$
        by \cref{intervalclosed,subseteq,real_leq_add,real_add_comm,real_leq_subst_right_local,real_add_closed,real_leq_trans,real_add_assoc,space_filling_three_in_naturalsplus,real_div_naturalsplus_real,real_div_leq_same_denom,space_filling_triple_third_real}.
    Therefore $J\subseteq \intervalclosed{a}{b}$ by \cref{subseteq}.
\end{proof}

% from §44 helper lemma 35: the right retained third lies in the parent interval
\begin{lemma}\label{space_filling_right_third_subset_parent}
    Suppose $a\in\reals$ and $b\in\reals$ and $a \leq b$.
    Let $J = \intervalclosed{(a + b + b) / (1 + 1 + 1)}{b}$.
    Then $J\subseteq \intervalclosed{a}{b}$.
\end{lemma}
\begin{proof}
    For all $x$ such that $x\in J$ we have $(a + b + b) / (1 + 1 + 1) \leq x$ and $x \leq b$
        by \cref{intervalclosed}.
    We have $a + a \leq b + a$ by \cref{real_leq_add}.
    Hence $a + a \leq a + b$ by \cref{real_add_comm,real_leq_subst_right_local}.
    We have $a + b \leq b + b$ by \cref{real_leq_add}.
    We have $a + a\in\reals$ and $a + b\in\reals$ and $b + b\in\reals$
        by \cref{real_add_closed}.
    Thus $a + a \leq b + b$ by \cref{real_leq_trans}.
    Hence $(a + a) + a \leq (b + b) + a$ by \cref{real_leq_add}.
    Therefore $a + a + a \leq a + b + b$ by \cref{real_add_assoc,real_add_comm}.
    $a + a + a\in\reals$ and $a + b + b\in\reals$ by \cref{real_add_closed}.
    We have $(a + a + a) / (1 + 1 + 1)\in\reals$ and $(a + b + b) / (1 + 1 + 1)\in\reals$
        by \cref{space_filling_three_in_naturalsplus,real_div_naturalsplus_real}.
    Thus $(a + a + a) / (1 + 1 + 1) \leq (a + b + b) / (1 + 1 + 1)$
        by \cref{space_filling_three_in_naturalsplus,real_div_leq_same_denom}.
    Therefore $a \leq (a + b + b) / (1 + 1 + 1)$
        by \cref{space_filling_triple_third_real,real_leq_subst_left_local}.
    Thus for all $x$ such that $x\in J$ we have $x\in\intervalclosed{a}{b}$
        by \cref{intervalclosed,real_leq_trans}.
    Hence $J\subseteq \intervalclosed{a}{b}$ by \cref{subseteq}.
\end{proof}

% from §44 helper lemma 36: the two retained-third endpoints are strictly ordered
\begin{lemma}\label{space_filling_left_third_lt_right_third}
    Suppose $a\in\reals$ and $b\in\reals$ and $a < b$.
    Then $(a + a + b) / (1 + 1 + 1) < (a + b + b) / (1 + 1 + 1)$.
\end{lemma}
\begin{proof}
    $a + (a + b) < b + (a + b)$ by \cref{real_add_closed,real_lt_add}.
    Therefore $a + a + b < a + b + b$
        by \cref{real_add_assoc,real_add_comm,real_add_eq_subst,real_lt_subst_right}.
    $a + a + b\in\reals$ and $a + b + b\in\reals$ by \cref{real_add_closed}.
    We have $\inv{(1 + 1 + 1)}\in\reals$ and $0 < \inv{(1 + 1 + 1)}$
        by \cref{space_filling_three_in_naturalsplus,real_naturals_subset_reals,subseteq,naturalsplus_positive_real,real_pos_nonzero,real_mul_inverse,real_inv_pos}.
    Hence $(a + a + b) \rmul \inv{(1 + 1 + 1)} < (a + b + b) \rmul \inv{(1 + 1 + 1)}$
        by \cref{real_lt_mul_pos}.
    Therefore $(a + a + b) / (1 + 1 + 1) < (a + b + b) / (1 + 1 + 1)$
        by \cref{real_div,real_lt_subst_left,real_lt_subst_right}.
\end{proof}

% from §44 helper lemma 37: the two retained thirds are disjoint
\begin{lemma}\label{space_filling_retained_thirds_disjoint}
    Suppose $a\in\reals$ and $b\in\reals$ and $a < b$.
    Let $J_0 = \intervalclosed{a}{(a + a + b) / (1 + 1 + 1)}$.
    Let $J_1 = \intervalclosed{(a + b + b) / (1 + 1 + 1)}{b}$.
    Then $J_0$ is disjoint from $J_1$.
\end{lemma}
\begin{proof}
    Suppose not.
    Take $x$ such that $x\in J_0, J_1$.
    We have $x \leq (a + a + b) / (1 + 1 + 1)$ and
        $(a + b + b) / (1 + 1 + 1) \leq x$
        by \cref{intervalclosed}.
    Hence $x\in\reals$ by \cref{intervalclosed,subseteq}.
    $a + a + b\in\reals$ and $a + b + b\in\reals$ by \cref{real_add_closed}.
    Thus $(a + a + b) / (1 + 1 + 1)\in\reals$ and
        $(a + b + b) / (1 + 1 + 1)\in\reals$
        by \cref{space_filling_three_in_naturalsplus,real_div_naturalsplus_real}.
    Hence $(a + b + b) / (1 + 1 + 1) \leq (a + a + b) / (1 + 1 + 1)$
        by \cref{real_leq_trans}.
    Therefore $(a + b + b) / (1 + 1 + 1) < (a + b + b) / (1 + 1 + 1)$
        by \cref{real_leq_lt_trans,space_filling_left_third_lt_right_third}.
    Contradiction by \cref{real_lt_irreflexive}.
\end{proof}

% from §44 helper lemma 38: a retained third can be chosen from a binary digit
\begin{lemma}\label{space_filling_choose_retained_third}
    Let $D = \{0,1\}$.
    Suppose $a\in\reals$ and $b\in\reals$ and $a \leq b$.
    Let $J = \intervalclosed{a}{b}$.
    Assume $d\in D$.
    Then there exists $H$ such that
        $H\subseteq J$ and
        if $d = 0$ then $H = \intervalclosed{a}{(a + a + b) / (1 + 1 + 1)}$ and
        if $d = 1$ then $H = \intervalclosed{(a + b + b) / (1 + 1 + 1)}{b}$.
\end{lemma}
\begin{proof}
    \begin{byCase}
    \caseOf{$d = 0$.}
        Let $H = \intervalclosed{a}{(a + a + b) / (1 + 1 + 1)}$.
        We have $H\subseteq J$ by \cref{space_filling_left_third_subset_parent}.
        Show that if $d = 0$ then $H = \intervalclosed{a}{(a + a + b) / (1 + 1 + 1)}$ and
            if $d = 1$ then $H = \intervalclosed{(a + b + b) / (1 + 1 + 1)}{b}$.
        \begin{subproof}
            Trivial.
        \end{subproof}
    \caseOf{$d = 1$.}
        Let $H = \intervalclosed{(a + b + b) / (1 + 1 + 1)}{b}$.
        We have $H\subseteq J$ by \cref{space_filling_right_third_subset_parent}.
        Show that if $d = 0$ then $H = \intervalclosed{a}{(a + a + b) / (1 + 1 + 1)}$ and
            if $d = 1$ then $H = \intervalclosed{(a + b + b) / (1 + 1 + 1)}{b}$.
        \begin{subproof}
            Trivial.
        \end{subproof}
    \end{byCase}
\end{proof}

% from §44 helper lemma 39: the left retained-third endpoint is above the left endpoint
\begin{lemma}\label{space_filling_left_endpoint_leq_left_third}
    Suppose $a\in\reals$ and $b\in\reals$ and $a \leq b$.
    Then $a \leq (a + a + b) / (1 + 1 + 1)$.
\end{lemma}
\begin{proof}
    Follows by \cref{real_leq_add,real_add_closed,real_add_assoc,real_add_comm,real_add_eq_subst,real_leq_subst_right_local,space_filling_three_in_naturalsplus,real_div_leq_same_denom,space_filling_triple_third_real,real_leq_subst_left_local}.
\end{proof}

% from §44 helper lemma 40: the right retained-third endpoint is below the right endpoint
\begin{lemma}\label{space_filling_right_third_leq_right_endpoint}
    Suppose $a\in\reals$ and $b\in\reals$ and $a \leq b$.
    Then $(a + b + b) / (1 + 1 + 1) \leq b$.
\end{lemma}
\begin{proof}
    Follows by \cref{real_add_closed,real_leq_add,real_add_assoc,real_leq_subst_right_local,space_filling_three_in_naturalsplus,real_div_leq_same_denom,space_filling_triple_third_real}.
\end{proof}

% from §44 helper lemma 41: a dyadic half can be chosen from a binary digit
\begin{lemma}\label{space_filling_choose_half_from_digit}
    Let $D = \{0,1\}$.
    Suppose $a\in\reals$ and $b\in\reals$ and $a \leq b$.
    Let $J = \intervalclosed{a}{b}$.
    Assume $d\in D$.
    Then there exists $H$ such that
        $H\subseteq J$ and
        if $d = 0$ then $H = \intervalclosed{a}{(a + b) / (1 + 1)}$ and
        if $d = 1$ then $H = \intervalclosed{(a + b) / (1 + 1)}{b}$.
\end{lemma}
\begin{proof}
    \begin{byCase}
    \caseOf{$d = 0$.}
        Let $H = \intervalclosed{a}{(a + b) / (1 + 1)}$.
        We have $H\subseteq J$ by \cref{space_filling_left_half_subset_parent}.
        Show that if $d = 0$ then $H = \intervalclosed{a}{(a + b) / (1 + 1)}$ and
            if $d = 1$ then $H = \intervalclosed{(a + b) / (1 + 1)}{b}$.
        \begin{subproof}
            Trivial.
        \end{subproof}
    \caseOf{$d = 1$.}
        Let $H = \intervalclosed{(a + b) / (1 + 1)}{b}$.
        We have $H\subseteq J$ by \cref{space_filling_right_half_subset_parent}.
        Show that if $d = 0$ then $H = \intervalclosed{a}{(a + b) / (1 + 1)}$ and
            if $d = 1$ then $H = \intervalclosed{(a + b) / (1 + 1)}{b}$.
        \begin{subproof}
            Trivial.
        \end{subproof}
    \end{byCase}
\end{proof}

% from §44 helper lemma 42: the unit interval belongs to the closed-interval family
\begin{lemma}\label{space_filling_closed_interval_family_root}
    Let $I = \intervalclosed{0}{1}$.
    Let $K = \{J\in\pow{I} \mid \text{there exist $a,b\in\reals$ such that
        $J = \intervalclosed{a}{b}$ and
        $a \leq b$}\}$.
    Then $I\in K$.
\end{lemma}
\begin{proof}
    Follows by \cref{real_add_ident,real_one_in_naturals,real_naturals_subset_reals,subseteq,real_zero_lt_one,real_lt_imp_leq,subseteq_refl,powerset_intro}.
\end{proof}

% from §44 helper lemma 43: a binary sequence selects dyadic child intervals statewise
\begin{lemma}\label{space_filling_dyadic_child_state_selector}
    Let $I = \intervalclosed{0}{1}$.
    Let $D = \{0,1\}$.
    Let $E = \{(n,D) \mid n\in\naturalsPlus\}$.
    Assume $x\in\productFamily{E}$.
    Let $K = \{J\in\pow{I} \mid \text{there exist $a,b\in\reals$ such that
        $J = \intervalclosed{a}{b}$ and
        $a \leq b$}\}$.
    Let $G = \naturalsPlus\times K$.
    Let $R = \{p\in G\times G \mid \text{there exist $n,J,H,a,b,d$ such that
        $n\in\naturalsPlus$ and
        $\fst{p} = (n,J)$ and
        $\snd{p} = (n + 1,H)$ and
        $d = x(n)$ and
        $a\in\reals$ and
        $b\in\reals$ and
        $a \leq b$ and
        $J = \intervalclosed{a}{b}$ and
        $H\subseteq J$ and
        $(d,H)\in\{(0,\intervalclosed{a}{(a + b) / (1 + 1)}),
            (1,\intervalclosed{(a + b) / (1 + 1)}{b})\}$}\}$.
    Then there exists $g$ such that
        $g\in\funs{G}{G}$ and
        for all $z\in G$ we have $z\mathrel{R}\apply{g}{z}$.
\end{lemma}
\begin{proof}
    For all $n,u,v$ such that $(n,u)\in E$ and $(n,v)\in E$ we have $u = v$ by \cref{pair_eq_iff}.
    We have $E$ is a function on $\naturalsPlus$
        by \cref{relation,rightunique,dom_intro,dom_iff,pair_eq_iff,subseteq,subseteq_antisymmetric}.
    Therefore $\dom{E} = \naturalsPlus$ by \cref{funs}.
    For all $n\in\naturalsPlus$ we have $E(n) = D$ by \cref{function_apply_intro}.
    Show that for all $z\in G$ there exists $w$ such that $z\mathrel{R} w$.
    \begin{subproof}
        Fix $z\in G$.
        Take $n,J$ such that $z = (n,J)$ and $n\in\naturalsPlus$ and $J\in K$
            by \cref{subseteq,times_elem_is_tuple}.
        Hence $x(n)\in D$ by \cref{indexed_product,subseteq}.
        Take $a,b$ such that
            $a\in\reals$ and
            $b\in\reals$ and
            $J = \intervalclosed{a}{b}$ and
            $a \leq b$
            by \cref{subseteq}.
        $x(n) = 0$ or $x(n) = 1$ by \cref{upair_iff}.
        \begin{byCase}
        \caseOf{$x(n) = 0$.}
            Let $H = \intervalclosed{a}{(a + b) / (1 + 1)}$.
            We have $H\subseteq J$ by \cref{space_filling_left_half_subset_parent}.
            Hence $(x(n),H)\in\{(0,\intervalclosed{a}{(a + b) / (1 + 1)}),
                (1,\intervalclosed{(a + b) / (1 + 1)}{b})\}$ by \cref{upair_intro_left,pair_eq_iff,subseteq}.
            $H\in K$ by \cref{subseteq,subseteq_transitive,powerset_elim,powerset_intro,real_add_closed,real_one_in_naturals,real_naturals_closed_succ,real_div_naturalsplus_real,space_filling_left_endpoint_leq_midpoint}.
            Hence $(n + 1,H)\in G$ by \cref{real_naturals_closed_succ,times_tuple_intro}.
            Take $w$ such that $w = (n + 1,H)$.
            Therefore $z\mathrel{R} w$ by \cref{relation,subseteq,times_tuple_intro,fst_eq,snd_eq,pair_eq_iff}.
        \caseOf{$x(n) = 1$.}
            Let $H = \intervalclosed{(a + b) / (1 + 1)}{b}$.
            We have $H\subseteq J$ by \cref{space_filling_right_half_subset_parent}.
            Hence $(x(n),H)\in\{(0,\intervalclosed{a}{(a + b) / (1 + 1)}),
                (1,\intervalclosed{(a + b) / (1 + 1)}{b})\}$ by \cref{upair_intro_right,pair_eq_iff,subseteq}.
            $H\in K$ by \cref{subseteq,subseteq_transitive,powerset_elim,powerset_intro,real_add_closed,real_one_in_naturals,real_naturals_closed_succ,real_div_naturalsplus_real,space_filling_midpoint_leq_right_endpoint}.
            Hence $(n + 1,H)\in G$ by \cref{real_naturals_closed_succ,times_tuple_intro}.
            Take $w$ such that $w = (n + 1,H)$.
            Therefore $z\mathrel{R} w$ by \cref{relation,subseteq,times_tuple_intro,fst_eq,snd_eq,pair_eq_iff}.
        \end{byCase}
    \end{subproof}
    Take $g$ such that
        $g$ is a function on $G$ and
        for all $z\in G$ we have $z\mathrel{R}\apply{g}{z}$
        by \cref{choice_function,relation,subseteq,times_elem_is_tuple}.
    Thus $g\in\funs{G}{G}$ by \cref{funs_intro,relation,subseteq,times_tuple_elim}.
\end{proof}

% from §44 helper lemma 44: a dyadic child selector has an iterated state orbit
\begin{lemma}\label{space_filling_dyadic_child_state_iteration}
    Let $I = \intervalclosed{0}{1}$.
    Let $K = \{J\in\pow{I} \mid \text{there exist $a,b\in\reals$ such that
        $J = \intervalclosed{a}{b}$ and
        $a \leq b$}\}$.
    Let $G = \naturalsPlus\times K$.
    Suppose $g\in\funs{G}{G}$.
    Then there exists $u$ such that
        $u$ is a sequence in $G$ and
        $u(1) = (1,I)$ and
        for all $n\in\naturalsPlus$ we have $u(n + 1) = \apply{g}{u(n)}$.
\end{lemma}
\begin{proof}
    Follows by \cref{space_filling_closed_interval_family_root,real_one_in_naturals,times_tuple_intro,iteration_sequence_naturalsplus}.
\end{proof}

% from §44 helper lemma 45: the first coordinate of a dyadic state orbit is the stage number
\begin{lemma}\label{space_filling_dyadic_child_state_orbit_index}
    Let $I = \intervalclosed{0}{1}$.
    Let $D = \{0,1\}$.
    Let $E = \{(n,D) \mid n\in\naturalsPlus\}$.
    Assume $x\in\productFamily{E}$.
    Let $K = \{J\in\pow{I} \mid \text{there exist $a,b\in\reals$ such that
        $J = \intervalclosed{a}{b}$ and
        $a \leq b$}\}$.
    Let $G = \naturalsPlus\times K$.
    Let $R = \{p\in G\times G \mid \text{there exist $n,J,H,a,b,d$ such that
        $n\in\naturalsPlus$ and
        $\fst{p} = (n,J)$ and
        $\snd{p} = (n + 1,H)$ and
        $d = x(n)$ and
        $a\in\reals$ and
        $b\in\reals$ and
        $a \leq b$ and
        $J = \intervalclosed{a}{b}$ and
        $H\subseteq J$ and
        $(d,H)\in\{(0,\intervalclosed{a}{(a + b) / (1 + 1)}),
            (1,\intervalclosed{(a + b) / (1 + 1)}{b})\}$}\}$.
    Suppose $g\in\funs{G}{G}$.
    Suppose for all $z\in G$ we have $z\mathrel{R}\apply{g}{z}$.
    Suppose $u$ is a sequence in $G$.
    Suppose $u(1) = (1,I)$.
    Suppose for all $n\in\naturalsPlus$ we have $u(n + 1) = \apply{g}{u(n)}$.
    Then for all $n\in\naturalsPlus$ we have $\fst{u(n)} = n$.
\end{lemma}
\begin{proof}
    Let $A_0 = \{n\in\naturalsPlus \mid \fst{u(n)} = n\}$.
    $A_0\subseteq\naturalsPlus$ by \cref{subseteq}.
    We have $1\in A_0$ by \cref{real_one_in_naturals,pair_eq_iff,fst_eq,subseteq}.
    Show that for all $n\in A_0$ we have $n + 1\in A_0$.
    \begin{subproof}
        Fix $n\in A_0$.
        We have $n\in\naturalsPlus$ and $\fst{u(n)} = n$ by \cref{subseteq}.
        Hence $(u(n),u(n + 1))\in R$ by \cref{sequence_in,funs_type_apply,relation,subseteq}.
        Take $k,J,H,a,b,d$ such that
            $(u(n),u(n + 1))\in G\times G$ and
            $k\in\naturalsPlus$ and
            $\fst{(u(n),u(n + 1))} = (k,J)$ and
            $\snd{(u(n),u(n + 1))} = (k + 1,H)$ and
            $a\in\reals$ and
            $b\in\reals$ and
            $a \leq b$ and
            $J = \intervalclosed{a}{b}$ and
            $H\subseteq J$ and
            $(d,H)\in\{(0,\intervalclosed{a}{(a + b) / (1 + 1)}),
                (1,\intervalclosed{(a + b) / (1 + 1)}{b})\}$
            by \cref{relation,subseteq}.
        Thus $k = n$ and $\fst{u(n + 1)} = n + 1$
            by \cref{fst_eq,snd_eq,pair_eq_iff,subseteq,real_add_eq_subst}.
        Hence $n + 1\in A_0$ by \cref{real_naturals_closed_succ,subseteq}.
    \end{subproof}
    Therefore for all $n\in\naturalsPlus$ we have $n\in A_0$ by \cref{real_naturals_induction}.
    Hence for all $n\in\naturalsPlus$ we have $\fst{u(n)} = n$ by \cref{subseteq}.
\end{proof}

% from §44 helper lemma 46: the second coordinate of a dyadic state orbit is nested
\begin{lemma}\label{space_filling_dyadic_child_state_orbit_nested}
    Let $I = \intervalclosed{0}{1}$.
    Let $D = \{0,1\}$.
    Let $E = \{(n,D) \mid n\in\naturalsPlus\}$.
    Assume $x\in\productFamily{E}$.
    Let $K = \{J\in\pow{I} \mid \text{there exist $a,b\in\reals$ such that
        $J = \intervalclosed{a}{b}$ and
        $a \leq b$}\}$.
    Let $G = \naturalsPlus\times K$.
    Let $R = \{p\in G\times G \mid \text{there exist $n,J,H,a,b,d$ such that
        $n\in\naturalsPlus$ and
        $\fst{p} = (n,J)$ and
        $\snd{p} = (n + 1,H)$ and
        $d = x(n)$ and
        $a\in\reals$ and
        $b\in\reals$ and
        $a \leq b$ and
        $J = \intervalclosed{a}{b}$ and
        $H\subseteq J$ and
        $(d,H)\in\{(0,\intervalclosed{a}{(a + b) / (1 + 1)}),
            (1,\intervalclosed{(a + b) / (1 + 1)}{b})\}$}\}$.
    Suppose $g\in\funs{G}{G}$.
    Suppose for all $z\in G$ we have $z\mathrel{R}\apply{g}{z}$.
    Suppose $u$ is a sequence in $G$.
    Suppose for all $n\in\naturalsPlus$ we have $u(n + 1) = \apply{g}{u(n)}$.
    Then for all $n\in\naturalsPlus$ we have $\snd{u(n + 1)}\subseteq \snd{u(n)}$.
\end{lemma}
\begin{proof}
    Follows by \cref{sequence_in,funs_type_apply,relation,subseteq,fst_eq,snd_eq,pair_eq_iff}.
\end{proof}

% from §44 helper lemma 47: the interval coordinates of a dyadic state orbit are monotone
\begin{lemma}\label{space_filling_dyadic_child_state_orbit_monotone}
    Let $I = \intervalclosed{0}{1}$.
    Let $D = \{0,1\}$.
    Let $E = \{(n,D) \mid n\in\naturalsPlus\}$.
    Assume $x\in\productFamily{E}$.
    Let $K = \{J\in\pow{I} \mid \text{there exist $a,b\in\reals$ such that
        $J = \intervalclosed{a}{b}$ and
        $a \leq b$}\}$.
    Let $G = \naturalsPlus\times K$.
    Let $R = \{p\in G\times G \mid \text{there exist $n,J,H,a,b,d$ such that
        $n\in\naturalsPlus$ and
        $\fst{p} = (n,J)$ and
        $\snd{p} = (n + 1,H)$ and
        $d = x(n)$ and
        $a\in\reals$ and
        $b\in\reals$ and
        $a \leq b$ and
        $J = \intervalclosed{a}{b}$ and
        $H\subseteq J$ and
        $(d,H)\in\{(0,\intervalclosed{a}{(a + b) / (1 + 1)}),
            (1,\intervalclosed{(a + b) / (1 + 1)}{b})\}$}\}$.
    Suppose $g\in\funs{G}{G}$.
    Suppose for all $z\in G$ we have $z\mathrel{R}\apply{g}{z}$.
    Suppose $u$ is a sequence in $G$.
    Suppose for all $n\in\naturalsPlus$ we have $u(n + 1) = \apply{g}{u(n)}$.
    Then for all $m,n\in\naturalsPlus$ if $m \leq n$ then $\snd{u(n)}\subseteq \snd{u(m)}$.
\end{lemma}
\begin{proof}
    Let $P = \{n\in\naturalsPlus \mid \text{for all $m\in\initialSegment{n}$ we have $\snd{u(n)}\subseteq \snd{u(m)}$}\}$.
    $P\subseteq\naturalsPlus$ by \cref{subseteq}.
    For all $m\in\initialSegment{1}$ we have $\snd{u(1)}\subseteq \snd{u(m)}$.
    Therefore $1\in P$ by \cref{real_one_in_naturals,subseteq}.
    Show that for all $n\in P$ we have $n + 1\in P$.
    \begin{subproof}
        Fix $n\in P$.
        Show that for all $m\in\initialSegment{n + 1}$ we have $\snd{u(n + 1)}\subseteq \snd{u(m)}$.
        \begin{subproof}
            Fix $m\in\initialSegment{n + 1}$.
            $\snd{u(n + 1)}\subseteq \snd{u(n)}$ by \cref{space_filling_dyadic_child_state_orbit_nested}.
            \begin{byCase}
            \caseOf{$m = n + 1$.}
                We have $\snd{u(n + 1)}\subseteq \snd{u(m)}$ by \cref{subseteq_refl}.
            \caseOf{$m \neq n + 1$.}
                We have $\snd{u(n + 1)}\subseteq \snd{u(m)}$
                    by \cref{subseteq,initial_segment_naturalsplus,real_naturals_leq_succ_elim,subseteq_transitive}.
            \end{byCase}
        \end{subproof}
        Thus $n + 1\in P$ by \cref{real_naturals_closed_succ,subseteq}.
    \end{subproof}
    Hence for all $n\in\naturalsPlus$ we have $n\in P$ by \cref{real_naturals_induction}.
    Fix $m$.
    Fix $n$ such that $m\in\naturalsPlus$ and $n\in\naturalsPlus$.
    Assume $m \leq n$.
    Therefore $\snd{u(n)}\subseteq \snd{u(m)}$ by \cref{subseteq,initial_segment_naturalsplus}.
\end{proof}

% from §44 helper lemma 48: every closed interval with ordered endpoints is inhabited
\begin{lemma}\label{space_filling_intervalclosed_inhabited}
    Suppose $a\in\reals$ and $b\in\reals$ and $a \leq b$.
    Then $\intervalclosed{a}{b}$ is inhabited.
\end{lemma}
\begin{proof}
    Follows by \cref{real_add_closed,real_one_in_naturals,real_naturals_closed_succ,real_div_naturalsplus_real,space_filling_left_endpoint_leq_midpoint,space_filling_midpoint_leq_right_endpoint,intervalclosed,nonempty_inhabited}.
\end{proof}

% from §44 helper lemma 49: every member of the dyadic closed-interval family is closed in the unit interval
\begin{lemma}\label{space_filling_closed_interval_family_member_closed}
    Let $I = \intervalclosed{0}{1}$.
    Let $T_I = \subspaceTopology{\standardTopologyR}{I}$.
    Let $K = \{J\in\pow{I} \mid \text{there exist $a,b\in\reals$ such that
        $J = \intervalclosed{a}{b}$ and
        $a \leq b$}\}$.
    Assume $J\in K$.
    Then $J$ is closed in $I$ with respect to $T_I$.
\end{lemma}
\begin{proof}
    Take $a,b$ such that
        $a\in\reals$ and
        $b\in\reals$ and
        $J = \intervalclosed{a}{b}$ and
        $a \leq b$
        by \cref{subseteq}.
    Therefore $J$ is closed in $I$ with respect to $T_I$
        by \cref{intervalclosed,subseteq,standard_basis_is_basis,basis_topology_is_topology,standard_topology_reals,subspace_of,real_intervalclosed_closed_standard,closed_in_subspace_iff,powerset_elim,inter_absorb_supseteq_right}.
\end{proof}

% from §44 helper lemma 50: the interval members of a dyadic state orbit have the finite intersection property
\begin{lemma}\label{space_filling_dyadic_child_state_orbit_fip}
    Let $I = \intervalclosed{0}{1}$.
    Let $D = \{0,1\}$.
    Let $E = \{(n,D) \mid n\in\naturalsPlus\}$.
    Assume $x\in\productFamily{E}$.
    Let $K = \{J\in\pow{I} \mid \text{there exist $a,b\in\reals$ such that
        $J = \intervalclosed{a}{b}$ and
        $a \leq b$}\}$.
    Let $G = \naturalsPlus\times K$.
    Let $R = \{p\in G\times G \mid \text{there exist $n,J,H,a,b,d$ such that
        $n\in\naturalsPlus$ and
        $\fst{p} = (n,J)$ and
        $\snd{p} = (n + 1,H)$ and
        $d = x(n)$ and
        $a\in\reals$ and
        $b\in\reals$ and
        $a \leq b$ and
        $J = \intervalclosed{a}{b}$ and
        $H\subseteq J$ and
        $(d,H)\in\{(0,\intervalclosed{a}{(a + b) / (1 + 1)}),
            (1,\intervalclosed{(a + b) / (1 + 1)}{b})\}$}\}$.
    Suppose $g\in\funs{G}{G}$.
    Suppose for all $z\in G$ we have $z\mathrel{R}\apply{g}{z}$.
    Suppose $u$ is a sequence in $G$.
    Suppose for all $n\in\naturalsPlus$ we have $u(n + 1) = \apply{g}{u(n)}$.
    Let $v = \{(n,\snd{u(n)}) \mid n\in\naturalsPlus\}$.
    Let $C = \img{v}{\naturalsPlus}$.
    Then $C$ has the finite intersection property.
\end{lemma}
\begin{proof}
    $u\in\funs{\naturalsPlus}{G}$ by \cref{sequence_in}.
    We have $\dom{v} = \naturalsPlus$ by \cref{relation,rightunique,pair_eq_iff,dom_intro,dom_iff,subseteq,subseteq_antisymmetric}.
    For all $n\in\dom{v}$ we have $v(n) = \snd{u(n)}$ by \cref{function_apply_intro,relation,rightunique,pair_eq_iff}.
    Thus $C$ is inhabited by \cref{real_one_in_naturals,function_apply_elim,subseteq,img_iff,nonempty_inhabited}.
    Show that for all $F$ such that $F\subseteq C$ and $F$ is finite and inhabited
    we have $\inters{F}\neq\emptyset$.
    \begin{subproof}
        Fix $F$ such that $F\subseteq C$ and $F$ is finite and inhabited.
        Let $Q = \{w\in F\times\naturalsPlus \mid \text{there exists $n\in\naturalsPlus$ such that
            $\snd{w} = n$ and
            $\fst{w} = v(n)$}\}$.
        For all $J\in F$ there exists $n$ such that $J\mathrel{Q} n$
            by \cref{img_iff,fst_eq,snd_eq,pair_eq_iff,relation,subseteq,times_tuple_intro}.
        Take $k$ such that
            $k$ is a function on $F$ and
            for all $J\in F$ we have $J\mathrel{Q}\apply{k}{J}$
            by \cref{choice_function,relation,subseteq,times_elem_is_tuple}.
        Let $M = \img{k}{F}$.
        $M$ is finite by \cref{finite_image_function}.
        Take $J_0$ such that $J_0\in F$ by \cref{nonempty_inhabited}.
        Therefore $M$ is inhabited by \cref{img_iff,function_apply_elim,nonempty_inhabited}.
        For all $m\in M$ we have $m\in\naturalsPlus$
            by \cref{img_iff,function_apply_intro,pair_eq_iff,relation,subseteq,times_tuple_elim}.
        Thus $M\subseteq\reals$ by \cref{real_naturals_subset_reals,subseteq}.
        Take $N\in M$ such that $N$ is a maximum of $M$ by \cref{inhabited_nonempty,finite_real_maximum}.
        Therefore $N\in\naturalsPlus$ by \cref{subseteq}.
        Thus $\snd{u(N)}\in K$
            by \cref{funs_type_apply,subseteq,times_elem_is_tuple,snd_eq,pair_eq_iff}.
        Take $a,b$ such that
            $a\in\reals$ and
            $b\in\reals$ and
            $\snd{u(N)} = \intervalclosed{a}{b}$ and
            $a \leq b$
            by \cref{subseteq}.
        Take $y$ such that $y = (a + b) / (1 + 1)$.
        Therefore $y\in\snd{u(N)}$
            by \cref{intervalclosed,subseteq,real_add_closed,real_one_in_naturals,real_naturals_closed_succ,real_div_naturalsplus_real,space_filling_left_endpoint_leq_midpoint,space_filling_midpoint_leq_right_endpoint}.
        Show that for all $J\in F$ we have $y\in J$.
        \begin{subproof}
            Fix $J\in F$.
            We have $(J,k(J))\in Q$ by \cref{relation}.
            Take $n$ such that
                $(J,k(J))\in F\times\naturalsPlus$ and
                $n\in\naturalsPlus$ and
                $\snd{(J,k(J))} = n$ and
                $\fst{(J,k(J))} = v(n)$
                by \cref{subseteq}.
            Therefore $J = v(k(J))$ and $k(J)\in\naturalsPlus$
                by \cref{fst_eq,snd_eq,pair_eq_iff,times_tuple_elim,real_lt_subst_right}.
            Thus $k(J)\in M$ by \cref{img_iff,function_apply_elim}.
            Hence $k(J)\leq N$ by \cref{real_maximum}.
            Since $k(J)\in\naturalsPlus$ and $N\in\naturalsPlus$ and $k(J)\leq N$,
            we have $\snd{u(N)}\subseteq \snd{u(k(J))}$
                by \cref{space_filling_dyadic_child_state_orbit_monotone}.
            Therefore $y\in J$ by \cref{function_apply_intro,subseteq,real_lt_subst_right}.
        \end{subproof}
        Thus $y\in\inters{F}$ by \cref{inters_intro}.
        Therefore $\inters{F}\neq\emptyset$ by \cref{emptyset}.
    \end{subproof}
    Therefore $C$ has the finite intersection property by \cref{finite_intersection_property}.
\end{proof}

% generic helper: a monotone closed-interval state orbit has the finite intersection property
\begin{lemma}\label{space_filling_closed_interval_state_orbit_fip_from_monotone}
    Let $I = \intervalclosed{0}{1}$.
    Let $K = \{J\in\pow{I} \mid \text{there exist $a,b\in\reals$ such that
        $J = \intervalclosed{a}{b}$ and
        $a \leq b$}\}$.
    Let $G = \naturalsPlus\times K$.
    Suppose $u$ is a sequence in $G$.
    Suppose for all $m,n\in\naturalsPlus$ if $m \leq n$ then $\snd{u(n)}\subseteq \snd{u(m)}$.
    Let $v = \{(n,\snd{u(n)}) \mid n\in\naturalsPlus\}$.
    Let $C = \img{v}{\naturalsPlus}$.
    Then $C$ has the finite intersection property.
\end{lemma}
\begin{proof}
    $u\in\funs{\naturalsPlus}{G}$ by \cref{sequence_in}.
    For all $n,J_1,J_2$ such that $(n,J_1)\in v$ and $(n,J_2)\in v$ we have $J_1 = J_2$ by \cref{pair_eq_iff}.
    Therefore $v$ is a function by \cref{relation,rightunique,pair_eq_iff}.
    Thus $\dom{v} = \naturalsPlus$ by \cref{dom_intro,dom_iff,pair_eq_iff,subseteq,subseteq_antisymmetric}.
    For all $n\in\dom{v}$ we have $v(n) = \snd{u(n)}$ by \cref{function_apply_intro}.
    Therefore $C$ is inhabited by \cref{real_one_in_naturals,function_apply_elim,subseteq,img_iff,nonempty_inhabited}.
    Show that for all $F$ such that $F\subseteq C$ and $F$ is finite and inhabited
    we have $\inters{F}\neq\emptyset$.
    \begin{subproof}
        Fix $F$ such that $F\subseteq C$ and $F$ is finite and inhabited.
        Let $Q = \{w\in F\times\naturalsPlus \mid \text{there exists $n\in\naturalsPlus$ such that
            $\snd{w} = n$ and
            $\fst{w} = v(n)$}\}$.
        For all $J\in F$ there exists $n$ such that $J\mathrel{Q} n$
            by \cref{img_iff,fst_eq,snd_eq,pair_eq_iff,relation,subseteq,times_tuple_intro}.
        Take $k$ such that
            $k$ is a function on $F$ and
            for all $J\in F$ we have $J\mathrel{Q}\apply{k}{J}$
            by \cref{choice_function,relation,subseteq,times_elem_is_tuple}.
        Let $M = \img{k}{F}$.
        $M$ is finite by \cref{finite_image_function}.
        Take $J_0$ such that $J_0\in F$ by \cref{nonempty_inhabited}.
        Thus $M$ is inhabited by \cref{img_iff,function_apply_elim,choice_function,nonempty_inhabited}.
        For all $m\in M$ we have $m\in\naturalsPlus$
            by \cref{img_iff,choice_function,function_apply_intro,pair_eq_iff,relation,subseteq,times_tuple_elim}.
        Therefore $M\subseteq\reals$ by \cref{real_naturals_subset_reals,subseteq}.
        Take $N\in M$ such that $N$ is a maximum of $M$ by \cref{inhabited_nonempty,finite_real_maximum}.
        Thus $N\in\naturalsPlus$ by \cref{subseteq}.
        Therefore $\snd{u(N)}\in K$
            by \cref{funs_type_apply,subseteq,times_elem_is_tuple,snd_eq,pair_eq_iff}.
        Take $a,b$ such that
            $a\in\reals$ and
            $b\in\reals$ and
            $\snd{u(N)} = \intervalclosed{a}{b}$ and
            $a \leq b$
            by \cref{subseteq}.
        Take $y$ such that $y = (a + b) / (1 + 1)$.
        Thus $y\in\snd{u(N)}$
            by \cref{intervalclosed,subseteq,real_add_closed,real_one_in_naturals,real_naturals_closed_succ,real_div_naturalsplus_real,space_filling_left_endpoint_leq_midpoint,space_filling_midpoint_leq_right_endpoint}.
        Show that for all $J\in F$ we have $y\in J$.
        \begin{subproof}
            Fix $J\in F$.
            We have $(J,k(J))\in Q$ by \cref{choice_function,relation}.
            Take $n$ such that
                $(J,k(J))\in F\times\naturalsPlus$ and
                $n\in\naturalsPlus$ and
                $\snd{(J,k(J))} = n$ and
                $\fst{(J,k(J))} = v(n)$
                by \cref{subseteq}.
            Thus $J = v(k(J))$ and $k(J)\in\naturalsPlus$
                by \cref{fst_eq,snd_eq,pair_eq_iff,times_tuple_elim,real_lt_subst_right}.
            Therefore $k(J)\in M$ by \cref{img_iff,function_apply_elim,choice_function}.
            Hence $k(J)\leq N$ by \cref{real_maximum}.
            Since $k(J)\in\naturalsPlus$ and $N\in\naturalsPlus$ and $k(J)\leq N$,
            we have $\snd{u(N)}\subseteq \snd{u(k(J))}$.
            Therefore $y\in J$ by \cref{function_apply_intro,subseteq,real_lt_subst_right}.
        \end{subproof}
        Thus $y\in\inters{F}$ by \cref{inters_intro}.
        Therefore $\inters{F}\neq\emptyset$ by \cref{emptyset}.
    \end{subproof}
    Thus $C$ has the finite intersection property by \cref{finite_intersection_property}.
\end{proof}

% generic helper: a closed-interval state orbit with the finite intersection property has a common point
\begin{lemma}\label{space_filling_closed_interval_state_orbit_common_point_from_fip}
    Let $I = \intervalclosed{0}{1}$.
    Let $K = \{J\in\pow{I} \mid \text{there exist $a,b\in\reals$ such that
        $J = \intervalclosed{a}{b}$ and
        $a \leq b$}\}$.
    Let $G = \naturalsPlus\times K$.
    Suppose $u$ is a sequence in $G$.
    Let $v = \{(n,\snd{u(n)}) \mid n\in\naturalsPlus\}$.
    Let $C = \img{v}{\naturalsPlus}$.
    Assume $C$ has the finite intersection property.
    Then there exists $y\in I$ such that
        for all $n\in\naturalsPlus$ we have $y\in\snd{u(n)}$.
\end{lemma}
\begin{proof}
    Let $T_I = \subspaceTopology{\standardTopologyR}{I}$.
    $u\in\funs{\naturalsPlus}{G}$ by \cref{sequence_in}.
    For all $n,J_1,J_2$ such that $(n,J_1)\in v$ and $(n,J_2)\in v$ we have $J_1 = J_2$ by \cref{pair_eq_iff}.
    Therefore $v$ is a function by \cref{relation,rightunique,pair_eq_iff}.
    Thus $\dom{v} = \naturalsPlus$ by \cref{dom_intro,dom_iff,pair_eq_iff,subseteq,subseteq_antisymmetric}.
    For all $n\in\dom{v}$ we have $v(n) = \snd{u(n)}$ by \cref{function_apply_intro}.
    For all $n\in\dom{v}$ we have $v(n)\in K$
        by \cref{funs_type_apply,subseteq,times_elem_is_tuple,function_apply_intro,pair_eq_iff,snd_eq}.
    Therefore $v\in\funs{\naturalsPlus}{K}$ by \cref{funs_intro}.

    $I\subseteq\reals$ by \cref{intervalclosed,subseteq}.
    $\standardTopologyR$ is a topology on $\reals$
        by \cref{standard_basis_is_basis,basis_topology_is_topology,standard_topology_reals}.
    Thus $I$ is a subspace of $\reals$ with respect to $\standardTopologyR$
        by \cref{subspace_of,subseteq}.
    For all $J$ such that $J\in C$ we have $J$ is closed in $I$ with respect to $T_I$
        by \cref{img_iff,subseteq,funs_type_apply,space_filling_closed_interval_family_member_closed}.
    Therefore $\inters{C}$ is inhabited
        by \cref{subspace_topology_is_topology,space_filling_domain_compact,space_filling_compact_fip_inhabited_intersection}.
    Take $y$ such that $y\in\inters{C}$ by \cref{nonempty_inhabited}.
    Take $J_0$ such that $J_0\in C$ by \cref{finite_intersection_property}.
    Therefore there exists $z\in I$ such that
        for all $n\in\naturalsPlus$ we have $z\in\snd{u(n)}$
        by \cref{subseteq,function_apply_elim,img_iff,inters_destr,function_apply_intro,funs_type_apply,powerset_elim}.
\end{proof}

% from §44 helper lemma 51: every dyadic state orbit has a common point
\begin{lemma}\label{space_filling_dyadic_child_state_orbit_common_point}
    Let $I = \intervalclosed{0}{1}$.
    Let $D = \{0,1\}$.
    Let $E = \{(n,D) \mid n\in\naturalsPlus\}$.
    Assume $x\in\productFamily{E}$.
    Let $K = \{J\in\pow{I} \mid \text{there exist $a,b\in\reals$ such that
        $J = \intervalclosed{a}{b}$ and
        $a \leq b$}\}$.
    Let $G = \naturalsPlus\times K$.
    Let $R = \{p\in G\times G \mid \text{there exist $n,J,H,a,b,d$ such that
        $n\in\naturalsPlus$ and
        $\fst{p} = (n,J)$ and
        $\snd{p} = (n + 1,H)$ and
        $d = x(n)$ and
        $a\in\reals$ and
        $b\in\reals$ and
        $a \leq b$ and
        $J = \intervalclosed{a}{b}$ and
        $H\subseteq J$ and
        $(d,H)\in\{(0,\intervalclosed{a}{(a + b) / (1 + 1)}),
            (1,\intervalclosed{(a + b) / (1 + 1)}{b})\}$}\}$.
    Suppose $g\in\funs{G}{G}$.
    Suppose for all $z\in G$ we have $z\mathrel{R}\apply{g}{z}$.
    Suppose $u$ is a sequence in $G$.
    Suppose for all $n\in\naturalsPlus$ we have $u(n + 1) = \apply{g}{u(n)}$.
    Then there exists $y\in I$ such that
        for all $n\in\naturalsPlus$ we have $y\in\snd{u(n)}$.
\end{lemma}
\begin{proof}
    Let $v = \{(n,\snd{u(n)}) \mid n\in\naturalsPlus\}$.
    Let $C = \img{v}{\naturalsPlus}$.
    Therefore there exists $z\in I$ such that
        for all $n\in\naturalsPlus$ we have $z\in\snd{u(n)}$
        by \cref{space_filling_dyadic_child_state_orbit_fip,space_filling_closed_interval_state_orbit_common_point_from_fip}.
\end{proof}

% from §44 helper lemma 52: a binary sequence selects retained-third child intervals statewise
\begin{lemma}\label{space_filling_retained_third_state_selector}
    Let $I = \intervalclosed{0}{1}$.
    Let $D = \{0,1\}$.
    Let $E = \{(n,D) \mid n\in\naturalsPlus\}$.
    Assume $x\in\productFamily{E}$.
    Let $K = \{J\in\pow{I} \mid \text{there exist $a,b\in\reals$ such that
        $J = \intervalclosed{a}{b}$ and
        $a \leq b$}\}$.
    Let $G = \naturalsPlus\times K$.
    Let $R = \{p\in G\times G \mid \text{there exist $n,J,H,a,b,d$ such that
        $n\in\naturalsPlus$ and
        $\fst{p} = (n,J)$ and
        $\snd{p} = (n + 1,H)$ and
        $d = x(n)$ and
        $a\in\reals$ and
        $b\in\reals$ and
        $a \leq b$ and
        $J = \intervalclosed{a}{b}$ and
        $H\subseteq J$ and
        $(d,H)\in\{(0,\intervalclosed{a}{(a + a + b) / (1 + 1 + 1)}),
            (1,\intervalclosed{(a + b + b) / (1 + 1 + 1)}{b})\}$}\}$.
    Then there exists $g$ such that
        $g\in\funs{G}{G}$ and
        for all $z\in G$ we have $z\mathrel{R}\apply{g}{z}$.
\end{lemma}
\begin{proof}
    $E$ is a function on $\naturalsPlus$
        by \cref{relation,rightunique,dom_intro,dom_iff,pair_eq_iff,subseteq,subseteq_antisymmetric}.
    $\dom{E} = \naturalsPlus$ by \cref{funs}.
    For all $n\in\naturalsPlus$ we have $E(n) = D$ by \cref{function_apply_intro}.
    $R$ is a relation by \cref{relation,subseteq,times_elem_is_tuple}.
    Show that for all $z\in G$ there exists $w$ such that $z\mathrel{R} w$.
    \begin{subproof}
        Fix $z\in G$.
        Take $n,J$ such that $z = (n,J)$ and $n\in\naturalsPlus$ and $J\in K$
            by \cref{subseteq,times_elem_is_tuple}.
        Take $a,b$ such that
            $a\in\reals$ and
            $b\in\reals$ and
            $J = \intervalclosed{a}{b}$ and
            $a \leq b$
            by \cref{subseteq}.
        $x(n) = 0$ or $x(n) = 1$ by \cref{indexed_product,subseteq,upair_iff}.
        \begin{byCase}
        \caseOf{$x(n) = 0$.}
            Let $H = \intervalclosed{a}{(a + a + b) / (1 + 1 + 1)}$.
            $H\subseteq J$ by \cref{space_filling_left_third_subset_parent}.
            We have $(x(n),H)\in\{(0,\intervalclosed{a}{(a + a + b) / (1 + 1 + 1)}),
                (1,\intervalclosed{(a + b + b) / (1 + 1 + 1)}{b})\}$ by \cref{upair_iff,pair_eq_iff,subseteq}.
            Hence $H\in K$
                by \cref{subseteq,powerset_elim,subseteq_transitive,powerset_intro,real_add_closed,space_filling_three_in_naturalsplus,real_div_naturalsplus_real,space_filling_left_endpoint_leq_left_third}.
            Thus $z\mathrel{R} (n + 1,H)$
                by \cref{relation,subseteq,real_naturals_closed_succ,times_tuple_intro,fst_eq,snd_eq,pair_eq_iff}.
        \caseOf{$x(n) = 1$.}
            Let $H = \intervalclosed{(a + b + b) / (1 + 1 + 1)}{b}$.
            $H\subseteq J$ by \cref{space_filling_right_third_subset_parent}.
            We have $(x(n),H)\in\{(0,\intervalclosed{a}{(a + a + b) / (1 + 1 + 1)}),
                (1,\intervalclosed{(a + b + b) / (1 + 1 + 1)}{b})\}$ by \cref{upair_iff,pair_eq_iff,subseteq}.
            Hence $H\in K$
                by \cref{subseteq,powerset_elim,subseteq_transitive,powerset_intro,real_add_closed,space_filling_three_in_naturalsplus,real_div_naturalsplus_real,space_filling_right_third_leq_right_endpoint}.
            Thus $z\mathrel{R} (n + 1,H)$
                by \cref{relation,subseteq,real_naturals_closed_succ,times_tuple_intro,fst_eq,snd_eq,pair_eq_iff}.
        \end{byCase}
    \end{subproof}
    Take $g$ such that
        $g$ is a function on $G$ and
        for all $z\in G$ we have $z\mathrel{R}\apply{g}{z}$
        by \cref{choice_function}.
    $g\in\funs{G}{G}$ by \cref{funs_intro,relation,subseteq,times_tuple_elim}.
\end{proof}

% from §44 helper lemma 53: the second coordinate of a retained-third state orbit is nested
\begin{lemma}\label{space_filling_retained_third_state_orbit_nested}
    Let $I = \intervalclosed{0}{1}$.
    Let $D = \{0,1\}$.
    Let $E = \{(n,D) \mid n\in\naturalsPlus\}$.
    Assume $x\in\productFamily{E}$.
    Let $K = \{J\in\pow{I} \mid \text{there exist $a,b\in\reals$ such that
        $J = \intervalclosed{a}{b}$ and
        $a \leq b$}\}$.
    Let $G = \naturalsPlus\times K$.
    Let $R = \{p\in G\times G \mid \text{there exist $n,J,H,a,b,d$ such that
        $n\in\naturalsPlus$ and
        $\fst{p} = (n,J)$ and
        $\snd{p} = (n + 1,H)$ and
        $d = x(n)$ and
        $a\in\reals$ and
        $b\in\reals$ and
        $a \leq b$ and
        $J = \intervalclosed{a}{b}$ and
        $H\subseteq J$ and
        $(d,H)\in\{(0,\intervalclosed{a}{(a + a + b) / (1 + 1 + 1)}),
            (1,\intervalclosed{(a + b + b) / (1 + 1 + 1)}{b})\}$}\}$.
    Suppose $g\in\funs{G}{G}$.
    Suppose for all $z\in G$ we have $z\mathrel{R}\apply{g}{z}$.
    Suppose $u$ is a sequence in $G$.
    Suppose for all $n\in\naturalsPlus$ we have $u(n + 1) = \apply{g}{u(n)}$.
    Then for all $n\in\naturalsPlus$ we have $\snd{u(n + 1)}\subseteq \snd{u(n)}$.
\end{lemma}
\begin{proof}
    Follows by \cref{sequence_in,funs_type_apply,relation,subseteq,fst_eq,snd_eq,pair_eq_iff}.
\end{proof}

% from §44 helper lemma 54: the interval coordinates of a retained-third state orbit are monotone
\begin{lemma}\label{space_filling_retained_third_state_orbit_monotone}
    Let $I = \intervalclosed{0}{1}$.
    Let $D = \{0,1\}$.
    Let $E = \{(n,D) \mid n\in\naturalsPlus\}$.
    Assume $x\in\productFamily{E}$.
    Let $K = \{J\in\pow{I} \mid \text{there exist $a,b\in\reals$ such that
        $J = \intervalclosed{a}{b}$ and
        $a \leq b$}\}$.
    Let $G = \naturalsPlus\times K$.
    Let $R = \{p\in G\times G \mid \text{there exist $n,J,H,a,b,d$ such that
        $n\in\naturalsPlus$ and
        $\fst{p} = (n,J)$ and
        $\snd{p} = (n + 1,H)$ and
        $d = x(n)$ and
        $a\in\reals$ and
        $b\in\reals$ and
        $a \leq b$ and
        $J = \intervalclosed{a}{b}$ and
        $H\subseteq J$ and
        $(d,H)\in\{(0,\intervalclosed{a}{(a + a + b) / (1 + 1 + 1)}),
            (1,\intervalclosed{(a + b + b) / (1 + 1 + 1)}{b})\}$}\}$.
    Suppose $g\in\funs{G}{G}$.
    Suppose for all $z\in G$ we have $z\mathrel{R}\apply{g}{z}$.
    Suppose $u$ is a sequence in $G$.
    Suppose for all $n\in\naturalsPlus$ we have $u(n + 1) = \apply{g}{u(n)}$.
    Then for all $m,n\in\naturalsPlus$ if $m \leq n$ then $\snd{u(n)}\subseteq \snd{u(m)}$.
\end{lemma}
\begin{proof}
    Let $P = \{n\in\naturalsPlus \mid \text{for all $m\in\initialSegment{n}$ we have $\snd{u(n)}\subseteq \snd{u(m)}$}\}$.
    $P\subseteq\naturalsPlus$ by \cref{subseteq}.
    Show that for all $m\in\initialSegment{1}$ we have $\snd{u(1)}\subseteq \snd{u(m)}$.
    \begin{subproof}
        Trivial.
    \end{subproof}
    $1\in P$ by \cref{real_one_in_naturals,subseteq}.
    Show that for all $n\in P$ we have $n + 1\in P$.
    \begin{subproof}
        Fix $n\in P$.
        Show that for all $m\in\initialSegment{n + 1}$ we have $\snd{u(n + 1)}\subseteq \snd{u(m)}$.
        \begin{subproof}
            Fix $m\in\initialSegment{n + 1}$.
            $\snd{u(n + 1)}\subseteq \snd{u(n)}$ by \cref{space_filling_retained_third_state_orbit_nested}.
            Hence $\snd{u(n + 1)}\subseteq \snd{u(m)}$
                by \cref{subseteq_refl,subseteq,initial_segment_naturalsplus,real_naturals_leq_succ_elim,subseteq_transitive}.
        \end{subproof}
        Thus $n + 1\in P$ by \cref{subseteq,real_naturals_closed_succ}.
    \end{subproof}
    Hence for all $n\in\naturalsPlus$ we have $n\in P$ by \cref{real_naturals_induction}.
    Fix $m$.
    Fix $n$ such that $m\in\naturalsPlus$ and $n\in\naturalsPlus$.
    Assume $m \leq n$.
    $\snd{u(n)}\subseteq \snd{u(m)}$ by \cref{subseteq,initial_segment_naturalsplus}.
\end{proof}

% from §44 helper lemma 55: the interval members of a retained-third state orbit have the finite intersection property
\begin{lemma}\label{space_filling_retained_third_state_orbit_fip}
    Let $I = \intervalclosed{0}{1}$.
    Let $D = \{0,1\}$.
    Let $E = \{(n,D) \mid n\in\naturalsPlus\}$.
    Assume $x\in\productFamily{E}$.
    Let $K = \{J\in\pow{I} \mid \text{there exist $a,b\in\reals$ such that
        $J = \intervalclosed{a}{b}$ and
        $a \leq b$}\}$.
    Let $G = \naturalsPlus\times K$.
    Let $R = \{p\in G\times G \mid \text{there exist $n,J,H,a,b,d$ such that
        $n\in\naturalsPlus$ and
        $\fst{p} = (n,J)$ and
        $\snd{p} = (n + 1,H)$ and
        $d = x(n)$ and
        $a\in\reals$ and
        $b\in\reals$ and
        $a \leq b$ and
        $J = \intervalclosed{a}{b}$ and
        $H\subseteq J$ and
        $(d,H)\in\{(0,\intervalclosed{a}{(a + a + b) / (1 + 1 + 1)}),
            (1,\intervalclosed{(a + b + b) / (1 + 1 + 1)}{b})\}$}\}$.
    Suppose $g\in\funs{G}{G}$.
    Suppose for all $z\in G$ we have $z\mathrel{R}\apply{g}{z}$.
    Suppose $u$ is a sequence in $G$.
    Suppose for all $n\in\naturalsPlus$ we have $u(n + 1) = \apply{g}{u(n)}$.
    Let $v = \{(n,\snd{u(n)}) \mid n\in\naturalsPlus\}$.
    Let $C = \img{v}{\naturalsPlus}$.
    Then $C$ has the finite intersection property.
\end{lemma}
\begin{proof}
    Follows by \cref{space_filling_retained_third_state_orbit_monotone,space_filling_closed_interval_state_orbit_fip_from_monotone}.
\end{proof}

% from §44 helper lemma 56: every retained-third state orbit has a common point
\begin{lemma}\label{space_filling_retained_third_state_orbit_common_point}
    Let $I = \intervalclosed{0}{1}$.
    Let $D = \{0,1\}$.
    Let $E = \{(n,D) \mid n\in\naturalsPlus\}$.
    Assume $x\in\productFamily{E}$.
    Let $K = \{J\in\pow{I} \mid \text{there exist $a,b\in\reals$ such that
        $J = \intervalclosed{a}{b}$ and
        $a \leq b$}\}$.
    Let $G = \naturalsPlus\times K$.
    Let $R = \{p\in G\times G \mid \text{there exist $n,J,H,a,b,d$ such that
        $n\in\naturalsPlus$ and
        $\fst{p} = (n,J)$ and
        $\snd{p} = (n + 1,H)$ and
        $d = x(n)$ and
        $a\in\reals$ and
        $b\in\reals$ and
        $a \leq b$ and
        $J = \intervalclosed{a}{b}$ and
        $H\subseteq J$ and
        $(d,H)\in\{(0,\intervalclosed{a}{(a + a + b) / (1 + 1 + 1)}),
            (1,\intervalclosed{(a + b + b) / (1 + 1 + 1)}{b})\}$}\}$.
    Suppose $g\in\funs{G}{G}$.
    Suppose for all $z\in G$ we have $z\mathrel{R}\apply{g}{z}$.
    Suppose $u$ is a sequence in $G$.
    Suppose for all $n\in\naturalsPlus$ we have $u(n + 1) = \apply{g}{u(n)}$.
    Then there exists $y\in I$ such that
        for all $n\in\naturalsPlus$ we have $y\in\snd{u(n)}$.
\end{lemma}
\begin{proof}
    Let $v = \{(n,\snd{u(n)}) \mid n\in\naturalsPlus\}$.
    Let $C = \img{v}{\naturalsPlus}$.
    There exists $z\in I$ such that
        for all $n\in\naturalsPlus$ we have $z\in\snd{u(n)}$
        by \cref{space_filling_retained_third_state_orbit_fip,space_filling_closed_interval_state_orbit_common_point_from_fip}.
\end{proof}

% from §44 helper lemma 57: the first coordinate of a retained-third state orbit is the stage number
\begin{lemma}\label{space_filling_retained_third_state_orbit_index}
    Let $I = \intervalclosed{0}{1}$.
    Let $D = \{0,1\}$.
    Let $E = \{(n,D) \mid n\in\naturalsPlus\}$.
    Assume $x\in\productFamily{E}$.
    Let $K = \{J\in\pow{I} \mid \text{there exist $a,b\in\reals$ such that
        $J = \intervalclosed{a}{b}$ and
        $a \leq b$}\}$.
    Let $G = \naturalsPlus\times K$.
    Let $R = \{p\in G\times G \mid \text{there exist $n,J,H,a,b,d$ such that
        $n\in\naturalsPlus$ and
        $\fst{p} = (n,J)$ and
        $\snd{p} = (n + 1,H)$ and
        $d = x(n)$ and
        $a\in\reals$ and
        $b\in\reals$ and
        $a \leq b$ and
        $J = \intervalclosed{a}{b}$ and
        $H\subseteq J$ and
        $(d,H)\in\{(0,\intervalclosed{a}{(a + a + b) / (1 + 1 + 1)}),
            (1,\intervalclosed{(a + b + b) / (1 + 1 + 1)}{b})\}$}\}$.
    Suppose $g\in\funs{G}{G}$.
    Suppose for all $z\in G$ we have $z\mathrel{R}\apply{g}{z}$.
    Suppose $u$ is a sequence in $G$.
    Suppose $u(1) = (1,I)$.
    Suppose for all $n\in\naturalsPlus$ we have $u(n + 1) = \apply{g}{u(n)}$.
    Then for all $n\in\naturalsPlus$ we have $\fst{u(n)} = n$.
\end{lemma}
\begin{proof}
    Let $A_0 = \{n\in\naturalsPlus \mid \fst{u(n)} = n\}$.
    $A_0\subseteq\naturalsPlus$ by \cref{subseteq}.
    Show that $\fst{u(1)} = 1$.
    \begin{subproof}
        Follows by \cref{pair_eq_iff,fst_eq,subseteq}.
    \end{subproof}
    $1\in A_0$ by \cref{real_one_in_naturals,subseteq}.
    Show that for all $n\in A_0$ we have $n + 1\in A_0$.
    \begin{subproof}
        Fix $n\in A_0$.
        We have $n\in\naturalsPlus$ and $\fst{u(n)} = n$ by \cref{subseteq}.
        Take $k,J,H,a,b,d$ such that
            $(u(n),u(n + 1))\in G\times G$ and
            $k\in\naturalsPlus$ and
            $\fst{(u(n),u(n + 1))} = (k,J)$ and
            $\snd{(u(n),u(n + 1))} = (k + 1,H)$ and
            $a\in\reals$ and
            $b\in\reals$ and
            $a \leq b$ and
            $J = \intervalclosed{a}{b}$ and
            $H\subseteq J$ and
            $(d,H)\in\{(0,\intervalclosed{a}{(a + a + b) / (1 + 1 + 1)}),
                (1,\intervalclosed{(a + b + b) / (1 + 1 + 1)}{b})\}$
            by \cref{sequence_in,funs_type_apply,relation,subseteq}.
        $k = n$ and $\fst{u(n + 1)} = n + 1$
            by \cref{fst_eq,snd_eq,pair_eq_iff,subseteq,real_add_eq_subst}.
        Hence $n + 1\in A_0$ by \cref{real_naturals_closed_succ,subseteq}.
    \end{subproof}
    Thus for all $n\in\naturalsPlus$ we have $\fst{u(n)} = n$
        by \cref{real_naturals_induction,subseteq}.
\end{proof}

% from §44 helper lemma 84: a dyadic child is unique for a fixed parent interval and digit
\begin{lemma}\label{space_filling_dyadic_child_unique}
    Let $D = \{0,1\}$.
    Suppose $a\in\reals$ and $b\in\reals$ and $a \leq b$.
    Let $J = \intervalclosed{a}{b}$.
    Assume $d\in D$.
    Assume $H_1\subseteq J$.
    Assume if $d = 0$ then $H_1 = \intervalclosed{a}{(a + b) / (1 + 1)}$.
    Assume if $d = 1$ then $H_1 = \intervalclosed{(a + b) / (1 + 1)}{b}$.
    Assume $H_2\subseteq J$.
    Assume if $d = 0$ then $H_2 = \intervalclosed{a}{(a + b) / (1 + 1)}$.
    Assume if $d = 1$ then $H_2 = \intervalclosed{(a + b) / (1 + 1)}{b}$.
    Then $H_1 = H_2$.
\end{lemma}
\begin{proof}
    Follows by \cref{upair_iff,subseteq,real_lt_subst_left}.
\end{proof}

% from §44 helper lemma 85: dyadic child choices with the same digit coincide
\begin{lemma}\label{space_filling_dyadic_child_choice_unique}
    Let $D = \{0,1\}$.
    Suppose $d\in D$.
    Suppose $(d,H_1)\in\{(0,\intervalclosed{a}{(a + b) / (1 + 1)}),
        (1,\intervalclosed{(a + b) / (1 + 1)}{b})\}$.
    Suppose $(d,H_2)\in\{(0,\intervalclosed{a}{(a + b) / (1 + 1)}),
        (1,\intervalclosed{(a + b) / (1 + 1)}{b})\}$.
    Then $H_1 = H_2$.
\end{lemma}
\begin{proof}
    Follows by \cref{upair_iff,subseteq,pair_eq_iff,real_add_ident,real_one_in_naturals,real_naturals_subset_reals,real_zero_lt_one,real_lt_irreflexive,real_lt_subst_left,real_lt_subst_right}.
\end{proof}

% from §44 helper lemma 86: dyadic state orbits agree on a common binary prefix
% from §44 helper lemma 59: ordered closed intervals have unique endpoints
\begin{lemma}\label{space_filling_intervalclosed_endpoints_unique}
    Suppose $a_1\in\reals$ and $b_1\in\reals$ and $a_1 \leq b_1$.
    Suppose $a_2\in\reals$ and $b_2\in\reals$ and $a_2 \leq b_2$.
    Suppose $\intervalclosed{a_1}{b_1} = \intervalclosed{a_2}{b_2}$.
    Then $a_1 = a_2$ and $b_1 = b_2$.
\end{lemma}
\begin{proof}
    Follows by \cref{intervalclosed,real_leq_antisym}.
\end{proof}

\begin{lemma}\label{space_filling_dyadic_state_orbit_prefix_agreement}
    Let $I = \intervalclosed{0}{1}$.
    Let $D = \{0,1\}$.
    Let $E = \{(n,D) \mid n\in\naturalsPlus\}$.
    Assume $x_1\in\productFamily{E}$.
    Assume $x_2\in\productFamily{E}$.
    Let $K = \{J\in\pow{I} \mid \text{there exist $a,b\in\reals$ such that
        $J = \intervalclosed{a}{b}$ and
        $a \leq b$}\}$.
    Let $G = \naturalsPlus\times K$.
    Let $R_1 = \{p\in G\times G \mid \text{there exist $n,J,H,a,b,d$ such that
        $n\in\naturalsPlus$ and
        $\fst{p} = (n,J)$ and
        $\snd{p} = (n + 1,H)$ and
        $d = x_1(n)$ and
        $a\in\reals$ and
        $b\in\reals$ and
        $a \leq b$ and
        $J = \intervalclosed{a}{b}$ and
        $H\subseteq J$ and
        $(d,H)\in\{(0,\intervalclosed{a}{(a + b) / (1 + 1)}),
            (1,\intervalclosed{(a + b) / (1 + 1)}{b})\}$}\}$.
    Let $R_2 = \{p\in G\times G \mid \text{there exist $n,J,H,a,b,d$ such that
        $n\in\naturalsPlus$ and
        $\fst{p} = (n,J)$ and
        $\snd{p} = (n + 1,H)$ and
        $d = x_2(n)$ and
        $a\in\reals$ and
        $b\in\reals$ and
        $a \leq b$ and
        $J = \intervalclosed{a}{b}$ and
        $H\subseteq J$ and
        $(d,H)\in\{(0,\intervalclosed{a}{(a + b) / (1 + 1)}),
            (1,\intervalclosed{(a + b) / (1 + 1)}{b})\}$}\}$.
    Suppose $g_1\in\funs{G}{G}$.
    Suppose for all $z\in G$ we have $z\mathrel{R_1}\apply{g_1}{z}$.
    Suppose $u_1$ is a sequence in $G$.
    Suppose $u_1(1) = (1,I)$.
    Suppose for all $n\in\naturalsPlus$ we have $u_1(n + 1) = \apply{g_1}{u_1(n)}$.
    Suppose $g_2\in\funs{G}{G}$.
    Suppose for all $z\in G$ we have $z\mathrel{R_2}\apply{g_2}{z}$.
    Suppose $u_2$ is a sequence in $G$.
    Suppose $u_2(1) = (1,I)$.
    Suppose for all $n\in\naturalsPlus$ we have $u_2(n + 1) = \apply{g_2}{u_2(n)}$.
    Assume $N\in\naturalsPlus$.
    Assume for all $i\in\initialSegment{N}$ we have $x_1(i) = x_2(i)$.
    Then for all $n\in\initialSegment{N}$ we have $\snd{u_1(n)} = \snd{u_2(n)}$.
\end{lemma}
\begin{proof}
    $E$ is a function on $\naturalsPlus$
        by \cref{relation,rightunique,dom_intro,dom_iff,pair_eq_iff,subseteq,subseteq_antisymmetric}.
    Hence $\dom{E} = \naturalsPlus$ by \cref{funs}.
    For all $n\in\naturalsPlus$ we have $E(n) = D$ by \cref{function_apply_intro}.
    $u_1\in\funs{\naturalsPlus}{G}$ and $u_2\in\funs{\naturalsPlus}{G}$ by \cref{sequence_in}.
    $R_1$ is a relation by \cref{relation,subseteq,times_elem_is_tuple}.
    $R_2$ is a relation by \cref{relation,subseteq,times_elem_is_tuple}.
    Let $A_0 = \{n\in\naturalsPlus \mid \text{if $n \leq N$ then $\snd{u_1(n)} = \snd{u_2(n)}$}\}$.
    $A_0\subseteq\naturalsPlus$ by \cref{subseteq}.
    Show that if $1 \leq N$ then $\snd{u_1(1)} = \snd{u_2(1)}$.
    \begin{subproof}
        Follows by \cref{subseteq,snd_eq,pair_eq_iff}.
    \end{subproof}
    $1\in A_0$ by \cref{real_one_in_naturals,subseteq}.
    Show that for all $n\in A_0$ we have $n + 1\in A_0$.
    \begin{subproof}
        Fix $n\in A_0$.
        We have $n\in\naturalsPlus$ by \cref{subseteq}.
        Show that if $n + 1 \leq N$ then $\snd{u_1(n + 1)} = \snd{u_2(n + 1)}$.
        \begin{subproof}
            Assume $n + 1 \leq N$.
            $n \leq n + 1$ by \cref{real_naturals_lt_succ_local,real_lt_imp_leq}.
            We have $n\in\reals$ and $N\in\reals$ and $n + 1\in\reals$
                by \cref{real_naturals_subset_reals,subseteq,real_one_in_naturals,real_add_closed}.
            Thus $n \leq N$ by \cref{real_leq_trans}.
            $\snd{u_1(n)} = \snd{u_2(n)}$ by \cref{subseteq}.
            We have $n\in\initialSegment{N}$ and $n + 1\in\initialSegment{N}$ by \cref{initial_segment_naturalsplus,real_naturals_closed_succ}.
            $x_1(n) = x_2(n)$ by \cref{subseteq}.
            We have $x_1(n)\in D$ and $x_2(n)\in D$ by \cref{indexed_product,subseteq}.

            We have $u_1(n)\in G$ and $u_2(n)\in G$ by \cref{funs_type_apply}.
            $u_1(n)\mathrel{R_1}\apply{g_1}{u_1(n)}$ and
                $u_2(n)\mathrel{R_2}\apply{g_2}{u_2(n)}$ by assumption.
            Thus $u_1(n)\mathrel{R_1}u_1(n + 1)$ and
                $u_2(n)\mathrel{R_2}u_2(n + 1)$ by \cref{subseteq}.
            $(u_1(n),u_1(n + 1))\in R_1$ and
                $(u_2(n),u_2(n + 1))\in R_2$ by \cref{relation}.

            Take $k_1,J_1,H_1,a_1,b_1,d_1$ such that
                $(u_1(n),u_1(n + 1))\in G\times G$ and
                $k_1\in\naturalsPlus$ and
                $\fst{(u_1(n),u_1(n + 1))} = (k_1,J_1)$ and
                $\snd{(u_1(n),u_1(n + 1))} = (k_1 + 1,H_1)$ and
                $d_1 = x_1(k_1)$ and
                $a_1\in\reals$ and
                $b_1\in\reals$ and
                $a_1 \leq b_1$ and
                $J_1 = \intervalclosed{a_1}{b_1}$ and
                $H_1\subseteq J_1$ and
                $(d_1,H_1)\in\{(0,\intervalclosed{a_1}{(a_1 + b_1) / (1 + 1)}),
                    (1,\intervalclosed{(a_1 + b_1) / (1 + 1)}{b_1})\}$
                by \cref{subseteq}.
            Take $k_2,J_2,H_2,a_2,b_2,d_2$ such that
                $(u_2(n),u_2(n + 1))\in G\times G$ and
                $k_2\in\naturalsPlus$ and
                $\fst{(u_2(n),u_2(n + 1))} = (k_2,J_2)$ and
                $\snd{(u_2(n),u_2(n + 1))} = (k_2 + 1,H_2)$ and
                $d_2 = x_2(k_2)$ and
                $a_2\in\reals$ and
                $b_2\in\reals$ and
                $a_2 \leq b_2$ and
                $J_2 = \intervalclosed{a_2}{b_2}$ and
                $H_2\subseteq J_2$ and
                $(d_2,H_2)\in\{(0,\intervalclosed{a_2}{(a_2 + b_2) / (1 + 1)}),
                    (1,\intervalclosed{(a_2 + b_2) / (1 + 1)}{b_2})\}$
                by \cref{subseteq}.

            We have $u_1(n) = (k_1,J_1)$ and $u_1(n + 1) = (k_1 + 1,H_1)$ and
                $u_2(n) = (k_2,J_2)$ and $u_2(n + 1) = (k_2 + 1,H_2)$
                by \cref{fst_eq,snd_eq,pair_eq_iff}.
            $\fst{u_1(n)} = k_1$ and $\fst{u_2(n)} = k_2$ and
                $\snd{u_1(n + 1)} = H_1$ and $\snd{u_2(n + 1)} = H_2$
                by \cref{fst_eq,snd_eq,pair_eq_iff}.
            Hence $\fst{u_1(n)} = n$ and $\fst{u_2(n)} = n$
                by \cref{space_filling_dyadic_child_state_orbit_index}.
            $k_1 = n$ and $k_2 = n$ by \cref{subseteq}.
            $d_1 = x_1(n)$ and $d_2 = x_2(n)$ and $d_1 = d_2$ by \cref{subseteq}.
            $\snd{u_1(n)} = J_1$ and $\snd{u_2(n)} = J_2$ by \cref{snd_eq,pair_eq_iff}.
            $J_1 = J_2$ by \cref{subseteq}.
            $a_1 = a_2$ and $b_1 = b_2$ by \cref{subseteq,space_filling_intervalclosed_endpoints_unique}.
            $(d_2,H_1)\in\{(0,\intervalclosed{a_2}{(a_2 + b_2) / (1 + 1)}),
                (1,\intervalclosed{(a_2 + b_2) / (1 + 1)}{b_2})\}$
                by \cref{subseteq,real_add_eq_subst}.
            Hence $H_1 = H_2$ by \cref{space_filling_dyadic_child_choice_unique,subseteq,real_add_eq_subst}.
            $\snd{u_1(n + 1)} = \snd{u_2(n + 1)}$ by \cref{subseteq}.
        \end{subproof}
        Thus $n + 1\in A_0$ by \cref{subseteq,real_naturals_closed_succ}.
    \end{subproof}
    Hence for all $n\in\naturalsPlus$ we have $n\in A_0$ by \cref{real_naturals_induction}.
    For all $n\in\initialSegment{N}$ we have $\snd{u_1(n)} = \snd{u_2(n)}$
        by \cref{initial_segment_naturalsplus,subseteq}.
\end{proof}

% from §44 helper definition 2: retained-third state selector
\begin{definition}\label{space_filling_retained_third_state_selector_on}
    $g$ is a retained-third state selector on $G$ and $R$ iff
        $g\in\funs{G}{G}$ and
        for all $z\in G$ we have $z\mathrel{R}\apply{g}{z}$.
\end{definition}

% from §44 helper definition 3: retained-third state orbit from the root interval
\begin{definition}\label{space_filling_retained_third_state_orbit_from}
    $u$ is a retained-third state orbit of $g$ on $G$ from $I$ iff
        $u$ is a sequence in $G$ and
        $u(1) = (1,I)$ and
        for all $n\in\naturalsPlus$ we have $u(n + 1) = \apply{g}{u(n)}$.
\end{definition}

% from §44 helper definition 4: retained-third code point of a binary sequence
\begin{definition}\label{space_filling_retained_third_code_point}
    $y$ is a retained-third code point of $x$ in $I$ and $E$ iff
        $x\in\productFamily{E}$ and
        $y\in I$ and
        there exist $K,G,R,g,u$ such that
            $K = \{J\in\pow{I} \mid \text{there exist $a,b\in\reals$ such that
                $J = \intervalclosed{a}{b}$ and
                $a \leq b$}\}$ and
            $G = \naturalsPlus\times K$ and
            $R = \{p\in G\times G \mid \text{there exist $n,J,H,a,b,d$ such that
                $n\in\naturalsPlus$ and
                $\fst{p} = (n,J)$ and
                $\snd{p} = (n + 1,H)$ and
                $d = x(n)$ and
                $a\in\reals$ and
                $b\in\reals$ and
                $a \leq b$ and
                $J = \intervalclosed{a}{b}$ and
                $H\subseteq J$ and
                $(d,H)\in\{(0,\intervalclosed{a}{(a + a + b) / (1 + 1 + 1)}),
                    (1,\intervalclosed{(a + b + b) / (1 + 1 + 1)}{b})\}$}\}$ and
            $g$ is a retained-third state selector on $G$ and $R$ and
            $u$ is a retained-third state orbit of $g$ on $G$ from $I$ and
            for all $n\in\naturalsPlus$ we have $y\in\snd{u(n)}$.
\end{definition}

% from §44 helper lemma 89: every binary sequence has a retained-third code point
\begin{lemma}\label{space_filling_retained_third_code_point_exists}
    Let $I = \intervalclosed{0}{1}$.
    Let $D = \{0,1\}$.
    Let $E = \{(n,D) \mid n\in\naturalsPlus\}$.
    Assume $x\in\productFamily{E}$.
    Then there exists $y$ such that $y$ is a retained-third code point of $x$ in $I$ and $E$.
\end{lemma}
\begin{proof}
    Let $K = \{J\in\pow{I} \mid \text{there exist $a,b\in\reals$ such that
        $J = \intervalclosed{a}{b}$ and
        $a \leq b$}\}$.
    Let $G = \naturalsPlus\times K$.
    Let $R = \{p\in G\times G \mid \text{there exist $n,J,H,a,b,d$ such that
        $n\in\naturalsPlus$ and
        $\fst{p} = (n,J)$ and
        $\snd{p} = (n + 1,H)$ and
        $d = x(n)$ and
        $a\in\reals$ and
        $b\in\reals$ and
        $a \leq b$ and
        $J = \intervalclosed{a}{b}$ and
        $H\subseteq J$ and
        $(d,H)\in\{(0,\intervalclosed{a}{(a + a + b) / (1 + 1 + 1)}),
            (1,\intervalclosed{(a + b + b) / (1 + 1 + 1)}{b})\}$}\}$.
    Take $g$ such that
        $g\in\funs{G}{G}$ and
        for all $z\in G$ we have $z\mathrel{R}\apply{g}{z}$
        by \cref{space_filling_retained_third_state_selector}.
    Take $u$ such that
        $u$ is a sequence in $G$ and
        $u(1) = (1,I)$ and
        for all $n\in\naturalsPlus$ we have $u(n + 1) = \apply{g}{u(n)}$
        by \cref{space_filling_dyadic_child_state_iteration}.
    Take $y\in I$ such that
        for all $n\in\naturalsPlus$ we have $y\in\snd{u(n)}$
        by \cref{space_filling_retained_third_state_orbit_common_point}.
    $y$ is a retained-third code point of $x$ in $I$ and $E$
        by \cref{space_filling_retained_third_state_selector_on,space_filling_retained_third_state_orbit_from,space_filling_retained_third_code_point}.
\end{proof}

% from §44 helper lemma 87: retained-third child choices with the same digit coincide
\begin{lemma}\label{space_filling_retained_third_child_choice_unique}
    Let $D = \{0,1\}$.
    Suppose $d\in D$.
    Suppose $(d,H_1)\in\{(0,\intervalclosed{a}{(a + a + b) / (1 + 1 + 1)}),
        (1,\intervalclosed{(a + b + b) / (1 + 1 + 1)}{b})\}$.
    Suppose $(d,H_2)\in\{(0,\intervalclosed{a}{(a + a + b) / (1 + 1 + 1)}),
        (1,\intervalclosed{(a + b + b) / (1 + 1 + 1)}{b})\}$.
    Then $H_1 = H_2$.
\end{lemma}
\begin{proof}
    Follows by \cref{upair_iff,subseteq,pair_eq_iff,real_add_ident,real_one_in_naturals,real_naturals_subset_reals,real_zero_lt_one,real_lt_irreflexive,real_lt_subst_left,real_lt_subst_right}.
\end{proof}

% from §44 helper lemma 88: retained-third state orbits agree on a common binary prefix
\begin{lemma}\label{space_filling_retained_third_state_orbit_prefix_agreement}
    Let $I = \intervalclosed{0}{1}$.
    Let $D = \{0,1\}$.
    Let $E = \{(n,D) \mid n\in\naturalsPlus\}$.
    Assume $x_1\in\productFamily{E}$.
    Assume $x_2\in\productFamily{E}$.
    Let $K = \{J\in\pow{I} \mid \text{there exist $a,b\in\reals$ such that
        $J = \intervalclosed{a}{b}$ and
        $a \leq b$}\}$.
    Let $G = \naturalsPlus\times K$.
    Let $R_1 = \{p\in G\times G \mid \text{there exist $n,J,H,a,b,d$ such that
        $n\in\naturalsPlus$ and
        $\fst{p} = (n,J)$ and
        $\snd{p} = (n + 1,H)$ and
        $d = x_1(n)$ and
        $a\in\reals$ and
        $b\in\reals$ and
        $a \leq b$ and
        $J = \intervalclosed{a}{b}$ and
        $H\subseteq J$ and
        $(d,H)\in\{(0,\intervalclosed{a}{(a + a + b) / (1 + 1 + 1)}),
            (1,\intervalclosed{(a + b + b) / (1 + 1 + 1)}{b})\}$}\}$.
    Let $R_2 = \{p\in G\times G \mid \text{there exist $n,J,H,a,b,d$ such that
        $n\in\naturalsPlus$ and
        $\fst{p} = (n,J)$ and
        $\snd{p} = (n + 1,H)$ and
        $d = x_2(n)$ and
        $a\in\reals$ and
        $b\in\reals$ and
        $a \leq b$ and
        $J = \intervalclosed{a}{b}$ and
        $H\subseteq J$ and
        $(d,H)\in\{(0,\intervalclosed{a}{(a + a + b) / (1 + 1 + 1)}),
            (1,\intervalclosed{(a + b + b) / (1 + 1 + 1)}{b})\}$}\}$.
    Suppose $g_1\in\funs{G}{G}$.
    Suppose for all $z\in G$ we have $z\mathrel{R_1}\apply{g_1}{z}$.
    Suppose $u_1$ is a sequence in $G$.
    Suppose $u_1(1) = (1,I)$.
    Suppose for all $n\in\naturalsPlus$ we have $u_1(n + 1) = \apply{g_1}{u_1(n)}$.
    Suppose $g_2\in\funs{G}{G}$.
    Suppose for all $z\in G$ we have $z\mathrel{R_2}\apply{g_2}{z}$.
    Suppose $u_2$ is a sequence in $G$.
    Suppose $u_2(1) = (1,I)$.
    Suppose for all $n\in\naturalsPlus$ we have $u_2(n + 1) = \apply{g_2}{u_2(n)}$.
    Assume $N\in\naturalsPlus$.
    Assume for all $i\in\initialSegment{N}$ we have $x_1(i) = x_2(i)$.
    Then for all $n\in\initialSegment{N}$ we have $\snd{u_1(n)} = \snd{u_2(n)}$.
\end{lemma}
\begin{proof}
    $E$ is a function on $\naturalsPlus$
        by \cref{relation,rightunique,dom_intro,dom_iff,pair_eq_iff,subseteq,subseteq_antisymmetric}.
    Hence $\dom{E} = \naturalsPlus$ by \cref{funs}.
    For all $n\in\naturalsPlus$ we have $E(n) = D$ by \cref{function_apply_intro}.
    $u_1\in\funs{\naturalsPlus}{G}$ and $u_2\in\funs{\naturalsPlus}{G}$ by \cref{sequence_in}.
    $R_1$ is a relation by \cref{relation,subseteq,times_elem_is_tuple}.
    $R_2$ is a relation by \cref{relation,subseteq,times_elem_is_tuple}.
    Let $A_0 = \{n\in\naturalsPlus \mid \text{if $n \leq N$ then $\snd{u_1(n)} = \snd{u_2(n)}$}\}$.
    $A_0\subseteq\naturalsPlus$ by \cref{subseteq}.
    Show that if $1 \leq N$ then $\snd{u_1(1)} = \snd{u_2(1)}$.
    \begin{subproof}
        Follows by \cref{subseteq,snd_eq,pair_eq_iff}.
    \end{subproof}
    $1\in A_0$ by \cref{real_one_in_naturals,subseteq}.
    Show that for all $n\in A_0$ we have $n + 1\in A_0$.
    \begin{subproof}
        Fix $n\in A_0$.
        We have $n\in\naturalsPlus$ by \cref{subseteq}.
        Show that if $n + 1 \leq N$ then $\snd{u_1(n + 1)} = \snd{u_2(n + 1)}$.
        \begin{subproof}
            Assume $n + 1 \leq N$.
            $n \leq n + 1$ by \cref{real_naturals_lt_succ_local,real_lt_imp_leq}.
            We have $n\in\reals$ and $N\in\reals$ and $n + 1\in\reals$
                by \cref{real_naturals_subset_reals,subseteq,real_one_in_naturals,real_add_closed}.
            Thus $n \leq N$ by \cref{real_leq_trans}.
            $\snd{u_1(n)} = \snd{u_2(n)}$ by \cref{subseteq}.
            We have $n\in\initialSegment{N}$ and $n + 1\in\initialSegment{N}$ by \cref{initial_segment_naturalsplus,real_naturals_closed_succ}.
            $x_1(n) = x_2(n)$ by \cref{subseteq}.
            We have $x_1(n)\in D$ and $x_2(n)\in D$ by \cref{indexed_product,subseteq}.

            We have $u_1(n)\in G$ and $u_2(n)\in G$ by \cref{funs_type_apply}.
            $u_1(n)\mathrel{R_1}\apply{g_1}{u_1(n)}$ and
                $u_2(n)\mathrel{R_2}\apply{g_2}{u_2(n)}$ by assumption.
            Thus $u_1(n)\mathrel{R_1}u_1(n + 1)$ and
                $u_2(n)\mathrel{R_2}u_2(n + 1)$ by \cref{subseteq}.
            $(u_1(n),u_1(n + 1))\in R_1$ and
                $(u_2(n),u_2(n + 1))\in R_2$ by \cref{relation}.

            Take $k_1,J_1,H_1,a_1,b_1,d_1$ such that
                $(u_1(n),u_1(n + 1))\in G\times G$ and
                $k_1\in\naturalsPlus$ and
                $\fst{(u_1(n),u_1(n + 1))} = (k_1,J_1)$ and
                $\snd{(u_1(n),u_1(n + 1))} = (k_1 + 1,H_1)$ and
                $d_1 = x_1(k_1)$ and
                $a_1\in\reals$ and
                $b_1\in\reals$ and
                $a_1 \leq b_1$ and
                $J_1 = \intervalclosed{a_1}{b_1}$ and
                $H_1\subseteq J_1$ and
                $(d_1,H_1)\in\{(0,\intervalclosed{a_1}{(a_1 + a_1 + b_1) / (1 + 1 + 1)}),
                    (1,\intervalclosed{(a_1 + b_1 + b_1) / (1 + 1 + 1)}{b_1})\}$
                by \cref{subseteq}.
            Take $k_2,J_2,H_2,a_2,b_2,d_2$ such that
                $(u_2(n),u_2(n + 1))\in G\times G$ and
                $k_2\in\naturalsPlus$ and
                $\fst{(u_2(n),u_2(n + 1))} = (k_2,J_2)$ and
                $\snd{(u_2(n),u_2(n + 1))} = (k_2 + 1,H_2)$ and
                $d_2 = x_2(k_2)$ and
                $a_2\in\reals$ and
                $b_2\in\reals$ and
                $a_2 \leq b_2$ and
                $J_2 = \intervalclosed{a_2}{b_2}$ and
                $H_2\subseteq J_2$ and
                $(d_2,H_2)\in\{(0,\intervalclosed{a_2}{(a_2 + a_2 + b_2) / (1 + 1 + 1)}),
                    (1,\intervalclosed{(a_2 + b_2 + b_2) / (1 + 1 + 1)}{b_2})\}$
                by \cref{subseteq}.

            We have $u_1(n) = (k_1,J_1)$ and $u_1(n + 1) = (k_1 + 1,H_1)$ and
                $u_2(n) = (k_2,J_2)$ and $u_2(n + 1) = (k_2 + 1,H_2)$
                by \cref{fst_eq,snd_eq,pair_eq_iff}.
            $\fst{u_1(n)} = k_1$ and $\fst{u_2(n)} = k_2$ and
                $\snd{u_1(n + 1)} = H_1$ and $\snd{u_2(n + 1)} = H_2$
                by \cref{fst_eq,snd_eq,pair_eq_iff}.
            Hence $\fst{u_1(n)} = n$ by \cref{space_filling_retained_third_state_orbit_index}.
            $\fst{u_2(n)} = n$ by \cref{space_filling_retained_third_state_orbit_index}.
            $k_1 = n$ and $k_2 = n$ by \cref{subseteq}.
            $d_1 = x_1(n)$ and $d_2 = x_2(n)$ and $d_1 = d_2$ by \cref{subseteq}.
            $\snd{u_1(n)} = J_1$ and $\snd{u_2(n)} = J_2$ by \cref{snd_eq,pair_eq_iff}.
            $J_1 = J_2$ by \cref{subseteq}.
            $a_1 = a_2$ and $b_1 = b_2$ by \cref{subseteq,space_filling_intervalclosed_endpoints_unique}.
            $(d_2,H_1)\in\{(0,\intervalclosed{a_2}{(a_2 + a_2 + b_2) / (1 + 1 + 1)}),
                (1,\intervalclosed{(a_2 + b_2 + b_2) / (1 + 1 + 1)}{b_2})\}$
                by \cref{subseteq,real_add_eq_subst}.
            Hence $H_1 = H_2$ by \cref{space_filling_retained_third_child_choice_unique,subseteq,real_add_eq_subst}.
            $\snd{u_1(n + 1)} = \snd{u_2(n + 1)}$ by \cref{subseteq}.
        \end{subproof}
        Thus $n + 1\in A_0$ by \cref{subseteq,real_naturals_closed_succ}.
    \end{subproof}
    Hence for all $n\in\naturalsPlus$ we have $n\in A_0$ by \cref{real_naturals_induction}.
    For all $n\in\initialSegment{N}$ we have $\snd{u_1(n)} = \snd{u_2(n)}$
        by \cref{initial_segment_naturalsplus,subseteq}.
\end{proof}

% from §44 helper lemma 58: a retained-third child is unique for a fixed parent interval and digit
\begin{lemma}\label{space_filling_retained_third_child_unique}
    Let $D = \{0,1\}$.
    Suppose $a\in\reals$ and $b\in\reals$ and $a \leq b$.
    Let $J = \intervalclosed{a}{b}$.
    Assume $d\in D$.
    Assume $H_1\subseteq J$.
    Assume if $d = 0$ then $H_1 = \intervalclosed{a}{(a + a + b) / (1 + 1 + 1)}$.
    Assume if $d = 1$ then $H_1 = \intervalclosed{(a + b + b) / (1 + 1 + 1)}{b}$.
    Assume $H_2\subseteq J$.
    Assume if $d = 0$ then $H_2 = \intervalclosed{a}{(a + a + b) / (1 + 1 + 1)}$.
    Assume if $d = 1$ then $H_2 = \intervalclosed{(a + b + b) / (1 + 1 + 1)}{b}$.
    Then $H_1 = H_2$.
\end{lemma}
\begin{proof}
    Follows by \cref{upair_iff,subseteq,real_lt_subst_left}.
\end{proof}

% from §44 helper lemma 60: the retained-third state relation is right-unique
\begin{lemma}\label{space_filling_retained_third_state_relation_rightunique}
    Let $I = \intervalclosed{0}{1}$.
    Let $D = \{0,1\}$.
    Let $E = \{(n,D) \mid n\in\naturalsPlus\}$.
    Assume $x\in\productFamily{E}$.
    Let $K = \{J\in\pow{I} \mid \text{there exist $a,b\in\reals$ such that
        $J = \intervalclosed{a}{b}$ and
        $a \leq b$}\}$.
    Let $G = \naturalsPlus\times K$.
    Let $R = \{p\in G\times G \mid \text{there exist $n,J,H,a,b,d$ such that
        $n\in\naturalsPlus$ and
        $\fst{p} = (n,J)$ and
        $\snd{p} = (n + 1,H)$ and
        $d = x(n)$ and
        $a\in\reals$ and
        $b\in\reals$ and
        $a \leq b$ and
        $J = \intervalclosed{a}{b}$ and
        $H\subseteq J$ and
        $(d,H)\in\{(0,\intervalclosed{a}{(a + a + b) / (1 + 1 + 1)}),
            (1,\intervalclosed{(a + b + b) / (1 + 1 + 1)}{b})\}$}\}$.
    Then $R$ is right-unique.
\end{lemma}
\begin{proof}
    $E$ is a function on $\naturalsPlus$
        by \cref{relation,rightunique,dom_intro,dom_iff,pair_eq_iff,subseteq,subseteq_antisymmetric}.
    Hence $\dom{E} = \naturalsPlus$ by \cref{funs}.
    For all $n\in\naturalsPlus$ we have $E(n) = D$ by \cref{function_apply_intro}.
    $R$ is a relation by \cref{relation,subseteq,times_elem_is_tuple}.
    Show that for all $z,w_1,w_2$ such that $(z,w_1)\in R$ and $(z,w_2)\in R$ we have $w_1 = w_2$.
    \begin{subproof}
        Fix $z$.
        Fix $w_1$.
        Fix $w_2$ such that $(z,w_1)\in R$ and $(z,w_2)\in R$.
        Take $n_1,J_1,H_1,a_1,b_1,d_1$ such that
            $(z,w_1)\in G\times G$ and
            $n_1\in\naturalsPlus$ and
            $\fst{(z,w_1)} = (n_1,J_1)$ and
            $\snd{(z,w_1)} = (n_1 + 1,H_1)$ and
            $d_1 = x(n_1)$ and
            $a_1\in\reals$ and
            $b_1\in\reals$ and
            $a_1 \leq b_1$ and
            $J_1 = \intervalclosed{a_1}{b_1}$ and
            $(d_1,H_1)\in\{(0,\intervalclosed{a_1}{(a_1 + a_1 + b_1) / (1 + 1 + 1)}),
                (1,\intervalclosed{(a_1 + b_1 + b_1) / (1 + 1 + 1)}{b_1})\}$
            by \cref{subseteq}.
        Take $n_2,J_2,H_2,a_2,b_2,d_2$ such that
            $(z,w_2)\in G\times G$ and
            $n_2\in\naturalsPlus$ and
            $\fst{(z,w_2)} = (n_2,J_2)$ and
            $\snd{(z,w_2)} = (n_2 + 1,H_2)$ and
            $d_2 = x(n_2)$ and
            $a_2\in\reals$ and
            $b_2\in\reals$ and
            $a_2 \leq b_2$ and
            $J_2 = \intervalclosed{a_2}{b_2}$ and
            $(d_2,H_2)\in\{(0,\intervalclosed{a_2}{(a_2 + a_2 + b_2) / (1 + 1 + 1)}),
                (1,\intervalclosed{(a_2 + b_2 + b_2) / (1 + 1 + 1)}{b_2})\}$
            by \cref{subseteq}.
        We have $z = (n_1,J_1)$ and $w_1 = (n_1 + 1,H_1)$ by \cref{fst_eq,snd_eq,pair_eq_iff}.
        $z = (n_2,J_2)$ and $w_2 = (n_2 + 1,H_2)$ by \cref{fst_eq,snd_eq,pair_eq_iff}.
        Therefore $n_1 = n_2$ and $J_1 = J_2$ by \cref{pair_eq_iff}.
        $d_1 = d_2$ by \cref{subseteq}.
        $a_1 = a_2$ and $b_1 = b_2$ by \cref{subseteq,space_filling_intervalclosed_endpoints_unique}.
        We have $d_1\in D$ by \cref{indexed_product,funs,subseteq}.
        $(d_2,H_1)\in\{(0,\intervalclosed{a_2}{(a_2 + a_2 + b_2) / (1 + 1 + 1)}),
            (1,\intervalclosed{(a_2 + b_2 + b_2) / (1 + 1 + 1)}{b_2})\}$
            by \cref{subseteq,real_add_eq_subst}.
        Thus $H_1 = H_2$ by \cref{space_filling_retained_third_child_choice_unique,subseteq,real_add_eq_subst}.
        Hence $w_1 = w_2$ by \cref{pair_eq_iff,real_add_eq_subst}.
    \end{subproof}
    Hence $R$ is right-unique by \cref{rightunique}.
\end{proof}

% from §44 helper lemma 61: points in an ordered closed interval differ by at most its width
\begin{lemma}\label{space_filling_interval_width_bounds_difference}
    Suppose $a\in\reals$ and $b\in\reals$ and $a \leq b$.
    Suppose $y_1\in\intervalclosed{a}{b}$.
    Suppose $y_2\in\intervalclosed{a}{b}$.
    Suppose $y_1 \leq y_2$.
    Then $y_2 - y_1 \leq b - a$.
\end{lemma}
\begin{proof}
    We have $y_1\in\reals$ and $a \leq y_1$ and $y_1 \leq b$ and
        $y_2\in\reals$ and $a \leq y_2$ and $y_2 \leq b$ by \cref{intervalclosed,subseteq}.
    $\neg{y_1}\in\reals$ and $y_2 - y_1\in\reals$ by \cref{real_add_inverse,real_sub,real_add_closed}.
    Thus $a + (y_2 - y_1) \leq y_2$
        by \cref{real_leq_add,real_add_sub_cancel_right,real_leq_subst_right_local}.
    $a + (y_2 - y_1) \leq b$ by \cref{real_add_closed,real_leq_trans}.
    Hence $(a + (y_2 - y_1)) + \neg{a} \leq b + \neg{a}$
        by \cref{real_add_inverse,real_sub,real_add_closed,real_leq_add}.
    $(a + (y_2 - y_1)) + \neg{a} = y_2 - y_1$
        by \cref{real_add_assoc,real_add_comm,real_add_inverse,real_add_ident,real_add_eq_subst}.
    Therefore $y_2 - y_1 \leq b - a$
        by \cref{real_sub,real_leq_subst_left_local,real_leq_subst_right_local}.
\end{proof}

% from §44 helper lemma 62: the left dyadic half has half the parent width
\begin{lemma}\label{space_filling_left_half_width_double}
    Suppose $a\in\reals$ and $b\in\reals$.
    Then $(((a + b) / (1 + 1)) - a) + (((a + b) / (1 + 1)) - a) = b - a$.
\end{lemma}
\begin{proof}
    $1 + 1\in\naturalsPlus$ by \cref{real_one_in_naturals,real_naturals_closed_succ}.
    Let $m = (a + b) / (1 + 1)$.
    We have $a + b\in\reals$ and $m\in\reals$ and $\neg{a}\in\reals$ and $m - a\in\reals$ and $b - a\in\reals$
        by \cref{real_add_closed,real_div_naturalsplus_real,real_add_inverse,real_sub}.
    $m + m = a + b$ by \cref{real_half_plus_half}.
    Therefore $(m - a) + (m - a) = (a + b) + \neg{(a + a)}$
        by \cref{real_sub,real_half_plus_half,real_neg_addition,real_add_assoc,real_add_comm,real_add_eq_subst}.
    Hence $(((a + b) / (1 + 1)) - a) + (((a + b) / (1 + 1)) - a) = b - a$
        by \cref{real_sub,real_add_assoc,real_add_comm,real_add_inverse,real_add_ident,real_add_eq_subst,subseteq}.
\end{proof}

% from §44 helper lemma 63: the midpoint splits the interval symmetrically
\begin{lemma}\label{space_filling_midpoint_symmetric_difference}
    Suppose $a\in\reals$ and $b\in\reals$.
    Then $((a + b) / (1 + 1)) - a = b - ((a + b) / (1 + 1))$.
\end{lemma}
\begin{proof}
    $1 + 1\in\naturalsPlus$ by \cref{real_one_in_naturals,real_naturals_closed_succ}.
    Let $m = (a + b) / (1 + 1)$.
    We have $a + b\in\reals$ and $m\in\reals$ and $\neg{a}\in\reals$ and $\neg{m}\in\reals$
        by \cref{real_add_closed,real_div_naturalsplus_real,real_add_inverse}.
    Thus $m - a\in\reals$ and $b - m\in\reals$ by \cref{real_sub,real_add_closed}.
    $m + (m - a) = b$
        by \cref{real_sub,real_half_plus_half,real_add_assoc,real_add_comm,real_add_inverse,real_add_ident,real_add_eq_subst}.
    We have $m + (b - m) = b$ by \cref{real_add_sub_cancel_right}.
    $(m - a) + m = b$ and $(b - m) + m = b$ by \cref{real_add_comm,real_add_eq_subst}.
    Therefore $((a + b) / (1 + 1)) - a = b - ((a + b) / (1 + 1))$
        by \cref{real_add_cancel,subseteq}.
\end{proof}

% from §44 helper lemma 64: the right dyadic half has half the parent width
\begin{lemma}\label{space_filling_right_half_width_double}
    Suppose $a\in\reals$ and $b\in\reals$.
    Then $(b - ((a + b) / (1 + 1))) + (b - ((a + b) / (1 + 1))) = b - a$.
\end{lemma}
\begin{proof}
    We have $(b - ((a + b) / (1 + 1))) + (b - ((a + b) / (1 + 1))) = b - a$
        by \cref{space_filling_midpoint_symmetric_difference,space_filling_left_half_width_double,real_add_eq_subst}.
\end{proof}

% from §44 helper lemma 65: a successor is bounded by the corresponding double
\begin{lemma}\label{space_filling_succ_leq_double}
    Suppose $n\in\naturalsPlus$.
    Then $n + 1 \leq n + n$.
\end{lemma}
\begin{proof}
    Follows by \cref{real_naturals_subset_reals,subseteq,real_one_in_naturals,real_one_leq_naturalsplus,real_leq_add,real_add_comm,real_leq_subst_left_local}.
\end{proof}

% from §44 helper lemma 66: reciprocal doubles are below reciprocal successors
\begin{lemma}\label{space_filling_double_reciprocal_leq_succ_reciprocal}
    Suppose $n\in\naturalsPlus$.
    Then $1 / (n + n) \leq 1 / (n + 1)$.
\end{lemma}
\begin{proof}
    Follows by \cref{real_naturals_closed_succ,naturalsplus_add_closed,space_filling_succ_leq_double,naturalsplus_reciprocal_antitone}.
\end{proof}

% from §44 helper lemma 67: a reciprocal is bounded by the doubled next reciprocal
\begin{lemma}\label{space_filling_reciprocal_leq_double_succ_reciprocal}
    Suppose $n\in\naturalsPlus$.
    Then $1 / n \leq 1 / (n + 1) + 1 / (n + 1)$.
\end{lemma}
\begin{proof}
    We have $n\in\reals$ and $n + 1\in\reals$ and $n + n\in\reals$
        by \cref{real_naturals_subset_reals,subseteq,real_naturals_closed_succ,real_add_closed}.
    $n \neq 0$ and $n + 1 \neq 0$
        by \cref{naturalsplus_positive_real,subseteq,real_naturals_closed_succ,real_pos_nonzero}.
    Hence $\inv{n}\in\reals$ and $\inv{(n + 1)}\in\reals$ by \cref{real_mul_inverse}.
    Let $c = \inv{n} \rmul \inv{(n + 1)}$.
    We have $c\in\reals$ and $0 < c$
        by \cref{naturalsplus_positive_real,subseteq,real_naturals_closed_succ,real_inv_pos,real_mul_closed,real_mul_pos_pos_gt}.
    Therefore $(n + 1) \rmul c \leq (n + n) \rmul c$
        by \cref{space_filling_succ_leq_double,real_zero_lt_one,real_pos_add,real_leq_split,real_lt_mul_pos,real_lt_imp_leq,real_mul_eq_subst,subseteq}.
    $(n + 1) \rmul c = 1 / n$
        by \cref{real_mul_assoc,real_mul_comm,real_mul_inverse,real_div,real_mul_ident}.
    Thus $(n + n) \rmul c = 1 / (n + 1) + 1 / (n + 1)$
        by \cref{real_right_distrib,real_mul_assoc,real_mul_inverse,real_div,real_mul_ident,real_add_eq_subst}.
    Hence $1 / n \leq 1 / (n + 1) + 1 / (n + 1)$
        by \cref{real_leq_subst_left_local,real_leq_subst_right_local}.
\end{proof}

% from §44 helper lemma 68: equal denominators cancel from doubled weak inequalities
\begin{lemma}\label{space_filling_double_leq_cancel}
    Suppose $d\in\reals$ and $r\in\reals$.
    Suppose $d + d \leq r + r$.
    Then $d \leq r$.
\end{lemma}
\begin{proof}
    Follows by \cref{real_one_in_naturals,real_naturals_closed_succ,real_add_closed,real_div_leq_same_denom,space_filling_double_half_real,real_leq_subst_left_local,real_leq_subst_right_local}.
\end{proof}

% from §44 helper lemma 69: a left dyadic child improves the reciprocal width bound
\begin{lemma}\label{space_filling_left_child_width_leq_succ_reciprocal}
    Suppose $n\in\naturalsPlus$.
    Suppose $a\in\reals$ and $b\in\reals$ and $a \leq b$.
    Suppose $b - a \leq 1 / n$.
    Then $((a + b) / (1 + 1)) - a \leq 1 / (n + 1)$.
\end{lemma}
\begin{proof}
    We have $1 / n\in\reals$ and $1 / (n + 1)\in\reals$
        by \cref{naturalsplus_reciprocal_real,real_naturals_closed_succ}.
    Let $d = ((a + b) / (1 + 1)) - a$.
    $d\in\reals$ and $d + d \leq 1 / n$
        by \cref{real_sub,real_add_inverse,real_add_closed,real_div_naturalsplus_real,real_one_in_naturals,real_naturals_closed_succ,naturalsplus_reciprocal_real,space_filling_left_half_width_double,real_leq_subst_left_local,real_leq_trans}.
    Thus $d + d \leq 1 / (n + 1) + 1 / (n + 1)$
        by \cref{real_add_closed,real_naturals_closed_succ,naturalsplus_reciprocal_real,real_leq_trans,space_filling_reciprocal_leq_double_succ_reciprocal}.
    Hence $((a + b) / (1 + 1)) - a \leq 1 / (n + 1)$ by \cref{space_filling_double_leq_cancel}.
\end{proof}

% from §44 helper lemma 70: a right dyadic child improves the reciprocal width bound
\begin{lemma}\label{space_filling_right_child_width_leq_succ_reciprocal}
    Suppose $n\in\naturalsPlus$.
    Suppose $a\in\reals$ and $b\in\reals$ and $a \leq b$.
    Suppose $b - a \leq 1 / n$.
    Then $b - ((a + b) / (1 + 1)) \leq 1 / (n + 1)$.
\end{lemma}
\begin{proof}
    Follows by \cref{space_filling_midpoint_symmetric_difference,space_filling_left_child_width_leq_succ_reciprocal,real_leq_subst_left_local}.
\end{proof}

% from §44 helper lemma 71: dyadic state-orbit widths are bounded by reciprocal stage numbers
\begin{lemma}\label{space_filling_dyadic_child_state_orbit_width_bound}
    Let $I = \intervalclosed{0}{1}$.
    Let $D = \{0,1\}$.
    Let $E = \{(n,D) \mid n\in\naturalsPlus\}$.
    Assume $x\in\productFamily{E}$.
    Let $K = \{J\in\pow{I} \mid \text{there exist $a,b\in\reals$ such that
        $J = \intervalclosed{a}{b}$ and
        $a \leq b$}\}$.
    Let $G = \naturalsPlus\times K$.
    Let $R = \{p\in G\times G \mid \text{there exist $n,J,H,a,b,d$ such that
        $n\in\naturalsPlus$ and
        $\fst{p} = (n,J)$ and
        $\snd{p} = (n + 1,H)$ and
        $d = x(n)$ and
        $a\in\reals$ and
        $b\in\reals$ and
        $a \leq b$ and
        $J = \intervalclosed{a}{b}$ and
        $H\subseteq J$ and
        $(d,H)\in\{(0,\intervalclosed{a}{(a + b) / (1 + 1)}),
            (1,\intervalclosed{(a + b) / (1 + 1)}{b})\}$}\}$.
    Suppose $g\in\funs{G}{G}$.
    Suppose for all $z\in G$ we have $z\mathrel{R}\apply{g}{z}$.
    Suppose $u$ is a sequence in $G$.
    Suppose $u(1) = (1,I)$.
    Suppose for all $n\in\naturalsPlus$ we have $u(n + 1) = \apply{g}{u(n)}$.
    Then for all $n,a,b$ if
        $n\in\naturalsPlus$ and
        $a\in\reals$ and
        $b\in\reals$ and
        $\snd{u(n)} = \intervalclosed{a}{b}$ and
        $a \leq b$
    then $b - a \leq 1 / n$.
\end{lemma}
\begin{proof}
    Let $A_0 = \{n\in\naturalsPlus \mid \text{for all $a,b$ if
        $a\in\reals$ and
        $b\in\reals$ and
        $\snd{u(n)} = \intervalclosed{a}{b}$ and
        $a \leq b$
    then $b - a \leq 1 / n$}\}$.
    $A_0\subseteq\naturalsPlus$ by \cref{subseteq}.
    Show that for all $a,b$ if
        $a\in\reals$ and
        $b\in\reals$ and
        $\snd{u(1)} = \intervalclosed{a}{b}$ and
        $a \leq b$
    then $b - a \leq 1 / 1$.
    \begin{subproof}
        Follows by \cref{snd_eq,pair_eq_iff,subseteq,real_add_ident,real_one_in_naturals,real_naturals_subset_reals,real_zero_lt_one,real_lt_imp_leq,space_filling_intervalclosed_endpoints_unique,real_sub,real_add_eq_subst,real_neg_zero,real_add_comm,real_one_div_one}.
    \end{subproof}
    $1\in A_0$ by \cref{subseteq,real_one_in_naturals}.
    Show that for all $n\in A_0$ we have $n + 1\in A_0$.
    \begin{subproof}
        Fix $n\in A_0$.
        We have $n\in\naturalsPlus$ by \cref{subseteq}.
        Show that for all $a,b$ if
            $a\in\reals$ and
            $b\in\reals$ and
            $\snd{u(n + 1)} = \intervalclosed{a}{b}$ and
            $a \leq b$
        then $b - a \leq 1 / (n + 1)$.
        \begin{subproof}
            Fix $a$.
            Fix $b$.
            Assume $a\in\reals$ and $b\in\reals$ and $\snd{u(n + 1)} = \intervalclosed{a}{b}$ and $a \leq b$.
            $(u(n),u(n + 1))\in R$ by \cref{sequence_in,funs_type_apply,relation,subseteq}.
            Take $k,J,H,c,d,e$ such that
                $(u(n),u(n + 1))\in G\times G$ and
                $k\in\naturalsPlus$ and
                $\fst{(u(n),u(n + 1))} = (k,J)$ and
                $\snd{(u(n),u(n + 1))} = (k + 1,H)$ and
                $e = x(k)$ and
                $c\in\reals$ and
                $d\in\reals$ and
                $c \leq d$ and
                $J = \intervalclosed{c}{d}$ and
                $(e,H)\in\{(0,\intervalclosed{c}{(c + d) / (1 + 1)}),
                    (1,\intervalclosed{(c + d) / (1 + 1)}{d})\}$
                by \cref{subseteq}.
            We have $\snd{u(n)} = J$ and $\snd{u(n + 1)} = H$ and $\fst{u(n)} = k$
                by \cref{fst_eq,snd_eq,pair_eq_iff}.
            Thus $k = n$ by \cref{space_filling_dyadic_child_state_orbit_index,subseteq}.
            $d - c \leq 1 / n$ by \cref{subseteq}.
            $(e,H) = (0,\intervalclosed{c}{(c + d) / (1 + 1)})$ or
                $(e,H) = (1,\intervalclosed{(c + d) / (1 + 1)}{d})$
                by \cref{upair_iff}.
            \begin{byCase}
            \caseOf{$(e,H) = (0,\intervalclosed{c}{(c + d) / (1 + 1)})$.}
                $b - a \leq 1 / (n + 1)$
                    by \cref{pair_eq_iff,subseteq,real_add_closed,real_one_in_naturals,real_naturals_closed_succ,real_div_naturalsplus_real,space_filling_left_endpoint_leq_midpoint,space_filling_intervalclosed_endpoints_unique,real_sub,real_add_eq_subst,space_filling_left_child_width_leq_succ_reciprocal}.
            \caseOf{$(e,H) = (1,\intervalclosed{(c + d) / (1 + 1)}{d})$.}
                $b - a \leq 1 / (n + 1)$
                    by \cref{pair_eq_iff,subseteq,real_add_closed,real_one_in_naturals,real_naturals_closed_succ,real_div_naturalsplus_real,space_filling_midpoint_leq_right_endpoint,space_filling_intervalclosed_endpoints_unique,real_sub,real_add_eq_subst,space_filling_right_child_width_leq_succ_reciprocal}.
            \end{byCase}
        \end{subproof}
        Hence $n + 1\in A_0$ by \cref{subseteq,real_naturals_closed_succ}.
    \end{subproof}
    Therefore for all $n\in\naturalsPlus$ we have $n\in A_0$ by \cref{real_naturals_induction}.
    For all $n,a,b$ if
        $n\in\naturalsPlus$ and
        $a\in\reals$ and
        $b\in\reals$ and
        $\snd{u(n)} = \intervalclosed{a}{b}$ and
        $a \leq b$
    then $b - a \leq 1 / n$ by \cref{subseteq}.
\end{proof}

% from §44 helper lemma 72: a dyadic state orbit has at most one common point
\begin{lemma}\label{space_filling_dyadic_child_state_orbit_common_point_unique}
    Let $I = \intervalclosed{0}{1}$.
    Let $D = \{0,1\}$.
    Let $E = \{(n,D) \mid n\in\naturalsPlus\}$.
    Assume $x\in\productFamily{E}$.
    Let $K = \{J\in\pow{I} \mid \text{there exist $a,b\in\reals$ such that
        $J = \intervalclosed{a}{b}$ and
        $a \leq b$}\}$.
    Let $G = \naturalsPlus\times K$.
    Let $R = \{p\in G\times G \mid \text{there exist $n,J,H,a,b,d$ such that
        $n\in\naturalsPlus$ and
        $\fst{p} = (n,J)$ and
        $\snd{p} = (n + 1,H)$ and
        $d = x(n)$ and
        $a\in\reals$ and
        $b\in\reals$ and
        $a \leq b$ and
        $J = \intervalclosed{a}{b}$ and
        $H\subseteq J$ and
        $(d,H)\in\{(0,\intervalclosed{a}{(a + b) / (1 + 1)}),
            (1,\intervalclosed{(a + b) / (1 + 1)}{b})\}$}\}$.
    Suppose $g\in\funs{G}{G}$.
    Suppose for all $z\in G$ we have $z\mathrel{R}\apply{g}{z}$.
    Suppose $u$ is a sequence in $G$.
    Suppose $u(1) = (1,I)$.
    Suppose for all $n\in\naturalsPlus$ we have $u(n + 1) = \apply{g}{u(n)}$.
    Suppose $y_1\in I$ and for all $n\in\naturalsPlus$ we have $y_1\in\snd{u(n)}$.
    Suppose $y_2\in I$ and for all $n\in\naturalsPlus$ we have $y_2\in\snd{u(n)}$.
    Then $y_1 = y_2$.
\end{lemma}
\begin{proof}
    $y_1\in\reals$ and $y_2\in\reals$ by \cref{intervalclosed,subseteq}.
    Suppose not.
    $y_1 < y_2$ or $y_2 < y_1$ by \cref{real_lt_trichotomy}.
    \begin{byCase}
    \caseOf{$y_1 < y_2$.}
        We have $y_2 - y_1\in\reals$ and $0 < y_2 - y_1$
            by \cref{real_sub,real_add_inverse,real_add_closed,real_lt_imp_pos_diff}.
        Take $N\in\naturalsPlus$ such that $0 < 1 / N$ and $1 / N < y_2 - y_1$
            by \cref{subseteq,real_small_reciprocal}.
        Take $m,J$ such that
            $u(N) = (m,J)$ and
            $J\in K$
            by \cref{sequence_in,funs_type_apply,subseteq,times_elem_is_tuple}.
        Take $a,b$ such that
            $a\in\reals$ and
            $b\in\reals$ and
            $J = \intervalclosed{a}{b}$ and
            $a \leq b$
            by \cref{subseteq}.
        $\snd{u(N)} = \intervalclosed{a}{b}$ by \cref{snd_eq,pair_eq_iff,subseteq}.
        $y_1\in\intervalclosed{a}{b}$ and $y_2\in\intervalclosed{a}{b}$ by \cref{subseteq}.
        $y_2 - y_1 \leq b - a$
            by \cref{real_lt_imp_leq,space_filling_interval_width_bounds_difference}.
        $b - a \leq 1 / N$ by \cref{space_filling_dyadic_child_state_orbit_width_bound}.
        We have $b - a\in\reals$ and $1 / N\in\reals$
            by \cref{real_sub,real_add_inverse,real_add_closed,naturalsplus_reciprocal_real}.
        $y_2 - y_1 \leq 1 / N$ by \cref{real_leq_trans}.
        $1 / N < 1 / N$ by \cref{real_leq_split,real_lt_trans,real_lt_subst_right}.
        Contradiction by \cref{real_lt_irreflexive}.
    \caseOf{$y_2 < y_1$.}
        We have $y_1 - y_2\in\reals$ and $0 < y_1 - y_2$
            by \cref{real_sub,real_add_inverse,real_add_closed,real_lt_imp_pos_diff}.
        Take $N\in\naturalsPlus$ such that $0 < 1 / N$ and $1 / N < y_1 - y_2$
            by \cref{subseteq,real_small_reciprocal}.
        Take $m,J$ such that
            $u(N) = (m,J)$ and
            $J\in K$
            by \cref{sequence_in,funs_type_apply,subseteq,times_elem_is_tuple}.
        Take $a,b$ such that
            $a\in\reals$ and
            $b\in\reals$ and
            $J = \intervalclosed{a}{b}$ and
            $a \leq b$
            by \cref{subseteq}.
        $\snd{u(N)} = \intervalclosed{a}{b}$ by \cref{snd_eq,pair_eq_iff,subseteq}.
        $y_1\in\intervalclosed{a}{b}$ and $y_2\in\intervalclosed{a}{b}$ by \cref{subseteq}.
        $y_1 - y_2 \leq b - a$
            by \cref{real_lt_imp_leq,space_filling_interval_width_bounds_difference}.
        $b - a \leq 1 / N$ by \cref{space_filling_dyadic_child_state_orbit_width_bound}.
        We have $b - a\in\reals$ and $1 / N\in\reals$
            by \cref{real_sub,real_add_inverse,real_add_closed,naturalsplus_reciprocal_real}.
        $y_1 - y_2 \leq 1 / N$ by \cref{real_leq_trans}.
        $1 / N < 1 / N$ by \cref{real_leq_split,real_lt_trans,real_lt_subst_right}.
        Contradiction by \cref{real_lt_irreflexive}.
    \end{byCase}
\end{proof}

% from §44 helper lemma 73: a successor is bounded by the corresponding triple
\begin{lemma}\label{space_filling_succ_leq_triple}
    Suppose $n\in\naturalsPlus$.
    Then $n + 1 \leq n + n + n$.
\end{lemma}
\begin{proof}
    Follows by \cref{real_naturals_subset_reals,subseteq,real_one_in_naturals,real_add_closed,space_filling_succ_leq_double,real_add_ident,naturalsplus_positive_real,real_lt_imp_leq,real_leq_add,real_add_assoc,real_leq_subst_right_local,real_leq_trans}.
\end{proof}

% from §44 helper lemma 74: reciprocal triples are below reciprocal successors
\begin{lemma}\label{space_filling_triple_reciprocal_leq_succ_reciprocal}
    Suppose $n\in\naturalsPlus$.
    Then $1 / (n + n + n) \leq 1 / (n + 1)$.
\end{lemma}
\begin{proof}
    Follows by \cref{real_naturals_closed_succ,naturalsplus_add_closed,space_filling_succ_leq_triple,naturalsplus_reciprocal_antitone}.
\end{proof}

% from §44 helper lemma 75: a reciprocal is bounded by the tripled next reciprocal
\begin{lemma}\label{space_filling_reciprocal_leq_triple_succ_reciprocal}
    Suppose $n\in\naturalsPlus$.
    Then $1 / n \leq 1 / (n + 1) + 1 / (n + 1) + 1 / (n + 1)$.
\end{lemma}
\begin{proof}
    We have $n\in\reals$ and $n + 1\in\reals$ and $n + n\in\reals$ and $n + n + n\in\reals$
        by \cref{real_naturals_subset_reals,subseteq,real_naturals_closed_succ,real_add_closed}.
    $n \neq 0$ and $n + 1 \neq 0$
        by \cref{naturalsplus_positive_real,subseteq,real_naturals_closed_succ,real_pos_nonzero}.
    $\inv{n}\in\reals$ and $\inv{(n + 1)}\in\reals$ by \cref{real_mul_inverse}.
    Let $c = \inv{n} \rmul \inv{(n + 1)}$.
    We have $c\in\reals$ and $0 < c$
        by \cref{naturalsplus_positive_real,subseteq,real_naturals_closed_succ,real_inv_pos,real_mul_closed,real_mul_pos_pos_gt}.
    $(n + 1) \rmul c \leq (n + n + n) \rmul c$
        by \cref{space_filling_succ_leq_triple,real_zero_lt_one,real_pos_add,real_leq_split,real_lt_mul_pos,real_lt_imp_leq,real_mul_eq_subst,subseteq}.
    $(n + 1) \rmul c = 1 / n$
        by \cref{real_mul_assoc,real_mul_comm,real_mul_inverse,real_div,real_mul_ident}.
    $((n + n + n) \rmul c) = ((n + n) \rmul c) + (n \rmul c)$ by \cref{real_right_distrib}.
    $(n + n + n) \rmul c = (1 / (n + 1) + 1 / (n + 1)) + 1 / (n + 1)$
        by \cref{real_right_distrib,real_mul_assoc,real_mul_inverse,real_div,real_mul_ident,real_add_eq_subst}.
    $1 / n \leq 1 / (n + 1) + 1 / (n + 1) + 1 / (n + 1)$
        by \cref{real_leq_subst_left_local,real_leq_subst_right_local}.
\end{proof}

% from §44 helper lemma 76: equal denominators cancel from tripled weak inequalities
\begin{lemma}\label{space_filling_triple_leq_cancel}
    Suppose $d\in\reals$ and $r\in\reals$.
    Suppose $d + d + d \leq r + r + r$.
    Then $d \leq r$.
\end{lemma}
\begin{proof}
    Follows by \cref{real_one_in_naturals,real_naturals_closed_succ,real_add_closed,real_div_leq_same_denom,space_filling_triple_third_real,real_leq_subst_left_local,real_leq_subst_right_local}.
\end{proof}

% from §44 helper lemma 77: the left retained third has one third of the parent width
\begin{lemma}\label{space_filling_left_third_width_triple}
    Suppose $a\in\reals$ and $b\in\reals$.
    Then $((((a + a + b) / (1 + 1 + 1)) - a) + (((a + a + b) / (1 + 1 + 1)) - a))
        + (((a + a + b) / (1 + 1 + 1)) - a) = b - a$.
\end{lemma}
\begin{proof}
    We have $1 + 1 + 1\in\naturalsPlus$ by \cref{real_one_in_naturals,real_naturals_closed_succ}.
    $\neg{a}\in\reals$ and $b + \neg{a}\in\reals$ and
        $a + a\in\reals$ and $a + a + b\in\reals$ and $(a + a + b) + (a + a + b)\in\reals$ and
        $((a + a + b) + (a + a + b)) + (a + a + b)\in\reals$ and $a + a + a\in\reals$
        by \cref{real_add_inverse,real_add_closed}.
    Let $m = (a + a + b) / (1 + 1 + 1)$.
    Let $d = m - a$.
    We have $m\in\reals$ by \cref{real_div_naturalsplus_real}.
    $d\in\reals$ and $d + d\in\reals$ and $(d + d) + d\in\reals$ and $d + a\in\reals$ and
        $b - a\in\reals$ and $m + m\in\reals$ and $m + m + m\in\reals$
        by \cref{real_sub,real_add_inverse,real_add_closed}.
    $d + a = m$ by \cref{real_add_sub_cancel_right,real_add_comm,real_add_eq_subst}.
    $(d + a) + (d + a) = m + m$ by \cref{real_add_eq_subst}.
    $((d + a) + (d + a)) + (d + a) = (m + m) + m$ by \cref{real_add_eq_subst}.
    $(d + d) + (a + a) = (d + a) + (d + a)$ by \cref{real_add_four_exchange}.
    $((d + d) + (a + a)) + (d + a) = (m + m) + m$ by \cref{real_add_eq_subst}.
    $(d + d) + ((a + a) + (d + a)) = (m + m) + m$ by \cref{real_add_assoc}.
    $(d + d) + ((a + d) + (a + a)) = (m + m) + m$ by \cref{real_add_four_exchange,real_add_eq_subst}.
    $(d + d) + ((d + a) + (a + a)) = (m + m) + m$
        by \cref{real_add_comm,real_add_eq_subst}.
    $((d + d) + (d + a)) + (a + a) = (m + m) + m$ by \cref{real_add_assoc}.
    $(((d + d) + d) + a) + (a + a) = (m + m) + m$ by \cref{real_add_assoc,real_add_eq_subst}.
    $((d + d) + d) + (a + (a + a)) = (m + m) + m$ by \cref{real_add_assoc}.
    $((d + d) + d) + ((a + a) + a) = (m + m) + m$ by \cref{real_add_assoc,real_add_eq_subst}.
    $(m + m) + m
        = (((a + a + b) / (1 + 1 + 1)) + ((a + a + b) / (1 + 1 + 1)))
            + ((a + a + b) / (1 + 1 + 1))$ by \cref{real_add_eq_subst}.
    $(m + m) + m
        = (((a + a + b) + (a + a + b)) / (1 + 1 + 1))
            + ((a + a + b) / (1 + 1 + 1))$ by \cref{real_div_add_same_denom,real_add_eq_subst}.
    $(m + m) + m
        = (((a + a + b) + (a + a + b)) + (a + a + b)) / (1 + 1 + 1)$
        by \cref{real_div_add_same_denom,real_add_eq_subst}.
    $((a + a + b) + (a + a + b)) + (a + a + b) = (a + a + b) + ((a + a + b) + (a + a + b))$
        by \cref{real_add_assoc}.
    $((a + a + b) + (a + a + b)) + (a + a + b) = (a + a + b) + (a + a + b) + (a + a + b)$
        by \cref{real_add_assoc,real_add_eq_subst}.
    $(m + m) + m = a + a + b$ by \cref{space_filling_triple_third_real,real_add_eq_subst}.
    $((d + d) + d) + ((a + a) + a) = a + a + b$ by \cref{real_add_eq_subst}.
    We have $b - a = b + \neg{a}$ by \cref{real_sub}.
    $(b - a) + ((a + a) + a) = (b + \neg{a}) + ((a + a) + a)$ by \cref{real_add_eq_subst}.
    $(b + \neg{a}) + ((a + a) + a) = b + (\neg{a} + ((a + a) + a))$ by \cref{real_add_assoc}.
    $\neg{a} + ((a + a) + a) = (\neg{a} + (a + a)) + a$ by \cref{real_add_assoc}.
    $(\neg{a} + (a + a)) + a = ((\neg{a} + a) + a) + a$ by \cref{real_add_assoc}.
    $\neg{a} + a = 0$ by \cref{real_add_comm,real_add_inverse,real_add_eq_subst}.
    $(\neg{a} + a) + a = 0 + a$ by \cref{real_add_eq_subst}.
    $((\neg{a} + a) + a) + a = (0 + a) + a$ by \cref{real_add_eq_subst}.
    $(0 + a) + a = a + a$ by \cref{real_add_ident}.
    $\neg{a} + ((a + a) + a) = a + a$ by \cref{real_add_eq_subst}.
    $b + (\neg{a} + ((a + a) + a)) = b + (a + a)$ by \cref{real_add_eq_subst}.
    $b + (a + a) = (b + a) + a$ by \cref{real_add_assoc}.
    $(b + a) + a = ((a + b) + a)$ by \cref{real_add_comm,real_add_eq_subst}.
    $(a + b) + a = a + (b + a)$ by \cref{real_add_assoc}.
    $a + (b + a) = a + (a + b)$ by \cref{real_add_comm,real_add_eq_subst}.
    $a + (a + b) = (a + a) + b$ by \cref{real_add_assoc}.
    $(b - a) + ((a + a) + a) = a + a + b$ by \cref{real_add_eq_subst}.
    $((d + d) + d) + ((a + a) + a) = (b - a) + ((a + a) + a)$ by \cref{real_add_eq_subst}.
    $(d + d) + d = b - a$ by \cref{real_add_cancel}.
\end{proof}

% from §44 helper lemma 78: the two retained-third widths are equal
\begin{lemma}\label{space_filling_retained_third_symmetric_difference}
    Suppose $a\in\reals$ and $b\in\reals$.
    Then $((a + a + b) / (1 + 1 + 1)) - a = b - ((a + b + b) / (1 + 1 + 1))$.
\end{lemma}
\begin{proof}
    $1 + 1 + 1\in\naturalsPlus$ by \cref{space_filling_three_in_naturalsplus}.
    We have $a + b\in\reals$ and $a + a + b\in\reals$ and $a + b + b\in\reals$ and $\neg{a}\in\reals$
        by \cref{real_add_closed,real_add_inverse}.
    Let $m = (a + a + b) / (1 + 1 + 1)$.
    Let $n = (a + b + b) / (1 + 1 + 1)$.
    $m\in\reals$ and $n\in\reals$ by \cref{real_div_naturalsplus_real}.
    $m - a\in\reals$ and $b - n\in\reals$ by \cref{real_sub,real_add_inverse,real_add_closed}.
    $m + n = ((a + a + b) + (a + b + b)) / (1 + 1 + 1)$ by \cref{real_div_add_same_denom}.
    $m + n = a + b$
        by \cref{real_div_add_same_denom,real_add_assoc,real_add_comm,space_filling_triple_third_real,real_add_eq_subst}.
    $n + (m - a) = b$
        by \cref{real_sub,real_add_assoc,real_add_comm,space_filling_triple_third_real,real_add_inverse,real_add_ident,real_add_eq_subst}.
    $(m - a) + n = b$ and $(b - n) + n = b$
        by \cref{real_add_sub_cancel_right,real_add_comm,real_add_eq_subst}.
    $((a + a + b) / (1 + 1 + 1)) - a = b - ((a + b + b) / (1 + 1 + 1))$ by \cref{real_add_cancel}.
\end{proof}

% from §44 helper lemma 79: the right retained third has one third of the parent width
\begin{lemma}\label{space_filling_right_third_width_triple}
    Suppose $a\in\reals$ and $b\in\reals$.
    Then $(b - ((a + b + b) / (1 + 1 + 1)) + (b - ((a + b + b) / (1 + 1 + 1))))
        + (b - ((a + b + b) / (1 + 1 + 1))) = b - a$.
\end{lemma}
\begin{proof}
    Follows by \cref{space_filling_retained_third_symmetric_difference,space_filling_left_third_width_triple,real_add_eq_subst}.
\end{proof}

% from §44 helper lemma 80: a left retained third improves the reciprocal width bound
\begin{lemma}\label{space_filling_left_third_width_leq_succ_reciprocal}
    Suppose $n\in\naturalsPlus$.
    Suppose $a\in\reals$ and $b\in\reals$ and $a \leq b$.
    Suppose $b - a \leq 1 / n$.
    Then $((a + a + b) / (1 + 1 + 1)) - a \leq 1 / (n + 1)$.
\end{lemma}
\begin{proof}
    We have $1 / n\in\reals$ and $1 / (n + 1)\in\reals$
        by \cref{naturalsplus_reciprocal_real,real_naturals_closed_succ}.
    Let $d = ((a + a + b) / (1 + 1 + 1)) - a$.
    $d\in\reals$ and $(d + d) + d \leq 1 / n$
        by \cref{real_sub,real_add_inverse,real_add_closed,real_div_naturalsplus_real,real_one_in_naturals,real_naturals_closed_succ,space_filling_left_third_width_triple,real_leq_subst_left_local,real_leq_trans}.
    $(d + d) + d \leq 1 / (n + 1) + 1 / (n + 1) + 1 / (n + 1)$
        by \cref{real_add_closed,real_leq_trans,space_filling_reciprocal_leq_triple_succ_reciprocal}.
    $((a + a + b) / (1 + 1 + 1)) - a \leq 1 / (n + 1)$ by \cref{space_filling_triple_leq_cancel}.
\end{proof}

% from §44 helper lemma 81: a right retained third improves the reciprocal width bound
\begin{lemma}\label{space_filling_right_third_width_leq_succ_reciprocal}
    Suppose $n\in\naturalsPlus$.
    Suppose $a\in\reals$ and $b\in\reals$ and $a \leq b$.
    Suppose $b - a \leq 1 / n$.
    Then $b - ((a + b + b) / (1 + 1 + 1)) \leq 1 / (n + 1)$.
\end{lemma}
\begin{proof}
    Follows by \cref{space_filling_retained_third_symmetric_difference,space_filling_left_third_width_leq_succ_reciprocal,real_leq_subst_left_local}.
\end{proof}

% from §44 helper lemma 82: retained-third state-orbit widths are bounded by reciprocal stage numbers
\begin{lemma}\label{space_filling_retained_third_state_orbit_width_bound}
    Let $I = \intervalclosed{0}{1}$.
    Let $D = \{0,1\}$.
    Let $E = \{(n,D) \mid n\in\naturalsPlus\}$.
    Assume $x\in\productFamily{E}$.
    Let $K = \{J\in\pow{I} \mid \text{there exist $a,b\in\reals$ such that
        $J = \intervalclosed{a}{b}$ and
        $a \leq b$}\}$.
    Let $G = \naturalsPlus\times K$.
    Let $R = \{p\in G\times G \mid \text{there exist $n,J,H,a,b,d$ such that
        $n\in\naturalsPlus$ and
        $\fst{p} = (n,J)$ and
        $\snd{p} = (n + 1,H)$ and
        $d = x(n)$ and
        $a\in\reals$ and
        $b\in\reals$ and
        $a \leq b$ and
        $J = \intervalclosed{a}{b}$ and
        $H\subseteq J$ and
        $(d,H)\in\{(0,\intervalclosed{a}{(a + a + b) / (1 + 1 + 1)}),
            (1,\intervalclosed{(a + b + b) / (1 + 1 + 1)}{b})\}$}\}$.
    Suppose $g\in\funs{G}{G}$.
    Suppose for all $z\in G$ we have $z\mathrel{R}\apply{g}{z}$.
    Suppose $u$ is a sequence in $G$.
    Suppose $u(1) = (1,I)$.
    Suppose for all $n\in\naturalsPlus$ we have $u(n + 1) = \apply{g}{u(n)}$.
    Then for all $n,a,b$ if
        $n\in\naturalsPlus$ and
        $a\in\reals$ and
        $b\in\reals$ and
        $\snd{u(n)} = \intervalclosed{a}{b}$ and
        $a \leq b$
    then $b - a \leq 1 / n$.
\end{lemma}
\begin{proof}
    Let $A_0 = \{n\in\naturalsPlus \mid \text{for all $a,b$ if
        $a\in\reals$ and
        $b\in\reals$ and
        $\snd{u(n)} = \intervalclosed{a}{b}$ and
        $a \leq b$
    then $b - a \leq 1 / n$}\}$.
    $A_0\subseteq\naturalsPlus$ by \cref{subseteq}.
    Show that for all $a,b$ if
        $a\in\reals$ and
        $b\in\reals$ and
        $\snd{u(1)} = \intervalclosed{a}{b}$ and
        $a \leq b$
    then $b - a \leq 1 / 1$.
    \begin{subproof}
        Follows by \cref{snd_eq,pair_eq_iff,subseteq,real_add_ident,real_one_in_naturals,real_naturals_subset_reals,real_zero_lt_one,real_lt_imp_leq,space_filling_intervalclosed_endpoints_unique,real_sub,real_add_eq_subst,real_neg_zero,real_add_comm,real_one_div_one}.
    \end{subproof}
    $1\in A_0$ by \cref{subseteq,real_one_in_naturals}.
    Show that for all $n\in A_0$ we have $n + 1\in A_0$.
    \begin{subproof}
        Fix $n\in A_0$.
        $n\in\naturalsPlus$ by \cref{subseteq}.
        Show that for all $a,b$ if
            $a\in\reals$ and
            $b\in\reals$ and
            $\snd{u(n + 1)} = \intervalclosed{a}{b}$ and
            $a \leq b$
        then $b - a \leq 1 / (n + 1)$.
        \begin{subproof}
            Fix $a$.
            Fix $b$.
            Assume $a\in\reals$ and $b\in\reals$ and $\snd{u(n + 1)} = \intervalclosed{a}{b}$ and $a \leq b$.
            We have $(u(n),u(n + 1))\in R$ by \cref{sequence_in,funs_type_apply,relation,subseteq}.
            Take $k,J,H,c,d,e$ such that
                $(u(n),u(n + 1))\in G\times G$ and
                $k\in\naturalsPlus$ and
                $\fst{(u(n),u(n + 1))} = (k,J)$ and
                $\snd{(u(n),u(n + 1))} = (k + 1,H)$ and
                $e = x(k)$ and
                $c\in\reals$ and
                $d\in\reals$ and
                $c \leq d$ and
                $J = \intervalclosed{c}{d}$ and
                $(e,H)\in\{(0,\intervalclosed{c}{(c + c + d) / (1 + 1 + 1)}),
                    (1,\intervalclosed{(c + d + d) / (1 + 1 + 1)}{d})\}$
                by \cref{subseteq}.
            $\snd{u(n)} = J$ and $\snd{u(n + 1)} = H$ and $\fst{u(n)} = k$
                by \cref{fst_eq,snd_eq,pair_eq_iff}.
            $k = n$ by \cref{space_filling_retained_third_state_orbit_index,subseteq}.
            $d - c \leq 1 / n$ by \cref{subseteq}.
            $(e,H) = (0,\intervalclosed{c}{(c + c + d) / (1 + 1 + 1)})$ or
                $(e,H) = (1,\intervalclosed{(c + d + d) / (1 + 1 + 1)}{d})$
                by \cref{upair_iff}.
            \begin{byCase}
                \caseOf{$(e,H) = (0,\intervalclosed{c}{(c + c + d) / (1 + 1 + 1)})$.}
                    We have $b - a \leq 1 / (n + 1)$
                        by \cref{pair_eq_iff,subseteq,real_add_closed,real_div_naturalsplus_real,real_one_in_naturals,real_naturals_closed_succ,space_filling_left_endpoint_leq_left_third,space_filling_intervalclosed_endpoints_unique,real_sub,real_add_eq_subst,space_filling_left_third_width_leq_succ_reciprocal}.
                \caseOf{$(e,H) = (1,\intervalclosed{(c + d + d) / (1 + 1 + 1)}{d})$.}
                    We have $b - a \leq 1 / (n + 1)$
                        by \cref{pair_eq_iff,subseteq,real_add_closed,real_div_naturalsplus_real,real_one_in_naturals,real_naturals_closed_succ,space_filling_right_third_leq_right_endpoint,space_filling_intervalclosed_endpoints_unique,real_sub,real_add_eq_subst,space_filling_right_third_width_leq_succ_reciprocal}.
            \end{byCase}
        \end{subproof}
        $n + 1\in A_0$ by \cref{subseteq,real_naturals_closed_succ}.
    \end{subproof}
    For all $n\in\naturalsPlus$ we have $n\in A_0$ by \cref{real_naturals_induction}.
    For all $n,a,b$ if
        $n\in\naturalsPlus$ and
        $a\in\reals$ and
        $b\in\reals$ and
        $\snd{u(n)} = \intervalclosed{a}{b}$ and
        $a \leq b$
    then $b - a \leq 1 / n$ by \cref{subseteq}.
\end{proof}

% from §44 helper lemma 83: a retained-third state orbit has at most one common point
\begin{lemma}\label{space_filling_retained_third_state_orbit_common_point_unique}
    Let $I = \intervalclosed{0}{1}$.
    Let $D = \{0,1\}$.
    Let $E = \{(n,D) \mid n\in\naturalsPlus\}$.
    Assume $x\in\productFamily{E}$.
    Let $K = \{J\in\pow{I} \mid \text{there exist $a,b\in\reals$ such that
        $J = \intervalclosed{a}{b}$ and
        $a \leq b$}\}$.
    Let $G = \naturalsPlus\times K$.
    Let $R = \{p\in G\times G \mid \text{there exist $n,J,H,a,b,d$ such that
        $n\in\naturalsPlus$ and
        $\fst{p} = (n,J)$ and
        $\snd{p} = (n + 1,H)$ and
        $d = x(n)$ and
        $a\in\reals$ and
        $b\in\reals$ and
        $a \leq b$ and
        $J = \intervalclosed{a}{b}$ and
        $H\subseteq J$ and
        $(d,H)\in\{(0,\intervalclosed{a}{(a + a + b) / (1 + 1 + 1)}),
            (1,\intervalclosed{(a + b + b) / (1 + 1 + 1)}{b})\}$}\}$.
    Suppose $g\in\funs{G}{G}$.
    Suppose for all $z\in G$ we have $z\mathrel{R}\apply{g}{z}$.
    Suppose $u$ is a sequence in $G$.
    Suppose $u(1) = (1,I)$.
    Suppose for all $n\in\naturalsPlus$ we have $u(n + 1) = \apply{g}{u(n)}$.
    Suppose $y_1\in I$ and for all $n\in\naturalsPlus$ we have $y_1\in\snd{u(n)}$.
    Suppose $y_2\in I$ and for all $n\in\naturalsPlus$ we have $y_2\in\snd{u(n)}$.
    Then $y_1 = y_2$.
\end{lemma}
\begin{proof}
    $y_1\in\reals$ and $y_2\in\reals$ by \cref{intervalclosed,subseteq}.
    Suppose not.
    $y_1 < y_2$ or $y_2 < y_1$ by \cref{real_lt_trichotomy}.
    \begin{byCase}
    \caseOf{$y_1 < y_2$.}
        We have $y_2 - y_1\in\reals$ and $0 < y_2 - y_1$
            by \cref{real_sub,real_add_inverse,real_add_closed,real_lt_imp_pos_diff}.
        Take $N\in\naturalsPlus$ such that $0 < 1 / N$ and $1 / N < y_2 - y_1$
            by \cref{subseteq,real_small_reciprocal}.
        Take $m,J$ such that
            $u(N) = (m,J)$ and
            $J\in K$
            by \cref{sequence_in,funs_type_apply,subseteq,times_elem_is_tuple}.
        Take $a,b$ such that
            $a\in\reals$ and
            $b\in\reals$ and
            $J = \intervalclosed{a}{b}$ and
            $a \leq b$
            by \cref{subseteq}.
        $\snd{u(N)} = \intervalclosed{a}{b}$ by \cref{snd_eq,pair_eq_iff,subseteq}.
        $y_1\in\intervalclosed{a}{b}$ and $y_2\in\intervalclosed{a}{b}$ by \cref{subseteq}.
        $y_2 - y_1 \leq b - a$
            by \cref{real_lt_imp_leq,space_filling_interval_width_bounds_difference}.
        $b - a \leq 1 / N$ by \cref{space_filling_retained_third_state_orbit_width_bound}.
        We have $b - a\in\reals$ and $1 / N\in\reals$
            by \cref{real_sub,real_add_inverse,real_add_closed,naturalsplus_reciprocal_real}.
        $y_2 - y_1 \leq 1 / N$ by \cref{real_leq_trans}.
        $1 / N < 1 / N$ by \cref{real_leq_split,real_lt_trans,real_lt_subst_right}.
        Contradiction by \cref{real_lt_irreflexive}.
    \caseOf{$y_2 < y_1$.}
        We have $y_1 - y_2\in\reals$ and $0 < y_1 - y_2$
            by \cref{real_sub,real_add_inverse,real_add_closed,real_lt_imp_pos_diff}.
        Take $N\in\naturalsPlus$ such that $0 < 1 / N$ and $1 / N < y_1 - y_2$
            by \cref{subseteq,real_small_reciprocal}.
        Take $m,J$ such that
            $u(N) = (m,J)$ and
            $J\in K$
            by \cref{sequence_in,funs_type_apply,subseteq,times_elem_is_tuple}.
        Take $a,b$ such that
            $a\in\reals$ and
            $b\in\reals$ and
            $J = \intervalclosed{a}{b}$ and
            $a \leq b$
            by \cref{subseteq}.
        $\snd{u(N)} = \intervalclosed{a}{b}$ by \cref{snd_eq,pair_eq_iff,subseteq}.
        $y_1\in\intervalclosed{a}{b}$ and $y_2\in\intervalclosed{a}{b}$ by \cref{subseteq}.
        $y_1 - y_2 \leq b - a$
            by \cref{real_lt_imp_leq,space_filling_interval_width_bounds_difference}.
        $b - a \leq 1 / N$ by \cref{space_filling_retained_third_state_orbit_width_bound}.
        We have $b - a\in\reals$ and $1 / N\in\reals$
            by \cref{real_sub,real_add_inverse,real_add_closed,naturalsplus_reciprocal_real}.
        $y_1 - y_2 \leq 1 / N$ by \cref{real_leq_trans}.
        $1 / N < 1 / N$ by \cref{real_leq_split,real_lt_trans,real_lt_subst_right}.
        Contradiction by \cref{real_lt_irreflexive}.
\end{byCase}
\end{proof}

% from §44 helper lemma 90: a binary sequence has at most one retained-third code point
\begin{lemma}\label{space_filling_retained_third_code_point_unique}
    Let $I = \intervalclosed{0}{1}$.
    Let $D = \{0,1\}$.
    Let $E = \{(n,D) \mid n\in\naturalsPlus\}$.
    Assume $x\in\productFamily{E}$.
    Suppose $y_1$ is a retained-third code point of $x$ in $I$ and $E$.
    Suppose $y_2$ is a retained-third code point of $x$ in $I$ and $E$.
    Then $y_1 = y_2$.
\end{lemma}
\begin{proof}
    We have $y_1\in I$ by \cref{space_filling_retained_third_code_point}.
    Take $K_1,G_1,R_1,g_1,u_1$ such that
        $K_1 = \{J\in\pow{I} \mid \text{there exist $a,b\in\reals$ such that
            $J = \intervalclosed{a}{b}$ and
            $a \leq b$}\}$ and
        $G_1 = \naturalsPlus\times K_1$ and
        $R_1 = \{p\in G_1\times G_1 \mid \text{there exist $n,J,H,a,b,d$ such that
            $n\in\naturalsPlus$ and
            $\fst{p} = (n,J)$ and
            $\snd{p} = (n + 1,H)$ and
            $d = x(n)$ and
            $a\in\reals$ and
            $b\in\reals$ and
            $a \leq b$ and
            $J = \intervalclosed{a}{b}$ and
            $H\subseteq J$ and
            $(d,H)\in\{(0,\intervalclosed{a}{(a + a + b) / (1 + 1 + 1)}),
                (1,\intervalclosed{(a + b + b) / (1 + 1 + 1)}{b})\}$}\}$ and
        $g_1$ is a retained-third state selector on $G_1$ and $R_1$ and
        $u_1$ is a retained-third state orbit of $g_1$ on $G_1$ from $I$ and
        for all $n\in\naturalsPlus$ we have $y_1\in\snd{u_1(n)}$
        by \cref{space_filling_retained_third_code_point}.
    $g_1\in\funs{G_1}{G_1}$ and
        for all $z\in G_1$ we have $z\mathrel{R_1}\apply{g_1}{z}$
        by \cref{space_filling_retained_third_state_selector_on}.
    $u_1$ is a sequence in $G_1$ and
        $u_1(1) = (1,I)$ and
        for all $n\in\naturalsPlus$ we have $u_1(n + 1) = \apply{g_1}{u_1(n)}$
        by \cref{space_filling_retained_third_state_orbit_from}.
    We have $y_2\in I$
        by \cref{space_filling_retained_third_code_point}.
    Take $K_2,G_2,R_2,g_2,u_2$ such that
        $K_2 = \{J\in\pow{I} \mid \text{there exist $a,b\in\reals$ such that
            $J = \intervalclosed{a}{b}$ and
            $a \leq b$}\}$ and
        $G_2 = \naturalsPlus\times K_2$ and
        $R_2 = \{p\in G_2\times G_2 \mid \text{there exist $n,J,H,a,b,d$ such that
            $n\in\naturalsPlus$ and
            $\fst{p} = (n,J)$ and
            $\snd{p} = (n + 1,H)$ and
            $d = x(n)$ and
            $a\in\reals$ and
            $b\in\reals$ and
            $a \leq b$ and
            $J = \intervalclosed{a}{b}$ and
            $H\subseteq J$ and
            $(d,H)\in\{(0,\intervalclosed{a}{(a + a + b) / (1 + 1 + 1)}),
                (1,\intervalclosed{(a + b + b) / (1 + 1 + 1)}{b})\}$}\}$ and
        $g_2$ is a retained-third state selector on $G_2$ and $R_2$ and
        $u_2$ is a retained-third state orbit of $g_2$ on $G_2$ from $I$ and
        for all $n\in\naturalsPlus$ we have $y_2\in\snd{u_2(n)}$
        by \cref{space_filling_retained_third_code_point}.
    $g_2\in\funs{G_2}{G_2}$ and
        for all $z\in G_2$ we have $z\mathrel{R_2}\apply{g_2}{z}$
        by \cref{space_filling_retained_third_state_selector_on}.
    $u_2$ is a sequence in $G_2$ and
        $u_2(1) = (1,I)$ and
        for all $n\in\naturalsPlus$ we have $u_2(n + 1) = \apply{g_2}{u_2(n)}$
        by \cref{space_filling_retained_third_state_orbit_from}.
    $K_1 = K_2$ and $G_1 = G_2$ and $g_2\in\funs{G_1}{G_1}$
        by \cref{subseteq,subseteq_antisymmetric}.
    $R_1\subseteq R_2$ and $R_2\subseteq R_1$ by \cref{subseteq}.
    $R_1 = R_2$ by \cref{subseteq_antisymmetric}.
    $u_2$ is a sequence in $G_1$ and
        $u_2(1) = (1,I)$ and
        for all $n\in\naturalsPlus$ we have $u_2(n + 1) = \apply{g_2}{u_2(n)}$
        by \cref{subseteq}.
    $R_1$ is right-unique by \cref{space_filling_retained_third_state_relation_rightunique}.
    $R_1$ is a relation by \cref{relation,subseteq,times_elem_is_tuple}.
    For all $z,w_1,w_2$ such that $(z,w_1)\in R_1$ and $(z,w_2)\in R_1$ we have $w_1 = w_2$
        by \cref{rightunique}.
    Let $A_0 = \{n\in\naturalsPlus \mid u_1(n) = u_2(n)\}$.
    $1\in A_0$ by \cref{real_one_in_naturals,pair_eq_iff,subseteq}.
    Show that for all $n\in A_0$ we have $n + 1\in A_0$.
    \begin{subproof}
        Fix $n\in A_0$.
        We have $n\in\naturalsPlus$ and $u_1(n) = u_2(n)$ by \cref{subseteq}.
        $u_1(n)\in G_1$ and $u_2(n)\in G_1$ by \cref{sequence_in,funs_type_apply}.
        $u_1(n)\mathrel{R_1}\apply{g_1}{u_1(n)}$ and $u_2(n)\mathrel{R_2}\apply{g_2}{u_2(n)}$
            by assumption.
        $u_2(n)\mathrel{R_1}\apply{g_2}{u_2(n)}$ by \cref{subseteq}.
        $(u_1(n),u_1(n + 1))\in R_1$ and $(u_1(n),u_2(n + 1))\in R_1$
            by \cref{relation,subseteq}.
        $u_1(n + 1) = u_2(n + 1)$ by \cref{subseteq}.
        $n + 1\in A_0$ by \cref{subseteq,real_naturals_closed_succ}.
    \end{subproof}
    For all $n\in\naturalsPlus$ we have $u_1(n) = u_2(n)$
        by \cref{real_naturals_induction,subseteq}.
    For all $n\in\naturalsPlus$ we have $y_2\in\snd{u_1(n)}$ by \cref{snd_eq,pair_eq_iff,subseteq}.
    $y_1 = y_2$ by \cref{space_filling_retained_third_state_orbit_common_point_unique}.
\end{proof}

% from §44 helper definition 5: dyadic code point of a binary sequence
\begin{definition}\label{space_filling_dyadic_code_point}
    $y$ is a dyadic code point of $x$ in $I$ and $E$ iff
        $x\in\productFamily{E}$ and
        $y\in I$ and
        there exist $K,G,R,g,u$ such that
            $K = \{J\in\pow{I} \mid \text{there exist $a,b\in\reals$ such that
                $J = \intervalclosed{a}{b}$ and
                $a \leq b$}\}$ and
            $G = \naturalsPlus\times K$ and
            $R = \{p\in G\times G \mid \text{there exist $n,J,H,a,b,d$ such that
                $n\in\naturalsPlus$ and
                $\fst{p} = (n,J)$ and
                $\snd{p} = (n + 1,H)$ and
                $d = x(n)$ and
                $a\in\reals$ and
                $b\in\reals$ and
                $a \leq b$ and
                $J = \intervalclosed{a}{b}$ and
                $H\subseteq J$ and
                $(d,H)\in\{(0,\intervalclosed{a}{(a + b) / (1 + 1)}),
                    (1,\intervalclosed{(a + b) / (1 + 1)}{b})\}$}\}$ and
            $g$ is a retained-third state selector on $G$ and $R$ and
            $u$ is a retained-third state orbit of $g$ on $G$ from $I$ and
            for all $n\in\naturalsPlus$ we have $y\in\snd{u(n)}$.
\end{definition}

% from §44 helper lemma 91: every binary sequence has a dyadic code point
\begin{lemma}\label{space_filling_dyadic_code_point_exists}
    Let $I = \intervalclosed{0}{1}$.
    Let $D = \{0,1\}$.
    Let $E = \{(n,D) \mid n\in\naturalsPlus\}$.
    Assume $x\in\productFamily{E}$.
    Then there exists $y$ such that $y$ is a dyadic code point of $x$ in $I$ and $E$.
\end{lemma}
\begin{proof}
    Let $K = \{J\in\pow{I} \mid \text{there exist $a,b\in\reals$ such that
        $J = \intervalclosed{a}{b}$ and
        $a \leq b$}\}$.
    Let $G = \naturalsPlus\times K$.
    Let $R = \{p\in G\times G \mid \text{there exist $n,J,H,a,b,d$ such that
        $n\in\naturalsPlus$ and
        $\fst{p} = (n,J)$ and
        $\snd{p} = (n + 1,H)$ and
        $d = x(n)$ and
        $a\in\reals$ and
        $b\in\reals$ and
        $a \leq b$ and
        $J = \intervalclosed{a}{b}$ and
        $H\subseteq J$ and
        $(d,H)\in\{(0,\intervalclosed{a}{(a + b) / (1 + 1)}),
            (1,\intervalclosed{(a + b) / (1 + 1)}{b})\}$}\}$.
    Take $g$ such that
        $g\in\funs{G}{G}$ and
        for all $z\in G$ we have $z\mathrel{R}\apply{g}{z}$
        by \cref{space_filling_dyadic_child_state_selector}.
    Take $u$ such that
        $u$ is a sequence in $G$ and
        $u(1) = (1,I)$ and
        for all $n\in\naturalsPlus$ we have $u(n + 1) = \apply{g}{u(n)}$
        by \cref{space_filling_dyadic_child_state_iteration}.
    Take $y\in I$ such that
        for all $n\in\naturalsPlus$ we have $y\in\snd{u(n)}$
        by \cref{space_filling_dyadic_child_state_orbit_common_point}.
    $y$ is a dyadic code point of $x$ in $I$ and $E$
        by \cref{space_filling_retained_third_state_selector_on,space_filling_retained_third_state_orbit_from,space_filling_dyadic_code_point}.
\end{proof}

% from §44 helper lemma 92: a binary sequence has at most one dyadic code point
\begin{lemma}\label{space_filling_dyadic_code_point_unique}
    Let $I = \intervalclosed{0}{1}$.
    Let $D = \{0,1\}$.
    Let $E = \{(n,D) \mid n\in\naturalsPlus\}$.
    Assume $x\in\productFamily{E}$.
    Suppose $y_1$ is a dyadic code point of $x$ in $I$ and $E$.
    Suppose $y_2$ is a dyadic code point of $x$ in $I$ and $E$.
    Then $y_1 = y_2$.
\end{lemma}
\begin{proof}
    We have $y_1\in I$ by \cref{space_filling_dyadic_code_point}.
    Take $K_1,G_1,R_1,g_1,u_1$ such that
        $K_1 = \{J\in\pow{I} \mid \text{there exist $a,b\in\reals$ such that
            $J = \intervalclosed{a}{b}$ and
            $a \leq b$}\}$ and
        $G_1 = \naturalsPlus\times K_1$ and
        $R_1 = \{p\in G_1\times G_1 \mid \text{there exist $n,J,H,a,b,d$ such that
            $n\in\naturalsPlus$ and
            $\fst{p} = (n,J)$ and
            $\snd{p} = (n + 1,H)$ and
            $d = x(n)$ and
            $a\in\reals$ and
            $b\in\reals$ and
            $a \leq b$ and
            $J = \intervalclosed{a}{b}$ and
            $H\subseteq J$ and
            $(d,H)\in\{(0,\intervalclosed{a}{(a + b) / (1 + 1)}),
                (1,\intervalclosed{(a + b) / (1 + 1)}{b})\}$}\}$ and
        $g_1$ is a retained-third state selector on $G_1$ and $R_1$ and
        $u_1$ is a retained-third state orbit of $g_1$ on $G_1$ from $I$ and
        for all $n\in\naturalsPlus$ we have $y_1\in\snd{u_1(n)}$
        by \cref{space_filling_dyadic_code_point}.
    $g_1\in\funs{G_1}{G_1}$ and
        for all $z\in G_1$ we have $z\mathrel{R_1}\apply{g_1}{z}$
        by \cref{space_filling_retained_third_state_selector_on}.
    $u_1$ is a sequence in $G_1$ and
        $u_1(1) = (1,I)$ and
        for all $n\in\naturalsPlus$ we have $u_1(n + 1) = \apply{g_1}{u_1(n)}$
        by \cref{space_filling_retained_third_state_orbit_from}.
    We have $y_2\in I$
        by \cref{space_filling_dyadic_code_point}.
    Take $K_2,G_2,R_2,g_2,u_2$ such that
        $K_2 = \{J\in\pow{I} \mid \text{there exist $a,b\in\reals$ such that
            $J = \intervalclosed{a}{b}$ and
            $a \leq b$}\}$ and
        $G_2 = \naturalsPlus\times K_2$ and
        $R_2 = \{p\in G_2\times G_2 \mid \text{there exist $n,J,H,a,b,d$ such that
            $n\in\naturalsPlus$ and
            $\fst{p} = (n,J)$ and
            $\snd{p} = (n + 1,H)$ and
            $d = x(n)$ and
            $a\in\reals$ and
            $b\in\reals$ and
            $a \leq b$ and
            $J = \intervalclosed{a}{b}$ and
            $H\subseteq J$ and
            $(d,H)\in\{(0,\intervalclosed{a}{(a + b) / (1 + 1)}),
                (1,\intervalclosed{(a + b) / (1 + 1)}{b})\}$}\}$ and
        $g_2$ is a retained-third state selector on $G_2$ and $R_2$ and
        $u_2$ is a retained-third state orbit of $g_2$ on $G_2$ from $I$ and
        for all $n\in\naturalsPlus$ we have $y_2\in\snd{u_2(n)}$
        by \cref{space_filling_dyadic_code_point}.
    $g_2\in\funs{G_2}{G_2}$ and
        for all $z\in G_2$ we have $z\mathrel{R_2}\apply{g_2}{z}$
        by \cref{space_filling_retained_third_state_selector_on}.
    $u_2$ is a sequence in $G_2$ and
        $u_2(1) = (1,I)$ and
        for all $n\in\naturalsPlus$ we have $u_2(n + 1) = \apply{g_2}{u_2(n)}$
        by \cref{space_filling_retained_third_state_orbit_from}.
    $K_1 = K_2$ and $G_1 = G_2$ and $g_2\in\funs{G_1}{G_1}$
        by \cref{subseteq,subseteq_antisymmetric}.
    $u_2$ is a sequence in $G_1$ and
        $u_2(1) = (1,I)$ and
        for all $n\in\naturalsPlus$ we have $u_2(n + 1) = \apply{g_2}{u_2(n)}$
        by \cref{subseteq}.
    For all $z\in G_1$ we have $z\mathrel{R_2}\apply{g_2}{z}$ by \cref{subseteq}.
    $R_1\subseteq R_2$ and $R_2\subseteq R_1$ by \cref{subseteq}.
    $R_1 = R_2$ by \cref{subseteq,subseteq_antisymmetric}.
    $R_1$ is a relation by \cref{relation,subseteq,times_elem_is_tuple}.
    $u_1\in\funs{\naturalsPlus}{G_1}$ and $u_2\in\funs{\naturalsPlus}{G_1}$ by \cref{sequence_in}.
    For all $n,u,v$ such that $(n,u)\in E$ and $(n,v)\in E$ we have $u = v$ by \cref{pair_eq_iff}.
    $E$ is a function on $\naturalsPlus$
        by \cref{relation,rightunique,dom_intro,dom_iff,pair_eq_iff,subseteq,subseteq_antisymmetric}.
    $\dom{E} = \naturalsPlus$ by \cref{funs}.
    For all $n\in\naturalsPlus$ we have $E(n) = D$ by \cref{function_apply_intro}.
    Show that for all $z\in G_1$ we have $\apply{g_1}{z} = \apply{g_2}{z}$.
    \begin{subproof}
        Fix $z\in G_1$.
        We have $(z,\apply{g_1}{z})\in R_1$ by \cref{subseteq,relation}.
        $(z,\apply{g_2}{z})\in R_2$ by \cref{subseteq,relation}.
        $(z,\apply{g_2}{z})\in R_1$ by \cref{subseteq}.
        Take $n_1,J_1,H_1,a_1,b_1,d_1$ such that
            $(z,\apply{g_1}{z})\in G_1\times G_1$ and
            $n_1\in\naturalsPlus$ and
            $\fst{(z,\apply{g_1}{z})} = (n_1,J_1)$ and
            $\snd{(z,\apply{g_1}{z})} = (n_1 + 1,H_1)$ and
            $d_1 = x(n_1)$ and
            $a_1\in\reals$ and
            $b_1\in\reals$ and
            $a_1 \leq b_1$ and
            $J_1 = \intervalclosed{a_1}{b_1}$ and
            $H_1\subseteq J_1$ and
            $(d_1,H_1)\in\{(0,\intervalclosed{a_1}{(a_1 + b_1) / (1 + 1)}),
                (1,\intervalclosed{(a_1 + b_1) / (1 + 1)}{b_1})\}$
            by \cref{subseteq}.
        Take $n_2,J_2,H_2,a_2,b_2,d_2$ such that
            $(z,\apply{g_2}{z})\in G_1\times G_1$ and
            $n_2\in\naturalsPlus$ and
            $\fst{(z,\apply{g_2}{z})} = (n_2,J_2)$ and
            $\snd{(z,\apply{g_2}{z})} = (n_2 + 1,H_2)$ and
            $d_2 = x(n_2)$ and
            $a_2\in\reals$ and
            $b_2\in\reals$ and
            $a_2 \leq b_2$ and
            $J_2 = \intervalclosed{a_2}{b_2}$ and
            $H_2\subseteq J_2$ and
            $(d_2,H_2)\in\{(0,\intervalclosed{a_2}{(a_2 + b_2) / (1 + 1)}),
                (1,\intervalclosed{(a_2 + b_2) / (1 + 1)}{b_2})\}$
            by \cref{subseteq}.
        $z = (n_1,J_1)$ and $\apply{g_1}{z} = (n_1 + 1,H_1)$
            by \cref{fst_eq,snd_eq,pair_eq_iff}.
        $z = (n_2,J_2)$ and $\apply{g_2}{z} = (n_2 + 1,H_2)$
            by \cref{fst_eq,snd_eq,pair_eq_iff}.
        $n_1 = n_2$ and $J_1 = J_2$ by \cref{pair_eq_iff}.
        $d_1 = d_2$ and $d_1\in D$
            by \cref{indexed_product,function_apply_intro,subseteq}.
        $a_1 = a_2$ and $b_1 = b_2$
            by \cref{subseteq,space_filling_intervalclosed_endpoints_unique}.
        $H_1 = H_2$ by \cref{space_filling_dyadic_child_choice_unique}.
        $\apply{g_1}{z} = \apply{g_2}{z}$ by \cref{pair_eq_iff,real_add_eq_subst}.
    \end{subproof}
    Let $A_0 = \{m\in\naturalsPlus \mid u_1(m) = u_2(m)\}$.
    $A_0\subseteq\naturalsPlus$ by \cref{subseteq}.
    $1\in A_0$ by \cref{real_one_in_naturals,pair_eq_iff,subseteq}.
    Show that for all $m\in A_0$ we have $m + 1\in A_0$.
    \begin{subproof}
        Follows by \cref{subseteq,funs_type_apply,real_naturals_closed_succ}.
    \end{subproof}
    For all $m\in\naturalsPlus$ we have $m\in A_0$ by \cref{real_naturals_induction}.
    For all $n\in\naturalsPlus$ we have $\snd{u_1(n)} = \snd{u_2(n)}$ by \cref{subseteq}.
    For all $n\in\naturalsPlus$ we have $y_2\in\snd{u_1(n)}$ by \cref{subseteq}.
    $y_1\in\reals$ and $y_2\in\reals$ by \cref{intervalclosed,subseteq}.
    Suppose not.
    $y_1 < y_2$ or $y_2 < y_1$ by \cref{real_lt_trichotomy}.
    \begin{byCase}
    \caseOf{$y_1 < y_2$.}
        We have $y_2 - y_1\in\reals$ and $0 < y_2 - y_1$
            by \cref{real_sub,real_add_inverse,real_add_closed,real_lt_imp_pos_diff}.
        Take $N\in\naturalsPlus$ such that $0 < 1 / N$ and $1 / N < y_2 - y_1$
            by \cref{subseteq,real_small_reciprocal}.
        Take $m,J$ such that
            $u_1(N) = (m,J)$ and
            $J\in K_1$
            by \cref{sequence_in,funs_type_apply,subseteq,times_elem_is_tuple}.
        Take $a,b$ such that
            $a\in\reals$ and
            $b\in\reals$ and
            $J = \intervalclosed{a}{b}$ and
            $a \leq b$
            by \cref{subseteq}.
        $\snd{u_1(N)} = \intervalclosed{a}{b}$ by \cref{snd_eq,pair_eq_iff,subseteq}.
        $y_1\in\intervalclosed{a}{b}$ and $y_2\in\intervalclosed{a}{b}$ by \cref{subseteq}.
        $y_2 - y_1 \leq b - a$
            by \cref{real_lt_imp_leq,space_filling_interval_width_bounds_difference}.
        $b - a \leq 1 / N$ by \cref{space_filling_dyadic_child_state_orbit_width_bound}.
        We have $b - a\in\reals$ and $1 / N\in\reals$
            by \cref{real_sub,real_add_inverse,real_add_closed,naturalsplus_reciprocal_real}.
        $y_2 - y_1 \leq 1 / N$ by \cref{real_leq_trans}.
        $1 / N < 1 / N$ by \cref{real_leq_split,real_lt_trans,real_lt_subst_right}.
        Contradiction by \cref{real_lt_irreflexive}.
    \caseOf{$y_2 < y_1$.}
        We have $y_1 - y_2\in\reals$ and $0 < y_1 - y_2$
            by \cref{real_sub,real_add_inverse,real_add_closed,real_lt_imp_pos_diff}.
        Take $N\in\naturalsPlus$ such that $0 < 1 / N$ and $1 / N < y_1 - y_2$
            by \cref{subseteq,real_small_reciprocal}.
        Take $m,J$ such that
            $u_1(N) = (m,J)$ and
            $J\in K_1$
            by \cref{sequence_in,funs_type_apply,subseteq,times_elem_is_tuple}.
        Take $a,b$ such that
            $a\in\reals$ and
            $b\in\reals$ and
            $J = \intervalclosed{a}{b}$ and
            $a \leq b$
            by \cref{subseteq}.
        $\snd{u_1(N)} = \intervalclosed{a}{b}$ by \cref{snd_eq,pair_eq_iff,subseteq}.
        $y_1\in\intervalclosed{a}{b}$ and $y_2\in\intervalclosed{a}{b}$ by \cref{subseteq}.
        $y_1 - y_2 \leq b - a$
            by \cref{real_lt_imp_leq,space_filling_interval_width_bounds_difference}.
        $b - a \leq 1 / N$ by \cref{space_filling_dyadic_child_state_orbit_width_bound}.
        We have $b - a\in\reals$ and $1 / N\in\reals$
            by \cref{real_sub,real_add_inverse,real_add_closed,naturalsplus_reciprocal_real}.
        $y_1 - y_2 \leq 1 / N$ by \cref{real_leq_trans}.
        $1 / N < 1 / N$ by \cref{real_leq_split,real_lt_trans,real_lt_subst_right}.
        Contradiction by \cref{real_lt_irreflexive}.
    \end{byCase}
\end{proof}

% from §44 helper lemma 93: the retained-third code-point assignment is a function
\begin{lemma}\label{space_filling_retained_third_code_map_exists}
    Let $I = \intervalclosed{0}{1}$.
    Let $D = \{0,1\}$.
    Let $E = \{(n,D) \mid n\in\naturalsPlus\}$.
    Then there exists $r\in\funs{\productFamily{E}}{I}$ such that
        for all $x\in\productFamily{E}$ we have $r(x)$ is a retained-third code point of $x$ in $I$ and $E$.
\end{lemma}
\begin{proof}
    Let $Q = \{w\in\productFamily{E}\times I \mid \text{$\snd{w}$ is a retained-third code point of $\fst{w}$ in $I$ and $E$}\}$.
    For all $x\in\productFamily{E}$ there exists $y$ such that $x\mathrel{Q} y$
        by \cref{space_filling_retained_third_code_point_exists,space_filling_retained_third_code_point,times_tuple_intro,fst_eq,snd_eq,pair_eq_iff,relation,subseteq}.
    Take $r$ such that
        $r$ is a function on $\productFamily{E}$ and
        for all $x\in\productFamily{E}$ we have $x\mathrel{Q}\apply{r}{x}$
        by \cref{choice_function,relation,subseteq,times_elem_is_tuple}.
    $r\in\funs{\productFamily{E}}{I}$ by \cref{choice_function,funs_intro,relation,subseteq,times_tuple_elim}.
    For all $x\in\productFamily{E}$ we have $r(x)$ is a retained-third code point of $x$ in $I$ and $E$
        by \cref{choice_function,relation,subseteq,fst_eq,snd_eq,pair_eq_iff}.
\end{proof}

% from §44 helper lemma 94: the dyadic code-point assignment is a function
\begin{lemma}\label{space_filling_dyadic_code_map_exists}
    Let $I = \intervalclosed{0}{1}$.
    Let $D = \{0,1\}$.
    Let $E = \{(n,D) \mid n\in\naturalsPlus\}$.
    Then there exists $r\in\funs{\productFamily{E}}{I}$ such that
        for all $x\in\productFamily{E}$ we have $r(x)$ is a dyadic code point of $x$ in $I$ and $E$.
\end{lemma}
\begin{proof}
    Let $Q = \{w\in\productFamily{E}\times I \mid \text{$\snd{w}$ is a dyadic code point of $\fst{w}$ in $I$ and $E$}\}$.
    For all $x\in\productFamily{E}$ there exists $y$ such that $x\mathrel{Q} y$
        by \cref{space_filling_dyadic_code_point_exists,space_filling_dyadic_code_point,times_tuple_intro,fst_eq,snd_eq,pair_eq_iff,relation,subseteq}.
    Take $r$ such that
        $r$ is a function on $\productFamily{E}$ and
        for all $x\in\productFamily{E}$ we have $x\mathrel{Q}\apply{r}{x}$
        by \cref{choice_function,relation,subseteq,times_elem_is_tuple}.
    $r\in\funs{\productFamily{E}}{I}$ by \cref{choice_function,funs_intro,relation,subseteq,times_tuple_elim}.
    For all $x\in\productFamily{E}$ we have $r(x)$ is a dyadic code point of $x$ in $I$ and $E$
        by \cref{choice_function,relation,subseteq,fst_eq,snd_eq,pair_eq_iff}.
\end{proof}

% from §44 helper lemma 95: every chosen dyadic child interval determines a binary digit
\begin{lemma}\label{space_filling_point_interval_child_has_digit}
    Let $I = \intervalclosed{0}{1}$.
    Let $D = \{0,1\}$.
    Assume $y\in I$.
    Let $K = \{J\in\pow{I} \mid \text{there exist $a,b\in\reals$ such that
        $J = \intervalclosed{a}{b}$ and
        $a \leq b$ and
        $y\in J$}\}$.
    Let $R = \{w\in K\times K \mid \text{there exist $a,b$ such that
        $a\in\reals$ and
        $b\in\reals$ and
        $a \leq b$ and
        $\fst{w} = \intervalclosed{a}{b}$ and
        $\snd{w}\subseteq\fst{w}$ and
        $\snd{w}\in\{\intervalclosed{a}{(a + b) / (1 + 1)}, \intervalclosed{(a + b) / (1 + 1)}{b}\}$}\}$.
    Suppose $J\in K$.
    Suppose $J\mathrel{R} H$.
    Then there exist $a,b,d$ such that
        $a\in\reals$ and
        $b\in\reals$ and
        $a \leq b$ and
        $J = \intervalclosed{a}{b}$ and
        $d\in D$ and
        $(d,H)\in\{(0,\intervalclosed{a}{(a + b) / (1 + 1)}),
            (1,\intervalclosed{(a + b) / (1 + 1)}{b})\}$.
\end{lemma}
\begin{proof}
    We have $(J,H)\in R$ by \cref{relation,subseteq,times_elem_is_tuple}.
    Take $a,b$ such that
        $(J,H)\in K\times K$ and
        $a\in\reals$ and
        $b\in\reals$ and
        $a \leq b$ and
        $\fst{(J,H)} = \intervalclosed{a}{b}$ and
        $\snd{(J,H)}\subseteq\fst{(J,H)}$ and
        $\snd{(J,H)}\in\{\intervalclosed{a}{(a + b) / (1 + 1)}, \intervalclosed{(a + b) / (1 + 1)}{b}\}$
        by \cref{subseteq}.
    $J = \intervalclosed{a}{b}$ and
        $H\in\{\intervalclosed{a}{(a + b) / (1 + 1)}, \intervalclosed{(a + b) / (1 + 1)}{b}\}$
        by \cref{fst_eq,snd_eq,pair_eq_iff}.
    $H = \intervalclosed{a}{(a + b) / (1 + 1)}$ or
        $H = \intervalclosed{(a + b) / (1 + 1)}{b}$
        by \cref{upair_iff}.
    There exists $d\in D$ such that
        $(d,H)\in\{(0,\intervalclosed{a}{(a + b) / (1 + 1)}),
            (1,\intervalclosed{(a + b) / (1 + 1)}{b})\}$
        by \cref{upair_intro_left,upair_intro_right,pair_eq_iff,subseteq,upair_iff}.
    There exist $c,e,f$ such that
        $c\in\reals$ and
        $e\in\reals$ and
        $c \leq e$ and
        $J = \intervalclosed{c}{e}$ and
        $f\in D$ and
        $(f,H)\in\{(0,\intervalclosed{c}{(c + e) / (1 + 1)}),
            (1,\intervalclosed{(c + e) / (1 + 1)}{e})\}$
        by \cref{subseteq}.
\end{proof}

% from §44 helper lemma 96: the left retained-third endpoint is strictly above the left endpoint
\begin{lemma}\label{space_filling_left_endpoint_lt_left_third}
    Suppose $a\in\reals$ and $b\in\reals$ and $a < b$.
    Then $a < (a + a + b) / (1 + 1 + 1)$.
\end{lemma}
\begin{proof}
    We have $a + a + a < a + a + b$
        by \cref{real_lt_add,real_add_closed,real_add_assoc,real_add_comm,real_lt_subst_right}.
    $a + a + a\in\reals$ and $a + a + b\in\reals$ by \cref{real_add_closed}.
    $1 + 1 + 1\in\reals$ and $0 < 1 + 1 + 1$ and $1 + 1 + 1\neq 0$ and
        $\inv{(1 + 1 + 1)}\in\reals$ and $0 < \inv{(1 + 1 + 1)}$
        by \cref{space_filling_three_in_naturalsplus,real_naturals_subset_reals,subseteq,naturalsplus_positive_real,real_pos_nonzero,real_mul_inverse,real_inv_pos}.
    $(a + a + a) \rmul \inv{(1 + 1 + 1)} < (a + a + b) \rmul \inv{(1 + 1 + 1)}$
        by \cref{real_lt_mul_pos}.
    $(a + a + a) / (1 + 1 + 1) < (a + a + b) / (1 + 1 + 1)$
        by \cref{real_div,real_lt_subst_left,real_lt_subst_right}.
    $a < (a + a + b) / (1 + 1 + 1)$ by \cref{space_filling_triple_third_real,real_lt_subst_left}.
\end{proof}

% from §44 helper lemma 97: the right retained-third endpoint is strictly below the right endpoint
\begin{lemma}\label{space_filling_right_third_lt_right_endpoint}
    Suppose $a\in\reals$ and $b\in\reals$ and $a < b$.
    Then $(a + b + b) / (1 + 1 + 1) < b$.
\end{lemma}
\begin{proof}
    We have $a + b + b < b + b + b$ by \cref{real_add_closed,real_lt_add,real_add_assoc,real_lt_subst_right}.
    $a + b + b\in\reals$ and $b + b + b\in\reals$ by \cref{real_add_closed}.
    $1 + 1 + 1\in\reals$ and $0 < 1 + 1 + 1$ and $1 + 1 + 1\neq 0$ and
        $\inv{(1 + 1 + 1)}\in\reals$ and $0 < \inv{(1 + 1 + 1)}$
        by \cref{space_filling_three_in_naturalsplus,real_naturals_subset_reals,subseteq,naturalsplus_positive_real,real_pos_nonzero,real_mul_inverse,real_inv_pos}.
    $(a + b + b) \rmul \inv{(1 + 1 + 1)} < (b + b + b) \rmul \inv{(1 + 1 + 1)}$
        by \cref{real_lt_mul_pos}.
    $(a + b + b) / (1 + 1 + 1) < (b + b + b) / (1 + 1 + 1)$
        by \cref{real_div,real_lt_subst_left,real_lt_subst_right}.
    $(a + b + b) / (1 + 1 + 1) < b$ by \cref{space_filling_triple_third_real,real_lt_subst_right}.
\end{proof}

% from §44 helper lemma 98: retained-third orbit intervals are nondegenerate
\begin{lemma}\label{space_filling_retained_third_state_orbit_strict_interval}
    Let $I = \intervalclosed{0}{1}$.
    Let $D = \{0,1\}$.
    Let $E = \{(n,D) \mid n\in\naturalsPlus\}$.
    Assume $x\in\productFamily{E}$.
    Let $K = \{J\in\pow{I} \mid \text{there exist $a,b\in\reals$ such that
        $J = \intervalclosed{a}{b}$ and
        $a \leq b$}\}$.
    Let $G = \naturalsPlus\times K$.
    Let $R = \{p\in G\times G \mid \text{there exist $n,J,H,a,b,d$ such that
        $n\in\naturalsPlus$ and
        $\fst{p} = (n,J)$ and
        $\snd{p} = (n + 1,H)$ and
        $d = x(n)$ and
        $a\in\reals$ and
        $b\in\reals$ and
        $a \leq b$ and
        $J = \intervalclosed{a}{b}$ and
        $H\subseteq J$ and
        $(d,H)\in\{(0,\intervalclosed{a}{(a + a + b) / (1 + 1 + 1)}),
            (1,\intervalclosed{(a + b + b) / (1 + 1 + 1)}{b})\}$}\}$.
    Suppose $g\in\funs{G}{G}$.
    Suppose for all $z\in G$ we have $z\mathrel{R}\apply{g}{z}$.
    Suppose $u$ is a sequence in $G$.
    Suppose $u(1) = (1,I)$.
    Suppose for all $n\in\naturalsPlus$ we have $u(n + 1) = \apply{g}{u(n)}$.
    Then for all $n,a,b$ if
        $n\in\naturalsPlus$ and
        $a\in\reals$ and
        $b\in\reals$ and
        $\snd{u(n)} = \intervalclosed{a}{b}$ and
        $a \leq b$
    then $a < b$.
\end{lemma}
\begin{proof}
    Let $A_0 = \{n\in\naturalsPlus \mid \text{for all $a,b$ if
        $a\in\reals$ and
        $b\in\reals$ and
        $\snd{u(n)} = \intervalclosed{a}{b}$ and
        $a \leq b$
    then $a < b$}\}$.
    $A_0\subseteq\naturalsPlus$ by \cref{subseteq}.
    Show that for all $a,b$ if
        $a\in\reals$ and
        $b\in\reals$ and
        $\snd{u(1)} = \intervalclosed{a}{b}$ and
        $a \leq b$
    then $a < b$.
    \begin{subproof}
        Follows by \cref{snd_eq,pair_eq_iff,subseteq,space_filling_intervalclosed_endpoints_unique,real_add_ident,real_one_in_naturals,real_naturals_subset_reals,real_zero_lt_one,real_lt_imp_leq}.
    \end{subproof}
    $1\in A_0$ by \cref{real_one_in_naturals,subseteq}.
    Show that for all $n\in A_0$ we have $n + 1\in A_0$.
    \begin{subproof}
        Fix $n\in A_0$.
        $n\in\naturalsPlus$ by \cref{subseteq}.
        Show that for all $a,b$ if
            $a\in\reals$ and
            $b\in\reals$ and
            $\snd{u(n + 1)} = \intervalclosed{a}{b}$ and
            $a \leq b$
        then $a < b$.
        \begin{subproof}
            Fix $a$.
            Fix $b$.
            Assume $a\in\reals$ and $b\in\reals$ and $\snd{u(n + 1)} = \intervalclosed{a}{b}$ and $a \leq b$.
            We have $(u(n),u(n + 1))\in R$ by \cref{sequence_in,funs_type_apply,subseteq,relation,times_elem_is_tuple}.
            Take $k,J,H,c,d,e$ such that
                $(u(n),u(n + 1))\in G\times G$ and
                $k\in\naturalsPlus$ and
                $\fst{(u(n),u(n + 1))} = (k,J)$ and
                $\snd{(u(n),u(n + 1))} = (k + 1,H)$ and
                $e = x(k)$ and
                $c\in\reals$ and
                $d\in\reals$ and
                $c \leq d$ and
                $J = \intervalclosed{c}{d}$ and
                $H\subseteq J$ and
                $(e,H)\in\{(0,\intervalclosed{c}{(c + c + d) / (1 + 1 + 1)}),
                    (1,\intervalclosed{(c + d + d) / (1 + 1 + 1)}{d})\}$
                by \cref{subseteq}.
            $\fst{u(n)} = k$ and $\snd{u(n)} = J$ and $\snd{u(n + 1)} = H$
                by \cref{fst_eq,snd_eq,pair_eq_iff}.
            $k = n$ by \cref{space_filling_retained_third_state_orbit_index,subseteq}.
            $c < d$ by \cref{subseteq}.
            $(e,H) = (0,\intervalclosed{c}{(c + c + d) / (1 + 1 + 1)})$ or
                $(e,H) = (1,\intervalclosed{(c + d + d) / (1 + 1 + 1)}{d})$
                by \cref{upair_iff}.
            \begin{byCase}
            \caseOf{$(e,H) = (0,\intervalclosed{c}{(c + c + d) / (1 + 1 + 1)})$.}
                We have $H = \intervalclosed{c}{(c + c + d) / (1 + 1 + 1)}$ by \cref{pair_eq_iff}.
                $\intervalclosed{a}{b} = \intervalclosed{c}{(c + c + d) / (1 + 1 + 1)}$ by \cref{subseteq}.
                $a = c$ and $b = (c + c + d) / (1 + 1 + 1)$
                    by \cref{space_filling_left_endpoint_lt_left_third,real_add_closed,real_div_naturalsplus_real,real_one_in_naturals,real_naturals_closed_succ,real_lt_imp_leq,space_filling_intervalclosed_endpoints_unique}.
                $a < b$ by \cref{space_filling_left_endpoint_lt_left_third,real_lt_subst_left,real_lt_subst_right}.
            \caseOf{$(e,H) = (1,\intervalclosed{(c + d + d) / (1 + 1 + 1)}{d})$.}
                We have $H = \intervalclosed{(c + d + d) / (1 + 1 + 1)}{d}$ by \cref{pair_eq_iff}.
                $\intervalclosed{a}{b} = \intervalclosed{(c + d + d) / (1 + 1 + 1)}{d}$ by \cref{subseteq}.
                $a = (c + d + d) / (1 + 1 + 1)$ and $b = d$
                    by \cref{space_filling_right_third_lt_right_endpoint,real_add_closed,real_div_naturalsplus_real,real_one_in_naturals,real_naturals_closed_succ,real_lt_imp_leq,space_filling_intervalclosed_endpoints_unique}.
                $a < b$ by \cref{space_filling_right_third_lt_right_endpoint,real_lt_subst_left,real_lt_subst_right}.
            \end{byCase}
        \end{subproof}
        $n + 1\in A_0$ by \cref{subseteq,real_naturals_closed_succ}.
    \end{subproof}
    For all $n\in\naturalsPlus$ we have $n\in A_0$ by \cref{real_naturals_induction}.
    For all $n,a,b$ if
        $n\in\naturalsPlus$ and
        $a\in\reals$ and
        $b\in\reals$ and
        $\snd{u(n)} = \intervalclosed{a}{b}$ and
        $a \leq b$
    then $a < b$ by \cref{subseteq}.
\end{proof}

% from §44 helper lemma 99: retained-third state orbits with the same prefix have the same next interval
\begin{lemma}\label{space_filling_retained_third_state_orbit_prefix_agreement_succ}
    Let $I = \intervalclosed{0}{1}$.
    Let $D = \{0,1\}$.
    Let $E = \{(n,D) \mid n\in\naturalsPlus\}$.
    Assume $x_1\in\productFamily{E}$.
    Assume $x_2\in\productFamily{E}$.
    Let $K = \{J\in\pow{I} \mid \text{there exist $a,b\in\reals$ such that
        $J = \intervalclosed{a}{b}$ and
        $a \leq b$}\}$.
    Let $G = \naturalsPlus\times K$.
    Let $R_1 = \{p\in G\times G \mid \text{there exist $n,J,H,a,b,d$ such that
        $n\in\naturalsPlus$ and
        $\fst{p} = (n,J)$ and
        $\snd{p} = (n + 1,H)$ and
        $d = x_1(n)$ and
        $a\in\reals$ and
        $b\in\reals$ and
        $a \leq b$ and
        $J = \intervalclosed{a}{b}$ and
        $H\subseteq J$ and
        $(d,H)\in\{(0,\intervalclosed{a}{(a + a + b) / (1 + 1 + 1)}),
            (1,\intervalclosed{(a + b + b) / (1 + 1 + 1)}{b})\}$}\}$.
    Let $R_2 = \{p\in G\times G \mid \text{there exist $n,J,H,a,b,d$ such that
        $n\in\naturalsPlus$ and
        $\fst{p} = (n,J)$ and
        $\snd{p} = (n + 1,H)$ and
        $d = x_2(n)$ and
        $a\in\reals$ and
        $b\in\reals$ and
        $a \leq b$ and
        $J = \intervalclosed{a}{b}$ and
        $H\subseteq J$ and
        $(d,H)\in\{(0,\intervalclosed{a}{(a + a + b) / (1 + 1 + 1)}),
            (1,\intervalclosed{(a + b + b) / (1 + 1 + 1)}{b})\}$}\}$.
    Suppose $g_1\in\funs{G}{G}$.
    Suppose for all $z\in G$ we have $z\mathrel{R_1}\apply{g_1}{z}$.
    Suppose $u_1$ is a sequence in $G$.
    Suppose $u_1(1) = (1,I)$.
    Suppose for all $n\in\naturalsPlus$ we have $u_1(n + 1) = \apply{g_1}{u_1(n)}$.
    Suppose for all $n\in\naturalsPlus$ we have $\fst{u_1(n)} = n$.
    Suppose $g_2\in\funs{G}{G}$.
    Suppose for all $z\in G$ we have $z\mathrel{R_2}\apply{g_2}{z}$.
    Suppose $u_2$ is a sequence in $G$.
    Suppose $u_2(1) = (1,I)$.
    Suppose for all $n\in\naturalsPlus$ we have $u_2(n + 1) = \apply{g_2}{u_2(n)}$.
    Suppose for all $n\in\naturalsPlus$ we have $\fst{u_2(n)} = n$.
    Assume $N\in\naturalsPlus$.
    Assume for all $i\in\initialSegment{N}$ we have $x_1(i) = x_2(i)$.
    Then $\snd{u_1(N + 1)} = \snd{u_2(N + 1)}$.
\end{lemma}
\begin{proof}
    We have $E$ is a function on $\naturalsPlus$
        by \cref{relation,rightunique,dom_intro,dom_iff,pair_eq_iff,subseteq,subseteq_antisymmetric}.
    Hence $\dom{E} = \naturalsPlus$ by \cref{funs}.
    For all $n\in\naturalsPlus$ we have $E(n) = D$ by \cref{function_apply_intro}.
    We have for all $n\in\initialSegment{N}$ we have $\snd{u_1(n)} = \snd{u_2(n)}$
        by \cref{space_filling_retained_third_state_orbit_prefix_agreement}.
    $N\in\initialSegment{N}$ by \cref{initial_segment_naturalsplus}.
    Hence $\snd{u_1(N)} = \snd{u_2(N)}$ by \cref{subseteq}.
    $u_1\in\funs{\naturalsPlus}{G}$ and $u_2\in\funs{\naturalsPlus}{G}$ by \cref{sequence_in}.
    We have $(u_1(N),u_1(N + 1))\in R_1$ and
        $(u_2(N),u_2(N + 1))\in R_2$
        by \cref{funs_type_apply,real_naturals_closed_succ,subseteq,relation}.
    Take $k_1,J_1,H_1,a_1,b_1,d_1$ such that
        $(u_1(N),u_1(N + 1))\in G\times G$ and
        $k_1\in\naturalsPlus$ and
        $\fst{(u_1(N),u_1(N + 1))} = (k_1,J_1)$ and
        $\snd{(u_1(N),u_1(N + 1))} = (k_1 + 1,H_1)$ and
        $d_1 = x_1(k_1)$ and
        $a_1\in\reals$ and
        $b_1\in\reals$ and
        $a_1 \leq b_1$ and
        $J_1 = \intervalclosed{a_1}{b_1}$ and
        $H_1\subseteq J_1$ and
        $(d_1,H_1)\in\{(0,\intervalclosed{a_1}{(a_1 + a_1 + b_1) / (1 + 1 + 1)}),
            (1,\intervalclosed{(a_1 + b_1 + b_1) / (1 + 1 + 1)}{b_1})\}$
        by \cref{subseteq}.
    Take $k_2,J_2,H_2,a_2,b_2,d_2$ such that
        $(u_2(N),u_2(N + 1))\in G\times G$ and
        $k_2\in\naturalsPlus$ and
        $\fst{(u_2(N),u_2(N + 1))} = (k_2,J_2)$ and
        $\snd{(u_2(N),u_2(N + 1))} = (k_2 + 1,H_2)$ and
        $d_2 = x_2(k_2)$ and
        $a_2\in\reals$ and
        $b_2\in\reals$ and
        $a_2 \leq b_2$ and
        $J_2 = \intervalclosed{a_2}{b_2}$ and
        $H_2\subseteq J_2$ and
        $(d_2,H_2)\in\{(0,\intervalclosed{a_2}{(a_2 + a_2 + b_2) / (1 + 1 + 1)}),
            (1,\intervalclosed{(a_2 + b_2 + b_2) / (1 + 1 + 1)}{b_2})\}$
        by \cref{subseteq}.
    We have $k_1 = N$ and $k_2 = N$ and
        $\snd{u_1(N)} = J_1$ and $\snd{u_2(N)} = J_2$ and
        $\snd{u_1(N + 1)} = H_1$ and $\snd{u_2(N + 1)} = H_2$
        by \cref{fst_eq,snd_eq,pair_eq_iff,subseteq}.
    $d_1 = x_1(N)$ and $d_2 = x_2(N)$ and $J_1 = J_2$ by \cref{subseteq}.
    We have $x_1(N)\in D$ and $x_2(N)\in D$ by \cref{indexed_product,function_apply_intro,subseteq}.
    Hence $d_1 = d_2$ and $d_2\in D$ by \cref{subseteq}.
    $\intervalclosed{a_1}{b_1} = \intervalclosed{a_2}{b_2}$ by \cref{subseteq}.
    Therefore $a_1 = a_2$ and $b_1 = b_2$ by \cref{space_filling_intervalclosed_endpoints_unique}.
    We have $(d_2,H_1)\in\{(0,\intervalclosed{a_2}{(a_2 + a_2 + b_2) / (1 + 1 + 1)}),
        (1,\intervalclosed{(a_2 + b_2 + b_2) / (1 + 1 + 1)}{b_2})\}$
        by \cref{subseteq,real_add_eq_subst}.
    Thus $H_1 = H_2$ by \cref{space_filling_retained_third_child_choice_unique,subseteq,real_add_eq_subst}.
    Therefore $\snd{u_1(N + 1)} = \snd{u_2(N + 1)}$ by \cref{subseteq}.
\end{proof}

% from §44 helper lemma 100: opposite retained-third digits separate a shared stage interval
\begin{lemma}\label{space_filling_retained_third_shared_interval_digit_separation}
    Let $I = \intervalclosed{0}{1}$.
    Let $D = \{0,1\}$.
    Let $E = \{(n,D) \mid n\in\naturalsPlus\}$.
    Assume $x_1\in\productFamily{E}$.
    Assume $x_2\in\productFamily{E}$.
    Let $K = \{J\in\pow{I} \mid \text{there exist $c,d\in\reals$ such that
        $J = \intervalclosed{c}{d}$ and
        $c \leq d$}\}$.
    Let $G = \naturalsPlus\times K$.
    Let $R_1 = \{p\in G\times G \mid \text{there exist $n,J,H,c,d,e$ such that
        $n\in\naturalsPlus$ and
        $\fst{p} = (n,J)$ and
        $\snd{p} = (n + 1,H)$ and
        $e = x_1(n)$ and
        $c\in\reals$ and
        $d\in\reals$ and
        $c \leq d$ and
        $J = \intervalclosed{c}{d}$ and
        $H\subseteq J$ and
        $(e,H)\in\{(0,\intervalclosed{c}{(c + c + d) / (1 + 1 + 1)}),
            (1,\intervalclosed{(c + d + d) / (1 + 1 + 1)}{d})\}$}\}$.
    Let $R_2 = \{p\in G\times G \mid \text{there exist $n,J,H,c,d,e$ such that
        $n\in\naturalsPlus$ and
        $\fst{p} = (n,J)$ and
        $\snd{p} = (n + 1,H)$ and
        $e = x_2(n)$ and
        $c\in\reals$ and
        $d\in\reals$ and
        $c \leq d$ and
        $J = \intervalclosed{c}{d}$ and
        $H\subseteq J$ and
        $(e,H)\in\{(0,\intervalclosed{c}{(c + c + d) / (1 + 1 + 1)}),
            (1,\intervalclosed{(c + d + d) / (1 + 1 + 1)}{d})\}$}\}$.
    Suppose $g_1\in\funs{G}{G}$.
    Suppose for all $z\in G$ we have $z\mathrel{R_1}\apply{g_1}{z}$.
    Suppose $u_1$ is a sequence in $G$.
    Suppose $u_1(1) = (1,I)$.
    Suppose for all $n\in\naturalsPlus$ we have $u_1(n + 1) = \apply{g_1}{u_1(n)}$.
    Suppose for all $n\in\naturalsPlus$ we have $\fst{u_1(n)} = n$.
    Suppose $g_2\in\funs{G}{G}$.
    Suppose for all $z\in G$ we have $z\mathrel{R_2}\apply{g_2}{z}$.
    Suppose $u_2$ is a sequence in $G$.
    Suppose $u_2(1) = (1,I)$.
    Suppose for all $n\in\naturalsPlus$ we have $u_2(n + 1) = \apply{g_2}{u_2(n)}$.
    Suppose for all $n\in\naturalsPlus$ we have $\fst{u_2(n)} = n$.
    Suppose $y\in I$ and for all $n\in\naturalsPlus$ we have $y\in\snd{u_1(n)}$.
    Suppose for all $n\in\naturalsPlus$ we have $y\in\snd{u_2(n)}$.
    Assume $N\in\naturalsPlus$.
    Assume $a\in\reals$ and $b\in\reals$ and $a < b$.
    Assume $\snd{u_1(N + 1)} = \intervalclosed{a}{b}$.
    Assume $\snd{u_2(N + 1)} = \intervalclosed{a}{b}$.
    Assume $x_1(N + 1) = 0$.
    Assume $x_2(N + 1) = 1$.
    Then $y\notin I$.
\end{lemma}
\begin{proof}
    We have $u_1\in\funs{\naturalsPlus}{G}$ and $u_2\in\funs{\naturalsPlus}{G}$ and
        $N + 1\in\naturalsPlus$ and $(N + 1) + 1\in\naturalsPlus$
        by \cref{sequence_in,real_naturals_closed_succ}.
    $(u_1(N + 1),u_1((N + 1) + 1))\in R_1$ and
        $(u_2(N + 1),u_2((N + 1) + 1))\in R_2$
        by \cref{funs_type_apply,subseteq,relation}.
    Take $k_1,J_1,H_1,a_1,b_1,d_1$ such that
        $(u_1(N + 1),u_1((N + 1) + 1))\in G\times G$ and
        $k_1\in\naturalsPlus$ and
        $\fst{(u_1(N + 1),u_1((N + 1) + 1))} = (k_1,J_1)$ and
        $\snd{(u_1(N + 1),u_1((N + 1) + 1))} = (k_1 + 1,H_1)$ and
        $d_1 = x_1(k_1)$ and
        $a_1\in\reals$ and
        $b_1\in\reals$ and
        $a_1 \leq b_1$ and
        $J_1 = \intervalclosed{a_1}{b_1}$ and
        $H_1\subseteq J_1$ and
        $(d_1,H_1)\in\{(0,\intervalclosed{a_1}{(a_1 + a_1 + b_1) / (1 + 1 + 1)}),
            (1,\intervalclosed{(a_1 + b_1 + b_1) / (1 + 1 + 1)}{b_1})\}$
        by \cref{subseteq}.
    Take $k_2,J_2,H_2,a_2,b_2,d_2$ such that
        $(u_2(N + 1),u_2((N + 1) + 1))\in G\times G$ and
        $k_2\in\naturalsPlus$ and
        $\fst{(u_2(N + 1),u_2((N + 1) + 1))} = (k_2,J_2)$ and
        $\snd{(u_2(N + 1),u_2((N + 1) + 1))} = (k_2 + 1,H_2)$ and
        $d_2 = x_2(k_2)$ and
        $a_2\in\reals$ and
        $b_2\in\reals$ and
        $a_2 \leq b_2$ and
        $J_2 = \intervalclosed{a_2}{b_2}$ and
        $H_2\subseteq J_2$ and
        $(d_2,H_2)\in\{(0,\intervalclosed{a_2}{(a_2 + a_2 + b_2) / (1 + 1 + 1)}),
            (1,\intervalclosed{(a_2 + b_2 + b_2) / (1 + 1 + 1)}{b_2})\}$
        by \cref{subseteq}.
    We have $k_1 = N + 1$ and $k_2 = N + 1$ and
        $\snd{u_1(N + 1)} = J_1$ and $\snd{u_2(N + 1)} = J_2$ and
        $\snd{u_1((N + 1) + 1)} = H_1$ and $\snd{u_2((N + 1) + 1)} = H_2$
        by \cref{fst_eq,snd_eq,pair_eq_iff,subseteq}.
    $d_1 = x_1(N + 1)$ and $d_2 = x_2(N + 1)$ by \cref{subseteq}.
    We have $d_1 = 0$ and $d_2 = 1$ by \cref{subseteq}.
    $\intervalclosed{a_1}{b_1} = \intervalclosed{a}{b}$ and
        $\intervalclosed{a_2}{b_2} = \intervalclosed{a}{b}$ by \cref{subseteq}.
    We have $a_1 = a$ and $b_1 = b$
        by \cref{real_lt_imp_leq,space_filling_intervalclosed_endpoints_unique}.
    $a_2 = a$ and $b_2 = b$
        by \cref{real_lt_imp_leq,space_filling_intervalclosed_endpoints_unique}.
    We have $(d_1,H_1)\in\{(0,\intervalclosed{a}{(a + a + b) / (1 + 1 + 1)}),
        (1,\intervalclosed{(a + b + b) / (1 + 1 + 1)}{b})\}$
        by \cref{subseteq,real_add_eq_subst}.
    $(d_2,H_2)\in\{(0,\intervalclosed{a}{(a + a + b) / (1 + 1 + 1)}),
        (1,\intervalclosed{(a + b + b) / (1 + 1 + 1)}{b})\}$
        by \cref{subseteq,real_add_eq_subst}.
    Let $L = \intervalclosed{a}{(a + a + b) / (1 + 1 + 1)}$.
    Let $R = \intervalclosed{(a + b + b) / (1 + 1 + 1)}{b}$.
    We have $d_1\in D$ and $d_2\in D$ by \cref{upair_intro_left,upair_intro_right,subseteq}.
    $(d_1,L)\in\{(0,\intervalclosed{a}{(a + a + b) / (1 + 1 + 1)}),
        (1,\intervalclosed{(a + b + b) / (1 + 1 + 1)}{b})\}$ and
        $(d_2,R)\in\{(0,\intervalclosed{a}{(a + a + b) / (1 + 1 + 1)}),
        (1,\intervalclosed{(a + b + b) / (1 + 1 + 1)}{b})\}$
        by \cref{upair_intro_left,upair_intro_right,pair_eq_iff,subseteq}.
    Hence $H_1 = L$ and $H_2 = R$
        by \cref{space_filling_retained_third_child_choice_unique,subseteq,real_add_eq_subst}.
    $L$ is disjoint from $R$ by \cref{space_filling_retained_thirds_disjoint}.
    Suppose not.
    We have $y\in L$ and $y\in R$ by \cref{subseteq}.
    Contradiction by \cref{disjoint}.
\end{proof}

% from §44 helper lemma 101: every point of the unit interval has a dyadic binary code
\begin{lemma}\label{space_filling_point_has_dyadic_code}
    Let $I = \intervalclosed{0}{1}$.
    Let $D = \{0,1\}$.
    Let $E = \{(n,D) \mid n\in\naturalsPlus\}$.
    Assume $y\in I$.
    Then there exists $x\in\productFamily{E}$ such that
        $y$ is a dyadic code point of $x$ in $I$ and $E$.
\end{lemma}
\begin{proof}
    Let $K_0 = \{J\in\pow{I} \mid \text{there exist $a,b\in\reals$ such that
        $J = \intervalclosed{a}{b}$ and
        $a \leq b$ and
        $y\in J$}\}$.
    Let $R_0 = \{w\in K_0\times K_0 \mid \text{there exist $a,b$ such that
        $a\in\reals$ and
        $b\in\reals$ and
        $a \leq b$ and
        $\fst{w} = \intervalclosed{a}{b}$ and
        $\snd{w}\subseteq\fst{w}$ and
        $\snd{w}\in\{\intervalclosed{a}{(a + b) / (1 + 1)}, \intervalclosed{(a + b) / (1 + 1)}{b}\}$}\}$.
    Take $g_0$ such that
        $g_0\in\funs{K_0}{K_0}$ and
        for all $J\in K_0$ we have $J\mathrel{R_0}\apply{g_0}{J}$
        by \cref{space_filling_point_interval_child_relation_selector}.
    Take $u_0$ such that
        $u_0$ is a sequence in $K_0$ and
        $u_0(1) = I$ and
        for all $n\in\naturalsPlus$ we have $u_0(n + 1) = \apply{g_0}{u_0(n)}$
        by \cref{space_filling_point_interval_family_iteration}.
    Let $Q = \{w\in\naturalsPlus\times D \mid \text{there exist $a,b$ such that
        $a\in\reals$ and
        $b\in\reals$ and
        $a \leq b$ and
        $u_0(\fst{w}) = \intervalclosed{a}{b}$ and
        $(\snd{w},u_0(\fst{w} + 1))\in\{(0,\intervalclosed{a}{(a + b) / (1 + 1)}),
            (1,\intervalclosed{(a + b) / (1 + 1)}{b})\}$}\}$.
    $Q$ is a relation by \cref{relation,subseteq,times_elem_is_tuple}.
    For all $n\in\naturalsPlus$ we have $u_0(n)\in K_0$ and $u_0(n)\mathrel{R_0}u_0(n + 1)$
        by \cref{sequence_in,funs_type_apply,real_naturals_closed_succ,subseteq}.
    Thus for all $n\in\naturalsPlus$ there exists $d$ such that $n\mathrel{Q} d$
        by \cref{space_filling_point_interval_child_has_digit,times_tuple_intro,fst_eq,snd_eq,pair_eq_iff,relation}.
    Take $x$ such that
        $x$ is a function on $\naturalsPlus$ and
        for all $n\in\naturalsPlus$ we have $n\mathrel{Q}\apply{x}{n}$
        by \cref{choice_function}.
    We have $E$ is a function on $\naturalsPlus$
        by \cref{relation,rightunique,dom_intro,dom_iff,pair_eq_iff,subseteq,subseteq_antisymmetric}.
    $\dom{E} = \naturalsPlus$ by \cref{funs}.
    For all $n\in\naturalsPlus$ we have $E(n) = D$ by \cref{function_apply_intro}.
    Hence $x\in\funs{\naturalsPlus}{D}$ by \cref{choice_function,relation,subseteq,times_tuple_elim,funs_intro}.
    For all $n\in\dom{E}$ we have $x(n)\in E(n)$ by \cref{choice_function,function_apply_intro,subseteq,funs_type_apply}.
    For all $n\in\dom{x}$ we have $x(n)\in\unions{\ran{E}}$
        by \cref{funs_elim,function_apply_elim,ran_intro,unions_iff,function_apply_intro,subseteq}.
    We have $x\in\funs{\dom{E}}{\unions{\ran{E}}}$ by \cref{funs_intro}.
    Hence $x\in\productFamily{E}$ by \cref{indexed_product}.
    Let $K = \{J\in\pow{I} \mid \text{there exist $a,b\in\reals$ such that
        $J = \intervalclosed{a}{b}$ and
        $a \leq b$}\}$.
    Let $G = \naturalsPlus\times K$.
    Let $R = \{p\in G\times G \mid \text{there exist $n,J,H,a,b,d$ such that
        $n\in\naturalsPlus$ and
        $\fst{p} = (n,J)$ and
        $\snd{p} = (n + 1,H)$ and
        $d = x(n)$ and
        $a\in\reals$ and
        $b\in\reals$ and
        $a \leq b$ and
        $J = \intervalclosed{a}{b}$ and
        $H\subseteq J$ and
        $(d,H)\in\{(0,\intervalclosed{a}{(a + b) / (1 + 1)}),
            (1,\intervalclosed{(a + b) / (1 + 1)}{b})\}$}\}$.
    Take $g$ such that
        $g\in\funs{G}{G}$ and
        for all $z\in G$ we have $z\mathrel{R}\apply{g}{z}$
        by \cref{space_filling_dyadic_child_state_selector}.
    Take $u$ such that
        $u$ is a sequence in $G$ and
        $u(1) = (1,I)$ and
        for all $n\in\naturalsPlus$ we have $u(n + 1) = \apply{g}{u(n)}$
        by \cref{space_filling_dyadic_child_state_iteration}.
    We have for all $n\in\naturalsPlus$ we have $\fst{u(n)} = n$
        by \cref{space_filling_dyadic_child_state_orbit_index}.
    Let $A_0 = \{n\in\naturalsPlus \mid \snd{u(n)} = u_0(n)\}$.
    $A_0\subseteq\naturalsPlus$ by \cref{subseteq}.
    $1\in A_0$ by \cref{snd_eq,pair_eq_iff,real_one_in_naturals,subseteq}.
    Show that for all $n\in A_0$ we have $n + 1\in A_0$.
    \begin{subproof}
        Fix $n\in A_0$.
        We have $n\in\naturalsPlus$ and $\snd{u(n)} = u_0(n)$ by \cref{subseteq}.
        $(u(n),u(n + 1))\in R$
            by \cref{sequence_in,funs_type_apply,real_naturals_closed_succ,relation,subseteq}.
        Take $k,J,H,a,b,d$ such that
            $(u(n),u(n + 1))\in G\times G$ and
            $k\in\naturalsPlus$ and
            $\fst{(u(n),u(n + 1))} = (k,J)$ and
            $\snd{(u(n),u(n + 1))} = (k + 1,H)$ and
            $d = x(k)$ and
            $a\in\reals$ and
            $b\in\reals$ and
            $a \leq b$ and
            $J = \intervalclosed{a}{b}$ and
            $H\subseteq J$ and
            $(d,H)\in\{(0,\intervalclosed{a}{(a + b) / (1 + 1)}),
                (1,\intervalclosed{(a + b) / (1 + 1)}{b})\}$
            by \cref{subseteq}.
        $\fst{u(n)} = k$ and $\snd{u(n)} = J$ and $\snd{u(n + 1)} = H$
            by \cref{fst_eq,snd_eq,pair_eq_iff}.
        We have $(n,x(n))\in Q$ by \cref{choice_function,relation}.
        Take $c,e$ such that
            $c\in\reals$ and
            $e\in\reals$ and
            $c \leq e$ and
            $u_0(n) = \intervalclosed{c}{e}$ and
            $(x(n),u_0(n + 1))\in\{(0,\intervalclosed{c}{(c + e) / (1 + 1)}),
                (1,\intervalclosed{(c + e) / (1 + 1)}{e})\}$
            by \cref{subseteq,fst_eq,snd_eq,pair_eq_iff}.
        $a = c$ and $b = e$ by \cref{subseteq,space_filling_intervalclosed_endpoints_unique}.
        We have $(x(n),u_0(n + 1))\in\{(0,\intervalclosed{a}{(a + b) / (1 + 1)}),
            (1,\intervalclosed{(a + b) / (1 + 1)}{b})\}$ by \cref{subseteq,real_add_eq_subst}.
        Hence $H = u_0(n + 1)$ by \cref{space_filling_dyadic_child_choice_unique,subseteq,real_add_eq_subst}.
        Therefore $n + 1\in A_0$ by \cref{subseteq,real_naturals_closed_succ}.
    \end{subproof}
    We have for all $n\in\naturalsPlus$ we have $n\in A_0$ by \cref{real_naturals_induction}.
    Therefore for all $n\in\naturalsPlus$ we have $y\in\snd{u(n)}$
        by \cref{sequence_in,funs_type_apply,subseteq}.
    $y$ is a dyadic code point of $x$ in $I$ and $E$
        by \cref{space_filling_retained_third_state_selector_on,space_filling_retained_third_state_orbit_from,space_filling_dyadic_code_point}.
    Therefore there exists $z\in\productFamily{E}$ such that
        $y$ is a dyadic code point of $z$ in $I$ and $E$ by \cref{subseteq}.
\end{proof}

% from §44 helper lemma 102: opposite first retained-third digits cannot share a code point
\begin{lemma}\label{space_filling_retained_third_first_digit_separation}
    Let $I = \intervalclosed{0}{1}$.
    Let $D = \{0,1\}$.
    Let $E = \{(n,D) \mid n\in\naturalsPlus\}$.
    Assume $x_1\in\productFamily{E}$.
    Assume $x_2\in\productFamily{E}$.
    Suppose $y$ is a retained-third code point of $x_1$ in $I$ and $E$.
    Suppose $y$ is a retained-third code point of $x_2$ in $I$ and $E$.
    Assume $x_1(1) = 0$.
    Assume $x_2(1) = 1$.
    Then $y\notin I$.
\end{lemma}
\begin{proof}
    Take $K_1,G_1,R_1,g_1,u_1$ such that
        $K_1 = \{J\in\pow{I} \mid \text{there exist $a,b\in\reals$ such that
            $J = \intervalclosed{a}{b}$ and
            $a \leq b$}\}$ and
        $G_1 = \naturalsPlus\times K_1$ and
        $R_1 = \{p\in G_1\times G_1 \mid \text{there exist $n,J,H,a,b,d$ such that
            $n\in\naturalsPlus$ and
            $\fst{p} = (n,J)$ and
            $\snd{p} = (n + 1,H)$ and
            $d = x_1(n)$ and
            $a\in\reals$ and
            $b\in\reals$ and
            $a \leq b$ and
            $J = \intervalclosed{a}{b}$ and
            $H\subseteq J$ and
            $(d,H)\in\{(0,\intervalclosed{a}{(a + a + b) / (1 + 1 + 1)}),
                (1,\intervalclosed{(a + b + b) / (1 + 1 + 1)}{b})\}$}\}$ and
        $g_1$ is a retained-third state selector on $G_1$ and $R_1$ and
        $u_1$ is a retained-third state orbit of $g_1$ on $G_1$ from $I$ and
        for all $n\in\naturalsPlus$ we have $y\in\snd{u_1(n)}$
        by \cref{space_filling_retained_third_code_point}.
    We have $g_1\in\funs{G_1}{G_1}$ and
        for all $z\in G_1$ we have $z\mathrel{R_1}\apply{g_1}{z}$
        by \cref{space_filling_retained_third_state_selector_on}.
    $u_1$ is a sequence in $G_1$ and
        $u_1(1) = (1,I)$ and
        for all $n\in\naturalsPlus$ we have $u_1(n + 1) = \apply{g_1}{u_1(n)}$
        by \cref{space_filling_retained_third_state_orbit_from}.
    We have $(u_1(1),\apply{g_1}{u_1(1)})\in R_1$
        by \cref{sequence_in,funs_type_apply,real_one_in_naturals,relation,subseteq,times_elem_is_tuple}.
    Take $k_1,J_1,H_1,a_1,b_1,d_1$ such that
        $(u_1(1),\apply{g_1}{u_1(1)})\in G_1\times G_1$ and
        $k_1\in\naturalsPlus$ and
        $\fst{(u_1(1),\apply{g_1}{u_1(1)})} = (k_1,J_1)$ and
        $\snd{(u_1(1),\apply{g_1}{u_1(1)})} = (k_1 + 1,H_1)$ and
        $d_1 = x_1(k_1)$ and
        $a_1\in\reals$ and
        $b_1\in\reals$ and
        $a_1 \leq b_1$ and
        $J_1 = \intervalclosed{a_1}{b_1}$ and
        $H_1\subseteq J_1$ and
        $(d_1,H_1)\in\{(0,\intervalclosed{a_1}{(a_1 + a_1 + b_1) / (1 + 1 + 1)}),
            (1,\intervalclosed{(a_1 + b_1 + b_1) / (1 + 1 + 1)}{b_1})\}$
        by \cref{subseteq}.
    $u_1(1) = (k_1,J_1)$ and $\snd{\apply{g_1}{u_1(1)}} = H_1$
        by \cref{fst_eq,snd_eq,pair_eq_iff}.
    We have $k_1 = 1$ and $J_1 = I$ by \cref{pair_eq_iff,subseteq}.
    $d_1 = 0$ by \cref{subseteq}.
    We have $a_1 = 0$ and $b_1 = 1$
        by \cref{real_add_ident,real_one_in_naturals,real_naturals_subset_reals,subseteq,real_zero_lt_one,real_lt_imp_leq,space_filling_intervalclosed_endpoints_unique}.
    $(d_1,H_1)\in\{(0,\intervalclosed{0}{(0 + 0 + 1) / (1 + 1 + 1)}),
        (1,\intervalclosed{(0 + 1 + 1) / (1 + 1 + 1)}{1})\}$
        by \cref{subseteq,real_add_eq_subst}.
    Let $L = \intervalclosed{0}{(0 + 0 + 1) / (1 + 1 + 1)}$.
    Hence $H_1 = L$
        by \cref{space_filling_retained_third_child_choice_unique,upair_intro_left,pair_eq_iff,subseteq,real_add_eq_subst}.
    Take $K_2,G_2,R_2,g_2,u_2$ such that
        $K_2 = \{J\in\pow{I} \mid \text{there exist $a,b\in\reals$ such that
            $J = \intervalclosed{a}{b}$ and
            $a \leq b$}\}$ and
        $G_2 = \naturalsPlus\times K_2$ and
        $R_2 = \{p\in G_2\times G_2 \mid \text{there exist $n,J,H,a,b,d$ such that
            $n\in\naturalsPlus$ and
            $\fst{p} = (n,J)$ and
            $\snd{p} = (n + 1,H)$ and
            $d = x_2(n)$ and
            $a\in\reals$ and
            $b\in\reals$ and
            $a \leq b$ and
            $J = \intervalclosed{a}{b}$ and
            $H\subseteq J$ and
            $(d,H)\in\{(0,\intervalclosed{a}{(a + a + b) / (1 + 1 + 1)}),
                (1,\intervalclosed{(a + b + b) / (1 + 1 + 1)}{b})\}$}\}$ and
        $g_2$ is a retained-third state selector on $G_2$ and $R_2$ and
        $u_2$ is a retained-third state orbit of $g_2$ on $G_2$ from $I$ and
        for all $n\in\naturalsPlus$ we have $y\in\snd{u_2(n)}$
        by \cref{space_filling_retained_third_code_point}.
    We have $g_2\in\funs{G_2}{G_2}$ and
        for all $z\in G_2$ we have $z\mathrel{R_2}\apply{g_2}{z}$
        by \cref{space_filling_retained_third_state_selector_on}.
    $u_2$ is a sequence in $G_2$ and
        $u_2(1) = (1,I)$ and
        for all $n\in\naturalsPlus$ we have $u_2(n + 1) = \apply{g_2}{u_2(n)}$
        by \cref{space_filling_retained_third_state_orbit_from}.
    We have $(u_2(1),\apply{g_2}{u_2(1)})\in R_2$
        by \cref{sequence_in,funs_type_apply,real_one_in_naturals,relation,subseteq,times_elem_is_tuple}.
    Take $k_2,J_2,H_2,a_2,b_2,d_2$ such that
        $(u_2(1),\apply{g_2}{u_2(1)})\in G_2\times G_2$ and
        $k_2\in\naturalsPlus$ and
        $\fst{(u_2(1),\apply{g_2}{u_2(1)})} = (k_2,J_2)$ and
        $\snd{(u_2(1),\apply{g_2}{u_2(1)})} = (k_2 + 1,H_2)$ and
        $d_2 = x_2(k_2)$ and
        $a_2\in\reals$ and
        $b_2\in\reals$ and
        $a_2 \leq b_2$ and
        $J_2 = \intervalclosed{a_2}{b_2}$ and
        $H_2\subseteq J_2$ and
        $(d_2,H_2)\in\{(0,\intervalclosed{a_2}{(a_2 + a_2 + b_2) / (1 + 1 + 1)}),
            (1,\intervalclosed{(a_2 + b_2 + b_2) / (1 + 1 + 1)}{b_2})\}$
        by \cref{subseteq}.
    $u_2(1) = (k_2,J_2)$ and $\snd{\apply{g_2}{u_2(1)}} = H_2$
        by \cref{fst_eq,snd_eq,pair_eq_iff}.
    We have $k_2 = 1$ and $J_2 = I$ by \cref{pair_eq_iff,subseteq}.
    $d_2 = 1$ by \cref{subseteq}.
    We have $a_2 = 0$ and $b_2 = 1$ by \cref{subseteq,space_filling_intervalclosed_endpoints_unique}.
    $(d_2,H_2)\in\{(0,\intervalclosed{0}{(0 + 0 + 1) / (1 + 1 + 1)}),
        (1,\intervalclosed{(0 + 1 + 1) / (1 + 1 + 1)}{1})\}$
        by \cref{subseteq,real_add_eq_subst}.
    Let $R = \intervalclosed{(0 + 1 + 1) / (1 + 1 + 1)}{1}$.
    Hence $H_2 = R$
        by \cref{space_filling_retained_third_child_choice_unique,upair_intro_right,pair_eq_iff,subseteq,real_add_eq_subst}.
    $L$ is disjoint from $R$ by \cref{space_filling_retained_thirds_disjoint,real_zero_lt_one}.
    Suppose not.
    We have $y\in\snd{u_1(1 + 1)}$ and $y\in\snd{u_2(1 + 1)}$ by \cref{subseteq,real_one_in_naturals,real_naturals_closed_succ}.
    $u_1(1 + 1) = \apply{g_1}{u_1(1)}$ and $u_2(1 + 1) = \apply{g_2}{u_2(1)}$
        by \cref{subseteq,real_one_in_naturals}.
    Hence $y\in L$ and $y\in R$ by \cref{subseteq}.
    Contradiction by \cref{disjoint}.
\end{proof}

% from §44 helper lemma 103: retained-third state orbits are unique for a fixed binary sequence
\begin{lemma}\label{space_filling_retained_third_state_orbit_unique}
    Let $I = \intervalclosed{0}{1}$.
    Let $D = \{0,1\}$.
    Let $E = \{(n,D) \mid n\in\naturalsPlus\}$.
    Assume $x\in\productFamily{E}$.
    Let $K = \{J\in\pow{I} \mid \text{there exist $a,b\in\reals$ such that
        $J = \intervalclosed{a}{b}$ and
        $a \leq b$}\}$.
    Let $G = \naturalsPlus\times K$.
    Let $R = \{p\in G\times G \mid \text{there exist $n,J,H,a,b,d$ such that
        $n\in\naturalsPlus$ and
        $\fst{p} = (n,J)$ and
        $\snd{p} = (n + 1,H)$ and
        $d = x(n)$ and
        $a\in\reals$ and
        $b\in\reals$ and
        $a \leq b$ and
        $J = \intervalclosed{a}{b}$ and
        $H\subseteq J$ and
        $(d,H)\in\{(0,\intervalclosed{a}{(a + a + b) / (1 + 1 + 1)}),
            (1,\intervalclosed{(a + b + b) / (1 + 1 + 1)}{b})\}$}\}$.
    Suppose $g_1\in\funs{G}{G}$.
    Suppose for all $z\in G$ we have $z\mathrel{R}\apply{g_1}{z}$.
    Suppose $u_1$ is a sequence in $G$.
    Suppose $u_1(1) = (1,I)$.
    Suppose for all $n\in\naturalsPlus$ we have $u_1(n + 1) = \apply{g_1}{u_1(n)}$.
    Suppose $g_2\in\funs{G}{G}$.
    Suppose for all $z\in G$ we have $z\mathrel{R}\apply{g_2}{z}$.
    Suppose $u_2$ is a sequence in $G$.
    Suppose $u_2(1) = (1,I)$.
    Suppose for all $n\in\naturalsPlus$ we have $u_2(n + 1) = \apply{g_2}{u_2(n)}$.
    Then for all $n\in\naturalsPlus$ we have $u_1(n) = u_2(n)$.
\end{lemma}
\begin{proof}
    We have $R$ is right-unique by \cref{space_filling_retained_third_state_relation_rightunique}.
    For all $z\in G$ we have $\apply{g_1}{z} = \apply{g_2}{z}$.
    Let $A_0 = \{m\in\naturalsPlus \mid u_1(m) = u_2(m)\}$.
    $1\in A_0$ by \cref{real_one_in_naturals,pair_eq_iff,subseteq}.
    Show that for all $m\in A_0$ we have $m + 1\in A_0$.
    \begin{subproof}
        Follows by \cref{subseteq,sequence_in,funs_type_apply,real_naturals_closed_succ}.
    \end{subproof}
    Therefore for all $n\in\naturalsPlus$ we have $u_1(n) = u_2(n)$
        by \cref{real_naturals_induction,subseteq}.
\end{proof}

% from §44 helper lemma 104: every retained-third code point lies on every standard orbit for its sequence
\begin{lemma}\label{space_filling_retained_third_code_point_on_standard_orbit}
    Let $I = \intervalclosed{0}{1}$.
    Let $D = \{0,1\}$.
    Let $E = \{(n,D) \mid n\in\naturalsPlus\}$.
    Assume $x\in\productFamily{E}$.
    Let $K = \{J\in\pow{I} \mid \text{there exist $a,b\in\reals$ such that
        $J = \intervalclosed{a}{b}$ and
        $a \leq b$}\}$.
    Let $G = \naturalsPlus\times K$.
    Let $R = \{p\in G\times G \mid \text{there exist $n,J,H,a,b,d$ such that
        $n\in\naturalsPlus$ and
        $\fst{p} = (n,J)$ and
        $\snd{p} = (n + 1,H)$ and
        $d = x(n)$ and
        $a\in\reals$ and
        $b\in\reals$ and
        $a \leq b$ and
        $J = \intervalclosed{a}{b}$ and
        $H\subseteq J$ and
        $(d,H)\in\{(0,\intervalclosed{a}{(a + a + b) / (1 + 1 + 1)}),
            (1,\intervalclosed{(a + b + b) / (1 + 1 + 1)}{b})\}$}\}$.
    Suppose $g\in\funs{G}{G}$.
    Suppose for all $z\in G$ we have $z\mathrel{R}\apply{g}{z}$.
    Suppose $u$ is a sequence in $G$.
    Suppose $u(1) = (1,I)$.
    Suppose for all $n\in\naturalsPlus$ we have $u(n + 1) = \apply{g}{u(n)}$.
    Suppose $y$ is a retained-third code point of $x$ in $I$ and $E$.
    Then for all $n\in\naturalsPlus$ we have $y\in\snd{u(n)}$.
\end{lemma}
\begin{proof}
    Take $K_1,G_1,R_1,g_1,u_1$ such that
        $K_1 = \{J\in\pow{I} \mid \text{there exist $a,b\in\reals$ such that
            $J = \intervalclosed{a}{b}$ and
            $a \leq b$}\}$ and
        $G_1 = \naturalsPlus\times K_1$ and
        $R_1 = \{p\in G_1\times G_1 \mid \text{there exist $n,J,H,a,b,d$ such that
            $n\in\naturalsPlus$ and
            $\fst{p} = (n,J)$ and
            $\snd{p} = (n + 1,H)$ and
            $d = x(n)$ and
            $a\in\reals$ and
            $b\in\reals$ and
            $a \leq b$ and
            $J = \intervalclosed{a}{b}$ and
            $H\subseteq J$ and
            $(d,H)\in\{(0,\intervalclosed{a}{(a + a + b) / (1 + 1 + 1)}),
                (1,\intervalclosed{(a + b + b) / (1 + 1 + 1)}{b})\}$}\}$ and
        $g_1$ is a retained-third state selector on $G_1$ and $R_1$ and
        $u_1$ is a retained-third state orbit of $g_1$ on $G_1$ from $I$ and
        for all $n\in\naturalsPlus$ we have $y\in\snd{u_1(n)}$
        by \cref{space_filling_retained_third_code_point}.
    We have $g_1\in\funs{G_1}{G_1}$ and
        for all $z\in G_1$ we have $z\mathrel{R_1}\apply{g_1}{z}$
        by \cref{space_filling_retained_third_state_selector_on}.
    $u_1$ is a sequence in $G_1$ and
        $u_1(1) = (1,I)$ and
        for all $n\in\naturalsPlus$ we have $u_1(n + 1) = \apply{g_1}{u_1(n)}$
        by \cref{space_filling_retained_third_state_orbit_from}.
    $K_1 = K$ by \cref{subseteq,subseteq_antisymmetric}.
    $G_1 = G$ by \cref{subseteq}.
    $R_1\subseteq R$ and $R\subseteq R_1$ by \cref{subseteq}.
    $R_1 = R$ by \cref{subseteq_antisymmetric}.
    Therefore $g_1\in\funs{G}{G}$ and for all $z\in G$ we have $z\mathrel{R}\apply{g_1}{z}$ by \cref{subseteq}.
    $u_1$ is a sequence in $G$ and
        $u_1(1) = (1,I)$ and
        for all $n\in\naturalsPlus$ we have $u_1(n + 1) = \apply{g_1}{u_1(n)}$
        by \cref{subseteq}.
    Hence for all $n\in\naturalsPlus$ we have $u_1(n) = u(n)$
        by \cref{space_filling_retained_third_state_orbit_unique}.
    Thus for all $n\in\naturalsPlus$ we have $y\in\snd{u(n)}$ by \cref{subseteq}.
\end{proof}


% from §44 helper lemma 105: the dyadic state relation is right-unique
\begin{lemma}\label{space_filling_dyadic_state_relation_rightunique}
    Let $I = \intervalclosed{0}{1}$.
    Let $D = \{0,1\}$.
    Let $E = \{(n,D) \mid n\in\naturalsPlus\}$.
    Assume $x\in\productFamily{E}$.
    Let $K = \{J\in\pow{I} \mid \text{there exist $a,b\in\reals$ such that
        $J = \intervalclosed{a}{b}$ and
        $a \leq b$}\}$.
    Let $G = \naturalsPlus\times K$.
    Let $R = \{p\in G\times G \mid \text{there exist $n,J,H,a,b,d$ such that
        $n\in\naturalsPlus$ and
        $\fst{p} = (n,J)$ and
        $\snd{p} = (n + 1,H)$ and
        $d = x(n)$ and
        $a\in\reals$ and
        $b\in\reals$ and
        $a \leq b$ and
        $J = \intervalclosed{a}{b}$ and
        $H\subseteq J$ and
        $(d,H)\in\{(0,\intervalclosed{a}{(a + b) / (1 + 1)}),
            (1,\intervalclosed{(a + b) / (1 + 1)}{b})\}$}\}$.
    Then $R$ is right-unique.
\end{lemma}
\begin{proof}
    We have $E$ is a function on $\naturalsPlus$
        by \cref{relation,rightunique,dom_intro,dom_iff,pair_eq_iff,subseteq,subseteq_antisymmetric}.
    $\dom{E} = \naturalsPlus$ by \cref{funs}.
    For all $n\in\naturalsPlus$ we have $E(n) = D$ by \cref{function_apply_intro}.
    Show that for all $z,w_1,w_2$ such that $(z,w_1)\in R$ and $(z,w_2)\in R$ we have $w_1 = w_2$.
    \begin{subproof}
        Follows by \cref{subseteq,fst_eq,snd_eq,pair_eq_iff,indexed_product,space_filling_intervalclosed_endpoints_unique,space_filling_dyadic_child_choice_unique,real_add_eq_subst}.
    \end{subproof}
    Therefore $R$ is right-unique by \cref{rightunique}.
\end{proof}

% from §44 helper lemma 106: dyadic state orbits are unique for a fixed binary sequence
\begin{lemma}\label{space_filling_dyadic_state_orbit_unique}
    Let $I = \intervalclosed{0}{1}$.
    Let $D = \{0,1\}$.
    Let $E = \{(n,D) \mid n\in\naturalsPlus\}$.
    Assume $x\in\productFamily{E}$.
    Let $K = \{J\in\pow{I} \mid \text{there exist $a,b\in\reals$ such that
        $J = \intervalclosed{a}{b}$ and
        $a \leq b$}\}$.
    Let $G = \naturalsPlus\times K$.
    Let $R = \{p\in G\times G \mid \text{there exist $n,J,H,a,b,d$ such that
        $n\in\naturalsPlus$ and
        $\fst{p} = (n,J)$ and
        $\snd{p} = (n + 1,H)$ and
        $d = x(n)$ and
        $a\in\reals$ and
        $b\in\reals$ and
        $a \leq b$ and
        $J = \intervalclosed{a}{b}$ and
        $H\subseteq J$ and
        $(d,H)\in\{(0,\intervalclosed{a}{(a + b) / (1 + 1)}),
            (1,\intervalclosed{(a + b) / (1 + 1)}{b})\}$}\}$.
    Suppose $g_1\in\funs{G}{G}$.
    Suppose for all $z\in G$ we have $z\mathrel{R}\apply{g_1}{z}$.
    Suppose $u_1$ is a sequence in $G$.
    Suppose $u_1(1) = (1,I)$.
    Suppose for all $n\in\naturalsPlus$ we have $u_1(n + 1) = \apply{g_1}{u_1(n)}$.
    Suppose $g_2\in\funs{G}{G}$.
    Suppose for all $z\in G$ we have $z\mathrel{R}\apply{g_2}{z}$.
    Suppose $u_2$ is a sequence in $G$.
    Suppose $u_2(1) = (1,I)$.
    Suppose for all $n\in\naturalsPlus$ we have $u_2(n + 1) = \apply{g_2}{u_2(n)}$.
    Then for all $n\in\naturalsPlus$ we have $u_1(n) = u_2(n)$.
\end{lemma}
\begin{proof}
    We have $R$ is right-unique by \cref{space_filling_dyadic_state_relation_rightunique}.
    For all $z\in G$ we have $\apply{g_1}{z} = \apply{g_2}{z}$.
    Let $A_0 = \{m\in\naturalsPlus \mid u_1(m) = u_2(m)\}$.
    $1\in A_0$ by \cref{real_one_in_naturals,pair_eq_iff,subseteq}.
    Show that for all $m\in A_0$ we have $m + 1\in A_0$.
    \begin{subproof}
        Follows by \cref{subseteq,sequence_in,funs_type_apply,real_naturals_closed_succ}.
    \end{subproof}
    Therefore for all $n\in\naturalsPlus$ we have $u_1(n) = u_2(n)$
        by \cref{real_naturals_induction,subseteq}.
\end{proof}

% from §44 helper lemma 107: every dyadic code point lies on every standard orbit for its sequence
\begin{lemma}\label{space_filling_dyadic_code_point_on_standard_orbit}
    Let $I = \intervalclosed{0}{1}$.
    Let $D = \{0,1\}$.
    Let $E = \{(n,D) \mid n\in\naturalsPlus\}$.
    Assume $x\in\productFamily{E}$.
    Let $K = \{J\in\pow{I} \mid \text{there exist $a,b\in\reals$ such that
        $J = \intervalclosed{a}{b}$ and
        $a \leq b$}\}$.
    Let $G = \naturalsPlus\times K$.
    Let $R = \{p\in G\times G \mid \text{there exist $n,J,H,a,b,d$ such that
        $n\in\naturalsPlus$ and
        $\fst{p} = (n,J)$ and
        $\snd{p} = (n + 1,H)$ and
        $d = x(n)$ and
        $a\in\reals$ and
        $b\in\reals$ and
        $a \leq b$ and
        $J = \intervalclosed{a}{b}$ and
        $H\subseteq J$ and
        $(d,H)\in\{(0,\intervalclosed{a}{(a + b) / (1 + 1)}),
            (1,\intervalclosed{(a + b) / (1 + 1)}{b})\}$}\}$.
    Suppose $g\in\funs{G}{G}$.
    Suppose for all $z\in G$ we have $z\mathrel{R}\apply{g}{z}$.
    Suppose $u$ is a sequence in $G$.
    Suppose $u(1) = (1,I)$.
    Suppose for all $n\in\naturalsPlus$ we have $u(n + 1) = \apply{g}{u(n)}$.
    Suppose $y$ is a dyadic code point of $x$ in $I$ and $E$.
    Then for all $n\in\naturalsPlus$ we have $y\in\snd{u(n)}$.
\end{lemma}
\begin{proof}
    Take $K_1,G_1,R_1,g_1,u_1$ such that
        $K_1 = \{J\in\pow{I} \mid \text{there exist $a,b\in\reals$ such that
            $J = \intervalclosed{a}{b}$ and
            $a \leq b$}\}$ and
        $G_1 = \naturalsPlus\times K_1$ and
        $R_1 = \{p\in G_1\times G_1 \mid \text{there exist $n,J,H,a,b,d$ such that
            $n\in\naturalsPlus$ and
            $\fst{p} = (n,J)$ and
            $\snd{p} = (n + 1,H)$ and
            $d = x(n)$ and
            $a\in\reals$ and
            $b\in\reals$ and
            $a \leq b$ and
            $J = \intervalclosed{a}{b}$ and
            $H\subseteq J$ and
            $(d,H)\in\{(0,\intervalclosed{a}{(a + b) / (1 + 1)}),
                (1,\intervalclosed{(a + b) / (1 + 1)}{b})\}$}\}$ and
        $g_1$ is a retained-third state selector on $G_1$ and $R_1$ and
        $u_1$ is a retained-third state orbit of $g_1$ on $G_1$ from $I$ and
        for all $n\in\naturalsPlus$ we have $y\in\snd{u_1(n)}$
        by \cref{space_filling_dyadic_code_point}.
    We have $g_1\in\funs{G_1}{G_1}$ and
        for all $z\in G_1$ we have $z\mathrel{R_1}\apply{g_1}{z}$
        by \cref{space_filling_retained_third_state_selector_on}.
    $u_1$ is a sequence in $G_1$ and
        $u_1(1) = (1,I)$ and
        for all $n\in\naturalsPlus$ we have $u_1(n + 1) = \apply{g_1}{u_1(n)}$
        by \cref{space_filling_retained_third_state_orbit_from}.
    $K_1 = K$ by \cref{subseteq,subseteq_antisymmetric}.
    $G_1 = G$ by \cref{subseteq}.
    $R_1\subseteq R$ and $R\subseteq R_1$ by \cref{subseteq}.
    $R_1 = R$ by \cref{subseteq_antisymmetric}.
    Therefore $g_1\in\funs{G}{G}$ and for all $z\in G$ we have $z\mathrel{R}\apply{g_1}{z}$ by \cref{subseteq}.
    $u_1$ is a sequence in $G$ and
        $u_1(1) = (1,I)$ and
        for all $n\in\naturalsPlus$ we have $u_1(n + 1) = \apply{g_1}{u_1(n)}$
        by \cref{subseteq}.
    Hence for all $n\in\naturalsPlus$ we have $u_1(n) = u(n)$
        by \cref{space_filling_dyadic_state_orbit_unique}.
    Thus for all $n\in\naturalsPlus$ we have $y\in\snd{u(n)}$ by \cref{subseteq}.
\end{proof}

% from §44 helper lemma 115: equal dyadic prefixes yield the same stage interval
\begin{lemma}\label{space_filling_dyadic_state_orbit_prefix_stage_local}
    Let $I = \intervalclosed{0}{1}$.
    Let $D = \{0,1\}$.
    Let $E = \{(n,D) \mid n\in\naturalsPlus\}$.
    Assume $x_1\in\productFamily{E}$.
    Assume $x_2\in\productFamily{E}$.
    Let $K = \{J\in\pow{I} \mid \text{there exist $a,b\in\reals$ such that
        $J = \intervalclosed{a}{b}$ and
        $a \leq b$}\}$.
    Let $G = \naturalsPlus\times K$.
    Let $R_1 = \{p\in G\times G \mid \text{there exist $n,J,H,a,b,d$ such that
        $n\in\naturalsPlus$ and
        $\fst{p} = (n,J)$ and
        $\snd{p} = (n + 1,H)$ and
        $d = x_1(n)$ and
        $a\in\reals$ and
        $b\in\reals$ and
        $a \leq b$ and
        $J = \intervalclosed{a}{b}$ and
        $H\subseteq J$ and
        $(d,H)\in\{(0,\intervalclosed{a}{(a + b) / (1 + 1)}),
            (1,\intervalclosed{(a + b) / (1 + 1)}{b})\}$}\}$.
    Let $R_2 = \{p\in G\times G \mid \text{there exist $n,J,H,a,b,d$ such that
        $n\in\naturalsPlus$ and
        $\fst{p} = (n,J)$ and
        $\snd{p} = (n + 1,H)$ and
        $d = x_2(n)$ and
        $a\in\reals$ and
        $b\in\reals$ and
        $a \leq b$ and
        $J = \intervalclosed{a}{b}$ and
        $H\subseteq J$ and
        $(d,H)\in\{(0,\intervalclosed{a}{(a + b) / (1 + 1)}),
            (1,\intervalclosed{(a + b) / (1 + 1)}{b})\}$}\}$.
    Suppose $g_1\in\funs{G}{G}$.
    Suppose for all $z\in G$ we have $z\mathrel{R_1}\apply{g_1}{z}$.
    Suppose $u_1$ is a sequence in $G$.
    Suppose $u_1(1) = (1,I)$.
    Suppose for all $n\in\naturalsPlus$ we have $u_1(n + 1) = \apply{g_1}{u_1(n)}$.
    Suppose $g_2\in\funs{G}{G}$.
    Suppose for all $z\in G$ we have $z\mathrel{R_2}\apply{g_2}{z}$.
    Suppose $u_2$ is a sequence in $G$.
    Suppose $u_2(1) = (1,I)$.
    Suppose for all $n\in\naturalsPlus$ we have $u_2(n + 1) = \apply{g_2}{u_2(n)}$.
    Assume $p\in\naturalsPlus$.
    Assume for all $i\in\initialSegment{p}$ we have $x_1(i) = x_2(i)$.
    Then $\snd{u_1(p)} = \snd{u_2(p)}$.
\end{lemma}
\begin{proof}
    We have for all $n\in\initialSegment{p}$ we have $\snd{u_1(n)} = \snd{u_2(n)}$
        by \cref{space_filling_dyadic_state_orbit_prefix_agreement}.
    Therefore $\snd{u_1(p)} = \snd{u_2(p)}$
        by \cref{initial_segment_naturalsplus,subseteq}.
\end{proof}

% from §44 helper lemma 116: a standard dyadic orbit has the expected stage index at any stage
\begin{lemma}\label{space_filling_dyadic_state_orbit_stage_index_local}
    Let $I = \intervalclosed{0}{1}$.
    Let $K = \{J\in\pow{I} \mid \text{there exist $a,b\in\reals$ such that
        $J = \intervalclosed{a}{b}$ and
        $a \leq b$}\}$.
    Let $G = \naturalsPlus\times K$.
    Let $R = \{p\in G\times G \mid \text{there exist $n,J,H,a,b,d$ such that
        $n\in\naturalsPlus$ and
        $\fst{p} = (n,J)$ and
        $\snd{p} = (n + 1,H)$ and
        $a\in\reals$ and
        $b\in\reals$ and
        $a \leq b$ and
        $J = \intervalclosed{a}{b}$ and
        $H\subseteq J$ and
        $(d,H)\in\{(0,\intervalclosed{a}{(a + b) / (1 + 1)}),
            (1,\intervalclosed{(a + b) / (1 + 1)}{b})\}$}\}$.
    Suppose $g\in\funs{G}{G}$.
    Suppose for all $z\in G$ we have $z\mathrel{R}\apply{g}{z}$.
    Suppose $u$ is a sequence in $G$.
    Suppose $u(1) = (1,I)$.
    Suppose for all $n\in\naturalsPlus$ we have $u(n + 1) = \apply{g}{u(n)}$.
    Assume $p\in\naturalsPlus$.
    Then $\fst{u(p)} = p$.
\end{lemma}
\begin{proof}
    Let $A_0 = \{n\in\naturalsPlus \mid \fst{u(n)} = n\}$.
    $1\in A_0$ by \cref{real_one_in_naturals,pair_eq_iff,fst_eq,subseteq}.
    Show that for all $n\in A_0$ we have $n + 1\in A_0$.
    \begin{subproof}
        Follows by \cref{subseteq,sequence_in,funs_type_apply,relation,fst_eq,snd_eq,pair_eq_iff,real_add_eq_subst,real_naturals_closed_succ}.
    \end{subproof}
    Therefore $\fst{u(p)} = p$ by \cref{real_naturals_induction,subseteq}.
\end{proof}

% from §44 helper lemma 108: initial segments are transitive
\begin{lemma}\label{space_filling_initial_segment_transitive}
    Suppose $k\in\naturalsPlus$.
    Suppose $m\in\initialSegment{k}$.
    Suppose $i\in\initialSegment{m}$.
    Then $i\in\initialSegment{k}$.
\end{lemma}
\begin{proof}
    We have $i\in\initialSegment{k}$
        by \cref{initial_segment_naturalsplus,real_naturals_subset_reals,real_leq_trans,subseteq}.
\end{proof}

% from §44 helper lemma 117: a standard dyadic stage interval has width at most 1/n
\begin{lemma}\label{space_filling_dyadic_stage_interval_width_leq_local}
    Let $I = \intervalclosed{0}{1}$.
    Let $D = \{0,1\}$.
    Let $E = \{(n,D) \mid n\in\naturalsPlus\}$.
    Assume $x\in\productFamily{E}$.
    Let $K = \{J\in\pow{I} \mid \text{there exist $a,b\in\reals$ such that
        $J = \intervalclosed{a}{b}$ and
        $a \leq b$}\}$.
    Let $G = \naturalsPlus\times K$.
    Let $R = \{p\in G\times G \mid \text{there exist $n,J,H,a,b,d$ such that
        $n\in\naturalsPlus$ and
        $\fst{p} = (n,J)$ and
        $\snd{p} = (n + 1,H)$ and
        $d = x(n)$ and
        $a\in\reals$ and
        $b\in\reals$ and
        $a \leq b$ and
        $J = \intervalclosed{a}{b}$ and
        $H\subseteq J$ and
        $(d,H)\in\{(0,\intervalclosed{a}{(a + b) / (1 + 1)}),
            (1,\intervalclosed{(a + b) / (1 + 1)}{b})\}$}\}$.
    Suppose $g\in\funs{G}{G}$.
    Suppose for all $z\in G$ we have $z\mathrel{R}\apply{g}{z}$.
    Suppose $u$ is a sequence in $G$.
    Suppose $u(1) = (1,I)$.
    Suppose for all $n\in\naturalsPlus$ we have $u(n + 1) = \apply{g}{u(n)}$.
    Assume $N\in\naturalsPlus$.
    Assume $p\in\reals$.
    Assume $q\in\reals$.
    Assume $\snd{u(N)} = \intervalclosed{p}{q}$.
    Assume $p \leq q$.
    Then $q - p \leq 1 / N$.
\end{lemma}
\begin{proof}
    We have $q - p \leq 1 / N$ by \cref{space_filling_dyadic_child_state_orbit_width_bound}.
\end{proof}

% from §44 helper lemma 109: a standard retained-third orbit has the expected stage index at any stage
\begin{lemma}\label{space_filling_retained_third_state_orbit_stage_index_local}
    Let $I = \intervalclosed{0}{1}$.
    Let $K = \{J\in\pow{I} \mid \text{there exist $a,b\in\reals$ such that
        $J = \intervalclosed{a}{b}$ and
        $a \leq b$}\}$.
    Let $G = \naturalsPlus\times K$.
    Let $R = \{p\in G\times G \mid \text{there exist $n,J,H,a,b,d$ such that
        $n\in\naturalsPlus$ and
        $\fst{p} = (n,J)$ and
        $\snd{p} = (n + 1,H)$ and
        $a\in\reals$ and
        $b\in\reals$ and
        $a \leq b$ and
        $J = \intervalclosed{a}{b}$ and
        $H\subseteq J$ and
        $(d,H)\in\{(0,\intervalclosed{a}{(a + a + b) / (1 + 1 + 1)}),
            (1,\intervalclosed{(a + b + b) / (1 + 1 + 1)}{b})\}$}\}$.
    Suppose $g\in\funs{G}{G}$.
    Suppose for all $z\in G$ we have $z\mathrel{R}\apply{g}{z}$.
    Suppose $u$ is a sequence in $G$.
    Suppose $u(1) = (1,I)$.
    Suppose for all $n\in\naturalsPlus$ we have $u(n + 1) = \apply{g}{u(n)}$.
    Assume $p\in\naturalsPlus$.
    Then $\fst{u(p)} = p$.
\end{lemma}
\begin{proof}
    Let $A_0 = \{n\in\naturalsPlus \mid \fst{u(n)} = n\}$.
    $1\in A_0$ by \cref{real_one_in_naturals,pair_eq_iff,fst_eq,subseteq}.
    Show that for all $n\in A_0$ we have $n + 1\in A_0$.
    \begin{subproof}
        Follows by \cref{subseteq,sequence_in,funs_type_apply,relation,fst_eq,snd_eq,pair_eq_iff,real_add_eq_subst,real_naturals_closed_succ}.
    \end{subproof}
    Therefore $\fst{u(p)} = p$ by \cref{real_naturals_induction,subseteq}.
\end{proof}

% from §44 helper lemma 110: equal retained-third prefixes yield the same stage interval
\begin{lemma}\label{space_filling_retained_third_state_orbit_prefix_stage_local}
    Let $I = \intervalclosed{0}{1}$.
    Let $D = \{0,1\}$.
    Let $E = \{(n,D) \mid n\in\naturalsPlus\}$.
    Assume $x_1\in\productFamily{E}$.
    Assume $x_2\in\productFamily{E}$.
    Let $K = \{J\in\pow{I} \mid \text{there exist $a,b\in\reals$ such that
        $J = \intervalclosed{a}{b}$ and
        $a \leq b$}\}$.
    Let $G = \naturalsPlus\times K$.
    Let $R_1 = \{p\in G\times G \mid \text{there exist $n,J,H,a,b,d$ such that
        $n\in\naturalsPlus$ and
        $\fst{p} = (n,J)$ and
        $\snd{p} = (n + 1,H)$ and
        $d = x_1(n)$ and
        $a\in\reals$ and
        $b\in\reals$ and
        $a \leq b$ and
        $J = \intervalclosed{a}{b}$ and
        $H\subseteq J$ and
        $(d,H)\in\{(0,\intervalclosed{a}{(a + a + b) / (1 + 1 + 1)}),
            (1,\intervalclosed{(a + b + b) / (1 + 1 + 1)}{b})\}$}\}$.
    Let $R_2 = \{p\in G\times G \mid \text{there exist $n,J,H,a,b,d$ such that
        $n\in\naturalsPlus$ and
        $\fst{p} = (n,J)$ and
        $\snd{p} = (n + 1,H)$ and
        $d = x_2(n)$ and
        $a\in\reals$ and
        $b\in\reals$ and
        $a \leq b$ and
        $J = \intervalclosed{a}{b}$ and
        $H\subseteq J$ and
        $(d,H)\in\{(0,\intervalclosed{a}{(a + a + b) / (1 + 1 + 1)}),
            (1,\intervalclosed{(a + b + b) / (1 + 1 + 1)}{b})\}$}\}$.
    Suppose $g_1\in\funs{G}{G}$.
    Suppose for all $z\in G$ we have $z\mathrel{R_1}\apply{g_1}{z}$.
    Suppose $u_1$ is a sequence in $G$.
    Suppose $u_1(1) = (1,I)$.
    Suppose for all $n\in\naturalsPlus$ we have $u_1(n + 1) = \apply{g_1}{u_1(n)}$.
    Suppose $g_2\in\funs{G}{G}$.
    Suppose for all $z\in G$ we have $z\mathrel{R_2}\apply{g_2}{z}$.
    Suppose $u_2$ is a sequence in $G$.
    Suppose $u_2(1) = (1,I)$.
    Suppose for all $n\in\naturalsPlus$ we have $u_2(n + 1) = \apply{g_2}{u_2(n)}$.
    Assume $p\in\naturalsPlus$.
    Assume for all $i\in\initialSegment{p}$ we have $x_1(i) = x_2(i)$.
    Then $\snd{u_1(p)} = \snd{u_2(p)}$.
\end{lemma}
\begin{proof}
    We have for all $n\in\initialSegment{p}$ we have $\snd{u_1(n)} = \snd{u_2(n)}$
        by \cref{space_filling_retained_third_state_orbit_prefix_agreement}.
    Therefore $\snd{u_1(p)} = \snd{u_2(p)}$
        by \cref{initial_segment_naturalsplus,subseteq}.
\end{proof}

% from §44 helper lemma 111: a retained-third successor interval is nondegenerate
\begin{lemma}\label{space_filling_retained_third_state_orbit_succ_interval_strict_local}
    Let $I = \intervalclosed{0}{1}$.
    Let $K = \{J\in\pow{I} \mid \text{there exist $a,b\in\reals$ such that
        $J = \intervalclosed{a}{b}$ and
        $a \leq b$}\}$.
    Let $G = \naturalsPlus\times K$.
    Let $R = \{p\in G\times G \mid \text{there exist $n,J,H,a,b,d$ such that
        $n\in\naturalsPlus$ and
        $\fst{p} = (n,J)$ and
        $\snd{p} = (n + 1,H)$ and
        $a\in\reals$ and
        $b\in\reals$ and
        $a \leq b$ and
        $J = \intervalclosed{a}{b}$ and
        $H\subseteq J$ and
        $(d,H)\in\{(0,\intervalclosed{a}{(a + a + b) / (1 + 1 + 1)}),
            (1,\intervalclosed{(a + b + b) / (1 + 1 + 1)}{b})\}$}\}$.
    Suppose $g$ is a retained-third state selector on $G$ and $R$.
    Suppose $u$ is a retained-third state orbit of $g$ on $G$ from $I$.
    Then for all $n,a,b$ if
        $n\in\naturalsPlus$ and
        $a\in\reals$ and
        $b\in\reals$ and
        $\snd{u(n)} = \intervalclosed{a}{b}$ and
        $a \leq b$
    then $a < b$.
\end{lemma}
\begin{proof}
    We have $g\in\funs{G}{G}$ and
        for all $z\in G$ we have $z\mathrel{R}\apply{g}{z}$
        by \cref{space_filling_retained_third_state_selector_on}.
    $u$ is a sequence in $G$ and
        $u(1) = (1,I)$ and
        for all $n\in\naturalsPlus$ we have $u(n + 1) = \apply{g}{u(n)}$
        by \cref{space_filling_retained_third_state_orbit_from}.
    Let $A_0 = \{n\in\naturalsPlus \mid \text{for all $a,b$ if
        $a\in\reals$ and
        $b\in\reals$ and
        $\snd{u(n)} = \intervalclosed{a}{b}$ and
        $a \leq b$
    then $a < b$}\}$.
    $A_0\subseteq\naturalsPlus$ by \cref{subseteq}.
    Show that for all $a,b$ if
        $a\in\reals$ and
        $b\in\reals$ and
        $\snd{u(1)} = \intervalclosed{a}{b}$ and
        $a \leq b$
    then $a < b$.
    \begin{subproof}
        Follows by \cref{snd_eq,pair_eq_iff,subseteq,space_filling_intervalclosed_endpoints_unique,real_add_ident,real_one_in_naturals,real_naturals_subset_reals,real_zero_lt_one,real_lt_imp_leq}.
    \end{subproof}
    Hence $1\in A_0$ by \cref{real_one_in_naturals,subseteq}.
    Show that for all $n\in A_0$ we have $n + 1\in A_0$.
    \begin{subproof}
        Fix $n\in A_0$.
        Hence $n\in\naturalsPlus$ by \cref{subseteq}.
        Show that for all $a,b$ if
            $a\in\reals$ and
            $b\in\reals$ and
            $\snd{u(n + 1)} = \intervalclosed{a}{b}$ and
            $a \leq b$
        then $a < b$.
        \begin{subproof}
            Fix $a$.
            Fix $b$.
            Assume $a\in\reals$ and $b\in\reals$ and $\snd{u(n + 1)} = \intervalclosed{a}{b}$ and $a \leq b$.
            Hence $(u(n),u(n + 1))\in R$ by \cref{sequence_in,funs_type_apply,subseteq,relation}.
            Take $k,J,H,c,d,e$ such that
                $(u(n),u(n + 1))\in G\times G$ and
                $k\in\naturalsPlus$ and
                $\fst{(u(n),u(n + 1))} = (k,J)$ and
                $\snd{(u(n),u(n + 1))} = (k + 1,H)$ and
                $c\in\reals$ and
                $d\in\reals$ and
                $c \leq d$ and
                $J = \intervalclosed{c}{d}$ and
                $H\subseteq J$ and
                $(e,H)\in\{(0,\intervalclosed{c}{(c + c + d) / (1 + 1 + 1)}),
                    (1,\intervalclosed{(c + d + d) / (1 + 1 + 1)}{d})\}$
                by \cref{subseteq}.
            Therefore $\fst{u(n)} = k$ and $\snd{u(n)} = J$ and $\snd{u(n + 1)} = H$
                by \cref{fst_eq,snd_eq,pair_eq_iff}.
            Hence $k = n$ by \cref{space_filling_retained_third_state_orbit_stage_index_local,subseteq}.
            Thus $c < d$ by \cref{subseteq}.
            $(e,H) = (0,\intervalclosed{c}{(c + c + d) / (1 + 1 + 1)})$ or
                $(e,H) = (1,\intervalclosed{(c + d + d) / (1 + 1 + 1)}{d})$
                by \cref{upair_iff}.
            \begin{byCase}
            \caseOf{$(e,H) = (0,\intervalclosed{c}{(c + c + d) / (1 + 1 + 1)})$.}
                We have $H = \intervalclosed{c}{(c + c + d) / (1 + 1 + 1)}$ by \cref{pair_eq_iff}.
                Therefore $\intervalclosed{a}{b} = \intervalclosed{c}{(c + c + d) / (1 + 1 + 1)}$ by \cref{subseteq}.
                Hence $a = c$ and $b = (c + c + d) / (1 + 1 + 1)$
                    by \cref{space_filling_left_endpoint_lt_left_third,real_add_closed,real_div_naturalsplus_real,real_one_in_naturals,real_naturals_closed_succ,real_lt_imp_leq,space_filling_intervalclosed_endpoints_unique}.
                Therefore $a < b$ by \cref{space_filling_left_endpoint_lt_left_third,real_lt_subst_left,real_lt_subst_right}.
            \caseOf{$(e,H) = (1,\intervalclosed{(c + d + d) / (1 + 1 + 1)}{d})$.}
                We have $H = \intervalclosed{(c + d + d) / (1 + 1 + 1)}{d}$ by \cref{pair_eq_iff}.
                Therefore $\intervalclosed{a}{b} = \intervalclosed{(c + d + d) / (1 + 1 + 1)}{d}$ by \cref{subseteq}.
                Hence $a = (c + d + d) / (1 + 1 + 1)$ and $b = d$
                    by \cref{space_filling_right_third_lt_right_endpoint,real_add_closed,real_div_naturalsplus_real,real_one_in_naturals,real_naturals_closed_succ,real_lt_imp_leq,space_filling_intervalclosed_endpoints_unique}.
                Therefore $a < b$ by \cref{space_filling_right_third_lt_right_endpoint,real_lt_subst_left,real_lt_subst_right}.
            \end{byCase}
        \end{subproof}
        Therefore $n + 1\in A_0$ by \cref{real_naturals_closed_succ,subseteq}.
    \end{subproof}
    Hence for all $n\in\naturalsPlus$ we have $n\in A_0$ by \cref{real_naturals_induction}.
    Therefore for all $n,a,b$ if
        $n\in\naturalsPlus$ and
        $a\in\reals$ and
        $b\in\reals$ and
        $\snd{u(n)} = \intervalclosed{a}{b}$ and
        $a \leq b$
    then $a < b$ by \cref{subseteq}.
\end{proof}

% from §44 helper lemma 112: local retained-third digit separation from a shared strict interval
\begin{lemma}\label{space_filling_retained_third_shared_interval_digit_separation_local}
    Let $I = \intervalclosed{0}{1}$.
    Let $K = \{J\in\pow{I} \mid \text{there exist $c,d\in\reals$ such that
        $J = \intervalclosed{c}{d}$ and
        $c \leq d$}\}$.
    Let $G = \naturalsPlus\times K$.
    Let $R_1 = \{p\in G\times G \mid \text{there exist $n,J,H,c,d,e$ such that
        $n\in\naturalsPlus$ and
        $\fst{p} = (n,J)$ and
        $\snd{p} = (n + 1,H)$ and
        $e = x_1(n)$ and
        $c\in\reals$ and
        $d\in\reals$ and
        $c \leq d$ and
        $J = \intervalclosed{c}{d}$ and
        $H\subseteq J$ and
        $(e,H)\in\{(0,\intervalclosed{c}{(c + c + d) / (1 + 1 + 1)}),
            (1,\intervalclosed{(c + d + d) / (1 + 1 + 1)}{d})\}$}\}$.
    Let $R_2 = \{p\in G\times G \mid \text{there exist $n,J,H,c,d,e$ such that
        $n\in\naturalsPlus$ and
        $\fst{p} = (n,J)$ and
        $\snd{p} = (n + 1,H)$ and
        $e = x_2(n)$ and
        $c\in\reals$ and
        $d\in\reals$ and
        $c \leq d$ and
        $J = \intervalclosed{c}{d}$ and
        $H\subseteq J$ and
        $(e,H)\in\{(0,\intervalclosed{c}{(c + c + d) / (1 + 1 + 1)}),
            (1,\intervalclosed{(c + d + d) / (1 + 1 + 1)}{d})\}$}\}$.
    Suppose $g_1\in\funs{G}{G}$.
    Suppose for all $z\in G$ we have $z\mathrel{R_1}\apply{g_1}{z}$.
    Suppose $u_1$ is a sequence in $G$.
    Suppose $u_1(1) = (1,I)$.
    Suppose for all $n\in\naturalsPlus$ we have $u_1(n + 1) = \apply{g_1}{u_1(n)}$.
    Suppose $\fst{u_1(N + 1)} = N + 1$.
    Suppose $g_2\in\funs{G}{G}$.
    Suppose for all $z\in G$ we have $z\mathrel{R_2}\apply{g_2}{z}$.
    Suppose $u_2$ is a sequence in $G$.
    Suppose $u_2(1) = (1,I)$.
    Suppose for all $n\in\naturalsPlus$ we have $u_2(n + 1) = \apply{g_2}{u_2(n)}$.
    Suppose $\fst{u_2(N + 1)} = N + 1$.
    Suppose $y\in I$ and for all $n\in\naturalsPlus$ we have $y\in\snd{u_1(n)}$.
    Suppose for all $n\in\naturalsPlus$ we have $y\in\snd{u_2(n)}$.
    Assume $N\in\naturalsPlus$.
    Assume $a\in\reals$ and $b\in\reals$ and $a < b$.
    Assume $\snd{u_1(N + 1)} = \intervalclosed{a}{b}$.
    Assume $\snd{u_2(N + 1)} = \intervalclosed{a}{b}$.
    Assume $x_1(N + 1) = 0$.
    Assume $x_2(N + 1) = 1$.
    Then $y\notin I$.
\end{lemma}
\begin{proof}
    We have $u_1\in\funs{\naturalsPlus}{G}$ and $u_2\in\funs{\naturalsPlus}{G}$ and
        $N + 1\in\naturalsPlus$ and $(N + 1) + 1\in\naturalsPlus$
        by \cref{sequence_in,real_naturals_closed_succ}.
    Therefore $(u_1(N + 1),u_1((N + 1) + 1))\in R_1$ and
        $(u_2(N + 1),u_2((N + 1) + 1))\in R_2$
        by \cref{funs_type_apply,subseteq,relation}.
    Take $k_1,J_1,H_1,a_1,b_1,d_1$ such that
        $(u_1(N + 1),u_1((N + 1) + 1))\in G\times G$ and
        $k_1\in\naturalsPlus$ and
        $\fst{(u_1(N + 1),u_1((N + 1) + 1))} = (k_1,J_1)$ and
        $\snd{(u_1(N + 1),u_1((N + 1) + 1))} = (k_1 + 1,H_1)$ and
        $d_1 = x_1(k_1)$ and
        $a_1\in\reals$ and
        $b_1\in\reals$ and
        $a_1 \leq b_1$ and
        $J_1 = \intervalclosed{a_1}{b_1}$ and
        $H_1\subseteq J_1$ and
        $(d_1,H_1)\in\{(0,\intervalclosed{a_1}{(a_1 + a_1 + b_1) / (1 + 1 + 1)}),
            (1,\intervalclosed{(a_1 + b_1 + b_1) / (1 + 1 + 1)}{b_1})\}$
        by \cref{subseteq}.
    Take $k_2,J_2,H_2,a_2,b_2,d_2$ such that
        $(u_2(N + 1),u_2((N + 1) + 1))\in G\times G$ and
        $k_2\in\naturalsPlus$ and
        $\fst{(u_2(N + 1),u_2((N + 1) + 1))} = (k_2,J_2)$ and
        $\snd{(u_2(N + 1),u_2((N + 1) + 1))} = (k_2 + 1,H_2)$ and
        $d_2 = x_2(k_2)$ and
        $a_2\in\reals$ and
        $b_2\in\reals$ and
        $a_2 \leq b_2$ and
        $J_2 = \intervalclosed{a_2}{b_2}$ and
        $H_2\subseteq J_2$ and
        $(d_2,H_2)\in\{(0,\intervalclosed{a_2}{(a_2 + a_2 + b_2) / (1 + 1 + 1)}),
            (1,\intervalclosed{(a_2 + b_2 + b_2) / (1 + 1 + 1)}{b_2})\}$
        by \cref{subseteq}.
    Hence $k_1 = N + 1$ and $k_2 = N + 1$ and
        $\snd{u_1(N + 1)} = J_1$ and $\snd{u_2(N + 1)} = J_2$ and
        $\snd{u_1((N + 1) + 1)} = H_1$ and $\snd{u_2((N + 1) + 1)} = H_2$
        by \cref{fst_eq,snd_eq,pair_eq_iff,subseteq}.
    Therefore $d_1 = x_1(N + 1)$ and $d_2 = x_2(N + 1)$ by \cref{subseteq}.
    $d_1 = 0$ and $d_2 = 1$ by \cref{subseteq}.
    Thus $\intervalclosed{a_1}{b_1} = \intervalclosed{a}{b}$ and
        $\intervalclosed{a_2}{b_2} = \intervalclosed{a}{b}$ by \cref{subseteq}.
    $a_1 = a$ and $b_1 = b$
        by \cref{real_lt_imp_leq,space_filling_intervalclosed_endpoints_unique}.
    We have $a_2 = a$ and $b_2 = b$
        by \cref{real_lt_imp_leq,space_filling_intervalclosed_endpoints_unique}.
    $(d_1,H_1)\in\{(0,\intervalclosed{a}{(a + a + b) / (1 + 1 + 1)}),
        (1,\intervalclosed{(a + b + b) / (1 + 1 + 1)}{b})\}$
        by \cref{subseteq,real_add_eq_subst}.
    We have $(d_2,H_2)\in\{(0,\intervalclosed{a}{(a + a + b) / (1 + 1 + 1)}),
        (1,\intervalclosed{(a + b + b) / (1 + 1 + 1)}{b})\}$
        by \cref{subseteq,real_add_eq_subst}.
    Let $L = \intervalclosed{a}{(a + a + b) / (1 + 1 + 1)}$.
    Let $R = \intervalclosed{(a + b + b) / (1 + 1 + 1)}{b}$.
    $(d_1,L)\in\{(0,\intervalclosed{a}{(a + a + b) / (1 + 1 + 1)}),
        (1,\intervalclosed{(a + b + b) / (1 + 1 + 1)}{b})\}$ and
        $(d_2,R)\in\{(0,\intervalclosed{a}{(a + a + b) / (1 + 1 + 1)}),
            (1,\intervalclosed{(a + b + b) / (1 + 1 + 1)}{b})\}$
        by \cref{upair_intro_left,upair_intro_right,pair_eq_iff,subseteq}.
    Therefore $H_1 = L$ and $H_2 = R$
        by \cref{space_filling_retained_third_child_choice_unique,upair_intro_left,upair_intro_right,pair_eq_iff,subseteq,real_add_eq_subst}.
    $L$ is disjoint from $R$ by \cref{space_filling_retained_thirds_disjoint}.
    Suppose not.
    Therefore $y\in L$ and $y\in R$ by \cref{subseteq}.
    Contradiction by \cref{disjoint}.
\end{proof}

% from §44 helper lemma 112b: the local shared-interval digit separation is symmetric
\begin{lemma}\label{space_filling_retained_third_shared_interval_digit_separation_local_rev}
    Let $I = \intervalclosed{0}{1}$.
    Let $K = \{J\in\pow{I} \mid \text{there exist $c,d\in\reals$ such that
        $J = \intervalclosed{c}{d}$ and
        $c \leq d$}\}$.
    Let $G = \naturalsPlus\times K$.
    Let $R_1 = \{p\in G\times G \mid \text{there exist $n,J,H,c,d,e$ such that
        $n\in\naturalsPlus$ and
        $\fst{p} = (n,J)$ and
        $\snd{p} = (n + 1,H)$ and
        $e = x_1(n)$ and
        $c\in\reals$ and
        $d\in\reals$ and
        $c \leq d$ and
        $J = \intervalclosed{c}{d}$ and
        $H\subseteq J$ and
        $(e,H)\in\{(0,\intervalclosed{c}{(c + c + d) / (1 + 1 + 1)}),
            (1,\intervalclosed{(c + d + d) / (1 + 1 + 1)}{d})\}$}\}$.
    Let $R_2 = \{p\in G\times G \mid \text{there exist $n,J,H,c,d,e$ such that
        $n\in\naturalsPlus$ and
        $\fst{p} = (n,J)$ and
        $\snd{p} = (n + 1,H)$ and
        $e = x_2(n)$ and
        $c\in\reals$ and
        $d\in\reals$ and
        $c \leq d$ and
        $J = \intervalclosed{c}{d}$ and
        $H\subseteq J$ and
        $(e,H)\in\{(0,\intervalclosed{c}{(c + c + d) / (1 + 1 + 1)}),
            (1,\intervalclosed{(c + d + d) / (1 + 1 + 1)}{d})\}$}\}$.
    Suppose $g_1\in\funs{G}{G}$.
    Suppose for all $z\in G$ we have $z\mathrel{R_1}\apply{g_1}{z}$.
    Suppose $u_1$ is a sequence in $G$.
    Suppose $u_1(1) = (1,I)$.
    Suppose for all $n\in\naturalsPlus$ we have $u_1(n + 1) = \apply{g_1}{u_1(n)}$.
    Suppose $\fst{u_1(N + 1)} = N + 1$.
    Suppose $g_2\in\funs{G}{G}$.
    Suppose for all $z\in G$ we have $z\mathrel{R_2}\apply{g_2}{z}$.
    Suppose $u_2$ is a sequence in $G$.
    Suppose $u_2(1) = (1,I)$.
    Suppose for all $n\in\naturalsPlus$ we have $u_2(n + 1) = \apply{g_2}{u_2(n)}$.
    Suppose $\fst{u_2(N + 1)} = N + 1$.
    Suppose $y\in I$ and for all $n\in\naturalsPlus$ we have $y\in\snd{u_1(n)}$.
    Suppose for all $n\in\naturalsPlus$ we have $y\in\snd{u_2(n)}$.
    Assume $N\in\naturalsPlus$.
    Assume $a\in\reals$ and $b\in\reals$ and $a < b$.
    Assume $\snd{u_1(N + 1)} = \intervalclosed{a}{b}$.
    Assume $\snd{u_2(N + 1)} = \intervalclosed{a}{b}$.
    Assume $x_1(N + 1) = 1$.
    Assume $x_2(N + 1) = 0$.
    Then $y\notin I$.
\end{lemma}
\begin{proof}
    Suppose not.
    Hence $y\in I$ by \cref{subseteq}.
    Therefore $y\notin I$ by \cref{space_filling_retained_third_shared_interval_digit_separation_local}.
    Contradiction by \cref{subseteq}.
\end{proof}

% from §44 helper lemma 113: the retained-third code-point assignment is injective
\begin{lemma}\label{space_filling_retained_third_code_map_injective}
    Let $I = \intervalclosed{0}{1}$.
    Let $D = \{0,1\}$.
    Let $E = \{(n,D) \mid n\in\naturalsPlus\}$.
    Suppose $r\in\funs{\productFamily{E}}{I}$.
    Suppose for all $x\in\productFamily{E}$ we have $r(x)$ is a retained-third code point of $x$ in $I$ and $E$.
    Then $r$ is injective.
\end{lemma}
\begin{proof}
    Show that for all $x_1,x_2$ such that
        $x_1\in\productFamily{E}$ and
        $x_2\in\productFamily{E}$ and
        $r(x_1) = r(x_2)$
    we have $x_1 = x_2$.
    \begin{subproof}
        Fix $x_1$.
        Fix $x_2$.
        Assume $x_1\in\productFamily{E}$ and
            $x_2\in\productFamily{E}$ and
            $r(x_1) = r(x_2)$.
        Let $y = r(x_1)$.
        Hence $y$ is a retained-third code point of $x_1$ in $I$ and $E$ and
            $y$ is a retained-third code point of $x_2$ in $I$ and $E$
            by \cref{subseteq}.
        Therefore $y\in I$ by \cref{space_filling_retained_third_code_point}.
        For all $n,u,v$ such that $(n,u)\in E$ and $(n,v)\in E$ we have $u = v$ by \cref{pair_eq_iff}.
        Therefore $E$ is a function on $\naturalsPlus$ and $\dom{E} = \naturalsPlus$
            by \cref{relation,rightunique,dom_intro,dom_iff,pair_eq_iff,subseteq,subseteq_antisymmetric,funs}.
        Suppose not.
        Take $w$ such that either $w\in x_1$ and $w\notin x_2$ or $w\notin x_1$ and $w\in x_2$
            by \cref{neq_witness}.
        We have $x_1$ is a function and $\dom{x_1} = \dom{E}$ and
            $x_2$ is a function and $\dom{x_2} = \dom{E}$ by \cref{indexed_product,funs_elim}.
        Therefore there exists $n\in\dom{E}$ such that $\apply{x_1}{n}\neq\apply{x_2}{n}$
            by \cref{function_member_elim,function_apply_elim,pair_eq_iff}.
        Take $k\in\dom{E}$ such that $\apply{x_1}{k}\neq\apply{x_2}{k}$ by \cref{subseteq}.
        Thus $k\in\naturalsPlus$ by \cref{subseteq}.
        Let $P = \{n\in\initialSegment{k} \mid x_1(n) \neq x_2(n)\}$.
        $k\in\initialSegment{k}$ by \cref{initial_segment_naturalsplus,real_one_leq_naturalsplus}.
        Hence $k\in P$ by \cref{subseteq}.
        Therefore $P$ is inhabited.
        $P\subseteq\initialSegment{k}$ by \cref{subseteq}.
        Therefore $P\subseteq\reals$
            by \cref{subseteq,initial_segment_naturalsplus,real_naturals_subset_reals,subseteq_transitive}.
        Thus $P$ is finite by \cref{initial_segment_finite,subseteq_finite_is_finite,subseteq}.
        Take $m$ such that $m$ is a smallest element of $P$ with respect to $\realOrderRel$
            by \cref{finite_inhabited_smallest_element,real_order_relation_simple}.
        Therefore $m\in P$ by \cref{smallest_element}.
        Hence $m\in\initialSegment{k}$ and $x_1(m) \neq x_2(m)$ by \cref{subseteq}.
        Thus $m\in\naturalsPlus$ by \cref{initial_segment_naturalsplus,subseteq}.
        Hence $x_1(m)\in E(m)$ and $x_2(m)\in E(m)$ by \cref{indexed_product,subseteq}.
        Therefore $x_1(m)\in D$ and $x_2(m)\in D$ by \cref{subseteq,function_apply_intro}.
        Thus $x_1(m) = 0$ or $x_1(m) = 1$ by \cref{upair_iff}.
        Therefore $x_2(m) = 0$ or $x_2(m) = 1$ by \cref{upair_iff}.
        Let $K = \{J\in\pow{I} \mid \text{there exist $a,b\in\reals$ such that
            $J = \intervalclosed{a}{b}$ and
            $a \leq b$}\}$.
        Let $G = \naturalsPlus\times K$.
        Let $R = \{p\in G\times G \mid \text{there exist $n,J,H,a,b,d$ such that
            $n\in\naturalsPlus$ and
            $\fst{p} = (n,J)$ and
            $\snd{p} = (n + 1,H)$ and
            $a\in\reals$ and
            $b\in\reals$ and
            $a \leq b$ and
            $J = \intervalclosed{a}{b}$ and
            $H\subseteq J$ and
            $(d,H)\in\{(0,\intervalclosed{a}{(a + a + b) / (1 + 1 + 1)}),
                (1,\intervalclosed{(a + b + b) / (1 + 1 + 1)}{b})\}$}\}$.
        Let $R_1 = \{p\in G\times G \mid \text{there exist $n,J,H,a,b,d$ such that
            $n\in\naturalsPlus$ and
            $\fst{p} = (n,J)$ and
            $\snd{p} = (n + 1,H)$ and
            $d = x_1(n)$ and
            $a\in\reals$ and
            $b\in\reals$ and
            $a \leq b$ and
            $J = \intervalclosed{a}{b}$ and
            $H\subseteq J$ and
            $(d,H)\in\{(0,\intervalclosed{a}{(a + a + b) / (1 + 1 + 1)}),
                (1,\intervalclosed{(a + b + b) / (1 + 1 + 1)}{b})\}$}\}$.
        Let $R_2 = \{p\in G\times G \mid \text{there exist $n,J,H,a,b,d$ such that
            $n\in\naturalsPlus$ and
            $\fst{p} = (n,J)$ and
            $\snd{p} = (n + 1,H)$ and
            $d = x_2(n)$ and
            $a\in\reals$ and
            $b\in\reals$ and
            $a \leq b$ and
            $J = \intervalclosed{a}{b}$ and
            $H\subseteq J$ and
            $(d,H)\in\{(0,\intervalclosed{a}{(a + a + b) / (1 + 1 + 1)}),
                (1,\intervalclosed{(a + b + b) / (1 + 1 + 1)}{b})\}$}\}$.
        Take $g_1$ such that
            $g_1\in\funs{G}{G}$ and
            for all $z\in G$ we have $z\mathrel{R_1}\apply{g_1}{z}$
            by \cref{space_filling_retained_third_state_selector}.
        Therefore $g_1$ is a retained-third state selector on $G$ and $R$
            by \cref{subseteq,space_filling_retained_third_state_selector_on}.
        Take $u_1$ such that
            $u_1$ is a sequence in $G$ and
            $u_1(1) = (1,I)$ and
            for all $n\in\naturalsPlus$ we have $u_1(n + 1) = \apply{g_1}{u_1(n)}$
            by \cref{space_filling_dyadic_child_state_iteration}.
        Thus $u_1$ is a retained-third state orbit of $g_1$ on $G$ from $I$
            by \cref{space_filling_retained_third_state_orbit_from}.
        Hence for all $n\in\naturalsPlus$ we have $y\in\snd{u_1(n)}$
            by \cref{space_filling_retained_third_code_point_on_standard_orbit}.
        Take $g_2$ such that
            $g_2\in\funs{G}{G}$ and
            for all $z\in G$ we have $z\mathrel{R_2}\apply{g_2}{z}$
            by \cref{space_filling_retained_third_state_selector}.
        Therefore $g_2$ is a retained-third state selector on $G$ and $R$
            by \cref{subseteq,space_filling_retained_third_state_selector_on}.
        Take $u_2$ such that
            $u_2$ is a sequence in $G$ and
            $u_2(1) = (1,I)$ and
            for all $n\in\naturalsPlus$ we have $u_2(n + 1) = \apply{g_2}{u_2(n)}$
            by \cref{space_filling_dyadic_child_state_iteration}.
        Thus $u_2$ is a retained-third state orbit of $g_2$ on $G$ from $I$
            by \cref{space_filling_retained_third_state_orbit_from}.
        Hence for all $n\in\naturalsPlus$ we have $y\in\snd{u_2(n)}$
            by \cref{space_filling_retained_third_code_point_on_standard_orbit}.
        $m = 1$ or there exists $p\in\naturalsPlus$ such that $m = p + 1$
            by \cref{real_naturals_predecessor}.
        \begin{byCase}
        \caseOf{$m = 1$.}
            We have $x_1(1) \neq x_2(1)$ by \cref{subseteq}.
            \begin{byCase}
            \caseOf{$x_1(1) = 0$.}
                We have $x_2(1) = 1$ by \cref{subseteq}.
                Therefore $y\notin I$ by \cref{space_filling_retained_third_first_digit_separation}.
                Contradiction by \cref{space_filling_retained_third_code_point}.
            \caseOf{$x_1(1) = 1$.}
                We have $x_2(1) = 0$ by \cref{subseteq}.
                Therefore $y\notin I$ by \cref{space_filling_retained_third_first_digit_separation}.
                Contradiction by \cref{space_filling_retained_third_code_point}.
            \end{byCase}
        \caseOf{there exists $p\in\naturalsPlus$ such that $m = p + 1$.}
            Take $p\in\naturalsPlus$ such that $m = p + 1$ by \cref{subseteq}.
            Show that for all $i\in\initialSegment{p}$ we have $x_1(i) = x_2(i)$.
            \begin{subproof}
                Fix $i\in\initialSegment{p}$.
                Hence $i\in\initialSegment{m}$ by \cref{initial_segment_subset_succ,subseteq}.
                Suppose not.
                Therefore $m\in\initialSegment{k}$ by \cref{subseteq}.
                Hence $i\in\initialSegment{k}$ by \cref{space_filling_initial_segment_transitive}.
                Thus $i\in P$ by \cref{subseteq}.
                Hence $m\mathrel{\realOrderRel} i$ or $m = i$ by \cref{smallest_element}.
                \begin{byCase}
                \caseOf{$m\mathrel{\realOrderRel} i$.}
                    We have $m < i$ by \cref{real_order_relation_elim}.
                    Hence $m \leq i$ by \cref{real_lt_imp_leq}.
                    Thus $i\in\naturalsPlus$ and $i \leq m$ by \cref{initial_segment_naturalsplus,subseteq}.
                    We have $m\in\reals$ and $i\in\reals$ by \cref{real_naturals_subset_reals,subseteq}.
                    Therefore $m = i$ by \cref{real_leq_antisym}.
                    $m\in\initialSegment{p}$ by \cref{subseteq}.
                    Therefore $m \leq p$ by \cref{initial_segment_naturalsplus,subseteq}.
                    Hence $p + 1 \leq p$ by \cref{subseteq}.
                    Thus $p < p + 1$ by \cref{real_naturals_lt_succ_local}.
                    We have $p\in\reals$ and $p + 1\in\reals$
                        by \cref{real_naturals_subset_reals,subseteq,real_one_in_naturals,real_add_closed}.
                    Therefore $p + 1 < p + 1$ by \cref{real_leq_lt_trans}.
                    Contradiction by \cref{real_lt_irreflexive}.
                \caseOf{$m = i$.}
                    $m\in\initialSegment{p}$ by \cref{subseteq}.
                    Therefore $m \leq p$ by \cref{initial_segment_naturalsplus,subseteq}.
                    $p + 1 \leq p$ by \cref{subseteq}.
                    Thus $p < p + 1$ by \cref{real_naturals_lt_succ_local}.
                    We have $p\in\reals$ and $p + 1\in\reals$
                        by \cref{real_naturals_subset_reals,subseteq,real_one_in_naturals,real_add_closed}.
                    Therefore $p + 1 < p + 1$ by \cref{real_leq_lt_trans}.
                    Contradiction by \cref{real_lt_irreflexive}.
                \end{byCase}
            \end{subproof}
            Hence for all $n\in\initialSegment{p}$ we have $\snd{u_1(n)} = \snd{u_2(n)}$
                by \cref{space_filling_retained_third_state_orbit_prefix_agreement}.
            We have $p\in\initialSegment{p}$ by \cref{initial_segment_naturalsplus}.
            Therefore $\snd{u_1(p)} = \snd{u_2(p)}$ by \cref{subseteq}.
            Let $A_1 = \{n\in\naturalsPlus \mid \fst{u_1(n)} = n\}$.
            Hence $1\in A_1$ by \cref{real_one_in_naturals,pair_eq_iff,fst_eq,subseteq}.
            Show that for all $n\in A_1$ we have $n + 1\in A_1$.
            \begin{subproof}
                Follows by \cref{subseteq,sequence_in,funs_type_apply,relation,fst_eq,snd_eq,pair_eq_iff,real_add_eq_subst,real_naturals_closed_succ}.
            \end{subproof}
            Therefore $p\in A_1$ by \cref{real_naturals_induction,subseteq}.
            Thus $\fst{u_1(p)} = p$ by \cref{subseteq}.
            Let $A_2 = \{n\in\naturalsPlus \mid \fst{u_2(n)} = n\}$.
            Hence $1\in A_2$ by \cref{real_one_in_naturals,pair_eq_iff,fst_eq,subseteq}.
            Show that for all $n\in A_2$ we have $n + 1\in A_2$.
            \begin{subproof}
                Follows by \cref{subseteq,sequence_in,funs_type_apply,relation,fst_eq,snd_eq,pair_eq_iff,real_add_eq_subst,real_naturals_closed_succ}.
            \end{subproof}
            Therefore $p\in A_2$ by \cref{real_naturals_induction,subseteq}.
            Thus $\fst{u_2(p)} = p$ by \cref{subseteq}.
            We have $u_1(p)\in G$ and $u_1(p + 1)\in G$ and
                $u_2(p)\in G$ and $u_2(p + 1)\in G$
                by \cref{sequence_in,funs_type_apply,real_naturals_closed_succ}.
            Therefore $(u_1(p),u_1(p + 1))\in R_1$ and
                $(u_2(p),u_2(p + 1))\in R_2$
                by \cref{subseteq,relation}.
            Take $k_1,J_1,H_1,a_1,b_1,d_1$ such that
                $(u_1(p),u_1(p + 1))\in G\times G$ and
                $k_1\in\naturalsPlus$ and
                $\fst{(u_1(p),u_1(p + 1))} = (k_1,J_1)$ and
                $\snd{(u_1(p),u_1(p + 1))} = (k_1 + 1,H_1)$ and
                $d_1 = x_1(k_1)$ and
                $a_1\in\reals$ and
                $b_1\in\reals$ and
                $a_1 \leq b_1$ and
                $J_1 = \intervalclosed{a_1}{b_1}$ and
                $H_1\subseteq J_1$ and
                $(d_1,H_1)\in\{(0,\intervalclosed{a_1}{(a_1 + a_1 + b_1) / (1 + 1 + 1)}),
                    (1,\intervalclosed{(a_1 + b_1 + b_1) / (1 + 1 + 1)}{b_1})\}$
                by \cref{subseteq}.
            Take $k_2,J_2,H_2,a_2,b_2,d_2$ such that
                $(u_2(p),u_2(p + 1))\in G\times G$ and
                $k_2\in\naturalsPlus$ and
                $\fst{(u_2(p),u_2(p + 1))} = (k_2,J_2)$ and
                $\snd{(u_2(p),u_2(p + 1))} = (k_2 + 1,H_2)$ and
                $d_2 = x_2(k_2)$ and
                $a_2\in\reals$ and
                $b_2\in\reals$ and
                $a_2 \leq b_2$ and
                $J_2 = \intervalclosed{a_2}{b_2}$ and
                $H_2\subseteq J_2$ and
                $(d_2,H_2)\in\{(0,\intervalclosed{a_2}{(a_2 + a_2 + b_2) / (1 + 1 + 1)}),
                    (1,\intervalclosed{(a_2 + b_2 + b_2) / (1 + 1 + 1)}{b_2})\}$
                by \cref{subseteq}.
            Therefore $u_1(p) = (k_1,J_1)$ and $u_1(p + 1) = (k_1 + 1,H_1)$ and
                $u_2(p) = (k_2,J_2)$ and $u_2(p + 1) = (k_2 + 1,H_2)$
                by \cref{fst_eq,snd_eq,pair_eq_iff}.
            Thus $\fst{u_1(p)} = k_1$ and $\fst{u_2(p)} = k_2$ and
                $\snd{u_1(p)} = J_1$ and $\snd{u_2(p)} = J_2$ and
                $\snd{u_1(p + 1)} = H_1$ and $\snd{u_2(p + 1)} = H_2$
                by \cref{fst_eq,snd_eq,pair_eq_iff}.
            Hence $k_1 = p$ and $k_2 = p$ by \cref{subseteq}.
            $d_1 = x_1(p)$ and $d_2 = x_2(p)$ by \cref{subseteq}.
            Therefore $x_1(p) = x_2(p)$ by \cref{initial_segment_naturalsplus,subseteq}.
            Hence $d_1 = d_2$ by \cref{subseteq}.
            Thus $J_1 = J_2$ by \cref{subseteq}.
            Hence $\intervalclosed{a_1}{b_1} = \intervalclosed{a_2}{b_2}$ by \cref{subseteq}.
            Thus $a_1 = a_2$ and $b_1 = b_2$ by \cref{space_filling_intervalclosed_endpoints_unique}.
            Therefore $(d_2,H_1)\in\{(0,\intervalclosed{a_2}{(a_2 + a_2 + b_2) / (1 + 1 + 1)}),
                (1,\intervalclosed{(a_2 + b_2 + b_2) / (1 + 1 + 1)}{b_2})\}$
                by \cref{subseteq,real_add_eq_subst}.
            Thus $d_2\in D$ by \cref{upair_iff,subseteq,pair_eq_iff}.
            Hence $H_1 = H_2$ by \cref{space_filling_retained_third_child_choice_unique,subseteq,real_add_eq_subst}.
            Therefore $\snd{u_1(p + 1)} = \snd{u_2(p + 1)}$ by \cref{subseteq}.
            Thus $u_1(p + 1)\in G$ by \cref{sequence_in,funs_type_apply,real_naturals_closed_succ}.
            Take $t,J$ such that
                $u_1(p + 1) = (t,J)$ and
                $t\in\naturalsPlus$ and
                $J\in K$
                by \cref{subseteq,times_elem_is_tuple}.
            Hence $u_1(p + 1) = (k_1 + 1,H_1)$ by \cref{subseteq}.
            Therefore $t = k_1 + 1$ and $J = H_1$ by \cref{pair_eq_iff}.
            Thus $t = p + 1$ by \cref{subseteq,real_add_eq_subst}.
            Take $a,b$ such that
                $a\in\reals$ and
                $b\in\reals$ and
                $J = \intervalclosed{a}{b}$ and
                $a \leq b$
                by \cref{subseteq}.
            Hence $u_1(p + 1) = (t,\intervalclosed{a}{b})$ by \cref{subseteq}.
            Therefore $\snd{u_1(p + 1)} = \intervalclosed{a}{b}$ by \cref{snd_eq,pair_eq_iff}.
            Thus $\snd{u_2(p + 1)} = \intervalclosed{a}{b}$ by \cref{subseteq}.
            Hence $u_1(p + 1) = (p + 1,\intervalclosed{a}{b})$ by \cref{subseteq}.
            Therefore $u_2(p + 1) = (p + 1,\intervalclosed{a}{b})$ by \cref{subseteq,real_add_eq_subst}.
            Thus $\fst{u_1(p + 1)} = p + 1$ and $\fst{u_2(p + 1)} = p + 1$
                by \cref{fst_eq,pair_eq_iff}.
            Hence $p + 1\in\naturalsPlus$ by \cref{real_naturals_closed_succ}.
            Therefore $a < b$ by \cref{space_filling_retained_third_state_orbit_succ_interval_strict_local,subseteq}.
            Thus $x_1(p + 1) \neq x_2(p + 1)$ by \cref{subseteq}.
            We have $u_1\in\funs{\naturalsPlus}{G}$ and $u_2\in\funs{\naturalsPlus}{G}$ and
                $(p + 1) + 1\in\naturalsPlus$
                by \cref{sequence_in,real_naturals_closed_succ}.
            Therefore $(u_1(p + 1),u_1((p + 1) + 1))\in R_1$ and
                $(u_2(p + 1),u_2((p + 1) + 1))\in R_2$
                by \cref{funs_type_apply,subseteq,relation}.
            Take $k_3,J_3,H_3,a_3,b_3,d_3$ such that
                $(u_1(p + 1),u_1((p + 1) + 1))\in G\times G$ and
                $k_3\in\naturalsPlus$ and
                $\fst{(u_1(p + 1),u_1((p + 1) + 1))} = (k_3,J_3)$ and
                $\snd{(u_1(p + 1),u_1((p + 1) + 1))} = (k_3 + 1,H_3)$ and
                $d_3 = x_1(k_3)$ and
                $a_3\in\reals$ and
                $b_3\in\reals$ and
                $a_3 \leq b_3$ and
                $J_3 = \intervalclosed{a_3}{b_3}$ and
                $H_3\subseteq J_3$ and
                $(d_3,H_3)\in\{(0,\intervalclosed{a_3}{(a_3 + a_3 + b_3) / (1 + 1 + 1)}),
                    (1,\intervalclosed{(a_3 + b_3 + b_3) / (1 + 1 + 1)}{b_3})\}$
                by \cref{subseteq}.
            Take $k_4,J_4,H_4,a_4,b_4,d_4$ such that
                $(u_2(p + 1),u_2((p + 1) + 1))\in G\times G$ and
                $k_4\in\naturalsPlus$ and
                $\fst{(u_2(p + 1),u_2((p + 1) + 1))} = (k_4,J_4)$ and
                $\snd{(u_2(p + 1),u_2((p + 1) + 1))} = (k_4 + 1,H_4)$ and
                $d_4 = x_2(k_4)$ and
                $a_4\in\reals$ and
                $b_4\in\reals$ and
                $a_4 \leq b_4$ and
                $J_4 = \intervalclosed{a_4}{b_4}$ and
                $H_4\subseteq J_4$ and
                $(d_4,H_4)\in\{(0,\intervalclosed{a_4}{(a_4 + a_4 + b_4) / (1 + 1 + 1)}),
                    (1,\intervalclosed{(a_4 + b_4 + b_4) / (1 + 1 + 1)}{b_4})\}$
                by \cref{subseteq}.
            Therefore $k_3 = p + 1$ and $k_4 = p + 1$ and
                $\snd{u_1(p + 1)} = J_3$ and $\snd{u_2(p + 1)} = J_4$ and
                $\snd{u_1((p + 1) + 1)} = H_3$ and $\snd{u_2((p + 1) + 1)} = H_4$
                by \cref{fst_eq,snd_eq,pair_eq_iff,subseteq}.
            Thus $d_3 = x_1(p + 1)$ and $d_4 = x_2(p + 1)$ by \cref{subseteq}.
            Hence $\intervalclosed{a_3}{b_3} = \intervalclosed{a}{b}$ and
                $\intervalclosed{a_4}{b_4} = \intervalclosed{a}{b}$ by \cref{subseteq}.
            Thus $a_3 = a$ and $b_3 = b$
                by \cref{real_lt_imp_leq,space_filling_intervalclosed_endpoints_unique}.
            Hence $a_4 = a$ and $b_4 = b$
                by \cref{real_lt_imp_leq,space_filling_intervalclosed_endpoints_unique}.
            Therefore $(d_3,H_3)\in\{(0,\intervalclosed{a}{(a + a + b) / (1 + 1 + 1)}),
                (1,\intervalclosed{(a + b + b) / (1 + 1 + 1)}{b})\}$ and
                $(d_4,H_4)\in\{(0,\intervalclosed{a}{(a + a + b) / (1 + 1 + 1)}),
                    (1,\intervalclosed{(a + b + b) / (1 + 1 + 1)}{b})\}$
                by \cref{subseteq,real_add_eq_subst}.
            We have $y\in\snd{u_1((p + 1) + 1)}$ and $y\in\snd{u_2((p + 1) + 1)}$ by \cref{subseteq}.
            Therefore $y\in H_3$ and $y\in H_4$ by \cref{subseteq}.
            Let $L = \intervalclosed{a}{(a + a + b) / (1 + 1 + 1)}$.
            Let $R_0 = \intervalclosed{(a + b + b) / (1 + 1 + 1)}{b}$.
            Thus $L$ is disjoint from $R_0$ by \cref{space_filling_retained_thirds_disjoint}.
            \begin{byCase}
            \caseOf{$x_1(m) = 0$.}
                We have $x_1(p + 1) = 0$ and $x_2(p + 1) = 1$ by \cref{subseteq}.
                Thus $d_3 = 0$ and $d_4 = 1$ by \cref{subseteq}.
                Therefore $(d_3,L)\in\{(0,\intervalclosed{a}{(a + a + b) / (1 + 1 + 1)}),
                    (1,\intervalclosed{(a + b + b) / (1 + 1 + 1)}{b})\}$ and
                    $(d_4,R_0)\in\{(0,\intervalclosed{a}{(a + a + b) / (1 + 1 + 1)}),
                        (1,\intervalclosed{(a + b + b) / (1 + 1 + 1)}{b})\}$
                    by \cref{upair_intro_left,upair_intro_right,pair_eq_iff,subseteq}.
                Hence $H_3 = L$ and $H_4 = R_0$
                    by \cref{space_filling_retained_third_child_choice_unique,upair_intro_left,upair_intro_right,pair_eq_iff,subseteq,real_add_eq_subst}.
                Therefore $y\in L$ and $y\in R_0$ by \cref{subseteq}.
                Contradiction by \cref{disjoint}.
            \caseOf{$x_1(m) = 1$.}
                We have $x_2(p + 1) = 0$ and $x_1(p + 1) = 1$ by \cref{subseteq}.
                Thus $d_4 = 0$ and $d_3 = 1$ by \cref{subseteq}.
                Therefore $(d_4,L)\in\{(0,\intervalclosed{a}{(a + a + b) / (1 + 1 + 1)}),
                    (1,\intervalclosed{(a + b + b) / (1 + 1 + 1)}{b})\}$ and
                    $(d_3,R_0)\in\{(0,\intervalclosed{a}{(a + a + b) / (1 + 1 + 1)}),
                        (1,\intervalclosed{(a + b + b) / (1 + 1 + 1)}{b})\}$
                    by \cref{upair_intro_left,upair_intro_right,pair_eq_iff,subseteq}.
                Hence $H_4 = L$ and $H_3 = R_0$
                    by \cref{space_filling_retained_third_child_choice_unique,upair_intro_left,upair_intro_right,pair_eq_iff,subseteq,real_add_eq_subst}.
                Therefore $y\in L$ and $y\in R_0$ by \cref{subseteq}.
                Contradiction by \cref{disjoint}.
            \end{byCase}
        \end{byCase}
    \end{subproof}
    We have $r$ is a function and $\dom{r} = \productFamily{E}$ by \cref{funs_elim,funs_is_function}.
    Therefore $r$ is injective by \cref{funs_is_function,injective_function}.
\end{proof}

% from §44 helper lemma 118: ordered dyadic difference bounds the standard metric
\begin{lemma}\label{space_filling_standard_metric_lt_from_ordered_difference_local}
    Suppose $d = \standardMetricR$.
    Suppose $x\in\reals$ and $y\in\reals$ and $\epsilon\in\reals$.
    Suppose $x < y$ and $y - x < \epsilon$.
    Then $d((x,y)) < \epsilon$.
\end{lemma}
\begin{proof}
    We have $y - x\in\reals$ by \cref{real_sub,real_add_inverse,real_add_closed}.
    Hence $0 < \epsilon$ by \cref{real_lt_imp_pos_diff,real_add_ident,real_lt_trans}.
    Thus $x - \epsilon < y$
        by \cref{real_sub,real_add_inverse,real_add_closed,real_sub_pos_lt,real_lt_trans}.
    Therefore $y < x + \epsilon$
        by \cref{real_lt_add,real_add_comm,real_lt_subst_left,real_lt_subst_right,real_add_sub_cancel_right}.
    Hence $\abs{x - y} < \epsilon$ by \cref{real_abs_lt_from_interval}.
    Therefore $d((x,y)) < \epsilon$ by \cref{standard_metric_value,real_lt_subst_left}.
\end{proof}

% from §44 helper lemma 114: a chosen dyadic code-point assignment is continuous
\begin{lemma}\label{space_filling_dyadic_code_map_continuous}
    Let $I = \intervalclosed{0}{1}$.
    Let $T_I = \subspaceTopology{\standardTopologyR}{I}$.
    Let $D = \{0,1\}$.
    Let $T_D = \subspaceTopology{T_I}{D}$.
    Let $E = \{(n,D) \mid n\in\naturalsPlus\}$.
    Let $U = \{(n,T_D) \mid n\in\naturalsPlus\}$.
    Let $P_B = \productTopologyIndexed{U}{E}$.
    Suppose $r\in\funs{\productFamily{E}}{I}$.
    Suppose for all $x\in\productFamily{E}$ we have $r(x)$ is a dyadic code point of $x$ in $I$ and $E$.
    Then $r$ is continuous from $\productFamily{E}$ to $I$ with respect to $P_B$ and $T_I$.
\end{lemma}
\begin{proof}
    $I\subseteq\reals$ by \cref{intervalclosed,subseteq}.
    $\standardTopologyR$ is a topology on $\reals$
        by \cref{standard_basis_is_basis,basis_topology_is_topology,standard_topology_reals}.
    Hence $I$ is a subspace of $\reals$ with respect to $\standardTopologyR$ by \cref{subspace_of}.
    Therefore $T_I$ is a topology on $I$ by \cref{subspace_topology_is_topology}.
    $\productFamily{E}$ is compact with respect to $P_B$ and
        $\productFamily{E}$ is Hausdorff with respect to $P_B$
        by \cref{space_filling_binary_product_compact_hausdorff}.
    Hence $P_B$ is a topology on $\productFamily{E}$ by \cref{hausdorff}.
    We have $E$ is a function on $\naturalsPlus$
        by \cref{relation,rightunique,dom_intro,dom_iff,pair_eq_iff,subseteq,subseteq_antisymmetric}.
    Therefore $\dom{E} = \naturalsPlus$ by \cref{funs}.
    For all $n\in\naturalsPlus$ we have $E(n) = D$ by \cref{function_apply_intro}.
    Hence $r\in\funs{\productFamily{E}}{\reals}$ by \cref{funs_weaken_codom,subseteq}.
    Therefore $r$ is a function and $\dom{r} = \productFamily{E}$ by \cref{funs_elim,funs_is_function}.

    Show that for all $x\in\productFamily{E}$ we have
        $r$ is continuous at $x$ from $\productFamily{E}$ to $\reals$ with respect to $P_B$ and $\standardTopologyR$.
    \begin{subproof}
        Fix $x\in\productFamily{E}$.
        Let $y = r(x)$.
        Hence $y$ is a dyadic code point of $x$ in $I$ and $E$ by \cref{subseteq}.
        Take $K_1,G_1,R_1,g_1,u_1$ such that
            $K_1 = \{J\in\pow{I} \mid \text{there exist $a,b\in\reals$ such that
                $J = \intervalclosed{a}{b}$ and
                $a \leq b$}\}$ and
            $G_1 = \naturalsPlus\times K_1$ and
            $R_1 = \{p\in G_1\times G_1 \mid \text{there exist $n,J,H,a,b,d$ such that
                $n\in\naturalsPlus$ and
                $\fst{p} = (n,J)$ and
                $\snd{p} = (n + 1,H)$ and
                $d = x(n)$ and
                $a\in\reals$ and
                $b\in\reals$ and
                $a \leq b$ and
                $J = \intervalclosed{a}{b}$ and
                $H\subseteq J$ and
                $(d,H)\in\{(0,\intervalclosed{a}{(a + b) / (1 + 1)}),
                    (1,\intervalclosed{(a + b) / (1 + 1)}{b})\}$}\}$ and
            $g_1$ is a retained-third state selector on $G_1$ and $R_1$ and
            $u_1$ is a retained-third state orbit of $g_1$ on $G_1$ from $I$ and
            for all $n\in\naturalsPlus$ we have $y\in\snd{u_1(n)}$
            by \cref{space_filling_dyadic_code_point}.
        Therefore $g_1\in\funs{G_1}{G_1}$ and
            for all $z\in G_1$ we have $z\mathrel{R_1}\apply{g_1}{z}$
            by \cref{space_filling_retained_third_state_selector_on}.
        Hence $u_1$ is a sequence in $G_1$ and
            $u_1(1) = (1,I)$ and
            for all $n\in\naturalsPlus$ we have $u_1(n + 1) = \apply{g_1}{u_1(n)}$
            by \cref{space_filling_retained_third_state_orbit_from}.

        Show that for all $V$ such that
            $V$ is open in $\reals$ with respect to $\standardTopologyR$ and
            $y\in V$
        there exists $W$ such that
            $W$ is open in $\productFamily{E}$ with respect to $P_B$ and
            $x\in W$ and
            $\img{r}{W}\subseteq V$.
        \begin{subproof}
            Fix $V$ such that
                $V$ is open in $\reals$ with respect to $\standardTopologyR$ and
                $y\in V$.
            Let $d = \standardMetricR$.
            $\metricTopology{d}{\reals} = \standardTopologyR$ by \cref{standard_metric_topology}.
            Therefore $V\subseteq\reals$ by \cref{open_in,topology_on,subseteq,pow_iff,standard_basis_is_basis,basis_topology_is_topology,standard_topology_reals}.
            $d$ is a metric on $\reals$ by \cref{standard_metric_is_metric}.
            Take $\epsilon\in\reals$ such that
                $0 < \epsilon$ and
                $\metricBall{d}{\reals}{y}{\epsilon}\subseteq V$
                by \cref{standard_metric_topology,metric_open_iff,neighborhood}.
            Take $N\in\naturalsPlus$ such that
                $0 < 1 / N$ and
                $1 / N < \epsilon$
                by \cref{real_small_reciprocal}.
            Hence $u_1(N)\in G_1$ by \cref{sequence_in,funs_type_apply}.
            Take $m_0,J_0$ such that
                $u_1(N) = (m_0,J_0)$ and
                $m_0\in\naturalsPlus$ and
                $J_0\in K_1$
                by \cref{subseteq,times_elem_is_tuple}.
            Take $p_0,q_0$ such that
                $p_0\in\reals$ and
                $q_0\in\reals$ and
                $J_0 = \intervalclosed{p_0}{q_0}$ and
                $p_0 \leq q_0$
                by \cref{subseteq}.
            Therefore $\snd{u_1(N)} = \intervalclosed{p_0}{q_0}$ by \cref{snd_eq,pair_eq_iff,subseteq}.
            Hence $q_0 - p_0 \leq 1 / N$ by \cref{space_filling_dyadic_stage_interval_width_leq_local}.
            Thus $y\in\intervalclosed{p_0}{q_0}$ by \cref{subseteq}.
            Therefore $x$ is a function and $\dom{x} = \dom{E}$ by \cref{indexed_product,funs_elim}.
            Let $a = \restrl{x}{\initialSegment{N}}$.
            $a$ is a function by \cref{restrl_of_function_is_function}.
            Hence $\dom{a} = \dom{x}\inter\initialSegment{N}$ by \cref{dom_of_restrl_eq_inter}.
            Therefore $\dom{x} = \naturalsPlus$ by \cref{subseteq}.
            Hence $\initialSegment{N}\subseteq\dom{x}$ by \cref{initial_segment_naturalsplus,subseteq}.
            Therefore $\dom{a} = \initialSegment{N}$ by \cref{inter_absorb_supseteq_left}.
            For all $i\in\dom{a}$ we have $i\in\initialSegment{N}$ by \cref{subseteq}.
            Hence for all $i\in\dom{a}$ we have $i\in\naturalsPlus$ by \cref{initial_segment_naturalsplus,subseteq}.
            Therefore for all $i\in\dom{a}$ we have $i\in\dom{E}$ by \cref{subseteq}.
            Hence for all $i\in\dom{a}$ we have $x(i)\in E(i)$ by \cref{indexed_product,subseteq}.
            For all $i\in\dom{a}$ we have $E(i) = D$ by \cref{function_apply_intro,subseteq}.
            Therefore for all $i\in\dom{a}$ we have $x(i)\in D$ by \cref{subseteq}.
            For all $i\in\dom{a}$ we have $a(i) = x(i)$ by \cref{restrl_of_function_apply,subseteq}.
            Hence for all $i\in\dom{a}$ we have $a(i)\in D$ by \cref{subseteq}.
            Therefore $a\in\funs{\initialSegment{N}}{D}$ by \cref{funs_intro}.
            Let $W = \{z\in\productFamily{E} \mid \text{for all $i\in\dom{a}$ we have $z(i) = a(i)$}\}$.
            Hence $W$ is the binary prefix cylinder at $a$ in $E$ by \cref{space_filling_binary_prefix_cylinder}.
            Therefore $W$ is open in $\productFamily{E}$ with respect to $P_B$
                by \cref{space_filling_binary_prefix_cylinder_open}.
            Hence $x\in W$ by \cref{subseteq,restrl_of_function_apply}.
            Show that for all $u$ such that $u\in\img{r}{W}$ we have $u\in V$.
            \begin{subproof}
                    Fix $u$ such that $u\in\img{r}{W}$.
                    Take $z\in\dom{r}\inter W$ such that $u = r(z)$ by \cref{img_of_function_elim}.
                    Therefore $z\in W$ by \cref{inter_elim_right}.
                    Hence $z\in\productFamily{E}$ and
                        for all $i\in\dom{a}$ we have $z(i) = a(i)$ by \cref{subseteq}.
                    Hence for all $i\in\initialSegment{N}$ we have $z(i) = x(i)$
                        by \cref{subseteq,restrl_of_function_apply}.
                    Therefore $u$ is a dyadic code point of $z$ in $I$ and $E$ by \cref{subseteq}.
                    Let $K_2 = \{J\in\pow{I} \mid \text{there exist $a_2,b_2\in\reals$ such that
                        $J = \intervalclosed{a_2}{b_2}$ and
                        $a_2 \leq b_2$}\}$.
                    Let $G_2 = \naturalsPlus\times K_2$.
                    Let $R_2 = \{p\in G_2\times G_2 \mid \text{there exist $n,J,H,a_2,b_2,d_2$ such that
                        $n\in\naturalsPlus$ and
                        $\fst{p} = (n,J)$ and
                        $\snd{p} = (n + 1,H)$ and
                        $d_2 = z(n)$ and
                        $a_2\in\reals$ and
                        $b_2\in\reals$ and
                        $a_2 \leq b_2$ and
                        $J = \intervalclosed{a_2}{b_2}$ and
                        $H\subseteq J$ and
                        $(d_2,H)\in\{(0,\intervalclosed{a_2}{(a_2 + b_2) / (1 + 1)}),
                            (1,\intervalclosed{(a_2 + b_2) / (1 + 1)}{b_2})\}$}\}$.
                    Take $g_2$ such that
                        $g_2\in\funs{G_2}{G_2}$ and
                        for all $w\in G_2$ we have $w\mathrel{R_2}\apply{g_2}{w}$
                        by \cref{space_filling_dyadic_child_state_selector}.
                    Take $u_2$ such that
                        $u_2$ is a sequence in $G_2$ and
                        $u_2(1) = (1,I)$ and
                        for all $n\in\naturalsPlus$ we have $u_2(n + 1) = \apply{g_2}{u_2(n)}$
                        by \cref{space_filling_dyadic_child_state_iteration}.
                    Take $K_3,G_3,R_5,g_3,v_3$ such that
                        $K_3 = \{J\in\pow{I} \mid \text{there exist $a_3,b_3\in\reals$ such that
                            $J = \intervalclosed{a_3}{b_3}$ and
                            $a_3 \leq b_3$}\}$ and
                        $G_3 = \naturalsPlus\times K_3$ and
                        $R_5 = \{p\in G_3\times G_3 \mid \text{there exist $n,J,H,a_3,b_3,d_3$ such that
                            $n\in\naturalsPlus$ and
                            $\fst{p} = (n,J)$ and
                            $\snd{p} = (n + 1,H)$ and
                            $d_3 = z(n)$ and
                            $a_3\in\reals$ and
                            $b_3\in\reals$ and
                            $a_3 \leq b_3$ and
                            $J = \intervalclosed{a_3}{b_3}$ and
                            $H\subseteq J$ and
                            $(d_3,H)\in\{(0,\intervalclosed{a_3}{(a_3 + b_3) / (1 + 1)}),
                                (1,\intervalclosed{(a_3 + b_3) / (1 + 1)}{b_3})\}$}\}$ and
                        $g_3$ is a retained-third state selector on $G_3$ and $R_5$ and
                        $v_3$ is a retained-third state orbit of $g_3$ on $G_3$ from $I$ and
                        for all $n\in\naturalsPlus$ we have $u\in\snd{v_3(n)}$
                        by \cref{space_filling_dyadic_code_point}.
                    Therefore $g_3\in\funs{G_3}{G_3}$ and
                        for all $w\in G_3$ we have $w\mathrel{R_5}\apply{g_3}{w}$
                        by \cref{space_filling_retained_third_state_selector_on}.
                    Hence $v_3$ is a sequence in $G_3$ and
                        $v_3(1) = (1,I)$ and
                        for all $n\in\naturalsPlus$ we have $v_3(n + 1) = \apply{g_3}{v_3(n)}$
                        by \cref{space_filling_retained_third_state_orbit_from}.
                    Therefore $K_3 = K_2$ by \cref{subseteq_antisymmetric,subseteq}.
                    Hence $G_3 = G_2$ by \cref{subseteq}.
                    $R_5\subseteq R_2$ and $R_2\subseteq R_5$ by \cref{subseteq}.
                    Therefore $R_5 = R_2$ by \cref{subseteq_antisymmetric}.
                    Hence $g_3\in\funs{G_2}{G_2}$ and
                        for all $w\in G_2$ we have $w\mathrel{R_2}\apply{g_3}{w}$
                        by \cref{subseteq}.
                    Therefore $v_3$ is a sequence in $G_2$ and
                        $v_3(1) = (1,I)$ and
                        for all $n\in\naturalsPlus$ we have $v_3(n + 1) = \apply{g_3}{v_3(n)}$
                        by \cref{subseteq}.
                    Hence for all $n\in\naturalsPlus$ we have $v_3(n) = u_2(n)$
                        by \cref{space_filling_dyadic_state_orbit_unique}.
                    Therefore for all $n\in\naturalsPlus$ we have $u\in\snd{u_2(n)}$
                        by \cref{subseteq}.
                    We have $K_2 = K_1$ by \cref{subseteq_antisymmetric,subseteq}.
                    Hence $G_2 = G_1$ by \cref{subseteq}.
                    Therefore $u_2$ is a sequence in $G_1$ and
                        $u_2(1) = (1,I)$ and
                        for all $n\in\naturalsPlus$ we have $u_2(n + 1) = \apply{g_2}{u_2(n)}$
                        by \cref{subseteq}.
                    Let $R_3 = \{p\in G_1\times G_1 \mid \text{there exist $n,J,H,a_3,b_3,d_3$ such that
                        $n\in\naturalsPlus$ and
                        $\fst{p} = (n,J)$ and
                        $\snd{p} = (n + 1,H)$ and
                        $d_3 = z(n)$ and
                        $a_3\in\reals$ and
                        $b_3\in\reals$ and
                        $a_3 \leq b_3$ and
                        $J = \intervalclosed{a_3}{b_3}$ and
                        $H\subseteq J$ and
                        $(d_3,H)\in\{(0,\intervalclosed{a_3}{(a_3 + b_3) / (1 + 1)}),
                            (1,\intervalclosed{(a_3 + b_3) / (1 + 1)}{b_3})\}$}\}$.
                    $R_2\subseteq R_3$ by \cref{subseteq}.
                    $R_3\subseteq R_2$ by \cref{subseteq}.
                    Therefore $R_2 = R_3$ by \cref{subseteq_antisymmetric}.
                    For all $w\in G_1$ we have $w\mathrel{R_2}\apply{g_2}{w}$ by \cref{subseteq}.
                    For all $w\in G_1$ we have $w\mathrel{R_3}\apply{g_2}{w}$ by \cref{subseteq}.
                    Let $R_4 = \{p\in G_1\times G_1 \mid \text{there exist $n,J,H,a_4,b_4,d_4$ such that
                        $n\in\naturalsPlus$ and
                        $\fst{p} = (n,J)$ and
                        $\snd{p} = (n + 1,H)$ and
                        $a_4\in\reals$ and
                        $b_4\in\reals$ and
                        $a_4 \leq b_4$ and
                        $J = \intervalclosed{a_4}{b_4}$ and
                        $H\subseteq J$ and
                        $(d_4,H)\in\{(0,\intervalclosed{a_4}{(a_4 + b_4) / (1 + 1)}),
                            (1,\intervalclosed{(a_4 + b_4) / (1 + 1)}{b_4})\}$}\}$.
                    For all $w\in G_1$ we have $w\mathrel{R_4}\apply{g_1}{w}$ by \cref{subseteq}.
                    For all $w\in G_1$ we have $w\mathrel{R_4}\apply{g_2}{w}$ by \cref{subseteq}.
                    Therefore for all $n\in\naturalsPlus$ we have $\fst{u_1(n)} = n$
                        by \cref{space_filling_dyadic_state_orbit_stage_index_local}.
                    Hence for all $n\in\naturalsPlus$ we have $\fst{u_2(n)} = n$
                        by \cref{space_filling_dyadic_state_orbit_stage_index_local}.
                    Let $A_0 = \{n\in\naturalsPlus \mid \text{if $n \leq N$ then $\snd{u_1(n)} = \snd{u_2(n)}$}\}$.
                    $A_0\subseteq\naturalsPlus$ by \cref{subseteq}.
                    Show that if $1 \leq N$ then $\snd{u_1(1)} = \snd{u_2(1)}$.
                    \begin{subproof}
                        Follows by \cref{snd_eq,pair_eq_iff,subseteq}.
                    \end{subproof}
                    Therefore $1\in A_0$ by \cref{real_one_in_naturals,subseteq}.
                    Show that for all $n\in A_0$ we have $n + 1\in A_0$.
                    \begin{subproof}
                        Fix $n\in A_0$.
                        Hence $n\in\naturalsPlus$ by \cref{subseteq}.
                        Show that if $n + 1 \leq N$ then $\snd{u_1(n + 1)} = \snd{u_2(n + 1)}$.
                        \begin{subproof}
                            Assume $n + 1 \leq N$.
                            Therefore $n < n + 1$ by \cref{real_naturals_lt_succ_local}.
                            Hence $n \leq n + 1$ by \cref{real_lt_imp_leq}.
                            Therefore $n\in\reals$ and $N\in\reals$
                                by \cref{real_naturals_subset_reals,subseteq}.
                            Hence $1\in\reals$ by \cref{real_one_in_naturals,real_naturals_subset_reals,subseteq}.
                            Therefore $n + 1\in\reals$ by \cref{real_add_closed}.
                            Hence $n \leq N$ by \cref{real_leq_trans}.
                            Therefore $\snd{u_1(n)} = \snd{u_2(n)}$ by \cref{subseteq}.
                            Hence $n\in\initialSegment{N}$ and $n + 1\in\initialSegment{N}$
                                by \cref{initial_segment_naturalsplus,real_naturals_closed_succ}.
                            Therefore $z(n) = x(n)$ by \cref{subseteq}.
                            Hence $x(n)\in E(n)$ and $z(n)\in E(n)$ by \cref{indexed_product,subseteq}.
                            Therefore $E(n) = D$ by \cref{function_apply_intro,subseteq}.
                            Hence $x(n)\in D$ and $z(n)\in D$ by \cref{subseteq}.

                            Therefore $u_1(n)\in G_1$ and $u_2(n)\in G_1$ and
                                $u_1(n + 1)\in G_1$ and $u_2(n + 1)\in G_1$
                                by \cref{funs_type_apply,real_naturals_closed_succ,sequence_in}.
                            Hence $u_1(n)\mathrel{R_1}\apply{g_1}{u_1(n)}$ and
                                $u_2(n)\mathrel{R_3}\apply{g_2}{u_2(n)}$ by \cref{subseteq}.
                            Therefore $u_1(n)\mathrel{R_1}u_1(n + 1)$ and
                                $u_2(n)\mathrel{R_3}u_2(n + 1)$ by \cref{subseteq}.
                            Hence $(u_1(n),u_1(n + 1))\in R_1$ and
                                $(u_2(n),u_2(n + 1))\in R_3$ by \cref{relation,subseteq,times_elem_is_tuple}.

                            Take $k_1,J_1,H_1,a_1,b_1,d_1$ such that
                                $(u_1(n),u_1(n + 1))\in G_1\times G_1$ and
                                $k_1\in\naturalsPlus$ and
                                $\fst{(u_1(n),u_1(n + 1))} = (k_1,J_1)$ and
                                $\snd{(u_1(n),u_1(n + 1))} = (k_1 + 1,H_1)$ and
                                $d_1 = x(k_1)$ and
                                $a_1\in\reals$ and
                                $b_1\in\reals$ and
                                $a_1 \leq b_1$ and
                                $J_1 = \intervalclosed{a_1}{b_1}$ and
                                $H_1\subseteq J_1$ and
                                $(d_1,H_1)\in\{(0,\intervalclosed{a_1}{(a_1 + b_1) / (1 + 1)}),
                                    (1,\intervalclosed{(a_1 + b_1) / (1 + 1)}{b_1})\}$
                                by \cref{subseteq}.
                            Take $k_3,J_3,H_3,a_3,b_3,d_3$ such that
                                $(u_2(n),u_2(n + 1))\in G_1\times G_1$ and
                                $k_3\in\naturalsPlus$ and
                                $\fst{(u_2(n),u_2(n + 1))} = (k_3,J_3)$ and
                                $\snd{(u_2(n),u_2(n + 1))} = (k_3 + 1,H_3)$ and
                                $d_3 = z(k_3)$ and
                                $a_3\in\reals$ and
                                $b_3\in\reals$ and
                                $a_3 \leq b_3$ and
                                $J_3 = \intervalclosed{a_3}{b_3}$ and
                                $H_3\subseteq J_3$ and
                                $(d_3,H_3)\in\{(0,\intervalclosed{a_3}{(a_3 + b_3) / (1 + 1)}),
                                    (1,\intervalclosed{(a_3 + b_3) / (1 + 1)}{b_3})\}$
                                by \cref{subseteq}.

                            Therefore $\fst{u_1(n)} = k_1$ and $\fst{u_2(n)} = k_3$ and
                                $\snd{u_1(n)} = J_1$ and $\snd{u_2(n)} = J_3$ and
                                $\snd{u_1(n + 1)} = H_1$ and $\snd{u_2(n + 1)} = H_3$
                                by \cref{fst_eq,snd_eq,pair_eq_iff,subseteq}.
                            Hence $\fst{u_1(n)} = n$ and $\fst{u_2(n)} = n$ by \cref{subseteq}.
                            Therefore $k_1 = n$ and $k_3 = n$ by \cref{subseteq}.
                            Hence $d_1 = x(n)$ and $d_3 = z(n)$ by \cref{subseteq}.
                            Therefore $d_3 = d_1$ by \cref{subseteq}.
                            Hence $J_1 = J_3$ by \cref{subseteq}.
                            Therefore $a_1 = a_3$ and $b_1 = b_3$
                                by \cref{space_filling_intervalclosed_endpoints_unique,subseteq}.
                            Hence $d_1\in D$ by \cref{subseteq}.
                            Therefore $(d_3,H_1)\in\{(0,\intervalclosed{a_3}{(a_3 + b_3) / (1 + 1)}),
                                (1,\intervalclosed{(a_3 + b_3) / (1 + 1)}{b_3})\}$
                                by \cref{subseteq,real_add_eq_subst}.
                            Hence $H_1 = H_3$ by \cref{space_filling_dyadic_child_choice_unique,subseteq,real_add_eq_subst}.
                            Therefore $\snd{u_1(n + 1)} = \snd{u_2(n + 1)}$ by \cref{subseteq}.
                        \end{subproof}
                        Therefore $n + 1\in A_0$ by \cref{real_naturals_closed_succ,subseteq}.
                    \end{subproof}
                    Hence for all $n\in\naturalsPlus$ we have $n\in A_0$ by \cref{real_naturals_induction}.
                    Therefore $N\in A_0$ by \cref{subseteq}.
                    Hence $\snd{u_1(N)} = \snd{u_2(N)}$ by \cref{subseteq}.
                    Therefore $\snd{u_2(N)} = J_0$ by \cref{snd_eq,pair_eq_iff,subseteq}.
                    Hence $u\in\intervalclosed{p_0}{q_0}$ by \cref{subseteq}.
                    Therefore $y\in\reals$ and $u\in\reals$ by \cref{intervalclosed,subseteq}.
                    $u < y$ or $u = y$ or $y < u$ by \cref{real_lt_trichotomy}.
                    \begin{byCase}
                    \caseOf{$u < y$.}
                        We have $u \leq y$ by \cref{real_lt_imp_leq}.
                        Thus $y - u \leq q_0 - p_0$ by \cref{space_filling_interval_width_bounds_difference}.
                        Hence $y - u\in\reals$ by \cref{real_sub,real_add_inverse,real_add_closed}.
                        Therefore $q_0 - p_0\in\reals$ and $1 / N\in\reals$
                            by \cref{real_sub,real_add_inverse,real_add_closed,naturalsplus_reciprocal_real}.
                        Thus $y - u \leq 1 / N$ by \cref{real_leq_trans}.
                        Hence $y - u < \epsilon$ by \cref{real_leq_lt_trans}.
                        Therefore $d((u,y)) < \epsilon$ by \cref{space_filling_standard_metric_lt_from_ordered_difference_local}.
                        Thus $u\in\metricBall{d}{\reals}{y}{\epsilon}$
                            by \cref{metric_symmetric,metric_on,real_lt_subst_left,metric_ball}.
                        Therefore $u\in V$ by \cref{subseteq}.
                    \caseOf{$u = y$.}
                        We have $d((y,u)) = 0$ by \cref{standard_metric_is_metric,metric_on,metric_zero_characterization}.
                        Thus $u\in\metricBall{d}{\reals}{y}{\epsilon}$ by \cref{real_lt_subst_left,metric_ball}.
                        Therefore $u\in V$ by \cref{subseteq}.
                    \caseOf{$y < u$.}
                        We have $y \leq u$ by \cref{real_lt_imp_leq}.
                        Therefore $u - y \leq q_0 - p_0$ by \cref{space_filling_interval_width_bounds_difference}.
                        Thus $u - y\in\reals$ by \cref{real_sub,real_add_inverse,real_add_closed}.
                        Hence $q_0 - p_0\in\reals$ and $1 / N\in\reals$
                            by \cref{real_sub,real_add_inverse,real_add_closed,naturalsplus_reciprocal_real}.
                        Therefore $u - y \leq 1 / N$ by \cref{real_leq_trans}.
                        Thus $u - y < \epsilon$ by \cref{real_leq_lt_trans}.
                        Hence $d((y,u)) < \epsilon$ by \cref{space_filling_standard_metric_lt_from_ordered_difference_local}.
                        Therefore $u\in\metricBall{d}{\reals}{y}{\epsilon}$ by \cref{metric_ball}.
                        Hence $u\in V$ by \cref{subseteq}.
                    \end{byCase}
            \end{subproof}
            Therefore there exists $W_0$ such that
                $W_0$ is open in $\productFamily{E}$ with respect to $P_B$ and
                $x\in W_0$ and
                $\img{r}{W_0}\subseteq V$ by \cref{subseteq}.
        \end{subproof}
        Therefore $r$ is continuous at $x$ from $\productFamily{E}$ to $\reals$ with respect to $P_B$ and $\standardTopologyR$
            by \cref{continuous_at_open,standard_basis_is_basis,basis_topology_is_topology,standard_topology_reals}.
    \end{subproof}
    Hence $r$ is continuous from $\productFamily{E}$ to $\reals$ with respect to $P_B$ and $\standardTopologyR$
        by \cref{continuous_at_implies_continuous,standard_basis_is_basis,basis_topology_is_topology,standard_topology_reals}.
    We have $\img{r}{\productFamily{E}}\subseteq I$ by \cref{img_subseteq_ran,funs_ran,subseteq_transitive}.
    Therefore $r$ is continuous from $\productFamily{E}$ to $I$ with respect to $P_B$ and $T_I$
        by \cref{continuous_subspace_range}.
\end{proof}

% from §44 helper lemma 119: every positive natural is 1, even, or odd
\begin{lemma}\label{space_filling_naturalsplus_parity_decomposition}
    Suppose $n\in\naturalsPlus$.
    Then $n = 1$ or there exists $k\in\naturalsPlus$ such that
        $n = k + k$ or
        $n = (k + k) + 1$.
\end{lemma}
\begin{proof}
    Let $A = \{m\in\naturalsPlus \mid \text{$m = 1$ or there exists $k\in\naturalsPlus$ such that
        $m = k + k$ or
        $m = (k + k) + 1$}\}$.
    $A\subseteq\naturalsPlus$ by \cref{subseteq}.
    $1\in A$ by \cref{real_one_in_naturals,subseteq}.
    Show that for all $m\in A$ we have $m + 1\in A$.
    \begin{subproof}
        Fix $m\in A$.
        $m = 1$ or there exists $k\in\naturalsPlus$ such that
            $m = k + k$ or
            $m = (k + k) + 1$ by \cref{subseteq}.
        \begin{byCase}
        \caseOf{$m = 1$.}
            $m + 1\in A$
                by \cref{real_one_in_naturals,subseteq,real_add_eq_subst,real_naturals_closed_succ}.
        \caseOf{there exists $k\in\naturalsPlus$ such that $m = k + k$ or $m = (k + k) + 1$.}
            Take $k\in\naturalsPlus$ such that $m = k + k$ or $m = (k + k) + 1$.
            \begin{byCase}
            \caseOf{$m = k + k$.}
                $m + 1\in A$
                    by \cref{naturalsplus_add_closed,real_naturals_closed_succ,real_add_eq_subst,subseteq}.
            \caseOf{$m = (k + k) + 1$.}
                We have $(k + 1) + (k + 1) = ((k + k) + 1) + 1$
                    by \cref{real_naturals_subset_reals,real_one_in_naturals,subseteq,real_add_closed,real_add_four_exchange,real_add_assoc,real_add_eq_subst}.
                $m + 1\in A$ by \cref{subseteq,real_add_eq_subst,real_naturals_closed_succ}.
            \end{byCase}
        \end{byCase}
    \end{subproof}
    For all $m\in\naturalsPlus$ we have $m\in A$ by \cref{real_naturals_induction}.
    $n = 1$ or there exists $k\in\naturalsPlus$ such that
        $n = k + k$ or
        $n = (k + k) + 1$ by \cref{subseteq}.
\end{proof}

% from §44 helper lemma 120: strict natural inequality yields a successor bound
\begin{lemma}\label{space_filling_naturalsplus_lt_imp_succ_leq}
    Suppose $m\in\naturalsPlus$ and $n\in\naturalsPlus$.
    Suppose $m < n$.
    Then $m + 1 \leq n$.
\end{lemma}
\begin{proof}
    We have $n = 1$ or there exists $k\in\naturalsPlus$ such that $n = k + 1$
        by \cref{real_naturals_predecessor}.
    \begin{byCase}
    \caseOf{$n = 1$.}
        Suppose not.
        $1 \leq m$ by \cref{real_one_leq_naturalsplus}.
        We have $1\in\reals$ and $m\in\reals$ by \cref{real_one_in_naturals,real_naturals_subset_reals,subseteq}.
        $1 < 1$ by \cref{subseteq,real_leq_lt_trans}.
        Contradiction by \cref{real_lt_irreflexive}.
    \caseOf{there exists $k\in\naturalsPlus$ such that $n = k + 1$.}
        Take $k\in\naturalsPlus$ such that $n = k + 1$.
        $m \leq k$ by \cref{subseteq,real_lt_imp_leq,real_lt_irreflexive,real_lt_subst_right,real_naturals_subset_reals,real_naturals_leq_succ_elim}.
        $m + 1 \leq k + 1$
            by \cref{real_naturals_subset_reals,real_one_in_naturals,subseteq,real_leq_add}.
        $m + 1 \leq n$ by \cref{subseteq}.
    \end{byCase}
\end{proof}

% from §44 helper lemma 121: equal doubles have equal halves
\begin{lemma}\label{space_filling_double_eq_halves_eq}
    Suppose $a\in\naturalsPlus$ and $b\in\naturalsPlus$.
    Suppose $a + a = b + b$.
    Then $a = b$.
\end{lemma}
\begin{proof}
    Suppose not.
    $a \neq b$.
    $a\in\reals$ and $b\in\reals$ by \cref{real_naturals_subset_reals,subseteq}.
    $a < b$ or $b < a$ by \cref{real_lt_trichotomy}.
    \begin{byCase}
    \caseOf{$a < b$.}
        $a + a < a + b$ by \cref{real_lt_add,real_add_comm,real_lt_subst_right}.
        $a + b < b + b$ by \cref{real_lt_add}.
        We have $a + a\in\reals$ and $a + b\in\reals$ and $b + b\in\reals$ by \cref{real_add_closed}.
        $a + a < b + b$ by \cref{real_lt_trans}.
        $a + a < a + a$ by \cref{subseteq,real_lt_subst_right}.
        Contradiction by \cref{real_lt_irreflexive}.
    \caseOf{$b < a$.}
        $b + b < b + a$ by \cref{real_lt_add,real_add_comm,real_lt_subst_right}.
        $b + a < a + a$ by \cref{real_lt_add}.
        We have $b + b\in\reals$ and $b + a\in\reals$ and $a + a\in\reals$ by \cref{real_add_closed}.
        $b + b < a + a$ by \cref{real_lt_trans}.
        $b + b < b + b$ by \cref{subseteq,real_lt_subst_right}.
        Contradiction by \cref{real_lt_irreflexive}.
    \end{byCase}
\end{proof}

% from §44 helper lemma 122: a double is never the successor of a double
\begin{lemma}\label{space_filling_double_neq_double_succ}
    Suppose $a\in\naturalsPlus$ and $b\in\naturalsPlus$.
    Then $a + a \neq (b + b) + 1$.
\end{lemma}
\begin{proof}
    Suppose not.
    $a + a = (b + b) + 1$.
    $a\in\reals$ and $b\in\reals$ by \cref{real_naturals_subset_reals,subseteq}.
    $a = b$ or $a < b$ or $b < a$ by \cref{real_lt_trichotomy}.
    \begin{byCase}
    \caseOf{$a = b$.}
        $a + a = (a + a) + 1$ by \cref{real_add_eq_subst}.
        $a + a\in\naturalsPlus$ by \cref{naturalsplus_add_closed}.
        We have $a + a\in\reals$ by \cref{real_naturals_subset_reals,subseteq}.
        $a + a < (a + a) + 1$ by \cref{real_naturals_lt_succ_local}.
        $a + a < a + a$ by \cref{subseteq,real_lt_subst_right}.
        Contradiction by \cref{real_lt_irreflexive}.
    \caseOf{$a < b$.}
        $a + a < a + b$ by \cref{real_lt_add,real_add_comm,real_lt_subst_right}.
        $a + b < b + b$ by \cref{real_lt_add}.
        We have $a + a\in\reals$ and $a + b\in\reals$ and $b + b\in\reals$ by \cref{real_add_closed}.
        $a + a < b + b$ by \cref{real_lt_trans}.
        $b + b < (b + b) + 1$
            by \cref{naturalsplus_add_closed,real_naturals_subset_reals,real_naturals_closed_succ,subseteq,real_naturals_lt_succ_local}.
        $a + a < (b + b) + 1$ by \cref{real_lt_trans}.
        $a + a < a + a$ by \cref{subseteq,real_lt_subst_right}.
        Contradiction by \cref{real_lt_irreflexive}.
    \caseOf{$b < a$.}
        $b + 1 \leq a$ by \cref{space_filling_naturalsplus_lt_imp_succ_leq}.
        We have $1\in\reals$ by \cref{real_one_in_naturals,real_naturals_subset_reals,subseteq}.
        $(b + 1) + b \leq a + b$
            by \cref{real_one_in_naturals,real_naturals_subset_reals,subseteq,real_add_closed,real_leq_add}.
        $(b + 1) + b = b + 1 + b$ by \cref{real_add_assoc}.
        We have $1 + b = b + 1$ by \cref{real_add_comm}.
        $b + 1 + b = b + (1 + b)$ by \cref{real_add_assoc}.
        $b + (1 + b) = b + (b + 1)$ by \cref{real_add_eq_subst}.
        $b + (b + 1) = (b + b) + 1$ by \cref{real_add_assoc}.
        $(b + b) + 1 \leq a + b$ by \cref{real_leq_subst_left_local}.
        $a + b < a + a$ by \cref{real_lt_add,real_add_comm,real_lt_subst_right}.
        We have $(b + b) + 1\in\reals$ and $a + b\in\reals$ and $a + a\in\reals$ by \cref{real_add_closed}.
        $(b + b) + 1 < a + a$ by \cref{real_leq_lt_trans}.
        $a + a < a + a$ by \cref{subseteq,real_lt_subst_left}.
        Contradiction by \cref{real_lt_irreflexive}.
    \end{byCase}
\end{proof}

% from §44 helper lemma 123: a chosen retained-third code-point assignment is continuous
\begin{lemma}\label{space_filling_retained_third_code_map_continuous}
    Let $I = \intervalclosed{0}{1}$.
    Let $T_I = \subspaceTopology{\standardTopologyR}{I}$.
    Let $D = \{0,1\}$.
    Let $T_D = \subspaceTopology{T_I}{D}$.
    Let $E = \{(n,D) \mid n\in\naturalsPlus\}$.
    Let $U = \{(n,T_D) \mid n\in\naturalsPlus\}$.
    Let $P_B = \productTopologyIndexed{U}{E}$.
    Suppose $r\in\funs{\productFamily{E}}{I}$.
    Suppose for all $x\in\productFamily{E}$ we have $r(x)$ is a retained-third code point of $x$ in $I$ and $E$.
    Then $r$ is continuous from $\productFamily{E}$ to $I$ with respect to $P_B$ and $T_I$.
\end{lemma}
\begin{proof}
    $I\subseteq\reals$ by \cref{intervalclosed,subseteq}.
    $\standardTopologyR$ is a topology on $\reals$
        by \cref{standard_basis_is_basis,basis_topology_is_topology,standard_topology_reals}.
    $I$ is a subspace of $\reals$ with respect to $\standardTopologyR$ by \cref{subspace_of}.
    $T_I$ is a topology on $I$ by \cref{subspace_topology_is_topology}.
    $\productFamily{E}$ is compact with respect to $P_B$ and
        $\productFamily{E}$ is Hausdorff with respect to $P_B$
        by \cref{space_filling_binary_product_compact_hausdorff}.
    $P_B$ is a topology on $\productFamily{E}$ by \cref{hausdorff}.
    We have $E$ is a function on $\naturalsPlus$
        by \cref{relation,rightunique,dom_intro,dom_iff,pair_eq_iff,subseteq,subseteq_antisymmetric}.
    $\dom{E} = \naturalsPlus$ by \cref{funs}.
    For all $n\in\naturalsPlus$ we have $E(n) = D$ by \cref{function_apply_intro}.
    $r\in\funs{\productFamily{E}}{\reals}$ by \cref{funs_weaken_codom,subseteq}.
    $r$ is a function and $\dom{r} = \productFamily{E}$ by \cref{funs_elim,funs_is_function}.

    Show that for all $x\in\productFamily{E}$ we have
        $r$ is continuous at $x$ from $\productFamily{E}$ to $\reals$ with respect to $P_B$ and $\standardTopologyR$.
    \begin{subproof}
        Fix $x\in\productFamily{E}$.
        Let $y = r(x)$.
        $y$ is a retained-third code point of $x$ in $I$ and $E$ by \cref{subseteq}.
        Take $K_1,G_1,R_1,g_1,u_1$ such that
            $K_1 = \{J\in\pow{I} \mid \text{there exist $a,b\in\reals$ such that
                $J = \intervalclosed{a}{b}$ and
                $a \leq b$}\}$ and
            $G_1 = \naturalsPlus\times K_1$ and
            $R_1 = \{p\in G_1\times G_1 \mid \text{there exist $n,J,H,a,b,d$ such that
                $n\in\naturalsPlus$ and
                $\fst{p} = (n,J)$ and
                $\snd{p} = (n + 1,H)$ and
                $d = x(n)$ and
                $a\in\reals$ and
                $b\in\reals$ and
                $a \leq b$ and
                $J = \intervalclosed{a}{b}$ and
                $H\subseteq J$ and
                $(d,H)\in\{(0,\intervalclosed{a}{(a + a + b) / (1 + 1 + 1)}),
                    (1,\intervalclosed{(a + b + b) / (1 + 1 + 1)}{b})\}$}\}$ and
            $g_1$ is a retained-third state selector on $G_1$ and $R_1$ and
            $u_1$ is a retained-third state orbit of $g_1$ on $G_1$ from $I$ and
            for all $n\in\naturalsPlus$ we have $y\in\snd{u_1(n)}$
            by \cref{space_filling_retained_third_code_point}.
        We have $g_1\in\funs{G_1}{G_1}$ and
            for all $w\in G_1$ we have $w\mathrel{R_1}\apply{g_1}{w}$
            by \cref{space_filling_retained_third_state_selector_on}.
        $u_1$ is a sequence in $G_1$ and
            $u_1(1) = (1,I)$ and
            for all $n\in\naturalsPlus$ we have $u_1(n + 1) = \apply{g_1}{u_1(n)}$
            by \cref{space_filling_retained_third_state_orbit_from}.

        Show that for all $V$ such that
            $V$ is open in $\reals$ with respect to $\standardTopologyR$ and
            $y\in V$
        there exists $W$ such that
            $W$ is open in $\productFamily{E}$ with respect to $P_B$ and
            $x\in W$ and
            $\img{r}{W}\subseteq V$.
        \begin{subproof}
            Fix $V$ such that
                $V$ is open in $\reals$ with respect to $\standardTopologyR$ and
                $y\in V$.
            Let $d = \standardMetricR$.
            $\metricTopology{d}{\reals} = \standardTopologyR$ by \cref{standard_metric_topology}.
            $V\subseteq\reals$ by \cref{open_in,topology_on,subseteq,pow_iff,standard_basis_is_basis,basis_topology_is_topology,standard_topology_reals}.
            $d$ is a metric on $\reals$ by \cref{standard_metric_is_metric}.
            Take $\epsilon\in\reals$ such that
                $0 < \epsilon$ and
                $\metricBall{d}{\reals}{y}{\epsilon}\subseteq V$
                by \cref{standard_metric_topology,metric_open_iff,neighborhood}.
            Take $N\in\naturalsPlus$ such that
                $0 < 1 / N$ and
                $1 / N < \epsilon$
                by \cref{real_small_reciprocal}.
            $u_1(N)\in G_1$ by \cref{sequence_in,funs_type_apply}.
            Take $m_0,J_0$ such that
                $u_1(N) = (m_0,J_0)$ and
                $m_0\in\naturalsPlus$ and
                $J_0\in K_1$
                by \cref{subseteq,times_elem_is_tuple}.
            Take $p_0,q_0$ such that
                $p_0\in\reals$ and
                $q_0\in\reals$ and
                $J_0 = \intervalclosed{p_0}{q_0}$ and
                $p_0 \leq q_0$
                by \cref{subseteq}.
            $\fst{u_1(N)} = m_0$ and $\snd{u_1(N)} = J_0$ by \cref{fst_eq,snd_eq,pair_eq_iff}.
            $m_0 = N$ by \cref{space_filling_retained_third_state_orbit_index,subseteq}.
            $\snd{u_1(N)} = \intervalclosed{p_0}{q_0}$ by \cref{subseteq}.
            $q_0 - p_0 \leq 1 / N$ by \cref{space_filling_retained_third_state_orbit_width_bound,subseteq}.
            $y\in\intervalclosed{p_0}{q_0}$ by \cref{subseteq}.
            We have $x$ is a function and $\dom{x} = \dom{E}$ by \cref{indexed_product,funs_elim}.
            Let $a = \restrl{x}{\initialSegment{N}}$.
            $a$ is a function by \cref{restrl_of_function_is_function}.
            $\dom{a} = \dom{x}\inter\initialSegment{N}$ by \cref{dom_of_restrl_eq_inter}.
            $\dom{x} = \naturalsPlus$ by \cref{subseteq}.
            $\initialSegment{N}\subseteq\dom{x}$ by \cref{initial_segment_naturalsplus,subseteq}.
            $\dom{a} = \initialSegment{N}$ by \cref{inter_absorb_supseteq_left}.
            $a\in\funs{\initialSegment{N}}{D}$
                by \cref{subseteq,initial_segment_naturalsplus,indexed_product,function_apply_intro,restrl_of_function_apply,funs_intro}.
            Let $W = \{z\in\productFamily{E} \mid \text{for all $i\in\dom{a}$ we have $z(i) = a(i)$}\}$.
            $W$ is the binary prefix cylinder at $a$ in $E$ by \cref{space_filling_binary_prefix_cylinder}.
            $W$ is open in $\productFamily{E}$ with respect to $P_B$
                by \cref{space_filling_binary_prefix_cylinder_open}.
            $x\in W$ by \cref{subseteq,restrl_of_function_apply}.
            Show that for all $u$ such that $u\in\img{r}{W}$ we have $u\in V$.
            \begin{subproof}
                    Fix $u$ such that $u\in\img{r}{W}$.
                    Take $z\in\dom{r}\inter W$ such that $u = r(z)$ by \cref{img_of_function_elim}.
                    $z\in W$ by \cref{inter_elim_right}.
                    $z\in\productFamily{E}$ and
                        for all $i\in\dom{a}$ we have $z(i) = a(i)$ by \cref{subseteq}.
                    For all $i\in\initialSegment{N}$ we have $z(i) = x(i)$
                        by \cref{subseteq,restrl_of_function_apply}.
                    $u$ is a retained-third code point of $z$ in $I$ and $E$ by \cref{subseteq}.
                    Take $K_2,G_2,R_2,g_2,u_2$ such that
                        $K_2 = \{J\in\pow{I} \mid \text{there exist $a_2,b_2\in\reals$ such that
                            $J = \intervalclosed{a_2}{b_2}$ and
                            $a_2 \leq b_2$}\}$ and
                        $G_2 = \naturalsPlus\times K_2$ and
                        $R_2 = \{p\in G_2\times G_2 \mid \text{there exist $n,J,H,a_2,b_2,d_2$ such that
                            $n\in\naturalsPlus$ and
                            $\fst{p} = (n,J)$ and
                            $\snd{p} = (n + 1,H)$ and
                            $d_2 = z(n)$ and
                            $a_2\in\reals$ and
                            $b_2\in\reals$ and
                            $a_2 \leq b_2$ and
                            $J = \intervalclosed{a_2}{b_2}$ and
                            $H\subseteq J$ and
                            $(d_2,H)\in\{(0,\intervalclosed{a_2}{(a_2 + a_2 + b_2) / (1 + 1 + 1)}),
                                (1,\intervalclosed{(a_2 + b_2 + b_2) / (1 + 1 + 1)}{b_2})\}$}\}$ and
                        $g_2$ is a retained-third state selector on $G_2$ and $R_2$ and
                        $u_2$ is a retained-third state orbit of $g_2$ on $G_2$ from $I$ and
                        for all $n\in\naturalsPlus$ we have $u\in\snd{u_2(n)}$
                        by \cref{space_filling_retained_third_code_point}.
                    We have $g_2\in\funs{G_2}{G_2}$ and
                        for all $w\in G_2$ we have $w\mathrel{R_2}\apply{g_2}{w}$
                        by \cref{space_filling_retained_third_state_selector_on}.
                    $u_2$ is a sequence in $G_2$ and
                        $u_2(1) = (1,I)$ and
                        for all $n\in\naturalsPlus$ we have $u_2(n + 1) = \apply{g_2}{u_2(n)}$
                        by \cref{space_filling_retained_third_state_orbit_from}.
                    $K_2 = K_1$ by \cref{subseteq_antisymmetric,subseteq}.
                    $G_2 = G_1$ by \cref{subseteq}.
                    $g_2\in\funs{G_1}{G_1}$ by \cref{subseteq}.
                    $u_2$ is a sequence in $G_1$ and
                        $u_2(1) = (1,I)$ and
                        for all $n\in\naturalsPlus$ we have $u_2(n + 1) = \apply{g_2}{u_2(n)}$
                        by \cref{subseteq}.
                    Let $B_1 = \{m\in\naturalsPlus \mid \fst{u_1(m)} = m\}$.
                    $B_1\subseteq\naturalsPlus$ by \cref{subseteq}.
                    $1\in B_1$ by \cref{real_one_in_naturals,pair_eq_iff,fst_eq,subseteq}.
                    Show that for all $m\in B_1$ we have $m + 1\in B_1$.
                    \begin{subproof}
                        Follows by \cref{sequence_in,funs_type_apply,relation,subseteq,fst_eq,snd_eq,pair_eq_iff,real_add_eq_subst,real_naturals_closed_succ}.
                    \end{subproof}
                    For all $m\in\naturalsPlus$ we have $m\in B_1$ by \cref{real_naturals_induction}.
                    Let $B_2 = \{m\in\naturalsPlus \mid \fst{u_2(m)} = m\}$.
                    $B_2\subseteq\naturalsPlus$ by \cref{subseteq}.
                    $1\in B_2$ by \cref{real_one_in_naturals,pair_eq_iff,fst_eq,subseteq}.
                    Show that for all $m\in B_2$ we have $m + 1\in B_2$.
                    \begin{subproof}
                        Follows by \cref{sequence_in,funs_type_apply,relation,subseteq,fst_eq,snd_eq,pair_eq_iff,real_add_eq_subst,real_naturals_closed_succ}.
                    \end{subproof}
                    For all $m\in\naturalsPlus$ we have $m\in B_2$ by \cref{real_naturals_induction}.
                    For all $n\in\naturalsPlus$ we have $\fst{u_1(n)} = n$ by \cref{subseteq}.
                    For all $n\in\naturalsPlus$ we have $\fst{u_2(n)} = n$ by \cref{subseteq}.
                    Let $A_0 = \{n\in\naturalsPlus \mid \text{if $n \leq N$ then $\snd{u_1(n)} = \snd{u_2(n)}$}\}$.
                    $A_0\subseteq\naturalsPlus$ by \cref{subseteq}.
                    Show that if $1 \leq N$ then $\snd{u_1(1)} = \snd{u_2(1)}$.
                    \begin{subproof}
                        Follows by \cref{subseteq,snd_eq,pair_eq_iff}.
                    \end{subproof}
                    $1\in A_0$ by \cref{real_one_in_naturals,subseteq}.
                    Show that for all $n\in A_0$ we have $n + 1\in A_0$.
                    \begin{subproof}
                        Fix $n\in A_0$.
                        $n\in\naturalsPlus$ by \cref{subseteq}.
                        Show that if $n + 1 \leq N$ then $\snd{u_1(n + 1)} = \snd{u_2(n + 1)}$.
                        \begin{subproof}
                            Assume $n + 1 \leq N$.
                            $n < n + 1$ by \cref{real_naturals_lt_succ_local}.
                            $n \leq n + 1$ by \cref{real_lt_imp_leq}.
                            $n\in\reals$ and $N\in\reals$
                                by \cref{real_naturals_subset_reals,subseteq}.
                            $1\in\reals$ by \cref{real_one_in_naturals,real_naturals_subset_reals,subseteq}.
                            $n + 1\in\reals$ by \cref{real_add_closed}.
                            $n \leq N$ by \cref{real_leq_trans}.
                            $\snd{u_1(n)} = \snd{u_2(n)}$ by \cref{subseteq}.
                            $n\in\initialSegment{N}$ and $n + 1\in\initialSegment{N}$
                                by \cref{initial_segment_naturalsplus,real_naturals_closed_succ}.
                            $z(n) = x(n)$ by \cref{subseteq}.
                            $x(n)\in E(n)$ and $z(n)\in E(n)$ by \cref{indexed_product,subseteq}.
                            $E(n) = D$ by \cref{function_apply_intro,subseteq}.
                            $x(n)\in D$ and $z(n)\in D$ by \cref{subseteq}.

                            $u_1(n)\in G_1$ and $u_2(n)\in G_1$ and
                                $u_1(n + 1)\in G_1$ and $u_2(n + 1)\in G_1$
                                by \cref{funs_type_apply,real_naturals_closed_succ,sequence_in}.
                            $u_1(n)\mathrel{R_1}\apply{g_1}{u_1(n)}$ and
                                $u_2(n)\mathrel{R_2}\apply{g_2}{u_2(n)}$ by \cref{subseteq}.
                            $u_1(n)\mathrel{R_1}u_1(n + 1)$ and
                                $u_2(n)\mathrel{R_2}u_2(n + 1)$ by \cref{subseteq}.
                            $(u_1(n),u_1(n + 1))\in R_1$ and
                                $(u_2(n),u_2(n + 1))\in R_2$ by \cref{relation}.

                            Take $k_1,J_1,H_1,a_1,b_1,d_1$ such that
                                $(u_1(n),u_1(n + 1))\in G_1\times G_1$ and
                                $k_1\in\naturalsPlus$ and
                                $\fst{(u_1(n),u_1(n + 1))} = (k_1,J_1)$ and
                                $\snd{(u_1(n),u_1(n + 1))} = (k_1 + 1,H_1)$ and
                                $d_1 = x(k_1)$ and
                                $a_1\in\reals$ and
                                $b_1\in\reals$ and
                                $a_1 \leq b_1$ and
                                $J_1 = \intervalclosed{a_1}{b_1}$ and
                                $H_1\subseteq J_1$ and
                                $(d_1,H_1)\in\{(0,\intervalclosed{a_1}{(a_1 + a_1 + b_1) / (1 + 1 + 1)}),
                                    (1,\intervalclosed{(a_1 + b_1 + b_1) / (1 + 1 + 1)}{b_1})\}$
                                by \cref{subseteq}.
                            Take $k_3,J_3,H_3,a_3,b_3,d_3$ such that
                                $(u_2(n),u_2(n + 1))\in G_1\times G_1$ and
                                $k_3\in\naturalsPlus$ and
                                $\fst{(u_2(n),u_2(n + 1))} = (k_3,J_3)$ and
                                $\snd{(u_2(n),u_2(n + 1))} = (k_3 + 1,H_3)$ and
                                $d_3 = z(k_3)$ and
                                $a_3\in\reals$ and
                                $b_3\in\reals$ and
                                $a_3 \leq b_3$ and
                                $J_3 = \intervalclosed{a_3}{b_3}$ and
                                $H_3\subseteq J_3$ and
                                $(d_3,H_3)\in\{(0,\intervalclosed{a_3}{(a_3 + a_3 + b_3) / (1 + 1 + 1)}),
                                    (1,\intervalclosed{(a_3 + b_3 + b_3) / (1 + 1 + 1)}{b_3})\}$
                                by \cref{subseteq}.

                            We have $\fst{u_1(n)} = k_1$ and $\fst{u_2(n)} = k_3$ and
                                $\snd{u_1(n)} = J_1$ and $\snd{u_2(n)} = J_3$ and
                                $\snd{u_1(n + 1)} = H_1$ and $\snd{u_2(n + 1)} = H_3$
                                by \cref{fst_eq,snd_eq,pair_eq_iff,subseteq}.
                            $\fst{u_1(n)} = n$ and $\fst{u_2(n)} = n$ by \cref{subseteq}.
                            $k_1 = n$ and $k_3 = n$ by \cref{subseteq}.
                            $d_1 = x(n)$ and $d_3 = z(n)$ by \cref{subseteq}.
                            $d_3 = d_1$ by \cref{subseteq}.
                            $J_1 = J_3$ by \cref{subseteq}.
                            $a_1 = a_3$ and $b_1 = b_3$
                                by \cref{space_filling_intervalclosed_endpoints_unique,subseteq}.
                            $d_1\in D$ by \cref{subseteq}.
                            $(d_3,H_1)\in\{(0,\intervalclosed{a_3}{(a_3 + a_3 + b_3) / (1 + 1 + 1)}),
                                (1,\intervalclosed{(a_3 + b_3 + b_3) / (1 + 1 + 1)}{b_3})\}$
                                by \cref{subseteq,real_add_eq_subst}.
                            $H_1 = H_3$ by \cref{space_filling_retained_third_child_choice_unique,subseteq,real_add_eq_subst}.
                            $\snd{u_1(n + 1)} = \snd{u_2(n + 1)}$ by \cref{subseteq}.
                        \end{subproof}
                        $n + 1\in A_0$ by \cref{real_naturals_closed_succ,subseteq}.
                    \end{subproof}
                    For all $n\in\naturalsPlus$ we have $n\in A_0$ by \cref{real_naturals_induction}.
                    $N\in A_0$ by \cref{subseteq}.
                    $\snd{u_1(N)} = \snd{u_2(N)}$ by \cref{subseteq}.
                    $\snd{u_2(N)} = J_0$ by \cref{snd_eq,pair_eq_iff,subseteq}.
                    $u\in\intervalclosed{p_0}{q_0}$ by \cref{subseteq}.
                    $y\in\reals$ and $u\in\reals$ by \cref{intervalclosed,subseteq}.
                    $u < y$ or $u = y$ or $y < u$ by \cref{real_lt_trichotomy}.
                    \begin{byCase}
                    \caseOf{$u < y$.}
                        $u \leq y$ by \cref{real_lt_imp_leq}.
                        $y - u \leq q_0 - p_0$ by \cref{space_filling_interval_width_bounds_difference}.
                        $y - u\in\reals$ by \cref{real_sub,real_add_inverse,real_add_closed}.
                        $q_0 - p_0\in\reals$ and $1 / N\in\reals$
                            by \cref{real_sub,real_add_inverse,real_add_closed,naturalsplus_reciprocal_real}.
                        $y - u \leq 1 / N$ by \cref{real_leq_trans}.
                        $y - u < \epsilon$ by \cref{real_leq_lt_trans}.
                        $d((u,y)) < \epsilon$ by \cref{space_filling_standard_metric_lt_from_ordered_difference_local}.
                        $u\in\metricBall{d}{\reals}{y}{\epsilon}$
                            by \cref{metric_symmetric,metric_on,real_lt_subst_left,metric_ball}.
                        $u\in V$ by \cref{subseteq}.
                    \caseOf{$u = y$.}
                        $d((y,u)) = 0$ by \cref{standard_metric_is_metric,metric_on,metric_zero_characterization}.
                        $u\in\metricBall{d}{\reals}{y}{\epsilon}$ by \cref{real_lt_subst_left,metric_ball}.
                        $u\in V$ by \cref{subseteq}.
                    \caseOf{$y < u$.}
                        $y \leq u$ by \cref{real_lt_imp_leq}.
                        $u - y \leq q_0 - p_0$ by \cref{space_filling_interval_width_bounds_difference}.
                        $u - y\in\reals$ by \cref{real_sub,real_add_inverse,real_add_closed}.
                        $q_0 - p_0\in\reals$ and $1 / N\in\reals$
                            by \cref{real_sub,real_add_inverse,real_add_closed,naturalsplus_reciprocal_real}.
                        $u - y \leq 1 / N$ by \cref{real_leq_trans}.
                        $u - y < \epsilon$ by \cref{real_leq_lt_trans}.
                        $d((y,u)) < \epsilon$ by \cref{space_filling_standard_metric_lt_from_ordered_difference_local}.
                        $u\in\metricBall{d}{\reals}{y}{\epsilon}$ by \cref{metric_ball}.
                        $u\in V$ by \cref{subseteq}.
                    \end{byCase}
            \end{subproof}
            There exists $W_0$ such that
                $W_0$ is open in $\productFamily{E}$ with respect to $P_B$ and
                $x\in W_0$ and
                $\img{r}{W_0}\subseteq V$ by \cref{subseteq}.
        \end{subproof}
        $r$ is continuous at $x$ from $\productFamily{E}$ to $\reals$ with respect to $P_B$ and $\standardTopologyR$
            by \cref{continuous_at_open,standard_basis_is_basis,basis_topology_is_topology,standard_topology_reals}.
    \end{subproof}
    $r$ is continuous from $\productFamily{E}$ to $\reals$ with respect to $P_B$ and $\standardTopologyR$
        by \cref{continuous_at_implies_continuous,standard_basis_is_basis,basis_topology_is_topology,standard_topology_reals}.
    We have $\img{r}{\productFamily{E}}\subseteq I$ by \cref{img_subseteq_ran,funs_ran,subseteq_transitive}.
    Therefore $r$ is continuous from $\productFamily{E}$ to $I$ with respect to $P_B$ and $T_I$
        by \cref{continuous_subspace_range}.
\end{proof}

% from §44 helper definition 6: even binary extractor on the binary product
\begin{definition}\label{space_filling_even_binary_extractor_on}
    $e$ is an even binary extractor on $E$ iff
        $e\in\funs{\productFamily{E}}{\productFamily{E}}$ and
        for all $w\in\productFamily{E}\times\naturalsPlus$ we have
            $\apply{\apply{e}{\fst{w}}}{\snd{w}} = \apply{\fst{w}}{(\snd{w} + \snd{w})}$.
\end{definition}

% from §44 helper lemma 124: the even-position extractor on binary sequences is continuous
\begin{lemma}\label{space_filling_even_binary_extractor}
    Let $I = \intervalclosed{0}{1}$.
    Let $T_I = \subspaceTopology{\standardTopologyR}{I}$.
    Let $D = \{0,1\}$.
    Let $T_D = \subspaceTopology{T_I}{D}$.
    Let $E = \{(n,D) \mid n\in\naturalsPlus\}$.
    Let $U = \{(n,T_D) \mid n\in\naturalsPlus\}$.
    Let $P_B = \productTopologyIndexed{U}{E}$.
    Then there exists $e$ such that
        $e$ is an even binary extractor on $E$ and
        $e$ is continuous from $\productFamily{E}$ to $\productFamily{E}$ with respect to $P_B$ and $P_B$.
\end{lemma}
\begin{proof}
    We have $E$ is a function on $\naturalsPlus$ and $U$ is a function on $\naturalsPlus$
        by \cref{relation,rightunique,dom_intro,dom_iff,pair_eq_iff,subseteq,subseteq_antisymmetric}.
    $\dom{E} = \naturalsPlus$ and $\dom{U} = \naturalsPlus$ by \cref{funs}.
    For all $n\in\naturalsPlus$ we have $E(n) = D$ and $U(n) = T_D$ by \cref{function_apply_intro}.
    $T_I$ is a topology on $I$
        by \cref{intervalclosed,subseteq,standard_basis_is_basis,basis_topology_is_topology,standard_topology_reals,subspace_of,subspace_topology_is_topology}.
    $0\in\reals$ and $1\in\reals$ and $0 < 1$
        by \cref{real_add_ident,real_one_in_naturals,real_naturals_subset_reals,subseteq,real_zero_lt_one}.
    $0\in I$ and $1\in I$ by \cref{intervalclosed,real_lt_imp_leq}.
    $D\subseteq I$ by \cref{upair_iff,subseteq}.
    $D$ is a subspace of $I$ with respect to $T_I$ by \cref{subspace_of}.
    $T_D$ is a topology on $D$ by \cref{subseteq,subspace_topology_is_topology}.
    $\productFamily{E}$ is compact with respect to $P_B$ and
        $\productFamily{E}$ is Hausdorff with respect to $P_B$
        by \cref{space_filling_binary_product_compact_hausdorff}.
    $P_B$ is a topology on $\productFamily{E}$ by \cref{hausdorff}.

    Let $Q = \{w\in\productFamily{E}\times\productFamily{E} \mid
        \text{for all $n\in\naturalsPlus$ we have $\apply{\snd{w}}{n} = \apply{\fst{w}}{(n + n)}$}\}$.
    $Q$ is a relation by \cref{relation,subseteq,times_elem_is_tuple}.
    Show that for all $x\in\productFamily{E}$ there exists $y$ such that $x\mathrel{Q} y$.
    \begin{subproof}
        Fix $x\in\productFamily{E}$.
        Let $R = \{p\in\naturalsPlus\times D \mid \snd{p} = x(\fst{p} + \fst{p})\}$.
        $R$ is a relation by \cref{relation,subseteq,times_elem_is_tuple}.
        For all $n\in\naturalsPlus$ we have $x(n + n)\in D$
            by \cref{naturalsplus_add_closed,indexed_product,function_apply_intro,subseteq}.
        For all $n\in\naturalsPlus$ there exists $d$ such that $n\mathrel{R} d$
            by \cref{relation,subseteq,times_tuple_intro,fst_eq,snd_eq,pair_eq_iff}.
        Take $y$ such that
            $y$ is a function on $\naturalsPlus$ and
            for all $n\in\naturalsPlus$ we have $n\mathrel{R}\apply{y}{n}$
            by \cref{choice_function}.
        $y\in\funs{\naturalsPlus}{D}$
            by \cref{choice_function,relation,subseteq,fst_eq,snd_eq,pair_eq_iff,indexed_product,naturalsplus_add_closed,function_apply_intro,funs_intro}.
        $y\in\funs{\dom{E}}{\unions{\ran{E}}}$
            by \cref{function_apply_intro,funs_type_apply,subseteq,function_apply_elim,ran_intro,unions_iff,upair_intro_left,funs_intro}.
        $y\in\productFamily{E}$ and $x\mathrel{Q} y$
            by \cref{indexed_product,subseteq,times_tuple_intro,choice_function,fst_eq,snd_eq,pair_eq_iff,relation}.
    \end{subproof}
    Take $e$ such that
        $e$ is a function on $\productFamily{E}$ and
        for all $x\in\productFamily{E}$ we have $x\mathrel{Q}\apply{e}{x}$
        by \cref{choice_function}.
    $e\in\funs{\productFamily{E}}{\productFamily{E}}$
        by \cref{choice_function,funs_intro,relation,subseteq,times_tuple_elim}.
    $e$ is an even binary extractor on $E$
        by \cref{space_filling_even_binary_extractor_on,times_elem_is_tuple,fst_eq,snd_eq,pair_eq_iff,choice_function,relation,subseteq}.

    Show that for all $n\in\dom{E}$ we have $\piIndex{E}{n}\circ e$ is continuous from $\productFamily{E}$ to $\apply{E}{n}$ with respect to $P_B$ and $\apply{U}{n}$.
    \begin{subproof}
        Fix $n\in\dom{E}$.
        Let $g = \piIndex{E}{n}\circ e$.
        $g\in\funs{\productFamily{E}}{E(n)}$ by \cref{projection_map_indexed_function,funs_circ}.
        $n + n\in\naturalsPlus$ by \cref{naturalsplus_add_closed,subseteq}.
        $\piIndex{E}{(n + n)}\in\funs{\productFamily{E}}{E(n + n)}$ by \cref{projection_map_indexed_function,subseteq}.
        For all $x\in\productFamily{E}$ we have $g(x) = \apply{\piIndex{E}{(n + n)}}{x}$
            by \cref{funs_elim,funs_is_function,funs_type_apply,projection_map_indexed_function,funs_ran,circ_apply,projection_map_indexed,function_apply_intro,choice_function,relation,subseteq,fst_eq,snd_eq,pair_eq_iff}.
        $g = \piIndex{E}{(n + n)}$ by \cref{functions_eq_iff}.
        $\piIndex{E}{(n + n)}$ is continuous from $\productFamily{E}$ to $E(n + n)$
            with respect to $P_B$ and $U(n + n)$ by \cref{projection_map_indexed_continuous,subseteq}.
        $g$ is continuous from $\productFamily{E}$ to $E(n)$ with respect to $P_B$ and $U(n)$
            by \cref{function_apply_intro,subseteq}.
    \end{subproof}
    $e$ is continuous from $\productFamily{E}$ to $\productFamily{E}$ with respect to $P_B$ and $P_B$
        by \cref{continuous_map_into_indexed_product,real_one_in_naturals,subseteq}.
\end{proof}

% from §44 helper definition 7: successor-of-even binary extractor on the binary product
\begin{definition}\label{space_filling_odd_binary_extractor_on}
    $o$ is a successor-of-even binary extractor on $E$ iff
        $o\in\funs{\productFamily{E}}{\productFamily{E}}$ and
        for all $w\in\productFamily{E}\times\naturalsPlus$ we have
            $\apply{\apply{o}{\fst{w}}}{\snd{w}} = \apply{\fst{w}}{((\snd{w} + \snd{w}) + 1)}$.
\end{definition}

% from §44 helper lemma 125: the successor-of-even-position extractor on binary sequences is continuous
\begin{lemma}\label{space_filling_odd_binary_extractor}
    Let $I = \intervalclosed{0}{1}$.
    Let $T_I = \subspaceTopology{\standardTopologyR}{I}$.
    Let $D = \{0,1\}$.
    Let $T_D = \subspaceTopology{T_I}{D}$.
    Let $E = \{(n,D) \mid n\in\naturalsPlus\}$.
    Let $U = \{(n,T_D) \mid n\in\naturalsPlus\}$.
    Let $P_B = \productTopologyIndexed{U}{E}$.
    Then there exists $o$ such that
        $o$ is a successor-of-even binary extractor on $E$ and
        $o$ is continuous from $\productFamily{E}$ to $\productFamily{E}$ with respect to $P_B$ and $P_B$.
\end{lemma}
\begin{proof}
    We have $E$ is a function on $\naturalsPlus$ and $U$ is a function on $\naturalsPlus$
        by \cref{relation,rightunique,dom_intro,dom_iff,pair_eq_iff,subseteq,subseteq_antisymmetric}.
    $\dom{E} = \naturalsPlus$ and $\dom{U} = \naturalsPlus$ by \cref{funs}.
    For all $n\in\naturalsPlus$ we have $E(n) = D$ and $U(n) = T_D$ by \cref{function_apply_intro}.
    $T_I$ is a topology on $I$
        by \cref{intervalclosed,subseteq,standard_basis_is_basis,basis_topology_is_topology,standard_topology_reals,subspace_of,subspace_topology_is_topology}.
    $0\in\reals$ and $1\in\reals$ and $0 < 1$
        by \cref{real_add_ident,real_one_in_naturals,real_naturals_subset_reals,subseteq,real_zero_lt_one}.
    $0\in I$ and $1\in I$ by \cref{intervalclosed,real_lt_imp_leq}.
    $D\subseteq I$ by \cref{upair_iff,subseteq}.
    $D$ is a subspace of $I$ with respect to $T_I$ by \cref{subspace_of}.
    $T_D$ is a topology on $D$ by \cref{subseteq,subspace_topology_is_topology}.
    $\productFamily{E}$ is compact with respect to $P_B$ and
        $\productFamily{E}$ is Hausdorff with respect to $P_B$
        by \cref{space_filling_binary_product_compact_hausdorff}.
    $P_B$ is a topology on $\productFamily{E}$ by \cref{hausdorff}.

    Let $Q = \{w\in\productFamily{E}\times\productFamily{E} \mid
        \text{for all $n\in\naturalsPlus$ we have $\apply{\snd{w}}{n} = \apply{\fst{w}}{((n + n) + 1)}$}\}$.
    $Q$ is a relation by \cref{relation,subseteq,times_elem_is_tuple}.
    Show that for all $x\in\productFamily{E}$ there exists $y$ such that $x\mathrel{Q} y$.
    \begin{subproof}
        Fix $x\in\productFamily{E}$.
        Let $R = \{p\in\naturalsPlus\times D \mid \snd{p} = x((\fst{p} + \fst{p}) + 1)\}$.
        $R$ is a relation by \cref{relation,subseteq,times_elem_is_tuple}.
        For all $n\in\naturalsPlus$ we have $x((n + n) + 1)\in D$
            by \cref{naturalsplus_add_closed,real_naturals_closed_succ,indexed_product,function_apply_intro,subseteq}.
        For all $n\in\naturalsPlus$ there exists $d$ such that $n\mathrel{R} d$
            by \cref{relation,subseteq,times_tuple_intro,fst_eq,snd_eq,pair_eq_iff}.
        Take $y$ such that
            $y$ is a function on $\naturalsPlus$ and
            for all $n\in\naturalsPlus$ we have $n\mathrel{R}\apply{y}{n}$
            by \cref{choice_function}.
        $y\in\funs{\naturalsPlus}{D}$
            by \cref{choice_function,relation,subseteq,fst_eq,snd_eq,pair_eq_iff,indexed_product,naturalsplus_add_closed,real_naturals_closed_succ,function_apply_intro,funs_intro}.
        $y\in\funs{\dom{E}}{\unions{\ran{E}}}$
            by \cref{function_apply_intro,funs_type_apply,subseteq,function_apply_elim,ran_intro,unions_iff,upair_intro_left,funs_intro}.
        $y\in\productFamily{E}$ and $x\mathrel{Q} y$
            by \cref{indexed_product,subseteq,times_tuple_intro,choice_function,fst_eq,snd_eq,pair_eq_iff,relation}.
    \end{subproof}
    Take $o$ such that
        $o$ is a function on $\productFamily{E}$ and
        for all $x\in\productFamily{E}$ we have $x\mathrel{Q}\apply{o}{x}$
        by \cref{choice_function}.
    $o\in\funs{\productFamily{E}}{\productFamily{E}}$
        by \cref{choice_function,funs_intro,relation,subseteq,times_tuple_elim}.
    $o$ is a successor-of-even binary extractor on $E$
        by \cref{space_filling_odd_binary_extractor_on,times_elem_is_tuple,fst_eq,snd_eq,pair_eq_iff,choice_function,relation,subseteq}.

    Show that for all $n\in\dom{E}$ we have $\piIndex{E}{n}\circ o$ is continuous from $\productFamily{E}$ to $\apply{E}{n}$ with respect to $P_B$ and $\apply{U}{n}$.
    \begin{subproof}
        Fix $n\in\dom{E}$.
        Let $g = \piIndex{E}{n}\circ o$.
        $g\in\funs{\productFamily{E}}{E(n)}$ by \cref{projection_map_indexed_function,funs_circ}.
        $n + n\in\naturalsPlus$ by \cref{naturalsplus_add_closed,subseteq}.
        $(n + n) + 1\in\naturalsPlus$ by \cref{real_naturals_closed_succ}.
        $\piIndex{E}{((n + n) + 1)}\in\funs{\productFamily{E}}{E((n + n) + 1)}$
            by \cref{projection_map_indexed_function,subseteq}.
        For all $x\in\productFamily{E}$ we have $g(x) = \apply{\piIndex{E}{((n + n) + 1)}}{x}$
            by \cref{funs_elim,funs_is_function,funs_type_apply,projection_map_indexed_function,funs_ran,circ_apply,projection_map_indexed,function_apply_intro,choice_function,relation,subseteq,fst_eq,snd_eq,pair_eq_iff}.
        $g = \piIndex{E}{((n + n) + 1)}$ by \cref{functions_eq_iff}.
        $\piIndex{E}{((n + n) + 1)}$ is continuous from $\productFamily{E}$ to $E((n + n) + 1)$
            with respect to $P_B$ and $U((n + n) + 1)$ by \cref{projection_map_indexed_continuous,subseteq}.
        $g$ is continuous from $\productFamily{E}$ to $E(n)$ with respect to $P_B$ and $U(n)$
            by \cref{function_apply_intro,subseteq}.
    \end{subproof}
    $o$ is continuous from $\productFamily{E}$ to $\productFamily{E}$ with respect to $P_B$ and $P_B$
        by \cref{continuous_map_into_indexed_product,real_one_in_naturals,subseteq}.
\end{proof}

% from §44 helper lemma 126: two binary sequences can be interleaved into one
\begin{lemma}\label{space_filling_binary_interleaving}
    Let $I = \intervalclosed{0}{1}$.
    Let $D = \{0,1\}$.
    Let $E = \{(n,D) \mid n\in\naturalsPlus\}$.
    Assume $x_1\in\productFamily{E}$.
    Assume $x_2\in\productFamily{E}$.
    Then there exists $x\in\productFamily{E}$ such that
        for all $n\in\naturalsPlus$ we have
            $x(n + n) = x_1(n)$ and
            $x((n + n) + 1) = x_2(n)$.
\end{lemma}
\begin{proof}
    We have $E$ is a function on $\naturalsPlus$
        by \cref{relation,rightunique,dom_intro,dom_iff,pair_eq_iff,subseteq,subseteq_antisymmetric}.
    $\dom{E} = \naturalsPlus$ by \cref{funs}.
    For all $n\in\naturalsPlus$ we have $E(n) = D$ by \cref{function_apply_intro}.
    Let $R_1 = \{p\in\naturalsPlus\times D \mid \text{$\fst{p} = 1$ and $\snd{p} = 0$}\}$.
    Let $R_2 = \{p\in\naturalsPlus\times D \mid \text{there exists $k\in\naturalsPlus$ such that
        $\fst{p} = k + k$ and
        $\snd{p} = x_1(k)$}\}$.
    Let $R_3 = \{p\in\naturalsPlus\times D \mid \text{there exists $k\in\naturalsPlus$ such that
        $\fst{p} = (k + k) + 1$ and
        $\snd{p} = x_2(k)$}\}$.
    Let $R = R_1\union(R_2\union R_3)$.
    $R_1$ is a relation and $R_2$ is a relation and $R_3$ is a relation
        by \cref{relation,subseteq,times_elem_is_tuple}.
    $R$ is a relation by \cref{relation,subseteq,union_iff}.
    $R\subseteq\naturalsPlus\times D$ by \cref{subseteq,union_iff}.
    For all $n\in\naturalsPlus$ we have $n = 1$ or there exists $k\in\naturalsPlus$ such that
        $n = k + k$ or
        $n = (k + k) + 1$ by \cref{space_filling_naturalsplus_parity_decomposition}.
    For all $n\in\naturalsPlus$ there exists $d$ such that $n\mathrel{R} d$
        by \cref{upair_intro_left,indexed_product,function_apply_intro,times_tuple_intro,subseteq,fst_eq,snd_eq,pair_eq_iff,union_intro_left,union_intro_right,relation}.
    Take $x$ such that
        $x$ is a function on $\naturalsPlus$ and
        for all $n\in\naturalsPlus$ we have $n\mathrel{R}\apply{x}{n}$
        by \cref{choice_function}.
    $x$ is a function and $\dom{x} = \naturalsPlus$ by \cref{choice_function}.
    $x\in\funs{\naturalsPlus}{D}$ by \cref{choice_function,relation,subseteq,times_tuple_elim,funs_intro}.
    For all $n\in\dom{E}$ we have $x(n)\in E(n)$ by \cref{function_apply_intro,funs_type_apply,subseteq}.
    $x\in\funs{\dom{E}}{\unions{\ran{E}}}$
        by \cref{subseteq,function_apply_elim,ran_intro,unions_iff,upair_intro_left,funs_intro}.
    $x\in\productFamily{E}$ by \cref{indexed_product,subseteq}.

    Show that for all $n\in\naturalsPlus$ we have
        $x(n + n) = x_1(n)$ and
        $x((n + n) + 1) = x_2(n)$.
    \begin{subproof}
        Fix $n\in\naturalsPlus$.
        $n + n\in\naturalsPlus$ by \cref{naturalsplus_add_closed}.
        $(n + n) + 1\in\naturalsPlus$ by \cref{real_naturals_closed_succ}.

        Show that $x(n + n) = x_1(n)$.
        \begin{subproof}
            We have $((n + n),x(n + n))\in R$ by \cref{choice_function,relation}.
            \begin{byCase}
            \caseOf{$((n + n),x(n + n))\in R_1$.}
                Suppose not.
                We have $\fst{((n + n),x(n + n))} = 1$ and
                    $\snd{((n + n),x(n + n))} = 0$
                    by \cref{subseteq}.
                $n + n = 1$ by \cref{fst_eq,snd_eq,pair_eq_iff}.
                $1 \leq n$ by \cref{real_one_leq_naturalsplus}.
                $n < n + 1$ by \cref{real_naturals_lt_succ_local}.
                $1\in\reals$ and $n\in\reals$ and $n + 1\in\reals$
                    by \cref{real_one_in_naturals,real_naturals_subset_reals,subseteq,real_add_closed}.
                $1 < n + 1$ by \cref{real_leq_lt_trans}.
                $n + 1 \leq n + n$ by \cref{space_filling_succ_leq_double}.
                $1 < n + n$ by \cref{real_lt_leq_trans_local}.
                $1 < 1$ by \cref{subseteq,real_lt_subst_right}.
                Contradiction by \cref{real_lt_irreflexive}.
            \caseOf{$((n + n),x(n + n))\in R_2\union R_3$.}
                $((n + n),x(n + n))\in R_2$ or $((n + n),x(n + n))\in R_3$ by \cref{union_iff}.
                \begin{byCase}
                \caseOf{$((n + n),x(n + n))\in R_2$.}
                    Take $k\in\naturalsPlus$ such that
                        $\fst{((n + n),x(n + n))} = k + k$ and
                        $\snd{((n + n),x(n + n))} = x_1(k)$
                        by \cref{subseteq}.
                    $n + n = k + k$ and $x(n + n) = x_1(k)$ by \cref{fst_eq,snd_eq,pair_eq_iff}.
                    $n = k$ by \cref{space_filling_double_eq_halves_eq}.
                    $x(n + n) = x_1(n)$ by \cref{subseteq}.
                \caseOf{$((n + n),x(n + n))\in R_3$.}
                    Suppose not.
                    Take $k\in\naturalsPlus$ such that
                        $\fst{((n + n),x(n + n))} = (k + k) + 1$ and
                        $\snd{((n + n),x(n + n))} = x_2(k)$
                        by \cref{subseteq}.
                    $n + n = (k + k) + 1$ by \cref{fst_eq,snd_eq,pair_eq_iff}.
                    $n + n \neq (k + k) + 1$ by \cref{space_filling_double_neq_double_succ}.
                    Contradiction by \cref{subseteq}.
                \end{byCase}
            \end{byCase}
        \end{subproof}

        Show that $x((n + n) + 1) = x_2(n)$.
        \begin{subproof}
            We have $(((n + n) + 1),x((n + n) + 1))\in R$ by \cref{choice_function,relation}.
            \begin{byCase}
            \caseOf{$(((n + n) + 1),x((n + n) + 1))\in R_1$.}
                Suppose not.
                We have $\fst{(((n + n) + 1),x((n + n) + 1))} = 1$ and
                    $\snd{(((n + n) + 1),x((n + n) + 1))} = 0$
                    by \cref{subseteq}.
                $(n + n) + 1 = 1$ by \cref{fst_eq,snd_eq,pair_eq_iff}.
                $n + n < (n + n) + 1$ by \cref{real_naturals_lt_succ_local}.
                $1 \leq n + n$ by \cref{real_one_leq_naturalsplus}.
                $1\in\reals$ and $n + n\in\reals$ and $(n + n) + 1\in\reals$
                    by \cref{real_one_in_naturals,real_naturals_subset_reals,naturalsplus_add_closed,subseteq,real_add_closed}.
                $1 < (n + n) + 1$ by \cref{real_leq_lt_trans}.
                $1 < 1$ by \cref{subseteq,real_lt_subst_right}.
                Contradiction by \cref{real_lt_irreflexive}.
            \caseOf{$(((n + n) + 1),x((n + n) + 1))\in R_2\union R_3$.}
                $(((n + n) + 1),x((n + n) + 1))\in R_2$ or
                    $(((n + n) + 1),x((n + n) + 1))\in R_3$ by \cref{union_iff}.
                \begin{byCase}
                \caseOf{$(((n + n) + 1),x((n + n) + 1))\in R_2$.}
                    Suppose not.
                    Take $k\in\naturalsPlus$ such that
                        $\fst{(((n + n) + 1),x((n + n) + 1))} = k + k$ and
                        $\snd{(((n + n) + 1),x((n + n) + 1))} = x_1(k)$
                        by \cref{subseteq}.
                    $(n + n) + 1 = k + k$ by \cref{fst_eq,snd_eq,pair_eq_iff}.
                    $k + k \neq (n + n) + 1$ by \cref{space_filling_double_neq_double_succ}.
                    Contradiction by \cref{subseteq}.
                \caseOf{$(((n + n) + 1),x((n + n) + 1))\in R_3$.}
                    Take $k\in\naturalsPlus$ such that
                        $\fst{(((n + n) + 1),x((n + n) + 1))} = (k + k) + 1$ and
                        $\snd{(((n + n) + 1),x((n + n) + 1))} = x_2(k)$
                        by \cref{subseteq}.
                    $(n + n) + 1 = (k + k) + 1$ and $x((n + n) + 1) = x_2(k)$
                        by \cref{fst_eq,snd_eq,pair_eq_iff}.
                    $n + n\in\reals$ and $k + k\in\reals$ and $1\in\reals$
                        by \cref{real_naturals_subset_reals,naturalsplus_add_closed,subseteq,real_one_in_naturals,real_add_closed}.
                    $n + n = k + k$ by \cref{real_add_cancel}.
                    $n = k$ by \cref{space_filling_double_eq_halves_eq}.
                    $x((n + n) + 1) = x_2(n)$ by \cref{subseteq}.
                \end{byCase}
            \end{byCase}
        \end{subproof}
        Therefore $x(n + n) = x_1(n)$ and $x((n + n) + 1) = x_2(n)$ by \cref{subseteq}.
    \end{subproof}
\end{proof}

% from §44 helper lemma 127: a chosen dyadic code-point assignment is surjective
\begin{lemma}\label{space_filling_dyadic_code_map_surjective}
    Let $I = \intervalclosed{0}{1}$.
    Let $D = \{0,1\}$.
    Let $E = \{(n,D) \mid n\in\naturalsPlus\}$.
    Suppose $r\in\funs{\productFamily{E}}{I}$.
    Suppose for all $x\in\productFamily{E}$ we have $r(x)$ is a dyadic code point of $x$ in $I$ and $E$.
    Then $r\in\surj{\productFamily{E}}{I}$.
\end{lemma}
\begin{proof}
    Follows by \cref{space_filling_point_has_dyadic_code,space_filling_dyadic_code_point_unique,subseteq,surj}.
\end{proof}

% from §44 helper lemma 128: any two interval coordinates determine a square point
\begin{lemma}\label{space_filling_square_point_exists}
    Let $I = \intervalclosed{0}{1}$.
    Let $m_2 = 1 + 1$.
    Let $C = \{(i,I) \mid i\in\initialSegment{m_2}\}$.
    Assume $a\in I$.
    Assume $b\in I$.
    Then there exists $y\in\productFamily{C}$ such that
        $y(1) = a$ and
        $y(m_2) = b$.
\end{lemma}
\begin{proof}
    Let $J = \initialSegment{m_2}$.
    We have $C$ is a function on $J$ and $\dom{C} = J$
        by \cref{relation,rightunique,dom_intro,dom_iff,pair_eq_iff,subseteq,subseteq_antisymmetric,funs}.
    $1\in J$ and $m_2\in J$
        by \cref{real_one_in_naturals,real_naturals_closed_succ,real_one_leq_naturalsplus,initial_segment_naturalsplus}.
    For all $i\in J$ we have $i = 1$ or $i = m_2$
        by \cref{initial_segment_naturalsplus,real_naturals_leq_succ_elim,real_one_leq_naturalsplus,real_naturals_subset_reals,real_one_in_naturals,real_leq_antisym,subseteq}.
    For all $i\in J$ we have $C(i) = I$ by \cref{function_apply_intro}.
    $0\in\reals$ and $1\in\reals$ and $m_2\in\reals$
        by \cref{real_add_ident,real_one_in_naturals,real_naturals_subset_reals,real_naturals_closed_succ,subseteq}.
    $1 < m_2$ by \cref{real_zero_lt_one,real_lt_add,real_add_ident,subseteq}.
    $1 \neq m_2$ by \cref{real_lt_irreflexive,real_lt_subst_right}.

    Let $R_1 = \{p\in J\times I \mid \text{$\fst{p} = 1$ and $\snd{p} = a$}\}$.
    Let $R_2 = \{p\in J\times I \mid \text{$\fst{p} = m_2$ and $\snd{p} = b$}\}$.
    Let $R = R_1\union R_2$.
    $R_1$ is a relation and $R_2$ is a relation by \cref{relation,subseteq,times_elem_is_tuple}.
    $R$ is a relation by \cref{relation,subseteq,union_iff}.
    $R\subseteq J\times I$ by \cref{subseteq,union_iff}.
    For all $i\in J$ there exists $t$ such that $i\mathrel{R} t$
        by \cref{subseteq,times_tuple_intro,fst_eq,snd_eq,pair_eq_iff,union_intro_left,union_intro_right,relation}.
    Take $y$ such that
        $y$ is a function on $J$ and
        for all $i\in J$ we have $i\mathrel{R}\apply{y}{i}$
        by \cref{choice_function}.
    $y$ is a function and $\dom{y} = J$ by \cref{choice_function}.
    $y\in\funs{J}{I}$ by \cref{choice_function,relation,subseteq,times_tuple_elim,funs_intro}.
    $y\in\productFamily{C}$ by \cref{funs_intro,indexed_product,funs_elim,function_apply_intro,ran_intro,unions_iff}.

    Show that $y(1) = a$.
    \begin{subproof}
        We have $(1,y(1))\in R_1$
            by \cref{choice_function,subseteq,relation,union_iff,fst_eq,snd_eq,pair_eq_iff}.
        Follows by \cref{fst_eq,snd_eq,pair_eq_iff,subseteq}.
    \end{subproof}
    Show that $y(m_2) = b$.
    \begin{subproof}
        We have $(m_2,y(m_2))\in R_2$
            by \cref{choice_function,subseteq,relation,union_iff,fst_eq,snd_eq,pair_eq_iff}.
        Follows by \cref{fst_eq,snd_eq,pair_eq_iff,subseteq}.
    \end{subproof}
    Therefore there exists $w\in\productFamily{C}$ such that
        $w(1) = a$ and
        $w(m_2) = b$ by \cref{subseteq}.
\end{proof}

% from §44 helper lemma 129: the even extractor recovers the even subsequence of an interleaving
\begin{lemma}\label{space_filling_even_extractor_interleaving_inverse}
    Let $I = \intervalclosed{0}{1}$.
    Let $D = \{0,1\}$.
    Let $E = \{(n,D) \mid n\in\naturalsPlus\}$.
    Suppose $e$ is an even binary extractor on $E$.
    Suppose $x\in\productFamily{E}$.
    Suppose $x_1\in\productFamily{E}$.
    Suppose for all $n\in\naturalsPlus$ we have $x(n + n) = x_1(n)$.
    Then $e(x) = x_1$.
\end{lemma}
\begin{proof}
    We have $E$ is a function on $\naturalsPlus$ and $\dom{E} = \naturalsPlus$
        by \cref{relation,rightunique,dom_intro,dom_iff,pair_eq_iff,subseteq,subseteq_antisymmetric,funs}.
    For all $n\in\naturalsPlus$ we have $\apply{\apply{e}{x}}{n} = x_1(n)$
        by \cref{times_tuple_intro,space_filling_even_binary_extractor_on,subseteq,fst_eq,snd_eq,pair_eq_iff}.
    Therefore $e(x) = x_1$
        by \cref{functions_eq_iff,funs_elim,funs_type_apply,indexed_product,subseteq,space_filling_even_binary_extractor_on}.
\end{proof}

% from §44 helper lemma 130: the odd extractor recovers the odd subsequence of an interleaving
\begin{lemma}\label{space_filling_odd_extractor_interleaving_inverse}
    Let $I = \intervalclosed{0}{1}$.
    Let $D = \{0,1\}$.
    Let $E = \{(n,D) \mid n\in\naturalsPlus\}$.
    Suppose $o$ is a successor-of-even binary extractor on $E$.
    Suppose $x\in\productFamily{E}$.
    Suppose $x_2\in\productFamily{E}$.
    Suppose for all $n\in\naturalsPlus$ we have $x((n + n) + 1) = x_2(n)$.
    Then $o(x) = x_2$.
\end{lemma}
\begin{proof}
    We have $E$ is a function on $\naturalsPlus$ and $\dom{E} = \naturalsPlus$
        by \cref{relation,rightunique,dom_intro,dom_iff,pair_eq_iff,subseteq,subseteq_antisymmetric,funs}.
    For all $n\in\naturalsPlus$ we have $\apply{\apply{o}{x}}{n} = x_2(n)$
        by \cref{times_tuple_intro,space_filling_odd_binary_extractor_on,subseteq,fst_eq,snd_eq,pair_eq_iff}.
    Therefore $o(x) = x_2$
        by \cref{functions_eq_iff,funs_elim,funs_type_apply,indexed_product,subseteq,space_filling_odd_binary_extractor_on}.
\end{proof}

% from §44 helper lemma 131: evaluating a composite at a point
\begin{lemma}\label{space_filling_composite_apply}
    Suppose $f = g\circ h$.
    Suppose $h\in\funs{X}{Y}$.
    Suppose $g\in\funs{Y}{Z}$.
    Suppose $x\in X$.
    Then $f(x) = g(h(x))$.
\end{lemma}
\begin{proof}
    Follows by \cref{funs,subseteq,funs_ran,funs_is_function,function_rightunique,funs_is_relation,circ_apply}.
\end{proof}

% from §44 helper lemma 132: the binary product continuously surjects onto the unit square
\begin{lemma}\label{space_filling_binary_square_surjection}
    Let $I = \intervalclosed{0}{1}$.
    Let $T_I = \subspaceTopology{\standardTopologyR}{I}$.
    Let $D = \{0,1\}$.
    Let $T_D = \subspaceTopology{T_I}{D}$.
    Let $E = \{(n,D) \mid n\in\naturalsPlus\}$.
    Let $U = \{(n,T_D) \mid n\in\naturalsPlus\}$.
    Let $P_B = \productTopologyIndexed{U}{E}$.
    Let $m_2 = 1 + 1$.
    Let $C = \{(i,I) \mid i\in\initialSegment{m_2}\}$.
    Let $S = \{(i,T_I) \mid i\in\initialSegment{m_2}\}$.
    Let $P = \productTopologyIndexed{S}{C}$.
    Then there exists $h\in\surj{\productFamily{E}}{\productFamily{C}}$ such that
        $h$ is continuous from $\productFamily{E}$ to $\productFamily{C}$ with respect to $P_B$ and $P$.
\end{lemma}
\begin{proof}
    Take $r\in\funs{\productFamily{E}}{I}$ such that
        for all $x\in\productFamily{E}$ we have $r(x)$ is a dyadic code point of $x$ in $I$ and $E$
        by \cref{space_filling_dyadic_code_map_exists}.
    $r$ is continuous from $\productFamily{E}$ to $I$ with respect to $P_B$ and $T_I$
        by \cref{space_filling_dyadic_code_map_continuous}.
    $r\in\surj{\productFamily{E}}{I}$ by \cref{space_filling_dyadic_code_map_surjective}.

    Take $e$ such that
        $e$ is an even binary extractor on $E$ and
        $e$ is continuous from $\productFamily{E}$ to $\productFamily{E}$ with respect to $P_B$ and $P_B$
        by \cref{space_filling_even_binary_extractor}.
    $e\in\funs{\productFamily{E}}{\productFamily{E}}$ by \cref{space_filling_even_binary_extractor_on}.
    Take $o$ such that
        $o$ is a successor-of-even binary extractor on $E$ and
        $o$ is continuous from $\productFamily{E}$ to $\productFamily{E}$ with respect to $P_B$ and $P_B$
        by \cref{space_filling_odd_binary_extractor}.
    $o\in\funs{\productFamily{E}}{\productFamily{E}}$ by \cref{space_filling_odd_binary_extractor_on}.

    Let $J = \initialSegment{m_2}$.
    We have $C$ is a function on $J$ and $S$ is a function on $J$
        by \cref{relation,rightunique,dom_intro,dom_iff,pair_eq_iff,subseteq,subseteq_antisymmetric}.
    $\dom{C} = J$ and $\dom{S} = J$ by \cref{funs}.
    $1\in J$ and $m_2\in J$
        by \cref{real_one_in_naturals,real_naturals_closed_succ,real_one_leq_naturalsplus,initial_segment_naturalsplus}.
    For all $i\in J$ we have $C(i) = I$ and $S(i) = T_I$ by \cref{function_apply_intro}.
    $I\subseteq\reals$ by \cref{intervalclosed,subseteq}.
    $\standardTopologyR$ is a topology on $\reals$
        by \cref{standard_basis_is_basis,basis_topology_is_topology,standard_topology_reals}.
    $I$ is a subspace of $\reals$ with respect to $\standardTopologyR$ by \cref{subspace_of}.
    $T_I$ is a topology on $I$ by \cref{subspace_topology_is_topology}.
    For all $i\in J$ we have $S(i)$ is a topology on $C(i)$ by \cref{function_apply_intro,subseteq}.
    $J$ is inhabited by \cref{real_one_in_naturals,real_naturals_closed_succ,real_one_leq_naturalsplus,initial_segment_naturalsplus}.
    $\productFamily{E}$ is compact with respect to $P_B$ and
        $\productFamily{E}$ is Hausdorff with respect to $P_B$
        by \cref{space_filling_binary_product_compact_hausdorff}.
    $P_B$ is a topology on $\productFamily{E}$ by \cref{hausdorff}.

    Let $f_1 = r\circ e$.
    Let $f_2 = r\circ o$.
    $f_1$ is continuous from $\productFamily{E}$ to $I$ with respect to $P_B$ and $T_I$ and
        $f_2$ is continuous from $\productFamily{E}$ to $I$ with respect to $P_B$ and $T_I$
        by \cref{continuous_composite}.

    Let $Q = \{w\in\productFamily{E}\times\productFamily{C} \mid
        \apply{\snd{w}}{1} = f_1(\fst{w}) \land \apply{\snd{w}}{m_2} = f_2(\fst{w})\}$.
    $Q$ is a relation by \cref{relation,subseteq,times_elem_is_tuple}.
    For all $x\in\productFamily{E}$ we have $f_1(x)\in I$ and $f_2(x)\in I$ by \cref{continuous,funs_type_apply}.
    For all $x\in\productFamily{E}$ there exists $y\in\productFamily{C}$ such that
        $y(1) = f_1(x)$ and
        $y(m_2) = f_2(x)$ by \cref{space_filling_square_point_exists}.
    For all $x\in\productFamily{E}$ there exists $y$ such that $x\mathrel{Q} y$
        by \cref{relation,times_tuple_intro,fst_eq,snd_eq,pair_eq_iff,subseteq}.
    Take $h$ such that
        $h$ is a function on $\productFamily{E}$ and
        for all $x\in\productFamily{E}$ we have $x\mathrel{Q}\apply{h}{x}$
        by \cref{choice_function}.
    $h\in\funs{\productFamily{E}}{\productFamily{C}}$
        by \cref{choice_function,funs_intro,relation,subseteq,times_tuple_elim}.

    Show that for all $i\in J$ we have $\piIndex{C}{i}\circ h$ is continuous from $\productFamily{E}$ to $C(i)$ with respect to $P_B$ and $S(i)$.
    \begin{subproof}
        Fix $i\in J$.
        Let $g = \piIndex{C}{i}\circ h$.
        $\piIndex{C}{i}\in\funs{\productFamily{C}}{C(i)}$ by \cref{projection_map_indexed_function,subseteq}.
        $g\in\funs{\productFamily{E}}{C(i)}$ by \cref{funs_circ}.
        $i = 1$ or $i = m_2$ by \cref{initial_segment_naturalsplus,real_naturals_leq_succ_elim,real_one_leq_naturalsplus,real_naturals_subset_reals,real_one_in_naturals,real_leq_antisym,subseteq}.
        \begin{byCase}
        \caseOf{$i = 1$.}
            $f_1\in\funs{\productFamily{E}}{C(i)}$ by \cref{continuous,function_apply_intro,subseteq}.
            For all $x\in\productFamily{E}$ we have $g(x) = f_1(x)$
                by \cref{funs_elim,funs_is_function,funs_type_apply,projection_map_indexed_function,funs_ran,circ_apply,projection_map_indexed,function_apply_intro,choice_function,relation,subseteq,fst_eq,snd_eq,pair_eq_iff}.
            $g = f_1$ by \cref{functions_eq_iff}.
            $g$ is continuous from $\productFamily{E}$ to $C(i)$ with respect to $P_B$ and $S(i)$
                by \cref{continuous,continuous_eq_subst,function_apply_intro,subseteq}.
            $\piIndex{C}{i}\circ h$ is continuous from $\productFamily{E}$ to $C(i)$ with respect to $P_B$ and $S(i)$
                by \cref{subseteq}.
        \caseOf{$i = m_2$.}
            $f_2\in\funs{\productFamily{E}}{C(i)}$ by \cref{continuous,function_apply_intro,subseteq}.
            For all $x\in\productFamily{E}$ we have $g(x) = f_2(x)$
                by \cref{funs_elim,funs_is_function,funs_type_apply,projection_map_indexed_function,funs_ran,circ_apply,projection_map_indexed,function_apply_intro,choice_function,relation,subseteq,fst_eq,snd_eq,pair_eq_iff}.
            $g = f_2$ by \cref{functions_eq_iff}.
            $g$ is continuous from $\productFamily{E}$ to $C(i)$ with respect to $P_B$ and $S(i)$
                by \cref{continuous,continuous_eq_subst,function_apply_intro,subseteq}.
            $\piIndex{C}{i}\circ h$ is continuous from $\productFamily{E}$ to $C(i)$ with respect to $P_B$ and $S(i)$
                by \cref{subseteq}.
        \end{byCase}
    \end{subproof}
    $h$ is continuous from $\productFamily{E}$ to $\productFamily{C}$ with respect to $P_B$ and $P$
        by \cref{continuous_map_into_indexed_product,subspace_topology_is_topology,funs,subseteq}.

    Show that for all $y\in\productFamily{C}$ there exists $x\in\productFamily{E}$ such that $h(x) = y$.
    \begin{subproof}
        Fix $y\in\productFamily{C}$.
        $y(1)\in I$ and $y(m_2)\in I$
            by \cref{indexed_product,function_apply_intro,subseteq}.
        Take $x_1\in\productFamily{E}$ such that $r(x_1) = y(1)$ by \cref{surj}.
        Take $x_2\in\productFamily{E}$ such that $r(x_2) = y(m_2)$ by \cref{surj}.
        Take $x\in\productFamily{E}$ such that
            for all $n\in\naturalsPlus$ we have
                $x(n + n) = x_1(n)$ and
                $x((n + n) + 1) = x_2(n)$
            by \cref{space_filling_binary_interleaving}.
        Show that $e(x) = x_1$.
        \begin{subproof}
            Follows by \cref{space_filling_even_extractor_interleaving_inverse}.
        \end{subproof}
        Show that $o(x) = x_2$.
        \begin{subproof}
            Follows by \cref{space_filling_odd_extractor_interleaving_inverse}.
        \end{subproof}
        Show that $f_1(x) = y(1)$.
        \begin{subproof}
            Follows by \cref{space_filling_composite_apply,subseteq}.
        \end{subproof}
        Show that $f_2(x) = y(m_2)$.
        \begin{subproof}
            Follows by \cref{space_filling_composite_apply,subseteq}.
        \end{subproof}
        $(x,h(x))\in Q$ by \cref{choice_function,relation}.
        $h(x)\in\productFamily{C}$ by \cref{funs_type_apply,choice_function,subseteq}.
        $(x,h(x))\in\productFamily{E}\times\productFamily{C}$ and
            $\apply{\snd{(x,h(x))}}{1} = f_1(\fst{(x,h(x))})$ and
            $\apply{\snd{(x,h(x))}}{m_2} = f_2(\fst{(x,h(x))})$
            by \cref{choice_function,relation,funs_type_apply,subseteq}.
        Show that $\apply{\apply{h}{x}}{1} = y(1)$.
        \begin{subproof}
            Follows by \cref{fst_eq,snd_eq,pair_eq_iff,subseteq}.
        \end{subproof}
        Show that $\apply{\apply{h}{x}}{m_2} = y(m_2)$.
        \begin{subproof}
            Follows by \cref{fst_eq,snd_eq,pair_eq_iff,subseteq}.
        \end{subproof}
        $h(x) = y$ by \cref{space_filling_square_point_eq}.
        There exists $z\in\productFamily{E}$ such that $h(z) = y$.
    \end{subproof}
    Therefore $h\in\surj{\productFamily{E}}{\productFamily{C}}$ by \cref{surj}.
\end{proof}

% from §44 helper lemma 6: classical closed-subspace surjection onto the unit square
\begin{lemma}\label{space_filling_closed_subspace_surjection}
    Let $I = \intervalclosed{0}{1}$.
    Let $T_I = \subspaceTopology{\standardTopologyR}{I}$.
    Let $C = \{(i,I) \mid i\in\initialSegment{1 + 1}\}$.
    Let $S = \{(i,T_I) \mid i\in\initialSegment{1 + 1}\}$.
    Let $P = \productTopologyIndexed{S}{C}$.
    Then there exist $A,h$ such that
        $A$ is closed in $I$ with respect to $T_I$ and
        $h$ is continuous from $A$ to $\productFamily{C}$ with respect to $\subspaceTopology{T_I}{A}$ and $P$ and
        $h$ is surjective on $\productFamily{C}$.
\end{lemma}
\begin{proof}
    Let $D = \{0,1\}$.
    Let $T_D = \subspaceTopology{T_I}{D}$.
    Let $E = \{(n,D) \mid n\in\naturalsPlus\}$.
    Let $U = \{(n,T_D) \mid n\in\naturalsPlus\}$.
    Let $P_B = \productTopologyIndexed{U}{E}$.

    Take $r\in\funs{\productFamily{E}}{I}$ such that
        for all $x\in\productFamily{E}$ we have $r(x)$ is a retained-third code point of $x$ in $I$ and $E$
        by \cref{space_filling_retained_third_code_map_exists}.
    We have $r$ is continuous from $\productFamily{E}$ to $I$ with respect to $P_B$ and $T_I$
        by \cref{space_filling_retained_third_code_map_continuous}.
    $r$ is injective by \cref{space_filling_retained_third_code_map_injective}.

    Let $A = \img{r}{\productFamily{E}}$.
    $A\subseteq I$ by \cref{img_subseteq_ran,funs_ran,subseteq_transitive,subseteq}.
    $I\subseteq\reals$ by \cref{intervalclosed,subseteq}.
    $\standardTopologyR$ is a topology on $\reals$
        by \cref{standard_basis_is_basis,basis_topology_is_topology,standard_topology_reals}.
    Hence $I$ is a subspace of $\reals$ with respect to $\standardTopologyR$ by \cref{subspace_of}.
    $T_I$ is a topology on $I$ by \cref{subspace_topology_is_topology}.
    $\reals$ is Hausdorff with respect to $\standardTopologyR$
        by \cref{standard_metric_is_metric,metric_space_hausdorff,standard_metric_topology}.
    Thus $I$ is Hausdorff with respect to $T_I$ by \cref{subspace_hausdorff}.
    $\productFamily{E}$ is compact with respect to $P_B$ and
        $\productFamily{E}$ is Hausdorff with respect to $P_B$
        by \cref{space_filling_binary_product_compact_hausdorff}.
    Therefore $A$ is compact with respect to $\subspaceTopology{T_I}{A}$ by \cref{continuous_image_compact}.
    Hence $A$ is closed in $I$ with respect to $T_I$ by \cref{compact_subspace_closed}.

    We have $r$ is continuous from $\productFamily{E}$ to $A$ with respect to $P_B$ and $\subspaceTopology{T_I}{A}$
        by \cref{continuous_restrict_range,subspace_of,subseteq}.
    Thus $r\in\surj{\productFamily{E}}{A}$ by \cref{continuous,funs_is_function,funs,funs_intro,function_apply_elim,ran_intro,img_dom,funs_surj_iff}.
    Hence $r$ is a bijection from $\productFamily{E}$ to $A$ by \cref{bijections,injective_function}.
    $A$ is Hausdorff with respect to $\subspaceTopology{T_I}{A}$ by \cref{subspace_hausdorff}.
    Therefore $r$ is a homeomorphism between $\productFamily{E}$ and $A$ with respect to
        $P_B$ and $\subspaceTopology{T_I}{A}$ by \cref{compact_hausdorff_bijection_homeomorphism}.
    Thus $\converse{r}$ is continuous from $A$ to $\productFamily{E}$ with respect to
        $\subspaceTopology{T_I}{A}$ and $P_B$ by \cref{homeomorphism}.
    $\converse{r}\in\funs{A}{\productFamily{E}}$ by \cref{bijective_converse_are_funs,bijections}.

    Take $g\in\surj{\productFamily{E}}{\productFamily{C}}$ such that
        $g$ is continuous from $\productFamily{E}$ to $\productFamily{C}$ with respect to $P_B$ and $P$
        by \cref{space_filling_binary_square_surjection}.
    Let $h = g\circ \converse{r}$.
    Hence $h$ is continuous from $A$ to $\productFamily{C}$ with respect to $\subspaceTopology{T_I}{A}$ and $P$
        by \cref{continuous_composite}.
    We have $\converse{r}\in\surj{A}{\productFamily{E}}$ by \cref{bijection_converse_is_bijection,bijections}.
    Thus $h\in\surj{A}{\productFamily{C}}$ by \cref{surjections_circ}.
    Hence $h$ is surjective on $\productFamily{C}$.
    Follows by definition.
\end{proof}
\begin{theorem}\label{space_filling_curve_dense_image}
    Let $I = \intervalclosed{0}{1}$.
    Let $T_I = \subspaceTopology{\standardTopologyR}{I}$.
    Let $C = \{(i,I) \mid i\in\initialSegment{1 + 1}\}$.
    Let $S = \{(i,T_I) \mid i\in\initialSegment{1 + 1}\}$.
    Let $P = \productTopologyIndexed{S}{C}$.
    Then there exists $f$ such that
        $f$ is continuous from $I$ to $\productFamily{C}$ with respect to $T_I$ and $P$ and
        $\img{f}{I}$ is dense in $\productFamily{C}$ with respect to $P$.
\end{theorem}
\begin{proof}
    Let $m_2 = 1 + 1$.
    Take $A,h$ such that
        $A$ is closed in $I$ with respect to $T_I$ and
        $h$ is continuous from $A$ to $\productFamily{C}$ with respect to $\subspaceTopology{T_I}{A}$ and $P$ and
        $h$ is surjective on $\productFamily{C}$
        by \cref{space_filling_closed_subspace_surjection}.
    Let $T_A = \subspaceTopology{T_I}{A}$.
    $I$ is normal with respect to $T_I$ by \cref{space_filling_domain_normal}.
    Hence $T_I$ is a topology on $I$ by \cref{normal_space}.
    $P$ is a topology on $\productFamily{C}$ by \cref{space_filling_target_hausdorff,hausdorff}.

    For all $i,u,v$ such that $(i,u)\in C$ and $(i,v)\in C$ we have $u = v$ by \cref{pair_eq_iff}.
    For all $i,u,v$ such that $(i,u)\in S$ and $(i,v)\in S$ we have $u = v$ by \cref{pair_eq_iff}.
    Hence $C$ is a function and $S$ is a function by \cref{relation,rightunique}.
    Thus $\dom{C} = \initialSegment{m_2}$ and $\dom{S} = \initialSegment{m_2}$
        by \cref{dom_intro,dom_iff,pair_eq_iff,subseteq,subseteq_antisymmetric}.
    Hence $1\in\initialSegment{m_2}$ and $m_2\in\initialSegment{m_2}$
        by \cref{real_one_in_naturals,real_one_leq_naturalsplus,real_naturals_closed_succ,initial_segment_naturalsplus}.
    For all $i\in\initialSegment{m_2}$ we have $\apply{C}{i} = I$ and $\apply{S}{i} = T_I$
        by \cref{function_apply_intro}.
    We have for all $i\in\dom{C}$ $\apply{S}{i}$ is a topology on $\apply{C}{i}$ by \cref{subseteq}.

    We have for all $i\in\dom{C}$ $\piIndex{C}{i}$ is continuous from $\productFamily{C}$ to $\apply{C}{i}$
        with respect to $P$ and $\apply{S}{i}$ by \cref{projection_map_indexed_continuous}.
    Therefore $\piIndex{C}{1}$ is continuous from $\productFamily{C}$ to $I$ with respect to $P$ and $T_I$ and
        $\piIndex{C}{m_2}$ is continuous from $\productFamily{C}$ to $I$ with respect to $P$ and $T_I$
        by \cref{function_apply_intro,subseteq}.

    Let $h_1 = \piIndex{C}{1}\circ h$.
    Let $h_2 = \piIndex{C}{m_2}\circ h$.
    Therefore $h_1$ is continuous from $A$ to $I$ with respect to $T_A$ and $T_I$ and
        $h_2$ is continuous from $A$ to $I$ with respect to $T_A$ and $T_I$
        by \cref{continuous_composite}.

    $0\in\reals$ and $1\in\reals$ and $0 < 1$
        by \cref{real_add_ident,real_one_in_naturals,real_naturals_subset_reals,subseteq,real_zero_lt_one}.
    Take $g_1$ such that
        $g_1$ is continuous from $I$ to $I$ with respect to $T_I$ and $T_I$ and
        for all $x\in A$ we have $g_1(x) = h_1(x)$
        by \cref{tietze_extension_theorem_closed_interval}.
    Take $g_2$ such that
        $g_2$ is continuous from $I$ to $I$ with respect to $T_I$ and $T_I$ and
        for all $x\in A$ we have $g_2(x) = h_2(x)$
        by \cref{tietze_extension_theorem_closed_interval}.
    Hence $g_1$ is a function and $\dom{g_1} = I$ and
        $g_2$ is a function and $\dom{g_2} = I$ by \cref{continuous,funs_is_function,funs}.

    Let $f = \{(x,z) \mid x\in I, z\in\productFamily{C} \mid z(1) = g_1(x) \land z(m_2) = g_2(x)\}$.
    $f$ is a relation by \cref{relation}.
    For all $x,z_1,z_2$ such that $(x,z_1)\in f$ and $(x,z_2)\in f$ we have $z_1 = z_2$
        by \cref{pair_eq_iff,space_filling_square_point_eq}.
    Hence $f$ is right-unique by \cref{rightunique}.
    Thus $f$ is a function by \cref{relation}.
    For all $x\in I$ we have $g_1(x)\in I$ and $g_2(x)\in I$ by \cref{continuous,funs_type_apply}.
    For all $x\in I$ there exists $z\in\productFamily{C}$ such that
        $z(1) = g_1(x)$ and
        $z(m_2) = g_2(x)$ by \cref{space_filling_square_point_exists}.
    Hence for all $x\in I$ there exists $z$ such that $(x,z)\in f$ by \cref{pair_eq_iff}.
    Therefore $\dom{f} = I$ by \cref{dom_iff,pair_eq_iff,subseteq,subseteq_antisymmetric}.
    For all $x\in\dom{f}$ we have $f(x)\in\productFamily{C}$ by \cref{function_apply_iff,pair_eq_iff}.
    Thus $f\in\funs{I}{\productFamily{C}}$ by \cref{funs_intro}.

    Show that for all $i\in\dom{C}$ we have $\piIndex{C}{i}\circ f$ is continuous from $I$ to $\apply{C}{i}$ with respect to $T_I$ and $\apply{S}{i}$.
    \begin{subproof}
        Fix $i\in\dom{C}$.
        Let $g = \piIndex{C}{i}\circ f$.
        $g\in\funs{I}{\apply{C}{i}}$ by \cref{projection_map_indexed_function,funs_circ,funs_elim}.
        We have $g$ is a function and $\dom{g} = I$ by \cref{funs_elim,funs}.
        We have $i\in\initialSegment{m_2}$ by \cref{subseteq}.
        Hence $i = 1$ or $i = m_2$.
        \begin{byCase}
        \caseOf{$i = 1$.}
            For all $x\in I$ we have $g(x) = g_1(x)$
                by \cref{function_apply_elim,pair_eq_iff,real_add_eq_subst,funs_type_apply,projection_map_indexed,circ_elem_intro,funs_elim,function_apply_intro}.
            Hence $g = g_1$ by \cref{fun_subseteq,subseteq_refl,subseteq_antisymmetric}.
            We have $\apply{C}{i} = I$ and $\apply{S}{i} = T_I$ by \cref{function_apply_intro,real_add_eq_subst}.
            Therefore $g$ is continuous from $I$ to $\apply{C}{i}$ with respect to $T_I$ and $\apply{S}{i}$
                by \cref{continuous_eq_subst}.
            Thus $\piIndex{C}{i}\circ f$ is continuous from $I$ to $\apply{C}{i}$ with respect to $T_I$ and $\apply{S}{i}$
                by \cref{subseteq}.
        \caseOf{$i = m_2$.}
            For all $x\in I$ we have $g(x) = g_2(x)$
                by \cref{function_apply_elim,pair_eq_iff,real_add_eq_subst,funs_type_apply,projection_map_indexed,circ_elem_intro,funs_elim,function_apply_intro}.
            Thus $g = g_2$ by \cref{fun_subseteq,subseteq_refl,subseteq_antisymmetric}.
            We have $\apply{C}{i} = I$ and $\apply{S}{i} = T_I$ by \cref{function_apply_intro}.
            Therefore $g$ is continuous from $I$ to $\apply{C}{i}$ with respect to $T_I$ and $\apply{S}{i}$
                by \cref{continuous_eq_subst}.
            Thus $\piIndex{C}{i}\circ f$ is continuous from $I$ to $\apply{C}{i}$ with respect to $T_I$ and $\apply{S}{i}$
                by \cref{subseteq}.
        \end{byCase}
    \end{subproof}
    Therefore $f$ is continuous from $I$ to $\productFamily{C}$ with respect to $T_I$ and $P$
        by \cref{continuous_map_into_indexed_product,subseteq}.

    Show that for all $y$ such that $y\in\productFamily{C}$ we have $y\in\img{f}{I}$.
    \begin{subproof}
        Fix $y$ such that $y\in\productFamily{C}$.
        $h$ is right-total on $\productFamily{C}$ by definition.
        Take $x$ such that $x\mathrel{h} y$ by \cref{righttotal}.
        Hence $x\in A$ by \cref{dom_intro,continuous,funs,subseteq}.
        Thus $x\in I$ by \cref{closed_in,subseteq}.
        $(y,\apply{y}{1})\in\piIndex{C}{1}$ and $(y,\apply{y}{m_2})\in\piIndex{C}{m_2}$
            by \cref{projection_map_indexed,function_apply_elim}.
        Hence $(x,\apply{y}{1})\in h_1$ and $(x,\apply{y}{m_2})\in h_2$
            by \cref{circ_elem_intro,continuous,funs_ran,function_apply_elim}.
        Show that $h_1(x) = y(1)$.
        \begin{subproof}
            Follows by \cref{function_apply_intro,continuous,funs_is_function}.
        \end{subproof}
        Show that $h_2(x) = y(m_2)$.
        \begin{subproof}
            Follows by \cref{function_apply_intro,continuous,funs_is_function}.
        \end{subproof}
        Show that $g_1(x) = y(1)$.
        \begin{subproof}
            Follows by \cref{subseteq}.
        \end{subproof}
        Show that $g_2(x) = y(m_2)$.
        \begin{subproof}
            Follows by \cref{subseteq}.
        \end{subproof}
        Hence $(x,y)\in f$ by \cref{pair_eq_iff}.
        Thus $y\in\img{f}{I}$ by \cref{img_iff}.
    \end{subproof}
    Therefore $\img{f}{I} = \productFamily{C}$
        by \cref{subseteq,continuous,funs,rels_ran_subseteq,img_subseteq_ran,subseteq_transitive,subseteq_antisymmetric}.
    Hence $\img{f}{I}$ is dense in $\productFamily{C}$ with respect to $P$
        by \cref{dense_in,closed_eq_closure,closed_whole_space}.
\end{proof}
% from §44 Item 1: space-filling curve
\begin{theorem}\label{space_filling_curve}
    Let $I = \intervalclosed{0}{1}$.
    Let $T_I = \subspaceTopology{\standardTopologyR}{I}$.
    Let $C = \{(i,I) \mid i\in\initialSegment{1 + 1}\}$.
    Let $S = \{(i,T_I) \mid i\in\initialSegment{1 + 1}\}$.
    Let $P = \productTopologyIndexed{S}{C}$.
    Then there exists $f$ such that
        $f$ is continuous from $I$ to $\productFamily{C}$ with respect to $T_I$ and $P$ and
        $f$ is surjective on $\productFamily{C}$.
\end{theorem}
\begin{proof}
    Take $f$ such that
        $f$ is continuous from $I$ to $\productFamily{C}$ with respect to $T_I$ and $P$ and
        $\img{f}{I}$ is dense in $\productFamily{C}$ with respect to $P$
        by \cref{space_filling_curve_dense_image}.
    Let $A = \img{f}{I}$.
    $I$ is compact with respect to $T_I$ by \cref{space_filling_domain_compact}.
    Hence $A$ is compact with respect to $\subspaceTopology{P}{A}$ by \cref{continuous_image_compact}.
    $A\subseteq\productFamily{C}$
        by \cref{continuous,funs,rels_ran_subseteq,img_subseteq_ran,subseteq_transitive}.
    $P$ is a topology on $\productFamily{C}$ by \cref{space_filling_target_hausdorff,hausdorff}.
    Therefore $A$ is closed in $\productFamily{C}$ with respect to $P$
        by \cref{space_filling_target_hausdorff,compact_subspace_closed}.
    Hence $A = \productFamily{C}$ by \cref{closed_eq_closure,dense_in}.
    Therefore $f$ is surjective on $\productFamily{C}$
        by \cref{continuous,funs_is_function,img_of_function_elim,inter_elim_left,function_apply_elim,righttotal}.
\end{proof}
\section{§45 Compactness in Metric Spaces}

% from §45 helper definition 1: totally bounded subset of a metric space
\begin{definition}\label{totally_bounded_subset_metric_space}
    $A$ is totally bounded in $X$ with respect to $d$ iff
        $d$ is a metric on $X$ and
        $A\subseteq X$ and
        for all $\epsilon\in\reals$ such that $0 < \epsilon$
        there exists $C$ such that
            $C$ is finite and
            $C$ is a covering of $A$ and
            for all $U\in C$ there exists $x\in X$ such that
                $U = A\inter\metricBall{d}{X}{x}{\epsilon}$.
\end{definition}

% from §45 Item 1: totally bounded metric space
\begin{definition}\label{totally_bounded_metric_space}
    $X$ is totally bounded with respect to $d$ iff
        $X$ is totally bounded in $X$ with respect to $d$.
\end{definition}

% from §45 Item 2: compact metric spaces are complete
\begin{theorem}\label{compact_metric_space_complete}
    Assume $d$ is a metric on $X$.
    Let $T = \metricTopology{d}{X}$.
    Suppose $X$ is compact with respect to $T$.
    Then $X$ is complete with respect to $d$.
\end{theorem}
\begin{proof}
    $X$ is sequentially compact with respect to $T$
        by \cref{metric_topology,metric_basis_is_basis,basis_for_topology,basis_topology_is_topology,metrizable,metrizable_compact_iff_limit_point_compact_iff_sequentially_compact}.
    Hence $X$ is complete with respect to $d$
        by \cref{cauchy_sequence,sequentially_compact,cauchy_subsequence_completeness_criterion}.
\end{proof}

% from §45 Item 2: compact metric spaces are totally bounded
\begin{theorem}\label{compact_metric_space_totally_bounded}
    Assume $d$ is a metric on $X$.
    Let $T = \metricTopology{d}{X}$.
    Suppose $X$ is compact with respect to $T$.
    Then $X$ is totally bounded with respect to $d$.
\end{theorem}
\begin{proof}
    $X$ is sequentially compact with respect to $T$
        by \cref{metric_topology,metric_basis_is_basis,basis_for_topology,basis_topology_is_topology,metrizable,metrizable_compact_iff_limit_point_compact_iff_sequentially_compact}.
    Thus for all $\epsilon\in\reals$ such that $0 < \epsilon$
        there exists $C$ such that
            $C$ is finite and
            $C$ is a covering of $X$ and
            for all $U\in C$ there exists $x\in X$ such that
                $U = X\inter\metricBall{d}{X}{x}{\epsilon}$
        by \cref{sequentially_compact_finite_ball_cover,metric_ball,inter_absorb_supseteq_left,subseteq}.
    Therefore $X$ is totally bounded with respect to $d$
        by \cref{subseteq_refl,totally_bounded_subset_metric_space,totally_bounded_metric_space}.
\end{proof}

% from §45 helper lemma 1: finite covers of an infinite indexed image have an infinitely visited member
\begin{lemma}\label{finite_cover_infinite_sequence_trace}
    Suppose $s$ is a sequence in $X$.
    Suppose $J\subseteq\naturalsPlus$.
    Suppose $J$ is infinite.
    Suppose $C$ is finite.
    Suppose $C$ is a covering of $\img{s}{J}$.
    Then there exists $U\in C$ such that
        $J\inter\preimg{s}{U}$ is infinite.
\end{lemma}
\begin{proof}
    Let $h = \{(U, J\inter\preimg{s}{U}) \mid U\in C\}$.
    For all $U,V_1,V_2$ such that $(U,V_1)\in h$ and $(U,V_2)\in h$ we have $V_1 = V_2$
        by \cref{pair_eq_iff}.
    Hence $h$ is right-unique and $h$ is a function by \cref{rightunique,relation}.
    Therefore $\dom{h} = C$ by \cref{pair_eq_iff,dom_intro,dom_iff,subseteq,subseteq_antisymmetric}.
    Let $F = \img{h}{C}$.
    We have $F$ is finite by \cref{finite_image_function}.
    For all $n$ such that $n\in J$ we have $n\in\dom{s}$ by \cref{sequence_in,funs,subseteq}.
    Hence for all $n$ such that $n\in J$ we have $\apply{s}{n}\in\img{s}{J}$
        by \cref{img_iff,function_apply_iff,sequence_in,funs_is_function}.
    Therefore for all $n$ such that $n\in J$ there exists $U\in C$ such that $\apply{s}{n}\in U$
        by \cref{covering,subseteq,unions_iff}.
    Thus for all $n$ such that $n\in J$ there exists $U\in C$ such that $n\in J\inter\preimg{s}{U}$
        by \cref{inter_intro,preimg_iff,function_apply_iff,sequence_in,funs_is_function}.
    Hence for all $n$ such that $n\in J$ we have $n\in\unions{F}$
        by \cref{subseteq,img_iff,function_apply_iff,unions_iff}.
    Therefore $\unions{F} = J$ by \cref{subseteq,unions_iff,img_iff,pair_eq_iff,inter,subseteq_antisymmetric}.
    Suppose not.
    For all $U\in C$ we have $J\inter\preimg{s}{U}$ is finite.
    For all $K\in F$ we have $K$ is finite by \cref{img_iff,function_apply_iff}.
    Hence $J$ is finite by \cref{subseteq_antisymmetric,finite_union_of_finite_family}.
    Contradiction.
\end{proof}

% from §45 helper lemma 2a: instantiate a totally bounded cover element
\begin{lemma}\label{totally_bounded_cover_member_ball}
    Assume for all $U\in C$ there exists $x\in X$ such that
        $U = X\inter\metricBall{d}{X}{x}{\epsilon}$.
    Suppose $U\in C$.
    Then there exists $x\in X$ such that
        $U = X\inter\metricBall{d}{X}{x}{\epsilon}$.
\end{lemma}
\begin{proof}
    Follows by definition.
\end{proof}

% from §45 helper lemma 2: total boundedness yields an infinite epsilon-ball trace
\begin{lemma}\label{totally_bounded_infinite_sequence_epsilon_ball}
    Assume $d$ is a metric on $X$.
    Assume $X$ is totally bounded with respect to $d$.
    Suppose $s$ is a sequence in $X$.
    Suppose $J\subseteq\naturalsPlus$.
    Suppose $J$ is infinite.
    Suppose $\epsilon\in\reals$.
    Suppose $0 < \epsilon$.
    Then there exist $K,U$ such that
        $K\subseteq J$ and
        $K$ is infinite and
        for all $m\in K$ we have $\apply{s}{m}\in U$ and
        there exists $x\in X$ such that
            $U = X\inter\metricBall{d}{X}{x}{\epsilon}$.
\end{lemma}
\begin{proof}
    Take $C$ such that
        $C$ is finite and
        $C$ is a covering of $X$ and
        for all $U\in C$ there exists $x\in X$ such that
            $U = X\inter\metricBall{d}{X}{x}{\epsilon}$
        by \cref{totally_bounded_metric_space,totally_bounded_subset_metric_space}.
    Hence $C$ is a covering of $\img{s}{J}$
        by \cref{covering,img_iff,subseteq,sequence_in,funs,funs_is_function,function_apply_intro,funs_type_apply,subseteq_transitive}.
    Take $U\in C$ such that $J\inter\preimg{s}{U}$ is infinite
        by \cref{finite_cover_infinite_sequence_trace}.
    Take $x\in X$ such that $U = X\inter\metricBall{d}{X}{x}{\epsilon}$
        by \cref{totally_bounded_cover_member_ball}.
    Let $K = J\inter\preimg{s}{U}$.
    Show that $K\subseteq J$.
    \begin{subproof}
        Follows by \cref{inter_lower_left}.
    \end{subproof}
    Show that $K$ is infinite.
    \begin{subproof}
        Follows by \cref{inter_lower_left}.
    \end{subproof}
    Show that for all $m\in K$ we have $\apply{s}{m}\in U$.
    \begin{subproof}
        Follows by \cref{inter,subseteq,sequence_in,funs,preimg_iff,function_apply_iff,funs_is_function}.
    \end{subproof}
    Show that there exists $z\in X$ such that $U = X\inter\metricBall{d}{X}{z}{\epsilon}$.
    \begin{subproof}
        Trivial.
    \end{subproof}
    Follows by definition.
\end{proof}

% from §45 helper lemma 3: a sequence of infinite subsets of naturalsPlus has an increasing diagonal
\begin{lemma}\label{infinite_naturalsplus_subset_sequence_diagonal}
    Suppose $j$ is a sequence in $\pow{\naturalsPlus}$.
    Suppose for all $n\in\naturalsPlus$ we have $\apply{j}{n}$ is infinite.
    Then there exists $u$ such that
        $u$ is strictly increasing on $\naturalsPlus$ and
        for all $n\in\naturalsPlus$ we have $u(n)\in\apply{j}{n}$.
\end{lemma}
\begin{proof}
    Let $D = \{p\in\naturalsPlus\times\naturalsPlus \mid \snd{p}\in\apply{j}{\fst{p}}\}$.
    Take $k$ such that $k\in\apply{j}{1}$ by \cref{real_one_in_naturals,emptyset_is_finite,nonempty_inhabited,subseteq}.
    Let $a_0 = (1,k)$.
    We have $a_0\in D$ and $\fst{a_0} = 1$ by \cref{real_one_in_naturals,times_tuple_intro,fst_eq,snd_eq,sequence_in,funs_type_apply,powerset_elim,subseteq}.

    Let $R = \{w\in D\times D \mid
        \fst{\snd{w}} = \fst{\fst{w}} + 1 \land \snd{\fst{w}} < \snd{\snd{w}}\}$.

    Show that for all $p\in D$ there exists $q$ such that $p\mathrel{R} q$.
    \begin{subproof}
        Fix $p\in D$.
        Take $n_0,m_0$ such that $p = (n_0,m_0)$ by \cref{times_elem_is_tuple}.
        $n_0\in\naturalsPlus$ and $m_0\in\naturalsPlus$ by \cref{times_tuple_elim}.
        Hence $n_0 + 1\in\naturalsPlus$ by \cref{real_naturals_closed_succ}.
        Thus $\apply{j}{(n_0 + 1)}$ is infinite and $\apply{j}{(n_0 + 1)}\subseteq\naturalsPlus$
            by \cref{subseteq,sequence_in,funs_type_apply,powerset_elim}.
        Take $w$ such that
            $w$ is strictly increasing on $\naturalsPlus$ and
            for all $n\in\naturalsPlus$ we have $w(n)\in\apply{j}{(n_0 + 1)}$
            by \cref{infinite_naturalsplus_has_strictly_increasing_sequence}.
        Let $k_1 = \apply{w}{(m_0 + 1)}$.
        We have $k_1\in\apply{j}{(n_0 + 1)}$ and $k_1\in\naturalsPlus$ by \cref{real_naturals_closed_succ,subseteq}.
        $m_0 + 1 \leq k_1$ by \cref{real_naturals_closed_succ,strictly_increasing_naturalsplus_geq_identity}.
        Hence $m_0 < k_1$
            by \cref{real_naturals_lt_succ_local,real_naturals_subset_reals,real_naturals_closed_succ,subseteq,real_lt_leq_trans_local}.
        Let $q = (n_0 + 1,k_1)$.
        $q\in\naturalsPlus\times\naturalsPlus$ by \cref{times_tuple_intro,subseteq,sequence_in,funs_type_apply,powerset_elim}.
        $\fst{q} = n_0 + 1$ and $\snd{q} = k_1$ by \cref{fst_eq,snd_eq}.
        $\snd{q}\in\apply{j}{\fst{q}}$ by \cref{real_lt_subst_right}.
        Hence $(p,q)\in D\times D$ by \cref{times_tuple_intro,subseteq}.
        $\fst{\snd{(p,q)}} = n_0 + 1$ and $\fst{\fst{(p,q)}} = n_0$ by \cref{fst_eq,snd_eq}.
        $\snd{\fst{(p,q)}} = m_0$ and $\snd{\snd{(p,q)}} = k_1$ by \cref{fst_eq,snd_eq}.
        Thus $\fst{\snd{(p,q)}} = \fst{\fst{(p,q)}} + 1$ and
            $\snd{\fst{(p,q)}} < \snd{\snd{(p,q)}}$
            by \cref{real_lt_subst_left,real_lt_subst_right,real_add_eq_subst}.
        Hence there exists $r$ such that $p\mathrel{R} r$ by \cref{subseteq}.
    \end{subproof}

    Take $f$ such that $f$ is a function on $D$ and for all $p\in D$ we have
        $p\mathrel{R}\apply{f}{p}$ by \cref{subseteq,times_elem_is_tuple,relation,choice_function}.
    Hence $f\in\funs{D}{D}$ by \cref{funs_intro,subseteq,times_tuple_elim}.
    Take $v$ such that
        $v$ is a sequence in $D$ and
        $v(1) = a_0$ and
        for all $n\in\naturalsPlus$ we have $v(n + 1) = \apply{f}{v(n)}$
        by \cref{iteration_sequence_naturalsplus}.

    Let $g = \{(p,\snd{p}) \mid p\in D\}$.
    For all $p,y_1,y_2$ such that $(p,y_1)\in g$ and $(p,y_2)\in g$ we have $y_1 = y_2$.
    Therefore $g$ is a function and $\dom{g} = D$
        by \cref{rightunique,relation,dom_iff,pair_eq_iff,subseteq,subseteq_antisymmetric}.
    For all $p\in\dom{g}$ we have $g(p)\in\naturalsPlus$
        by \cref{pair_eq_iff,function_apply_intro,subseteq,times_elem_is_tuple,snd_eq}.
    Hence $g\in\funs{D}{\naturalsPlus}$ by \cref{funs_intro}.

    Let $u = g\circ v$.
    Hence $u$ is a sequence in $\naturalsPlus$ by \cref{sequence_in,funs_circ}.

    For all $n\in\naturalsPlus$ we have $u(n) = \snd{v(n)}$
        by \cref{funs,sequence_in,funs_is_function,function_apply_elim,funs_type_apply,circ_iff,function_apply_intro}.

    Let $P = \{n\in\naturalsPlus \mid \fst{v(n)} = n\}$.
    $P\subseteq\naturalsPlus$ by \cref{subseteq}.
    $1\in P$ by \cref{real_one_in_naturals,fst_eq,pair_eq_pair_of_proj}.
    For all $n\in P$ we have $n + 1\in P$
        by \cref{real_naturals_closed_succ,iteration_sequence_naturalsplus,sequence_in,funs_type_apply,fst_eq,snd_eq,real_add_eq_subst,subseteq}.
    Hence for all $n\in\naturalsPlus$ we have $n\in P$ by \cref{real_naturals_induction}.
    For all $n\in\naturalsPlus$ we have $u(n)\in\apply{j}{n}$
        by \cref{funs_type_apply,sequence_in,snd_eq,subseteq}.

    For all $n\in\naturalsPlus$ we have $v(n)\mathrel{R} v(n + 1)$
        by \cref{iteration_sequence_naturalsplus,sequence_in,funs_type_apply,function_apply_elim}.
    Therefore for all $n\in\naturalsPlus$ we have $u(n) < u(n + 1)$
        by \cref{real_naturals_closed_succ,snd_eq,real_lt_subst_left,real_lt_subst_right,fst_eq}.

    Let $T_0 = \{n\in\naturalsPlus \mid \text{for all $m\in\naturalsPlus$ such that $m < n$ we have $u(m) < u(n)$}\}$.
    $T_0\subseteq\naturalsPlus$ by \cref{subseteq}.
    $1\in T_0$ by \cref{real_one_in_naturals,real_one_leq_naturalsplus,real_naturals_subset_reals,subseteq,real_lt_subst_right,real_lt_trans,real_lt_irreflexive}.
    Show that for all $n\in T_0$ we have $n + 1\in T_0$.
    \begin{subproof}
        Fix $n$ such that $n\in T_0$.
        It suffices to show that for all $m\in\naturalsPlus$ such that $m < n + 1$ we have $u(m) < u(n + 1)$.
        Fix $m$ such that $m\in\naturalsPlus$ and $m < n + 1$.
        Suppose not.
        Thus $m < n$ or $m = n$ or $m = n + 1$ by \cref{real_succ_discrete}.
        \begin{byCase}
        \caseOf{$m = n$.}
            Contradiction by \cref{real_lt_subst_left}.
        \caseOf{$m < n$.}
        We have $u(m) < u(n)$ by \cref{subseteq}.
        Hence $u(m) < u(n + 1)$
            by \cref{sequence_in,real_naturals_closed_succ,real_naturals_subset_reals,subseteq,funs_type_apply,real_lt_trans}.
            Contradiction.
        \caseOf{$m = n + 1$.}
            Contradiction
                by \cref{real_lt_subst_left,real_naturals_subset_reals,subseteq,real_lt_irreflexive}.
        \end{byCase}
    \end{subproof}
    Hence for all $n\in\naturalsPlus$ we have $n\in T_0$ by \cref{real_naturals_induction}.
    Therefore for all $m,n\in\naturalsPlus$ such that $m < n$ we have $u(m) < u(n)$ by \cref{subseteq}.
    Hence $u$ is strictly increasing on $\naturalsPlus$ by \cref{sequence_in,strictly_increasing_on_set}.
    Follows by definition.
\end{proof}

% helper lemma: points close to a common center are pairwise close
\begin{lemma}\label{metric_common_center_double_bound}
    Assume $d$ is a metric on $X$.
    Suppose $x,y,z\in X$.
    Suppose $\alpha\in\reals$.
    Suppose $d((x,z)) < \alpha$.
    Suppose $d((y,z)) < \alpha$.
    Then $d((x,y)) < \alpha + \alpha$.
\end{lemma}
\begin{proof}
    We have $d((x,z))\in\reals$ and $d((y,z))\in\reals$ and
        $d((z,y))\in\reals$ and $d((x,y))\in\reals$
        by \cref{times_tuple_intro,metric_on,funs_type_apply}.
    Thus $d((z,y)) < \alpha$ by \cref{times_tuple_intro,metric_on,metric_symmetric,real_lt_subst_left}.
    Therefore $d((x,z)) + d((z,y)) < \alpha + \alpha$ by \cref{real_lt_add_two}.
    We have $d((x,y)) \leq d((x,z)) + d((z,y))$ by \cref{times_tuple_intro,metric_on,metric_triangle_inequality}.
    Hence $d((x,y)) < \alpha + \alpha$ by \cref{real_add_closed,real_leq_lt_trans}.
\end{proof}

\begin{lemma}\label{one_plus_one_reciprocal_positive}
    Suppose $m_2 = 1 + 1$.
    Then $m_2\in\reals$ and $0 < m_2$ and $m_2\neq 0$ and $\inv{m_2}\in\reals$ and $0 < \inv{m_2}$.
\end{lemma}
\begin{proof}
    $1 + 1\in\reals$ and $0 < 1 + 1$
        by \cref{real_naturals_subset_reals,real_one_in_naturals,subseteq,real_add_ident,real_add_closed,real_zero_lt_one,real_lt_add,real_lt_trans}.
    Show that $m_2\in\reals$.
    \begin{subproof}
        Follows by \cref{real_add_closed}.
    \end{subproof}
    Show that $0 < m_2$.
    \begin{subproof}
        Follows by \cref{real_lt_subst_right}.
    \end{subproof}
    Show that $m_2\neq 0$.
    \begin{subproof}
        Follows by \cref{real_pos_nonzero}.
    \end{subproof}
    Show that $\inv{m_2}\in\reals$.
    \begin{subproof}
        Follows by \cref{real_mul_inverse}.
    \end{subproof}
    Show that $0 < \inv{m_2}$.
    \begin{subproof}
        Follows by \cref{real_inv_pos}.
    \end{subproof}
\end{proof}

% from §45 Item 2: complete and totally bounded metric spaces are compact
\begin{theorem}\label{complete_totally_bounded_metric_space_compact}
    Assume $d$ is a metric on $X$.
    Assume $X$ is complete with respect to $d$.
    Assume $X$ is totally bounded with respect to $d$.
    Let $T = \metricTopology{d}{X}$.
    Then $X$ is compact with respect to $T$.
\end{theorem}
\begin{proof}
    $T$ is a topology on $X$ by \cref{metric_topology,metric_basis_is_basis,basis_topology_is_topology,basis_for_topology}.
    Show that for all $s$ such that $s$ is a sequence in $X$
        there exist $t,x$ such that
            $t$ is a sequence in $X$ and
            $t$ is a subsequence of $s$ and
            $x\in X$ and
            $t$ converges to $x$ in $X$ with respect to $T$.
    \begin{subproof}
            Fix $s$ such that $s$ is a sequence in $X$.
            Take $J_1,U_1$ such that
                $J_1\subseteq\naturalsPlus$ and
                $J_1$ is infinite and
                for all $m\in J_1$ we have $\apply{s}{m}\in U_1$ and
                there exists $x\in X$ such that
                    $U_1 = X\inter\metricBall{d}{X}{x}{1 / 1}$
                by \cref{subseteq_refl,naturalsplus_infinite,real_one_in_naturals,naturalsplus_reciprocal_real,positive_reciprocal_naturalsplus,totally_bounded_infinite_sequence_epsilon_ball}.

            Let $A = \{p\in\naturalsPlus\times\pow{\naturalsPlus} \mid \text{$\snd{p}$ is infinite}\}$.
            Let $a_0 = (1,J_1)$.
            We have $a_0\in A$ and $\fst{a_0} = 1$
                by \cref{real_one_in_naturals,times_tuple_intro,powerset_intro,snd_eq,fst_eq,subseteq}.

            Let $R = \{w\in A\times A \mid \text{$\fst{\snd{w}} = \fst{\fst{w}} + 1$ and
                $\snd{\snd{w}}\subseteq\snd{\fst{w}}$ and
                there exists $U$ such that
                    for all $m\in\snd{\snd{w}}$ we have $\apply{s}{m}\in U$ and
                    there exists $x\in X$ such that
                        $U = X\inter\metricBall{d}{X}{x}{1 / (\fst{\snd{w}})}$}\}$.

            Show that for all $p\in A$ there exists $q$ such that $p\mathrel{R} q$.
            \begin{subproof}
                Fix $p\in A$.
                Take $n_0,J$ such that $p = (n_0,J)$ by \cref{times_elem_is_tuple}.
                $n_0\in\naturalsPlus$ and $J\in\pow{\naturalsPlus}$ by \cref{times_tuple_elim}.
                Thus $J$ is infinite and $J\subseteq\naturalsPlus$ by \cref{powerset_elim,fst_eq,snd_eq,subseteq}.
                Take $K,U$ such that
                    $K\subseteq J$ and
                    $K$ is infinite and
                    for all $m\in K$ we have $\apply{s}{m}\in U$ and
                    there exists $x\in X$ such that
                        $U = X\inter\metricBall{d}{X}{x}{1 / (n_0 + 1)}$
                    by \cref{real_naturals_closed_succ,naturalsplus_reciprocal_real,positive_reciprocal_naturalsplus,totally_bounded_infinite_sequence_epsilon_ball}.
                Let $q = (n_0 + 1,K)$.
                Hence $(p,q)\in A\times A$ by \cref{times_tuple_intro,powerset_intro,snd_eq,subseteq,real_naturals_closed_succ}.
                $\fst{\snd{(p,q)}} = n_0 + 1$ and $\fst{\fst{(p,q)}} = n_0$
                    by \cref{fst_eq,snd_eq}.
                $\snd{\snd{(p,q)}} = K$ and $\snd{\fst{(p,q)}} = J$
                    by \cref{fst_eq,snd_eq}.
                Thus $\fst{\snd{(p,q)}} = \fst{\fst{(p,q)}} + 1$ and
                    $\snd{\snd{(p,q)}}\subseteq\snd{\fst{(p,q)}}$
                    by \cref{real_add_eq_subst}.
                Hence there exists $r$ such that $p\mathrel{R} r$
                    by \cref{real_add_eq_subst,real_lt_subst_left,real_lt_subst_right}.
            \end{subproof}

            Take $f$ such that $f$ is a function on $A$ and for all $p\in A$ we have
                $p\mathrel{R}\apply{f}{p}$ by \cref{subseteq,times_elem_is_tuple,relation,choice_function}.
            Hence $f\in\funs{A}{A}$ by \cref{funs_intro,subseteq,times_tuple_elim}.
            Take $v$ such that
                $v$ is a sequence in $A$ and
                $v(1) = a_0$ and
                for all $n\in\naturalsPlus$ we have $v(n + 1) = \apply{f}{v(n)}$
                by \cref{iteration_sequence_naturalsplus}.

            Let $g = \{(p,\snd{p}) \mid p\in A\}$.
            For all $p,y_1,y_2$ such that $(p,y_1)\in g$ and $(p,y_2)\in g$ we have $y_1 = y_2$.
            Therefore $g$ is a function and $\dom{g} = A$
                by \cref{rightunique,relation,dom_iff,pair_eq_iff,subseteq,subseteq_antisymmetric}.
            For all $p\in\dom{g}$ we have $g(p)\in\pow{\naturalsPlus}$
                by \cref{pair_eq_iff,function_apply_intro,subseteq,times_elem_is_tuple,snd_eq}.
            Hence $g\in\funs{A}{\pow{\naturalsPlus}}$ by \cref{funs_intro}.

            Let $j = g\circ v$.
            We have $j$ is a sequence in $\pow{\naturalsPlus}$ by \cref{sequence_in,funs_circ}.
            For all $n\in\naturalsPlus$ we have $j(n) = \snd{v(n)}$
                by \cref{funs,sequence_in,funs_is_function,function_apply_elim,funs_type_apply,circ_iff,function_apply_intro}.

            Let $P = \{n\in\naturalsPlus \mid \fst{v(n)} = n\}$.
            $P\subseteq\naturalsPlus$ by \cref{subseteq}.
            Show that $1\in P$.
            \begin{subproof}
                Follows by \cref{real_one_in_naturals,fst_eq,pair_eq_pair_of_proj}.
            \end{subproof}
            Show that for all $n\in P$ we have $n + 1\in P$.
            \begin{subproof}
                Follows by \cref{real_naturals_closed_succ,iteration_sequence_naturalsplus,sequence_in,funs_type_apply,subseteq,fst_eq,snd_eq,real_add_eq_subst}.
            \end{subproof}
            Hence for all $n\in\naturalsPlus$ we have $n\in P$ by \cref{real_naturals_induction}.

            For all $n\in\naturalsPlus$ we have $\apply{j}{n}$ is infinite
                by \cref{funs_type_apply,sequence_in,snd_eq,subseteq}.

            Take $u$ such that
                $u$ is strictly increasing on $\naturalsPlus$ and
                for all $n\in\naturalsPlus$ we have $u(n)\in\apply{j}{n}$
                by \cref{infinite_naturalsplus_subset_sequence_diagonal}.
            Hence $u$ is a sequence in $\naturalsPlus$ by \cref{sequence_in,strictly_increasing_on_set}.

            Let $t = s\circ u$.
            Thus $t$ is a sequence in $X$ by \cref{sequence_in,funs_circ}.
            For all $n\in\naturalsPlus$ we have $t(n) = \apply{s}{\apply{u}{n}}$
                by \cref{funs,sequence_in,funs_is_function,function_apply_elim,funs_type_apply,circ_iff,function_apply_intro}.
            Hence $t$ is a subsequence of $s$ by \cref{subsequence}.

            For all $n\in\naturalsPlus$ we have $v(n)\mathrel{R} v(n + 1)$
                by \cref{iteration_sequence_naturalsplus,sequence_in,funs_type_apply,function_apply_elim}.
            Hence for all $n\in\naturalsPlus$ we have $j(n + 1)\subseteq j(n)$
                by \cref{real_naturals_closed_succ,snd_eq,subseteq,fst_eq}.

            Let $A_0 = \{m\in\naturalsPlus \mid \text{for all $n\in\naturalsPlus$ if $n \leq m$ then
                $j(m)\subseteq j(n)$}\}$.
            $A_0\subseteq\naturalsPlus$ by \cref{subseteq}.
            Show that $1\in A_0$.
            \begin{subproof}
                Follows by \cref{real_one_in_naturals,real_naturals_subset_reals,subseteq,real_one_leq_naturalsplus,real_leq_antisym,subseteq_refl}.
            \end{subproof}
            Show that for all $k\in A_0$ we have $k + 1\in A_0$.
            \begin{subproof}
                Follows by \cref{subseteq,real_naturals_closed_succ,real_naturals_leq_succ_elim,subseteq_refl,subseteq_transitive}.
            \end{subproof}
            Hence for all $m\in\naturalsPlus$ we have $m\in A_0$ by \cref{real_naturals_induction}.

            Let $Q = \{n\in\naturalsPlus \mid \text{there exists $U$ such that
                for all $m\in j(n)$ we have $\apply{s}{m}\in U$ and
                there exists $x\in X$ such that
                    $U = X\inter\metricBall{d}{X}{x}{1 / n}$}\}$.
            $Q\subseteq\naturalsPlus$ by \cref{subseteq}.
            For all $m\in j(1)$ we have $\apply{s}{m}\in U_1$
                by \cref{real_one_in_naturals,snd_eq,pair_eq_pair_of_proj,subseteq}.
            Take $x_1\in X$ such that $U_1 = X\inter\metricBall{d}{X}{x_1}{1 / 1}$.
            Show that $1\in Q$.
            \begin{subproof}
                Follows by \cref{real_one_in_naturals,subseteq}.
            \end{subproof}
            For all $n\in Q$ there exists $U$ such that
                for all $m\in j(n + 1)$ we have $\apply{s}{m}\in U$ and
                there exists $x\in X$ such that
                    $U = X\inter\metricBall{d}{X}{x}{1 / (\fst{v(n + 1)})}$
                by \cref{iteration_sequence_naturalsplus,sequence_in,funs_type_apply,function_apply_elim,fst_eq,snd_eq,subseteq}.
            Show that for all $n\in Q$ we have $n + 1\in Q$.
            \begin{subproof}
                Follows by \cref{real_naturals_closed_succ,subseteq,pair_eq_iff,fst_eq,snd_eq}.
            \end{subproof}
            Hence for all $n\in\naturalsPlus$ we have $n\in Q$ by \cref{real_naturals_induction}.

            Show that for all $\epsilon\in\reals$ such that $0 < \epsilon$
                there exists $N\in\naturalsPlus$ such that
                    for all $n,m\in\naturalsPlus$ if $N \leq n$ and $N \leq m$ then
                        $d((t(n),t(m))) < \epsilon$.
            \begin{subproof}
                    Fix $\epsilon$ such that $\epsilon\in\reals$ and $0 < \epsilon$.
                    Let $m_2 = 1 + 1$.
                    Let $\delta = \epsilon / m_2$.
                    We have $m_2\in\reals$ and $0 < m_2$ and $m_2\neq 0$ and $\inv{m_2}\in\reals$ and
                        $0 < \inv{m_2}$
                        by \cref{one_plus_one_reciprocal_positive}.
                    $0\in\reals$ by \cref{real_add_ident,real_one_in_naturals,real_naturals_subset_reals,subseteq}.
                    $\delta = \epsilon \rmul \inv{m_2}$ by \cref{real_div}.
                    Hence $\delta\in\reals$ and $0 < \delta$
                        by \cref{real_mul_closed,real_lt_mul_pos,real_mul_zero,real_lt_subst_left}.
                    Take $N\in\naturalsPlus$ such that $0 < 1 / N$ and $1 / N < \delta$ by \cref{real_small_reciprocal}.
                    Take $U$ such that
                        for all $m\in j(N)$ we have $\apply{s}{m}\in U$ and
                        there exists $x\in X$ such that
                            $U = X\inter\metricBall{d}{X}{x}{1 / N}$ by \cref{subseteq}.
                    Take $x\in X$ such that $U = X\inter\metricBall{d}{X}{x}{1 / N}$.
                    It suffices to show that for all $n\in\naturalsPlus$ we have for all $m\in\naturalsPlus$ if $N \leq n$ and $N \leq m$ then
                        $d((t(n),t(m))) < \epsilon$.
                    Fix $n\in\naturalsPlus$.
                    Fix $m\in\naturalsPlus$.
                    Assume $N \leq n$ and $N \leq m$.
                    $u(n)\in j(n)$ and $u(m)\in j(m)$ by \cref{subseteq}.
                    Hence $j(n)\subseteq j(N)$ and $j(m)\subseteq j(N)$ by \cref{subseteq}.
                    Thus $u(n)\in j(N)$ and $u(m)\in j(N)$ by \cref{subseteq}.
                    Therefore $t(n)\in U$ and $t(m)\in U$ by \cref{subseteq}.
                    $t(n)\in X\inter\metricBall{d}{X}{x}{1 / N}$ and
                        $t(m)\in X\inter\metricBall{d}{X}{x}{1 / N}$.
                    $t(n)\in X$ and $t(m)\in X$ by \cref{inter}.
                    Thus $t(n)\in\metricBall{d}{X}{x}{1 / N}$ and
                        $t(m)\in\metricBall{d}{X}{x}{1 / N}$ by \cref{inter}.
                    Hence $d((t(n),x)) < 1 / N$ and $d((t(m),x)) < 1 / N$
                        by \cref{metric_ball,metric_on,metric_symmetric,real_lt_subst_left}.
                    $1 / N\in\reals$ by \cref{naturalsplus_reciprocal_real}.
                    Thus $d((t(n),t(m))) < 1 / N + 1 / N$
                        by \cref{metric_common_center_double_bound}.
                    $d((t(n),t(m)))\in\reals$ by \cref{times_tuple_intro,metric_on,funs_type_apply}.
                    We have $1 / N + 1 / N\in\reals$ and
                        $\delta + \delta\in\reals$ by \cref{real_add_closed}.
                    Therefore $1 / N + 1 / N < \delta + \delta$ by \cref{real_lt_add_two}.
                    Hence $d((t(n),t(m))) < \delta + \delta$ by \cref{real_lt_trans}.
                    Thus $\delta + \delta = \epsilon$ by \cref{real_half_plus_half,real_div,real_add_eq_subst}.
                    Therefore $d((t(n),t(m))) < \epsilon$ by \cref{real_lt_subst_right}.
                    Follows by definition.
            \end{subproof}
            Hence $t$ is a Cauchy sequence in $X$ with respect to $d$ by \cref{cauchy_sequence}.

            Take $x\in X$ such that $t$ converges to $x$ in $X$ with respect to $T$
                by \cref{complete_metric_space}.
            Follows by definition.
    \end{subproof}
    Hence $X$ is sequentially compact with respect to $T$ by \cref{sequentially_compact}.
    Show that for all $A$ if $A$ is an open covering of $X$ with respect to $T$ then
        there exists $B$ such that $B\subseteq A$ and $B$ is finite and $B$ is a covering of $X$.
    \begin{subproof}
        Fix $A$ such that $A$ is an open covering of $X$ with respect to $T$.
        \begin{byCase}
        \caseOf{$X = \emptyset$.}
            There exists $C$ such that $C\subseteq A$ and $C$ is finite and $C$ is a covering of $X$
                by \cref{emptyset_subseteq,emptyset_is_finite,unions_emptyset,covering}.
        \caseOf{$X\neq\emptyset$.}
            Take $\delta\in\reals$ such that $0 < \delta$ and
                for all $B$ such that $B\subseteq X$ and
                    for all $x,y\in B$ we have $d((x,y)) < \delta$
                there exists $U\in A$ such that $B\subseteq U$
                by \cref{sequentially_compact_lebesgue}.
            Let $m_2 = 1 + 1$.
            Let $\epsilon = \delta / m_2$.
            We have $m_2\in\reals$ and $0 < m_2$ and $m_2\neq 0$ and $\inv{m_2}\in\reals$ and
                $0 < \inv{m_2}$
                by \cref{one_plus_one_reciprocal_positive}.
            $\epsilon = \delta \rmul \inv{m_2}$ by \cref{real_div}.
            Hence $\epsilon\in\reals$ and $0 < \epsilon$
                by \cref{real_mul_closed,real_add_ident,real_lt_mul_pos,real_mul_zero,real_lt_subst_left}.
            Take $B$ such that
                $B$ is finite and
                $B$ is a covering of $X$ and
                for all $U\in B$ there exists $x\in X$ such that $U = \metricBall{d}{X}{x}{\epsilon}$
                by \cref{sequentially_compact_finite_ball_cover}.

            We have for all $V\in B$ there exists $U\in A$ such that $V\subseteq U$
                by \cref{metric_ball,subseteq,metric_on,metric_symmetric,real_lt_subst_left,times_tuple_intro,funs_type_apply,metric_common_center_double_bound,real_half_plus_half,real_div,real_add_eq_subst,real_lt_subst_right}.

            Let $R = \{w\in B\times A \mid \fst{w}\subseteq \snd{w}\}$.
            For all $V\in B$ there exists $U$ such that $V\mathrel{R} U$
                by \cref{subseteq,times_tuple_intro,fst_eq,snd_eq}.
            Take $h$ such that $h$ is a function on $B$ and for all $V\in B$ we have $V\mathrel{R}\apply{h}{V}$
                by \cref{subseteq,times_elem_is_tuple,relation,choice_function}.
            Let $C = \img{h}{B}$.
            Hence $C$ is finite by \cref{finite_image_function}.
            Therefore $C\subseteq A$ by \cref{img_iff,function_apply_intro,subseteq,times_tuple_elim}.
            For all $x$ such that $x\in X$ we have $x\in\unions{C}$
                by \cref{covering,subseteq,unions_iff,fst_eq,snd_eq,function_apply_elim,img_iff}.
            Hence $C$ is a covering of $X$ by \cref{subseteq,covering}.
            Follows by definition.
        \end{byCase}
    \end{subproof}
    Hence $X$ is compact with respect to $T$ by \cref{compact}.
\end{proof}

% from §45 Item 3: equicontinuous at a point
\begin{definition}\label{equicontinuous_at_point}
    $F$ is equicontinuous at $x_0$ from $X$ to $Y$ with respect to $T$ and $d$ iff
        $T$ is a topology on $X$ and
        $x_0\in X$ and
        $d$ is a metric on $Y$ and
        $F\subseteq\funs{X}{Y}$ and
        for all $\epsilon\in\reals$ such that $0 < \epsilon$
        there exists $U$ such that
            $U$ is a neighborhood of $x_0$ in $X$ with respect to $T$ and
            for all $x, f$ such that $x\in U$ and $f\in F$ we have
                $d((\apply{f}{x},\apply{f}{x_0})) < \epsilon$.
\end{definition}

% from §45 Item 3: equicontinuous family
\begin{definition}\label{equicontinuous_family}
    $F$ is equicontinuous from $X$ to $Y$ with respect to $T$ and $d$ iff
        $T$ is a topology on $X$ and
        $d$ is a metric on $Y$ and
        $F\subseteq\funs{X}{Y}$ and
        for all $x\in X$ we have
            $F$ is equicontinuous at $x$ from $X$ to $Y$ with respect to $T$ and $d$.
\end{definition}

% from §45 helper lemma 4: finite intersections of neighborhoods are neighborhoods
\begin{lemma}\label{finite_neighborhood_intersection}
    Assume $T$ is a topology on $X$.
    Suppose $A\subseteq T$.
    Suppose $A$ is finite.
    Suppose $A$ is inhabited.
    Suppose for all $U\in A$ we have $x\in U$.
    Then $\inters{A}$ is a neighborhood of $x$ in $X$ with respect to $T$.
\end{lemma}
\begin{proof}
    Follows by \cref{finite_intersections_open_in_topology,open_in,inters_iff_forall,neighborhood}.
\end{proof}

% from §45 helper lemma 5: small uniform distance gives pointwise control
\begin{lemma}\label{uniform_metric_small_implies_pointwise_small}
    Assume $d$ is a metric on $Y$.
    Assume $\rho$ is the uniform metric on functions from $X$ to $Y$ with respect to $d$.
    Suppose $f,g\in\funs{X}{Y}$.
    Suppose $\epsilon\in\reals$.
    Suppose $0 < \epsilon$.
    Suppose $\epsilon < 1$.
    Suppose $\rho((f,g)) < \epsilon$.
    Then for all $x\in X$ we have
        $d((f(x),g(x))) < \epsilon$.
\end{lemma}
\begin{proof}
    Let $e = \boundedMetric{d}$.
    We have $e$ is a metric on $Y$ and $\rho((f,g))\in\reals$
        by \cref{bounded_metric_is_metric,uniform_metric_on_function_space,times_tuple_intro,metric_on,funs_type_apply}.
    For all $x\in X$ we have $e((f(x),g(x))) \leq \rho((f,g))$ and $e((f(x),g(x)))\in\reals$
        by \cref{uniform_function_metric_coordinate_leq,times_tuple_intro,funs_type_apply,metric_on}.
    Show that for all $x\in X$ we have $e((f(x),g(x))) < \epsilon$.
    \begin{subproof}
        Follows by \cref{real_leq_lt_trans}.
    \end{subproof}
    Show that for all $x\in X$ we have $e((f(x),g(x))) = d((f(x),g(x)))$.
    \begin{subproof}
        Follows by \cref{funs_type_apply,times_tuple_intro,bounded_metric_apply_value_lt_one,real_naturals_subset_reals,real_one_in_naturals,subseteq,real_lt_trans}.
    \end{subproof}
    Show that for all $x\in X$ we have $d((f(x),g(x))) < \epsilon$.
    \begin{subproof}
        Follows by \cref{real_lt_subst_left}.
    \end{subproof}
\end{proof}

% from §45 helper lemma 6: totally bounded subsets admit finite internal epsilon-nets
\begin{lemma}\label{totally_bounded_subset_finite_internal_net}
    Assume $\rho$ is a metric on $Z$.
    Suppose $F\subseteq Z$.
    Suppose $F$ is totally bounded in $Z$ with respect to $\rho$.
    Suppose $F$ is inhabited.
    Suppose $\epsilon\in\reals$.
    Suppose $0 < \epsilon$.
    Then there exists $H$ such that
        $H\subseteq F$ and
        $H$ is finite and
        for all $f\in F$ there exists $h\in H$ such that
            $\rho((f,h)) < \epsilon$.
\end{lemma}
\begin{proof}
    Let $m_2 = 1 + 1$.
    We have $m_2\in\reals$ and $0 < m_2$ and $m_2\neq 0$ and $\inv{m_2}\in\reals$ and
        $0 < \inv{m_2}$
        by \cref{one_plus_one_reciprocal_positive}.
    $0\in\reals$ by \cref{real_add_ident,real_one_in_naturals,real_naturals_subset_reals,subseteq}.
    Let $\delta = \epsilon / m_2$.
    $\delta = \epsilon \rmul \inv{m_2}$ by \cref{real_div}.
    Hence $\delta\in\reals$ by \cref{real_mul_closed}.
    Therefore $0 < \delta$ by \cref{real_lt_mul_pos,real_mul_zero,real_lt_subst_left}.
    Take $C$ such that
        $C$ is finite and
        $C$ is a covering of $F$ and
        for all $U\in C$ there exists $z\in Z$ such that
            $U = F\inter\metricBall{\rho}{Z}{z}{\delta}$
        by \cref{totally_bounded_subset_metric_space}.
    Let $D = \{U\in C \mid \exists f. f\in U\}$.
    $D$ is finite by \cref{subseteq,subseteq_finite_is_finite}.
    Take $f_0$ such that $f_0\in F$.
    $f_0\in\unions{C}$ by \cref{covering,subseteq}.
    Take $U_0\in C$ such that $f_0\in U_0$ by \cref{unions_iff}.
    Hence $U_0\in D$ by \cref{subseteq}.
    Let $R = \{w\in D\times F \mid \snd{w}\in\fst{w}\}$.
    For all $U\in D$ there exists $f$ such that $U\mathrel{R} f$
        by \cref{subseteq,inter_elim_left,real_lt_subst_left,times_tuple_intro,fst_eq,snd_eq}.
    Take $g$ such that
        $g$ is a function on $D$ and
        for all $U\in D$ we have $U\mathrel{R}\apply{g}{U}$
        by \cref{subseteq,times_elem_is_tuple,relation,choice_function}.
    Let $H = \img{g}{D}$.
    Show that $H$ is finite.
    \begin{subproof}
        Follows by \cref{finite_image_function}.
    \end{subproof}
    Show that $H\subseteq F$.
    \begin{subproof}
        Follows by \cref{img_iff,function_apply_intro,function_apply_elim,times_tuple_elim,subseteq}.
    \end{subproof}
    Show that for all $f\in F$ there exists $h\in H$ such that
        $\rho((f,h)) < \epsilon$.
    \begin{subproof}
        Fix $f$ such that $f\in F$.
        Take $U\in C$ such that $f\in U$ by \cref{covering,subseteq,unions_iff}.
        Therefore $U\in D$ by \cref{subseteq}.
        Let $h = g(U)$.
        We have $h\in H$ and $U\mathrel{R} h$ by \cref{img_iff,function_apply_elim}.
        Thus $h\in U$ by \cref{snd_eq,fst_eq}.
        Take $z\in Z$ such that $U = F\inter\metricBall{\rho}{Z}{z}{\delta}$ by \cref{subseteq}.
        Hence $f\in F\inter\metricBall{\rho}{Z}{z}{\delta}$ and
            $h\in F\inter\metricBall{\rho}{Z}{z}{\delta}$ by \cref{real_lt_subst_left,real_lt_subst_right}.
        $f\in Z$ and $h\in Z$ by \cref{subseteq,inter_elim_left}.
        Therefore $f\in\metricBall{\rho}{Z}{z}{\delta}$ and
            $h\in\metricBall{\rho}{Z}{z}{\delta}$ by \cref{inter_elim_right}.
        Hence $\rho((z,f)) < \delta$ and $\rho((z,h)) < \delta$ by \cref{metric_ball}.
        Thus $\rho((f,z)) < \delta$ and $\rho((h,z)) < \delta$
            by \cref{metric_on,metric_symmetric,real_lt_subst_left}.
        Therefore $\rho((f,h)) < \epsilon$
            by \cref{metric_common_center_double_bound,real_half_plus_half,real_div,real_add_eq_subst,real_lt_subst_right}.
        Follows by definition.
    \end{subproof}
\end{proof}

% from §45 helper lemma 7: finite continuous families admit a common epsilon-neighborhood
\begin{lemma}\label{finite_continuous_family_common_neighborhood}
    Assume $T$ is a topology on $X$.
    Assume $d$ is a metric on $Y$.
    Assume $C$ is the set of continuous functions from $X$ to $Y$ with respect to $T$ and $\metricTopology{d}{Y}$.
    Suppose $H\subseteq C$.
    Suppose $H$ is finite.
    Suppose $x_0\in X$.
    Suppose $\epsilon\in\reals$.
    Suppose $0 < \epsilon$.
    Then there exists $U$ such that
        $U$ is a neighborhood of $x_0$ in $X$ with respect to $T$ and
        for all $x, h$ such that $x\in U$ and $h\in H$ we have
            $d((h(x),h(x_0))) < \epsilon$.
\end{lemma}
\begin{proof}
    \begin{byCase}
    \caseOf{$H = \emptyset$.}
            Show that there exists $U$ such that
                $U$ is a neighborhood of $x_0$ in $X$ with respect to $T$ and
                for all $x, h$ such that $x\in U$ and $h\in H$ we have
                    $d((h(x),h(x_0))) < \epsilon$.
            \begin{subproof}
                Follows by \cref{topology_on,open_in,neighborhood,real_lt_subst_right,emptyset}.
            \end{subproof}
    \caseOf{$H\neq\emptyset$.}
        $H$ is inhabited by \cref{nonempty_inhabited}.
        Let $R = \{w\in H\times T \mid \text{$\snd{w}$ is a neighborhood of $x_0$ in $X$ with respect to $T$ and
            for all $x\in\snd{w}$ we have
                $d((\apply{\fst{w}}{x},\apply{\fst{w}}{x_0})) < \epsilon$}\}$.
        Show that for all $h\in H$ there exists $U$ such that $h\mathrel{R} U$.
        \begin{subproof}
            Fix $h$ such that $h\in H$.
            We have $h$ is continuous from $X$ to $Y$ with respect to $T$ and $\metricTopology{d}{Y}$
                by \cref{subseteq,continuous_function_set}.
            $h(x_0)\in Y$ by \cref{continuous,funs_type_apply}.
            Therefore $h$ is continuous at $x_0$ from $X$ to $Y$ with respect to $T$ and $\metricTopology{d}{Y}$
                by \cref{continuous_implies_continuous_at}.
            Let $V = \metricBall{d}{Y}{h(x_0)}{\epsilon}$.
            Hence $V$ is a neighborhood of $h(x_0)$ in $Y$ with respect to $\metricTopology{d}{Y}$
                by \cref{metric_ball,subseteq,powerset_intro,metric_basis,metric_basis_is_basis,basis_elem_open_in_generated,metric_topology,open_in,basis_topology_is_topology,basis_for_topology,metric_zero_characterization,real_lt_subst_left,metric_on,neighborhood}.
            Take $U$ such that
                $U$ is a neighborhood of $x_0$ in $X$ with respect to $T$ and
                $\img{h}{U}\subseteq V$
                by \cref{continuous_at}.
            Thus for all $x\in U$ we have $d((h(x),h(x_0))) < \epsilon$
                by \cref{neighborhood,open_in,topology_on,powerset_elim,subseteq,continuous,funs,inter_intro,funs_is_function,img_of_function_intro,metric_ball,metric_on,metric_symmetric,real_lt_subst_left,funs_type_apply}.
            Hence $(h,U)\in H\times T$ by \cref{neighborhood,open_in,times_tuple_intro}.
            Thus there exists $W_0$ such that $h\mathrel{R} W_0$ by \cref{subseteq,fst_eq,snd_eq}.
        \end{subproof}
        Take $u$ such that
            $u$ is a function on $H$ and
            for all $h\in H$ we have $h\mathrel{R}\apply{u}{h}$
            by \cref{subseteq,times_elem_is_tuple,relation,choice_function}.
        Let $A = \img{u}{H}$.
        Hence $A$ is finite and $A$ is inhabited
            by \cref{finite_image_function,function_apply_elim,img_iff,nonempty_inhabited}.
        Let $W = \inters{A}$.
        Show that $W$ is a neighborhood of $x_0$ in $X$ with respect to $T$.
        \begin{subproof}
            Follows by \cref{img_iff,function_apply_intro,function_apply_elim,subseteq,times_tuple_elim,snd_eq,neighborhood,finite_neighborhood_intersection}.
        \end{subproof}
        Show that for all $x, h$ such that $x\in W$ and $h\in H$ we have $d((h(x),h(x_0))) < \epsilon$.
        \begin{subproof}
            Follows by \cref{function_apply_elim,img_iff,inters_iff_forall,subseteq,fst_eq,snd_eq}.
        \end{subproof}
        Follows by definition.
    \end{byCase}
\end{proof}

% from §45 helper lemma 8: every positive real admits a small triple-sum bound below it
\begin{lemma}\label{positive_real_small_triple_bound}
    Suppose $\epsilon\in\reals$.
    Suppose $0 < \epsilon$.
    Then there exists $\delta\in\reals$ such that
        $0 < \delta$ and
        $\delta < 1$ and
        $\delta + \delta + \delta < \epsilon$.
\end{lemma}
\begin{proof}
    $1\in\reals$ by \cref{real_naturals_subset_reals,real_one_in_naturals,subseteq}.
    $1 + 1\in\reals$ and $0 < 1 + 1$ and $1 + 1\neq 0$ and $\inv{(1 + 1)}\in\reals$ and
        $0 < \inv{(1 + 1)}$ by \cref{one_plus_one_reciprocal_positive}.
    $0 < 1$ by \cref{real_zero_lt_one}.
    Hence $\epsilon < 1$ or $\epsilon = 1$ or $1 < \epsilon$ by \cref{real_lt_trichotomy}.
    \begin{byCase}
    \caseOf{$\epsilon < 1$.}
        Let $\eta = \epsilon / (1 + 1)$.
        $\eta\in\reals$ and $0 < \eta$
            by \cref{real_div,real_mul_inverse,real_mul_closed,real_add_ident,real_lt_mul_pos,real_inv_pos,real_mul_zero,real_lt_subst_left}.
        We have $\eta + \eta = \epsilon$ by \cref{real_half_plus_half}.
        Hence $\eta < 1$ by \cref{real_half_plus_half,real_add_ident,real_lt_add,real_lt_subst_left,real_lt_subst_right,real_lt_trans}.
        Let $\delta = \eta / (1 + 1)$.
        $\delta\in\reals$ and $0 < \delta$
            by \cref{real_div,real_mul_inverse,real_mul_closed,real_add_ident,real_lt_mul_pos,real_inv_pos,real_mul_zero,real_lt_subst_left}.
        We have $\delta + \delta = \eta$ by \cref{real_half_plus_half}.
        Hence $\delta < \eta$ by \cref{real_add_ident,real_lt_add,real_lt_subst_left,real_lt_subst_right}.
        Therefore $\delta < 1$ by \cref{real_half_plus_half,real_add_ident,real_lt_add,real_lt_subst_left,real_lt_subst_right,real_lt_trans}.
        $\delta + \eta < \eta + \eta$ by \cref{real_lt_add}.
        Thus $\delta + \delta + \delta < \epsilon$
            by \cref{real_add_comm,real_add_assoc,real_half_plus_half,real_lt_subst_left,real_lt_subst_right}.
        Follows by definition.
    \caseOf{$\epsilon = 1$.}
        Let $\eta = \epsilon / (1 + 1)$.
        $\eta\in\reals$ and $0 < \eta$
            by \cref{real_div,real_mul_inverse,real_mul_closed,real_add_ident,real_lt_mul_pos,real_inv_pos,real_mul_zero,real_lt_subst_left}.
        We have $\eta + \eta = \epsilon$ by \cref{real_half_plus_half}.
        Hence $\eta < 1$ by \cref{real_half_plus_half,real_add_ident,real_lt_add,real_lt_subst_left,real_lt_subst_right}.
        Let $\delta = \eta / (1 + 1)$.
        $\delta\in\reals$ and $0 < \delta$
            by \cref{real_div,real_mul_inverse,real_mul_closed,real_add_ident,real_lt_mul_pos,real_inv_pos,real_mul_zero,real_lt_subst_left}.
        We have $\delta + \delta = \eta$ by \cref{real_half_plus_half}.
        Hence $\delta < \eta$ by \cref{real_add_ident,real_lt_add,real_lt_subst_left,real_lt_subst_right}.
        Therefore $\delta < 1$ by \cref{real_half_plus_half,real_add_ident,real_lt_add,real_lt_subst_left,real_lt_subst_right,real_lt_trans}.
        $\delta + \eta < \eta + \eta$ by \cref{real_lt_add}.
        Thus $\delta + \delta + \delta < \epsilon$
            by \cref{real_add_comm,real_add_assoc,real_half_plus_half,real_lt_subst_left,real_lt_subst_right}.
        Follows by definition.
    \caseOf{$1 < \epsilon$.}
        Let $\eta = 1 / (1 + 1)$.
        $\eta\in\reals$ and $0 < \eta$
            by \cref{real_div,real_mul_inverse,real_mul_closed,real_add_ident,real_lt_mul_pos,real_inv_pos,real_mul_zero,real_lt_subst_left}.
        We have $\eta + \eta = 1$ by \cref{real_half_plus_half}.
        Hence $\eta < 1$ by \cref{real_add_ident,real_lt_add,real_lt_subst_left}.
        Let $\delta = \eta / (1 + 1)$.
        $\delta\in\reals$ and $0 < \delta$
            by \cref{real_div,real_mul_inverse,real_mul_closed,real_add_ident,real_lt_mul_pos,real_inv_pos,real_mul_zero,real_lt_subst_left}.
        We have $\delta + \delta = \eta$ by \cref{real_half_plus_half}.
        Hence $\delta < \eta$ by \cref{real_add_ident,real_lt_add,real_lt_subst_left,real_lt_subst_right}.
        Therefore $\delta < 1$ by \cref{real_lt_trans}.
        $\delta + \eta < \eta + \eta$ by \cref{real_lt_add}.
        Hence $\eta + \delta < \epsilon$
            by \cref{real_add_closed,real_add_comm,real_half_plus_half,real_lt_trans,real_lt_subst_left,real_lt_subst_right}.
        Thus $\delta + \delta + \delta < \epsilon$
            by \cref{real_add_comm,real_add_assoc,real_half_plus_half,real_lt_subst_left}.
        Follows by definition.
    \end{byCase}
\end{proof}

% from §45 Item 4: total boundedness implies equicontinuity
\begin{lemma}\label{totally_bounded_function_family_equicontinuous}
    Assume $T$ is a topology on $X$.
    Assume $d$ is a metric on $Y$.
    Assume $\rho$ is the uniform metric on functions from $X$ to $Y$ with respect to $d$.
    Assume $C$ is the set of continuous functions from $X$ to $Y$ with respect to $T$ and $\metricTopology{d}{Y}$.
    Assume $F\subseteq C$.
    Suppose $F$ is totally bounded in $\funs{X}{Y}$ with respect to $\rho$.
    Then $F$ is equicontinuous from $X$ to $Y$ with respect to $T$ and $d$.
\end{lemma}
\begin{proof}
    $F\subseteq\funs{X}{Y}$ by \cref{subseteq,continuous_function_set}.
    \begin{byCase}
    \caseOf{$F = \emptyset$.}
        Show that for all $x_0\in X$ we have
            $F$ is equicontinuous at $x_0$ from $X$ to $Y$ with respect to $T$ and $d$.
        \begin{subproof}
            Follows by \cref{topology_on,open_in,neighborhood,real_lt_subst_right,emptyset,equicontinuous_at_point}.
        \end{subproof}
        Hence $F$ is equicontinuous from $X$ to $Y$ with respect to $T$ and $d$
            by \cref{equicontinuous_family}.
    \caseOf{$F\neq\emptyset$.}
        $F$ is inhabited by \cref{nonempty_inhabited}.
        Show that for all $x_0\in X$ we have
            $F$ is equicontinuous at $x_0$ from $X$ to $Y$ with respect to $T$ and $d$.
        \begin{subproof}
            Fix $x_0$ such that $x_0\in X$.
            Show that for all $\epsilon\in\reals$ such that $0 < \epsilon$
                there exists $U$ such that
                    $U$ is a neighborhood of $x_0$ in $X$ with respect to $T$ and
                    for all $x, f$ such that $x\in U$ and $f\in F$ we have
                        $d((\apply{f}{x},\apply{f}{x_0})) < \epsilon$.
            \begin{subproof}
                Fix $\epsilon$ such that $\epsilon\in\reals$ and $0 < \epsilon$.
                Take $\delta\in\reals$ such that
                    $0 < \delta$ and
                    $\delta < 1$ and
                    $\delta + \delta + \delta < \epsilon$
                    by \cref{positive_real_small_triple_bound}.
                Take $H$ such that
                    $H\subseteq F$ and
                    $H$ is finite and
                    for all $f\in F$ there exists $h\in H$ such that
                        $\rho((f,h)) < \delta$
                    by \cref{uniform_metric_on_function_space,totally_bounded_subset_finite_internal_net}.
                Hence $H\subseteq C$ by \cref{subseteq_transitive}.
                Take $U$ such that
                    $U$ is a neighborhood of $x_0$ in $X$ with respect to $T$ and
                    for all $x, h$ such that $x\in U$ and $h\in H$ we have
                        $d((h(x),h(x_0))) < \delta$
                    by \cref{finite_continuous_family_common_neighborhood}.
                Show that for all $x, f$ such that $x\in U$ and $f\in F$ we have
                    $d((\apply{f}{x},\apply{f}{x_0})) < \epsilon$.
                \begin{subproof}
                    Fix $x$.
                    Fix $f$.
                    Assume $x\in U$ and $f\in F$.
                    Take $h\in H$ such that $\rho((f,h)) < \delta$ by \cref{subseteq}.
                    We have $x\in X$ and $f\in\funs{X}{Y}$ and $h\in\funs{X}{Y}$
                        by \cref{neighborhood,open_in,topology_on,powerset_elim,subseteq,subseteq_transitive}.
                    Hence $d((f(x),h(x))) < \delta$ by \cref{uniform_metric_small_implies_pointwise_small}.
                    Therefore $d((h(x_0),f(x_0))) < \delta$
                        by \cref{uniform_metric_small_implies_pointwise_small,metric_on,metric_symmetric,funs_type_apply,real_lt_subst_left}.
                    We have $d((h(x),h(x_0))) < \delta$ by \cref{subseteq}.
                    Let $a = d((f(x),h(x)))$.
                    Let $b = d((h(x),h(x_0)))$.
                    Let $c = d((h(x_0),f(x_0)))$.
                    We have $a\in\reals$ and $b\in\reals$ and $c\in\reals$ and
                        $a < \delta$ and $b < \delta$ and $c < \delta$
                        by \cref{metric_on,times_tuple_intro,funs_type_apply,real_lt_subst_left,real_lt_subst_right}.
                    Hence $a + b < \delta + \delta$ by \cref{real_lt_add_two}.
                    Thus $(a + b) + c < \epsilon$
                        by \cref{real_lt_add_two,real_add_closed,real_add_assoc,real_lt_trans,real_lt_subst_right}.
                    We have $d((f(x),f(x_0))) \leq a + d((h(x),f(x_0)))$
                        by \cref{metric_on,metric_triangle_inequality,funs_type_apply}.
                    We have $d((h(x),f(x_0))) \leq b + c$
                        by \cref{metric_on,metric_triangle_inequality,funs_type_apply}.
                    We have $d((h(x),f(x_0)))\in\reals$ and $d((f(x),f(x_0)))\in\reals$
                        by \cref{metric_on,funs_type_apply,times_tuple_intro}.
                    $a + d((h(x),f(x_0)))\in\reals$ and $(a + b) + c\in\reals$ and $b + c\in\reals$
                        by \cref{real_add_closed}.
                    $a + d((h(x),f(x_0))) < (a + b) + c$ or
                        $a + d((h(x),f(x_0))) = (a + b) + c$
                        by \cref{real_leq_split,real_lt_add,real_add_comm,real_lt_subst_left,real_lt_subst_right,real_add_assoc,real_add_eq_subst}.
                    Therefore $a + d((h(x),f(x_0))) < \epsilon$ by \cref{real_lt_trans,real_lt_subst_left}.
                    Hence $d((f(x),f(x_0))) < \epsilon$
                        by \cref{real_leq_split,real_lt_trans,real_lt_subst_left}.
                \end{subproof}
                Follows by definition.
            \end{subproof}
            Hence $F$ is equicontinuous at $x_0$ from $X$ to $Y$ with respect to $T$ and $d$
                by \cref{equicontinuous_at_point}.
        \end{subproof}
        Hence $F$ is equicontinuous from $X$ to $Y$ with respect to $T$ and $d$
            by \cref{equicontinuous_family}.
    \end{byCase}
\end{proof}

% from §45 Item 5: equicontinuity implies total boundedness in the uniform metric
\begin{lemma}\label{common_center_small_ball_distance_bound}
    Assume $d$ is a metric on $Y$.
    Suppose $y_1,y_2,z\in Y$.
    Suppose $\alpha\in\reals$.
    Suppose $d((y_1,z)) < \alpha$.
    Suppose $d((y_2,z)) < \alpha$.
    Then $d((y_1,y_2)) < \alpha + \alpha$.
\end{lemma}
\begin{proof}
    Follows by \cref{metric_common_center_double_bound}.
\end{proof}

% from §45 helper lemma 10: every positive real has a positive half
\begin{lemma}\label{positive_real_half_decomposition}
    Suppose $\delta\in\reals$.
    Suppose $0 < \delta$.
    Then there exists $\alpha\in\reals$ such that
        $0 < \alpha$ and
        $\alpha + \alpha = \delta$ and
        $\alpha < \delta$.
\end{lemma}
\begin{proof}
    Let $\alpha = \delta / (1 + 1)$.
    Show that $\alpha\in\reals$ and $0 < \alpha$ and $\alpha + \alpha = \delta$ and $\alpha < \delta$.
    \begin{subproof}
        Follows by \cref{real_half_of_positive,real_add_ident,real_lt_add,real_lt_subst_left,real_lt_subst_right}.
    \end{subproof}
    Follows by definition.
\end{proof}

% from §45 helper lemma 11: Cartesian products of finite sets are finite
\begin{lemma}\label{finite_cartesian_product}
    Suppose $A$ is finite.
    Suppose $B$ is finite.
    Then $A\times B$ is finite.
\end{lemma}
\begin{proof}
    \begin{byCase}
    \caseOf{$A$ is empty.}
        Show that for all $p$ we have $p\notin A\times B$.
        \begin{subproof}
            Fix $p$.
            Suppose not.
            Then $p\in A\times B$.
            Take $a,b$ such that $p = (a,b)$ and $a\in A$ and $b\in B$ by \cref{times}.
            Contradiction.
        \end{subproof}
        Hence $A\times B$ is empty.
        Thus $A\times B$ is finite by \cref{finite}.
    \caseOf{$A$ is not empty.}
    Take $n\in\naturalsPlus$ such that there exists a bijection from $\initialSegment{n}$ to $A$
        by \cref{finite}.
    Take $f$ such that $f$ is a bijection from $\initialSegment{n}$ to $A$.
    We have $f\in\surj{\initialSegment{n}}{A}$ and $f\in\funs{\initialSegment{n}}{A}$ by \cref{bijections,bijections_to_funs}.
    $\initialSegment{n}$ is finite by \cref{initial_segment_finite}.
    Let $g = \{(i,\{f(i)\}\times B) \mid i\in \initialSegment{n}\}$.
    For all $i,U_1,U_2$ such that $(i,U_1)\in g$ and $(i,U_2)\in g$ we have $U_1 = U_2$
        by \cref{pair_eq_iff}.
    Hence $g$ is a function on $\initialSegment{n}$
        by \cref{rightunique,relation,dom_intro,dom_iff,pair_eq_iff,subseteq,subseteq_antisymmetric}.
    Let $H = \img{g}{\initialSegment{n}}$.
    Hence $H$ is finite by \cref{finite_image_function}.
    Show that for all $C\in H$ we have $C$ is finite.
    \begin{subproof}
        Fix $C$ such that $C\in H$.
        Take $i\in\initialSegment{n}$ such that $(i,C)\in g$ by \cref{img_iff}.
        We have $C = \{f(i)\}\times B$ by \cref{function_apply_intro}.
        Let $h = \{(b,(f(i),b)) \mid b\in B\}$.
        For all $b,y_1,y_2$ such that $(b,y_1)\in h$ and $(b,y_2)\in h$ we have $y_1 = y_2$
            by \cref{pair_eq_iff}.
        We have $h$ is a function by \cref{rightunique,relation}.
        Therefore $\dom{h} = B$ by
            \cref{dom_intro,dom_iff,pair_eq_iff,subseteq,subseteq_antisymmetric}.
        For all $y$ such that $y\in\ran{h}$ we have $y\in C$
            by \cref{ran_iff,pair_eq_iff,singleton_intro,times_tuple_intro,real_lt_subst_right}.
        We have $\ran{h}\subseteq C$ by \cref{subseteq}.
        For all $y$ such that $y\in C$ we have $y\in\ran{h}$
            by \cref{real_lt_subst_left,times,singleton_elim,pair_eq_iff,ran_intro}.
        We have $C\subseteq\ran{h}$ by \cref{subseteq}.
        Hence $\ran{h} = C$ by \cref{subseteq_antisymmetric}.
        Therefore $h\in\funs{B}{C}$ by \cref{function_apply_iff,ran_intro,subseteq,funs_intro,subseteq_antisymmetric}.
        For all $y\in C$ there exists $b\in B$ such that $h(b) = y$.
        We have $h\in\surj{B}{C}$ and $h$ is injective
            by \cref{surj,function_apply_iff,pair_eq_iff,injective_function}.
        Hence $h$ is a bijection from $B$ to $C$ by \cref{surj,bijections}.
        Therefore $C$ is finite by \cref{finite_bijection}.
    \end{subproof}
    Hence $\unions{H}$ is finite by \cref{finite_union_of_finite_family}.
    Show that for all $p$ such that $p\in A\times B$ we have $p\in\unions{H}$.
    \begin{subproof}
        Follows by \cref{times,surj,function_apply_intro,img_iff,function_apply_iff,singleton_intro,times_tuple_intro,real_lt_subst_left,real_lt_subst_right,unions_iff}.
    \end{subproof}
    We have $A\times B\subseteq\unions{H}$ by \cref{subseteq}.
    Show that for all $p$ such that $p\in\unions{H}$ we have $p\in A\times B$.
    \begin{subproof}
        Follows by \cref{unions_iff,img_iff,function_apply_intro,real_lt_subst_right,times,singleton_elim,bijections_to_funs,funs_type_apply,subseteq,real_lt_subst_left,times_tuple_intro}.
    \end{subproof}
    We have $\unions{H}\subseteq A\times B$ by \cref{subseteq}.
    Thus $\unions{H} = A\times B$ by \cref{subseteq_antisymmetric}.
    Therefore $A\times B$ is finite.
    \end{byCase}
\end{proof}

% from §45 helper lemma 12a: singleton graphs on the first initial segment are typed functions
\begin{lemma}\label{singleton_initial_segment_function}
    Suppose $r\in\naturalsPlus$.
    Suppose $a\in\initialSegment{r}$.
    Then $\{(1,a)\}\in\funs{\initialSegment{1}}{\initialSegment{r}}$.
\end{lemma}
\begin{proof}
    $\{(1,a)\}$ is a function by \cref{singleton_elim,pair_eq_iff,rightunique,relation}.
    Show that $\dom{\{(1,a)\}} = \initialSegment{1}$.
    \begin{subproof}
        Follows by \cref{dom_iff,singleton_elim,pair_eq_iff,real_one_in_naturals,initial_segment_naturalsplus,real_naturals_subset_reals,subseteq,real_one_leq_naturalsplus,real_leq_antisym,singleton_intro,dom_intro,subseteq_antisymmetric}.
    \end{subproof}
    Show that for all $x$ such that $x\in\dom{\{(1,a)\}}$ we have $\apply{\{(1,a)\}}{x}\in\initialSegment{r}$.
    \begin{subproof}
        Follows by \cref{dom_iff,singleton_elim,pair_eq_iff,singleton_intro,function_apply_intro,rightunique,relation,real_lt_subst_left}.
    \end{subproof}
    Hence $\{(1,a)\}\in\funs{\initialSegment{1}}{\initialSegment{r}}$ by \cref{funs_intro}.
\end{proof}

% from §45 helper lemma 12: finite grids of index assignments are finite
\begin{lemma}\label{finite_initial_segment_function_space}
    Suppose $n\in\naturalsPlus$.
    Suppose $m\in\naturalsPlus$.
    Then $\funs{\initialSegment{n}}{\initialSegment{m}}$ is finite.
\end{lemma}
\begin{proof}
    Let $A = \{k\in\naturalsPlus \mid \text{for all $m\in\naturalsPlus$ we have $\funs{\initialSegment{k}}{\initialSegment{m}}$ is finite}\}$.
    $A\subseteq\naturalsPlus$ by \cref{subseteq}.
        Show that for all $r\in\naturalsPlus$ we have $\funs{\initialSegment{1}}{\initialSegment{r}}$ is finite.
        \begin{subproof}
        Fix $r$ such that $r\in\naturalsPlus$.
        Let $g = \{(a,\{(1,a)\}) \mid a\in\initialSegment{r}\}$.
        For all $a,b_1,b_2$ such that $(a,b_1)\in g$ and $(a,b_2)\in g$ we have $b_1 = b_2$
            by \cref{pair_eq_iff}.
        Hence $g$ is a function by \cref{rightunique,relation}.
            For all $a$ such that $a\in\dom{g}$ we have $g(a)\in\funs{\initialSegment{1}}{\initialSegment{r}}$
                by \cref{dom_iff,pair_eq_iff,subseteq,function_apply_intro,singleton_initial_segment_function,real_lt_subst_left}.
            We have $\dom{g} = \initialSegment{r}$ by
                \cref{dom_intro,dom_iff,pair_eq_iff,subseteq,subseteq_antisymmetric}.
        Hence $g\in\funs{\initialSegment{r}}{\funs{\initialSegment{1}}{\initialSegment{r}}}$ by \cref{funs_intro}.
        For all $a_1,a_2$ such that $a_1\in\dom{g}$ and $a_2\in\dom{g}$ and $g(a_1) = g(a_2)$ we have $a_1 = a_2$
            by \cref{subseteq,function_apply_intro,real_lt_subst_left,real_lt_subst_right,singleton_iff,pair_eq_iff}.
        Show that $g$ is injective.
        \begin{subproof}
            Follows by \cref{injective_function,funs_is_function}.
        \end{subproof}
            Show that for all $u\in\funs{\initialSegment{1}}{\initialSegment{r}}$ there exists $a\in\initialSegment{r}$ such that $g(a) = u$.
            \begin{subproof}
                Fix $u$ such that $u\in\funs{\initialSegment{1}}{\initialSegment{r}}$.
                Let $a = u(1)$.
                We have $a\in\initialSegment{r}$ by \cref{real_one_in_naturals,initial_segment_naturalsplus,funs_type_apply}.
                We have $g(a) = \{(1,a)\}$ by \cref{subseteq,function_apply_intro}.
                    $u$ is a function and $\dom{u} = \initialSegment{1}$ by \cref{funs_elim}.
                    For all $x\in\initialSegment{1}$ we have $x = 1$
                        by \cref{initial_segment_naturalsplus,real_one_in_naturals,real_naturals_subset_reals,subseteq,real_one_leq_naturalsplus,real_leq_antisym}.
                    Hence for all $x\in\initialSegment{1}$ we have $\apply{\{(1,a)\}}{x} = u(x)$
                        by \cref{singleton_initial_segment_function,funs_is_function,funs_is_relation,funs_elim,singleton_intro,pair_eq_iff,function_apply_intro,real_lt_subst_left,real_lt_subst_right}.
                    Hence $\{(1,a)\} = u$
                        by \cref{singleton_initial_segment_function,funs_is_function,funs_is_relation,funs_elim,functions_eq_iff}.
                Follows by definition.
        \end{subproof}
        Show that $g\in\surj{\initialSegment{r}}{\funs{\initialSegment{1}}{\initialSegment{r}}}$.
        \begin{subproof}
            Follows by \cref{surj}.
        \end{subproof}
        Hence $g$ is a bijection from $\initialSegment{r}$ to $\funs{\initialSegment{1}}{\initialSegment{r}}$ by \cref{bijections,surj}.
        Therefore $\funs{\initialSegment{1}}{\initialSegment{r}}$ is finite by \cref{initial_segment_finite,finite_bijection}.
        \end{subproof}
        Therefore $1\in A$ by \cref{real_one_in_naturals,subseteq}.
    Show that for all $k\in A$ we have $k + 1\in A$.
    \begin{subproof}
        Fix $k$ such that $k\in A$.
        Show that for all $r\in\naturalsPlus$ we have $\funs{\initialSegment{k + 1}}{\initialSegment{r}}$ is finite.
        \begin{subproof}
        Fix $r$ such that $r\in\naturalsPlus$.
        We have $\funs{\initialSegment{k}}{\initialSegment{r}}\times\initialSegment{r}$ is finite
            by \cref{subseteq,initial_segment_finite,finite_cartesian_product}.
        Let $g = \{(p,\fst{p}\union\{(k + 1,\snd{p})\}) \mid p\in\funs{\initialSegment{k}}{\initialSegment{r}}\times\initialSegment{r}\}$.
        For all $p,u_1,u_2$ such that $(p,u_1)\in g$ and $(p,u_2)\in g$ we have $u_1 = u_2$
            by \cref{pair_eq_iff}.
        Hence $g$ is a function by \cref{rightunique,relation}.
            Show that for all $p$ such that $p\in\dom{g}$ we have $g(p)\in\funs{\initialSegment{k + 1}}{\initialSegment{r}}$.
            \begin{subproof}
                Fix $p$ such that $p\in\dom{g}$.
                Take $u$ such that $(p,u)\in g$ by \cref{dom_iff}.
                We have $p\in\funs{\initialSegment{k}}{\initialSegment{r}}\times\initialSegment{r}$ by \cref{pair_eq_iff,subseteq}.
                Take $v,b$ such that $p = (v,b)$ by \cref{times_elem_is_tuple}.
                We have $v\in\funs{\initialSegment{k}}{\initialSegment{r}}$ and $b\in\initialSegment{r}$ by \cref{times_tuple_elim}.
                Hence $(p,\fst{p}\union\{(k + 1,\snd{p})\})\in g$ by \cref{subseteq}.
                We have $g(p) = v\union\{(k + 1,b)\}$ by \cref{fst_eq,snd_eq,function_apply_intro}.
                Let $s = \{(k + 1,b)\}$.
                For all $x$ such that $x\in\dom{v}\inter\dom{s}$ we have $v(x) = s(x)$
                    by \cref{inter_elim_left,inter_elim_right,funs_elim,dom_iff,singleton_elim,pair_eq_iff,initial_segment_naturalsplus,real_naturals_subset_reals,subseteq,real_naturals_closed_succ,real_add_closed,real_one_in_naturals,real_naturals_lt_succ_local,real_lt_subst_right,real_lt_trans,real_lt_irreflexive}.
                Hence $v\union s$ is a function by \cref{funs_is_function,singleton_elim,pair_eq_iff,rightunique,relation,union_compatible_functions}.
                We have $\dom{s} = \{k + 1\}$
                    by \cref{subseteq_antisymmetric,subseteq,dom_iff,singleton_elim,pair_eq_iff,dom_intro,singleton_intro}.
                Hence $\dom{(v\union s)} = \initialSegment{k}\union\{k + 1\}$ by \cref{dom_union,funs_elim,real_lt_subst_right}.
                For all $x$ such that $x\in\initialSegment{k}\union\{k + 1\}$ we have $x\in\initialSegment{k + 1}$
                    by \cref{subseteq,real_naturals_closed_succ,union_iff,initial_segment_naturalsplus,real_naturals_subset_reals,real_one_in_naturals,real_add_ident,real_add_closed,real_zero_lt_one,real_lt_add,real_add_comm,real_lt_subst_left,real_lt_trans,singleton_elim}.
                We have $\dom{(v\union s)}\subseteq\initialSegment{k + 1}$ by \cref{subseteq}.
                For all $x$ such that $x\in\initialSegment{k + 1}$ we have $x\in\dom{(v\union s)}$
                    by \cref{singleton_intro,real_lt_subst_right,union_iff,real_naturals_leq_succ_elim,initial_segment_naturalsplus,funs_elim,dom_union}.
                We have $\initialSegment{k + 1}\subseteq\dom{(v\union s)}$ by \cref{subseteq}.
                Thus $\dom{(v\union s)} = \initialSegment{k + 1}$ by \cref{subseteq_antisymmetric,subseteq}.
                For all $x$ such that $x\in\dom{(v\union s)}$ we have $(v\union s)(x)\in\initialSegment{r}$
                    by \cref{dom_union,union_iff,funs_elim,function_apply_elim,function_apply_intro,funs_type_apply,real_lt_subst_left,dom_iff,singleton_elim,pair_eq_iff,singleton_intro}.
                Hence $v\union s\in\funs{\initialSegment{k + 1}}{\initialSegment{r}}$ by \cref{funs_intro}.
                Therefore $g(p)\in\funs{\initialSegment{k + 1}}{\initialSegment{r}}$ by \cref{real_lt_subst_left}.
            \end{subproof}
            Hence $\dom{g} = \funs{\initialSegment{k}}{\initialSegment{r}}\times\initialSegment{r}$ by
                \cref{dom_intro,dom_iff,pair_eq_iff,subseteq,subseteq_antisymmetric}.
            Therefore $g\in\funs{\funs{\initialSegment{k}}{\initialSegment{r}}\times\initialSegment{r}}{\funs{\initialSegment{k + 1}}{\initialSegment{r}}}$ by \cref{funs_intro}.
        We have $g$ is a function and $g$ is a relation and
            $\dom{g} = \funs{\initialSegment{k}}{\initialSegment{r}}\times\initialSegment{r}$
            by \cref{funs_is_function,funs_is_relation,funs}.
        Show that for all $p_1,p_2$ such that $p_1\in\funs{\initialSegment{k}}{\initialSegment{r}}\times\initialSegment{r}$ and
            $p_2\in\funs{\initialSegment{k}}{\initialSegment{r}}\times\initialSegment{r}$ and
            $g(p_1) = g(p_2)$ we have $p_1 = p_2$.
        \begin{subproof}
            Fix $p_1$.
            Fix $p_2$ such that $p_1\in\funs{\initialSegment{k}}{\initialSegment{r}}\times\initialSegment{r}$ and
                $p_2\in\funs{\initialSegment{k}}{\initialSegment{r}}\times\initialSegment{r}$ and
                $g(p_1) = g(p_2)$.
            Take $v_1,b_1$ such that $p_1 = (v_1,b_1)$ by \cref{times_elem_is_tuple}.
            Take $v_2,b_2$ such that $p_2 = (v_2,b_2)$ by \cref{times_elem_is_tuple}.
            Hence $v_1\in\funs{\initialSegment{k}}{\initialSegment{r}}$ and $b_1\in\initialSegment{r}$ and
                $v_2\in\funs{\initialSegment{k}}{\initialSegment{r}}$ and $b_2\in\initialSegment{r}$ by \cref{times_tuple_elim}.
            We have $g(p_1) = v_1\union\{(k + 1,b_1)\}$ and $g(p_2) = v_2\union\{(k + 1,b_2)\}$
                by \cref{fst_eq,snd_eq,function_apply_intro}.
            Hence $v_1\union\{(k + 1,b_1)\} = v_2\union\{(k + 1,b_2)\}$ by \cref{real_lt_subst_left,real_lt_subst_right}.
            Show that for all $x\in\initialSegment{k}$ we have $v_1(x) = v_2(x)$.
            \begin{subproof}
                Fix $x$ such that $x\in\initialSegment{k}$.
                We have $x\in\dom{v_1}$ and $x\in\dom{v_2}$ by \cref{funs_elim}.
                $v_1$ is a function and $v_1$ is a relation and $v_1$ is right-unique
                    by \cref{funs_is_function,funs_is_relation,function_rightunique}.
                $v_2$ is a function and $v_2$ is a relation and $v_2$ is right-unique
                    by \cref{funs_is_function,funs_is_relation,function_rightunique}.
                Let $u_2 = v_2\union\{(k + 1,b_2)\}$.
                We have $u_2$ is a function and $u_2$ is a relation and $u_2$ is right-unique
                    by \cref{real_lt_subst_right,funs_type_apply,funs_is_function,funs_is_relation,function_rightunique}.
                We have $(x,v_1(x))\in v_1$ by \cref{function_apply_elim}.
                Hence $(x,v_1(x))\in v_1\union\{(k + 1,b_1)\}$ by \cref{union_iff}.
                We have $(x,v_1(x))\in u_2$ by \cref{real_lt_subst_left,real_lt_subst_right}.
                Hence $u_2(x) = v_1(x)$ by \cref{function_apply_intro}.
                We have $(x,v_2(x))\in v_2$ by \cref{function_apply_elim}.
                Thus $(x,v_2(x))\in u_2$ by \cref{union_iff}.
                Therefore $u_2(x) = v_2(x)$ by \cref{function_apply_intro}.
                Thus $v_1(x) = v_2(x)$ by \cref{real_lt_subst_left,real_lt_subst_right}.
            \end{subproof}
            Hence $v_1 = v_2$ by \cref{functions_eq_iff}.
            We have $(k + 1,b_1)\in v_2\union\{(k + 1,b_2)\}$ by \cref{union_iff,singleton_intro,real_lt_subst_left}.
            Show that $(k + 1,b_1)\notin v_2$.
            \begin{subproof}
                Suppose not.
                We have $k + 1\in\dom{v_2}$ by \cref{dom_iff}.
                Hence $k + 1\in\initialSegment{k}$ by \cref{funs_elim}.
                $k + 1 < k$ or $k + 1 = k$ by \cref{initial_segment_naturalsplus}.
                $k\in\naturalsPlus$ by \cref{subseteq}.
                We have $k\in\reals$ and $1\in\reals$ and $0\in\reals$
                    by \cref{real_naturals_subset_reals,subseteq,real_one_in_naturals,real_add_ident}.
                $k < k + 1$ by \cref{real_zero_lt_one,real_lt_add,real_add_ident,real_add_comm,real_lt_subst_left}.
                Hence $k + 1\in\reals$ by \cref{real_add_closed}.
                \begin{byCase}
                \caseOf{$k + 1 < k$.}
                    Contradiction by \cref{real_lt_irreflexive,real_lt_trans}.
                \caseOf{$k + 1 = k$.}
                    We have $k < k$ by \cref{real_lt_subst_left}.
                    Contradiction by \cref{real_lt_irreflexive}.
                \end{byCase}
            \end{subproof}
            $(k + 1,b_1)\in\{(k + 1,b_2)\}$ by \cref{union_iff}.
            Hence $(k + 1,b_1) = (k + 1,b_2)$ by \cref{singleton_elim}.
            Hence $b_1 = b_2$ by \cref{pair_eq_iff}.
            Thus $p_1 = p_2$ by \cref{pair_eq_iff}.
        \end{subproof}
        Show that $g$ is injective.
        \begin{subproof}
            Follows by \cref{injective_function,funs}.
        \end{subproof}
            Show that for all $u\in\funs{\initialSegment{k + 1}}{\initialSegment{r}}$ there exists $p\in\funs{\initialSegment{k}}{\initialSegment{r}}\times\initialSegment{r}$ such that $g(p) = u$.
            \begin{subproof}
                Fix $u$ such that $u\in\funs{\initialSegment{k + 1}}{\initialSegment{r}}$.
                Let $v = \restrl{u}{\initialSegment{k}}$.
                $k\in\naturalsPlus$ by \cref{subseteq}.
                We have $u$ is a function and $\dom{u} = \initialSegment{k + 1}$ by \cref{funs_elim}.
                $v$ is a function by \cref{restrl_of_function_is_function}.
                $\dom{v} = \initialSegment{k}$
                    by \cref{initial_segment_subset_succ,subseteq,elem_dom_and_restr_implies_elem_of_restr,restrl_dom,subseteq_antisymmetric}.
                Therefore $v\in\funs{\initialSegment{k}}{\initialSegment{r}}$
                    by \cref{funs_elim,initial_segment_subset_succ,subseteq,restrl_of_function_apply,funs_intro,elem_dom_of_restrl_implies_elem_dom_and_restr}.
                Let $b = u(k + 1)$.
                We have $b\in\initialSegment{r}$ by \cref{real_naturals_closed_succ,initial_segment_naturalsplus,funs_type_apply}.
                Let $p = (v,b)$.
                We have $p\in\funs{\initialSegment{k}}{\initialSegment{r}}\times\initialSegment{r}$ by \cref{times_tuple_intro}.
                We have $g(p) = v\union\{(k + 1,b)\}$ by \cref{fst_eq,snd_eq,function_apply_intro}.
                    $v\union\{(k + 1,b)\}$ is a function and $\dom{(v\union\{(k + 1,b)\})} = \initialSegment{k + 1}$
                        by \cref{funs_type_apply,real_lt_subst_left,funs_elim}.
                    For all $x\in\initialSegment{k + 1}$ we have $(v\union\{(k + 1,b)\})(x) = u(x)$
                        by \cref{union_iff,singleton_intro,function_apply_intro,real_lt_subst_left,real_naturals_leq_succ_elim,initial_segment_naturalsplus,funs_elim,function_apply_elim,initial_segment_subset_succ,subseteq,real_lt_subst_right,restrl_of_function_apply}.
                    Hence $v\union\{(k + 1,b)\} = u$ by \cref{functions_eq_iff}.
                Follows by definition.
            \end{subproof}
            Show that $g\in\surj{\funs{\initialSegment{k}}{\initialSegment{r}}\times\initialSegment{r}}{\funs{\initialSegment{k + 1}}{\initialSegment{r}}}$.
            \begin{subproof}
                Follows by \cref{surj}.
            \end{subproof}
        Hence $g$ is a bijection from $\funs{\initialSegment{k}}{\initialSegment{r}}\times\initialSegment{r}$ to $\funs{\initialSegment{k + 1}}{\initialSegment{r}}$
            by \cref{bijections,surj}.
        Therefore $\funs{\initialSegment{k + 1}}{\initialSegment{r}}$ is finite by \cref{finite_bijection}.
        \end{subproof}
        Therefore $k + 1\in A$ by \cref{subseteq,real_naturals_closed_succ}.
    \end{subproof}
    Hence $n\in A$ by \cref{real_naturals_induction}.
    Therefore $\funs{\initialSegment{n}}{\initialSegment{m}}$ is finite by \cref{subseteq}.
\end{proof}

% from §45 helper lemma 13: function spaces over finite sets are finite
\begin{lemma}\label{finite_function_space}
    Suppose $A$ is finite.
    Suppose $B$ is finite.
    Then $\funs{A}{B}$ is finite.
\end{lemma}
\begin{proof}
    \begin{byCase}
    \caseOf{$A$ is empty.}
        For all $f$ such that $f\in\funs{A}{B}$ we have $f\in\{\emptyset\}$.
        We have $\funs{A}{B}\subseteq\{\emptyset\}$ by \cref{subseteq}.
        Therefore $\funs{A}{B}$ is finite by \cref{singleton_subset_intro,upair_intro_left,pair_is_finite,subseteq_finite_is_finite}.
    \caseOf{$A$ is not empty.}
        Take $n\in\naturalsPlus$ such that there exists a bijection from $\initialSegment{n}$ to $A$
            by \cref{finite}.
        Take $f$ such that $f$ is a bijection from $\initialSegment{n}$ to $A$.
        \begin{byCase}
        \caseOf{$B$ is empty.}
            Show that $\funs{A}{B}$ is empty.
            \begin{subproof}
                Show that for all $u$ we have $u\notin\funs{A}{B}$.
                \begin{subproof}
                    Fix $u$.
                    Suppose not.
                    Then $u\in\funs{A}{B}$.
                    Take $a$ such that $a\in A$ by \cref{nonempty_inhabited}.
                    Then $u(a)\in B$ by \cref{funs_type_apply}.
                    Contradiction.
                \end{subproof}
            \end{subproof}
            Thus $\funs{A}{B}$ is finite by \cref{finite}.
        \caseOf{$B$ is not empty.}
            Take $m\in\naturalsPlus$ such that there exists a bijection from $\initialSegment{m}$ to $B$
                by \cref{finite}.
            Take $e$ such that $e$ is a bijection from $\initialSegment{m}$ to $B$.
            Let $R = \funs{\initialSegment{n}}{\initialSegment{m}}$.
            We have $R$ is finite by \cref{finite_initial_segment_function_space}.
            $\converse{f}\in\funs{A}{\initialSegment{n}}$ and
                $\converse{e}\in\funs{B}{\initialSegment{m}}$
                by \cref{bijective_converse_are_funs}.
            $\converse{f}$ is a function and $\dom{\converse{f}} = A$ and
                $\converse{e}$ is a function and $\dom{\converse{e}} = B$
                by \cref{funs_elim}.
            $f\in\funs{\initialSegment{n}}{A}$ and $e\in\funs{\initialSegment{m}}{B}$
                by \cref{bijections_to_funs}.
            Let $q = \{(u,e\circ (u\circ \converse{f})) \mid u\in R\}$.
            For all $u,v_1,v_2$ such that $(u,v_1)\in q$ and $(u,v_2)\in q$ we have $v_1 = v_2$
                by \cref{pair_eq_iff}.
            Hence $q$ is a function on $R$
                by \cref{rightunique,relation,dom_intro,dom_iff,pair_eq_iff,subseteq,subseteq_antisymmetric}.
            Let $H = \img{q}{R}$.
            Hence $H$ is finite by \cref{finite_image_function}.
            Show that $\funs{A}{B}\subseteq H$.
            \begin{subproof}
                Show that for all $v$ such that $v\in\funs{A}{B}$ we have $v\in H$.
                \begin{subproof}
                    Fix $v$ such that $v\in\funs{A}{B}$.
                    Let $u = \converse{e}\circ (v\circ f)$.
                    We have $v\circ f\in\funs{\initialSegment{n}}{B}$ by \cref{funs_circ}.
                    Hence $u\in R$ by \cref{funs_circ}.
                    We have $q(u) = e\circ (u\circ \converse{f})$ by \cref{function_apply_intro}.
                    Show that $e\circ (u\circ \converse{f}) = v$.
                    \begin{subproof}
                        We have $e\circ (u\circ \converse{f})\in\funs{A}{B}$ by \cref{funs_circ}.
                        Show that for all $x\in A$ we have $\apply{e\circ (u\circ \converse{f})}{x} = v(x)$.
                        \begin{subproof}
                            Fix $x\in A$.
                            Let $x_0 = (\converse{f})(x)$.
                            $x_0\in\initialSegment{n}$ by \cref{funs_type_apply}.
                            $(x,x_0)\in\converse{f}$ by \cref{function_apply_elim}.
                            Hence $(x_0,x)\in f$ by \cref{converse_iff}.
                            $f(x_0) = x$ by \cref{function_apply_intro,bijections_to_funs,funs_is_function}.
                            $v(x)\in B$ by \cref{funs_type_apply}.
                            Let $y_0 = (\converse{e})(v(x))$.
                            $y_0\in\initialSegment{m}$ by \cref{funs_type_apply}.
                            $(v(x),y_0)\in\converse{e}$ by \cref{function_apply_elim}.
                            Hence $(y_0,v(x))\in e$ by \cref{converse_iff}.
                            $e(y_0) = v(x)$ by \cref{function_apply_intro,bijections_to_funs,funs_is_function}.
                            We have $u\in\funs{\initialSegment{n}}{\initialSegment{m}}$ by \cref{subseteq}.
                            $u$ is a function and $\dom{u} = \initialSegment{n}$ by \cref{funs_elim}.
                            $v$ is a function and $\dom{v} = A$ by \cref{funs_elim}.
                            $f$ is a function and $\dom{f} = \initialSegment{n}$ by \cref{bijections_to_funs,funs_elim}.
                            $e$ is a function and $\dom{e} = \initialSegment{m}$ by \cref{bijections_to_funs,funs_elim}.
                            We have $u\circ\converse{f}\in\funs{A}{\initialSegment{m}}$ by \cref{funs_circ}.
                            $(x_0,u(x_0))\in u$ by \cref{function_apply_elim}.
                            $(x,u(x_0))\in u\circ\converse{f}$ by \cref{circ_elem_intro}.
                            $\apply{u\circ\converse{f}}{x} = u(x_0)$
                                by \cref{function_apply_intro,funs_is_function}.
                            $(x_0,f(x_0))\in f$ by \cref{function_apply_elim}.
                            $(f(x_0),v(f(x_0)))\in v$ by \cref{function_apply_elim,real_lt_subst_right}.
                            $(x_0,v(f(x_0)))\in v\circ f$ by \cref{circ_elem_intro}.
                            $\apply{v\circ f}{x_0} = v(f(x_0))$
                                by \cref{function_apply_intro,funs_is_function,funs_circ}.
                            Hence $\apply{v\circ f}{x_0} = v(x)$ by \cref{real_lt_subst_right}.
                            $(x_0,\apply{v\circ f}{x_0})\in v\circ f$ by \cref{function_apply_elim,funs_circ,funs}.
                            $(\apply{v\circ f}{x_0},\apply{\converse{e}}{\apply{v\circ f}{x_0}})\in\converse{e}$
                                by \cref{function_apply_elim,funs}.
                            $(x_0,\apply{\converse{e}}{\apply{v\circ f}{x_0}})\in \converse{e}\circ (v\circ f)$
                                by \cref{circ_elem_intro}.
                            $(x_0,\apply{\converse{e}}{\apply{v\circ f}{x_0}})\in u$ by \cref{real_lt_subst_left}.
                            $u(x_0) = \apply{\converse{e}}{\apply{v\circ f}{x_0}}$
                                by \cref{function_apply_intro,funs_is_function}.
                            Hence $u(x_0) = y_0$ by \cref{real_lt_subst_left,real_lt_subst_right}.
                            Thus $\apply{u\circ\converse{f}}{x} = y_0$ by \cref{real_lt_subst_right}.
                            $(x,\apply{u\circ\converse{f}}{x})\in u\circ\converse{f}$ by \cref{function_apply_elim,funs}.
                            $(\apply{u\circ\converse{f}}{x},e(\apply{u\circ\converse{f}}{x}))\in e$
                                by \cref{function_apply_elim,bijections_to_funs,funs,funs_type_apply}.
                            $(x,e(\apply{u\circ\converse{f}}{x}))\in e\circ (u\circ\converse{f})$
                                by \cref{circ_elem_intro}.
                            $\apply{e\circ (u\circ \converse{f})}{x} = e(\apply{u\circ\converse{f}}{x})$
                                by \cref{function_apply_intro,funs_is_function}.
                            Thus $\apply{e\circ (u\circ \converse{f})}{x} = v(x)$
                                by \cref{real_lt_subst_left,real_lt_subst_right}.
                        \end{subproof}
                        Therefore $e\circ (u\circ \converse{f}) = v$ by \cref{functions_eq_iff}.
                    \end{subproof}
                    Thus $q(u) = v$ by \cref{real_lt_subst_left}.
                    Therefore $v\in H$ by \cref{img_iff,function_apply_iff}.
                \end{subproof}
            \end{subproof}
            Therefore $\funs{A}{B}$ is finite by \cref{subseteq_finite_is_finite}.
        \end{byCase}
    \end{byCase}
\end{proof}

% from §45 helper lemma 14: equicontinuous families admit finite uniform epsilon-nets
\begin{lemma}\label{equicontinuous_function_family_uniform_net}
    Assume $T$ is a topology on $X$.
    Assume $X$ is compact with respect to $T$.
    Assume $d$ is a metric on $Y$.
    Assume $Y$ is compact with respect to $\metricTopology{d}{Y}$.
    Assume $\rho$ is the uniform metric on functions from $X$ to $Y$ with respect to $d$.
    Assume $C$ is the set of continuous functions from $X$ to $Y$ with respect to $T$ and $\metricTopology{d}{Y}$.
    Assume $F\subseteq C$.
    Suppose $F$ is equicontinuous from $X$ to $Y$ with respect to $T$ and $d$.
    Then for all $\epsilon\in\reals$ such that $0 < \epsilon$
        there exists $H$ such that
            $H$ is finite and
            $H\subseteq\funs{X}{Y}$ and
            for all $f\in F$ there exists $g\in H$ such that
                $\rho((f,g)) < \epsilon$.
\end{lemma}
\begin{proof}
    $\rho$ is a metric on $\funs{X}{Y}$ by \cref{uniform_metric_on_function_space}.
    $F\subseteq\funs{X}{Y}$ by \cref{subseteq,continuous_function_set}.
    We have $Y$ is totally bounded with respect to $d$
        by \cref{metric_topology,metric_basis_is_basis,basis_for_topology,basis_topology_is_topology,compact_metric_space_totally_bounded}.
    Fix $\epsilon$ such that $\epsilon\in\reals$ and $0 < \epsilon$.
        Take $\delta\in\reals$ such that
            $0 < \delta$ and
            $\delta < 1$ and
            $\delta + \delta + \delta < \epsilon$
            by \cref{positive_real_small_triple_bound}.

        Take $C_Y$ such that
            $C_Y$ is finite and
            $C_Y$ is a covering of $Y$ and
            for all $V\in C_Y$ there exists $y\in Y$ such that
                $V = Y\inter\metricBall{d}{Y}{y}{\delta}$
            by \cref{totally_bounded_metric_space,totally_bounded_subset_metric_space}.

        Let $R_Y = \{w\in C_Y\times Y \mid \fst{w} = Y\inter\metricBall{d}{Y}{\snd{w}}{\delta}\}$.
        For all $V\in C_Y$ there exists $y$ such that $V\mathrel{R_Y} y$
            by \cref{subseteq,times_tuple_intro,fst_eq,snd_eq}.
        Take $p$ such that
            $p$ is a function on $C_Y$ and
            for all $V\in C_Y$ we have $V\mathrel{R_Y}\apply{p}{V}$
            by \cref{subseteq,times_elem_is_tuple,relation,choice_function}.
        We have $p\in\funs{C_Y}{Y}$
            by \cref{funs_elim,function_apply_elim,subseteq,times_tuple_elim,funs_intro}.
        For all $V\in C_Y$ we have $V = Y\inter\metricBall{d}{Y}{p(V)}{\delta}$
            by \cref{subseteq,fst_eq,snd_eq}.

        Let $N = \{U\in T \mid \text{there exists $a\in X$ such that
            $U$ is a neighborhood of $a$ in $X$ with respect to $T$ and
            for all $x,f$ such that $x\in U$ and $f\in F$ we have $d((\apply{f}{x},\apply{f}{a})) < \delta$}\}$.
        For all $x$ such that $x\in X$ we have $x\in\unions{N}$
            by \cref{equicontinuous_family,equicontinuous_at_point,neighborhood,open_in,subseteq,unions_iff}.
        Hence $N$ is an open covering of $X$ with respect to $T$
            by \cref{subseteq,covering,open_in,open_covering}.
        Take $K$ such that
            $K\subseteq N$ and
            $K$ is finite and
            $K$ is a covering of $X$
            by \cref{compact}.

        Let $R_X = \{w\in K\times X \mid \text{$\fst{w}$ is a neighborhood of $\snd{w}$ in $X$ with respect to $T$ and
            for all $x,f$ such that $x\in\fst{w}$ and $f\in F$ we have
                $d((\apply{f}{x},\apply{f}{\snd{w}})) < \delta$}\}$.
        For all $U\in K$ there exists $a$ such that $U\mathrel{R_X} a$
            by \cref{subseteq,times_tuple_intro,fst_eq,snd_eq}.
        Take $c$ such that
            $c$ is a function on $K$ and
            for all $U\in K$ we have $U\mathrel{R_X}\apply{c}{U}$
            by \cref{subseteq,times_elem_is_tuple,relation,choice_function}.
        We have $c\in\funs{K}{X}$
            by \cref{funs_elim,function_apply_elim,subseteq,times_tuple_elim,funs_intro}.
        For all $U\in K$ we have
            $U$ is a neighborhood of $c(U)$ in $X$ with respect to $T$ and
            for all $x,f$ such that $x\in U$ and $f\in F$ we have
                $d((\apply{f}{x},\apply{f}{c(U)})) < \delta$
            by \cref{subseteq,fst_eq,snd_eq}.

        Let $R_u = \{w\in X\times K \mid \fst{w}\in\snd{w}\}$.
        For all $x\in X$ there exists $U$ such that $x\mathrel{R_u} U$
            by \cref{covering,subseteq,unions_iff,times_tuple_intro,fst_eq,snd_eq}.
        Take $u$ such that
            $u$ is a function on $X$ and
            for all $x\in X$ we have $x\mathrel{R_u}\apply{u}{x}$
            by \cref{subseteq,times_elem_is_tuple,relation,choice_function}.
        We have $u\in\funs{X}{K}$
            by \cref{funs_elim,function_apply_elim,subseteq,times_tuple_elim,funs_intro}.

        Let $A = \funs{K}{C_Y}$.
        We have $A$ is finite by \cref{finite_function_space}.
        Let $h = \{(\alpha,p\circ\alpha\circ u) \mid \alpha\in A\}$.
        For all $\alpha,g_1,g_2$ such that $(\alpha,g_1)\in h$ and $(\alpha,g_2)\in h$ we have $g_1 = g_2$
            by \cref{pair_eq_iff}.
        Hence $h$ is a function on $A$
            by \cref{rightunique,relation,subseteq_antisymmetric,subseteq,dom_intro,pair_eq_iff,dom_iff}.
        Let $H = \img{h}{A}$.
        Therefore $H$ is finite by \cref{finite_image_function}.
        For all $g$ such that $g\in H$ we have $g\in\funs{X}{Y}$
            by \cref{img_iff,function_apply_intro,real_lt_subst_left,funs_circ,circ_assoc}.
        Hence $H\subseteq\funs{X}{Y}$ by \cref{subseteq}.

        Show that for all $f\in F$ there exists $g\in H$ such that $\rho((f,g)) < \epsilon$.
        \begin{subproof}
            Fix $f$ such that $f\in F$.
            We have $f\in\funs{X}{Y}$ by \cref{subseteq,continuous_function_set}.
            Let $R_f = \{w\in K\times C_Y \mid \apply{f}{\apply{c}{\fst{w}}}\in\snd{w}\}$.
            For all $U\in K$ there exists $V$ such that $U\mathrel{R_f} V$
                by \cref{funs_type_apply,covering,subseteq,unions_iff,times_tuple_intro,fst_eq,snd_eq}.
            Take $\beta$ such that
                $\beta$ is a function on $K$ and
                for all $U\in K$ we have $U\mathrel{R_f}\apply{\beta}{U}$
                by \cref{subseteq,times_elem_is_tuple,relation,choice_function}.
                We have $\beta\in A$
                    by \cref{funs_elim,function_apply_elim,subseteq,times_tuple_elim,funs_intro}.
            For all $U\in K$ we have $f(c(U))\in\beta(U)$ by \cref{subseteq,fst_eq,snd_eq}.
            $(\beta,p\circ\beta\circ u)\in h$ by \cref{subseteq}.
            Hence $h(\beta) = p\circ\beta\circ u$ by \cref{function_apply_intro}.
            Let $g = h(\beta)$.
            We have $g\in H$ by \cref{img_iff,function_apply_iff}.
            Therefore $g = p\circ\beta\circ u$ by \cref{real_lt_subst_left}.
            We have $g\in\funs{X}{Y}$ by \cref{subseteq}.

            Show that for all $x\in X$ we have $d((f(x),g(x))) < \delta + \delta$.
            \begin{subproof}
                Fix $x$ such that $x\in X$.
                We have $x\in u(x)$ by \cref{function_apply_elim,fst_eq,snd_eq}.
                We have $u(x)\in K$ and $c(u(x))\in X$ by \cref{funs_type_apply}.
                $u(x)$ is a neighborhood of $c(u(x))$ in $X$ with respect to $T$ and
                    for all $y,h_0$ such that $y\in u(x)$ and $h_0\in F$ we have
                        $d((\apply{h_0}{y},\apply{h_0}{c(u(x))})) < \delta$ by \cref{subseteq}.
                Hence $d((f(x),f(c(u(x))))) < \delta$ by \cref{subseteq}.
                We have $f(c(u(x)))\in\beta(u(x))$ by \cref{subseteq}.
                Hence $\beta(u(x)) = Y\inter\metricBall{d}{Y}{p(\beta(u(x)))}{\delta}$ by \cref{funs_type_apply,subseteq}.
                $f(c(u(x)))\in\metricBall{d}{Y}{p(\beta(u(x)))}{\delta}$
                    by \cref{real_lt_subst_left,inter_elim_right}.
                Hence $d((p(\beta(u(x))),f(c(u(x))))) < \delta$ by \cref{metric_ball}.
                We have $d((f(c(u(x))),p(\beta(u(x))))) < \delta$
                    by \cref{metric_on,metric_symmetric,funs_type_apply,real_lt_subst_left,times_tuple_intro}.
                $\beta$ is a function and $\beta$ is a relation and $\dom{\beta} = K$ and $\ran{\beta}\subseteq C_Y$
                    by \cref{funs_is_function,funs_is_relation,funs_elim,funs_ran,subseteq}.
                $u$ is a function and $u$ is a relation and $\dom{u} = X$ and $\ran{u}\subseteq K$
                    by \cref{funs_is_function,funs_is_relation,funs_elim,funs_ran,subseteq}.
                $p$ is a function and $p$ is a relation and $\dom{p} = C_Y$
                    by \cref{funs_is_function,funs_is_relation,funs_elim}.
                Hence $\ran{\beta}\subseteq\dom{p}$ by \cref{real_lt_subst_right}.
                We have $(p\circ\beta)(u(x)) = p(\beta(u(x)))$ by \cref{circ_apply}.
                $p\circ\beta$ is a function and $p\circ\beta$ is a relation and $\dom{(p\circ\beta)} = K$
                    by \cref{funs_circ,funs_is_function,funs_is_relation,funs_elim}.
                Hence $g(x) = p(\beta(u(x)))$
                    by \cref{circ_apply,real_lt_subst_left,real_lt_subst_right}.
                We have $d((f(c(u(x))),g(x))) < \delta$ by \cref{real_lt_subst_right}.
                Hence $d((f(x),g(x))) < \delta + \delta$
                    by \cref{common_center_small_ball_distance_bound,funs_type_apply}.
            \end{subproof}

            Take $E$ such that
                $E = \{r\in\reals \mid \exists \alpha\in X.
                    ((r = d((\apply{f}{\alpha},\apply{g}{\alpha})) \land d((\apply{f}{\alpha},\apply{g}{\alpha})) < 1) \lor
                    (r = 1 \land 1 \leq d((\apply{f}{\alpha},\apply{g}{\alpha}))))\}$ and
                $\rho((f,g))$ is a supremum of $E$
                by \cref{uniform_metric_on_function_space}.
            We have $\rho((f,g))$ is an upper bound of $E$ by \cref{real_supremum}.
                Show that for all $r$ such that $r\in E$ we have $r \leq \delta + \delta$.
                \begin{subproof}
                    Fix $r$ such that $r\in E$.
                    Take $x\in X$ such that
                        $((r = d((\apply{f}{x},\apply{g}{x})) \land d((\apply{f}{x},\apply{g}{x})) < 1) \lor
                        (r = 1 \land 1 \leq d((\apply{f}{x},\apply{g}{x}))))$
                        by \cref{subseteq}.
                    \begin{byCase}
                    \caseOf{$r = d((\apply{f}{x},\apply{g}{x}))$ and $d((\apply{f}{x},\apply{g}{x})) < 1$.}
                        We have $r \leq \delta + \delta$ by \cref{subseteq,real_lt_subst_left,real_lt_imp_leq}.
                    \caseOf{$r = 1$ and $1 \leq d((\apply{f}{x},\apply{g}{x}))$.}
                            We have $1 < \delta + \delta$
                                by \cref{metric_on,funs_type_apply,times_tuple_intro,real_naturals_subset_reals,real_one_in_naturals,subseteq,real_add_closed,real_leq_split,real_lt_trans,real_lt_subst_left}.
                        Hence $r \leq \delta + \delta$ by \cref{real_lt_subst_left,real_lt_imp_leq}.
                    \end{byCase}
                \end{subproof}
                $E\subseteq\reals$ and $\delta + \delta\in\reals$ by \cref{subseteq,real_add_closed}.
                Hence $\delta + \delta$ is an upper bound of $E$ by \cref{real_upper_bound}.
            Hence $\rho((f,g)) \leq \delta + \delta$ by \cref{real_supremum}.
            Thus $\delta + \delta < \epsilon$
                by \cref{real_add_closed,real_add_ident,real_lt_add,real_add_comm,real_add_assoc,real_lt_trans,real_lt_subst_left,real_lt_subst_right}.
            Therefore $\rho((f,g)) < \epsilon$
                by \cref{metric_on,funs_type_apply,times_tuple_intro,real_leq_split,real_lt_trans,real_lt_subst_left}.
        \end{subproof}

    Follows by definition.
\end{proof}

% from §45 helper lemma 15: finite epsilon-nets imply total boundedness
\begin{lemma}\label{finite_metric_subset_totally_bounded_from_net}
    Assume $\rho$ is a metric on $Z$.
    Assume $F\subseteq Z$.
    Suppose for all $\epsilon\in\reals$ such that $0 < \epsilon$
        there exists $H$ such that
            $H$ is finite and
            $H\subseteq Z$ and
            for all $f\in F$ there exists $g\in H$ such that
                $\rho((f,g)) < \epsilon$.
    Then $F$ is totally bounded in $Z$ with respect to $\rho$.
\end{lemma}
\begin{proof}
    Show that for all $\epsilon\in\reals$ such that $0 < \epsilon$
        there exists $D$ such that
            $D$ is finite and
            $D$ is a covering of $F$ and
            for all $U\in D$ there exists $z\in Z$ such that
                $U = F\inter\metricBall{\rho}{Z}{z}{\epsilon}$.
    \begin{subproof}
        Fix $\epsilon$ such that $\epsilon\in\reals$ and $0 < \epsilon$.
        Take $H$ such that
            $H$ is finite and
            $H\subseteq Z$ and
            for all $f\in F$ there exists $g\in H$ such that
                $\rho((f,g)) < \epsilon$.
        Let $h = \{(g,F\inter\metricBall{\rho}{Z}{g}{\epsilon}) \mid g\in H\}$.
        For all $g,U_1,U_2$ such that $(g,U_1)\in h$ and $(g,U_2)\in h$ we have $U_1 = U_2$
            by \cref{pair_eq_iff}.
        Hence $h$ is a function on $H$
            by \cref{rightunique,relation,subseteq_antisymmetric,subseteq,dom_intro,pair_eq_iff,dom_iff}.
        Let $D = \img{h}{H}$.
        Therefore $D$ is finite by \cref{finite_image_function}.
        For all $f$ such that $f\in F$ we have $f\in\unions{D}$
            by \cref{subseteq,metric_on,metric_symmetric,times_tuple_intro,real_lt_subst_left,metric_ball,inter_intro,img_iff,function_apply_iff,unions_iff}.
        Hence $D$ is a covering of $F$ by \cref{subseteq,covering}.
        For all $U\in D$ there exists $z\in Z$ such that
            $U = F\inter\metricBall{\rho}{Z}{z}{\epsilon}$
            by \cref{img_iff,subseteq,function_apply_intro,real_lt_subst_right}.
        Follows by definition.
    \end{subproof}
    Hence $F$ is totally bounded in $Z$ with respect to $\rho$
        by \cref{totally_bounded_subset_metric_space}.
\end{proof}

% from §45 helper lemma 16: equicontinuous families admit finite uniform epsilon-covers
\begin{lemma}\label{equicontinuous_function_family_uniform_cover}
    Assume $T$ is a topology on $X$.
    Assume $X$ is compact with respect to $T$.
    Assume $d$ is a metric on $Y$.
    Assume $Y$ is compact with respect to $\metricTopology{d}{Y}$.
    Assume $\rho$ is the uniform metric on functions from $X$ to $Y$ with respect to $d$.
    Assume $C$ is the set of continuous functions from $X$ to $Y$ with respect to $T$ and $\metricTopology{d}{Y}$.
    Assume $F\subseteq C$.
    Suppose $F$ is equicontinuous from $X$ to $Y$ with respect to $T$ and $d$.
    Then for all $\epsilon\in\reals$ such that $0 < \epsilon$
        there exists $D$ such that
            $D$ is finite and
            $D$ is a covering of $F$ and
            for all $U\in D$ there exists $g\in\funs{X}{Y}$ such that
                $U = F\inter\metricBall{\rho}{\funs{X}{Y}}{g}{\epsilon}$.
\end{lemma}
\begin{proof}
    Follows by \cref{uniform_metric_on_function_space,subseteq,continuous_function_set,finite_metric_subset_totally_bounded_from_net,equicontinuous_function_family_uniform_net,totally_bounded_subset_metric_space}.
\end{proof}

\begin{lemma}\label{equicontinuous_function_family_totally_bounded_uniform}
    Assume $T$ is a topology on $X$.
    Assume $X$ is compact with respect to $T$.
    Assume $d$ is a metric on $Y$.
    Assume $Y$ is compact with respect to $\metricTopology{d}{Y}$.
    Assume $\rho$ is the uniform metric on functions from $X$ to $Y$ with respect to $d$.
    Assume $C$ is the set of continuous functions from $X$ to $Y$ with respect to $T$ and $\metricTopology{d}{Y}$.
    Assume $F\subseteq C$.
    Suppose $F$ is equicontinuous from $X$ to $Y$ with respect to $T$ and $d$.
    Then $F$ is totally bounded in $\funs{X}{Y}$ with respect to $\rho$.
\end{lemma}
\begin{proof}
    Follows by \cref{equicontinuous_function_family_uniform_cover,uniform_metric_on_function_space,subseteq,continuous_function_set,totally_bounded_subset_metric_space}.
\end{proof}

% from §45 helper lemma 16: a three-step triangle bound
\begin{lemma}\label{metric_three_step_lt_sum}
    Assume $d$ is a metric on $Y$.
    Suppose $x_1,x_2,x_3,x_4\in Y$.
    Suppose $a,b,c\in\reals$.
    Suppose $d((x_1,x_2)) < a$.
    Suppose $d((x_2,x_3)) < b$.
    Suppose $d((x_3,x_4)) < c$.
    Then $d((x_1,x_4)) < a + b + c$.
\end{lemma}
\begin{proof}
    We have $d((x_2,x_4)) < b + c$
        by \cref{real_lt_add_two,metric_on,metric_triangle_inequality,funs_type_apply,times_tuple_intro,real_add_closed,real_leq_split,real_lt_trans,real_lt_subst_left}.
    Hence $d((x_1,x_4)) \leq d((x_1,x_2)) + d((x_2,x_4))$
        by \cref{metric_on,metric_triangle_inequality,times_tuple_intro}.
    Thus $d((x_1,x_2)) + d((x_2,x_4)) < a + (b + c)$
        by \cref{real_lt_add_two,metric_on,funs_type_apply,times_tuple_intro,real_add_closed}.
    Hence $d((x_1,x_4)) < a + b + c$
        by \cref{metric_on,funs_type_apply,times_tuple_intro,real_add_closed,real_add_assoc,real_leq_split,real_lt_trans,real_lt_subst_left}.
\end{proof}

% from §45 helper lemma 17: a pointwise bound gives a sup-metric upper bound
\begin{lemma}\label{sup_metric_upper_bound_from_pointwise}
    Assume $X$ is inhabited.
    Assume $d$ is a metric on $Y$.
    Assume $C\subseteq\funs{X}{Y}$.
    Assume $\sigma$ is the sup metric on $C$ with respect to $X$ and $Y$ and $d$.
    Suppose $f,g\in C$.
    Suppose $M\in\reals$.
    Suppose for all $x\in X$ we have $d((f(x),g(x))) < M$.
    Then $\sigma((f,g)) \leq M$.
\end{lemma}
\begin{proof}
    Follows by \cref{sup_metric_on_bounded_function_space,function_distance_set,subseteq,real_lt_subst_left,real_lt_imp_leq,funs_type_apply,metric_on,times_tuple_intro,real_upper_bound,real_supremum}.
\end{proof}

% from §45 Item 5: equicontinuity implies total boundedness in the sup metric
\begin{lemma}\label{equicontinuous_function_family_totally_bounded_sup}
    Assume $T$ is a topology on $X$.
    Assume $X$ is compact with respect to $T$.
    Assume $X$ is inhabited.
    Assume $d$ is a metric on $Y$.
    Assume $Y$ is compact with respect to $\metricTopology{d}{Y}$.
    Assume $C$ is the set of continuous functions from $X$ to $Y$ with respect to $T$ and $\metricTopology{d}{Y}$.
    Assume $C$ is inhabited.
    Assume $\sigma$ is the sup metric on $C$ with respect to $X$ and $Y$ and $d$.
    Assume $F\subseteq C$.
    Suppose $F$ is equicontinuous from $X$ to $Y$ with respect to $T$ and $d$.
    Then $F$ is totally bounded in $C$ with respect to $\sigma$.
\end{lemma}
\begin{proof}
    $\sigma$ is a metric on $C$ by \cref{sup_metric_on_bounded_function_space}.
    Show that for all $\epsilon\in\reals$ such that $0 < \epsilon$
        there exists $H$ such that
            $H$ is finite and
            $H\subseteq C$ and
            for all $f\in F$ there exists $g\in H$ such that
                $\sigma((f,g)) < \epsilon$.
    \begin{subproof}
        Fix $\epsilon$ such that $\epsilon\in\reals$ and $0 < \epsilon$.
        Take $\delta\in\reals$ such that
            $0 < \delta$ and
            $\delta < 1$ and
            $\delta + \delta + \delta < \epsilon$
            by \cref{positive_real_small_triple_bound}.
        Take $\alpha\in\reals$ such that
            $0 < \alpha$ and
            $\alpha + \alpha = \delta$ and
            $\alpha < \delta$
            by \cref{positive_real_half_decomposition}.
        We have $Y$ is totally bounded with respect to $d$ by \cref{compact_metric_space_totally_bounded}.
        Take $C_Y$ such that
            $C_Y$ is finite and
            $C_Y$ is a covering of $Y$ and
            for all $V\in C_Y$ there exists $y\in Y$ such that
                $V = Y\inter\metricBall{d}{Y}{y}{\alpha}$
            by \cref{totally_bounded_metric_space,totally_bounded_subset_metric_space}.
        Let $R_Y = \{w\in C_Y\times Y \mid \fst{w} = Y\inter\metricBall{d}{Y}{\snd{w}}{\alpha}\}$.
        For all $V\in C_Y$ there exists $y$ such that $V\mathrel{R_Y} y$
            by \cref{subseteq,times_tuple_intro,fst_eq,snd_eq}.
        Take $p$ such that
            $p$ is a function on $C_Y$ and
            for all $V\in C_Y$ we have $V\mathrel{R_Y}\apply{p}{V}$
            by \cref{subseteq,times_elem_is_tuple,relation,choice_function}.
        We have $p\in\funs{C_Y}{Y}$
            by \cref{funs_elim,function_apply_elim,subseteq,times_tuple_elim,funs_intro}.
        For all $V\in C_Y$ we have $V = Y\inter\metricBall{d}{Y}{p(V)}{\alpha}$
            by \cref{subseteq,fst_eq,snd_eq}.

        Let $N = \{U\in T \mid \text{there exists $a\in X$ such that
            $U$ is a neighborhood of $a$ in $X$ with respect to $T$ and
            for all $x,f$ such that $x\in U$ and $f\in F$ we have $d((\apply{f}{x},\apply{f}{a})) < \delta$}\}$.
        For all $x$ such that $x\in X$ we have $x\in\unions{N}$
            by \cref{equicontinuous_family,equicontinuous_at_point,neighborhood,open_in,subseteq,unions_iff}.
        Hence $N$ is a covering of $X$ by \cref{subseteq,covering}.
        For all $U$ such that $U\in N$ we have $U$ is open in $X$ with respect to $T$
            by \cref{subseteq,open_in}.
        Hence $N$ is an open covering of $X$ with respect to $T$ by \cref{open_covering}.
        Take $K$ such that
            $K\subseteq N$ and
            $K$ is finite and
            $K$ is a covering of $X$
            by \cref{compact}.
        Let $R_X = \{w\in K\times X \mid \text{$\fst{w}$ is a neighborhood of $\snd{w}$ in $X$ with respect to $T$ and
            for all $x,f$ such that $x\in\fst{w}$ and $f\in F$ we have
                $d((\apply{f}{x},\apply{f}{\snd{w}})) < \delta$}\}$.
        For all $U\in K$ there exists $a$ such that $U\mathrel{R_X} a$
            by \cref{subseteq,times_tuple_intro,fst_eq,snd_eq}.
        Take $c$ such that
            $c$ is a function on $K$ and
            for all $U\in K$ we have $U\mathrel{R_X}\apply{c}{U}$
            by \cref{subseteq,times_elem_is_tuple,relation,choice_function}.
        We have $c\in\funs{K}{X}$
            by \cref{funs_elim,function_apply_elim,subseteq,times_tuple_elim,funs_intro}.
        For all $U\in K$ we have
            $U$ is a neighborhood of $c(U)$ in $X$ with respect to $T$ and
            for all $x,f$ such that $x\in U$ and $f\in F$ we have
                $d((\apply{f}{x},\apply{f}{c(U)})) < \delta$
            by \cref{subseteq,fst_eq,snd_eq}.
        Let $R_u = \{w\in X\times K \mid \fst{w}\in\snd{w}\}$.
        For all $x\in X$ there exists $U$ such that $x\mathrel{R_u} U$
            by \cref{covering,subseteq,unions_iff,times_tuple_intro,fst_eq,snd_eq}.
        Take $u$ such that
            $u$ is a function on $X$ and
            for all $x\in X$ we have $x\mathrel{R_u}\apply{u}{x}$
            by \cref{subseteq,times_elem_is_tuple,relation,choice_function}.
        We have $u\in\funs{X}{K}$
            by \cref{funs_elim,function_apply_elim,subseteq,times_tuple_elim,funs_intro}.
        Let $A = \funs{K}{C_Y}$.
        We have $A$ is finite by \cref{finite_function_space}.
        Let $A_0 = \{\beta\in A \mid \exists f\in F. \forall U\in K. \apply{f}{\apply{c}{U}}\in\apply{\beta}{U}\}$.
        Hence $A_0$ is finite by \cref{subseteq,subseteq_finite_is_finite}.
        Let $R_q = \{w\in A_0\times F \mid \forall U\in K. \apply{\snd{w}}{\apply{c}{U}}\in\apply{\fst{w}}{U}\}$.
        For all $\beta\in A_0$ there exists $f$ such that $\beta\mathrel{R_q} f$
            by \cref{subseteq,times_tuple_intro,fst_eq,snd_eq}.
        Take $q$ such that
            $q$ is a function on $A_0$ and
            for all $\beta\in A_0$ we have $\beta\mathrel{R_q}\apply{q}{\beta}$
            by \cref{subseteq,times_elem_is_tuple,relation,choice_function}.
        For all $\beta$ such that $\beta\in\dom{q}$ we have $q(\beta)\in F$
            by \cref{funs_elim,function_apply_elim,subseteq,times_tuple_elim}.
        Hence $q\in\funs{A_0}{F}$ by \cref{funs_elim,function_apply_elim,subseteq,times_tuple_elim,funs_intro}.
        For all $\beta\in A_0$ we have
            for all $U\in K$ we have $\apply{\apply{q}{\beta}}{\apply{c}{U}}\in\apply{\beta}{U}$
            by \cref{subseteq,fst_eq,snd_eq}.
        Let $H = \img{q}{A_0}$.
        We have $H$ is finite by \cref{finite_image_function}.
        Hence $H\subseteq C$ by \cref{img_subseteq_ran,funs_ran,subseteq_transitive}.
        Show that for all $f\in F$ there exists $g\in H$ such that
            $\sigma((f,g)) < \epsilon$.
        \begin{subproof}
            Fix $f$ such that $f\in F$.
            We have $f\in\funs{X}{Y}$ by \cref{subseteq,continuous_function_set}.
            Let $R_f = \{w\in K\times C_Y \mid \apply{f}{\apply{c}{\fst{w}}}\in\snd{w}\}$.
            For all $U\in K$ there exists $V$ such that $U\mathrel{R_f} V$
                by \cref{funs_type_apply,covering,subseteq,unions_iff,times_tuple_intro,fst_eq,snd_eq}.
            Take $\beta$ such that
                $\beta$ is a function on $K$ and
                for all $U\in K$ we have $U\mathrel{R_f}\apply{\beta}{U}$
                by \cref{subseteq,times_elem_is_tuple,relation,choice_function}.
            We have $\beta\in A$
                by \cref{funs_elim,function_apply_elim,subseteq,times_tuple_elim,funs_intro}.
            For all $U\in K$ we have $\apply{f}{\apply{c}{U}}\in\apply{\beta}{U}$
                by \cref{subseteq,fst_eq,snd_eq}.
            Hence $\beta\in A_0$ by \cref{subseteq}.
            Let $g = \apply{q}{\beta}$.
            Show that $g\in H$ and $g\in C$ and $g\in\funs{X}{Y}$.
            \begin{subproof}
                Follows by \cref{img_iff,function_apply_iff,real_lt_subst_left,subseteq,continuous_function_set}.
            \end{subproof}

            Show that for all $x\in X$ we have $d((f(x),g(x))) < \delta + \delta + \delta$.
            \begin{subproof}
                Fix $x$ such that $x\in X$.
                We have $x\in u(x)$ by \cref{function_apply_elim,fst_eq,snd_eq}.
                We have $u(x)\in K$ by \cref{funs_type_apply}.
                We have $c(u(x))\in X$ by \cref{funs_type_apply}.
                $u(x)$ is a neighborhood of $c(u(x))$ in $X$ with respect to $T$ and
                    for all $y,h_0$ such that $y\in u(x)$ and $h_0\in F$ we have
                        $d((\apply{h_0}{y},\apply{h_0}{c(u(x))})) < \delta$ by \cref{subseteq}.
                Hence $d((f(x),f(c(u(x))))) < \delta$ by \cref{subseteq}.
                We have $d((g(x),g(c(u(x))))) < \delta$ by \cref{subseteq}.
                Thus $d((g(c(u(x))),g(x))) < \delta$
                    by \cref{metric_on,metric_symmetric,funs_type_apply,real_lt_subst_left,times_tuple_intro}.
                We have $f(c(u(x)))\in\beta(u(x))$ and $g(c(u(x)))\in\beta(u(x))$
                    by \cref{subseteq,real_lt_subst_left}.
                Thus $f(c(u(x)))\in Y\inter\metricBall{d}{Y}{p(\beta(u(x)))}{\alpha}$ and
                    $g(c(u(x)))\in Y\inter\metricBall{d}{Y}{p(\beta(u(x)))}{\alpha}$
                    by \cref{subseteq,funs_type_apply,real_lt_subst_left}.
                Therefore $d((f(c(u(x))),p(\beta(u(x))))) < \alpha$ and
                    $d((g(c(u(x))),p(\beta(u(x))))) < \alpha$
                    by \cref{inter_elim_right,metric_ball,metric_on,metric_symmetric,funs_type_apply,real_lt_subst_left,times_tuple_intro}.
                Hence $d((f(c(u(x))),g(c(u(x))))) < \alpha + \alpha$
                    by \cref{common_center_small_ball_distance_bound,funs_type_apply}.
                Thus $d((f(c(u(x))),g(c(u(x))))) < \delta$ by \cref{real_lt_subst_right}.
                Therefore $d((f(x),g(x))) < \delta + \delta + \delta$
                    by \cref{metric_three_step_lt_sum,funs_type_apply}.
            \end{subproof}
            We have $f\in C$ by \cref{subseteq}.
            $\delta + \delta + \delta\in\reals$ by \cref{real_add_closed}.
            Hence $\sigma((f,g)) \leq \delta + \delta + \delta$
                by \cref{sup_metric_upper_bound_from_pointwise,subseteq,continuous_function_set,real_add_closed}.
            $\sigma((f,g))\in\reals$ by \cref{metric_on,funs_type_apply,times_tuple_intro}.
            Therefore $\sigma((f,g)) < \epsilon$
                by \cref{real_leq_split,real_lt_trans,real_lt_subst_left}.
        \end{subproof}
        Follows by definition.
    \end{subproof}
    Hence $F$ is totally bounded in $C$ with respect to $\sigma$
        by \cref{finite_metric_subset_totally_bounded_from_net}.
\end{proof}

% from §45 helper lemma 18: compact metric subspaces are totally bounded in the ambient metric
\begin{lemma}\label{compact_metric_subspace_ambient_totally_bounded}
    Assume $d$ is a metric on $X$.
    Let $T = \metricTopology{d}{X}$.
    Suppose $A\subseteq X$.
    Suppose $A$ is compact with respect to $\subspaceTopology{T}{A}$.
    Then $A$ is totally bounded in $X$ with respect to $d$.
\end{lemma}
\begin{proof}
    $T$ is a topology on $X$
        by \cref{metric_topology,metric_basis_is_basis,basis_topology_is_topology,basis_for_topology}.
    Show that for all $\epsilon\in\reals$ such that $0 < \epsilon$
        there exists $D$ such that
            $D$ is finite and
            $D$ is a covering of $A$ and
            for all $U\in D$ there exists $x\in X$ such that
                $U = A\inter\metricBall{d}{X}{x}{\epsilon}$.
    \begin{subproof}
        Fix $\epsilon$ such that $\epsilon\in\reals$ and $0 < \epsilon$.
        Let $C = \{\metricBall{d}{X}{x}{\epsilon} \mid x\in A\}$.
        We have $C$ is a covering of $A$
            by \cref{subseteq,metric_ball,metric_on,metric_zero_characterization,real_lt_subst_left,unions_iff,covering}.
        For all $U$ such that $U\in C$ we have $U$ is open in $X$ with respect to $T$
            by \cref{subseteq,metric_ball,metric_basis,powerset_intro,real_lt_subst_left,metric_topology,metric_basis_is_basis,basis_topology_is_topology,basis_for_topology,basis_elem_open_in_generated,open_in}.
        Take $B$ such that
            $B\subseteq C$ and
            $B$ is finite and
            $B$ is a covering of $A$
            by \cref{subspace_of,compact_subspace_open_covering}.
        Let $h = \{(U, A\inter U) \mid U\in B\}$.
        We have $h$ is a function on $B$ by \cref{rightunique,relation,dom_intro,dom_iff,pair_eq_iff,subseteq,subseteq_antisymmetric}.
        Let $D = \img{h}{B}$.
        $D$ is finite by \cref{finite_image_function}.
        $D$ is a covering of $A$
            by \cref{covering,subseteq,unions_iff,inter_intro,img_iff,function_apply_iff}.
        For all $U\in D$ there exists $x\in X$ such that
            $U = A\inter\metricBall{d}{X}{x}{\epsilon}$
            by \cref{img_iff,subseteq,function_apply_intro,real_lt_subst_right}.
        Follows by definition.
    \end{subproof}
    Therefore $A$ is totally bounded in $X$ with respect to $d$
        by \cref{totally_bounded_subset_metric_space}.
\end{proof}

% from §45 helper lemma 19: finite subsets of metric spaces are bounded
\begin{lemma}\label{finite_metric_subset_bounded}
    Assume $d$ is a metric on $X$.
    Suppose $A\subseteq X$.
    Suppose $A$ is finite.
    Then $A$ is bounded in $X$ with respect to $d$.
\end{lemma}
\begin{proof}
    \begin{byCase}
    \caseOf{$A = \emptyset$.}
        Let $M = 0$.
        For all $a_1,a_2\in A$ we have $d((a_1,a_2)) \leq M$
            by \cref{subseteq_emptyset_iff,elem_subseteq,emptyset}.
        Therefore $A$ is bounded in $X$ with respect to $d$ by \cref{real_add_ident,metric_bounded}.
    \caseOf{$A\neq\emptyset$.}
        Take $a_0$ such that $a_0\in A$ by \cref{nonempty_inhabited}.
        Let $f = \{(a,d((a,a_0))) \mid a\in A\}$.
        We have $f$ is a function on $A$ by \cref{rightunique,relation,dom_intro,dom_iff,pair_eq_iff,subseteq,subseteq_antisymmetric}.
        Let $D = \img{f}{A}$.
        $D$ is finite by \cref{finite_image_function}.
        For all $r\in D$ we have $r\in\reals$
            by \cref{img_iff,subseteq,times_tuple_intro,function_apply_intro,metric_on,funs_type_apply,real_lt_subst_left}.
        $D\subseteq\reals$ by \cref{subseteq}.
        We have $D\neq\emptyset$
            by \cref{subseteq,function_apply_iff,img_iff,emptyset}.
        Take $M\in D$ such that $M$ is a maximum of $D$ by \cref{finite_real_maximum}.
        Show that for all $x,y\in A$ we have $d((x,y)) \leq M + M$.
        \begin{subproof}
            Fix $x$.
            Fix $y$ such that $x\in A$ and $y\in A$.
            We have $x\in X$ and $y\in X$ and $a_0\in X$ by \cref{subseteq}.
            Hence $(x,y)\in X\times X$ and $(x,a_0)\in X\times X$ and $(a_0,y)\in X\times X$
                by \cref{times_tuple_intro}.
            Therefore $d((x,y)) \leq d((x,a_0)) + d((a_0,y))$
                by \cref{metric_on,metric_triangle_inequality}.
            Hence $d((x,a_0)) \leq M$ and $d((a_0,y)) \leq M$
                by \cref{subseteq,function_apply_iff,img_iff,function_apply_intro,real_lt_subst_left,real_lt_subst_right,real_maximum,metric_on,metric_symmetric}.
            We have $d((x,a_0))\in\reals$ and $d((a_0,y))\in\reals$ and $d((x,y))\in\reals$ and $M\in\reals$
                by \cref{metric_on,funs_type_apply,real_maximum}.
            Hence $d((x,a_0)) + d((a_0,y)) \leq M + d((x,a_0))$ by \cref{real_leq_add,real_add_comm}.
            $M + d((x,a_0)) \leq M + M$ by \cref{real_leq_add,real_add_comm}.
            We have $d((x,a_0)) + d((a_0,y))\in\reals$ and $M + d((x,a_0))\in\reals$ and $M + M\in\reals$
                by \cref{real_add_closed}.
            Thus $d((x,a_0)) + d((a_0,y)) \leq M + M$ by \cref{real_leq_trans}.
            Therefore $d((x,y)) \leq M + M$ by \cref{real_leq_trans}.
        \end{subproof}
        Therefore $A$ is bounded in $X$ with respect to $d$ by \cref{real_maximum,real_add_closed,metric_bounded}.
    \end{byCase}
\end{proof}

% from §45 helper lemma 20: closure of bounded metric subsets is bounded
\begin{lemma}\label{closure_metric_bounded}
    Assume $d$ is a metric on $X$.
    Let $T = \metricTopology{d}{X}$.
    Suppose $A\subseteq X$.
    Suppose $A$ is bounded in $X$ with respect to $d$.
    Let $B = \closure{A}{X}{T}$.
    Then $B$ is bounded in $X$ with respect to $d$.
\end{lemma}
\begin{proof}
    $T$ is a topology on $X$
        by \cref{metric_topology,metric_basis_is_basis,basis_topology_is_topology,basis_for_topology}.
    Hence $B\subseteq X$ by \cref{closure_closed_in,closed_in,subseteq}.
    Take $M\in\reals$ such that for all $a_1,a_2\in A$ we have $d((a_1,a_2)) \leq M$
        by \cref{metric_bounded}.
    Show that for all $x,z\in B$ we have $d((x,z)) \leq 1 + (M + 1) + 1$.
    \begin{subproof}
        Fix $x$.
        Fix $z$ such that $x\in B$ and $z\in B$.
        We have $x\in X$ and $z\in X$ by \cref{subseteq}.
        $1\in\reals$ and $0 < 1$
            by \cref{real_naturals_subset_reals,real_one_in_naturals,subseteq,real_zero_lt_one}.
        Hence $\metricBall{d}{X}{x}{1}$ is open in $X$ with respect to $T$
            by \cref{metric_ball,subseteq,metric_basis,metric_basis_is_basis,basis_elem_open_in_generated,metric_topology,basis_for_topology,powerset_intro,open_in}.
        Hence $\metricBall{d}{X}{x}{1}$ intersects $A$
            by \cref{metric_ball,metric_on,metric_zero_characterization,real_lt_subst_left,closure_iff_intersects_open}.
        Take $a_1$ such that $a_1\in\metricBall{d}{X}{x}{1}\inter A$ by \cref{intersects,nonempty_inhabited}.
        We have $\metricBall{d}{X}{z}{1}$ is open in $X$ with respect to $T$
            by \cref{metric_ball,subseteq,metric_basis,metric_basis_is_basis,basis_elem_open_in_generated,metric_topology,basis_for_topology,powerset_intro,open_in,real_lt_subst_left}.
        Hence $\metricBall{d}{X}{z}{1}$ intersects $A$
            by \cref{metric_ball,metric_on,metric_zero_characterization,real_lt_subst_left,closure_iff_intersects_open}.
        Take $a_2$ such that $a_2\in\metricBall{d}{X}{z}{1}\inter A$ by \cref{intersects,nonempty_inhabited}.
        We have $a_1\in A$ and $a_1\in X$ and $d((x,a_1)) < 1$
            by \cref{inter_elim_left,inter_elim_right,metric_ball}.
        We have $a_2\in A$ and $a_2\in X$ and $d((z,a_2)) < 1$
            by \cref{inter_elim_left,inter_elim_right,metric_ball}.
        We have $d((a_1,a_2)) \leq M$.
        Hence $d((a_2,z)) < 1$ by \cref{metric_on,metric_symmetric,real_lt_subst_left,times_tuple_intro}.
        We have $(x,a_1)\in X\times X$ and $(a_1,a_2)\in X\times X$ and $(a_2,z)\in X\times X$
            by \cref{times_tuple_intro}.
        We have $d((x,a_1))\in\reals$ and $d((a_1,a_2))\in\reals$ and $d((a_2,z))\in\reals$ and
            $M\in\reals$ by \cref{metric_on,funs_type_apply}.
        We have $M + 1\in\reals$ by \cref{real_add_closed}.
        Therefore $M < M + 1$ by \cref{real_lt_add,real_add_ident,real_add_comm,real_lt_subst_left}.
        Thus $d((a_1,a_2)) < M + 1$ by \cref{real_leq_split,real_lt_trans,real_lt_subst_left}.
        Hence $d((x,z)) < 1 + (M + 1) + 1$
            by \cref{metric_three_step_lt_sum}.
        Therefore $d((x,z)) \leq 1 + (M + 1) + 1$ by \cref{real_lt_imp_leq}.
    \end{subproof}
    Therefore $B$ is bounded in $X$ with respect to $d$
        by \cref{real_add_closed,real_naturals_subset_reals,real_one_in_naturals,subseteq,metric_bounded}.
\end{proof}

% from §45 helper lemma 22: bounded subsets of euclidean spaces admit finite internal epsilon-nets
\begin{lemma}\label{bounded_reals_power_subset_finite_internal_net}
    Assume $n\in\naturalsPlus$.
    Let $J = \initialSegment{n}$.
    Let $Y = \realsPower{J}$.
    Assume $d$ is a metric on $Y$.
    Assume $d = \squareMetric{n}$ or $d = \euclideanMetric{n}$.
    Suppose $A\subseteq Y$.
    Suppose $A$ is bounded in $Y$ with respect to $d$.
    Suppose $\epsilon\in\reals$.
    Suppose $0 < \epsilon$.
    Then there exists $H$ such that
        $H$ is finite and
        $H\subseteq Y$ and
        for all $a\in A$ there exists $h\in H$ such that
            $d((a,h)) < \epsilon$.
\end{lemma}
\begin{proof}
    Let $S = \metricTopology{d}{Y}$.
    Let $B = \closure{A}{Y}{S}$.
    $S$ is a topology on $Y$
        by \cref{metric_topology,metric_basis_is_basis,basis_for_topology,basis_topology_is_topology}.
    We have $B\subseteq Y$ by \cref{closure,subseteq}.
    We have $B$ is bounded in $Y$ with respect to $d$ by \cref{closure_metric_bounded}.
    Let $d_0 = \squareMetric{n}$.
    We have $S = \metricTopology{d_0}{Y}$
        by \cref{square_metric_topology_equals_euclidean,real_lt_subst_left}.
    Hence $B$ is closed in $Y$ with respect to $\metricTopology{d_0}{Y}$
        by \cref{closure_closed_in,real_lt_subst_right}.
    We have $B$ is bounded in $Y$ with respect to $\euclideanMetric{n}$ or
        $B$ is bounded in $Y$ with respect to $d_0$ by \cref{real_lt_subst_right}.
    Thus $B$ is compact with respect to $\subspaceTopology{\metricTopology{d_0}{Y}}{B}$
        by \cref{compact_reals_power_iff_closed_bounded}.
    Hence $B$ is totally bounded in $Y$ with respect to $d$
        by \cref{real_lt_subst_left,compact_metric_subspace_ambient_totally_bounded}.
    We have $A\subseteq B$ by \cref{subseteq_closure}.
    \begin{byCase}
    \caseOf{$B = \emptyset$.}
        Show that there exists $H$ such that
            $H$ is finite and
            $H\subseteq Y$ and
            for all $a\in A$ there exists $h\in H$ such that
                $d((a,h)) < \epsilon$.
        \begin{subproof}
            Follows by \cref{subseteq_emptyset_iff,subseteq_transitive,emptyset_is_finite,emptyset_subseteq,subseteq,real_lt_subst_left,emptyset}.
        \end{subproof}
        Follows by definition.
    \caseOf{$B\neq\emptyset$.}
        Show that there exists $H$ such that
            $H$ is finite and
            $H\subseteq Y$ and
            for all $a\in A$ there exists $h\in H$ such that
                $d((a,h)) < \epsilon$.
        \begin{subproof}
            Follows by \cref{nonempty_inhabited,totally_bounded_subset_finite_internal_net,subseteq_transitive,subseteq}.
        \end{subproof}
        Follows by definition.
    \end{byCase}
\end{proof}

% from §45 helper lemma 21: equicontinuous pointwise bounded families are totally bounded in the uniform metric
\begin{lemma}\label{equicontinuous_pointwise_bounded_function_family_totally_bounded_uniform}
    Assume $n\in\naturalsPlus$.
    Let $J = \initialSegment{n}$.
    Let $Y = \realsPower{J}$.
    Assume $d$ is a metric on $Y$.
    Assume $d = \squareMetric{n}$ or $d = \euclideanMetric{n}$.
    Assume $T$ is a topology on $X$.
    Assume $X$ is compact with respect to $T$.
    Assume $\rho$ is the uniform metric on functions from $X$ to $Y$ with respect to $d$.
    Assume $C$ is the set of continuous functions from $X$ to $Y$ with respect to $T$ and $\metricTopology{d}{Y}$.
    Assume $F\subseteq C$.
    Suppose $F$ is equicontinuous from $X$ to $Y$ with respect to $T$ and $d$.
    Suppose for all $x\in X$ there exists $A$ such that
        $A = \{\apply{f}{x} \mid f\in F\}$ and
        $A$ is bounded in $Y$ with respect to $d$.
    Then $F$ is totally bounded in $\funs{X}{Y}$ with respect to $\rho$.
\end{lemma}
\begin{proof}
    Let $Z = \funs{X}{Y}$.
    $\rho$ is a metric on $Z$ by \cref{uniform_metric_on_function_space}.
    $F\subseteq Z$ by \cref{subseteq,continuous_function_set}.
    Show that for all $\epsilon\in\reals$ such that $0 < \epsilon$
        there exists $H$ such that
            $H$ is finite and
            $H\subseteq Z$ and
            for all $f\in F$ there exists $g\in H$ such that
                $\rho((f,g)) < \epsilon$.
    \begin{subproof}
        Fix $\epsilon$ such that $\epsilon\in\reals$ and $0 < \epsilon$.
        Take $\delta\in\reals$ such that
            $0 < \delta$ and
            $\delta < 1$ and
            $\delta + \delta + \delta < \epsilon$
            by \cref{positive_real_small_triple_bound}.

        Let $N = \{U\in T \mid \text{there exists $a\in X$ such that
            $U$ is a neighborhood of $a$ in $X$ with respect to $T$ and
            for all $x,f$ such that $x\in U$ and $f\in F$ we have $d((\apply{f}{x},\apply{f}{a})) < \delta$}\}$.
        For all $x$ such that $x\in X$ there exists $U$ such that
            $U$ is a neighborhood of $x$ in $X$ with respect to $T$ and
            for all $y,f$ such that $y\in U$ and $f\in F$ we have
                $d((\apply{f}{y},\apply{f}{x})) < \delta$
            by \cref{equicontinuous_family,equicontinuous_at_point}.
        Hence $N$ is a covering of $X$
            by \cref{neighborhood,open_in,subseteq,unions_iff,covering}.
        For all $U$ such that $U\in N$ we have $U$ is open in $X$ with respect to $T$
            by \cref{subseteq,open_in}.
        Therefore $N$ is an open covering of $X$ with respect to $T$ by \cref{open_covering}.
        Take $K$ such that
            $K\subseteq N$ and
            $K$ is finite and
            $K$ is a covering of $X$
            by \cref{compact}.

        Let $R_c = \{w\in K\times X \mid \text{$\fst{w}$ is a neighborhood of $\snd{w}$ in $X$ with respect to $T$ and
            for all $x,f$ such that $x\in\fst{w}$ and $f\in F$ we have
                $d((\apply{f}{x},\apply{f}{\snd{w}})) < \delta$}\}$.
        For all $U\in K$ there exists $a$ such that $U\mathrel{R_c} a$
            by \cref{subseteq,times_tuple_intro,fst_eq,snd_eq}.
        Take $c$ such that
            $c$ is a function on $K$ and
            for all $U\in K$ we have $U\mathrel{R_c}\apply{c}{U}$
            by \cref{subseteq,times_elem_is_tuple,relation,choice_function}.
        We have $c\in\funs{K}{X}$ by \cref{funs_elim,function_apply_elim,subseteq,times_tuple_elim,funs_intro}.
        For all $U\in K$ we have
            $U$ is a neighborhood of $c(U)$ in $X$ with respect to $T$ and
            for all $x,f$ such that $x\in U$ and $f\in F$ we have
                $d((\apply{f}{x},\apply{f}{c(U)})) < \delta$
            by \cref{subseteq,fst_eq,snd_eq}.

        For all $U\in K$ there exists $A$ such that
            $A$ is finite and
            $A\subseteq Y$ and
            for all $f\in F$ there exists $y\in A$ such that
                $d((\apply{f}{c(U)},y)) < \delta$
            by \cref{funs_type_apply,subseteq,metric_bounded,bounded_reals_power_subset_finite_internal_net}.

        Let $R_m = \{w\in K\times\pow{Y} \mid \text{$\snd{w}$ is finite and
            $\snd{w}\subseteq Y$ and
            for all $f\in F$ there exists $y\in\snd{w}$ such that
                $d((\apply{f}{c(\fst{w})},y)) < \delta$}\}$.
        For all $U\in K$ there exists $A$ such that $U\mathrel{R_m} A$
            by \cref{subseteq,times_tuple_intro,powerset_intro,fst_eq,snd_eq}.
        Take $m$ such that
            $m$ is a function on $K$ and
            for all $U\in K$ we have $U\mathrel{R_m}\apply{m}{U}$
            by \cref{subseteq,times_elem_is_tuple,relation,choice_function}.
        For all $U\in K$ we have
            $m(U)$ is finite and
            $m(U)\subseteq Y$ and
            for all $f\in F$ there exists $y\in m(U)$ such that
                $d((\apply{f}{c(U)},y)) < \delta$
            by \cref{subseteq,fst_eq,snd_eq}.

        Let $M_0 = \img{m}{K}$.
        We have $M_0$ is finite by \cref{finite_image_function}.
        For all $A$ such that $A\in M_0$ we have $A$ is finite
            by \cref{img_iff,function_apply_intro,subseteq}.
        Let $D = \unions{M_0}$.
        Hence $D$ is finite by \cref{finite_union_of_finite_family}.
        $D\subseteq Y$ by \cref{subseteq,unions_iff,img_iff,function_apply_intro}.
        For all $U\in K$ we have $m(U)\subseteq D$
            by \cref{subseteq,img_iff,function_apply_iff,unions_iff}.

        Let $R_u = \{w\in X\times K \mid \fst{w}\in\snd{w}\}$.
        For all $x\in X$ there exists $U$ such that $x\mathrel{R_u} U$
            by \cref{covering,subseteq,unions_iff,times_tuple_intro,fst_eq,snd_eq}.
        Take $u$ such that
            $u$ is a function on $X$ and
            for all $x\in X$ we have $x\mathrel{R_u}\apply{u}{x}$
            by \cref{subseteq,times_elem_is_tuple,relation,choice_function}.
        Hence $u\in\funs{X}{K}$ by \cref{funs_elim,function_apply_elim,subseteq,times_tuple_elim,funs_intro}.

        Let $A = \funs{K}{D}$.
        We have $A$ is finite by \cref{finite_function_space}.
        Let $h = \{(\alpha,\alpha\circ u) \mid \alpha\in A\}$.
        We have $h$ is a function on $A$ by \cref{rightunique,relation,dom_intro,dom_iff,pair_eq_iff,subseteq,subseteq_antisymmetric}.
        Let $H = \img{h}{A}$.
        Therefore $H$ is finite by \cref{finite_image_function}.
        $H\subseteq Z$
            by \cref{subseteq,img_iff,function_apply_intro,real_lt_subst_left,funs_circ,funs_weaken_codom}.

        Show that for all $f\in F$ there exists $g\in H$ such that $\rho((f,g)) < \epsilon$.
        \begin{subproof}
            Fix $f$ such that $f\in F$.
            We have $f\in Z$ by \cref{subseteq,continuous_function_set}.
            Let $R_f = \{w\in K\times D \mid \snd{w}\in m(\fst{w}) \land d((\apply{f}{c(\fst{w})},\snd{w})) < \delta\}$.
            For all $U\in K$ there exists $y$ such that $U\mathrel{R_f} y$
                by \cref{subseteq,times_tuple_intro,fst_eq,snd_eq}.
            Take $\beta$ such that
                $\beta$ is a function on $K$ and
                for all $U\in K$ we have $U\mathrel{R_f}\apply{\beta}{U}$
                by \cref{subseteq,times_elem_is_tuple,relation,choice_function}.
            We have $\beta\in A$
                by \cref{funs_elim,function_apply_elim,subseteq,times_tuple_elim,funs_intro}.
            For all $U\in K$ we have
                $\beta(U)\in m(U)$ and
                $d((\apply{f}{c(U)},\beta(U))) < \delta$
                by \cref{subseteq,fst_eq,snd_eq}.
            Let $g = h(\beta)$.
            Show that $g\in H$ and $g = \beta\circ u$.
            \begin{subproof}
                Follows by \cref{subseteq,function_apply_intro,img_iff,function_apply_iff,real_lt_subst_left}.
            \end{subproof}
            We have $g\in Z$ by \cref{subseteq}.

            Show that for all $x\in X$ we have $d((f(x),g(x))) < \delta + \delta$.
            \begin{subproof}
                Fix $x$ such that $x\in X$.
                We have $x\in u(x)$ by \cref{function_apply_elim,fst_eq,snd_eq}.
                We have $u(x)\in K$ by \cref{funs_type_apply}.
                We have $u(x)$ is a neighborhood of $c(u(x))$ in $X$ with respect to $T$ and
                    for all $y,h_0$ such that $y\in u(x)$ and $h_0\in F$ we have
                        $d((\apply{h_0}{y},\apply{h_0}{c(u(x))})) < \delta$ by \cref{subseteq}.
                Hence $d((f(x),f(c(u(x))))) < \delta$ by \cref{subseteq}.
                We have $\beta(u(x))\in m(u(x))$ and $d((f(c(u(x))),\beta(u(x)))) < \delta$ by \cref{subseteq}.
                We have $g(x) = \beta(u(x))$
                    by \cref{circ_apply,real_lt_subst_left,funs_is_function,funs_is_relation,funs_elim,funs_ran,subseteq}.
                Therefore $d((f(c(u(x))),g(x))) < \delta$ by \cref{real_lt_subst_right}.
                We have $f(x)\in Y$ and $f(c(u(x)))\in Y$ and $g(x)\in Y$
                    by \cref{funs_type_apply,subseteq}.
                Hence $d((g(x),f(c(u(x))))) < \delta$
                    by \cref{metric_on,metric_symmetric,real_lt_subst_left,times_tuple_intro}.
                Therefore $d((f(x),g(x))) < \delta + \delta$
                    by \cref{common_center_small_ball_distance_bound}.
            \end{subproof}

            Let $e = \boundedMetric{d}$.
            Show that for all $x\in X$ we have $e((f(x),g(x))) \leq \delta + \delta$.
            \begin{subproof}
                Fix $x$ such that $x\in X$.
                We have $f(x)\in Y$ and $g(x)\in Y$ by \cref{subseteq,funs_type_apply}.
                Let $p = (f(x),g(x))$.
                We have $p\in Y\times Y$ by \cref{times_tuple_intro}.
                Hence $d(p) < \delta + \delta$ by \cref{subseteq,real_lt_subst_left}.
                $e$ is a metric on $Y$ by \cref{bounded_metric_is_metric}.
                We have $e\in\funs{Y\times Y}{\reals}$ by \cref{metric_on}.
                \begin{byCase}
                \caseOf{$\delta + \delta < 1$.}
                    We have $d(p)\in\reals$ and $\delta + \delta\in\reals$ and $1\in\reals$
                        by \cref{metric_on,funs_type_apply,real_add_closed,real_one_in_naturals,real_naturals_subset_reals,subseteq}.
                    Hence $d(p) < 1$ by \cref{real_lt_trans}.
                    Therefore $e(p) = d(p)$ by \cref{bounded_metric_apply_below_one}.
                    Hence $e((f(x),g(x))) \leq \delta + \delta$ by \cref{real_lt_imp_leq,real_lt_subst_left}.
                \caseOf{it is wrong that $\delta + \delta < 1$.}
                    We have $\delta + \delta\in\reals$ and $1\in\reals$
                        by \cref{real_add_closed,real_one_in_naturals,real_naturals_subset_reals,subseteq}.
                    Hence $1 \leq \delta + \delta$ by \cref{real_lt_trichotomy}.
                    Therefore $e(p) \leq 1$ by \cref{bounded_metric_leq_one}.
                    We have $e(p)\in\reals$ by \cref{metric_on,funs_type_apply}.
                    Thus $e((f(x),g(x))) \leq \delta + \delta$ by \cref{real_leq_trans,real_lt_subst_left}.
                \end{byCase}
            \end{subproof}
            We have $\delta + \delta\in\reals$ by \cref{real_add_closed}.
            Hence $\rho((f,g)) \leq \delta + \delta$
                by \cref{uniform_metric_leq_from_coordinate_bounds}.
            We have $\delta + \delta < \delta + \delta + \delta$
                by \cref{real_lt_add,real_add_ident,real_add_assoc,real_add_comm,real_lt_subst_left,real_lt_subst_right}.
            We have $\delta + \delta\in\reals$ and $\delta + \delta + \delta\in\reals$ and $\epsilon\in\reals$
                by \cref{real_add_closed}.
            Hence $\delta + \delta < \epsilon$ by \cref{real_lt_trans}.
            We have $\rho((f,g))\in\reals$ by \cref{metric_on,funs_type_apply,times_tuple_intro}.
            Therefore $\rho((f,g)) < \epsilon$
                by \cref{real_leq_split,real_lt_trans,real_lt_subst_left}.
        \end{subproof}

        Show that $H$ is finite and
            $H\subseteq Z$ and
            for all $f\in F$ there exists $g\in H$ such that
                $\rho((f,g)) < \epsilon$.
        \begin{subproof}
            Trivial.
        \end{subproof}
        Follows by definition.
    \end{subproof}
    Hence $F$ is totally bounded in $Z$ with respect to $\rho$
        by \cref{finite_metric_subset_totally_bounded_from_net}.
\end{proof}

% from §45 Item 7: compact closure implies equicontinuity
\begin{theorem}\label{ascoli_compact_closure_implies_equicontinuous}
    Assume $n\in\naturalsPlus$.
    Let $J = \initialSegment{n}$.
    Let $Y = \realsPower{J}$.
    Assume $d$ is a metric on $Y$.
    Assume $d = \squareMetric{n}$ or $d = \euclideanMetric{n}$.
    Assume $T$ is a topology on $X$.
    Assume $X$ is compact with respect to $T$.
    Assume $\rho$ is the uniform metric on functions from $X$ to $Y$ with respect to $d$.
    Assume $C$ is the set of continuous functions from $X$ to $Y$ with respect to $T$ and $\metricTopology{d}{Y}$.
    Assume $F\subseteq C$.
    Let $G = \closure{F}{\funs{X}{Y}}{\metricTopology{\rho}{\funs{X}{Y}}}$.
    Suppose $G$ is compact with respect to $\subspaceTopology{\metricTopology{\rho}{\funs{X}{Y}}}{G}$.
    Then $F$ is equicontinuous from $X$ to $Y$ with respect to $T$ and $d$.
\end{theorem}
\begin{proof}
    Let $Z = \funs{X}{Y}$.
    Let $T_Z = \metricTopology{\rho}{Z}$.
    $G\subseteq Z$ by \cref{closure,subseteq}.
    $T_Z$ is a topology on $Z$
        by \cref{metric_topology,metric_basis_is_basis,basis_topology_is_topology,basis_for_topology,uniform_metric_on_function_space}.
    $C$ is closed in $Z$ with respect to $T_Z$ by \cref{continuous_function_space_closed_uniform}.
    $F\subseteq Z$ by \cref{subseteq,continuous_function_set}.
    We have $G\subseteq C$ by \cref{closure_subseteq_closed,subseteq}.
        Hence $G$ is totally bounded in $Z$ with respect to $\rho$
        by \cref{closure,subseteq,uniform_metric_on_function_space,compact_metric_subspace_ambient_totally_bounded}.
    Therefore $G$ is equicontinuous from $X$ to $Y$ with respect to $T$ and $d$
        by \cref{totally_bounded_function_family_equicontinuous}.
    Hence $F$ is equicontinuous from $X$ to $Y$ with respect to $T$ and $d$
        by \cref{equicontinuous_family,equicontinuous_at_point,subseteq_closure,subseteq}.
\end{proof}

% from §45 Item 6: pointwise bounded family
% from §45 Item 7: compact closure implies pointwise boundedness
\begin{theorem}\label{ascoli_compact_closure_implies_pointwise_bounded}
    Assume $n\in\naturalsPlus$.
    Let $J = \initialSegment{n}$.
    Let $Y = \realsPower{J}$.
    Assume $d$ is a metric on $Y$.
    Assume $d = \squareMetric{n}$ or $d = \euclideanMetric{n}$.
    Assume $T$ is a topology on $X$.
    Assume $X$ is compact with respect to $T$.
    Assume $\rho$ is the uniform metric on functions from $X$ to $Y$ with respect to $d$.
    Assume $C$ is the set of continuous functions from $X$ to $Y$ with respect to $T$ and $\metricTopology{d}{Y}$.
    Assume $F\subseteq C$.
    Let $G = \closure{F}{\funs{X}{Y}}{\metricTopology{\rho}{\funs{X}{Y}}}$.
    Suppose $G$ is compact with respect to $\subspaceTopology{\metricTopology{\rho}{\funs{X}{Y}}}{G}$.
    Then for all $x\in X$ there exists $A$ such that
        $A = \{\apply{f}{x} \mid f\in F\}$ and
        $A$ is bounded in $Y$ with respect to $d$.
\end{theorem}
\begin{proof}
    Let $Z = \funs{X}{Y}$.
    Let $T_Z = \metricTopology{\rho}{Z}$.
    $T_Z$ is a topology on $Z$
        by \cref{metric_topology,metric_basis_is_basis,basis_topology_is_topology,basis_for_topology,uniform_metric_on_function_space}.
    $F\subseteq Z$ by \cref{subseteq,continuous_function_set}.
    $G\subseteq Z$ by \cref{closure,subseteq}.
    $\rho$ is a metric on $Z$ by \cref{uniform_metric_on_function_space}.
    We have $G$ is totally bounded in $Z$ with respect to $\rho$
        by \cref{real_lt_subst_left,real_lt_subst_right,compact_metric_subspace_ambient_totally_bounded}.
    Show that for all $x\in X$ there exists $A$ such that
        $A = \{\apply{f}{x} \mid f\in F\}$ and
        $A$ is bounded in $Y$ with respect to $d$.
    \begin{subproof}
        Fix $x$ such that $x\in X$.
        Let $A = \{\apply{f}{x} \mid f\in F\}$.
        \begin{byCase}
        \caseOf{$F = \emptyset$.}
            We have $A$ is bounded in $Y$ with respect to $d$
                by \cref{subseteq,emptyset,emptyset_subseteq,subseteq_antisymmetric,subseteq_emptyset_iff,emptyset_is_finite,finite_metric_subset_bounded}.
            Follows by definition.
        \caseOf{$F\neq\emptyset$.}
            $1\in\reals$ and $0 < 1$
                by \cref{real_naturals_subset_reals,real_one_in_naturals,subseteq,real_zero_lt_one}.
            Take $\alpha\in\reals$ such that
                $0 < \alpha$ and
                $\alpha + \alpha = 1$ and
                $\alpha < 1$
                by \cref{positive_real_half_decomposition}.
            Take $H$ such that
                $H\subseteq G$ and
                $H$ is finite and
                for all $g\in G$ there exists $h\in H$ such that
                    $\rho((g,h)) < \alpha$
                by \cref{subseteq_closure,subseteq,nonempty_inhabited,totally_bounded_subset_finite_internal_net}.
            Let $e = \{(h,\apply{h}{x}) \mid h\in H\}$.
            We have $e$ is a function on $H$ by \cref{rightunique,relation,dom_intro,dom_iff,pair_eq_iff,subseteq,subseteq_antisymmetric}.
            Let $B = \img{e}{H}$.
            $B$ is finite by \cref{finite_image_function}.
            Therefore $B$ is bounded in $Y$ with respect to $d$
                by \cref{img_iff,subseteq,closure,funs_type_apply,function_apply_intro,real_lt_subst_left,finite_metric_subset_bounded}.
            Take $M\in\reals$ such that for all $b_1,b_2\in B$ we have $d((b_1,b_2)) \leq M$
                by \cref{metric_bounded}.
            Show that for all $y_1,y_2\in A$ we have $d((y_1,y_2)) \leq 1 + (M + 1)$.
            \begin{subproof}
                Fix $y_1$.
                Fix $y_2$ such that $y_1\in A$ and $y_2\in A$.
                Take $f_1\in F$ such that $y_1 = f_1(x)$ by \cref{subseteq}.
                Take $f_2\in F$ such that $y_2 = f_2(x)$ by \cref{subseteq}.
                We have $f_1\in G$ and $f_2\in G$ by \cref{subseteq_closure,subseteq}.
                Take $h_1\in H$ such that $\rho((f_1,h_1)) < \alpha$ by \cref{subseteq}.
                Take $h_2\in H$ such that $\rho((f_2,h_2)) < \alpha$ by \cref{subseteq}.
                We have $f_1\in Z$ and $f_2\in Z$ and $h_1\in Z$ and $h_2\in Z$ by \cref{subseteq,closure}.
                Hence $d((f_1(x),h_1(x))) < \alpha$ and $d((f_2(x),h_2(x))) < \alpha$
                    by \cref{uniform_metric_small_implies_pointwise_small}.
                Therefore $d((y_1,h_1(x))) < \alpha$ and $d((y_2,h_2(x))) < \alpha$
                    by \cref{real_lt_subst_left,real_lt_subst_right}.
                Hence $d((h_2(x),y_2)) < \alpha$ by \cref{metric_on,metric_symmetric,real_lt_subst_left,funs_type_apply}.
                Thus $h_1(x)\in B$ and $h_2(x)\in B$
                    by \cref{subseteq,function_apply_iff,img_iff,function_apply_intro,real_lt_subst_left,real_lt_subst_right}.
                Therefore $d((h_1(x),h_2(x))) \leq M$.
                We have $y_1\in Y$ and $y_2\in Y$ and $h_1(x)\in Y$ and $h_2(x)\in Y$ by \cref{subseteq,funs_type_apply}.
                We have $(y_1,h_1(x))\in Y\times Y$ and $(h_1(x),h_2(x))\in Y\times Y$ and
                    $(h_2(x),y_2)\in Y\times Y$ by \cref{times_tuple_intro}.
                We have $d((y_1,h_1(x)))\in\reals$ and $d((h_1(x),h_2(x)))\in\reals$ and
                    $d((h_2(x),y_2))\in\reals$ by \cref{metric_on,funs_type_apply}.
                $1\in\reals$ by \cref{real_naturals_subset_reals,real_one_in_naturals,subseteq}.
                We have $M + 1\in\reals$ by \cref{real_add_closed}.
                Therefore $M < M + 1$ by \cref{real_lt_add,real_add_ident,real_add_comm,real_lt_subst_left}.
                Thus $d((h_1(x),h_2(x))) < M + 1$ by \cref{real_leq_split,real_lt_trans,real_lt_subst_left}.
                We have $d((y_1,h_1(x))) < 1$ by \cref{real_lt_trans}.
                Therefore $d((h_2(x),y_2)) < 1$ by \cref{real_lt_trans}.
                Thus $d((y_1,y_2)) < \alpha + (M + 1) + \alpha$ by \cref{metric_three_step_lt_sum}.
                Hence $d((y_1,y_2)) < 1 + (M + 1)$
                    by \cref{real_lt_subst_right,real_add_assoc,real_add_comm,real_lt_subst_left}.
                Therefore $d((y_1,y_2)) \leq 1 + (M + 1)$ by \cref{real_lt_imp_leq}.
            \end{subproof}
            Therefore $A$ is bounded in $Y$ with respect to $d$
                by \cref{real_add_closed,real_naturals_subset_reals,real_one_in_naturals,subseteq,continuous_function_set,funs_type_apply,real_lt_subst_left,metric_bounded}.
            Follows by definition.
        \end{byCase}
    \end{subproof}
\end{proof}

% from §45 helper lemma 23: closures of totally bounded subsets are totally bounded
\begin{lemma}\label{closure_totally_bounded}
    Assume $d$ is a metric on $X$.
    Let $T = \metricTopology{d}{X}$.
    Suppose $A\subseteq X$.
    Suppose $A$ is totally bounded in $X$ with respect to $d$.
    Let $B = \closure{A}{X}{T}$.
    Then $B$ is totally bounded in $X$ with respect to $d$.
\end{lemma}
\begin{proof}
    $T$ is a topology on $X$
        by \cref{metric_topology,metric_basis_is_basis,basis_topology_is_topology,basis_for_topology}.
    We have $B\subseteq X$ by \cref{closure,subseteq}.
    Show that for all $\epsilon\in\reals$ such that $0 < \epsilon$
        there exists $H$ such that
            $H$ is finite and
            $H\subseteq X$ and
            for all $b\in B$ there exists $h\in H$ such that
                $d((b,h)) < \epsilon$.
    \begin{subproof}
        Fix $\epsilon$ such that $\epsilon\in\reals$ and $0 < \epsilon$.
        \begin{byCase}
        \caseOf{$A = \emptyset$.}
            There exists $H$ such that
                $H$ is finite and
                $H\subseteq X$ and
                for all $b\in B$ there exists $h\in H$ such that $d((b,h)) < \epsilon$
                by \cref{closed_emptyset,real_lt_subst_left,closed_eq_closure,real_lt_subst_right,emptyset_is_finite,emptyset_subseteq,subseteq,emptyset,subseteq_emptyset_iff}.
            Follows by definition.
        \caseOf{$A\neq\emptyset$.}
            Let $\delta = \epsilon / (1 + 1)$.
            We have $\delta\in\reals$ and $0 < \delta$ and $\delta + \delta = \epsilon$
                by \cref{real_half_of_positive}.
            Take $H$ such that
                $H\subseteq A$ and
                $H$ is finite and
                for all $a\in A$ there exists $h\in H$ such that
                    $d((a,h)) < \delta$
                by \cref{nonempty_inhabited,totally_bounded_subset_finite_internal_net}.
            Show that for all $b\in B$ there exists $h\in H$ such that
                $d((b,h)) < \epsilon$.
            \begin{subproof}
                Fix $b$ such that $b\in B$.
                We have $b\in X$ by \cref{closure,subseteq}.
                Hence $\metricBall{d}{X}{b}{\delta}$ is open in $X$ with respect to $T$
                    by \cref{metric_ball,subseteq,metric_basis,metric_basis_is_basis,basis_elem_open_in_generated,metric_topology,basis_for_topology,powerset_intro,open_in}.
                Therefore $\metricBall{d}{X}{b}{\delta}$ intersects $A$
                    by \cref{metric_zero_characterization,metric_on,real_lt_subst_left,metric_ball,closure_iff_intersects_open}.
                Take $a$ such that $a\in\metricBall{d}{X}{b}{\delta}\inter A$
                    by \cref{intersects,nonempty_inhabited}.
                We have $a\in A$ and $a\in X$ and $d((b,a)) < \delta$ by \cref{subseteq,inter,metric_ball}.
                Take $h\in H$ such that $d((a,h)) < \delta$ by \cref{subseteq}.
                We have $h\in X$ and $d((h,a)) < \delta$
                    by \cref{subseteq,subseteq_transitive,metric_on,metric_symmetric,real_lt_subst_left,times_tuple_intro}.
                Hence $d((b,h)) < \delta + \delta$ by \cref{common_center_small_ball_distance_bound}.
                Therefore $d((b,h)) < \epsilon$ by \cref{real_lt_subst_right}.
            \end{subproof}
            We have $H$ is finite and
                $H\subseteq X$ and
                for all $b\in B$ there exists $h\in H$ such that
                    $d((b,h)) < \epsilon$.
            Follows by definition.
        \end{byCase}
    \end{subproof}
    Therefore $B$ is totally bounded in $X$ with respect to $d$
        by \cref{finite_metric_subset_totally_bounded_from_net}.
\end{proof}

% from §45 Item 7: equicontinuity and pointwise boundedness imply compact closure
\begin{theorem}\label{ascoli_equicontinuous_pointwise_bounded_implies_compact_closure}
    Assume $n\in\naturalsPlus$.
    Let $J = \initialSegment{n}$.
    Let $Y = \realsPower{J}$.
    Assume $d$ is a metric on $Y$.
    Assume $d = \squareMetric{n}$ or $d = \euclideanMetric{n}$.
    Assume $T$ is a topology on $X$.
    Assume $X$ is compact with respect to $T$.
    Assume $\rho$ is the uniform metric on functions from $X$ to $Y$ with respect to $d$.
    Assume $C$ is the set of continuous functions from $X$ to $Y$ with respect to $T$ and $\metricTopology{d}{Y}$.
    Assume $F\subseteq C$.
    Suppose $F$ is equicontinuous from $X$ to $Y$ with respect to $T$ and $d$.
    Suppose for all $x\in X$ there exists $A$ such that
        $A = \{\apply{f}{x} \mid f\in F\}$ and
        $A$ is bounded in $Y$ with respect to $d$.
    Let $G = \closure{F}{\funs{X}{Y}}{\metricTopology{\rho}{\funs{X}{Y}}}$.
    Then $G$ is compact with respect to $\subspaceTopology{\metricTopology{\rho}{\funs{X}{Y}}}{G}$.
\end{theorem}
\begin{proof}
    Let $Z = \funs{X}{Y}$.
    Let $T_Z = \metricTopology{\rho}{Z}$.
    Let $e = \restrl{\rho}{G\times G}$.
    $\rho$ is a metric on $Z$ by \cref{uniform_metric_on_function_space}.
    $T_Z$ is a topology on $Z$
        by \cref{metric_topology,metric_basis_is_basis,basis_topology_is_topology,basis_for_topology,uniform_metric_on_function_space}.
    $F\subseteq Z$ by \cref{subseteq,continuous_function_set}.
    $G\subseteq Z$ by \cref{closure,subseteq}.
    $G$ is closed in $Z$ with respect to $T_Z$ by \cref{closure_closed_in}.
    $\rho$ is a function and $\rho$ is right-unique and $\rho$ is a relation and $\dom{\rho} = Z\times Z$
        by \cref{metric_on,funs_is_function,function_rightunique,funs_is_relation,funs_elim}.
    Hence $G\times G\subseteq\dom{\rho}$
        by \cref{times_subseteq_left,times_subseteq_right,subseteq_transitive,real_lt_subst_right}.
    $F$ is totally bounded in $Z$ with respect to $\rho$
        by \cref{equicontinuous_pointwise_bounded_function_family_totally_bounded_uniform}.
    Therefore $G$ is totally bounded in $Z$ with respect to $\rho$
        by \cref{closure_totally_bounded}.
    $Y$ is complete with respect to $d$
        by \cref{euclidean_space_square_metric_complete,euclidean_space_euclidean_metric_complete,complete_metric_space,real_lt_subst_right}.
    Hence $Z$ is complete with respect to $\rho$ by \cref{uniform_function_space_complete}.
    Therefore $G$ is complete with respect to $e$ by \cref{closed_subset_complete_restricted_metric}.
    Thus $e$ is a metric on $G$ by \cref{complete_metric_space}.
    $G\subseteq G$ by \cref{subseteq_refl}.
    Show that $G$ is totally bounded with respect to $e$.
        \begin{subproof}
            \begin{byCase}
            \caseOf{$G = \emptyset$.}
                We have $G$ is totally bounded with respect to $e$
                    by \cref{finite_metric_subset_totally_bounded_from_net,emptyset_is_finite,emptyset_subseteq,real_lt_subst_left,subseteq,emptyset,subseteq_emptyset_iff,totally_bounded_metric_space}.
            \caseOf{$G\neq\emptyset$.}
                We have $G$ is totally bounded with respect to $e$
                    by \cref{nonempty_inhabited,finite_metric_subset_totally_bounded_from_net,totally_bounded_subset_finite_internal_net,subseteq,times_tuple_intro,restrl_of_function_apply,real_lt_subst_left,totally_bounded_metric_space}.
            \end{byCase}
        \end{subproof}
    Let $T_G = \metricTopology{e}{G}$.
    Hence $G$ is compact with respect to $T_G$
        by \cref{complete_totally_bounded_metric_space_compact}.
    Show that for all $A$ such that
        $A$ is a covering of $G$ and
        for all $U$ such that $U\in A$ we have $U$ is open in $Z$ with respect to $T_Z$
        there exists $B$ such that
            $B\subseteq A$ and
            $B$ is finite and
            $B$ is a covering of $G$.
    \begin{subproof}
        Fix $A$ such that
            $A$ is a covering of $G$ and
            for all $U$ such that $U\in A$ we have $U$ is open in $Z$ with respect to $T_Z$.
        $T_G$ is a topology on $G$
            by \cref{metric_topology,metric_basis_is_basis,basis_topology_is_topology,basis_for_topology}.
        Let $h = \{(U,G\inter U) \mid U\in A\}$.
        For all $U,V_1,V_2$ such that $(U,V_1)\in h$ and $(U,V_2)\in h$ we have $V_1 = V_2$
            by \cref{pair_eq_iff}.
        We have $h$ is right-unique by \cref{rightunique}.
        Show that for all $U$ we have $U\in\dom{h}$ iff $U\in A$.
        \begin{subproof}
            Follows by \cref{dom_intro,dom_iff,pair_eq_iff}.
        \end{subproof}
        Therefore $h$ is a function on $A$ by \cref{rightunique,relation,dom_intro,dom_iff,pair_eq_iff,subseteq,subseteq_antisymmetric}.
        Let $D = \img{h}{A}$.
        We have $D$ is a covering of $G$
            by \cref{covering,subseteq,unions_iff,inter_intro,img_iff,function_apply_iff}.
        Show that for all $W$ such that $W\in D$ we have $W$ is open in $G$ with respect to $T_G$.
        \begin{subproof}
            Fix $W$ such that $W\in D$.
            Take $U\in A$ such that $(U,W)\in h$ by \cref{img_iff}.
            We have $W = G\inter U$ by \cref{function_apply_intro}.
            Show that for all $g\in W$ there exists $\delta\in\reals$ such that
                $0 < \delta$ and
                $\metricBall{e}{G}{g}{\delta}\subseteq W$.
            \begin{subproof}
                Fix $g$ such that $g\in W$.
                We have $g\in U$ and $g\in G$ by \cref{inter}.
                Take $\delta\in\reals$ such that
                    $0 < \delta$ and
                    $\metricBall{\rho}{Z}{g}{\delta}\subseteq U$
                    by \cref{open_in,topology_on,subseteq,powerset_elim,metric_open_iff}.
                We have $\metricBall{e}{G}{g}{\delta}\subseteq W$
                    by \cref{metric_ball,subseteq,times_tuple_intro,restrl_of_function_apply,real_lt_subst_right,inter_intro,real_lt_subst_left}.
                Follows by definition.
            \end{subproof}
            Therefore $W$ is open in $G$ with respect to $T_G$ by \cref{inter_lower_left,metric_open_iff}.
        \end{subproof}
        Hence $D$ is an open covering of $G$ with respect to $T_G$ by \cref{open_covering}.
        Take $B_0$ such that
            $B_0\subseteq D$ and
            $B_0$ is finite and
            $B_0$ is a covering of $G$
            by \cref{compact}.
        Let $S = \{w\in B_0\times A \mid \fst{w} = G\inter\snd{w}\}$.
        For all $W\in B_0$ there exists $U$ such that $W\mathrel{S} U$
            by \cref{subseteq,img_iff,function_apply_intro,times_tuple_intro,fst_eq,snd_eq}.
        Take $k$ such that
            $k$ is a function on $B_0$ and
            for all $W\in B_0$ we have $W\mathrel{S}\apply{k}{W}$
            by \cref{subseteq,times_elem_is_tuple,relation,choice_function}.
        Let $B = \img{k}{B_0}$.
        We have $B\subseteq A$ by \cref{img_iff,function_apply_intro,subseteq,times_tuple_elim}.
        $B$ is finite by \cref{finite_image_function}.
        We have $B$ is a covering of $G$
            by \cref{covering,subseteq,unions_iff,fst_eq,snd_eq,inter,img_iff,function_apply_iff}.
        Follows by definition.
    \end{subproof}
    Therefore $G$ is compact with respect to $\subspaceTopology{T_Z}{G}$
        by \cref{subspace_of,compact_subspace_open_covering}.
\end{proof}

% from §45 helper lemma 24: totally bounded subsets of metric spaces are bounded
\begin{lemma}\label{totally_bounded_subset_bounded}
    Assume $d$ is a metric on $X$.
    Suppose $A\subseteq X$.
    Suppose $A$ is totally bounded in $X$ with respect to $d$.
    Then $A$ is bounded in $X$ with respect to $d$.
\end{lemma}
\begin{proof}
    \begin{byCase}
    \caseOf{$A = \emptyset$.}
        Let $M = 0$.
        For all $x,z\in A$ we have $d((x,z)) \leq M$
            by \cref{subseteq_emptyset_iff,elem_subseteq,emptyset}.
        Therefore $A$ is bounded in $X$ with respect to $d$ by \cref{real_add_ident,metric_bounded}.
    \caseOf{$A\neq\emptyset$.}
        $1\in\reals$ and $0 < 1$
            by \cref{real_naturals_subset_reals,real_one_in_naturals,subseteq,real_zero_lt_one}.
        Take $H$ such that
            $H\subseteq A$ and
            $H$ is finite and
            for all $a\in A$ there exists $h\in H$ such that
                $d((a,h)) < 1$
            by \cref{nonempty_inhabited,totally_bounded_subset_finite_internal_net}.
        Take $M\in\reals$ such that for all $h_1,h_2\in H$ we have $d((h_1,h_2)) \leq M$
            by \cref{subseteq_transitive,finite_metric_subset_bounded,metric_bounded}.
        Show that for all $x,z\in A$ we have $d((x,z)) \leq 1 + (M + 1) + 1$.
        \begin{subproof}
            Fix $x$.
            Fix $z$ such that $x\in A$ and $z\in A$.
            Take $h_1\in H$ such that $d((x,h_1)) < 1$ by \cref{subseteq}.
            Take $h_2\in H$ such that $d((z,h_2)) < 1$ by \cref{subseteq}.
            We have $x\in X$ and $z\in X$ and $h_1\in X$ and $h_2\in X$ by \cref{subseteq}.
            Hence $d((h_1,h_2)) < M + 1$
                by \cref{metric_on,funs_type_apply,times_tuple_intro,real_add_closed,real_naturals_subset_reals,real_one_in_naturals,subseteq,real_lt_add,real_add_ident,real_add_comm,real_lt_subst_left,real_leq_lt_trans}.
            Therefore $d((h_2,z)) < 1$ by \cref{metric_on,metric_symmetric,real_lt_subst_left,times_tuple_intro}.
            Thus $d((x,z)) < 1 + (M + 1) + 1$
                by \cref{metric_three_step_lt_sum,real_add_closed,real_naturals_subset_reals,real_one_in_naturals,subseteq}.
            Therefore $d((x,z)) \leq 1 + (M + 1) + 1$ by \cref{real_lt_imp_leq}.
        \end{subproof}
        Therefore $A$ is bounded in $X$ with respect to $d$
            by \cref{real_add_closed,real_naturals_subset_reals,real_one_in_naturals,subseteq,metric_bounded}.
    \end{byCase}
\end{proof}

% from §45 helper lemma 25: continuous function spaces into finite real powers are inhabited
\begin{lemma}\label{continuous_reals_power_function_set_inhabited}
    Assume $n\in\naturalsPlus$.
    Let $J = \initialSegment{n}$.
    Let $Y = \realsPower{J}$.
    Assume $T$ is a topology on $X$.
    Assume $d$ is a metric on $Y$.
    Assume $C$ is the set of continuous functions from $X$ to $Y$ with respect to $T$ and $\metricTopology{d}{Y}$.
    Then $C$ is inhabited.
\end{lemma}
\begin{proof}
    Let $y_0 = \{(j,0) \mid j\in J\}$.
    We have $y_0\in Y$
        by \cref{real_add_ident,constant_map_function,funs,tuple_power,reals_power,real_lt_subst_right}.
    Let $f_0 = \{(x,y_0) \mid x\in X\}$.
    We have $f_0\in\funs{X}{Y}$ by \cref{constant_map_function,funs}.
    Hence $f_0$ is continuous from $X$ to $Y$ with respect to $T$ and $\metricTopology{d}{Y}$
        by \cref{metric_topology,metric_basis_is_basis,basis_topology_is_topology,basis_for_topology,funs_is_function,function_rightunique,funs_is_relation,function_apply_intro,constant_continuous}.
    Hence $C$ is inhabited by \cref{continuous_function_set,inhabited_nonempty,nonempty_inhabited}.
\end{proof}

% from §45 helper lemma 26: each coordinate distance is bounded by the sup metric
\begin{lemma}\label{sup_metric_coordinate_leq_general}
    Assume $d$ is a metric on $Y$.
    Assume $B\subseteq\funs{X}{Y}$.
    Assume $\rho$ is the sup metric on $B$ with respect to $X$ and $Y$ and $d$.
    Suppose $f\in B$.
    Suppose $g\in B$.
    Suppose $x\in X$.
    Then $d((f(x),g(x))) \leq \rho((f,g))$.
\end{lemma}
\begin{proof}
    We have $d((f(x),g(x))) \leq \rho((f,g))$
        by \cref{sup_metric_on_bounded_function_space,subseteq,function_distance_set,metric_on,funs_type_apply,real_supremum,real_upper_bound}.
\end{proof}

% from §45 helper lemma 27: bounded subsets in the sup metric are pointwise bounded
\begin{lemma}\label{bounded_sup_implies_pointwise_bounded}
    Assume $d$ is a metric on $Y$.
    Assume $T$ is a topology on $X$.
    Assume $C$ is the set of continuous functions from $X$ to $Y$ with respect to $T$ and $\metricTopology{d}{Y}$.
    Assume $\sigma$ is the sup metric on $C$ with respect to $X$ and $Y$ and $d$.
    Assume $F\subseteq C$.
    Suppose $F$ is bounded in $C$ with respect to $\sigma$.
    Then for all $x\in X$ there exists $A$ such that
        $A = \{\apply{f}{x} \mid f\in F\}$ and
        $A$ is bounded in $Y$ with respect to $d$.
\end{lemma}
\begin{proof}
    Take $M\in\reals$ such that for all $f_1,f_2\in F$ we have $\sigma((f_1,f_2)) \leq M$
        by \cref{metric_bounded}.
    Show that for all $x\in X$ there exists $A$ such that
        $A = \{\apply{f}{x} \mid f\in F\}$ and
        $A$ is bounded in $Y$ with respect to $d$.
    \begin{subproof}
        Fix $x$ such that $x\in X$.
        Let $A = \{\apply{f}{x} \mid f\in F\}$.
        Show that for all $y_1,y_2\in A$ we have $d((y_1,y_2)) \leq M$.
        \begin{subproof}
            Fix $y_1$.
            Fix $y_2$ such that $y_1\in A$ and $y_2\in A$.
            Take $f_1\in F$ such that $y_1 = f_1(x)$ by \cref{subseteq}.
            Take $f_2\in F$ such that $y_2 = f_2(x)$ by \cref{subseteq}.
            Hence $d((y_1,y_2)) \leq M$
                by \cref{sup_metric_coordinate_leq_general,metric_bounded,subseteq,metric_on,continuous_function_set,funs_type_apply,times_tuple_intro,real_leq_trans,real_lt_subst_left,real_lt_subst_right}.
        \end{subproof}
        Therefore $A$ is bounded in $Y$ with respect to $d$
            by \cref{subseteq,continuous_function_set,funs_type_apply,real_lt_subst_left,metric_bounded}.
        Follows by definition.
    \end{subproof}
    Follows by definition.
\end{proof}

% from §45 Item 8: compact function families are closed
\begin{corollary}\label{ascoli_compact_function_family_closed}
    Assume $n\in\naturalsPlus$.
    Let $J = \initialSegment{n}$.
    Let $Y = \realsPower{J}$.
    Assume $d$ is a metric on $Y$.
    Assume $d = \squareMetric{n}$ or $d = \euclideanMetric{n}$.
    Assume $T$ is a topology on $X$.
    Assume $X$ is compact with respect to $T$.
    Assume $\rho$ is the uniform metric on functions from $X$ to $Y$ with respect to $d$.
    Assume $C$ is the set of continuous functions from $X$ to $Y$ with respect to $T$ and $\metricTopology{d}{Y}$.
    Assume $F\subseteq C$.
    Suppose $F$ is compact with respect to $\subspaceTopology{\metricTopology{\rho}{\funs{X}{Y}}}{F}$.
    Then $F$ is closed in $\funs{X}{Y}$ with respect to $\metricTopology{\rho}{\funs{X}{Y}}$.
\end{corollary}
\begin{proof}
    Let $Z = \funs{X}{Y}$.
    Let $T_Z = \metricTopology{\rho}{Z}$.
    $\rho$ is a metric on $Z$ by \cref{uniform_metric_on_function_space}.
    $T_Z$ is a topology on $Z$
        by \cref{metric_topology,metric_basis_is_basis,basis_topology_is_topology,basis_for_topology,uniform_metric_on_function_space}.
    Hence $F$ is closed in $Z$ with respect to $T_Z$
        by \cref{metric_space_hausdorff,compact_subspace_closed,subseteq,continuous_function_set,subseteq_transitive}.
\end{proof}

% from §45 Item 8: compact function families are bounded in the sup metric
\begin{corollary}\label{ascoli_compact_function_family_bounded_sup}
    Assume $n\in\naturalsPlus$.
    Let $J = \initialSegment{n}$.
    Let $Y = \realsPower{J}$.
    Assume $d$ is a metric on $Y$.
    Assume $d = \squareMetric{n}$ or $d = \euclideanMetric{n}$.
    Assume $T$ is a topology on $X$.
    Assume $X$ is compact with respect to $T$.
    Assume $X$ is inhabited.
    Assume $C$ is the set of continuous functions from $X$ to $Y$ with respect to $T$ and $\metricTopology{d}{Y}$.
    Assume $\sigma$ is the sup metric on $C$ with respect to $X$ and $Y$ and $d$.
    Assume $F\subseteq C$.
    Suppose $F$ is compact with respect to $\subspaceTopology{\metricTopology{\sigma}{C}}{F}$.
    Then $F$ is bounded in $C$ with respect to $\sigma$.
\end{corollary}
\begin{proof}
    We have $F$ is bounded in $C$ with respect to $\sigma$
        by \cref{continuous_reals_power_function_set_inhabited,sup_metric_on_bounded_function_space,compact_metric_subspace_ambient_totally_bounded,totally_bounded_subset_bounded}.
\end{proof}

% from §45 Item 8: compact function families are equicontinuous
\begin{corollary}\label{ascoli_compact_function_family_equicontinuous}
    Assume $n\in\naturalsPlus$.
    Let $J = \initialSegment{n}$.
    Let $Y = \realsPower{J}$.
    Assume $d$ is a metric on $Y$.
    Assume $d = \squareMetric{n}$ or $d = \euclideanMetric{n}$.
    Assume $T$ is a topology on $X$.
    Assume $X$ is compact with respect to $T$.
    Assume $\rho$ is the uniform metric on functions from $X$ to $Y$ with respect to $d$.
    Assume $C$ is the set of continuous functions from $X$ to $Y$ with respect to $T$ and $\metricTopology{d}{Y}$.
    Assume $F\subseteq C$.
    Suppose $F$ is compact with respect to $\subspaceTopology{\metricTopology{\rho}{\funs{X}{Y}}}{F}$.
    Then $F$ is equicontinuous from $X$ to $Y$ with respect to $T$ and $d$.
\end{corollary}
\begin{proof}
    Let $Z = \funs{X}{Y}$.
    Let $T_Z = \metricTopology{\rho}{Z}$.
    Let $G = \closure{F}{Z}{T_Z}$.
    $T_Z$ is a topology on $Z$
        by \cref{metric_topology,metric_basis_is_basis,basis_topology_is_topology,basis_for_topology,uniform_metric_on_function_space}.
    We have $G = F$
        by \cref{ascoli_compact_function_family_closed,closure_subseteq_closed,subseteq,continuous_function_set,subseteq_closure,subseteq_antisymmetric}.
    Therefore $G$ is compact with respect to $\subspaceTopology{T_Z}{G}$
        by \cref{real_lt_subst_left,real_lt_subst_right}.
    Hence $F$ is equicontinuous from $X$ to $Y$ with respect to $T$ and $d$
        by \cref{ascoli_compact_closure_implies_equicontinuous}.
\end{proof}

% from §45 Item 8: closed bounded equicontinuous function families are compact
\begin{corollary}\label{ascoli_closed_bounded_equicontinuous_implies_compact}
    Assume $n\in\naturalsPlus$.
    Let $J = \initialSegment{n}$.
    Let $Y = \realsPower{J}$.
    Assume $d$ is a metric on $Y$.
    Assume $d = \squareMetric{n}$ or $d = \euclideanMetric{n}$.
    Assume $T$ is a topology on $X$.
    Assume $X$ is compact with respect to $T$.
    Assume $\rho$ is the uniform metric on functions from $X$ to $Y$ with respect to $d$.
    Assume $C$ is the set of continuous functions from $X$ to $Y$ with respect to $T$ and $\metricTopology{d}{Y}$.
    Assume $\sigma$ is the sup metric on $C$ with respect to $X$ and $Y$ and $d$.
    Assume $F\subseteq C$.
    Suppose $F$ is closed in $\funs{X}{Y}$ with respect to $\metricTopology{\rho}{\funs{X}{Y}}$.
    Suppose $F$ is bounded in $C$ with respect to $\sigma$.
    Suppose $F$ is equicontinuous from $X$ to $Y$ with respect to $T$ and $d$.
    Then $F$ is compact with respect to $\subspaceTopology{\metricTopology{\rho}{\funs{X}{Y}}}{F}$.
\end{corollary}
\begin{proof}
    Let $Z = \funs{X}{Y}$.
    Let $T_Z = \metricTopology{\rho}{Z}$.
    Let $G = \closure{F}{Z}{T_Z}$.
    $T_Z$ is a topology on $Z$
        by \cref{metric_topology,metric_basis_is_basis,basis_topology_is_topology,basis_for_topology,uniform_metric_on_function_space}.
    We have for all $x\in X$ there exists $A$ such that
        $A = \{\apply{f}{x} \mid f\in F\}$ and
        $A$ is bounded in $Y$ with respect to $d$
        by \cref{bounded_sup_implies_pointwise_bounded}.
    Therefore $G$ is compact with respect to $\subspaceTopology{T_Z}{G}$
        by \cref{ascoli_equicontinuous_pointwise_bounded_implies_compact_closure}.
    We have $G = F$
        by \cref{closure_subseteq_closed,subseteq,continuous_function_set,subseteq_closure,subseteq_antisymmetric}.
    Hence $F$ is compact with respect to $\subspaceTopology{T_Z}{F}$
        by \cref{real_lt_subst_left,real_lt_subst_right}.
\end{proof}



\section{§46 Pointwise and Compact Convergence}

% from §46 Item 1: point-open set
\begin{definition}\label{point_open_set}
    $A$ is the point-open set at $x$ with value set $V$ from $X$ to $Y$ iff
        $x\in X$ and
        $V\subseteq Y$ and
        $A = \{f\in\funs{X}{Y} \mid \apply{f}{x}\in V\}$.
\end{definition}

% from §46 helper definition 1: pointwise subbasis on a function space
\begin{definition}\label{pointwise_function_space_subbasis}
    $S$ is the pointwise subbasis on functions from $X$ to $Y$ with respect to $T$ iff
        $T$ is a topology on $Y$ and
        $S = \{A\in\pow{\funs{X}{Y}} \mid \text{there exist $x,V$ such that
            $x\in X$ and
            $V\in T$ and
            $A$ is the point-open set at $x$ with value set $V$ from $X$ to $Y$}\}$.
\end{definition}

% from §46 Item 1: topology of pointwise convergence
\begin{definition}\label{pointwise_convergence_topology}
    $T_0$ is the topology of pointwise convergence on functions from $X$ to $Y$ with respect to $T$ iff
        $T$ is a topology on $Y$ and
        $T_0$ is a topology on $\funs{X}{Y}$ and
        there exists $S$ such that
            $S$ is the pointwise subbasis on functions from $X$ to $Y$ with respect to $T$ and
            $T_0 = \subbasisTopology{S}{\funs{X}{Y}}$.
\end{definition}

% from §46 helper lemma 2: point-open sets belong to a pointwise subbasis
\begin{lemma}\label{pointwise_subbasis_contains_point_open_set}
    Assume $T$ is a topology on $Y$.
    Assume $S$ is the pointwise subbasis on functions from $X$ to $Y$ with respect to $T$.
    Assume $x\in X$.
    Assume $V\in T$.
    Let $A = \{g\in\funs{X}{Y} \mid \apply{g}{x}\in V\}$.
    Then $A\in S$.
\end{lemma}
\begin{proof}
    We have $A\in S$
        by \cref{topology_on,subseteq,pow_iff,point_open_set,pointwise_function_space_subbasis}.
\end{proof}

% from §46 helper definition 2: coordinate sequence of a function sequence at a point
\begin{definition}\label{function_sequence_coordinate_sequence}
    $t$ is the coordinate sequence of $s$ at $x$ from $X$ to $Y$ iff
        $s$ is a sequence in $\funs{X}{Y}$ and
        $x\in X$ and
        $t$ is a sequence in $Y$ and
        for all $n\in\naturalsPlus$ we have $t(n) = \apply{\apply{s}{n}}{x}$.
\end{definition}

% from §46 helper lemma 1: coordinate sequences exist
\begin{lemma}\label{coordinate_sequence_exists}
    Assume $s$ is a sequence in $\funs{X}{Y}$.
    Assume $x\in X$.
    Then there exists $t$ such that $t$ is the coordinate sequence of $s$ at $x$ from $X$ to $Y$.
\end{lemma}
\begin{proof}
    Let $R = \{w\in\naturalsPlus\times Y \mid \snd{w} = \apply{\apply{s}{\fst{w}}}{x}\}$.
    $R\subseteq\naturalsPlus\times Y$ by \cref{subseteq}.
    We have $R$ is a relation by \cref{relation,subseteq,times_elem_is_tuple}.
    For all $n\in\naturalsPlus$ we have $s(n)\in\funs{X}{Y}$ by \cref{sequence_in,funs_type_apply}.
    For all $n\in\naturalsPlus$ we have $\apply{\apply{s}{n}}{x}\in Y$ by \cref{funs_type_apply}.
    Show that for all $n\in\naturalsPlus$ there exists $y$ such that $(n,y)\in R$.
    \begin{subproof}
        Fix $n\in\naturalsPlus$.
        Let $y = \apply{\apply{s}{n}}{x}$.
        We have $y\in Y$ by \cref{subseteq}.
        $(n,y)\in\naturalsPlus\times Y$ by \cref{times_tuple_intro}.
        Therefore $(n,y)\in R$ by \cref{fst_eq,snd_eq,subseteq}.
    \end{subproof}
    Take $t$ such that $t$ is a function on $\naturalsPlus$ and
        for all $n\in\naturalsPlus$ we have $n\mathrel{R} t(n)$
        by \cref{choice_function}.
    Therefore $t$ is the coordinate sequence of $s$ at $x$ from $X$ to $Y$
        by \cref{choice_function,funs_intro,sequence_in,function_apply_elim,subseteq,times_tuple_elim,fst_eq,snd_eq,function_sequence_coordinate_sequence}.
\end{proof}

% from §46 helper lemma 1a: coordinate sequences are indexed projections of function sequences
\begin{lemma}\label{coordinate_sequence_is_projection_composition}
    Let $X_0 = \{(x,Y) \mid x\in X\}$.
    Assume $s$ is a sequence in $\funs{X}{Y}$.
    Suppose $x\in X$.
    Suppose $t$ is the coordinate sequence of $s$ at $x$ from $X$ to $Y$.
    Then $t = \piIndex{X_0}{x}\circ s$.
\end{lemma}
\begin{proof}
    For all $x_1,y_1,y_2$ such that $(x_1,y_1)\in X_0$ and $(x_1,y_2)\in X_0$ we have $y_1 = y_2$.
    We have $X_0$ is a function on $X$
        by \cref{rightunique,relation,dom_intro,dom_iff,pair_eq_iff,subseteq,subseteq_antisymmetric}.
    $\dom{X_0} = X$ by \cref{funs}.
    Hence $X_0(x) = Y$ by \cref{function_apply_intro}.
    $\funs{X}{Y}\subseteq\productFamily{X_0}$
        by \cref{subseteq,funs_is_function,funs_elim,funs_type_apply,relation,ran_intro,unions_iff,funs_intro,function_apply_intro,indexed_product}.
    We have $s$ is a sequence in $\productFamily{X_0}$ by \cref{sequence_in,funs_weaken_codom}.
    Hence $\piIndex{X_0}{x}\circ s\in\funs{\naturalsPlus}{Y}$
        by \cref{projection_map_indexed_function,funs_elim,funs_circ,sequence_in,subseteq,function_apply_intro}.
    Therefore $t = \piIndex{X_0}{x}\circ s$
        by \cref{functions_eq_iff,function_sequence_coordinate_sequence,sequence_in,funs_type_apply,projection_map_indexed,function_apply_elim,circ_iff,function_apply_intro,projection_map_indexed_function,funs_is_function,subseteq,funs}.
\end{proof}

% from §46 Item 2: pointwise convergence is coordinatewise convergence
\begin{theorem}\label{pointwise_convergence_iff_coordinatewise}
    Assume $T$ is a topology on $Y$.
    Assume $T_0$ is the topology of pointwise convergence on functions from $X$ to $Y$ with respect to $T$.
    Assume $s$ is a sequence in $\funs{X}{Y}$.
    Assume $f\in\funs{X}{Y}$.
    Then $s$ converges to $f$ in $\funs{X}{Y}$ with respect to $T_0$ iff
        for all $x\in X$ there exists $t$ such that
            $t$ is the coordinate sequence of $s$ at $x$ from $X$ to $Y$ and
            $t$ converges to $\apply{f}{x}$ in $Y$ with respect to $T$.
\end{theorem}
\begin{proof}
    Let $X_0 = \{(x,Y) \mid x\in X\}$.
    Let $T_1 = \{(x,T) \mid x\in X\}$.
    For all $x,y_1,y_2$ such that $(x,y_1)\in X_0$ and $(x,y_2)\in X_0$ we have $y_1 = y_2$.
    Hence $X_0$ is a function on $X$
        by \cref{rightunique,relation,dom_intro,dom_iff,pair_eq_iff,subseteq,subseteq_antisymmetric}.
    Therefore $\dom{X_0} = X$ by \cref{funs}.
    For all $x,U_1,U_2$ such that $(x,U_1)\in T_1$ and $(x,U_2)\in T_1$ we have $U_1 = U_2$.
    We have $T_1$ is a function on $X$
        by \cref{rightunique,relation,dom_intro,dom_iff,pair_eq_iff,subseteq,subseteq_antisymmetric}.
    Hence $\dom{T_1} = X$ by \cref{funs}.
    For all $x\in X$ we have $X_0(x) = Y$ and $T_1(x) = T$ by \cref{function_apply_intro}.
    Show that $\productFamily{X_0}\subseteq\funs{X}{Y}$.
    \begin{subproof}
        Follows by \cref{subseteq,indexed_product,funs_is_function,funs_elim,function_apply_intro,funs_intro}.
    \end{subproof}
    Show that $\funs{X}{Y}\subseteq\productFamily{X_0}$.
    \begin{subproof}
        Follows by \cref{subseteq,funs_is_function,funs_elim,funs_type_apply,relation,ran_intro,unions_iff,funs_intro,function_apply_intro,indexed_product}.
    \end{subproof}
    Hence $\productFamily{X_0} = \funs{X}{Y}$ by \cref{subseteq,subseteq_antisymmetric}.
    For all $x\in\dom{X_0}$ we have $T_1(x)$ is a topology on $X_0(x)$
        by \cref{function_apply_intro,pair_eq_iff}.
    Let $T_2 = \productTopologyIndexed{T_1}{X_0}$.
    Take $S$ such that
        $S$ is the pointwise subbasis on functions from $X$ to $Y$ with respect to $T$ and
        $T_0 = \subbasisTopology{S}{\funs{X}{Y}}$ by \cref{pointwise_convergence_topology}.
    Show that for all $A$ if $A\in S$ then $A\in\productSubbasis{T_1}{X_0}$.
    \begin{subproof}
        Fix $A$.
        Assume $A\in S$.
        Take $x,V$ such that
            $x\in X$ and
            $V\in T$ and
            $A$ is the point-open set at $x$ with value set $V$ from $X$ to $Y$
            by \cref{pointwise_function_space_subbasis}.
        We have $A = \{g\in\funs{X}{Y} \mid \apply{g}{x}\in V\}$ by \cref{point_open_set}.
        Show that $A\subseteq \preimg{\piIndex{X_0}{x}}{V}$.
        \begin{subproof}
            Follows by \cref{subseteq,inters_iff_forall,projection_map_indexed,function_apply_elim,preimg_iff}.
        \end{subproof}
        Show that $\preimg{\piIndex{X_0}{x}}{V}\subseteq A$.
        \begin{subproof}
            Follows by \cref{subseteq,preimg_iff,projection_map_indexed,pair_eq_iff,inters_iff_forall}.
        \end{subproof}
        Therefore $A = \preimg{\piIndex{X_0}{x}}{V}$ by \cref{subseteq,subseteq_antisymmetric}.
        We have $A\in\pow{\productFamily{X_0}}$
            by \cref{projection_map_indexed_function,funs_elim,preimg_subseteq_dom,subseteq_transitive,pow_iff}.
        Hence $A\in\productSubbasisAt{T_1}{X_0}{x}$ by \cref{product_subbasis_indexed,function_apply_intro}.
        Therefore $A\in\productSubbasis{T_1}{X_0}$ by \cref{product_subbasis_union_indexed}.
    \end{subproof}
    We have $S\subseteq\productSubbasis{T_1}{X_0}$ by \cref{subseteq}.
    Show that for all $A$ if $A\in\productSubbasis{T_1}{X_0}$ then $A\in S$.
    \begin{subproof}
        Fix $A$.
        Assume $A\in\productSubbasis{T_1}{X_0}$.
        Take $x\in\dom{X_0}$ such that $A\in\productSubbasisAt{T_1}{X_0}{x}$ by \cref{product_subbasis_union_indexed}.
        Take $V\in T_1(x)$ such that $A = \preimg{\piIndex{X_0}{x}}{V}$ by \cref{product_subbasis_indexed}.
        Let $B = \{g\in\funs{X}{Y} \mid \apply{g}{x}\in V\}$.
        Show that $A\subseteq B$.
        \begin{subproof}
            Follows by \cref{subseteq,pair_eq_iff,preimg_iff,projection_map_indexed,inters_iff_forall}.
        \end{subproof}
        Show that $B\subseteq A$.
        \begin{subproof}
            Follows by \cref{subseteq,inters_iff_forall,projection_map_indexed,function_apply_elim,preimg_iff,pair_eq_iff}.
        \end{subproof}
        Hence $A\in S$ by \cref{subseteq,subseteq_antisymmetric,pointwise_subbasis_contains_point_open_set,function_apply_intro,pair_eq_iff}.
    \end{subproof}
    We have $\productSubbasis{T_1}{X_0}\subseteq S$ by \cref{subseteq}.
    Therefore $S = \productSubbasis{T_1}{X_0}$ by \cref{subseteq,subseteq_antisymmetric}.
    Hence $T_2 = \subbasisTopology{S}{\productFamily{X_0}}$ by \cref{product_topology_indexed}.
    Thus $T_2 = T_0$ by \cref{topology_generated_by_subbasis}.
    Show that if $s$ converges to $f$ in $\funs{X}{Y}$ with respect to $T_0$ then
        for all $x\in X$ there exists $t$ such that
            $t$ is the coordinate sequence of $s$ at $x$ from $X$ to $Y$ and
            $t$ converges to $\apply{f}{x}$ in $Y$ with respect to $T$.
    \begin{subproof}
        Assume $s$ converges to $f$ in $\funs{X}{Y}$ with respect to $T_0$.
        Hence $s$ converges to $f$ in $\productFamily{X_0}$ with respect to $T_2$ by \cref{converges_to}.
        For all $x\in\dom{X_0}$ we have
            $\piIndex{X_0}{x}\circ s$ converges to $\apply{f}{x}$ in $\apply{X_0}{x}$ with respect to $\apply{T_1}{x}$
            by \cref{product_sequence_converges_iff_coordinatewise,subseteq}.
        For all $x\in X$ there exists $t$ such that
            $t$ is the coordinate sequence of $s$ at $x$ from $X$ to $Y$ and
            $t$ converges to $\apply{f}{x}$ in $Y$ with respect to $T$
            by \cref{coordinate_sequence_exists,coordinate_sequence_is_projection_composition,subseteq,function_apply_intro}.
    \end{subproof}
    Show that if for all $x\in X$ there exists $t$ such that
        $t$ is the coordinate sequence of $s$ at $x$ from $X$ to $Y$ and
        $t$ converges to $\apply{f}{x}$ in $Y$ with respect to $T$
    then $s$ converges to $f$ in $\funs{X}{Y}$ with respect to $T_0$.
    \begin{subproof}
        Assume for all $x\in X$ there exists $t$ such that
            $t$ is the coordinate sequence of $s$ at $x$ from $X$ to $Y$ and
            $t$ converges to $\apply{f}{x}$ in $Y$ with respect to $T$.
        For all $x\in\dom{X_0}$ we have
            $\piIndex{X_0}{x}\circ s$ converges to $\apply{f}{x}$ in $\apply{X_0}{x}$ with respect to $\apply{T_1}{x}$
            by \cref{coordinate_sequence_exists,coordinate_sequence_is_projection_composition,converges_to,function_apply_intro}.
        Hence $s$ converges to $f$ in $\productFamily{X_0}$ with respect to $T_2$
            by \cref{product_sequence_converges_iff_coordinatewise,subseteq}.
        Therefore $s$ converges to $f$ in $\funs{X}{Y}$ with respect to $T_0$ by \cref{converges_to}.
    \end{subproof}
    Therefore $s$ converges to $f$ in $\funs{X}{Y}$ with respect to $T_0$ iff
        for all $x\in X$ there exists $t$ such that
            $t$ is the coordinate sequence of $s$ at $x$ from $X$ to $Y$ and
            $t$ converges to $\apply{f}{x}$ in $Y$ with respect to $T$.
\end{proof}

% from §46 Item 3: compact-convergence ball
\begin{definition}\label{compact_convergence_ball}
    $A$ is the compact-convergence ball about $f$ over $C$ with radius $\epsilon$
        from $X$ to $Y$ with respect to $T$ and $d$ iff
        $T$ is a topology on $X$ and
        $d$ is a metric on $Y$ and
        $f\in\funs{X}{Y}$ and
        $C$ is a compact subspace of $X$ with respect to $T$ and
        $\epsilon\in\reals$ and
        $0 < \epsilon$ and
        there exists $\rho$ such that
            $\rho$ is the uniform metric on functions from $C$ to $Y$ with respect to $d$ and
            $A = \{g\in\funs{X}{Y} \mid \rho((\restrl{f}{C},\restrl{g}{C})) < \epsilon\}$.
\end{definition}

% from §46 helper definition 3: compact-convergence basis on a function space
\begin{definition}\label{compact_convergence_function_space_basis}
    $B$ is the compact-convergence basis on functions from $X$ to $Y$ with respect to $T$ and $d$ iff
        $T$ is a topology on $X$ and
        $d$ is a metric on $Y$ and
        $B = \{A\in\pow{\funs{X}{Y}} \mid \text{there exist $f, C, \epsilon$ such that
            $A$ is the compact-convergence ball about $f$ over $C$ with radius $\epsilon$
                from $X$ to $Y$ with respect to $T$ and $d$}\}$.
\end{definition}

% from §46 Item 3: topology of compact convergence
\begin{definition}\label{compact_convergence_topology}
    $T_0$ is the topology of compact convergence on functions from $X$ to $Y$ with respect to $T$ and $d$ iff
        $T$ is a topology on $X$ and
        $d$ is a metric on $Y$ and
        $T_0$ is a topology on $\funs{X}{Y}$ and
        there exists $B$ such that
            $B$ is the compact-convergence basis on functions from $X$ to $Y$ with respect to $T$ and $d$ and
            $T_0 = \basisTopology{B}{\funs{X}{Y}}$.
\end{definition}

% from §46 helper definition 4: restriction sequence on a subset
\begin{definition}\label{restriction_sequence_on_subset}
    $t$ is the restriction sequence of $s$ to $C$ from $X$ to $Y$ iff
        $s$ is a sequence in $\funs{X}{Y}$ and
        $C\subseteq X$ and
        $t$ is a sequence in $\funs{C}{Y}$ and
        for all $n\in\naturalsPlus$ we have $t(n) = \restrl{s(n)}{C}$.
\end{definition}

% from §46 helper lemma 3: restricting a typed function preserves codomain
\begin{lemma}\label{restrl_typed_function}
    Assume $f\in\funs{X}{Y}$.
    Assume $C\subseteq X$.
    Then $\restrl{f}{C}\in\funs{C}{Y}$.
\end{lemma}
\begin{proof}
    We have $f$ is a function and $\dom{f} = X$ by \cref{funs_elim}.
    Hence $\restrl{f}{C}\in\funs{C}{Y}$
        by \cref{restrl_of_function_is_function,subseteq,elem_dom_and_restr_implies_elem_of_restr,restrl_dom,subseteq_antisymmetric,elem_dom_of_restrl_implies_elem_dom_and_restr,funs_elim,restrl_of_function_apply,funs_intro}.
\end{proof}

% from §46 helper lemma 4: restriction sequences exist
\begin{lemma}\label{restriction_sequence_exists}
    Assume $s$ is a sequence in $\funs{X}{Y}$.
    Assume $C\subseteq X$.
    Then there exists $t$ such that $t$ is the restriction sequence of $s$ to $C$ from $X$ to $Y$.
\end{lemma}
\begin{proof}
    Let $t = \{(n,\restrl{s(n)}{C}) \mid n\in\naturalsPlus\}$.
    We have $t$ is a relation by \cref{relation}.
    For all $n,u,v$ such that $(n,u)\in t$ and $(n,v)\in t$ we have $u = v$ by \cref{pair_eq_iff}.
    Therefore $t$ is a function by \cref{rightunique}.
    Therefore $\dom{t} = \naturalsPlus$
        by \cref{dom_iff,pair_eq_iff,relation,dom_intro,subseteq,subseteq_antisymmetric}.
    For all $n$ such that $n\in\dom{t}$ we have $t(n)\in\funs{C}{Y}$
        by \cref{function_apply_elim,pair_eq_iff,sequence_in,funs_type_apply,restrl_typed_function}.
    Hence $t$ is the restriction sequence of $s$ to $C$ from $X$ to $Y$
        by \cref{funs_intro,sequence_in,function_apply_elim,pair_eq_iff,restriction_sequence_on_subset}.
\end{proof}

% from §46 helper lemma 5: compact-convergence balls refine around their points
\begin{lemma}\label{compact_convergence_ball_refinement}
    Assume $A$ is the compact-convergence ball about $f$ over $C$ with radius $\epsilon$
        from $X$ to $Y$ with respect to $T$ and $d$.
    Assume $h\in A$.
    Then there exist $D,\delta$ such that
        $\delta\in\reals$ and
        $0 < \delta$ and
        $D$ is the compact-convergence ball about $h$ over $C$ with radius $\delta$
            from $X$ to $Y$ with respect to $T$ and $d$ and
        $D\subseteq A$.
\end{lemma}
\begin{proof}
    Let $F = \funs{X}{Y}$.
    Take $\rho$ such that
        $\rho$ is the uniform metric on functions from $C$ to $Y$ with respect to $d$ and
        $A = \{g\in F \mid \rho((\restrl{f}{C},\restrl{g}{C})) < \epsilon\}$
        by \cref{compact_convergence_ball}.
    Let $F_C = \funs{C}{Y}$.
    $f\in F$ by \cref{compact_convergence_ball}.
    $C\subseteq X$ by \cref{compact_convergence_ball,compact_subspace}.
    $\rho$ is a metric on $F_C$ by \cref{uniform_metric_on_function_space}.
    Let $U = \metricBall{\rho}{F_C}{\restrl{f}{C}}{\epsilon}$.
    $U\subseteq F_C$ by \cref{metric_ball,subseteq}.
    $\restrl{f}{C}\in F_C$ by \cref{restrl_typed_function}.
    $\metricTopology{\rho}{F_C}$ is a topology on $F_C$
        by \cref{metric_topology,metric_basis_is_basis,basis_topology_is_topology,basis_for_topology}.
    We have $U\in\metricBasis{\rho}{F_C}$ by \cref{compact_convergence_ball,metric_basis,powerset_intro}.
    Therefore $U$ is open in $F_C$ with respect to $\metricTopology{\rho}{F_C}$
        by \cref{metric_basis_is_basis,basis_elem_open_in_generated,metric_topology,open_in}.
    We have $\restrl{h}{C}\in U$
        by \cref{inters_iff_forall,compact_convergence_ball,compact_subspace,restrl_typed_function,metric_ball}.
    Take $\delta\in\reals$ such that $0 < \delta$ and
        $\metricBall{\rho}{F_C}{\restrl{h}{C}}{\delta} \subseteq U$
        by \cref{metric_open_iff}.
    Let $D = \{g\in F \mid \rho((\restrl{h}{C},\restrl{g}{C})) < \delta\}$.
    We have $D$ is the compact-convergence ball about $h$ over $C$ with radius $\delta$
        from $X$ to $Y$ with respect to $T$ and $d$ by \cref{compact_convergence_ball}.
    Therefore $D\subseteq A$ by \cref{inters_iff_forall,restrl_typed_function,metric_ball,subseteq}.
    Follows by definition.
\end{proof}

% from §46 helper lemma 6: the center belongs to each compact-convergence ball
\begin{lemma}\label{compact_convergence_ball_center}
    Assume $A$ is the compact-convergence ball about $f$ over $C$ with radius $\epsilon$
        from $X$ to $Y$ with respect to $T$ and $d$.
    Then $f\in A$.
\end{lemma}
\begin{proof}
    Let $F = \funs{X}{Y}$.
    Let $F_C = \funs{C}{Y}$.
    Take $\rho$ such that
        $\rho$ is the uniform metric on functions from $C$ to $Y$ with respect to $d$ and
        $A = \{g\in F \mid \rho((\restrl{f}{C},\restrl{g}{C})) < \epsilon\}$
        by \cref{compact_convergence_ball}.
    We have $\restrl{f}{C}\in F_C$ by \cref{compact_convergence_ball,compact_subspace,restrl_typed_function}.
    Therefore $\rho((\restrl{f}{C},\restrl{f}{C})) = 0$
        by \cref{uniform_metric_on_function_space,metric_on,metric_zero_characterization,times_tuple_intro}.
    Thus $\rho((\restrl{f}{C},\restrl{f}{C})) < \epsilon$
        by \cref{compact_convergence_ball,real_lt_subst_left}.
    Therefore $f\in A$ by \cref{compact_convergence_ball,inters_iff_forall}.
\end{proof}

% from §46 helper lemma 7: compact-convergence balls are open
\begin{lemma}\label{compact_convergence_ball_open}
    Assume $A$ is the compact-convergence ball about $f$ over $C$ with radius $\epsilon$
        from $X$ to $Y$ with respect to $T$ and $d$.
    Assume $T_0$ is the topology of compact convergence on functions from $X$ to $Y$ with respect to $T$ and $d$.
    Then $A$ is open in $\funs{X}{Y}$ with respect to $T_0$.
\end{lemma}
\begin{proof}
    Let $F = \funs{X}{Y}$.
    Take $B$ such that
        $B$ is the compact-convergence basis on functions from $X$ to $Y$ with respect to $T$ and $d$ and
        $T_0 = \basisTopology{B}{F}$ by \cref{compact_convergence_topology}.
    For all $h\in A$ there exists $D\in B$ such that $h\in D$ and $D\subseteq A$
        by \cref{compact_convergence_ball_refinement,compact_convergence_ball,subseteq,compact_convergence_function_space_basis,powerset_intro,compact_convergence_ball_center}.
    Therefore $A$ is open in $F$ with respect to $T_0$
        by \cref{compact_convergence_ball,subseteq,powerset_intro,topology_generated_by_basis,compact_convergence_topology,basis_topology_is_topology,basis_for_topology,open_in}.
\end{proof}

% from §46 helper lemma 8: positive reals have smaller positive reals below one
\begin{lemma}\label{positive_real_has_smaller_positive_below_one}
    Suppose $\epsilon\in\reals$.
    Suppose $0 < \epsilon$.
    Then there exists $\delta$ such that
        $\delta\in\reals$ and
        $0 < \delta$ and
        $\delta < 1$ and
        $\delta < \epsilon$.
\end{lemma}
\begin{proof}
    \begin{byCase}
    \caseOf{$\epsilon < 1$.}
        Let $\delta = \epsilon / (1 + 1)$.
        $\delta\in\reals$ and $0 < \delta$ and $\delta + \delta = \epsilon$
            by \cref{real_one_in_naturals,real_naturals_subset_reals,subseteq,real_zero_lt_one,real_half_of_positive}.
        Hence $\delta < \epsilon$ by \cref{real_add_ident,real_lt_add,real_lt_subst_left,real_lt_subst_right}.
        Therefore $\delta < 1$ by \cref{real_one_in_naturals,real_naturals_subset_reals,subseteq,real_lt_trans}.
        Follows by definition.
    \caseOf{it is wrong that $\epsilon < 1$.}
        Let $\delta = 1 / (1 + 1)$.
        $1\in\reals$ and $0 < 1$ by \cref{real_one_in_naturals,real_naturals_subset_reals,subseteq,real_zero_lt_one}.
        Hence $\delta\in\reals$ and $0 < \delta$ and $\delta + \delta = 1$ by \cref{real_half_of_positive}.
        Therefore $\delta < 1$ by \cref{real_add_ident,real_lt_add,real_lt_subst_left,real_lt_subst_right}.
        $1 < \epsilon$ or $1 = \epsilon$ by \cref{real_lt_trichotomy}.
        Hence $\delta < \epsilon$ by \cref{real_lt_trans,real_lt_subst_right}.
        Follows by definition.
    \end{byCase}
\end{proof}

% from §46 helper lemma 9: equal sets preserve supremum witnesses
\begin{lemma}\label{equal_set_supremum_has_witness}
    Assume $A = B$.
    Assume $r$ is a supremum of $A$.
    Then there exists $C$ such that
        $C = B$ and
        $r$ is a supremum of $C$.
\end{lemma}
\begin{proof}
    There exists $C$ such that $C = B$ and $r$ is a supremum of $C$ by \cref{pair_eq_iff}.
\end{proof}

% from §46 helper lemma 10: inhabited domains admit uniform metrics on function spaces
\begin{lemma}\label{uniform_metric_on_function_space_exists_inhabited}
    Assume $X$ is inhabited.
    Assume $d$ is a metric on $Y$.
    Then there exists $\rho$ such that
        $\rho$ is the uniform metric on functions from $X$ to $Y$ with respect to $d$.
\end{lemma}
\begin{proof}
    Let $F = \funs{X}{Y}$.
    Let $e = \boundedMetric{d}$.
    $e$ is a metric on $Y$ by \cref{bounded_metric_is_metric}.
    Take $x_0$ such that $x_0\in X$.
    Let $R = \{w\in(F\times F)\times\reals \mid
        \text{there exists $A$ such that
            $A$ is the distance set of $\fst{\fst{w}}$ and $\snd{\fst{w}}$ on $X$ with respect to $e$ and
            $\snd{w}$ is a supremum of $A$}\}$.
    $R$ is a relation by \cref{relation,subseteq,times_elem_is_tuple}.
    Show that for all $p\in F\times F$ there exists $r$ such that $p\mathrel{R} r$.
    \begin{subproof}
        Fix $p$ such that $p\in F\times F$.
        Take $f, g$ such that $p = (f,g)$ and $f\in F$ and $g\in F$ by \cref{times}.
        Let $A = \{\apply{e}{(f(x),g(x))} \mid x\in X\}$.
        $A$ is the distance set of $f$ and $g$ on $X$ with respect to $e$ by \cref{function_distance_set}.
        $A\subseteq\reals$
            by \cref{function_distance_set,metric_on,funs_type_apply,times_tuple_intro,pair_eq_iff,subseteq}.
        We have $A\neq\emptyset$
            by \cref{funs_type_apply,times_tuple_intro,function_distance_set,inhabited_nonempty,nonempty_inhabited}.
        For all $r$ if $r\in A$ then $r \leq 1$
            by \cref{function_distance_set,bounded_metric_leq_one,pair_eq_iff,metric_on,funs_type_apply,times_tuple_intro}.
        We have $1$ is an upper bound of $A$ by \cref{real_one_in_naturals,real_naturals_subset_reals,subseteq,real_upper_bound}.
        Hence $A$ is bounded above by \cref{real_bounded_above}.
        Take $r$ such that $r$ is a supremum of $A$ by \cref{real_completeness_axiom}.
        $r\in\reals$ by \cref{real_supremum,real_upper_bound}.
        There exists $s$ such that $p\mathrel{R} s$
            by \cref{times_tuple_intro,fst_eq,snd_eq,relation}.
    \end{subproof}
    Take $\rho$ such that
        $\rho$ is a function on $F\times F$ and
        for all $p\in F\times F$ we have $p\mathrel{R}\rho(p)$ by \cref{choice_function}.
    For all $p\in F\times F$ we have $\rho(p)\in\reals$
        by \cref{choice_function,relation,subseteq,times_tuple_elim}.
    We have $\rho$ is a function to $\reals$ by \cref{choice_function}.
    $\dom{\rho} = F\times F$ by \cref{choice_function}.
    Hence $\rho\in\funs{F\times F}{\reals}$ by \cref{funs_intro}.
    For all $f,g$ such that $f\in F$ and $g\in F$ we have $\rho((f,g)) \geq 0$
        by \cref{times_tuple_intro,choice_function,relation,fst_eq,snd_eq,funs_type_apply,function_distance_set,metric_on,real_supremum,real_upper_bound,metric_nonnegative,real_add_ident,real_leq_trans}.
    Hence $\rho$ is nonnegative on $F$ by \cref{metric_nonnegative}.
    Show that for all $f,g$ such that $f\in F$ and $g\in F$ we have $\rho((f,g)) = 0$ iff $f = g$.
    \begin{subproof}
        Fix $f$.
        Fix $g$ such that $f\in F$ and $g\in F$.
        We have $f\in\funs{X}{Y}$ and $g\in\funs{X}{Y}$ by \cref{subseteq}.
        $f$ is a function and $\dom{f} = X$ and $g$ is a function and $\dom{g} = X$
            by \cref{funs_is_function,funs_elim}.
        Show that if $\rho((f,g)) = 0$ then $f = g$.
        \begin{subproof}
            Assume $\rho((f,g)) = 0$.
            Let $p = (f,g)$.
            We have $(p,\rho(p))\in R$ by \cref{times_tuple_intro,choice_function,relation}.
            Take $A$ such that
                $A$ is the distance set of $\fst{p}$ and $\snd{p}$ on $X$ with respect to $e$ and
                $\rho(p)$ is a supremum of $A$ by \cref{fst_eq,snd_eq,relation}.
            For all $x\in X$ we have $f(x) = g(x)$
                by \cref{fst_eq,snd_eq,function_distance_set,real_supremum,real_upper_bound,funs_type_apply,times_tuple_intro,metric_on,metric_nonnegative,real_add_ident,real_leq_antisym,metric_zero_characterization}.
            Therefore $f = g$ by \cref{pair_eq_iff,functions_eq_iff}.
        \end{subproof}
        Show that if $f = g$ then $\rho((f,g)) = 0$.
        \begin{subproof}
            Assume $f = g$.
            Let $p = (f,g)$.
            We have $(p,\rho(p))\in R$ by \cref{times_tuple_intro,choice_function,relation}.
            Take $A$ such that
                $A$ is the distance set of $\fst{p}$ and $\snd{p}$ on $X$ with respect to $e$ and
                $\rho(p)$ is a supremum of $A$ by \cref{fst_eq,snd_eq,relation}.
            We have $A\subseteq\reals$
                by \cref{fst_eq,snd_eq,function_distance_set,metric_on,funs_type_apply,times_tuple_intro,pair_eq_iff,subseteq}.
            For all $r$ if $r\in A$ then $r = 0$
                by \cref{fst_eq,snd_eq,function_distance_set,metric_on,metric_zero_characterization,funs_type_apply,times_tuple_intro,pair_eq_iff,subseteq}.
            $0\in\reals$ by \cref{real_add_ident}.
            We have $0$ is an upper bound of $A$ by \cref{real_upper_bound,subseteq}.
            We have $e((f(x_0),g(x_0))) = 0$
                by \cref{metric_on,metric_zero_characterization,funs_type_apply,times_tuple_intro,pair_eq_iff,subseteq}.
            Therefore $0\in A$ by \cref{fst_eq,snd_eq,function_distance_set,pair_eq_iff,subseteq}.
            We have $0 \leq \rho(p)$ by \cref{real_supremum_pred,real_supremum,real_upper_bound}.
            Hence $\rho(p) \leq 0$ by \cref{real_supremum,real_upper_bound}.
            We have $\rho(p)\in\reals$ by \cref{real_supremum,real_upper_bound}.
            Thus $\rho(p) = 0$ by \cref{real_leq_antisym,real_add_ident}.
        \end{subproof}
    \end{subproof}
    Therefore $\rho$ is zero-characterized on $F$ by \cref{metric_zero_characterization}.
    Show that for all $f,g$ such that $f\in F$ and $g\in F$ we have $\rho((f,g)) = \rho((g,f))$.
    \begin{subproof}
        Fix $f$.
        Fix $g$ such that $f\in F$ and $g\in F$.
        We have $f\in\funs{X}{Y}$ and $g\in\funs{X}{Y}$ by \cref{subseteq}.
        Let $p = (f,g)$.
        Let $q = (g,f)$.
        We have $(p,\rho(p))\in R$ and $(q,\rho(q))\in R$ by \cref{times_tuple_intro,choice_function,relation}.
        Take $A$ such that
            $A$ is the distance set of $\fst{p}$ and $\snd{p}$ on $X$ with respect to $e$ and
            $\rho(p)$ is a supremum of $A$ by \cref{fst_eq,snd_eq,relation}.
        Take $C$ such that
            $C$ is the distance set of $\fst{q}$ and $\snd{q}$ on $X$ with respect to $e$ and
            $\rho(q)$ is a supremum of $C$ by \cref{fst_eq,snd_eq,relation}.
        $\fst{p} = f$ and $\snd{p} = g$ and $\fst{q} = g$ and $\snd{q} = f$ by \cref{fst_eq,snd_eq}.
        Show that for all $r$ if $r\in A$ then $r\in C$.
        \begin{subproof}
            Follows by \cref{function_distance_set,metric_on,metric_symmetric,times_tuple_intro,funs_type_apply,pair_eq_iff}.
        \end{subproof}
        We have $A\subseteq C$ by \cref{subseteq}.
        Show that for all $r$ if $r\in C$ then $r\in A$.
        \begin{subproof}
            Fix $r$.
            Assume $r\in C$.
            Take $x\in X$ such that $e((g(x),f(x))) = r$ by \cref{function_distance_set}.
            We have $g(x)\in Y$ and $f(x)\in Y$ by \cref{funs_type_apply}.
            Hence $(g(x),f(x))\in Y\times Y$ and $(f(x),g(x))\in Y\times Y$ by \cref{times_tuple_intro}.
            Therefore $e((f(x),g(x))) = r$ by \cref{metric_on,metric_symmetric,pair_eq_iff,real_add_ident,real_lt_subst_left,real_lt_subst_right}.
            Hence $r\in A$ by \cref{function_distance_set}.
        \end{subproof}
        Thus $C\subseteq A$ by \cref{subseteq}.
        Therefore $A = C$ by \cref{subseteq_antisymmetric}.
        Hence $\rho((f,g)) = \rho((g,f))$ by \cref{real_supremum_unique}.
    \end{subproof}
    Therefore $\rho$ is symmetric on $F$ by \cref{metric_symmetric}.
    Show that for all $f,g,h$ such that $f\in F$ and $g\in F$ and $h\in F$ we have
        $\rho((f,g)) + \rho((g,h)) \geq \rho((f,h))$.
    \begin{subproof}
        Fix $f$.
        Fix $g$.
        Fix $h$ such that $f\in F$ and $g\in F$ and $h\in F$.
        Let $p = (f,h)$.
        Let $q = (f,g)$.
        Let $s = (g,h)$.
        We have $(p,\rho(p))\in R$ and $(q,\rho(q))\in R$ and $(s,\rho(s))\in R$
            by \cref{times_tuple_intro,choice_function,relation}.
        Take $A$ such that
            $A$ is the distance set of $\fst{p}$ and $\snd{p}$ on $X$ with respect to $e$ and
            $\rho(p)$ is a supremum of $A$ by \cref{fst_eq,snd_eq,relation}.
        Take $C$ such that
            $C$ is the distance set of $\fst{q}$ and $\snd{q}$ on $X$ with respect to $e$ and
            $\rho(q)$ is a supremum of $C$ by \cref{fst_eq,snd_eq,relation}.
        Take $D$ such that
            $D$ is the distance set of $\fst{s}$ and $\snd{s}$ on $X$ with respect to $e$ and
            $\rho(s)$ is a supremum of $D$ by \cref{fst_eq,snd_eq,relation}.
        Hence $\fst{p} = f$ and $\snd{p} = h$ and $\fst{q} = f$ and $\snd{q} = g$ and $\fst{s} = g$ and $\snd{s} = h$
            by \cref{fst_eq,snd_eq}.
        $A\subseteq\reals$
            by \cref{function_distance_set,metric_on,funs_type_apply,times_tuple_intro,pair_eq_iff,subseteq}.
        $(f,g)\in F\times F$ and $(g,h)\in F\times F$ by \cref{times_tuple_intro}.
        We have $\rho((f,g))\in\reals$ and $\rho((g,h))\in\reals$ by \cref{funs_type_apply}.
        Show that for all $r\in A$ we have $r \leq \rho((f,g)) + \rho((g,h))$.
        \begin{subproof}
            Fix $r$ such that $r\in A$.
            Take $x\in X$ such that $r = \apply{e}{(f(x),h(x))}$ by \cref{function_distance_set}.
            We have $f(x)\in Y$ and $g(x)\in Y$ and $h(x)\in Y$ by \cref{funs_type_apply}.
            Hence $(f(x),g(x))\in Y\times Y$ and $(g(x),h(x))\in Y\times Y$ and $(f(x),h(x))\in Y\times Y$
                by \cref{times_tuple_intro}.
            We have $\apply{e}{(f(x),g(x))}\in\reals$ and $\apply{e}{(g(x),h(x))}\in\reals$ and $\apply{e}{(f(x),h(x))}\in\reals$
                by \cref{metric_on,funs_type_apply}.
            Therefore $\apply{e}{(f(x),h(x))} \leq \apply{e}{(f(x),g(x))} + \apply{e}{(g(x),h(x))}$
                by \cref{metric_on,metric_triangle_inequality}.
        We have $\apply{e}{(f(x),g(x))}\in C$ and $\apply{e}{(g(x),h(x))}\in D$ by \cref{function_distance_set}.
            We have $\rho((f,g))$ is an upper bound of $C$ and $\rho((g,h))$ is an upper bound of $D$ by \cref{real_supremum}.
            Therefore $\apply{e}{(f(x),g(x))} \leq \rho((f,g))$ and $\apply{e}{(g(x),h(x))} \leq \rho((g,h))$
                by \cref{real_upper_bound}.
            Hence $\apply{e}{(f(x),g(x))} + \apply{e}{(g(x),h(x))} \leq \rho((f,g)) + \apply{e}{(g(x),h(x))}$
                by \cref{real_leq_add}.
            We have $\apply{e}{(g(x),h(x))} + \rho((f,g)) \leq \rho((g,h)) + \rho((f,g))$ by \cref{real_leq_add}.
            Therefore $\rho((f,g)) + \apply{e}{(g(x),h(x))} \leq \rho((f,g)) + \rho((g,h))$ by \cref{real_add_comm}.
            We have $\apply{e}{(f(x),g(x))} + \apply{e}{(g(x),h(x))}\in\reals$ and
                $\rho((f,g)) + \apply{e}{(g(x),h(x))}\in\reals$ and
                $\rho((f,g)) + \rho((g,h))\in\reals$ by \cref{real_add_closed}.
            Hence $\apply{e}{(f(x),g(x))} + \apply{e}{(g(x),h(x))} \leq \rho((f,g)) + \rho((g,h))$ by \cref{real_leq_trans}.
            Hence $\apply{e}{(f(x),h(x))} \leq \rho((f,g)) + \rho((g,h))$ by \cref{real_leq_trans}.
            Thus $r \leq \rho((f,g)) + \rho((g,h))$ by \cref{pair_eq_iff}.
        \end{subproof}
        We have $\rho((f,g)) + \rho((g,h))$ is an upper bound of $A$ by \cref{real_upper_bound,real_add_closed}.
        Hence $\rho((f,h)) \leq \rho((f,g)) + \rho((g,h))$ by \cref{real_supremum}.
        Therefore $\rho((f,g)) + \rho((g,h)) \geq \rho((f,h))$.
    \end{subproof}
    Therefore $\rho$ satisfies the triangle inequality on $F$ by \cref{metric_triangle_inequality}.
    Hence $\rho$ is a metric on $F$ by \cref{metric_on}.
    Show that for all $f,g\in F$ there exists $A$ such that
        $A = \{r\in\reals \mid \exists \alpha\in X.
            ((r = d((\apply{f}{\alpha},\apply{g}{\alpha})) \land d((\apply{f}{\alpha},\apply{g}{\alpha})) < 1) \lor
            (r = 1 \land 1 \leq d((\apply{f}{\alpha},\apply{g}{\alpha}))))\}$ and
        $\rho((f,g))$ is a supremum of $A$.
    \begin{subproof}
        Fix $f$.
        Fix $g$ such that $f\in F$ and $g\in F$.
        Let $p = (f,g)$.
        We have $(p,\rho(p))\in R$ by \cref{times_tuple_intro,choice_function,relation}.
        Take $D$ such that
            $D$ is the distance set of $\fst{p}$ and $\snd{p}$ on $X$ with respect to $e$ and
            $\rho(p)$ is a supremum of $D$ by \cref{fst_eq,snd_eq,relation}.
        $\fst{p} = f$ and $\snd{p} = g$ by \cref{fst_eq,snd_eq}.
        Let $B = \{r\in\reals \mid \exists \alpha\in X.
            ((r = d((\apply{f}{\alpha},\apply{g}{\alpha})) \land d((\apply{f}{\alpha},\apply{g}{\alpha})) < 1) \lor
            (r = 1 \land 1 \leq d((\apply{f}{\alpha},\apply{g}{\alpha}))))\}$.
        Show that for all $r$ if $r\in D$ then $r\in B$.
        \begin{subproof}
                Fix $r$.
                Assume $r\in D$.
                Take $\alpha\in X$ such that $r = \apply{e}{(f(\alpha),g(\alpha))}$ by \cref{function_distance_set}.
                We have $f(\alpha)\in Y$ and $g(\alpha)\in Y$ by \cref{funs_type_apply}.
                Hence $(f(\alpha),g(\alpha))\in Y\times Y$ by \cref{times_tuple_intro}.
                We have $e\in\funs{Y\times Y}{\reals}$ by \cref{metric_on}.
                \begin{byCase}
                \caseOf{$d((f(\alpha),g(\alpha))) < 1$.}
                    We have $\apply{e}{(f(\alpha),g(\alpha))} = d((f(\alpha),g(\alpha)))$
                        by \cref{bounded_metric_apply_below_one}.
                    Hence $r = d((f(\alpha),g(\alpha)))$ by \cref{real_lt_subst_right}.
                    Therefore $r = d((\apply{f}{\alpha},\apply{g}{\alpha})) \land d((\apply{f}{\alpha},\apply{g}{\alpha})) < 1$.
                    We have $d((f(\alpha),g(\alpha)))\in\reals$ by \cref{metric_on,funs_type_apply}.
                    Hence $r\in\reals$ by \cref{real_lt_subst_right}.
                    Therefore $r\in B$ by \cref{subseteq}.
                \caseOf{it is wrong that $d((f(\alpha),g(\alpha))) < 1$.}
                    We have $\apply{e}{(f(\alpha),g(\alpha))} = 1$
                        by \cref{bounded_metric_apply_geq_one}.
                    Hence $r = 1$ by \cref{real_lt_subst_right}.
                    $1\in\reals$ by \cref{real_one_in_naturals,real_naturals_subset_reals,subseteq}.
                    We have $d((f(\alpha),g(\alpha)))\in\reals$ by \cref{metric_on,funs_type_apply}.
                    Thus $1 = d((f(\alpha),g(\alpha)))$ or $1 < d((f(\alpha),g(\alpha)))$ by \cref{real_lt_trichotomy}.
                    Therefore $1 \leq d((f(\alpha),g(\alpha)))$ by \cref{pair_eq_iff,real_lt_imp_leq}.
                    We have $r = 1 \land 1 \leq d((\apply{f}{\alpha},\apply{g}{\alpha}))$.
                    Hence $r\in\reals$ by \cref{real_lt_subst_right}.
                    Therefore $r\in B$ by \cref{subseteq}.
                \end{byCase}
        \end{subproof}
        We have $D\subseteq B$ by \cref{subseteq}.
        Show that for all $r$ if $r\in B$ then $r\in D$.
        \begin{subproof}
                Fix $r$.
                Assume $r\in B$.
                Take $\alpha\in X$ such that
                    $(r = d((\apply{f}{\alpha},\apply{g}{\alpha})) \land d((\apply{f}{\alpha},\apply{g}{\alpha})) < 1) \lor
                    (r = 1 \land 1 \leq d((\apply{f}{\alpha},\apply{g}{\alpha})))$ by \cref{subseteq}.
                We have $f(\alpha)\in Y$ and $g(\alpha)\in Y$ by \cref{funs_type_apply}.
                Hence $(f(\alpha),g(\alpha))\in Y\times Y$ by \cref{times_tuple_intro}.
                We have $e\in\funs{Y\times Y}{\reals}$ by \cref{metric_on}.
                \begin{byCase}
                \caseOf{$r = d((\apply{f}{\alpha},\apply{g}{\alpha})) \land d((\apply{f}{\alpha},\apply{g}{\alpha})) < 1$.}
                    We have $\apply{e}{(f(\alpha),g(\alpha))} = d((f(\alpha),g(\alpha)))$ by \cref{bounded_metric_apply_below_one}.
                    Hence $r = \apply{e}{(f(\alpha),g(\alpha))}$ by \cref{real_lt_subst_left}.
                    Thus $r\in D$ by \cref{function_distance_set}.
                \caseOf{$r = 1 \land 1 \leq d((\apply{f}{\alpha},\apply{g}{\alpha}))$.}
                    $1\in\reals$ by \cref{real_one_in_naturals,real_naturals_subset_reals,subseteq}.
                    We have $d((f(\alpha),g(\alpha)))\in\reals$ by \cref{metric_on,funs_type_apply}.
                    Hence it is wrong that $d((f(\alpha),g(\alpha))) < 1$ by \cref{real_lt_trans,real_lt_irreflexive}.
                    Therefore $\apply{e}{(f(\alpha),g(\alpha))} = 1$ by \cref{bounded_metric_apply_geq_one}.
                    Hence $r = \apply{e}{(f(\alpha),g(\alpha))}$ by \cref{pair_eq_iff}.
                    Thus $r\in D$ by \cref{function_distance_set}.
                \end{byCase}
        \end{subproof}
        We have $B\subseteq D$ by \cref{subseteq}.
        Hence $D = B$ by \cref{subseteq_antisymmetric}.
        We have $\rho((f,g))$ is a supremum of $B$ by \cref{real_add_eq_subst}.
        Thus there exists $A$ such that
            $A = \{r\in\reals \mid \exists \alpha\in X.
                ((r = d((\apply{f}{\alpha},\apply{g}{\alpha})) \land d((\apply{f}{\alpha},\apply{g}{\alpha})) < 1) \lor
                (r = 1 \land 1 \leq d((\apply{f}{\alpha},\apply{g}{\alpha}))))\}$ and
            $\rho((f,g))$ is a supremum of $A$ by \cref{pair_eq_iff}.
        Follows by definition.
    \end{subproof}
    Therefore there exists $\tau$ such that $\tau$ is the uniform metric on functions from $X$ to $Y$
        with respect to $d$ by \cref{pair_eq_iff,uniform_metric_on_function_space}.
\end{proof}

% from §46 Item 4: compact convergence is uniform convergence on compact subspaces
\begin{theorem}\label{compact_convergence_iff_uniform_on_compact_subspaces}
    Assume $T$ is a topology on $X$.
    Assume $d$ is a metric on $Y$.
    Assume $T_0$ is the topology of compact convergence on functions from $X$ to $Y$ with respect to $T$ and $d$.
    Assume $s$ is a sequence in $\funs{X}{Y}$.
    Assume $f\in\funs{X}{Y}$.
    Then $s$ converges to $f$ in $\funs{X}{Y}$ with respect to $T_0$ iff
        for all $C$ if
            $C$ is a compact subspace of $X$ with respect to $T$
        then there exists $t$ such that
            $t$ is the restriction sequence of $s$ to $C$ from $X$ to $Y$ and
            $t$ converges uniformly to $\restrl{f}{C}$ from $C$ to $Y$ with respect to $d$.
\end{theorem}
\begin{proof}
    Let $F = \funs{X}{Y}$.
    Show that if $s$ converges to $f$ in $F$ with respect to $T_0$ then
        for all $C$ if
            $C$ is a compact subspace of $X$ with respect to $T$
        then there exists $t$ such that
            $t$ is the restriction sequence of $s$ to $C$ from $X$ to $Y$ and
            $t$ converges uniformly to $\restrl{f}{C}$ from $C$ to $Y$ with respect to $d$.
    \begin{subproof}
        Assume $s$ converges to $f$ in $F$ with respect to $T_0$.
        Show that for all $C$ if
            $C$ is a compact subspace of $X$ with respect to $T$
        then there exists $t$ such that
            $t$ is the restriction sequence of $s$ to $C$ from $X$ to $Y$ and
            $t$ converges uniformly to $\restrl{f}{C}$ from $C$ to $Y$ with respect to $d$.
        \begin{subproof}
            Fix $C$ such that $C$ is a compact subspace of $X$ with respect to $T$.
            We have $C\subseteq X$ by \cref{compact_subspace}.
            Take $t$ such that $t$ is the restriction sequence of $s$ to $C$ from $X$ to $Y$
                by \cref{restriction_sequence_exists}.
            Show that $t$ converges uniformly to $\restrl{f}{C}$ from $C$ to $Y$ with respect to $d$.
            \begin{subproof}
                \begin{byCase}
                \caseOf{$C = \emptyset$.}
                    Show that for all $\epsilon$ if $\epsilon\in\reals$ and $0 < \epsilon$ then
                        there exists $N\in\naturalsPlus$ such that for all $n\in\naturalsPlus$ if $N\leq n$ then
                            for all $x\in C$ we have $d((\apply{\apply{t}{n}}{x}, \apply{\restrl{f}{C}}{x})) < \epsilon$.
                    \begin{subproof}
                        Trivial.
                    \end{subproof}
                    Hence $t$ converges uniformly to $\restrl{f}{C}$ from $C$ to $Y$ with respect to $d$
                        by \cref{uniform_converges_to,restriction_sequence_on_subset,restrl_typed_function}.
                \caseOf{$C \neq \emptyset$.}
                    Take $\rho$ such that
                        $\rho$ is the uniform metric on functions from $C$ to $Y$ with respect to $d$
                        by \cref{nonempty_inhabited,uniform_metric_on_function_space_exists_inhabited}.
                    We have $\rho$ is a metric on $\funs{C}{Y}$ by \cref{uniform_metric_on_function_space}.
                    Show that for all $\epsilon$ if $\epsilon\in\reals$ and $0 < \epsilon$ then
                        there exists $N\in\naturalsPlus$ such that for all $n\in\naturalsPlus$ if $N\leq n$ then
                            for all $x\in C$ we have $d((\apply{\apply{t}{n}}{x}, \apply{\restrl{f}{C}}{x})) < \epsilon$.
                    \begin{subproof}
                        Fix $\epsilon$ such that $\epsilon\in\reals$ and $0 < \epsilon$.
                        Take $\delta$ such that
                            $\delta\in\reals$ and
                            $0 < \delta$ and
                            $\delta < 1$ and
                            $\delta < \epsilon$ by \cref{positive_real_has_smaller_positive_below_one}.
                        Let $A = \{g\in F \mid \rho((\restrl{f}{C},\restrl{g}{C})) < \delta\}$.
                        We have $A$ is the compact-convergence ball about $f$ over $C$ with radius $\delta$
                            from $X$ to $Y$ with respect to $T$ and $d$ by \cref{compact_convergence_ball}.
                        Hence $A$ is a neighborhood of $f$ in $F$ with respect to $T_0$
                            by \cref{compact_convergence_ball_open,compact_convergence_ball_center,neighborhood}.
                        Take $N\in\naturalsPlus$ such that for all $n$ such that $n\in\naturalsPlus$ if $N\leq n$ then $s(n)\in A$
                            by \cref{converges_to}.
                        Show that for all $n\in\naturalsPlus$ if $N\leq n$ then
                            for all $x\in C$ we have $d((\apply{\apply{t}{n}}{x}, \apply{\restrl{f}{C}}{x})) < \epsilon$.
                        \begin{subproof}
                            Fix $n\in\naturalsPlus$.
                            Assume $N\leq n$.
                            We have $s(n)\in A$.
                            Hence $\rho((\restrl{f}{C},\restrl{\apply{s}{n}}{C})) < \delta$ by \cref{subseteq}.
                            We have $t(n) = \restrl{s(n)}{C}$ by \cref{restriction_sequence_on_subset}.
                            We have $\restrl{s(n)}{C}\in\funs{C}{Y}$
                                by \cref{sequence_in,funs_type_apply,restrl_typed_function}.
                            Hence $\restrl{f}{C}\in\funs{C}{Y}$ by \cref{restrl_typed_function,compact_subspace,subseteq}.
                            We have $t(n)\in\funs{C}{Y}$ by \cref{real_lt_subst_right}.
                            Therefore $\rho((\restrl{\apply{s}{n}}{C},\restrl{f}{C})) = \rho((\restrl{f}{C},\restrl{\apply{s}{n}}{C}))$
                                by \cref{uniform_metric_on_function_space,metric_on,metric_symmetric,times_tuple_intro}.
                            Hence $\rho((t(n),\restrl{f}{C})) < \delta$ by \cref{real_lt_subst_right,real_lt_subst_left}.
                            Therefore for all $x\in C$ we have $d((\apply{\apply{t}{n}}{x}, \apply{\restrl{f}{C}}{x})) < \epsilon$
                                by \cref{uniform_metric_small_implies_pointwise_small,times_tuple_intro,funs_type_apply,metric_on,real_lt_trans}.
                            Follows by definition.
                        \end{subproof}
                        Follows by definition.
                    \end{subproof}
                    Therefore $t$ converges uniformly to $\restrl{f}{C}$ from $C$ to $Y$ with respect to $d$
                        by \cref{uniform_converges_to,restriction_sequence_on_subset,restrl_typed_function}.
                \end{byCase}
            \end{subproof}
            Follows by definition.
        \end{subproof}
        Follows by definition.
    \end{subproof}
    Show that if
        for all $C$ if
            $C$ is a compact subspace of $X$ with respect to $T$
        then there exists $t$ such that
            $t$ is the restriction sequence of $s$ to $C$ from $X$ to $Y$ and
            $t$ converges uniformly to $\restrl{f}{C}$ from $C$ to $Y$ with respect to $d$
    then $s$ converges to $f$ in $F$ with respect to $T_0$.
    \begin{subproof}
        Assume for all $C$ if
            $C$ is a compact subspace of $X$ with respect to $T$
        then there exists $t$ such that
            $t$ is the restriction sequence of $s$ to $C$ from $X$ to $Y$ and
            $t$ converges uniformly to $\restrl{f}{C}$ from $C$ to $Y$ with respect to $d$.
        Show that for all $U$ such that $U$ is a neighborhood of $f$ in $F$ with respect to $T_0$
            there exists $N\in\naturalsPlus$ such that for all $n$ such that $n\in\naturalsPlus$ if $N\leq n$ then $s(n)\in U$.
        \begin{subproof}
            Fix $U$ such that $U$ is a neighborhood of $f$ in $F$ with respect to $T_0$.
            We have $U$ is open in $F$ with respect to $T_0$ and $f\in U$ by \cref{neighborhood}.
            Take $B$ such that
                $B$ is the compact-convergence basis on functions from $X$ to $Y$ with respect to $T$ and $d$ and
                $T_0 = \basisTopology{B}{F}$ by \cref{compact_convergence_topology}.
            Take $A\in B$ such that $f\in A$ and $A\subseteq U$
                by \cref{basis_for_topology,open_in,topology_generated_by_basis}.
            Take $h, C, \epsilon$ such that
                $A$ is the compact-convergence ball about $h$ over $C$ with radius $\epsilon$
                    from $X$ to $Y$ with respect to $T$ and $d$ by \cref{compact_convergence_function_space_basis}.
            Take $\delta, D$ such that
                $\delta\in\reals$ and
                $0 < \delta$ and
                $D$ is the compact-convergence ball about $f$ over $C$ with radius $\delta$
                    from $X$ to $Y$ with respect to $T$ and $d$ and
                $D\subseteq A$ by \cref{compact_convergence_ball_refinement}.
            We have $C$ is a compact subspace of $X$ with respect to $T$ by \cref{compact_convergence_ball}.
            Take $t$ such that
                $t$ is the restriction sequence of $s$ to $C$ from $X$ to $Y$ and
                $t$ converges uniformly to $\restrl{f}{C}$ from $C$ to $Y$ with respect to $d$ by \cref{subseteq}.
            Take $\rho$ such that
                $\rho$ is the uniform metric on functions from $C$ to $Y$ with respect to $d$ and
                $D = \{g\in F \mid \rho((\restrl{f}{C},\restrl{g}{C})) < \delta\}$ by \cref{compact_convergence_ball}.
            We have $\rho$ is a metric on $\funs{C}{Y}$ by \cref{uniform_metric_on_function_space}.
            Take $\eta$ such that
                $\eta\in\reals$ and
                $0 < \eta$ and
                $\eta < 1$ and
                $\eta < \delta$ by \cref{positive_real_has_smaller_positive_below_one}.
            Take $N\in\naturalsPlus$ such that for all $n\in\naturalsPlus$ if $N\leq n$ then
                for all $x\in C$ we have $d((\apply{\apply{t}{n}}{x}, \apply{\restrl{f}{C}}{x})) < \eta$
                by \cref{uniform_converges_to}.
            Show that for all $n$ such that $n\in\naturalsPlus$ if $N\leq n$ then $s(n)\in U$.
            \begin{subproof}
                Fix $n$ such that $n\in\naturalsPlus$.
                Assume $N\leq n$.
                We have $t(n)\in\funs{C}{Y}$ by \cref{restriction_sequence_on_subset,sequence_in,funs_type_apply}.
                $\restrl{f}{C}\in\funs{C}{Y}$ by \cref{restrl_typed_function,compact_subspace,subseteq}.
                For all $x\in C$ we have $d((\apply{\apply{t}{n}}{x}, \apply{\restrl{f}{C}}{x})) < \eta$
                    by \cref{uniform_converges_to}.
                Hence for all $x\in C$ we have $\apply{\boundedMetric{d}}{(\apply{\apply{t}{n}}{x}, \apply{\restrl{f}{C}}{x})} \leq \eta$
                    by \cref{times_tuple_intro,funs_type_apply,metric_on,real_one_in_naturals,real_naturals_subset_reals,subseteq,real_lt_trans,bounded_metric_apply_below_one,real_lt_subst_left,real_lt_imp_leq}.
                We have $\rho((t(n),\restrl{f}{C})) \leq \eta$ by \cref{uniform_metric_leq_from_coordinate_bounds}.
                Hence $\rho((\restrl{f}{C},t(n))) = \rho((t(n),\restrl{f}{C}))$
                    by \cref{uniform_metric_on_function_space,metric_on,metric_symmetric,times_tuple_intro}.
                We have $\rho((\restrl{f}{C},t(n)))\in\reals$
                    by \cref{metric_on,funs_type_apply,times_tuple_intro}.
                Therefore $\rho((\restrl{f}{C},t(n))) < \delta$
                    by \cref{real_leq_split,real_lt_trans,real_lt_subst_left}.
                We have $t(n) = \restrl{s(n)}{C}$ by \cref{restriction_sequence_on_subset}.
                Hence $\rho((\restrl{f}{C},\restrl{s(n)}{C})) < \delta$ by \cref{real_lt_subst_right}.
                Therefore $s(n)\in D$ by \cref{sequence_in,funs_type_apply,subseteq}.
                Hence $s(n)\in U$ by \cref{subseteq,subseteq_transitive}.
            \end{subproof}
            Follows by definition.
        \end{subproof}
        Thus $s$ converges to $f$ in $F$ with respect to $T_0$ by \cref{converges_to}.
    \end{subproof}
    Therefore $s$ converges to $f$ in $F$ with respect to $T_0$ iff
        for all $C$ if
            $C$ is a compact subspace of $X$ with respect to $T$
        then there exists $t$ such that
            $t$ is the restriction sequence of $s$ to $C$ from $X$ to $Y$ and
            $t$ converges uniformly to $\restrl{f}{C}$ from $C$ to $Y$ with respect to $d$.
\end{proof}

% from §46 Item 5: compactly generated space
\begin{definition}\label{compactly_generated_space}
    $X$ is compactly generated with respect to $T$ iff
        $T$ is a topology on $X$ and
        for all $A\subseteq X$ if
            for all $C$ if
                $C$ is a compact subspace of $X$ with respect to $T$
            then $A\inter C$ is open in $C$ with respect to $\subspaceTopology{T}{C}$
        then $A$ is open in $X$ with respect to $T$.
\end{definition}

% from §46 helper lemma 11: neighborhood criterion for openness
\begin{lemma}\label{open_from_local_neighborhoods}
    Assume $T$ is a topology on $X$.
    Suppose $A\subseteq X$.
    Suppose for all $x\in A$ there exists $U$ such that
        $U$ is a neighborhood of $x$ in $X$ with respect to $T$ and
        $U\subseteq A$.
    Then $A$ is open in $X$ with respect to $T$.
\end{lemma}
\begin{proof}
    Let $E = \{U\in\pow{X} \mid \text{$U$ is open in $X$ with respect to $T$ and $U\subseteq A$}\}$.
    We have for all $x$ such that $x\in A$ we have $x\in\unions{E}$
        by \cref{neighborhood,open_in,topology_on,subseteq,pow_iff,powerset_intro,unions_iff}.
    $\unions{E}\subseteq A$ by \cref{unions_iff,subseteq}.
    Therefore $A = \unions{E}$ by \cref{subseteq,unions_iff,subseteq_antisymmetric}.
    Hence $A$ is open in $X$ with respect to $T$
        by \cref{subseteq,unions_iff,subseteq_antisymmetric,open_in,topology_on}.
\end{proof}

% from §46 Item 6: locally compact spaces are compactly generated
\begin{lemma}\label{locally_compact_implies_compactly_generated}
    Assume $T$ is a topology on $X$.
    Suppose $X$ is locally compact with respect to $T$.
    Then $X$ is compactly generated with respect to $T$.
\end{lemma}
\begin{proof}
    Show that for all $A\subseteq X$ if
        for all $C$ if
            $C$ is a compact subspace of $X$ with respect to $T$
        then $A\inter C$ is open in $C$ with respect to $\subspaceTopology{T}{C}$
    then $A$ is open in $X$ with respect to $T$.
    \begin{subproof}
        Fix $A$ such that $A\subseteq X$.
        Assume for all $C$ if
            $C$ is a compact subspace of $X$ with respect to $T$
        then $A\inter C$ is open in $C$ with respect to $\subspaceTopology{T}{C}$.
        Show that for all $x\in A$ there exists $U$ such that
            $U$ is a neighborhood of $x$ in $X$ with respect to $T$ and
            $U\subseteq A$.
        \begin{subproof}
                Fix $x$ such that $x\in A$.
                Hence $X$ is locally compact at $x$ with respect to $T$ by \cref{subseteq,locally_compact}.
                Take $C$ such that
                    $C$ is a compact subspace of $X$ with respect to $T$ and
                    there exists $U$ such that
                        $U$ is a neighborhood of $x$ in $X$ with respect to $T$ and
                        $U\subseteq C$ by \cref{locally_compact_at_point}.
                Take $U$ such that
                    $U$ is a neighborhood of $x$ in $X$ with respect to $T$ and
                    $U\subseteq C$ by \cref{locally_compact_at_point}.
                We have $A\inter C$ is open in $C$ with respect to $\subspaceTopology{T}{C}$ by assumption.
                Take $V\in T$ such that $A\inter C = C\inter V$ by \cref{open_in,subspace_topology}.
                We have $x\in A\inter C$ by \cref{neighborhood,subseteq,inter_intro}.
                Hence $x\in V$ by \cref{inter}.
                We have $U\inter V$ is open in $X$ with respect to $T$
                    by \cref{neighborhood,open_in,topology_on,open_in}.
                Thus $U\inter V$ is a neighborhood of $x$ in $X$ with respect to $T$ by \cref{inter_intro,neighborhood}.
                Show that $U\inter V\subseteq A$.
                \begin{subproof}
                    Follows by \cref{inter,subseteq,inter_intro}.
                \end{subproof}
                Follows by definition.
        \end{subproof}
        Hence $A$ is open in $X$ with respect to $T$ by \cref{open_from_local_neighborhoods}.
    \end{subproof}
    Therefore $X$ is compactly generated with respect to $T$ by \cref{compactly_generated_space}.
\end{proof}

% from §46 helper lemma 12: a convergent sequence together with its limit is compact
\begin{lemma}\label{convergent_sequence_adjoin_limit_compact}
    Assume $T$ is a topology on $X$.
    Assume $s$ is a sequence in $X$.
    Assume $x\in X$.
    Suppose $s$ converges to $x$ in $X$ with respect to $T$.
    Let $C = \{x\}\union\img{s}{\naturalsPlus}$.
    Then $C$ is a compact subspace of $X$ with respect to $T$.
\end{lemma}
\begin{proof}
    We have $\img{s}{\naturalsPlus}\subseteq X$ by \cref{sequence_in,funs_ran,img_subseteq_ran,subseteq_transitive}.
    We have $C\subseteq X$ by \cref{singleton_subset_intro,union_subsets_is_subset}.
    Hence $C$ is a subspace of $X$ with respect to $T$ by \cref{subspace_of}.
    Show that for all $A$ such that
        $A$ is a covering of $C$ and
        for all $U$ such that $U\in A$ we have $U$ is open in $X$ with respect to $T$
    there exists $B$ such that
        $B\subseteq A$ and
        $B$ is finite and
        $B$ is a covering of $C$.
    \begin{subproof}
        Fix $A$ such that
            $A$ is a covering of $C$ and
            for all $U$ such that $U\in A$ we have $U$ is open in $X$ with respect to $T$.
        We have $x\in\unions{A}$ by \cref{singleton_intro,union_intro_left,covering,subseteq}.
        Take $U\in A$ such that $x\in U$ by \cref{unions_iff}.
        Take $N\in\naturalsPlus$ such that for all $n\in\naturalsPlus$ if $N\leq n$ then $s(n)\in U$
            by \cref{neighborhood,converges_to}.
        Let $I = \initialSegment{N}$.
        We have $I$ is finite by \cref{initial_segment_finite}.
        Let $R = \{w\in I\times A \mid \apply{s}{\fst{w}}\in\snd{w}\}$.
        We have $R$ is a relation by \cref{relation,subseteq,times_elem_is_tuple}.
        For all $n\in I$ we have $s(n)\in C$
            by \cref{initial_segment_naturalsplus,sequence_in,funs_is_function,funs,subseteq,inter_intro,img_of_function_intro,union_intro_right}.
        For all $n\in I$ there exists $V$ such that $n\mathrel{R} V$
            by \cref{covering,subseteq,unions_iff,times_tuple_intro,fst_eq,snd_eq,relation}.
        Take $g$ such that $g$ is a function on $I$ and for all $n\in I$ we have $n\mathrel{R}\apply{g}{n}$ by \cref{choice_function}.
        Let $D = \img{g}{I}$.
        Let $E = \{U,U\}$.
        Let $B = D\union E$.
        $D\subseteq A$
            by \cref{img_iff,choice_function,function_apply_iff,relation,subseteq,fst_eq,snd_eq,times,real_lt_subst_left}.
        We have $D$ is finite by \cref{finite_image_function}.
        Hence $E$ is finite by \cref{pair_is_finite}.
        Therefore $B$ is finite by \cref{union_of_finite_is_finite}.
        $E\subseteq A$ by \cref{upair_iff,subseteq}.
        We have $B\subseteq A$ by \cref{union_subsets_is_subset}.
        Show that for all $y$ such that $y\in C$ we have $y\in\unions{B}$.
        \begin{subproof}
            Fix $y$ such that $y\in C$.
            We have $y = x$ or $y\in\img{s}{\naturalsPlus}$ by \cref{union_iff,singleton_iff}.
            \begin{byCase}
            \caseOf{$y = x$.}
                Therefore $y\in U$ by \cref{real_lt_subst_left}.
                Hence $y\in\unions{B}$ by \cref{union_intro_right,upair_intro_left,unions_iff}.
            \caseOf{$y\in\img{s}{\naturalsPlus}$.}
                Take $n\in\naturalsPlus$ such that $(n,y)\in s$ by \cref{img_iff}.
                We have $s(n) = y$ by \cref{sequence_in,funs_is_function,function_apply_iff}.
                $n\in\reals$ and $N\in\reals$ by \cref{real_naturals_subset_reals,subseteq}.
                Hence $n < N$ or $N \leq n$ by \cref{real_lt_trichotomy,real_lt_imp_leq}.
                \begin{byCase}
                \caseOf{$n < N$.}
                    We have $n \leq N$ by \cref{real_lt_imp_leq}.
                    Hence $n\in I$ by \cref{initial_segment_naturalsplus}.
                    Therefore $n\mathrel{R}\apply{g}{n}$ by \cref{choice_function}.
                    We have $s(n)\in g(n)$ by \cref{relation,fst_eq,snd_eq}.
                    We have $n\in\dom{g}$ by \cref{choice_function}.
                    Hence $n\in\dom{g}\inter I$ by \cref{inter_intro}.
                    Therefore $g(n)\in D$ by \cref{img_of_function_intro}.
                    Hence $g(n)\in B$ by \cref{union_intro_left}.
                    Thus $y\in\unions{B}$ by \cref{real_lt_subst_right,unions_iff}.
                \caseOf{$N \leq n$.}
                    Therefore $s(n)\in U$ by \cref{converges_to}.
                    Hence $y\in\unions{B}$ by \cref{real_lt_subst_right,union_intro_right,upair_intro_left,unions_iff}.
                \end{byCase}
            \end{byCase}
        \end{subproof}
        Hence $B$ is a covering of $C$ by \cref{subseteq,covering}.
        Follows by definition.
    \end{subproof}
    Therefore $C$ is compact with respect to $\subspaceTopology{T}{C}$ by \cref{compact_subspace_open_covering_backward}.
    Hence $C$ is a compact subspace of $X$ with respect to $T$ by \cref{compact_subspace}.
\end{proof}

% from §46 Item 6: first-countable spaces are compactly generated
\begin{lemma}\label{first_countable_implies_compactly_generated}
    Assume $T$ is a topology on $X$.
    Suppose $X$ satisfies the first countability axiom with respect to $T$.
    Then $X$ is compactly generated with respect to $T$.
\end{lemma}
\begin{proof}
    Show that for all $A\subseteq X$ if
        for all $C$ if
            $C$ is a compact subspace of $X$ with respect to $T$
        then $A\inter C$ is open in $C$ with respect to $\subspaceTopology{T}{C}$
    then $A$ is open in $X$ with respect to $T$.
    \begin{subproof}
        Fix $A$ such that $A\subseteq X$.
        Assume for all $C$ if
            $C$ is a compact subspace of $X$ with respect to $T$
        then $A\inter C$ is open in $C$ with respect to $\subspaceTopology{T}{C}$.
        Show that for all $x\in A$ there exists $U$ such that
            $U$ is a neighborhood of $x$ in $X$ with respect to $T$ and
            $U\subseteq A$.
        \begin{subproof}
            Fix $x\in A$.
            Hence $x\in X$ by \cref{subseteq}.
            Suppose not.
            Hence for all $U$ if
                $U$ is a neighborhood of $x$ in $X$ with respect to $T$
            then it is wrong that $U\subseteq A$.
            We have $X\setminus A\subseteq X$ by \cref{setminus_subseteq}.
            Show that for all $U$ such that
                $U$ is open in $X$ with respect to $T$ and
                $x\in U$
            we have $U$ intersects $X\setminus A$.
            \begin{subproof}
                Fix $U$ such that
                    $U$ is open in $X$ with respect to $T$ and
                    $x\in U$.
                Hence it is wrong that $U\subseteq A$ by \cref{neighborhood}.
                Suppose not.
                Show that for all $y$ such that $y\in U$ we have $y\in A$.
                \begin{subproof}
                    Fix $y$ such that $y\in U$.
                    Suppose not.
                    Hence $y\in X\setminus A$ by \cref{open_in,topology_on,subseteq,pow_iff,subseteq,setminus_intro}.
                    Therefore $y\in U\inter (X\setminus A)$ by \cref{inter_intro}.
                    Contradiction.
                \end{subproof}
                Hence $U\subseteq A$ by \cref{subseteq}.
                Contradiction.
                Therefore $U$ intersects $X\setminus A$.
            \end{subproof}
            Therefore $x\in\closure{X\setminus A}{X}{T}$ by \cref{closure_iff_intersects_open}.
            Take $s$ such that
                $s$ is a sequence in $X\setminus A$ and
                $s$ converges to $x$ in $X$ with respect to $T$
                by \cref{closure_implies_sequence_first_countable}.
            Let $C = \{x\}\union\img{s}{\naturalsPlus}$.
            Hence $C$ is a compact subspace of $X$ with respect to $T$
                by \cref{sequence_in,funs_weaken_codom,setminus_subseteq,convergent_sequence_adjoin_limit_compact}.
            Therefore $A\inter C$ is open in $C$ with respect to $\subspaceTopology{T}{C}$ by assumption.
            Take $V\in T$ such that $A\inter C = C\inter V$ by \cref{open_in,subspace_topology}.
            We have $x\in A\inter C$ by \cref{singleton_intro,union_intro_left,inter_intro}.
            Take $N\in\naturalsPlus$ such that for all $n\in\naturalsPlus$ if $N\leq n$ then $s(n)\in V$
                by \cref{inter,open_in,neighborhood,converges_to}.
            We have $N \leq N$.
            Therefore $s(N)\in V$ by \cref{subseteq}.
            We have $s(N)\in X\setminus A$ by \cref{sequence_in,funs_type_apply}.
            $s$ is a function and $\dom{s} = \naturalsPlus$ by \cref{sequence_in,funs_is_function,funs}.
            Thus $N\in\dom{s}\inter\naturalsPlus$ by \cref{inter_intro}.
            Hence $s(N)\in\img{s}{\naturalsPlus}$ by \cref{img_of_function_intro}.
            Therefore $s(N)\in C\inter V$ by \cref{union_intro_right,inter_intro}.
            Hence $s(N)\in A\inter C$ by \cref{inter}.
            Thus $s(N)\in A$ by \cref{inter}.
            Contradiction.
        \end{subproof}
        Hence $A$ is open in $X$ with respect to $T$ by \cref{open_from_local_neighborhoods}.
    \end{subproof}
    Therefore $X$ is compactly generated with respect to $T$ by \cref{compactly_generated_space}.
\end{proof}

% from §46 Item 7: continuity from continuity on compact subspaces
\begin{lemma}\label{compactly_generated_continuity_from_compact_subspaces}
    Assume $X$ is compactly generated with respect to $T$.
    Assume $f\in\funs{X}{Y}$.
    Suppose for all $C$ if
        $C$ is a compact subspace of $X$ with respect to $T$
    then $\restrl{f}{C}$ is continuous from $C$ to $Y$ with respect to $\subspaceTopology{T}{C}$ and $U$.
    Then $f$ is continuous from $X$ to $Y$ with respect to $T$ and $U$.
\end{lemma}
\begin{proof}
    We have $T$ is a topology on $X$ by \cref{compactly_generated_space}.
    Let $T_E = \subspaceTopology{T}{\emptyset}$.
    Hence $T_E$ is a topology on $\emptyset$ by \cref{emptyset_subseteq,subspace_topology_is_topology}.
    Therefore $\emptyset$ is a compact subspace of $X$ with respect to $T$
        by \cref{emptyset_subseteq,emptyset_compact,compact_subspace}.
    Thus $\restrl{f}{\emptyset}$ is continuous from $\emptyset$ to $Y$ with respect to $T_E$ and $U$ by assumption.
    Therefore $U$ is a topology on $Y$ by \cref{continuous}.
    Show that for all $V$ such that $V$ is open in $Y$ with respect to $U$ we have
        $\preimg{f}{V}$ is open in $X$ with respect to $T$.
    \begin{subproof}
        Fix $V$ such that $V$ is open in $Y$ with respect to $U$.
        Let $A = \preimg{f}{V}$.
        Show that for all $C$ if
            $C$ is a compact subspace of $X$ with respect to $T$
        then $A\inter C$ is open in $C$ with respect to $\subspaceTopology{T}{C}$.
        \begin{subproof}
            Fix $C$ such that $C$ is a compact subspace of $X$ with respect to $T$.
            Let $g = \restrl{f}{C}$.
            Hence $g$ is continuous from $C$ to $Y$ with respect to $\subspaceTopology{T}{C}$ and $U$ by assumption.
            Therefore $\preimg{g}{V}$ is open in $C$ with respect to $\subspaceTopology{T}{C}$ by \cref{continuous}.
            Thus $\preimg{g}{V} = A\inter C$
                by \cref{preimg_iff,restrl_iff,inter,inter_intro,subseteq,subseteq_antisymmetric}.
            Hence $A\inter C$ is open in $C$ with respect to $\subspaceTopology{T}{C}$ by \cref{pair_eq_iff}.
        \end{subproof}
        Therefore $\preimg{f}{V}$ is open in $X$ with respect to $T$
            by \cref{preimg_subseteq_dom,funs,subseteq,compactly_generated_space,pair_eq_iff}.
    \end{subproof}
    Hence $f$ is continuous from $X$ to $Y$ with respect to $T$ and $U$ by \cref{funs,continuous}.
\end{proof}

% from §46 Item 8: continuous functions are closed in the compact-convergence topology
\begin{theorem}\label{continuous_function_space_closed_compact_convergence}
    Assume $X$ is compactly generated with respect to $T$.
    Assume $d$ is a metric on $Y$.
    Assume $C$ is the set of continuous functions from $X$ to $Y$ with respect to $T$ and $\metricTopology{d}{Y}$.
    Assume $T_0$ is the topology of compact convergence on functions from $X$ to $Y$ with respect to $T$ and $d$.
    Then $C$ is closed in $\funs{X}{Y}$ with respect to $T_0$.
\end{theorem}
\begin{proof}
    Let $F = \funs{X}{Y}$.
    Let $U = \metricTopology{d}{Y}$.
    We have $T$ is a topology on $X$ by \cref{compactly_generated_space}.
    Therefore $T_0$ is a topology on $F$ by \cref{compact_convergence_topology}.
    $U$ is a topology on $Y$ by \cref{metric_basis_is_basis,basis_topology_is_topology,metric_topology,basis_for_topology}.
    $C\subseteq F$ by \cref{continuous_function_set,subseteq}.
    Show that for all $h$ if $h\in\closure{C}{F}{T_0}$ then $h\in C$.
    \begin{subproof}
        Fix $h$ such that $h\in\closure{C}{F}{T_0}$.
        Hence $h\in F$ by \cref{closure,subseteq}.
        Show that for all $K$ if
            $K$ is a compact subspace of $X$ with respect to $T$
        then $\restrl{h}{K}$ is continuous from $K$ to $Y$ with respect to $\subspaceTopology{T}{K}$ and $U$.
        \begin{subproof}
            Fix $K$ such that $K$ is a compact subspace of $X$ with respect to $T$.
            Let $T_K = \subspaceTopology{T}{K}$.
            Hence $K\subseteq X$ by \cref{compact_subspace}.
            Therefore $T_K$ is a topology on $K$ by \cref{subspace_topology_is_topology}.
            Let $g = \restrl{h}{K}$.
            Hence $g\in\funs{K}{Y}$ by \cref{restrl_typed_function}.
            \begin{byCase}
            \caseOf{$K = \emptyset$.}
                Show that for all $V$ such that $V$ is open in $Y$ with respect to $U$ we have
                    $\preimg{g}{V}$ is open in $K$ with respect to $T_K$.
                \begin{subproof}
                    Follows by \cref{preimg_subseteq_dom,funs,subseteq,pair_eq_iff,emptyset_subseteq,subseteq_antisymmetric,topology_on,open_in}.
                \end{subproof}
                Therefore $g$ is continuous from $K$ to $Y$ with respect to $T_K$ and $U$ by \cref{funs,continuous}.
            \caseOf{$K \neq \emptyset$.}
                Take $\rho$ such that
                    $\rho$ is the uniform metric on functions from $K$ to $Y$ with respect to $d$
                    by \cref{nonempty_inhabited,uniform_metric_on_function_space_exists_inhabited}.
                Hence $\rho$ is a metric on $\funs{K}{Y}$ by \cref{uniform_metric_on_function_space}.
                Let $D = \{u\in\funs{K}{Y} \mid \text{$u$ is continuous from $K$ to $Y$ with respect to $T_K$ and $U$}\}$.
                Therefore $D$ is closed in $\funs{K}{Y}$ with respect to $\metricTopology{\rho}{\funs{K}{Y}}$
                    by \cref{continuous_function_set,continuous_function_space_closed_uniform}.
                Show that for all $W$ such that
                    $W$ is open in $\funs{K}{Y}$ with respect to $\metricTopology{\rho}{\funs{K}{Y}}$ and
                    $g\in W$
                we have $W$ intersects $D$.
                \begin{subproof}
                    Fix $W$ such that
                        $W$ is open in $\funs{K}{Y}$ with respect to $\metricTopology{\rho}{\funs{K}{Y}}$ and
                        $g\in W$.
                    Hence $W\subseteq\funs{K}{Y}$ by \cref{open_in,topology_on,subseteq,pow_iff}.
                    Take $\epsilon\in\reals$ such that
                        $0 < \epsilon$ and
                        $\metricBall{\rho}{\funs{K}{Y}}{g}{\epsilon} \subseteq W$
                        by \cref{metric_open_iff}.
                    Let $A = \{u\in F \mid \rho((\restrl{h}{K},\restrl{u}{K})) < \epsilon\}$.
                    Hence $A$ is the compact-convergence ball about $h$ over $K$ with radius $\epsilon$
                        from $X$ to $Y$ with respect to $T$ and $d$ by \cref{compact_convergence_ball}.
                    Therefore $A$ intersects $C$
                        by \cref{compact_convergence_ball_open,compact_convergence_ball_center,closure_iff_intersects_open}.
                    Take $u$ such that $u\in A\inter C$ by \cref{intersects,nonempty_inhabited}.
                    Hence $\restrl{u}{K}\in D$
                        by \cref{inter,continuous_function_set,compact_subspace,subspace_of,continuous_restrict_domain,restrl_typed_function}.
                    We have $\rho((\restrl{h}{K},\restrl{u}{K})) < \epsilon$ by \cref{inter,subseteq}.
                    Hence $\restrl{u}{K}\in W\inter D$ by \cref{inter_intro,metric_ball,subseteq}.
                    Therefore $W$ intersects $D$ by \cref{intersects,emptyset}.
                \end{subproof}
                $\metricTopology{\rho}{\funs{K}{Y}}$ is a topology on $\funs{K}{Y}$
                    by \cref{metric_basis_is_basis,basis_topology_is_topology,metric_topology,basis_for_topology}.
                Therefore $g\in\closure{D}{\funs{K}{Y}}{\metricTopology{\rho}{\funs{K}{Y}}}$ by
                    \cref{continuous_function_set,subseteq,closure_iff_intersects_open}.
                Thus $D = \closure{D}{\funs{K}{Y}}{\metricTopology{\rho}{\funs{K}{Y}}}$
                    by \cref{continuous_function_set,subseteq,closed_eq_closure}.
                We have $g\in D$ by \cref{subseteq}.
                Therefore $g$ is continuous from $K$ to $Y$ with respect to $T_K$ and $U$
                    by \cref{continuous_function_set}.
            \end{byCase}
        \end{subproof}
        Hence $h$ is continuous from $X$ to $Y$ with respect to $T$ and $U$
            by \cref{compactly_generated_continuity_from_compact_subspaces}.
        Therefore $h\in C$ by \cref{continuous_function_set}.
    \end{subproof}
    Hence $\closure{C}{F}{T_0}\subseteq C$ by \cref{subseteq}.
    Therefore $C$ is closed in $F$ with respect to $T_0$ by
        \cref{subseteq,subseteq_closure,subseteq_antisymmetric,closure_closed_in}.
\end{proof}

% from §46 Item 9: compact-convergence limits of continuous functions are continuous
\begin{corollary}\label{compact_convergence_limit_continuous}
    Assume $X$ is compactly generated with respect to $T$.
    Assume $d$ is a metric on $Y$.
    Assume $T_0$ is the topology of compact convergence on functions from $X$ to $Y$ with respect to $T$ and $d$.
    Assume $s$ is a sequence in $\funs{X}{Y}$.
    Assume $f\in\funs{X}{Y}$.
    Suppose for all $n\in\naturalsPlus$ we have
        $s(n)$ is continuous from $X$ to $Y$ with respect to $T$ and $\metricTopology{d}{Y}$.
    Suppose $s$ converges to $f$ in $\funs{X}{Y}$ with respect to $T_0$.
    Then $f$ is continuous from $X$ to $Y$ with respect to $T$ and $\metricTopology{d}{Y}$.
\end{corollary}
\begin{proof}
    Let $F = \funs{X}{Y}$.
    Let $C = \{g\in F \mid \text{$g$ is continuous from $X$ to $Y$ with respect to $T$ and $\metricTopology{d}{Y}$}\}$.
    Hence $C$ is closed in $F$ with respect to $T_0$
        by \cref{continuous_function_set,continuous_function_space_closed_compact_convergence}.
    Therefore $T_0$ is a topology on $F$ by \cref{compact_convergence_topology}.
    We have $C\subseteq F$ by \cref{continuous_function_set,subseteq}.
    $s$ is a function and $\dom{s} = \naturalsPlus$ by \cref{sequence_in,funs_is_function,funs}.
    Show that for all $n$ such that $n\in\dom{s}$ we have $s(n)\in C$.
    \begin{subproof}
        Follows by \cref{sequence_in,funs,subseteq,continuous_function_set,funs_type_apply}.
    \end{subproof}
    Hence $s$ is a sequence in $C$ by \cref{function_rightunique,funs_intro,sequence_in}.
    Therefore $f\in\closure{C}{F}{T_0}$ by \cref{sequence_implies_closure}.
    Hence $f$ is continuous from $X$ to $Y$ with respect to $T$ and $\metricTopology{d}{Y}$
        by \cref{closed_eq_closure,subseteq,continuous_function_set}.
\end{proof}

% from §46 Item 10: uniform topology is finer than compact convergence
\begin{theorem}\label{uniform_topology_finer_than_compact_convergence}
    Assume $T$ is a topology on $X$.
    Assume $d$ is a metric on $Y$.
    Assume $\rho$ is the uniform metric on functions from $X$ to $Y$ with respect to $d$.
    Assume $T_0$ is the topology of compact convergence on functions from $X$ to $Y$ with respect to $T$ and $d$.
    Then $\metricTopology{\rho}{\funs{X}{Y}}$ is finer than $T_0$ on $\funs{X}{Y}$.
\end{theorem}
\begin{proof}
    Let $F = \funs{X}{Y}$.
    Take $B$ such that
        $B$ is the compact-convergence basis on functions from $X$ to $Y$ with respect to $T$ and $d$ and
        $T_0 = \basisTopology{B}{F}$ by \cref{compact_convergence_topology}.
    Show that for all $h,A$ such that $h\in F$ and $A\in B$ and $h\in A$
        there exists $M\in\metricBasis{\rho}{F}$ such that $h\in M$ and $M\subseteq A$.
    \begin{subproof}
        Fix $h$.
        Fix $A$ such that $h\in F$ and $A\in B$ and $h\in A$.
        Take $f, C, \epsilon$ such that
            $A$ is the compact-convergence ball about $f$ over $C$ with radius $\epsilon$
                from $X$ to $Y$ with respect to $T$ and $d$
                by \cref{compact_convergence_function_space_basis}.
        Take $D,\delta$ such that
            $\delta\in\reals$ and
            $0 < \delta$ and
            $D$ is the compact-convergence ball about $h$ over $C$ with radius $\delta$
                from $X$ to $Y$ with respect to $T$ and $d$ and
            $D\subseteq A$ by \cref{compact_convergence_ball_refinement}.
        Take $\eta$ such that
            $\eta\in\reals$ and
            $0 < \eta$ and
            $\eta < 1$ and
            $\eta < \delta$ by \cref{positive_real_has_smaller_positive_below_one}.
        Let $M = \metricBall{\rho}{F}{h}{\eta}$.
        Take $\sigma$ such that
            $\sigma$ is the uniform metric on functions from $C$ to $Y$ with respect to $d$ and
            $D = \{g\in F \mid \sigma((\restrl{h}{C},\restrl{g}{C})) < \delta\}$ by \cref{compact_convergence_ball}.
        Show that for all $g$ if $g\in M$ then $g\in D$.
        \begin{subproof}
            Fix $g$ such that $g\in M$.
            Hence $g\in F$ and $\rho((h,g)) < \eta$ by \cref{metric_ball}.
            Therefore $\restrl{h}{C}\in\funs{C}{Y}$ and $\restrl{g}{C}\in\funs{C}{Y}$
                by \cref{compact_convergence_ball,compact_subspace,subseteq,restrl_typed_function}.
            For all $x\in C$ we have $\apply{\boundedMetric{d}}{(\apply{\restrl{h}{C}}{x}, \apply{\restrl{g}{C}}{x})} \leq \eta$ by \cref{compact_convergence_ball,compact_subspace,subseteq,uniform_metric_small_implies_pointwise_small,times_tuple_intro,funs_type_apply,metric_on,real_one_in_naturals,real_naturals_subset_reals,real_lt_trans,bounded_metric_apply_below_one,real_lt_subst_left,real_lt_imp_leq,restrl_of_function_apply,funs_is_function,funs_is_relation,funs_elim,pair_eq_iff}.
            Hence $\sigma((\restrl{h}{C},\restrl{g}{C})) \leq \eta$ by \cref{uniform_metric_leq_from_coordinate_bounds}.
            We have $\sigma((\restrl{h}{C},\restrl{g}{C}))\in\reals$
                by \cref{uniform_metric_on_function_space,metric_on,funs_type_apply,times_tuple_intro}.
            $\sigma((\restrl{h}{C},\restrl{g}{C})) < \eta$ or $\sigma((\restrl{h}{C},\restrl{g}{C})) = \eta$
                by \cref{real_leq_split}.
            Hence $\sigma((\restrl{h}{C},\restrl{g}{C})) < \delta$ by \cref{real_lt_trans,real_lt_subst_left}.
            Therefore $g\in D$ by \cref{subseteq}.
        \end{subproof}
        Hence $M\subseteq D$ by \cref{subseteq}.
        Therefore $M\subseteq A$ by \cref{subseteq_transitive}.
        Thus $M\in\metricBasis{\rho}{F}$ by \cref{metric_ball,subseteq,metric_basis,powerset_intro}.
        Therefore $h\in M$ by \cref{metric_ball,metric_on,uniform_metric_on_function_space,metric_zero_characterization,times_tuple_intro,real_lt_subst_left}.
        Follows by definition.
    \end{subproof}
    Show that for all $U$ if $U$ is open in $F$ with respect to $T_0$ then
        $U$ is open in $F$ with respect to $\metricTopology{\rho}{F}$.
    \begin{subproof}
        Fix $U$ such that $U$ is open in $F$ with respect to $T_0$.
        We have $U\subseteq F$ by \cref{open_in,topology_generated_by_basis,pow_iff,subseteq}.
        Hence $\rho$ is a metric on $F$ by \cref{uniform_metric_on_function_space}.
        Therefore $\metricTopology{\rho}{F}$ is a topology on $F$
            by \cref{metric_basis_is_basis,basis_topology_is_topology,metric_topology,basis_for_topology}.
        Show that for all $h\in U$ there exists $M$ such that
            $M$ is a neighborhood of $h$ in $F$ with respect to $\metricTopology{\rho}{F}$ and
            $M\subseteq U$.
        \begin{subproof}
            Fix $h\in U$.
            Hence $h\in F$ by \cref{subseteq}.
            Take $A\in B$ such that $h\in A$ and $A\subseteq U$
                by \cref{open_in,topology_generated_by_basis}.
            Take $M\in\metricBasis{\rho}{F}$ such that $h\in M$ and $M\subseteq A$ by assumption.
            Therefore $M$ is a neighborhood of $h$ in $F$ with respect to $\metricTopology{\rho}{F}$ and $M\subseteq U$
                by \cref{metric_basis_is_basis,basis_elem_open_in_generated,metric_topology,open_in,neighborhood,subseteq_transitive}.
            Follows by definition.
        \end{subproof}
        Therefore $U$ is open in $F$ with respect to $\metricTopology{\rho}{F}$ by \cref{open_from_local_neighborhoods}.
    \end{subproof}
    We have $T_0$ is a topology on $F$ by \cref{compact_convergence_topology}.
    Therefore $\metricTopology{\rho}{F}$ is a topology on $F$
        by \cref{uniform_metric_on_function_space,metric_basis_is_basis,basis_topology_is_topology,metric_topology,basis_for_topology}.
    $T_0\subseteq\metricTopology{\rho}{F}$ by \cref{open_in,subseteq}.
    Hence $\metricTopology{\rho}{F}$ is finer than $T_0$ on $F$ by \cref{finer_topology}.
\end{proof}

% from §46 helper lemma 13: singleton subsets are compact subspaces
\begin{lemma}\label{singleton_compact_subspace}
    Assume $T$ is a topology on $X$.
    Assume $x\in X$.
    Then $\{x\}$ is a compact subspace of $X$ with respect to $T$.
\end{lemma}
\begin{proof}
    We have $\{x\}$ is a subspace of $X$ with respect to $T$ by \cref{singleton_subset_intro,subspace_of}.
    Show that for all $A$ such that
        $A$ is a covering of $\{x\}$ and
        for all $U$ such that $U\in A$ we have $U$ is open in $X$ with respect to $T$
    there exists $B$ such that
        $B\subseteq A$ and
        $B$ is finite and
        $B$ is a covering of $\{x\}$.
    \begin{subproof}
        Fix $A$ such that
            $A$ is a covering of $\{x\}$ and
            for all $U$ such that $U\in A$ we have $U$ is open in $X$ with respect to $T$.
        Hence $x\in\unions{A}$ by \cref{singleton_intro,covering,subseteq}.
        Take $U\in A$ such that $x\in U$ by \cref{unions_iff}.
        Let $B = \{U\}$.
        We have $B\subseteq A$ by \cref{singleton_subset_intro}.
        Therefore $B$ is finite by \cref{upair_intro_left,singleton_subset_intro,pair_is_finite,subseteq_finite_is_finite}.
        Hence $B$ is a covering of $\{x\}$ by \cref{singleton_intro,unions_iff,singleton_subset_intro,covering}.
    \end{subproof}
    Hence $\{x\}$ is compact with respect to $\subspaceTopology{T}{\{x\}}$
        by \cref{compact_subspace_open_covering_backward}.
    Therefore $\{x\}$ is a compact subspace of $X$ with respect to $T$ by \cref{compact_subspace,singleton_subset_intro}.
\end{proof}

% from §46 helper lemma 14: restricting a typed function to its whole domain does nothing
\begin{lemma}\label{restrl_self_typed_function}
    Assume $f\in\funs{X}{Y}$.
    Then $\restrl{f}{X} = f$.
\end{lemma}
\begin{proof}
    $\restrl{f}{X}\subseteq f$ by \cref{restrl_subseteq}.
    Show that for all $w$ if $w\in f$ then $w\in\restrl{f}{X}$.
    \begin{subproof}
        Follows by \cref{funs_elim,relation,dom_intro,pair_eq_iff,restrl}.
    \end{subproof}
    Therefore $\restrl{f}{X} = f$ by \cref{subseteq,subseteq_antisymmetric}.
\end{proof}

% from §46 Item 10: compact convergence is finer than pointwise convergence
\begin{theorem}\label{compact_convergence_finer_than_pointwise}
    Assume $T$ is a topology on $X$.
    Assume $d$ is a metric on $Y$.
    Let $U = \metricTopology{d}{Y}$.
    Assume $T_0$ is the topology of compact convergence on functions from $X$ to $Y$ with respect to $T$ and $d$.
    Assume $T_1$ is the topology of pointwise convergence on functions from $X$ to $Y$ with respect to $U$.
    Then $T_0$ is finer than $T_1$ on $\funs{X}{Y}$.
\end{theorem}
\begin{proof}
    Let $F = \funs{X}{Y}$.
    Take $B$ such that
        $B$ is the compact-convergence basis on functions from $X$ to $Y$ with respect to $T$ and $d$ and
        $T_0 = \basisTopology{B}{F}$ by \cref{compact_convergence_topology}.
    Take $S$ such that
        $S$ is the pointwise subbasis on functions from $X$ to $Y$ with respect to $U$ and
        $T_1 = \subbasisTopology{S}{F}$ by \cref{pointwise_convergence_topology}.
    We have $T_0$ is a topology on $F$ by \cref{compact_convergence_topology}.
    $T_1$ is a topology on $F$ by \cref{pointwise_convergence_topology}.
    We have $U$ is a topology on $Y$ by \cref{metric_topology,metric_basis_is_basis,basis_topology_is_topology}.
    Show that for all $A$ if $A\in S$ then $A\in T_0$.
    \begin{subproof}
        Fix $A$ such that $A\in S$.
        Take $x,V$ such that
            $x\in X$ and
            $V\in U$ and
            $A$ is the point-open set at $x$ with value set $V$ from $X$ to $Y$
            by \cref{pointwise_function_space_subbasis}.
        Hence $A = \{g\in F \mid \apply{g}{x}\in V\}$ by \cref{point_open_set}.
        Therefore $A\subseteq F$ by \cref{point_open_set,subseteq}.
        Show that for all $h\in A$ there exists $M$ such that
            $M$ is a neighborhood of $h$ in $F$ with respect to $T_0$ and
            $M\subseteq A$.
        \begin{subproof}
            Fix $h\in A$.
            Hence $h(x)\in V$ by \cref{inters_iff_forall}.
            Take $\epsilon$ such that
                $\epsilon\in\reals$ and
                $0 < \epsilon$ and
                $\metricBall{d}{Y}{h(x)}{\epsilon}\subseteq V$ by \cref{open_in,topology_on,subseteq,pow_iff,metric_open_iff}.
            Take $\eta$ such that
                $\eta\in\reals$ and
                $0 < \eta$ and
                $\eta < 1$ and
                $\eta < \epsilon$ by \cref{positive_real_has_smaller_positive_below_one}.
            Hence $\{x\}$ is a compact subspace of $X$ with respect to $T$ by \cref{singleton_compact_subspace}.
            Take $\sigma$ such that
                $\sigma$ is the uniform metric on functions from $\{x\}$ to $Y$ with respect to $d$
                by \cref{singleton_intro,emptyset,nonempty_inhabited,uniform_metric_on_function_space_exists_inhabited}.
            Let $M = \{g\in F \mid \sigma((\restrl{h}{\{x\}},\restrl{g}{\{x\}})) < \eta\}$.
            Hence $M$ is the compact-convergence ball about $h$ over $\{x\}$ with radius $\eta$
                from $X$ to $Y$ with respect to $T$ and $d$ by \cref{compact_convergence_ball}.
            Therefore $M$ is a neighborhood of $h$ in $F$ with respect to $T_0$
                by \cref{compact_convergence_ball_open,compact_convergence_ball_center,neighborhood}.
            Show that for all $g$ if $g\in M$ then $g\in A$.
            \begin{subproof}
                Fix $g$ such that $g\in M$.
                Hence $g\in F$ and $\sigma((\restrl{h}{\{x\}},\restrl{g}{\{x\}})) < \eta$ by \cref{subseteq}.
                Therefore $d((h(x),g(x))) < \eta$
                    by \cref{subseteq,restrl_typed_function,singleton_intro,uniform_metric_small_implies_pointwise_small,restrl_of_function_apply,funs_is_function,funs_is_relation,funs_elim,singleton_subset_intro,pair_eq_iff}.
                Thus $d((h(x),g(x))) < \epsilon$ by \cref{metric_on,funs_type_apply,times_tuple_intro,real_lt_trans}.
                Therefore $g\in A$ by \cref{metric_ball,times_tuple_intro,funs_type_apply,subseteq,inters_iff_forall}.
            \end{subproof}
            Hence $M\subseteq A$ by \cref{subseteq}.
            Follows by definition.
        \end{subproof}
        Therefore $A\in T_0$ by \cref{open_from_local_neighborhoods,open_in}.
    \end{subproof}
    Hence $S\subseteq T_0$ by \cref{subseteq}.
    Show that for all $W$ if $W\in T_1$ then $W\in T_0$.
    \begin{subproof}
        Fix $W$ such that $W\in T_1$.
        Hence $W\in\basisTopology{\finiteIntersections{S}}{F}$ by \cref{topology_generated_by_subbasis}.
        Therefore $W\subseteq F$ by \cref{topology_on,subseteq,pow_iff}.
        Show that for all $h\in W$ there exists $M$ such that
            $M$ is a neighborhood of $h$ in $F$ with respect to $T_0$ and
            $M\subseteq W$.
        \begin{subproof}
            Fix $h\in W$.
            Take $C\in\finiteIntersections{S}$ such that $h\in C$ and $C\subseteq W$
                by \cref{topology_generated_by_basis}.
            Take $G$ such that
                $G\subseteq S$ and
                $G$ is finite and
                $G$ is inhabited and
                $C = \inters{G}$ by \cref{finite_intersections}.
            Show that $C$ is a neighborhood of $h$ in $F$ with respect to $T_0$.
            \begin{subproof}
                Follows by \cref{subseteq_transitive,inters_iff_forall,pair_eq_iff,finite_neighborhood_intersection}.
            \end{subproof}
            Follows by definition.
        \end{subproof}
        Therefore $W\in T_0$ by \cref{open_from_local_neighborhoods,open_in}.
    \end{subproof}
    Hence $T_1\subseteq T_0$ by \cref{subseteq}.
    Therefore $T_0$ is finer than $T_1$ on $F$ by \cref{finer_topology}.
\end{proof}

% from §46 Item 10: compactness of X identifies uniform and compact-convergence topologies
\begin{theorem}\label{compact_space_uniform_equals_compact_convergence}
    Assume $T$ is a topology on $X$.
    Assume $X$ is compact with respect to $T$.
    Assume $d$ is a metric on $Y$.
    Assume $\rho$ is the uniform metric on functions from $X$ to $Y$ with respect to $d$.
    Assume $T_0$ is the topology of compact convergence on functions from $X$ to $Y$ with respect to $T$ and $d$.
    Then $\metricTopology{\rho}{\funs{X}{Y}} = T_0$.
\end{theorem}
\begin{proof}
    Let $F = \funs{X}{Y}$.
    Let $T_1 = \metricTopology{\rho}{F}$.
    We have $T_1$ is finer than $T_0$ on $F$ by \cref{uniform_topology_finer_than_compact_convergence}.
    Hence $T_0$ is a topology on $F$ by \cref{compact_convergence_topology}.
    Therefore $\rho$ is a metric on $F$ by \cref{uniform_metric_on_function_space}.
    Hence $T_1$ is a topology on $F$ by \cref{metric_topology,metric_basis_is_basis,basis_topology_is_topology}.
    Let $T_X = \subspaceTopology{T}{X}$.
    Hence $T_X = T$ by \cref{topology_on,powerset_elim,subseteq,inter_absorb_supseteq_left,subspace_topology,subseteq_antisymmetric}.
    Therefore $X$ is a compact subspace of $X$ with respect to $T$ by \cref{subseteq_refl,compact_subspace}.
    Show that for all $U$ if $U$ is open in $F$ with respect to $T_1$ then
        $U$ is open in $F$ with respect to $T_0$.
    \begin{subproof}
        Fix $U$ such that $U$ is open in $F$ with respect to $T_1$.
        Hence $U\subseteq F$ by \cref{open_in,topology_on,subseteq,pow_iff}.
        Show that for all $h\in U$ there exists $A$ such that
            $A$ is a neighborhood of $h$ in $F$ with respect to $T_0$ and
            $A\subseteq U$.
        \begin{subproof}
            Fix $h\in U$.
            Hence $h\in F$ by \cref{subseteq}.
            Take $\epsilon$ such that
                $\epsilon\in\reals$ and
                $0 < \epsilon$ and
                $\metricBall{\rho}{F}{h}{\epsilon}\subseteq U$ by \cref{metric_open_iff}.
            Let $A = \{g\in F \mid \rho((\restrl{h}{X},\restrl{g}{X})) < \epsilon\}$.
            Hence $A$ is the compact-convergence ball about $h$ over $X$ with radius $\epsilon$
                from $X$ to $Y$ with respect to $T$ and $d$ by \cref{compact_convergence_ball}.
            Therefore $A$ is a neighborhood of $h$ in $F$ with respect to $T_0$
                by \cref{compact_convergence_ball_open,compact_convergence_ball_center,neighborhood}.
            Show that $A\subseteq U$.
            \begin{subproof}
                Follows by \cref{subseteq,restrl_self_typed_function,real_lt_subst_left,real_lt_subst_right,metric_ball}.
            \end{subproof}
            Follows by definition.
        \end{subproof}
        Therefore $U$ is open in $F$ with respect to $T_0$ by \cref{open_from_local_neighborhoods}.
    \end{subproof}
    Hence $T_1\subseteq T_0$ by \cref{open_in,subseteq}.
    Thus $T_0\subseteq T_1$ by \cref{finer_topology}.
    Therefore $T_1 = T_0$ by \cref{subseteq_antisymmetric}.
\end{proof}

% from §46 helper lemma 15: compact subspaces of a discrete space are finite
\begin{lemma}\label{compact_subspace_of_discrete_is_finite}
    Assume $T = \discrete{X}$.
    Assume $C$ is a compact subspace of $X$ with respect to $T$.
    Then $C$ is finite.
\end{lemma}
\begin{proof}
    Let $d = \discreteMetric{X}$.
    We have $d$ is a metric on $X$ by \cref{discrete_metric_is_metric}.
    Therefore $T = \metricTopology{d}{X}$ by \cref{discrete_metric_topology,pair_eq_iff}.
    Hence $T$ is a topology on $X$ by \cref{metric_topology,metric_basis_is_basis,basis_topology_is_topology,pair_eq_iff}.
    We have $C\subseteq X$ by \cref{compact_subspace}.
    Let $S = \{\{x\} \mid x\in C\}$.
    Hence $S$ is a covering of $C$ by \cref{relation,singleton_intro,unions_iff,subseteq,covering}.
    For all $U$ such that $U\in S$ we have $U$ is open in $X$ with respect to $T$
        by \cref{relation,singleton_subset_intro,subseteq,powerset_intro,discrete_topology,pair_eq_iff,open_in}.
    Let $T_C = \subspaceTopology{T}{C}$.
    Hence $C$ is a subspace of $X$ with respect to $T$ by \cref{subspace_of,compact_subspace}.
    Therefore $C$ is compact with respect to $T_C$ by \cref{compact_subspace}.
    Take $B$ such that $B\subseteq S$ and $B$ is finite and $B$ is a covering of $C$
        by \cref{compact_subspace_open_covering}.
    Show that for all $U$ such that $U\in B$ we have $U$ is finite.
    \begin{subproof}
        Follows by \cref{subseteq,relation,upair_intro_left,singleton_subset_intro,pair_eq_iff,pair_is_finite,subseteq_finite_is_finite}.
    \end{subproof}
    Hence $\unions{B}$ is finite by \cref{finite_union_of_finite_family}.
    We have $C\subseteq\unions{B}$ by \cref{covering}.
    Therefore $C$ is finite by \cref{subseteq_finite_is_finite}.
\end{proof}

% from §46 helper definition 6: point-open metric neighborhood set
\begin{definition}\label{point_open_metric_neighborhood_set}
    $E$ is the point-open metric neighborhood set at $x$ around $h$ with radius $\eta$ from $X$ to $Y$ with respect to $d$ iff
        $h\in\funs{X}{Y}$ and
        $x\in X$ and
        $d$ is a metric on $Y$ and
        $\eta\in\reals$ and
        $0 < \eta$ and
        $E$ is the point-open set at $x$ with value set $\metricBall{d}{Y}{h(x)}{\eta}$ from $X$ to $Y$.
\end{definition}

% from §46 helper lemma 16: pointwise subbasis is a subbasis on inhabited domains
\begin{lemma}\label{pointwise_subbasis_is_subbasis_inhabited_domain}
    Assume $x_0\in X$.
    Assume $d$ is a metric on $Y$.
    Let $U = \metricTopology{d}{Y}$.
    Assume $S$ is the pointwise subbasis on functions from $X$ to $Y$ with respect to $U$.
    Then $S$ is a subbasis on $\funs{X}{Y}$.
\end{lemma}
\begin{proof}
    Let $F = \funs{X}{Y}$.
    $S\subseteq\pow{F}$ by \cref{pointwise_function_space_subbasis,subseteq}.
    Let $E = \{u\in F \mid \apply{u}{x_0}\in Y\}$.
    Hence $E\in S$ by \cref{metric_topology,metric_basis_is_basis,basis_topology_is_topology,open_in,topology_on,pointwise_subbasis_contains_point_open_set}.
    $F\subseteq\unions{S}$ by \cref{funs_type_apply,subseteq,unions_iff}.
    Therefore $\unions{S} = F$ by \cref{subseteq,pow_iff,unions_subseteq_of_powerset_is_subseteq,subseteq_antisymmetric}.
    Therefore $S$ is a subbasis on $F$ by \cref{subbasis_on}.
\end{proof}

% from §46 helper lemma 17: finite point-open metric neighborhoods form an open intersection
\begin{lemma}\label{finite_point_open_metric_neighborhoods_open}
    Assume $C\subseteq X$.
    Assume $C$ is finite.
    Assume $C$ is inhabited.
    Assume $h\in\funs{X}{Y}$.
    Assume $d$ is a metric on $Y$.
    Let $U = \metricTopology{d}{Y}$.
    Assume $S$ is the pointwise subbasis on functions from $X$ to $Y$ with respect to $U$.
    Assume $S$ is a subbasis on $\funs{X}{Y}$.
    Assume $\eta\in\reals$.
    Assume $0 < \eta$.
    Then there exists $M$ such that
        $M$ is open in $\funs{X}{Y}$ with respect to $\subbasisTopology{S}{\funs{X}{Y}}$ and
        $h\in M$ and
        for all $g$ if $g\in M$ then
            $g\in\funs{X}{Y}$ and
            for all $x\in C$ we have $g(x)\in\metricBall{d}{Y}{h(x)}{\eta}$.
\end{lemma}
\begin{proof}
    Let $F = \funs{X}{Y}$.
    Let $R = \{w\in C\times\pow{F} \mid \text{$\snd{w}$ is the point-open metric neighborhood set
        at $\fst{w}$ around $h$ with radius $\eta$ from $X$ to $Y$ with respect to $d$}\}$.
    Hence $R$ is a relation by \cref{times_elem_is_tuple,relation}.
    Show that for all $x\in C$ there exists $A$ such that $x\mathrel{R} A$.
    \begin{subproof}
        Fix $x\in C$.
        We have $x\in X$ by \cref{subseteq}.
        Let $A = \{g\in F \mid \apply{g}{x}\in\metricBall{d}{Y}{h(x)}{\eta}\}$.
        Hence $A$ is the point-open set at $x$ with value set $\metricBall{d}{Y}{h(x)}{\eta}$ from $X$ to $Y$
            by \cref{metric_ball,subseteq,point_open_set}.
        Therefore $A$ is the point-open metric neighborhood set at $x$ around $h$ with radius $\eta$
            from $X$ to $Y$ with respect to $d$ by \cref{point_open_metric_neighborhood_set}.
        Thus $(x,A)\in C\times\pow{F}$
            by \cref{subseteq,point_open_set,point_open_metric_neighborhood_set,pow_iff,times_tuple_intro}.
        Therefore $(x,A)\in R$ by \cref{fst_eq,snd_eq}.
    \end{subproof}
    Take $e$ such that $e$ is a function on $C$ and for all $x\in C$ we have $x\mathrel{R}\apply{e}{x}$
        by \cref{choice_function}.
    Let $G = \img{e}{C}$.
    Hence $G$ is finite by \cref{finite_image_function}.
    Take $x_0$ such that $x_0\in C$ by assumption.
    Therefore $G$ is inhabited by \cref{function_apply_elim,img_iff,emptyset,nonempty_inhabited}.
    We have $U$ is a topology on $Y$ by \cref{metric_topology,metric_basis_is_basis,basis_topology_is_topology}.
    Show that for all $A$ if $A\in G$ then $A\in S$.
    \begin{subproof}
        Fix $A$.
        Assume $A\in G$.
        Take $x\in C$ such that $e(x) = A$ by \cref{img_iff,function_apply_intro}.
        Hence $x\mathrel{R} A$ by \cref{choice_function,pair_eq_iff}.
        Therefore $A$ is the point-open metric neighborhood set at $x$ around $h$ with radius $\eta$
            from $X$ to $Y$ with respect to $d$ by \cref{fst_eq,snd_eq}.
        Thus $\metricBall{d}{Y}{h(x)}{\eta}\in U$ by
            \cref{subseteq,funs_type_apply,metric_ball,subseteq,powerset_intro,metric_basis,metric_basis_is_basis,basis_elem_open_in_generated,metric_topology}.
        Therefore $A\in S$ by \cref{point_open_metric_neighborhood_set,pointwise_subbasis_contains_point_open_set,point_open_set}.
    \end{subproof}
    Hence $G\subseteq S$ by \cref{subseteq}.
    Let $M = \inters{G}$.
    Show that for all $g$ if $g\in M$ then $g\in\unions{S}$.
    \begin{subproof}
        Follows by \cref{inters_iff_forall,subseteq,function_apply_elim,img_iff,unions_iff}.
    \end{subproof}
    Therefore $M\subseteq\unions{S}$ by \cref{subseteq}.
    Show that $M$ is open in $F$ with respect to $\subbasisTopology{S}{F}$.
    \begin{subproof}
        Follows by
            \cref{powerset_intro,finite_intersections,subbasis_finite_intersections_basis,basis_elem_open_in_generated,topology_generated_by_subbasis,basis_for_topology,basis_topology_is_topology,open_in}.
    \end{subproof}
    Show that for all $A\in G$ we have $h\in A$.
    \begin{subproof}
        Fix $A\in G$.
        Take $x\in C$ such that $e(x) = A$ by \cref{img_iff,function_apply_intro}.
        Hence $x\mathrel{R} A$ by \cref{choice_function,pair_eq_iff}.
        Therefore $A$ is the point-open metric neighborhood set at $x$ around $h$ with radius $\eta$
            from $X$ to $Y$ with respect to $d$ by \cref{fst_eq,snd_eq}.
        We have $h(x)\in Y$ by \cref{subseteq,funs_type_apply}.
        Therefore $d((h(x),h(x))) = 0$ by \cref{metric_on,metric_zero_characterization,times_tuple_intro,funs_type_apply}.
        Hence $d((h(x),h(x))) < \eta$ by \cref{real_add_ident,real_lt_subst_left}.
        Therefore $h(x)\in\metricBall{d}{Y}{h(x)}{\eta}$ by \cref{metric_ball}.
        Therefore $h\in A$ by \cref{point_open_metric_neighborhood_set,point_open_set}.
    \end{subproof}
    Hence $h\in M$ by \cref{inters_iff_forall}.
    Show that for all $g$ if $g\in M$ then
        $g\in\funs{X}{Y}$ and
        for all $x\in C$ we have $g(x)\in\metricBall{d}{Y}{h(x)}{\eta}$.
    \begin{subproof}
        Follows by \cref{function_apply_elim,img_iff,inters_iff_forall,choice_function,fst_eq,snd_eq,point_open_metric_neighborhood_set,point_open_set}.
    \end{subproof}
    Follows by definition.
\end{proof}

% from §46 helper lemma 18: compact-convergence balls over the empty subspace are the whole function space
\begin{lemma}\label{compact_convergence_ball_over_empty_is_whole_space}
    Assume $A$ is the compact-convergence ball about $f$ over $\emptyset$ with radius $\epsilon$
        from $X$ to $Y$ with respect to $T$ and $d$.
    Then $A = \funs{X}{Y}$.
\end{lemma}
\begin{proof}
    Let $F = \funs{X}{Y}$.
    Take $\rho$ such that
        $\rho$ is the uniform metric on functions from $\emptyset$ to $Y$ with respect to $d$ and
        $A = \{g\in F \mid \rho((\restrl{f}{\emptyset},\restrl{g}{\emptyset})) < \epsilon\}$ by \cref{compact_convergence_ball}.
    $A\subseteq F$ by \cref{subseteq}.
    Show that for all $g$ if $g\in F$ then $g\in A$.
    \begin{subproof}
        Follows by \cref{compact_convergence_ball,restrl_typed_function,emptyset_subseteq,functions_eq_iff,emptyset,uniform_metric_on_function_space,metric_on,metric_zero_characterization,times_tuple_intro,real_add_ident,real_lt_subst_left,subseteq}.
    \end{subproof}
    Hence $F\subseteq A$ by \cref{subseteq}.
    Therefore $A = F$ by \cref{subseteq_antisymmetric}.
\end{proof}

% from §46 Item 10: discreteness of X identifies compact-convergence and pointwise topologies
\begin{theorem}\label{discrete_space_compact_convergence_equals_pointwise}
    Assume $T = \discrete{X}$.
    Assume $d$ is a metric on $Y$.
    Let $U = \metricTopology{d}{Y}$.
    Assume $T_0$ is the topology of compact convergence on functions from $X$ to $Y$ with respect to $T$ and $d$.
    Assume $T_1$ is the topology of pointwise convergence on functions from $X$ to $Y$ with respect to $U$.
    Then $T_0 = T_1$.
\end{theorem}
\begin{proof}
    Let $F = \funs{X}{Y}$.
    Take $B$ such that
        $B$ is the compact-convergence basis on functions from $X$ to $Y$ with respect to $T$ and $d$ and
        $T_0 = \basisTopology{B}{F}$ by \cref{compact_convergence_topology}.
    Take $S$ such that
        $S$ is the pointwise subbasis on functions from $X$ to $Y$ with respect to $U$ and
        $T_1 = \subbasisTopology{S}{F}$ by \cref{pointwise_convergence_topology}.
    We have $T = \pow{X}$ by \cref{discrete_topology}.
    Therefore $T\subseteq\pow{X}$ by \cref{subseteq_refl}.
    $\emptyset\in T$ and $X\in T$ by \cref{discrete_topology,powerset_bottom,powerset_top}.
    Hence $T$ is closed under arbitrary unions by \cref{subseteq_transitive,powerset_closed_unions,discrete_topology}.
    Therefore $T$ is closed under binary intersections by \cref{discrete_topology,subseteq,inter_powerset}.
    Hence $T$ is a topology on $X$ by \cref{topology_on}.
    Therefore $T_0$ is finer than $T_1$ on $F$ by \cref{compact_convergence_finer_than_pointwise}.
    We have $T_0$ is a topology on $F$ by \cref{compact_convergence_topology}.
    $T_1$ is a topology on $F$ by \cref{pointwise_convergence_topology}.
    Show that for all $A\in B$ we have $A$ is open in $F$ with respect to $T_0$.
    \begin{subproof}
        Follows by \cref{compact_convergence_function_space_basis,compact_convergence_ball_open}.
    \end{subproof}
    Show that for all $W,h$ such that $W\in T_0$ and $h\in W$ there exists $A\in B$ such that $h\in A$ and $A\subseteq W$.
    \begin{subproof}
        Follows by \cref{topology_generated_by_basis}.
    \end{subproof}
    Hence $B$ is a basis for $T_0$ on $F$ by \cref{open_collection_basis}.
    Show that for all $A$ if $A\in B$ then $A\in T_1$.
    \begin{subproof}
        Fix $A$.
        Assume $A\in B$.
        Take $f, C, \epsilon$ such that
            $A$ is the compact-convergence ball about $f$ over $C$ with radius $\epsilon$
                from $X$ to $Y$ with respect to $T$ and $d$ by \cref{compact_convergence_function_space_basis}.
        Hence $A\subseteq F$ by \cref{compact_convergence_ball,subseteq}.
        \begin{byCase}
        \caseOf{$C = \emptyset$.}
            Hence $A = F$ by \cref{compact_convergence_ball_over_empty_is_whole_space}.
            Therefore $A\in T_1$ by \cref{open_in,topology_on,pair_eq_iff}.
        \caseOf{$C\neq\emptyset$.}
                We have $C$ is a compact subspace of $X$ with respect to $T$ by \cref{compact_convergence_ball}.
                Therefore $C\subseteq X$ by \cref{compact_subspace,subseteq}.
                Hence $C$ is finite by \cref{compact_subspace_of_discrete_is_finite}.
                Hence $C$ is inhabited by \cref{nonempty_inhabited}.
                Take $x_0$ such that $x_0\in C$ by assumption.
                Thus $S$ is a subbasis on $F$ by \cref{subseteq,pointwise_subbasis_is_subbasis_inhabited_domain}.
                Show that for all $h\in A$ there exists $M$ such that
                    $M$ is a neighborhood of $h$ in $F$ with respect to $T_1$ and
                    $M\subseteq A$.
                \begin{subproof}
                    Fix $h\in A$.
                    Hence $h\in F$ by \cref{subseteq}.
                    Take $\delta, D$ such that
                        $\delta\in\reals$ and
                        $0 < \delta$ and
                        $D$ is the compact-convergence ball about $h$ over $C$ with radius $\delta$
                            from $X$ to $Y$ with respect to $T$ and $d$ and
                        $D\subseteq A$ by \cref{compact_convergence_ball_refinement}.
                    Take $\eta$ such that
                        $\eta\in\reals$ and
                        $0 < \eta$ and
                        $\eta < 1$ and
                        $\eta < \delta$ by \cref{positive_real_has_smaller_positive_below_one}.
                    Take $M$ such that
                        $M$ is open in $F$ with respect to $\subbasisTopology{S}{F}$ and
                        $h\in M$ and
                        for all $g$ if $g\in M$ then
                            $g\in F$ and
                            for all $x\in C$ we have $g(x)\in\metricBall{d}{Y}{h(x)}{\eta}$
                        by \cref{finite_point_open_metric_neighborhoods_open}.
                    Hence $M$ is open in $F$ with respect to $T_1$ by \cref{pair_eq_iff}.
                    Take $\sigma$ such that
                        $\sigma$ is the uniform metric on functions from $C$ to $Y$ with respect to $d$ and
                        $D = \{g\in F \mid \sigma((\restrl{h}{C},\restrl{g}{C})) < \delta\}$ by \cref{compact_convergence_ball}.
                    Show that for all $g$ if $g\in M$ then $g\in D$.
                    \begin{subproof}
                        Fix $g$.
                        Assume $g\in M$.
                        Hence $g\in F$ and for all $x\in C$ we have $g(x)\in\metricBall{d}{Y}{h(x)}{\eta}$ by \cref{subseteq}.
                        Therefore $\restrl{h}{C}\in\funs{C}{Y}$ and $\restrl{g}{C}\in\funs{C}{Y}$
                            by \cref{restrl_typed_function,compact_subspace,subseteq}.
                        Show that for all $x\in C$ we have $\apply{\boundedMetric{d}}{(\apply{\restrl{h}{C}}{x}, \apply{\restrl{g}{C}}{x})} \leq \eta$.
                        \begin{subproof}
                            Follows by \cref{subseteq,metric_ball,funs_type_apply,times_tuple_intro,metric_on,real_one_in_naturals,real_naturals_subset_reals,real_lt_trans,bounded_metric_apply_below_one,real_lt_subst_left,real_lt_imp_leq,restrl_of_function_apply,compact_subspace,funs_is_function,funs_is_relation,funs_elim,pair_eq_iff}.
                        \end{subproof}
                        Hence $\sigma((\restrl{h}{C},\restrl{g}{C})) \leq \eta$ by \cref{uniform_metric_leq_from_coordinate_bounds}.
                        Therefore $\sigma((\restrl{h}{C},\restrl{g}{C})) < \delta$
                            by \cref{uniform_metric_on_function_space,metric_on,funs_type_apply,times_tuple_intro,real_leq_split,real_lt_trans,real_lt_subst_left}.
                        Hence $g\in D$ by \cref{subseteq}.
                    \end{subproof}
                    Hence $M\subseteq D$ by \cref{subseteq}.
                    Therefore $M\subseteq A$ by \cref{subseteq_transitive}.
                    Hence $M$ is a neighborhood of $h$ in $F$ with respect to $T_1$ by \cref{neighborhood}.
                    Follows by definition.
                \end{subproof}
            Therefore $A\in T_1$ by \cref{open_from_local_neighborhoods,open_in}.
        \end{byCase}
    \end{subproof}
    Hence $B\subseteq T_1$ by \cref{subseteq}.
    Show that for all $W$ if $W\in T_0$ then $W\in T_1$.
    \begin{subproof}
        Follows by \cref{basis_topology_unions,basis_for_topology,subseteq_transitive,topology_on,subseteq}.
    \end{subproof}
    Therefore $T_0\subseteq T_1$ by \cref{subseteq}.
    Thus $T_1\subseteq T_0$ by \cref{finer_topology}.
    Therefore $T_0 = T_1$ by \cref{subseteq_antisymmetric}.
\end{proof}

% from §46 Item 11: compact-open set
\begin{definition}\label{compact_open_set}
    $A$ is the set of continuous functions from $X$ to $Y$ that carry $C$ into $V$ with respect to $T$ and $U$ iff
        $C$ is a compact subspace of $X$ with respect to $T$ and
        $V$ is open in $Y$ with respect to $U$ and
        there exists $F$ such that
            $F$ is the set of continuous functions from $X$ to $Y$ with respect to $T$ and $U$ and
            $A = \{f\in F \mid \img{f}{C}\subseteq V\}$.
\end{definition}

% from §46 helper definition 5: compact-open subbasis on a continuous function space
\begin{definition}\label{compact_open_function_space_subbasis}
    $S$ is the compact-open subbasis on continuous functions from $X$ to $Y$ with respect to $T$ and $U$ iff
        $T$ is a topology on $X$ and
        $U$ is a topology on $Y$ and
        $S = \{A\in\pow{\funs{X}{Y}} \mid \text{there exist $C_0, V$ such that
            $A$ is the set of continuous functions from $X$ to $Y$ that carry $C_0$ into $V$
                with respect to $T$ and $U$}\}$.
\end{definition}

% from §46 Item 11: compact-open topology
\begin{definition}\label{compact_open_topology}
    $K$ is the compact open topology on continuous functions from $X$ to $Y$ with respect to $T$ and $U$ iff
        $T$ is a topology on $X$ and
        $U$ is a topology on $Y$ and
        there exists $F, S$ such that
            $F$ is the set of continuous functions from $X$ to $Y$ with respect to $T$ and $U$ and
            $S$ is the compact-open subbasis on continuous functions from $X$ to $Y$ with respect to $T$ and $U$ and
            $K = \subbasisTopology{S}{F}$.
\end{definition}

% from §46 helper definition 7: closed metric ball
\begin{definition}\label{closed_metric_ball}
    $A$ is the closed metric ball in $X$ about $x$ with radius $\epsilon$ with respect to $d$ iff
        $d$ is a metric on $X$ and
        $x\in X$ and
        $\epsilon\in\reals$ and
        $0 \leq \epsilon$ and
        $A = \{y\in X \mid d((x,y)) \leq \epsilon\}$.
\end{definition}

% from §46 helper lemma 19: closed metric balls lie in the ambient space
\begin{lemma}\label{closed_metric_ball_subseteq}
    Assume $A$ is the closed metric ball in $X$ about $x$ with radius $\epsilon$ with respect to $d$.
    Then $A\subseteq X$.
\end{lemma}
\begin{proof}
    Follows by \cref{closed_metric_ball,subseteq}.
\end{proof}

% from §46 helper lemma 20: open metric balls lie in the corresponding closed balls
\begin{lemma}\label{metric_ball_subset_closed_metric_ball}
    Assume $A$ is the closed metric ball in $X$ about $x$ with radius $\epsilon$ with respect to $d$.
    Assume $y\in\metricBall{d}{X}{x}{\epsilon}$.
    Then $y\in A$.
\end{lemma}
\begin{proof}
    Follows by \cref{metric_ball,real_lt_imp_leq,closed_metric_ball}.
\end{proof}

% from §46 helper lemma 21: closed metric balls are closed
\begin{lemma}\label{closed_metric_ball_closed}
    Assume $A$ is the closed metric ball in $X$ about $x$ with radius $\epsilon$ with respect to $d$.
    Let $U = \metricTopology{d}{X}$.
    Then $A$ is closed in $X$ with respect to $U$.
\end{lemma}
\begin{proof}
    We have $d$ is a metric on $X$ and $x\in X$ and $\epsilon\in\reals$ and $0 \leq \epsilon$ by \cref{closed_metric_ball}.
    Therefore $U$ is a topology on $X$ by \cref{metric_basis_is_basis,basis_topology_is_topology,metric_topology,basis_for_topology}.
    Let $B = X\setminus A$.
    Hence $B\subseteq X$ by \cref{setminus_subseteq}.
    Show that for all $z$ if $z\in B$ then there exists $V$ such that
        $V$ is a neighborhood of $z$ in $X$ with respect to $U$ and
        $V\subseteq B$.
    \begin{subproof}
        Fix $z$.
        Assume $z\in B$.
        Hence $z\in X$ and $z\notin A$ by \cref{setminus_elim_left,setminus_elim_right}.
        Therefore $d((x,z))\in\reals$ by \cref{metric_on,times_tuple_intro,funs_type_apply}.
        Hence $\epsilon < d((x,z))$ by \cref{closed_metric_ball,real_lt_trichotomy}.
        Let $\delta = d((x,z)) - \epsilon$.
        We have $\delta\in\reals$ by \cref{metric_on,times_tuple_intro,funs_type_apply,real_sub,real_add_inverse,real_add_closed}.
        Therefore $0 < \delta$ by \cref{real_lt_imp_pos_diff}.
        Let $V = \metricBall{d}{X}{z}{\delta}$.
        Hence $V$ is open in $X$ with respect to $U$ by
            \cref{metric_ball,subseteq,powerset_intro,metric_basis,metric_basis_is_basis,basis_elem_open_in_generated,metric_topology,open_in}.
        Therefore $z\in V$ by \cref{metric_zero_characterization,metric_on,real_lt_subst_left,metric_ball}.
        Show that for all $y$ if $y\in V$ then $y\in B$.
        \begin{subproof}
            Fix $y$.
            Assume $y\in V$.
            Hence $y\in X$ and $d((z,y)) < \delta$ by \cref{metric_ball}.
            Therefore $d((x,y))\in\reals$ and $d((z,y))\in\reals$ by \cref{metric_on,times_tuple_intro,funs_type_apply}.
            Suppose not.
            Hence $y\in A$ by \cref{setminus}.
            Therefore $\epsilon + d((z,y)) < d((x,z))$ by
                \cref{metric_on,times_tuple_intro,funs_type_apply,real_sub,real_add_inverse,real_add_closed,real_lt_subst_right,real_lt_add,real_add_assoc,real_add_comm,real_add_ident,real_lt_subst_left}.
            Hence $d((x,y)) + d((z,y)) < d((x,z))$ by
                \cref{closed_metric_ball,real_add_closed,real_leq_split,real_lt_add,real_add_eq_subst,real_lt_trans,real_lt_subst_left}.
            Hence $d((y,z)) = d((z,y))$ by \cref{metric_on,metric_symmetric}.
            Therefore $d((x,y)) + d((y,z)) < d((x,z))$ by \cref{real_add_eq_subst}.
            We have $d((x,z)) \leq d((x,y)) + d((y,z))$ by \cref{metric_on,metric_triangle_inequality}.
            Hence $d((x,z)) < d((x,z))$ by
                \cref{real_add_closed,real_leq_split,real_lt_trans,real_lt_subst_left}.
            Contradiction by \cref{real_lt_irreflexive}.
            Thus $y\in B$.
        \end{subproof}
        Hence $V\subseteq B$ by \cref{subseteq}.
        Therefore $V$ is a neighborhood of $z$ in $X$ with respect to $U$ by \cref{neighborhood}.
    \end{subproof}
    Hence $B$ is open in $X$ with respect to $U$ by \cref{open_from_local_neighborhoods}.
    Therefore $A$ is closed in $X$ with respect to $U$ by \cref{closed_metric_ball_subseteq,closed_in}.
\end{proof}

% from §46 helper lemma 22: compact-open subbasis is a subbasis
\begin{lemma}\label{compact_open_subbasis_is_subbasis}
    Assume $T$ is a topology on $X$.
    Assume $U$ is a topology on $Y$.
    Assume $C$ is the set of continuous functions from $X$ to $Y$ with respect to $T$ and $U$.
    Assume $S$ is the compact-open subbasis on continuous functions from $X$ to $Y$ with respect to $T$ and $U$.
    Then $S$ is a subbasis on $C$.
\end{lemma}
\begin{proof}
    We have $S\subseteq\pow{C}$ by
        \cref{compact_open_function_space_subbasis,compact_open_set,continuous_function_set,subseteq,subseteq_antisymmetric,powerset_intro}.
    Let $E = \{f\in C \mid \img{f}{\emptyset}\subseteq Y\}$.
    Hence $\emptyset$ is a compact subspace of $X$ with respect to $T$
        by \cref{emptyset_subseteq,subspace_topology_is_topology,emptyset_compact,compact_subspace}.
    Therefore $E$ is the set of continuous functions from $X$ to $Y$ that carry $\emptyset$ into $Y$ with respect to $T$ and $U$
        by \cref{open_in,topology_on,compact_open_set}.
    Hence $E\in S$ by \cref{subseteq,continuous_function_set,powerset_intro,compact_open_function_space_subbasis}.
    Show that for all $f$ if $f\in C$ then $f\in\unions{S}$.
    \begin{subproof}
        Trivial.
    \end{subproof}
    Hence $C\subseteq\unions{S}$ by \cref{subseteq}.
    Therefore $\unions{S} = C$ by \cref{subseteq,pow_iff,unions_subseteq_of_powerset_is_subseteq,subseteq_antisymmetric}.
    Hence $S$ is a subbasis on $C$ by \cref{subbasis_on}.
\end{proof}

% from §46 helper lemma 23: half-radius closed balls lie in full-radius open balls
\begin{lemma}\label{half_radius_closed_ball_subset_open_ball}
    Assume $A$ is the closed metric ball in $Y$ about $y$ with radius $\delta$ with respect to $d$.
    Assume $\epsilon\in\reals$.
    Assume $0 < \delta$.
    Assume $\delta + \delta = \epsilon$.
    Assume $z\in A$.
    Then $z\in\metricBall{d}{Y}{y}{\epsilon}$.
\end{lemma}
\begin{proof}
    We have $d$ is a metric on $Y$ and $y\in Y$ and $\delta\in\reals$ and $z\in Y$ and $d((y,z)) \leq \delta$
        by \cref{closed_metric_ball}.
    Therefore $d((y,z))\in\reals$ and $\delta + \delta\in\reals$ and $\delta < \delta + \delta$ by
        \cref{times_tuple_intro,metric_on,funs_type_apply,real_add_closed,real_add_ident,real_lt_add,real_lt_subst_left}.
    Hence $z\in\metricBall{d}{Y}{y}{\epsilon}$
        by \cref{real_leq_split,real_lt_trans,real_lt_subst_left,real_lt_subst_right,metric_ball}.
\end{proof}

% from §46 helper lemma 24: compact-open subbasis elements are open
\begin{lemma}\label{compact_open_subbasis_element_open}
    Assume $T$ is a topology on $X$.
    Assume $U$ is a topology on $Y$.
    Assume $C$ is the set of continuous functions from $X$ to $Y$ with respect to $T$ and $U$.
    Assume $S$ is the compact-open subbasis on continuous functions from $X$ to $Y$ with respect to $T$ and $U$.
    Assume $K$ is the compact open topology on continuous functions from $X$ to $Y$ with respect to $T$ and $U$.
    Assume $A\in S$.
    Then $A$ is open in $C$ with respect to $K$.
\end{lemma}
\begin{proof}
    Take $F, S_0$ such that
        $F$ is the set of continuous functions from $X$ to $Y$ with respect to $T$ and $U$ and
        $S_0$ is the compact-open subbasis on continuous functions from $X$ to $Y$ with respect to $T$ and $U$ and
        $K = \subbasisTopology{S_0}{F}$ by \cref{compact_open_topology}.
    Hence $K = \subbasisTopology{S}{C}$
        by \cref{continuous_function_set,setext,compact_open_function_space_subbasis,setext,pair_eq_iff}.
    Therefore $S$ is a subbasis on $C$ by \cref{compact_open_subbasis_is_subbasis}.
    Hence $A\in K$ by \cref{subbasis_on,subseteq,pow_iff,singleton_subset_intro,upair_iff,pair_is_finite,subseteq_finite_is_finite,singleton_intro,inters_singleton,finite_intersections,subbasis_finite_intersections_basis,basis_elem_open_in_generated,basis_for_topology,topology_generated_by_subbasis}.
    Therefore $A$ is open in $C$ with respect to $K$ by
        \cref{subbasis_finite_intersections_basis,basis_for_topology,basis_topology_is_topology,topology_generated_by_subbasis,open_in}.
\end{proof}

% from §46 helper lemma 25: compact-open neighborhoods are locally compact-convergence-open
\begin{lemma}\label{compact_open_local_compact_convergence}
    Assume $T$ is a topology on $X$.
    Assume $d$ is a metric on $Y$.
    Let $U = \metricTopology{d}{Y}$.
    Assume $C$ is the set of continuous functions from $X$ to $Y$ with respect to $T$ and $U$.
    Assume $T_0$ is the topology of compact convergence on functions from $X$ to $Y$ with respect to $T$ and $d$.
    Assume $A$ is the set of continuous functions from $X$ to $Y$ that carry $C_0$ into $V$ with respect to $T$ and $U$.
    Assume $h\in A$.
    Then there exists $M$ such that
        $M$ is open in $\funs{X}{Y}$ with respect to $T_0$ and
        $h\in M$ and
        for all $g$ if $g\in M\inter C$ then $g\in A$.
\end{lemma}
\begin{proof}
    Take $F$ such that
        $F$ is the set of continuous functions from $X$ to $Y$ with respect to $T$ and $U$ and
        $A = \{f\in F \mid \img{f}{C_0}\subseteq V\}$ by \cref{compact_open_set}.
    We have $F = C$ by \cref{continuous_function_set,setext}.
    $h\in F$ and $\img{h}{C_0}\subseteq V$ by \cref{inters_iff_forall}.
    $h\in\funs{X}{Y}$ and $h$ is continuous from $X$ to $Y$ with respect to $T$ and $U$
        by \cref{pair_eq_iff,continuous_function_set}.
    \begin{byCase}
    \caseOf{$C_0 = \emptyset$.}
        Let $M = \funs{X}{Y}$.
        $M$ is open in $\funs{X}{Y}$ with respect to $T_0$
            by \cref{compact_convergence_topology,open_in,topology_on,pair_eq_iff}.
        $h\in M$ by \cref{subseteq}.
        Show that for all $g$ if $g\in M\inter C$ then $g\in A$.
        \begin{subproof}
            Trivial.
        \end{subproof}
        Follows by definition.
    \caseOf{$C_0 \neq \emptyset$.}
        $C_0$ is a compact subspace of $X$ with respect to $T$ and
            $V$ is open in $Y$ with respect to $U$ by \cref{compact_open_set}.
        $V\subseteq Y$ by \cref{open_in,topology_on,subseteq,powerset_elim}.
        $C_0\subseteq X$ by \cref{compact_subspace}.
        $C_0$ is a subspace of $X$ with respect to $T$ by \cref{subspace_of}.
        Let $T_C = \subspaceTopology{T}{C_0}$.
        $T_C$ is a topology on $C_0$ by \cref{subspace_topology_is_topology}.
        $C_0$ is compact with respect to $T_C$ by \cref{compact_subspace}.
        $\restrl{h}{C_0}$ is continuous from $C_0$ to $Y$ with respect to $T_C$ and $U$
            by \cref{continuous_restrict_domain}.
        Let $D = \{N\in\pow{C_0} \mid \text{there exist $x, \epsilon, \delta$ such that
            $x\in C_0$ and
            $\epsilon\in\reals$ and
            $\delta\in\reals$ and
            $0 < \epsilon$ and
            $0 < \delta$ and
            $\delta < 1$ and
            $\delta + \delta = \epsilon$ and
            $\metricBall{d}{Y}{h(x)}{\epsilon}\subseteq V$ and
            $N = \preimg{\restrl{h}{C_0}}{\metricBall{d}{Y}{h(x)}{\delta}}$}\}$.
        Show that for all $x\in C_0$ we have $x\in\unions{D}$.
        \begin{subproof}
            Fix $x\in C_0$.
            $x\in X$ by \cref{subseteq}.
            $x\in\dom{h}$ by \cref{continuous,funs_elim,subseteq}.
            We have $x\in\dom{h}\inter C_0$ by \cref{inter_intro}.
            $h(x)\in\img{h}{C_0}$ by \cref{continuous,funs_is_function,img_of_function_intro}.
            Thus $h(x)\in V$ by \cref{subseteq}.
            Take $\epsilon_0\in\reals$ such that $0 < \epsilon_0$ and
                $\metricBall{d}{Y}{h(x)}{\epsilon_0}\subseteq V$ by \cref{metric_open_iff}.
            Take $\epsilon$ such that
                $\epsilon\in\reals$ and
                $0 < \epsilon$ and
                $\epsilon < 1$ and
                $\epsilon < \epsilon_0$ by \cref{positive_real_has_smaller_positive_below_one}.
            Let $\delta = \epsilon / (1 + 1)$.
            $\delta\in\reals$ and $0 < \delta$ and $\delta + \delta = \epsilon$
                by \cref{real_one_in_naturals,real_naturals_subset_reals,subseteq,real_zero_lt_one,real_half_of_positive}.
            We have $\delta < \epsilon$ by \cref{real_add_ident,real_lt_add,real_lt_subst_left,real_lt_subst_right}.
            $\delta < 1$ by \cref{real_one_in_naturals,real_naturals_subset_reals,subseteq,real_lt_trans}.
            $h(x)\in Y$ by \cref{funs_type_apply}.
            $\metricBall{d}{Y}{h(x)}{\epsilon}\subseteq\metricBall{d}{Y}{h(x)}{\epsilon_0}$ by \cref{funs_type_apply,metric_ball_mono}.
            Therefore $\metricBall{d}{Y}{h(x)}{\epsilon}\subseteq V$ by \cref{subseteq_transitive}.
            Let $N = \preimg{\restrl{h}{C_0}}{\metricBall{d}{Y}{h(x)}{\delta}}$.
            We have $N\in D$ by \cref{preimg_subseteq_dom,restrl_dom,subseteq,powerset_intro}.
            $d((h(x),h(x))) = 0$ by \cref{metric_on,metric_zero_characterization,times_tuple_intro,funs_type_apply}.
            $d((h(x),h(x))) < \delta$ by \cref{real_lt_subst_left}.
            Thus $h(x)\in\metricBall{d}{Y}{h(x)}{\delta}$ by \cref{metric_ball}.
            $(x,h(x))\in h$ by \cref{function_apply_elim,funs_is_function}.
            $(x,h(x))\in\restrl{h}{C_0}$ by \cref{restrl_iff}.
            $x\in N$ by \cref{preimg_iff}.
            Therefore $x\in\unions{D}$ by \cref{unions_iff}.
        \end{subproof}
        $C_0\subseteq\unions{D}$ by \cref{subseteq}.
        Therefore $D$ is a covering of $C_0$ by \cref{covering}.
        Show that for all $N$ if $N\in D$ then $N$ is open in $C_0$ with respect to $T_C$.
        \begin{subproof}
            Follows by \cref{subseteq,funs_type_apply,metric_ball,powerset_intro,metric_basis,metric_basis_is_basis,basis_elem_open_in_generated,metric_topology,open_in,basis_topology_is_topology,basis_for_topology,continuous}.
        \end{subproof}
        Hence $D$ is an open covering of $C_0$ with respect to $T_C$ by \cref{open_covering}.
        Take $G$ such that $G\subseteq D$ and $G$ is finite and $G$ is a covering of $C_0$
            by \cref{compact}.
        Let $R = \{w\in G\times\pow{\funs{X}{Y}} \mid \text{there exist $N, M, x, \epsilon, \delta, B_0$ such that
            $N\in G$ and
            $M\in\pow{\funs{X}{Y}}$ and
            $x\in C_0$ and
            $\epsilon\in\reals$ and
            $\delta\in\reals$ and
            $w = (N,M)$ and
            $0 < \epsilon$ and
            $0 < \delta$ and
            $\delta < 1$ and
            $\delta + \delta = \epsilon$ and
            $B_0$ is the closed metric ball in $Y$ about $h(x)$ with radius $\delta$ with respect to $d$ and
            $\metricBall{d}{Y}{h(x)}{\epsilon}\subseteq V$ and
            $N = \preimg{\restrl{h}{C_0}}{\metricBall{d}{Y}{h(x)}{\delta}}$ and
            $M$ is the compact-convergence ball about $h$ over $\preimg{\restrl{h}{C_0}}{B_0}$ with radius $\delta$
                from $X$ to $Y$ with respect to $T$ and $d$}\}$.
        We have $R$ is a relation by \cref{relation,subseteq,times_elem_is_tuple}.
        Show that for all $N\in G$ there exists $M$ such that $N\mathrel{R} M$.
        \begin{subproof}
            Fix $N\in G$.
            Take $x, \epsilon, \delta$ such that
                $x\in C_0$ and
                $\epsilon\in\reals$ and
                $\delta\in\reals$ and
                $0 < \epsilon$ and
                $0 < \delta$ and
                $\delta < 1$ and
                $\delta + \delta = \epsilon$ and
                $\metricBall{d}{Y}{h(x)}{\epsilon}\subseteq V$ and
                $N = \preimg{\restrl{h}{C_0}}{\metricBall{d}{Y}{h(x)}{\delta}}$ by \cref{subseteq}.
            $h(x)\in Y$ by \cref{subseteq,funs_type_apply}.
            Let $B_0 = \{y\in Y \mid d((h(x),y)) \leq \delta\}$.
            $B_0$ is the closed metric ball in $Y$ about $h(x)$ with radius $\delta$ with respect to $d$
                by \cref{subseteq,funs_type_apply,closed_metric_ball}.
            Let $K = \preimg{\restrl{h}{C_0}}{B_0}$.
            $B_0$ is closed in $Y$ with respect to $U$ by \cref{closed_metric_ball_closed}.
            $\preimg{\restrl{h}{C_0}}{Y\setminus B_0}$ is open in $C_0$ with respect to $T_C$ by \cref{continuous,closed_in}.
            We have $\preimg{\restrl{h}{C_0}}{Y\setminus B_0} = C_0\setminus K$ by
                \cref{continuous,closed_in,setminus_subseteq,preimg_setminus_function,pair_eq_iff}.
            Therefore $C_0\setminus K$ is open in $C_0$ with respect to $T_C$ by \cref{pair_eq_iff}.
            We have $K\subseteq C_0$ by \cref{preimg_subseteq_dom,restrl_dom,subseteq}.
            $K$ is compact with respect to $\subspaceTopology{T}{K}$
                by \cref{closed_in,closed_subspace_compact,subspace_subspace_topology,subseteq}.
            $K\subseteq X$ by \cref{subseteq_transitive}.
            Therefore $K$ is a compact subspace of $X$ with respect to $T$ by \cref{compact_subspace}.
            $x\in\dom{h}$ by \cref{continuous,funs_elim,subseteq}.
            $d((h(x),h(x))) = 0$ by \cref{metric_on,metric_zero_characterization,times_tuple_intro,funs_type_apply}.
            $h(x)\in B_0$ by \cref{closed_metric_ball,real_lt_subst_left}.
            $(x,h(x))\in h$ by \cref{function_apply_elim,funs_is_function}.
            $(x,h(x))\in\restrl{h}{C_0}$ by \cref{restrl_iff}.
            Hence $x\in K$ by \cref{preimg_iff}.
            Take $\rho$ such that $\rho$ is the uniform metric on functions from $K$ to $Y$ with respect to $d$
                by \cref{uniform_metric_on_function_space_exists_inhabited,nonempty_inhabited,metric_ball_subset_closed_metric_ball,subseteq}.
            Let $M = \{g\in\funs{X}{Y} \mid \rho((\restrl{h}{K},\restrl{g}{K})) < \delta\}$.
            $M$ is the compact-convergence ball about $h$ over $K$ with radius $\delta$
                from $X$ to $Y$ with respect to $T$ and $d$ by \cref{compact_convergence_ball}.
            $M$ is the compact-convergence ball about $h$ over $\preimg{\restrl{h}{C_0}}{B_0}$ with radius $\delta$
                from $X$ to $Y$ with respect to $T$ and $d$ by \cref{pair_eq_iff}.
            Let $w = (N,M)$.
            $w\in G\times\pow{\funs{X}{Y}}$ by \cref{times_tuple_intro,subseteq,powerset_intro}.
            We have $N\in G$ and $M\in\pow{\funs{X}{Y}}$ by \cref{fst_eq,snd_eq,times_tuple_elim}.
            $w\in R$ by \cref{subseteq}.
            Therefore $N\mathrel{R} M$ by \cref{relation}.
        \end{subproof}
        Take $m$ such that $m$ is a function on $G$ and for all $N\in G$ we have $N\mathrel{R}\apply{m}{N}$ by \cref{choice_function}.
        Let $H = \img{m}{G}$.
        Let $W = \inters{H}$.
        We have $H\subseteq T_0$ by \cref{img_iff,subseteq,function_apply_iff,relation,pair_eq_iff,compact_convergence_ball_open,compact_convergence_topology,open_in,topology_on}.
        $H$ is finite by \cref{pair_eq_iff,finite_image_function}.
        $G$ is inhabited by \cref{covering,subseteq,unions_iff,nonempty_inhabited}.
        $H$ is inhabited by \cref{pair_eq_iff,function_apply_elim,img_iff,nonempty_inhabited}.
        Hence $W$ is open in $\funs{X}{Y}$ with respect to $T_0$
            by \cref{compact_convergence_topology,finite_intersections_open_in_topology,open_in}.
        Show that for all $P$ if $P\in H$ then $h\in P$.
        \begin{subproof}
            Follows by \cref{img_iff,subseteq,function_apply_iff,relation,pair_eq_iff,compact_convergence_ball_center}.
        \end{subproof}
        $h\in W$ by \cref{inters_iff_forall}.
        Show that for all $g$ if $g\in W\inter C$ then $g\in A$.
        \begin{subproof}
            Fix $g$.
            Assume $g\in W\inter C$.
            We have $g\in W$ and $g\in C$ by \cref{inter_elim_left,inter_elim_right}.
            $g\in F$ by \cref{pair_eq_iff}.
            Show that for all $z$ if $z\in\img{g}{C_0}$ then $z\in V$.
            \begin{subproof}
                    Fix $z$.
                    Assume $z\in\img{g}{C_0}$.
                    $g\in\funs{X}{Y}$ by \cref{continuous_function_set}.
                    Take $x\in\dom{g}\inter C_0$ such that $z = g(x)$ by \cref{funs_is_function,img_of_function_elim}.
                    $x\in C_0$ by \cref{inter_elim_right}.
                    Take $N\in G$ such that $x\in N$ by \cref{covering,subseteq,unions_iff}.
                    $N\in\dom{m}$ by \cref{subseteq}.
                    $N\mathrel{R}\apply{m}{N}$ by \cref{subseteq}.
                    Take $x_0, \epsilon, \delta, B_0$ such that
                        $x_0\in C_0$ and
                        $\epsilon\in\reals$ and
                        $\delta\in\reals$ and
                        $0 < \epsilon$ and
                        $0 < \delta$ and
                        $\delta < 1$ and
                        $\delta + \delta = \epsilon$ and
                        $B_0$ is the closed metric ball in $Y$ about $h(x_0)$ with radius $\delta$ with respect to $d$ and
                        $\metricBall{d}{Y}{h(x_0)}{\epsilon}\subseteq V$ and
                        $N = \preimg{\restrl{h}{C_0}}{\metricBall{d}{Y}{h(x_0)}{\delta}}$ and
                        $\apply{m}{N}$ is the compact-convergence ball about $h$ over $\preimg{\restrl{h}{C_0}}{B_0}$ with radius $\delta$
                            from $X$ to $Y$ with respect to $T$ and $d$
                        by \cref{relation,pair_eq_iff}.
                    Take $y\in\metricBall{d}{Y}{h(x_0)}{\delta}$ such that $(x,y)\in\restrl{h}{C_0}$ by \cref{preimg_iff}.
                    $x\in X$ by \cref{subseteq}.
                    $x\in\dom{h}$ by \cref{continuous,funs_elim,subseteq}.
                    $(x,h(x))\in h$ by \cref{function_apply_elim,funs_is_function}.
                    $(x,h(x))\in\restrl{h}{C_0}$ by \cref{restrl_iff}.
                    $\restrl{h}{C_0}\in\funs{C_0}{Y}$ by \cref{continuous,funs}.
                    Therefore $y = h(x)$ by \cref{function_rightunique,funs_is_function}.
                    $h(x)\in\metricBall{d}{Y}{h(x_0)}{\delta}$ by \cref{pair_eq_iff}.
                    $x\in\preimg{\restrl{h}{C_0}}{B_0}$ by \cref{metric_ball_subset_closed_metric_ball,preimg_iff}.
                    Let $K = \preimg{\restrl{h}{C_0}}{B_0}$.
                    $x\in K$ by \cref{pair_eq_iff}.
                    $\apply{m}{N}\in H$ by \cref{pair_eq_iff,function_apply_elim,img_iff}.
                    $g\in\apply{m}{N}$ by \cref{inters_iff_forall}.
                    Take $\rho$ such that
                        $\rho$ is the uniform metric on functions from $K$ to $Y$ with respect to $d$ and
                        $\apply{m}{N} = \{u\in\funs{X}{Y} \mid \rho((\restrl{h}{K},\restrl{u}{K})) < \delta\}$
                        by \cref{compact_convergence_ball}.
                    We have $g\in\funs{X}{Y}$ and $\rho((\restrl{h}{K},\restrl{g}{K})) < \delta$ by \cref{inters_iff_forall}.
                    $K\subseteq X$ by \cref{preimg_subseteq_dom,restrl_dom,subseteq}.
                    $\restrl{h}{K}\in\funs{K}{Y}$ and $\restrl{g}{K}\in\funs{K}{Y}$ by \cref{restrl_typed_function}.
                    $d((\apply{\restrl{h}{K}}{x}, \apply{\restrl{g}{K}}{x})) < \delta$
                        by \cref{uniform_metric_small_implies_pointwise_small}.
                    $\apply{\restrl{h}{K}}{x} = h(x)$ and $\apply{\restrl{g}{K}}{x} = g(x)$
                        by \cref{restrl_of_function_apply,subseteq,funs_is_function,funs_is_relation,funs_elim}.
                    Hence $d((h(x),g(x))) < \delta$ by \cref{pair_eq_iff}.
                    $h(x_0)\in Y$ by \cref{funs_type_apply,subseteq}.
                    $h(x)\in Y$ by \cref{metric_ball,subseteq}.
                    $g(x)\in Y$ by \cref{funs_type_apply}.
                    $d((h(x_0),h(x))) < \delta$ by \cref{metric_ball}.
                    $(h(x_0),h(x))\in Y\times Y$ and $(h(x),g(x))\in Y\times Y$ and $(h(x_0),g(x))\in Y\times Y$
                        by \cref{times_tuple_intro}.
                    $d((h(x_0),g(x))) \leq d((h(x_0),h(x))) + d((h(x),g(x)))$
                        by \cref{metric_on,metric_triangle_inequality}.
                    $d((h(x_0),h(x)))\in\reals$ and $d((h(x),g(x)))\in\reals$ and $d((h(x_0),g(x)))\in\reals$
                        by \cref{metric_on,funs_type_apply}.
                    $d((h(x_0),h(x))) + d((h(x),g(x))) < \delta + \delta$ by \cref{real_lt_add_two}.
                    $d((h(x_0),g(x))) < \delta + \delta$
                        by \cref{real_leq_split,real_lt_trans,real_lt_subst_left,real_add_closed}.
                    Hence $d((h(x_0),g(x))) < \epsilon$ by \cref{real_lt_subst_right}.
                    $g(x)\in\metricBall{d}{Y}{h(x_0)}{\epsilon}$ by \cref{metric_ball}.
                    Thus $g(x)\in V$ by \cref{subseteq}.
                    Therefore $z\in V$ by \cref{pair_eq_iff}.
            \end{subproof}
            $\img{g}{C_0}\subseteq V$ by \cref{subseteq}.
            Therefore $g\in A$ by \cref{inters_iff_forall}.
        \end{subproof}
        Follows by definition.
    \end{byCase}
\end{proof}

% from §46 helper lemma 26: preimages of closed metric balls on compact subspaces are compact
\begin{lemma}\label{closed_ball_preimage_compact_subspace}
    Assume $T$ is a topology on $X$.
    Assume $d$ is a metric on $Y$.
    Let $U = \metricTopology{d}{Y}$.
    Assume $h$ is continuous from $X$ to $Y$ with respect to $T$ and $U$.
    Assume $C_0$ is a compact subspace of $X$ with respect to $T$.
    Assume $B_0$ is the closed metric ball in $Y$ about $y$ with radius $\delta$ with respect to $d$.
    Let $Q = \preimg{\restrl{h}{C_0}}{B_0}$.
    Then $Q$ is a compact subspace of $X$ with respect to $T$.
\end{lemma}
\begin{proof}
    $C_0\subseteq X$ by \cref{compact_subspace}.
    Let $T_C = \subspaceTopology{T}{C_0}$.
    $T_C$ is a topology on $C_0$ by \cref{compact_subspace,subseteq,subspace_topology_is_topology}.
    $C_0$ is compact with respect to $T_C$ by \cref{compact_subspace}.
    $\restrl{h}{C_0}$ is continuous from $C_0$ to $Y$ with respect to $T_C$ and $U$
        by \cref{subspace_of,compact_subspace,continuous_restrict_domain}.
    $B_0$ is closed in $Y$ with respect to $U$ by \cref{closed_metric_ball_closed}.
    $\preimg{\restrl{h}{C_0}}{Y\setminus B_0}$ is open in $C_0$ with respect to $T_C$ by \cref{continuous,closed_in}.
    We have $\preimg{\restrl{h}{C_0}}{Y\setminus B_0} = C_0\setminus Q$ by
        \cref{continuous,funs,closed_in,setminus_subseteq,preimg_setminus_function,pair_eq_iff}.
    Therefore $C_0\setminus Q$ is open in $C_0$ with respect to $T_C$ by \cref{pair_eq_iff}.
    $Q\subseteq C_0$ by \cref{preimg_subseteq_dom,restrl_dom,subseteq}.
    $Q$ is compact with respect to $\subspaceTopology{T}{Q}$
        by \cref{closed_in,closed_subspace_compact,subspace_subspace_topology,subseteq}.
    Hence $Q$ is a compact subspace of $X$ with respect to $T$ by \cref{compact_subspace,subseteq_transitive}.
\end{proof}

% from §46 helper lemma 27: compact-convergence neighborhoods are locally compact-open
\begin{lemma}\label{compact_convergence_local_compact_open}
    Assume $T$ is a topology on $X$.
    Assume $d$ is a metric on $Y$.
    Let $U = \metricTopology{d}{Y}$.
    Assume $C$ is the set of continuous functions from $X$ to $Y$ with respect to $T$ and $U$.
    Assume $K$ is the compact open topology on continuous functions from $X$ to $Y$ with respect to $T$ and $U$.
    Assume $A$ is the compact-convergence ball about $h$ over $C_0$ with radius $\epsilon$
        from $X$ to $Y$ with respect to $T$ and $d$.
    Assume $h\in C$.
    Then there exists $M$ such that
        $M$ is open in $C$ with respect to $K$ and
        $h\in M$ and
        for all $g$ if $g\in M$ then $g\in A$.
\end{lemma}
\begin{proof}
    Let $F = \funs{X}{Y}$.
    Take $F_0, S$ such that
        $F_0$ is the set of continuous functions from $X$ to $Y$ with respect to $T$ and $U$ and
        $S$ is the compact-open subbasis on continuous functions from $X$ to $Y$ with respect to $T$ and $U$ and
        $K = \subbasisTopology{S}{F_0}$ by \cref{compact_open_topology}.
    We have $K = \subbasisTopology{S}{C}$ by \cref{continuous_function_set,setext,pair_eq_iff}.
    $K$ is a topology on $C$ by \cref{compact_open_subbasis_is_subbasis,compact_open_topology,subbasis_finite_intersections_basis,basis_for_topology,basis_topology_is_topology,topology_generated_by_subbasis}.
    Take $\rho$ such that
        $\rho$ is the uniform metric on functions from $C_0$ to $Y$ with respect to $d$ and
        $A = \{g\in F \mid \rho((\restrl{h}{C_0},\restrl{g}{C_0})) < \epsilon\}$ by \cref{compact_convergence_ball}.
    $C\subseteq F$ by \cref{continuous_function_set,subseteq}.
    $h\in F$ and $h$ is continuous from $X$ to $Y$ with respect to $T$ and $U$
        by \cref{continuous_function_set}.
    $U$ is a topology on $Y$ by \cref{metric_basis_is_basis,basis_topology_is_topology,metric_topology,basis_for_topology}.
    $C_0$ is a compact subspace of $X$ with respect to $T$ and $\epsilon\in\reals$ and $0 < \epsilon$
        by \cref{compact_convergence_ball}.
    $C_0\subseteq X$ by \cref{compact_subspace}.
    $C_0$ is a subspace of $X$ with respect to $T$ by \cref{subspace_of}.
    Let $T_C = \subspaceTopology{T}{C_0}$.
    $T_C$ is a topology on $C_0$ by \cref{subspace_topology_is_topology}.
    $C_0$ is compact with respect to $T_C$ by \cref{compact_subspace}.
    $\restrl{h}{C_0}$ is continuous from $C_0$ to $Y$ with respect to $T_C$ and $U$
        by \cref{continuous_restrict_domain}.
    Take $\alpha$ such that
        $\alpha\in\reals$ and
        $0 < \alpha$ and
        $\alpha < 1$ and
        $\alpha < \epsilon$ by \cref{positive_real_has_smaller_positive_below_one}.
    Let $\eta = \alpha / (1 + 1)$.
    $\eta\in\reals$ and $0 < \eta$ and $\eta + \eta = \alpha$
        by \cref{real_one_in_naturals,real_naturals_subset_reals,subseteq,real_zero_lt_one,real_half_of_positive}.
    Let $\lambda = \eta / (1 + 1)$.
    $\lambda\in\reals$ and $0 < \lambda$ and $\lambda + \lambda = \eta$
        by \cref{real_one_in_naturals,real_naturals_subset_reals,subseteq,real_zero_lt_one,real_half_of_positive}.
    Let $D = \{N\in\pow{C_0} \mid \text{there exists $x$ such that
        $x\in C_0$ and
        $N = \preimg{\restrl{h}{C_0}}{\metricBall{d}{Y}{h(x)}{\lambda}}$}\}$.
    Show that for all $x\in C_0$ we have $x\in\unions{D}$.
    \begin{subproof}
        Fix $x\in C_0$.
        $x\in\dom{h}$ by \cref{compact_subspace,continuous,funs_elim,subseteq}.
        $h(x)\in Y$ by \cref{compact_subspace,subseteq,funs_type_apply}.
        $d((h(x),h(x))) = 0$ by \cref{metric_on,metric_zero_characterization,times_tuple_intro,funs_type_apply}.
        $h(x)\in\metricBall{d}{Y}{h(x)}{\lambda}$ by \cref{real_lt_subst_left,metric_ball}.
        Let $N = \preimg{\restrl{h}{C_0}}{\metricBall{d}{Y}{h(x)}{\lambda}}$.
        We have $N\in D$ by \cref{preimg_subseteq_dom,restrl_dom,subseteq,powerset_intro}.
        $x\in N$ by \cref{function_apply_elim,funs_is_function,restrl_iff,preimg_iff}.
        Therefore $x\in\unions{D}$ by \cref{unions_iff}.
    \end{subproof}
    $D$ is a covering of $C_0$ by \cref{subseteq,covering}.
    Show that for all $N$ if $N\in D$ then $N$ is open in $C_0$ with respect to $T_C$.
    \begin{subproof}
        Follows by \cref{subseteq,funs_type_apply,metric_ball,powerset_intro,metric_basis,metric_basis_is_basis,basis_elem_open_in_generated,metric_topology,open_in,basis_topology_is_topology,basis_for_topology,continuous}.
    \end{subproof}
    Hence $D$ is an open covering of $C_0$ with respect to $T_C$ by \cref{open_covering}.
    Take $G$ such that $G\subseteq D$ and $G$ is finite and $G$ is a covering of $C_0$
        by \cref{compact}.
    Let $R = \{w\in G\times\pow{F} \mid \text{there exist $N, P, x, B_0, L, V$ such that
        $N\in G$ and
        $P\in\pow{F}$ and
        $x\in C_0$ and
        $w = (N,P)$ and
        $B_0$ is the closed metric ball in $Y$ about $h(x)$ with radius $\lambda$ with respect to $d$ and
        $L = \preimg{\restrl{h}{C_0}}{B_0}$ and
        $V = \metricBall{d}{Y}{h(x)}{\eta}$ and
        $N = \preimg{\restrl{h}{C_0}}{\metricBall{d}{Y}{h(x)}{\lambda}}$ and
        $L$ is a compact subspace of $X$ with respect to $T$ and
        $P$ is the set of continuous functions from $X$ to $Y$ that carry $L$ into $V$ with respect to $T$ and $U$}\}$.
    We have $R$ is a relation by \cref{relation,subseteq,times_elem_is_tuple}.
    Show that for all $N\in G$ there exists $P$ such that $N\mathrel{R} P$.
    \begin{subproof}
        Fix $N\in G$.
        Take $x$ such that
            $x\in C_0$ and
            $N = \preimg{\restrl{h}{C_0}}{\metricBall{d}{Y}{h(x)}{\lambda}}$ by \cref{subseteq}.
        $h(x)\in Y$ by \cref{compact_subspace,subseteq,funs_type_apply}.
        Let $B_0 = \{y\in Y \mid d((h(x),y)) \leq \lambda\}$.
        $B_0$ is the closed metric ball in $Y$ about $h(x)$ with radius $\lambda$ with respect to $d$
            by \cref{closed_metric_ball}.
        Let $L = \preimg{\restrl{h}{C_0}}{B_0}$.
        $L$ is a compact subspace of $X$ with respect to $T$ by \cref{closed_ball_preimage_compact_subspace}.
        Let $V = \metricBall{d}{Y}{h(x)}{\eta}$.
        $V$ is open in $Y$ with respect to $U$ by
            \cref{metric_ball,subseteq,powerset_intro,metric_basis,metric_basis_is_basis,basis_elem_open_in_generated,basis_for_topology,metric_topology,basis_topology_is_topology,open_in}.
        Let $P = \{g\in C \mid \img{g}{L}\subseteq V\}$.
        $P$ is the set of continuous functions from $X$ to $Y$ that carry $L$ into $V$ with respect to $T$ and $U$
            by \cref{compact_open_set}.
        $P\in\pow{F}$ by \cref{subseteq,continuous_function_set,subseteq_transitive,powerset_intro}.
        Let $w = (N,P)$.
        $w\in R$ by \cref{times_tuple_intro,subseteq}.
        Therefore $N\mathrel{R} P$ by \cref{relation}.
    \end{subproof}
    Take $m$ such that $m$ is a function on $G$ and for all $N\in G$ we have $N\mathrel{R}\apply{m}{N}$ by \cref{choice_function}.
    Let $H_0 = \img{m}{G}$.
    Let $H = H_0\union\{C,C\}$.
    Let $M = \inters{H}$.
    We have $H_0\subseteq K$ by \cref{img_iff,subseteq,function_apply_iff,relation,pair_eq_iff,compact_open_function_space_subbasis,compact_open_subbasis_element_open,open_in,topology_on}.
    $H\subseteq K$ by \cref{union_iff,subseteq,upair_iff,pair_eq_iff,topology_on}.
    $H_0$ is finite by \cref{pair_eq_iff,finite_image_function}.
    $H$ is finite by \cref{pair_is_finite,union_of_finite_is_finite}.
    $C\in H$ by \cref{upair_intro_left,union_intro_right}.
    Thus $M\in K$ by \cref{finite_intersections_open_in_topology}.
    $M$ is open in $C$ with respect to $K$ by \cref{open_in}.
    Show that for all $P$ if $P\in H$ then $h\in P$.
    \begin{subproof}
        Fix $P$.
        Assume $P\in H$.
        $P\in H_0$ or $P\in\{C,C\}$ by \cref{union_iff}.
        \begin{byCase}
        \caseOf{$P\in H_0$.}
            Take $N\in G$ such that $(N,P)\in m$ by \cref{img_iff}.
            $N\in\dom{m}$ by \cref{subseteq}.
            $N\mathrel{R}\apply{m}{N}$ by \cref{subseteq}.
            $P = \apply{m}{N}$ by \cref{function_apply_iff}.
            Take $N_0, P_0, x_0, B_0, L, V$ such that
                $(N,\apply{m}{N}) = (N_0,P_0)$ and
                $N_0\in G$ and
                $P_0\in\pow{F}$ and
                $x_0\in C_0$ and
                $B_0$ is the closed metric ball in $Y$ about $h(x_0)$ with radius $\lambda$ with respect to $d$ and
                $L = \preimg{\restrl{h}{C_0}}{B_0}$ and
                $V = \metricBall{d}{Y}{h(x_0)}{\eta}$ and
                $N_0 = \preimg{\restrl{h}{C_0}}{\metricBall{d}{Y}{h(x_0)}{\lambda}}$ and
                $L$ is a compact subspace of $X$ with respect to $T$ and
                $P_0$ is the set of continuous functions from $X$ to $Y$ that carry $L$ into $V$ with respect to $T$ and $U$
                by \cref{relation,pair_eq_iff}.
            $P$ is the set of continuous functions from $X$ to $Y$ that carry $L$ into $V$ with respect to $T$ and $U$
                by \cref{pair_eq_iff}.
            Show that for all $z$ if $z\in\img{h}{L}$ then $z\in V$.
            \begin{subproof}
                Follows by \cref{funs_is_function,img_of_function_elim,inter_elim_right,preimg_iff,preimg_subseteq_dom,restrl_dom,subseteq,continuous,funs_elim,function_apply_elim,restrl_iff,function_rightunique,pair_eq_iff,half_radius_closed_ball_subset_open_ball}.
            \end{subproof}
            $\img{h}{L}\subseteq V$ by \cref{subseteq}.
            Take $Q$ such that
                $Q$ is the set of continuous functions from $X$ to $Y$ with respect to $T$ and $U$ and
                $P = \{f\in Q \mid \img{f}{L}\subseteq V\}$ by \cref{compact_open_set}.
            $h\in Q$ by \cref{continuous_function_set,setext,pair_eq_iff}.
            Therefore $h\in P$ by \cref{inters_iff_forall}.
        \caseOf{$P\in\{C,C\}$.}
            $h\in P$ by \cref{upair_iff,pair_eq_iff}.
        \end{byCase}
    \end{subproof}
    $h\in M$ by \cref{inters_iff_forall}.
    Show that for all $g$ if $g\in M$ then $g\in A$.
    \begin{subproof}
        Fix $g$.
        Assume $g\in M$.
        We have $g\in F$ by \cref{open_in,topology_on,subseteq,powerset_elim,continuous_function_set}.
        $\restrl{h}{C_0}\in\funs{C_0}{Y}$ and $\restrl{g}{C_0}\in\funs{C_0}{Y}$
            by \cref{restrl_typed_function,continuous_function_set,compact_subspace,subseteq}.
        Show that for all $y\in C_0$ we have $\apply{\boundedMetric{d}}{(\apply{\restrl{h}{C_0}}{y}, \apply{\restrl{g}{C_0}}{y})} \leq \alpha$.
        \begin{subproof}
            Fix $y$.
            Assume $y\in C_0$.
            Take $N\in G$ such that $y\in N$ by \cref{covering,subseteq,unions_iff}.
            $N\in\dom{m}$ by \cref{subseteq}.
            $N\mathrel{R}\apply{m}{N}$ by \cref{subseteq}.
            Take $x, B_0, L, V$ such that
                $x\in C_0$ and
                $B_0$ is the closed metric ball in $Y$ about $h(x)$ with radius $\lambda$ with respect to $d$ and
                $L = \preimg{\restrl{h}{C_0}}{B_0}$ and
                $V = \metricBall{d}{Y}{h(x)}{\eta}$ and
                $N = \preimg{\restrl{h}{C_0}}{\metricBall{d}{Y}{h(x)}{\lambda}}$ and
                $L$ is a compact subspace of $X$ with respect to $T$ and
                $\apply{m}{N}$ is the set of continuous functions from $X$ to $Y$ that carry $L$ into $V$ with respect to $T$ and $U$
                by \cref{relation,pair_eq_iff}.
            Take $u\in\metricBall{d}{Y}{h(x)}{\lambda}$ such that $(y,u)\in\restrl{h}{C_0}$ by \cref{preimg_iff}.
            $y\in\dom{h}$ by \cref{compact_subspace,continuous,funs_elim,subseteq}.
            $(y,h(y))\in h$ by \cref{function_apply_elim,funs_is_function}.
            $(y,h(y))\in\restrl{h}{C_0}$ by \cref{restrl_iff}.
            $u = h(y)$ by \cref{function_rightunique,funs_is_function,restrl_typed_function,continuous_function_set,compact_subspace,subseteq}.
            $h(y)\in\metricBall{d}{Y}{h(x)}{\lambda}$ by \cref{pair_eq_iff}.
            $y\in L$ by \cref{metric_ball_subset_closed_metric_ball,preimg_iff,pair_eq_iff}.
            $g\in\apply{m}{N}$ by \cref{pair_eq_iff,function_apply_elim,img_iff,union_intro_left,inters_iff_forall}.
            Take $Q$ such that
                $Q$ is the set of continuous functions from $X$ to $Y$ with respect to $T$ and $U$ and
                $\apply{m}{N} = \{f\in Q \mid \img{f}{L}\subseteq V\}$ by \cref{compact_open_set}.
            $g\in Q$ by \cref{continuous_function_set,setext,pair_eq_iff}.
            $\img{g}{L}\subseteq V$ by \cref{inters_iff_forall}.
            $y\in\dom{g}\inter L$ by \cref{continuous_function_set,pair_eq_iff,funs_elim,subseteq,inter_intro}.
            $g(y)\in\img{g}{L}$ by \cref{continuous_function_set,funs_is_function,img_of_function_intro}.
            $g(y)\in V$ by \cref{subseteq}.
            $h(y)\in V$ by \cref{half_radius_closed_ball_subset_open_ball,metric_ball_subset_closed_metric_ball,preimg_iff,pair_eq_iff}.
            $h(x)\in Y$ by \cref{compact_subspace,subseteq,funs_type_apply}.
            $h(y)\in Y$ and $g(y)\in Y$ by \cref{metric_ball,subseteq}.
            $(h(y),g(y))\in Y\times Y$ and $(h(y),h(x))\in Y\times Y$ and $(h(x),g(y))\in Y\times Y$
                by \cref{times_tuple_intro}.
            $d((h(y),h(x))) < \eta$ and $d((h(x),g(y))) < \eta$ by \cref{metric_ball,metric_on,metric_symmetric,real_lt_subst_left}.
            $d((h(y),g(y))) \leq d((h(y),h(x))) + d((h(x),g(y)))$
                by \cref{metric_on,metric_triangle_inequality}.
            $d((h(y),h(x)))\in\reals$ and $d((h(x),g(y)))\in\reals$ and $d((h(y),g(y)))\in\reals$
                by \cref{metric_on,funs_type_apply}.
            $d((h(y),h(x))) + d((h(x),g(y))) < \eta + \eta$ by \cref{real_lt_add_two}.
            $d((h(y),g(y))) < \eta + \eta$ by \cref{real_leq_split,real_lt_trans,real_lt_subst_left,real_add_closed}.
            Hence $d((h(y),g(y))) < \alpha$ by \cref{real_lt_subst_right}.
            $d((h(y),g(y))) < 1$
                by \cref{metric_on,times_tuple_intro,funs_type_apply,real_one_in_naturals,real_naturals_subset_reals,subseteq,real_lt_trans}.
            Therefore $\apply{\boundedMetric{d}}{(h(y),g(y))} \leq \alpha$
                by \cref{bounded_metric_apply_below_one,real_lt_subst_left,real_lt_imp_leq}.
            $\apply{\restrl{h}{C_0}}{y} = h(y)$ and $\apply{\restrl{g}{C_0}}{y} = g(y)$
                by \cref{restrl_of_function_apply,compact_subspace,subseteq,continuous_function_set,funs_is_function,funs_is_relation,funs_elim}.
            $\apply{\boundedMetric{d}}{(\apply{\restrl{h}{C_0}}{y}, \apply{\restrl{g}{C_0}}{y})}
                = \apply{\boundedMetric{d}}{(h(y),g(y))}$ by \cref{pair_eq_iff}.
            Hence $\apply{\boundedMetric{d}}{(\apply{\restrl{h}{C_0}}{y}, \apply{\restrl{g}{C_0}}{y})} \leq \alpha$
                by \cref{pair_eq_iff}.
        \end{subproof}
        $\rho((\restrl{h}{C_0},\restrl{g}{C_0})) \leq \alpha$ by \cref{uniform_metric_leq_from_coordinate_bounds}.
        $\rho((\restrl{h}{C_0},\restrl{g}{C_0}))\in\reals$
            by \cref{uniform_metric_on_function_space,metric_on,funs_type_apply,times_tuple_intro}.
        Hence $\rho((\restrl{h}{C_0},\restrl{g}{C_0})) < \epsilon$ by
            \cref{real_leq_split,real_lt_trans,real_lt_subst_left}.
        Therefore $g\in A$ by \cref{subseteq}.
    \end{subproof}
    Follows by definition.
\end{proof}

\begin{lemma}\label{continuous_function_sets_equal}
    Assume $C_1$ is the set of continuous functions from $X$ to $Y$ with respect to $T$ and $U$.
    Assume $C_2$ is the set of continuous functions from $X$ to $Y$ with respect to $T$ and $U$.
    Then $C_1 = C_2$.
\end{lemma}

% from §46 Item 12: compact-open and compact-convergence topologies agree on continuous functions
\begin{theorem}\label{compact_open_equals_compact_convergence_on_continuous_functions}
    Assume $T$ is a topology on $X$.
    Assume $d$ is a metric on $Y$.
    Assume $C$ is the set of continuous functions from $X$ to $Y$ with respect to $T$ and $\metricTopology{d}{Y}$.
    Assume $K$ is the compact open topology on continuous functions from $X$ to $Y$
        with respect to $T$ and $\metricTopology{d}{Y}$.
    Assume $T_0$ is the topology of compact convergence on functions from $X$ to $Y$ with respect to $T$ and $d$.
    Then $K = \subspaceTopology{T_0}{C}$.
\end{theorem}
\begin{proof}
    Let $F = \funs{X}{Y}$.
    Let $U = \metricTopology{d}{Y}$.
    Take $F_0, S$ such that
        $F_0$ is the set of continuous functions from $X$ to $Y$ with respect to $T$ and $U$ and
        $S$ is the compact-open subbasis on continuous functions from $X$ to $Y$ with respect to $T$ and $U$ and
        $K = \subbasisTopology{S}{F_0}$ by \cref{compact_open_topology}.
    We have $K = \subbasisTopology{S}{C}$ by \cref{continuous_function_sets_equal,pair_eq_iff}.
    $U$ is a topology on $Y$ by \cref{metric_basis_is_basis,basis_topology_is_topology,metric_topology,basis_for_topology}.
    $S$ is a subbasis on $C$ by \cref{compact_open_subbasis_is_subbasis}.
    $\finiteIntersections{S}$ is a basis for $K$ on $C$
        by \cref{subbasis_finite_intersections_basis,basis_for_topology,topology_generated_by_subbasis}.
    $K$ is a topology on $C$ by \cref{basis_for_topology,basis_topology_is_topology}.
    Take $B$ such that
        $B$ is the compact-convergence basis on functions from $X$ to $Y$ with respect to $T$ and $d$ and
        $T_0 = \basisTopology{B}{F}$ by \cref{compact_convergence_topology}.
    Let $T_1 = \subspaceTopology{T_0}{C}$.
    $C\subseteq F$ by \cref{continuous_function_set,subseteq}.
    $T_1$ is a topology on $C$ by \cref{compact_convergence_topology,subspace_topology_is_topology}.
    Show that for all $A$ if $A\in S$ then $A\in T_1$.
    \begin{subproof}
        Fix $A$.
        Assume $A\in S$.
        $A\subseteq C$ by \cref{compact_open_subbasis_is_subbasis,subbasis_on,subseteq,pow_iff}.
        Show that for all $h$ if $h\in A$ then there exists $P$ such that
            $P$ is a neighborhood of $h$ in $C$ with respect to $T_1$ and
            $P\subseteq A$.
        \begin{subproof}
            Fix $h$.
            Assume $h\in A$.
            Take $C_0, V$ such that
                $A$ is the set of continuous functions from $X$ to $Y$ that carry $C_0$ into $V$
                    with respect to $T$ and $U$ by \cref{compact_open_function_space_subbasis}.
            Take $M$ such that
                $M$ is open in $F$ with respect to $T_0$ and
                $h\in M$ and
                for all $g$ if $g\in M\inter C$ then $g\in A$
                by \cref{compact_open_local_compact_convergence}.
            Let $N = C\inter M$.
            $N$ is a neighborhood of $h$ in $C$ with respect to $T_1$
                by \cref{subspace_topology,open_in,compact_open_subbasis_is_subbasis,subbasis_on,subseteq,pow_iff,inter_intro,neighborhood}.
            $N\subseteq A$ by \cref{inter,subseteq}.
            Follows by definition.
        \end{subproof}
        Therefore $A\in T_1$ by \cref{open_from_local_neighborhoods,open_in}.
    \end{subproof}
    $S\subseteq T_1$ by \cref{subseteq}.
    Show that for all $W$ if $W\in\finiteIntersections{S}$ then $W\in T_1$.
    \begin{subproof}
        Follows by \cref{finite_intersections,subseteq_transitive,finite_intersections_open_in_topology,pair_eq_iff}.
    \end{subproof}
    $\finiteIntersections{S}\subseteq T_1$ by \cref{subseteq}.
    Show that for all $W$ if $W\in K$ then $W\in T_1$.
    \begin{subproof}
        Follows by \cref{basis_topology_unions,subseteq_transitive,topology_on}.
    \end{subproof}
    $K\subseteq T_1$ by \cref{subseteq}.
    Show that for all $W$ if $W\in T_1$ then $W\in K$.
    \begin{subproof}
        Fix $W$.
        Assume $W\in T_1$.
        $W\subseteq C$ by \cref{subspace_topology,inter_lower_left,pow_iff,subseteq}.
        Show that for all $h$ if $h\in W$ then there exists $P$ such that
            $P$ is a neighborhood of $h$ in $C$ with respect to $K$ and
            $P\subseteq W$.
        \begin{subproof}
            Fix $h$.
            Assume $h\in W$.
            Take $V\in T_0$ such that $W = C\inter V$ by \cref{subspace_topology}.
            $h\in C$ and $h\in V$ by \cref{inter}.
            Take $A\in B$ such that $h\in A$ and $A\subseteq V$
                by \cref{open_in,pair_eq_iff,topology_generated_by_basis}.
            Take $h_0, C_0, \epsilon$ such that
                $A$ is the compact-convergence ball about $h_0$ over $C_0$ with radius $\epsilon$
                    from $X$ to $Y$ with respect to $T$ and $d$
                by \cref{compact_convergence_function_space_basis}.
            Take $\delta, D$ such that
                $\delta\in\reals$ and
                $0 < \delta$ and
                $D$ is the compact-convergence ball about $h$ over $C_0$ with radius $\delta$
                    from $X$ to $Y$ with respect to $T$ and $d$ and
                $D\subseteq A$ by \cref{compact_convergence_ball_refinement}.
            Take $N$ such that
                $N$ is open in $C$ with respect to $K$ and
                $h\in N$ and
                for all $g$ if $g\in N$ then $g\in D$
                by \cref{compact_convergence_local_compact_open}.
            $N$ is a neighborhood of $h$ in $C$ with respect to $K$ by \cref{neighborhood}.
            $N\subseteq W$
                by \cref{subseteq,open_in,topology_on,pow_iff,inter_intro,pair_eq_iff}.
            Follows by definition.
        \end{subproof}
        Hence $W\in K$ by \cref{open_from_local_neighborhoods,open_in}.
    \end{subproof}
    $T_1\subseteq K$ by \cref{subseteq}.
    Therefore $K = T_1$ by \cref{subseteq_antisymmetric}.
\end{proof}

% from §46 Item 13: compact-convergence topology on continuous functions is metric-independent
\begin{corollary}\label{compact_convergence_on_continuous_functions_metric_independent}
    Assume $T$ is a topology on $X$.
    Assume $d_1$ is a metric on $Y$.
    Assume $d_2$ is a metric on $Y$.
    Assume $\metricTopology{d_1}{Y} = \metricTopology{d_2}{Y}$.
    Assume $C$ is the set of continuous functions from $X$ to $Y$ with respect to $T$ and $\metricTopology{d_1}{Y}$.
    Assume $T_1$ is the topology of compact convergence on functions from $X$ to $Y$ with respect to $T$ and $d_1$.
    Assume $T_2$ is the topology of compact convergence on functions from $X$ to $Y$ with respect to $T$ and $d_2$.
    Then $\subspaceTopology{T_1}{C} = \subspaceTopology{T_2}{C}$.
\end{corollary}
\begin{proof}
    Let $U = \metricTopology{d_1}{Y}$.
    Let $S = \{A\in\pow{\funs{X}{Y}} \mid \text{there exist $C_0, V$ such that
        $A$ is the set of continuous functions from $X$ to $Y$ that carry $C_0$ into $V$
            with respect to $T$ and $U$}\}$.
    Let $K = \subbasisTopology{S}{C}$.
    $U$ is a topology on $Y$ by \cref{metric_basis_is_basis,basis_topology_is_topology,metric_topology,basis_for_topology}.
    $S$ is the compact-open subbasis on continuous functions from $X$ to $Y$ with respect to $T$ and $U$
        by \cref{compact_open_function_space_subbasis}.
    $K$ is the compact open topology on continuous functions from $X$ to $Y$ with respect to $T$ and $U$
        by \cref{compact_open_topology}.
    $K = \subspaceTopology{T_1}{C}$ by \cref{compact_open_equals_compact_convergence_on_continuous_functions}.
    $C$ is the set of continuous functions from $X$ to $Y$ with respect to $T$ and $\metricTopology{d_2}{Y}$
        by \cref{pair_eq_iff}.
    $S$ is the compact-open subbasis on continuous functions from $X$ to $Y$ with respect to $T$ and $\metricTopology{d_2}{Y}$
        by \cref{pair_eq_iff}.
    $K$ is the compact open topology on continuous functions from $X$ to $Y$ with respect to $T$ and $\metricTopology{d_2}{Y}$
        by \cref{compact_open_topology,pair_eq_iff}.
    $K = \subspaceTopology{T_2}{C}$ by \cref{compact_open_equals_compact_convergence_on_continuous_functions}.
    Therefore $\subspaceTopology{T_1}{C} = \subspaceTopology{T_2}{C}$ by \cref{pair_eq_iff}.
\end{proof}

% from §46 helper lemma 28: on compact domains, the uniform metric topology is the compact-convergence topology
\begin{lemma}\label{compact_domain_uniform_metric_topology_is_compact_convergence}
    Assume $T$ is a topology on $X$.
    Assume $X$ is compact with respect to $T$.
    Assume $d$ is a metric on $Y$.
    Assume $\rho$ is the uniform metric on functions from $X$ to $Y$ with respect to $d$.
    Then $\metricTopology{\rho}{\funs{X}{Y}}$ is the topology of compact convergence on functions from $X$ to $Y$ with respect to $T$ and $d$.
\end{lemma}
\begin{proof}
    Let $F = \funs{X}{Y}$.
    Let $B = \{A\in\pow{F} \mid \text{there exist $f, C, \epsilon$ such that
        $A$ is the compact-convergence ball about $f$ over $C$ with radius $\epsilon$
            from $X$ to $Y$ with respect to $T$ and $d$}\}$.
    $B$ is the compact-convergence basis on functions from $X$ to $Y$ with respect to $T$ and $d$
        by \cref{compact_convergence_function_space_basis}.
    Let $T_0 = \basisTopology{B}{F}$.
    Let $T_X = \subspaceTopology{T}{X}$.
    We have $T_X = T$ by \cref{topology_on,powerset_elim,subseteq,inter_absorb_supseteq_left,subspace_topology,subseteq_antisymmetric}.

    Show that for all $U$ if $U\in\metricTopology{\rho}{F}$ then $U\in T_0$.
    \begin{subproof}
        Fix $U$.
        Assume $U\in\metricTopology{\rho}{F}$.
        $U\subseteq F$ by \cref{uniform_metric_on_function_space,metric_topology,topology_generated_by_basis,pow_iff,subseteq}.
        Show that for all $h\in U$ there exists $A\in B$ such that $h\in A$ and $A\subseteq U$.
        \begin{subproof}
            Fix $h\in U$.
            Take $M\in\metricBasis{\rho}{F}$ such that $h\in M$ and $M\subseteq U$
                by \cref{metric_topology,metric_basis_is_basis,basis_for_topology,topology_generated_by_basis}.
            Take $f,\epsilon$ such that
                $f\in F$ and
                $\epsilon\in\reals$ and
                $0 < \epsilon$ and
                $M = \metricBall{\rho}{F}{f}{\epsilon}$ by \cref{metric_basis}.
            Let $A = \{g\in F \mid \rho((\restrl{f}{X},\restrl{g}{X})) < \epsilon\}$.
            We have $M = A$
                by \cref{metric_ball,restrl_self_typed_function,real_lt_subst_left,real_lt_subst_right,inters_iff_forall,subseteq,subseteq_antisymmetric}.
            $M$ is the compact-convergence ball about $f$ over $X$ with radius $\epsilon$
                from $X$ to $Y$ with respect to $T$ and $d$
                by \cref{subseteq_refl,compact_subspace,compact_convergence_ball,pair_eq_iff}.
            $M\in B$ by \cref{metric_ball,subseteq,compact_convergence_function_space_basis,powerset_intro}.
            Follows by definition.
        \end{subproof}
        Therefore $U\in T_0$ by \cref{powerset_intro,topology_generated_by_basis}.
    \end{subproof}
    $\metricTopology{\rho}{F}\subseteq T_0$ by \cref{subseteq}.

    Show that for all $U$ if $U\in T_0$ then $U\in\metricTopology{\rho}{F}$.
    \begin{subproof}
        Fix $U$.
        Assume $U\in T_0$.
        $U\subseteq F$ by \cref{topology_generated_by_basis,pow_iff,subseteq}.
        $\metricTopology{\rho}{F}$ is a topology on $F$
            by \cref{uniform_metric_on_function_space,metric_basis_is_basis,basis_topology_is_topology,metric_topology,basis_for_topology}.
        Show that for all $h\in U$ there exists $N$ such that
            $N$ is a neighborhood of $h$ in $F$ with respect to $\metricTopology{\rho}{F}$ and
            $N\subseteq U$.
        \begin{subproof}
            Fix $h\in U$.
            $h\in F$ by \cref{subseteq}.
            Take $A\in B$ such that $h\in A$ and $A\subseteq U$ by \cref{topology_generated_by_basis}.
            Take $f, C, \epsilon$ such that
                $A$ is the compact-convergence ball about $f$ over $C$ with radius $\epsilon$
                    from $X$ to $Y$ with respect to $T$ and $d$
                by \cref{compact_convergence_function_space_basis}.
            Take $D,\delta$ such that
                $\delta\in\reals$ and
                $0 < \delta$ and
                $D$ is the compact-convergence ball about $h$ over $C$ with radius $\delta$
                    from $X$ to $Y$ with respect to $T$ and $d$ and
                $D\subseteq A$ by \cref{compact_convergence_ball_refinement}.
            Take $\eta$ such that
                $\eta\in\reals$ and
                $0 < \eta$ and
                $\eta < 1$ and
                $\eta < \delta$ by \cref{positive_real_has_smaller_positive_below_one}.
            Let $M = \metricBall{\rho}{F}{h}{\eta}$.
            Take $\sigma$ such that
                $\sigma$ is the uniform metric on functions from $C$ to $Y$ with respect to $d$ and
                $D = \{g\in F \mid \sigma((\restrl{h}{C},\restrl{g}{C})) < \delta\}$ by \cref{compact_convergence_ball}.
            Show that for all $g$ if $g\in M$ then $g\in D$.
            \begin{subproof}
                Fix $g$.
                Assume $g\in M$.
                We have $g\in F$ and $\rho((h,g)) < \eta$ by \cref{metric_ball}.
                $\restrl{h}{C}\in\funs{C}{Y}$ and $\restrl{g}{C}\in\funs{C}{Y}$
                    by \cref{compact_convergence_ball,compact_subspace,subseteq,restrl_typed_function}.
                Show that for all $x\in C$ we have
                    $\apply{\boundedMetric{d}}{(\apply{\restrl{h}{C}}{x}, \apply{\restrl{g}{C}}{x})} \leq \eta$.
                \begin{subproof}
                    Follows by \cref{compact_convergence_ball,compact_subspace,subseteq,uniform_metric_small_implies_pointwise_small,times_tuple_intro,funs_type_apply,metric_on,real_one_in_naturals,real_naturals_subset_reals,real_lt_trans,bounded_metric_apply_below_one,real_lt_subst_left,real_lt_imp_leq,restrl_of_function_apply,funs_is_function,funs_is_relation,funs_elim,pair_eq_iff}.
                \end{subproof}
                $\sigma((\restrl{h}{C},\restrl{g}{C})) \leq \eta$ by \cref{uniform_metric_leq_from_coordinate_bounds}.
                    $\sigma((\restrl{h}{C},\restrl{g}{C}))\in\reals$
                        by \cref{uniform_metric_on_function_space,metric_on,funs_type_apply,times_tuple_intro}.
                Hence $\sigma((\restrl{h}{C},\restrl{g}{C})) < \delta$
                    by \cref{real_leq_split,real_lt_trans,real_lt_subst_left}.
                Therefore $g\in D$ by \cref{subseteq}.
            \end{subproof}
            $M\subseteq A$ by \cref{subseteq,subseteq_transitive}.
            $M\in\metricBasis{\rho}{F}$ by \cref{metric_ball,subseteq,metric_basis,powerset_intro}.
            $M$ is a neighborhood of $h$ in $F$ with respect to $\metricTopology{\rho}{F}$
                by \cref{metric_topology,metric_basis_is_basis,basis_elem_open_in_generated,open_in,metric_ball,metric_on,uniform_metric_on_function_space,metric_zero_characterization,times_tuple_intro,real_lt_subst_left,neighborhood}.
            Hence there exists $N$ such that
                $N$ is a neighborhood of $h$ in $F$ with respect to $\metricTopology{\rho}{F}$ and
                $N\subseteq U$ by \cref{subseteq,pair_eq_iff}.
        \end{subproof}
        Therefore $U\in\metricTopology{\rho}{F}$ by \cref{open_from_local_neighborhoods,open_in}.
    \end{subproof}
    $T_0\subseteq\metricTopology{\rho}{F}$ by \cref{subseteq}.

    Thus $\metricTopology{\rho}{F} = T_0$ by \cref{subseteq_antisymmetric}.
    $\metricTopology{\rho}{F}$ is a topology on $F$
        by \cref{metric_basis_is_basis,basis_topology_is_topology,metric_topology,basis_for_topology,uniform_metric_on_function_space}.
    Hence $\metricTopology{\rho}{F}$ is the topology of compact convergence on functions from $X$ to $Y$ with respect to $T$ and $d$
        by \cref{compact_convergence_topology}.
\end{proof}

% from §46 Item 13: uniform topology on continuous functions is metric-independent for compact domains
\begin{corollary}\label{uniform_topology_on_continuous_functions_metric_independent_for_compact_domain}
    Assume $T$ is a topology on $X$.
    Assume $X$ is compact with respect to $T$.
    Assume $d_1$ is a metric on $Y$.
    Assume $d_2$ is a metric on $Y$.
    Assume $\metricTopology{d_1}{Y} = \metricTopology{d_2}{Y}$.
    Assume $\rho_1$ is the uniform metric on functions from $X$ to $Y$ with respect to $d_1$.
    Assume $\rho_2$ is the uniform metric on functions from $X$ to $Y$ with respect to $d_2$.
    Assume $C$ is the set of continuous functions from $X$ to $Y$ with respect to $T$ and $\metricTopology{d_1}{Y}$.
    Then $\subspaceTopology{\metricTopology{\rho_1}{\funs{X}{Y}}}{C}
        = \subspaceTopology{\metricTopology{\rho_2}{\funs{X}{Y}}}{C}$.
\end{corollary}
\begin{proof}
    Follows by \cref{pair_eq_iff,compact_domain_uniform_metric_topology_is_compact_convergence,compact_convergence_on_continuous_functions_metric_independent}.
\end{proof}

% from §46 helper lemma 29: compact-domain uniform subspaces agree with compact-convergence subspaces
\begin{lemma}\label{compact_domain_uniform_subspace_equals_compact_convergence_subspace}
    Assume $T$ is a topology on $X$.
    Assume $X$ is compact with respect to $T$.
    Assume $d$ is a metric on $Y$.
    Assume $\rho$ is the uniform metric on functions from $X$ to $Y$ with respect to $d$.
    Assume $C$ is the set of continuous functions from $X$ to $Y$ with respect to $T$ and $\metricTopology{d}{Y}$.
    Assume $T_0$ is the topology of compact convergence on functions from $X$ to $Y$ with respect to $T$ and $d$.
    Then $\subspaceTopology{\metricTopology{\rho}{\funs{X}{Y}}}{C} = \subspaceTopology{T_0}{C}$.
\end{lemma}

% from §46 helper lemma 30: compact-domain uniform subspaces are metric-independent once compact-convergence topologies are given
\begin{lemma}\label{compact_domain_uniform_subspace_metric_independence_from_compact_convergence}
    Assume $T$ is a topology on $X$.
    Assume $X$ is compact with respect to $T$.
    Assume $d_1$ is a metric on $Y$.
    Assume $d_2$ is a metric on $Y$.
    Assume $\metricTopology{d_1}{Y} = \metricTopology{d_2}{Y}$.
    Assume $\rho_1$ is the uniform metric on functions from $X$ to $Y$ with respect to $d_1$.
    Assume $\rho_2$ is the uniform metric on functions from $X$ to $Y$ with respect to $d_2$.
    Assume $C$ is the set of continuous functions from $X$ to $Y$ with respect to $T$ and $\metricTopology{d_1}{Y}$.
    Assume $T_1$ is the topology of compact convergence on functions from $X$ to $Y$ with respect to $T$ and $d_1$.
    Assume $T_2$ is the topology of compact convergence on functions from $X$ to $Y$ with respect to $T$ and $d_2$.
    Then $\subspaceTopology{\metricTopology{\rho_1}{\funs{X}{Y}}}{C}
        = \subspaceTopology{\metricTopology{\rho_2}{\funs{X}{Y}}}{C}$.
\end{lemma}

% from §46 Item 14: evaluation map
\begin{definition}\label{evaluation_map}
    $e$ is the evaluation map from $X$ and $C$ to $Y$ iff
        $C\subseteq\funs{X}{Y}$ and
        $e\in\funs{X\times C}{Y}$ and
        for all $x, f$ such that $x\in X$ and $f\in C$ we have
            $\apply{e}{(x,f)} = \apply{f}{x}$.
\end{definition}

% from §46 Item 14: continuity of the evaluation map
\begin{theorem}\label{evaluation_map_continuous_compact_open}
    Assume $T$ is a topology on $X$.
    Assume $U$ is a topology on $Y$.
    Assume $X$ is locally compact with respect to $T$.
    Assume $X$ is Hausdorff with respect to $T$.
    Assume $C$ is the set of continuous functions from $X$ to $Y$ with respect to $T$ and $U$.
    Assume $K$ is the compact open topology on continuous functions from $X$ to $Y$ with respect to $T$ and $U$.
    Assume $e$ is the evaluation map from $X$ and $C$ to $Y$.
    Then $e$ is continuous from $X\times C$ to $Y$ with respect to $\productTopology{T}{K}{X}{C}$ and $U$.
\end{theorem}
\begin{proof}
    Let $P = \productTopology{T}{K}{X}{C}$.
    $e\in\funs{X\times C}{Y}$ by \cref{evaluation_map}.
    Take $F_0, S$ such that
        $F_0$ is the set of continuous functions from $X$ to $Y$ with respect to $T$ and $U$ and
        $S$ is the compact-open subbasis on continuous functions from $X$ to $Y$ with respect to $T$ and $U$ and
        $K = \subbasisTopology{S}{F_0}$ by \cref{compact_open_topology}.
    We have $K = \subbasisTopology{S}{C}$ by \cref{continuous_function_sets_equal,pair_eq_iff}.
    $S$ is a subbasis on $C$ by \cref{compact_open_subbasis_is_subbasis}.
    $K$ is a topology on $C$
        by \cref{subbasis_finite_intersections_basis,basis_for_topology,topology_generated_by_subbasis,basis_topology_is_topology}.
    $P$ is a topology on $X\times C$
        by \cref{product_basis_is_basis,basis_topology_is_topology,product_topology}.
    Show that for all $p\in X\times C$ we have
        $e$ is continuous at $p$ from $X\times C$ to $Y$ with respect to $P$ and $U$.
    \begin{subproof}
        Fix $p\in X\times C$.
        Take $x, g$ such that $p = (x,g)$ and $x\in X$ and $g\in C$ by \cref{times_elem_is_tuple}.
        $g\in\funs{X}{Y}$ and $g$ is continuous from $X$ to $Y$ with respect to $T$ and $U$
            by \cref{continuous_function_set}.
        $\apply{e}{p} = \apply{g}{x}$ by \cref{evaluation_map,pair_eq_iff}.
        Show that for all $V$ such that $V$ is a neighborhood of $\apply{e}{p}$ in $Y$ with respect to $U$
            there exists $N$ such that
                $N$ is a neighborhood of $p$ in $X\times C$ with respect to $P$ and
                $\img{e}{N}\subseteq V$.
        \begin{subproof}
            Fix $V$ such that $V$ is a neighborhood of $\apply{e}{p}$ in $Y$ with respect to $U$.
            $V$ is open in $Y$ with respect to $U$ and $\apply{g}{x}\in V$
                by \cref{neighborhood,pair_eq_iff}.
            Let $W = \preimg{g}{V}$.
            $W$ is a neighborhood of $x$ in $X$ with respect to $T$
                by \cref{continuous,pair_eq_iff,funs,subseteq,function_apply_elim,funs_is_function,preimg_iff,neighborhood}.
            Take $U_0$ such that
                $U_0$ is a neighborhood of $x$ in $X$ with respect to $T$ and
                $\closure{U_0}{X}{T}$ is a compact subspace of $X$ with respect to $T$ and
                $\closure{U_0}{X}{T}\subseteq W$
                by \cref{locally_compact_neighborhood_characterization}.
            Let $D = \closure{U_0}{X}{T}$.
            $D$ is a compact subspace of $X$ with respect to $T$ and $D\subseteq W$ by \cref{subseteq}.
            Let $A = \{f\in C \mid \img{f}{D}\subseteq V\}$.
            $A$ is the set of continuous functions from $X$ to $Y$ that carry $D$ into $V$
                with respect to $T$ and $U$ by \cref{compact_open_set}.
            $A\in S$ by \cref{subseteq,continuous_function_set,subseteq_transitive,powerset_intro,compact_open_function_space_subbasis}.
            $A$ is open in $C$ with respect to $K$
                by \cref{compact_open_subbasis_element_open}.
            Show that for all $y$ if $y\in\img{g}{D}$ then $y\in V$.
            \begin{subproof}
                Fix $y$.
                Assume $y\in\img{g}{D}$.
                Take $x_0\in\dom{g}\inter D$ such that $y = g(x_0)$ by \cref{funs_is_function,img_of_function_elim}.
                Therefore $y\in V$ by \cref{inter_elim_right,subseteq,preimg_iff,function_apply_iff,funs_is_function,pair_eq_iff}.
            \end{subproof}
            $\img{g}{D}\subseteq V$ by \cref{subseteq}.
            $g\in A$ by \cref{inters_iff_forall}.
            Let $N = U_0\times A$.
            $U_0$ is open in $X$ with respect to $T$ and $x\in U_0$ by \cref{neighborhood}.
            $N$ is a neighborhood of $p$ in $X\times C$ with respect to $P$ by
                \cref{open_in,topology_on,subseteq,pow_iff,times_subseteq_left,times_subseteq_right,product_basis,product_basis_is_basis,product_topology,basis_elem_open_in_generated,times_tuple_intro,pair_eq_iff,neighborhood}.
            Show that for all $y$ if $y\in\img{e}{N}$ then $y\in V$.
            \begin{subproof}
                Fix $y$.
                Assume $y\in\img{e}{N}$.
                Take $q\in\dom{e}\inter N$ such that $y = e(q)$ by \cref{evaluation_map,funs_is_function,img_of_function_elim}.
                Take $x_0, h$ such that $q = (x_0,h)$ and $x_0\in U_0$ and $h\in A$ by \cref{inter_elim_right,times_elem_is_tuple}.
                $x_0\in D$ by \cref{neighborhood,open_in,topology_on,subseteq,pow_iff,subseteq_closure}.
                We have $h\in C$ and $\img{h}{D}\subseteq V$ by \cref{inters_iff_forall}.
                $h\in\funs{X}{Y}$ by \cref{continuous_function_set,pair_eq_iff}.
                $x_0\in\dom{h}\inter D$
                    by \cref{funs,inter_intro,compact_subspace,subseteq}.
                $h(x_0)\in\img{h}{D}$ by \cref{funs_is_function,img_of_function_intro}.
                $h(x_0)\in V$ by \cref{subseteq}.
                Therefore $y\in V$ by \cref{evaluation_map,function_apply_elim,compact_subspace,subseteq,pair_eq_iff}.
            \end{subproof}
            $\img{e}{N}\subseteq V$ by \cref{subseteq}.
            Follows by definition.
        \end{subproof}
        Hence $e$ is continuous at $p$ from $X\times C$ to $Y$ with respect to $P$ and $U$
            by \cref{continuous_at}.
    \end{subproof}
    Therefore $e$ is continuous from $X\times C$ to $Y$ with respect to $P$ and $U$
        by \cref{continuous_at,continuous_at_implies_continuous}.
\end{proof}

% from §46 Item 15: induced map
\begin{definition}\label{induced_map_into_function_space}
    $F$ is the map induced by $f$ from $Z$ to $C$ with respect to $X$ and $Y$ iff
        $f\in\funs{X\times Z}{Y}$ and
        $C\subseteq\funs{X}{Y}$ and
        $F\in\funs{Z}{C}$ and
        for all $z, x$ such that $z\in Z$ and $x\in X$ we have
            $\apply{\apply{F}{z}}{x} = \apply{f}{(x,z)}$.
\end{definition}

% from §46 helper lemma 30: induced maps are locally inside compact-open subbasic sets
\begin{lemma}\label{induced_map_compact_open_local_subbasic}
    Assume $T$ is a topology on $X$.
    Assume $S$ is a topology on $Z$.
    Assume $U$ is a topology on $Y$.
    Assume $C$ is the set of continuous functions from $X$ to $Y$ with respect to $T$ and $U$.
    Assume $f$ is continuous from $X\times Z$ to $Y$ with respect to $\productTopology{T}{S}{X}{Z}$ and $U$.
    Assume $F$ is the map induced by $f$ from $Z$ to $C$ with respect to $X$ and $Y$.
    Assume $A$ is the set of continuous functions from $X$ to $Y$ that carry $C_0$ into $V$ with respect to $T$ and $U$.
    Assume $z_0\in Z$.
    Suppose $\apply{F}{z_0}\in A$.
    Then there exists $W$ such that
        $W$ is open in $Z$ with respect to $S$ and
        $z_0\in W$ and
        $W\subseteq\preimg{F}{A}$.
\end{lemma}
\begin{proof}
    Take $Q$ such that
        $Q$ is the set of continuous functions from $X$ to $Y$ with respect to $T$ and $U$ and
        $A = \{g\in Q \mid \img{g}{C_0}\subseteq V\}$ by \cref{compact_open_set}.
    $Q = C$ by \cref{continuous_function_sets_equal}.
    $C_0$ is a compact subspace of $X$ with respect to $T$ and
        $V$ is open in $Y$ with respect to $U$ by \cref{compact_open_set}.
    $C_0\subseteq X$ by \cref{compact_subspace}.
    $\apply{F}{z_0}\in Q$ and $\img{\apply{F}{z_0}}{C_0}\subseteq V$ by \cref{inters_iff_forall}.
    $\apply{F}{z_0}\in\funs{X}{Y}$ by \cref{pair_eq_iff,continuous_function_set}.

    Let $T_C = \subspaceTopology{T}{C_0}$.
    $T_C$ is a topology on $C_0$ and $C_0$ is compact with respect to $T_C$
        by \cref{subspace_topology_is_topology,compact_subspace}.
    $C_0$ is a subspace of $X$ with respect to $T$ by \cref{subspace_of,compact_subspace}.
    Let $j = \identity{C_0}$.
    $j$ is continuous from $C_0$ to $X$ with respect to $T_C$ and $T$
        by \cref{inclusion_continuous}.

    Let $P = \productTopology{S}{T_C}{Z}{C_0}$.
    Let $g_1 = j\circ \piTwo{Z}{C_0}$.
    $\piTwo{Z}{C_0}$ is continuous from $Z\times C_0$ to $C_0$ with respect to $P$ and $T_C$
        by \cref{projection_two_continuous}.
    $g_1$ is continuous from $Z\times C_0$ to $X$ with respect to $P$ and $T$
        by \cref{continuous_composite}.
    Let $g_2 = \piOne{Z}{C_0}$.
    $g_2$ is continuous from $Z\times C_0$ to $Z$ with respect to $P$ and $S$
        by \cref{projection_one_continuous}.
    Let $G = \{(p,(g_1(p),g_2(p))) \mid p\in Z\times C_0\}$.
    $G$ is continuous from $Z\times C_0$ to $X\times Z$ with respect to $P$ and
        $\productTopology{T}{S}{X}{Z}$ by \cref{continuous_map_into_product_backward}.
    Let $h = f\circ G$.
    $h$ is continuous from $Z\times C_0$ to $Y$ with respect to $P$ and $U$
        by \cref{continuous_composite}.

    Let $N = \preimg{h}{V}$.
    $N$ is open in $Z\times C_0$ with respect to $P$ by \cref{continuous}.
    Show that for all $p$ if $p\in\{z_0\}\times C_0$ then $p\in N$.
    \begin{subproof}
        Fix $p$.
        Assume $p\in\{z_0\}\times C_0$.
        Take $z, x$ such that $p = (z,x)$ and $z\in\{z_0\}$ and $x\in C_0$ by \cref{times_elem_is_tuple}.
        $z = z_0$ by \cref{singleton_iff}.
        $p\in Z\times C_0$ by \cref{times_tuple_intro,subseteq}.
        $\apply{g_2}{p} = z$ by \cref{projection_first,continuous,funs_is_function,function_apply_intro}.
        $\apply{\piTwo{Z}{C_0}}{p} = x$
            by \cref{projection_second,projection_two_continuous,continuous,funs_is_function,function_apply_intro}.
        $\apply{j}{x} = x$ by \cref{id_apply,id_elem_funs}.
        $\piTwo{Z}{C_0}$ is a function and $j$ is a function
            by \cref{projection_two_continuous,continuous,funs_is_function}.
        $\ran{\piTwo{Z}{C_0}}\subseteq C_0$ by \cref{projection_two_continuous,continuous,funs_ran,subseteq}.
        $\ran{\piTwo{Z}{C_0}}\subseteq \dom{j}$ by \cref{id_dom,pair_eq_iff,subseteq}.
        $p\in\dom{\piTwo{Z}{C_0}}$ by \cref{projection_two_continuous,continuous,funs}.
        $\apply{g_1}{p} = \apply{j}{\apply{\piTwo{Z}{C_0}}{p}}$
            by \cref{circ_apply}.
        $\apply{g_1}{p} = x$ by \cref{pair_eq_iff}.
        $\apply{G}{p} = (\apply{g_1}{p},\apply{g_2}{p})$ by \cref{continuous,funs_is_function,function_apply_intro}.
        $\apply{G}{p} = (x,z_0)$ by \cref{pair_eq_iff}.
        $G$ is a function and $f$ is a function by \cref{continuous,funs_is_function}.
        $\ran{G}\subseteq X\times Z$ by \cref{continuous,funs_ran,subseteq}.
        $\ran{G}\subseteq \dom{f}$ by \cref{continuous,funs,pair_eq_iff,subseteq}.
        $p\in\dom{G}$ by \cref{continuous,funs}.
        $\apply{h}{p} = \apply{f}{\apply{G}{p}}$ by \cref{circ_apply}.
        $\apply{h}{p} = \apply{f}{(x,z_0)}$ by \cref{pair_eq_iff}.
        $x\in\dom{\apply{F}{z_0}}$ by \cref{funs,subseteq}.
        $x\in\dom{\apply{F}{z_0}}\inter C_0$ by \cref{inter_intro}.
        $\apply{\apply{F}{z_0}}{x}\in\img{\apply{F}{z_0}}{C_0}$
            by \cref{funs_is_function,img_of_function_intro,continuous_function_set,pair_eq_iff}.
        $\apply{\apply{F}{z_0}}{x}\in V$ by \cref{subseteq}.
        $\apply{\apply{F}{z_0}}{x} = \apply{f}{(x,z_0)}$ by \cref{induced_map_into_function_space,subseteq}.
        $h(p)\in V$ by \cref{pair_eq_iff}.
        $p\in\dom{h}$ by \cref{continuous,funs}.
        $(p,\apply{h}{p})\in h$ by \cref{function_apply_elim,continuous,funs_is_function}.
        $p\in N$ by \cref{preimg_iff,pair_eq_iff}.
    \end{subproof}
    $\{z_0\}\times C_0\subseteq N$ by \cref{subseteq}.
    Take $W$ such that
        $W$ is open in $Z$ with respect to $S$ and
        $z_0\in W$ and
        $W\times C_0\subseteq N$ by \cref{tube_lemma}.

    Show that for all $z$ if $z\in W$ then $z\in\preimg{F}{A}$.
    \begin{subproof}
        Fix $z$.
        Assume $z\in W$.
        $z\in Z$ by \cref{open_in,topology_on,subseteq,pow_iff}.
        $\apply{F}{z}\in C$ by \cref{induced_map_into_function_space,funs_type_apply}.
        $\apply{F}{z}\in\funs{X}{Y}$ by \cref{continuous_function_set,subseteq}.
        Show that for all $y$ if $y\in\img{\apply{F}{z}}{C_0}$ then $y\in V$.
        \begin{subproof}
            Fix $y$.
            Assume $y\in\img{\apply{F}{z}}{C_0}$.
            Take $x\in\dom{\apply{F}{z}}\inter C_0$ such that $y = \apply{\apply{F}{z}}{x}$
                by \cref{funs_is_function,img_of_function_elim,induced_map_into_function_space,funs_type_apply}.
            $x\in C_0$ by \cref{inter_elim_right}.
            $(z,x)\in W\times C_0$ by \cref{times_tuple_intro}.
            $(z,x)\in N$ by \cref{subseteq}.
            Let $p = (z,x)$.
            $p\in Z\times C_0$ by \cref{times_tuple_intro}.
            $\apply{g_2}{p} = z$ by \cref{projection_first,continuous,funs_is_function,function_apply_intro,pair_eq_iff}.
            $\apply{\piTwo{Z}{C_0}}{p} = x$
                by \cref{projection_second,projection_two_continuous,continuous,funs_is_function,function_apply_intro,pair_eq_iff}.
            $\apply{j}{x} = x$ by \cref{id_apply,id_elem_funs}.
            $\piTwo{Z}{C_0}$ is a function and $j$ is a function
                by \cref{projection_two_continuous,continuous,funs_is_function}.
            $\ran{\piTwo{Z}{C_0}}\subseteq C_0$ by \cref{projection_two_continuous,continuous,funs_ran,subseteq}.
            $\ran{\piTwo{Z}{C_0}}\subseteq \dom{j}$ by \cref{id_dom,pair_eq_iff,subseteq}.
            $p\in\dom{\piTwo{Z}{C_0}}$ by \cref{projection_two_continuous,continuous,funs}.
            $\apply{g_1}{p} = \apply{j}{\apply{\piTwo{Z}{C_0}}{p}}$
                by \cref{circ_apply}.
            $\apply{g_1}{p} = x$ by \cref{pair_eq_iff}.
            $\apply{G}{p} = (\apply{g_1}{p},\apply{g_2}{p})$ by \cref{continuous,funs_is_function,function_apply_intro}.
            $\apply{G}{p} = (x,z)$ by \cref{pair_eq_iff}.
            $G$ is a function and $f$ is a function by \cref{continuous,funs_is_function}.
            $\ran{G}\subseteq X\times Z$ by \cref{continuous,funs_ran,subseteq}.
            $\ran{G}\subseteq \dom{f}$ by \cref{continuous,funs,pair_eq_iff,subseteq}.
            $p\in\dom{G}$ by \cref{continuous,funs}.
            $\apply{h}{p} = \apply{f}{\apply{G}{p}}$ by \cref{circ_apply}.
            $\apply{h}{p} = \apply{f}{(x,z)}$ by \cref{pair_eq_iff}.
            $p\in\dom{h}$ by \cref{continuous,funs}.
            Take $u\in V$ such that $(p,u)\in h$ by \cref{preimg_iff,pair_eq_iff}.
            $h$ is a function by \cref{continuous,funs_is_function}.
            $\apply{h}{p} = u$ by \cref{function_apply_iff}.
            $\apply{f}{(x,z)} = u$ by \cref{pair_eq_iff}.
            $\apply{\apply{F}{z}}{x} = \apply{f}{(x,z)}$ by \cref{induced_map_into_function_space,subseteq}.
            $y\in V$ by \cref{pair_eq_iff}.
        \end{subproof}
        $\img{\apply{F}{z}}{C_0}\subseteq V$ by \cref{subseteq}.
        $\apply{F}{z}\in A$ by \cref{pair_eq_iff,inters_iff_forall}.
        $(z,\apply{F}{z})\in F$
            by \cref{induced_map_into_function_space,funs_is_function,funs,function_apply_elim}.
        $z\in\preimg{F}{A}$ by \cref{preimg_iff}.
    \end{subproof}
    $W\subseteq\preimg{F}{A}$ by \cref{subseteq}.
    Follows by definition.
\end{proof}

% from §46 Item 15: continuity of an induced map into the compact-open function space
\begin{theorem}\label{induced_map_continuous_compact_open}
    Assume $T$ is a topology on $X$.
    Assume $S$ is a topology on $Z$.
    Assume $U$ is a topology on $Y$.
    Assume $C$ is the set of continuous functions from $X$ to $Y$ with respect to $T$ and $U$.
    Assume $K$ is the compact open topology on continuous functions from $X$ to $Y$ with respect to $T$ and $U$.
    Assume $f$ is continuous from $X\times Z$ to $Y$ with respect to $\productTopology{T}{S}{X}{Z}$ and $U$.
    Assume $F$ is the map induced by $f$ from $Z$ to $C$ with respect to $X$ and $Y$.
    Then $F$ is continuous from $Z$ to $C$ with respect to $S$ and $K$.
\end{theorem}
\begin{proof}
    We have $F$ is a function from $Z$ to $C$ by \cref{induced_map_into_function_space,funs}.
    Take $F_0, S_0$ such that
        $F_0$ is the set of continuous functions from $X$ to $Y$ with respect to $T$ and $U$ and
        $S_0$ is the compact-open subbasis on continuous functions from $X$ to $Y$ with respect to $T$ and $U$ and
        $K = \subbasisTopology{S_0}{F_0}$ by \cref{compact_open_topology}.
    $K = \subbasisTopology{S_0}{C}$ by \cref{continuous_function_sets_equal,pair_eq_iff}.
    $S_0$ is a subbasis on $C$ by \cref{compact_open_subbasis_is_subbasis}.
    $\finiteIntersections{S_0}$ is a basis for $K$ on $C$
        by \cref{subbasis_finite_intersections_basis,basis_for_topology,topology_generated_by_subbasis}.
    $K$ is a topology on $C$ by \cref{basis_for_topology,basis_topology_is_topology}.

    Show that for all $W$ such that $W$ is open in $C$ with respect to $K$
        we have $\preimg{F}{W}$ is open in $Z$ with respect to $S$.
    \begin{subproof}
        Fix $W$ such that $W$ is open in $C$ with respect to $K$.
        $\preimg{F}{W}\subseteq Z$ by \cref{preimg_subseteq_dom,induced_map_into_function_space,funs,subseteq}.
        Show that for all $z\in\preimg{F}{W}$ there exists $P$ such that
            $P$ is a neighborhood of $z$ in $Z$ with respect to $S$ and
            $P\subseteq\preimg{F}{W}$.
        \begin{subproof}
            Fix $z$ such that $z\in\preimg{F}{W}$.
            $\apply{F}{z}\in W$
                by \cref{preimg_subseteq_dom,subseteq,preimg_iff,funs_is_function,function_apply_iff,pair_eq_iff}.
            Take $B$ such that $B\subseteq\finiteIntersections{S_0}$ and $W = \unions{B}$
                by \cref{basis_topology_unions,basis_for_topology,open_in}.
            Take $M\in B$ such that $\apply{F}{z}\in M$ by \cref{unions_iff,pair_eq_iff}.
            $M\subseteq W$ by \cref{subseteq,unions_iff,pair_eq_iff}.
            Take $G$ such that $G\subseteq S_0$ and $G$ is finite and inhabited and $M = \inters{G}$
                by \cref{subseteq,finite_intersections}.

            Let $R = \{w\in G\times S \mid \text{there exist $A, P$ such that
                $w = (A,P)$ and
                $z\in P$ and
                $P\subseteq\preimg{F}{A}$}\}$.
            $R$ is a relation by \cref{relation,subseteq,times_elem_is_tuple}.
            For all $A\in G$ there exists $P$ such that $A\mathrel{R} P$ by
                \cref{subseteq,compact_open_function_space_subbasis,inters_iff_forall,induced_map_compact_open_local_subbasic,open_in,times_tuple_intro,pair_eq_iff}.
            Take $m$ such that $m$ is a function on $G$ and for all $A\in G$ we have $A\mathrel{R}\apply{m}{A}$
                by \cref{choice_function}.
            Let $H_0 = \img{m}{G}$.
            Let $H = H_0\union\{Z,Z\}$.
            Let $P = \inters{H}$.

            For all $Q$ if $Q\in H_0$ then $Q\in S$
                by \cref{img_iff,subseteq,function_apply_iff,relation,times_tuple_elim}.
            $H_0\subseteq S$ by \cref{subseteq}.
            For all $Q$ if $Q\in H$ then $Q\in S$
                by \cref{union_iff,subseteq,upair_iff,pair_eq_iff,topology_on}.
            $H\subseteq S$ by \cref{subseteq}.
            $H$ is finite by \cref{pair_eq_iff,finite_image_function,pair_is_finite,union_of_finite_is_finite}.
            $Z\in H$ by \cref{upair_intro_left,union_intro_right}.
            $P$ is open in $Z$ with respect to $S$ by \cref{finite_intersections_open_in_topology,open_in}.

            For all $Q$ if $Q\in H$ then $z\in Q$
                by \cref{union_iff,img_iff,subseteq,function_apply_iff,relation,pair_eq_iff,upair_iff,preimg_subseteq_dom}.
            $P$ is a neighborhood of $z$ in $Z$ with respect to $S$ by \cref{inters_iff_forall,pair_eq_iff,neighborhood}.

            Let $E = \preimg{F}{W}$.
            Show that for all $y$ such that $y\in P$ we have $y\in E$.
            \begin{subproof}
                Fix $y$ such that $y\in P$.
                $y\in Z$ by \cref{open_in,topology_on,subseteq,pow_iff}.
                For all $A$ such that $A\in G$ we have $\apply{F}{y}\in A$ by
                    \cref{subseteq,relation,function_apply_elim,img_iff,union_intro_left,inters_iff_forall,pair_eq_iff,preimg_subseteq_dom,funs_is_function,preimg_iff,function_apply_iff}.
                $\apply{F}{y}\in M$ by \cref{inters_iff_forall,pair_eq_iff}.
                $\apply{F}{y}\in W$ by \cref{pair_eq_iff,subseteq}.
                $(y,\apply{F}{y})\in F$
                    by \cref{induced_map_into_function_space,funs,subseteq,function_apply_elim,funs_is_function}.
                $y\in E$ by \cref{preimg_iff,pair_eq_iff}.
            \end{subproof}
            There exists $Q$ such that
                $Q$ is a neighborhood of $z$ in $Z$ with respect to $S$ and
                $Q\subseteq\preimg{F}{W}$ by \cref{subseteq,pair_eq_iff}.
        \end{subproof}
        $\preimg{F}{W}$ is open in $Z$ with respect to $S$ by \cref{open_from_local_neighborhoods}.
    \end{subproof}
    $F$ is continuous from $Z$ to $C$ with respect to $S$ and $K$
        by \cref{open_from_local_neighborhoods,continuous}.
\end{proof}

% from §46 helper lemma 31: evaluation maps exist
\begin{lemma}\label{evaluation_map_exists}
    Assume $C$ is the set of continuous functions from $X$ to $Y$ with respect to $T$ and $U$.
    Then there exists $e$ such that $e$ is the evaluation map from $X$ and $C$ to $Y$.
\end{lemma}
\begin{proof}
    We have $C\subseteq\funs{X}{Y}$ by \cref{continuous_function_set,subseteq}.
    Let $e = \{((x,h),h(x)) \mid x\in X, h\in C\}$.
    $e$ is a relation by \cref{relation,pair_eq_iff}.
    $e$ is a function by \cref{pair_eq_iff,rightunique,relation}.
    Show that for all $p$ if $p\in\dom{e}$ then $p\in X\times C$.
    \begin{subproof}
        Fix $p$.
        Assume $p\in\dom{e}$.
        Take $x,h,y$ such that $x\in X$ and $h\in C$ and $((x,h),h(x)) = (p,y)$ by \cref{dom_iff,pair_eq_iff}.
        $p\in X\times C$ by \cref{times_tuple_intro,pair_eq_iff}.
    \end{subproof}
    $\dom{e}\subseteq X\times C$ by \cref{subseteq}.
    Show that for all $p$ if $p\in X\times C$ then $p\in\dom{e}$.
    \begin{subproof}
        Fix $p$.
        Assume $p\in X\times C$.
        Take $x,h$ such that $p = (x,h)$ and $x\in X$ and $h\in C$ by \cref{times_elem_is_tuple}.
        $p\in\dom{e}$ by \cref{pair_eq_iff,dom_intro}.
    \end{subproof}
    $X\times C\subseteq\dom{e}$ by \cref{subseteq}.
    $\dom{e} = X\times C$ by \cref{subseteq_antisymmetric}.
    Show that for all $p\in\dom{e}$ we have $e(p)\in Y$.
    \begin{subproof}
        Follows by \cref{function_apply_elim,pair_eq_iff,funs_type_apply,subseteq}.
    \end{subproof}
    $e\in\funs{X\times C}{Y}$ by \cref{funs_intro}.
    Show that for all $x,h$ such that $x\in X$ and $h\in C$ we have $\apply{e}{(x,h)} = \apply{h}{x}$.
    \begin{subproof}
        Follows by \cref{times_tuple_intro,pair_eq_iff,function_apply_elim}.
    \end{subproof}
    There exists $f$ such that $f$ is the evaluation map from $X$ and $C$ to $Y$ by \cref{evaluation_map,pair_eq_iff}.
\end{proof}

% from §46 helper lemma 32: evaluation composed with the induced graph map recovers the original map
\begin{lemma}\label{evaluation_composite_with_induced_graph}
    Assume $T$ is a topology on $X$.
    Assume $S$ is a topology on $Z$.
    Assume $K$ is a topology on $C$.
    Assume $e$ is the evaluation map from $X$ and $C$ to $Y$.
    Assume $f\in\funs{X\times Z}{Y}$.
    Assume $F$ is the map induced by $f$ from $Z$ to $C$ with respect to $X$ and $Y$.
    Assume $F$ is continuous from $Z$ to $C$ with respect to $S$ and $K$.
    Let $P = \productTopology{T}{S}{X}{Z}$.
    Let $Q = \productTopology{T}{K}{X}{C}$.
    Let $g_1 = \piOne{X}{Z}$.
    Let $g_2 = F\circ \piTwo{X}{Z}$.
    Let $G = \{(p,(g_1(p),g_2(p))) \mid p\in X\times Z\}$.
    Let $h = e\circ G$.
    Then $h = f$.
\end{lemma}
\begin{proof}
    We have $G$ is continuous from $X\times Z$ to $X\times C$ with respect to $P$ and $Q$
        by \cref{projection_one_continuous,projection_two_continuous,continuous_composite,continuous_map_into_product_backward}.
    Show that for all $p\in X\times Z$ we have $h(p) = f(p)$.
    \begin{subproof}
        Fix $p\in X\times Z$.
        Take $x,z$ such that $p = (x,z)$ and $x\in X$ and $z\in Z$ by \cref{times_elem_is_tuple}.
        $\apply{G}{p} = (x,\apply{F}{z})$
            by \cref{projection_first,projection_one_continuous,projection_second,projection_two_continuous,circ_apply,continuous,funs_is_function,function_apply_intro,funs_ran,subseteq,funs,pair_eq_iff}.
        $h(p) = \apply{\apply{F}{z}}{x}$
            by \cref{circ_apply,continuous,funs_is_function,funs_ran,subseteq,evaluation_map,funs,pair_eq_iff,funs_type_apply}.
        $h(p) = f(p)$ by \cref{induced_map_into_function_space,pair_eq_iff}.
    \end{subproof}
    $h\in\funs{X\times Z}{Y}$ by \cref{evaluation_map,continuous,funs_circ}.
    $h = f$ by \cref{funs_elim,subseteq,fun_subseteq,subseteq_antisymmetric}.
\end{proof}

% from §46 helper lemma 33: continuity follows from the continuous induced map via evaluation
\begin{lemma}\label{induced_map_continuity_via_evaluation}
    Assume $T$ is a topology on $X$.
    Assume $S$ is a topology on $Z$.
    Assume $U$ is a topology on $Y$.
    Assume $X$ is locally compact with respect to $T$.
    Assume $X$ is Hausdorff with respect to $T$.
    Assume $C$ is the set of continuous functions from $X$ to $Y$ with respect to $T$ and $U$.
    Assume $K$ is the compact open topology on continuous functions from $X$ to $Y$ with respect to $T$ and $U$.
    Assume $f\in\funs{X\times Z}{Y}$.
    Assume $F$ is the map induced by $f$ from $Z$ to $C$ with respect to $X$ and $Y$.
    Assume $F$ is continuous from $Z$ to $C$ with respect to $S$ and $K$.
    Then $f$ is continuous from $X\times Z$ to $Y$ with respect to $\productTopology{T}{S}{X}{Z}$ and $U$.
\end{lemma}
\begin{proof}
    Let $P = \productTopology{T}{S}{X}{Z}$.
    Let $Q = \productTopology{T}{K}{X}{C}$.
    Take $e$ such that $e$ is the evaluation map from $X$ and $C$ to $Y$ by \cref{evaluation_map_exists}.
    Take $F_0, S_0$ such that
        $F_0$ is the set of continuous functions from $X$ to $Y$ with respect to $T$ and $U$ and
        $S_0$ is the compact-open subbasis on continuous functions from $X$ to $Y$ with respect to $T$ and $U$ and
        $K = \subbasisTopology{S_0}{F_0}$ by \cref{compact_open_topology}.
    $K = \subbasisTopology{S_0}{C}$
        by \cref{compact_open_topology,continuous_function_sets_equal,pair_eq_iff}.
    $S_0$ is a subbasis on $C$ by \cref{compact_open_subbasis_is_subbasis}.
    $K$ is a topology on $C$
        by \cref{subbasis_finite_intersections_basis,basis_for_topology,topology_generated_by_subbasis,basis_topology_is_topology}.
    $e$ is continuous from $X\times C$ to $Y$ with respect to $Q$ and $U$
        by \cref{evaluation_map_continuous_compact_open}.
    Let $g_1 = \piOne{X}{Z}$.
    Let $g_2 = F\circ \piTwo{X}{Z}$.
    Let $G = \{(p,(g_1(p),g_2(p))) \mid p\in X\times Z\}$.
    $G$ is continuous from $X\times Z$ to $X\times C$ with respect to $P$ and $Q$
        by \cref{projection_one_continuous,projection_two_continuous,continuous_composite,continuous_map_into_product_backward}.
    Let $h = e\circ G$.
    $h = f$ by \cref{evaluation_composite_with_induced_graph}.
    $f$ is continuous from $X\times Z$ to $Y$ with respect to $P$ and $U$
        by \cref{continuous_composite,pair_eq_iff}.
\end{proof}

% from §46 Item 15: continuity of an induced map characterizes continuity for locally compact Hausdorff domains
\begin{theorem}\label{induced_map_continuity_characterization_compact_open}
    Assume $T$ is a topology on $X$.
    Assume $S$ is a topology on $Z$.
    Assume $U$ is a topology on $Y$.
    Assume $X$ is locally compact with respect to $T$.
    Assume $X$ is Hausdorff with respect to $T$.
    Assume $C$ is the set of continuous functions from $X$ to $Y$ with respect to $T$ and $U$.
    Assume $K$ is the compact open topology on continuous functions from $X$ to $Y$ with respect to $T$ and $U$.
    Assume $f\in\funs{X\times Z}{Y}$.
    Assume $F$ is the map induced by $f$ from $Z$ to $C$ with respect to $X$ and $Y$.
    Suppose $F$ is continuous from $Z$ to $C$ with respect to $S$ and $K$.
    Then $f$ is continuous from $X\times Z$ to $Y$ with respect to $\productTopology{T}{S}{X}{Z}$ and $U$.
\end{theorem}
\begin{proof}
    Follows by \cref{induced_map_continuity_via_evaluation}.
\end{proof}

% from §46 Item 16: homotopy
\begin{definition}\label{homotopy_between_maps}
    $h$ is a homotopy between $f$ and $g$ from $X$ to $Y$ with respect to $T$ and $U$ iff
        $f$ is continuous from $X$ to $Y$ with respect to $T$ and $U$ and
        $g$ is continuous from $X$ to $Y$ with respect to $T$ and $U$ and
        $h\in\funs{X\times\intervalclosed{0}{1}}{Y}$ and
        $h$ is continuous from $X\times\intervalclosed{0}{1}$ to $Y$ with respect to
            $\productTopology{T}{\subspaceTopology{\standardTopologyR}{\intervalclosed{0}{1}}}{X}{\intervalclosed{0}{1}}$
            and $U$ and
        for all $x\in X$ we have
            $\apply{h}{(x,0)} = \apply{f}{x}$ and
            $\apply{h}{(x,1)} = \apply{g}{x}$.
\end{definition}

% from §46 Item 16: homotopic maps
\begin{definition}\label{homotopic_maps}
    $f$ is homotopic to $g$ from $X$ to $Y$ with respect to $T$ and $U$ iff
        there exists $h$ such that
            $h$ is a homotopy between $f$ and $g$ from $X$ to $Y$ with respect to $T$ and $U$.
\end{definition}
\section{§47 Ascoli's Theorem}

% from §47 helper definition 1: pointwise compact closure of a function family
\begin{definition}\label{pointwise_compact_closure_function_family}
    $F$ has pointwise compact closure from $X$ to $Y$ with respect to $d$ iff
        $d$ is a metric on $Y$ and
        $F\subseteq\funs{X}{Y}$ and
        for all $a\in X$ there exists $A$ such that
            $A = \{\apply{f}{a} \mid f\in F\}$ and
            $\closure{A}{Y}{\metricTopology{d}{Y}}$ is compact with respect to
                $\subspaceTopology{\metricTopology{d}{Y}}{\closure{A}{Y}{\metricTopology{d}{Y}}}$.
\end{definition}

% from §47 helper lemma 1: finer ambient topologies induce finer subspace topologies
\begin{lemma}\label{subspace_finer_from_ambient_finer_47}
    Assume $T_1$ is finer than $T$ on $X$.
    Assume $A\subseteq X$.
    Then $\subspaceTopology{T_1}{A}$ is finer than $\subspaceTopology{T}{A}$ on $A$.
\end{lemma}
\begin{proof}
    Follows by \cref{finer_topology,subspace_topology_is_topology,subspace_topology,subseteq}.
\end{proof}

% from §47 helper lemma 2: constant indexed products are function spaces
\begin{lemma}\label{constant_indexed_product_function_space_47}
    Let $R = \{(a,Y) \mid a\in X\}$.
    Then $\productFamily{R} = \funs{X}{Y}$.
\end{lemma}
\begin{proof}
    For all $a,b_1,b_2$ such that $(a,b_1)\in R$ and $(a,b_2)\in R$ we have $b_1 = b_2$
        by \cref{pair_eq_iff}.
    $R$ is a function on $X$ by \cref{relation,rightunique,dom_iff,dom_intro,pair_eq_iff,subseteq,subseteq_antisymmetric}.
    Show that for all $f$ such that $f\in\productFamily{R}$ we have $f\in\funs{X}{Y}$.
    \begin{subproof}
        Follows by \cref{indexed_product,funs_elim,subseteq,function_apply_intro,funs_intro}.
    \end{subproof}
    $\productFamily{R}\subseteq\funs{X}{Y}$ by \cref{subseteq}.
    Show that for all $f$ such that $f\in\funs{X}{Y}$ we have $f\in\productFamily{R}$.
    \begin{subproof}
        Follows by \cref{funs,funs_is_function,subseteq,funs_type_apply,ran_intro,unions_iff,funs_intro,function_apply_intro,indexed_product}.
    \end{subproof}
    $\funs{X}{Y}\subseteq\productFamily{R}$ by \cref{subseteq}.
    $\productFamily{R} = \funs{X}{Y}$ by \cref{subseteq_antisymmetric}.
\end{proof}

% from §47 helper lemma 3: equicontinuous families consist of continuous functions
\begin{lemma}\label{equicontinuous_family_continuous_functions_47}
    Assume $T$ is a topology on $X$.
    Assume $d$ is a metric on $Y$.
    Let $U = \metricTopology{d}{Y}$.
    Assume $H\subseteq\funs{X}{Y}$.
    Assume $C$ is the set of continuous functions from $X$ to $Y$ with respect to $T$ and $U$.
    Suppose $H$ is equicontinuous from $X$ to $Y$ with respect to $T$ and $d$.
    Then $H\subseteq C$.
\end{lemma}
\begin{proof}
    $U$ is a topology on $Y$
        by \cref{metric_topology,metric_basis_is_basis,basis_for_topology,basis_topology_is_topology}.
    Show that for all $h$ such that $h\in H$ we have $h\in C$.
    \begin{subproof}
        Fix $h\in H$.
        $h\in\funs{X}{Y}$ by \cref{subseteq}.
        Show that for all $x\in X$ we have
            $h$ is continuous at $x$ from $X$ to $Y$ with respect to $T$ and $U$.
        \begin{subproof}
            Fix $x\in X$.
            Show that for all $V$ such that $V$ is a neighborhood of $h(x)$ in $Y$ with respect to $U$
                there exists $W$ such that
                    $W$ is a neighborhood of $x$ in $X$ with respect to $T$ and
                    $\img{h}{W}\subseteq V$.
            \begin{subproof}
                Fix $V$ such that $V$ is a neighborhood of $h(x)$ in $Y$ with respect to $U$.
                Take $\epsilon\in\reals$ such that $0 < \epsilon$ and
                    $\metricBall{d}{Y}{h(x)}{\epsilon}\subseteq V$
                    by \cref{neighborhood,open_in,topology_on,subseteq,pow_iff,metric_open_iff}.
                Take $W$ such that
                    $W$ is a neighborhood of $x$ in $X$ with respect to $T$ and
                    for all $y, f$ such that $y\in W$ and $f\in H$ we have
                        $d((\apply{f}{y},\apply{f}{x})) < \epsilon$
                    by \cref{equicontinuous_family,equicontinuous_at_point}.
                For all $z$ such that $z\in\img{h}{W}$ we have $z\in V$
                    by \cref{img_iff,funs_is_function,function_apply_intro,neighborhood,open_in,topology_on,subseteq,pow_iff,funs_type_apply,equicontinuous_family,equicontinuous_at_point,metric_on,metric_symmetric,real_lt_subst_left,metric_ball}.
                $\img{h}{W}\subseteq V$ by \cref{subseteq}.
                Follows by definition.
            \end{subproof}
            $h$ is continuous at $x$ from $X$ to $Y$ with respect to $T$ and $U$
                by \cref{continuous_at}.
        \end{subproof}
        $h\in C$ by \cref{continuous_at_implies_continuous,continuous_function_set}.
    \end{subproof}
    $H\subseteq C$ by \cref{subseteq}.
\end{proof}

% from §47 helper lemma 4: constant indexed families are functions with the expected values
\begin{lemma}\label{constant_indexed_family_47}
    Let $Q = \{(a,B) \mid a\in X\}$.
    Then $Q$ is a function and
        $\dom{Q} = X$ and
        for all $a\in X$ we have $\apply{Q}{a} = B$.
\end{lemma}
\begin{proof}
    For all $a,b_1,b_2$ such that $(a,b_1)\in Q$ and $(a,b_2)\in Q$ we have $b_1 = b_2$
        by \cref{pair_eq_iff}.
    Show that $Q$ is a function and
        $\dom{Q} = X$ and
        for all $a\in X$ we have $\apply{Q}{a} = B$.
    \begin{subproof}
        Follows by \cref{relation,rightunique,dom_iff,dom_intro,pair_eq_iff,subseteq,subseteq_antisymmetric,function_apply_intro}.
    \end{subproof}
\end{proof}

% from §47 helper lemma 5: product topology on coordinate subspaces is finer than the ambient subspace topology
\begin{lemma}\label{indexed_product_subspace_finer_47}
    Assume $U$ is a topology on $Y$.
    Let $R = \{(a,Y) \mid a\in X\}$.
    Let $S = \{(a,U) \mid a\in X\}$.
    Assume $A$ is a function on $X$.
    Assume $X$ is inhabited.
    Assume for all $a\in X$ we have $\apply{A}{a}\subseteq Y$.
    Let $T_1 = \{(a,\subspaceTopology{U}{\apply{A}{a}}) \mid a\in X\}$.
    Let $P = \productFamily{A}$.
    Let $T_P = \productTopologyIndexed{T_1}{A}$.
    Let $T_0 = \subspaceTopology{\productTopologyIndexed{S}{R}}{P}$.
    Then $T_P$ is finer than $T_0$ on $P$.
\end{lemma}
\begin{proof}
    Show that $R$ is a function and $\dom{R} = X$ and for all $a\in X$ we have $\apply{R}{a} = Y$.
    \begin{subproof}
        Follows by \cref{constant_indexed_family_47}.
    \end{subproof}
    Show that $S$ is a function and $\dom{S} = X$ and for all $a\in X$ we have $\apply{S}{a} = U$.
    \begin{subproof}
        Follows by \cref{constant_indexed_family_47}.
    \end{subproof}
    For all $a,b_1,b_2$ such that $(a,b_1)\in T_1$ and $(a,b_2)\in T_1$ we have $b_1 = b_2$
        by \cref{pair_eq_iff}.
    Show that $T_1$ is a function and
        $\dom{T_1} = X$ and
        for all $a\in X$ we have $\apply{T_1}{a} = \subspaceTopology{U}{\apply{A}{a}}$.
    \begin{subproof}
        Follows by \cref{relation,rightunique,dom_iff,dom_intro,pair_eq_iff,subseteq,subseteq_antisymmetric,function_apply_intro}.
    \end{subproof}
    For all $a\in X$ we have $\apply{T_1}{a}$ is a topology on $\apply{A}{a}$
        by \cref{subseteq,subspace_topology_is_topology,function_apply_intro}.
    $A\in\funs{X}{\pow{Y}}$ by \cref{funs,subseteq,powerset_intro,funs_intro}.
    For all $p$ such that $p\in P$ we have $p\in\productFamily{R}$
        by \cref{indexed_product,funs_elim,funs,subseteq,funs_intro,constant_indexed_product_function_space_47}.
    $P\subseteq\productFamily{R}$ by \cref{subseteq}.
    Let $B_0 = \finiteIntersections{\productSubbasis{S}{R}}$.
    $B_0$ is a basis for $\productTopologyIndexed{S}{R}$ on $\productFamily{R}$
        by \cref{product_subbasis_finite_intersections_basis}.
    Let $C = \{D\in\pow{P} \mid \exists E\in B_0. D = E\inter P\}$.
    $C$ is a basis for $T_0$ on $P$ by \cref{subspace_basis}.
    Let $B_1 = \finiteIntersections{\productSubbasis{T_1}{A}}$.
    $B_1$ is a basis for $T_P$ on $P$ by \cref{product_subbasis_finite_intersections_basis}.
    $\productSubbasis{T_1}{A}$ is a subbasis on $P$ by \cref{product_subbasis_is_subbasis}.
    Show that for all $D$ such that $D\in C$ we have $D\in T_P$.
    \begin{subproof}
        Fix $D$ such that $D\in C$.
        Take $E\in B_0$ such that $D = E\inter P$ by \cref{subseteq}.
        Take $F\subseteq\productSubbasis{S}{R}$ such that $F$ is finite and inhabited and $E = \inters{F}$
            by \cref{finite_intersections}.
        Let $g = \{(V,V\inter P) \mid V\in F\}$.
        For all $V,W_1,W_2$ such that $(V,W_1)\in g$ and $(V,W_2)\in g$ we have $W_1 = W_2$
            by \cref{pair_eq_iff}.
        $g$ is a function and $\dom{g} = F$
            by \cref{rightunique,relation,dom_iff,dom_intro,pair_eq_iff,subseteq,subseteq_antisymmetric}.
        Let $H = \img{g}{F}$.
        $H$ is finite by \cref{finite_image_function}.
        Take $V_0$ such that $V_0\in F$ by \cref{nonempty_inhabited}.
        $H$ is inhabited by \cref{subseteq,img_iff,function_apply_intro,nonempty_inhabited,inhabited_nonempty}.
        Show that for all $W$ such that $W\in H$ we have $W\in\productSubbasis{T_1}{A}$.
        \begin{subproof}
            Fix $W$ such that $W\in H$.
            Take $V\in F$ such that $(V,W)\in g$ by \cref{img_iff}.
            $W = V\inter P$ by \cref{function_apply_intro}.
            Take $a\in X$ such that $V\in\productSubbasisAt{S}{R}{a}$ by \cref{subseteq,product_subbasis_union_indexed}.
            Take $Z\in\apply{S}{a}$ such that $V = \preimg{\piIndex{R}{a}}{Z}$ by \cref{product_subbasis_indexed}.
            Let $Z_0 = \apply{A}{a}\inter Z$.
            $Z_0\in\subspaceTopology{U}{\apply{A}{a}}$ by \cref{subspace_topology,function_apply_intro}.
            Show that for all $p$ such that $p\in W$ we have $p\in\preimg{\piIndex{A}{a}}{Z_0}$.
            \begin{subproof}
                Fix $p$ such that $p\in W$.
                $p\in V$ and $p\in P$ by \cref{inter}.
                Take $z$ such that $z\in Z$ and $(p,z)\in\piIndex{R}{a}$ by \cref{preimg_iff}.
                $p(a)\in Z_0$
                    by \cref{projection_map_indexed,pair_eq_iff,indexed_product,funs_elim,inter_intro}.
                $(p,p(a))\in\piIndex{A}{a}$ by \cref{projection_map_indexed,pair_eq_iff}.
                $p\in\preimg{\piIndex{A}{a}}{Z_0}$ by \cref{preimg_iff}.
            \end{subproof}
            $W\subseteq \preimg{\piIndex{A}{a}}{Z_0}$ by \cref{subseteq}.
            Show that for all $p$ such that $p\in\preimg{\piIndex{A}{a}}{Z_0}$ we have $p\in W$.
            \begin{subproof}
                Follows by \cref{preimg_iff,projection_map_indexed,pair_eq_iff,inter_elim_right,subseteq,inter_intro}.
            \end{subproof}
            $\preimg{\piIndex{A}{a}}{Z_0}\subseteq W$ by \cref{subseteq}.
            $W = \preimg{\piIndex{A}{a}}{Z_0}$ by \cref{subseteq_antisymmetric}.
            $W\in\pow{P}$ by \cref{inter_lower_right,powerset_intro}.
            $W\in\productSubbasisAt{T_1}{A}{a}$ by \cref{product_subbasis_indexed,function_apply_intro}.
            $W\in\productSubbasis{T_1}{A}$ by \cref{product_subbasis_union_indexed}.
        \end{subproof}
        $H\subseteq\productSubbasis{T_1}{A}$ by \cref{subseteq}.
        For all $p$ such that $p\in\inters{H}$ we have $p\in P$
            by \cref{nonempty_inhabited,img_iff,function_apply_intro,inters_iff_forall,inter_elim_right}.
        $\inters{H}\subseteq P$ by \cref{subseteq}.
        $\unions{\productSubbasis{T_1}{A}} = P$ by \cref{subbasis_on}.
        $\inters{H}\in\pow{\unions{\productSubbasis{T_1}{A}}}$
            by \cref{powerset_intro,subseteq_antisymmetric}.
        $\inters{H}\in T_P$
            by \cref{finite_intersections,basis_elem_open_in_generated,basis_for_topology}.
        For all $p$ such that $p\in D$ we have $p\in\inters{H}$
            by \cref{inter,img_iff,function_apply_intro,inters_iff_forall,inter_intro,inters_intro}.
        $D\subseteq \inters{H}$ by \cref{subseteq}.
        For all $p$ such that $p\in\inters{H}$ we have $p\in D$
            by \cref{nonempty_inhabited,subseteq,img_iff,function_apply_intro,inters_iff_forall,inter_elim_right,inter_elim_left,inters_intro,inter_intro}.
        $\inters{H}\subseteq D$ by \cref{subseteq}.
        $D\in T_P$ by \cref{subseteq_antisymmetric}.
    \end{subproof}
    $C\subseteq T_P$ by \cref{subseteq}.
    $T_0$ is a topology on $P$ and $T_P$ is a topology on $P$
        by \cref{subspace_topology_is_topology,basis_topology_is_topology,basis_for_topology}.
    For all $p,D$ such that $p\in P$ and $D\in C$ and $p\in D$
        there exists $B\in B_1$ such that $p\in B$ and $B\subseteq D$
        by \cref{subseteq,basis_for_topology,topology_generated_by_basis}.
    $T_P$ is finer than $T_0$ on $P$ by \cref{basis_finer_criterion}.
\end{proof}

% from §47 helper lemma 6: constant indexed product topology is pointwise convergence
\begin{lemma}\label{constant_indexed_product_topology_is_pointwise_47}
    Assume $U$ is a topology on $Y$.
    Assume $X$ is inhabited.
    Let $R = \{(a,Y) \mid a\in X\}$.
    Let $S = \{(a,U) \mid a\in X\}$.
    Let $T_1 = \productTopologyIndexed{S}{R}$.
    Then $T_1$ is the topology of pointwise convergence on functions from $X$ to $Y$ with respect to $U$.
\end{lemma}
\begin{proof}
    $R$ is a function and $\dom{R} = X$ and for all $a\in X$ we have $\apply{R}{a} = Y$
        by \cref{constant_indexed_family_47}.
    $S$ is a function and $\dom{S} = X$ and for all $a\in X$ we have $\apply{S}{a} = U$
        by \cref{constant_indexed_family_47}.
    Let $P = \productSubbasis{S}{R}$.
    Show that for all $A$ if $A\in P$ then there exist $x,V$ such that
        $x\in X$ and
        $V\in U$ and
        $A$ is the point-open set at $x$ with value set $V$ from $X$ to $Y$.
    \begin{subproof}
        Fix $A\in P$.
        Take $x\in X$ such that $A\in\productSubbasisAt{S}{R}{x}$
            by \cref{product_subbasis_union_indexed}.
        Take $V\in\apply{S}{x}$ such that $A = \preimg{\piIndex{R}{x}}{V}$
            by \cref{product_subbasis_indexed}.
        $V\subseteq Y$ by \cref{function_apply_intro,topology_on,subseteq,pow_iff}.
        Let $B = \{f\in\funs{X}{Y} \mid \apply{f}{x}\in V\}$.
        Show that for all $f$ if $f\in A$ then $f\in B$.
        \begin{subproof}
            Fix $f$.
            Assume $f\in A$.
            Take $y\in V$ such that $(f,y)\in\piIndex{R}{x}$ by \cref{preimg_iff}.
            Take $g\in\productFamily{R}$ such that $(g,\apply{g}{x}) = (f,y)$ by \cref{projection_map_indexed}.
            $g\in\funs{X}{Y}$ by \cref{constant_indexed_product_function_space_47,subseteq}.
            $f = g$ and $y = g(x)$ by \cref{pair_eq_iff}.
            $f\in B$ by \cref{pair_eq_iff}.
        \end{subproof}
        $A\subseteq B$ by \cref{subseteq}.
        Show that for all $f$ if $f\in B$ then $f\in A$.
        \begin{subproof}
            Fix $f$.
            Assume $f\in B$.
            $f\in\funs{X}{Y}$ and $f(x)\in V$ by \cref{inters_iff_forall}.
            $f\in\productFamily{R}$ by \cref{constant_indexed_product_function_space_47,subseteq}.
            $(f,\apply{f}{x})\in\piIndex{R}{x}$ by \cref{projection_map_indexed,pair_eq_iff}.
            $f\in A$ by \cref{preimg_iff}.
        \end{subproof}
        $B\subseteq A$ by \cref{subseteq}.
        $A = B$ by \cref{subseteq_antisymmetric}.
        $A$ is the point-open set at $x$ with value set $V$ from $X$ to $Y$ by \cref{point_open_set}.
    \end{subproof}
    For all $A\in P$ we have $A\in\pow{\funs{X}{Y}}$ by \cref{point_open_set,subseteq,powerset_intro}.
    Show that for all $A$ if
        there exist $x,V$ such that
            $x\in X$ and
            $V\in U$ and
            $A$ is the point-open set at $x$ with value set $V$ from $X$ to $Y$
    then $A\in P$.
    \begin{subproof}
        Fix $A$.
        Assume there exist $x,V$ such that
            $x\in X$ and
            $V\in U$ and
            $A$ is the point-open set at $x$ with value set $V$ from $X$ to $Y$.
        Take $x,V$ such that
            $x\in X$ and
            $V\in U$ and
            $A$ is the point-open set at $x$ with value set $V$ from $X$ to $Y$.
        Show that for all $f$ such that $f\in A$ we have $f\in\preimg{\piIndex{R}{x}}{V}$.
        \begin{subproof}
            Follows by \cref{point_open_set,constant_indexed_product_function_space_47,projection_map_indexed,function_apply_elim,preimg_iff}.
        \end{subproof}
        $A\subseteq \preimg{\piIndex{R}{x}}{V}$ by \cref{subseteq}.
        Show that for all $f$ such that $f\in\preimg{\piIndex{R}{x}}{V}$ we have $f\in A$.
        \begin{subproof}
            Follows by \cref{preimg_iff,projection_map_indexed,pair_eq_iff,constant_indexed_product_function_space_47,point_open_set}.
        \end{subproof}
        $\preimg{\piIndex{R}{x}}{V}\subseteq A$ by \cref{subseteq}.
        $A = \preimg{\piIndex{R}{x}}{V}$ by \cref{subseteq_antisymmetric}.
        $A\in\pow{\productFamily{R}}$
            by \cref{projection_map_indexed_function,funs_elim,preimg_subseteq_dom,subseteq_transitive,pow_iff}.
        $A\in\productSubbasisAt{S}{R}{x}$ by \cref{function_apply_intro,product_subbasis_indexed}.
        $A\in P$ by \cref{product_subbasis_union_indexed}.
    \end{subproof}
    $P$ is the pointwise subbasis on functions from $X$ to $Y$ with respect to $U$
        by \cref{pointwise_function_space_subbasis}.
    $T_1$ is a topology on $\productFamily{R}$
        by \cref{product_subbasis_finite_intersections_basis,basis_for_topology,basis_topology_is_topology}.
    $\productFamily{R} = \funs{X}{Y}$ by \cref{constant_indexed_product_function_space_47}.
    $T_1 = \subbasisTopology{P}{\funs{X}{Y}}$ by \cref{product_topology_indexed,pair_eq_iff}.
    $T_1$ is the topology of pointwise convergence on functions from $X$ to $Y$ with respect to $U$
        by \cref{pointwise_convergence_topology}.
\end{proof}

% from §47 helper lemma 7: compactness passes to coarser topologies
\begin{lemma}\label{compactness_from_finer_topology_47}
    Assume $T_1$ is finer than $T$ on $X$.
    Suppose $X$ is compact with respect to $T_1$.
    Then $X$ is compact with respect to $T$.
\end{lemma}
\begin{proof}
    Follows by \cref{finer_topology,open_covering,open_in,subseteq,compact}.
\end{proof}

% from §47 helper lemma 8: coordinate projection on a constant indexed product evaluates a function
\begin{lemma}\label{constant_indexed_projection_apply_47}
    Assume $a\in X$.
    Let $R = \{(x,Y) \mid x\in X\}$.
    Assume $f\in\funs{X}{Y}$.
    Then $\apply{\piIndex{R}{a}}{f} = f(a)$.
\end{lemma}
\begin{proof}
    Follows by \cref{constant_indexed_family_47,constant_indexed_product_function_space_47,projection_map_indexed_function,funs_is_function,projection_map_indexed,function_apply_elim,function_apply_intro}.
\end{proof}

% from §47 helper lemma 9: pointwise closure of an equicontinuous family is equicontinuous
\begin{lemma}\label{pointwise_closure_equicontinuous_47}
    Assume $T$ is a topology on $X$.
    Assume $d$ is a metric on $Y$.
    Let $U = \metricTopology{d}{Y}$.
    Assume $T_1$ is the topology of pointwise convergence on functions from $X$ to $Y$ with respect to $U$.
    Assume $F\subseteq\funs{X}{Y}$.
    Suppose $F$ is equicontinuous from $X$ to $Y$ with respect to $T$ and $d$.
    Let $G = \closure{F}{\funs{X}{Y}}{T_1}$.
    Then $G$ is equicontinuous from $X$ to $Y$ with respect to $T$ and $d$.
\end{lemma}
\begin{proof}
    $T_1$ is a topology on $\funs{X}{Y}$ by \cref{pointwise_convergence_topology}.
    Show that for all $x_0\in X$ we have
        $G$ is equicontinuous at $x_0$ from $X$ to $Y$ with respect to $T$ and $d$.
    \begin{subproof}
        Fix $x_0\in X$.
        Take $S$ such that
            $S$ is the pointwise subbasis on functions from $X$ to $Y$ with respect to $U$ and
            $T_1 = \subbasisTopology{S}{\funs{X}{Y}}$ by \cref{pointwise_convergence_topology}.
        Show that for all $\epsilon\in\reals$ such that $0 < \epsilon$
            there exists $W$ such that
                $W$ is a neighborhood of $x_0$ in $X$ with respect to $T$ and
                for all $x, g$ such that $x\in W$ and $g\in G$ we have
                    $d((\apply{g}{x},\apply{g}{x_0})) < \epsilon$.
        \begin{subproof}
            Fix $\epsilon$ such that $\epsilon\in\reals$ and $0 < \epsilon$.
            Take $\delta\in\reals$ such that
                $0 < \delta$ and
                $\delta < 1$ and
                $\delta + \delta + \delta < \epsilon$
                by \cref{positive_real_small_triple_bound}.
            Take $W$ such that
                $W$ is a neighborhood of $x_0$ in $X$ with respect to $T$ and
                for all $x, f$ such that $x\in W$ and $f\in F$ we have
                    $d((\apply{f}{x},\apply{f}{x_0})) < \delta$
                by \cref{equicontinuous_family,equicontinuous_at_point}.
            Show that for all $x, g$ such that $x\in W$ and $g\in G$ we have
                $d((\apply{g}{x},\apply{g}{x_0})) < \epsilon$.
            \begin{subproof}
                Fix $x$.
                Fix $g$.
                Assume $x\in W$ and $g\in G$.
                $g\in\funs{X}{Y}$ by \cref{closure,subseteq}.
                $x\in X$ by \cref{neighborhood,open_in,topology_on,subseteq,powerset_elim}.
                Let $C = \{x_0,x\}$.
                $C\subseteq X$ and $C$ is finite and inhabited
                    by \cref{upair_iff,subseteq,pair_is_finite,upair_intro_left,nonempty_inhabited,emptyset}.
                Take $M$ such that
                    $M$ is open in $\funs{X}{Y}$ with respect to $\subbasisTopology{S}{\funs{X}{Y}}$ and
                    $g\in M$ and
                    for all $h$ if $h\in M$ then
                        $h\in\funs{X}{Y}$ and
                        for all $z\in C$ we have $h(z)\in\metricBall{d}{Y}{g(z)}{\delta}$
                    by \cref{nonempty_inhabited,emptyset,pointwise_subbasis_is_subbasis_inhabited_domain,finite_point_open_metric_neighborhoods_open}.
                $M$ intersects $F$ by \cref{closure_iff_intersects_open,pair_eq_iff}.
                Take $f$ such that $f\in M\inter F$ by \cref{intersects,nonempty_inhabited}.
                $f\in M$ and $f\in F$ by \cref{inter}.
                $f(x)\in\metricBall{d}{Y}{g(x)}{\delta}$ and
                    $f(x_0)\in\metricBall{d}{Y}{g(x_0)}{\delta}$ by \cref{subseteq,upair_intro_right,upair_intro_left}.
                $d((g(x),f(x))) < \delta$ and $d((g(x_0),f(x_0))) < \delta$ by \cref{metric_ball}.
                $d((f(x_0),g(x_0))) < \delta$
                    by \cref{metric_on,metric_symmetric,funs_type_apply,real_lt_subst_left}.
                $d((f(x),f(x_0))) < \delta$ by \cref{equicontinuous_at_point}.
                $f(x)\in Y$ and $f(x_0)\in Y$ and $g(x)\in Y$ and $g(x_0)\in Y$
                    by \cref{funs_type_apply}.
                $d((g(x),g(x_0))) < \delta + \delta + \delta$
                    by \cref{metric_three_step_lt_sum}.
                $d((g(x),g(x_0))) < \epsilon$
                    by \cref{metric_on,funs_type_apply,times_tuple_intro,real_add_closed,real_lt_trans}.
            \end{subproof}
            Follows by definition.
        \end{subproof}
        $G$ is equicontinuous at $x_0$ from $X$ to $Y$ with respect to $T$ and $d$
            by \cref{closure,subseteq,equicontinuous_at_point}.
    \end{subproof}
    $G$ is equicontinuous from $X$ to $Y$ with respect to $T$ and $d$
        by \cref{closure,subseteq,equicontinuous_family}.
\end{proof}

% from §47 helper lemma 10: product-closure of an equicontinuous family is equicontinuous
\begin{lemma}\label{coordinatewise_product_closure_equicontinuous_47}
    Assume $T$ is a topology on $X$.
    Assume $d$ is a metric on $Y$.
    Let $U = \metricTopology{d}{Y}$.
    Assume $A$ is a function on $X$.
    Assume for all $a\in X$ we have $\apply{A}{a}\subseteq Y$.
    Let $T_A = \{(a,\subspaceTopology{U}{\apply{A}{a}}) \mid a\in X\}$.
    Let $P = \productFamily{A}$.
    Let $T_P = \productTopologyIndexed{T_A}{A}$.
    Assume $F\subseteq P$.
    Suppose $F$ is equicontinuous from $X$ to $Y$ with respect to $T$ and $d$.
    Let $G = \closure{F}{P}{T_P}$.
    Then $G$ is equicontinuous from $X$ to $Y$ with respect to $T$ and $d$.
\end{lemma}
\begin{proof}
    For all $a,b_1,b_2$ such that $(a,b_1)\in T_A$ and $(a,b_2)\in T_A$ we have $b_1 = b_2$
        by \cref{pair_eq_iff}.
    Show that $T_A$ is a function and
        $\dom{T_A} = X$ and
        for all $a\in X$ we have $\apply{T_A}{a} = \subspaceTopology{U}{\apply{A}{a}}$.
    \begin{subproof}
        Follows by \cref{relation,rightunique,dom_iff,dom_intro,pair_eq_iff,subseteq,subseteq_antisymmetric,function_apply_intro}.
    \end{subproof}
    For all $a\in X$ we have $\apply{T_A}{a}$ is a topology on $\apply{A}{a}$
        by \cref{metric_topology,metric_basis_is_basis,basis_topology_is_topology,basis_for_topology,subseteq,subspace_topology_is_topology,function_apply_intro}.
    For all $a,p$ such that $a\in X$ and $p\in P$ we have $p(a)\in Y$
        by \cref{indexed_product,funs_elim,pair_eq_iff,subseteq}.
    $P\subseteq\funs{X}{Y}$ by \cref{indexed_product,funs_elim,pair_eq_iff,funs_intro,subseteq}.
    \begin{byCase}
    \caseOf{$X = \emptyset$.}
        $G\subseteq\funs{X}{Y}$ by \cref{closure,subseteq,subseteq_transitive}.
        For all $x\in X$ we have
            $G$ is equicontinuous at $x$ from $X$ to $Y$ with respect to $T$ and $d$
            by \cref{emptyset}.
        $G$ is equicontinuous from $X$ to $Y$ with respect to $T$ and $d$
            by \cref{equicontinuous_family}.
    \caseOf{$X \neq \emptyset$.}
        $X$ is inhabited by \cref{nonempty_inhabited}.
        Let $B_P = \finiteIntersections{\productSubbasis{T_A}{A}}$.
        $B_P$ is a basis for $T_P$ on $P$ by \cref{product_subbasis_finite_intersections_basis}.
        $T_P$ is a topology on $P$ by \cref{basis_for_topology,basis_topology_is_topology}.
        Show that for all $x_0\in X$ we have
            $G$ is equicontinuous at $x_0$ from $X$ to $Y$ with respect to $T$ and $d$.
        \begin{subproof}
            Fix $x_0\in X$.
            Show that for all $\epsilon\in\reals$ such that $0 < \epsilon$
                there exists $W$ such that
                    $W$ is a neighborhood of $x_0$ in $X$ with respect to $T$ and
                    for all $x, g$ such that $x\in W$ and $g\in G$ we have
                        $d((\apply{g}{x},\apply{g}{x_0})) < \epsilon$.
            \begin{subproof}
                Fix $\epsilon$ such that $\epsilon\in\reals$ and $0 < \epsilon$.
                Take $\delta\in\reals$ such that
                    $0 < \delta$ and
                    $\delta < 1$ and
                    $\delta + \delta + \delta < \epsilon$
                    by \cref{positive_real_small_triple_bound}.
                Take $W$ such that
                    $W$ is a neighborhood of $x_0$ in $X$ with respect to $T$ and
                    for all $x, f$ such that $x\in W$ and $f\in F$ we have
                        $d((\apply{f}{x},\apply{f}{x_0})) < \delta$
                    by \cref{equicontinuous_family,equicontinuous_at_point}.
                Show that for all $x, g$ such that $x\in W$ and $g\in G$ we have
                    $d((\apply{g}{x},\apply{g}{x_0})) < \epsilon$.
                \begin{subproof}
                    Fix $x$.
                    Fix $g$.
                    Assume $x\in W$ and $g\in G$.
                    $g\in P$ by \cref{closure,subseteq}.
                    $g$ is a function and $\dom{g} = X$ by \cref{indexed_product,funs_elim,pair_eq_iff}.
                    $x\in X$ by \cref{neighborhood,open_in,topology_on,subseteq,powerset_elim}.
                    $g(x_0)\in\apply{A}{x_0}$ and $g(x)\in\apply{A}{x}$ by \cref{indexed_product}.
                    $g(x_0)\in Y$ and $g(x)\in Y$ by \cref{subseteq}.
                    Let $Z_0 = \apply{A}{x_0}\inter\metricBall{d}{Y}{g(x_0)}{\delta}$.
                    Let $Z_1 = \apply{A}{x}\inter\metricBall{d}{Y}{g(x)}{\delta}$.
                    $\metricBall{d}{Y}{g(x_0)}{\delta}\in U$ and
                        $\metricBall{d}{Y}{g(x)}{\delta}\in U$
                        by \cref{metric_ball,subseteq,powerset_intro,metric_basis,metric_basis_is_basis,basis_elem_open_in_generated,metric_topology,open_in}.
                    $Z_0\in\subspaceTopology{U}{\apply{A}{x_0}}$ and
                        $Z_1\in\subspaceTopology{U}{\apply{A}{x}}$ by \cref{subspace_topology}.
                    Let $N_0 = \preimg{\piIndex{A}{x_0}}{Z_0}$.
                    Let $N_1 = \preimg{\piIndex{A}{x}}{Z_1}$.
                    $Z_0\in\apply{T_A}{x_0}$ and $Z_1\in\apply{T_A}{x}$ by \cref{function_apply_intro}.
                    $N_0\subseteq P$ and $N_1\subseteq P$
                        by \cref{projection_map_indexed_function,funs_elim,preimg_subseteq_dom,subseteq_transitive}.
                    $N_0\in\pow{P}$ and $N_1\in\pow{P}$ by \cref{powerset_intro}.
                    $N_0\in\productSubbasisAt{T_A}{A}{x_0}$ and
                        $N_1\in\productSubbasisAt{T_A}{A}{x}$ by \cref{product_subbasis_indexed}.
                    $N_0\in\productSubbasis{T_A}{A}$ and
                        $N_1\in\productSubbasis{T_A}{A}$ by \cref{product_subbasis_union_indexed}.
                    Let $H = \{N_0,N_1\}$.
                    $H\subseteq\productSubbasis{T_A}{A}$ by \cref{upair_iff,subseteq}.
                    $H$ is finite and inhabited by \cref{pair_is_finite,upair_intro_left}.
                    Let $M = \inters{H}$.
                    For all $p$ such that $p\in M$ we have $p\in P$ by \cref{inters_destr,upair_intro_left,subseteq}.
                    $M\subseteq P$ by \cref{subseteq}.
                    $M\in\pow{\unions{\productSubbasis{T_A}{A}}}$
                        by \cref{product_subbasis_is_subbasis,subseteq,subbasis_on,powerset_intro,subseteq_antisymmetric}.
                    $M$ is open in $P$ with respect to $T_P$
                        by \cref{finite_intersections,basis_elem_open_in_generated,basis_for_topology,open_in}.
                    $(g,g(x_0))\in\piIndex{A}{x_0}$ by \cref{projection_map_indexed,function_apply_elim}.
                    $d((g(x_0),g(x_0))) < \delta$
                        by \cref{metric_on,metric_zero_characterization,real_lt_subst_left}.
                    $g(x_0)\in\metricBall{d}{Y}{g(x_0)}{\delta}$ by \cref{metric_ball}.
                    $g(x_0)\in Z_0$ by \cref{inter_intro}.
                    $g\in N_0$ by \cref{preimg_iff}.
                    $(g,g(x))\in\piIndex{A}{x}$ by \cref{projection_map_indexed,function_apply_elim}.
                    $d((g(x),g(x))) < \delta$
                        by \cref{metric_on,metric_zero_characterization,real_lt_subst_left}.
                    $g(x)\in\metricBall{d}{Y}{g(x)}{\delta}$ by \cref{metric_ball}.
                    $g(x)\in Z_1$ by \cref{inter_intro}.
                    $g\in N_1$ by \cref{preimg_iff}.
                    For all $B$ such that $B\in H$ we have $g\in B$ by \cref{upair_iff,real_lt_subst_left,real_lt_subst_right}.
                    $g\in M$ by \cref{inters_intro}.
                    $M$ intersects $F$ by \cref{closure_iff_intersects_open}.
                    Take $f$ such that $f\in M\inter F$ by \cref{intersects,nonempty_inhabited}.
                    $f\in M$ and $f\in F$ by \cref{inter}.
                    $f\in P$ by \cref{subseteq}.
                    $f$ is a function and $\dom{f} = X$ by \cref{indexed_product,funs_elim,pair_eq_iff}.
                    $f(x_0)\in\apply{A}{x_0}$ and $f(x)\in\apply{A}{x}$ by \cref{indexed_product}.
                    $f(x_0)\in Y$ and $f(x)\in Y$ by \cref{subseteq}.
                    $f\in N_0$ and $f\in N_1$ by \cref{inters_iff_forall,upair_intro_left,upair_intro_right}.
                    Take $z_0$ such that $z_0\in Z_0$ and $(f,z_0)\in\piIndex{A}{x_0}$ by \cref{preimg_iff}.
                    Take $z_1$ such that $z_1\in Z_1$ and $(f,z_1)\in\piIndex{A}{x}$ by \cref{preimg_iff}.
                    $z_0 = f(x_0)$ and $z_1 = f(x)$ by \cref{projection_map_indexed,pair_eq_iff}.
                    $f(x_0)\in Z_0$ and $f(x)\in Z_1$ by \cref{projection_map_indexed,pair_eq_iff}.
                    $f(x_0)\in\metricBall{d}{Y}{g(x_0)}{\delta}$ and
                        $f(x)\in\metricBall{d}{Y}{g(x)}{\delta}$ by \cref{inter_elim_right}.
                    $d((g(x_0),f(x_0))) < \delta$ and $d((g(x),f(x))) < \delta$ by \cref{metric_ball}.
                    $d((f(x_0),g(x_0))) < \delta$
                        by \cref{metric_on,metric_symmetric,funs_type_apply,real_lt_subst_left}.
                    $d((f(x),f(x_0))) < \delta$ by \cref{equicontinuous_at_point}.
                    $d((g(x),g(x_0))) < \delta + \delta + \delta$
                        by \cref{metric_three_step_lt_sum}.
                    $d((g(x),g(x_0))) < \epsilon$
                        by \cref{metric_on,funs_type_apply,times_tuple_intro,real_add_closed,real_lt_trans}.
                \end{subproof}
                Follows by definition.
            \end{subproof}
            $G$ is equicontinuous at $x_0$ from $X$ to $Y$ with respect to $T$ and $d$
                by \cref{closure,subseteq,subseteq_transitive,equicontinuous_at_point}.
        \end{subproof}
        $G$ is equicontinuous from $X$ to $Y$ with respect to $T$ and $d$
            by \cref{closure,subseteq,subseteq_transitive,equicontinuous_family}.
    \end{byCase}
\end{proof}

% from §47 helper lemma 11: indexed products of coordinate subsets are function spaces
\begin{lemma}\label{indexed_product_subset_function_space_47}
    Assume $A$ is a function on $X$.
    Assume for all $a\in X$ we have $\apply{A}{a}\subseteq Y$.
    Let $P = \productFamily{A}$.
    Then $P\subseteq\funs{X}{Y}$.
\end{lemma}
\begin{proof}
    Follows by \cref{indexed_product,funs_elim,pair_eq_iff,funs_intro,subseteq}.
\end{proof}

% from §47 helper lemma 12: product topology is finer than compact convergence on equicontinuous subfamilies
\begin{lemma}\label{equicontinuous_product_topology_finer_than_compact_convergence_47}
    Assume $T$ is a topology on $X$.
    Assume $X$ is inhabited.
    Assume $d$ is a metric on $Y$.
    Let $U = \metricTopology{d}{Y}$.
    Assume $T_0$ is the topology of compact convergence on functions from $X$ to $Y$ with respect to $T$ and $d$.
    Assume $A$ is a function on $X$.
    Assume for all $a\in X$ we have $\apply{A}{a}\subseteq Y$.
    Let $T_A = \{(a,\subspaceTopology{U}{\apply{A}{a}}) \mid a\in X\}$.
    Let $P = \productFamily{A}$.
    Let $T_P = \productTopologyIndexed{T_A}{A}$.
    Assume $G\subseteq P$.
    Suppose $G$ is equicontinuous from $X$ to $Y$ with respect to $T$ and $d$.
    Then $\subspaceTopology{T_P}{G}$ is finer than $\subspaceTopology{T_0}{G}$ on $G$.
\end{lemma}
\begin{proof}
    Let $F = \funs{X}{Y}$.
    Let $R = \{(a,Y) \mid a\in X\}$.
    Let $S = \{(a,U) \mid a\in X\}$.
    Let $T_1 = \productTopologyIndexed{S}{R}$.
    $U$ is a topology on $Y$
        by \cref{metric_topology,metric_basis_is_basis,basis_topology_is_topology,basis_for_topology}.
    $G\subseteq F$ by \cref{indexed_product_subset_function_space_47,subseteq_transitive}.
    $T_1$ is the topology of pointwise convergence on functions from $X$ to $Y$ with respect to $U$
        by \cref{constant_indexed_product_topology_is_pointwise_47}.
    $T_P$ is finer than $\subspaceTopology{T_1}{P}$ on $P$
        by \cref{indexed_product_subspace_finer_47}.
    Show that $T_P$ is a topology on $P$ and $\subspaceTopology{T_1}{P}$ is a topology on $P$.
    \begin{subproof}
        Follows by \cref{finer_topology}.
    \end{subproof}
    $\subspaceTopology{T_P}{G}$ is a topology on $G$
        by \cref{subspace_topology_is_topology,subseteq}.
    $\subspaceTopology{T_0}{G}$ is a topology on $G$
        by \cref{subspace_topology_is_topology,compact_convergence_topology,subseteq}.
    Take $B$ such that
        $B$ is the compact-convergence basis on functions from $X$ to $Y$ with respect to $T$ and $d$ and
        $T_0 = \basisTopology{B}{F}$ by \cref{compact_convergence_topology}.
    $T_0$ is a topology on $F$ by \cref{compact_convergence_topology}.
    For all $E\in B$ we have $E$ is open in $F$ with respect to $T_0$
        by \cref{compact_convergence_function_space_basis,compact_convergence_ball_open}.
    For all $W,h$ such that $W\in T_0$ and $h\in W$ there exists $E\in B$ such that
        $h\in E$ and $E\subseteq W$ by \cref{topology_generated_by_basis}.
    $B$ is a basis for $T_0$ on $F$ by \cref{open_collection_basis}.
    Let $C = \{D\in\pow{G} \mid \exists E\in B. D = E\inter G\}$.
    $C$ is a basis for $\subspaceTopology{T_0}{G}$ on $G$ by \cref{subspace_basis}.
    Show that for all $D$ such that $D\in C$ we have
        $D$ is open in $G$ with respect to $\subspaceTopology{T_P}{G}$.
    \begin{subproof}
        Fix $D$ such that $D\in C$.
        Take $E\in B$ such that $D = E\inter G$ by \cref{subseteq}.
        Show that for all $g\in D$ there exists $N$ such that
            $N$ is a neighborhood of $g$ in $G$ with respect to $\subspaceTopology{T_P}{G}$ and
            $N\subseteq D$.
        \begin{subproof}
            Fix $g\in D$.
            $g\in G$ and $g\in E$ by \cref{inter}.
            Take $f_0, K, \epsilon$ such that
                $E$ is the compact-convergence ball about $f_0$ over $K$ with radius $\epsilon$
                    from $X$ to $Y$ with respect to $T$ and $d$
                by \cref{compact_convergence_function_space_basis}.
            Take $H,\delta$ such that
                $\delta\in\reals$ and
                $0 < \delta$ and
                $H$ is the compact-convergence ball about $g$ over $K$ with radius $\delta$
                    from $X$ to $Y$ with respect to $T$ and $d$ and
                $H\subseteq E$ by \cref{compact_convergence_ball_refinement}.
            \begin{byCase}
            \caseOf{$K = \emptyset$.}
                $H = F$ by \cref{compact_convergence_ball_over_empty_is_whole_space}.
                Let $N = G$.
                $N$ is open in $G$ with respect to $\subspaceTopology{T_P}{G}$ by \cref{open_in,topology_on}.
                $N$ is a neighborhood of $g$ in $G$ with respect to $\subspaceTopology{T_P}{G}$
                    by \cref{neighborhood,subseteq_refl}.
                $G\subseteq H$ by \cref{subseteq,real_lt_subst_left}.
                $H\inter G = G$ by \cref{inter_absorb_supseteq_left,subseteq}.
                $N\subseteq H\inter G$ by \cref{pair_eq_iff,subseteq_refl}.
                For all $u$ such that $u\in H\inter G$ we have $u\in E\inter G$ by \cref{inter,subseteq,inter_intro}.
                $H\inter G\subseteq E\inter G$ by \cref{subseteq}.
                $N\subseteq D$ by \cref{subseteq_transitive,pair_eq_iff}.
                Follows by definition.
            \caseOf{$K \neq \emptyset$.}
                $K$ is a compact subspace of $X$ with respect to $T$ by \cref{compact_convergence_ball}.
                $K\subseteq X$ by \cref{compact_subspace}.
                $K$ is inhabited by \cref{nonempty_inhabited}.
                Take $\theta\in\reals$ such that
                    $0 < \theta$ and
                    $\theta < 1$ and
                    $\theta < \delta$ by \cref{positive_real_has_smaller_positive_below_one}.
                Take $\eta\in\reals$ such that
                    $0 < \eta$ and
                    $\eta < 1$ and
                    $\eta + \eta + \eta < \theta$
                    by \cref{positive_real_small_triple_bound}.
                $\eta + \eta + \eta < \delta$ by \cref{real_add_closed,real_lt_trans}.
                Let $N_0 = \{V\in T \mid \text{there exists $a\in K$ such that
                    $V$ is a neighborhood of $a$ in $X$ with respect to $T$ and
                    for all $x,h$ such that $x\in V$ and $h\in G$ we have
                        $d((\apply{h}{x},\apply{h}{a})) < \eta$}\}$.
                For all $x$ such that $x\in K$ we have $x\in\unions{N_0}$
                    by \cref{subseteq,equicontinuous_family,equicontinuous_at_point,neighborhood,open_in,unions_iff}.
                $N_0$ is a covering of $K$ by \cref{subseteq,covering}.
                For all $V$ such that $V\in N_0$ we have $V$ is open in $X$ with respect to $T$ by \cref{subseteq,open_in}.
                $K$ is a subspace of $X$ with respect to $T$ by \cref{subspace_of}.
                $K$ is compact with respect to $\subspaceTopology{T}{K}$ by \cref{compact_subspace}.
                Take $L$ such that
                    $L\subseteq N_0$ and
                    $L$ is finite and
                    $L$ is a covering of $K$
                    by \cref{compact_subspace_open_covering}.
                Let $Q = \{w\in L\times K \mid \text{$\fst{w}$ is a neighborhood of $\snd{w}$ in $X$ with respect to $T$ and
                    for all $x,h$ such that $x\in\fst{w}$ and $h\in G$ we have
                        $d((\apply{h}{x},\apply{h}{\snd{w}})) < \eta$}\}$.
                $Q$ is a relation by \cref{times_elem_is_tuple,relation}.
                For all $V\in L$ there exists $a$ such that $V\mathrel{Q} a$
                    by \cref{subseteq,times_tuple_intro,fst_eq,snd_eq}.
                Take $e$ such that $e$ is a function on $L$ and for all $V\in L$ we have $V\mathrel{Q}\apply{e}{V}$
                    by \cref{choice_function}.
                Let $J = \img{e}{L}$.
                $J$ is finite by \cref{finite_image_function}.
                For all $a$ such that $a\in J$ we have $a\in X$
                    by \cref{img_iff,function_apply_intro,choice_function,pair_eq_iff,subseteq,times_tuple_elim}.
                $J\subseteq X$ by \cref{subseteq}.
                $L$ is inhabited by \cref{covering,subseteq,unions_iff,nonempty_inhabited}.
                $J$ is inhabited by \cref{choice_function,subseteq,function_apply_iff,img_iff,nonempty_inhabited}.
                Take $S_0$ such that
                    $S_0$ is the pointwise subbasis on functions from $X$ to $Y$ with respect to $U$ and
                    $T_1 = \subbasisTopology{S_0}{F}$ by \cref{pointwise_convergence_topology}.
                $g\in F$ by \cref{subseteq}.
                Take $M_0$ such that
                    $M_0$ is open in $F$ with respect to $\subbasisTopology{S_0}{F}$ and
                    $g\in M_0$ and
                    for all $h$ if $h\in M_0$ then
                        $h\in F$ and
                        for all $a\in J$ we have $h(a)\in\metricBall{d}{Y}{g(a)}{\eta}$
                    by \cref{pointwise_subbasis_is_subbasis_inhabited_domain,finite_point_open_metric_neighborhoods_open}.
                Let $M = P\inter M_0$.
                $M$ is open in $P$ with respect to $\subspaceTopology{T_1}{P}$
                    by \cref{subspace_topology,open_in,pair_eq_iff}.
                $M$ is open in $P$ with respect to $T_P$ by \cref{finer_topology,open_in,subseteq}.
                $g\in M$ by \cref{subseteq,inter_intro}.
                Let $N = G\inter M$.
                $N$ is open in $G$ with respect to $\subspaceTopology{T_P}{G}$
                    by \cref{subspace_topology,open_in}.
                $N$ is a neighborhood of $g$ in $G$ with respect to $\subspaceTopology{T_P}{G}$
                    by \cref{neighborhood,inter_intro}.
                Take $\sigma$ such that
                    $\sigma$ is the uniform metric on functions from $K$ to $Y$ with respect to $d$ and
                    $H = \{h\in F \mid \sigma((\restrl{g}{K},\restrl{h}{K})) < \delta\}$
                    by \cref{compact_convergence_ball}.
                Show that for all $h$ such that $h\in N$ we have $h\in H\inter G$.
                \begin{subproof}
                    Fix $h\in N$.
                    $h\in M$ and $h\in G$ by \cref{inter}.
                    $h\in P$ and $h\in M_0$ by \cref{inter}.
                    $h\in F$ and $g\in F$ by \cref{subseteq_transitive}.
                    $\restrl{g}{K}\in\funs{K}{Y}$ and $\restrl{h}{K}\in\funs{K}{Y}$
                        by \cref{restrl_typed_function,compact_subspace,subseteq}.
                    Show that for all $x\in K$ we have
                        $d((\apply{h}{x},\apply{g}{x})) < \eta + \eta + \eta$.
                    \begin{subproof}
                        Fix $x\in K$.
                        $x\in\unions{L}$ by \cref{covering,subseteq}.
                        Take $V\in L$ such that $x\in V$ by \cref{unions_iff}.
                        $(V,\apply{e}{V})\in Q$ by \cref{choice_function,pair_eq_iff}.
                        $\apply{e}{V}\in K$ and
                            $V$ is a neighborhood of $\apply{e}{V}$ in $X$ with respect to $T$ and
                            for all $y,u$ such that $y\in V$ and $u\in G$ we have
                                $d((\apply{u}{y},\apply{u}{\apply{e}{V}})) < \eta$
                            by \cref{fst_eq,snd_eq,subseteq,times_tuple_elim}.
                        $e(V)\in J$ by \cref{function_apply_elim,img_iff}.
                        $d((\apply{h}{x},\apply{h}{e(V)})) < \eta$ and
                            $d((\apply{g}{x},\apply{g}{e(V)})) < \eta$ by \cref{subseteq}.
                        $h(e(V))\in\metricBall{d}{Y}{g(e(V))}{\eta}$ by \cref{subseteq}.
                        $d((\apply{g}{e(V)},\apply{h}{e(V)})) < \eta$ by \cref{metric_ball}.
                        $h(e(V))\in Y$ and $g(e(V))\in Y$ and $h(x)\in Y$ and $g(x)\in Y$
                            by \cref{funs_type_apply,subseteq}.
                        $d((\apply{h}{e(V)},\apply{g}{e(V)})) < \eta$ and
                            $d((\apply{g}{e(V)},\apply{g}{x})) < \eta$
                            by \cref{metric_on,metric_symmetric,real_lt_subst_left}.
                        $d((\apply{h}{x},\apply{g}{x})) < \eta + \eta + \eta$
                            by \cref{metric_three_step_lt_sum}.
                    \end{subproof}
                    $\eta + \eta + \eta\in\reals$ by \cref{real_add_closed}.
                    Show that for all $x\in K$ we have
                        $\apply{\boundedMetric{d}}{(\apply{\restrl{g}{K}}{x}, \apply{\restrl{h}{K}}{x})}
                            \leq \eta + \eta + \eta$.
                    \begin{subproof}
                        Fix $x\in K$.
                        $x\in X$ by \cref{subseteq}.
                        $d((\apply{h}{x},\apply{g}{x})) < \eta + \eta + \eta$ by \cref{subseteq}.
                        $h(x)\in Y$ and $g(x)\in Y$ by \cref{funs_type_apply}.
                        $d((h(x),g(x)))\in\reals$ by \cref{metric_on,funs_type_apply,times_tuple_intro}.
                        $1\in\reals$ by \cref{real_one_in_naturals,real_naturals_subset_reals,subseteq}.
                        $\eta + \eta + \eta < 1$ by \cref{real_lt_trans}.
                        $d((h(x),g(x))) < 1$ by \cref{real_lt_trans}.
                        $(g(x),h(x))\in Y\times Y$ by \cref{times_tuple_intro}.
                        $d((g(x),h(x))) < 1$ by \cref{metric_on,metric_symmetric,real_lt_subst_left}.
                        $\apply{\boundedMetric{d}}{(g(x),h(x))} = d((g(x),h(x)))$
                            by \cref{bounded_metric_apply_below_one}.
                        $d((g(x),h(x))) < \eta + \eta + \eta$
                            by \cref{metric_on,metric_symmetric,real_lt_subst_left}.
                        $\apply{\boundedMetric{d}}{(g(x),h(x))} < \eta + \eta + \eta$
                            by \cref{real_lt_subst_left}.
                        $\apply{\boundedMetric{d}}{(g(x),h(x))}
                            \leq \eta + \eta + \eta$ by \cref{real_lt_imp_leq}.
                        $\apply{\restrl{g}{K}}{x} = g(x)$ and $\apply{\restrl{h}{K}}{x} = h(x)$
                            by \cref{restrl_of_function_apply,compact_subspace,subseteq,funs_is_function,funs_is_relation,funs_elim}.
                        $\apply{\boundedMetric{d}}{(\apply{\restrl{g}{K}}{x}, \apply{\restrl{h}{K}}{x})}
                            = \apply{\boundedMetric{d}}{(g(x),h(x))}$ by \cref{pair_eq_iff}.
                        $\apply{\boundedMetric{d}}{(\apply{\restrl{g}{K}}{x}, \apply{\restrl{h}{K}}{x})}
                            \leq \eta + \eta + \eta$ by \cref{pair_eq_iff}.
                    \end{subproof}
                    $\sigma((\restrl{g}{K},\restrl{h}{K})) \leq \eta + \eta + \eta$
                        by \cref{uniform_metric_leq_from_coordinate_bounds}.
                    $(\restrl{g}{K},\restrl{h}{K})\in\funs{K}{Y}\times\funs{K}{Y}$
                        by \cref{times_tuple_intro}.
                    $\sigma\in\funs{\funs{K}{Y}\times\funs{K}{Y}}{\reals}$
                        by \cref{uniform_metric_on_function_space,metric_on}.
                    $\sigma((\restrl{g}{K},\restrl{h}{K})) < \delta$
                        by \cref{funs_type_apply,real_leq_lt_trans}.
                    $h\in H$ by \cref{subseteq}.
                    $h\in H\inter G$ by \cref{inter_intro}.
                \end{subproof}
                $N\subseteq H\inter G$ by \cref{subseteq}.
                For all $u$ such that $u\in H\inter G$ we have $u\in E\inter G$ by \cref{inter,subseteq,inter_intro}.
                $H\inter G\subseteq E\inter G$ by \cref{subseteq}.
                $N\subseteq D$ by \cref{pair_eq_iff,subseteq_transitive}.
                Follows by definition.
            \end{byCase}
        \end{subproof}
        $D$ is open in $G$ with respect to $\subspaceTopology{T_P}{G}$
            by \cref{pow_iff,open_from_local_neighborhoods}.
    \end{subproof}
    For all $W$ such that $W\in\subspaceTopology{T_0}{G}$ we have
        $W\in\subspaceTopology{T_P}{G}$
        by \cref{basis_topology_unions,basis_topology_is_topology,basis_for_topology,open_in,subseteq,topology_on}.
    $\subspaceTopology{T_0}{G}\subseteq\subspaceTopology{T_P}{G}$ by \cref{subseteq}.
    $\subspaceTopology{T_P}{G}$ is finer than $\subspaceTopology{T_0}{G}$ on $G$
        by \cref{finer_topology}.
\end{proof}

% from §47 helper lemma 13: pointwise value sets of functions lie in the codomain
\begin{lemma}\label{pointwise_value_set_subset_codomain_47}
    Assume $F\subseteq\funs{X}{Y}$.
    Assume $a\in X$.
    Assume $B = \{\apply{f}{a} \mid f\in F\}$.
    Then $B\subseteq Y$.
\end{lemma}

% from §47 Item 1: equicontinuity and pointwise compact closure imply compact containment
\begin{theorem}\label{ascoli_general_compact_containment}
    Assume $T$ is a topology on $X$.
    Assume $d$ is a metric on $Y$.
    Assume $C$ is the set of continuous functions from $X$ to $Y$ with respect to $T$ and $\metricTopology{d}{Y}$.
    Assume $T_0$ is the topology of compact convergence on functions from $X$ to $Y$ with respect to $T$ and $d$.
    Assume $F\subseteq C$.
    Suppose $F$ is equicontinuous from $X$ to $Y$ with respect to $T$ and $d$.
    Suppose $F$ has pointwise compact closure from $X$ to $Y$ with respect to $d$.
    Then there exists $G$ such that
        $F\subseteq G$ and
        $G\subseteq C$ and
        $G$ is compact with respect to $\subspaceTopology{T_0}{G}$.
\end{theorem}
\begin{proof}
    Let $U = \metricTopology{d}{Y}$.
    $C\subseteq\funs{X}{Y}$ by \cref{continuous_function_set,subseteq}.
    \begin{byCase}
    \caseOf{$X = \emptyset$.}
        Let $G = F$.
        Show that $F\subseteq G$ and $G\subseteq C$ and
            $G$ is compact with respect to $\subspaceTopology{T_0}{G}$.
        \begin{subproof}
            $F\subseteq G$ by \cref{subseteq_refl}.
            $G\subseteq C$ by \cref{subseteq_refl}.
            $T_0$ is a topology on $\funs{X}{Y}$ by \cref{compact_convergence_topology}.
            \begin{byCase}
            \caseOf{$G = \emptyset$.}
                $\subspaceTopology{T_0}{G}$ is a topology on $G$ by \cref{subspace_topology_is_topology,emptyset_subseteq}.
                $G$ is compact with respect to $\subspaceTopology{T_0}{G}$ by \cref{emptyset_compact}.
            \caseOf{$G \neq \emptyset$.}
                Take $f$ such that $f\in G$ by \cref{nonempty_inhabited}.
                For all $g$ such that $g\in G$ we have $g\in\{\emptyset\}$
                    by \cref{subseteq,continuous_function_set,funs_elim,function_from_emptyset_is_empty,singleton_intro}.
                $G\subseteq\{\emptyset\}$ by \cref{subseteq}.
                $\emptyset\in G$
                    by \cref{real_lt_subst_left,subseteq,continuous_function_set,funs_elim,function_from_emptyset_is_empty}.
                $G = \{\emptyset\}$ by \cref{subseteq_antisymmetric,singleton_subset_intro}.
                $\{\emptyset\}$ is a compact subspace of $\funs{X}{Y}$ with respect to $T_0$
                    by \cref{continuous_function_set,subseteq,real_lt_subst_left,singleton_compact_subspace,compact_convergence_topology}.
                $G$ is compact with respect to $\subspaceTopology{T_0}{G}$ by \cref{compact_subspace,real_lt_subst_left}.
            \end{byCase}
        \end{subproof}
        There exists $H$ such that
            $F\subseteq H$ and
            $H\subseteq C$ and
            $H$ is compact with respect to $\subspaceTopology{T_0}{H}$.
    \caseOf{$X \neq \emptyset$.}
        $X$ is inhabited by \cref{nonempty_inhabited}.
        Let $Q = \{w\in X\times\pow{Y} \mid \text{$\snd{w} = \{\apply{f}{\fst{w}} \mid f\in F\}$ and
            $\closure{\snd{w}}{Y}{U}$ is compact with respect to
                $\subspaceTopology{U}{\closure{\snd{w}}{Y}{U}}$}\}$.
        $Q$ is a relation by \cref{times_elem_is_tuple,relation}.
        $F\subseteq\funs{X}{Y}$ by \cref{pointwise_compact_closure_function_family}.
        For all $a\in X$ there exists $B$ such that $a\mathrel{Q} B$
            by \cref{pointwise_compact_closure_function_family,pointwise_value_set_subset_codomain_47,powerset_intro,times_tuple_intro,fst_eq,snd_eq}.
        Take $e$ such that $e$ is a function on $X$ and for all $a\in X$ we have $a\mathrel{Q}\apply{e}{a}$
            by \cref{choice_function}.
        Let $A = \{(a,\closure{e(a)}{Y}{U}) \mid a\in X\}$.
        For all $a,b_1,b_2$ such that $(a,b_1)\in A$ and $(a,b_2)\in A$ we have $b_1 = b_2$
            by \cref{pair_eq_iff}.
        For all $w\in A$ there exist $a,b$ such that $w = (a,b)$.
        $A$ is a function by \cref{relation,rightunique}.
        Show that $\dom{A} = X$ and for all $a\in X$ we have
            $\apply{A}{a} = \closure{e(a)}{Y}{U}$ and
            $\apply{A}{a}$ is compact with respect to $\subspaceTopology{U}{\apply{A}{a}}$ and
            $\apply{A}{a}\subseteq Y$.
        \begin{subproof}
            Follows by \cref{dom_iff,dom_intro,pair_eq_iff,subseteq,subseteq_antisymmetric,function_apply_intro,choice_function,fst_eq,snd_eq,closure}.
        \end{subproof}
        Let $T_A = \{(a,\subspaceTopology{U}{\apply{A}{a}}) \mid a\in X\}$.
        For all $a,b_1,b_2$ such that $(a,b_1)\in T_A$ and $(a,b_2)\in T_A$ we have $b_1 = b_2$
            by \cref{pair_eq_iff}.
        For all $w\in T_A$ there exist $a,b$ such that $w = (a,b)$.
        Show that $T_A$ is a function and $\dom{T_A} = X$.
        \begin{subproof}
            Follows by \cref{relation,rightunique,dom_iff,dom_intro,pair_eq_iff,subseteq,subseteq_antisymmetric}.
        \end{subproof}
        Let $P = \productFamily{A}$.
        Let $T_P = \productTopologyIndexed{T_A}{A}$.
        For all $f$ such that $f\in F$ we have
            $f\in\funs{X}{Y}$ and
            for all $a$ such that $a\in X$ we have $f(a)\in\apply{A}{a}$
            by \cref{subseteq,continuous_function_set,choice_function,pointwise_compact_closure_function_family,fst_eq,snd_eq,pointwise_value_set_subset_codomain_47,real_lt_subst_left,metric_basis_is_basis,basis_topology_is_topology,metric_topology,subseteq_closure,function_apply_intro}.
        For all $f$ such that $f\in F$ we have $f\in\funs{\dom{A}}{\unions{\ran{A}}}$
            by \cref{funs_is_function,funs,funs_intro,real_lt_subst_left,subseteq,function_apply_elim,ran_intro,unions_iff}.
        $F\subseteq P$ by \cref{indexed_product,subseteq}.
        For all $a$ such that $a\in\dom{A}$ we have
            $\apply{T_A}{a}$ is a topology on $\apply{A}{a}$ and
            $\apply{A}{a}$ is compact with respect to $\apply{T_A}{a}$
            by \cref{real_lt_subst_left,subseteq,function_apply_intro,subspace_topology_is_topology,metric_topology,metric_basis_is_basis,basis_topology_is_topology,basis_for_topology}.
        $P$ is compact with respect to $T_P$ by \cref{real_lt_subst_left,tychonoff_theorem}.
        $T_P$ is a topology on $P$ by
            \cref{product_subbasis_finite_intersections_basis,product_topology_indexed,basis_topology_is_topology,basis_for_topology}.
        Let $G = \closure{F}{P}{T_P}$.
        $F\subseteq G$ by \cref{subseteq_closure}.
        $G$ is closed in $P$ with respect to $T_P$ by \cref{closure_closed_in}.
        $G\subseteq P$ by \cref{real_lt_subst_left,closure,subseteq}.
        $G$ is compact with respect to $\subspaceTopology{T_P}{G}$ by \cref{closed_subspace_compact}.
        $G$ is equicontinuous from $X$ to $Y$ with respect to $T$ and $d$
            by \cref{coordinatewise_product_closure_equicontinuous_47}.
        $G\subseteq\funs{X}{Y}$
            by \cref{indexed_product_subset_function_space_47,closure,subseteq,subseteq_transitive}.
        $G\subseteq C$ by \cref{equicontinuous_family_continuous_functions_47}.
        $\subspaceTopology{T_P}{G}$ is finer than $\subspaceTopology{T_0}{G}$ on $G$
            by \cref{equicontinuous_product_topology_finer_than_compact_convergence_47}.
        $G$ is compact with respect to $\subspaceTopology{T_0}{G}$
            by \cref{compactness_from_finer_topology_47}.
        There exists $K$ such that
            $F\subseteq K$ and
            $K\subseteq C$ and
            $K$ is compact with respect to $\subspaceTopology{T_0}{K}$.
    \end{byCase}
\end{proof}

% from §47 helper lemma 14: restriction of a restricted typed function equals the smaller restriction
\begin{lemma}\label{restrl_restrl_typed_function_47}
    Assume $f\in\funs{X}{Y}$.
    Assume $C\subseteq A\subseteq X$.
    Then $\restrl{\restrl{f}{A}}{C} = \restrl{f}{C}$.
\end{lemma}

% from §47 helper lemma 15: naming a restriction preserves its smaller restrictions
\begin{lemma}\label{named_restrl_restrl_typed_function_47}
    Assume $p = \restrl{f}{A}$.
    Assume $f\in\funs{X}{Y}$.
    Assume $C\subseteq A\subseteq X$.
    Then $\restrl{p}{C} = \restrl{f}{C}$.
\end{lemma}

% from §47 helper lemma 15a: the graph of the restriction map is a typed function
\begin{lemma}\label{restriction_graph_typed_function_47}
    Assume $A\subseteq X$.
    Let $F = \funs{X}{Y}$.
    Let $F_A = \funs{A}{Y}$.
    Let $r = \{(f,\restrl{f}{A}) \mid f\in F\}$.
    Then $r\in\funs{F}{F_A}$ and $\dom{r} = F$ and for all $f\in F$ we have $r(f) = \restrl{f}{A}$.
\end{lemma}
\begin{proof}
    For all $f,g_1,g_2$ such that $(f,g_1)\in r$ and $(f,g_2)\in r$ we have $g_1 = g_2$
        by \cref{pair_eq_iff}.
    For all $w\in r$ there exist $f,g$ such that $w = (f,g)$.
    $r$ is a function by \cref{relation,rightunique}.
    Show that $r\in\funs{F}{F_A}$ and $\dom{r} = F$ and for all $f\in F$ we have $r(f) = \restrl{f}{A}$.
    \begin{subproof}
        Follows by \cref{function_apply_intro,restrl_typed_function,subseteq,real_lt_subst_left,funs_intro,dom_iff,pair_eq_iff,dom_intro,subseteq_antisymmetric,funs}.
    \end{subproof}
\end{proof}

% from §47 helper lemma 16: restriction maps are continuous for compact-convergence topologies on subspaces
\begin{lemma}\label{restriction_map_continuous_compact_convergence_47}
    Assume $T$ is a topology on $X$.
    Assume $A$ is a subspace of $X$ with respect to $T$.
    Let $T_A = \subspaceTopology{T}{A}$.
    Assume $d$ is a metric on $Y$.
    Assume $T_0$ is the topology of compact convergence on functions from $X$ to $Y$ with respect to $T$ and $d$.
    Assume $T_1$ is the topology of compact convergence on functions from $A$ to $Y$ with respect to $T_A$ and $d$.
    Let $r = \{(f,\restrl{f}{A}) \mid f\in\funs{X}{Y}\}$.
    Then $r$ is continuous from $\funs{X}{Y}$ to $\funs{A}{Y}$ with respect to $T_0$ and $T_1$.
\end{lemma}
\begin{proof}
    Let $F = \funs{X}{Y}$.
    Let $F_A = \funs{A}{Y}$.
    Show that $r\in\funs{F}{F_A}$ and $\dom{r} = F$ and for all $f\in F$ we have $r(f) = \restrl{f}{A}$.
    \begin{subproof}
        Follows by \cref{restriction_graph_typed_function_47,subspace_of}.
    \end{subproof}
    Show that $r$ is right-unique and $r$ is a relation.
    \begin{subproof}
        Follows by \cref{funs,funs_is_function,function_rightunique,funs_is_relation}.
    \end{subproof}
    Show that $T_0$ is a topology on $F$ and $T_1$ is a topology on $F_A$.
    \begin{subproof}
        Follows by \cref{compact_convergence_topology}.
    \end{subproof}
    Show that for all $U$ such that $U$ is open in $F_A$ with respect to $T_1$ we have
        $\preimg{r}{U}$ is open in $F$ with respect to $T_0$.
    \begin{subproof}
        Fix $U$ such that $U$ is open in $F_A$ with respect to $T_1$.
        Take $B$ such that
            $B$ is the compact-convergence basis on functions from $A$ to $Y$ with respect to $T_A$ and $d$ and
            $T_1 = \basisTopology{B}{F_A}$ by \cref{compact_convergence_topology}.
        $U\in T_1$ by \cref{open_in}.
        $\preimg{r}{U}\subseteq F$ by \cref{preimg_subseteq_dom,subseteq}.
        Show that for all $f\in\preimg{r}{U}$ there exists $E$ such that
            $E$ is a neighborhood of $f$ in $F$ with respect to $T_0$ and
            $E\subseteq \preimg{r}{U}$.
        \begin{subproof}
            Fix $f\in\preimg{r}{U}$.
            $f\in F$ by \cref{subseteq}.
            Take $u$ such that $u\in U$ and $(f,u)\in r$ by \cref{preimg_iff}.
            Take $V\in B$ such that $u\in V$ and $V\subseteq U$
                by \cref{topology_generated_by_basis}.
            Take $h, C, \epsilon$ such that
                $V$ is the compact-convergence ball about $h$ over $C$ with radius $\epsilon$
                    from $A$ to $Y$ with respect to $T_A$ and $d$
                by \cref{compact_convergence_function_space_basis}.
            $r(f)\in V$ by \cref{function_apply_intro,real_lt_subst_left}.
            $r(f) = \restrl{f}{A}$ by \cref{subseteq}.
            Take $D,\delta$ such that
                $\delta\in\reals$ and
                $0 < \delta$ and
                $D$ is the compact-convergence ball about $r(f)$ over $C$ with radius $\delta$
                    from $A$ to $Y$ with respect to $T_A$ and $d$ and
                $D\subseteq V$ by \cref{compact_convergence_ball_refinement}.
            Take $\rho$ such that
                $\rho$ is the uniform metric on functions from $C$ to $Y$ with respect to $d$ and
                $D = \{k\in F_A \mid \rho((\restrl{r(f)}{C},\restrl{k}{C})) < \delta\}$
                by \cref{compact_convergence_ball}.
            $C$ is a compact subspace of $A$ with respect to $T_A$ by \cref{compact_convergence_ball}.
            $\subspaceTopology{T_A}{C} = \subspaceTopology{T}{C}$
                by \cref{subspace_subspace_topology,compact_subspace,subspace_of,subseteq}.
            $C$ is compact with respect to $\subspaceTopology{T}{C}$ by \cref{compact_subspace,real_lt_subst_right}.
            $C$ is a compact subspace of $X$ with respect to $T$
                by \cref{compact_subspace,subspace_of,subseteq_transitive,real_lt_subst_right}.
            Let $E = \{g\in F \mid \rho((\restrl{f}{C},\restrl{g}{C})) < \delta\}$.
            $E$ is the compact-convergence ball about $f$ over $C$ with radius $\delta$
                from $X$ to $Y$ with respect to $T$ and $d$ by \cref{compact_convergence_ball}.
            $E$ is open in $F$ with respect to $T_0$ by \cref{compact_convergence_ball_open}.
            For all $g$ such that $g\in E$ we have $g\in\preimg{r}{U}$
                by \cref{compact_convergence_ball,subseteq,funs_type_apply,compact_subspace,subspace_of,named_restrl_restrl_typed_function_47,real_lt_subst_left,real_lt_subst_right,funs,function_apply_elim,preimg_iff}.
            There exists $H$ such that
                $H$ is a neighborhood of $f$ in $F$ with respect to $T_0$ and
                $H\subseteq \preimg{r}{U}$
                by \cref{subseteq,neighborhood,compact_convergence_ball_open,compact_convergence_ball_center}.
        \end{subproof}
        $\preimg{r}{U}$ is open in $F$ with respect to $T_0$ by \cref{open_from_local_neighborhoods}.
    \end{subproof}
    $r$ is continuous from $F$ to $F_A$ with respect to $T_0$ and $T_1$ by \cref{continuous}.
\end{proof}

% from §47 helper lemma 17: restriction maps into uniform metric function spaces are continuous on compact subspaces
\begin{lemma}\label{restriction_map_continuous_uniform_metric_47}
    Assume $T$ is a topology on $X$.
    Assume $A$ is a compact subspace of $X$ with respect to $T$.
    Assume $d$ is a metric on $Y$.
    Assume $\rho$ is the uniform metric on functions from $A$ to $Y$ with respect to $d$.
    Assume $T_0$ is the topology of compact convergence on functions from $X$ to $Y$ with respect to $T$ and $d$.
    Let $r = \{(f,\restrl{f}{A}) \mid f\in\funs{X}{Y}\}$.
    Then $r$ is continuous from $\funs{X}{Y}$ to $\funs{A}{Y}$ with respect to $T_0$ and
        $\metricTopology{\rho}{\funs{A}{Y}}$.
\end{lemma}
\begin{proof}
    Let $F = \funs{X}{Y}$.
    Let $F_A = \funs{A}{Y}$.
    Let $T_1 = \metricTopology{\rho}{F_A}$.
    Show that $r\in\funs{F}{F_A}$ and $\dom{r} = F$ and for all $f\in F$ we have $r(f) = \restrl{f}{A}$.
    \begin{subproof}
        Follows by \cref{restriction_graph_typed_function_47,compact_subspace}.
    \end{subproof}
    Show that $r$ is right-unique and $r$ is a relation.
    \begin{subproof}
        Follows by \cref{funs,funs_is_function,function_rightunique,funs_is_relation}.
    \end{subproof}
    Show that $T_0$ is a topology on $F$ and $T_1$ is a topology on $F_A$.
    \begin{subproof}
        Follows by \cref{compact_convergence_topology,metric_topology,metric_basis_is_basis,basis_topology_is_topology,basis_for_topology,uniform_metric_on_function_space}.
    \end{subproof}
    Show that for all $U$ such that $U$ is open in $F_A$ with respect to $T_1$ we have
        $\preimg{r}{U}$ is open in $F$ with respect to $T_0$.
    \begin{subproof}
        Fix $U$ such that $U$ is open in $F_A$ with respect to $T_1$.
        $U\subseteq F_A$ by \cref{open_in,topology_on,subseteq,pow_iff}.
        $\preimg{r}{U}\subseteq F$ by \cref{preimg_subseteq_dom,subseteq}.
        Show that for all $f\in\preimg{r}{U}$ there exists $E$ such that
            $E$ is a neighborhood of $f$ in $F$ with respect to $T_0$ and
            $E\subseteq \preimg{r}{U}$.
        \begin{subproof}
            Fix $f\in\preimg{r}{U}$.
            $f\in F$ by \cref{subseteq}.
            Take $u$ such that $u\in U$ and $(f,u)\in r$ by \cref{preimg_iff}.
            $r(f)\in U$ by \cref{subseteq,function_apply_intro,real_lt_subst_left}.
            Take $\epsilon\in\reals$ such that
                $0 < \epsilon$ and
                $\metricBall{\rho}{F_A}{r(f)}{\epsilon}\subseteq U$
                by \cref{metric_open_iff,uniform_metric_on_function_space}.
            Let $E = \{g\in F \mid \rho((\restrl{f}{A},\restrl{g}{A})) < \epsilon\}$.
            $E$ is the compact-convergence ball about $f$ over $A$ with radius $\epsilon$
                from $X$ to $Y$ with respect to $T$ and $d$ by \cref{compact_convergence_ball}.
            $E$ is open in $F$ with respect to $T_0$ by \cref{compact_convergence_ball_open}.
            For all $g$ such that $g\in E$ we have $g\in\preimg{r}{U}$
                by \cref{compact_convergence_ball,subseteq,funs_type_apply,real_lt_subst_left,real_lt_subst_right,metric_ball,funs,function_apply_elim,preimg_iff}.
            There exists $H$ such that
                $H$ is a neighborhood of $f$ in $F$ with respect to $T_0$ and
                $H\subseteq \preimg{r}{U}$
                by \cref{subseteq,neighborhood,compact_convergence_ball_open,compact_convergence_ball_center}.
        \end{subproof}
        $\preimg{r}{U}$ is open in $F$ with respect to $T_0$ by \cref{open_from_local_neighborhoods}.
    \end{subproof}
    $r$ is continuous from $F$ to $F_A$ with respect to $T_0$ and $T_1$ by \cref{continuous}.
\end{proof}

% from §47 Item 1: compact containment implies equicontinuity for locally compact Hausdorff domains
\begin{theorem}\label{ascoli_general_compact_containment_implies_equicontinuous}
    Assume $T$ is a topology on $X$.
    Assume $X$ is locally compact with respect to $T$.
    Assume $X$ is Hausdorff with respect to $T$.
    Assume $d$ is a metric on $Y$.
    Assume $C$ is the set of continuous functions from $X$ to $Y$ with respect to $T$ and $\metricTopology{d}{Y}$.
    Assume $T_0$ is the topology of compact convergence on functions from $X$ to $Y$ with respect to $T$ and $d$.
    Assume $F\subseteq C$.
    Suppose there exists $G$ such that
        $F\subseteq G$ and
        $G\subseteq C$ and
        $G$ is compact with respect to $\subspaceTopology{T_0}{G}$.
    Then $F$ is equicontinuous from $X$ to $Y$ with respect to $T$ and $d$.
\end{theorem}
\begin{proof}
    Let $U = \metricTopology{d}{Y}$.
    Let $F_X = \funs{X}{Y}$.
    Take $G$ such that
        $F\subseteq G$ and
        $G\subseteq C$ and
        $G$ is compact with respect to $\subspaceTopology{T_0}{G}$.
    $F\subseteq\funs{X}{Y}$ by \cref{continuous_function_set,subseteq,subseteq_transitive}.
    Show that for all $a\in X$ we have
        $G$ is equicontinuous at $a$ from $X$ to $Y$ with respect to $T$ and $d$.
    \begin{subproof}
        Fix $a\in X$.
        $X$ is a neighborhood of $a$ in $X$ with respect to $T$ by \cref{topology_on,open_in,neighborhood}.
        Take $V$ such that
            $V$ is a neighborhood of $a$ in $X$ with respect to $T$ and
            $\closure{V}{X}{T}$ is a compact subspace of $X$ with respect to $T$ and
            $\closure{V}{X}{T}\subseteq X$
            by \cref{locally_compact_neighborhood_characterization}.
        Let $A = \closure{V}{X}{T}$.
        Let $T_A = \subspaceTopology{T}{A}$.
        $A$ is a subspace of $X$ with respect to $T$ by \cref{compact_subspace,subspace_of}.
        $V$ is open in $X$ with respect to $T$ and $a\in V$ by \cref{neighborhood}.
        $V\in T$ by \cref{open_in}.
        $V\subseteq A$ by \cref{topology_on,pow_iff,subseteq,subseteq_closure,subspace_of}.
        $a\in A$ by \cref{neighborhood,subseteq}.
        $A$ is inhabited by \cref{nonempty_inhabited}.
        Let $C_A = \{h\in\funs{A}{Y} \mid \text{$h$ is continuous from $A$ to $Y$ with respect to $T_A$ and $U$}\}$.
        $T_A$ is a topology on $A$ by \cref{subspace_of,subspace_topology_is_topology}.
        $C_A$ is the set of continuous functions from $A$ to $Y$ with respect to $T_A$ and $U$
            by \cref{continuous_function_set}.
        Take $\rho$ such that
            $\rho$ is the uniform metric on functions from $A$ to $Y$ with respect to $d$
            by \cref{uniform_metric_on_function_space_exists_inhabited}.
        Let $T_1 = \metricTopology{\rho}{\funs{A}{Y}}$.
        Let $q = \{(f,\restrl{f}{A}) \mid f\in F_X\}$.
        $q$ is continuous from $F_X$ to $\funs{A}{Y}$ with respect to $T_0$ and $T_1$
            by \cref{restriction_map_continuous_uniform_metric_47}.
        $q\in\funs{F_X}{\funs{A}{Y}}$ by \cref{continuous}.
        $q$ is right-unique and $q$ is a relation
            by \cref{funs,funs_is_function,function_rightunique,funs_is_relation}.
        $G\subseteq F_X$ by \cref{continuous_function_set,subseteq,subseteq_transitive}.
        $T_0$ is a topology on $F_X$ by \cref{compact_convergence_topology}.
        $G$ is a subspace of $F_X$ with respect to $T_0$ by \cref{subspace_of}.
        Let $r = \restrl{q}{G}$.
        $r$ is continuous from $G$ to $\funs{A}{Y}$ with respect to $\subspaceTopology{T_0}{G}$ and $T_1$
            by \cref{continuous_restrict_domain}.
        Let $R = \img{r}{G}$.
        For all $h$ such that $h\in R$ we have $h\in C_A$
            by \cref{img_iff,subseteq,continuous_function_set,restrl_iff,function_apply_intro,continuous_restrict_domain,funs_type_apply,real_lt_subst_left,real_lt_subst_right}.
        $R\subseteq C_A$ by \cref{subseteq}.
        $R$ is compact with respect to $\subspaceTopology{T_1}{R}$ by \cref{continuous_image_compact}.
        $R$ is compact with respect to $\subspaceTopology{\metricTopology{\rho}{\funs{A}{Y}}}{R}$
            by \cref{real_lt_subst_left}.
        $\rho$ is a metric on $\funs{A}{Y}$ by \cref{uniform_metric_on_function_space}.
        $R\subseteq \funs{A}{Y}$ by \cref{continuous_function_set,subseteq}.
        $R$ is totally bounded in $\funs{A}{Y}$ with respect to $\rho$
            by \cref{compact_metric_subspace_ambient_totally_bounded}.
        $R$ is equicontinuous from $A$ to $Y$ with respect to $T_A$ and $d$
            by \cref{totally_bounded_function_family_equicontinuous}.
        $R$ is equicontinuous at $a$ from $A$ to $Y$ with respect to $T_A$ and $d$
            by \cref{equicontinuous_family}.
        Show that for all $\epsilon\in\reals$ such that $0 < \epsilon$
            there exists $W$ such that
                $W$ is a neighborhood of $a$ in $X$ with respect to $T$ and
                for all $x, g$ such that $x\in W$ and $g\in G$ we have
                    $d((\apply{g}{x},\apply{g}{a})) < \epsilon$.
        \begin{subproof}
            Fix $\epsilon$ such that $\epsilon\in\reals$ and $0 < \epsilon$.
            Take $B$ such that
                $B$ is a neighborhood of $a$ in $A$ with respect to $T_A$ and
                for all $x, h$ such that $x\in B$ and $h\in R$ we have
                    $d((\apply{h}{x},\apply{h}{a})) < \epsilon$
                by \cref{equicontinuous_at_point}.
            $B$ is open in $A$ with respect to $T_A$ and $a\in B$ by \cref{neighborhood}.
            Take $W_0\in T$ such that $B = A\inter W_0$ by \cref{subspace_topology,open_in,subspace_of}.
            Let $W = V\inter W_0$.
            $W\in T$ by \cref{topology_on}.
            $W$ is open in $X$ with respect to $T$ by \cref{open_in}.
            $a\in W$ by \cref{neighborhood,real_lt_subst_right,inter,inter_intro}.
            Show that for all $x, g$ such that $x\in W$ and $g\in G$ we have
                $d((\apply{g}{x},\apply{g}{a})) < \epsilon$.
            \begin{subproof}
                Fix $x$.
                Fix $g$.
                Assume $x\in W$ and $g\in G$.
                $x\in V$ and $x\in W_0$ by \cref{inter}.
                $x\in B$ by \cref{inter_intro,subseteq,real_lt_subst_left}.
                $(g,\restrl{g}{A})\in q$ by \cref{subseteq}.
                $(g,\restrl{g}{A})\in r$ by \cref{restrl_iff,subseteq}.
                $\restrl{g}{A}\in R$ by \cref{img_iff,function_apply_intro}.
                $d((\apply{\restrl{g}{A}}{x},\apply{\restrl{g}{A}}{a})) < \epsilon$ by \cref{subseteq}.
                $g\in\funs{X}{Y}$ by \cref{continuous_function_set,subseteq}.
                $\dom{g} = X$ by \cref{funs}.
                $g$ is right-unique and $g$ is a relation
                    by \cref{funs,funs_is_function,function_rightunique,funs_is_relation}.
                $A\subseteq\dom{g}$ by \cref{subspace_of,real_lt_subst_right}.
                $\apply{\restrl{g}{A}}{x} = g(x)$ by \cref{restrl_of_function_apply,subseteq}.
                $\apply{\restrl{g}{A}}{a} = g(a)$ by \cref{restrl_of_function_apply,subseteq}.
                $d((\apply{g}{x},\apply{g}{a})) < \epsilon$ by \cref{real_lt_subst_left,real_lt_subst_right}.
            \end{subproof}
            There exists $Z$ such that
                $Z$ is a neighborhood of $a$ in $X$ with respect to $T$ and
                for all $x, g$ such that $x\in Z$ and $g\in G$ we have
                    $d((\apply{g}{x},\apply{g}{a})) < \epsilon$
                by \cref{neighborhood}.
        \end{subproof}
        $G$ is equicontinuous at $a$ from $X$ to $Y$ with respect to $T$ and $d$
            by \cref{equicontinuous_at_point}.
    \end{subproof}
    $G\subseteq\funs{X}{Y}$ by \cref{continuous_function_set,subseteq}.
    $G$ is equicontinuous from $X$ to $Y$ with respect to $T$ and $d$
        by \cref{equicontinuous_family}.
    For all $a\in X$ we have
        $F$ is equicontinuous at $a$ from $X$ to $Y$ with respect to $T$ and $d$
        by \cref{equicontinuous_family,equicontinuous_at_point,subseteq}.
    $F$ is equicontinuous from $X$ to $Y$ with respect to $T$ and $d$
        by \cref{equicontinuous_family}.
\end{proof}

% from §47 Item 1: compact containment implies pointwise compact closure for locally compact Hausdorff domains
\begin{theorem}\label{ascoli_general_compact_containment_implies_pointwise_compact_closure}
    Assume $T$ is a topology on $X$.
    Assume $X$ is locally compact with respect to $T$.
    Assume $X$ is Hausdorff with respect to $T$.
    Assume $d$ is a metric on $Y$.
    Assume $C$ is the set of continuous functions from $X$ to $Y$ with respect to $T$ and $\metricTopology{d}{Y}$.
    Assume $T_0$ is the topology of compact convergence on functions from $X$ to $Y$ with respect to $T$ and $d$.
    Assume $F\subseteq C$.
    Suppose there exists $G$ such that
        $F\subseteq G$ and
        $G\subseteq C$ and
        $G$ is compact with respect to $\subspaceTopology{T_0}{G}$.
    Then $F$ has pointwise compact closure from $X$ to $Y$ with respect to $d$.
\end{theorem}
\begin{proof}
    Let $U = \metricTopology{d}{Y}$.
    $C\subseteq\funs{X}{Y}$ by \cref{continuous_function_set,subseteq}.
    Take $G$ such that
        $F\subseteq G$ and
        $G\subseteq C$ and
        $G$ is compact with respect to $\subspaceTopology{T_0}{G}$.
    Show that for all $a\in X$ there exists $A$ such that
        $A = \{\apply{f}{a} \mid f\in F\}$ and
        $\closure{A}{Y}{U}$ is compact with respect to
            $\subspaceTopology{U}{\closure{A}{Y}{U}}$.
    \begin{subproof}
            Fix $a$ such that $a\in X$.
                Let $R = \{(x,Y) \mid x\in X\}$.
                Let $S = \{(x,U) \mid x\in X\}$.
                Let $T_1 = \productTopologyIndexed{S}{R}$.
                $U$ is a topology on $Y$
                    by \cref{metric_topology,metric_basis_is_basis,basis_topology_is_topology,basis_for_topology}.
                $T_1$ is the topology of pointwise convergence on functions from $X$ to $Y$ with respect to $U$
                    by \cref{constant_indexed_product_topology_is_pointwise_47,nonempty_inhabited,emptyset}.
                $T_0$ is finer than $T_1$ on $\funs{X}{Y}$ by \cref{compact_convergence_finer_than_pointwise}.
                $G\subseteq\funs{X}{Y}$ by \cref{subseteq_transitive,subseteq}.
                $\subspaceTopology{T_0}{G}$ is finer than $\subspaceTopology{T_1}{G}$ on $G$
                    by \cref{subspace_finer_from_ambient_finer_47}.
                $G$ is compact with respect to $\subspaceTopology{T_1}{G}$
                    by \cref{compactness_from_finer_topology_47}.
                Show that $R$ is a function and $\dom{R} = X$ and for all $x\in X$ we have $\apply{R}{x} = Y$.
                \begin{subproof}
                    Follows by \cref{constant_indexed_family_47}.
                \end{subproof}
                Show that $S$ is a function and $\dom{S} = X$ and for all $x\in X$ we have $\apply{S}{x} = U$.
                \begin{subproof}
                    Follows by \cref{constant_indexed_family_47}.
                \end{subproof}
                $\piIndex{R}{a}$ is continuous from $\productFamily{R}$ to $\apply{R}{a}$
                    with respect to $T_1$ and $\apply{S}{a}$ by \cref{projection_map_indexed_continuous,function_apply_intro}.
                $\piIndex{R}{a}$ is continuous from $\funs{X}{Y}$ to $Y$ with respect to $T_1$ and $U$
                    by \cref{constant_indexed_product_function_space_47,function_apply_intro}.
                Show that $\subspaceTopology{T_1}{G}$ is a topology on $G$ and $G$ is a subspace of $\funs{X}{Y}$ with respect to $T_1$.
                \begin{subproof}
                    Follows by \cref{subspace_topology_is_topology,subspace_of,pointwise_convergence_topology}.
                \end{subproof}
                Let $p = \restrl{\piIndex{R}{a}}{G}$.
                $p$ is continuous from $G$ to $Y$ with respect to $\subspaceTopology{T_1}{G}$ and $U$
                    by \cref{continuous_restrict_domain}.
                Let $B = \img{p}{G}$.
                $B$ is compact with respect to $\subspaceTopology{U}{B}$ by \cref{continuous_image_compact}.
                $B\subseteq Y$ by \cref{continuous,funs,rels_ran_subseteq,img_subseteq_ran,subseteq_transitive}.
                $B$ is closed in $Y$ with respect to $U$ by \cref{compact_subspace_closed,metric_space_hausdorff}.
                Let $A = \{\apply{f}{a} \mid f\in F\}$.
                Show that $\piIndex{R}{a}$ is a function and $\dom{\piIndex{R}{a}} = \productFamily{R}$.
                \begin{subproof}
                    Follows by \cref{projection_map_indexed_function,funs_is_function,funs}.
                \end{subproof}
                $G\subseteq \dom{\piIndex{R}{a}}$ by \cref{constant_indexed_product_function_space_47,subseteq}.
                For all $z$ such that $z\in A$ we have $z\in B$
                    by \cref{pair_eq_iff,subseteq,constant_indexed_product_function_space_47,projection_map_indexed,function_apply_elim,restrl_iff,img_iff}.
                $A\subseteq B$ by \cref{subseteq}.
                $A\subseteq Y$ by \cref{subseteq_transitive}.
                $\closure{A}{Y}{U}\subseteq B$ by \cref{subseteq_transitive,closure_subseteq_closed}.
                $\closure{A}{Y}{U}$ is closed in $Y$ with respect to $U$ by \cref{closure_closed_in}.
                $B$ is a subspace of $Y$ with respect to $U$
                    by \cref{subspace_of,continuous,funs,rels_ran_subseteq,img_subseteq_ran,subseteq_transitive}.
                $\closure{A}{Y}{U}\inter B = \closure{A}{Y}{U}$
                    by \cref{inter_absorb_supseteq_right,subseteq}.
                There exists $D$ such that
                    $D$ is closed in $Y$ with respect to $U$ and
                    $\closure{A}{Y}{U} = D\inter B$.
                $\closure{A}{Y}{U}$ is closed in $B$ with respect to $\subspaceTopology{U}{B}$
                    by \cref{closed_in_subspace_iff}.
                $\closure{A}{Y}{U}$ is compact with respect to
                    $\subspaceTopology{U}{\closure{A}{Y}{U}}$
                    by \cref{closed_in_subspace_iff,closed_subspace_compact,subspace_subspace_topology,subseteq,subspace_topology_is_topology,subspace_of}.
                There exists $E$ such that
                    $E = \{\apply{f}{a} \mid f\in F\}$ and
                    $\closure{E}{Y}{U}$ is compact with respect to
                        $\subspaceTopology{U}{\closure{E}{Y}{U}}$.
    \end{subproof}
    $F$ has pointwise compact closure from $X$ to $Y$ with respect to $d$
        by \cref{pointwise_compact_closure_function_family,continuous_function_set,subseteq,subseteq_transitive}.
\end{proof}

% from §47 helper lemma 18: compactly contained continuous families are equicontinuous
\begin{lemma}\label{compact_contained_family_equicontinuous_47}
    Assume $T$ is a topology on $X$.
    Assume $X$ is locally compact with respect to $T$.
    Assume $X$ is Hausdorff with respect to $T$.
    Assume $d$ is a metric on $Y$.
    Assume $C$ is the set of continuous functions from $X$ to $Y$ with respect to $T$ and $\metricTopology{d}{Y}$.
    Assume $T_0$ is the topology of compact convergence on functions from $X$ to $Y$ with respect to $T$ and $d$.
    Assume $G\subseteq C$.
    Suppose $G$ is compact with respect to $\subspaceTopology{T_0}{G}$.
    Then $G$ is equicontinuous from $X$ to $Y$ with respect to $T$ and $d$.
\end{lemma}
\begin{proof}
    Follows by \cref{ascoli_general_compact_containment_implies_equicontinuous,subseteq_refl,real_lt_subst_left,real_lt_subst_right}.
\end{proof}
\section{§48 Baire Spaces}

% from §48 Item 1: empty interior
\begin{definition}\label{empty_interior}
    $A$ has empty interior in $X$ with respect to $T$ iff
        $A\subseteq X$ and
        $\interior{A}{X}{T} = \emptyset$.
\end{definition}

% from §48 Item 1: empty interior is equivalent to dense complement
\begin{lemma}\label{empty_interior_iff_dense_complement}
    Assume $T$ is a topology on $X$.
    Assume $A\subseteq X$.
    Then $A$ has empty interior in $X$ with respect to $T$ iff
        $X\setminus A$ is dense in $X$ with respect to $T$.
\end{lemma}
\begin{proof}
    Show that if $A$ has empty interior in $X$ with respect to $T$ then
        $X\setminus A$ is dense in $X$ with respect to $T$.
    \begin{subproof}
        Assume $A$ has empty interior in $X$ with respect to $T$.
        Show that for all $x\in X$ we have $x\in\closure{X\setminus A}{X}{T}$.
        \begin{subproof}
            Fix $x\in X$.
            Show that for all $U$ such that $U$ is open in $X$ with respect to $T$ and $x\in U$
                we have $U$ intersects $X\setminus A$.
            \begin{subproof}
                Fix $U$ such that $U$ is open in $X$ with respect to $T$ and $x\in U$.
                Suppose not.
                $U\inter (X\setminus A) = \emptyset$ by \cref{intersects}.
                $U\subseteq A$ by \cref{open_in,topology_on,pow_iff,elem_subseteq,subseteq_setminus,double_relative_complement,subseteq,subseteq_transitive}.
                $x\in\interior{A}{X}{T}$ by \cref{open_in,interior}.
                Contradiction by \cref{empty_interior,emptyset}.
            \end{subproof}
            $x\in\closure{X\setminus A}{X}{T}$ by \cref{closure_iff_intersects_open,setminus_subseteq}.
        \end{subproof}
        $\closure{X\setminus A}{X}{T} = X$
            by \cref{closure_closed_in,closed_in,setminus_subseteq,subseteq,subseteq_antisymmetric}.
        $X\setminus A$ is dense in $X$ with respect to $T$ by \cref{dense_in}.
    \end{subproof}
    Show that if $X\setminus A$ is dense in $X$ with respect to $T$ then
        $A$ has empty interior in $X$ with respect to $T$.
    \begin{subproof}
        Assume $X\setminus A$ is dense in $X$ with respect to $T$.
        Show that $\interior{A}{X}{T} = \emptyset$.
        \begin{subproof}
            Suppose not.
            Take $x$ such that $x\in\interior{A}{X}{T}$ by \cref{emptyset,nonempty_inhabited}.
            Take $U\in T$ such that $x\in U$ and $U\subseteq A$ by \cref{interior}.
            $\closure{X\setminus A}{X}{T} = X$ by \cref{dense_in}.
            $x\in\closure{X\setminus A}{X}{T}$ by \cref{subseteq}.
            $U$ intersects $X\setminus A$
                by \cref{closure_iff_intersects_open,setminus_subseteq,open_in,topology_on}.
            Take $y$ such that $y\in U\inter (X\setminus A)$ by \cref{intersects,nonempty_inhabited}.
            $y\in A$ and $y\notin A$ by \cref{inter,setminus_elim_right,subseteq}.
            Contradiction.
        \end{subproof}
        $A$ has empty interior in $X$ with respect to $T$ by \cref{empty_interior}.
    \end{subproof}
\end{proof}

% from §48 helper lemma 1: supersets of dense subsets are dense
\begin{lemma}\label{dense_subseteq_superset}
    Assume $T$ is a topology on $X$.
    Assume $A\subseteq B$.
    Assume $B\subseteq X$.
    Suppose $A$ is dense in $X$ with respect to $T$.
    Then $B$ is dense in $X$ with respect to $T$.
\end{lemma}
\begin{proof}
    We have $\closure{A}{X}{T}\subseteq\closure{B}{X}{T}$ by \cref{subseteq_transitive,subseteq_closure,closure_subseteq_closed,closure_closed_in}.
    $B$ is dense in $X$ with respect to $T$ by \cref{dense_in,real_lt_subst_left,closure_closed_in,closed_in,subseteq,subseteq_antisymmetric}.
\end{proof}

% from §48 helper lemma 2: subsets of empty-interior sets have empty interior
\begin{lemma}\label{empty_interior_subset}
    Assume $T$ is a topology on $X$.
    Assume $A\subseteq B$.
    Assume $B\subseteq X$.
    Suppose $B$ has empty interior in $X$ with respect to $T$.
    Then $A$ has empty interior in $X$ with respect to $T$.
\end{lemma}
\begin{proof}
    We have $A$ has empty interior in $X$ with respect to $T$
        by \cref{empty_interior_iff_dense_complement,dense_subseteq_superset,setminus_subseteq,subseteq_transitive,subseteq_implies_setminus_supseteq,subseteq}.
\end{proof}

% from §48 helper lemma 3: complements of a powerset-valued sequence form a powerset-valued sequence
\begin{lemma}\label{complement_sequence_in_powerset_48}
    Assume $s$ is a sequence in $\pow{X}$.
    Let $t = \{(n, X\setminus s(n)) \mid n\in\naturalsPlus\}$.
    Then $t$ is a sequence in $\pow{X}$.
\end{lemma}
\begin{proof}
    For all $w\in t$ there exist $n,U$ such that $w = (n,U)$ by \cref{pair_eq_iff}.
    For all $n,U_1,U_2$ such that $(n,U_1)\in t$ and $(n,U_2)\in t$ we have $U_1 = U_2$
        by \cref{pair_eq_iff}.
    $t$ is a function by \cref{relation,rightunique}.
    For all $n\in\dom{t}$ we have $t(n)\in\pow{X}$
        by \cref{dom_iff,pair_eq_iff,setminus_subseteq,powerset_intro,function_apply_iff}.
    $t$ is a sequence in $\pow{X}$ by \cref{pair_eq_iff,dom_intro,dom_iff,subseteq,subseteq_antisymmetric,funs_intro,sequence_in}.
\end{proof}

% from §48 Item 2: Baire space
\begin{definition}\label{baire_space}
    $X$ is a Baire space with respect to $T$ iff
        $T$ is a topology on $X$ and
        for all $s$ if
            $s$ is a sequence in $\pow{X}$ and
            for all $n\in\naturalsPlus$ we have
                $s(n)$ is closed in $X$ with respect to $T$ and
                $s(n)$ has empty interior in $X$ with respect to $T$
        then $\unions{\img{s}{\naturalsPlus}}$ has empty interior in $X$ with respect to $T$.
\end{definition}

% from §48 Item 3: open-set characterization of a Baire space
\begin{lemma}\label{baire_space_open_set_characterization}
    Assume $T$ is a topology on $X$.
    Then $X$ is a Baire space with respect to $T$ iff
        for all $s$ if
            $s$ is a sequence in $\pow{X}$ and
            for all $n\in\naturalsPlus$ we have
                $s(n)$ is open in $X$ with respect to $T$ and
                $s(n)$ is dense in $X$ with respect to $T$
        then $\inters{\img{s}{\naturalsPlus}}$ is dense in $X$ with respect to $T$.
\end{lemma}
\begin{proof}
    Show that if $X$ is a Baire space with respect to $T$ then
        for all $s$ if
            $s$ is a sequence in $\pow{X}$ and
            for all $n\in\naturalsPlus$ we have
                $s(n)$ is open in $X$ with respect to $T$ and
                $s(n)$ is dense in $X$ with respect to $T$
        then $\inters{\img{s}{\naturalsPlus}}$ is dense in $X$ with respect to $T$.
    \begin{subproof}
        Assume $X$ is a Baire space with respect to $T$.
        Fix $s$ such that
            $s$ is a sequence in $\pow{X}$ and
            for all $n\in\naturalsPlus$ we have
                $s(n)$ is open in $X$ with respect to $T$ and
                $s(n)$ is dense in $X$ with respect to $T$.
        Let $t = \{(n, X\setminus s(n)) \mid n\in\naturalsPlus\}$.
        $t$ is a sequence in $\pow{X}$ by \cref{complement_sequence_in_powerset_48}.
        For all $n\in\naturalsPlus$ we have $s(n)\subseteq X$ by \cref{sequence_in,funs_elim,powerset_elim}.
        For all $n\in\naturalsPlus$ we have $t(n) = X\setminus s(n)$
            by \cref{pair_eq_iff,sequence_in,funs_elim,function_apply_intro}.
        For all $n\in\naturalsPlus$ we have $t(n)$ is closed in $X$ with respect to $T$
            by \cref{sequence_in,funs_elim,powerset_elim,setminus_subseteq,double_relative_complement,closed_in}.
        For all $n\in\naturalsPlus$ we have
            $t(n)$ has empty interior in $X$ with respect to $T$
            by \cref{function_apply_iff,closed_in,double_relative_complement,empty_interior_iff_dense_complement}.
        $\unions{\img{t}{\naturalsPlus}}$ has empty interior in $X$ with respect to $T$ by \cref{baire_space}.
        $X\setminus \unions{\img{t}{\naturalsPlus}}$ is dense in $X$ with respect to $T$
            by \cref{empty_interior_iff_dense_complement,unions_subseteq_of_powerset_is_subseteq,sequence_in,funs_elim,function_to_ran,img_subseteq_ran,subseteq_transitive}.
        For all $x$ such that $x\in X\setminus \unions{\img{t}{\naturalsPlus}}$
            we have $x\in \inters{\img{s}{\naturalsPlus}}$
            by \cref{img_iff,sequence_in,funs_elim,function_apply_iff,setminus_elim_left,setminus_elim_right,pair_eq_iff,function_apply_intro,setminus_intro,function_apply_elim,unions_iff,real_one_in_naturals,inters_intro}.
        $X\setminus \unions{\img{t}{\naturalsPlus}}\subseteq \inters{\img{s}{\naturalsPlus}}$ by \cref{subseteq}.
        $\inters{\img{s}{\naturalsPlus}}\subseteq X$
            by \cref{real_one_in_naturals,sequence_in,funs_elim,function_apply_elim,img_iff,powerset_elim,inters_subseteq_elem,subseteq_transitive}.
        $\inters{\img{s}{\naturalsPlus}}$ is dense in $X$ with respect to $T$
            by \cref{dense_subseteq_superset}.
    \end{subproof}
    Show that if
        for all $s$ if
            $s$ is a sequence in $\pow{X}$ and
            for all $n\in\naturalsPlus$ we have
                $s(n)$ is open in $X$ with respect to $T$ and
                $s(n)$ is dense in $X$ with respect to $T$
        then $\inters{\img{s}{\naturalsPlus}}$ is dense in $X$ with respect to $T$
    then $X$ is a Baire space with respect to $T$.
    \begin{subproof}
        Assume for all $s$ if
            $s$ is a sequence in $\pow{X}$ and
            for all $n\in\naturalsPlus$ we have
                $s(n)$ is open in $X$ with respect to $T$ and
                $s(n)$ is dense in $X$ with respect to $T$
        then $\inters{\img{s}{\naturalsPlus}}$ is dense in $X$ with respect to $T$.
        Show that for all $s$ if
            $s$ is a sequence in $\pow{X}$ and
            for all $n\in\naturalsPlus$ we have
                $s(n)$ is closed in $X$ with respect to $T$ and
                $s(n)$ has empty interior in $X$ with respect to $T$
        then $\unions{\img{s}{\naturalsPlus}}$ has empty interior in $X$ with respect to $T$.
        \begin{subproof}
            Fix $s$ such that
                $s$ is a sequence in $\pow{X}$ and
                for all $n\in\naturalsPlus$ we have
                    $s(n)$ is closed in $X$ with respect to $T$ and
                    $s(n)$ has empty interior in $X$ with respect to $T$.
            Let $t = \{(n, X\setminus s(n)) \mid n\in\naturalsPlus\}$.
            $t$ is a sequence in $\pow{X}$ by \cref{complement_sequence_in_powerset_48}.
            For all $n\in\naturalsPlus$ we have $s(n)\subseteq X$ by \cref{sequence_in,funs_elim,powerset_elim}.
            For all $n\in\naturalsPlus$ we have $t(n) = X\setminus s(n)$
                by \cref{pair_eq_iff,sequence_in,funs_elim,function_apply_intro}.
            For all $n\in\naturalsPlus$ we have $t(n)$ is open in $X$ with respect to $T$ by \cref{closed_in}.
            For all $n\in\naturalsPlus$ we have $t(n)$ is dense in $X$ with respect to $T$
                by \cref{function_apply_iff,open_in,topology_on,pow_iff,elem_subseteq,double_relative_complement,empty_interior_iff_dense_complement}.
            $\inters{\img{t}{\naturalsPlus}}$ is dense in $X$ with respect to $T$.
            Let $B = X\setminus \inters{\img{t}{\naturalsPlus}}$.
            $B$ has empty interior in $X$ with respect to $T$
                by \cref{dense_in,setminus_subseteq,double_relative_complement,empty_interior_iff_dense_complement,real_one_in_naturals,sequence_in,funs_elim,function_apply_elim,img_iff,powerset_elim,inters_subseteq_elem,subseteq_transitive}.
            For all $x$ such that $x\in\unions{\img{s}{\naturalsPlus}}$ we have $x\in B$
                by \cref{unions_iff,img_iff,sequence_in,funs_elim,function_apply_iff,powerset_elim,subseteq,function_apply_elim,inters_subseteq_elem,pair_eq_iff,function_apply_intro,setminus_elim_right,setminus_intro}.
            $\unions{\img{s}{\naturalsPlus}}\subseteq B$ by \cref{subseteq}.
            $\unions{\img{s}{\naturalsPlus}}$ has empty interior in $X$ with respect to $T$
                by \cref{empty_interior_subset,setminus_subseteq}.
        \end{subproof}
        $X$ is a Baire space with respect to $T$ by \cref{baire_space}.
    \end{subproof}
\end{proof}

% from §48 helper lemma 3: a nested sequence is monotone under the natural order
\begin{lemma}\label{nested_sequence_monotone}
    Assume $s$ is a sequence in $\pow{X}$.
    Suppose for all $n\in\naturalsPlus$ we have $s(n + 1)\subseteq s(n)$.
    Then for all $n,m\in\naturalsPlus$ if $n \leq m$ then $s(m)\subseteq s(n)$.
\end{lemma}
\begin{proof}
    Let $P = \{m\in\naturalsPlus \mid \text{for all $n\in\naturalsPlus$ if $n \leq m$ then $s(m)\subseteq s(n)$}\}$.
    $P\subseteq\naturalsPlus$ by \cref{subseteq}.
    For all $n\in\naturalsPlus$ if $n \leq 1$ then $s(1)\subseteq s(n)$
        by \cref{real_naturals_subset_reals,real_one_in_naturals,real_one_leq_naturalsplus,real_leq_antisym,subseteq_refl,subseteq}.
    $1\in P$ by \cref{real_one_in_naturals}.
    Show that for all $k\in P$ we have $k + 1\in P$.
    \begin{subproof}
        Fix $k$ such that $k\in P$.
        $k + 1\in\naturalsPlus$ by \cref{subseteq,real_naturals_closed_succ}.
        Show that for all $n\in\naturalsPlus$ if $n \leq k + 1$ then $s(k + 1)\subseteq s(n)$.
        \begin{subproof}
            Fix $n\in\naturalsPlus$.
            Assume $n \leq k + 1$.
            \begin{byCase}
            \caseOf{$n = k + 1$.}
                $s(k + 1)\subseteq s(n)$ by \cref{subseteq_refl}.
            \caseOf{$n \neq k + 1$.}
                $s(k + 1)\subseteq s(n)$ by \cref{real_naturals_leq_succ_elim,subseteq,subseteq_transitive}.
            \end{byCase}
        \end{subproof}
        $k + 1\in P$.
    \end{subproof}
    For all $m\in\naturalsPlus$ we have $m\in P$ by \cref{real_naturals_induction}.
    For all $n,m\in\naturalsPlus$ if $n \leq m$ then $s(m)\subseteq s(n)$ by \cref{subseteq}.
\end{proof}

% from §48 helper lemma 4: the image of a nested nonempty sequence has the finite intersection property
\begin{lemma}\label{nested_sequence_image_fip}
    Assume $s$ is a sequence in $\pow{X}$.
    Suppose for all $n\in\naturalsPlus$ we have $s(n)\neq\emptyset$.
    Suppose for all $n\in\naturalsPlus$ we have $s(n + 1)\subseteq s(n)$.
    Then $\img{s}{\naturalsPlus}$ has the finite intersection property.
\end{lemma}
\begin{proof}
    Show that $\img{s}{\naturalsPlus}$ is inhabited.
    \begin{subproof}
        Follows by \cref{real_one_in_naturals,sequence_in,funs,funs_is_function,function_apply_iff,img_iff,nonempty_inhabited,inhabited_nonempty}.
    \end{subproof}
    Show that for all $F$ such that $F\subseteq\img{s}{\naturalsPlus}$ and $F$ is finite and inhabited
        we have $\inters{F}\neq\emptyset$.
    \begin{subproof}
        Fix $F$ such that $F\subseteq\img{s}{\naturalsPlus}$ and $F$ is finite and inhabited.
        Let $R = \{(A,n) \mid A\in F, n\in\naturalsPlus \mid A = s(n)\}$.
        $R$ is a relation by \cref{relation,pair_eq_iff}.
        For all $A\in F$ there exists $n$ such that $A\mathrel{R} n$ by
            \cref{subseteq,img_iff,sequence_in,funs_is_function,function_apply_iff,pair_eq_iff}.
        Take $q$ such that $q$ is a function on $F$ and for all $A\in F$ we have $A\mathrel{R}\apply{q}{A}$
            by \cref{choice_function}.
        Let $J = \img{q}{F}$.
        We have $J$ is finite and $J$ is inhabited and $J\subseteq\naturalsPlus$
            by \cref{finite_image_function,choice_function,subseteq,function_apply_iff,img_iff,pair_eq_iff,nonempty_inhabited}.
        Take $m$ such that $m$ is a largest element of $J$ with respect to $\realOrderRel$
            by \cref{real_naturals_subset_reals,subseteq_transitive,subseteq,finite_inhabited_largest_element,real_order_relation_simple}.
        For all $A\in F$ we have $s(m)\subseteq A$ by \cref{choice_function,pair_eq_iff,function_apply_elim,img_iff,largest_element,subseteq,real_order_relation_elim,nested_sequence_monotone,real_lt_subst_right}.
        $s(m)\subseteq\inters{F}$ by \cref{subseteq_inters_iff,subseteq}.
        Take $z$ such that $z\in s(m)$ by \cref{subseteq,largest_element,nonempty_inhabited}.
        Hence $\inters{F}\neq\emptyset$ by \cref{subseteq,nonempty_inhabited,inhabited_nonempty}.
    \end{subproof}
    Therefore $\img{s}{\naturalsPlus}$ has the finite intersection property by \cref{finite_intersection_property}.
\end{proof}

% from §48 Item 4: compact Hausdorff spaces are Baire
\begin{theorem}\label{compact_hausdorff_space_baire}
    Assume $T$ is a topology on $X$.
    Assume $X$ is compact with respect to $T$.
    Assume $X$ is Hausdorff with respect to $T$.
    Then $X$ is a Baire space with respect to $T$.
\end{theorem}
\begin{proof}
    We have $X$ is regular with respect to $T$ by \cref{compact_hausdorff_normal,normal_implies_regular}.
    Show that for all $s$ if
        $s$ is a sequence in $\pow{X}$ and
        for all $n\in\naturalsPlus$ we have
            $s(n)$ is open in $X$ with respect to $T$ and
            $s(n)$ is dense in $X$ with respect to $T$
    then $\inters{\img{s}{\naturalsPlus}}$ is dense in $X$ with respect to $T$.
    \begin{subproof}
        Fix $s$ such that
            $s$ is a sequence in $\pow{X}$ and
            for all $n\in\naturalsPlus$ we have
                $s(n)$ is open in $X$ with respect to $T$ and
                $s(n)$ is dense in $X$ with respect to $T$.
        We have $\inters{\img{s}{\naturalsPlus}}\subseteq X$ by
            \cref{unions_subseteq_of_powerset_is_subseteq,inters,elem_subseteq,subseteq,img_subseteq_ran,sequence_in,funs_ran,subseteq_transitive}.
        Show that for all $x\in X$ we have $x\in\closure{\inters{\img{s}{\naturalsPlus}}}{X}{T}$.
        \begin{subproof}
            Fix $x\in X$.
            Show that for all $U$ such that
                $U$ is open in $X$ with respect to $T$ and
                $x\in U$
                we have $U$ intersects $\inters{\img{s}{\naturalsPlus}}$.
            \begin{subproof}
                Fix $U$ such that
                    $U$ is open in $X$ with respect to $T$ and
                    $x\in U$.
                We have $s(1)$ is open in $X$ with respect to $T$ and $s(1)$ is dense in $X$ with respect to $T$
                    by \cref{real_one_in_naturals,subseteq}.
                $s(1)\subseteq X$ by \cref{open_in,topology_on,subseteq,pow_iff}.
                $x\in\closure{s(1)}{X}{T}$ by \cref{dense_in,real_lt_subst_right}.
                Hence $U$ intersects $s(1)$ by \cref{closure_iff_intersects_open}.
                Take $y_0$ such that $y_0\in U\inter s(1)$ by \cref{intersects,nonempty_inhabited}.
                Let $W_0 = U\inter s(1)$.
                $W_0$ is open in $X$ with respect to $T$ and $y_0\in W_0$
                    by \cref{inter,open_in,topology_on,neighborhood}.
                $y_0\in X$ by \cref{open_in,topology_on,subseteq,pow_iff,subseteq}.
                $W_0$ is a neighborhood of $y_0$ in $X$ with respect to $T$ by \cref{neighborhood}.
                Take $V_0$ such that
                    $V_0$ is a neighborhood of $y_0$ in $X$ with respect to $T$ and
                    $\closure{V_0}{X}{T}\subseteq W_0$
                    by \cref{regular_closure_neighborhood}.
                $V_0$ is open in $X$ with respect to $T$ and $y_0\in V_0$ by \cref{neighborhood}.
                Let $a_0 = (1,V_0)$.
                Let $A = \{p\in\naturalsPlus\times\pow{X} \mid
                    \snd{p}\in T \land
                    \snd{p}\neq\emptyset \land
                    \closure{\snd{p}}{X}{T}\subseteq s(\fst{p})\}$.
                $a_0\in A$
                    by \cref{real_one_in_naturals,subseteq,neighborhood,open_in,topology_on,emptyset,pow_iff,powerset_intro,fst_eq,snd_eq,times_tuple_intro,pair_eq_iff,inter_lower_right,subseteq_transitive}.
                Let $R = \{w\in A\times A \mid \fst{\snd{w}} = \fst{\fst{w}} + 1 \land \closure{\snd{\snd{w}}}{X}{T}\subseteq \snd{\fst{w}}\}$.
                $R$ is a relation by \cref{relation,subseteq,times_elem_is_tuple}.
                Show that for all $p\in A$ there exists $q$ such that $p\mathrel{R} q$.
                \begin{subproof}
                    Fix $p\in A$.
                    We have $p\in\naturalsPlus\times\pow{X}$ and
                        $\snd{p}\in T$ and
                        $\snd{p}\neq\emptyset$ and
                        $\closure{\snd{p}}{X}{T}\subseteq s(\fst{p})$ by \cref{subseteq}.
                    Take $n,W$ such that
                        $p = (n,W)$ and
                        $n\in\naturalsPlus$ and
                        $W\in\pow{X}$ by \cref{times}.
                    $W = \snd{p}$ and $n = \fst{p}$ by \cref{fst_eq,snd_eq}.
                    $W\in T$ and $W\neq\emptyset$ and $\closure{W}{X}{T}\subseteq s(n)$
                        by \cref{real_lt_subst_left,real_lt_subst_right}.
                    $W$ is open in $X$ with respect to $T$ by \cref{open_in}.
                    Take $z$ such that $z\in W$ by \cref{nonempty_inhabited}.
                    $s(n + 1)$ is open in $X$ with respect to $T$ and $s(n + 1)$ is dense in $X$ with respect to $T$
                        by \cref{real_naturals_closed_succ,subseteq}.
                    $s(n + 1)\subseteq X$ by \cref{open_in,topology_on,subseteq,pow_iff}.
                    $z\in X$ by \cref{subseteq,open_in,topology_on,pow_iff}.
                    Hence $z\in\closure{s(n + 1)}{X}{T}$ by \cref{dense_in,real_lt_subst_right}.
                    Therefore $W$ intersects $s(n + 1)$ by \cref{closure_iff_intersects_open}.
                    Take $y$ such that $y\in W\inter s(n + 1)$ by \cref{intersects,nonempty_inhabited}.
                    Let $B = W\inter s(n + 1)$.
                    $B$ is open in $X$ with respect to $T$ and $y\in B$
                        by \cref{inter,open_in,topology_on,neighborhood}.
                    $y\in X$ by \cref{open_in,topology_on,subseteq,pow_iff,subseteq}.
                    $B$ is a neighborhood of $y$ in $X$ with respect to $T$ by \cref{neighborhood}.
                    Take $V$ such that
                        $V$ is a neighborhood of $y$ in $X$ with respect to $T$ and
                        $\closure{V}{X}{T}\subseteq B$
                        by \cref{regular_closure_neighborhood}.
                    $V$ is open in $X$ with respect to $T$ and $y\in V$ by \cref{neighborhood}.
                    Let $q = (n + 1,V)$.
                    $q\in A$
                        by \cref{real_naturals_closed_succ,subseteq,neighborhood,open_in,topology_on,emptyset,pow_iff,powerset_intro,fst_eq,snd_eq,times_tuple_intro,pair_eq_iff,inter_lower_right,subseteq_transitive}.
                    $\closure{\snd{q}}{X}{T}\subseteq \snd{p}$ by \cref{fst_eq,snd_eq,inter_lower_left,subseteq_transitive}.
                    Hence $p\mathrel{R} q$ by \cref{fst_eq,snd_eq,times_tuple_intro,relation}.
                \end{subproof}
                Take $c$ such that $c$ is a function on $A$ and for all $p\in A$ we have $p\mathrel{R}\apply{c}{p}$
                    by \cref{choice_function}.
                We have $c\in\funs{A}{A}$ by \cref{choice_function,relation,times_tuple_elim,funs_intro}.
                Take $u$ such that
                    $u$ is a sequence in $A$ and
                    $u(1) = a_0$ and
                    for all $n\in\naturalsPlus$ we have $u(n + 1) = \apply{c}{u(n)}$
                    by \cref{iteration_sequence_naturalsplus}.
                Let $P = \{n\in\naturalsPlus \mid \fst{u(n)} = n\}$.
                $P\subseteq\naturalsPlus$ by \cref{subseteq}.
                $1\in P$ by \cref{fst_eq,real_one_in_naturals}.
                For all $n\in P$ we have $n + 1\in P$ by \cref{subseteq,sequence_in,funs_type_apply,choice_function,real_lt_subst_right,relation,fst_eq,snd_eq,real_naturals_closed_succ}.
                Therefore for all $n\in\naturalsPlus$ we have $n\in P$ by \cref{real_naturals_induction}.
                For all $n\in\naturalsPlus$ we have $\fst{u(n)} = n$ by \cref{subseteq}.
                Let $h = \{(n,\closure{\snd{u(n)}}{X}{T}) \mid n\in\naturalsPlus\}$.
                For all $a,b_1,b_2$ such that $(a,b_1)\in h$ and $(a,b_2)\in h$ we have $b_1 = b_2$
                    by \cref{pair_eq_iff}.
                $h$ is a function by \cref{pair_eq_iff,relation,rightunique}.
                We have $\dom{h} = \naturalsPlus$ by \cref{dom_intro,pair_eq_iff,dom_iff,subseteq,subseteq_antisymmetric}.
                For all $n$ such that $n\in\dom{h}$ we have $h(n)\in\pow{X}$
                    by \cref{subseteq,pair_eq_iff,function_apply_iff,real_lt_subst_right,closure,powerset_intro}.
                Therefore $h\in\funs{\naturalsPlus}{\pow{X}}$ by \cref{funs_intro}.
                $h$ is a sequence in $\pow{X}$ by \cref{sequence_in}.
                Show that for all $n\in\naturalsPlus$ we have $h(n)\neq\emptyset$.
                \begin{subproof}
                    Follows by \cref{sequence_in,funs_type_apply,subseteq,times,snd_eq,real_lt_subst_left,powerset_elim,pair_eq_iff,function_apply_iff,nonempty_inhabited,real_lt_subst_right,subseteq_closure,inhabited_nonempty}.
                \end{subproof}
                Show that for all $n\in\naturalsPlus$ we have $h(n + 1)\subseteq h(n)$.
                \begin{subproof}
                    Follows by \cref{sequence_in,funs_type_apply,real_naturals_closed_succ,subseteq,choice_function,real_lt_subst_right,relation,fst_eq,snd_eq,times,powerset_elim,subseteq_closure,subseteq_transitive,function_apply_intro}.
                \end{subproof}
                Therefore $\img{h}{\naturalsPlus}$ has the finite intersection property by \cref{nested_sequence_image_fip}.
                Show that for all $C$ such that $C\in\img{h}{\naturalsPlus}$ we have
                    $C$ is closed in $X$ with respect to $T$.
                \begin{subproof}
                    Follows by \cref{img_iff,sequence_in,funs_is_function,function_apply_iff,sequence_in,funs_type_apply,subseteq,times,snd_eq,powerset_elim,real_lt_subst_right,function_apply_intro,closure_closed_in,real_lt_subst_left}.
                \end{subproof}
                Hence $\inters{\img{h}{\naturalsPlus}}\neq\emptyset$ by \cref{compact_iff_finite_intersection_property}.
                Take $p$ such that $p\in\inters{\img{h}{\naturalsPlus}}$ by \cref{nonempty_inhabited}.
                $p\in h(1)$ by \cref{sequence_in,function_apply_elim,img_iff,real_one_in_naturals,inters_destr}.
                $h(1) = \closure{\snd{u(1)}}{X}{T}$ by \cref{real_one_in_naturals,function_apply_intro}.
                $\snd{u(1)} = V_0$ by \cref{subseteq,snd_eq}.
                $p\in\closure{V_0}{X}{T}$ by \cref{real_lt_subst_left,real_lt_subst_right}.
                $\closure{V_0}{X}{T}\subseteq W_0$ by \cref{subseteq}.
                $p\in U$ by \cref{subseteq,inter}.
                Show that for all $A_1$ such that $A_1\in\img{s}{\naturalsPlus}$ we have
                    $p\in A_1$.
                \begin{subproof}
                    Follows by \cref{img_iff,sequence_in,funs_is_function,function_apply_iff,function_apply_elim,inters_destr,sequence_in,funs_type_apply,fst_eq,snd_eq,pair_eq_iff,function_apply_intro,subseteq,real_lt_subst_right,real_lt_subst_left}.
                \end{subproof}
                Thus $p\in\inters{\img{s}{\naturalsPlus}}$
                    by \cref{sequence_in,funs,funs_is_function,function_apply_iff,img_iff,real_one_in_naturals,inters_intro}.
                Hence $U$ intersects $\inters{\img{s}{\naturalsPlus}}$ by \cref{inter_intro,emptyset,intersects,inter}.
            \end{subproof}
            Thus $x\in\closure{\inters{\img{s}{\naturalsPlus}}}{X}{T}$ by \cref{closure_iff_intersects_open}.
        \end{subproof}
        $X\subseteq\closure{\inters{\img{s}{\naturalsPlus}}}{X}{T}$ by \cref{subseteq}.
        Hence $\closure{\inters{\img{s}{\naturalsPlus}}}{X}{T} = X$
            by \cref{closure,subseteq,subseteq_antisymmetric}.
        $\inters{\img{s}{\naturalsPlus}}$ is dense in $X$ with respect to $T$ by \cref{dense_in}.
    \end{subproof}
    Therefore $X$ is a Baire space with respect to $T$ by \cref{baire_space_open_set_characterization}.
\end{proof}

% from §48 helper lemma 5: two points in the same half-ball are epsilon-close
\begin{lemma}\label{metric_half_ball_pair_bound}
    Assume $d$ is a metric on $X$.
    Suppose $a,x,y\in X$.
    Suppose $\epsilon\in\reals$ and $0 < \epsilon$.
    Suppose $x\in\metricBall{d}{X}{a}{\epsilon / (1 + 1)}$.
    Suppose $y\in\metricBall{d}{X}{a}{\epsilon / (1 + 1)}$.
    Then $d((x,y)) < \epsilon$.
\end{lemma}
\begin{proof}
    We have $d((x,y)) < \epsilon$ by
        \cref{real_half_of_positive,metric_ball,metric_on,metric_symmetric,metric_common_center_double_bound,real_lt_subst_left,real_lt_subst_right}.
\end{proof}

% from §48 helper lemma 5a: half a positive reciprocal is positive and real
\begin{lemma}\label{half_reciprocal_naturalsplus_real_positive}
    Suppose $n\in\naturalsPlus$.
    Then $(1 / n) / (1 + 1)\in\reals$ and $0 < (1 / n) / (1 + 1)$.
\end{lemma}
\begin{proof}
    Follows by \cref{naturalsplus_reciprocal_real,positive_reciprocal_naturalsplus,real_half_of_positive}.
\end{proof}

% from §48 helper definition 1: a set fits into a half-reciprocal ball at a stage
\begin{definition}\label{metric_stage_small_set}
    $W$ is stage-small at $n$ in $X$ with respect to $d$ iff
        there exists $a\in X$ such that
            $W\subseteq\metricBall{d}{X}{a}{(1 / n) / (1 + 1)}$.
\end{definition}

% from §48 helper lemma 6: an eventual closed-set constraint passes to the limit
\begin{lemma}\label{eventually_in_closed_limit}
    Assume $T$ is a topology on $X$.
    Assume $A$ is closed in $X$ with respect to $T$.
    Assume $s$ is a sequence in $X$.
    Assume $x\in X$.
    Suppose $s$ converges to $x$ in $X$ with respect to $T$.
    Suppose there exists $N\in\naturalsPlus$ such that for all $n\in\naturalsPlus$ if $N \leq n$ then $s(n)\in A$.
    Then $x\in A$.
\end{lemma}
\begin{proof}
    Take $N\in\naturalsPlus$ such that for all $n\in\naturalsPlus$ if $N \leq n$ then $s(n)\in A$.
    Show that for all $O$ such that
        $O$ is open in $X$ with respect to $T$ and
        $x\in O$
        we have $O$ intersects $A$.
    \begin{subproof}
        Follows by \cref{converges_to,neighborhood,naturalsplus_add_closed,real_naturals_subset_reals,subseteq,naturalsplus_positive_real,real_add_ident,real_lt_add,real_add_comm,real_lt_subst_left,real_lt_subst_right,real_lt_imp_leq,inter_intro,emptyset,intersects,inter}.
    \end{subproof}
    Therefore $x\in A$ by \cref{closed_in,closure_iff_intersects_open,closure_subseteq_closed,subseteq_refl,subseteq}.
\end{proof}

% from §48 helper lemma 8: an explicit witness yields the Cauchy existence clause
\begin{lemma}\label{cauchy_witness_intro}
    Suppose $L\in\naturalsPlus$.
    Suppose for all $n,m\in\naturalsPlus$ if $L \leq n$ and $L \leq m$ then $d((s(n),s(m))) < \epsilon$.
    Then there exists $N\in\naturalsPlus$ such that
        for all $n,m\in\naturalsPlus$ if $N \leq n$ and $N \leq m$ then $d((s(n),s(m))) < \epsilon$.
\end{lemma}
\begin{proof}
    Take $N\in\naturalsPlus$ such that
        $N = L$ and
        for all $n,m\in\naturalsPlus$ if $N \leq n$ and $N \leq m$ then $d((s(n),s(m))) < \epsilon$.
\end{proof}

% from §48 helper lemma 9: a diameter bounds all pairwise distances
\begin{lemma}\label{metric_diameter_bound}
    Assume $d$ is a metric on $X$.
    Assume $r$ is a diameter of $A$ in $X$ with respect to $d$.
    Suppose $x\in A$ and $y\in A$.
    Then $d((x,y)) \leq r$.
\end{lemma}
\begin{proof}
    Follows by \cref{metric_diameter,metric_distance_set,real_supremum,real_upper_bound}.
\end{proof}

% from §48 helper lemma 10: nested members are eventually pairwise close
\begin{lemma}\label{nested_member_sequence_small_tail}
    Assume $s$ is a sequence in $\pow{X}$.
    Assume for all $n\in\naturalsPlus$ we have $s(n+1)\subseteq s(n)$.
    Assume $N\in\naturalsPlus$.
    Assume for all $n\in\naturalsPlus$ if $N \leq n$ then
        for all $x,y\in s(n)$ we have $d((x,y)) < \epsilon$.
    Assume $t$ is a sequence in $X$.
    Assume for all $n\in\naturalsPlus$ we have $t(n)\in s(n)$.
    Then for all $n\in\naturalsPlus$ we have
        for all $m\in\naturalsPlus$ if $N \leq n$ and $N \leq m$ then
            $d((t(n),t(m))) < \epsilon$.
\end{lemma}
\begin{proof}
    Show that for all $n\in\naturalsPlus$ we have
        for all $m\in\naturalsPlus$ if $N \leq n$ and $N \leq m$ then
            $d((t(n),t(m))) < \epsilon$.
    \begin{subproof}
        Follows by \cref{subseteq,nested_sequence_monotone}.
    \end{subproof}
\end{proof}

% from §48 helper lemma 11: a naturalsPlus sequence has a strictly increasing majorant
\begin{lemma}\label{naturalsplus_sequence_strictly_increasing_majorant}
    Assume $v$ is a sequence in $\naturalsPlus$.
    Then there exists $u$ such that
        $u$ is strictly increasing on $\naturalsPlus$ and
        for all $n\in\naturalsPlus$ we have $v(n) \leq u(n)$.
\end{lemma}
\begin{proof}
    Let $D = \{p\in\naturalsPlus\times\naturalsPlus \mid v(\fst{p}) \leq \snd{p}\}$.
    Let $a_0 = (1,v(1))$.
    $v(1)\in\naturalsPlus$ by \cref{sequence_in,funs_type_apply,real_one_in_naturals}.
    $v(1)\in\reals$ by \cref{real_naturals_subset_reals,subseteq}.
    $v(1) \leq v(1)$.
    $a_0\in D$ by \cref{real_one_in_naturals,times_tuple_intro,fst_eq,snd_eq}.

    Let $R = \{w\in D\times D \mid
        \fst{\snd{w}} = \fst{\fst{w}} + 1 \land \snd{\fst{w}} < \snd{\snd{w}}\}$.

    Show that for all $p\in D$ there exists $q$ such that $p\mathrel{R} q$.
    \begin{subproof}
        Fix $p\in D$.
        Take $n_0,m_0$ such that $p = (n_0,m_0)$ by \cref{times_elem_is_tuple}.
        We have $n_0\in\naturalsPlus$ and $m_0\in\naturalsPlus$ by \cref{times_tuple_elim}.
        $n_0 + 1\in\naturalsPlus$ by \cref{real_naturals_closed_succ}.
        $v(n_0 + 1)\in\naturalsPlus$ by \cref{sequence_in,funs_type_apply}.
        $v(n_0 + 1)\in\reals$ and $m_0\in\reals$ by \cref{real_naturals_subset_reals,subseteq}.
        $m_0 < v(n_0 + 1)$ or $m_0 = v(n_0 + 1)$ or $v(n_0 + 1) < m_0$ by \cref{real_lt_trichotomy}.
        \begin{byCase}
        \caseOf{$m_0 < v(n_0 + 1)$.}
            Let $k_1 = v(n_0 + 1)$.
            $k_1\in\naturalsPlus$ and $k_1\in\reals$
                by \cref{real_lt_subst_left,real_naturals_subset_reals,subseteq}.
            $v(n_0 + 1) \leq k_1$.
            Let $q = (n_0 + 1,k_1)$.
            $q\in D$ by \cref{real_lt_subst_left,real_lt_subst_right,times_tuple_intro,fst_eq,snd_eq}.
            Hence $p\mathrel{R} q$
                by \cref{times_tuple_intro,fst_eq,snd_eq,relation,real_lt_subst_left,real_lt_subst_right,real_add_eq_subst}.
            Follows by definition.
        \caseOf{$m_0 = v(n_0 + 1)$.}
            Let $k_1 = m_0 + 1$.
            $k_1\in\naturalsPlus$ by \cref{real_naturals_closed_succ}.
            $m_0 < k_1$ by \cref{real_naturals_lt_succ_local}.
            $v(n_0 + 1) \leq k_1$ by \cref{real_lt_subst_right,real_lt_imp_leq}.
            Let $q = (n_0 + 1,k_1)$.
            $q\in D$ by \cref{times_tuple_intro,fst_eq,snd_eq,real_lt_subst_left,real_lt_subst_right}.
            Thus $p\mathrel{R} q$
                by \cref{times_tuple_intro,fst_eq,snd_eq,relation,real_lt_subst_left,real_lt_subst_right,real_add_eq_subst}.
            Follows by definition.
        \caseOf{$v(n_0 + 1) < m_0$.}
            Let $k_1 = m_0 + 1$.
            $k_1\in\naturalsPlus$ by \cref{real_naturals_closed_succ}.
            $k_1\in\reals$ by \cref{real_naturals_subset_reals,subseteq}.
            $m_0 < k_1$ by \cref{real_naturals_lt_succ_local}.
            $v(n_0 + 1) \leq k_1$ by \cref{real_naturals_lt_succ_local,real_lt_imp_leq,real_leq_trans}.
            Let $q = (n_0 + 1,k_1)$.
            $q\in D$ by \cref{times_tuple_intro,fst_eq,snd_eq,real_lt_subst_left,real_lt_subst_right}.
            Therefore $p\mathrel{R} q$
                by \cref{times_tuple_intro,fst_eq,snd_eq,relation,real_lt_subst_left,real_lt_subst_right,real_add_eq_subst}.
            Follows by definition.
        \end{byCase}
    \end{subproof}

    Take $f$ such that $f$ is a function on $D$ and for all $p\in D$ we have
        $p\mathrel{R}\apply{f}{p}$ by \cref{subseteq,times_elem_is_tuple,relation,choice_function}.
    Therefore $f\in\funs{D}{D}$ by \cref{funs_intro,subseteq,times_tuple_elim}.
    Take $w$ such that
        $w$ is a sequence in $D$ and
        $w(1) = a_0$ and
        for all $n\in\naturalsPlus$ we have $w(n + 1) = \apply{f}{w(n)}$
        by \cref{iteration_sequence_naturalsplus}.

    Let $g = \{(p,\snd{p}) \mid p\in D\}$.
    We have $g$ is a function by \cref{pair_eq_iff,rightunique,relation}.
    $\dom{g} = D$ by \cref{dom_iff,pair_eq_iff,subseteq,subseteq_antisymmetric}.
    For all $p\in\dom{g}$ we have $g(p)\in\naturalsPlus$
        by \cref{pair_eq_iff,function_apply_intro,subseteq,times_elem_is_tuple,snd_eq}.
    Thus $g\in\funs{D}{\naturalsPlus}$ by \cref{funs_intro}.

    Let $u = g\circ w$.
    We have $u$ is a sequence in $\naturalsPlus$ by \cref{sequence_in,funs_circ}.
    For all $n\in\naturalsPlus$ we have $u(n) = \snd{w(n)}$
        by \cref{funs,sequence_in,funs_is_function,function_apply_elim,funs_type_apply,circ_iff,function_apply_intro}.

    Let $P = \{n\in\naturalsPlus \mid \fst{w(n)} = n\}$.
    $P\subseteq\naturalsPlus$ by \cref{subseteq}.
    $1\in P$ by \cref{real_one_in_naturals,fst_eq,pair_eq_pair_of_proj}.
    For all $n\in P$ we have $n + 1\in P$
        by \cref{subseteq,sequence_in,funs_type_apply,choice_function,fst_eq,snd_eq,real_naturals_closed_succ,iteration_sequence_naturalsplus,real_lt_subst_left,real_lt_subst_right}.
    Therefore for all $n\in\naturalsPlus$ we have $n\in P$ by \cref{real_naturals_induction}.

    For all $n\in\naturalsPlus$ we have $u(n) < u(n + 1)$ by \cref{iteration_sequence_naturalsplus,sequence_in,funs_type_apply,times_elem_is_tuple,times_tuple_elim,fst_eq,snd_eq,pair_eq_iff,real_lt_subst_left,real_lt_subst_right,real_add_eq_subst,real_naturals_closed_succ}.

    Let $T_0 = \{n\in\naturalsPlus \mid \text{for all $m\in\naturalsPlus$ such that $m < n$ we have $u(m) < u(n)$}\}$.
    $T_0\subseteq\naturalsPlus$ by \cref{subseteq}.
    $1\in T_0$ by \cref{real_one_in_naturals,real_one_leq_naturalsplus,real_naturals_subset_reals,subseteq,real_lt_subst_right,real_lt_trans,real_lt_irreflexive}.
    Show that for all $n\in T_0$ we have $n + 1\in T_0$.
    \begin{subproof}
        Fix $n$ such that $n\in T_0$.
        Show that for all $m\in\naturalsPlus$ such that $m < n + 1$ we have $u(m) < u(n + 1)$.
        \begin{subproof}
            Fix $m$ such that $m\in\naturalsPlus$ and $m < n + 1$.
            Suppose not.
            We have $m < n$ or $m = n$ or $m = n + 1$ by \cref{real_succ_discrete}.
            \begin{byCase}
            \caseOf{$m = n$.}
                Contradiction by \cref{real_lt_subst_left}.
            \caseOf{$m < n$.}
                $u(m) < u(n)$ by \cref{subseteq}.
                Hence $u(m) < u(n + 1)$
                    by \cref{sequence_in,real_naturals_closed_succ,funs_type_apply,real_naturals_subset_reals,subseteq,real_lt_trans}.
                Contradiction.
            \caseOf{$m = n + 1$.}
                Contradiction by \cref{real_lt_subst_left,real_naturals_subset_reals,subseteq,real_lt_irreflexive}.
            \end{byCase}
        \end{subproof}
        $n + 1\in T_0$.
    \end{subproof}
    Therefore for all $n\in\naturalsPlus$ we have $n\in T_0$ by \cref{real_naturals_induction}.
    For all $m,n\in\naturalsPlus$ such that $m < n$ we have $u(m) < u(n)$ by \cref{subseteq}.
    Thus $u$ is strictly increasing on $\naturalsPlus$ by \cref{sequence_in,strictly_increasing_on_set}.

    Take $w_0$ such that
        $w_0 = u$ and
        $w_0$ is strictly increasing on $\naturalsPlus$ and
        for all $n\in\naturalsPlus$ we have $v(n) \leq w_0(n)$
        by \cref{subseteq,sequence_in,funs_type_apply,real_lt_subst_left,real_lt_subst_right}.
\end{proof}

% from §48 helper lemma 13: eventual small tails make a selected sequence Cauchy
\begin{lemma}\label{nested_member_sequence_cauchy}
    Assume $d$ is a metric on $X$.
    Assume $s$ is a sequence in $\pow{X}$.
    Assume for all $n\in\naturalsPlus$ we have $s(n+1)\subseteq s(n)$.
    Assume $t$ is a sequence in $X$.
    Assume for all $n\in\naturalsPlus$ we have $t(n)\in s(n)$.
    Suppose for all $\epsilon\in\reals$ such that $0 < \epsilon$
        there exists $N\in\naturalsPlus$ such that
            for all $n\in\naturalsPlus$ if $N \leq n$ then
                for all $x,y\in s(n)$ we have $d((x,y)) < \epsilon$.
    Then $t$ is a Cauchy sequence in $X$ with respect to $d$.
\end{lemma}
\begin{proof}
    We have $t$ is a Cauchy sequence in $X$ with respect to $d$
        by \cref{subseteq,nested_member_sequence_small_tail,cauchy_witness_intro,cauchy_sequence}.
\end{proof}

% from §48 helper lemma 12: nested closed sets with eventually small tails meet
\begin{lemma}\label{nested_closed_sets_eventually_small_nonempty_intersection}
    Assume $d$ is a metric on $X$.
    Assume $X$ is complete with respect to $d$.
    Assume $s$ is a sequence in $\pow{X}$.
    Suppose for all $n\in\naturalsPlus$ we have
        $s(n)$ is closed in $X$ with respect to $\metricTopology{d}{X}$ and
        $s(n)\neq\emptyset$.
    Suppose for all $n\in\naturalsPlus$ we have
        $s(n+1)\subseteq s(n)$.
    Suppose for all $\epsilon\in\reals$ such that $0 < \epsilon$
        there exists $N\in\naturalsPlus$ such that
            for all $n\in\naturalsPlus$ if $N \leq n$ then
                for all $x,y\in s(n)$ we have $d((x,y)) < \epsilon$.
    Then $\inters{\img{s}{\naturalsPlus}}\neq\emptyset$.
\end{lemma}
\begin{proof}
    Let $R = \{(n,x) \mid n\in\naturalsPlus, x\in X \mid x\in s(n)\}$.
    $R$ is a relation by \cref{relation,pair_eq_iff}.
    For all $n\in\naturalsPlus$ there exists $x$ such that $n\mathrel{R} x$ by
        \cref{subseteq,nonempty_inhabited,closed_in,pair_eq_iff}.
    Take $t$ such that $t$ is a function on $\naturalsPlus$ and
        for all $n\in\naturalsPlus$ we have $n\mathrel{R} \apply{t}{n}$ by \cref{choice_function}.
    Hence $t$ is a sequence in $X$ by \cref{choice_function,pair_eq_iff,subseteq,closed_in,funs_intro,sequence_in}.
    For all $n\in\naturalsPlus$ we have $t(n)\in s(n)$ by \cref{choice_function,pair_eq_iff}.

    Therefore $t$ is a Cauchy sequence in $X$ with respect to $d$
        by \cref{nested_member_sequence_cauchy,subseteq}.

    Take $p\in X$ such that $t$ converges to $p$ in $X$ with respect to $\metricTopology{d}{X}$ by \cref{complete_metric_space}.
    For all $n\in\naturalsPlus$ we have $p\in s(n)$ by
        \cref{subseteq,metric_on,metric_basis_is_basis,basis_topology_is_topology,metric_topology,nested_sequence_monotone,eventually_in_closed_limit}.
    Show that for all $A$ such that $A\in\img{s}{\naturalsPlus}$ we have $p\in A$.
    \begin{subproof}
        Follows by \cref{img_iff,sequence_in,funs_is_function,function_apply_iff,subseteq,real_lt_subst_right}.
    \end{subproof}
    Hence $p\in\inters{\img{s}{\naturalsPlus}}$ by
        \cref{sequence_in,funs,funs_is_function,function_apply_iff,img_iff,real_one_in_naturals,inters_intro}.
    Therefore $\inters{\img{s}{\naturalsPlus}}\neq\emptyset$ by \cref{nonempty_inhabited,inhabited_nonempty}.
\end{proof}

% from §48 Item 4: complete metric spaces are Baire
\begin{theorem}\label{complete_metric_space_baire}
    Assume $d$ is a metric on $X$.
    Assume $X$ is complete with respect to $d$.
    Then $X$ is a Baire space with respect to $\metricTopology{d}{X}$.
\end{theorem}
\begin{proof}
    Let $T = \metricTopology{d}{X}$.
    We have $T$ is a topology on $X$ by \cref{metric_on,metric_basis_is_basis,basis_topology_is_topology,metric_topology}.
    Hence $X$ is regular with respect to $T$ by \cref{metrizable,metrizable_normal,normal_implies_regular}.
    Show that for all $s$ if
        $s$ is a sequence in $\pow{X}$ and
        for all $n\in\naturalsPlus$ we have
            $s(n)$ is open in $X$ with respect to $T$ and
            $s(n)$ is dense in $X$ with respect to $T$
    then $\inters{\img{s}{\naturalsPlus}}$ is dense in $X$ with respect to $T$.
    \begin{subproof}
        Fix $s$ such that
            $s$ is a sequence in $\pow{X}$ and
            for all $n\in\naturalsPlus$ we have
                $s(n)$ is open in $X$ with respect to $T$ and
                $s(n)$ is dense in $X$ with respect to $T$.
        $\img{s}{\naturalsPlus}\subseteq\pow{X}$ by \cref{img_subseteq_ran,sequence_in,funs_ran,subseteq_transitive}.
        We have $\inters{\img{s}{\naturalsPlus}}\subseteq X$ by
            \cref{unions_subseteq_of_powerset_is_subseteq,inters,elem_subseteq,subseteq}.
        Show that for all $x\in X$ we have $x\in\closure{\inters{\img{s}{\naturalsPlus}}}{X}{T}$.
        \begin{subproof}
            Fix $x\in X$.
            Show that for all $U$ such that
                $U$ is open in $X$ with respect to $T$ and
                $x\in U$
                we have $U$ intersects $\inters{\img{s}{\naturalsPlus}}$.
            \begin{subproof}
                Fix $U$ such that
                    $U$ is open in $X$ with respect to $T$ and
                    $x\in U$.
                We have $s(1)$ is open in $X$ with respect to $T$ and $s(1)$ is dense in $X$ with respect to $T$
                    by \cref{real_one_in_naturals,subseteq}.
                $s(1)\subseteq X$ by \cref{open_in,topology_on,subseteq,pow_iff}.
                $x\in\closure{s(1)}{X}{T}$ by \cref{dense_in,real_lt_subst_right}.
                Hence $U$ intersects $s(1)$ by \cref{closure_iff_intersects_open}.
                Take $y_0$ such that $y_0\in U\inter s(1)$ by \cref{intersects,nonempty_inhabited}.
                $y_0\in X$ by \cref{inter_elim_left,open_in,topology_on,subseteq,pow_iff}.
                $(1 / 1) / (1 + 1)\in\reals$ and $0 < (1 / 1) / (1 + 1)$
                    by \cref{real_one_in_naturals,half_reciprocal_naturalsplus_real_positive}.
                Let $C_0 = s(1)\inter\metricBall{d}{X}{y_0}{(1 / 1) / (1 + 1)}$.
                $\metricBall{d}{X}{y_0}{(1 / 1) / (1 + 1)}$ is open in $X$ with respect to $T$
                    by \cref{metric_ball,subseteq,powerset_intro,metric_basis,metric_on,metric_basis_is_basis,basis_elem_open_in_generated,metric_topology,basis_for_topology,open_in}.
                $y_0\in\metricBall{d}{X}{y_0}{(1 / 1) / (1 + 1)}$
                    by \cref{metric_on,metric_zero_characterization,metric_ball,real_lt_subst_left}.
                $C_0$ is open in $X$ with respect to $T$ and $y_0\in C_0$
                    by \cref{inter,open_in,topology_on,neighborhood,metric_ball}.
                Let $B_0 = U\inter C_0$.
                $B_0$ is open in $X$ with respect to $T$ and $y_0\in B_0$
                    by \cref{inter,open_in,topology_on,neighborhood}.
                Take $V_0$ such that
                    $V_0$ is a neighborhood of $y_0$ in $X$ with respect to $T$ and
                    $\closure{V_0}{X}{T}\subseteq B_0$
                    by \cref{inter,open_in,topology_on,neighborhood,regular_closure_neighborhood}.
                $V_0$ is open in $X$ with respect to $T$ and $y_0\in V_0$ by \cref{neighborhood}.
                Let $a_0 = (1,V_0)$.
                Let $A = \{p\in\naturalsPlus\times\pow{X} \mid \text{$\snd{p}\in T$ and
                    $\snd{p}\neq\emptyset$ and
                    $\closure{\snd{p}}{X}{T}\subseteq s(\fst{p})$ and
                    $\closure{\snd{p}}{X}{T}$ is stage-small at $\fst{p}$ in $X$ with respect to $d$}\}$.
                $B_0\subseteq C_0$ by \cref{inter_lower_right}.
                $C_0\subseteq s(1)$ and $C_0\subseteq\metricBall{d}{X}{y_0}{(1 / 1) / (1 + 1)}$
                    by \cref{inter_lower_left,inter_lower_right}.
                $\closure{V_0}{X}{T}\subseteq s(1)$ by \cref{subseteq_transitive}.
                Thus $\closure{V_0}{X}{T}$ is stage-small at $1$ in $X$ with respect to $d$
                    by \cref{real_one_in_naturals,subseteq,metric_stage_small_set}.
                $a_0\in A$ by
                    \cref{real_one_in_naturals,subseteq,neighborhood,open_in,topology_on,emptyset,pow_iff,powerset_intro,fst_eq,snd_eq,times_tuple_intro,pair_eq_iff,inter_lower_right,subseteq_transitive,metric_stage_small_set}.
                Let $R = \{w\in A\times A \mid \fst{\snd{w}} = \fst{\fst{w}} + 1 \land \closure{\snd{\snd{w}}}{X}{T}\subseteq \snd{\fst{w}}\}$.
                $R$ is a relation by \cref{relation,subseteq,times_elem_is_tuple}.
                Show that for all $p\in A$ there exists $q$ such that $p\mathrel{R} q$.
                \begin{subproof}
                    Fix $p\in A$.
                    We have $p\in\naturalsPlus\times\pow{X}$ and
                        $\snd{p}\in T$ and
                        $\snd{p}\neq\emptyset$ and
                        $\closure{\snd{p}}{X}{T}\subseteq s(\fst{p})$ by \cref{subseteq}.
                    Take $n,W$ such that
                        $p = (n,W)$ and
                        $n\in\naturalsPlus$ and
                        $W\in\pow{X}$ by \cref{times}.
                    $W = \snd{p}$ and $n = \fst{p}$ by \cref{fst_eq,snd_eq}.
                    $W\in T$ and $W\neq\emptyset$ and $\closure{W}{X}{T}\subseteq s(n)$
                        by \cref{real_lt_subst_left,real_lt_subst_right}.
                    $W$ is open in $X$ with respect to $T$ by \cref{open_in}.
                    Take $z$ such that $z\in W$ by \cref{nonempty_inhabited}.
                    $s(n + 1)$ is open in $X$ with respect to $T$ and $s(n + 1)$ is dense in $X$ with respect to $T$
                        by \cref{real_naturals_closed_succ,subseteq}.
                    $s(n + 1)\subseteq X$ by \cref{open_in,topology_on,subseteq,pow_iff}.
                    $W\subseteq X$ by \cref{open_in,topology_on,subseteq,pow_iff}.
                    $z\in X$ by \cref{subseteq}.
                    Hence $z\in\closure{s(n + 1)}{X}{T}$ by \cref{dense_in,real_lt_subst_right}.
                    Therefore $W$ intersects $s(n + 1)$ by \cref{closure_iff_intersects_open}.
                    Take $y$ such that $y\in W\inter s(n + 1)$ by \cref{intersects,nonempty_inhabited}.
                    $y\in X$ by \cref{inter_elim_left,subseteq}.
                    $(1 / (n + 1)) / (1 + 1)\in\reals$ and $0 < (1 / (n + 1)) / (1 + 1)$
                        by \cref{real_naturals_closed_succ,half_reciprocal_naturalsplus_real_positive}.
                    Let $C = s(n + 1)\inter\metricBall{d}{X}{y}{(1 / (n + 1)) / (1 + 1)}$.
                    $\metricBall{d}{X}{y}{(1 / (n + 1)) / (1 + 1)}$ is open in $X$ with respect to $T$
                        by \cref{metric_ball,subseteq,powerset_intro,metric_basis,metric_on,metric_basis_is_basis,basis_elem_open_in_generated,metric_topology,basis_for_topology,open_in}.
                    $y\in\metricBall{d}{X}{y}{(1 / (n + 1)) / (1 + 1)}$
                        by \cref{metric_on,metric_zero_characterization,metric_ball,real_lt_subst_left}.
                    $C$ is open in $X$ with respect to $T$ and $y\in C$
                        by \cref{inter,open_in,topology_on,neighborhood,metric_ball}.
                    Let $B = W\inter C$.
                    $B$ is open in $X$ with respect to $T$ and $y\in B$
                        by \cref{inter,open_in,topology_on,neighborhood}.
                    Take $V$ such that
                        $V$ is a neighborhood of $y$ in $X$ with respect to $T$ and
                        $\closure{V}{X}{T}\subseteq B$
                        by \cref{inter,open_in,topology_on,neighborhood,regular_closure_neighborhood}.
                    $V$ is open in $X$ with respect to $T$ and $y\in V$ by \cref{neighborhood}.
                    Let $q = (n + 1,V)$.
                    $B\subseteq C$ by \cref{inter_lower_right}.
                    $C\subseteq s(n + 1)$ and $C\subseteq\metricBall{d}{X}{y}{(1 / (n + 1)) / (1 + 1)}$
                        by \cref{inter_lower_left,inter_lower_right}.
                    $\closure{V}{X}{T}\subseteq s(n + 1)$ by \cref{subseteq_transitive}.
                    Thus $\closure{V}{X}{T}$ is stage-small at $n + 1$ in $X$ with respect to $d$
                        by \cref{real_naturals_closed_succ,subseteq,metric_stage_small_set}.
                    $q\in A$ by
                        \cref{real_naturals_closed_succ,subseteq,neighborhood,open_in,topology_on,emptyset,pow_iff,powerset_intro,fst_eq,snd_eq,times_tuple_intro,pair_eq_iff,inter_lower_right,subseteq_transitive,metric_stage_small_set}.
                    $\closure{\snd{q}}{X}{T}\subseteq \snd{p}$ by \cref{fst_eq,snd_eq,inter_lower_left,subseteq_transitive}.
                    Hence $p\mathrel{R} q$ by \cref{fst_eq,snd_eq,times_tuple_intro,relation}.
                    Follows by definition.
                \end{subproof}
                Take $c$ such that $c$ is a function on $A$ and for all $p\in A$ we have $p\mathrel{R}\apply{c}{p}$
                    by \cref{choice_function}.
                We have $c\in\funs{A}{A}$ by \cref{choice_function,relation,times_tuple_elim,funs_intro}.
                Take $u$ such that
                    $u$ is a sequence in $A$ and
                    $u(1) = a_0$ and
                    for all $n\in\naturalsPlus$ we have $u(n + 1) = \apply{c}{u(n)}$
                    by \cref{iteration_sequence_naturalsplus}.
                Let $P = \{n\in\naturalsPlus \mid \fst{u(n)} = n\}$.
                $P\subseteq\naturalsPlus$ by \cref{subseteq}.
                $1\in P$ by \cref{fst_eq,real_one_in_naturals}.
                For all $n\in P$ we have $n + 1\in P$ by \cref{subseteq,sequence_in,funs_type_apply,choice_function,real_lt_subst_right,relation,fst_eq,snd_eq,real_naturals_closed_succ}.
                Therefore for all $n\in\naturalsPlus$ we have $n\in P$ by \cref{real_naturals_induction}.
                For all $n\in\naturalsPlus$ we have $\fst{u(n)} = n$ by \cref{subseteq}.
                Let $h = \{(n,\closure{\snd{u(n)}}{X}{T}) \mid n\in\naturalsPlus\}$.
                $h$ is a function by \cref{pair_eq_iff,rightunique,relation}.
                For all $n\in\naturalsPlus$ we have $n\in\dom{h}$ by \cref{dom_intro,pair_eq_iff}.
                For all $n$ such that $n\in\dom{h}$ we have $n\in\naturalsPlus$ by \cref{dom_iff,pair_eq_iff}.
                We have $\dom{h} = \naturalsPlus$ by \cref{subseteq,subseteq_antisymmetric}.
                For all $n$ such that $n\in\dom{h}$ we have $h(n)\in\pow{X}$
                    by \cref{subseteq,pair_eq_iff,function_apply_iff,real_lt_subst_right,closure,powerset_intro}.
                Therefore $h\in\funs{\naturalsPlus}{\pow{X}}$ by \cref{funs_intro}.
                $h$ is a sequence in $\pow{X}$ by \cref{sequence_in}.
                Show that for all $n\in\naturalsPlus$ we have
                    $h(n)$ is closed in $X$ with respect to $T$ and
                    $h(n)\neq\emptyset$.
                \begin{subproof}
                    Follows by \cref{sequence_in,funs_type_apply,subseteq,pair_eq_iff,function_apply_iff,nonempty_inhabited,times,snd_eq,powerset_elim,real_lt_subst_right,subseteq_closure,closure_closed_in,inhabited_nonempty}.
                \end{subproof}
                Show that for all $n\in\naturalsPlus$ we have $h(n + 1)\subseteq h(n)$.
                \begin{subproof}
                    Follows by \cref{sequence_in,funs_type_apply,real_naturals_closed_succ,subseteq,choice_function,real_lt_subst_right,relation,fst_eq,snd_eq,times,powerset_elim,subseteq_closure,subseteq_transitive,function_apply_intro}.
                \end{subproof}
                Show that for all $\epsilon\in\reals$ such that $0 < \epsilon$
                    there exists $N\in\naturalsPlus$ such that
                        for all $n\in\naturalsPlus$ if $N \leq n$ then
                            for all $x_1,y_1\in h(n)$ we have $d((x_1,y_1)) < \epsilon$.
                \begin{subproof}
                    Fix $\epsilon$ such that $\epsilon\in\reals$ and $0 < \epsilon$.
                    Take $N\in\naturalsPlus$ such that $0 < 1 / N$ and $1 / N < \epsilon$ by \cref{real_small_reciprocal}.
                    Show that for all $n\in\naturalsPlus$ if $N \leq n$ then
                        for all $x_1,y_1\in h(n)$ we have $d((x_1,y_1)) < \epsilon$.
                    \begin{subproof}
                        Fix $n\in\naturalsPlus$.
                        Assume $N \leq n$.
                        We have $u(n)\in A$ by \cref{sequence_in,funs_type_apply}.
                        $\closure{\snd{u(n)}}{X}{T}$ is stage-small at $n$ in $X$ with respect to $d$
                            by \cref{subseteq,real_lt_subst_right}.
                        $(n,\closure{\snd{u(n)}}{X}{T})\in h$ by \cref{pair_eq_iff}.
                        Hence $h(n) = \closure{\snd{u(n)}}{X}{T}$ by \cref{function_apply_iff}.
                        Take $a\in X$ such that $h(n)\subseteq\metricBall{d}{X}{a}{(1 / n) / (1 + 1)}$
                            by \cref{real_lt_subst_right,metric_stage_small_set}.
                        For all $x_1\in h(n)$ we have
                            for all $y_1\in h(n)$ we have $d((x_1,y_1)) < \epsilon$ by \cref{subseteq,sequence_in,funs_type_apply,powerset_elim,naturalsplus_reciprocal_real,positive_reciprocal_naturalsplus,metric_half_ball_pair_bound,naturalsplus_reciprocal_antitone,metric_on,times_tuple_intro,real_leq_split,real_lt_subst_right,real_lt_trans}.
                    \end{subproof}
                    Follows by definition.
                \end{subproof}
                Therefore $\inters{\img{h}{\naturalsPlus}}\neq\emptyset$
                    by \cref{nested_closed_sets_eventually_small_nonempty_intersection}.
                Take $p$ such that $p\in\inters{\img{h}{\naturalsPlus}}$ by \cref{nonempty_inhabited}.
                $p\in U$ by \cref{sequence_in,function_apply_elim,img_iff,real_one_in_naturals,inters_destr,function_apply_intro,subseteq,snd_eq,real_lt_subst_left,real_lt_subst_right,inter}.
                Show that for all $A_1$ such that $A_1\in\img{s}{\naturalsPlus}$ we have
                    $p\in A_1$.
                \begin{subproof}
                    Follows by \cref{img_iff,sequence_in,funs_is_function,function_apply_iff,function_apply_elim,inters_destr,funs_type_apply,subseteq,function_apply_intro,real_lt_subst_right,real_lt_subst_left}.
                \end{subproof}
                Hence $p\in\inters{\img{s}{\naturalsPlus}}$ by
                    \cref{sequence_in,funs,funs_is_function,function_apply_iff,img_iff,real_one_in_naturals,inters_intro}.
                $p\in U\inter\inters{\img{s}{\naturalsPlus}}$ by \cref{inter_intro}.
                Therefore $U$ intersects $\inters{\img{s}{\naturalsPlus}}$ by \cref{emptyset,intersects,inter}.
            \end{subproof}
            Thus $x\in\closure{\inters{\img{s}{\naturalsPlus}}}{X}{T}$ by \cref{closure_iff_intersects_open}.
        \end{subproof}
        $X\subseteq\closure{\inters{\img{s}{\naturalsPlus}}}{X}{T}$ by \cref{subseteq}.
        Hence $\closure{\inters{\img{s}{\naturalsPlus}}}{X}{T} = X$ by \cref{closure,subseteq,subseteq_antisymmetric}.
        $\inters{\img{s}{\naturalsPlus}}$ is dense in $X$ with respect to $T$ by \cref{dense_in}.
    \end{subproof}
    Therefore $X$ is a Baire space with respect to $T$ by \cref{baire_space_open_set_characterization}.
\end{proof}

% from §48 Item 5: nested closed sets with shrinking diameters have nonempty intersection
\begin{lemma}\label{nested_closed_sets_shrinking_diameters_nonempty_intersection}
    Assume $d$ is a metric on $X$.
    Assume $X$ is complete with respect to $d$.
    Assume $s$ is a sequence in $\pow{X}$.
    Suppose for all $n\in\naturalsPlus$ we have
        $s(n)$ is closed in $X$ with respect to $\metricTopology{d}{X}$ and
        $s(n)\neq\emptyset$.
    Suppose for all $n\in\naturalsPlus$ we have
        $s(n+1)\subseteq s(n)$.
    Suppose for all $\epsilon\in\reals$ such that $0 < \epsilon$
        there exists $N\in\naturalsPlus$ such that
            for all $n\in\naturalsPlus$ if $N \leq n$ then
                there exists $r$ such that
                    $r$ is a diameter of $s(n)$ in $X$ with respect to $d$ and
                    $r < \epsilon$.
    Then $\inters{\img{s}{\naturalsPlus}}\neq\emptyset$.
\end{lemma}
\begin{proof}
    Show that for all $\epsilon\in\reals$ such that $0 < \epsilon$
        there exists $N\in\naturalsPlus$ such that
            for all $n\in\naturalsPlus$ if $N \leq n$ then
                for all $x,y\in s(n)$ we have $d((x,y)) < \epsilon$.
    \begin{subproof}
        Fix $\epsilon$ such that $\epsilon\in\reals$ and $0 < \epsilon$.
        Take $N\in\naturalsPlus$ such that
            for all $n\in\naturalsPlus$ if $N \leq n$ then
                there exists $r$ such that
                    $r$ is a diameter of $s(n)$ in $X$ with respect to $d$ and
                    $r < \epsilon$ by \cref{subseteq}.
        Show that for all $n\in\naturalsPlus$ if $N \leq n$ then
            for all $x,y\in s(n)$ we have $d((x,y)) < \epsilon$.
        \begin{subproof}
            Fix $n\in\naturalsPlus$.
            Assume $N \leq n$.
            Take $r$ such that
                $r$ is a diameter of $s(n)$ in $X$ with respect to $d$ and
                $r < \epsilon$ by \cref{subseteq}.
            Show that for all $x,y$ such that $x\in s(n)$ and $y\in s(n)$ we have
                $d((x,y)) < \epsilon$.
            \begin{subproof}
                Fix $x$.
                Fix $y$.
                Assume $x\in s(n)$ and $y\in s(n)$.
                We have $r\in\reals$ by \cref{metric_diameter,real_supremum,real_upper_bound}.
                $s(n)\subseteq X$ by \cref{sequence_in,funs_type_apply,powerset_elim}.
                $x\in X$ and $y\in X$ by \cref{subseteq}.
                $d((x,y))\in\reals$ by \cref{metric_on,times_tuple_intro,funs_type_apply}.
                $d((x,y)) \leq r$ by \cref{metric_diameter_bound,real_leq_split}.
                Hence $d((x,y)) < \epsilon$ by \cref{real_lt_trans,real_lt_subst_left}.
            \end{subproof}
        \end{subproof}
        Follows by definition.
    \end{subproof}
    Therefore $\inters{\img{s}{\naturalsPlus}}\neq\emptyset$
        by \cref{nested_closed_sets_eventually_small_nonempty_intersection}.
\end{proof}

% from §48 Item 6: open subspaces of Baire spaces are Baire
\begin{lemma}\label{open_subspace_of_baire_space}
    Assume $X$ is a Baire space with respect to $T$.
    Assume $Y$ is open in $X$ with respect to $T$.
    Then $Y$ is a Baire space with respect to $\subspaceTopology{T}{Y}$.
\end{lemma}
\begin{proof}
    Let $T_Y = \subspaceTopology{T}{Y}$.
    We have $T$ is a topology on $X$ by \cref{baire_space}.
    $Y\subseteq X$ by \cref{open_in,topology_on,subseteq,pow_iff}.
    $T_Y$ is a topology on $Y$ by \cref{subspace_topology_is_topology,subspace_of}.
    Show that for all $s$ if
        $s$ is a sequence in $\pow{Y}$ and
        for all $n\in\naturalsPlus$ we have
            $s(n)$ is open in $Y$ with respect to $T_Y$ and
            $s(n)$ is dense in $Y$ with respect to $T_Y$
    then $\inters{\img{s}{\naturalsPlus}}$ is dense in $Y$ with respect to $T_Y$.
    \begin{subproof}
        Fix $s$ such that
            $s$ is a sequence in $\pow{Y}$ and
            for all $n\in\naturalsPlus$ we have
                $s(n)$ is open in $Y$ with respect to $T_Y$ and
                $s(n)$ is dense in $Y$ with respect to $T_Y$.
        Let $t = \{(n, s(n)\union(X\setminus\closure{Y}{X}{T})) \mid n\in\naturalsPlus\}$.
        We have $t$ is a function by \cref{pair_eq_iff,rightunique,relation}.
        For all $n$ such that $n\in\naturalsPlus$ we have $n\in\dom{t}$
            by \cref{pair_eq_iff,dom_intro}.
        For all $n$ such that $n\in\dom{t}$ we have $n\in\naturalsPlus$
            by \cref{dom_iff,pair_eq_iff}.
        $\dom{t} = \naturalsPlus$ by \cref{subseteq,subseteq_antisymmetric}.
    Hence $t$ is a sequence in $\pow{X}$
            by \cref{dom_iff,pair_eq_iff,sequence_in,funs_elim,powerset_elim,subseteq_transitive,setminus_subseteq,union_subsets_is_subset,function_apply_intro,real_lt_subst_left,powerset_intro,funs_intro,sequence_in}.
        Show that for all $n\in\naturalsPlus$ we have
            $t(n)$ is open in $X$ with respect to $T$ and
            $t(n)$ is dense in $X$ with respect to $T$.
        \begin{subproof}
            Fix $n\in\naturalsPlus$.
            We have $s(n)$ is open in $Y$ with respect to $T_Y$ and $s(n)$ is dense in $Y$ with respect to $T_Y$
                by \cref{subseteq}.
            $s(n)$ is open in $X$ with respect to $T$ by \cref{open_in_subspace_open_in_space,subspace_of}.
            $X\setminus\closure{Y}{X}{T}$ is open in $X$ with respect to $T$ by \cref{closure_closed_in,closed_in}.
            $(n, s(n)\union(X\setminus\closure{Y}{X}{T}))\in t$ by \cref{pair_eq_iff}.
            $t(n) = s(n)\union(X\setminus\closure{Y}{X}{T})$ by \cref{sequence_in,funs_elim,function_apply_intro}.
            Let $M = \{s(n), X\setminus\closure{Y}{X}{T}\}$.
            Show that $t(n)$ is open in $X$ with respect to $T$.
            \begin{subproof}
                Follows by \cref{upair_elim,open_in,subseteq,topology_on,real_lt_subst_left,union_as_unions,real_lt_subst_right}.
            \end{subproof}
            Show that for all $x\in X$ we have $x\in\closure{t(n)}{X}{T}$.
            \begin{subproof}
                Fix $x\in X$.
                Show that for all $U$ such that
                    $U$ is open in $X$ with respect to $T$ and
                    $x\in U$
                    we have $U$ intersects $t(n)$.
                \begin{subproof}
                    Fix $U$ such that
                        $U$ is open in $X$ with respect to $T$ and
                        $x\in U$.
                    \begin{byCase}
                    \caseOf{$U$ intersects $Y$.}
                        Take $z$ such that $z\in U\inter Y$ by \cref{intersects,nonempty_inhabited}.
                        $U\inter Y$ is open in $Y$ with respect to $T_Y$
                            by \cref{open_in,subspace_topology,inter_comm}.
                        $z\in Y$ by \cref{inter}.
                        Hence $z\in\closure{s(n)}{Y}{T_Y}$ by \cref{dense_in,real_lt_subst_left}.
                        $U\inter Y$ intersects $s(n)$ by \cref{sequence_in,funs_elim,powerset_elim,closure_iff_intersects_open}.
                        Take $w$ such that $w\in (U\inter Y)\inter s(n)$ by \cref{intersects,nonempty_inhabited}.
                        $w\in U$ and $w\in s(n)$ by \cref{inter}.
                        $w\in t(n)$ by \cref{real_lt_subst_right,union_iff}.
                        $w\in U\inter t(n)$ by \cref{inter_intro}.
                        Therefore $U$ intersects $t(n)$ by \cref{emptyset,intersects,inter}.
                    \caseOf{$U$ does not intersect $Y$.}
                        $x\notin\closure{Y}{X}{T}$ by \cref{closure_iff_intersects_open}.
                        $x\in X\setminus\closure{Y}{X}{T}$ by \cref{setminus_intro}.
                        $x\in t(n)$ by \cref{real_lt_subst_right,union_iff}.
                        $x\in U\inter t(n)$ by \cref{inter_intro}.
                        Hence $U$ intersects $t(n)$ by \cref{emptyset,intersects,inter}.
                    \end{byCase}
                \end{subproof}
                Hence $x\in\closure{t(n)}{X}{T}$ by \cref{sequence_in,funs_elim,powerset_elim,closure_iff_intersects_open}.
            \end{subproof}
            $X\subseteq\closure{t(n)}{X}{T}$ by \cref{subseteq}.
            Therefore $\closure{t(n)}{X}{T} = X$ by \cref{closure,subseteq,subseteq_antisymmetric}.
            $t(n)$ is dense in $X$ with respect to $T$ by \cref{dense_in}.
        \end{subproof}
        Hence $\inters{\img{t}{\naturalsPlus}}$ is dense in $X$ with respect to $T$
            by \cref{baire_space_open_set_characterization,baire_space}.
        $\img{s}{\naturalsPlus}\subseteq\pow{Y}$ by \cref{img_subseteq_ran,sequence_in,funs_ran,subseteq_transitive}.
        $\inters{\img{s}{\naturalsPlus}}\subseteq Y$ by
            \cref{unions_subseteq_of_powerset_is_subseteq,inters,elem_subseteq,subseteq}.
        Show that for all $y\in Y$ we have $y\in\closure{\inters{\img{s}{\naturalsPlus}}}{Y}{T_Y}$.
        \begin{subproof}
            Fix $y\in Y$.
            Show that for all $W$ such that
                $W$ is open in $Y$ with respect to $T_Y$ and
                $y\in W$
                we have $W$ intersects $\inters{\img{s}{\naturalsPlus}}$.
            \begin{subproof}
                Fix $W$ such that
                    $W$ is open in $Y$ with respect to $T_Y$ and
                    $y\in W$.
                $\img{t}{\naturalsPlus}\subseteq\pow{X}$ by \cref{img_subseteq_ran,sequence_in,funs_ran,subseteq_transitive}.
                $\inters{\img{t}{\naturalsPlus}}\subseteq X$ by
                    \cref{unions_subseteq_of_powerset_is_subseteq,inters,elem_subseteq,subseteq}.
                We have $W$ is open in $X$ with respect to $T$ and $y\in X$ by \cref{open_in_subspace_open_in_space,subspace_of,subseteq}.
        $y\in\closure{\inters{\img{t}{\naturalsPlus}}}{X}{T}$ by \cref{dense_in,real_lt_subst_left}.
                Therefore $W$ intersects $\inters{\img{t}{\naturalsPlus}}$ by \cref{closure_iff_intersects_open}.
                Take $p$ such that $p\in W\inter\inters{\img{t}{\naturalsPlus}}$ by \cref{intersects,nonempty_inhabited}.
                $p\in W$ and $p\in\inters{\img{t}{\naturalsPlus}}$ by \cref{inter}.
                $p\in Y$ by \cref{open_in,topology_on,subseteq,pow_iff,inter}.
                Show that for all $A$ such that $A\in\img{s}{\naturalsPlus}$ we have $p\in A$.
                \begin{subproof}
                    Follows by \cref{img_iff,sequence_in,funs_is_function,function_apply_iff,pair_eq_iff,funs_elim,function_apply_intro,function_apply_elim,inters_destr,subseteq_closure,not_in_subseteq,setminus_elim_right,union_iff,real_lt_subst_left}.
                \end{subproof}
                Hence $p\in\inters{\img{s}{\naturalsPlus}}$ by
                    \cref{sequence_in,funs,funs_is_function,function_apply_iff,img_iff,real_one_in_naturals,inters_intro}.
                $p\in W\inter\inters{\img{s}{\naturalsPlus}}$ by \cref{inter_intro}.
                Therefore $W$ intersects $\inters{\img{s}{\naturalsPlus}}$ by \cref{emptyset,intersects,inter}.
            \end{subproof}
            Thus $y\in\closure{\inters{\img{s}{\naturalsPlus}}}{Y}{T_Y}$ by \cref{closure_iff_intersects_open}.
        \end{subproof}
        $Y\subseteq\closure{\inters{\img{s}{\naturalsPlus}}}{Y}{T_Y}$ by \cref{subseteq}.
        Hence $\closure{\inters{\img{s}{\naturalsPlus}}}{Y}{T_Y} = Y$ by \cref{closure,subseteq,subseteq_antisymmetric}.
        $\inters{\img{s}{\naturalsPlus}}$ is dense in $Y$ with respect to $T_Y$ by \cref{dense_in}.
    \end{subproof}
    Therefore $Y$ is a Baire space with respect to $T_Y$ by \cref{baire_space_open_set_characterization}.
\end{proof}

% from §48 helper definition 2: Cauchy stage set for a function sequence
\begin{definition}\label{function_sequence_cauchy_stage_set}
    $A$ is the $k$-th Cauchy stage set of $s$ at $N$ from $X$ to $Y$ with respect to $d$ iff
        $s$ is a sequence in $\funs{X}{Y}$ and
        $k\in\naturalsPlus$ and
        $N\in\naturalsPlus$ and
        $d$ is a metric on $Y$ and
        $A = \{x\in X \mid \forall n,m\in\naturalsPlus. (N \leq n \land N \leq m) \implies
            d((\apply{\apply{s}{n}}{x},\apply{\apply{s}{m}}{x})) \leq 1 / k\}$.
\end{definition}

% from §48 helper lemma 14: reverse triangle inequality for two pairs
\begin{lemma}\label{metric_reverse_triangle_two_pairs}
    Assume $d$ is a metric on $X$.
    Assume $x_0,x_1,y_0,y_1\in X$.
    Then $\apply{\standardMetricR}{(d((x_0,y_0)),d((x_1,y_1)))} \leq d((x_0,x_1)) + d((y_0,y_1))$.
\end{lemma}
\begin{proof}
    Let $a = d((x_0,x_1))$.
    Let $b = d((y_0,y_1))$.
    Let $s = a + b$.
    Let $u = d((x_0,y_0))$.
    Let $v = d((x_1,y_1))$.
    $(x_0,x_1)\in X\times X$ and $(y_0,y_1)\in X\times X$ and
        $(x_0,y_0)\in X\times X$ and $(x_1,y_1)\in X\times X$
        by \cref{times_tuple_intro}.
    We have $(x_1,y_0)\in X\times X$ and $(x_0,y_1)\in X\times X$ by \cref{times_tuple_intro}.
    $d\in\funs{X\times X}{\reals}$ by \cref{metric_on}.
    $a\in\reals$ and $b\in\reals$ and $u\in\reals$ and $v\in\reals$ by \cref{funs_type_apply}.
    $d((x_1,y_0))\in\reals$ and $d((x_0,y_1))\in\reals$ by \cref{funs_type_apply}.
    $s\in\reals$ by \cref{real_add_closed}.
    $d((x_0,y_0)) \leq d((x_0,x_1)) + d((x_1,y_0))$ by \cref{metric_on,metric_triangle_inequality}.
    $u \leq a + d((x_1,y_0))$.
    $d((y_1,y_0)) = d((y_0,y_1))$ by \cref{metric_on,metric_symmetric}.
    Hence $d((x_1,y_0)) \leq d((x_1,y_1)) + b$ by \cref{metric_on,metric_triangle_inequality}.
    $d((x_1,y_0)) \leq v + b$.
    $v + b\in\reals$ by \cref{real_add_closed}.
    $d((x_1,y_0)) + a \leq (v + b) + a$ by \cref{real_leq_add}.
    $d((x_1,y_0)) + a = a + d((x_1,y_0))$ and $(v + b) + a = a + (v + b)$
        by \cref{real_add_comm}.
    $a + d((x_1,y_0))\in\reals$ and $a + (v + b)\in\reals$ by \cref{real_add_closed}.
    $a + d((x_1,y_0)) \leq a + (v + b)$ by \cref{real_leq_trans}.
    Thus $u \leq a + (v + b)$ by \cref{real_leq_trans}.
    $a + (v + b) = (a + b) + v$ by \cref{real_add_assoc,real_add_comm}.
    $a + (v + b) = s + v$ by \cref{real_lt_subst_left}.
    $s + v\in\reals$ and $\neg{v}\in\reals$ by \cref{real_add_closed,real_add_inverse}.
    $u \leq s + v$ by \cref{real_leq_trans}.
    $u + \neg{v} \leq (s + v) + \neg{v}$ by \cref{real_add_inverse,real_add_closed,real_leq_add}.
    Hence $u - v \leq s$ by \cref{real_add_assoc,real_add_inverse,real_lt_subst_right,real_add_ident,real_add_comm,real_sub,real_leq_trans}.

    $d((x_1,y_1)) \leq d((x_1,x_0)) + d((x_0,y_1))$ by \cref{metric_on,metric_triangle_inequality}.
    $d((x_1,x_0)) = d((x_0,x_1))$ by \cref{metric_on,metric_symmetric}.
    $v \leq a + d((x_0,y_1))$.
    $d((x_0,y_1)) \leq d((x_0,y_0)) + d((y_0,y_1))$ by \cref{metric_on,metric_triangle_inequality}.
    Thus $d((x_0,y_1)) \leq u + b$.
    $u + b\in\reals$ by \cref{real_add_closed}.
    $d((x_0,y_1)) + a \leq (u + b) + a$ by \cref{real_leq_add}.
    $d((x_0,y_1)) + a = a + d((x_0,y_1))$ and $(u + b) + a = a + (u + b)$
        by \cref{real_add_comm}.
    $a + d((x_0,y_1))\in\reals$ and $a + (u + b)\in\reals$ by \cref{real_add_closed}.
    $a + d((x_0,y_1)) \leq a + (u + b)$ by \cref{real_leq_trans}.
    Hence $v \leq a + (u + b)$ by \cref{real_leq_trans}.
    $a + (u + b) = (a + b) + u$ by \cref{real_add_assoc,real_add_comm}.
    $a + (u + b) = s + u$ by \cref{real_lt_subst_left}.
    $s + u\in\reals$ and $\neg{u}\in\reals$ by \cref{real_add_closed,real_add_inverse}.
    $v \leq s + u$ by \cref{real_leq_trans}.
    $v + \neg{u} \leq (s + u) + \neg{u}$ by \cref{real_add_inverse,real_add_closed,real_leq_add}.
    Thus $v - u \leq s$ by \cref{real_add_assoc,real_add_inverse,real_lt_subst_right,real_add_ident,real_add_comm,real_sub,real_leq_trans}.

    $u + \neg{v}\in\reals$ and $v + \neg{u}\in\reals$ by \cref{real_add_closed}.
    $u - v = u + \neg{v}$ and $v - u = v + \neg{u}$ by \cref{real_sub}.
    $u - v\in\reals$ and $v - u\in\reals$ by \cref{real_lt_subst_left}.
    $v - u = \neg{(u - v)}$ by \cref{real_sub_swap_neg}.
    Therefore $\abs{(u - v)} \leq s$ by \cref{real_abs_leq_if_pm_leq}.
    Let $p = (u,v)$.
    $p\in\reals\times\reals$ by \cref{times_tuple_intro}.
    $\fst{p} = u$ and $\snd{p} = v$ by \cref{fst_eq,snd_eq}.
    $\abs{(u - v)}\in\reals$ by \cref{real_abs_in_reals}.
    $\abs{(\fst{p} - \snd{p})} = \abs{(u - v)}$.
    $(p,\abs{(u - v)})\in\standardMetricR$ by \cref{standard_metric_r}.
    Hence $\apply{\standardMetricR}{(u,v)} = \abs{(u - v)}$
        by \cref{standard_metric_is_metric,metric_on,funs_elim,function_apply_intro}.
    Thus $\apply{\standardMetricR}{(u,v)} \leq a + b$ by \cref{real_lt_subst_right}.
    Therefore $\apply{\standardMetricR}{(d((x_0,y_0)),d((x_1,y_1)))} \leq d((x_0,x_1)) + d((y_0,y_1))$
        by \cref{real_lt_subst_left,real_lt_subst_right}.
\end{proof}

% from §48 helper lemma 15: pointwise distance of continuous maps is continuous
\begin{lemma}\label{continuous_maps_pointwise_distance}
    Assume $T$ is a topology on $X$.
    Assume $d$ is a metric on $Y$.
    Assume $f$ is continuous from $X$ to $Y$ with respect to $T$ and $\metricTopology{d}{Y}$.
    Assume $g$ is continuous from $X$ to $Y$ with respect to $T$ and $\metricTopology{d}{Y}$.
    Let $h = \{(x, d((f(x),g(x)))) \mid x\in X\}$.
    Then $h$ is continuous from $X$ to $\reals$ with respect to $T$ and $\standardTopologyR$.
\end{lemma}
\begin{proof}
    Let $U = \metricTopology{d}{Y}$.
    $f$ is a function and $g$ is a function by \cref{continuous,funs_is_function}.
    Hence $h$ is a function by \cref{pair_eq_iff,rightunique,relation}.
    For all $x,z$ such that $x\in\dom{h}$ and $(x,z)\in h$ we have $x\in X$ and $z = d((f(x),g(x)))$
        by \cref{pair_eq_iff}.
    For all $x\in\dom{h}$ we have $h(x)\in\reals$
        by \cref{dom_iff,pair_eq_iff,continuous,funs_type_apply,times_tuple_intro,metric_on,funs_is_function,function_apply_intro,real_lt_subst_left,real_lt_subst_right}.
    For all $x\in\dom{h}$ we have $x\in X$ and for all $x\in X$ we have $(x,d((f(x),g(x))))\in h$.
    Therefore $h\in\funs{X}{\reals}$ by \cref{subseteq,dom_intro,subseteq_antisymmetric,funs_intro}.
    Show that for all $x\in X$ we have $h$ is continuous at $x$ from $X$ to $\reals$ with respect to $T$ and $\standardTopologyR$.
    \begin{subproof}
        Fix $x\in X$.
        $f$ is continuous at $x$ from $X$ to $Y$ with respect to $T$ and $U$ and
            $g$ is continuous at $x$ from $X$ to $Y$ with respect to $T$ and $U$
            by \cref{continuous_implies_continuous_at}.
        Show that for all $V$ such that $V$ is open in $\reals$ with respect to $\standardTopologyR$ and $h(x)\in V$
            there exists $W$ such that $W$ is open in $X$ with respect to $T$ and $x\in W$ and $\img{h}{W}\subseteq V$.
        \begin{subproof}
            Fix $V$ such that $V$ is open in $\reals$ with respect to $\standardTopologyR$ and $h(x)\in V$.
            Let $T_R = \metricTopology{\standardMetricR}{\reals}$.
            We have $\standardMetricR$ is a metric on $\reals$ by \cref{standard_metric_is_metric}.
            $V$ is open in $\reals$ with respect to $T_R$
                by \cref{standard_metric_is_metric,standard_metric_topology,real_lt_subst_right}.
            $V\subseteq\reals$ by
                \cref{open_in,basis_topology_is_topology,metric_topology,metric_basis_is_basis,basis_for_topology,topology_on,subseteq,pow_iff}.
            Take $\epsilon\in\reals$ such that $0 < \epsilon$ and
                $\metricBall{\standardMetricR}{\reals}{h(x)}{\epsilon} \subseteq V$
                by \cref{subseteq,metric_open_iff}.
            Let $\epsilon_0 = \epsilon / (1 + 1)$.
            $\epsilon_0\in\reals$ and $0 < \epsilon_0$ and $\epsilon_0 + \epsilon_0 = \epsilon$
                by \cref{real_half_of_positive}.
            Let $B_f = \metricBall{d}{Y}{f(x)}{\epsilon_0}$.
            Let $B_g = \metricBall{d}{Y}{g(x)}{\epsilon_0}$.
            $f(x)\in Y$ and $g(x)\in Y$ by \cref{continuous,funs_type_apply}.
            $B_f$ is open in $Y$ with respect to $U$ and $B_g$ is open in $Y$ with respect to $U$
                by \cref{metric_ball,subseteq,powerset_intro,metric_basis,metric_basis_is_basis,basis_elem_open_in_generated,basis_for_topology,metric_topology,basis_topology_is_topology,open_in}.
            $f(x)\in B_f$ and $g(x)\in B_g$
                by \cref{metric_zero_characterization,metric_on,continuous,funs_type_apply,real_lt_subst_left,metric_ball}.
            Hence $B_f$ is a neighborhood of $f(x)$ in $Y$ with respect to $U$ and
                $B_g$ is a neighborhood of $g(x)$ in $Y$ with respect to $U$ by \cref{neighborhood}.
            Take $U_f$ such that
                $U_f$ is a neighborhood of $x$ in $X$ with respect to $T$ and
                $\img{f}{U_f}\subseteq B_f$ by \cref{continuous_at}.
            Take $U_g$ such that
                $U_g$ is a neighborhood of $x$ in $X$ with respect to $T$ and
                $\img{g}{U_g}\subseteq B_g$ by \cref{continuous_at}.
            $U_f$ is open in $X$ with respect to $T$ and $x\in U_f$ and
                $U_g$ is open in $X$ with respect to $T$ and $x\in U_g$ by \cref{neighborhood}.
            Let $W = U_f\inter U_g$.
            $W$ is open in $X$ with respect to $T$ and $x\in W$ by \cref{open_in,topology_on,inter_intro}.
            Show that for all $z$ such that $z\in\img{h}{W}$ we have $z\in V$.
            \begin{subproof}
                Fix $z$ such that $z\in\img{h}{W}$.
                Take $y\in\dom{h}\inter W$ such that $z = h(y)$ by \cref{img_of_function_elim}.
                We have $y\in X$ and $y\in U_f$ and $y\in U_g$ by \cref{funs,inter}.
                Therefore $f(y)\in B_f$ and $g(y)\in B_g$
                    by \cref{continuous,funs,inter_intro,img_of_function_intro,subseteq}.
                $d((f(x),f(y))) < \epsilon_0$ and $d((g(x),g(y))) < \epsilon_0$ by \cref{metric_ball}.
                $h(x) = d((f(x),g(x)))$ and $h(y) = d((f(y),g(y)))$ by \cref{funs_is_function,function_apply_intro}.
                Hence $\apply{\standardMetricR}{(h(x),h(y))} \leq d((f(x),f(y))) + d((g(x),g(y)))$
                    by \cref{continuous,funs_type_apply,metric_reverse_triangle_two_pairs}.
                $(f(x),f(y))\in Y\times Y$ and $(g(x),g(y))\in Y\times Y$
                    by \cref{continuous,funs_type_apply,times_tuple_intro}.
                $d((f(x),f(y)))\in\reals$ and $d((g(x),g(y)))\in\reals$ and
                    $h(x)\in\reals$ and $h(y)\in\reals$ by \cref{metric_on,funs_type_apply}.
                $\apply{\standardMetricR}{(h(x),h(y))}\in\reals$ and
                    $d((f(x),f(y))) + d((g(x),g(y)))\in\reals$ and $\epsilon_0 + \epsilon_0\in\reals$
                    by \cref{standard_metric_is_metric,metric_on,funs_type_apply,times_tuple_intro,real_add_closed}.
                $\apply{\standardMetricR}{(h(x),h(y))} < \epsilon_0 + \epsilon_0$
                    by \cref{real_lt_add_two,real_leq_split,real_lt_trans,real_lt_subst_left}.
                Therefore $\apply{\standardMetricR}{(h(x),h(y))} < \epsilon$ by \cref{real_lt_subst_right}.
                Hence $z\in\metricBall{\standardMetricR}{\reals}{h(x)}{\epsilon}$
                    by \cref{funs_type_apply,metric_ball,real_lt_subst_left}.
                $z\in V$ by \cref{subseteq}.
            \end{subproof}
            $\img{h}{W}\subseteq V$ by \cref{subseteq}.
            There exists $U_0$ such that
                $U_0$ is open in $X$ with respect to $T$ and
                $x\in U_0$ and
                $\img{h}{U_0}\subseteq V$.
        \end{subproof}
        Hence $h$ is continuous at $x$ from $X$ to $\reals$ with respect to $T$ and $\standardTopologyR$
            by \cref{continuous_at_open,standard_basis_is_basis,basis_topology_is_topology,standard_topology_reals}.
    \end{subproof}
    We have $\standardTopologyR$ is a topology on $\reals$ by \cref{standard_basis_is_basis,basis_topology_is_topology,standard_topology_reals}.
    Therefore $h$ is continuous from $X$ to $\reals$ with respect to $T$ and $\standardTopologyR$
        by \cref{continuous_at_implies_continuous}.
\end{proof}

% from §48 helper lemma 16: convergent metric sequences are Cauchy
\begin{lemma}\label{metric_convergent_sequence_cauchy}
    Assume $d$ is a metric on $X$.
    Assume $s$ is a sequence in $X$.
    Assume $x\in X$.
    Suppose $s$ converges to $x$ in $X$ with respect to $\metricTopology{d}{X}$.
    Then $s$ is a Cauchy sequence in $X$ with respect to $d$.
\end{lemma}
\begin{proof}
    Show that for all $\epsilon\in\reals$ such that $0 < \epsilon$
        there exists $N\in\naturalsPlus$ such that
            for all $n,m\in\naturalsPlus$ if $N \leq n$ and $N \leq m$ then
                $d((s(n),s(m))) < \epsilon$.
    \begin{subproof}
        Fix $\epsilon$ such that $\epsilon\in\reals$ and $0 < \epsilon$.
        Let $\delta = \epsilon / (1 + 1)$.
        $\delta\in\reals$ and $0 < \delta$ and $\delta + \delta = \epsilon$ by \cref{real_half_of_positive}.
        Let $B = \metricBall{d}{X}{x}{\delta}$.
        $B$ is a neighborhood of $x$ in $X$ with respect to $\metricTopology{d}{X}$
            by \cref{metric_ball,subseteq,powerset_intro,metric_basis,metric_basis_is_basis,basis_elem_open_in_generated,basis_for_topology,metric_topology,basis_topology_is_topology,open_in,metric_zero_characterization,metric_on,real_lt_subst_left,neighborhood}.
        Take $N\in\naturalsPlus$ such that for all $n\in\naturalsPlus$ if $N \leq n$ then $s(n)\in B$
            by \cref{converges_to}.
        For all $n\in\naturalsPlus$ if $N \leq n$ then
            for all $m\in\naturalsPlus$ if $N \leq m$ then
                $d((s(n),s(m))) < \epsilon$ by \cref{subseteq,sequence_in,funs_type_apply,metric_half_ball_pair_bound}.
        Follows by definition.
    \end{subproof}
    Therefore $s$ is a Cauchy sequence in $X$ with respect to $d$ by \cref{cauchy_sequence}.
\end{proof}

% from §48 helper lemma 17: closed intervals in the standard topology are closed
\begin{lemma}\label{standard_closed_interval_closed}
    Assume $r\in\reals$.
    Then $\intervalclosed{0}{r}$ is closed in $\reals$ with respect to $\standardTopologyR$.
\end{lemma}
\begin{proof}
    Let $I = \intervalclosed{0}{r}$.
    $\realOrderRel$ is a simple order relation on $\reals$ and $\reals$ has more than one element
        by \cref{reals_linear_continuum,linear_continuum}.
    Hence $\realOrderTopology$ is the order topology on $\reals$ with respect to $\realOrderRel$
        by \cref{real_order_topology,order_topology}.
    $0\in\reals$ by \cref{real_add_ident}.
    Show that $I\subseteq \intervalclosedrel{\realOrderRel}{\reals}{0}{r}$.
    \begin{subproof}
        Follows by \cref{intervalclosed,real_add_ident,real_order_relation_intro,intervalclosed_rel,subseteq}.
    \end{subproof}
    Show that $\intervalclosedrel{\realOrderRel}{\reals}{0}{r}\subseteq I$.
    \begin{subproof}
        Follows by \cref{intervalclosed_rel,real_order_relation_elim,intervalclosed,subseteq}.
    \end{subproof}
    Therefore $\intervalclosed{0}{r}$ is closed in $\reals$ with respect to $\standardTopologyR$
        by \cref{subseteq_antisymmetric,intervalclosedrel_closed_in_order_topology,real_lt_subst_left,real_order_topology_eq_standard,real_lt_subst_right}.
\end{proof}

% from §48 helper lemma 18: closed metric balls are closed
\begin{lemma}\label{metric_closed_ball_closed}
    Assume $d$ is a metric on $Y$.
    Assume $a\in Y$.
    Assume $r\in\reals$.
    Let $B = \{y\in Y \mid d((y,a)) \leq r\}$.
    Then $B$ is closed in $Y$ with respect to $\metricTopology{d}{Y}$.
\end{lemma}
\begin{proof}
    Let $T = \metricTopology{d}{Y}$.
    Let $j = \identity{Y}$.
    $T$ is a topology on $Y$ by \cref{metric_topology,metric_basis_is_basis,basis_for_topology,basis_topology_is_topology}.
    $j$ is continuous from $Y$ to $Y$ with respect to $T$ and $T$ by \cref{identity_continuous}.
    Let $c = \{(y,a) \mid y\in Y\}$.
    $c$ is a function from $Y$ to $Y$ by \cref{constant_map_function}.
    For all $y\in Y$ we have $c(y) = a$ by \cref{funs,funs_is_function,function_apply_intro}.
    Therefore $c$ is continuous from $Y$ to $Y$ with respect to $T$ and $T$ by \cref{constant_continuous}.
    Let $h = \{(y, d((j(y),c(y)))) \mid y\in Y\}$.
    $h$ is continuous from $Y$ to $\reals$ with respect to $T$ and $\standardTopologyR$
        by \cref{continuous_maps_pointwise_distance}.
    $h\in\funs{Y}{\reals}$ by \cref{continuous}.
    $\standardTopologyR$ is a topology on $\reals$
        by \cref{standard_basis_is_basis,basis_topology_is_topology,standard_topology_reals}.
    Let $I = \intervalclosed{0}{r}$.
    $I$ is closed in $\reals$ with respect to $\standardTopologyR$ by \cref{standard_closed_interval_closed}.
    For all $A$ such that $A\subseteq Y$ we have
        $\img{h}{\closure{A}{Y}{T}}\subseteq \closure{\img{h}{A}}{\reals}{\standardTopologyR}$
        by \cref{continuous_closure_image_subset}.
    Therefore $\preimg{h}{I}$ is closed in $Y$ with respect to $T$ by \cref{closure_image_subset_implies_preimg_closed}.
    Show that $\preimg{h}{I}\subseteq B$.
    \begin{subproof}
        Follows by \cref{preimg_subseteq_dom,subseteq,funs,preimg_iff,funs_is_function,function_apply_intro,id_iff,continuous,real_lt_subst_left,real_lt_subst_right,intervalclosed,real_leq_split}.
    \end{subproof}
    Show that $B\subseteq\preimg{h}{I}$.
    \begin{subproof}
        Follows by \cref{subseteq,times_tuple_intro,metric_on,funs_type_apply,metric_nonnegative,real_leq_split,intervalclosed,id_iff,continuous,funs_is_function,function_apply_intro,real_lt_subst_left,real_lt_subst_right,elem_subseteq,preimg_iff}.
    \end{subproof}
    Therefore $B$ is closed in $Y$ with respect to $T$
        by \cref{subseteq_antisymmetric,real_lt_subst_right}.
\end{proof}

% from §48 helper lemma 19: Cauchy stage sets are closed
\begin{lemma}\label{function_sequence_cauchy_stage_set_closed}
    Assume $T$ is a topology on $X$.
    Assume $d$ is a metric on $Y$.
    Assume $s$ is a sequence in $\funs{X}{Y}$.
    Assume $k\in\naturalsPlus$.
    Assume $N\in\naturalsPlus$.
    Suppose for all $n\in\naturalsPlus$ we have
        $s(n)$ is continuous from $X$ to $Y$ with respect to $T$ and $\metricTopology{d}{Y}$.
    Assume $A$ is the $k$-th Cauchy stage set of $s$ at $N$ from $X$ to $Y$ with respect to $d$.
    Then $A$ is closed in $X$ with respect to $T$.
\end{lemma}
\begin{proof}
    Let $T_Y = \metricTopology{d}{Y}$.
    $A = \{x\in X \mid \forall n,m\in\naturalsPlus. (N \leq n \land N \leq m) \implies
        d((\apply{\apply{s}{n}}{x},\apply{\apply{s}{m}}{x})) \leq 1 / k\}$ by \cref{function_sequence_cauchy_stage_set}.
    $A\subseteq X$ by \cref{subseteq}.
    Show that for all $x$ such that $x\in\closure{A}{X}{T}$ we have $x\in A$.
    \begin{subproof}
        Fix $x$ such that $x\in\closure{A}{X}{T}$.
        $x\in X$ by \cref{closure}.
        Show that for all $n\in\naturalsPlus$ we have
            for all $m\in\naturalsPlus$ if $N \leq n$ and $N \leq m$ then
                $d((\apply{\apply{s}{n}}{x},\apply{\apply{s}{m}}{x})) \leq 1 / k$.
        \begin{subproof}
            Fix $n\in\naturalsPlus$.
            Show that for all $m\in\naturalsPlus$ if $N \leq n$ and $N \leq m$ then
                $d((\apply{\apply{s}{n}}{x},\apply{\apply{s}{m}}{x})) \leq 1 / k$.
            \begin{subproof}
                Fix $m\in\naturalsPlus$.
                Assume $N \leq n$ and $N \leq m$.
                Let $h = \{(u, d((\apply{\apply{s}{n}}{u},\apply{\apply{s}{m}}{u}))) \mid u\in X\}$.
                $h$ is continuous from $X$ to $\reals$ with respect to $T$ and $\standardTopologyR$
                    by \cref{subseteq,continuous_maps_pointwise_distance}.
                $h$ is a function by \cref{continuous,funs_is_function}.
                $h\in\funs{X}{\reals}$ by \cref{continuous}.
                Let $I = \intervalclosed{0}{1 / k}$.
                $I$ is closed in $\reals$ with respect to $\standardTopologyR$
                    by \cref{naturalsplus_reciprocal_real,standard_closed_interval_closed}.
                Show that for all $z$ such that $z\in\img{h}{A}$ we have $z\in I$.
                \begin{subproof}
                    Fix $z$ such that $z\in\img{h}{A}$.
                    Take $y\in\dom{h}\inter A$ such that $z = h(y)$ by \cref{img_of_function_elim}.
                    $y\in X$ and $y\in A$ by \cref{funs,inter}.
                    $d((\apply{\apply{s}{n}}{y},\apply{\apply{s}{m}}{y})) \leq 1 / k$ by \cref{subseteq}.
                    $\apply{\apply{s}{n}}{y}\in Y$ and $\apply{\apply{s}{m}}{y}\in Y$ by \cref{sequence_in,funs_type_apply}.
                    $d((\apply{\apply{s}{n}}{y},\apply{\apply{s}{m}}{y}))\in\reals$
                        by \cref{sequence_in,funs_type_apply,times_tuple_intro,metric_on}.
                    $0 < d((\apply{\apply{s}{n}}{y},\apply{\apply{s}{m}}{y}))$ or
                        $d((\apply{\apply{s}{n}}{y},\apply{\apply{s}{m}}{y})) = 0$
                        by \cref{metric_on,metric_nonnegative,real_leq_split}.
                    $d((\apply{\apply{s}{n}}{y},\apply{\apply{s}{m}}{y})) < 1 / k$ or
                        $d((\apply{\apply{s}{n}}{y},\apply{\apply{s}{m}}{y})) = 1 / k$ by \cref{real_leq_split}.
                    $h(y) = z$ and $h(y) = d((\apply{\apply{s}{n}}{y},\apply{\apply{s}{m}}{y}))$
                        by \cref{funs_is_function,function_apply_intro,real_lt_subst_right}.
                    Hence $z\in I$ by \cref{intervalclosed,real_lt_subst_left,real_lt_subst_right}.
                \end{subproof}
                We have $\img{h}{A}\subseteq I$ by \cref{subseteq}.
                Therefore $\closure{\img{h}{A}}{\reals}{\standardTopologyR}\subseteq I$
                    by \cref{img_subseteq_ran,continuous,funs_ran,subseteq_transitive,closure_subseteq_closed}.
                $h(x)\in\img{h}{\closure{A}{X}{T}}$ by \cref{funs,subseteq,inter_intro,img_of_function_intro}.
                $h(x)\in\closure{\img{h}{A}}{\reals}{\standardTopologyR}$ by \cref{continuous_closure_image_subset,subseteq}.
                $h(x)\in I$ by \cref{elem_subseteq}.
                Hence $d((\apply{\apply{s}{n}}{x},\apply{\apply{s}{m}}{x})) \leq 1 / k$
                    by \cref{intervalclosed,funs_is_function,function_apply_intro,real_leq_split,real_lt_subst_left,real_lt_subst_right}.
            \end{subproof}
        \end{subproof}
        Hence $x\in A$ by \cref{subseteq}.
    \end{subproof}
    Therefore $A$ is closed in $X$ with respect to $T$
        by \cref{subseteq,subseteq_closure,subseteq_antisymmetric,closure_closed_in}.
\end{proof}

% from §48 helper lemma 20: stage bounds pass to the pointwise limit
\begin{lemma}\label{function_sequence_cauchy_stage_limit_bound}
    Assume $d$ is a metric on $Y$.
    Assume $s$ is a sequence in $\funs{X}{Y}$.
    Assume $f\in\funs{X}{Y}$.
    Assume $k\in\naturalsPlus$.
    Assume $N\in\naturalsPlus$.
    Assume $A$ is the $k$-th Cauchy stage set of $s$ at $N$ from $X$ to $Y$ with respect to $d$.
    Assume $x\in A$.
    Suppose $t$ is the coordinate sequence of $s$ at $x$ from $X$ to $Y$.
    Suppose $t$ converges to $\apply{f}{x}$ in $Y$ with respect to $\metricTopology{d}{Y}$.
    Then $d((\apply{\apply{s}{N}}{x},\apply{f}{x})) \leq 1 / k$.
\end{lemma}
\begin{proof}
    $A = \{u\in X \mid \forall n,m\in\naturalsPlus. (N \leq n \land N \leq m) \implies
        d((\apply{\apply{s}{n}}{u},\apply{\apply{s}{m}}{u})) \leq 1 / k\}$ by \cref{function_sequence_cauchy_stage_set}.
    Let $a_0 = \apply{\apply{s}{N}}{x}$.
    $a_0\in Y$ by \cref{subseteq,sequence_in,funs_type_apply}.
    Let $r_0 = 1 / k$.
    Let $B = \{y\in Y \mid d((y,a_0)) \leq r_0\}$.
    $B$ is closed in $Y$ with respect to $\metricTopology{d}{Y}$
        by \cref{naturalsplus_reciprocal_real,metric_closed_ball_closed}.
    $\metricTopology{d}{Y}$ is a topology on $Y$
        by \cref{metric_on,metric_basis_is_basis,basis_topology_is_topology,metric_topology}.
    $t$ is a sequence in $Y$ by \cref{function_sequence_coordinate_sequence}.
    $\apply{f}{x}\in Y$ by \cref{funs_type_apply}.
    For all $n\in\naturalsPlus$ if $N \leq n$ then $t(n)\in B$
        by \cref{subseteq,function_sequence_coordinate_sequence,sequence_in,funs_type_apply,times_tuple_intro,metric_on,metric_symmetric,real_lt_subst_left,real_lt_subst_right}.
    There exists $L\in\naturalsPlus$ such that for all $n\in\naturalsPlus$ if $L \leq n$ then $t(n)\in B$.
    Therefore $\apply{f}{x}\in B$ by \cref{eventually_in_closed_limit}.
    $d((\apply{f}{x},a_0)) \leq r_0$ by \cref{subseteq}.
    We have $d((a_0,\apply{f}{x})) \leq r_0$
        by \cref{funs_type_apply,times_tuple_intro,metric_on,metric_symmetric,real_lt_subst_left}.
    Hence $d((\apply{\apply{s}{N}}{x},\apply{f}{x})) \leq 1 / k$ by \cref{real_lt_subst_left,real_lt_subst_right}.
\end{proof}

% from §48 helper definition 3: union of interiors of Cauchy stage sets
\begin{definition}\label{function_sequence_cauchy_stage_union}
    $U$ is the $k$-th Cauchy stage union of $s$ from $X$ to $Y$ with respect to $T$ and $d$ iff
        for all $x$ we have
            $x\in U$ iff
                $x\in X$ and
                there exists $N\in\naturalsPlus$ such that
                    there exists $A$ such that
                        $A$ is the $k$-th Cauchy stage set of $s$ at $N$ from $X$ to $Y$ with respect to $d$ and
                        $x\in\interior{A}{X}{T}$.
\end{definition}

% from §48 helper lemma 21: a four-point metric upper bound
\begin{lemma}\label{metric_four_point_upper_bound}
    Assume $d$ is a metric on $Y$.
    Assume $x_0,x_1,y_0,y_1\in Y$.
    Then $d((x_0,y_0)) \leq d((x_0,x_1)) + d((x_1,y_1)) + d((y_1,y_0))$.
\end{lemma}
\begin{proof}
    $d((x_0,y_0)) \leq d((x_0,x_1)) + d((x_1,y_0))$
        by \cref{times_tuple_intro,metric_on,metric_triangle_inequality}.
    $d((x_1,y_0)) \leq d((x_1,y_1)) + d((y_1,y_0))$
        by \cref{times_tuple_intro,metric_on,metric_triangle_inequality}.
    $d((x_0,y_0))\in\reals$ and $d((x_0,x_1))\in\reals$ and $d((x_1,y_0))\in\reals$ and
        $d((x_1,y_1))\in\reals$ and $d((y_1,y_0))\in\reals$ by \cref{times_tuple_intro,metric_on,funs_type_apply}.
    $d((x_0,x_1)) + d((x_1,y_0))\in\reals$ and
        $d((x_1,y_1)) + d((y_1,y_0))\in\reals$ and
        $d((x_0,x_1)) + (d((x_1,y_1)) + d((y_1,y_0)))\in\reals$ by \cref{real_add_closed}.
    Therefore $d((x_1,y_0)) + d((x_0,x_1)) \leq (d((x_1,y_1)) + d((y_1,y_0))) + d((x_0,x_1))$
        by \cref{real_leq_add}.
    $d((x_1,y_0)) + d((x_0,x_1)) = d((x_0,x_1)) + d((x_1,y_0))$ by \cref{real_add_comm}.
    $(d((x_1,y_1)) + d((y_1,y_0))) + d((x_0,x_1))
        = d((x_0,x_1)) + (d((x_1,y_1)) + d((y_1,y_0)))$ by \cref{real_add_comm}.
    $d((x_0,x_1)) + d((x_1,y_0)) \leq d((x_0,x_1)) + (d((x_1,y_1)) + d((y_1,y_0)))$
        by \cref{real_lt_subst_left,real_lt_subst_right}.
    Hence $d((x_0,y_0)) \leq d((x_0,x_1)) + (d((x_1,y_1)) + d((y_1,y_0)))$ by \cref{real_leq_trans}.
    $d((x_0,x_1)) + (d((x_1,y_1)) + d((y_1,y_0)))
        = d((x_0,x_1)) + d((x_1,y_1)) + d((y_1,y_0))$ by \cref{real_add_assoc}.
    Therefore $d((x_0,y_0)) \leq d((x_0,x_1)) + d((x_1,y_1)) + d((y_1,y_0))$ by \cref{real_lt_subst_right}.
\end{proof}

% from §48 helper lemma 22: Cauchy stage unions are open
\begin{lemma}\label{function_sequence_cauchy_stage_union_open}
    Assume $T$ is a topology on $X$.
    Assume $U$ is the $k$-th Cauchy stage union of $s$ from $X$ to $Y$ with respect to $T$ and $d$.
    Then $U$ is open in $X$ with respect to $T$.
\end{lemma}
\begin{proof}
    Let $M = \{V\in T \mid \text{there exists $N\in\naturalsPlus$ such that
        there exists $A$ such that
            $A$ is the $k$-th Cauchy stage set of $s$ at $N$ from $X$ to $Y$ with respect to $d$ and
            $V = \interior{A}{X}{T}$}\}$.
    $\unions{M}\in T$ by \cref{subseteq,topology_on}.
    Show that for all $x$ such that $x\in U$ we have $x\in\unions{M}$.
    \begin{subproof}
        Fix $x$ such that $x\in U$.
        $x\in X$ and there exists $N\in\naturalsPlus$ such that
            there exists $A$ such that
                $A$ is the $k$-th Cauchy stage set of $s$ at $N$ from $X$ to $Y$ with respect to $d$ and
                $x\in\interior{A}{X}{T}$ by \cref{function_sequence_cauchy_stage_union}.
        Take $N\in\naturalsPlus$ such that
            there exists $A$ such that
                $A$ is the $k$-th Cauchy stage set of $s$ at $N$ from $X$ to $Y$ with respect to $d$ and
                $x\in\interior{A}{X}{T}$.
        Take $A$ such that
            $A$ is the $k$-th Cauchy stage set of $s$ at $N$ from $X$ to $Y$ with respect to $d$ and
            $x\in\interior{A}{X}{T}$.
        $\interior{A}{X}{T}$ is open in $X$ with respect to $T$ by \cref{interior_open_in}.
        $\interior{A}{X}{T}\in T$ by \cref{open_in}.
        We have $\interior{A}{X}{T}\in M$.
        Therefore $x\in\unions{M}$ by \cref{unions_iff}.
    \end{subproof}
    $U\subseteq\unions{M}$ by \cref{subseteq}.
    Show that for all $x$ such that $x\in\unions{M}$ we have $x\in U$.
    \begin{subproof}
        Fix $x$ such that $x\in\unions{M}$.
        Take $V\in M$ such that $x\in V$ by \cref{unions_iff}.
        Take $N\in\naturalsPlus$ such that
            there exists $A$ such that
                $A$ is the $k$-th Cauchy stage set of $s$ at $N$ from $X$ to $Y$ with respect to $d$ and
                $V = \interior{A}{X}{T}$ by \cref{subseteq}.
        Take $A$ such that
            $A$ is the $k$-th Cauchy stage set of $s$ at $N$ from $X$ to $Y$ with respect to $d$ and
            $V = \interior{A}{X}{T}$.
        $V\subseteq X$ by \cref{topology_on,subseteq,pow_iff}.
        $x\in X$ and $x\in\interior{A}{X}{T}$ by \cref{subseteq,real_lt_subst_left}.
        Therefore $x\in U$ by \cref{function_sequence_cauchy_stage_union}.
    \end{subproof}
    We have $\unions{M}\subseteq U$ by \cref{subseteq}.
    $U = \unions{M}$ by \cref{subseteq_antisymmetric}.
    Therefore $U$ is open in $X$ with respect to $T$ by \cref{open_in}.
\end{proof}

% from §48 helper definition 4: intersection of an open set with a Cauchy stage set
\begin{definition}\label{function_sequence_cauchy_stage_slice}
    $B$ is the $k$-th Cauchy stage slice of $s$ at $N$ over $V$ from $X$ to $Y$ with respect to $d$ iff
        there exists $A$ such that
            $A$ is the $k$-th Cauchy stage set of $s$ at $N$ from $X$ to $Y$ with respect to $d$ and
            $B = V\inter A$.
\end{definition}

% from §48 helper lemma 24: Cauchy stage sets at a fixed stage are unique
\begin{lemma}\label{function_sequence_cauchy_stage_set_unique}
    Assume $A$ is the $k$-th Cauchy stage set of $s$ at $N$ from $X$ to $Y$ with respect to $d$.
    Assume $B$ is the $k$-th Cauchy stage set of $s$ at $N$ from $X$ to $Y$ with respect to $d$.
    Then $A = B$.
\end{lemma}
\begin{proof}
    Follows by \cref{function_sequence_cauchy_stage_set,subseteq,subseteq_antisymmetric}.
\end{proof}

% from §48 helper lemma 23: Cauchy stage unions are dense
\begin{lemma}\label{function_sequence_cauchy_stage_union_dense}
    Assume $T$ is a topology on $X$.
    Assume $d$ is a metric on $Y$.
    Assume $s$ is a sequence in $\funs{X}{Y}$.
    Assume $f\in\funs{X}{Y}$.
    Assume $k\in\naturalsPlus$.
    Assume $U$ is the $k$-th Cauchy stage union of $s$ from $X$ to $Y$ with respect to $T$ and $d$.
    Assume $X$ is a Baire space with respect to $T$.
    Suppose for all $n\in\naturalsPlus$ we have
        $s(n)$ is continuous from $X$ to $Y$ with respect to $T$ and $\metricTopology{d}{Y}$.
    Suppose for all $x\in X$ there exists $t$ such that
        $t$ is the coordinate sequence of $s$ at $x$ from $X$ to $Y$ and
        $t$ converges to $\apply{f}{x}$ in $Y$ with respect to $\metricTopology{d}{Y}$.
    Then $U$ is dense in $X$ with respect to $T$.
\end{lemma}
\begin{proof}
    Show that for all $x\in X$ we have $x\in\closure{U}{X}{T}$.
    \begin{subproof}
        Fix $x\in X$.
        Show that for all $V$ such that
            $V$ is open in $X$ with respect to $T$ and
            $x\in V$
            we have $V$ intersects $U$.
        \begin{subproof}
            Fix $V$ such that
                $V$ is open in $X$ with respect to $T$ and
                $x\in V$.
            Suppose not.
            $V$ does not intersect $U$.
            Let $T_V = \subspaceTopology{T}{V}$.
            $V\subseteq X$ by \cref{open_in,topology_on,subseteq,pow_iff}.
            $T_V$ is a topology on $V$ by \cref{subspace_topology_is_topology,subspace_of}.
            $V$ is a Baire space with respect to $T_V$ by \cref{open_subspace_of_baire_space}.
            Let $R = \{w\in\naturalsPlus\times\pow{X} \mid
                \text{$\snd{w}$ is the $k$-th Cauchy stage set of $s$ at $\fst{w}$ from $X$ to $Y$ with respect to $d$}\}$.
            $R$ is a relation by \cref{times_elem_is_tuple,relation}.
            Show that for all $n\in\naturalsPlus$ there exists $A$ such that $n\mathrel{R} A$.
            \begin{subproof}
                Fix $n\in\naturalsPlus$.
                Let $A = \{y\in X \mid \forall p,q\in\naturalsPlus. (n \leq p \land n \leq q) \implies
                    d((\apply{\apply{s}{p}}{y},\apply{\apply{s}{q}}{y})) \leq 1 / k\}$.
                $A$ is the $k$-th Cauchy stage set of $s$ at $n$ from $X$ to $Y$ with respect to $d$
                    by \cref{function_sequence_cauchy_stage_set}.
                $(n,A)\in R$
                    by \cref{function_sequence_cauchy_stage_set,subseteq,powerset_intro,times_tuple_intro,fst_eq,snd_eq}.
                There exists $B$ such that $n\mathrel{R} B$.
            \end{subproof}
            Take $a$ such that $a$ is a function on $\naturalsPlus$ and
                for all $n\in\naturalsPlus$ we have $n\mathrel{R} a(n)$ by \cref{choice_function}.
            $a$ is a sequence in $\pow{X}$
                by \cref{choice_function,function_apply_elim,subseteq,times_tuple_elim,funs_intro,sequence_in}.
            Let $b = \{(n, V\inter a(n)) \mid n\in\naturalsPlus\}$.
            For all $n,B_1,B_2$ such that $(n,B_1)\in b$ and $(n,B_2)\in b$ we have $B_1 = B_2$
                by \cref{pair_eq_iff}.
            $b$ is right-unique by \cref{rightunique}.
            $b$ is a function.
            For all $n$ such that $n\in\naturalsPlus$ we have $n\in\dom{b}$
                by \cref{pair_eq_iff,dom_intro}.
            For all $n$ such that $n\in\dom{b}$ we have $n\in\naturalsPlus$
                by \cref{dom_iff,pair_eq_iff}.
            $\dom{b} = \naturalsPlus$ by \cref{subseteq,subseteq_antisymmetric}.
            $b$ is a sequence in $\pow{V}$
                by \cref{dom_iff,pair_eq_iff,function_apply_intro,inter_lower_left,real_lt_subst_left,powerset_intro,funs_intro,sequence_in}.
            Show that for all $n\in\naturalsPlus$ we have
                $b(n)$ is closed in $V$ with respect to $T_V$ and
                $b(n)$ has empty interior in $V$ with respect to $T_V$.
            \begin{subproof}
                Fix $n\in\naturalsPlus$.
                $a(n)$ is the $k$-th Cauchy stage set of $s$ at $n$ from $X$ to $Y$ with respect to $d$
                    by \cref{choice_function,function_apply_elim,fst_eq,snd_eq}.
                $a(n)$ is closed in $X$ with respect to $T$ by \cref{function_sequence_cauchy_stage_set_closed}.
                $b(n) = V\inter a(n)$ by \cref{pair_eq_iff,sequence_in,funs_elim,function_apply_intro}.
                $b(n)$ is closed in $V$ with respect to $T_V$
                    by \cref{inter_comm,real_lt_subst_left,closed_in_subspace_iff,subspace_of}.
                Suppose not.
                $\interior{b(n)}{V}{T_V}\neq\emptyset$
                    by \cref{sequence_in,funs_elim,powerset_elim,empty_interior}.
                Take $y$ such that $y\in\interior{b(n)}{V}{T_V}$ by \cref{emptyset,nonempty_inhabited}.
                Take $W\in T_V$ such that $y\in W$ and $W\subseteq b(n)$ by \cref{interior}.
                $W$ is open in $V$ with respect to $T_V$ by \cref{open_in}.
                $W$ is open in $X$ with respect to $T$ by \cref{open_in_subspace_open_in_space,subspace_of}.
                $V\inter a(n)\subseteq a(n)$ by \cref{inter_lower_right}.
                $b(n)\subseteq a(n)$ by \cref{real_lt_subst_left}.
                $W\subseteq a(n)$ by \cref{subseteq_transitive}.
                $W\in T$ by \cref{open_in}.
                $W\subseteq X$ by \cref{topology_on,subseteq,pow_iff}.
                $y\in X$ by \cref{subseteq}.
                $y\in\interior{a(n)}{X}{T}$ by \cref{interior}.
                $y\in V$ by \cref{interior}.
                $y\in U$ by \cref{function_sequence_cauchy_stage_union}.
                $y\in V\inter U$ by \cref{inter_intro}.
                $V$ intersects $U$ by \cref{emptyset,intersects,inter}.
                Contradiction.
                $b(n)$ has empty interior in $V$ with respect to $T_V$ by \cref{empty_interior}.
            \end{subproof}
            $\unions{\img{b}{\naturalsPlus}}$ has empty interior in $V$ with respect to $T_V$ by \cref{baire_space}.
            Show that for all $y$ such that $y\in V$ we have $y\in\unions{\img{b}{\naturalsPlus}}$.
            \begin{subproof}
                Fix $y$ such that $y\in V$.
                $y\in X$ by \cref{subseteq}.
                Take $t$ such that
                    $t$ is the coordinate sequence of $s$ at $y$ from $X$ to $Y$ and
                    $t$ converges to $\apply{f}{y}$ in $Y$ with respect to $\metricTopology{d}{Y}$ by \cref{subseteq}.
                $t$ is a sequence in $Y$ by \cref{function_sequence_coordinate_sequence}.
                $\apply{f}{y}\in Y$ by \cref{funs_type_apply}.
                $t$ is a Cauchy sequence in $Y$ with respect to $d$ by \cref{metric_convergent_sequence_cauchy}.
                $1 / k\in\reals$ and $0 < 1 / k$
                    by \cref{naturalsplus_reciprocal_real,positive_reciprocal_naturalsplus}.
                Take $N\in\naturalsPlus$ such that
                    for all $p,q\in\naturalsPlus$ if $N \leq p$ and $N \leq q$ then
                        $d((t(p),t(q))) < 1 / k$ by \cref{cauchy_sequence}.
                Let $A = \{z\in X \mid \forall p,q\in\naturalsPlus. (N \leq p \land N \leq q) \implies
                    d((\apply{\apply{s}{p}}{z},\apply{\apply{s}{q}}{z})) \leq 1 / k\}$.
                $A$ is the $k$-th Cauchy stage set of $s$ at $N$ from $X$ to $Y$ with respect to $d$
                    by \cref{function_sequence_cauchy_stage_set}.
                For all $p,q$ such that
                    $p\in\naturalsPlus$ and
                    $q\in\naturalsPlus$ and
                    $N \leq p$ and
                    $N \leq q$
                    we have $d((\apply{\apply{s}{p}}{y},\apply{\apply{s}{q}}{y})) \leq 1 / k$
                    by \cref{subseteq,function_sequence_coordinate_sequence,real_lt_subst_left,real_lt_subst_right,real_lt_imp_leq}.
                $y\in A$ by \cref{subseteq}.
                $a(N)$ is the $k$-th Cauchy stage set of $s$ at $N$ from $X$ to $Y$ with respect to $d$
                    by \cref{choice_function,function_apply_elim,fst_eq,snd_eq}.
                $A = a(N)$ by \cref{function_sequence_cauchy_stage_set_unique}.
                $b(N) = V\inter a(N)$ by \cref{pair_eq_iff,sequence_in,funs_elim,function_apply_intro}.
                $y\in V\inter a(N)$ by \cref{inter_intro,real_lt_subst_left}.
                $y\in b(N)$ by \cref{real_lt_subst_right}.
                $b(N)\in\img{b}{\naturalsPlus}$ by \cref{sequence_in,funs_elim,function_apply_elim,img_iff}.
                $y\in\unions{\img{b}{\naturalsPlus}}$ by \cref{unions_iff}.
            \end{subproof}
            $V\subseteq\unions{\img{b}{\naturalsPlus}}$ by \cref{subseteq}.
            $\img{b}{\naturalsPlus}\subseteq\pow{V}$ by \cref{img_subseteq_ran,sequence_in,funs_ran,subseteq_transitive}.
            $\unions{\img{b}{\naturalsPlus}}\subseteq V$ by \cref{unions_subseteq_of_powerset_is_subseteq}.
            $\unions{\img{b}{\naturalsPlus}} = V$ by \cref{subseteq_antisymmetric}.
            $V$ has empty interior in $V$ with respect to $T_V$ by \cref{real_lt_subst_left}.
            $x\in\interior{V}{V}{T_V}$ by \cref{topology_on,subseteq_refl,interior}.
            $\interior{V}{V}{T_V} = \emptyset$ by \cref{empty_interior}.
            Contradiction by \cref{emptyset}.
            $V$ intersects $U$.
        \end{subproof}
        $U\subseteq X$ by \cref{function_sequence_cauchy_stage_union,subseteq}.
        $x\in\closure{U}{X}{T}$ by \cref{closure_iff_intersects_open}.
    \end{subproof}
    $X\subseteq\closure{U}{X}{T}$ by \cref{subseteq}.
    $\closure{U}{X}{T}\subseteq X$ by \cref{closure,subseteq}.
    $\closure{U}{X}{T} = X$ by \cref{subseteq_antisymmetric}.
    $U$ is dense in $X$ with respect to $T$ by \cref{dense_in}.
\end{proof}

% from §48 Item 7: pointwise limits of continuous maps are continuous on a dense set
\begin{theorem}\label{pointwise_limit_continuous_on_dense_set}
    Assume $T$ is a topology on $X$.
    Assume $d$ is a metric on $Y$.
    Assume $s$ is a sequence in $\funs{X}{Y}$.
    Assume $f\in\funs{X}{Y}$.
    Assume $X$ is a Baire space with respect to $T$.
    Suppose for all $n\in\naturalsPlus$ we have
        $s(n)$ is continuous from $X$ to $Y$ with respect to $T$ and $\metricTopology{d}{Y}$.
    Suppose for all $x\in X$ there exists $t$ such that
        $t$ is the coordinate sequence of $s$ at $x$ from $X$ to $Y$ and
        $t$ converges to $\apply{f}{x}$ in $Y$ with respect to $\metricTopology{d}{Y}$.
    Then there exists $C$ such that
        $C\subseteq X$ and
        $C$ is dense in $X$ with respect to $T$ and
        for all $x\in C$ we have
            $f$ is continuous at $x$ from $X$ to $Y$ with respect to $T$ and $\metricTopology{d}{Y}$.
\end{theorem}
\begin{proof}
    Let $T_Y = \metricTopology{d}{Y}$.
    $T$ is a topology on $X$ by \cref{baire_space}.
    $T_Y$ is a topology on $Y$
        by \cref{metric_on,metric_basis_is_basis,basis_topology_is_topology,metric_topology}.
    Let $R = \{w\in\naturalsPlus\times\pow{X} \mid
        \text{$\snd{w}$ is the $\fst{w}$-th Cauchy stage union of $s$ from $X$ to $Y$ with respect to $T$ and $d$}\}$.
    $R$ is a relation by \cref{times_elem_is_tuple,relation}.
    Show that for all $k\in\naturalsPlus$ there exists $U$ such that $k\mathrel{R} U$.
        \begin{subproof}
            Fix $k\in\naturalsPlus$.
            Let $S = \{w\in\naturalsPlus\times\pow{X} \mid
                \text{$\snd{w}$ is the $k$-th Cauchy stage set of $s$ at $\fst{w}$ from $X$ to $Y$ with respect to $d$}\}$.
            $S$ is a relation by \cref{times_elem_is_tuple,relation}.
            Let $U = \{x\in X \mid \exists N\in\naturalsPlus. \exists A\in\pow{X}. (N,A)\in S \land x\in\interior{A}{X}{T}\}$.
            For all $x$ we have
                $x\in U$ iff
                    $x\in X$ and
                    there exists $N\in\naturalsPlus$ such that
                        there exists $A$ such that
                            $A$ is the $k$-th Cauchy stage set of $s$ at $N$ from $X$ to $Y$ with respect to $d$ and
                            $x\in\interior{A}{X}{T}$
                by \cref{subseteq,function_sequence_cauchy_stage_set,powerset_intro,times_tuple_intro,fst_eq,snd_eq}.
            $U$ is the $k$-th Cauchy stage union of $s$ from $X$ to $Y$ with respect to $T$ and $d$
                by \cref{function_sequence_cauchy_stage_union}.
        $(k,U)\in R$
            by \cref{function_sequence_cauchy_stage_union,subseteq,powerset_intro,times_tuple_intro,fst_eq,snd_eq}.
        There exists $V$ such that $k\mathrel{R} V$.
    \end{subproof}
    Take $u$ such that $u$ is a function on $\naturalsPlus$ and
        for all $k\in\naturalsPlus$ we have $k\mathrel{R} u(k)$ by \cref{choice_function}.
    $u$ is a sequence in $\pow{X}$
        by \cref{choice_function,function_apply_elim,subseteq,times_tuple_elim,funs_intro,sequence_in}.
    For all $k\in\naturalsPlus$ we have
        $u(k)$ is open in $X$ with respect to $T$ and
        $u(k)$ is dense in $X$ with respect to $T$
        by \cref{choice_function,function_apply_elim,fst_eq,snd_eq,function_sequence_cauchy_stage_union_open,function_sequence_cauchy_stage_union_dense}.
    Let $C = \inters{\img{u}{\naturalsPlus}}$.
    $C$ is dense in $X$ with respect to $T$ by \cref{baire_space_open_set_characterization,baire_space}.
    $\img{u}{\naturalsPlus}\subseteq\pow{X}$ by \cref{img_subseteq_ran,sequence_in,funs_ran,subseteq_transitive}.
    $C\subseteq X$ by \cref{unions_subseteq_of_powerset_is_subseteq,inters,elem_subseteq,subseteq}.
    Show that for all $x\in C$ we have
        $f$ is continuous at $x$ from $X$ to $Y$ with respect to $T$ and $T_Y$.
        \begin{subproof}
            Fix $x\in C$.
            Show that for all $V$ such that $V$ is open in $Y$ with respect to $T_Y$ and $\apply{f}{x}\in V$
                there exists $W$ such that $W$ is open in $X$ with respect to $T$ and $x\in W$ and $\img{f}{W}\subseteq V$.
        \begin{subproof}
            Fix $V$ such that $V$ is open in $Y$ with respect to $T_Y$ and $\apply{f}{x}\in V$.
            $x\in X$ by \cref{subseteq}.
            $\apply{f}{x}\in Y$ by \cref{funs_type_apply}.
            $V\subseteq Y$ by \cref{open_in,topology_on,subseteq,pow_iff}.
            Take $\epsilon\in\reals$ such that $0 < \epsilon$ and
                $\metricBall{d}{Y}{\apply{f}{x}}{\epsilon}\subseteq V$ by \cref{metric_open_iff}.
            Let $\delta = \epsilon / (1 + 1)$.
            $\delta\in\reals$ and $0 < \delta$ and $\delta + \delta = \epsilon$ by \cref{real_half_of_positive}.
            Let $\eta = \delta / (1 + 1)$.
            $\eta\in\reals$ and $0 < \eta$ and $\eta + \eta = \delta$ by \cref{real_half_of_positive}.
            Take $k\in\naturalsPlus$ such that $0 < 1 / k$ and $1 / k < \eta$ by \cref{real_small_reciprocal}.
            $1 / k\in\reals$ by \cref{naturalsplus_reciprocal_real}.
            $u(k)\in\img{u}{\naturalsPlus}$ by \cref{choice_function,function_apply_elim,img_iff}.
            $x\in u(k)$ by \cref{inters_destr}.
            $(k,u(k))\in R$ by \cref{choice_function}.
            $u(k)$ is the $k$-th Cauchy stage union of $s$ from $X$ to $Y$ with respect to $T$ and $d$
                by \cref{fst_eq,snd_eq}.
            $x\in X$ and
                there exists $N\in\naturalsPlus$ such that
                    there exists $A$ such that
                        $A$ is the $k$-th Cauchy stage set of $s$ at $N$ from $X$ to $Y$ with respect to $d$ and
                        $x\in\interior{A}{X}{T}$ by \cref{function_sequence_cauchy_stage_union}.
            Take $N\in\naturalsPlus$ such that
                there exists $A$ such that
                    $A$ is the $k$-th Cauchy stage set of $s$ at $N$ from $X$ to $Y$ with respect to $d$ and
                    $x\in\interior{A}{X}{T}$ by \cref{subseteq}.
            Take $A$ such that
                $A$ is the $k$-th Cauchy stage set of $s$ at $N$ from $X$ to $Y$ with respect to $d$ and
                $x\in\interior{A}{X}{T}$ by \cref{subseteq}.
            $s(N)$ is continuous at $x$ from $X$ to $Y$ with respect to $T$ and $T_Y$
                by \cref{subseteq,continuous_implies_continuous_at}.
            Let $B = \metricBall{d}{Y}{\apply{\apply{s}{N}}{x}}{\delta}$.
            $B\subseteq Y$ by \cref{metric_ball,subseteq}.
            $\apply{\apply{s}{N}}{x}\in Y$ by \cref{sequence_in,funs_type_apply}.
            $B\in\metricBasis{d}{Y}$ by \cref{metric_ball,subseteq,powerset_intro,metric_basis}.
            $B\in T_Y$ by \cref{metric_basis_is_basis,basis_elem_open_in_generated,basis_for_topology,metric_topology}.
            $B$ is open in $Y$ with respect to $T_Y$ by \cref{open_in,topology_on}.
            $d((\apply{\apply{s}{N}}{x},\apply{\apply{s}{N}}{x})) = 0$ by \cref{metric_zero_characterization,metric_on,sequence_in,funs_type_apply}.
            $\apply{\apply{s}{N}}{x}\in B$ by \cref{real_lt_subst_left,metric_ball}.
            $B$ is a neighborhood of $\apply{\apply{s}{N}}{x}$ in $Y$ with respect to $T_Y$ by \cref{neighborhood}.
            Take $W_0$ such that
                $W_0$ is a neighborhood of $x$ in $X$ with respect to $T$ and
                $\img{\apply{s}{N}}{W_0}\subseteq B$ by \cref{continuous_at}.
            $W_0$ is open in $X$ with respect to $T$ and $x\in W_0$ by \cref{neighborhood}.
            Let $W = W_0\inter\interior{A}{X}{T}$.
            $W$ is open in $X$ with respect to $T$ and $x\in W$
                by \cref{interior_open_in,open_in,topology_on,inter_intro}.
            Show that for all $z$ such that $z\in\img{f}{W}$ we have $z\in V$.
            \begin{subproof}
                Fix $z$ such that $z\in\img{f}{W}$.
                $f$ is right-unique and $f$ is a relation and $\dom{f} = X$
                    by \cref{funs,funs_is_function,function_rightunique,funs_is_relation}.
                Take $y\in\dom{f}\inter W$ such that $z = f(y)$ by \cref{img_of_function_elim}.
                $y\in X$ and $y\in W_0$ and $y\in\interior{A}{X}{T}$ by \cref{funs,inter}.
                $y\in A$ and $x\in A$ by \cref{interior_subseteq,subseteq}.
                $\apply{s}{N}$ is right-unique and $\apply{s}{N}$ is a relation and $\dom{\apply{s}{N}} = X$
                    by \cref{continuous,funs,funs_is_function,function_rightunique,funs_is_relation}.
                $y\in\dom{\apply{s}{N}}$ by \cref{real_lt_subst_right}.
                $\apply{\apply{s}{N}}{y}\in\img{\apply{s}{N}}{W_0}$ by \cref{inter_intro,img_of_function_intro}.
                $\apply{\apply{s}{N}}{y}\in B$ by \cref{subseteq}.
                $\apply{\apply{s}{N}}{y}\in Y$ by \cref{subseteq}.
                $d((\apply{\apply{s}{N}}{x},\apply{\apply{s}{N}}{y})) < \delta$ by \cref{metric_ball}.
                $d((\apply{\apply{s}{N}}{y},\apply{\apply{s}{N}}{x})) < \delta$
                    by \cref{metric_on,metric_symmetric,real_lt_subst_left}.
                Take $t_y$ such that
                    $t_y$ is the coordinate sequence of $s$ at $y$ from $X$ to $Y$ and
                    $t_y$ converges to $\apply{f}{y}$ in $Y$ with respect to $T_Y$ by \cref{subseteq}.
                Take $t_x$ such that
                    $t_x$ is the coordinate sequence of $s$ at $x$ from $X$ to $Y$ and
                    $t_x$ converges to $\apply{f}{x}$ in $Y$ with respect to $T_Y$ by \cref{subseteq}.
                $d((\apply{\apply{s}{N}}{y},\apply{f}{y})) \leq 1 / k$ by \cref{function_sequence_cauchy_stage_limit_bound}.
                $d((\apply{\apply{s}{N}}{x},\apply{f}{x})) \leq 1 / k$ by \cref{function_sequence_cauchy_stage_limit_bound}.
                $\apply{f}{y}\in Y$ and $\apply{f}{x}\in Y$ and
                    $\apply{\apply{s}{N}}{y}\in Y$ and $\apply{\apply{s}{N}}{x}\in Y$
                    by \cref{funs_type_apply,sequence_in}.
                $(\apply{\apply{s}{N}}{y},\apply{f}{y})\in Y\times Y$ and
                    $(\apply{f}{y},\apply{\apply{s}{N}}{y})\in Y\times Y$ and
                    $(\apply{\apply{s}{N}}{x},\apply{f}{x})\in Y\times Y$ and
                    $(\apply{\apply{s}{N}}{y},\apply{\apply{s}{N}}{x})\in Y\times Y$
                    by \cref{times_tuple_intro}.
                $d((\apply{f}{y},\apply{\apply{s}{N}}{y})) \leq 1 / k$
                    by \cref{metric_on,metric_symmetric,real_lt_subst_left}.
                Let $a = d((\apply{f}{y},\apply{\apply{s}{N}}{y}))$.
                Let $b = d((\apply{\apply{s}{N}}{y},\apply{\apply{s}{N}}{x}))$.
                Let $c = d((\apply{\apply{s}{N}}{x},\apply{f}{x}))$.
                $a\in\reals$ and $b\in\reals$ and $c\in\reals$ by \cref{metric_on,funs_type_apply}.
                $a \leq 1 / k$ and $b < \delta$ and $c \leq 1 / k$
                    by \cref{real_lt_subst_left,real_lt_subst_right}.
                $a < 1 / k$ or $a = 1 / k$ by \cref{real_leq_split}.
                $a < \eta$ by \cref{real_lt_trans,real_lt_subst_left}.
                $c < 1 / k$ or $c = 1 / k$ by \cref{real_leq_split}.
                $c < \eta$ by \cref{real_lt_trans,real_lt_subst_left}.
                $a + c\in\reals$ and $\eta + \eta\in\reals$ by \cref{real_add_closed}.
                $a + c < \eta + \eta$ by \cref{real_lt_add_two}.
                $a + c < \delta$ by \cref{real_lt_subst_right}.
                $b + (a + c)\in\reals$ and $\delta + \delta\in\reals$ by \cref{real_add_closed}.
                $b + (a + c) < \delta + \delta$ by \cref{real_lt_add_two}.
                $b + (a + c) = a + b + c$ by \cref{real_add_assoc,real_add_comm}.
                $a + b + c < \epsilon$ by \cref{real_lt_subst_left,real_lt_subst_right}.
                $d((\apply{f}{y},\apply{f}{x})) \leq a + b + c$
                    by \cref{metric_four_point_upper_bound,real_lt_subst_left,real_lt_subst_right}.
                $d((\apply{f}{y},\apply{f}{x})) < a + b + c$ or
                    $d((\apply{f}{y},\apply{f}{x})) = a + b + c$
                    by \cref{real_leq_split}.
                $d((\apply{f}{y},\apply{f}{x}))\in\reals$ by \cref{times_tuple_intro,metric_on,funs_type_apply}.
                $d((\apply{f}{y},\apply{f}{x})) < \epsilon$
                    by \cref{real_lt_trans,real_lt_subst_left}.
                $d((\apply{f}{x},\apply{f}{y})) < \epsilon$
                    by \cref{times_tuple_intro,metric_on,metric_symmetric,real_lt_subst_left}.
                $\apply{f}{y}\in\metricBall{d}{Y}{\apply{f}{x}}{\epsilon}$ by \cref{metric_ball}.
                $z\in\metricBall{d}{Y}{\apply{f}{x}}{\epsilon}$ by \cref{real_lt_subst_left}.
                $z\in V$ by \cref{subseteq}.
            \end{subproof}
            $\img{f}{W}\subseteq V$ by \cref{subseteq}.
            There exists $U_0$ such that
                $U_0$ is open in $X$ with respect to $T$ and
                $x\in U_0$ and
                $\img{f}{U_0}\subseteq V$.
        \end{subproof}
        $x\in X$ by \cref{subseteq}.
        $f$ is continuous at $x$ from $X$ to $Y$ with respect to $T$ and $T_Y$
            by \cref{continuous_at_open}.
    \end{subproof}
    There exists $D$ such that
        $D\subseteq X$ and
        $D$ is dense in $X$ with respect to $T$ and
        for all $x\in D$ we have
            $f$ is continuous at $x$ from $X$ to $Y$ with respect to $T$ and $T_Y$.
\end{proof}
\section{§49 A Nowhere-Differentiable Function}

% from §49 helper definition 1: differentiability at a point on a real subset
\begin{definition}\label{differentiable_at_point_on_real_subset}
    $f$ is differentiable at $x$ on $I$ iff
        $I\subseteq\reals$ and
        $f\in\funs{I}{\reals}$ and
        $x\in I$ and
        there exists $l\in\reals$ such that
            for all $\epsilon\in\reals$ such that $0 < \epsilon$
            there exists $\delta\in\reals$ such that
                $0 < \delta$ and
                for all $y\in I$ if
                    $y\neq x$ and
                    $\abs{(y - x)} < \delta$
                then $\abs{(((\apply{f}{y} - \apply{f}{x}) / (y - x)) - l)} < \epsilon$.
\end{definition}

% from §49 helper definition 2: nowhere differentiable on a real subset
\begin{definition}\label{nowhere_differentiable_on_real_subset}
    $f$ is nowhere differentiable on $I$ iff
        $I\subseteq\reals$ and
        $f\in\funs{I}{\reals}$ and
        for all $x\in I$ there exists no $l\in\reals$ such that
            for all $\epsilon\in\reals$ such that $0 < \epsilon$
            there exists $\delta\in\reals$ such that
                $0 < \delta$ and
                for all $y\in I$ if
                    $y\neq x$ and
                    $\abs{(y - x)} < \delta$
                then $\abs{(((\apply{f}{y} - \apply{f}{x}) / (y - x)) - l)} < \epsilon$.
\end{definition}

% from §49 helper lemma 1: distance realized by a closest point
\begin{lemma}\label{point_set_distance_realized_by_closest_point}
    Assume $d$ is a metric on $X$.
    Assume $A\subseteq X$.
    Assume $x\in X$.
    Assume $a\in A$.
    Assume $r$ is the distance from $x$ to $A$ in $X$ with respect to $d$.
    Suppose for all $b\in A$ we have $d((x,a)) \leq d((x,b))$.
    Then $r = d((x,a))$.
\end{lemma}
\begin{proof}
    Let $S = \pointDistanceSet{d}{A}{x}$.
    $a\in X$ and $(x,a)\in X\times X$ and $d((x,a))\in\reals$
        by \cref{subseteq,times_tuple_intro,metric_on,funs_type_apply}.
    $S\subseteq\reals$ by \cref{point_distance_set,subseteq,times_tuple_intro,metric_on,funs_type_apply}.
    $d((x,a))\in S$ by \cref{point_distance_set}.
    $r \leq d((x,a))$ and $d((x,a)) \leq r$
        by \cref{point_set_distance,real_infimum,real_lower_bound,point_distance_set}.
    $r = d((x,a))$ by \cref{point_set_distance,real_leq_antisym}.
\end{proof}

% from §49 helper definition 3: large fixed-scale secants on a real subset
\begin{definition}\label{large_fixed_scale_secants_on_real_subset}
    $f$ has secants above $n$ at scale $h$ on $I$ iff
        $I\subseteq\reals$ and
        $f\in\funs{I}{\reals}$ and
        $n\in\reals$ and
        $h\in\reals$ and
        $0 < h$ and
        for all $x\in I$ there exists $y\in I$ such that
            either $y = x + h$ and
                $n < \abs{((\apply{f}{y} - \apply{f}{x}) / (y - x))}$
            or $y = x - h$ and
                $n < \abs{((\apply{f}{y} - \apply{f}{x}) / (y - x))}$.
\end{definition}

% from §49 helper definition 4: equal unit mesh
\begin{definition}\label{unit_interval_equal_mesh_49}
    $A$ is the equal mesh of order $m$ on the unit interval iff
        $m\in\naturalsPlus$ and
        $A = \{x\in\reals \mid \text{$x = 0$ or there exists $p\in\initialSegment{m}$ such that $x = p / m$}\}$.
\end{definition}

% from §49 helper lemma 2: continuity respects equal domain topology
\begin{lemma}\label{continuous_eq_domain}
    Suppose $T = T_0$.
    Suppose $f$ is continuous from $X$ to $Y$ with respect to $T$ and $U$.
    Then $f$ is continuous from $X$ to $Y$ with respect to $T_0$ and $U$.
\end{lemma}
\begin{proof}
    Trivial.
\end{proof}

% from §49 helper lemma 3: compactness respects equal topologies
\begin{lemma}\label{compact_eq_topology}
    Suppose $T = T_0$.
    Suppose $X$ is compact with respect to $T$.
    Then $X$ is compact with respect to $T_0$.
\end{lemma}
\begin{proof}
    Trivial.
\end{proof}

% from §49 helper lemma 4: the restricted metric induces the subspace topology
\begin{lemma}\label{restricted_metric_induces_subspace_topology}
    Assume $d$ is a metric on $X$.
    Let $T = \metricTopology{d}{X}$.
    Assume $Y$ is a subspace of $X$ with respect to $T$.
    Let $e = \restrl{d}{Y\times Y}$.
    Let $T_Y = \subspaceTopology{T}{Y}$.
    Let $T_e = \metricTopology{e}{Y}$.
    Then $T_Y = T_e$.
\end{lemma}
\begin{proof}
    $T$ is a topology on $X$ and $Y\subseteq X$ by \cref{metric_basis_is_basis,basis_topology_is_topology,metric_topology,subspace_of}.
    Show that $e$ is a metric on $Y$.
    \begin{subproof}
        Show that $e\in\funs{Y\times Y}{\reals}$.
        \begin{subproof}
            $d\in\funs{X\times X}{\reals}$ and $d$ is a function and $d$ is a relation and $\dom{d} = X\times X$
                by \cref{metric_on,funs_elim,funs_is_relation}.
            $Y\times Y\subseteq X\times X$ by \cref{times_subseteq_left,times_subseteq_right,subseteq_transitive}.
            $e$ is a function by \cref{restrl_of_function_is_function}.
            $\dom{e} = (X\times X)\inter(Y\times Y)$ by \cref{dom_of_restrl_eq_inter,pair_eq_iff}.
            $\dom{e} = Y\times Y$ by \cref{subseteq,inter_intro,restrl_dom,subseteq_antisymmetric}.
            For all $p\in\dom{e}$ we have $e(p)\in\reals$
                by \cref{inter_elim_left,dom_of_restrl_eq_inter,funs_type_apply,restrl_of_function_apply}.
            $e\in\funs{Y\times Y}{\reals}$ by \cref{funs_intro}.
        \end{subproof}
        $d$ is a function and $d$ is a relation and $\dom{d} = X\times X$
            by \cref{metric_on,funs_elim,funs_is_relation}.
        $Y\times Y\subseteq X\times X$ by \cref{times_subseteq_left,times_subseteq_right,subseteq_transitive}.

        For all $x,y$ such that $x\in Y$ and $y\in Y$ we have $e((x,y)) = d((x,y))$
            by \cref{times_tuple_intro,subseteq,restrl_of_function_apply}.

        Show that $e$ is nonnegative on $Y$.
        \begin{subproof}
            For all $x,y$ such that $x\in Y$ and $y\in Y$ we have $e((x,y)) \geq 0$
                by \cref{subseteq,metric_on,metric_nonnegative}.
            $e$ is nonnegative on $Y$ by \cref{metric_nonnegative}.
        \end{subproof}

        Show that $e$ is zero-characterized on $Y$.
        \begin{subproof}
            For all $x,y$ such that $x\in Y$ and $y\in Y$ we have
                $e((x,y)) = 0$ iff $x = y$
                by \cref{subseteq,metric_on,metric_zero_characterization}.
            $e$ is zero-characterized on $Y$ by \cref{metric_zero_characterization}.
        \end{subproof}

        Show that $e$ is symmetric on $Y$.
        \begin{subproof}
            For all $x,y$ such that $x\in Y$ and $y\in Y$ we have $e((x,y)) = e((y,x))$
                by \cref{subseteq,metric_on,metric_symmetric}.
            $e$ is symmetric on $Y$ by \cref{metric_symmetric}.
        \end{subproof}

        Show that $e$ satisfies the triangle inequality on $Y$.
        \begin{subproof}
            For all $x,y,z$ such that $x\in Y$ and $y\in Y$ and $z\in Y$
                we have $e((x,y)) + e((y,z)) \geq e((x,z))$
                by \cref{subseteq,metric_on,metric_triangle_inequality}.
            $e$ satisfies the triangle inequality on $Y$ by \cref{metric_triangle_inequality}.
        \end{subproof}
        $e$ is a metric on $Y$ by \cref{metric_on}.
    \end{subproof}

    $T_Y$ is a topology on $Y$ and $T_e$ is a topology on $Y$
        by \cref{subspace_topology_is_topology,metric_basis_is_basis,basis_topology_is_topology,metric_topology,subspace_of}.
    $e\in\funs{Y\times Y}{\reals}$ and $\dom{e} = Y\times Y$ by \cref{metric_on,funs_elim}.
    $d\in\funs{X\times X}{\reals}$ and $d$ is a function and $d$ is a relation and $\dom{d} = X\times X$
        by \cref{metric_on,funs_elim,funs_is_relation}.
    $Y\times Y\subseteq\dom{d}$ by \cref{times_subseteq_left,times_subseteq_right,subseteq_transitive,pair_eq_iff}.

    Show that for all $U$ such that $U$ is open in $Y$ with respect to $T_Y$ we have $U$ is open in $Y$ with respect to $T_e$.
    \begin{subproof}
        Fix $U$ such that $U$ is open in $Y$ with respect to $T_Y$.
        $U\subseteq Y$ by \cref{open_in,topology_on,subseteq,pow_iff}.
        Take $V\in T$ such that $U = Y\inter V$ by \cref{subspace_topology,open_in}.
        $V$ is open in $X$ with respect to $T$ by \cref{open_in}.
        $V\subseteq X$ by \cref{open_in,topology_on,subseteq,pow_iff}.
        Show that for all $y\in U$ there exists $\delta\in\reals$ such that $0 < \delta$ and $\metricBall{e}{Y}{y}{\delta}\subseteq U$.
        \begin{subproof}
            Fix $y\in U$.
            $y\in Y$ and $y\in V$ by \cref{inter}.
            Take $\delta\in\reals$ such that $0 < \delta$ and $\metricBall{d}{X}{y}{\delta}\subseteq V$
                by \cref{metric_open_iff}.
            Show that for all $z$ such that $z\in\metricBall{e}{Y}{y}{\delta}$ we have $z\in U$.
            \begin{subproof}
                Fix $z$ such that $z\in\metricBall{e}{Y}{y}{\delta}$.
                $z\in Y$ and $e((y,z)) < \delta$ by \cref{metric_ball}.
                $e((y,z)) = d((y,z))$ by \cref{times_tuple_intro,subseteq,restrl_of_function_apply}.
                $d((y,z)) < \delta$ by \cref{real_lt_subst_left}.
                $z\in\metricBall{d}{X}{y}{\delta}$ by \cref{metric_ball,subseteq}.
                $z\in U$ by \cref{subseteq,inter_intro}.
            \end{subproof}
            $\metricBall{e}{Y}{y}{\delta}\subseteq U$ by \cref{subseteq}.
        \end{subproof}
        $U$ is open in $Y$ with respect to $T_e$ by \cref{metric_open_iff,open_in}.
    \end{subproof}

    Show that for all $U$ such that $U$ is open in $Y$ with respect to $T_e$ we have $U$ is open in $Y$ with respect to $T_Y$.
    \begin{subproof}
        Fix $U$ such that $U$ is open in $Y$ with respect to $T_e$.
        $U\subseteq Y$ by \cref{open_in,topology_on,subseteq,pow_iff}.
        Let $B = \metricBasis{e}{Y}$.
        $B$ is a basis on $Y$ by \cref{metric_basis_is_basis}.
        Show that for all $W\in B$ we have $W$ is open in $Y$ with respect to $T_Y$.
        \begin{subproof}
            Fix $W\in B$.
            Take $y\in Y$ such that there exists $\delta\in\reals$ such that $0 < \delta$ and $W = \metricBall{e}{Y}{y}{\delta}$
                by \cref{metric_basis}.
            Take $\delta\in\reals$ such that $0 < \delta$ and $W = \metricBall{e}{Y}{y}{\delta}$ by \cref{metric_basis}.
            $y\in X$ by \cref{subseteq}.
            $\metricBall{d}{X}{y}{\delta}\subseteq X$ by \cref{metric_ball,subseteq}.
            Show that for all $z$ such that $z\in W$ we have $z\in Y\inter\metricBall{d}{X}{y}{\delta}$.
            \begin{subproof}
                Fix $z$ such that $z\in W$.
                $z\in\metricBall{e}{Y}{y}{\delta}$ by \cref{pair_eq_iff}.
                $z\in Y$ and $e((y,z)) < \delta$ by \cref{metric_ball}.
                $e((y,z)) = d((y,z))$ by \cref{times_tuple_intro,subseteq,restrl_of_function_apply}.
                $d((y,z)) < \delta$ by \cref{real_lt_subst_left}.
                $z\in\metricBall{d}{X}{y}{\delta}$ by \cref{metric_ball,subseteq}.
                $z\in Y\inter\metricBall{d}{X}{y}{\delta}$ by \cref{inter_intro}.
            \end{subproof}
            $W\subseteq Y\inter\metricBall{d}{X}{y}{\delta}$ by \cref{subseteq}.
            Show that for all $z$ such that $z\in Y\inter\metricBall{d}{X}{y}{\delta}$ we have $z\in W$.
            \begin{subproof}
                Fix $z$ such that $z\in Y\inter\metricBall{d}{X}{y}{\delta}$.
                $z\in Y$ and $z\in\metricBall{d}{X}{y}{\delta}$ by \cref{inter}.
                $d((y,z)) < \delta$ by \cref{metric_ball}.
                $e((y,z)) = d((y,z))$ by \cref{times_tuple_intro,subseteq,restrl_of_function_apply}.
                $e((y,z)) < \delta$ by \cref{real_lt_subst_left}.
                $z\in\metricBall{e}{Y}{y}{\delta}$ by \cref{metric_ball}.
                $z\in W$ by \cref{pair_eq_iff}.
            \end{subproof}
            $Y\inter\metricBall{d}{X}{y}{\delta}\subseteq W$ by \cref{subseteq}.
            $W = Y\inter\metricBall{d}{X}{y}{\delta}$ by \cref{subseteq_antisymmetric}.
            $\metricBall{d}{X}{y}{\delta}\in\metricBasis{d}{X}$ by \cref{metric_ball,powerset_intro,metric_basis}.
            $\metricBall{d}{X}{y}{\delta}$ is open in $X$ with respect to $T$
                by \cref{metric_topology,metric_basis_is_basis,basis_elem_open_in_generated,open_in}.
            $W$ is open in $Y$ with respect to $T_Y$ by \cref{subspace_topology,open_in}.
        \end{subproof}

        Let $M = \{W\in B \mid W\subseteq U\}$.
        Show that for all $y$ such that $y\in U$ we have $y\in\unions{M}$.
        \begin{subproof}
            Fix $y$ such that $y\in U$.
            Take $\delta\in\reals$ such that $0 < \delta$ and $\metricBall{e}{Y}{y}{\delta}\subseteq U$
                by \cref{metric_open_iff}.
            Let $W = \metricBall{e}{Y}{y}{\delta}$.
            $y\in Y$ by \cref{subseteq}.
            $W\in B$ by \cref{metric_ball,subseteq,powerset_intro,metric_basis,pair_eq_iff}.
            $y\in\unions{M}$
                by \cref{metric_ball,real_lt_subst_left,metric_on,metric_zero_characterization,pair_eq_iff,unions_intro}.
        \end{subproof}
        $U\subseteq\unions{M}$ by \cref{subseteq}.
        Show that for all $y$ such that $y\in\unions{M}$ we have $y\in U$.
        \begin{subproof}
            Trivial.
        \end{subproof}
        $\unions{M}\subseteq U$ by \cref{subseteq}.
        $U = \unions{M}$ by \cref{subseteq_antisymmetric}.
        $M\subseteq T_Y$ by \cref{pair_eq_iff,subseteq,open_in}.
        $\unions{M}\in T_Y$ by \cref{topology_on,subseteq,pow_iff}.
        $U$ is open in $Y$ with respect to $T_Y$ by \cref{subseteq_antisymmetric,open_in}.
    \end{subproof}

    Show that $T_Y\subseteq T_e$.
    \begin{subproof}
        Follows by \cref{subseteq,open_in}.
    \end{subproof}
    Show that $T_e\subseteq T_Y$.
    \begin{subproof}
        Follows by \cref{subseteq,open_in}.
    \end{subproof}
    $T_Y = T_e$ by \cref{subseteq_antisymmetric}.
\end{proof}

% from §49 helper lemma 5: differentiability provides real-subset data
\begin{lemma}\label{differentiable_at_point_local_secant_bound_49}
    Assume $f$ is differentiable at $x$ on $I$.
    Then $I\subseteq\reals$ and $f\in\funs{I}{\reals}$ and $x\in I$.
\end{lemma}
\begin{proof}
    Follows by \cref{differentiable_at_point_on_real_subset}.
\end{proof}

% from §49 helper lemma 6: arbitrarily large secants at arbitrarily small fixed scales prevent differentiability
\begin{lemma}\label{arbitrarily_large_fixed_scale_secants_not_differentiable_49}
    Assume $I\subseteq\reals$.
    Assume $f\in\funs{I}{\reals}$.
    Assume $x\in I$.
    Suppose for all $n\in\naturalsPlus$ there exists $h\in\reals$ such that
        $h \leq 1 / n$ and
        $f$ has secants above $n$ at scale $h$ on $I$.
    Then $f$ is not differentiable at $x$ on $I$.
\end{lemma}
\begin{proof}
    Suppose not.
    $f$ is differentiable at $x$ on $I$.
    $I\subseteq\reals$ and $x\in I$ by \cref{differentiable_at_point_local_secant_bound_49}.
    Take $l\in\reals$ such that
        for all $\epsilon\in\reals$ such that $0 < \epsilon$
        there exists $\delta\in\reals$ such that
            $0 < \delta$ and
            for all $y\in I$ if
                $y\neq x$ and
                $\abs{(y - x)} < \delta$
            then $\abs{(((\apply{f}{y} - \apply{f}{x}) / (y - x)) - l)} < \epsilon$
        by \cref{differentiable_at_point_on_real_subset}.
    Take $\delta_0\in\reals$ such that
        $0 < \delta_0$ and
        for all $y\in I$ if
            $y\neq x$ and
            $\abs{(y - x)} < \delta_0$
        then $\abs{(((\apply{f}{y} - \apply{f}{x}) / (y - x)) - l)} < 1$
        by \cref{real_one_in_naturals,real_naturals_subset_reals,subseteq,real_zero_lt_one}.
    Take $m\in\naturalsPlus$ such that
        $0 < 1 / m$ and
        $1 / m < \delta_0$ by \cref{real_small_reciprocal}.
    Let $b = \abs{l} + 1$.
    $\abs{l}\in\reals$ by \cref{real_abs_in_reals}.
    $1\in\reals$ by \cref{real_one_in_naturals,real_naturals_subset_reals,subseteq}.
    $b\in\reals$ by \cref{real_add_closed}.
    $0 \leq \abs{l}$ by \cref{real_abs_nonnegative}.
    $0\in\reals$ by \cref{real_add_ident}.
    $0 + 1 \leq \abs{l} + 1$ by \cref{real_add_ident,real_leq_add}.
    $0 + 1 = 1$ by \cref{real_add_ident}.
    $1 \leq b$ by \cref{real_lt_subst_right}.
    $0 < b$ by \cref{real_zero_lt_one,real_leq_split,real_lt_trans,real_lt_subst_right}.
    Let $r = b + m$.
    $m\in\reals$ by \cref{real_naturals_subset_reals,subseteq}.
    $r\in\reals$ by \cref{real_naturals_subset_reals,subseteq,real_add_closed}.
    Take $n\in\naturalsPlus$ such that $r < n$ by \cref{real_archimedean_property}.
    $0 < m$ by \cref{naturalsplus_positive_real}.
    $n\in\reals$ by \cref{real_naturals_subset_reals,subseteq}.
    $0 + m < b + m$ by \cref{real_lt_add}.
    $0 + m = m$ by \cref{real_add_ident}.
    $b + m = r$ by \cref{real_lt_subst_right}.
    $m < r$ by \cref{real_lt_subst_left}.
    $m < n$ by \cref{real_lt_trans}.
    $m \leq n$ by \cref{real_lt_imp_leq}.
    $0 + b < m + b$ by \cref{real_lt_add}.
    $0 + b = b$ by \cref{real_add_ident}.
    $m + b = r$ by \cref{real_add_comm,real_lt_subst_right}.
    $b < r$ by \cref{real_lt_subst_left}.
    $b < n$ by \cref{real_lt_trans}.
    Take $h\in\reals$ such that
        $h \leq 1 / n$ and
        $f$ has secants above $n$ at scale $h$ on $I$ by \cref{subseteq}.
    $1 / n \leq 1 / m$ by \cref{naturalsplus_reciprocal_antitone}.
    $1 / n\in\reals$ and $1 / m\in\reals$ by \cref{naturalsplus_reciprocal_real}.
    $h \leq 1 / m$ by \cref{real_leq_trans}.
    $0 < h$ by \cref{large_fixed_scale_secants_on_real_subset}.
    $h < 1 / m$ or $h = 1 / m$ by \cref{real_leq_split}.
    $h < \delta_0$ by \cref{real_lt_trans,real_lt_subst_left}.
    Take $y\in I$ such that
        either $y = x + h$ and
            $n < \abs{((\apply{f}{y} - \apply{f}{x}) / (y - x))}$
        or $y = x - h$ and
            $n < \abs{((\apply{f}{y} - \apply{f}{x}) / (y - x))}$ by \cref{large_fixed_scale_secants_on_real_subset}.
    $y\in\reals$ and $x\in\reals$ by \cref{subseteq}.
    Show that $y\neq x$ and $\abs{(y - x)} < \delta_0$.
    \begin{subproof}
        \begin{byCase}
        \caseOf{$y = x + h$ and $n < \abs{((\apply{f}{y} - \apply{f}{x}) / (y - x))}$.}
            $y - x = h$ by \cref{real_sub,real_add_assoc,real_add_inverse,real_add_ident,real_add_comm}.
            Suppose not.
            $y = x$.
            $y - x = 0$ by \cref{real_sub,real_add_inverse,real_add_closed,real_add_ident}.
            $h = 0$ by \cref{real_lt_subst_left}.
            Contradiction by \cref{real_pos_nonzero}.
            $y\neq x$.
            $\abs{(y - x)} = \abs{h}$ by \cref{real_lt_subst_left}.
            $\abs{(y - x)} = h$ by \cref{real_abs_nonneg,real_lt_imp_leq}.
            $y\neq x$ and $\abs{(y - x)} < \delta_0$ by \cref{real_lt_subst_right}.
        \caseOf{$y = x - h$ and $n < \abs{((\apply{f}{y} - \apply{f}{x}) / (y - x))}$.}
            $y - x = \neg{h}$ by \cref{real_sub,real_add_assoc,real_add_inverse,real_add_ident,real_add_comm}.
            Suppose not.
            $y = x$.
            $y - x = 0$ by \cref{real_sub,real_add_inverse,real_add_closed,real_add_ident}.
            $\neg{h} = 0$ by \cref{real_lt_subst_left}.
            $h = 0$ by \cref{real_add_inverse,real_add_ident,real_add_cancel}.
            Contradiction by \cref{real_pos_nonzero}.
            $y\neq x$.
            $\abs{(y - x)} = \abs{\neg{h}}$ by \cref{real_lt_subst_left}.
            $\neg{h}\in\reals$ by \cref{real_add_inverse}.
            $\neg{h} < 0$ by \cref{real_lt_neg,real_neg_zero,real_lt_subst_right}.
            $\abs{\neg{h}} = \neg{(\neg{h})}$ by \cref{real_abs_neg}.
            $\neg{(\neg{h})} = h$ by \cref{real_neg_neg}.
            $\abs{(y - x)} = h$ by \cref{real_lt_subst_right}.
            $y\neq x$ and $\abs{(y - x)} < \delta_0$ by \cref{real_lt_subst_right}.
        \end{byCase}
    \end{subproof}
    $n < \abs{((\apply{f}{y} - \apply{f}{x}) / (y - x))}$.
    $\abs{(((\apply{f}{y} - \apply{f}{x}) / (y - x)) - l)} < 1$ by \cref{subseteq}.
    $\apply{f}{y}\in\reals$ and $\apply{f}{x}\in\reals$ by \cref{funs_type_apply}.
    $y - x\in\reals$ by \cref{real_sub,real_add_inverse,real_add_closed}.
    $y - x\neq 0$.
    $(\apply{f}{y} - \apply{f}{x}) / (y - x)\in\reals$
        by \cref{real_sub,real_add_inverse,real_add_closed,real_div,real_mul_inverse,real_mul_closed}.
    $\abs{(((\apply{f}{y} - \apply{f}{x}) / (y - x)))} < \abs{l} + 1$
        by \cref{real_abs_sub_lt_bound}.
    $\abs{(((\apply{f}{y} - \apply{f}{x}) / (y - x)))} < b$ by \cref{real_lt_subst_right}.
    $\abs{(((\apply{f}{y} - \apply{f}{x}) / (y - x)))}\in\reals$ by \cref{real_abs_in_reals}.
    $n < b$ by \cref{real_lt_trans}.
    Contradiction by \cref{real_lt_trans,real_lt_irreflexive}.
\end{proof}

% from §49 helper lemma 7: equal meshes cover the unit interval
\begin{lemma}\label{unit_interval_equal_mesh_cover_49}
    Assume $m\in\naturalsPlus$.
    Assume $x\in\intervalclosed{0}{1}$.
    Then there exists $p$ such that
        either $p = 0$ and
            $x \leq 1 / m$
        or $p\in\naturalsPlus$ and
            $p + 1 \leq m$ and
            $p / m \leq x$ and
            $x \leq (p + 1) / m$.
\end{lemma}
\begin{proof}
    Let $S = \{a\in\initialSegment{m} \mid x \leq a / m\}$.
    $S\subseteq\initialSegment{m}$ and $\initialSegment{m}\subseteq\naturalsPlus$
        by \cref{subseteq,initial_segment_naturalsplus}.
    $m\in\initialSegment{m}$ by \cref{initial_segment_naturalsplus}.
    $x\in\reals$ and $x \leq 1$ by \cref{intervalclosed}.
    $m\in\reals$ and $0 < m$ by \cref{real_naturals_subset_reals,subseteq,naturalsplus_positive_real}.
    Hence $m \neq 0$ by \cref{real_pos_nonzero}.
    $m / m = m \rmul \inv{m}$ by \cref{real_div}.
    $m / m = 1$ by \cref{real_mul_inverse,real_mul_comm,real_lt_subst_left}.
    $x \leq m / m$ by \cref{real_lt_subst_right}.
    Thus $m\in S$ by \cref{initial_segment_naturalsplus}.
    Therefore $S\neq\emptyset$ by \cref{nonempty_inhabited,inhabited_nonempty}.
    Take $k\in S$ such that for all $a\in S$ we have $k \leq a$ by \cref{subseteq_transitive,real_naturals_least_element}.
    We have $k\in\naturalsPlus$ and $k \leq m$ and $x \leq k / m$ by \cref{subseteq,initial_segment_naturalsplus}.
    $k = 1$ or there exists $p\in\naturalsPlus$ such that $k = p + 1$ by \cref{real_naturals_predecessor}.
    \begin{byCase}
    \caseOf{$k = 1$.}
        $k / m = 1 / m$ by \cref{real_one_in_naturals,real_naturals_subset_reals,subseteq,real_lt_subst_left}.
        Hence $x \leq 1 / m$ by \cref{real_lt_subst_right}.
        Let $p = 0$.
        $p = 0$ by \cref{real_add_ident}.
        It is wrong that $p\in\naturalsPlus$
            by \cref{real_add_ident,naturalsplus_positive_real,real_lt_subst_left,real_lt_irreflexive}.
        Show that either $p = 0$ and
            $x \leq 1 / m$
        or $p\in\naturalsPlus$ and
            $p + 1 \leq m$ and
            $p / m \leq x$ and
            $x \leq (p + 1) / m$.
        \begin{subproof}
            Trivial.
        \end{subproof}
    \caseOf{There exists $p\in\naturalsPlus$ such that $k = p + 1$.}
        Take $p\in\naturalsPlus$ such that $k = p + 1$ by \cref{real_naturals_predecessor}.
        $p + 1 \leq m$ by \cref{real_lt_subst_left}.
        $p \leq m$ by \cref{real_naturals_succ_bound_elim_local}.
        Hence $p\in\initialSegment{m}$ by \cref{initial_segment_naturalsplus}.
        Show that $p / m \leq x$.
        \begin{subproof}
            Suppose not.
            It is wrong that $p / m \leq x$.
            $p\in\reals$ by \cref{real_naturals_subset_reals,subseteq}.
            $p / m\in\reals$ by \cref{real_div_naturalsplus_real}.
            $p / m < x$ or $p / m = x$ or $x < p / m$ by \cref{real_lt_trichotomy}.
            \begin{byCase}
            \caseOf{$p / m < x$.}
                $p / m \leq x$ by \cref{real_lt_imp_leq}.
                Contradiction.
            \caseOf{$p / m = x$.}
                Contradiction.
            \caseOf{$x < p / m$.}
                $x \leq p / m$ by \cref{real_lt_imp_leq}.
                Hence $p\in S$ by \cref{initial_segment_naturalsplus}.
                Thus $k \leq p$.
                $1\in\reals$ and $0\in\reals$ by \cref{real_one_in_naturals,real_naturals_subset_reals,subseteq,real_add_ident}.
                $0 < 1$ by \cref{real_zero_lt_one}.
                $0 + p < 1 + p$ by \cref{real_lt_add}.
                $0 + p = p$ by \cref{real_add_ident}.
                $1 + p = p + 1$ by \cref{real_add_comm}.
                $p + 1 = k$ by \cref{real_lt_subst_right}.
                Hence $p < k$ by \cref{real_lt_subst_left}.
                $p \leq k$ by \cref{real_lt_imp_leq}.
                $k\in\reals$ by \cref{real_naturals_subset_reals,subseteq}.
                Therefore $k = p$ by \cref{real_leq_antisym}.
                Contradiction by \cref{real_lt_irreflexive,real_lt_subst_right}.
            \end{byCase}
        \end{subproof}
        $k / m = (p + 1) / m$ by \cref{real_lt_subst_left}.
        $x \leq (p + 1) / m$ by \cref{real_lt_subst_right}.
        It is wrong that $p = 0$
            by \cref{real_naturals_subset_reals,subseteq,naturalsplus_positive_real,real_pos_nonzero}.
        Show that either $p = 0$ and
            $x \leq 1 / m$
        or $p\in\naturalsPlus$ and
            $p + 1 \leq m$ and
            $p / m \leq x$ and
            $x \leq (p + 1) / m$.
        \begin{subproof}
            Trivial.
        \end{subproof}
    \end{byCase}
\end{proof}

% from §49 helper lemma 8: uniform-metric bounds below one control pointwise real distances
\begin{lemma}\label{uniform_metric_pointwise_lt_below_one_49}
    Assume $\rho$ is the uniform metric on functions from $I$ to $\reals$ with respect to $\standardMetricR$.
    Assume $f\in\funs{I}{\reals}$.
    Assume $g\in\funs{I}{\reals}$.
    Assume $x\in I$.
    Assume $\delta\in\reals$.
    Suppose $\rho((f,g)) < \delta$.
    Suppose $\delta \leq 1$.
    Then $\apply{\standardMetricR}{(\apply{f}{x},\apply{g}{x})} < \delta$.
\end{lemma}
\begin{proof}
    Let $e = \boundedMetric{\standardMetricR}$.
    $\standardMetricR$ is a metric on $\reals$ by \cref{standard_metric_is_metric}.
    We have $\apply{e}{(\apply{f}{x},\apply{g}{x})} \leq \rho((f,g))$
        by \cref{uniform_function_metric_coordinate_leq}.
    $\rho((f,g))\in\reals$ by \cref{uniform_metric_on_function_space,metric_on,funs_type_apply,times_tuple_intro}.
    $\apply{f}{x}\in\reals$ and $\apply{g}{x}\in\reals$ by \cref{funs_type_apply}.
    $(\apply{f}{x},\apply{g}{x})\in\reals\times\reals$ by \cref{times_tuple_intro}.
    $\apply{e}{(\apply{f}{x},\apply{g}{x})}\in\reals$
        by \cref{bounded_metric_is_metric,metric_on,funs_type_apply,times_tuple_intro}.
    Hence $\apply{e}{(\apply{f}{x},\apply{g}{x})} < \delta$ by \cref{real_leq_lt_trans}.
    Thus $\apply{e}{(\apply{f}{x},\apply{g}{x})} < 1$
        by \cref{real_one_in_naturals,real_naturals_subset_reals,subseteq,real_lt_leq_trans_local}.
    Hence $\apply{e}{(\apply{f}{x},\apply{g}{x})}
        = \apply{\standardMetricR}{(\apply{f}{x},\apply{g}{x})}$ by
        \cref{bounded_metric_apply_value_lt_one,bounded_metric_is_metric,standard_metric_is_metric}.
    Therefore $\apply{\standardMetricR}{(\apply{f}{x},\apply{g}{x})} < \delta$
        by \cref{real_lt_subst_left}.
\end{proof}

% from §49 helper lemma 9: the restricted standard metric is a metric on the unit interval
\begin{lemma}\label{unit_interval_restricted_standard_metric_49}
    Let $I = \intervalclosed{0}{1}$.
    Let $e = \restrl{\standardMetricR}{I\times I}$.
    Then $e$ is a metric on $I$.
\end{lemma}
\begin{proof}
    $0\in\reals$ and $1\in\reals$ and $0 < 1$
        by \cref{real_add_ident,real_one_in_naturals,real_naturals_subset_reals,subseteq,real_zero_lt_one}.
    $I\subseteq\reals$ by \cref{intervalclosed,subseteq,real_lt_imp_leq}.
    Hence $I$ is compact with respect to $\subspaceTopology{\standardTopologyR}{I}$
        by \cref{real_intervalclosed_compact_standard}.
    $\standardMetricR$ is a metric on $\reals$ by \cref{standard_metric_is_metric}.
    $\standardTopologyR$ is a topology on $\reals$
        by \cref{standard_topology_reals,standard_basis_is_basis,basis_topology_is_topology}.
    Thus $\reals$ is Hausdorff with respect to $\standardTopologyR$
        by \cref{metric_space_hausdorff,standard_metric_topology}.
    Therefore $I$ is closed in $\reals$ with respect to $\standardTopologyR$
        by \cref{compact_subspace_closed,metric_space_hausdorff}.
    Hence $I$ is closed in $\reals$ with respect to $\metricTopology{\standardMetricR}{\reals}$
        by \cref{standard_metric_topology}.
    Thus $I$ is complete with respect to $e$ by \cref{reals_complete_standard_metric,closed_subset_complete_restricted_metric}.
    Therefore $e$ is a metric on $I$ by \cref{complete_metric_space}.
\end{proof}

% from §49 helper lemma 10: the restricted standard metric induces the standard subspace topology on the unit interval
\begin{lemma}\label{unit_interval_restricted_standard_metric_topology_49}
    Let $I = \intervalclosed{0}{1}$.
    Let $T_I = \subspaceTopology{\standardTopologyR}{I}$.
    Let $e = \restrl{\standardMetricR}{I\times I}$.
    Let $T_e = \metricTopology{e}{I}$.
    Then $T_I = T_e$.
\end{lemma}
\begin{proof}
    Follows by \cref{standard_metric_is_metric,standard_topology_reals,standard_basis_is_basis,basis_topology_is_topology,intervalclosed,subseteq,subspace_of,standard_metric_topology,restricted_metric_induces_subspace_topology}.
\end{proof}

% from §49 helper lemma 11: the unit interval is compact for the restricted standard metric
\begin{lemma}\label{unit_interval_compact_restricted_standard_metric_49}
    Let $I = \intervalclosed{0}{1}$.
    Let $T_I = \subspaceTopology{\standardTopologyR}{I}$.
    Let $e = \restrl{\standardMetricR}{I\times I}$.
    Let $T_e = \metricTopology{e}{I}$.
    Then $I$ is compact with respect to $T_e$.
\end{lemma}
\begin{proof}
    Follows by \cref{real_add_ident,real_one_in_naturals,real_naturals_subset_reals,subseteq,real_zero_lt_one,real_intervalclosed_compact_standard,real_lt_imp_leq,unit_interval_restricted_standard_metric_topology_49,compact_eq_topology}.
\end{proof}

% from §49 helper lemma 12: equal meshes are nonempty subsets of the unit interval
\begin{lemma}\label{unit_interval_equal_mesh_basic_49}
    Assume $A$ is the equal mesh of order $m$ on the unit interval.
    Then $A\subseteq\intervalclosed{0}{1}$ and $A\neq\emptyset$.
\end{lemma}
\begin{proof}
    Follows by \cref{unit_interval_equal_mesh_49,subseteq,intervalclosed,real_add_ident,real_zero_lt_one,real_lt_imp_leq,real_lt_subst_left,initial_segment_naturalsplus,naturalsplus_positive_real,real_naturals_subset_reals,real_div_naturalsplus_nonnegative,real_div_leq_same_denom,real_pos_nonzero,real_div,real_mul_inverse,real_mul_comm,real_lt_subst_right,nonempty_inhabited,inhabited_nonempty}.
\end{proof}

% from §49 helper lemma 13: distance to an equal mesh is continuous on the unit interval
\begin{lemma}\label{unit_interval_equal_mesh_distance_continuous_49}
    Assume $A$ is the equal mesh of order $m$ on the unit interval.
    Let $I = \intervalclosed{0}{1}$.
    Let $T_I = \subspaceTopology{\standardTopologyR}{I}$.
    Then there exists $q$ such that
        $q\in\funs{I}{\reals}$ and
        for all $x\in I$ we have $q(x)$ is the distance from $x$ to $A$ in $\reals$ with respect to $\standardMetricR$ and
        $q$ is continuous from $I$ to $\reals$ with respect to $T_I$ and $\standardTopologyR$.
\end{lemma}
\begin{proof}
    We have $A\subseteq I$ and $A\neq\emptyset$ by \cref{unit_interval_equal_mesh_basic_49}.
    Thus $A\subseteq\reals$ by \cref{intervalclosed,subseteq,subseteq_transitive}.
    Take $f$ such that
        $f$ is a function from $\reals$ to $\reals$ and
        for all $x\in\reals$ we have
            $f(x)$ is the distance from $x$ to $A$ in $\reals$ with respect to $\standardMetricR$ and
            $f$ is continuous from $\reals$ to $\reals$ with respect to $\standardTopologyR$ and $\standardTopologyR$
        by \cref{standard_metric_is_metric,standard_metric_topology,point_set_distance_continuous}.
    $f$ is continuous from $\reals$ to $\reals$ with respect to $\standardTopologyR$ and $\standardTopologyR$
        by \cref{real_add_ident}.
    Let $q = \restrl{f}{I}$.
    Thus $I$ is a subspace of $\reals$ with respect to $\standardTopologyR$
        by \cref{intervalclosed,subseteq,standard_topology_reals,standard_basis_is_basis,basis_topology_is_topology,subspace_of}.
    Hence $q$ is continuous from $I$ to $\reals$ with respect to $T_I$ and $\standardTopologyR$
        by \cref{subspace_of,continuous_restrict_domain}.
    For all $x\in I$ we have $q(x)$ is the distance from $x$ to $A$ in $\reals$ with respect to $\standardMetricR$
        by \cref{intervalclosed,subseteq,restrl_of_function_apply,real_lt_subst_left,funs_elim,funs_is_relation}.
    Therefore there exists $g$ such that
        $g\in\funs{I}{\reals}$ and
        for all $x\in I$ we have $g(x)$ is the distance from $x$ to $A$ in $\reals$ with respect to $\standardMetricR$ and
        $g$ is continuous from $I$ to $\reals$ with respect to $T_I$ and $\standardTopologyR$.
\end{proof}

% from §49 helper lemma 14: large fixed-scale secants are monotone in the threshold
\begin{lemma}\label{large_fixed_scale_secants_threshold_monotone_49}
    Assume $f$ has secants above $r$ at scale $h$ on $I$.
    Assume $n\in\reals$.
    Suppose $n < r$.
    Then $f$ has secants above $n$ at scale $h$ on $I$.
\end{lemma}
\begin{proof}
    We have $I\subseteq\reals$ and
        $f\in\funs{I}{\reals}$ and
        $r\in\reals$ and
        $h\in\reals$ and
        $0 < h$ by \cref{large_fixed_scale_secants_on_real_subset}.
    Show that for all $x\in I$ there exists $y\in I$ such that
        either $y = x + h$ and
            $n < \abs{((\apply{f}{y} - \apply{f}{x}) / (y - x))}$
        or $y = x - h$ and
            $n < \abs{((\apply{f}{y} - \apply{f}{x}) / (y - x))}$.
    \begin{subproof}
        Fix $x\in I$.
        Take $y\in I$ such that
            either $y = x + h$ and
                $r < \abs{((\apply{f}{y} - \apply{f}{x}) / (y - x))}$
            or $y = x - h$ and
                $r < \abs{((\apply{f}{y} - \apply{f}{x}) / (y - x))}$ by \cref{large_fixed_scale_secants_on_real_subset}.
        \begin{byCase}
        \caseOf{$y = x + h$ and $r < \abs{((\apply{f}{y} - \apply{f}{x}) / (y - x))}$.}
            Thus there exists $z\in I$ such that
                either $z = x + h$ and
                    $n < \abs{((\apply{f}{z} - \apply{f}{x}) / (z - x))}$
                or $z = x - h$ and
                    $n < \abs{((\apply{f}{z} - \apply{f}{x}) / (z - x))}$
                by \cref{funs_type_apply,subseteq,real_sub,real_add_assoc,real_add_inverse,real_add_ident,real_add_comm,real_pos_nonzero,real_lt_subst_left,real_add_closed,real_div,real_mul_inverse,real_mul_closed,real_lt_trans,real_abs_in_reals}.
        \caseOf{$y = x - h$ and $r < \abs{((\apply{f}{y} - \apply{f}{x}) / (y - x))}$.}
            Thus there exists $z\in I$ such that
                either $z = x + h$ and
                    $n < \abs{((\apply{f}{z} - \apply{f}{x}) / (z - x))}$
                or $z = x - h$ and
                    $n < \abs{((\apply{f}{z} - \apply{f}{x}) / (z - x))}$
                by \cref{funs_type_apply,subseteq,real_sub,real_add_assoc,real_add_inverse,real_add_ident,real_add_comm,real_lt_subst_left,real_add_cancel,real_pos_nonzero,real_add_closed,real_div,real_mul_inverse,real_mul_closed,real_lt_trans,real_abs_in_reals}.
        \end{byCase}
    \end{subproof}
    Therefore $f$ has secants above $n$ at scale $h$ on $I$ by \cref{large_fixed_scale_secants_on_real_subset}.
\end{proof}

% from §49 helper lemma 15: lower bounds on absolute values survive small perturbations
\begin{lemma}\label{real_abs_perturbation_lower_bound_49}
    Assume $a\in\reals$.
    Assume $b\in\reals$.
    Assume $n\in\reals$.
    Assume $s\in\reals$.
    Assume $r\in\reals$.
    Suppose $\abs{(a - b)} < s$.
    Suppose $n + s < r$.
    Suppose $r < \abs{a}$.
    Then $n < \abs{b}$.
\end{lemma}
\begin{proof}
    Follows by \cref{real_abs_in_reals,real_add_closed,real_lt_trichotomy,real_lt_imp_leq,real_lt_subst_right,real_leq_add,real_abs_sub_lt_bound,real_lt_leq_trans_local,real_lt_trans,real_lt_irreflexive}.
\end{proof}

% from §49 helper lemma 16a: subtracting zero from a real number does nothing
\begin{lemma}\label{real_sub_zero_local_49}
    Suppose $x\in\reals$.
    Then $x - 0 = x$.
\end{lemma}
\begin{proof}
    Follows by \cref{real_sub,real_neg_zero,real_add_ident,real_add_comm}.
\end{proof}

% from §49 helper lemma 16b: the standard metric to zero equals the coordinate on the nonnegative reals
\begin{lemma}\label{standard_metric_zero_value_nonnegative_49}
    Suppose $x\in\reals$ and $0 \leq x$.
    Then $\apply{\standardMetricR}{(x,0)} = x$.
\end{lemma}
\begin{proof}
    Follows by \cref{real_add_ident,standard_metric_value,real_sub_zero_local_49,real_abs_nonneg,real_lt_subst_right}.
\end{proof}

% from §49 helper lemma 16: distance to a set is bounded by the distance to any chosen point
\begin{lemma}\label{point_set_distance_leq_point_distance_49}
    Assume $d$ is a metric on $X$.
    Assume $A\subseteq X$.
    Assume $x\in X$.
    Assume $a\in A$.
    Assume $r$ is the distance from $x$ to $A$ in $X$ with respect to $d$.
    Then $r \leq d((x,a))$.
\end{lemma}
\begin{proof}
    Follows by \cref{subseteq,times_tuple_intro,metric_on,funs_type_apply,point_distance_set,point_set_distance,real_infimum,real_lower_bound}.
\end{proof}

% from §49 helper lemma 17: equal-mesh distance is bounded by one mesh step
\begin{lemma}\label{unit_interval_equal_mesh_distance_bound_49}
    Assume $A$ is the equal mesh of order $m$ on the unit interval.
    Assume $q\in\funs{\intervalclosed{0}{1}}{\reals}$.
    Assume for all $x\in\intervalclosed{0}{1}$ we have
        $q(x)$ is the distance from $x$ to $A$ in $\reals$ with respect to $\standardMetricR$.
    Assume $x\in\intervalclosed{0}{1}$.
    Then $q(x) \leq 1 / m$.
\end{lemma}
\begin{proof}
    We have $m\in\naturalsPlus$ and
        $A = \{y\in\reals \mid \text{$y = 0$ or there exists $p\in\initialSegment{m}$ such that $y = p / m$}\}$
        by \cref{unit_interval_equal_mesh_49}.
    Hence $A\subseteq\reals$ by \cref{subseteq}.
    $x\in\reals$ by \cref{intervalclosed}.
    Thus $q(x)$ is the distance from $x$ to $A$ in $\reals$ with respect to $\standardMetricR$ by \cref{subseteq}.
    Take $p$ such that
        either $p = 0$ and
            $x \leq 1 / m$
        or $p\in\naturalsPlus$ and
            $p + 1 \leq m$ and
            $p / m \leq x$ and
            $x \leq (p + 1) / m$ by \cref{unit_interval_equal_mesh_cover_49}.
    \begin{byCase}
    \caseOf{$p = 0$ and $x \leq 1 / m$.}
        $0\in A$ by \cref{real_add_ident,subseteq}.
        Hence $q(x) \leq \apply{\standardMetricR}{(x,0)}$ by \cref{standard_metric_is_metric,point_set_distance_leq_point_distance_49}.
        $\apply{\standardMetricR}{(x,0)} = x$ by \cref{intervalclosed,standard_metric_zero_value_nonnegative_49}.
        Therefore $q(x) \leq 1 / m$
            by \cref{real_leq_trans,real_lt_subst_right,funs_type_apply,naturalsplus_reciprocal_real}.
    \caseOf{$p\in\naturalsPlus$ and $p + 1 \leq m$ and $p / m \leq x$ and $x \leq (p + 1) / m$.}
        $p \leq m$ by \cref{real_naturals_succ_bound_elim_local}.
        Hence $p\in\initialSegment{m}$ by \cref{initial_segment_naturalsplus}.
        $p\in\reals$ by \cref{real_naturals_subset_reals,subseteq}.
        Thus $p / m\in A$ by \cref{subseteq,real_div_naturalsplus_real}.
        $q(x) \leq \apply{\standardMetricR}{(x,p / m)}$ by \cref{standard_metric_is_metric,point_set_distance_leq_point_distance_49}.
        $\apply{\standardMetricR}{(x,p / m)} = \abs{(x - p / m)}$ by \cref{standard_metric_value}.
        $\neg{(p / m)}\in\reals$ by \cref{real_add_inverse}.
        $(p / m) + \neg{(p / m)} \leq x + \neg{(p / m)}$ by \cref{real_leq_add}.
        $(p / m) + \neg{(p / m)} = 0$ by \cref{real_add_inverse}.
        $x + \neg{(p / m)} = x - p / m$ by \cref{real_sub}.
        $0 \leq x - p / m$ by \cref{real_lt_subst_left,real_lt_subst_right}.
        $x - p / m\in\reals$ by \cref{real_sub,real_add_inverse,real_add_closed}.
        $\abs{(x - p / m)} = x - p / m$ by \cref{real_abs_nonneg}.
        $m\in\reals$ and $1\in\reals$ and $1 / m\in\reals$ and $0\in\reals$
            by \cref{real_naturals_subset_reals,subseteq,real_one_in_naturals,naturalsplus_reciprocal_real,real_add_ident}.
        $p + 1\in\reals$ and $(p + 1) / m\in\reals$
            by \cref{real_add_closed,real_div_naturalsplus_real,real_naturals_closed_succ}.
        $(p + 1) / m = (p / m) + (1 / m)$ by \cref{real_div_add_same_denom}.
        $x + \neg{(p / m)} \leq ((p / m) + (1 / m)) + \neg{(p / m)}$ by \cref{real_leq_add}.
        $x + \neg{(p / m)} = x - (p / m)$ by \cref{real_sub}.
        Hence $x - (p / m) \leq ((p / m) + (1 / m)) + \neg{(p / m)}$ by \cref{real_lt_subst_left}.
        $(p / m) + (1 / m) = (1 / m) + (p / m)$ by \cref{real_add_comm}.
        $((p / m) + (1 / m)) + \neg{(p / m)} = ((1 / m) + (p / m)) + \neg{(p / m)}$
            by \cref{real_lt_subst_left}.
        $((1 / m) + (p / m)) + \neg{(p / m)} = (1 / m) + ((p / m) + \neg{(p / m)})$
            by \cref{real_add_assoc}.
        $(p / m) + \neg{(p / m)} = 0$ by \cref{real_add_inverse}.
        $((1 / m) + (p / m)) + \neg{(p / m)} = (1 / m) + 0$ by \cref{real_lt_subst_right}.
        $(1 / m) + 0 = 0 + (1 / m)$ by \cref{real_add_comm}.
        $0 + (1 / m) = 1 / m$ by \cref{real_add_ident}.
        $x - (p / m) \leq 1 / m$ by \cref{real_lt_subst_right}.
        $q(x) \leq x - (p / m)$ by \cref{real_lt_subst_right}.
        Therefore $q(x) \leq 1 / m$ by \cref{real_leq_trans,funs_type_apply}.
    \end{byCase}
\end{proof}

% from §49 helper lemma 18: the distance from a point to any set containing it is zero
\begin{lemma}\label{point_set_distance_zero_when_member_49}
    Assume $d$ is a metric on $X$.
    Assume $A\subseteq X$.
    Assume $x\in A$.
    Assume $r$ is the distance from $x$ to $A$ in $X$ with respect to $d$.
    Then $r = 0$.
\end{lemma}
\begin{proof}
    Follows by \cref{subseteq,times_tuple_intro,metric_on,metric_zero_characterization,metric_nonnegative,real_lt_subst_left,point_set_distance_realized_by_closest_point,real_lt_subst_right}.
\end{proof}

% from §49 helper definition 5: large-secant continuous functions on the unit interval
\begin{definition}\label{unit_interval_large_secant_class_49}
    $U$ is the large-secant class above $n$ on $I$ with respect to $T_I$ iff
        $U = \{f\in\funs{I}{\reals} \mid \text{$f$ is continuous from $I$ to $\reals$ with respect to $T_I$ and $\standardTopologyR$ and
            there exists $h\in\reals$ such that
                $h \leq 1 / n$ and
                $f$ has secants above $n$ at scale $h$ on $I$}\}$.
\end{definition}

% from §49 helper definition 6: robust large-secant continuous functions on the unit interval
\begin{definition}\label{unit_interval_robust_large_secant_class_49}
    $U$ is the robust secant class above $n$ on $I$ with respect to $T_I$ and $\rho$ iff
        $U = \{g\in\funs{I}{\reals} \mid \text{$g$ is continuous from $I$ to $\reals$ with respect to $T_I$ and $\standardTopologyR$ and
            there exist $f,h$ such that
                $f$ is continuous from $I$ to $\reals$ with respect to $T_I$ and $\standardTopologyR$ and
                $h\in\reals$ and
                $h \leq 1 / n$ and
                $f$ has secants above $n + n$ at scale $h$ on $I$ and
                $\rho((g,f)) < (((h \rmul n) / (1 + 1)) / (1 + 1)) / (1 + 1)$}\}$.
\end{definition}

% from §49 helper lemma 18a: ordered left endpoints are closer in the standard metric
\begin{lemma}\label{standard_metric_ordered_left_endpoint_leq_49}
    Assume $c\in\reals$.
    Assume $a\in\reals$.
    Assume $x\in\reals$.
    Suppose $c \leq a$ and $a \leq x$.
    Then $\apply{\standardMetricR}{(x,a)} \leq \apply{\standardMetricR}{(x,c)}$.
\end{lemma}
\begin{proof}
    Follows by \cref{real_add_inverse,real_sub,real_add_closed,real_leq_trans,real_leq_add,real_add_ident,real_lt_subst_left,real_lt_subst_right,real_abs_nonneg,real_leq_neg,real_add_comm,standard_metric_value}.
\end{proof}

% from §49 helper lemma 18b: ordered right endpoints are closer in the standard metric
\begin{lemma}\label{standard_metric_ordered_right_endpoint_leq_49}
    Assume $x\in\reals$.
    Assume $b\in\reals$.
    Assume $c\in\reals$.
    Suppose $x \leq b$ and $b \leq c$.
    Then $\apply{\standardMetricR}{(x,b)} \leq \apply{\standardMetricR}{(x,c)}$.
\end{lemma}
\begin{proof}
    Follows by \cref{real_add_inverse,real_sub,real_add_closed,real_leq_trans,real_leq_add,real_add_ident,real_lt_subst_left,real_lt_subst_right,real_abs_nonneg,standard_metric_value,real_abs_sub_swap}.
\end{proof}

% from §49 helper lemma 19: a separated left endpoint realizes the distance
\begin{lemma}\label{point_set_distance_left_endpoint_realized_49}
    Assume $A\subseteq\reals$.
    Assume $a\in A$.
    Assume $b\in A$.
    Assume $x\in\reals$.
    Assume $a \leq x$ and $x \leq b$.
    Assume $\apply{\standardMetricR}{(x,a)} \leq \apply{\standardMetricR}{(x,b)}$.
    Assume $r$ is the distance from $x$ to $A$ in $\reals$ with respect to $\standardMetricR$.
    Suppose for all $c\in A$ if $c\neq a$ and $c\neq b$ then $c \leq a$ or $b \leq c$.
    Then $r = \apply{\standardMetricR}{(x,a)}$.
\end{lemma}
\begin{proof}
    Follows by \cref{standard_metric_is_metric,subseteq,real_lt_subst_right,standard_metric_ordered_left_endpoint_leq_49,standard_metric_ordered_right_endpoint_leq_49,metric_on,funs_type_apply,times_tuple_intro,real_leq_trans,point_set_distance_realized_by_closest_point}.
\end{proof}

% from §49 helper lemma 20: a separated right endpoint realizes the distance
\begin{lemma}\label{point_set_distance_right_endpoint_realized_49}
    Assume $A\subseteq\reals$.
    Assume $a\in A$.
    Assume $b\in A$.
    Assume $x\in\reals$.
    Assume $a \leq x$ and $x \leq b$.
    Assume $\apply{\standardMetricR}{(x,b)} \leq \apply{\standardMetricR}{(x,a)}$.
    Assume $r$ is the distance from $x$ to $A$ in $\reals$ with respect to $\standardMetricR$.
    Suppose for all $c\in A$ if $c\neq a$ and $c\neq b$ then $c \leq a$ or $b \leq c$.
    Then $r = \apply{\standardMetricR}{(x,b)}$.
\end{lemma}
\begin{proof}
    Follows by \cref{standard_metric_is_metric,subseteq,real_lt_subst_right,standard_metric_ordered_left_endpoint_leq_49,standard_metric_ordered_right_endpoint_leq_49,metric_on,funs_type_apply,times_tuple_intro,real_leq_trans,point_set_distance_realized_by_closest_point}.
\end{proof}

% from §49 helper lemma 21: nonzero equal-mesh points lie beyond the first step
\begin{lemma}\label{unit_interval_equal_mesh_first_step_separation_49}
    Assume $A$ is the equal mesh of order $m$ on the unit interval.
    Assume $c\in A$.
    Suppose $c \neq 0$.
    Then $1 / m \leq c$.
\end{lemma}
\begin{proof}
    Follows by \cref{unit_interval_equal_mesh_49,subseteq,initial_segment_naturalsplus,real_one_leq_naturalsplus,real_one_in_naturals,real_naturals_subset_reals,naturalsplus_positive_real,real_div_leq_same_denom,real_lt_subst_right}.
\end{proof}

% from §49 helper lemma 22: adjacent positive mesh points separate all other mesh points
\begin{lemma}\label{unit_interval_equal_mesh_adjacent_separation_49}
    Assume $A$ is the equal mesh of order $m$ on the unit interval.
    Assume $p\in\naturalsPlus$.
    Assume $p + 1 \leq m$.
    Assume $c\in A$.
    Suppose $c \neq p / m$ and $c \neq (p + 1) / m$.
    Then $c \leq p / m$ or $(p + 1) / m \leq c$.
\end{lemma}
\begin{proof}
    We have $m\in\naturalsPlus$ and
        $A = \{x\in\reals \mid \text{$x = 0$ or there exists $q\in\initialSegment{m}$ such that $x = q / m$}\}$
        by \cref{unit_interval_equal_mesh_49}.
    \begin{byCase}
    \caseOf{$c = 0$.}
        $c \leq p / m$
            by \cref{real_div_naturalsplus_nonnegative,real_lt_imp_leq,real_add_ident,real_naturals_subset_reals,subseteq,naturalsplus_positive_real,real_lt_subst_left}.
    \caseOf{$c \neq 0$.}
        Take $q\in\initialSegment{m}$ such that $c = q / m$ by \cref{subseteq}.
        $q\in\naturalsPlus$ and $q \leq m$ by \cref{initial_segment_naturalsplus}.
        $q\in\reals$ and $p\in\reals$ and $m\in\reals$
            by \cref{real_naturals_subset_reals,subseteq}.
        $q < p$ or $p < q$ by \cref{real_lt_trichotomy,real_div_naturalsplus_real,real_lt_subst_left,real_lt_subst_right}.
        \begin{byCase}
        \caseOf{$q < p$.}
            $q / m \leq p / m$
                by \cref{real_lt_imp_leq,real_div_leq_same_denom,real_naturals_subset_reals,subseteq,naturalsplus_positive_real}.
            Hence $c \leq p / m$ by \cref{real_lt_subst_left}.
        \caseOf{$p < q$.}
            $p + 1 \leq q$ by \cref{space_filling_naturalsplus_lt_imp_succ_leq}.
            $(p + 1) / m \leq q / m$
                by \cref{real_div_leq_same_denom,unit_interval_equal_mesh_49,real_one_in_naturals,real_naturals_subset_reals,subseteq,real_add_closed,naturalsplus_positive_real}.
            Thus $(p + 1) / m \leq c$ by \cref{real_lt_subst_right}.
        \end{byCase}
    \end{byCase}
\end{proof}

% from §49 helper lemma 23: functions in every large-secant class are nowhere differentiable
\begin{lemma}\label{unit_interval_large_secant_classes_nowhere_differentiable_49}
    Assume $C$ is the set of continuous functions from $I$ to $\reals$ with respect to $T_I$ and $\standardTopologyR$.
    Assume $f\in C$.
    Suppose for all $n\in\naturalsPlus$ there exists $U$ such that
        $U$ is the large-secant class above $n$ on $I$ with respect to $T_I$ and
        $f\in U$.
    Then $f$ is nowhere differentiable on $I$.
\end{lemma}
\begin{proof}
    Follows by \cref{continuous_function_set,continuous,funs,real_one_in_naturals,subseteq,unit_interval_large_secant_class_49,large_fixed_scale_secants_on_real_subset,arbitrarily_large_fixed_scale_secants_not_differentiable_49,differentiable_at_point_on_real_subset,nowhere_differentiable_on_real_subset}.
\end{proof}

% from §49 helper lemma 24: robust large-secant classes are open in the restricted uniform metric
\begin{lemma}\label{unit_interval_robust_large_secant_class_open_49}
    Assume $C$ is the set of continuous functions from $I$ to $\reals$ with respect to $T_I$ and $\standardTopologyR$.
    Assume $\rho$ is the uniform metric on functions from $I$ to $\reals$ with respect to $\standardMetricR$.
    Let $e = \restrl{\rho}{C\times C}$.
    Assume $e$ is a metric on $C$.
    Assume $n\in\naturalsPlus$.
    Assume $V$ is the robust secant class above $n$ on $I$ with respect to $T_I$ and $\rho$.
    Then $V$ is open in $C$ with respect to $\metricTopology{e}{C}$.
\end{lemma}
\begin{proof}
    $C\subseteq\funs{I}{\reals}$ and $\rho$ is a metric on $\funs{I}{\reals}$
        by \cref{continuous_function_set,subseteq,uniform_metric_on_function_space}.
    Show that for all $g\in V$ there exists $\delta\in\reals$ such that $0 < \delta$ and $\metricBall{e}{C}{g}{\delta}\subseteq V$.
    \begin{subproof}
        Fix $g\in V$.
        Take $f,h$ such that
            $f$ is continuous from $I$ to $\reals$ with respect to $T_I$ and $\standardTopologyR$ and
            $h\in\reals$ and
            $h \leq 1 / n$ and
            $f$ has secants above $n + n$ at scale $h$ on $I$ and
            $\rho((g,f)) < (((h \rmul n) / (1 + 1)) / (1 + 1)) / (1 + 1)$
            by \cref{unit_interval_robust_large_secant_class_49,subseteq}.
        $f\in\funs{I}{\reals}$ and $f\in C$
            by \cref{continuous,funs,continuous_function_set}.
        $g\in\funs{I}{\reals}$ and $g\in C$
            by \cref{unit_interval_robust_large_secant_class_49,continuous_function_set}.
        $0 < h$ by \cref{large_fixed_scale_secants_on_real_subset}.
        $n\in\reals$ and $0 < n$ by \cref{real_naturals_subset_reals,subseteq,naturalsplus_positive_real}.
        $h \rmul n\in\reals$ and $0 < h \rmul n$ by \cref{real_mul_closed,real_mul_pos_pos_gt}.
        Let $a = (h \rmul n) / (1 + 1)$.
        $a\in\reals$ and $0 < a$ and $a + a = h \rmul n$ by \cref{real_half_of_positive}.
        Let $b = a / (1 + 1)$.
        $b\in\reals$ and $0 < b$ and $b + b = a$ by \cref{real_half_of_positive}.
        Let $r = b / (1 + 1)$.
        $r\in\reals$ and $0 < r$ and $r + r = b$ by \cref{real_half_of_positive}.
        $b = ((h \rmul n) / (1 + 1)) / (1 + 1)$ by \cref{real_lt_subst_left}.
        $r = (((h \rmul n) / (1 + 1)) / (1 + 1)) / (1 + 1)$ by \cref{real_lt_subst_left}.
        $\rho((g,f)) < r$ by \cref{real_lt_subst_right}.
        $\rho((g,f))\in\reals$ by \cref{metric_on,funs_type_apply,times_tuple_intro}.
        Let $\delta = r - \rho((g,f))$.
        $\delta\in\reals$ by \cref{real_sub,real_add_inverse,real_add_closed}.
        $0 < \delta$ by \cref{real_lt_imp_pos_diff}.
        Show that for all $u$ such that $u\in\metricBall{e}{C}{g}{\delta}$ we have $u\in V$.
        \begin{subproof}
            Fix $u$ such that $u\in\metricBall{e}{C}{g}{\delta}$.
            $u\in C$ and $e((u,g)) < \delta$ by \cref{metric_ball,metric_on,metric_symmetric,times_tuple_intro,real_lt_subst_left}.
            $e((u,g)) = \rho((u,g))$ and $e((g,f)) = \rho((g,f))$ and $e((u,f)) = \rho((u,f))$
                by \cref{restrl_of_function_apply,subseteq,times_subseteq_left,times_subseteq_right,subseteq_transitive,pair_eq_iff,metric_on,funs_elim,funs_is_function,times_tuple_intro}.
            $e((u,f)) \leq e((u,g)) + e((g,f))$ by \cref{metric_on,metric_triangle_inequality,times_tuple_intro}.
            $\rho((u,g)) + \rho((g,f)) < \delta + \rho((g,f))$
                by \cref{real_lt_subst_left,real_lt_subst_right,real_lt_add,metric_on,funs_type_apply,times_tuple_intro,real_add_closed}.
            $\delta + \rho((g,f)) = r$ by
                \cref{real_sub,real_add_inverse,real_add_closed,real_add_assoc,real_add_comm,real_add_ident}.
            $e((u,g)) + e((g,f)) < r$ by \cref{real_lt_subst_left,real_lt_subst_right}.
            Thus $e((u,f)) < r$ by \cref{real_leq_lt_trans,metric_on,funs_type_apply,times_tuple_intro,real_add_closed}.
            Therefore $u\in V$ by \cref{real_lt_subst_left,real_lt_subst_right,unit_interval_robust_large_secant_class_49,subseteq,continuous_function_set}.
        \end{subproof}
        Therefore there exists $\epsilon\in\reals$ such that $0 < \epsilon$ and $\metricBall{e}{C}{g}{\epsilon}\subseteq V$ by \cref{subseteq}.
    \end{subproof}
    Therefore $V$ is open in $C$ with respect to $\metricTopology{e}{C}$ by \cref{metric_open_iff,unit_interval_robust_large_secant_class_49,subseteq,continuous,funs,continuous_function_set}.
\end{proof}

% from §49 helper lemma 25: pointwise bounds below one control the uniform metric from above
\begin{lemma}\label{uniform_metric_leq_from_pointwise_bound_49}
    Assume $\rho$ is the uniform metric on functions from $I$ to $\reals$ with respect to $\standardMetricR$.
    Assume $f\in\funs{I}{\reals}$.
    Assume $g\in\funs{I}{\reals}$.
    Assume $\delta\in\reals$.
    Suppose $\delta \leq 1$.
    Suppose for all $x\in I$ we have $\apply{\standardMetricR}{(\apply{f}{x},\apply{g}{x})} \leq \delta$.
    Then $\rho((f,g)) \leq \delta$.
\end{lemma}
\begin{proof}
    Follows by \cref{standard_metric_is_metric,uniform_metric_on_function_space,subseteq,real_lt_subst_left,real_leq_trans,metric_on,funs_type_apply,times_tuple_intro,real_upper_bound,real_supremum}.
\end{proof}

% from §49 helper lemma 25a: left-half points are closer to the left endpoint
\begin{lemma}\label{real_left_half_endpoint_gap_49}
    Assume $a,x,t\in\reals$.
    Suppose $t + t = a$ and $x \leq t$.
    Then $x \leq a - x$.
\end{lemma}
\begin{proof}
    Follows by \cref{real_leq_add,real_add_inverse,real_add_assoc,real_add_ident,real_add_comm,real_sub,real_lt_subst_left,real_lt_subst_right,real_add_closed,real_leq_trans}.
\end{proof}

% from §49 helper lemma 25b: right-half points are closer to the right endpoint
\begin{lemma}\label{real_right_half_endpoint_gap_49}
    Assume $a,x,t\in\reals$.
    Suppose $t + t = a$ and $t \leq x$.
    Then $a - x \leq x$.
\end{lemma}
\begin{proof}
    Follows by \cref{real_leq_add,real_add_inverse,real_add_assoc,real_add_ident,real_add_comm,real_sub,real_lt_subst_left,real_lt_subst_right,real_add_closed,real_leq_trans}.
\end{proof}

% from §49 helper lemma 26: on the left half of the first mesh interval the distance equals the coordinate
\begin{lemma}\label{unit_interval_equal_mesh_left_half_distance_49}
    Assume $A$ is the equal mesh of order $m$ on the unit interval.
    Assume $q\in\funs{\intervalclosed{0}{1}}{\reals}$.
    Assume for all $x\in\intervalclosed{0}{1}$ we have
        $q(x)$ is the distance from $x$ to $A$ in $\reals$ with respect to $\standardMetricR$.
    Assume $x\in\intervalclosed{0}{1}$.
    Suppose $x \leq (1 / m) / (1 + 1)$.
    Then $q(x) = x$.
\end{lemma}
\begin{proof}
    We have $m\in\naturalsPlus$ and
        $A = \{y\in\reals \mid \text{$y = 0$ or there exists $p\in\initialSegment{m}$ such that $y = p / m$}\}$
        by \cref{unit_interval_equal_mesh_49}.
    $x\in\reals$ and $0 \leq x$ by \cref{intervalclosed}.
    $1 / m\in\reals$ and $0 < 1 / m$ by \cref{naturalsplus_reciprocal_real,positive_reciprocal_naturalsplus}.
    Let $t = (1 / m) / (1 + 1)$.
    $t\in\reals$ and $0 < t$ and $t + t = 1 / m$ by \cref{real_half_of_positive}.
    $t < 1 / m$ by \cref{real_add_ident,real_lt_add,real_lt_subst_left,real_lt_subst_right}.
    $x \leq t$ by \cref{real_lt_subst_right}.
    $x \leq (1 / m) - x$ by \cref{real_left_half_endpoint_gap_49}.
    $x \leq 1 / m$ by \cref{real_leq_trans}.
    $0\in A$ by \cref{real_add_ident,subseteq}.
    $1 / m\in A$
        by \cref{real_one_in_naturals,real_one_leq_naturalsplus,initial_segment_naturalsplus,subseteq,naturalsplus_reciprocal_real}.
    $\apply{\standardMetricR}{(x,1 / m)} = \abs{(x - 1 / m)}$ by \cref{standard_metric_value}.
    $\apply{\standardMetricR}{(x,0)} = x$ by \cref{standard_metric_zero_value_nonnegative_49}.
    $(1 / m) - x\in\reals$ by \cref{real_sub,real_add_inverse,real_add_closed}.
    $0 \leq (1 / m) - x$ by \cref{real_leq_trans}.
    $\abs{((1 / m) - x)} = (1 / m) - x$
        by \cref{real_abs_nonneg,real_sub,real_add_inverse,real_add_closed}.
    $\abs{(x - 1 / m)} = \abs{((1 / m) - x)}$ by \cref{real_abs_sub_swap}.
    $\abs{(x - 1 / m)} = (1 / m) - x$ by \cref{real_lt_subst_right}.
    $\apply{\standardMetricR}{(x,1 / m)} = (1 / m) - x$ by \cref{real_lt_subst_right}.
    $\apply{\standardMetricR}{(x,0)} \leq \apply{\standardMetricR}{(x,1 / m)}$ by \cref{real_lt_subst_left,real_lt_subst_right}.
    Let $r = q(x)$.
    $A\subseteq\reals$ by \cref{subseteq}.
    $r$ is the distance from $x$ to $A$ in $\reals$ with respect to $\standardMetricR$ by \cref{subseteq}.
    For all $c\in A$ if $c\neq 0$ and $c\neq 1 / m$ then $c \leq 0$ or $1 / m \leq c$
        by \cref{unit_interval_equal_mesh_first_step_separation_49}.
    $r = \apply{\standardMetricR}{(x,0)}$
        by \cref{point_set_distance_left_endpoint_realized_49,standard_metric_is_metric,real_lt_subst_right}.
    Hence $q(x) = x$ by \cref{real_lt_subst_left,real_lt_subst_right}.
\end{proof}

% from §49 helper lemma 27: on the right half of the first mesh interval the distance equals the gap to the first mesh point
\begin{lemma}\label{unit_interval_equal_mesh_right_half_distance_49}
    Assume $A$ is the equal mesh of order $m$ on the unit interval.
    Assume $q\in\funs{\intervalclosed{0}{1}}{\reals}$.
    Assume for all $x\in\intervalclosed{0}{1}$ we have
        $q(x)$ is the distance from $x$ to $A$ in $\reals$ with respect to $\standardMetricR$.
    Assume $x\in\intervalclosed{0}{1}$.
    Suppose $(1 / m) / (1 + 1) \leq x$ and $x \leq 1 / m$.
    Then $q(x) = (1 / m) - x$.
\end{lemma}
\begin{proof}
    $m\in\naturalsPlus$ by \cref{unit_interval_equal_mesh_49}.
    $x\in\reals$ and $0 \leq x$ by \cref{intervalclosed}.
    $1 / m\in\reals$ and $0 < 1 / m$ by \cref{naturalsplus_reciprocal_real,positive_reciprocal_naturalsplus}.
    Let $t = (1 / m) / (1 + 1)$.
    $t\in\reals$ and $0 < t$ and $t + t = 1 / m$ by \cref{real_half_of_positive}.
    $t \leq x$ by \cref{real_lt_subst_left}.
    $(1 / m) - x \leq x$ by \cref{real_right_half_endpoint_gap_49}.
    $0\in A$ by \cref{unit_interval_equal_mesh_49,real_add_ident,subseteq}.
    $1 / m\in A$
        by \cref{unit_interval_equal_mesh_49,real_one_in_naturals,real_one_leq_naturalsplus,initial_segment_naturalsplus,subseteq,naturalsplus_reciprocal_real}.
    $\apply{\standardMetricR}{(x,1 / m)} = \abs{(x - 1 / m)}$ by \cref{standard_metric_value}.
    $\apply{\standardMetricR}{(x,0)} = x$ by \cref{standard_metric_zero_value_nonnegative_49}.
    $0 \leq (1 / m) - x$ by \cref{real_leq_add,real_add_inverse,real_add_ident,real_sub,real_lt_subst_left,real_lt_subst_right}.
    $\abs{((1 / m) - x)} = (1 / m) - x$
        by \cref{real_abs_nonneg,real_sub,real_add_inverse,real_add_closed}.
    $\abs{(x - 1 / m)} = \abs{((1 / m) - x)}$ by \cref{real_abs_sub_swap}.
    $\abs{(x - 1 / m)} = (1 / m) - x$ by \cref{real_lt_subst_right}.
    $\apply{\standardMetricR}{(x,1 / m)} = (1 / m) - x$ by \cref{real_lt_subst_right}.
    $\apply{\standardMetricR}{(x,1 / m)} \leq \apply{\standardMetricR}{(x,0)}$ by \cref{real_lt_subst_left,real_lt_subst_right}.
    Let $r = q(x)$.
    $A\subseteq\reals$ by \cref{unit_interval_equal_mesh_49,subseteq}.
    $r$ is the distance from $x$ to $A$ in $\reals$ with respect to $\standardMetricR$ by \cref{subseteq}.
    Show that for all $c\in A$ if $c\neq 0$ and $c\neq 1 / m$ then $c \leq 0$ or $1 / m \leq c$.
    \begin{subproof}
        Follows by \cref{unit_interval_equal_mesh_first_step_separation_49}.
    \end{subproof}
    $r = \apply{\standardMetricR}{(x,1 / m)}$
        by \cref{point_set_distance_right_endpoint_realized_49,standard_metric_is_metric,real_lt_subst_right}.
    Hence $q(x) = (1 / m) - x$ by \cref{real_lt_subst_left,real_lt_subst_right}.
\end{proof}

% from §49 helper lemma 28: dense subsets of metric spaces meet every positive ball
\begin{lemma}\label{dense_metric_ball_intersection_49}
    Assume $d$ is a metric on $X$.
    Assume $A\subseteq X$.
    Assume $A$ is dense in $X$ with respect to $\metricTopology{d}{X}$.
    Assume $x\in X$.
    Assume $\epsilon\in\reals$ and $0 < \epsilon$.
    Then there exists $y\in A$ such that $d((x,y)) < \epsilon$.
\end{lemma}
\begin{proof}
    Follows by \cref{metric_basis_is_basis,basis_topology_is_topology,metric_topology,basis_for_topology,metric_ball_neighborhood_metric_topology,neighborhood,metric_ball_center,dense_in,real_lt_subst_right,closure_iff_intersects_open,intersects,nonempty_inhabited,inter,metric_ball}.
\end{proof}

% from §49 helper lemma 28a: the unit-interval continuous function space is complete
\begin{lemma}\label{unit_interval_continuous_function_space_complete_49}
    Let $I = \intervalclosed{0}{1}$.
    Let $T_I = \subspaceTopology{\standardTopologyR}{I}$.
    Assume $C$ is the set of continuous functions from $I$ to $\reals$ with respect to $T_I$ and $\standardTopologyR$.
    Assume $\rho$ is the uniform metric on functions from $I$ to $\reals$ with respect to $\standardMetricR$.
    Let $e = \restrl{\rho}{C\times C}$.
    Then $C$ is complete with respect to $e$.
\end{lemma}
\begin{proof}
    Follows by \cref{real_add_ident,real_one_in_naturals,real_naturals_subset_reals,subseteq,real_zero_lt_one,intervalclosed,standard_metric_is_metric,reals_complete_standard_metric,standard_metric_topology,standard_topology_reals,standard_basis_is_basis,basis_topology_is_topology,subspace_topology_is_topology,continuous_function_set,continuous_function_space_complete_uniform}.
\end{proof}

% from §49 helper lemma 29: the unit-interval continuous function space is Baire
\begin{lemma}\label{unit_interval_continuous_function_space_baire_49}
    Let $I = \intervalclosed{0}{1}$.
    Let $T_I = \subspaceTopology{\standardTopologyR}{I}$.
    Assume $C$ is the set of continuous functions from $I$ to $\reals$ with respect to $T_I$ and $\standardTopologyR$.
    Assume $\rho$ is the uniform metric on functions from $I$ to $\reals$ with respect to $\standardMetricR$.
    Let $e = \restrl{\rho}{C\times C}$.
    Then $C$ is a Baire space with respect to $\metricTopology{e}{C}$.
\end{lemma}
\begin{proof}
    Follows by \cref{unit_interval_continuous_function_space_complete_49,complete_metric_space,complete_metric_space_baire}.
\end{proof}

% from §49 helper lemma 30: continuous functions belong to the continuous function space
\begin{lemma}\label{continuous_function_set_intro_49}
    Assume $C$ is the set of continuous functions from $I$ to $\reals$ with respect to $T_I$ and $\standardTopologyR$.
    Assume $f$ is continuous from $I$ to $\reals$ with respect to $T_I$ and $\standardTopologyR$.
    Then $f\in C$.
\end{lemma}
\begin{proof}
    Follows by \cref{continuous_function_set,continuous,funs}.
\end{proof}

% from §49 helper lemma 31: continuous function space members are continuous
\begin{lemma}\label{continuous_function_set_elim_49}
    Assume $C$ is the set of continuous functions from $I$ to $\reals$ with respect to $T_I$ and $\standardTopologyR$.
    Suppose $f\in C$.
    Then $f$ is continuous from $I$ to $\reals$ with respect to $T_I$ and $\standardTopologyR$.
\end{lemma}
\begin{proof}
    Follows by \cref{continuous_function_set,subseteq}.
\end{proof}

% from §49 helper lemma 32: the restricted uniform metric is a metric on the continuous function space
\begin{lemma}\label{unit_interval_continuous_function_space_metric_49}
    Let $I = \intervalclosed{0}{1}$.
    Let $T_I = \subspaceTopology{\standardTopologyR}{I}$.
    Assume $C$ is the set of continuous functions from $I$ to $\reals$ with respect to $T_I$ and $\standardTopologyR$.
    Assume $\rho$ is the uniform metric on functions from $I$ to $\reals$ with respect to $\standardMetricR$.
    Let $e = \restrl{\rho}{C\times C}$.
    Then $e$ is a metric on $C$.
\end{lemma}
\begin{proof}
    Follows by \cref{unit_interval_continuous_function_space_complete_49,complete_metric_space}.
\end{proof}

% from §49 helper lemma 33: the restricted uniform metric topology is a topology
\begin{lemma}\label{unit_interval_continuous_function_space_metric_topology_49}
    Let $I = \intervalclosed{0}{1}$.
    Let $T_I = \subspaceTopology{\standardTopologyR}{I}$.
    Assume $C$ is the set of continuous functions from $I$ to $\reals$ with respect to $T_I$ and $\standardTopologyR$.
    Assume $\rho$ is the uniform metric on functions from $I$ to $\reals$ with respect to $\standardMetricR$.
    Let $e = \restrl{\rho}{C\times C}$.
    Then $\metricTopology{e}{C}$ is a topology on $C$.
\end{lemma}
\begin{proof}
    Follows by \cref{unit_interval_continuous_function_space_metric_49,metric_basis_is_basis,basis_topology_is_topology,metric_topology,basis_for_topology}.
\end{proof}

% from §49 helper lemma 34: the restricted uniform metric agrees with the uniform metric on continuous functions
\begin{lemma}\label{restricted_uniform_metric_apply_49}
    Assume $C$ is the set of continuous functions from $I$ to $\reals$ with respect to $T_I$ and $\standardTopologyR$.
    Assume $\rho$ is the uniform metric on functions from $I$ to $\reals$ with respect to $\standardMetricR$.
    Let $e = \restrl{\rho}{C\times C}$.
    Suppose $f\in C$ and $g\in C$.
    Then $e((f,g)) = \rho((f,g))$.
\end{lemma}
\begin{proof}
    Follows by \cref{continuous_function_set,subseteq,uniform_metric_on_function_space,metric_on,funs_elim,funs_is_function,times_subseteq_left,times_subseteq_right,subseteq_transitive,pair_eq_iff,times_tuple_intro,restrl_of_function_apply}.
\end{proof}

% from §49 helper lemma 35: restricted uniform metric balls give uniform metric bounds
\begin{lemma}\label{restricted_uniform_metric_ball_elim_49}
    Assume $C$ is the set of continuous functions from $I$ to $\reals$ with respect to $T_I$ and $\standardTopologyR$.
    Assume $\rho$ is the uniform metric on functions from $I$ to $\reals$ with respect to $\standardMetricR$.
    Let $e = \restrl{\rho}{C\times C}$.
    Assume $h\in C$.
    Assume $g\in\metricBall{e}{C}{h}{\epsilon}$.
    Then $g\in C$ and $\rho((h,g)) < \epsilon$.
\end{lemma}
\begin{proof}
    Follows by \cref{metric_ball,restricted_uniform_metric_apply_49,real_lt_subst_left}.
\end{proof}

% from §49 helper lemma 36: uniform metric bounds give pointwise absolute-value bounds
\begin{lemma}\label{uniform_metric_pointwise_abs_lt_49}
    Assume $\rho$ is the uniform metric on functions from $I$ to $\reals$ with respect to $\standardMetricR$.
    Assume $f\in\funs{I}{\reals}$.
    Assume $g\in\funs{I}{\reals}$.
    Assume $x\in I$.
    Assume $\epsilon\in\reals$.
    Suppose $\rho((f,g)) < \epsilon$.
    Suppose $\epsilon \leq 1$.
    Then $\abs{(\apply{f}{x} - \apply{g}{x})} < \epsilon$.
\end{lemma}
\begin{proof}
    Follows by \cref{uniform_metric_pointwise_lt_below_one_49,standard_metric_value,funs_type_apply,real_lt_subst_left}.
\end{proof}

% from §49 helper lemma 37: robust large-secant classes lie in the continuous function space
\begin{lemma}\label{unit_interval_robust_large_secant_class_subset_continuous_49}
    Assume $C$ is the set of continuous functions from $I$ to $\reals$ with respect to $T_I$ and $\standardTopologyR$.
    Assume $V$ is the robust secant class above $n$ on $I$ with respect to $T_I$ and $\rho$.
    Then $V\subseteq C$.
\end{lemma}
\begin{proof}
    Follows by \cref{unit_interval_robust_large_secant_class_49,subseteq,continuous_function_set_intro_49}.
\end{proof}

% from §49 helper lemma 38: large-secant classes lie in the continuous function space
\begin{lemma}\label{unit_interval_large_secant_class_subset_continuous_49}
    Assume $C$ is the set of continuous functions from $I$ to $\reals$ with respect to $T_I$ and $\standardTopologyR$.
    Assume $U$ is the large-secant class above $n$ on $I$ with respect to $T_I$.
    Then $U\subseteq C$.
\end{lemma}
\begin{proof}
    Follows by \cref{unit_interval_large_secant_class_49,subseteq,continuous_function_set_intro_49}.
\end{proof}

% from §49 helper lemma 39: robust large-secant classes are open in the continuous function space
\begin{lemma}\label{unit_interval_robust_large_secant_class_open_auto_49}
    Let $I = \intervalclosed{0}{1}$.
    Let $T_I = \subspaceTopology{\standardTopologyR}{I}$.
    Assume $C$ is the set of continuous functions from $I$ to $\reals$ with respect to $T_I$ and $\standardTopologyR$.
    Assume $\rho$ is the uniform metric on functions from $I$ to $\reals$ with respect to $\standardMetricR$.
    Let $e = \restrl{\rho}{C\times C}$.
    Assume $n\in\naturalsPlus$.
    Assume $V$ is the robust secant class above $n$ on $I$ with respect to $T_I$ and $\rho$.
    Then $V$ is open in $C$ with respect to $\metricTopology{e}{C}$.
\end{lemma}
\begin{proof}
    Follows by \cref{unit_interval_continuous_function_space_metric_49,unit_interval_robust_large_secant_class_open_49}.
\end{proof}

% from §49 helper lemma 40: dense subsets of the continuous function space give uniform approximants
\begin{lemma}\label{dense_continuous_function_space_uniform_approximant_49}
    Let $I = \intervalclosed{0}{1}$.
    Let $T_I = \subspaceTopology{\standardTopologyR}{I}$.
    Assume $C$ is the set of continuous functions from $I$ to $\reals$ with respect to $T_I$ and $\standardTopologyR$.
    Assume $\rho$ is the uniform metric on functions from $I$ to $\reals$ with respect to $\standardMetricR$.
    Let $e = \restrl{\rho}{C\times C}$.
    Assume $D\subseteq C$.
    Assume $D$ is dense in $C$ with respect to $\metricTopology{e}{C}$.
    Assume $h\in C$.
    Assume $\epsilon\in\reals$ and $0 < \epsilon$.
    Then there exists $g\in D$ such that $\rho((h,g)) < \epsilon$.
\end{lemma}
\begin{proof}
    Follows by \cref{unit_interval_continuous_function_space_metric_49,dense_metric_ball_intersection_49,subseteq,restricted_uniform_metric_apply_49,real_lt_subst_left}.
\end{proof}

% from §49 helper lemma 41: positive radii can be reduced below one
\begin{lemma}\label{positive_radius_below_one_49}
    Assume $\epsilon\in\reals$ and $0 < \epsilon$.
    Then there exists $\eta\in\reals$ such that $0 < \eta$ and $\eta \leq 1$ and $\eta \leq \epsilon$.
\end{lemma}
\begin{proof}
    Follows by \cref{real_one_in_naturals,real_naturals_subset_reals,subseteq,real_zero_lt_one,real_lt_trichotomy,real_lt_imp_leq,real_lt_subst_left}.
\end{proof}

% from §49 helper lemma 42: dense nowhere-differentiable sets yield pointwise approximations
\begin{lemma}\label{dense_nowhere_differentiable_set_gives_approximation_49}
    Let $I = \intervalclosed{0}{1}$.
    Let $T_I = \subspaceTopology{\standardTopologyR}{I}$.
    Assume $C$ is the set of continuous functions from $I$ to $\reals$ with respect to $T_I$ and $\standardTopologyR$.
    Assume $\rho$ is the uniform metric on functions from $I$ to $\reals$ with respect to $\standardMetricR$.
    Let $e = \restrl{\rho}{C\times C}$.
    Assume $D\subseteq C$.
    Assume $D$ is dense in $C$ with respect to $\metricTopology{e}{C}$.
    Suppose for all $f\in D$ we have $f$ is nowhere differentiable on $I$.
    Assume $h$ is continuous from $I$ to $\reals$ with respect to $T_I$ and $\standardTopologyR$.
    Assume $\epsilon\in\reals$ and $0 < \epsilon$.
    Then there exists $g$ such that
        $g$ is continuous from $I$ to $\reals$ with respect to $T_I$ and $\standardTopologyR$ and
        $g$ is nowhere differentiable on $I$ and
        for all $x\in I$ we have $\abs{(\apply{h}{x} - \apply{g}{x})} < \epsilon$.
\end{lemma}
\begin{proof}
    $h\in C$ by \cref{continuous_function_set_intro_49}.
    Take $\eta\in\reals$ such that $0 < \eta$ and $\eta \leq 1$ and $\eta \leq \epsilon$
        by \cref{positive_radius_below_one_49}.
    Take $g\in D$ such that $\rho((h,g)) < \eta$
        by \cref{dense_continuous_function_space_uniform_approximant_49}.
    Show that $g$ is continuous from $I$ to $\reals$ with respect to $T_I$ and $\standardTopologyR$.
    \begin{subproof}
        Follows by \cref{subseteq,continuous_function_set_elim_49}.
    \end{subproof}
    Show that $g$ is nowhere differentiable on $I$.
    \begin{subproof}
        Follows by \cref{subseteq}.
    \end{subproof}
    Show that for all $x\in I$ we have $\abs{(\apply{h}{x} - \apply{g}{x})} < \epsilon$.
    \begin{subproof}
        Follows by \cref{continuous,funs,uniform_metric_pointwise_abs_lt_49,funs_type_apply,real_sub,real_add_inverse,real_add_closed,real_abs_in_reals,real_lt_leq_trans_local}.
    \end{subproof}
\end{proof}

% from §49 helper lemma 43: metric ball intersections imply density
\begin{lemma}\label{metric_ball_intersections_imply_dense_49}
    Assume $d$ is a metric on $X$.
    Assume $A\subseteq X$.
    Suppose for all $x\in X$ we have for all $\epsilon\in\reals$ if $0 < \epsilon$ then
        $\metricBall{d}{X}{x}{\epsilon}$ intersects $A$.
    Then $A$ is dense in $X$ with respect to $\metricTopology{d}{X}$.
\end{lemma}
\begin{proof}
    Let $T = \metricTopology{d}{X}$.
    Show that for all $x$ such that $x\in X$ we have for all $U$ such that $U$ is open in $X$ with respect to $T$ and $x\in U$
        there exists $\delta\in\reals$ such that $0 < \delta$ and $\metricBall{d}{X}{x}{\delta}\subseteq U$.
    \begin{subproof}
        Follows by \cref{open_in,topology_on,subseteq,powerset_elim,metric_open_iff}.
    \end{subproof}
    Show that for all $x$ such that $x\in X$ we have for all $U$ such that $U$ is open in $X$ with respect to $T$ and $x\in U$
        $U$ intersects $A$.
    \begin{subproof}
        Follows by \cref{subseteq,intersects,nonempty_inhabited,inter_elim_left,inter_elim_right,inter_intro,inhabited_nonempty}.
    \end{subproof}
    Show that for all $x$ such that $x\in X$ we have $x\in\closure{A}{X}{T}$.
    \begin{subproof}
        Follows by \cref{metric_basis_is_basis,basis_topology_is_topology,metric_topology,closure_iff_intersects_open}.
    \end{subproof}
    Hence $\closure{A}{X}{T} = X$ by \cref{subseteq,closure,subseteq_antisymmetric}.
    Therefore $A$ is dense in $X$ with respect to $T$ by \cref{metric_basis_is_basis,basis_topology_is_topology,metric_topology,dense_in}.
\end{proof}

% from §49 helper lemma 44: ball intersections imply density in the continuous function space
\begin{lemma}\label{continuous_function_space_ball_intersections_imply_dense_49}
    Let $I = \intervalclosed{0}{1}$.
    Let $T_I = \subspaceTopology{\standardTopologyR}{I}$.
    Assume $C$ is the set of continuous functions from $I$ to $\reals$ with respect to $T_I$ and $\standardTopologyR$.
    Assume $\rho$ is the uniform metric on functions from $I$ to $\reals$ with respect to $\standardMetricR$.
    Let $e = \restrl{\rho}{C\times C}$.
    Assume $A\subseteq C$.
    Suppose for all $f\in C$ we have for all $\epsilon\in\reals$ if $0 < \epsilon$ then
        $\metricBall{e}{C}{f}{\epsilon}$ intersects $A$.
    Then $A$ is dense in $C$ with respect to $\metricTopology{e}{C}$.
\end{lemma}
\begin{proof}
    Follows by \cref{unit_interval_continuous_function_space_metric_49,metric_ball_intersections_imply_dense_49}.
\end{proof}

% from §49 helper lemma 45: choose a positive real below two positive bounds
\begin{lemma}\label{positive_radius_below_two_bounds_49}
    Assume $a\in\reals$ and $0 < a$.
    Assume $b\in\reals$ and $0 < b$.
    Then there exists $r\in\reals$ such that $0 < r$ and $r \leq a$ and $r \leq b$.
\end{lemma}
\begin{proof}
    Follows by \cref{real_lt_trichotomy,real_lt_imp_leq,real_lt_subst_left}.
\end{proof}

% from §49 helper lemma 46: choose a positive real below three positive bounds
\begin{lemma}\label{positive_radius_below_three_bounds_49}
    Assume $a\in\reals$ and $0 < a$.
    Assume $b\in\reals$ and $0 < b$.
    Assume $c\in\reals$ and $0 < c$.
    Then there exists $r\in\reals$ such that $0 < r$ and $r \leq a$ and $r \leq b$ and $r \leq c$.
\end{lemma}
\begin{proof}
    Follows by \cref{positive_radius_below_two_bounds_49,real_leq_trans}.
\end{proof}

% from §49 helper lemma 47: choose a positive real below four positive bounds
\begin{lemma}\label{positive_radius_below_four_bounds_49}
    Assume $a\in\reals$ and $0 < a$.
    Assume $b\in\reals$ and $0 < b$.
    Assume $c\in\reals$ and $0 < c$.
    Assume $d\in\reals$ and $0 < d$.
    Then there exists $r\in\reals$ such that $0 < r$ and $r \leq a$ and $r \leq b$ and $r \leq c$ and $r \leq d$.
\end{lemma}
\begin{proof}
    Take $s\in\reals$ such that $0 < s$ and $s \leq a$ and $s \leq b$ and $s \leq c$
        by \cref{positive_radius_below_three_bounds_49}.
    Take $r\in\reals$ such that $0 < r$ and $r \leq s$ and $r \leq d$
        by \cref{positive_radius_below_two_bounds_49}.
    Follows by \cref{real_leq_trans}.
\end{proof}

% from §49 helper lemma 48: choose a positive mesh step below a positive bound
\begin{lemma}\label{positive_mesh_step_below_bound_49}
    Assume $a\in\reals$ and $0 < a$.
    Then there exists $m\in\naturalsPlus$ such that $0 < 1 / m$ and $1 / m \leq a$.
\end{lemma}
\begin{proof}
    Follows by \cref{real_small_reciprocal,real_lt_imp_leq}.
\end{proof}

% from §49 helper lemma 49: a quarter mesh step is positive and bounded by the mesh step
\begin{lemma}\label{quarter_mesh_step_positive_bounded_49}
    Assume $m\in\naturalsPlus$.
    Let $h = ((1 / m) / (1 + 1)) / (1 + 1)$.
    Then $h\in\reals$ and $0 < h$ and $h \leq 1 / m$.
\end{lemma}
\begin{proof}
    Follows by \cref{naturalsplus_reciprocal_real,positive_reciprocal_naturalsplus,real_half_of_positive,real_add_ident,real_lt_add,real_lt_subst_left,real_lt_subst_right,real_lt_trans,real_lt_imp_leq}.
\end{proof}

% from §49 helper lemma 50: a quarter mesh step inherits an upper bound
\begin{lemma}\label{quarter_mesh_step_leq_bound_49}
    Assume $m\in\naturalsPlus$.
    Assume $a\in\reals$.
    Suppose $1 / m \leq a$.
    Let $h = ((1 / m) / (1 + 1)) / (1 + 1)$.
    Then $h \leq a$.
\end{lemma}
\begin{proof}
    Follows by \cref{quarter_mesh_step_positive_bounded_49,naturalsplus_reciprocal_real,real_leq_trans}.
\end{proof}

% from §49 helper lemma 51: scalar multiples of continuous real-valued functions are continuous
\begin{lemma}\label{real_function_scalar_multiple_continuous_49}
    Assume $T$ is a topology on $X$.
    Assume $f$ is continuous from $X$ to $\reals$ with respect to $T$ and $\standardTopologyR$.
    Assume $a\in\reals$.
    Let $c = \{(x,a) \mid x\in X\}$.
    Let $g = \{(x,c(x) \rmul f(x)) \mid x\in X\}$.
    Then $g$ is continuous from $X$ to $\reals$ with respect to $T$ and $\standardTopologyR$.
\end{lemma}
\begin{proof}
    $\standardTopologyR$ is a topology on $\reals$
        by \cref{standard_topology_reals,standard_basis_is_basis,basis_topology_is_topology}.
    $c$ is a function from $X$ to $\reals$ by \cref{constant_map_function}.
    For all $x\in X$ we have $c(x) = a$ by \cref{constant_map_function,funs_is_function,function_apply_intro}.
    $c$ is continuous from $X$ to $\reals$ with respect to $T$ and $\standardTopologyR$ by \cref{constant_continuous}.
    Therefore $g$ is continuous from $X$ to $\reals$ with respect to $T$ and $\standardTopologyR$
        by \cref{real_function_multiplication_continuous}.
\end{proof}

% from §49 helper lemma 52: adding a scalar multiple preserves continuity
\begin{lemma}\label{real_function_add_scalar_multiple_continuous_49}
    Assume $T$ is a topology on $X$.
    Assume $f$ is continuous from $X$ to $\reals$ with respect to $T$ and $\standardTopologyR$.
    Assume $q$ is continuous from $X$ to $\reals$ with respect to $T$ and $\standardTopologyR$.
    Assume $a\in\reals$.
    Let $c = \{(x,a) \mid x\in X\}$.
    Let $s = \{(x,c(x) \rmul q(x)) \mid x\in X\}$.
    Let $g = \{(x,f(x) + s(x)) \mid x\in X\}$.
    Then $g$ is continuous from $X$ to $\reals$ with respect to $T$ and $\standardTopologyR$.
\end{lemma}
\begin{proof}
    Follows by \cref{real_function_scalar_multiple_continuous_49,real_function_addition_continuous}.
\end{proof}

% from §49 helper lemma 53: a witness inside a metric ball gives intersection
\begin{lemma}\label{metric_ball_intersection_witness_49}
    Assume $d$ is a metric on $X$.
    Assume $A\subseteq X$.
    Assume $x\in X$.
    Assume $y\in A$.
    Assume $\epsilon\in\reals$.
    Suppose $d((x,y)) < \epsilon$.
    Then $\metricBall{d}{X}{x}{\epsilon}$ intersects $A$.
\end{lemma}
\begin{proof}
    Follows by \cref{subseteq,metric_ball,inter_intro,intersects,inhabited_nonempty}.
\end{proof}

% from §49 helper lemma 54: choose a positive real below five positive bounds
\begin{lemma}\label{positive_radius_below_five_bounds_49}
    Assume $a\in\reals$ and $0 < a$.
    Assume $b\in\reals$ and $0 < b$.
    Assume $c\in\reals$ and $0 < c$.
    Assume $d\in\reals$ and $0 < d$.
    Assume $e\in\reals$ and $0 < e$.
    Then there exists $r\in\reals$ such that $0 < r$ and $r \leq a$ and $r \leq b$ and $r \leq c$ and $r \leq d$ and $r \leq e$.
\end{lemma}
\begin{proof}
    Take $s\in\reals$ such that $0 < s$ and $s \leq a$ and $s \leq b$ and $s \leq c$ and $s \leq d$
        by \cref{positive_radius_below_four_bounds_49}.
    Take $r\in\reals$ such that $0 < r$ and $r \leq s$ and $r \leq e$
        by \cref{positive_radius_below_two_bounds_49}.
    Thus $r\in\reals$ and $0 < r$ and $r \leq a$ and $r \leq b$ and $r \leq c$ and $r \leq d$ and $r \leq e$
        by \cref{real_leq_trans}.
\end{proof}

% from §49 helper lemma 55: choose a positive real below six positive bounds
\begin{lemma}\label{positive_radius_below_six_bounds_49}
    Assume $a\in\reals$ and $0 < a$.
    Assume $b\in\reals$ and $0 < b$.
    Assume $c\in\reals$ and $0 < c$.
    Assume $d\in\reals$ and $0 < d$.
    Assume $e\in\reals$ and $0 < e$.
    Assume $p\in\reals$ and $0 < p$.
    Then there exists $r\in\reals$ such that $0 < r$ and $r \leq a$ and $r \leq b$ and $r \leq c$ and $r \leq d$ and $r \leq e$ and $r \leq p$.
\end{lemma}
\begin{proof}
    Take $s\in\reals$ such that $0 < s$ and $s \leq a$ and $s \leq b$ and $s \leq c$ and $s \leq d$ and $s \leq e$
        by \cref{positive_radius_below_five_bounds_49}.
    Take $r\in\reals$ such that $0 < r$ and $r \leq s$ and $r \leq p$
        by \cref{positive_radius_below_two_bounds_49}.
    Thus $r\in\reals$ and $0 < r$ and $r \leq a$ and $r \leq b$ and $r \leq c$ and $r \leq d$ and $r \leq e$ and $r \leq p$
        by \cref{real_leq_trans}.
\end{proof}

% from §49 helper lemma 56: choose a positive radius below one and another positive bound
\begin{lemma}\label{positive_radius_below_one_and_bound_49}
    Assume $a\in\reals$ and $0 < a$.
    Then there exists $r\in\reals$ such that $0 < r$ and $r \leq 1$ and $r \leq a$.
\end{lemma}
\begin{proof}
    Follows by \cref{positive_radius_below_one_49}.
\end{proof}

% from §49 helper lemma 57: choose a mesh step below two positive bounds
\begin{lemma}\label{positive_mesh_step_below_two_bounds_49}
    Assume $a\in\reals$ and $0 < a$.
    Assume $b\in\reals$ and $0 < b$.
    Then there exists $m\in\naturalsPlus$ such that $0 < 1 / m$ and $1 / m \leq a$ and $1 / m \leq b$.
\end{lemma}
\begin{proof}
    Follows by \cref{positive_radius_below_two_bounds_49,positive_mesh_step_below_bound_49,real_leq_trans,naturalsplus_reciprocal_real}.
\end{proof}

% from §49 helper lemma 58: choose a mesh step below three positive bounds
\begin{lemma}\label{positive_mesh_step_below_three_bounds_49}
    Assume $a\in\reals$ and $0 < a$.
    Assume $b\in\reals$ and $0 < b$.
    Assume $c\in\reals$ and $0 < c$.
    Then there exists $m\in\naturalsPlus$ such that $0 < 1 / m$ and $1 / m \leq a$ and $1 / m \leq b$ and $1 / m \leq c$.
\end{lemma}
\begin{proof}
    Follows by \cref{positive_radius_below_three_bounds_49,positive_mesh_step_below_bound_49,real_leq_trans,naturalsplus_reciprocal_real}.
\end{proof}

% from §49 helper lemma 59: choose a mesh step below four positive bounds
\begin{lemma}\label{positive_mesh_step_below_four_bounds_49}
    Assume $a\in\reals$ and $0 < a$.
    Assume $b\in\reals$ and $0 < b$.
    Assume $c\in\reals$ and $0 < c$.
    Assume $d\in\reals$ and $0 < d$.
    Then there exists $m\in\naturalsPlus$ such that $0 < 1 / m$ and $1 / m \leq a$ and $1 / m \leq b$ and $1 / m \leq c$ and $1 / m \leq d$.
\end{lemma}
\begin{proof}
    Take $r\in\reals$ such that $0 < r$ and $r \leq a$ and $r \leq b$ and $r \leq c$ and $r \leq d$
        by \cref{positive_radius_below_four_bounds_49}.
    Take $m\in\naturalsPlus$ such that $0 < 1 / m$ and $1 / m \leq r$
        by \cref{positive_mesh_step_below_bound_49}.
    Thus there exists $k\in\naturalsPlus$ such that $0 < 1 / k$ and $1 / k \leq a$ and $1 / k \leq b$ and $1 / k \leq c$ and $1 / k \leq d$
        by \cref{real_leq_trans,naturalsplus_reciprocal_real}.
\end{proof}

% from §49 helper lemma 60: choose a mesh step below five positive bounds
\begin{lemma}\label{positive_mesh_step_below_five_bounds_49}
    Assume $a\in\reals$ and $0 < a$.
    Assume $b\in\reals$ and $0 < b$.
    Assume $c\in\reals$ and $0 < c$.
    Assume $d\in\reals$ and $0 < d$.
    Assume $e\in\reals$ and $0 < e$.
    Then there exists $m\in\naturalsPlus$ such that $0 < 1 / m$ and $1 / m \leq a$ and $1 / m \leq b$ and $1 / m \leq c$ and $1 / m \leq d$ and $1 / m \leq e$.
\end{lemma}
\begin{proof}
    Take $r\in\reals$ such that $0 < r$ and $r \leq a$ and $r \leq b$ and $r \leq c$ and $r \leq d$ and $r \leq e$
        by \cref{positive_radius_below_five_bounds_49}.
    Take $m\in\naturalsPlus$ such that $0 < 1 / m$ and $1 / m \leq r$
        by \cref{positive_mesh_step_below_bound_49}.
    Thus there exists $k\in\naturalsPlus$ such that $0 < 1 / k$ and $1 / k \leq a$ and $1 / k \leq b$ and $1 / k \leq c$ and $1 / k \leq d$ and $1 / k \leq e$
        by \cref{real_leq_trans,naturalsplus_reciprocal_real}.
\end{proof}

% from §49 helper lemma 61: choose a mesh step below six positive bounds
\begin{lemma}\label{positive_mesh_step_below_six_bounds_49}
    Assume $a\in\reals$ and $0 < a$.
    Assume $b\in\reals$ and $0 < b$.
    Assume $c\in\reals$ and $0 < c$.
    Assume $d\in\reals$ and $0 < d$.
    Assume $e\in\reals$ and $0 < e$.
    Assume $p\in\reals$ and $0 < p$.
    Then there exists $m\in\naturalsPlus$ such that $0 < 1 / m$ and $1 / m \leq a$ and $1 / m \leq b$ and $1 / m \leq c$ and $1 / m \leq d$ and $1 / m \leq e$ and $1 / m \leq p$.
\end{lemma}
\begin{proof}
    Take $r\in\reals$ such that $0 < r$ and $r \leq a$ and $r \leq b$ and $r \leq c$ and $r \leq d$ and $r \leq e$ and $r \leq p$
        by \cref{positive_radius_below_six_bounds_49}.
    Take $m\in\naturalsPlus$ such that $0 < 1 / m$ and $1 / m \leq r$
        by \cref{positive_mesh_step_below_bound_49}.
    Thus there exists $k\in\naturalsPlus$ such that $0 < 1 / k$ and $1 / k \leq a$ and $1 / k \leq b$ and $1 / k \leq c$ and $1 / k \leq d$ and $1 / k \leq e$ and $1 / k \leq p$
        by \cref{real_leq_trans,naturalsplus_reciprocal_real}.
\end{proof}

% from §49 helper lemma 62: choose a quarter mesh scale below two positive bounds
\begin{lemma}\label{quarter_mesh_step_below_two_bounds_49}
    Assume $a\in\reals$ and $0 < a$.
    Assume $b\in\reals$ and $0 < b$.
    Then there exist $m,h$ such that $m\in\naturalsPlus$ and $h\in\reals$ and $0 < h$ and $h \leq a$ and $h \leq b$ and $h = ((1 / m) / (1 + 1)) / (1 + 1)$.
\end{lemma}
\begin{proof}
    Follows by \cref{positive_mesh_step_below_two_bounds_49,quarter_mesh_step_positive_bounded_49,quarter_mesh_step_leq_bound_49}.
\end{proof}

% from §49 helper lemma 63: choose a quarter mesh scale below three positive bounds
\begin{lemma}\label{quarter_mesh_step_below_three_bounds_49}
    Assume $a\in\reals$ and $0 < a$.
    Assume $b\in\reals$ and $0 < b$.
    Assume $c\in\reals$ and $0 < c$.
    Then there exist $m,h$ such that $m\in\naturalsPlus$ and $h\in\reals$ and $0 < h$ and $h \leq a$ and $h \leq b$ and $h \leq c$ and $h = ((1 / m) / (1 + 1)) / (1 + 1)$.
\end{lemma}
\begin{proof}
    Follows by \cref{positive_mesh_step_below_three_bounds_49,quarter_mesh_step_positive_bounded_49,quarter_mesh_step_leq_bound_49}.
\end{proof}

% from §49 helper lemma 64: choose a quarter mesh scale below four positive bounds
\begin{lemma}\label{quarter_mesh_step_below_four_bounds_49}
    Assume $a\in\reals$ and $0 < a$.
    Assume $b\in\reals$ and $0 < b$.
    Assume $c\in\reals$ and $0 < c$.
    Assume $d\in\reals$ and $0 < d$.
    Then there exist $m,h$ such that $m\in\naturalsPlus$ and $h\in\reals$ and $0 < h$ and $h \leq a$ and $h \leq b$ and $h \leq c$ and $h \leq d$ and $h = ((1 / m) / (1 + 1)) / (1 + 1)$.
\end{lemma}
\begin{proof}
    Follows by \cref{positive_mesh_step_below_four_bounds_49,quarter_mesh_step_positive_bounded_49,quarter_mesh_step_leq_bound_49}.
\end{proof}

% from §49 helper lemma 65: choose a quarter mesh scale below five positive bounds
\begin{lemma}\label{quarter_mesh_step_below_five_bounds_49}
    Assume $a\in\reals$ and $0 < a$.
    Assume $b\in\reals$ and $0 < b$.
    Assume $c\in\reals$ and $0 < c$.
    Assume $d\in\reals$ and $0 < d$.
    Assume $e\in\reals$ and $0 < e$.
    Then there exist $m,h$ such that $m\in\naturalsPlus$ and $h\in\reals$ and $0 < h$ and $h \leq a$ and $h \leq b$ and $h \leq c$ and $h \leq d$ and $h \leq e$ and $h = ((1 / m) / (1 + 1)) / (1 + 1)$.
\end{lemma}
\begin{proof}
    Take $m\in\naturalsPlus$ such that $0 < 1 / m$ and $1 / m \leq a$ and $1 / m \leq b$ and $1 / m \leq c$ and $1 / m \leq d$ and $1 / m \leq e$
        by \cref{positive_mesh_step_below_five_bounds_49}.
    Let $h = ((1 / m) / (1 + 1)) / (1 + 1)$.
    Thus there exist $k,r$ such that $k\in\naturalsPlus$ and $r\in\reals$ and $0 < r$ and $r \leq a$ and $r \leq b$ and $r \leq c$ and $r \leq d$ and $r \leq e$ and $r = ((1 / k) / (1 + 1)) / (1 + 1)$
        by \cref{quarter_mesh_step_positive_bounded_49,quarter_mesh_step_leq_bound_49}.
\end{proof}

% from §49 helper lemma 66: choose a quarter mesh scale below six positive bounds
\begin{lemma}\label{quarter_mesh_step_below_six_bounds_49}
    Assume $a\in\reals$ and $0 < a$.
    Assume $b\in\reals$ and $0 < b$.
    Assume $c\in\reals$ and $0 < c$.
    Assume $d\in\reals$ and $0 < d$.
    Assume $e\in\reals$ and $0 < e$.
    Assume $p\in\reals$ and $0 < p$.
    Then there exist $m,h$ such that $m\in\naturalsPlus$ and $h\in\reals$ and $0 < h$ and $h \leq a$ and $h \leq b$ and $h \leq c$ and $h \leq d$ and $h \leq e$ and $h \leq p$ and $h = ((1 / m) / (1 + 1)) / (1 + 1)$.
\end{lemma}
\begin{proof}
    Take $m\in\naturalsPlus$ such that $0 < 1 / m$ and $1 / m \leq a$ and $1 / m \leq b$ and $1 / m \leq c$ and $1 / m \leq d$ and $1 / m \leq e$ and $1 / m \leq p$
        by \cref{positive_mesh_step_below_six_bounds_49}.
    Let $h = ((1 / m) / (1 + 1)) / (1 + 1)$.
    Thus there exist $k,r$ such that $k\in\naturalsPlus$ and $r\in\reals$ and $0 < r$ and $r \leq a$ and $r \leq b$ and $r \leq c$ and $r \leq d$ and $r \leq e$ and $r \leq p$ and $r = ((1 / k) / (1 + 1)) / (1 + 1)$
        by \cref{quarter_mesh_step_positive_bounded_49,quarter_mesh_step_leq_bound_49}.
\end{proof}

% from §49 helper lemma 67: choose a quarter mesh scale below one and another positive bound
\begin{lemma}\label{quarter_mesh_step_below_one_and_bound_49}
    Assume $a\in\reals$ and $0 < a$.
    Then there exist $m,h$ such that $m\in\naturalsPlus$ and $h\in\reals$ and $0 < h$ and $h \leq 1$ and $h \leq a$ and $h = ((1 / m) / (1 + 1)) / (1 + 1)$.
\end{lemma}
\begin{proof}
    Follows by \cref{real_one_in_naturals,real_naturals_subset_reals,subseteq,real_zero_lt_one,quarter_mesh_step_below_two_bounds_49}.
\end{proof}

% from §49 helper lemma 68: choose a quarter mesh scale below one and two positive bounds
\begin{lemma}\label{quarter_mesh_step_below_one_and_two_bounds_49}
    Assume $a\in\reals$ and $0 < a$.
    Assume $b\in\reals$ and $0 < b$.
    Then there exist $m,h$ such that $m\in\naturalsPlus$ and $h\in\reals$ and $0 < h$ and $h \leq 1$ and $h \leq a$ and $h \leq b$ and $h = ((1 / m) / (1 + 1)) / (1 + 1)$.
\end{lemma}
\begin{proof}
    Follows by \cref{real_one_in_naturals,real_naturals_subset_reals,subseteq,real_zero_lt_one,quarter_mesh_step_below_three_bounds_49}.
\end{proof}

% from §49 helper lemma 69: choose a quarter mesh scale below one and three positive bounds
\begin{lemma}\label{quarter_mesh_step_below_one_and_three_bounds_49}
    Assume $a\in\reals$ and $0 < a$.
    Assume $b\in\reals$ and $0 < b$.
    Assume $c\in\reals$ and $0 < c$.
    Then there exist $m,h$ such that $m\in\naturalsPlus$ and $h\in\reals$ and $0 < h$ and $h \leq 1$ and $h \leq a$ and $h \leq b$ and $h \leq c$ and $h = ((1 / m) / (1 + 1)) / (1 + 1)$.
\end{lemma}
\begin{proof}
    Follows by \cref{real_one_in_naturals,real_naturals_subset_reals,subseteq,real_zero_lt_one,quarter_mesh_step_below_four_bounds_49}.
\end{proof}

% from §49 helper lemma 70: choose a quarter mesh scale below one and four positive bounds
\begin{lemma}\label{quarter_mesh_step_below_one_and_four_bounds_49}
    Assume $a\in\reals$ and $0 < a$.
    Assume $b\in\reals$ and $0 < b$.
    Assume $c\in\reals$ and $0 < c$.
    Assume $d\in\reals$ and $0 < d$.
    Then there exist $m,h$ such that $m\in\naturalsPlus$ and $h\in\reals$ and $0 < h$ and $h \leq 1$ and $h \leq a$ and $h \leq b$ and $h \leq c$ and $h \leq d$ and $h = ((1 / m) / (1 + 1)) / (1 + 1)$.
\end{lemma}
\begin{proof}
    Follows by \cref{real_one_in_naturals,real_naturals_subset_reals,subseteq,real_zero_lt_one,quarter_mesh_step_below_five_bounds_49}.
\end{proof}

% from §49 helper lemma 71: choose a quarter mesh scale below one and five positive bounds
\begin{lemma}\label{quarter_mesh_step_below_one_and_five_bounds_49}
    Assume $a\in\reals$ and $0 < a$.
    Assume $b\in\reals$ and $0 < b$.
    Assume $c\in\reals$ and $0 < c$.
    Assume $d\in\reals$ and $0 < d$.
    Assume $e\in\reals$ and $0 < e$.
    Then there exist $m,h$ such that $m\in\naturalsPlus$ and $h\in\reals$ and $0 < h$ and $h \leq 1$ and $h \leq a$ and $h \leq b$ and $h \leq c$ and $h \leq d$ and $h \leq e$ and $h = ((1 / m) / (1 + 1)) / (1 + 1)$.
\end{lemma}
\begin{proof}
    $1\in\reals$ and $0 < 1$ by \cref{real_one_in_naturals,real_naturals_subset_reals,subseteq,real_zero_lt_one}.
    Take $m\in\naturalsPlus$ such that $0 < 1 / m$ and $1 / m \leq 1$ and $1 / m \leq a$ and $1 / m \leq b$ and $1 / m \leq c$ and $1 / m \leq d$ and $1 / m \leq e$
        by \cref{positive_mesh_step_below_six_bounds_49}.
    Let $h = ((1 / m) / (1 + 1)) / (1 + 1)$.
    Thus there exist $k,r$ such that $k\in\naturalsPlus$ and $r\in\reals$ and $0 < r$ and $r \leq 1$ and $r \leq a$ and $r \leq b$ and $r \leq c$ and $r \leq d$ and $r \leq e$ and $r = ((1 / k) / (1 + 1)) / (1 + 1)$
        by \cref{quarter_mesh_step_positive_bounded_49,quarter_mesh_step_leq_bound_49}.
\end{proof}

% from §49 helper lemma 72: choose a mesh step below one and another positive bound
\begin{lemma}\label{positive_mesh_step_below_one_and_bound_49}
    Assume $a\in\reals$ and $0 < a$.
    Then there exists $m\in\naturalsPlus$ such that $0 < 1 / m$ and $1 / m \leq 1$ and $1 / m \leq a$.
\end{lemma}
\begin{proof}
    Follows by \cref{real_one_in_naturals,real_naturals_subset_reals,subseteq,real_zero_lt_one,positive_mesh_step_below_two_bounds_49}.
\end{proof}

% from §49 helper lemma 73: choose a mesh step below one and two positive bounds
\begin{lemma}\label{positive_mesh_step_below_one_and_two_bounds_49}
    Assume $a\in\reals$ and $0 < a$.
    Assume $b\in\reals$ and $0 < b$.
    Then there exists $m\in\naturalsPlus$ such that $0 < 1 / m$ and $1 / m \leq 1$ and $1 / m \leq a$ and $1 / m \leq b$.
\end{lemma}
\begin{proof}
    Follows by \cref{real_one_in_naturals,real_naturals_subset_reals,subseteq,real_zero_lt_one,positive_mesh_step_below_three_bounds_49}.
\end{proof}

% from §49 helper lemma 74: a strong secant function is in its robust class with itself as witness
\begin{lemma}\label{unit_interval_robust_class_self_witness_49}
    Let $I = \intervalclosed{0}{1}$.
    Let $T_I = \subspaceTopology{\standardTopologyR}{I}$.
    Assume $\rho$ is the uniform metric on functions from $I$ to $\reals$ with respect to $\standardMetricR$.
    Assume $n\in\naturalsPlus$.
    Assume $V$ is the robust secant class above $n$ on $I$ with respect to $T_I$ and $\rho$.
    Assume $g$ is continuous from $I$ to $\reals$ with respect to $T_I$ and $\standardTopologyR$.
    Assume $h\in\reals$.
    Suppose $h \leq 1 / n$.
    Suppose $g$ has secants above $n + n$ at scale $h$ on $I$.
    Then $g\in V$.
\end{lemma}
\begin{proof}
    $g\in\funs{I}{\reals}$ by \cref{continuous,funs}.
    $\rho$ is a metric on $\funs{I}{\reals}$ by \cref{uniform_metric_on_function_space}.
    $0 < h$ by \cref{large_fixed_scale_secants_on_real_subset}.
    $n\in\reals$ and $0 < n$ by \cref{real_naturals_subset_reals,subseteq,naturalsplus_positive_real}.
    $h \rmul n\in\reals$ and $0 < h \rmul n$ by \cref{real_mul_closed,real_mul_pos_pos_gt}.
    Let $a = (h \rmul n) / (1 + 1)$.
    $a\in\reals$ and $0 < a$ by \cref{real_half_of_positive}.
    Let $b = a / (1 + 1)$.
    $b\in\reals$ and $0 < b$ by \cref{real_half_of_positive}.
    Let $r = b / (1 + 1)$.
    $r\in\reals$ and $0 < r$ by \cref{real_half_of_positive}.
    $(g,g)\in\funs{I}{\reals}\times\funs{I}{\reals}$ by \cref{times_tuple_intro}.
    $\rho((g,g)) = 0$ by \cref{metric_on,metric_zero_characterization}.
    $\rho((g,g)) < r$ by \cref{real_lt_subst_left}.
    $r = (((h \rmul n) / (1 + 1)) / (1 + 1)) / (1 + 1)$ by \cref{real_lt_subst_left}.
    Hence $\rho((g,g)) < (((h \rmul n) / (1 + 1)) / (1 + 1)) / (1 + 1)$ by \cref{real_lt_subst_right}.
    Therefore $g\in V$ by \cref{unit_interval_robust_large_secant_class_49,subseteq}.
\end{proof}

% from §49 helper lemma 75: pointwise values of adding a scalar multiple
\begin{lemma}\label{real_function_add_scalar_multiple_apply_49}
    Assume $T$ is a topology on $X$.
    Assume $f$ is continuous from $X$ to $\reals$ with respect to $T$ and $\standardTopologyR$.
    Assume $q$ is continuous from $X$ to $\reals$ with respect to $T$ and $\standardTopologyR$.
    Assume $a\in\reals$.
    Let $c = \{(x,a) \mid x\in X\}$.
    Let $s = \{(x,c(x) \rmul q(x)) \mid x\in X\}$.
    Let $g = \{(x,f(x) + s(x)) \mid x\in X\}$.
    Assume $x\in X$.
    Then $g(x) = f(x) + a \rmul q(x)$.
\end{lemma}
\begin{proof}
    $c$ is a function from $X$ to $\reals$ by \cref{constant_map_function}.
    $c(x) = a$ by \cref{constant_map_function,funs_is_function,function_apply_intro}.
    $s$ is continuous from $X$ to $\reals$ with respect to $T$ and $\standardTopologyR$
        by \cref{real_function_scalar_multiple_continuous_49}.
    $g$ is continuous from $X$ to $\reals$ with respect to $T$ and $\standardTopologyR$
        by \cref{real_function_add_scalar_multiple_continuous_49}.
    $q(x)\in\reals$ and $f(x)\in\reals$ and $s(x)\in\reals$ by \cref{continuous,funs_type_apply}.
    $s(x) = c(x) \rmul q(x)$ by \cref{continuous,funs_is_function,function_apply_intro}.
    Hence $s(x) = a \rmul q(x)$ by \cref{real_mul_eq_subst,real_lt_subst_left,real_lt_subst_right}.
    $g(x) = f(x) + s(x)$ by \cref{continuous,funs_is_function,function_apply_intro}.
    Thus $g(x) = f(x) + a \rmul q(x)$ by \cref{real_add_eq_subst,real_lt_subst_left,real_lt_subst_right}.
\end{proof}

% from §49 helper lemma 76: multiplying a weak inequality by a positive real
\begin{lemma}\label{real_leq_mul_pos_left_49}
    Assume $a\in\reals$.
    Assume $b\in\reals$.
    Assume $c\in\reals$.
    Suppose $0 < a$.
    Suppose $b \leq c$.
    Then $a \rmul b \leq a \rmul c$.
\end{lemma}
\begin{proof}
    Follows by \cref{real_leq_split,real_lt_mul_pos,real_mul_comm,real_lt_imp_leq}.
\end{proof}

% from §49 helper lemma 77: absolute value is invariant under negation before the density proof
\begin{lemma}\label{real_abs_neg_value_before_density_49}
    Assume $x\in\reals$.
    Then $\abs{\neg{x}} = \abs{x}$.
\end{lemma}
\begin{proof}
    Follows by \cref{real_add_ident,real_add_inverse,real_sub,real_sub_zero_local_49,real_abs_sub_swap}.
\end{proof}

% from §49 helper lemma 78: scalar perturbation pointwise distance bound
\begin{lemma}\label{real_function_add_scalar_multiple_pointwise_bound_49}
    Assume $T$ is a topology on $X$.
    Assume $f$ is continuous from $X$ to $\reals$ with respect to $T$ and $\standardTopologyR$.
    Assume $q$ is continuous from $X$ to $\reals$ with respect to $T$ and $\standardTopologyR$.
    Assume $a\in\reals$.
    Assume $b\in\reals$.
    Suppose $0 < a$.
    Let $c = \{(x,a) \mid x\in X\}$.
    Let $s = \{(x,c(x) \rmul q(x)) \mid x\in X\}$.
    Let $g = \{(x,f(x) + s(x)) \mid x\in X\}$.
    Assume $x\in X$.
    Suppose $0 \leq q(x)$.
    Suppose $q(x) \leq b$.
    Then $\apply{\standardMetricR}{(f(x),g(x))} \leq a \rmul b$.
\end{lemma}
\begin{proof}
    $q(x)\in\reals$ and $f(x)\in\reals$ and $a \rmul q(x)\in\reals$
        by \cref{continuous,funs,funs_type_apply,real_mul_closed}.
    $g(x) = f(x) + a \rmul q(x)$ by \cref{real_function_add_scalar_multiple_apply_49}.
    $g(x) - f(x) = a \rmul q(x)$
        by \cref{real_sub,real_add_assoc,real_add_inverse,real_add_ident,real_add_comm,real_lt_subst_left}.
    $\abs{a \rmul q(x)} = a \rmul q(x)$
        by \cref{real_abs_mul,real_abs_nonneg,real_lt_imp_leq,real_mul_eq_subst}.
    $\abs{(g(x) - f(x))} = a \rmul q(x)$ by \cref{real_lt_subst_left,real_lt_subst_right}.
    $a \rmul q(x) \leq a \rmul b$ by \cref{real_leq_mul_pos_left_49}.
    $g$ is continuous from $X$ to $\reals$ with respect to $T$ and $\standardTopologyR$
        by \cref{real_function_add_scalar_multiple_continuous_49}.
    $g(x)\in\reals$ by \cref{continuous,funs,funs_type_apply}.
    $\apply{\standardMetricR}{(g(x),f(x))} = \abs{(g(x) - f(x))}$ by \cref{standard_metric_value}.
    $\apply{\standardMetricR}{(f(x),g(x))} = \apply{\standardMetricR}{(g(x),f(x))}$
        by \cref{standard_metric_is_metric,metric_on,metric_symmetric,times_tuple_intro}.
    Therefore $\apply{\standardMetricR}{(f(x),g(x))} \leq a \rmul b$ by \cref{real_lt_subst_left}.
\end{proof}

% from §49 helper lemma 79: scalar perturbation uniform distance bound
\begin{lemma}\label{real_function_add_scalar_multiple_uniform_bound_49}
    Assume $T$ is a topology on $X$.
    Assume $\rho$ is the uniform metric on functions from $X$ to $\reals$ with respect to $\standardMetricR$.
    Assume $f$ is continuous from $X$ to $\reals$ with respect to $T$ and $\standardTopologyR$.
    Assume $q$ is continuous from $X$ to $\reals$ with respect to $T$ and $\standardTopologyR$.
    Assume $a\in\reals$.
    Assume $b\in\reals$.
    Assume $\delta\in\reals$.
    Suppose $0 < a$.
    Suppose $a \rmul b \leq \delta$.
    Suppose $\delta \leq 1$.
    Suppose for all $x\in X$ we have $0 \leq q(x)$ and $q(x) \leq b$.
    Let $c = \{(x,a) \mid x\in X\}$.
    Let $s = \{(x,c(x) \rmul q(x)) \mid x\in X\}$.
    Let $g = \{(x,f(x) + s(x)) \mid x\in X\}$.
    Then $\rho((f,g)) \leq \delta$.
\end{lemma}
\begin{proof}
    $g$ is continuous from $X$ to $\reals$ with respect to $T$ and $\standardTopologyR$
        by \cref{real_function_add_scalar_multiple_continuous_49}.
    $f\in\funs{X}{\reals}$ and $g\in\funs{X}{\reals}$ and $a \rmul b\in\reals$
        by \cref{continuous,funs,real_mul_closed}.
    Show that for all $x\in X$ we have $\apply{\standardMetricR}{(f(x),g(x))} \leq a \rmul b$.
    \begin{subproof}
        Follows by \cref{subseteq,real_function_add_scalar_multiple_pointwise_bound_49}.
    \end{subproof}
    Show that for all $x\in X$ we have $\apply{\standardMetricR}{(f(x),g(x))} \leq \delta$.
    \begin{subproof}
        Follows by \cref{standard_metric_is_metric,metric_on,times_tuple_intro,funs_type_apply,real_leq_trans}.
    \end{subproof}
    Therefore $\rho((f,g)) \leq \delta$ by \cref{uniform_metric_leq_from_pointwise_bound_49}.
\end{proof}

% from §49 helper lemma 80: equal-mesh distance values are nonnegative
\begin{lemma}\label{unit_interval_equal_mesh_distance_nonnegative_49}
    Assume $A$ is the equal mesh of order $m$ on the unit interval.
    Assume $q\in\funs{\intervalclosed{0}{1}}{\reals}$.
    Assume for all $x\in\intervalclosed{0}{1}$ we have
        $q(x)$ is the distance from $x$ to $A$ in $\reals$ with respect to $\standardMetricR$.
    Assume $x\in\intervalclosed{0}{1}$.
    Then $0 \leq q(x)$.
\end{lemma}
\begin{proof}
    We have $A\subseteq\intervalclosed{0}{1}$ and $A\neq\emptyset$ by \cref{unit_interval_equal_mesh_basic_49}.
    $A\subseteq\reals$ and $x\in\reals$ by \cref{intervalclosed,subseteq,subseteq_transitive}.
    Let $S = \pointDistanceSet{\standardMetricR}{A}{x}$.
    $q(x)$ is an infimum of $S$ by \cref{point_set_distance}.
    $S\subseteq\reals$ by \cref{point_distance_set,standard_metric_is_metric,metric_on,funs_type_apply,subseteq,times_tuple_intro}.
    $0$ is a lower bound of $S$
        by \cref{point_distance_set,subseteq,standard_metric_is_metric,metric_on,metric_nonnegative,real_lt_subst_right,real_add_ident,real_lower_bound}.
    Therefore $0 \leq q(x)$ by \cref{real_infimum}.
\end{proof}

% from §49 helper lemma 81: scalar mesh-distance perturbations are uniformly small
\begin{lemma}\label{unit_interval_equal_mesh_scalar_perturbation_uniform_bound_49}
    Let $I = \intervalclosed{0}{1}$.
    Let $T_I = \subspaceTopology{\standardTopologyR}{I}$.
    Assume $A$ is the equal mesh of order $m$ on the unit interval.
    Assume $\rho$ is the uniform metric on functions from $I$ to $\reals$ with respect to $\standardMetricR$.
    Assume $f$ is continuous from $I$ to $\reals$ with respect to $T_I$ and $\standardTopologyR$.
    Assume $q\in\funs{I}{\reals}$.
    Assume for all $x\in I$ we have
        $q(x)$ is the distance from $x$ to $A$ in $\reals$ with respect to $\standardMetricR$.
    Assume $q$ is continuous from $I$ to $\reals$ with respect to $T_I$ and $\standardTopologyR$.
    Assume $a\in\reals$.
    Assume $\delta\in\reals$.
    Suppose $0 < a$.
    Suppose $a \rmul (1 / m) \leq \delta$.
    Suppose $\delta \leq 1$.
    Let $c = \{(x,a) \mid x\in I\}$.
    Let $s = \{(x,c(x) \rmul q(x)) \mid x\in I\}$.
    Let $g = \{(x,f(x) + s(x)) \mid x\in I\}$.
    Then $\rho((f,g)) \leq \delta$.
\end{lemma}
\begin{proof}
    Follows by \cref{standard_topology_reals,standard_basis_is_basis,basis_topology_is_topology,intervalclosed,subseteq,subspace_topology_is_topology,unit_interval_equal_mesh_49,naturalsplus_reciprocal_real,unit_interval_equal_mesh_distance_nonnegative_49,unit_interval_equal_mesh_distance_bound_49,real_function_add_scalar_multiple_uniform_bound_49}.
\end{proof}

% from §49 helper lemma 82: continuous functions on the unit interval are uniformly continuous
\begin{lemma}\label{unit_interval_continuous_uniformly_continuous_49}
    Let $I = \intervalclosed{0}{1}$.
    Let $T_I = \subspaceTopology{\standardTopologyR}{I}$.
    Let $d_I = \restrl{\standardMetricR}{I\times I}$.
    Assume $f$ is continuous from $I$ to $\reals$ with respect to $T_I$ and $\standardTopologyR$.
    Then $f$ is uniformly continuous from $I$ to $\reals$ with respect to $d_I$ and $\standardMetricR$.
\end{lemma}
\begin{proof}
    Follows by \cref{unit_interval_restricted_standard_metric_topology_49,continuous_eq_domain,standard_metric_topology,unit_interval_restricted_standard_metric_49,standard_metric_is_metric,unit_interval_compact_restricted_standard_metric_49,uniform_continuity_theorem}.
\end{proof}

% from §49 helper lemma 83: adjacent left endpoint realizes equal-mesh distance
\begin{lemma}\label{unit_interval_equal_mesh_adjacent_left_realized_49}
    Assume $A$ is the equal mesh of order $m$ on the unit interval.
    Assume $q\in\funs{\intervalclosed{0}{1}}{\reals}$.
    Assume for all $x\in\intervalclosed{0}{1}$ we have
        $q(x)$ is the distance from $x$ to $A$ in $\reals$ with respect to $\standardMetricR$.
    Assume $p\in\naturalsPlus$.
    Assume $p + 1 \leq m$.
    Assume $x\in\intervalclosed{0}{1}$.
    Suppose $p / m \leq x$ and $x \leq (p + 1) / m$.
    Suppose $\apply{\standardMetricR}{(x,p / m)} \leq \apply{\standardMetricR}{(x,(p + 1) / m)}$.
    Then $q(x) = \apply{\standardMetricR}{(x,p / m)}$.
\end{lemma}
\begin{proof}
    Follows by \cref{unit_interval_equal_mesh_49,real_naturals_succ_bound_elim_local,initial_segment_naturalsplus,real_naturals_subset_reals,subseteq,real_naturals_closed_succ,real_div_naturalsplus_real,intervalclosed,unit_interval_equal_mesh_adjacent_separation_49,point_set_distance_left_endpoint_realized_49}.
\end{proof}

% from §49 helper lemma 84: adjacent right endpoint realizes equal-mesh distance
\begin{lemma}\label{unit_interval_equal_mesh_adjacent_right_realized_49}
    Assume $A$ is the equal mesh of order $m$ on the unit interval.
    Assume $q\in\funs{\intervalclosed{0}{1}}{\reals}$.
    Assume for all $x\in\intervalclosed{0}{1}$ we have
        $q(x)$ is the distance from $x$ to $A$ in $\reals$ with respect to $\standardMetricR$.
    Assume $p\in\naturalsPlus$.
    Assume $p + 1 \leq m$.
    Assume $x\in\intervalclosed{0}{1}$.
    Suppose $p / m \leq x$ and $x \leq (p + 1) / m$.
    Suppose $\apply{\standardMetricR}{(x,(p + 1) / m)} \leq \apply{\standardMetricR}{(x,p / m)}$.
    Then $q(x) = \apply{\standardMetricR}{(x,(p + 1) / m)}$.
\end{lemma}
\begin{proof}
    Follows by \cref{unit_interval_equal_mesh_49,real_naturals_succ_bound_elim_local,initial_segment_naturalsplus,real_naturals_subset_reals,subseteq,real_naturals_closed_succ,real_div_naturalsplus_real,intervalclosed,unit_interval_equal_mesh_adjacent_separation_49,point_set_distance_right_endpoint_realized_49}.
\end{proof}

% from §49 helper lemma 85: adjacent left equal-mesh distance formula
\begin{lemma}\label{unit_interval_equal_mesh_adjacent_left_formula_49}
    Assume $A$ is the equal mesh of order $m$ on the unit interval.
    Assume $q\in\funs{\intervalclosed{0}{1}}{\reals}$.
    Assume for all $x\in\intervalclosed{0}{1}$ we have
        $q(x)$ is the distance from $x$ to $A$ in $\reals$ with respect to $\standardMetricR$.
    Assume $p\in\naturalsPlus$.
    Assume $p + 1 \leq m$.
    Assume $x\in\intervalclosed{0}{1}$.
    Suppose $p / m \leq x$ and $x \leq (p + 1) / m$.
    Suppose $\apply{\standardMetricR}{(x,p / m)} \leq \apply{\standardMetricR}{(x,(p + 1) / m)}$.
    Then $q(x) = x - p / m$.
\end{lemma}
\begin{proof}
    Follows by \cref{unit_interval_equal_mesh_adjacent_left_realized_49,unit_interval_equal_mesh_49,real_naturals_subset_reals,subseteq,real_div_naturalsplus_real,intervalclosed,real_leq_add,real_sub,real_add_inverse,real_add_ident,real_lt_subst_left,real_lt_subst_right,standard_metric_value,real_abs_nonneg,real_add_closed}.
\end{proof}

% from §49 helper lemma 86: adjacent right equal-mesh distance formula
\begin{lemma}\label{unit_interval_equal_mesh_adjacent_right_formula_49}
    Assume $A$ is the equal mesh of order $m$ on the unit interval.
    Assume $q\in\funs{\intervalclosed{0}{1}}{\reals}$.
    Assume for all $x\in\intervalclosed{0}{1}$ we have
        $q(x)$ is the distance from $x$ to $A$ in $\reals$ with respect to $\standardMetricR$.
    Assume $p\in\naturalsPlus$.
    Assume $p + 1 \leq m$.
    Assume $x\in\intervalclosed{0}{1}$.
    Suppose $p / m \leq x$ and $x \leq (p + 1) / m$.
    Suppose $\apply{\standardMetricR}{(x,(p + 1) / m)} \leq \apply{\standardMetricR}{(x,p / m)}$.
    Then $q(x) = (p + 1) / m - x$.
\end{lemma}
\begin{proof}
    $q(x) = \apply{\standardMetricR}{(x,(p + 1) / m)}$
        by \cref{unit_interval_equal_mesh_adjacent_right_realized_49}.
    $m\in\naturalsPlus$ by \cref{unit_interval_equal_mesh_49}.
    $p\in\reals$ and $p + 1\in\reals$ and $m\in\reals$
        by \cref{real_naturals_subset_reals,subseteq,real_naturals_closed_succ}.
    $(p + 1) / m\in\reals$ by \cref{real_div_naturalsplus_real}.
    $x\in\reals$ by \cref{intervalclosed}.
    $0 \leq (p + 1) / m - x$
        by \cref{real_leq_add,real_sub,real_add_inverse,real_add_ident,real_lt_subst_left,real_lt_subst_right}.
    $\apply{\standardMetricR}{(x,(p + 1) / m)} = \abs{(x - (p + 1) / m)}$ by \cref{standard_metric_value}.
    $\abs{(x - (p + 1) / m)} = \abs{((p + 1) / m - x)}$ by \cref{real_abs_sub_swap}.
    $\abs{((p + 1) / m - x)} = (p + 1) / m - x$
        by \cref{real_abs_nonneg,real_sub,real_add_inverse,real_add_closed}.
    Therefore $q(x) = (p + 1) / m - x$ by \cref{real_lt_subst_left,real_lt_subst_right}.
\end{proof}

% from §49 helper lemma 87: forward linear secant quotient equals one
\begin{lemma}\label{forward_linear_secant_quotient_one_49}
    Assume $a\in\reals$.
    Assume $x\in\reals$.
    Assume $h\in\reals$.
    Suppose $0 < h$.
    Let $y = x + h$.
    Assume $q_x = x - a$.
    Assume $q_y = y - a$.
    Then $\abs{((q_y - q_x) / (y - x))} = 1$.
\end{lemma}
\begin{proof}
    Follows by \cref{real_add_closed,real_sub,real_add_assoc,real_add_inverse,real_add_ident,real_add_comm,real_sub_swap_neg,real_sub_add_chain,real_pos_nonzero,real_div,real_mul_inverse,real_mul_comm,real_abs_nonneg,real_zero_lt_one,real_lt_imp_leq,real_one_in_naturals,real_naturals_subset_reals,subseteq}.
\end{proof}

% from §49 helper lemma 88: backward linear secant quotient has absolute value one
\begin{lemma}\label{backward_linear_secant_quotient_one_49}
    Assume $a\in\reals$.
    Assume $x\in\reals$.
    Assume $h\in\reals$.
    Suppose $0 < h$.
    Let $y = x - h$.
    Assume $q_x = x - a$.
    Assume $q_y = y - a$.
    Then $\abs{((q_y - q_x) / (y - x))} = 1$.
\end{lemma}
\begin{proof}
    Follows by \cref{real_sub,real_add_inverse,real_add_closed,real_add_assoc,real_add_ident,real_add_comm,real_sub_swap_neg,real_sub_add_chain,real_add_cancel,real_pos_nonzero,real_neg_zero,real_div,real_mul_inverse,real_mul_comm,real_abs_nonneg,real_zero_lt_one,real_lt_imp_leq,real_one_in_naturals,real_naturals_subset_reals,subseteq}.
\end{proof}

% from §49 helper lemma 89: forward descending linear secant quotient has absolute value one
\begin{lemma}\label{forward_descending_linear_secant_quotient_one_49}
    Assume $a\in\reals$.
    Assume $x\in\reals$.
    Assume $h\in\reals$.
    Suppose $0 < h$.
    Let $y = x + h$.
    Assume $q_x = a - x$.
    Assume $q_y = a - y$.
    Then $\abs{((q_y - q_x) / (y - x))} = 1$.
\end{lemma}
\begin{proof}
    $y\in\reals$ by \cref{real_add_closed}.
    $y - x = h$ by \cref{real_sub,real_add_assoc,real_add_inverse,real_add_ident,real_add_comm}.
    $q_y - q_x = \neg{h}$
        by \cref{real_sub,real_lt_subst_left,real_lt_subst_right,real_sub_swap_neg,real_add_eq_subst,real_add_inverse,real_add_closed,real_add_comm,real_sub_add_chain,real_add_assoc,real_add_ident}.
    $h\neq 0$ by \cref{real_pos_nonzero}.
    $(q_y - q_x) / (y - x) = \neg{1}$
        by \cref{real_lt_subst_left,real_lt_subst_right,real_mul_inverse,real_add_inverse,real_div,real_mul_comm,real_mul_neg_right}.
    Thus $\abs{((q_y - q_x) / (y - x))} = 1$
        by \cref{real_abs_neg_value_before_density_49,real_abs_nonneg,real_zero_lt_one,real_lt_imp_leq,real_one_in_naturals,real_naturals_subset_reals,subseteq}.
\end{proof}

% from §49 helper lemma 90: inverse of a negative real before the density proof
\begin{lemma}\label{real_inv_neg_value_before_density_49}
    Assume $h\in\reals$.
    Suppose $h\neq 0$.
    Then $\inv{\neg{h}} = \neg{\inv{h}}$.
\end{lemma}
\begin{proof}
    Follows by \cref{real_add_inverse,real_neg_neg,real_neg_zero,real_mul_inverse,real_mul_neg_right,real_mul_comm,real_mul_inverse_unique}.
\end{proof}

% from §49 helper lemma 91: backward descending linear secant quotient has absolute value one
\begin{lemma}\label{backward_descending_linear_secant_quotient_one_49}
    Assume $a\in\reals$.
    Assume $x\in\reals$.
    Assume $h\in\reals$.
    Suppose $0 < h$.
    Let $y = x - h$.
    Assume $q_x = a - x$.
    Assume $q_y = a - y$.
    Then $\abs{((q_y - q_x) / (y - x))} = 1$.
\end{lemma}
\begin{proof}
    Follows by \cref{real_sub,real_add_inverse,real_add_closed,real_add_assoc,real_add_ident,real_add_comm,real_lt_subst_left,real_lt_subst_right,real_sub_swap_neg,real_add_eq_subst,real_sub_add_chain,real_pos_imp_neg,real_neg_nonzero,real_pos_nonzero,real_neg_zero,real_div,real_mul_inverse,real_inv_neg_value_before_density_49,real_mul_eq_subst,real_mul_neg_right,real_mul_comm,real_abs_neg_value_before_density_49,real_abs_nonneg,real_zero_lt_one,real_lt_imp_leq,real_one_in_naturals,real_naturals_subset_reals,subseteq}.
\end{proof}

% from §49 helper lemma 92: a large scalar sawtooth quotient dominates a small quotient
\begin{lemma}\label{large_scalar_sawtooth_quotient_dominates_49}
    Assume $u\in\reals$.
    Assume $v\in\reals$.
    Assume $a\in\reals$.
    Assume $n\in\reals$.
    Suppose $0 < n$.
    Suppose $\abs{u} < 1 / (1 + 1)$.
    Suppose $\abs{v} = 1$.
    Suppose $n + 1 < a$.
    Then $n < \abs{(u + a \rmul v)}$.
\end{lemma}
\begin{proof}
    $1\in\reals$ and $0 < 1$ and $0\in\reals$
        by \cref{real_one_in_naturals,real_naturals_subset_reals,subseteq,real_zero_lt_one,real_add_ident}.
    $1 / (1 + 1)\in\reals$ and $0 < 1 / (1 + 1)$ by \cref{real_half_of_positive}.
    $a \rmul v\in\reals$ and $u + a \rmul v\in\reals$ by \cref{real_mul_closed,real_add_closed}.
    $1 < n + 1$ and $n + 1\in\reals$ and $0 < n + 1$
        by \cref{real_lt_add,real_add_ident,real_lt_subst_left,real_add_closed,real_zero_lt_one,real_lt_trans}.
    $0 < a$ by \cref{real_lt_trans}.
    $\abs{a} = a$ by \cref{real_abs_nonneg,real_lt_imp_leq}.
    $\abs{a \rmul v} = a$ by \cref{real_abs_mul,real_mul_eq_subst,real_mul_comm,real_mul_ident}.
    $\abs{((a \rmul v) - (u + a \rmul v))} < 1 / (1 + 1)$
        by \cref{real_sub,real_add_assoc,real_add_inverse,real_add_ident,real_add_comm,real_abs_sub_swap,real_lt_subst_left}.
    Let $r = n + 1$.
    $r\in\reals$ by \cref{real_add_closed}.
    $1 / (1 + 1) < 1$ by \cref{real_half_plus_half,real_lt_add,real_add_ident,real_lt_subst_left}.
    $n + (1 / (1 + 1)) < r$ by \cref{real_lt_add,real_add_comm,real_lt_subst_left,real_lt_subst_right}.
    $r < a$ by \cref{real_lt_subst_left}.
    $r < \abs{a \rmul v}$ by \cref{real_lt_subst_right}.
    Therefore $n < \abs{(u + a \rmul v)}$ by \cref{real_abs_perturbation_lower_bound_49}.
\end{proof}

% from §49 helper lemma 93: a large scalar quotient dominates a bounded quotient
\begin{lemma}\label{large_scalar_quotient_dominates_bound_49}
    Assume $u\in\reals$.
    Assume $v\in\reals$.
    Assume $a\in\reals$.
    Assume $n\in\reals$.
    Assume $s\in\reals$.
    Assume $r\in\reals$.
    Suppose $0 < n$.
    Suppose $0 < r$.
    Suppose $\abs{u} < s$.
    Suppose $\abs{v} = 1$.
    Suppose $n + s < r$.
    Suppose $r < a$.
    Then $n < \abs{(u + a \rmul v)}$.
\end{lemma}
\begin{proof}
    $1\in\reals$ and $0\in\reals$ and $a \rmul v\in\reals$ and $u + a \rmul v\in\reals$
        by \cref{real_one_in_naturals,real_naturals_subset_reals,subseteq,real_add_ident,real_mul_closed,real_add_closed}.
    $0 < a$ by \cref{real_lt_trans}.
    $\abs{a} = a$ by \cref{real_abs_nonneg,real_lt_imp_leq}.
    $\abs{a \rmul v} = a$ by \cref{real_abs_mul,real_mul_eq_subst,real_mul_comm,real_mul_ident}.
    $\abs{((a \rmul v) - (u + a \rmul v))} < s$
        by \cref{real_sub,real_add_assoc,real_add_inverse,real_add_ident,real_add_comm,real_abs_sub_swap,real_lt_subst_left}.
    $r < \abs{a \rmul v}$ by \cref{real_lt_subst_right}.
    Therefore $n < \abs{(u + a \rmul v)}$ by \cref{real_abs_perturbation_lower_bound_49}.
\end{proof}

% from §49 helper lemma 94: secant quotients distribute over scalar perturbations
\begin{lemma}\label{real_secant_add_scalar_multiple_quotient_49}
    Assume $f_y\in\reals$.
    Assume $f_x\in\reals$.
    Assume $q_y\in\reals$.
    Assume $q_x\in\reals$.
    Assume $a\in\reals$.
    Assume $d\in\reals$.
    Suppose $d\neq 0$.
    Then $(((f_y + a \rmul q_y) - (f_x + a \rmul q_x)) / d)
        = ((f_y - f_x) / d) + a \rmul ((q_y - q_x) / d)$.
\end{lemma}
\begin{proof}
    $f_y - f_x\in\reals$ and $q_y - q_x\in\reals$
        by \cref{real_sub,real_add_inverse,real_add_closed}.
    $a \rmul q_y\in\reals$ and $a \rmul q_x\in\reals$ by \cref{real_mul_closed}.
    $f_y + a \rmul q_y\in\reals$ and $f_x + a \rmul q_x\in\reals$ by \cref{real_add_closed}.
    $f_y + \neg{f_x}\in\reals$ and $a \rmul q_y + \neg{(a \rmul q_x)}\in\reals$
        by \cref{real_add_inverse,real_add_closed}.
    $\neg{(f_x + a \rmul q_x)} = \neg{f_x} + \neg{(a \rmul q_x)}$ by \cref{real_neg_addition}.
    $(f_y + a \rmul q_y) - (f_x + a \rmul q_x)
        = (f_y + a \rmul q_y) + \neg{(f_x + a \rmul q_x)}$ by \cref{real_sub}.
    Hence $(f_y + a \rmul q_y) - (f_x + a \rmul q_x)
        = (f_y + a \rmul q_y) + (\neg{f_x} + \neg{(a \rmul q_x)})$ by \cref{real_add_eq_subst}.
    $\neg{f_x}\in\reals$ and $\neg{(a \rmul q_x)}\in\reals$ by \cref{real_add_inverse}.
    $((f_y + a \rmul q_y) + \neg{f_x}) + \neg{(a \rmul q_x)}
        = (f_y + a \rmul q_y) + (\neg{f_x} + \neg{(a \rmul q_x)})$ by \cref{real_add_assoc}.
    $(f_y + a \rmul q_y) + \neg{f_x} = f_y + (a \rmul q_y + \neg{f_x})$ by \cref{real_add_assoc}.
    $a \rmul q_y + \neg{f_x} = \neg{f_x} + a \rmul q_y$ by \cref{real_add_comm}.
    $f_y + (a \rmul q_y + \neg{f_x}) = f_y + (\neg{f_x} + a \rmul q_y)$ by \cref{real_add_eq_subst}.
    $f_y + (\neg{f_x} + a \rmul q_y) = (f_y + \neg{f_x}) + a \rmul q_y$ by \cref{real_add_assoc}.
    Hence $(f_y + a \rmul q_y) + \neg{f_x} = (f_y + \neg{f_x}) + a \rmul q_y$ by \cref{real_lt_subst_left,real_lt_subst_right}.
    $((f_y + a \rmul q_y) + \neg{f_x}) + \neg{(a \rmul q_x)}
        = ((f_y + \neg{f_x}) + a \rmul q_y) + \neg{(a \rmul q_x)}$ by \cref{real_add_eq_subst}.
    $((f_y + \neg{f_x}) + a \rmul q_y) + \neg{(a \rmul q_x)}
        = (f_y + \neg{f_x}) + (a \rmul q_y + \neg{(a \rmul q_x)})$ by \cref{real_add_assoc}.
    Hence $(f_y + a \rmul q_y) + (\neg{f_x} + \neg{(a \rmul q_x)})
        = (f_y + \neg{f_x}) + (a \rmul q_y + \neg{(a \rmul q_x)})$
        by \cref{real_lt_subst_left,real_lt_subst_right}.
    $f_y + \neg{f_x} = f_y - f_x$ by \cref{real_sub}.
    $a \rmul q_y + \neg{(a \rmul q_x)} = a \rmul q_y + a \rmul \neg{q_x}$ by \cref{real_mul_neg_right,real_add_eq_subst}.
    $a \rmul q_y = q_y \rmul a$ and $a \rmul \neg{q_x} = \neg{q_x} \rmul a$ by \cref{real_mul_comm,real_add_inverse}.
    $a \rmul q_y + a \rmul \neg{q_x} = q_y \rmul a + \neg{q_x} \rmul a$ by \cref{real_add_eq_subst}.
    $(q_y + \neg{q_x}) \rmul a = q_y \rmul a + \neg{q_x} \rmul a$ by \cref{real_right_distrib,real_add_inverse}.
    $(q_y + \neg{q_x}) \rmul a = (q_y - q_x) \rmul a$ by \cref{real_sub,real_mul_eq_subst}.
    $(q_y - q_x) \rmul a = a \rmul (q_y - q_x)$ by \cref{real_mul_comm}.
    $a \rmul q_y + \neg{(a \rmul q_x)} = a \rmul (q_y - q_x)$ by \cref{real_lt_subst_left,real_lt_subst_right}.
    $(f_y + a \rmul q_y) - (f_x + a \rmul q_x) = (f_y - f_x) + a \rmul (q_y - q_x)$
        by \cref{real_lt_subst_left,real_lt_subst_right,real_add_eq_subst}.
    $((f_y + a \rmul q_y) - (f_x + a \rmul q_x)) / d
        = ((f_y - f_x) + a \rmul (q_y - q_x)) / d$
        by \cref{real_lt_subst_left}.
    $\inv{d}\in\reals$ by \cref{real_mul_inverse}.
    $a \rmul (q_y - q_x)\in\reals$ and $(f_y - f_x) + a \rmul (q_y - q_x)\in\reals$
        by \cref{real_mul_closed,real_add_closed}.
    $((f_y - f_x) + a \rmul (q_y - q_x)) / d
        = ((f_y - f_x) \rmul \inv{d}) + ((a \rmul (q_y - q_x)) \rmul \inv{d})$
        by \cref{real_div,real_right_distrib}.
    $((a \rmul (q_y - q_x)) \rmul \inv{d}) = a \rmul ((q_y - q_x) \rmul \inv{d})$
        by \cref{real_mul_assoc}.
    $(f_y - f_x) / d = (f_y - f_x) \rmul \inv{d}$ by \cref{real_div}.
    $a \rmul ((q_y - q_x) / d) = a \rmul ((q_y - q_x) \rmul \inv{d})$ by \cref{real_div,real_mul_eq_subst}.
    Therefore $(((f_y + a \rmul q_y) - (f_x + a \rmul q_x)) / d)
        = ((f_y - f_x) / d) + a \rmul ((q_y - q_x) / d)$
        by \cref{real_lt_subst_left,real_lt_subst_right,real_add_eq_subst}.
\end{proof}

% from §49 helper lemma 95: scalar perturbations inherit large secants from a unit-slope perturbation
\begin{lemma}\label{scalar_perturbation_large_secants_from_unit_quotients_49}
    Assume $I\subseteq\reals$.
    Assume $T$ is a topology on $I$.
    Assume $f$ is continuous from $I$ to $\reals$ with respect to $T$ and $\standardTopologyR$.
    Assume $q$ is continuous from $I$ to $\reals$ with respect to $T$ and $\standardTopologyR$.
    Assume $a\in\reals$.
    Assume $n\in\reals$.
    Assume $b\in\reals$.
    Assume $r\in\reals$.
    Assume $h\in\reals$.
    Suppose $0 < n$.
    Suppose $0 < h$.
    Suppose $0 < r$.
    Suppose $n + b < r$.
    Suppose $r < a$.
    Suppose for all $x\in I$ there exists $y\in I$ such that
        either $y = x + h$ and
            $\abs{((\apply{f}{y} - \apply{f}{x}) / (y - x))} < b$ and
            $\abs{((\apply{q}{y} - \apply{q}{x}) / (y - x))} = 1$
        or $y = x - h$ and
            $\abs{((\apply{f}{y} - \apply{f}{x}) / (y - x))} < b$ and
            $\abs{((\apply{q}{y} - \apply{q}{x}) / (y - x))} = 1$.
    Let $c = \{(x,a) \mid x\in I\}$.
    Let $w = \{(x,c(x) \rmul q(x)) \mid x\in I\}$.
    Let $g = \{(x,f(x) + w(x)) \mid x\in I\}$.
    Then $g$ has secants above $n$ at scale $h$ on $I$.
\end{lemma}
\begin{proof}
    $f\in\funs{I}{\reals}$ and $q\in\funs{I}{\reals}$ by \cref{continuous,funs}.
    $g$ is continuous from $I$ to $\reals$ with respect to $T$ and $\standardTopologyR$
        by \cref{real_function_add_scalar_multiple_continuous_49}.
    $g\in\funs{I}{\reals}$ by \cref{continuous,funs}.
    Show that for all $x\in I$ there exists $y\in I$ such that
        either $y = x + h$ and
            $n < \abs{((\apply{g}{y} - \apply{g}{x}) / (y - x))}$
        or $y = x - h$ and
            $n < \abs{((\apply{g}{y} - \apply{g}{x}) / (y - x))}$.
    \begin{subproof}
        Fix $x\in I$.
        Take $y\in I$ such that
            either $y = x + h$ and
                $\abs{((\apply{f}{y} - \apply{f}{x}) / (y - x))} < b$ and
                $\abs{((\apply{q}{y} - \apply{q}{x}) / (y - x))} = 1$
            or $y = x - h$ and
                $\abs{((\apply{f}{y} - \apply{f}{x}) / (y - x))} < b$ and
                $\abs{((\apply{q}{y} - \apply{q}{x}) / (y - x))} = 1$ by \cref{subseteq}.
        $\apply{f}{y}\in\reals$ and $\apply{f}{x}\in\reals$ and
            $\apply{q}{y}\in\reals$ and $\apply{q}{x}\in\reals$ and
            $\apply{g}{y}\in\reals$ and $\apply{g}{x}\in\reals$ by \cref{funs_type_apply}.
        $x\in\reals$ and $y\in\reals$ by \cref{subseteq}.
        $\apply{g}{y} = \apply{f}{y} + a \rmul \apply{q}{y}$ and
            $\apply{g}{x} = \apply{f}{x} + a \rmul \apply{q}{x}$
            by \cref{real_function_add_scalar_multiple_apply_49}.
        Let $d = y - x$.
        $d\in\reals$ by \cref{real_sub,real_add_inverse,real_add_closed}.
        $d = h$ or $d = \neg{h}$
            by \cref{real_sub,real_add_assoc,real_add_inverse,real_add_ident,real_add_comm,real_lt_subst_left}.
        $d\neq 0$ by \cref{real_pos_nonzero,real_pos_imp_neg,real_neg_nonzero,real_lt_subst_left}.
        Let $u = (\apply{f}{y} - \apply{f}{x}) / (y - x)$.
        Let $v = (\apply{q}{y} - \apply{q}{x}) / (y - x)$.
        $u\in\reals$ and $v\in\reals$
            by \cref{real_sub,real_add_inverse,real_add_closed,real_div,real_mul_inverse,real_mul_closed}.
        $((\apply{g}{y} - \apply{g}{x}) / (y - x)) = u + a \rmul v$
            by \cref{real_secant_add_scalar_multiple_quotient_49,real_lt_subst_left,real_lt_subst_right}.
        $\abs{u} < b$ and $\abs{v} = 1$ by \cref{real_lt_subst_left}.
        $n < \abs{(u + a \rmul v)}$ by \cref{large_scalar_quotient_dominates_bound_49}.
        $n < \abs{((\apply{g}{y} - \apply{g}{x}) / (y - x))}$ by \cref{real_lt_subst_right}.
        Show that there exists $z\in I$ such that
            either $z = x + h$ and
                $n < \abs{((\apply{g}{z} - \apply{g}{x}) / (z - x))}$
            or $z = x - h$ and
                $n < \abs{((\apply{g}{z} - \apply{g}{x}) / (z - x))}$.
        \begin{subproof}
            Trivial.
        \end{subproof}
    \end{subproof}
    Therefore $g$ has secants above $n$ at scale $h$ on $I$ by \cref{large_fixed_scale_secants_on_real_subset}.
\end{proof}

% from §49 helper lemma 96: scalar perturbation data give a robust-class point in a ball
\begin{lemma}\label{robust_class_ball_witness_from_scalar_perturbation_49}
    Let $I = \intervalclosed{0}{1}$.
    Let $T_I = \subspaceTopology{\standardTopologyR}{I}$.
    Assume $C$ is the set of continuous functions from $I$ to $\reals$ with respect to $T_I$ and $\standardTopologyR$.
    Assume $\rho$ is the uniform metric on functions from $I$ to $\reals$ with respect to $\standardMetricR$.
    Let $e = \restrl{\rho}{C\times C}$.
    Assume $n\in\naturalsPlus$.
    Assume $V$ is the robust secant class above $n$ on $I$ with respect to $T_I$ and $\rho$.
    Assume $f\in C$.
    Assume $\epsilon\in\reals$.
    Assume $0 < \epsilon$.
    Assume $q$ is continuous from $I$ to $\reals$ with respect to $T_I$ and $\standardTopologyR$.
    Assume $a\in\reals$.
    Assume $b\in\reals$.
    Assume $r\in\reals$.
    Assume $h\in\reals$.
    Suppose $0 < h$.
    Suppose $0 < r$.
    Suppose $h \leq 1 / n$.
    Suppose $(n + n) + b < r$.
    Suppose $r < a$.
    Suppose for all $x\in I$ there exists $y\in I$ such that
        either $y = x + h$ and
            $\abs{((\apply{f}{y} - \apply{f}{x}) / (y - x))} < b$ and
            $\abs{((\apply{q}{y} - \apply{q}{x}) / (y - x))} = 1$
        or $y = x - h$ and
            $\abs{((\apply{f}{y} - \apply{f}{x}) / (y - x))} < b$ and
            $\abs{((\apply{q}{y} - \apply{q}{x}) / (y - x))} = 1$.
    Let $c = \{(x,a) \mid x\in I\}$.
    Let $w = \{(x,c(x) \rmul q(x)) \mid x\in I\}$.
    Let $g = \{(x,f(x) + w(x)) \mid x\in I\}$.
    Suppose $\rho((f,g)) < \epsilon$.
    Then $\metricBall{e}{C}{f}{\epsilon}$ intersects $V$.
\end{lemma}
\begin{proof}
    $I\subseteq\reals$ by \cref{intervalclosed,subseteq}.
    $f$ is continuous from $I$ to $\reals$ with respect to $T_I$ and $\standardTopologyR$
        by \cref{continuous_function_set_elim_49}.
    $T_I$ is a topology on $I$ by \cref{subspace_topology_is_topology,standard_topology_reals,standard_basis_is_basis,basis_topology_is_topology}.
    $n\in\reals$ and $0 < n$ by \cref{real_naturals_subset_reals,subseteq,naturalsplus_positive_real}.
    $n + n\in\reals$ and $0 < n + n$ by \cref{real_add_closed,real_pos_add}.
    $g$ is continuous from $I$ to $\reals$ with respect to $T_I$ and $\standardTopologyR$ and
        $g$ has secants above $n + n$ at scale $h$ on $I$
        by \cref{real_function_add_scalar_multiple_continuous_49,scalar_perturbation_large_secants_from_unit_quotients_49}.
    $g\in V$ by \cref{unit_interval_robust_class_self_witness_49}.
    $V\subseteq C$ by \cref{unit_interval_robust_large_secant_class_subset_continuous_49}.
    $g\in C$ by \cref{subseteq}.
    $e((f,g)) < \epsilon$ by \cref{restricted_uniform_metric_apply_49,real_lt_subst_left}.
    $e$ is a metric on $C$ by \cref{unit_interval_continuous_function_space_metric_49}.
    Therefore $\metricBall{e}{C}{f}{\epsilon}$ intersects $V$ by \cref{metric_ball_intersection_witness_49}.
\end{proof}

% from §49 helper lemma 97: continuous real functions on the unit interval have bounded image
\begin{lemma}\label{unit_interval_continuous_image_bounded_49}
    Let $I = \intervalclosed{0}{1}$.
    Let $T_I = \subspaceTopology{\standardTopologyR}{I}$.
    Suppose $f$ is continuous from $I$ to $\reals$ with respect to $T_I$ and $\standardTopologyR$.
    Then $\img{f}{I}$ is bounded in $\reals$ with respect to $\standardMetricR$.
\end{lemma}
\begin{proof}
    Let $A = \img{f}{I}$.
    $I$ is compact with respect to $T_I$
        by \cref{unit_interval_compact_restricted_standard_metric_49,unit_interval_restricted_standard_metric_topology_49,compact_eq_topology}.
    $A$ is compact with respect to $\subspaceTopology{\standardTopologyR}{A}$ by \cref{continuous_image_compact}.
    $A$ is totally bounded in $\reals$ with respect to $\standardMetricR$
        by \cref{continuous,funs,funs_ran,img_subseteq_ran,subseteq_transitive,standard_metric_is_metric,standard_metric_topology,compact_eq_topology,compact_metric_subspace_ambient_totally_bounded}.
    Therefore $A$ is bounded in $\reals$ with respect to $\standardMetricR$
        by \cref{continuous,funs,funs_ran,img_subseteq_ran,subseteq_transitive,standard_metric_is_metric,totally_bounded_subset_bounded}.
\end{proof}

% from §49 helper lemma 98: bounded image gives a uniform bound on value distances
\begin{lemma}\label{unit_interval_continuous_value_distance_bound_49}
    Let $I = \intervalclosed{0}{1}$.
    Let $T_I = \subspaceTopology{\standardTopologyR}{I}$.
    Suppose $f$ is continuous from $I$ to $\reals$ with respect to $T_I$ and $\standardTopologyR$.
    Then there exists $M\in\reals$ such that for all $x,y\in I$ we have
        $\apply{\standardMetricR}{(\apply{f}{x},\apply{f}{y})} \leq M$.
\end{lemma}
\begin{proof}
    $\img{f}{I}$ is bounded in $\reals$ with respect to $\standardMetricR$
        by \cref{unit_interval_continuous_image_bounded_49}.
    Take $M\in\reals$ such that for all $a_1,a_2\in\img{f}{I}$ we have
        $\apply{\standardMetricR}{(a_1,a_2)} \leq M$ by \cref{metric_bounded}.
    Thus there exists $B\in\reals$ such that for all $x,y\in I$ we have
        $\apply{\standardMetricR}{(\apply{f}{x},\apply{f}{y})} \leq B$
        by \cref{continuous,funs,funs_is_function,function_apply_iff,img_iff}.
\end{proof}

% from §49 helper lemma 99: numerator bounds give quotient bounds by a positive scale before density
\begin{lemma}\label{real_abs_div_lt_from_mul_bound_before_density_49}
    Assume $a\in\reals$.
    Assume $h\in\reals$.
    Assume $n\in\reals$.
    Suppose $0 < h$.
    Suppose $0 < n$.
    Suppose $\abs{a} < h \rmul n$.
    Then $\abs{a / h} < n$.
\end{lemma}
\begin{proof}
    $h\neq 0$ and $\inv{h}\in\reals$ and $0 < \inv{h}$ by \cref{real_pos_nonzero,real_mul_inverse,real_inv_pos}.
    Therefore $\abs{a / h} < n$
        by \cref{real_div,real_abs_mul,real_abs_nonneg,real_lt_imp_leq,real_mul_eq_subst,real_mul_closed,real_abs_in_reals,real_lt_mul_pos,real_mul_assoc,real_mul_comm,real_mul_inverse,real_mul_ident,real_lt_subst_right,real_lt_subst_left}.
\end{proof}

% from §49 helper lemma 100: numerator bounds give quotient bounds by a negative scale before density
\begin{lemma}\label{real_abs_div_neg_lt_from_mul_bound_before_density_49}
    Assume $a\in\reals$.
    Assume $h\in\reals$.
    Assume $n\in\reals$.
    Suppose $0 < h$.
    Suppose $0 < n$.
    Suppose $\abs{a} < h \rmul n$.
    Then $\abs{a / (\neg{h})} < n$.
\end{lemma}
\begin{proof}
    Therefore $\abs{a / (\neg{h})} < n$
        by \cref{real_pos_nonzero,real_add_inverse,real_mul_inverse,real_inv_neg_value_before_density_49,real_div,real_mul_eq_subst,real_mul_neg_right,real_lt_subst_left,real_lt_subst_right,real_abs_neg_value_before_density_49,real_mul_closed,real_abs_div_lt_from_mul_bound_before_density_49}.
\end{proof}

% from §49 helper lemma 103: strict value bounds give bounded fixed-scale secant quotients
\begin{lemma}\label{unit_interval_strict_fixed_scale_quotient_bound_49}
    Let $I = \intervalclosed{0}{1}$.
    Let $T_I = \subspaceTopology{\standardTopologyR}{I}$.
    Suppose $f$ is continuous from $I$ to $\reals$ with respect to $T_I$ and $\standardTopologyR$.
    Assume $h\in\reals$.
    Suppose $0 < h$.
    Assume $\eta\in\reals$.
    Suppose $0 < \eta$.
    Suppose for all $x,y\in I$ if $y = x + h$ or $y = x - h$ then
        $\apply{\standardMetricR}{(\apply{f}{y},\apply{f}{x})} < \eta$.
    Let $b = \eta / h$.
    Then $b\in\reals$ and $0 < b$ and
        for all $x,y\in I$ if $y = x + h$ or $y = x - h$ then
            $\abs{((\apply{f}{y} - \apply{f}{x}) / (y - x))} < b$.
\end{lemma}
\begin{proof}
    $h\neq 0$ and $b\in\reals$ and $\inv{h}\in\reals$ and $0 < \inv{h}$ and $0 < b$
        by \cref{real_pos_nonzero,real_div,real_mul_inverse,real_mul_closed,real_inv_pos,real_add_ident,real_lt_mul_pos,real_mul_zero,real_lt_subst_left,real_lt_subst_right}.
    For all $x,y\in I$ we have $\apply{f}{x}\in\reals$ and $\apply{f}{y}\in\reals$ and
        $\apply{f}{y} - \apply{f}{x}\in\reals$
        by \cref{continuous,funs,funs_type_apply,real_sub,real_add_inverse,real_add_closed}.
    For all $x,y\in I$ if $y = x + h$ or $y = x - h$ then
        $\abs{(\apply{f}{y} - \apply{f}{x})} < \eta$
        by \cref{intervalclosed,subseteq,standard_metric_value,real_lt_subst_left}.
    For all $x,y\in I$ if $y = x + h$ or $y = x - h$ then
        $\abs{(\apply{f}{y} - \apply{f}{x})} < h \rmul b$
        by \cref{real_div,real_mul_assoc,real_mul_inverse,real_mul_comm,real_mul_ident,real_pos_nonzero,real_lt_subst_right}.
    For all $x,y\in I$ if $y = x + h$ or $y = x - h$ then $y - x = h$ or $y - x = \neg{h}$
        by \cref{intervalclosed,real_sub,real_add_assoc,real_add_inverse,real_add_ident,real_add_comm}.
    Show that $b\in\reals$ and $0 < b$ and
        for all $x,y\in I$ if $y = x + h$ or $y = x - h$ then
            $\abs{((\apply{f}{y} - \apply{f}{x}) / (y - x))} < b$.
    \begin{subproof}
        Follows by \cref{real_abs_div_lt_from_mul_bound_before_density_49,real_abs_div_neg_lt_from_mul_bound_before_density_49,real_lt_subst_left,real_lt_subst_right}.
    \end{subproof}
\end{proof}

% from §49 helper lemma 101: continuous functions have bounded fixed-scale secant quotients
\begin{lemma}\label{unit_interval_continuous_fixed_scale_quotient_bound_49}
    Let $I = \intervalclosed{0}{1}$.
    Let $T_I = \subspaceTopology{\standardTopologyR}{I}$.
    Suppose $f$ is continuous from $I$ to $\reals$ with respect to $T_I$ and $\standardTopologyR$.
    Assume $h\in\reals$.
    Suppose $0 < h$.
    Assume $M\in\reals$.
    Suppose for all $x,y\in I$ we have
        $\apply{\standardMetricR}{(\apply{f}{x},\apply{f}{y})} \leq M$.
    Let $B = M + 1$.
    Let $b = B / h$.
    Then $b\in\reals$ and $0 < b$ and
    for all $x,y\in I$ if $y = x + h$ or $y = x - h$ then
            $\abs{((\apply{f}{y} - \apply{f}{x}) / (y - x))} < b$.
\end{lemma}
\begin{proof}
    $0 \leq M$
        by \cref{intervalclosed,real_add_ident,real_zero_lt_one,real_lt_imp_leq,subseteq,continuous,funs,funs_type_apply,standard_metric_is_metric,metric_on,metric_zero_characterization,real_lt_subst_left}.
    $1\in\reals$ and $0 < 1$ by \cref{real_one_in_naturals,real_naturals_subset_reals,subseteq,real_zero_lt_one}.
    $B\in\reals$ and $0 < B$ by \cref{real_add_closed,real_leq_add,real_add_ident,real_zero_lt_one,real_leq_split,real_lt_trans,real_lt_subst_right}.
    For all $x,y\in I$ if $y = x + h$ or $y = x - h$ then
        $\apply{\standardMetricR}{(\apply{f}{y},\apply{f}{x})} < B$
        by \cref{real_lt_add,real_add_ident,real_add_comm,real_lt_subst_left,real_lt_subst_right,continuous,funs,funs_type_apply,real_sub,real_add_inverse,real_add_closed,standard_metric_value,real_abs_in_reals,real_leq_lt_trans,subseteq}.
    Show that $b\in\reals$ and $0 < b$ and
        for all $x,y\in I$ if $y = x + h$ or $y = x - h$ then
            $\abs{((\apply{f}{y} - \apply{f}{x}) / (y - x))} < b$.
    \begin{subproof}
        Follows by \cref{unit_interval_strict_fixed_scale_quotient_bound_49}.
    \end{subproof}
\end{proof}

% from §49 helper lemma 104: uniform-continuity witnesses control a fixed step
\begin{lemma}\label{unit_interval_uniform_continuity_step_value_bound_49}
    Let $I = \intervalclosed{0}{1}$.
    Let $T_I = \subspaceTopology{\standardTopologyR}{I}$.
    Let $d_I = \restrl{\standardMetricR}{I\times I}$.
    Suppose $f$ is continuous from $I$ to $\reals$ with respect to $T_I$ and $\standardTopologyR$.
    Assume $\eta\in\reals$.
    Suppose $0 < \eta$.
    Assume $\delta\in\reals$.
    Suppose $0 < \delta$.
    Suppose for all $x_0,x_1\in I$ such that $d_I((x_0,x_1)) < \delta$ we have
        $\apply{\standardMetricR}{(\apply{f}{x_0},\apply{f}{x_1})} < \eta$.
    Assume $h\in\reals$.
    Suppose $0 < h$.
    Suppose $h < \delta$.
    Then for all $x,y\in I$ if $y = x + h$ or $y = x - h$ then
        $\apply{\standardMetricR}{(\apply{f}{y},\apply{f}{x})} < \eta$.
\end{lemma}
\begin{proof}
    Show that for all $x,y\in I$ if $y = x + h$ then $d_I((y,x)) < \delta$.
    \begin{subproof}
        Follows by \cref{intervalclosed,times_tuple_intro,standard_metric_is_metric,metric_on,funs_elim,funs_is_function,subseteq,times_subseteq_left,times_subseteq_right,subseteq_transitive,pair_eq_iff,restrl_of_function_apply,standard_metric_value,real_sub,real_add_assoc,real_add_inverse,real_add_ident,real_add_comm,real_lt_subst_left,real_abs_neg_value_before_density_49,real_abs_nonneg,real_lt_imp_leq,real_lt_subst_right}.
    \end{subproof}
    Show that for all $x,y\in I$ if $y = x - h$ then $d_I((y,x)) < \delta$.
    \begin{subproof}
        Follows by \cref{intervalclosed,times_tuple_intro,standard_metric_is_metric,metric_on,funs_elim,funs_is_function,subseteq,times_subseteq_left,times_subseteq_right,subseteq_transitive,pair_eq_iff,restrl_of_function_apply,standard_metric_value,real_sub,real_add_assoc,real_add_inverse,real_add_ident,real_add_comm,real_lt_subst_left,real_abs_neg_value_before_density_49,real_abs_nonneg,real_lt_imp_leq,real_lt_subst_right}.
    \end{subproof}
    Show that for all $x,y\in I$ if $y = x + h$ or $y = x - h$ then
        $\apply{\standardMetricR}{(\apply{f}{y},\apply{f}{x})} < \eta$.
    \begin{subproof}
        Follows by \cref{subseteq}.
    \end{subproof}
\end{proof}

% from §49 helper lemma 105: reciprocals of positive natural numbers are at most one
\begin{lemma}\label{naturalsplus_reciprocal_leq_one_49}
    Assume $m\in\naturalsPlus$.
    Then $1 / m \leq 1$.
\end{lemma}
\begin{proof}
    Thus $1 / m \leq 1$
        by \cref{real_one_in_naturals,subseteq,real_one_leq_naturalsplus,naturalsplus_reciprocal_antitone,real_one_div_one,real_lt_subst_right}.
\end{proof}

% from §49 helper lemma 106: first mesh interval forward unit quotient on the initial quarter
\begin{lemma}\label{unit_interval_equal_mesh_first_left_forward_unit_quotient_49}
    Assume $A$ is the equal mesh of order $m$ on the unit interval.
    Assume $q\in\funs{\intervalclosed{0}{1}}{\reals}$.
    Assume for all $x\in\intervalclosed{0}{1}$ we have
        $q(x)$ is the distance from $x$ to $A$ in $\reals$ with respect to $\standardMetricR$.
    Let $h = ((1 / m) / (1 + 1)) / (1 + 1)$.
    Assume $x\in\intervalclosed{0}{1}$.
    Suppose $x \leq h$.
    Let $y = x + h$.
    Then $y\in\intervalclosed{0}{1}$ and
        $y = x + h$ and
        $\abs{((q(y) - q(x)) / (y - x))} = 1$.
\end{lemma}
\begin{proof}
    $m\in\naturalsPlus$ and $1 / m\in\reals$ and $0 < 1 / m$ by \cref{unit_interval_equal_mesh_49,naturalsplus_reciprocal_real,positive_reciprocal_naturalsplus}.
    Let $t = (1 / m) / (1 + 1)$.
    $t\in\reals$ and $0 < t$ and $t + t = 1 / m$ by \cref{real_half_of_positive}.
    $h\in\reals$ and $0 < h$ and $h + h = t$ by \cref{real_half_of_positive}.
    $x\in\reals$ and $y\in\reals$ by \cref{intervalclosed,real_add_closed}.
    $x \leq t$ by \cref{real_leq_trans,real_lt_imp_leq,real_add_ident,real_lt_add,real_lt_subst_left,real_lt_subst_right}.
    $y\leq t$ by \cref{real_leq_add,real_add_comm,real_leq_subst_right_local,real_lt_subst_left,real_lt_subst_right}.
    $t < 1 / m$ and $y < 1 / m$ by \cref{real_add_ident,real_lt_add,real_lt_subst_left,real_lt_subst_right,real_leq_lt_trans}.
    $0\leq x$ and $0\in\reals$ by \cref{intervalclosed,real_add_ident}.
    $0 \leq y$
        by \cref{real_lt_imp_leq,real_leq_add,real_add_ident,real_lt_subst_left,real_lt_subst_right,real_leq_trans}.
    Thus $y\in\intervalclosed{0}{1}$ by \cref{naturalsplus_reciprocal_leq_one_49,real_one_in_naturals,real_naturals_subset_reals,subseteq,real_lt_imp_leq,real_leq_trans,intervalclosed}.
    $q(x) = x - 0$ and $q(y) = y - 0$
        by \cref{unit_interval_equal_mesh_left_half_distance_49,real_sub_zero_local_49,real_lt_subst_right}.
    $\abs{((q(y) - q(x)) / (y - x))} = 1$
        by \cref{forward_linear_secant_quotient_one_49}.
\end{proof}

% from §49 helper lemma 107: first mesh interval backward unit quotient on the second quarter
\begin{lemma}\label{unit_interval_equal_mesh_first_left_backward_unit_quotient_49}
    Assume $A$ is the equal mesh of order $m$ on the unit interval.
    Assume $q\in\funs{\intervalclosed{0}{1}}{\reals}$.
    Assume for all $x\in\intervalclosed{0}{1}$ we have
        $q(x)$ is the distance from $x$ to $A$ in $\reals$ with respect to $\standardMetricR$.
    Let $h = ((1 / m) / (1 + 1)) / (1 + 1)$.
    Let $t = (1 / m) / (1 + 1)$.
    Assume $x\in\intervalclosed{0}{1}$.
    Suppose $h \leq x$.
    Suppose $x \leq t$.
    Let $y = x - h$.
    Then $y\in\intervalclosed{0}{1}$ and
        $y = x - h$ and
        $\abs{((q(y) - q(x)) / (y - x))} = 1$.
\end{lemma}
\begin{proof}
    $m\in\naturalsPlus$ and $1 / m\in\reals$ and $0 < 1 / m$ by \cref{unit_interval_equal_mesh_49,naturalsplus_reciprocal_real,positive_reciprocal_naturalsplus}.
    $t\in\reals$ and $0 < t$ and $t + t = 1 / m$ by \cref{real_half_of_positive}.
    $h\in\reals$ and $0 < h$ and $h + h = t$ by \cref{real_half_of_positive}.
    $x\in\reals$ and $y\in\reals$ by \cref{intervalclosed,real_sub,real_add_inverse,real_add_closed}.
    $0 \leq y$ by \cref{real_leq_add,real_sub,real_add_inverse,real_add_ident,real_lt_subst_left,real_lt_subst_right}.
    $y \leq x$ by \cref{real_sub,real_add_inverse,real_add_closed,real_lt_imp_leq,real_add_ident,real_lt_add,real_add_comm,real_leq_subst_left_local}.
    $y \leq t$ by \cref{real_leq_trans}.
    $t < 1 / m$ and $y < 1 / m$ by \cref{real_add_ident,real_lt_add,real_lt_subst_left,real_lt_subst_right,real_leq_lt_trans}.
    Thus $y\in\intervalclosed{0}{1}$
        by \cref{naturalsplus_reciprocal_leq_one_49,real_one_in_naturals,real_naturals_subset_reals,subseteq,real_lt_imp_leq,real_leq_trans,intervalclosed}.
    $0\in\reals$ by \cref{real_add_ident}.
    $q(x) = x - 0$ and $q(y) = y - 0$
        by \cref{unit_interval_equal_mesh_left_half_distance_49,real_sub_zero_local_49,real_lt_subst_right}.
    $\abs{((q(y) - q(x)) / (y - x))} = 1$
        by \cref{backward_linear_secant_quotient_one_49}.
\end{proof}

% from §49 helper lemma 108: first mesh interval forward unit quotient on the third quarter
\begin{lemma}\label{unit_interval_equal_mesh_first_right_forward_unit_quotient_49}
    Assume $A$ is the equal mesh of order $m$ on the unit interval.
    Assume $q\in\funs{\intervalclosed{0}{1}}{\reals}$.
    Assume for all $x\in\intervalclosed{0}{1}$ we have
        $q(x)$ is the distance from $x$ to $A$ in $\reals$ with respect to $\standardMetricR$.
    Let $h = ((1 / m) / (1 + 1)) / (1 + 1)$.
    Let $t = (1 / m) / (1 + 1)$.
    Assume $x\in\intervalclosed{0}{1}$.
    Suppose $t \leq x$.
    Suppose $x \leq (1 / m) - h$.
    Let $y = x + h$.
    Then $y\in\intervalclosed{0}{1}$ and
        $y = x + h$ and
        $\abs{((q(y) - q(x)) / (y - x))} = 1$.
\end{lemma}
\begin{proof}
    $m\in\naturalsPlus$ and $1 / m\in\reals$ and $0 < 1 / m$ by \cref{unit_interval_equal_mesh_49,naturalsplus_reciprocal_real,positive_reciprocal_naturalsplus}.
    $t\in\reals$ and $0 < t$ and $t + t = 1 / m$ by \cref{real_half_of_positive}.
    $h\in\reals$ and $0 < h$ and $h + h = t$ by \cref{real_half_of_positive}.
    $x\in\reals$ and $y\in\reals$ by \cref{intervalclosed,real_add_closed}.
    $1 / m - h\in\reals$ by \cref{real_sub,real_add_inverse,real_add_closed}.
    $y \leq 1 / m$
        by \cref{real_leq_add,real_sub,real_add_assoc,real_add_inverse,real_add_ident,real_add_comm,real_lt_subst_left,real_lt_subst_right}.
    $0\in\reals$ and $0 < x$ by \cref{real_add_ident,real_leq_lt_trans}.
    $0 < y$ by \cref{real_pos_add,real_lt_subst_right}.
    $0 \leq y$ by \cref{real_lt_imp_leq}.
    Thus $y\in\intervalclosed{0}{1}$
        by \cref{naturalsplus_reciprocal_leq_one_49,real_one_in_naturals,real_naturals_subset_reals,subseteq,real_leq_trans,intervalclosed}.
    $(1 / m) - h < 1 / m$ and $(1 / m) - h \leq 1 / m$ and $x \leq 1 / m$
        by \cref{real_add_inverse,real_pos_imp_neg,real_lt_add,real_add_comm,real_sub,real_add_ident,real_lt_subst_left,real_lt_subst_right,real_lt_imp_leq,real_leq_trans}.
    $t \leq y$
        by \cref{real_leq_add,real_lt_imp_leq,real_add_ident,real_lt_add,real_add_comm,real_lt_subst_left,real_leq_trans,real_lt_subst_right,real_add_closed}.
    $q(x) = (1 / m) - x$ and $q(y) = (1 / m) - y$ by \cref{unit_interval_equal_mesh_right_half_distance_49}.
    $\abs{((q(y) - q(x)) / (y - x))} = 1$
        by \cref{forward_descending_linear_secant_quotient_one_49}.
\end{proof}

% from §49 helper lemma 109: first mesh interval backward unit quotient on the last quarter
\begin{lemma}\label{unit_interval_equal_mesh_first_right_backward_unit_quotient_49}
    Assume $A$ is the equal mesh of order $m$ on the unit interval.
    Assume $q\in\funs{\intervalclosed{0}{1}}{\reals}$.
    Assume for all $x\in\intervalclosed{0}{1}$ we have
        $q(x)$ is the distance from $x$ to $A$ in $\reals$ with respect to $\standardMetricR$.
    Let $h = ((1 / m) / (1 + 1)) / (1 + 1)$.
    Let $t = (1 / m) / (1 + 1)$.
    Assume $x\in\intervalclosed{0}{1}$.
    Suppose $(1 / m) - h \leq x$.
    Suppose $x \leq 1 / m$.
    Let $y = x - h$.
    Then $y\in\intervalclosed{0}{1}$ and
        $y = x - h$ and
        $\abs{((q(y) - q(x)) / (y - x))} = 1$.
\end{lemma}
\begin{proof}
    $m\in\naturalsPlus$ and $1 / m\in\reals$ and $0 < 1 / m$ by \cref{unit_interval_equal_mesh_49,naturalsplus_reciprocal_real,positive_reciprocal_naturalsplus}.
    $t\in\reals$ and $0 < t$ and $t + t = 1 / m$ by \cref{real_half_of_positive}.
    $h\in\reals$ and $0 < h$ and $h + h = t$ by \cref{real_half_of_positive}.
    $x\in\reals$ and $y\in\reals$ by \cref{intervalclosed,real_sub,real_add_inverse,real_add_closed}.
    $1 / m - h\in\reals$ by \cref{real_sub,real_add_inverse,real_add_closed}.
    $\neg{h}\in\reals$ by \cref{real_add_inverse}.
    $t \leq y$
        by \cref{real_sub,real_add_assoc,real_add_inverse,real_add_ident,real_add_comm,real_add_eq_subst,real_leq_add,real_lt_subst_left,real_lt_subst_right}.
    $0\in\reals$ and $0 < y$ and $0 \leq y$ by \cref{real_add_ident,real_leq_lt_trans,real_lt_imp_leq}.
    $y < x$ and $y \leq x$
        by \cref{real_pos_imp_neg,real_lt_add,real_add_comm,real_add_ident,real_sub,real_lt_subst_left,real_lt_subst_right,real_lt_imp_leq}.
    $y \leq 1 / m$ by \cref{real_leq_trans}.
    Thus $y\in\intervalclosed{0}{1}$
        by \cref{naturalsplus_reciprocal_leq_one_49,real_one_in_naturals,real_naturals_subset_reals,subseteq,real_leq_trans,intervalclosed}.
    $t \leq x$
        by \cref{real_sub,real_add_assoc,real_add_inverse,real_add_ident,real_add_comm,real_add_eq_subst,real_lt_subst_right,real_lt_imp_leq,real_lt_add,real_lt_subst_left,real_leq_trans}.
    $q(x) = (1 / m) - x$ and $q(y) = (1 / m) - y$ by \cref{unit_interval_equal_mesh_right_half_distance_49}.
    $\abs{((q(y) - q(x)) / (y - x))} = 1$
        by \cref{backward_descending_linear_secant_quotient_one_49}.
\end{proof}

% from §49 helper lemma 110: first mesh interval has a quarter-step unit quotient
\begin{lemma}\label{unit_interval_equal_mesh_first_interval_unit_quotient_49}
    Assume $A$ is the equal mesh of order $m$ on the unit interval.
    Assume $q\in\funs{\intervalclosed{0}{1}}{\reals}$.
    Assume for all $x\in\intervalclosed{0}{1}$ we have
        $q(x)$ is the distance from $x$ to $A$ in $\reals$ with respect to $\standardMetricR$.
    Let $h = ((1 / m) / (1 + 1)) / (1 + 1)$.
    Let $t = (1 / m) / (1 + 1)$.
    Assume $x\in\intervalclosed{0}{1}$.
    Suppose $x \leq 1 / m$.
    Then there exists $y\in\intervalclosed{0}{1}$ such that
        either $y = x + h$ and
            $\abs{((q(y) - q(x)) / (y - x))} = 1$
        or $y = x - h$ and
            $\abs{((q(y) - q(x)) / (y - x))} = 1$.
\end{lemma}
\begin{proof}
    $m\in\naturalsPlus$ and $1 / m\in\reals$ and $0 < 1 / m$ by \cref{unit_interval_equal_mesh_49,naturalsplus_reciprocal_real,positive_reciprocal_naturalsplus}.
    $t\in\reals$ and $0 < t$ and $t + t = 1 / m$ by \cref{real_half_of_positive}.
    $h\in\reals$ and $0 < h$ and $h + h = t$ by \cref{real_half_of_positive}.
    $x\in\reals$ by \cref{intervalclosed}.
    $1 / m - h\in\reals$ by \cref{real_sub,real_add_inverse,real_add_closed}.
    \begin{byCase}
    \caseOf{$x \leq h$.}
        Let $y = x + h$.
        $y\in\intervalclosed{0}{1}$ and $y = x + h$ and
            $\abs{((q(y) - q(x)) / (y - x))} = 1$
            by \cref{unit_interval_equal_mesh_first_left_forward_unit_quotient_49}.
        Show that there exists $z\in\intervalclosed{0}{1}$ such that either $z = x + h$ and
            $\abs{((q(z) - q(x)) / (z - x))} = 1$ or $z = x - h$ and
            $\abs{((q(z) - q(x)) / (z - x))} = 1$.
        \begin{subproof}
            Trivial.
        \end{subproof}
    \caseOf{it is wrong that $x \leq h$.}
        $h < x$ and $h \leq x$ by \cref{real_lt_trichotomy,real_lt_imp_leq,real_lt_subst_right}.
        \begin{byCase}
        \caseOf{$x \leq t$.}
            Let $y = x - h$.
            $y\in\intervalclosed{0}{1}$ and $y = x - h$ and
                $\abs{((q(y) - q(x)) / (y - x))} = 1$
                by \cref{unit_interval_equal_mesh_first_left_backward_unit_quotient_49}.
            Show that there exists $z\in\intervalclosed{0}{1}$ such that either $z = x + h$ and
                $\abs{((q(z) - q(x)) / (z - x))} = 1$ or $z = x - h$ and
                $\abs{((q(z) - q(x)) / (z - x))} = 1$.
            \begin{subproof}
                Trivial.
            \end{subproof}
        \caseOf{it is wrong that $x \leq t$.}
            $t < x$ and $t \leq x$ by \cref{real_lt_trichotomy,real_lt_imp_leq,real_lt_subst_right}.
            \begin{byCase}
            \caseOf{$x \leq (1 / m) - h$.}
                Let $y = x + h$.
                $y\in\intervalclosed{0}{1}$ and $y = x + h$ and
                    $\abs{((q(y) - q(x)) / (y - x))} = 1$
                    by \cref{unit_interval_equal_mesh_first_right_forward_unit_quotient_49}.
                Show that there exists $z\in\intervalclosed{0}{1}$ such that either $z = x + h$ and
                    $\abs{((q(z) - q(x)) / (z - x))} = 1$ or $z = x - h$ and
                    $\abs{((q(z) - q(x)) / (z - x))} = 1$.
                \begin{subproof}
                    Trivial.
                \end{subproof}
            \caseOf{it is wrong that $x \leq (1 / m) - h$.}
                $(1 / m) - h < x$ and $(1 / m) - h \leq x$ by \cref{real_lt_trichotomy,real_lt_imp_leq,real_lt_subst_right}.
                Let $y = x - h$.
                $y\in\intervalclosed{0}{1}$ and $y = x - h$ and
                    $\abs{((q(y) - q(x)) / (y - x))} = 1$
                    by \cref{unit_interval_equal_mesh_first_right_backward_unit_quotient_49}.
                Show that there exists $z\in\intervalclosed{0}{1}$ such that either $z = x + h$ and
                    $\abs{((q(z) - q(x)) / (z - x))} = 1$ or $z = x - h$ and
                    $\abs{((q(z) - q(x)) / (z - x))} = 1$.
                \begin{subproof}
                    Trivial.
                \end{subproof}
            \end{byCase}
        \end{byCase}
    \end{byCase}
\end{proof}

% from §49 helper lemma 111: adjacent left-half equal-mesh distance formula
\begin{lemma}\label{unit_interval_equal_mesh_adjacent_left_half_formula_49}
    Assume $A$ is the equal mesh of order $m$ on the unit interval.
    Assume $q\in\funs{\intervalclosed{0}{1}}{\reals}$.
    Assume for all $x\in\intervalclosed{0}{1}$ we have
        $q(x)$ is the distance from $x$ to $A$ in $\reals$ with respect to $\standardMetricR$.
    Assume $p\in\naturalsPlus$.
    Assume $p + 1 \leq m$.
    Assume $x\in\intervalclosed{0}{1}$.
    Suppose $p / m \leq x$.
    Suppose $x \leq p / m + (1 / m) / (1 + 1)$.
    Then $q(x) = x - p / m$.
\end{lemma}
\begin{proof}
    $m\in\naturalsPlus$ and $p\in\reals$ and $m\in\reals$ and $p + 1\in\reals$ by \cref{unit_interval_equal_mesh_49,real_naturals_subset_reals,subseteq,real_naturals_closed_succ}.
    $1\in\reals$ by \cref{real_one_in_naturals,real_naturals_subset_reals,subseteq}.
    $1 / m\in\reals$ and $0 < 1 / m$ by \cref{naturalsplus_reciprocal_real,positive_reciprocal_naturalsplus}.
    Let $t = (1 / m) / (1 + 1)$.
    $t\in\reals$ and $0 < t$ and $t + t = 1 / m$ by \cref{real_half_of_positive}.
    $x\in\reals$ by \cref{intervalclosed}.
    $x \leq p / m + t$ by \cref{real_lt_subst_right}.
    $p / m\in\reals$ and $(p + 1) / m\in\reals$ by \cref{real_div_naturalsplus_real,real_naturals_closed_succ}.
    $p / m + t\in\reals$ and $x + t\in\reals$ and $((p / m) + t) + t\in\reals$ by \cref{real_add_closed}.
    $(p + 1) / m = (p / m) + (1 / m)$ by \cref{real_div_add_same_denom}.
    $(p / m) + (1 / m) = (p / m) + (t + t)$ by \cref{real_add_eq_subst}.
    $(p / m) + (t + t) = ((p / m) + t) + t$ by \cref{real_add_assoc}.
    $x + t \leq ((p / m) + t) + t$ by \cref{real_leq_add}.
    $(p + 1) / m = ((p / m) + t) + t$ by \cref{real_lt_subst_left,real_lt_subst_right}.
    $x + t \leq (p + 1) / m$ by \cref{real_lt_subst_right}.
    $x \leq x + t$ by \cref{real_lt_imp_leq,real_add_ident,real_lt_add,real_add_comm,real_lt_subst_left}.
    $x \leq (p + 1) / m$ by \cref{real_leq_trans,real_add_closed}.
    $\neg{(p / m)}\in\reals$ by \cref{real_add_inverse}.
    $x + \neg{(p / m)} \leq ((p / m) + t) + \neg{(p / m)}$ by \cref{real_leq_add}.
    $x + \neg{(p / m)} = x - p / m$ by \cref{real_sub}.
    $((p / m) + t) + \neg{(p / m)} = t$
        by \cref{real_add_assoc,real_add_inverse,real_add_ident,real_add_comm}.
    $x - p / m \leq t$ by \cref{real_lt_subst_left,real_lt_subst_right}.
    $\neg{x}\in\reals$ by \cref{real_add_inverse}.
    $(x + t) + \neg{x} \leq ((p + 1) / m) + \neg{x}$ by \cref{real_leq_add}.
    $(x + t) + \neg{x} = t$ by \cref{real_add_assoc,real_add_inverse,real_add_ident,real_add_comm}.
    $((p + 1) / m) + \neg{x} = (p + 1) / m - x$ by \cref{real_sub}.
    $t \leq (p + 1) / m - x$ by \cref{real_lt_subst_left,real_lt_subst_right}.
    $x - p / m \leq (p + 1) / m - x$ by \cref{real_leq_trans,real_sub,real_add_inverse,real_add_closed}.
    $\apply{\standardMetricR}{(x,p / m)} = \abs{(x - p / m)}$ by \cref{standard_metric_value}.
    $0 \leq x - p / m$ by \cref{real_leq_add,real_sub,real_add_inverse,real_add_ident,real_lt_subst_left,real_lt_subst_right}.
    $\abs{(x - p / m)} = x - p / m$ by \cref{real_abs_nonneg,real_sub,real_add_inverse,real_add_closed}.
    $\apply{\standardMetricR}{(x,p / m)} = x - p / m$ by \cref{real_lt_subst_right}.
    $\apply{\standardMetricR}{(x,(p + 1) / m)} = \abs{(x - (p + 1) / m)}$ by \cref{standard_metric_value}.
    $0 \leq (p + 1) / m - x$ by \cref{real_leq_add,real_sub,real_add_inverse,real_add_ident,real_lt_subst_left,real_lt_subst_right}.
    $\abs{(x - (p + 1) / m)} = \abs{((p + 1) / m - x)}$ by \cref{real_abs_sub_swap}.
    $\abs{((p + 1) / m - x)} = (p + 1) / m - x$ by \cref{real_abs_nonneg,real_sub,real_add_inverse,real_add_closed}.
    $\apply{\standardMetricR}{(x,(p + 1) / m)} = (p + 1) / m - x$ by \cref{real_lt_subst_left,real_lt_subst_right}.
    $\apply{\standardMetricR}{(x,p / m)} \leq \apply{\standardMetricR}{(x,(p + 1) / m)}$
        by \cref{real_lt_subst_left,real_lt_subst_right}.
    Therefore $q(x) = x - p / m$ by \cref{unit_interval_equal_mesh_adjacent_left_formula_49}.
\end{proof}

% from §49 helper lemma 112: adjacent right-half equal-mesh distance formula
\begin{lemma}\label{unit_interval_equal_mesh_adjacent_right_half_formula_49}
    Assume $A$ is the equal mesh of order $m$ on the unit interval.
    Assume $q\in\funs{\intervalclosed{0}{1}}{\reals}$.
    Assume for all $x\in\intervalclosed{0}{1}$ we have
        $q(x)$ is the distance from $x$ to $A$ in $\reals$ with respect to $\standardMetricR$.
    Assume $p\in\naturalsPlus$.
    Assume $p + 1 \leq m$.
    Assume $x\in\intervalclosed{0}{1}$.
    Suppose $p / m + (1 / m) / (1 + 1) \leq x$.
    Suppose $x \leq (p + 1) / m$.
    Then $q(x) = (p + 1) / m - x$.
\end{lemma}
\begin{proof}
    $m\in\naturalsPlus$ and $p\in\reals$ and $m\in\reals$ and $p + 1\in\reals$ by \cref{unit_interval_equal_mesh_49,real_naturals_subset_reals,subseteq,real_naturals_closed_succ}.
    $1\in\reals$ by \cref{real_one_in_naturals,real_naturals_subset_reals,subseteq}.
    $1 / m\in\reals$ and $0 < 1 / m$ by \cref{naturalsplus_reciprocal_real,positive_reciprocal_naturalsplus}.
    Let $t = (1 / m) / (1 + 1)$.
    $t\in\reals$ and $0 < t$ and $t + t = 1 / m$ by \cref{real_half_of_positive}.
    $x\in\reals$ by \cref{intervalclosed}.
    $p / m + t \leq x$ by \cref{real_lt_subst_left}.
    $p / m\in\reals$ and $(p + 1) / m\in\reals$ by \cref{real_div_naturalsplus_real,real_naturals_closed_succ}.
    $p / m + t\in\reals$ and $x + t\in\reals$ and $((p / m) + t) + t\in\reals$ by \cref{real_add_closed}.
    $(p + 1) / m = (p / m) + (1 / m)$ by \cref{real_div_add_same_denom}.
    $(p / m) + (1 / m) = (p / m) + (t + t)$ by \cref{real_add_eq_subst}.
    $(p / m) + (t + t) = ((p / m) + t) + t$ by \cref{real_add_assoc}.
    $(p + 1) / m = ((p / m) + t) + t$ by \cref{real_lt_subst_left,real_lt_subst_right}.
    $p / m \leq p / m + t$ by \cref{real_lt_imp_leq,real_add_ident,real_lt_add,real_add_comm,real_lt_subst_left}.
    $p / m \leq x$ by \cref{real_leq_trans,real_add_closed}.
    $\neg{x}\in\reals$ by \cref{real_add_inverse}.
    $((p / m) + t) + t \leq x + t$ by \cref{real_leq_add}.
    $(p + 1) / m \leq x + t$ by \cref{real_lt_subst_left}.
    $((p + 1) / m) + \neg{x} \leq (x + t) + \neg{x}$ by \cref{real_leq_add}.
    $((p + 1) / m) + \neg{x} = (p + 1) / m - x$ by \cref{real_sub}.
    $(x + t) + \neg{x} = t$ by \cref{real_add_assoc,real_add_inverse,real_add_ident,real_add_comm}.
    $(p + 1) / m - x \leq t$ by \cref{real_lt_subst_left,real_lt_subst_right}.
    $\neg{(p / m)}\in\reals$ by \cref{real_add_inverse}.
    $((p / m) + t) + \neg{(p / m)} \leq x + \neg{(p / m)}$ by \cref{real_leq_add}.
    $((p / m) + t) + \neg{(p / m)} = t$
        by \cref{real_add_assoc,real_add_inverse,real_add_ident,real_add_comm}.
    $x + \neg{(p / m)} = x - p / m$ by \cref{real_sub}.
    $t \leq x - p / m$ by \cref{real_lt_subst_left,real_lt_subst_right}.
    $(p + 1) / m - x \leq x - p / m$ by \cref{real_leq_trans,real_sub,real_add_inverse,real_add_closed}.
    $\apply{\standardMetricR}{(x,p / m)} = \abs{(x - p / m)}$ by \cref{standard_metric_value}.
    $0 \leq x - p / m$ by \cref{real_leq_add,real_sub,real_add_inverse,real_add_ident,real_lt_subst_left,real_lt_subst_right}.
    $\abs{(x - p / m)} = x - p / m$ by \cref{real_abs_nonneg,real_sub,real_add_inverse,real_add_closed}.
    $\apply{\standardMetricR}{(x,p / m)} = x - p / m$ by \cref{real_lt_subst_right}.
    $\apply{\standardMetricR}{(x,(p + 1) / m)} = \abs{(x - (p + 1) / m)}$ by \cref{standard_metric_value}.
    $0 \leq (p + 1) / m - x$ by \cref{real_leq_add,real_sub,real_add_inverse,real_add_ident,real_lt_subst_left,real_lt_subst_right}.
    $\abs{(x - (p + 1) / m)} = \abs{((p + 1) / m - x)}$ by \cref{real_abs_sub_swap}.
    $\abs{((p + 1) / m - x)} = (p + 1) / m - x$ by \cref{real_abs_nonneg,real_sub,real_add_inverse,real_add_closed}.
    $\apply{\standardMetricR}{(x,(p + 1) / m)} = (p + 1) / m - x$ by \cref{real_lt_subst_left,real_lt_subst_right}.
    $\apply{\standardMetricR}{(x,(p + 1) / m)} \leq \apply{\standardMetricR}{(x,p / m)}$
        by \cref{real_lt_subst_left,real_lt_subst_right}.
    Therefore $q(x) = (p + 1) / m - x$ by \cref{unit_interval_equal_mesh_adjacent_right_formula_49}.
\end{proof}

% from §49 helper lemma 113: adjacent mesh interval forward unit quotient on the initial quarter
\begin{lemma}\label{unit_interval_equal_mesh_adjacent_left_forward_unit_quotient_49}
    Assume $A$ is the equal mesh of order $m$ on the unit interval.
    Assume $q\in\funs{\intervalclosed{0}{1}}{\reals}$.
    Assume for all $x\in\intervalclosed{0}{1}$ we have
        $q(x)$ is the distance from $x$ to $A$ in $\reals$ with respect to $\standardMetricR$.
    Assume $p\in\naturalsPlus$.
    Assume $p + 1 \leq m$.
    Let $h = ((1 / m) / (1 + 1)) / (1 + 1)$.
    Let $t = (1 / m) / (1 + 1)$.
    Assume $x\in\intervalclosed{0}{1}$.
    Suppose $p / m \leq x$.
    Suppose $x \leq p / m + h$.
    Let $y = x + h$.
    Then $y\in\intervalclosed{0}{1}$ and
        $y = x + h$ and
        $\abs{((q(y) - q(x)) / (y - x))} = 1$.
\end{lemma}
\begin{proof}
    $m\in\naturalsPlus$ and $p\in\reals$ and $m\in\reals$ and $p + 1\in\reals$ by \cref{unit_interval_equal_mesh_49,real_naturals_subset_reals,subseteq,real_naturals_closed_succ}.
    $1\in\reals$ by \cref{real_one_in_naturals,real_naturals_subset_reals,subseteq}.
    $1 / m\in\reals$ and $0 < 1 / m$ by \cref{naturalsplus_reciprocal_real,positive_reciprocal_naturalsplus}.
    $t\in\reals$ and $0 < t$ and $t + t = 1 / m$ by \cref{real_half_of_positive}.
    $h\in\reals$ and $0 < h$ and $h + h = t$ by \cref{real_half_of_positive}.
    $p / m\in\reals$ and $(p + 1) / m\in\reals$ by \cref{real_div_naturalsplus_real,real_naturals_closed_succ}.
    $x\in\reals$ and $y\in\reals$ by \cref{intervalclosed,real_add_closed}.
    $p / m + h\in\reals$ and $p / m + t\in\reals$ by \cref{real_add_closed}.
    $h < t$ and $h \leq t$ by \cref{real_add_ident,real_lt_add,real_lt_subst_left,real_lt_subst_right,real_lt_imp_leq}.
    $p / m + h \leq p / m + t$ by \cref{real_leq_add,real_add_comm,real_leq_subst_right_local}.
    $x \leq p / m + t$ by \cref{real_leq_trans}.
    $x + h \leq (p / m + h) + h$ by \cref{real_leq_add}.
    $(p / m + h) + h = p / m + (h + h)$ by \cref{real_add_assoc}.
    $p / m + (h + h) = p / m + t$ by \cref{real_add_eq_subst}.
    $y \leq p / m + t$ by \cref{real_lt_subst_left,real_lt_subst_right}.
    $x \leq y$ by \cref{real_lt_imp_leq,real_add_ident,real_lt_add,real_add_comm,real_lt_subst_left}.
    $p / m \leq y$ by \cref{real_leq_trans}.
    $(p + 1) / m = (p / m) + (1 / m)$ by \cref{real_div_add_same_denom}.
    $(p / m) + (1 / m) = (p / m) + (t + t)$ by \cref{real_add_eq_subst}.
    $(p / m) + (t + t) = (p / m + t) + t$ by \cref{real_add_assoc}.
    $(p + 1) / m = (p / m + t) + t$ by \cref{real_lt_subst_left,real_lt_subst_right}.
    $p / m + t \leq (p / m + t) + t$ by \cref{real_lt_imp_leq,real_add_ident,real_lt_add,real_add_comm,real_lt_subst_left}.
    $p / m + t \leq (p + 1) / m$ by \cref{real_lt_subst_right}.
    $y \leq (p + 1) / m$ by \cref{real_leq_trans}.
    $0\leq x$ by \cref{intervalclosed}.
    $0 \leq y$ by \cref{real_leq_split,real_pos_add,real_lt_imp_leq,real_lt_subst_right,pair_eq_iff,real_add_ident,real_lt_subst_left}.
    Thus $y\in\intervalclosed{0}{1}$
        by \cref{real_div_leq_same_denom,naturalsplus_positive_real,real_pos_nonzero,real_div,real_mul_inverse,real_mul_comm,real_lt_subst_left,real_lt_subst_right,real_leq_trans,intervalclosed}.
    $q(x) = x - p / m$ and $q(y) = y - p / m$ by \cref{unit_interval_equal_mesh_adjacent_left_half_formula_49}.
    $\abs{((q(y) - q(x)) / (y - x))} = 1$
        by \cref{forward_linear_secant_quotient_one_49}.
\end{proof}

% from §49 helper lemma 114: adjacent mesh interval backward unit quotient on the second quarter
\begin{lemma}\label{unit_interval_equal_mesh_adjacent_left_backward_unit_quotient_49}
    Assume $A$ is the equal mesh of order $m$ on the unit interval.
    Assume $q\in\funs{\intervalclosed{0}{1}}{\reals}$.
    Assume for all $x\in\intervalclosed{0}{1}$ we have
        $q(x)$ is the distance from $x$ to $A$ in $\reals$ with respect to $\standardMetricR$.
    Assume $p\in\naturalsPlus$.
    Assume $p + 1 \leq m$.
    Let $h = ((1 / m) / (1 + 1)) / (1 + 1)$.
    Let $t = (1 / m) / (1 + 1)$.
    Assume $x\in\intervalclosed{0}{1}$.
    Suppose $p / m + h \leq x$.
    Suppose $x \leq p / m + t$.
    Let $y = x - h$.
    Then $y\in\intervalclosed{0}{1}$ and
        $y = x - h$ and
        $\abs{((q(y) - q(x)) / (y - x))} = 1$.
\end{lemma}
\begin{proof}
    $m\in\naturalsPlus$ and $p\in\reals$ and $m\in\reals$ and $p + 1\in\reals$ by \cref{unit_interval_equal_mesh_49,real_naturals_subset_reals,subseteq,real_naturals_closed_succ}.
    $1\in\reals$ by \cref{real_one_in_naturals,real_naturals_subset_reals,subseteq}.
    $1 / m\in\reals$ and $0 < 1 / m$ by \cref{naturalsplus_reciprocal_real,positive_reciprocal_naturalsplus}.
    $t\in\reals$ and $0 < t$ and $t + t = 1 / m$ by \cref{real_half_of_positive}.
    $h\in\reals$ and $0 < h$ and $h + h = t$ by \cref{real_half_of_positive}.
    $p / m\in\reals$ by \cref{real_div_naturalsplus_real}.
    $x\in\reals$ and $y\in\reals$ by \cref{intervalclosed,real_sub,real_add_inverse,real_add_closed}.
    $p / m + h\in\reals$ and $p / m + t\in\reals$ by \cref{real_add_closed}.
    $0\in\reals$ by \cref{real_add_ident}.
    $p / m \leq p / m + h$ by \cref{real_lt_imp_leq,real_add_ident,real_lt_add,real_add_comm,real_lt_subst_left}.
    $p / m \leq x$ by \cref{real_leq_trans}.
    $\neg{h}\in\reals$ by \cref{real_add_inverse}.
    $(p / m + h) + \neg{h} \leq x + \neg{h}$ by \cref{real_leq_add}.
    $(p / m + h) + \neg{h} = p / m$ by \cref{real_add_assoc,real_add_inverse,real_add_ident,real_add_comm}.
    $x + \neg{h} = x - h$ by \cref{real_sub}.
    $p / m \leq y$ by \cref{real_lt_subst_left,real_lt_subst_right}.
    $\neg{h} < 0$ by \cref{real_pos_imp_neg}.
    $\neg{h} + x < 0 + x$ by \cref{real_lt_add}.
    $\neg{h} + x = x + \neg{h}$ by \cref{real_add_comm}.
    $0 + x = x$ by \cref{real_add_ident}.
    $y < x$ and $y \leq x$ by \cref{real_lt_subst_left,real_lt_subst_right,real_lt_imp_leq}.
    $y \leq p / m + t$ by \cref{real_leq_trans}.
    $0 \leq p / m$ by \cref{real_div_naturalsplus_nonnegative,real_lt_imp_leq,real_add_ident,naturalsplus_positive_real}.
    $0 \leq y$ by \cref{real_leq_trans}.
    Thus $y\in\intervalclosed{0}{1}$ by \cref{intervalclosed,real_leq_trans}.
    $q(x) = x - p / m$ and $q(y) = y - p / m$ by \cref{unit_interval_equal_mesh_adjacent_left_half_formula_49}.
    $\abs{((q(y) - q(x)) / (y - x))} = 1$
        by \cref{backward_linear_secant_quotient_one_49}.
\end{proof}

% from §49 helper lemma 115: adjacent mesh interval forward unit quotient on the third quarter
\begin{lemma}\label{unit_interval_equal_mesh_adjacent_right_forward_unit_quotient_49}
    Assume $A$ is the equal mesh of order $m$ on the unit interval.
    Assume $q\in\funs{\intervalclosed{0}{1}}{\reals}$.
    Assume for all $x\in\intervalclosed{0}{1}$ we have
        $q(x)$ is the distance from $x$ to $A$ in $\reals$ with respect to $\standardMetricR$.
    Assume $p\in\naturalsPlus$.
    Assume $p + 1 \leq m$.
    Let $h = ((1 / m) / (1 + 1)) / (1 + 1)$.
    Let $t = (1 / m) / (1 + 1)$.
    Assume $x\in\intervalclosed{0}{1}$.
    Suppose $p / m + t \leq x$.
    Suppose $x \leq (p + 1) / m - h$.
    Let $y = x + h$.
    Then $y\in\intervalclosed{0}{1}$ and
        $y = x + h$ and
        $\abs{((q(y) - q(x)) / (y - x))} = 1$.
\end{lemma}
\begin{proof}
    $m\in\naturalsPlus$ and $p\in\reals$ and $m\in\reals$ and $p + 1\in\reals$ by \cref{unit_interval_equal_mesh_49,real_naturals_subset_reals,subseteq,real_naturals_closed_succ}.
    $1\in\reals$ by \cref{real_one_in_naturals,real_naturals_subset_reals,subseteq}.
    $1 / m\in\reals$ and $0 < 1 / m$ by \cref{naturalsplus_reciprocal_real,positive_reciprocal_naturalsplus}.
    $t\in\reals$ and $0 < t$ and $t + t = 1 / m$ by \cref{real_half_of_positive}.
    $h\in\reals$ and $0 < h$ and $h + h = t$ by \cref{real_half_of_positive}.
    $p / m\in\reals$ and $(p + 1) / m\in\reals$ by \cref{real_div_naturalsplus_real,real_naturals_closed_succ}.
    $x\in\reals$ and $y\in\reals$ and $0\in\reals$ by \cref{intervalclosed,real_add_closed,real_add_ident}.
    $(p + 1) / m - h\in\reals$ by \cref{real_sub,real_add_inverse,real_add_closed}.
    $x + h \leq ((p + 1) / m - h) + h$ by \cref{real_leq_add}.
    $((p + 1) / m - h) + h = (p + 1) / m$
        by \cref{real_sub,real_add_assoc,real_add_inverse,real_add_ident,real_add_comm}.
    $y \leq (p + 1) / m$ by \cref{real_lt_subst_left,real_lt_subst_right}.
    $x \leq y$ by \cref{real_lt_imp_leq,real_add_ident,real_lt_add,real_add_comm,real_lt_subst_left}.
    $p / m + t \leq y$ by \cref{real_leq_trans,real_add_closed}.
    $\neg{h}\in\reals$ by \cref{real_add_inverse}.
    $\neg{h} < 0$ by \cref{real_pos_imp_neg}.
    $\neg{h} + (p + 1) / m < 0 + (p + 1) / m$ by \cref{real_lt_add}.
    $\neg{h} + (p + 1) / m = (p + 1) / m + \neg{h}$ by \cref{real_add_comm}.
    $0 + (p + 1) / m = (p + 1) / m$ by \cref{real_add_ident}.
    $(p + 1) / m + \neg{h} = (p + 1) / m - h$ by \cref{real_sub}.
    $(p + 1) / m - h < (p + 1) / m$ by \cref{real_lt_subst_left,real_lt_subst_right}.
    $(p + 1) / m - h \leq (p + 1) / m$ by \cref{real_lt_imp_leq}.
    $x \leq (p + 1) / m$ by \cref{real_leq_trans}.
    $0\leq x$ by \cref{intervalclosed}.
    $0 \leq y$ by \cref{real_leq_split,real_pos_add,real_lt_imp_leq,real_lt_subst_right,pair_eq_iff,real_add_ident,real_lt_subst_left}.
    Thus $y\in\intervalclosed{0}{1}$
        by \cref{real_div_leq_same_denom,naturalsplus_positive_real,real_pos_nonzero,real_div,real_mul_inverse,real_mul_comm,real_lt_subst_left,real_lt_subst_right,real_leq_trans,intervalclosed}.
    $q(x) = (p + 1) / m - x$ and $q(y) = (p + 1) / m - y$ by \cref{unit_interval_equal_mesh_adjacent_right_half_formula_49}.
    $\abs{((q(y) - q(x)) / (y - x))} = 1$
        by \cref{forward_descending_linear_secant_quotient_one_49}.
\end{proof}

% from §49 helper lemma 116: adjacent mesh interval forward unit quotient before the last quarter
\begin{lemma}\label{unit_interval_equal_mesh_adjacent_right_forward_sum_unit_quotient_49}
    Assume $A$ is the equal mesh of order $m$ on the unit interval.
    Assume $q\in\funs{\intervalclosed{0}{1}}{\reals}$.
    Assume for all $x\in\intervalclosed{0}{1}$ we have
        $q(x)$ is the distance from $x$ to $A$ in $\reals$ with respect to $\standardMetricR$.
    Assume $p\in\naturalsPlus$.
    Assume $p + 1 \leq m$.
    Let $h = ((1 / m) / (1 + 1)) / (1 + 1)$.
    Let $t = (1 / m) / (1 + 1)$.
    Assume $x\in\intervalclosed{0}{1}$.
    Suppose $p / m + t \leq x$.
    Suppose $x \leq (p / m + t) + h$.
    Let $y = x + h$.
    Then $y\in\intervalclosed{0}{1}$ and
        $y = x + h$ and
        $\abs{((q(y) - q(x)) / (y - x))} = 1$.
\end{lemma}
\begin{proof}
    $m\in\naturalsPlus$ and $p\in\reals$ and $m\in\reals$ and $p + 1\in\reals$ by \cref{unit_interval_equal_mesh_49,real_naturals_subset_reals,subseteq,real_naturals_closed_succ}.
    $1\in\reals$ by \cref{real_one_in_naturals,real_naturals_subset_reals,subseteq}.
    $1 / m\in\reals$ and $0 < 1 / m$ by \cref{naturalsplus_reciprocal_real,positive_reciprocal_naturalsplus}.
    $t\in\reals$ and $0 < t$ and $t + t = 1 / m$ by \cref{real_half_of_positive}.
    $h\in\reals$ and $0 < h$ and $h + h = t$ by \cref{real_half_of_positive}.
    $p / m\in\reals$ and $(p + 1) / m\in\reals$ by \cref{real_div_naturalsplus_real,real_naturals_closed_succ}.
    $x\in\reals$ and $y\in\reals$ by \cref{intervalclosed,real_add_closed}.
    $p / m + t\in\reals$ and $(p / m + t) + h\in\reals$ by \cref{real_add_closed}.
    $x + h \leq ((p / m + t) + h) + h$ by \cref{real_leq_add}.
    $((p / m + t) + h) + h = (p / m + t) + (h + h)$ by \cref{real_add_assoc}.
    $(p / m + t) + (h + h) = (p / m + t) + t$ by \cref{real_add_eq_subst}.
    $(p / m + t) + t = p / m + (t + t)$ by \cref{real_add_assoc}.
    $p / m + (t + t) = p / m + (1 / m)$ by \cref{real_add_eq_subst}.
    $(p + 1) / m = (p / m) + (1 / m)$ by \cref{real_div_add_same_denom}.
    $y \leq (p + 1) / m$ by \cref{real_lt_subst_left,real_lt_subst_right}.
    $x \leq y$ by \cref{real_lt_imp_leq,real_add_ident,real_lt_add,real_add_comm,real_lt_subst_left}.
    $p / m + t \leq y$ by \cref{real_leq_trans,real_add_closed}.
    $x \leq (p + 1) / m$ by \cref{real_leq_trans}.
    $0\leq x$ by \cref{intervalclosed}.
    $0 \leq y$ by \cref{real_leq_split,real_pos_add,real_lt_imp_leq,real_lt_subst_right,pair_eq_iff,real_add_ident,real_lt_subst_left}.
    Thus $y\in\intervalclosed{0}{1}$
        by \cref{real_div_leq_same_denom,naturalsplus_positive_real,real_pos_nonzero,real_div,real_mul_inverse,real_mul_comm,real_lt_subst_left,real_lt_subst_right,real_leq_trans,intervalclosed}.
    $q(x) = (p + 1) / m - x$ and $q(y) = (p + 1) / m - y$ by \cref{unit_interval_equal_mesh_adjacent_right_half_formula_49}.
    $\abs{((q(y) - q(x)) / (y - x))} = 1$
        by \cref{forward_descending_linear_secant_quotient_one_49}.
\end{proof}

% from §49 helper lemma 117: adjacent mesh interval backward unit quotient on the last quarter
\begin{lemma}\label{unit_interval_equal_mesh_adjacent_right_backward_sum_unit_quotient_49}
    Assume $A$ is the equal mesh of order $m$ on the unit interval.
    Assume $q\in\funs{\intervalclosed{0}{1}}{\reals}$.
    Assume for all $x\in\intervalclosed{0}{1}$ we have
        $q(x)$ is the distance from $x$ to $A$ in $\reals$ with respect to $\standardMetricR$.
    Assume $p\in\naturalsPlus$.
    Assume $p + 1 \leq m$.
    Let $h = ((1 / m) / (1 + 1)) / (1 + 1)$.
    Let $t = (1 / m) / (1 + 1)$.
    Assume $x\in\intervalclosed{0}{1}$.
    Suppose $(p / m + t) + h \leq x$.
    Suppose $x \leq (p + 1) / m$.
    Let $y = x - h$.
    Then $y\in\intervalclosed{0}{1}$ and
        $y = x - h$ and
        $\abs{((q(y) - q(x)) / (y - x))} = 1$.
\end{lemma}
\begin{proof}
    $m\in\naturalsPlus$ and $p\in\reals$ and $m\in\reals$ and $p + 1\in\reals$ by \cref{unit_interval_equal_mesh_49,real_naturals_subset_reals,subseteq,real_naturals_closed_succ}.
    $0\in\reals$ and $1\in\reals$ by \cref{real_add_ident,real_one_in_naturals,real_naturals_subset_reals,subseteq}.
    $1 / m\in\reals$ and $0 < 1 / m$ by \cref{naturalsplus_reciprocal_real,positive_reciprocal_naturalsplus}.
    $t\in\reals$ and $0 < t$ and $t + t = 1 / m$ by \cref{real_half_of_positive}.
    $h\in\reals$ and $0 < h$ and $h + h = t$ by \cref{real_half_of_positive}.
    $p / m\in\reals$ and $(p + 1) / m\in\reals$ by \cref{real_div_naturalsplus_real,real_naturals_closed_succ}.
    $x\in\reals$ and $y\in\reals$ by \cref{intervalclosed,real_sub,real_add_inverse,real_add_closed}.
    $p / m + t\in\reals$ and $(p / m + t) + h\in\reals$ by \cref{real_add_closed}.
    $\neg{h}\in\reals$ by \cref{real_add_inverse}.
    $p / m + t \leq (p / m + t) + h$ by \cref{real_lt_imp_leq,real_add_ident,real_lt_add,real_add_comm,real_lt_subst_left}.
    $p / m + t \leq x$ by \cref{real_leq_trans}.
    $((p / m + t) + h) + \neg{h} \leq x + \neg{h}$ by \cref{real_leq_add}.
    $((p / m + t) + h) + \neg{h} = p / m + t$
        by \cref{real_add_assoc,real_add_inverse,real_add_ident,real_add_comm}.
    $x + \neg{h} = x - h$ by \cref{real_sub}.
    $p / m + t \leq y$ by \cref{real_lt_subst_left,real_lt_subst_right}.
    $\neg{h} < 0$ by \cref{real_pos_imp_neg}.
    $\neg{h} + x < 0 + x$ by \cref{real_lt_add}.
    $\neg{h} + x = x + \neg{h}$ by \cref{real_add_comm}.
    $0 + x = x$ by \cref{real_add_ident}.
    $y < x$ and $y \leq x$ by \cref{real_lt_subst_left,real_lt_subst_right,real_lt_imp_leq}.
    $y \leq (p + 1) / m$ by \cref{real_leq_trans}.
    $0 \leq p / m$ by \cref{real_div_naturalsplus_nonnegative,real_lt_imp_leq,real_add_ident,naturalsplus_positive_real}.
    $p / m \leq p / m + t$ by \cref{real_lt_imp_leq,real_add_ident,real_lt_add,real_add_comm,real_lt_subst_left}.
    $0 \leq p / m + t$ by \cref{real_leq_trans}.
    $0 \leq y$ by \cref{real_leq_trans}.
    Thus $y\in\intervalclosed{0}{1}$ by \cref{intervalclosed,real_leq_trans}.
    $q(x) = (p + 1) / m - x$ and $q(y) = (p + 1) / m - y$ by \cref{unit_interval_equal_mesh_adjacent_right_half_formula_49}.
    $\abs{((q(y) - q(x)) / (y - x))} = 1$
        by \cref{backward_descending_linear_secant_quotient_one_49}.
\end{proof}

% from §49 helper lemma 119: positive forward and backward steps are distinct
\begin{lemma}\label{real_forward_backward_points_distinct_49}
    Assume $x\in\reals$.
    Assume $h\in\reals$.
    Suppose $0 < h$.
    Then $x + h \neq x - h$.
\end{lemma}
\begin{proof}
    Suppose not.
    $0\in\reals$ and $\neg{h}\in\reals$ and $x + h = x + \neg{h}$
        by \cref{real_add_ident,real_add_inverse,real_sub,real_lt_subst_right}.
    $h = \neg{h}$ and $h + h = 0$
        by \cref{real_add_comm,real_lt_subst_left,real_lt_subst_right,real_add_cancel,real_add_eq_subst,real_add_inverse}.
    Contradiction by \cref{real_pos_add,real_lt_irreflexive,real_lt_subst_right}.
\end{proof}

% from §49 helper lemma 118: adjacent mesh intervals have a quarter-step unit quotient
\begin{lemma}\label{unit_interval_equal_mesh_adjacent_interval_unit_quotient_49}
    Assume $A$ is the equal mesh of order $m$ on the unit interval.
    Assume $q\in\funs{\intervalclosed{0}{1}}{\reals}$.
    Assume for all $x\in\intervalclosed{0}{1}$ we have
        $q(x)$ is the distance from $x$ to $A$ in $\reals$ with respect to $\standardMetricR$.
    Assume $p\in\naturalsPlus$.
    Assume $p + 1 \leq m$.
    Let $h = ((1 / m) / (1 + 1)) / (1 + 1)$.
    Let $t = (1 / m) / (1 + 1)$.
    Assume $x\in\intervalclosed{0}{1}$.
    Suppose $p / m \leq x$.
    Suppose $x \leq (p + 1) / m$.
    Then there exists $y\in\intervalclosed{0}{1}$ such that
        either $y = x + h$ and
            $\abs{((q(y) - q(x)) / (y - x))} = 1$
        or $y = x - h$ and
            $\abs{((q(y) - q(x)) / (y - x))} = 1$.
\end{lemma}
\begin{proof}
    $m\in\naturalsPlus$ and $1 / m\in\reals$ and $0 < 1 / m$
        by \cref{unit_interval_equal_mesh_49,naturalsplus_reciprocal_real,positive_reciprocal_naturalsplus}.
    $t\in\reals$ and $0 < t$ and $t + t = 1 / m$ by \cref{real_half_of_positive}.
    $h\in\reals$ and $0 < h$ and $h + h = t$ by \cref{real_half_of_positive}.
    $p\in\reals$ and $p + 1\in\reals$ and $p / m\in\reals$ by
        \cref{real_naturals_subset_reals,subseteq,real_naturals_closed_succ,real_div_naturalsplus_real}.
    $x\in\reals$ by \cref{intervalclosed}.
    $p / m + h\in\reals$ and $p / m + t\in\reals$ and $(p / m + t) + h\in\reals$ by \cref{real_add_closed}.
    \begin{byCase}
    \caseOf{$x \leq p / m + h$.}
        Let $y = x + h$.
        $y\in\intervalclosed{0}{1}$ and $y = x + h$ and
            $\abs{((q(y) - q(x)) / (y - x))} = 1$
            by \cref{unit_interval_equal_mesh_adjacent_left_forward_unit_quotient_49}.
        Show that there exists $z\in\intervalclosed{0}{1}$ such that either $z = x + h$ and
            $\abs{((q(z) - q(x)) / (z - x))} = 1$ or $z = x - h$ and
            $\abs{((q(z) - q(x)) / (z - x))} = 1$.
        \begin{subproof}
            Trivial.
        \end{subproof}
    \caseOf{it is wrong that $x \leq p / m + h$.}
        $p / m + h < x$ and $p / m + h \leq x$
            by \cref{real_lt_trichotomy,real_lt_imp_leq,real_lt_subst_right}.
        \begin{byCase}
        \caseOf{$x \leq p / m + t$.}
            Let $y = x - h$.
            $y\in\intervalclosed{0}{1}$ and $y = x - h$ and
                $\abs{((q(y) - q(x)) / (y - x))} = 1$
                by \cref{unit_interval_equal_mesh_adjacent_left_backward_unit_quotient_49}.
            Show that there exists $z\in\intervalclosed{0}{1}$ such that either $z = x + h$ and
                $\abs{((q(z) - q(x)) / (z - x))} = 1$ or $z = x - h$ and
                $\abs{((q(z) - q(x)) / (z - x))} = 1$.
            \begin{subproof}
                Trivial.
            \end{subproof}
        \caseOf{it is wrong that $x \leq p / m + t$.}
            $p / m + t < x$ and $p / m + t \leq x$
                by \cref{real_lt_trichotomy,real_lt_imp_leq,real_lt_subst_right}.
            \begin{byCase}
            \caseOf{$x \leq (p / m + t) + h$.}
                Let $y = x + h$.
                $y\in\intervalclosed{0}{1}$ and $y = x + h$ and
                    $\abs{((q(y) - q(x)) / (y - x))} = 1$
                    by \cref{unit_interval_equal_mesh_adjacent_right_forward_sum_unit_quotient_49}.
                Show that there exists $z\in\intervalclosed{0}{1}$ such that either $z = x + h$ and
                    $\abs{((q(z) - q(x)) / (z - x))} = 1$ or $z = x - h$ and
                    $\abs{((q(z) - q(x)) / (z - x))} = 1$.
                \begin{subproof}
                    Trivial.
                \end{subproof}
            \caseOf{it is wrong that $x \leq (p / m + t) + h$.}
                $(p / m + t) + h < x$ and $(p / m + t) + h \leq x$
                    by \cref{real_lt_trichotomy,real_lt_imp_leq,real_lt_subst_right}.
                Let $y = x - h$.
                $y\in\intervalclosed{0}{1}$ and $y = x - h$ and
                    $\abs{((q(y) - q(x)) / (y - x))} = 1$
                    by \cref{unit_interval_equal_mesh_adjacent_right_backward_sum_unit_quotient_49}.
                Show that there exists $z\in\intervalclosed{0}{1}$ such that either $z = x + h$ and
                    $\abs{((q(z) - q(x)) / (z - x))} = 1$ or $z = x - h$ and
                    $\abs{((q(z) - q(x)) / (z - x))} = 1$.
                \begin{subproof}
                    Trivial.
                \end{subproof}
            \end{byCase}
        \end{byCase}
    \end{byCase}
\end{proof}

% from §49 helper lemma 120: equal mesh distance has a quarter-step unit quotient everywhere
\begin{lemma}\label{unit_interval_equal_mesh_global_unit_quotient_49}
    Assume $A$ is the equal mesh of order $m$ on the unit interval.
    Assume $q\in\funs{\intervalclosed{0}{1}}{\reals}$.
    Assume for all $x\in\intervalclosed{0}{1}$ we have
        $q(x)$ is the distance from $x$ to $A$ in $\reals$ with respect to $\standardMetricR$.
    Let $h = ((1 / m) / (1 + 1)) / (1 + 1)$.
    Let $t = (1 / m) / (1 + 1)$.
    Assume $x\in\intervalclosed{0}{1}$.
    Then there exists $y\in\intervalclosed{0}{1}$ such that
        either $y = x + h$ and
            $\abs{((q(y) - q(x)) / (y - x))} = 1$
        or $y = x - h$ and
            $\abs{((q(y) - q(x)) / (y - x))} = 1$.
\end{lemma}
\begin{proof}
    $m\in\naturalsPlus$ by \cref{unit_interval_equal_mesh_49}.
    Take $p$ such that
        either $p = 0$ and
            $x \leq 1 / m$
        or $p\in\naturalsPlus$ and
            $p + 1 \leq m$ and
            $p / m \leq x$ and
            $x \leq (p + 1) / m$ by \cref{unit_interval_equal_mesh_cover_49}.
    Thus there exists $z\in\intervalclosed{0}{1}$ such that
        either $z = x + h$ and
            $\abs{((q(z) - q(x)) / (z - x))} = 1$
        or $z = x - h$ and
            $\abs{((q(z) - q(x)) / (z - x))} = 1$
        by \cref{unit_interval_equal_mesh_first_interval_unit_quotient_49,unit_interval_equal_mesh_adjacent_interval_unit_quotient_49}.
\end{proof}

% from §49 helper lemma 121: four quarter mesh steps equal one mesh width
\begin{lemma}\label{quarter_mesh_step_fourfold_width_49}
    Assume $m\in\naturalsPlus$.
    Let $h = ((1 / m) / (1 + 1)) / (1 + 1)$.
    Then $(h + h) + (h + h) = 1 / m$.
\end{lemma}
\begin{proof}
    Follows by \cref{naturalsplus_reciprocal_real,positive_reciprocal_naturalsplus,real_half_of_positive,real_add_eq_subst}.
\end{proof}

% from §49 helper lemma 122: quotient by a quarter mesh step times the mesh width
\begin{lemma}\label{quarter_mesh_step_quotient_times_width_49}
    Assume $m\in\naturalsPlus$.
    Assume $\eta\in\reals$.
    Let $h = ((1 / m) / (1 + 1)) / (1 + 1)$.
    Let $b = \eta / h$.
    Then $b \rmul (1 / m) = (\eta + \eta) + (\eta + \eta)$.
\end{lemma}
\begin{proof}
    $1 / m\in\reals$ and $0 < 1 / m$ by \cref{naturalsplus_reciprocal_real,positive_reciprocal_naturalsplus}.
    $h\in\reals$ and $0 < h$ by \cref{quarter_mesh_step_positive_bounded_49}.
    $h\neq 0$ and $b\in\reals$ and $\inv{h}\in\reals$
        by \cref{real_pos_nonzero,real_div,real_mul_inverse,real_mul_closed}.
    $b = \eta \rmul \inv{h}$ and $b \rmul h = (\eta \rmul \inv{h}) \rmul h$
        by \cref{real_div,real_mul_eq_subst}.
    $(\eta \rmul \inv{h}) \rmul h = \eta \rmul (\inv{h} \rmul h)$ by \cref{real_mul_assoc}.
    $\inv{h} \rmul h = h \rmul \inv{h}$ and $h \rmul \inv{h} = 1$
        by \cref{real_mul_comm,real_mul_inverse}.
    $\eta \rmul 1 = \eta$ by \cref{real_mul_comm,real_mul_ident}.
    $b \rmul h = \eta$ by \cref{real_lt_subst_left,real_lt_subst_right,real_mul_eq_subst}.
    $(h + h) + (h + h) = 1 / m$ and $b \rmul (1 / m) = b \rmul ((h + h) + (h + h))$
        by \cref{quarter_mesh_step_fourfold_width_49,real_mul_eq_subst}.
    $b \rmul ((h + h) + (h + h)) = (b \rmul (h + h)) + (b \rmul (h + h))$ and
        $b \rmul (h + h) = (b \rmul h) + (b \rmul h)$ by \cref{real_distrib,real_add_closed}.
    Therefore $b \rmul (1 / m) = (\eta + \eta) + (\eta + \eta)$ by \cref{real_lt_subst_left,real_lt_subst_right,real_add_eq_subst}.
\end{proof}

% from §49 helper lemma 123: four eighths of a positive real are below the real
\begin{lemma}\label{real_four_eighths_lt_self_49}
    Assume $\epsilon\in\reals$.
    Suppose $0 < \epsilon$.
    Let $\eta = ((\epsilon / (1 + 1)) / (1 + 1)) / (1 + 1)$.
    Then $(\eta + \eta) + (\eta + \eta) < \epsilon$.
\end{lemma}
\begin{proof}
    Follows by \cref{real_half_of_positive,real_add_eq_subst,real_lt_subst_right,real_add_ident,real_lt_add}.
\end{proof}

% from §49 helper lemma 124: adding a strict and a weak inequality gives a strict inequality
\begin{lemma}\label{real_lt_add_leq_two_49}
    Assume $a\in\reals$.
    Assume $b\in\reals$.
    Assume $c\in\reals$.
    Assume $d\in\reals$.
    Suppose $a < b$.
    Suppose $c \leq d$.
    Then $a + c < b + d$.
\end{lemma}
\begin{proof}
    Follows by \cref{real_leq_split,real_lt_add_two,real_lt_add,real_add_eq_subst}.
\end{proof}

% from §49 helper lemma 125: adding two weak inequalities gives a weak inequality
\begin{lemma}\label{real_leq_add_two_49}
    Assume $a\in\reals$.
    Assume $b\in\reals$.
    Assume $c\in\reals$.
    Assume $d\in\reals$.
    Suppose $a \leq b$.
    Suppose $c \leq d$.
    Then $a + c \leq b + d$.
\end{lemma}
\begin{proof}
    Follows by \cref{real_leq_add,real_add_comm,real_leq_subst_left_local,real_leq_subst_right_local,real_leq_trans,real_add_closed}.
\end{proof}

% from §49 helper lemma 126: five eighths of a positive real are below the real
\begin{lemma}\label{real_five_eighths_lt_self_49}
    Assume $\epsilon\in\reals$.
    Suppose $0 < \epsilon$.
    Let $\eta = ((\epsilon / (1 + 1)) / (1 + 1)) / (1 + 1)$.
    Then $((\eta + \eta) + (\eta + \eta)) + \eta < \epsilon$.
\end{lemma}
\begin{proof}
    Follows by \cref{real_half_of_positive,real_add_eq_subst,real_lt_subst_right,real_lt_trans,real_add_ident,real_lt_add,real_add_comm}.
\end{proof}

% from §49 helper lemma 127: the chosen scalar perturbation amplitude is small
\begin{lemma}\label{quarter_mesh_scalar_amplitude_small_49}
    Assume $n\in\naturalsPlus$.
    Assume $\epsilon\in\reals$.
    Suppose $0 < \epsilon$.
    Assume $m\in\naturalsPlus$.
    Let $\eta = ((\epsilon / (1 + 1)) / (1 + 1)) / (1 + 1)$.
    Let $h = ((1 / m) / (1 + 1)) / (1 + 1)$.
    Let $b = \eta / h$.
    Let $K = (n + n) + 1$.
    Let $a = K + b$.
    Suppose $K \rmul (1 / m) \leq \eta$.
    Then $a \rmul (1 / m) < \epsilon$.
\end{lemma}
\begin{proof}
    $n\in\reals$ and $0 < n$ and $n + n\in\reals$ and $K\in\reals$
        by \cref{real_naturals_subset_reals,subseteq,naturalsplus_positive_real,real_add_closed,real_one_in_naturals}.
    $\eta\in\reals$ and $0 < \eta$ by \cref{real_half_of_positive}.
    $h\in\reals$ and $0 < h$ and $b\in\reals$
        by \cref{quarter_mesh_step_positive_bounded_49,real_div,real_mul_inverse,real_mul_closed,real_pos_nonzero}.
    $a\in\reals$ and $1 / m\in\reals$ and $K \rmul (1 / m)\in\reals$
        by \cref{real_add_closed,naturalsplus_reciprocal_real,real_mul_closed}.
    $a \rmul (1 / m) = (K + b) \rmul (1 / m)$ and $(K + b) \rmul (1 / m) = (K \rmul (1 / m)) + (b \rmul (1 / m))$ by \cref{real_mul_eq_subst,real_right_distrib}.
    $b \rmul (1 / m) = (\eta + \eta) + (\eta + \eta)$ and
        $(\eta + \eta) + (\eta + \eta)\in\reals$
        by \cref{quarter_mesh_step_quotient_times_width_49,real_add_closed}.
    $a \rmul (1 / m) = (K \rmul (1 / m)) + ((\eta + \eta) + (\eta + \eta))$ and
        $(K \rmul (1 / m)) + ((\eta + \eta) + (\eta + \eta)) = ((\eta + \eta) + (\eta + \eta)) + (K \rmul (1 / m))$
        by \cref{real_lt_subst_left,real_lt_subst_right,real_add_eq_subst,real_add_comm}.
    $((\eta + \eta) + (\eta + \eta)) \leq ((\eta + \eta) + (\eta + \eta))$ and $((\eta + \eta) + (\eta + \eta)) + (K \rmul (1 / m)) \leq ((\eta + \eta) + (\eta + \eta)) + \eta$
        by \cref{real_lt_subst_left,real_leq_add_two_49}.
    $((\eta + \eta) + (\eta + \eta)) + \eta < \epsilon$ by \cref{real_five_eighths_lt_self_49}.
    $a \rmul (1 / m) < \epsilon$ by \cref{real_leq_lt_trans,real_lt_subst_left,real_add_closed}.
\end{proof}

% from §49 helper lemma 128: choose a mesh width with a scaled reciprocal bound
\begin{lemma}\label{positive_mesh_step_scaled_bound_49}
    Assume $K\in\reals$.
    Suppose $0 < K$.
    Assume $\eta\in\reals$.
    Suppose $0 < \eta$.
    Then there exists $m\in\naturalsPlus$ such that $0 < 1 / m$ and $K \rmul (1 / m) \leq \eta$.
\end{lemma}
\begin{proof}
    $K\neq 0$ and $0\in\reals$ and $\eta / K\in\reals$ and $\inv{K}\in\reals$ and $0 < \inv{K}$
        by \cref{real_pos_nonzero,real_add_ident,real_div,real_mul_inverse,real_mul_closed,real_inv_pos}.
    $\eta / K = \eta \rmul \inv{K}$ and $0 \rmul \inv{K} < \eta \rmul \inv{K}$ and $0 \rmul \inv{K} = 0$
        by \cref{real_div,real_lt_mul_pos,real_mul_zero}.
    $0 < \eta / K$ by \cref{real_lt_subst_left,real_lt_subst_right}.
    Take $m\in\naturalsPlus$ such that $0 < 1 / m$ and $1 / m \leq \eta / K$
        by \cref{positive_mesh_step_below_bound_49}.
    $K \rmul (1 / m) \leq K \rmul (\eta / K)$ by \cref{real_leq_mul_pos_left_49,naturalsplus_reciprocal_real}.
    $K \rmul (\eta / K) = K \rmul (\eta \rmul \inv{K})$ and $K \rmul (\eta \rmul \inv{K}) = (K \rmul \inv{K}) \rmul \eta$
        by \cref{real_mul_eq_subst,real_mul_assoc,real_mul_comm}.
    $K \rmul \inv{K} = 1$ and $K \rmul (\eta / K) = \eta$
        by \cref{real_mul_inverse,real_mul_ident,real_lt_subst_left,real_lt_subst_right,real_mul_eq_subst}.
    $K \rmul (1 / m) \leq \eta$ by \cref{real_leq_subst_right_local}.
    Thus there exists $l\in\naturalsPlus$ such that $0 < 1 / l$ and $K \rmul (1 / l) \leq \eta$.
\end{proof}

% from §49 helper lemma 129: choose a mesh width satisfying one scaled and two direct bounds
\begin{lemma}\label{positive_mesh_step_scaled_and_two_bounds_49}
    Assume $K\in\reals$.
    Suppose $0 < K$.
    Assume $\eta\in\reals$.
    Suppose $0 < \eta$.
    Assume $c\in\reals$.
    Suppose $0 < c$.
    Assume $d\in\reals$.
    Suppose $0 < d$.
    Then there exists $m\in\naturalsPlus$ such that
        $0 < 1 / m$ and
        $K \rmul (1 / m) \leq \eta$ and
        $1 / m \leq c$ and
        $1 / m \leq d$.
\end{lemma}
\begin{proof}
    $K\neq 0$ and $0\in\reals$ and $\eta / K\in\reals$ and $\inv{K}\in\reals$ and $0 < \inv{K}$
        by \cref{real_pos_nonzero,real_add_ident,real_div,real_mul_inverse,real_mul_closed,real_inv_pos}.
    $\eta / K = \eta \rmul \inv{K}$ and $0 \rmul \inv{K} < \eta \rmul \inv{K}$ and $0 \rmul \inv{K} = 0$
        by \cref{real_div,real_lt_mul_pos,real_mul_zero}.
    $0 < \eta / K$ by \cref{real_lt_subst_left,real_lt_subst_right}.
    Take $m\in\naturalsPlus$ such that $0 < 1 / m$ and $1 / m \leq \eta / K$ and $1 / m \leq c$ and $1 / m \leq d$
        by \cref{positive_mesh_step_below_three_bounds_49}.
    $K \rmul (1 / m) \leq K \rmul (\eta / K)$ by \cref{real_leq_mul_pos_left_49,naturalsplus_reciprocal_real}.
    $K \rmul (\eta / K) = K \rmul (\eta \rmul \inv{K})$ and $K \rmul (\eta \rmul \inv{K}) = (K \rmul \inv{K}) \rmul \eta$
        by \cref{real_mul_eq_subst,real_mul_assoc,real_mul_comm}.
    $K \rmul \inv{K} = 1$ and $K \rmul (\eta / K) = \eta$
        by \cref{real_mul_inverse,real_mul_ident,real_lt_subst_left,real_lt_subst_right,real_mul_eq_subst}.
    $K \rmul (1 / m) \leq \eta$ by \cref{real_leq_subst_right_local}.
    Thus there exists $l\in\naturalsPlus$ such that $0 < 1 / l$ and $K \rmul (1 / l) \leq \eta$ and $1 / l \leq c$ and $1 / l \leq d$.
\end{proof}

% from §49 helper lemma 130: a quarter mesh step is strictly below the mesh width
\begin{lemma}\label{quarter_mesh_step_lt_width_49}
    Assume $m\in\naturalsPlus$.
    Let $h = ((1 / m) / (1 + 1)) / (1 + 1)$.
    Then $h < 1 / m$.
\end{lemma}
\begin{proof}
    Follows by \cref{naturalsplus_reciprocal_real,positive_reciprocal_naturalsplus,real_half_of_positive,real_add_ident,real_lt_add,real_lt_trans}.
\end{proof}

% from §49 helper lemma 131: robust scalar parameters from a positive quotient bound
\begin{lemma}\label{robust_scalar_parameters_from_bound_49}
    Assume $n\in\naturalsPlus$.
    Assume $b\in\reals$.
    Suppose $0 < b$.
    Let $K = (n + n) + 1$.
    Let $a = K + b$.
    Let $r = ((n + n) + b) + (1 / (1 + 1))$.
    Then $r\in\reals$ and $0 < r$ and $(n + n) + b < r$ and $r < a$.
\end{lemma}
\begin{proof}
    Follows by \cref{real_naturals_subset_reals,subseteq,naturalsplus_positive_real,real_one_in_naturals,real_zero_lt_one,real_add_closed,real_pos_add,real_half_of_positive,real_add_ident,real_lt_add,real_lt_subst_left,real_lt_subst_right,real_add_comm,real_add_assoc}.
\end{proof}

% from §49 helper lemma 132: scalar mesh-distance perturbations are strictly uniformly small
\begin{lemma}\label{unit_interval_equal_mesh_scalar_perturbation_uniform_strict_bound_49}
    Let $I = \intervalclosed{0}{1}$.
    Let $T_I = \subspaceTopology{\standardTopologyR}{I}$.
    Assume $A$ is the equal mesh of order $m$ on the unit interval.
    Assume $\rho$ is the uniform metric on functions from $I$ to $\reals$ with respect to $\standardMetricR$.
    Assume $f$ is continuous from $I$ to $\reals$ with respect to $T_I$ and $\standardTopologyR$.
    Assume $q\in\funs{I}{\reals}$.
    Assume for all $x\in I$ we have
        $q(x)$ is the distance from $x$ to $A$ in $\reals$ with respect to $\standardMetricR$.
    Assume $q$ is continuous from $I$ to $\reals$ with respect to $T_I$ and $\standardTopologyR$.
    Assume $a\in\reals$.
    Assume $\epsilon\in\reals$.
    Suppose $0 < a$.
    Suppose $a \rmul (1 / m) < \epsilon$.
    Suppose $\epsilon \leq 1$.
    Let $c = \{(x,a) \mid x\in I\}$.
    Let $s = \{(x,c(x) \rmul q(x)) \mid x\in I\}$.
    Let $g = \{(x,f(x) + s(x)) \mid x\in I\}$.
    Then $\rho((f,g)) < \epsilon$.
\end{lemma}
\begin{proof}
    Let $\delta = a \rmul (1 / m)$.
    $1 / m\in\reals$ and $1\in\reals$ by \cref{unit_interval_equal_mesh_49,naturalsplus_reciprocal_real,real_one_in_naturals,real_naturals_subset_reals,subseteq}.
    $\delta\in\reals$ and $a \rmul (1 / m) \leq \delta$ and $\delta < \epsilon$
        by \cref{real_mul_closed,real_lt_subst_left}.
    $\delta < 1$ and $\delta \leq 1$ by \cref{real_lt_leq_trans_local,real_lt_imp_leq}.
    $\rho((f,g)) \leq \delta$ by \cref{unit_interval_equal_mesh_scalar_perturbation_uniform_bound_49}.
    $g$ is continuous from $I$ to $\reals$ with respect to $T_I$ and $\standardTopologyR$ and
        $f\in\funs{I}{\reals}$ and $g\in\funs{I}{\reals}$
        by \cref{standard_topology_reals,standard_basis_is_basis,basis_topology_is_topology,intervalclosed,subseteq,subspace_topology_is_topology,real_function_add_scalar_multiple_continuous_49,continuous,funs}.
    Therefore $\rho((f,g)) < \epsilon$ by \cref{real_leq_lt_trans,uniform_metric_on_function_space,metric_on,funs_type_apply,times_tuple_intro}.
\end{proof}

% from §49 helper lemma 133: distance to an equal mesh with separated continuity witness
\begin{lemma}\label{unit_interval_equal_mesh_distance_continuous_separated_49}
    Assume $A$ is the equal mesh of order $m$ on the unit interval.
    Let $I = \intervalclosed{0}{1}$.
    Let $T_I = \subspaceTopology{\standardTopologyR}{I}$.
    Then there exists $q$ such that
        $q\in\funs{I}{\reals}$ and
        $q$ is continuous from $I$ to $\reals$ with respect to $T_I$ and $\standardTopologyR$ and
        for all $x\in I$ we have $q(x)$ is the distance from $x$ to $A$ in $\reals$ with respect to $\standardMetricR$.
\end{lemma}
\begin{proof}
    We have $A\subseteq I$ and $A\neq\emptyset$ by \cref{unit_interval_equal_mesh_basic_49}.
    $I\subseteq\reals$ by \cref{intervalclosed,subseteq}.
    Thus $A\subseteq\reals$ by \cref{subseteq_transitive}.
    Take $f$ such that
        $f$ is a function from $\reals$ to $\reals$ and
        for all $x\in\reals$ we have
            $f(x)$ is the distance from $x$ to $A$ in $\reals$ with respect to $\standardMetricR$ and
            $f$ is continuous from $\reals$ to $\reals$ with respect to $\standardTopologyR$ and $\standardTopologyR$
        by \cref{standard_metric_is_metric,standard_metric_topology,point_set_distance_continuous}.
    $f$ is continuous from $\reals$ to $\reals$ with respect to $\standardTopologyR$ and $\standardTopologyR$
        by \cref{real_add_ident,subseteq}.
    $f$ is a function and $f$ is a relation and $\dom{f} = \reals$
        by \cref{funs_elim,funs_is_relation}.
    Let $q = \restrl{f}{I}$.
    We have $\standardTopologyR$ is a topology on $\reals$
        by \cref{standard_topology_reals,standard_basis_is_basis,basis_topology_is_topology}.
    Thus $I$ is a subspace of $\reals$ with respect to $\standardTopologyR$ by \cref{subspace_of}.
    Show that $q$ is continuous from $I$ to $\reals$ with respect to $T_I$ and $\standardTopologyR$.
    \begin{subproof}
        Follows by \cref{subspace_of,continuous_restrict_domain}.
    \end{subproof}
    Show that $q\in\funs{I}{\reals}$.
    \begin{subproof}
        Follows by \cref{continuous,funs}.
    \end{subproof}
    Show that for all $x\in I$ we have $q(x)$ is the distance from $x$ to $A$ in $\reals$ with respect to $\standardMetricR$.
    \begin{subproof}
        Follows by \cref{intervalclosed,subseteq,restrl_of_function_apply,real_lt_subst_left}.
    \end{subproof}
    Therefore there exists $g$ such that $g\in\funs{I}{\reals}$ and $g$ is continuous from $I$ to $\reals$ with respect to $T_I$ and $\standardTopologyR$ and for all $x\in I$ we have $g(x)$ is the distance from $x$ to $A$ in $\reals$ with respect to $\standardMetricR$.
\end{proof}

% from §49 helper lemma 102: robust large-secant classes are dense in the continuous function space
\begin{lemma}\label{unit_interval_robust_large_secant_class_dense_49}
    Let $I = \intervalclosed{0}{1}$.
    Let $T_I = \subspaceTopology{\standardTopologyR}{I}$.
    Assume $C$ is the set of continuous functions from $I$ to $\reals$ with respect to $T_I$ and $\standardTopologyR$.
    Assume $\rho$ is the uniform metric on functions from $I$ to $\reals$ with respect to $\standardMetricR$.
    Let $e = \restrl{\rho}{C\times C}$.
    Assume $n\in\naturalsPlus$.
    Assume $V$ is the robust secant class above $n$ on $I$ with respect to $T_I$ and $\rho$.
    Then $V$ is dense in $C$ with respect to $\metricTopology{e}{C}$.
\end{lemma}
\begin{proof}
    $V\subseteq C$ by \cref{unit_interval_robust_large_secant_class_subset_continuous_49}.
    Show that for all $f\in C$ we have for all $\epsilon\in\reals$ if $0 < \epsilon$ then $\metricBall{e}{C}{f}{\epsilon}$ intersects $V$.
    \begin{subproof}
        Fix $f\in C$.
        Fix $\epsilon\in\reals$.
        Assume $0 < \epsilon$.
        Take $\epsilon_0\in\reals$ such that $0 < \epsilon_0$ and $\epsilon_0 \leq 1$ and $\epsilon_0 \leq \epsilon$ by \cref{positive_radius_below_one_49}.
        $f$ is continuous from $I$ to $\reals$ with respect to $T_I$ and $\standardTopologyR$ by \cref{continuous_function_set_elim_49}.
        Let $\eta = ((\epsilon_0 / (1 + 1)) / (1 + 1)) / (1 + 1)$.
        Let $d_I = \restrl{\standardMetricR}{I\times I}$.
        $\eta\in\reals$ and $0 < \eta$ by \cref{real_half_of_positive}.
        $f$ is uniformly continuous from $I$ to $\reals$ with respect to $d_I$ and $\standardMetricR$ by \cref{unit_interval_continuous_uniformly_continuous_49}.
        Take $\delta\in\reals$ such that $0 < \delta$ and for all $x_0,x_1\in I$ such that $d_I((x_0,x_1)) < \delta$ we have $\apply{\standardMetricR}{(\apply{f}{x_0},\apply{f}{x_1})} < \eta$ by \cref{uniformly_continuous}.
        Let $K = (n + n) + 1$.
        $n\in\reals$ and $0 < n$ and $1\in\reals$ and $0 < 1$ and $n + n\in\reals$ and $0 < n + n$ and $K\in\reals$ and $0 < K$
            by \cref{real_naturals_subset_reals,subseteq,naturalsplus_positive_real,real_one_in_naturals,real_zero_lt_one,real_add_closed,real_pos_add}.
        $1 / n\in\reals$ and $0 < 1 / n$ by \cref{naturalsplus_reciprocal_real,positive_reciprocal_naturalsplus}.
        Take $m\in\naturalsPlus$ such that $0 < 1 / m$ and $K \rmul (1 / m) \leq \eta$ and $1 / m \leq \delta$ and $1 / m \leq 1 / n$ by \cref{positive_mesh_step_scaled_and_two_bounds_49}.
        $1 / m\in\reals$ by \cref{naturalsplus_reciprocal_real}.
        Let $h = ((1 / m) / (1 + 1)) / (1 + 1)$.
        $h\in\reals$ and $0 < h$ and $h \leq 1 / m$ by \cref{quarter_mesh_step_positive_bounded_49}.
        $h \leq 1 / n$ by \cref{real_leq_trans,naturalsplus_reciprocal_real}.
        $h < 1 / m$ and $h < \delta$ by \cref{quarter_mesh_step_lt_width_49,real_lt_leq_trans_local}.
        Let $A = \{y\in\reals \mid \text{$y = 0$ or there exists $p\in\initialSegment{m}$ such that $y = p / m$}\}$.
        $A$ is the equal mesh of order $m$ on the unit interval by \cref{unit_interval_equal_mesh_49}.
        Take $q$ such that $q\in\funs{I}{\reals}$ and $q$ is continuous from $I$ to $\reals$ with respect to $T_I$ and $\standardTopologyR$ and for all $x\in I$ we have $q(x)$ is the distance from $x$ to $A$ in $\reals$ with respect to $\standardMetricR$ by \cref{unit_interval_equal_mesh_distance_continuous_separated_49}.
        For all $x,y\in I$ if $y = x + h$ or $y = x - h$ then $\apply{\standardMetricR}{(\apply{f}{y},\apply{f}{x})} < \eta$ by \cref{unit_interval_uniform_continuity_step_value_bound_49}.
        Let $b = \eta / h$.
        $b\in\reals$ and $0 < b$ and
            for all $x,y\in I$ if $y = x + h$ or $y = x - h$ then
                $\abs{((\apply{f}{y} - \apply{f}{x}) / (y - x))} < b$
            by \cref{unit_interval_strict_fixed_scale_quotient_bound_49}.
        For all $x\in I$ there exists $y\in I$ such that
            either $y = x + h$ and
                $\abs{((\apply{f}{y} - \apply{f}{x}) / (y - x))} < b$ and
                $\abs{((q(y) - q(x)) / (y - x))} = 1$
            or $y = x - h$ and
                $\abs{((\apply{f}{y} - \apply{f}{x}) / (y - x))} < b$ and
                $\abs{((q(y) - q(x)) / (y - x))} = 1$
            by \cref{unit_interval_equal_mesh_global_unit_quotient_49,subseteq}.
        Let $a = K + b$.
        Let $r = ((n + n) + b) + (1 / (1 + 1))$.
        $r\in\reals$ and $0 < r$ and $(n + n) + b < r$ and $r < a$
            by \cref{robust_scalar_parameters_from_bound_49}.
        Let $c = \{(x,a) \mid x\in I\}$.
        Let $w = \{(x,c(x) \rmul q(x)) \mid x\in I\}$.
        Let $g = \{(x,f(x) + w(x)) \mid x\in I\}$.
        $a\in\reals$ and $0 < a$ by \cref{real_add_closed,real_pos_add}.
        $a \rmul (1 / m) < \epsilon_0$ by \cref{quarter_mesh_scalar_amplitude_small_49}.
        $\rho((f,g)) < \epsilon_0$ by \cref{unit_interval_equal_mesh_scalar_perturbation_uniform_strict_bound_49}.
        $\standardTopologyR$ is a topology on $\reals$ and $I\subseteq\reals$ and $T_I$ is a topology on $I$
            by \cref{standard_topology_reals,standard_basis_is_basis,basis_topology_is_topology,intervalclosed,subseteq,subspace_topology_is_topology}.
        $g$ is continuous from $I$ to $\reals$ with respect to $T_I$ and $\standardTopologyR$
            by \cref{real_function_add_scalar_multiple_continuous_49}.
        $f\in\funs{I}{\reals}$ and $g\in\funs{I}{\reals}$ by \cref{continuous,funs}.
        $\rho((f,g))\in\reals$ by \cref{uniform_metric_on_function_space,metric_on,funs_type_apply,times_tuple_intro}.
        $\rho((f,g)) < \epsilon$ by \cref{real_lt_leq_trans_local}.
        Therefore $\metricBall{e}{C}{f}{\epsilon}$ intersects $V$
            by \cref{robust_class_ball_witness_from_scalar_perturbation_49}.
    \end{subproof}
    Therefore $V$ is dense in $C$ with respect to $\metricTopology{e}{C}$
        by \cref{continuous_function_space_ball_intersections_imply_dense_49}.
\end{proof}

% from §49 helper lemma 46: absolute value is invariant under negation
\begin{lemma}\label{real_abs_neg_value_49}
    Assume $x\in\reals$.
    Then $\abs{\neg{x}} = \abs{x}$.
\end{lemma}

% from §49 helper lemma 47: the robust secant perturbation radius is at most one
\begin{lemma}\label{unit_interval_robust_radius_leq_one_49}
    Assume $n\in\naturalsPlus$.
    Assume $h\in\reals$.
    Suppose $0 < h$.
    Suppose $h \leq 1 / n$.
    Then $(((h \rmul n) / (1 + 1)) / (1 + 1)) / (1 + 1) \leq 1$.
\end{lemma}
\begin{proof}
    $n\in\reals$ and $0 < n$ and $h \rmul n\in\reals$ and $0 < h \rmul n$
        by \cref{real_naturals_subset_reals,subseteq,naturalsplus_positive_real,real_mul_closed,real_mul_pos_pos_gt}.
    $1 / n\in\reals$ by \cref{naturalsplus_reciprocal_real}.
    $n \rmul h \leq n \rmul (1 / n)$ by \cref{real_leq_mul_pos_left_49,naturalsplus_reciprocal_real}.
    $n \rmul h = h \rmul n$ and $n \rmul (1 / n) = (1 / n) \rmul n$ by \cref{real_mul_comm}.
    $h \rmul n \leq (1 / n) \rmul n$
        by \cref{real_leq_subst_left_local,real_leq_subst_right_local}.
    $(1 / n) \rmul n = 1$
        by \cref{real_div,real_mul_assoc,real_mul_inverse,real_mul_comm,real_mul_ident,real_pos_nonzero}.
    $h \rmul n \leq 1$ by \cref{real_leq_subst_right_local}.
    $(((h \rmul n) / (1 + 1)) / (1 + 1)) / (1 + 1) \leq h \rmul n$
        by \cref{real_half_of_positive,real_add_ident,real_lt_add,real_lt_subst_left,real_lt_subst_right,real_lt_trans,real_lt_imp_leq}.
    $(((h \rmul n) / (1 + 1)) / (1 + 1)) / (1 + 1)\in\reals$ and $1\in\reals$
        by \cref{real_half_of_positive,real_one_in_naturals,real_naturals_subset_reals,subseteq}.
    Therefore $(((h \rmul n) / (1 + 1)) / (1 + 1)) / (1 + 1) \leq 1$ by \cref{real_leq_trans}.
\end{proof}

% from §49 helper lemma 48: differences of two small errors are bounded by the doubled radius
\begin{lemma}\label{real_abs_difference_lt_double_radius_49}
    Assume $a\in\reals$.
    Assume $b\in\reals$.
    Assume $r\in\reals$.
    Suppose $\abs{a} < r$.
    Suppose $\abs{b} < r$.
    Then $\abs{a - b} < r + r$.
\end{lemma}
\begin{proof}
    Follows by \cref{real_add_inverse,real_sub,real_add_closed,real_lt_subst_left,real_triangle_inequality,real_abs_neg_value_49,real_add_eq_subst,real_abs_in_reals,real_lt_add_two,real_leq_lt_trans}.
\end{proof}

% from §49 helper lemma 49: numerator bounds give quotient bounds by a positive scale
\begin{lemma}\label{real_abs_div_lt_from_mul_bound_49}
    Assume $a\in\reals$.
    Assume $h\in\reals$.
    Assume $n\in\reals$.
    Suppose $0 < h$.
    Suppose $0 < n$.
    Suppose $\abs{a} < h \rmul n$.
    Then $\abs{a / h} < n$.
\end{lemma}

% from §49 helper lemma 50: inverse of a negative real
\begin{lemma}\label{real_inv_neg_value_49}
    Assume $h\in\reals$.
    Suppose $h\neq 0$.
    Then $\inv{\neg{h}} = \neg{\inv{h}}$.
\end{lemma}

% from §49 helper lemma 51: numerator bounds give quotient bounds by a negative positive scale
\begin{lemma}\label{real_abs_div_neg_lt_from_mul_bound_49}
    Assume $a\in\reals$.
    Assume $h\in\reals$.
    Assume $n\in\reals$.
    Suppose $0 < h$.
    Suppose $0 < n$.
    Suppose $\abs{a} < h \rmul n$.
    Then $\abs{a / (\neg{h})} < n$.
\end{lemma}

% from §49 helper lemma 52: subtraction of quotients with a common denominator
\begin{lemma}\label{real_div_sub_same_denominator_49}
    Assume $a\in\reals$.
    Assume $b\in\reals$.
    Assume $h\in\reals$.
    Suppose $h\neq 0$.
    Then $(a / h) - (b / h) = (a - b) / h$.
\end{lemma}
\begin{proof}
    Follows by \cref{real_div,real_mul_inverse,real_sub,real_mul_neg_right,real_mul_comm,real_add_eq_subst,real_add_inverse,real_right_distrib,real_mul_eq_subst}.
\end{proof}

% from §49 helper lemma 53: secant numerator perturbation rearrangement
\begin{lemma}\label{secant_numerator_perturbation_rearrange_49}
    Assume $g_y\in\reals$.
    Assume $g_x\in\reals$.
    Assume $f_y\in\reals$.
    Assume $f_x\in\reals$.
    Then $(g_y - g_x) - (f_y - f_x) = (g_y - f_y) - (g_x - f_x)$.
\end{lemma}
\begin{proof}
    Follows by \cref{real_add_inverse,real_sub_swap_neg,real_sub,real_add_closed,real_add_assoc,real_add_comm,real_add_eq_subst}.
\end{proof}

% from §49 helper lemma 54: quotient differences at a positive scale are controlled by endpoint errors
\begin{lemma}\label{positive_scale_secant_quotient_difference_bound_49}
    Assume $g_y\in\reals$.
    Assume $g_x\in\reals$.
    Assume $f_y\in\reals$.
    Assume $f_x\in\reals$.
    Assume $h\in\reals$.
    Assume $n\in\reals$.
    Suppose $0 < h$.
    Suppose $0 < n$.
    Suppose $\abs{((g_y - f_y) - (g_x - f_x))} < h \rmul n$.
    Then $\abs{(((g_y - g_x) / h) - ((f_y - f_x) / h))} < n$.
\end{lemma}
\begin{proof}
    Follows by \cref{real_sub,real_add_inverse,real_add_closed,real_pos_nonzero,real_div_sub_same_denominator_49,secant_numerator_perturbation_rearrange_49,real_abs_div_lt_from_mul_bound_49}.
\end{proof}

% from §49 helper lemma 55: quotient differences at a negative scale are controlled by endpoint errors
\begin{lemma}\label{negative_scale_secant_quotient_difference_bound_49}
    Assume $g_y\in\reals$.
    Assume $g_x\in\reals$.
    Assume $f_y\in\reals$.
    Assume $f_x\in\reals$.
    Assume $h\in\reals$.
    Assume $n\in\reals$.
    Suppose $0 < h$.
    Suppose $0 < n$.
    Suppose $\abs{((g_y - f_y) - (g_x - f_x))} < h \rmul n$.
    Then $\abs{(((g_y - g_x) / (\neg{h})) - ((f_y - f_x) / (\neg{h})))} < n$.
\end{lemma}
\begin{proof}
    Follows by \cref{real_sub,real_add_inverse,real_add_closed,real_pos_nonzero,real_neg_neg,real_neg_zero,real_div_sub_same_denominator_49,secant_numerator_perturbation_rearrange_49,real_abs_div_neg_lt_from_mul_bound_49}.
\end{proof}

% from §49 helper lemma 56: lower bounds on absolute values survive small perturbations with a weak middle bound
\begin{lemma}\label{real_abs_perturbation_lower_bound_leq_49}
    Assume $a\in\reals$.
    Assume $b\in\reals$.
    Assume $n\in\reals$.
    Assume $s\in\reals$.
    Assume $r\in\reals$.
    Suppose $\abs{(a - b)} < s$.
    Suppose $n + s \leq r$.
    Suppose $r < \abs{a}$.
    Then $n < \abs{b}$.
\end{lemma}
\begin{proof}
    Follows by \cref{real_abs_in_reals,real_add_closed,real_lt_trichotomy,real_lt_imp_leq,real_lt_subst_right,real_leq_add,real_abs_sub_lt_bound,real_lt_leq_trans_local,real_lt_irreflexive,real_lt_trans}.
\end{proof}

% from §49 helper lemma 57: twice the robust secant radius is below the scale product
\begin{lemma}\label{unit_interval_robust_double_radius_lt_product_49}
    Assume $n\in\naturalsPlus$.
    Assume $h\in\reals$.
    Suppose $0 < h$.
    Then $((((h \rmul n) / (1 + 1)) / (1 + 1)) / (1 + 1)) + ((((h \rmul n) / (1 + 1)) / (1 + 1)) / (1 + 1)) < h \rmul n$.
\end{lemma}
\begin{proof}
    Follows by \cref{real_naturals_subset_reals,subseteq,naturalsplus_positive_real,real_mul_closed,real_mul_pos_pos_gt,real_half_of_positive,real_add_ident,real_lt_add,real_lt_trans}.
\end{proof}

% from §49 helper lemma 58: robust endpoint errors are below the scale product
\begin{lemma}\label{unit_interval_robust_endpoint_errors_product_bound_49}
    Assume $\rho$ is the uniform metric on functions from $I$ to $\reals$ with respect to $\standardMetricR$.
    Assume $f\in\funs{I}{\reals}$.
    Assume $g\in\funs{I}{\reals}$.
    Assume $n\in\naturalsPlus$.
    Assume $h\in\reals$.
    Assume $x\in I$.
    Assume $y\in I$.
    Suppose $0 < h$.
    Suppose $h \leq 1 / n$.
    Suppose $\rho((g,f)) < (((h \rmul n) / (1 + 1)) / (1 + 1)) / (1 + 1)$.
    Then $\abs{((\apply{g}{y} - \apply{f}{y}) - (\apply{g}{x} - \apply{f}{x}))} < h \rmul n$.
\end{lemma}
\begin{proof}
    Follows by \cref{real_naturals_subset_reals,subseteq,naturalsplus_positive_real,real_mul_closed,real_mul_pos_pos_gt,real_half_of_positive,unit_interval_robust_radius_leq_one_49,uniform_metric_pointwise_abs_lt_49,funs_type_apply,real_sub,real_add_inverse,real_add_closed,real_abs_difference_lt_double_radius_49,unit_interval_robust_double_radius_lt_product_49,real_abs_in_reals,real_lt_trans}.
\end{proof}

% from §49 helper lemma 59: positive robust secants transfer to nearby functions
\begin{lemma}\label{unit_interval_positive_robust_secant_transfer_49}
    Assume $I\subseteq\reals$.
    Assume $\rho$ is the uniform metric on functions from $I$ to $\reals$ with respect to $\standardMetricR$.
    Assume $f\in\funs{I}{\reals}$.
    Assume $g\in\funs{I}{\reals}$.
    Assume $n\in\naturalsPlus$.
    Assume $h\in\reals$.
    Assume $x\in I$.
    Assume $y\in I$.
    Suppose $0 < h$.
    Suppose $h \leq 1 / n$.
    Suppose $\rho((g,f)) < (((h \rmul n) / (1 + 1)) / (1 + 1)) / (1 + 1)$.
    Suppose $y = x + h$.
    Suppose $n + n < \abs{((\apply{f}{y} - \apply{f}{x}) / (y - x))}$.
    Then $n < \abs{((\apply{g}{y} - \apply{g}{x}) / (y - x))}$.
\end{lemma}
\begin{proof}
    $n\in\reals$ and $0 < n$ and $h \rmul n\in\reals$ by \cref{real_naturals_subset_reals,subseteq,naturalsplus_positive_real,real_mul_closed}.
    $\apply{g}{y}\in\reals$ and $\apply{g}{x}\in\reals$ and
    $\apply{f}{y}\in\reals$ and $\apply{f}{x}\in\reals$ by \cref{funs_type_apply}.
    $y\in\reals$ and $x\in\reals$ and $y - x = h$
        by \cref{subseteq,real_sub,real_add_assoc,real_add_inverse,real_add_ident,real_add_comm}.
    $y - x\neq 0$ by \cref{real_pos_nonzero,real_lt_subst_left}.
    $\abs{((\apply{g}{y} - \apply{f}{y}) - (\apply{g}{x} - \apply{f}{x}))} < h \rmul n$
        by \cref{unit_interval_robust_endpoint_errors_product_bound_49}.
    $\abs{(((\apply{g}{y} - \apply{g}{x}) / h) - ((\apply{f}{y} - \apply{f}{x}) / h))} < n$
        by \cref{positive_scale_secant_quotient_difference_bound_49}.
    $\abs{(((\apply{g}{y} - \apply{g}{x}) / (y - x)) - ((\apply{f}{y} - \apply{f}{x}) / (y - x)))} < n$
        by \cref{real_lt_subst_left,real_lt_subst_right}.
    $\apply{g}{y} - \apply{g}{x}\in\reals$ and $\apply{f}{y} - \apply{f}{x}\in\reals$
        by \cref{real_sub,real_add_inverse,real_add_closed}.
    Let $A = (\apply{f}{y} - \apply{f}{x}) / (y - x)$.
    Let $B = (\apply{g}{y} - \apply{g}{x}) / (y - x)$.
    $A\in\reals$ and $B\in\reals$ and $\abs{(A - B)} = \abs{(B - A)}$
        by \cref{real_div,real_mul_inverse,real_mul_closed,real_abs_sub_swap}.
    $\abs{(A - B)} < n$ by \cref{real_lt_subst_left}.
    $n + n\in\reals$ and $n + n \leq n + n$ by \cref{real_add_closed}.
    Therefore $n < \abs{B}$ by \cref{real_abs_perturbation_lower_bound_leq_49}.
\end{proof}

% from §49 helper lemma 60: negative robust secants transfer to nearby functions
\begin{lemma}\label{unit_interval_negative_robust_secant_transfer_49}
    Assume $I\subseteq\reals$.
    Assume $\rho$ is the uniform metric on functions from $I$ to $\reals$ with respect to $\standardMetricR$.
    Assume $f\in\funs{I}{\reals}$.
    Assume $g\in\funs{I}{\reals}$.
    Assume $n\in\naturalsPlus$.
    Assume $h\in\reals$.
    Assume $x\in I$.
    Assume $y\in I$.
    Suppose $0 < h$.
    Suppose $h \leq 1 / n$.
    Suppose $\rho((g,f)) < (((h \rmul n) / (1 + 1)) / (1 + 1)) / (1 + 1)$.
    Suppose $y = x - h$.
    Suppose $n + n < \abs{((\apply{f}{y} - \apply{f}{x}) / (y - x))}$.
    Then $n < \abs{((\apply{g}{y} - \apply{g}{x}) / (y - x))}$.
\end{lemma}
\begin{proof}
    $n\in\reals$ and $0 < n$ and $h \rmul n\in\reals$ by \cref{real_naturals_subset_reals,subseteq,naturalsplus_positive_real,real_mul_closed}.
    $\apply{g}{y}\in\reals$ and $\apply{g}{x}\in\reals$ and
        $\apply{f}{y}\in\reals$ and $\apply{f}{x}\in\reals$ by \cref{funs_type_apply}.
    $y\in\reals$ and $x\in\reals$ and $y - x = \neg{h}$
        by \cref{subseteq,real_sub,real_add_assoc,real_add_inverse,real_add_ident,real_add_comm}.
    $y - x\in\reals$ by \cref{real_lt_subst_left,real_add_inverse}.
    $y - x\neq 0$ by \cref{real_lt_subst_left,real_add_inverse,real_add_ident,real_add_cancel,real_pos_nonzero}.
    $\abs{((\apply{g}{y} - \apply{f}{y}) - (\apply{g}{x} - \apply{f}{x}))} < h \rmul n$
        by \cref{unit_interval_robust_endpoint_errors_product_bound_49}.
    $\abs{(((\apply{g}{y} - \apply{g}{x}) / (\neg{h})) - ((\apply{f}{y} - \apply{f}{x}) / (\neg{h})))} < n$
        by \cref{negative_scale_secant_quotient_difference_bound_49}.
    $\abs{(((\apply{g}{y} - \apply{g}{x}) / (y - x)) - ((\apply{f}{y} - \apply{f}{x}) / (y - x)))} < n$
        by \cref{real_lt_subst_left,real_lt_subst_right}.
    $\apply{g}{y} - \apply{g}{x}\in\reals$ and $\apply{f}{y} - \apply{f}{x}\in\reals$
        by \cref{real_sub,real_add_inverse,real_add_closed}.
    Let $A = (\apply{f}{y} - \apply{f}{x}) / (y - x)$.
    Let $B = (\apply{g}{y} - \apply{g}{x}) / (y - x)$.
    $A\in\reals$ and $B\in\reals$ and $\abs{(A - B)} = \abs{(B - A)}$
        by \cref{real_div,real_mul_inverse,real_mul_closed,real_abs_sub_swap}.
    $\abs{(A - B)} < n$ by \cref{real_lt_subst_left}.
    $n + n\in\reals$ and $n + n \leq n + n$ by \cref{real_add_closed}.
    Therefore $n < \abs{B}$ by \cref{real_abs_perturbation_lower_bound_leq_49}.
\end{proof}

% from §49 helper lemma 61: robust fixed-scale secants transfer to nearby functions
\begin{lemma}\label{unit_interval_robust_fixed_scale_secants_transfer_49}
    Assume $I\subseteq\reals$.
    Assume $\rho$ is the uniform metric on functions from $I$ to $\reals$ with respect to $\standardMetricR$.
    Assume $g\in\funs{I}{\reals}$.
    Assume $n\in\naturalsPlus$.
    Assume $h\in\reals$.
    Suppose $h \leq 1 / n$.
    Assume $f$ has secants above $n + n$ at scale $h$ on $I$.
    Suppose $\rho((g,f)) < (((h \rmul n) / (1 + 1)) / (1 + 1)) / (1 + 1)$.
    Then $g$ has secants above $n$ at scale $h$ on $I$.
\end{lemma}
\begin{proof}
    Follows by \cref{large_fixed_scale_secants_on_real_subset,real_naturals_subset_reals,subseteq,unit_interval_positive_robust_secant_transfer_49,unit_interval_negative_robust_secant_transfer_49}.
\end{proof}

% from §49 helper lemma 62: robust large-secant classes are contained in large-secant classes
\begin{lemma}\label{unit_interval_robust_large_secant_class_subset_large_49}
    Assume $C$ is the set of continuous functions from $I$ to $\reals$ with respect to $T_I$ and $\standardTopologyR$.
    Assume $\rho$ is the uniform metric on functions from $I$ to $\reals$ with respect to $\standardMetricR$.
    Assume $n\in\naturalsPlus$.
    Assume $V$ is the robust secant class above $n$ on $I$ with respect to $T_I$ and $\rho$.
    Assume $U$ is the large-secant class above $n$ on $I$ with respect to $T_I$.
    Then $V\subseteq U$.
\end{lemma}
\begin{proof}
    Follows by \cref{unit_interval_robust_large_secant_class_49,subseteq,large_fixed_scale_secants_on_real_subset,unit_interval_robust_fixed_scale_secants_transfer_49,unit_interval_large_secant_class_49}.
\end{proof}

% from §49 helper lemma 63: the half of a positive real is strictly smaller
\begin{lemma}\label{real_half_lt_self_positive_49}
    Assume $a\in\reals$.
    Suppose $0 < a$.
    Then $a / (1 + 1) < a$.
\end{lemma}

% from §49 helper lemma 64: the quarter of a positive real is strictly smaller
\begin{lemma}\label{real_quarter_lt_self_positive_49}
    Assume $a\in\reals$.
    Suppose $0 < a$.
    Then $(a / (1 + 1)) / (1 + 1) < a$.
\end{lemma}

% from §49 helper lemma 65: the eighth of a positive real is strictly smaller
\begin{lemma}\label{real_eighth_lt_self_positive_49}
    Assume $a\in\reals$.
    Suppose $0 < a$.
    Then $((a / (1 + 1)) / (1 + 1)) / (1 + 1) < a$.
\end{lemma}

% from §49 helper lemma 66: the half of a positive real is bounded by the real
\begin{lemma}\label{real_half_leq_self_positive_49}
    Assume $a\in\reals$.
    Suppose $0 < a$.
    Then $a / (1 + 1) \leq a$.
\end{lemma}
% from §49 helper lemma 67: the quarter of a positive real is bounded by the real
\begin{lemma}\label{real_quarter_leq_self_positive_49}
    Assume $a\in\reals$.
    Suppose $0 < a$.
    Then $(a / (1 + 1)) / (1 + 1) \leq a$.
\end{lemma}
% from §49 helper lemma 68: the eighth of a positive real is bounded by the real
\begin{lemma}\label{real_eighth_leq_self_positive_49}
    Assume $a\in\reals$.
    Suppose $0 < a$.
    Then $((a / (1 + 1)) / (1 + 1)) / (1 + 1) \leq a$.
\end{lemma}
% from §49 helper lemma 69: the quarter of a positive real is positive
\begin{lemma}\label{real_quarter_positive_49}
    Assume $a\in\reals$.
    Suppose $0 < a$.
    Then $0 < (a / (1 + 1)) / (1 + 1)$.
\end{lemma}
% from §49 helper lemma 70: the eighth of a positive real is positive
\begin{lemma}\label{real_eighth_positive_49}
    Assume $a\in\reals$.
    Suppose $0 < a$.
    Then $0 < ((a / (1 + 1)) / (1 + 1)) / (1 + 1)$.
\end{lemma}
% from §49 helper lemma 71: the quarter of a positive real is real
\begin{lemma}\label{real_quarter_in_reals_positive_49}
    Assume $a\in\reals$.
    Suppose $0 < a$.
    Then $(a / (1 + 1)) / (1 + 1)\in\reals$.
\end{lemma}
% from §49 helper lemma 72: the eighth of a positive real is real
\begin{lemma}\label{real_eighth_in_reals_positive_49}
    Assume $a\in\reals$.
    Suppose $0 < a$.
    Then $((a / (1 + 1)) / (1 + 1)) / (1 + 1)\in\reals$.
\end{lemma}
% from §49 helper lemma 73: two quarters make a half
\begin{lemma}\label{real_quarter_plus_quarter_eq_half_49}
    Assume $a\in\reals$.
    Suppose $0 < a$.
    Then $((a / (1 + 1)) / (1 + 1)) + ((a / (1 + 1)) / (1 + 1)) = a / (1 + 1)$.
\end{lemma}
% from §49 helper lemma 74: two eighths make a quarter
\begin{lemma}\label{real_eighth_plus_eighth_eq_quarter_49}
    Assume $a\in\reals$.
    Suppose $0 < a$.
    Then $(((a / (1 + 1)) / (1 + 1)) / (1 + 1)) + (((a / (1 + 1)) / (1 + 1)) / (1 + 1)) = (a / (1 + 1)) / (1 + 1)$.
\end{lemma}
% from §49 helper lemma 75: two quarters are below a positive real
\begin{lemma}\label{real_quarter_plus_quarter_lt_self_49}
    Assume $a\in\reals$.
    Suppose $0 < a$.
    Then $((a / (1 + 1)) / (1 + 1)) + ((a / (1 + 1)) / (1 + 1)) < a$.
\end{lemma}
% from §49 helper lemma 76: two eighths are below a positive real
\begin{lemma}\label{real_eighth_plus_eighth_lt_self_49}
    Assume $a\in\reals$.
    Suppose $0 < a$.
    Then $(((a / (1 + 1)) / (1 + 1)) / (1 + 1)) + (((a / (1 + 1)) / (1 + 1)) / (1 + 1)) < a$.
\end{lemma}
% from §49 helper lemma 77: two quarters are bounded by a positive real
\begin{lemma}\label{real_quarter_plus_quarter_leq_self_49}
    Assume $a\in\reals$.
    Suppose $0 < a$.
    Then $((a / (1 + 1)) / (1 + 1)) + ((a / (1 + 1)) / (1 + 1)) \leq a$.
\end{lemma}
% from §49 helper lemma 78: two eighths are bounded by a positive real
\begin{lemma}\label{real_eighth_plus_eighth_leq_self_49}
    Assume $a\in\reals$.
    Suppose $0 < a$.
    Then $(((a / (1 + 1)) / (1 + 1)) / (1 + 1)) + (((a / (1 + 1)) / (1 + 1)) / (1 + 1)) \leq a$.
\end{lemma}
% from §49 helper lemma 79: the half of a positive real below one is below one
\begin{lemma}\label{real_half_leq_one_from_leq_one_49}
    Assume $a\in\reals$.
    Suppose $0 < a$.
    Suppose $a \leq 1$.
    Then $a / (1 + 1) \leq 1$.
\end{lemma}
% from §49 helper lemma 80: the quarter of a positive real below one is below one
\begin{lemma}\label{real_quarter_leq_one_from_leq_one_49}
    Assume $a\in\reals$.
    Suppose $0 < a$.
    Suppose $a \leq 1$.
    Then $(a / (1 + 1)) / (1 + 1) \leq 1$.
\end{lemma}
% from §49 helper lemma 81: the eighth of a positive real below one is below one
\begin{lemma}\label{real_eighth_leq_one_from_leq_one_49}
    Assume $a\in\reals$.
    Suppose $0 < a$.
    Suppose $a \leq 1$.
    Then $((a / (1 + 1)) / (1 + 1)) / (1 + 1) \leq 1$.
\end{lemma}
% from §49 helper lemma 82: a quarter is below the corresponding half
\begin{lemma}\label{real_quarter_lt_half_positive_49}
    Assume $a\in\reals$.
    Suppose $0 < a$.
    Then $(a / (1 + 1)) / (1 + 1) < a / (1 + 1)$.
\end{lemma}

% from §49 helper lemma 83: an eighth is below the corresponding quarter
\begin{lemma}\label{real_eighth_lt_quarter_positive_49}
    Assume $a\in\reals$.
    Suppose $0 < a$.
    Then $((a / (1 + 1)) / (1 + 1)) / (1 + 1) < (a / (1 + 1)) / (1 + 1)$.
\end{lemma}

% from §49 helper lemma 84: an eighth is below the corresponding half
\begin{lemma}\label{real_eighth_lt_half_positive_49}
    Assume $a\in\reals$.
    Suppose $0 < a$.
    Then $((a / (1 + 1)) / (1 + 1)) / (1 + 1) < a / (1 + 1)$.
\end{lemma}

% from §49 helper lemma 85: a quarter is bounded by the corresponding half
\begin{lemma}\label{real_quarter_leq_half_positive_49}
    Assume $a\in\reals$.
    Suppose $0 < a$.
    Then $(a / (1 + 1)) / (1 + 1) \leq a / (1 + 1)$.
\end{lemma}

% from §49 helper lemma 86: an eighth is bounded by the corresponding quarter
\begin{lemma}\label{real_eighth_leq_quarter_positive_49}
    Assume $a\in\reals$.
    Suppose $0 < a$.
    Then $((a / (1 + 1)) / (1 + 1)) / (1 + 1) \leq (a / (1 + 1)) / (1 + 1)$.
\end{lemma}

% from §49 helper lemma 87: an eighth is bounded by the corresponding half
\begin{lemma}\label{real_eighth_leq_half_positive_49}
    Assume $a\in\reals$.
    Suppose $0 < a$.
    Then $((a / (1 + 1)) / (1 + 1)) / (1 + 1) \leq a / (1 + 1)$.
\end{lemma}

% from §49 helper lemma 88: the robust-class intersection is dense and nowhere differentiable
\begin{lemma}\label{unit_interval_dense_nowhere_differentiable_set_49}
    Let $I = \intervalclosed{0}{1}$.
    Let $T_I = \subspaceTopology{\standardTopologyR}{I}$.
    Assume $C$ is the set of continuous functions from $I$ to $\reals$ with respect to $T_I$ and $\standardTopologyR$.
    Assume $\rho$ is the uniform metric on functions from $I$ to $\reals$ with respect to $\standardMetricR$.
    Let $e = \restrl{\rho}{C\times C}$.
    Then there exists $D$ such that
        $D\subseteq C$ and
        $D$ is dense in $C$ with respect to $\metricTopology{e}{C}$ and
        for all $f\in D$ we have $f$ is nowhere differentiable on $I$.
\end{lemma}
\begin{proof}
    Let $T_C = \metricTopology{e}{C}$.
    $T_C$ is a topology on $C$ and $C$ is a Baire space with respect to $T_C$
        by \cref{unit_interval_continuous_function_space_metric_topology_49,unit_interval_continuous_function_space_baire_49}.
    Let $R = \{w\in\naturalsPlus\times\pow{C} \mid
        \text{$\snd{w}$ is the robust secant class above $\fst{w}$ on $I$ with respect to $T_I$ and $\rho$}\}$.
    $R$ is a relation by \cref{times_elem_is_tuple,relation}.
    Show that for all $n\in\naturalsPlus$ there exists $V$ such that $n\mathrel{R} V$.
    \begin{subproof}
        Fix $n\in\naturalsPlus$.
        Let $V = \{g\in\funs{I}{\reals} \mid \text{$g$ is continuous from $I$ to $\reals$ with respect to $T_I$ and $\standardTopologyR$ and
            there exist $f,h$ such that
                $f$ is continuous from $I$ to $\reals$ with respect to $T_I$ and $\standardTopologyR$ and
                $h\in\reals$ and
                $h \leq 1 / n$ and
                $f$ has secants above $n + n$ at scale $h$ on $I$ and
                $\rho((g,f)) < (((h \rmul n) / (1 + 1)) / (1 + 1)) / (1 + 1)$}\}$.
        $V$ is the robust secant class above $n$ on $I$ with respect to $T_I$ and $\rho$
            by \cref{unit_interval_robust_large_secant_class_49}.
        Follows by \cref{unit_interval_robust_large_secant_class_subset_continuous_49,times_tuple_intro,powerset_intro,fst_eq,snd_eq}.
    \end{subproof}
    Take $s$ such that $s$ is a function on $\naturalsPlus$ and
        for all $n\in\naturalsPlus$ we have $n\mathrel{R} s(n)$ by \cref{choice_function}.
    $s$ is a sequence in $\pow{C}$
        by \cref{choice_function,function_apply_elim,subseteq,times_tuple_elim,funs_intro,sequence_in}.
    For all $n\in\naturalsPlus$ we have
        $s(n)$ is open in $C$ with respect to $T_C$ and
        $s(n)$ is dense in $C$ with respect to $T_C$
        by \cref{choice_function,function_apply_elim,fst_eq,snd_eq,unit_interval_robust_large_secant_class_open_auto_49,unit_interval_robust_large_secant_class_dense_49}.
    Let $D = \inters{\img{s}{\naturalsPlus}}$.
    Show that $D$ is dense in $C$ with respect to $T_C$.
    \begin{subproof}
        Follows by \cref{baire_space_open_set_characterization}.
    \end{subproof}
    $\img{s}{\naturalsPlus}\subseteq\pow{C}$ by \cref{img_subseteq_ran,sequence_in,funs_ran,subseteq_transitive}.
    Show that $D\subseteq C$.
    \begin{subproof}
        Follows by \cref{real_one_in_naturals,sequence_in,funs_elim,function_apply_elim,img_iff,powerset_elim,inters_subseteq_elem,subseteq_transitive}.
    \end{subproof}
    Show that for all $f\in D$ we have $f$ is nowhere differentiable on $I$.
    \begin{subproof}
        Fix $f\in D$.
            Show that for all $n\in\naturalsPlus$ there exists $U$ such that
                $U$ is the large-secant class above $n$ on $I$ with respect to $T_I$ and
                $f\in U$.
            \begin{subproof}
                Fix $n\in\naturalsPlus$.
                Let $V = s(n)$.
                $V\in\img{s}{\naturalsPlus}$ by \cref{choice_function,function_apply_elim,img_iff}.
                $f\in V$ by \cref{inters_destr}.
                $V$ is the robust secant class above $n$ on $I$ with respect to $T_I$ and $\rho$
                    by \cref{choice_function,function_apply_elim,fst_eq,snd_eq}.
                Let $U = \{g\in\funs{I}{\reals} \mid \text{$g$ is continuous from $I$ to $\reals$ with respect to $T_I$ and $\standardTopologyR$ and
                    there exists $h\in\reals$ such that
                        $h \leq 1 / n$ and
                        $g$ has secants above $n$ at scale $h$ on $I$}\}$.
                $U$ is the large-secant class above $n$ on $I$ with respect to $T_I$
                    by \cref{unit_interval_large_secant_class_49}.
                Follows by \cref{unit_interval_robust_large_secant_class_subset_large_49,subseteq}.
            \end{subproof}
        Follows by \cref{subseteq,unit_interval_large_secant_classes_nowhere_differentiable_49}.
    \end{subproof}
\end{proof}

% from §49 Item 1: approximation by a continuous nowhere-differentiable function
\begin{theorem}\label{continuous_nowhere_differentiable_approximation_on_unit_interval}
    Let $I = \intervalclosed{0}{1}$.
    Let $T_I = \subspaceTopology{\standardTopologyR}{I}$.
    Assume $h$ is continuous from $I$ to $\reals$ with respect to $T_I$ and $\standardTopologyR$.
    Assume $\epsilon\in\reals$ and $0 < \epsilon$.
    Then there exists $g$ such that
        $g$ is continuous from $I$ to $\reals$ with respect to $T_I$ and $\standardTopologyR$ and
        $g$ is nowhere differentiable on $I$ and
        for all $x\in I$ we have $\abs{(\apply{h}{x} - \apply{g}{x})} < \epsilon$.
\end{theorem}
\begin{proof}
    Let $C = \{f\in\funs{I}{\reals} \mid \text{$f$ is continuous from $I$ to $\reals$ with respect to $T_I$ and $\standardTopologyR$}\}$.
    $C$ is the set of continuous functions from $I$ to $\reals$ with respect to $T_I$ and $\standardTopologyR$
        by \cref{continuous_function_set}.
    $I$ is inhabited and $\standardMetricR$ is a metric on $\reals$
        by \cref{intervalclosed,real_add_ident,real_zero_lt_one,real_lt_imp_leq,nonempty_inhabited,standard_metric_is_metric}.
    Take $\rho$ such that $\rho$ is the uniform metric on functions from $I$ to $\reals$ with respect to $\standardMetricR$
        by \cref{uniform_metric_on_function_space_exists_inhabited}.
    Let $e = \restrl{\rho}{C\times C}$.
    Take $D$ such that
        $D\subseteq C$ and
        $D$ is dense in $C$ with respect to $\metricTopology{e}{C}$ and
        for all $f\in D$ we have $f$ is nowhere differentiable on $I$
        by \cref{unit_interval_dense_nowhere_differentiable_set_49}.
    Follows by \cref{dense_nowhere_differentiable_set_gives_approximation_49}.
\end{proof}

\section{§50 Introduction to Dimension Theory}

% from §50 helper definition 1: at most finitely many members through a point
\begin{definition}\label{point_multiplicity_at_most_50}
    $x$ has multiplicity at most $m$ in $A$ iff
        $m\in\naturalsPlus$ and
        there exists no $F$ such that
            $F\subseteq A$ and
            there exists $g$ such that
                $g$ is a bijection from $\initialSegment{m + 1}$ to $F$ and
                for all $B\in F$ we have $x\in B$.
\end{definition}

% from §50 helper definition 2: at least finitely many members through a point
\begin{definition}\label{point_multiplicity_at_least_50}
    $x$ has multiplicity at least $m$ in $A$ iff
        $m\in\naturalsPlus$ and
        there exists $F$ such that
            $F\subseteq A$ and
            there exists $g$ such that
                $g$ is a bijection from $\initialSegment{m}$ to $F$ and
                for all $B\in F$ we have $x\in B$.
\end{definition}

% from §50 helper definition 3: collection has order at most a positive integer
\begin{definition}\label{covering_order_at_most_50}
    $A$ has order at most $m$ in $X$ iff
        $m\in\naturalsPlus$ and
        $A\subseteq\pow{X}$ and
        for all $x\in X$ we have $x$ has multiplicity at most $m$ in $A$.
\end{definition}

% from §50 Item 1: collection of subsets has order m
\begin{definition}\label{covering_order_50}
    $A$ has order $m$ in $X$ iff
        $m\in\naturalsPlus$ and
        $A\subseteq\pow{X}$ and
        for all $x\in X$ we have $x$ has multiplicity at most $m$ in $A$ and
        there exists $x\in X$ such that $x$ has multiplicity at least $m$ in $A$.
\end{definition}

% from §50 helper definition 4: order at most at points of a subspace
\begin{definition}\label{covering_order_at_most_on_subset_50}
    $A$ has order at most $m$ at points of $Y$ in $X$ iff
        $m\in\naturalsPlus$ and
        $Y\subseteq X$ and
        $A\subseteq\pow{X}$ and
        for all $y\in Y$ we have $y$ has multiplicity at most $m$ in $A$.
\end{definition}

% from §50 helper definition 5: topological dimension bound
\begin{definition}\label{topological_dimension_bound_50}
    $m$ is a topological dimension bound for $X$ with respect to $T$ iff
        $m\in\naturals$ and
        $T$ is a topology on $X$ and
        for all $A$ if
            $A$ is an open covering of $X$ with respect to $T$
        then there exists $B$ such that
            $B$ is an open refinement of $A$ on $X$ with respect to $T$ and
            $B$ has order at most $m + 1$ in $X$.
\end{definition}

% from §50 Item 2: finite dimensional space
\begin{definition}\label{finite_dimensional_space_50}
    $X$ is finite dimensional with respect to $T$ iff
        $T$ is a topology on $X$ and
        there exists $m\in\naturals$ such that $m$ is a topological dimension bound for $X$ with respect to $T$.
\end{definition}

% from §50 Item 3: topological dimension
\begin{definition}\label{topological_dimension_50}
    $m$ is the topological dimension of $X$ with respect to $T$ iff
        $m$ is a topological dimension bound for $X$ with respect to $T$ and
        for all $n$ such that $n$ is a topological dimension bound for $X$ with respect to $T$
            we have $m \leq n$.
\end{definition}

% from §50 helper lemma 1: intersecting a collection with a subspace does not increase its order
\begin{lemma}\label{subspace_intersections_preserve_order_bound_50}
    Assume $Y\subseteq X$.
    Suppose $B$ has order at most $m$ in $X$.
    Let $C = \{U\in\pow{Y} \mid \exists V\in B. U = Y\inter V\}$.
    Then $C$ has order at most $m$ in $Y$.
\end{lemma}
\begin{proof}
    $m\in\naturalsPlus$ and $B\subseteq\pow{X}$
        by \cref{covering_order_at_most_50}.
    $C\subseteq\pow{Y}$ by \cref{subseteq}.
    Show that for all $y\in Y$ we have $y$ has multiplicity at most $m$ in $C$.
    \begin{subproof}
        Fix $y\in Y$.
        Suppose not.
        Take $F,h$ such that
            $F\subseteq C$ and
            $h$ is a bijection from $\initialSegment{m + 1}$ to $F$ and
            for all $U\in F$ we have $y\in U$
            by \cref{point_multiplicity_at_most_50}.
        Let $R = \{w\in F\times B \mid
            \exists U\in F. \exists V\in B. w = (U,V) \land U = Y\inter V\}$.
        $R$ is a relation by \cref{relation,subseteq,times_elem_is_tuple}.
        For all $U\in F$ there exists $V$ such that $U\mathrel{R}V$
            by \cref{subseteq,times_tuple_intro}.
        Take $c$ such that
            $c$ is a function on $F$ and
            for all $U\in F$ we have $U\mathrel{R}\apply{c}{U}$
            by \cref{choice_function}.
        For all $U\in F$ we have
            $\apply{c}{U}\in B$ and $U = Y\inter\apply{c}{U}$
            by \cref{pair_eq_iff}.
        Show that $c$ is injective.
        \begin{subproof}
            It suffices to show that for all $U,V$ such that
                $U\in\dom{c}$ and $V\in\dom{c}$ and $\apply{c}{U} = \apply{c}{V}$
                we have $U = V$
                by \cref{injective_function}.
            Fix $U$.
            Fix $V$.
            Assume $U\in\dom{c}$ and $V\in\dom{c}$ and $\apply{c}{U} = \apply{c}{V}$.
            $U,V\in F$.
            $U = Y\inter\apply{c}{U}$ and $V = Y\inter\apply{c}{V}$.
            Thus $U = V$.
        \end{subproof}
        Let $H = \img{c}{F}$.
        $H\subseteq B$ by \cref{subseteq,img_iff,function_apply_iff}.
        For all $U\in\dom{c}$ we have $\apply{c}{U}\in H$
            by \cref{img_of_function_intro,inter_intro}.
        $c$ is a function to $H$.
        $c\in\funs{F}{H}$ by \cref{funs_intro}.
        $c\in\surj{F}{H}$ by \cref{funs_surj_iff,img_dom}.
        Thus $c$ is a bijection from $F$ to $H$ by \cref{bijections}.
        Let $q = c\circ h$.
        $q$ is a bijection from $\initialSegment{m + 1}$ to $H$
            by \cref{bijection_circ}.
        For all $V\in H$ we have $y\in V$
            by \cref{img_iff,function_apply_iff,inter_elim_right}.
        $y\in X$ by \cref{subseteq}.
        $y$ has multiplicity at most $m$ in $B$
            by \cref{covering_order_at_most_50}.
        Contradiction by \cref{point_multiplicity_at_most_50}.
    \end{subproof}
    Follows by \cref{covering_order_at_most_50}.
\end{proof}

% from §50 helper lemma 2: a subcollection does not have larger order
\begin{lemma}\label{subcollection_preserves_order_bound_50}
    Assume $C\subseteq B$.
    Suppose $B$ has order at most $m$ in $X$.
    Then $C$ has order at most $m$ in $X$.
\end{lemma}
\begin{proof}
    $m\in\naturalsPlus$ and $B\subseteq\pow{X}$
        by \cref{covering_order_at_most_50}.
    $C\subseteq\pow{X}$ by \cref{subseteq_transitive}.
    Show that for all $x\in X$ we have $x$ has multiplicity at most $m$ in $C$.
    \begin{subproof}
        Fix $x\in X$.
        Suppose not.
        $x$ has multiplicity at most $m$ in $B$
            by \cref{covering_order_at_most_50}.
        Take $F,g$ such that
            $F\subseteq C$ and
            $g$ is a bijection from $\initialSegment{m + 1}$ to $F$ and
            for all $A\in F$ we have $x\in A$
            by \cref{point_multiplicity_at_most_50}.
        $F\subseteq B$ by \cref{subseteq_transitive}.
        Contradiction by \cref{point_multiplicity_at_most_50}.
    \end{subproof}
    Follows by \cref{covering_order_at_most_50}.
\end{proof}

% from §50 Item 4: closed subspaces of finite dimensional spaces are finite dimensional
\begin{theorem}\label{closed_subspace_finite_dimensional_50}
    Assume $X$ is finite dimensional with respect to $T$.
    Assume $Y$ is a subspace of $X$ with respect to $T$.
    Assume $Y$ is closed in $X$ with respect to $T$.
    Let $T_Y = \subspaceTopology{T}{Y}$.
    Then $Y$ is finite dimensional with respect to $T_Y$.
\end{theorem}
\begin{proof}
    $T$ is a topology on $X$ and $Y\subseteq X$
        by \cref{finite_dimensional_space_50,subspace_of}.
    $T_Y$ is a topology on $Y$
        by \cref{subspace_topology_is_topology}.
    Take $m\in\naturals$ such that
        $m$ is a topological dimension bound for $X$ with respect to $T$
        by \cref{finite_dimensional_space_50}.
    Show that $m$ is a topological dimension bound for $Y$ with respect to $T_Y$.
    \begin{subproof}
        Show that for all $A$ if
            $A$ is an open covering of $Y$ with respect to $T_Y$
        then there exists $C$ such that
            $C$ is an open refinement of $A$ on $Y$ with respect to $T_Y$ and
            $C$ has order at most $m + 1$ in $Y$.
        \begin{subproof}
            Fix $A$.
            Assume $A$ is an open covering of $Y$ with respect to $T_Y$.
            $A\subseteq\pow{Y}$
                by \cref{open_covering,open_in,topology_on,subseteq,pow_iff,powerset_intro}.
            Let $E = \{V\in T \mid Y\inter V\in A\}$.
            Let $D = E\union\{X\setminus Y\}$.
            Show that for all $x\in X$ we have $x\in\unions{D}$.
            \begin{subproof}
                Fix $x\in X$.
                \begin{byCase}
                \caseOf{$x\in Y$.}
                    Take $U\in A$ such that $x\in U$
                        by \cref{open_covering,covering,subseteq,unions_iff}.
                    Take $V\in T$ such that $U = Y\inter V$
                        by \cref{open_covering,open_in,subspace_topology}.
                    $V\in E$ and $V\in D$.
                    $x\in V$ by \cref{inter_elim_right}.
                    Thus $x\in\unions{D}$ by \cref{unions_iff}.
                \caseOf{$x\notin Y$.}
                    $X\setminus Y\in D$ and $x\in X\setminus Y$
                        by \cref{union_intro_right,singleton_intro,setminus_intro}.
                    Thus $x\in\unions{D}$ by \cref{unions_iff}.
                \end{byCase}
            \end{subproof}
            $D$ is a covering of $X$ by \cref{subseteq,covering}.
            For all $V\in D$ we have $V$ is open in $X$ with respect to $T$
                by \cref{union_iff,singleton_iff,subseteq,open_in,closed_in}.
            $D$ is an open covering of $X$ with respect to $T$
                by \cref{open_covering}.
            Take $B$ such that
                $B$ is an open refinement of $D$ on $X$ with respect to $T$ and
                $B$ has order at most $m + 1$ in $X$
                by \cref{topological_dimension_bound_50}.
            Let $C_0 = \{U\in\pow{Y} \mid \exists V\in B. U = Y\inter V\}$.
            $C_0$ has order at most $m + 1$ in $Y$
                by \cref{subspace_intersections_preserve_order_bound_50}.
            Let $C = C_0\setminus\{\emptyset\}$.
            $C\subseteq C_0$ by \cref{setminus_subseteq}.
            $C$ has order at most $m + 1$ in $Y$
                by \cref{subcollection_preserves_order_bound_50}.
            $C\subseteq\pow{Y}$ by \cref{subseteq,covering_order_at_most_50}.
            Show that for all $U\in C$ there exists $W\in A$ such that $U\subseteq W$.
            \begin{subproof}
                Fix $U\in C$.
                $U\in C_0$ and $U\neq\emptyset$
                    by \cref{setminus_elim_left,setminus_elim_right,singleton_iff}.
                Take $V\in B$ such that $U = Y\inter V$.
                Take $W\in D$ such that $V\subseteq W$
                    by \cref{open_refinement_on,refinement_on}.
                \begin{byCase}
                \caseOf{$W\in E$.}
                    $Y\inter W\in A$.
                    $U\subseteq Y\inter W$
                        by \cref{inter_lower_left,inter_lower_right,subseteq_transitive,subseteq_inter_iff}.
                \caseOf{$W\notin E$.}
                    Show that there exists $G\in A$ such that $U\subseteq G$.
                    \begin{subproof}
                        Suppose not.
                        $W = X\setminus Y$ by \cref{union_iff,singleton_iff}.
                        $U\subseteq Y\inter(X\setminus Y)$
                            by \cref{inter_lower_left,inter_lower_right,subseteq_transitive,subseteq_inter_iff}.
                        $Y\inter(X\setminus Y) = \emptyset$ by \cref{setminus_disjoint}.
                        Thus $U = \emptyset$ by \cref{subseteq_emptyset_iff}.
                        Contradiction.
                    \end{subproof}
                \end{byCase}
            \end{subproof}
            $C$ is a refinement of $A$ on $Y$
                by \cref{refinement_on}.
            Show that for all $y\in Y$ we have $y\in\unions{C}$.
            \begin{subproof}
                Fix $y\in Y$.
                $y\in X$ by \cref{subseteq}.
                Take $V\in B$ such that $y\in V$
                    by \cref{open_refinement_on,covering,subseteq,unions_iff}.
                Let $U = Y\inter V$.
                $U\subseteq Y$ by \cref{inter_lower_left}.
                $U\in\pow{Y}$ by \cref{powerset_intro}.
                $U\in C_0$.
                $y\in U$ by \cref{inter_intro}.
                $U\neq\emptyset$ by \cref{emptyset}.
                $U\in C$ by \cref{setminus_intro,singleton_iff}.
                Thus $y\in\unions{C}$ by \cref{unions_iff}.
            \end{subproof}
            $C$ is a covering of $Y$ by \cref{subseteq,covering}.
            For all $U\in C$ we have $U$ is open in $Y$ with respect to $T_Y$
                by \cref{setminus_elim_left,subseteq,open_refinement_on,open_in,subspace_topology}.
            $C$ is an open refinement of $A$ on $Y$ with respect to $T_Y$
                by \cref{open_refinement_on}.
        \end{subproof}
        Follows by \cref{topological_dimension_bound_50}.
    \end{subproof}
    Follows by \cref{finite_dimensional_space_50}.
\end{proof}

% from §50 helper lemma 3: an ambient dimension bound is inherited by a closed subspace
\begin{lemma}\label{closed_subspace_inherits_dimension_bound_50}
    Assume $Y$ is a subspace of $X$ with respect to $T$.
    Assume $Y$ is closed in $X$ with respect to $T$.
    Let $T_Y = \subspaceTopology{T}{Y}$.
    Suppose $m$ is a topological dimension bound for $X$ with respect to $T$.
    Then $m$ is a topological dimension bound for $Y$ with respect to $T_Y$.
\end{lemma}
\begin{proof}
    $T$ is a topology on $X$ and $Y\subseteq X$ by \cref{subspace_of}.
    $T_Y$ is a topology on $Y$ by \cref{subspace_topology_is_topology}.
    Show that for all $A$ if
        $A$ is an open covering of $Y$ with respect to $T_Y$
    then there exists $C$ such that
        $C$ is an open refinement of $A$ on $Y$ with respect to $T_Y$ and
        $C$ has order at most $m + 1$ in $Y$.
    \begin{subproof}
        Fix $A$.
        Assume $A$ is an open covering of $Y$ with respect to $T_Y$.
        $A\subseteq\pow{Y}$
            by \cref{open_covering,open_in,topology_on,subseteq,pow_iff,powerset_intro}.
        Let $E = \{V\in T \mid Y\inter V\in A\}$.
        Let $D = E\union\{X\setminus Y\}$.
        Show that for all $x\in X$ we have $x\in\unions{D}$.
        \begin{subproof}
            Fix $x\in X$.
            \begin{byCase}
            \caseOf{$x\in Y$.}
                Take $U\in A$ such that $x\in U$
                    by \cref{open_covering,covering,subseteq,unions_iff}.
                Take $V\in T$ such that $U = Y\inter V$
                    by \cref{open_covering,open_in,subspace_topology}.
                $V\in E$ and $V\in D$.
                $x\in V$ by \cref{inter_elim_right}.
                Thus $x\in\unions{D}$ by \cref{unions_iff}.
            \caseOf{$x\notin Y$.}
                $X\setminus Y\in D$ and $x\in X\setminus Y$
                    by \cref{union_intro_right,singleton_intro,setminus_intro}.
                Thus $x\in\unions{D}$ by \cref{unions_iff}.
            \end{byCase}
        \end{subproof}
        $D$ is a covering of $X$ by \cref{subseteq,covering}.
        For all $V\in D$ we have $V$ is open in $X$ with respect to $T$
            by \cref{union_iff,singleton_iff,subseteq,open_in,closed_in}.
        $D$ is an open covering of $X$ with respect to $T$ by \cref{open_covering}.
        Take $B$ such that
            $B$ is an open refinement of $D$ on $X$ with respect to $T$ and
            $B$ has order at most $m + 1$ in $X$
            by \cref{topological_dimension_bound_50}.
        Let $C_0 = \{U\in\pow{Y} \mid \exists V\in B. U = Y\inter V\}$.
        $C_0$ has order at most $m + 1$ in $Y$
            by \cref{subspace_intersections_preserve_order_bound_50}.
        Let $C = C_0\setminus\{\emptyset\}$.
        $C\subseteq C_0$ by \cref{setminus_subseteq}.
        $C$ has order at most $m + 1$ in $Y$
            by \cref{subcollection_preserves_order_bound_50}.
        $C\subseteq\pow{Y}$ by \cref{subseteq,covering_order_at_most_50}.
        Show that for all $U\in C$ there exists $W\in A$ such that $U\subseteq W$.
        \begin{subproof}
            Fix $U\in C$.
            $U\in C_0$ and $U\neq\emptyset$
                by \cref{setminus_elim_left,setminus_elim_right,singleton_iff}.
            Take $V\in B$ such that $U = Y\inter V$.
            Take $W\in D$ such that $V\subseteq W$
                by \cref{open_refinement_on,refinement_on}.
            \begin{byCase}
            \caseOf{$W\in E$.}
                $Y\inter W\in A$.
                $U\subseteq Y\inter W$
                    by \cref{inter_lower_left,inter_lower_right,subseteq_transitive,subseteq_inter_iff}.
            \caseOf{$W\notin E$.}
                Show that there exists $G\in A$ such that $U\subseteq G$.
                \begin{subproof}
                    Suppose not.
                    $W = X\setminus Y$ by \cref{union_iff,singleton_iff}.
                    $U\subseteq Y\inter(X\setminus Y)$
                        by \cref{inter_lower_left,inter_lower_right,subseteq_transitive,subseteq_inter_iff}.
                    $Y\inter(X\setminus Y) = \emptyset$ by \cref{setminus_disjoint}.
                    Thus $U = \emptyset$ by \cref{subseteq_emptyset_iff}.
                    Contradiction.
                \end{subproof}
            \end{byCase}
        \end{subproof}
        $C$ is a refinement of $A$ on $Y$ by \cref{refinement_on}.
        Show that for all $y\in Y$ we have $y\in\unions{C}$.
        \begin{subproof}
            Fix $y\in Y$.
            $y\in X$ by \cref{subseteq}.
            Take $V\in B$ such that $y\in V$
                by \cref{open_refinement_on,covering,subseteq,unions_iff}.
            Let $U = Y\inter V$.
            $U\subseteq Y$ by \cref{inter_lower_left}.
            $U\in\pow{Y}$ by \cref{powerset_intro}.
            $U\in C_0$.
            $y\in U$ by \cref{inter_intro}.
            $U\neq\emptyset$ by \cref{emptyset}.
            $U\in C$ by \cref{setminus_intro,singleton_iff}.
            Thus $y\in\unions{C}$ by \cref{unions_iff}.
        \end{subproof}
        $C$ is a covering of $Y$ by \cref{subseteq,covering}.
        For all $U\in C$ we have $U$ is open in $Y$ with respect to $T_Y$
            by \cref{setminus_elim_left,subseteq,open_refinement_on,open_in,subspace_topology}.
        $C$ is an open refinement of $A$ on $Y$ with respect to $T_Y$
            by \cref{open_refinement_on}.
    \end{subproof}
    Follows by \cref{topological_dimension_bound_50}.
\end{proof}

% from §50 Item 4: dimension of a closed subspace is bounded by the ambient dimension
\begin{theorem}\label{closed_subspace_dimension_leq_50}
    Assume $X$ is finite dimensional with respect to $T$.
    Assume $Y$ is a subspace of $X$ with respect to $T$.
    Assume $Y$ is closed in $X$ with respect to $T$.
    Let $T_Y = \subspaceTopology{T}{Y}$.
    Suppose $m$ is the topological dimension of $X$ with respect to $T$.
    Suppose $n$ is the topological dimension of $Y$ with respect to $T_Y$.
    Then $n \leq m$.
\end{theorem}
\begin{proof}
    $m$ is a topological dimension bound for $X$ with respect to $T$
        by \cref{topological_dimension_50}.
    $m$ is a topological dimension bound for $Y$ with respect to $T_Y$
        by \cref{closed_subspace_inherits_dimension_bound_50}.
    Follows by \cref{topological_dimension_50}.
\end{proof}

% from §50 helper lemma 4: the successor of a natural number is positive natural
\begin{lemma}\label{natural_successor_positive_50}
    Suppose $m\in\naturals$.
    Then $m + 1\in\naturalsPlus$.
\end{lemma}
\begin{proof}
    \begin{byCase}
    \caseOf{$m\in\naturalsPlus$.}
        Follows by \cref{real_naturals_closed_succ}.
    \caseOf{$m\notin\naturalsPlus$.}
        $m = 0$ by \cref{naturals,union_iff,singleton_iff}.
        $1\in\reals$ by \cref{real_one_in_naturals,real_naturals_subset_reals,subseteq}.
        $m + 1 = 1$ by \cref{real_add_ident}.
        Follows by \cref{real_one_in_naturals}.
    \end{byCase}
\end{proof}

% from §50 helper lemma 5: monotonicity of positive-natural initial segments
\begin{lemma}\label{initial_segment_monotone_50}
    Suppose $a\in\naturalsPlus$ and $b\in\naturalsPlus$.
    Suppose $a\leq b$.
    Then $\initialSegment{a}\subseteq\initialSegment{b}$.
\end{lemma}
\begin{proof}
    Show that for all $i\in\initialSegment{a}$ we have $i\in\initialSegment{b}$.
    \begin{subproof}
        Fix $i\in\initialSegment{a}$.
        $i\in\naturalsPlus$ and $i\leq a$
            by \cref{initial_segment_naturalsplus}.
        $i,a,b\in\reals$ by \cref{real_naturals_subset_reals,subseteq}.
        $i\leq b$ by \cref{real_leq_trans}.
        Follows by \cref{initial_segment_naturalsplus}.
    \end{subproof}
    Follows by \cref{subseteq}.
\end{proof}

% from §50 helper lemma 6: monotonicity of point multiplicity bounds
\begin{lemma}\label{point_multiplicity_bound_monotone_50}
    Suppose $a\in\naturalsPlus$ and $b\in\naturalsPlus$.
    Suppose $a\leq b$.
    Suppose $x$ has multiplicity at most $a$ in $A$.
    Then $x$ has multiplicity at most $b$ in $A$.
\end{lemma}
\begin{proof}
    Suppose not.
    Take $F,h$ such that
        $F\subseteq A$ and
        $h$ is a bijection from $\initialSegment{b + 1}$ to $F$ and
        for all $C\in F$ we have $x\in C$
        by \cref{point_multiplicity_at_most_50}.
    $a + 1\in\naturalsPlus$ and $b + 1\in\naturalsPlus$
        by \cref{real_naturals_closed_succ}.
    $a,b,1\in\reals$
        by \cref{real_naturals_subset_reals,real_one_in_naturals,subseteq}.
    $a + 1\leq b + 1$ by \cref{real_leq_add}.
    Let $I = \initialSegment{a + 1}$.
    Let $J = \initialSegment{b + 1}$.
    $I\subseteq J$ by \cref{initial_segment_monotone_50}.
    $h\in\funs{J}{F}$ and $h$ is injective
        by \cref{bijections,bijections_to_funs}.
    $h$ is a function and $\dom{h} = J$ by \cref{funs_elim}.
    Let $r = \restrl{h}{I}$.
    $r$ is a function by \cref{restrl_of_function_is_function}.
    $\dom{r} = \dom{h}\inter I$
        by \cref{dom_of_restrl_eq_inter}.
    $\dom{r} = I$ by \cref{inter_absorb_supseteq_left}.
    $r$ is a function on $I$.
    $r$ is injective by \cref{restrl_injective}.
    Let $H = \img{r}{I}$.
    Show that $H\subseteq F$.
    \begin{subproof}
        It suffices to show that for all $C\in H$ we have $C\in F$.
        Fix $C\in H$.
        Take $i$ such that $i\in\dom{r}\inter I$ and $C = \apply{r}{i}$
            by \cref{img_of_function_elim}.
        $i\in I$ by \cref{inter_elim_right}.
        $\apply{r}{i} = C$.
        $i\in J$ by \cref{subseteq}.
        $\apply{r}{i} = \apply{h}{i}$
            by \cref{restrl_of_function_apply}.
        $\apply{h}{i}\in F$ by \cref{funs_type_apply}.
        Thus $C\in F$.
    \end{subproof}
    For all $i\in\dom{r}$ we have $\apply{r}{i}\in H$
        by \cref{img_of_function_intro,inter_intro}.
    $r$ is a function to $H$.
    $r\in\funs{I}{H}$ by \cref{funs_intro}.
    $r\in\surj{I}{H}$ by \cref{funs_surj_iff,img_dom}.
    Thus $r$ is a bijection from $I$ to $H$ by \cref{bijections}.
    $H\subseteq A$ by \cref{subseteq_transitive}.
    For all $C\in H$ we have $x\in C$ by \cref{subseteq}.
    Contradiction by \cref{point_multiplicity_at_most_50}.
\end{proof}

% from §50 helper lemma 7: monotonicity of collection order bounds
\begin{lemma}\label{covering_order_bound_monotone_50}
    Suppose $a\in\naturalsPlus$ and $b\in\naturalsPlus$.
    Suppose $a\leq b$.
    Suppose $A$ has order at most $a$ in $X$.
    Then $A$ has order at most $b$ in $X$.
\end{lemma}
\begin{proof}
    Follows by \cref{covering_order_at_most_50,point_multiplicity_bound_monotone_50}.
\end{proof}

% from §50 helper lemma 8: monotonicity of topological dimension bounds
\begin{lemma}\label{topological_dimension_bound_monotone_50}
    Suppose $a\in\naturals$ and $b\in\naturals$.
    Suppose $a\leq b$.
    Suppose $a$ is a topological dimension bound for $X$ with respect to $T$.
    Then $b$ is a topological dimension bound for $X$ with respect to $T$.
\end{lemma}
\begin{proof}
    Follows by
        \cref{topological_dimension_bound_50,natural_successor_positive_50,naturals,union_iff,singleton_iff,real_naturals_subset_reals,subseteq,real_add_ident,real_one_in_naturals,real_leq_add,covering_order_bound_monotone_50}.
\end{proof}

% from §50 helper lemma 9: controlled images preserve order at points of a subspace
\begin{lemma}\label{controlled_image_preserves_order_on_subset_50}
    Assume $Y\subseteq X$.
    Suppose $m\in\naturalsPlus$.
    Assume $s$ is a function on $B$.
    Assume for all $V\in B$ we have $\apply{s}{V}\subseteq X$.
    Assume for all $V,y$ such that
        $V\in B$ and $y\in Y$ and $y\in\apply{s}{V}$
        we have $y\in V$.
    Suppose $B$ has order at most $m$ in $Y$.
    Then $\img{s}{B}$ has order at most $m$ at points of $Y$ in $X$.
\end{lemma}
\begin{proof}
    $\img{s}{B}\subseteq\pow{X}$
        by \cref{img_iff,function_apply_iff,powerset_intro,subseteq}.
    Show that for all $y\in Y$ we have
        $y$ has multiplicity at most $m$ in $\img{s}{B}$.
    \begin{subproof}
        Fix $y\in Y$.
        Suppose not.
        Take $F,h$ such that
            $F\subseteq\img{s}{B}$ and
            $h$ is a bijection from $\initialSegment{m + 1}$ to $F$ and
            for all $U\in F$ we have $y\in U$
            by \cref{point_multiplicity_at_most_50}.
        Let $I = \initialSegment{m + 1}$.
        $h\in\funs{I}{F}$ and $h$ is injective
            by \cref{bijections,bijections_to_funs}.
        Let $R = \{w\in I\times B \mid \apply{s}{\snd{w}} = \apply{h}{\fst{w}}\}$.
        $R$ is a relation by \cref{relation,subseteq,times_elem_is_tuple}.
        For all $i\in I$ there exists $V$ such that $i\mathrel{R}V$
            by \cref{funs_type_apply,subseteq,img_of_function_elim,inter_elim_right,times_tuple_intro,fst_eq,snd_eq}.
        Take $g$ such that
            $g$ is a function on $I$ and
            for all $i\in I$ we have $i\mathrel{R}\apply{g}{i}$
            by \cref{choice_function}.
        For all $i\in I$ we have
            $\apply{g}{i}\in B$ and
            $\apply{s}{\apply{g}{i}} = \apply{h}{i}$
            by \cref{choice_function,times_tuple_elim,fst_eq,snd_eq}.
        $g$ is injective
            by \cref{injective_function,choice_function,funs_elim,funs_is_function}.
        Let $H = \img{g}{I}$.
        $g$ is a bijection from $I$ to $H$
            by \cref{img_of_function_intro,choice_function,inter_intro,funs_intro,funs_surj_iff,img_dom,bijections}.
        $H\subseteq B$
            by \cref{img_iff,function_apply_iff,subseteq}.
        For all $V\in H$ we have $y\in V$
            by \cref{img_of_function_elim,inter_elim_right,funs_type_apply,subseteq,choice_function,times_tuple_elim,fst_eq,snd_eq}.
        Contradiction by
            \cref{covering_order_at_most_50,point_multiplicity_at_most_50}.
    \end{subproof}
    Follows by \cref{covering_order_at_most_on_subset_50}.
\end{proof}

% from §50 helper lemma 10: bounded-order refinement at a closed subspace
\begin{lemma}\label{closed_subspace_pointwise_refinement_50}
    Assume $Y$ is a subspace of $X$ with respect to $T$.
    Assume $Y$ is closed in $X$ with respect to $T$.
    Let $T_Y = \subspaceTopology{T}{Y}$.
    Suppose $m$ is a topological dimension bound for $Y$ with respect to $T_Y$.
    Suppose $A$ is an open covering of $X$ with respect to $T$.
    Then there exists $C$ such that
        $C$ is an open refinement of $A$ on $X$ with respect to $T$ and
        $C$ has order at most $m + 1$ at points of $Y$ in $X$.
\end{lemma}
\begin{proof}
    $T$ is a topology on $X$ and $Y\subseteq X$
        by \cref{subspace_of}.
    $T_Y$ is a topology on $Y$
        by \cref{subspace_topology_is_topology}.
    $A\subseteq\pow{X}$
        by \cref{open_covering,open_in,topology_on,subseteq,pow_iff,powerset_intro}.
    Let $A_Y = \{V\in\pow{Y} \mid \exists U\in A. V = Y\inter U\}$.
    For all $y\in Y$ we have $y\in\unions{A_Y}$
        by \cref{subseteq,open_covering,covering,unions_iff,inter_lower_left,powerset_intro,inter_intro}.
    $A_Y$ is an open covering of $Y$ with respect to $T_Y$
        by \cref{subseteq,covering,open_covering,open_in,subspace_topology}.
    Take $B$ such that
        $B$ is an open refinement of $A_Y$ on $Y$ with respect to $T_Y$ and
        $B$ has order at most $m + 1$ in $Y$
        by \cref{topological_dimension_bound_50}.
    Let $E = \{W\in T \mid \exists U\in A. W\subseteq U\}$.
    Let $R = \{w\in B\times E \mid
        \exists V\in B.\exists W\in E.
        w = (V,W)\land V = Y\inter W\}$.
    $R$ is a relation by \cref{relation,subseteq,times_elem_is_tuple}.
    Show that for all $V\in B$ there exists $W$ such that $V\mathrel{R}W$.
    \begin{subproof}
        Fix $V\in B$.
        Take $D\in A_Y$ such that $V\subseteq D$
            by \cref{open_refinement_on,refinement_on}.
        Take $U_0\in A$ such that $D = Y\inter U_0$.
        Take $U\in T$ such that $V = Y\inter U$
            by \cref{open_refinement_on,open_in,subspace_topology}.
        Let $W = U\inter U_0$.
        $W\in T$ by \cref{open_covering,open_in,topology_on}.
        $Y\inter W = V$
            by \cref{inter_lower_right,subseteq_transitive,inter_assoc,inter_absorb_supseteq_right}.
        Thus $V\mathrel{R}W$
            by \cref{inter_lower_right,times_tuple_intro,fst_eq,snd_eq}.
    \end{subproof}
    Take $s$ such that
        $s$ is a function on $B$ and
        for all $V\in B$ we have $V\mathrel{R}\apply{s}{V}$
        by \cref{choice_function}.
    For all $V\in B$ we have
        $\apply{s}{V}\in E$ and
        $V = Y\inter\apply{s}{V}$
        by \cref{choice_function,times_tuple_elim,fst_eq,snd_eq}.
    For all $V\in B$ we have
        $\apply{s}{V}\in T$ and
        there exists $U\in A$ such that $\apply{s}{V}\subseteq U$
        by \cref{subseteq}.
    For all $V\in B$ we have $\apply{s}{V}\subseteq X$
        by \cref{topology_on,subseteq,powerset_elim}.
    For all $V,y$ such that
        $V\in B$ and $y\in Y$ and $y\in\apply{s}{V}$
        we have $y\in V$
        by \cref{choice_function,times_tuple_elim,fst_eq,snd_eq,inter_intro}.
    $m + 1\in\naturalsPlus$
        by \cref{topological_dimension_bound_50,natural_successor_positive_50}.
    Let $C_0 = \img{s}{B}$.
    $C_0$ has order at most $m + 1$ at points of $Y$ in $X$
        by \cref{controlled_image_preserves_order_on_subset_50}.
    Let $Q = \{W\in\pow{X} \mid
        \exists U\in A. W = U\inter(X\setminus Y)\}$.
    Let $C = C_0\union Q$.
    $C\subseteq\pow{X}$
        by \cref{covering_order_at_most_on_subset_50,subseteq,union_subsets_is_subset}.
    Show that for all $y\in Y$ we have
        $y$ has multiplicity at most $m + 1$ in $C$.
    \begin{subproof}
        Fix $y\in Y$.
        Suppose not.
        Take $F,h$ such that
            $F\subseteq C$ and
            $h$ is a bijection from $\initialSegment{m + 1 + 1}$ to $F$ and
            for all $W\in F$ we have $y\in W$
            by \cref{point_multiplicity_at_most_50}.
        $F\subseteq C_0$
            by \cref{subseteq,union_iff,inter_elim_right,setminus_elim_right}.
        Contradiction by
            \cref{covering_order_at_most_on_subset_50,point_multiplicity_at_most_50}.
    \end{subproof}
    $C$ has order at most $m + 1$ at points of $Y$ in $X$
        by \cref{covering_order_at_most_on_subset_50}.
    Show that for all $x\in X$ we have $x\in\unions{C}$.
    \begin{subproof}
        Fix $x\in X$.
        \begin{byCase}
        \caseOf{$x\in Y$.}
            Follows by
                \cref{open_refinement_on,covering,subseteq,unions_iff,choice_function,times_tuple_elim,fst_eq,snd_eq,inter_elim_right,img_of_function_intro,inter_intro,union_intro_left}.
        \caseOf{$x\notin Y$.}
            Follows by
                \cref{open_covering,covering,subseteq,unions_iff,powerset_elim,inter_lower_left,subseteq_transitive,powerset_intro,union_intro_right,inter_intro,setminus_intro}.
        \end{byCase}
    \end{subproof}
    For all $W\in C$ we have
        $W$ is open in $X$ with respect to $T$
        by \cref{union_iff,img_of_function_elim,inter_elim_right,choice_function,times_tuple_elim,fst_eq,snd_eq,open_in,open_covering,closed_in,topology_on}.
    $C$ is an open refinement of $A$ on $X$ with respect to $T$
        by \cref{union_iff,img_of_function_elim,inter_elim_right,choice_function,times_tuple_elim,fst_eq,snd_eq,inter_lower_left,subseteq,covering,refinement_on,open_refinement_on}.
\end{proof}

% from §50 helper lemma 11: grouping a refinement preserves split order bounds
\begin{lemma}\label{grouped_refinement_split_order_50}
    Assume $X = Y\union Z$.
    Suppose $m\in\naturalsPlus$.
    Suppose $B$ has order at most $m$ at points of $Y$ in $X$.
    Suppose $C$ has order at most $m$ at points of $Z$ in $X$.
    Assume $f$ is a function on $C$.
    Assume for all $W\in C$ we have
        $\apply{f}{W}\in B$ and $W\subseteq\apply{f}{W}$.
    Assume $p$ is a function on $B$.
    Assume for all $V\in B$ we have $\apply{p}{V}\subseteq C$.
    Assume for all $V,W$ such that
        $V\in B$ and $W\in C$
        we have $W\in\apply{p}{V}$ iff $\apply{f}{W} = V$.
    Let $D = \{U\in\pow{X} \mid
        \exists V\in B.
        U = \unions{\apply{p}{V}}\}$.
    Then $D$ has order at most $m$ in $X$.
\end{lemma}
\begin{proof}
    $D\subseteq\pow{X}$.
    Show that for all $x\in X$ we have $x$ has multiplicity at most $m$ in $D$.
    \begin{subproof}
        Fix $x\in X$.
        Suppose not.
        Take $F,h$ such that
            $F\subseteq D$ and
            $h$ is a bijection from $\initialSegment{m + 1}$ to $F$ and
            for all $U\in F$ we have $x\in U$
            by \cref{point_multiplicity_at_most_50}.
        Let $I = \initialSegment{m + 1}$.
        $h\in\funs{I}{F}$ and $h$ is injective
            by \cref{bijections,bijections_to_funs}.
        Let $R = \{w\in I\times B \mid
            \exists i\in I.\exists V\in B.
            w = (i,V)\land
            \apply{h}{i} = \unions{\apply{p}{V}}\}$.
        $R$ is a relation by \cref{relation,subseteq,times_elem_is_tuple}.
        For all $i\in I$ there exists $V$ such that $i\mathrel{R}V$
            by \cref{funs_type_apply,subseteq,times_tuple_intro,fst_eq,snd_eq}.
        Take $g$ such that
            $g$ is a function on $I$ and
            for all $i\in I$ we have $i\mathrel{R}\apply{g}{i}$
            by \cref{choice_function}.
        For all $i\in I$ we have
            $\apply{g}{i}\in B$ and
            $\apply{h}{i} = \unions{\apply{p}{\apply{g}{i}}}$
            by \cref{choice_function,times_tuple_elim,fst_eq,snd_eq}.
        $g$ is injective
            by \cref{injective_function,choice_function,funs_elim,funs_is_function}.
        Let $H = \img{g}{I}$.
        $g$ is a bijection from $I$ to $H$
            by \cref{img_of_function_intro,choice_function,inter_intro,funs_intro,funs_surj_iff,img_dom,bijections}.
        $H\subseteq B$
            by \cref{img_iff,function_apply_iff,subseteq}.
        Let $S = \{w\in I\times C \mid
            \exists i\in I.\exists W\in C.
            w = (i,W)\land
            x\in W\land
            \apply{f}{W} = \apply{g}{i}\}$.
        $S$ is a relation by \cref{relation,subseteq,times_elem_is_tuple}.
        For all $i\in I$ there exists $W$ such that $i\mathrel{S}W$
            by \cref{funs_type_apply,subseteq,choice_function,times_tuple_elim,fst_eq,snd_eq,unions_iff,times_tuple_intro}.
        Take $r$ such that
            $r$ is a function on $I$ and
            for all $i\in I$ we have $i\mathrel{S}\apply{r}{i}$
            by \cref{choice_function}.
        For all $i\in I$ we have
            $\apply{r}{i}\in C$ and
            $x\in\apply{r}{i}$ and
            $\apply{f}{\apply{r}{i}} = \apply{g}{i}$
            by \cref{choice_function,times_tuple_elim,fst_eq,snd_eq}.
        $r$ is injective
            by \cref{injective_function,choice_function}.
        Let $K = \img{r}{I}$.
        $r$ is a bijection from $I$ to $K$
            by \cref{img_of_function_intro,choice_function,inter_intro,funs_intro,funs_surj_iff,img_dom,bijections}.
        $K\subseteq C$
            by \cref{img_iff,function_apply_iff,subseteq}.
        For all $W\in K$ we have $x\in W$
            by \cref{img_iff,function_apply_iff,subseteq,choice_function,times_tuple_elim,fst_eq,snd_eq}.
        Contradiction by
            \cref{img_of_function_elim,inter_elim_right,choice_function,times_tuple_elim,fst_eq,snd_eq,subseteq,union_iff,covering_order_at_most_on_subset_50,point_multiplicity_at_most_50}.
    \end{subproof}
    Follows by \cref{covering_order_at_most_50}.
\end{proof}

% from §50 helper lemma 12: merge two pointwise bounded refinements
\begin{lemma}\label{merge_split_order_refinements_50}
    Assume $X = Y\union Z$.
    Assume $T$ is a topology on $X$.
    Suppose $m\in\naturalsPlus$.
    Suppose $B$ has order at most $m$ at points of $Y$ in $X$.
    Suppose $C$ has order at most $m$ at points of $Z$ in $X$.
    Suppose $C$ is an open refinement of $B$ on $X$ with respect to $T$.
    Then there exists $D$ such that
        $D$ is an open refinement of $B$ on $X$ with respect to $T$ and
        $D$ has order at most $m$ in $X$.
\end{lemma}
\begin{proof}
    $B\subseteq\pow{X}$ and $C\subseteq\pow{X}$
        by \cref{open_refinement_on,refinement_on}.
    Let $R = \{w\in C\times B \mid
        \exists W\in C.\exists V\in B.
        w = (W,V)\land W\subseteq V\}$.
    $R$ is a relation by \cref{relation,subseteq,times_elem_is_tuple}.
    For all $W\in C$ there exists $V$ such that $W\mathrel{R}V$
        by \cref{open_refinement_on,refinement_on,times_tuple_intro,fst_eq,snd_eq}.
    Take $f$ such that
        $f$ is a function on $C$ and
        for all $W\in C$ we have $W\mathrel{R}\apply{f}{W}$
        by \cref{choice_function}.
    For all $W\in C$ we have
        $\apply{f}{W}\in B$ and $W\subseteq\apply{f}{W}$
        by \cref{choice_function,times_tuple_elim,fst_eq,snd_eq}.
    Let $P = \{w\in B\times\pow{C} \mid
        \forall W. (W\in\snd{w}\iff
        (W\in C\land\apply{f}{W} = \fst{w}))\}$.
    $P$ is a relation by \cref{relation,subseteq,times_elem_is_tuple}.
    Show that for all $V\in B$ there exists $E$ such that $V\mathrel{P}E$.
    \begin{subproof}
        Fix $V\in B$.
        Let $E = \{W\in C \mid \apply{f}{W} = V\}$.
        Thus $V\mathrel{P}E$
            by \cref{subseteq,powerset_intro,times_tuple_intro,fst_eq,snd_eq}.
    \end{subproof}
    Take $p$ such that
        $p$ is a function on $B$ and
        for all $V\in B$ we have $V\mathrel{P}\apply{p}{V}$
        by \cref{choice_function}.
    For all $V\in B$ we have $\apply{p}{V}\subseteq C$
        by \cref{choice_function,times_tuple_elim,fst_eq,snd_eq,powerset_elim}.
    For all $V,W$ such that $V\in B$ and $W\in C$
        we have $W\in\apply{p}{V}$ iff $\apply{f}{W} = V$
        by \cref{choice_function,times_tuple_elim,fst_eq,snd_eq}.
    Let $D = \{U\in\pow{X} \mid
        \exists V\in B. U = \unions{\apply{p}{V}}\}$.
    $D$ has order at most $m$ in $X$
        by \cref{grouped_refinement_split_order_50}.
    For all $x\in X$ we have $x\in\unions{D}$
        by \cref{open_refinement_on,covering,subseteq,unions_iff,choice_function,times_tuple_elim,fst_eq,snd_eq,powerset_elim,subseteq_transitive,unions_family,powerset_intro}.
    For all $U\in D$ there exists $V\in B$ such that $U\subseteq V$
        by \cref{unions_iff,powerset_elim,subseteq,choice_function,times_tuple_elim,fst_eq,snd_eq}.
    For all $U\in D$ we have
        $U$ is open in $X$ with respect to $T$
        by \cref{powerset_elim,subseteq,open_refinement_on,open_in,topology_on}.
    $D$ is an open refinement of $B$ on $X$ with respect to $T$
        by \cref{covering_order_at_most_50,subseteq,covering,refinement_on,open_refinement_on}.
\end{proof}

% from §50 helper lemma 13: transitivity of open refinements
\begin{lemma}\label{open_refinement_transitive_50}
    Suppose $B$ is an open refinement of $A$ on $X$ with respect to $T$.
    Suppose $C$ is an open refinement of $B$ on $X$ with respect to $T$.
    Then $C$ is an open refinement of $A$ on $X$ with respect to $T$.
\end{lemma}
\begin{proof}
    Follows by
        \cref{open_refinement_on,refinement_on,subseteq_transitive}.
\end{proof}

% from §50 Item 5: dimension of a union of two closed finite dimensional subspaces
\begin{theorem}\label{closed_union_two_dimension_50}
    Assume $X = Y\union Z$.
    Assume $Y$ is a subspace of $X$ with respect to $T$.
    Assume $Z$ is a subspace of $X$ with respect to $T$.
    Assume $Y$ is closed in $X$ with respect to $T$.
    Assume $Z$ is closed in $X$ with respect to $T$.
    Let $T_Y = \subspaceTopology{T}{Y}$.
    Let $T_Z = \subspaceTopology{T}{Z}$.
    Suppose $p$ is the topological dimension of $Y$ with respect to $T_Y$.
    Suppose $q$ is the topological dimension of $Z$ with respect to $T_Z$.
    Suppose $m$ is a maximum of $\{p,q\}$.
    Then $m$ is the topological dimension of $X$ with respect to $T$.
\end{theorem}
\begin{proof}
    $T$ is a topology on $X$ and $Y\subseteq X$ and $Z\subseteq X$
        by \cref{subspace_of}.
    $m\in\naturals$
        by \cref{topological_dimension_50,topological_dimension_bound_50,real_maximum,upair_elim}.
    $m$ is a topological dimension bound for $Y$ with respect to $T_Y$ and
        $m$ is a topological dimension bound for $Z$ with respect to $T_Z$
        by \cref{upair_intro_left,upair_intro_right,real_maximum,topological_dimension_50,topological_dimension_bound_50,topological_dimension_bound_monotone_50}.
    Show that for all $A$ if
        $A$ is an open covering of $X$ with respect to $T$
        then there exists $D$ such that
            $D$ is an open refinement of $A$ on $X$ with respect to $T$ and
            $D$ has order at most $m + 1$ in $X$.
    \begin{subproof}
        Fix $A$.
        Assume $A$ is an open covering of $X$ with respect to $T$.
        Take $B$ such that
            $B$ is an open refinement of $A$ on $X$ with respect to $T$ and
            $B$ has order at most $m + 1$ at points of $Y$ in $X$
            by \cref{closed_subspace_pointwise_refinement_50}.
        Take $C$ such that
            $C$ is an open refinement of $B$ on $X$ with respect to $T$ and
            $C$ has order at most $m + 1$ at points of $Z$ in $X$
            by \cref{open_refinement_on,open_covering,closed_subspace_pointwise_refinement_50}.
        Take $D$ such that
            $D$ is an open refinement of $A$ on $X$ with respect to $T$ and
            $D$ has order at most $m + 1$ in $X$
            by \cref{natural_successor_positive_50,merge_split_order_refinements_50,open_refinement_transitive_50}.
    \end{subproof}
    $m$ is a topological dimension bound for $X$ with respect to $T$
        by \cref{topological_dimension_bound_50}.
    For all $n$ such that
        $n$ is a topological dimension bound for $X$ with respect to $T$
        we have $m\leq n$
        by \cref{closed_subspace_inherits_dimension_bound_50,topological_dimension_50,real_maximum,upair_elim}.
    Follows by \cref{topological_dimension_50}.
\end{proof}

% from §50 helper lemma 14: a nonempty set of natural numbers has a least element
\begin{lemma}\label{natural_set_least_element_50}
    Suppose $A\subseteq\naturals$ and $A\neq\emptyset$.
    Then there exists $m\in A$ such that for all $a\in A$ we have $m\leq a$.
\end{lemma}
\begin{proof}
    Follows by
        \cref{subseteq,naturals,union_iff,singleton_iff,real_add_ident,real_one_in_naturals,real_naturals_subset_reals,real_zero_lt_one,real_one_leq_naturalsplus,real_lt_leq_trans_local,real_naturals_least_element}.
\end{proof}

% from §50 helper lemma 15: finite-dimensional spaces have a topological dimension
\begin{lemma}\label{finite_dimensional_dimension_exists_50}
    Suppose $X$ is finite dimensional with respect to $T$.
    Then there exists $m\in\naturals$ such that
        $m$ is the topological dimension of $X$ with respect to $T$.
\end{lemma}
\begin{proof}
    Let $A = \{n\in\naturals \mid
        \text{$n$ is a topological dimension bound for $X$ with respect to $T$}\}$.
    Take $m\in A$ such that for all $a\in A$ we have $m\leq a$
        by \cref{subseteq,finite_dimensional_space_50,inhabited_nonempty,natural_set_least_element_50}.
    Thus there exists $d\in\naturals$ such that
        $d$ is the topological dimension of $X$ with respect to $T$
        by \cref{topological_dimension_bound_50,topological_dimension_50}.
\end{proof}

% from §50 helper lemma 16: a common dimension bound is preserved by a closed union
\begin{lemma}\label{closed_union_preserves_common_dimension_bound_50}
    Assume $T$ is a topology on $X$.
    Assume $Y\subseteq X$ and $Z\subseteq X$.
    Assume $Y$ is closed in $X$ with respect to $T$.
    Assume $Z$ is closed in $X$ with respect to $T$.
    Let $P = Y\union Z$.
    Let $T_P = \subspaceTopology{T}{P}$.
    Suppose $m$ is a topological dimension bound for $Y$ with respect to $\subspaceTopology{T}{Y}$.
    Suppose $m$ is a topological dimension bound for $Z$ with respect to $\subspaceTopology{T}{Z}$.
    Then $m$ is a topological dimension bound for $P$ with respect to $T_P$.
\end{lemma}
\begin{proof}
    $P\subseteq X$ by \cref{union_subsets_is_subset}.
    $T_P$ is a topology on $P$ by \cref{subspace_topology_is_topology}.
    $Y\subseteq P$ and $Z\subseteq P$
        by \cref{union_upper_left,union_upper_right}.
    $Y$ is a subspace of $P$ with respect to $T_P$ and
        $Z$ is a subspace of $P$ with respect to $T_P$
        by \cref{subspace_of}.
    $Y$ is closed in $P$ with respect to $T_P$ and
        $Z$ is closed in $P$ with respect to $T_P$
        by \cref{inter_absorb_supseteq_right,closed_in_subspace_iff,subspace_of}.
    $m$ is a topological dimension bound for $Y$ with respect to $\subspaceTopology{T_P}{Y}$ and
        $m$ is a topological dimension bound for $Z$ with respect to $\subspaceTopology{T_P}{Z}$
        by \cref{subspace_subspace_topology}.
    Show that for all $A$ if
        $A$ is an open covering of $P$ with respect to $T_P$
        then there exists $D$ such that
            $D$ is an open refinement of $A$ on $P$ with respect to $T_P$ and
            $D$ has order at most $m + 1$ in $P$.
    \begin{subproof}
        Fix $A$.
        Assume $A$ is an open covering of $P$ with respect to $T_P$.
        Take $B$ such that
            $B$ is an open refinement of $A$ on $P$ with respect to $T_P$ and
            $B$ has order at most $m + 1$ at points of $Y$ in $P$
            by \cref{closed_subspace_pointwise_refinement_50}.
        Take $C$ such that
            $C$ is an open refinement of $B$ on $P$ with respect to $T_P$ and
            $C$ has order at most $m + 1$ at points of $Z$ in $P$
            by \cref{open_refinement_on,open_covering,closed_subspace_pointwise_refinement_50}.
        Take $D$ such that
            $D$ is an open refinement of $A$ on $P$ with respect to $T_P$ and
            $D$ has order at most $m + 1$ in $P$
            by \cref{topological_dimension_bound_50,natural_successor_positive_50,merge_split_order_refinements_50,open_refinement_transitive_50}.
    \end{subproof}
    Follows by \cref{topological_dimension_bound_50}.
\end{proof}

% from §50 helper lemma 17: an initial finite union of closed indexed subspaces is closed
\begin{lemma}\label{indexed_initial_union_closed_50}
    Assume $T$ is a topology on $X$.
    Assume $k\in\naturalsPlus$ and $n\in\naturalsPlus$.
    Assume $n\leq k$.
    Assume $Y$ is a function on $\initialSegment{k}$.
    Assume for all $i\in\initialSegment{k}$ we have
        $\apply{Y}{i}$ is closed in $X$ with respect to $T$.
    Let $P = \unions{\img{Y}{\initialSegment{n}}}$.
    Then $P$ is closed in $X$ with respect to $T$.
\end{lemma}
\begin{proof}
    Let $F = \img{Y}{\initialSegment{n}}$.
    $F$ is finite by
        \cref{initial_segment_monotone_50,initial_segment_finite,restrl_of_function_is_function,dom_of_restrl_eq_inter,inter_absorb_supseteq_left,subseteq,inter_absorb_supseteq_right,subseteq_refl,restrl_img,finite_image_function}.
    Follows by
        \cref{initial_segment_monotone_50,img_of_function_elim,inter_elim_right,subseteq,closed_in,powerset_intro,closed_finite_unions}.
\end{proof}

% from §50 helper lemma 18: split an indexed initial union at its last member
\begin{lemma}\label{indexed_initial_union_successor_50}
    Assume $k\in\naturalsPlus$ and $n\in\naturalsPlus$.
    Assume $n + 1\leq k$.
    Assume $Y$ is a function on $\initialSegment{k}$.
    Then
        $\unions{\img{Y}{\initialSegment{n + 1}}}
        =
        \unions{\img{Y}{\initialSegment{n}}}\union\apply{Y}{n + 1}$.
\end{lemma}
\begin{proof}
    $\unions{\img{Y}{\initialSegment{n + 1}}}
        \subseteq
        \unions{\img{Y}{\initialSegment{n}}}\union\apply{Y}{n + 1}$
        by \cref{real_naturals_closed_succ,initial_segment_naturalsplus,initial_segment_subset_succ,subseteq,unions_iff,img_of_function_elim,inter_elim_right,initial_segment_succ_split,img_of_function_intro,inter_intro,inter_elim_left,union_intro_left,union_intro_right}.
    $\unions{\img{Y}{\initialSegment{n}}}\union\apply{Y}{n + 1}
        \subseteq
        \unions{\img{Y}{\initialSegment{n + 1}}}$
        by \cref{real_naturals_closed_succ,initial_segment_naturalsplus,initial_segment_subset_succ,subseteq,union_iff,unions_iff,img_of_function_elim,inter_elim_right,img_of_function_intro,inter_intro,inter_elim_left}.
    Follows by \cref{subseteq_antisymmetric}.
\end{proof}

% from §50 helper lemma 19: a common bound is preserved by a finite indexed closed union
\begin{lemma}\label{finite_closed_family_common_dimension_bound_50}
    Assume $T$ is a topology on $X$.
    Assume $k\in\naturalsPlus$.
    Assume $Y$ is a function on $\initialSegment{k}$.
    Assume for all $i\in\initialSegment{k}$ we have
        $\apply{Y}{i}\subseteq X$ and
        $\apply{Y}{i}$ is closed in $X$ with respect to $T$ and
        $m$ is a topological dimension bound for $\apply{Y}{i}$ with respect to
            $\subspaceTopology{T}{\apply{Y}{i}}$.
    Let $P = \unions{\img{Y}{\initialSegment{k}}}$.
    Then $m$ is a topological dimension bound for $P$ with respect to
        $\subspaceTopology{T}{P}$.
\end{lemma}
\begin{proof}
    Let $S = \{n\in\naturalsPlus \mid
        \text{if $n\leq k$ then $m$ is a topological dimension bound for
        $\unions{\img{Y}{\initialSegment{n}}}$ with respect to
        $\subspaceTopology{T}{\unions{\img{Y}{\initialSegment{n}}}}$}\}$.
    $S\subseteq\naturalsPlus$.
    $1\in\initialSegment{k}$
        by \cref{real_one_leq_naturalsplus,real_one_in_naturals,initial_segment_naturalsplus}.
    $m$ is a topological dimension bound for
        $\unions{\img{Y}{\initialSegment{1}}}$ with respect to
        $\subspaceTopology{T}{\unions{\img{Y}{\initialSegment{1}}}}$
        by \cref{unions_iff,img_of_function_elim,inter_elim_right,initial_segment_one,singleton_iff,subseteq,img_of_function_intro,inter_intro,subseteq_antisymmetric}.
    $1\in S$.
    For all $n\in S$ if $n + 1\leq k$ then
            $m$ is a topological dimension bound for
                $\unions{\img{Y}{\initialSegment{n + 1}}}$ with respect to
                $\subspaceTopology{T}{\unions{\img{Y}{\initialSegment{n + 1}}}}$
            by \cref{subseteq,real_naturals_closed_succ,real_naturals_succ_bound_elim_local,indexed_initial_union_closed_50,closed_in,initial_segment_naturalsplus,closed_union_preserves_common_dimension_bound_50,indexed_initial_union_successor_50}.
    Thus for all $n\in S$ we have $n + 1\in S$
        by \cref{subseteq,real_naturals_closed_succ}.
    Follows by \cref{real_one_in_naturals,real_naturals_induction,subseteq}.
\end{proof}

% from §50 helper lemma 20: the whole-space subspace topology is unchanged
\begin{lemma}\label{whole_subspace_topology_50}
    Assume $T$ is a topology on $X$.
    Then $\subspaceTopology{T}{X} = T$.
\end{lemma}
\begin{proof}
    Follows by
        \cref{subspace_topology,topology_on,pow_iff,subseteq,inter_absorb_supseteq_left,subseteq_antisymmetric}.
\end{proof}

% from §50 Item 6: dimension of a finite union of closed finite dimensional subspaces
\begin{corollary}\label{closed_finite_union_dimension_50}
    Assume $k\in\naturalsPlus$.
    Assume $Y$ is a function on $\initialSegment{k}$.
    Assume $X = \unions{\img{Y}{\initialSegment{k}}}$.
    Suppose for all $i\in\initialSegment{k}$ we have
        $\apply{Y}{i}$ is a subspace of $X$ with respect to $T$ and
        $\apply{Y}{i}$ is closed in $X$ with respect to $T$ and
        $\apply{Y}{i}$ is finite dimensional with respect to $\subspaceTopology{T}{\apply{Y}{i}}$.
    Let $D = \{d\in\naturals \mid \text{there exists $i\in\initialSegment{k}$ such that $d$ is the topological dimension of $\apply{Y}{i}$ with respect to $\subspaceTopology{T}{\apply{Y}{i}}$}\}$.
    Suppose $m$ is a maximum of $D$.
    Then $m$ is the topological dimension of $X$ with respect to $T$.
\end{corollary}
\begin{proof}
    $T$ is a topology on $X$
        by \cref{real_one_leq_naturalsplus,real_one_in_naturals,initial_segment_naturalsplus,subspace_of}.
    $m\in D$ by \cref{real_maximum}.
    $m\in\naturals$.
    For all $i\in\initialSegment{k}$ we have
        $\apply{Y}{i}\subseteq X$ and
        $\apply{Y}{i}$ is closed in $X$ with respect to $T$ and
        $m$ is a topological dimension bound for $\apply{Y}{i}$ with respect to
            $\subspaceTopology{T}{\apply{Y}{i}}$
        by \cref{subspace_of,finite_dimensional_dimension_exists_50,real_maximum,topological_dimension_50,topological_dimension_bound_monotone_50}.
    $m$ is a topological dimension bound for $X$ with respect to $T$
        by \cref{finite_closed_family_common_dimension_bound_50,whole_subspace_topology_50}.
    Follows by
        \cref{closed_subspace_inherits_dimension_bound_50,topological_dimension_50}.
\end{proof}

% from §50 helper definition 6: arc
\begin{definition}\label{arc_50}
    $A$ is an arc in $X$ with respect to $T$ iff
        $A$ is a subspace of $X$ with respect to $T$ and
        there exists $h$ such that
            $h$ is a homeomorphism between $\intervalclosed{0}{1}$ and $A$ with respect to
                $\subspaceTopology{\standardTopologyR}{\intervalclosed{0}{1}}$ and $\subspaceTopology{T}{A}$.
\end{definition}

% from §50 helper definition 7: endpoint of an arc
\begin{definition}\label{arc_endpoint_50}
    $p$ is an endpoint of $A$ in $X$ with respect to $T$ iff
        $A$ is an arc in $X$ with respect to $T$ and
        $p\in A$ and
        $A\setminus\{p\}$ is connected with respect to $\subspaceTopology{T}{A\setminus\{p\}}$.
\end{definition}

% from §50 helper definition 8: family of arcs
\begin{definition}\label{family_of_arcs_50}
    $E$ is a family of arcs in $G$ with respect to $T$ iff
        for all $A\in E$ we have $A$ is an arc in $G$ with respect to $T$.
\end{definition}

% from §50 helper definition 9: linear graph intersection condition
\begin{definition}\label{endpoint_intersection_50}
    $A$ has endpoint intersection with $B$ in $G$ with respect to $T$ iff
        there exists $p$ such that
            $p$ is an endpoint of $A$ in $G$ with respect to $T$ and
            $p$ is an endpoint of $B$ in $G$ with respect to $T$ and
            $A\inter B = \{p\}$.
\end{definition}

% from §50 helper definition 10: linear graph intersection condition
\begin{definition}\label{linear_graph_intersections_50}
    $E$ has linear graph intersections in $G$ with respect to $T$ iff
        for all $A,B\in E$ if $A\neq B$ then
            $A\inter B = \emptyset$ or
            $A$ has endpoint intersection with $B$ in $G$ with respect to $T$.
\end{definition}

% from §50 helper definition 11: finite linear graph decomposition
\begin{definition}\label{finite_linear_graph_decomposition_50}
    $E$ is a finite linear graph decomposition of $G$ with respect to $T$ iff
        $G$ is Hausdorff with respect to $T$ and
        $E$ is finite and
        $E$ is a covering of $G$ and
        $E$ is a family of arcs in $G$ with respect to $T$ and
        $E$ has linear graph intersections in $G$ with respect to $T$.
\end{definition}

% from §50 helper definition 12: finite linear graph
\begin{definition}\label{finite_linear_graph_50}
    $G$ is a finite linear graph with respect to $T$ iff
        there exists $E$ such that $E$ is a finite linear graph decomposition of $G$ with respect to $T$.
\end{definition}

% from §50 helper definition 13: edge of a finite linear graph
\begin{definition}\label{edge_of_finite_linear_graph_50}
    $e$ is an edge of $G$ with respect to $T$ iff
        there exists $E$ such that
            $E$ is a finite linear graph decomposition of $G$ with respect to $T$ and
            $e\in E$.
\end{definition}

% from §50 helper definition 14: vertex of a finite linear graph
\begin{definition}\label{vertex_of_finite_linear_graph_50}
    $v$ is a vertex of $G$ with respect to $T$ iff
        there exists $e$ such that
            $e$ is an edge of $G$ with respect to $T$ and
            $v$ is an endpoint of $e$ in $G$ with respect to $T$.
\end{definition}

% from §50 helper definition 15: finite set has at most a positive number of elements
\begin{definition}\label{finite_set_has_at_most_50}
    $A$ has at most $m$ elements iff
        $m\in\naturalsPlus$ and
        $A$ is finite and
        there exists no $B$ such that
            $B\subseteq A$ and
            there exists $g$ such that $g$ is a bijection from $\initialSegment{m + 1}$ to $B$.
\end{definition}

% from §50 helper definition 16: zero vector in a finite euclidean space
\begin{definition}\label{zero_vector_50}
    $\zeroVector{N} = \{(i,0) \mid i\in\initialSegment{N}\}$.
\end{definition}

% from §50 helper definition 17: affine combination of a finite list of points
\begin{definition}\label{affine_combination_value_50}
    $y$ is an affine combination of $s$ with coefficients $a$ in dimension $N$ iff
        there exists $n\in\naturalsPlus$ such that
            $s\in\funs{\initialSegment{n}}{\realsPower{\initialSegment{N}}}$ and
            $a\in\funs{\initialSegment{n}}{\reals}$ and
            $y\in\realsPower{\initialSegment{N}}$ and
            for all $j\in\initialSegment{N}$ there exists $b$ such that
                $b\in\funs{\initialSegment{n}}{\reals}$ and
                for all $i\in\initialSegment{n}$ we have $b(i) = a(i) \rmul \apply{\apply{s}{i}}{j}$ and
                $\sumFiniteSequence{b}{\apply{y}{j}}$.
\end{definition}

% from §50 Item 7: geometrically independent finite set of points
\begin{definition}\label{geometrically_independent_sequence_50}
    $s$ is a geometrically independent sequence in dimension $N$ iff
        $N\in\naturalsPlus$ and
        there exists $n\in\naturalsPlus$ such that
            $s\in\funs{\initialSegment{n}}{\realsPower{\initialSegment{N}}}$ and
            for all $a$ if
                $a\in\funs{\initialSegment{n}}{\reals}$ and
                $\zeroVector{N}$ is an affine combination of $s$ with coefficients $a$ in dimension $N$ and
                $\sumFiniteSequence{a}{0}$
            then for all $i\in\initialSegment{n}$ we have $a(i) = 0$.
\end{definition}

% from §50 helper definition 18: geometrically independent set of points
\begin{definition}\label{geometrically_independent_set_50}
    $A$ is a geometrically independent set in dimension $N$ iff
        $A\subseteq\realsPower{\initialSegment{N}}$ and
        $A$ is finite and
        either $A = \emptyset$ or
        there exist $n,s$ such that
            $n\in\naturalsPlus$ and
            $s$ is a bijection from $\initialSegment{n}$ to $A$ and
            $s$ is a geometrically independent sequence in dimension $N$.
\end{definition}

% from §50 Item 8: plane determined by geometrically independent points
\begin{definition}\label{plane_determined_by_points_50}
    $P$ is the plane determined by $A$ in dimension $N$ iff
        $A$ is a geometrically independent set in dimension $N$ and
        $P = \{y\in\realsPower{\initialSegment{N}} \mid \text{there exist $n,s,a$ such that
            $n\in\naturalsPlus$ and
            $s$ is a bijection from $\initialSegment{n}$ to $A$ and
            $a\in\funs{\initialSegment{n}}{\reals}$ and
            $\sumFiniteSequence{a}{1}$ and
            $y$ is an affine combination of $s$ with coefficients $a$ in dimension $N$}\}$.
\end{definition}

% from §50 helper definition 19: k-plane
\begin{definition}\label{k_plane_50}
    $P$ is a $k$-plane in dimension $N$ iff
        $k\in\naturals$ and
        $N\in\naturalsPlus$ and
        there exists $A$ such that
            $A$ is a geometrically independent set in dimension $N$ and
            there exists $g$ such that
                $g$ is a bijection from $\initialSegment{k + 1}$ to $A$ and
                $P$ is the plane determined by $A$ in dimension $N$.
\end{definition}

% from §50 Item 9: general position in euclidean space
\begin{definition}\label{general_position_50}
    $A$ is in general position in dimension $N$ iff
        $N\in\naturalsPlus$ and
        $A\subseteq\realsPower{\initialSegment{N}}$ and
        for all $B$ such that
            $B\subseteq A$ and
            $B$ has at most $N + 1$ elements
        we have $B$ is a geometrically independent set in dimension $N$.
\end{definition}

% from §50 helper lemma 21: the empty set is geometrically independent
\begin{lemma}\label{empty_geometrically_independent_50}
    Then $\emptyset$ is a geometrically independent set in dimension $N$.
\end{lemma}
\begin{proof}
    Follows by
        \cref{real_one_in_naturals,real_one_leq_naturalsplus,initial_segment_naturalsplus,bijections_to_funs,funs_type_apply,emptyset,emptyset_subseteq,emptyset_is_finite,geometrically_independent_set_50}.
\end{proof}

% from §50 helper lemma 22: a singleton is geometrically independent
\begin{lemma}\label{singleton_geometrically_independent_50}
    Assume $N\in\naturalsPlus$.
    Assume $y\in\realsPower{\initialSegment{N}}$.
    Then $\{y\}$ is a geometrically independent set in dimension $N$.
\end{lemma}
\begin{proof}
    Let $I = \initialSegment{1}$.
    Let $s = \{(i,y) \mid i\in I\}$.
    For all $i\in I$ we have $\apply{s}{i} = y$
        by \cref{relation,rightunique,pair_eq_iff,dom_intro,dom_iff,subseteq,subseteq_antisymmetric,function_apply_iff}.
    $s$ is a bijection from $I$ to $\{y\}$
        by \cref{relation,rightunique,pair_eq_iff,dom_intro,dom_iff,subseteq_antisymmetric,funs_intro,injective_function,bijections,surj,initial_segment_one,singleton_iff,subseteq}.
    Thus $s$ is a geometrically independent sequence in dimension $N$
        by \cref{relation,rightunique,pair_eq_iff,dom_intro,dom_iff,subseteq,subseteq_antisymmetric,funs_intro,sum_finite_sequence_one_term,sum_finite_sequence_unique,initial_segment_one,singleton_iff,geometrically_independent_sequence_50,real_one_in_naturals}.
    Follows by
        \cref{singleton_subset_intro,pair_is_finite,upair_intro_left,subseteq_finite_is_finite,emptyset,singleton_intro,geometrically_independent_set_50,real_one_in_naturals}.
\end{proof}

% from §50 helper lemma 23: a nonzero last coefficient places the last point in the preceding affine plane
\begin{lemma}\label{affine_relation_last_point_in_plane_50}
    Assume $N,n\in\naturalsPlus$.
    Let $I = \initialSegment{n}$.
    Let $I_1 = \initialSegment{n + 1}$.
    Let $E = \realsPower{\initialSegment{N}}$.
    Assume $A$ is a geometrically independent set in dimension $N$.
    Assume $s$ is a bijection from $I$ to $A$.
    Assume $P$ is the plane determined by $A$ in dimension $N$.
    Assume $y\in E$.
    Assume $t\in\funs{I_1}{E}$.
    Assume $\restrl{t}{I} = s$ and $t(n + 1) = y$.
    Assume $a\in\funs{I_1}{\reals}$.
    Assume $\zeroVector{N}$ is an affine combination of $t$ with coefficients $a$ in dimension $N$.
    Assume $\sumFiniteSequence{a}{0}$.
    Suppose $a(n + 1)\neq 0$.
    Then $y\in P$.
\end{lemma}
\begin{proof}
    Let $c = a(n + 1)$.
    Let $a_0 = \restrl{a}{I}$.
    $a_0\in\funs{I}{\reals}$ by
        \cref{sum_finite_sequence_split_last}.
    Take $r\in\reals$ such that
        $\sumFiniteSequence{a_0}{r}$ and $0 = r + c$
        by \cref{sum_finite_sequence_split_last}.
    $c\in\reals$ by
        \cref{real_naturals_closed_succ,initial_segment_naturalsplus,funs_type_apply}.
    $\inv{c}\in\reals$ and $\neg{\inv{c}}\in\reals$
        by \cref{real_mul_inverse,real_add_inverse}.
    Let $q = \neg{\inv{c}}$.
    Let $d = \coordScalarSequence{n}{q}{a_0}$.
    $d\in\funs{I}{\reals}$ by
        \cref{coord_scalar_sequence_function}.
    $\sumFiniteSequence{d}{q \rmul r}$
        by \cref{sum_finite_sequence_scalar}.
    Hence $\sumFiniteSequence{d}{1}$ by
        \cref{real_add_inverse,real_add_comm,real_lt_subst_left,real_lt_subst_right,real_add_cancel,real_mul_neg_right,real_mul_comm,real_mul_closed,real_neg_neg,real_mul_inverse}.
    Take $k$ such that
        $k\in\naturalsPlus$ and
        $t\in\funs{\initialSegment{k}}{E}$ and
        $a\in\funs{\initialSegment{k}}{\reals}$ and
        $\zeroVector{N}\in E$ and
        for all $j\in\initialSegment{N}$ there exists $e$ such that
            $e\in\funs{\initialSegment{k}}{\reals}$ and
            for all $i\in\initialSegment{k}$ we have
                $e(i) = a(i) \rmul \apply{\apply{t}{i}}{j}$ and
                $\sumFiniteSequence{e}{\apply{\zeroVector{N}}{j}}$
        by \cref{affine_combination_value_50}.
    $\initialSegment{k} = I_1$ by
        \cref{funs_elim}.
    $k = n + 1$ by
        \cref{real_naturals_closed_succ,initial_segment_naturalsplus,real_lt_subst_left,real_lt_subst_right,real_naturals_subset_reals,subseteq,real_leq_antisym}.
    $n + 1\in\initialSegment{k}$ by
        \cref{real_naturals_closed_succ,initial_segment_naturalsplus,real_lt_subst_left,real_lt_subst_right}.
    Show that for all $j\in\initialSegment{N}$ there exists $b$ such that
            $b\in\funs{I}{\reals}$ and
            for all $i\in I$ we have
                $b(i) = d(i) \rmul \apply{\apply{s}{i}}{j}$ and
            $\sumFiniteSequence{b}{\apply{y}{j}}$.
        \begin{subproof}
            Fix $j\in\initialSegment{N}$.
            Take $e$ such that
                $e\in\funs{\initialSegment{k}}{\reals}$ and
                for all $i\in\initialSegment{k}$ we have
                    $e(i) = a(i) \rmul \apply{\apply{t}{i}}{j}$ and
                $\sumFiniteSequence{e}{\apply{\zeroVector{N}}{j}}$
                by \cref{real_lt_subst_right}.
            $\apply{\zeroVector{N}}{j} = 0$
                by \cref{reals_power,tuple_power,real_lt_subst_left,zero_vector_50,pair_eq_iff,funs_is_function,function_apply_intro}.
            Let $e_0 = \restrl{e}{I}$.
            $e_0\in\funs{I}{\reals}$ by
                \cref{real_lt_subst_left,real_lt_subst_right,sum_finite_sequence_split_last}.
            Take $u\in\reals$ such that
                $\sumFiniteSequence{e_0}{u}$ and
                $0 = u + e(n + 1)$
                by \cref{real_lt_subst_right,sum_finite_sequence_split_last}.
            $e(n + 1) = c \rmul \apply{y}{j}$
                by \cref{real_naturals_closed_succ,initial_segment_naturalsplus,real_lt_subst_left,real_lt_subst_right}.
            Let $b = \coordScalarSequence{n}{q}{e_0}$.
            $b\in\funs{I}{\reals}$ by
                \cref{coord_scalar_sequence_function}.
            $\sumFiniteSequence{b}{\apply{y}{j}}$
                by \cref{sum_finite_sequence_scalar,reals_power,tuple_power,funs_type_apply,real_mul_closed,real_add_inverse,real_add_comm,real_add_cancel,real_lt_subst_left,real_lt_subst_right,real_mul_neg_right,real_mul_comm,real_mul_assoc,real_neg_neg,real_mul_inverse,real_mul_ident}.
            For all $i\in I$ we have
                $b(i) = d(i) \rmul \apply{\apply{s}{i}}{j}$
                by \cref{initial_segment_subset_succ,subseteq,funs_elim,restrl_of_function_apply,funs_is_function,real_lt_subst_left,coord_scalar_sequence,function_apply_intro,funs_type_apply,reals_power,tuple_power,real_mul_assoc,real_mul_eq_subst,real_lt_subst_right}.
            Follows.
        \end{subproof}
    Follows by
        \cref{geometrically_independent_set_50,bijections_to_funs,funs_elim,funs_type_apply,subseteq,funs_intro,affine_combination_value_50,plane_determined_by_points_50}.
\end{proof}

% from §50 helper lemma 24: a zero last coefficient restricts an affine zero relation
\begin{lemma}\label{affine_zero_relation_restrict_last_zero_50}
    Assume $N,n\in\naturalsPlus$.
    Let $I = \initialSegment{n}$.
    Let $I_1 = \initialSegment{n + 1}$.
    Let $E = \realsPower{\initialSegment{N}}$.
    Assume $s\in\funs{I}{E}$.
    Assume $t\in\funs{I_1}{E}$.
    Assume $\restrl{t}{I} = s$.
    Assume $a\in\funs{I_1}{\reals}$.
    Assume $a(n + 1) = 0$.
    Assume $\zeroVector{N}$ is an affine combination of $t$ with coefficients $a$ in dimension $N$.
    Let $a_0 = \restrl{a}{I}$.
    Then $\zeroVector{N}$ is an affine combination of $s$ with coefficients $a_0$ in dimension $N$.
\end{lemma}
\begin{proof}
    $n + 1\in\naturalsPlus$ by
        \cref{real_naturals_closed_succ}.
    $a_0\in\funs{I}{\reals}$ by
        \cref{sum_finite_sequence_exists,sum_finite_sequence_split_last}.
    Take $k$ such that
        $k\in\naturalsPlus$ and
        $t\in\funs{\initialSegment{k}}{E}$ and
        $a\in\funs{\initialSegment{k}}{\reals}$ and
        $\zeroVector{N}\in E$ and
        for all $j\in\initialSegment{N}$ there exists $e$ such that
            $e\in\funs{\initialSegment{k}}{\reals}$ and
            for all $i\in\initialSegment{k}$ we have
                $e(i) = a(i) \rmul \apply{\apply{t}{i}}{j}$ and
                $\sumFiniteSequence{e}{\apply{\zeroVector{N}}{j}}$
        by \cref{affine_combination_value_50}.
    $\initialSegment{k} = I_1$ by \cref{funs_elim}.
    $k = n + 1$ by
        \cref{initial_segment_naturalsplus,real_lt_subst_left,real_lt_subst_right,real_naturals_subset_reals,subseteq,real_leq_antisym}.
    $n + 1\in\initialSegment{k}$ by
        \cref{initial_segment_naturalsplus,real_lt_subst_left,real_lt_subst_right}.
    Show that for all $j\in\initialSegment{N}$ there exists $b$ such that
        $b\in\funs{I}{\reals}$ and
        for all $i\in I$ we have
            $b(i) = a_0(i) \rmul \apply{\apply{s}{i}}{j}$ and
            $\sumFiniteSequence{b}{\apply{\zeroVector{N}}{j}}$.
    \begin{subproof}
        Fix $j\in\initialSegment{N}$.
        Take $e$ such that
            $e\in\funs{\initialSegment{k}}{\reals}$ and
            for all $i\in\initialSegment{k}$ we have
                $e(i) = a(i) \rmul \apply{\apply{t}{i}}{j}$ and
                $\sumFiniteSequence{e}{\apply{\zeroVector{N}}{j}}$.
        $\apply{\zeroVector{N}}{j} = 0$
            by \cref{reals_power,tuple_power,real_lt_subst_left,zero_vector_50,pair_eq_iff,funs_is_function,function_apply_intro}.
        Let $b = \restrl{e}{I}$.
        $b\in\funs{I}{\reals}$ by
            \cref{real_lt_subst_left,real_lt_subst_right,sum_finite_sequence_split_last}.
        Take $u\in\reals$ such that
            $\sumFiniteSequence{b}{u}$ and
            $0 = u + e(n + 1)$
            by \cref{real_lt_subst_right,sum_finite_sequence_split_last}.
        $\sumFiniteSequence{b}{0}$ by
            \cref{real_naturals_closed_succ,initial_segment_naturalsplus,funs_type_apply,reals_power,tuple_power,real_mul_zero,real_add_ident,real_add_comm,real_lt_subst_left,real_lt_subst_right}.
        For all $i\in I$ we have
            $b(i) = a_0(i) \rmul \apply{\apply{s}{i}}{j}$
            by \cref{initial_segment_subset_succ,subseteq,funs_elim,restrl_of_function_apply,funs_is_function,real_lt_subst_left,real_lt_subst_right}.
        Follows by \cref{real_lt_subst_right}.
    \end{subproof}
    Follows by \cref{affine_combination_value_50}.
\end{proof}

% from §50 helper lemma 25: adjoining a point outside the affine plane preserves independence
\begin{lemma}\label{affine_independence_extension_outside_plane_50}
    Assume $N,n\in\naturalsPlus$.
    Let $I = \initialSegment{n}$.
    Let $I_1 = \initialSegment{n + 1}$.
    Let $E = \realsPower{\initialSegment{N}}$.
    Assume $A$ is a geometrically independent set in dimension $N$.
    Assume $s$ is a bijection from $I$ to $A$.
    Assume $s$ is a geometrically independent sequence in dimension $N$.
    Assume $P$ is the plane determined by $A$ in dimension $N$.
    Assume $y\in E\setminus P$.
    Assume $t\in\funs{I_1}{E}$.
    Assume $\restrl{t}{I} = s$ and $t(n + 1) = y$.
    Then $t$ is a geometrically independent sequence in dimension $N$.
\end{lemma}
\begin{proof}
    $n + 1\in\naturalsPlus$ by
        \cref{real_naturals_closed_succ}.
    Show that for all $a$ if
        $a\in\funs{I_1}{\reals}$ and
        $\zeroVector{N}$ is an affine combination of $t$ with coefficients $a$ in dimension $N$ and
        $\sumFiniteSequence{a}{0}$
        then for all $i\in I_1$ we have $a(i) = 0$.
    \begin{subproof}
        Fix $a$.
        Assume $a\in\funs{I_1}{\reals}$ and
            $\zeroVector{N}$ is an affine combination of $t$ with coefficients $a$ in dimension $N$ and
            $\sumFiniteSequence{a}{0}$.
        \begin{byCase}
        \caseOf{$a(n + 1) = 0$.}
            Let $a_0 = \restrl{a}{I}$.
            $a_0\in\funs{I}{\reals}$ by
                \cref{sum_finite_sequence_split_last}.
            $\sumFiniteSequence{a_0}{0}$ by
                \cref{sum_finite_sequence_split_last,real_add_ident,real_add_comm,real_lt_subst_left,real_lt_subst_right}.
            $s\in\funs{I}{E}$ by
                \cref{geometrically_independent_set_50,bijections_to_funs,funs_elim,funs_type_apply,subseteq,funs_intro}.
            $\zeroVector{N}$ is an affine combination of $s$ with coefficients $a_0$ in dimension $N$
                by \cref{affine_zero_relation_restrict_last_zero_50}.
            For all $i\in I_1$ we have $a(i) = 0$
                by \cref{geometrically_independent_sequence_50,initial_segment_naturalsplus,real_naturals_subset_reals,subseteq,real_leq_antisym,initial_segment_succ_split,funs_elim,initial_segment_subset_succ,restrl_of_function_apply,funs_is_function,real_lt_subst_left,real_lt_subst_right}.
        \caseOf{$a(n + 1)\neq 0$.}
            Show that for all $i\in I_1$ we have $a(i) = 0$.
            \begin{subproof}
                Suppose not.
                Contradiction by
                    \cref{setminus_elim_left,affine_relation_last_point_in_plane_50,setminus_elim_right}.
            \end{subproof}
        \end{byCase}
    \end{subproof}
    Follows by \cref{geometrically_independent_sequence_50}.
\end{proof}

% from §50 helper lemma 26: a finite sequence supported at its last index sums to its last value
\begin{lemma}\label{finite_sum_supported_at_last_50}
    Assume $n\in\naturalsPlus$.
    Let $I = \initialSegment{n}$.
    Assume $a\in\funs{I}{\reals}$.
    Assume for all $i\in I$ if $i\neq n$ then $a(i) = 0$.
    Then $\sumFiniteSequence{a}{a(n)}$.
\end{lemma}
\begin{proof}
    $n = 1$ or there exists $m\in\naturalsPlus$ such that $n = m + 1$
        by \cref{real_naturals_predecessor}.
    \begin{byCase}
    \caseOf{$n = 1$.}
        Follows by \cref{sum_finite_sequence_one_term,real_lt_subst_left,real_lt_subst_right}.
    \caseOf{$n\neq 1$.}
        Take $m\in\naturalsPlus$ such that $n = m + 1$.
        Let $J = \initialSegment{m}$.
        Let $a_0 = \restrl{a}{J}$.
        $a_0\in\funs{J}{\reals}$ by
            \cref{real_lt_subst_left,real_lt_subst_right,sum_finite_sequence_exists,sum_finite_sequence_split_last}.
        For all $i\in J$ we have $a_0(i) = 0$
            by \cref{initial_segment_subset_succ,subseteq,real_lt_subst_left,real_lt_subst_right,initial_segment_naturalsplus,real_naturals_subset_reals,real_naturals_lt_succ_local,real_lt_trans,real_lt_irreflexive,funs_elim,restrl_of_function_apply,funs_is_function}.
        $\sumFiniteSequence{a_0}{0}$ by
            \cref{sum_finite_sequence_exists,sum_finite_sequence_all_zero}.
        Follows by
            \cref{real_naturals_closed_succ,initial_segment_naturalsplus,funs_type_apply,real_add_eq_subst,real_add_ident,real_lt_subst_left,real_lt_subst_right,sum_finite_sequence_append_last}.
    \end{byCase}
\end{proof}

% from §50 helper lemma 27: every generator belongs to its affine plane
\begin{lemma}\label{generator_belongs_to_affine_plane_50}
    Assume $N\in\naturalsPlus$.
    Assume $A$ is a geometrically independent set in dimension $N$.
    Assume $P$ is the plane determined by $A$ in dimension $N$.
    Suppose $z\in A$.
    Then $z\in P$.
\end{lemma}
\begin{proof}
    $A$ is finite by \cref{geometrically_independent_set_50}.
    Take $n,s$ such that
        $n\in\naturalsPlus$ and
        $s$ is a bijection from $\initialSegment{n}$ to $A$ and
        $s(n) = z$
        by \cref{finite_inhabited_initial_segment_enumeration_last}.
    Let $I = \initialSegment{n}$.
    Let $E = \realsPower{\initialSegment{N}}$.
    Let $a = \{(i,r) \mid i\in I, r\in\reals \mid
        (i=n\land r=1)\lor(i\neq n\land r=0)\}$.
    $a\in\funs{I}{\reals}$ by
        \cref{relation,pair_eq_iff,rightunique,subseteq,dom_iff,real_add_ident,real_one_in_naturals,real_naturals_subset_reals,subseteq_antisymmetric,function_apply_iff,funs_intro}.
    $a(n)=1$ by
        \cref{initial_segment_naturalsplus,funs_is_function,real_one_in_naturals,real_naturals_subset_reals,subseteq,pair_eq_iff,function_apply_intro}.
    For all $i\in I$ if $i\neq n$ then $a(i)=0$ by
        \cref{funs_is_function,real_add_ident,pair_eq_iff,function_apply_intro}.
    $s\in\funs{I}{E}$ by
        \cref{geometrically_independent_set_50,bijections_to_funs,funs_elim,funs_type_apply,subseteq,funs_intro}.
    Show that for all $j\in\initialSegment{N}$ there exists $b$ such that
        $b\in\funs{I}{\reals}$ and
        for all $i\in I$ we have
            $b(i)=a(i)\rmul\apply{\apply{s}{i}}{j}$ and
            $\sumFiniteSequence{b}{\apply{z}{j}}$.
    \begin{subproof}
            Fix $j\in\initialSegment{N}$.
            Let $b =
                \{(i,a(i)\rmul\apply{\apply{s}{i}}{j}) \mid i\in I\}$.
            $b$ is a function by
                \cref{relation,pair_eq_iff,rightunique}.
            $b\in\funs{I}{\reals}$ by
                \cref{dom_intro,dom_iff,pair_eq_iff,subseteq,subseteq_antisymmetric,function_apply_iff,funs_type_apply,reals_power,tuple_power,real_mul_closed,funs_intro}.
            For all $i\in I$ we have
                $b(i)=a(i)\rmul\apply{\apply{s}{i}}{j}$
                by \cref{funs_is_function,pair_eq_iff,function_apply_intro}.
            For all $i\in I$ if $i\neq n$ then $b(i)=0$ by
                \cref{subseteq,funs_type_apply,reals_power,tuple_power,real_mul_zero,real_lt_subst_left,real_lt_subst_right}.
            $\sumFiniteSequence{b}{\apply{z}{j}}$ by
                \cref{geometrically_independent_set_50,initial_segment_naturalsplus,subseteq,real_lt_subst_left,reals_power,tuple_power,funs_type_apply,real_mul_ident,real_mul_comm,finite_sum_supported_at_last_50,real_lt_subst_right}.
    \end{subproof}
    Follows by
        \cref{finite_sum_supported_at_last_50,real_lt_subst_right,geometrically_independent_set_50,subseteq,affine_combination_value_50,plane_determined_by_points_50}.
\end{proof}

% from §50 helper lemma 28: adjoining a point outside a generated plane preserves set independence
\begin{lemma}\label{independent_set_extension_outside_plane_50}
    Assume $N\in\naturalsPlus$.
    Let $E=\realsPower{\initialSegment{N}}$.
    Assume $A$ is a geometrically independent set in dimension $N$.
    Assume $A$ is inhabited.
    Assume $P$ is the plane determined by $A$ in dimension $N$.
    Assume $y\in E\setminus P$.
    Then $A\union\{y\}$ is a geometrically independent set in dimension $N$.
\end{lemma}
\begin{proof}
    Take $n,s$ such that
        $n\in\naturalsPlus$ and
        $s$ is a bijection from $\initialSegment{n}$ to $A$ and
        $s$ is a geometrically independent sequence in dimension $N$
        by \cref{inhabited_nonempty,geometrically_independent_set_50}.
    Let $I=\initialSegment{n}$.
    Let $I_1=\initialSegment{n+1}$.
    $y\notin A$ by
        \cref{generator_belongs_to_affine_plane_50,setminus_elim_right}.
    Let $r=\{(n+1,y)\}$.
    $s$ is a function and $\dom{s}=I$ by \cref{bijection_is_function,bijections_dom}.
    $r$ is a function and $\dom{r}=\{n+1\}$ by
        \cref{singleton_elim,pair_eq_iff,relation,rightunique,dom_iff,singleton_intro,subseteq,subseteq_antisymmetric}.
    $\dom{s}\inter\dom{r}=\emptyset$ by
        \cref{initial_segment_naturalsplus,real_naturals_subset_reals,real_naturals_closed_succ,real_naturals_lt_succ_local,real_lt_trans,real_lt_irreflexive,real_lt_subst_left,real_lt_subst_right,singleton_iff,subseteq,inter,emptyset,emptyset_subseteq,subseteq_antisymmetric}.
    Let $t=s\union r$.
    $\dom{t}=I_1$ by
        \cref{dom_union,real_lt_subst_left,real_lt_subst_right,union_iff,initial_segment_subset_succ,subseteq,singleton_iff,real_naturals_closed_succ,initial_segment_naturalsplus,initial_segment_succ_split,subseteq_antisymmetric}.
    $t\in\funs{I_1}{A\union\{y\}}$ by
        \cref{union_compatible_functions,emptyset,inter,setminus_elim_left,bijections_to_funs,funs,function_apply_elim,union_iff,dom_intro,bijections_dom,subseteq,funs_type_apply,funs_is_function,function_apply_intro,real_lt_subst_left,union_intro_left,singleton_elim,pair_eq_iff,singleton_iff,union_intro_right,funs_intro}.
    Show that $t$ is injective.
    \begin{subproof}
        It suffices to show that for all $i,j$ such that
            $i\in\dom{t}$ and $j\in\dom{t}$ and $t(i)=t(j)$
        we have $i=j$
            by \cref{funs_is_function,injective_function}.
        Fix $i$.
        Fix $j$.
        Assume $i\in\dom{t}$ and $j\in\dom{t}$ and $t(i)=t(j)$.
        $i\in I$ or $i=n+1$ by \cref{real_lt_subst_left,initial_segment_succ_split}.
        $j\in I$ or $j=n+1$ by \cref{real_lt_subst_left,initial_segment_succ_split}.
        \begin{byCase}
        \caseOf{$i=n+1$.}
            \begin{byCase}
            \caseOf{$j=n+1$.}
                We have $i=j$.
            \caseOf{$j\neq n+1$.}
                Suppose not.
                Contradiction by
                    \cref{funs_is_function,bijections_to_funs,bijections_dom,subseteq,function_apply_elim,union_iff,function_apply_intro,funs_type_apply,singleton_intro,union_intro_right,real_lt_subst_left,real_lt_subst_right}.
            \end{byCase}
        \caseOf{$i\neq n+1$.}
            \begin{byCase}
            \caseOf{$j=n+1$.}
                Suppose not.
                Contradiction by
                    \cref{funs_is_function,bijections_to_funs,bijections_dom,subseteq,function_apply_elim,union_iff,function_apply_intro,funs_type_apply,singleton_intro,union_intro_right,real_lt_subst_left,real_lt_subst_right}.
            \caseOf{$j\neq n+1$.}
                Follows by
                    \cref{funs_is_function,bijections_dom,subseteq,function_apply_elim,union_iff,function_apply_intro,real_lt_subst_left,real_lt_subst_right,bijections,injective_function}.
            \end{byCase}
        \end{byCase}
    \end{subproof}
    $t$ is a bijection from $I_1$ to $A\union\{y\}$ by
        \cref{funs_is_function,union_iff,bijections,surj,initial_segment_subset_succ,subseteq,bijections_dom,function_apply_elim,function_apply_intro,real_lt_subst_right,singleton_iff,real_naturals_closed_succ,initial_segment_naturalsplus,singleton_intro,union_intro_right}.
    $\restrl{t}{I}=s$ by
        \cref{initial_segment_subset_succ,real_lt_subst_left,restrl_of_function_is_function,dom_of_restrl_eq_inter,inter_absorb_supseteq_left,bijection_is_function,bijections_dom,restrl_of_function_apply,funs_is_function,subseteq,function_apply_elim,union_iff,function_apply_intro,fun_subseteq,subseteq_refl,subseteq_antisymmetric}.
    $t(n+1)=y$ by
        \cref{funs_is_function,singleton_intro,union_intro_right,function_apply_intro}.
    $A\union\{y\}\subseteq E$ by
        \cref{setminus_elim_left,geometrically_independent_set_50,union_iff,singleton_iff,subseteq}.
    $t$ is a geometrically independent sequence in dimension $N$ by
        \cref{funs_is_function,funs_type_apply,subseteq,funs_intro,affine_independence_extension_outside_plane_50}.
    Follows by
        \cref{singleton_subset_intro,upair_intro_left,pair_is_finite,subseteq_finite_is_finite,union_of_finite_is_finite,singleton_intro,union_intro_right,emptyset,real_naturals_closed_succ,geometrically_independent_set_50,real_lt_subst_left,real_lt_subst_right}.
\end{proof}

% from §50 helper lemma 29: bijective pullbacks preserve covering-order bounds
\begin{lemma}\label{bijective_pullback_preserves_order_bound_50}
    Assume $f$ is a bijection from $X$ to $Y$.
    Suppose $C$ has order at most $m$ in $Y$.
    Let $D=\{U\in\pow{X}\mid\exists V\in C. U=\preimg{f}{V}\}$.
    Then $D$ has order at most $m$ in $X$.
\end{lemma}
\begin{proof}
    Show that for all $x\in X$ we have $x$ has multiplicity at most $m$ in $D$.
    \begin{subproof}
        Fix $x\in X$.
        Suppose not.
        Take $F,h$ such that
            $F\subseteq D$ and
            $h$ is a bijection from $\initialSegment{m+1}$ to $F$ and
            for all $U\in F$ we have $x\in U$
            by \cref{covering_order_at_most_50,point_multiplicity_at_most_50}.
        Let $R=\{w\in F\times C\mid
            \exists U\in F.\exists V\in C.
                w=(U,V)\land U=\preimg{f}{V}\}$.
        Take $c$ such that
            $c$ is a function on $F$ and
            for all $U\in F$ we have $U\mathrel{R}\apply{c}{U}$
            by \cref{relation,subseteq,times_elem_is_tuple,times_tuple_intro,choice_function}.
        For all $V\in\img{c}{F}$ we have $f(x)\in V$ by
            \cref{bijections_to_funs,img_of_function_elim,inter_elim_right,subseteq,choice_function,pair_eq_iff,preimg_iff,funs_is_function,function_apply_iff,real_lt_subst_left,real_lt_subst_right}.
        Contradiction by
            \cref{injective_function,img_of_function_intro,inter_intro,funs_intro,funs_surj_iff,img_dom,bijections,bijection_circ,img_iff,function_apply_iff,choice_function,pair_eq_iff,subseteq,bijections_to_funs,funs_type_apply,covering_order_at_most_50,point_multiplicity_at_most_50}.
    \end{subproof}
    Follows by
        \cref{bijections_to_funs,preimg_subseteq_dom,funs_elim,subseteq,covering_order_at_most_50}.
\end{proof}

% from §50 helper lemma 30: topological dimension bounds are preserved by homeomorphism
\begin{lemma}\label{homeomorphism_preserves_dimension_bound_50}
    Assume $f$ is a homeomorphism between $X$ and $Y$ with respect to $T$ and $S$.
    Suppose $m$ is a topological dimension bound for $Y$ with respect to $S$.
    Then $m$ is a topological dimension bound for $X$ with respect to $T$.
\end{lemma}
\begin{proof}
    $f$ is a bijection from $X$ to $Y$ and
        $f$ is continuous from $X$ to $Y$ with respect to $T$ and $S$ and
        $\converse{f}$ is continuous from $Y$ to $X$ with respect to $S$ and $T$
        by \cref{homeomorphism}.
    Let $g=\converse{f}$.
    It suffices to show that for all $A$ if
        $A$ is an open covering of $X$ with respect to $T$
        then there exists $B$ such that
            $B$ is an open refinement of $A$ on $X$ with respect to $T$ and
            $B$ has order at most $m+1$ in $X$
        by \cref{continuous,topological_dimension_bound_50}.
        Fix $A$.
        Suppose $A$ is an open covering of $X$ with respect to $T$.
        Let $C=\{V\in\pow{Y}\mid\exists U\in A. V=\preimg{g}{U}\}$.
        $C$ is an open covering of $Y$ with respect to $S$ by
            \cref{bijection_converse_is_bijection,open_covering,continuous,real_lt_subst_left,bijections_to_funs,funs_type_apply,covering,subseteq,unions_iff,bijections_dom,preimg_subseteq_dom,real_lt_subst_right,powerset_intro,funs_is_function,funs,function_apply_elim,preimg_iff}.
        Take $D$ such that
            $D$ is an open refinement of $C$ on $Y$ with respect to $S$ and
            $D$ has order at most $m+1$ in $Y$
            by \cref{topological_dimension_bound_50}.
        Let $B=\{U\in\pow{X}\mid\exists V\in D. U=\preimg{f}{V}\}$.
        $B$ has order at most $m+1$ in $X$ by
            \cref{bijective_pullback_preserves_order_bound_50}.
        Show that $B$ is an open refinement of $A$ on $X$ with respect to $T$.
        \begin{subproof}
            For all $W\in B$ there exists $U\in A$ such that $W\subseteq U$ by
                \cref{open_refinement_on,refinement_on,subseteq,powerset_elim,real_lt_subst_left,continuous,funs_is_function,function_apply_iff,preimg_iff,bijections_to_funs,funs,function_apply_elim,converse_iff,homeomorphism,bijections,bijective_converse_are_funs,function_apply_intro}.
            $B$ is a refinement of $A$ on $X$ by
                \cref{open_covering,open_in,topology_on,powerset_elim,continuous,preimg_subseteq_dom,powerset_intro,subseteq,refinement_on}.
            For all $x\in X$ we have $x\in\unions{B}$ by
                \cref{continuous,funs_type_apply,open_refinement_on,covering,subseteq,unions_iff,funs,preimg_subseteq_dom,real_lt_subst_right,powerset_intro,funs_is_function,function_apply_elim,preimg_iff}.
            For all $W\in B$ we have $W$ is open in $X$ with respect to $T$
                by \cref{continuous,open_refinement_on}.
            Follows by \cref{subseteq,covering,open_refinement_on}.
        \end{subproof}
\end{proof}

% from §50 helper lemma 31: compact restrictions of homeomorphisms are homeomorphisms onto their images
\begin{lemma}\label{compact_homeomorphism_restriction_50}
    Assume $f$ is a homeomorphism between $X$ and $Y$ with respect to $T$ and $S$.
    Assume $Y$ is Hausdorff with respect to $S$.
    Assume $A$ is a subspace of $X$ with respect to $T$.
    Let $T_A=\subspaceTopology{T}{A}$.
    Assume $A$ is compact with respect to $T_A$.
    Let $g=\restrl{f}{A}$.
    Let $B=\img{f}{A}$.
    Let $S_B=\subspaceTopology{S}{B}$.
    Then $g$ is a homeomorphism between $A$ and $B$ with respect to $T_A$ and $S_B$.
\end{lemma}
\begin{proof}
    Follows by
        \cref{homeomorphism,continuous_restrict_domain,subspace_of,bijections_to_funs,funs,restrl_img,inter_idempotent,continuous,rels_ran_subseteq,img_subseteq_ran,subseteq_transitive,continuous_restrict_range,subseteq_refl,real_lt_subst_left,bijections,restrl_injective,funs_surj_iff,img_dom,subspace_hausdorff,compact_hausdorff_bijection_homeomorphism}.
\end{proof}

% from §50 helper lemma 32: the empty space has every natural dimension bound
\begin{lemma}\label{empty_space_dimension_bound_50}
    Assume $T$ is a topology on $X$.
    Assume $X=\emptyset$.
    Assume $m\in\naturals$.
    Then $m$ is a topological dimension bound for $X$ with respect to $T$.
\end{lemma}
\begin{proof}
    It suffices to show that for all $A$ if
        $A$ is an open covering of $X$ with respect to $T$
    then there exists $B$ such that
        $B$ is an open refinement of $A$ on $X$ with respect to $T$ and
        $B$ has order at most $m+1$ in $X$
        by \cref{topological_dimension_bound_50}.
    Fix $A$.
    Assume $A$ is an open covering of $X$ with respect to $T$.
    Let $B=\emptyset$.
    $B$ is an open refinement of $A$ on $X$ with respect to $T$ and
    $B$ has order at most $m+1$ in $X$ by
        \cref{open_covering,open_in,topology_on,powerset_elim,powerset_intro,refinement_on,covering,open_refinement_on,naturals,union_iff,singleton_iff,real_one_in_naturals,real_naturals_subset_reals,real_add_ident,real_lt_subst_left,real_naturals_closed_succ,emptyset_subseteq,emptyset,subseteq,covering_order_at_most_50}.
\end{proof}

% from §50 helper lemma 33: the image of a finite euclidean sequence consists of euclidean points
\begin{lemma}\label{finite_euclidean_sequence_image_subset_50}
    Assume $N,n\in\naturalsPlus$.
    Let $E=\realsPower{\initialSegment{N}}$.
    Assume $x\in\funs{\initialSegment{n}}{E}$.
    Then $\img{x}{\initialSegment{n}}\subseteq E$.
\end{lemma}
\begin{proof}
    Follows by
        \cref{funs,rels_ran_subseteq,img_subseteq_ran,subseteq_transitive}.
\end{proof}

% from §50 helper lemma 34: the image of a finite sequence is finite
\begin{lemma}\label{finite_sequence_image_finite_50}
    Assume $n\in\naturalsPlus$.
    Assume $x$ is a function on $\initialSegment{n}$.
    Then $\img{x}{\initialSegment{n}}$ is finite.
\end{lemma}
\begin{proof}
    Follows by \cref{initial_segment_finite,finite_image_function}.
\end{proof}

% from §50 helper definition 20: injectivity deviation set
\begin{definition}\label{injectivity_deviation_set_50}
    $\injectivityDeviationSet{f}{X}{d} =
        \{s\in\reals \mid \text{there exists $z\in\img{f}{X}$ such that $s$ is a diameter of $\preimg{f}{\{z\}}$ in $X$ with respect to $d$}\}$.
\end{definition}

% from §50 helper definition 21: deviation from injectivity
\begin{definition}\label{deviation_from_injectivity_50}
    $r$ is the deviation from injectivity of $f$ on $X$ with respect to $d$ iff
        $f$ is a function on $X$ and
        $d$ is a metric on $X$ and
        $r$ is a supremum of $\injectivityDeviationSet{f}{X}{d}$.
\end{definition}

% from §50 helper definition 22: maps with small deviation from injectivity
\begin{definition}\label{small_deviation_class_50}
    $U$ is the small-deviation class below $\epsilon$ on $X$ to $Y$ with respect to $T$ and $S$ and $d$ iff
        $U = \{f\in\funs{X}{Y} \mid \text{$f$ is continuous from $X$ to $Y$ with respect to $T$ and $S$ and
            there exists $r\in\reals$ such that $r$ is the deviation from injectivity of $f$ on $X$ with respect to $d$ and $r < \epsilon$}\}$.
\end{definition}

% from §50 helper definition 23: maps separating all pairs at a fixed metric scale
\begin{definition}\label{metric_scale_separating_class_50}
    $U$ is the separating class at scale $\epsilon$ on $X$ to $Y$ with respect to $T$ and $S$ and $d$ iff
        $U=\{f\in\funs{X}{Y}\mid
            \text{$f$ is continuous from $X$ to $Y$ with respect to $T$ and $S$ and
            for all $x,y\in X$ if $\epsilon\leq d((x,y))$ then $f(x)\neq f(y)$}\}$.
\end{definition}

% from §50 helper lemma 36: separation at every reciprocal scale implies injectivity
\begin{lemma}\label{reciprocal_scale_separation_implies_injective_50}
    Assume $d$ is a metric on $X$.
    Assume $f\in\funs{X}{Y}$.
    Suppose for all $n\in\naturalsPlus$ we have
        for all $x,y\in X$ if $1/n\leq d((x,y))$ then $f(x)\neq f(y)$.
    Then $f$ is injective.
\end{lemma}
\begin{proof}
    It suffices to show that for all $x,y$ such that
        $x\in\dom{f}$ and $y\in\dom{f}$ and $f(x)=f(y)$
    we have $x=y$ by
        \cref{funs_is_function,injective_function}.
    Fix $x$.
    Fix $y$.
    Assume $x\in\dom{f}$ and $y\in\dom{f}$ and $f(x)=f(y)$.
    Suppose not.
    Take $n\in\naturalsPlus$ such that $1/n<d((x,y))$ by
        \cref{funs,metric_on,metric_zero_characterization,metric_nonnegative,real_leq_split,funs_type_apply,times_tuple_intro,real_small_reciprocal}.
    Contradiction by \cref{funs,real_lt_imp_leq}.
\end{proof}

% from §50 helper lemma 37: a finite sum supported at one index equals its supported value
\begin{lemma}\label{finite_sum_single_support_50}
    Assume $n\in\naturalsPlus$.
    Let $I=\initialSegment{n}$.
    Assume $k\in I$.
    Assume $a\in\funs{I}{\reals}$.
    Suppose for all $i\in I$ if $i\neq k$ then $a(i)=0$.
    Then $\sumFiniteSequence{a}{a(k)}$.
\end{lemma}
\begin{proof}
    Let $A=\{r\in\naturalsPlus\mid
        \text{for all $q\in\initialSegment{r}$ we have for all
        $c\in\funs{\initialSegment{r}}{\reals}$ if
        for all $i\in\initialSegment{r}$ if $i\neq q$ then $c(i)=0$,
        then $\sumFiniteSequence{c}{c(q)}$}\}$.
    Show that for all $r\in A$ we have $r+1\in A$.
    \begin{subproof}
        Fix $r\in A$.
        Let $J=\initialSegment{r}$.
        Let $J_1=\initialSegment{r+1}$.
        Show that for all $q\in J_1$ we have for all
            $c\in\funs{J_1}{\reals}$ if
            for all $i\in J_1$ if $i\neq q$ then $c(i)=0$,
        then $\sumFiniteSequence{c}{c(q)}$.
        \begin{subproof}
            Fix $q\in J_1$.
            Fix $c\in\funs{J_1}{\reals}$.
            Assume for all $i\in J_1$ if $i\neq q$ then $c(i)=0$.
            \begin{byCase}
            \caseOf{$q=r+1$.}
                Follows by
                    \cref{real_naturals_closed_succ,real_lt_subst_right,finite_sum_supported_at_last_50}.
            \caseOf{$q\neq r+1$.}
                Let $c_0=\restrl{c}{J}$.
                $c_0\in\funs{J}{\reals}$ by
                    \cref{funs_elim,restrl_of_function_is_function,dom_of_restrl_eq_inter,initial_segment_subset_succ,inter_absorb_supseteq_left,subseteq,funs_type_apply,restrl_of_function_apply,funs_intro}.
                $\sumFiniteSequence{c_0}{c_0(q)}$ by
                    \cref{initial_segment_naturalsplus,real_naturals_leq_succ_elim,funs_is_function,funs,initial_segment_subset_succ,subseteq,restrl_of_function_apply}.
                Follows by
                    \cref{real_naturals_leq_succ_elim,funs_is_function,funs,initial_segment_subset_succ,restrl_of_function_apply,real_naturals_closed_succ,initial_segment_naturalsplus,sum_finite_sequence_append_zero,real_lt_subst_left}.
            \end{byCase}
        \end{subproof}
        Follows by \cref{real_naturals_closed_succ}.
    \end{subproof}
    Follows by
        \cref{real_one_in_naturals,initial_segment_one,singleton_iff,sum_finite_sequence_one_term,real_lt_subst_right,subseteq,real_naturals_induction}.
\end{proof}

% from §50 helper definition 24: the zero-or-coordinate point of the standard affine frame
\begin{definition}\label{standard_affine_frame_point_50}
    $\standardAffineFramePoint{N}{i}=
        \{(j,r)\mid j\in\initialSegment{N},r\in\reals\mid
            \text{either $i=1$ and $r=0$ or
            $i\neq 1$ and either $j+1=i$ and $r=1$ or
            $j+1\neq i$ and $r=0$}\}$.
\end{definition}

% from §50 helper definition 25: the standard affine frame sequence
\begin{definition}\label{standard_affine_frame_sequence_50}
    $\standardAffineFrame{N}=
        \{(i,\standardAffineFramePoint{N}{i})\mid
            i\in\initialSegment{N+1}\}$.
\end{definition}

% from §50 helper lemma 38: standard affine frame points are euclidean points
\begin{lemma}\label{standard_affine_frame_point_coordinates_50}
    Assume $N\in\naturalsPlus$.
    Let $I=\initialSegment{N}$.
    Let $I_1=\initialSegment{N+1}$.
    Assume $i\in I_1$.
    Let $v=\standardAffineFramePoint{N}{i}$.
    Then $v\in\realsPower{I}$.
\end{lemma}
\begin{proof}
    Follows by
        \cref{standard_affine_frame_point_50,relation,pair_eq_iff,rightunique,subseteq,dom_iff,real_add_ident,real_one_in_naturals,real_naturals_subset_reals,subseteq_antisymmetric,function_apply_iff,funs_intro,reals_power,tuple_power}.
\end{proof}

% from §50 helper lemma 39: the first standard affine frame point is zero
\begin{lemma}\label{standard_affine_frame_point_zero_50}
    Assume $N\in\naturalsPlus$.
    Let $I=\initialSegment{N}$.
    Let $v=\standardAffineFramePoint{N}{1}$.
    Then for all $j\in I$ we have $v(j)=0$.
\end{lemma}
\begin{proof}
    Follows by
        \cref{real_one_in_naturals,real_one_leq_naturalsplus,initial_segment_naturalsplus,real_naturals_closed_succ,standard_affine_frame_point_coordinates_50,reals_power,tuple_power,funs_is_function,funs,function_apply_elim,standard_affine_frame_point_50,pair_eq_iff}.
\end{proof}

% from §50 helper lemma 40: a matching standard affine frame coordinate is one
\begin{lemma}\label{standard_affine_frame_point_one_50}
    Assume $N\in\naturalsPlus$.
    Let $I=\initialSegment{N}$.
    Assume $j\in I$.
    Let $v=\standardAffineFramePoint{N}{j+1}$.
    Then $v(j)=1$.
\end{lemma}
\begin{proof}
    Follows by
        \cref{initial_segment_naturalsplus,real_naturals_subset_reals,subseteq,real_one_in_naturals,real_leq_add,real_naturals_closed_succ,standard_affine_frame_point_coordinates_50,reals_power,tuple_power,funs_is_function,real_one_leq_naturalsplus,real_naturals_lt_succ_local,real_lt_trans,real_lt_irreflexive,real_lt_subst_left,real_lt_subst_right,standard_affine_frame_point_50,pair_eq_iff,function_apply_intro}.
\end{proof}

% from §50 helper lemma 41: a nonmatching nonzero frame coordinate is zero
\begin{lemma}\label{standard_affine_frame_point_off_coordinate_50}
    Assume $N\in\naturalsPlus$.
    Let $I=\initialSegment{N}$.
    Let $I_1=\initialSegment{N+1}$.
    Assume $i\in I_1$.
    Assume $i\neq 1$.
    Assume $j\in I$ and $j+1\neq i$.
    Let $v=\standardAffineFramePoint{N}{i}$.
    Then $v(j)=0$.
\end{lemma}
\begin{proof}
    Follows by
        \cref{standard_affine_frame_point_coordinates_50,reals_power,tuple_power,funs_is_function,real_add_ident,standard_affine_frame_point_50,pair_eq_iff,function_apply_intro}.
\end{proof}

% from §50 helper lemma 42: the standard affine frame is a euclidean sequence
\begin{lemma}\label{standard_affine_frame_is_function_50}
    Assume $N\in\naturalsPlus$.
    Let $I=\initialSegment{N}$.
    Let $I_1=\initialSegment{N+1}$.
    Let $E=\realsPower{I}$.
    Let $s=\standardAffineFrame{N}$.
    Then $s\in\funs{I_1}{E}$ and
        for all $i\in I_1$ we have
            $s(i)=\standardAffineFramePoint{N}{i}$.
\end{lemma}
\begin{proof}
    Follows by
        \cref{standard_affine_frame_sequence_50,relation,pair_eq_iff,rightunique,function_apply_intro,dom_intro,dom_iff,subseteq,subseteq_antisymmetric,standard_affine_frame_point_coordinates_50,funs_intro}.
\end{proof}

% from §50 helper lemma 43: the standard affine frame is geometrically independent
\begin{lemma}\label{standard_affine_frame_independent_50}
    Assume $N\in\naturalsPlus$.
    Let $I=\initialSegment{N}$.
    Let $I_1=\initialSegment{N+1}$.
    Let $E=\realsPower{I}$.
    Let $s=\standardAffineFrame{N}$.
    Then $s$ is a geometrically independent sequence in dimension $N$.
\end{lemma}
\begin{proof}
    Show that for all $a$ if
        $a\in\funs{I_1}{\reals}$ and
        $\zeroVector{N}$ is an affine combination of $s$ with coefficients $a$ in dimension $N$ and
        $\sumFiniteSequence{a}{0}$
        then for all $i\in I_1$ we have $a(i)=0$.
    \begin{subproof}
        Fix $a$.
        Assume $a\in\funs{I_1}{\reals}$ and
            $\zeroVector{N}$ is an affine combination of $s$ with coefficients $a$ in dimension $N$ and
            $\sumFiniteSequence{a}{0}$.
        Take $k$ such that
            $k\in\naturalsPlus$ and
            $s\in\funs{\initialSegment{k}}{E}$ and
            $a\in\funs{\initialSegment{k}}{\reals}$ and
            $\zeroVector{N}\in E$ and
            for all $j\in I$ there exists $b$ such that
                $b\in\funs{\initialSegment{k}}{\reals}$ and
                for all $i\in\initialSegment{k}$ we have
                    $b(i)=a(i)\rmul\apply{\apply{s}{i}}{j}$ and
                $\sumFiniteSequence{b}{\apply{\zeroVector{N}}{j}}$
            by \cref{affine_combination_value_50}.
        $k=N+1$ by
            \cref{real_naturals_closed_succ,funs_elim,initial_segment_naturalsplus,real_lt_subst_left,real_lt_subst_right,real_naturals_subset_reals,subseteq,real_leq_antisym}.
        Show that for all $j\in I$ we have $a(j+1)=0$.
        \begin{subproof}
            Fix $j\in I$.
            Take $b$ such that
                $b\in\funs{\initialSegment{k}}{\reals}$ and
                for all $i\in\initialSegment{k}$ we have
                    $b(i)=a(i)\rmul\apply{\apply{s}{i}}{j}$ and
                $\sumFiniteSequence{b}{\apply{\zeroVector{N}}{j}}$
                by \cref{subseteq}.
            For all $i\in I_1$ if $i\neq j+1$ then $b(i)=0$ by
                \cref{standard_affine_frame_is_function_50,standard_affine_frame_point_zero_50,standard_affine_frame_point_off_coordinate_50,funs_type_apply,reals_power,tuple_power,subseteq,real_mul_zero,real_mul_comm,real_lt_subst_left,real_lt_subst_right}.
            $b(j+1)=0$ by
                \cref{initial_segment_naturalsplus,real_naturals_subset_reals,subseteq,real_one_in_naturals,real_leq_add,real_naturals_closed_succ,reals_power,tuple_power,funs_is_function,zero_vector_50,pair_eq_iff,function_apply_intro,real_lt_subst_left,real_lt_subst_right,finite_sum_single_support_50,sum_finite_sequence_unique}.
            Thus $a(j+1)=0$ by
                \cref{initial_segment_naturalsplus,real_naturals_subset_reals,subseteq,real_one_in_naturals,real_leq_add,real_naturals_closed_succ,standard_affine_frame_is_function_50,standard_affine_frame_point_one_50,funs_type_apply,real_mul_ident,real_mul_comm,real_lt_subst_left,real_lt_subst_right}.
        \end{subproof}
        For all $i\in I_1$ if $i\neq 1$ then $a(i)=0$ by
            \cref{initial_segment_naturalsplus,real_naturals_predecessor,real_naturals_subset_reals,subseteq,real_lt_trichotomy,real_one_in_naturals,real_lt_add,real_lt_subst_left,real_add_closed,real_lt_trans,real_lt_irreflexive,real_lt_subst_right}.
        $a(1)=0$ by
            \cref{finite_sum_single_support_50,real_one_in_naturals,real_one_leq_naturalsplus,initial_segment_naturalsplus,sum_finite_sequence_unique}.
        Follows by \cref{real_lt_subst_left,subseteq}.
    \end{subproof}
    Follows by
        \cref{standard_affine_frame_is_function_50,real_naturals_closed_succ,geometrically_independent_sequence_50}.
\end{proof}

% from §50 helper lemma 52: package a controlled indexed cover into existential form
\begin{lemma}\label{controlled_finite_indexed_cover_packaging_50}
    Assume $n\in\naturalsPlus$.
    Assume $U$ is a bijection from $\initialSegment{n}$ to
        $\img{U}{\initialSegment{n}}$.
    Assume $\img{U}{\initialSegment{n}}$ is a covering of $X$.
    Assume $\img{U}{\initialSegment{n}}$ has order at most $m+1$ in $X$.
    Assume for all $i\in\initialSegment{n}$ we have
        $U(i)$ is open in $X$ with respect to $T$ and $U(i)$ is inhabited and
        for all $x,y\in U(i)$ we have
            $d((x,y))<\epsilon$ and $e((f(x),f(y)))<\delta$.
    Then there exist $n_0,U_0$ such that
        $n_0\in\naturalsPlus$ and
        $U_0$ is a bijection from $\initialSegment{n_0}$ to
            $\img{U_0}{\initialSegment{n_0}}$ and
        $\img{U_0}{\initialSegment{n_0}}$ is a covering of $X$ and
        $\img{U_0}{\initialSegment{n_0}}$ has order at most $m+1$ in $X$ and
        for all $i\in\initialSegment{n_0}$ we have
            $U_0(i)$ is open in $X$ with respect to $T$ and $U_0(i)$ is inhabited and
            for all $x,y\in U_0(i)$ we have
                $d((x,y))<\epsilon$ and $e((f(x),f(y)))<\delta$.
\end{lemma}
\begin{proof}
    Follows by \cref{subseteq_refl}.
\end{proof}

% from §50 helper lemma 53: noninitial affine-frame indices have coordinate predecessors
\begin{lemma}\label{affine_frame_index_predecessor_50}
    Assume $N\in\naturalsPlus$.
    Assume $i\in\initialSegment{N+1}$ and $i\neq 1$.
    Then there exists $j\in\initialSegment{N}$ such that $i=j+1$.
\end{lemma}
\begin{proof}
    Follows.
\end{proof}

% from §50 helper lemma 54: the standard affine frame has distinct points
\begin{lemma}\label{standard_affine_frame_injective_50}
    Assume $N\in\naturalsPlus$.
    Let $I=\initialSegment{N}$.
    Let $I_1=\initialSegment{N+1}$.
    Let $E=\realsPower{I}$.
    Let $s=\standardAffineFrame{N}$.
    Then $s$ is injective.
\end{lemma}
\begin{proof}
    It suffices to show that for all $i,k$ such that
        $i\in\dom{s}$ and $k\in\dom{s}$ and $s(i)=s(k)$
        we have $i=k$ by
        \cref{standard_affine_frame_is_function_50,funs_is_function,injective_function}.
    Fix $i$.
    Fix $k$.
    Assume $i\in\dom{s}$ and $k\in\dom{s}$ and $s(i)=s(k)$.
    $i,k\in I_1$ by
        \cref{standard_affine_frame_is_function_50,funs}.
    Follows by
        \cref{affine_frame_index_predecessor_50,funs_type_apply,reals_power,tuple_power,real_add_ident,real_one_in_naturals,real_naturals_subset_reals,subseteq,standard_affine_frame_is_function_50,standard_affine_frame_point_zero_50,standard_affine_frame_point_one_50,standard_affine_frame_point_off_coordinate_50,real_zero_lt_one,real_lt_irreflexive,real_lt_subst_left,real_lt_subst_right}.
\end{proof}

% from §50 helper lemma 55: the standard affine-frame image is independent
\begin{lemma}\label{standard_affine_frame_image_independent_50}
    Assume $N\in\naturalsPlus$.
    Let $I=\initialSegment{N}$.
    Let $I_1=\initialSegment{N+1}$.
    Let $E=\realsPower{I}$.
    Let $s=\standardAffineFrame{N}$.
    Let $B=\img{s}{I_1}$.
    Then $B$ is a geometrically independent set in dimension $N$.
\end{lemma}
\begin{proof}
    Follows by
        \cref{real_naturals_closed_succ,standard_affine_frame_is_function_50,standard_affine_frame_injective_50,funs_elim,img_of_function_intro,inter_intro,funs_intro,funs_surj_iff,img_dom,bijections,img_iff,function_apply_iff,funs_type_apply,subseteq,real_one_in_naturals,real_one_leq_naturalsplus,initial_segment_naturalsplus,emptyset,initial_segment_finite,finite_image_function,standard_affine_frame_independent_50,geometrically_independent_set_50,real_lt_subst_left,real_lt_subst_right}.
\end{proof}

% from §50 helper lemma 56: finite real sums are invariant under enumeration
\begin{lemma}\label{finite_sum_reindex_bijection_50}
    Assume $A$ is finite and $A$ is inhabited.
    Assume $n,m\in\naturalsPlus$.
    Assume $g$ is a bijection from $\initialSegment{n}$ to $A$.
    Assume $h$ is a bijection from $\initialSegment{m}$ to $A$.
    Assume $c\in\funs{A}{\reals}$.
    Assume $a\in\funs{\initialSegment{n}}{\reals}$ and
        $\sumFiniteSequence{a}{r}$ and
        for all $i\in\initialSegment{n}$ we have $a(i)=c(g(i))$.
    Assume $b\in\funs{\initialSegment{m}}{\reals}$ and
        $\sumFiniteSequence{b}{s}$ and
        for all $i\in\initialSegment{m}$ we have $b(i)=c(h(i))$.
    Then $r=s$.
\end{lemma}
\begin{proof}
    Let $F=\{(\alpha,\{(0,c(\alpha))\})\mid \alpha\in A\}$.
    $F$ is a function by
        \cref{relation,pair_eq_iff,rightunique}.
    $\dom{F}=A$ by
        \cref{dom_intro,dom_iff,pair_eq_iff,subseteq,subseteq_antisymmetric}.
    For all $\alpha\in A$ we have
        $F(\alpha)\in\funs{\{0\}}{\reals}$ and
        $\apply{F(\alpha)}{0}=c(\alpha)$ by
        \cref{pair_eq_iff,function_apply_intro,relation,singleton_elim,rightunique,dom_intro,dom_iff,singleton_intro,subseteq,subseteq_antisymmetric,funs_type_apply,singleton_iff,funs_intro,real_lt_subst_left,real_lt_subst_right}.
    Take $R$ such that $R$ is a well order relation on $A$ by
        \cref{well_ordering_theorem_local_finiteness}.
    Take $k,u$ such that
        $k\in\naturalsPlus$ and
        $u$ is a bijection from $\initialSegment{k}$ to $A$ and
        $u$ is increasing on $\initialSegment{k}$ with respect to $R$
        by \cref{subseteq_refl,finite_well_ordered_enumeration}.
    $u$ is a well-ordered enumeration of $A$ in $A$ with respect to $R$ and $k$
        by \cref{well_ordered_enumeration_of_subset,subseteq_refl}.
    Let $d=\{(i,c(u(i)))\mid i\in\initialSegment{k}\}$.
    $d\in\funs{\initialSegment{k}}{\reals}$ by
        \cref{relation,pair_eq_iff,rightunique,dom_intro,dom_iff,subseteq,subseteq_antisymmetric,bijections,surj,funs_type_apply,funs_is_function,function_apply_intro,real_lt_subst_left,funs_intro}.
    For all $i\in\initialSegment{k}$ we have $d(i)=c(u(i))$ by
        \cref{funs_is_function,pair_eq_iff,function_apply_intro}.
    Take $t\in\reals$ such that $\sumFiniteSequence{d}{t}$ by
        \cref{sum_finite_sequence_exists}.
    $a$ is a real sum sequence of $F$ at $0$ along $g$
        with respect to $n$ and result $r$ by
        \cref{bijections,surj,funs_type_apply,subseteq,real_lt_subst_left,real_lt_subst_right,real_sum_sequence_along_enumeration}.
    $b$ is a real sum sequence of $F$ at $0$ along $h$
        with respect to $m$ and result $s$ by
        \cref{bijections,surj,funs_type_apply,subseteq,real_lt_subst_left,real_lt_subst_right,real_sum_sequence_along_enumeration}.
    $d$ is a real sum sequence of $F$ at $0$ along $u$
        with respect to $k$ and result $t$ by
        \cref{bijections,surj,funs_type_apply,subseteq,real_lt_subst_left,real_lt_subst_right,real_sum_sequence_along_enumeration}.
    Follows by
        \cref{well_ordered_same_set_sum_comparison,real_lt_subst_left,real_lt_subst_right}.
\end{proof}

% from §50 helper lemma 57: affine coefficients transfer between enumerations
\begin{lemma}\label{coefficient_sequence_reindex_50}
    Assume $A$ is finite and $A$ is inhabited.
    Assume $n,m\in\naturalsPlus$.
    Assume $g$ is a bijection from $\initialSegment{n}$ to $A$.
    Assume $h$ is a bijection from $\initialSegment{m}$ to $A$.
    Assume $a\in\funs{\initialSegment{n}}{\reals}$ and
        $\sumFiniteSequence{a}{r}$.
    Then there exist $c,b$ such that
        $c\in\funs{A}{\reals}$ and
        $b\in\funs{\initialSegment{m}}{\reals}$ and
        $\sumFiniteSequence{b}{r}$ and
        for all $i\in\initialSegment{n}$ we have $a(i)=c(g(i))$ and
        for all $i\in\initialSegment{m}$ we have $b(i)=c(h(i))$.
\end{lemma}
\begin{proof}
    Let $I=\initialSegment{n}$.
    Let $J=\initialSegment{m}$.
    Let $c=\{(g(i),a(i))\mid i\in I\}$.
    $c\in\funs{A}{\reals}$ by
        \cref{relation,pair_eq_iff,rightunique,bijections_dom,subseteq,bijections,injective_function,real_lt_subst_left,real_lt_subst_right,dom_iff,surj,funs_type_apply,dom_intro,subseteq_antisymmetric,funs_is_function,function_apply_intro,funs_intro}.
    For all $i\in I$ we have $a(i)=c(g(i))$ by
        \cref{funs_is_function,pair_eq_iff,function_apply_intro}.
    Let $b=\{(i,c(h(i)))\mid i\in J\}$.
    $b\in\funs{J}{\reals}$ by
        \cref{relation,pair_eq_iff,rightunique,dom_intro,dom_iff,subseteq,subseteq_antisymmetric,bijections,surj,funs_type_apply,funs_is_function,function_apply_intro,real_lt_subst_left,funs_intro}.
    For all $i\in J$ we have $b(i)=c(h(i))$ by
        \cref{funs_is_function,pair_eq_iff,function_apply_intro}.
    Take $q\in\reals$ such that $\sumFiniteSequence{b}{q}$ by
        \cref{sum_finite_sequence_exists}.
    $\sumFiniteSequence{b}{r}$ by
        \cref{finite_sum_reindex_bijection_50,real_lt_subst_right}.
\end{proof}

% from §50 helper lemma 58: normalize affine-plane coordinates to a fixed enumeration
\begin{lemma}\label{plane_coordinates_fixed_enumeration_50}
    Assume $N,n\in\naturalsPlus$.
    Let $I=\initialSegment{n}$.
    Let $E=\realsPower{\initialSegment{N}}$.
    Assume $A$ is a geometrically independent set in dimension $N$.
    Assume $A$ is inhabited.
    Assume $g$ is a bijection from $I$ to $A$.
    Assume $P$ is the plane determined by $A$ in dimension $N$.
    Assume $y\in P$.
    Then there exists $a$ such that
        $a\in\funs{I}{\reals}$ and
        $\sumFiniteSequence{a}{1}$ and
        $y$ is an affine combination of $g$ with coefficients $a$ in dimension $N$.
\end{lemma}
\begin{proof}
    $A$ is finite and $A\subseteq E$ by
        \cref{geometrically_independent_set_50}.
    Take $k,s,a_0$ such that
        $k\in\naturalsPlus$ and
        $s$ is a bijection from $\initialSegment{k}$ to $A$ and
        $a_0\in\funs{\initialSegment{k}}{\reals}$ and
        $\sumFiniteSequence{a_0}{1}$ and
        $y$ is an affine combination of $s$ with coefficients $a_0$ in dimension $N$
        by \cref{plane_determined_by_points_50}.
    Take $c,a$ such that
        $c\in\funs{A}{\reals}$ and
        $a\in\funs{I}{\reals}$ and
        $\sumFiniteSequence{a}{1}$ and
        for all $i\in\initialSegment{k}$ we have $a_0(i)=c(s(i))$ and
        for all $i\in I$ we have $a(i)=c(g(i))$
        by \cref{coefficient_sequence_reindex_50}.
    $g\in\funs{I}{E}$ by
        \cref{bijections_to_funs,funs_elim,funs_type_apply,subseteq,funs_intro}.
    Take $l$ such that
        $l\in\naturalsPlus$ and
        $s\in\funs{\initialSegment{l}}{E}$ and
        $a_0\in\funs{\initialSegment{l}}{\reals}$ and
        $y\in E$ and
        for all $j\in\initialSegment{N}$ there exists $e$ such that
            $e\in\funs{\initialSegment{l}}{\reals}$ and
            for all $i\in\initialSegment{l}$ we have
                $e(i)=a_0(i)\rmul\apply{\apply{s}{i}}{j}$ and
            $\sumFiniteSequence{e}{\apply{y}{j}}$
        by \cref{affine_combination_value_50}.
    $\initialSegment{l}=\initialSegment{k}$ by
        \cref{bijections_to_funs,funs_elim}.
    Show that for all $j\in\initialSegment{N}$ there exists $d$ such that
        $d\in\funs{I}{\reals}$ and
        for all $i\in I$ we have
            $d(i)=a(i)\rmul\apply{\apply{g}{i}}{j}$ and
        $\sumFiniteSequence{d}{\apply{y}{j}}$.
    \begin{subproof}
        Fix $j\in\initialSegment{N}$.
        Take $e$ such that
            $e\in\funs{\initialSegment{l}}{\reals}$ and
            for all $i\in\initialSegment{l}$ we have
                $e(i)=a_0(i)\rmul\apply{\apply{s}{i}}{j}$ and
            $\sumFiniteSequence{e}{\apply{y}{j}}$.
        $\sumFiniteSequence{e}{\apply{y}{j}}$ by
            \cref{real_one_in_naturals,real_one_leq_naturalsplus,initial_segment_naturalsplus,subseteq}.
        For all $i\in\initialSegment{k}$ we have
            $e(i)=a_0(i)\rmul\apply{\apply{s}{i}}{j}$ by
            \cref{real_lt_subst_left,real_lt_subst_right}.
        Let $v=\{(z,c(z)\rmul z(j))\mid z\in A\}$.
        $v\in\funs{A}{\reals}$ by
            \cref{relation,pair_eq_iff,rightunique,dom_intro,dom_iff,subseteq,subseteq_antisymmetric,funs_type_apply,reals_power,tuple_power,real_mul_closed,function_apply_intro,real_lt_subst_left,funs_intro}.
        For all $i\in\initialSegment{k}$ we have $e(i)=v(s(i))$ by
            \cref{bijections,surj,funs_type_apply,funs_is_function,pair_eq_iff,function_apply_intro,subseteq,real_lt_subst_left,real_lt_subst_right}.
        Let $d=\{(i,v(g(i)))\mid i\in I\}$.
        $d\in\funs{I}{\reals}$ by
            \cref{relation,pair_eq_iff,rightunique,dom_intro,dom_iff,subseteq,subseteq_antisymmetric,bijections,surj,funs_type_apply,funs_is_function,function_apply_intro,real_lt_subst_left,funs_intro}.
        For all $i\in I$ we have $d(i)=v(g(i))$ by
            \cref{funs_is_function,pair_eq_iff,function_apply_intro}.
        Take $q\in\reals$ such that $\sumFiniteSequence{d}{q}$ by
            \cref{sum_finite_sequence_exists}.
        $\sumFiniteSequence{d}{\apply{y}{j}}$ by
            \cref{real_lt_subst_left,finite_sum_reindex_bijection_50,real_lt_subst_right}.
        For all $i\in I$ we have
            $d(i)=a(i)\rmul\apply{\apply{g}{i}}{j}$ by
            \cref{bijections,surj,funs_type_apply,funs_is_function,pair_eq_iff,function_apply_intro,subseteq,real_lt_subst_left,real_lt_subst_right}.
    \end{subproof}
    $y$ is an affine combination of $g$ with coefficients $a$ in dimension $N$ by
        \cref{affine_combination_value_50}.
\end{proof}

% from §50 helper lemma 59: affine planes are closed under binary affine combinations
\begin{lemma}\label{affine_plane_binary_combination_50}
    Assume $N,n\in\naturalsPlus$.
    Let $I=\initialSegment{n}$.
    Let $J=\initialSegment{N}$.
    Let $E=\realsPower{J}$.
    Assume $A$ is a geometrically independent set in dimension $N$.
    Assume $A$ is inhabited.
    Assume $g$ is a bijection from $I$ to $A$.
    Assume $P$ is the plane determined by $A$ in dimension $N$.
    Assume $u,v\in P$.
    Assume $t\in\reals$.
    Let $\alpha=1+\neg{t}$.
    Let $w=\{(j,\alpha\rmul\apply{u}{j}+t\rmul\apply{v}{j})\mid j\in J\}$.
    Then $w\in P$.
\end{lemma}
\begin{proof}
    $\neg{t}\in\reals$ and $\alpha\in\reals$ by
        \cref{real_one_in_naturals,real_naturals_subset_reals,subseteq,real_add_inverse,real_add_closed}.
    Take $a$ such that
        $a\in\funs{I}{\reals}$ and
        $\sumFiniteSequence{a}{1}$ and
        $u$ is an affine combination of $g$ with coefficients $a$ in dimension $N$
        by \cref{plane_coordinates_fixed_enumeration_50}.
    Take $b$ such that
        $b\in\funs{I}{\reals}$ and
        $\sumFiniteSequence{b}{1}$ and
        $v$ is an affine combination of $g$ with coefficients $b$ in dimension $N$
        by \cref{plane_coordinates_fixed_enumeration_50}.
    Let $a_1=\coordScalarSequence{n}{\alpha}{a}$.
    Let $b_1=\coordScalarSequence{n}{t}{b}$.
    Let $c=\coordSumSequence{n}{a_1}{b_1}$.
    $a_1,b_1,c\in\funs{I}{\reals}$ by
        \cref{coord_scalar_sequence_function,coord_sum_sequence_function}.
    $\sumFiniteSequence{c}{1}$ by
        \cref{sum_finite_sequence_scalar,sum_finite_sequence_add,real_mul_ident,real_mul_comm,real_lt_subst_left,real_add_assoc,real_add_inverse,real_add_comm,real_add_eq_subst,real_add_ident,real_lt_subst_right}.
    $w\in E$ by
        \cref{relation,pair_eq_iff,rightunique,dom_intro,dom_iff,subseteq,subseteq_antisymmetric,plane_determined_by_points_50,reals_power,tuple_power,funs_type_apply,real_mul_closed,real_add_closed,function_apply_intro,real_lt_subst_left,funs_intro}.
    $A\subseteq E$ by \cref{geometrically_independent_set_50}.
    $g\in\funs{I}{E}$ by
        \cref{bijections_to_funs,funs_elim,funs_type_apply,subseteq,funs_intro}.
    Show that for all $j\in J$ there exists $r$ such that
            $r\in\funs{I}{\reals}$ and
            for all $i\in I$ we have
                $r(i)=c(i)\rmul\apply{\apply{g}{i}}{j}$ and
            $\sumFiniteSequence{r}{\apply{w}{j}}$.
    \begin{subproof}
            Fix $j\in J$.
            Take $k$ such that
                $k\in\naturalsPlus$ and
                $g\in\funs{\initialSegment{k}}{E}$ and
                $a\in\funs{\initialSegment{k}}{\reals}$ and
                $u\in E$ and
                for all $q\in J$ there exists $p$ such that
                    $p\in\funs{\initialSegment{k}}{\reals}$ and
                    for all $i\in\initialSegment{k}$ we have
                        $p(i)=a(i)\rmul\apply{\apply{g}{i}}{q}$ and
                    $\sumFiniteSequence{p}{\apply{u}{q}}$
                by \cref{affine_combination_value_50}.
            $\initialSegment{k}=I$ by \cref{funs_elim}.
            Take $p$ such that
                $p\in\funs{\initialSegment{k}}{\reals}$ and
                for all $i\in\initialSegment{k}$ we have
                    $p(i)=a(i)\rmul\apply{\apply{g}{i}}{j}$ and
                $\sumFiniteSequence{p}{\apply{u}{j}}$.
            $p\in\funs{I}{\reals}$ and
                $\sumFiniteSequence{p}{\apply{u}{j}}$ by
                \cref{real_one_in_naturals,real_one_leq_naturalsplus,initial_segment_naturalsplus,real_lt_subst_left,real_lt_subst_right,subseteq}.
            For all $i\in I$ we have
                $p(i)=a(i)\rmul\apply{\apply{g}{i}}{j}$ by
                \cref{real_lt_subst_left,real_lt_subst_right}.
            Take $l$ such that
                $l\in\naturalsPlus$ and
                $g\in\funs{\initialSegment{l}}{E}$ and
                $b\in\funs{\initialSegment{l}}{\reals}$ and
                $v\in E$ and
                for all $q\in J$ there exists $q_0$ such that
                    $q_0\in\funs{\initialSegment{l}}{\reals}$ and
                    for all $i\in\initialSegment{l}$ we have
                        $q_0(i)=b(i)\rmul\apply{\apply{g}{i}}{q}$ and
                    $\sumFiniteSequence{q_0}{\apply{v}{q}}$
                by \cref{affine_combination_value_50}.
            $\initialSegment{l}=I$ by \cref{funs_elim}.
            Take $q$ such that
                $q\in\funs{\initialSegment{l}}{\reals}$ and
                for all $i\in\initialSegment{l}$ we have
                    $q(i)=b(i)\rmul\apply{\apply{g}{i}}{j}$ and
                $\sumFiniteSequence{q}{\apply{v}{j}}$.
            $q\in\funs{I}{\reals}$ and
                $\sumFiniteSequence{q}{\apply{v}{j}}$ by
                \cref{real_one_in_naturals,real_one_leq_naturalsplus,initial_segment_naturalsplus,real_lt_subst_left,real_lt_subst_right,subseteq}.
            For all $i\in I$ we have
                $q(i)=b(i)\rmul\apply{\apply{g}{i}}{j}$ by
                \cref{real_lt_subst_left,real_lt_subst_right}.
            Let $p_1=\coordScalarSequence{n}{\alpha}{p}$.
            Let $q_1=\coordScalarSequence{n}{t}{q}$.
            Let $r=\coordSumSequence{n}{p_1}{q_1}$.
            $p_1,q_1,r\in\funs{I}{\reals}$ by
                \cref{coord_scalar_sequence_function,coord_sum_sequence_function}.
            $\sumFiniteSequence{r}{\apply{w}{j}}$ by
                \cref{sum_finite_sequence_scalar,sum_finite_sequence_add,relation,pair_eq_iff,rightunique,function_apply_intro,real_lt_subst_right}.
            Show that for all $i\in I$ we have
                $r(i)=c(i)\rmul\apply{\apply{g}{i}}{j}$.
            \begin{subproof}
                Fix $i\in I$.
                Let $z=\apply{\apply{g}{i}}{j}$.
                $a_1(i)=\alpha\rmul a(i)$ and
                    $b_1(i)=t\rmul b(i)$ by
                    \cref{coord_scalar_sequence,funs_is_function,function_apply_intro}.
                $c(i)=a_1(i)+b_1(i)$ by
                    \cref{coord_sum_sequence,funs_is_function,function_apply_intro}.
                $p_1(i)=\alpha\rmul p(i)$ and
                    $q_1(i)=t\rmul q(i)$ by
                    \cref{coord_scalar_sequence,funs_is_function,function_apply_intro}.
                $r(i)=p_1(i)+q_1(i)$ by
                    \cref{coord_sum_sequence,funs_is_function,function_apply_intro}.
                $p(i)=a(i)\rmul z$ and $q(i)=b(i)\rmul z$ by
                    \cref{subseteq}.
                Follows by
                    \cref{funs_type_apply,reals_power,tuple_power,real_mul_assoc,real_mul_comm,real_mul_closed,real_add_closed,real_distrib,real_add_eq_subst,real_lt_subst_left,real_lt_subst_right}.
            \end{subproof}
    \end{subproof}
    $w$ is an affine combination of $g$ with coefficients $c$ in dimension $N$
        by \cref{affine_combination_value_50}.
    Follows by \cref{plane_determined_by_points_50}.
\end{proof}

% from §50 helper definition 26: dot product of two finite real sequences
\begin{definition}\label{finite_real_dot_product_50}
    $r$ is the dot product of $a$ and $b$ of length $n$ iff
        $n\in\naturalsPlus$ and
        $a,b\in\funs{\initialSegment{n}}{\reals}$ and
        there exists $c$ such that
            $c=\coordProductSequence{n}{a}{b}$ and
            $\sumFiniteSequence{c}{r}$.
\end{definition}

% from §50 helper lemma 60: finite real dot products exist
\begin{lemma}\label{finite_real_dot_product_exists_50}
    Assume $n\in\naturalsPlus$.
    Assume $a,b\in\funs{\initialSegment{n}}{\reals}$.
    Then there exists $r\in\reals$ such that
        $r$ is the dot product of $a$ and $b$ of length $n$.
\end{lemma}
\begin{proof}
    Follows by
        \cref{coord_product_sequence_function,sum_finite_sequence_exists,finite_real_dot_product_50}.
\end{proof}

% from §50 helper lemma 61: finite real dot products are unique
\begin{lemma}\label{finite_real_dot_product_unique_50}
    Assume $r$ is the dot product of $a$ and $b$ of length $n$.
    Assume $s$ is the dot product of $a$ and $b$ of length $n$.
    Then $r=s$.
\end{lemma}
\begin{proof}
    Follows by
        \cref{finite_real_dot_product_50,sum_finite_sequence_unique}.
\end{proof}

% from §50 helper lemma 62: scaling one factor scales a finite dot product
\begin{lemma}\label{finite_real_dot_product_scalar_50}
    Assume $r$ is the dot product of $a$ and $b$ of length $n$.
    Assume $q\in\reals$.
    Let $d=\coordScalarSequence{n}{q}{a}$.
    Then $q\rmul r$ is the dot product of $d$ and $b$ of length $n$.
\end{lemma}
\begin{proof}
    $n\in\naturalsPlus$ and
        $a,b\in\funs{\initialSegment{n}}{\reals}$ by
        \cref{finite_real_dot_product_50}.
    Let $I=\initialSegment{n}$.
    Take $c$ such that
        $c=\coordProductSequence{n}{a}{b}$ and
        $\sumFiniteSequence{c}{r}$ by
        \cref{finite_real_dot_product_50}.
    Let $e=\coordProductSequence{n}{d}{b}$.
    Let $h=\coordScalarSequence{n}{q}{c}$.
    $e=h$ by
        \cref{coord_product_sequence_function,coord_scalar_sequence_function,coord_scalar_sequence,coord_product_sequence,funs_is_function,function_apply_intro,funs_type_apply,real_mul_assoc,real_lt_subst_left,real_lt_subst_right,functions_eq_iff}.
    Follows by
        \cref{coord_product_sequence_function,coord_scalar_sequence_function,sum_finite_sequence_scalar,real_lt_subst_left,finite_real_dot_product_50}.
\end{proof}

% from §50 helper lemma 63: adding first factors adds finite dot products
\begin{lemma}\label{finite_real_dot_product_add_50}
    Assume $r$ is the dot product of $a$ and $b$ of length $n$.
    Assume $s$ is the dot product of $d$ and $b$ of length $n$.
    Let $e=\coordSumSequence{n}{a}{d}$.
    Then $r+s$ is the dot product of $e$ and $b$ of length $n$.
\end{lemma}
\begin{proof}
    $n\in\naturalsPlus$ and
        $a,b,d\in\funs{\initialSegment{n}}{\reals}$ by
        \cref{finite_real_dot_product_50}.
    Let $I=\initialSegment{n}$.
    Take $c$ such that
        $c=\coordProductSequence{n}{a}{b}$ and
        $\sumFiniteSequence{c}{r}$ by
        \cref{finite_real_dot_product_50}.
    Take $h$ such that
        $h=\coordProductSequence{n}{d}{b}$ and
        $\sumFiniteSequence{h}{s}$ by
        \cref{finite_real_dot_product_50}.
    Let $p=\coordProductSequence{n}{e}{b}$.
    Let $q=\coordSumSequence{n}{c}{h}$.
    $p=q$ by
        \cref{coord_product_sequence_function,coord_sum_sequence_function,coord_sum_sequence,coord_product_sequence,funs_is_function,function_apply_intro,funs_type_apply,real_add_closed,real_mul_comm,real_distrib,real_add_eq_subst,real_lt_subst_left,real_lt_subst_right,functions_eq_iff}.
    Follows by
        \cref{coord_product_sequence_function,coord_sum_sequence_function,sum_finite_sequence_add,real_lt_subst_left,finite_real_dot_product_50}.
\end{proof}

% from §50 helper definition 27: tail coordinates of a successor-length sequence
\begin{definition}\label{successor_tail_sequence_50}
    $\successorTailSequence{n}{a}=
        \{(i,a(i+1))\mid i\in\initialSegment{n}\}$.
\end{definition}

% from §50 helper lemma 64: a successor-length real sequence has a real tail
\begin{lemma}\label{successor_tail_sequence_function_50}
    Assume $n\in\naturalsPlus$.
    Assume $a\in\funs{\initialSegment{n+1}}{\reals}$.
    Then $\successorTailSequence{n}{a}\in
        \funs{\initialSegment{n}}{\reals}$.
\end{lemma}
\begin{proof}
    Let $I=\initialSegment{n}$.
    Let $d=\successorTailSequence{n}{a}$.
    $d$ is a function by
        \cref{successor_tail_sequence_50,relation,rightunique,pair_eq_iff}.
    $\dom{d}=I$ by
        \cref{successor_tail_sequence_50,dom_intro,dom_iff,pair_eq_iff,subseteq,subseteq_antisymmetric}.
    Follows by
        \cref{initial_segment_naturalsplus,real_naturals_subset_reals,real_one_in_naturals,subseteq,real_leq_add,real_naturals_closed_succ,funs_type_apply,successor_tail_sequence_50,pair_eq_iff,function_apply_intro,real_lt_subst_left,funs_intro}.
\end{proof}

% from §50 helper lemma 65: split a successor-length dot product at its last coordinate
\begin{lemma}\label{finite_real_dot_product_split_last_50}
    Assume $n\in\naturalsPlus$.
    Let $I=\initialSegment{n}$.
    Let $I_1=\initialSegment{n+1}$.
    Assume $r$ is the dot product of $a$ and $b$ of length $n+1$.
    Let $a_0=\restrl{a}{I}$.
    Let $b_0=\restrl{b}{I}$.
    Assume $s$ is the dot product of $a_0$ and $b_0$ of length $n$.
    Then $r=s+(a(n+1)\rmul b(n+1))$.
\end{lemma}
\begin{proof}
    $n+1\in\naturalsPlus$ by
        \cref{real_naturals_closed_succ}.
    $a,b\in\funs{I_1}{\reals}$ by
        \cref{finite_real_dot_product_50}.
    Take $c$ such that
        $c=\coordProductSequence{n+1}{a}{b}$ and
        $\sumFiniteSequence{c}{r}$ by
        \cref{finite_real_dot_product_50}.
    $c\in\funs{I_1}{\reals}$ by
        \cref{coord_product_sequence_function}.
    Let $c_0=\restrl{c}{I}$.
    $c_0\in\funs{I}{\reals}$ by
        \cref{sum_finite_sequence_split_last}.
    Take $t\in\reals$ such that
        $\sumFiniteSequence{c_0}{t}$ and
        $r=t+c(n+1)$ by
        \cref{sum_finite_sequence_split_last}.
    Take $e$ such that
        $e=\coordProductSequence{n}{a_0}{b_0}$ and
        $\sumFiniteSequence{e}{s}$ by
        \cref{finite_real_dot_product_50}.
    $a_0,b_0\in\funs{I}{\reals}$ by
        \cref{finite_real_dot_product_50}.
    $e\in\funs{I}{\reals}$ by
        \cref{coord_product_sequence_function}.
    $t=s$ by
        \cref{initial_segment_subset_succ,subseteq,funs_elim,restrl_of_function_apply,funs_is_function,coord_product_sequence,function_apply_intro,functions_eq_iff,real_lt_subst_left,real_lt_subst_right,sum_finite_sequence_unique}.
    Follows by
        \cref{initial_segment_naturalsplus,real_naturals_closed_succ,coord_product_sequence,funs_is_function,function_apply_intro,real_lt_subst_left,real_lt_subst_right}.
\end{proof}

% from §50 helper definition 28: nontrivial kernel property for a finite number of equations
\begin{definition}\label{finite_homogeneous_kernel_property_50}
    $\finiteHomogeneousKernelProperty{m}$ iff
        $m\in\naturalsPlus$ and
        for all $N\in\naturalsPlus$ such that $m<N$ we have
            for all $M\in\funs{\initialSegment{m}}{
                \realsPower{\initialSegment{N}}}$
            there exists $z$ such that
                $z\in\realsPower{\initialSegment{N}}$ and
                $z\neq\zeroVector{N}$ and
                for all $i\in\initialSegment{m}$ we have
                    $0$ is the dot product of $M(i)$ and $z$ of length $N$.
\end{definition}

% from §50 helper lemma 66: a zero last coefficient is annihilated by the last coordinate unit
\begin{lemma}\label{last_coordinate_unit_annihilates_50}
    Assume $N\in\naturalsPlus$.
    Let $I=\initialSegment{N}$.
    Let $a\in\realsPower{I}$.
    Assume $a(N)=0$.
    Let $z=\standardAffineFramePoint{N}{N+1}$.
    Then $0$ is the dot product of $a$ and $z$ of length $N$.
\end{lemma}
\begin{proof}
    $N+1\in\initialSegment{N+1}$ by
        \cref{real_naturals_closed_succ,initial_segment_naturalsplus}.
    $z\in\realsPower{I}$ by
        \cref{standard_affine_frame_point_coordinates_50}.
    $a,z\in\funs{I}{\reals}$ by
        \cref{reals_power,tuple_power}.
    Let $c=\coordProductSequence{N}{a}{z}$.
    $c\in\funs{I}{\reals}$ by
        \cref{coord_product_sequence_function}.
    For all $i\in I$ if $i\neq N$ then $c(i)=0$ by
        \cref{initial_segment_naturalsplus,real_naturals_subset_reals,real_one_in_naturals,subseteq,real_add_cancel,real_one_leq_naturalsplus,real_naturals_lt_succ_local,real_add_closed,real_leq_lt_trans,real_lt_irreflexive,standard_affine_frame_point_off_coordinate_50,coord_product_sequence,funs_is_function,function_apply_intro,funs_type_apply,real_mul_zero,real_mul_comm,real_lt_subst_left,real_lt_subst_right}.
    $N\in I$ by \cref{initial_segment_naturalsplus}.
    $\sumFiniteSequence{c}{0}$ by
        \cref{finite_sum_single_support_50,standard_affine_frame_point_one_50,funs_type_apply,coord_product_sequence,funs_is_function,function_apply_intro,real_mul_zero,real_lt_subst_left,real_lt_subst_right}.
    Follows by \cref{finite_real_dot_product_50}.
\end{proof}

% from §50 helper lemma 67: restrict a successor-length sequence to its initial part
\begin{lemma}\label{successor_sequence_restriction_50}
    Assume $n\in\naturalsPlus$.
    Assume $a\in\funs{\initialSegment{n+1}}{X}$.
    Then $\restrl{a}{\initialSegment{n}}\in
        \funs{\initialSegment{n}}{X}$.
\end{lemma}
\begin{proof}
    Follows by
        \cref{funs_elim,restrl_of_function_is_function,dom_of_restrl_eq_inter,initial_segment_subset_succ,inter_absorb_supseteq_left,subseteq,funs_type_apply,restrl_of_function_apply,funs_intro}.
\end{proof}

% from §50 helper lemma 68: one nonzero equation coefficient admits a two-coordinate kernel vector
\begin{lemma}\label{one_equation_nonzero_last_kernel_50}
    Assume $N\in\naturalsPlus$.
    Assume $1<N$.
    Let $I=\initialSegment{N}$.
    Assume $a\in\realsPower{I}$.
    Assume $a(N)\neq 0$.
    Then there exists $z$ such that
        $z\in\realsPower{I}$ and
        $z\neq\zeroVector{N}$ and
        $0$ is the dot product of $a$ and $z$ of length $N$.
\end{lemma}
\begin{proof}
    $N\neq 1$ by
        \cref{real_one_in_naturals,real_naturals_subset_reals,subseteq,real_lt_irreflexive,real_lt_subst_right}.
    Take $k\in\naturalsPlus$ such that $N=k+1$ by
        \cref{real_naturals_predecessor}.
    Let $K=\initialSegment{k}$.
    $a\in\funs{I}{\reals}$ by \cref{reals_power,tuple_power}.
    $1,N\in I$ by
        \cref{real_one_in_naturals,real_one_leq_naturalsplus,initial_segment_naturalsplus}.
    $a(1),a(N)\in\reals$ by \cref{funs_type_apply}.
    Let $q=\neg{(\inv{a(N)}\rmul a(1))}$.
    $q\in\reals$ by
        \cref{real_mul_inverse,real_mul_closed,real_add_inverse}.
    Let $z=\{(i,r)\mid i\in I,r\in\reals\mid
        (i=1\land r=1)\lor
        (i=N\land r=q)\lor
        (i\neq 1\land i\neq N\land r=0)\}$.
    $z\in\funs{I}{\reals}$ by
        \cref{relation,rightunique,pair_eq_iff,dom_iff,subseteq,real_add_ident,real_one_in_naturals,real_naturals_subset_reals,dom_intro,subseteq_antisymmetric,function_apply_intro,real_lt_subst_left,funs_intro}.
    $z\neq\zeroVector{N}$ by
        \cref{funs_is_function,real_one_in_naturals,real_naturals_subset_reals,subseteq,pair_eq_iff,function_apply_intro,zero_vector_50,real_zero_lt_one,real_lt_irreflexive,real_lt_subst_left,real_lt_subst_right}.
    Let $c=\coordProductSequence{N}{a}{z}$.
    $c\in\funs{I}{\reals}$ by
        \cref{coord_product_sequence_function}.
    Let $c_0=\restrl{c}{K}$.
    $c_0\in\funs{K}{\reals}$ by
        \cref{successor_sequence_restriction_50,real_lt_subst_left,real_lt_subst_right}.
    For all $i\in K$ if $i\neq 1$ then $c_0(i)=0$ by
        \cref{initial_segment_subset_succ,subseteq,initial_segment_naturalsplus,real_naturals_lt_succ_local,real_naturals_subset_reals,real_leq_lt_trans,real_lt_irreflexive,real_add_ident,pair_eq_iff,function_apply_intro,funs_elim,restrl_of_function_apply,funs_is_function,coord_product_sequence,funs_type_apply,real_mul_zero,real_mul_comm,real_lt_subst_left,real_lt_subst_right}.
    $\sumFiniteSequence{c}{c_0(1)+c(k+1)}$ by
        \cref{real_one_in_naturals,real_one_leq_naturalsplus,initial_segment_naturalsplus,finite_sum_single_support_50,sum_finite_sequence_append_last}.
    $c_0(1)+c(k+1)=0$ by
        \cref{real_mul_inverse,real_mul_assoc,real_mul_neg_right,real_mul_closed,real_one_in_naturals,real_one_leq_naturalsplus,initial_segment_naturalsplus,funs_elim,initial_segment_subset_succ,restrl_of_function_apply,coord_product_sequence,funs_is_function,real_naturals_subset_reals,subseteq,pair_eq_iff,function_apply_intro,real_mul_ident,real_mul_comm,real_add_eq_subst,real_add_inverse,real_lt_subst_left,real_lt_subst_right}.
    Thus there exists $w$ such that
        $w\in\realsPower{I}$ and
        $w\neq\zeroVector{N}$ and
        $0$ is the dot product of $a$ and $w$ of length $N$ by
        \cref{real_lt_subst_right,reals_power,tuple_power,finite_real_dot_product_50}.
\end{proof}

% from §50 helper lemma 71: the last coordinate unit is nonzero
\begin{lemma}\label{last_coordinate_unit_nonzero_50}
    Assume $N\in\naturalsPlus$.
    Let $I=\initialSegment{N}$.
    Let $z=\standardAffineFramePoint{N}{N+1}$.
    Then $z\in\realsPower{I}\setminus\{\zeroVector{N}\}$.
\end{lemma}
\begin{proof}
    $N+1\in\initialSegment{N+1}$ by
        \cref{real_naturals_closed_succ,initial_segment_naturalsplus}.
    $z\in\realsPower{I}$ by
        \cref{standard_affine_frame_point_coordinates_50}.
    Follows by
        \cref{initial_segment_naturalsplus,standard_affine_frame_point_one_50,reals_power,tuple_power,funs_elim,function_apply_iff,zero_vector_50,pair_eq_iff,real_zero_lt_one,real_add_ident,real_one_in_naturals,real_naturals_subset_reals,subseteq,real_lt_irreflexive,real_lt_subst_left,real_lt_subst_right,singleton_iff,setminus_intro}.
\end{proof}

% from §50 helper lemma 69: one homogeneous equation has a nontrivial kernel in at least two variables
\begin{lemma}\label{finite_homogeneous_kernel_base_50}
    $\finiteHomogeneousKernelProperty{1}$.
\end{lemma}
\begin{proof}
    $1\in\naturalsPlus$ by \cref{real_one_in_naturals}.
    Show that for all $N\in\naturalsPlus$ such that $1<N$ we have
        for all $M\in\funs{\initialSegment{1}}{
            \realsPower{\initialSegment{N}}}$
        there exists $z$ such that
            $z\in\realsPower{\initialSegment{N}}$ and
            $z\neq\zeroVector{N}$ and
            for all $i\in\initialSegment{1}$ we have
                $0$ is the dot product of $M(i)$ and $z$ of length $N$.
    \begin{subproof}
        Fix $N\in\naturalsPlus$.
        Assume $1<N$.
        Fix $M\in\funs{\initialSegment{1}}{
            \realsPower{\initialSegment{N}}}$.
        Let $I=\initialSegment{N}$.
        Let $a=M(1)$.
        $1\in\initialSegment{1}$ by
            \cref{initial_segment_one,singleton_iff}.
        $a\in\realsPower{I}$ by \cref{funs_type_apply}.
        \begin{byCase}
        \caseOf{$a(N)=0$.}
            Thus there exists $w$ such that
                $w\in\realsPower{I}$ and
                $w\neq\zeroVector{N}$ and
                for all $i\in\initialSegment{1}$ we have
                    $0$ is the dot product of $M(i)$ and $w$ of length $N$ by
                    \cref{last_coordinate_unit_nonzero_50,setminus_elim_left,setminus_elim_right,last_coordinate_unit_annihilates_50,initial_segment_one,singleton_iff,real_lt_subst_left}.
        \caseOf{$a(N)\neq 0$.}
            Thus there exists $w$ such that
                $w\in\realsPower{I}$ and
                $w\neq\zeroVector{N}$ and
                for all $i\in\initialSegment{1}$ we have
                    $0$ is the dot product of $M(i)$ and $w$ of length $N$ by
                    \cref{one_equation_nonzero_last_kernel_50,initial_segment_one,singleton_iff,real_lt_subst_left}.
        \end{byCase}
    \end{subproof}
    Follows by \cref{finite_homogeneous_kernel_property_50}.
\end{proof}

% from §50 helper lemma 70: enumerate a successor initial segment with one prescribed point removed
\begin{lemma}\label{initial_segment_remove_point_enumeration_50}
    Assume $m\in\naturalsPlus$.
    Let $J=\initialSegment{m}$.
    Let $J_1=\initialSegment{m+1}$.
    Assume $p\in J_1$.
    Then there exists $h$ such that
        $h$ is a bijection from $J$ to $J_1\setminus\{p\}$.
\end{lemma}
\begin{proof}
    $m+1\in\naturalsPlus$ by
        \cref{real_naturals_closed_succ}.
    $m+1\in J_1$ by \cref{initial_segment_naturalsplus}.
    Let $g=\{(i,j)\mid i\in J_1,j\in J_1\mid
        (i=p\land j=m+1)\lor
        (i=m+1\land j=p)\lor
        (i\neq p\land i\neq m+1\land j=i)\}$.
    $g\in\funs{J_1}{J_1}$ by
        \cref{relation,rightunique,pair_eq_iff,dom_iff,subseteq,dom_intro,subseteq_antisymmetric,function_apply_iff,funs_intro}.
    For all $i\in J_1$ we have $g(g(i))=i$ by
        \cref{funs_is_function,pair_eq_iff,function_apply_intro,real_lt_subst_left,real_lt_subst_right}.
    $g$ is a bijection from $J_1$ to $J_1$ by
        \cref{funs_elim,funs_is_function,funs_type_apply,surj,injective_function,subseteq,bijections}.
    $g(m+1)=p$ by
        \cref{funs_is_function,pair_eq_iff,function_apply_intro}.
    Let $h=\restrl{g}{J}$.
    Thus there exists $q$ such that
        $q$ is a bijection from $J$ to $J_1\setminus\{p\}$ by
        \cref{finite_enumeration_remove_last}.
\end{proof}

% from §50 helper definition 29: eliminate a final coordinate from one row using a pivot row
\begin{definition}\label{reduced_homogeneous_row_50}
    $\reducedHomogeneousRow{k}{M}{r}{p}=
        \{(j,\apply{\apply{M}{r}}{j}+
            \neg{(\apply{\apply{M}{r}}{k+1}
                \rmul\inv{\apply{\apply{M}{p}}{k+1}})}
                \rmul\apply{\apply{M}{p}}{j})\mid
            j\in\initialSegment{k}\}$.
\end{definition}

% from §50 helper definition 30: matrix of rows after final-coordinate elimination
\begin{definition}\label{reduced_homogeneous_matrix_50}
    $\reducedHomogeneousMatrix{m}{k}{M}{h}{p}=
        \{(u,\reducedHomogeneousRow{k}{M}{\apply{h}{u}}{p})\mid
            u\in\initialSegment{m}\}$.
\end{definition}

% from §50 helper lemma 72: pivot elimination produces a smaller real matrix
\begin{lemma}\label{reduced_homogeneous_matrix_function_50}
    Assume $m,k\in\naturalsPlus$.
    Let $L=\initialSegment{m}$.
    Let $R=\initialSegment{m+1}$.
    Let $K=\initialSegment{k}$.
    Let $C=\initialSegment{k+1}$.
    Assume $M\in\funs{R}{\realsPower{C}}$.
    Assume $p\in R$ and $\apply{\apply{M}{p}}{k+1}\neq 0$.
    Assume $h\in\funs{L}{R}$.
    Then $\reducedHomogeneousMatrix{m}{k}{M}{h}{p}\in
        \funs{L}{\realsPower{K}}$.
\end{lemma}
\begin{proof}
    $m+1,k+1\in\naturalsPlus$ by
        \cref{real_naturals_closed_succ}.
    $k+1\in C$ by \cref{initial_segment_naturalsplus}.
    $\apply{M}{p}\in\funs{C}{\reals}$ by
        \cref{funs_type_apply,reals_power,tuple_power}.
    $\inv{\apply{\apply{M}{p}}{k+1}}\in\reals$ by
        \cref{funs_type_apply,real_mul_inverse}.
    Let $B=\reducedHomogeneousMatrix{m}{k}{M}{h}{p}$.
    $B$ is a function by
        \cref{reduced_homogeneous_matrix_50,relation,rightunique,pair_eq_iff}.
    $\dom{B}=L$ by
        \cref{reduced_homogeneous_matrix_50,dom_intro,dom_iff,pair_eq_iff,subseteq,subseteq_antisymmetric}.
    Show that for all $u\in\dom{B}$ we have $B(u)\in\realsPower{K}$.
    \begin{subproof}
        Fix $u\in\dom{B}$.
        $\apply{M}{h(u)}\in\funs{C}{\reals}$ by
            \cref{funs_type_apply,reals_power,tuple_power}.
        $\neg{(\apply{\apply{M}{h(u)}}{k+1}
            \rmul\inv{\apply{\apply{M}{p}}{k+1}})}\in\reals$ by
            \cref{funs_type_apply,real_mul_closed,real_add_inverse}.
        Let $d=\reducedHomogeneousRow{k}{M}{h(u)}{p}$.
        $d$ is a function by
            \cref{reduced_homogeneous_row_50,relation,rightunique,pair_eq_iff}.
        $\dom{d}=K$ by
            \cref{reduced_homogeneous_row_50,dom_intro,dom_iff,pair_eq_iff,subseteq,subseteq_antisymmetric}.
        $d\in\funs{K}{\reals}$ by
            \cref{initial_segment_subset_succ,subseteq,funs_type_apply,real_mul_closed,real_add_closed,reduced_homogeneous_row_50,pair_eq_iff,function_apply_intro,real_lt_subst_left,funs_intro}.
        Follows by
            \cref{reals_power,tuple_power,reduced_homogeneous_matrix_50,pair_eq_iff,function_apply_intro,real_lt_subst_left}.
    \end{subproof}
    Follows by \cref{funs_intro}.
\end{proof}

% from §50 helper lemma 73: the finite zero vector is a euclidean point
\begin{lemma}\label{zero_vector_euclidean_50}
    Assume $N\in\naturalsPlus$.
    Then $\zeroVector{N}\in
        \realsPower{\initialSegment{N}}$.
\end{lemma}
\begin{proof}
    Follows by
        \cref{zero_vector_50,relation,pair_eq_iff,rightunique,dom_intro,dom_iff,subseteq,subseteq_antisymmetric,real_add_ident,function_apply_intro,funs_intro,reals_power,tuple_power}.
\end{proof}

% from §50 helper lemma 74: pivot-elimination scalars give the same correction term
\begin{lemma}\label{pivot_elimination_scalar_identity_50}
    Assume $a,b,s\in\reals$.
    Assume $a\neq 0$.
    Let $q=\neg{(b\rmul\inv{a})}$.
    Let $t=\neg{(\inv{a}\rmul s)}$.
    Then $b\rmul t=q\rmul s$.
\end{lemma}
\begin{proof}
    Follows by
        \cref{real_mul_inverse,real_mul_closed,real_mul_neg_right,real_mul_assoc,real_add_inverse,real_mul_comm,real_lt_subst_left,real_lt_subst_right}.
\end{proof}

% from §50 helper lemma 75: the recovered pivot coordinate cancels the pivot tail
\begin{lemma}\label{pivot_recovery_cancels_50}
    Assume $a,s\in\reals$.
    Assume $a\neq 0$.
    Let $t=\neg{(\inv{a}\rmul s)}$.
    Then $s+(a\rmul t)=0$.
\end{lemma}
\begin{proof}
    Follows by
        \cref{real_mul_inverse,real_mul_closed,real_mul_assoc,real_mul_ident,real_mul_neg_right,real_add_inverse,real_lt_subst_left,real_lt_subst_right}.
\end{proof}

% from §50 helper lemma 76: lift a kernel vector after pivot elimination
\begin{lemma}\label{pivot_elimination_kernel_lift_50}
    Assume $m,k\in\naturalsPlus$.
    Let $L=\initialSegment{m}$.
    Let $R=\initialSegment{m+1}$.
    Let $K=\initialSegment{k}$.
    Let $C=\initialSegment{k+1}$.
    Assume $M\in\funs{R}{\realsPower{C}}$.
    Assume $p\in R$ and $\apply{\apply{M}{p}}{k+1}\neq 0$.
    Assume $h$ is a bijection from $L$ to $R\setminus\{p\}$.
    Let $B=\reducedHomogeneousMatrix{m}{k}{M}{h}{p}$.
    Assume $z_0\in\realsPower{K}$ and
        $z_0\neq\zeroVector{k}$.
    Assume for all $u\in L$ we have
        $0$ is the dot product of $B(u)$ and $z_0$ of length $k$.
    Then there exists $z$ such that
        $z\in\realsPower{C}$ and
        $z\neq\zeroVector{k+1}$ and
        for all $i\in R$ we have
            $0$ is the dot product of $M(i)$ and $z$ of length $k+1$.
\end{lemma}
\begin{proof}
    $m+1,k+1\in\naturalsPlus$ by
        \cref{real_naturals_closed_succ}.
    $k+1\in C$ by \cref{initial_segment_naturalsplus}.
    $h\in\funs{L}{R}$ by
        \cref{bijections_to_funs,funs_type_apply,setminus_elim_left,funs_elim,subseteq,funs_intro}.
    $B\in\funs{L}{\realsPower{K}}$ by
        \cref{reduced_homogeneous_matrix_function_50}.
    Let $P=M(p)$.
    $P\in\realsPower{C}$ by \cref{funs_type_apply}.
    $P\in\funs{C}{\reals}$ by \cref{reals_power,tuple_power}.
    Let $P_0=\restrl{P}{K}$.
    $P_0\in\funs{K}{\reals}$ by
        \cref{successor_sequence_restriction_50}.
    $z_0\in\funs{K}{\reals}$ by \cref{reals_power,tuple_power}.
    Take $s\in\reals$ such that
        $s$ is the dot product of $P_0$ and $z_0$ of length $k$ by
        \cref{finite_real_dot_product_exists_50}.
    Let $a=P(k+1)$.
    $a\in\reals$ by \cref{funs_type_apply}.
    $a\neq 0$ by \cref{real_lt_subst_left}.
    $\inv{a}\in\reals$ by \cref{real_mul_inverse}.
    Let $t=\neg{(\inv{a}\rmul s)}$.
    $t\in\reals$ by
        \cref{real_mul_closed,real_add_inverse}.
    Let $r=\{(k+1,t)\}$.
    $r$ is a function by
        \cref{relation,rightunique,singleton_elim,pair_eq_iff}.
    $\dom{r}=\{k+1\}$ by
        \cref{dom_intro,dom_iff,singleton_elim,pair_eq_iff,singleton_intro,subseteq,subseteq_antisymmetric}.
    $k,k+1\in\reals$ by
        \cref{real_naturals_closed_succ,real_naturals_subset_reals,subseteq}.
    $k+1\notin K$ by
        \cref{initial_segment_naturalsplus,real_naturals_lt_succ_local,real_lt_trans,real_lt_irreflexive}.
    $\dom{z_0}=K$ by \cref{funs_elim}.
    $\dom{z_0}\inter\dom{r}=\emptyset$ by
        \cref{inter,emptyset,singleton_iff,subseteq,emptyset_subseteq,subseteq_antisymmetric}.
    Let $z=z_0\union r$.
    $z$ is a function by
        \cref{funs_is_function,union_compatible_functions,emptyset,inter}.
    $\dom{z}=K\union\{k+1\}$ by
        \cref{dom_union,real_lt_subst_left,real_lt_subst_right}.
    $K\union\{k+1\}=C$ by
        \cref{union_iff,singleton_iff,initial_segment_subset_succ,subseteq,real_naturals_closed_succ,initial_segment_naturalsplus,initial_segment_succ_split,union_intro_left,union_intro_right,singleton_intro,subseteq_antisymmetric}.
    $\dom{z}=C$ by
        \cref{real_lt_subst_left,real_lt_subst_right}.
    For all $j\in\dom{z}$ we have $z(j)\in\reals$ by
        \cref{dom_iff,dom_intro,funs_is_function,function_apply_intro,funs_type_apply,real_lt_subst_left,real_lt_subst_right,union_iff,singleton_iff,pair_eq_iff}.
    $z\in\funs{C}{\reals}$ by \cref{funs_intro}.
    $z\in\realsPower{C}$ by \cref{reals_power,tuple_power}.
    $z(k+1)=t$ by
        \cref{singleton_intro,union_intro_right,function_apply_intro}.
    $\restrl{z}{K}=z_0$ by
        \cref{successor_sequence_restriction_50,funs_elim,funs_is_function,function_apply_iff,union_intro_left,function_apply_intro,initial_segment_subset_succ,restrl_of_function_apply,functions_eq_iff}.
    $z\neq\zeroVector{k+1}$ by
        \cref{zero_vector_euclidean_50,real_naturals_closed_succ,reals_power,tuple_power,initial_segment_subset_succ,funs_elim,restrl_of_function_apply,funs_is_function,real_lt_subst_left,real_lt_subst_right,zero_vector_50,pair_eq_iff,subseteq,function_apply_intro,functions_eq_iff}.
    Show that for all $i\in R$ we have
        $0$ is the dot product of $M(i)$ and $z$ of length $k+1$.
    \begin{subproof}
        Fix $i\in R$.
        Let $A=M(i)$.
        $A\in\realsPower{C}$ by \cref{funs_type_apply}.
        $A\in\funs{C}{\reals}$ by \cref{reals_power,tuple_power}.
        Let $A_0=\restrl{A}{K}$.
        $A_0\in\funs{K}{\reals}$ by
            \cref{successor_sequence_restriction_50}.
        \begin{byCase}
        \caseOf{$i=p$.}
            $s$ is the dot product of $A_0$ and
                $\restrl{z}{K}$ of length $k$ by
                \cref{real_lt_subst_left,real_lt_subst_right}.
            Take $v\in\reals$ such that
                $v$ is the dot product of $A$ and $z$ of length $k+1$ by
                \cref{finite_real_dot_product_exists_50}.
            Follows by
                \cref{finite_real_dot_product_split_last_50,pivot_recovery_cancels_50,real_lt_subst_left,real_lt_subst_right}.
        \caseOf{$i\neq p$.}
            $i\in R\setminus\{p\}$ by
                \cref{setminus_intro,singleton_iff}.
            Take $u\in L$ such that $h(u)=i$ by
                \cref{bijections,surj}.
            Let $q=\neg{(A(k+1)\rmul\inv{a})}$.
            $A(k+1)\in\reals$ by \cref{funs_type_apply}.
            Let $D=\coordScalarSequence{k}{q}{P_0}$.
            Let $E=\coordSumSequence{k}{A_0}{D}$.
            $D,E\in\funs{K}{\reals}$ by
                \cref{real_mul_closed,real_add_inverse,coord_scalar_sequence_function,coord_sum_sequence_function}.
            Let $d=\reducedHomogeneousRow{k}{M}{h(u)}{p}$.
            $d=B(u)$ by
                \cref{reduced_homogeneous_matrix_50,pair_eq_iff,funs_is_function,function_apply_intro}.
            $d=E$ by
                \cref{funs_type_apply,reals_power,tuple_power,initial_segment_subset_succ,subseteq,funs_elim,restrl_of_function_apply,funs_is_function,coord_scalar_sequence,function_apply_intro,coord_sum_sequence,real_lt_subst_left,real_lt_subst_right,reduced_homogeneous_row_50,pair_eq_iff,functions_eq_iff}.
            $0$ is the dot product of $E$ and $z_0$ of length $k$ by
                \cref{subseteq,real_lt_subst_left,real_lt_subst_right}.
            Take $r_0\in\reals$ such that
                $r_0$ is the dot product of $A_0$ and $z_0$ of length $k$ by
                \cref{finite_real_dot_product_exists_50}.
            $r_0+q\rmul s=0$ by
                \cref{real_mul_closed,real_add_inverse,finite_real_dot_product_scalar_50,finite_real_dot_product_add_50,finite_real_dot_product_unique_50}.
            $r_0$ is the dot product of $A_0$ and
                $\restrl{z}{K}$ of length $k$ by
                \cref{real_lt_subst_right}.
            Take $v\in\reals$ such that
                $v$ is the dot product of $A$ and $z$ of length $k+1$ by
                \cref{finite_real_dot_product_exists_50}.
            Follows by
                \cref{finite_real_dot_product_split_last_50,pivot_elimination_scalar_identity_50,real_add_eq_subst,real_lt_subst_left,real_lt_subst_right}.
        \end{byCase}
    \end{subproof}
    Follows.
\end{proof}

% from §50 helper lemma 77: the finite homogeneous-kernel property is closed under successors
\begin{lemma}\label{finite_homogeneous_kernel_successor_50}
    Assume $\finiteHomogeneousKernelProperty{m}$.
    Then $\finiteHomogeneousKernelProperty{m+1}$.
\end{lemma}
\begin{proof}
    $m\in\naturalsPlus$ by
        \cref{finite_homogeneous_kernel_property_50}.
    It suffices to show that for all $N\in\naturalsPlus$ such that $m+1<N$ we have
        for all $M\in\funs{\initialSegment{m+1}}{
            \realsPower{\initialSegment{N}}}$
        there exists $z$ such that
            $z\in\realsPower{\initialSegment{N}}$ and
            $z\neq\zeroVector{N}$ and
            for all $i\in\initialSegment{m+1}$ we have
                $0$ is the dot product of $M(i)$ and $z$ of length $N$
        by \cref{real_naturals_closed_succ,finite_homogeneous_kernel_property_50}.
        Fix $N\in\naturalsPlus$.
        Assume $m+1<N$.
        Fix $M\in\funs{\initialSegment{m+1}}{
            \realsPower{\initialSegment{N}}}$.
        Let $R=\initialSegment{m+1}$.
        Take $k\in\naturalsPlus$ such that $N=k+1$ by
            \cref{real_naturals_closed_succ,real_one_leq_naturalsplus,real_one_in_naturals,real_naturals_subset_reals,subseteq,real_leq_lt_trans,real_lt_irreflexive,real_lt_subst_right,real_naturals_predecessor}.
        Let $F=\{i\in R\mid \apply{\apply{M}{i}}{N}\neq 0\}$.
        \begin{byCase}
        \caseOf{$F=\emptyset$.}
            Follows by
                \cref{pair_eq_iff,real_lt_subst_left,emptyset,last_coordinate_unit_nonzero_50,setminus_elim_left,setminus_elim_right,singleton_iff,funs_type_apply,subseteq,last_coordinate_unit_annihilates_50}.
        \caseOf{$F\neq\emptyset$.}
            Take $p$ such that $p\in R$ and
                $\apply{\apply{M}{p}}{N}\neq 0$ by
                \cref{nonempty_inhabited,pair_eq_iff}.
            Take $h$ such that
                $h$ is a bijection from $\initialSegment{m}$ to
                    $R\setminus\{p\}$ by
                \cref{initial_segment_remove_point_enumeration_50}.
            Let $B=\reducedHomogeneousMatrix{m}{k}{M}{h}{p}$.
            Follows by
                \cref{real_naturals_subset_reals,subseteq,real_lt_trichotomy,real_one_in_naturals,real_leq_add,real_add_closed,real_lt_trans,real_lt_irreflexive,bijections_to_funs,funs_elim,funs_type_apply,setminus_elim_left,funs_intro,reduced_homogeneous_matrix_function_50,finite_homogeneous_kernel_property_50,pivot_elimination_kernel_lift_50,real_lt_subst_left,real_lt_subst_right}.
        \end{byCase}
\end{proof}

% from §50 helper lemma 78: the finite homogeneous-kernel property holds in every positive dimension
\begin{lemma}\label{finite_homogeneous_kernel_all_dimensions_50}
    For all $m\in\naturalsPlus$ we have
        $\finiteHomogeneousKernelProperty{m}$.
\end{lemma}
\begin{proof}
    Let $A=\{m\in\naturalsPlus\mid
        \finiteHomogeneousKernelProperty{m}\}$.
    Follows by
        \cref{subseteq,real_one_in_naturals,finite_homogeneous_kernel_base_50,pair_eq_iff,finite_homogeneous_kernel_successor_50,real_naturals_closed_succ,real_naturals_induction}.
\end{proof}

% from §50 helper lemma 79: fewer homogeneous equations than variables have a nonzero solution
\begin{lemma}\label{finite_homogeneous_kernel_50}
    Assume $m,N\in\naturalsPlus$.
    Assume $m<N$.
    Let $M\in\funs{\initialSegment{m}}{
        \realsPower{\initialSegment{N}}}$.
    Then there exists $z$ such that
        $z\in\realsPower{\initialSegment{N}}$ and
        $z\neq\zeroVector{N}$ and
        for all $i\in\initialSegment{m}$ we have
            $0$ is the dot product of $M(i)$ and $z$ of length $N$.
\end{lemma}
\begin{proof}
    Follows by
        \cref{finite_homogeneous_kernel_all_dimensions_50,finite_homogeneous_kernel_property_50}.
\end{proof}

% from §50 helper definition 32: dot products commute with finite coordinatewise linear combinations
\begin{definition}\label{finite_affine_dot_property_50}
    $\finiteAffineDotProperty{n}$ iff
        $n\in\naturalsPlus$ and
        for all $N\in\naturalsPlus$ we have
        for all $s\in\funs{\initialSegment{n}}{
            \realsPower{\initialSegment{N}}}$ we have
        for all $a\in\funs{\initialSegment{n}}{\reals}$ we have
        for all $y\in\realsPower{\initialSegment{N}}$ we have
        for all $c\in\realsPower{\initialSegment{N}}$ if
            $y$ is an affine combination of $s$ with coefficients $a$
                in dimension $N$
        then for all $q\in\funs{\initialSegment{n}}{\reals}$ we have
        for all $b\in\funs{\initialSegment{n}}{\reals}$ we have
        for all $r\in\reals$ we have
        for all $t\in\reals$ if
            for all $i\in\initialSegment{n}$ we have
                $q(i)$ is the dot product of $s(i)$ and $c$ of length $N$ and
            for all $i\in\initialSegment{n}$ we have
                $b(i)=a(i)\rmul q(i)$ and
            $r$ is the dot product of $y$ and $c$ of length $N$ and
            $\sumFiniteSequence{b}{t}$
        then $r=t$.
\end{definition}

% from §50 helper lemma 80: dot products commute with one-term linear combinations
\begin{lemma}\label{finite_affine_dot_base_50}
    $\finiteAffineDotProperty{1}$.
\end{lemma}
\begin{proof}
    Let $I=\initialSegment{1}$.
    It suffices to show that for all $N\in\naturalsPlus$ we have
        for all $s\in\funs{I}{\realsPower{\initialSegment{N}}}$ we have
        for all $a\in\funs{I}{\reals}$ we have
        for all $y\in\realsPower{\initialSegment{N}}$ we have
        for all $c\in\realsPower{\initialSegment{N}}$ if
            $y$ is an affine combination of $s$ with coefficients $a$
                in dimension $N$
        then for all $q\in\funs{I}{\reals}$ we have
        for all $b\in\funs{I}{\reals}$ we have
        for all $r\in\reals$ we have
        for all $t\in\reals$ if
            for all $i\in I$ we have
                $q(i)$ is the dot product of $s(i)$ and $c$ of length $N$ and
            for all $i\in I$ we have $b(i)=a(i)\rmul q(i)$ and
            $r$ is the dot product of $y$ and $c$ of length $N$ and
            $\sumFiniteSequence{b}{t}$
        then $r=t$ by
        \cref{finite_affine_dot_property_50,real_one_in_naturals}.
        Fix $N\in\naturalsPlus$.
        Fix $s\in\funs{I}{\realsPower{\initialSegment{N}}}$.
        Fix $a\in\funs{I}{\reals}$.
        Fix $y\in\realsPower{\initialSegment{N}}$.
        Fix $c\in\realsPower{\initialSegment{N}}$.
        Assume $y$ is an affine combination of $s$ with coefficients $a$
            in dimension $N$.
        Fix $q\in\funs{I}{\reals}$.
        Fix $b\in\funs{I}{\reals}$.
        Fix $r\in\reals$.
        Fix $t\in\reals$.
        Assume for all $i\in I$ we have
                $q(i)$ is the dot product of $s(i)$ and $c$ of length $N$ and
            for all $i\in I$ we have $b(i)=a(i)\rmul q(i)$ and
            $r$ is the dot product of $y$ and $c$ of length $N$ and
            $\sumFiniteSequence{b}{t}$.
        Let $v=\coordScalarSequence{N}{a(1)}{s(1)}$.
        $y=v$ by
            \cref{affine_combination_value_50,funs_elim,real_one_in_naturals,initial_segment_one,singleton_iff,funs_type_apply,reals_power,tuple_power,coord_scalar_sequence_function,real_lt_subst_right,subseteq,real_lt_subst_left,sum_finite_sequence_one_term,sum_finite_sequence_unique,coord_scalar_sequence,funs_is_function,function_apply_intro,functions_eq_iff}.
        Follows by
            \cref{initial_segment_one,singleton_iff,subseteq,funs_type_apply,finite_real_dot_product_scalar_50,real_lt_subst_left,real_lt_subst_right,finite_real_dot_product_unique_50,sum_finite_sequence_one_term,sum_finite_sequence_unique}.
\end{proof}

% from §50 helper lemma 81: the finite affine-dot property is closed under successors
\begin{lemma}\label{finite_affine_dot_successor_50}
    Assume $\finiteAffineDotProperty{n}$.
    Then $\finiteAffineDotProperty{n+1}$.
\end{lemma}
\begin{proof}
    $n\in\naturalsPlus$ by \cref{finite_affine_dot_property_50}.
    Let $I=\initialSegment{n}$.
    Let $I_1=\initialSegment{n+1}$.
    It suffices to show that for all $N\in\naturalsPlus$ we have
        for all $s\in\funs{I_1}{\realsPower{\initialSegment{N}}}$ we have
        for all $a\in\funs{I_1}{\reals}$ we have
        for all $y\in\realsPower{\initialSegment{N}}$ we have
        for all $c\in\realsPower{\initialSegment{N}}$ if
            $y$ is an affine combination of $s$ with coefficients $a$
                in dimension $N$
        then for all $q\in\funs{I_1}{\reals}$ we have
        for all $b\in\funs{I_1}{\reals}$ we have
        for all $r\in\reals$ we have
        for all $t\in\reals$ if
            for all $i\in I_1$ we have
                $q(i)$ is the dot product of $s(i)$ and $c$ of length $N$ and
            for all $i\in I_1$ we have $b(i)=a(i)\rmul q(i)$ and
            $r$ is the dot product of $y$ and $c$ of length $N$ and
            $\sumFiniteSequence{b}{t}$
        then $r=t$ by
        \cref{real_naturals_closed_succ,finite_affine_dot_property_50}.
        Fix $N\in\naturalsPlus$.
        Fix $s\in\funs{I_1}{\realsPower{\initialSegment{N}}}$.
        Fix $a\in\funs{I_1}{\reals}$.
        Fix $y\in\realsPower{\initialSegment{N}}$.
        Fix $c\in\realsPower{\initialSegment{N}}$.
        Assume $y$ is an affine combination of $s$ with coefficients $a$
            in dimension $N$.
        Fix $q\in\funs{I_1}{\reals}$.
        Fix $b\in\funs{I_1}{\reals}$.
        Fix $r\in\reals$.
        Fix $t\in\reals$.
        Assume for all $i\in I_1$ we have
                $q(i)$ is the dot product of $s(i)$ and $c$ of length $N$ and
            for all $i\in I_1$ we have $b(i)=a(i)\rmul q(i)$ and
            $r$ is the dot product of $y$ and $c$ of length $N$ and
            $\sumFiniteSequence{b}{t}$.
        Let $J=\initialSegment{N}$.
        $n+1\in I_1$ by
            \cref{real_naturals_closed_succ,initial_segment_naturalsplus}.
        $q(n+1)$ is the dot product of $s(n+1)$ and $c$ of length $N$ and
            for all $i\in I_1$ we have $b(i)=a(i)\rmul q(i)$ and
            $r$ is the dot product of $y$ and $c$ of length $N$ and
            $\sumFiniteSequence{b}{t}$ by \cref{subseteq}.
        Let $s_0=\restrl{s}{I}$.
        Let $a_0=\restrl{a}{I}$.
        Let $q_0=\restrl{q}{I}$.
        Let $b_0=\restrl{b}{I}$.
        $s_0\in\funs{I}{\realsPower{J}}$ and
            $a_0,q_0,b_0\in\funs{I}{\reals}$ by
            \cref{successor_sequence_restriction_50}.
        Let $y_0=\{(j,u)\mid j\in J,u\in\reals\mid
            \text{there exists $e\in\funs{I}{\reals}$ such that
            for all $i\in I$ we have
                $e(i)=a_0(i)\rmul\apply{\apply{s_0}{i}}{j}$ and
            $\sumFiniteSequence{e}{u}$}\}$.
        $y_0$ is a function by
            \cref{relation,rightunique,pair_eq_iff,real_one_in_naturals,real_one_leq_naturalsplus,initial_segment_naturalsplus,subseteq,real_lt_subst_left,real_lt_subst_right,funs_is_function,functions_eq_iff,sum_finite_sequence_unique}.
        $\dom{y_0}\subseteq J$ by
                \cref{subseteq,dom_iff,pair_eq_iff}.
        $\dom{y_0}=J$ by
            \cref{affine_combination_value_50,funs_elim,real_lt_subst_left,subseteq,successor_sequence_restriction_50,sum_finite_sequence_split_last,initial_segment_subset_succ,restrl_of_function_apply,funs_is_function,real_lt_subst_right,pair_eq_iff,dom_intro,subseteq_antisymmetric}.
        $y_0\in\funs{J}{\reals}$ by
            \cref{dom_iff,pair_eq_iff,function_apply_intro,real_lt_subst_left,funs_intro}.
        $y_0\in\realsPower{J}$ by
            \cref{reals_power,tuple_power}.
        $y_0$ is an affine combination of $s_0$ with
            coefficients $a_0$ in dimension $N$ by
            \cref{funs_is_function,function_apply_elim,subseteq,pair_eq_iff,affine_combination_value_50}.
        For all $i\in I$ we have
            $q_0(i)$ is the dot product of $s_0(i)$ and $c$ of length $N$ by
            \cref{initial_segment_subset_succ,subseteq,funs_elim,restrl_of_function_apply,funs_is_function,real_lt_subst_left,real_lt_subst_right}.
        For all $i\in I$ we have
            $b_0(i)=a_0(i)\rmul q_0(i)$ by
            \cref{initial_segment_subset_succ,subseteq,funs_elim,restrl_of_function_apply,funs_is_function,real_lt_subst_left,real_lt_subst_right}.
        Take $r_0\in\reals$ such that
            $r_0$ is the dot product of $y_0$ and $c$ of length $N$ by
            \cref{reals_power,tuple_power,real_lt_subst_left,real_lt_subst_right,finite_real_dot_product_exists_50}.
        Take $t_0\in\reals$ such that
            $\sumFiniteSequence{b_0}{t_0}$ and
            $t=t_0+b(n+1)$ by
            \cref{sum_finite_sequence_split_last}.
        $r_0=t_0$ by \cref{finite_affine_dot_property_50}.
        Let $v=\coordScalarSequence{N}{a(n+1)}{s(n+1)}$.
        $a(n+1)\in\reals$ and
            $s(n+1)\in\realsPower{J}$ by \cref{funs_type_apply}.
        $v\in\funs{J}{\reals}$ by
            \cref{reals_power,tuple_power,coord_scalar_sequence_function}.
        Let $w=\coordSumSequence{N}{y_0}{v}$.
        $w\in\funs{J}{\reals}$ by
            \cref{coord_sum_sequence_function}.
        $w\in\realsPower{J}$ by \cref{reals_power,tuple_power}.
        Show that for all $j\in J$ we have $y(j)=w(j)$.
        \begin{subproof}
            Fix $j\in J$.
            Take $d$ such that
                $d\in\funs{I_1}{\reals}$ and
                for all $i\in I_1$ we have
                    $d(i)=a(i)\rmul\apply{\apply{s}{i}}{j}$ and
                $\sumFiniteSequence{d}{y(j)}$ by
                \cref{affine_combination_value_50,funs_elim,real_lt_subst_left,real_lt_subst_right}.
            Let $e=\restrl{d}{I}$.
            Take $u\in\reals$ such that
                $\sumFiniteSequence{e}{u}$ and
                $y(j)=u+d(n+1)$ by
                \cref{sum_finite_sequence_split_last}.
            Follows by
                \cref{reals_power,tuple_power,initial_segment_subset_succ,subseteq,funs_elim,restrl_of_function_apply,successor_sequence_restriction_50,coord_scalar_sequence,funs_is_function,function_apply_intro,pair_eq_iff,coord_sum_sequence,real_lt_subst_left,real_lt_subst_right}.
        \end{subproof}
        $y=w$ by
            \cref{reals_power,tuple_power,funs_is_function,functions_eq_iff}.
        Follows by
            \cref{finite_real_dot_product_scalar_50,finite_real_dot_product_add_50,real_lt_subst_left,real_lt_subst_right,finite_real_dot_product_unique_50}.
\end{proof}

% from §50 helper lemma 82: the finite affine-dot property holds in every positive length
\begin{lemma}\label{finite_affine_dot_all_lengths_50}
    For all $n\in\naturalsPlus$ we have
        $\finiteAffineDotProperty{n}$.
\end{lemma}
\begin{proof}
    Let $A=\{n\in\naturalsPlus\mid\finiteAffineDotProperty{n}\}$.
    Follows by
        \cref{subseteq,real_one_in_naturals,finite_affine_dot_base_50,pair_eq_iff,finite_affine_dot_successor_50,real_naturals_closed_succ,real_naturals_induction}.
\end{proof}

% from §50 helper lemma 83: a finite linear combination commutes with a dot product
\begin{lemma}\label{finite_affine_combination_dot_50}
    Assume $n,N\in\naturalsPlus$.
    Let $I=\initialSegment{n}$.
    Let $J=\initialSegment{N}$.
    Assume $s\in\funs{I}{\realsPower{J}}$.
    Assume $a\in\funs{I}{\reals}$.
    Assume $y,c\in\realsPower{J}$.
    Assume $y$ is an affine combination of $s$ with coefficients $a$
        in dimension $N$.
    Assume $q,b\in\funs{I}{\reals}$.
    Assume for all $i\in I$ we have
        $q(i)$ is the dot product of $s(i)$ and $c$ of length $N$.
    Assume for all $i\in I$ we have $b(i)=a(i)\rmul q(i)$.
    Assume $r,t\in\reals$.
    Assume $r$ is the dot product of $y$ and $c$ of length $N$.
    Assume $\sumFiniteSequence{b}{t}$.
    Then $r=t$.
\end{lemma}
\begin{proof}
    Follows by
        \cref{finite_affine_dot_all_lengths_50,real_one_in_naturals,real_one_leq_naturalsplus,initial_segment_naturalsplus,subseteq,finite_affine_dot_property_50}.
\end{proof}

% from §50 helper definition 33: append a final coordinate one to a euclidean point
\begin{definition}\label{augmented_affine_row_50}
    $\augmentedAffineRow{N}{x}=
        \{(j,r)\mid j\in\initialSegment{N+1},r\in\reals\mid
            \text{either $j=N+1$ and $r=1$ or
            $j\neq N+1$ and $r=x(j)$}\}$.
\end{definition}

% from §50 helper lemma 84: augmented affine rows are real vectors
\begin{lemma}\label{augmented_affine_row_function_50}
    Assume $N\in\naturalsPlus$.
    Let $I=\initialSegment{N}$.
    Let $I_1=\initialSegment{N+1}$.
    Assume $x\in\realsPower{I}$.
    Let $u=\augmentedAffineRow{N}{x}$.
    Then $u\in\realsPower{I_1}$ and
        $\restrl{u}{I}=x$ and
        $u(N+1)=1$.
\end{lemma}
\begin{proof}
    $u\in\realsPower{I_1}$ by
        \cref{reals_power,tuple_power,augmented_affine_row_50,relation,rightunique,pair_eq_iff,subseteq,dom_iff,initial_segment_succ_split,real_one_in_naturals,real_naturals_subset_reals,funs_type_apply,dom_intro,subseteq_antisymmetric,function_apply_intro,real_lt_subst_left,funs_intro}.
    $u(N+1)=1$ by
        \cref{reals_power,tuple_power,real_naturals_closed_succ,initial_segment_naturalsplus,real_one_in_naturals,real_naturals_subset_reals,subseteq,augmented_affine_row_50,pair_eq_iff,funs_is_function,function_apply_intro}.
    $\restrl{u}{I}=x$ by
        \cref{reals_power,tuple_power,successor_sequence_restriction_50,initial_segment_subset_succ,subseteq,initial_segment_naturalsplus,real_naturals_lt_succ_local,real_naturals_subset_reals,real_add_closed,real_one_in_naturals,real_leq_lt_trans,real_lt_irreflexive,real_lt_subst_left,funs_type_apply,augmented_affine_row_50,pair_eq_iff,function_apply_intro,restrl_of_function_apply,funs_is_function,funs_elim,functions_eq_iff}.
\end{proof}

% from §50 helper definition 34: matrix of augmented affine rows
\begin{definition}\label{augmented_affine_matrix_50}
    $\augmentedAffineMatrix{n}{N}{g}=
        \{(i,\augmentedAffineRow{N}{g(i)})\mid
            i\in\initialSegment{n}\}$.
\end{definition}

% from §50 helper lemma 85: an augmented affine matrix is a real matrix
\begin{lemma}\label{augmented_affine_matrix_function_50}
    Assume $n,N\in\naturalsPlus$.
    Let $I=\initialSegment{n}$.
    Let $J=\initialSegment{N}$.
    Assume $g\in\funs{I}{\realsPower{J}}$.
    Let $M=\augmentedAffineMatrix{n}{N}{g}$.
    Then $M\in\funs{I}{\realsPower{\initialSegment{N+1}}}$ and
        for all $i\in I$ we have
            $\restrl{M(i)}{J}=g(i)$ and
            $\apply{\apply{M}{i}}{N+1}=1$.
\end{lemma}
\begin{proof}
    Follows by
        \cref{augmented_affine_matrix_50,relation,rightunique,pair_eq_iff,dom_intro,dom_iff,subseteq,subseteq_antisymmetric,function_apply_intro,funs_type_apply,augmented_affine_row_function_50,real_lt_subst_left,funs_intro,funs_is_function,real_lt_subst_right}.
\end{proof}

% from §50 helper lemma 86: every vector has zero dot product with the zero vector
\begin{lemma}\label{finite_dot_zero_vector_50}
    Assume $N\in\naturalsPlus$.
    Let $J=\initialSegment{N}$.
    Assume $x\in\realsPower{J}$.
    Then $0$ is the dot product of $x$ and $\zeroVector{N}$ of length $N$.
\end{lemma}
\begin{proof}
    Follows by
        \cref{reals_power,tuple_power,zero_vector_euclidean_50,coord_product_sequence_function,zero_vector_50,pair_eq_iff,funs_is_function,function_apply_intro,coord_product_sequence,funs_type_apply,real_add_ident,real_mul_zero,real_mul_comm,real_lt_subst_left,real_lt_subst_right,sum_finite_sequence_exists,sum_finite_sequence_all_zero,finite_real_dot_product_50}.
\end{proof}

% from §50 helper lemma 87: at most N points admit a nonconstant vanishing affine functional
\begin{lemma}\label{finite_points_affine_functional_50}
    Assume $n,N\in\naturalsPlus$.
    Assume $n\leq N$.
    Let $I=\initialSegment{n}$.
    Let $J=\initialSegment{N}$.
    Assume $g\in\funs{I}{\realsPower{J}}$.
    Then there exist $c,d$ such that
        $c\in\realsPower{J}$ and
        $c\neq\zeroVector{N}$ and
        $d\in\reals$ and
        for all $i\in I$ there exists $q\in\reals$ such that
            $q$ is the dot product of $g(i)$ and $c$ of length $N$ and
            $q+d=0$.
\end{lemma}
\begin{proof}
    $n<N+1$ and $N+1\in\naturalsPlus$ by
        \cref{real_naturals_subset_reals,subseteq,real_one_in_naturals,real_add_closed,real_naturals_lt_succ_local,real_leq_lt_trans,real_naturals_closed_succ}.
    Let $M=\augmentedAffineMatrix{n}{N}{g}$.
    $M\in\funs{I}{\realsPower{\initialSegment{N+1}}}$ and
        for all $i\in I$ we have
            $\restrl{M(i)}{J}=g(i)$ and
            $\apply{\apply{M}{i}}{N+1}=1$ by
        \cref{augmented_affine_matrix_function_50}.
    Take $z$ such that
        $z\in\realsPower{\initialSegment{N+1}}$ and
        $z\neq\zeroVector{N+1}$ and
        for all $i\in I$ we have
            $0$ is the dot product of $M(i)$ and $z$ of length $N+1$
        by \cref{finite_homogeneous_kernel_50}.
    Let $c=\restrl{z}{J}$.
    Let $d=z(N+1)$.
    $c\in\realsPower{J}$ and $d\in\reals$ by
        \cref{reals_power,tuple_power,successor_sequence_restriction_50,real_naturals_closed_succ,initial_segment_naturalsplus,funs_type_apply}.
    For all $i\in I$ there exists $q\in\reals$ such that
        $q$ is the dot product of $g(i)$ and $c$ of length $N$ and
        $q+d=0$ by
        \cref{funs_type_apply,reals_power,tuple_power,finite_real_dot_product_exists_50,subseteq,real_lt_subst_left,real_lt_subst_right,finite_real_dot_product_split_last_50,real_one_in_naturals,real_naturals_subset_reals,real_mul_ident}.
    $c\neq\zeroVector{N}$ by
        \cref{real_one_in_naturals,real_one_leq_naturalsplus,initial_segment_naturalsplus,subseteq,funs_type_apply,finite_dot_zero_vector_50,real_lt_subst_left,real_lt_subst_right,finite_real_dot_product_unique_50,real_add_ident,zero_vector_euclidean_50,reals_power,tuple_power,zero_vector_50,pair_eq_iff,funs_is_function,function_apply_intro,initial_segment_succ_split,funs_elim,initial_segment_subset_succ,restrl_of_function_apply,functions_eq_iff}.
\end{proof}

% from §50 helper lemma 88: an affine functional vanishing on generators vanishes on their plane
\begin{lemma}\label{affine_plane_functional_50}
    Assume $n,N\in\naturalsPlus$.
    Assume $n\leq N$.
    Let $I=\initialSegment{n}$.
    Let $J=\initialSegment{N}$.
    Let $E=\realsPower{J}$.
    Assume $A$ is a geometrically independent set in dimension $N$.
    Assume $A$ is inhabited.
    Assume $g$ is a bijection from $I$ to $A$.
    Assume $P$ is the plane determined by $A$ in dimension $N$.
    Then there exist $c,d$ such that
        $c\in E$ and
        $c\neq\zeroVector{N}$ and
        $d\in\reals$ and
        for all $y\in P$ there exists $r\in\reals$ such that
            $r$ is the dot product of $y$ and $c$ of length $N$ and
            $r+d=0$.
\end{lemma}
\begin{proof}
    $g\in\funs{I}{E}$ by
        \cref{geometrically_independent_set_50,bijections_to_funs,funs_elim,funs_type_apply,subseteq,funs_intro}.
    Take $c,d$ such that
        $c\in E$ and
        $c\neq\zeroVector{N}$ and
        $d\in\reals$ and
        for all $i\in I$ there exists $q\in\reals$ such that
            $q$ is the dot product of $g(i)$ and $c$ of length $N$ and
            $q+d=0$
        by \cref{finite_points_affine_functional_50}.
    Let $q=\{(i,u)\mid i\in I,u\in\reals\mid
        \text{$u$ is the dot product of $g(i)$ and $c$ of length $N$}\}$.
    $q\in\funs{I}{\reals}$ by
        \cref{relation,rightunique,pair_eq_iff,finite_real_dot_product_unique_50,subseteq,dom_iff,dom_intro,subseteq_antisymmetric,function_apply_intro,real_lt_subst_left,funs_intro}.
    For all $i\in I$ we have
        $q(i)$ is the dot product of $g(i)$ and $c$ of length $N$ and
        $q(i)+d=0$ by
        \cref{funs_is_function,funs_elim,function_apply_elim,subseteq,pair_eq_iff,finite_real_dot_product_unique_50,real_lt_subst_left}.
    Show that for all $y\in P$ there exists $r\in\reals$ such that
        $r$ is the dot product of $y$ and $c$ of length $N$ and
        $r+d=0$.
    \begin{subproof}
        Fix $y\in P$.
        Take $a$ such that
            $a\in\funs{I}{\reals}$ and
            $\sumFiniteSequence{a}{1}$ and
            $y$ is an affine combination of $g$ with coefficients $a$
                in dimension $N$
            by \cref{plane_coordinates_fixed_enumeration_50}.
        Let $b=\coordProductSequence{n}{a}{q}$.
        Let $h=\coordScalarSequence{n}{\neg{d}}{a}$.
        $\sumFiniteSequence{b}{\neg{d}}$ by
            \cref{coord_product_sequence_function,coord_product_sequence,real_add_inverse,real_add_comm,real_add_cancel,coord_scalar_sequence_function,coord_scalar_sequence,funs_is_function,function_apply_intro,funs_type_apply,real_mul_comm,real_lt_subst_left,real_lt_subst_right,subseteq,functions_eq_iff,sum_finite_sequence_scalar,real_mul_ident}.
        Thus there exists $u\in\reals$ such that
            $u$ is the dot product of $y$ and $c$ of length $N$ and
            $u+d=0$ by
            \cref{plane_determined_by_points_50,finite_real_dot_product_exists_50,reals_power,tuple_power,coord_product_sequence_function,coord_product_sequence,funs_is_function,function_apply_intro,real_lt_subst_left,real_add_inverse,finite_affine_combination_dot_50,real_add_comm}.
    \end{subproof}
\end{proof}

% from §50 helper lemma 89: finite real dot products are symmetric
\begin{lemma}\label{finite_real_dot_product_symmetric_50}
    Assume $N\in\naturalsPlus$.
    Assume $r$ is the dot product of $a$ and $b$ of length $N$.
    Then $r$ is the dot product of $b$ and $a$ of length $N$.
\end{lemma}
\begin{proof}
    Follows by
        \cref{coord_product_sequence_function,coord_product_sequence,real_lt_subst_left,funs_is_function,function_apply_intro,funs_type_apply,real_mul_comm,real_lt_subst_right,functions_eq_iff,finite_real_dot_product_50}.
\end{proof}

% from §50 helper lemma 90: a coordinate unit extracts the matching coordinate by a dot product
\begin{lemma}\label{coordinate_unit_dot_50}
    Assume $N\in\naturalsPlus$.
    Let $J=\initialSegment{N}$.
    Assume $j\in J$.
    Assume $c\in\realsPower{J}$.
    Let $u=\standardAffineFramePoint{N}{j+1}$.
    Then $c(j)$ is the dot product of $u$ and $c$ of length $N$.
\end{lemma}
\begin{proof}
    $u,c\in\funs{J}{\reals}$ by
        \cref{standard_affine_frame_point_coordinates_50,initial_segment_naturalsplus,real_naturals_closed_succ,real_leq_add,real_one_in_naturals,real_naturals_subset_reals,subseteq,reals_power,tuple_power}.
    Let $b=\coordProductSequence{N}{u}{c}$.
    Follows by
        \cref{coord_product_sequence_function,initial_segment_naturalsplus,real_naturals_subset_reals,real_one_in_naturals,subseteq,real_add_cancel,real_one_leq_naturalsplus,real_naturals_lt_succ_local,real_leq_lt_trans,real_add_closed,real_lt_irreflexive,real_lt_subst_right,real_leq_add,real_naturals_closed_succ,standard_affine_frame_point_off_coordinate_50,coord_product_sequence,funs_is_function,function_apply_intro,real_add_ident,funs_type_apply,real_mul_zero,real_lt_subst_left,standard_affine_frame_point_one_50,real_mul_ident,finite_sum_single_support_50,finite_real_dot_product_50}.
\end{proof}

% from §50 helper lemma 91: an affine plane generated by at most N points is proper
\begin{lemma}\label{small_affine_plane_proper_50}
    Assume $n,N\in\naturalsPlus$.
    Assume $n\leq N$.
    Let $I=\initialSegment{n}$.
    Let $J=\initialSegment{N}$.
    Let $E=\realsPower{J}$.
    Assume $A$ is a geometrically independent set in dimension $N$.
    Assume $A$ is inhabited.
    Assume $g$ is a bijection from $I$ to $A$.
    Assume $P$ is the plane determined by $A$ in dimension $N$.
    Then $P\neq E$.
\end{lemma}
\begin{proof}
    Suppose not.
    Take $c,d$ such that
        $c\in E$ and
        $c\neq\zeroVector{N}$ and
        $d\in\reals$ and
        for all $y\in P$ there exists $r\in\reals$ such that
            $r$ is the dot product of $y$ and $c$ of length $N$ and
            $r+d=0$
        by \cref{affine_plane_functional_50}.
    Take $j\in J$ such that $c(j)\neq 0$ by
        \cref{reals_power,tuple_power,zero_vector_euclidean_50,zero_vector_50,pair_eq_iff,funs_is_function,function_apply_intro,real_lt_subst_right,functions_eq_iff}.
    \begin{byCase}
    \caseOf{$d\neq 0$.}
        Let $w=\zeroVector{N}$.
        $w\in E,P$ by \cref{zero_vector_euclidean_50,real_lt_subst_left}.
        Contradiction by
            \cref{subseteq,finite_dot_zero_vector_50,finite_real_dot_product_symmetric_50,finite_real_dot_product_unique_50,real_lt_subst_left,real_lt_subst_right,real_add_ident}.
    \caseOf{$d=0$.}
        Let $w=\standardAffineFramePoint{N}{j+1}$.
        $j,N,1\in\reals$ and $j\leq N$ by
            \cref{initial_segment_naturalsplus,real_naturals_subset_reals,real_one_in_naturals,subseteq}.
        $j+1\leq N+1$ and $j+1\in\initialSegment{N+1}$ by
            \cref{real_leq_add,real_naturals_closed_succ,initial_segment_naturalsplus}.
        $w\in E,P$ by
            \cref{standard_affine_frame_point_coordinates_50,real_lt_subst_left}.
        Contradiction by
            \cref{subseteq,coordinate_unit_dot_50,finite_real_dot_product_unique_50,real_lt_subst_left,real_lt_subst_right,real_add_ident,real_add_comm}.
    \end{byCase}
\end{proof}

% from §50 helper lemma 92: a single-coordinate perturbation has controlled metric and dot product
\begin{lemma}\label{coordinate_perturbation_50}
    Assume $N\in\naturalsPlus$.
    Let $J=\initialSegment{N}$.
    Let $E=\realsPower{J}$.
    Assume $j\in J$.
    Assume $y,c\in E$.
    Assume $t\in\reals$ and $0<t$.
    Assume $r$ is the dot product of $y$ and $c$ of length $N$.
    Let $u=\standardAffineFramePoint{N}{j+1}$.
    Let $v=\coordScalarSequence{N}{t}{u}$.
    Let $w=\coordSumSequence{N}{y}{v}$.
    Then $w\in E$ and
        $\apply{\squareMetric{N}}{(y,w)}=t$ and
        $r+(t\rmul c(j))$ is the dot product of $w$ and $c$ of length $N$.
\end{lemma}
\begin{proof}
    $u\in E$ by
        \cref{standard_affine_frame_point_coordinates_50,initial_segment_naturalsplus,real_naturals_closed_succ,real_leq_add,real_one_in_naturals,real_naturals_subset_reals,subseteq}.
    $j,N,1\in\reals$ and $j\leq N$ and $j+1\leq N+1$ and
        $j+1\in\initialSegment{N+1}$ by
        \cref{initial_segment_naturalsplus,real_naturals_closed_succ,real_leq_add,real_one_in_naturals,real_naturals_subset_reals,subseteq}.
    $u,y,c\in\funs{J}{\reals}$ by \cref{reals_power,tuple_power}.
    $0,\neg{t}\in\reals$ and $v,w\in\funs{J}{\reals}$ and $w\in E$ by
        \cref{real_add_ident,real_add_inverse,coord_scalar_sequence_function,coord_sum_sequence_function,reals_power,tuple_power}.
    $c(j)$ is the dot product of $u$ and $c$ of length $N$ by
        \cref{coordinate_unit_dot_50}.
    $t\rmul c(j)$ is the dot product of $v$ and $c$ of length $N$ by
        \cref{finite_real_dot_product_scalar_50}.
    $r+(t\rmul c(j))$ is the dot product of $w$ and $c$ of length $N$ by
        \cref{finite_real_dot_product_add_50}.
    Take $m\in\reals$ such that
        $((y,w),m)\in\squareMetric{N}$ and
        for all $i\in J$ we have $\abs{y(i)-w(i)}\leq m$
        by \cref{square_metric_value_exists}.
    $\apply{\squareMetric{N}}{(y,w)}=m$ by
        \cref{square_metric_function,function_apply_intro}.
    For all $i\in J$ we have
        $\abs{y(i)-w(i)}\leq t$ by
        \cref{coord_scalar_sequence,funs_is_function,function_apply_intro,coord_sum_sequence,funs_type_apply,standard_affine_frame_point_one_50,real_mul_ident,real_mul_comm,real_lt_subst_left,real_lt_subst_right,real_sub,real_neg_neg,real_add_eq_subst,real_sub_sub_cancel,real_lt_neg,real_neg_zero,real_abs_neg,real_lt_imp_leq,real_abs_nonneg,initial_segment_naturalsplus,real_naturals_subset_reals,real_one_in_naturals,subseteq,real_add_cancel,real_one_leq_naturalsplus,real_naturals_lt_succ_local,real_leq_lt_trans,real_add_closed,real_lt_irreflexive,real_leq_add,real_naturals_closed_succ,standard_affine_frame_point_off_coordinate_50,real_add_ident,real_mul_zero,real_add_comm,real_add_inverse}.
    $\apply{\squareMetric{N}}{(y,w)}=t$ by
        \cref{square_metric,pair_eq_iff,coord_diff_set,real_lt_subst_left,coord_scalar_sequence,funs_is_function,function_apply_intro,standard_affine_frame_point_one_50,real_mul_ident,real_mul_comm,real_lt_subst_right,coord_sum_sequence,funs_type_apply,real_sub,real_neg_neg,real_add_eq_subst,real_sub_sub_cancel,real_lt_neg,real_neg_zero,real_abs_neg,real_lt_imp_leq,real_abs_nonneg,subseteq,real_leq_antisym}.
\end{proof}

% from §50 helper definition 35: one weighted coordinate function
\begin{definition}\label{fixed_dot_coordinate_function_50}
    $\fixedDotCoordinateFunction{R}{E}{c}{j}=
        \{(y,c(j)\rmul\apply{\piIndex{R}{j}}{y})\mid y\in E\}$.
\end{definition}

% from §50 helper lemma 93: a fixed finite dot-product functional is continuous
\begin{lemma}\label{finite_dot_product_function_continuous_50}
    Assume $N\in\naturalsPlus$.
    Let $J=\initialSegment{N}$.
    Let $E=\realsPower{J}$.
    Let $S=\metricTopology{\squareMetric{N}}{E}$.
    Assume $c\in E$.
    Let $h=\{(y,r)\mid y\in E,r\in\reals\mid
        \text{$r$ is the dot product of $y$ and $c$ of length $N$}\}$.
    Then $h$ is continuous from $E$ to $\reals$ with respect to
        $S$ and $\standardTopologyR$.
\end{lemma}
\begin{proof}
    $c\in\funs{J}{\reals}$ by \cref{reals_power,tuple_power}.
    $h\in\funs{E}{\reals}$ by
        \cref{relation,rightunique,pair_eq_iff,finite_real_dot_product_unique_50,subseteq,dom_iff,reals_power,tuple_power,finite_real_dot_product_exists_50,dom_intro,subseteq_antisymmetric,function_apply_intro,real_lt_subst_left,funs_intro}.
    Let $R_0=\{(j,\reals)\mid j\in J\}$.
    Let $T_0=\{(j,\standardTopologyR)\mid j\in J\}$.
    Let $P=\productTopologyIndexed{T_0}{R_0}$.
    $R_0$ is a function on $J$ and
        $T_0$ is a function on $J$ and for all $j\in J$ we have
        $R_0(j)=\reals$ and $T_0(j)=\standardTopologyR$ by
        \cref{relation,rightunique,dom_intro,dom_iff,pair_eq_iff,subseteq,subseteq_antisymmetric,funs,function_apply_intro}.
    For all $j\in\dom{R_0}$ we have
        $T_0(j)$ is a topology on $R_0(j)$ and
        $P$ is a topology on $\productFamily{R_0}$ by
        \cref{standard_topology_reals,standard_basis_is_basis,basis_topology_is_topology,real_lt_subst_left,real_lt_subst_right,subseteq,real_one_in_naturals,real_one_leq_naturalsplus,initial_segment_naturalsplus,product_subbasis_finite_intersections_basis,product_topology_indexed,basis_for_topology}.
    $\productFamily{R_0}=E$ and $P=S$ by
        \cref{real_one_in_naturals,real_one_leq_naturalsplus,initial_segment_naturalsplus,product_family_reals_power,euclidean_metric_product_topology,square_metric_topology_equals_euclidean,real_lt_subst_left,real_lt_subst_right}.
    Let $p=\{(j,\fixedDotCoordinateFunction{R_0}{E}{c}{j})
        \mid j\in J\}$.
    $p$ is a function on $J$ by
        \cref{relation,rightunique,dom_intro,dom_iff,pair_eq_iff,subseteq,subseteq_antisymmetric,funs}.
    Show that for all $j\in J$ we have
        $p(j)$ is continuous from $E$ to $\reals$ with respect to
            $P$ and $\standardTopologyR$.
    \begin{subproof}
        Fix $j\in J$.
        Let $f=\piIndex{R_0}{j}$.
        $f$ is continuous from $E$ to $\reals$ with respect to
            $P$ and $\standardTopologyR$ by
            \cref{projection_map_indexed_continuous,real_lt_subst_left,real_lt_subst_right}.
        Let $k=\{(y,c(j))\mid y\in E\}$.
        Let $g_0=\{(y,k(y)\rmul f(y))\mid y\in E\}$.
        $g_0$ is continuous from $E$ to $\reals$ with respect to
            $P$ and $\standardTopologyR$ by
            \cref{funs_type_apply,real_function_scalar_multiple_continuous_49}.
        $\fixedDotCoordinateFunction{R_0}{E}{c}{j}\subseteq g_0$ by
            \cref{subseteq,fixed_dot_coordinate_function_50,funs_type_apply,constant_map_function,pair_eq_iff,funs_is_function,function_apply_intro,real_lt_subst_left,real_lt_subst_right}.
        $g_0\subseteq
            \fixedDotCoordinateFunction{R_0}{E}{c}{j}$ by
            \cref{subseteq,funs_type_apply,constant_map_function,pair_eq_iff,funs_is_function,function_apply_intro,real_lt_subst_left,real_lt_subst_right,fixed_dot_coordinate_function_50}.
        $\fixedDotCoordinateFunction{R_0}{E}{c}{j}=g_0$ by
            \cref{subseteq_antisymmetric}.
        Follows.
    \end{subproof}
    Take $g$ such that
        $g$ is a finite real sum witness for $p$ on $E$ with respect to
            $P$ and $N$ by
        \cref{finite_real_sum_continuous}.
    $g$ is continuous from $E$ to $\reals$ with respect to
        $P$ and $\standardTopologyR$ and
        for all $y\in E$ there exists $a$ such that
            $a\in\funs{J}{\reals}$ and
            $\sumFiniteSequence{a}{g(y)}$ and
            for all $j\in J$ we have
                $a(j)=\apply{\apply{p}{j}}{y}$ by
        \cref{finite_real_sum_witness}.
    Show that for all $y\in E$ we have $h(y)=g(y)$.
    \begin{subproof}
        Fix $y\in E$.
        Let $b=\coordProductSequence{N}{y}{c}$.
        $\sumFiniteSequence{b}{h(y)}$ by
            \cref{funs_is_function,funs_elim,function_apply_elim,subseteq,pair_eq_iff,finite_real_dot_product_50}.
        Take $a$ such that
            $a\in\funs{J}{\reals}$ and
            $\sumFiniteSequence{a}{g(y)}$ and
            for all $j\in J$ we have
                $a(j)=\apply{\apply{p}{j}}{y}$ by \cref{subseteq}.
        For all $j\in J$ we have $a(j)=b(j)$ by
            \cref{pair_eq_iff,coord_product_sequence_function,funs_is_function,function_apply_intro,fixed_dot_coordinate_function_50,continuous,real_lt_subst_left,projection_map_indexed_function,projection_map_indexed,coord_product_sequence,reals_power,tuple_power,funs_type_apply,real_mul_comm,real_lt_subst_right,subseteq}.
        Follows by
            \cref{coord_product_sequence_function,reals_power,tuple_power,funs_is_function,functions_eq_iff,real_lt_subst_left,sum_finite_sequence_unique}.
    \end{subproof}
    $h$ is continuous from $E$ to $\reals$ with respect to
        $S$ and $\standardTopologyR$ by
        \cref{continuous,funs_is_function,functions_eq_iff,real_lt_subst_left,real_lt_subst_right}.
\end{proof}

% from §50 helper lemma 94: affine dot-product level sets are closed
\begin{lemma}\label{affine_dot_level_set_closed_50}
    Assume $N\in\naturalsPlus$.
    Let $J=\initialSegment{N}$.
    Let $E=\realsPower{J}$.
    Let $S=\metricTopology{\squareMetric{N}}{E}$.
    Assume $c\in E$ and $d\in\reals$.
    Let $H=\{y\in E\mid\text{there exists $r\in\reals$ such that
        $r$ is the dot product of $y$ and $c$ of length $N$ and
        $r+d=0$}\}$.
    Then $H$ is closed in $E$ with respect to $S$.
\end{lemma}
\begin{proof}
    Let $h=\{(y,r)\mid y\in E,r\in\reals\mid
        \text{$r$ is the dot product of $y$ and $c$ of length $N$}\}$.
    $h$ is continuous from $E$ to $\reals$ with respect to
        $S$ and $\standardTopologyR$ by
        \cref{finite_dot_product_function_continuous_50}.
    Let $a=\neg{d}$.
    $a\in\reals$ by \cref{real_add_inverse}.
    $\reals$ is Hausdorff with respect to $\standardTopologyR$ by
        \cref{standard_metric_is_metric,metric_space_hausdorff,standard_metric_topology}.
    $\{a\}$ is closed in $\reals$ with respect to $\standardTopologyR$ by
        \cref{singleton_subset_intro,pair_is_finite,upair_intro_left,subseteq_finite_is_finite,finite_point_set_closed_hausdorff}.
    $\standardTopologyR$ is a topology on $\reals$ by
        \cref{standard_basis_is_basis,basis_topology_is_topology,standard_topology_reals}.
    For all $B$ such that $B\subseteq E$ we have
        $\img{h}{\closure{B}{E}{S}}
            \subseteq\closure{\img{h}{B}}{\reals}{\standardTopologyR}$ by
        \cref{continuous_closure_image_subset}.
    $\preimg{h}{\{a\}}$ is closed in $E$ with respect to $S$ by
        \cref{continuous,closure_image_subset_implies_preimg_closed}.
    $H=\preimg{h}{\{a\}}$ by
        \cref{subseteq_antisymmetric,subseteq,pair_eq_iff,continuous,funs_is_function,function_apply_intro,real_add_inverse,real_add_comm,real_lt_subst_left,real_lt_subst_right,real_add_cancel,singleton_iff,function_apply_elim,preimg_iff,preimg_subseteq_dom,funs}.
    Follows by \cref{real_lt_subst_left}.
\end{proof}

% from §50 helper lemma 95: nonconstant affine dot-product level sets have empty interior
\begin{lemma}\label{affine_dot_level_set_empty_interior_50}
    Assume $N\in\naturalsPlus$.
    Let $J=\initialSegment{N}$.
    Let $E=\realsPower{J}$.
    Let $S=\metricTopology{\squareMetric{N}}{E}$.
    Assume $c\in E$ and $c\neq\zeroVector{N}$ and $d\in\reals$.
    Let $H=\{y\in E\mid\text{there exists $r\in\reals$ such that
        $r$ is the dot product of $y$ and $c$ of length $N$ and
        $r+d=0$}\}$.
    Then $H$ has empty interior in $E$ with respect to $S$.
\end{lemma}
\begin{proof}
    $H\subseteq E$ by \cref{subseteq}.
    Suppose not.
    $\interior{H}{E}{S}\neq\emptyset$ by
        \cref{empty_interior}.
    Take $y$ such that $y\in\interior{H}{E}{S}$ by
        \cref{nonempty_inhabited}.
    Take $U\in S$ such that $y\in U$ and $U\subseteq H$ by
        \cref{interior}.
    $\squareMetric{N}$ is a metric on $E$ by
        \cref{square_metric_is_metric}.
    $S$ is a topology on $E$ by
        \cref{metric_topology,metric_basis_is_basis,basis_topology_is_topology}.
    $U$ is open in $E$ with respect to $S$ by
        \cref{open_in}.
    $U\subseteq E$ by
        \cref{topology_on,pow_iff,subseteq}.
    $y\in E$ by \cref{subseteq}.
    Take $\epsilon\in\reals$ such that $0<\epsilon$ and
        $\metricBall{\squareMetric{N}}{E}{y}{\epsilon}\subseteq U$ by
        \cref{metric_open_iff}.
    Let $t=\epsilon/(1+1)$.
    $t\in\reals$ and $0<t$ and $t+t=\epsilon$ by
        \cref{real_half_of_positive}.
    $0\in\reals$ by \cref{real_add_ident}.
    $0+t<t+t$ by \cref{real_lt_add}.
    $0+t=t$ by \cref{real_add_ident}.
    $t<\epsilon$ by
        \cref{real_lt_subst_left,real_lt_subst_right}.
    There exists $i\in J$ such that $c(i)\neq 0$ by
        \cref{reals_power,tuple_power,zero_vector_euclidean_50,zero_vector_50,pair_eq_iff,funs_is_function,function_apply_intro,real_lt_subst_right,functions_eq_iff}.
    Take $j\in J$ such that $c(j)\neq 0$.
    Take $r\in\reals$ such that
        $r$ is the dot product of $y$ and $c$ of length $N$ by
        \cref{finite_real_dot_product_exists_50,reals_power,tuple_power}.
    Let $u=\standardAffineFramePoint{N}{j+1}$.
    Let $v=\coordScalarSequence{N}{t}{u}$.
    Let $w=\coordSumSequence{N}{y}{v}$.
    $w\in E$ and
        $\apply{\squareMetric{N}}{(y,w)}=t$ and
        $r+(t\rmul c(j))$ is the dot product of $w$ and $c$ of length $N$ by
        \cref{coordinate_perturbation_50}.
    $\apply{\squareMetric{N}}{(y,w)}<\epsilon$ by
        \cref{real_lt_subst_left}.
    $w\in\metricBall{\squareMetric{N}}{E}{y}{\epsilon}$ by
        \cref{metric_ball}.
    $w\in U$ by \cref{subseteq}.
    $w\in H$ by \cref{subseteq}.
    Take $q\in\reals$ such that
        $q$ is the dot product of $w$ and $c$ of length $N$ and
        $q+d=0$ by \cref{subseteq}.
    $q=r+(t\rmul c(j))$ by
        \cref{finite_real_dot_product_unique_50}.
    $y\in H$ by \cref{interior,subseteq}.
    Take $s\in\reals$ such that
        $s$ is the dot product of $y$ and $c$ of length $N$ and
        $s+d=0$ by \cref{subseteq}.
    $s=r$ by \cref{finite_real_dot_product_unique_50}.
    $r+d=0$ by \cref{real_lt_subst_left}.
    Let $p=t\rmul c(j)$.
    $c(j)\in\reals$ by
        \cref{reals_power,tuple_power,funs_type_apply}.
    $p\in\reals$ by \cref{real_mul_closed}.
    $(r+p)+d=0$ by
        \cref{real_lt_subst_left,real_lt_subst_right}.
    $(r+p)+d=(r+d)+p$ by
        \cref{real_add_assoc,real_add_comm}.
    $(r+d)+p=0+p$ by \cref{real_add_eq_subst}.
    $0+p=p$ by \cref{real_add_ident}.
    $p=0$ by
        \cref{real_lt_subst_left,real_lt_subst_right}.
    $t=0$ or $c(j)=0$ by
        \cref{real_mul_zero_product}.
    Contradiction by
        \cref{real_lt_subst_right,real_lt_irreflexive}.
\end{proof}

% from §50 helper lemma 96: a finite union of closed empty-interior sets has empty interior
\begin{lemma}\label{finite_closed_empty_interior_union_50}
    Assume $T$ is a topology on $X$.
    Assume $X$ is a Baire space with respect to $T$.
    Assume $k\in\naturalsPlus$.
    Let $I=\initialSegment{k}$.
    Assume $h$ is a bijection from $I$ to $F$.
    Assume $F\subseteq\pow{X}$.
    Assume for all $H\in F$ we have
        $H$ is closed in $X$ with respect to $T$ and
        $H$ has empty interior in $X$ with respect to $T$.
    Then $\unions{F}$ has empty interior in $X$ with respect to $T$.
\end{lemma}
\begin{proof}
    Let $s=\{(n,U)\mid n\in\naturalsPlus,U\in\pow{X}\mid
        (n\in I\land U=h(n))\lor(n\notin I\land U=\emptyset)\}$.
    $s$ is a relation by \cref{relation,pair_eq_iff}.
    $s$ is right-unique by
        \cref{rightunique,pair_eq_iff}.
    $s$ is a function by \cref{relation,rightunique}.
    $\dom{s}=\naturalsPlus$ by
        \cref{subseteq,dom_iff,pair_eq_iff,bijections_to_funs,funs_type_apply,powerset_bottom,dom_intro,subseteq_antisymmetric}.
    For all $n\in\dom{s}$ we have $s(n)\in\pow{X}$ by
        \cref{function_apply_iff,pair_eq_iff}.
    $s\in\funs{\naturalsPlus}{\pow{X}}$ by
        \cref{funs_intro}.
    $s$ is a sequence in $\pow{X}$ by
        \cref{sequence_in}.
    $\emptyset$ is closed in $X$ with respect to $T$ by
        \cref{closed_emptyset}.
    $\interior{\emptyset}{X}{T}=\emptyset$ by
        \cref{interior,subseteq,emptyset,emptyset_subseteq,subseteq_antisymmetric}.
    $\emptyset$ has empty interior in $X$ with respect to $T$ by
        \cref{empty_interior,emptyset_subseteq}.
    For all $n\in\naturalsPlus$ we have
        $s(n)$ is closed in $X$ with respect to $T$ and
        $s(n)$ has empty interior in $X$ with respect to $T$ by
        \cref{bijections_to_funs,funs_type_apply,subseteq,pair_eq_iff,funs_is_function,function_apply_intro,real_lt_subst_left,powerset_bottom}.
    $\unions{\img{s}{\naturalsPlus}}$ has empty interior in $X$
        with respect to $T$ by \cref{baire_space}.
    For all $x\in\unions{\img{s}{\naturalsPlus}}$ we have
        $x\in\unions{F}$ by
        \cref{unions_iff,img_of_function_elim,inter_elim_right,bijections_to_funs,funs_type_apply,subseteq,pair_eq_iff,funs_is_function,function_apply_intro,real_lt_subst_left,real_lt_subst_right,powerset_bottom,emptyset}.
    Thus $\unions{\img{s}{\naturalsPlus}}\subseteq\unions{F}$ by
        \cref{subseteq}.
    For all $x\in\unions{F}$ we have
        $x\in\unions{\img{s}{\naturalsPlus}}$ by
        \cref{unions_iff,bijections,surj,initial_segment_naturalsplus,bijections_to_funs,funs_type_apply,subseteq,pair_eq_iff,funs_is_function,function_apply_intro,function_apply_elim,img_iff,real_lt_subst_left,real_lt_subst_right}.
    Thus $\unions{F}\subseteq\unions{\img{s}{\naturalsPlus}}$ by
        \cref{subseteq}.
    $\unions{\img{s}{\naturalsPlus}}=\unions{F}$ by
        \cref{subseteq_antisymmetric}.
    Follows by \cref{real_lt_subst_left}.
\end{proof}

% from §50 helper lemma 97: append a new point to a finite bijective enumeration
\begin{lemma}\label{finite_bijection_append_new_point_50}
    Assume $n\in\naturalsPlus$.
    Let $I=\initialSegment{n}$.
    Let $I_1=\initialSegment{n+1}$.
    Assume $g$ is a bijection from $I$ to $A$.
    Assume $y\notin A$.
    Then there exists $h$ such that
        $h$ is a bijection from $I_1$ to $A\union\{y\}$ and
        $\restrl{h}{I}=g$ and $h(n+1)=y$.
\end{lemma}
\begin{proof}
    Let $r=\{(n+1,y)\}$.
    Let $h=g\union r$.
    $g$ is a function and $\dom{g}=I$ by
        \cref{bijection_is_function,bijections_dom}.
    $r$ is a function by
        \cref{singleton_elim,pair_eq_iff,relation,rightunique}.
    $\dom{r}=\{n+1\}$ by
        \cref{dom_iff,singleton_elim,pair_eq_iff,singleton_intro,subseteq,subseteq_antisymmetric}.
    $n+1\notin I$ by
        \cref{initial_segment_naturalsplus,real_naturals_subset_reals,subseteq,real_naturals_closed_succ,real_naturals_lt_succ_local,real_lt_trans,real_lt_irreflexive}.
    $\dom{g}\inter\dom{r}=\emptyset$ by
        \cref{inter,subseteq,emptyset,subseteq_antisymmetric,emptyset_subseteq,singleton_iff,real_lt_subst_left,real_lt_subst_right}.
    $h$ is a function by
        \cref{union_compatible_functions,emptyset,inter}.
    $\dom{h}=I\union\{n+1\}$ by
        \cref{dom_union,real_lt_subst_left,real_lt_subst_right}.
    $\dom{h}=I_1$ by
        \cref{real_lt_subst_left,union_iff,initial_segment_subset_succ,subseteq,singleton_iff,real_naturals_closed_succ,initial_segment_naturalsplus,initial_segment_succ_split,subseteq_antisymmetric}.
    For all $i\in\dom{h}$ we have $h(i)\in A\union\{y\}$ by
        \cref{funs,function_apply_elim,union_iff,dom_intro,bijections_to_funs,funs_type_apply,funs_is_function,function_apply_intro,real_lt_subst_left,union_intro_left,singleton_elim,pair_eq_iff,singleton_iff,union_intro_right}.
    $h\in\funs{I_1}{A\union\{y\}}$ by
        \cref{funs_intro}.
    For all $i,j$ such that
        $i\in\dom{h}$ and $j\in\dom{h}$ and $h(i)=h(j)$
    we have $i=j$ by
        \cref{real_lt_subst_left,initial_segment_succ_split,subseteq,function_apply_elim,union_iff,function_apply_intro,bijections_to_funs,funs_type_apply,singleton_intro,union_intro_right,funs_is_function,real_lt_subst_right,bijections,injective_function}.
    Thus $h$ is injective by \cref{injective_function}.
    For all $z\in A\union\{y\}$ there exists $i\in I_1$ such that
        $h(i)=z$ by
        \cref{union_iff,bijections,surj,initial_segment_subset_succ,subseteq,bijections_dom,function_apply_elim,function_apply_intro,real_lt_subst_right,singleton_iff,real_naturals_closed_succ,initial_segment_naturalsplus,singleton_intro,union_intro_right,funs_is_function,real_lt_subst_left}.
    $h\in\surj{I_1}{A\union\{y\}}$ by \cref{surj}.
    $h$ is a bijection from $I_1$ to $A\union\{y\}$ by
        \cref{bijections}.
    Show that $\restrl{h}{I}=g$.
    \begin{subproof}
        Let $u=\restrl{h}{I}$.
        $I\subseteq\dom{h}$ by
            \cref{initial_segment_subset_succ,real_lt_subst_left}.
        $u$ is a function by
            \cref{restrl_of_function_is_function}.
        $\dom{u}=I$ by
            \cref{dom_of_restrl_eq_inter,inter_absorb_supseteq_left}.
        $g$ is a function and $\dom{g}=I$ by
            \cref{bijection_is_function,bijections_dom}.
        For all $i\in I$ we have $u(i)=g(i)$ by
            \cref{restrl_of_function_apply,funs_is_function,bijections_dom,subseteq,function_apply_elim,union_iff,function_apply_intro}.
        $u=g$ by
            \cref{fun_subseteq,subseteq_refl,subseteq,real_lt_subst_left,subseteq_antisymmetric}.
        Follows by \cref{real_lt_subst_left}.
    \end{subproof}
    $(n+1,y)\in h$ by
        \cref{singleton_intro,union_intro_right}.
    $h(n+1)=y$ by
        \cref{funs_is_function,function_apply_intro}.
    Follows by \cref{subseteq_refl}.
\end{proof}

% from §50 helper lemma 98: the finite at-most bound passes to subsets
\begin{lemma}\label{finite_at_most_subset_50}
    Assume $A$ has at most $m$ elements.
    Assume $B\subseteq A$.
    Then $B$ has at most $m$ elements.
\end{lemma}
\begin{proof}
    Follows by
        \cref{finite_set_has_at_most_50,subseteq_finite_is_finite,subseteq_transitive}.
\end{proof}

% from §50 helper lemma 99: deleting a point lowers a positive finite at-most bound
\begin{lemma}\label{finite_at_most_remove_point_50}
    Assume $m\in\naturalsPlus$.
    Assume $D$ has at most $m+1$ elements.
    Assume $y\in D$.
    Let $B=D\setminus\{y\}$.
    Then $B$ has at most $m$ elements.
\end{lemma}
\begin{proof}
    $m+1\in\naturalsPlus$ by
        \cref{real_naturals_closed_succ}.
    $D$ is finite by \cref{finite_set_has_at_most_50}.
    $B\subseteq D$ and $B$ is finite by
        \cref{setminus_subseteq,subseteq_finite_is_finite}.
    There exists no $C$ such that
        $C\subseteq B$ and
        there exists $g$ such that
            $g$ is a bijection from $\initialSegment{m+1}$ to $C$ by
        \cref{setminus_elim_right,singleton_intro,not_in_subseteq,finite_bijection_append_new_point_50,real_naturals_closed_succ,union_iff,singleton_iff,subseteq,finite_set_has_at_most_50}.
    Follows by \cref{finite_set_has_at_most_50}.
\end{proof}

% from §50 helper lemma 100: an enumeration of an at-most set has bounded length
\begin{lemma}\label{finite_at_most_enumeration_length_50}
    Assume $N,n\in\naturalsPlus$.
    Assume $A$ has at most $N$ elements.
    Assume $g$ is a bijection from $\initialSegment{n}$ to $A$.
    Then $n\leq N$.
\end{lemma}
\begin{proof}
    Suppose not.
    $N,n\in\reals$ by
        \cref{real_naturals_subset_reals,subseteq}.
    $N<n$ by \cref{real_lt_trichotomy}.
    $N+1,n\in\reals$ by
        \cref{real_naturals_closed_succ,real_naturals_subset_reals,real_add_closed,real_one_in_naturals,subseteq}.
    $N+1\leq n$ by
        \cref{real_lt_trichotomy,real_lt_imp_leq,real_lt_subst_right,real_lt_irreflexive,real_naturals_leq_succ_elim,real_leq_split,real_lt_trans}.
    Let $L=\initialSegment{N+1}$.
    Let $I=\initialSegment{n}$.
    $L\subseteq I$ by
        \cref{initial_segment_naturalsplus,subseteq,real_naturals_closed_succ,real_naturals_subset_reals,real_leq_trans}.
    Let $q=\restrl{g}{L}$.
    $g$ is a function and $\dom{g}=I$ by
        \cref{bijection_is_function,bijections_dom}.
    $q$ is a function by
        \cref{restrl_of_function_is_function}.
    $\dom{q}=L$ by
        \cref{dom_of_restrl_eq_inter,inter_absorb_supseteq_left}.
    For all $i\in\dom{q}$ we have $q(i)\in A$ by
        \cref{real_lt_subst_left,subseteq,restrl_of_function_apply,bijections_to_funs,funs_type_apply}.
    $q\in\funs{L}{A}$ by \cref{funs_intro}.
    $q$ is injective by
        \cref{restrl_of_function_apply,bijections,injective_function,subseteq,real_lt_subst_left}.
    Let $C=\img{q}{L}$.
    $C\subseteq A$ by
        \cref{img_subseteq_ran,funs_ran,subseteq_transitive}.
    $q$ is a bijection from $L$ to $C$ by
        \cref{real_lt_subst_right,img_dom,real_lt_subst_left,img_of_function_intro,inter_intro,funs_intro,funs_surj_iff,bijections}.
    Contradiction by \cref{finite_set_has_at_most_50}.
\end{proof}

% from §50 helper definition 36: binary indicator of a subset
\begin{definition}\label{binary_subset_indicator_50}
    $f$ is a binary indicator of $B$ on $A$ iff
        $f\in\funs{A}{\{0,1\}}$ and
        $B=\preimg{f}{\{1\}}$.
\end{definition}

% from §50 helper lemma 101: the powerset of a finite set is finite
\begin{lemma}\label{finite_powerset_50}
    Assume $A$ is finite.
    Then $\pow{A}$ is finite.
\end{lemma}
\begin{proof}
    Let $Q=\{0,1\}$.
    $Q$ is finite by \cref{pair_is_finite}.
    Let $F=\funs{A}{Q}$.
    $F$ is finite by \cref{finite_function_space}.
    Let $R=\{w\in\pow{A}\times F\mid
        \text{$\snd{w}$ is a binary indicator of $\fst{w}$ on $A$}\}$.
    $R$ is a relation by
        \cref{relation,times_elem_is_tuple}.
    Show that for all $B\in\pow{A}$ there exists $f$ such that
        $B\mathrel{R}f$.
    \begin{subproof}
        Fix $B\in\pow{A}$.
        Let $f=\{(a,r)\mid a\in A,r\in Q\mid
            (a\in B\land r=1)\lor(a\notin B\land r=0)\}$.
        $f$ is a relation by \cref{relation,pair_eq_iff}.
        $f$ is right-unique by
            \cref{rightunique,pair_eq_iff}.
        $f$ is a function by \cref{relation,rightunique}.
        $\dom{f}=A$ by
            \cref{subseteq,dom_iff,pair_eq_iff,upair_intro_right,dom_intro,upair_intro_left,subseteq_antisymmetric}.
        For all $a\in\dom{f}$ we have $f(a)\in Q$ by
            \cref{function_apply_iff,pair_eq_iff}.
        $f\in\funs{A}{Q}$ by \cref{funs_intro}.
        $B=\preimg{f}{\{1\}}$ by
            \cref{powerset_elim,subseteq,upair_intro_right,pair_eq_iff,singleton_intro,preimg_iff,real_add_ident,real_zero_lt_one,real_lt_subst_right,real_lt_irreflexive,singleton_iff,real_lt_subst_left,subseteq_antisymmetric}.
        $f$ is a binary indicator of $B$ on $A$ by
            \cref{binary_subset_indicator_50}.
        $(B,f)\in\pow{A}\times F$ by
            \cref{times_tuple_intro}.
        $(B,f)\in R$ by
            \cref{fst_eq,snd_eq}.
        Follows.
    \end{subproof}
    Take $c$ such that $c$ is a function on $\pow{A}$ and
        for all $B\in\pow{A}$ we have
            $B\mathrel{R}c(B)$ by
        \cref{choice_function}.
    For all $B\in\dom{c}$ we have $c(B)\in F$ by
        \cref{choice_function,subseteq,times_tuple_elim}.
    $c\in\funs{\pow{A}}{F}$ by
        \cref{choice_function,funs_intro}.
    For all $B,C$ such that
        $B\in\dom{c}$ and $C\in\dom{c}$ and $c(B)=c(C)$
    we have $B=C$ by
        \cref{choice_function,subseteq,fst_eq,snd_eq,binary_subset_indicator_50,real_lt_subst_left,real_lt_subst_right}.
    Thus $c$ is injective by \cref{injective_function}.
    Let $D=\img{c}{\pow{A}}$.
    $D\subseteq F$ by
        \cref{img_subseteq_ran,funs_ran,subseteq_transitive}.
    $D$ is finite by
        \cref{subseteq_finite_is_finite}.
    For all $B\in\dom{c}$ we have $c(B)\in D$ by
        \cref{img_of_function_intro,inter_intro}.
    $c\in\funs{\pow{A}}{D}$ by
        \cref{choice_function,funs_intro}.
    $c\in\surj{\pow{A}}{D}$ by
        \cref{funs_surj_iff,img_dom}.
    $c$ is a bijection from $\pow{A}$ to $D$ by
        \cref{bijections}.
    $\converse{c}$ is a bijection from $D$ to $\pow{A}$ by
        \cref{bijection_converse_is_bijection}.
    Follows by \cref{finite_bijection}.
\end{proof}

% from §50 helper lemma 102: an at-most bound remains true after increasing the bound
\begin{lemma}\label{finite_at_most_successor_50}
    Assume $A$ has at most $m$ elements.
    Then $A$ has at most $m+1$ elements.
\end{lemma}
\begin{proof}
    $m\in\naturalsPlus$ and $A$ is finite by
        \cref{finite_set_has_at_most_50}.
    $m+1\in\naturalsPlus$ by
        \cref{real_naturals_closed_succ}.
    Show that there exists no $C$ such that
        $C\subseteq A$ and
        there exists $g$ such that
            $g$ is a bijection from
                $\initialSegment{(m+1)+1}$ to $C$.
    \begin{subproof}
        Suppose not.
        Take $C,g$ such that
            $C\subseteq A$ and
            $g$ is a bijection from
                $\initialSegment{(m+1)+1}$ to $C$.
        Let $L=\initialSegment{m+1}$.
        Let $I=\initialSegment{(m+1)+1}$.
        $L\subseteq I$ by
            \cref{initial_segment_subset_succ}.
        Let $q=\restrl{g}{L}$.
        $g$ is a function and $\dom{g}=I$ by
            \cref{bijection_is_function,bijections_dom}.
        $q$ is a function by
            \cref{restrl_of_function_is_function}.
        $\dom{q}=L$ by
            \cref{dom_of_restrl_eq_inter,inter_absorb_supseteq_left}.
        For all $i\in\dom{q}$ we have $q(i)\in A$ by
            \cref{real_lt_subst_left,subseteq,restrl_of_function_apply,bijections_to_funs,funs_type_apply}.
        $q\in\funs{L}{A}$ by \cref{funs_intro}.
        $q$ is injective by
            \cref{restrl_of_function_apply,bijections,injective_function,subseteq,real_lt_subst_left}.
        Let $D=\img{q}{L}$.
        $D\subseteq A$ by
            \cref{img_subseteq_ran,funs_ran,subseteq_transitive}.
        $q$ is a bijection from $L$ to $D$ by
            \cref{real_lt_subst_right,img_dom,real_lt_subst_left,img_of_function_intro,inter_intro,funs_intro,funs_surj_iff,bijections}.
        Contradiction by \cref{finite_set_has_at_most_50}.
    \end{subproof}
    Follows by \cref{finite_set_has_at_most_50}.
\end{proof}

% from §50 helper lemma 103: a singleton obeys every positive at-most bound
\begin{lemma}\label{singleton_at_most_positive_50}
    Assume $N\in\naturalsPlus$.
    Then $\{a\}$ has at most $N$ elements.
\end{lemma}
\begin{proof}
    Let $P=\{a,a\}$.
    $P$ is finite by \cref{pair_is_finite}.
    $a\in P$ by \cref{upair_intro_left}.
    $\{a\}\subseteq P$ by \cref{singleton_subset_intro}.
    $\{a\}$ is finite by
        \cref{subseteq_finite_is_finite}.
    There exists no $C$ such that
        $C\subseteq\{a\}$ and
        there exists $g$ such that
            $g$ is a bijection from $\initialSegment{N+1}$ to $C$ by
        \cref{real_naturals_closed_succ,real_one_in_naturals,real_one_leq_naturalsplus,initial_segment_naturalsplus,bijections_to_funs,funs_type_apply,subseteq,singleton_iff,real_lt_subst_left,real_lt_subst_right,bijection_is_function,bijections,injective_function,bijections_dom,real_naturals_subset_reals,real_naturals_lt_succ_local,real_leq_lt_trans,real_lt_irreflexive}.
    Follows by \cref{finite_set_has_at_most_50}.
\end{proof}

% from §50 helper lemma 104: a finite general-position set admits a nearby one-point extension
\begin{lemma}\label{general_position_one_point_extension_50}
    Assume $N\in\naturalsPlus$.
    Let $J=\initialSegment{N}$.
    Let $E=\realsPower{J}$.
    Assume $A$ is finite and $A$ is inhabited.
    Assume $A$ is in general position in dimension $N$.
    Assume $x\in E$.
    Assume $\delta\in\reals$ and $0<\delta$.
    Then there exists $y\in E$ such that
        $\apply{\squareMetric{N}}{(x,y)}<\delta$ and
        $y\notin A$ and
        $A\union\{y\}$ is in general position in dimension $N$.
\end{lemma}
\begin{proof}
    Let $S=\metricTopology{\squareMetric{N}}{E}$.
    Let $C=\{B\in\pow{A}\mid
        \text{$B$ is inhabited and $B$ has at most $N$ elements}\}$.
    $C\subseteq\pow{A}$ by \cref{subseteq}.
    $\pow{A}$ is finite by \cref{finite_powerset_50}.
    $C$ is finite by \cref{subseteq_finite_is_finite}.
    Take $a$ such that $a\in A$ by assumption.
    $\{a\}\subseteq A$ by \cref{singleton_subset_intro}.
    $\{a\}\in\pow{A}$ by \cref{powerset_intro}.
    $\{a\}$ is inhabited by
        \cref{singleton_intro}.
    $\{a\}$ has at most $N$ elements by
        \cref{singleton_at_most_positive_50}.
    $\{a\}\in C$ by \cref{subseteq}.
    $C$ is inhabited.
    Let $R=\{w\in C\times\pow{E}\mid
        \text{there exists $P$ such that
            $P$ is the plane determined by $\fst{w}$ in dimension $N$ and
            $P\subseteq\snd{w}$ and
            $\snd{w}$ is closed in $E$ with respect to $S$ and
            $\snd{w}$ has empty interior in $E$ with respect to $S$}\}$.
    $R$ is a relation by
        \cref{relation,subseteq,times_elem_is_tuple}.
    Show that for all $B\in C$ there exists $H$ such that
        $B\mathrel{R}H$.
    \begin{subproof}
        Fix $B\in C$.
        $B\in\pow{A}$ and $B$ is inhabited and
            $B$ has at most $N$ elements by \cref{subseteq}.
        $B\subseteq A$ by \cref{powerset_elim}.
        $B$ has at most $N+1$ elements by
            \cref{finite_at_most_successor_50}.
        $B$ is a geometrically independent set in dimension $N$ by
            \cref{general_position_50}.
        $B$ is finite by \cref{finite_set_has_at_most_50}.
        Take $n,g$ such that
            $n\in\naturalsPlus$ and
            $g$ is a bijection from $\initialSegment{n}$ to $B$ by
            \cref{finite_inhabited_initial_segment_enumeration}.
        $n\leq N$ by
            \cref{finite_at_most_enumeration_length_50}.
        Let $P=\{z\in E\mid
            \text{there exist $m,s,b$ such that
                $m\in\naturalsPlus$ and
                $s$ is a bijection from $\initialSegment{m}$ to $B$ and
                $b\in\funs{\initialSegment{m}}{\reals}$ and
                $\sumFiniteSequence{b}{1}$ and
                $z$ is an affine combination of $s$ with coefficients $b$
                    in dimension $N$}\}$.
        $P$ is the plane determined by $B$ in dimension $N$ by
            \cref{plane_determined_by_points_50}.
        Take $c,d$ such that
            $c\in E$ and
            $c\neq\zeroVector{N}$ and
            $d\in\reals$ and
            for all $z\in P$ there exists $r\in\reals$ such that
                $r$ is the dot product of $z$ and $c$ of length $N$ and
                $r+d=0$ by
            \cref{affine_plane_functional_50}.
        Let $H=\{z\in E\mid
            \text{there exists $r\in\reals$ such that
                $r$ is the dot product of $z$ and $c$ of length $N$ and
                $r+d=0$}\}$.
        $P\subseteq H$ by \cref{subseteq}.
        $H\subseteq E$ by \cref{subseteq}.
        $H\in\pow{E}$ by \cref{powerset_intro}.
        $H$ is closed in $E$ with respect to $S$ by
            \cref{affine_dot_level_set_closed_50}.
        $H$ has empty interior in $E$ with respect to $S$ by
            \cref{affine_dot_level_set_empty_interior_50}.
        $(B,H)\in C\times\pow{E}$ by
            \cref{times_tuple_intro}.
        $(B,H)\in R$ by
            \cref{fst_eq,snd_eq,subseteq}.
        Follows by \cref{relation}.
    \end{subproof}
    Take $h$ such that
        $h$ is a function on $C$ and
        for all $B\in C$ we have $B\mathrel{R}h(B)$ by
        \cref{choice_function}.
    Let $F=\img{h}{C}$.
    $F$ is finite by \cref{finite_image_function}.
    $F$ is inhabited by
        \cref{choice_function,function_apply_elim,img_iff,nonempty_inhabited}.
    $F\subseteq\pow{E}$ by
        \cref{subseteq,img_of_function_elim,inter_elim_right,choice_function,times_tuple_elim,snd_eq,real_lt_subst_left}.
    For all $H\in F$ we have
        $H$ is closed in $E$ with respect to $S$ and
        $H$ has empty interior in $E$ with respect to $S$ by
        \cref{img_of_function_elim,inter_elim_right,choice_function,subseteq,fst_eq,snd_eq,real_lt_subst_left}.
    Take $k,u$ such that
        $k\in\naturalsPlus$ and
        $u$ is a bijection from $\initialSegment{k}$ to $F$ by
        \cref{finite_inhabited_initial_segment_enumeration}.
    $\squareMetric{N}$ is a metric on $E$ by
        \cref{square_metric_is_metric}.
    $S$ is a topology on $E$ by
        \cref{metric_topology,metric_basis_is_basis,basis_topology_is_topology}.
    $E$ is complete with respect to $\squareMetric{N}$ by
        \cref{euclidean_space_square_metric_complete}.
    $E$ is a Baire space with respect to $S$ by
        \cref{complete_metric_space_baire}.
    Let $Z=\unions{F}$.
    $Z$ has empty interior in $E$ with respect to $S$ by
        \cref{finite_closed_empty_interior_union_50}.
    $Z\subseteq E$ by
        \cref{unions_iff,powerset_elim,subseteq}.
    $E\setminus Z$ is dense in $E$ with respect to $S$ by
        \cref{empty_interior_iff_dense_complement}.
    Let $U=\metricBall{\squareMetric{N}}{E}{x}{\delta}$.
    $U$ is a neighborhood of $x$ in $E$ with respect to $S$ by
        \cref{metric_ball_neighborhood_metric_topology}.
    $U$ is open in $E$ with respect to $S$ and $x\in U$ by
        \cref{neighborhood}.
    $x\in\closure{E\setminus Z}{E}{S}$ by
        \cref{dense_in,real_lt_subst_right}.
    $U$ intersects $E\setminus Z$ by
        \cref{closure_iff_intersects_open,setminus_subseteq}.
    Take $y$ such that $y\in U\inter(E\setminus Z)$ by
        \cref{intersects,nonempty_inhabited}.
    $y\in U$ and $y\in E\setminus Z$ by \cref{inter}.
    $y\in E$ and $y\notin Z$ by
        \cref{setminus_elim_left,setminus_elim_right}.
    $\apply{\squareMetric{N}}{(x,y)}<\delta$ by
        \cref{metric_ball}.
    $y\notin A$ by
        \cref{singleton_subset_intro,powerset_intro,singleton_intro,singleton_at_most_positive_50,subseteq,choice_function,fst_eq,snd_eq,singleton_geometrically_independent_50,generator_belongs_to_affine_plane_50,function_apply_elim,img_iff,unions_iff}.
    Show that $A\union\{y\}$ is in general position in dimension $N$.
    \begin{subproof}
        $A\subseteq E$ by \cref{general_position_50}.
        $A\union\{y\}\subseteq E$ by
            \cref{union_iff,singleton_iff,subseteq}.
        Show that for all $D$ such that
            $D\subseteq A\union\{y\}$ and
            $D$ has at most $N+1$ elements
        we have $D$ is a geometrically independent set in dimension $N$.
        \begin{subproof}
            Fix $D$.
            Assume $D\subseteq A\union\{y\}$ and
                $D$ has at most $N+1$ elements.
            \begin{byCase}
            \caseOf{$y\notin D$.}
                $D\subseteq A$ by
                    \cref{subseteq,union_iff,singleton_iff,real_lt_subst_left}.
                Follows by \cref{general_position_50}.
            \caseOf{$y\in D$.}
                Let $B=D\setminus\{y\}$.
                $B$ has at most $N$ elements by
                    \cref{finite_at_most_remove_point_50}.
                $B\subseteq A$ by
                    \cref{subseteq,setminus_elim_left,setminus_elim_right,union_iff}.
                \begin{byCase}
                \caseOf{$B=\emptyset$.}
                    $D=\{y\}$ by
                        \cref{subseteq,singleton_iff,setminus_intro,real_lt_subst_left,emptyset,singleton_subset_intro,subseteq_antisymmetric}.
                    $\{y\}$ is a geometrically independent set in dimension $N$ by
                        \cref{singleton_geometrically_independent_50}.
                    Follows by \cref{real_lt_subst_left}.
                \caseOf{$B\neq\emptyset$.}
                    $B$ is inhabited by \cref{nonempty_inhabited}.
                    $B\in\pow{A}$ by \cref{powerset_intro}.
                    $B\in C$ by \cref{subseteq}.
                    $B\mathrel{R}h(B)$ by \cref{choice_function}.
                    Take $P$ such that
                        $P$ is the plane determined by $B$ in dimension $N$ and
                        $P\subseteq h(B)$ and
                        $h(B)$ is closed in $E$ with respect to $S$ and
                        $h(B)$ has empty interior in $E$ with respect to $S$ by
                        \cref{subseteq,fst_eq,snd_eq}.
                    $h(B)\in F$ by
                        \cref{choice_function,function_apply_elim,img_iff}.
                    $y\notin h(B)$ by \cref{unions_iff}.
                    $y\notin P$ by \cref{not_in_subseteq}.
                    $y\in E\setminus P$ by \cref{setminus_intro}.
                    $B$ has at most $N+1$ elements by
                        \cref{finite_at_most_successor_50}.
                    $B$ is a geometrically independent set in dimension $N$ by
                        \cref{general_position_50}.
                    $B\union\{y\}$ is a geometrically independent set in dimension $N$ by
                        \cref{independent_set_extension_outside_plane_50}.
                    $D=B\union\{y\}$ by
                        \cref{subseteq,singleton_iff,setminus_intro,union_intro_right,union_intro_left,union_iff,setminus_elim_left,real_lt_subst_left,subseteq_antisymmetric}.
                    Follows by \cref{real_lt_subst_left}.
                \end{byCase}
            \end{byCase}
        \end{subproof}
        Follows by \cref{general_position_50}.
    \end{subproof}
    Therefore there exists $y_0\in E$ such that
        $\apply{\squareMetric{N}}{(x,y_0)}<\delta$ and
        $y_0\notin A$ and
        $A\union\{y_0\}$ is in general position in dimension $N$.
\end{proof}

% from §50 helper lemma 105: append one value to a finite function and compute its image
\begin{lemma}\label{finite_function_append_value_50}
    Assume $n\in\naturalsPlus$.
    Let $I=\initialSegment{n}$.
    Let $I_1=\initialSegment{n+1}$.
    Assume $g\in\funs{I}{E}$.
    Assume $y\in E$.
    Then there exists $h$ such that
        $h\in\funs{I_1}{E}$ and
        $\restrl{h}{I}=g$ and
        $h(n+1)=y$ and
        $\img{h}{I_1}=\img{g}{I}\union\{y\}$.
\end{lemma}
\begin{proof}
    Let $r=\{(n+1,y)\}$.
    Let $h=g\union r$.
    $g$ is a function and $\dom{g}=I$ by
        \cref{funs_elim}.
    $r$ is a function by
        \cref{singleton_elim,pair_eq_iff,relation,rightunique}.
    $\dom{r}=\{n+1\}$ by
        \cref{dom_iff,singleton_elim,pair_eq_iff,singleton_intro,subseteq,subseteq_antisymmetric}.
    $n+1\notin I$ by
        \cref{initial_segment_naturalsplus,real_naturals_subset_reals,subseteq,real_naturals_closed_succ,real_naturals_lt_succ_local,real_lt_trans,real_lt_irreflexive}.
    $\dom{g}\inter\dom{r}=\emptyset$ by
        \cref{inter,subseteq,emptyset,subseteq_antisymmetric,emptyset_subseteq,singleton_iff,real_lt_subst_left,real_lt_subst_right}.
    $h$ is a function by
        \cref{union_compatible_functions,emptyset,inter}.
    $\dom{h}=I\union\{n+1\}$ by
        \cref{dom_union,real_lt_subst_left,real_lt_subst_right}.
    $\dom{h}=I_1$ by
        \cref{union_iff,initial_segment_subset_succ,subseteq,singleton_iff,real_naturals_closed_succ,initial_segment_naturalsplus,initial_segment_succ_split,subseteq_antisymmetric}.
    For all $i\in\dom{h}$ we have $h(i)\in E$ by
        \cref{funs,function_apply_elim,union_iff,dom_intro,funs_type_apply,funs_is_function,function_apply_intro,singleton_elim,pair_eq_iff,real_lt_subst_left}.
    $h\in\funs{I_1}{E}$ by \cref{funs_intro}.
    $\restrl{h}{I}=g$ by
        \cref{initial_segment_subset_succ,real_lt_subst_left,restrl_of_function_is_function,dom_of_restrl_eq_inter,inter_absorb_supseteq_left,restrl_of_function_apply,funs_is_function,funs_elim,subseteq,function_apply_elim,union_iff,function_apply_intro,fun_subseteq,subseteq_refl,subseteq_antisymmetric}.
    $(n+1,y)\in h$ by
        \cref{singleton_intro,union_intro_right}.
    $h(n+1)=y$ by
        \cref{funs_is_function,function_apply_intro}.
    $\img{h}{I_1}=\img{g}{I}\union\{y\}$ by
        \cref{subseteq,img_of_function_elim,inter_elim_right,initial_segment_succ_split,funs_elim,function_apply_elim,union_iff,function_apply_intro,img_of_function_intro,inter_intro,real_lt_subst_left,real_lt_subst_right,union_intro_left,singleton_iff,union_intro_right,initial_segment_subset_succ,inter_elim_left,funs_is_function,img_iff,real_naturals_closed_succ,initial_segment_naturalsplus,subseteq_antisymmetric}.
    Therefore there exists $h_0$ such that
        $h_0\in\funs{I_1}{E}$ and
        $\restrl{h_0}{I}=g$ and
        $h_0(n+1)=y$ and
        $\img{h_0}{I_1}=\img{g}{I}\union\{y\}$.
\end{proof}

% from §50 helper lemma 106: a singleton of a euclidean space is in general position
\begin{lemma}\label{singleton_general_position_50}
    Assume $N\in\naturalsPlus$.
    Let $E=\realsPower{\initialSegment{N}}$.
    Assume $y\in E$.
    Then $\{y\}$ is in general position in dimension $N$.
\end{lemma}
\begin{proof}
    Follows by
        \cref{singleton_subset_intro,empty_geometrically_independent_50,real_lt_subst_left,nonempty_inhabited,subseteq,singleton_iff,subseteq_antisymmetric,singleton_geometrically_independent_50,general_position_50}.
\end{proof}

% from §50 helper definition 37: approximation property at a finite sequence length
\begin{definition}\label{finite_general_position_approximation_property_50}
    $k$ has the finite general-position approximation property iff
        $k\in\naturalsPlus$ and
        for all $N,x,\delta$ such that
            $N\in\naturalsPlus$ and
            $x\in\funs{\initialSegment{k}}
                {\realsPower{\initialSegment{N}}}$ and
            $\delta\in\reals$ and $0<\delta$
        there exists $y$ such that
            $y\in\funs{\initialSegment{k}}
                {\realsPower{\initialSegment{N}}}$ and
            $\img{y}{\initialSegment{k}}$ is in general position
                in dimension $N$ and
            for all $i\in\initialSegment{k}$ we have
                $\apply{\squareMetric{N}}{(x(i),y(i))}<\delta$.
\end{definition}

% from §50 helper lemma 107: the approximation property starts at length one
\begin{lemma}\label{finite_general_position_approximation_property_one_50}
    Then $1$ has the finite general-position approximation property.
\end{lemma}
\begin{proof}
    $1\in\naturalsPlus$ by \cref{real_one_in_naturals}.
    Show that for all $N,x,\delta$ such that
        $N\in\naturalsPlus$ and
        $x\in\funs{\initialSegment{1}}
            {\realsPower{\initialSegment{N}}}$ and
        $\delta\in\reals$ and $0<\delta$
    there exists $y$ such that
        $y\in\funs{\initialSegment{1}}
            {\realsPower{\initialSegment{N}}}$ and
        $\img{y}{\initialSegment{1}}$ is in general position
            in dimension $N$ and
        for all $i\in\initialSegment{1}$ we have
            $\apply{\squareMetric{N}}{(x(i),y(i))}<\delta$.
    \begin{subproof}
        Fix $N$.
        Fix $x$.
        Fix $\delta$.
        Assume $N\in\naturalsPlus$ and
            $x\in\funs{\initialSegment{1}}
                {\realsPower{\initialSegment{N}}}$ and
            $\delta\in\reals$ and $0<\delta$.
        Let $I=\initialSegment{1}$.
        Let $E=\realsPower{\initialSegment{N}}$.
        Let $y=x$.
        $y\in\funs{I}{E}$ by \cref{real_lt_subst_left}.
        $y$ is a function by \cref{funs_is_function}.
        $\dom{y}=I$ by \cref{funs_elim}.
        $1\in I$ by
            \cref{real_one_in_naturals,initial_segment_naturalsplus}.
        $y(1)\in E$ by \cref{funs_type_apply}.
        $\img{y}{I}=\{y(1)\}$ by
            \cref{subseteq,img_of_function_elim,inter_elim_right,initial_segment_one,singleton_iff,real_lt_subst_left,real_lt_subst_right,funs_is_function,function_apply_elim,img_iff,singleton_subset_intro,subseteq_antisymmetric}.
        $\{y(1)\}$ is in general position in dimension $N$ by
            \cref{singleton_general_position_50}.
        $\img{y}{I}$ is in general position in dimension $N$ by
            \cref{real_lt_subst_left}.
        $\squareMetric{N}$ is a metric on $E$ by
            \cref{square_metric_is_metric}.
        For all $i\in I$ we have
            $\apply{\squareMetric{N}}{(x(i),y(i))}<\delta$ by
            \cref{funs_type_apply,real_lt_subst_right,metric_zero_characterization,metric_on,real_lt_subst_left}.
        Therefore there exists $y_0$ such that
            $y_0\in\funs{\initialSegment{1}}
                {\realsPower{\initialSegment{N}}}$ and
            $\img{y_0}{\initialSegment{1}}$ is in general position
                in dimension $N$ and
            for all $i\in\initialSegment{1}$ we have
                $\apply{\squareMetric{N}}{(x(i),y_0(i))}<\delta$.
    \end{subproof}
    Follows by
        \cref{finite_general_position_approximation_property_50}.
\end{proof}

% from §50 helper lemma 108: the approximation property is closed under successors
\begin{lemma}\label{finite_general_position_approximation_property_successor_50}
    Assume $k$ has the finite general-position approximation property.
    Then $k+1$ has the finite general-position approximation property.
\end{lemma}
\begin{proof}
    $k\in\naturalsPlus$ by
        \cref{finite_general_position_approximation_property_50}.
    $k+1\in\naturalsPlus$ by
        \cref{real_naturals_closed_succ}.
    Show that for all $N,x,\delta$ such that
        $N\in\naturalsPlus$ and
        $x\in\funs{\initialSegment{k+1}}
            {\realsPower{\initialSegment{N}}}$ and
        $\delta\in\reals$ and $0<\delta$
    there exists $y$ such that
        $y\in\funs{\initialSegment{k+1}}
            {\realsPower{\initialSegment{N}}}$ and
        $\img{y}{\initialSegment{k+1}}$ is in general position
            in dimension $N$ and
        for all $i\in\initialSegment{k+1}$ we have
            $\apply{\squareMetric{N}}{(x(i),y(i))}<\delta$.
    \begin{subproof}
        Fix $N$.
        Fix $x$.
        Fix $\delta$.
        Assume $N\in\naturalsPlus$ and
            $x\in\funs{\initialSegment{k+1}}
                {\realsPower{\initialSegment{N}}}$ and
            $\delta\in\reals$ and $0<\delta$.
        Let $I=\initialSegment{k}$.
        Let $I_1=\initialSegment{k+1}$.
        Let $E=\realsPower{\initialSegment{N}}$.
        Let $x_0=\restrl{x}{I}$.
        $I\subseteq I_1$ by
            \cref{initial_segment_subset_succ}.
        $x$ is a function and $\dom{x}=I_1$ by
            \cref{funs_elim}.
        $x_0$ is a function by
            \cref{restrl_of_function_is_function}.
        $\dom{x_0}=I$ by
            \cref{dom_of_restrl_eq_inter,inter_absorb_supseteq_left}.
        For all $i\in\dom{x_0}$ we have $x_0(i)\in E$ by
            \cref{real_lt_subst_left,subseteq,restrl_of_function_apply,funs_type_apply}.
        $x_0\in\funs{I}{E}$ by \cref{funs_intro}.
        Take $g$ such that
            $g\in\funs{I}{E}$ and
            $\img{g}{I}$ is in general position in dimension $N$ and
            for all $i\in I$ we have
                $\apply{\squareMetric{N}}{(x_0(i),g(i))}<\delta$ by
            \cref{finite_general_position_approximation_property_50}.
        Let $A=\img{g}{I}$.
        $g$ is a function and $\dom{g}=I$ by
            \cref{funs_elim}.
        $A$ is finite by
            \cref{finite_sequence_image_finite_50}.
        $1\in I$ by
            \cref{real_one_in_naturals,real_one_leq_naturalsplus,initial_segment_naturalsplus}.
        $g(1)\in A$ by
            \cref{function_apply_elim,img_iff}.
        $A$ is inhabited.
        $x(k+1)\in E$ by
            \cref{real_naturals_closed_succ,initial_segment_naturalsplus,funs_type_apply}.
        Take $z\in E$ such that
            $\apply{\squareMetric{N}}{(x(k+1),z)}<\delta$ and
            $A\union\{z\}$ is in general position in dimension $N$ by
            \cref{general_position_one_point_extension_50}.
        Take $h$ such that
            $h\in\funs{I_1}{E}$ and
            $\restrl{h}{I}=g$ and
            $h(k+1)=z$ and
            $\img{h}{I_1}=A\union\{z\}$ by
            \cref{finite_function_append_value_50}.
        $h$ is a function by \cref{funs_is_function}.
        $\img{h}{I_1}$ is in general position in dimension $N$ by
            \cref{real_lt_subst_left}.
        For all $i\in I_1$ we have
            $\apply{\squareMetric{N}}{(x(i),h(i))}<\delta$ by
            \cref{initial_segment_succ_split,funs_elim,initial_segment_subset_succ,restrl_of_function_apply,real_lt_subst_left,subseteq,real_lt_subst_right}.
        Therefore there exists $y_0$ such that
            $y_0\in\funs{\initialSegment{k+1}}
                {\realsPower{\initialSegment{N}}}$ and
            $\img{y_0}{\initialSegment{k+1}}$ is in general position
                in dimension $N$ and
            for all $i\in\initialSegment{k+1}$ we have
                $\apply{\squareMetric{N}}{(x(i),y_0(i))}<\delta$.
    \end{subproof}
    Follows by
        \cref{finite_general_position_approximation_property_50}.
\end{proof}

% from §50 helper lemma 109: every positive finite length has the approximation property
\begin{lemma}\label{finite_general_position_approximation_property_all_50}
    Assume $k\in\naturalsPlus$.
    Then $k$ has the finite general-position approximation property.
\end{lemma}
\begin{proof}
    Let $Q=\{m\in\naturalsPlus\mid
        \text{$m$ has the finite general-position approximation property}\}$.
    Follows by
        \cref{subseteq,finite_general_position_approximation_property_one_50,real_one_in_naturals,real_naturals_closed_succ,finite_general_position_approximation_property_successor_50,real_naturals_induction}.
\end{proof}

% from §50 helper definition 38: injective approximation property at a finite sequence length
\begin{definition}\label{finite_injective_general_position_approximation_property_50}
    $k$ has the finite injective general-position approximation property iff
        $k\in\naturalsPlus$ and
        for all $N,x,\delta$ such that
            $N\in\naturalsPlus$ and
            $x\in\funs{\initialSegment{k}}
                {\realsPower{\initialSegment{N}}}$ and
            $\delta\in\reals$ and $0<\delta$
        there exists $y$ such that
            $y\in\funs{\initialSegment{k}}
                {\realsPower{\initialSegment{N}}}$ and
            $y$ is injective and
            $\img{y}{\initialSegment{k}}$ is in general position
                in dimension $N$ and
            for all $i\in\initialSegment{k}$ we have
                $\apply{\squareMetric{N}}{(x(i),y(i))}<\delta$.
\end{definition}

% from §50 helper lemma 110: injective approximation starts at length one
\begin{lemma}\label{finite_injective_general_position_approximation_property_one_50}
    Then $1$ has the finite injective general-position approximation property.
\end{lemma}
\begin{proof}
    Follows by
        \cref{real_one_in_naturals,finite_general_position_approximation_property_one_50,finite_general_position_approximation_property_50,funs_elim,injective_function,real_lt_subst_left,initial_segment_one,singleton_iff,real_lt_subst_right,finite_injective_general_position_approximation_property_50}.
\end{proof}

% from §50 helper lemma 111: injective approximation is closed under successors
\begin{lemma}\label{finite_injective_general_position_approximation_property_successor_50}
    Assume $k$ has the finite injective general-position approximation property.
    Then $k+1$ has the finite injective general-position approximation property.
\end{lemma}
\begin{proof}
    $k\in\naturalsPlus$ by
        \cref{finite_injective_general_position_approximation_property_50}.
    $k+1\in\naturalsPlus$ by
        \cref{real_naturals_closed_succ}.
    Show that for all $N,x,\delta$ such that
        $N\in\naturalsPlus$ and
        $x\in\funs{\initialSegment{k+1}}
            {\realsPower{\initialSegment{N}}}$ and
        $\delta\in\reals$ and $0<\delta$
    there exists $y$ such that
        $y\in\funs{\initialSegment{k+1}}
            {\realsPower{\initialSegment{N}}}$ and
        $y$ is injective and
        $\img{y}{\initialSegment{k+1}}$ is in general position
            in dimension $N$ and
        for all $i\in\initialSegment{k+1}$ we have
            $\apply{\squareMetric{N}}{(x(i),y(i))}<\delta$.
    \begin{subproof}
        Fix $N$.
        Fix $x$.
        Fix $\delta$.
        Assume $N\in\naturalsPlus$ and
            $x\in\funs{\initialSegment{k+1}}
                {\realsPower{\initialSegment{N}}}$ and
            $\delta\in\reals$ and $0<\delta$.
        Let $I=\initialSegment{k}$.
        Let $I_1=\initialSegment{k+1}$.
        Let $E=\realsPower{\initialSegment{N}}$.
        Let $x_0=\restrl{x}{I}$.
        $I\subseteq I_1$ by
            \cref{initial_segment_subset_succ}.
        $x$ is a function and $\dom{x}=I_1$ by
            \cref{funs_elim}.
        $x_0$ is a function by
            \cref{restrl_of_function_is_function}.
        $\dom{x_0}=I$ by
            \cref{dom_of_restrl_eq_inter,inter_absorb_supseteq_left}.
        For all $i\in\dom{x_0}$ we have $x_0(i)\in E$ by
            \cref{real_lt_subst_left,subseteq,restrl_of_function_apply,funs_type_apply}.
        $x_0\in\funs{I}{E}$ by \cref{funs_intro}.
        Take $g$ such that
            $g\in\funs{I}{E}$ and
            $g$ is injective and
            $\img{g}{I}$ is in general position in dimension $N$ and
            for all $i\in I$ we have
                $\apply{\squareMetric{N}}{(x_0(i),g(i))}<\delta$ by
            \cref{finite_injective_general_position_approximation_property_50}.
        Let $A=\img{g}{I}$.
        $g$ is a function and $\dom{g}=I$ by
            \cref{funs_elim}.
        $A$ is finite by
            \cref{finite_sequence_image_finite_50}.
        $1\in I$ by
            \cref{real_one_in_naturals,real_one_leq_naturalsplus,initial_segment_naturalsplus}.
        $g(1)\in A$ by
            \cref{function_apply_elim,img_iff}.
        $A$ is inhabited.
        $x(k+1)\in E$ by
            \cref{real_naturals_closed_succ,initial_segment_naturalsplus,funs_type_apply}.
        Take $z\in E$ such that
            $\apply{\squareMetric{N}}{(x(k+1),z)}<\delta$ and
            $z\notin A$ and
            $A\union\{z\}$ is in general position in dimension $N$ by
            \cref{general_position_one_point_extension_50}.
        $A=\img{g}{\dom{g}}$ by \cref{real_lt_subst_right}.
        $A=\ran{g}$ by \cref{img_dom,real_lt_subst_left}.
        For all $i\in\dom{g}$ we have $g(i)\in A$ by
            \cref{real_lt_subst_left,img_of_function_intro,inter_intro}.
        $g\in\funs{I}{A}$ by \cref{funs_intro}.
        $g\in\surj{I}{A}$ by \cref{funs_surj_iff}.
        $g$ is a bijection from $I$ to $A$ by
            \cref{bijections}.
        Take $h$ such that
            $h$ is a bijection from $I_1$ to $A\union\{z\}$ and
            $\restrl{h}{I}=g$ and
            $h(k+1)=z$ by
            \cref{finite_bijection_append_new_point_50}.
        $h\in\funs{I_1}{A\union\{z\}}$ by
            \cref{bijections_to_funs}.
        $A\union\{z\}\subseteq E$ by
            \cref{general_position_50}.
        $h$ is a function and $\dom{h}=I_1$ by
            \cref{funs_elim}.
        For all $i\in\dom{h}$ we have $h(i)\in E$ by
            \cref{funs_type_apply,subseteq}.
        $h\in\funs{I_1}{E}$ by \cref{funs_intro}.
        $h$ is injective by \cref{bijections}.
        $h\in\surj{I_1}{A\union\{z\}}$ by \cref{bijections}.
        $\img{h}{I_1}=A\union\{z\}$ by
            \cref{funs_surj_iff,img_dom,real_lt_subst_left,real_lt_subst_right}.
        $\img{h}{I_1}$ is in general position in dimension $N$ by
            \cref{real_lt_subst_left}.
        For all $i\in I_1$ we have
            $\apply{\squareMetric{N}}{(x(i),h(i))}<\delta$ by
            \cref{initial_segment_succ_split,funs_elim,initial_segment_subset_succ,restrl_of_function_apply,real_lt_subst_left,subseteq,real_lt_subst_right}.
        Therefore there exists $y_0$ such that
            $y_0\in\funs{\initialSegment{k+1}}
                {\realsPower{\initialSegment{N}}}$ and
            $y_0$ is injective and
            $\img{y_0}{\initialSegment{k+1}}$ is in general position
                in dimension $N$ and
            for all $i\in\initialSegment{k+1}$ we have
                $\apply{\squareMetric{N}}{(x(i),y_0(i))}<\delta$.
    \end{subproof}
    Follows by
        \cref{finite_injective_general_position_approximation_property_50}.
\end{proof}

% from §50 helper lemma 112: every positive length has injective approximation
\begin{lemma}\label{finite_injective_general_position_approximation_property_all_50}
    Assume $k\in\naturalsPlus$.
    Then $k$ has the finite injective general-position approximation property.
\end{lemma}
\begin{proof}
    Let $Q=\{m\in\naturalsPlus\mid
        \text{$m$ has the finite injective general-position approximation property}\}$.
    Follows by
        \cref{subseteq,finite_injective_general_position_approximation_property_one_50,real_one_in_naturals,real_naturals_closed_succ,finite_injective_general_position_approximation_property_successor_50,real_naturals_induction}.
\end{proof}

% from §50 helper lemma 113: finite euclidean sequences admit injective general-position approximations
\begin{lemma}\label{injective_general_position_approximation_50}
    Assume $N,n\in\naturalsPlus$.
    Let $E=\realsPower{\initialSegment{N}}$.
    Assume $x\in\funs{\initialSegment{n}}{E}$.
    Assume $\delta\in\reals$ and $0<\delta$.
    Then there exists $y$ such that
        $y\in\funs{\initialSegment{n}}{E}$ and
        $y$ is injective and
        $\img{y}{\initialSegment{n}}$ is in general position in dimension $N$ and
        for all $i\in\initialSegment{n}$ we have
            $\apply{\squareMetric{N}}{(x(i),y(i))}<\delta$.
\end{lemma}
\begin{proof}
    Follows by
        \cref{finite_injective_general_position_approximation_property_all_50,finite_injective_general_position_approximation_property_50}.
\end{proof}

% from §50 helper lemma 35: finite euclidean sequences admit general-position approximations
\begin{lemma}\label{general_position_approximation_principle_50}
    Assume $N\in\naturalsPlus$.
    Assume $n\in\naturalsPlus$.
    Let $E = \realsPower{\initialSegment{N}}$.
    Suppose $x\in\funs{\initialSegment{n}}{E}$.
    Suppose $\delta\in\reals$ and $0 < \delta$.
    Then there exists $y$ such that
        $y\in\funs{\initialSegment{n}}{E}$ and
        $\img{y}{\initialSegment{n}}$ is in general position in dimension $N$ and
        for all $i\in\initialSegment{n}$ we have
            $\apply{\squareMetric{N}}{(\apply{x}{i},\apply{y}{i})} < \delta$.
\end{lemma}
\begin{proof}
    $n$ has the finite general-position approximation property by
        \cref{finite_general_position_approximation_property_all_50}.
    Follows by
        \cref{finite_general_position_approximation_property_50}.
\end{proof}

% from §50 Item 10: approximation by points in general position
\begin{lemma}\label{general_position_approximation_50}
    Assume $N\in\naturalsPlus$.
    Assume $n\in\naturalsPlus$.
    Let $E = \realsPower{\initialSegment{N}}$.
    Suppose $x\in\funs{\initialSegment{n}}{E}$.
    Suppose $\delta\in\reals$ and $0 < \delta$.
    Then there exists $y$ such that
        $y\in\funs{\initialSegment{n}}{E}$ and
        $\img{y}{\initialSegment{n}}$ is in general position in dimension $N$ and
        for all $i\in\initialSegment{n}$ we have $\apply{\squareMetric{N}}{(\apply{x}{i},\apply{y}{i})} < \delta$.
\end{lemma}
\begin{proof}
    Follows by \cref{general_position_approximation_principle_50}.
\end{proof}

% from §50 helper lemma 46: a positive continuous function on a compact space has a positive lower bound
\begin{lemma}\label{compact_positive_continuous_lower_bound_50}
    Assume $T$ is a topology on $X$.
    Assume $X$ is compact with respect to $T$.
    Assume $X\neq\emptyset$.
    Assume $h$ is continuous from $X$ to $\reals$ with respect to $T$ and $\standardTopologyR$.
    Suppose for all $x\in X$ we have $0<h(x)$.
    Then there exists $r\in\reals$ such that
        $0<r$ and
        for all $x\in X$ we have $r\leq h(x)$.
\end{lemma}
\begin{proof}
    $\realOrderRel$ is a simple order relation on $\reals$ and
        $\reals$ has more than one element
        by \cref{reals_linear_continuum,linear_continuum}.
    $\realOrderTopology$ is the order topology on $\reals$ with respect to $\realOrderRel$
        by \cref{real_order_topology,order_topology}.
    $\standardTopologyR$ is the order topology on $\reals$ with respect to $\realOrderRel$
        by \cref{real_order_topology_eq_standard,real_lt_subst_left}.
    Take $c,d\in X$ such that
        $h(c)$ is a smallest element of $\img{h}{X}$ with respect to $\realOrderRel$ and
        $h(d)$ is a largest element of $\img{h}{X}$ with respect to $\realOrderRel$
        by \cref{extreme_value_theorem}.
    Let $r=h(c)$.
    $r\in\reals$ and $0<r$ by \cref{continuous,funs_type_apply,subseteq}.
    For all $x\in X$ we have $r\leq h(x)$ by
        \cref{continuous,funs_is_function,funs,inter_intro,img_of_function_intro,smallest_element,real_order_relation_elim,real_leq_split}.
    Follows by \cref{subseteq}.
\end{proof}

% from §50 helper lemma 47: pairs separated by a fixed metric scale form a compact subspace
\begin{lemma}\label{metric_large_pair_set_compact_50}
    Assume $d$ is a metric on $X$.
    Let $T=\metricTopology{d}{X}$.
    Assume $X$ is compact with respect to $T$.
    Assume $\epsilon\in\reals$.
    Let $P=\productTopology{T}{T}{X}{X}$.
    Let $A=\{p\in X\times X\mid \epsilon\leq d(p)\}$.
    Then $A$ is compact with respect to $\subspaceTopology{P}{A}$.
\end{lemma}
\begin{proof}
    $T$ is a topology on $X$ by
        \cref{metric_topology,metric_basis_is_basis,basis_for_topology,basis_topology_is_topology}.
    $P$ is a topology on $X\times X$ by
        \cref{product_basis_is_basis,basis_topology_is_topology,product_topology}.
    $X\times X$ is compact with respect to $P$ by \cref{product_compact}.
    Let $p_1=\piOne{X}{X}$.
    Let $p_2=\piTwo{X}{X}$.
    $p_1$ is continuous from $X\times X$ to $X$ with respect to $P$ and $T$
        by \cref{projection_one_continuous}.
    $p_2$ is continuous from $X\times X$ to $X$ with respect to $P$ and $T$
        by \cref{projection_two_continuous}.
    Let $h=\{(p,d((p_1(p),p_2(p))))\mid p\in X\times X\}$.
    $h$ is continuous from $X\times X$ to $\reals$ with respect to $P$ and $\standardTopologyR$
        by \cref{continuous_maps_pointwise_distance}.
    $h\in\funs{X\times X}{\reals}$ by \cref{continuous}.
    $\standardTopologyR$ is a topology on $\reals$ by
        \cref{standard_basis_is_basis,basis_topology_is_topology,standard_topology_reals}.
    Let $C=\closedrayright{\realOrderRel}{\reals}{\epsilon}$.
    $\realOrderRel$ is a simple order relation on $\reals$ and
        $\reals$ has more than one element
        by \cref{reals_linear_continuum,linear_continuum}.
    $\realOrderTopology$ is the order topology on $\reals$ with respect to $\realOrderRel$
        by \cref{real_order_topology,order_topology}.
    $C$ is closed in $\reals$ with respect to $\realOrderTopology$
        by \cref{closedrayright_closed_in_order_topology}.
    $C$ is closed in $\reals$ with respect to $\standardTopologyR$
        by \cref{real_order_topology_eq_standard,real_lt_subst_right}.
    For all $B$ such that $B\subseteq X\times X$ we have
        $\img{h}{\closure{B}{X\times X}{P}}
            \subseteq\closure{\img{h}{B}}{\reals}{\standardTopologyR}$
        by \cref{continuous_closure_image_subset}.
    $\preimg{h}{C}$ is closed in $X\times X$ with respect to $P$
        by \cref{closure_image_subset_implies_preimg_closed}.
    Show that $\preimg{h}{C}=A$.
    \begin{subproof}
        Show that $\preimg{h}{C}\subseteq A$.
        \begin{subproof}
            It suffices to show that for all $p\in\preimg{h}{C}$ we have $p\in A$
                by \cref{subseteq}.
            Fix $p\in\preimg{h}{C}$.
            $p\in X\times X$ by
                \cref{preimg_subseteq_dom,continuous,funs,subseteq}.
            Take $x,y$ such that $p=(x,y)$ and $x,y\in X$
                by \cref{times_elem_is_tuple}.
            $p_1(p)=x$ and $p_2(p)=y$ by
                \cref{projection_first,projection_second,continuous,funs_is_function,function_apply_intro,pair_eq_iff}.
            $h(p)=d((x,y))$ by
                \cref{continuous,funs_is_function,function_apply_intro,real_lt_subst_left,real_lt_subst_right}.
            Take $z\in C$ such that $(p,z)\in h$ by \cref{preimg_iff}.
            $h(p)=z$ by \cref{continuous,funs_is_function,function_apply_intro}.
            $h(p)\in C$ by \cref{real_lt_subst_left}.
            $\epsilon<h(p)$ or $h(p)=\epsilon$ by
                \cref{closedray_right,real_order_relation_elim}.
            $\epsilon\leq d(p)$ by
                \cref{real_leq_split,real_lt_subst_left,real_lt_subst_right,pair_eq_iff}.
            Follows by \cref{subseteq}.
        \end{subproof}
        Show that $A\subseteq\preimg{h}{C}$.
        \begin{subproof}
            It suffices to show that for all $p\in A$ we have $p\in\preimg{h}{C}$
                by \cref{subseteq}.
            Fix $p\in A$.
            $p\in X\times X$ and $\epsilon\leq d(p)$ by \cref{subseteq}.
            Take $x,y$ such that $p=(x,y)$ and $x,y\in X$
                by \cref{times_elem_is_tuple}.
            $p_1(p)=x$ and $p_2(p)=y$ by
                \cref{projection_first,projection_second,continuous,funs_is_function,function_apply_intro,pair_eq_iff}.
            $h(p)=d((x,y))$ by
                \cref{continuous,funs_is_function,function_apply_intro,real_lt_subst_left,real_lt_subst_right}.
            $d(p)=d((x,y))$ by \cref{pair_eq_iff}.
            $d(p)\in\reals$ and $h(p)\in\reals$ by
                \cref{metric_on,funs_type_apply,continuous}.
            $\epsilon<d(p)$ or $\epsilon=d(p)$ by \cref{real_leq_split}.
            $\epsilon<h(p)$ or $\epsilon=h(p)$ by
                \cref{real_lt_subst_right,real_lt_subst_left}.
            $h(p)\in C$ by
                \cref{closedray_right,real_order_relation_intro}.
            Follows by \cref{preimg_iff}.
        \end{subproof}
        Follows by \cref{subseteq_antisymmetric}.
    \end{subproof}
    $A$ is closed in $X\times X$ with respect to $P$.
    $A\subseteq X\times X$ by \cref{closed_in}.
    Follows by \cref{closed_subspace_compact}.
\end{proof}

% from §50 helper lemma 48: uniform bounds below one control pointwise distances in any metric target
\begin{lemma}\label{uniform_metric_pointwise_lt_below_one_50}
    Assume $e$ is a metric on $Y$.
    Assume $\rho$ is the uniform metric on functions from $X$ to $Y$ with respect to $e$.
    Assume $f\in\funs{X}{Y}$.
    Assume $g\in\funs{X}{Y}$.
    Assume $x\in X$.
    Assume $\delta\in\reals$.
    Suppose $\rho((f,g))<\delta$.
    Suppose $\delta\leq 1$.
    Then $e((f(x),g(x)))<\delta$.
\end{lemma}
\begin{proof}
    Let $b=\boundedMetric{e}$.
    $b((f(x),g(x)))\leq\rho((f,g))$
        by \cref{uniform_function_metric_coordinate_leq}.
    $\rho((f,g))\in\reals$ by
        \cref{uniform_metric_on_function_space,metric_on,funs_type_apply,times_tuple_intro}.
    $f(x),g(x)\in Y$ by \cref{funs_type_apply}.
    $(f(x),g(x))\in Y\times Y$ by \cref{times_tuple_intro}.
    $b((f(x),g(x)))\in\reals$ by
        \cref{bounded_metric_is_metric,metric_on,funs_type_apply,times_tuple_intro}.
    $b((f(x),g(x)))<\delta$ by \cref{real_leq_lt_trans}.
    $b((f(x),g(x)))<1$ by
        \cref{real_one_in_naturals,real_naturals_subset_reals,subseteq,real_lt_leq_trans_local}.
    $b((f(x),g(x)))=e((f(x),g(x)))$ by
        \cref{bounded_metric_apply_value_lt_one,bounded_metric_is_metric}.
    Follows by \cref{real_lt_subst_left}.
\end{proof}

% from §50 helper lemma 49: fixed-scale separating classes are open in the uniform topology
\begin{lemma}\label{metric_scale_separating_class_open_50}
    Assume $d$ is a metric on $X$.
    Let $T=\metricTopology{d}{X}$.
    Assume $X$ is compact with respect to $T$.
    Assume $e$ is a metric on $Y$.
    Let $S=\metricTopology{e}{Y}$.
    Assume $C$ is the set of continuous functions from $X$ to $Y$ with respect to $T$ and $S$.
    Assume $\rho$ is the uniform metric on functions from $X$ to $Y$ with respect to $e$.
    Let $q=\restrl{\rho}{C\times C}$.
    Assume $q$ is a metric on $C$.
    Assume $\epsilon\in\reals$.
    Assume $U$ is the separating class at scale $\epsilon$ on $X$ to $Y$
        with respect to $T$ and $S$ and $d$.
    Then $U$ is open in $C$ with respect to $\metricTopology{q}{C}$.
\end{lemma}
\begin{proof}
    $T$ is a topology on $X$ by
        \cref{metric_topology,metric_basis_is_basis,basis_for_topology,basis_topology_is_topology}.
    $S$ is a topology on $Y$ by
        \cref{metric_topology,metric_basis_is_basis,basis_for_topology,basis_topology_is_topology}.
    $C\subseteq\funs{X}{Y}$ by \cref{continuous_function_set,subseteq}.
    $\rho$ is a metric on $\funs{X}{Y}$ by
        \cref{uniform_metric_on_function_space}.
    $U\subseteq C$ by
        \cref{metric_scale_separating_class_50,continuous_function_set,subseteq}.
    Let $P=\productTopology{T}{T}{X}{X}$.
    $P$ is a topology on $X\times X$ by
        \cref{product_basis_is_basis,basis_topology_is_topology,product_topology}.
    Let $A=\{p\in X\times X\mid \epsilon\leq d(p)\}$.
    $A$ is compact with respect to $\subspaceTopology{P}{A}$
        by \cref{metric_large_pair_set_compact_50}.
    Show that for all $f\in U$ there exists $\delta\in\reals$ such that
        $0<\delta$ and $\metricBall{q}{C}{f}{\delta}\subseteq U$.
    \begin{subproof}
        Fix $f\in U$.
        $f\in C$ and $f\in\funs{X}{Y}$ by \cref{subseteq}.
        \begin{byCase}
        \caseOf{$A=\emptyset$.}
            $1\in\reals$ and $0<1$ by
                \cref{real_one_in_naturals,real_naturals_subset_reals,subseteq,real_zero_lt_one}.
            $\metricBall{q}{C}{f}{1}\subseteq U$ by
                \cref{subseteq,metric_ball,continuous_function_set,times_tuple_intro,emptyset,metric_scale_separating_class_50,continuous,funs}.
            Follows by \cref{subseteq}.
        \caseOf{$A\neq\emptyset$.}
            Let $p_1=\piOne{X}{X}$.
            Let $p_2=\piTwo{X}{X}$.
            Let $f_1=f\circ p_1$.
            Let $f_2=f\circ p_2$.
            $p_1$ is continuous from $X\times X$ to $X$ with respect to $P$ and $T$
                by \cref{projection_one_continuous}.
            $p_2$ is continuous from $X\times X$ to $X$ with respect to $P$ and $T$
                by \cref{projection_two_continuous}.
            $f$ is continuous from $X$ to $Y$ with respect to $T$ and $S$
                by \cref{metric_scale_separating_class_50,subseteq}.
            $f_1$ is continuous from $X\times X$ to $Y$ with respect to $P$ and $S$
                by \cref{continuous_composite}.
            $f_2$ is continuous from $X\times X$ to $Y$ with respect to $P$ and $S$
                by \cref{continuous_composite}.
            $f_1$ is continuous from $X\times X$ to $Y$ with respect to $P$ and
                $\metricTopology{e}{Y}$ by \cref{real_lt_subst_right}.
            $f_2$ is continuous from $X\times X$ to $Y$ with respect to $P$ and
                $\metricTopology{e}{Y}$ by \cref{real_lt_subst_right}.
            Let $h=\{(p,e((f_1(p),f_2(p))))\mid p\in X\times X\}$.
            $h$ is continuous from $X\times X$ to $\reals$ with respect to $P$ and
                $\standardTopologyR$ by
                \cref{continuous_maps_pointwise_distance}.
            Let $T_A=\subspaceTopology{P}{A}$.
            Let $k=\restrl{h}{A}$.
            $A\subseteq X\times X$ by \cref{subseteq}.
            $A$ is a subspace of $X\times X$ with respect to $P$ by
                \cref{subspace_of}.
            $T_A$ is a topology on $A$ by
                \cref{subspace_topology_is_topology}.
            $k$ is continuous from $A$ to $\reals$ with respect to $T_A$ and
                $\standardTopologyR$ by \cref{continuous_restrict_domain}.
            Show that for all $p\in A$ we have $0<k(p)$.
            \begin{subproof}
                Fix $p\in A$.
                $p\in X\times X$ and $\epsilon\leq d(p)$ by \cref{subseteq}.
                Take $x,y$ such that $p=(x,y)$ and $x,y\in X$
                    by \cref{times_elem_is_tuple}.
                $p_1(p)=x$ and $p_2(p)=y$ by
                    \cref{projection_first,projection_second,continuous,funs_is_function,function_apply_intro,pair_eq_iff}.
                $f_1(p)=f(x)$ and $f_2(p)=f(y)$ by
                    \cref{circ_apply,continuous,funs,funs_is_function,funs_ran,subseteq,pair_eq_iff}.
                $h(p)=e((f(x),f(y)))$ by
                    \cref{continuous,funs_is_function,function_apply_intro,real_lt_subst_left,real_lt_subst_right}.
                $k(p)=h(p)$ by
                    \cref{continuous,funs_is_function,funs,restrl_of_function_apply,subseteq}.
                $\epsilon\leq d((x,y))$ by \cref{pair_eq_iff}.
                $f(x)\neq f(y)$ by
                    \cref{metric_scale_separating_class_50,subseteq}.
                $f(x),f(y)\in Y$ by \cref{continuous,funs_type_apply}.
                $e((f(x),f(y)))\in\reals$ and
                    $0\leq e((f(x),f(y)))$ by
                    \cref{metric_on,funs_type_apply,times_tuple_intro,metric_nonnegative}.
                $e((f(x),f(y)))\neq 0$ by
                    \cref{metric_on,metric_zero_characterization}.
                $0<e((f(x),f(y)))$ by \cref{real_leq_split}.
                Follows by \cref{real_lt_subst_left,real_lt_subst_right}.
            \end{subproof}
            Take $r\in\reals$ such that $0<r$ and
                for all $p\in A$ we have $r\leq k(p)$
                by \cref{compact_positive_continuous_lower_bound_50}.
            Let $a=r/(1+1)$.
            $a\in\reals$ and $0<a$ and $a+a=r$ by
                \cref{real_half_of_positive}.
            $a<r$ by
                \cref{real_add_ident,real_lt_add,real_lt_subst_left,real_lt_subst_right}.
            Let $b=a/(1+1)$.
            $b\in\reals$ and $0<b$ and $b+b=a$ by
                \cref{real_half_of_positive}.
            $1\in\reals$ and $0<1$ by
                \cref{real_one_in_naturals,real_naturals_subset_reals,subseteq,real_zero_lt_one}.
            Take $\delta\in\reals$ such that
                $0<\delta$ and $\delta\leq b$ and $\delta\leq 1$
                by \cref{positive_radius_below_two_bounds_49}.
            $\delta+\delta\leq b+b$ by
                \cref{real_leq_add_two_49}.
            $\delta+\delta,b+b\in\reals$ by \cref{real_add_closed}.
            $\delta+\delta\leq a$ by \cref{real_lt_subst_right}.
            $\delta+\delta<r$ by
                \cref{real_leq_lt_trans}.
            Show that $\metricBall{q}{C}{f}{\delta}\subseteq U$.
            \begin{subproof}
                It suffices to show that for all $g\in\metricBall{q}{C}{f}{\delta}$
                    we have $g\in U$ by \cref{subseteq}.
                Fix $g\in\metricBall{q}{C}{f}{\delta}$.
                $g\in C$ and $q((f,g))<\delta$ by \cref{metric_ball}.
                $g\in\funs{X}{Y}$ and
                    $g$ is continuous from $X$ to $Y$ with respect to $T$ and $S$
                    by \cref{continuous_function_set,subseteq}.
                $q((f,g))=\rho((f,g))$ by
                    \cref{continuous_function_set,subseteq,uniform_metric_on_function_space,metric_on,funs_elim,funs_is_function,times_subseteq_left,times_subseteq_right,subseteq_transitive,pair_eq_iff,times_tuple_intro,restrl_of_function_apply}.
                $\rho((f,g))<\delta$ by \cref{real_lt_subst_left}.
                Show that for all $x\in X$ we have for all $y\in X$ if $\epsilon\leq d((x,y))$
                    then $g(x)\neq g(y)$.
                \begin{subproof}
                    Fix $x\in X$.
                    Fix $y\in X$.
                    Assume $\epsilon\leq d((x,y))$.
                    Let $p=(x,y)$.
                    $p\in A$ by \cref{times_tuple_intro,subseteq}.
                    $p_1(p)=x$ and $p_2(p)=y$ by
                        \cref{projection_first,projection_second,continuous,funs_is_function,function_apply_intro,fst_eq,snd_eq}.
                    $f_1(p)=f(x)$ and $f_2(p)=f(y)$ by
                        \cref{circ_apply,continuous,funs,funs_is_function,funs_ran,subseteq,pair_eq_iff}.
                    $h(p)=e((f(x),f(y)))$ by
                        \cref{continuous,funs_is_function,function_apply_intro,real_lt_subst_left,real_lt_subst_right}.
                    $k(p)=h(p)$ by
                        \cref{continuous,funs_is_function,funs,restrl_of_function_apply,subseteq}.
                    $r\leq k(p)$ by \cref{subseteq}.
                    $e((f(x),g(x)))<\delta$ and
                        $e((f(y),g(y)))<\delta$ by
                        \cref{uniform_metric_pointwise_lt_below_one_50}.
                    Suppose not.
                    $g(x)=g(y)$.
                    $f(x),f(y),g(x),g(y)\in Y$ by
                        \cref{funs_type_apply}.
                    $e((g(x),f(y)))=e((g(y),f(y)))$ by
                        \cref{real_lt_subst_left}.
                    $e((g(y),f(y)))=e((f(y),g(y)))$ by
                        \cref{metric_on,metric_symmetric}.
                    $e((g(x),f(y)))<\delta$ by
                        \cref{real_lt_subst_left,real_lt_subst_right}.
                    $e((f(x),f(y)))\leq
                        e((f(x),g(x)))+e((g(x),f(y)))$ by
                        \cref{metric_on,metric_triangle_inequality}.
                    $e((f(x),g(x)))+e((g(x),f(y)))<\delta+\delta$ by
                        \cref{metric_on,funs_type_apply,times_tuple_intro,real_lt_add_two}.
                    $e((f(x),f(y)))\in\reals$ and
                        $e((f(x),g(x)))+e((g(x),f(y)))\in\reals$ and
                        $\delta+\delta\in\reals$ by
                        \cref{metric_on,funs_type_apply,times_tuple_intro,real_add_closed}.
                    $e((f(x),f(y)))<\delta+\delta$ by
                        \cref{real_leq_lt_trans}.
                    $e((f(x),f(y)))<r$ by \cref{real_lt_trans}.
                    $r\leq e((f(x),f(y)))$ by
                        \cref{real_lt_subst_left,real_lt_subst_right}.
                    Contradiction by
                        \cref{real_leq_split,real_lt_trans,real_lt_irreflexive}.
                \end{subproof}
                Follows by
                    \cref{metric_scale_separating_class_50,continuous,funs}.
            \end{subproof}
            Follows by \cref{subseteq}.
        \end{byCase}
    \end{subproof}
    Follows by \cref{metric_open_iff}.
\end{proof}

% from §50 helper lemma 50: pointwise control bounds the uniform metric
\begin{lemma}\label{uniform_metric_leq_from_pointwise_bound_50}
    Assume $e$ is a metric on $Y$.
    Assume $\rho$ is the uniform metric on functions from $X$ to $Y$ with respect to $e$.
    Assume $f,g\in\funs{X}{Y}$.
    Assume $\delta\in\reals$ and $\delta\leq 1$.
    Assume for all $x\in X$ we have $e((f(x),g(x)))\leq\delta$.
    Then $\rho((f,g))\leq\delta$.
\end{lemma}
\begin{proof}
    Follows by \cref{uniform_metric_on_function_space,subseteq,real_lt_subst_left,real_leq_trans,metric_on,funs_type_apply,times_tuple_intro,real_upper_bound,real_supremum}.
\end{proof}

% from §50 helper lemma 51: finite dimension gives a controlled finite indexed cover
\begin{lemma}\label{controlled_finite_indexed_cover_50}
    Assume $m\in\naturals$.
    Assume $m$ is a topological dimension bound for $X$ with respect to $T$.
    Assume $d$ is a metric on $X$ and $T=\metricTopology{d}{X}$.
    Assume $X$ is compact with respect to $T$.
    Assume $X$ is inhabited.
    Assume $e$ is a metric on $Y$.
    Assume $f$ is uniformly continuous from $X$ to $Y$ with respect to $d$ and $e$.
    Assume $\epsilon,\delta\in\reals$ and $0<\epsilon$ and $0<\delta$.
    Then there exist $n,U$ such that
        $n\in\naturalsPlus$ and
        $U$ is a bijection from $\initialSegment{n}$ to $\img{U}{\initialSegment{n}}$ and
        $\img{U}{\initialSegment{n}}$ is a covering of $X$ and
        $\img{U}{\initialSegment{n}}$ has order at most $m+1$ in $X$ and
        for all $i\in\initialSegment{n}$ we have
            $U(i)$ is open in $X$ with respect to $T$ and $U(i)$ is inhabited and
            for all $x,y\in U(i)$ we have
                $d((x,y))<\epsilon$ and $e((f(x),f(y)))<\delta$.
\end{lemma}
\begin{proof}
    Take $\eta\in\reals$ such that
        $0<\eta$ and
        for all $x,y\in X$ if $d((x,y))<\eta$ then $e((f(x),f(y)))<\delta$
        by \cref{uniformly_continuous}.
    Let $a=\epsilon/(1+1)$.
    $a\in\reals$ and $0<a$ and $a+a=\epsilon$ by
        \cref{real_half_of_positive}.
    Let $b=\eta/(1+1)$.
    $b\in\reals$ and $0<b$ and $b+b=\eta$ by
        \cref{real_half_of_positive}.
    Take $r\in\reals$ such that $0<r$ and $r\leq a$ and $r\leq b$
        by \cref{positive_radius_below_two_bounds_49}.
    $X$ is totally bounded with respect to $d$ by
        \cref{compact_metric_space_totally_bounded}.
    Take $C$ such that
        $C$ is finite and
        $C$ is a covering of $X$ and
        for all $V\in C$ there exists $c\in X$ such that
            $V=X\inter\metricBall{d}{X}{c}{r}$
        by \cref{totally_bounded_metric_space,totally_bounded_subset_metric_space}.
    For all $V\in C$ we have $V$ is open in $X$ with respect to $T$ by
        \cref{metric_ball,subseteq,inter_absorb_supseteq_left,metric_basis,powerset_intro,metric_basis_is_basis,metric_topology,basis_elem_open_in_generated,basis_topology_is_topology,basis_for_topology,open_in}.
    $C$ is an open covering of $X$ with respect to $T$ by
        \cref{open_covering}.
    Take $B$ such that
        $B$ is an open refinement of $C$ on $X$ with respect to $T$ and
        $B$ has order at most $m+1$ in $X$
        by \cref{topological_dimension_bound_50}.
    $B$ is an open covering of $X$ with respect to $T$ by
        \cref{open_refinement_on,open_covering}.
    Take $D$ such that
        $D\subseteq B$ and $D$ is finite and $D$ is a covering of $X$
        by \cref{compact}.
    Let $F=D\setminus\{\emptyset\}$.
    $F\subseteq D$ and $F$ is finite by
        \cref{setminus_subseteq,subseteq_finite_is_finite}.
    $F$ is a covering of $X$ by
        \cref{covering,subseteq,unions_iff,emptyset,setminus_intro,singleton_iff}.
    $F$ is inhabited by
        \cref{covering,subseteq,unions_iff,nonempty_inhabited}.
    Take $n,U$ such that
        $n\in\naturalsPlus$ and
        $U$ is a bijection from $\initialSegment{n}$ to $F$
        by \cref{finite_inhabited_initial_segment_enumeration}.
    $U\in\surj{\initialSegment{n}}{F}$ by \cref{bijections}.
    $\img{U}{\initialSegment{n}}=F$ by
        \cref{surj,funs_elim,funs_surj_iff,img_dom}.
    $U$ is a bijection from $\initialSegment{n}$ to
        $\img{U}{\initialSegment{n}}$ by \cref{real_lt_subst_right}.
    For all $i\in\initialSegment{n}$ we have
        $U(i)$ is open in $X$ with respect to $T$ and $U(i)$ is inhabited by
        \cref{bijections_to_funs,funs_type_apply,setminus_elim_left,setminus_elim_right,singleton_iff,subseteq,open_refinement_on,nonempty_inhabited}.
    $\img{U}{\initialSegment{n}}$ is a covering of $X$ by
        \cref{real_lt_subst_left}.
    $F\subseteq B$ by \cref{subseteq_transitive}.
    $F$ has order at most $m+1$ in $X$ by
        \cref{subcollection_preserves_order_bound_50}.
    $\img{U}{\initialSegment{n}}$ has order at most $m+1$ in $X$ by
        \cref{real_lt_subst_left}.
    Show that for all $i\in\initialSegment{n}$ for all $x,y\in U(i)$ we have
        $d((x,y))<\epsilon$ and $e((f(x),f(y)))<\delta$.
    \begin{subproof}
        Fix $i\in\initialSegment{n}$.
        Fix $x$.
        Fix $y$.
        Assume $x\in U(i)$ and $y\in U(i)$.
        $U(i)\in F$ by \cref{bijections_to_funs,funs_type_apply}.
        $U(i)\in B$ by
            \cref{setminus_elim_left,subseteq}.
        Take $V\in C$ such that $U(i)\subseteq V$ by
            \cref{open_refinement_on,refinement_on}.
        Take $c\in X$ such that
            $V=X\inter\metricBall{d}{X}{c}{r}$.
        $x,y\in V$ by \cref{subseteq}.
        $x,y\in\metricBall{d}{X}{c}{r}$ by
            \cref{inter_elim_right,real_lt_subst_left}.
        $d((c,x))<r$ and $d((c,y))<r$ by \cref{metric_ball}.
        $x,y,c\in X$ by \cref{metric_ball}.
        $d((x,c))=d((c,x))$ by
            \cref{metric_on,metric_symmetric}.
        $d((x,c))<r$ by \cref{real_lt_subst_left}.
        $d((x,c)),d((c,y)),d((x,y)),d((x,c))+d((c,y)),r+r,a+a,b+b\in\reals$ by
            \cref{metric_on,funs_type_apply,times_tuple_intro,real_add_closed}.
        $d((x,y))\leq d((x,c))+d((c,y))$ by
            \cref{metric_on,metric_triangle_inequality}.
        $d((x,c))+d((c,y))<r+r$ by
            \cref{real_lt_add_two}.
        $r+r\leq a+a$ and $r+r\leq b+b$ by \cref{real_leq_add_two_49}.
        $d((x,y))<\epsilon$ and $d((x,y))<\eta$ by
            \cref{real_leq_lt_trans,real_lt_leq_trans_local,real_lt_subst_right}.
        $e((f(x),f(y)))<\delta$ by \cref{real_lt_subst_left}.
    \end{subproof}
    For all $i\in\initialSegment{n}$ we have
        $U(i)$ is open in $X$ with respect to $T$ and $U(i)$ is inhabited and
        for all $x,y\in U(i)$ we have
            $d((x,y))<\epsilon$ and $e((f(x),f(y)))<\delta$.
    Follows by \cref{controlled_finite_indexed_cover_packaging_50}.
\end{proof}

% from §50 helper lemma 114: every bijective enumeration of an independent set is independent
\begin{lemma}\label{independent_set_bijective_enumeration_50}
    Assume $N,n\in\naturalsPlus$.
    Let $I=\initialSegment{n}$.
    Let $E=\realsPower{\initialSegment{N}}$.
    Assume $A$ is a geometrically independent set in dimension $N$.
    Assume $A$ is inhabited.
    Assume $g$ is a bijection from $I$ to $A$.
    Then $g$ is a geometrically independent sequence in dimension $N$.
\end{lemma}
\begin{proof}
    $A$ is finite and $A\subseteq E$ by
        \cref{geometrically_independent_set_50}.
    $A\neq\emptyset$ by \cref{inhabited_nonempty}.
    Take $m,s$ such that
        $m\in\naturalsPlus$ and
        $s$ is a bijection from $\initialSegment{m}$ to $A$ and
        $s$ is a geometrically independent sequence in dimension $N$ by
        \cref{geometrically_independent_set_50}.
    Let $K=\initialSegment{m}$.
    $g\in\funs{I}{A}$ by \cref{bijections_to_funs}.
    $g$ is a function and $\dom{g}=I$ by \cref{funs_elim}.
    For all $i\in\dom{g}$ we have $g(i)\in E$ by
        \cref{funs_type_apply,subseteq}.
    $g\in\funs{I}{E}$ by \cref{funs_intro}.
    Show that for all $a$ if
        $a\in\funs{I}{\reals}$ and
        $\zeroVector{N}$ is an affine combination of $g$ with coefficients $a$
            in dimension $N$ and
        $\sumFiniteSequence{a}{0}$
    then for all $i\in I$ we have $a(i)=0$.
    \begin{subproof}
        Fix $a$.
        Assume $a\in\funs{I}{\reals}$ and
            $\zeroVector{N}$ is an affine combination of $g$ with coefficients $a$
                in dimension $N$ and
            $\sumFiniteSequence{a}{0}$.
        Take $c,b$ such that
            $c\in\funs{A}{\reals}$ and
            $b\in\funs{K}{\reals}$ and
            $\sumFiniteSequence{b}{0}$ and
            for all $i\in I$ we have $a(i)=c(g(i))$ and
            for all $i\in K$ we have $b(i)=c(s(i))$ by
            \cref{coefficient_sequence_reindex_50}.
        $s\in\funs{K}{A}$ by \cref{bijections_to_funs}.
        $s$ is a function and $\dom{s}=K$ by \cref{funs_elim}.
        For all $i\in\dom{s}$ we have $s(i)\in E$ by
            \cref{funs_type_apply,subseteq}.
        $s\in\funs{K}{E}$ by \cref{funs_intro}.
        $\zeroVector{N}\in E$ by \cref{zero_vector_euclidean_50}.
        Take $l$ such that
            $l\in\naturalsPlus$ and
            $g\in\funs{\initialSegment{l}}{E}$ and
            $a\in\funs{\initialSegment{l}}{\reals}$ and
            for all $j\in\initialSegment{N}$ there exists $e$ such that
                $e\in\funs{\initialSegment{l}}{\reals}$ and
                for all $i\in\initialSegment{l}$ we have
                    $e(i)=a(i)\rmul\apply{\apply{g}{i}}{j}$ and
                $\sumFiniteSequence{e}{\apply{\zeroVector{N}}{j}}$ by
            \cref{affine_combination_value_50}.
        $\initialSegment{l}=I$ by \cref{funs_elim}.
        Show that for all $j\in\initialSegment{N}$ there exists $d$ such that
            $d\in\funs{K}{\reals}$ and
            for all $i\in K$ we have
                $d(i)=b(i)\rmul\apply{\apply{s}{i}}{j}$ and
            $\sumFiniteSequence{d}{\apply{\zeroVector{N}}{j}}$.
        \begin{subproof}
            Fix $j\in\initialSegment{N}$.
            Take $e$ such that
                $e\in\funs{\initialSegment{l}}{\reals}$ and
                for all $i\in\initialSegment{l}$ we have
                    $e(i)=a(i)\rmul\apply{\apply{g}{i}}{j}$ and
                $\sumFiniteSequence{e}{\apply{\zeroVector{N}}{j}}$.
            $1\in\initialSegment{l}$ by
                \cref{real_one_in_naturals,real_one_leq_naturalsplus,initial_segment_naturalsplus}.
            $\sumFiniteSequence{e}{\apply{\zeroVector{N}}{j}}$ by
                \cref{subseteq}.
            $e\in\funs{I}{\reals}$ by
                \cref{real_lt_subst_left}.
            Let $v=\{(z,c(z)\rmul z(j))\mid z\in A\}$.
            $v$ is a function by \cref{relation,rightunique,pair_eq_iff}.
            $\dom{v}=A$ by
                \cref{dom_intro,dom_iff,pair_eq_iff,subseteq,subseteq_antisymmetric}.
            For all $z\in\dom{v}$ we have $v(z)\in\reals$ by
                \cref{subseteq,reals_power,tuple_power,funs_type_apply,real_mul_closed,pair_eq_iff,function_apply_intro,real_lt_subst_left}.
            $v\in\funs{A}{\reals}$ by \cref{funs_intro}.
            For all $i\in I$ we have $e(i)=v(g(i))$ by
                \cref{bijections_to_funs,funs_type_apply,pair_eq_iff,function_apply_intro,subseteq,real_lt_subst_left,real_lt_subst_right}.
            Let $d=\{(i,v(s(i)))\mid i\in K\}$.
            $d$ is a function by \cref{relation,rightunique,pair_eq_iff}.
            $\dom{d}=K$ by
                \cref{dom_intro,dom_iff,pair_eq_iff,subseteq,subseteq_antisymmetric}.
            For all $i\in\dom{d}$ we have $d(i)\in\reals$ by
                \cref{bijections_to_funs,funs_type_apply,pair_eq_iff,function_apply_intro,real_lt_subst_left}.
            $d\in\funs{K}{\reals}$ by \cref{funs_intro}.
            For all $i\in K$ we have $d(i)=v(s(i))$ by
                \cref{pair_eq_iff,function_apply_intro}.
            Take $q\in\reals$ such that $\sumFiniteSequence{d}{q}$ by
                \cref{sum_finite_sequence_exists}.
            $\apply{\zeroVector{N}}{j}=q$ by
                \cref{finite_sum_reindex_bijection_50}.
            $\sumFiniteSequence{d}{\apply{\zeroVector{N}}{j}}$ by
                \cref{real_lt_subst_right}.
            For all $i\in K$ we have
                $d(i)=b(i)\rmul\apply{\apply{s}{i}}{j}$ by
                \cref{bijections_to_funs,funs_type_apply,pair_eq_iff,function_apply_intro,subseteq,real_lt_subst_left,real_lt_subst_right}.
            Follows.
        \end{subproof}
        $\zeroVector{N}$ is an affine combination of $s$ with coefficients $b$
            in dimension $N$ by
            \cref{affine_combination_value_50}.
        Take $q$ such that
            $q\in\naturalsPlus$ and
            $s\in\funs{\initialSegment{q}}{E}$ and
            for all $d$ if
                $d\in\funs{\initialSegment{q}}{\reals}$ and
                $\zeroVector{N}$ is an affine combination of $s$ with coefficients $d$
                    in dimension $N$ and
                $\sumFiniteSequence{d}{0}$
            then for all $i\in\initialSegment{q}$ we have $d(i)=0$ by
            \cref{geometrically_independent_sequence_50}.
        $\initialSegment{q}=K$ by \cref{funs_elim}.
        For all $i\in K$ we have $b(i)=0$ by
            \cref{real_lt_subst_left,real_lt_subst_right}.
        For all $i\in I$ we have $a(i)=0$ by
            \cref{bijections_to_funs,funs_type_apply,bijections,surj,subseteq,real_lt_subst_left,real_lt_subst_right}.
    \end{subproof}
    Follows by \cref{geometrically_independent_sequence_50}.
\end{proof}

% from §50 helper lemma 115: affine coordinates in an independent sequence are unique
\begin{lemma}\label{independent_affine_coordinates_unique_50}
    Assume $N,n\in\naturalsPlus$.
    Let $I=\initialSegment{n}$.
    Let $E=\realsPower{\initialSegment{N}}$.
    Assume $s$ is a geometrically independent sequence in dimension $N$.
    Assume $s\in\funs{I}{E}$.
    Assume $a,b\in\funs{I}{\reals}$.
    Assume $\sumFiniteSequence{a}{1}$ and $\sumFiniteSequence{b}{1}$.
    Assume $y$ is an affine combination of $s$ with coefficients $a$ in dimension $N$.
    Assume $y$ is an affine combination of $s$ with coefficients $b$ in dimension $N$.
    Then $a=b$.
\end{lemma}
\begin{proof}
    $y\in E$ by \cref{affine_combination_value_50}.
    $1,\neg{1}\in\reals$ by
        \cref{real_one_in_naturals,real_naturals_subset_reals,subseteq,real_add_inverse}.
    Let $b_0=\coordScalarSequence{n}{\neg{1}}{b}$.
    Let $c=\coordSumSequence{n}{a}{b_0}$.
    $b_0,c\in\funs{I}{\reals}$ by
        \cref{coord_scalar_sequence_function,coord_sum_sequence_function}.
    $\sumFiniteSequence{b_0}{\neg{1}\rmul 1}$ by
        \cref{sum_finite_sequence_scalar}.
    $\neg{1}\rmul 1=\neg{1}$ by
        \cref{real_mul_ident,real_mul_comm}.
    $\sumFiniteSequence{b_0}{\neg{1}}$ by
        \cref{real_lt_subst_right}.
    $\sumFiniteSequence{c}{1+\neg{1}}$ by
        \cref{sum_finite_sequence_add}.
    $1+\neg{1}=0$ by
        \cref{real_add_inverse}.
    $\sumFiniteSequence{c}{0}$ by
        \cref{real_lt_subst_right}.
    Show that $\zeroVector{N}$ is an affine combination of $s$ with coefficients $c$
        in dimension $N$.
    \begin{subproof}
        $\zeroVector{N}\in E$ by
            \cref{zero_vector_euclidean_50}.
        Show that for all $j\in\initialSegment{N}$ there exists $q$ such that
            $q\in\funs{I}{\reals}$ and
            for all $i\in I$ we have
                $q(i)=c(i)\rmul\apply{\apply{s}{i}}{j}$ and
            $\sumFiniteSequence{q}{\apply{\zeroVector{N}}{j}}$.
        \begin{subproof}
            Fix $j\in\initialSegment{N}$.
            Take $u$ such that
                $u\in\funs{I}{\reals}$ and
                for all $i\in I$ we have
                    $u(i)=a(i)\rmul\apply{\apply{s}{i}}{j}$ and
                $\sumFiniteSequence{u}{y(j)}$ by
                \cref{affine_combination_value_50,funs_elim,real_lt_subst_left,real_lt_subst_right}.
            Take $v$ such that
                $v\in\funs{I}{\reals}$ and
                for all $i\in I$ we have
                    $v(i)=b(i)\rmul\apply{\apply{s}{i}}{j}$ and
                $\sumFiniteSequence{v}{y(j)}$ by
                \cref{affine_combination_value_50,funs_elim,real_lt_subst_left,real_lt_subst_right}.
            $1\in I$ by
                \cref{real_one_in_naturals,real_one_leq_naturalsplus,initial_segment_naturalsplus}.
            $\sumFiniteSequence{u}{y(j)}$ and
                $\sumFiniteSequence{v}{y(j)}$ by \cref{subseteq}.
            Let $v_0=\coordScalarSequence{n}{\neg{1}}{v}$.
            Let $q=\coordSumSequence{n}{u}{v_0}$.
            $v_0,q\in\funs{I}{\reals}$ by
                \cref{coord_scalar_sequence_function,coord_sum_sequence_function}.
            $\sumFiniteSequence{v_0}{\neg{1}\rmul y(j)}$ by
                \cref{sum_finite_sequence_scalar}.
            $y(j)\in\reals$ by
                \cref{reals_power,tuple_power,funs_type_apply}.
            $\neg{1}\rmul y(j)=\neg{y(j)}$ by
                \cref{real_neg_as_mul}.
            $\sumFiniteSequence{v_0}{\neg{y(j)}}$ by
                \cref{real_lt_subst_right}.
            $\sumFiniteSequence{q}{y(j)+\neg{y(j)}}$ by
                \cref{sum_finite_sequence_add}.
            $y(j)+\neg{y(j)}=0$ by
                \cref{reals_power,tuple_power,funs_type_apply,real_add_inverse}.
            $\sumFiniteSequence{q}{0}$ by
                \cref{real_lt_subst_right}.
            $(j,0)\in\zeroVector{N}$ by
                \cref{zero_vector_50,pair_eq_iff}.
            $\apply{\zeroVector{N}}{j}=0$ by
                \cref{zero_vector_euclidean_50,reals_power,tuple_power,funs_is_function,function_apply_intro}.
            $\sumFiniteSequence{q}{\apply{\zeroVector{N}}{j}}$ by
                \cref{real_lt_subst_right}.
            For all $i\in I$ we have
                $q(i)=c(i)\rmul\apply{\apply{s}{i}}{j}$ by
                \cref{subseteq,coord_scalar_sequence,coord_sum_sequence,funs_is_function,function_apply_intro,funs_type_apply,reals_power,tuple_power,real_right_distrib,real_add_inverse,real_mul_closed,real_mul_assoc,real_lt_subst_left,real_lt_subst_right}.
            Follows.
        \end{subproof}
        Follows by \cref{affine_combination_value_50}.
    \end{subproof}
    Take $m$ such that
        $m\in\naturalsPlus$ and
        $s\in\funs{\initialSegment{m}}{E}$ and
        for all $d$ if
            $d\in\funs{\initialSegment{m}}{\reals}$ and
            $\zeroVector{N}$ is an affine combination of $s$ with coefficients $d$
                in dimension $N$ and
            $\sumFiniteSequence{d}{0}$
        then for all $i\in\initialSegment{m}$ we have $d(i)=0$ by
        \cref{geometrically_independent_sequence_50}.
    $\initialSegment{m}=I$ by \cref{funs_elim}.
    For all $i\in I$ we have $c(i)=0$ by
        \cref{real_lt_subst_left,real_lt_subst_right}.
    For all $i\in I$ we have $a(i)=b(i)$ by
        \cref{coord_scalar_sequence,coord_sum_sequence,funs_is_function,function_apply_intro,funs_type_apply,real_neg_as_mul,real_lt_subst_left,real_lt_subst_right,real_add_inverse,real_add_cancel}.
    Follows by \cref{funs_is_function,functions_eq_iff}.
\end{proof}

% from §50 helper lemma 116: finite partitions of unity produce continuous barycentric maps
\begin{lemma}\label{finite_partition_barycentric_map_50}
    Assume $N,n\in\naturalsPlus$.
    Let $I=\initialSegment{n}$.
    Let $J=\initialSegment{N}$.
    Let $E=\realsPower{J}$.
    Let $S=\metricTopology{\squareMetric{N}}{E}$.
    Assume $T$ is a topology on $X$.
    Assume $p$ is a partition of unity on $X$ dominated by $U$
        with respect to $T$ and $n$.
    Assume $z\in\funs{I}{E}$.
    Then there exists $g$ such that
        $g$ is continuous from $X$ to $E$ with respect to $T$ and $S$ and
        for all $x\in X$ there exists $a$ such that
            $a\in\funs{I}{\reals}$ and
            for all $i\in I$ we have $a(i)=\apply{\apply{p}{i}}{x}$ and
            $\sumFiniteSequence{a}{1}$ and
            $g(x)$ is an affine combination of $z$ with coefficients $a$
                in dimension $N$.
\end{lemma}
\begin{proof}
    Let $I_0=\intervalclosed{0}{1}$.
    Let $j_0=\identity{I_0}$.
    Let $R=\{w\in J\times\pow{X\times\reals}\mid
        \text{$\snd{w}$ is continuous from $X$ to $\reals$
            with respect to $T$ and $\standardTopologyR$ and
            for all $x\in X$ there exists $b$ such that
                $b\in\funs{I}{\reals}$ and
                $\sumFiniteSequence{b}{\apply{\snd{w}}{x}}$ and
                for all $i\in I$ we have
                    $b(i)=\apply{\apply{p}{i}}{x}
                        \rmul\apply{\apply{z}{i}}{\fst{w}}$}\}$.
    $R$ is a relation by
        \cref{relation,subseteq,times_elem_is_tuple}.
    Show that for all $j\in J$ there exists $r$ such that $j\mathrel{R}r$.
    \begin{subproof}
        Fix $j\in J$.
        Let $Q=\{w\in I\times\pow{X\times\reals}\mid
            \text{there exist $c,f$ such that
                $c=\{(x,\apply{\apply{z}{\fst{w}}}{j})\mid x\in X\}$ and
                $f=j_0\circ\apply{p}{\fst{w}}$ and
                $\snd{w}=\{(x,c(x)\rmul f(x))\mid x\in X\}$ and
                $\snd{w}$ is continuous from $X$ to $\reals$
                    with respect to $T$ and $\standardTopologyR$}\}$.
        $Q$ is a relation by
            \cref{relation,subseteq,times_elem_is_tuple}.
        Show that for all $i\in I$ there exists $r$ such that $i\mathrel{Q}r$.
        \begin{subproof}
            Fix $i\in I$.
            $\apply{p}{i}$ is continuous from $X$ to $I_0$ with respect to
                $T$ and $\subspaceTopology{\standardTopologyR}{I_0}$ by
                \cref{finite_partition_of_unity,finite_partition_component_condition}.
            Let $f=j_0\circ\apply{p}{i}$.
            $f$ is continuous from $X$ to $\reals$ with respect to
                $T$ and $\standardTopologyR$ by
                \cref{continuous_interval_expand_range}.
            $\apply{z}{i}\in E$ by \cref{funs_type_apply}.
            $\apply{\apply{z}{i}}{j}\in\reals$ by
                \cref{reals_power,tuple_power,funs_type_apply}.
            Let $c=\{(x,\apply{\apply{z}{i}}{j})\mid x\in X\}$.
            Let $r=\{(x,c(x)\rmul f(x))\mid x\in X\}$.
            $r$ is continuous from $X$ to $\reals$ with respect to
                $T$ and $\standardTopologyR$ by
                \cref{real_function_scalar_multiple_continuous_49}.
            $r\subseteq X\times\reals$ by
                \cref{continuous,funs,funs_subseteq_rels,rels_ran_subseteq,rels_elim,subseteq}.
            $r\in\pow{X\times\reals}$ by \cref{powerset_intro}.
            $(i,r)\in I\times\pow{X\times\reals}$ by
                \cref{times_tuple_intro}.
            $(i,r)\in Q$ by
                \cref{fst_eq,snd_eq,subseteq}.
            Follows by \cref{relation}.
        \end{subproof}
        Take $q$ such that
            $q$ is a function on $I$ and
            for all $i\in I$ we have $i\mathrel{Q}q(i)$ by
            \cref{choice_function}.
        For all $i\in I$ we have
            $q(i)$ is continuous from $X$ to $\reals$ with respect to
                $T$ and $\standardTopologyR$ by
            \cref{choice_function,subseteq,fst_eq,snd_eq}.
        Take $r$ such that
            $r$ is a finite real sum witness for $q$ on $X$ with respect to
                $T$ and $n$ by
            \cref{finite_real_sum_continuous}.
        $r$ is continuous from $X$ to $\reals$ with respect to
            $T$ and $\standardTopologyR$ by
            \cref{finite_real_sum_witness}.
        $r\subseteq X\times\reals$ by
            \cref{continuous,funs,funs_subseteq_rels,rels_ran_subseteq,rels_elim,subseteq}.
        $r\in\pow{X\times\reals}$ by \cref{powerset_intro}.
        $(j,r)\in J\times\pow{X\times\reals}$ by
            \cref{times_tuple_intro}.
        Show that for all $x\in X$ there exists $b$ such that
            $b\in\funs{I}{\reals}$ and
            $\sumFiniteSequence{b}{r(x)}$ and
            for all $i\in I$ we have
                $b(i)=\apply{\apply{p}{i}}{x}
                    \rmul\apply{\apply{z}{i}}{j}$.
        \begin{subproof}
            Fix $x\in X$.
            Take $b$ such that
                $b\in\funs{I}{\reals}$ and
                $\sumFiniteSequence{b}{r(x)}$ and
                for all $i\in I$ we have
                    $b(i)=\apply{\apply{q}{i}}{x}$ by
                \cref{finite_real_sum_witness}.
            Show that for all $i\in I$ we have
                $b(i)=\apply{\apply{p}{i}}{x}
                    \rmul\apply{\apply{z}{i}}{j}$.
            \begin{subproof}
                Fix $i\in I$.
                $i\mathrel{Q}q(i)$ by \cref{choice_function}.
                Take $c,f$ such that
                    $c=\{(u,\apply{\apply{z}{i}}{j})\mid u\in X\}$ and
                    $f=j_0\circ\apply{p}{i}$ and
                    $q(i)=\{(u,c(u)\rmul f(u))\mid u\in X\}$ and
                    $q(i)$ is continuous from $X$ to $\reals$
                        with respect to $T$ and $\standardTopologyR$ by
                    \cref{subseteq,fst_eq,snd_eq}.
                $c$ is a function from $X$ to $\reals$ by
                    \cref{constant_map_function,reals_power,tuple_power,funs_type_apply}.
                $c(x)=\apply{\apply{z}{i}}{j}$ by
                    \cref{constant_map_function,funs_is_function,function_apply_intro}.
                $\apply{p}{i}$ is continuous from $X$ to $I_0$ with respect to
                    $T$ and $\subspaceTopology{\standardTopologyR}{I_0}$ by
                    \cref{finite_partition_of_unity,finite_partition_component_condition}.
                $f(x)=\apply{\apply{p}{i}}{x}$ by
                    \cref{continuous,funs,funs_is_function,function_rightunique,funs_is_relation,funs_subseteq_rels,subseteq,rels_ran_subseteq,id_dom,id_rightunique,id_is_relation,circ_apply,funs_type_apply,id_apply}.
                $(x,c(x)\rmul f(x))\in q(i)$ by
                    \cref{real_lt_subst_left,pair_eq_iff}.
                $\apply{\apply{q}{i}}{x}=c(x)\rmul f(x)$ by
                    \cref{continuous,funs_is_function,function_apply_intro}.
                $b(i)=\apply{\apply{q}{i}}{x}$ by \cref{subseteq}.
                $\apply{\apply{p}{i}}{x},\apply{\apply{z}{i}}{j}\in\reals$ by
                    \cref{continuous,funs_type_apply,intervalclosed,subseteq,reals_power,tuple_power}.
                Follows by
                    \cref{real_mul_comm,real_lt_subst_left,real_lt_subst_right}.
            \end{subproof}
            Follows.
        \end{subproof}
        $(j,r)\in R$ by
            \cref{fst_eq,snd_eq,subseteq}.
        Follows by \cref{relation}.
    \end{subproof}
    Take $h$ such that
        $h$ is a function on $J$ and
        for all $j\in J$ we have $j\mathrel{R}h(j)$ by
        \cref{choice_function}.
    For all $j\in J$ we have
        $h(j)$ is continuous from $X$ to $\reals$ with respect to
            $T$ and $\standardTopologyR$ by
        \cref{choice_function,subseteq,fst_eq,snd_eq}.
    Let $R_0=\{(j,\reals)\mid j\in J\}$.
    Let $S_0=\{(j,\standardTopologyR)\mid j\in J\}$.
    Let $P=\productTopologyIndexed{S_0}{R_0}$.
    Let $g=\{(x,v)\mid x\in X,v\in\productFamily{R_0}\mid
        \text{for all $j\in J$ we have $v(j)=\apply{\apply{h}{j}}{x}$}\}$.
    $g$ is a function from $X$ to $\productFamily{R_0}$ and
        $g$ is continuous from $X$ to $\productFamily{R_0}$ with respect to
            $T$ and $P$ and
        for all $x\in X$ for all $j\in J$ we have
            $\apply{\apply{g}{x}}{j}=\apply{\apply{h}{j}}{x}$ by
        \cref{finite_real_family_evaluation_map}.
    $J$ is inhabited by
        \cref{initial_segment_inhabited}.
    $\productFamily{R_0}=E$ by
        \cref{product_family_reals_power}.
    $\metricTopology{\euclideanMetric{N}}{E}=P$ by
        \cref{euclidean_metric_product_topology}.
    $\metricTopology{\euclideanMetric{N}}{E}=S$ by
        \cref{square_metric_topology_equals_euclidean}.
    $P=S$.
    $g$ is continuous from $X$ to $E$ with respect to $T$ and $S$ by
        \cref{real_lt_subst_left,real_lt_subst_right}.
    Show that for all $x\in X$ there exists $a$ such that
        $a\in\funs{I}{\reals}$ and
        for all $i\in I$ we have $a(i)=\apply{\apply{p}{i}}{x}$ and
        $\sumFiniteSequence{a}{1}$ and
        $g(x)$ is an affine combination of $z$ with coefficients $a$
            in dimension $N$.
    \begin{subproof}
        Fix $x\in X$.
        Take $a$ such that
            $a\in\funs{I}{\reals}$ and
            for all $i\in I$ we have $a(i)=\apply{\apply{p}{i}}{x}$ and
            $\sumFiniteSequence{a}{1}$ by
            \cref{finite_partition_of_unity_sum_witness}.
        $g(x)\in E$ by \cref{continuous,funs_type_apply}.
        Show that for all $j\in J$ there exists $b$ such that
            $b\in\funs{I}{\reals}$ and
            for all $i\in I$ we have
                $b(i)=a(i)\rmul\apply{\apply{z}{i}}{j}$ and
            $\sumFiniteSequence{b}{\apply{\apply{g}{x}}{j}}$.
        \begin{subproof}
            Fix $j\in J$.
            $j\mathrel{R}h(j)$ by \cref{choice_function}.
            Take $b$ such that
                $b\in\funs{I}{\reals}$ and
                $\sumFiniteSequence{b}{\apply{\apply{h}{j}}{x}}$ and
                for all $i\in I$ we have
                    $b(i)=\apply{\apply{p}{i}}{x}
                        \rmul\apply{\apply{z}{i}}{j}$ by
                \cref{subseteq,fst_eq,snd_eq}.
            $\apply{\apply{g}{x}}{j}=\apply{\apply{h}{j}}{x}$ by
                \cref{subseteq}.
            $\sumFiniteSequence{b}{\apply{\apply{g}{x}}{j}}$ by
                \cref{real_lt_subst_right}.
            For all $i\in I$ we have
                $b(i)=a(i)\rmul\apply{\apply{z}{i}}{j}$ by
                \cref{subseteq,real_lt_subst_left}.
            Follows.
        \end{subproof}
        $g(x)$ is an affine combination of $z$ with coefficients $a$
            in dimension $N$ by
            \cref{affine_combination_value_50}.
        Follows.
    \end{subproof}
    Follows.
\end{proof}

% from §50 helper lemma 117: coordinatewise control bounds the square metric
\begin{lemma}\label{square_metric_from_coordinate_bounds_50}
    Assume $N\in\naturalsPlus$.
    Let $J=\initialSegment{N}$.
    Let $E=\realsPower{J}$.
    Assume $x,y\in E$.
    Assume $r\in\reals$.
    Assume for all $j\in J$ we have $\abs{x(j)-y(j)}\leq r$.
    Then $\apply{\squareMetric{N}}{(x,y)}\leq r$.
\end{lemma}
\begin{proof}
    Let $e=\squareMetric{N}$.
    $e$ is a metric on $E$ by \cref{square_metric_is_metric}.
    $e((x,y))\in\reals$ by
        \cref{metric_on,funs_type_apply,times_tuple_intro}.
    $(x,y)\in E\times E$ by \cref{times_tuple_intro}.
    $(x,y)\in\dom{e}$ by \cref{metric_on,funs}.
    $((x,y),e((x,y)))\in e$ by
        \cref{metric_on,funs_is_function,function_apply_elim}.
    $e((x,y))\in\coordDiffSet{N}{x}{y}$ by
        \cref{square_metric,pair_eq_iff}.
    Take $j\in J$ such that
        $e((x,y))=\abs{x(j)-y(j)}$ by
        \cref{coord_diff_set}.
    Follows by \cref{real_leq_subst_left_local}.
\end{proof}

% from §50 helper lemma 118: translated two-sided bounds give an absolute difference bound
\begin{lemma}\label{real_between_implies_abs_difference_50}
    Assume $a,b,r\in\reals$.
    Assume $0\leq r$.
    Assume $b-r\leq a$.
    Assume $a\leq b+r$.
    Then $\abs{a-b}\leq r$.
\end{lemma}
\begin{proof}
    $a-b\in\reals$ by
        \cref{real_sub,real_add_inverse,real_add_closed}.
    $\neg{r}\leq a-b$ by
        \cref{real_sub,real_add_inverse,real_add_closed,real_add_assoc,real_add_ident,real_add_comm,real_leq_add,real_leq_subst_left_local,real_leq_subst_right_local}.
    $a-b\leq r$ by
        \cref{real_sub,real_add_inverse,real_add_closed,real_add_assoc,real_add_ident,real_add_comm,real_leq_add,real_leq_subst_left_local,real_leq_subst_right_local}.
    Follows by \cref{real_abs_leq_from_bounds}.
\end{proof}

% from §50 helper lemma 119: nonnegative affine averages preserve square-metric bounds
\begin{lemma}\label{barycentric_square_metric_bound_50}
    Assume $N,n\in\naturalsPlus$.
    Let $I=\initialSegment{n}$.
    Let $J=\initialSegment{N}$.
    Let $E=\realsPower{J}$.
    Assume $z\in\funs{I}{E}$.
    Assume $a\in\funs{I}{\reals}$.
    Assume $x,y\in E$.
    Assume $r\in\reals$ and $0\leq r$.
    Assume $\sumFiniteSequence{a}{1}$.
    Assume for all $i\in I$ we have $0\leq a(i)$.
    Assume $y$ is an affine combination of $z$ with coefficients $a$
        in dimension $N$.
    Assume for all $i\in I$ if $0<a(i)$ then
        $\apply{\squareMetric{N}}{(z(i),x)}\leq r$.
    Then $\apply{\squareMetric{N}}{(y,x)}\leq r$.
\end{lemma}
\begin{proof}
    Show that for all $j\in J$ we have
        $\abs{y(j)-x(j)}\leq r$.
    \begin{subproof}
        Fix $j\in J$.
        Take $b$ such that
            $b\in\funs{I}{\reals}$ and
            for all $i\in I$ we have
                $b(i)=a(i)\rmul\apply{z(i)}{j}$ and
            $\sumFiniteSequence{b}{y(j)}$ by
            \cref{affine_combination_value_50,funs_elim,real_lt_subst_left,real_lt_subst_right}.
        $1\in I$ by
            \cref{real_one_in_naturals,real_one_leq_naturalsplus,initial_segment_naturalsplus}.
        $\sumFiniteSequence{b}{y(j)}$ by \cref{subseteq}.
        $x(j),y(j)\in\reals$ by
            \cref{reals_power,tuple_power,funs_type_apply}.
        $x(j)-r,x(j)+r\in\reals$ by
            \cref{real_sub,real_add_inverse,real_add_closed}.
        Let $l=\coordScalarSequence{n}{x(j)-r}{a}$.
        Let $u=\coordScalarSequence{n}{x(j)+r}{a}$.
        $l,u\in\funs{I}{\reals}$ by
            \cref{coord_scalar_sequence_function}.
        $\sumFiniteSequence{l}{(x(j)-r)\rmul 1}$ and
            $\sumFiniteSequence{u}{(x(j)+r)\rmul 1}$ by
            \cref{sum_finite_sequence_scalar}.
        $(x(j)-r)\rmul 1=x(j)-r$ and
            $(x(j)+r)\rmul 1=x(j)+r$ by
            \cref{real_mul_ident,real_mul_comm}.
        $\sumFiniteSequence{l}{x(j)-r}$ and
            $\sumFiniteSequence{u}{x(j)+r}$ by
            \cref{real_lt_subst_right}.
        Show that for all $i\in I$ we have
            $l(i)\leq b(i)$ and $b(i)\leq u(i)$.
        \begin{subproof}
            Fix $i\in I$.
            $a(i),b(i),l(i),u(i)\in\reals$ by
                \cref{funs_type_apply}.
            $\apply{z(i)}{j}\in\reals$ by
                \cref{funs_type_apply,reals_power,tuple_power}.
            $l(i)=(x(j)-r)\rmul a(i)$ and
                $u(i)=(x(j)+r)\rmul a(i)$ by
                \cref{coord_scalar_sequence,funs_is_function,function_apply_intro}.
            $b(i)=a(i)\rmul\apply{z(i)}{j}$ by \cref{subseteq}.
            \begin{byCase}
            \caseOf{$0<a(i)$.}
                Take $s\in\reals$ such that
                    $((z(i),x),s)\in\squareMetric{N}$ and
                    for all $k\in J$ we have
                        $\abs{\apply{z(i)}{k}-x(k)}\leq s$ by
                    \cref{square_metric_value_exists,funs_type_apply}.
                $s=\apply{\squareMetric{N}}{(z(i),x)}$ by
                    \cref{square_metric_function,function_apply_intro}.
                $s\leq r$ by \cref{real_leq_subst_left_local}.
                $\abs{\apply{z(i)}{j}-x(j)}\in\reals$ by
                    \cref{real_sub,real_add_inverse,real_add_closed,real_abs_in_reals}.
                $\abs{\apply{z(i)}{j}-x(j)}\leq r$ by
                    \cref{subseteq,real_leq_trans}.
                $x(j)-r\leq\apply{z(i)}{j}$ by
                    \cref{real_abs_sub_leq_imp_left_bound}.
                $\apply{z(i)}{j}\leq x(j)+r$ by
                    \cref{real_abs_sub_leq_imp_right_bound}.
                $a(i)\rmul(x(j)-r)\leq
                    a(i)\rmul\apply{z(i)}{j}$ and
                    $a(i)\rmul\apply{z(i)}{j}\leq
                    a(i)\rmul(x(j)+r)$ by
                    \cref{real_leq_mul_pos_left_49}.
                $(x(j)-r)\rmul a(i)=a(i)\rmul(x(j)-r)$ and
                    $(x(j)+r)\rmul a(i)=a(i)\rmul(x(j)+r)$ by
                    \cref{real_mul_comm}.
                Follows by
                    \cref{real_leq_subst_left_local,real_leq_subst_right_local}.
            \caseOf{$0\not<a(i)$.}
                $a(i)=0$ by \cref{real_leq_split}.
                $l(i)=0$ and $b(i)=0$ and $u(i)=0$ by
                    \cref{real_mul_zero,real_mul_comm,real_lt_subst_left,real_lt_subst_right}.
                Follows.
            \end{byCase}
        \end{subproof}
        $x(j)-r\leq y(j)$ by
            \cref{sum_finite_sequence_coordwise_leq}.
        $y(j)\leq x(j)+r$ by
            \cref{sum_finite_sequence_coordwise_leq}.
        Follows by \cref{real_between_implies_abs_difference_50}.
    \end{subproof}
    Follows by \cref{square_metric_from_coordinate_bounds_50}.
\end{proof}

% from §50 helper lemma 120: indices active at a point inherit the covering-order bound
\begin{lemma}\label{cover_point_index_set_at_most_50}
    Assume $m,n\in\naturalsPlus$.
    Let $I=\initialSegment{n}$.
    Assume $U$ is a bijection from $I$ to $\img{U}{I}$.
    Assume $\img{U}{I}$ has order at most $m$ in $X$.
    Assume $x\in X$.
    Let $A=\{i\in I\mid x\in U(i)\}$.
    Then $A$ has at most $m$ elements.
\end{lemma}
\begin{proof}
    $I$ is finite by \cref{initial_segment_finite}.
    $A\subseteq I$ by \cref{subseteq}.
    $A$ is finite by \cref{subseteq_finite_is_finite}.
    Show that there exists no $C$ such that
        $C\subseteq A$ and
        there exists $g$ such that
            $g$ is a bijection from $\initialSegment{m+1}$ to $C$.
    \begin{subproof}
        Suppose not.
        Take $C,g$ such that
            $C\subseteq A$ and
            $g$ is a bijection from $\initialSegment{m+1}$ to $C$.
        $C\subseteq I$ by \cref{subseteq_transitive}.
        Let $v=\restrl{U}{C}$.
        $U$ is a function and $\dom{U}=I$ by
            \cref{bijection_is_function,bijections_dom}.
        $v$ is a function by
            \cref{restrl_of_function_is_function}.
        $\dom{v}=C$ by
            \cref{dom_of_restrl_eq_inter,inter_absorb_supseteq_left}.
        Let $F=\img{U}{C}$.
        For all $i\in\dom{v}$ we have $v(i)\in F$ by
            \cref{real_lt_subst_left,restrl_of_function_apply,subseteq,img_of_function_intro,inter_intro}.
        $v\in\funs{C}{F}$ by \cref{funs_intro}.
        $v$ is injective by
            \cref{restrl_of_function_apply,bijections,injective_function,subseteq}.
        $\img{v}{C}=F$ by
            \cref{restrl_img,inter_idempotent}.
        $\img{v}{C}=\ran{v}$ by
            \cref{img_dom,real_lt_subst_left}.
        $v\in\surj{C}{F}$ by
            \cref{funs_surj_iff,real_lt_subst_left}.
        $v$ is a bijection from $C$ to $F$ by
            \cref{bijections}.
        Let $q=v\circ g$.
        $q$ is a bijection from $\initialSegment{m+1}$ to $F$ by
            \cref{bijection_circ}.
        $F\subseteq\img{U}{I}$ by
            \cref{img_subseteq}.
        For all $B\in F$ we have $x\in B$ by
            \cref{img_of_function_elim,inter_elim_right,subseteq}.
        $x$ has multiplicity at most $m$ in $\img{U}{I}$ by
            \cref{covering_order_at_most_50}.
        Contradiction by \cref{point_multiplicity_at_most_50}.
    \end{subproof}
    Follows by \cref{finite_set_has_at_most_50}.
\end{proof}

% from §50 helper lemma 121: adjoining one point raises a finite at-most bound by at most one
\begin{lemma}\label{finite_at_most_from_deleted_point_50}
    Assume $m\in\naturalsPlus$.
    Assume $D$ is finite.
    Assume $y\in D$.
    Let $B=D\setminus\{y\}$.
    Assume $B$ has at most $m$ elements.
    Then $D$ has at most $m+1$ elements.
\end{lemma}
\begin{proof}
    $m+1\in\naturalsPlus$ by
        \cref{real_naturals_closed_succ}.
    Show that there exists no $C$ such that
        $C\subseteq D$ and
        there exists $g$ such that
            $g$ is a bijection from
                $\initialSegment{(m+1)+1}$ to $C$.
    \begin{subproof}
        Suppose not.
        Take $C,g$ such that
            $C\subseteq D$ and
            $g$ is a bijection from
                $\initialSegment{(m+1)+1}$ to $C$.
        Let $L=\initialSegment{m+1}$.
        Let $K=\initialSegment{(m+1)+1}$.
        $(m+1)+1\in K$ by
            \cref{real_naturals_closed_succ,initial_segment_naturalsplus}.
        \begin{byCase}
        \caseOf{$y\in C$.}
            Take $p\in K$ such that $g(p)=y$ by
                \cref{bijections,surj}.
            Take $h$ such that
                $h$ is a bijection from $L$ to $K\setminus\{p\}$ by
                \cref{initial_segment_remove_point_enumeration_50,real_naturals_closed_succ}.
            $p\notin K\setminus\{p\}$ by
                \cref{setminus_elim_right,singleton_intro}.
            Take $t$ such that
                $t$ is a bijection from
                    $\initialSegment{(m+1)+1}$ to
                    $(K\setminus\{p\})\union\{p\}$ and
                $\restrl{t}{L}=h$ and
                $t((m+1)+1)=p$ by
                \cref{finite_bijection_append_new_point_50,real_naturals_closed_succ}.
            $(K\setminus\{p\})\union\{p\}\subseteq K$ by
                \cref{setminus_subseteq,singleton_subset_intro,union_subsets_is_subset}.
            Show that $K\subseteq(K\setminus\{p\})\union\{p\}$.
                \begin{subproof}
                    It suffices to show that for all $q\in K$ we have
                        $q\in(K\setminus\{p\})\union\{p\}$ by \cref{subseteq}.
                    Fix $q\in K$.
                    \begin{byCase}
                    \caseOf{$q=p$.}
                        $q\in\{p\}$ by \cref{singleton_iff}.
                        Follows by \cref{union_intro_right}.
                    \caseOf{$q\neq p$.}
                        $q\in K\setminus\{p\}$ by
                            \cref{setminus_intro,singleton_iff}.
                        Follows by \cref{union_intro_left}.
                    \end{byCase}
            \end{subproof}
            $(K\setminus\{p\})\union\{p\}=K$ by
                \cref{subseteq_antisymmetric}.
            $t$ is a bijection from $K$ to $K$ by
                \cref{real_lt_subst_right}.
            Let $q=g\circ t$.
            $q$ is a bijection from $K$ to $C$ by
                \cref{bijection_circ}.
            $q((m+1)+1)=g(p)$ by
                \cref{circ_apply,bijection_is_function,bijections_dom,bijections,ran_of_surj,subseteq}.
            $q((m+1)+1)=y$ by
                \cref{real_lt_subst_right}.
            Let $q_0=\restrl{q}{L}$.
            $q_0$ is a bijection from $L$ to $C\setminus\{y\}$ by
                \cref{finite_enumeration_remove_last,real_naturals_closed_succ}.
            $C\setminus\{y\}\subseteq B$ by
                \cref{setminus_elim_left,setminus_elim_right,setminus_intro,subseteq}.
            Contradiction by \cref{finite_set_has_at_most_50}.
        \caseOf{$y\notin C$.}
            Let $\beta=g((m+1)+1)$.
            $\beta\in C$ by
                \cref{bijections_to_funs,funs_type_apply}.
            Let $q_0=\restrl{g}{L}$.
            $q_0$ is a bijection from $L$ to $C\setminus\{\beta\}$ by
                \cref{finite_enumeration_remove_last,real_naturals_closed_succ}.
            $C\setminus\{\beta\}\subseteq B$ by
                \cref{setminus_elim_left,setminus_intro,singleton_iff,subseteq}.
            Contradiction by \cref{finite_set_has_at_most_50}.
        \end{byCase}
    \end{subproof}
    Follows by \cref{finite_set_has_at_most_50}.
\end{proof}

% from §50 helper lemma 122: adjoining one point preserves the successor at-most bound
\begin{lemma}\label{finite_at_most_adjoin_point_50}
    Assume $A$ has at most $m$ elements.
    Let $D=A\union\{y\}$.
    Then $D$ has at most $m+1$ elements.
\end{lemma}
\begin{proof}
    $m\in\naturalsPlus$ and $A$ is finite by
        \cref{finite_set_has_at_most_50}.
    $\{y\}$ is finite by
        \cref{singleton_at_most_positive_50,finite_set_has_at_most_50}.
    $D$ is finite by \cref{union_of_finite_is_finite}.
    $y\in D$ by
        \cref{singleton_intro,union_intro_right}.
    Let $B=D\setminus\{y\}$.
    $B\subseteq A$ by
        \cref{setminus_elim_left,setminus_elim_right,union_iff,singleton_iff,subseteq}.
    $B$ has at most $m$ elements by
        \cref{finite_at_most_subset_50}.
    Follows by \cref{finite_at_most_from_deleted_point_50}.
\end{proof}

% from §50 helper lemma 123: the empty set obeys every positive at-most bound
\begin{lemma}\label{empty_at_most_positive_50}
    Assume $m\in\naturalsPlus$.
    Then $\emptyset$ has at most $m$ elements.
\end{lemma}
\begin{proof}
    $\emptyset$ is finite by \cref{emptyset_is_finite}.
    Show that there exists no $C$ such that
        $C\subseteq\emptyset$ and
        there exists $g$ such that
            $g$ is a bijection from $\initialSegment{m+1}$ to $C$.
    \begin{subproof}
        Suppose not.
        Take $C,g$ such that
            $C\subseteq\emptyset$ and
            $g$ is a bijection from $\initialSegment{m+1}$ to $C$.
        $C=\emptyset$ by \cref{subseteq_emptyset_iff}.
        $m+1\in\naturalsPlus$ by
            \cref{real_naturals_closed_succ}.
        $1\in\initialSegment{m+1}$ by
            \cref{real_one_in_naturals,real_one_leq_naturalsplus,initial_segment_naturalsplus}.
        $g(1)\in C$ by
            \cref{bijections_to_funs,funs_type_apply}.
        Contradiction by \cref{emptyset}.
    \end{subproof}
    Follows by \cref{finite_set_has_at_most_50}.
\end{proof}

% from §50 helper lemma 124: an inhabited set with at most one element is a singleton
\begin{lemma}\label{inhabited_at_most_one_singleton_50}
    Assume $B$ has at most $1$ elements.
    Assume $y\in B$.
    Then $B=\{y\}$.
\end{lemma}
\begin{proof}
    Show that $B\subseteq\{y\}$.
    \begin{subproof}
        It suffices to show that for all $b\in B$ we have
            $b\in\{y\}$ by \cref{subseteq}.
        Fix $b\in B$.
        Suppose not.
        $b\neq y$ by \cref{singleton_iff}.
        Let $I=\initialSegment{1}$.
        Let $h=\{(i,y)\mid i\in I\}$.
        $y\in\{y\}$ by \cref{singleton_intro}.
        $h$ is a function from $I$ to $\{y\}$ by
            \cref{constant_map_function}.
        $h$ is a function by \cref{funs_is_function}.
        Show that $h$ is injective.
        \begin{subproof}
            It suffices to show that for all $i,j$ such that
                $i\in\dom{h}$ and $j\in\dom{h}$ and $h(i)=h(j)$
            we have $i=j$ by \cref{injective_function}.
            Fix $i$.
            Fix $j$.
            Assume $i\in\dom{h}$ and $j\in\dom{h}$ and $h(i)=h(j)$.
            Follows by
                \cref{funs,initial_segment_one,singleton_iff,real_lt_subst_left,real_lt_subst_right}.
        \end{subproof}
        $1\in I$ by
            \cref{initial_segment_one,singleton_iff}.
        $h(1)=y$ by
            \cref{constant_map_function,funs_is_function,function_apply_intro}.
        $h\in\surj{I}{\{y\}}$ by
            \cref{surj,singleton_iff}.
        $h$ is a bijection from $I$ to $\{y\}$ by
            \cref{bijections}.
        $b\notin\{y\}$ by \cref{singleton_iff}.
        Take $q$ such that
            $q$ is a bijection from $\initialSegment{1+1}$ to
                $\{y\}\union\{b\}$ and
            $\restrl{q}{I}=h$ and
            $q(1+1)=b$ by
            \cref{finite_bijection_append_new_point_50,real_one_in_naturals}.
        $\{y\}\union\{b\}\subseteq B$ by
            \cref{singleton_subset_intro,union_subsets_is_subset}.
        Contradiction by \cref{finite_set_has_at_most_50}.
    \end{subproof}
    $\{y\}\subseteq B$ by \cref{singleton_subset_intro}.
    Follows by \cref{subseteq_antisymmetric}.
\end{proof}

% from §50 helper lemma 125: finite at-most bounds add under unions
\begin{lemma}\label{finite_at_most_union_sum_50}
    Assume $m,k\in\naturalsPlus$.
    Assume $A$ has at most $m$ elements.
    Assume $B$ has at most $k$ elements.
    Then $A\union B$ has at most $m+k$ elements.
\end{lemma}
\begin{proof}
    Let $S=\{q\in\naturalsPlus\mid
        \text{for all $r,C,D$ if
            $r\in\naturalsPlus$ and
            $C$ has at most $r$ elements and
            $D$ has at most $q$ elements
        then $C\union D$ has at most $r+q$ elements}\}$.
    $S\subseteq\naturalsPlus$ by \cref{subseteq}.
    $1\in\naturalsPlus$ by \cref{real_one_in_naturals}.
    Show that for all $r,C,D$ if
            $r\in\naturalsPlus$ and
            $C$ has at most $r$ elements and
            $D$ has at most $1$ elements
        then $C\union D$ has at most $r+1$ elements.
        \begin{subproof}
            Fix $r$.
            Fix $C$.
            Fix $D$.
            Assume $r\in\naturalsPlus$ and
                $C$ has at most $r$ elements and
                $D$ has at most $1$ elements.
            \begin{byCase}
            \caseOf{$D=\emptyset$.}
                $C\union D=C$ by \cref{union_emptyset}.
                $C$ has at most $r+1$ elements by
                    \cref{finite_at_most_successor_50}.
                Follows by \cref{real_lt_subst_left}.
            \caseOf{$D\neq\emptyset$.}
                Take $y$ such that $y\in D$ by \cref{nonempty_inhabited}.
                $D=\{y\}$ by
                    \cref{inhabited_at_most_one_singleton_50}.
                $C\union\{y\}$ has at most $r+1$ elements by
                    \cref{finite_at_most_adjoin_point_50}.
                Follows by \cref{real_lt_subst_left}.
            \end{byCase}
    \end{subproof}
    $1\in S$ by \cref{subseteq}.
    Show that for all $q\in S$ we have $q+1\in S$.
    \begin{subproof}
        Fix $q\in S$.
        $q\in\naturalsPlus$ by \cref{subseteq}.
        $q+1\in\naturalsPlus$ by
            \cref{real_naturals_closed_succ}.
        Show that for all $r,C,D$ if
            $r\in\naturalsPlus$ and
            $C$ has at most $r$ elements and
            $D$ has at most $q+1$ elements
        then $C\union D$ has at most $r+(q+1)$ elements.
        \begin{subproof}
            Fix $r$.
            Fix $C$.
            Fix $D$.
            Assume $r\in\naturalsPlus$ and
                $C$ has at most $r$ elements and
                $D$ has at most $q+1$ elements.
            $r+q\in\naturalsPlus$ by
                \cref{naturalsplus_add_closed}.
            \begin{byCase}
            \caseOf{$D=\emptyset$.}
                $\emptyset$ has at most $q$ elements by
                    \cref{empty_at_most_positive_50}.
                $C\union\emptyset$ has at most $r+q$ elements by
                    \cref{subseteq}.
                $C\union\emptyset$ has at most $(r+q)+1$ elements by
                    \cref{finite_at_most_successor_50}.
                $(r+q)+1=r+(q+1)$ by
                    \cref{real_naturals_subset_reals,real_one_in_naturals,subseteq,real_add_assoc}.
                Follows by
                    \cref{real_lt_subst_left,real_lt_subst_right}.
            \caseOf{$D\neq\emptyset$.}
                Take $y$ such that $y\in D$ by \cref{nonempty_inhabited}.
                Let $D_0=D\setminus\{y\}$.
                $D_0$ has at most $q$ elements by
                    \cref{finite_at_most_remove_point_50}.
                $C\union D_0$ has at most $r+q$ elements by
                    \cref{subseteq}.
                Let $P=(C\union D_0)\union\{y\}$.
                $P$ has at most $(r+q)+1$ elements by
                    \cref{finite_at_most_adjoin_point_50}.
                $D_0\union\{y\}\subseteq D$ by
                    \cref{setminus_subseteq,singleton_subset_intro,union_subsets_is_subset}.
                Show that $D\subseteq D_0\union\{y\}$.
                    \begin{subproof}
                        It suffices to show that for all $z\in D$ we have
                            $z\in D_0\union\{y\}$ by \cref{subseteq}.
                        Fix $z\in D$.
                        \begin{byCase}
                        \caseOf{$z=y$.}
                            $z\in\{y\}$ by \cref{singleton_iff}.
                            Follows by \cref{union_intro_right}.
                        \caseOf{$z\neq y$.}
                            $z\in D_0$ by
                                \cref{setminus_intro,singleton_iff}.
                            Follows by \cref{union_intro_left}.
                        \end{byCase}
                \end{subproof}
                $D_0\union\{y\}=D$ by
                    \cref{subseteq_antisymmetric}.
                $P=C\union D$ by
                    \cref{union_assoc,real_lt_subst_left,real_lt_subst_right}.
                $(r+q)+1=r+(q+1)$ by
                    \cref{real_naturals_subset_reals,real_one_in_naturals,subseteq,real_add_assoc}.
                Follows by
                    \cref{real_lt_subst_left,real_lt_subst_right}.
            \end{byCase}
        \end{subproof}
        Follows by \cref{subseteq}.
    \end{subproof}
    For all $q\in\naturalsPlus$ we have $q\in S$ by
        \cref{real_naturals_induction,real_one_in_naturals}.
    $k\in S$ by \cref{subseteq}.
    Follows by \cref{subseteq}.
\end{proof}

% from §50 helper lemma 126: bijections preserve finite at-most bounds
\begin{lemma}\label{bijection_preserves_finite_at_most_50}
    Assume $A$ has at most $m$ elements.
    Assume $f$ is a bijection from $A$ to $B$.
    Then $B$ has at most $m$ elements.
\end{lemma}
\begin{proof}
    $m\in\naturalsPlus$ and $A$ is finite by
        \cref{finite_set_has_at_most_50}.
    $B$ is finite by \cref{finite_bijection}.
    Show that there exists no $C$ such that
        $C\subseteq B$ and
        there exists $g$ such that
            $g$ is a bijection from $\initialSegment{m+1}$ to $C$.
    \begin{subproof}
        Suppose not.
        Take $C,g$ such that
            $C\subseteq B$ and
            $g$ is a bijection from $\initialSegment{m+1}$ to $C$.
        $\converse{f}$ is a bijection from $B$ to $A$ by
            \cref{bijection_converse_is_bijection}.
        Let $v=\restrl{\converse{f}}{C}$.
        $\converse{f}$ is a function and $\dom{\converse{f}}=B$ by
            \cref{bijection_is_function,bijections_dom}.
        $v$ is a function by
            \cref{restrl_of_function_is_function}.
        $\dom{v}=C$ by
            \cref{dom_of_restrl_eq_inter,inter_absorb_supseteq_left}.
        Let $D=\img{\converse{f}}{C}$.
        For all $c\in\dom{v}$ we have $v(c)\in D$ by
            \cref{real_lt_subst_left,restrl_of_function_apply,subseteq,img_of_function_intro,inter_intro}.
        $v\in\funs{C}{D}$ by \cref{funs_intro}.
        $v$ is injective by
            \cref{restrl_of_function_apply,bijections,injective_function,subseteq}.
        $\img{v}{C}=D$ by
            \cref{restrl_img,inter_idempotent}.
        $\img{v}{C}=\ran{v}$ by
            \cref{img_dom,real_lt_subst_left}.
        $v\in\surj{C}{D}$ by
            \cref{funs_surj_iff,real_lt_subst_left}.
        $v$ is a bijection from $C$ to $D$ by
            \cref{bijections}.
        Let $q=v\circ g$.
        $q$ is a bijection from $\initialSegment{m+1}$ to $D$ by
            \cref{bijection_circ}.
        $D\subseteq A$ by
            \cref{img_subseteq_ran,bijections_to_funs,funs_ran,subseteq_transitive}.
        Contradiction by \cref{finite_set_has_at_most_50}.
    \end{subproof}
    Follows by \cref{finite_set_has_at_most_50}.
\end{proof}

% from §50 helper lemma 127: equal affine values give a zero affine difference relation
\begin{lemma}\label{affine_combination_difference_zero_50}
    Assume $N,n\in\naturalsPlus$.
    Let $I=\initialSegment{n}$.
    Let $E=\realsPower{\initialSegment{N}}$.
    Assume $z\in\funs{I}{E}$.
    Assume $a,b\in\funs{I}{\reals}$.
    Assume $\sumFiniteSequence{a}{1}$ and
        $\sumFiniteSequence{b}{1}$.
    Assume $y$ is an affine combination of $z$ with coefficients $a$
        in dimension $N$.
    Assume $y$ is an affine combination of $z$ with coefficients $b$
        in dimension $N$.
    Then there exists $d$ such that
        $d\in\funs{I}{\reals}$ and
        for all $i\in I$ we have $d(i)=a(i)+\neg{b(i)}$ and
        $\sumFiniteSequence{d}{0}$ and
        $\zeroVector{N}$ is an affine combination of $z$
            with coefficients $d$ in dimension $N$.
\end{lemma}
\begin{proof}
    $y\in E$ by \cref{affine_combination_value_50}.
    $1,\neg{1}\in\reals$ by
        \cref{real_one_in_naturals,real_naturals_subset_reals,subseteq,real_add_inverse}.
    Let $b_0=\coordScalarSequence{n}{\neg{1}}{b}$.
    Let $d=\coordSumSequence{n}{a}{b_0}$.
    $b_0,d\in\funs{I}{\reals}$ by
        \cref{coord_scalar_sequence_function,coord_sum_sequence_function}.
    $\sumFiniteSequence{b_0}{\neg{1}\rmul 1}$ by
        \cref{sum_finite_sequence_scalar}.
    $\neg{1}\rmul 1=\neg{1}$ by
        \cref{real_mul_ident,real_mul_comm}.
    $\sumFiniteSequence{b_0}{\neg{1}}$ by
        \cref{real_lt_subst_right}.
    $\sumFiniteSequence{d}{1+\neg{1}}$ by
        \cref{sum_finite_sequence_add}.
    $1+\neg{1}=0$ by \cref{real_add_inverse}.
    $\sumFiniteSequence{d}{0}$ by
        \cref{real_lt_subst_right}.
    For all $i\in I$ we have
        $d(i)=a(i)+\neg{b(i)}$ by
        \cref{coord_scalar_sequence,coord_sum_sequence,funs_is_function,function_apply_intro,funs_type_apply,real_neg_as_mul,real_lt_subst_left,real_lt_subst_right}.
    $\zeroVector{N}\in E$ by
        \cref{zero_vector_euclidean_50}.
    Show that for all $j\in\initialSegment{N}$ there exists $q$ such that
            $q\in\funs{I}{\reals}$ and
            for all $i\in I$ we have
                $q(i)=d(i)\rmul\apply{z(i)}{j}$ and
            $\sumFiniteSequence{q}{\apply{\zeroVector{N}}{j}}$.
        \begin{subproof}
            Fix $j\in\initialSegment{N}$.
            Take $u$ such that
                $u\in\funs{I}{\reals}$ and
                for all $i\in I$ we have
                    $u(i)=a(i)\rmul\apply{z(i)}{j}$ and
                $\sumFiniteSequence{u}{y(j)}$ by
                \cref{affine_combination_value_50,funs_elim,real_lt_subst_left,real_lt_subst_right}.
            Take $v$ such that
                $v\in\funs{I}{\reals}$ and
                for all $i\in I$ we have
                    $v(i)=b(i)\rmul\apply{z(i)}{j}$ and
                $\sumFiniteSequence{v}{y(j)}$ by
                \cref{affine_combination_value_50,funs_elim,real_lt_subst_left,real_lt_subst_right}.
            $1\in I$ by
                \cref{real_one_in_naturals,real_one_leq_naturalsplus,initial_segment_naturalsplus}.
            $\sumFiniteSequence{u}{y(j)}$ and
                $\sumFiniteSequence{v}{y(j)}$ by \cref{subseteq}.
            Let $v_0=\coordScalarSequence{n}{\neg{1}}{v}$.
            Let $q=\coordSumSequence{n}{u}{v_0}$.
            $v_0,q\in\funs{I}{\reals}$ by
                \cref{coord_scalar_sequence_function,coord_sum_sequence_function}.
            $\sumFiniteSequence{v_0}{\neg{1}\rmul y(j)}$ by
                \cref{sum_finite_sequence_scalar}.
            $y(j)\in\reals$ by
                \cref{reals_power,tuple_power,funs_type_apply}.
            $\neg{1}\rmul y(j)=\neg{y(j)}$ by
                \cref{real_neg_as_mul}.
            $\sumFiniteSequence{v_0}{\neg{y(j)}}$ by
                \cref{real_lt_subst_right}.
            $\sumFiniteSequence{q}{y(j)+\neg{y(j)}}$ by
                \cref{sum_finite_sequence_add}.
            $y(j)+\neg{y(j)}=0$ by \cref{real_add_inverse}.
            $\sumFiniteSequence{q}{0}$ by
                \cref{real_lt_subst_right}.
            $(j,0)\in\zeroVector{N}$ by
                \cref{zero_vector_50,pair_eq_iff}.
            $\apply{\zeroVector{N}}{j}=0$ by
                \cref{zero_vector_euclidean_50,reals_power,tuple_power,funs_is_function,function_apply_intro}.
            $\sumFiniteSequence{q}{\apply{\zeroVector{N}}{j}}$ by
                \cref{real_lt_subst_right}.
            Show that for all $i\in I$ we have
                $q(i)=d(i)\rmul\apply{z(i)}{j}$.
            \begin{subproof}
                Fix $i\in I$.
                $u(i)=a(i)\rmul\apply{z(i)}{j}$ and
                    $v(i)=b(i)\rmul\apply{z(i)}{j}$ by \cref{subseteq}.
                $b_0(i)=\neg{1}\rmul b(i)$ and
                    $d(i)=a(i)+b_0(i)$ and
                    $v_0(i)=\neg{1}\rmul v(i)$ and
                    $q(i)=u(i)+v_0(i)$ by
                    \cref{coord_scalar_sequence,coord_sum_sequence,funs_is_function,function_apply_intro}.
                $a(i),b(i),\apply{z(i)}{j}\in\reals$ by
                    \cref{funs_type_apply,reals_power,tuple_power}.
                $(a(i)+(\neg{1}\rmul b(i)))\rmul\apply{z(i)}{j}
                    =(a(i)\rmul\apply{z(i)}{j})
                        +((\neg{1}\rmul b(i))\rmul\apply{z(i)}{j})$ by
                    \cref{real_right_distrib,real_add_inverse,real_mul_closed}.
                $(\neg{1}\rmul b(i))\rmul\apply{z(i)}{j}
                    =\neg{1}\rmul(b(i)\rmul\apply{z(i)}{j})$ by
                    \cref{real_mul_assoc}.
                Follows by
                    \cref{real_lt_subst_left,real_lt_subst_right}.
            \end{subproof}
            Follows.
    \end{subproof}
    $\zeroVector{N}$ is an affine combination of $z$
        with coefficients $d$ in dimension $N$ by
        \cref{affine_combination_value_50}.
    Follows.
\end{proof}

% from §50 helper lemma 128: a zero affine relation restricts to its nonzero support
\begin{lemma}\label{zero_affine_relation_support_restriction_50}
    Assume $N,n\in\naturalsPlus$.
    Let $I=\initialSegment{n}$.
    Let $E=\realsPower{\initialSegment{N}}$.
    Assume $z\in\funs{I}{E}$.
    Assume $d\in\funs{I}{\reals}$.
    Assume $\sumFiniteSequence{d}{0}$.
    Assume $\zeroVector{N}$ is an affine combination of $z$
        with coefficients $d$ in dimension $N$.
    Let $K=\{i\in I\mid d(i)\neq 0\}$.
    Assume $K$ is inhabited.
    Then there exist $k,h,s,c$ such that
        $k\in\naturalsPlus$ and
        $h$ is a bijection from $\initialSegment{k}$ to $K$ and
        $s\in\funs{\initialSegment{k}}{E}$ and
        $c\in\funs{\initialSegment{k}}{\reals}$ and
        for all $t\in\initialSegment{k}$ we have
            $s(t)=z(h(t))$ and $c(t)=d(h(t))$ and
        $\sumFiniteSequence{c}{0}$ and
        $\zeroVector{N}$ is an affine combination of $s$
            with coefficients $c$ in dimension $N$.
\end{lemma}
\begin{proof}
    $I$ is finite and $I$ is inhabited by
        \cref{initial_segment_finite,initial_segment_inhabited}.
    $K\subseteq I$ by \cref{subseteq}.
    $K$ is finite by \cref{subseteq_finite_is_finite}.
    Take $R$ such that $R$ is a well order relation on $I$ by
        \cref{well_ordering_theorem_local_finiteness}.
    $R$ is a simple order relation on $I$ by
        \cref{well_order_relation}.
    Take $n_0,g$ such that
        $n_0\in\naturalsPlus$ and
        $g$ is a bijection from $\initialSegment{n_0}$ to $I$ and
        $g$ is increasing on $\initialSegment{n_0}$ with respect to $R$ by
        \cref{subseteq_refl,finite_well_ordered_enumeration}.
    $g$ is a well-ordered enumeration of $I$ in $I$ with respect to
        $R$ and $n_0$ by
        \cref{well_ordered_enumeration_of_subset,subseteq_refl}.
    Take $k,h$ such that
        $k\in\naturalsPlus$ and
        $h$ is a bijection from $\initialSegment{k}$ to $K$ and
        $h$ is increasing on $\initialSegment{k}$ with respect to $R$ by
        \cref{subseteq,finite_well_ordered_enumeration}.
    Let $L=\initialSegment{k}$.
    $h$ is a well-ordered enumeration of $K$ in $I$ with respect to
        $R$ and $k$ by
        \cref{well_ordered_enumeration_of_subset}.
    Let $s=\{(t,z(h(t)))\mid t\in L\}$.
    Let $c=\{(t,d(h(t)))\mid t\in L\}$.
    $s$ is a relation and $s$ is right-unique by
        \cref{relation,rightunique,pair_eq_iff}.
    $s$ is a function by \cref{relation,rightunique}.
    $\dom{s}=L$ by
        \cref{dom_intro,dom_iff,pair_eq_iff,subseteq,subseteq_antisymmetric}.
    For all $t\in\dom{s}$ we have $s(t)\in E$ by
        \cref{real_lt_subst_left,bijections_to_funs,funs_type_apply,subseteq,pair_eq_iff,function_apply_intro}.
    $s\in\funs{L}{E}$ by \cref{funs_intro}.
    $c$ is a relation and $c$ is right-unique by
        \cref{relation,rightunique,pair_eq_iff}.
    $c$ is a function by \cref{relation,rightunique}.
    $\dom{c}=L$ by
        \cref{dom_intro,dom_iff,pair_eq_iff,subseteq,subseteq_antisymmetric}.
    For all $t\in\dom{c}$ we have $c(t)\in\reals$ by
        \cref{real_lt_subst_left,bijections_to_funs,funs_type_apply,subseteq,pair_eq_iff,function_apply_intro}.
    $c\in\funs{L}{\reals}$ by \cref{funs_intro}.
    For all $t\in L$ we have
        $s(t)=z(h(t))$ and $c(t)=d(h(t))$ by
        \cref{pair_eq_iff,funs_is_function,function_apply_intro}.
    Let $F=\{(\alpha,\{(0,d(\alpha))\})\mid\alpha\in I\}$.
    $F$ is a relation and $F$ is right-unique by
        \cref{relation,rightunique,pair_eq_iff}.
    $F$ is a function by \cref{relation,rightunique}.
    $\dom{F}=I$ by
        \cref{dom_intro,dom_iff,pair_eq_iff,subseteq,subseteq_antisymmetric}.
    Show that for all $\alpha\in I$ we have
        $F(\alpha)\in\funs{\{0\}}{\reals}$ and
        $\apply{F(\alpha)}{0}=d(\alpha)$.
    \begin{subproof}
        Fix $\alpha\in I$.
        $F(\alpha)=\{(0,d(\alpha))\}$ by
            \cref{pair_eq_iff,function_apply_intro}.
        Let $q=\{(0,d(\alpha))\}$.
        $q$ is a relation and $q$ is right-unique by
            \cref{relation,rightunique,singleton_elim,pair_eq_iff}.
        $q$ is a function by \cref{relation,rightunique}.
        $\dom{q}=\{0\}$ by
            \cref{dom_intro,dom_iff,singleton_elim,pair_eq_iff,singleton_intro,subseteq,subseteq_antisymmetric}.
        $d(\alpha)\in\reals$ by \cref{funs_type_apply}.
        For all $u\in\dom{q}$ we have $q(u)\in\reals$ by
            \cref{singleton_iff,function_apply_intro}.
        $q\in\funs{\{0\}}{\reals}$ by \cref{funs_intro}.
        Follows by
            \cref{singleton_intro,function_apply_intro,real_lt_subst_left,real_lt_subst_right}.
    \end{subproof}
    Let $d_g=\{(t,d(g(t)))\mid t\in\initialSegment{n_0}\}$.
    $d_g$ is a relation and $d_g$ is right-unique by
        \cref{relation,rightunique,pair_eq_iff}.
    $d_g$ is a function by \cref{relation,rightunique}.
    $\dom{d_g}=\initialSegment{n_0}$ by
        \cref{dom_intro,dom_iff,pair_eq_iff,subseteq,subseteq_antisymmetric}.
    For all $t\in\dom{d_g}$ we have $d_g(t)\in\reals$ by
        \cref{real_lt_subst_left,bijections_to_funs,funs_type_apply,pair_eq_iff,function_apply_intro}.
    $d_g\in\funs{\initialSegment{n_0}}{\reals}$ by
        \cref{funs_intro}.
    Take $r\in\reals$ such that $\sumFiniteSequence{d_g}{r}$ by
        \cref{sum_finite_sequence_exists}.
    Let $j_0=\identity{I}$.
    $j_0$ is a bijection from $I$ to $I$ by
        \cref{id_elem_bijections}.
    For all $i\in I$ we have $d(i)=d(j_0(i))$ by
        \cref{id_apply}.
    For all $t\in\initialSegment{n_0}$ we have
        $d_g(t)=d(g(t))$ by
        \cref{pair_eq_iff,function_apply_intro}.
    $0=r$ by
        \cref{finite_sum_reindex_bijection_50}.
    $\sumFiniteSequence{d_g}{0}$ by
        \cref{real_lt_subst_right}.
    For all $t\in\initialSegment{n_0}$ we have
        $d_g(t)=\apply{\apply{F}{g(t)}}{0}$ by
        \cref{bijections_to_funs,funs_type_apply,subseteq,pair_eq_iff,function_apply_intro,real_lt_subst_left,real_lt_subst_right}.
    $d_g$ is a real sum sequence of $F$ at $0$ along $g$
        with respect to $n_0$ and result $0$ by
        \cref{real_sum_sequence_along_enumeration}.
    Take $r_0\in\reals$ such that $\sumFiniteSequence{c}{r_0}$ by
        \cref{sum_finite_sequence_exists}.
    For all $t\in L$ we have
        $c(t)=\apply{\apply{F}{h(t)}}{0}$ by
        \cref{bijections_to_funs,funs_type_apply,subseteq,pair_eq_iff,function_apply_intro,real_lt_subst_left,real_lt_subst_right}.
    $c$ is a real sum sequence of $F$ at $0$ along $h$
        with respect to $k$ and result $r_0$ by
        \cref{real_sum_sequence_along_enumeration}.
    For all $\alpha\in I\setminus K$ we have
        $\apply{\apply{F}{\alpha}}{0}=0$ by
        \cref{setminus_elim_left,setminus_elim_right,subseteq,real_lt_subst_right}.
    $0=r_0$ by
        \cref{well_ordered_zero_extension_sum_comparison}.
    $\sumFiniteSequence{c}{0}$ by
        \cref{real_lt_subst_right}.
    $\zeroVector{N}\in E$ by
        \cref{zero_vector_euclidean_50}.
    Show that for all $j\in\initialSegment{N}$ there exists $w$ such that
            $w\in\funs{L}{\reals}$ and
            for all $t\in L$ we have
                $w(t)=c(t)\rmul\apply{s(t)}{j}$ and
            $\sumFiniteSequence{w}{\apply{\zeroVector{N}}{j}}$.
        \begin{subproof}
            Fix $j\in\initialSegment{N}$.
            Take $q$ such that
                $q\in\funs{I}{\reals}$ and
                for all $i\in I$ we have
                    $q(i)=d(i)\rmul\apply{z(i)}{j}$ and
                $\sumFiniteSequence{q}{\apply{\zeroVector{N}}{j}}$ by
                \cref{affine_combination_value_50,funs_elim,real_lt_subst_left,real_lt_subst_right}.
            $1\in I$ by
                \cref{real_one_in_naturals,real_one_leq_naturalsplus,initial_segment_naturalsplus}.
            $\sumFiniteSequence{q}{\apply{\zeroVector{N}}{j}}$ by
                \cref{subseteq}.
            $\apply{\zeroVector{N}}{j}=0$ by
                \cref{zero_vector_euclidean_50,zero_vector_50,reals_power,tuple_power,funs_is_function,function_apply_intro,pair_eq_iff}.
            $\sumFiniteSequence{q}{0}$ by
                \cref{real_lt_subst_right}.
            Let $q_g=\{(t,q(g(t)))\mid t\in\initialSegment{n_0}\}$.
            $q_g$ is a relation and $q_g$ is right-unique by
                \cref{relation,rightunique,pair_eq_iff}.
            $q_g$ is a function by \cref{relation,rightunique}.
            $\dom{q_g}=\initialSegment{n_0}$ by
                \cref{dom_intro,dom_iff,pair_eq_iff,subseteq,subseteq_antisymmetric}.
            For all $t\in\dom{q_g}$ we have $q_g(t)\in\reals$ by
                \cref{real_lt_subst_left,bijections_to_funs,funs_type_apply,pair_eq_iff,function_apply_intro}.
            $q_g\in\funs{\initialSegment{n_0}}{\reals}$ by
                \cref{funs_intro}.
            Take $v\in\reals$ such that $\sumFiniteSequence{q_g}{v}$ by
                \cref{sum_finite_sequence_exists}.
            For all $i\in I$ we have $q(i)=q(j_0(i))$ by
                \cref{id_apply}.
            For all $t\in\initialSegment{n_0}$ we have
                $q_g(t)=q(g(t))$ by
                \cref{pair_eq_iff,function_apply_intro}.
            $0=v$ by
                \cref{finite_sum_reindex_bijection_50}.
            $\sumFiniteSequence{q_g}{0}$ by
                \cref{real_lt_subst_right}.
            Let $w=\{(t,q(h(t)))\mid t\in L\}$.
            $w$ is a relation and $w$ is right-unique by
                \cref{relation,rightunique,pair_eq_iff}.
            $w$ is a function by \cref{relation,rightunique}.
            $\dom{w}=L$ by
                \cref{dom_intro,dom_iff,pair_eq_iff,subseteq,subseteq_antisymmetric}.
            For all $t\in\dom{w}$ we have $w(t)\in\reals$ by
                \cref{real_lt_subst_left,bijections_to_funs,funs_type_apply,subseteq,pair_eq_iff,function_apply_intro}.
            $w\in\funs{L}{\reals}$ by \cref{funs_intro}.
            For all $t\in L$ we have
                $w(t)=c(t)\rmul\apply{s(t)}{j}$ by
                \cref{bijections_to_funs,funs_type_apply,subseteq,pair_eq_iff,funs_is_function,function_apply_intro,real_lt_subst_left,real_lt_subst_right}.
            Take $v_0\in\reals$ such that $\sumFiniteSequence{w}{v_0}$ by
                \cref{sum_finite_sequence_exists}.
            Let $F_j=\{(\alpha,\{(0,q(\alpha))\})\mid\alpha\in I\}$.
            $F_j$ is a relation and $F_j$ is right-unique by
                \cref{relation,rightunique,pair_eq_iff}.
            $F_j$ is a function by \cref{relation,rightunique}.
            $\dom{F_j}=I$ by
                \cref{dom_intro,dom_iff,pair_eq_iff,subseteq,subseteq_antisymmetric}.
            Show that for all $\alpha\in I$ we have
                $F_j(\alpha)\in\funs{\{0\}}{\reals}$ and
                $\apply{F_j(\alpha)}{0}=q(\alpha)$.
            \begin{subproof}
                Fix $\alpha\in I$.
                $F_j(\alpha)=\{(0,q(\alpha))\}$ by
                    \cref{pair_eq_iff,function_apply_intro}.
                Let $e_0=\{(0,q(\alpha))\}$.
                $e_0$ is a relation and $e_0$ is right-unique by
                    \cref{relation,rightunique,singleton_elim,pair_eq_iff}.
                $e_0$ is a function by \cref{relation,rightunique}.
                $\dom{e_0}=\{0\}$ by
                    \cref{dom_intro,dom_iff,singleton_elim,pair_eq_iff,singleton_intro,subseteq,subseteq_antisymmetric}.
                $q(\alpha)\in\reals$ by \cref{funs_type_apply}.
                For all $u\in\dom{e_0}$ we have $e_0(u)\in\reals$ by
                    \cref{singleton_iff,function_apply_intro}.
                $e_0\in\funs{\{0\}}{\reals}$ by \cref{funs_intro}.
                Follows by
                    \cref{singleton_intro,function_apply_intro,real_lt_subst_left,real_lt_subst_right}.
            \end{subproof}
            For all $t\in\initialSegment{n_0}$ we have
                $q_g(t)=\apply{\apply{F_j}{g(t)}}{0}$ by
                \cref{bijections_to_funs,funs_type_apply,pair_eq_iff,function_apply_intro,subseteq,real_lt_subst_left,real_lt_subst_right}.
            $q_g$ is a real sum sequence of $F_j$ at $0$ along $g$
                with respect to $n_0$ and result $0$ by
                \cref{real_sum_sequence_along_enumeration}.
            For all $t\in L$ we have
                $w(t)=\apply{\apply{F_j}{h(t)}}{0}$ by
                \cref{bijections_to_funs,funs_type_apply,subseteq,pair_eq_iff,function_apply_intro,real_lt_subst_left,real_lt_subst_right}.
            $w$ is a real sum sequence of $F_j$ at $0$ along $h$
                with respect to $k$ and result $v_0$ by
                \cref{real_sum_sequence_along_enumeration}.
            For all $\alpha\in I\setminus K$ we have
                $\apply{\apply{F_j}{\alpha}}{0}=0$ by
                \cref{setminus_elim_left,setminus_elim_right,subseteq,funs_type_apply,reals_power,tuple_power,real_mul_zero,real_lt_subst_left,real_lt_subst_right}.
            $0=v_0$ by
                \cref{well_ordered_zero_extension_sum_comparison}.
            $\sumFiniteSequence{w}{0}$ by
                \cref{real_lt_subst_right}.
            $\sumFiniteSequence{w}{\apply{\zeroVector{N}}{j}}$ by
                \cref{real_lt_subst_right}.
            Follows.
    \end{subproof}
    $\zeroVector{N}$ is an affine combination of $s$
        with coefficients $c$ in dimension $N$ by
        \cref{affine_combination_value_50}.
    For all $t\in L$ we have
        $s(t)=z(h(t))$ and $c(t)=d(h(t))$ and
        $\sumFiniteSequence{c}{0}$ and
        $\zeroVector{N}$ is an affine combination of $s$
            with coefficients $c$ in dimension $N$ by \cref{subseteq}.
    Follows by
        \cref{real_lt_subst_left,real_lt_subst_right}.
\end{proof}

% from §50 helper lemma 129: sparse affine coordinates on general-position points are unique
\begin{lemma}\label{sparse_general_position_affine_coordinates_unique_50}
    Assume $N,n\in\naturalsPlus$.
    Let $I=\initialSegment{n}$.
    Let $E=\realsPower{\initialSegment{N}}$.
    Assume $z\in\funs{I}{E}$.
    Assume $z$ is injective.
    Assume $\img{z}{I}$ is in general position in dimension $N$.
    Assume $a,b\in\funs{I}{\reals}$.
    Assume $\sumFiniteSequence{a}{1}$ and
        $\sumFiniteSequence{b}{1}$.
    Assume $y$ is an affine combination of $z$ with coefficients $a$
        in dimension $N$.
    Assume $y$ is an affine combination of $z$ with coefficients $b$
        in dimension $N$.
    Let $D=\{i\in I\mid a(i)\neq b(i)\}$.
    Assume $D$ has at most $N+1$ elements.
    Then $a=b$.
\end{lemma}
\begin{proof}
    Take $d$ such that
        $d\in\funs{I}{\reals}$ and
        for all $i\in I$ we have $d(i)=a(i)+\neg{b(i)}$ and
        $\sumFiniteSequence{d}{0}$ and
        $\zeroVector{N}$ is an affine combination of $z$
            with coefficients $d$ in dimension $N$ by
        \cref{affine_combination_difference_zero_50}.
    $1\in I$ by
        \cref{real_one_in_naturals,real_one_leq_naturalsplus,initial_segment_naturalsplus}.
    $\sumFiniteSequence{d}{0}$ and
        $\zeroVector{N}$ is an affine combination of $z$
            with coefficients $d$ in dimension $N$ by
        \cref{subseteq}.
    Let $K=\{i\in I\mid d(i)\neq 0\}$.
    $K\subseteq D$ by
        \cref{subseteq,funs_type_apply,real_add_inverse,real_lt_subst_left,real_lt_subst_right}.
    $K$ has at most $N+1$ elements by
        \cref{finite_at_most_subset_50}.
    \begin{byCase}
    \caseOf{$K=\emptyset$.}
        For all $i\in I$ we have $a(i)=b(i)$ by
            \cref{emptyset,subseteq,funs_type_apply,real_add_inverse,real_lt_subst_left,real_lt_subst_right,real_add_cancel}.
        Follows by \cref{functions_eq_iff}.
    \caseOf{$K\neq\emptyset$.}
        Suppose not.
        $K$ is inhabited by \cref{nonempty_inhabited}.
        Take $k,h,s,c$ such that
            $k\in\naturalsPlus$ and
            $h$ is a bijection from $\initialSegment{k}$ to $K$ and
            $s\in\funs{\initialSegment{k}}{E}$ and
            $c\in\funs{\initialSegment{k}}{\reals}$ and
            for all $t\in\initialSegment{k}$ we have
                $s(t)=z(h(t))$ and $c(t)=d(h(t))$ and
                $\sumFiniteSequence{c}{0}$ and
                $\zeroVector{N}$ is an affine combination of $s$
                    with coefficients $c$ in dimension $N$ by
            \cref{zero_affine_relation_support_restriction_50}.
        Let $L=\initialSegment{k}$.
        $1\in L$ by
            \cref{real_one_in_naturals,real_one_leq_naturalsplus,initial_segment_naturalsplus}.
        $\sumFiniteSequence{c}{0}$ and
            $\zeroVector{N}$ is an affine combination of $s$
                with coefficients $c$ in dimension $N$ by
            \cref{subseteq}.
        $s$ is injective by
            \cref{funs_elim,real_lt_subst_left,bijections_to_funs,funs_type_apply,subseteq,real_lt_subst_right,injective_function,bijections}.
        Let $A=\img{s}{L}$.
        $s$ is a bijection from $L$ to $A$ by
            \cref{funs_elim,real_lt_subst_left,inter_intro,img_of_function_intro,function_apply_elim,funs_intro,img_dom,real_lt_subst_right,funs_surj_iff,bijections}.
        $\converse{h}$ is a bijection from $K$ to $L$ by
            \cref{bijection_converse_is_bijection}.
        Let $r=s\circ\converse{h}$.
        $r$ is a bijection from $K$ to $A$ by
            \cref{bijection_circ}.
        $A$ has at most $N+1$ elements by
            \cref{bijection_preserves_finite_at_most_50}.
        $A\subseteq\img{z}{I}$ by
            \cref{subseteq,img_of_function_elim,inter_elim_right,bijections_to_funs,funs_type_apply,funs_elim,inter_intro,img_of_function_intro,real_lt_subst_left,real_lt_subst_right}.
        $A$ is a geometrically independent set in dimension $N$ by
            \cref{general_position_50}.
        $s(1)\in A$ by
            \cref{bijections_to_funs,funs_elim,inter_intro,img_of_function_intro}.
        $A$ is inhabited.
        $s$ is a geometrically independent sequence in dimension $N$ by
            \cref{independent_set_bijective_enumeration_50}.
        Take $l$ such that
            $l\in\naturalsPlus$ and
            $s\in\funs{\initialSegment{l}}{E}$ and
            for all $e$ if
                $e\in\funs{\initialSegment{l}}{\reals}$ and
                $\zeroVector{N}$ is an affine combination of $s$
                    with coefficients $e$ in dimension $N$ and
                $\sumFiniteSequence{e}{0}$
            then for all $t\in\initialSegment{l}$ we have $e(t)=0$ by
            \cref{geometrically_independent_sequence_50}.
        $\initialSegment{l}=L$ by \cref{funs_elim}.
        For all $t\in L$ we have $c(t)=0$ by
            \cref{real_lt_subst_left,real_lt_subst_right}.
        Contradiction by
            \cref{bijections_to_funs,funs_type_apply,subseteq,real_lt_subst_left,real_lt_subst_right}.
    \end{byCase}
\end{proof}

% from §50 helper lemma 130: partition coefficients are nonnegative
\begin{lemma}\label{finite_partition_coefficients_nonnegative_50}
    Assume $n\in\naturalsPlus$.
    Let $I=\initialSegment{n}$.
    Assume $T$ is a topology on $X$.
    Assume $p$ is a partition of unity on $X$ dominated by $U$
        with respect to $T$ and $n$.
    Assume $x\in X$.
    Assume $a\in\funs{I}{\reals}$.
    Assume for all $i\in I$ we have
        $a(i)=\apply{\apply{p}{i}}{x}$.
    Then for all $i\in I$ we have $0\leq a(i)$.
\end{lemma}
\begin{proof}
    Follows by
        \cref{finite_partition_of_unity,finite_partition_component_condition,continuous,funs_type_apply,intervalclosed,subseteq,real_lt_subst_left,real_lt_subst_right}.
\end{proof}

% from §50 helper lemma 131: partition coefficients vanish outside their dominating set
\begin{lemma}\label{finite_partition_coefficient_zero_outside_50}
    Assume $n\in\naturalsPlus$.
    Let $I=\initialSegment{n}$.
    Assume $T$ is a topology on $X$.
    Assume $p$ is a partition of unity on $X$ dominated by $U$
        with respect to $T$ and $n$.
    Assume $i\in I$ and $x\in X$.
    Assume $x\notin U(i)$.
    Then $\apply{\apply{p}{i}}{x}=0$.
\end{lemma}
\begin{proof}
    Follows by
        \cref{finite_partition_of_unity,finite_partition_component_condition,continuous,funs_type_apply,intervalclosed,real_leq_split,positive_partition_component_in_support,subseteq}.
\end{proof}

% from §50 helper lemma 132: differing partition coordinates occur only in the two active index sets
\begin{lemma}\label{partition_coefficient_difference_support_50}
    Assume $n\in\naturalsPlus$.
    Let $I=\initialSegment{n}$.
    Assume $T$ is a topology on $X$.
    Assume $p$ is a partition of unity on $X$ dominated by $U$
        with respect to $T$ and $n$.
    Assume $x,y\in X$.
    Assume $a,b\in\funs{I}{\reals}$.
    Assume for all $i\in I$ we have
        $a(i)=\apply{\apply{p}{i}}{x}$ and
        $b(i)=\apply{\apply{p}{i}}{y}$.
    Let $A=\{i\in I\mid x\in U(i)\}$.
    Let $B=\{i\in I\mid y\in U(i)\}$.
    Let $D=\{i\in I\mid a(i)\neq b(i)\}$.
    Then $D\subseteq A\union B$.
\end{lemma}
\begin{proof}
    Follows by
        \cref{subseteq,union_iff,finite_partition_coefficient_zero_outside_50,real_lt_subst_left,real_lt_subst_right}.
\end{proof}

% from §50 helper lemma 133: the Menger--Nobeling target dimension is positive
\begin{lemma}\label{menger_nobeling_target_dimension_positive_50}
    Assume $m\in\naturals$.
    Let $N=(1+1)\rmul m+1$.
    Then $N\in\naturalsPlus$.
\end{lemma}
\begin{proof}
    Follows by
        \cref{naturals,union_iff,singleton_iff,real_one_in_naturals,real_naturals_subset_reals,subseteq,real_add_closed,real_add_ident,real_mul_zero,real_mul_comm,real_lt_subst_left,naturalsplus_add_closed,naturalsplus_mul_closed,real_naturals_closed_succ}.
\end{proof}

% from §50 helper lemma 134: two cover-order bounds add to the affine-independence threshold
\begin{lemma}\label{menger_nobeling_order_arithmetic_50}
    Assume $m\in\naturals$.
    Let $N=(1+1)\rmul m+1$.
    Then $(m+1)+(m+1)=N+1$.
\end{lemma}
\begin{proof}
    Follows by
        \cref{naturals,union_iff,singleton_iff,real_add_ident,real_naturals_subset_reals,real_one_in_naturals,subseteq,real_add_closed,real_distrib,real_mul_ident,real_mul_comm,real_lt_subst_left,real_lt_subst_right,real_add_assoc,real_add_comm}.
\end{proof}

% from §50 helper lemma 135: choose one point from every member of a finite indexed family
\begin{lemma}\label{finite_indexed_inhabited_family_choice_50}
    Assume $n\in\naturalsPlus$.
    Let $I=\initialSegment{n}$.
    Assume $U$ is a function on $I$.
    Assume for all $i\in I$ we have
        $U(i)\subseteq X$ and $U(i)$ is inhabited.
    Then there exists $c$ such that
        $c\in\funs{I}{X}$ and
        for all $i\in I$ we have $c(i)\in U(i)$.
\end{lemma}
\begin{proof}
    Let $R=\{w\in I\times X\mid \snd{w}\in U(\fst{w})\}$.
    $R$ is a relation by
        \cref{relation,subseteq,times_elem_is_tuple}.
    For all $i\in I$ there exists $x$ such that $i\mathrel{R}x$ by
        \cref{subseteq,times_tuple_intro,fst_eq,snd_eq,relation}.
    Take $c$ such that
        $c$ is a function on $I$ and
        for all $i\in I$ we have $i\mathrel{R}c(i)$ by
        \cref{choice_function}.
    $c\in\funs{I}{X}$ and
        for all $i\in I$ we have $c(i)\in U(i)$ by
        \cref{real_lt_subst_left,subseteq,fst_eq,snd_eq,times,funs_intro}.
    Follows.
\end{proof}

% from §50 helper lemma 136: every uniform ball contains a fixed-scale separating map
\begin{lemma}\label{menger_nobeling_separating_ball_intersection_50}
    Assume $m\in\naturals$.
    Assume $m$ is a topological dimension bound for $X$ with respect to $T$.
    Assume $d$ is a metric on $X$ and $T=\metricTopology{d}{X}$.
    Assume $X$ is compact with respect to $T$ and $X$ is inhabited.
    Let $N=(1+1)\rmul m+1$.
    Let $E=\realsPower{\initialSegment{N}}$.
    Let $e=\squareMetric{N}$.
    Let $S=\metricTopology{e}{E}$.
    Assume $C$ is the set of continuous functions from $X$ to $E$
        with respect to $T$ and $S$.
    Assume $\rho$ is the uniform metric on functions from $X$ to $E$
        with respect to $e$.
    Let $q=\restrl{\rho}{C\times C}$.
    Assume $q$ is a metric on $C$.
    Assume $\epsilon,\eta\in\reals$ and
        $0<\epsilon$ and $0<\eta$.
    Assume $W$ is the separating class at scale $\epsilon$ on $X$ to $E$
        with respect to $T$ and $S$ and $d$.
    Assume $f\in C$.
    Then $\metricBall{q}{C}{f}{\eta}$ intersects $W$.
\end{lemma}
\begin{proof}
    $N\in\naturalsPlus$ by
        \cref{menger_nobeling_target_dimension_positive_50}.
    $e$ is a metric on $E$ by \cref{square_metric_is_metric}.
    $f$ is continuous from $X$ to $E$ with respect to $T$ and $S$ by
        \cref{continuous_function_set,subseteq}.
    $f$ is uniformly continuous from $X$ to $E$ with respect to $d$ and $e$ by
        \cref{uniform_continuity_theorem}.
    Let $a=\eta/(1+1)$.
    $a\in\reals$ and $0<a$ and $a+a=\eta$ by
        \cref{real_half_of_positive}.
    $a<\eta$ by
        \cref{real_add_ident,real_lt_add,real_lt_subst_left,real_lt_subst_right}.
    $1\in\reals$ and $0<1$ by
        \cref{real_one_in_naturals,real_naturals_subset_reals,subseteq,real_zero_lt_one}.
    Take $r\in\reals$ such that
        $0<r$ and $r\leq a$ and $r\leq 1$ by
        \cref{positive_radius_below_two_bounds_49}.
    $r<\eta$ by \cref{real_leq_lt_trans}.
    Let $b=r/(1+1)$.
    $b\in\reals$ and $0<b$ and $b+b=r$ by
        \cref{real_half_of_positive}.
    Take $n,U$ such that
        $n\in\naturalsPlus$ and
        $U$ is a bijection from $\initialSegment{n}$ to
            $\img{U}{\initialSegment{n}}$ and
        $\img{U}{\initialSegment{n}}$ is a covering of $X$ and
        $\img{U}{\initialSegment{n}}$ has order at most $m+1$ in $X$ and
        for all $i\in\initialSegment{n}$ we have
            $U(i)$ is open in $X$ with respect to $T$ and
            $U(i)$ is inhabited and
            for all $x,y\in U(i)$ we have
                $d((x,y))<\epsilon$ and
                $e((f(x),f(y)))<b$ by
        \cref{controlled_finite_indexed_cover_50}.
    Let $I=\initialSegment{n}$.
    $U$ is a function on $I$ by
        \cref{bijections_to_funs,funs_elim}.
    For all $i\in I$ we have
        $U(i)\subseteq X$ and $U(i)$ is inhabited by
        \cref{subseteq,open_in,topology_on,powerset_elim}.
    Take $c$ such that
        $c\in\funs{I}{X}$ and
        for all $i\in I$ we have $c(i)\in U(i)$ by
        \cref{finite_indexed_inhabited_family_choice_50}.
    Let $v=f\circ c$.
    $f\in\funs{X}{E}$ by \cref{continuous}.
    $v\in\funs{I}{E}$ by \cref{funs_circ}.
    For all $i\in I$ we have $v(i)=f(c(i))$ by
        \cref{funs_elim,funs_ran,circ_apply}.
    Take $z$ such that
        $z\in\funs{I}{E}$ and
        $z$ is injective and
        $\img{z}{I}$ is in general position in dimension $N$ and
        for all $i\in I$ we have
            $e((v(i),z(i)))<b$ by
        \cref{injective_general_position_approximation_50}.
    $T$ is a topology on $X$ by
        \cref{topological_dimension_bound_50}.
    $X$ is Hausdorff with respect to $T$ by
        \cref{metric_space_hausdorff}.
    $X$ is normal with respect to $T$ by
        \cref{compact_hausdorff_normal}.
    For all $i\in I$ we have
        $U(i)$ is open in $X$ with respect to $T$ by \cref{subseteq}.
    Take $p$ such that
        $p$ is a partition of unity on $X$ dominated by $U$
            with respect to $T$ and $n$ by
        \cref{finite_partition_of_unity_existence}.
    Take $g$ such that
        $g$ is continuous from $X$ to $E$ with respect to $T$ and $S$ and
        for all $x\in X$ there exists $A_x$ such that
            $A_x\in\funs{I}{\reals}$ and
            for all $i\in I$ we have
                $A_x(i)=\apply{\apply{p}{i}}{x}$ and
                $\sumFiniteSequence{A_x}{1}$ and
                $g(x)$ is an affine combination of $z$
                    with coefficients $A_x$ in dimension $N$ by
        \cref{finite_partition_barycentric_map_50}.
    $g\in C$ by
        \cref{continuous_function_set,continuous,funs,subseteq}.
    Show that for all $x\in X$ we have
        $e((f(x),g(x)))\leq r$.
    \begin{subproof}
        Fix $x\in X$.
        Take $A_x$ such that
            $A_x\in\funs{I}{\reals}$ and
            for all $i\in I$ we have
                $A_x(i)=\apply{\apply{p}{i}}{x}$ and
                $\sumFiniteSequence{A_x}{1}$ and
                $g(x)$ is an affine combination of $z$
                    with coefficients $A_x$ in dimension $N$ by
            \cref{subseteq}.
        $1\in I$ by
            \cref{real_one_in_naturals,real_one_leq_naturalsplus,initial_segment_naturalsplus}.
        $\sumFiniteSequence{A_x}{1}$ and
            $g(x)$ is an affine combination of $z$
                with coefficients $A_x$ in dimension $N$ by
            \cref{subseteq}.
        For all $i\in I$ we have $0\leq A_x(i)$ by
            \cref{finite_partition_coefficients_nonnegative_50}.
        Show that for all $i\in I$ if $0<A_x(i)$ then
            $e((z(i),f(x)))\leq r$.
        \begin{subproof}
            Fix $i\in I$.
            Assume $0<A_x(i)$.
            $A_x(i)=\apply{\apply{p}{i}}{x}$ by \cref{subseteq}.
            $0<\apply{\apply{p}{i}}{x}$ by
                \cref{real_lt_subst_right}.
            $x\in\supp{\apply{p}{i}}{X}{T}$ by
                \cref{positive_partition_component_in_support}.
            $\supp{\apply{p}{i}}{X}{T}\subseteq U(i)$ by
                \cref{finite_partition_of_unity,finite_partition_component_condition}.
            $x\in U(i)$ by \cref{subseteq}.
            $c(i)\in U(i)$ by \cref{subseteq}.
            $e((f(c(i)),f(x)))<b$ by \cref{subseteq}.
            $v(i)=f(c(i))$ by \cref{subseteq}.
            $e((v(i),f(x)))<b$ by
                \cref{real_lt_subst_left}.
            $e((v(i),z(i)))<b$ by \cref{subseteq}.
            $v(i),z(i),f(x)\in E$ by
                \cref{funs_type_apply,continuous}.
            Follows by
                \cref{metric_on,metric_symmetric,real_lt_subst_left,metric_triangle_inequality,funs_type_apply,times_tuple_intro,real_add_closed,real_lt_add_two,real_leq_lt_trans,real_lt_subst_right,real_lt_imp_leq}.
        \end{subproof}
        Follows by
            \cref{continuous,funs_type_apply,real_lt_imp_leq,barycentric_square_metric_bound_50,metric_on,metric_symmetric,real_leq_subst_left_local}.
    \end{subproof}
    $f,g\in\funs{X}{E}$ by
        \cref{continuous}.
    $\rho((f,g))\leq r$ by
        \cref{uniform_metric_leq_from_pointwise_bound_50}.
    $(f,g)\in C\times C$ by \cref{times_tuple_intro}.
    $q((f,g))=\rho((f,g))$ by
        \cref{continuous_function_set,subseteq,times_subseteq_left,times_subseteq_right,subseteq_transitive,uniform_metric_on_function_space,metric_on,funs_elim,restrl_of_function_apply}.
    $q((f,g))\leq r$ by
        \cref{real_leq_subst_left_local}.
    $q((f,g))<\eta$ by
        \cref{metric_on,funs_type_apply,times_tuple_intro,real_leq_lt_trans}.
    $g\in\metricBall{q}{C}{f}{\eta}$ by
        \cref{metric_ball}.
    Show that for all $x\in X$ we have for all $y\in X$ if
        $\epsilon\leq d((x,y))$
    then $g(x)\neq g(y)$.
    \begin{subproof}
        Fix $x\in X$.
        Fix $y\in X$.
        Assume $\epsilon\leq d((x,y))$.
        Suppose not.
        $g(x)=g(y)$.
        Take $A_x$ such that
            $A_x\in\funs{I}{\reals}$ and
            for all $i\in I$ we have
                $A_x(i)=\apply{\apply{p}{i}}{x}$ and
                $\sumFiniteSequence{A_x}{1}$ and
                $g(x)$ is an affine combination of $z$
                    with coefficients $A_x$ in dimension $N$ by
            \cref{subseteq}.
        Take $A_y$ such that
            $A_y\in\funs{I}{\reals}$ and
            for all $i\in I$ we have
                $A_y(i)=\apply{\apply{p}{i}}{y}$ and
                $\sumFiniteSequence{A_y}{1}$ and
                $g(y)$ is an affine combination of $z$
                    with coefficients $A_y$ in dimension $N$ by
            \cref{subseteq}.
        $1\in I$ by
            \cref{real_one_in_naturals,real_one_leq_naturalsplus,initial_segment_naturalsplus}.
        $\sumFiniteSequence{A_x}{1}$ and
            $g(x)$ is an affine combination of $z$
                with coefficients $A_x$ in dimension $N$ by
            \cref{subseteq}.
        $\sumFiniteSequence{A_y}{1}$ and
            $g(y)$ is an affine combination of $z$
                with coefficients $A_y$ in dimension $N$ by
            \cref{subseteq}.
        $g(x)$ is an affine combination of $z$
            with coefficients $A_y$ in dimension $N$ by
            \cref{real_lt_subst_left,real_lt_subst_right}.
        Let $P_x=\{i\in I\mid x\in U(i)\}$.
        Let $P_y=\{i\in I\mid y\in U(i)\}$.
        $m+1\in\naturalsPlus$ by
            \cref{natural_successor_positive_50}.
        $P_x$ has at most $m+1$ elements and
            $P_y$ has at most $m+1$ elements by
            \cref{cover_point_index_set_at_most_50}.
        Let $D=\{i\in I\mid A_x(i)\neq A_y(i)\}$.
        $D$ has at most $N+1$ elements by
            \cref{finite_at_most_union_sum_50,menger_nobeling_order_arithmetic_50,real_lt_subst_right,partition_coefficient_difference_support_50,finite_at_most_subset_50}.
        $A_x=A_y$ by
            \cref{sparse_general_position_affine_coordinates_unique_50}.
        For all $j\in I$ we have $0\leq A_x(j)$ by
            \cref{finite_partition_coefficients_nonnegative_50}.
        Take $i\in I$ such that $0<A_x(i)$ by
            \cref{finite_real_sum_one_positive_term}.
        Contradiction by
            \cref{real_lt_subst_left,real_lt_subst_right,subseteq,positive_partition_component_in_support,finite_partition_of_unity,finite_partition_component_condition,metric_on,funs_type_apply,times_tuple_intro,real_leq_split,real_lt_trans,real_lt_irreflexive}.
    \end{subproof}
    $g\in W$ by
        \cref{metric_scale_separating_class_50,continuous,funs}.
    $g\in\metricBall{q}{C}{f}{\eta}\inter W$ by
        \cref{inter_intro}.
    $\metricBall{q}{C}{f}{\eta}\inter W$ is inhabited.
    Follows by \cref{intersects,inhabited_nonempty}.
\end{proof}

% from §50 helper lemma 137: each fixed-scale separating class is dense
\begin{lemma}\label{menger_nobeling_separating_class_dense_50}
    Assume $m\in\naturals$.
    Assume $m$ is a topological dimension bound for $X$ with respect to $T$.
    Assume $d$ is a metric on $X$ and $T=\metricTopology{d}{X}$.
    Assume $X$ is compact with respect to $T$ and $X$ is inhabited.
    Let $N=(1+1)\rmul m+1$.
    Let $E=\realsPower{\initialSegment{N}}$.
    Let $e=\squareMetric{N}$.
    Let $S=\metricTopology{e}{E}$.
    Assume $C$ is the set of continuous functions from $X$ to $E$
        with respect to $T$ and $S$.
    Assume $\rho$ is the uniform metric on functions from $X$ to $E$
        with respect to $e$.
    Let $q=\restrl{\rho}{C\times C}$.
    Assume $q$ is a metric on $C$.
    Assume $\epsilon\in\reals$ and $0<\epsilon$.
    Assume $W$ is the separating class at scale $\epsilon$ on $X$ to $E$
        with respect to $T$ and $S$ and $d$.
    Then $W$ is dense in $C$ with respect to $\metricTopology{q}{C}$.
\end{lemma}
\begin{proof}
    Follows by
        \cref{subseteq,metric_scale_separating_class_50,continuous_function_set,continuous,funs,menger_nobeling_separating_ball_intersection_50,metric_ball_intersections_imply_dense_49}.
\end{proof}

% from §50 helper lemma 44: Menger--Nobeling for a prescribed dimension bound
\begin{lemma}\label{menger_nobeling_dimension_bound_principle_50}
    Assume $m\in\naturals$.
    Assume $m$ is a topological dimension bound for $X$ with respect to $T$.
    Assume $X$ is compact with respect to $T$.
    Assume $X$ is metrizable with respect to $T$.
    Let $N=(1+1)\rmul m+1$.
    Let $E=\realsPower{\initialSegment{N}}$.
    Let $S=\metricTopology{\squareMetric{N}}{E}$.
    Then there exists $f$ such that
        $f$ is a topological embedding from $X$ to $E$ with respect to $T$ and $S$.
\end{lemma}
\begin{proof}
    $T$ is a topology on $X$ by
        \cref{topological_dimension_bound_50}.
    $N\in\naturalsPlus$ by
        \cref{menger_nobeling_target_dimension_positive_50}.
    Let $e=\squareMetric{N}$.
    $e$ is a metric on $E$ by \cref{square_metric_is_metric}.
    $S$ is a topology on $E$ by
        \cref{metric_topology,metric_basis_is_basis,basis_topology_is_topology}.
    \begin{byCase}
    \caseOf{$X=\emptyset$.}
        Let $f=\emptyset$.
        Thus there exists $g$ such that
            $g$ is a topological embedding from $X$ to $E$ with respect to
                $T$ and $S$ by \cref{empty_domain_topological_embedding}.
    \caseOf{$X\neq\emptyset$.}
        $X$ is inhabited by \cref{nonempty_inhabited}.
        Take $d$ such that
            $d$ is a metric on $X$ and $T=\metricTopology{d}{X}$ by
            \cref{metrizable}.
        Take $\rho$ such that
            $\rho$ is the uniform metric on functions from $X$ to $E$
                with respect to $e$ by
            \cref{uniform_metric_on_function_space_exists_inhabited}.
        Let $C=\{f\in\funs{X}{E}\mid
            \text{$f$ is continuous from $X$ to $E$ with respect to $T$ and $S$}\}$.
        $C$ is the set of continuous functions from $X$ to $E$ with respect to
            $T$ and $S$ by \cref{continuous_function_set}.
        Let $q=\restrl{\rho}{C\times C}$.
        $E$ is complete with respect to $e$ by
            \cref{euclidean_space_square_metric_complete}.
        $C$ is complete with respect to $q$ by
            \cref{continuous_function_space_complete_uniform}.
        $q$ is a metric on $C$ by \cref{complete_metric_space}.
        Let $T_C=\metricTopology{q}{C}$.
        $T_C$ is a topology on $C$ by
            \cref{metric_topology,metric_basis_is_basis,basis_topology_is_topology}.
        $C$ is a Baire space with respect to $T_C$ by
            \cref{complete_metric_space_baire}.
        Let $R=\{w\in\naturalsPlus\times\pow{C}\mid
            \text{$\snd{w}$ is the separating class at scale $1/(\fst{w})$ on
                $X$ to $E$ with respect to $T$ and $S$ and $d$}\}$.
        $R$ is a relation by
            \cref{times_elem_is_tuple,relation}.
        Show that for all $n\in\naturalsPlus$ there exists $W$ such that
            $n\mathrel{R}W$.
        \begin{subproof}
            Fix $n\in\naturalsPlus$.
            Let $W=\{f\in\funs{X}{E}\mid
                \text{$f$ is continuous from $X$ to $E$ with respect to
                    $T$ and $S$ and
                    for all $x,y\in X$ if $1/n\leq d((x,y))$
                    then $f(x)\neq f(y)$}\}$.
            Follows by
                \cref{metric_scale_separating_class_50,subseteq,continuous_function_set,continuous,funs,powerset_intro,times_tuple_intro,fst_eq,snd_eq,relation}.
        \end{subproof}
        Take $s$ such that
            $s$ is a function on $\naturalsPlus$ and
            for all $n\in\naturalsPlus$ we have
                $n\mathrel{R}s(n)$ by \cref{choice_function}.
        $s$ is a sequence in $\pow{C}$ by
            \cref{choice_function,function_apply_elim,subseteq,times_tuple_elim,funs_intro,sequence_in}.
        For all $n\in\naturalsPlus$ we have
            $s(n)$ is open in $C$ with respect to $T_C$ and
            $s(n)$ is dense in $C$ with respect to $T_C$ by
            \cref{choice_function,function_apply_elim,fst_eq,snd_eq,naturalsplus_reciprocal_real,positive_reciprocal_naturalsplus,metric_scale_separating_class_open_50,menger_nobeling_separating_class_dense_50}.
        Let $D=\inters{\img{s}{\naturalsPlus}}$.
        $D$ is dense in $C$ with respect to $T_C$ by
            \cref{baire_space_open_set_characterization}.
        $\img{s}{\naturalsPlus}\subseteq\pow{C}$ by
            \cref{img_subseteq_ran,sequence_in,funs_ran,subseteq_transitive}.
        $D\subseteq C$ by
            \cref{real_one_in_naturals,sequence_in,funs_elim,function_apply_elim,img_iff,powerset_elim,inters_subseteq_elem,subseteq_transitive}.
        $C$ is inhabited by
            \cref{continuous_reals_power_function_set_inhabited}.
        Take $c$ such that $c\in C$.
        $1\in\reals$ and $0<1$ by
            \cref{real_one_in_naturals,real_naturals_subset_reals,subseteq,real_zero_lt_one}.
        Take $f\in D$ such that $q((c,f))<1$ by
            \cref{dense_metric_ball_intersection_49}.
        $f\in C$ by \cref{subseteq}.
        $f$ is continuous from $X$ to $E$ with respect to $T$ and $S$ by
            \cref{continuous_function_set,subseteq}.
        For all $n\in\naturalsPlus$ we have
            for all $x,y\in X$ if $1/n\leq d((x,y))$
            then $f(x)\neq f(y)$ by
            \cref{sequence_in,funs_is_function,function_apply_elim,img_iff,inters_destr,choice_function,fst_eq,snd_eq,metric_scale_separating_class_50,subseteq}.
        $f\in\funs{X}{E}$ by \cref{continuous}.
        $f$ is injective by
            \cref{reciprocal_scale_separation_implies_injective_50}.
        Let $A=\img{f}{X}$.
        $f$ is a function and $\dom{f}=X$ by
            \cref{continuous,funs_elim}.
        $f\in\funs{X}{A}$ by
            \cref{funs_intro,function_apply_elim,ran_intro,img_dom}.
        $f\in\surj{X}{A}$ by
            \cref{funs_surj_iff,img_dom}.
        $f$ is a bijection from $X$ to $A$ by
            \cref{bijections,injective_function}.
        $A\subseteq E$ by
            \cref{continuous,funs_ran,img_subseteq_ran,subseteq_transitive}.
        $A$ is a subspace of $E$ with respect to $S$ by
            \cref{subspace_of}.
        $E$ is Hausdorff with respect to $S$ by
            \cref{metric_space_hausdorff}.
        $A$ is Hausdorff with respect to $\subspaceTopology{S}{A}$ by
            \cref{subspace_hausdorff}.
        $f$ is continuous from $X$ to $A$ with respect to $T$ and
            $\subspaceTopology{S}{A}$ by
            \cref{continuous_restrict_range,subseteq_refl}.
        $f$ is a homeomorphism between $X$ and $A$ with respect to $T$ and
            $\subspaceTopology{S}{A}$ by
            \cref{compact_hausdorff_bijection_homeomorphism}.
        $f$ is a topological embedding from $X$ to $E$ with respect to
            $T$ and $S$ by \cref{topological_embedding}.
        Thus there exists $g$ such that
            $g$ is a topological embedding from $X$ to $E$ with respect to
                $T$ and $S$.
    \end{byCase}
\end{proof}

% from §50 Item 11: the Menger-Nobeling imbedding theorem
\begin{theorem}\label{menger_nobeling_embedding_theorem_50}
    Assume $m\in\naturals$.
    Assume $m$ is the topological dimension of $X$ with respect to $T$.
    Assume $X$ is compact with respect to $T$.
    Assume $X$ is metrizable with respect to $T$.
    Let $N = (1 + 1) \rmul m + 1$.
    Let $E = \realsPower{\initialSegment{N}}$.
    Let $S = \metricTopology{\squareMetric{N}}{E}$.
    Then there exists $f$ such that
        $f$ is a topological embedding from $X$ to $E$ with respect to $T$ and $S$.
\end{theorem}
\begin{proof}
    Follows by
        \cref{topological_dimension_50,menger_nobeling_dimension_bound_principle_50}.
\end{proof}

% from §50 helper lemma 138: fine bounded-order covers give a compact metric dimension bound
\begin{lemma}\label{fine_ordered_metric_covers_dimension_bound_50}
    Assume $m\in\naturals$.
    Assume $d$ is a metric on $X$.
    Let $T=\metricTopology{d}{X}$.
    Assume $X$ is compact with respect to $T$.
    Suppose for all $\epsilon\in\reals$ such that $0<\epsilon$
        there exists $B$ such that
            $B$ is an open covering of $X$ with respect to $T$ and
            $B$ has order at most $m+1$ in $X$ and
            for all $V\in B$ for all $x,y\in V$ we have
                $d((x,y))<\epsilon$.
    Then $m$ is a topological dimension bound for $X$ with respect to $T$.
\end{lemma}
\begin{proof}
    $T$ is a topology on $X$ by
        \cref{metric_topology,metric_basis_is_basis,basis_topology_is_topology}.
    Show that for all $A$ if
        $A$ is an open covering of $X$ with respect to $T$
    then there exists $B$ such that
        $B$ is an open refinement of $A$ on $X$ with respect to $T$ and
        $B$ has order at most $m+1$ in $X$.
    \begin{subproof}
        Fix $A$.
        Assume $A$ is an open covering of $X$ with respect to $T$.
        $A\subseteq\pow{X}$ by
            \cref{open_covering,open_in,topology_on,subseteq,pow_iff,powerset_intro}.
        \begin{byCase}
        \caseOf{$X=\emptyset$.}
            Let $B=\emptyset$.
            Thus there exists $C$ such that
                $C$ is an open refinement of $A$ on $X$ with respect to $T$ and
                $C$ has order at most $m+1$ in $X$ by
                \cref{emptyset_subseteq,emptyset,covering,subseteq,unions_iff,refinement_on,open_refinement_on,natural_successor_positive_50,covering_order_at_most_50}.
        \caseOf{$X\neq\emptyset$.}
            Take $\epsilon\in\reals$ such that
                $0<\epsilon$ and
                for all $D$ such that
                    $D\subseteq X$ and
                    for all $x,y\in D$ we have $d((x,y))<\epsilon$
                there exists $U\in A$ such that $D\subseteq U$ by
                \cref{lebesgue_number_lemma}.
            Take $B$ such that
                $B$ is an open covering of $X$ with respect to $T$ and
                $B$ has order at most $m+1$ in $X$ and
                for all $V\in B$ for all $x,y\in V$ we have
                    $d((x,y))<\epsilon$.
            $B\subseteq\pow{X}$ by
                \cref{open_covering,open_in,topology_on,subseteq,pow_iff,powerset_intro}.
            For all $V\in B$ there exists $U\in A$ such that
                $V\subseteq U$ by \cref{powerset_elim,subseteq}.
            $B$ is an open refinement of $A$ on $X$ with respect to $T$ by
                \cref{refinement_on,open_covering,open_refinement_on}.
            Thus there exists $C$ such that
                $C$ is an open refinement of $A$ on $X$ with respect to $T$ and
                $C$ has order at most $m+1$ in $X$.
        \end{byCase}
    \end{subproof}
    Follows by \cref{topological_dimension_bound_50}.
\end{proof}

% from §50 helper lemma 139: an initial segment has its evident finite bound
\begin{lemma}\label{initial_segment_at_most_self_50}
    Assume $n\in\naturalsPlus$.
    Then $\initialSegment{n}$ has at most $n$ elements.
\end{lemma}
\begin{proof}
    Let $A=\{k\in\naturalsPlus\mid
        \text{$\initialSegment{k}$ has at most $k$ elements}\}$.
    $A\subseteq\naturalsPlus$.
    $1\in A$ by
        \cref{initial_segment_one,singleton_at_most_positive_50,real_one_in_naturals}.
    Show that for all $k\in A$ we have $k+1\in A$.
    \begin{subproof}
        Fix $k\in A$.
        Let $D=\initialSegment{k}\union\{k+1\}$.
        $D$ has at most $k+1$ elements by
            \cref{subseteq,finite_at_most_adjoin_point_50}.
        For all $i\in D$ we have $i\in\initialSegment{k+1}$ by
            \cref{union_iff,singleton_iff,initial_segment_subset_succ,subseteq,real_naturals_closed_succ,initial_segment_naturalsplus}.
        $D\subseteq\initialSegment{k+1}$ by \cref{subseteq}.
        For all $i\in\initialSegment{k+1}$ we have $i\in D$ by
            \cref{initial_segment_succ_split,union_intro_left,union_intro_right,singleton_intro}.
        $\initialSegment{k+1}\subseteq D$ by \cref{subseteq}.
        $D=\initialSegment{k+1}$ by \cref{subseteq_antisymmetric}.
        Follows by
            \cref{real_lt_subst_left,real_naturals_closed_succ,subseteq}.
    \end{subproof}
    For all $k\in\naturalsPlus$ we have $k\in A$ by
        \cref{real_naturals_induction}.
    Follows by \cref{subseteq}.
\end{proof}

% from §50 helper lemma 140: every coordinate subset has a rank bounded by the ambient dimension
\begin{lemma}\label{coordinate_subset_rank_50}
    Assume $N\in\naturalsPlus$.
    Let $I=\initialSegment{N}$.
    Assume $J\subseteq I$.
    Then there exists $m\in\naturals$ such that
        $m\leq N$ and
        either $m=0$ and $J=\emptyset$ or
        $m\in\naturalsPlus$ and
            there exists $g$ such that
                $g$ is a bijection from $\initialSegment{m}$ to $J$.
\end{lemma}
\begin{proof}
    \begin{byCase}
    \caseOf{$J=\emptyset$.}
        $0\in\naturals$ by
            \cref{naturals,union_intro_right,singleton_intro}.
        $0\leq N$ by
            \cref{real_add_ident,real_naturals_subset_reals,subseteq,naturalsplus_positive_real,real_lt_imp_leq}.
        Thus there exists $m\in\naturals$ such that
            $m\leq N$ and
            either $m=0$ and $J=\emptyset$ or
            $m\in\naturalsPlus$ and
                there exists $g$ such that
                    $g$ is a bijection from $\initialSegment{m}$ to $J$.
    \caseOf{$J\neq\emptyset$.}
        Take $m,g$ such that
            $m\in\naturalsPlus$ and
            $g$ is a bijection from $\initialSegment{m}$ to $J$ by
            \cref{initial_segment_finite,subseteq_finite_is_finite,nonempty_inhabited,finite_inhabited_initial_segment_enumeration}.
        $m\leq N$ by
            \cref{initial_segment_at_most_self_50,finite_at_most_subset_50,finite_at_most_enumeration_length_50}.
        Thus there exists $n\in\naturals$ such that
            $n\leq N$ and
            either $n=0$ and $J=\emptyset$ or
            $n\in\naturalsPlus$ and
                there exists $h$ such that
                    $h$ is a bijection from $\initialSegment{n}$ to $J$ by
            \cref{naturals,union_intro_left}.
    \end{byCase}
\end{proof}

% from §50 helper lemma 141: distinct integers are separated by a full unit
\begin{lemma}\label{integer_successor_gap_50}
    Assume $a,b\in\integers$.
    Suppose $a<b$.
    Then $a+1\leq b$.
\end{lemma}
\begin{proof}
    For all $p,q$ such that
        $p\in\naturalsPlus$ and $q\in\naturalsPlus$ and $p<q$
    we have $p+1\leq q$ by
        \cref{real_naturals_subset_reals,subseteq,real_naturals_closed_succ,real_lt_trichotomy,real_succ_discrete,real_lt_trans,real_lt_irreflexive,real_lt_subst_right,real_lt_subst_left}.
    $a,b\in\reals$ by \cref{real_integers,subseteq}.
    $0\in\reals$ by \cref{real_add_ident}.
    $1,a+1\in\reals$ by
        \cref{real_one_in_naturals,real_naturals_subset_reals,subseteq,real_add_closed}.
    $a\in\naturalsPlus$ or $a=0$ or $\neg{a}\in\naturalsPlus$ by
        \cref{real_integers,subseteq}.
    $b\in\naturalsPlus$ or $b=0$ or $\neg{b}\in\naturalsPlus$ by
        \cref{real_integers,subseteq}.
    \begin{byCase}
    \caseOf{$a\in\naturalsPlus$.}
        Follows by
            \cref{naturalsplus_positive_real,real_lt_subst_left,real_add_ident,real_add_inverse,real_lt_neg,real_neg_neg,real_lt_trans,real_lt_irreflexive}.
    \caseOf{$a=0$.}
        $b\in\naturalsPlus$ by
            \cref{real_lt_subst_left,real_lt_subst_right,real_lt_irreflexive,naturalsplus_positive_real,real_add_ident,real_add_inverse,real_lt_neg,real_neg_neg,real_lt_trans}.
        Follows by
            \cref{real_one_leq_naturalsplus,real_lt_imp_leq,real_one_in_naturals,real_naturals_subset_reals,subseteq,real_add_ident,real_add_comm,real_lt_subst_left}.
    \caseOf{$\neg{a}\in\naturalsPlus$.}
        Let $p=\neg{a}$.
        $p\in\naturalsPlus$.
        $1<p$ or $1=p$ by
            \cref{real_one_leq_naturalsplus}.
        $1\leq p$ by \cref{real_lt_imp_leq}.
        $\neg{p}\leq\neg{1}$ by
            \cref{real_naturals_subset_reals,real_one_in_naturals,subseteq,real_leq_neg}.
        $\neg{p}+1\leq\neg{1}+1$ by
            \cref{real_add_inverse,real_one_in_naturals,real_naturals_subset_reals,subseteq,real_leq_add}.
        $\neg{1}+1=0$ by
            \cref{real_one_in_naturals,real_naturals_subset_reals,subseteq,real_add_inverse,real_add_comm}.
        $a=\neg{p}$ by
            \cref{real_add_inverse,real_neg_neg}.
        $a+1\leq0$ by
            \cref{real_lt_subst_left,real_lt_subst_right}.
        \begin{byCase}
        \caseOf{$b\in\naturalsPlus$.}
            Follows by
                \cref{naturalsplus_positive_real,real_lt_imp_leq,real_add_ident,real_one_in_naturals,real_naturals_subset_reals,subseteq,real_leq_trans}.
        \caseOf{$b=0$.}
            Follows by \cref{real_lt_subst_right}.
        \caseOf{$\neg{b}\in\naturalsPlus$.}
            Let $q=\neg{b}$.
            $q\in\naturalsPlus$.
            $q<p$ by
                \cref{real_lt_neg,real_neg_neg,real_lt_subst_left,real_lt_subst_right}.
            $\neg{p}\leq\neg{(q+1)}$ by
                \cref{real_naturals_subset_reals,subseteq,real_naturals_closed_succ,real_leq_neg}.
            $\neg{p}+1\leq\neg{(q+1)}+1$ by
                \cref{real_add_inverse,real_one_in_naturals,real_naturals_subset_reals,subseteq,real_naturals_closed_succ,real_leq_add}.
            $\neg{(q+1)}+1=\neg{q}$ by
                \cref{real_naturals_subset_reals,real_one_in_naturals,subseteq,real_neg_addition,real_add_inverse,real_add_assoc,real_add_ident}.
            Follows by
                \cref{real_add_inverse,real_neg_neg,real_lt_subst_left,real_lt_subst_right}.
        \end{byCase}
    \end{byCase}
\end{proof}

% from §50 helper definition 39: agreement at integer coordinates
\begin{definition}\label{cubically_integer_compatible_50}
    $x$ is cubically integer-compatible with $y$ in dimension $N$ iff
        $N\in\naturalsPlus$ and
        $x,y\in\realsPower{\initialSegment{N}}$ and
        for all $i\in\initialSegment{N}$ if
            $x(i)\in\integers$
        then $y(i)=x(i)$.
\end{definition}

% from §50 helper definition 40: agreement inside noninteger unit intervals
\begin{definition}\label{cubically_noninteger_compatible_50}
    $x$ is cubically noninteger-compatible with $y$ in dimension $N$ iff
        $N\in\naturalsPlus$ and
        $x,y\in\realsPower{\initialSegment{N}}$ and
        for all $i\in\initialSegment{N}$ if
            $x(i)\notin\integers$
        then
                there exists $n\in\integers$ such that
                    $n<x(i)$ and $x(i)<n+1$ and
                    $n<y(i)$ and $y(i)<n+1$.
\end{definition}

% from §50 helper definition 41: equivalence inside the integer cubical decomposition
\begin{definition}\label{cubically_equivalent_50}
    $x$ is cubically equivalent to $y$ in dimension $N$ iff
        $x$ is cubically integer-compatible with $y$ in dimension $N$ and
        $x$ is cubically noninteger-compatible with $y$ in dimension $N$.
\end{definition}

% from §50 helper definition 42: a cell in the integer cubical decomposition
\begin{definition}\label{integer_cubical_cell_50}
    $C$ is an integer cubical cell in dimension $N$ iff
        there exists $x\in\realsPower{\initialSegment{N}}$ such that
            $C=\{y\in\realsPower{\initialSegment{N}}\mid
                \text{$x$ is cubically equivalent to $y$ in dimension $N$}\}$.
\end{definition}

% from §50 helper definition 43: the positive level coding a cubical cell's dimension
\begin{definition}\label{integer_cubical_level_50}
    $C$ has cubical level $k$ in dimension $N$ iff
        $k\in\naturalsPlus$ and
        $C$ is an integer cubical cell in dimension $N$ and
        there exist $x,J$ such that
            $x\in C$ and
            $J=\{i\in\initialSegment{N}\mid x(i)\notin\integers\}$ and
            either $k=1$ and $J=\emptyset$ or
                there exists $m\in\naturalsPlus$ such that
                    $k=m+1$ and
                    there exists $g$ such that
                        $g$ is a bijection from $\initialSegment{m}$ to $J$.
\end{definition}

% from §50 helper lemma 142: every euclidean point lies in its generated cubical cell
\begin{lemma}\label{generated_integer_cubical_cell_50}
    Assume $N\in\naturalsPlus$.
    Let $E=\realsPower{\initialSegment{N}}$.
    Assume $x\in E$.
    Let $C=\{y\in E\mid
        \text{$x$ is cubically equivalent to $y$ in dimension $N$}\}$.
    Then $x\in C$ and $C$ is an integer cubical cell in dimension $N$.
\end{lemma}
\begin{proof}
    Let $I=\initialSegment{N}$.
    For all $i\in I$ if $x(i)\in\integers$ then $x(i)=x(i)$.
    For all $i\in I$ if $x(i)\notin\integers$ then
        there exists $n\in\integers$ such that
            $n<x(i)$ and $x(i)<n+1$ and
            $n<x(i)$ and $x(i)<n+1$ by
        \cref{reals_power,tuple_power,funs_type_apply,real_consecutive_integers,real_lt_subst_right,real_leq_split}.
    $x$ is cubically equivalent to $x$ in dimension $N$ by
        \cref{cubically_integer_compatible_50,cubically_noninteger_compatible_50,cubically_equivalent_50}.
    $x\in C$ by \cref{subseteq}.
    $C$ is an integer cubical cell in dimension $N$ by
        \cref{integer_cubical_cell_50}.
\end{proof}

% from §50 helper lemma 143: an open unit interval between integers contains no integer
\begin{lemma}\label{no_integer_between_successive_integers_50}
    Assume $n\in\integers$ and $r\in\reals$.
    Suppose $n<r$ and $r<n+1$.
    Then $r\notin\integers$.
\end{lemma}
\begin{proof}
    Suppose not.
    $n+1\in\reals$ by
        \cref{real_integers,subseteq,real_one_in_naturals,real_naturals_subset_reals,real_add_closed}.
    $n+1<n+1$ by
        \cref{integer_successor_gap_50,real_leq_lt_trans}.
    Contradiction by \cref{real_lt_irreflexive}.
\end{proof}

% from §50 helper lemma 144: cubically equivalent points have the same nonintegral coordinates
\begin{lemma}\label{cubical_noninteger_coordinates_equal_50}
    Assume $x$ is cubically equivalent to $y$ in dimension $N$.
    Let $I=\initialSegment{N}$.
    Let $J_x=\{i\in I\mid x(i)\notin\integers\}$.
    Let $J_y=\{i\in I\mid y(i)\notin\integers\}$.
    Then $J_x=J_y$.
\end{lemma}
\begin{proof}
    $N\in\naturalsPlus$ and
        $x,y\in\realsPower{I}$ by
        \cref{cubically_equivalent_50,cubically_integer_compatible_50}.
    For all $i\in J_x$ we have $i\in J_y$ by
        \cref{subseteq,cubically_equivalent_50,cubically_noninteger_compatible_50,reals_power,tuple_power,funs_type_apply,no_integer_between_successive_integers_50}.
    $J_x\subseteq J_y$ by \cref{subseteq}.
    For all $i\in J_y$ we have $i\in J_x$ by
        \cref{subseteq,cubically_equivalent_50,cubically_integer_compatible_50,real_lt_subst_left}.
    $J_y\subseteq J_x$ by \cref{subseteq}.
    Follows by \cref{subseteq_antisymmetric}.
\end{proof}

% from §50 helper lemma 145: every cubical cell has a level among the first N+1 levels
\begin{lemma}\label{integer_cubical_cell_level_exists_50}
    Assume $N\in\naturalsPlus$.
    Assume $C$ is an integer cubical cell in dimension $N$.
    Then there exists $k\in\initialSegment{N+1}$ such that
        $C$ has cubical level $k$ in dimension $N$.
\end{lemma}
\begin{proof}
    Let $I=\initialSegment{N}$.
    Let $E=\realsPower{I}$.
    Take $x\in E$ such that
        $C=\{y\in E\mid
            \text{$x$ is cubically equivalent to $y$ in dimension $N$}\}$ by
        \cref{integer_cubical_cell_50}.
    $x\in C$ by
        \cref{generated_integer_cubical_cell_50,real_lt_subst_left}.
    Let $J=\{i\in I\mid x(i)\notin\integers\}$.
    $J\subseteq I$.
    Take $m\in\naturals$ such that
        $m\leq N$ and
        either $m=0$ and $J=\emptyset$ or
        $m\in\naturalsPlus$ and
            there exists $g$ such that
                $g$ is a bijection from $\initialSegment{m}$ to $J$ by
        \cref{coordinate_subset_rank_50}.
    \begin{byCase}
    \caseOf{$m=0$ and $J=\emptyset$.}
        $1\in\naturalsPlus$ by \cref{real_one_in_naturals}.
        $N+1\in\naturalsPlus$ by
            \cref{real_naturals_closed_succ}.
        $1\leq N+1$ by
            \cref{real_one_leq_naturalsplus,real_lt_imp_leq}.
        $1\in\initialSegment{N+1}$ by
            \cref{initial_segment_naturalsplus}.
        There exists no $r\in\naturalsPlus$ such that
            $1=r+1$ and
            there exists $g$ such that
                $g$ is a bijection from $\initialSegment{r}$ to $J$ by
            \cref{real_one_in_naturals,real_one_leq_naturalsplus,initial_segment_naturalsplus,bijections_to_funs,funs_type_apply,emptyset}.
        $C$ has cubical level $1$ in dimension $N$ by
            \cref{integer_cubical_level_50}.
        Thus there exists $k\in\initialSegment{N+1}$ such that
            $C$ has cubical level $k$ in dimension $N$.
    \caseOf{$m\in\naturalsPlus$ and
            there exists $g$ such that
                $g$ is a bijection from $\initialSegment{m}$ to $J$.}
        $J\neq\emptyset$ by
            \cref{real_one_in_naturals,real_one_leq_naturalsplus,initial_segment_naturalsplus,bijections_to_funs,funs_type_apply,emptyset}.
        $m+1\in\naturalsPlus$ by
            \cref{real_naturals_closed_succ}.
        $m+1\in\initialSegment{N+1}$ by
            \cref{real_naturals_closed_succ,real_naturals_subset_reals,real_one_in_naturals,subseteq,real_leq_add,initial_segment_naturalsplus}.
        $C$ has cubical level $m+1$ in dimension $N$ by
            \cref{integer_cubical_level_50}.
        Thus there exists $k\in\initialSegment{N+1}$ such that
            $C$ has cubical level $k$ in dimension $N$.
    \end{byCase}
\end{proof}

% from §50 helper lemma 146: a cubical cell has only one level
\begin{lemma}\label{integer_cubical_cell_level_unique_50}
    Assume $C$ has cubical level $k$ in dimension $N$.
    Assume $C$ has cubical level $l$ in dimension $N$.
    Then $k=l$.
\end{lemma}
\begin{proof}
    $k,l\in\naturalsPlus$ and
        $C$ is an integer cubical cell in dimension $N$ by
        \cref{integer_cubical_level_50}.
    Take $x,J$ such that
        $x\in C$ and
        $J=\{i\in\initialSegment{N}\mid x(i)\notin\integers\}$ and
        either $k=1$ and $J=\emptyset$ or
            there exists $m\in\naturalsPlus$ such that
                $k=m+1$ and
                there exists $g$ such that
                    $g$ is a bijection from $\initialSegment{m}$ to $J$ by
        \cref{integer_cubical_level_50}.
    Take $y,K$ such that
        $y\in C$ and
        $K=\{i\in\initialSegment{N}\mid y(i)\notin\integers\}$ and
        either $l=1$ and $K=\emptyset$ or
            there exists $n\in\naturalsPlus$ such that
                $l=n+1$ and
                there exists $h$ such that
                    $h$ is a bijection from $\initialSegment{n}$ to $K$ by
        \cref{integer_cubical_level_50}.
    Let $I=\initialSegment{N}$.
    Let $E=\realsPower{I}$.
    Take $z\in E$ such that
        $C=\{w\in E\mid
            \text{$z$ is cubically equivalent to $w$ in dimension $N$}\}$ by
        \cref{integer_cubical_cell_50}.
    $z$ is cubically equivalent to $x$ in dimension $N$ and
        $z$ is cubically equivalent to $y$ in dimension $N$ by
        \cref{subseteq}.
    Let $L=\{i\in I\mid z(i)\notin\integers\}$.
    $L=J$ and $L=K$ by
        \cref{cubical_noninteger_coordinates_equal_50}.
    $J=K$ by \cref{real_lt_subst_left,real_lt_subst_right}.
    \begin{byCase}
    \caseOf{$J=\emptyset$.}
        $K=\emptyset$ by
            \cref{real_lt_subst_left,real_lt_subst_right}.
        $k=1$ by
            \cref{real_one_in_naturals,real_one_leq_naturalsplus,initial_segment_naturalsplus,bijections_to_funs,funs_type_apply,emptyset}.
        $l=1$ by
            \cref{real_one_in_naturals,real_one_leq_naturalsplus,initial_segment_naturalsplus,bijections_to_funs,funs_type_apply,emptyset}.
        Follows by \cref{real_lt_subst_left,real_lt_subst_right}.
    \caseOf{$J\neq\emptyset$.}
        $K\neq\emptyset$ by
            \cref{real_lt_subst_left,real_lt_subst_right}.
        Take $m\in\naturalsPlus$ such that
            $k=m+1$ and
            there exists $g$ such that
                $g$ is a bijection from $\initialSegment{m}$ to $J$ by
            \cref{emptyset}.
        Take $n\in\naturalsPlus$ such that
            $l=n+1$ and
            there exists $h$ such that
                $h$ is a bijection from $\initialSegment{n}$ to $K$ by
            \cref{emptyset}.
        $J$ has at most $m$ elements and
            $K$ has at most $n$ elements by
            \cref{initial_segment_at_most_self_50,bijection_preserves_finite_at_most_50}.
        $m\leq n$ and $n\leq m$ by
            \cref{real_lt_subst_left,real_lt_subst_right,finite_at_most_enumeration_length_50}.
        $m=n$ by
            \cref{real_naturals_subset_reals,subseteq,real_leq_antisym}.
        Follows by
            \cref{real_lt_subst_left,real_lt_subst_right}.
    \end{byCase}
\end{proof}

% from §50 helper lemma 147: two successive-integer intervals containing a point coincide
\begin{lemma}\label{successive_integer_interval_unique_50}
    Assume $n,q\in\integers$ and $r\in\reals$.
    Assume $n<r$ and $r<n+1$.
    Assume $q<r$ and $r<q+1$.
    Then $n=q$.
\end{lemma}
\begin{proof}
    Follows by
        \cref{real_integers,subseteq,real_lt_trichotomy,integer_successor_gap_50,real_one_in_naturals,real_naturals_subset_reals,real_add_closed,real_leq_lt_trans,real_lt_trans,real_lt_irreflexive}.
\end{proof}

% from §50 helper lemma 148: cubical equivalence is symmetric
\begin{lemma}\label{cubical_equivalence_symmetric_50}
    Assume $x$ is cubically equivalent to $y$ in dimension $N$.
    Then $y$ is cubically equivalent to $x$ in dimension $N$.
\end{lemma}
\begin{proof}
    Let $I=\initialSegment{N}$.
    Let $E=\realsPower{I}$.
    $N\in\naturalsPlus$ and $x,y\in E$ by
        \cref{cubically_equivalent_50,cubically_integer_compatible_50}.
    Let $J_x=\{i\in I\mid x(i)\notin\integers\}$.
    Let $J_y=\{i\in I\mid y(i)\notin\integers\}$.
    $J_x=J_y$ by
        \cref{cubical_noninteger_coordinates_equal_50}.
    For all $i\in I$ if $y(i)\in\integers$ then $x(i)=y(i)$ by
        \cref{subseteq,real_lt_subst_left,cubically_equivalent_50,cubically_integer_compatible_50}.
    For all $i\in I$ if $y(i)\notin\integers$ then
        there exists $n\in\integers$ such that
            $n<y(i)$ and $y(i)<n+1$ and
            $n<x(i)$ and $x(i)<n+1$ by
        \cref{subseteq,real_lt_subst_left,cubically_equivalent_50,cubically_noninteger_compatible_50}.
    Follows by
        \cref{cubically_integer_compatible_50,cubically_noninteger_compatible_50,cubically_equivalent_50}.
\end{proof}

% from §50 helper lemma 149: cubical equivalence is transitive
\begin{lemma}\label{cubical_equivalence_transitive_50}
    Assume $x$ is cubically equivalent to $y$ in dimension $N$.
    Assume $y$ is cubically equivalent to $z$ in dimension $N$.
    Then $x$ is cubically equivalent to $z$ in dimension $N$.
\end{lemma}
\begin{proof}
    Let $I=\initialSegment{N}$.
    Let $E=\realsPower{I}$.
    $N\in\naturalsPlus$ and $x,y,z\in E$ by
        \cref{cubically_equivalent_50,cubically_integer_compatible_50}.
    For all $i\in I$ if $x(i)\in\integers$ then $z(i)=x(i)$ by
        \cref{cubically_equivalent_50,cubically_integer_compatible_50,real_lt_subst_left}.
    For all $i\in I$ if $x(i)\notin\integers$ then
        there exists $n\in\integers$ such that
            $n<x(i)$ and $x(i)<n+1$ and
            $n<z(i)$ and $z(i)<n+1$ by
        \cref{cubically_equivalent_50,cubically_noninteger_compatible_50,reals_power,tuple_power,funs_type_apply,no_integer_between_successive_integers_50,successive_integer_interval_unique_50}.
    Follows by
        \cref{cubically_integer_compatible_50,cubically_noninteger_compatible_50,cubically_equivalent_50}.
\end{proof}

% from §50 helper lemma 150: two cubical cells containing one point are equal
\begin{lemma}\label{integer_cubical_cells_common_point_equal_50}
    Assume $C$ is an integer cubical cell in dimension $N$.
    Assume $D$ is an integer cubical cell in dimension $N$.
    Assume $x\in C\inter D$.
    Then $C=D$.
\end{lemma}
\begin{proof}
    Let $I=\initialSegment{N}$.
    Let $E=\realsPower{I}$.
    Take $a\in E$ such that
        $C=\{y\in E\mid
            \text{$a$ is cubically equivalent to $y$ in dimension $N$}\}$ by
        \cref{integer_cubical_cell_50}.
    Take $b\in E$ such that
        $D=\{y\in E\mid
            \text{$b$ is cubically equivalent to $y$ in dimension $N$}\}$ by
        \cref{integer_cubical_cell_50}.
    $x\in C$ and $x\in D$ by \cref{inter_elim_left,inter_elim_right}.
    $a$ is cubically equivalent to $x$ in dimension $N$ and
        $b$ is cubically equivalent to $x$ in dimension $N$ by
        \cref{subseteq}.
    $C\subseteq D$ by
        \cref{subseteq,cubical_equivalence_symmetric_50,cubical_equivalence_transitive_50}.
    $D\subseteq C$ by
        \cref{subseteq,cubical_equivalence_symmetric_50,cubical_equivalence_transitive_50}.
    Follows by \cref{subseteq_antisymmetric}.
\end{proof}

% from §50 helper lemma 151: finite subsets of the same positive rank are equal
\begin{lemma}\label{equal_finite_rank_subset_equal_50}
    Assume $m\in\naturalsPlus$.
    Assume $f$ is a bijection from $\initialSegment{m}$ to $A$.
    Assume $g$ is a bijection from $\initialSegment{m}$ to $B$.
    Assume $A\subseteq B$.
    Then $A=B$.
\end{lemma}
\begin{proof}
    Suppose not.
    Take $y\in B$ such that $y\notin A$ by
        \cref{subseteq,subseteq_antisymmetric}.
    $\initialSegment{m}$ has at most $m$ elements by
        \cref{initial_segment_at_most_self_50}.
    $B$ has at most $m$ elements by
        \cref{bijection_preserves_finite_at_most_50}.
    \begin{byCase}
    \caseOf{$m=1$.}
        $y\in A$ by
            \cref{initial_segment_one,singleton_iff,real_lt_subst_right,bijections_to_funs,funs_type_apply,subseteq,inhabited_at_most_one_singleton_50,real_lt_subst_left}.
        Contradiction.
    \caseOf{$m\neq 1$.}
        Take $q\in\naturalsPlus$ such that $m=q+1$ by
            \cref{real_naturals_predecessor}.
        Let $D=B\setminus\{y\}$.
        $m\leq q$ by
            \cref{finite_at_most_remove_point_50,real_lt_subst_left,setminus_intro,singleton_iff,subseteq,finite_at_most_subset_50,finite_at_most_enumeration_length_50}.
        $m\in\reals$ by
            \cref{real_naturals_subset_reals,subseteq}.
        $m<m$ by
            \cref{real_naturals_subset_reals,real_one_in_naturals,subseteq,real_zero_lt_one,real_lt_add,real_add_ident,real_add_comm,real_lt_subst_right,real_leq_lt_trans}.
        Contradiction by \cref{real_lt_irreflexive}.
    \end{byCase}
\end{proof}

% from §50 helper lemma 152: a sufficiently small perturbation stays in one unit interval
\begin{lemma}\label{small_perturbation_stays_in_integer_interval_50}
    Assume $n\in\integers$ and $x,y,r\in\reals$.
    Assume $0<r$.
    Assume $n<x$ and $x<n+1$.
    Assume $r<x-n$ and $r<n+1-x$.
    Assume $\abs{x-y}<r$.
    Then $n<y$ and $y<n+1$.
\end{lemma}
\begin{proof}
    $\abs{y-x}\leq r$ by
        \cref{real_lt_imp_leq,real_abs_sub_swap,real_lt_subst_left}.
    $x-r\leq y$ and $y\leq x+r$ by
        \cref{real_lt_imp_leq,real_abs_sub_leq_imp_left_bound,real_abs_sub_leq_imp_right_bound}.
    $n,r,x\in\reals$ by
        \cref{real_integers,subseteq}.
    $x-r,x+r,n+1\in\reals$ by
        \cref{real_sub,real_add_inverse,real_add_closed,real_one_in_naturals,real_naturals_subset_reals,subseteq}.
    $n<x-r$ by
        \cref{real_sub,real_add_inverse,real_add_closed,real_lt_add,real_add_assoc,real_add_comm,real_add_inverse,real_add_ident,real_lt_subst_left,real_lt_subst_right}.
    $n<y$ by \cref{real_lt_leq_trans_local}.
    $x+r<n+1$ by
        \cref{real_sub,real_add_inverse,real_add_closed,real_lt_add,real_add_assoc,real_add_comm,real_add_inverse,real_add_ident,real_lt_subst_left,real_lt_subst_right}.
    $y<n+1$ by \cref{real_leq_lt_trans}.
\end{proof}

% from §50 helper lemma 153: points of one cubical cell have the same noninteger coordinates
\begin{lemma}\label{same_cubical_cell_noninteger_coordinates_50}
    Assume $C$ is an integer cubical cell in dimension $N$.
    Assume $x,y\in C$.
    Let $I=\initialSegment{N}$.
    Let $J_x=\{i\in I\mid x(i)\notin\integers\}$.
    Let $J_y=\{i\in I\mid y(i)\notin\integers\}$.
    Then $J_x=J_y$.
\end{lemma}
\begin{proof}
    Let $E=\realsPower{\initialSegment{N}}$.
    Take $z\in E$ such that
        $C=\{w\in E\mid
            \text{$z$ is cubically equivalent to $w$ in dimension $N$}\}$ by
        \cref{integer_cubical_cell_50}.
    $z$ is cubically equivalent to $x$ in dimension $N$ and
        $z$ is cubically equivalent to $y$ in dimension $N$ by
        \cref{subseteq}.
    $x$ is cubically equivalent to $y$ in dimension $N$ by
        \cref{cubical_equivalence_symmetric_50,cubical_equivalence_transitive_50}.
    Follows by \cref{cubical_noninteger_coordinates_equal_50}.
\end{proof}

% from §50 helper lemma 154: inclusion of coordinate patterns at one cubical level is equality
\begin{lemma}\label{same_level_noninteger_coordinate_inclusion_equal_50}
    Assume $C$ has cubical level $k$ in dimension $N$.
    Assume $D$ has cubical level $k$ in dimension $N$.
    Assume $x\in C$ and $y\in D$.
    Let $I=\initialSegment{N}$.
    Let $J_x=\{i\in I\mid x(i)\notin\integers\}$.
    Let $J_y=\{i\in I\mid y(i)\notin\integers\}$.
    Assume $J_x\subseteq J_y$.
    Then $J_x=J_y$.
\end{lemma}
\begin{proof}
    $k\in\naturalsPlus$ and
        $C$ is an integer cubical cell in dimension $N$ and
        $D$ is an integer cubical cell in dimension $N$ by
        \cref{integer_cubical_level_50}.
    Take $a,J$ such that
        $a\in C$ and
        $J=\{i\in\initialSegment{N}\mid a(i)\notin\integers\}$ and
        either $k=1$ and $J=\emptyset$ or
            there exists $m\in\naturalsPlus$ such that
                $k=m+1$ and
                there exists $f$ such that
                    $f$ is a bijection from $\initialSegment{m}$ to $J$ by
        \cref{integer_cubical_level_50}.
    Take $b,K$ such that
        $b\in D$ and
        $K=\{i\in\initialSegment{N}\mid b(i)\notin\integers\}$ and
        either $k=1$ and $K=\emptyset$ or
            there exists $n\in\naturalsPlus$ such that
                $k=n+1$ and
                there exists $g$ such that
                    $g$ is a bijection from $\initialSegment{n}$ to $K$ by
        \cref{integer_cubical_level_50}.
    $J=J_x$ and $K=J_y$ by
        \cref{same_cubical_cell_noninteger_coordinates_50}.
    \begin{byCase}
    \caseOf{$k=1$ and $J=\emptyset$.}
        $K=\emptyset$ by
            \cref{naturalsplus_positive_real,real_add_ident,real_one_in_naturals,real_naturals_subset_reals,subseteq,real_naturals_closed_succ,real_lt_add,real_add_comm,real_lt_subst_left,real_lt_subst_right,real_lt_irreflexive}.
        $J_x=\emptyset$ and $J_y=\emptyset$ by
            \cref{real_lt_subst_left,real_lt_subst_right}.
        Follows by
            \cref{real_lt_subst_left,real_lt_subst_right}.
    \caseOf{There exists $m\in\naturalsPlus$ such that
            $k=m+1$ and
            there exists $f$ such that
                $f$ is a bijection from $\initialSegment{m}$ to $J$.}
        Take $m\in\naturalsPlus$ such that
            $k=m+1$ and
            there exists $f$ such that
                $f$ is a bijection from $\initialSegment{m}$ to $J$.
        There exists $p\in\naturalsPlus$ such that
            $k=p+1$ and
            there exists $g$ such that
                $g$ is a bijection from $\initialSegment{p}$ to $K$ by
            \cref{naturalsplus_positive_real,real_add_ident,real_one_in_naturals,real_naturals_subset_reals,subseteq,real_naturals_closed_succ,real_lt_add,real_add_comm,real_lt_subst_left,real_lt_subst_right,real_lt_irreflexive}.
        Take $n\in\naturalsPlus$ such that
            $k=n+1$ and
            there exists $g$ such that
                $g$ is a bijection from $\initialSegment{n}$ to $K$.
        $m+1=n+1$ by
            \cref{real_lt_subst_left,real_lt_subst_right}.
        $m=n$ by
            \cref{real_naturals_subset_reals,real_one_in_naturals,subseteq,real_add_cancel}.
        $J\subseteq K$ by
            \cref{real_lt_subst_left,real_lt_subst_right}.
        $J=K$ by
            \cref{equal_finite_rank_subset_equal_50,real_lt_subst_left,real_lt_subst_right}.
        Follows by
            \cref{real_lt_subst_left,real_lt_subst_right}.
    \end{byCase}
\end{proof}

% from §50 helper lemma 155: integers less than one unit apart are equal
\begin{lemma}\label{integers_strictly_within_unit_equal_50}
    Assume $a,b\in\integers$.
    Assume $\abs{a-b}<1$.
    Then $a=b$.
\end{lemma}
\begin{proof}
    $a,b,a-b,b-a,1\in\reals$ by
        \cref{real_integers,subseteq,real_sub,real_add_inverse,real_add_closed,real_one_in_naturals,real_naturals_subset_reals}.
    $a<b$ or $a=b$ or $b<a$ by \cref{real_lt_trichotomy}.
    \begin{byCase}
    \caseOf{$a<b$.}
        Suppose not.
        Contradiction by
            \cref{integer_successor_gap_50,real_sub,real_add_inverse,real_add_closed,real_leq_add,real_add_assoc,real_add_comm,real_add_ident,real_leq_subst_left_local,real_leq_subst_right_local,real_leq_abs,real_abs_in_reals,real_abs_sub_swap,real_lt_subst_left,real_leq_lt_trans,real_lt_irreflexive}.
    \caseOf{$a=b$.}
        Follows.
    \caseOf{$b<a$.}
        Suppose not.
        Contradiction by
            \cref{integer_successor_gap_50,real_sub,real_add_inverse,real_add_closed,real_leq_add,real_add_assoc,real_add_comm,real_add_ident,real_leq_subst_left_local,real_leq_subst_right_local,real_leq_abs,real_abs_in_reals,real_leq_lt_trans,real_lt_irreflexive}.
    \end{byCase}
\end{proof}

% from §50 helper definition 44: a radius excluding all other cells at one level
\begin{definition}\label{cubical_protective_radius_50}
    $r$ is a cubical protective radius for $x$ at level $k$ in dimension $N$ iff
        $N,k\in\naturalsPlus$ and
        $x\in\realsPower{\initialSegment{N}}$ and
        $r\in\reals$ and $0<r$ and $r\leq1$ and
        for all $y\in\realsPower{\initialSegment{N}}$ if
            $\apply{\squareMetric{N}}{(x,y)}<r$
        then for all $D$ if
            $D$ has cubical level $k$ in dimension $N$ and $y\in D$
        then $x\in D$.
\end{definition}

% from §50 helper lemma 156: finitely many noninteger coordinates admit a common safe radius
\begin{lemma}\label{noninteger_coordinates_common_safe_radius_50}
    Assume $N\in\naturalsPlus$.
    Let $I=\initialSegment{N}$.
    Let $E=\realsPower{I}$.
    Assume $x\in E$.
    Let $J=\{i\in I\mid x(i)\notin\integers\}$.
    Then there exists $r\in\reals$ such that
        $0<r$ and $r\leq1$ and
        for all $i\in J$ there exists $n\in\integers$ such that
            $n<x(i)$ and $x(i)<n+1$ and
            $r<x(i)-n$ and $r<n+1-x(i)$.
\end{lemma}
\begin{proof}
    \begin{byCase}
    \caseOf{$J=\emptyset$.}
        $1\in\reals$ and $0<1$ and $1\leq1$ by
            \cref{real_one_in_naturals,real_naturals_subset_reals,subseteq,real_zero_lt_one}.
        Follows by \cref{emptyset}.
    \caseOf{$J\neq\emptyset$.}
        $I$ is finite by \cref{initial_segment_finite}.
        $J\subseteq I$.
        $J$ is finite by \cref{subseteq_finite_is_finite}.
        Let $R=\{w\in J\times\reals\mid
            \text{there exist $i,t,n$ such that
            $i\in J$ and $t\in\reals$ and $n\in\integers$ and
            $w=(i,t)$ and
            $n<x(i)$ and $x(i)<n+1$ and
            $0<t$ and $t\leq1$ and
            $t\leq x(i)-n$ and $t\leq n+1-x(i)$}\}$.
        For all $w\in R$ there exist $i,t$ such that $w=(i,t)$.
        $R$ is a relation by \cref{relation}.
        Show that for all $i\in J$ there exists $t$ such that
            $i\mathrel{R}t$.
        \begin{subproof}
            Fix $i\in J$.
            $i\in I$ and $x(i)\notin\integers$ by \cref{subseteq}.
            $x(i)\in\reals$ by
                \cref{reals_power,tuple_power,funs_type_apply}.
            Take $n\in\integers$ such that
                $n\leq x(i)$ and $x(i)<n+1$ by
                \cref{real_consecutive_integers}.
            $n\neq x(i)$ by \cref{real_lt_subst_right}.
            $n<x(i)$ by \cref{real_leq_split}.
            $n\in\reals$ by \cref{real_integers,subseteq}.
            $x(i)-n,n+1-x(i),1\in\reals$ by
                \cref{real_sub,real_add_inverse,real_add_closed,real_one_in_naturals,real_naturals_subset_reals,subseteq}.
            $0<x(i)-n$ and $0<n+1-x(i)$ by
                \cref{real_sub,real_add_inverse,real_add_closed,real_add_ident,real_lt_add,real_add_comm,real_add_assoc,real_add_inverse,real_lt_subst_left,real_lt_subst_right}.
            Take $t\in\reals$ such that
                $0<t$ and $t\leq1$ and
                $t\leq x(i)-n$ and $t\leq n+1-x(i)$ by
                \cref{real_zero_lt_one,positive_radius_below_three_bounds_49}.
            $(i,t)\in R$ by
                \cref{subseteq,times_tuple_intro,pair_eq_iff,real_lt_subst_left,real_lt_subst_right}.
            Follows by \cref{relation}.
        \end{subproof}
        Take $h$ such that
            $h$ is a function on $J$ and
            for all $i\in J$ we have $i\mathrel{R}h(i)$ by
            \cref{choice_function}.
        Let $H=\img{h}{J}$.
        $H$ is finite by \cref{finite_image_function}.
        Take $i_0$ such that $i_0\in J$ by
            \cref{nonempty_inhabited}.
        $h(i_0)\in H$ by
            \cref{function_apply_elim,img_iff}.
        $H\neq\emptyset$ by \cref{emptyset}.
        For all $t\in H$ we have
            $t\in\reals$ and $0<t$ and $t\leq1$ by
            \cref{img_iff,function_apply_intro,subseteq,relation,pair_eq_iff,real_lt_subst_left,real_lt_subst_right}.
        $H\subseteq\reals$ by \cref{subseteq}.
        Take $s\in H$ such that $s$ is a minimum of $H$ by
            \cref{finite_real_minimum}.
        $s\in\reals$ and $0<s$ and $s\leq1$ by \cref{subseteq}.
        Let $r=s/(1+1)$.
        $r\in\reals$ and $0<r$ and $r+r=s$ by
            \cref{real_half_of_positive}.
        $r<s$ by \cref{real_half_lt_self_positive_49}.
        $1\in\reals$ by
            \cref{real_one_in_naturals,real_naturals_subset_reals,subseteq}.
        $r\leq s$ by \cref{real_lt_imp_leq}.
        $r\leq1$ by
            \cref{real_leq_trans}.
        For all $i\in J$ there exists $n\in\integers$ such that
            $n<x(i)$ and $x(i)<n+1$ and
            $r<x(i)-n$ and $r<n+1-x(i)$ by
            \cref{subseteq,relation,pair_eq_iff,function_apply_elim,img_iff,real_minimum,real_lt_leq_trans_local,real_integers,reals_power,tuple_power,funs_type_apply,real_sub,real_add_inverse,real_add_closed,real_one_in_naturals,real_naturals_subset_reals,real_lt_subst_left,real_lt_subst_right}.
        Follows.
    \end{byCase}
\end{proof}

% from §50 helper lemma 157: every point of a cell has a protective radius at its level
\begin{lemma}\label{cubical_protective_radius_exists_50}
    Assume $C$ has cubical level $k$ in dimension $N$.
    Assume $x\in C$.
    Then there exists $r\in\reals$ such that
        $r$ is a cubical protective radius for $x$ at level $k$ in dimension $N$.
\end{lemma}
\begin{proof}
    Let $I=\initialSegment{N}$.
    Let $E=\realsPower{I}$.
    Let $e=\squareMetric{N}$.
    $k\in\naturalsPlus$ and
        $C$ is an integer cubical cell in dimension $N$ by
        \cref{integer_cubical_level_50}.
    Take $a\in E$ such that
        $C=\{z\in E\mid
            \text{$a$ is cubically equivalent to $z$ in dimension $N$}\}$ by
        \cref{integer_cubical_cell_50}.
    $x\in E$ and
        $a$ is cubically equivalent to $x$ in dimension $N$ by
        \cref{subseteq}.
    $N\in\naturalsPlus$ by
        \cref{cubically_equivalent_50,cubically_integer_compatible_50}.
    Let $J_x=\{i\in I\mid x(i)\notin\integers\}$.
    Take $r\in\reals$ such that
        $0<r$ and $r\leq1$ and
        for all $i\in J_x$ there exists $n\in\integers$ such that
            $n<x(i)$ and $x(i)<n+1$ and
            $r<x(i)-n$ and $r<n+1-x(i)$ by
        \cref{noninteger_coordinates_common_safe_radius_50}.
    Show that for all $y\in E$ if $e((x,y))<r$ then
        for all $D$ if
            $D$ has cubical level $k$ in dimension $N$ and $y\in D$
        then $x\in D$.
    \begin{subproof}
        Fix $y\in E$.
        Assume $e((x,y))<r$.
        Fix $D$.
        Assume $D$ has cubical level $k$ in dimension $N$ and $y\in D$.
        $D$ is an integer cubical cell in dimension $N$ by
            \cref{integer_cubical_level_50}.
        Take $q\in\reals$ such that
            $((x,y),q)\in e$ and
            for all $i\in I$ we have $\abs{x(i)-y(i)}\leq q$ by
            \cref{square_metric_value_exists}.
        $q=e((x,y))$ by
            \cref{square_metric_function,function_apply_intro}.
        $q<r$ by \cref{real_lt_subst_left}.
        Let $J_y=\{i\in I\mid y(i)\notin\integers\}$.
        For all $i\in J_x$ we have $i\in J_y$ by
            \cref{subseteq,reals_power,tuple_power,funs_type_apply,real_sub,real_add_inverse,real_add_closed,real_abs_in_reals,real_leq_lt_trans,small_perturbation_stays_in_integer_interval_50,no_integer_between_successive_integers_50}.
        $J_x\subseteq J_y$ by \cref{subseteq}.
        $J_x=J_y$ by
            \cref{same_level_noninteger_coordinate_inclusion_equal_50}.
        For all $i\in I$ if $x(i)\in\integers$ then
            $y(i)=x(i)$ by
            \cref{subseteq,real_lt_subst_left,reals_power,tuple_power,funs_type_apply,real_sub,real_add_inverse,real_add_closed,real_abs_in_reals,real_leq_lt_trans,real_one_in_naturals,real_naturals_subset_reals,real_lt_leq_trans_local,integers_strictly_within_unit_equal_50}.
        $x$ is cubically integer-compatible with $y$ in dimension $N$ by
            \cref{cubically_integer_compatible_50}.
        For all $i\in I$ if $x(i)\notin\integers$ then
            there exists $n\in\integers$ such that
                $n<x(i)$ and $x(i)<n+1$ and
                $n<y(i)$ and $y(i)<n+1$ by
            \cref{subseteq,reals_power,tuple_power,funs_type_apply,real_sub,real_add_inverse,real_add_closed,real_abs_in_reals,real_leq_lt_trans,small_perturbation_stays_in_integer_interval_50}.
        $x$ is cubically noninteger-compatible with $y$ in dimension $N$ by
            \cref{cubically_noninteger_compatible_50}.
        $x$ is cubically equivalent to $y$ in dimension $N$ by
            \cref{cubically_equivalent_50}.
        $a$ is cubically equivalent to $y$ in dimension $N$ by
            \cref{cubical_equivalence_transitive_50}.
        $y\in C$ by \cref{subseteq}.
        $y\in C\inter D$ by \cref{inter_intro}.
        $C=D$ by
            \cref{integer_cubical_cells_common_point_equal_50}.
        Follows by
            \cref{real_lt_subst_left,real_lt_subst_right}.
    \end{subproof}
    $r$ is a cubical protective radius for $x$ at level $k$ in dimension $N$ by
        \cref{cubical_protective_radius_50}.
    Follows.
\end{proof}

% from §50 helper definition 45: the open expansion of a cubical cell
\begin{definition}\label{cubical_cell_expansion_50}
    $U$ is the cubical expansion of $C$ at level $k$ in dimension $N$ iff
        $C$ has cubical level $k$ in dimension $N$ and
        $U=\{z\in\realsPower{\initialSegment{N}}\mid
            \text{there exist $x,r$ such that $x\in C$ and $r\in\reals$ and
            $r$ is a cubical protective radius for $x$ at level $k$ in dimension $N$
            and
            $\apply{\squareMetric{N}}{(x,z)}<(r/(1+1))/(1+1)$}\}$.
\end{definition}

% from §50 helper lemma 158: every leveled cubical cell has an expansion
\begin{lemma}\label{cubical_cell_expansion_exists_50}
    Assume $C$ has cubical level $k$ in dimension $N$.
    Then there exists $U$ such that
        $U$ is the cubical expansion of $C$ at level $k$ in dimension $N$.
\end{lemma}
\begin{proof}
    Let $E=\realsPower{\initialSegment{N}}$.
    Let $U=\{z\in E\mid
        \text{there exist $x,r$ such that $x\in C$ and $r\in\reals$ and
        $r$ is a cubical protective radius for $x$ at level $k$ in dimension $N$
        and
        $\apply{\squareMetric{N}}{(x,z)}<(r/(1+1))/(1+1)$}\}$.
    Follows by \cref{cubical_cell_expansion_50}.
\end{proof}

% from §50 helper lemma 159: a cubical expansion contains its cell
\begin{lemma}\label{cubical_cell_subset_expansion_50}
    Assume $U$ is the cubical expansion of $C$ at level $k$ in dimension $N$.
    Then $C\subseteq U$.
\end{lemma}
\begin{proof}
    It suffices to show that for all $x\in C$ we have $x\in U$ by
        \cref{subseteq}.
    Fix $x\in C$.
    $C$ has cubical level $k$ in dimension $N$ by
        \cref{cubical_cell_expansion_50}.
    $C$ is an integer cubical cell in dimension $N$ by
        \cref{integer_cubical_level_50}.
    Let $I=\initialSegment{N}$.
    Let $E=\realsPower{I}$.
    Let $e=\squareMetric{N}$.
    Take $a\in E$ such that
        $C=\{z\in E\mid
            \text{$a$ is cubically equivalent to $z$ in dimension $N$}\}$ by
        \cref{integer_cubical_cell_50}.
    $x\in E$ by \cref{subseteq}.
    $N\in\naturalsPlus$ by
        \cref{integer_cubical_level_50,cubically_equivalent_50,cubically_integer_compatible_50,subseteq}.
    $e$ is a metric on $E$ by \cref{square_metric_is_metric}.
    Take $r\in\reals$ such that
        $r$ is a cubical protective radius for $x$ at level $k$ in dimension $N$ by
        \cref{cubical_protective_radius_exists_50}.
    $0<r$ by \cref{cubical_protective_radius_50}.
    Let $h=r/(1+1)$.
    Let $q=h/(1+1)$.
    $h\in\reals$ and $0<h$ by
        \cref{real_half_of_positive}.
    $q\in\reals$ and $0<q$ by
        \cref{real_half_of_positive}.
    $e((x,x))=0$ by
        \cref{metric_on,metric_zero_characterization}.
    $e((x,x))<q$ by
        \cref{real_lt_subst_left}.
    $x\in U$ by
        \cref{cubical_cell_expansion_50,subseteq}.
\end{proof}

% from §50 helper lemma 160: cubical expansions are open
\begin{lemma}\label{cubical_cell_expansion_open_50}
    Assume $U$ is the cubical expansion of $C$ at level $k$ in dimension $N$.
    Let $E=\realsPower{\initialSegment{N}}$.
    Let $S=\metricTopology{\squareMetric{N}}{E}$.
    Then $U$ is open in $E$ with respect to $S$.
\end{lemma}
\begin{proof}
    Let $I=\initialSegment{N}$.
    Let $e=\squareMetric{N}$.
    $C$ has cubical level $k$ in dimension $N$ by
        \cref{cubical_cell_expansion_50}.
    $C$ is an integer cubical cell in dimension $N$ by
        \cref{integer_cubical_level_50}.
    Take $a\in E$ such that
        $C=\{z\in E\mid
            \text{$a$ is cubically equivalent to $z$ in dimension $N$}\}$ by
        \cref{integer_cubical_cell_50}.
    Take $c$ such that $c\in C$ by
        \cref{integer_cubical_level_50}.
    $a$ is cubically equivalent to $c$ in dimension $N$ by
        \cref{subseteq}.
    $N\in\naturalsPlus$ by
        \cref{cubically_equivalent_50,cubically_integer_compatible_50}.
    $e$ is a metric on $E$ by \cref{square_metric_is_metric}.
    $U\subseteq E$ by
        \cref{cubical_cell_expansion_50,subseteq}.
    Show that for all $z\in U$ there exists $\delta\in\reals$ such that
        $0<\delta$ and $\metricBall{e}{E}{z}{\delta}\subseteq U$.
    \begin{subproof}
        Fix $z\in U$.
        Take $x,r$ such that
            $x\in C$ and $r\in\reals$ and
            $r$ is a cubical protective radius for $x$ at level $k$ in dimension $N$ and
            $e((x,z))<(r/(1+1))/(1+1)$ by
            \cref{cubical_cell_expansion_50,subseteq}.
        $x,z\in E$ by
            \cref{cubical_protective_radius_50,subseteq}.
        $0<r$ by \cref{cubical_protective_radius_50}.
        Let $h=r/(1+1)$.
        Let $q=h/(1+1)$.
        $h\in\reals$ and $0<h$ by
            \cref{real_half_of_positive}.
        $q\in\reals$ and $0<q$ by
            \cref{real_half_of_positive}.
        $e((x,z))\in\reals$ by
            \cref{metric_on,times_tuple_intro,funs_type_apply}.
        Let $\delta=q-e((x,z))$.
        $\delta\in\reals$ and $0<\delta$ by
            \cref{real_sub,real_add_inverse,real_add_closed,real_lt_imp_pos_diff}.
        Show that $\metricBall{e}{E}{z}{\delta}\subseteq U$.
        \begin{subproof}
            It suffices to show that for all
                $w\in\metricBall{e}{E}{z}{\delta}$ we have $w\in U$ by
                \cref{subseteq}.
            Fix $w\in\metricBall{e}{E}{z}{\delta}$.
            $w\in E$ and $e((z,w))<\delta$ by
                \cref{metric_ball}.
            $e((z,w))\in\reals$ by
                \cref{metric_on,times_tuple_intro,funs_type_apply}.
            $e((x,z))+e((z,w))<q$ by
                \cref{real_sub,real_add_inverse,real_add_closed,real_lt_add,real_add_assoc,real_add_comm,real_add_inverse,real_add_ident,real_lt_subst_left,real_lt_subst_right}.
            $e((x,w))\leq e((x,z))+e((z,w))$ by
                \cref{metric_on,metric_triangle_inequality}.
            $e((x,w))<q$ by
                \cref{metric_on,times_tuple_intro,funs_type_apply,real_add_closed,real_leq_lt_trans}.
            Follows by
                \cref{cubical_cell_expansion_50,subseteq}.
        \end{subproof}
        Follows.
    \end{subproof}
    Follows by \cref{metric_open_iff}.
\end{proof}

% from §50 helper lemma 161: distinct cells at one level have disjoint expansions
\begin{lemma}\label{same_level_cubical_expansions_disjoint_50}
    Assume $U$ is the cubical expansion of $C$ at level $k$ in dimension $N$.
    Assume $V$ is the cubical expansion of $D$ at level $k$ in dimension $N$.
    Assume $C\neq D$.
    Then $U$ is disjoint from $V$.
\end{lemma}
\begin{proof}
    Suppose not.
    Take $z$ such that $z\in U$ and $z\in V$ by
        \cref{disjoint}.
    Let $I=\initialSegment{N}$.
    Let $E=\realsPower{I}$.
    Let $e=\squareMetric{N}$.
    Take $x,r$ such that
        $x\in C$ and $r\in\reals$ and
        $r$ is a cubical protective radius for $x$ at level $k$ in dimension $N$ and
        $e((x,z))<(r/(1+1))/(1+1)$ by
        \cref{cubical_cell_expansion_50,subseteq}.
    Take $y,s$ such that
        $y\in D$ and $s\in\reals$ and
        $s$ is a cubical protective radius for $y$ at level $k$ in dimension $N$ and
        $e((y,z))<(s/(1+1))/(1+1)$ by
        \cref{cubical_cell_expansion_50,subseteq}.
    $N,k\in\naturalsPlus$ and $x\in E$ and
        $0<r$ and $r\leq1$ by
        \cref{cubical_protective_radius_50}.
    $y\in E$ and $0<s$ and $s\leq1$ by
        \cref{cubical_protective_radius_50}.
    $z\in E$ by
        \cref{cubical_cell_expansion_50,subseteq}.
    $e$ is a metric on $E$ by \cref{square_metric_is_metric}.
    $C$ has cubical level $k$ in dimension $N$ and
        $D$ has cubical level $k$ in dimension $N$ by
        \cref{cubical_cell_expansion_50}.
    $C$ is an integer cubical cell in dimension $N$ and
        $D$ is an integer cubical cell in dimension $N$ by
        \cref{integer_cubical_level_50}.
    Let $h_r=r/(1+1)$.
    Let $q_r=h_r/(1+1)$.
    Let $h_s=s/(1+1)$.
    Let $q_s=h_s/(1+1)$.
    $h_r,q_r,h_s,q_s\in\reals$ and
        $0<h_r$ and $0<q_r$ and $0<h_s$ and $0<q_s$ by
        \cref{real_half_of_positive}.
    $e((z,y))=e((y,z))$ by
        \cref{metric_on,metric_symmetric}.
    $e((z,y))<q_s$ by \cref{real_lt_subst_left}.
    $e((x,z)),e((z,y)),e((x,y))\in\reals$ by
        \cref{metric_on,times_tuple_intro,funs_type_apply}.
    $e((x,z))+e((z,y))<q_r+e((z,y))$ by
        \cref{real_lt_add}.
    $q_r+e((z,y))<q_r+q_s$ by
        \cref{real_lt_add,real_add_comm,real_lt_subst_left,real_lt_subst_right}.
    $e((x,z))+e((z,y)),q_r+e((z,y)),q_r+q_s\in\reals$ by
        \cref{real_add_closed}.
    $e((x,z))+e((z,y))<q_r+q_s$ by
        \cref{real_lt_trans}.
    $e((x,y))\leq e((x,z))+e((z,y))$ by
        \cref{metric_on,metric_triangle_inequality}.
    $e((x,y))<q_r+q_s$ by
        \cref{real_add_closed,real_leq_lt_trans}.
    \begin{byCase}
    \caseOf{$r\leq s$.}
        $1+1\in\naturalsPlus$ by
            \cref{real_one_in_naturals,real_naturals_closed_succ}.
        $h_r\leq h_s$ by
            \cref{real_div_leq_same_denom}.
        $q_r\leq q_s$ by
            \cref{real_div_leq_same_denom}.
        $q_r+q_s\leq q_s+q_s$ by
            \cref{real_leq_add}.
        $q_s+q_s=h_s$ by \cref{real_half_plus_half}.
        $h_s<s$ by \cref{real_half_lt_self_positive_49}.
        $q_r+q_s<s$ by
            \cref{real_lt_subst_left,real_leq_lt_trans}.
        $e((x,y))<s$ by \cref{real_lt_trans}.
        $e((y,x))=e((x,y))$ by
            \cref{metric_on,metric_symmetric}.
        $e((y,x))<s$ by \cref{real_lt_subst_left}.
        $y\in C$ by
            \cref{cubical_protective_radius_50}.
        $y\in C\inter D$ by \cref{inter_intro}.
        $C=D$ by
            \cref{integer_cubical_cells_common_point_equal_50}.
        Contradiction.
    \caseOf{$s<r$.}
        $1+1\in\naturalsPlus$ by
            \cref{real_one_in_naturals,real_naturals_closed_succ}.
        $h_s\leq h_r$ by
            \cref{real_lt_imp_leq,real_div_leq_same_denom}.
        $q_s\leq q_r$ by
            \cref{real_div_leq_same_denom}.
        $q_s+q_r\leq q_r+q_r$ by
            \cref{real_leq_add}.
        $q_r+q_s=q_s+q_r$ by \cref{real_add_comm}.
        $q_r+q_s\leq q_r+q_r$ by
            \cref{real_leq_subst_left_local}.
        $q_r+q_r=h_r$ by \cref{real_half_plus_half}.
        $h_r<r$ by \cref{real_half_lt_self_positive_49}.
        $q_r+q_s<r$ by
            \cref{real_lt_subst_left,real_leq_lt_trans}.
        $e((x,y))<r$ by \cref{real_lt_trans}.
        $x\in D$ by
            \cref{cubical_protective_radius_50}.
        $x\in C\inter D$ by \cref{inter_intro}.
        $C=D$ by
            \cref{integer_cubical_cells_common_point_equal_50}.
        Contradiction.
    \end{byCase}
\end{proof}

% from §50 helper lemma 162: two points in one open unit interval differ by at most one
\begin{lemma}\label{unit_integer_interval_absolute_difference_bound_50}
    Assume $n\in\integers$ and $x,y\in\reals$.
    Assume $n<x$ and $x<n+1$.
    Assume $n<y$ and $y<n+1$.
    Then $\abs{x-y}\leq1$.
\end{lemma}
\begin{proof}
    $n,1,n+1\in\reals$ by
        \cref{real_integers,real_one_in_naturals,real_naturals_subset_reals,subseteq,real_add_closed}.
    $y-1,y+1\in\reals$ by
        \cref{real_sub,real_add_inverse,real_add_closed}.
    $y-1<n$ by
        \cref{real_sub,real_add_inverse,real_add_closed,real_lt_add,real_add_assoc,real_add_comm,real_add_inverse,real_add_ident,real_lt_subst_left,real_lt_subst_right}.
    $y-1<x$ by \cref{real_lt_trans}.
    $y-1\leq x$ by \cref{real_lt_imp_leq}.
    $n+1<y+1$ by \cref{real_lt_add}.
    $x<y+1$ by \cref{real_lt_trans}.
    $x\leq y+1$ by \cref{real_lt_imp_leq}.
    $0\leq1$ by
        \cref{real_zero_lt_one,real_lt_imp_leq}.
    Follows by \cref{real_between_implies_abs_difference_50}.
\end{proof}

% from §50 helper lemma 163: an integer cubical cell has square-metric diameter at most one
\begin{lemma}\label{integer_cubical_cell_metric_bound_50}
    Assume $C$ is an integer cubical cell in dimension $N$.
    Assume $x,y\in C$.
    Then $\apply{\squareMetric{N}}{(x,y)}\leq1$.
\end{lemma}
\begin{proof}
    Let $I=\initialSegment{N}$.
    Let $E=\realsPower{I}$.
    Take $a\in E$ such that
        $C=\{z\in E\mid
            \text{$a$ is cubically equivalent to $z$ in dimension $N$}\}$ by
        \cref{integer_cubical_cell_50}.
    $x,y\in E$ and
        $a$ is cubically equivalent to $x$ in dimension $N$ and
        $a$ is cubically equivalent to $y$ in dimension $N$ by
        \cref{subseteq}.
    $N\in\naturalsPlus$ by
        \cref{cubically_equivalent_50,cubically_integer_compatible_50}.
    $1\in\reals$ by
        \cref{real_one_in_naturals,real_naturals_subset_reals,subseteq}.
    For all $i\in I$ we have $\abs{x(i)-y(i)}\leq1$ by
        \cref{cubically_equivalent_50,cubically_integer_compatible_50,cubically_noninteger_compatible_50,real_lt_subst_left,real_lt_subst_right,reals_power,tuple_power,funs_type_apply,real_sub,real_add_inverse,real_add_eq_subst,real_add_ident,real_abs_nonneg,real_zero_lt_one,real_lt_imp_leq,real_leq_subst_left_local,successive_integer_interval_unique_50,unit_integer_interval_absolute_difference_bound_50}.
    Follows by \cref{square_metric_from_coordinate_bounds_50}.
\end{proof}

% from §50 helper lemma 164: a cubical expansion has uniformly bounded diameter
\begin{lemma}\label{cubical_cell_expansion_metric_bound_50}
    Assume $U$ is the cubical expansion of $C$ at level $k$ in dimension $N$.
    Assume $z,w\in U$.
    Then $\apply{\squareMetric{N}}{(z,w)}<(1+1)+1$.
\end{lemma}
\begin{proof}
    Let $I=\initialSegment{N}$.
    Let $E=\realsPower{I}$.
    Let $e=\squareMetric{N}$.
    Take $x,r$ such that
        $x\in C$ and $r\in\reals$ and
        $r$ is a cubical protective radius for $x$ at level $k$ in dimension $N$ and
        $e((x,z))<(r/(1+1))/(1+1)$ by
        \cref{cubical_cell_expansion_50,subseteq}.
    Take $y,s$ such that
        $y\in C$ and $s\in\reals$ and
        $s$ is a cubical protective radius for $y$ at level $k$ in dimension $N$ and
        $e((y,w))<(s/(1+1))/(1+1)$ by
        \cref{cubical_cell_expansion_50,subseteq}.
    $N,k\in\naturalsPlus$ and $x,y,z,w\in E$ and
        $0<r$ and $r\leq1$ and $0<s$ and $s\leq1$ by
        \cref{cubical_protective_radius_50,cubical_cell_expansion_50,subseteq}.
    $e$ is a metric on $E$ by \cref{square_metric_is_metric}.
    $C$ is an integer cubical cell in dimension $N$ by
        \cref{cubical_cell_expansion_50,integer_cubical_level_50}.
    Let $h_r=r/(1+1)$.
    Let $q_r=h_r/(1+1)$.
    Let $h_s=s/(1+1)$.
    Let $q_s=h_s/(1+1)$.
    $h_r,q_r,h_s,q_s\in\reals$ and
        $0<h_r$ and $0<q_r$ and $0<h_s$ and $0<q_s$ by
        \cref{real_half_of_positive}.
    $q_r<h_r$ and $h_r<r$ and $q_s<h_s$ and $h_s<s$ by
        \cref{real_half_lt_self_positive_49}.
    $e((z,x)),e((x,y)),e((z,y)),e((y,w)),e((z,w))\in\reals$ by
        \cref{metric_on,times_tuple_intro,funs_type_apply}.
    $1\in\reals$ and $1+1\in\reals$ by
        \cref{real_one_in_naturals,real_naturals_subset_reals,subseteq,real_add_closed}.
    $e((z,x))=e((x,z))$ by
        \cref{metric_on,metric_symmetric}.
    $e((x,z))<q_r$ and $e((y,w))<q_s$ by
        \cref{real_lt_subst_right}.
    $e((z,x))<q_r$ by
        \cref{real_lt_subst_left}.
    $e((z,x))<h_r$ and $e((z,x))<r$ and $e((z,x))<1$ by
        \cref{real_lt_trans,real_lt_leq_trans_local}.
    $e((y,w))<h_s$ and $e((y,w))<s$ and $e((y,w))<1$ by
        \cref{real_lt_trans,real_lt_leq_trans_local}.
    $e((x,y))\leq1$ by
        \cref{integer_cubical_cell_metric_bound_50}.
    $e((z,x))+e((x,y))<1+1$ by
        \cref{real_lt_add_leq_two_49}.
    $e((z,y))\leq e((z,x))+e((x,y))$ by
        \cref{metric_on,metric_triangle_inequality}.
    $e((z,y))<1+1$ by
        \cref{real_add_closed,real_leq_lt_trans}.
    $e((z,y))+e((y,w))<(1+1)+1$ by
        \cref{real_lt_add_two}.
    $e((z,w))\leq e((z,y))+e((y,w))$ by
        \cref{metric_on,metric_triangle_inequality}.
    Follows by
        \cref{real_add_closed,real_leq_lt_trans}.
\end{proof}

% from §50 helper lemma 165: integer cubical expansions form a bounded-order cover
\begin{lemma}\label{integer_cubical_expansion_cover_50}
    Assume $N\in\naturalsPlus$.
    Let $E=\realsPower{\initialSegment{N}}$.
    Let $S=\metricTopology{\squareMetric{N}}{E}$.
    Then there exists $A$ such that
        $A$ is an open covering of $E$ with respect to $S$ and
        $A$ has order at most $N+1$ in $E$ and
        for all $U\in A$ for all $x,y\in U$ we have
            $\apply{\squareMetric{N}}{(x,y)}<(1+1)+1$.
\end{lemma}
\begin{proof}
    Let $J=\initialSegment{N+1}$.
    Let $A=\{U\in\pow{E}\mid
        \text{there exist $C,k$ such that $k\in J$ and
            $U$ is the cubical expansion of $C$ at level $k$ in dimension $N$}\}$.
    $A\subseteq\pow{E}$ by \cref{subseteq}.
    Show that for all $x\in E$ we have $x\in\unions{A}$.
    \begin{subproof}
        Fix $x\in E$.
        Let $C=\{y\in E\mid
            \text{$x$ is cubically equivalent to $y$ in dimension $N$}\}$.
        $x\in C$ and $C$ is an integer cubical cell in dimension $N$ by
            \cref{generated_integer_cubical_cell_50}.
        Take $k\in J$ such that
            $C$ has cubical level $k$ in dimension $N$ by
            \cref{integer_cubical_cell_level_exists_50}.
        Take $U$ such that
            $U$ is the cubical expansion of $C$ at level $k$ in dimension $N$ by
            \cref{cubical_cell_expansion_exists_50}.
        Follows by
            \cref{cubical_cell_subset_expansion_50,cubical_cell_expansion_50,powerset_intro,subseteq,unions_iff}.
    \end{subproof}
    $A$ is a covering of $E$ by \cref{covering,subseteq}.
    For all $U\in A$ we have
        $U$ is open in $E$ with respect to $S$ by
        \cref{subseteq,cubical_cell_expansion_open_50}.
    $A$ is an open covering of $E$ with respect to $S$ by
        \cref{open_covering}.
    $N+1\in\naturalsPlus$ by
        \cref{real_naturals_closed_succ}.
    Show that for all $z\in E$ we have
        $z$ has multiplicity at most $N+1$ in $A$.
    \begin{subproof}
        Fix $z\in E$.
        Suppose not.
        Take $F,g$ such that
            $F\subseteq A$ and
            $g$ is a bijection from $\initialSegment{(N+1)+1}$ to $F$ and
            for all $U\in F$ we have $z\in U$ by
            \cref{point_multiplicity_at_most_50}.
        Let $R=\{w\in F\times J\mid
            \text{there exists $C$ such that
                $\fst{w}$ is the cubical expansion of $C$ at level $\snd{w}$
                    in dimension $N$}\}$.
        $R$ is a relation by
            \cref{relation,subseteq,times_elem_is_tuple}.
        For all $U\in F$ there exists $k$ such that
            $U\mathrel{R}k$ by
            \cref{subseteq,times_tuple_intro,fst_eq,snd_eq}.
        Take $h$ such that
            $h$ is a function on $F$ and
            for all $U\in F$ we have $U\mathrel{R}h(U)$ by
            \cref{choice_function}.
        For all $U\in F$ we have $h(U)\in J$ by
            \cref{choice_function,times_tuple_elim}.
        Show that $h$ is injective.
        \begin{subproof}
            It suffices to show that for all $U,V$ such that
                $U\in\dom{h}$ and $V\in\dom{h}$ and $h(U)=h(V)$
            we have $U=V$ by \cref{injective_function}.
            Fix $U$.
            Fix $V$.
            Assume $U\in\dom{h}$ and $V\in\dom{h}$ and $h(U)=h(V)$.
            $U,V\in F$ by \cref{choice_function}.
            $U\mathrel{R}h(U)$ and $V\mathrel{R}h(V)$ by
                \cref{choice_function}.
            Take $C$ such that
                $U$ is the cubical expansion of $C$ at level $h(U)$
                    in dimension $N$ by
                \cref{subseteq,fst_eq,snd_eq}.
            Take $D$ such that
                $V$ is the cubical expansion of $D$ at level $h(V)$
                    in dimension $N$ by
                \cref{subseteq,fst_eq,snd_eq}.
            $z\in U$ and $z\in V$ by \cref{subseteq}.
            $C=D$ by
                \cref{same_level_cubical_expansions_disjoint_50,real_lt_subst_right,disjoint}.
            $U\subseteq V$ and $V\subseteq U$ by
                \cref{cubical_cell_expansion_50,subseteq,real_lt_subst_left,real_lt_subst_right}.
            Follows by \cref{subseteq_antisymmetric}.
        \end{subproof}
        Let $H=\img{h}{F}$.
        $H\subseteq J$ by
            \cref{img_iff,function_apply_iff,subseteq}.
        For all $U\in\dom{h}$ we have $h(U)\in H$ by
            \cref{img_of_function_intro,choice_function,inter_intro}.
        $h\in\funs{F}{H}$ by
            \cref{choice_function,funs_intro}.
        $h\in\surj{F}{H}$ by
            \cref{funs_surj_iff,img_dom}.
        $h$ is a bijection from $F$ to $H$ by
            \cref{bijections}.
        Let $q=h\circ g$.
        $q$ is a bijection from $\initialSegment{(N+1)+1}$ to $H$ by
            \cref{bijection_circ}.
        $J$ has at most $N+1$ elements by
            \cref{initial_segment_at_most_self_50}.
        Contradiction by
            \cref{finite_set_has_at_most_50}.
    \end{subproof}
    $A$ has order at most $N+1$ in $E$ by
        \cref{covering_order_at_most_50}.
    For all $U\in A$ for all $x,y\in U$ we have
        $\apply{\squareMetric{N}}{(x,y)}<(1+1)+1$ by
        \cref{subseteq,cubical_cell_expansion_metric_bound_50}.
    Follows.
\end{proof}

% from §50 helper definition 46: coordinatewise scalar dilation of finite euclidean space
\begin{definition}\label{euclidean_scalar_dilation_50}
    $f$ is the scalar dilation by $c$ in dimension $N$ iff
        $c\in\reals$ and
        $f=\{w\in
            \realsPower{\initialSegment{N}}\times
            \realsPower{\initialSegment{N}}\mid
            \text{there exists $x\in\realsPower{\initialSegment{N}}$ such that
                $w=(x,\coordScalarSequence{N}{c}{x})$}\}$.
\end{definition}

% from §50 helper lemma 166: every positive scalar induces a euclidean bijection
\begin{lemma}\label{positive_euclidean_scalar_dilation_bijection_50}
    Assume $N\in\naturalsPlus$.
    Assume $c\in\reals$ and $0<c$.
    Let $E=\realsPower{\initialSegment{N}}$.
    Then there exists $f$ such that
        $f$ is the scalar dilation by $c$ in dimension $N$ and
        $f$ is a bijection from $E$ to $E$.
\end{lemma}
\begin{proof}
    Let $I=\initialSegment{N}$.
    Let $f=\{w\in E\times E\mid
        \text{there exists $x\in E$ such that
            $w=(x,\coordScalarSequence{N}{c}{x})$}\}$.
    $f$ is the scalar dilation by $c$ in dimension $N$ by
        \cref{euclidean_scalar_dilation_50}.
    $f$ is a relation by
        \cref{relation,subseteq,times_elem_is_tuple}.
    For all $x,y,z$ such that
        $(x,y)\in f$ and $(x,z)\in f$ we have $y=z$ by
        \cref{subseteq,pair_eq_iff,real_lt_subst_left,real_lt_subst_right}.
    $f$ is right-unique by \cref{rightunique}.
    $f$ is a function by \cref{relation,rightunique}.
    $\dom{f}\subseteq E$ by
        \cref{dom_iff,subseteq,times_tuple_elim}.
    $E\subseteq\dom{f}$ by
        \cref{coord_scalar_sequence_function,reals_power,tuple_power,times_tuple_intro,subseteq,dom_intro}.
    Thus $\dom{f}=E$ by \cref{subseteq_antisymmetric}.
    For all $x\in\dom{f}$ we have $f(x)\in E$ by
        \cref{function_apply_elim,subseteq,times_tuple_elim}.
    $f\in\funs{E}{E}$ by \cref{funs_intro}.
    For all $x\in E$ we have
        $f(x)=\coordScalarSequence{N}{c}{x}$ by
        \cref{coord_scalar_sequence_function,reals_power,tuple_power,times_tuple_intro,subseteq,funs_is_function,function_apply_intro}.
    Show that $f$ is injective.
    \begin{subproof}
        It suffices to show that for all $x,y$ such that
            $x\in\dom{f}$ and $y\in\dom{f}$ and $f(x)=f(y)$
        we have $x=y$ by \cref{injective_function}.
        Fix $x$.
        Fix $y$.
        Assume $x\in\dom{f}$ and $y\in\dom{f}$ and $f(x)=f(y)$.
        $x,y\in E$ by \cref{funs}.
        Show that for all $i\in I$ we have $x(i)=y(i)$.
        \begin{subproof}
            Fix $i\in I$.
            $f(x),f(y)\in E$ by \cref{funs_type_apply}.
            $\apply{\apply{f}{x}}{i}=c\rmul x(i)$ and
                $\apply{\apply{f}{y}}{i}=c\rmul y(i)$ by
                \cref{subseteq,coord_scalar_sequence,real_lt_subst_left,real_lt_subst_right,reals_power,tuple_power,funs_is_function,function_apply_intro}.
            $c\rmul x(i)=c\rmul y(i)$ by
                \cref{real_lt_subst_left,real_lt_subst_right}.
            $x(i),y(i)\in\reals$ by
                \cref{reals_power,tuple_power,funs_type_apply}.
            Follows by \cref{positive_real_factor_cancel}.
        \end{subproof}
        Follows by
            \cref{reals_power,tuple_power,funs_is_function,funs,functions_eq_iff}.
    \end{subproof}
    $\img{f}{E}\subseteq E$ by
        \cref{funs_ran,img_subseteq_ran,subseteq_transitive}.
    Show that for all $y\in E$ we have $y\in\img{f}{E}$.
    \begin{subproof}
            Fix $y\in E$.
            Let $x=\coordScalarSequence{N}{\inv{c}}{y}$.
            $c\neq0$ and $\inv{c}\in\reals$ by
                \cref{real_pos_nonzero,real_mul_inverse}.
            $x\in E$ by
                \cref{coord_scalar_sequence_function,reals_power,tuple_power}.
            $f(x),y\in E$ by \cref{funs_type_apply}.
                Show that for all $i\in I$ we have
                    $\apply{\apply{f}{x}}{i}=y(i)$.
                \begin{subproof}
                    Fix $i\in I$.
                    $(i,\inv{c}\rmul y(i))\in x$ by
                        \cref{coord_scalar_sequence}.
                    $x(i)=\inv{c}\rmul y(i)$ by
                        \cref{reals_power,tuple_power,funs_is_function,function_apply_intro}.
                    $\apply{\apply{f}{x}}{i}=c\rmul x(i)$ by
                        \cref{subseteq,coord_scalar_sequence,real_lt_subst_left,real_lt_subst_right,reals_power,tuple_power,funs_is_function,function_apply_intro}.
                    $y(i)\in\reals$ by
                        \cref{reals_power,tuple_power,funs_type_apply}.
                    Follows by
                        \cref{real_mul_assoc,real_mul_inverse,real_mul_ident,real_lt_subst_left,real_lt_subst_right}.
                \end{subproof}
            Thus $f(x)=y$ by
                \cref{reals_power,tuple_power,funs_is_function,funs,functions_eq_iff}.
            Follows by
                \cref{img_of_function_intro,inter_intro}.
    \end{subproof}
    Thus $E\subseteq\img{f}{E}$ by \cref{subseteq}.
    Thus $\img{f}{E}=E$ by \cref{subseteq_antisymmetric}.
    $f\in\surj{E}{E}$ by
        \cref{funs_surj_iff,img_dom}.
    $f$ is a bijection from $E$ to $E$ by
        \cref{bijections}.
    Thus there exists $g$ such that
        $g$ is the scalar dilation by $c$ in dimension $N$ and
        $g$ is a bijection from $E$ to $E$.
\end{proof}

% from §50 helper lemma 167: coordinate formula for a scalar dilation
\begin{lemma}\label{euclidean_scalar_dilation_coordinate_50}
    Assume $N\in\naturalsPlus$.
    Assume $c\in\reals$.
    Let $E=\realsPower{\initialSegment{N}}$.
    Assume $f$ is the scalar dilation by $c$ in dimension $N$.
    Assume $f\in\funs{E}{E}$.
    Assume $x\in E$ and $i\in\initialSegment{N}$.
    Then $\apply{\apply{f}{x}}{i}=c\rmul x(i)$.
\end{lemma}
\begin{proof}
    $\coordScalarSequence{N}{c}{x}\in E$ by
        \cref{coord_scalar_sequence_function,reals_power,tuple_power}.
    $(x,\coordScalarSequence{N}{c}{x})\in f$ by
        \cref{euclidean_scalar_dilation_50,times_tuple_intro,subseteq}.
    $f(x)=\coordScalarSequence{N}{c}{x}$ by
        \cref{funs_is_function,function_apply_intro}.
    $(i,c\rmul x(i))\in\coordScalarSequence{N}{c}{x}$ by
        \cref{coord_scalar_sequence}.
    $(i,c\rmul x(i))\in f(x)$ by
        \cref{real_lt_subst_left,real_lt_subst_right}.
    $f(x)\in E$ by \cref{funs_type_apply}.
    Follows by
        \cref{reals_power,tuple_power,funs_is_function,function_apply_intro}.
\end{proof}

% from §50 helper lemma 168: positive dilations scale the square metric
\begin{lemma}\label{euclidean_scalar_dilation_metric_50}
    Assume $N\in\naturalsPlus$.
    Assume $c\in\reals$ and $0<c$.
    Let $I=\initialSegment{N}$.
    Let $E=\realsPower{I}$.
    Let $e=\squareMetric{N}$.
    Assume $f$ is the scalar dilation by $c$ in dimension $N$.
    Assume $f\in\funs{E}{E}$.
    Assume $x,y\in E$.
    Then $e((f(x),f(y)))=c\rmul e((x,y))$.
\end{lemma}
\begin{proof}
    Let $a=e((x,y))$.
    Let $b=e((f(x),f(y)))$.
    $e$ is a metric on $E$ by \cref{square_metric_is_metric}.
    $f(x),f(y)\in E$ by \cref{funs_type_apply}.
    $a,b\in\reals$ by
        \cref{metric_on,times_tuple_intro,funs_type_apply}.
    $c\rmul a\in\reals$ by \cref{real_mul_closed}.
    Take $r\in\reals$ such that
        $((x,y),r)\in e$ and
        for all $i\in I$ we have $\abs{x(i)-y(i)}\leq r$ by
        \cref{square_metric_value_exists}.
    $r=a$ by
        \cref{metric_on,funs_is_function,function_apply_intro}.
    Show that for all $i\in I$ we have
        $\abs{\apply{\apply{f}{x}}{i}-\apply{\apply{f}{y}}{i}}
            =c\rmul\abs{x(i)-y(i)}$.
    \begin{subproof}
        Fix $i\in I$.
        $\apply{\apply{f}{x}}{i}=c\rmul x(i)$ and
            $\apply{\apply{f}{y}}{i}=c\rmul y(i)$ by
            \cref{euclidean_scalar_dilation_coordinate_50}.
        $x(i),y(i)\in\reals$ and
            $\apply{\apply{f}{x}}{i},\apply{\apply{f}{y}}{i}\in\reals$ by
            \cref{reals_power,tuple_power,funs_type_apply}.
        $x(i)-y(i)\in\reals$ and
            $\apply{\apply{f}{x}}{i}-\apply{\apply{f}{y}}{i}\in\reals$ by
            \cref{real_sub,real_add_inverse,real_add_closed}.
        $(c\rmul x(i))-(c\rmul y(i))=c\rmul(x(i)-y(i))$ by
            \cref{real_sub,real_add_inverse,real_mul_neg_right,real_distrib}.
        $\abs{c}=c$ by
            \cref{real_lt_imp_leq,real_abs_nonneg}.
        $\abs{\apply{\apply{f}{x}}{i}-\apply{\apply{f}{y}}{i}}
            =\abs{(c\rmul x(i))-(c\rmul y(i))}$ by
            \cref{real_lt_subst_left,real_lt_subst_right}.
        $\abs{(c\rmul x(i))-(c\rmul y(i))}
            =\abs{c\rmul(x(i)-y(i))}$ by
            \cref{real_lt_subst_left,real_lt_subst_right}.
        $\abs{c\rmul(x(i)-y(i))}
            =\abs{c}\rmul\abs{x(i)-y(i)}$ by
            \cref{real_abs_mul}.
        $\abs{c}\rmul\abs{x(i)-y(i)}
            =c\rmul\abs{x(i)-y(i)}$ by
            \cref{real_abs_in_reals,real_mul_eq_subst}.
        Follows by
            \cref{real_lt_subst_left,real_lt_subst_right}.
    \end{subproof}
    For all $i\in I$ we have
        $\abs{\apply{\apply{f}{x}}{i}-\apply{\apply{f}{y}}{i}}
            \leq c\rmul a$ by
        \cref{real_leq_subst_right_local,reals_power,tuple_power,funs_type_apply,real_sub,real_add_inverse,real_add_closed,real_abs_in_reals,real_leq_mul_pos_left_49,real_leq_subst_left_local}.
    $b\leq c\rmul a$ by
        \cref{square_metric_from_coordinate_bounds_50}.
    $(x,y)\in E\times E$ and $(x,y)\in\dom{e}$ by
        \cref{times_tuple_intro,metric_on,funs}.
    $((x,y),a)\in e$ by
        \cref{metric_on,funs_is_function,function_apply_elim}.
    $a\in\coordDiffSet{N}{x}{y}$ by
        \cref{square_metric,pair_eq_iff}.
    Take $j\in I$ such that $a=\abs{x(j)-y(j)}$ by
        \cref{coord_diff_set}.
    Take $s\in\reals$ such that
        $((f(x),f(y)),s)\in e$ and
        for all $i\in I$ we have
            $\abs{\apply{\apply{f}{x}}{i}-\apply{\apply{f}{y}}{i}}\leq s$ by
        \cref{square_metric_value_exists}.
    $s=b$ by
        \cref{metric_on,funs_is_function,function_apply_intro}.
    $\abs{\apply{\apply{f}{x}}{j}-\apply{\apply{f}{y}}{j}}
        =c\rmul\abs{x(j)-y(j)}$ by \cref{subseteq}.
    $c\rmul a\leq b$ by
        \cref{real_lt_subst_left,real_lt_subst_right,real_leq_subst_left_local,real_leq_subst_right_local}.
    Follows by \cref{real_leq_antisym}.
\end{proof}

% from §50 helper lemma 169: positive euclidean scalar dilations are continuous
\begin{lemma}\label{positive_euclidean_scalar_dilation_continuous_50}
    Assume $N\in\naturalsPlus$.
    Assume $c\in\reals$ and $0<c$.
    Let $E=\realsPower{\initialSegment{N}}$.
    Let $e=\squareMetric{N}$.
    Let $S=\metricTopology{e}{E}$.
    Assume $f$ is the scalar dilation by $c$ in dimension $N$.
    Assume $f$ is a bijection from $E$ to $E$.
    Then $f$ is continuous from $E$ to $E$ with respect to $S$ and $S$.
\end{lemma}
\begin{proof}
    $e$ is a metric on $E$ by \cref{square_metric_is_metric}.
    $f\in\funs{E}{E}$ by \cref{bijections_to_funs}.
    Show that for all $x\in E$ we have
        for all $\epsilon\in\reals$ such that $0<\epsilon$
        there exists $\delta\in\reals$ such that $0<\delta$ and
        for all $y\in E$ if $e((x,y))<\delta$
        then $e((f(x),f(y)))<\epsilon$.
    \begin{subproof}
        Fix $x\in E$.
        Fix $\epsilon\in\reals$.
        Assume $0<\epsilon$.
        Let $\delta=\epsilon/c$.
        $c\neq0$ and $\inv{c}\in\reals$ and $0<\inv{c}$ by
            \cref{real_pos_nonzero,real_mul_inverse,real_inv_pos}.
        $\delta=\epsilon\rmul\inv{c}$ and $\delta\in\reals$ by
            \cref{real_div,real_mul_closed}.
        $0<\delta$ by
            \cref{real_add_ident,real_lt_mul_pos,real_mul_zero,real_lt_subst_left,real_lt_subst_right}.
        For all $y\in E$ if $e((x,y))<\delta$
            then $e((f(x),f(y)))<\epsilon$ by
            \cref{metric_on,times_tuple_intro,funs_type_apply,real_lt_mul_pos,real_mul_comm,real_lt_subst_left,real_lt_subst_right,real_div,real_mul_assoc,real_mul_inverse,real_mul_ident,euclidean_scalar_dilation_metric_50}.
        Thus there exists $\eta\in\reals$ such that $0<\eta$ and
            for all $y\in E$ if $e((x,y))<\eta$
            then $e((f(x),f(y)))<\epsilon$.
    \end{subproof}
    Follows by \cref{metric_continuity_eps_delta}.
\end{proof}

% from §50 helper lemma 170: euclidean space has arbitrarily fine bounded-order covers
\begin{lemma}\label{fine_integer_cubical_expansion_cover_50}
    Assume $N\in\naturalsPlus$.
    Let $E=\realsPower{\initialSegment{N}}$.
    Let $e=\squareMetric{N}$.
    Let $S=\metricTopology{e}{E}$.
    Assume $\epsilon\in\reals$ and $0<\epsilon$.
    Then there exists $B$ such that
        $B$ is an open covering of $E$ with respect to $S$ and
        $B$ has order at most $N+1$ in $E$ and
        for all $V\in B$ for all $x,y\in V$ we have
            $e((x,y))<\epsilon$.
\end{lemma}
\begin{proof}
    Let $K=(1+1)+1$.
    $1\in\reals$ and $0<1$ by
        \cref{real_one_in_naturals,real_naturals_subset_reals,subseteq,real_zero_lt_one}.
    $1+1,K\in\reals$ and $0<1+1$ and $0<K$ by
        \cref{real_add_closed,real_pos_add}.
    Take $m\in\naturalsPlus$ such that
        $0<1/m$ and $K\rmul(1/m)\leq\epsilon$ by
        \cref{positive_mesh_step_scaled_bound_49}.
    $m\in\reals$ and $0<m$ by
        \cref{real_naturals_subset_reals,subseteq,naturalsplus_positive_real}.
    Take $f$ such that
        $f$ is the scalar dilation by $m$ in dimension $N$ and
        $f$ is a bijection from $E$ to $E$ by
        \cref{positive_euclidean_scalar_dilation_bijection_50}.
    $f\in\funs{E}{E}$ by \cref{bijections_to_funs}.
    $f$ is continuous from $E$ to $E$ with respect to $S$ and $S$ by
        \cref{positive_euclidean_scalar_dilation_continuous_50}.
    Take $A$ such that
        $A$ is an open covering of $E$ with respect to $S$ and
        $A$ has order at most $N+1$ in $E$ and
        for all $U\in A$ for all $u,v\in U$ we have
            $e((u,v))<K$ by
        \cref{integer_cubical_expansion_cover_50}.
    Let $B=\{V\in\pow{E}\mid
        \text{there exists $U\in A$ such that $V=\preimg{f}{U}$}\}$.
    $B$ has order at most $N+1$ in $E$ by
        \cref{bijective_pullback_preserves_order_bound_50}.
    $B\subseteq\pow{E}$ by \cref{subseteq}.
    For all $x\in E$ we have $x\in\unions{B}$ by
        \cref{funs_type_apply,open_covering,covering,subseteq,unions_iff,preimg_subseteq_dom,funs,powerset_intro,funs_is_function,real_lt_subst_left,real_lt_subst_right,function_apply_elim,preimg_iff}.
    Thus $B$ is a covering of $E$ by \cref{covering,subseteq}.
    For all $V\in B$ we have
        $V$ is open in $E$ with respect to $S$ by
        \cref{subseteq,open_covering,continuous,real_lt_subst_left,real_lt_subst_right}.
    $B$ is an open covering of $E$ with respect to $S$ by
        \cref{open_covering}.
    Show that for all $V\in B$ for all $x,y$ such that
        $x\in V$ and $y\in V$ we have $e((x,y))<\epsilon$.
    \begin{subproof}
        Fix $V\in B$.
        Fix $x$.
        Fix $y$.
        Assume $x\in V$ and $y\in V$.
        Take $U\in A$ such that $V=\preimg{f}{U}$ by
            \cref{subseteq}.
        $x,y\in E$ by
            \cref{powerset_elim,subseteq}.
        $f(x),f(y)\in U$ by
            \cref{preimg_iff,funs_is_function,function_apply_iff,real_lt_subst_left,real_lt_subst_right}.
        $m\rmul e((x,y))<K$ by
            \cref{subseteq,euclidean_scalar_dilation_metric_50,real_lt_subst_left}.
        $e((x,y))\in\reals$ and $m\rmul e((x,y))\in\reals$ by
            \cref{square_metric_is_metric,metric_on,times_tuple_intro,funs_type_apply,real_mul_closed}.
        $(m\rmul e((x,y)))\rmul(1/m)<K\rmul(1/m)$ by
            \cref{naturalsplus_reciprocal_real,real_lt_mul_pos,real_lt_subst_left}.
        $m\neq0$ and $\inv{m}\in\reals$ and $1/m=\inv{m}$ by
            \cref{real_pos_nonzero,real_mul_inverse,real_div,real_mul_ident}.
        $(m\rmul e((x,y)))\rmul(1/m)=e((x,y))$ by
            \cref{real_mul_assoc,real_mul_comm,real_mul_inverse,real_mul_ident,real_lt_subst_left,real_lt_subst_right}.
        $e((x,y))<K\rmul(1/m)$ by
            \cref{real_lt_subst_left}.
        Follows by
            \cref{real_mul_closed,naturalsplus_reciprocal_real,real_lt_leq_trans_local}.
    \end{subproof}
    Thus there exists $C$ such that
        $C$ is an open covering of $E$ with respect to $S$ and
        $C$ has order at most $N+1$ in $E$ and
        for all $V\in C$ for all $x,y\in V$ we have
            $e((x,y))<\epsilon$.
\end{proof}

% from §50 helper lemma 171: restricting the square metric to a subset gives a metric
\begin{lemma}\label{restricted_square_metric_is_metric_50}
    Assume $N\in\naturalsPlus$.
    Let $E=\realsPower{\initialSegment{N}}$.
    Let $e=\squareMetric{N}$.
    Assume $X\subseteq E$.
    Let $d=\restrl{e}{X\times X}$.
    Then $d$ is a metric on $X$.
\end{lemma}
\begin{proof}
    $e$ is a metric on $E$ by \cref{square_metric_is_metric}.
    $e\in\funs{E\times E}{\reals}$ and
        $e$ is a function and $\dom{e}=E\times E$ by
        \cref{metric_on,funs_elim,funs_is_function}.
    $X\times X\subseteq E\times E$ by
        \cref{times_subseteq_left,times_subseteq_right,subseteq_transitive}.
    $d$ is a function by
        \cref{restrl_of_function_is_function}.
    $\dom{d}=(E\times E)\inter(X\times X)$ by
        \cref{dom_of_restrl_eq_inter,pair_eq_iff}.
    $\dom{d}=X\times X$ by
        \cref{subseteq,inter_intro,restrl_dom,subseteq_antisymmetric}.
    For all $p\in\dom{d}$ we have $d(p)\in\reals$ by
        \cref{inter_elim_left,dom_of_restrl_eq_inter,funs_type_apply,restrl_of_function_apply}.
    $d\in\funs{X\times X}{\reals}$ by \cref{funs_intro}.
    For all $x,y$ such that $x\in X$ and $y\in X$ we have
        $d((x,y))=e((x,y))$ by
        \cref{times_tuple_intro,subseteq,restrl_of_function_apply}.
    $d$ is nonnegative on $X$ by
        \cref{subseteq,metric_on,metric_nonnegative}.
    $d$ is zero-characterized on $X$ by
        \cref{subseteq,metric_on,metric_zero_characterization}.
    $d$ is symmetric on $X$ by
        \cref{subseteq,metric_on,metric_symmetric}.
    $d$ satisfies the triangle inequality on $X$ by
        \cref{subseteq,metric_on,metric_triangle_inequality}.
    Follows by \cref{metric_on}.
\end{proof}

% from §50 helper lemma 45: cubical refinements bound compact euclidean dimension
\begin{lemma}\label{compact_euclidean_dimension_bound_principle_50}
    Assume $N\in\naturalsPlus$.
    Let $E=\realsPower{\initialSegment{N}}$.
    Let $S=\metricTopology{\squareMetric{N}}{E}$.
    Assume $X$ is a subspace of $E$ with respect to $S$.
    Let $T=\subspaceTopology{S}{X}$.
    Suppose $X$ is compact with respect to $T$.
    Then $N$ is a topological dimension bound for $X$ with respect to $T$.
\end{lemma}
\begin{proof}
    Let $e=\squareMetric{N}$.
    Let $d=\restrl{e}{X\times X}$.
    Let $T_d=\metricTopology{d}{X}$.
    $e$ is a metric on $E$ by \cref{square_metric_is_metric}.
    $X\subseteq E$ by \cref{subspace_of}.
    $e$ is a function and $\dom{e}=E\times E$ by
        \cref{metric_on,funs_elim,funs_is_function}.
    $X\times X\subseteq\dom{e}$ by
        \cref{times_subseteq_left,times_subseteq_right,subseteq_transitive,real_lt_subst_left,real_lt_subst_right}.
    $S$ is a topology on $E$ and $T$ is a topology on $X$ by
        \cref{metric_topology,metric_basis_is_basis,basis_topology_is_topology,subspace_topology_is_topology,subspace_of}.
    $d$ is a metric on $X$ by
        \cref{restricted_square_metric_is_metric_50}.
    $T=T_d$ by
        \cref{restricted_metric_induces_subspace_topology}.
    $X$ is compact with respect to $T_d$ by
        \cref{compact_eq_topology}.
    $N\in\naturals$ by
        \cref{naturals,union_intro_left}.
    Show that for all $\epsilon\in\reals$ such that $0<\epsilon$
        there exists $C$ such that
            $C$ is an open covering of $X$ with respect to $T_d$ and
            $C$ has order at most $N+1$ in $X$ and
            for all $U\in C$ for all $x,y\in U$ we have
                $d((x,y))<\epsilon$.
    \begin{subproof}
        Fix $\epsilon\in\reals$.
        Assume $0<\epsilon$.
        Take $B$ such that
            $B$ is an open covering of $E$ with respect to $S$ and
            $B$ has order at most $N+1$ in $E$ and
            for all $V\in B$ for all $x,y\in V$ we have
                $e((x,y))<\epsilon$ by
            \cref{fine_integer_cubical_expansion_cover_50}.
        Let $C=\{U\in\pow{X}\mid
            \text{there exists $V\in B$ such that $U=X\inter V$}\}$.
        $C\subseteq\pow{X}$ by \cref{subseteq}.
        $C$ has order at most $N+1$ in $X$ by
            \cref{subspace_intersections_preserve_order_bound_50}.
        For all $x\in X$ we have $x\in\unions{C}$ by
            \cref{subseteq,open_covering,covering,unions_iff,inter_lower_left,powerset_intro,inter_intro}.
        Thus $C$ is a covering of $X$ by \cref{covering,subseteq}.
        For all $U\in C$ we have
            $U$ is open in $X$ with respect to $T_d$ by
            \cref{subseteq,open_covering,open_in,subspace_topology,real_lt_subst_left,real_lt_subst_right}.
        $C$ is an open covering of $X$ with respect to $T_d$ by
            \cref{open_covering}.
        For all $U\in C$ for all $x,y$ such that
            $x\in U$ and $y\in U$ we have $d((x,y))<\epsilon$ by
            \cref{subseteq,real_lt_subst_left,real_lt_subst_right,inter_elim_left,inter_elim_right,times_tuple_intro,metric_on,funs_is_function,restrl_of_function_apply}.
        Thus there exists $D$ such that
            $D$ is an open covering of $X$ with respect to $T_d$ and
            $D$ has order at most $N+1$ in $X$ and
            for all $U\in D$ for all $x,y\in U$ we have
                $d((x,y))<\epsilon$.
    \end{subproof}
    $N$ is a topological dimension bound for $X$ with respect to $T_d$ by
        \cref{fine_ordered_metric_covers_dimension_bound_50}.
    Follows by \cref{real_lt_subst_left,real_lt_subst_right}.
\end{proof}

% from §50 Item 12: compact subspaces of euclidean space have bounded topological dimension
\begin{theorem}\label{compact_subspace_euclidean_dimension_bound_50}
    Assume $N\in\naturalsPlus$.
    Let $E = \realsPower{\initialSegment{N}}$.
    Let $S = \metricTopology{\squareMetric{N}}{E}$.
    Assume $X$ is a subspace of $E$ with respect to $S$.
    Let $T = \subspaceTopology{S}{X}$.
    Suppose $X$ is compact with respect to $T$.
    Then $N$ is a topological dimension bound for $X$ with respect to $T$.
\end{theorem}
\begin{proof}
    Follows by \cref{compact_euclidean_dimension_bound_principle_50}.
\end{proof}

% from §50 Item 13: compact manifolds have dimension at most their manifold dimension
\begin{corollary}\label{compact_manifold_dimension_bound_50}
    Assume $m\in\naturalsPlus$.
    Assume $X$ is a manifold of dimension $m$ with respect to $T$.
    Assume $X$ is compact with respect to $T$.
    Then $m$ is a topological dimension bound for $X$ with respect to $T$.
\end{corollary}
\begin{proof}
    $X$ is Hausdorff with respect to $T$ by \cref{m_manifold}.
    $T$ is a topology on $X$ by \cref{hausdorff}.
    $m\in\naturals$ by
        \cref{naturals,union_intro_left}.
    \begin{byCase}
    \caseOf{$X=\emptyset$.}
        Follows by \cref{empty_space_dimension_bound_50}.
    \caseOf{$X\neq\emptyset$.}
        $X$ is inhabited by \cref{nonempty_inhabited}.
        Let $I=\initialSegment{m}$.
        Let $R_0=\{(j,\reals)\mid j\in I\}$.
        Let $S_0=\{(j,\standardTopologyR)\mid j\in I\}$.
        Let $P_0=\productTopologyIndexed{S_0}{R_0}$.
        Let $E=\realsPower{I}$.
        Let $S=\metricTopology{\squareMetric{m}}{E}$.
        $I$ is inhabited by
            \cref{real_one_in_naturals,real_one_leq_naturalsplus,initial_segment_naturalsplus}.
        $\productFamily{R_0}=E$ by
            \cref{product_family_reals_power}.
        $\metricTopology{\euclideanMetric{m}}{E}=P_0$ by
            \cref{euclidean_metric_product_topology}.
        $\metricTopology{\euclideanMetric{m}}{E}=S$ by
            \cref{square_metric_topology_equals_euclidean}.
        $P_0=S$.
        $E$ is Hausdorff with respect to $S$ by
            \cref{square_metric_is_metric,metric_space_hausdorff}.
        $S$ is a topology on $E$ by \cref{hausdorff}.
        Take $n,U$ such that
            $n\in\naturalsPlus$ and
            $U$ is a function on $\initialSegment{n}$ and
            for all $i\in\initialSegment{n}$ there exist $W,h$ such that
                $\apply{U}{i}$ is open in $X$ with respect to $T$ and
                $W$ is open in $\productFamily{R_0}$ with respect to $P_0$ and
                $h$ is a homeomorphism between $\apply{U}{i}$ and $W$ with respect to
                    $\subspaceTopology{T}{\apply{U}{i}}$ and
                    $\subspaceTopology{P_0}{W}$ and
            $\img{U}{\initialSegment{n}}$ is a covering of $X$
            by \cref{compact_manifold_finite_chart_cover}.
        $X$ is normal with respect to $T$ by
            \cref{compact_hausdorff_normal}.
        Take $V$ such that
            $V$ is a function on $\initialSegment{n}$ and
            for all $i\in\initialSegment{n}$ we have
                $\apply{V}{i}$ is open in $X$ with respect to $T$ and
                $\closure{\apply{V}{i}}{X}{T}\subseteq\apply{U}{i}$ and
            $\img{V}{\initialSegment{n}}$ is a covering of $X$
            by \cref{finite_shrinking_cover}.
        Let $J=\initialSegment{n}$.
        Let $F=\{(i,\closure{\apply{V}{i}}{X}{T})\mid i\in J\}$.
        $F$ is a function by
            \cref{relation,rightunique,pair_eq_iff}.
        $\dom{F}=J$ by
            \cref{dom_intro,dom_iff,pair_eq_iff,subseteq,subseteq_antisymmetric}.
        $F$ is a function on $J$.
        For all $i\in J$ we have
            $\apply{F}{i}=\closure{\apply{V}{i}}{X}{T}$
            by \cref{function_apply_iff,pair_eq_iff}.
        Show that for all $i\in J$ we have
            $\apply{F}{i}\subseteq X$ and
            $\apply{F}{i}$ is closed in $X$ with respect to $T$ and
            $m$ is a topological dimension bound for $\apply{F}{i}$ with respect to
                $\subspaceTopology{T}{\apply{F}{i}}$.
        \begin{subproof}
            Fix $i\in J$.
            Let $A=\apply{F}{i}$.
            $A=\closure{\apply{V}{i}}{X}{T}$.
            $\apply{V}{i}\subseteq X$ by
                \cref{open_in,topology_on,subseteq,powerset_elim}.
            $A$ is closed in $X$ with respect to $T$ by
                \cref{closure_closed_in,real_lt_subst_left}.
            $A\subseteq X$ by \cref{closed_in}.
            $A$ is a subspace of $X$ with respect to $T$ by
                \cref{subspace_of}.
            $A$ is compact with respect to $\subspaceTopology{T}{A}$ by
                \cref{closed_subspace_compact}.
            $A\subseteq\apply{U}{i}$ by assumption.
            Take $W,h$ such that
                $\apply{U}{i}$ is open in $X$ with respect to $T$ and
                $W$ is open in $\productFamily{R_0}$ with respect to $P_0$ and
                $h$ is a homeomorphism between $\apply{U}{i}$ and $W$ with respect to
                    $\subspaceTopology{T}{\apply{U}{i}}$ and
                    $\subspaceTopology{P_0}{W}$ and
                $\img{U}{J}$ is a covering of $X$.
            Let $T_U=\subspaceTopology{T}{\apply{U}{i}}$.
            Let $S_W=\subspaceTopology{P_0}{W}$.
            $T_U$ is a topology on $\apply{U}{i}$ by
                \cref{homeomorphism,continuous,real_lt_subst_left}.
            $A$ is a subspace of $\apply{U}{i}$ with respect to $T_U$ by
                \cref{subspace_of}.
            $\apply{U}{i}\subseteq X$ by
                \cref{open_in,topology_on,subseteq,powerset_elim}.
            $\subspaceTopology{T_U}{A}=\subspaceTopology{T}{A}$ by
                \cref{subspace_subspace_topology}.
            $W\subseteq E$ by
                \cref{open_in,topology_on,subseteq,powerset_elim,real_lt_subst_right}.
            $W$ is Hausdorff with respect to $S_W$ by
                \cref{subspace_hausdorff}.
            Let $g=\restrl{h}{A}$.
            Let $B=\img{h}{A}$.
            $g$ is a homeomorphism between $A$ and $B$ with respect to
                $\subspaceTopology{T_U}{A}$ and
                $\subspaceTopology{S_W}{B}$
                by \cref{compact_homeomorphism_restriction_50,real_lt_subst_left}.
            $B\subseteq W$ by
                \cref{homeomorphism,bijections_to_funs,funs,rels_ran_subseteq,img_subseteq_ran,subseteq_transitive}.
            $B\subseteq E$ by \cref{subseteq_transitive}.
            $B$ is a subspace of $E$ with respect to $S$ by
                \cref{subspace_of}.
            $\subspaceTopology{S_W}{B}=\subspaceTopology{S}{B}$ by
                \cref{subspace_subspace_topology,real_lt_subst_left,real_lt_subst_right}.
            $g$ is a homeomorphism between $A$ and $B$ with respect to
                $\subspaceTopology{T}{A}$ and $\subspaceTopology{S}{B}$.
            $g$ is continuous from $A$ to $B$ with respect to
                $\subspaceTopology{T}{A}$ and $\subspaceTopology{S}{B}$
                by \cref{homeomorphism}.
            Let $C=\img{g}{A}$.
            Let $S_C=\subspaceTopology{\subspaceTopology{S}{B}}{C}$.
            $C$ is compact with respect to $S_C$ by
                \cref{continuous_image_compact}.
            $C=B$ by
                \cref{homeomorphism,bijections,funs_surj_iff,bijections_to_funs,funs,img_dom}.
            $S_C=\subspaceTopology{S}{B}$ by
                \cref{whole_subspace_topology_50,subspace_topology_is_topology,real_lt_subst_left}.
            $B$ is compact with respect to $\subspaceTopology{S}{B}$.
            $m$ is a topological dimension bound for $B$ with respect to
                $\subspaceTopology{S}{B}$
                by \cref{compact_subspace_euclidean_dimension_bound_50}.
            Follows by
                \cref{homeomorphism_preserves_dimension_bound_50}.
        \end{subproof}
        Let $Q=\unions{\img{F}{J}}$.
        $m$ is a topological dimension bound for $Q$ with respect to
            $\subspaceTopology{T}{Q}$
            by \cref{finite_closed_family_common_dimension_bound_50}.
        $1\in J$ by
            \cref{real_one_in_naturals,real_one_leq_naturalsplus,initial_segment_naturalsplus}.
        $\img{V}{J}$ is a covering of $X$ by assumption.
        $Q\subseteq X$ by
            \cref{unions_iff,img_iff,function_apply_iff,subseteq,closure_closed_in,closed_in}.
        For all $x\in X$ we have $x\in Q$ by
            \cref{covering,subseteq,unions_iff,img_of_function_elim,inter_elim_right,img_of_function_intro,inter_intro,real_lt_subst_left,open_in,topology_on,powerset_elim,subseteq_closure,real_lt_subst_right}.
        Thus $X\subseteq Q$ by \cref{subseteq}.
        Thus $Q=X$ by \cref{subseteq_antisymmetric}.
        $\subspaceTopology{T}{Q}=T$ by
            \cref{whole_subspace_topology_50,real_lt_subst_left}.
        Follows.
    \end{byCase}
\end{proof}

% from §50 Item 14: compact manifolds embed in euclidean space of dimension 2m+1
\begin{corollary}\label{compact_manifold_embeds_in_dimension_2m_plus_1_50}
    Assume $m\in\naturalsPlus$.
    Assume $X$ is a manifold of dimension $m$ with respect to $T$.
    Assume $X$ is compact with respect to $T$.
    Let $N = (1 + 1) \rmul m + 1$.
    Let $E = \realsPower{\initialSegment{N}}$.
    Let $S = \metricTopology{\squareMetric{N}}{E}$.
    Then there exists $f$ such that
        $f$ is a topological embedding from $X$ to $E$ with respect to $T$ and $S$.
\end{corollary}
\begin{proof}
    $m\in\naturals$ by
        \cref{naturals,union_intro_left}.
    $m$ is a topological dimension bound for $X$ with respect to $T$ by
        \cref{compact_manifold_dimension_bound_50}.
    $X$ is Hausdorff with respect to $T$ and
        $X$ satisfies the second countability axiom with respect to $T$ by
        \cref{m_manifold}.
    $X$ is normal with respect to $T$ by
        \cref{compact_hausdorff_normal}.
    $X$ is regular with respect to $T$ by
        \cref{normal_implies_regular}.
    $X$ is metrizable with respect to $T$ by
        \cref{urysohn_metrization_theorem}.
    Follows by
        \cref{menger_nobeling_dimension_bound_principle_50}.
\end{proof}

% from §50 helper definition 31: embeddable in some euclidean space
\begin{definition}\label{embeddable_in_some_euclidean_space_50}
    $X$ is embeddable in some euclidean space with respect to $T$ iff
        there exist $N,f$ such that
            $N\in\naturalsPlus$ and
            $f$ is a topological embedding from $X$ to $\realsPower{\initialSegment{N}}$ with respect to
                $T$ and $\metricTopology{\squareMetric{N}}{\realsPower{\initialSegment{N}}}$.
\end{definition}

% from §50 Item 15: compact metrizable spaces embed in euclidean space iff finite dimensional
\begin{corollary}\label{compact_metrizable_euclidean_embedding_iff_finite_dimensional_50}
    Assume $X$ is compact with respect to $T$.
    Assume $X$ is metrizable with respect to $T$.
    Then $X$ is embeddable in some euclidean space with respect to $T$ iff
        $X$ is finite dimensional with respect to $T$.
\end{corollary}
\begin{proof}
    Show that if $X$ is embeddable in some euclidean space with respect to $T$
        then $X$ is finite dimensional with respect to $T$.
    \begin{subproof}
        Assume $X$ is embeddable in some euclidean space with respect to $T$.
        Take $N,f$ such that
            $N\in\naturalsPlus$ and
            $f$ is a topological embedding from $X$ to
                $\realsPower{\initialSegment{N}}$ with respect to
                $T$ and
                $\metricTopology{\squareMetric{N}}{\realsPower{\initialSegment{N}}}$
            by \cref{embeddable_in_some_euclidean_space_50}.
        Let $E=\realsPower{\initialSegment{N}}$.
        Let $S=\metricTopology{\squareMetric{N}}{E}$.
        Let $Y=\img{f}{X}$.
        $f$ is continuous from $X$ to $E$ with respect to $T$ and $S$ and
            $f$ is a homeomorphism between $X$ and $Y$ with respect to
                $T$ and $\subspaceTopology{S}{Y}$
            by \cref{topological_embedding}.
        $Y$ is a subspace of $E$ with respect to $S$ by
            \cref{continuous,funs,rels_ran_subseteq,img_subseteq_ran,subseteq_transitive,subspace_of}.
        $Y$ is compact with respect to $\subspaceTopology{S}{Y}$ by
            \cref{continuous_image_compact}.
        $N$ is a topological dimension bound for $Y$ with respect to
            $\subspaceTopology{S}{Y}$
            by \cref{compact_subspace_euclidean_dimension_bound_50}.
        $N$ is a topological dimension bound for $X$ with respect to $T$ by
            \cref{homeomorphism_preserves_dimension_bound_50}.
        Follows by
            \cref{finite_dimensional_space_50,topological_dimension_bound_50}.
    \end{subproof}
    Show that if $X$ is finite dimensional with respect to $T$
        then $X$ is embeddable in some euclidean space with respect to $T$.
    \begin{subproof}
        Assume $X$ is finite dimensional with respect to $T$.
        Take $m\in\naturals$ such that
            $m$ is the topological dimension of $X$ with respect to $T$
            by \cref{finite_dimensional_dimension_exists_50}.
        Let $N=(1+1)\rmul m+1$.
        Let $E=\realsPower{\initialSegment{N}}$.
        Let $S=\metricTopology{\squareMetric{N}}{E}$.
        Take $f$ such that
            $f$ is a topological embedding from $X$ to $E$ with respect to
                $T$ and $S$
            by \cref{menger_nobeling_embedding_theorem_50}.
        $N\in\naturalsPlus$ by
            \cref{menger_nobeling_target_dimension_positive_50}.
        Follows by \cref{embeddable_in_some_euclidean_space_50}.
    \end{subproof}
\end{proof}
